\chapter{A weakening of Vandiver's conjecture}\label{chapter-tate-shafarevich}

The window bounds of Chapter~\ref{chapter-host-fields} need the Tate-Shafarevich groups of the Tate twists $\mathbf Q_\ell/\mathbf Z_\ell(-k)$ over the host fields to be small: on most windows of scales $x$, the sum over primes $\ell$ of $\log\ell$ times their total length for $1\le k<x^{401}$ is at most $e^{x^{\varepsilon}/2}$.
This chapter proves that bound, assuming the Galois representations of Chapter~\ref{chapter-galois-representations}.
For odd $\ell$ the groups split into eigenspaces for complex conjugation: the minus part is bounded by generalized Bernoulli numbers, and the plus part, which Vandiver's conjecture predicts to vanish, by Eisenstein congruences for $U(3)$ and $GSp_4$.
The prime $2$ and finitely many bad primes are handled by Iwasawa theory in cyclotomic towers, the Ferrero-Washington theorem, a bound of Soulé type for $H^2$, and a pigeonhole argument over windows of scales.

\section{Tate-Shafarevich groups of Tate twists}\label{section-sha-groups}

\begin{definition}\label{sha-tate-shafarevich-group}
Let $E$ be a number field with algebraic closure $\bar E$, and let $\ell$ be a prime.
Let $S_\ell$ be the set of places of $E$ above $\ell$, let $E^{S_\ell}\subset\bar E$ be the maximal extension of $E$ unramified at every finite place outside $S_\ell$ (no condition is imposed at the archimedean places), and put $G_{E,S_\ell}=\operatorname{Gal}(E^{S_\ell}/E)$.
For each $v\in S_\ell$ choose a place of $\bar E$ above $v$; its decomposition group maps to $G_{E,S_\ell}$, and we write $G_{E_v}$ for it.
Let $\chi_{\mathrm{cyc}}:G_{E,S_\ell}\to\mathbf Z_\ell^\times$ be the $\ell$-adic cyclotomic character.
For an integer $k$, let $\mathbf Q_\ell/\mathbf Z_\ell(-k)$ be the discrete module $\mathbf Q_\ell/\mathbf Z_\ell$ on which $\sigma$ acts by multiplication by $\chi_{\mathrm{cyc}}(\sigma)^{-k}$.
The \emph{Tate-Shafarevich group} of the twist $k$ at $\ell$ is
\[
 \text{Ш}^1_\ell(E,k)=\ker\Big(H^1\big(G_{E,S_\ell},\mathbf Q_\ell/\mathbf Z_\ell(-k)\big)\to\prod_{v\in S_\ell}H^1\big(G_{E_v},\mathbf Q_\ell/\mathbf Z_\ell(-k)\big)\Big),
\]
with continuous cochain cohomology.
It is a $\mathbf Z_\ell$-module, and $\operatorname{len}\text{Ш}^1_\ell(E,k)\in\{0,1,2,\dots\}\cup\{\infty\}$ denotes its length.
Another choice of places conjugates each $G_{E_v}$ by an element of $G_{E,S_\ell}$, which acts trivially on $H^1(G_{E,S_\ell},-)$, so the group does not depend on the choices.
\end{definition}
\leanref{Formalism/GaloisCohomology/Statement.lean}{Formalism.GaloisCohomology.LData}
\leanref{Formalism/GaloisCohomology/Statement.lean}{Formalism.GaloisCohomology.Sha}
\leanref{Formalism/GaloisCohomology/Statement.lean}{Formalism.GaloisCohomology.shaLen}

\begin{definition}\label{sha-plus-minus}
Let $E$ be a CM field with complex conjugation $c\in\operatorname{Aut}(E)$, let $\ell$ be a prime and let $k\ge0$ be an integer.
Let $\sigma$ be an automorphism of $\bar E$ restricting to $c$.
Conjugation by $\sigma$ preserves the kernel of $\operatorname{Gal}(\bar E/E)\to G_{E,S_\ell}$, and $\sigma$ acts on $H^1(G_{E,S_\ell},\mathbf Q_\ell/\mathbf Z_\ell(-k))$ by
\[
 (\sigma\cdot z)(g)=\chi_{\mathrm{cyc}}(\sigma)^{-k}\,z(\sigma^{-1}g\sigma).
\]
This action does not depend on the choice of $\sigma$, it is an involution, and it preserves $\text{Ш}^1_\ell(E,k)$.
The \emph{plus part} $\text{Ш}^+_\ell(E,k)$ is the $(-1)^k$-eigenspace of this involution on $\text{Ш}^1_\ell(E,k)$, and the \emph{minus part} $\text{Ш}^-_\ell(E,k)$ is the $(-1)^{k+1}$-eigenspace.
\end{definition}
\leanref{Development/AutomorphicArithmetic/AutomorphicVarieties/LengthEqAddEigen.lean}{UpperTheory.V.ShaPlus}
\leanref{Development/AutomorphicArithmetic/AutomorphicVarieties/LengthEqAddEigen.lean}{UpperTheory.V.ShaMinus}
\leanref{Development/AutomorphicArithmetic/AutomorphicVarieties/PartsConj.lean}{UpperTheory.V.galoisActionOnCohomology}

\begin{lemma}\label{sha-plus-minus-length}
Let $E$ be a CM field, let $\ell$ be an odd prime and let $k\ge0$ be an integer.
Then
\[
 \operatorname{len}\text{Ш}^1_\ell(E,k)=\operatorname{len}\text{Ш}^+_\ell(E,k)+\operatorname{len}\text{Ш}^-_\ell(E,k).
\]
\end{lemma}
\begin{proof}
Let $f$ be the involution of Definition~\ref{sha-plus-minus} and $s=(-1)^k$.
Since $\ell$ is odd, $2$ is a unit in $\mathbf Z_\ell$.
Every $x\in\text{Ш}^1_\ell(E,k)$ is $\tfrac12(x+sf(x))+\tfrac12(x-sf(x))$, and the two summands lie in the $s$- and $(-s)$-eigenspaces.
So $\text{Ш}^1_\ell(E,k)$ is the direct sum of its plus and minus parts, and length is additive.
\end{proof}
\leanref{Development/AutomorphicArithmetic/AutomorphicVarieties/LengthEqAddEigen.lean}{UpperTheory.V.shaLen_eq_add}
\leanref{Development/AutomorphicArithmetic/AutomorphicVarieties/LengthEqAddEigen.lean}{UpperTheory.V.length_eq_add_eigen}

\begin{definition}\label{sha-bounded-field}
Let $E$ be a number field, let $t\ge0$ be an integer and let $B$ be a real number.
\begin{enumerate}
\item $E$ satisfies the \emph{Tate-Shafarevich bound} $(t,B)$ if $\text{Ш}^1_\ell(E,k)$ has finite length for every prime $\ell$ and every $0<k<t$, and
\[
 \sum_{\ell\in S}\log\ell\sum_{k=1}^{t-1}\operatorname{len}\text{Ш}^1_\ell(E,k)\le B
\]
for every finite set $S$ of primes.
\item For a prime $\ell$, $E$ satisfies the \emph{$\ell$-bound} $(t,B)$ if $\text{Ш}^1_\ell(E,k)$ has finite length for $0<k<t$ and $\sum_{k=1}^{t-1}\operatorname{len}\text{Ш}^1_\ell(E,k)\le B$.
\item Suppose that $E$ is a CM field. Then $E$ satisfies the \emph{plus bound} $(t,B)$ if $\text{Ш}^+_\ell(E,k)$ has finite length for every odd prime $\ell$ and every $0<k<t$, and $\sum_{\ell\in S}\log\ell\sum_{k=1}^{t-1}\operatorname{len}\text{Ш}^+_\ell(E,k)\le B$ for every finite set $S$ of odd primes. The \emph{minus bound} $(t,B)$ is defined in the same way with $\text{Ш}^-$.
\end{enumerate}
\end{definition}
\leanref{Development/AutomorphicArithmetic/CharacterBoxes/ExistsPageN.lean}{UpperTheory.Top.VVle}
\leanref{Development/AutomorphicArithmetic/AutomorphicVarieties/LengthEqAddEigen.lean}{UpperTheory.V.VVat}
\leanref{Development/AutomorphicArithmetic/AutomorphicVarieties/LengthEqAddEigen.lean}{UpperTheory.V.VVplusLe}
\leanref{Development/AutomorphicArithmetic/AutomorphicVarieties/LengthEqAddEigen.lean}{UpperTheory.V.VVminusLe}

\begin{lemma}\label{sha-bound-from-parts}
Let $E$ be a CM field, let $t\ge0$ be an integer and let $B_2,B_-,B_+$ be real numbers.
Suppose that $E$ satisfies the $2$-bound $(t,B_2)$, the minus bound $(t,B_-)$ and the plus bound $(t,B_+)$ (Definition~\ref{sha-bounded-field}).
Then $E$ satisfies the Tate-Shafarevich bound $(t,B_2\log2+B_-+B_+)$.
\end{lemma}
\begin{proof}
Finiteness at $\ell=2$ is part of the $2$-bound, and for odd $\ell$ it follows from Lemma~\ref{sha-plus-minus-length}.
Split a finite set $S$ of primes into $S\cap\{2\}$ and its odd part.
The term at $2$ is at most $B_2\log2$.
On the odd part, Lemma~\ref{sha-plus-minus-length} writes each length as the sum of the lengths of the plus and minus parts, whose weighted sums are at most $B_+$ and $B_-$.
\end{proof}
\leanref{Development/AutomorphicArithmetic/AutomorphicVarieties/LengthEqAddEigen.lean}{UpperTheory.V.vvle_of_parts}

\begin{definition}\label{sha-vandiver-weakening}
Let $\psi$ be a primitive Dirichlet character and let $C$ be a real number.
The character $\psi$ satisfies the \emph{weakening of Vandiver's conjecture with exponent $C$} if the following holds for every $0<\varepsilon<1/52$ and every function $e:\mathbf R\to\mathbf Z_{\ge0}$ with $x^{2+8\varepsilon}\le2^{e(x)}\le2x^{2+8\varepsilon}$ for all $x\ge1$.
There are $J\ge1$ and $X_0$ such that for every $X\ge X_0$ and every pair of primes $(q,d)$ admissible at $X^{1/J}$ (Definition~\ref{hf-admissible-pair}) there is $Y$ with $X^{1/J}\le Y$ and $Y^2\le X$ such that, for every $x\in[Y,Y^2]$, every host field $E$ of $(\psi,e(x),q,d)$ (Definition~\ref{hf-host-field-of-character}) satisfies the Tate-Shafarevich bound $\big(\lfloor x^{C}\rfloor,\sqrt{\exp(x^{\varepsilon})}\big)$ of Definition~\ref{sha-bounded-field}.
The \emph{weakening of Vandiver's conjecture} is the statement that every primitive Dirichlet character satisfies it with exponent $401$.
\end{definition}
\leanref{Development/AutomorphicArithmetic/CharacterBoxes/ExistsPageN.lean}{UpperTheory.Top.PageN}
\leanref{Development/AutomorphicArithmetic/CharacterBoxes/ExistsPageN.lean}{UpperTheory.Top.VAt}
\leanref{Development/AutomorphicArithmetic/CharacterBoxes/ExistsPageN.lean}{UpperTheory.Top.VWinAt}
\leanref{Development/AutomorphicArithmetic/CharacterBoxes/ExistsPageN.lean}{UpperTheory.Top.VWin}

Vandiver's conjecture asserts that $\ell$ does not divide the class number of $\mathbf Q(\zeta_\ell)^+$; equivalently, the plus parts above vanish when $E=\mathbf Q(\zeta_\ell)$.
The statement of Definition~\ref{sha-vandiver-weakening} allows every Tate-Shafarevich group to be nonzero, and bounds a weighted sum of their lengths for host fields of degree about $x^{2+8\varepsilon}$ and twists up to $x^{401}$.
Its proof occupies the rest of the chapter and ends in Theorem~\ref{sha-weakening-from-galois-representations}.

\section{Abelian fields and their composita}\label{section-sha-composita}

\begin{definition}\label{sha-abelian-field}
An \emph{abelian field datum} is a pair $M=(m,H)$ of an integer $m\ge1$ and a subgroup $H\subset(\mathbf Z/m)^\times$; its \emph{degree} is $\deg M=[(\mathbf Z/m)^\times:H]$.
\begin{enumerate}
\item A number field $L$ is a \emph{field of $M$} if for some $N\ge1$ divisible by $m$ it is isomorphic to the subfield of $\mathbf Q(\zeta_N)$ fixed by the automorphisms $\zeta_N\mapsto\zeta_N^{a}$ with $a\bmod m\in H$.
\item $M$ \emph{contains $\mu_\ell$}, for a prime $\ell$, if $\ell\mid m$ and every integer $a$ with $a\bmod m\in H$ satisfies $a\equiv1\pmod\ell$.
\item A \emph{prime pair} for $M$ is a pair of primes $(q,d)$ with $q\equiv3\pmod8$, $d\equiv7\pmod{16}$ and $qd$ prime to $m$.
\item Let $e\ge0$ be an integer and $q,d$ primes. A number field $E$ is a \emph{compositum} of $(M,e,q,d)$ if for some $N\ge1$ divisible by $2^{e+2}$, $m$, $q$ and $d$ it is isomorphic to the subfield of $\mathbf Q(\zeta_N)$ fixed by the automorphisms $\zeta_N\mapsto\zeta_N^{a}$ with $a\equiv\pm1\pmod{2^{e+2}}$, $a\bmod m\in H$ and $(a/q)=(a/d)=1$.
\item The \emph{conductor} $N_E$ of a number field $E$ is the least $N\ge1$ such that $E$ embeds in $\mathbf Q(\zeta_N)$, and $0$ if there is none.
\end{enumerate}
For a primitive Dirichlet character $\psi$ of conductor $N_\psi$, the datum of $\psi$ is $M_\psi=(N_\psi,\ker\psi)$. For $m\ge1$, the cyclotomic datum $(m,\{1\})$ has field $\mathbf Q(\zeta_m)$.
\end{definition}
\leanref{Development/AutomorphicArithmetic/AutomorphicVarieties/StmtFields.lean}{UpperTheory.V.AbF}
\leanref{Development/AutomorphicArithmetic/AutomorphicVarieties/StmtFields.lean}{UpperTheory.V.IsAbField}
\leanref{Development/AutomorphicArithmetic/AutomorphicVarieties/StmtFields.lean}{UpperTheory.V.AbF.HasMu}
\leanref{Development/AutomorphicArithmetic/AutomorphicVarieties/PrimeSetBoundDataOfPrimitiveChar.lean}{UpperTheory.V.PrimePair}
\leanref{Development/AutomorphicArithmetic/AutomorphicVarieties/StmtFields.lean}{UpperTheory.V.IsCompField}
\leanref{Development/AutomorphicArithmetic/AutomorphicVarieties/StmtFields.lean}{UpperTheory.V.condF}
\leanref{Development/AutomorphicArithmetic/AutomorphicVarieties/StmtFields.lean}{UpperTheory.V.AbF.ofChar}
\leanref{Development/AutomorphicArithmetic/AutomorphicVarieties/StmtFields.lean}{UpperTheory.V.AbF.cyc}

\begin{lemma}\label{sha-composita}
Let $\psi$ be a primitive Dirichlet character of conductor $N_\psi$, let $e\ge0$ be an integer and let $q,d$ be primes.
\begin{enumerate}
\item A number field is a compositum of $(M_\psi,e,q,d)$ if and only if it is a host field of $(\psi,e,q,d)$ (Definition~\ref{hf-host-field-of-character}).
\item If $q\equiv3\pmod4$, every host field $E$ of $(\psi,e,q,d)$ is a CM field, and $1\le N_E\le2^{e+2}N_\psi qd$ and $[E:\mathbf Q]\le N_E$.
\item Let $M$ be an abelian field datum, let $(q,d)$ be a prime pair for $M$, and let $E$ be a compositum of $(M,e,q,d)$. Every odd prime dividing $[E:\mathbf Q]$ divides $\deg M$.
\end{enumerate}
\end{lemma}
\begin{proof}
Part (a) holds because $a\bmod N_\psi\in\ker\psi$ if and only if $\psi(a)=1$.

For (b), $E$ is abelian over $\mathbf Q$, and $a=-1$ does not fix it because $(-1/q)=-1$.
So complex conjugation is nontrivial on $E$, and an abelian field with nontrivial complex conjugation is a CM field.
The fixing subgroup contains the kernel of reduction to $L=\operatorname{lcm}(2^{e+2},N_\psi,q,d)$, so $E$ embeds in $\mathbf Q(\zeta_L)$ and $N_E\le L\le2^{e+2}N_\psi qd$.
Then $[E:\mathbf Q]\le\varphi(N_E)\le N_E$.

For (c), let $\tau$ be an automorphism of $\mathbf Q(\zeta_N)$, $\zeta_N\mapsto\zeta_N^{a}$.
Then $a^{2^{e+1}\deg M}$ is $\equiv1\pmod{2^{e+2}}$, lies in $H$ modulo $m$, and is a square modulo $q$ and modulo $d$.
So $\tau^{2^{e+1}\deg M}$ fixes $E$, and the exponent of $\operatorname{Gal}(E/\mathbf Q)$ divides $2^{e+1}\deg M$.
An odd prime dividing $[E:\mathbf Q]$ divides this exponent, hence $\deg M$.
\end{proof}
\leanref{Development/AutomorphicArithmetic/AutomorphicVarieties/StmtFields.lean}{UpperTheory.V.isCompField_ofChar}
\leanref{Development/AutomorphicArithmetic/AutomorphicVarieties/CMFieldExponentialBounds/IsCMFieldOfIsChiField.lean}{UpperTheory.V.isCMField_of_isChiField}
\leanref{Development/AutomorphicArithmetic/AutomorphicVarieties/CMFieldExponentialBounds/ConductorBoundsOfIsChiField.lean}{UpperTheory.V.conductor_bounds_of_isChiField}
\leanref{Development/AutomorphicArithmetic/AutomorphicVarieties/OddResidueBounds/DegExistsLift.lean}{UpperTheory.V.Odd.deg}

\section{The minus part}\label{section-sha-minus}

For a Dirichlet character $\chi$ modulo $N$ with values in a field of characteristic zero, let $\chi^0$ be the primitive character inducing it, of conductor $f$, and let
\[
 B_{k,\chi}=f^{k-1}\sum_{a=1}^{f}\chi^0(a)B_k(a/f),\qquad \alpha_k(\chi)=N\Big(\prod_{p\le k+1}p\Big)B_{k,\chi},
\]
where $B_k(X)$ is the $k$-th Bernoulli polynomial and the product runs over primes.

\begin{theorem}[\cite{MW84,Kol90,Rub90}]\label{sha-bernoulli-bound}
Let $E$ be an abelian CM field with algebraic closure $\bar{\mathbf Q}$, let $N\ge1$, and let $\rho:(\mathbf Z/N)^\times\to\operatorname{Gal}(E/\mathbf Q)$ be a homomorphism such that $\rho(b)$ is the restriction of every automorphism of $\bar{\mathbf Q}$ that acts on $N$-th roots of unity by $\zeta\mapsto\zeta^{b}$.
Let $\ell$ be an odd prime not dividing $[E:\mathbf Q]$, and let $\mathcal O\subset\bar{\mathbf Q}_\ell$ be a subring that is a discrete valuation ring finite over $\mathbf Z_\ell$.
Let $k\ge1$, let $\psi:\operatorname{Gal}(E/\mathbf Q)\to\mathcal O^\times$ be a character, and let $\chi_\psi=\psi\circ\rho$, a Dirichlet character modulo $N$; suppose that $\chi_\psi(-1)=(-1)^{k+1}$.
Let $\text{Ш}_\ell(E,\mathcal O(-k))$ be the subgroup of $H^1\big(\operatorname{Gal}(\bar{\mathbf Q}/E),\mathcal O\otimes\mathbf Q_\ell/\mathbf Z_\ell(-k)\big)$ of classes that vanish on the decomposition groups at the places above $\ell$ and on the inertia groups at all places, with its action of $\operatorname{Gal}(\bar{\mathbf Q}/\mathbf Q)$, and let $\text{Ш}_\ell(E,\mathcal O(-k))^{\psi}$ be the part on which $\operatorname{Gal}(\bar{\mathbf Q}/\mathbf Q)$ acts through $\psi$.
Then for every $a\in\mathcal O$ equal to $\alpha_{k+1}(\chi_\psi)$,
\[
 \operatorname{len}_{\mathcal O}\text{Ш}_\ell(E,\mathcal O(-k))^{\psi}\le v_{\mathcal O}(a).
\]
\end{theorem}
\leanref{Development/AutomorphicArithmetic/AutomorphicVarieties/OddResidueBounds/ZeroMemUnrHom.lean}{UpperTheory.V.Odd.PsiBound}
\leanref{Development/AutomorphicArithmetic/AutomorphicVarieties/PrimitiveFrobeniusFixed/Iv.lean}{UpperTheory.V.ML.iv}

For class groups of abelian fields, bounds of this kind by generalized Bernoulli numbers are the theorem of Mazur and Wiles.
Rubin's proof uses Kolyvagin's Euler system of Gauss sums \cite{MW84,Kol90,Rub90}.

\begin{lemma}\label{sha-minus-prime-to-degree}
Let $E$ be a CM field that embeds in $\mathbf Q(\zeta_N)$ for some $N\ge1$, and let $k\ge1$.
Then $\text{Ш}^-_\ell(E,k)$ has finite length for every odd prime $\ell$ not dividing $[E:\mathbf Q]$, and
\[
 \sum_{\ell\in S}\log\ell\cdot\operatorname{len}\text{Ш}^-_\ell(E,k)\le[E:\mathbf Q]\log\big(16N^{k+2}(k+1)!\big)
\]
for every finite set $S$ of odd primes not dividing $[E:\mathbf Q]$.
\end{lemma}
\begin{proof}
Restriction from $\mathbf Q(\zeta_N)$ gives a surjection $\rho:(\mathbf Z/N)^\times\to\operatorname{Gal}(E/\mathbf Q)$ as in Theorem~\ref{sha-bernoulli-bound}, with kernel of index $[E:\mathbf Q]$, and $\rho(-1)$ is complex conjugation.

Fix an odd prime $\ell\nmid[E:\mathbf Q]$, and a discrete valuation ring $\mathcal O\subset\bar{\mathbf Q}_\ell$ finite over $\mathbf Z_\ell$ that contains the values of all characters of $\operatorname{Gal}(E/\mathbf Q)$.
Inflation identifies $\text{Ш}^1_\ell(E,k)\otimes\mathcal O$ with $\text{Ш}_\ell(E,\mathcal O(-k))$, and since $\ell\nmid[E:\mathbf Q]$ this group is the direct sum of its $\psi$-parts.
Complex conjugation acts on the $\psi$-part by $\psi(\rho(-1))$, so the minus part is the sum of the $\psi$-parts with $\chi_\psi(-1)=(-1)^{k+1}$.
Theorem~\ref{sha-bernoulli-bound} bounds each of them, which gives finiteness and
\[
 \operatorname{len}\text{Ш}^-_\ell(E,k)\cdot\log\ell\le\sum_{\chi\in T}-\log|\alpha_{k+1}(\chi)|_\ell
\]
for a set $T$ of characters modulo $N$ trivial on $\ker\rho$ with $\chi(-1)=(-1)^{k+1}$.

The characters $\chi$ modulo $N$ trivial on $\ker\rho$ with $\chi(-1)=(-1)^{k+1}$ form a set stable under $\operatorname{Gal}(\bar{\mathbf Q}/\mathbf Q)$, of size at most $[E:\mathbf Q]$.
The numbers $\alpha_{k+1}(\chi)$ are nonzero algebraic integers with complex absolute value at most $16N^{k+2}(k+1)!$, so their product over this set is a nonzero integer $M_0$ with $\log|M_0|\le[E:\mathbf Q]\log(16N^{k+2}(k+1)!)$.
By the ultrametric inequality, $\sum_{\chi\in T}-\log|\alpha_{k+1}(\chi)|_\ell\le v_\ell(M_0)\log\ell$.
Summing over $\ell\in S$ gives at most $\log|M_0|$.
\end{proof}
\leanref{Development/AutomorphicArithmetic/AutomorphicVarieties/OddResidueBounds/ZeroMemUnrHom.lean}{UpperTheory.V.Odd.PrimeBound}
\leanref{Development/AutomorphicArithmetic/AutomorphicVarieties/OddResidueBounds/ThmBcOfPsiBound.lean}{UpperTheory.V.Odd.primeBound_of}
\leanref{Development/AutomorphicArithmetic/AutomorphicVarieties/OddResidueBounds/AtPrime.lean}{UpperTheory.V.Odd.atPrime}
\leanref{Development/AutomorphicArithmetic/AutomorphicVarieties/OddResidueBounds/ProdVal.lean}{UpperTheory.V.Odd.prodVal}

Let $E$ be an abelian number field, let $M\ge1$ and let $\rho:(\mathbf Z/M)^\times\to\operatorname{Gal}(E/\mathbf Q)$ be a homomorphism such that $\rho(b)$ is the restriction of every automorphism of $\bar{\mathbf Q}$ acting on $M$-th roots of unity by $\zeta\mapsto\zeta^{b}$; write $\sigma_b=\rho(b)$.
The group $\operatorname{Gal}(E/\mathbf Q)$ acts on $H^1(G_{E,S_\ell},\mathbf Q_\ell/\mathbf Z_\ell(-k))$: for $\tau\in\operatorname{Gal}(E/\mathbf Q)$ choose an automorphism $\sigma$ of $\bar E$ restricting to $\tau$ and use the formula of Definition~\ref{sha-plus-minus}.
As there, the action does not depend on $\sigma$ and preserves $\text{Ш}^1_\ell(E,k)$, so $\text{Ш}^1_\ell(E,k)$ is a module over $\mathbf Z_\ell[\operatorname{Gal}(E/\mathbf Q)]$.

\begin{definition}\label{sha-stickelberger-elements}
Let $M\ge1$, let $\Gamma$ be a finite abelian group, and let $b\mapsto\sigma_b$ be a homomorphism $(\mathbf Z/M)^\times\to\Gamma$.
For $b\in\mathbf Z/M$ let $\langle b\rangle\in\{0,1,\dots,M-1\}$ be its representative.
Let $k\ge1$ and let $a\ge1$ be an integer prime to $M$.
\begin{enumerate}
\item The \emph{higher Stickelberger element} is
\[
 y_k(a)=(1-a^{k}\sigma_a)\,\theta_k\in\mathbf Q[\Gamma],\qquad\theta_k=-\sum_{b\in(\mathbf Z/M)^\times}\frac{M^{k-1}B_k(\langle b\rangle/M)}{k}\,\sigma_b.
\]
\item Suppose that $a$ is odd, and let $e\ge0$. The \emph{twisted Stickelberger element} of level $M$ is
\[
 \tilde y_e(a)=\sum_{b\in(\mathbf Z/M)^\times}\langle b\rangle^{e}\Big(\Big\lfloor\frac{a\langle a^{-1}b\rangle}{M}\Big\rfloor-\frac{a-1}{2}\Big)\sigma_b\in\mathbf Z[\Gamma].
\]
\end{enumerate}
\end{definition}
\leanref{Development/AutomorphicArithmetic/AutomorphicVarieties/StickelbergerElements/CoefMapSingle.lean}{UpperTheory.V.B4.thetaEl}
\leanref{Development/AutomorphicArithmetic/AutomorphicVarieties/StickelbergerElements/CoefMapSingle.lean}{UpperTheory.V.B4.yEl}
\leanref{Development/AutomorphicArithmetic/AutomorphicVarieties/StickelbergerElements/CoefMapSingle.lean}{UpperTheory.V.B4.ytEl}

\begin{lemma}\label{sha-stickelberger-annihilator}
Let $E$ be a number field that embeds in $\mathbf Q(\zeta_N)$, let $\ell$ be a prime dividing $N$, and suppose that $\ell$ is odd or that $\chi_{\mathrm{cyc}}\equiv1\pmod4$ on $\operatorname{Gal}(\bar E/E)$.
Let $\rho:(\mathbf Z/N)^\times\to\operatorname{Gal}(E/\mathbf Q)$ be the restriction map, as above, and form the elements of Definition~\ref{sha-stickelberger-elements} for $b\mapsto\rho(b)$.
Let $k\ge1$ and let $a$ be an odd integer prime to $N$.
Then the coefficients of $y_{k+1}(a)$ lie in $\mathbf Z_{(\ell)}$, and its image $z\in\mathbf Z_\ell[\operatorname{Gal}(E/\mathbf Q)]$ satisfies
\[
 z\cdot\varphi(N)x=0\qquad\text{for every }x\in\text{Ш}^1_\ell(E,k).
\]
\end{lemma}
\begin{proof}
Let $U_\infty\subset\operatorname{Gal}(\bar E/E)$ be the subgroup fixing all $\ell$-power roots of unity.
It acts trivially on $\mathbf Q_\ell/\mathbf Z_\ell(-k)$, so the restriction of a class $x$ to $U_\infty$ is a homomorphism $x|_{U_\infty}:U_\infty\to\mathbf Q_\ell/\mathbf Z_\ell$.
For $x\in\text{Ш}^1_\ell(E,k)$ it vanishes on the inertia groups at all places and on the decomposition groups at the places above $\ell$.

\emph{Step 1: restriction to $U_\infty$ is injective.}
The kernel of restriction is $H^1$ of $\operatorname{Gal}(E(\mu_{\ell^\infty})/E)$ with coefficients $\mathbf Q_\ell/\mathbf Z_\ell(-k)$.
Under the hypothesis on $\ell$ this group is procyclic, and a topological generator $\gamma$ acts on the coefficients by $u=\chi_{\mathrm{cyc}}(\gamma)^{-k}\ne1$.
Its $H^1$ is the module of coinvariants $(\mathbf Q_\ell/\mathbf Z_\ell)/(u-1)$, which is $0$ because $\mathbf Q_\ell/\mathbf Z_\ell$ is divisible.

\emph{Step 2: Stickelberger's theorem.}
Let $x\in\text{Ш}^1_\ell(E,k)$ and let $\ell^{r}x=0$.
There is $m_0$ such that for $m\ge m_0$ and $M=N\ell^m$, with $\tilde y_k(a)$ formed at level $M$ for $b\mapsto\rho(b\bmod N)$, the class $\tilde y_k(a)\cdot x$ vanishes on $U_\infty\cap\operatorname{Gal}(\bar E/E(\mu_N))$.
To see this, let $b\in(\mathbf Z/M)^\times$ and let $\sigma'_b$ be an automorphism of $\bar{\mathbf Q}$ acting on $M$-th roots of unity by $\zeta\mapsto\zeta^{b}$.
Then $\chi_{\mathrm{cyc}}(\sigma'_b)\equiv\langle b\rangle\pmod{\ell^m}$, so for $m\ge r$ the factor $\langle b\rangle^{k}$ in $\tilde y_k(a)$ cancels the factor $\chi_{\mathrm{cyc}}(\sigma'_b)^{-k}$ in the action of $\sigma_b$ on $\ell^r$-torsion.
Hence on $u\in U_\infty\cap\operatorname{Gal}(\bar E/E(\mu_M))$,
\[
 (\tilde y_k(a)\cdot x)(u)=\sum_{b}\Big(\Big\lfloor\frac{a\langle a^{-1}b\rangle}{M}\Big\rfloor-\frac{a-1}{2}\Big)\,x\big(\sigma'^{-1}_bu\,\sigma'_b\big).
\]
The homomorphism $x|_{U_\infty}$ extends, for $m$ large, to a homomorphism on $\operatorname{Gal}(\bar E/E(\mu_{M}))=\operatorname{Gal}(\bar{\mathbf Q}/\mathbf Q(\zeta_M))$ that is unramified everywhere, so by class field theory it factors through the class group of $\mathbf Q(\zeta_M)$, and the sum above is the action of $\sum_b\big(\lfloor a\langle a^{-1}b\rangle/M\rfloor-\frac{a-1}2\big)\sigma_b^{-1}$ on the class group.
Substituting $b=ab'$ shows that this element is $\sigma_a^{-1}(a-\sigma_a)\vartheta-\frac{a-1}{2}\sum_b\sigma_b^{-1}$, where $\vartheta=\sum_b(\langle b\rangle/M)\,\sigma_b^{-1}$ is the Stickelberger element of $\mathbf Q(\zeta_M)$.
The element $(a-\sigma_a)\vartheta$ kills the class group by Stickelberger's theorem \cite[Theorem 6.10]{Was97}, and the norm element $\sum_b\sigma_b^{-1}$ kills it because $\mathbf Q$ has class number $1$.
The elements $u$ of $U_\infty\cap\operatorname{Gal}(\bar E/E(\mu_N))$ already fix $\mu_M$, which gives the claim.

\emph{Step 3: a congruence.}
The coefficients of $y_{k+1}(a)$ lie in $\mathbf Z_{(\ell)}$, and for every $r$ there is $m_0$ such that for $m\ge m_0$ the image of $\tilde y_k(a)$ formed at level $N\ell^m$ is congruent to $y_{k+1}(a)$ modulo $\ell^r\mathbf Z_{(\ell)}[\operatorname{Gal}(E/\mathbf Q)]$.
This is a Kummer congruence for the Bernoulli polynomial $B_{k+1}$: for each $g\in\operatorname{Gal}(E/\mathbf Q)$, the coefficient of $g$ in $\tilde y_k(a)$ formed at level $N\ell^m$ differs from its coefficient in $y_{k+1}(a)$ by an element of $\ell^{m-1-c}\mathbf Z_{(\ell)}$, where $c$ is the $\ell$-adic valuation of $2(k+1)\prod_{p\le k+2}p$.
It uses that $a$ is prime to $N\ell$ and that $\ell\mid N$; one takes $m_0=r+c+1$.

\emph{Step 4.}
Let $x\in\text{Ш}^1_\ell(E,k)$ with $\ell^rx=0$, and choose $m$ for which Steps 2 and 3 apply.
Then $z\cdot x=\tilde y_k(a)\cdot x$, so $z\cdot x$ vanishes on $U_\infty\cap\operatorname{Gal}(\bar E/E(\mu_N))$.
For $u\in U_\infty$ the element $u^{\varphi(N)}$ fixes $\mu_N$, and $(z\cdot x)|_{U_\infty}$ is a homomorphism, so $\varphi(N)\,(z\cdot x)$ vanishes on $U_\infty$.
By Step 1, $\varphi(N)\,z\cdot x=0$.
\end{proof}
\leanref{Development/AutomorphicArithmetic/AutomorphicVarieties/OddResidueBounds/ZKill.lean}{UpperTheory.V.Odd.zKill}
\leanref{Development/AutomorphicArithmetic/AutomorphicVarieties/GaloisLibrary/KummerTheory/ResInfInj.lean}{UpperTheory.V.Kum.resInfInj}
\leanref{Development/AutomorphicArithmetic/AutomorphicVarieties/OddResidueBounds/ResKill.lean}{UpperTheory.V.Odd.resKill}
\leanref{Development/AutomorphicArithmetic/AutomorphicVarieties/OddResidueBounds/HomKill.lean}{UpperTheory.V.Odd.homKill}
\leanref{Development/AutomorphicArithmetic/AutomorphicVarieties/OddResidueBounds/ClKill.lean}{UpperTheory.V.Odd.clKill}
\leanref{Development/AutomorphicArithmetic/AutomorphicVarieties/OddResidueBounds/NearLevel.lean}{UpperTheory.V.Odd.congr}
\leanref{Development/AutomorphicArithmetic/AutomorphicVarieties/OddResidueBounds/TotKill.lean}{UpperTheory.V.Odd.totKill}

\begin{lemma}\label{sha-stickelberger-determinant}
Let $A$ be a commutative ring, let $\Gamma$ be a finite abelian group, let $c\in\Gamma$ with $c^2=1\ne c$, and let $\varepsilon\in\{1,-1\}$.
For $w\in A[\Gamma]$ let $\det_{c,\varepsilon}(w)\in A$ be the determinant of multiplication by $w$ on the free $A$-module $A[\Gamma]/(c-\varepsilon)$, with basis the images of a set of representatives of $\Gamma/\langle c\rangle$.
\begin{enumerate}
\item $\det_{c,\varepsilon}(w)$ lies in the ideal of $A[\Gamma]$ generated by $c-\varepsilon$ and $w$. So if $X$ is an $A[\Gamma]$-module and $x\in X$ satisfies $cx=\varepsilon x$ and $wx=0$, then $\det_{c,\varepsilon}(w)\,x=0$.
\item If $A\subset\mathbf C$, then $\det_{c,\varepsilon}(w)=\prod_\chi\chi(w)$ over the $|\Gamma|/2$ characters $\chi:\Gamma\to\mathbf C^\times$ with $\chi(c)=\varepsilon$.
\item Let $M\ge2$, let $b\mapsto\sigma_b$ be a surjection $(\mathbf Z/M)^\times\to\Gamma$ with $c=\sigma_{-1}$, let $k\ge1$, let $a\ge2$ be prime to $M$, and let $\varepsilon=(-1)^{k+1}$. Put $d=M(k+1)(k+2)!$ and $s=2a^{k+1}M^{k+1}(k+1)!$. Then $D=d^{|\Gamma|/2}\det_{c,\varepsilon}(y_{k+1}(a))$ is a nonzero integer with $|D|\le(ds)^{|\Gamma|/2}$.
\end{enumerate}
\end{lemma}
\begin{proof}
For (a), multiplication by $w$ on $A[\Gamma]/(c-\varepsilon)$ has a matrix $W$ in the given basis, and $\operatorname{adj}(W)W=\det(W)\cdot1$.
Lift the vector $\operatorname{adj}(W)$ applied to the coordinates of $1$ to an element $u\in A[\Gamma]$; then $wu\equiv\det_{c,\varepsilon}(w)$ modulo $c-\varepsilon$.
If $cx=\varepsilon x$ and $wx=0$, then $\det_{c,\varepsilon}(w)x=uwx=0$.

For (b), over $\mathbf C$ the algebra $\mathbf C[\Gamma]/(c-\varepsilon)$ is the product of the copies of $\mathbf C$ indexed by the characters with $\chi(c)=\varepsilon$, and $w$ acts on the copy of $\chi$ by $\chi(w)$.

For (c), $d\,y_{k+1}(a)$ has integer coefficients, so $D$ is an integer.
By (b), $\det_{c,\varepsilon}(y_{k+1}(a))$ is the product of $\chi(y_{k+1}(a))$ over the characters with $\chi(\sigma_{-1})=(-1)^{k+1}$.
For such $\chi$, $\chi(\theta_{k+1})$ is $-\tfrac1{k+1}$ times the generalized Bernoulli number of $\chi\circ\sigma$ at level $M$, which is nonzero because the parity of $\chi$ is that of $k+1$; and $1-a^{k+1}\chi(\sigma_a)\ne0$ because $a\ge2$.
So $D\ne0$.
Each factor has absolute value at most $s$, which gives $|D|\le(ds)^{|\Gamma|/2}$.
\end{proof}
\leanref{Development/AutomorphicArithmetic/AutomorphicVarieties/OddResidueBounds/AdjKill.lean}{UpperTheory.V.Odd.adjKill}
\leanref{Development/AutomorphicArithmetic/AutomorphicVarieties/StickelbergerCharacterEval/AlgebraMapSignedInvolutionQuotientDetMemSpan.lean}{UpperTheory.V.B5.algebraMap_signedInvolutionQuotientDet_mem_span}
\leanref{Development/AutomorphicArithmetic/AutomorphicVarieties/StickelbergerCharacterEval/Detprod.lean}{UpperTheory.V.B5.det_mul_quotientByOrderTwo_eq_prod_charEval}
\leanref{Development/AutomorphicArithmetic/AutomorphicVarieties/StickelbergerCharacterEval/Detint.lean}{UpperTheory.V.B5.higherStickelberger_det_exists_int_bounded}
\leanref{Development/AutomorphicArithmetic/AutomorphicVarieties/StickelbergerCharacterEval/GroupRingCharEvalStickelbergerYNeZeroOfParity.lean}{UpperTheory.V.B5.groupRingCharEval_stickelbergerY_ne_zero_of_parity}
\leanref{Development/AutomorphicArithmetic/AutomorphicVarieties/StickelbergerCharacterEval/NormCharEvalHigherStickelbergerLe.lean}{UpperTheory.V.B5.norm_charEval_higherStickelberger_le}

\begin{lemma}\label{sha-torsion-length}
Let $E$ be a number field that embeds in $\mathbf Q(\zeta_N)$ for some $N\ge1$, let $\ell$ be a prime and let $k$ be an integer.
Then the $\ell$-torsion of $H^1(G_{E,S_\ell},\mathbf Q_\ell/\mathbf Z_\ell(-k))$ has length at most $8\ell\,[E:\mathbf Q]\log(\ell N)$.
\end{lemma}
\begin{proof}
Let $F=E(\mu_\ell)$ and $m=[F:\mathbf Q]$.
Then $[F:E]$ divides $\ell-1$, so $m\le\ell[E:\mathbf Q]$, and $F$ embeds in $\mathbf Q(\zeta_{\ell N})$.
A class killed by $\ell$ is represented by a cocycle with values in the $\ell$-torsion $\mathbf Z/\ell$ of the coefficients, on which $\operatorname{Gal}(\bar F/F)$ acts trivially.
Its restriction to $\operatorname{Gal}(\bar F/F)$ is a homomorphism to $\mathbf Z/\ell$ that is unramified outside $\ell$, and restriction is injective on these classes because $[F:E]$ is prime to $\ell$.
By Kummer theory such homomorphisms inject into the group of $\alpha\in F^\times$ with $v(\alpha)\equiv0\pmod\ell$ at every finite place $v\nmid\ell$, modulo $F^{\times\ell}$.
That group has order at most $\ell^{2m}h_F$: its subgroup of $\ell$-units modulo $\ell$-th powers has order at most $\ell^{2m}$ by Dirichlet's unit theorem, and its quotient is contained in the $\ell$-torsion of the class group.
Every ideal class contains an integral ideal of norm at most $\sqrt{|d_F|}$, and counting these ideals gives $\log h_F\le\frac12\log|d_F|+m\log\big(1+\frac12\log|d_F|\big)$.
The conductor-discriminant formula gives $|d_F|\le(\ell N)^{m}$, and $m\le\ell N$, so $\log h_F\le3m\log(\ell N)$.
So a length $g$ of the $\ell$-torsion satisfies $g\log\ell\le2m\log\ell+3m\log(\ell N)$, and since $\log\ell\ge\log2>5/8$ this gives $g\le8\ell[E:\mathbf Q]\log(\ell N)$.
\end{proof}
\leanref{Development/AutomorphicArithmetic/AutomorphicVarieties/OddResidueBounds/Imports.lean}{UpperTheory.V.Odd.gen}
\leanref{Development/AutomorphicArithmetic/AutomorphicVarieties/OddResidueBounds/UnrInj.lean}{UpperTheory.V.Odd.unrInj}
\leanref{Development/AutomorphicArithmetic/AutomorphicVarieties/OddResidueBounds/UnrSel.lean}{UpperTheory.V.Odd.unrSel}
\leanref{Development/AutomorphicArithmetic/AutomorphicVarieties/OddResidueBounds/SelCard.lean}{UpperTheory.V.Odd.selCard}
\leanref{Development/AutomorphicArithmetic/AutomorphicVarieties/OddResidueBounds/LogH.lean}{UpperTheory.V.Odd.logH}

\begin{lemma}\label{sha-minus-exponent}
There is $C\ge0$ with the following property.
Let $E$ be a CM field that embeds in $\mathbf Q(\zeta_N)$ for some $N\ge1$, let $\ell$ be a prime, and suppose that $\ell$ is odd or that $\chi_{\mathrm{cyc}}\equiv1\pmod4$ on $\operatorname{Gal}(\bar E/E)$.
Let $k\ge1$.
Then $\text{Ш}^-_\ell(E,k)$ has finite length and
\[
 \operatorname{len}\text{Ш}^-_\ell(E,k)\cdot\log\ell\le C\,\ell\,[E:\mathbf Q]^2(k+1)\log(k\ell N)^2.
\]
One can take $C=240$.
\end{lemma}
\begin{proof}
Let $n=[E:\mathbf Q]$ and $M=N\ell$, so that $E$ embeds in $\mathbf Q(\zeta_M)$ and $\ell\mid M$.
Let $\rho:(\mathbf Z/M)^\times\to\operatorname{Gal}(E/\mathbf Q)$ be the restriction map, which is surjective; $\sigma_{-1}$ is complex conjugation, which is not the identity since $E$ is a CM field.
Put $a=2M+1$, which is odd and prime to $M$, and $\varepsilon=(-1)^{k+1}$.

\emph{An exponent for the minus part.}
Lemma~\ref{sha-stickelberger-annihilator} gives $z\in\mathbf Z_\ell[\operatorname{Gal}(E/\mathbf Q)]$, the image of $y_{k+1}(a)$, with $z\cdot\varphi(M)x=0$ for $x\in\text{Ш}^1_\ell(E,k)$.
Let $x\in\text{Ш}^-_\ell(E,k)$. Complex conjugation acts on $\varphi(M)x$ by $\varepsilon$, so Lemma~\ref{sha-stickelberger-determinant}(a) gives $\det_{\sigma_{-1},\varepsilon}(z)\,\varphi(M)x=0$.
Let $D$ be the integer of Lemma~\ref{sha-stickelberger-determinant}(c) and $r_0=v_\ell(D)$.
In $\mathbf Z_\ell$ we have $D=d^{n/2}\det_{\sigma_{-1},\varepsilon}(z)$, so $\det_{\sigma_{-1},\varepsilon}(z)$ divides $\ell^{r_0}$; and $\ell^{r_0}\le|D|\le(ds)^{n/2}$.
Let $r_1=v_\ell(\varphi(M))$, so that $\varphi(M)$ divides $\ell^{r_1}$ in $\mathbf Z_\ell$ and $\ell^{r_1}\le M$.
Hence $\ell^{r}$ kills $\text{Ш}^-_\ell(E,k)$ for $r=r_0+r_1$, and
\[
 r\log\ell\le n\log(ds)+\log M\le n\log(dsM)\le30\,n(k+1)\log(k\ell N),
\]
where the last step is the estimate $dsM\le(3k\ell N)^{10(k+1)}$ for $d=M(k+1)(k+2)!$, $s=2a^{k+1}M^{k+1}(k+1)!$ and $a=2M+1$.

\emph{The length.}
Let $g$ be the length of the $\ell$-torsion of $H^1(G_{E,S_\ell},\mathbf Q_\ell/\mathbf Z_\ell(-k))$.
A submodule killed by $\ell^r$ has a filtration by the kernels of $\ell,\ell^2,\dots,\ell^{r}$ whose $r$ graded pieces embed in that $\ell$-torsion.
So $\operatorname{len}\text{Ш}^-_\ell(E,k)\le rg$, which is finite.
By Lemma~\ref{sha-torsion-length}, $g\le8\ell n\log(\ell N)\le8\ell n\log(k\ell N)$, and so
\[
 \operatorname{len}\text{Ш}^-_\ell(E,k)\cdot\log\ell\le(r\log\ell)\,g\le240\,\ell\,n^2(k+1)\log(k\ell N)^2.
\]
\end{proof}
\leanref{Development/AutomorphicArithmetic/AutomorphicVarieties/OddResidueBounds/ZeroMemUnrHom.lean}{UpperTheory.V.Odd.StickBound}
\leanref{Development/AutomorphicArithmetic/AutomorphicVarieties/OddResidueBounds/ThmBcOfPsiBound.lean}{UpperTheory.V.Odd.stickBound}
\leanref{Development/AutomorphicArithmetic/AutomorphicVarieties/OddResidueBounds/KillPow.lean}{UpperTheory.V.Odd.killPow}
\leanref{Development/AutomorphicArithmetic/AutomorphicVarieties/OddResidueBounds/Imports.lean}{UpperTheory.V.Odd.gen}
\leanref{Development/AutomorphicArithmetic/AutomorphicVarieties/OddResidueBounds/LenPow.lean}{UpperTheory.V.Odd.lenPow}

\begin{theorem}\label{sha-minus-part}
There is $C$ with the following property.
Let $M$ be an abelian field datum of degree $o$, let $e,t\ge2$ be integers, let $(q,d)$ be a prime pair for $M$, and let $E$ be a CM field that is a compositum of $(M,e,q,d)$ (Definition~\ref{sha-abelian-field}), of degree $n$ and conductor $N$.
Then $E$ satisfies the minus bound (Definition~\ref{sha-bounded-field})
\[
 \Big(t,\;3nt^2\log(tN)+C\sum_{\substack{\ell\mid o\\ \ell\ \text{odd prime}}}\ell\,n^2t^2\log(t\ell N)^2\Big).
\]
\end{theorem}
\begin{proof}
Let $C$ be the constant of Lemma~\ref{sha-minus-exponent}.
The field $E$ embeds in $\mathbf Q(\zeta_N)$.
For finiteness, let $\ell$ be odd and $0<k<t$: if $\ell\mid n$ use Lemma~\ref{sha-minus-exponent}, and otherwise Lemma~\ref{sha-minus-prime-to-degree}.

Let $S$ be a finite set of odd primes, and split it into the primes $S_1$ not dividing $n$ and the primes $S_2$ dividing $n$.
For each $1\le k\le t-1$, Lemma~\ref{sha-minus-prime-to-degree} gives
\[
 \sum_{\ell\in S_1}\log\ell\cdot\operatorname{len}\text{Ш}^-_\ell(E,k)\le n\log\big(16N^{k+2}(k+1)!\big)\le n(2t+5)\log(tN),
\]
because $16\le(tN)^4$, $N^{k+2}\le(tN)^{t+1}$ and $(k+1)!\le(tN)^t$.
Summing over the $t-1$ values of $k$ and using $(t-1)(2t+5)\le3t^2$ bounds the contribution of $S_1$ by $3nt^2\log(tN)$.

Each $\ell\in S_2$ is odd and divides $n$, so it divides $o$ by Lemma~\ref{sha-composita}.
Lemma~\ref{sha-minus-exponent} gives $\log\ell\cdot\operatorname{len}\text{Ш}^-_\ell(E,k)\le C\ell n^2(k+1)\log(k\ell N)^2\le C\ell n^2t\log(t\ell N)^2$ for $k<t$.
Summing over $k$ gives at most $C\ell n^2t^2\log(t\ell N)^2$ for each $\ell\in S_2$, and $S_2$ is contained in the set of odd primes dividing $o$.
\end{proof}
\leanref{Development/AutomorphicArithmetic/AutomorphicVarieties/PrimeSetBoundDataOfPrimitiveChar.lean}{UpperTheory.V.ThmBcFree}
\leanref{Development/AutomorphicArithmetic/AutomorphicVarieties/StmtFields.lean}{UpperTheory.V.minusBound}
\leanref{Development/AutomorphicArithmetic/AutomorphicVarieties/OddResidueBounds/ThmBcOfPsiBound.lean}{UpperTheory.V.Odd.thmBcFree_of}

\section{The prime 2 and the bad primes}\label{section-sha-two}

\begin{theorem}[\cite{FW79}]\label{sha-ferrero-washington}
Let $\ell$ be a prime and let $K$ be an abelian number field that contains a primitive $\ell$-th root of unity, and a primitive fourth root of unity if $\ell=2$.
Then there is $r$ such that for every $n\ge0$ and every field $K'$ that is the $\ell^n$-th cyclotomic extension of $K$, the quotient $\operatorname{Cl}(K')/\operatorname{Cl}(K')^{\ell}$ has at most $\ell^r$ elements.
\end{theorem}
\leanref{Development/AutomorphicArithmetic/AutomorphicVarieties/PublishedTheorems/FerreroWashington.lean}{UpperTheory.V.Published.FerreroWashington}
\leanref{Development/AutomorphicArithmetic/AutomorphicVarieties/GaloisLibrary/FerreroWashingtonIndex/FerreroWashington.lean}{UpperTheory.V.FW.ferreroWashington}

Bounded $\ell$-rank of the class groups along the cyclotomic $\mathbf Z_\ell$-extension is equivalent to the vanishing of its Iwasawa $\mu$-invariant.
For abelian fields this vanishing is the theorem of Ferrero and Washington \cite{FW79}; see also \cite[Chapter 7]{Was97}.

\begin{definition}\label{sha-tower-data}
Let $\ell$ be a prime and let $L_0$ be a number field.
\begin{enumerate}
\item A \emph{layer of degree $\ell^{j}$} over $L_0$ is an extension $L_1/L_0$ that is the $\ell^{n}$-th cyclotomic extension of $L_0$ for some $n$, with $[L_1:L_0]=\ell^{j}$.
\item Let $R\ge0$ and $j_0=\lfloor\log_\ell R\rfloor+3$. The field $L_0$ has a \emph{tower bound with $R$ generators} if the following two conditions hold.
\begin{enumerate}
\item There are a field $K$ with an ultrametric absolute value extending that of $\mathbf Q_\ell$, integers $c\ge0$ and $R'\le R$, and elements $s_1,\dots,s_{R'}\in K$ such that $L_0$ contains a primitive $\ell^{c}$-th root of unity and: for every $k\ge1$, every layer $L_1$ of degree $\ell^{m}$ over $L_0$ and every integer $B\ge0$ with $|s_i-k|\ge\ell^{-B}$ for all $i$,
\[
 \operatorname{len}\text{Ш}^1_\ell(L_1,k)\le R'\big(\max(m,j_0)+c+B\big).
\]
\item For every $k\ge1$, every layer $L_1$ of degree $\ell^{m}$ and every layer $L_2$ of degree $\ell^{j_0}$ over $L_0$,
\[
 \operatorname{len}\text{Ш}^1_\ell(L_1,k)\le\operatorname{len}H^2\big(G_{L_2,S_\ell},\mathbf Z_\ell(k+1)\big)+R\max(m-j_0,0).
\]
\end{enumerate}
\item Let $M=(m,H)$ be an abelian field datum and $(q,d)$ a prime pair for $M$. A number field $L_0$ is a \emph{$2$-base} for $(M,q,d)$ if it contains a primitive fourth root of unity, $[L_0:\mathbf Q]\le8\deg M$, $L_0$ embeds in $\mathbf Q(\zeta_{4mqd})$, and for every $e\ge2$ and every compositum $E$ of $(M,e,q,d)$ there is a layer $L_1$ of degree $2^{j}$ over $L_0$, for $\ell=2$, with $j\le e+2$ and $\operatorname{len}\text{Ш}^1_2(E,k)\le1+\operatorname{len}\text{Ш}^1_2(L_1,k)$ for all $k\ge1$.
\end{enumerate}
\end{definition}
\leanref{Development/AutomorphicArithmetic/AutomorphicVarieties/StmtFields.lean}{UpperTheory.V.IsLayer}
\leanref{Development/AutomorphicArithmetic/AutomorphicVarieties/LayeredIsomorphisms/SecondBound.lean}{UpperTheory.V.Win.ABound}
\leanref{Development/AutomorphicArithmetic/AutomorphicVarieties/LayeredIsomorphisms/SecondBound.lean}{UpperTheory.V.Win.SecondBound}
\leanref{Development/AutomorphicArithmetic/AutomorphicVarieties/LayeredIsomorphisms/SecondBound.lean}{UpperTheory.V.Win.TwoBase}

\begin{definition}\label{sha-iwasawa-module}
Let $\ell$ be a prime and let $L$ be a number field that contains a primitive $\ell$-th root of unity, and a primitive fourth root of unity if $\ell=2$.
Let $L_\infty=L(\mu_{\ell^\infty})$ and $U_\infty=\operatorname{Gal}(\bar L/L_\infty)$.
Let $D(L)$ be the group of continuous homomorphisms $U_\infty\to\mathbf Q_\ell/\mathbf Z_\ell$ that vanish on the inertia groups at all places and on the decomposition groups at the places above $\ell$, and let $X'_\infty(L)=\operatorname{Hom}(D(L),\mathbf Q_\ell/\mathbf Z_\ell)$ be its Pontryagin dual.
Equivalently, $X'_\infty(L)$ is the Galois group over $L_\infty$ of its maximal abelian pro-$\ell$ extension that is unramified everywhere and in which the places above $\ell$ split completely.
The group $\operatorname{Gal}(\bar L/L)$ acts on $D(L)$ by $(\sigma\varphi)(u)=\varphi(\sigma^{-1}u\sigma)$, and on $X'_\infty(L)$ by the dual action; $U_\infty$ acts trivially.
\end{definition}
\leanref{Development/AutomorphicArithmetic/AutomorphicVarieties/NumberFieldTowers/XTwApplyApply.lean}{UpperTheory.V.Tower.XP}
\leanref{Development/AutomorphicArithmetic/AutomorphicVarieties/NumberFieldTowers/XTwApplyApply.lean}{UpperTheory.V.Tower.xTw}
\leanref{Development/AutomorphicArithmetic/AutomorphicVarieties/NumberFieldTowers/XTwApplyApply.lean}{UpperTheory.V.Tower.IsTower}

\begin{lemma}\label{sha-iwasawa-descent}
Let $\ell$ and $L$ be as in Definition~\ref{sha-iwasawa-module}, let $\gamma\in\operatorname{Gal}(\bar L/L)$ restrict to a topological generator of $\operatorname{Gal}(L_\infty/L)$, and put $u=\chi_{\mathrm{cyc}}(\gamma)$.
Let $k\ge1$ and let $L_1$ be a layer of degree $\ell^{m}$ over $L$ (Definition~\ref{sha-tower-data}).
Then $\operatorname{len}\text{Ш}^1_\ell(L_1,k)$ is the length of the cokernel of $(u^{k}\gamma)^{\ell^m}-1$ on $X'_\infty(L)$.
\end{lemma}
\begin{proof}
The field $L_1(\mu_{\ell^\infty})$ is $L_\infty$, and $\operatorname{Gal}(L_\infty/L_1)$ is topologically generated by $\gamma^{\ell^m}$.
Restriction to $U_\infty$ sends $\text{Ш}^1_\ell(L_1,k)$ into $D(L)$, since $U_\infty$ acts trivially on the coefficients.
It is injective by the argument of Lemma~\ref{sha-stickelberger-annihilator}, Step 1: the kernel is $H^1$ of the procyclic group $\operatorname{Gal}(L_\infty/L_1)$ with coefficients in the divisible module $\mathbf Q_\ell/\mathbf Z_\ell(-k)$, on which a generator acts by a scalar different from $1$.
Its image consists of the $\varphi\in D(L)$ fixed by $\gamma^{\ell^m}$ acting by $\varphi\mapsto\chi_{\mathrm{cyc}}(\gamma)^{-k\ell^m}\gamma^{\ell^m}\varphi$: such a $\varphi$ extends to a class on $\operatorname{Gal}(\bar L/L_1)$ because $H^2$ of a procyclic pro-$\ell$ group with torsion coefficients vanishes, and the extension satisfies the local conditions because the same vanishing of $H^1$ holds for the decomposition groups of $L_\infty/L_1$, which are infinite.
By Pontryagin duality, the kernel of $\chi_{\mathrm{cyc}}(\gamma)^{-k\ell^m}\gamma^{\ell^m}-1$ on $D(L)$ has the same length as the cokernel of the dual endomorphism $(u^k\gamma)^{\ell^m}-1$ on $X'_\infty(L)$.
\end{proof}
\leanref{Development/AutomorphicArithmetic/AutomorphicVarieties/AdditiveValuationNorms/Descends.lean}{UpperTheory.V.PropA.descends}
\leanref{Development/AutomorphicArithmetic/AutomorphicVarieties/NumberFieldTowers/Imports.lean}{UpperTheory.V.Tower.desc}
\leanref{Development/AutomorphicArithmetic/AutomorphicVarieties/NumberFieldTowers/ShaDGen.lean}{UpperTheory.V.Tower.shaDGen}
\leanref{Development/AutomorphicArithmetic/AutomorphicVarieties/NumberFieldTowers/CokerDual.lean}{UpperTheory.V.Tower.cokerDual}

\begin{lemma}\label{sha-cokernel-growth}
Let $\ell$ be a prime, let $R\ge1$, let $X$ be a $\mathbf Z_\ell$-module generated by $R$ elements, and let $g$ be an endomorphism of $X$ with $(g-1)^{n}X\subset\ell X$ for some $n$.
Put $j_0=\lfloor\log_\ell R\rfloor+3$.
Then for every $m\ge0$,
\[
 \operatorname{len}\operatorname{coker}\big(g^{\ell^m}-1\big)\le\operatorname{len}\operatorname{coker}\big(g^{\ell^{j_0}}-1\big)+R\max(m-j_0,0).
\]
\end{lemma}
\begin{proof}
The module $X/\ell X$ has length at most $R$ and $g-1$ is nilpotent on it, so $(g-1)^RX\subset\ell X$.
Let $j=\lfloor\log_\ell R\rfloor+1$, so that $R\le\ell^{j}$ and $(g-1)^{\ell^{j}}X\subset\ell X$.
Since $\ell$ divides the binomial coefficients $\binom{\ell^j}{i}$ for $0<i<\ell^j$, the endomorphism $g^{\ell^{j}}-1$ is congruent to $(g-1)^{\ell^j}$ modulo $\ell\operatorname{End}(X)$, so $g^{\ell^{j}}-1$ maps $X$ into $\ell X$.
If $w-1$ maps $X$ into $\ell^{i}X$ with $i\ge1$, then $w^\ell-1$ maps $X$ into $\ell^{i+1}X$; so $w=g^{\ell^{n}}$ satisfies $(w-1)X\subset\ell^2X$ for every $n\ge j+1$.
For such $w$ write $w^\ell-1=\big(\sum_{i<\ell}w^i\big)(w-1)$.
The first factor is $\ell$ times an endomorphism congruent to $1$ modulo $\ell$, hence $\ell$ times an automorphism, so its cokernel is $X/\ell X$, of length at most $R$.
The length of the cokernel of a composite is at most the sum of the lengths of the cokernels, so $\operatorname{len}\operatorname{coker}(w^\ell-1)\le\operatorname{len}\operatorname{coker}(w-1)+R$.
Applying this with $w=g^{\ell^{n}}$ for $n=j_0,j_0+1,\dots,m-1$, which all satisfy $n\ge j+1$, gives the claim for $m\ge j_0$; for $m<j_0$ the module $\operatorname{coker}(g^{\ell^m}-1)$ is a quotient of $\operatorname{coker}(g^{\ell^{j_0}}-1)$, because $g^{\ell^{j_0}}-1$ is a multiple of $g^{\ell^m}-1$.
\end{proof}
\leanref{Development/AutomorphicArithmetic/AutomorphicVarieties/AdditiveValuationNorms/Growth.lean}{UpperTheory.V.PropA.growth}
\leanref{Development/AutomorphicArithmetic/AutomorphicVarieties/AdditiveValuationNorms/UnipPow.lean}{UpperTheory.V.PropA.unipPow}
\leanref{Development/AutomorphicArithmetic/AutomorphicVarieties/AdditiveValuationNorms/ImportsMore.lean}{UpperTheory.V.PropA.step}

\begin{lemma}\label{sha-exceptional-centres}
Let $\ell$ be a prime, let $R\ge0$, let $X$ be a $\mathbf Z_\ell$-module generated by $R$ elements, let $\gamma$ be an endomorphism of $X$, and let $u\in\mathbf Z_\ell$ with $v_\ell(u-1)=c$, where $c\ge1$, and $c\ge2$ if $\ell=2$.
Let $j\ge0$.
Then there are $s_1,\dots,s_R\in\mathbf Z_\ell$ such that for every integer $k\ge0$,
\[
 \operatorname{len}\operatorname{coker}\big((u^{k}\gamma)^{\ell^j}-1\big)\le\sum_{i=1}^{R}\big(j+c+v_\ell(s_i-k)\big).
\]
\end{lemma}
\begin{proof}
Choose a surjection $\pi:\mathbf Z_\ell^R\to X$ and a matrix $S\in M_R(\mathbf Z_\ell)$ with $\pi\circ S=\gamma^{\ell^j}\circ\pi$.
For $a\in\mathbf Z_\ell$, $\pi$ induces a surjection from the cokernel of $aS-1$ on $\mathbf Z_\ell^R$ onto the cokernel of $a\gamma^{\ell^j}-1$ on $X$, and the first has length $v_\ell\det(aS-1)$.
Let $z_1,\dots,z_R\in\mathbf C_\ell$ be the eigenvalues of $S$, so that $\det(aS-1)=\prod_i(az_i-1)$.
Put $a_k=u^{k\ell^j}$.
By the binomial expansion of $(1+(u-1))^{\ell^j(k-k')}$, $|a_k-a_{k'}|=\ell^{-(j+c)}|k-k'|$ for all $k,k'\ge0$; this uses $c\ge1$, and $c\ge2$ if $\ell=2$.
Fix $i$ and put $f(k)=\ell^{j+c}|a_kz_i-1|$.
For units $\alpha,\beta$ of $\mathbf Z_\ell$, $|\alpha-\beta|\le\max(|\alpha z_i-1|,|\beta z_i-1|)$, so $|k-k'|\le\max(f(k),f(k'))$.
Thus any two of the closed balls of radius $f(k)$ about the integers $k$ meet, and by the ultrametric property and the compactness of $\mathbf Z_\ell$ they have a common point $s_i\in\mathbf Z_\ell$.
Then $\ell^{-(j+c)}|s_i-k|\le|a_kz_i-1|$ for every $k$, and taking the product over $i$ gives $v_\ell\det(a_kS-1)\le\sum_i\big(j+c+v_\ell(s_i-k)\big)$.
Finally $(u^k\gamma)^{\ell^j}=a_k\gamma^{\ell^j}$.
\end{proof}
\leanref{Development/AutomorphicArithmetic/AutomorphicVarieties/AdditiveValuationNorms/ExistsCentresLevel.lean}{UpperTheory.V.PropA.exists_centres_level}
\leanref{Development/AutomorphicArithmetic/AutomorphicVarieties/AdditiveValuationNorms/ImportsAlt.lean}{UpperTheory.V.PropA.lift}
\leanref{Development/AutomorphicArithmetic/AutomorphicVarieties/AdditiveValuationNorms/DetBound.lean}{UpperTheory.V.PropA.detBound}
\leanref{Development/AutomorphicArithmetic/AutomorphicVarieties/AdditiveValuationNorms/PadicUnitsDistanceLeMax.lean}{UpperTheory.V.PropA.ball}
\leanref{Development/AutomorphicArithmetic/AutomorphicVarieties/AdditiveValuationNorms/Helly.lean}{UpperTheory.V.PropA.helly}

\begin{lemma}\label{sha-tower-from-module}
Let $\ell$ be a prime, let $R\ge1$, and let $L_0$ be a totally complex number field that contains a primitive $\ell$-th root of unity, and a primitive fourth root of unity if $\ell=2$.
Suppose that $X'_\infty(L_0)$ (Definition~\ref{sha-iwasawa-module}) is generated by $R$ elements.
Then $L_0$ has a tower bound with $R$ generators (Definition~\ref{sha-tower-data}).
\end{lemma}
\begin{proof}
Let $\ell^{c}$ be the number of $\ell$-power roots of unity in $L_0$; then $c\ge1$, and $c\ge2$ if $\ell=2$.
The group $\chi_{\mathrm{cyc}}(\operatorname{Gal}(\bar L_0/L_0))$ is $1+\ell^{c}\mathbf Z_\ell$, so there is $\gamma$ restricting to a topological generator of $\operatorname{Gal}(L_{0,\infty}/L_0)$ with $u=\chi_{\mathrm{cyc}}(\gamma)$ of valuation $v_\ell(u-1)=c$.
Put $X=X'_\infty(L_0)$ and $j_0=\lfloor\log_\ell R\rfloor+3$.
The finite $\ell$-group $X/\ell X$ carries an action of the pro-$\ell$ group generated by $\gamma$, so $\gamma-1$ is nilpotent on it; since $u\equiv1\pmod\ell$, the same holds for $u^k\gamma$ for every $k$.
By Lemma~\ref{sha-iwasawa-descent} and Lemma~\ref{sha-cokernel-growth} applied to $g=u^k\gamma$, for every $k\ge1$ and every layer $L_1$ of degree $\ell^m$,
\[
 \operatorname{len}\text{Ш}^1_\ell(L_1,k)\le\operatorname{len}\operatorname{coker}\big((u^k\gamma)^{\ell^{j_0}}-1\big)+R\max(m-j_0,0).\tag{$*$}
\]
For condition (i), take $K=\mathbf Q_\ell$, $R'=R$, and the elements $s_1,\dots,s_R$ of Lemma~\ref{sha-exceptional-centres} for $j=j_0$.
If $|s_i-k|\ge\ell^{-B}$, that is $v_\ell(s_i-k)\le B$, for all $i$, then ($*$) and Lemma~\ref{sha-exceptional-centres} give $\operatorname{len}\text{Ш}^1_\ell(L_1,k)\le R(j_0+c+B)+R\max(m-j_0,0)=R(\max(m,j_0)+c+B)$.
For condition (ii), let $L_2$ be a layer of degree $\ell^{j_0}$. Lemma~\ref{sha-iwasawa-descent} identifies the first term of ($*$) with $\operatorname{len}\text{Ш}^1_\ell(L_2,k)$.
The field $L_2$ is totally complex, and by Poitou-Tate duality \cite[Chapter VIII, \S6]{NSW08} $\text{Ш}^1_\ell(L_2,k)$ is dual to the kernel of $H^2(G_{L_2,S_\ell},\mathbf Z_\ell(k+1))\to\prod_{v\mid\ell}H^2(G_{L_{2,v}},\mathbf Z_\ell(k+1))$, so its length is at most that of $H^2(G_{L_2,S_\ell},\mathbf Z_\ell(k+1))$.
\end{proof}
\leanref{Development/AutomorphicArithmetic/AutomorphicVarieties/AdditiveValuationNorms/ABoundOfIsTower.lean}{UpperTheory.V.PropA.aBound_of_isTower}
\leanref{Development/AutomorphicArithmetic/AutomorphicVarieties/AdditiveValuationNorms/ABoundOfDescends.lean}{UpperTheory.V.PropA.aBound_of_descends}
\leanref{Development/AutomorphicArithmetic/AutomorphicVarieties/AdditiveValuationNorms/GenLevel.lean}{UpperTheory.V.PropA.genLevel}
\leanref{Development/AutomorphicArithmetic/AutomorphicVarieties/AdditiveValuationNorms/ShaLeV.lean}{UpperTheory.V.PropA.shaLeV}

\begin{proposition}\label{sha-tower-bounds}
\begin{enumerate}
\item Let $\ell$ be an odd prime, let $A$ be an abelian field datum containing $\mu_\ell$, and let $L_0$ be a field of $A$ (Definition~\ref{sha-abelian-field}). Then $L_0$ has a tower bound with $R$ generators for some $R$ (Definition~\ref{sha-tower-data}).
\item Let $M=(m,H)$ be an abelian field datum. There is $R$ such that for every prime pair $(q,d)$ for $M$ there is a $2$-base $L_0$ for $(M,q,d)$ that has a tower bound for $\ell=2$ with $R$ generators.
\end{enumerate}
\end{proposition}
\begin{proof}
By Lemma~\ref{sha-tower-from-module} it suffices, in each case, to show that $L_0$ is totally complex, contains $\mu_\ell$ (and $\mu_4$ if $\ell=2$), and that $X'_\infty(L_0)$ is generated by $R$ elements.
By Nakayama's lemma, $X'_\infty(L_0)$ is generated by $R$ elements if $X'_\infty(L_0)/\ell$ has at most $\ell^R$ elements, that is, if the group of continuous homomorphisms $\operatorname{Gal}(\bar L_0/L_{0,\infty})\to\mathbf Z/\ell$ that are unramified everywhere and vanish on the decomposition groups above $\ell$ has at most $\ell^R$ elements.

(a) A field $L_0$ of $A$ contains $\mu_\ell$, and is totally complex because $\ell$ is odd.
By Theorem~\ref{sha-ferrero-washington} there is $r$ such that $\operatorname{Cl}(L_0(\zeta_{\ell^n}))/\ell$ has at most $\ell^{r}$ elements for every $n$.
A homomorphism $\operatorname{Gal}(\bar L_0/L_{0,\infty})\to\mathbf Z/\ell$ that is unramified everywhere comes from an unramified homomorphism on $\operatorname{Gal}(\bar L_0/L_0(\zeta_{\ell^n}))$ for some $n$, and by class field theory a set of at most $\ell^{r}$ such homomorphisms at each level exhausts them; so there are at most $\ell^{r}$.
Take $R=\max(1,r)$. Any two fields of $A$ are isomorphic, so $R$ depends only on $A$.

(b) For a prime $p$ let $Q_p=(p,(\mathbf Z/p)^{\times2})$, whose field is the quadratic subfield of $\mathbf Q(\zeta_p)$; for $p\equiv3\pmod4$ it is $\mathbf Q(\sqrt{-p})$.
For abelian field data $A=(m_A,H_A)$ and $B=(m_B,H_B)$ let $A\vee B=(m_Am_B,H)$, with $H$ the set of classes whose reductions lie in $H_A$ and in $H_B$; its fields are the composita of fields of $A$ and of $B$.
Let $L_0$ be a field of $M\vee(4,\{1\})\vee Q_q\vee Q_d$, the compositum of a field of $M$, $\mathbf Q(i)$, $\mathbf Q(\sqrt{-q})$ and $\mathbf Q(\sqrt{-d})$.
It contains $\mu_4$ and is therefore totally complex, $[L_0:\mathbf Q]\le2\cdot2\cdot2\deg M$, and $L_0$ embeds in $\mathbf Q(\zeta_{4mqd})$.

\emph{The comparison.} Let $e\ge2$ and let $E$ be a compositum of $(M,e,q,d)$, and let $L_1=L_0(\mu_{2^{e+2}})$, a layer of degree $2^{j}$ over $L_0$ with $j\le e+2$.
Comparing the congruence conditions that define $E$ and $L_1$ shows that $E\subset L_1$ and $[L_1:E]\mid2$.
The kernel of restriction from $H^1(G_{E,S_2},\mathbf Q_2/\mathbf Z_2(-k))$ to $H^1(G_{L_1,S_2},\mathbf Q_2/\mathbf Z_2(-k))$ is $H^1$ of $\operatorname{Gal}(L_1/E)$, of order at most $2$, with coefficients in a submodule of $\mathbf Q_2/\mathbf Z_2$, which has length at most $1$.
Restriction maps $\text{Ш}^1_2(E,k)$ into $\text{Ш}^1_2(L_1,k)$, so $\operatorname{len}\text{Ш}^1_2(E,k)\le1+\operatorname{len}\text{Ш}^1_2(L_1,k)$.
So $L_0$ is a $2$-base for $(M,q,d)$.

\emph{The generators.} Let $K$ be a field of $M\vee(4,\{1\})$, an abelian field containing $\mu_4$, and let $V=\operatorname{Gal}(\bar{\mathbf Q}/K(\mu_{2^\infty}))$.
As in (a), Theorem~\ref{sha-ferrero-washington} gives $g$, depending only on $M$, such that at most $2^{g}$ continuous homomorphisms $V\to\mathbf Z/2$ are unramified everywhere.
Allow ramification at $q$: a homomorphism $V\to\mathbf Z/2$ unramified outside $q$ is determined, modulo those unramified everywhere, by its restrictions to the inertia groups at the primes above $q$; each restriction factors through the procyclic tame inertia, and the primes of $K(\mu_{2^\infty})$ above $q$ are at most $2\deg M$ in number because $q\equiv3\pmod8$.
So at most $2^{g+2\deg M}$ homomorphisms are unramified outside $q$; in the same way, since $d\equiv7\pmod{16}$ gives at most $4\deg M$ primes above $d$, at most $2^{g+6\deg M}$ are unramified outside $qd$.
Passing from $V$ to the subgroups $V\cap\operatorname{Gal}(\bar{\mathbf Q}/\mathbf Q(\sqrt{-q}))$ and then $\operatorname{Gal}(\bar{\mathbf Q}/L_{0,\infty})$, each of index at most $2$, at most squares the number of homomorphisms unramified outside $qd$.
So there are at most $2^{4(g+6\deg M)}$ homomorphisms $\operatorname{Gal}(\bar L_0/L_{0,\infty})\to\mathbf Z/2$ unramified outside $qd$, and in particular at most that many unramified everywhere and split above $2$.
Take $R=4(g+6\deg M)$, which depends only on $M$.
\end{proof}
\leanref{Development/AutomorphicArithmetic/AutomorphicVarieties/LayeredIsomorphisms/SecondBound.lean}{UpperTheory.V.Win.FixedIn}
\leanref{Development/AutomorphicArithmetic/AutomorphicVarieties/LayeredIsomorphisms/SecondBound.lean}{UpperTheory.V.Win.TwoIn}
\leanref{Development/AutomorphicArithmetic/AutomorphicVarieties/QuadraticExtensions/IsAlgClosureRat.lean}{UpperTheory.V.Two.fixedIn}
\leanref{Development/AutomorphicArithmetic/AutomorphicVarieties/QuadraticExtensions/IsAlgClosureRat.lean}{UpperTheory.V.Two.twoIn}
\leanref{Development/AutomorphicArithmetic/AutomorphicVarieties/AdditiveValuationNorms/ABoundOfIsTower.lean}{UpperTheory.V.PropA.aBound_of_isTower}

\section{A bound of Soulé type}\label{section-sha-soule}

The bound for $H^2\big(G_{L,S_\ell},\mathbf Z_\ell(k+1)\big)$ that the prime $2$ and the bad primes need follows Soulé's method \cite{Sou79,Sou99}.
A class of the Tate-Shafarevich group gives homomorphisms from class groups of Kummer layers, and these are counted through elements whose twisted norms are detected, at primes of degree one, by cycles in the homology of $GL_N(\mathcal O_L)$.
The section first sets up the chains, then constructs the cycles, and ends with the bound.

\begin{definition}\label{sha-kummer-layers}
Let $\ell$ be a prime and let $L$ be a number field.
\begin{enumerate}
\item Let $1\le c\le\nu$ be integers, with $c\ge2$ if $\ell=2$. A \emph{Kummer layer} of type $(\ell,c,\nu)$ over $L$ is a Galois extension $F/L$ with a generator $\gamma$ of the cyclic group $\operatorname{Gal}(F/L)$, of order $\ell^{\nu-c}$, and a primitive $\ell^\nu$-th root of unity $\zeta\in\mathcal O_F^\times$ with $\gamma(\zeta)=\zeta^{1+\ell^{c}}$.
\item For such $F$ and integers $a,q\ge1$, let $\operatorname{Hom}_a(F,q)$ be the set of functions $\varphi$ from the set of nonzero primes of $\mathcal O_F$ to $\mathbf Z/q$ such that $\sum_P\operatorname{ord}_P(\alpha)\varphi(P)=0$ for every $\alpha\in F^\times$ and $\varphi(\gamma^{-1}P)=a\varphi(P)$ for every $P$. Equivalently, $\operatorname{Hom}_a(F,q)$ is the group of homomorphisms $\operatorname{Cl}(F)\to\mathbf Z/q$ on which $\gamma$ acts by $a$.
\item The \emph{twisted norm} of $\alpha\in F^\times$ is $N_a(\alpha)=\prod_{j=0}^{\ell^{\nu-c}-1}\gamma^{j}(\alpha)^{a^{j}}$.
\item For integers $s,q\ge1$ and $B\ge0$, the layer has \emph{few twisted norms} $(a,s,q,B)$ if every finite set $S\subset F^\times$ with the following two properties has at most $B$ elements: each $\alpha\in S$ has divisor $(\gamma-a)\mathfrak d$ for a divisor $\mathfrak d$ of $F$, that is $\operatorname{ord}_P(\alpha)=\operatorname{ord}_{\gamma P}(\mathfrak d)-a\operatorname{ord}_P(\mathfrak d)$ for every $P$; and two elements $\alpha,\alpha'\in S$ for which $N_a(\alpha)^{s}/N_a(\alpha')^{s}$ is a $q$-th power in $F^\times$ are equal.
\end{enumerate}
For a number field $L$, $r_2(L)$ is its number of complex places, $h_L$ its class number, $\operatorname{gen}(L)=2[L:\mathbf Q]+\lfloor\log_2h_L\rfloor+2$, and $t(L,N,m)$ the order of the torsion subgroup of the homology group $H_m(GL_N(\mathcal O_L),\mathbf Z)$.
\end{definition}
\leanref{Development/AutomorphicArithmetic/AutomorphicVarieties/ModularityArgument/ClassCounting/HasVal.lean}{UpperTheory.V.ZK.Lay}
\leanref{Development/AutomorphicArithmetic/AutomorphicVarieties/ModularityArgument/ClassCounting/HasVal.lean}{UpperTheory.V.ZK.EqHom}
\leanref{Development/AutomorphicArithmetic/AutomorphicVarieties/ModularityArgument/ClassCounting/HasVal.lean}{UpperTheory.V.ZK.nrm}
\leanref{Development/AutomorphicArithmetic/AutomorphicVarieties/ModularityArgument/ClassCounting/LemmaKOfSteps.lean}{UpperTheory.V.ZK.Few}
\leanref{Development/AutomorphicArithmetic/AutomorphicVarieties/ModularityArgument/ClassCounting/LemmaKOfSteps.lean}{UpperTheory.V.ZK.gens}
\leanref{Development/AutomorphicArithmetic/AutomorphicVarieties/ModularityArgument/ClassCounting/LemmaKOfSteps.lean}{UpperTheory.V.ZK.tc}

\begin{definition}\label{sha-bar-chains}
Let $A$ be a commutative ring and let $G$ be a group.
For $m\ge0$ let $C_m(G;A)$ be the free $A$-module on $G^m$, with basis elements written $[g_1|\dots|g_m]$.
\begin{enumerate}
\item The \emph{boundary} $d:C_{m+1}(G;A)\to C_m(G;A)$ is
\[
 d[g_1|\dots|g_{m+1}]=[g_2|\dots|g_{m+1}]+\sum_{j=1}^{m}(-1)^j[g_1|\dots|g_jg_{j+1}|\dots|g_{m+1}]+(-1)^{m+1}[g_1|\dots|g_m].
\]
These are the inhomogeneous chains that compute the group homology $H_*(G;A)$ with trivial coefficients. An element $z$ of $C_{m+1}(G;A)$ with $dz=0$ is a \emph{cycle}, and an element of $d\,C_{m+1}(G;A)$ is a \emph{boundary}.
\item A homomorphism $f:G\to H$ induces $f_*[g_1|\dots|g_m]=[f(g_1)|\dots|f(g_m)]$.
\item For $k\in G$ let $c_k:G\to G$ be $g\mapsto k^{-1}gk$, and let the \emph{prism operator} $h_k:C_m(G;A)\to C_{m+1}(G;A)$ be
\[
 h_k[g_1|\dots|g_m]=\sum_{i=0}^{m}(-1)^i[g_1|\dots|g_i|k|c_k(g_{i+1})|\dots|c_k(g_m)].
\]
\item Let $q\ge1$ and $e\ge0$, and write elements of the group $\mathbf Z/q$ additively inside the brackets. Put
\[
 z_{q,e}=\sum_{j\in(\mathbf Z/q)^{e+1}}[j_0|1|j_1|1|\dots|j_e|1]\in C_{2e+2}(\mathbf Z/q;\mathbf Z/q),
\]
and for an integer $t$ let $\Delta_t:\mathbf Z/q\to\mathbf Z/q\times\mathbf Z/q$ be $a\mapsto(a,ta)$.
\end{enumerate}
\end{definition}
\leanref{Development/AutomorphicArithmetic/AutomorphicVarieties/ModularityArgument/ClassCounting/HasVal.lean}{UpperTheory.V.ZK.Ch}
\leanref{Development/AutomorphicArithmetic/AutomorphicVarieties/ModularityArgument/ClassCounting/HasVal.lean}{UpperTheory.V.ZK.bd}
\leanref{Development/AutomorphicArithmetic/AutomorphicVarieties/ModularityArgument/ClassCounting/HasVal.lean}{UpperTheory.V.ZK.cm}
\leanref{Development/AutomorphicArithmetic/AutomorphicVarieties/ModularityArgument/ClassCounting/HasVal.lean}{UpperTheory.V.ZK.prism}
\leanref{Development/AutomorphicArithmetic/AutomorphicVarieties/ModularityArgument/ClassCounting/HasVal.lean}{UpperTheory.V.ZK.zc}
\leanref{Development/AutomorphicArithmetic/AutomorphicVarieties/ModularityArgument/ClassCounting/HasVal.lean}{UpperTheory.V.ZK.dt}

\begin{lemma}\label{sha-bar-identities}
Let $A$ be a commutative ring, let $G$ and $H$ be groups, let $f:G\to H$ be a homomorphism and let $k,k'\in G$.
\begin{enumerate}
\item $d\circ d=0$, $df_*=f_*d$, $(g\circ f)_*=g_*f_*$ for every homomorphism $g$ with source $H$, and $f_*h_k=h_{f(k)}f_*$. If $G$ is trivial, every cycle of positive degree is a boundary.
\item On $C_m(G;A)$ with $m\ge1$, $dh_k+h_kd=(c_k)_*-1$.
\item For every cycle $z$ of positive degree, $h_{kk'}z-h_kz-h_{k'}(c_k)_*z$ is a boundary.
\item For $q\ge1$ and $e\ge0$, $z_{q,e}$ is a cycle.
\end{enumerate}
\end{lemma}
\begin{proof}
(a) The complex $(C_*(G;A),d)$ is the complex of inhomogeneous chains of $G$ with coefficients in the trivial module $A$, so $d\circ d=0$.
The other identities hold on basis elements because $f$ is a homomorphism.
If $G$ is trivial, then $C_m(G;A)=A$ and $d:C_{m+1}\to C_m$ is multiplication by $1+\sum_{j=1}^{m+1}(-1)^j$, which is $0$ for $m$ even and $1$ for $m$ odd.
So a cycle of degree $m+1$ is $0$ if $m$ is odd, and is the boundary of the same element of $C_{m+2}=A$ if $m$ is even.

(b) Let $\tilde C_m$ be the free $A$-module on $G^{m+1}$, including the empty tuple for $m=-1$, with the boundary $\partial(x_0,\dots,x_m)=\sum_i(-1)^i(x_0,\dots,\widehat{x_i},\dots,x_m)$, and let $\mathrm{pr}:\tilde C_m\to C_m(G;A)$ send $(x_0,\dots,x_m)$ to $[x_0^{-1}x_1|\dots|x_{m-1}^{-1}x_m]$; it is surjective and commutes with the boundaries.
Let $R_k$ multiply every entry on the right by $k$; then $\mathrm{pr}\circ R_k=(c_k)_*\circ\mathrm{pr}$.
With the cone $\mathrm{cone}_a(x_0,\dots,x_m)=(a,x_0,\dots,x_m)$, define $\tilde h_k$ recursively by $\tilde h_k(x)=\mathrm{cone}_{x_0}\big(R_kx-\tilde h_k(x_1,\dots,x_m)\big)$, with $\tilde h_k$ zero on the empty tuple; explicitly $\tilde h_k(x_0,\dots,x_m)=\sum_i(-1)^i(x_0,\dots,x_i,x_ik,\dots,x_mk)$.
Induction on $m$ gives $\partial\tilde h_k+\tilde h_k\partial=R_k-1$ on $\tilde C_m$ for $m\ge0$, and $\mathrm{pr}\circ\tilde h_k=h_k\circ\mathrm{pr}$; applying $\mathrm{pr}$ gives (b).

(c) Define $\tilde d$ recursively by $\tilde d(x)=\mathrm{cone}_{x_0}\big(\tilde h_{k'}(R_kx)+\tilde d(x_1,\dots,x_m)\big)$, with $\tilde d$ zero on the empty tuple.
Induction on $m$ gives $\partial\tilde d+\tilde d\partial=-\big(\tilde h_{kk'}-\tilde h_k-\tilde h_{k'}R_k\big)$ on $\tilde C_m$ for $m\ge0$.
The chain $\mathrm{pr}(\tilde d\,y)$ depends only on $\mathrm{pr}(y)$, as one checks on the section $[g_1|\dots|g_m]\mapsto(1,g_1,g_1g_2,\dots,g_1\cdots g_m)$.
For a cycle $z=\mathrm{pr}(y)$ this gives $\mathrm{pr}(\tilde d\partial y)=0$, so $h_{kk'}z-h_kz-h_{k'}(c_k)_*z=-d\,\mathrm{pr}(\tilde dy)$.

(d) The first face, which deletes $j_0$, summed over $j$ gives every term $q$ times, hence $0$ in $\mathbf Z/q$.
For $0\le s\le e$, the face that merges the entries $j_s$ and $1$ in the term of $j$ equals the face that merges the entries $1$ and $j_{s+1}$ (or deletes the last entry $1$, if $s=e$) in the term of $j+\epsilon_s-\epsilon_{s+1}$, where the $\epsilon_s$ are the unit vectors and $\epsilon_{e+1}=0$.
These faces have opposite signs, so after reindexing the sum over $j$ they cancel in pairs.
\end{proof}
\leanref{Development/AutomorphicArithmetic/AutomorphicVarieties/ModularityArgument/ClassCounting/LemmaKOfSteps.lean}{UpperTheory.V.ZK.SChain}
\leanref{Development/AutomorphicArithmetic/AutomorphicVarieties/ModularityArgument/ClassCounting/SChain.lean}{UpperTheory.V.ZK.sChain}
\leanref{Development/AutomorphicArithmetic/AutomorphicVarieties/ModularityArgument/ClassCounting/SPrism.lean}{UpperTheory.V.ZK.sPrism}
\leanref{Development/AutomorphicArithmetic/AutomorphicVarieties/ModularityArgument/ClassCounting/SPrismMul.lean}{UpperTheory.V.ZK.sPrismMul}
\leanref{Development/AutomorphicArithmetic/AutomorphicVarieties/ModularityArgument/ClassCounting/FirstFace.lean}{UpperTheory.V.ZK.sZc}

\begin{lemma}\label{sha-bar-relation}
Let $\ell$ be a prime, let $\nu,e\ge0$ and let $q=\ell^\nu$.
Then there are $n_0,\dots,n_{2e+3}\in\mathbf Z/q$ and an index $t_0$ with $n_{t_0}$ a unit such that $\sum_{t=0}^{2e+3}n_t(\Delta_t)_*z_{q,e}$ is a boundary in $C_{2e+2}(\mathbf Z/q\times\mathbf Z/q;\mathbf Z/q)$.
\end{lemma}
\begin{proof}
Say that an abelian group is generated by $r$ elements if it is a quotient of $\mathbf Z^r$, and let $A=\mathbf Z/q$.

\emph{Step 1.} Let $G$ be a finite group and $p:G\to\mathbf Z/q$ a surjection with kernel $K$.
The sequence of $A[G]$-modules $0\to A\to A^{\mathbf Z/q}\to A^{\mathbf Z/q}\to A\to0$, given by the constant functions, $F\mapsto F(\cdot+1)-F$ and the sum, with $G$ acting on $A^{\mathbf Z/q}$ through translation by $p$, is exact.
By Shapiro's lemma $H_n(G;A^{\mathbf Z/q})\cong H_n(K;A)$, so the long exact sequences of its two short exact pieces show: if $H_{n+2}(K;A)$, $H_{n+1}(K;A)$ and $H_n(G;A)$ are generated by $a$, $b$ and $c$ elements, then $H_{n+2}(G;A)$ is generated by $a+b+c$ elements.

\emph{Step 2.} With $p$ the identity of $\mathbf Z/q$, so that $K$ is trivial, Step 1 and induction show that each $H_n(\mathbf Z/q;A)$ is generated by one element.
With $p$ the first projection of $\mathbf Z/q\times\mathbf Z/q$, whose kernel is isomorphic to $\mathbf Z/q$, Step 1 and induction show that $H_{2e+2}(\mathbf Z/q\times\mathbf Z/q;A)$ is generated by $2e+3$ elements.

\emph{Step 3.} The chains $(\Delta_t)_*z_{q,e}$ for $0\le t\le 2e+3$ are cycles by Lemma~\ref{sha-bar-identities}(a) and (d), and give $2e+4$ classes in a group generated by $2e+3$ elements.
Lift them to $\mathbf Z^{2e+3}$; the lifts are linearly dependent over $\mathbf Z$, and dividing a dependence by the largest power of $\ell$ dividing all its coefficients gives a relation with a coefficient prime to $\ell$.
A relation among homology classes says that the corresponding combination of cycles is a boundary.
\end{proof}
\leanref{Development/AutomorphicArithmetic/AutomorphicVarieties/ModularityArgument/ClassCounting/LemmaKOfSteps.lean}{UpperTheory.V.ZK.SRel}
\leanref{Development/EtaleCohomology/CyclotomicDiscriminants/FinSRel.lean}{LateralTheory.Fin.SRel.intersectionNumber_baseChange}
\leanref{Development/EtaleCohomology/CyclotomicDiscriminants/SRelTPf.lean}{LateralTheory.Q.SRelT.SRel_of}
\leanref{Development/EtaleCohomology/CyclotomicDiscriminants/SRelTPf.lean}{LateralTheory.Q.SRelT.step}
\leanref{Development/EtaleCohomology/CyclotomicDiscriminants/SRelTPf.lean}{LateralTheory.Q.SRelT.gen_two}

\begin{lemma}\label{sha-eta}
Let $\ell$ be a prime, let $1\le c\le\nu$ and $M<\ell^c$ be integers, let $q=\ell^\nu$ and $\kappa=1+\ell^c$, and write $\prod_{j=0}^{M-1}(X-\kappa^j)=\sum_{i=0}^{M}p_iX^i$.
For an integer $t$ let $\pi_t:(\mathbf Z/q)^2\to\mathbf Z/q$ be $(x,y)\mapsto x+ty$, and for a finitely supported function $\eta$ on $(\mathbf Z/q)^2$ let $(\pi_t)_*\eta(a)=\sum_{\pi_t(u)=a}\eta(u)$.
Let $0\le t_0\le M$.
Then there is $\eta:(\mathbf Z/q)^2\to\mathbf Z$ with $(\pi_t)_*\eta=0$ for $0\le t\le M$, $t\ne t_0$, and $(\pi_{t_0})_*\eta(a)=\sum_{i:\,\kappa^i=a}p_i$ for every $a\in\mathbf Z/q$, the sum over $0\le i\le M$.
\end{lemma}
\begin{proof}
Work in the group rings $R_2=\mathbf Z[(\mathbf Z/q)^2]$ and $R_1=\mathbf Z[\mathbf Z/q]$, with group elements written $[u]$; functions with finite support are elements of these rings, and $\pi_t$ induces a ring homomorphism $R_2\to R_1$.
Let $\iota:R_1\to R_2$ be induced by $a\mapsto(a,0)$, so that $\pi_t\circ\iota$ is the identity.
Put $w=\ell^c$ and $F=\sum_ip_i[\kappa^i]\in R_1$.

\emph{Step 1: $([w]-1)^M$ divides $F$.}
Let $E$ be the ring endomorphism $X\mapsto X^{1+w}$ of $\mathbf Z[X]$.
For every $g$, $X^w-1$ divides $E(g)-g$; and if $(X^w-1)^j$ divides $x$, then $(X^w-1)^{j+1}$ divides $E(x)-(1+w)^jx$, because $E(X^w-1)=(X^w-1)\sum_{i\le w}X^{wi}$ and the last sum is $\equiv1+w$ modulo $X^w-1$.
By induction on $M$, $(X^w-1)^M$ divides $\prod_{j<M}(E-\kappa^j)(X)=\sum_ip_iX^{\kappa^i}$.
Mapping $X$ to $[1]\in R_1$ gives Step 1.

\emph{Step 2.} For $t\ne t_0$ in $\{0,\dots,M\}$ we have $0<|t_0-t|<\ell^c$, so $t_0-t=\ell^au$ with $a<c$ and $u$ prime to $\ell$, and $t_0-t$ divides $w$ in $\mathbf Z/q$.
If $a$ divides $b$ in $\mathbf Z/q$, then $[a]-1$ divides $[b]-1$ in $R_1$.
So $\prod_{t\ne t_0}([t_0-t]-1)$ divides $([w]-1)^M$, and by Step 1 there is $Q\in R_1$ with $F=Q\prod_{t\ne t_0}([t_0-t]-1)$.

\emph{Step 3.} Put
\[
 \eta=\prod_{t'\ne t_0}\big([(0,1)]-[(t',0)]\big)\cdot\iota\Big(\Big[-\sum_{t'\ne t_0}t'\Big]Q\Big).
\]
For $t\ne t_0$ the factor with $t'=t$ maps to $[t]-[t]=0$ under $\pi_t$, so $\pi_t(\eta)=0$.
Under $\pi_{t_0}$ the factor of $t'$ maps to $[t_0]-[t']=[t']([t_0-t']-1)$, and the product of the $[t']$ cancels $[-\sum t']$, so $\pi_{t_0}(\eta)=Q\prod_{t'\ne t_0}([t_0-t']-1)=F$.
\end{proof}
\leanref{Development/AutomorphicArithmetic/AutomorphicVarieties/ModularityArgument/ClassCounting/LemmaKOfSteps.lean}{UpperTheory.V.ZK.SEta}
\leanref{Development/AutomorphicArithmetic/AutomorphicVarieties/ModularityArgument/ClassCounting/ExpStep.lean}{UpperTheory.V.ZK.sEta}

\begin{definition}\label{sha-additive-family}
Let $k$ be a finite field, let $A$ be a commutative ring and let $m\ge0$.
For a $k$-vector space $V$ let $GL(V)$ be its group of $k$-linear automorphisms.
\begin{enumerate}
\item An \emph{additive family} of degree $m+1$ over $k$ with coefficients in $A$ assigns to every finite-dimensional $k$-vector space $V$ an $A$-linear map $\theta_V:C_{m+1}(GL(V);A)\to A$ such that: $\theta_V$ vanishes on boundaries; for every isomorphism $\varphi:V\to W$ and every cycle $z$, $\theta_W((\varphi\,\cdot\,\varphi^{-1})_*z)=\theta_V(z)$; and for every group $G$, all homomorphisms $\rho:G\to GL(V)$ and $\sigma:G\to GL(W)$, and every cycle $z$ of $C_{m+1}(G;A)$, $\theta_{V\oplus W}((\rho\oplus\sigma)_*z)=\theta_V(\rho_*z)+\theta_W(\sigma_*z)$.
\item Let $q\ge1$ and $e\ge0$, let $\zeta\in k^\times$ with $\zeta^q=1$, and let $v:k^\times\to\mathbf Z/q$ be a homomorphism. For an integer $x\ge0$ let $\chi_x:\mathbf Z/q\to k^\times$ be $a\mapsto\zeta^{xa}$, and for a vector space $V$ let $s_V:k^\times\to GL(V)$ be the scalars. An additive family $\theta$ of degree $2e+3$ with coefficients in $\mathbf Z/q$ has \emph{value $v$} if
\[
 \theta_V\big(h_g\,(s_V\chi_x)_*z_{q,e}\big)=x^{e+1}\,v(\det g)
\]
for every $V$, every $x\ge0$ and every $g\in GL(V)$.
\item Let $R$ be a commutative ring and $B$ an $A$-module. A \emph{primitive family} of degree $m+1$ over $R$ with values in $B$ consists of $A$-linear maps $\Phi_n:C_{m+1}(GL_n(R);A)\to B$ for $n\ge0$ that vanish on boundaries and satisfy, for every cycle $z$: $\Phi_{n'}((i_{n,n'})_*z)=\Phi_n(z)$ for $n\le n'$, where $i_{n,n'}(g)=\operatorname{diag}(g,1_{n'-n})$; and $\Phi_{dn}((\psi_d)_*z)=d\,\Phi_n(z)$ for $d\ge1$, where $\psi_d(g)=\operatorname{diag}(g,\dots,g)$ with $d$ blocks.
\end{enumerate}
\end{definition}
\leanref{Development/AutomorphicArithmetic/AutomorphicVarieties/ModularityArgument/ClassCounting/HasVal.lean}{UpperTheory.V.ZK.Fam}
\leanref{Development/AutomorphicArithmetic/AutomorphicVarieties/ModularityArgument/ClassCounting/HasVal.lean}{UpperTheory.V.ZK.IsAdd}
\leanref{Development/AutomorphicArithmetic/AutomorphicVarieties/ModularityArgument/ClassCounting/HasVal.lean}{UpperTheory.V.ZK.HasVal}
\leanref{Development/AutomorphicArithmetic/AutomorphicVarieties/ModularityArgument/ClassCounting/HasVal.lean}{UpperTheory.V.ZK.chi}
\leanref{Development/AutomorphicArithmetic/AutomorphicVarieties/ModularityArgument/ClassCounting/HasVal.lean}{UpperTheory.V.ZK.psi}
\leanref{Development/AutomorphicArithmetic/AutomorphicVarieties/ModularityArgument/ClassCounting/LemmaKOfSteps.lean}{UpperTheory.V.ZK.IsPrimFam}

\begin{lemma}\label{sha-transfer}
Let $k$ be a finite field, let $A$ be a commutative ring in which $|k|=1$, let $m\ge0$, and let $t:C_{m+1}(k^\times;A)\to A$ be an $A$-linear map that vanishes on boundaries.
Then there is an additive family $\theta$ of degree $m+1$ over $k$ with coefficients in $A$ such that $\theta_k((s_k)_*z)=t(z)$ for every cycle $z$ of $C_{m+1}(k^\times;A)$.
\end{lemma}
\begin{proof}
\emph{The transfer.} Let $V$ be a finite-dimensional $k$-vector space and $P(V)$ its set of lines.
Choose $\lambda:V\smallsetminus\{0\}\to k^\times$ with $\lambda(uv)=u\lambda(v)$, by choosing a value on one vector of each line.
For a group $G$ and $\rho:G\to GL(V)$, let $G$ act on $P(V)$ on the right by $x\cdot g=\rho(g)^{-1}x$, and put $c(g,x)=\lambda(\tilde x)/\lambda(\rho(g)^{-1}\tilde x)$ for any nonzero $\tilde x\in x$.
Then $c(gh,x)=c(g,x)\,c(h,x\cdot g)$, and
\[
 T_\rho[g_1|\dots|g_m]=\sum_{x\in P(V)}\big[c(g_1,x)\,\big|\,c(g_2,x\cdot g_1)\,\big|\dots\big|\,c(g_m,x\cdot g_1\cdots g_{m-1})\big]
\]
defines a map $C_m(G;A)\to C_m(k^\times;A)$ that commutes with $d$, by the cocycle identity.
Another choice $\lambda'$ changes $c$ by the coboundary of $x\mapsto\lambda'(\tilde x)/\lambda(\tilde x)$, and changes $T_\rho(z)$ by a boundary when $z$ is a cycle.

\emph{The family.} Put $\theta_V=t\circ T_{\mathrm{id}}$ for the identity of $GL(V)$.
It vanishes on boundaries because $T$ commutes with $d$.
For $\rho:G\to GL(V)$ and a cycle $z$, $\theta_V(\rho_*z)=t(T_\rho z)$, since $T_{\mathrm{id}}\rho_*$ and $T_\rho$ are given by the same cocycle.
An isomorphism $\varphi:V\to W$ maps lines to lines and transports $\lambda$ to a homogeneous function on $W$, so isomorphism invariance follows from the independence of $\lambda$.
For $V=k$ the set $P(k)$ is a point and $c(u,\mathrm{pt})=u$, so $T_{s_k}$ is the identity and $\theta_k((s_k)_*z)=t(z)$.

\emph{Additivity.} Let $\rho:G\to GL(V)$ and $\sigma:G\to GL(W)$.
Map $P(V\oplus W)$ to $P(V)\sqcup P(W)$ by $[v,w]\mapsto[v]$ if $v\ne0$ and $[0,w]\mapsto[w]$; this is compatible with the actions.
Choose $\lambda(v,w)=\lambda_V(v)$ for $v\ne0$ and $\lambda(0,w)=\lambda_W(w)$.
Then the cocycle on $P(V\oplus W)$ is pulled back from those on $P(V)$ and $P(W)$, the fibre over a point of $P(V)$ has $|W|=|k|^{\dim W}$ elements, and the fibre over a point of $P(W)$ has one.
So $T_{\rho\oplus\sigma}(z)=|W|\,T_\rho(z)+T_\sigma(z)=T_\rho(z)+T_\sigma(z)$ in $C_{m+1}(k^\times;A)$, because $|k|=1$ in $A$.
\end{proof}
\leanref{Development/AutomorphicArithmetic/AutomorphicVarieties/ModularityArgument/ClassCounting/LemmaKOfSteps.lean}{UpperTheory.V.ZK.STrans}
\leanref{Development/AutomorphicArithmetic/AutomorphicVarieties/ModularityArgument/ClassCounting/STrans.lean}{UpperTheory.V.ZK.sTrans}

\begin{lemma}\label{sha-cocycle}
Let $k$ be a finite field, let $q\ge1$ and $e\ge0$, let $\zeta\in k^\times$ be a primitive $q$-th root of unity, and let $v:k^\times\to\mathbf Z/q$ be a homomorphism.
Then there is a $\mathbf Z/q$-linear map $t:C_{2e+3}(k^\times;\mathbf Z/q)\to\mathbf Z/q$ that vanishes on boundaries and satisfies $t\big(h_g(\chi_x)_*z_{q,e}\big)=x^{e+1}v(g)$ for every integer $x\ge0$ and every $g\in k^\times$.
\end{lemma}
\begin{proof}
\emph{A carry cocycle.} The group $k^\times$ is cyclic.
Let $y:k^\times\to\{0,\dots,|k^\times|-1\}$ be the discrete logarithm for a generator, let $r=\gcd(|k^\times|,y(\zeta))$, so that $|k^\times|=qr$, and write $r=a|k^\times|+b\,y(\zeta)$ with integers $a,b$.
Then $\lambda=b\,y$ satisfies $\lambda(g)+\lambda(g')-\lambda(gg')\in qr\mathbf Z$ for all $g,g'$, and $\lambda(\zeta^x)\equiv rx\pmod{qr}$.
Put $u(g,g')=(\lambda(g)+\lambda(g')-\lambda(gg'))/(qr)$ modulo $q$, a symmetric $2$-cocycle on $k^\times$.

\emph{The cochain.} Define $t$ on basis elements by
\[
 t[g_1|\dots|g_{2e+3}]=u(g_1,g_2)\,u(g_3,g_4)\cdots u(g_{2e+1},g_{2e+2})\,v(g_{2e+3}),
\]
the cup product of $e+1$ copies of $u$ with $v$.
It vanishes on boundaries: by induction on the degree, using the cocycle identity of $u$ and the additivity of $v$.
Since $k^\times$ is abelian, $c_g$ is the identity and $h_g$ inserts $g$ in every position with alternating signs; by the symmetry of $u$, the alternating sum of $t$ over the insertions of $g$ into $[g_1|\dots|g_{2e+2}]$ is $u(g_1,g_2)\cdots u(g_{2e+1},g_{2e+2})\,v(g)$.

\emph{The value.} Hence
\[
 t\big(h_g(\chi_x)_*z_{q,e}\big)=\sum_{j\in(\mathbf Z/q)^{e+1}}\prod_{s=0}^{e}u(\zeta^{xj_s},\zeta^x)\,v(g)=\Big(\sum_{a\in\mathbf Z/q}u(\zeta^{xa},\zeta^x)\Big)^{e+1}v(g).
\]
The sum $\sum_a\big(\lambda(\zeta^{xa})+\lambda(\zeta^x)-\lambda(\zeta^{x(a+1)})\big)$ telescopes to $q\lambda(\zeta^x)$, so $\sum_au(\zeta^{xa},\zeta^x)=\lambda(\zeta^x)/r\equiv x\pmod q$.
\end{proof}
\leanref{Development/AutomorphicArithmetic/AutomorphicVarieties/ModularityArgument/ClassCounting/LemmaKOfSteps.lean}{UpperTheory.V.ZK.SCoc}
\leanref{Development/AutomorphicArithmetic/AutomorphicVarieties/ModularityArgument/ClassCounting/FProd.lean}{UpperTheory.V.ZK.sCoc}

\begin{lemma}\label{sha-theta}
Let $k$ be a finite field, let $q\ge1$ and $e\ge0$, let $\zeta\in k^\times$ be a primitive $q$-th root of unity, and let $v:k^\times\to\mathbf Z/q$ be a homomorphism.
Then there is an additive family of degree $2e+3$ over $k$ with coefficients in $\mathbf Z/q$ that has value $v$ (Definition~\ref{sha-additive-family}).
\end{lemma}
\begin{proof}
Since $q$ divides $|k^\times|$, $|k|=1$ in $\mathbf Z/q$.
Lemma~\ref{sha-transfer} applied to the map $t$ of Lemma~\ref{sha-cocycle} gives an additive family $\theta$ with $\theta_k((s_k)_*z)=t(z)$ for cycles $z$.
If $q=1$ there is nothing more to prove, so let $q>1$; then $\zeta\ne0,1$ in $k$.

Fix $x\ge0$, and for a vector space $V$ put $w_V=(s_V\chi_x)_*z_{q,e}$, a cycle by Lemma~\ref{sha-bar-identities}(a),(d).
For $g\in GL(V)$, $c_g$ fixes the scalars, so Lemma~\ref{sha-bar-identities}(b) shows that $h_gw_V$ is a cycle.
Put $F_V(g)=\theta_V(h_gw_V)$.
By Lemma~\ref{sha-bar-identities}(c) and the vanishing of $\theta_V$ on boundaries, $F_V:GL(V)\to\mathbf Z/q$ is a homomorphism.
Isomorphism invariance and additivity of $\theta$, with Lemma~\ref{sha-bar-identities}(a), give $F_W(\varphi g\varphi^{-1})=F_V(g)$ and $F_{V\oplus W}(g\oplus g')=F_V(g)+F_W(g')$.
For $V=k$ and $g=u\in k^\times$, $h_uw_k=(s_k)_*h_u(\chi_x)_*z_{q,e}$, so $F_k(u)=t(h_u(\chi_x)_*z_{q,e})=x^{e+1}v(u)$.

A homomorphism from $GL_n(k)$ to an abelian group kills every elementary matrix $1+be_{ij}$: conjugation by $\operatorname{diag}(\zeta,1,\dots,1)$ or its analogue multiplies $b$ by $\zeta$, so the image of $1+b(\zeta-1)e_{ij}$ is trivial, and $\zeta-1$ is a unit.
Since $GL_n(k)$ is generated by elementary matrices and $\operatorname{diag}(u,1,\dots,1)$, $F_V(g)=F_V(\operatorname{diag}(\det g,1,\dots,1))=F_k(\det g)$ in a basis.
\end{proof}
\leanref{Development/AutomorphicArithmetic/AutomorphicVarieties/ModularityArgument/ClassCounting/LemmaKOfSteps.lean}{UpperTheory.V.ZK.STheta}
\leanref{Development/AutomorphicArithmetic/AutomorphicVarieties/ModularityArgument/ClassCounting/SThetaOf.lean}{UpperTheory.V.ZK.sThetaOf}

\begin{lemma}\label{sha-primitive}
Let $k$ be a finite field, let $A$ be a commutative ring, let $m\ge0$, and let $\theta$ be an additive family of degree $m+1$ over $k$ with coefficients in $A$.
Let $R$ be a commutative ring and $r:R\to k$ a ring homomorphism, and let $r_n:GL_n(R)\to GL_n(k)$ be the induced maps.
Then $\Phi_n=\theta_{k^n}\circ(r_n)_*$ is a primitive family of degree $m+1$ over $R$, with coefficients and values in $A$.
\end{lemma}
\begin{proof}
$\Phi_n$ vanishes on boundaries because $(r_n)_*$ commutes with $d$.
The decomposition $k^{n'}=k^n\oplus k^{n'-n}$ carries $r_{n'}\circ i_{n,n'}$ to $r_n\oplus1$, and $\theta$ vanishes on the cycles pushed forward by a trivial representation, since these come from the trivial group (Lemma~\ref{sha-bar-identities}(a)); isomorphism invariance and additivity give $\Phi_{n'}((i_{n,n'})_*z)=\Phi_n(z)$.
The decomposition $k^{dn}=(k^n)^{d}$ carries $r_{dn}\circ\psi_d$ to the sum of $d$ copies of $r_n$, and additivity gives $\Phi_{dn}((\psi_d)_*z)=d\,\Phi_n(z)$.
\end{proof}
\leanref{Development/AutomorphicArithmetic/AutomorphicVarieties/ModularityArgument/ClassCounting/LemmaKOfSteps.lean}{UpperTheory.V.ZK.SPhi}
\leanref{Development/AutomorphicArithmetic/AutomorphicVarieties/ModularityArgument/ClassCounting/PiRepApply.lean}{UpperTheory.V.ZK.sPhi}

\begin{lemma}\label{sha-steinitz}
Let $F$ be a number field, let $J$ be a finite set, and for $j\in J$ let $\mathfrak a_j,\mathfrak b_j$ be divisors of $F$; write $\Lambda(\mathfrak a)\subset F^J$ for the set of $x$ with $\operatorname{ord}_P(x_j)\ge\mathfrak a_j(P)$ for all $j$ and all primes $P$.
Let $\beta\in F^\times$ with $\operatorname{div}\beta=\sum_j(\mathfrak b_j-\mathfrak a_j)$, and suppose that $J$ is nonempty or $\beta=1$.
Then there is an $F$-linear automorphism $g$ of $F^J$ with $\det g=\beta$ and $g\Lambda(\mathfrak a)=\Lambda(\mathfrak b)$.
\end{lemma}
\begin{proof}
If $J$ is empty, take $g=1$. Otherwise fix $j_0\in J$.
Multiplying the coordinate $j_0$ by $\beta$ has determinant $\beta$ and carries $\Lambda(\mathfrak a)$ to $\Lambda(\mathfrak a')$, where $\mathfrak a'$ differs from $\mathfrak a$ only by $\mathfrak a'_{j_0}=\mathfrak a_{j_0}+\operatorname{div}\beta$; then $\sum_j\mathfrak a'_j=\sum_j\mathfrak b_j$.
For $i\ne j$ and a divisor $\delta$, there is a matrix of determinant $1$ acting on the coordinates $i,j$ that carries $\Lambda(\mathfrak a')$ to the lattice with $\mathfrak a'_i+\delta$ and $\mathfrak a'_j-\delta$ in those coordinates: choose $p,r$ with $\operatorname{ord}_P$ prescribed at the finitely many primes involved and complete $(p,r)$ to a matrix $\begin{pmatrix}p&s\\r&t\end{pmatrix}$ of determinant $1$ with entries in the appropriate fractional ideals, by the Chinese remainder theorem.
Moving divisors one coordinate at a time carries $\mathfrak a'$ to $\mathfrak b$.
\end{proof}
\leanref{Development/AutomorphicArithmetic/AutomorphicVarieties/ModularityArgument/ClassCounting/LemmaKOfSteps.lean}{UpperTheory.V.ZK.SStein}
\leanref{Development/EtaleCohomology/CyclotomicDiscriminants/FinSStein.lean}{LateralTheory.Fin.SStein.intersectionNumber_baseChange}
\leanref{Development/EtaleCohomology/CyclotomicDiscriminants/SSteinTPf.lean}{LateralTheory.Q.SSteinT.SStein_of}

\begin{lemma}\label{sha-detect}
Let $L\subset F$ be number fields, let $\ell$ be a prime and $\nu\ge1$, and suppose that $\mathcal O_F$ contains a primitive $\ell^\nu$-th root of unity.
Let $\beta\in F^\times$ and let $b\in\mathcal O_F$ be nonzero.
Suppose that for every prime $P$ of $F$ with $b\notin P$ and $\mathcal O_L\to\mathcal O_F/P$ surjective, the residue of $\beta$ at $P$ is defined and is an $\ell^\nu$-th power in $(\mathcal O_F/P)^\times$.
Then $\beta$ is an $\ell^\nu$-th power in $F^\times$.
\end{lemma}
\begin{proof}
Let $x\in\bar F$ with $x^{\ell^\nu}=\beta$, and let $\sigma\in\operatorname{Gal}(\bar F/F)$.
The set of $\tau$ with $\tau(x)=\sigma(x)$ is open, so by Theorem~\ref{lit-chebotarev-degree-one} it contains a Frobenius element at a place $v$ of $F$ of absolute residue degree one, prime to $b\ell$.
Its residue field has prime order, so $\mathcal O_L\to\mathcal O_F/P_v$ is surjective, and $\beta\equiv s^{\ell^\nu}$ at $v$ for some $s$.
The Frobenius $\tau$ satisfies $\tau(x)=\xi x$ for an $\ell^\nu$-th root of unity $\xi$, and $\tau(x)\equiv x^{p}$ modulo the place above $v$, where $p$ is the residue characteristic; since $x\equiv\xi'\tilde s$ for a root of unity $\xi'$ and a lift $\tilde s\in\mathcal O_F$ of $s$, this gives $\xi\equiv1$, and distinct $\ell^\nu$-th roots of unity remain distinct modulo a place prime to $\ell$.
So $\xi=1$ and $\sigma(x)=\tau(x)=x$.
Hence $x\in F$.
\end{proof}
\leanref{Development/AutomorphicArithmetic/AutomorphicVarieties/ModularityArgument/ClassCounting/LemmaKOfSteps.lean}{UpperTheory.V.ZK.SDetect}
\leanref{Development/EtaleCohomology/CyclotomicDiscriminants/FinSDetect.lean}{LateralTheory.Fin.SDetect.intersectionNumber_baseChange}
\leanref{Development/EtaleCohomology/CyclotomicDiscriminants/SDetectCut.lean}{LateralTheory.Q.SDetectCut.goal}

\begin{proposition}\label{sha-class-cycle}
Let $L\subset F$ be number fields, let $\ell$ be a prime, let $1\le c\le\nu$ and $e\ge0$ be integers with $2e+3<\ell^c$, and let $F/L$ with $\gamma$ and $\zeta\in\mathcal O_F^\times$ be a Kummer layer of type $(\ell,c,\nu)$ (Definition~\ref{sha-kummer-layers}).
Put $\kappa=1+\ell^c$, $a=\kappa^{e+1}$ and
\[
 p_e=\prod_{0\le j\le2e+2,\ j\ne e+1}\big(\kappa^{e+1}-\kappa^{j}\big)\in\mathbf Z.
\]
Let $\alpha\in F^\times$ have divisor $(\gamma-a)\mathfrak d$ for a divisor $\mathfrak d$ of $F$, that is $\operatorname{ord}_P(\alpha)=\mathfrak d(\gamma P)-a\,\mathfrak d(P)$ for every prime $P$.
Then there are $N\ge0$, a cycle $y\in C_{2e+3}(GL_N(\mathcal O_L);\mathbf Z/\ell^\nu)$ and a nonzero $b\in\mathcal O_F$ with the following property.
Let $k$ be a finite field and $r:\mathcal O_F\to k$ a ring homomorphism with $r(\mathcal O_L)=k$ and $r(b)\ne0$; let $v:k^\times\to\mathbf Z/\ell^\nu$ be a homomorphism, and let $\theta$ be an additive family of degree $2e+3$ over $k$ with coefficients in $\mathbf Z/\ell^\nu$ that has value $v$ relative to $r(\zeta)$ (Definition~\ref{sha-additive-family}).
Then $N_a(\alpha)$ can be written $a'/b'$ with $a',b'\in\mathcal O_F$, $r(b')\ne0$ and $r(a')\ne0$, and with $u=r(a')/r(b')$,
\[
 \theta_{k^N}\big((r_N)_*y\big)=p_e\,v(u),
\]
where $r_N:GL_N(\mathcal O_L)\to GL_N(k)$ is induced by $r$.
\end{proposition}
\begin{proof}
Write $q=\ell^\nu$, $\Gamma=\operatorname{Gal}(F/L)$, $J=(\mathbf Z/q)^2$ and $G=GL_L(F^J)$, the $L$-linear automorphisms of $F^J$.
For divisors $\mathfrak a=(\mathfrak a_j)_{j\in J}$ let $\Lambda(\mathfrak a)$ be as in Lemma~\ref{sha-steinitz}.
For nonzero $b\in\mathcal O_L$ let $G_b\subset G$ be the subgroup of $g$ such that $g$ and $g^{-1}$ preserve the set of vectors whose coordinates are integral at every prime not containing $b$, and let $H_b\subset G_b$ be the stabilizer of $\Lambda(0)=\mathcal O_F^J$.
Let $\rho:\mathbf Z/q\times\mathbf Z/q\to G_b$ act on the coordinate $j=(j_1,j_2)$ by $\zeta^{u_1j_1+u_2j_2}$ at $(u_1,u_2)$; then $\rho\circ\Delta_t$ acts on the coordinate $j$ through the character $\zeta^{\pi_t(j)}$, with $\pi_t$ as in Lemma~\ref{sha-eta}.
For $\sigma\in\Gamma$ let $\kappa(\sigma)$ be an integer with $\sigma(\zeta)=\zeta^{\kappa(\sigma)}$; then $\kappa(\gamma^i)\equiv\kappa^i\pmod q$.

\emph{Step 1: the element $b$.}
Since $F/L$ is Galois, there is a nonzero $b_0\in\mathcal O_L$ such that every prime of $L$ not containing $b_0$ is unramified in $F$; then for every $r$ as in the statement with $r(b_0)\ne0$, the prime $\mathfrak p=\ker(r)\cap\mathcal O_L$ splits completely in $F$ and $x\mapsto(r(\sigma x))_{\sigma\in\Gamma}$ is a surjection $\mathcal O_F\to k^\Gamma$ with kernel $\mathfrak p\mathcal O_F$.
Choose a nonzero $b\in\mathcal O_L$ divisible by $b_0$ and contained in every prime in the support of $\mathfrak d$; its image in $\mathcal O_F$ is the element $b$ of the statement.

\emph{Step 2: the combinatorial data.}
Lemma~\ref{sha-bar-relation} gives $n_t\in\mathbf Z/q$ for $0\le t\le2e+3$ and $t_0$ with $n_{t_0}$ a unit and $\sum_tn_t(\Delta_t)_*z_{q,e}$ a boundary.
Lemma~\ref{sha-eta} with $M=2e+3<\ell^c$ and this $t_0$ gives $\eta:J\to\mathbf Z$.
Write $\prod_{0\le j\le2e+2,\,j\ne e+1}(X-\kappa^j)=\sum_ie_iX^i$ and put $\beta=\prod_i\gamma^{-i}(\alpha)^{e_i}$.
Since $\operatorname{div}\alpha=(\gamma-a)\mathfrak d$ and $\operatorname{ord}_P(\gamma^{-i}\alpha)=\operatorname{ord}_{\gamma^iP}(\alpha)$, we get $\operatorname{ord}_P(\beta)=\sum_ip_i\,\mathfrak d(\gamma^iP)$, where $\sum_ip_iX^i=\prod_{j=0}^{2e+2}(X-\kappa^j)$.

\emph{Step 3: lattice elements.}
For each $t$ there is $g_t\in G$ with $g_t\Lambda(0)=\Lambda(\eta\mathfrak d)$, where $(\eta\mathfrak d)_j=\eta(j)\mathfrak d$, in the following \emph{block form}: there are $s:J\to\Gamma$ and $\beta_t:J\to\mathbf Z/q$ with $s_j(\zeta)^{\pi_t(j)}=\zeta^{\beta_t(j)}$, and $F$-linear automorphisms $g'_x$ of $F^{\beta_t^{-1}(x)}$ for $x\in\mathbf Z/q$, such that $s_j\big((g_tw)_j\big)=g'_{\beta_t(j)}\big((s_i(w_i))_i\big)_j$ for all $w$ and $j$, and $\det g'_x=\delta_t(x)$, where $\delta_t(x)=\beta$ if $t=t_0$ and $x=1$, and $\delta_t(x)=1$ otherwise.
For $t\ne t_0$ take $s=1$ and $\beta_t=\pi_t$: on each block $\pi_t^{-1}(x)$ the divisors $\eta(j)\mathfrak d$ add up to $(\pi_t)_*\eta(x)\,\mathfrak d=0$, and Lemma~\ref{sha-steinitz} with determinant $1$ gives $g'_x$ carrying $\mathcal O_F^{\pi_t^{-1}(x)}$ to the lattice of these divisors.
For $t=t_0$, on a coordinate $j$ with $\pi_{t_0}(j)=\kappa^i$ for some $0\le i\le2e+3$ take $s_j=\gamma^{-i}$ and $\beta_t(j)=1$; this is consistent because $\kappa^i=\kappa^{i'}$ in $\mathbf Z/q$ implies $\gamma^i=\gamma^{i'}$, and $\gamma^{-i}(\zeta)^{\kappa^i}=\zeta$.
On the other coordinates take $s_j=1$ and $\beta_t(j)=\pi_{t_0}(j)$.
Put $\mathfrak c_j(P)=\eta(j)\,\mathfrak d(s_j^{-1}P)$.
On the block $\beta_t^{-1}(1)$ the $\mathfrak c_j$ add up to $\sum_ip_i\,(\mathfrak d\circ\gamma^i)=\operatorname{div}\beta$, by Lemma~\ref{sha-eta} and Step 2, and on every other block they add up to $0$, because $(\pi_{t_0})_*\eta$ vanishes off the powers of $\kappa$.
Lemma~\ref{sha-steinitz} gives $g'_x$ with determinant $\delta_t(x)$ carrying $\mathcal O_F$ on the block to the lattice of the $\mathfrak c_j$, and the blocks assemble to an $L$-linear $g_t$; since $s_j$ carries the lattice of $\eta(j)\mathfrak d$ onto that of $\mathfrak c_j$, $g_t\Lambda(0)=\Lambda(\eta\mathfrak d)$.
A map in block form commutes with $\rho\circ\Delta_t$: after the twists $s_j$, $\rho\Delta_t(\tau)$ multiplies the coordinate $j$ by $\zeta^{\tau\beta_t(j)}$, a scalar on each block, which commutes with the $F$-linear $g'_x$.

\emph{Step 4: the subgroup $H_b$.}
Since every $\eta(j)\mathfrak d$ is supported on primes containing $b$, each $g_t$ lies in $G_b$.
The diagonal maps $\rho(u_1,u_2)$ preserve every $\Lambda(\mathfrak a)$, so they lie in $H_b$, and $g_{t_0}^{-1}\rho(u_1,u_2)g_{t_0}$ preserves $\Lambda(0)$; and $g_t^{-1}g_{t_0}$ preserves $\Lambda(0)$ because both $g_t$ and $g_{t_0}$ carry $\Lambda(0)$ to $\Lambda(\eta\mathfrak d)$.

\emph{Step 5: a cycle in $H_b$.}
Put $\rho'=c_{g_{t_0}}\circ\rho:\mathbf Z/q\times\mathbf Z/q\to H_b$ and $k_t=g_{t_0}^{-1}g_t\in H_b$, and let $w$ be a chain with $dw=\sum_tn_t(\Delta_t)_*z_{q,e}$.
Then
\[
 y_0=\rho'_*w-\rho_*w+\sum_tn_t\,h_{k_t}(\rho'\Delta_t)_*z_{q,e}\in C_{2e+3}(H_b;\mathbf Z/q)
\]
is a cycle: by Lemma~\ref{sha-bar-identities}(b), $dh_{k_t}(\rho'\Delta_t)_*z_{q,e}=(\rho\Delta_t)_*z_{q,e}-(\rho'\Delta_t)_*z_{q,e}$, because $c_{k_t}\circ\rho'\circ\Delta_t=\rho\circ\Delta_t$ as $g_t$ commutes with $\rho\Delta_t$.
By Lemma~\ref{sha-bar-identities}(b),(c), applied in $G_b$ to $h_{g_{t_0}}\rho_*w$ and to $h_{g_{t_0}k_t}$, the image of $y_0$ in $C_{2e+3}(G_b;\mathbf Z/q)$ differs from $\sum_tn_t\,h_{g_t}(\rho\Delta_t)_*z_{q,e}$ by a boundary.

\emph{Step 6: passage to $GL_N(\mathcal O_L)$.}
The $\mathcal O_L$-module $\mathcal O_F^J$ is finitely generated and projective, and $H_b$ acts on it $\mathcal O_L$-linearly.
Choose a module $P'$ with $\mathcal O_F^J\oplus P'\cong\mathcal O_L^N$, let $\lambda:H_b\to GL_N(\mathcal O_L)$ act on the first summand, and put $y=\lambda_*(n_{t_0}^{-1}y_0)$, a cycle.

\emph{Step 7: evaluation.}
Let $r$, $k$, $v$, $\theta$ be as in the statement; then $r(b_0)\ne0$ and Step 1 applies.
Reduction defines a homomorphism $\mu:G_b\to GL_k(k^{\Gamma\times J})$ such that, if the coordinates of $x\in F^J$ are integral at the primes not containing $b$, the vector $\big(r(\sigma(x_j))\big)_{(\sigma,j)}$ of residues of $gx$ is $\mu(g)$ applied to that of $x$.
On $\mathcal O_F^J$ this reduction is surjective with kernel $\mathfrak p\mathcal O_F^J$, so $r_N\circ\lambda$ is isomorphic to $\mu|_{H_b}\oplus1$, and additivity of $\theta$ with Lemma~\ref{sha-bar-identities}(a) gives $\theta_{k^N}((r_N)_*y)=n_{t_0}^{-1}\theta(\mu_*y_0)$.
By Step 5 and the vanishing of $\theta$ on boundaries, $\theta(\mu_*y_0)=\sum_tn_t\,\theta\big(h_{\mu(g_t)}(\mu\rho\Delta_t)_*z_{q,e}\big)$.

For each $t$, $k^{\Gamma\times J}$ is the direct sum of the blocks $B_{\sigma,x}$ spanned by the coordinates $(\sigma s_j,j)$ with $\beta_t(j)=x$.
On $B_{\sigma,x}$, $\mu\rho\Delta_t(\tau)$ is the scalar $r(\zeta)^{\kappa(\sigma)x\tau}$, and $\mu(g_t)$ acts by the reduction of $\sigma(g'_x)$, whose determinant is the residue of $\sigma(\delta_t(x))$.
Additivity of $\theta$ and its value $v$ give
\[
 \theta\big(h_{\mu(g_t)}(\mu\rho\Delta_t)_*z_{q,e}\big)=\sum_{\sigma\in\Gamma}\sum_{x\in\mathbf Z/q}\big(\kappa(\sigma)x\big)^{e+1}f\big(\sigma(\delta_t(x))\big),
\]
where $f(\xi)=v(\text{residue of }\xi)$ for $\xi\in F^\times$ that are units at all primes not containing $b$; $f$ is additive on this group, which is stable under $\Gamma$, and $f(1)=0$.
For $t\ne t_0$ all $\delta_t(x)$ are $1$ and the sum is $0$.
For $t=t_0$ only $x=1$ contributes, and since $\kappa(\gamma^i)\equiv\kappa^i$ and $\gamma$ has order $\ell^{\nu-c}$, which is the order of $\kappa$ modulo $q$,
\[
 \sum_{\sigma\in\Gamma}\kappa(\sigma)^{e+1}f(\sigma\beta)=\sum_{n}a^{n}\sum_ie_i\,f(\gamma^{n-i}\alpha)=\Big(\sum_ie_ia^i\Big)\sum_na^nf(\gamma^n\alpha)=p_e\,f\big(N_a(\alpha)\big),
\]
with $n$ running over $\mathbf Z/\ell^{\nu-c}$, because $\sum_ie_ia^i=\prod_{j\ne e+1}(\kappa^{e+1}-\kappa^j)=p_e$.
Here $\alpha$ is a unit at the primes not containing $b$, since its divisor is supported on primes containing $b$, so $N_a(\alpha)$ has a residue $u\in k^\times$ as in the statement and $f(N_a(\alpha))=v(u)$.
Altogether $\theta_{k^N}((r_N)_*y)=n_{t_0}^{-1}n_{t_0}\,p_e\,v(u)=p_e\,v(u)$.
\end{proof}
\leanref{Development/AutomorphicArithmetic/AutomorphicVarieties/ModularityArgument/ClassCounting/LemmaKOfSteps.lean}{UpperTheory.V.ZK.SClass}
\leanref{Development/AutomorphicArithmetic/AutomorphicVarieties/ModularityArgument/ClassCounting/SClassOf.lean}{UpperTheory.V.ZK.sClassOf}
\leanref{Development/AutomorphicArithmetic/AutomorphicVarieties/ModularityArgument/ClassCounting/LatticeBlockClasses/CJoin.lean}{UpperTheory.V.ZK.Cl.cJoin}
\leanref{Development/AutomorphicArithmetic/AutomorphicVarieties/ModularityArgument/ClassCounting/LatticeBlockClasses/CChain.lean}{UpperTheory.V.ZK.Cl.cChain}
\leanref{Development/AutomorphicArithmetic/AutomorphicVarieties/ModularityArgument/ClassCounting/LatticeBlockClasses/GtEqZeta.lean}{UpperTheory.V.ZK.Cl.cGtEq}
\leanref{Development/AutomorphicArithmetic/AutomorphicVarieties/ModularityArgument/ClassCounting/LatticeBlockClasses/CGtNe.lean}{UpperTheory.V.ZK.Cl.cGtNe}
\leanref{Development/AutomorphicArithmetic/AutomorphicVarieties/ModularityArgument/ClassCounting/LatticeBlockClasses/BlkEquiv.lean}{UpperTheory.V.ZK.Cl.cVal}
\leanref{Development/AutomorphicArithmetic/AutomorphicVarieties/ModularityArgument/ClassCounting/LatticeBlockClasses/CBlockVal.lean}{UpperTheory.V.ZK.Cl.cBlockVal}
\leanref{Development/AutomorphicArithmetic/AutomorphicVarieties/ModularityArgument/ClassCounting/LatticeBlockClasses/CBetaVal.lean}{UpperTheory.V.ZK.Cl.cBetaVal}
\leanref{Development/AutomorphicArithmetic/AutomorphicVarieties/ModularityArgument/ClassCounting/LatticeBlockClasses/RedApplyMap.lean}{UpperTheory.V.ZK.Cl.cFree}

\begin{theorem}[\cite{Bor74}]\label{sha-borel-rank}
Let $L$ be a totally complex number field with $r_2(L)$ complex places, let $k\ge0$ be even with $k+1<r_2(L)$, and let $l\ge1$ and $n\ge2k+6$ be integers.
For $m\ge1$ let $i_m:GL_m(\mathcal O_L)\to GL_{lm}(\mathcal O_L)$ and $\psi_l:GL_m(\mathcal O_L)\to GL_{lm}(\mathcal O_L)$ be $g\mapsto\operatorname{diag}(g,1)$ and $g\mapsto\operatorname{diag}(g,\dots,g)$, and for $x\in H_{k+1}(GL_n(\mathcal O_L);\mathbf Z)$ define $z_t(x)\in H_{k+1}(GL_{l^tn}(\mathcal O_L);\mathbf Z)$ by $z_0(x)=x$ and $z_{t+1}(x)=(\psi_l)_*z_t(x)-l^{t+2}(i_{l^tn})_*z_t(x)$.
Then the subgroup of $H_{k+1}(GL_{l^kn}(\mathcal O_L);\mathbf Z)$ generated by the classes $z_k(x)$ has rank at most $r_2(L)$.
\end{theorem}
\leanref{Development/AutomorphicArithmetic/AutomorphicVarieties/ModularityArgument/ClassCounting/LemmaKOfSteps.lean}{UpperTheory.V.ZK.BOdd}
\leanref{Development/AutomorphicArithmetic/AutomorphicVarieties/ModularityArgument/ClassCounting/HasVal.lean}{UpperTheory.V.ZK.zSeq}
\leanref{Development/AutomorphicArithmetic/AutomorphicVarieties/LiteratureResults/KTheoryStability/ReductionBoundsTorsion/TorsionIntegers/Bridge.lean}{PositivityTheory.KB.bOdd_zk}

Borel computed the stable real cohomology of $GL_n(\mathcal O_L)$ \cite{Bor74}: for $L$ totally complex, $\varinjlim_nH^*(GL_n(\mathcal O_L);\mathbf R)$ is an exterior algebra on primitive generators, of which there are $r_2(L)$ in each odd degree at least $3$, $r_2(L)-1$ in degree $1$, and none in even degrees.
Theorem~\ref{sha-borel-rank} follows from this computation.
By homological stability \cite{vdK80}, $H_{k+1}(GL_m(\mathcal O_L);\mathbf Q)$ is the stable group for $m\ge2k+6$, so the rank in question is the dimension of the span of the classes $z_k(x)$ in the stable rational homology, in which the maps $(i_m)_*$ are the identity.
Block sum makes the stable rational homology a connected, commutative and cocommutative graded Hopf algebra, which is therefore the free graded-commutative algebra on its primitive elements \cite{MM65}.
The map $(\psi_l)_*$ is the $l$-th convolution power of the identity, so it multiplies a product of $j$ primitive classes by $l^j$.
Hence $z_{t+1}(x)$ is the image of $z_t(x)$ under $(\psi_l)_*-l^{t+2}$, which kills the products of $t+2$ primitive classes.
A class of degree $k+1$ is a sum of products of between $1$ and $k+1$ primitive classes, so $z_k(x)$ is primitive, and the primitive part of degree $k+1$ has dimension at most $r_2(L)$.

\begin{lemma}\label{sha-few-values}
Let $L$ be a totally complex number field, let $k\ge0$ be even with $k+1<r_2(L)$, let $\ell$ be a prime and $\nu\ge0$, let $B$ be a $\mathbf Z/\ell^\nu$-module, and let $\Phi$ be a primitive family of degree $k+1$ over $\mathcal O_L$, with coefficients in $A=\mathbf Z/\ell^\nu$ and values in $B$ (Definition~\ref{sha-additive-family}).
Then every finite set of elements of $B$ of the form $\ell^k\Phi_n(y)$, for $n\ge0$ and cycles $y\in C_{k+1}(GL_n(\mathcal O_L);\mathbf Z/\ell^\nu)$, has at most $(\ell^\nu)^{r_2(L)}\,t(L,2k+6,k+1)\,t(L,2k+6,k)$ elements, with $t$ as in Definition~\ref{sha-kummer-layers}.
\end{lemma}
\begin{proof}
Write $q=\ell^\nu$.

\emph{Step 1: homological stability.} For $m\ge0$ and $n'\ge n\ge2m+4$, $H_m(GL_n(\mathcal O_L);\mathbf Z)\to H_m(GL_{n'}(\mathcal O_L);\mathbf Z)$ is an isomorphism, by van der Kallen's stability theorem for rings of stable rank $2$ \cite{vdK80}; so the torsion of $H_m(GL_{n'}(\mathcal O_L);\mathbf Z)$ is finite of order $t(L,n,m)$.

\emph{Step 2: descent to homology.} For an integral cycle $z$ let $\bar z$ be its reduction modulo $q$.
The maps $z\mapsto\Phi_n(\bar z)$ vanish on integral boundaries and define homomorphisms $\varphi_n:H_{k+1}(GL_n(\mathcal O_L);\mathbf Z)\to B$ with $\varphi_{n'}\circ(i_{n,n'})_*=\varphi_n$ and $\varphi_{dn}\circ(\psi_d)_*=d\varphi_n$.
By induction on $t$, $\varphi_{l^tn}(z_t(x))=\ell^t(1-\ell c_t)\varphi_n(x)$ for an integer $c_t$, with $l=\ell$ in Theorem~\ref{sha-borel-rank}; since $1-\ell c_t$ is invertible modulo $q$, there is an integer $u$ with $\ell^k\varphi_n(x)=\varphi_{\ell^kn}(u\,z_k(x))$ for all $x$.

\emph{Step 3: counting.} Let $S$ be a finite set of values $\ell^k\Phi_{n_s}(y_s)$.
Choose $n\ge2k+6$ with $n\ge n_s$ for all $s$, and replace $y_s$ by its image in $GL_n$, which does not change the value.
Let $Y$ be the group of cycles of $C_{k+1}(GL_n(\mathcal O_L);\mathbf Z/q)$ and $\beta:Y\to T$ the Bockstein map to the torsion $T$ of $H_k(GL_n(\mathcal O_L);\mathbf Z)$; a cycle with $\beta(y)=0$ lifts to an integral cycle, and $|T|=t(L,2k+6,k)$ by Step 1.
For $y$ with $\beta(y)=0$, lifted to $z$, Step 2 gives $\ell^k\Phi_n(y)=\varphi_{\ell^kn}(u\,z_k([z]))$, a value of $\varphi_{\ell^kn}$ on the subgroup $J$ generated by the $z_k(x)$.
By Theorem~\ref{sha-borel-rank} $J$ has rank at most $r_2(L)$, and the torsion of $H_{k+1}(GL_{\ell^kn}(\mathcal O_L);\mathbf Z)$ has order $t(L,2k+6,k+1)$ by Step 1.
A homomorphism to a group killed by $q$ takes at most $q^{r}|\operatorname{tors}|$ values on a subgroup of rank $r$, so the values on $\ker\beta$ number at most $q^{r_2(L)}t(L,2k+6,k+1)$.
The values on each fibre of $\beta$ are translates of a subset of the values on $\ker\beta$, so there are at most $|T|$ times as many values in all.
\end{proof}
\leanref{Development/AutomorphicArithmetic/AutomorphicVarieties/ModularityArgument/ClassCounting/LemmaKOfSteps.lean}{UpperTheory.V.ZK.SUpper}
\leanref{Development/AutomorphicArithmetic/AutomorphicVarieties/ModularityArgument/ClassCounting/SUpperOf.lean}{UpperTheory.V.ZK.sUpperOf}
\leanref{Development/AutomorphicArithmetic/AutomorphicVarieties/ModularityArgument/ClassCounting/TorsionRankCounts/TorsStab.lean}{UpperTheory.V.ZK.Up.torsStab}
\leanref{Development/AutomorphicArithmetic/AutomorphicVarieties/ModularityArgument/ClassCounting/TorsionRankCounts/Descend.lean}{UpperTheory.V.ZK.Up.descend}
\leanref{Development/AutomorphicArithmetic/AutomorphicVarieties/ModularityArgument/ClassCounting/TorsionRankCounts/PhiZSeq.lean}{UpperTheory.V.ZK.Up.zVal}
\leanref{Development/AutomorphicArithmetic/AutomorphicVarieties/ModularityArgument/ClassCounting/TorsionRankCounts/BockAbs.lean}{UpperTheory.V.ZK.Up.bock}

\begin{lemma}\label{sha-few-norms}
Let $L$ be a totally complex number field, let $\ell$ be a prime, let $1\le c\le\nu$ and $e\ge0$ be integers with $2e+3<\ell^c$ and $2e+3<r_2(L)$, and let $F/L$ with $\gamma$ and $\zeta$ be a Kummer layer of type $(\ell,c,\nu)$ (Definition~\ref{sha-kummer-layers}).
Then the layer has few twisted norms $\big((1+\ell^c)^{e+1},\ell^{b},\ell^\nu,B\big)$, where $b=2e+2+(2e+3)(c+2e+3)$ and
\[
 B=(\ell^\nu)^{r_2(L)}\,t(L,4e+10,2e+3)\,t(L,4e+10,2e+2).
\]
\end{lemma}
\begin{proof}
Write $q=\ell^\nu$, $a=(1+\ell^c)^{e+1}$ and $s=\ell^b$, and let $S\subset F^\times$ be a finite set as in Definition~\ref{sha-kummer-layers}(d).
For $\alpha\in S$, Proposition~\ref{sha-class-cycle} gives $N_\alpha$, a cycle $y_\alpha$ and $b_\alpha\ne0$.
Let $I$ be the set of primes $P$ of $F$ with $\ell\prod_\alpha b_\alpha\notin P$ and $\mathcal O_L\to\mathcal O_F/P$ surjective.

For $P\in I$ let $k_P=\mathcal O_F/P$ and $r_P:\mathcal O_F\to k_P$ the reduction.
Since $P\nmid\ell$, $r_P(\zeta)$ is a primitive $q$-th root of unity, so $q$ divides $|k_P^\times|$ and there is a homomorphism $v_P:k_P^\times\to\mathbf Z/q$ whose kernel is the group of $q$-th powers.
Lemma~\ref{sha-theta} gives an additive family $\theta_P$ of degree $2e+3$ over $k_P$ with value $v_P$ relative to $r_P(\zeta)$, and Lemma~\ref{sha-primitive} gives the primitive family $\Phi_P=\theta_P\circ(r_P)_*$ over $\mathcal O_L$.
Then $\Phi=(\Phi_P)_{P\in I}$ is a primitive family of degree $2e+3$ with values in $(\mathbf Z/q)^I$.

The map $\alpha\mapsto\ell^{2e+2}\Phi(y_\alpha)$ is injective on $S$.
Suppose that $\ell^{2e+2}\Phi_P(y_\alpha)=\ell^{2e+2}\Phi_P(y_{\alpha'})$ for every $P\in I$.
By Proposition~\ref{sha-class-cycle}, $\ell^{2e+2}p_e\big(v_P(u_\alpha)-v_P(u_{\alpha'})\big)=0$, where $u_\alpha$ is the residue of $N_a(\alpha)$ at $P$.
Each factor $\kappa^{e+1}-\kappa^j$ of $p_e$, with $\kappa=1+\ell^c$, $0\le j\le2e+2$ and $j\ne e+1$, is $\pm\kappa^{\min(j,e+1)}(\kappa^{m}-1)$ with $0<m=|e+1-j|<2e+3<\ell^c$, and $\ell^{c+2e+3}$ does not divide $\kappa^m-1$, because $\kappa^m-1=\ell^c(m+\ell^cx)$ for an integer $x$ and $v_\ell(m)<\min(c,2e+3)$.
So $p_e$ divides $\ell^{(2e+3)(c+2e+3)}$ in $\mathbf Z/q$, and $s\big(v_P(u_\alpha)-v_P(u_{\alpha'})\big)=0$: the residue of $N_a(\alpha)^s/N_a(\alpha')^s$ at $P$ is a $q$-th power.
By Lemma~\ref{sha-detect}, applied with $b=\ell\prod_\alpha b_\alpha$, the quotient $N_a(\alpha)^s/N_a(\alpha')^s$ is a $q$-th power in $F^\times$, and the second property of $S$ gives $\alpha=\alpha'$.

Finally $k=2e+2$ is even and $k+1=2e+3<r_2(L)$, so Lemma~\ref{sha-few-values} bounds the number of values $\ell^{2e+2}\Phi(y_\alpha)$ by $q^{r_2(L)}t(L,2k+6,k+1)t(L,2k+6,k)=B$, since $2k+6=4e+10$.
\end{proof}
\leanref{Development/AutomorphicArithmetic/AutomorphicVarieties/ModularityArgument/ClassCounting/LemmaKOfSteps.lean}{UpperTheory.V.ZK.SFew}
\leanref{Development/EtaleCohomology/CyclotomicDiscriminants/FinSFewOf.lean}{LateralTheory.Fin.SFewOf.intersectionNumber_baseChange}
\leanref{Development/EtaleCohomology/CyclotomicDiscriminants/SFewOfCut.lean}{LateralTheory.Q.SFewOfCut.sFewOf}

\begin{lemma}\label{sha-count-homs}
Let $L$ be a totally complex number field, let $\ell$ be a prime, let $F/L$ with $\gamma$ and $\zeta$ be a Kummer layer of type $(\ell,c,\nu)$, let $E\ge1$, $b\ge0$ and $B\ge0$ be integers, and put $a=(1+\ell^c)^{E}$.
If the layer has few twisted norms $(a,\ell^{b},\ell^{\nu},B)$, then $\operatorname{Hom}_a(F,\ell^\nu)$ is finite and
\[
 \#\operatorname{Hom}_a(F,\ell^\nu)\cdot(\ell^{\nu})^{r_2(L)}\le B\cdot\ell^{(b+v_\ell(E))\operatorname{gen}(L)+c+v_\ell(E)}.
\]
\end{lemma}
\begin{proof}
Write $n=\ell^{\nu-c}$ for the order of $\gamma$, $q=\ell^\nu$, $s=\ell^b$ and $\operatorname{Div}F$ for the group of divisors; let $g=\gamma^{-1}$ act on $F^\times$ and on $\operatorname{Div}F$, and let $\mu(F)$ be the roots of unity of $F$.
Since $c\ge1$, and $c\ge2$ if $\ell=2$, the $\ell$-adic valuations are $v_\ell(a-1)=c+v_\ell(E)$ and $v_\ell(a^n-1)=\nu+v_\ell(E)$; in particular $a\ge2$.

\emph{The groups.} Let $U\subset F^\times$ be the group of $\alpha$ whose divisor is $(\gamma-a)\mathfrak d$ for some divisor $\mathfrak d$, and put
\[
 W=U/\big(U\cap(g-a)F^\times\mu(F)\big),\quad E_Q=\mathcal O_F^\times/(g-a)\mathcal O_F^\times\mu(F),\quad C_Q=\operatorname{Div}F/\big(\operatorname{div}F^\times+(g-a)\operatorname{Div}F\big),
\]
written additively.
$C_Q$ is a quotient of the class group of $F$, hence finite.
Since $L$ is totally complex, every infinite place of $L$ splits completely in $F$, and by Dirichlet's unit theorem $(\mathcal O_F^\times/\mu(F))\otimes\mathbf Q$ is isomorphic, as a module over $\operatorname{Gal}(F/L)$, to $\mathbf Q[\operatorname{Gal}(F/L)]^{r_2(L)}$ modulo one trivial summand.
The determinant of $g-a$ on $\mathbf Q[\operatorname{Gal}(F/L)]$ is $\pm(a^n-1)$ and on the trivial summand $1-a$, so $E_Q$ is finite and $\#E_Q\cdot(a-1)=(a^n-1)^{r_2(L)}$.
Since $a\ge2$ and $g$ has finite order, $g-a$ is injective on $\operatorname{Div}F$, so $\mathfrak d$ is determined by $\alpha$.
The map $\alpha\mapsto$ (class of $\mathfrak d$) gives an exact sequence $0\to E_Q\to W\to\operatorname{Cl}(F)[\gamma-a]\to0$, where $\operatorname{Cl}(F)[\gamma-a]$ is the kernel of $\gamma-a$ on the class group; on the finite group $\operatorname{Cl}(F)$ this kernel has the order of the cokernel, which is $C_Q$.
So $\#W=\#E_Q\cdot\#C_Q$.

\emph{Homomorphisms.} An element of $\operatorname{Hom}_a(F,q)$ is a function on primes that vanishes on principal divisors and satisfies $\varphi(\gamma^{-1}P)=a\varphi(P)$, that is a homomorphism $C_Q\to\mathbf Z/q$.
So $\#\operatorname{Hom}_a(F,q)\le\ell^{v_\ell(\#C_Q)}$.

\emph{Twisted norms.} Write $a^n-1=q\,\ell^{v_\ell(E)}m$ with $\ell\nmid m$.
The twisted norm satisfies $N_a\circ(g-a)=(1-a^n)g$ on $F^\times$, so $f(\alpha)=N_a(\alpha)^s$ defines a homomorphism $f:W\to F^\times/F^{\times q}\mu(F)$, and if $f(w)=0$ then $(\ell^{v_\ell(E)}ms)w=0$, because $F^\times/\mu(F)$ is torsion-free.
Distinct elements of the image of $f$ come from elements of $U$ whose twisted norms satisfy the second condition of Definition~\ref{sha-kummer-layers}(d), so the image of $f$ has at most $B$ elements.

\emph{The $\ell$-rank of $W$.} The $\ell$-torsion of $W$ injects into the set of classes $y\in F^\times/F^{\times\ell}\mu(F)$ with $\gamma y\equiv y$ and $\operatorname{div}y\in\ell\operatorname{Div}F$, since $\ell\mid a-1$.
Such a class has a representative with $\gamma y=y\,w^\ell\rho^{r}$, where $\rho$ generates $\mu(F)$, $0\le r<\ell$ and $N_{F/L}(w)^\ell=N_{F/L}(\rho)^{-r}$; two representatives with the same $r$ and the same $N_{F/L}(w)$ differ by an element of $L^\times$ whose valuations at the primes not above $\ell$ are divisible by $\ell$, by Hilbert's Theorem 90.
The classes of such elements of $L^\times$ modulo $\ell$-th powers number at most $\ell^{2[L:\mathbf Q]+\lfloor\log_2h_L\rfloor}$, from the units, the primes above $\ell$ and the $\ell$-torsion of the class group.
So the $\ell$-torsion of $W$ has at most $\ell^2\cdot\ell^{2[L:\mathbf Q]+\lfloor\log_2h_L\rfloor}=\ell^{\operatorname{gen}(L)}$ elements.

\emph{Conclusion.} The kernel of $f$ has $\ell$-part killed by $\ell^{b+v_\ell(E)}$ and of $\ell$-rank at most $\operatorname{gen}(L)$, so $\ell^{v_\ell(\#W)}\le B\cdot\ell^{(b+v_\ell(E))\operatorname{gen}(L)}$.
From the valuations above, $v_\ell(\#E_Q)+c+v_\ell(E)=r_2(L)\,(\nu+v_\ell(E))\ge r_2(L)\nu$.
Hence
\[
 \#\operatorname{Hom}_a(F,q)\,q^{r_2(L)}\le\ell^{v_\ell(\#C_Q)+v_\ell(\#E_Q)+c+v_\ell(E)}=\ell^{v_\ell(\#W)}\ell^{c+v_\ell(E)}\le B\,\ell^{(b+v_\ell(E))\operatorname{gen}(L)+c+v_\ell(E)}.
\]
\end{proof}
\leanref{Development/AutomorphicArithmetic/AutomorphicVarieties/ModularityArgument/ClassCounting/LemmaKOfSteps.lean}{UpperTheory.V.ZK.SCount}
\leanref{Development/AutomorphicArithmetic/AutomorphicVarieties/ModularityArgument/ClassCounting/SelmerOrderCounts.lean}{UpperTheory.V.ZK.sCount}
\leanref{Development/AutomorphicArithmetic/AutomorphicVarieties/ModularityArgument/ClassCounting/SelmerOrderCounts/CardEQ.lean}{UpperTheory.V.ZK.Ct.cardEQ}
\leanref{Development/AutomorphicArithmetic/AutomorphicVarieties/ModularityArgument/ClassCounting/SelmerOrderCounts/KerFDelNm.lean}{UpperTheory.V.ZK.Ct.kerF}
\leanref{Development/AutomorphicArithmetic/AutomorphicVarieties/ModularityArgument/ClassCounting/SelmerOrderCounts/SelInvLe.lean}{UpperTheory.V.ZK.Ct.selInvLe}
\leanref{Development/AutomorphicArithmetic/AutomorphicVarieties/ModularityArgument/ClassCounting/SelmerOrderCounts/SelLCard.lean}{UpperTheory.V.ZK.Ct.selLCard}

\begin{lemma}\label{sha-kummer-descent}
Let $L$ be a totally complex number field, let $\ell$ be a prime, and suppose that the roots of unity of $\ell$-power order in $L$ are those of order dividing $\ell^c$, with $c\ge1$, and $c\ge2$ if $\ell=2$.
Let $k\ge1$ and $B\ge0$, and suppose that for every $\nu\ge c$ and every Kummer layer $F$ of type $(\ell,c,\nu)$ over $L$, the set $\operatorname{Hom}_{(1+\ell^c)^k}(F,\ell^\nu)$ is finite with at most $B$ elements.
Then $\text{Ш}^1_\ell(L,k)$ is finite and $\ell^{\operatorname{len}\text{Ш}^1_\ell(L,k)}\le B$.
\end{lemma}
\begin{proof}
Let $U_\infty=\operatorname{Gal}(\bar L/L(\mu_{\ell^\infty}))$.
By the argument of Lemma~\ref{sha-stickelberger-annihilator}, Step 1, restriction to $U_\infty$ is injective on $H^1(G_{L,S_\ell},\mathbf Q_\ell/\mathbf Z_\ell(-k))$, and it maps $\text{Ш}^1_\ell(L,k)$ into the group of continuous homomorphisms $\varphi:U_\infty\to\mathbf Q_\ell/\mathbf Z_\ell$ that are unramified everywhere, trivial on the decomposition groups above $\ell$, and satisfy $\varphi(\sigma u\sigma^{-1})=\chi_{\mathrm{cyc}}(\sigma)^{-k}\varphi(u)$.

For $\nu\ge c$ let $F_\nu=L(\mu_{\ell^\nu})$.
Since $\chi_{\mathrm{cyc}}(\operatorname{Gal}(\bar L/L))=1+\ell^c\mathbf Z_\ell$, there are $\sigma\in\operatorname{Gal}(\bar L/L)$ with $\chi_{\mathrm{cyc}}(\sigma)\equiv1+\ell^c\pmod{\ell^\nu}$, and its restriction $\gamma$ generates $\operatorname{Gal}(F_\nu/L)$, of order $\ell^{\nu-c}$; with a primitive $\ell^\nu$-th root of unity $\zeta$, $F_\nu$ is a Kummer layer of type $(\ell,c,\nu)$.
A homomorphism $\varphi$ as above factors, for $\nu$ large, through a homomorphism on $\operatorname{Gal}(\bar L/F_\nu)$ that is unramified everywhere, trivial above $\ell$ and killed by $\ell^\nu$; by class field theory this is a homomorphism $\operatorname{Cl}(F_\nu)\to\mathbf Z/\ell^\nu$, and the equivariance says $\varphi(\gamma^{-1}P)=(1+\ell^c)^k\varphi(P)$ for its values on primes, since $\operatorname{Frob}_{\gamma^{-1}P}=\sigma^{-1}\operatorname{Frob}_P\sigma$.
This gives an injection, for each finite set of such $\varphi$ and $\nu$ large, into $\operatorname{Hom}_{(1+\ell^c)^k}(F_\nu,\ell^\nu)$.
So every finite subset of the image of $\text{Ш}^1_\ell(L,k)$ has at most $B$ elements; the image is finite of order at most $B$, and its order is $\ell^{\operatorname{len}\text{Ш}^1_\ell(L,k)}$.
\end{proof}
\leanref{Development/AutomorphicArithmetic/AutomorphicVarieties/ModularityArgument/ClassCounting/LemmaKOfSteps.lean}{UpperTheory.V.ZK.SSha}
\leanref{Development/AutomorphicArithmetic/AutomorphicVarieties/ModularityArgument/ClassCounting/RingEquivalences/OrderOfOneAddPrimePowZmodPrimePow.lean}{PositivityTheory.QuaternionShaBounds.sSha}
\leanref{Development/AutomorphicArithmetic/AutomorphicVarieties/ModularityArgument/ClassCounting/RingEquivalences/OrderOfOneAddPrimePowZmodPrimePow.lean}{PositivityTheory.QuaternionShaBounds.LayEx_holds}
\leanref{Development/AutomorphicArithmetic/AutomorphicVarieties/ModularityArgument/ClassCounting/RingEquivalences/OrderOfOneAddPrimePowZmodPrimePow.lean}{PositivityTheory.QuaternionShaBounds.ClInj_holds}

\begin{lemma}\label{sha-enlarging}
Let $L$ be a number field and $\ell$ a prime such that $L$ contains a primitive $\ell$-th root of unity, and a primitive fourth root of unity if $\ell=2$; let $n=[L:\mathbf Q]$ and $k\ge1$.
Then there are a totally complex number field $L'\supset L$ and integers $c,d\ge1$ such that: the roots of unity of $\ell$-power order in $L'$ are those of order dividing $\ell^{c}$, with $c\ge2$ if $\ell=2$; $2k+1<\ell^{c}\le2[L':\mathbf Q]$; $2k+1<r_2(L')$; $[L':\mathbf Q]=dn$ with $d\le(n+1)(4k+4)$; $|d_{L'}|\le|d_L|^{d}(\ell^{c})^{[L':\mathbf Q]}$; and $\operatorname{len}\text{Ш}^1_\ell(L,k)\le\operatorname{len}\text{Ш}^1_\ell(L',k)$.
\end{lemma}
\begin{proof}
Let $\ell^{c_0}$ be the number of $\ell$-power roots of unity in $L$, let $m$ be the least integer with $\ell^m\ge4k+4$, so that $\ell^m\le\ell(4k+4)$, and let $L'=L(\mu_{\ell^{c_0+m}})$, a layer of degree $\ell^m$ over $L$ (Definition~\ref{sha-tower-data}).
Then $L'$ is totally complex, $d=\ell^m$, and $\ell\le n+1$ because $\mu_\ell\subset L$, so $d\le(n+1)(4k+4)$.
The number $\ell^c$ of $\ell$-power roots of unity in $L'$ is at least $\ell^{c_0+m}\ge\ell^m\ge4k+4>2k+1$, and at most $2[L':\mathbf Q]$ because $\varphi(\ell^c)\le[L':\mathbf Q]$.
Also $r_2(L')=[L':\mathbf Q]/2\ge\ell^m/2>2k+1$.
For the discriminant, $L'=L(\zeta)$ with $\zeta^{N}=1$ for $N=\ell^{c_0+m}$, so $N$ lies in the different of $L'/L$, whose norm divides $N^{[L':\mathbf Q]}$; this gives $|d_{L'}|\le|d_L|^{d}N^{[L':\mathbf Q]}\le|d_L|^d(\ell^c)^{[L':\mathbf Q]}$.
Finally, by Lemma~\ref{sha-iwasawa-descent} the lengths are those of the cokernels of $g-1$ and $g^{\ell^m}-1$ on $X'_\infty(L)$, for $g=u^k\gamma$, and the image of $g^{\ell^m}-1$ is contained in that of $g-1$.
\end{proof}
\leanref{Development/AutomorphicArithmetic/AutomorphicVarieties/ModularityArgument/ClassCounting/LemmaKOfSteps.lean}{UpperTheory.V.ZK.SLayer}
\leanref{Development/EtaleCohomology/CyclotomicDiscriminants/FinSLayer.lean}{LateralTheory.Fin.SLayer.intersectionNumber_baseChange}
\leanref{Development/EtaleCohomology/CyclotomicDiscriminants/SLayerCut.lean}{LateralTheory.Q.SLayerCut.goal_of}

\begin{lemma}\label{sha-convex-nerve}
Let $F$ be a real normed vector space, let $X\subset F$ be a nonempty convex subset, and let $(A_i)_{i\in I}$ be a family of closed convex subsets of $F$ that cover $X$, such that the convex hull of every finite subset of $X$ meets only finitely many of the $A_i$.
For $p\ge0$ let $N_p(X)$ be the set of tuples $(i_0,\dots,i_p)\in I^{p+1}$ with $X\cap A_{i_0}\cap\dots\cap A_{i_p}\ne\emptyset$.
Let $\partial:\mathbf Z[N_p(X)]\to\mathbf Z[N_{p-1}(X)]$ be $\partial(i_0,\dots,i_p)=\sum_{j=0}^{p}(-1)^j(i_0,\dots,\widehat{i_j},\dots,i_p)$ for $p\ge1$, and $\partial(i_0)=1\in\mathbf Z$ for $p=0$.
Then the \emph{augmented nerve complex}
\[
 \cdots\to\mathbf Z[N_2(X)]\to\mathbf Z[N_1(X)]\to\mathbf Z[N_0(X)]\to\mathbf Z\to0
\]
is exact.
\end{lemma}
\begin{proof}
Write $N(X)$ for the family of the sets $N_p(X)$, and call such a family of sets of tuples \emph{acyclic} if its augmented complex is exact.

\emph{Step 1: cones.} Let $(N_p)_{p\ge0}$ be sets of tuples closed under deleting entries, and suppose there is $a\in I$ with $(a)\in N_0$ and $(a,i_0,\dots,i_p)\in N_{p+1}$ for every $(i_0,\dots,i_p)\in N_p$.
Put $\kappa(1)=(a)$ in degree $-1$ and $\kappa(i_0,\dots,i_p)=(a,i_0,\dots,i_p)$.
Then $\partial\kappa=1$ in degree $-1$ and $\partial\kappa+\kappa\partial=1$ in degrees $p\ge0$, so the family is acyclic.

\emph{Step 2: $X$ compact and meeting only finitely many $A_i$.}
Then the set $\Sigma(X)$ of finite sets $\sigma\subset I$ with $X\cap\bigcap_{j\in\sigma}A_j\ne\emptyset$ is finite.
We argue by induction on $\#\Sigma(X)$, for all compact convex nonempty $X$ covered by the $A_i$ at once.
Suppose first that some $x\in X$ lies in every $A_i$ that meets $X$, and choose $a$ with $x\in A_a$.
Then $x$ shows that $(a,i_0,\dots,i_p)\in N_{p+1}(X)$ whenever $(i_0,\dots,i_p)\in N_p(X)$, and Step 1 applies.

Otherwise let $\sigma\in\Sigma(X)$ be maximal for inclusion.
Some $A_i$ that meets $X$ does not contain all of $X\cap\bigcap_{j\in\sigma}A_j$, so $i\notin\sigma$, and by maximality $X\cap A_i$ and $X\cap\bigcap_{j\in\sigma}A_j$ are disjoint.
These are nonempty compact convex sets, so there are a continuous linear form $f$ on $F$ and $u\in\mathbf R$ with $f<u$ on the first and $f>u$ on the second.
Put $X_{\le}=X\cap\{f\le u\}$, $X_{\ge}=X\cap\{f\ge u\}$ and $X_{=}=X\cap\{f=u\}$.
They are compact and convex, and nonempty because $X$ is connected and $f$ takes values below and above $u$ on it.
The sets $\Sigma(X_\le)$ and $\Sigma(X_=)$ do not contain $\sigma$, and $\Sigma(X_\ge)$ does not contain $\{i\}$, so all three are smaller than $\Sigma(X)$, and $N(X_\le)$, $N(X_\ge)$ and $N(X_=)$ are acyclic by induction.
A tuple of $N_p(X)$ with a witness $x$ lies in $N_p(X_\le)$ if $f(x)\le u$ and in $N_p(X_\ge)$ otherwise.
If it lies in both, with witnesses $x$ and $x'$, the segment from $x$ to $x'$ lies in $X$ and in the sets $A_{i_j}$ and meets $\{f=u\}$.
So $N_p(X)=N_p(X_\le)\cup N_p(X_\ge)$ and $N_p(X_=)=N_p(X_\le)\cap N_p(X_\ge)$.
The augmented complexes form an exact sequence $0\to\mathbf Z[N(X_=)]\to\mathbf Z[N(X_\le)]\oplus\mathbf Z[N(X_\ge)]\to\mathbf Z[N(X)]\to0$, and its long exact sequence shows that $N(X)$ is acyclic.

\emph{Step 3: the general case.}
An element of $\mathbf Z$ in degree $-1$, or a cycle of the augmented complex of $X$, involves finitely many tuples.
Choose a point of $X$ witnessing each of them, and let $K$ be the convex hull of these points and one further point of $X$.
Then $K\subset X$ is compact, convex, nonempty and covered by the $A_i$, and by hypothesis it meets only finitely many of them.
The tuples involved lie in $N(K)$, and $N_p(K)\subset N_p(X)$ for every $p$.
By Step 2 the element is a boundary in the complex of $K$, hence in that of $X$.
\end{proof}
\leanref{Development/AutomorphicArithmetic/AutomorphicVarieties/ConvexNerveComplexes/NerveConvex.lean}{UpperTheory.V.ConeExtensionAcyclic.nerve_convex}
\leanref{Development/AutomorphicArithmetic/AutomorphicVarieties/ConvexNerveComplexes/NerveCompact.lean}{UpperTheory.V.ConeExtensionAcyclic.nerve_compact}
\leanref{Development/AutomorphicArithmetic/AutomorphicVarieties/ConvexNerveComplexes/AcycCone.lean}{UpperTheory.V.ConeExtensionAcyclic.acyc_cone}
\leanref{Development/AutomorphicArithmetic/AutomorphicVarieties/ConvexNerveComplexes/NerveSep.lean}{UpperTheory.V.ConeExtensionAcyclic.nerve_sep}
\leanref{Development/AutomorphicArithmetic/AutomorphicVarieties/ConvexNerveComplexes/NerveCut.lean}{UpperTheory.V.ConeExtensionAcyclic.nerve_cut}
\leanref{Development/AutomorphicArithmetic/AutomorphicVarieties/ConvexNerveComplexes/AcycUnion.lean}{UpperTheory.V.ConeExtensionAcyclic.acyc_union}
\leanref{Development/AutomorphicArithmetic/AutomorphicVarieties/ConvexNerveComplexes/AcycDirected.lean}{UpperTheory.V.ConeExtensionAcyclic.acyc_directed}

\begin{lemma}\label{sha-polyhedron-nerve}
Let $E$ be a finite-dimensional real vector space with dual $E^\vee$, let $T\subset E$ be a closed subset with $cT\subset T$ for every real $c>0$, and let $(q_y)_{y\in Y}$ be a family of linear forms on $E$.
Put
\[
 P=\{h\in E:\ q_y(h)\ge1\text{ for all }y\in Y\},\qquad T^*=\{x\in E^\vee:\ x(h)>0\text{ for all }h\in T\smallsetminus\{0\}\}.
\]
Suppose that $P$ is nonempty and contained in $T$, that $T^*$ is nonempty, and that $P$ is \emph{locally finite}: every $h\in P$ has a neighbourhood $U$ and a finite set $S\subset Y$ such that $q_y>1$ on $U$ for every $y\notin S$.
A \emph{vertex} of $P$ is an $h\in P$ such that $0$ is the only $k\in E$ with $q_y(k)=0$ for all $y$ with $q_y(h)=1$, and the \emph{normal cone} of $h$ is $C(h)=\{x\in E^\vee:\ x(h)\le x(h')\text{ for all }h'\in P\}$.
Then $X=T^*$, with the family of normal cones $C(v)$ indexed by the vertices $v$ of $P$, satisfies the hypotheses of Lemma~\ref{sha-convex-nerve}.
\end{lemma}
\begin{proof}
The set $T^*$ is convex, and each $C(v)$ is closed and convex, as an intersection of closed half-spaces.

\emph{Compact sublevel sets.} Let $x\in T^*$ and fix a norm on $E$.
The elements of $T$ of norm $1$ form a compact set on which $x$ is positive, so there is $c>0$ with $x(h)\ge c\|h\|$ for $h\in T$.
Since $P\subset T$ is closed, every set $\{h\in P:\ x(h)\le b\}$ is compact.

\emph{Vertices are isolated.} Let $h\in P$, with $U$ and $S$ as in the definition of local finiteness.
For a vertex $v\in U$ the set $S_v=\{y:\ q_y(v)=1\}$ is contained in $S$, and $v$ is the only solution $k$ of the equations $q_y(k)=1$ for $y\in S_v$, since two solutions differ by an element on which every $q_y$ with $y\in S_v$ vanishes.
So $U$ contains at most $2^{\#S}$ vertices, and a compact subset of $P$ contains only finitely many vertices.

\emph{Covering.} Let $x\in T^*$.
By the first step $x$ attains its minimum $\mu$ on $P$, and the set $Z$ of points where it does is compact, so it is covered by finitely many of the neighbourhoods $U$ above.
Hence the sets $S_h=\{y:\ q_y(h)=1\}$ for $h\in Z$ have bounded size; choose $h\in Z$ with $\#S_h$ largest.
Suppose that $h$ is not a vertex, and let $k\ne0$ with $q_y(k)=0$ for $y\in S_h$.
By local finiteness at $h$, $h+tk\in P$ for $|t|$ small, so $x(k)=0$, and $h+tk\in Z$ whenever $h+tk\in P$.
As $Z$ is compact there is a largest $t_0>0$ with $h+tk\in P$ for $0\le t\le t_0$.
At $h'=h+t_0k$ every $q_y$ with $y\in S_h$ still takes the value $1$, and by local finiteness at $h'$ some further $q_y$ takes the value $1$, for otherwise $h'+tk\in P$ for small $t>0$.
So $h'\in Z$ and $\#S_{h'}>\#S_h$, a contradiction.
Hence $h$ is a vertex and $x\in C(h)$.

\emph{Finiteness.} Let $s\subset T^*$ be finite, fix $h_0\in P$, and let $K$ be the union over $x\in s$ of the compact sets $\{h\in P:\ x(h)\le x(h_0)\}$.
If a convex combination $x'=\sum_{x\in s}\lambda_xx$ lies in $C(v)$, then $x'(v)\le x'(h_0)$, so $x(v)\le x(h_0)$ for some $x\in s$, and $v\in K$.
By the second step only finitely many vertices lie in $K$.
\end{proof}
\leanref{Development/AutomorphicArithmetic/AutomorphicVarieties/ConvexNerveComplexes/AcycSimp.lean}{UpperTheory.V.ConeExtensionAcyclic.acyc_simp}
\leanref{Development/AutomorphicArithmetic/AutomorphicVarieties/ConvexNerveComplexes/Sublevel.lean}{UpperTheory.V.ConeExtensionAcyclic.sublevel}
\leanref{Development/AutomorphicArithmetic/AutomorphicVarieties/ConvexNerveComplexes/VertFinite.lean}{UpperTheory.V.ConeExtensionAcyclic.vert_finite}
\leanref{Development/AutomorphicArithmetic/AutomorphicVarieties/ConvexNerveComplexes/ExistsVert.lean}{UpperTheory.V.ConeExtensionAcyclic.exists_vert}
\leanref{Development/AutomorphicArithmetic/AutomorphicVarieties/ConvexNerveComplexes/ConeFinite.lean}{UpperTheory.V.ConeExtensionAcyclic.cone_finite}

\begin{lemma}\label{sha-hermitian-polyhedron}
Let $L$ be a totally complex number field of degree $n$ and discriminant $d_L$, let $W$ be its set of complex places, each given by an embedding $w:L\to\mathbf C$, and let $N\ge1$ and $D=nN$.
Let $E$ be the real vector space of tuples $h=(h_w)_{w\in W}$ of Hermitian $N\times N$ complex matrices, and $T\subset E$ the set of tuples of positive semidefinite matrices.
For nonzero $y\in\mathcal O_L^N$ put $q_y(h)=\sum_{w\in W}w(y)^*\,h_w\,w(y)$, and let $P$, $T^*$, the vertices of $P$ and their normal cones be as in Lemma~\ref{sha-polyhedron-nerve}.
A \emph{minimal vector} of $h\in P$ is a $y$ with $q_y(h)=1$.
For $p\ge0$ let $X_p$ be the set of $(p+1)$-tuples of vertices of $P$ whose normal cones have a common point in $T^*$.
The group $GL_N(\mathcal O_L)$ acts on $E$ by $(g\cdot h)_w=w(g)^{*-1}\,h_w\,w(g)^{-1}$; since $q_{gy}(g\cdot h)=q_y(h)$, it preserves $P$, permutes its vertices, and acts on each $X_p$.
\begin{enumerate}
\item The hypotheses of Lemma~\ref{sha-polyhedron-nerve} hold. Hence the augmented nerve complex $\cdots\to\mathbf Z[X_1]\to\mathbf Z[X_0]\to\mathbf Z\to0$ of Lemma~\ref{sha-convex-nerve} is exact.
\item The stabilizer in $GL_N(\mathcal O_L)$ of every element of $X_p$ is finite, of order at most $3^{DN}$.
\item For every $p\ge0$, the set $X_p$ has at most $(2|d_L|D)^{12(p+1)D^3N}$ orbits under $GL_N(\mathcal O_L)$.
\end{enumerate}
\end{lemma}
\begin{proof}
Let $V=\prod_{w\in W}\mathbf C^N$, in which $y\mapsto(w(y))_w$ embeds $\mathcal O_L^N$ as a lattice, and for $h\in E$ let $Q_h(z)=\sum_wz_w^*h_wz_w$, a real quadratic form on $V$ with $Q_h(y)=q_y(h)$.

(a) The set $T$ is closed and stable under positive scalars, and $T^*$ contains the form $h\mapsto\sum_w\operatorname{tr}h_w$.
The set $P$ is nonempty: for $h_w=1$ and a nonzero coordinate $y_a$ of $y$, $q_y(h)\ge\sum_w|w(y_a)|^2\ge\#W\cdot|N_{L/\mathbf Q}(y_a)|^{1/\#W}\ge1$ by the inequality of arithmetic and geometric means, since $\prod_w|w(y_a)|^2=|N_{L/\mathbf Q}(y_a)|$.

Every $h\in P$ is a tuple of positive definite matrices.
Indeed $Q_h\ge1$ on the nonzero lattice vectors, so $Q_h>0$ on their rational multiples, and $Q_h\ge0$ on $V$ by density.
If $Q_h(z)=0$ with $z\ne0$, then $Q_h$ vanishes on $\mathbf Rz$ and is less than $1$ on a thin cylinder around this line; by Minkowski's theorem a long enough cylinder contains a nonzero lattice vector, a contradiction.
So $Q_h$ is positive definite, and so is each $h_w$.

For local finiteness, let $h\in P$.
Since $Q_h$ is positive definite, there are a neighbourhood $U$ of $h$ and $c>0$ with $Q_k(z)\ge c\sum_w|z_w|^2$ for $k\in U$ and $z\in V$.
The set $S$ of lattice vectors $y$ with $\sum_w|w(y)|^2\le1/c$ is finite, and $q_y>1$ on $U$ for $y\notin S$.

(b) An element fixing a tuple fixes its first vertex $v$, which is positive definite by (a).
The minimal vectors of $v$ are finitely many, by local finiteness, and they span $L^N$: if they all lay in a proper subspace $W'\subset L^N$, the tuple $k$ of orthogonal projections onto the orthogonal complements of the spans of $w(W')$ would be nonzero with $q_y(k)=0$ for every minimal vector $y$, and $v$ would not be a vertex.
An element of the stabilizer permutes these minimal vectors, and is determined by this permutation, so the stabilizer is finite.
A finite group of units of the ring $M_N(\mathcal O_L)$, which is free of rank $DN$ over $\mathbf Z$, injects into $M_N(\mathcal O_L)/3M_N(\mathcal O_L)$, because no unit $g\equiv1$ modulo $3$ has prime order.
Indeed, let $g\ne1$ have prime order $p$, and write $g=1+3^ah$ with $a\ge1$ and $h\notin3M_N(\mathcal O_L)$.
Expanding $(1+3^ah)^p=1$ gives $p3^ah\equiv0$ modulo $3^{a+1}$ if $p\ne3$, and $3^{a+1}h\equiv0$ modulo $3^{a+2}$ if $p=3$; both contradict $h\notin3M_N(\mathcal O_L)$.
So the stabilizer has at most $3^{DN}$ elements.

(c) This is a counting computation; we give its steps and its result.
Put $I_0=|d_L|^ND^D$ and $R_0=N^NI_0^2$.
The vertices of an element of $X_p$ minimize a common $x\in T^*$ on $P$.
Let $h$ be a point in the relative interior of the face of $P$ on which $x$ is minimal; its minimal vectors are minimal vectors of every point of the face, and they span $L^N$.
Otherwise the tuple $k$ of (b) for their span would satisfy $x(k)>0$, while $h-tk\in P$ for small $t>0$ by local finiteness, and $x(h-tk)<x(h)$.
So there is $A\in M_N(\mathcal O_L)$ with $\det A\ne0$ whose columns are minimal vectors of every vertex of the tuple.
For such $A$ and a vertex $v$, the tuple $(w(A)^*v_ww(A))_w$ has diagonal entries summing over $w$ to $1$, so Hadamard's inequality and a lower bound for $\prod_w\det v_w$ from Minkowski's theorem give $|N_{L/\mathbf Q}(\det A)|\le I_0$.
By Hermite normal forms over $\mathcal O_L$, the matrices $A$ with $\det A\ne0$ and $|N_{L/\mathbf Q}(\det A)|=I$ fall into at most $I^{DN}$ classes under left multiplication by $GL_N(\mathcal O_L)$, so those with $|N_{L/\mathbf Q}(\det A)|\le I_0$ fall into at most $I_0^{1+DN}$ classes.
For a fixed $A$, the minimal vectors of the vertices having the columns of $A$ among their minimal vectors lie in a set of at most $(2R_0+1)^D$ vectors, and each such vertex is determined by at most $DN$ of its minimal vectors, because $\dim_{\mathbf R}E\le DN$; so there are at most $\big((2R_0+1)^D+1\big)^{DN}$ such vertices.
Translating a tuple by $GL_N(\mathcal O_L)$ so that its matrix $A$ is a chosen representative of its class, we find at most
\[
 I_0^{1+DN}\big((2R_0+1)^D+1\big)^{DN(p+1)}\le(2|d_L|D)^{12(p+1)D^3N}
\]
orbits in $X_p$, the inequality being a numerical estimate.
\end{proof}
\leanref{Development/AutomorphicArithmetic/AutomorphicVarieties/HermiteNormalForms/Imports.lean}{UpperTheory.V.KT.good}
\leanref{Development/AutomorphicArithmetic/AutomorphicVarieties/HermiteNormalForms/Psdcone.lean}{UpperTheory.V.KT.psdcone}
\leanref{Development/AutomorphicArithmetic/AutomorphicVarieties/HermiteNormalForms/RynePlace.lean}{UpperTheory.V.KT.ryne}
\leanref{Development/AutomorphicArithmetic/AutomorphicVarieties/HermiteNormalForms/Rypos.lean}{UpperTheory.V.KT.rypos}
\leanref{Development/AutomorphicArithmetic/AutomorphicVarieties/HermiteNormalForms/Locfin.lean}{UpperTheory.V.KT.locfin}
\leanref{Development/AutomorphicArithmetic/AutomorphicVarieties/HermiteNormalForms/Assembly.lean}{UpperTheory.V.KT.sx_exact}
\leanref{Development/AutomorphicArithmetic/AutomorphicVarieties/HermiteNormalForms/Stabfin.lean}{UpperTheory.V.KT.stabfin}
\leanref{Development/AutomorphicArithmetic/AutomorphicVarieties/HermiteNormalForms/ImportsAlt.lean}{UpperTheory.V.KT.stab}
\leanref{Development/AutomorphicArithmetic/AutomorphicVarieties/HermiteNormalForms/Mink3.lean}{UpperTheory.V.KT.mink3}
\leanref{Development/AutomorphicArithmetic/AutomorphicVarieties/HermiteNormalForms/ColsMinSmul.lean}{UpperTheory.V.KT.orb}
\leanref{Development/AutomorphicArithmetic/AutomorphicVarieties/HermiteNormalForms/DetPull.lean}{UpperTheory.V.KT.nrmle}
\leanref{Development/AutomorphicArithmetic/AutomorphicVarieties/HermiteNormalForms/Hnf.lean}{UpperTheory.V.KT.hnf}
\leanref{Development/AutomorphicArithmetic/AutomorphicVarieties/HermiteNormalForms/Nearvx.lean}{UpperTheory.V.KT.nearvx}
\leanref{Development/AutomorphicArithmetic/AutomorphicVarieties/HermiteNormalForms/Orbnum.lean}{UpperTheory.V.KT.orbnum}

\begin{lemma}\label{sha-torsion-from-action}
Let $G$ be a group, let $X_0,X_1,\dots$ be sets on which $G$ acts, let $M$ be a $\mathbf Z[G]$-module, and let
\[
 \cdots\to\mathbf Z[X_2]\xrightarrow{d_1}\mathbf Z[X_1]\xrightarrow{d_0}\mathbf Z[X_0]\xrightarrow{\varepsilon}M\to0
\]
be maps of $\mathbf Z[G]$-modules.
Let $m\ge0$, and suppose that $\varepsilon$ is surjective and that the sequence is exact at $\mathbf Z[X_0],\dots,\mathbf Z[X_m]$.
Let $A\ge0$ and $s,b\ge1$ be integers such that, for $p\le m$, $X_p$ has at most $A$ orbits and every stabilizer has at most $s$ elements, and such that $d_m$ maps each element of $X_{m+1}$ to an element of $\ell^1$-norm at most $b$.
Then the torsion subgroup of $H_m(G;M)$ is finite, of order at most $b^{A}s^{mAs^{m}}$.
\end{lemma}
\begin{proof}
\emph{Permutation modules.} Let $X$ be a $G$-set with at most $A$ orbits and stabilizers of order at most $s$, and let $j\ge1$.
By Shapiro's lemma $H_j(G;\mathbf Z[X])$ is the sum over the orbits of the groups $H_j(G_x;\mathbf Z)$ for stabilizers $G_x$.
For a finite group $H$ of order $h\le s$, $H_j(H;\mathbf Z)$ is killed by $h$ and is a quotient of a subgroup of $\mathbf Z[H^j]$, so it is a quotient of $(\mathbf Z/h)^{r}$ with $r\le h^j$, and has at most $s^{s^j}$ elements.
So $H_j(G;\mathbf Z[X])$ has at most $s^{As^j}$ elements.

\emph{The case $m=0$.} Then $H_0(G;M)=M_G$ is the quotient of $\mathbf Z[X_0]_G=\mathbf Z[X_0/G]$, free of rank $A'\le A$, by the subgroup $\Lambda$ generated by the images of the elements of $X_1$, which have $\ell^1$-norm at most $b$.
Let $\Lambda'\supset\Lambda$ be the elements of $\mathbf Z[X_0/G]$ having a multiple in $\Lambda$, so that the torsion of $M_G$ is $\Lambda'/\Lambda$, and let $\Lambda_0\subset\Lambda$ be spanned by $r=\operatorname{rank}\Lambda$ linearly independent generators.
Then $\#(\Lambda'/\Lambda)\le\#(\Lambda'/\Lambda_0)\le\operatorname{covol}(\Lambda_0)\le b^{r}\le b^{A}$ by Hadamard's inequality, since $\operatorname{covol}(\Lambda')\ge1$ and the $\ell^2$-norm is at most the $\ell^1$-norm.

\emph{Induction.} Let $m\ge0$, assume the statement for $m$, and suppose the hypotheses hold for $m+1$.
Let $K=\ker\varepsilon$, the image of $d_0$.
The sequence $\cdots\to\mathbf Z[X_2]\to\mathbf Z[X_1]\to K\to0$ satisfies the hypotheses for $m$, with the sets $X_{p+1}$, so the torsion of $H_m(G;K)$ has order at most $b^As^{mAs^m}$.
The exact sequence $0\to K\to\mathbf Z[X_0]\to M\to0$ gives an exact sequence $H_{m+1}(G;\mathbf Z[X_0])\to H_{m+1}(G;M)\to H_m(G;K)$, whose first term has at most $s^{As^{m+1}}$ elements.
If $A_1\to B_1\to C_1$ is exact, $A_1$ is finite and the torsion of $C_1$ is finite, then the torsion of $B_1$ is finite of order at most $\#A_1$ times the order of the torsion of $C_1$.
So the torsion of $H_{m+1}(G;M)$ has order at most $s^{As^{m+1}}b^As^{mAs^m}\le b^As^{(m+1)As^{m+1}}$.
\end{proof}
\leanref{Development/AutomorphicArithmetic/AutomorphicVarieties/TorsionBoundsChains/LogLeExp.lean}{UpperTheory.V.KH.main}
\leanref{Development/AutomorphicArithmetic/AutomorphicVarieties/TorsionBoundsChains/LogLeExp.lean}{UpperTheory.V.KH.shift}
\leanref{Development/AutomorphicArithmetic/AutomorphicVarieties/TorsionBoundsChains/Imports.lean}{UpperTheory.V.KH.base}
\leanref{Development/AutomorphicArithmetic/AutomorphicVarieties/TorsionBoundsChains/PermCard.lean}{UpperTheory.V.KH.permCard}
\leanref{Development/AutomorphicArithmetic/AutomorphicVarieties/TorsionBoundsChains/PermTrans.lean}{UpperTheory.V.KH.permTrans}
\leanref{Development/AutomorphicArithmetic/AutomorphicVarieties/TorsionBoundsChains/TorsStep.lean}{UpperTheory.V.KH.torsStep}

\begin{lemma}\label{sha-gl-torsion}
There is $C_0$ such that for every totally complex number field $L$ of degree $n$ and discriminant $d_L$, every $E\ge1$ and every $0\le m\le2E+1$,
\[
 \log t(L,4E+6,m)\le\exp\big((2En)^{C_0}\log(2|d_L|)\big).
\]
One can take $C_0=24$.
\end{lemma}
\begin{proof}
Let $N=4E+6$ and $D=nN$, the rank of $\mathcal O_L^N$ over $\mathbf Z$, and let $X_p$ be the sets of tuples of vertices of Lemma~\ref{sha-hermitian-polyhedron}, with their action of $G=GL_N(\mathcal O_L)$.

\emph{The torsion bound.} By Lemma~\ref{sha-hermitian-polyhedron}(a) the augmented nerve complex $\mathbf Z[X_*]\to\mathbf Z$ is exact.
By Lemma~\ref{sha-hermitian-polyhedron}(b) and (c), for $p\le m$ every stabilizer has at most $s=3^{DN}$ elements and $X_p$ has at most $A=(2|d_L|D)^{12(m+1)D^3N}$ orbits.
The boundary of an element of $X_{m+1}$ is an alternating sum of its $m+2$ faces, so it has $\ell^1$-norm at most $b=m+2$.
Lemma~\ref{sha-torsion-from-action}, with $M=\mathbf Z$ and the augmentation as $\varepsilon$, shows that the torsion of $H_m(G;\mathbf Z)$ is finite of order $t(L,N,m)\le b^As^{mAs^m}$.

\emph{The estimate.} Using $\log(m+2)\le m+1$, $\log s\le s$, $s\le e^{2DN}$ and $2m+1\le e^{2m}$,
\[
 \log t(L,N,m)\le A\log(m+2)+mAs^m\log s\le(2m+1)As^{m+1}\le\exp\big(12(m+1)D^3N\log(2|d_L|D)+2(m+1)DN+2m\big).
\]
Since $N\le4E+6$ and $m\le2E+1$, the exponent is at most $(2En)^{24}\log(2|d_L|)$, by a numerical estimate.
\end{proof}
\leanref{Development/AutomorphicArithmetic/AutomorphicVarieties/ModularityArgument/ClassCounting/LemmaKOfSteps.lean}{UpperTheory.V.ZK.STors}
\leanref{Development/AutomorphicArithmetic/AutomorphicVarieties/ModularityArgument/ClassCounting/ExponentLe.lean}{UpperTheory.V.ZK.sTors}
\leanref{Development/AutomorphicArithmetic/AutomorphicVarieties/ModularityArgument/ClassCounting/ExponentLe.lean}{UpperTheory.V.ZK.torsAll}
\leanref{Development/AutomorphicArithmetic/AutomorphicVarieties/HermiteNormalForms/Assembly.lean}{UpperTheory.V.KT.tors_le}
\leanref{Development/AutomorphicArithmetic/AutomorphicVarieties/TorsionBoundsChains/LogLeExp.lean}{UpperTheory.V.KH.mainFace}
\leanref{Development/AutomorphicArithmetic/AutomorphicVarieties/TorsionBoundsChains/LogLeExp.lean}{UpperTheory.V.KH.TorsLe.log_le_exp}
\leanref{Development/AutomorphicArithmetic/AutomorphicVarieties/ModularityArgument/ClassCounting/ExponentLe.lean}{UpperTheory.V.ZK.TorsionExponentBound.exponent_le}

\begin{theorem}\label{sha-soule-bound}
There is a constant $C_K$ with the following property.
Let $\ell$ be a prime and let $L$ be a totally complex number field that contains a primitive $\ell$-th root of unity, and a primitive fourth root of unity if $\ell=2$.
Let $k\ge1$.
Then $H^2(G_{L,S_\ell},\mathbf Z_\ell(k+1))$ is finite and
\[
 \log\#H^2\big(G_{L,S_\ell},\mathbf Z_\ell(k+1)\big)\le\exp\Big(\big(2k[L:\mathbf Q]\big)^{C_K}\log\big(2|d_L|\big)\Big),
\]
where $d_L$ is the discriminant of $L$ and $G_{L,S_\ell}$ is as in Definition~\ref{sha-tate-shafarevich-group}.
\end{theorem}
\begin{proof}
Write $n=[L:\mathbf Q]$.

\emph{Step 1: the bound for $\text{Ш}^1$.}
There is $C_1$ with $\operatorname{len}\text{Ш}^1_\ell(L,k)\log\ell\le\exp\big((2kn)^{C_1}\log(2|d_L|)\big)$ for all $\ell$, $L$ and $k$ as in the statement.
Let $L'$, $c$ and $d$ be as in Lemma~\ref{sha-enlarging}, let $\nu\ge c$, and let $F$ be a Kummer layer of type $(\ell,c,\nu)$ over $L'$.
Lemma~\ref{sha-few-norms} with $e=k-1$, which applies because $2k+1<\ell^c$ and $2k+1<r_2(L')$, and Lemma~\ref{sha-count-homs} with $E=k$ give
\[
 \#\operatorname{Hom}_{(1+\ell^c)^k}(F,\ell^\nu)\le t(L',4k+6,2k+1)\,t(L',4k+6,2k)\,\ell^{X},\qquad X=(b+v_\ell(k))\operatorname{gen}(L')+c+v_\ell(k),
\]
with $b=2k+(2k+1)(c+2k+1)$; the factor $(\ell^\nu)^{r_2(L')}$ appears on both sides and cancels, so the bound does not depend on $\nu$.
By Lemma~\ref{sha-kummer-descent}, $\operatorname{len}\text{Ш}^1_\ell(L',k)\log\ell\le\log t(L',4k+6,2k+1)+\log t(L',4k+6,2k)+X\log\ell$.
Lemma~\ref{sha-gl-torsion} with $E=k$ bounds the first two terms by $\exp\big((2k[L':\mathbf Q])^{24}\log(2|d_{L'}|)\big)$.
For the third, $\ell^c\le2[L':\mathbf Q]$, $v_\ell(k)\le\log_2k$, and $\log h_{L'}$ is bounded in terms of $[L':\mathbf Q]$ and $|d_{L'}|$ as in the proof of Lemma~\ref{sha-torsion-length}.
The bounds $[L':\mathbf Q]\le(n+1)(4k+4)n$ and $|d_{L'}|\le|d_L|^d(2[L':\mathbf Q])^{[L':\mathbf Q]}$ of Lemma~\ref{sha-enlarging} express everything in terms of $n$, $k$ and $|d_L|$, and $\operatorname{len}\text{Ш}^1_\ell(L,k)\le\operatorname{len}\text{Ш}^1_\ell(L',k)$.
This gives Step 1 with $C_1=5\max(24,9)+21$.

\emph{Step 2: from $\text{Ш}^1$ to $H^2$.}
By Poitou-Tate duality \cite[Chapter VIII, \S6]{NSW08}, the kernel of $H^2(G_{L,S_\ell},\mathbf Z_\ell(k+1))\to\prod_{v\mid\ell}H^2(G_{L_v},\mathbf Z_\ell(k+1))$ is dual to $\text{Ш}^1_\ell(L,k)$; $L$ is totally complex, so no real places intervene.
By local Tate duality the local group $H^2(G_{L_v},\mathbf Z_\ell(k+1))$ is dual to $H^0(G_{L_v},\mathbf Q_\ell/\mathbf Z_\ell(-k))$, and the total length of these local groups is at most $n(2+\log_\ell n+k)$.
So $H^2(G_{L,S_\ell},\mathbf Z_\ell(k+1))$ has length at most $\operatorname{len}\text{Ш}^1_\ell(L,k)+n(2+\log_\ell n+k)$, and is finite by Step 1.
Since $\mu_\ell\subset L$ we have $\ell\le n+1$, and Step 1 gives the bound of the statement for a suitable $C_K$.
\end{proof}
\leanref{Development/AutomorphicArithmetic/AutomorphicVarieties/GeneratorCountTuples/LemmaK.lean}{UpperTheory.V.LemmaK}
\leanref{Development/AutomorphicArithmetic/AutomorphicVarieties/GeneratorCountTuples/LemmaK.lean}{UpperTheory.V.kBound}
\leanref{Development/AutomorphicArithmetic/AutomorphicVarieties/ModularityArgument/ClassCounting/LemmaKOfLitN.lean}{UpperTheory.V.ZK.lemmaK}
\leanref{Development/AutomorphicArithmetic/AutomorphicVarieties/ModularityArgument/ClassCounting/LemmaKOfSteps.lean}{UpperTheory.V.ZK.lemmaK_of_steps}
\leanref{Development/AutomorphicArithmetic/AutomorphicVarieties/ModularityArgument/ClassCounting/NumThree.lean}{UpperTheory.V.ZK.sNum}
\leanref{Development/AutomorphicArithmetic/AutomorphicVarieties/ModularityArgument/ClassCounting/SFinal.lean}{UpperTheory.V.ZK.sFinal}

\section{Windows of scales}\label{section-sha-windows}

\begin{lemma}\label{sha-window-pigeonhole}
Let $R\ge0$ be an integer and let $\eta,C_t,X$ be real numbers with $0<\eta\le C_t$ and $X>1$.
Let $\mathrm{Exc}$ be a set of at most $R$ non-negative integers, and put $J=2^{R+2}(C_t/\eta)^{R}$.
Then there is $Y$ with $X^{1/J}\le Y$ and $Y^2\le X$ such that for every $x\in[Y,Y^2]$, every $k\in\mathrm{Exc}$ with $k\le x^{C_t}$ satisfies $k\le x^{\eta}$.
\end{lemma}
\begin{proof}
Put $U=\log X>0$, $b=C_t/\eta\ge1$ and $T=\{\log k/C_t:k\in\mathrm{Exc},\ k\ge1\}$; the integer $k=0$ satisfies the conclusion for every $x>0$.
We show by induction on $R$: for every $U>0$ and every set $T$ of at most $R$ reals there is $y$ with $U/(2^{R+2}b^{R})\le y$, $2y\le U$, and $ba\le u$ for all $a\in T$ and $u\in[y,2y]$ with $a\le u$.
If no $a\in T$ lies in $(U/2b,U]$, take $y=U/2$: an $a\in T$ with $a\le u\le U$ satisfies $a\le U/2b$, so $ba\le U/2\le u$.
Otherwise remove one such $a_0$ from $T$ and apply the induction hypothesis with $U/2b$ in place of $U$.
The resulting $y$ satisfies $U/(2^{R+2}b^{R})\le y$ and $2y\le U/2b\le U$.
For $u\in[y,2y]$ we have $u\le U/2b<a_0$, so $a_0$ imposes no condition, and the other elements of $T$ are handled by the induction hypothesis.

Now let $Y=e^{y}$. Then $X^{1/J}\le Y$ and $Y^2\le X$.
For $x\in[Y,Y^2]$ put $u=\log x\in[y,2y]$.
If $k\in\mathrm{Exc}$ with $k\ge1$ and $k\le x^{C_t}$, then $a=\log k/C_t\le u$, so $(C_t/\eta)a\le u$, that is $k\le x^{\eta}$.
\end{proof}
\leanref{Development/AutomorphicArithmetic/AutomorphicVarieties/LayeredIsomorphisms/Window.lean}{UpperTheory.V.Win.window}
\leanref{Development/AutomorphicArithmetic/AutomorphicVarieties/LayeredIsomorphisms/Window1.lean}{UpperTheory.V.Win.window1}
\leanref{Development/AutomorphicArithmetic/AutomorphicVarieties/LayeredIsomorphisms/SecondBound.lean}{UpperTheory.V.Win.Good}

\begin{definition}\label{sha-window-constants}
\begin{enumerate}
\item A \emph{window constant} $K$ consists of a function $R_0$ from abelian field data to $\mathbf Z_{\ge0}$, functions $R_\ell$ from abelian field data to $\mathbf Z_{\ge0}$ for every integer $\ell$, and a real number $C_K>0$. For an abelian field datum $M$, a finite set $B'$ of primes, abelian field data $L(\ell)$ for $\ell\in B'$, and real numbers $C_t,\varepsilon$, put
\[
 R=R_0(M)+\sum_{\ell\in B'}R_\ell(L(\ell)),\qquad J=2^{R+2}\Big(\frac{2C_KC_t}{\varepsilon}\Big)^{R}.
\]
\item Let $A_1,C_t$ be real numbers. Functions $e,t:\mathbf R\to\mathbf Z_{\ge0}$ are \emph{admitted} for $(A_1,C_t)$ if for every $x\ge2$: $2\le e(x)$, $2^{e(x)}\le x^{A_1}$ and $2\le t(x)\le x^{C_t}$.
\item In the situation of (a), let $(e,t)$ be functions, let $q,d$ be primes and let $X,Y$ be real numbers. Then $Y$ is a \emph{controlled window} for these data if $X^{1/J}\le Y$, $Y^2\le X$, and for every $x\in[Y,Y^2]$:
\begin{enumerate}
\item every compositum $E$ of $(M,e(x),q,d)$ satisfies the $2$-bound $\big(t(x),\exp(x^{4\varepsilon/5})\big)$;
\item for every $\ell\in B'$, every field $L_0$ of $L(\ell)$ and every layer $L_1$ of degree $\ell^{\lfloor\log_\ell R_\ell(L(\ell))\rfloor+3}$ over $L_0$, the field $L_1$ satisfies the $\ell$-bound $\big(t(x),\exp(x^{4\varepsilon/5})\big)$.
\end{enumerate}
\end{enumerate}
\end{definition}
\leanref{Development/AutomorphicArithmetic/AutomorphicVarieties/PrimeSetBoundDataOfPrimitiveChar.lean}{UpperTheory.V.BConst}
\leanref{Development/AutomorphicArithmetic/AutomorphicVarieties/PrimeSetBoundDataOfPrimitiveChar.lean}{UpperTheory.V.BConst.J}
\leanref{Development/AutomorphicArithmetic/AutomorphicVarieties/PrimeSetBoundDataOfPrimitiveChar.lean}{UpperTheory.V.admitted}
\leanref{Development/AutomorphicArithmetic/AutomorphicVarieties/PrimeSetBoundDataOfPrimitiveChar.lean}{UpperTheory.V.IsBWindow}

\begin{theorem}\label{sha-prime-two}
There is a window constant $K$ such that for every $C_t\ge1$ the following holds.
Let $M$ be an abelian field datum, let $B'$ be a finite set of odd primes, and for $\ell\in B'$ let $L(\ell)$ be an abelian field datum containing $\mu_\ell$.
Let $0<\varepsilon<1/52$ and let $A_1$ be real.
Then there is $X_0$ such that for all functions $(e,t)$ admitted for $(A_1,C_t)$, every $X\ge X_0$ and every prime pair $(q,d)$ for $M$ with $\log(qd)\le X^{\varepsilon/(4J)}$, there is a controlled window $Y$ for $(K,C_t,\varepsilon,M,B',L,e,t,q,d,X)$ (Definition~\ref{sha-window-constants}).
\end{theorem}
\begin{proof}
Let $C_K^0$ be the constant of Theorem~\ref{sha-soule-bound}, and put $C_K=\max(C_K^0,1)$; the bound of Theorem~\ref{sha-soule-bound} still holds with $C_K$.
Let $R_0(M)$ be the integer of Proposition~\ref{sha-tower-bounds}(b) for $M$.
For an odd prime $\ell$ and an abelian field datum $A$ containing $\mu_\ell$ that has a field, let $R_\ell(A)$ be an integer given by Proposition~\ref{sha-tower-bounds}(a); otherwise let $R_\ell(A)=0$.
Fix $C_t\ge1$ and the data of the statement, and put $\eta=\varepsilon/(2C_K)$, so that $0<\eta\le C_t$.

\emph{The prime $2$.} There are $x_A$ and, for every admitted $(e,t)$, every $X$ and every prime pair $(q,d)$ with $\log(qd)\le X^{\varepsilon/(4J)}$, a set $\mathrm{Exc}_A$ of at most $R_0(M)$ integers with the following property: if $x\ge x_A$, $X^{1/J}\le x\le X$, and every $k\in\mathrm{Exc}_A$ with $k\le x^{C_t}$ satisfies $k\le x^{\eta}$, then condition (i) of a controlled window holds at $x$.
To see this, take the $2$-base $L_0$ of Proposition~\ref{sha-tower-bounds}(b), with its elements $s_i$ and its integer $c\le16\deg M$, and put $Z=\lfloor X^{C_t}\rfloor$.
Condition (i) of Definition~\ref{sha-tower-data} gives a set $\mathrm{Exc}_A$ of at most $R_0(M)$ integers such that $\operatorname{len}\text{Ш}^1_2(L_1,k)\le R_0(M)\big(\max(m,j_0)+c+\log_2Z\big)$ for every layer $L_1$ of degree $2^{m}$ and every $1\le k\le Z$ outside $\mathrm{Exc}_A$: for each $i$, at most one $k\le Z$ has $|s_i-k|_2<2^{-\lfloor\log_2Z\rfloor}$.
For a compositum $E$ of $(M,e(x),q,d)$ the $2$-base gives a layer $L_1$ of degree $2^{m}$ with $m\le e(x)+2\le(2|A_1|+2)(1+\log x)$.
For $k<t(x)$ not in $\mathrm{Exc}_A$, this bounds $\operatorname{len}\text{Ш}^1_2(L_1,k)$ by a multiple of $1+\log x$, since $\log X\le J\log x$.
For $k\in\mathrm{Exc}_A$ with $k<t(x)$, we have $k\le x^{\eta}$, and condition (ii) bounds $\operatorname{len}\text{Ш}^1_2(L_1,k)$ by $\operatorname{len}H^2(G_{L_2,S_2},\mathbf Z_2(k+1))$ plus a multiple of $1+\log x$, for the layer $L_2$ of degree $2^{j_0}$.
The degree of $L_2$ is bounded in terms of $M$, and $\log(2|d_{L_2}|)$ is at most a constant times $\log(qd)+1\le x^{\varepsilon/4}+1$.
So Theorem~\ref{sha-soule-bound} gives $\log\#H^2\le\exp\big(C x^{\eta C_K}x^{\varepsilon/4}\big)=\exp(Cx^{3\varepsilon/4})$ for a constant $C$.
Each term is therefore at most $C(1+\log x)+2\exp(Cx^{3\varepsilon/4})$, and summing over $k<t(x)\le x^{C_t}$ gives at most $\exp(x^{4\varepsilon/5})$ for $x$ large.

\emph{The primes of $B'$.} For each $\ell\in B'$ the same argument, with the field $L_0$ of $L(\ell)$, the elements $s_i$ of Proposition~\ref{sha-tower-bounds}(a), and $L_1=L_2$ the layer of degree $\ell^{j_0}$ itself, gives $x_\ell$ and sets $\mathrm{Exc}_\ell$ of at most $R_\ell(L(\ell))$ integers with the corresponding property for condition (ii) of a controlled window.
Here the discriminant of $L_2$ does not depend on $x$, and any two fields of $L(\ell)$ and their layers of equal degree are isomorphic.

\emph{The window.} The set $\mathrm{Exc}=\mathrm{Exc}_A\cup\bigcup_{\ell\in B'}\mathrm{Exc}_\ell$ has at most $R$ elements.
Lemma~\ref{sha-window-pigeonhole}, with $\eta$ and $C_t$, gives $Y$ with $X^{1/J}\le Y$, $Y^2\le X$, and the condition on $\mathrm{Exc}$ at every $x\in[Y,Y^2]$; here $2^{R+2}(C_t/\eta)^{R}=J$.
Let $x_1=\max(1,x_A,\sum_\ell|x_\ell|)$ and $X_0=\max(2,x_1^{J})$.
For $X\ge X_0$ and $x\in[Y,Y^2]$ we have $x\ge X^{1/J}\ge x_1$ and $x\le X$, so both parts apply.
\end{proof}
\leanref{Development/AutomorphicArithmetic/AutomorphicVarieties/PrimeSetBoundDataOfPrimitiveChar.lean}{UpperTheory.V.ThmBab}
\leanref{Development/AutomorphicArithmetic/AutomorphicVarieties/LayeredIsomorphisms/ThmBabOfLemmaK.lean}{UpperTheory.V.Win.thmBab_of_lemmaK}
\leanref{Development/AutomorphicArithmetic/AutomorphicVarieties/LayeredIsomorphisms/ThmBabOf.lean}{UpperTheory.V.Win.thmBab_of}
\leanref{Development/AutomorphicArithmetic/AutomorphicVarieties/LayeredIsomorphisms/PartA.lean}{UpperTheory.V.Win.partA}
\leanref{Development/AutomorphicArithmetic/AutomorphicVarieties/LayeredIsomorphisms/PartB.lean}{UpperTheory.V.Win.partB}

\section{The plus part}\label{section-sha-plus}

\subsection{Character pieces and Eisenstein data}\label{subsection-sha-character-pieces}

\begin{definition}\label{sha-character-parts}
Let $\ell$ be a prime and let $\mathcal O\subset\bar{\mathbf Q}_\ell$ be a subring that is a discrete valuation ring finite over $\mathbf Z_\ell$, with ramification index $e_{\mathcal O}=\operatorname{len}_{\mathcal O}(\mathcal O/\ell\mathcal O)$.
Let $\theta:\operatorname{Gal}(\bar{\mathbf Q}/\mathbf Q)\to\mathcal O^\times$ be a character, and let $k\ge0$ and $n\ge1$ be integers.
Let $\mathcal O\otimes\mathbf Q_\ell/\mathbf Z_\ell(\theta^{-1}\chi_{\mathrm{cyc}}^{-k})$ be the discrete module on which $\sigma$ acts by $\theta(\sigma)^{-1}\chi_{\mathrm{cyc}}(\sigma)^{-k}$.
Let $V_{\ell,n}(k,\theta)$ be the $\mathcal O$-length of the subgroup of $H^1\big(\operatorname{Gal}(\bar{\mathbf Q}/\mathbf Q),\mathcal O\otimes\mathbf Q_\ell/\mathbf Z_\ell(\theta^{-1}\chi_{\mathrm{cyc}}^{-k})\big)$ of classes that vanish on the decomposition groups at the primes dividing $n$ and on the inertia groups at all primes.
Put $V_\ell(k,\theta)=V_{\ell,\ell}(k,\theta)$.
A Dirichlet character $\psi$ modulo $N$ with values in $\mathcal O$ is viewed as a character of $\operatorname{Gal}(\bar{\mathbf Q}/\mathbf Q)$ through the action on $N$-th roots of unity, and we put
\[
 V'_\ell(k,\psi)=V_{\ell,\ell N}(k,\psi),\qquad \bar V_\ell(k,\psi)=V'_\ell(k,\psi)/e_{\mathcal O}.
\]
\end{definition}
\leanref{Development/AutomorphicArithmetic/AutomorphicVarieties/CharacterLevelBounds/ShaFamMono.lean}{UpperTheory.V.Chr.VS}
\leanref{Development/AutomorphicArithmetic/AutomorphicVarieties/CharacterLevelBounds/ShaFamMono.lean}{UpperTheory.V.Chr.V}
\leanref{Development/AutomorphicArithmetic/AutomorphicVarieties/HeightAndVanishing/DirichletPrimeToLPartMulLPowerPart.lean}{UpperTheory.V.HV.Vlen}
\leanref{Development/AutomorphicArithmetic/AutomorphicVarieties/HeightAndVanishing/DirichletPrimeToLPartMulLPowerPart.lean}{UpperTheory.V.HV.Vn}

\begin{lemma}\label{sha-plus-by-characters}
Let $N\ge1$, let $H\subset(\mathbf Z/N)^\times$ be a subgroup, and let $E$ be a CM field isomorphic to the subfield of $\mathbf Q(\zeta_N)$ fixed by the automorphisms $\zeta_N\mapsto\zeta_N^{u}$ with $u\in H$.
Let $\ell$ be an odd prime, let $\mathcal O$ be as in Definition~\ref{sha-character-parts} and containing the $\varphi(N)$-th roots of unity, and let $k\ge1$.
Let $X_k$ be the set of characters $\psi$ of $(\mathbf Z/N)^\times$ with values in $\mathcal O$, trivial on $H$, with $\psi(-1)=(-1)^k$.
\begin{enumerate}
\item If $\ell\nmid[E:\mathbf Q]$, then $e_{\mathcal O}\operatorname{len}\text{Ш}^+_\ell(E,k)\le\sum_{\psi\in X_k}V'_\ell(k,\psi)$.
\item In general, $e_{\mathcal O}\operatorname{len}\text{Ш}^+_\ell(E,k)\le\sum_{\psi\in X_k}\big(V'_\ell(k,\psi)+\lambda e_{\mathcal O}\big)$, where, with $n=[E:\mathbf Q]$ and $\omega(N)$ the number of primes dividing $N$,
\[
 \lambda=(n+n^2)\,\nu+\omega(N)\,(n+1)\,v_\ell(n),\qquad\nu=\begin{cases}\max(v_\ell(N),1)+v_\ell(k)&\text{if }\ell\mid N\text{ or }\ell-1\mid k,\\0&\text{otherwise.}\end{cases}
\]
If $N\ge2$, then $\lambda\le(Nk)^6$.
\end{enumerate}
\end{lemma}
\begin{proof}
Let $\Gamma=\operatorname{Gal}(E/\mathbf Q)$, a quotient of $(\mathbf Z/N)^\times$ of order $n$, so that $\operatorname{Gal}(\bar{\mathbf Q}/E)$ is the kernel of $\operatorname{Gal}(\bar{\mathbf Q}/\mathbf Q)\to\Gamma$; and $E\subset\mathbf Q(\zeta_N)$.
Let $A=\mathcal O\otimes\mathbf Q_\ell/\mathbf Z_\ell(-k)$, and for a character $\psi$ of $\Gamma$ let $A(\psi)$ be $\mathcal O\otimes\mathbf Q_\ell/\mathbf Z_\ell$ with $\sigma$ acting by $\psi(\sigma)^{-1}\chi_{\mathrm{cyc}}(\sigma)^{-k}$, as in Definition~\ref{sha-character-parts}.
For a closed subgroup $G$ of $\operatorname{Gal}(\bar{\mathbf Q}/\mathbf Q)$ and a coefficient module, the \emph{Selmer group} below is the subgroup of $H^1(G,-)$ of classes that vanish on $G\cap I$ for every inertia group $I$ and on $G\cap D$ for the decomposition groups $D$ at $\ell$.

\emph{Step 1.} Let $\text{Ш}_{\mathcal O}$ be the Selmer group of $\operatorname{Gal}(\bar{\mathbf Q}/E)$ with coefficients $A$, which is $\text{Ш}^1_\ell(E,k)\otimes\mathcal O$, with the action of $\Gamma$.
Then $e_{\mathcal O}\operatorname{len}\text{Ш}^+_\ell(E,k)$ is the $\mathcal O$-length of the $(-1)^k$-eigenspace of complex conjugation on $\text{Ш}_{\mathcal O}$.
Since $2$ is a unit and $\mathcal O$ contains the values of the characters of $\Gamma$, this length is at most $\sum_{\psi\in X_k}\operatorname{len}\text{Ш}_{\mathcal O}^{\psi}$, where $\text{Ш}_{\mathcal O}^{\psi}$ is the subgroup on which $\tau\in\Gamma$ acts by $\psi(\tau)$.

\emph{Step 2: comparison with the group over $\mathbf Q$.}
Since $\psi$ is trivial on $\operatorname{Gal}(\bar{\mathbf Q}/E)$, the module $A(\psi)$ restricted to $\operatorname{Gal}(\bar{\mathbf Q}/E)$ is $A$, and $\text{Ш}_{\mathcal O}^{\psi}$ is the subgroup of $\Gamma$-invariants of the Selmer group of $\operatorname{Gal}(\bar{\mathbf Q}/E)$ with coefficients $A(\psi)$.
The group defining $V_\ell(k,\psi)$ is the Selmer group of $\operatorname{Gal}(\bar{\mathbf Q}/\mathbf Q)$ with coefficients $A(\psi)$, and restriction maps it to $\text{Ш}_{\mathcal O}^{\psi}$.
The length of $(A(\psi))^{\operatorname{Gal}(\bar{\mathbf Q}/\mathbf Q(\zeta_N))}$, and of its invariants under the decomposition groups at $\ell$ of $\mathbf Q(\zeta_N)$, is at most $\nu e_{\mathcal O}$: these invariants are $0$ unless $\chi_{\mathrm{cyc}}^{k}\equiv1\pmod\ell$ on these groups, that is, unless $\ell\mid N$ or $\ell-1\mid k$; and they are killed by $x^k-1$ for every $x$ in the image of $\chi_{\mathrm{cyc}}$, which contains an element with $v_\ell(x^{k}-1)=\max(v_\ell(N),1)+v_\ell(k)$.
A class in $\text{Ш}_{\mathcal O}^{\psi}$ is the restriction of a class over $\mathbf Q$, up to the cokernel of $H^1(\operatorname{Gal}(\bar{\mathbf Q}/\mathbf Q),A(\psi))\to H^1(\operatorname{Gal}(\bar{\mathbf Q}/E),A(\psi))^{\Gamma}$, which has length at most $n^2\nu e_{\mathcal O}$.
That class satisfies the local conditions over $\mathbf Q$ up to: the kernel of restriction from the decomposition group at $\ell$ to its intersection with $\operatorname{Gal}(\bar{\mathbf Q}/E)$, of length at most $n\nu e_{\mathcal O}$; and, for each prime $p\mid N$ other than $\ell$, the kernel of restriction from the inertia group at $p$ to its intersection with $\operatorname{Gal}(\bar{\mathbf Q}/E)$, of length at most $n\,v_\ell(n)e_{\mathcal O}$.
At the primes $p\nmid N\ell$ the inertia groups lie in $\operatorname{Gal}(\bar{\mathbf Q}/E)$, so there is no condition to recover.
Hence
\[
 \operatorname{len}\text{Ш}_{\mathcal O}^{\psi}\le V_\ell(k,\psi)+\big((n+n^2)\nu+\omega(N)\,n\,v_\ell(n)\big)e_{\mathcal O}.
\]
If $\ell\nmid n$, the cohomology of $\Gamma$ with coefficients in $\ell$-primary modules vanishes, restriction identifies the two groups, and $\operatorname{len}\text{Ш}_{\mathcal O}^{\psi}=V_\ell(k,\psi)$.

\emph{Step 3: the condition at the primes dividing $N$.}
A class in the group defining $V_\ell(k,\psi)$ is unramified at each prime $p\mid N$ with $p\ne\ell$, and its restriction to the decomposition group at $p$ lies in the kernel of restriction to the inertia group.
That kernel has length at most $v_\ell(n)e_{\mathcal O}$, since $\psi^{n}=1$ and Frobenius at $p$ acts on the inertia invariants of $A(\psi)$ by $\psi(p)^{-1}p^{-k}\ne1$; it is $0$ if $\ell\nmid n$.
So $V_\ell(k,\psi)\le V'_\ell(k,\psi)+\omega(N)\,v_\ell(n)e_{\mathcal O}$, with equality of $V_\ell$ and $V'_\ell$ if $\ell\nmid n$.

Steps 1 to 3 give (a) and (b).
For the bound on $\lambda$, let $N\ge2$; since $n\le N$, each of $v_\ell(N)$, $v_\ell(n)$, $v_\ell(k)$, $\omega(N)$ and $n$ is at most $Nk$, $\nu\le2Nk$, and the terms of $\lambda$ add up to at most $(Nk)^6$.
\end{proof}
\leanref{Development/AutomorphicArithmetic/AutomorphicVarieties/HeightAndVanishing/RedPrime.lean}{UpperTheory.V.HV.redPrime}
\leanref{Development/AutomorphicArithmetic/AutomorphicVarieties/HeightAndVanishing/RedAll.lean}{UpperTheory.V.HV.redAll}
\leanref{Development/AutomorphicArithmetic/AutomorphicVarieties/HeightAndVanishing/LayerLe.lean}{UpperTheory.V.HV.lossAll}
\leanref{Development/AutomorphicArithmetic/AutomorphicVarieties/HeightAndVanishing/LossLe.lean}{UpperTheory.V.HV.lossLe}

\begin{lemma}\label{sha-control}
Let $F$ be a number field with algebraic closure $\bar F$, let $\ell$ be an odd prime, let $\mathcal O$ be as in Definition~\ref{sha-character-parts}, let $n\ne0$ be an integer, and let $\psi:\operatorname{Gal}(\bar F/F)\to\mathcal O^\times$ be a continuous character.
Let $A$ be $\mathcal O\otimes\mathbf Q_\ell/\mathbf Z_\ell$ with $\sigma$ acting by $\psi(\sigma)\chi_{\mathrm{cyc}}(\sigma)^{n}$.
Let $U$ be an open normal subgroup of $\operatorname{Gal}(\bar F/F)$ on which $\psi$ is trivial, and let $U_\infty=U\cap\ker\chi_{\mathrm{cyc}}$.
Then restriction $H^1(U,A)\to H^1(U_\infty,A)$ is injective.
\end{lemma}
\begin{proof}
Let $U'=U\cap\ker(\chi_{\mathrm{cyc}}\bmod\ell)$, of index dividing $\ell-1$ in $U$.
Corestriction composed with restriction is multiplication by this index, which is prime to $\ell$, so restriction from $U$ to $U'$ is injective on the $\ell$-primary group $H^1(U,A)$.
The group $U'/U_\infty$ is isomorphic through $\chi_{\mathrm{cyc}}$ to an open subgroup of $1+\ell\mathbf Z_\ell$, so it is procyclic with a generator $\gamma$, and $\gamma$ acts on $A$ by $\chi_{\mathrm{cyc}}(\gamma)^{n}\ne1$.
By inflation-restriction, the kernel of restriction from $U'$ to $U_\infty$ is $H^1(U'/U_\infty,A)=A/(\chi_{\mathrm{cyc}}(\gamma)^n-1)A$, which is $0$ because $A$ is divisible.
\end{proof}
\leanref{Development/AutomorphicArithmetic/AutomorphicVarieties/GaloisLibrary/SelmerTransport/Control.lean}{UpperTheory.V.Base.control}

\begin{theorem}\label{sha-trivial-character}
Let $\ell\ge5$ be a prime, let $\mathcal O$ be as in Definition~\ref{sha-character-parts}, and let $k\ge2$ be even with $\ell-1\mid2k$.
Then $V_\ell(k,1)=0$: every class in $H^1\big(\operatorname{Gal}(\bar{\mathbf Q}/\mathbf Q),\mathcal O\otimes\mathbf Q_\ell/\mathbf Z_\ell(-k)\big)$ that is unramified at every prime and vanishes on the decomposition group at $\ell$ is zero.
\end{theorem}
\begin{proof}
Let $\text{Ш}$ be the group of such classes.
It is $\ell$-power torsion, so it suffices to show that it has no nonzero element killed by $\ell$.

Suppose first that $\ell-1\mid k$.
Then $\chi_{\mathrm{cyc}}^{-k}\equiv1\pmod\ell$, so $\operatorname{Gal}(\bar{\mathbf Q}/\mathbf Q)$ acts trivially on the $\ell$-torsion of the coefficients.
Let $x\in\text{Ш}$ with $\ell x=0$, represented by a cocycle $z$. Then $\ell z$ is the coboundary of some $m$; choose $m'$ with $\ell m'=m$ and replace $z$ by $z-dm'$, which represents $x$ and takes values in this $\ell$-torsion.
Since the action on the $\ell$-torsion is trivial, $z$ is a continuous homomorphism, and it vanishes on the inertia groups at the primes $p\ne\ell$ because $x$ is unramified there and the inertia groups at $p\ne\ell$ act trivially on the coefficients.
An abelian extension of $\mathbf Q$ unramified outside $\ell$ lies in $\mathbf Q(\mu_{\ell^\infty})$ by the Kronecker-Weber theorem, so $z$ vanishes on $U_\infty=\operatorname{Gal}(\bar{\mathbf Q}/\mathbf Q(\mu_{\ell^\infty}))$, and the restriction of $x$ to $U_\infty$ is $0$.
By Lemma~\ref{sha-control}, with $F=\mathbf Q$, $\psi=1$, $n=-k$ and $U=\operatorname{Gal}(\bar{\mathbf Q}/\mathbf Q)$, the class $x$ is $0$.

Suppose now that $\ell-1\nmid k$.
Since $k$ is even and $\ell-1\mid2k$, we get $\ell\equiv1\pmod4$ and $(\ell-1)/2\mid k$.
Let $K=\mathbf Q(\sqrt\ell)$, the real quadratic field of discriminant $\ell$.
The character $\chi_{\mathrm{cyc}}^{(\ell-1)/2}\bmod\ell$ is the quadratic character of $K$, so $\chi_{\mathrm{cyc}}^{-k}\equiv1\pmod\ell$ on $\operatorname{Gal}(\bar{\mathbf Q}/K)$.
Choose $u\in\mathbf Z_\ell^\times$ with $u^k\not\equiv1\pmod\ell$, which exists because $\ell-1\nmid k$, and $\delta$ in the inertia group at $\ell$ with $\chi_{\mathrm{cyc}}(\delta)=u$; then $\delta\notin\operatorname{Gal}(\bar{\mathbf Q}/K)$, and $\delta$ acts on the coefficients by $\alpha(\delta)=u^{-k}$ with $\alpha(\delta)-1$ a unit.
Let $x\in\text{Ш}$ with $\ell x=0$, represented by a cocycle $z$ with values in the $\ell$-torsion of the coefficients; adding a coboundary, we may assume $z(\delta)=0$.
The group $\operatorname{Gal}(\bar{\mathbf Q}/K)$ acts trivially on the $\ell$-torsion of the coefficients, so $z$ restricts to a homomorphism on it.
It vanishes on the inertia groups at the primes $p\ne\ell$.
At a place above $\ell$, $z$ is the coboundary of some $m$ on the decomposition group; evaluating at an element $t$ of the inertia group with $\chi_{\mathrm{cyc}}(t)=u$ shows that $\ell m=0$, so $z(g)=(\alpha(g)-1)m=0$ for $g$ in the decomposition group and in $\operatorname{Gal}(\bar{\mathbf Q}/K)$.
So $z$ restricts to a homomorphism on $\operatorname{Gal}(\bar{\mathbf Q}/K)$ that is unramified everywhere and trivial on the decomposition groups above $\ell$; it factors through the class group of $K$ and is killed by $h_K$.
The class number $h_K$ is less than $\ell$, hence prime to $\ell$, so $z$ vanishes on $\operatorname{Gal}(\bar{\mathbf Q}/K)$.
Every other $g$ is $\delta g'$ with $g'\in\operatorname{Gal}(\bar{\mathbf Q}/K)$, and $z(g)=z(\delta)+\alpha(\delta)z(g')=0$.
So $x=0$.
\end{proof}
\leanref{Development/AutomorphicArithmetic/AutomorphicVarieties/HeightAndVanishing/Counts.lean}{UpperTheory.V.HV.PropZ}
\leanref{Development/AutomorphicArithmetic/AutomorphicVarieties/LevelCohomology/PropZ.lean}{UpperTheory.V.X.propZ}
\leanref{Development/AutomorphicArithmetic/AutomorphicVarieties/LevelCohomology/TorsClass.lean}{UpperTheory.V.X.zOne}
\leanref{Development/AutomorphicArithmetic/AutomorphicVarieties/LevelCohomology/ZqField.lean}{UpperTheory.V.X.zQuad}

The remaining characters are treated by deformation theory.
For a ring $A$, let $\mathcal G_3(A)=GL_3(A)\rtimes\{\pm1\}$, where $-1$ acts by $g\mapsto J\,{}^tg^{-1}J$ for the antidiagonal matrix $J$ with entries $1$; this is the group $\mathcal G_3$ of Clozel, Harris and Taylor \cite[\S2.1]{CHT08} without its similitude factor.
For a character $\eta$ with values in $\{\pm1\}$ and a character $\Psi$, let $r_0(\eta,\Psi)(g)=\big(\operatorname{diag}(\Psi(g),1,\Psi(g)^{-1}),\eta(g)\big)$.

\begin{definition}\label{sha-ordinary-deformation}
Let $\ell$ be an odd prime, $\mathcal O$ as in Definition~\ref{sha-character-parts}, $k\ge1$, $N\ge1$ even, $\psi$ a Dirichlet character modulo $N$ with values in $\mathcal O$, and $K$ an imaginary quadratic field.
For a prime $p$ write $\psi_p$ for the component of $\psi$ at $p$, and $\psi'$ for the part of $\psi$ of order prime to $\ell$; let $\chi'_\ell(\psi)=\psi'\cdot(\psi'_\ell)^{-1}$.
\begin{enumerate}
\item The quadruple $(\ell,k,\psi,K)$ is an \emph{Eisenstein datum} if $\ell\nmid\operatorname{ord}\psi$, $\psi_\ell^{2}=1$, $\psi(-1)=(-1)^k$, $\ell\nmid d_Kh_K$, $d_K$ is prime to $N$, $K$ is not contained in the field cut out by $\chi_{\mathrm{cyc}}\bmod\ell$ and $\chi'_\ell(\psi)$, and $\ell-1\nmid2k$.
\item \emph{Deformation data} for an Eisenstein datum $(\ell,k,\psi,K)$ consist of the character $\eta$ of $\operatorname{Gal}(\bar{\mathbf Q}/\mathbf Q)$ with kernel $\operatorname{Gal}(\bar{\mathbf Q}/K)$, a complex conjugation $c$, a decomposition group at $\ell$, and the character $\Psi=\chi_{\mathrm{cyc}}^{k}\psi$; on the inertia group at $\ell$, $\Psi=\chi_{\mathrm{cyc}}^{k}\psi_\ell$, and $\eta$ is unramified at $\ell$.
\item Let $g_1\in\operatorname{Gal}(\bar{\mathbf Q}/\mathbf Q)$ and let $\mathfrak B$ be a place above $\ell$. Let $A$ be a complete local $\mathcal O$-algebra with maximal ideal $\mathfrak m_A$. An \emph{ordinary point} with values in $A$ is a homomorphism $r:\operatorname{Gal}(\bar{\mathbf Q}/\mathbf Q)\to\mathcal G_3(A)$ that has multiplier $\eta$, reduces to $r_0(\eta,\Psi)$ modulo $\mathfrak m_A$, is continuous for the $\mathfrak m_A$-adic topology, is normalized by requiring the matrix part of $r(c)$ to be $1$ and by the vanishing of the $(1,2)$ and $(2,1)$ entries of $r(g_1)-r(cg_1^{-1}c)$, satisfies $r(s)=1$ for every $s$ in an inertia group at a prime $p\ne\ell$ with $\eta(s)=\Psi(s)=1$, and is ordinary at $\mathfrak B$: after conjugation by a matrix congruent to $1$ modulo $\mathfrak m_A$, its restriction to the decomposition group at $\mathfrak B$ is upper triangular on the kernel of $\eta$, with diagonal entries $\Psi,1,\Psi^{-1}$ on inertia.
\end{enumerate}
\end{definition}
\leanref{Development/AutomorphicArithmetic/AutomorphicVarieties/HeightAndVanishing/Counts.lean}{UpperTheory.V.HV.GD}
\leanref{Development/AutomorphicArithmetic/AutomorphicVarieties/UnitaryThreeGroups/RhoKApply.lean}{UpperTheory.V.U3.Dat}
\leanref{Development/AutomorphicArithmetic/AutomorphicVarieties/CurveDVRSpecialization/RatTorsionFree.lean}{UpperTheory.V.Pr.OrdPt}
\leanref{Development/AutomorphicArithmetic/AutomorphicVarieties/UnitaryThreeGroups/UnitaryThreeReflectionRight.lean}{UpperTheory.V.U3.r0}
\leanref{Development/AutomorphicArithmetic/AutomorphicVarieties/UnitaryThreeGroups/UnitaryThreeReflectionRight.lean}{UpperTheory.V.U3.Normalised}

\subsection{Ordinary determinants and Hecke algebras}\label{subsection-sha-ordinary-determinants}

\begin{situation}\label{sha-hecke-situation}
Let $\ell$ be an odd prime and $L$ a CM field, with maximal totally real subfield $L^+$ and a complex conjugation $c$, such that $L/L^+$ is unramified at every finite place and $L$ contains no primitive $\ell$-th root of unity.
Let $\mathcal O$ be the ring of integers of a finite extension of $\mathbf Q_\ell$ that contains the images of all embeddings of $\mathcal O_L$, with residue field $k$.
Let $S$ be a finite set of places of $L^+$, all split in $L$, consisting of the places above $\ell$, a set $S_{\mathrm a}$ of auxiliary places prime to $\ell$, and a place $v_0\notin S_{\mathrm a}$ with residue field $\mathbf F_p$ for an odd prime $p\ne\ell$ unramified in $L$; for $v\in S$ fix a place $\tilde v$ of $L$ above $v$.
Let $\Gamma$ be the Galois group of the maximal extension of $L^+$ unramified outside $S$ and the infinite places, $\Gamma_0\subset\Gamma$ the subgroup fixing $L$, $\eta$ the character of $\Gamma$ with kernel $\Gamma_0$, $\varepsilon$ the $\ell$-adic cyclotomic character and $\Gamma'=\Gamma_0\cap\ker\varepsilon$.
For $v\mid\ell$ let $D_v$ be the decomposition group of $\operatorname{Gal}(\bar L/L)$ at a place above $\tilde v$, with its map to $\Gamma_0$, $I_v\subset D_v$ its inertia group, and $\mu_v\subset D_v$ the inverse image of the $\ell$-primary part of the abelianization of $D_v$; let $I_u\subset\Gamma_0$, for $u\in S_{\mathrm a}$, and $I_{v_0}\subset D_{v_0}\subset\Gamma_0$ be inertia and decomposition groups at $\tilde u$ and $\tilde v_0$.
Let $n\ge1$, and for an $\mathcal O$-algebra $R$ and $g\in\Gamma_0$ let $g^*=\varepsilon(g)^{1-n}(cgc^{-1})^{-1}\in R[\Gamma_0]$.
A \emph{residual datum} of rank $n$ consists of pairwise distinct characters $\bar\chi_1,\dots,\bar\chi_n:\Gamma_0\to k^\times$ with open kernels, such that the determinant $\bar D$ of $\bar\chi_1\oplus\dots\oplus\bar\chi_n$ takes the same value at $y$ and at the image of $y$ under $g\mapsto g^*$ for all $y$; characters $\bar\psi_{v,i}:D_v\to k^\times$ with open kernels, for $v\mid\ell$ and $1\le i\le n$; and characters $\theta_{v,i}:\mu_v\to\mathcal O^\times$ trivial on the closed commutator subgroup of $D_v$.
For a complete local $\mathcal O$-algebra $R$ with residue field $k$, an \emph{ordinary pair} over $R$ is a continuous determinant $D$ of $\Gamma_0$ of dimension $n$ with values in $R$ \cite{Che14}, reducing to $\bar D$, together with continuous characters $\psi_{v,i}:D_v\to R^\times$ lifting $\bar\psi_{v,i}$ and equal to $\theta_{v,i}$ on $\mu_v$, such that:
\begin{enumerate}
\item the coefficients of the characteristic polynomials of $D$ at $g$ and at $g^*$ agree, for $g\in\Gamma_0$;
\item (\emph{ordinary relations}) $\operatorname{tr}_D\big(h\prod_{i=1}^n(g_i-\psi_{v,i}(g_i))\big)=0$ for $v\mid\ell$, $g_1,\dots,g_n\in D_v$ and $h\in\Gamma_0$;
\item $\operatorname{tr}_D\big(h(g-1)\big)=0$ for $u\in S_{\mathrm a}$, $g\in I_u\cap\bigcap_i\ker\bar\chi_i$ and $h\in\Gamma_0$.
\end{enumerate}
Let $B$ be the \emph{universal ordinary determinant ring}: a complete local $\mathcal O$-algebra with an ordinary pair $(D^{\mathrm u},\psi^{\mathrm u})$ from which every ordinary pair over every $R$ is obtained by a unique $\mathcal O$-algebra map $B\to R$.
Let $\Lambda$ be the \emph{Iwasawa algebra}: the universal ring of families $(\psi_{v,i}|_{I_v})$ of restrictions to inertia of characters as above, with the map $\Lambda\to B$ given by $\psi^{\mathrm u}$.
\end{situation}
\leanref{Development/AutomorphicArithmetic/AutomorphicVarieties/WeilGroupCharacters/ContinuousResP.lean}{UpperTheory.V.W.HidaSetting}
\leanref{Development/AutomorphicArithmetic/AutomorphicVarieties/WeilGroupCharacters/UniversalOrdinaryDeterminantRing.lean}{UpperTheory.V.W.Resid}
\leanref{Development/AutomorphicArithmetic/AutomorphicVarieties/WeilGroupCharacters/UniversalOrdinaryDeterminantRing.lean}{UpperTheory.V.W.UniversalOrdinaryDeterminantRing}
\leanref{Development/AutomorphicArithmetic/AutomorphicVarieties/WeilGroupCharacters/UniversalOrdinaryDeterminantRing.lean}{UpperTheory.V.W.LamObj}

\begin{definition}\label{sha-hecke-arcs}
In Situation~\ref{sha-hecke-situation}:
\begin{enumerate}
\item A \emph{weight} is a family $h=(h_{v,j,\tau})$ of integers, for $v\mid\ell$, $1\le j\le n$ and embeddings $\tau:L_{\tilde v}\to\bar{\mathbf Q}_\ell$; it is \emph{increasing} if $j\mapsto h_{v,j,\tau}$ is strictly increasing for all $v,\tau$. A \emph{point} is an $\mathcal O$-algebra map $x:B\to\bar{\mathbf Q}_\ell$, with characters $\psi_{x,v,j}=x\circ\psi^{\mathrm u}_{v,j}$. For an increasing weight $h$, $x$ is \emph{arithmetic of weight $h$} if for all $v,j$ there is a character $\phi_{v,j}$ of $I_v$ with open kernel such that $\psi_{x,v,j}(\sigma)\phi_{v,j}(\sigma)=\prod_\tau\tau(\operatorname{Art}^{-1}\sigma)^{-h_{v,j,\tau}}$ for $\sigma\in I_v$, where $\operatorname{Art}$ is the local Artin map; it is \emph{of distinct type} if moreover $\phi_{v,1},\dots,\phi_{v,n}$ can be chosen pairwise distinct for each $v$.
\item An \emph{arc} consists of a finite field $\kappa$ containing $k$, a continuous representation $\rho:\Gamma_0\to GL_n(\kappa[[T]])$ with a symmetric invertible matrix $P$ such that ${}^t\rho(g)P\rho(cgc^{-1})=\varepsilon(g)^{1-n}P$ for $g\in\Gamma_0$, whose determinant reduces modulo $T$ to $\bar D$ and which is trivial on $I_u\cap\bigcap_i\ker\bar\chi_i$ for $u\in S_{\mathrm a}$, and continuous characters $\psi_{v,i}:D_v\to\kappa[[T]]^\times$ lifting $\bar\psi_{v,i}$ such that $\rho|_{D_v}$ is conjugate over the fraction field $\mathcal K$ of $\kappa[[T]]$ to an upper triangular representation with diagonal $\psi_{v,1},\dots,\psi_{v,n}$. It is \emph{attached} to an $\mathcal O$-algebra map $f:B\to\kappa[[T]]$ if $f$ carries $(D^{\mathrm u},\psi^{\mathrm u})$ to the determinant of $\rho$ and the $\psi_{v,i}$.
\item An arc is \emph{nice of type I}, respectively \emph{of type E}, if:
(N1) for each $v\mid\ell$, $\psi_{v,i}\ne\psi_{v,j}$ for $i\ne j$, and $\psi_{v,i}\ne\psi_{v,j}\varepsilon$ for $i<j$;
(N2) $\rho\otimes\bar{\mathcal K}$ is irreducible on $\Gamma'$; respectively, $\rho\otimes\mathcal K$ is conjugate to a direct sum $\rho_1\oplus\rho_2$ of representations of dimensions $a,b\ge1$, each with a matrix $P_i$ as in (2), each irreducible on $\Gamma'$ over $\bar{\mathcal K}$, with no nonzero $\Gamma'$-equivariant map from $\rho_1\otimes\bar{\mathcal K}$ to $\rho_2\otimes\bar{\mathcal K}$;
(N3) every open subset of $\Gamma_0$ that meets $\Gamma'$ contains some $\sigma\in\Gamma'$ at which the characteristic polynomial of $\rho(\sigma)$ is separable;
(N4) the only matrix $X$ over $\mathcal K$ with $\varepsilon(g)\rho(g)X=X\rho(g)$ for all $g\in D_{v_0}$ is $X=0$.
\item A prime $\mathfrak q$ of $B$ is \emph{stable} if there are $g_1,\dots,g_{n^2},h_1,\dots,h_{n^2}\in\Gamma_0$ with $\det\big(\operatorname{tr}_{D^{\mathrm u}}(g_sh_{s'})\big)_{1\le s,s'\le n^2}\notin\mathfrak q$. By Remark~\ref{sha-patching-chain-burnside}, this holds if and only if $D^{\mathrm u}\bmod\mathfrak q$, over an algebraic closure of the residue field of $\mathfrak q$, is the determinant of an absolutely irreducible representation. Let $E_0$ be the quotient of $B[\Gamma_0]$ by the closure of the two-sided ideal generated by the Cayley-Hamilton ideal of $D^{\mathrm u}$, the elements $\prod_i(g_i-\psi^{\mathrm u}_{v,i}(g_i))$ and $g-1$ of (2) and (3) of Situation~\ref{sha-hecke-situation}, and their images under $g\mapsto g^*$. A prime $\mathfrak q$ is \emph{thick} if it is not stable and $E_0\otimes_B\bar\kappa(\mathfrak q)$ is not semisimple, for an algebraic closure $\bar\kappa(\mathfrak q)$ of the residue field of $\mathfrak q$.
\end{enumerate}
\end{definition}
\leanref{Development/AutomorphicArithmetic/AutomorphicVarieties/WeilGroupCharacters/ArithmeticPointClassicality.lean}{UpperTheory.V.W.UniversalOrdinaryDeterminantRing.Pt.Arith}
\leanref{Development/AutomorphicArithmetic/AutomorphicVarieties/WeilGroupCharacters/ArithmeticPointClassicality.lean}{UpperTheory.V.W.UniversalOrdinaryDeterminantRing.Pt.Good}
\leanref{Development/AutomorphicArithmetic/AutomorphicVarieties/WeilGroupCharacters/UniversalOrdinaryDeterminantRing.lean}{UpperTheory.V.W.Point}
\leanref{Development/AutomorphicArithmetic/AutomorphicVarieties/WeilGroupCharacters/UniversalOrdinaryDeterminantRing.lean}{UpperTheory.V.W.Point.NiceI}
\leanref{Development/AutomorphicArithmetic/AutomorphicVarieties/WeilGroupCharacters/UniversalOrdinaryDeterminantRing.lean}{UpperTheory.V.W.Point.NiceE}
\leanref{Development/AutomorphicArithmetic/AutomorphicVarieties/PatchingModules/PairedDualFirstOrderDualMatrixSpace.lean}{UpperTheory.V.Pa.DiscSt}
\leanref{Development/AutomorphicArithmetic/AutomorphicVarieties/WeilGroupCharacters/UniversalOrdinaryDeterminantRing.lean}{UpperTheory.V.W.UniversalOrdinaryDeterminantRing.IsThick}

\begin{lemma}\label{sha-hecke-data}
Assume that Galois representations exist, in the sense of Definition~\ref{gal-galois-representations}, relative to the maps satisfying the laws of Definition~\ref{gal-llc-laws}. In Situation~\ref{sha-hecke-situation}, let $U(n)$ be the unitary group over $L^+$ of the Hermitian form $\sum_ix_i\bar x_i$ on $L^n$, and let $\iota:\bar{\mathbf Q}_\ell\cong\mathbf C$. Then there are:
\begin{enumerate}
\item a representation $\tau_0$ of $K_{\mathrm a}=\prod_{u\in S_{\mathrm a}}GL_n(\mathcal O_{L_{\tilde u}})$ on $\mathcal O^m$, with open kernel, with the following property. Let $(t_i)$ be a family of Teichm\"uller lifts of the $\bar\chi_i$: characters $t_i:\Gamma_0\to\mathcal O^\times$ with $t_i\bmod\mathfrak m_{\mathcal O}=\bar\chi_i$ and $t_i^{q-1}=1$, where $q$ is the number of elements of $k$. Let $M\ne0$ be a complex representation of $\prod_uGL_n(L_{\tilde u})$, and for each $u$ let $\sigma_u$ be an irreducible smooth representation of $GL_n(L_{\tilde u})$ such that $M$, restricted to the factor at $u$, is a sum of copies of $\sigma_u$. Say that $M$ \emph{contains} $\iota\tau_0$ if there is a nonzero $F=(F_1,\dots,F_m)\in M^m$ with $\sum_j\iota\tau_0(k)_{ij}\,kF_j=F_i$ for all $k\in K_{\mathrm a}$ and $1\le i\le m$. Then $M$ contains $\iota\tau_0$ if and only if, for every $u$ and every map $\operatorname{rc}$ that satisfies the laws (Definition~\ref{gal-llc-laws}), the Weil-Deligne representation $\operatorname{rc}(\sigma_u\nu^{(1-n)/2})$ has monodromy $0$ and restriction $\bigoplus_i\iota t_i$ to the inertia group at a place of $\bar L$ above $\tilde u$;
\item a character $\theta$ of the diagonal torus $\prod_{v\mid\ell}(\mathcal O_{L_{\tilde v}}^\times)^n$ with values in $\Lambda^\times$ such that $\kappa_{v,i}(\sigma)\varepsilon(\sigma)^{i-1}\theta(e_{v,i}(\operatorname{Art}^{-1}\sigma))=1$ for $\sigma\in I_v$, where $\kappa_{v,i}$ is the universal character of $I_v$ over $\Lambda$ and $e_{v,i}$ places a unit in the $(v,i)$ coordinate; $\theta$ is continuous, and for every level $\mathcal K$ at the places of $S_{\mathrm a}$ and $v_0$, every pair $(k,b)$, with $k\in\mathcal K$ and $b$ in the Iwahori subgroup at the places above $\ell$, that fixes a point of $U(n)(L^+)\backslash U(n)(\mathbf A_{L^+}^\infty)$ has $b$ central in the torus and $\theta(b)\tau_0(k)=1$;
\item elements $a_{w,j}$ and $a_{v,j}$ of $\mathcal O$, for the split places $w$ of $L$ outside $S$, $v\mid\ell$ and $1\le j\le n$, such that the Hecke polynomial at $w$ with coefficients $a_{w,j}\bmod\mathfrak m_{\mathcal O}$ is the characteristic polynomial of $\bar D$ at the inverse of a Frobenius element at $w$, and $a_{v,j}\bmod\mathfrak m_{\mathcal O}$ is $\prod_{i\le j}\bar\psi_{v,i}\bar\varepsilon^{\,i-1}$ at the image of a uniformizer under the Artin map.
\end{enumerate}
\end{lemma}
\leanref{Development/AutomorphicArithmetic/AutomorphicVarieties/ModularityArgument/ThresholdArguments/DenseThresholdCases/HdEx.lean}{UpperTheory.V.Hi.nonempty_unitaryHeckeData}
\leanref{Development/AutomorphicArithmetic/AutomorphicVarieties/ModularityArgument/ThresholdArguments/DenseThresholdCases/UnitaryHeckeData.lean}{UpperTheory.V.Hi.CaseU.UnitaryHeckeData}
\leanref{Development/AutomorphicArithmetic/AutomorphicVarieties/ModularityArgument/ThresholdArguments/DenseThresholdCases/TTypeEx.lean}{UpperTheory.V.Hi.exists_integral_type_for_residual_rep}
\leanref{Development/AutomorphicArithmetic/AutomorphicVarieties/ModularityArgument/ThresholdArguments/DenseThresholdCases/TThetaEx.lean}{UpperTheory.V.Hi.exists_theta_is_weight_of_lam_obj}
\begin{proof}
(1) If there is no family of Teichm\"uller lifts, take $m=0$. Otherwise the family $(t_i)$ is unique, as $q-1$ is prime to $\ell$, so that distinct $(q-1)$-th roots of unity in $\mathcal O$ have distinct reductions. Let $u\in S_{\mathrm a}$. By the inertial local Langlands correspondence of Schneider and Zink in the case of monodromy $0$ \cite[Theorem~3.7]{CEGGPS16}, \cite[\S6]{SZ99}, there is an irreducible smooth representation $\xi_u$ of $GL_n(\mathcal O_{L_{\tilde u}})$ such that an irreducible smooth representation $\sigma_u$ of $GL_n(L_{\tilde u})$ contains $\xi_u$ if and only if, for every $\operatorname{rc}$ that satisfies the laws, $\operatorname{rc}(\sigma_u\nu^{(1-n)/2})$ has monodromy $0$ and restriction $\bigoplus_i\iota t_i$ to inertia. This condition on $\sigma_u$ is preserved when $\sigma_u$ is transported by an automorphism of $\mathbf C$ fixing $\iota(\mathcal O)$, which fixes the values of the $\iota t_i$. As $\xi_u$ factors through a finite quotient of $GL_n(\mathcal O_{L_{\tilde u}})$, this gives a representation $\tau_u$ of $GL_n(\mathcal O_{L_{\tilde u}})$ on some $\mathcal O^{m_u}$, with open kernel, such that $\sigma_u$ contains $\iota\tau_u$, in the sense of (1) for the single place $u$, if and only if it contains $\xi_u$. Let $\tau_0=\bigotimes_u\tau_u$, on $\mathcal O^m$ with $m=\prod_um_u$; its kernel contains the product of the kernels of the $\tau_u$, so it is open. For $M$ and the $\sigma_u$ as in (1), $M$ contains $\iota\tau_0$ if and only if every $\sigma_u$ contains $\iota\tau_u$. Indeed, the equivariant maps from $M$ to $\sigma_u$ for the factor at $u$ separate the vectors of $M$, as $M$ is a sum of copies of $\sigma_u$ there; composing a vector $F$ for $\tau_0$ with one that is nonzero on a coordinate of $F$ gives a vector for $\tau_u$. Conversely, adding the places one at a time, the vectors for the tensor product over the places already added form a nonzero subspace stable under the factor at the next place $u$, which therefore contains a copy of $\sigma_u$, and the vector for $\tau_u$ in that copy extends the construction.
% (1): Hi.exists_integral_type_for_residual_rep = NonzeroTypeTheory.exists_type_for_residual_rep' = Z.Desc.P.exists_integral_type_for_residual_rep (PTypeEx.lean): by_cases on a Teichmuller family (m = 0 otherwise), te_root for its uniqueness; the local type te_loc from NonzeroTypes (nonzero_types_of_literal_automorphy, proved in the Lean in PublishedInputs/NonzeroTypes) with the integral model Hi.typeLoc (exists_integral_rep_detecting_invariant_subspace, le_span_integralValued_of_galoisStable) and the Galois stability Hi.typeConj; the product is piFamilyProdHom_bijective_characterization (private main, through isotypic_separates_and_embeds_into_invariant_subspace).
(2) By local class field theory each $\kappa_{v,i}\varepsilon^{i-1}$, inverted, is a continuous character of $\mathcal O_{L_{\tilde v}}^\times$ with values in $\Lambda^\times$; together they define $\theta$ on the torus. The levels are neat, so the only pairs $(k,b)$ fixing a point are trivial, and the last condition holds.
(3) Since the Hecke polynomial at $w$ is a monic polynomial whose coefficients determine it, lift the coefficients of the characteristic polynomial of $\bar D$ at the Frobenius element, and the residual values of the $U$-operators, to $\mathcal O$.
\end{proof}

\begin{lemma}\label{sha-hecke-algebra}
Assume that Galois representations exist, in the sense of Definition~\ref{gal-galois-representations}, relative to the maps satisfying the laws of Definition~\ref{gal-llc-laws}. In Situation~\ref{sha-hecke-situation}, let $\iota$, $U(n)$ and data $\tau_0$, $\theta$ and $a=(a_{w,j},a_{v,j})$ be as in Lemma~\ref{sha-hecke-data}, and let $\mathcal K_0$ be a compact open level at the places of $S_{\mathrm a}$ and $v_0$, with the principal congruence subgroup of level $\tilde v_0$. Let $M$ be the $\Lambda$-module of Hida's ordinary $\Lambda$-adic forms on $U(n)$ of type $\tau_0$, weight $\theta$ and level $\mathcal K_0$, localized at the maximal ideal of the eigenvalues $a$: the common generalized eigenspace of the Hecke operators $T_{w,j}$ at the split places $w\notin S$ and the operators $U_{v,j}$ at $v\mid\ell$ for the eigenvalues $a_{w,j}$ and $a_{v,j}$. Suppose $M\ne0$, and let $\mathbf T$ be the $\Lambda$-algebra of endomorphisms of $M$ generated by these operators. Then $\mathbf T$ is a complete local $\mathcal O$-algebra, finite over $\Lambda$ and torsion free over $\Lambda$; $\Lambda$ is a domain; each generator is congruent to its eigenvalue in $a$ modulo the maximal ideal; the $U_{v,j}$ are units; and for every weight an element of $\mathbf T$ that vanishes at all points of distinct type of that weight is $0$.
\end{lemma}
\leanref{Development/AutomorphicArithmetic/AutomorphicVarieties/HeckeAlgebras/OrdinaryHeckeCharacters/HiAlgUnit.lean}{UpperTheory.V.Hi.hiAlg}
\leanref{Development/AutomorphicArithmetic/AutomorphicVarieties/HeckeAlgebras/OrdinaryHeckeCharacters/OrdFree.lean}{UpperTheory.V.Hi.ordFree}
\leanref{Development/AutomorphicArithmetic/AutomorphicVarieties/ModularityArgument/ThresholdArguments/DenseThresholdCases/TDenseU.lean}{UpperTheory.V.Hi.isDomain_and_denseWt_of_isTheta}
\begin{proof}
The ordinary part of the $\Lambda$-adic forms is finite and free over $\Lambda$, and the operators $U_{v,j}$ act bijectively on it; this is Hida's theorem for definite unitary groups \cite{Ger19}. The localization $M$ is a direct summand, so $\mathbf T$ is finite over $\Lambda$, hence complete and semilocal, and local because the eigenvalues are fixed modulo the maximal ideal. A nonzero $\lambda\in\Lambda$ kills no nonzero element of the free module $M$, so it kills no nonzero element of $\mathbf T\subset\operatorname{End}_\Lambda M$.
By local class field theory $\Lambda$ is a power series ring over $\mathcal O$ in variables indexed by $v\mid\ell$, $1\le i\le n$ and a basis of the free part of $\mathcal O_{L_{\tilde v}}^\times\otimes\mathbf Z_\ell$, the variables being $c\,\kappa_{v,i}(\gamma)-1$ for suitable $c\in\mathcal O^\times$ and $\gamma\in I_v$; so $\Lambda$ is a domain. A point of $\Lambda$ is determined by the values of the $\kappa_{v,i}$ at the $\gamma$, and the points of distinct type of a given weight give values $\chi(\gamma)$ for characters $\chi$ ranging over a set whose values are Zariski dense in each variable; a power series vanishing at all of them is $0$. For the last statement, let $t\in\mathbf T$ vanish at all points of distinct type of weight $h$. At such a point of $\Lambda$ the specialization of $M$ is a space of classical ordinary forms, by Hida's control theorem, on which the Hecke operators act through the eigenvalues of automorphic representations of $U(n)$, so $t$ acts on it as $0$. Hence for $m\in M$ the coordinates of $tm$ in a basis of $M$ over $\Lambda$ vanish at all these points, so they are $0$, and $t=0$.
\end{proof}

\begin{lemma}\label{sha-hecke-glue}
Let $A$ be a compact topological ring, $G$ a topological group, $n\ge0$, and $(K_x)_{x\in X}$ a family of Hausdorff topological rings with continuous homomorphisms $f_x:A\to K_x$ such that an element $a\in A$ with $f_x(a)=0$ for all $x$ is $0$. For each $x$ let $\rho_x:G\to M_n(K_x)$ be a continuous multiplicative map with $\rho_x(1)=1$, and let $\Sigma\subset G$ be a dense subset such that for every $g\in\Sigma$ there is $p\in A[X]$ whose image under each $f_x$ is the characteristic polynomial of $\rho_x(g)$. Then there is a determinant $D$ of $G$ of dimension $n$ with values in $A$ \cite{Che14}, whose characteristic polynomial coefficients at the elements of $G$ are continuous functions on $G$, such that $f_x(D)$ is the determinant of $\rho_x$ for every $x$.
\end{lemma}
\leanref{Development/AutomorphicArithmetic/AutomorphicVarieties/AlgebraTools/DeterminantRepresentations/GlueRep.lean}{UpperTheory.V.AlgebraTools.Det.glueRep}
\begin{proof}
Let $\iota=(f_x)_x:A\to A'=\prod_xK_x$ and $R=(\rho_x)_x:G\to M_n(A')$, and let $D'$ be the determinant of $R$. The map $\iota$ is continuous and injective, and $A$ is compact and $A'$ Hausdorff, so $\iota$ is a closed embedding. The coefficients $\lambda_i(g)$ of the characteristic polynomial of $R(g)$ are continuous functions of $g$, and they lie in $\iota(A)$ for $g\in\Sigma$ by hypothesis. Since $\iota(A)$ is closed and $\Sigma$ is dense, they lie in $\iota(A)$ for all $g\in G$. By Amitsur's formula, the coefficients of the values of $D'$ on generic elements $\sum_jX_jg_j$ of the group algebra lie in the subring generated by the $\lambda_i(g)$, $g\in G$ \cite{Che14}, so they lie in $\iota(A)$, and $D'$ descends to a unique determinant $D$ over $A$ with $\iota(D)=D'$. Its coefficients at the elements of $G$ are continuous because $\iota$ is a closed embedding.
% Det.glueRep: the product representation Det.piRep, then Det.glue (closed embedding; the closed set of g with lam_i in the range), Det.descendLam = Det.descend with Det.amitsurSub (Amitsur's formula, through CA.d_coef_mem).
\end{proof}

\begin{definition}\label{sha-hecke-local-ordinary}
Let $K$ be a nonarchimedean local field, with ring of integers $\mathcal O_K$, uniformizer $\varpi$ and residue field of $q$ elements, and let $n\ge1$. In $GL_n(K)$, let $N_0$ be the group of upper triangular unipotent matrices with entries in $\mathcal O_K$; for $c\ge0$ let $\bar N_c$ be the group of lower triangular unipotent matrices congruent to $1$ modulo $\varpi^{c+1}$; and let $T_0$ be the group of diagonal matrices with entries in $\mathcal O_K^\times$, with $e_i(u)\in T_0$ the matrix with $u\in\mathcal O_K^\times$ in place $i$ and $1$ in the other places of the diagonal. For $0\le j\le n$ let $\alpha_j$ be the diagonal matrix whose first $j$ entries are $\varpi$ and whose other entries are $1$, and let $z=\alpha_1\alpha_2\cdots\alpha_{n-1}=\operatorname{diag}(\varpi^{n-1},\dots,\varpi,1)$.
\begin{enumerate}
\item Let $\pi$ be a smooth representation of $GL_n(K)$ on a complex vector space $V$. For $\alpha=\operatorname{diag}(\varpi^{a_1},\dots,\varpi^{a_n})$ with $a_1\ge\dots\ge a_n$, the operator $U_\alpha$ on the vectors fixed by $N_0$ is $U_\alpha v=\sum_x\pi(x\alpha)v$, the sum over representatives $x$ of $N_0/\alpha N_0\alpha^{-1}$.
\item For $c\ge0$ and a character $\phi$ of $T_0$, $V(c,\phi)$ is the space of vectors fixed by $N_0$ and $\bar N_c$ on which every $t\in T_0$ acts by $\phi(t)^{-1}$, and its \emph{ordinary part} is $V^{\mathrm{ord}}(c,\phi)=\bigcap_{k\ge0}U_z^k\,V(c,\phi)$.
\item For smooth characters $\omega=(\omega_1,\dots,\omega_n)$ of $K^\times$, $\mathrm{PS}(\omega)$ is the normalized parabolic induction, from the Borel subgroup of upper triangular matrices, of the character $\operatorname{diag}(t_1,\dots,t_n)\mapsto\prod_i\omega_i(t_i)|t_i|^{(n-1)/2}$; and $E_j(\omega)=\prod_{i=1}^j\omega_i(\varpi)q^{1-i}$ for $0\le j\le n$.
\end{enumerate}
\end{definition}
\leanref{Development/AutomorphicArithmetic/AutomorphicVarieties/HeckeAlgebras/TcAnti.lean}{UpperTheory.V.Hi.glData}
\leanref{Development/AutomorphicArithmetic/AutomorphicVarieties/HeckeAlgebras/ContinuousTorusCharacter.lean}{UpperTheory.V.Hi.Iwa.uOp}
\leanref{Development/AutomorphicArithmetic/AutomorphicVarieties/HeckeAlgebras/ContinuousTorusCharacter.lean}{UpperTheory.V.Hi.Iwa.eig}
\leanref{Development/AutomorphicArithmetic/AutomorphicVarieties/HeckeAlgebras/Distinct.lean}{UpperTheory.V.Hi.invPart}
\leanref{Development/AutomorphicArithmetic/AutomorphicVarieties/HeckeAlgebras/Distinct.lean}{UpperTheory.V.Hi.PS}
\leanref{Development/AutomorphicArithmetic/AutomorphicVarieties/HeckeAlgebras/Distinct.lean}{UpperTheory.V.Hi.psEig}

\begin{lemma}\label{sha-hecke-local-principal}
Let $K$, $q$, $\varpi$, $n$ and the groups and operators be as in Definition~\ref{sha-hecke-local-ordinary}. Let $\pi$ be an irreducible smooth representation of $GL_n(K)$ on $V$, let $c\ge0$, and let $\phi$ be a character of $T_0$, trivial on the matrices congruent to $1$ modulo $\varpi^{c+1}$, such that for $i\ne j$ the characters $u\mapsto\phi(e_i(u))$ and $u\mapsto\phi(e_j(u))$ of $\mathcal O_K^\times$ are distinct. If $V^{\mathrm{ord}}(c,\phi)\ne0$, then there are smooth characters $\omega_1,\dots,\omega_n$ of $K^\times$ with $\omega_i(u)=\phi(e_i(u))^{-1}$ for $u\in\mathcal O_K^\times$, such that $\pi\cong\mathrm{PS}(\omega)$, $V^{\mathrm{ord}}(c,\phi)$ has dimension $1$, and $U_{\alpha_j}$ acts on $V^{\mathrm{ord}}(c,\phi)$ by $E_j(\omega)$ for $0\le j\le n$.
\end{lemma}
\leanref{Development/AutomorphicArithmetic/AutomorphicVarieties/ModularityArgument/ThresholdArguments/DenseThresholdCases/ExistsForm.lean}{UpperTheory.V.Hi.exists_principal_series_of_distinct_inv_ne_bot}
\begin{proof}
Let $N$ be the group of upper triangular unipotent matrices, $T$ the diagonal torus, $\delta$ the modulus character of the upper triangular Borel subgroup, and $p:V\to V_N$ the projection to the Jacquet module, on which $t\in T$ acts by $p(v)\mapsto\delta(t)^{-1/2}p(\pi(t)v)$.
Since $\pi$ is admissible \cite{BZ77}, $V(c,\phi)$ is finite-dimensional. The operator $U_z$ preserves $V(c,\phi)$, so $V(c,\phi)$ is the direct sum of the generalized kernel of $U_z$ and of $V^{\mathrm{ord}}(c,\phi)$.
For $v\in V(c,\phi)$ and $m\ge0$, $U_z^mv=\pi(z^m)\sum_x\pi(x)v$, the sum over $x\in z^{-m}N_0z^m/N_0$, and the groups $z^{-m}N_0z^m$ exhaust $N$; so $p(v)=0$ if and only if $U_z^mv=0$ for some $m$ \cite{BZ77}. Hence $p$ is injective on $V^{\mathrm{ord}}(c,\phi)$. Since $\delta$ is trivial on $T_0$, $p$ maps $V^{\mathrm{ord}}(c,\phi)$ into the subspace $E\subset V_N$ on which every $t\in T_0$ acts by $\phi(t)^{-1}$, and $E$ is stable under $T$.
The Jacquet module $V_N$ has finite length \cite{BZ77}. So there are a smooth character $\psi=\psi_1\otimes\dots\otimes\psi_n$ of $T$ and a nonzero $T$-equivariant map from $V_N$ to $\psi$ that does not vanish on $p(V^{\mathrm{ord}}(c,\phi))$. Comparing the actions of $T_0$ gives $\psi_i(u)=\phi(e_i(u))^{-1}$ for $u\in\mathcal O_K^\times$, and by Frobenius reciprocity $\pi$ is a subquotient of the normalized induction of $\psi$.
By the geometric lemma \cite{BZ77}, for every smooth character $\chi$ of $T$ the generalized $\chi$-eigenspace of $T$ in $V_N$ has dimension at most the number of permutations $w$ of $\{1,\dots,n\}$ with $\chi=\psi\circ w$. The $\psi_i$ are pairwise distinct on $\mathcal O_K^\times$, so the only $w$ with $\psi\circ w=\psi$ on $T_0$ is $w=1$. Every character of $T$ occurring in $E$ agrees with $\psi$ on $T_0$, hence equals $\psi$. So $E$ lies in the generalized $\psi$-eigenspace, $E$ has dimension at most $1$, $T$ acts on $E$ by $\psi$, and $V^{\mathrm{ord}}(c,\phi)$ has dimension $1$.
Put $\omega_i=\psi_i|\cdot|^{-(n-1)/2}$, so that $\omega_i(u)=\phi(e_i(u))^{-1}$. The $\omega_i$ are pairwise distinct on $\mathcal O_K^\times$, so no two of them differ by $|\cdot|^{\pm1}$; hence $\mathrm{PS}(\omega)$ is irreducible and is its only irreducible subquotient \cite{BZ77,Zel80}, and $\pi\cong\mathrm{PS}(\omega)$.
Finally, let $v\in V^{\mathrm{ord}}(c,\phi)$ and $0\le j\le n$. The operator $U_{\alpha_j}$ preserves $V(c,\phi)$ and commutes with $U_z$, so $U_{\alpha_j}v\in V^{\mathrm{ord}}(c,\phi)$. For $x\in N$ the vector $\pi(x\alpha_j)v-\pi(\alpha_j)v$ lies in the kernel of $p$, and $[N_0:\alpha_jN_0\alpha_j^{-1}]=\delta(\alpha_j)^{-1}$, so
\[
p(U_{\alpha_j}v)=\delta(\alpha_j)^{-1}p(\pi(\alpha_j)v)=\delta(\alpha_j)^{-1/2}\psi(\alpha_j)\,p(v)=E_j(\omega)\,p(v),
\]
the last equality being the computation $\delta(\alpha_j)^{-1/2}\prod_{i\le j}\omega_i(\varpi)q^{-(n-1)/2}=\prod_{i\le j}\omega_i(\varpi)q^{1-i}$. Since $p$ is injective on $V^{\mathrm{ord}}(c,\phi)$, $U_{\alpha_j}v=E_j(\omega)v$.
\end{proof}

\begin{lemma}\label{sha-hecke-local-rec}
In Situation~\ref{gal-local-situation}, let $a$ be a local Artin map (Definition~\ref{gal-local-artin}) and let $\operatorname{rc}$ satisfy the laws (Definition~\ref{gal-llc-laws}). Let $n\ge1$, let $\sigma$ be an irreducible smooth representation of $GL_n(E_w)$, and let $\omega_1,\dots,\omega_n$ be smooth characters of $E_w^\times$, with $\mathrm{PS}(\omega)$ as in Definition~\ref{sha-hecke-local-ordinary}.
\begin{enumerate}
\item $\sigma$ is a subquotient of $\mathrm{PS}(\omega)$ if and only if, for every $s\in W_A$, $\operatorname{rc}(\sigma\nu^{(1-n)/2})(s)$ has characteristic polynomial $\prod_j(X-\omega_j(a(s)))$.
\item If $\sigma\cong\mathrm{PS}(\omega)$ and the $\omega_i$ are pairwise distinct on $\mathcal O_w^\times$, then $\operatorname{rc}(\sigma\nu^{(1-n)/2})=\bigoplus_j\omega_j\circ a$; in particular it is abelian and monodromy-free (Definition~\ref{gal-wd-classes}).
\end{enumerate}
\end{lemma}
\leanref{Development/AutomorphicArithmetic/AutomorphicVarieties/ModularityArgument/ThresholdArguments/DenseThresholdCases/TPsRec.lean}{UpperTheory.V.Hi.psRec_s_ty}
\leanref{Development/AutomorphicArithmetic/AutomorphicVarieties/ModularityArgument/ThresholdArguments/DenseThresholdCases/IsSubquotientOfHasSupport.lean}{UpperTheory.V.Z.Win.T2.psRec_alt}
\leanref{Development/AutomorphicArithmetic/AutomorphicVarieties/ModularityArgument/AutomorphyLifting/UnitaryGaloisFamilies/PlainGd.lean}{UpperTheory.V.Z.Win.Gu.plainGd}
\begin{proof}
By~(4) of Definition~\ref{gal-local-artin}, two local Artin maps agree after composition with a global reciprocity map, which is injective on the image of $E_w^\times$ by Chevalley's theorem on congruence subgroups of units \cite{Che51}; so $\operatorname{rc}$ satisfies the laws with respect to $a$.
(1) By Definition~\ref{gal-irreducible-classes}(3), $\sigma$ is a subquotient of $\mathrm{PS}(\omega)$ if and only if its supercuspidal support is $\{\omega_j\nu^{(n-1)/2}\}$, that is, if and only if the supercuspidal support of $\sigma\nu^{(1-n)/2}$ is $\{\omega_1,\dots,\omega_n\}$ \cite{BZ77,Zel80}. By laws~(3) and~(5), the characteristic polynomial of $\operatorname{rc}(\pi)(s)$ is the product of those of $\operatorname{rc}(\rho)(s)$ over the supercuspidal support of $\pi$, counted with multiplicity, and $\operatorname{rc}(\chi)=\chi\circ a$ for a character $\chi$ by law~(2); this gives the implication from left to right. Conversely, suppose the characteristic polynomials are $\prod_j(X-\omega_j(a(s)))$. For each $\rho$ in the supercuspidal support of $\sigma\nu^{(1-n)/2}$, $\operatorname{rc}(\rho)$ is irreducible by law~(4), and its characteristic polynomials divide these products, so it is a character $\chi\circ a$; then $\rho=\chi$ by laws~(1), (2) and~(10). Comparing the roots of the characteristic polynomials at all $s$ identifies the multiset of these characters with $\{\omega_1,\dots,\omega_n\}$.
(2) The characters $\omega_j\nu^{(n-1)/2}$ are pairwise distinct on $\mathcal O_w^\times$, so no two of them differ by $\nu^{\pm1}$, and $\sigma$ is the Langlands quotient of the segments $(\omega_j\nu^{(n-1)/2},1)$ \cite{Zel80}. By laws~(3) and~(2), $\operatorname{rc}(\sigma)=\bigoplus_j(\omega_j|\cdot|^{(n-1)/2})\circ a$, and by law~(5), $\operatorname{rc}(\sigma\nu^{(1-n)/2})=\bigoplus_j\omega_j\circ a$, a sum of characters with monodromy $0$.
\end{proof}

\begin{lemma}\label{sha-hecke-local-global}
In Situation~\ref{gal-local-situation}, let $\ell$ be the prime under $w$, $\iota:\bar{\mathbf Q}_\ell\to\mathbf C$ a field isomorphism, $n\ge1$, $r:G_E\to GL_n(\bar{\mathbf Q}_\ell)$ a continuous homomorphism, $\sigma$ an irreducible smooth representation of $GL_n(E_w)$, and $\mathrm{HT}_\tau$ a multiset of integers for each embedding $\tau:E\to\bar{\mathbf Q}_\ell$ that extends continuously to $E_w$. Suppose that $r$ is \emph{compatible with $\sigma$ above $\ell$}: for every $\operatorname{rc}$ that satisfies the laws (Definition~\ref{gal-llc-laws}) and for which $y=\operatorname{rc}(\sigma\nu^{(1-n)/2})$ is abelian and monodromy-free, $r|_{D_A}$ is de Rham and potentially semistable, $\mathrm{HT}_\tau(r|_{D_A})=\mathrm{HT}_\tau$ for every $\tau$, and $\mathrm{WD}_\iota(r|_{D_A})=y$ (Definition~\ref{gal-p-adic-hodge}). If $\sigma\cong\mathrm{PS}(\omega)$ for smooth characters $\omega_1,\dots,\omega_n$ of $E_w^\times$ that are pairwise distinct on $\mathcal O_w^\times$, then $r|_{D_A}$ is de Rham, $\mathrm{HT}_\tau(r|_{D_A})=\mathrm{HT}_\tau$ for every $\tau$, and for every local Artin map $a$ and every $s\in W_A$ the characteristic polynomial of $\mathrm{WD}_\iota(r|_{D_A})(s)$ is $\prod_j(X-\omega_j(a(s)))$.
\end{lemma}
\leanref{Development/AutomorphicArithmetic/AutomorphicVarieties/ModularityArgument/AutomorphyLifting/GalUnitaryT.lean}{UpperTheory.V.Z.Win.CmpT}
\leanref{Development/AutomorphicArithmetic/AutomorphicVarieties/ModularityArgument/AutomorphyLifting/GalUnitaryT.lean}{UpperTheory.V.Z.Win.LocGlobT}
\leanref{Development/AutomorphicArithmetic/AutomorphicVarieties/ModularityArgument/OrdinaryJacquetModules/AtLOf.lean}{UpperTheory.V.Z.Ord.AtLT}
\leanref{Development/AutomorphicArithmetic/AutomorphicVarieties/ModularityArgument/ThresholdArguments/DenseThresholdCases/OrdAtLOf.lean}{UpperTheory.V.Z.Ord.at_l_of_literal_automorphy_n}
\begin{proof}
By Theorem~\ref{gal-llc-exists} there is a map $\operatorname{rc}$ that satisfies the laws, and by the proof of Lemma~\ref{sha-hecke-local-rec} it satisfies them with respect to $a$. By Lemma~\ref{sha-hecke-local-rec}(2), $y=\operatorname{rc}(\sigma\nu^{(1-n)/2})=\bigoplus_j\omega_j\circ a$ is abelian and monodromy-free. By the compatibility, $r|_{D_A}$ is de Rham, $\mathrm{HT}_\tau(r|_{D_A})=\mathrm{HT}_\tau$, and $\mathrm{WD}_\iota(r|_{D_A})=y$, whose characteristic polynomial at $s$ is $\prod_j(X-\omega_j(a(s)))$.
\end{proof}

\begin{lemma}\label{sha-hecke-local-characters}
In Situation~\ref{gal-local-situation}, let $\ell$ be the prime under $w$, $\iota:\bar{\mathbf Q}_\ell\to\mathbf C$ a field isomorphism, $\varepsilon$ the $\ell$-adic cyclotomic character, $a$ a local Artin map and $\varpi$ a uniformizer of $E_w$. Let $n\ge1$. For $1\le i\le n$ let $h_{i,\tau}$ be integers, indexed by the continuous embeddings $\tau:E_w\to\bar{\mathbf Q}_\ell$, let $\omega_i$ be a smooth character of $E_w^\times$, $\psi_i:W_A\to\bar{\mathbf Q}_\ell^\times$ a homomorphism and $b_i:\mathcal O_w^\times\to\bar{\mathbf Q}_\ell^\times$ a map; and for $1\le j\le n$ let $u_j\in\bar{\mathbf Q}_\ell$. Suppose that:
\begin{enumerate}
\item $\psi_i(s)\varepsilon(s)^{i-1}b_i(a(s))=1$ for $s\in W_A$ with $a(s)\in\mathcal O_w^\times$;
\item $\prod_{i\le j}\psi_i(s)\varepsilon(s)^{i-1}=u_j$ for $s\in W_A$ with $a(s)=\varpi$;
\item $\omega_i(x)\,\iota\big(b_i(x)\prod_\tau\tau(x)^{i-1-h_{i,\tau}}\big)=1$ for $x\in\mathcal O_w^\times$;
\item $\iota(u_j)=\prod_{i\le j}\iota\big(\prod_\tau\tau(\varpi)^{i-1-h_{i,\tau}}\big)\,\omega_i(\varpi)\,q_w^{1-i}$.
\end{enumerate}
Then $\psi_i(s)=\iota^{-1}\big(\omega_i(a(s))\big)\prod_\tau\tau(a(s))^{-h_{i,\tau}}$ for every $i$ and every $s\in W_A$.
\end{lemma}
\leanref{Development/AutomorphicArithmetic/AutomorphicVarieties/ModularityArgument/ThresholdArguments/DenseThresholdCases/TCoreEig.lean}{UpperTheory.V.Hi.core_eigen_of_weight_char_hecke_level}
\leanref{Development/AutomorphicArithmetic/AutomorphicVarieties/PadicHodgeTheory/EpsArt.lean}{UpperTheory.V.Loc.epsArt}
\begin{proof}
The cyclotomic character is the $\ell$-adic character of the absolute value \cite{Ser68}, so $\varepsilon(s)=q_w^{d(s)}\prod_\tau\tau(a(s))$ for $s\in W_A$. Both sides of the claimed identity are homomorphisms of $W_A$, and $W_A$ is generated by the $s$ with $a(s)\in\mathcal O_w^\times$ together with one $s_0$ with $a(s_0)=\varpi$, since $a$ is surjective and $v_w(a(s))=-d(s)$.
If $a(s)=x\in\mathcal O_w^\times$, then $d(s)=0$ and $\varepsilon(s)=\prod_\tau\tau(x)$, so (1) and (3) give $\psi_i(s)=\prod_\tau\tau(x)^{1-i}b_i(x)^{-1}=\iota^{-1}(\omega_i(x))\prod_\tau\tau(x)^{-h_{i,\tau}}$.
If $a(s)=\varpi$, then $d(s)=-1$ and $\varepsilon(s)=q_w^{-1}\prod_\tau\tau(\varpi)$. By (2) and (4), for every $j$ the product over $i\le j$ of $\psi_i(s)\varepsilon(s)^{i-1}$ equals the product over $i\le j$ of $\iota^{-1}(\omega_i(\varpi))\prod_\tau\tau(\varpi)^{i-1-h_{i,\tau}}q_w^{1-i}$, so by induction on $i$ the factors agree. Substituting $\varepsilon(s)$ gives $\psi_i(s)=\iota^{-1}(\omega_i(\varpi))\prod_\tau\tau(\varpi)^{-h_{i,\tau}}$.
\end{proof}

\begin{lemma}\label{sha-hecke-local-flag}
In Situation~\ref{gal-local-situation}, let $\ell$ be the prime under $w$, assumed odd, $\iota:\bar{\mathbf Q}_\ell\to\mathbf C$ a field isomorphism, $a$ a local Artin map, $n\ge0$ and $r:D_A\to GL_n(\bar{\mathbf Q}_\ell)$ a continuous homomorphism. For $1\le j\le n$ let $h_{j,\tau}$ be integers, indexed by the continuous embeddings $\tau:E_w\to\bar{\mathbf Q}_\ell$, with $j\mapsto h_{j,\tau}$ strictly increasing for every $\tau$; let $\omega_j$ be smooth characters of $E_w^\times$, pairwise distinct on $\mathcal O_w^\times$; and let $\psi_j:D_A\to\bar{\mathbf Q}_\ell^\times$ be continuous characters with $\psi_j(s)=\iota^{-1}(\omega_j(a(s)))\prod_\tau\tau(a(s))^{-h_{j,\tau}}$ for $s\in W_A$. Suppose that $r$ is de Rham, with $\mathrm{HT}_\tau(r)=\{h_{1,\tau},\dots,h_{n,\tau}\}$ for every $\tau$, and that for every $s\in W_A$ the characteristic polynomial of $\mathrm{WD}_\iota(r)(s)$ is $\prod_j(X-\omega_j(a(s)))$. Then $r$ is potentially crystalline, and there is $g\in GL_n(\bar{\mathbf Q}_\ell)$ such that $gr(s)g^{-1}$ is upper triangular with diagonal $\psi_1(s),\dots,\psi_n(s)$ for every $s\in D_A$.
\end{lemma}
\leanref{Development/AutomorphicArithmetic/AutomorphicVarieties/HeckeAlgebras/Distinct.lean}{UpperTheory.V.Hi.OrdFlag}
\leanref{Development/AutomorphicArithmetic/AutomorphicVarieties/ModularityArgument/ThresholdArguments/DenseThresholdCases/LocOrdFlag.lean}{UpperTheory.V.Loc.ordFlag_alt}
\begin{proof}
The proof is by induction on $n$; for $n=0$ there is nothing to prove. Let $n\ge1$.
The character $\psi_1^{-1}$ satisfies the same relation with $\omega_1^{-1}$ and $-h_{1,\tau}$. A character with such a relation is the restriction to $D_A$ of the $\ell$-adic character attached to an algebraic Hecke character, so it is de Rham, with Hodge-Tate multiset $\{-h_{1,\tau}\}$ and $\mathrm{WD}_\iota$ equal to $\omega_1^{-1}\circ a$ \cite{Ser68,Fon94}. Hence $r'=\psi_1^{-1}\otimes r$ is de Rham, with $\mathrm{HT}_\tau(r')=\{h_{j,\tau}-h_{1,\tau}\}$, all of whose elements are $\ge0$, and the characteristic polynomial of $\mathrm{WD}_\iota(r')(s)$ is $\prod_j(X-\mu_j(s))$ with $\mu_j=(\omega_j\omega_1^{-1})\circ a$.
By Berger's theorem \cite{Ber02}, $r'$ is potentially semistable; let $(\rho',N')$ be the Weil-Deligne representation of $D_{\mathrm{pst}}(r')$, transported by $\iota$. The group $\rho'(I_A)$ is finite, and the $\mu_j$ are pairwise distinct on $I_A$, because $a(I_A)=\mathcal O_w^\times$. So the eigenspaces of $I_A$ are lines; each is stable under $W_A$, which acts on it by one of the $\mu_j$, since $I_A$ is normal in $W_A$ and the $\mu_j$ are characters of $W_A$. The operator $N'$ commutes with $\rho'(I_A)$, so it preserves each line, and being nilpotent it is $0$. Hence $r'$, and with it $r$, is potentially crystalline \cite{Fon94}.
The line on which $W_A$ acts through $\mu_1=1$ gives a nonzero element of $D_{\mathrm{pst}}(r')$ fixed by $I_A$, and from it a nonzero $x\in D_{\mathrm{pst}}(r')$ fixed by $D_A$ with $\varphi(x)=x$. Since every element of $\mathrm{HT}_\tau(r')$ is $\ge0$, the filtration of $D_{\mathrm{dR}}(r')$ has $\operatorname{Fil}^0=D_{\mathrm{dR}}(r')$, and the fundamental exact sequence $0\to\mathbf Q_\ell\to B_{\mathrm{cris}}^{\varphi=1}\to B_{\mathrm{dR}}/B_{\mathrm{dR}}^+\to0$ \cite{Fon94} shows that $x$ comes from a nonzero vector fixed by $r'$. So there is $v\ne0$ with $r(s)v=\psi_1(s)v$ for $s\in D_A$.
In a basis beginning with $v$, $r$ is block upper triangular, with $\psi_1$ in the first place and a quotient $r''$ of dimension $n-1$. Since $D_{\mathrm{dR}}$ and $D_{\mathrm{pst}}$ are exact on de Rham representations \cite{Fon94}, $r''$ is de Rham, with $\mathrm{HT}_\tau(r'')=\{h_{2,\tau},\dots,h_{n,\tau}\}$, and the characteristic polynomial of $\mathrm{WD}_\iota(r'')(s)$ is $\prod_{j\ge2}(X-\omega_j(a(s)))$. By induction, $r''$ is conjugate to an upper triangular representation with diagonal $\psi_2,\dots,\psi_n$, and so $r$ is conjugate to one with diagonal $\psi_1,\dots,\psi_n$.
% ordFlag_alt: the induction is the private lemma flag (ordChar_inv; twR with UpperTheory.deRham_rational_twist_by_character and V.Loc.hecke_char_de_rham_of_local_artin_above_l; UpperTheory.exists_common_eigenline; deRhamRational_of_conj and deRham_rational_of_extension_by_character); the invariant vector and potential crystallinity are the private lemma core (Berger through hL.aut.berger; core0 with inertia_eigenvector_extends_of_distinct_characters, exists_invariant_vector_of_finite_image_averaging and exists_frobenius_fixed_period; PstLaws.wdParameterAt_monodromy_zero_iff_potCrystalline; exists_invariant_of_dpst_phi_fixed_of_nonneg_weights).
\end{proof}

\begin{lemma}\label{sha-hecke-determinant}
Assume that Galois representations exist, in the sense of Definition~\ref{gal-galois-representations}, relative to the maps satisfying the laws of Definition~\ref{gal-llc-laws}. In Situation~\ref{sha-hecke-situation}, let $M\ne0$ and $\mathbf T$ be as in Lemma~\ref{sha-hecke-algebra}, for data $\tau_0$, $\theta$ and $a$ as in Lemma~\ref{sha-hecke-data} with Teichm\"uller lifts $t_i$ of the $\bar\chi_i$. Then there are a continuous determinant $D_{\mathbf T}$ of $\Gamma_0$ of dimension $n$ with values in $\mathbf T$ and continuous characters $\psi_{\mathbf T,v,i}:D_v\to\mathbf T^\times$ such that:
\begin{enumerate}
\item the characteristic polynomial of $D_{\mathbf T}$ at the inverse of a Frobenius element at a split place $w\notin S$ is the Hecke polynomial with coefficients $T_{w,j}$;
\item $\psi_{\mathbf T,v,i}(\sigma)\varepsilon(\sigma)^{i-1}\theta(e_{v,i}(\operatorname{Art}^{-1}\sigma))=1$ for $\sigma\in I_v$, and $\prod_{i\le j}\psi_{\mathbf T,v,i}\varepsilon^{i-1}$ at the image of a uniformizer is $U_{v,j}$;
\item for every $\mathcal O$-algebra map $y:\mathbf T\to\bar{\mathbf Q}_\ell$ that is arithmetic of some weight, in the sense of Definition~\ref{sha-hecke-arcs}(1) for the characters $y\circ\psi_{\mathbf T,v,i}$, there is a continuous semisimple $\rho:\Gamma_0\to GL_n(\bar{\mathbf Q}_\ell)$ with $y(D_{\mathbf T})$ the determinant of $\rho$, polarized, of inertial type $\bigoplus_it_i$ at the places of $S_{\mathrm a}$, with the properties of the automorphic representations of $U(n)$ of type $\tau_0$; and if $y$ is of distinct type, $\rho|_{D_v}$ is triangular with diagonal $y(\psi_{\mathbf T,v,i})$ and potentially crystalline, for $v\mid\ell$.
\end{enumerate}
\end{lemma}
\leanref{Development/AutomorphicArithmetic/AutomorphicVarieties/ModularityArgument/ThresholdArguments/DenseThresholdCases/TCoreDet.lean}{UpperTheory.V.Hi.coreDet_alt}
\leanref{Development/AutomorphicArithmetic/AutomorphicVarieties/ModularityArgument/ThresholdArguments/DenseThresholdCases/OrdCoreH.lean}{UpperTheory.V.Z.Ord.core_hecke_of_automorphic_hecke_dense_weight}
\begin{proof}
At an arithmetic point $y$ of $\mathbf T$, the eigensystem is that of an automorphic representation of $U(n)$ of type $\tau_0$ by Hida's control theorem, whose base change to $GL_n(\mathbf A_L)$ \cite{Lab11} has a Galois representation by Definition~\ref{gal-galois-representations}, compatible with the local Langlands correspondence at the split places; this gives $\rho$ in (3), with the inertial type at $S_{\mathrm a}$ from Lemma~\ref{sha-hecke-data}(1). The flags at $\ell$ at the points of distinct type are proved in the last paragraph.
Fix a weight. Apply Lemma~\ref{sha-hecke-glue} with $A=\mathbf T$, which is compact since it is a complete local noetherian ring with finite residue field; with $X$ the set of points $y$ of distinct type of that weight, $K_y=\bar{\mathbf Q}_\ell$ and $f_y=y$, which is continuous; with $\rho_y$ the representation attached to $y$ above; and with $\Sigma$ the set of inverses of the Frobenius elements at the split places $w\notin S$, which is dense in $\Gamma_0$, and $p$ the Hecke polynomial with coefficients $T_{w,j}$, whose image under $y$ is the characteristic polynomial of $\rho_y$ at the Frobenius element by the compatibility at the split places. The points of distinct type of the weight separate the elements of $\mathbf T$ (Lemma~\ref{sha-hecke-algebra}). This gives $D_{\mathbf T}$ with (1). The $U_{v,j}$ are units, so the ratios $U_{v,j}U_{v,j-1}^{-1}\varepsilon^{1-j}$ at the uniformizer and the characters of (2) on inertia define the $\psi_{\mathbf T,v,j}$ by local class field theory.

\emph{The flags at $\ell$.} Let $y$ be of distinct type and of weight $h$, let $\rho$ be as in the first paragraph, and let $v\mid\ell$. Write $K=L_{\tilde v}$, with uniformizer $\varpi$ and residue field of $q$ elements, $W_v\subset D_v$ for the Weil group and $a$ for a local Artin map (Definition~\ref{gal-local-artin}); identify $U(n)(L^+_v)$ with $GL_n(K)$ through $\tilde v$, and use the notation of Definition~\ref{sha-hecke-local-ordinary} for $GL_n(K)$. For $u\in\mathcal O_K^\times$ put $\chi_{v,i}(u)=\iota\big(y(\theta(e_{v,i}(u)))\prod_\tau\tau(u)^{i-1-h_{v,i,\tau}}\big)$, the product over the embeddings $\tau:K\to\bar{\mathbf Q}_\ell$, and put $c_{v,j}=\iota\big(\prod_{i\le j}\prod_\tau\tau(\varpi)^{h_{v,i,\tau}-i+1}\big)$.
Hida's control theorem, as in the proof of Lemma~\ref{sha-hecke-algebra}, gives the following: an integer $e\ge0$ such that each $\chi_{v,i}$ is trivial on $1+\varpi^{e+1}\mathcal O_K$; the finite-dimensional space $S$ of classical forms of weight $h$, type $\tau_0$ and level $\mathcal K_0$ that are fixed by $N_0$ and $\bar N_e$ at each $v\mid\ell$, on which $e_{v,i}(u)$ acts by $\chi_{v,i}(u)^{-1}$, and which lie in the generalized eigenspaces of $U_z$ for eigenvalues $\mu$ with $|\iota^{-1}\mu|=1$; a specialization map from the ordinary part of the $\Lambda$-adic forms to $S$, whose image spans $S$ and which carries $U_{v,j}$ to $c_{v,j}^{-1}U_{\alpha_j}$; and an automorphic representation $\sigma$ of $U(n)$, whose Galois representation is $\rho$, with a nonzero $F\in\sigma\cap S$ on which $c_{v,j}^{-1}U_{\alpha_j}$ acts by $\iota y(U_{v,j})$ for every $j$.
First, $U_{\alpha_j}$ maps $\sigma\cap S$ onto itself. It maps $S$ onto $S$, because $U_{v,j}$ is bijective on the ordinary part of the $\Lambda$-adic forms (Lemma~\ref{sha-hecke-algebra}) and the specialization is compatible with it. It maps $\sigma$ into $\sigma$, being a sum of translates. Since $S$ is finite-dimensional, $U_{\alpha_j}$ is injective on $S$, hence bijective on $\sigma\cap S$. As a representation of $GL_n(K)$, $\sigma$ is a sum of copies of an irreducible smooth representation $\sigma_v$. Since $U_z=U_{\alpha_1}\cdots U_{\alpha_{n-1}}$, $U_z$ maps $\sigma\cap S$ onto itself. Every $GL_n(K)$-equivariant map $p:\sigma\to\sigma_v$ maps $\sigma\cap S$ into $V(e,\phi_v)$, where $\phi_v(\operatorname{diag}(t_1,\dots,t_n))=\prod_i\chi_{v,i}(t_i)$, and commutes with $U_z$, so it sends $F$ into the ordinary part $V^{\mathrm{ord}}(e,\phi_v)$ of $\sigma_v$.
Second, the $\chi_{v,i}$ are pairwise distinct, because $y$ is of distinct type: by (2), $\iota^{-1}\circ\chi_{v,i}\circ a$ is the character $\phi_{v,i}$ of Definition~\ref{sha-hecke-arcs}(1) for $y\circ\psi_{\mathbf T,v,i}$ on $I_v$. By Lemma~\ref{sha-hecke-local-principal}, applied to $\sigma_v$, $c=e$ and $\phi_v$, there are smooth characters $\omega_{v,1},\dots,\omega_{v,n}$ of $K^\times$ with $\omega_{v,i}(u)\chi_{v,i}(u)=1$ for $u\in\mathcal O_K^\times$ and $\sigma_v\cong\mathrm{PS}(\omega_v)$, and $U_{\alpha_j}$ acts on $V^{\mathrm{ord}}(e,\phi_v)$ by $E_j(\omega_v)$. The maps $p$ commute with $U_{\alpha_j}$ and together separate the vectors of $\sigma$, so $U_{\alpha_j}F=E_j(\omega_v)F$, and $\iota y(U_{v,j})=c_{v,j}^{-1}E_j(\omega_v)$.
Third, by (2), the restrictions of the characters $y\circ\psi_{\mathbf T,v,i}$ to $W_v$ satisfy (1) and (2) of Lemma~\ref{sha-hecke-local-characters} with $b_i(u)=y(\theta(e_{v,i}(u)))$ and $u_j=y(U_{v,j})$, and (3) and (4) of that lemma are the relations of the second step. So $y(\psi_{\mathbf T,v,i}(s))=\iota^{-1}(\omega_{v,i}(a(s)))\prod_\tau\tau(a(s))^{-h_{v,i,\tau}}$ for $s\in W_v$.
Fourth, by the first paragraph, $\rho$ is compatible above $\ell$, in the sense of Lemma~\ref{sha-hecke-local-global}, with the local component $\sigma_v$ of $\sigma$ at the split place $v$, with $\mathrm{HT}_\tau=\{h_{v,1,\tau},\dots,h_{v,n,\tau}\}$ given by the weight $h$ of $\sigma$: this is the form of (5)(b) and (3) of Definition~\ref{gal-attached} at $\tilde v$, the local-global compatibility at $\ell$ \cite[Theorem~2.1.1]{BLGGT14}, \cite{Car14} used here. The $\omega_{v,i}$ are pairwise distinct on $\mathcal O_K^\times$, so by Lemma~\ref{sha-hecke-local-global}, $\rho|_{D_v}$ is de Rham with these Hodge-Tate multisets, and the characteristic polynomial of $\mathrm{WD}_\iota(\rho|_{D_v})(s)$ is $\prod_j(X-\omega_{v,j}(a(s)))$ for $s\in W_v$.
Fifth, $h$ is increasing, so $j\mapsto h_{v,j,\tau}$ is strictly increasing, and the $y\circ\psi_{\mathbf T,v,j}$ are continuous. Since $\ell$ is odd, Lemma~\ref{sha-hecke-local-flag}, applied to $\rho|_{D_v}$ with $\omega_j=\omega_{v,j}$, $h_{j,\tau}=h_{v,j,\tau}$ and $\psi_j=y\circ\psi_{\mathbf T,v,j}$, shows that $\rho|_{D_v}$ is potentially crystalline and conjugate to an upper triangular representation with diagonal $y(\psi_{\mathbf T,v,1}),\dots,y(\psi_{\mathbf T,v,n})$.
% Gluing: Hi.coreDet_alt applies cd_glue (= AlgebraTools.Det.glueRep with cd_compact, cd_cont, cd_pol) to H.Good points, separated by eq_zero_of_vanish_at_good_points_dense_weight, with the representations Y.r and Y.frob from partAct0_alt (spec); cd_tel and cd_psi (cd_xi, Loc.charExt) give the psi.
% Flags at l: Z.Ord.core_hecke_of_automorphic_hecke_dense_weight. Hida control is spec and partAct0_alt (Y, F, e; HSet.IsSpec, OrdAt, ClAt, finPart, lowVal = c_{v,j}, opC); steps 1-2 are goodScalB_alt (private local_step: onto_inf, uOp_factor, into_inv, isoSep, exists_principal_series_of_distinct_inv_ne_bot); step 3 is core_eigen_of_weight_char_hecke_level (eig_ord, eig_loc, Loc.epsArt); step 4 is Y.atL (from Z.Desc.P.partX_alt: Z.Win.Gu.galUnitaryT_of_printsH gives LocGlobT and HtGT, then Z.Desc.P.ordinary_at_l_of_local_global_hodge_tate_weights) used through Z.Ord.useAtL' with Z.Win.exRec'_alt, Z.Win.T2.psRec_alt and Lev.ssCalc; step 5 is Loc.ordFlag_alt.
\end{proof}

\subsection{Galois representations at arithmetic points}\label{subsection-sha-arithmetic-points}

\begin{lemma}\label{sha-hecke-ideal-speh}
Assume that Galois representations exist, in the sense of Definition~\ref{gal-galois-representations}, relative to the maps satisfying the laws of Definition~\ref{gal-llc-laws}. Let $L$ be a CM field such that $L/L^+$ is unramified at every finite place, let $\ell$ be an odd prime, $\iota:\bar{\mathbf Q}_\ell\to\mathbf C$ a field isomorphism and $n\ge1$. Let $\sigma$ be an automorphic representation of the unitary group $U(n)$ of Lemma~\ref{sha-hecke-data}, and let $\rho:\operatorname{Gal}(\bar L/L)\to GL_n(\bar{\mathbf Q}_\ell)$ be continuous and semisimple and compatible with $\sigma$ at the split places: for every finite place $w$ of $L$ split over $L^+$, every local component $\tau$ of $\sigma$ at $w$, and every $\operatorname{rc}$ that satisfies the laws and for which $y=\operatorname{rc}(\tau\nu^{(1-n)/2})$ is abelian and monodromy-free, $\mathrm{WD}_\iota(\rho|_{W_w})=y$ if $w\nmid\ell$, and $\rho|_{D_w}$ is potentially semistable with $\mathrm{WD}_\iota(\rho|_{D_w})=y$ if $w\mid\ell$. Then there are continuous representations $\rho'_1,\dots,\rho'_a$ of $\operatorname{Gal}(\bar L/L)$ over $\bar{\mathbf Q}_\ell$, of dimensions $d_i\ge1$, and integers $m_1,\dots,m_a\ge1$, such that
\[
\det(X-\rho(g))=\prod_{i=1}^a\prod_{0\le j<m_i}\det\big(X-\rho'_i(g)\varepsilon(g)^{-j}\big)
\]
for every $g$. If all $m_i=1$, then for all but finitely many places $w$ of $L$, any two Frobenius eigenvalues $z,z'$ of $\rho$ at $w$ satisfy $|\iota z|<\operatorname{N}w\,|\iota z'|$.
\end{lemma}
\leanref{Development/AutomorphicArithmetic/AutomorphicVarieties/ModularityArgument/ThresholdArguments/DenseThresholdCases/GuWrap.lean}{UpperTheory.V.Z.Win.Gu.endShape'_alt}
\leanref{Development/AutomorphicArithmetic/AutomorphicVarieties/ModularityArgument/AutomorphyLifting/UnitaryGaloisFamilies/MemOfProdEq.lean}{UpperTheory.V.Z.Win.Gu.endShape_of}
\begin{proof}
By base change from $U(n)$ to $GL_n(\mathbf A_L)$ (Theorem~\ref{gal-base-change-definite}) and the description of the discrete spectrum of $GL_N$ by Moeglin and Waldspurger \cite{MW89}, there are cuspidal automorphic representations $P_i$ of $GL_{d_i}(\mathbf A_L)$ with $P_i^\vee\cong P_i^c$, and integers $m_i\ge1$ with $\sum_im_id_i=n$, such that at every split place $w$ the local component of $\sigma$ is the product of the local Speh representations attached to $P_{i,w}$ and $m_i$.
The weight of $\sigma$ makes the infinitesimal characters of the $P_i$, shifted by a half-integer $s_i$ that depends only on $n$ and $m_i$, regular and integral, with $s_i=(n-1)/2$ when $m_i=1$. So by Definition~\ref{gal-galois-representations}, applied to the twist of $P_i$ by $|\det|^{-s_i}$, there are continuous representations $\rho'_i$ of $\operatorname{Gal}(\bar L/L)$ whose Frobenius eigenvalues at the places $w\nmid\ell$ where $P_i$ is unramified are, through $\iota$, the Satake parameters of $P_{i,w}$ multiplied by $\operatorname{N}w^{s_i}$.
Let $w$ be a place of $L$ of degree one over $\mathbf Q$, not above $\ell$ and outside the finite set where some $P_i$ is ramified. Since $L/L^+$ is unramified, $w$ is split over $L^+$. By the compatibility of $\rho$ with $\sigma$ at $w$ and the laws of Definition~\ref{gal-llc-laws} for the product of Speh representations, the Frobenius eigenvalues of $\rho$ at $w$ are the union, over $i$ and $0\le j<m_i$, of those of $\rho'_i\varepsilon^{-j}$. The Frobenius elements at these places are dense in $\operatorname{Gal}(\bar L/L)$ (Theorem~\ref{lit-chebotarev-degree-one}), and the coefficients of characteristic polynomials are continuous, so the characteristic polynomials agree at every $g$.
Suppose all $m_i=1$. Then the characteristic polynomial of $\rho(g)$ is the product of those of the $\rho'_i(g)$. Each $P_i$ is conjugate self-dual, hence has unitary central character, so by the Jacquet-Shalika bound (Theorem~\ref{lit-jacquet-shalika}(2)) its Satake parameters $\alpha$ at $w$ satisfy $\operatorname{N}w^{-1/2}<|\alpha|<\operatorname{N}w^{1/2}$. Since $s_i=(n-1)/2$, every Frobenius eigenvalue $z$ of $\rho$ at almost every $w$ satisfies $\operatorname{N}w^{(n-2)/2}<|\iota z|<\operatorname{N}w^{n/2}$, so $|\iota z|<\operatorname{N}w\,|\iota z'|$ for any two of them.
% endShape_of: BcSpeh (spehOf_alt), speh_parameter_strings (hint), HalfG (half_alt, from RowT = rowOf'), charpoly_eq_prod (ChebotarevDegreeOne) with satake_galois_link_of_printed_hida, recTotT_of_printsH, frobTotT_alt and frobTwT; the bound is GalN.satake_window (JacquetShalikaLocal.jacquetShalikaBound with CsdUnitary.isUnitary_of_isCsd).
\end{proof}

\begin{lemma}\label{sha-hecke-ideal-inertia}
In Situation~\ref{sha-hecke-situation}, let $x$ be a point, let $\rho:\operatorname{Gal}(\bar L/L)\to GL_n(\bar{\mathbf Q}_\ell)$ be a continuous semisimple representation with $x(D^{\mathrm u})$ the determinant of $\rho$, let $u\in S_{\mathrm a}$, let $I_{\tilde u}\subset\operatorname{Gal}(\bar L/L)$ be an inertia group at a place above $\tilde u$, and let $t_1,\dots,t_n:\Gamma_0\to\mathcal O^\times$ be Teichm\"uller lifts of $\bar\chi_1,\dots,\bar\chi_n$: $t_i\bmod\mathfrak m_{\mathcal O}=\bar\chi_i$ and $t_i^{q-1}=1$, where $q$ is the number of elements of $k$.
\begin{enumerate}
\item Each $t_i$ is trivial on $\ker\bar\chi_i$, and there is $P\in GL_n(\bar{\mathbf Q}_\ell)$ with $P\rho(g)P^{-1}=\operatorname{diag}(t_1(g),\dots,t_n(g))$ for every $g\in I_{\tilde u}$.
\item Let $\iota:\bar{\mathbf Q}_\ell\to\mathbf C$ be a field isomorphism, $\sigma_u$ an irreducible smooth representation of $GL_n(L_{\tilde u})$ and $\operatorname{rc}$ a map that satisfies the laws (Definition~\ref{gal-llc-laws}), with $\mathrm{WD}_\iota(\rho|_{W_{\tilde u}})=\operatorname{rc}(\sigma_u\nu^{(1-n)/2})$. Then $\operatorname{rc}(\sigma_u\nu^{(1-n)/2})$ has monodromy $0$, and its restriction to inertia is $\bigoplus_i\iota t_i$.
\end{enumerate}
\end{lemma}
\leanref{Development/AutomorphicArithmetic/AutomorphicVarieties/HeckeAlgebras/UpOperatorEigenvalues/WdCharpoly.lean}{UpperTheory.V.Hi.Up.diag_pt}
\leanref{Development/AutomorphicArithmetic/AutomorphicVarieties/HeckeAlgebras/UpOperatorEigenvalues/DiagOfCong.lean}{UpperTheory.V.Hi.Up.diag_of_cong}
\leanref{Development/AutomorphicArithmetic/AutomorphicVarieties/HeckeAlgebras/UpOperatorEigenvalues/WdCharpoly.lean}{UpperTheory.V.Hi.Up.wd_type}
\begin{proof}
(1) The integer $q-1$ is prime to $\ell$, so distinct $(q-1)$-th roots of unity in $\bar{\mathbf Q}_\ell$ remain distinct modulo the maximal ideal of its ring of integers. For $g\in\ker\bar\chi_i$, $t_i(g)$ is a $(q-1)$-th root of unity congruent to $1$, so $t_i(g)=1$.
Let $I'\subset I_{\tilde u}$ be the subgroup of the $g$ with $\bar\chi_i(g)=1$ for all $i$, so that $t_i(g)=1$ for all $i$. For $g\in I'$ and every $h$, relation (3) of Situation~\ref{sha-hecke-situation}, carried to $\rho$ by $x$, gives $\operatorname{tr}\rho(h)\rho(g)=\operatorname{tr}\rho(h)$. So $\rho(g)-1$ is orthogonal, for the trace form, to the algebra spanned by the image of $\rho$. This algebra is semisimple because $\rho$ is, so its trace form is nondegenerate in characteristic $0$, and $\rho(g)=1$.
Let $\delta=\operatorname{diag}(t_1,\dots,t_n)$ on $I_{\tilde u}$. Both $\rho|_{I_{\tilde u}}$ and $\delta$ are trivial on $I'$, so they factor through the finite group $I_{\tilde u}/I'$ and are semisimple.
For $g\in I_{\tilde u}$, $g^{q-1}\in I'$, so the eigenvalues of $\rho(g)$, like the $t_i(g)$, are $(q-1)$-th roots of unity. The characteristic polynomial of $\rho(g)$ is the image under $x$ of that of $D^{\mathrm u}$ at $g$, which reduces to that of $\bar D$, $\prod_i(X-\bar\chi_i(g))$; so its coefficients are congruent to those of $\prod_i(X-t_i(g))$. Two monic polynomials whose roots are $(q-1)$-th roots of unity and whose coefficients are congruent have the same roots, with multiplicities. So $\rho|_{I_{\tilde u}}$ and $\delta$ have the same characteristic polynomials, and they are conjugate by the Brauer-Nesbitt theorem \cite[(30.16)]{CR62}.
(2) After conjugation by $P$, $\rho$ is trivial on the open subgroup $I'$ of $I_{\tilde u}$. So the monodromy operator of $\mathrm{WD}_\iota(\rho|_{W_{\tilde u}})$ is $0$, and the restriction of this Weil-Deligne representation to inertia is $\iota\circ\rho|_{I_{\tilde u}}$, which is conjugate to $\bigoplus_i\iota t_i$. The claim follows from $\mathrm{WD}_\iota(\rho|_{W_{\tilde u}})=\operatorname{rc}(\sigma_u\nu^{(1-n)/2})$.
% diag_pt: diag_of_cong with N = |k| - 1 (cong_roots, Brauer-Nesbitt through SumRepresentations.brauerNesbitt_matrix and isSSRep_of_finite), the trace lemma eq_one_of_trace_translation_invariant_semisimple with B.u.I, and chi_cong. The kernel statement of (1) is in the private Up.types of UpUW.lean. (2) is wd_type (IsWDOf.of_conj, monodromy 0 by isMonodromy_zero_iff), whose conclusion is IsTypeWD.
\end{proof}

\begin{lemma}\label{sha-hecke-ideal-forms}
Assume that Galois representations exist, in the sense of Definition~\ref{gal-galois-representations}, relative to the maps satisfying the laws of Definition~\ref{gal-llc-laws}. In Situation~\ref{sha-hecke-situation}, let $\iota$, $U(n)$ and $\tau_0$ be as in Lemma~\ref{sha-hecke-data}, let $\mathcal K_0$ be the level of Lemma~\ref{sha-hecke-algebra}, let $h$ be an increasing weight, and let $\sigma$ be an automorphic representation of $U(n)$ such that:
\begin{enumerate}
\item $\sigma$ has weight $h$ at infinity;
\item $\sigma$ is unramified at the places of $L$ above the places of $L^+$ outside $S$;
\item the local component of $\sigma$ at $\tilde v_0$ has nonzero vectors fixed by the principal congruence subgroup of level $\tilde v_0$;
\item for every family $(t_i)$ of Teichm\"uller lifts of the $\bar\chi_i$, every $u\in S_{\mathrm a}$ and every map $\operatorname{rc}$ that satisfies the laws (Definition~\ref{gal-llc-laws}), the local component $\sigma_u$ of $\sigma$ at $\tilde u$ has $\operatorname{rc}(\sigma_u\nu^{(1-n)/2})$ with monodromy $0$ and restriction $\bigoplus_i\iota t_i$ to inertia.
\end{enumerate}
Then the space of classical forms of weight $h$, type $\tau_0$ and level $\mathcal K_0$ whose associated automorphic forms lie in $\sigma$ is nonzero.
\end{lemma}
\leanref{Development/AutomorphicArithmetic/AutomorphicVarieties/ModularityArgument/ThresholdArguments/DenseThresholdCases/AbsoluteIrreducibleRankFour.lean}{UpperTheory.V.Hi.nonempty_unitary_hecke_parts}
\begin{proof}
By (4) and Lemma~\ref{sha-hecke-data}(1), the representation $\bigotimes_u\sigma_u$ of $\prod_uGL_n(L_{\tilde u})$ contains $\iota\tau_0$. By (2) and (3), $\sigma$ has a nonzero vector fixed by the factors of $\mathcal K_0$ away from $S_{\mathrm a}$. Together with the copy of $\iota\tau_0$ this gives a nonzero vector of $\sigma^m$ fixed by those factors and transforming by $\iota\tau_0$ under $K_{\mathrm a}$. By (1), pairing it with the representation of weight $h$ of the compact group $U(n)(L^+\otimes\mathbf R)$ gives a nonzero classical form of weight $h$, type $\tau_0$ and level $\mathcal K_0$. This is the description of the algebraic modular forms on a definite group by automorphic representations \cite{Gro99}.
% nonempty_unitary_hecke_parts: exists_nonzero_form_fixed_by_level_of_unramified (the vector fixed away from S_a), ne_type with the TypeForAlt of R (the copy of tau_0), exists_tensor_dual_embedding_of_hasWeightAt (the weight), formU (the passage to the forms of hset). The conclusion is YOf ... sigma /= bot.
\end{proof}

\begin{lemma}\label{sha-hecke-ideal-unitary}
Assume that Galois representations exist, in the sense of Definition~\ref{gal-galois-representations}, relative to the maps satisfying the laws of Definition~\ref{gal-llc-laws}. Let $L$ be a CM field such that $L/L^+$ is unramified at every finite place, let $\ell$ be an odd prime, $\iota:\bar{\mathbf Q}_\ell\to\mathbf C$ a field isomorphism and $n\ge1$, and let $\sigma$ be an automorphic representation of the unitary group $U(n)$ of Lemma~\ref{sha-hecke-data}. Then there is a continuous semisimple representation $\rho_\sigma:\operatorname{Gal}(\bar L/L)\to GL_n(\bar{\mathbf Q}_\ell)$ such that:
\begin{enumerate}
\item $\rho_\sigma$ is unramified at every finite place $w\nmid\ell$ of $L$ at which $\sigma$ is unramified;
\item for every finite place $w$ of $L$ split over $L^+$, every valuation ring $A$ of $\bar L$ over $w$, every local component $\tau$ of $\sigma$ at $w$ and every $\operatorname{rc}$ that satisfies the laws, with $y=\operatorname{rc}(\tau\nu^{(1-n)/2})$: if $w\nmid\ell$, then $\rho_\sigma|_{W_A}$ has a Weil-Deligne representation, that is, the data of Definition~\ref{gal-wd-galois} exist, and $\mathrm{WD}_\iota(\rho_\sigma|_{W_A})=y$ if $y$ is abelian and monodromy-free; if $w\mid\ell$ and $y$ is abelian and monodromy-free, then $\rho_\sigma|_{D_A}$ is potentially semistable and $\mathrm{WD}_\iota(\rho_\sigma|_{D_A})=y$ (Definition~\ref{gal-p-adic-hodge}). In particular $\rho_\sigma$ is compatible with $\sigma$ at the split places, in the sense of Lemma~\ref{sha-hecke-ideal-speh};
\item at every place $w\nmid\ell$ of $L$, split over $L^+$, at which the local component of $\sigma$ is spherical with Satake parameters $B$, $\rho_\sigma$ is unramified and the eigenvalues of $\rho_\sigma$ at the inverse of a Frobenius element are, through $\iota$, the numbers $\operatorname{N}w^{(n-1)/2}\alpha$ for $\alpha\in B$;
\item there are multisets $\mathrm{HT}_\tau$ of integers, for the embeddings $\tau:L\to\bar{\mathbf Q}_\ell$, such that $\mathrm{HT}_\tau=\{k-\mu_k:0\le k<n\}$ whenever $\sigma$ has weight $\mu=(\mu_0\ge\dots\ge\mu_{n-1})$ at $\iota\circ\tau$ (Definition~\ref{gal-definite-unitary}(4));
\item with the multisets of (4): if $w\mid\ell$ is a finite place of $L$ split over $L^+$ and some local component of $\sigma$ at $w$ is plain (Definition~\ref{gal-plain}), then for every valuation ring $A$ of $\bar L$ over $w$, $\rho_\sigma|_{D_A}$ is de Rham and $\mathrm{HT}_\tau(\rho_\sigma|_{D_A})=\mathrm{HT}_\tau$ for every $\tau$ that extends continuously to $L_w$.
\end{enumerate}
\end{lemma}
\leanref{Development/AutomorphicArithmetic/AutomorphicVarieties/ModularityArgument/AutomorphyLifting/GalUnitaryT.lean}{UpperTheory.V.Z.Win.GalUnitaryT}
\leanref{Development/AutomorphicArithmetic/AutomorphicVarieties/ModularityArgument/AutomorphyLifting/GalUnitaryT.lean}{UpperTheory.V.Z.Win.HasHTAt}
\leanref{Development/AutomorphicArithmetic/AutomorphicVarieties/ModularityArgument/ThresholdArguments/DenseThresholdCases/GuWrap.lean}{UpperTheory.V.Z.Win.Gu.galUnitaryT_of_printsH}
\leanref{Development/AutomorphicArithmetic/AutomorphicVarieties/ModularityArgument/ThresholdArguments/DenseThresholdCases/TSatFrobT.lean}{UpperTheory.V.Z.Desc.W.satFrobT_alt}
\begin{proof}
(1), (2), (4) and (5) are the existence of Galois representations for automorphic representations of definite unitary groups, with their local-global compatibility and their Hodge-Tate weights \cite{CH13,BLGGT14}. Alpöge builds $\rho_\sigma$ as the direct sum of the representations given by Definition~\ref{gal-galois-representations} for the twists, by powers of $|\det|$, of the cuspidal representations $P_i$ in the base change of $\sigma$ to $GL_n(\mathbf A_L)$ (Theorem~\ref{gal-base-change-definite}), one for each segment of their Speh representations, as in the proof of Lemma~\ref{sha-hecke-ideal-speh}. The multiset $\mathrm{HT}_\tau$ is the sum of the multisets that Definition~\ref{gal-attached} gives for these summands. Alpöge derives (4) from (3) of that definition and the infinitesimal characters of Theorem~\ref{gal-base-change-definite}(6), and (5) from (5)(b) of that definition, as $\operatorname{rc}(\tau\nu^{(1-n)/2})$ is abelian and monodromy-free when $\tau$ is plain, as in the proof of Lemma~\ref{sha-hecke-ideal-attached}(2).
(3) Let $w$ be such a place and $\operatorname{rc}$ a map that satisfies the laws. By the laws, $\operatorname{rc}(\sigma_w\nu^{(1-n)/2})$ is the unramified sum of characters whose values at a Frobenius element are the $\operatorname{N}w^{(n-1)/2}\alpha$ for $\alpha\in B$; it is abelian and monodromy-free. So by (2), $\mathrm{WD}_\iota(\rho_\sigma|_{W_w})$ equals it, which gives (3).
% GalUnitaryT is proved by Z.Win.Gu.galUnitaryT_of_printsH_of_bcSpeh_of_rT (assemble_s with sumFamG, half_alt, speh_parameter_strings, recTotT_of_printsH, frobTotT_alt and irrCaseT_alt) from HypothesesH'; GlobalTheory.BcSpeh.BcSpeh is the Speh decomposition of the base change. (2) is its field LocGlobT (CmpT: WDT away from l, IsPstAt and wdAt above l); (4) is its field on HT (from hStr, speh_parameter_strings, and hw0 of half_alt, with HT the sum of the HT0 i k), and (5) is the second part of the local-global field, HasHTAt r HT w for a PlainTor local component (hr.ht of sumFamG with plainGd). It also gives the polarization and, when rho_sigma is irreducible, a cuspidal representation; these are used through Z.Desc.P.partX_alt. (3) is satFrobT_alt (rec_algNorm_diag); the descent of rho_sigma to Gamma_0 is exists_descent_to_level_group_of_unramified.
\end{proof}

\begin{lemma}\label{sha-hecke-ideal-ordinary-vector}
Let $K$, $q$, $\varpi$, $n$ and the groups and operators be as in Definition~\ref{sha-hecke-local-ordinary}, and let $c\ge0$. Let $\omega_1,\dots,\omega_n$ be smooth characters of $K^\times$, trivial on $1+\varpi^{c+1}\mathcal O_K$, and let $\phi$ be the character of $T_0$ with $\phi(e_i(u))=\omega_i(u)^{-1}$. Let $\pi$ be a smooth representation of $GL_n(K)$ on $V$ isomorphic to $\mathrm{PS}(\omega)$. Then there is a nonzero $v\in V(c,\phi)$ with $U_{\alpha_j}v=E_j(\omega)v$ for $0\le j\le n$; in particular $v\in V^{\mathrm{ord}}(c,\phi)$.
\end{lemma}
\leanref{Development/AutomorphicArithmetic/AutomorphicVarieties/ModularityArgument/ThresholdArguments/DenseThresholdCases/OrdA.lean}{UpperTheory.V.Hi.exists_ne_zero_mem_eig_uOp_eq_psEig_smul_of_equiv_ps}
\begin{proof}
Let $\delta$ be the modulus character of the upper triangular Borel subgroup $B$, and let $\chi$ be the character of $B$, trivial on its unipotent radical $N$, with $\chi(\operatorname{diag}(t_1,\dots,t_n))=\delta^{1/2}(t)\prod_i\omega_i(t_i)|t_i|^{(n-1)/2}$. Identify $V$ with $\mathrm{PS}(\omega)$, the space of smooth $f$ with $f(bg)=\chi(b)f(g)$, and let $l(f)=f(1)$. Then $l(\pi(b)f)=\chi(b)l(f)$ for $b\in B$.
Let $I_c$ be the open subgroup generated by $N_0$, $\bar N_c$ and the diagonal matrices congruent to $1$ modulo $\varpi^{c+1}$. The character $\chi$ is trivial on $B\cap I_c$, so the function $f_0$ supported on $BI_c$ with $f_0(bk)=\chi(b)$ for $k\in I_c$ lies in $V$, is fixed by $I_c$ and has $l(f_0)=1$.
The space $W$ of vectors fixed by $I_c$ is finite-dimensional, since $\mathrm{PS}(\omega)$ is admissible \cite{BZ77}. It is stable under $T_0$ and under the operators $U_a$ for $a=\operatorname{diag}(\varpi^{a_1},\dots,\varpi^{a_n})$ with $a_1\ge\dots\ge a_n$, and all these operators commute on $W$. For $f\in W$, $l(\pi(t)f)=\chi(t)l(f)$ for $t\in T_0$, and $l(U_af)=[N_0:aN_0a^{-1}]\chi(a)l(f)$, since $l(\pi(xa)f)=\chi(xa)l(f)=\chi(a)l(f)$ for $x\in N_0$.
So $l$ is a common eigenvector of the transposes of these operators on $W$, and $l(f_0)\ne0$. The projection of $f_0$ to the common generalized eigenspace for these eigenvalues has the same value under $l$, so this generalized eigenspace is nonzero, and a minimal nonzero subspace of it stable under the operators consists of common eigenvectors. This gives a nonzero $v\in W$ with $\pi(t)v=\chi(t)v$ for $t\in T_0$ and $U_av=[N_0:aN_0a^{-1}]\chi(a)v$.
Since $v$ is fixed by $N_0$ and $\bar N_c$, and $t\in T_0$ acts by $\chi(t)=\prod_i\omega_i(t_i)=\phi(t)^{-1}$, $v\in V(c,\phi)$. For $a=\alpha_j$, $[N_0:\alpha_jN_0\alpha_j^{-1}]=q^{j(n-j)}$ and $\chi(\alpha_j)=\prod_{i\le j}\omega_i(\varpi)q^{-(n-i)}$, whose product is $\prod_{i\le j}\omega_i(\varpi)q^{1-i}=E_j(\omega)$. Since $U_z$ acts on $v$ by a nonzero scalar, $v\in V^{\mathrm{ord}}(c,\phi)$.
% The Lean: indVec (f_0), ParIndAdm.isAdmissible_parInd (finite dimension), the private ordA_abs and ordA_lin of OrdA.lean (the common eigenvector, with propCp for the value of l on U_a), and the private ordA_val (the computation of E_j).
\end{proof}

\begin{lemma}\label{sha-hecke-ideal-eigenform}
Assume that Galois representations exist, in the sense of Definition~\ref{gal-galois-representations}, relative to the maps satisfying the laws of Definition~\ref{gal-llc-laws}. In Situation~\ref{sha-hecke-situation}, let $\iota$, $U(n)$, $\tau_0$, $\theta$ and $a=(a_{w,j},a_{v,j})$ be as in Lemma~\ref{sha-hecke-data}, and let $\mathcal K_0$ and $M$ be as in Lemma~\ref{sha-hecke-algebra}. Let $x$ be a point, arithmetic of an increasing weight $h$, and let $e$, $S$, $\chi_{v,i}$ and $c_{v,j}$ be defined as in the flags at $\ell$ in the proof of Lemma~\ref{sha-hecke-determinant}, with $x$ in place of $y$. Let $\sigma$ be an automorphic representation of $U(n)$, and $Y\ne0$ the space of classical forms of weight $h$, type $\tau_0$ and level $\mathcal K_0$ in $\sigma$. Suppose that:
\begin{enumerate}
\item for every split place $w\notin S$ and $1\le j\le n$, $T_{w,j}$ acts on $Y$ by a scalar $\mu_{w,j}$, with $\iota^{-1}\mu_{w,j}-a_{w,j}$ in the maximal ideal of the ring of integers of $\bar{\mathbf Q}_\ell$;
\item for $v\mid\ell$, the local component of $\sigma$ at $v$ is $\mathrm{PS}(\omega_v)$ for smooth characters $\omega_{v,i}$ of $L_{\tilde v}^\times$ with $\omega_{v,i}\chi_{v,i}=1$ on $\mathcal O_{L_{\tilde v}}^\times$;
\item for $v\mid\ell$ and $1\le j\le n$, $\mu_{v,j}=c_{v,j}^{-1}E_j(\omega_v)$ has $\iota^{-1}\mu_{v,j}$ a unit, congruent to $a_{v,j}$ modulo the maximal ideal.
\end{enumerate}
Then $M\ne0$, and for the algebra $\mathbf T$ of Lemma~\ref{sha-hecke-algebra} there is an $\mathcal O$-algebra map $y:\mathbf T\to\bar{\mathbf Q}_\ell$, equal to $x$ on $\Lambda$, with $\iota y(T_{w,j})=\mu_{w,j}$ and $\iota y(U_{v,j})=\mu_{v,j}$.
\end{lemma}
\leanref{Development/AutomorphicArithmetic/AutomorphicVarieties/ModularityArgument/ThresholdArguments/DenseThresholdCases/TCoreUp.lean}{UpperTheory.V.Hi.coreUp_alt}
\leanref{Development/AutomorphicArithmetic/AutomorphicVarieties/HeckeAlgebras/UpOperatorEigenvalues/MemOrdC.lean}{UpperTheory.V.Hi.Up.meet}
\leanref{Development/AutomorphicArithmetic/AutomorphicVarieties/HeckeAlgebras/UpOperatorEigenvalues/EigChar.lean}{UpperTheory.V.Hi.Up.eigPt}
\begin{proof}
For $v\mid\ell$ the characters $\omega_{v,i}=\chi_{v,i}^{-1}$ are trivial on $1+\varpi_v^{e+1}\mathcal O_{L_{\tilde v}}$, since the $\chi_{v,i}$ are. Let $\phi_v$ be the character of $T_0$ at $v$ with $\phi_v(e_i(u))=\omega_{v,i}(u)^{-1}=\chi_{v,i}(u)$, and let $V_v\subset Y$ be the subspace of the forms fixed by $N_0$ and $\bar N_e$ at $v$, on which $T_0$ at $v$ acts through $\phi_v^{-1}$ and $U_{\alpha_j}$ at $v$ acts by $E_j(\omega_v)$ for $0\le j\le n$.

\emph{An ordinary vector.} There is a nonzero $F\in Y$ that lies in $V_v$ for every $v\mid\ell$. The actions of $GL_n(L_{\tilde v})$ for different places $v\mid\ell$ commute and preserve $Y$, so each $V_v$ is stable under the groups at the other places. By induction on the number of places, it suffices to show that every nonzero subspace $Z\subset Y$ stable under $GL_n(L_{\tilde v})$ meets $V_v$. As a representation of $GL_n(L_{\tilde v})$, $Y$ is a sum of copies of the local component $\mathrm{PS}(\omega_v)$, so $Z$ contains the image of an injective equivariant map from $\mathrm{PS}(\omega_v)$, and the image of the vector of Lemma~\ref{sha-hecke-ideal-ordinary-vector}, for $c=e$ and $\phi=\phi_v$, is a nonzero element of $Z\cap V_v$.

\emph{It lies in $S$.} The form $F$ has weight $h$, type $\tau_0$ and level $\mathcal K_0$; it is fixed by $N_0$ and $\bar N_e$ at each $v\mid\ell$; and $e_{v,i}(u)$ acts on it by $\phi_v(e_{v,i}(u))^{-1}=\chi_{v,i}(u)^{-1}$. The normalized operators $U_{v,j}=c_{v,j}^{-1}U_{\alpha_j}$ act on $F$ by $\mu_{v,j}$, and $U_z$, normalized by the $c_{v,j}$, acts by the product of the $\mu_{v,j}$ over $v\mid\ell$ and $j<n$, whose image under $\iota^{-1}$ is a unit by (3). So $F\in S$, and $F$ is a common eigenvector of the $T_{w,j}$ and the $U_{v,j}$ with eigenvalues $\mu_{w,j}$ and $\mu_{v,j}$.

\emph{A point of $\mathbf T$.} Let $Y_\Lambda$ be the $\Lambda$-module of the ordinary $\Lambda$-adic forms of type $\tau_0$, weight $\theta$ and level $\mathcal K_0$, inside the module of all $\Lambda$-adic forms of that type, weight and level. It is finite and free over $\Lambda$, and the operators $T_{w,j}$ and $U_{v,j}$ map it into itself and commute on it (proof of Lemma~\ref{sha-hecke-algebra}). The module $M$ is the localization of $Y_\Lambda$ at these operators for the eigenvalues $a$ (Lemma~\ref{sha-hecke-ideal-tw-augmented-idempotent}), which they map into itself, and $\mathbf T$ is the $\Lambda$-subalgebra of $\operatorname{End}_\Lambda(M)$ that they generate. By Hida's control theorem, as in the proof of Lemma~\ref{sha-hecke-algebra}, the specialization $\operatorname{sp}$ at $x$, from the $\Lambda$-adic forms to the classical forms of weight $h$, is additive on $Y_\Lambda$, satisfies $\operatorname{sp}(bf)=\iota x(b)\operatorname{sp}(f)$ for $b\in\Lambda$ and $f\in Y_\Lambda$, carries every basis of $Y_\Lambda$ over $\Lambda$ to linearly independent forms, carries the operators $T_{w,j}$ and $U_{v,j}$ on $Y_\Lambda$ to the corresponding operators on classical forms, and maps $Y_\Lambda$ onto a set that spans $S$. By Lemma~\ref{sha-automorphy-unramified-near-tangent-points}, $|x(b)|\le1$ for $b\in\Lambda$ and $|x(b)|<1$ for $b\in\mathfrak m_\Lambda$; and $|\iota^{-1}\mu_{w,j}-a_{w,j}|<1$ and $|\iota^{-1}\mu_{v,j}-a_{v,j}|<1$, by (1) and (3). So Lemma~\ref{sha-automorphy-unramified-hecke-eigenspace-localization-eigenvector}, applied to $A=\Lambda$, $Y=Y_\Lambda$, the operators $T_{w,j}$ and $U_{v,j}$ with the eigenvalues $a$, the restriction of $x$ to $\Lambda$, $\operatorname{sp}$, the operators on classical forms, the eigenvector $F\in S$ and the eigenvalues $\mu_{w,j}$ and $\mu_{v,j}$, shows that $M\ne0$ and gives a ring map $y:\mathbf T\to\bar{\mathbf Q}_\ell$, equal to $x$ on $\Lambda$, with $\iota y(T_{w,j})=\mu_{w,j}$ and $\iota y(U_{v,j})=\mu_{v,j}$. As $y$ equals $x$ on $\Lambda$, it is an $\mathcal O$-algebra map.
% A point of T is Up.eigPt (Lemma sha-automorphy-unramified-hecke-eigenspace-localization-eigenvector) with Y = H.X.Ord (H.tauA A) theta (H.lev K), op = H.op A, the fields add, smul, indep, op and span of HSet.IsSpec, hx1 and hx2 (Up.normLt and the integrality of the values on O), and the instance hypothesis Fact (CommOn (H.MK theta a K) (H.op A)) of coreUp_alt for the localization.
% coreUp_alt: up_vec (Up.meet over V + Q, with Up.isoSub, Up.ordMap, Up.fixMap and exists_ne_zero_mem_eig_uOp_eq_psEig_smul_of_equiv_ps; Up.exInv for phi_v; the Q-part is empty here), up_ord (ClAt and OrdAt through uOpList and uOpPi; F lies in the eigenspace of uzC with eigenvalue the product of lowVal^{-1} psEig), then Up.eigPt with the fields span, indep, op, add and smul of HSet.IsSpec (eigChar, then heckeModel and heckeRestr for the localization).
\end{proof}

\begin{lemma}\label{sha-hecke-ideal-attached}
Assume that Galois representations exist, in the sense of Definition~\ref{gal-galois-representations}, relative to the maps satisfying the laws of Definition~\ref{gal-llc-laws}. Let $L$ be a CM field, $\ell$ a prime, $\iota:\bar{\mathbf Q}_\ell\to\mathbf C$ a field isomorphism and $n\ge1$. Let $\pi$ be a regular algebraic cuspidal automorphic representation of $GL_n(\mathbf A_L)$ such that $(\pi,\chi)$ is polarized (Definition~\ref{gal-polarized}) for some Hecke character $\chi$ of $L^+$, and let $r:\operatorname{Gal}(\bar L/L)\to GL_n(\bar{\mathbf Q}_\ell)$ be continuous and weakly attached to $\pi$ (Definition~\ref{gal-weakly-attached}). For an embedding $\tau:L\to\bar{\mathbf Q}_\ell$ let $\mathrm{HT}_\tau=\{\frac{n-1}2-z:z\in Z_{\iota\circ\tau}(\pi)\}$. Let $w$ be a finite place of $L$, $A$ a valuation ring over $w$ and $\pi_w$ a local component of $\pi$ at $w$.
\begin{enumerate}
\item Let $\operatorname{rc}$ satisfy the laws, with $y=\operatorname{rc}(\pi_w\nu^{(1-n)/2})$ abelian and monodromy-free. If $w\nmid\ell$, then $\mathrm{WD}_\iota(r|_{W_A})=y$; if $w\mid\ell$, then $r|_{D_A}$ is potentially semistable and $\mathrm{WD}_\iota(r|_{D_A})=y$.
\item If $w\mid\ell$ and $\pi_w$ is plain (Definition~\ref{gal-plain}), then $r|_{D_A}$ is de Rham and $\mathrm{HT}_\tau(r|_{D_A})=\mathrm{HT}_\tau$ for every $\tau$ that extends continuously to $L_w$.
\end{enumerate}
\end{lemma}
\leanref{Development/AutomorphicArithmetic/AutomorphicVarieties/ModularityArgument/AutomorphyLifting/GalUnitaryT.lean}{UpperTheory.V.Z.Win.CompatT}
\leanref{Development/AutomorphicArithmetic/AutomorphicVarieties/ModularityArgument/AutomorphyLifting/PolarizedLinkage/CompatTOf.lean}{UpperTheory.V.Z.Win.compatTOf}
\begin{proof}
By Definition~\ref{gal-galois-representations} there is a triple $(R,m,\mathrm{HT})$ attached to $\pi$ relative to the maps satisfying the laws (Definition~\ref{gal-attached}); by (3) of that definition, $\mathrm{HT}$ is the family of multisets of the statement. Both $R$ and $r$ are continuous and semisimple. At all but finitely many places $w'\nmid\ell$, $\pi$ is spherical with Satake parameters $B$, and $r$ has Frobenius eigenvalues $q_{w'}^{(n-1)/2}B$ because it is weakly attached to $\pi$. So does $R$, by (2) of Definition~\ref{gal-attached} and the laws, which make $\operatorname{rc}(\pi_{w'}\nu^{(1-n)/2})$ the unramified sum of characters with these values at a Frobenius element. By the Chebotarev density theorem (Theorem~\ref{lit-chebotarev-degree-one}), $R$ and $r$ have the same characteristic polynomials everywhere, so they are conjugate by the Brauer-Nesbitt theorem \cite[(30.16)]{CR62}.
(1) For $R$, this is (2) of Definition~\ref{gal-attached} if $w\nmid\ell$, the compatibility there being an equality because $y$ is abelian and monodromy-free, and (5)(b) if $w\mid\ell$. Both statements are invariant under conjugation, so they hold for $r$.
(2) If $\pi_w$ is plain, its supercuspidal support consists of characters, no two of which differ by $\nu^{\pm1}$; so $\pi_w$ is their product, and by the laws $\operatorname{rc}(\pi_w\nu^{(1-n)/2})$ is the sum of the corresponding characters, which is abelian and monodromy-free. Then (5)(b) of Definition~\ref{gal-attached} gives the claim for $R$, and so for $r$.
% compatTOf: RowT gives R (P.cont, P.ss, P.away, P.above, P.weights); GlobalTheory.IsoOfFrobFree.main (frobEig_of, rec_eq_sat, Z.Gam.suppFacts, plainIff, sphUnl) gives R conjugate to r; WDT.of_conj, GlobalTheory.PstConj.main and GlobalTheory.ConjLocal.local_conj transfer the properties; gd_of with plain_all gives Gd for a plain local component.
\end{proof}

\begin{lemma}\label{sha-hecke-ideal-spherical}
In Situation~\ref{gal-local-situation}, let $\ell$ be a prime not under $w$, $\iota:\bar{\mathbf Q}_\ell\to\mathbf C$ a field isomorphism, $n\ge1$ and $r:G_E\to GL_n(\bar{\mathbf Q}_\ell)$ a continuous homomorphism, trivial on $I_A$. Let $\sigma$ be an irreducible smooth representation of $GL_n(E_w)$ that is plain (Definition~\ref{gal-plain}), and $\operatorname{rc}$ a map that satisfies the laws (Definition~\ref{gal-llc-laws}) with $\mathrm{WD}_\iota(r|_{W_A})=\operatorname{rc}(\sigma\nu^{(1-n)/2})$. Then $\sigma$ has nonzero vectors fixed by $GL_n(\mathcal O_w)$.
\end{lemma}
\leanref{Development/AutomorphicArithmetic/AutomorphicVarieties/ModularityArgument/ThresholdArguments/DenseThresholdCases/WStmtx.lean}{UpperTheory.V.Z.Desc.W.SphTAlt}
\leanref{Development/AutomorphicArithmetic/AutomorphicVarieties/ModularityArgument/ThresholdArguments/DenseThresholdCases/EqOneOfUnr.lean}{UpperTheory.V.Z.Desc.W.sphT_alt}
\begin{proof}
Since $r$ is trivial on $I_A$, $\operatorname{rc}(\sigma\nu^{(1-n)/2})=\mathrm{WD}_\iota(r|_{W_A})$ is unramified, and so is $\operatorname{rc}(\sigma)$ by law~(5). Let $\xi_1,\dots,\xi_n$ be the characters of the supercuspidal support of $\sigma$. No two of them differ by $\nu^{\pm1}$, so $\sigma\cong\xi_1\times\dots\times\xi_n$ \cite{Zel80}, and by laws~(2) and~(3), $\operatorname{rc}(\sigma)=\bigoplus_i\xi_i\circ a$ for the local Artin map $a$. Hence each $\xi_i$ is unramified. By the Iwasawa decomposition $GL_n(E_w)=B\,GL_n(\mathcal O_w)$, for the upper triangular Borel subgroup $B$, the function $f(bk)=\chi(b)$, with $\chi$ the character of $B$ from which $\xi_1\times\dots\times\xi_n$ is induced, is a well defined nonzero vector of $\xi_1\times\dots\times\xi_n$ fixed by $GL_n(\mathcal O_w)$.
% sphT_alt: hN.unr_wdOf_iff (unramified), RoadU.plainTor_exists_chars and full_chars (sigma is the product of its characters, in Langlands quotient form hL), the laws langlands, one and nu (hr), eq_one_of_unr and exists_unr (unramified characters), then PrinSat.hasSatGL_prinSeries and hasSatGL_of_equiv (the spherical vector).
\end{proof}

\begin{lemma}\label{sha-hecke-ideal-above}
In Situation~\ref{gal-local-situation}, let $\ell$ be the prime under $w$, assumed odd, $\iota:\bar{\mathbf Q}_\ell\to\mathbf C$ a field isomorphism, $a$ a local Artin map, $\varepsilon$ the $\ell$-adic cyclotomic character, $n\ge1$ and $r:D_A\to GL_n(\bar{\mathbf Q}_\ell)$ a continuous homomorphism. For $1\le j\le n$ let $h_{j,\tau}$ be integers, indexed by the continuous embeddings $\tau:E_w\to\bar{\mathbf Q}_\ell$, with $j\mapsto h_{j,\tau}$ strictly increasing for every $\tau$, and let $\psi_j:D_A\to\bar{\mathbf Q}_\ell^\times$ be continuous characters such that $\psi_j(s)\phi_j(s)=\prod_\tau\tau(a(s))^{-h_{j,\tau}}$ for $s\in I_A$, for a character $\phi_j$ of $I_A$ with open kernel. Suppose that:
\begin{enumerate}
\item $r$ is conjugate to an upper triangular representation with diagonal $\psi_1,\dots,\psi_n$, and $\psi_i\ne\psi_j\varepsilon$ for all $i,j$;
\item $r$ is de Rham, with Hodge-Tate multisets $\mathrm{HT}_\tau$;
\item $\sigma$ is an irreducible smooth representation of $GL_n(E_w)$ that is plain (Definition~\ref{gal-plain}), and for a map $\operatorname{rc}$ that satisfies the laws, $y=\operatorname{rc}(\sigma\nu^{(1-n)/2})$ is abelian and monodromy-free, $r$ is potentially semistable, and $\mathrm{WD}_\iota(r)=y$.
\end{enumerate}
Then there are smooth characters $\omega_1,\dots,\omega_n$ of $E_w^\times$ with $\psi_j(s)=\iota^{-1}(\omega_j(a(s)))\prod_\tau\tau(a(s))^{-h_{j,\tau}}$ for $s\in W_A$, such that $\sigma\cong\mathrm{PS}(\omega)$; and $\mathrm{HT}_\tau=\{h_{1,\tau},\dots,h_{n,\tau}\}$ for every $\tau$.
\end{lemma}
\leanref{Development/AutomorphicArithmetic/AutomorphicVarieties/PadicHodgeTheory/OrdCharIff.lean}{UpperTheory.V.Loc.ordCharIff}
\leanref{Development/AutomorphicArithmetic/AutomorphicVarieties/ModularityArgument/ThresholdArguments/DenseThresholdCases/LocFlagWD.lean}{UpperTheory.V.Loc.flag_weil_deligne_of_literal_automorphy_n}
\leanref{Development/AutomorphicArithmetic/AutomorphicVarieties/ModularityArgument/ThresholdArguments/DenseThresholdCases/TPsIrr.lean}{UpperTheory.V.Hi.prin_series_irreducible_and_unique_subquotient}
\begin{proof}
\emph{The characters $\omega_j$.} On $W_A$, $s\mapsto\psi_j(s)\prod_\tau\tau(a(s))^{h_{j,\tau}}$ is a character that factors through $a$, by local class field theory, and on $I_A$ it is $\phi_j^{-1}$, which has open kernel. So it is $\iota^{-1}\circ\omega_j\circ a$ for a smooth character $\omega_j$ of $E_w^\times$.

\emph{The Weil-Deligne representation.} As in the proof of Lemma~\ref{sha-hecke-local-flag}, each $\psi_j$ is de Rham, with Hodge-Tate multisets $\{h_{j,\tau}\}$ and Weil-Deligne representation $\omega_j\circ a$ \cite{Ser68,Fon94}. Since $D_{\mathrm{dR}}$ and $D_{\mathrm{pst}}$ are exact on de Rham representations \cite{Fon94}, the flag of (1) gives $\mathrm{HT}_\tau=\{h_{1,\tau},\dots,h_{n,\tau}\}$, and the characteristic polynomial of $\mathrm{WD}_\iota(r)(s)$ is $\prod_j(X-\omega_j(a(s)))$ for $s\in W_A$. By (3), this is the characteristic polynomial of $\operatorname{rc}(\sigma\nu^{(1-n)/2})(s)$, so $\sigma$ is a subquotient of $\mathrm{PS}(\omega)$ by Lemma~\ref{sha-hecke-local-rec}(1).

\emph{No two $\omega_j$ are linked.} Let $i\ne j$. A continuous character of the compact group $D_A$ takes unit values, while $|\tau(\varpi)|<1$ for a uniformizer $\varpi$ and every $\tau$; so if $\prod_\tau\tau(a(s))^{d_\tau}$ is the restriction to $W_A$ of a continuous character of $D_A$, then $\sum_\tau d_\tau=0$. If $\omega_i=\omega_j$, then $\psi_i\psi_j^{-1}(s)=\prod_\tau\tau(a(s))^{h_{j,\tau}-h_{i,\tau}}$ on $W_A$, and the exponents are all nonzero of the same sign, which is impossible. Suppose $\omega_i=\omega_j|\cdot|$. Since $\varepsilon(s)=q_w^{d(s)}\prod_\tau\tau(a(s))$ (proof of Lemma~\ref{sha-hecke-local-characters}) and $\iota^{-1}|a(s)|=q_w^{d(s)}$, $\psi_i(\psi_j\varepsilon)^{-1}(s)=\prod_\tau\tau(a(s))^{h_{j,\tau}-h_{i,\tau}-1}$ on $W_A$. If $i>j$ the exponents are all negative, which is impossible. If $i<j$ they are all nonnegative with sum $0$, hence $0$; then $\psi_i=\psi_j\varepsilon$ on $W_A$, hence on $D_A$ since $W_A$ is dense, contradicting (1).

\emph{Conclusion.} So no two of the $\omega_j$ differ by $|\cdot|^{\pm1}$, $\mathrm{PS}(\omega)$ is irreducible, and it is its only irreducible subquotient \cite{BZ77,Zel80}. Hence $\sigma\cong\mathrm{PS}(\omega)$.
% The private above of DenseThresholdCases/DualOfWeights.lean: Loc.ordCharIff (omega_j), Loc.flag_weil_deligne_of_literal_automorphy_n (with Loc.PstExact.flag and Loc.hecke_char_de_rham_of_local_artin_above_l) and jumps_eq (HT), Z.Desc.W.pseudo_local_langlands_descent (the subquotient, the form of Lemma sha-hecke-local-rec(1) used here), the private not_linked (sum_zero, norm_one, Loc.epsArt, weil_group_dense_in_galois_group), prin_series_irreducible_and_unique_subquotient.
\end{proof}

\begin{lemma}\label{sha-hecke-ideal-twist-local}
In Situation~\ref{gal-local-situation}, let $\ell$ be a prime, $\iota:\bar{\mathbf Q}_\ell\to\mathbf C$ a field isomorphism and $n\ge1$, and let $\pi$ and $\pi'$ be irreducible smooth representations of $GL_n(E_w)$.
\begin{enumerate}
\item Let $\chi_1,\dots,\chi_n:G_E\to\mathbf C^\times$ be homomorphisms, and let $\operatorname{rc}$ be a map that satisfies the laws (Definition~\ref{gal-llc-laws}) such that $\operatorname{rc}(\pi\nu^{(1-n)/2})$ has monodromy $0$ and restriction $\bigoplus_i\chi_i|_{I_A}$ to inertia. Then $\pi$ is good (Definition~\ref{gal-good}), and $\operatorname{rc}(\pi\nu^{(1-n)/2})$ is abelian and monodromy-free (Definition~\ref{gal-wd-classes}).
\item Let $\kappa$ be a Hecke character of $E$, with local component $\kappa_w$ at $w$, such that $\pi'\cong\pi\otimes(\kappa_w\circ\det)$. If $\pi'$ is spherical with Satake parameters $B$, then $\pi$ is good, and $\operatorname{rc}(\pi\nu^{(1-n)/2})$ is abelian and monodromy-free for every map $\operatorname{rc}$ that satisfies the laws.
\item Suppose that $E$ is totally complex and $w\nmid\ell$. Let $\kappa$ be as in (2), with $\pi'\cong\pi\otimes(\kappa_w\circ\det)$, and let $\mu:G_E\to\bar{\mathbf Q}_\ell^\times$ be a continuous character with $\mu(g)=\iota^{-1}(\kappa(x))$ whenever a reciprocity map of global class field theory sends the idele $x$ to the image of $g$ in $G_E^{\mathrm{ab}}$. Let $\pi$ be good, let $r:G_E\to GL_n(\bar{\mathbf Q}_\ell)$ be continuous, and let $\operatorname{rc}$ be a map that satisfies the laws with $\mathrm{WD}_\iota(r|_{W_A})=\operatorname{rc}(\pi\nu^{(1-n)/2})$ (Definition~\ref{gal-wd-galois}). Then $\mathrm{WD}_\iota((r\otimes\mu)|_{W_A})=\operatorname{rc}(\pi'\nu^{(1-n)/2})$.
\item Suppose that $E$ is totally complex, and let $\kappa$ and $\mu$ be as in (3), with $\pi'\cong\pi\otimes(\kappa_w\circ\det)$. Let $a$ be a local Artin map (Definition~\ref{gal-local-artin}), and for $1\le i\le n$ let $\omega_i$ be a smooth character of $E_w^\times$, $h_{i,\tau}$ integers indexed by the continuous embeddings $\tau:E_w\to\bar{\mathbf Q}_\ell$, and $\psi_i:D_A\to\bar{\mathbf Q}_\ell^\times$ a continuous character with $\psi_i(s)=\iota^{-1}(\omega_i(a(s)))\prod_\tau\tau(a(s))^{-h_{i,\tau}}$ for $s\in W_A$. If $\pi\cong\mathrm{PS}(\omega)$ (Definition~\ref{sha-hecke-local-ordinary}), then $\pi'\cong\mathrm{PS}(\omega')$ for $\omega'_i=\omega_i\kappa_w$, and $\psi_i(s)\mu(s)=\iota^{-1}(\omega'_i(a(s)))\prod_\tau\tau(a(s))^{-h_{i,\tau}}$ for $s\in W_A$.
\end{enumerate}
\end{lemma}
\leanref{Development/AutomorphicArithmetic/AutomorphicVarieties/ModularityArgument/CuspidalDescent/SphericalTwists/HodgeTateFromSandwich.lean}{UpperTheory.V.Z.Desc.W.TypeGd}
\leanref{Development/AutomorphicArithmetic/AutomorphicVarieties/ModularityArgument/CuspidalDescent/SphericalTwists/UnramifiedTwists/TypeGd.lean}{UpperTheory.V.Z.Desc.W.typeGd}
\leanref{Development/AutomorphicArithmetic/AutomorphicVarieties/ModularityArgument/ThresholdArguments/DenseThresholdCases/SphUntw.lean}{UpperTheory.V.Z.Desc.W.SphUntwAlt}
\leanref{Development/AutomorphicArithmetic/AutomorphicVarieties/ModularityArgument/ThresholdArguments/DenseThresholdCases/TSphUntw.lean}{UpperTheory.V.Z.Desc.W.sphUntw_alt}
\leanref{Development/AutomorphicArithmetic/AutomorphicVarieties/ModularityArgument/ThresholdArguments/DenseThresholdCases/SphUntw.lean}{UpperTheory.V.Z.Desc.W.TwAwayAlt}
\leanref{Development/AutomorphicArithmetic/AutomorphicVarieties/ModularityArgument/ThresholdArguments/DenseThresholdCases/TTwAway.lean}{UpperTheory.V.Z.Desc.W.twist_away_statement}
\leanref{Development/AutomorphicArithmetic/AutomorphicVarieties/ModularityArgument/CuspidalDescent/SphericalTwists/HodgeTateFromSandwich.lean}{UpperTheory.V.Z.Desc.W.OrdTw}
\leanref{Development/AutomorphicArithmetic/AutomorphicVarieties/ModularityArgument/CuspidalDescent/SphericalTwists/UnramifiedTwists/PrinTw.lean}{UpperTheory.V.Z.Desc.W.ordTw}
\begin{proof}
(1) Put $x=\pi\nu^{(1-n)/2}$; the supercuspidal support of $\pi$ is that of $x$ twisted by $\nu^{(n-1)/2}$. The representation $x$ is the Langlands quotient of a multiset of segments \cite{Zel80}; let $(\rho,l)$ be one of them. By laws~(1), (3) and~(4), $\operatorname{rc}(\rho)$ is irreducible, of dimension the degree of $\rho$, and its characteristic polynomial at every element of $W_A$ divides that of $\operatorname{rc}(x)$. Let $s\in W_A$ and $t\in I_A$. The commutator $g=sts^{-1}t^{-1}$ lies in $I_A$, and $\chi_i(g)=1$ for every $i$, so the characteristic polynomial of $\operatorname{rc}(x)(g)$, and with it that of $\operatorname{rc}(\rho)(g)$, is a power of $X-1$. Write $\operatorname{rc}(\rho)=(V,\rho_0,N)$. Since $\rho_0(g)$ is semisimple, $\rho_0(g)=1$. So $\rho_0(t)$ commutes with $\rho_0(W_A)$, and it is a scalar by Schur's lemma, $\rho_0$ being irreducible. Since $W_A$ is generated by $I_A$ and a Frobenius element, an eigenvector of $\rho_0$ of that element spans a line stable under $W_A$, and $V$ has dimension $1$. So every segment of $x$ is made of characters, the supercuspidal supports of $x$ and of $\pi$ consist of characters, and $\pi$ is good. Let $a$ be a local Artin map with respect to which $\operatorname{rc}$ satisfies the laws. By laws~(3) and~(5), the characteristic polynomial of $\operatorname{rc}(x)(s)$, for $s\in W_A$, is the product of those of $\operatorname{rc}(\xi)(s)=\xi(a(s))$ (law~(2)) over the members $\xi$ of the supercuspidal support of $x$. At a commutator of two elements of $W_A$ each factor is $X-1$, so $\operatorname{rc}(x)$ is abelian. Its monodromy is $0$ by hypothesis.

(2) Since $\pi'$ is spherical, it is the Langlands quotient of segments $(\chi_1,1),\dots,(\chi_n,1)$ for unramified characters $\chi_i$ of $E_w^\times$, and its supercuspidal support consists of the $\chi_i$ \cite{Car79,Zel80}. The supercuspidal support of $\pi\cong\pi'\otimes(\kappa_w^{-1}\circ\det)$ consists of the characters $\chi_i\kappa_w^{-1}$ \cite{BZ77}. So $\pi$ and $\pi'$ are good. Let $\operatorname{rc}$ satisfy the laws with respect to a local Artin map $a$. By laws~(2) and~(3), $\operatorname{rc}(\pi')=\bigoplus_i\chi_i\circ a$, which is abelian and monodromy-free. By law~(5), law~(7) for the good representation $\pi'$, and law~(2),
\[
\operatorname{rc}(\pi\nu^{(1-n)/2})=\operatorname{rc}(\pi')\otimes(\kappa_w^{-1}\circ a)\otimes(|\cdot|^{(1-n)/2}\circ a),
\]
which is again abelian and monodromy-free.

(3) Let $\operatorname{rc}$ satisfy the laws with respect to a local Artin map $a$. By law~(5), law~(7) for the good representation $\pi$, and law~(2),
\[
\operatorname{rc}(\pi'\nu^{(1-n)/2})=\operatorname{rc}(\pi\otimes\kappa_w)\otimes(|\cdot|^{(1-n)/2}\circ a)=\operatorname{rc}(\pi\nu^{(1-n)/2})\otimes(\kappa_w\circ a).
\]
By (4) of Definition~\ref{gal-local-artin}, a reciprocity map sends the idele equal to $a(s)$ at $w$ and to $1$ elsewhere to the image of $s\in W_A$, so $\mu(s)=\iota^{-1}(\kappa_w(a(s)))$ for $s\in W_A$. The construction of Definition~\ref{gal-wd-galois} commutes with twisting by the continuous character $\mu$, so $\mathrm{WD}_\iota((r\otimes\mu)|_{W_A})=\mathrm{WD}_\iota(r|_{W_A})\otimes(\iota\circ\mu|_{W_A})=\operatorname{rc}(\pi\nu^{(1-n)/2})\otimes(\kappa_w\circ a)$, which is $\operatorname{rc}(\pi'\nu^{(1-n)/2})$.

(4) The map $f\mapsto(\kappa_w\circ\det)f$ is an isomorphism from $\mathrm{PS}(\omega)\otimes(\kappa_w\circ\det)$ to $\mathrm{PS}(\omega')$, since the determinant of an upper triangular matrix is the product of its diagonal entries. So $\pi'\cong\mathrm{PS}(\omega')$. As in (3), $\mu(s)=\iota^{-1}(\kappa_w(a(s)))$ for $s\in W_A$, and multiplying the relation for $\psi_i$ by it gives the relation for $\psi_i\mu$.
% (1) is typeGd: Zero.LLC.abConv (Langlands classification Pr.exh, the laws langlands, cusp and dim, the private dim_one of AbConv.lean: commutators unipotent and semisimple, Schur, a Frobenius eigenvector), then Toral, good_chars, and the supp and one laws of Z.Win.toV' for the commutators. (2) is sphUntw_alt: T.supercuspidal_support_characters, hasSupport_rep_tw, good_chars; the sph law of toV' (sat_law: a spherical representation is the Langlands quotient of its unramified characters) and Gu.gd_unr; the twist and nu laws and gd_lpTwist. (3) is twist_away_statement: the nu and twist laws (hG), the one law, wdOf_twGL_det, and IsLAdicOf.loc_away with IsLocalArtin.glob. (4) is ordTw: ps_tw, and loc_away for infinity type 0.
\end{proof}

\begin{lemma}\label{sha-hecke-ideal-twist-global}
Let $L$ be a CM field with maximal totally real subfield $L^+$, $c\in\operatorname{Gal}(\bar L/L^+)$ a complex conjugation, $\ell$ a prime, $\iota:\bar{\mathbf Q}_\ell\to\mathbf C$ a field isomorphism and $n\ge1$. Let $\sigma$ be an automorphic representation of the unitary group $U(n)$ of Lemma~\ref{sha-hecke-data} (Definition~\ref{gal-definite-unitary}), and let $\beta:\operatorname{Gal}(\bar L/L)\to\bar{\mathbf Q}_\ell^\times$ be a continuous character of odd order with $\beta(cgc^{-1})=\beta(g)^{-1}$ for every $g$. Then there are a Hecke character $\kappa$ of $L$, with $\beta(g)=\iota^{-1}(\kappa(x))$ whenever a reciprocity map of global class field theory sends the idele $x$ to the image of $g$ in $\operatorname{Gal}(\bar L/L)^{\mathrm{ab}}$, and an automorphic representation $\sigma'$ of $U(n)$, such that:
\begin{enumerate}
\item for every embedding $\phi:L\to\mathbf C$ and every weight $\lambda$, if $\sigma$ has weight $\lambda$ at $\phi$, so does $\sigma'$;
\item for every finite place $w$ of $L$ such that $\beta$ is trivial on the inertia groups at $w$, if $\sigma$ is unramified at $w$, so is $\sigma'$;
\item for every place $w$ of $L$ split over $L^+$ and every local component $\tau$ of $\sigma$ at $w$, there is a local component $\tau'$ of $\sigma'$ at $w$ with $\tau'\cong\tau\otimes(\kappa_w\circ\det)$; if moreover $\beta$ is trivial on the inertia groups at $w$, then every subgroup of $GL_n(\mathcal O_w)$ that has nonzero fixed vectors in $\tau$ has nonzero fixed vectors in $\tau'$;
\item for every place $w$ of $L$ split over $L^+$ and every local component $\tau'$ of $\sigma'$ at $w$, there is a local component $\tau$ of $\sigma$ at $w$ with $\tau'\cong\tau\otimes(\kappa_w\circ\det)$.
\end{enumerate}
\end{lemma}
\leanref{Development/AutomorphicArithmetic/AutomorphicVarieties/BuildingsResidueChars/ExistsAutomorphicRepTwistByCharacter.lean}{UpperTheory.exists_automorphicRep_twist_by_character}
\leanref{Development/AutomorphicArithmetic/AutomorphicVarieties/BuildingsResidueChars/DetCharacterEqLocalDetSplitNonsplit.lean}{UpperTheory.detCharacter_eq_localDet_split_nonsplit}
\leanref{Development/AutomorphicArithmetic/AutomorphicVarieties/BuildingsResidueChars/ThetaCharEqSelfOfConjugateDual.lean}{UpperTheory.thetaChar_eq_self_of_conjugate_dual}
\leanref{Development/AutomorphicArithmetic/AutomorphicVarieties/BuildingsResidueChars/DetCharacterKerOpenTrivialRationalArchimedean.lean}{UpperTheory.detCharacter_ker_open_trivial_rational_archimedean}
\begin{proof}
Let $2N+1$ be the order of $\beta$. By class field theory \cite[Chapter~VI]{Neu99} there is a Hecke character $\nu$ of $L$ with $\nu^{2N+1}=1$ whose $\ell$-adic character through $\iota$, in the sense of the statement, is $\beta^{N+1}$. Put $\kappa=\nu^2$; its $\ell$-adic character is $\beta^{2N+2}=\beta$. At a finite place $w$, $\kappa$ is unramified if and only if $\beta$ is trivial on the inertia groups at $w$, and then $\nu$ is unramified at $w$ too.

\emph{Conjugation.} $\nu(c(x))=\nu(x)^{-1}$ for every idele $x$. Let $v$ be a finite place of $L$ at which $\nu$ and $\nu\circ c$ are unramified, and $g$ a Frobenius element at a place above $v$. Then $cgc^{-1}$ is a Frobenius element at a place above $c(v)$, and $\beta^{N+1}(cgc^{-1})=\beta^{N+1}(g)^{-1}$; so the values of $\nu$ at a uniformizer $\varpi$ of $L_v$ and at $c(\varpi)$ are inverse to each other. The Hecke character $x\mapsto\nu(c(x))\nu(x)$ is therefore trivial on the local components at all but finitely many finite places. It is continuous and trivial on $L^\times$, so it is trivial, by weak approximation.

\emph{A character of $U(n)$.} The group $U(n)(\mathbf A_{L^+})$ is a subgroup of $GL_n(\mathbf A_L)$; let $\chi(g)=\nu(\det g)$. It has open kernel, since $\nu$ is continuous of finite order; it is trivial on $U(n)(L^+)$, since $\nu$ is trivial on $L^\times$; and it is trivial on the archimedean factor, since every element of $(L\otimes\mathbf R)^\times$ is a $(2N+1)$-th power. Let $w$ be a place of $L$ split over $L^+$, with conjugate place $w^c$, and identify the factor of $U(n)$ at the place under $w$ with $GL_n(L_w)$ through the component at $w$. Since the Hermitian form is the identity, an element $g$ of that factor has component ${}^tc(g_w)^{-1}$ at $w^c$. So $\nu_{w^c}(\det g_{w^c})=\nu(c(\det g_w))^{-1}=\nu_w(\det g_w)$ by the last step, and $\chi(g)=\kappa_w(\det g_w)$. At a place $w$ not split over $L^+$, $\chi(g)=\nu_w(\det g_w)$ on the factor at $w$.

\emph{The twist.} Let $\sigma'$ be the space of the functions $\varphi\chi$ for $\varphi\in\sigma$. It is an automorphic representation of $U(n)$, since $\chi$ is continuous, trivial on $U(n)(L^+)$ and at infinity, and right translation by $g$ sends $\varphi\chi$ to $\chi(g)$ times the product of $\chi$ with the translate of $\varphi$. So a representation $\tau$ of a local factor is a local component of $\sigma$ if and only if its twist by the restriction of $\chi$ is a local component of $\sigma'$. At a split place $w$ this twist is $\tau\otimes(\kappa_w\circ\det)$, which gives the first part of (3), and (4) with $\tau=\tau'\otimes(\kappa_w^{-1}\circ\det)$. Since $\chi$ is trivial on the archimedean factor, (1) holds. If $\beta$ is trivial on the inertia groups at $w$, then $\kappa$ and $\nu$ are unramified at $w$, so $\chi$ is trivial on the elements of the factor at $w$ that lie in $GL_n(\mathcal O_w)$, and a nonzero vector of $\sigma$ fixed by these elements gives one of $\sigma'$; this is (2). In the same case $\kappa_w\circ\det$ is trivial on $GL_n(\mathcal O_w)$, so a vector of $\tau$ fixed by a subgroup of $GL_n(\mathcal O_w)$ is fixed by it in $\tau\otimes(\kappa_w\circ\det)$, which gives the rest of (3).
% The private tw_aut of DenseThresholdCases/WJTwist.lean: P.aut.art and HeckeGal.exists_unique_hchar_of_finite_order (nu for beta^(N+1)); hkappa for nu*nu; thetaChar_eq_self_of_conjugate_dual (conjugation, through frob_unr, FrobDense.eq_one_of_away and eq_one_of_unif); detCharacter_ker_open_trivial_rational_archimedean (g1, g2, g3); detCharacter_eq_localDet_split_nonsplit (g4, with SplitPts.comp_eq); exists_automorphicRep_twist_by_character (sigma' = sigma.comap (mulCh chi).symm, a2, a3); IsLAdicOf.unr_iff and locDet_KGL for the unramified cases.
\end{proof}

\begin{lemma}\label{sha-hecke-ideal-plain}
Let $E$ be a CM field, $n\ge1$ and $\pi$ a cuspidal automorphic representation of $GL_n(\mathbf A_E)$. Say that $\pi$ is \emph{plain at} a finite place $w$ if its local components at $w$ are plain (Definition~\ref{gal-plain}). For a finite place $w$, let $\varpi_w$ be a uniformizer of $E_w$.
\begin{enumerate}
\item There is $t\in\mathbf R$ such that at every finite place $w$ at which $\pi$ is plain, every character $\xi$ in the supercuspidal support of a local component of $\pi$ at $w$ satisfies $q_w^{t-1/2}<|\xi(\varpi_w)|<q_w^{t+1/2}$.
\item If $\pi$ has Satake parameters at a finite place $w$, then $\pi$ is plain at $w$.
\item Suppose $\pi^\vee\cong\pi^c$, and let $w,w'$ be finite places of $E$ above the same place of its maximal totally real subfield. If $\pi$ is plain at $w$, then it is plain at $w'$.
\item Let $\ell$ be a prime, $\iota:\bar{\mathbf Q}_\ell\to\mathbf C$ a field isomorphism, $\pi'$ a second cuspidal automorphic representation of $GL_n(\mathbf A_E)$, and $\rho:\operatorname{Gal}(\bar E/E)\to GL_n(\bar{\mathbf Q}_\ell)$ weakly attached through $\iota$ to both $\pi$ and $\pi'$ (Definition~\ref{gal-weakly-attached}). If $\pi$ is plain at $w$, so is $\pi'$.
\item Let $w$ be a finite place, $\tau$ a local component of $\pi$ at $w$, $\chi_1,\dots,\chi_n:\operatorname{Gal}(\bar E/E)\to\mathbf C^\times$ homomorphisms, and $\operatorname{rc}$ a map that satisfies the laws (Definition~\ref{gal-llc-laws}), for a valuation ring $A$ over $w$, such that $\operatorname{rc}(\tau\nu^{(1-n)/2})$ has monodromy $0$ and restriction $\bigoplus_i\chi_i|_{I_A}$ to inertia. Then $\tau$ is plain.
\end{enumerate}
\end{lemma}
\leanref{Development/AutomorphicArithmetic/AutomorphicVarieties/ModularityArgument/LeviSmoothLocalization/WinUpDown.lean}{UpperTheory.V.Z.ImaginaryQuadraticEtaCharacter.WinOf}
\leanref{Development/AutomorphicArithmetic/AutomorphicVarieties/ModularityArgument/LeviSmoothLocalization/PrincipalSeriesMatching/NormAbsPowUnif.lean}{UpperTheory.V.Z.ImaginaryQuadraticEtaCharacter.winOf}
\leanref{Development/AutomorphicArithmetic/AutomorphicVarieties/ModularityArgument/LeviSmoothLocalization/PrincipalSeriesMatching/SphUnl.lean}{UpperTheory.V.Z.ImaginaryQuadraticEtaCharacter.sphUnl}
\leanref{Development/AutomorphicArithmetic/AutomorphicVarieties/ModularityArgument/OrdinaryJacquetModules/CoreOrdAlt.lean}{UpperTheory.V.Z.Ord.ConjTor}
\leanref{Development/AutomorphicArithmetic/AutomorphicVarieties/ModularityArgument/ThresholdArguments/DenseThresholdCases/OrdConjTor.lean}{UpperTheory.V.Z.Ord.conj_torsion_of_printed_results}
\leanref{Development/AutomorphicArithmetic/Foundations/GaloisRepresentations/PlainTorsionPairs/EndoscopyRank2Ordinary.lean}{UpperTheory.Gal.projectedFormsSystemPart.tu_same}
\leanref{Development/AutomorphicArithmetic/Foundations/GaloisRepresentations/PlainTorsionPairs/PlainTorOfMono0.lean}{UpperTheory.Gal.projectedFormsSystemPart.typePlain}
\begin{proof}
(1) There are $t_0\in\mathbf R$ and a cuspidal automorphic representation $\pi_0$ with unitary central character such that $\pi\cong\pi_0\otimes|\det|^{t_0}$: twist $\pi$ by the power of $|\det|$ that makes the absolute value of its central character trivial. The local components of $\pi_0$ are unitarizable, since $\pi_0$ is a space of square-integrable cusp forms, and generic \cite{Sha74}. An irreducible unitarizable generic representation of $GL_n(E_w)$ is a product $\delta_1\nu^{s_1}\times\dots\times\delta_k\nu^{s_k}$ of square-integrable $\delta_j$ twisted by real numbers $|s_j|<1/2$ \cite{Tad86}, and its supercuspidal support is the union of the segments of the $\delta_j\nu^{s_j}$ \cite{Zel80}. Let $w$ be a place at which $\pi$ is plain. Then so is $\pi_0$, and the members of the supercuspidal support of a local component of $\pi_0$ at $w$ are characters, no two linked. A segment with two members would contain two linked members, so each segment has one member, which is then $\delta_j\nu^{s_j}$ with $\delta_j$ a square-integrable, hence unitary, character. So each member $\xi_0$ has $|\xi_0(\varpi_w)|=q_w^{-s_j}$ for some $|s_j|<1/2$, and the members $\xi=\xi_0|\cdot|^{t_0}$ for $\pi$ satisfy the bound with $t=-t_0$.

(2) The supercuspidal support of a local component of $\pi$ at $w$ consists of the unramified characters whose values at $\varpi_w$ are the Satake parameters $\alpha$. With $\pi_0$ and $t_0$ as in (1), Theorem~\ref{lit-jacquet-shalika}(2), applied to $\pi_0$, gives $q_w^{-t_0-1/2}<|\alpha|<q_w^{-t_0+1/2}$ for every $\alpha$. So no two Satake parameters have ratio $q_w$, and no two members of the support are linked.

(3) We may assume $w'\ne w$, so $w'=w^c$. Since $\pi^\vee\cong\pi^c$, every local component of $\pi$ at $w^c$ is isomorphic, through $c:E_w\to E_{w^c}$, to the contragredient of a local component at $w$. The supercuspidal support transforms in the same way, member by member, and taking contragredients and composing with $c$ sends characters to characters and preserves linking.

(4) At all but finitely many places, $\pi$ and $\pi'$ have Satake parameters, determined through $\iota$ by the Frobenius eigenvalues of $\rho$; so they have the same Satake parameters there, and $\pi'\cong\pi$ by Theorem~\ref{lit-jacquet-shalika}(1). So $\pi'$ has the local components of $\pi$.

(5) As in the proof of Lemma~\ref{sha-hecke-ideal-twist-local}(1), the supercuspidal support of $x=\tau\nu^{(1-n)/2}$ consists of characters. The local component $\tau$ of the cuspidal representation $\pi$ is generic \cite{Sha74}, and so is $x$, and $\operatorname{rc}(x)$ is monodromy-free; so the supercuspidal support of $x$ is unlinked, by Lemma~\ref{gal-generic-monodromy-free}(1). Hence $x$ is plain, and so is $\tau$, whose support is that of $x$ twisted by $\nu^{(n-1)/2}$.
% (1) is winOf: CG.exists_isUnitary_twist_absPow, ResW.isUnitarisable_of_isLocComp_cusp and CuspRep.Generic.isGeneric_of_isLocComp, suppTw, and the private exponent_of_unl (JsR.localExpBound_holds, the Lean's proof of the bound |s_j| < 1/2 for unitarizable generic representations; the segments have length one by unlinkedness). (2) is sphUnl: suppFacts, LFun.SphComp.hasSatake_of_isLocComp and ConvolutionTheory.Sph.unlinked (RgJS.unit_exponent_holds and JacquetShalikaLocal.jacquetShalikaBound). (3) is conj_torsion_of_printed_results (OneExDatum.exists_equiv_of_dualIsConj, hasSupport_of_invT, DualCusp.dualCl_isCusp, unl_of_map). (4) is tu_same, through winOf, suppSmo (strongMultOneC) and winUnl. (5) is typePlain: Zero.LLC.abConv for the characters and plainTor_of_mono0 (genericUnlinkedGL_outright, the langlands and cusp laws and SegLen.not_mono0_lpChar_lpSp).
\end{proof}

\begin{lemma}\label{sha-hecke-ideal-cuspidal}
Assume that Galois representations exist, in the sense of Definition~\ref{gal-galois-representations}, relative to the maps satisfying the laws of Definition~\ref{gal-llc-laws}. Let $L$ be a CM field such that $L/L^+$ is unramified at every finite place, let $\ell$ be an odd prime, $\iota:\bar{\mathbf Q}_\ell\to\mathbf C$ a field isomorphism and $n\ge1$, and let $\sigma$ be an automorphic representation of the unitary group $U(n)$ of Lemma~\ref{sha-hecke-data}. Let $\rho:\operatorname{Gal}(\bar L/L)\to GL_n(\bar{\mathbf Q}_\ell)$ be continuous and irreducible, and suppose that for all but finitely many finite places $w$ of $L$ split over $L^+$, and every local component of $\sigma$ at $w$ that is spherical with Satake parameters $B$, $\rho$ has Frobenius eigenvalues $\operatorname{N}w^{(n-1)/2}B$ at $w$ (Definition~\ref{gal-weakly-attached}(1)). Then there is a regular algebraic cuspidal automorphic representation $\pi$ of $GL_n(\mathbf A_L)$ with $\pi^\vee\cong\pi^c$ such that $\rho$ is weakly attached to $\pi$ through $\iota$ (Definition~\ref{gal-weakly-attached}), $\pi$ has Satake parameters at every finite place at which $\sigma$ is unramified, and at every place $w$ split over $L^+$ the local components of $\pi$ at $w$ are those of $\sigma$.
\end{lemma}
\leanref{Development/AutomorphicArithmetic/AutomorphicVarieties/ModularityArgument/ThresholdArguments/DenseThresholdCases/Link2.lean}{UpperTheory.Gal.projectedFormsSystemPart.link2_alt}
\leanref{Development/AutomorphicArithmetic/AutomorphicVarieties/ModularityArgument/ThresholdArguments/DenseThresholdCases/GuIrrCaseT.lean}{UpperTheory.V.Z.Win.Gu.irrCaseT_alt}
\begin{proof}
\emph{The representation $\rho_\sigma$.} Let $\rho_\sigma$ be the representation of Lemma~\ref{sha-hecke-ideal-unitary}. The automorphic representation $\sigma$ is unramified at all but finitely many places. Let $w$ be a place of $L$ of degree one over $\mathbf Q$, not above $\ell$, at which $\sigma$ is unramified and the hypothesis on $\rho$ holds; all but finitely many places of degree one are such. Since $L/L^+$ is unramified at $w$ and $w$ has degree one, $w$ is split over $L^+$, and the local component of $\sigma$ at $w$ is spherical, with Satake parameters $B$. By Lemma~\ref{sha-hecke-ideal-unitary}(3) and the hypothesis, $\rho_\sigma$ and $\rho$ both have Frobenius eigenvalues $\operatorname{N}w^{(n-1)/2}B$ at $w$. These Frobenius elements are dense (Theorem~\ref{lit-chebotarev-degree-one}), and $\rho_\sigma$ and $\rho$ are continuous, so they have the same characteristic polynomials. Both are semisimple, $\rho$ because it is irreducible, so they are conjugate by the Brauer-Nesbitt theorem \cite[(30.16)]{CR62}. In particular $\rho_\sigma$ is irreducible.

\emph{The cuspidal representation.} Let $P_i$, of $GL_{d_i}(\mathbf A_L)$, and $m_i$ be the cuspidal representations and multiplicities of the Speh decomposition of the base change of $\sigma$, as in the proof of Lemma~\ref{sha-hecke-ideal-speh} (Theorem~\ref{gal-base-change-definite}, \cite{MW89}). By its construction (proof of Lemma~\ref{sha-hecke-ideal-unitary}), $\rho_\sigma$ is the direct sum, over $i$ and over the $m_i$ segments of the Speh representation of $P_i$, of representations of dimension $d_i$ attached to twists of $P_i$. Since $\rho_\sigma$ is irreducible, this sum has one term: there is one $P_1=\pi$, with $m_1=1$ and $d_1=n$. So the base change of $\sigma$ is the cuspidal representation $\pi$ of $GL_n(\mathbf A_L)$, with $\pi^c\cong\pi^\vee$, and $\rho_\sigma$ is the representation that Definition~\ref{gal-galois-representations} attaches to $\pi$, so it is weakly attached to $\pi$. So is $\rho$, which is conjugate to $\rho_\sigma$. The infinitesimal characters of $\pi$ are those given by the weight of $\sigma$, which are regular and integral, so $\pi$ is regular algebraic. At a place where $\sigma$ is unramified, $\pi$ is spherical and has Satake parameters (Theorem~\ref{gal-base-change-definite}(3)).

\emph{The split places.} Let $w$ be split over $L^+$. A local component of $\pi$ at $w$ is a local component of the base change at $w$, which is a local component of $\sigma$ at $w$. Conversely, $\pi$ has a local component at $w$, and any two local components of $\sigma$ at the split place $w$ are isomorphic, so every local component of $\sigma$ at $w$ is a local component of $\pi$.
% link2_alt: r is the representation of Gu.galUnitaryT_of_printsH; hagree (ev_offL, unram_ev, isSplitAt_of_unramified_of_inertiaDeg_eq_one, SplitComp.exists_comp_and_unique, SphBC.isSpherical_of_isCompAlong_split, satake_local_of_printed_results, frobLG'_alt with the compatibility hLG) and iso_deg_one (Af.Es.charpoly_eq_prod at the degree-one places, SumRepresentations.brauerNesbitt_matrix) give r conjugate to rho; then the field hcusp of galUnitaryT_of_printsH for the irreducible r (assemble_s with hi from sumFamG: an irreducible direct sum has one block, c = 1 and d = 1; Gu.irrCaseT_alt: P 0, isRegAlg_iff with T.exists_algebraicWeight_of_automorphic_definiteUnitary, dualIsConj_of_isCsd, hun, oneFactor and LocUniq.isLocComp_unique), the transfer of the attachment along the conjugation, and at the split places local_comp_complex_of_printed_results, building_component_uniqueness and BCT.isLocComp_iff_equiv. Proposition sha-hecke-ideal applies it, through the private part2 of ProjectedFormsLinkOfPrintedHida.lean and AutUOT.lean, to rho = rho_sigma with the Frobenius eigenvalues from the compatibility.
\end{proof}

\begin{lemma}\label{sha-hecke-ideal-jacquet}
Let $K$, $q$, $\varpi$, $n$ and the groups and operators be as in Definition~\ref{sha-hecke-local-ordinary}. Let $\pi$ be an irreducible smooth representation of $GL_n(K)$ on $V$, let $c\ge0$, let $\phi$ be a character of $T_0$, trivial on the matrices congruent to $1$ modulo $\varpi^{c+1}$, and let $v\in V(c,\phi)$ be nonzero, with $U_{\alpha_j}v=a_jv$ for nonzero complex numbers $a_1,\dots,a_n$. Then there are smooth characters $\omega_1,\dots,\omega_n$ of $K^\times$ with $\omega_i(u)=\phi(e_i(u))^{-1}$ for $u\in\mathcal O_K^\times$, such that $\pi$ is a subquotient of $\mathrm{PS}(\omega)$ and $E_j(\omega)=a_j$ for $1\le j\le n$.
\end{lemma}
\leanref{Development/AutomorphicArithmetic/AutomorphicVarieties/ModularityArgument/OrdinaryJacquetModules/CoreOrdAlt.lean}{UpperTheory.V.Z.Ord.OrdJac}
\leanref{Development/AutomorphicArithmetic/AutomorphicVarieties/ModularityArgument/ThresholdArguments/DenseThresholdCases/OrdOrdJac.lean}{UpperTheory.V.Z.Ord.ordinary_jacquet_of_literal_automorphy_n}
\begin{proof}
Let $N$, $T$, $\delta$ and $p:V\to V_N$ be as in the proof of Lemma~\ref{sha-hecke-local-principal}, with $t\in T$ acting on the Jacquet module $V_N$ by $p(w)\mapsto\delta(t)^{-1/2}p(\pi(t)w)$.

\emph{The vector $v$ is ordinary.} Since $\alpha_0=1$ and $U_z=U_{\alpha_1}\cdots U_{\alpha_{n-1}}$, $U_zv=bv$ with $b=a_1\cdots a_{n-1}\ne0$. So $v=U_z^k(b^{-k}v)$ lies in $U_z^kV(c,\phi)$ for every $k\ge0$, and $v\in V^{\mathrm{ord}}(c,\phi)$. As in the proof of Lemma~\ref{sha-hecke-local-principal}, $V(c,\phi)$ is finite-dimensional and $p$ is injective on $V^{\mathrm{ord}}(c,\phi)$, so $p(v)\ne0$.

\emph{A character of $T$.} The Jacquet module $V_N$ is a smooth representation of $T$ of finite length \cite{BZ77}. Let $N_1\subset V_N$ be a $T$-stable subspace that is maximal among those not containing $p(v)$, and let $P\supset N_1$ be a proper $T$-stable subspace that is maximal among those containing $N_1$. Then $V_N/P$ is irreducible, so it is a smooth character $\mu$ of the commutative group $T$, and the projection $l:V_N\to V_N/P\cong\mathbf C$ is a nonzero form with $l(t\,x)=\mu(t)l(x)$. Let $t\in T$ and $c'\in\mathbf C$ with $t\,p(v)=c'p(v)$, and put $E=t-c'$. The subspaces $E^{-k}(N_1)$ are $T$-stable and increase with $k$, so they are equal for $k\ge k_0$. If $E^{k_0+1}V_N\not\subset N_1$, then the $T$-stable subspace $N_1+E^{k_0+1}V_N$ contains $p(v)$, say $p(v)=x+E^{k_0+1}y$ with $x\in N_1$. Applying $E$ gives $E^{k_0+2}y=-Ex\in N_1$, so $E^{k_0+1}y\in N_1$ and $p(v)\in N_1$, which is impossible. So $E^{k_0+1}$ maps $V_N$ into $N_1\subset P$, hence $(\mu(t)-c')^{k_0+1}=0$ and $\mu(t)=c'$.

\emph{The principal series.} By Frobenius reciprocity, the form $l$ makes $\pi$ a subquotient of the normalized parabolic induction of $\mu$ from the upper triangular Borel subgroup \cite{BZ77}. Write $\mu(\operatorname{diag}(t_1,\dots,t_n))=\prod_i\mu_i(t_i)$ and $\omega_i=\mu_i|\cdot|^{-(n-1)/2}$. The $\omega_i$ are smooth, and $\pi$ is a subquotient of $\mathrm{PS}(\omega)$. For $u\in\mathcal O_K^\times$, $\delta(e_i(u))=1$ and $e_i(u)\,p(v)=p(\pi(e_i(u))v)=\phi(e_i(u))^{-1}p(v)$, so $\omega_i(u)=\mu_i(u)=\phi(e_i(u))^{-1}$.

\emph{The eigenvalues.} Let $1\le j\le n$. For $x\in N$ the vector $\pi(x\alpha_j)v-\pi(\alpha_j)v$ lies in the kernel of $p$, and $[N_0:\alpha_jN_0\alpha_j^{-1}]=\delta(\alpha_j)^{-1}$, so $a_jp(v)=p(U_{\alpha_j}v)=\delta(\alpha_j)^{-1}p(\pi(\alpha_j)v)$. Hence $\alpha_j\,p(v)=\delta(\alpha_j)^{1/2}a_j\,p(v)$ and $\mu(\alpha_j)=\delta(\alpha_j)^{1/2}a_j$. As in the proof of Lemma~\ref{sha-hecke-local-principal}, $\delta(\alpha_j)^{-1/2}\mu(\alpha_j)=E_j(\omega)$, so $E_j(\omega)=a_j$.
% ordinary_jacquet_of_literal_automorphy_n: glBridge; hlist with uOpMul and hz (U_z v = b v); hvi (v in invPart); OrdB.inv_inj (p(v) /= 0); the finite length of jacB from H.loc.bz; OrdB.exists_form (through the private exists_nil of ExistsNil.lean: N_1 and P, then TorusEigenform.exists_char_of_isIrreducible); SatakeExists.isSubquotient_of_eigenform; hpsiphi with UnramFixed.deltaHalf_torusO; propCp with hidx and OrdB.num for E_j.
\end{proof}

\begin{lemma}\label{sha-hecke-ideal-ordinary-point}
Assume that Galois representations exist, in the sense of Definition~\ref{gal-galois-representations}, relative to the maps satisfying the laws of Definition~\ref{gal-llc-laws}. In Situation~\ref{sha-hecke-situation}, let $\iota$, $U(n)$ and data $\tau_0$, $\theta$ and $(a_{w,j},a_{v,j})$ be as in Lemma~\ref{sha-hecke-data}, let $M\ne0$, $\mathcal K_0$ and $\mathbf T$ be as in Lemma~\ref{sha-hecke-algebra}, and let $D_{\mathbf T}$ and $\psi_{\mathbf T,v,i}$ be as in Lemma~\ref{sha-hecke-determinant}. Let $y:\mathbf T\to\bar{\mathbf Q}_\ell$ be an $\mathcal O$-algebra map that is arithmetic of a weight $h$, in the sense of Definition~\ref{sha-hecke-arcs}(1) for the characters $y\circ\psi_{\mathbf T,v,i}$. Then there is an automorphic representation $\sigma$ of $U(n)$ such that:
\begin{enumerate}
\item the space of classical forms of weight $h$, type $\tau_0$ and level $\mathcal K_0$ in $\sigma$ is nonzero;
\item the representation $\rho_\sigma$ of Lemma~\ref{sha-hecke-ideal-unitary} factors through $\Gamma_0$, and $y(D_{\mathbf T})$ is its determinant;
\item for every $v\mid\ell$, with the Weil group $W_v\subset D_v$ and a local Artin map $a$ as in the flags at $\ell$ in the proof of Lemma~\ref{sha-hecke-determinant}, there are smooth characters $\omega_{v,1},\dots,\omega_{v,n}$ of $L_{\tilde v}^\times$ such that the local component $\sigma_v$ of $\sigma$ at $\tilde v$ is a subquotient of $\mathrm{PS}(\omega_v)$ and $y(\psi_{\mathbf T,v,i}(s))=\iota^{-1}(\omega_{v,i}(a(s)))\prod_\tau\tau(a(s))^{-h_{v,i,\tau}}$ for $s\in W_v$, the product being over the embeddings $\tau:L_{\tilde v}\to\bar{\mathbf Q}_\ell$.
\end{enumerate}
\end{lemma}
\leanref{Development/AutomorphicArithmetic/AutomorphicVarieties/ModularityArgument/ThresholdArguments/DenseThresholdCases/OrdDefs.lean}{UpperTheory.V.Z.Ord.CoreOrdAlt}
\leanref{Development/AutomorphicArithmetic/AutomorphicVarieties/ModularityArgument/ThresholdArguments/DenseThresholdCases/OrdCoreOrd.lean}{UpperTheory.V.Z.Ord.core_ordinary_of_literal_automorphy_n}
\leanref{Development/AutomorphicArithmetic/AutomorphicVarieties/ModularityArgument/ThresholdArguments/DenseThresholdCases/OrdDefs.lean}{UpperTheory.V.Z.Ord.OrdScalAlt}
\leanref{Development/AutomorphicArithmetic/AutomorphicVarieties/ModularityArgument/ThresholdArguments/DenseThresholdCases/OrdOrdScal.lean}{UpperTheory.V.Z.Ord.ordinary_scalar_of_literal_automorphy_n}
\begin{proof}
Let $e$, $S$, $\chi_{v,i}$ and $c_{v,j}$ be defined as in the flags at $\ell$ in the proof of Lemma~\ref{sha-hecke-determinant}. Hida's control theorem, in the form used there, which holds for every arithmetic $y$, gives an automorphic representation $\sigma$ of $U(n)$ whose space of classical forms of weight $h$, type $\tau_0$ and level $\mathcal K_0$ is nonzero, and a nonzero form $F\in\sigma\cap S$ that is an eigenvector of the operators $T_{w,j}$, for the split places $w\notin S$, with eigenvalues $\iota y(T_{w,j})$, and of the operators $c_{v,j}^{-1}U_{\alpha_j}$ at the places $v\mid\ell$, with eigenvalues $\iota y(U_{v,j})$. This gives (1).

(2) The forms of level $\mathcal K_0$ are fixed by the maximal compact subgroups at the places above the places of $L^+$ outside $S$, so $\sigma$ is unramified there, and by Lemma~\ref{sha-hecke-ideal-unitary}(1) $\rho_\sigma$ factors through $\Gamma_0$. At a split place $w\notin S$, Lemma~\ref{sha-hecke-ideal-unitary}(3) and the Satake isomorphism make the characteristic polynomial of $\rho_\sigma$ at the inverse of a Frobenius element, carried by $\iota$, the Hecke polynomial with coefficients $\iota y(T_{w,j})$, which is the image under $\iota y$ of that of $D_{\mathbf T}$ by Lemma~\ref{sha-hecke-determinant}(1). These inverse Frobenius elements are dense in $\Gamma_0$, and the coefficients of the characteristic polynomials of both determinants are continuous, so $y(D_{\mathbf T})$ is the determinant of $\rho_\sigma$.

(3) Let $v\mid\ell$ and $K=L_{\tilde v}$, identify $U(n)(L^+_v)$ with $GL_n(K)$ through $\tilde v$, and use the notation of Definition~\ref{sha-hecke-local-ordinary} for $GL_n(K)$. Let $\phi_v$ be the character of $T_0$ with $\phi_v(\operatorname{diag}(t_1,\dots,t_n))=\prod_i\chi_{v,i}(t_i)$; it is trivial on the matrices congruent to $1$ modulo $\varpi^{e+1}$. As a representation of $GL_n(K)$, $\sigma$ is a sum of copies of $\sigma_v$, so there is a $GL_n(K)$-equivariant map $p:\sigma\to\sigma_v$ with $p(F)\ne0$. Since $F\in S$, $p(F)\in V(e,\phi_v)$, and $U_{\alpha_j}p(F)=p(U_{\alpha_j}F)=c_{v,j}\,\iota y(U_{v,j})\,p(F)$, where $\iota y(U_{v,j})\ne0$ because $U_{v,j}$ is a unit of $\mathbf T$ (Lemma~\ref{sha-hecke-algebra}). Lemma~\ref{sha-hecke-ideal-jacquet}, applied to $\sigma_v$, $p(F)$, $c=e$, $\phi=\phi_v$ and $a_j=c_{v,j}\,\iota y(U_{v,j})$, gives smooth characters $\omega_{v,i}$ with $\omega_{v,i}(u)\chi_{v,i}(u)=1$ for $u\in\mathcal O_K^\times$, such that $\sigma_v$ is a subquotient of $\mathrm{PS}(\omega_v)$ and $\iota y(U_{v,j})=c_{v,j}^{-1}E_j(\omega_v)$. As in the third step of the flags at $\ell$ in the proof of Lemma~\ref{sha-hecke-determinant}, Lemma~\ref{sha-hecke-local-characters}, with $b_i(u)=y(\theta(e_{v,i}(u)))$ and $u_j=y(U_{v,j})$, gives the relation of (3).
% core_ordinary_of_literal_automorphy_n: Hi.spec and ordinary_pt_of_literal_automorphy_n (PartPtAlt: the part Y with F in Y and in OrdAt, an eigenvector of every opC with eigenvalue iota y(gen)); core_eigen_of_weight_char_hecke_level gives (2) (eig_det with hU.frob, eig_pol, Y.frob, Y.sph) and the relation of (3); hne (the U_{v,j} are units); ordinary_scalar_of_literal_automorphy_n is the local step: the private local_exp of OrdOrdScal.lean (char_at for phi_v, isoSep for p, uOpNat) with ordinary_jacquet_of_literal_automorphy_n. Proposition sha-hecke-ideal uses it through the private core of IsoOfDet.lean and the private pt_part of WJTwist.lean, with P the cuspidal representation of autUO_of_printsH or the PartOfAlt data of autP.
\end{proof}

\begin{lemma}\label{sha-hecke-ideal-triangular}
Let $K$ be a field, $H$ a group, $n\ge0$, $r:H\to GL_n(K)$ a homomorphism and $\psi_1,\dots,\psi_n:H\to K^\times$ homomorphisms. Suppose that $\prod_{i=1}^n\big(r(g_i)-\psi_i(g_i)\big)=0$ for all $g_1,\dots,g_n\in H$, the product taken in increasing order of $i$; that the characteristic polynomial of $r(t)$ is $\prod_i(X-\psi_i(t))$ for every $t\in H$; and that $\psi_1(s),\dots,\psi_n(s)$ are pairwise distinct for some $s\in H$. Then there is $g\in GL_n(K)$ such that $gr(t)g^{-1}$ is upper triangular with diagonal $\psi_1(t),\dots,\psi_n(t)$ for every $t\in H$.
\end{lemma}
\leanref{Development/AutomorphicArithmetic/AutomorphicVarieties/AlgebraTools/LagTriangular.lean}{UpperTheory.V.AlgebraTools.lag_triangular}
\leanref{Development/AutomorphicArithmetic/AutomorphicVarieties/AlgebraTools/OrdExt.lean}{UpperTheory.V.AlgebraTools.ord_ext}
\begin{proof}
Extend $r$ and the $\psi_i$ linearly to the group algebra $K[H]$. Both sides of the relation are multilinear in $(g_1,\dots,g_n)$, so $\prod_i\big(r(y_i)-\psi_i(y_i)\big)=0$ for all $y_1,\dots,y_n\in K[H]$. Put $\lambda_a=\psi_a(s)$ and $e_a=\prod_{b\ne a}\big(r(s)-\lambda_b\big)/(\lambda_a-\lambda_b)$. With all $y_i=s$ the relation gives $\prod_i(r(s)-\lambda_i)=0$, so by Lagrange interpolation the $e_a$ are orthogonal idempotents with sum $1$ and $r(s)e_a=e_ar(s)=\lambda_ae_a$. Each $e_a$ is nonzero, since $\lambda_a$ is an eigenvalue of $r(s)$ and $e_a$ is the identity on its eigenvectors. So each $e_a$ has rank $1$, and there is $g\in GL_n(K)$ with $ge_ag^{-1}$ the matrix unit $E_{aa}$ for every $a$.
Let $y\in H$, and let $c,d,i$ be indices with $d\le i\le c$. Apply the relation with $y_j=s$ for $j\ne i$ and $y_i=y$, and multiply it by $e_c$ on the left and by $e_d$ on the right. For $j<i$, $e_c(r(s)-\lambda_j)=(\lambda_c-\lambda_j)e_c$ with $\lambda_c\ne\lambda_j$; for $j>i$, $(r(s)-\lambda_j)e_d=(\lambda_d-\lambda_j)e_d$ with $\lambda_d\ne\lambda_j$. Dividing by these scalars gives $e_c\big(r(y)-\psi_i(y)\big)e_d=0$. With $c=i$ and $d<i$ this is $e_ir(y)e_d=0$, and with $c=d=i$ it is $e_ir(y)e_i=\psi_i(y)e_i$. So the entry $(a,b)$ of $gr(y)g^{-1}=\sum_{a,b}g\,e_ar(y)e_b\,g^{-1}$ is $0$ for $b<a$ and $\psi_a(y)$ for $b=a$.
% The private flag_of_rel of OrdFlagJ.lean: ord_ext, lag_triangular (lag_coi for the Lagrange idempotents, the private key and mid_zero), lag_ne_zero with the characteristic polynomial, and coi_conj_single.
\end{proof}

\begin{lemma}\label{sha-hecke-ideal-flag}
In Situation~\ref{sha-hecke-situation}, let $J\subset B$ be an ideal such that for every $v\mid\ell$ and $g\in D_v$ the characteristic polynomial of $D_J=D^{\mathrm u}\bmod J$ at $g$ is $\prod_j(X-\psi_{J,v,j}(g))$, where $\psi_{J,v,j}=\psi^{\mathrm u}_{v,j}\bmod J$. Let $x$ be a point with $J\subset\ker x$, arithmetic of some weight, and let $\rho:\operatorname{Gal}(\bar L/L)\to GL_n(\bar{\mathbf Q}_\ell)$ be irreducible, with $x(D^{\mathrm u})$ the determinant of $\rho$. Then for every $v\mid\ell$, $\rho|_{D_v}$ is conjugate to an upper triangular representation with diagonal $\psi_{x,v,1},\dots,\psi_{x,v,n}$.
\end{lemma}
\leanref{Development/AutomorphicArithmetic/AutomorphicVarieties/ModularityArgument/OrdinaryJacquetModules/FlagJ.lean}{UpperTheory.V.Z.Ord.FlagJ}
\leanref{Development/AutomorphicArithmetic/AutomorphicVarieties/ModularityArgument/ThresholdArguments/DenseThresholdCases/OrdFlagJ.lean}{UpperTheory.V.Z.Ord.flagJ_alt}
\leanref{Development/AutomorphicArithmetic/AutomorphicVarieties/WeakShimuraTheory/EigenbasisLifts/RelKill.lean}{UpperTheory.V.Ws.Eb.relKill}
\leanref{Development/AutomorphicArithmetic/AutomorphicVarieties/ModularityArgument/ThresholdArguments/DenseThresholdCases/NmRegArith.lean}{UpperTheory.V.Th.Nm.regular_arithmetic_hida_of_literal_automorphy}
\begin{proof}
Let $v\mid\ell$ and $h$ be the weight of $x$. Since $\rho$ is irreducible, the image of $\rho$ spans $M_n(\bar{\mathbf Q}_\ell)$ (Burnside's theorem), so a matrix $X$ with $\operatorname{tr}(\rho(h')X)=0$ for all $h'$ is $0$, the trace form being nondegenerate.

\emph{The ordinary relations.} Let $g_1,\dots,g_n\in D_v$. Relation (2) of Situation~\ref{sha-hecke-situation}, carried to $\rho$ by $x$, gives $\operatorname{tr}\big(\rho(h')\prod_i(\rho(g_i)-\psi_{x,v,i}(g_i))\big)=0$ for all $h'$, so $\prod_i(\rho(g_i)-\psi_{x,v,i}(g_i))=0$.

\emph{The characteristic polynomials.} For $t\in D_v$ the characteristic polynomial of $\rho(t)$ is the image under $x$ of that of $D_J$ at $t$, which is $\prod_j(X-\psi_{x,v,j}(t))$.

\emph{A regular element.} Let $\phi_{v,j}$ be characters of $I_v$ with open kernels such that $\psi_{x,v,j}\phi_{v,j}=\prod_\tau\tau(\operatorname{Art}^{-1}\cdot)^{-h_{v,j,\tau}}$ on $I_v$ (Definition~\ref{sha-hecke-arcs}(1)), and let $U\subset I_v$ be the intersection of their kernels, a normal subgroup of finite index $m$. Let $s\in I_v$ with $\operatorname{Art}^{-1}(s)=1+\ell\in\mathcal O_{L_{\tilde v}}^\times$. Then $s^m\in U$, and $\psi_{x,v,j}(s^m)=\prod_\tau\tau\big((1+\ell)^m\big)^{-h_{v,j,\tau}}=Q^{-\sum_\tau h_{v,j,\tau}}$ with $Q=(1+\ell)^m>1$. There is at least one embedding $\tau$, and $j\mapsto h_{v,j,\tau}$ is strictly increasing for each $\tau$, so the sums $\sum_\tau h_{v,j,\tau}$ are pairwise distinct, and so are the $\psi_{x,v,j}(s^m)$.

Lemma~\ref{sha-hecke-ideal-triangular}, applied to $\rho|_{D_v}$, the $\psi_{x,v,j}$ and $s^m$, gives the claim.
% flagJ_alt: AlgebraTools.Det.irrIffSpan (the span), the private descend of OrdFlagJ.lean (rho factors through Gamma_0, by Ws.Eb.trPlus), Ws.Eb.relKill with B.u.Ord (the ordinary relations), Th.Nm.regular_arithmetic_hida_of_literal_automorphy (the private main of NmRegArith.lean, with emb_ne for the embeddings), hJ.O1 with AlgebraTools.Det.mapApi (the characteristic polynomials), and the private flag_of_rel.
\end{proof}

\begin{lemma}\label{sha-hecke-ideal-conjugate}
Let $E$ be a CM field with maximal totally real subfield $E^+$, $\ell$ a prime, $\iota:\bar{\mathbf Q}_\ell\to\mathbf C$ a field isomorphism, $\varepsilon$ the $\ell$-adic cyclotomic character and $n\ge1$. Let $\pi$ be a cuspidal automorphic representation of $GL_n(\mathbf A_E)$ with $\pi^\vee\cong\pi^c$, and let $\rho:\operatorname{Gal}(\bar E/E)\to GL_n(\bar{\mathbf Q}_\ell)$ be continuous and weakly attached to $\pi$ through $\iota$ (Definition~\ref{gal-weakly-attached}). Let $w$ and $w'$ be finite places of $E$ above the same place of $E^+$, let $A$ be a valuation ring of $\bar E$ over $w$, and let $\psi_1,\dots,\psi_n$ be characters of $D_A$ such that $\rho|_{D_A}$ is conjugate to an upper triangular representation with diagonal $\psi_1,\dots,\psi_n$.
\begin{enumerate}
\item Suppose that at every valuation ring $A'$ over $w'$, $\rho|_{D_{A'}}$ is conjugate to an upper triangular representation with diagonal characters $\kappa_1,\dots,\kappa_n$ such that $\kappa_a\ne\kappa_b\varepsilon$ for all $a,b$. Then $\psi_i\ne\psi_j\varepsilon$ for all $i,j$.
\item Suppose that $\psi_i^m\ne(\psi_j\varepsilon)^m$ for all $i,j$ and $m\ge1$. Then at every valuation ring $A'$ over $w'$, $\rho|_{D_{A'}}$ is conjugate to an upper triangular representation with diagonal characters $\kappa_1,\dots,\kappa_n$ such that $\kappa_a^m\ne(\kappa_b\varepsilon)^m$ for all $a,b$ and $m\ge1$.
\end{enumerate}
\end{lemma}
\leanref{Development/AutomorphicArithmetic/AutomorphicVarieties/ModularityArgument/OrdinaryJacquetModules/CoreOrdAlt.lean}{UpperTheory.V.Z.Ord.ConjUnl}
\leanref{Development/AutomorphicArithmetic/AutomorphicVarieties/ModularityArgument/ThresholdArguments/DenseThresholdCases/OrdConjUnl.lean}{UpperTheory.V.Z.Ord.conj_unlinked_of_printed_results}
\leanref{Development/AutomorphicArithmetic/AutomorphicVarieties/BuildingsResidueChars/TwistedDualFrobeniusCompatibility.lean}{UpperTheory.twistedDual_frobenius_compatibility}
\begin{proof}
\emph{The polarization.} Let $c$ be a complex conjugation in $\operatorname{Gal}(\bar E/E^+)$. Since $\pi^\vee\cong\pi^c$, the Satake parameters of $\pi$ at $c(u)$ are the inverses of those at $u$ for almost all $u$. So $g\mapsto\rho(cgc^{-1})$ and $g\mapsto\varepsilon(g)^{1-n}\,{}^t\rho(g)^{-1}$ have the same Frobenius eigenvalues at almost all places. Both are continuous and semisimple, so by the Chebotarev density theorem and the Brauer-Nesbitt theorem \cite[(30.16)]{CR62} they are conjugate: there is an invertible $P$ with ${}^t\rho(g)P\rho(cgc^{-1})=\varepsilon(g)^{1-n}P$ for every $g$.

\emph{Flags at the places over $w'$.} Let $A'$ be a valuation ring over $w'$. Since $A$ and $A'$ lie over the same place of $E^+$, there is $g\in\operatorname{Gal}(\bar E/E^+)$ with $g(A)=A'$, and $f(t)=g^{-1}tg$ is a surjection from $D_{A'}$ onto $D_A$; also $\varepsilon\circ f=\varepsilon$. If $g$ fixes $E$, then $\rho(t)=\rho(g)\rho(f(t))\rho(g)^{-1}$ for $t\in D_{A'}$, so $\rho|_{D_{A'}}$ is triangular with diagonal $\kappa_i=\psi_i\circ f$. Otherwise $u=gc^{-1}$ fixes $E$, and $u^{-1}tu=cf(t)c^{-1}$; the polarization gives $\rho(t)=\varepsilon(f(t))^{1-n}\,Q\,{}^t\rho(f(t))^{-1}Q^{-1}$ with $Q=\rho(u)P^{-1}$. The contragredient of an upper triangular representation with diagonal $\psi_1,\dots,\psi_n$ is conjugate, by transposition and the reversal of the basis, to an upper triangular one with diagonal $\psi_n^{-1},\dots,\psi_1^{-1}$; so $\rho|_{D_{A'}}$ is triangular with diagonal $\kappa_i=\psi_{n+1-i}^{-1}\circ f\cdot\mu$, where $\mu=\varepsilon^{1-n}\circ f$. In both cases, for every $m\ge1$, some $\psi_i^m=(\psi_j\varepsilon)^m$ if and only if some $\kappa_a^m=(\kappa_b\varepsilon)^m$: in the second case $\kappa_a^m=(\kappa_b\varepsilon)^m$ is equivalent to $\psi_{n+1-b}^m=(\psi_{n+1-a}\varepsilon)^m$, since $f$ is surjective.

\emph{Diagonal characters of two flags.} If $\rho|_{D_{A'}}$ is triangular with diagonals $\kappa$ and $\kappa'$, then $\sum_a\kappa_a=\sum_a\kappa'_a$ as functions on $D_{A'}$, both being the trace; by the linear independence of characters, every $\kappa_a$ is some $\kappa'_{a'}$.

(1) Choose $A'$ over $w'$, with $\kappa'$ as in the hypothesis. If $\psi_i=\psi_j\varepsilon$, the last two steps with $m=1$ give $a,b$ with $\kappa_a=\kappa_b\varepsilon$, and $a',b'$ with $\kappa'_{a'}=\kappa_a$ and $\kappa'_{b'}=\kappa_b$; then $\kappa'_{a'}=\kappa'_{b'}\varepsilon$, contrary to the hypothesis.
(2) At every $A'$ over $w'$ the characters $\kappa$ of the second step have the property.
% conj_unlinked_of_printed_results: CsdBack.isCsd_of_paired; the polarization hp from twistedDual_frobenius_compatibility (IsoOfFrob, with satake_parameters_complex_conjugate); key: CU.isPlaceOver_real, exists_smul_eq_of_isPlaceOver, CU.exists_rel (decMap, decMap_surj, cycBar_conjGal, gal_cases, rel_one and rel_conj) and CU.abs_main (flag_comp, flag_dual, flag_smul, aux1, aux2); CU.flag_range (trace_of_flag, mem_range_of_sum with linearIndependent_monoidHom).
\end{proof}

\begin{lemma}\label{sha-hecke-ideal-exponents}
In Situation~\ref{gal-local-situation}, let $\ell$ be the prime under $w$, $\iota:\bar{\mathbf Q}_\ell\to\mathbf C$ a field isomorphism, $a$ a local Artin map (Definition~\ref{gal-local-artin}), $\varepsilon$ the $\ell$-adic cyclotomic character, $\varpi$ a uniformizer of $E_w$ and $n\ge1$. For $1\le j\le n$ let $h_{j,\tau}$ be integers, indexed by the continuous embeddings $\tau:E_w\to\bar{\mathbf Q}_\ell$, with $j\mapsto h_{j,\tau}$ strictly increasing for every $\tau$, let $\omega_j$ be a smooth character of $E_w^\times$, and let $\psi_j:D_A\to\bar{\mathbf Q}_\ell^\times$ be a continuous character with $\psi_j(s)=\iota^{-1}(\omega_j(a(s)))\prod_\tau\tau(a(s))^{-h_{j,\tau}}$ for $s\in W_A$. Let $\sigma$ be an irreducible smooth representation of $GL_n(E_w)$ that is a subquotient of $\mathrm{PS}(\omega)$ (Definition~\ref{sha-hecke-local-ordinary}).
\begin{enumerate}
\item If $\psi_i\ne\psi_j\varepsilon$ for all $i,j$, then $\sigma$ is plain (Definition~\ref{gal-plain}).
\item If there is $t\in\mathbf R$ such that every character $\xi$ in the supercuspidal support of $\sigma$ satisfies $q_w^{t-1/2}<|\xi(\varpi)|<q_w^{t+1/2}$, then $\psi_i^m\ne(\psi_j\varepsilon)^m$ for all $i,j$ and $m\ge1$.
\end{enumerate}
\end{lemma}
\leanref{Development/AutomorphicArithmetic/AutomorphicVarieties/ModularityArgument/OrdinaryJacquetModules/CoreOrdAlt.lean}{UpperTheory.V.Z.Ord.ExpTor}
\leanref{Development/AutomorphicArithmetic/AutomorphicVarieties/ModularityArgument/ThresholdArguments/DenseThresholdCases/OrdExpTor.lean}{UpperTheory.V.Z.Ord.exp_torsion_of_literal_automorphy_n}
\begin{proof}
By Definition~\ref{gal-irreducible-classes}(3), the supercuspidal support of $\sigma$ consists of the characters $\xi_j=\omega_j\nu^{(n-1)/2}$ \cite{BZ77,Zel80}. Let $i,j$ be indices, put $\chi=\psi_i(\psi_j\varepsilon)^{-1}$, and let $d(s)$ be as in the proof of Lemma~\ref{sha-hecke-local-characters}, so that $\varepsilon(s)=q_w^{d(s)}\prod_\tau\tau(a(s))$ and $\iota^{-1}|a(s)|=q_w^{d(s)}$ for $s\in W_A$. Then for $s\in W_A$
\[
\chi(s)\,\iota^{-1}\big(\omega_j(a(s))\big)\,q_w^{d(s)}=\iota^{-1}\big(\omega_i(a(s))\big)\prod_\tau\tau(a(s))^{h_{j,\tau}-h_{i,\tau}-1}.
\]
Two facts about the exponents $d_\tau=h_{j,\tau}-h_{i,\tau}-1$ follow from the monotonicity: if $\sum_\tau d_\tau=0$, then all $d_\tau=0$ (if $i<j$ all $d_\tau\ge0$; if $i\ge j$ all $d_\tau<0$, and the sum cannot vanish); and the sum is $0$ if $\prod_\tau\tau(a(s))^{d_\tau}$ is the restriction to $W_A$ of a continuous character of the compact group $D_A$, by its value at an $s$ with $a(s)=\ell$, since a continuous character of $D_A$ has unit values and $|\ell|<1$.

(1) Suppose that $\xi_i=\xi_j\nu$, that is, $\omega_i=\omega_j|\cdot|$. Then $\iota^{-1}\omega_i(a(s))=\iota^{-1}\omega_j(a(s))q_w^{d(s)}$, so $\chi(s)=\prod_\tau\tau(a(s))^{d_\tau}$ on $W_A$. Hence all $d_\tau=0$, $\chi=1$ on $W_A$, and $\chi=1$ on $D_A$ since $W_A$ is dense. This contradicts the hypothesis, so no two members of the support are linked, and $\sigma$ is plain.

(2) Suppose that $\psi_i^m=(\psi_j\varepsilon)^m$, so that $\chi^m=1$ and $|\iota\chi(s)|=1$ for every $s$. The kernels of $\omega_i$ and $\omega_j$ are open, so for $k$ large the integer $N=1+\ell^k\ge2$ is a unit of $\mathcal O_w$ in both kernels; let $s\in W_A$ with $a(s)=N$, so $d(s)=0$. The displayed identity gives $\chi(s)=N^{\sum_\tau d_\tau}$, so $\sum_\tau d_\tau=0$ and all $d_\tau=0$. Now let $s_0\in W_A$ with $a(s_0)=\varpi$, so $d(s_0)=-1$. The identity gives $\chi(s_0)\,\iota^{-1}\omega_j(\varpi)\,q_w^{-1}=\iota^{-1}\omega_i(\varpi)$, and taking absolute values after $\iota$, $|\omega_j(\varpi)|=q_w|\omega_i(\varpi)|$, hence $|\xi_j(\varpi)|=q_w|\xi_i(\varpi)|$. With the bounds of the hypothesis, $|\xi_j(\varpi)|>q_w\cdot q_w^{t-1/2}=q_w^{t+1/2}$, which is impossible.
% exp_torsion_of_literal_automorphy_n: supp_ps (hasSupport_of_isSubquotient_prinSeries), twist_val (with Loc.epsArt and algPow_add'), abs_val; (1) is the private link (sum_eq_zero through unit_norm and algPow_nat, gap_zero, weil_group_dense_in_galois_group), applied with i and j exchanged after Unl; (2) is the private tors (near_one for N = 1 + l^M, ord_unif and abs_unif). The Lean states it for a cuspidal pi with ExpOf at w, with TorAt, Unl psi -> TU pi w and WinAt -> the power condition; the book states it for one local component.
\end{proof}

\subsection{Taylor-Wiles bases and level modules}\label{subsection-sha-taylor-wiles-bases}

\begin{definition}\label{sha-hecke-ideal-tw-base}
In Situation~\ref{sha-hecke-situation}, let $J\subset B$ be an ideal, $D\ge0$ an integer and $\Lambda_D=\mathcal O[[X_1,\dots,X_D]]$. A \emph{Taylor-Wiles base} for $J$ consists of places $v_k\mid\ell$ of $L^+$, indices $1\le i_k\le n$, elements $\gamma_k\in D_{v_k}$ and $c_k\in\mathcal O$ whose residues satisfy $\bar c_k\bar\psi_{v_k,i_k}(\gamma_k)=1$, for $1\le k\le D$; an $\mathcal O$-algebra map $\lambda:\Lambda_D\to B$ with $\lambda(X_k)=c_k\psi^{\mathrm u}_{v_k,i_k}(\gamma_k)-1$ for every $k$; and a $B$-module $M_0$, finite and free over $\Lambda_D$ through $\lambda$, whose support in $\operatorname{Spec}B$ is the zero locus of $J$.
\end{definition}
\leanref{Development/AutomorphicArithmetic/AutomorphicVarieties/PatchingModules/ThetaCoefficientMapEq.lean}{UpperTheory.V.Pa.WtVars}
\leanref{Development/AutomorphicArithmetic/AutomorphicVarieties/PatchingModules/ThetaCoefficientMapEq.lean}{UpperTheory.V.Pa.TWBase}

\begin{definition}\label{sha-hecke-ideal-tw-iwahori}
Let $F$ be a nonarchimedean local field with ring of integers $\mathcal O_F$, uniformizer $\varpi$ and residue field $k_F$, and let $n\ge1$.
\begin{enumerate}
\item For a subgroup $C\subset k_F^\times$, $\mathrm{Iw}_C\subset GL_n(F)$ is the subgroup generated by the matrices in $GL_n(\mathcal O_F)$ congruent to $1$ modulo $\varpi$, the upper triangular unipotent matrices with entries in $\mathcal O_F$, and the diagonal matrices whose entries are units of $\mathcal O_F$ with residues in $C$. For a prime $\ell$, $C_{\ell'}\subset k_F^\times$ is the subgroup of the elements of order prime to $\ell$. Write $\mathrm{Iw}=\mathrm{Iw}_{k_F^\times}$, the Iwahori subgroup, and $\mathrm{Iw}_{\ell'}=\mathrm{Iw}_{C_{\ell'}}$.
\item $T(F)\subset GL_n(F)$ is the diagonal torus. An element $\operatorname{diag}(t_1,\dots,t_n)$ of $T(F)$ is \emph{dominant} if $|t_1|\le|t_2|\le\dots\le|t_n|$. For $a\in F^\times$ and $1\le i\le n$, $e_i(a)\in T(F)$ has $a$ in the $i$-th place and $1$ in the others.
\item Let $A$ be a ring and $\rho$ a representation of a group $G$ on an $A$-module $P$. For subgroups $K,K'$ of $G$ and $g\in G$ such that $KgK'/K'$ is finite, the \emph{double coset operator} $[KgK']:P^{K'}\to P^K$ sends $m$ to $\sum_h\rho(h)m$, the sum over representatives $h$ of $KgK'/K'$.
\end{enumerate}
\end{definition}
\leanref{Development/AutomorphicArithmetic/AutomorphicVarieties/AlgebraTools/HeckePairTop.lean}{UpperTheory.V.AlgebraTools.iwC}
\leanref{Development/AutomorphicArithmetic/AutomorphicVarieties/PatchingModules/Classical.lean}{UpperTheory.V.Pa.ellC}
\leanref{Development/AutomorphicArithmetic/AutomorphicVarieties/AlgebraTools/HeckePairTop.lean}{UpperTheory.V.AlgebraTools.Dom}
\leanref{Development/AutomorphicArithmetic/AutomorphicVarieties/HeckeAlgebras/ContinuousTorusCharacter.lean}{UpperTheory.V.Hi.dc}
\leanref{Development/AutomorphicArithmetic/AutomorphicVarieties/AlgebraTools/HeckePairTop.lean}{UpperTheory.V.AlgebraTools.hk}

\begin{lemma}\label{sha-hecke-ideal-tw-torus}
Let $F_1,\dots,F_q$ be nonarchimedean local fields, with residue fields $k_{F_j}$, let $n\ge1$, let $A$ be a ring in which every $|k_{F_j}|$ is a unit, and let $\rho$ be a representation of $G=\prod_jGL_n(F_j)$ on an $A$-module $P$. Let $C_j\subset k_{F_j}^\times$ be subgroups and $K=\prod_j\mathrm{Iw}_{C_j}$ (Definition~\ref{sha-hecke-ideal-tw-iwahori}). Then for each $j$ there is a unique homomorphism $\Theta_j$ from $T(F_j)$ to the group of $A$-linear automorphisms of $P^K$ such that, for every dominant $t$, $\Theta_j(t)$ is the double coset operator $[\mathrm{Iw}_{C_j}t\,\mathrm{Iw}_{C_j}]$ of the $j$-th factor of $G$, restricted to $P^K$. Moreover $\Theta_j(t)$ and $\Theta_{j'}(t')$ commute for all $j,j'$, $t\in T(F_j)$ and $t'\in T(F_{j'})$.
\end{lemma}
\leanref{Development/AutomorphicArithmetic/AutomorphicVarieties/PatchingModules/NilpotentLevelLowering/ThetaQ.lean}{UpperTheory.V.Pa.Lev.thetaQ}
\leanref{Development/AutomorphicArithmetic/AutomorphicVarieties/AlgebraTools/IwThetaEx.lean}{UpperTheory.V.AlgebraTools.thetaEx}
\begin{proof}
On the monoid of dominant elements the double coset operators are multiplicative, and they are invertible because $|k_{F_j}|$ is a unit in $A$. Every element of $T(F_j)$ is a quotient of two dominant ones. This gives $\Theta_j$ and its uniqueness. The operators of two different factors commute because the factors commute in $G$, and the $\Theta_j(t)$ for one $j$ commute because $T(F_j)$ is commutative.
% thetaEx: domMul (hmul, and hunit from IsUnit (resCard k)), monoidExt on the submonoid of Dom elements, thetaEx_split for t = a b^{-1}. thetaQ: one_place for each j, comm_dc and comm_at for j /= j' on dominant elements, then dom_gen.
\end{proof}

\begin{lemma}\label{sha-hecke-ideal-tw-diamond}
Let $\ell$ be a prime, $n\ge1$, and $F_1,\dots,F_q$ nonarchimedean local fields with residue fields $k_{F_j}$. For each $j$ let $a_j$ be the exponent of $\ell$ in $|k_{F_j}|-1$, and let $\delta_j$ be a unit of $\mathcal O_{F_j}$ whose residue generates $k_{F_j}^\times$ modulo $C_{\ell'}$. Let $G=\prod_jGL_n(F_j)$, with the subgroups $\mathrm{Iw}_G=\prod_j\mathrm{Iw}$ and $\mathrm{Iw}_{\ell',G}=\prod_j\mathrm{Iw}_{\ell'}$ (Definition~\ref{sha-hecke-ideal-tw-iwahori}), and let $\Delta=\prod_{j,i}\mathbf Z/\ell^{a_j}$, over $1\le j\le q$ and $1\le i\le n$.
\begin{enumerate}
\item $\mathrm{Iw}_{\ell',G}$ is normal in $\mathrm{Iw}_G$, and there is an isomorphism $\Delta\cong\mathrm{Iw}_G/\mathrm{Iw}_{\ell',G}$ sending the generator of the factor $(j,i)$ to the class of $e_i(\delta_j)$, placed in the $j$-th factor of $G$.
\item Let $A$ be a ring in which the $|k_{F_j}|$ are units, $\rho$ a representation of $G$ on an $A$-module $P$, and $\Theta^1_j$ and $\Theta^{\mathrm{Iw}}_j$ the homomorphisms of Lemma~\ref{sha-hecke-ideal-tw-torus} on $P^{\mathrm{Iw}_{\ell',G}}$ and on $P^{\mathrm{Iw}_G}$. Then $\Theta^1_j(e_i(\delta_j))$ is the action of $e_i(\delta_j)$, placed in the $j$-th factor, on $P^{\mathrm{Iw}_{\ell',G}}$. The double coset operator $\mathcal N=[\mathrm{Iw}_G\,1\,\mathrm{Iw}_{\ell',G}]:P^{\mathrm{Iw}_{\ell',G}}\to P^{\mathrm{Iw}_G}$, $m\mapsto\sum_{g\in\mathrm{Iw}_G/\mathrm{Iw}_{\ell',G}}\rho(g)m$, satisfies $\mathcal N\circ\Theta^1_j(t)=\Theta^{\mathrm{Iw}}_j(t)\circ\mathcal N$ for all $j$ and $t\in T(F_j)$.
\end{enumerate}
\end{lemma}
\leanref{Development/AutomorphicArithmetic/AutomorphicVarieties/PatchingModules/NilpotentLevelLowering/DelIso.lean}{UpperTheory.V.Pa.Lev.delIso}
\leanref{Development/AutomorphicArithmetic/AutomorphicVarieties/PatchingModules/NilpotentLevelLowering/TrOfCommdc.lean}{UpperTheory.V.Pa.Lev.trOf}
\begin{proof}
(1) In each factor, $\mathrm{Iw}_{\ell'}$ is the inverse image of $C_{\ell'}^n$ under the homomorphism from $\mathrm{Iw}$ to $(k_{F_j}^\times)^n$ that reduces the diagonal entries, and $k_{F_j}^\times/C_{\ell'}$ is cyclic of order $\ell^{a_j}$, generated by the residue of $\delta_j$.
(2) The element $e_i(\delta_j)$ is dominant and lies in $\mathrm{Iw}$, which normalizes $\mathrm{Iw}_{\ell'}$, so its double coset is a single coset and its operator is the action of $e_i(\delta_j)$. The cosets of $\mathrm{Iw}_G/\mathrm{Iw}_{\ell',G}$ have diagonal representatives, which commute with $t$; this gives the commutation for dominant $t$, and so for all $t$.
% delIso: di_loc for each place (iwQuot: the reduction of the diagonal with kernel iwC C; di_cyc: the cyclic quotient generated by delta from IsGenMod), di_pi for the product. trOf: trOf_gen (e_i(delta) is Dom and lies in iwC top), trOf_unit, trOf_fix; hk2 for the norm map; trOf_main (trOf_reps: diagonal representatives of iwC top / iwC C; trOf_comm; trOf_split and trOf_ext for general t) with trOf_smul for the scalars q^i.
\end{proof}

\begin{lemma}\label{sha-hecke-ideal-tw-dimension}
Let $F$ be a nonarchimedean local field, $n\ge1$, $C\subset k_F^\times$ a subgroup and $\sigma$ an irreducible smooth representation of $GL_n(F)$ on a complex vector space. Then $\sigma^{\mathrm{Iw}_C}$ is finite-dimensional, of dimension at most $n!$.
\end{lemma}
\leanref{Development/AutomorphicArithmetic/AutomorphicVarieties/AlgebraTools/IwJqLine.lean}{UpperTheory.V.AlgebraTools.jqLine}
\begin{proof}
The invariants of $\mathrm{Iw}_C$ map injectively to the Jacquet module of $\sigma$ for the upper triangular unipotent subgroup, by the theorem of Borel and Casselman \cite{Bor76}. If $\sigma^{\mathrm{Iw}_C}\ne0$, then $\sigma$ is a subquotient of a principal series, and by the geometric lemma of Bernstein and Zelevinsky its Jacquet module has dimension at most $n!$ \cite{BZ77}.
% jqLine: finiteness from BernsteinZelevinskyFiniteness (admissibility), jqInj (injectivity into the Jacquet module), jqEmbed (a principal series containing sigma as a subquotient, from a common eigenvector of thetaEx), jqMult (GeometricLemma: the Jacquet module has dimension at most n!).
\end{proof}

\begin{lemma}\label{sha-hecke-ideal-tw-local}
Let $E$ be a number field, $\ell$ a prime, $w$ a finite place of $E$ not above $\ell$, with residue field $k_w$ of $q_w$ elements, and $A$ a valuation ring of $\bar E$ over $w$. Let $\phi$ be an element of the decomposition group $D_A$ whose inverse is a Frobenius element, and $\tau$ an element of the inertia group $I_A$ that topologically generates the $\ell$-part of tame inertia. Then there are a uniformizer $x$ of $E_w$ and a unit $\delta$ of $\mathcal O_{E_w}$ whose residue generates $k_w^\times$ modulo $C_{\ell'}$ with the following property. Let $n\ge1$, $\iota:\bar{\mathbf Q}_\ell\to\mathbf C$ a field isomorphism, $r:\operatorname{Gal}(\bar E/E)\to GL_n(\bar{\mathbf Q}_\ell)$ a homomorphism and $\sigma$ an irreducible smooth representation of $GL_n(E_w)$ such that, for every $\operatorname{rc}$ that satisfies the laws (Definition~\ref{gal-llc-laws}), $r|_{W_A}$ has a Weil-Deligne representation $\mathrm{WD}_\iota(r|_{W_A})$, equal to $\operatorname{rc}(\sigma\nu^{(1-n)/2})$ when the latter is abelian and monodromy-free. Let $C\subset k_w^\times$ be a subgroup containing $C_{\ell'}$, let $\Theta$ be the homomorphism of Lemma~\ref{sha-hecke-ideal-tw-torus} for $\sigma$ and $\mathrm{Iw}_C$, and let $\chi$ be a character of $T(E_w)$ such that the common generalized eigenspace of the $\Theta(t)$ for the eigenvalues $\chi(t)$ is nonzero. Put $\alpha_i=q_w^{i-1}\chi(e_i(x))$, and call $\alpha$ \emph{generic} if $\alpha_i\ne\alpha_{i'}$ and $\alpha_i\ne q_w\alpha_{i'}$ for all $i\ne i'$. Then:
\begin{enumerate}
\item if $\sigma$ has nonzero vectors fixed by $GL_n(\mathcal O_{E_w})$, or $\alpha$ is generic, the characteristic polynomial of $\iota r(\phi)$ is $\prod_i(X-\alpha_i)$;
\item if $\alpha$ is generic, then $\prod_{i'\ne i}\big(\iota r(\phi)-\alpha_{i'}\big)\cdot\big(\iota r(\tau)-\chi(e_i(\delta))\big)=0$ for every $i$.
\end{enumerate}
\end{lemma}
\leanref{Development/AutomorphicArithmetic/AutomorphicVarieties/ModularityArgument/ThresholdArguments/DenseThresholdCases/LocDataT.lean}{UpperTheory.V.Z.Win.LocDataTAlt}
\leanref{Development/AutomorphicArithmetic/AutomorphicVarieties/ModularityArgument/AutomorphyLifting/LocDataT.lean}{UpperTheory.V.Z.Win.LocRelT}
\leanref{Development/AutomorphicArithmetic/AutomorphicVarieties/ModularityArgument/ThresholdArguments/DenseThresholdCases/TLocDataT.lean}{UpperTheory.V.Z.Win.locDataT_alt}
\begin{proof}
\emph{The elements $x$ and $\delta$.} Let $\operatorname{rc}_0$ satisfy the laws (Theorem~\ref{gal-llc-exists}) with respect to a local Artin map $a:W_A\to E_w^\times$ (Definition~\ref{gal-local-artin}). The elements $\phi$ and $\tau$ lie in $W_A$. By local class field theory and the structure of the tame quotient of $I_A$, $x=a(\phi)$ is a uniformizer, $\delta=a(\tau)$ is a unit whose residue generates $k_w^\times$ modulo $C_{\ell'}$, $a$ maps the kernel of the map from $I_A$ to its tame $\ell$-quotient into the units with residues in $C_{\ell'}$, and $g_0=\tau^{-q_w}\phi^{-1}\tau\phi$ lies in that kernel.

\emph{The characters $c_i$.} Let $n$, $\iota$, $r$, $\sigma$, $C$, $\Theta$ and $\chi$ be as in the statement. Since the common generalized eigenspace for $\chi$ is nonzero, $\sigma$ is a subquotient of a principal series $\mathrm{PS}(\psi_1,\dots,\psi_n)$ of smooth characters $\psi_i$ of $E_w^\times$ that are trivial on the units with residues in $C$, with $\prod_i\psi_i(t_i)=\chi(t)\,\delta_B^{1/2}(t)$ for $t=\operatorname{diag}(t_1,\dots,t_n)$, where $\delta_B$ is the modulus character of the upper triangular Borel subgroup \cite{BZ77}. So $\psi_i(z)=q_w^{(i-1)\operatorname{ord}z}\chi(e_i(z))|z|^{(n-1)/2}$. By the laws for $\operatorname{rc}_0$, for $s\in W_A$ the characteristic polynomial of $\operatorname{rc}_0(\sigma\nu^{(1-n)/2})(s)$ is $\prod_i(X-c_i(s))$, where $c_i(s)=q_w^{(i-1)\operatorname{ord}a(s)}\chi(e_i(a(s)))$ is a character of $W_A$. As $\operatorname{ord}x=1$, $c_i(\phi)=\alpha_i$; as $\tau\in I_A$, $c_i(\tau)=\chi(e_i(\delta))$.

\emph{Abelian and monodromy-free.} If $\sigma$ has nonzero vectors fixed by $GL_n(\mathcal O_{E_w})$, then $\operatorname{rc}_0(\sigma)$ is unramified, by the laws, hence abelian and monodromy-free, and so is its twist $\operatorname{rc}_0(\sigma\nu^{(1-n)/2})$. If $\alpha$ is generic, then no two of the $\psi_i$ differ by $|\cdot|$: $\psi_i=\psi_{i'}|\cdot|$ would give $\alpha_{i'}=q_w\alpha_i$ at $z=x$, which is excluded for $i\ne i'$, and gives $q_w=1$ for $i=i'$. So $\sigma$ is plain (Definition~\ref{gal-plain}), and $\operatorname{rc}_0(\sigma\nu^{(1-n)/2})$ is again abelian and monodromy-free.

\emph{(1).} In both cases the hypothesis gives $\mathrm{WD}_\iota(r|_{W_A})=\operatorname{rc}_0(\sigma\nu^{(1-n)/2})$, so the characteristic polynomial of $\iota r(s)$ is $\prod_i(X-c_i(s))$ for $s\in W_A$. At $s=\phi$ this is (1).

\emph{(2).} Let $\alpha$ be generic. The element $a(g_0)$ is a unit with residue in $C_{\ell'}\subset C$, so the $\psi_i$ and $\delta_B$ are trivial on it, and $c_i(g_0)=1$ for every $i$. So the characteristic polynomial of $\iota r(g_0)$ is $(X-1)^n$. As $g_0$ is in the kernel of the tame $\ell$-quotient of $I_A$, the Weil-Deligne representation of $r|_{W_A}$ shows that $r(g_0)$ has finite order; being unipotent, $r(g_0)=1$. Hence $R=\iota r|_{W_A}$ satisfies $R(\tau)R(\phi)=R(\phi)R(\tau)^{q_w}$, and $c_i(\tau)^{q_w-1}=1$ for every $i$. With the characteristic polynomials $\prod_i(X-c_i(s))$ of the $R(s)$, which give the traces of the $R(\phi)^bR(\tau)^a$, and the genericity of the $c_i(\phi)=\alpha_i$, a computation in linear algebra gives $\prod_{i'\ne i}(R(\phi)-\alpha_{i'})(R(\tau)-c_i(\tau))=0$, which is (2).
% locDataT_alt: recTExists and extraV (the laws LawsV for the Artin map a of IsRecT); artinEl (from the published hypothesis TameLocal: resCard = normV, x = a(fr) IsUnif, delta = a(tau) IsGenMod, a(ineKer) in unitsC ellC, and the tame relation); jqEmbed (BernsteinZelevinskyFiniteness); nuCharV (with jqRecV from the support law, L.nu_eq, psi_val); hord1; hsph (L.sph and Gu.gdLaws); hgen (unlinked_of_gen and Gu.plainGd); core: wdChar (charpoly of r(s) from IsWDOf; part 2: on ineKer the monodromy representation has finite image, so a unipotent r(g) is 1), chi_unit, and the private rel_of_chars of LLocData.lean (prod_sub_lam_mul_sub_mu_eq_zero with the relation R(t)R(f) = R(f)R(t)^q, the semisimplicity of R(f) from wdb_ofReal_eq_conj_wd_ofReal, (R(t)^(q-1) - 1)^n = 0 and the traces).
\end{proof}

\begin{definition}\label{sha-hecke-ideal-tw-level-module}
In Situation~\ref{sha-hecke-situation}, let $(\lambda,M_0)$ be a Taylor-Wiles base for an ideal $J\subset B$ (Definition~\ref{sha-hecke-ideal-tw-base}), with $\Lambda_D$, $c_k$, $v_k$, $i_k$ and $\gamma_k$ as there. Let $S'\supset S$ be a finite set of places of $L^+$, $\Gamma'$ the Galois group of the maximal extension of $L^+$ unramified outside $S'$ and the infinite places, $p:\Gamma'\to\Gamma$ the projection and $\Gamma'_0=p^{-1}(\Gamma_0)$. An \emph{ordinary pair for $\Gamma'$} over a complete local $\mathcal O$-algebra $R$ with residue field $k$ is a pair $(D_R,\psi_R)$ that satisfies the conditions on an ordinary pair of Situation~\ref{sha-hecke-situation}, with $\Gamma'_0$ in place of $\Gamma_0$ and the residual datum $(\bar\chi_i\circ p,\bar\psi_{v,i},\theta_{v,i})$. Let $w_1,\dots,w_q$ be finite places of $L$, $G_w=\prod_jGL_n(L_{w_j})$, and $K_w=\prod_jGL_n(\mathcal O_{L_{w_j}})$, $\mathrm{Iw}_w=\prod_j\mathrm{Iw}$ and $\mathrm{Iw}_{\ell',w}=\prod_j\mathrm{Iw}_{\ell'}$ its subgroups (Definition~\ref{sha-hecke-ideal-tw-iwahori}).
\begin{enumerate}
\item A \emph{module at level $\Gamma'$ for the places $w_j$} consists of: a complete local noetherian $\mathcal O$-algebra $R$ with residue field $k$, an ordinary pair $(D_R,\psi_R)$ for $\Gamma'$ over $R$, an $\mathcal O$-algebra map $\Lambda_D\to R$ sending each $X_k$ to $c_k\psi_{R,v_k,i_k}(\gamma_k)-1$, and an $\mathcal O$-algebra map $\beta:R\to B$ carrying $(D_R,\psi_R)$ to $(D^{\mathrm u}\circ p,\psi^{\mathrm u})$; an $R$-module $P$ with an $R$-linear action $\rho$ of $G_w$, such that $P^{\mathrm{Iw}_{\ell',w}}$ is finite over $\Lambda_D$, $\mathrm{Iw}_{\ell',w}$ is normal in $\mathrm{Iw}_w$, and $P^{\mathrm{Iw}_{\ell',w}}$ is projective over the group ring $\Lambda_D[\mathrm{Iw}_w/\mathrm{Iw}_{\ell',w}]$; and an additive isomorphism $\epsilon:P^{K_w}\to M_0$ with $\epsilon(tm)=\beta(t)\epsilon(m)$ for $t\in R$.
\item Let $\iota:\bar{\mathbf Q}_\ell\to\mathbf C$ be a field isomorphism. A \emph{point} $y$ of such a module consists of: a ring map $\theta_y:R\to\bar{\mathbf Q}_\ell$; a continuous homomorphism $\rho_y:\Gamma'_0\to GL_n(\bar{\mathbf Q}_\ell)$ such that $\theta_y(D_R)$ is the determinant of $\rho_y$; a representation $\pi_y$ of $G_w$ on a complex vector space $V_y$ and an additive map $\operatorname{sp}_y:P\to V_y$ with $\operatorname{sp}_y(tm)=\iota\theta_y(t)\operatorname{sp}_y(m)$ and $\operatorname{sp}_y(\rho(g)m)=\pi_y(g)\operatorname{sp}_y(m)$; and irreducible smooth representations $\sigma_{y,j}$ of $GL_n(L_{w_j})$ such that $V_y$, as a representation of the $j$-th factor of $G_w$, is a sum of copies of $\sigma_{y,j}$. View $\rho_y$ as a representation of $\operatorname{Gal}(\bar L/L)$ through $\Gamma'_0$. It is required that for every $j$, every valuation ring $A$ of $\bar L$ over $w_j$ and every $\operatorname{rc}$ that satisfies the laws (Definition~\ref{gal-llc-laws}), $\rho_y|_{W_A}$ has a Weil-Deligne representation $\mathrm{WD}_\iota(\rho_y|_{W_A})$, equal to $\operatorname{rc}(\sigma_{y,j}\nu^{(1-n)/2})$ when the latter is abelian and monodromy-free.
\item The module is \emph{separated by its points} for $\iota$ if every $m\in P$ with $\operatorname{sp}_y(m)=0$ for every point $y$ is $0$.
\item Suppose that no $w_j$ lies above $\ell$, and let $q_{w_j}$ be the number of elements of the residue field $k_{w_j}$. Let $x_j$ be uniformizers of $L_{w_j}$, $\delta_j$ units of $\mathcal O_{L_{w_j}}$ whose residues generate $k_{w_j}^\times$ modulo $C_{\ell'}$, and $\phi_j,\tau_j\in\Gamma'_0$. The module satisfies the \emph{local conditions} for the $x_j$, $\delta_j$, $\phi_j$, $\tau_j$ and $\iota$ if it is separated by its points for $\iota$ and, for every point $y$, every $j$, every subgroup $C\supset C_{\ell'}$ of $k_{w_j}^\times$ and every character $\chi$ as in Lemma~\ref{sha-hecke-ideal-tw-local}, the conclusions (1) and (2) of that lemma hold for $\sigma_{y,j}$, $x_j$, $\delta_j$, $q_{w_j}$ and the matrices $\iota\rho_y(\phi_j)$ and $\iota\rho_y(\tau_j)$ in place of $\iota r(\phi)$ and $\iota r(\tau)$.
\end{enumerate}
\end{definition}
\leanref{Development/AutomorphicArithmetic/AutomorphicVarieties/PatchingModules/Classical.lean}{UpperTheory.V.Pa.LevMod}
\leanref{Development/AutomorphicArithmetic/AutomorphicVarieties/ModularityArgument/ThresholdArguments/DenseThresholdCases/PaLevel.lean}{UpperTheory.V.Pa.LevPtAlt}
\leanref{Development/AutomorphicArithmetic/AutomorphicVarieties/ModularityArgument/ThresholdArguments/DenseThresholdCases/PaLevel.lean}{UpperTheory.V.Pa.LevMod.ClassicalAlt}
\leanref{Development/AutomorphicArithmetic/AutomorphicVarieties/ModularityArgument/ThresholdArguments/DenseThresholdCases/LocOk.lean}{UpperTheory.V.Pa.Lev.LocalLevelConditions}

\begin{lemma}\label{sha-hecke-ideal-tw-kill}
In Situation~\ref{sha-hecke-situation}, let $(\lambda,M_0)$ be a Taylor-Wiles base for an ideal $J\subset B$, $\Gamma'$ and $w_1,\dots,w_q$ as in Definition~\ref{sha-hecke-ideal-tw-level-module}, with residue fields of $q_{w_j}$ elements, $\iota:\bar{\mathbf Q}_\ell\to\mathbf C$ a field isomorphism, and $(R,P)$ a module at level $\Gamma'$ for the places $w_j$ that is separated by its points for $\iota$. Suppose that the $q_{w_j}$ are units in $R$, and let $\Theta^1_j$ be the homomorphisms of Lemma~\ref{sha-hecke-ideal-tw-torus} on $P^{\mathrm{Iw}_{\ell',w}}$. Suppose that no $w_j$ lies above $\ell$. Let $x_j$ be uniformizers of $L_{w_j}$ and $\delta_j$ units of $\mathcal O_{L_{w_j}}$ whose residues generate $k_{w_j}^\times$ modulo $C_{\ell'}$; these three hypotheses, on the $w_j$, the $x_j$ and the $\delta_j$, are not used in the proof. Put $U_{j,i}=q_{w_j}^{i-1}\Theta^1_j(e_i(x_j))$ and $Y_{j,i}=\Theta^1_j(e_i(\delta_j))-1$. Fix $j$ and a polynomial $F\in R[u_1,\dots,u_n,y_1,\dots,y_n]$. Suppose that for every point $y$ and every character $\chi$ of $T(L_{w_j})$ such that the homomorphism $\Theta$ of Lemma~\ref{sha-hecke-ideal-tw-torus} for $\sigma_{y,j}$ and $\mathrm{Iw}_{\ell'}$ has a nonzero common generalized eigenspace for the eigenvalues $\chi(t)$, the image of $F$ under the ring map $R[u,y]\to\mathbf C$ that is $\iota\theta_y$ on $R$ and sends $u_i$ to $q_{w_j}^{i-1}\chi(e_i(x_j))$ and $y_i$ to $\chi(e_i(\delta_j))-1$ is $0$. Then $F(U_{j,1},\dots,U_{j,n},Y_{j,1},\dots,Y_{j,n})^{n!}=0$ on $P^{\mathrm{Iw}_{\ell',w}}$.
\end{lemma}
\leanref{Development/AutomorphicArithmetic/AutomorphicVarieties/ModularityArgument/ThresholdArguments/DenseThresholdCases/LevRelOf.lean}{UpperTheory.V.Pa.Lev.relOf_alt}
\leanref{Development/AutomorphicArithmetic/AutomorphicVarieties/PatchingModules/NilpotentLevelLowering/NilKill.lean}{UpperTheory.V.Pa.Lev.nilKill}
\begin{proof}
Let $m\in P^{\mathrm{Iw}_{\ell',w}}$ and $m'=F(U_{j,\cdot},Y_{j,\cdot})^{n!}m$. Since the module is separated by its points, it suffices to show $\operatorname{sp}_y(m')=0$ for every point $y$.

\emph{Transfer to the point.} Fix $y$. The integers $q_{w_j}$ are units in $\mathbf C$, so Lemma~\ref{sha-hecke-ideal-tw-torus} gives homomorphisms $\Theta^y_j$ on $V_y^{\mathrm{Iw}_{\ell',w}}$. The map $\operatorname{sp}_y$ carries $P^{\mathrm{Iw}_{\ell',w}}$ into $V_y^{\mathrm{Iw}_{\ell',w}}$, because it intertwines $\rho$ and $\pi_y$. For the same reason it intertwines the double coset operators of the dominant elements, hence $\Theta^1_j(t)$ and $\Theta^y_j(t)$ for all $t$. Let $F_y\in\mathbf C[u,y]$ be the image of $F$ under $\iota\theta_y$ on the coefficients, and $G=F_y\big(q_{w_j}^{i-1}\Theta^y_j(e_i(x_j)),\Theta^y_j(e_i(\delta_j))-1\big)$. Since $\operatorname{sp}_y(tm)=\iota\theta_y(t)\operatorname{sp}_y(m)$, $\operatorname{sp}_y(m')=G^{n!}\operatorname{sp}_y(m)$.

\emph{The operator on $\sigma_{y,j}$.} Let $\Theta$ be the homomorphism of Lemma~\ref{sha-hecke-ideal-tw-torus} for $\sigma=\sigma_{y,j}$ and $\mathrm{Iw}_{\ell'}$, and $G_\sigma$ the same expression as $G$ with $\Theta$ in place of $\Theta^y_j$. By Lemma~\ref{sha-hecke-ideal-tw-dimension}, $\sigma^{\mathrm{Iw}_{\ell'}}$ has dimension at most $n!$. Suppose $G_\sigma^{n!}\ne0$. Its image is a nonzero finite-dimensional subspace stable under the commuting automorphisms $\Theta(t)$, so it contains a common eigenvector $v=G_\sigma^{n!}u\ne0$, with eigenvalues $\chi(t)$ for a character $\chi$ of $T(L_{w_j})$. Then $\chi$ has a nonzero common generalized eigenspace, so by hypothesis $G_\sigma v$ is the image of $F$ at $\chi$ times $v$, which is $0$. Hence $u$ lies in the kernel of $G_\sigma^{n!+1}$, which is the kernel of $G_\sigma^{n!}$ since the dimension is at most $n!$; so $v=0$, a contradiction. Thus $G_\sigma^{n!}=0$.

\emph{Back to $V_y$.} Every vector of $V_y^{\mathrm{Iw}_{\ell',w}}$ is fixed by $\mathrm{Iw}_{\ell'}$ in the $j$-th factor. As $V_y$ is a sum of copies of $\sigma_{y,j}$ for this factor, such a vector is a sum of images of vectors of $\sigma^{\mathrm{Iw}_{\ell'}}$ under equivariant maps $\sigma\to V_y$. These maps intertwine the double coset operators of the dominant elements, hence $\Theta(t)$ with the homomorphism of Lemma~\ref{sha-hecke-ideal-tw-torus} for the $j$-th factor of $V_y$ and $\mathrm{Iw}_{\ell'}$, which agrees with $\Theta^y_j$ on $V_y^{\mathrm{Iw}_{\ell',w}}$. So $G^{n!}=0$ on $V_y^{\mathrm{Iw}_{\ell',w}}$, and $\operatorname{sp}_y(m')=0$.
% relOf_alt: thetaQ for y.pi, relOf_act, thetaMap (sp intertwines Theta1 and Theta for y.pi), nilKill, relOf_uy (with relOf_int). nilKill: thetaEx for pi at the j-th factor, nilKill_loc (thetaEx for sigma, jqLine for the dimension bound, nilKill_nil with nilKill_eig, isoFix for the sum of copies, nilKill_dc, nilKill_split and nilKill_ext for the intertwining), thetaLoc for the agreement with Theta_j. The hypotheses w_j not above l, x_j a uniformizer and delta_j as in IsGenMod are not used.
\end{proof}

\begin{definition}\label{sha-hecke-ideal-tw-operator-module}
In Situation~\ref{sha-hecke-situation}, let $(\lambda,M_0)$ be a Taylor-Wiles base for an ideal $J\subset B$, with $\Lambda_D$, $c_k$, $v_k$, $i_k$ and $\gamma_k$ as in Definition~\ref{sha-hecke-ideal-tw-base}, and let $\Gamma'$, $p$ and $\Gamma'_0$ be as in Definition~\ref{sha-hecke-ideal-tw-level-module}. Let $q,N\ge0$.
\begin{enumerate}
\item An \emph{operator module of level $N$} for $\Gamma'$ consists of: a complete local noetherian $\mathcal O$-algebra $R$ with residue field $k$, an ordinary pair $(D_R,\psi_R)$ for $\Gamma'$ over $R$, an $\mathcal O$-algebra map $\Lambda_D\to R$ sending each $X_k$ to $c_k\psi_{R,v_k,i_k}(\gamma_k)-1$, and an $\mathcal O$-algebra map $\beta:R\to B$ carrying $(D_R,\psi_R)$ to $(D^{\mathrm u}\circ p,\psi^{\mathrm u})$; an $R$-module $M_1$, finite over $\Lambda_D$, with commuting $R$-linear endomorphisms $U_{j,i}$ and $Y_{j,i}$, for $1\le j\le q$ and $1\le i\le n$, and an $R$-module $M_w$ with commuting $R$-linear endomorphisms $U^w_{j,i}$; integers $a_j\ge N$ and a $\Lambda_D$-linear action of $\Delta=\prod_{j,i}\mathbf Z/\ell^{a_j}$ on $M_1$ in which the generator of the factor $(j,i)$ acts as $1+Y_{j,i}$, such that $M_1$ is projective over $\Lambda_D[\Delta]$; and a surjective $R$-linear map $\mathcal N:M_1\to M_w$ with kernel $\sum_{j,i}Y_{j,i}M_1$ and $\mathcal N\circ U_{j,i}=U^w_{j,i}\circ\mathcal N$.
\item Let $b_1,\dots,b_q$ be integers, $\phi_j,\tau_j\in\Gamma'_0$ and $e\ge0$. Let $P_j$ be the characteristic polynomial of $D^{\mathrm u}$ at $p(\phi_j)$, let $B[U]$ be the quotient of the polynomial ring $B[U_{j,i}]$ by the coefficients of the polynomials $\prod_i(X-U_{j,i})-P_j$, and let $f=\prod_j\prod_i\prod_{i'\ne i}(U_{j,i}-U_{j,i'})(U_{j,i}-b_jU_{j,i'})$. For $t$ in the polynomial ring $R[U_{j,i}]$, write $\beta(t)$ for its image in $B[U]$ and $t(U^w)$ for the endomorphism of $M_w$ obtained by substituting the $U^w_{j,i}$. The operator module \emph{satisfies the relations of exponent $e$} for the $b_j$, $\phi_j$ and $\tau_j$ if:
\begin{enumerate}
\item there are additive maps $\varphi_U:B[U]\otimes_BM_0\to M_w$ and $\psi_U:M_w\to B[U]\otimes_BM_0$ with $\varphi_U(\beta(t)z)=t(U^w)\varphi_U(z)$ and $\psi_U(t(U^w)v)=\beta(t)\psi_U(v)$ for $t\in R[U_{j,i}]$, such that $\varphi_U\psi_U$ is $f(U^w)$ and $\psi_U\varphi_U$ is multiplication by $\beta(f)$;
\item in the polynomial ring $R[U_{j,i},Y_{j,i}]$, with $\operatorname{tr}$ the trace of $D_R$ extended linearly to $R[U_{j,i},Y_{j,i}][\Gamma'_0]$, the $e$-th powers of $f\operatorname{tr}\big(h\prod_i(\phi_j-U_{j,i})\big)$ and of $f\operatorname{tr}\big(h\prod_{k\ne i}(\phi_j-U_{j,k})\,(\tau_j-1-Y_{j,i})\big)$, for all $j$, $i$ and $h\in\Gamma'_0$, act on $M_1$ as $0$ when each variable acts as the endomorphism of the same name.
\end{enumerate}
\end{enumerate}
\end{definition}
\leanref{Development/AutomorphicArithmetic/AutomorphicVarieties/PatchingModules/NilpotentLevelLowering/LevelModule.lean}{UpperTheory.V.Pa.Lev.OpMod}
\leanref{Development/AutomorphicArithmetic/AutomorphicVarieties/PatchingModules/NilpotentLevelLowering/LevelModule.lean}{UpperTheory.V.Pa.Lev.TwoMaps}
\leanref{Development/AutomorphicArithmetic/AutomorphicVarieties/PatchingModules/NilpotentLevelLowering/LevelModule.lean}{UpperTheory.V.Pa.Lev.OpRel}

\begin{definition}\label{sha-hecke-ideal-tw-divided}
Let $R$ be a commutative ring, $n\ge1$ and $q\in R$. Let the symmetric group $S_n$ act on $R[X_1,\dots,X_n]$ by $w(X_i)=X_{w(i)}$, and let $\Delta=\prod_{i<j}(X_j-X_i)$ and $N_q=\prod_{i<j}(X_j-qX_i)$, the products over $1\le i<j\le n$. For a polynomial $x$ let $\operatorname{alt}(x)=\sum_{w\in S_n}\operatorname{sgn}(w)\,w(x)$. The \emph{divided difference} $\delta(x)$ is a polynomial $y$ with $\Delta y=\operatorname{alt}(x)$ if there is one, and $0$ otherwise; when $R$ is a domain, $y$ is unique. Two families $(a_\kappa)_{\kappa\in\Sigma}$ and $(a'_\kappa)_{\kappa\in\Sigma}$ of polynomials, indexed by a finite set $\Sigma$, are \emph{dual} if for every $w\in S_n$ the sum $\sum_\kappa a_\kappa\,w(a'_\kappa)$ is $\Delta$ when $w=1$ and $0$ otherwise.
\end{definition}
\leanref{Development/AutomorphicArithmetic/AutomorphicVarieties/AlgebraTools/HkDefs.lean}{UpperTheory.V.AlgebraTools.Hk.De}
\leanref{Development/AutomorphicArithmetic/AutomorphicVarieties/AlgebraTools/HkDefs.lean}{UpperTheory.V.AlgebraTools.Hk.NN}
\leanref{Development/AutomorphicArithmetic/AutomorphicVarieties/AlgebraTools/HkDefs.lean}{UpperTheory.V.AlgebraTools.Hk.dd}
\leanref{Development/AutomorphicArithmetic/AutomorphicVarieties/AlgebraTools/HkDefs.lean}{UpperTheory.V.AlgebraTools.Hk.IsDual}

\begin{lemma}\label{sha-hecke-ideal-tw-dual}
Let $R$ be a domain and $n\ge1$, with the notation of Definition~\ref{sha-hecke-ideal-tw-divided}, and let $\Pi$ be the set of pairs $(i,j)$ with $1\le i<j\le n$.
\begin{enumerate}
\item The families $a_t=\prod_{(i,j)\in t}X_j$ and $a'_t=\prod_{(i,j)\in\Pi\smallsetminus t}(-X_i)$, indexed by the subsets $t\subset\Pi$, are dual.
\item For every $x$, $\Delta\,\delta(x)=\operatorname{alt}(x)$. The map $\delta$ is additive, $\delta(zx)=z\,\delta(x)$ for $z$ symmetric, and $\delta(x)$ is symmetric.
\item If $(a_\kappa)$ and $(a'_\kappa)$ are dual, then $\sum_\kappa a_\kappa\,\delta(a'_\kappa x)=x$ and $\sum_\kappa a'_\kappa\,\delta(a_\kappa x)=x$ for every $x$.
\end{enumerate}
\end{lemma}
\leanref{Development/AutomorphicArithmetic/AutomorphicVarieties/AlgebraTools/HkDualEx.lean}{UpperTheory.V.AlgebraTools.hkDualEx}
\leanref{Development/AutomorphicArithmetic/AutomorphicVarieties/AlgebraTools/HkDdSpec.lean}{UpperTheory.V.AlgebraTools.hkDdSpec}
\leanref{Development/AutomorphicArithmetic/AutomorphicVarieties/AlgebraTools/HkDdProps.lean}{UpperTheory.V.AlgebraTools.hkDdProps}
\leanref{Development/AutomorphicArithmetic/AutomorphicVarieties/AlgebraTools/HkDualExp.lean}{UpperTheory.V.AlgebraTools.hkDualExp}
\begin{proof}
(1) Expanding the product, $\sum_ta_t\,w(a'_t)=\prod_{(i,j)\in\Pi}(X_j-X_{w(i)})$. For $w=1$ this is $\Delta$. For $w\ne1$ let $i$ be the least index with $w(i)\ne i$; then $w(i)>i$, and the factor for the pair $(i,w(i))$ is $0$.

(2) Let $i<j$. Substituting $X_i$ for $X_j$ in $\operatorname{alt}(x)$ gives $0$, since the terms for $w$ and $(i\;j)w$ cancel; so the prime element $X_j-X_i$ divides $\operatorname{alt}(x)$. These elements are pairwise non-associate, so $\Delta$ divides $\operatorname{alt}(x)$, and $\Delta\,\delta(x)=\operatorname{alt}(x)$. As $\Delta$ is not a zero divisor, the remaining statements follow from the additivity of $\operatorname{alt}$, from $\operatorname{alt}(zx)=z\operatorname{alt}(x)$ for $z$ symmetric, and from $w(\Delta)=\operatorname{sgn}(w)\Delta$ and $w(\operatorname{alt}(x))=\operatorname{sgn}(w)\operatorname{alt}(x)$.

(3) By (2), $\Delta\sum_\kappa a_\kappa\delta(a'_\kappa x)=\sum_\kappa a_\kappa\operatorname{alt}(a'_\kappa x)=\sum_w\operatorname{sgn}(w)\big(\sum_\kappa a_\kappa\,w(a'_\kappa)\big)w(x)=\Delta x$, and $\Delta$ is not a zero divisor. The families $(a'_\kappa)$ and $(a_\kappa)$ are also dual, as one sees by applying $w$ to the identity for $w^{-1}$; this gives the second identity.
% hkDualEx: the expansion hsum. hkDdSpec: dvd_iff_sb (X_j - X_i divides y iff the substitution sb i j kills y), alt_sb, prime_sb and finset_prod_dvd_of_prime_of_dvd_imp_eq. hkDdProps: hadd, hmul, hsym through eq_of_De_mul_eq_De_mul and rename_alt. hkDualExp: dual_exp and dual_symm.
\end{proof}

\begin{lemma}\label{sha-hecke-ideal-tw-bernstein}
Let $F$ be a nonarchimedean local field whose residue field has $q$ elements, $x$ a uniformizer of $F$, $n\ge1$, and $A$ a commutative ring in which $q$ is a unit. Let $H_A$ be the Iwahori-Hecke algebra over $A$: the $A$-algebra of endomorphisms of $A[GL_n(F)/\mathrm{Iw}]^{\mathrm{Iw}}$ generated by the double coset operators $[\mathrm{Iw}\,g\,\mathrm{Iw}]$ (Definition~\ref{sha-hecke-ideal-tw-iwahori}). For $w\in S_n$ let $p_w$ be its permutation matrix and $T_w=[\mathrm{Iw}\,p_w\mathrm{Iw}]$; for $0\le m\le n$ let $t_m\in T(F)$ have $x$ in the first $m$ places and $1$ in the others.
\begin{enumerate}
\item There is an $A$-algebra map $\theta:A[X_1,\dots,X_n]\to H_A$ with $\theta(X_i)\,[\mathrm{Iw}\,t_{i-1}\mathrm{Iw}]=q^{i-1}\,[\mathrm{Iw}\,t_i\mathrm{Iw}]$ for $1\le i\le n$, such that for $w\in S_n$, $1\le a<n$ and $s=(a\;a{+}1)$: $T_1=1$; $T_{sw}=T_sT_w$ if $w^{-1}(a)<w^{-1}(a+1)$; $T_s^2=(q-1)T_s+q$; $T_s\theta(X_j)=\theta(X_j)T_s$ for $j\ne a,a+1$; $T_s\theta(X_a)=\theta(X_{a+1})T_s+(q-1)\theta(X_a)$; and $T_s\theta(X_{a+1})=\theta(X_a)T_s-(q-1)\theta(X_a)$.
\item Let $e=\sum_{w\in S_n}T_w$. If $y\in A[X_1,\dots,X_n]$ and $\theta(y)\,e=0$, then $y=0$.
\item For a representation $\rho$ of $GL_n(F)$ on an $A$-module $M$, there is an $A$-algebra map $\mathrm{ev}_\rho:H_A\to\operatorname{End}_A(M^{\mathrm{Iw}})$ that sends each $[\mathrm{Iw}\,g\,\mathrm{Iw}]$ to its double coset operator on $M^{\mathrm{Iw}}$.
\end{enumerate}
\end{lemma}
\leanref{Development/AutomorphicArithmetic/AutomorphicVarieties/AlgebraTools/HeckePairTop.lean}{UpperTheory.V.AlgebraTools.HkAlg}
\leanref{Development/AutomorphicArithmetic/AutomorphicVarieties/AlgebraTools/HkDefs.lean}{UpperTheory.V.AlgebraTools.Hk.Bern}
\leanref{Development/AutomorphicArithmetic/AutomorphicVarieties/AlgebraTools/IwBernHk.lean}{UpperTheory.V.AlgebraTools.bernHk}
\leanref{Development/AutomorphicArithmetic/AutomorphicVarieties/AlgebraTools/IwBernFaith.lean}{UpperTheory.V.AlgebraTools.bernFaith}
\leanref{Development/AutomorphicArithmetic/AutomorphicVarieties/AlgebraTools/IwHkEv.lean}{UpperTheory.V.AlgebraTools.hkEv}
\begin{proof}
(1) is the Bernstein presentation of the Iwahori-Hecke algebra of $GL_n(F)$, whose finite part is the presentation of Iwahori and Matsumoto \cite{IM65,Lus89}. (2) follows from the basis of $H_A$ over $A$ formed by the elements $\theta(X^\lambda)T_w$, for the monomials $X^\lambda$ with exponents in $\mathbf Z$ and $w\in S_n$ \cite{Lus89}, since $\theta(y)e=\sum_w\theta(y)T_w$. (3) is the action of the Hecke algebra on the vectors fixed by $\mathrm{Iw}$ \cite{Bor76}.
% bernHk constructs the presentation (the Lean's Hk.Bern, with th(X_i) * hT(pw x (fund n i)) = q^i hT(pw x (fund n (i+1))) for 0-based i); bernFaith is (2), Hk.Bern.Faithful; hkEv is (3).
\end{proof}

\begin{lemma}\label{sha-hecke-ideal-tw-identities}
Let $A$ be a domain, $n\ge1$, $q\in A$, $H$ an $A$-algebra, $T_w\in H$ for $w\in S_n$, and $\theta:A[X_1,\dots,X_n]\to H$ an $A$-algebra map, which satisfy the relations of Lemma~\ref{sha-hecke-ideal-tw-bernstein}(1) and the condition of Lemma~\ref{sha-hecke-ideal-tw-bernstein}(2), with $e=\sum_wT_w$. Then, with the notation of Definition~\ref{sha-hecke-ideal-tw-divided}:
\begin{enumerate}
\item $\theta(z)$ commutes with every $T_w$, for $z$ symmetric;
\item $\sum_\kappa\theta(a_\kappa)\,e\,\theta(a'_\kappa)=\theta(N_q)$ for every dual pair of families $(a_\kappa)$ and $(a'_\kappa)$;
\item $e\,\theta(y)\,e=\theta\big(\delta(N_qy)\big)\,e$ for every $y\in A[X_1,\dots,X_n]$.
\end{enumerate}
\end{lemma}
\leanref{Development/AutomorphicArithmetic/AutomorphicVarieties/AlgebraTools/HkBern.lean}{UpperTheory.V.AlgebraTools.hkBern}
\leanref{Development/AutomorphicArithmetic/AutomorphicVarieties/AlgebraTools/HkSum.lean}{UpperTheory.V.AlgebraTools.hkSum}
\leanref{Development/AutomorphicArithmetic/AutomorphicVarieties/AlgebraTools/HkSandwich.lean}{UpperTheory.V.AlgebraTools.hkSandwich}
\begin{proof}
(1) Let $s=(a\;a{+}1)$. By induction on polynomials, the relations give, for every $x$, a polynomial $y$ with $(X_a-X_{a+1})y=x-s(x)$ and $T_s\theta(x)=\theta(s(x))T_s+(q-1)\theta(X_ay)$. For $z$ symmetric, $(X_a-X_{a+1})y=0$, so $y=0$ and $T_s$ commutes with $\theta(z)$. Every $T_w$ is a product of the $T_s$, by the relation $T_{sw}=T_sT_w$ and induction on the length of $w$. (2) and (3) are computations in $H$ from the relations and the condition of Lemma~\ref{sha-hecke-ideal-tw-bernstein}(2).
% hkBern: main (by MvPolynomial.induction_on, with relA, relB and comm for the variables), then hkInd over the length.
\end{proof}

\begin{lemma}\label{sha-hecke-ideal-tw-trace}
Let $F$ be a nonarchimedean local field whose residue field has $q$ elements, $x$ a uniformizer of $F$, $n\ge1$ and $K=GL_n(\mathcal O_F)$. Let $O$ be a domain in which $q$ is a unit, $A$ an $O$-algebra, and $\rho$ a representation of $GL_n(F)$ on an $A$-module $M$ in which no nonzero element is killed by a nonzero integer. Let $V$ be an $A$-module with an injective $A$-linear map $\jmath:V\to M^{\mathrm{Iw}}$; let $\Theta$ be a homomorphism from $T(F)$ to the group of $A$-linear automorphisms of $V$ with $\jmath\Theta(t)=[\mathrm{Iw}\,t\,\mathrm{Iw}]\jmath$ for every dominant $t$; let $E$ be an $A$-linear endomorphism of $V$ with $\jmath E=[K\,1\,\mathrm{Iw}]\jmath$; and let $\varphi:O[X_1,\dots,X_n]\to\operatorname{End}_A(V)$ be a ring map that is multiplication on $O$ and sends $X_i$ to $q^{i-1}\Theta(e_i(x))$. Then, with the notation of Definition~\ref{sha-hecke-ideal-tw-divided} over $O$:
\begin{enumerate}
\item $E$ commutes with $\varphi(z)$ for every symmetric $z$;
\item $\sum_\kappa\varphi(a_\kappa)\,E\,\varphi(a'_\kappa)=\varphi(N_q)$ for every dual pair of families $(a_\kappa)$ and $(a'_\kappa)$;
\item $E\varphi(y)v=\varphi\big(\delta(N_qy)\big)v$ for every $y$ and every $v\in V$ with $\jmath(v)\in M^K$.
\end{enumerate}
\end{lemma}
\leanref{Development/AutomorphicArithmetic/AutomorphicVarieties/PatchingModules/NilpotentLevelLowering/HkHypOf.lean}{UpperTheory.V.Pa.Lev.hkHypOf}
\leanref{Development/AutomorphicArithmetic/AutomorphicVarieties/AlgebraTools/IwHkTrace.lean}{UpperTheory.V.AlgebraTools.hkTrace}
\begin{proof}
Regard $M$ as an $O$-module through $A$. Let $H_O$, $\theta$, $T_w$, $e$ and $t_m$ be as in Lemma~\ref{sha-hecke-ideal-tw-bernstein} for $F$, $x$ and $O$, and $\mathrm{ev}=\mathrm{ev}_\rho:H_O\to\operatorname{End}_O(M^{\mathrm{Iw}})$. By (1) and (2) of that lemma, Lemma~\ref{sha-hecke-ideal-tw-identities} applies to $H_O$.

\emph{The element $e$.} By the Bruhat decomposition, $K$ is the disjoint union of the double cosets $\mathrm{Iw}\,p_w\mathrm{Iw}$ \cite{IM65}. Hence $\sum_w[\mathrm{Iw}\,p_w\mathrm{Iw}]$ is $[K\,1\,\mathrm{Iw}]$ on $M^{\mathrm{Iw}}$, followed by the inclusion of $M^K$, and $[K\,1\,\mathrm{Iw}]m=[K:\mathrm{Iw}]\,m$ for $m\in M^K$. So $\mathrm{ev}(e)\jmath=\jmath E$.

\emph{The elements $\theta(X_i)$.} The $t_m$ are dominant, and $t_i=e_i(x)t_{i-1}$. Let $v\in V$, and write $v=\Theta(t_{i-1})u$. Then $\mathrm{ev}(\theta(X_i))\jmath(v)=\mathrm{ev}(\theta(X_i))[\mathrm{Iw}\,t_{i-1}\mathrm{Iw}]\jmath(u)=q^{i-1}[\mathrm{Iw}\,t_i\mathrm{Iw}]\jmath(u)=q^{i-1}\jmath\Theta(t_i)u=\jmath\varphi(X_i)v$. By induction on polynomials, $\mathrm{ev}(\theta(z))\jmath=\jmath\varphi(z)$ for every $z\in O[X_1,\dots,X_n]$.

\emph{(1) and (2).} They are the images under $\mathrm{ev}$ of (1) and (2) of Lemma~\ref{sha-hecke-ideal-tw-identities}, carried to $V$ through the injective map $\jmath$.

\emph{(3).} Let $v\in V$ with $\jmath(v)\in M^K$. Then $\jmath(Ev)=[K:\mathrm{Iw}]\jmath(v)$, so $Ev=[K:\mathrm{Iw}]v$. By (3) of Lemma~\ref{sha-hecke-ideal-tw-identities}, $E\varphi(y)Ev=\varphi(\delta(N_qy))Ev$, that is $[K:\mathrm{Iw}]\big(E\varphi(y)v-\varphi(\delta(N_qy))v\big)=0$. The index $[K:\mathrm{Iw}]$ is a nonzero integer, so (3) follows from the hypothesis on $M$ and the injectivity of $\jmath$.
% The private locA of HkHypOf.lean, through locO: bernHk, bernFaith, hkEv, hkTrace (with bruhatKO, heckePairIwC and iw_le), bern_ev (hkBern, hkSum, hkSandwich), ev_rel; then je, jTheta and jX (fund_split, t_m = pw x (fund n m), dom_fund), jphi by induction; (3) from dc_K, relIndex_ne and the torsion hypothesis tf.
\end{proof}

\begin{lemma}\label{sha-hecke-ideal-tw-nil}
Let $F_1,\dots,F_q$ be nonarchimedean local fields, $n\ge1$, $G=\prod_jGL_n(F_j)$, $C_j\subset k_{F_j}^\times$ subgroups and $K_C=\prod_j\mathrm{Iw}_{C_j}$ (Definition~\ref{sha-hecke-ideal-tw-iwahori}). Let $\pi$ be a representation of $G$ on a complex vector space $Y$, and $\Theta_j$ the homomorphisms of Lemma~\ref{sha-hecke-ideal-tw-torus} on $Y^{K_C}$. Fix $j$, and let $\sigma$ be an irreducible smooth representation of $GL_n(F_j)$ such that $Y$, as a representation of the $j$-th factor of $G$, is a sum of copies of $\sigma$; let $\Theta_\sigma$ be the homomorphism of Lemma~\ref{sha-hecke-ideal-tw-torus} for $\sigma$ and $\mathrm{Iw}_{C_j}$. Let $(t_a)_{a\in\Sigma}$ be a family of elements of $T(F_j)$ and $F\in\mathbf C[X_a:a\in\Sigma]$, and suppose that $F\big((\chi(t_a))_a\big)=0$ for every character $\chi$ of $T(F_j)$ for which the common generalized eigenspace of the $\Theta_\sigma(t)$ for the eigenvalues $\chi(t)$ is nonzero. Then $F\big((\Theta_j(t_a))_a\big)^{n!}=0$ on $Y^{K_C}$.
\end{lemma}
\leanref{Development/AutomorphicArithmetic/AutomorphicVarieties/PatchingModules/NilpotentLevelLowering/NilKill.lean}{UpperTheory.V.Pa.Lev.nilKill}
\begin{proof}
\emph{On $\sigma$.} By Lemma~\ref{sha-hecke-ideal-tw-dimension}, $W=\sigma^{\mathrm{Iw}_{C_j}}$ has dimension at most $n!$. Let $G_\sigma=F((\Theta_\sigma(t_a))_a)$, and suppose that $G_\sigma^{n!}\ne0$. Its image is a nonzero subspace of $W$ stable under the commuting automorphisms $\Theta_\sigma(t)$, so it contains a common eigenvector $v=G_\sigma^{n!}u\ne0$, with eigenvalues $\chi(t)$ for a character $\chi$ of $T(F_j)$. Then $\chi$ has a nonzero common generalized eigenspace, so $G_\sigma v=F((\chi(t_a))_a)v=0$. So $u$ lies in the kernel of $G_\sigma^{n!+1}$, which is that of $G_\sigma^{n!}$ as $\dim W\le n!$; hence $v=0$, a contradiction. Thus $G_\sigma^{n!}=0$.

\emph{On $Y$.} Let $\Theta_Y$ be the homomorphism of Lemma~\ref{sha-hecke-ideal-tw-torus} for $Y$, as a representation of the $j$-th factor, and $\mathrm{Iw}_{C_j}$. A vector of $Y^{\mathrm{Iw}_{C_j}}$ is a sum of images of vectors of $W$ under equivariant maps $f:\sigma\to Y$. Such an $f$ carries the double coset operators of the dominant elements on $W$ to those on $Y^{\mathrm{Iw}_{C_j}}$. Every $t\in T(F_j)$ is a quotient $(tb)b^{-1}$ of dominant elements $tb$ and $b$, so $f\Theta_\sigma(t)=\Theta_Y(t)f$ for all $t$, and $F((\Theta_Y(t_a))_a)^{n!}=0$ on $Y^{\mathrm{Iw}_{C_j}}$. Finally $Y^{K_C}\subset Y^{\mathrm{Iw}_{C_j}}$, and $\Theta_Y(t)$ restricts to $\Theta_j(t)$ on $Y^{K_C}$.
% nilKill: thetaEx for the j-th factor, nilKill_loc (thetaEx for sigma, jqLine for the dimension bound, nilKill_nil with nilKill_eig, isoFix for the sum of copies, nilKill_dc, and nilKill_split with nilKill_ext for t = (tb) b^{-1} with b = diag(c^(n-i))), then thetaLoc and nilKill_int on Y^{K_C}. It assumes the published GeometricLemma and BernsteinZelevinskyFiniteness, through jqLine.
\end{proof}

\begin{lemma}\label{sha-hecke-ideal-tw-scalar}
Let $F_1,\dots,F_q$ be nonarchimedean local fields, $n\ge1$, $G=\prod_jGL_n(F_j)$ and $\mathrm{Iw}_G=\prod_j\mathrm{Iw}$. Let $\pi$ be a representation of $G$ on a complex vector space $Y$, and $\Theta_j$ the homomorphisms of Lemma~\ref{sha-hecke-ideal-tw-torus} on $Y^{\mathrm{Iw}_G}$. Fix $j$, let $q_j$ be the number of elements of the residue field of $F_j$ and $x$ a uniformizer of $F_j$, and let $\sigma$ be an irreducible smooth representation of $GL_n(F_j)$ such that $Y$, as a representation of the $j$-th factor of $G$, is a sum of copies of $\sigma$; let $\Theta_\sigma$ be the homomorphism of Lemma~\ref{sha-hecke-ideal-tw-torus} for $\sigma$ and $\mathrm{Iw}$. Let $z\in\mathbf C[X_1,\dots,X_n]$ be symmetric and $c\in\mathbf C$, and suppose that $z(\alpha_1,\dots,\alpha_n)=c$, where $\alpha_i=q_j^{i-1}\chi(e_i(x))$, for every character $\chi$ of $T(F_j)$ for which the common generalized eigenspace of the $\Theta_\sigma(t)$ for the eigenvalues $\chi(t)$ is nonzero. Then $z\big((q_j^{i-1}\Theta_j(e_i(x)))_{1\le i\le n}\big)$ acts as $c$ on the vectors of $Y^{\mathrm{Iw}_G}$ fixed by $GL_n(\mathcal O_{F_j})$ in the $j$-th factor.
\end{lemma}
\leanref{Development/AutomorphicArithmetic/AutomorphicVarieties/PatchingModules/NilpotentLevelLowering/SphScalar.lean}{UpperTheory.V.Pa.Lev.sphScalar}
\leanref{Development/AutomorphicArithmetic/AutomorphicVarieties/AlgebraTools/IwSphEq.lean}{UpperTheory.V.AlgebraTools.sphEq}
\leanref{Development/AutomorphicArithmetic/AutomorphicVarieties/PatchingModules/NilpotentLevelLowering/SphLine.lean}{UpperTheory.V.Pa.Lev.sphLine}
\begin{proof}
Write $K_j=GL_n(\mathcal O_{F_j})$ and $\varphi(z)=z\big((q_j^{i-1}\Theta_j(e_i(x)))_i\big)$. Let $H_{\mathbf C}$, $\theta$, $T_w$ and $\mathrm{ev}_\tau$ be as in Lemma~\ref{sha-hecke-ideal-tw-bernstein} for $F_j$, $x$ and $\mathbf C$. For a representation $\tau$ of $GL_n(F_j)$, with the homomorphism $\Theta_\tau$ of Lemma~\ref{sha-hecke-ideal-tw-torus} for $\tau$ and $\mathrm{Iw}$, the computation in the proof of Lemma~\ref{sha-hecke-ideal-tw-trace} gives $\mathrm{ev}_\tau(\theta(X_i))=q_j^{i-1}\Theta_\tau(e_i(x))$.

\emph{A scalar on $\sigma$.} A vector $v\in\sigma^{\mathrm{Iw}}$ is fixed by $K_j$ if and only if $T_sv=q_jv$ for every simple transposition $s$. Indeed, $\mathrm{Iw}\,p_s\mathrm{Iw}$ is the union of $q_j$ cosets of $\mathrm{Iw}$, which gives one direction. Conversely, if $T_sv=q_jv$ for all simple $s$, then $T_wv=q_j^{\ell(w)}v$ by induction on the length $\ell(w)$, so $[K_j\,1\,\mathrm{Iw}]v=\sum_wT_wv=[K_j:\mathrm{Iw}]v$, which is fixed by $K_j$; and $[K_j:\mathrm{Iw}]$ is invertible in $\mathbf C$. By Lemma~\ref{sha-hecke-ideal-tw-identities}(1), $\theta(z)$ commutes with the $T_s$, so $\mathrm{ev}_\sigma(\theta(z))$ preserves $\sigma^{K_j}$. The space $\sigma^{K_j}$ has dimension at most one \cite{Car79}, so $\mathrm{ev}_\sigma(\theta(z))$ acts on it as a scalar $c'$.

\emph{The same scalar on $Y$.} Let $Y_j$ be $Y$ as a representation of the $j$-th factor. A vector of $Y_j^{K_j}$ is a sum of images of vectors of $\sigma^{K_j}$ under equivariant maps $f:\sigma\to Y_j$, and $f\circ\mathrm{ev}_\sigma(h)=\mathrm{ev}_{Y_j}(h)\circ f$ for $h\in H_{\mathbf C}$. So $\mathrm{ev}_{Y_j}(\theta(z))$ acts on $Y_j^{K_j}$ as $c'$. On $Y^{\mathrm{Iw}_G}\subset Y_j^{\mathrm{Iw}}$, $\Theta_{Y_j}$ restricts to $\Theta_j$, so $\varphi(z)v=c'v$ for every $v\in Y^{\mathrm{Iw}_G}$ fixed by $K_j$.

\emph{The scalar is $c$.} The polynomial $F=z(X_1,q_jX_2,\dots,q_j^{n-1}X_n)-c$ vanishes at $(\chi(e_i(x)))_i$ for every character $\chi$ as in the hypothesis. By Lemma~\ref{sha-hecke-ideal-tw-nil}, for the $C_{j'}=k_{F_{j'}}^\times$ and the family $(e_i(x))_i$, $(\varphi(z)-c)^{n!}=0$ on $Y^{\mathrm{Iw}_G}$. If $v\ne0$ is fixed by $K_j$, then $0=(\varphi(z)-c)^{n!}v=(c'-c)^{n!}v$, so $c'=c$; if $v=0$ there is nothing to show.
% sphScalar: sphScalar_loc (thetaEx; bernHk and hkEv; thetaTie for ev(theta(X_i)); sphScalar_gen with sphEq for the K-fixed vectors (hk_fixed and permCard one way, lenStep, permMul and hkTrace the other), hkBern, sphLine (from sphComm and BernsteinZelevinskyFiniteness) for the line, isoFix and sphScalar_ev for the sum of copies); sphScalar_tr with thetaLoc; sphScalar_fin and nilKill; sphScalar_nil.
\end{proof}

\begin{lemma}\label{sha-hecke-ideal-tw-symmetric}
In Situation~\ref{sha-hecke-situation}, let $(\lambda,M_0)$ be a Taylor-Wiles base for an ideal $J\subset B$, let $\Gamma'$ and places $w_1,\dots,w_q$ of $L$ not above $\ell$ be as in Definition~\ref{sha-hecke-ideal-tw-level-module}, with residue fields of $q_{w_j}$ elements, and let $(R,P)$ be a module at level $\Gamma'$ for the places $w_j$ that satisfies the local conditions for elements $x_j$, $\delta_j$, $\phi_j$, $\tau_j$ and a field isomorphism $\iota$ as in (4) of that definition. Let $\Theta^{\mathrm{Iw}}_j$ be the homomorphisms of Lemma~\ref{sha-hecke-ideal-tw-torus} on $P^{\mathrm{Iw}_w}$ and $U^w_{j,i}=q_{w_j}^{i-1}\Theta^{\mathrm{Iw}}_j(e_i(x_j))$, and let $\lambda_{j,k}\in R$ be the coefficients of the characteristic polynomial $X^n-\lambda_{j,1}X^{n-1}+\dots+(-1)^n\lambda_{j,n}$ of $D_R$ at $\phi_j$, with $\lambda_{j,0}=1$. Then $e_k(U^w_{j,1},\dots,U^w_{j,n})\,m=\lambda_{j,k}m$ for every $j$, every $0\le k\le n$ and every $m\in P^{K_w}$, where $e_k$ is the $k$-th elementary symmetric polynomial.
\end{lemma}
\leanref{Development/AutomorphicArithmetic/AutomorphicVarieties/ModularityArgument/ThresholdArguments/DenseThresholdCases/SymOfVieta.lean}{UpperTheory.V.Z.Win.T2.symOf_alt}
\begin{proof}
As the $w_j$ do not lie above $\ell$, the $q_{w_j}$ are units in $R$, so the $\Theta^{\mathrm{Iw}}_j$ are defined. Let $m\in P^{K_w}$ and $m'=e_k(U^w_{j,1},\dots,U^w_{j,n})m-\lambda_{j,k}m$. The module is separated by its points, so it suffices to show that $\operatorname{sp}_y(m')=0$ for every point $y$.

Fix $y$, let $\Theta^y_j$ be the homomorphisms of Lemma~\ref{sha-hecke-ideal-tw-torus} on $V_y^{\mathrm{Iw}_w}$, and $\varphi_y(z)=z\big((q_{w_j}^{i-1}\Theta^y_j(e_i(x_j)))_i\big)$ for $z\in\mathbf C[X_1,\dots,X_n]$. As in the proof of Lemma~\ref{sha-hecke-ideal-tw-kill}, $\operatorname{sp}_y$ carries $P^{\mathrm{Iw}_w}$ into $V_y^{\mathrm{Iw}_w}$ and intertwines $\Theta^{\mathrm{Iw}}_j(t)$ with $\Theta^y_j(t)$; with $\operatorname{sp}_y(tm)=\iota\theta_y(t)\operatorname{sp}_y(m)$, this gives $\operatorname{sp}_y(m')=\varphi_y(e_k)v-\iota\theta_y(\lambda_{j,k})v$ for $v=\operatorname{sp}_y(m)$. The vector $v$ is fixed by $GL_n(\mathcal O_{L_{w_j}})$ in the $j$-th factor. If $v=0$, there is nothing to show. Otherwise $\sigma_{y,j}$ has nonzero vectors fixed by $GL_n(\mathcal O_{L_{w_j}})$, since $V_y$ is a sum of copies of $\sigma_{y,j}$ for the $j$-th factor and its vectors fixed by this group are sums of images of such vectors of $\sigma_{y,j}$ under equivariant maps.

Apply Lemma~\ref{sha-hecke-ideal-tw-scalar} to $V_y$, $\sigma_{y,j}$, $x_j$, $z=e_k$ and $c=\iota\theta_y(\lambda_{j,k})$. Let $\chi$ be a character as there and $\alpha_i=q_{w_j}^{i-1}\chi(e_i(x_j))$. By the local conditions, for $C=k_{w_j}^\times$, conclusion (1) of Lemma~\ref{sha-hecke-ideal-tw-local} holds in its first case: the characteristic polynomial of $\iota\rho_y(\phi_j)$ is $\prod_i(X-\alpha_i)$. Since $\theta_y(D_R)$ is the determinant of $\rho_y$, the characteristic polynomial of $\rho_y(\phi_j)$ is the image under $\theta_y$ of that of $D_R$ at $\phi_j$, so $e_k(\alpha)=\iota\theta_y(\lambda_{j,k})$. The lemma gives $\varphi_y(e_k)v=\iota\theta_y(\lambda_{j,k})v$, so $\operatorname{sp}_y(m')=0$.
% symOf_alt (statement symOf_s_ty): thetaQ for y.pi with C = top, symOf_act, thetaMap (hTheta g), symOf_sym through symOf_int; hfix; the case g m = 0; symOf_sph (isoFix) for a spherical sigma_{y,j}; sphScalar with relT' (the first conclusion of LocRelT, Or.inl hsph) and symOf_vieta (with symOf_chi). The Lean takes the Theta of a given ThData; the hypotheses w_j not above l and IsGenMod delta_j are not used in it.
\end{proof}

\begin{lemma}\label{sha-hecke-ideal-tw-split}
Let $O$ be a domain, $R$ an $O$-algebra, $B$ a commutative ring, $\beta:R\to B$ a ring homomorphism, $q,n\ge1$ and $b_1,\dots,b_q$ integers. Let $V$ be an $R$-module with commuting $R$-linear endomorphisms $U_{j,i}$, for $1\le j\le q$ and $1\le i\le n$; for $z\in O[X_1,\dots,X_n]$ write $z(U_j)=z(U_{j,1},\dots,U_{j,n})$, and for $t$ in the polynomial ring $R[U_{j,i}]$ write $t(U)$ for the endomorphism obtained by substituting the $U_{j,i}$. Let $E_1,\dots,E_q$ be $R$-linear endomorphisms of $V$ and $W_1,\dots,W_q\subset V$ submodules such that, with the notation of Definition~\ref{sha-hecke-ideal-tw-divided} over $O$:
\begin{enumerate}
\item[(a)] the $E_j$ commute with each other, $E_j$ commutes with the $U_{j',i}$ for $j'\ne j$, and with $z(U_j)$ for $z$ symmetric;
\item[(b)] $\sum_\kappa a_\kappa(U_j)\,E_j\,a'_\kappa(U_j)=N_{b_j}(U_j)$ for every $j$ and every dual pair of families $(a_\kappa)$ and $(a'_\kappa)$;
\item[(c)] $W_j$ is stable under the $U_{j',i}$ and under $E_{j'}$, for $j'\ne j$;
\item[(d)] $E_j\big(y(U_j)v\big)=\delta(N_{b_j}y)(U_j)\,v$ for every $j$, $y\in O[X_1,\dots,X_n]$ and $v\in W_j$.
\end{enumerate}
Let $M$ be a $B$-module, $P_1,\dots,P_q\in B[X]$ monic of degree $n$, written $P_j=\sum_{k=0}^n(-1)^kp_{j,k}X^{n-k}$, and $B[U]$ the quotient of $B[U_{j,i}]$ by the coefficients of the polynomials $\prod_i(X-U_{j,i})-P_j$, with $\beta(t)$ the image in $B[U]$ of $t\in R[U_{j,i}]$. Let $\jmath:M\to V$ and $\operatorname{tr}:V\to M$ be additive maps, with $\jmath$ injective, such that $\jmath(\beta(t)m)=t\,\jmath(m)$ for $t\in R$, $\jmath\circ\operatorname{tr}=E_1\cdots E_q$, $\jmath(M)\subset W_j$ for every $j$, and $e_k(U_j)\,\jmath(m)=\jmath(p_{j,k}m)$ for every $j$, $0\le k\le n$ and $m\in M$. Let $f=\prod_j\prod_i\prod_{i'\ne i}(U_{j,i}-U_{j,i'})(U_{j,i}-b_jU_{j,i'})$. Then there are additive maps $\varphi_U:B[U]\otimes_BM\to V$ and $\psi_U:V\to B[U]\otimes_BM$ with $\varphi_U(\beta(t)z)=t(U)\varphi_U(z)$ and $\psi_U(t(U)v)=\beta(t)\psi_U(v)$ for $t\in R[U_{j,i}]$, such that $\varphi_U\psi_U=f(U)$ and $\psi_U\varphi_U$ is multiplication by $\beta(f)$.
\end{lemma}
\leanref{Development/AutomorphicArithmetic/AutomorphicVarieties/PatchingModules/NilpotentLevelLowering/MultiTwo.lean}{UpperTheory.V.Pa.Lev.multiTwo}
\leanref{Development/AutomorphicArithmetic/AutomorphicVarieties/PatchingModules/NilpotentLevelLowering/TwoMul.lean}{UpperTheory.V.Pa.Lev.twoMul}
\leanref{Development/AutomorphicArithmetic/AutomorphicVarieties/PatchingModules/NilpotentLevelLowering/ActUVariableEqPow.lean}{UpperTheory.V.Pa.Lev.HkHyp}
\begin{proof}
Let $f_N=\prod_jN_{b_j}(U_j)$.

\emph{The map $\varphi_U$.} For $b\in B$, a monomial $U^a$ in the $U_{j,i}$ and $m\in M$ put $G(bU^a,m)=U^a\jmath(bm)$, and extend $G$ additively to $B[U_{j,i}]\times M$. Then $G(\beta(t)h,m)=t(U)G(h,m)$ for $t\in R[U_{j,i}]$, and $G(bh,m)=G(h,bm)$ for $b\in B$. For $k\le n$ the coefficient of $X^k$ in $\prod_i(X-U_{j,i})-P_j$ is $g=(-1)^{n-k}\big(e_{n-k}(U_j)-p_{j,n-k}\big)$, and $G(g,m')=(-1)^{n-k}\big(e_{n-k}(U_j)\jmath(m')-\jmath(p_{j,n-k}m')\big)=0$; for $k>n$ the coefficient is $0$. By the two rules, $G(hg,m)=0$ for every $h$, every such $g$ and every $m$. So $G$ induces a $B$-balanced map on $B[U]\times M$, hence $\varphi_U$, with $\varphi_U(1\otimes m)=\jmath(m)$ and $\varphi_U(\beta(t)z)=t(U)\varphi_U(z)$.

\emph{Symmetric functions through $\operatorname{tr}$.} Since $\jmath$ is injective and $\jmath\operatorname{tr}=E_1\cdots E_q$, $\operatorname{tr}(tv)=\beta(t)\operatorname{tr}(v)$ for $t\in R$. For $z$ symmetric, $z(U_j)$ commutes with every $E_{j'}$ by (a), so $\jmath\operatorname{tr}(z(U_j)v)=z(U_j)\jmath\operatorname{tr}(v)$; for $z=e_k$ this is $\jmath(p_{j,k}\operatorname{tr}(v))$. In $B[U]$ the image of $e_k(U_j)$ is $p_{j,k}$, and the symmetric polynomials are polynomials in the $e_k$. Hence $y\otimes\operatorname{tr}(z(U_j)v)=y\beta(z(U_j))\otimes\operatorname{tr}(v)$ in $B[U]\otimes_BM$, for $y\in B[U]$ and $z$ symmetric.

\emph{The map $\psi_U$.} Let $(a_\kappa)_{\kappa\in\Sigma}$ and $(a'_\kappa)_{\kappa\in\Sigma}$ be dual families over $O$ (Lemma~\ref{sha-hecke-ideal-tw-dual}(1)). For a map $r:\{1,\dots,q\}\to\Sigma$ let $a_r=\prod_ja_{r(j)}(U_j)$ and $a'_r=\prod_ja'_{r(j)}(U_j)$, in $R[U_{j,i}]$, and put
\[
\psi_U(v)=\sum_r\beta(a_r)\otimes\operatorname{tr}\big(a'_r(U)v\big).
\]
It suffices to check $\psi_U(t(U)v)=\beta(t)\psi_U(v)$ for $t\in R$, which follows from the semilinearity of $\operatorname{tr}$, and for $t=U_{j,i}$. Fix $j$, let $x\in O[X_1,\dots,X_n]$, and for $s,s'\in R[U_{j,i}]$ let $\Phi(s,s')=\beta(s)\otimes\operatorname{tr}(s'(U)v)$. By Lemma~\ref{sha-hecke-ideal-tw-dual}(3), $a'_\kappa x=\sum_\mu a'_\mu\,\delta(a_\mu a'_\kappa x)$, and $\delta(a_\mu a'_\kappa x)$ is symmetric, so by the last step $\Phi(a_\kappa(U_j)s,a'_\kappa(U_j)x(U_j)s')=\sum_\mu\Phi\big(a_\kappa(U_j)\delta(a_\mu a'_\kappa x)(U_j)s,a'_\mu(U_j)s'\big)$. Summing over $\kappa$ and using $\sum_\kappa a_\kappa\delta(a'_\kappa xa_\mu)=xa_\mu$, again by Lemma~\ref{sha-hecke-ideal-tw-dual}(3),
\[
\sum_\kappa\Phi\big(a_\kappa(U_j)s,a'_\kappa(U_j)x(U_j)s'\big)=\sum_\mu\Phi\big(x(U_j)a_\mu(U_j)s,a'_\mu(U_j)s'\big).
\]
Applied to the factor of index $j$ of each term of $\psi_U$, with $x=X_i$, this gives $\psi_U(U_{j,i}v)=\beta(U_{j,i})\psi_U(v)$.

\emph{The composite $\varphi_U\psi_U$.} $\varphi_U\psi_U(v)=\sum_ra_r(U)\,E_1\cdots E_q\,a'_r(U)v$. By (a), $E_j$ commutes with the polynomials in the $U_{j',i}$, $j'\ne j$, so the sum is $\prod_j\sum_\kappa a_\kappa(U_j)E_ja'_\kappa(U_j)=f_N(U)$, by (b).

\emph{The composite $\psi_U\varphi_U$.} Let $m\in M$; then $\jmath(m)\in W_j$ for every $j$. Using (a) to move each $E_j$ past the factors of the other indices and (c) to stay in $W_j$, (d) gives $E_1\cdots E_q\,a'_r(U)\jmath(m)=\prod_j\delta(N_{b_j}a'_{r(j)})(U_j)\,\jmath(m)$. The $\delta(N_{b_j}a'_{r(j)})$ are symmetric, so this is $\jmath(c_rm)$ for an element $c_r\in B$ that is the image of $\prod_j\delta(N_{b_j}a'_{r(j)})(U_j)$ in $B[U]$. Hence
\[
\psi_U(\jmath(m))=\sum_r\beta(a_r)c_r\otimes m=\beta\Big(\prod_j\sum_\kappa a_\kappa\,\delta(a'_\kappa N_{b_j})\,(U_j)\Big)\otimes m=\beta(f_N)\otimes m,
\]
by Lemma~\ref{sha-hecke-ideal-tw-dual}(3). The elements $\beta(t)(1\otimes m)$ span $B[U]\otimes_BM$, and $\psi_U\varphi_U(\beta(t)(1\otimes m))=\psi_U(t(U)\jmath(m))=\beta(t)\beta(f_N)(1\otimes m)$.

\emph{The element $f$.} For $i<i'$ the factor $U_{j,i'}-b_jU_{j,i}$ of $f_N$ is the second factor of the term of $f$ for the pair $(i',i)$, so $f=f_Ng$ for some $g\in R[U_{j,i}]$. Replacing $\psi_U$ by $\psi_U\circ g(U)$ gives maps for $f$.
% multiTwo: multiPhi (exists_mvPolynomial_action_extending_inc, prod_coeff, hgen, the quotient by mSplitIdeal and TensorProduct.liftAddHom), hkDualEx, htrS and htrE (comm_sym), multiPsi (bal through esymmAlgHom_surjective and esymm_root; one_place and many_places with hkDualExp and hkDdProps), hphipsi (sum_all with sum_list and sum_ofFn_prod, and H.sum), hpsiinc (tr_list with H.tr, comm_list, mem_list; hTsym; hpoly from hkDualExp), and the induction over the tensor product. twoMul: prod_pairs_dvd and twoMaps_of_dvd. The hypotheses (a) to (d) are the Lean's HkHyp.
\end{proof}

\begin{lemma}\label{sha-hecke-ideal-tw-comparison}
In Situation~\ref{sha-hecke-situation}, let $(\lambda,M_0)$ be a Taylor-Wiles base for an ideal $J\subset B$, let $\Gamma'$ and places $w_1,\dots,w_q$ of $L$ not above $\ell$ be as in Definition~\ref{sha-hecke-ideal-tw-level-module}, with residue fields of $q_{w_j}$ elements, and let $(R,P)$ be a module at level $\Gamma'$ for the places $w_j$. Let $x_j$, $\delta_j$, $\phi_j$, $\tau_j$ and $\iota$ be as in (4) of that definition, and suppose that $(R,P)$ satisfies the local conditions for them. Let $\Theta^{\mathrm{Iw}}_j$ be the homomorphisms of Lemma~\ref{sha-hecke-ideal-tw-torus} on $P^{\mathrm{Iw}_w}$ and $U^w_{j,i}=q_{w_j}^{i-1}\Theta^{\mathrm{Iw}}_j(e_i(x_j))$. Then there are maps $\varphi_U$ and $\psi_U$ between $B[U]\otimes_BM_0$ and $M_w=P^{\mathrm{Iw}_w}$ as in Definition~\ref{sha-hecke-ideal-tw-operator-module}(2)(a), for $b_j=q_{w_j}$ and the $\phi_j$.
\end{lemma}
\leanref{Development/AutomorphicArithmetic/AutomorphicVarieties/ModularityArgument/ThresholdArguments/DenseThresholdCases/LevTwoOf.lean}{UpperTheory.V.Pa.Lev.twoOf_alt}
\begin{proof}
\emph{Units and torsion.} As $w_j\nmid\ell$, the number $q_{w_j}$ is prime to $\ell$, so it is a unit in $\mathcal O$, whose residue field has characteristic $\ell$, and in $R$. If $cm=0$ for a nonzero integer $c$ and $m\in P$, then $c\operatorname{sp}_y(m)=0$ in the complex vector space $V_y$, so $\operatorname{sp}_y(m)=0$, for every point $y$; as the module is separated by its points, $m=0$.

\emph{The operators $E_j$.} Let $K_j=GL_n(\mathcal O_{L_{w_j}})$, $V=P^{\mathrm{Iw}_w}$, and $E_j$ the restriction to $V$ of the double coset operator $[K_j\,1\,\mathrm{Iw}]$ of the $j$-th factor of $G_w$. It maps $V$ into $V$: its values are fixed by $K_j\supset\mathrm{Iw}$ in the $j$-th factor, and it commutes with the action of the other factors. Let $W_j\subset V$ be the submodule of the vectors fixed by $K_j$ in the $j$-th factor. Lemma~\ref{sha-hecke-ideal-tw-trace}, applied to $L_{w_j}$, $x_j$, $O=\mathcal O$, $A=R$, $P$ as a representation of the $j$-th factor, $V$ with its inclusion, $\Theta^{\mathrm{Iw}}_j$, $E_j$ and $\varphi(z)=z(U^w_{j,1},\dots,U^w_{j,n})$, gives (b), (d) and the last part of (a) of Lemma~\ref{sha-hecke-ideal-tw-split}, for the $U^w_{j,i}$ and $b_j=q_{w_j}$. The $E_j$ commute with each other, as the factors of $G_w$ commute. For $j'\ne j$, $E_j$ commutes with the double coset operators of the $j'$-th factor, so with $\Theta^{\mathrm{Iw}}_{j'}(t)$ for dominant $t$; hence with $\Theta^{\mathrm{Iw}}_{j'}(e_i(x_{j'}))=\Theta^{\mathrm{Iw}}_{j'}(t_i)\Theta^{\mathrm{Iw}}_{j'}(t_{i-1})^{-1}$, for the dominant $t_m$ of Lemma~\ref{sha-hecke-ideal-tw-bernstein} built from $x_{j'}$, and with the $U^w_{j',i}$. This gives (a). An endomorphism $h$ of $V$ that commutes with $E_j$ preserves $W_j$: for $v\in W_j$, $E_jv=[K_j:\mathrm{Iw}]v$, so $[K_j:\mathrm{Iw}]h(v)=E_jh(v)$ is fixed by $K_j$, and so is $h(v)$, by the first step. This gives (c).

\emph{The product of the $E_j$.} $E_1\cdots E_q$ maps $V$ into $P^{K_w}$: for each $j$, $E_j$ maps $V$ into $W_j$, and the $E_{j'}$, $j'\ne j$, preserve $W_j$.

\emph{The maps.} Let $\epsilon:P^{K_w}\to M_0$ be the isomorphism of the module $(R,P)$, $\jmath:M_0\to V$ the inclusion of $P^{K_w}$ composed with $\epsilon^{-1}$, and $\operatorname{tr}=\epsilon\circ E_1\cdots E_q$. Then $\jmath$ is injective, $\jmath(\beta(t)m)=t\jmath(m)$ for $t\in R$, $\jmath\circ\operatorname{tr}=E_1\cdots E_q$, and $\jmath(M_0)=P^{K_w}\subset W_j$ for every $j$. With $\lambda_{j,k}$ as in Lemma~\ref{sha-hecke-ideal-tw-symmetric}, that lemma gives $e_k(U^w_{j,1},\dots,U^w_{j,n})\jmath(m)=\lambda_{j,k}\jmath(m)=\jmath(\beta(\lambda_{j,k})m)$. As $\beta$ carries $D_R$ to $D^{\mathrm u}\circ p$, the $\beta(\lambda_{j,k})$ are the coefficients, with the signs of Lemma~\ref{sha-hecke-ideal-tw-split}, of the characteristic polynomial $P_j$ of $D^{\mathrm u}$ at $p(\phi_j)$, which is monic of degree $n$. That lemma, for $O=\mathcal O$, $\beta$, $M=M_0$, $V$, the $U^w_{j,i}$, the $E_j$ and $W_j$, $b_j=q_{w_j}$, $\jmath$ and $\operatorname{tr}$, gives $\varphi_U$ and $\psi_U$.
% twoOf_alt: unit_card (resCard_eq_normV and CharP of the residue field of O), the torsion hypothesis from hiota, hkHypOf (e_j = dc of the j-th factor restricted by pres; e_e and e_X with comm_dd and comm_unit; locA for esym, sum and tr; fixed_vectors_stable_under_hecke_action for WX and We; e_prod with lprod_mem), symOf_alt, and two_gen (inc from fix_le and eM, tr = eM o t0, fix_at for hW, hc from lam_map, beta_pair and lam_comap, chi_monic) through multiTwo and twoMul. The hypotheses IsGenMod delta_j are not used.
\end{proof}

\begin{lemma}\label{sha-hecke-ideal-tw-operators}
In Situation~\ref{sha-hecke-situation}, let $(\lambda,M_0)$ be a Taylor-Wiles base for an ideal $J\subset B$, let $\Gamma'$ and places $w_1,\dots,w_q$ of $L$ not above $\ell$ be as in Definition~\ref{sha-hecke-ideal-tw-level-module}, with residue fields of $q_{w_j}$ elements, and let $(R,P)$ be a module at level $\Gamma'$ for the places $w_j$ that satisfies the local conditions for elements $x_j$, $\delta_j$, $\phi_j$, $\tau_j$ and a field isomorphism $\iota$ as in (4) of that definition. Let $N\ge0$ and suppose $q_{w_j}\equiv1\pmod{\ell^N}$ for every $j$. Then there is an operator module of level $N$ for $\Gamma'$ (Definition~\ref{sha-hecke-ideal-tw-operator-module}) that satisfies the relations of exponent $n!$ for $b_j=q_{w_j}$ and the $\phi_j$ and $\tau_j$, with: the ring, pair and maps of the module $(R,P)$; $M_1=P^{\mathrm{Iw}_{\ell',w}}$ and $M_w=P^{\mathrm{Iw}_w}$; $U_{j,i}=q_{w_j}^{i-1}\Theta^1_j(e_i(x_j))$, $Y_{j,i}=\Theta^1_j(e_i(\delta_j))-1$ and $U^w_{j,i}=q_{w_j}^{i-1}\Theta^{\mathrm{Iw}}_j(e_i(x_j))$, for the homomorphisms $\Theta^1_j$ and $\Theta^{\mathrm{Iw}}_j$ of Lemma~\ref{sha-hecke-ideal-tw-torus} on $P^{\mathrm{Iw}_{\ell',w}}$ and $P^{\mathrm{Iw}_w}$; $a_j$ the exponent of $\ell$ in $q_{w_j}-1$; and the double coset operator $\mathcal N=[\mathrm{Iw}_w\,1\,\mathrm{Iw}_{\ell',w}]$.
\end{lemma}
\leanref{Development/AutomorphicArithmetic/AutomorphicVarieties/ModularityArgument/ThresholdArguments/DenseThresholdCases/LevOpOf.lean}{UpperTheory.V.Pa.Lev.opOf_alt}
\leanref{Development/AutomorphicArithmetic/AutomorphicVarieties/ModularityArgument/ThresholdArguments/DenseThresholdCases/LocOkLevelModuleOfThetaData.lean}{UpperTheory.V.Z.Win.T2.level_relation_one_two}
\begin{proof}
\emph{The operators.} Since $w_j\nmid\ell$, $q_{w_j}$ is prime to $\ell$, so it is a unit in $k$, hence in $\mathcal O$ and in $R$. Lemma~\ref{sha-hecke-ideal-tw-torus} gives the $\Theta^1_j$ and $\Theta^{\mathrm{Iw}}_j$, and the $U_{j,i}$, $Y_{j,i}$ commute, as do the $U^w_{j,i}$. $M_1$ is finite over $\Lambda_D$ by Definition~\ref{sha-hecke-ideal-tw-level-module}(1).

\emph{The group $\Delta$.} By Lemma~\ref{sha-hecke-ideal-tw-diamond}(1), $\Delta\cong\mathrm{Iw}_w/\mathrm{Iw}_{\ell',w}$, and $\Delta$ acts on $M_1$ through this isomorphism; $M_1$ is projective over $\Lambda_D[\Delta]$, being projective over the group ring of $\mathrm{Iw}_w/\mathrm{Iw}_{\ell',w}$. By Lemma~\ref{sha-hecke-ideal-tw-diamond}(2), the generator of the factor $(j,i)$ acts as $e_i(\delta_j)$, that is as $\Theta^1_j(e_i(\delta_j))=1+Y_{j,i}$. As $q_{w_j}\equiv1\pmod{\ell^N}$, $a_j\ge N$.

\emph{The map $\mathcal N$.} By Lemma~\ref{sha-hecke-ideal-tw-diamond}(2), $\mathcal N\circ U_{j,i}=U^w_{j,i}\circ\mathcal N$. For a finite group $Q$ and a projective $\Lambda_D[Q]$-module $M$, every $Q$-invariant element of $M$ is $\sum_{g\in Q}gm'$ for some $m'$, and an element $m$ with $\sum_{g\in Q}gm=0$ lies in $\sum_{s\in\Sigma}(s-1)M$ for every generating set $\Sigma$ of $Q$ (Higman's criterion). Apply this to $Q=\mathrm{Iw}_w/\mathrm{Iw}_{\ell',w}$ and $M=M_1$, whose $Q$-invariants are $P^{\mathrm{Iw}_w}$, with $\Sigma$ the images of the generators of $\Delta$, which act as $1+Y_{j,i}$. So $\mathcal N$ is surjective, and its kernel is contained in $\sum_{j,i}Y_{j,i}M_1$; conversely $\mathcal N\big((1+Y_{j,i})m\big)=\mathcal N(m)$, as the sum defining $\mathcal N$ is invariant under translation by $Q$.

\emph{Relation (a).} This is Lemma~\ref{sha-hecke-ideal-tw-comparison}.

\emph{Relations (b).} Fix $j$, $i$ and $h\in\Gamma'_0$, and write $f=f_jf'_j$, where $f_j=\prod_k\prod_{k'\ne k}(U_{j,k}-U_{j,k'})(U_{j,k}-q_{w_j}U_{j,k'})$ is the factor of index $j$. It suffices to show that $F^{n!}$ acts as $0$ for $F=f_j\operatorname{tr}\big(h\prod_k(\phi_j-U_{j,k})\big)$ and for $F=f_j\operatorname{tr}\big(h\prod_{k\ne i}(\phi_j-U_{j,k})(\tau_j-1-Y_{j,i})\big)$, polynomials in the variables of index $j$. We apply Lemma~\ref{sha-hecke-ideal-tw-kill}. Let $y$ be a point and $\chi$ a character as there, and put $\alpha_k=q_{w_j}^{k-1}\chi(e_k(x_j))$. The image of $f_j$ is $\prod_k\prod_{k'\ne k}(\alpha_k-\alpha_{k'})(\alpha_k-q_{w_j}\alpha_{k'})$. If it is $0$, so is the image of $F$. Otherwise $\alpha$ is generic. Since $\iota\theta_y(D_R)$ is the determinant of $\iota\rho_y$, the image of $\operatorname{tr}(z)$ is the trace of the matrix $\iota\rho_y(z)$, for $z$ in the group ring. For the first $F$, conclusion (1) of Lemma~\ref{sha-hecke-ideal-tw-local} makes $\prod_k(X-\alpha_k)$ the characteristic polynomial of $\iota\rho_y(\phi_j)$, so $\prod_k(\iota\rho_y(\phi_j)-\alpha_k)=0$ by the Cayley-Hamilton theorem, and the image of $F$ is the trace of $\iota\rho_y(h)\cdot0$. For the second, $y_i$ maps to $\chi(e_i(\delta_j))-1$, so $\tau_j-1-Y_{j,i}$ maps to $\iota\rho_y(\tau_j)-\chi(e_i(\delta_j))$, and conclusion (2) makes the matrix $\prod_{k\ne i}(\iota\rho_y(\phi_j)-\alpha_k)(\iota\rho_y(\tau_j)-\chi(e_i(\delta_j)))$ vanish. In both cases the image of $F$ is $0$, so $F^{n!}$ acts as $0$ by Lemma~\ref{sha-hecke-ideal-tw-kill}, and so does $(f\operatorname{tr}(\dots))^{n!}=f'^{\,n!}_jF^{n!}$.
% opOf_alt: opOf_cop and hq (q_w a unit), thEx (ThData from thetaQ for C = ellC and C = top, with act_uVariable_eq_pow, act_yVariable, actw_u), A.proj and opOf_proj (projectivity transported along delIso), delIso, trOf (the y-relation, the norm map hk2 and its commutation with U), opOf_tr with projNorm (Higman) and opOf_gen, ha (padicValNat l (q - 1) >= N), level_relation_one_two (reduction_of_polynomial_element, relOf_alt, trace_vanishing_of_evaluation, rep_first_relation_eq_zero and rep_second_relation_eq_zero, with relT' for the local relation LocRelT), twoOf_alt.
\end{proof}

\begin{lemma}\label{sha-hecke-ideal-tw-nilpotent}
Let $\Lambda$ be a complete noetherian local ring, $R$ a $\Lambda$-algebra, $M$ an $R$-module that is finite over $\Lambda$, and $(T_s)_{s\in\Sigma}$ a finite family of commuting $R$-linear endomorphisms of $M$. Let $M^{\mathrm{nil}}\subset M$ be the set of $m$ such that for every $s$ and every $k\ge0$, $T_s^dm\in\mathfrak m_\Lambda^kM$ for all large $d$. Then there is an idempotent $R$-linear endomorphism $\pi$ of $M$ with image $M^{\mathrm{nil}}$ that commutes with every $\Lambda$-linear endomorphism of $M$ commuting with the $T_s$. The $T_s$ preserve $M^{\mathrm{nil}}$, which is finite over $\Lambda$ and is a module over the power series ring $R[[Z_s:s\in\Sigma]]$ in which $R$ acts as on $M$ and $Z_s$ acts as $T_s$.
\end{lemma}
\leanref{Development/AutomorphicArithmetic/AutomorphicVarieties/AlgebraTools/TnIdem.lean}{UpperTheory.V.AlgebraTools.exists_idem_topNil}
\leanref{Development/AutomorphicArithmetic/AutomorphicVarieties/PatchingModules/NilpotentLevelLowering/OpParts.lean}{UpperTheory.V.Pa.Lev.OpP.nilPart}
\begin{proof}
The $\Lambda$-subalgebra of $\operatorname{End}_\Lambda M$ generated by the $T_s$ is commutative and finite over $\Lambda$, and it contains an idempotent $\pi$ with image $M^{\mathrm{nil}}$. An endomorphism that commutes with the $T_s$ commutes with this subalgebra, so with $\pi$; in particular $\pi$ is $R$-linear. The $T_s$ act topologically nilpotently on the finite $\Lambda$-module $M^{\mathrm{nil}}$, which gives the action of the power series ring.
% nilPart: exists_idem_topNil (topNil_eq_idem in the finite commutative algebra generated by the T_s), hcomm, heR, heT, the submodule Ns = range e, htop and psModule for the action of MvPowerSeries.
\end{proof}

\begin{lemma}\label{sha-hecke-ideal-tw-free}
Let $\ell$ be a prime and $\Lambda$ a complete noetherian local ring with $\ell\in\mathfrak m_\Lambda$. Let $a_1,\dots,a_q\ge0$, $n\ge1$, $\Delta=\prod_{j,i}\mathbf Z/\ell^{a_j}$ over $1\le j\le q$ and $1\le i\le n$, and $M$ a $\Lambda$-module, finite over $\Lambda$, with a $\Lambda$-linear action of $\Delta$ that makes it projective over $\Lambda[\Delta]$. Let $M=N\oplus C$ with $N$ and $C$ stable under $\Delta$, and suppose that the power series ring $\Lambda[[y_{j,i}]]$ acts $\Lambda$-linearly on $N$ with $1+y_{j,i}$ acting as the generator of the factor $(j,i)$ of $\Delta$. Then $N\cong\big(\Lambda[[y_{j,i}]]/((1+y_{j,i})^{\ell^{a_j}}-1)_{j,i}\big)^r$ as $\Lambda[[y_{j,i}]]$-modules, for some $r\ge0$.
\end{lemma}
\leanref{Development/AutomorphicArithmetic/AutomorphicVarieties/PatchingModules/NilpotentLevelLowering/FreeSummand.lean}{UpperTheory.V.Pa.Lev.freeSummand}
\begin{proof}
Sending $y_{j,i}$ to the generator of the factor $(j,i)$ minus $1$ gives a surjection $\Lambda[[y_{j,i}]]\to\Lambda[\Delta]$ with kernel generated by the $(1+y_{j,i})^{\ell^{a_j}}-1$, and $\Lambda[\Delta]$ is local since $\ell\in\mathfrak m_\Lambda$. The summand $N$ is finite and projective over the local ring $\Lambda[\Delta]$, hence free. The two actions of $\Lambda[[y_{j,i}]]$ on $N$, the given one and the one through $\Lambda[\Delta]$, agree: the elements on which they agree form a subalgebra that contains the $1+y_{j,i}$ and so maps onto $\Lambda[\Delta]$, and the given action kills the $(1+y_{j,i})^{\ell^{a_j}}-1$, which generate the kernel.
% freeSummand: exists_groupRing_surj and groupRing_ker (the surjection and its kernel, with the local ring), fs_abs (a finite projective direct summand over the local group ring is free), fsAgree, hsurjA, hkill and hagree for the agreement of the two actions. freePart applies it to the range of i with complement ker pr.
\end{proof}

\begin{lemma}\label{sha-hecke-ideal-tw-quotient}
Let $\Lambda$ be a complete noetherian local ring, $R$ a $\Lambda$-algebra, $\Sigma$ a finite set and $\tilde u_s\in R$ for $s\in\Sigma$. Let $M$ be an $R$-module that is finite over $\Lambda$, with commuting $R$-linear endomorphisms $U_s$ and $Y_s$, and $W$ an $R$-module with commuting $R$-linear endomorphisms $U^w_s$. Let $\mathcal N:M\to W$ be a surjective $R$-linear map with kernel $\sum_sY_sM$ and $\mathcal N\circ U_s=U^w_s\circ\mathcal N$. Let $M^{\mathrm{nil}}$ and $\pi$ be as in Lemma~\ref{sha-hecke-ideal-tw-nilpotent} for the endomorphisms $U_s-\tilde u_s$ and $Y_s$ of $M$, and $W^{\mathrm{nil}}$ as there for the $U^w_s-\tilde u_s$ on $W$. Then $\mathcal N$ induces an isomorphism $\jmath$ from $M^{\mathrm{nil}}/\sum_sY_sM^{\mathrm{nil}}$ onto $W^{\mathrm{nil}}$, carrying $U_s$ to $U^w_s$, and there is an additive map $\kappa:W\to M^{\mathrm{nil}}/\sum_sY_sM^{\mathrm{nil}}$ with $\kappa\circ\jmath=1$ and $\kappa\circ U^w_s=U_s\circ\kappa$.
\end{lemma}
\leanref{Development/AutomorphicArithmetic/AutomorphicVarieties/PatchingModules/NilpotentLevelLowering/OpParts.lean}{UpperTheory.V.Pa.Lev.OpP.quotPart}
\begin{proof}
The map $\mathcal N$ kills $Y_sM$, so it induces $\jmath$ on $M^{\mathrm{nil}}/\sum_sY_sM^{\mathrm{nil}}$. The projection $\pi$ commutes with the $Y_s$, so it carries $\ker\mathcal N=\sum_sY_sM$ into $\sum_sY_sM^{\mathrm{nil}}$; hence $\mathcal N(m)\mapsto\pi(m)$ is a well defined map $\kappa$, and $\kappa\circ\jmath=1$ because $\pi$ is the identity on $M^{\mathrm{nil}}$. Both maps commute with the $U$'s because $\mathcal N$ and $\pi$ do. Finally, the image of $M^{\mathrm{nil}}$ under the surjective map $\mathcal N$, which carries $U_s-\tilde u_s$ to $U^w_s-\tilde u_s$ and $Y_s$ to $0$, is $W^{\mathrm{nil}}$.
% quotPart: hkill, htr, hJkill, hprJ, exLift for j and kappa, topNil_map for the image.
\end{proof}

\begin{lemma}\label{sha-hecke-ideal-tw-splitting}
Let $\Lambda\to B$ be a local homomorphism of complete noetherian local rings, $M_0$ a $B$-module that is finite over $\Lambda$, and $P_1,\dots,P_q\in B[X]$ monic polynomials of degree $n$. Let $B[U]$ be the quotient of $B[U_{j,i}]$ by the coefficients of the $\prod_i(X-U_{j,i})-P_j$, let $\tilde u_{j,i}\in B$, and let $B[u]$ be the quotient of $B[[y_{j,i}]]$ by the coefficients of the $\prod_i(X-u_{j,i})-P_j$, where $u_{j,i}=\tilde u_{j,i}+y_{j,i}$. Let $R$ be a $\Lambda$-algebra with a homomorphism $\beta:R\to B$ of $\Lambda$-algebras, $\tilde u'_{j,i}\in R$ with $\beta(\tilde u'_{j,i})=\tilde u_{j,i}$, $W$ an $R$-module with commuting $R$-linear endomorphisms $U^w_{j,i}$, and $f\in R[U_{j,i}]$. Suppose that there are additive maps $\varphi_U:B[U]\otimes_BM_0\to W$ and $\psi_U:W\to B[U]\otimes_BM_0$ with $\varphi_U(\beta(t)z)=t(U^w)\varphi_U(z)$, $\psi_U(t(U^w)v)=\beta(t)\psi_U(v)$, $\varphi_U\psi_U=f(U^w)$ and $\psi_U\varphi_U=\beta(f)$, in the notation of Definition~\ref{sha-hecke-ideal-tw-operator-module}(2). Let $\eta:R[U_{j,i}]\to B[u]$ be $\beta$ on $R$ and send $U_{j,i}$ to $u_{j,i}$. Then $B[u]\otimes_BM_0$ is finite over $\Lambda$, the $u_{j,i}-\tilde u_{j,i}$ act on it topologically nilpotently, and there are additive maps $\varphi_u:B[u]\otimes_BM_0\to W$ and $\psi_u:W\to B[u]\otimes_BM_0$ with $\varphi_u(\eta(t)z)=t(U^w)\varphi_u(z)$, $\psi_u(t(U^w)v)=\eta(t)\psi_u(v)$ and $\psi_u\varphi_u=\eta(f)$, such that $\varphi_u$ takes values in $W^{\mathrm{nil}}$, for the endomorphisms $U^w_{j,i}-\tilde u'_{j,i}$ as in Lemma~\ref{sha-hecke-ideal-tw-nilpotent}, and $\varphi_u\psi_u=f(U^w)$ on $W^{\mathrm{nil}}$.
\end{lemma}
\leanref{Development/AutomorphicArithmetic/AutomorphicVarieties/PatchingModules/NilpotentLevelLowering/SpLocIdem.lean}{UpperTheory.V.Pa.Lev.spLoc}
\leanref{Development/AutomorphicArithmetic/AutomorphicVarieties/PatchingModules/NilpotentLevelLowering/OpParts.lean}{UpperTheory.V.Pa.Lev.OpP.spPart}
\begin{proof}
The $B$-algebra map $\varpi_u:B[U]\to B[u]$ sending $U_{j,i}$ to $u_{j,i}$ is surjective, and its kernel is generated by $1-\epsilon$ for an idempotent $\epsilon$ of $B[U]$; on $B[U]\otimes_BM_0$, $\epsilon$ acts as the identity on the elements on which the $U_{j,i}-\tilde u_{j,i}$ act topologically nilpotently, and only on those. So $\varpi_u\otimes1$ restricts to an isomorphism from $\epsilon(B[U]\otimes_BM_0)$ onto $B[u]\otimes_BM_0$; let $s$ be its inverse, and put $\varphi_u=\varphi_U\circ s$ and $\psi_u=(\varpi_u\otimes1)\circ\psi_U$. As $B[U]$ is finite over $B$, so is $B[u]$, and $B[u]\otimes_BM_0$ is finite over $\Lambda$. The compatibilities follow from those of $\varphi_U$ and $\psi_U$, as $\eta=\varpi_u\circ\beta$. The elements $s(z)$ are killed by high powers of the $U_{j,i}-\tilde u_{j,i}$ modulo $\mathfrak m_\Lambda^k$, so $\varphi_u(z)=\varphi_U(s(z))\in W^{\mathrm{nil}}$. If $v\in W^{\mathrm{nil}}$, then $\psi_U(v)$ is fixed by $\epsilon$, so $\varphi_u\psi_u(v)=\varphi_U\psi_U(v)=f(U^w)v$.
% spPart: spLoc (mSplitAlg_basis, MSplitAlg.lift with Core.spRoots, spLoc_alg for the surjection and the idempotent, topNil_eq_idem and spLoc_top for the characterization), finTens, tensSec for the section iota, nilMap for the nilpotence statements.
\end{proof}

\begin{lemma}\label{sha-hecke-ideal-tw-patching}
In Situation~\ref{sha-hecke-situation}, let $(\lambda,M_0)$ be a Taylor-Wiles base for an ideal $J\subset B$, $\Gamma'$, $p$ and $\Gamma'_0$ as in Definition~\ref{sha-hecke-ideal-tw-level-module}, $q,N\ge0$, $b_1,\dots,b_q$ integers, $\phi_j,\tau_j\in\Gamma'_0$, $e\ge0$, and an operator module of level $N$ for $\Gamma'$ that satisfies the relations of exponent $e$ for the $b_j$, $\phi_j$ and $\tau_j$ (Definition~\ref{sha-hecke-ideal-tw-operator-module}), with integers $a_j$. Let $\tilde u_{j,i}\in\mathcal O$, for $1\le j\le q$ and $1\le i\le n$, let $S=\Lambda_D[[y_{j,i}]]$, and let $B[u]$ be the quotient of $B[[y_{j,i}]]$ by the coefficients of the polynomials $\prod_i(X-u_{j,i})-P_j$, where $u_{j,i}=\tilde u_{j,i}+y_{j,i}$ and $P_j$ is the characteristic polynomial of $D^{\mathrm u}$ at $p(\phi_j)$. Then there are:
\begin{enumerate}
\item a complete local noetherian $\mathcal O$-algebra $R'$ with residue field $k$, with an ordinary pair $(D',\psi')$ for $\Gamma'$, elements $U'_{j,i}\equiv\tilde u_{j,i}$ modulo $\mathfrak m_{R'}$, and an $\mathcal O$-algebra map $s':S\to R'$ with $s'(X_k)=c_k\psi'_{v_k,i_k}(\gamma_k)-1$;
\item an $R'$-module $M'$ with an additive isomorphism $M'\cong\big(S/((1+y_{j,i})^{\ell^{a_j}}-1)_{j,i}\big)^r$, for some $r$, compatible with the action of $S$ through $s'$;
\item an $\mathcal O$-algebra map $\theta:R'\to B[u]$ carrying $(D',\psi')$ to $(D^{\mathrm u}\circ p,\psi^{\mathrm u})$, $U'_{j,i}$ to $u_{j,i}$, and $s'(z)$ to $\lambda(z_0)$ for $z\in S$ with constant term $z_0\in\Lambda_D$ in the $y_{j,i}$;
\item additive maps $\varphi_X$ from $B[u]\otimes_BM_0$ to $X'=M'/\sum_{j,i}s'(y_{j,i})M'$ and $\psi_X$ back, with $\varphi_X(\theta(t)z)=t\varphi_X(z)$ and $\psi_X(tz)=\theta(t)\psi_X(z)$ for $t\in R'$, such that $\varphi_X\psi_X$ and $\psi_X\varphi_X$ are multiplication by $f'$ and $\theta(f')$, where $f'=\prod_j\prod_i\prod_{i'\ne i}(U'_{j,i}-U'_{j,i'})(U'_{j,i}-b_jU'_{j,i'})$;
\end{enumerate}
such that the $e$-th powers of $f'\operatorname{tr}_{D'}\big(h\prod_i(\phi_j-U'_{j,i})\big)$ and of $f'\operatorname{tr}_{D'}\big(h\prod_{k\ne i}(\phi_j-U'_{j,k})(\tau_j-1-s'(y_{j,i}))\big)$ kill $M'$, for all $j$, $i$ and $h\in\Gamma'_0$.
\end{lemma}
\leanref{Development/AutomorphicArithmetic/AutomorphicVarieties/PatchingModules/NilpotentLevelLowering/OpToTW.lean}{UpperTheory.V.Pa.Lev.opToTW}
\begin{proof}
Let $R$, $(D_R,\psi_R)$, $\beta$, $M_1$, $U_{j,i}$, $Y_{j,i}$, $M_w$, $U^w_{j,i}$, $\Delta$ and $\mathcal N$ be the data of the operator module. The composite of $\Lambda_D\to R$ with $\beta$ is $\lambda$: both are $\mathcal O$-algebra maps from the power series ring $\Lambda_D$ to the complete local ring $B$ with the same values $c_k\psi^{\mathrm u}_{v_k,i_k}(\gamma_k)-1$ at the $X_k$, which lie in $\mathfrak m_B$.

\emph{The ring.} Let $R'=R[[u'_{j,i},y_{j,i}]]$, a power series ring over $R$ in $2nq$ variables; it is a complete local noetherian $\mathcal O$-algebra with residue field $k$. Let $(D',\psi')$ be the image of $(D_R,\psi_R)$ under $R\to R'$, $U'_{j,i}=\tilde u_{j,i}+u'_{j,i}$, and $s':S\to R'$ the map that is $\Lambda_D\to R$ on coefficients and sends $y_{j,i}$ to $y_{j,i}$. This gives (1). Let $\mathrm{ev}:R[U_{j,i},Y_{j,i}]\to R'$ be the $R$-algebra map sending $U_{j,i}$ to $U'_{j,i}$ and $Y_{j,i}$ to $y_{j,i}$.

\emph{The module.} Apply Lemma~\ref{sha-hecke-ideal-tw-nilpotent} to $M_1$, over $\Lambda=\Lambda_D$, and the endomorphisms $U_{j,i}-\tilde u_{j,i}$ and $Y_{j,i}$; let $M'=M_1^{\mathrm{nil}}$, an $R'$-module in which $u'_{j,i}$ acts as $U_{j,i}-\tilde u_{j,i}$ and $y_{j,i}$ as $Y_{j,i}$, so that $\mathrm{ev}(t)$ acts on $M'$ as $t$ acts on $M_1$, and $\pi(tm)=\mathrm{ev}(t)\pi(m)$. In particular, if $t$ acts on $M_1$ as $0$, then $\mathrm{ev}(t)$ kills $M'$. The generators of $\Delta$ act as $1+Y_{j,i}$, which commute with the $U_{j,i}$ and $Y_{j,i}$, so $\Delta$ commutes with $\pi$ and preserves $M'$ and $\ker\pi$. Lemma~\ref{sha-hecke-ideal-tw-free}, with $1+y_{j,i}$ acting on $M'$ through $s'$ as $1+Y_{j,i}$, gives (2).

\emph{The quotient.} Lemma~\ref{sha-hecke-ideal-tw-quotient}, applied to $M_1$, $M_w$, $\mathcal N$ and the $\tilde u_{j,i}$, gives $\jmath:X'\to M_w$ and $\kappa:M_w\to X'$, where $X'=M'/\sum s'(y_{j,i})M'$, compatible with $U'_{j,i}$ on $X'$ and $U^w_{j,i}$ on $M_w$, with $\kappa\jmath=1$ and with image $\jmath(X')=M_w^{\mathrm{nil}}$.

\emph{The splitting algebra.} The $P_j$ are monic of degree $n$. Relation (a) and Lemma~\ref{sha-hecke-ideal-tw-splitting}, applied with $\Lambda_D\to B$ through $\lambda$, $R$, $\beta$, the images of the $\tilde u_{j,i}$, $W=M_w$ and $f$, give $\eta$, $\varphi_u$ and $\psi_u$.

\emph{The map $\theta$.} Let $\theta:R'\to B[u]$ kill the $y_{j,i}$, apply $\beta$ to the coefficients, and send $u'_{j,i}$ to $u_{j,i}-\tilde u_{j,i}$. It carries $(D',\psi')$ to $(D^{\mathrm u}\circ p,\psi^{\mathrm u})$ because $\beta$ does, $U'_{j,i}$ to $u_{j,i}$, and $s'(z)$ to $\beta(z_0)=\lambda(z_0)$. Also $\theta(\mathrm{ev}(t))=\eta(t)$ for $t\in R[U_{j,i}]$. This gives (3).

\emph{The maps.} Put $\varphi_X=\kappa\circ\varphi_u$ and $\psi_X=\psi_u\circ\jmath$. For $t$ a constant or a variable $u'_{j,i}$ of $R'$, $t=\mathrm{ev}(t_0)$ with $t_0\in R[U_{j,i}]$ and $\theta(t)=\eta(t_0)$, so $\varphi_X(\theta(t)z)=\kappa(t_0(U^w)\varphi_u(z))=t\varphi_X(z)$, and similarly for $\psi_X$; for $t=y_{j,i}$ both sides vanish, as $\theta(y_{j,i})=0$ and $y_{j,i}$ kills $X'$. Both $B[u]\otimes_BM_0$ and $X'$ are finite over $\Lambda_D$, and the variables of $R'$ act on them topologically nilpotently: on $X'$ by the definition of $M'$, and on $B[u]\otimes_BM_0$ by Lemma~\ref{sha-hecke-ideal-tw-splitting}. So the compatibility with constants and variables extends to all power series in $R'$. For $x\in X'$, $\jmath(x)\in M_w^{\mathrm{nil}}$, so $\varphi_X\psi_X(x)=\kappa(f(U^w)\jmath(x))=\mathrm{ev}(f)\kappa\jmath(x)=f'x$. For $z\in B[u]\otimes_BM_0$, $\varphi_u(z)\in M_w^{\mathrm{nil}}=\jmath(X')$, so $\jmath\kappa\varphi_u(z)=\varphi_u(z)$ and $\psi_X\varphi_X(z)=\psi_u\varphi_u(z)=\eta(f)z=\theta(f')z$. This gives (4).

\emph{The relations.} By relations (b) and the module step, the images under $\mathrm{ev}$ of the elements of relations (b) kill $M'$. The map $\mathrm{ev}$ carries $f$ to $f'$, the trace of $D_R$ over $R[U_{j,i},Y_{j,i}]$ to that of $D'$, and $h\prod_i(\phi_j-U_{j,i})$ and $h\prod_{k\ne i}(\phi_j-U_{j,k})(\tau_j-1-Y_{j,i})$ to the corresponding elements with $U'_{j,i}$ and $y_{j,i}=s'(y_{j,i})$.
% opToTW: co_mvps (R'), CA.pair_map (the pair over R'), hbeta (co_ev uniqueness for beta o (Lambda -> R) = lam), opCore: nilPart, del_gen and freePart, quotPart, spPart (with Det.chiApi: chi monic of degree n), the map theta (killCompl, MvPowerSeries.map beta, Sp.mk), psLin' for the extension from constants and variables (hnS', hnX), and for the relations hk with fEl_map, htr (CA.d_gmap) and the q1El, q2El images.
\end{proof}

\subsection{Ordinary forms at augmented levels}\label{subsection-sha-augmented-levels}

\begin{lemma}\label{sha-hecke-ideal-tw-augmented-idempotent}
Let $A$ be a complete local noetherian ring with maximal ideal $\mathfrak m$, let $X$ be an $A$-module with $A$-linear endomorphisms $(t_i)_{i\in I}$, let $(a_i)_{i\in I}$ be elements of $A$ and $Y\subset X$ a submodule. Let the \emph{localization} $Y_{(a)}$ of $Y$ at the $t_i$ for the eigenvalues $a_i$ be the set of $y\in Y$ such that for every $i$ and every $k\ge0$ some power of $t_i-a_i$ maps $y$ into $\mathfrak m^kY$. Suppose that $Y$ is finite over $A$ and that the $t_i$ map $Y$ into itself and commute on it, and let $T\subset\operatorname{End}_A(Y)$ be the $A$-subalgebra generated by the $t_i$.
\begin{enumerate}
\item There is an idempotent $e\in T$ such that for every $i$ some power of $t_i-a_i$, multiplied by $e$, lies in $\mathfrak mT$, and $1-e$ lies in the ideal $\mathfrak mT+\sum_i(t_i-a_i)T$.
\item Let $e\in T$ be an idempotent and $(u_j)_{j\in J}$ elements of $T$ such that for every $j$ some $u_j^Ne$ lies in $\mathfrak mT$, and $1-e\in\mathfrak mT+\sum_ju_jT$. Then an element $y\in Y$ satisfies $ey=y$ if and only if for every $j$ and every $k\ge0$ some power of $u_j$ maps $y$ into $\mathfrak m^kY$.
\item For $e$ as in (1), $Y_{(a)}=eY$, and $e$ commutes with every endomorphism of $Y$ that commutes with all the $t_i$.
\end{enumerate}
\end{lemma}
\leanref{Development/AutomorphicArithmetic/AutomorphicVarieties/HeckeAlgebras/ContinuousTorusCharacter.lean}{UpperTheory.V.Hi.locPart}
\leanref{Development/AutomorphicArithmetic/AutomorphicVarieties/AlgebraTools/HidaCompleteRings/LocAbs.lean}{UpperTheory.V.AlgebraTools.AnnihilatorQuotientFinite.locAbs}
\leanref{Development/AutomorphicArithmetic/AutomorphicVarieties/AlgebraTools/HidaCompleteRings/LocChar.lean}{UpperTheory.V.AlgebraTools.AnnihilatorQuotientFinite.locChar}
\leanref{Development/AutomorphicArithmetic/AutomorphicVarieties/AlgebraTools/HidaCompleteRings/LocIdem.lean}{UpperTheory.V.AlgebraTools.AnnihilatorQuotientFinite.locIdem}
\begin{proof}
(1) The $A$-module $T$ is a submodule of the finite $A$-module $\operatorname{End}_A(Y)$, so it is finite, hence complete for the $\mathfrak mT$-adic topology \cite[\S8]{Mat86}, and $\bar T=T/\mathfrak mT$ is a finite-dimensional algebra over $A/\mathfrak m$, hence artinian. Let $J=\mathfrak mT+\sum_i(t_i-a_i)T$. If $J=T$, take $e=0$. Otherwise let $\bar J$ be the image of $J$ in $\bar T$. The powers of $\bar J$ stabilize, so $\bar J^N=\bar J^{N+1}=\bar J^{2N}$ for some $N\ge1$. The finitely generated ideal $\bar J^N$ equals its square, so by Nakayama's lemma it is generated by an idempotent $\bar f$. Then $\bar f\in\bar J^{N+1}\subset\bar J$, and $\bar s^N(1-\bar f)=0$ for $\bar s\in\bar J$, since $\bar s^N\in\bar T\bar f$. As $T$ is $\mathfrak mT$-adically complete, the idempotent $1-\bar f$ lifts to an idempotent $e\in T$, by Newton's iteration for $X^2-X$, whose derivative $2X-1$ is a unit at an idempotent. Then $(t_i-a_i)^Ne\in\mathfrak mT$ for every $i$, and $1-e\in J$, because its image $\bar f$ lies in $\bar J$ and $J\supset\mathfrak mT$.

(2) The elements of $\mathfrak mT$ map $\mathfrak m^kY$ into $\mathfrak m^{k+1}Y$. Suppose $ey=y$, and let $u_j^Ne\in\mathfrak mT$. Then $u_j^{Nk}y=(u_j^Ne)^ky\in\mathfrak m^kY$. Conversely, suppose that the condition on the powers of the $u_j$ holds, and let $y'=(1-e)y$. Fix $k$. The $t\in T$ with $ty'\in\mathfrak m^kY$ form an ideal $S_k$ of $T$. Its radical contains $\mathfrak mT$, since $\mu^ky'\in\mathfrak m^kY$ for $\mu\in\mathfrak mT$, and it contains every $u_j$, since $u_j^ny'=(1-e)u_j^ny\in\mathfrak m^kY$ for a suitable $n$. So it contains $1-e$, and $(1-e)^Ky'\in\mathfrak m^kY$ for some $K$. As $(1-e)^Ky'=y'$, $y'\in\mathfrak m^kY$. This holds for every $k$, so $y'=0$ by Krull's intersection theorem \cite[\S8]{Mat86}, and $ey=y$.

(3) By (1) and (2), with $u_i=t_i-a_i$, $Y_{(a)}$ is the set of $y$ with $ey=y$, which is $eY$ as $e$ is idempotent. Since $e$ lies in the algebra generated by the $t_i$, it commutes with every endomorphism that commutes with them.
% locAbs is Reformulated.B.Qe2.locAbs (PrecompleteOfFinite.lean): precomplete_of_finite, artin_idem (the stabilized power of J/IT generated by an idempotent), lift_idem (Hensel for X^2 - X); the case J = T gives e = 0. locChar is followed in (2), with the radical of the ideal S of t with t y in m^k N and IsHausdorff for Krull. locIdem: heckeModel, locAbs, locChar.
\end{proof}

\begin{lemma}\label{sha-hecke-ideal-tw-augmented-fewer}
Let $\ell$ be a prime, $\mathcal O$ the ring of integers of a finite extension of $\mathbf Q_\ell$, with residue field $k$, and $|\cdot|$ the absolute value of $\bar{\mathbf Q}_\ell$. Let $A$ be a complete local noetherian $\mathcal O$-algebra with residue field $k$ that is a domain, $Y$ a finite free $A$-module, $(t_i)_{i\in I}$ commuting $A$-linear endomorphisms of $Y$, $T\subset\operatorname{End}_A(Y)$ the $A$-subalgebra they generate, $(a_i)_{i\in I}$ elements of $\mathcal O$, $I'\subset I$ a subset and $x:A\to\bar{\mathbf Q}_\ell$ an $\mathcal O$-algebra map. Suppose that every ring map $z:T\to\bar{\mathbf Q}_\ell$ that equals $x$ on $A$ and satisfies $|z(t_i)-a_i|<1$ for all $i\in I'$ satisfies $|z(t_i)-a_i|<1$ for all $i\in I$. Then the localization of $Y$ at the $t_i$, $i\in I'$, and its localization at the $t_i$, $i\in I$, for the eigenvalues $a_i$ (Lemma~\ref{sha-hecke-ideal-tw-augmented-idempotent}), are equal.
\end{lemma}
\leanref{Development/AutomorphicArithmetic/AutomorphicVarieties/HeckeAlgebras/OrdinaryHeckeCharacters/LevelUR2LocSub.lean}{UpperTheory.V.Hi.levelUR2LocSub}
\begin{proof}
The algebra $T$ is finite over $A$, as in the proof of Lemma~\ref{sha-hecke-ideal-tw-augmented-idempotent}(1), and torsion free: if $b\in A$ is nonzero and $bt=0$ for some $t\in T$, then $t=0$, since $Y$ is free over the domain $A$. Also $|x(b)|<1$ for $b\in\mathfrak m_A$: otherwise $1/x(b)$ is integral over $\mathcal O$, and clearing denominators in its minimal polynomial gives $R\in\mathcal O[X]$ with $R(0)=1$ and $R(x(b))=0$, while $R(b)$ is a unit of $A$. As every element of $A$ is $c+b$ with $c\in\mathcal O$ and $b\in\mathfrak m_A$, $|x|\le1$ on $A$. Let $e\in T$ be an idempotent as in Lemma~\ref{sha-hecke-ideal-tw-augmented-idempotent}(1) for the $t_i$, $i\in I$. Let $T_0\subset T$ be the $A$-subalgebra generated by the $t_i$, $i\in I'$, and $e_0\in T_0$ an idempotent as in that lemma for $T_0$ and these $t_i$. Then $f=e_0(1-e)$ is an idempotent of $T$.

\emph{$f=0$.} Suppose $f\ne0$. A nonzero idempotent is not nilpotent, so $f$ lies outside some prime of $T$, and outside a minimal prime $Q$ contained in it. An element of $A$ that lies in the minimal prime $Q$ is a zero divisor of $T$, so it is $0$, as $T$ is torsion free; thus $Q\cap A=0$. Since $T$ is integral over $A$ and $\bar{\mathbf Q}_\ell$ is algebraically closed, $x$ extends to a ring map $z:T\to\bar{\mathbf Q}_\ell$ with $z(Q)=0$, by lying over \cite[Theorem~9.3]{Mat86}. Every $s\in T$ is a root of a monic polynomial with coefficients in $A$, and $|x|\le1$ on $A$, so $|z(s)|\le1$. Hence the $s\in T$ with $|z(s)|<1$ form an ideal $J_z$, which contains $\mathfrak m_AT$. As $f(1-f)=0\in Q$ and $f\notin Q$, $1-f\in Q$, so $z(f)=1$, and therefore $z(e_0)=1$ and $z(1-e)=1$. For $i\in I'$ some $(t_i-a_i)^Ne_0$ lies in $\mathfrak m_AT_0$, so $|z(t_i)-a_i|^N<1$ and $|z(t_i)-a_i|<1$. By hypothesis $|z(t_i)-a_i|<1$ for every $i\in I$, so $J_z$ contains $\mathfrak m_AT+\sum_{i\in I}(t_i-a_i)T$, hence $1-e$. This contradicts $z(1-e)=1$.

\emph{The localizations.} Since $f=0$, $1-e=(1-e_0)(1-e)$, which lies in $\mathfrak m_AT+\sum_{i\in I'}(t_i-a_i)T$ because $1-e_0$ does. Also every $t_i-a_i$, $i\in I$, has a power whose product with $e$ lies in $\mathfrak m_AT$. By Lemma~\ref{sha-hecke-ideal-tw-augmented-idempotent}(2), applied to $e$ with the family $(t_i-a_i)_{i\in I'}$ and with the family $(t_i-a_i)_{i\in I}$, both localizations are the set of $y\in Y$ with $ey=y$.
% levelUR2LocSub: htf (torsion free), the private ls_core (locAbs for T and for T0 = adjoin of the t_(j i); f = e0 (1 - e); a minimal prime Q avoiding f by Ideal.exists_minimalPrimes_le; hQA through disjoint_nonZeroDivisors_of_mem_minimalPrimes; AnnihilatorQuotientFinite.extendPt; ls_le and ls_small for the ideal of |z| < 1; Up.pt_le and Up.pt_lt for |x|), then locChar twice through ls_mem. The Lean indexes the smaller family by a map j : I' -> I.
\end{proof}

\begin{lemma}\label{sha-hecke-ideal-tw-augmented-isotypic}
\begin{enumerate}
\item Let $G$ be a group, $Y''$ a representation of $G$ on a complex vector space, $(\sigma_i)_{i\in I}$ irreducible representations of $G$ on complex vector spaces such that $Y''$ is the sum of the images of the $G$-equivariant maps $\sigma_i\to Y''$, $i\in I$, and $Y\subset Y''$ a subrepresentation. Then $Y$ is the sum of the images of the $G$-equivariant maps $\sigma_i\to Y$, $i\in I$.
\item Let $H_1,\dots,H_q$ be groups, $H=\prod_jH_j$, $\pi$ a representation of $H$ on a complex vector space $Y$, and $(\sigma_{i,j})_{i\in I}$, for each $j$, irreducible representations of $H_j$ on complex vector spaces, such that for every $j$, $Y$ is the sum of the images of the $H_j$-equivariant maps $\sigma_{i,j}\to Y$, $i\in I$. Then for every nonzero $y\in Y$ there are an $H$-equivariant surjection $p:Y\to Y'$ onto a representation $Y'$ of $H$ with $p(y)\ne0$, and indices $c_1,\dots,c_q\in I$, such that for every $j$, $Y'$ is the sum of the images of the $H_j$-equivariant maps $\sigma_{c_j,j}\to Y'$.
\end{enumerate}
\end{lemma}
\leanref{Development/AutomorphicArithmetic/Foundations/GaloisRepresentations/AdmissibleResidues/IsoSub.lean}{UpperTheory.GalN.V.isoSub}
\leanref{Development/AutomorphicArithmetic/Foundations/GaloisRepresentations/AdmissibleResidues/IsoPick.lean}{UpperTheory.GalN.V.isoPick}
\leanref{Development/AutomorphicArithmetic/Foundations/GaloisRepresentations/AdmissibleResidues/IsoRefine.lean}{UpperTheory.GalN.V.isoRefine}
\begin{proof}
(1) The image of an equivariant map from an irreducible representation is $0$ or irreducible, so $Y''$ is a sum of irreducible subrepresentations, each isomorphic to some $\sigma_i$. Hence $Y''$ is semisimple, and each of its irreducible subrepresentations is isomorphic to some $\sigma_i$. The subrepresentation $Y$ is semisimple, so it is the sum of its irreducible subrepresentations, and these are images of equivariant maps from the $\sigma_i$.

(2) \emph{One factor.} Fix $j$ and a nonzero $y\in Y$. As in (1), $Y$ is a semisimple representation of $H_j$, the direct sum of its isotypic components, each of the type of some $\sigma_{i,j}$. Choose $i_0$ such that the component of $y$ in the isotypic component of type $\sigma_{i_0,j}$ is nonzero, and let $N_j$ be the sum of the images of the equivariant maps from the $\sigma_{i,j}$ that are not isomorphic to $\sigma_{i_0,j}$. Then $N_j$ is contained in the sum of the other isotypic components, so $y\notin N_j$. The subspace $N_j$ is stable under $H$: under $H_j$ because it is a sum of images of $H_j$-equivariant maps, and under the other factors $H_{j'}$ because they commute with $H_j$, so that the composite of such a map with $\pi(h')$, $h'\in H_{j'}$, is again $H_j$-equivariant. In $Y/N_j$, the image of every equivariant map from a $\sigma_{i,j}$ not isomorphic to $\sigma_{i_0,j}$ is $0$, and the others are images of maps from $\sigma_{i_0,j}$; so $Y/N_j$ is the sum of the images of the $H_j$-equivariant maps $\sigma_{i_0,j}\to Y/N_j$.

\emph{All factors.} The two properties pass to quotients: the image in a quotient of a sum of images of equivariant maps from given representations is again such a sum. Apply the first step for $j=1,\dots,q$ in turn, each time to the quotient obtained so far and the image of $y$, which stays nonzero; let $c_j$ be the index $i_0$ chosen at step $j$, and let $p$ be the composite of the quotient maps.
% isoSub is irreducible_subrep_isotypic (MemIsotypicOfInjective: mem_isotypic_of_injective, iso_sub). isoPick is isoPick_gen. isoRefine: induction over Finset Q, stab_N (N is H-stable), iso_quot, iso_push and span_push for the quotient. The Lean's representations sigma are SmIrr (smooth irreducible); only irreducibility is used.
\end{proof}

\begin{lemma}\label{sha-hecke-ideal-tw-augmented-projection}
Let $V$ be a complex vector space, $(T_t)_{t\in R}$ and $(\pi_g)_{g\in G}$ families of linear endomorphisms of $V$, and $E\subset V$ a subspace with $\pi_gE\subset E$ and $T_t\pi_gF=\pi_gT_tF$ for all $t$, $g$ and $F\in E$. Let $(W_j)_{j\in J}$ be subspaces of $V$ with $E=\sum_jW_j$, and $\chi_j:R\to\mathbf C$ maps such that $T_t$ acts on $W_j$ by multiplication by $\chi_j(t)$. Let $F_0\in E$ be nonzero. Then there are $j_0\in J$ with $W_{j_0}\ne0$, a subspace $Y\subset\sum_{j:\chi_j=\chi_{j_0}}W_j$ with $\pi_gY\subset Y$ for all $g$, and a linear map $p:V\to Y$ with $p(F_0)\ne0$, such that $p(T_tF)=\chi_{j_0}(t)p(F)$ and $p(\pi_gF)=\pi_gp(F)$ for all $t$, $g$ and $F\in E$.
\end{lemma}
\leanref{Development/AutomorphicArithmetic/AutomorphicVarieties/GaloisLibrary/MemIsotypicOfInjective.lean}{UpperTheory.V.PiLin.exists_proj}
\begin{proof}
For a map $c:R\to\mathbf C$ let $E_c=\sum_{\chi_j=c}W_j$ and $E^c=\sum_{\chi_j\ne c}W_j$, so that $E=E_c+E^c$. Each $T_t$ acts on $E_c$ by $c(t)$ and maps $E^c$ into itself. An element of $E^c$ on which every $T_t$ acts by $c(t)$ is $0$: writing it as a finite sum of elements of $W_j$ with $\chi_j\ne c$, this follows by induction on the number of terms, applying $T_t-\chi_{j_1}(t)$ for a $t$ with $\chi_{j_1}(t)\ne c(t)$. So $E=E_c\oplus E^c$, and an element of $E$ on which every $T_t$ acts by $c(t)$ lies in $E_c$. Hence $\pi_g$ maps $E_c$ into $E_c$: for $F\in E_c$, $\pi_gF\in E$ and $T_t\pi_gF=\pi_gT_tF=c(t)\pi_gF$. Similarly $\pi_g$ maps each $W_j$, which lies in $E_{\chi_j}$, into $E_{\chi_j}$, so it maps $E^c$ into $E^c$.

Write $F_0=\sum_{j\in J_0}F_j$ with $J_0$ finite and $F_j\in W_j$. Grouping the terms by the value of $\chi_j$, some group has a nonzero sum, since $F_0\ne0$; choose $j_0$ in it with $F_{j_0}\ne0$, and let $c=\chi_{j_0}$ and $Y=E_c$. Let $C$ be a complement of $E$ in $V$, and $p:V=E_c\oplus(E^c\oplus C)\to E_c$ the projection. Then $p(F_0)$ is the sum of the $F_j$ with $\chi_j=c$, which is nonzero. For $F=F'+F''\in E$ with $F'\in E_c$ and $F''\in E^c$, $p(T_tF)=T_tF'=c(t)p(F)$ and $p(\pi_gF)=\pi_gF'=\pi_gp(F)$, since $T_t$ and $\pi_g$ preserve $E_c$ and $E^c$.
% exists_proj: hYe, hKz (the private exists_equivariant_projection_onto_invariant_submodule), hsplit, hEY, hpiY, hpiK, then Submodule.exists_isCompl and projectionOnto.
\end{proof}

\begin{situation}\label{sha-hecke-ideal-tw-augmented-setup}
In Situation~\ref{sha-hecke-situation}, assume that Galois representations exist, in the sense of Definition~\ref{gal-galois-representations}, relative to the maps satisfying the laws of Definition~\ref{gal-llc-laws}. Let $\iota:\bar{\mathbf Q}_\ell\cong\mathbf C$, and let $U(n)$, $\tau_0$, $\theta$, $a$ and $\mathcal K_0$ be as in Lemma~\ref{sha-hecke-algebra}, with $\tau_0$ acting on $\mathcal O^m$. Let $S'\supset S$ be a finite set of places of $L^+$ whose elements outside $S$ are split in $L$, and $\Gamma'$, $p:\Gamma'\to\Gamma$ and $\Gamma'_0$ as in Definition~\ref{sha-hecke-ideal-tw-level-module}. Let $v_1,\dots,v_q$ be the places of $S'\setminus S$, $w_j$ a place of $L$ above $v_j$, $G_w=\prod_jGL_n(L_{w_j})$ and $K_w=\prod_jGL_n(\mathcal O_{L_{w_j}})$.
\begin{enumerate}
\item (Groups.) Write $U(n)(\mathbf A^\infty_{L^+})=G^{(\ell)}\times G_\ell$, where $G^{(\ell)}$ is the group of points over the finite adeles away from $\ell$ and $G_\ell=\prod_{v\mid\ell}GL_n(L_{\tilde v})$, through the places $\tilde v$. In $G_\ell$ use the notation $N_0$, $\bar N_c$, $T_0$, $e_{v,i}$, $\alpha_j$, $z$ and $U_\alpha$ of Definition~\ref{sha-hecke-local-ordinary}, formed in each factor and multiplied over $v\mid\ell$, and let $I_\ell=N_0T_0\bar N_0$. Let $G^{(\ell)}_0\subset G^{(\ell)}$ be the subgroup of elements whose components at the places of $S_{\mathrm a}$ are integral and whose component at $v_0$ lies in the principal congruence subgroup of level $\tilde v_0$; it acts on $\mathcal O^m$ through $\tau_0$ on these components, and it contains $G_w$, through the $w_j$, and $GL_n(L_w)$, through $w$, for every split place $w$ of $L$ above a place outside $S$.
\item (Forms.) A \emph{$\Lambda$-adic form} is a function $f:G^{(\ell)}\times G_\ell\to\Lambda^m$ such that $f(\gamma g)=f(g)$ for $\gamma\in U(n)(L^+)$; $f(gn)=f(g)$ for $n\in N_0$; $f(gt)=\theta(t)^{-1}f(g)$ for $t\in T_0$; and for every $k\ge0$ there is $c\ge0$ such that $f(gy)-f(g)$ has entries in $\mathfrak m_\Lambda^k$ for all $g$ and $y\in\bar N_c$. It is \emph{ordinary} if for every $k\ge0$ there is $N\ge0$ such that $(U_z^{N+1}-1)f$ has values in $\mathfrak m_\Lambda^k\Lambda^m$. Let $\mathcal F\supset\mathcal F^{\mathrm{ord}}$ be the $\Lambda$-modules of $\Lambda$-adic forms and of ordinary ones. The group $G^{(\ell)}_0$ acts on functions by $(u\cdot f)(g)=\tau_0(u)f(gu)$; for a subgroup $U\subset G^{(\ell)}_0$, $\mathcal F(U)$ and $\mathcal F^{\mathrm{ord}}(U)$ are the forms fixed by $U$, and $\rho(g)f=g\cdot f$ for $g\in G_w$, so that $(\rho(g)f)(h)=f(hg)$.
\item (Levels.) A subgroup $U\subset G^{(\ell)}_0$ is \emph{neat} if $\gamma g=g(u,b)$, with $\gamma\in U(n)(L^+)$, $g\in G^{(\ell)}\times G_\ell$, $u\in U$ and $b\in I_\ell$, implies $u=1$ and $b=1$; it has \emph{finite level} if $U(n)(L^+)\backslash(G^{(\ell)}\times G_\ell)/(U\times I_\ell)$ is finite. For a family $K=(K_j)$ of compact open subgroups $K_j\subset GL_n(L_{w_j})$, the \emph{level} $\mathcal K_0K$ is the subgroup of $G^{(\ell)}_0$ of the elements whose component at each $v_j$ lies in $K_j$, through $w_j$, and whose components at every place $v\notin S'$ lie in $GL_n(\mathcal O_{L_{w'}})$ for the places $w'$ of $L$ above $v$. Write $K'\subset K$ if $K'_j\subset K_j$ for every $j$, and $K\cap K'$ for $(K_j\cap K'_j)$.
\item (Operators.) For a split place $w$ of $L$ above a place of $L^+$ outside $S'$ and $1\le j\le n$, $T_{w,j}$ is the double coset operator $[GL_n(\mathcal O_{L_w})\alpha_jGL_n(\mathcal O_{L_w})]$ of the action of $GL_n(L_w)$ on functions (Definition~\ref{sha-hecke-ideal-tw-iwahori}(3)), with $\alpha_j$ formed with a uniformizer of $L_w$; for $v\mid\ell$, $U_{v,j}=U_{\alpha_j}$ in the factor at $v$. Their \emph{eigenvalues} are the $a_{w,j}$ and $a_{v,j}$ of $a$.
\item (Localized modules.) For a level $\mathcal K_0K$, $M_K$ is the localization of $\mathcal F^{\mathrm{ord}}(\mathcal K_0K)$ at the operators $T_{w,j}$ and $U_{v,j}$ for their eigenvalues (Lemma~\ref{sha-hecke-ideal-tw-augmented-idempotent}), and, when $M_K\ne0$, $\mathbf T_K$ is the $\Lambda$-algebra of endomorphisms of $M_K$ generated by these operators, as in Lemma~\ref{sha-hecke-algebra}. Let $P=\sum_KM_K\subset\mathcal F$, the sum over all families $K$. For $S'=S$ there are no places $w_j$, and $M_{()}$ is the module $M$ of Lemma~\ref{sha-hecke-algebra}.
\end{enumerate}
\end{situation}
\leanref{Development/AutomorphicArithmetic/AutomorphicVarieties/HeckeAlgebras/CoefficientRepTrivialOnUnramified.lean}{UpperTheory.V.Hi.HSet}
\leanref{Development/AutomorphicArithmetic/AutomorphicVarieties/HeckeAlgebras/IsTheta.lean}{UpperTheory.V.Hi.CaseU.hset}
\leanref{Development/AutomorphicArithmetic/AutomorphicVarieties/HeckeAlgebras/IsTheta.lean}{UpperTheory.V.Hi.CaseU.Ub}
\leanref{Development/AutomorphicArithmetic/AutomorphicVarieties/HeckeAlgebras/IsTheta.lean}{UpperTheory.V.Hi.CaseU.UQ}
\leanref{Development/AutomorphicArithmetic/AutomorphicVarieties/HeckeAlgebras/ContinuousTorusCharacter.lean}{UpperTheory.V.Hi.Frame.Amb}
\leanref{Development/AutomorphicArithmetic/AutomorphicVarieties/HeckeAlgebras/ContinuousTorusCharacter.lean}{UpperTheory.V.Hi.Frame.Ord}
\leanref{Development/AutomorphicArithmetic/AutomorphicVarieties/HeckeAlgebras/CoefficientRepTrivialOnUnramified.lean}{UpperTheory.V.Hi.HSet.lev}
\leanref{Development/AutomorphicArithmetic/AutomorphicVarieties/HeckeAlgebras/CoefficientRepTrivialOnUnramified.lean}{UpperTheory.V.Hi.HSet.op}
\leanref{Development/AutomorphicArithmetic/AutomorphicVarieties/HeckeAlgebras/CoefficientRepTrivialOnUnramified.lean}{UpperTheory.V.Hi.HSet.MK}

\begin{definition}\label{sha-hecke-ideal-tw-augmented-parts}
In Situation~\ref{sha-hecke-ideal-tw-augmented-setup}:
\begin{enumerate}
\item A \emph{weight} is a family $\lambda=(\lambda_{\tau,1}\ge\dots\ge\lambda_{\tau,n})$ of nonincreasing $n$-tuples of integers indexed by the embeddings $\tau:L\to\bar{\mathbf Q}_\ell$. An $\mathcal O$-algebra map $x:\Lambda\to\bar{\mathbf Q}_\ell$ is an \emph{arithmetic point of $\Lambda$ of weight $\lambda$} if there is $e\ge0$ such that for all $v\mid\ell$ and $1\le i\le n$ the character $\chi_{x,v,i}(u)=x(\theta(e_{v,i}(u)))\prod_\tau\tau(u)^{-\lambda_{\tau,n+1-i}}$ of $\mathcal O_{L_{\tilde v}}^\times$, the product over the embeddings $\tau:L_{\tilde v}\to\bar{\mathbf Q}_\ell$ and $\lambda_\tau$ taken for the restriction of $\tau$ to $L$, is trivial on the units congruent to $1$ modulo $\varpi^e$. It is a \emph{point of $\Lambda$ of distinct type of weight $\lambda$} if moreover for each $v$ the characters $\chi_{x,v,1},\dots,\chi_{x,v,n}$ are pairwise distinct.
\item Let $\lambda$ be a weight and $\xi$ a representation of $G_\ell$ on $\mathbf C^d$ that is \emph{algebraic of weight $\lambda$ through $\iota$}: isomorphic to the representation $g\mapsto\bigotimes_{(v,\tau)}\rho_{\lambda_\tau}(\iota\tau(g_v))$ on $\bigotimes_{(v,\tau)}\nabla(\lambda_\tau)$, the tensor products over the pairs of a place $v\mid\ell$ and a continuous embedding $\tau:L_{\tilde v}\to\bar{\mathbf Q}_\ell$, with $\lambda_\tau$ taken for the restriction of $\tau$ to $L$, $g_v$ the component of $g$ at $v$, and $\nabla$ and $\rho$ as in Definition~\ref{is-polynomial-model} over $\mathbf C$. Let $V^\xi$ be the space of functions $G^{(\ell)}\times G_\ell\to(\mathbf C^d)^m$, on which $G^{(\ell)}_0$ acts by $(u\cdot F)(g)=\iota\tau_0(u)F(gu)$ and $G_w$ by $\pi_\xi(g)F=g\cdot F$, and let $T^\xi_{w,j}$ be defined on $V^\xi$ as $T_{w,j}$ is on functions with values in $\Lambda^m$. An \emph{automorphic part} of weight $\xi$ is a nonzero subspace $Y\subset V^\xi$ of functions $F$ with $F(\gamma g)=\xi(\gamma_\ell)F(g)$ for $\gamma\in U(n)(L^+)$, with $\gamma_\ell$ its component in $G_\ell$, and fixed by the elements of $\mathcal K_0K_w$ whose components at the $v_j$ are $1$, together with:
\begin{enumerate}
\item irreducible smooth representations $\sigma_{Y,j}$ of $GL_n(L_{w_j})$, for $1\le j\le q$, and $\sigma_{Y,v}$ of $GL_n(L_{\tilde v})$, for $v\mid\ell$, such that $Y$ is the sum of the images of the equivariant maps from $\sigma_{Y,j}$ to $Y$ under right translation by $GL_n(L_{w_j})$, and similarly for the $\sigma_{Y,v}$;
\item complex numbers $a_{Y,w,k}$, for the split places $w$ of $L$ above places outside $S'$ and $0\le k\le n$, with $a_{Y,w,0}=1$, such that $T^\xi_{w,k}$ acts on $Y$ by multiplication by $a_{Y,w,k}$ for $k\ge1$;
\item a continuous semisimple representation $\rho_Y:\Gamma'_0\to GL_n(\bar{\mathbf Q}_\ell)$ such that for every such $w$ the characteristic polynomial of $\rho_Y$ at the inverse of a Frobenius element at $w$, carried by $\iota$, is the Hecke polynomial with coefficients $a_{Y,w,k}$;
\item for every $j$, every valuation ring $A$ of $\bar L$ over $w_j$ and every $\operatorname{rc}$ that satisfies the laws (Definition~\ref{gal-llc-laws}), $\rho_Y|_{W_A}$, viewed through $\operatorname{Gal}(\bar L/L)\to\Gamma'_0$, has a Weil-Deligne representation $\mathrm{WD}_\iota(\rho_Y|_{W_A})$, equal to $\operatorname{rc}(\sigma_{Y,j}\nu^{(1-n)/2})$ when the latter is abelian and monodromy-free;
\item for every $v\mid\ell$: if $\sigma_{Y,v}\cong\mathrm{PS}(\omega)$ for smooth characters $\omega_1,\dots,\omega_n$ that are pairwise distinct on $\mathcal O_{L_{\tilde v}}^\times$, then for every $\operatorname{rc}$ that satisfies the laws with $\operatorname{rc}(\sigma_{Y,v}\nu^{(1-n)/2})$ abelian, $\rho_Y|_{D_v}$ is de Rham, with $\mathrm{HT}_\tau(\rho_Y|_{D_v})=\{\lambda_{\tau,i}+n-i:1\le i\le n\}$ for every embedding $\tau:L\to\bar{\mathbf Q}_\ell$ that extends continuously to $L_{\tilde v}$ (Definition~\ref{gal-p-adic-hodge}), and its Weil-Deligne representation has the characteristic polynomials of $\operatorname{rc}(\sigma_{Y,v}\nu^{(1-n)/2})$.
\end{enumerate}
\item Let $S''\supset S$ be a finite set of places of $L^+$, $\Gamma''_0$ the subgroup fixing $L$ of the Galois group of the maximal extension of $L^+$ unramified outside $S''$ and the infinite places, and $T$ a complete local noetherian $\mathcal O$-algebra with residue field $k$ and an $\mathcal O$-algebra map $\Lambda\to T$, with elements $t_{w,j}$, for the split places $w$ of $L$ above places outside $S''$ and $1\le j\le n$, and $u_{v,j}$, for $v\mid\ell$ and $1\le j\le n$. A pair $(D,\psi)$ of a continuous determinant $D$ of $\Gamma''_0$ of dimension $n$ with values in $T$ and continuous characters $\psi_{v,i}:D_v\to T^\times$ is \emph{attached} to these elements if (1) and (2) of Lemma~\ref{sha-hecke-determinant} hold with $t_{w,j}$ and $u_{v,j}$ in place of $T_{w,j}$ and $U_{v,j}$, for the split places $w$ above places outside $S''$.
\end{enumerate}
\end{definition}
\leanref{Development/AutomorphicArithmetic/AutomorphicVarieties/HeckeAlgebras/CoefficientRepTrivialOnUnramified.lean}{UpperTheory.V.Hi.HSet.Arith}
\leanref{Development/AutomorphicArithmetic/AutomorphicVarieties/HeckeAlgebras/CoefficientRepTrivialOnUnramified.lean}{UpperTheory.V.Hi.IsAlgWt}
\leanref{Development/AutomorphicArithmetic/AutomorphicVarieties/HeckeAlgebras/CoefficientRepTrivialOnUnramified.lean}{UpperTheory.V.Hi.algRep}
\leanref{Development/AutomorphicArithmetic/AutomorphicVarieties/HeckeAlgebras/CoefficientRepTrivialOnUnramified.lean}{UpperTheory.V.Hi.HSet.Good}
\leanref{Development/AutomorphicArithmetic/AutomorphicVarieties/ModularityArgument/ThresholdArguments/DenseThresholdCases/UnitaryPointProperties.lean}{UpperTheory.V.Hi.HSet.PartAltAlt}
\leanref{Development/AutomorphicArithmetic/AutomorphicVarieties/HeckeAlgebras/CoefficientRepTrivialOnUnramified.lean}{UpperTheory.V.Hi.HSet.IsUniv}

\begin{lemma}\label{sha-hecke-ideal-tw-augmented-torsion}
Let $L$ be a number field, $p$ an odd prime unramified in $L$, $w$ a place of $L$ above $p$ and $n\ge1$. Then the only element of finite order of the principal congruence subgroup $\ker\big(GL_n(\mathcal O_{L_w})\to GL_n(\mathcal O_{L_w}/w)\big)$ is $1$.
\end{lemma}
\leanref{Development/AutomorphicArithmetic/AutomorphicVarieties/HeckeAlgebras/OrdinaryHeckeCharacters/NtExpand.lean}{nt_free}
\begin{proof}
As $p$ is unramified, the maximal ideal of $\mathcal O_{L_w}$ is $p\mathcal O_{L_w}$, so an element of the principal congruence subgroup is $1+pz$ with $z$ a matrix over $\mathcal O_{L_w}$. Suppose $(1+pz)^m=1$ with $m\ge1$. We show $pz=0$, by induction on the number of prime factors of $m$; for $m=1$ it is clear.

\emph{A prime exponent.} Let $q$ be a prime and $(1+p^{k+1}y)^q=1$ for a matrix $y$ and some $k\ge0$. Expanding, $0=qp^{k+1}y+\binom q2p^{2k+2}y^2+p^{3k+3}y^3y'$ for a matrix $y'$ over $\mathcal O_{L_w}$. If $q=p$, then $p$ divides $\binom p2$, as $p$ is odd, and dividing by $p^{k+2}$ shows that $y$ is $p$ times a matrix. If $q\ne p$, then $q$ is a unit, and dividing by $p^{k+1}$ shows the same. Starting from $(1+pz)^q=1$ and applying this repeatedly, $pz$ is $p^{k+1}$ times a matrix for every $k$, so $pz=0$.

\emph{The induction.} Write $m=qa$ with $q$ prime. Then $(1+pz)^a=1+pz'$ for a matrix $z'$, and $(1+pz')^q=1$, so $pz'=0$ by the prime case. Hence $(1+pz)^a=1$, and $pz=0$ by induction.
% nt_free: nt_sq (p is not in w^2, from p not dividing the discriminant), hdv (the entries of x - 1 divisible by p), nt_mat (induction_on_primes, with nt_expand for the binomial expansion), nt_prime and nt_step (the cases q = p, through Nat.Prime.dvd_choose_self, and q a unit).
\end{proof}

\begin{lemma}\label{sha-hecke-ideal-tw-augmented-finite-order}
Let $L$ be a CM field with maximal totally real subfield $L^+$ and complex conjugation $c$, let $n\ge1$, and let $M\in GL_n(L)$ with ${}^tc(M)M=1$. Suppose that for every finite place $v$ of $L^+$ there are a place $w$ of $L$ above $v$ and $g\in GL_n(L_w)$ with $g^{-1}Mg\in GL_n(\mathcal O_{L_w})$. Then $M$ has finite order.
\end{lemma}
\leanref{Development/AutomorphicArithmetic/AutomorphicVarieties/HeckeAlgebras/OrdinaryHeckeCharacters/NtExpand.lean}{nt_ord}
\begin{proof}
\emph{Integrality.} Let $a$ be a coefficient of the characteristic polynomial of $M$, let $v$ be a finite place of $L^+$ and $w$ and $g$ as in the hypothesis. As $g^{-1}Mg$ has entries in $\mathcal O_{L_w}$, $a$ is integral at $w$. The matrix $M^{-1}={}^tc(M)$ is also conjugate by $g$ into $GL_n(\mathcal O_{L_w})$, and its characteristic polynomial has the coefficients of that of $c(M)$; so $c(a)$ is integral at $w$. Hence $a+c(a)$ and $ac(a)$, which lie in $L^+$, are integral at $v$. As $v$ is arbitrary, they lie in $\mathcal O_{L^+}$, and $a$, a root of $X^2-(a+c(a))X+ac(a)$, is integral.

\emph{Finiteness.} So the characteristic polynomial of $M$ has integral coefficients, and by the Cayley-Hamilton theorem the $\mathcal O_L$-module generated by the powers of $M$ is finitely generated. Hence there is a nonzero $d\in\mathcal O_L$ such that $d$ times every entry of every $M^k$ is integral. For an embedding $\phi:L\to\mathbf C$, $\phi\circ c$ is complex conjugation composed with $\phi$, as $L$ is CM, so $\phi(M^k)$ is a unitary matrix and its entries have absolute value at most $1$. So the entries of the $dM^k$ are integers of $L$ whose images under all embeddings are bounded by $\sum_\phi|\phi(d)|$. There are finitely many such integers, so $M$ has finitely many distinct powers, and it has finite order.
% nt_ord: nt_char (the coefficients of the characteristic polynomial integral at T v, through Matrix.charpoly_units_conj) for M and M inverse, the latter with charpoly_transpose and charpoly_map; nt_int (a + c a and a c a in the integers of L^+ by mem_integers_of_valuation_le_one, valuation_liesOver); nt_fin (IsIntegral of M from the charpoly, fg_adjoin_singleton, isFractional_of_fg, entry_norm_bound_of_unitary, Embeddings.finite_of_norm_le, finite_powers).
\end{proof}

\begin{lemma}\label{sha-hecke-ideal-tw-augmented-uop}
Let $K_1,\dots,K_r$ be nonarchimedean local fields, $n\ge1$ and $G=\prod_sGL_n(K_s)$, and form $N_0$, $\bar N_c$, $T_0$ and $z$ of Definition~\ref{sha-hecke-local-ordinary} in each factor and multiply them over the factors. Call $\alpha\in G$ \emph{dominant} if each of its components is $\operatorname{diag}(\varpi^{a_1},\dots,\varpi^{a_n})$ with $a_1\ge\dots\ge a_n$, as in Definition~\ref{sha-hecke-ideal-tw-iwahori}(2). For $c\ge0$ let $T_0^{(c)}\subset T_0$ be the subgroup of the matrices congruent to $1$ modulo $\varpi^{c+1}$ in each factor, and $I_1(c)$ the subgroup generated by $N_0$, $T_0^{(c)}$ and $\bar N_c$. Let $\pi$ be a representation of $G$ on a module $P$ over a commutative ring, and for dominant $\alpha$ let $U_\alpha v=\sum_x\pi(x\alpha)v$, the sum over representatives $x$ of $N_0/\alpha N_0\alpha^{-1}$. Let $\alpha$ and $\beta$ be dominant.
\begin{enumerate}
\item If $v\in P$ is fixed by $N_0$, then $U_\alpha v$ is fixed by $N_0$, $U_\alpha U_\beta v=U_{\alpha\beta}v$, $U_1v=v$, and $\pi(t)U_\alpha v=U_\alpha\pi(t)v$ for $t\in T_0$.
\item If $v$ is fixed by $I_1(c)$, then so is $U_\alpha v$.
\item If $v$ is fixed by $N_0$, $T_0^{(c)}$ and $\bar N_{c+1}$, then $U_zv$ is fixed by $I_1(c)$.
\end{enumerate}
\end{lemma}
\leanref{Development/AutomorphicArithmetic/AutomorphicVarieties/HeckeAlgebras/OrdinaryHeckeCharacters/UOpMul.lean}{UpperTheory.V.Hi.uOpMul}
\leanref{Development/AutomorphicArithmetic/AutomorphicVarieties/HeckeAlgebras/OrdinaryHeckeCharacters/UOpLevel.lean}{UpperTheory.V.Hi.uOpLevel}
\begin{proof}
(1) If $x$ runs over representatives of $N_0/\alpha N_0\alpha^{-1}$, so do $nx$ for $n\in N_0$ and $txt^{-1}$ for $t\in T_0$, which commutes with $\alpha$ and normalizes $N_0$; this gives the invariance under $N_0$ and the commutation with $T_0$. If $y$ runs over representatives of $N_0/\beta N_0\beta^{-1}$, the products $x\cdot\alpha y\alpha^{-1}$ run over representatives of $N_0/\alpha\beta N_0(\alpha\beta)^{-1}$, which gives $U_\alpha U_\beta v=U_{\alpha\beta}v$. For $\alpha=1$ there is one representative.

(2) and (3) Let $v$ be fixed by $N_0$ and $T_0^{(c)}$, and by $\alpha^{-1}y\alpha$ for every $y\in\bar N_c$; this holds in (2), as $\alpha^{-1}\bar N_c\alpha\subset\bar N_c$, and in (3) for $\alpha=z$, as $z^{-1}\bar N_cz\subset\bar N_{c+1}$. By (1), $U_\alpha v$ is fixed by $N_0$ and $T_0^{(c)}$. Let $y\in\bar N_c$. For each representative $x$, write $yx=x't\,y'$ with $x'\in N_0$, $t\in T_0^{(c)}$ and $y'\in\bar N_c$, and let $x''$ be the representative of the class of $x'$. As $t$ commutes with $\alpha$, $yx\alpha=x''\alpha\cdot k$ with $k=(\alpha^{-1}x''^{-1}x'\alpha)\,t\,(\alpha^{-1}y'\alpha)$, and $k$ fixes $v$, its first factor lying in $N_0$. The map $x\mapsto x''$ is injective: if $x_1''=x_2''$, then $\alpha^{-1}x_1^{-1}x_2\alpha=k_1^{-1}k_2$ is an upper triangular unipotent element of $N_0T_0\bar N_c$, so it lies in $N_0$, and $x_1$ and $x_2$ have the same class. So $x\mapsto x''$ permutes the representatives, and $\pi(y)U_\alpha v=U_\alpha v$. Hence $U_\alpha v$ is fixed by $N_0$, $T_0^{(c)}$ and $\bar N_c$, which generate $I_1(c)$.
% uOpMul: uOpMul_sum and uOpMul_ex (uOpSum) for the coset representatives; fix (n times the representatives), the product set s x r with p.1 * (a * p.2 * a^-1), and the conjugated representatives t u t^-1. uOpLevel: fixed_pro_unipotent_iff, uOpLevel_core (Iwa.swap for y x = x' t y', uniqueness from iwMem and Iwa.uniq), with Iwa.P_Nm and Iwa.z_Nm for the conjugates of Nbar_c.
\end{proof}

\begin{lemma}\label{sha-hecke-ideal-tw-augmented-forms}
In Situation~\ref{sha-hecke-ideal-tw-augmented-setup}, let $\alpha\in G_\ell$ be dominant and $T_0^{(c)}\subset T_0$ as in Lemma~\ref{sha-hecke-ideal-tw-augmented-uop}, and let $U\subset G^{(\ell)}_0$ be a subgroup.
\begin{enumerate}
\item $U_\alpha$ maps $\Lambda$-adic forms to $\Lambda$-adic forms, commutes with $f\mapsto u\cdot f$ for $u\in G^{(\ell)}_0$, and maps $\mathcal F(U)$ and $\mathcal F^{\mathrm{ord}}(U)$ into themselves.
\item If $f_0,f_1,\dots\in\mathcal F(U)$ and $f_{k+1}-f_k$ has values in $\mathfrak m_\Lambda^k\Lambda^m$ for every $k$, then there is $f\in\mathcal F(U)$ such that $f-f_k$ has values in $\mathfrak m_\Lambda^k\Lambda^m$ for every $k$.
\item Let $f$ be a $\Lambda$-adic form and $k,c\ge0$ with $\theta(t)-1\in\mathfrak m_\Lambda^k$ for all $t\in T_0^{(c)}$. Then there is $n_0$ such that $(U_z^nf)(gy)-(U_z^nf)(g)$ has entries in $\mathfrak m_\Lambda^k$ for all $n\ge n_0$, $g\in G^{(\ell)}\times G_\ell$ and $y\in\bar N_c$.
\end{enumerate}
\end{lemma}
\leanref{Development/AutomorphicArithmetic/AutomorphicVarieties/HeckeAlgebras/OrdinaryHeckeCharacters/SpStab.lean}{UpperTheory.V.Hi.spStab}
\leanref{Development/AutomorphicArithmetic/AutomorphicVarieties/BuildingsResidueChars/ExistsAdicLimitMemAutomorphicSpace.lean}{UpperTheory.exists_adic_limit_mem_automorphicSpace}
\leanref{Development/AutomorphicArithmetic/AutomorphicVarieties/BuildingsResidueChars/UOperatorPowEventuallyInvariantModPow.lean}{UpperTheory.uOperator_pow_eventually_invariant_mod_pow}
\begin{proof}
(1) The operator $U_\alpha$ is a finite sum of right translations by elements of $G_\ell$. So it commutes with left multiplication by $U(n)(L^+)$ and with $f\mapsto u\cdot f$, which acts on the variable in $G^{(\ell)}$ and on the values. Let $f$ be a $\Lambda$-adic form. By Lemma~\ref{sha-hecke-ideal-tw-augmented-uop}(1), $U_\alpha f$ is fixed by $N_0$ and $(U_\alpha f)(gt)=\theta(t)^{-1}(U_\alpha f)(g)$ for $t\in T_0$. For the continuity, fix $k$, choose $c_1$ with $f(gy)-f(g)$ in $\mathfrak m_\Lambda^k\Lambda^m$ for $y\in\bar N_{c_1}$, and $c_2$ with $\theta(t)-1\in\mathfrak m_\Lambda^k$ for $t\in T_0^{(c_2)}$, by the continuity of $\theta$; let $c=\max(c_1,c_2)$. The reduction $\bar f$ of $f$ modulo $\mathfrak m_\Lambda^k$ is fixed by $N_0$, $T_0^{(c)}$ and $\bar N_c$, hence by $I_1(c)$. By Lemma~\ref{sha-hecke-ideal-tw-augmented-uop}(2) so is $U_\alpha\bar f$, the reduction of $U_\alpha f$. So $U_\alpha f$ is a $\Lambda$-adic form, and it is fixed by $U$ if $f$ is. If $f$ is ordinary and $(U_z^{N+1}-1)f$ has values in $\mathfrak m_\Lambda^k\Lambda^m$, then so does $(U_z^{N+1}-1)U_\alpha f=U_\alpha(U_z^{N+1}-1)f$, since $U_\alpha U_z=U_{\alpha z}=U_{z\alpha}=U_zU_\alpha$ on functions fixed by $N_0$, by Lemma~\ref{sha-hecke-ideal-tw-augmented-uop}(1), and $U_\alpha$ preserves values in $\mathfrak m_\Lambda^k\Lambda^m$.

(2) As $\Lambda$ is complete, the $f_k(g)$ converge to an $f(g)$ with $f(g)-f_k(g)\in\mathfrak m_\Lambda^k\Lambda^m$ for every $k$. The conditions defining $\mathcal F(U)$ are equalities between values of $f$, possibly multiplied by $\theta(t)^{-1}$ or by a matrix $\tau_0(u)$, which pass to the limit as $\bigcap_k\mathfrak m_\Lambda^k=0$, and the continuity condition, which passes to the limit because $f(gy)-f(g)\equiv f_k(gy)-f_k(g)$ modulo $\mathfrak m_\Lambda^k$.

(3) Choose $n_0$ with $f(gy)-f(g)$ in $\mathfrak m_\Lambda^k\Lambda^m$ for $y\in\bar N_{n_0}$. The reduction $\bar f$ of $f$ modulo $\mathfrak m_\Lambda^k$ is fixed by $N_0$, by $T_0^{(c)}$, as $\theta\equiv1$ modulo $\mathfrak m_\Lambda^k$ there, and by $\bar N_{c+n}$ for $n\ge n_0$. If a function $h$ is fixed by $N_0$, $T_0^{(c)}$ and $\bar N_{d+1}$ for some $d\ge c$, then $U_zh$ is fixed by $N_0$ and $T_0^{(c)}$, by Lemma~\ref{sha-hecke-ideal-tw-augmented-uop}(1), and by $\bar N_d$, by Lemma~\ref{sha-hecke-ideal-tw-augmented-uop}(3) for $d$, as $T_0^{(d)}\subset T_0^{(c)}$. Applying this $n$ times, $U_z^n\bar f$ is fixed by $\bar N_c$, and it is the reduction of $U_z^nf$.
% spStab: spStab_U (U_a as a finite sum of rp), spStab_comm and spStab_commU, spStab_ambU (uOpMul, then uOpLevel (1) on the reduction modulo m^k by mapFn, with uOpNat), the Ord part through uOpMul in both orders and Iwa.comm. exists_adic_limit_mem_automorphicSpace: spC_lim, spC_eq and the key extension of linear maps preserving valIn. uOperator_pow_eventually_invariant_mod_pow: lower_one and lower_many (uOpLevel (2)), uOpNat for the reduction.
\end{proof}

\begin{lemma}\label{sha-hecke-ideal-tw-augmented-sections}
In Situation~\ref{sha-hecke-ideal-tw-augmented-setup}, let $U\subset G^{(\ell)}_0$ be a subgroup of finite level that satisfies the \emph{freeness condition}: if $\gamma g=g(u,b)$ with $\gamma\in U(n)(L^+)$, $g\in G^{(\ell)}\times G_\ell$, $u\in U$ and $b\in I_\ell$, then $b$ lies in $T_0$ and in the center of $G_\ell$, and $\theta(b)\tau_0(u)=1$. Let $c\ge0$, $I(c)=N_0T_0\bar N_c\subset I_\ell$ and $T_0^{(c)}$ as in Lemma~\ref{sha-hecke-ideal-tw-augmented-uop}.
\begin{enumerate}
\item There are finitely many $x_1,\dots,x_N\in G^{(\ell)}\times G_\ell$ such that every $g$ is $\gamma x_i(u,b)$ with $\gamma\in U(n)(L^+)$, $u\in U$ and $b\in I(c)$, for exactly one $i$.
\item For such $x_i$ there is a $\Lambda$-linear map $s:(\Lambda^m)^N\to\mathcal F(U)$ with $s(v)(x_i)=v_i$ for all $v$ and $i$, such that if $\theta(t)-1\in\mathfrak m_\Lambda^k$ for all $t\in T_0^{(c)}$, then $s(v)(gy)-s(v)(g)$ has entries in $\mathfrak m_\Lambda^k$ for all $v$, $g$ and $y\in\bar N_c$.
\item Let $f\in\mathcal F(U)$ and $k\ge0$ be such that $f(gy)-f(g)$ has entries in $\mathfrak m_\Lambda^k$ for all $g$ and $y\in\bar N_c$, and $f(x_i)$ has entries in $\mathfrak m_\Lambda^k$ for all $i$. Then $f$ has values in $\mathfrak m_\Lambda^k\Lambda^m$.
\end{enumerate}
\end{lemma}
\leanref{Development/AutomorphicArithmetic/AutomorphicVarieties/BuildingsResidueChars/ExistsFiniteIwahoriDoubleCosetRepresentatives.lean}{UpperTheory.exists_finite_iwahori_doubleCoset_representatives}
\leanref{Development/AutomorphicArithmetic/AutomorphicVarieties/BuildingsResidueChars/ExistsSectionAutomorphicSpaceWithPrescribedValue.lean}{UpperTheory.exists_section_automorphicSpace_with_prescribed_values}
\leanref{Development/AutomorphicArithmetic/AutomorphicVarieties/BuildingsResidueChars/MemValInOfInvariantAndValuesMemPow.lean}{UpperTheory.mem_valIn_of_invariant_and_values_mem_pow}
\begin{proof}
(1) Let $s$ be a finite set of representatives of $U(n)(L^+)\backslash(G^{(\ell)}\times G_\ell)/(U\times I_\ell)$ and $R$ one of $I_\ell/I(c)$, which is finite. Every $g$ is $\gamma x(u,b)$ with $x\in s$, $u\in U$ and $b\in I_\ell$, and $b=rb'$ with $r\in R$ and $b'\in I(c)$, so $g=\gamma\,x(1,r)\,(u,b')$. Hence the classes of $U(n)(L^+)\backslash(G^{(\ell)}\times G_\ell)/(U\times I(c))$ are among those of the finitely many $x(1,r)$; let the $x_i$ be one representative of each class.

(2) Every $b\in I_\ell$ factors as $b=y\,t\,x$ with $y\in\bar N_0$, $t\in T_0$ and $x\in N_0$, and $t=t(b)$ is unique. Then $t(1)=1$, $t(sb)=st(b)$ and $t(bs)=t(b)s$ for $s\in T_0$, $t(bx)=t(b)$ for $x\in N_0$, and $t(by)\in t(b)T_0^{(c)}$ for $y\in\bar N_c$. Define
\[
s(v)\big(\gamma x_i(u,b)\big)=\theta\big(t(b)\big)^{-1}\tau_0(u)^{-1}v_i\qquad(\gamma\in U(n)(L^+),\ u\in U,\ b\in I(c)).
\]
This is well defined. If $\gamma x_i(u,b)=\gamma'x_{i'}(u',b')$, then $i'=i$ by (1), and $\gamma'^{-1}\gamma\,x_i=x_i(u'u^{-1},b'b^{-1})$; by the freeness condition, $s_0=b'b^{-1}$ lies in $T_0$ and is central, and $\theta(s_0)\tau_0(u'u^{-1})=1$. Then $t(b')=s_0t(b)$ and $\tau_0(u')^{-1}=\theta(s_0)\tau_0(u)^{-1}$, so the two values agree. The function $s(v)$ is invariant under $U(n)(L^+)$ on the left, satisfies $u\cdot s(v)=s(v)$ for $u\in U$, is fixed by $N_0$, and satisfies $s(v)(gt)=\theta(t)^{-1}s(v)(g)$ for $t\in T_0$, by the properties of $t(b)$. For $y\in\bar N_c$, $s(v)(gy)=\theta(s')^{-1}s(v)(g)$ with $s'\in T_0^{(c)}$, which gives the congruence when $\theta\equiv1$ modulo $\mathfrak m_\Lambda^k$ on $T_0^{(c)}$; as $\theta$ is continuous, this also gives the continuity condition, so $s(v)\in\mathcal F(U)$. Finally $s(v)(x_i)=v_i$, as $t(1)=1$.

(3) Let $g=\gamma x_i(u,b)$ as in (1), and write $b=xty$ with $x\in N_0$, $t\in T_0$ and $y\in\bar N_c$. As $f$ is fixed by $u^{-1}$, $f(x_i(u,1))=\tau_0(u)^{-1}f(x_i)$, which has entries in $\mathfrak m_\Lambda^k$. As $f$ is fixed by $N_0$ and $f(ht)=\theta(t)^{-1}f(h)$, the same holds for $f(x_i(u,1)(1,xt))$. By the congruence for $y$, it holds for $f(x_i(u,1)(1,xty))$, which is $f(g)$ by the invariance of $f$ under $U(n)(L^+)$.
% exists_finite_iwahori_doubleCoset_representatives: the setoid repsRel, iwFin for the finite index of Iw c in Iw 0, Finite.exists_equiv_fin. exists_section_automorphicSpace_with_prescribed_values: exists_iwahori_torus_projection (torus_fn, torus_uniq, torus_ex, Iwa.swap for t(b y)), exists_linear_section_automorphicFunctions_of_free (well defined by Frame.Free: hi, hs, he, ht, htau), and hdiff for the congruences. mem_valIn_of_invariant_and_values_mem_pow: iwMem for b = x t y, then a1 to a4.
\end{proof}

\begin{lemma}\label{sha-hecke-ideal-tw-augmented-projector}
Let $A$ be a complete local noetherian ring with maximal ideal $\mathfrak m$ and finite residue field, $M$ an $A$-module, $T$ an $A$-linear endomorphism of $M$, and $V_0\supset V_1\supset\dots$ submodules of $M$. Suppose:
\begin{enumerate}
\item $\mathfrak m^kM\subset V_k$, $\mathfrak m^kV_1\subset V_{k+1}$ and $T(V_k)\subset V_k$ for every $k$, and $\bigcap_kV_k=0$;
\item for every sequence $(a_k)$ in $M$ with $a_{k+1}-a_k\in V_k$ for all $k$ there is $a\in M$ with $a-a_k\in V_k$ for all $k$;
\item for every $k$ there are a finite set $I$ and $A$-linear maps $s:A^I\to M$ and $\operatorname{ev}:M\to A^I$, with $\operatorname{ev}(V_j)\subset(\mathfrak m^j)^I$ for all $j$, such that for every $f\in M$ there is $n_0$ with $T^nf-s(\operatorname{ev}(T^nf))\in V_k$ for all $n\ge n_0$.
\end{enumerate}
Let $O$ be the set of $f\in M$ such that for every $k$ there is $n\ge1$ with $T^nf-f\in V_k$. Then $O$ is a finite free $A$-submodule of $M$, $T$ maps $O$ bijectively onto itself, $O\cap V_j\subset\mathfrak m^jO$ for every $j$, and there is an $A$-linear map $e:M\to O$, equal to the identity on $O$, such that for every $f\in M$ and $k$ there is $J$ with $T^{j!}f-ef\in V_k$ for all $j\ge J$.
\end{lemma}
\leanref{Development/AutomorphicArithmetic/AutomorphicVarieties/HeckeAlgebras/OrdinaryHeckeCharacters/OrdFree.lean}{UpperTheory.V.Hi.ordAbs}
\leanref{Development/AutomorphicArithmetic/AutomorphicVarieties/BuildingsResidueChars/ExistsOrdinaryProjectorOfEventuallyPeriodic.lean}{UpperTheory.exists_ordinaryProjector_of_eventually_periodic}
\leanref{Development/AutomorphicArithmetic/AutomorphicVarieties/BuildingsResidueChars/ExistsLinearPreimageOfSuccessiveApproximation.lean}{UpperTheory.exists_linear_preimage_of_successive_approximation}
\begin{proof}
\emph{Periodicity.} Let $f\in M$ and $k\ge0$, and take $I$, $s$, $\operatorname{ev}$ and $n_0$ as in (3). The ring $A/\mathfrak m^k$ is finite, so there are $x<y$ with $\operatorname{ev}(T^{n_0+y}f)\equiv\operatorname{ev}(T^{n_0+x}f)$ modulo $\mathfrak m^k$. As $s((\mathfrak m^k)^I)\subset\mathfrak m^kM\subset V_k$,
\[
T^{n_0+y}f-T^{n_0+x}f=\big(T^{n_0+y}f-s\operatorname{ev}T^{n_0+y}f\big)-\big(T^{n_0+x}f-s\operatorname{ev}T^{n_0+x}f\big)+s\operatorname{ev}\big(T^{n_0+y}f-T^{n_0+x}f\big)\in V_k.
\]
So there are $a\ge0$ and $p\ge1$ with $T^{a+p}f-T^af\in V_k$. As $T$ preserves $V_k$, then $T^{b+pc}f-T^bf\in V_k$ for all $b\ge a$ and $c\ge0$, and so $T^{i!}f-T^{ap}f\in V_k$ for $i\ge ap+p$, since then $p$ divides $i!$ and $i!\ge ap$.

\emph{The projector.} By the periodicity, for each $k$ there is $J_k$ such that the $T^{j!}f$, $j\ge J_k$, agree modulo $V_k$; take $J_k$ increasing in $k$. By (2) there is $ef\in M$ with $ef-T^{j!}f\in V_k$ for $j\ge J_k$, unique since $\bigcap_kV_k=0$; uniqueness makes $f\mapsto ef$ $A$-linear. With $a$ and $p$ as above for $k$, $T^p(ef)-ef\equiv T^{j!+p}f-T^{j!}f\equiv0$ modulo $V_k$ for $j!\ge a$, so $ef\in O$. If $f\in V_k$ then $ef\in V_k$. If $T^nf-f\in V_k$ with $n\ge1$, then $T^{j!}f-f\in V_k$ for $j\ge n$, as $n$ divides $j!$, so $ef-f\in V_k$. Hence $e$ is the identity on $O$.

\emph{Approximation.} For $f\in O$ and $k$, with $I$, $s$ and $\operatorname{ev}$ as in (3) for $k$, $f-e(s(\operatorname{ev}f))\in V_k$. Indeed, for $j$ large, $T^{j!}f-f$ and $T^{j!}f-s\operatorname{ev}T^{j!}f$ lie in $V_k$, and $s\operatorname{ev}(T^{j!}f-f)\in\mathfrak m^kM\subset V_k$; so $f-s\operatorname{ev}f\in V_k$, and applying $e$ gives the claim.

\emph{Generation.} Fix $I_1$, $s_1$ and $\operatorname{ev}_1$ as in (3) for $k=1$, and let $\sigma=e\circ s_1:A^{I_1}\to O$. Let $f\in O\cap V_j$. Take $I'$, $s'$ and $\operatorname{ev}'$ for $k=j+1$, so $\operatorname{ev}'f\in(\mathfrak m^j)^{I'}$, and let $\delta_i$ be the standard basis of $A^{I'}$. Put $a=\sum_i(\operatorname{ev}'f)_i\operatorname{ev}_1(es'\delta_i)\in(\mathfrak m^j)^{I_1}$. Then
\[
f-\sigma(a)=\big(f-es'\operatorname{ev}'f\big)+\sum_i(\operatorname{ev}'f)_i\big(es'\delta_i-\sigma(\operatorname{ev}_1(es'\delta_i))\big),
\]
where the first term lies in $V_{j+1}$ and each bracket in $V_1$, by the approximation, so $f-\sigma(a)\in V_{j+1}$. Repeating, with $r_0=f$ and $r_{t+1}=r_t-\sigma(b_t)$ for such $b_t\in(\mathfrak m^{j+t})^{I_1}$, the series $\sum_tb_t$ converges in $A^{I_1}$, as $A$ is complete, to some $b\in(\mathfrak m^j)^{I_1}$, and $f-\sigma(b)\in V_N+\mathfrak m^NM\subset V_N$ for every $N$, so $f=\sigma(b)$. Hence $O=\sigma(A^{I_1})$ is finite, taking $j=0$, and $O\cap V_j\subset\mathfrak m^jO$.

\emph{Freeness.} Let $\iota:O\to A^{I_1}$ be the restriction of $\operatorname{ev}_1$. For $f\in O$, $f-\sigma\iota f\in O\cap V_1\subset\mathfrak mO$ by the approximation. By Nakayama's lemma $\sigma\circ\iota$ is surjective, and a surjective endomorphism of a finite module is bijective. So $O$ is a direct summand of $A^{I_1}$, hence projective, hence free over the local ring $A$.

\emph{The operator.} If $f\in O$, then $Tf\in O$, as $T^n(Tf)-Tf=T(T^nf-f)$. For $f\in O$ choose $n\ge1$ with $T^nf-f\in V_1$; then $f-T(T^{n-1}f)\in O\cap V_1\subset\mathfrak mO$. So $T$ is surjective on $O$ by Nakayama's lemma, hence bijective.
% ordAbs: exists_eventually_periodic_mod_V (Finite.exists_ne_map_eq_of_infinite on the classes modulo m^k, Ideal.finite_quotient_pow); exists_ordinaryProjector_of_eventually_periodic (periodic_difference_shift, periodic_difference_mul_period, factorial_powers_cauchy, the limit L with Nm f k the sums of the J f i, huniq, hPadd, hPsmul); heO, heV, heid; happ; hstep with linearMap_pi_eq_sum_single; exists_linear_preimage_of_successive_approximation (the remainders r n, IsPrecomplete for the sum of the b n, hsep); ordFree_bij (le_of_le_smul_of_le_jacobson_bot, OrzechProperty.injective_of_surjective_endomorphism); Module.Projective.of_split and Module.free_of_flat_of_isLocalRing; hmaps and hTO.
\end{proof}

\begin{lemma}\label{sha-hecke-ideal-tw-augmented-ordinary}
In Situation~\ref{sha-hecke-ideal-tw-augmented-setup}, let $U\subset G^{(\ell)}_0$ be a subgroup of finite level that satisfies the freeness condition of Lemma~\ref{sha-hecke-ideal-tw-augmented-sections}. Then:
\begin{enumerate}
\item $\mathcal F^{\mathrm{ord}}(U)$ is finite and free over $\Lambda$, and $U_z$ maps it bijectively onto itself;
\item for $j\ge0$, every element of $\mathcal F^{\mathrm{ord}}(U)$ with values in $\mathfrak m_\Lambda^j\Lambda^m$ lies in $\mathfrak m_\Lambda^j\mathcal F^{\mathrm{ord}}(U)$;
\item there is a $\Lambda$-linear map $e_U:\mathcal F(U)\to\mathcal F^{\mathrm{ord}}(U)$, equal to the identity on $\mathcal F^{\mathrm{ord}}(U)$, such that for every $f\in\mathcal F(U)$ and $k\ge0$ there is $N'$ such that $U_z^{N!}f-e_Uf$ has values in $\mathfrak m_\Lambda^k\Lambda^m$ for all $N\ge N'$.
\end{enumerate}
\end{lemma}
\leanref{Development/AutomorphicArithmetic/AutomorphicVarieties/HeckeAlgebras/OrdinaryHeckeCharacters/OrdFree.lean}{UpperTheory.V.Hi.ordStruct}
\begin{proof}
Apply Lemma~\ref{sha-hecke-ideal-tw-augmented-projector} to $A=\Lambda$, which is a power series ring over $\mathcal O$ (proof of Lemma~\ref{sha-hecke-algebra}), so a complete local noetherian ring with finite residue field $k$; to $M=\mathcal F(U)$; to $T=U_z$; and to the submodules $V_k$ of the elements of $\mathcal F(U)$ with values in $\mathfrak m_\Lambda^k\Lambda^m$. The set $O$ of that lemma is $\mathcal F^{\mathrm{ord}}(U)$, by the definition of ordinary forms, so the lemma gives (1), (2) and (3) once its hypotheses hold.

(1) The first two inclusions are clear. The operator $U_z$ maps $\mathcal F(U)$ into itself by Lemma~\ref{sha-hecke-ideal-tw-augmented-forms}(1), and $V_k$ into itself, as it is a sum of right translations. By Krull's intersection theorem $\bigcap_k\mathfrak m_\Lambda^k=0$, so $\bigcap_kV_k=0$.

(2) This is Lemma~\ref{sha-hecke-ideal-tw-augmented-forms}(2).

(3) Fix $k$. As $\theta$ is continuous, there is $c$ with $\theta(t)-1\in\mathfrak m_\Lambda^k$ for $t\in T_0^{(c)}$. Let $x_1,\dots,x_N$ and $s$ be as in Lemma~\ref{sha-hecke-ideal-tw-augmented-sections}(1) and (2) for $c$, let $I=\{1,\dots,N\}\times\{1,\dots,m\}$, and let $\operatorname{ev}(f)=(f(x_i))_i$; then $\operatorname{ev}(V_j)\subset(\mathfrak m_\Lambda^j)^I$. Let $f\in\mathcal F(U)$. By Lemma~\ref{sha-hecke-ideal-tw-augmented-forms}(3) there is $n_0$ such that $(U_z^nf)(gy)-(U_z^nf)(g)$ has entries in $\mathfrak m_\Lambda^k$ for $n\ge n_0$ and $y\in\bar N_c$. For such $n$, $h=U_z^nf-s(\operatorname{ev}(U_z^nf))$ lies in $\mathcal F(U)$, satisfies the same congruences, by Lemma~\ref{sha-hecke-ideal-tw-augmented-sections}(2), and vanishes at the $x_i$. By Lemma~\ref{sha-hecke-ideal-tw-augmented-sections}(3), $h\in V_k$.
% ordStruct: S = Sp U with T the restriction of U_z (spStab), V k = valIn k, hanti, hm, hm1, hTV (ordFree_valU), hsep (Ideal.iInf_pow_eq_bot_of_isLocalRing), hcomp (exists_adic_limit_mem_automorphicSpace), hlev (exists_finite_iwahori_doubleCoset_representatives, exists_section_automorphicSpace_with_prescribed_values, uOperator_pow_eventually_invariant_mod_pow, mem_valIn_of_invariant_and_values_mem_pow), hO through mem_ord_iff; then ordAbs and the transport along comapSubtypeEquivOfLe. The Lean states (2) as hfilt; ordFree adds that every U_a, a in P, is bijective on Ord U.
\end{proof}

\begin{lemma}\label{sha-hecke-ideal-tw-augmented-spherical}
Let $F$ be a nonarchimedean local field, $n\ge1$, $K=GL_n(\mathcal O_F)$, and $\pi$ a representation of $GL_n(F)$ on a module $P$ over a commutative ring. Then $KxK/K$ is finite for every $x\in GL_n(F)$, and the double coset operators $[KxK]$ and $[KyK]$ of Definition~\ref{sha-hecke-ideal-tw-iwahori}(3) commute on $P^K$ for all $x,y\in GL_n(F)$.
\end{lemma}
\leanref{Development/AutomorphicArithmetic/AutomorphicVarieties/HeckeAlgebras/OrdinaryHeckeCharacters/SphComm.lean}{UpperTheory.V.Hi.sphComm}
\leanref{Development/AutomorphicArithmetic/AutomorphicVarieties/HeckeAlgebras/OrdinaryHeckeCharacters/DcGelfand.lean}{UpperTheory.V.Hi.dcGelfand}
\begin{proof}
The set $KxK/K$ is the image of the compact set $Kx$ in the discrete space $GL_n(F)/K$, so it is finite. For $v\in P^K$, $[KxK][KyK]v=\sum_{p,q}\pi(\tilde p\tilde q)v$, the sum over $p\in KxK/K$ and $q\in KyK/K$ with representatives $\tilde p$ and $\tilde q$; so it is $\sum_rN_{x,y}(r)\pi(\tilde r)v$, the sum over $r\in GL_n(F)/K$, where $N_{x,y}(r)$ is the number of pairs $(p,q)$ with $\tilde pq=r$. The transpose $g\mapsto{}^tg$ reverses products, preserves $K$, and preserves every double coset $KgK$, by the elementary divisor theorem. Let $r=\zeta K$ and write ${}^t\zeta=k\zeta k'$ with $k,k'\in K$. For a pair $(p,q)$ for $(x,y)$ with $\tilde pq=r$, let $p'=k^{-1}\,{}^t(\tilde p^{-1}\zeta)K$ and $q'=\tilde p'^{-1}r$. Then $p'\in KyK/K$, as $\tilde p^{-1}\zeta\in KyK$, and $q'\in KxK/K$, so $(p',q')$ is a pair for $(y,x)$ with $\tilde p'q'=r$; and $p'$ determines $p$, since ${}^t(\tilde p_2^{-1}\tilde p_1)\in K$ forces $p_1=p_2$, while $q=\tilde p^{-1}r$. Hence $N_{x,y}(r)\le N_{y,x}(r)$, and by symmetry the two are equal.
% sphComm: sphC_fin; dcGelfand (dcG_two for the double sum, dcG_le for the injection phi p = k^-1 s (p.out^-1 z), dcG_sum) with s x = invT (x^-1), the transpose; sphC_invT and sphC_cartan (Smith normal form, KStar.Selmer.exists_smith) for the double cosets. The cardinality hypothesis of dcGelfand, supplied through sphC_le, is not used in its proof.
\end{proof}

\subsection{Hida's control theorem}\label{subsection-sha-hida-control}

\begin{lemma}[Finiteness of class numbers {\cite[Theorem~5.1]{Bor63}}]\label{sha-hecke-ideal-tw-augmented-hida-classes}
Let $L$ be a CM field with maximal totally real subfield $L^+$, let $n\ge1$, let $U(n)$ be the unitary group over $L^+$ of the Hermitian form $\sum_ix_i\bar x_i$ on $L^n$, and let $C$ be an open subgroup of $U(n)(\mathbf A^\infty_{L^+})$. Then there is a finite set $\Sigma\subset U(n)(\mathbf A_{L^+})$ such that every element of $U(n)(\mathbf A_{L^+})$ is $\gamma sxk$ with $\gamma\in U(n)(L^+)$, $s\in\Sigma$, $x\in U(n)(L^+\otimes\mathbf R)$ and $k\in C$.
\end{lemma}
\leanref{Development/AbelianVarietyTheory/Proofs/IwasawaDecomposition/DefiniteUnitaryClassSetFinite.lean}{GlobalTheory.DefU.ClassFin.classSet_finite}

\begin{definition}\label{sha-hecke-ideal-tw-augmented-hida-weight}
In Situation~\ref{sha-hecke-ideal-tw-augmented-setup}, let $T\supset T_0$ be the group of diagonal matrices of $G_\ell$ and $N\supset N_0$ the group of its upper triangular unipotent matrices. Let $\mathcal O'\supset\mathcal O$ be the ring of integers of a finite extension of $\mathbf Q_\ell$ in $\bar{\mathbf Q}_\ell$, with maximal ideal $\mathfrak m'$, viewed in $\mathbf C$ through $\iota$.
\begin{enumerate}
\item A \emph{weight datum} over $\mathcal O'$ is a triple $(\xi,\chi,c_0)$ of a representation $\xi$ of $G_\ell$ on $\mathbf C^{d+1}$ for some $d\ge0$, with coordinates indexed by $0,\dots,d$, a character $\chi:T\to\mathbf C^\times$ and an integer $c_0\ge0$, such that:
\begin{enumerate}
\item $\xi(y)$ has entries in $\mathcal O'$ for $y$ in $N_0$, in $T_0$ or in $\bar N_{c_0}$, and so has $\chi(\alpha)^{-1}\xi(\alpha)$ for every dominant $\alpha$ (Lemma~\ref{sha-hecke-ideal-tw-augmented-uop});
\item the row $0$ of $\xi(y)$ is $(1,0,\dots,0)$ for $y\in N$ and $(\chi(t),0,\dots,0)$ for $y=t\in T$, and $\xi(t)=\chi(t)$ for $t\in T$ central in $G_\ell$;
\item the entries of $\chi(z)^{-1}\xi(z)$ outside the column $0$ lie in $\mathfrak m'$;
\item for every $k\ge0$ there is $c\ge0$ such that $\xi(y)-1$ has entries in $\mathfrak m'^k$ for $y\in\bar N_c$ and $\chi(t)-1\in\mathfrak m'^k$ for $t\in T_0^{(c)}$.
\end{enumerate}
\item Let $(\xi,\chi,c_0)$ be a weight datum, $V^\xi$ the space of Definition~\ref{sha-hecke-ideal-tw-augmented-parts}(2) formed for $\xi$, $U\subset G^{(\ell)}_0$ a subgroup, $c\ge0$, and $\phi$ a character of $T_0$ with values in $\mathbf C^\times$. The space $\mathrm{Cl}(U,c,\phi)$ of \emph{classical forms} is the space of the $F\in V^\xi$ such that $F(\gamma g)=\xi(\gamma_\ell)F(g)$ for $\gamma\in U(n)(L^+)$, $u\cdot F=F$ for $u\in U$, $F(gy)=F(g)$ for $y\in N_0$ and for $y\in\bar N_c$, and $F(gt)=\phi(t)^{-1}F(g)$ for $t\in T_0$. For dominant $\alpha$ let $U_\alpha$ be the operator of Lemma~\ref{sha-hecke-ideal-tw-augmented-uop} on $V^\xi$, for right translation on $G_\ell$, and $U'_\alpha=\chi(\alpha)^{-1}U_\alpha$. The \emph{ordinary part} $\mathrm{Cl}^{\mathrm{ord}}(U,c,\phi)$ is the sum, over the nonzero $\mu\in\mathbf C$ such that $\mu$ and $\mu^{-1}$ are integral over $\mathcal O'$, of the generalized $\mu$-eigenspaces of $U'_z$ in $\mathrm{Cl}(U,c,\phi)$.
\item The \emph{control map} takes $F\in V^\xi$, with components $F_1,\dots,F_m$, to the function $\operatorname{ctrl}(F):G^{(\ell)}\times G_\ell\to\mathbf C^m$ whose $i$-th entry at $g$ is the coordinate $0$ of $\xi(g_\ell)^{-1}F_i(g)$.
\item For a continuous character $\theta'$ of $T_0$ with values in $\mathcal O'^\times$ and a subgroup $U\subset G^{(\ell)}_0$, $\mathcal F^{\mathrm{ord}}_{\mathcal O'}(U,\theta')$ is the $\mathcal O'$-module defined as $\mathcal F^{\mathrm{ord}}(U)$ in Situation~\ref{sha-hecke-ideal-tw-augmented-setup}(2), with $\mathcal O'$, $\theta'$ and $\tau_0$ composed with $\mathcal O\subset\mathcal O'$ in place of $\Lambda$, $\theta$ and $\tau_0$.
\end{enumerate}
\end{definition}
\leanref{Development/AutomorphicArithmetic/AutomorphicVarieties/HeckeAlgebras/ContinuousTorusCharacter.lean}{UpperTheory.V.Hi.WtRep}
\leanref{Development/AutomorphicArithmetic/AutomorphicVarieties/HeckeAlgebras/ContinuousTorusCharacter.lean}{UpperTheory.V.Hi.Frame.Cl}
\leanref{Development/AutomorphicArithmetic/AutomorphicVarieties/HeckeAlgebras/ContinuousTorusCharacter.lean}{UpperTheory.V.Hi.Frame.UF}
\leanref{Development/AutomorphicArithmetic/AutomorphicVarieties/HeckeAlgebras/ContinuousTorusCharacter.lean}{UpperTheory.V.Hi.Frame.ClOrd}
\leanref{Development/AutomorphicArithmetic/AutomorphicVarieties/HeckeAlgebras/ContinuousTorusCharacter.lean}{UpperTheory.V.Hi.Frame.ctrl}

\begin{lemma}\label{sha-hecke-ideal-tw-augmented-hida-classical}
In Situation~\ref{sha-hecke-ideal-tw-augmented-setup}, let $(\xi,\chi,c_0)$ be a weight datum over $\mathcal O'$ (Definition~\ref{sha-hecke-ideal-tw-augmented-hida-weight}), let $U\subset G^{(\ell)}_0$ be a subgroup of finite level, let $c\ge c_0$, and let $\phi$ be a character of $T_0$ with values in $\mathcal O'^\times$.
\begin{enumerate}
\item If $\phi$ is trivial on $T_0^{(c)}$, then $\mathrm{Cl}(U,c,\phi)$ is finite-dimensional, and for dominant $\alpha$ the operator $U'_\alpha$ maps it into itself, with eigenvalues integral over $\mathcal O'$.
\item Let $u_1,\dots,u_r\in G^{(\ell)}_0$, and let $F\in\mathrm{Cl}(U,c,\phi)$ be nonzero with $\sum_qu_q\cdot F=\mu F$ for some $\mu\in\mathbf C$. Then $\mu$ is integral over $\mathcal O'$.
\end{enumerate}
\end{lemma}
\leanref{Development/AutomorphicArithmetic/AutomorphicVarieties/HeckeAlgebras/OrdinaryHeckeCharacters/ClFinMulVec.lean}{UpperTheory.V.Hi.clFin}
\leanref{Development/AutomorphicArithmetic/AutomorphicVarieties/HeckeAlgebras/HeckeFixedLevelSpaces/FreeMap.lean}{UpperTheory.V.Hi.Sp.hecke_eigenvalue_isIntegral}
\begin{proof}
For $F\in V^\xi$, $g\in G^{(\ell)}\times G_\ell$ and $1\le i\le m$ let $F^\flat_i(g)=\xi(g_\ell)^{-1}F_i(g)\in\mathbf C^{d+1}$, and for an $\mathcal O'$-submodule $L\subset\mathbf C$ say that $F^\flat(g)$ lies in $L$ if all the coordinates of all the $F^\flat_i(g)$ do. For $F\in\mathrm{Cl}(U,c,\phi)$ and $\gamma\in U(n)(L^+)$, $F^\flat(\gamma g)=F^\flat(g)$, and for $u\in U$, $F^\flat_i(g(u,1))=\sum_j\tau_0(u^{-1})_{ij}F^\flat_j(g)$, as $u^{-1}\cdot F=F$; the entries of $\tau_0$ lie in $\mathcal O$.

(1) \emph{Values at finitely many points.} Let $F\in\mathrm{Cl}(U,c,\phi)$. Every element of $I_1(c)$ is a product $nty$ with $n\in N_0$, $t\in T_0^{(c)}$ and $y\in\bar N_c$, and $\phi(t)=1$, so $F$ is invariant under right translation by $I_1(c)$; and $\xi$ of such an element has entries in $\mathcal O'$, as $\bar N_c\subset\bar N_{c_0}$. Let $s$ be a finite set as in the definition of finite level, $R$ a finite set of representatives of $I_\ell/I_1(c)$, and $\Sigma$ the finite set of the $x(1,r)$ with $x\in s$ and $r\in R$. We show that if $F^\flat(y)$ lies in $L$ for every $y\in\Sigma$, then $F^\flat(g)$ lies in $L$ for every $g$. Write $g=\gamma x(u,b)$ with $\gamma\in U(n)(L^+)$, $x\in s$, $u\in U$ and $b\in I_\ell$, and $b=re^{-1}$ with $r\in R$ and $e\in I_1(c)$. Then $g=\gamma\,x(1,r)\,(u,1)\,(1,e^{-1})$, and $F_i(g)=\xi(\gamma_\ell)\sum_j\tau_0(u^{-1})_{ij}F_j(x(1,r))$, so $F^\flat_i(g)=\xi(e)\sum_j\tau_0(u^{-1})_{ij}F^\flat_j(x(1,r))$, which has coordinates in $L$.

\emph{Finite dimension.} Taking $L=0$, an $F\in\mathrm{Cl}(U,c,\phi)$ that vanishes on $\Sigma$ is $0$. So $F\mapsto(F(y))_{y\in\Sigma}$ is injective, and $\mathrm{Cl}(U,c,\phi)$ is finite-dimensional.

\emph{The operators.} The operator $U_\alpha$ is a finite sum of right translations by elements $x\alpha$ of $G_\ell$. So it commutes with left multiplication by $U(n)(L^+)$ and with $F\mapsto u\cdot F$, and it preserves the condition $F(\gamma g)=\xi(\gamma_\ell)F(g)$. Let $F\in\mathrm{Cl}(U,c,\phi)$. As $F$ is fixed by $N_0$ and by $I_1(c)$, and $T_0$ acts on it by $\phi^{-1}$, so does $U_\alpha F$, by Lemma~\ref{sha-hecke-ideal-tw-augmented-uop}(1) and (2). Hence $U'_\alpha F\in\mathrm{Cl}(U,c,\phi)$.

\emph{Integral eigenvalues.} Let $F\in\mathrm{Cl}(U,c,\phi)$ be nonzero with $U'_\alpha F=\mu F$, and let $L$ be the $\mathcal O'$-span of the coordinates of the $F^\flat_i(y)$, $y\in\Sigma$. It is finitely generated, and nonzero by the finite dimension. For each representative $x$ of $N_0/\alpha N_0\alpha^{-1}$, $\xi(g_\ell)^{-1}=\xi(x)\xi(\alpha)\xi(g_\ell x\alpha)^{-1}$, so
\[
\mu F^\flat_i(g)=\sum_x\xi(x)\,\chi(\alpha)^{-1}\xi(\alpha)\,F^\flat_i(g(1,x\alpha)),
\]
whose coordinates lie in $L$ by the first step and Definition~\ref{sha-hecke-ideal-tw-augmented-hida-weight}(1)(a). So $\mu L\subset L$, and $\mu$ is integral over $\mathcal O'$, as $L$ is a nonzero finitely generated $\mathcal O'$-submodule of the field $\mathbf C$.

(2) For $g\in G^{(\ell)}\times G_\ell$ let $N(g)$ be the $\mathcal O'$-span of the coordinates of the $F^\flat_i(g)$. Then $N(\gamma g)=N(g)$ for $\gamma\in U(n)(L^+)$, and $N(g(u,1))\subset N(g)$ for $u\in U$. For $y$ in $N_0$, in $\bar N_c$ or in $T_0$, $F_i(g(1,y))$ is $F_i(g)$ or $\phi(y)^{-1}F_i(g)$, so $F^\flat_i(g(1,y))$ is $\xi(y)^{-1}F^\flat_i(g)$, up to the factor $\phi(y)^{-1}\in\mathcal O'$, and $N(g(1,y))\subset N(g)$. Let $R'$ be a finite set of representatives of $\bar N_0/\bar N_c$. Every $b\in I_\ell$ has $b^{-1}=nty$ with $n\in N_0$, $t\in T_0$ and $y\in\bar N_0$, and $y^{-1}=ry'$ with $r\in R'$ and $y'\in\bar N_c$. So $g=\gamma x(u,b)=\gamma\,x(1,r)(1,y')(1,t^{-1})(1,n^{-1})(u,1)$, and $N(g)\subset N(x(1,r))$. Hence the sum $N_0$ of all the $N(g)$ is the sum of the finitely many $N(x(1,r))$ with $x\in s$ and $r\in R'$, so it is finitely generated, and it is nonzero as $F\ne0$. As
\[
\mu F^\flat_i(g)=\sum_q\sum_j\tau_0(u_q)_{ij}F^\flat_j(g(u_q,1)),
\]
$\mu N_0\subset N_0$, and $\mu$ is integral over $\mathcal O'$.
% clFin: clFin_key (the finite set S of the x (1, r), with iwFin for the finite index of I1 c in Iw 0, hfix through iwMem, hint, e1 to e4), clFin_zero, the evaluation ev at S for the finite dimension; uOpSum and uOpLevel (2) for the operators; the O'-module L, hmuL (int_N and int_P), isIntegral_of_smul_mem_submodule. The Lean's hypotheses also include theta, the relation theta = phi chi and the freeness condition, which the proof does not use. hecke_eigenvalue_isIntegral (FreeMap.lean): Np, mvG, mvU, mvN, mvNm, mvT (with W.int_N, int_Nm, integral_torusCharacter), the cover through iwahori_decomposition_N_T_Nm and exists_finite_coset_reps_Nm, hfg, hgen, hstep.
\end{proof}

\begin{lemma}\label{sha-hecke-ideal-tw-augmented-hida-base-change}
In Situation~\ref{sha-hecke-ideal-tw-augmented-setup}, let $y:\Lambda\to\mathcal O'$ be an $\mathcal O$-algebra map with $y(\mathfrak m_\Lambda)\subset\mathfrak m'$, with $\mathcal O'$ as in Definition~\ref{sha-hecke-ideal-tw-augmented-hida-weight}, and let $U\subset G^{(\ell)}_0$ be a subgroup of finite level that satisfies the freeness condition of Lemma~\ref{sha-hecke-ideal-tw-augmented-sections}. Then $f\mapsto y\circ f$ maps $\mathcal F^{\mathrm{ord}}(U)$ into $\mathcal F^{\mathrm{ord}}_{\mathcal O'}(U,y\circ\theta)$, and it induces an isomorphism $\mathcal O'\otimes_\Lambda\mathcal F^{\mathrm{ord}}(U)\cong\mathcal F^{\mathrm{ord}}_{\mathcal O'}(U,y\circ\theta)$.
\end{lemma}
\leanref{Development/AutomorphicArithmetic/AutomorphicVarieties/HeckeAlgebras/OrdinaryHeckeCharacters/OrdBCMapOrd.lean}{UpperTheory.V.Hi.ordBC}
\begin{proof}
Let $\theta'=y\circ\theta$, a continuous character of $T_0$ with values in $\mathcal O'^\times$, as $y(\mathfrak m_\Lambda^k)\subset\mathfrak m'^k$. Applying $y$ to $\theta(b)\tau_0(u)=1$ shows that $U$ satisfies the freeness condition for $\theta'$. Lemmas~\ref{sha-hecke-ideal-tw-augmented-forms}, \ref{sha-hecke-ideal-tw-augmented-sections} and~\ref{sha-hecke-ideal-tw-augmented-ordinary} hold, with the same proofs, for $\mathcal O'$ and $\theta'$ in place of $\Lambda$ and $\theta$: they use only that $\Lambda$ is a complete local noetherian ring with finite residue field and that $\theta$ is continuous. Write $\mathcal M=\mathcal F^{\mathrm{ord}}(U)$ and $\mathcal M'=\mathcal F^{\mathrm{ord}}_{\mathcal O'}(U,\theta')$, finite over $\Lambda$ and over $\mathcal O'$ by Lemma~\ref{sha-hecke-ideal-tw-augmented-ordinary}(1).

\emph{The map.} The conditions defining $\mathcal F(U)$ and $\mathcal M$ are equalities between values of a function, possibly multiplied by $\theta(t)^{-1}$ or by a matrix $\tau_0(u)$, and conditions that functions built linearly from it, such as $f(gy')-f(g)$ and $(U_z^{N+1}-1)f$, have values in $\mathfrak m_\Lambda^k$. Applying $y$, which commutes with $U_z$ and maps $\mathfrak m_\Lambda^k$ into $\mathfrak m'^k$, gives them for $y\circ f$, with $\theta'$ in place of $\theta$. So $y\circ f$ lies in $\mathcal F_{\mathcal O'}(U)$ for $f\in\mathcal F(U)$, and in $\mathcal M'$ for $f\in\mathcal M$.

\emph{Approximation.} Let $c$ be such that $\theta(t)-1\in\mathfrak m_\Lambda$ for $t\in T_0^{(c)}$; then $\theta'(t)-1\in\mathfrak m'$ for these $t$. Let $x_1,\dots,x_N$ and $s$ be as in Lemma~\ref{sha-hecke-ideal-tw-augmented-sections}(1) and (2) for $c$, and $P=\{1,\dots,N\}\times\{1,\dots,m\}$. For $p\in P$ let $s_p=s(\delta_p)$, for the standard basis $(\delta_p)$ of $(\Lambda^m)^N$, and $w_p=e_U(s_p)\in\mathcal M$, with $e_U$ as in Lemma~\ref{sha-hecke-ideal-tw-augmented-ordinary}(3). For a function $f$ and $p=(i,j)$ let $\operatorname{ev}_p(f)$ be the entry $j$ of $f(x_i)$; so $\operatorname{ev}_{p'}(s_p)$ is $1$ if $p'=p$ and $0$ otherwise. We show that $f-\sum_p\operatorname{ev}_p(f)\,y\circ w_p\in\mathfrak m'\mathcal M'$ for $f\in\mathcal M'$. As $f$ is ordinary, $f$ is fixed by the map $e'_U$ of Lemma~\ref{sha-hecke-ideal-tw-augmented-ordinary}(3) for $\mathcal O'$, so $U_z^{n!}f-f$ has values in $\mathfrak m'$ for $n$ large. By Lemma~\ref{sha-hecke-ideal-tw-augmented-forms}(3), $(U_z^{n!}f)(gy')-(U_z^{n!}f)(g)$ has entries in $\mathfrak m'$ for $n$ large and $y'\in\bar N_c$; and $U_z^{n!}s_p-w_p$ has values in $\mathfrak m_\Lambda$ for $n$ large and every $p$. Fix such an $n$, and let $h=\sum_p\operatorname{ev}_p(f)\,y\circ s_p$. Then $f-h\in\mathcal F_{\mathcal O'}(U)$ vanishes at the $x_i$, and $(f-h)(gy')-(f-h)(g)$ has entries in $\mathfrak m'$ for $y'\in\bar N_c$: for $f$ by the two congruences above, and for $h$ by the congruence of Lemma~\ref{sha-hecke-ideal-tw-augmented-sections}(2). By Lemma~\ref{sha-hecke-ideal-tw-augmented-sections}(3), $f-h$ has values in $\mathfrak m'$. Let $g_f=\sum_p\operatorname{ev}_p(f)\,y\circ w_p$. Then
\[
f-g_f=U_z^{n!}(f-h)-(U_z^{n!}f-f)+\sum_p\operatorname{ev}_p(f)\,y\circ(U_z^{n!}s_p-w_p)
\]
has values in $\mathfrak m'$, as $U_z$ preserves values in $\mathfrak m'$. As $f-g_f\in\mathcal M'$, it lies in $\mathfrak m'\mathcal M'$ by Lemma~\ref{sha-hecke-ideal-tw-augmented-ordinary}(2). The same argument, with the identity of $\Lambda$ in place of $y$, gives $f-\sum_p\operatorname{ev}_p(f)w_p\in\mathfrak m_\Lambda\mathcal M$ for $f\in\mathcal M$.

\emph{The isomorphism.} Let $\psi:\mathcal M\to\mathcal M$ be $f\mapsto\sum_p\operatorname{ev}_p(f)w_p$. As $f-\psi(f)\in\mathfrak m_\Lambda\mathcal M$, $\psi$ is surjective by Nakayama's lemma, hence bijective, as $\mathcal M$ is finite. Let $L:\mathcal O'\otimes_\Lambda\mathcal M\to\mathcal M'$ be $b\otimes f\mapsto b\,(y\circ f)$, and $R:\mathcal M'\to\mathcal O'\otimes_\Lambda\mathcal M$ be $f\mapsto\sum_p\operatorname{ev}_p(f)\otimes w_p$. As $\operatorname{ev}_p(y\circ f)=y(\operatorname{ev}_p(f))$, $R\circ L=1\otimes\psi$, which is bijective; so $L$ is injective. As $f-L(R(f))\in\mathfrak m'\mathcal M'$ for $f\in\mathcal M'$, $L$ is surjective by Nakayama's lemma. So $L$ is an isomorphism.
% ordBC (OrdBCMapOrd.lean, 232 lines): ordBC_mapOrd (ordBC_mapSp, ordBC_mapU through uOpNat, ordBC_mapval through ordBC_pow), ordBC_cont and ordBC_free; ordStruct for both rings (the Lean's Frame lemmas hold for every complete local noetherian ring with finite residue field); exists_finite_iwahori_doubleCoset_representatives and exists_section_automorphicSpace_with_prescribed_values for x and s, with w p = e (s (Pi.single p 1)); ordBC_key (he1 and he2 for f, uOperator_pow_eventually_invariant_mod_pow, the J of the w p, mem_valIn_of_invariant_and_values_mem_pow for b1, b2, ordBC_valU, and h4 = hfilt), applied to B and to A itself; ordBC_abs (psi surjective by Submodule.le_of_le_smul_of_le_jacobson_bot and bijective by OrzechProperty, L = liftBaseChange, R with hRL, IsBaseChange.of_equiv).
\end{proof}

\begin{lemma}\label{sha-hecke-ideal-tw-augmented-hida-inputs-separation}
Let $H$ be a group, $\sigma$ an irreducible representation of $H$ on a complex vector space $V$, $\pi$ a representation of $H$ on a complex vector space $M$, and $Y\subset M$ a subspace that is the sum of the images of the $H$-equivariant maps $V\to M$ with image in $Y$. Then $Y$ is stable under $\pi$, and for every nonzero $F\in Y$ there is a linear map $p:M\to V$ with $p(F)\ne0$ and $p(\pi(h)y)=\sigma(h)p(y)$ for all $h\in H$ and $y\in Y$.
\end{lemma}
\leanref{Development/AutomorphicArithmetic/AutomorphicVarieties/HeckeAlgebras/OrdinaryHeckeCharacters/IsoSep.lean}{UpperTheory.V.Hi.isoSep}
\leanref{Development/AutomorphicArithmetic/IntegralCasts/CapelliIdentities/IrreducibleProjectionExists.lean}{UpperTheory.irreducible_projection_exists}
\begin{proof}
The image of an equivariant map is stable under $\pi$, so $Y$ is. Regard $V$, $M$ and $Y$ as modules over the group algebra $\mathbf C[H]$. The image of an equivariant map $V\to Y$ is $0$ or isomorphic to $V$, as its kernel is a submodule of the simple module $V$. So $Y$ is a sum of simple submodules isomorphic to $V$, hence semisimple and isotypic of type $V$, and there is an isomorphism of $\mathbf C[H]$-modules $e:Y\cong\bigoplus_{i\in I}V$ for some set $I$. As $F\ne0$, some coordinate $e(F)_i$ is nonzero. Let $p_0:Y\to V$ be $y\mapsto e(y)_i$, which is equivariant with $p_0(F)\ne0$, and let $p:M\to V$ be a linear map that extends $p_0$.
% isoSep is irreducible_projection_exists: hinv (the stability of Y), the intertwining maps psi i from sigma to Y, htop (their ranges span Y), sepCore (isSemisimpleModule_of_isSemisimpleModule_submodule', IsIsotypicOfType.of_isotypicComponent_eq_top, linearEquiv_finsupp and a coordinate Finsupp.lapply i), then LinearMap.exists_extend.
\end{proof}

\begin{lemma}\label{sha-hecke-ideal-tw-augmented-hida-inputs-embedding}
Let $L$ be a number field, $\ell$ a prime, $w$ a place of $L$ above $\ell$ with ring of integers $\mathcal O_{L_w}$ in the completion $L_w$, and $\tau:L_w\to\bar{\mathbf Q}_\ell$ a continuous ring homomorphism. Then for every ring of integers $\mathcal O_1$ of a finite extension of $\mathbf Q_\ell$ in $\bar{\mathbf Q}_\ell$ there is such a ring $\mathcal O'\supset\mathcal O_1$ with $\tau(\mathcal O_{L_w})\subset\mathcal O'$. For every such $\mathcal O'$ with $\tau(\mathcal O_{L_w})\subset\mathcal O'$, $\tau$ maps the maximal ideal of $\mathcal O_{L_w}$ into the maximal ideal of $\mathcal O'$.
\end{lemma}
\leanref{Development/AutomorphicArithmetic/AutomorphicVarieties/BuildingsResidueChars/ExpCharIntegralityCriterion.lean}{UpperTheory.expChar_integrality_criterion}
\begin{proof}
Let $K\subset\bar{\mathbf Q}_\ell$ be the extension of $\mathbf Q_\ell$ generated by the images of a basis of $L$ over $\mathbf Q$. It is finite over $\mathbf Q_\ell$, hence closed, and $\tau(L)\subset K$; as $L$ is dense in $L_w$ and $\tau$ is continuous, $\tau(L_w)\subset K$. For $x\in\mathcal O_{L_w}$ the powers of $\tau(x)$ are bounded, as $\tau(\mathcal O_{L_w})$ is compact, so $|\tau(x)|\le1$, and $\tau(x)$ is integral over $\mathbf Z_\ell$, as its minimal polynomial over $\mathbf Q_\ell$ then has coefficients in $\mathbf Z_\ell$. Let $\mathcal O'$ be the ring of integers of the extension generated by $K$ and the fraction field of $\mathcal O_1$; it contains $\mathcal O_1$ and $\tau(\mathcal O_{L_w})$. Now let $\mathcal O'\supset\tau(\mathcal O_{L_w})$ be any such ring and $x$ in the maximal ideal of $\mathcal O_{L_w}$. Then $x^k\to0$, so $\tau(x)^k\to0$ and $|\tau(x)|<1$. If $y\tau(x)=1$ with $y\in\mathcal O'$, then $|y|\le1$, as $y$ is integral over $\mathbf Z_\ell$, and $|y\tau(x)|<1$, a contradiction. So $\tau(x)$ lies in the maximal ideal of $\mathcal O'$.
% expChar_integrality_criterion: ec_field (Reformulated.B.Qh13.ec_field: K adjoins the images of Module.finBasis Q L, closed as finite dimensional, denseRange_algebraMap), HV.coefSup (a common Coef for c and the integral closure of Z_l in K, dvr_intClosure), ec_le (compactness), ec_intIff (spectralValue_le_one_iff), ec_lt (Valued.tendsto_zero_pow_of_v_lt_one).
\end{proof}

\begin{lemma}\label{sha-hecke-ideal-tw-augmented-hida-inputs-antidiagonal}
Let $R$ be a commutative ring and $n\ge0$. For $k=1,\dots,K$ let $0\le s_k\le n$ and $I_k(1),\dots,I_k(s_k)\in\{1,\dots,n\}$, and let $\beta=\prod_k\Delta_k$ be the bideterminant in $R[x_{ab}:1\le a,b\le n]$ whose factor $\Delta_k$ is the minor of $x$ on the rows $I_k(1),\dots,I_k(s_k)$ and the first $s_k$ columns. Let $\nu_t$ be the number of $k$ with $s_k\ge t$, and $J$ the matrix with $J_{ab}=1$ if $a+b=n+1$ and $J_{ab}=0$ otherwise. Then:
\begin{enumerate}
\item $\beta({}^tg\,J)=\prod_{i=1}^ng_{ii}^{\nu_{n+1-i}}\,\beta(J)$ for every upper triangular matrix $g$ over $R$;
\item if $I_k(p)+s_k=n+p$ for all $k$ and $p$, then $\beta(J)$ is a unit;
\item if each $I_k$ is strictly increasing and $\beta(J)\ne0$, then $I_k(p)+s_k=n+p$ for all $k$ and $p$.
\end{enumerate}
\end{lemma}
\leanref{Development/AutomorphicArithmetic/AutomorphicVarieties/BuildingsResidueChars/BideterminantAntidiagEvalWeightUnitSupport.lean}{UpperTheory.bideterminant_antidiag_eval_weight_unit_support}

\begin{lemma}\label{sha-hecke-ideal-tw-augmented-hida-inputs-dominance}
Let $n$, the $s_k$, the $I_k$ and $\nu$ be as in Lemma~\ref{sha-hecke-ideal-tw-augmented-hida-inputs-antidiagonal}, with each $I_k$ strictly increasing, and let $c_i$ be the number of pairs $(k,p)$ with $I_k(p)=i$. Then:
\begin{enumerate}
\item $\sum_ia_i\nu_{n+1-i}\le\sum_ia_ic_i$ for every nonincreasing $a\in\mathbf Z^n$;
\item if $\sum_i(n-i)c_i\le\sum_i(n-i)\nu_{n+1-i}$, then $I_k(p)+s_k=n+p$ for all $k$ and $p$;
\item if $I_k(p)+s_k=n+p$ for all $k$ and $p$, then $c_i=\nu_{n+1-i}$ for every $i$, and every family $(s'_k,I'_k)$ with the same property and the same numbers $\nu_t$ gives the same multiset of pairs $(I'_k(p),p)$.
\end{enumerate}
\end{lemma}
\leanref{Development/AutomorphicArithmetic/AutomorphicVarieties/BuildingsResidueChars/BideterminantRowWeightDominance.lean}{UpperTheory.bideterminant_row_weight_dominance}

\begin{lemma}\label{sha-hecke-ideal-tw-augmented-hida-inputs-tableaux}
Let $n\ge1$, $\lambda$ a dominant weight of length $n$ (Definition~\ref{is-polynomial-model}), $\nabla(\lambda)$ and $\rho_\lambda$ as there over $\mathbf C$, and $\mathcal O\subset\mathbf C$ a subring. Then there are a basis $(b_j)_{j\in\kappa}$ of $\nabla(\lambda)$, an index $k_0\in\kappa$ and $\mu_j\in\mathbf Z^n$ for $j\in\kappa$ such that:
\begin{enumerate}
\item for every nonzero $c\in\mathcal O$ and every $g\in GL_n(\mathcal O)$ with $g\equiv1$ modulo $c$, the matrix of $\rho_\lambda(g)$ in the basis $b$ is $1+cX$ for a matrix $X$ over $\mathcal O$;
\item for every upper triangular $g\in GL_n(\mathbf C)$ and every $j$, the coordinate $k_0$ of $\rho_\lambda(g)b_j$ is $\prod_ig_{ii}^{\mu_{k_0,i}}$ if $j=k_0$ and $0$ otherwise;
\item $\rho_\lambda(\operatorname{diag}(t_1,\dots,t_n))b_j=\prod_it_i^{\mu_{j,i}}\,b_j$ for all $j$ and $t_i\in\mathbf C^\times$;
\item $\mu_{k_0,i}=\lambda_{n+1-i}$ for every $i$;
\item $\sum_ia_i\mu_{k_0,i}\le\sum_ia_i\mu_{j,i}$ for every $j$ and every nonincreasing $a\in\mathbf Z^n$;
\item $\sum_i(n-i)\mu_{k_0,i}<\sum_i(n-i)\mu_{j,i}$ for every $j\ne k_0$.
\end{enumerate}
\end{lemma}
\leanref{Development/AutomorphicArithmetic/AutomorphicVarieties/HeckeAlgebras/OrdinaryHeckeCharacters/SqBasis.lean}{UpperTheory.V.Hi.sqBasis}
\begin{proof}
Let $\nu=\lambda^\natural$ and $s=\operatorname{sh}(\lambda)$, so that $\nabla(\lambda)=\nabla(\nu)$ and $(\rho_\lambda(g)f)(x)=\det(g)^sf({}^tg\,x)$. For a domain $A$ let $\nabla_A(\nu)\subset A[x_{ab}]$ be defined by the conditions of Definition~\ref{is-polynomial-model}(2). By the standard basis theorem for bitableaux over $\mathcal O$ \cite{DEP80} (Lemma~\ref{is-standard-basis}(1)), the bideterminants $b_j$ of shape $\nu$ indexed by the semistandard tableaux $j$, with strictly increasing rows $I_k$, form an $\mathcal O$-basis of $\nabla_{\mathcal O}(\nu)$; they have integer coefficients, and they form a basis of $\nabla(\nu)$ over $\mathbf C$. Each $b_j$ is homogeneous of degree $c_{j,i}$ in the entries of the row $i$, with $c_{j,i}$ the number of pairs $(k,p)$ with $I_k(p)=i$. Let $\mu_{j,i}=c_{j,i}+s$.

(3) $b_j(\operatorname{diag}(t)x)=\prod_it_i^{c_{j,i}}b_j(x)$, and $\det(\operatorname{diag}(t))^s=\prod_it_i^s$.

(1) Let $g\in GL_n(\mathcal O)$ with $g\equiv1$ modulo $c$, and $f=\rho_\lambda(g)b_j-b_j$. Reduction modulo $c\mathcal O$ commutes with the action, and $g$ reduces to $1$; so the coefficients of $f$ lie in $c\mathcal O$, and $f=ch$ with $h\in\mathcal O[x_{ab}]$. As $\rho_\lambda(g)b_j$ and $b_j$ lie in $\nabla_{\mathcal O}(\nu)$, so does $ch$, and so does $h$, as the conditions defining $\nabla_{\mathcal O}(\nu)$ are linear and $\mathcal O[x_{ab}]$ is a domain. So $h=\sum_iy_ib_i$ with $y_i\in\mathcal O$, and the coordinate $i$ of $\rho_\lambda(g)b_j$ is $\delta_{ij}+cy_i$.

\emph{The index $k_0$.} Let $\ell(f)=f(J)$, with $J$ as in Lemma~\ref{sha-hecke-ideal-tw-augmented-hida-inputs-antidiagonal}. The bideterminant of shape $\nu$ with $I_k(p)=n-s_k+p$ lies in $\nabla(\nu)$ and has $\ell\ne0$ by Lemma~\ref{sha-hecke-ideal-tw-augmented-hida-inputs-antidiagonal}(2); so $\ell(b_{k_0})\ne0$ for some $k_0$. If $\ell(b_j)\ne0$, then the rows of $j$ satisfy $I_k(p)+s_k=n+p$ by Lemma~\ref{sha-hecke-ideal-tw-augmented-hida-inputs-antidiagonal}(3), and then $j=k_0$ by Lemma~\ref{sha-hecke-ideal-tw-augmented-hida-inputs-dominance}(3), as a semistandard tableau is determined by its pairs $(I_k(p),p)$. So the coordinate $k_0$ of $f\in\nabla(\nu)$ in the basis $b$ is $\ell(f)/\ell(b_{k_0})$.

(4) By Lemma~\ref{sha-hecke-ideal-tw-augmented-hida-inputs-dominance}(3), $c_{k_0,i}=\nu_{n+1-i}$, so $\mu_{k_0,i}=\nu_{n+1-i}+s=\lambda_{n+1-i}$.

(2) Let $g$ be upper triangular, so that $\det g=\prod_ig_{ii}$. By Lemma~\ref{sha-hecke-ideal-tw-augmented-hida-inputs-antidiagonal}(1), $\ell(\rho_\lambda(g)b_j)=\prod_ig_{ii}^{s}\prod_ig_{ii}^{\nu_{n+1-i}}\ell(b_j)=\prod_ig_{ii}^{\mu_{k_0,i}}\ell(b_j)$, using the shape of $b_j$ and (4). Dividing by $\ell(b_{k_0})$ gives (2), as $\ell(b_j)=0$ for $j\ne k_0$.

(5) and (6) By Lemma~\ref{sha-hecke-ideal-tw-augmented-hida-inputs-dominance}(1), $\sum_ia_i\nu_{n+1-i}\le\sum_ia_ic_{j,i}$; adding $s\sum_ia_i$ gives (5). If $\sum_i(n-i)\mu_{j,i}\le\sum_i(n-i)\mu_{k_0,i}$, then $\sum_i(n-i)c_{j,i}\le\sum_i(n-i)\nu_{n+1-i}$, so the rows of $j$ satisfy $I_k(p)+s_k=n+p$ by Lemma~\ref{sha-hecke-ideal-tw-augmented-hida-inputs-dominance}(2), and $j=k_0$ as above.
% sqBasis: sSBasis over O (the book's Lemma is-standard-basis (1), cited to DEP80), ldlat_linInd_base and span_map_weightBasis_eq_dualWeyl for the basis over K, hbc and hbK (the same bideterminants over K); (1) is cong (map_sqAct, dualWeylMap, C_dvd_iff_dvd_coeff and dualWeylSat for the saturation); (3) is rowActDiagWeight with hb.2.2; the linear form ell is evaluation at J (aeval J), hne and hk0, hbot (bideterminant_antidiag_eval_weight_unit_support (3)), huniq and hw0 (bideterminant_row_weight_dominance (3)), hco; (2) through (1) of the antidiagonal lemma and detPow_eq; (5) and (6) through (1) and (2) of the dominance lemma. The weights are mu j i = rowWt + detShift.
\end{proof}

\begin{lemma}\label{sha-hecke-ideal-tw-augmented-hida-inputs-matrix}
In Situation~\ref{sha-hecke-ideal-tw-augmented-setup}, let $T$, $N$ and $\mathcal O'$ be as in Definition~\ref{sha-hecke-ideal-tw-augmented-hida-weight}. Let $\rho$ be a representation of $G_\ell$ on a complex vector space $M$ with a finite basis $(b_j)_{j\in\kappa}$, let $k_0\in\kappa$, and let $A(g)$ be the matrix of $\rho(g)$ in this basis. Suppose:
\begin{enumerate}
\item $A(y)$ has entries in $\mathcal O'$ for $y$ in $N_0$, in $T_0$ or in $\bar N_0$;
\item for dominant $\alpha$, $A(\alpha)_{k_0k_0}^{-1}A(\alpha)$ has entries in $\mathcal O'$;
\item $A(y)_{k_0j}$ is $1$ for $j=k_0$ and $0$ otherwise, for $y\in N$; and $A(t)_{k_0j}=0$ for $t\in T$ and $j\ne k_0$;
\item the entries of $A(z)_{k_0k_0}^{-1}A(z)$ outside the column $k_0$ lie in $\mathfrak m'$;
\item for every $k\ge0$ there is $c\ge0$ such that $A(y)-1$ has entries in $\mathfrak m'^k$ for $y$ in $\bar N_c$ or in $T_0^{(c)}$;
\item $A(t)=A(t)_{k_0k_0}\cdot1$ for $t\in T$ central in $G_\ell$.
\end{enumerate}
Then there is a weight datum $(\xi,\chi,0)$ over $\mathcal O'$ (Definition~\ref{sha-hecke-ideal-tw-augmented-hida-weight}) with $\xi$ isomorphic to $\rho$ and $\chi(t)=A(t)_{k_0k_0}$ for $t\in T$.
\end{lemma}
\leanref{Development/AutomorphicArithmetic/AutomorphicVarieties/HeckeAlgebras/OrdinaryHeckeCharacters/ProdZpowConst.lean}{ofBasis}
\begin{proof}
Number $\kappa$ as $0,\dots,d$ with $k_0$ as $0$, and let $\xi(g)=A(g)$ in this numbering, so that $\xi$ is isomorphic to $\rho$ through the basis. For $s,t\in T$, by (3), $A(st)_{k_0k_0}=\sum_jA(s)_{k_0j}A(t)_{jk_0}=A(s)_{k_0k_0}A(t)_{k_0k_0}$; so $\chi(t)=A(t)_{k_0k_0}$ is a character of $T$. Then (1) and (2) give (1)(a) of Definition~\ref{sha-hecke-ideal-tw-augmented-hida-weight} with $c_0=0$; (3) and (6) give (1)(b); (4) gives (1)(c); and (5), at the entry $(k_0,k_0)$ for $T_0^{(c)}$, gives (1)(d).
% ofBasis: sigma moves k0 to 0 (Equiv.swap), b' = b.reindex, xi through LinearMap.toMatrixAlgEquiv, chi0 multiplicative by hT, then the fields int_N, integral_torusCharacter and int_Nm (hK), int_P (hP), row_N (hN), torusAction (hT), nilpotent_central_action (hz), cont_Nm and continuous_torusCharacter (hc), central (hZ), with c0 = 0.
\end{proof}

\begin{definition}\label{sha-hecke-ideal-tw-augmented-hida-inputs-lattices}
In Situation~\ref{sha-hecke-ideal-tw-augmented-setup}, let $(\xi,\chi,c_0)$ be a weight datum over $\mathcal O'$ (Definition~\ref{sha-hecke-ideal-tw-augmented-hida-weight}), with $\xi$ acting on $\mathbf C^{d+1}$, and regard $\mathcal O'$ as a subring of $\mathbf C$ through $\iota$. For $F\in V^\xi$ with components $F_1,\dots,F_m$, $g\in G^{(\ell)}\times G_\ell$ and $1\le i\le m$, let $F^\flat_i(g)=\xi(g_\ell)^{-1}F_i(g)\in\mathbf C^{d+1}$; so $\operatorname{ctrl}(F)_i(g)$ is the coordinate $0$ of $F^\flat_i(g)$. For an $\mathcal O'$-submodule $\Lambda\subset\mathbf C$, $\mathrm{Lat}(\Lambda)$ is the set of the $F\in V^\xi$ such that every coordinate of every $F^\flat_i(g)$ lies in $\Lambda$.
\end{definition}
\leanref{Development/AutomorphicArithmetic/AutomorphicVarieties/HeckeAlgebras/CtrlDefs.lean}{UpperTheory.V.Hi.Frame.unt}
\leanref{Development/AutomorphicArithmetic/AutomorphicVarieties/HeckeAlgebras/CtrlDefs.lean}{UpperTheory.V.Hi.Frame.Lat}
\leanref{Development/AutomorphicArithmetic/AutomorphicVarieties/HeckeAlgebras/CtrlDefs.lean}{UpperTheory.V.Hi.mIm}

\begin{lemma}\label{sha-hecke-ideal-tw-augmented-hida-inputs-lattice}
In Situation~\ref{sha-hecke-ideal-tw-augmented-setup}, let $(\xi,\chi,c_0)$ be a weight datum over $\mathcal O'$, and $\Lambda,\Lambda'\subset\mathbf C$ $\mathcal O'$-submodules with $\mathfrak m'\Lambda\subset\Lambda'$. Let $F\in\mathrm{Lat}(\Lambda)$ (Definition~\ref{sha-hecke-ideal-tw-augmented-hida-inputs-lattices}) be such that $\operatorname{ctrl}(F)$ has values in $\Lambda'^m$. Then $U'_zF\in\mathrm{Lat}(\Lambda')$. In particular, if $\mathfrak m'\Lambda\subset\Lambda$, then $U'_z$ maps $\mathrm{Lat}(\Lambda)$ into itself.
\end{lemma}
\leanref{Development/AutomorphicArithmetic/AutomorphicVarieties/BuildingsResidueChars/UpOperatorMapsLatticeOfControlValues.lean}{UpperTheory.Up_operator_maps_lattice_of_control_values}
\begin{proof}
For each representative $x$ of $N_0/zN_0z^{-1}$, $\xi(g_\ell)^{-1}=\xi(x)\xi(z)\xi(g_\ell xz)^{-1}$, so
\[
(U'_zF)^\flat_i(g)=\sum_x\xi(x)\,\chi(z)^{-1}\xi(z)\,F^\flat_i(g(1,xz)).
\]
Let $w=F^\flat_i(g(1,xz))$. The coordinate $l$ of $\chi(z)^{-1}\xi(z)w$ is $\sum_j(\chi(z)^{-1}\xi(z))_{lj}w_j$. The term $j=0$ lies in $\mathcal O'\Lambda'\subset\Lambda'$, as $w_0$ is a value of $\operatorname{ctrl}(F)$ and $\chi(z)^{-1}\xi(z)$ has entries in $\mathcal O'$ (Definition~\ref{sha-hecke-ideal-tw-augmented-hida-weight}(1)(a)); the terms $j\ne0$ lie in $\mathfrak m'\Lambda\subset\Lambda'$, by (1)(c). As $\xi(x)$ has entries in $\mathcal O'$, by (1)(a), the coordinates of $(U'_zF)^\flat_i(g)$ lie in $\Lambda'$. For the last statement, take $\Lambda'=\Lambda$: the values of $\operatorname{ctrl}(F)$ are coordinates of the $F^\flat_i(g)$.
% Up_operator_maps_lattice_of_control_values: hF' (F = xi(g) times unt F), hM (integral matrices preserve Lambda'), main by finsum_induction over the uOp sum with c = chi(z)^-1, using W.int_N, W.int_P at z and W.nilpotent_central_action. The last statement is ctrlOrd_stab of CtrlOrd.lean.
\end{proof}

\begin{lemma}\label{sha-hecke-ideal-tw-augmented-hida-inputs-levels}
In Situation~\ref{sha-hecke-ideal-tw-augmented-setup}, let $(\xi,\chi,c_0)$ be a weight datum over $\mathcal O'$, $U\subset G^{(\ell)}_0$ a subgroup, $c\ge0$ and $\phi$ a character of $T_0$, trivial on $T_0^{(c)}$. Then:
\begin{enumerate}
\item $\mathrm{Cl}(U,c,\phi)\subset\mathrm{Cl}(U,c',\phi)$ for $c'\ge c$;
\item $U'_z$ maps $\mathrm{Cl}(U,c+1,\phi)$ into $\mathrm{Cl}(U,c,\phi)$, and $\mathrm{Cl}(U,c,\phi)$ into itself;
\item $U_z'^n$ maps $\mathrm{Cl}(U,c+e,\phi)$ into $\mathrm{Cl}(U,c,\phi)$ for $0\le e\le n$.
\end{enumerate}
\end{lemma}
\leanref{Development/AutomorphicArithmetic/AutomorphicVarieties/BuildingsResidueChars/ClassicalFormsLevelMonotoneAndUpStable.lean}{UpperTheory.classicalForms_level_monotone_and_Up_stable}
\begin{proof}
(1) $\bar N_{c'}\subset\bar N_c$.

(2) Let $F\in\mathrm{Cl}(U,c+1,\phi)$. The operator $U'_z$ is a combination of right translations by elements of $G_\ell$, which preserve the condition $F(\gamma g)=\xi(\gamma_\ell)F(g)$ and commute with $F\mapsto u\cdot F$; so $U'_zF$ satisfies the first two conditions defining $\mathrm{Cl}(U,c,\phi)$. As $F$ is fixed by $N_0$ and by $\bar N_{c+1}$, and by $T_0^{(c)}$, on which $\phi$ is trivial, $U_zF$ is fixed by $I_1(c)$, hence by $N_0$ and $\bar N_c$, by Lemma~\ref{sha-hecke-ideal-tw-augmented-uop}(3); and $T_0$ acts on it by $\phi^{-1}$, by Lemma~\ref{sha-hecke-ideal-tw-augmented-uop}(1). So $U'_zF\in\mathrm{Cl}(U,c,\phi)$. With (1), $U'_z$ maps $\mathrm{Cl}(U,c,\phi)\subset\mathrm{Cl}(U,c+1,\phi)$ into $\mathrm{Cl}(U,c,\phi)$.

(3) As $\phi$ is trivial on $T_0^{(c+e')}\subset T_0^{(c)}$, (2) for $c+e-1,\dots,c$ shows that $U_z'^e$ maps $\mathrm{Cl}(U,c+e,\phi)$ into $\mathrm{Cl}(U,c,\phi)$, which $U_z'^{n-e}$ maps into itself.
% classicalForms_level_monotone_and_Up_stable: mono, step from the private lower_step (uOpSum, ctrlEq (3) for the commutation with raF, uOpLevel (2.2) for the eigen-conditions), down by induction on e, stay by induction on k.
\end{proof}

\begin{lemma}\label{sha-hecke-ideal-tw-augmented-hida-inputs-values}
In Situation~\ref{sha-hecke-ideal-tw-augmented-setup}, let $(\xi,\chi,c_0)$ be a weight datum over $\mathcal O'$, $c\ge0$, $\theta'$ and $\phi$ characters of $T_0$ with values in $\mathcal O'^\times$, with $\phi$ trivial on $T_0^{(c)}$ and $\theta'(t)=\phi(t)\chi(t)$ for $t\in T_0$, and $U\subset G^{(\ell)}_0$ a subgroup of finite level that satisfies the freeness condition of Lemma~\ref{sha-hecke-ideal-tw-augmented-sections} for $\theta'$. Let $x_1,\dots,x_N$ be as in Lemma~\ref{sha-hecke-ideal-tw-augmented-sections}(1) for $c$. Then:
\begin{enumerate}
\item for all $v_{n,i}\in\mathbf C^{d+1}$, $1\le n\le N$, $1\le i\le m$, there is $F\in\mathrm{Cl}(U,c,\phi)$ with $F^\flat_i(x_n)=v_{n,i}$ for all $n$ and $i$ (Definition~\ref{sha-hecke-ideal-tw-augmented-hida-inputs-lattices});
\item if $c\ge c_0$, $F\in\mathrm{Cl}(U,c,\phi)$, and $\Lambda\subset\mathbf C$ is an $\mathcal O'$-submodule containing every coordinate of every $F^\flat_i(x_n)$, then $F\in\mathrm{Lat}(\Lambda)$.
\end{enumerate}
\end{lemma}
\leanref{Development/AutomorphicArithmetic/AutomorphicVarieties/BuildingsResidueChars/ExistsAutomorphicFormWithPrescribedIwahoriValues.lean}{UpperTheory.exists_automorphicForm_with_prescribed_iwahori_values}
\leanref{Development/AutomorphicArithmetic/AutomorphicVarieties/BuildingsResidueChars/ExistsIwahoriCharacterExtendingTorusCharacter.lean}{UpperTheory.exists_iwahori_character_extending_torus_character}
\leanref{Development/AutomorphicArithmetic/AutomorphicVarieties/BuildingsResidueChars/ExistsAutomorphicFunctionWithValuesOnCosetReps.lean}{UpperTheory.exists_automorphic_function_with_values_on_coset_reps}
\leanref{Development/AutomorphicArithmetic/AutomorphicVarieties/BuildingsResidueChars/MemLatticeOfValuesAtDoubleCosetReps.lean}{UpperTheory.mem_lattice_of_values_at_double_coset_reps}
\begin{proof}
Let $I(c)=N_0T_0\bar N_c$ and $t(b)\in T_0$ for $b\in I_\ell$ be as in the proof of Lemma~\ref{sha-hecke-ideal-tw-augmented-sections}(2).

(1) \emph{A character of $I(c)$.} Let $\psi(b)=\phi(t(b))$ for $b\in I(c)$. For $b'=xty$ with $x\in N_0$, $t\in T_0$ and $y\in\bar N_c$, the properties of $t$ listed there give $t(bb')\in t(b)\,t\,T_0^{(c)}$ and $t(b')\in t\,T_0^{(c)}$; as $\phi$ is trivial on $T_0^{(c)}$, $\psi(bb')=\psi(b)\psi(b')$. So $\psi$ is a character of $I(c)$, trivial on $N_0$ and $\bar N_c$, and equal to $\phi$ on $T_0$.

\emph{A function on double cosets.} For $u\in U$ and $b\in I(c)$, let $\rho(u,b)$ act on the $m\times(d+1)$ matrices by $\rho(u,b)w=\psi(b)\,\tau_0(u)\,w\,{}^t\xi(b)$. If $\gamma x_n=x_n(u,b)$ with $\gamma\in U(n)(L^+)$, then $\rho(u,b)=1$: by the freeness condition, $b$ lies in $T_0$ and in the center of $G_\ell$, and $\theta'(b)\tau_0(u)=1$; so $\psi(b)=\phi(b)$, $\xi(b)=\chi(b)\cdot1$ by Definition~\ref{sha-hecke-ideal-tw-augmented-hida-weight}(1)(b), and $\rho(u,b)w=\phi(b)\chi(b)\tau_0(u)w=w$. Let $V_n$ be the matrix with rows $v_{n,1},\dots,v_{n,m}$. By the uniqueness in Lemma~\ref{sha-hecke-ideal-tw-augmented-sections}(1), $\Phi(\gamma x_n(u,b))=\rho(u,b)^{-1}V_n$ is a well defined function on $G^{(\ell)}\times G_\ell$: if $\gamma x_n(u,b)=\gamma'x_n(u',b')$, then $\gamma'^{-1}\gamma x_n=x_n(u'u^{-1},b'b^{-1})$, so $\rho(u'u^{-1},b'b^{-1})=1$ and $\rho(u,b)^{-1}=\rho(u',b')^{-1}$. It satisfies $\Phi(\gamma g)=\Phi(g)$, $\Phi(g(u,b))=\rho(u,b)^{-1}\Phi(g)$ and $\Phi(x_n)=V_n$.

\emph{The form.} Let $F_i(g)=\xi(g_\ell)\Phi(g)_i$, with $\Phi(g)_i$ the row $i$ of $\Phi(g)$ as a column vector. Then $F_i(\gamma g)=\xi(\gamma_\ell)F_i(g)$. As $\Phi(g(u,1))=\tau_0(u)^{-1}\Phi(g)$, $u\cdot F=F$ for $u\in U$. As $\Phi(g(1,b))=\psi(b)^{-1}\Phi(g)\,{}^t\xi(b)^{-1}$, $F_i(g(1,b))=\psi(b)^{-1}F_i(g)$ for $b\in I(c)$; so $F$ is fixed by $N_0$ and $\bar N_c$, and $T_0$ acts on it by $\phi^{-1}$. Hence $F\in\mathrm{Cl}(U,c,\phi)$, and $F^\flat_i(x_n)=\Phi(x_n)_i=v_{n,i}$.

(2) Let $g\in G^{(\ell)}\times G_\ell$. By Lemma~\ref{sha-hecke-ideal-tw-augmented-sections}(1), $g=\gamma x_n(u,b)$ with $b\in I(c)$; write $b=yty'$ with $y\in N_0$, $t\in T_0$ and $y'\in\bar N_c$, so that $g=\gamma\,x_n(u,1)(1,y)(1,t)(1,y')$. For $\gamma\in U(n)(L^+)$, $F^\flat_i(\gamma h)=F^\flat_i(h)$. For $u\in U$, $F^\flat_i(h(u,1))=\sum_j\tau_0(u^{-1})_{ij}F^\flat_j(h)$, as $u^{-1}\cdot F=F$, and $\tau_0$ has entries in $\mathcal O$. For $b'$ in $N_0$, in $T_0$ or in $\bar N_c$, $F(h(1,b'))=sF(h)$ with $s$ equal to $1$ or to $\phi(b')^{-1}\in\mathcal O'$, so $F^\flat_i(h(1,b'))=s\,\xi(b'^{-1})F^\flat_i(h)$, and $\xi(b'^{-1})$ has entries in $\mathcal O'$ by Definition~\ref{sha-hecke-ideal-tw-augmented-hida-weight}(1)(a), as $\bar N_c\subset\bar N_{c_0}$. Each of these steps keeps the coordinates in $\Lambda$, starting from those of $F^\flat_i(x_n)$.
% exists_automorphicForm_with_prescribed_iwahori_values: exists_iwahori_character_extending_torus_character (psi b = phi (tt b), exists_iwahori_torus_projection, key and key1), rhoT, exists_automorphic_function_with_values_on_coset_reps (key: equal representations give equal values, through the stabilizer hstab supplied by Frame.Free, W.central and hθ), h2', hp, then the fields of Cl and the values at x n. mem_lattice_of_values_at_double_coset_reps: stepG, stepU, stepN (int_N), stepT (integral_torusCharacter), stepM (int_Nm with Nm_anti for c0 at most c), iwMem for b = y1 t y2.
\end{proof}

\begin{lemma}\label{sha-hecke-ideal-tw-augmented-hida-inputs-lift}
In Situation~\ref{sha-hecke-ideal-tw-augmented-setup}, let $(\xi,\chi,c_0)$ be a weight datum over $\mathcal O'$, $c\ge0$, $\theta'$ and $\phi$ characters of $T_0$ with values in $\mathcal O'^\times$ with $\theta'(t)=\phi(t)\chi(t)$ for $t\in T_0$, $U\subset G^{(\ell)}_0$ a subgroup, and $F\in\mathrm{Cl}(U,c,\phi)\cap\mathrm{Lat}(\mathcal O')$. Then $f=\operatorname{ctrl}(F)$ has values in $\mathcal O'^m$, and it lies in the $\mathcal O'$-module $\mathcal F_{\mathcal O'}(U,\theta')$ of the functions $G^{(\ell)}\times G_\ell\to\mathcal O'^m$ that satisfy the conditions defining $\mathcal F(U)$ in Situation~\ref{sha-hecke-ideal-tw-augmented-setup}(2) with $\mathcal O'$ and $\theta'$ in place of $\Lambda$ and $\theta$. If $k\ge0$ and $c'\ge c$ are such that $\xi(y)-1$ has entries in $\mathfrak m'^k$ for $y\in\bar N_{c'}$, then $f(g(1,y))-f(g)$ has entries in $\mathfrak m'^k$ for all $g$ and $y\in\bar N_{c'}$.
\end{lemma}
\leanref{Development/AutomorphicArithmetic/AutomorphicVarieties/BuildingsResidueChars/ExistsIntegralFormLiftingControlMap.lean}{UpperTheory.exists_integral_form_lifting_controlMap}
\begin{proof}
As $F\in\mathrm{Lat}(\mathcal O')$, $f$ has values in $\mathcal O'^m$. Let $b\in G_\ell$ and $a\in\mathbf C$ be such that $F(g(1,b))=aF(g)$ for all $g$. As $\xi(g_\ell b)^{-1}=\xi(b^{-1})\xi(g_\ell)^{-1}$,
\[
f_i(g(1,b))=a\sum_j\xi(b^{-1})_{0j}\,F^\flat_i(g)_j.
\]
For $b\in N_0$, $a=1$ and the row $0$ of $\xi(b^{-1})$ is $(1,0,\dots,0)$ by Definition~\ref{sha-hecke-ideal-tw-augmented-hida-weight}(1)(b); so $f(g(1,b))=f(g)$. For $t\in T_0$, $a=\phi(t)^{-1}$ and the row $0$ of $\xi(t^{-1})$ is $(\chi(t)^{-1},0,\dots,0)$; so $f(g(1,t))=\phi(t)^{-1}\chi(t)^{-1}f(g)=\theta'(t)^{-1}f(g)$. For $y\in\bar N_{c'}$ with $k$ and $c'$ as in the statement, $a=1$, as $\bar N_{c'}\subset\bar N_c$, and $\xi(y^{-1})_{0j}=\delta_{0j}+a_j$ with $a_j\in\mathfrak m'^k$; so $f_i(g(1,y))-f_i(g)=\sum_ja_jF^\flat_i(g)_j\in\mathfrak m'^k$. By Definition~\ref{sha-hecke-ideal-tw-augmented-hida-weight}(1)(d), for every $k$ there is $c_1$ such that $\xi(y)-1$ has entries in $\mathfrak m'^k$ for $y\in\bar N_{c_1}$; with $c'=\max(c,c_1)$ this gives the continuity condition. As $F(\gamma g)=\xi(\gamma_\ell)F(g)$, $f(\gamma g)=f(g)$ for $\gamma\in U(n)(L^+)$. Finally $\operatorname{ctrl}(u\cdot F)=u\cdot\operatorname{ctrl}(F)$, as $\xi(g_\ell)$ does not change under $g\mapsto g(u,1)$; so $f$ is fixed by $U$.
% exists_integral_form_lifting_controlMap: v the O'-coordinates of unt F, f g i = v i g 0, control_map_translate_trace for hs, key for the congruences on Nm c' (with hxi at y^-1), then the fields of Amb (hA, row_N, torusAction with hθ, cont_Nm for c1) and the U-invariance through ctrlEq (1).
\end{proof}

\begin{definition}\label{sha-hecke-ideal-tw-augmented-hida-inputs-fitting-projector}
Let $R$ be a commutative ring, $I\subset R$ an ideal, $M$ an $R$-module and $U$ an $R$-linear endomorphism of $M$. An \emph{ordinary projector} of $U$ for $I$ is an idempotent $R$-linear endomorphism $e$ of $M$ that commutes with $U$, such that $U$ maps $e(M)$ bijectively onto itself and for every $k\ge0$ there is $n\ge0$ with $U^n(x-e(x))\in I^kM$ for all $x\in M$.
\end{definition}
\leanref{Development/AutomorphicArithmetic/AutomorphicVarieties/AlgebraTools/HidaCompleteRings/IsOrdProjZeroOfIsNilpotent.lean}{UpperTheory.V.AlgebraTools.AnnihilatorQuotientFinite.IsOrdProj}

\begin{lemma}\label{sha-hecke-ideal-tw-augmented-hida-inputs-fitting-naturality}
Let $R$ be a commutative ring, $I\subset R$ an ideal, $M$ and $M'$ $R$-modules, $U$ and $U'$ $R$-linear endomorphisms of $M$ and of $M'$, and $e$ and $e'$ ordinary projectors of $U$ and of $U'$ for $I$ (Definition~\ref{sha-hecke-ideal-tw-augmented-hida-inputs-fitting-projector}).
\begin{enumerate}
\item If $z\in M$ satisfies $e(z)=z$ and $U^nz\in JM$ for an ideal $J\subset R$ and some $n\ge0$, then $z\in JM$.
\item If $\bigcap_kI^kM'=0$ and $f:M\to M'$ is an $R$-linear map with $f\circ U=U'\circ f$, then $f\circ e=e'\circ f$.
\item If $\bigcap_kI^kM=0$, then $e(M)=\bigcap_{n\ge0}U^n(M)$.
\end{enumerate}
\end{lemma}
\leanref{Development/AutomorphicArithmetic/AutomorphicVarieties/AlgebraTools/HidaCompleteRings/OrdBasic.lean}{UpperTheory.V.AlgebraTools.AnnihilatorQuotientFinite.ordCancel}
\leanref{Development/AutomorphicArithmetic/AutomorphicVarieties/AlgebraTools/HidaCompleteRings/OrdBasic.lean}{UpperTheory.V.AlgebraTools.AnnihilatorQuotientFinite.ordNatLin}
\leanref{Development/AutomorphicArithmetic/AutomorphicVarieties/AlgebraTools/HidaCompleteRings/OrdBasic.lean}{UpperTheory.V.AlgebraTools.AnnihilatorQuotientFinite.ordNat}
\leanref{Development/AutomorphicArithmetic/AutomorphicVarieties/AlgebraTools/HidaCompleteRings/OrdBasic.lean}{UpperTheory.V.AlgebraTools.AnnihilatorQuotientFinite.ordRange}
\begin{proof}
(1) Let $V$ be the endomorphism of $M$ that is $e$ followed by the inverse of the bijection $U:e(M)\to e(M)$. If $e(y)=y$, then $e(Uy)=Ue(y)=Uy$, so $V(Uy)=y$. For $n=0$, (1) is the hypothesis. If $U^{n+1}z\in JM$, then $Uz$ satisfies $e(Uz)=Uz$ and $U^n(Uz)\in JM$, so $Uz\in JM$ by induction on $n$, and $z=V(Uz)\in JM$, as $V$ is $R$-linear.

(2) Let $x\in M$. First, $e'(f(x-e(x)))=0$. Indeed, let $k\ge0$, and take $n$ with $U^n(x-e(x))\in I^kM$. As $e'$ commutes with $U'$ and $f\circ U^n=U'^n\circ f$, $U'^n\big(e'(f(x-e(x)))\big)=e'\big(f(U^n(x-e(x)))\big)$, which lies in $I^kM'$, as $f$ and $e'$ are $R$-linear. By (1) for $e'$, $e'(f(x-e(x)))\in I^kM'$, for every $k$, so it is $0$. Second, $f(e(x))=e'(f(e(x)))$. Indeed, let $k\ge0$, and take $n$ with $U'^n(y'-e'(y'))\in I^kM'$ for all $y'\in M'$. As $U^n$ maps $e(M)$ onto itself, $e(x)=U^ny$ for some $y\in e(M)$, and $f(e(x))-e'(f(e(x)))=U'^n\big(f(y)-e'(f(y))\big)\in I^kM'$, for every $k$. Hence $f(e(x))=e'(f(e(x)))=e'(f(x))$, by the first step.

(3) As $U$ maps $e(M)$ onto itself, so does every $U^n$, and $e(M)\subset\bigcap_nU^n(M)$. Conversely, let $x\in\bigcap_nU^n(M)$ and $k\ge0$. Take $n$ with $U^n(y-e(y))\in I^kM$ for all $y\in M$, and $y$ with $x=U^ny$. Then $x-e(x)=U^n(y-e(y))\in I^kM$, as $e$ commutes with $U$. So $x-e(x)\in\bigcap_kI^kM=0$, and $x=e(x)\in e(M)$.
% ordCancel (V = subtype o E.symm o e.rangeRestrict with E the bijection U on range e, hV, induction on n through pow_succ); ordNatLin is ordNat with l = k (h1 through ordCancel for e' and hc', h2 through (h.bij.iterate n).surjOn and pow_intertwine); (3) is the first conjunct h1 of ordRange, which also proves ker e = nilPart I U and the complement, not used here.
\end{proof}

\begin{lemma}\label{sha-hecke-ideal-tw-augmented-hida-inputs-fitting-length}
Let $R$ be a commutative ring, $M$ an $R$-module of finite length and $U$ an $R$-linear endomorphism of $M$. Then there are an $R$-linear endomorphism $e$ of $M$ and $n\ge0$ such that $e$ is an ordinary projector of $U$ for every ideal of $R$ (Definition~\ref{sha-hecke-ideal-tw-augmented-hida-inputs-fitting-projector}) and $U^n(x-e(x))=0$ for all $x\in M$.
\end{lemma}
\leanref{Development/AutomorphicArithmetic/AutomorphicVarieties/AlgebraTools/HidaCompleteRings/OrdBasic.lean}{UpperTheory.V.AlgebraTools.AnnihilatorQuotientFinite.ordFinLen}
\begin{proof}
As $M$ has finite length, it is noetherian and artinian. By Fitting's lemma \cite[pp.~113--114]{Jac09}, there is $n\ge0$ with $M=\ker U^k\oplus U^k(M)$ for every $k\ge n$. As $M$ is artinian, there is $m\ge0$ with $U^k(M)=\bigcap_{j\ge0}U^j(M)$ for every $k\ge m$. Let $N=n+m$. Then $M=\ker U^N\oplus U^N(M)$, $M=\ker U^{N+1}\oplus U^{N+1}(M)$ and $U^{N+1}(M)=U^N(M)$.

Let $e$ be the projection of $M=\ker U^N\oplus U^N(M)$ onto $U^N(M)$. It is idempotent with image $U^N(M)$, and $x-e(x)\in\ker U^N$, so $U^N(x-e(x))=0$ for all $x$; this gives the last property of an ordinary projector, for every ideal, with $n=N$. As $U$ commutes with $U^N$, it maps $U^N(M)$ and $\ker U^N$ into themselves. So for $x\in M$, $U(x-e(x))\in\ker U^N$ and $U(e(x))\in U^N(M)$, whence $e(U(x-e(x)))=0$ and $e(U(e(x)))=U(e(x))$; adding, $e(Ux)=U(e(x))$. The operator $U$ maps $U^N(M)$ injectively into itself: if $x,y\in U^N(M)$ and $Ux=Uy$, then $x-y$ lies in $\ker U^{N+1}$ and in $U^N(M)=U^{N+1}(M)$, so $x=y$, by the decomposition at $N+1$. And it maps $U^N(M)$ onto itself: $x\in U^N(M)=U^{N+1}(M)$ is $U^{N+1}z=U(U^Nz)$.
% ordFinLen: isFiniteLength_iff_isNoetherian_isArtinian; n from LinearMap.eventually_isCompl_ker_pow_range_pow (Fitting's lemma, in Mathlib, cited here to Jacobson) and m from LinearMap.eventually_iInf_range_pow_eq, N = n + m, hc (the complement at N) and hr (the ranges at N and N + 1 agree); e = Submodule.projection with he1 to he4 and hre; the commutation through he2 at U (x - e x) (hU2) and he1 at U (e x) (hU1); the injectivity through Submodule.disjoint_def for the complement at N + 1, and the surjectivity through hr.
\end{proof}

\begin{lemma}\label{sha-hecke-ideal-tw-augmented-hida-inputs-fitting-factorial}
Let $R$ be a commutative ring, $M$ an $R$-module with finitely many elements and $U$ an $R$-linear endomorphism of $M$. Then there is $N\ge0$ such that for every $j\ge N$, $U^{j!}=U^{N!}$ and $U^{j!}$ is an ordinary projector of $U$ for every ideal of $R$ (Definition~\ref{sha-hecke-ideal-tw-augmented-hida-inputs-fitting-projector}).
\end{lemma}
\leanref{Development/AutomorphicArithmetic/AutomorphicVarieties/AlgebraTools/HidaCompleteRings/OrdBasic.lean}{UpperTheory.V.AlgebraTools.AnnihilatorQuotientFinite.factIdem}
\leanref{Development/AutomorphicArithmetic/AutomorphicVarieties/AlgebraTools/HidaCompleteRings/OrdBasic.lean}{UpperTheory.V.AlgebraTools.AnnihilatorQuotientFinite.ordFactFin}
\begin{proof}
The endomorphisms of $M$ form a finite monoid, so there are $s\ge0$ and $p>0$ with $U^{s+p}=U^s$. Then $U^{m+pq}=U^m$ for all $m\ge s$ and $q\ge0$, by induction on $q$. Let $N=s+p$. Then $s\le N!$, and $N!=pc$ for some $c$, as $p\le N$. So $U^{N!}U^{N!}=U^{N!+pc}=U^{N!}$, and $E=U^{N!}$ is idempotent. For $j\ge N$, $j!=N!(d+1)$ for some $d\ge0$, so $j!=N!+p\,cd$ and $U^{j!}=U^{N!}=E$.

The operator $E$ commutes with $U$, so $U$ maps $E(M)$ into itself. It is injective on $E(M)$: write $N!=m+1$; if $U(Ex)=U(Ey)$, then applying $U^m$ gives $E(Ex)=E(Ey)$, that is $Ex=Ey$. As $E(M)$ is finite, $U$ maps it bijectively onto itself. Finally $U^{N!}(x-Ex)=Ex-E^2x=0$ for all $x$. So $E$ is an ordinary projector of $U$ for every ideal.
% factIdem in the finite monoid Module.End R M (s and p from Finite.exists_ne_map_eq_of_infinite, h1, N = s + p, Nat.dvd_factorial and Nat.factorial_dvd_factorial); ordFactFin (hi, hc, hmap, hinj with N! = m + 1, then Finite.injOn_iff_bijOn_of_mapsTo, and the nil condition with n = N!).
\end{proof}

\begin{lemma}\label{sha-hecke-ideal-tw-augmented-hida-inputs-fitting}
Let $\mathcal O$ be a complete local noetherian domain with maximal ideal $\mathfrak m$ and finite residue field, contained in an algebraically closed field $\Omega$, let $\Sigma$ be a finite set, $T_0$ an $\mathcal O$-linear endomorphism of $\mathcal O^\Sigma$, $T$ the $\Omega$-linear endomorphism of $\Omega^\Sigma$ that extends it, and $P=\bigcap_{n\ge0}T_0^n(\mathcal O^\Sigma)$. Then:
\begin{enumerate}
\item there is an idempotent $\mathcal O$-linear endomorphism $e$ of $\mathcal O^\Sigma$ with image $P$ that commutes with $T_0$, such that $T_0$ maps $P$ bijectively onto itself and for every $k\ge0$ there is $n\ge0$ with $T_0^n(v-e(v))\in\mathfrak m^k\mathcal O^\Sigma$ for all $v\in\mathcal O^\Sigma$;
\item for every $k\ge0$ there is $N$ such that $T_0^{j!}v-e(v)\in\mathfrak m^k\mathcal O^\Sigma$ for all $j\ge N$ and $v\in\mathcal O^\Sigma$;
\item $P$ is free over $\mathcal O$, and $v\in\mathcal O^\Sigma$ lies in $P$ if and only if for every $k\ge0$ there is $n\ge0$ with $T_0^{n+1}v-v\in\mathfrak m^k\mathcal O^\Sigma$;
\item the $\Omega$-span of $P$ is the sum of the generalized $\mu$-eigenspaces of $T$ over the nonzero $\mu\in\Omega$ such that $\mu$ and $\mu^{-1}$ are integral over $\mathcal O$.
\end{enumerate}
\end{lemma}
\leanref{Development/AutomorphicArithmetic/AutomorphicVarieties/AlgebraTools/HidaCompleteRings/OrdBasic.lean}{UpperTheory.V.AlgebraTools.AnnihilatorQuotientFinite.ordExists}
\leanref{Development/AutomorphicArithmetic/AutomorphicVarieties/AlgebraTools/HidaCompleteRings/OrdBasic.lean}{UpperTheory.V.AlgebraTools.AnnihilatorQuotientFinite.ordFactLim}
\leanref{Development/AutomorphicArithmetic/AutomorphicVarieties/AlgebraTools/HidaCompleteRings/OrdBasic.lean}{UpperTheory.V.AlgebraTools.AnnihilatorQuotientFinite.ordProjective}
\leanref{Development/AutomorphicArithmetic/AutomorphicVarieties/AlgebraTools/HidaCompleteRings/OrdBasic.lean}{UpperTheory.V.AlgebraTools.AnnihilatorQuotientFinite.ordPer}
\leanref{Development/AutomorphicArithmetic/AutomorphicVarieties/AlgebraTools/HidaCompleteRings/EigOrdInt.lean}{UpperTheory.V.AlgebraTools.AnnihilatorQuotientFinite.eigOrdInt}
\leanref{Development/AutomorphicArithmetic/AutomorphicVarieties/BuildingsResidueChars/SpanOrdinaryPartEqSupUnitGeneralizedEigenspaces.lean}{UpperTheory.span_ordinaryPart_eq_sup_unit_generalized_eigenspaces}
\begin{proof}
Let $M=\mathcal O^\Sigma$, and for $k\ge0$ let $M_k=M/\mathfrak m^kM$, $r_k:M\to M_k$ the reduction and $U_k$ the endomorphism of $M_k$ induced by $T_0$, so that $r_kT_0=U_kr_k$. As $\mathcal O$ is complete and separated for the $\mathfrak m$-adic topology, so is $M$. Each $M_k=(\mathcal O/\mathfrak m^k)^\Sigma$ is finite, as $\mathcal O/\mathfrak m^k$ is: $\mathcal O$ is noetherian with finite residue field, so each $\mathfrak m^i/\mathfrak m^{i+1}$ is finite. In particular $M_k$ has finite length.

(1) \emph{Projectors modulo $\mathfrak m^k$.} For each $k$ let $e_k$ and $n_k$ be as in Lemma~\ref{sha-hecke-ideal-tw-augmented-hida-inputs-fitting-length} for $U_k$. For $k\le k'$ the reduction $r:M_{k'}\to M_k$ satisfies $rU_{k'}=U_kr$, and $e_{k'}$ and $e_k$ are ordinary projectors of $U_{k'}$ and $U_k$ for the zero ideal (Definition~\ref{sha-hecke-ideal-tw-augmented-hida-inputs-fitting-projector}), for which $\bigcap_j0^jM_k=0$. By Lemma~\ref{sha-hecke-ideal-tw-augmented-hida-inputs-fitting-naturality}(2), $re_{k'}=e_kr$. So for $v\in M$ the elements $e_k(r_kv)$ form a compatible family, and, as $M$ is complete and separated, there is a unique $e(v)\in M$ with $r_k(e(v))=e_k(r_kv)$ for all $k$. The map $e$ is $\mathcal O$-linear.

\emph{The map $e$ is an ordinary projector of $T_0$ for $\mathfrak m$.} Two elements of $M$ with the same image under every $r_k$ are equal. As each $e_k$ is idempotent and commutes with $U_k$, $e$ is idempotent and commutes with $T_0$; so $T_0$ maps $e(M)$ into itself. If $T_0e(x)=T_0e(y)$, then $U_ke_k(r_kx)=U_ke_k(r_ky)$ for every $k$, so $e_k(r_kx)=e_k(r_ky)$, as $U_k$ is injective on $e_k(M_k)$; hence $e(x)=e(y)$. Let $x\in M$. For each $k$, as $U_k$ maps $e_k(M_k)$ onto itself, there is $y_k\in M$ with $U_ke_k(r_ky_k)=e_k(r_kx)$, that is $r_k(T_0e(y_k))=r_k(e(x))$. For $a\le b$, $r_a(T_0e(y_a))=r_a(T_0e(y_b))$, so $r_a(e(y_a))=r_a(e(y_b))$, by the injectivity of $U_a$ on $e_a(M_a)$. As $M$ is complete, there is $y\in M$ with $r_k(y)=r_k(e(y_k))$ for all $k$. Then $r_k(e(y))=e_k(r_ke(y_k))=r_k(e(y_k))=r_k(y)$ for all $k$, so $e(y)=y$; and $r_k(T_0y)=r_k(T_0e(y_k))=r_k(e(x))$ for all $k$, so $T_0y=e(x)$. Hence $T_0$ maps $e(M)$ onto itself. Finally, for $k\ge0$ and $x\in M$, $r_k\big(T_0^{n_k}(x-e(x))\big)=U_k^{n_k}\big(r_kx-e_k(r_kx)\big)=0$, so $T_0^{n_k}(x-e(x))\in\mathfrak m^kM$.

\emph{Conclusion.} By Lemma~\ref{sha-hecke-ideal-tw-augmented-hida-inputs-fitting-naturality}(3), $e(M)=\bigcap_nT_0^n(M)=P$. So $e$ is an idempotent with image $P$ that commutes with $T_0$, $T_0$ maps $P$ bijectively onto itself, and the last condition of (1) is the last property of an ordinary projector.

(2) Fix $k$. By Lemma~\ref{sha-hecke-ideal-tw-augmented-hida-inputs-fitting-factorial}, applied to $U_k$ on the finite module $M_k$, there is $N$ such that $U_k^{j!}$ is an ordinary projector of $U_k$ for $\mathfrak m$ for every $j\ge N$. As $\mathfrak m^kM_k=0$, $\bigcap_i\mathfrak m^iM_k=0$, and Lemma~\ref{sha-hecke-ideal-tw-augmented-hida-inputs-fitting-naturality}(2), applied to $r_k$, $e$ and $U_k^{j!}$, gives $r_k(e(v))=U_k^{j!}(r_kv)=r_k(T_0^{j!}v)$ for all $v$. So $T_0^{j!}v-e(v)\in\mathfrak m^k\mathcal O^\Sigma$ for $j\ge N$ and all $v$.
% (1) is ordExists: OrdAdm (O^Sigma) m from ordAdm_maximalIdeal_and_finite_quotients (AdmLocal.lean: completeness through tateH1FinitenessAtPlacesOver and completeFin, and the finite quotients through Submodule.finite_quotient_smul); ordLimProj (Uk = mapQ of U, ek and nk from ordFinLen, hc through ordNatLin for the maps factorPow at the zero ideal, e = IsAdicComplete.lift); ordLimBij (hz, hi, hc, the injectivity through (hp a).bij.injOn, the surjectivity through IsPrecomplete.prec applied to the e (y k), the last condition through pow_intertwine); then isOrdProj_ordProj, as the Lean names the chosen ordProj m U, which is an ordinary projector because e is, and ordRange for its image. (2) is ordFactLim (the direction h1): ordFactFin for Uk on the finite quotient, its IsHausdorff instance, and ordNatLin with mkQ; the converse, that a factorial limit is ordProj, is not stated. Lemma sha-hecke-ideal-tw-augmented-projector proves the analogous facts for the periodic part, by a different argument (ordAbs).

(3) As $e$ is an idempotent with image $P$, $P$ is a direct summand of $\mathcal O^\Sigma$, so it is finite and projective, hence free over the local ring $\mathcal O$. Let $v\in P$, so $e(v)=v$, and $k\ge0$; with $N$ as in (2), $T_0^{N!}v-v\in\mathfrak m^k\mathcal O^\Sigma$, and $n=N!-1$ will do. Conversely, let $v$ be periodic as in (3), and $k\ge0$; take $n$ for $k$ and $N$ as in (2). Then $T_0^{(n+1)q}v-v\in\mathfrak m^k\mathcal O^\Sigma$ for every $q\ge0$, as $T_0$ preserves $\mathfrak m^k\mathcal O^\Sigma$, and $n+1$ divides $(N+n+1)!$; so $e(v)-v\in\mathfrak m^k\mathcal O^\Sigma$. As $\bigcap_k\mathfrak m^k\mathcal O^\Sigma=0$, $v=e(v)\in P$.

(4) Let $e_\Omega$ be the extension of $e$ to $\Omega^\Sigma$. It is an idempotent that commutes with $T$, and its image is the $\Omega$-span of $P$. Let $y\in\Omega^\Sigma$ be nonzero with $Ty=\mu y$.

\emph{If $e_\Omega y=y$, then $\mu\ne0$ and $\mu$ and $\mu^{-1}$ are integral over $\mathcal O$.} Let $V$ be the endomorphism of $\mathcal O^\Sigma$ that is the inverse of $T_0$ on $P$ composed with $e$, so that $VT_0e=e$. Then $\mu V_\Omega y=y$, so $\mu\ne0$ and $V_\Omega y=\mu^{-1}y$. By the Cayley-Hamilton theorem a monic polynomial $p$ over $\mathcal O$ has $p(T_0)=0$, so $p(\mu)y=p(T)y=0$ and $\mu$ is integral over $\mathcal O$; in the same way $\mu^{-1}$ is.

\emph{If $e_\Omega y=0$ and $\mu\ne0$, then $\mu^{-1}$ is not integral over $\mathcal O$.} By (1), $S=T_0^n(1-e)$ maps $\mathcal O^\Sigma$ into $\mathfrak m\mathcal O^\Sigma$ for some $n$, and $S_\Omega y=\mu^ny$. By the Cayley-Hamilton theorem for an endomorphism with image in $\mathfrak m\mathcal O^\Sigma$, a monic polynomial $X^r+\sum_{i<r}a_iX^i$ with $a_i\in\mathfrak m$ kills $S$; so $\mu^{nr}+\sum_ia_i\mu^{ni}=0$, and with $\nu=\mu^{-n}$, $1=-\sum_ia_i\nu^{r-i}$. If $\mu^{-1}$ were integral over $\mathcal O$, so would be $\nu$, and $1$ would lie in $\mathfrak m\mathcal O[\nu]$, which is impossible as $\mathcal O[\nu]$ is integral over $\mathcal O$.

\emph{Conclusion.} The space $\Omega^\Sigma$ is the direct sum of the generalized eigenspaces $W_\mu$ of $T$, as $\Omega$ is algebraically closed, and $e_\Omega$ maps each $W_\mu$ into itself. Let $\mu$ be such that $\mu\ne0$ and $\mu$ and $\mu^{-1}$ are integral. If $(1-e_\Omega)W_\mu\ne0$, then it contains an eigenvector $y$ of $T$, with eigenvalue $\mu$ and $e_\Omega y=0$, contradicting the second step; so $W_\mu$ lies in the image of $e_\Omega$. For any other $\mu$, if $e_\Omega W_\mu\ne0$, it contains an eigenvector $y$ with eigenvalue $\mu$ and $e_\Omega y=y$, contradicting the first step; so $e_\Omega W_\mu=0$. Hence the image of $e_\Omega$ is the sum of the $W_\mu$ over the $\mu$ of the statement.
% (3) is ordProjective (Module.Projective.of_split) with Module.free_of_flat_of_isLocalRing, and ordPer (ordFactLim both ways, IsHausdorff). (4) is span_ordinaryPart_eq_sup_unit_generalized_eigenspaces (TensorProduct.piScalarRight for the base change) and eigOrdInt: eigOrdInt_inv for V, eigUnitInt (eigenvalue_integral_of_finite, LinearMap.exists_monic_and_aeval_eq_zero), eigNilInt (LinearMap.exists_monic_and_natDegree_eq_and_coeff_mem_pow_and_aeval_eq_zero, Ideal.map_eq_top_iff_of_ker_le) and eigSum (exists_eigenvector_fixed_idempotent, idempotent_kills_generalized_eigenspace, iSup_maxGenEigenspace_eq_top).
\end{proof}

\begin{lemma}\label{sha-hecke-ideal-tw-augmented-hida-inputs-ordinary-lattice}
Let $\mathcal O$ be a complete local noetherian domain with maximal ideal $\mathfrak m$ and finite residue field, contained in an algebraically closed field $\Omega$. Let $V$ be an $\Omega$-vector space, $C\subset V$ a subspace, $T$ an endomorphism of $V$ that maps $C$ into itself, $\Sigma$ a finite set and $\operatorname{ev}:V\to\Omega^\Sigma$ a linear map that maps $C$ bijectively onto $\Omega^\Sigma$. Suppose that for $F\in C$ with $\operatorname{ev}(F)\in\mathcal O^\Sigma$, $\operatorname{ev}(TF)\in\mathcal O^\Sigma$. Let $C^{\mathrm{ord}}$ be the sum, over the nonzero $\mu\in\Omega$ such that $\mu$ and $\mu^{-1}$ are integral over $\mathcal O$, of the generalized $\mu$-eigenspaces of $T$ in $C$. Then there is an $\mathcal O$-submodule $S\subset C$, free over $\mathcal O$, such that:
\begin{enumerate}
\item the $\Omega$-span of $S$ is $C^{\mathrm{ord}}$;
\item every $F\in S$ has $\operatorname{ev}(F)\in\mathcal O^\Sigma$, and for every $k\ge0$ there is $n\ge0$ with $\operatorname{ev}(T^{n+1}F-F)\in(\mathfrak m^k)^\Sigma$;
\item for every $F\in C$ with $\operatorname{ev}(F)\in\mathcal O^\Sigma$ there is $F'\in S$ such that for every $k\ge0$ there is $J$ with $\operatorname{ev}(T^{j!}F-F')\in(\mathfrak m^k)^\Sigma$ for all $j\ge J$.
\end{enumerate}
\end{lemma}
\leanref{Development/AutomorphicArithmetic/AutomorphicVarieties/BuildingsResidueChars/ExistsFreeLatticeSpanningOrdinaryPart.lean}{UpperTheory.exists_free_lattice_spanning_ordinary_part}
\begin{proof}
Let $g:\Omega^\Sigma\to C$ be the inverse of $\operatorname{ev}$ on $C$, and $T'=\operatorname{ev}\circ T\circ g$, an endomorphism of $\Omega^\Sigma$ with $g\circ T'=T\circ g$. By hypothesis $T'$ maps $\mathcal O^\Sigma$ into itself; let $T_0$ be its restriction, and let $P$ and $e$ be as in Lemma~\ref{sha-hecke-ideal-tw-augmented-hida-inputs-fitting} for $T_0$. Let $S=g(P)$, which is free over $\mathcal O$, by Lemma~\ref{sha-hecke-ideal-tw-augmented-hida-inputs-fitting}(3), as $g$ is injective.

(1) As $g$ intertwines $T'$ and $T|_C$, it maps the generalized $\mu$-eigenspace of $T'$ onto the generalized $\mu$-eigenspace of $T$ in $C$. So, by Lemma~\ref{sha-hecke-ideal-tw-augmented-hida-inputs-fitting}(4), the $\Omega$-span of $S$, which is $g$ of the $\Omega$-span of $P$, is $C^{\mathrm{ord}}$.

(2) For $F=g(v)$ with $v\in P$, $\operatorname{ev}(F)=v\in\mathcal O^\Sigma$ and $\operatorname{ev}(T^{n+1}F-F)=T_0^{n+1}v-v$, which lies in $(\mathfrak m^k)^\Sigma$ for some $n$, by Lemma~\ref{sha-hecke-ideal-tw-augmented-hida-inputs-fitting}(3).

(3) Let $v=\operatorname{ev}(F)$, so $F=g(v)$, and $F'=g(e(v))\in S$. Then $\operatorname{ev}(T^{j!}F-F')=T_0^{j!}v-e(v)$, which lies in $(\mathfrak m^k)^\Sigma$ for $j$ large, by Lemma~\ref{sha-hecke-ideal-tw-augmented-hida-inputs-fitting}(2).
% exists_free_lattice_spanning_ordinary_part: e (LinearEquiv.ofBijective evC), g, T', jL, T0 from hint', key (span_ordinaryPart_eq_sup_unit_generalized_eigenspaces), L = g jL, S = (ordPart T0).map L; free by ordProjective and Module.free_of_flat_of_isLocalRing; the span through heig (pow_intertwine); hSprop through ordPer; the approximation through ordExists and ordFactLim with ordProj.
\end{proof}

\begin{lemma}[{\cite{GGPS69,BJ79}}]\label{sha-hecke-ideal-tw-augmented-hida-inputs-automorphic-sum}
Let $E$ be a CM field with maximal totally real subfield $F$, let $n\ge1$, and let $U_n$ be as in Definition~\ref{gal-definite-unitary}. Let $\varphi:U_n\to\mathbf C$ be a continuous function, left invariant under $U_n(F)$, invariant on the right under an open subgroup of the finite part of $U_n$, whose right translates under the part of $U_n$ at infinity span a finite-dimensional space (an \emph{automorphic form} on $U_n$). Then $\varphi$ lies in a finite sum of automorphic representations of $U_n$ (Definition~\ref{gal-definite-unitary}(2)).
\end{lemma}
\leanref{Development/AutomorphicArithmetic/AutomorphicVarieties/BuildingsResidueChars/MemISupAutomorphicRepsOfInvariantForm.lean}{UpperTheory.mem_iSup_automorphicReps_of_invariant_form}
\leanref{Development/AutomorphicArithmetic/AutomorphicVarieties/BuildingsResidueChars/ExistsNondegenerateRightInvariantFormOfCompactQu.lean}{UpperTheory.exists_nondegenerate_right_invariant_form_of_compact_quotient}
\leanref{Development/AutomorphicArithmetic/AutomorphicVarieties/BuildingsResidueChars/ExistsCompactFundamentalSetUnitary.lean}{UpperTheory.exists_compact_fundamental_set_unitary}

\begin{lemma}\label{sha-hecke-ideal-tw-augmented-hida-inputs-automorphic-projections}
Let $E$, $F$, $n$ and $U_n$ be as in Lemma~\ref{sha-hecke-ideal-tw-augmented-hida-inputs-automorphic-sum}, and let $\mathcal A$ be the space of automorphic forms on $U_n$. Then $\mathcal A$ is stable under right translation by $U_n$, and for every finite-dimensional subspace $M\subset\mathcal A$ there are automorphic representations $\sigma_1,\dots,\sigma_r$ of $U_n$ and linear maps $P_1,\dots,P_r:\mathcal A\to\mathcal A$ that commute with right translation by $U_n$, such that $P_t(\varphi)\in\sigma_t$ for every $t$ and $\varphi\in\mathcal A$, and $\sum_tP_t(\varphi)=\varphi$ for $\varphi\in M$.
\end{lemma}
\leanref{Development/AutomorphicArithmetic/AutomorphicVarieties/HeckeAlgebras/OrdinaryHeckeCharacters/DecompU.lean}{autProj}
\leanref{Development/AutomorphicArithmetic/AutomorphicVarieties/BuildingsResidueChars/ExistsEquivariantProjectionsOfSumIrreducible.lean}{UpperTheory.exists_equivariant_projections_of_sum_irreducible}
\begin{proof}
Write $R(x)$ for right translation by $x\in U_n$, and $x=x_\infty x_{\mathrm f}$ for its parts at infinity and finite, which commute with each other. Let $\varphi\in\mathcal A$, fixed on the right by an open subgroup $S$ of the finite part. The function $R(x)\varphi$ is continuous and left invariant under $U_n(F)$. For $k\in x_{\mathrm f}Sx_{\mathrm f}^{-1}$, an open subgroup, $R(k)R(x)\varphi=R(x)R(x^{-1}kx)\varphi=R(x)\varphi$, as $x^{-1}kx=x_{\mathrm f}^{-1}kx_{\mathrm f}\in S$. For $y$ in the part at infinity, $R(y)R(x)\varphi=R(x)R(x_\infty^{-1}yx_\infty)\varphi$, so the right translates of $R(x)\varphi$ under the part at infinity lie in the image under $R(x)$ of the finite-dimensional span of those of $\varphi$. So $R(x)\varphi\in\mathcal A$.

Every automorphic representation of $U_n$ is a subspace of $\mathcal A$, stable and irreducible under $R$. By Lemma~\ref{sha-hecke-ideal-tw-augmented-hida-inputs-automorphic-sum}, $\mathcal A$ is the sum of such subspaces, so it is a semisimple representation of $U_n$, a direct sum $\mathcal A=\bigoplus_{\alpha\in I}W_\alpha$ of irreducible subrepresentations. Each $W_\alpha$ is a nonzero space of automorphic forms, stable and irreducible under right translation, so an automorphic representation. Let $\varphi_1,\dots,\varphi_s$ span $M$, let $\alpha_1,\dots,\alpha_r$ be the finitely many indices at which some $\varphi_k$ has a nonzero component, and let $P_t$ be the projection of $\mathcal A$ onto $W_{\alpha_t}$ along the other summands. It commutes with $R$, as every $W_\alpha$ is stable under $R$; and for $\varphi\in M$ the components of $\varphi$ at the other indices vanish, so $\sum_tP_t(\varphi)=\varphi$.
% autProj: unitary_adelic_t2_locallyCompact_secondCountable, exists_compact_fundamental_set_unitary and exists_nondegenerate_right_invariant_form_of_compact_quotient (the form B for mem_iSup_automorphicReps_of_invariant_form), then exists_equivariant_projections_of_sum_irreducible for autRep on autU (ss_proj: IsSemisimpleModule.exists_linearEquiv_dfinsupp, the supports of a finite spanning set of M, the projections Q onto the summands), and the passage between the subrepresentations of autU and the IsAutRepU subspaces of functions. The stability of autU is RU_mem_autU.
\end{proof}

\begin{lemma}\label{sha-hecke-ideal-tw-augmented-hida-inputs-automorphic-forms}
In Situation~\ref{sha-hecke-ideal-tw-augmented-setup}, let $\lambda$ be a weight, $\xi$ a representation of $G_\ell$ on $\mathbf C^d$, algebraic of weight $\lambda$ through $\iota$ (Definition~\ref{sha-hecke-ideal-tw-augmented-parts}), $X=\bigotimes_{(v,\tau)}\nabla(\lambda_\tau)$, and $e:\mathbf C^d\cong X$ a linear isomorphism that carries $\xi$ to the representation $g\mapsto\bigotimes_{(v,\tau)}\rho_{\lambda_\tau}(\iota\tau(g_v))$. For $g\in U(n)(\mathbf A_{L^+})$ let $g^\infty=(g^{(\ell)},g_\ell)\in G^{(\ell)}\times G_\ell$ be its finite part, through Situation~\ref{sha-hecke-ideal-tw-augmented-setup}(1), and $\rho_\infty(g)=\bigotimes_{(v,\tau)}\rho_{\lambda_\tau}(g_{\iota\tau})$, where $g_{\iota\tau}\in GL_n(\mathbf C)$ is the component of $g$ at the infinite place given by the embedding $\iota\tau|_L$. For a function $F:G^{(\ell)}\times G_\ell\to\mathbf C^d$ and a linear form $l$ on $X$, let $\varphi_{F,l}(g)=l\big(\rho_\infty(g)^{-1}e(F(g^\infty))\big)$ for $g\in U(n)(\mathbf A_{L^+})$. Then:
\begin{enumerate}
\item if $\varphi_{F,l}=0$ for every $l$, then $F=0$;
\item if $F(\gamma h)=\xi(\gamma_\ell)F(h)$ for $\gamma\in U(n)(L^+)$ and $h\in G^{(\ell)}\times G_\ell$, and the $k\in U(n)(\mathbf A^\infty_{L^+})$ with $F(hk^\infty)=F(h)$ for all $h$ form an open subgroup, then every $\varphi_{F,l}$ is an automorphic form on $U(n)$ (Lemma~\ref{sha-hecke-ideal-tw-augmented-hida-inputs-automorphic-sum});
\item for $x\in U(n)(\mathbf A_{L^+})$, the function $g\mapsto\varphi_{F,l}(gx)$ is $\varphi_{F',l'}$, with $F'(h)=F(hx^\infty)$ and $l'=l\circ\rho_\infty(x)^{-1}$;
\item if $\Phi$ is a linear map from the dual of $X$ to the space of automorphic forms on $U(n)$ such that $\Phi(l)(gx)=\Phi(l\circ\rho_\infty(x)^{-1})(g)$ for all $l$, all $g$ and all $x\in U(n)(L^+\otimes\mathbf R)$, then there is $F:G^{(\ell)}\times G_\ell\to\mathbf C^d$ with $F(\gamma h)=\xi(\gamma_\ell)F(h)$ for $\gamma\in U(n)(L^+)$ and $h\in G^{(\ell)}\times G_\ell$, and $\varphi_{F,l}=\Phi(l)$ for every $l$.
\end{enumerate}
\end{lemma}
\leanref{Development/AutomorphicArithmetic/AutomorphicVarieties/HeckeAlgebras/OrdinaryHeckeCharacters/PlaceBetween.lean}{UpperTheory.V.Hi.formU}
\leanref{Development/AutomorphicArithmetic/AutomorphicVarieties/HeckeAlgebras/IsTheta.lean}{UpperTheory.V.Hi.CaseU.toForm}
\leanref{Development/AutomorphicArithmetic/AutomorphicVarieties/HeckeAlgebras/IsTheta.lean}{UpperTheory.V.Hi.CaseU.infRep}
\begin{proof}
\emph{Preliminaries.} The map $g\mapsto g^\infty$ is a homomorphism with kernel $U(n)(L^+\otimes\mathbf R)$, and it maps $U(n)(\mathbf A^\infty_{L^+})$ bijectively onto $G^{(\ell)}\times G_\ell$. The map $\rho_\infty$ is a homomorphism, trivial on $U(n)(\mathbf A^\infty_{L^+})$, and $g\mapsto l(\rho_\infty(g)u)$ is continuous for $u\in X$ and every $l$, as the matrix coefficients of $\rho_{\lambda_\tau}$ are polynomials in the entries and a power of the determinant. Let $\gamma\in U(n)(L^+)$, a matrix in $GL_n(L)$. Its components at the infinite place given by $\iota\tau|_L$ and at $\tilde v$ are images of the same matrix, so $\gamma_{\iota\tau}=\iota\tau(\gamma_v)$, and $\rho_\infty(\gamma)e=e\,\xi(\gamma_\ell)$ by the choice of $e$. Write $F_e=e\circ F$, so that $\varphi_{F,l}(g)=l(\rho_\infty(g)^{-1}F_e(g^\infty))$.

(1) Every element of $G^{(\ell)}\times G_\ell$ is $g^\infty$ for some $g$. If every $l$ vanishes at $\rho_\infty(g)^{-1}F_e(g^\infty)$, then this vector is $0$, and so is $F_e(g^\infty)$.

(2) Let $\varphi=\varphi_{F,l}$ and $S$ the open subgroup of the statement. For $\gamma\in U(n)(L^+)$, $F_e(\gamma^\infty g^\infty)=e\,\xi(\gamma_\ell)F(g^\infty)=\rho_\infty(\gamma)F_e(g^\infty)$, so $\varphi(\gamma g)=l\big(\rho_\infty(g)^{-1}\rho_\infty(\gamma)^{-1}\rho_\infty(\gamma)F_e(g^\infty)\big)=\varphi(g)$. For $k\in S$, $\varphi(gk)=l\big(\rho_\infty(g)^{-1}\rho_\infty(k)^{-1}F_e(g^\infty k^\infty)\big)=\varphi(g)$, as $\rho_\infty(k)=1$. Let $g_0\in U(n)(\mathbf A_{L^+})$. For $y$ in a neighbourhood of $1$, the finite part of $y$ lies in $S$, so $F_e((g_0y)^\infty)=F_e(g_0^\infty)$, and $\varphi(g_0y)=l\big(\rho_\infty(g_0y)^{-1}F_e(g_0^\infty)\big)$, which is continuous in $y$; so $\varphi$ is continuous at $g_0$. For $y\in U(n)(L^+\otimes\mathbf R)$, $y^\infty=1$ and $\varphi(gy)=(l\circ\rho_\infty(y)^{-1})\big(\rho_\infty(g)^{-1}F_e(g^\infty)\big)$. So the right translates of $\varphi$ under $U(n)(L^+\otimes\mathbf R)$ lie in the image of the linear map $\mu\mapsto\big(g\mapsto\mu(\rho_\infty(g)^{-1}F_e(g^\infty))\big)$ on the dual of $X$, which is finite-dimensional.

(3) $\varphi_{F,l}(gx)=l\big(\rho_\infty(x)^{-1}\rho_\infty(g)^{-1}F_e(g^\infty x^\infty)\big)=\varphi_{F',l'}(g)$.

(4) As $X$ is finite-dimensional, for each $g$ there is a unique $\Psi(g)\in X$ with $l(\Psi(g))=\Phi(l)(g)$ for every $l$. For $x\in U(n)(L^+\otimes\mathbf R)$, $l(\Psi(gx))=\Phi(l\circ\rho_\infty(x)^{-1})(g)=l\big(\rho_\infty(x)^{-1}\Psi(g)\big)$, so $\Psi(gx)=\rho_\infty(x)^{-1}\Psi(g)$ and $\rho_\infty(gx)\Psi(gx)=\rho_\infty(g)\Psi(g)$. As $g\mapsto g^\infty$ is onto with kernel $U(n)(L^+\otimes\mathbf R)$, there is a function $F_e:G^{(\ell)}\times G_\ell\to X$ with $F_e(g^\infty)=\rho_\infty(g)\Psi(g)$ for all $g$. Each $\Phi(l)$ is left invariant under $U(n)(L^+)$, so $\Psi(\gamma g)=\Psi(g)$ and $F_e(\gamma^\infty g^\infty)=\rho_\infty(\gamma)F_e(g^\infty)$ for $\gamma\in U(n)(L^+)$. Let $F=e^{-1}\circ F_e$. Then $F(\gamma h)=\xi(\gamma_\ell)F(h)$, by the preliminaries, and $\varphi_{F,l}(g)=l(\Psi(g))=\Phi(l)(g)$.
% formU (PlaceBetween.lean, 352 lines): J' = (away, cLoc at the s.w v) with hsurj (exists_adjust, the surjectivity on finU) and hker (cLoc_conj, mem_inf_of, cLoc_inf: the kernel is infU); hfin and hcont for infRep (cont_multi, cont_sq, cont_cEmb); hrat (cEmb_rat) and hAF, the AF condition through IsBasis; then the private core: (1) through Module.forall_dual_apply_eq_zero_iff, (2) through aut_coef (continuity with the continuous finite part exists_fin, left invariance, the open level by isOpen_stab_of_mem_nhds, kfin through the map C on the dual), (3) by computation, and (4) through Psi = (Module.evalEquiv C X).symm and key. The book's identification of U(n)(A^infinity) with G^(l) x G_l (Situation (1)) gives the surjectivity on the finite part.
\end{proof}

\begin{lemma}\label{sha-hecke-ideal-tw-augmented-hida-inputs-automorphic-near}
In Situation~\ref{sha-hecke-ideal-tw-augmented-setup}, let $K=(K_j)$ be a family of compact open subgroups $K_j\subset GL_n(L_{w_j})$ and $b\ge0$. Then there is a neighbourhood $W$ of $1$ in $U(n)(\mathbf A^\infty_{L^+})$ such that every $k\in W$ has $k^\infty=(u,y)$, in the notation of Lemma~\ref{sha-hecke-ideal-tw-augmented-hida-inputs-automorphic-forms}, with $u\in\mathcal K_0K$, $\tau_0(u)=1$ and $y\in I_1(b)$ (Lemma~\ref{sha-hecke-ideal-tw-augmented-uop}).
\end{lemma}
\leanref{Development/AutomorphicArithmetic/AutomorphicVarieties/BuildingsResidueChars/EventuallyMemHeckeLevelAndIwahoriCongruence.lean}{UpperTheory.eventually_mem_heckeLevel_and_iwahoriCongruence}
\begin{proof}
The elements of $U(n)(\mathbf A^\infty_{L^+})$ whose components at all finite places are integral form an open subgroup, on which the component at each place is continuous. The following conditions on such an element are open, as they concern finitely many places and open subgroups, and $1$ satisfies them: the components at the places of $S_{\mathrm a}$ lie in the kernel of $\tau_0$, which is open (Lemma~\ref{sha-hecke-data}(1)); the component at $\tilde v_0$ lies in the principal congruence subgroup of level $\tilde v_0$; the component at each $w_j$ lies in $K_j$; and the component at each $\tilde v$, $v\mid\ell$, is congruent to $1$ modulo $\varpi^{b+1}$. Let $W$ be the set of the integral elements that satisfy them, and $k\in W$ with $k^\infty=(u,y)$. Then $u\in G^{(\ell)}_0$, $u\in\mathcal K_0K$ by the description of the level in Situation~\ref{sha-hecke-ideal-tw-augmented-setup}(3), and $\tau_0(u)=1$. Each component of $y$ is a matrix over $\mathcal O_{L_{\tilde v}}$ congruent to $1$ modulo $\varpi^{b+1}$, so by the Iwahori factorization it is a product of an element of $N_0$, an element of $T_0^{(b)}$ and an element of $\bar N_b$; so $y\in I_1(b)$.
% eventually_mem_heckeLevel_and_iwahoriCongruence: unitaryAdelic_localComponents_continuous_levels_open (hc, hK, hint), h1 (the kernel of tau0 through isOpen_induced_iff and cmp_cont), h2 (Kc 1 at s.w s.v0), h3 (the K q), h4 (Kc (b + 1) at the s.w v); mem_heckeLevel_of_local_components_mem for the level, and glBridge (Kc (c + 1) <= I1 c) with pi_mem_first_iwahori_filtration for the Iwahori part.
\end{proof}

\begin{lemma}[{\cite{Fla79,Car79}}]\label{sha-hecke-ideal-tw-augmented-hida-inputs-automorphic-local}
Let $E$, $F$, $n$ and $U_n$ be as in Lemma~\ref{sha-hecke-ideal-tw-augmented-hida-inputs-automorphic-sum}, $\sigma$ an automorphic representation of $U_n$, and $w$ a finite place of $E$ split over $F$, with $U_{n,w}$ and $c_w$ as in Definition~\ref{gal-definite-unitary}(3), and $q_w$ the norm of $w$. Then:
\begin{enumerate}
\item $\sigma$ has a local component at $w$, which is spherical if and only if $\sigma$ is unramified at $w$;
\item for every local component $\tau$ of $\sigma$ at $w$, $\sigma$ is the sum of the images of the linear maps $f:\tau\to\sigma$ with $f(\tau(c_w(x))u)=R(x)f(u)$ for $x\in U_{n,w}$, where $R$ is right translation;
\item an irreducible smooth representation $\tau$ of $GL_n(E_w)$ that is spherical has Satake parameters $B$, a multiset of $n$ nonzero complex numbers, and for $0\le k\le n$ the double coset operator $[GL_n(\mathcal O_w)\alpha_kGL_n(\mathcal O_w)]$, with $\alpha_k$ as in Definition~\ref{sha-hecke-local-ordinary}, acts on the vectors of $\tau$ fixed by $GL_n(\mathcal O_w)$ by $q_w^{k(n-k)/2}e_k(B)$.
\end{enumerate}
\end{lemma}
\leanref{Development/AbelianVarietyTheory/Proofs/IwasawaDecomposition/ExistsLocCompSplit.lean}{GlobalTheory.DefU.exists_locComp_split}
\leanref{Development/AutomorphicArithmetic/AutomorphicVarieties/BuildingsResidueChars/IsoSumRestrictionOfLocalComponentUnitary.lean}{UpperTheory.isoSum_restriction_of_local_component_unitary}
\leanref{Development/AutomorphicArithmetic/AutomorphicVarieties/ModularityArgument/ThresholdArguments/DenseThresholdCases/TSatLoc.lean}{UpperTheory.V.Af.satake_local_of_printed_results}
\leanref{Development/AutomorphicArithmetic/AutomorphicVarieties/HeckeAlgebras/OrdinaryHeckeCharacters/SatTUnif.lean}{UpperTheory.V.Hi.satT}

\begin{lemma}\label{sha-hecke-ideal-tw-augmented-hida-inputs-conjugate-dual}
Let $n\ge1$, $\lambda$ a dominant weight of length $n$, $\lambda^\vee=(-\lambda_n\ge\dots\ge-\lambda_1)$, and $\nabla$, $\rho$ and the Fock pairing $\langle\ ,\ \rangle_F$ as in Definition~\ref{is-polynomial-model} over $\mathbf C$. Let $U(n)=\{h\in GL_n(\mathbf C):{}^t\bar hh=1\}$, where $\bar h$ is the matrix of the complex conjugates of the entries of $h$; for $h\in U(n)$, $\bar h\in U(n)$. Then:
\begin{enumerate}
\item there is a linear isomorphism $\Psi:\nabla(\lambda)\to\nabla(\lambda^\vee)$ with $\Psi(\rho_\lambda(\bar h)f)=\rho_{\lambda^\vee}(h)\Psi(f)$ for all $h\in U(n)$ and $f\in\nabla(\lambda)$;
\item the map $f\mapsto\langle f,\cdot\rangle_F$ is a linear isomorphism from $\nabla(\lambda)$ onto its dual, and $\langle\rho_\lambda(\bar h)f,f'\rangle_F=\langle f,\rho_\lambda(h^{-1})f'\rangle_F$ for all $h\in U(n)$ and $f,f'\in\nabla(\lambda)$.
\end{enumerate}
In particular, the map $e:\Psi(f)\mapsto\langle f,\cdot\rangle_F$ is a linear isomorphism from $\nabla(\lambda^\vee)$ onto the dual of $\nabla(\lambda)$ with $e(\rho_{\lambda^\vee}(h)x)=e(x)\circ\rho_\lambda(h)^{-1}$ for $h\in U(n)$ and $x\in\nabla(\lambda^\vee)$.
\end{lemma}
\leanref{Development/AbelianVarietyTheory/Proofs/TwistedOperators/UnitaryIrrepCompConjIsoDualStatement.lean}{GlobalTheory.T.unitaryIrrep_comp_conj_iso_dual}
\leanref{Development/AbelianVarietyTheory/Proofs/TwistedOperators/IsoUnitaryIrrepConjDualStatement.lean}{GlobalTheory.T.iso_unitaryIrrep_conj_dual}
\leanref{Development/AbelianVarietyTheory/Proofs/SatakeParameters/DualIsoDetGenMNeZero.lean}{GlobalTheory.ValCmp.exists_posOk_dual}
\leanref{Development/AbelianVarietyTheory/Proofs/SatakeParameters/ExistsCofactorTwistDualWeylStatement.lean}{GlobalTheory.T.exists_cofactorTwist_dualWeyl}
\leanref{Development/AbelianVarietyTheory/Proofs/SatakeParameters/CofactorTwistRelationRowActEquivariantStatement.lean}{GlobalTheory.T.cofactorTwistRelation_rowAct_equivariant}
\leanref{Development/AbelianVarietyTheory/Proofs/SatakeParameters/CofactorTwistRelationDualWeightInvolutiveStatement.lean}{GlobalTheory.T.cofactorTwistRelation_dualWeight_involutive}
\leanref{Formalism/GradedWeylDivisibility/DualWeylDominanceOf.lean}{Formalism.GradedWeylDivisibility.fockRowAdj}
\begin{proof}
Let $\nu=\lambda^\natural$ and $s=\operatorname{sh}(\lambda)$, so that $\nu_n=0$ and $\lambda_1=s+\nu_1$; then $(\lambda^\vee)^\natural=\nu^*$ with $\nu^*_i=\nu_1-\nu_{n+1-i}$, and $\operatorname{sh}(\lambda^\vee)=-\lambda_1$. For $h\in U(n)$, ${}^t\bar h=h^{-1}$ and $\det\bar h=(\det h)^{-1}$, so $(\rho_\lambda(\bar h)f)(x)=\det(h)^{-s}f(h^{-1}x)$.

(1) \emph{The cofactor substitution.} Let $x=(x_{ab})$ be the matrix of the variables, $\operatorname{adj}(x)$ its adjugate, $\tilde x$ the matrix with entries $\tilde x_{ab}=\operatorname{adj}(x)_{n+1-b,a}$, and $|\nu|=\nu_1+\dots+\nu_n$. For $f\in\nabla(\nu)$ there is a polynomial $\Psi(f)\in\nabla(\nu^*)$ with $\det(x)^{|\nu|-\nu_1}\Psi(f)=f(\tilde x)$. It is unique, as $\det(x)$ is not a zero divisor, so $\Psi$ is linear. The same construction for $\nu^*$ carries $\Psi(f)$ to $\varepsilon f$, with $\varepsilon=\pm1$ independent of $f$; so $\Psi$ is injective and maps $\nabla(\nu)$ onto $\nabla(\nu^*)$. For $h\in GL_n(\mathbf C)$ and $f\in\nabla(\nu)$, the polynomial $f(h^{-1}x)$ lies in $\nabla(\nu)$, and $\Psi$ carries it to $\det(h)^{-\nu_1}\Psi(f)({}^thx)$. For $h\in U(n)$ this gives $\Psi(\rho_\lambda(\bar h)f)(x)=\det(h)^{-s-\nu_1}\Psi(f)({}^thx)=(\rho_{\lambda^\vee}(h)\Psi(f))(x)$, as $-s-\nu_1=\operatorname{sh}(\lambda^\vee)$.

(2) \emph{The Fock pairing.} For every $A\in M_n(\mathbf C)$ and all polynomials $f,f'$, $\langle f(Ax),f'\rangle_F=\langle f,f'({}^tAx)\rangle_F$. For $h\in U(n)$, taking $A=h^{-1}$ gives $\langle\rho_\lambda(\bar h)f,f'\rangle_F=\det(h)^{-s}\langle f,f'({}^th^{-1}x)\rangle_F=\langle f,\rho_\lambda(h^{-1})f'\rangle_F$. Let $f\in\nabla(\lambda)$ be nonzero, and $\bar f$ the polynomial of the complex conjugates of its coefficients. As the conditions defining $\nabla(\nu)$ have rational coefficients, $\bar f\in\nabla(\lambda)$, and $\langle f,\bar f\rangle_F=\sum_m|f_m|^2m!>0$. So $f\mapsto\langle f,\cdot\rangle_F$ is injective, and it is onto the dual, which has the same dimension.

For the last claim, $e(\rho_{\lambda^\vee}(h)\Psi(f))=e(\Psi(\rho_\lambda(\bar h)f))=\langle\rho_\lambda(\bar h)f,\cdot\rangle_F=\langle f,\rho_\lambda(h)^{-1}(\cdot)\rangle_F$ by (1) and (2).
\end{proof}
% iso_irr_comp_conjU (UnitaryIrrepCompConjIsoDual.lean, 102 lines) is (1): Psi from ValCmp.exists_posOk_dual with e = detShift and c = 0, that is dualIso_exists_map (the cofactor twist cofAct, x_{a,s} -> adj(x)_{rev s, a}, through CofRel), injective and onto by cofactorTwistRelation_dualWeight_involutive and exists_cofactorTwist_dualWeyl, equivariant by cofactorTwistRelation_rowAct_equivariant; sqAct_eq_rnd_invTranspose and hg (the inverse transpose of conj g is g) give the action. ConjDual.iso_conj_dual (IsoUnitaryIrrepConjDual.lean, 79 lines) is (2): bil f f' = fockF f (conj f'), the bilinear Fock pairing; bil_rowAct from fockRowAdj, detPow_conjU, and toDual_injective from bil_conj, fockPos and dualWeylMap. The cofactor lemmas are about 360 lines in SatakeParameters.

\begin{lemma}\label{sha-hecke-ideal-tw-augmented-hida-inputs-automorphic-weight}
In the notation of Lemma~\ref{sha-hecke-ideal-tw-augmented-hida-inputs-automorphic-forms}, let $\sigma$ be an automorphic representation of $U(n)$ (Definition~\ref{gal-definite-unitary}) and $F:G^{(\ell)}\times G_\ell\to\mathbf C^d$ a nonzero function with $\varphi_{F,l}\in\sigma$ for every linear form $l$ on $X$. Then for every pair $(v,\tau)$ of a place $v\mid\ell$ and a continuous embedding $\tau:L_{\tilde v}\to\bar{\mathbf Q}_\ell$, $\sigma$ has weight $(-\lambda_{\tau,n}\ge\dots\ge-\lambda_{\tau,1})$ at $\iota\circ\tau|_L$ (Definition~\ref{gal-definite-unitary}(4)).
\end{lemma}
\leanref{Development/AutomorphicArithmetic/AutomorphicVarieties/HeckeAlgebras/OrdinaryHeckeCharacters/WtOfForm.lean}{UpperTheory.V.Hi.partUR2Wt}
\leanref{Development/AutomorphicArithmetic/AutomorphicVarieties/BuildingsResidueChars/HasWeightAtOfTensorDualIntertwiner.lean}{UpperTheory.hasWeightAt_of_tensor_dual_intertwiner}
\leanref{Development/AutomorphicArithmetic/AutomorphicVarieties/BuildingsResidueChars/InjectiveInfinitePlaceOfEmbAtSplitPlaces.lean}{UpperTheory.injective_infinitePlace_of_emb_at_split_places}
\begin{proof}
For a pair $y=(v,\tau)$ write $\phi_y=\iota\circ\tau|_L$ and $\lambda_y=\lambda_\tau$, so that $X=\bigotimes_y\nabla(\lambda_y)$ and $\rho_\infty(g)=\bigotimes_y\rho_{\lambda_y}(g_{\phi_y})$, where $g_{\phi_y}$ is the component of $g$ at the infinite place given by $\phi_y$.

\emph{Distinct infinite places.} If $y=(v,\tau)$ and $y'=(v',\tau')$ are distinct, then $\phi_y$ and $\phi_{y'}$ define distinct infinite places of $L$. Indeed, suppose they define the same place. Then $\tau'|_L=\tau|_L$ or $\tau'|_L=\tau|_L\circ c$, as $\iota$ is injective. As $\tau$ is continuous, an element $x\in\mathcal O_L$ lies in the prime ideal of $\tilde v$ if and only if $|\tau(x)|<1$, and similarly for $\tau'$ and $\tilde v'$. As $c$ is the identity on $L^+$, in both cases the prime ideals of $\mathcal O_{L^+}$ under $\tilde v$ and $\tilde v'$ agree, so $v=v'$. If $\tau'|_L=\tau|_L$, then $\tau'=\tau$, as both are continuous and $L$ is dense in $L_{\tilde v}$; this is excluded. If $\tau'|_L=\tau|_L\circ c$, then $x\in\mathcal O_L$ lies in the prime ideal of $\tilde v$ if and only if $c(x)$ does, so $c(\tilde v)=\tilde v$, which is impossible as $v$ splits in $L$.

\emph{The forms.} The map $\Phi:l\mapsto\varphi_{F,l}$ is linear, from the dual $X^*$ of $X$ to $\sigma$, and it is nonzero by Lemma~\ref{sha-hecke-ideal-tw-augmented-hida-inputs-automorphic-forms}(1), as $F\ne0$. For $x\in U(n)(L^+\otimes\mathbf R)$, $x^\infty=1$, so Lemma~\ref{sha-hecke-ideal-tw-augmented-hida-inputs-automorphic-forms}(3) gives $R(x)\Phi(l)=\Phi(l\circ\rho_\infty(x)^{-1})$, where $R$ is right translation.

\emph{One factor.} Fix $y$, and let $U_y$ be the subgroup of the $x\in U(n)(L^+\otimes\mathbf R)$ whose components at the infinite places other than the one given by $\phi_y$ are $1$. It is the factor at $\phi_y$ of Definition~\ref{gal-definite-unitary}(4), and $x\mapsto x_{\phi_y}$ identifies it with the compact unitary group $\{h\in GL_n(\mathbf C):{}^t\bar hh=1\}$. For $x\in U_y$ and $y'\ne y$, $x_{\phi_{y'}}=1$ by the first step, so $\rho_\infty(x)$ acts by $\rho_{\lambda_y}(x_{\phi_y})$ on the factor $\nabla(\lambda_y)$ and as the identity on the others. As $\Phi\ne0$ and the products $\bigotimes_{y'}f_{y'}$ of linear forms $f_{y'}$ on the factors span $X^*$, there are such $f_{y'}$ with $\Phi(\bigotimes_{y'}f_{y'})\ne0$. For a linear form $m$ on $\nabla(\lambda_y)$ let $E(m)\in X^*$ be the product of $m$ at $y$ and of the $f_{y'}$ at the $y'\ne y$. Then $\Phi\circ E\ne0$, and $E(m)\circ\rho_\infty(x)^{-1}=E\big(m\circ\rho_{\lambda_y}(x_{\phi_y})^{-1}\big)$ for $x\in U_y$. So $\Phi\circ E$ is a nonzero linear map from the contragredient of $\nabla(\lambda_y)$ to $\sigma$ that carries the action of $x_{\phi_y}$ to $R(x)$, for $x\in U_y$.

\emph{The contragredient.} Let $\lambda_y^\vee=(-\lambda_{\tau,n}\ge\dots\ge-\lambda_{\tau,1})$, and let $e$ be the isomorphism of Lemma~\ref{sha-hecke-ideal-tw-augmented-hida-inputs-conjugate-dual} from $\nabla(\lambda_y^\vee)$ onto the dual of $\nabla(\lambda_y)$, which carries $\rho_{\lambda_y^\vee}(h)$ to $m\mapsto m\circ\rho_{\lambda_y}(h)^{-1}$ for $h$ in the compact unitary group. Then $\Phi\circ E\circ e$ is a nonzero linear map from $\nabla(\lambda_y^\vee)$ to $\sigma$ that intertwines its restriction to the compact unitary group with right translation by $U_y$, which is the claim.
% partUR2Wt (WtOfForm.lean): exists_gen gives F in YOf with a component F i /= 0, formU gives h3 (the translation, the book's Lemma sha-hecke-ideal-tw-augmented-hida-inputs-automorphic-forms (3)) and h4 (the surjectivity on finU), then the private wt_of_form, which is this lemma for one function: Phi = toForm iota e F, nonzero through Module.forall_dual_apply_eq_zero_iff and the inverse of infRep, the translation by x in infU through cLoc_inf (J x = 1). hasWeightAt_of_tensor_dual_intertwiner: E from exists_factor_dual_embedding_of_ne_zero_tensor (dualDistrib of a pure tensor of forms f with Phi /= 0 on it), u from EmbUnitary.exists_equiv_unitaryGroup, hT (the other factors act trivially, through the definition of embU and the injectivity), and the isomorphism e of nabla(lambda dual) with the dual of nabla(lambda) on U(n), from T.unitaryIrrep_comp_conj_iso_dual and T.iso_unitaryIrrep_conj_dual, which compare both with nabla(lambda) composed with complex conjugation (Lemma sha-hecke-ideal-tw-augmented-hida-inputs-conjugate-dual). injective_infinitePlace_of_emb_at_split_places: InfinitePlace.mk_eq_iff with complexEmbedding_complexConj, w_under and mem_iff for v = v', denseRange_algebraMap for tau = tau', and not_stable for the conjugate case. DomWeight.dual lam is i -> -lam.w (rev i), the book's (-lambda_{tau,n}, ..., -lambda_{tau,1}).
\end{proof}

\begin{lemma}\label{sha-hecke-ideal-tw-augmented-hida-inputs}
In Situation~\ref{sha-hecke-ideal-tw-augmented-setup}, with the notation of Definition~\ref{sha-hecke-ideal-tw-augmented-hida-weight}:
\begin{enumerate}
\item for every weight $\lambda$ and every ring of integers $\mathcal O_1\supset\mathcal O$ of a finite extension of $\mathbf Q_\ell$ in $\bar{\mathbf Q}_\ell$, there are $\mathcal O'\supset\mathcal O_1$ and a weight datum $(\xi,\chi,c_0)$ over $\mathcal O'$ such that $\xi$ is algebraic of weight $\lambda$ through $\iota$ and $\chi(t)=\iota\big(\prod_{v\mid\ell}\prod_{i=1}^n\prod_\tau\tau(t_{v,i})^{\lambda_{\tau,n+1-i}}\big)$ for $t\in T$ with diagonal entries $t_{v,i}$, the last product over the embeddings $\tau:L_{\tilde v}\to\bar{\mathbf Q}_\ell$, with $\lambda_\tau$ taken for the restriction of $\tau$ to $L$;
\item let $(\xi,\chi,c_0)$ be a weight datum over $\mathcal O'$, $c\ge c_0$, $\theta'$ and $\phi$ characters of $T_0$ with values in $\mathcal O'^\times$, with $\phi$ trivial on $T_0^{(c)}$ and $\theta'(t)=\phi(t)\chi(t)$ for $t\in T_0$, and $U\subset G^{(\ell)}_0$ a subgroup of finite level that satisfies the freeness condition of Lemma~\ref{sha-hecke-ideal-tw-augmented-sections} for $\theta'$; then the control map is injective on $\mathrm{Cl}^{\mathrm{ord}}(U,c,\phi)$, and it maps $\mathrm{Cl}^{\mathrm{ord}}(U,c,\phi)$ onto the $\mathbf C$-span of $\mathcal F^{\mathrm{ord}}_{\mathcal O'}(U,\theta')$;
\item let $\lambda$ be a weight and $\xi$ a representation of $G_\ell$ on some $\mathbf C^d$, algebraic of weight $\lambda$ through $\iota$; then every $F\in V^\xi$ with $F(\gamma g)=\xi(\gamma_\ell)F(g)$ for $\gamma\in U(n)(L^+)$, fixed by the elements of $\mathcal K_0K_w$ whose components at the $v_j$ are $1$, fixed by some level $\mathcal K_0K$, and invariant under right translation by $I_1(e)$ for some $e\ge0$, lies in the sum of the automorphic parts of weight $\xi$ (Definition~\ref{sha-hecke-ideal-tw-augmented-parts});
\item let $K$, $q$, $\varpi$, $n$ and the groups and operators be as in Definition~\ref{sha-hecke-local-ordinary}, $\pi$ a representation of $GL_n(K)$ on a complex vector space $M$, $Y\subset M$ a subspace stable under $\pi$ that is the sum of the images of the equivariant maps from an irreducible smooth representation $\sigma$ to $M$ with image in $Y$, and $E\subset Y$ a nonzero subspace; let $c\ge0$ and $f_1,\dots,f_n:\mathcal O_K^\times\to\mathbf C^\times$ maps such that $f_i(u)=1$ for $u\equiv1$ modulo $\varpi^{c+1}$ and $f_i\ne f_j$ for $i\ne j$; suppose that $E$ is fixed by $N_0$ and $\bar N_c$, that $\pi(e_i(u))F=f_i(u)^{-1}F$ for $F\in E$, $1\le i\le n$ and $u\in\mathcal O_K^\times$, and that $U_{\alpha_j}$ maps $E$ onto $E$ for $1\le j\le n$; then there are smooth characters $\omega_1,\dots,\omega_n$ of $K^\times$ with $\omega_i(u)f_i(u)=1$ for $u\in\mathcal O_K^\times$, such that $\sigma\cong\mathrm{PS}(\omega)$ and $U_{\alpha_j}$ acts on $E$ by $E_j(\omega)$ for $0\le j\le n$.
\end{enumerate}
\end{lemma}
\leanref{Development/AutomorphicArithmetic/AutomorphicVarieties/HeckeAlgebras/OrdinaryHeckeCharacters/ProdZpowConst.lean}{UpperTheory.V.Hi.wtRep}
\leanref{Development/AutomorphicArithmetic/AutomorphicVarieties/HeckeAlgebras/OrdinaryHeckeCharacters/CtrlOrd.lean}{UpperTheory.V.Hi.ctrlOrd}
\leanref{Development/AutomorphicArithmetic/AutomorphicVarieties/ModularityArgument/ThresholdArguments/DenseThresholdCases/AuTeichEx.lean}{UpperTheory.V.Hi.automorphicHeckeSet_unitaryPointProperties_of_typeFor}
\leanref{Development/AutomorphicArithmetic/AutomorphicVarieties/ModularityArgument/ThresholdArguments/DenseThresholdCases/CharOfScalar.lean}{local_step}
\begin{proof}
(1) Let $E$ be the finite set of pairs $e=(v,\tau)$ of a place $v\mid\ell$ and a continuous embedding $\tau:L_{\tilde v}\to\bar{\mathbf Q}_\ell$. Applying Lemma~\ref{sha-hecke-ideal-tw-augmented-hida-inputs-embedding} to the pairs one at a time gives $\mathcal O'\supset\mathcal O_1$ with $\tau(\mathcal O_{L_{\tilde v}})\subset\mathcal O'$ for every $(v,\tau)\in E$, and $\tau$ maps the maximal ideal of $\mathcal O_{L_{\tilde v}}$ into $\mathfrak m'$. Let $\psi_e=\iota\circ\tau$. For each $e$ let $(b_{e,j})_{j\in\kappa_e}$, $k_{0,e}$ and $\mu_{e,j}$ be as in Lemma~\ref{sha-hecke-ideal-tw-augmented-hida-inputs-tableaux} for $\lambda_\tau$ and $\mathcal O'$. Let $\rho$ be the representation $g\mapsto\bigotimes_e\rho_{\lambda_\tau}(\psi_e(g_v))$ on $M=\bigotimes_e\nabla(\lambda_\tau)$, which is algebraic of weight $\lambda$ through $\iota$, with the basis $b_j=\bigotimes_eb_{e,j_e}$ indexed by $j\in\kappa=\prod_e\kappa_e$, and $k_0=(k_{0,e})_e$. The matrix of $\rho(g)$ in this basis has entries $A(g)_{ij}=\prod_eA_e(\psi_e(g_v))_{i_ej_e}$, where $A_e(h)$ is the matrix of $\rho_{\lambda_\tau}(h)$ in the basis $b_e$. These matrices satisfy the hypotheses (1) to (6) of Lemma~\ref{sha-hecke-ideal-tw-augmented-hida-inputs-matrix}.

\emph{Congruences.} Let $\varpi'$ be a uniformizer of $\mathcal O'$, $k\ge0$, and $g\in G_\ell$ with components in $GL_n(\mathcal O_{L_{\tilde v}})$ congruent to $1$ modulo $\varpi^k$. As $\tau$ maps $\varpi^k\mathcal O_{L_{\tilde v}}$ into $\mathfrak m'^k=\varpi'^k\mathcal O'$, $\psi_e(g_v)$ is the image of a matrix over $\mathcal O'$ congruent to $1$ modulo $\varpi'^k$. By Lemma~\ref{sha-hecke-ideal-tw-augmented-hida-inputs-tableaux}(1) with $c=\varpi'^k$, $A_e(\psi_e(g_v))\equiv1$ modulo $\mathfrak m'^k$ for every $e$, and so $A(g)_{ij}\equiv\delta_{ij}$ modulo $\mathfrak m'^k$. For $k=0$ this gives hypothesis (1) of Lemma~\ref{sha-hecke-ideal-tw-augmented-hida-inputs-matrix}, and for the elements of $\bar N_k$ and $T_0^{(k)}$, which are congruent to $1$ modulo $\varpi^{k+1}$, it gives (5) with $c=k$.

\emph{Diagonal elements.} For $t\in T$ with diagonal entries $t_{v,i}$, Lemma~\ref{sha-hecke-ideal-tw-augmented-hida-inputs-tableaux}(3) gives $A(t)_{ij}=\delta_{ij}\prod_e\prod_i\psi_e(t_{v,i})^{\mu_{e,j_e,i}}$. This gives the second half of (3), and, by Lemma~\ref{sha-hecke-ideal-tw-augmented-hida-inputs-tableaux}(4), $A(t)_{k_0k_0}=\iota\big(\prod_{v\mid\ell}\prod_{i=1}^n\prod_\tau\tau(t_{v,i})^{\lambda_{\tau,n+1-i}}\big)$.

\emph{Dominant elements.} Let $\alpha$ be dominant with components $\operatorname{diag}(\varpi^{a_{v,1}},\dots,\varpi^{a_{v,n}})$. Then $A(\alpha)_{k_0k_0}^{-1}A(\alpha)_{ij}=\delta_{ij}\prod_e\psi_e(\varpi)^{m_e}$ with $m_e=\sum_ia_{v,i}(\mu_{e,j_e,i}-\mu_{e,k_{0,e},i})\ge0$, by Lemma~\ref{sha-hecke-ideal-tw-augmented-hida-inputs-tableaux}(5). As $\psi_e(\varpi)\in\iota(\mathfrak m')$, this lies in $\mathcal O'$, which is (2), and in $\mathfrak m'$ if some $m_e>0$. For $\alpha=z$, $a_{v,i}=n-i$; if $j\ne k_0$, some $j_e\ne k_{0,e}$, and $m_e>0$ by Lemma~\ref{sha-hecke-ideal-tw-augmented-hida-inputs-tableaux}(6). This gives (4).

\emph{Unipotent elements.} For $y\in N$, each $\psi_e(y_v)$ is upper unitriangular, so by Lemma~\ref{sha-hecke-ideal-tw-augmented-hida-inputs-tableaux}(2) the entry $(k_{0,e},j_e)$ of $A_e(\psi_e(y_v))$ is $1$ if $j_e=k_{0,e}$ and $0$ otherwise. Taking the product over $e$ gives the first half of (3).

\emph{Central elements.} Let $t\in T$ be central in $G_\ell$, so that each component is a scalar $s_v$. Then $A(t)_{jj}=\prod_e\psi_e(s_v)^{\sum_i\mu_{e,j_e,i}}$. By Lemma~\ref{sha-hecke-ideal-tw-augmented-hida-inputs-tableaux}(5), applied to $a=(1,\dots,1)$ and to $a=(-1,\dots,-1)$, $\sum_i\mu_{e,j,i}=\sum_i\mu_{e,k_{0,e},i}$ for every $j$. So $A(t)=A(t)_{k_0k_0}\cdot1$, which is (6).

Lemma~\ref{sha-hecke-ideal-tw-augmented-hida-inputs-matrix} now gives a weight datum $(\xi,\chi,0)$ over $\mathcal O'$ with $\xi$ isomorphic to $\rho$, hence algebraic of weight $\lambda$ through $\iota$, and $\chi(t)=A(t)_{k_0k_0}$, which is the stated formula.
% (1) is wtRep (ProdZpowConst.lean): the Coef c' by induction over Finset.univ of Emb l S with expChar_integrality_criterion (Lemma sha-hecke-ideal-tw-augmented-hida-inputs-embedding), phi e and hphim; sqBasis (Lemma sha-hecke-ideal-tw-augmented-hida-inputs-tableaux) for each e with c'.O; A g i j the product over e and hA (Basis.piTensorProduct, PiTensorProduct.map_tprod); hK (the uniformizer p of c'.O, sqBasis (1) with c = p^k, prod_delta), hDiag (sqBasis (3)), hPw (sqBasis (5), the uniformizers p v, prod_zpow_const); then ofBasis (Lemma sha-hecke-ideal-tw-augmented-hida-inputs-matrix) with hK at k = 0 (GlAx.Kc_zero, T_le_KO, Kc_le), hPw, the row through sqBasis (2), hDiag, hPw at z with sqBasis (6), hK with Kc_anti, and the central elements through mem_center_iff_val_mem_range_scalar and sqBasis (5) at the constants 1 and -1; the final formula through hDiag, Fintype.prod_sigma and algPow.

(2) Regard $\mathcal O'$ as a subring of $\mathbf C$ through $\iota$, and use the notation of Definition~\ref{sha-hecke-ideal-tw-augmented-hida-inputs-lattices}. Write $\mathcal M'=\mathcal F^{\mathrm{ord}}_{\mathcal O'}(U,\theta')$ and $C=\mathrm{Cl}(U,c,\phi)$.

\emph{Preliminaries.} The character $\theta'$ is continuous: for $k\ge0$, Definition~\ref{sha-hecke-ideal-tw-augmented-hida-weight}(1)(d) gives $c_3$ with $\chi(t)-1\in\mathfrak m'^k$ for $t\in T_0^{(c_3)}$, and for $t\in T_0^{(\max(c,c_3))}$, $\theta'(t)=\chi(t)$. So Lemma~\ref{sha-hecke-ideal-tw-augmented-ordinary}(1) and (2), and Lemma~\ref{sha-hecke-ideal-tw-augmented-sections}, hold for $\mathcal O'$ and $\theta'$, as in the proof of Lemma~\ref{sha-hecke-ideal-tw-augmented-hida-base-change}: $\mathcal M'$ is finite over $\mathcal O'$, and an element of $\mathcal M'$ with values in $\mathfrak m'$ lies in $\mathfrak m'\mathcal M'$. As in the proof of Lemma~\ref{sha-hecke-ideal-tw-augmented-hida-control}, $\operatorname{ctrl}(U'_zF)=U_z\operatorname{ctrl}(F)$ for $F\in V^\xi$; so if $\operatorname{ctrl}(F)=h$ with $h$ a function with values in $\mathcal O'^m$, then $\operatorname{ctrl}(U_z'^nF)=U_z^nh$ for all $n$. The operator $U_z$ preserves functions with values in $\mathfrak m'^k$.

\emph{Coordinates.} Let $x_1,\dots,x_N$ be as in Lemma~\ref{sha-hecke-ideal-tw-augmented-sections}(1) for $c$, $\Sigma_0=\{1,\dots,N\}\times\{1,\dots,m\}\times\{0,\dots,d\}$, and $\operatorname{ev}(F)\in\mathbf C^{\Sigma_0}$ the family of the coordinates of the $F^\flat_i(x_n)$. By Lemma~\ref{sha-hecke-ideal-tw-augmented-hida-inputs-values}, $\operatorname{ev}$ maps $C$ bijectively onto $\mathbf C^{\Sigma_0}$: it is onto by (1), and injective by (2) with $\Lambda=0$, as $F^\flat=0$ forces $F=0$. For $F\in C$ and an $\mathcal O'$-submodule $\Lambda$, $\operatorname{ev}(F)\in\Lambda^{\Sigma_0}$ if and only if $F\in\mathrm{Lat}(\Lambda)$, by Lemma~\ref{sha-hecke-ideal-tw-augmented-hida-inputs-values}(2). By Lemma~\ref{sha-hecke-ideal-tw-augmented-hida-inputs-levels}(2), $U'_z$ maps $C$ into itself, and by Lemma~\ref{sha-hecke-ideal-tw-augmented-hida-inputs-lattice} it maps $C\cap\mathrm{Lat}(\mathcal O')$ into itself. Lemma~\ref{sha-hecke-ideal-tw-augmented-hida-inputs-ordinary-lattice}, applied to $\mathcal O'$, $\Omega=\mathbf C$, $V=V^\xi$, $C$, $T=U'_z$ and $\operatorname{ev}$, gives a free $\mathcal O'$-submodule $S\subset C\cap\mathrm{Lat}(\mathcal O')$ whose $\mathbf C$-span is $\mathrm{Cl}^{\mathrm{ord}}(U,c,\phi)$, such that every $F\in S$ has, for each $k$, some $n$ with $U_z'^{n+1}F-F\in\mathrm{Lat}(\mathfrak m'^k)$, and with the approximation property (3) of that lemma.

\emph{Lifts of $S$.} For $F\in S$ there is $f\in\mathcal M'$ with $f=\operatorname{ctrl}(F)$. Indeed, $f=\operatorname{ctrl}(F)$ lies in $\mathcal F_{\mathcal O'}(U,\theta')$ by Lemma~\ref{sha-hecke-ideal-tw-augmented-hida-inputs-lift}. For $k\ge0$ take $n$ as above; then $U_z^{n+1}f-f=\operatorname{ctrl}(U_z'^{n+1}F-F)$ has values in $\mathfrak m'^k$, so $f$ is ordinary.

\emph{Injectivity on $S$.} Let $F\in S$ with $\operatorname{ctrl}(F)=0$. We show $F\in\mathrm{Lat}(\mathfrak m'^k)$ for every $k$, by induction; for $k=0$ it holds. If $F\in\mathrm{Lat}(\mathfrak m'^k)$, then $U'_zF\in\mathrm{Lat}(\mathfrak m'^{k+1})$ by Lemma~\ref{sha-hecke-ideal-tw-augmented-hida-inputs-lattice}, as $\operatorname{ctrl}(F)=0$, and so $U_z'^{n+1}F\in\mathrm{Lat}(\mathfrak m'^{k+1})$ for every $n$, by the last statement of that lemma. Taking $n$ with $U_z'^{n+1}F-F\in\mathrm{Lat}(\mathfrak m'^{k+1})$ gives $F\in\mathrm{Lat}(\mathfrak m'^{k+1})$. As $\bigcap_k\mathfrak m'^k=0$, $F^\flat=0$ and $F=0$.

\emph{Injectivity.} Let $(s_r)$ be a basis of $S$ over $\mathcal O'$. The functions $\operatorname{ctrl}(s_r)$ have values in $\mathcal O'^m$, and they are linearly independent over $\mathcal O'$, as $\operatorname{ctrl}$ is injective on $S$. So they are linearly independent over $\mathbf C$, as in the proof of Lemma~\ref{sha-hecke-ideal-tw-augmented-hida-control}. As $\mathrm{Cl}^{\mathrm{ord}}(U,c,\phi)$ is spanned by the $s_r$, $\operatorname{ctrl}$ is injective on it.

\emph{The image is in the span.} The image of $\mathrm{Cl}^{\mathrm{ord}}(U,c,\phi)$ is spanned by the $\operatorname{ctrl}(F)$, $F\in S$, which lie in $\mathcal M'$ by the lifts of $S$.

\emph{An approximation.} Let $g\in\mathcal M'$. There are $G\in C\cap\mathrm{Lat}(\mathcal O')$ and a function $h$ with values in $\mathcal O'^m$ such that $\operatorname{ctrl}(G)=h$ and $h-g$ has values in $\mathfrak m'$. Indeed, choose $c_2$ such that $g(x(1,y))-g(x)$ has entries in $\mathfrak m'$ for all $x\in G^{(\ell)}\times G_\ell$ and $y\in\bar N_{c_2}$, $c_1$ such that $\xi(y)-1$ has entries in $\mathfrak m'$ for $y\in\bar N_{c_1}$, by Definition~\ref{sha-hecke-ideal-tw-augmented-hida-weight}(1)(d), and $n$ with $U_z^{n+1}g-g$ with values in $\mathfrak m'$. Let $e=c_1+c_2$, $c'=c+e$ and $M=(n+1)(e+1)\ge e$, and let $x'_1,\dots,x'_{N'}$ be as in Lemma~\ref{sha-hecke-ideal-tw-augmented-sections}(1) for $c'$. As $\phi$ is trivial on $T_0^{(c')}$, Lemma~\ref{sha-hecke-ideal-tw-augmented-hida-inputs-values}(1) gives $F'\in\mathrm{Cl}(U,c',\phi)$ with $F'^\flat_i(x'_r)=(g(x'_r)_i,0,\dots,0)$, and $F'\in\mathrm{Lat}(\mathcal O')$ by (2) of that lemma. By Lemma~\ref{sha-hecke-ideal-tw-augmented-hida-inputs-lift}, $h'=\operatorname{ctrl}(F')$ lies in $\mathcal F_{\mathcal O'}(U,\theta')$, and $h'(x(1,y))-h'(x)$ has entries in $\mathfrak m'$ for all $x$ and $y\in\bar N_{c'}$, as $c'\ge c_1$. So $h'-g\in\mathcal F_{\mathcal O'}(U,\theta')$ has this property too, as $c'\ge c_2$, and vanishes at the $x'_r$. By Lemma~\ref{sha-hecke-ideal-tw-augmented-sections}(3) for $\mathcal O'$ and $c'$, $h'-g$ has values in $\mathfrak m'$. Let $G=U_z'^MF'$ and $h=U_z^Mh'$. Then $G\in C$ by Lemma~\ref{sha-hecke-ideal-tw-augmented-hida-inputs-levels}(3), $G\in\mathrm{Lat}(\mathcal O')$ by Lemma~\ref{sha-hecke-ideal-tw-augmented-hida-inputs-lattice}, $\operatorname{ctrl}(G)=h$, and $h-g=U_z^M(h'-g)+(U_z^Mg-g)$ has values in $\mathfrak m'$, as $n+1$ divides $M$.

\emph{The span is in the image.} Let $L$ be the $\mathcal O'$-submodule of the $g\in\mathcal M'$ that lie in the $\mathbf C$-span of $\operatorname{ctrl}(S)$. Let $g\in\mathcal M'$, and $G$ and $h$ as in the approximation. By Lemma~\ref{sha-hecke-ideal-tw-augmented-hida-inputs-ordinary-lattice}(3) there are $F'\in S$ and $J$ with $U_z'^{j!}G-F'\in\mathrm{Lat}(\mathfrak m')$ for $j\ge J$, as this element lies in $C$. Let $g'\in\mathcal M'$ with $g'=\operatorname{ctrl}(F')$, by the lifts of $S$, so $g'\in L$. Take $n$ with $U_z^{n+1}g-g$ with values in $\mathfrak m'$, and $j=\max(J,n+1)$, so that $n+1$ divides $j!$. Then
\[
g-g'=\big(U_z^{j!}h-g'\big)-U_z^{j!}(h-g)-\big(U_z^{j!}g-g\big),
\]
where $U_z^{j!}h-g'=\operatorname{ctrl}(U_z'^{j!}G-F')$, and all three terms have values in $\mathfrak m'$. So $g-g'\in\mathfrak m'\mathcal M'$, and $\mathcal M'=L+\mathfrak m'\mathcal M'$. By Nakayama's lemma, as $\mathcal M'$ is finite, $L=\mathcal M'$. Hence the span of $\mathcal M'$ is the image of $\mathrm{Cl}^{\mathrm{ord}}(U,c,\phi)$.
% (2) is ctrlOrd (CtrlOrd.lean): hcont, ordStruct over O' (hfin, hdiv), exists_finite_iwahori_doubleCoset_representatives, ev and hbij (exists_automorphicForm_with_prescribed_iwahori_values and mem_lattice_of_values_at_double_coset_reps, Lemma sha-hecke-ideal-tw-augmented-hida-inputs-values), hTC (Lemma sha-hecke-ideal-tw-augmented-hida-inputs-levels), hint (ctrlOrd_stab, Lemma sha-hecke-ideal-tw-augmented-hida-inputs-lattice), exists_free_lattice_spanning_ordinary_part (Lemma sha-hecke-ideal-tw-augmented-hida-inputs-ordinary-lattice), hSper, hSord (exists_integral_form_lifting_controlMap, Lemma sha-hecke-ideal-tw-augmented-hida-inputs-lift, with ctrlOrd_nat from ctrlEq (2) and uOpNat, and ctrlOrd_map); the injectivity is injective_on_span_of_integral_free_submodule (fnIndep) with ctrlOrd_inj (hKrull, hstep through Up_operator_maps_lattice_of_control_values and ctrlOrd_per); the image through Submodule.map_span, hSord, and for the converse L with hle from ctrlOrd_approx (cont_Nm, exists_automorphicForm_with_prescribed_iwahori_values at c', exists_integral_form_lifting_controlMap, mem_valIn_of_invariant_and_values_mem_pow, classicalForms_level_monotone_and_Up_stable (3), hper1), hSapprox, e1, e2 and e3, then le_of_le_smul_of_le_jacobson_bot.

(3) As $\xi$ is algebraic of weight $\lambda$ through $\iota$, fix a linear isomorphism $e:\mathbf C^d\cong\bigotimes_{(v,\tau)}\nabla(\lambda_\tau)$ that carries $\xi$ to the representation of Definition~\ref{sha-hecke-ideal-tw-augmented-parts}(2). For $g\in U(n)(\mathbf A_{L^+})$ let $g^{(\ell)}$ and $g_\ell$ be its components in $G^{(\ell)}$ and $G_\ell$, and $\rho_\infty(g)=\bigotimes_{(v,\tau)}\rho_{\lambda_\tau}(g_{\iota\tau})$, where $g_{\iota\tau}\in GL_n(\mathbf C)$ is the component of $g$ at the infinite place given by the embedding $\iota\tau|_L$. For an automorphic representation $\sigma$ of $U(n)$ (Definition~\ref{gal-definite-unitary}), let $Y(\sigma)$ be the space of the $F\in V^\xi$ with $F(\gamma g)=\xi(\gamma_\ell)F(g)$ for $\gamma\in U(n)(L^+)$, fixed by the elements of $\mathcal K_0K_w$ whose components at the $v_j$ are $1$, that lie in the span of the $F'\in V^\xi$ such that, for every $i$ and every linear form $l$ on $\bigotimes_{(v,\tau)}\nabla(\lambda_\tau)$, the function $g\mapsto l\big(\rho_\infty(g)^{-1}e(F'_i(g^{(\ell)},g_\ell))\big)$ on $U(n)(\mathbf A_{L^+})$ lies in $\sigma$. Two facts about these spaces are used:
\begin{enumerate}
\item[(A)] every $F$ as in (3) lies in the sum of the $Y(\sigma)$ over the automorphic representations $\sigma$ of $U(n)$;
\item[(B)] if $Y(\sigma)\ne0$, then $Y(\sigma)$, with suitable data (a) to (e), is an automorphic part of weight $\xi$.
\end{enumerate}
Given these, let $F$ be as in (3). By (A), $F$ is a finite sum of elements of spaces $Y(\sigma)$; the terms in spaces $Y(\sigma)=0$ vanish, and the others lie in automorphic parts of weight $\xi$ by (B). So $F$ lies in the sum of the automorphic parts of weight $\xi$.

\emph{Proof of (A).} Use the notation of Lemma~\ref{sha-hecke-ideal-tw-augmented-hida-inputs-automorphic-forms}, so that the function in the definition of $Y(\sigma)$ is $\varphi_{F'_i,l}$; note that $\varphi_{F,l}$ is linear in $F$ and in $l$. Let $F$ be as in (3), with components $F_1,\dots,F_m$, fixed by a level $U=\mathcal K_0K$ and invariant under right translation by $I_1(b)$, and let $U_Q$ be the group of the elements of $\mathcal K_0K_w$ whose components at the $v_j$ are $1$. For $u\in G^{(\ell)}_0$, the condition $u\cdot F=F$ reads $\sum_j\iota\tau_0(u)_{ij}F_j(h(u,1))=F_i(h)$ for all $i$ and $h$.

\emph{Open stabilizers.} For each $i$, the $k\in U(n)(\mathbf A^\infty_{L^+})$ with $F_i(hk^\infty)=F_i(h)$ for all $h$ form a subgroup $S_i$. It contains the neighbourhood $W$ of Lemma~\ref{sha-hecke-ideal-tw-augmented-hida-inputs-automorphic-near} for $K$ and $b$, so it is open. Indeed, for $k\in W$, $k^\infty=(u,1)(1,y)$ with $u\in U$, $\tau_0(u)=1$ and $y\in I_1(b)$; so $F_i(h(u,1))=F_i(h)$, as $u\cdot F=F$, and $F_i(h(1,y))=F_i(h)$.

\emph{Projections.} By Lemma~\ref{sha-hecke-ideal-tw-augmented-hida-inputs-automorphic-forms}(2), the $\varphi_{F_i,l}$ are automorphic forms on $U(n)$. They span a finite-dimensional space $M$, as there are finitely many $i$ and the dual of $\bigotimes_{(v,\tau)}\nabla(\lambda_\tau)$ is finite-dimensional. Let $\sigma_1,\dots,\sigma_r$ and $P_1,\dots,P_r$ be as in Lemma~\ref{sha-hecke-ideal-tw-augmented-hida-inputs-automorphic-projections} for $M$. Fix $t$ and $i$. The map $l\mapsto P_t(\varphi_{F_i,l})$ is linear, with values automorphic forms, and for $x\in U(n)(L^+\otimes\mathbf R)$, $P_t(\varphi_{F_i,l})(gx)=P_t(\varphi_{F_i,l\circ\rho_\infty(x)^{-1}})(g)$, by Lemma~\ref{sha-hecke-ideal-tw-augmented-hida-inputs-automorphic-forms}(3), as $x^\infty=1$ and $P_t$ commutes with right translation. By Lemma~\ref{sha-hecke-ideal-tw-augmented-hida-inputs-automorphic-forms}(4) there is $G_{t,i}:G^{(\ell)}\times G_\ell\to\mathbf C^d$ with $G_{t,i}(\gamma h)=\xi(\gamma_\ell)G_{t,i}(h)$ for $\gamma\in U(n)(L^+)$, and $\varphi_{G_{t,i},l}=P_t(\varphi_{F_i,l})\in\sigma_t$ for every $l$. Let $G_t\in V^\xi$ have components $G_{t,1},\dots,G_{t,m}$.

\emph{The sum.} For every $i$ and $l$, $\varphi_{\sum_tG_{t,i}-F_i,l}=\sum_tP_t(\varphi_{F_i,l})-\varphi_{F_i,l}=0$, as $\varphi_{F_i,l}\in M$. So $\sum_tG_t=F$, by Lemma~\ref{sha-hecke-ideal-tw-augmented-hida-inputs-automorphic-forms}(1).

\emph{Invariance.} Let $u\in U_Q$, and $k\in U(n)(\mathbf A^\infty_{L^+})$ with $k^\infty=(u,1)$. As $\rho_\infty(k)=1$, Lemma~\ref{sha-hecke-ideal-tw-augmented-hida-inputs-automorphic-forms}(3) gives $\varphi_{H',l}=R(k)\varphi_{H,l}$ for every function $H:G^{(\ell)}\times G_\ell\to\mathbf C^d$ and $H'(h)=H(h(u,1))$, where $R(k)$ is right translation by $k$. As $P_t$ commutes with $R(k)$,
\[
\varphi_{(u\cdot G_t)_i-G_{t,i},l}=\sum_j\iota\tau_0(u)_{ij}R(k)P_t(\varphi_{F_j,l})-P_t(\varphi_{F_i,l})=P_t\big(\varphi_{(u\cdot F)_i-F_i,l}\big)=0
\]
for every $i$ and $l$. So $u\cdot G_t=G_t$, by Lemma~\ref{sha-hecke-ideal-tw-augmented-hida-inputs-automorphic-forms}(1).

\emph{Conclusion.} Each $G_t$ satisfies $G_t(\gamma h)=\xi(\gamma_\ell)G_t(h)$ for $\gamma\in U(n)(L^+)$, is fixed by $U_Q$, and has $\varphi_{G_{t,i},l}\in\sigma_t$ for all $i$ and $l$. So $G_t\in Y(\sigma_t)$, and $F=\sum_tG_t$ lies in the sum of the $Y(\sigma)$.
% (A) is decompU (DecompU.lean, 208 lines): J = (away, cLoc at the s.w v), formU (Lemma sha-hecke-ideal-tw-augmented-hida-inputs-automorphic-forms), tf = toForm as a bilinear map, hJinf and hRfin, then core2: hopen (the stabilizer subgroup S, open by isOpen_of_mem_nhds with hev = eventually_mem_heckeLevel_and_iwahoriCongruence, hra for the action raF, h3 and h5), k u a preimage of (u, 1) under J, and core (autProj with M the supremum of the ranges of the Phi i, hex through hsur for G, the relations through hinj, hT and hPe, the sum through hPs), then the membership of each G t in YOf (sigma t) through Submodule.subset_span.
% (3) is automorphicHeckeSet_unitaryPointProperties_of_typeFor (AuTeichEx.lean): au_isTeich for a Teichmuller family t, au_open for the openness of ker tau0 (here Lemma sha-hecke-data (1)), decompU for (A), and, for each sigma with YOf sigma /= bot, Z.Desc.P.partU_alt for (B), with Y.Y = YOf. The book's Y(sigma) is YOf Q iota tau0 e xi sigma: Big inter the span of the F with toForm iota e (F i) l in sigma for all i and l (IsTheta.lean), toForm F l g = l (infRep iota lam g^-1 (e (F (away g, cLoc g)))). partU_alt is partX_alt (PartU.lean, 156 lines); its parts that UnitaryPointProperties needs (partUR2Type_alt, partUR2V0, hpair, hcusp) are not used for (B).

\emph{Proof of (B): the local components and the Hecke eigenvalues.} Let $\sigma$ be an automorphic representation of $U(n)$ with $Y(\sigma)\ne0$, and keep the notation of Lemma~\ref{sha-hecke-ideal-tw-augmented-hida-inputs-automorphic-forms}. The condition that $\varphi_{F_i,l}\in\sigma$ for all $i$ and $l$ is linear in $F$, so every $F\in Y(\sigma)$ satisfies it.

\emph{Translations.} Let $w$ be a place of $L$ split over $L^+$ that is one of the $w_j$, one of the $\tilde v$ with $v\mid\ell$, or a place above a place of $L^+$ outside $S'$, and let $U_{n,w}$ and $c_w$ be as in Definition~\ref{gal-definite-unitary}(3). For $h\in GL_n(L_w)$ let $h\cdot F$ be the action of $h$ on $V^\xi$, through $G_w$, $G_\ell$ or $GL_n(L_w)\subset G^{(\ell)}_0$; it is right translation by $h$, as $\tau_0$ involves only the components at the places of $S_{\mathrm a}$. Let $x\in U_{n,w}$ with $c_w(x)=h$. Then $x\in U(n)(\mathbf A^\infty_{L^+})$, $x^\infty$ is $h$ placed at $w$, and $\rho_\infty(x)=1$; so $\varphi_{(h\cdot F)_i,l}=R(x)\varphi_{F_i,l}$, by Lemma~\ref{sha-hecke-ideal-tw-augmented-hida-inputs-automorphic-forms}(3), where $R$ is right translation. For $w=w_j$ or $w=\tilde v$, $Y(\sigma)$ is stable under the action of $GL_n(L_w)$: for $F\in Y(\sigma)$, $h\cdot F$ satisfies $(h\cdot F)(\gamma g)=\xi(\gamma_\ell)(h\cdot F)(g)$; it is fixed by the elements of $\mathcal K_0K_w$ whose components at the $v_j$ are $1$, which commute with $h$; and $\varphi_{(h\cdot F)_i,l}=R(x)\varphi_{F_i,l}\in\sigma$.

\emph{The local components (a).} Let $w$ be $w_j$ or $\tilde v$, and $\pi_w$ a local component of $\sigma$ at $w$, by Lemma~\ref{sha-hecke-ideal-tw-augmented-hida-inputs-automorphic-local}(1). Let $(l_k)$ be a basis of the dual of $\bigotimes_{(v,\tau)}\nabla(\lambda_\tau)$. The maps $F\mapsto\varphi_{F_i,l_k}$, for $1\le i\le m$ and all $k$, are linear maps $Y(\sigma)\to\sigma$ that carry the action of $h\in GL_n(L_w)$ to $R(x)$ for $x\in U_{n,w}$ with $c_w(x)=h$, by the translations; and an $F$ that all of them kill is $0$, by Lemma~\ref{sha-hecke-ideal-tw-augmented-hida-inputs-automorphic-forms}(1). So $Y(\sigma)$ is isomorphic, as a representation of $GL_n(L_w)$, to a subrepresentation of a finite direct sum of copies of $\sigma$, on which $GL_n(L_w)$ acts through $c_w^{-1}$ and $R$. By Lemma~\ref{sha-hecke-ideal-tw-augmented-hida-inputs-automorphic-local}(2), this direct sum is the sum of the images of the equivariant maps from $\pi_w$, and so is $Y(\sigma)$, by Lemma~\ref{sha-hecke-ideal-tw-augmented-isotypic}(1). Let $\sigma_{Y,j}$ and $\sigma_{Y,v}$ be these local components at $w_j$ and at $\tilde v$.

\emph{Unramified places.} The representation $\sigma$ is unramified at every place $w$ of $L$ above a place $v$ of $L^+$ outside $S'$. Indeed, let $F\in Y(\sigma)$ be nonzero, and $i$ and $l$ with $\varphi_{F_i,l}\ne0$, by Lemma~\ref{sha-hecke-ideal-tw-augmented-hida-inputs-automorphic-forms}(1). Let $x\in U_{n,w}$ with $c_w(x)\in GL_n(\mathcal O_{L_w})$. Then $x\in U(n)(\mathbf A^\infty_{L^+})$, $\rho_\infty(x)=1$, and $x^\infty=(u,1)$, where $u\in G^{(\ell)}_0$ has components $1$ at the places of $L^+$ other than $v$ and integral components at the places of $L$ above $v$. So $u$ lies in $\mathcal K_0K_w$, with components $1$ at the $v_j$, and $\tau_0(u)=1$; as $u\cdot F=F$, $F_i(h(u,1))=F_i(h)$ for all $h$, and $R(x)\varphi_{F_i,l}=\varphi_{F_i,l}$ by Lemma~\ref{sha-hecke-ideal-tw-augmented-hida-inputs-automorphic-forms}(3). So $\varphi_{F_i,l}$ is a nonzero element of $\sigma$ fixed by $c_w^{-1}(GL_n(\mathcal O_{L_w}))$.

\emph{The Hecke eigenvalues (b).} Let $w$ be a split place of $L$ above a place of $L^+$ outside $S'$, of norm $q_w$, and $K_w=GL_n(\mathcal O_{L_w})$. Every $F\in Y(\sigma)$ is fixed by $K_w$, which lies in $\mathcal K_0K_w$ and has components $1$ at the $v_j$. As $\sigma$ is unramified at $w$, Lemma~\ref{sha-hecke-ideal-tw-augmented-hida-inputs-automorphic-local}(1) and (3) show that $\sigma$ has a spherical local component $\pi_w$ at $w$, with Satake parameters $B_w$. Let $a_{Y,w,k}=q_w^{k(n-k)/2}e_k(B_w)$ for $0\le k\le n$; then $a_{Y,w,0}=1$.

Let $1\le k\le n$, $F\in Y(\sigma)$, $1\le i\le m$ and $l$ a linear form. Let $h$ run through representatives of $K_w\alpha_kK_w/K_w$, so that $T^\xi_{w,k}F=\sum_hh\cdot F$, and let $x_h\in U_{n,w}$ with $c_w(x_h)=h$. By the translations, $\varphi_{(T^\xi_{w,k}F)_i,l}=\sum_hR(x_h)\varphi$, with $\varphi=\varphi_{F_i,l}$. The function $\varphi\in\sigma$ is fixed by $c_w^{-1}(K_w)$, as $F$ is fixed by $K_w$. By Lemma~\ref{sha-hecke-ideal-tw-augmented-hida-inputs-automorphic-local}(2), $\varphi=\sum_rf_r(u_r)$ for finitely many equivariant maps $f_r:\pi_w\to\sigma$ and vectors $u_r$ of $\pi_w$. For a vector $u$ of $\pi_w$, fixed by an open subgroup $K'\subset K_w$, of finite index as $K_w$ is compact, let $\operatorname{av}(u)=[K_w:K']^{-1}\sum_{k'}\pi_w(k')u$, the sum over representatives of $K_w/K'$; it is fixed by $K_w$, and $f_r(\operatorname{av}(u))$ is the same average of the $R(x)f_r(u)$ over $c_w^{-1}(K_w)$. As $\varphi$ is fixed by $c_w^{-1}(K_w)$, $\varphi=\sum_rf_r(\operatorname{av}(u_r))$. As $f_r$ is equivariant, $\sum_hR(x_h)f_r(\operatorname{av}(u_r))=f_r\big(\sum_h\pi_w(h)\operatorname{av}(u_r)\big)=a_{Y,w,k}f_r(\operatorname{av}(u_r))$, by Lemma~\ref{sha-hecke-ideal-tw-augmented-hida-inputs-automorphic-local}(3). So $\varphi_{(T^\xi_{w,k}F-a_{Y,w,k}F)_i,l}=0$ for all $i$ and $l$, and $T^\xi_{w,k}F=a_{Y,w,k}F$, by Lemma~\ref{sha-hecke-ideal-tw-augmented-hida-inputs-automorphic-forms}(1).
% (a) and (b) follow partX_alt: loc w from GlobalTheory.DefU.exists_locComp_split with hfix (unramified gives a fixed vector); isoL = partUR2IsoL and isoQ = partUR2IsoQ (partUR2Mem for the linearity, partUR2AtV and partUR2AtW for the translations, with sec and toKT_sec for tau0 = 1, frame_stab and big_atW for the stability, isoSum_restriction_of_local_component_unitary, then isoSum_of_injective_equivariant_family with the family formL against the dual basis and formU (1)); a from satake_local_of_printed_results and satT (a 0 = 1, sphOp acting by a j on the KO-fixed vectors); sph = partUR2Sph (W_le of HSet.Ax for the KO-invariance, push through partUR2Abs, the averaging of RepScalar.repScalar); partUR2Unr through partUR2Fix for the unramifiedness.

\emph{The Galois representation (c).} Let $\rho_\sigma:\operatorname{Gal}(\bar L/L)\to GL_n(\bar{\mathbf Q}_\ell)$ be the representation of Lemma~\ref{sha-hecke-ideal-unitary} for $\sigma$, whose hypotheses hold in Situation~\ref{sha-hecke-situation}. The places of $L$ above places of $L^+$ outside $S'$ do not lie above $\ell$, as the places above $\ell$ lie in $S$; $\sigma$ is unramified at them, by the unramified places; so $\rho_\sigma$ is unramified at them, by Lemma~\ref{sha-hecke-ideal-unitary}(1). As $L/L^+$ is unramified at the finite places, the restriction map $\operatorname{Gal}(\bar L/L)\to\Gamma'_0$ is onto, and its kernel is the closed normal subgroup generated by the inertia groups at these places. So $\rho_\sigma=\rho_Y\circ p'$, with $p'$ this map, for a homomorphism $\rho_Y:\Gamma'_0\to GL_n(\bar{\mathbf Q}_\ell)$, continuous as $p'$ is a continuous surjection of compact groups, and semisimple as $\rho_\sigma$ is and $p'$ is onto. Let $w$ be a split place of $L$ above a place of $L^+$ outside $S'$, and $\pi_w$ and $B_w$ as in (b). By Lemma~\ref{sha-hecke-ideal-unitary}(3), the eigenvalues of $\rho_\sigma$ at the inverse of a Frobenius element at $w$ are, through $\iota$, the $q_w^{(n-1)/2}\alpha$ for $\alpha\in B_w$. So its characteristic polynomial, carried by $\iota$, is
\[
\prod_{\alpha\in B_w}\big(X-q_w^{(n-1)/2}\alpha\big)=\sum_{k=0}^n(-1)^kq_w^{k(n-1)/2}e_k(B_w)X^{n-k}=\sum_{k=0}^n(-1)^kq_w^{k(k-1)/2}a_{Y,w,k}X^{n-k},
\]
as $k(n-1)/2=k(k-1)/2+k(n-k)/2$; this is the Hecke polynomial with coefficients $a_{Y,w,k}$.

\emph{The places $w_j$ (d).} Let $1\le j\le q$, $A$ a valuation ring of $\bar L$ over $w_j$, and $\operatorname{rc}$ a map that satisfies the laws. The place $v_j$ lies in $S'\setminus S$, so $w_j$ is split over $L^+$ and does not lie above $\ell$, and $\rho_Y|_{W_A}$, viewed through $p'$, is $\rho_\sigma|_{W_A}$. The representation $\sigma_{Y,j}$ is a local component of $\sigma$ at $w_j$. So by Lemma~\ref{sha-hecke-ideal-unitary}(2), $\rho_\sigma|_{W_A}$ has a Weil-Deligne representation, and $\mathrm{WD}_\iota(\rho_\sigma|_{W_A})=\operatorname{rc}(\sigma_{Y,j}\nu^{(1-n)/2})$ when the latter is abelian and monodromy-free.
% (c) and (d) follow partX_alt: R from Z.Win.Gu.galUnitaryT_of_printsH (Lemma sha-hecke-ideal-unitary), hRu with partUR2Unr (unr_sum, partUR2Fix, mem_maximalCompact_of_under_eq for the places above v other than w), exists_descent_to_level_group_of_unramified (ker_pr_le for the kernel of V.pr, liftOfSurjective, the quotient map V.pr for the continuity), ss from hRs through pr surjective, frob from satT (its characteristic polynomial identity satT_pol) with satFrobT_alt (Lemma sha-hecke-ideal-unitary (3)) and lgAwayT; atQ = lgAwayT at the w_q, with HSet.Ax.Q_away for w_q not above l: the first clause of CmpT in the field LocGlobT of galUnitaryT_of_printsH (hLG), whose WDT gives both the Weil-Deligne representation (IsWDOf) and its equality with rc(algNorm sigma_q) under Gd, which is Lemma sha-hecke-ideal-unitary (2) away from l.

\emph{The places above $\ell$ (e).} Let $v\mid\ell$, and let $A$ be the valuation ring of $\bar L$ at the place above $\tilde v$ that defines $D_v$, so that $\rho_Y|_{D_v}$, viewed through $p'$, is $\rho_\sigma|_{D_A}$. The place $\tilde v$ is split over $L^+$, as $v\in S$, and $\sigma_{Y,v}$ is a local component of $\sigma$ at $\tilde v$. Let $\omega_1,\dots,\omega_n$ be smooth characters of $L_{\tilde v}^\times$, pairwise distinct on $\mathcal O_{L_{\tilde v}}^\times$, with $\sigma_{Y,v}\cong\mathrm{PS}(\omega)$, let $\operatorname{rc}$ satisfy the laws, and let $y=\operatorname{rc}(\sigma_{Y,v}\nu^{(1-n)/2})$.

\emph{Plainness.} The representation $\sigma_{Y,v}$ is plain (Definition~\ref{gal-plain}). Indeed, $\mathrm{PS}(\omega)$ is the normalized parabolic induction of the characters $\omega_i\nu^{(n-1)/2}$ (Definition~\ref{sha-hecke-local-ordinary}(3)), so these characters form the supercuspidal support of $\sigma_{Y,v}$ (Definition~\ref{gal-irreducible-classes}(3)). No two of them differ by $\nu$: an equality $\omega_i\nu^{(n-1)/2}=\omega_j\nu^{(n+1)/2}$ is impossible for $i=j$, as $\nu(\varpi)=|\varpi|\ne1$ for a uniformizer $\varpi$, and for $i\ne j$, as $\omega_i$ and $\omega_j$ differ at some unit, at which $\nu$ is $1$.

\emph{The parameter.} By Lemma~\ref{sha-hecke-local-rec}(2), applied with a local Artin map (Proposition~\ref{gal-local-artin-unique}), $y$ is abelian and monodromy-free.

\emph{De Rham and the Hodge-Tate multisets.} By Lemma~\ref{sha-hecke-ideal-unitary}(5) at $\tilde v$, as $\sigma_{Y,v}$ is plain, $\rho_\sigma|_{D_A}$ is de Rham, and $\mathrm{HT}_\tau(\rho_\sigma|_{D_A})=\mathrm{HT}_\tau$ for every embedding $\tau:L\to\bar{\mathbf Q}_\ell$ that extends continuously to $L_{\tilde v}$. Let $\tau$ be such an embedding, the restriction of a continuous $\tau':L_{\tilde v}\to\bar{\mathbf Q}_\ell$. As $Y(\sigma)\ne0$, there are $F\in Y(\sigma)$ and an index $i_0$ with $F_{i_0}\ne0$, and $\varphi_{F_{i_0},l}\in\sigma$ for every $l$. By Lemma~\ref{sha-hecke-ideal-tw-augmented-hida-inputs-automorphic-weight} for the pair $(v,\tau')$, $\sigma$ has weight $\mu$ at $\iota\circ\tau$, with $\mu_k=-\lambda_{\tau,n-k}$ for $0\le k<n$. By Lemma~\ref{sha-hecke-ideal-unitary}(4),
\[
\mathrm{HT}_\tau=\{k+\lambda_{\tau,n-k}:0\le k<n\}=\{\lambda_{\tau,i}+n-i:1\le i\le n\}.
\]

\emph{The Weil-Deligne representation.} By Lemma~\ref{sha-hecke-ideal-unitary}(2) at $\tilde v$, as $y$ is abelian and monodromy-free, $\rho_\sigma|_{D_A}$ is potentially semistable and $\mathrm{WD}_\iota(\rho_\sigma|_{D_A})=y$. So the characteristic polynomials of the Weil-Deligne representation of $\rho_Y|_{D_v}$ are those of $y$.

\emph{Conclusion.} With the data (a) to (e), $Y(\sigma)$ is an automorphic part of weight $\xi$: it is nonzero, and by its definition its elements $F$ satisfy $F(\gamma g)=\xi(\gamma_\ell)F(g)$ for $\gamma\in U(n)(L^+)$ and are fixed by the elements of $\mathcal K_0K_w$ whose components at the $v_j$ are $1$.
% (e) follows partX_alt: atL = ordinary_at_l_of_local_global_hodge_tate_weights (MkAtLT.lean), which is Z.Ord.at_l_of_literal_automorphy_n (OrdAtLOf.lean; its hypothesis LiteratureRankNHypothesis is not used), applied with hHG (HtGT, the book's Lemma sha-hecke-ideal-unitary (5)), hLG (LocGlobT at the s.w v, (2) of that lemma above l) and hW ((4) of that lemma) at the weights hwt = partUR2Wt (Lemma sha-hecke-ideal-tw-augmented-hida-inputs-automorphic-weight); its argument hall is not used. at_l_of_literal_automorphy_n: the private plainTor_of_ps (hasSupport_of_isSubquotient_prinSeries for the support, the characters omega_i absPow((n-1)/2), unlinked through absPow_ofReal_eq_one and isUnram_absPow), Z.Win.Gu.plainGd for Gd (Tor and mono0) of rc(algNorm sigma), which the book cites as Lemma sha-hecke-local-rec (2), hHT for IsdRAt and IsJumpsAt, the private ex_emb for the extension of tau to L_w, the multiset identity e from hw (infOfWeight of DomWeight.dual, htW lam j tau = lambda_{tau,n-j} + j for j in Fin n), and the second clause of CmpT for IsPstAt and wdAt = rc(algNorm sigma). The hypothesis AbW of AtLT is not used. The conclusion uses Y.le = inf_le_left (YOf is below Big).

(4) Write $V_\sigma$ for the space of $\sigma$, and form the operators $U_\alpha$ of Definition~\ref{sha-hecke-local-ordinary}(1) for $\pi$ on the vectors of $M$ fixed by $N_0$ and for $\sigma$ on those of $V_\sigma$. Call a linear map $p:M\to V_\sigma$ \emph{equivariant on $Y$} if $p(\pi(g)y)=\sigma(g)p(y)$ for all $g\in GL_n(K)$ and $y\in Y$. By Lemma~\ref{sha-hecke-ideal-tw-augmented-hida-inputs-separation}, $Y$ is stable under $\pi$.

\emph{The character $\phi$.} Every $t\in T_0$ is a product $e_1(u_1)\cdots e_n(u_n)$ with $u_i\in\mathcal O_K^\times$, so it acts on $E$ by a scalar, the product of the $f_i(u_i)^{-1}$. As $E\ne0$, this scalar $a(t)$ is determined by $t$; so $a(st)=a(s)a(t)$ and $a(1)=1$, comparing the actions on a nonzero element of $E$. Hence $\phi(t)=a(t)^{-1}$ is a character of $T_0$, every $t\in T_0$ acts on $E$ by $\phi(t)^{-1}$, and $\phi(e_i(u))=f_i(u)$ for all $i$ and $u$. So the characters $u\mapsto\phi(e_i(u))$ of $\mathcal O_K^\times$ are pairwise distinct, and $\phi$ is trivial on the matrices of $T_0$ congruent to $1$ modulo $\varpi^{c+1}$, which are the products of the $e_i(u_i)$ with $u_i\equiv1$ modulo $\varpi^{c+1}$.

\emph{$U_z$ maps $E$ onto $E$.} By Lemma~\ref{sha-hecke-ideal-tw-augmented-uop}(1), $U_\alpha U_\beta=U_{\alpha\beta}$ on the vectors fixed by $N_0$, and $U_{\alpha_0}=U_1$ is the identity on them. If $U_\alpha$ and $U_\beta$ map $E$ onto $E$, so does $U_{\alpha\beta}$: for $F\in E$ write $F=U_\alpha F_1$ and $F_1=U_\beta F_2$ with $F_1,F_2\in E$; then $U_{\alpha\beta}F_2=U_\alpha U_\beta F_2=F$, as $F_2$ is fixed by $N_0$. As $z=\alpha_0\alpha_1\cdots\alpha_{n-1}$ and each $U_{\alpha_j}$ with $1\le j\le n$ maps $E$ onto $E$, so does $U_z$.

\emph{Into the ordinary part.} Let $p$ be equivariant on $Y$ and $F\in E$. Then $p(F)$ lies in the ordinary part $V_\sigma^{\mathrm{ord}}(c,\phi)$ of Definition~\ref{sha-hecke-local-ordinary}(2). Indeed, $p(F)$ is fixed by $N_0$ and $\bar N_c$, and every $t\in T_0$ acts on it by $\phi(t)^{-1}$, as $F$ is and $p$ is equivariant on $Y$; so $p(F)\in V_\sigma(c,\phi)$, and the same holds for every element of $E$. For $k\ge0$, the previous step gives $F_k\in E$ with $F=U_z^kF_k$. The operator $U_z$ is a finite sum of translations $\pi(xz)$, and $F_k$ and its images under them lie in $Y$, so $p(U_z^kF_k)=U_z^kp(F_k)$. Hence $p(F)\in U_z^kV_\sigma(c,\phi)$ for every $k$.

\emph{The representation $\sigma$.} Let $F_0\in E$ be nonzero. By Lemma~\ref{sha-hecke-ideal-tw-augmented-hida-inputs-separation} there is $p_0$ equivariant on $Y$ with $p_0(F_0)\ne0$, and $p_0(F_0)\in V_\sigma^{\mathrm{ord}}(c,\phi)$ by the previous step. So $V_\sigma^{\mathrm{ord}}(c,\phi)\ne0$. By Lemma~\ref{sha-hecke-local-principal}, applied to $\sigma$, $c$ and $\phi$, there are smooth characters $\omega_1,\dots,\omega_n$ of $K^\times$ with $\omega_i(u)=\phi(e_i(u))^{-1}$ for $u\in\mathcal O_K^\times$, such that $\sigma\cong\mathrm{PS}(\omega)$ and $U_{\alpha_j}$ acts on $V_\sigma^{\mathrm{ord}}(c,\phi)$ by $E_j(\omega)$ for $0\le j\le n$. Then $\omega_i(u)f_i(u)=\phi(e_i(u))^{-1}\phi(e_i(u))=1$.

\emph{The eigenvalues.} Let $F\in E$, $0\le j\le n$ and $D=U_{\alpha_j}F-E_j(\omega)F$. Then $D\in Y$, as $Y$ is stable under $\pi$ and $U_{\alpha_j}$ is a finite sum of translations. Suppose $D\ne0$. By Lemma~\ref{sha-hecke-ideal-tw-augmented-hida-inputs-separation} there is $p$ equivariant on $Y$ with $p(D)\ne0$. But $p(U_{\alpha_j}F)=U_{\alpha_j}p(F)=E_j(\omega)p(F)$, as $p(F)\in V_\sigma^{\mathrm{ord}}(c,\phi)$; so $p(D)=0$, a contradiction. Hence $U_{\alpha_j}$ acts on $E$ by $E_j(\omega)$.
% (4) is local_step (CharOfScalar.lean): glBridge for hz (z is the product of the al j, j < n) and hgen (T0 generated by the eT i y); char_of_scalar (the scalar a t, multiplicative by smul_left_injective at F0) and hval (phi (eT i y) = f i y) build the character phi of T0 from the scalars f i y, through the product form of hgen; the Lean's f takes values in C with 0^-1 = 0, and a value 0 is excluded as eT i y acts invertibly and E contains F0 /= 0; ordB's hypotheses on phi are stated on the eT i y (htriv on Uc (e + 1), hdis); hzo through onto_prod (uOpMul, al_zero); hinv through into_inv (Iwa.eig of sigma, uOpNat); isoSep (Lemma sha-hecke-ideal-tw-augmented-hida-inputs-separation) for p0; exists_principal_series_of_distinct_inv_ne_bot (Lemma sha-hecke-local-principal) with h2, h3, h4; the eigenvalues by contradiction with isoSep, uOp_stable and uOpNat.
% goodScalB_alt applies local_step at each v with E = Y meet OrdAt. The parts of (3) also satisfy UnitaryPointProperties, not used here.
\end{proof}

\begin{lemma}\label{sha-hecke-ideal-tw-augmented-hida-control}
In Situation~\ref{sha-hecke-ideal-tw-augmented-setup}, let $\mathcal U$ be a set of subgroups of $G^{(\ell)}_0$ of finite level that satisfy the freeness condition of Lemma~\ref{sha-hecke-ideal-tw-augmented-sections}, such that any two members of $\mathcal U$ contain a third. Let $x:\Lambda\to\mathcal O'$ be an $\mathcal O$-algebra map with $x(\mathfrak m_\Lambda)\subset\mathfrak m'$, $(\xi,\chi,c_0)$ a weight datum over $\mathcal O'$ (Definition~\ref{sha-hecke-ideal-tw-augmented-hida-weight}), $c\ge c_0$, and $\phi$ a character of $T_0$ with values in $\mathcal O'^\times$, trivial on $T_0^{(c)}$, such that $x(\theta(t))=\phi(t)\chi(t)$ for $t\in T_0$. Then:
\begin{enumerate}
\item for $U\in\mathcal U$ and $e\ge c$, $\mathrm{Cl}(U,e,\phi)$ is finite-dimensional; for dominant $\alpha$, $U'_\alpha$ maps it into itself, with eigenvalues integral over $\mathcal O'$; and for $u_1,\dots,u_r\in G^{(\ell)}_0$ the eigenvalues of $F\mapsto\sum_qu_q\cdot F$ on its nonzero elements are integral over $\mathcal O'$;
\item there is a map $\operatorname{sp}:\mathcal F\to V^\xi$ such that for all $U,U'\in\mathcal U$: $\operatorname{sp}$ is additive on $\mathcal F^{\mathrm{ord}}(U)$ and $\operatorname{sp}(bf)=\iota x(b)\operatorname{sp}(f)$ for $b\in\Lambda$ and $f\in\mathcal F^{\mathrm{ord}}(U)$; the span of $\operatorname{sp}(\mathcal F^{\mathrm{ord}}(U))$ is $\mathrm{Cl}^{\mathrm{ord}}(U,e,\phi)$ for every $e\ge c$; the images under $\operatorname{sp}$ of a basis of $\mathcal F^{\mathrm{ord}}(U)$ over $\Lambda$ are linearly independent over $\mathbf C$; $\operatorname{sp}(U_\alpha f)=U'_\alpha\operatorname{sp}(f)$ for dominant $\alpha$ and $f\in\mathcal F^{\mathrm{ord}}(U)$; and if $u_1,\dots,u_r\in G^{(\ell)}_0$ are such that $F\mapsto\sum_qu_q\cdot F$ maps the elements of $V^\xi$ fixed by $U$ to elements fixed by $U'$, then $\operatorname{sp}(\sum_qu_q\cdot f)=\sum_qu_q\cdot\operatorname{sp}(f)$ for $f\in\mathcal F^{\mathrm{ord}}(U)$.
\end{enumerate}
\end{lemma}
\leanref{Development/AutomorphicArithmetic/AutomorphicVarieties/HeckeAlgebras/HeckeFixedLevelSpaces/FreeMap.lean}{UpperTheory.V.Hi.Sp.absSpec}
\begin{proof}
Let $\theta'=x\circ\theta$, a continuous character of $T_0$ with values in $\mathcal O'^\times$. Applying $x$ to $\theta(b)\tau_0(u)=1$ shows that every $U\in\mathcal U$ satisfies the freeness condition for $\theta'$. Write $\iota x\circ f$ for the function $g\mapsto\iota x(f(g))$.

(1) This is Lemma~\ref{sha-hecke-ideal-tw-augmented-hida-classical} for $U$ and $e$, as $\phi$ is trivial on $T_0^{(e)}\subset T_0^{(c)}$.

\emph{Monotonicity.} If $U_2\subset U_1$ and $e_1\le e_2$, then $\mathrm{Cl}(U_1,e_1,\phi)\subset\mathrm{Cl}(U_2,e_2,\phi)$, as $\bar N_{e_2}\subset\bar N_{e_1}$, and so $\mathrm{Cl}^{\mathrm{ord}}(U_1,e_1,\phi)\subset\mathrm{Cl}^{\mathrm{ord}}(U_2,e_2,\phi)$.

\emph{The map.} Call $F\in V^\xi$ a \emph{classical specialization} of $f\in\mathcal F$ if $F\in\mathrm{Cl}^{\mathrm{ord}}(U,e,\phi)$ for some $U\in\mathcal U$ and $e\ge c$, and $\operatorname{ctrl}(F)=\iota x\circ f$. There is at most one. Indeed, let $F_1$ and $F_2$ be classical specializations of $f$, for $(U_1,e_1)$ and $(U_2,e_2)$, and let $U_3\in\mathcal U$ lie in $U_1$ and $U_2$. By the monotonicity, $F_1-F_2\in\mathrm{Cl}^{\mathrm{ord}}(U_3,\max(e_1,e_2),\phi)$, and $\operatorname{ctrl}(F_1-F_2)=0$; so $F_1=F_2$ by Lemma~\ref{sha-hecke-ideal-tw-augmented-hida-inputs}(2). For $U\in\mathcal U$, $e\ge c$ and $f\in\mathcal F^{\mathrm{ord}}(U)$ there is one in $\mathrm{Cl}^{\mathrm{ord}}(U,e,\phi)$: $x\circ f\in\mathcal F^{\mathrm{ord}}_{\mathcal O'}(U,\theta')$ by Lemma~\ref{sha-hecke-ideal-tw-augmented-hida-base-change}, so $\iota x\circ f$ is the control map of an element of $\mathrm{Cl}^{\mathrm{ord}}(U,e,\phi)$ by Lemma~\ref{sha-hecke-ideal-tw-augmented-hida-inputs}(2). Let $\operatorname{sp}(f)$ be the classical specialization of $f$ if there is one, and $0$ otherwise.

\emph{Linearity.} For $f,f'\in\mathcal F^{\mathrm{ord}}(U)$ and $b\in\Lambda$, $\operatorname{sp}(f)+\operatorname{sp}(f')$ and $\iota x(b)\operatorname{sp}(f)$ lie in $\mathrm{Cl}^{\mathrm{ord}}(U,c,\phi)$, with control maps $\iota x\circ(f+f')$ and $\iota x\circ(bf)$. So they are $\operatorname{sp}(f+f')$ and $\operatorname{sp}(bf)$.

\emph{The span.} Let $e\ge c$. The span of $\operatorname{sp}(\mathcal F^{\mathrm{ord}}(U))$ lies in $\mathrm{Cl}^{\mathrm{ord}}(U,e,\phi)$. Conversely, let $F\in\mathrm{Cl}^{\mathrm{ord}}(U,e,\phi)$. By Lemma~\ref{sha-hecke-ideal-tw-augmented-hida-inputs}(2), $\operatorname{ctrl}(F)$ is a $\mathbf C$-linear combination of elements of $\mathcal F^{\mathrm{ord}}_{\mathcal O'}(U,\theta')$, which by Lemma~\ref{sha-hecke-ideal-tw-augmented-hida-base-change} are $\mathcal O'$-linear combinations of the $x\circ f$, $f\in\mathcal F^{\mathrm{ord}}(U)$. As $\operatorname{ctrl}(\operatorname{sp}(f))=\iota x\circ f$, $\operatorname{ctrl}(F)=\operatorname{ctrl}(F_0)$ for some $F_0$ in the span of $\operatorname{sp}(\mathcal F^{\mathrm{ord}}(U))$. Then $F-F_0\in\mathrm{Cl}^{\mathrm{ord}}(U,e,\phi)$ has control map $0$, so $F=F_0$.

\emph{Independence.} Let $(b_k)$ be a basis of $\mathcal F^{\mathrm{ord}}(U)$ over $\Lambda$. By Lemma~\ref{sha-hecke-ideal-tw-augmented-hida-base-change}, the $x\circ b_k$ form a basis of $\mathcal F^{\mathrm{ord}}_{\mathcal O'}(U,\theta')$ over $\mathcal O'$, so they are linearly independent over $\mathcal O'$ as functions. Functions with values in $\mathcal O'$ that are linearly independent over $\mathcal O'$ are linearly independent over the fraction field $F'$ of $\mathcal O'$, by clearing denominators, and then over $\mathbf C$, by writing the coefficients in a basis of $\mathbf C$ over $F'$. As $\operatorname{ctrl}(\operatorname{sp}(b_k))=\iota x\circ b_k$, the $\operatorname{sp}(b_k)$ are linearly independent over $\mathbf C$.

\emph{The operators $U_\alpha$.} By (1), $U'_\alpha$ maps $\mathrm{Cl}(U,c,\phi)$ into itself. On vectors fixed by $N_0$ it commutes with $U'_z$, by Lemma~\ref{sha-hecke-ideal-tw-augmented-uop}(1), as $\alpha z=z\alpha$; so it maps every generalized eigenspace of $U'_z$ in $\mathrm{Cl}(U,c,\phi)$ into itself, and $\mathrm{Cl}^{\mathrm{ord}}(U,c,\phi)$ into itself. For $y\in N$ and $t\in T$ the row $0$ of $\xi(yt)$ is $(\chi(t),0,\dots,0)$, by Definition~\ref{sha-hecke-ideal-tw-augmented-hida-weight}(1)(b). So for each representative $y$ of $N_0/\alpha N_0\alpha^{-1}$ the coordinate $0$ of $\xi(g_\ell)^{-1}F_i(g(1,y\alpha))=\xi(y\alpha)\,\xi(g_\ell y\alpha)^{-1}F_i(g(1,y\alpha))$ is $\chi(\alpha)\operatorname{ctrl}(F)_i(g(1,y\alpha))$, and $\operatorname{ctrl}(U'_\alpha F)=U_\alpha\operatorname{ctrl}(F)$. Hence for $f\in\mathcal F^{\mathrm{ord}}(U)$, $U'_\alpha\operatorname{sp}(f)\in\mathrm{Cl}^{\mathrm{ord}}(U,c,\phi)$ has control map $U_\alpha(\iota x\circ f)=\iota x\circ U_\alpha f$, and it is $\operatorname{sp}(U_\alpha f)$.

\emph{Sums of translations.} Let $S=\sum_qu_q\cdot$, which maps the elements of $V^\xi$ fixed by $U$ to elements fixed by $U'$. It commutes with right translation by $G_\ell$, as the $u_q$ act on the variable in $G^{(\ell)}$, and it preserves the condition $F(\gamma g)=\xi(\gamma_\ell)F(g)$. So it maps $\mathrm{Cl}(U,c,\phi)$ into $\mathrm{Cl}(U',c,\phi)$ and commutes with $U'_z$, hence maps $\mathrm{Cl}^{\mathrm{ord}}(U,c,\phi)$ into $\mathrm{Cl}^{\mathrm{ord}}(U',c,\phi)$. As $\xi(g_\ell)$ does not change under $g\mapsto g(u_q,1)$, $\operatorname{ctrl}(u_q\cdot F)=u_q\cdot\operatorname{ctrl}(F)$. So for $f\in\mathcal F^{\mathrm{ord}}(U)$, $S\operatorname{sp}(f)\in\mathrm{Cl}^{\mathrm{ord}}(U',c,\phi)$ has control map $S(\iota x\circ f)=\iota x\circ Sf$, and it is $\operatorname{sp}(Sf)$.
% absSpec (FreeMap.lean, 589 lines): free_map; hCF from clFin (Lemma sha-hecke-ideal-tw-augmented-hida-classical) and hecke_eigenvalue_isIntegral; hCO from ctrlOrd (Lemma sha-hecke-ideal-tw-augmented-hida-inputs (2)) and hBC from ordBC (Lemma sha-hecke-ideal-tw-augmented-hida-base-change); hClOmono through cl_mono; the predicate Good, huniq, hex, sp through Classical.choose, hsp and hspec; add and smul by hsp; span through hsub (IsBaseChange.inductionOn); indep through hbc.basis and fnIndep (FnIndep.lean); the U operators through UF_maps_clOrd_of_maps_cl (UF_comm_on_N_fixed), ctrl_UF_eq_U_ctrl (wtRep_inv_mulVec_zero_eq_chi_inv) and mapFn_of_mem_P (uOpNat); the sums through sum_raF_maps_clOrd_of_maps_fixd (raF_mem_pi_AF, sum_raF_rpF_comm, uOpNat), ctrl_raF_comm and mapFn_ra. The Lean's levels are a family U : Lambda -> Subgroup with hdir.
\end{proof}

\begin{lemma}\label{sha-hecke-ideal-tw-augmented-hida-decomposition}
In Situation~\ref{sha-hecke-ideal-tw-augmented-setup}, let $\lambda$ be a weight, $(\xi,\chi,c_0)$ a weight datum over $\mathcal O'$ (Definition~\ref{sha-hecke-ideal-tw-augmented-hida-weight}) with $\xi$ algebraic of weight $\lambda$ through $\iota$, $c\ge c_0$, $\phi$ a character of $T_0$ with values in $\mathcal O'^\times$, trivial on $T_0^{(c)}$, and $U=\mathcal K_0K$ a level. Then $\mathrm{Cl}(U,c,\phi)$ is the sum of its intersections with the automorphic parts of weight $\xi$ (Definition~\ref{sha-hecke-ideal-tw-augmented-parts}), and so is $\mathrm{Cl}^{\mathrm{ord}}(U,c,\phi)$.
\end{lemma}
\leanref{Development/AutomorphicArithmetic/AutomorphicVarieties/ModularityArgument/ThresholdArguments/DenseThresholdCases/TClParts.lean}{UpperTheory.V.Hi.clParts_alt}
\leanref{Development/AutomorphicArithmetic/AutomorphicVarieties/ModularityArgument/ThresholdArguments/DenseThresholdCases/TClParts.lean}{classical_in_iwahori_parts}
\begin{proof}
\emph{Averaging.} Let a group $H$ act linearly on a complex vector space $M$ by $\rho$, let $K\subset H$ be a subgroup, and let $(Y_i)$ be subspaces of $M$ stable under $K$, such that every element of every $Y_i$ is fixed by a subgroup of $H$ whose intersection with $K$ has finite index in $K$. For $F\in M$ let $E(K,F)$ be the subspace of the $y\in M$ such that $\rho(k)y=ay$ whenever $k\in K$, $a\in\mathbf C$ and $\rho(k)F=aF$. We show that if $F\in\sum_iY_i$, then $F\in\sum_i(Y_i\cap E(K,F))$. We may assume $F\ne0$. The $k\in K$ with $\rho(k)F\in\mathbf CF$ form a subgroup $K_F$, and $\rho(k)F=\chi_F(k)F$ for a character $\chi_F$ of $K_F$. Write $F=\sum_if_i$ with finitely many $f_i\in Y_i$, and let $K'\subset K_F$ be a subgroup of finite index in $K_F$ that fixes every $f_i$. Then $\chi_F$ is trivial on $K'$. The operator $P=[K_F:K']^{-1}\sum_s\chi_F(s)^{-1}\rho(s)$, the sum over representatives $s$ of $K_F/K'$, is well defined on the vectors fixed by $K'$, and maps them to vectors $w$ with $\rho(k)w=\chi_F(k)w$ for $k\in K_F$. It fixes $F$, and it maps $f_i$ into $Y_i$. So $F=\sum_iP(f_i)$, with $P(f_i)\in Y_i\cap E(K,F)$, since a $k\in K$ with $\rho(k)F=aF$ lies in $K_F$ and has $\chi_F(k)=a$.

\emph{Several groups.} Let groups $H_1,\dots,H_r$ act on $M$ by $\rho_1,\dots,\rho_r$, with $\rho_j(h)\rho_{j'}(h')=\rho_{j'}(h')\rho_j(h)$ for $j\ne j'$, let $K_j\subset H_j$ be subgroups, and let the $Y_i$ satisfy the hypothesis of the averaging for each $(H_j,K_j)$. Then $F\in\sum_iY_i$ implies $F\in\sum_i\big(Y_i\cap\bigcap_jE_j(K_j,F)\big)$, by the averaging applied successively for $j=1,\dots,r$ to the subspaces $Y_i\cap\bigcap_{j'<j}E_{j'}(K_{j'},F)$, which are stable under $K_j$, as $\rho_j$ commutes with the $\rho_{j'}$.

\emph{Classical forms.} Let $F\in\mathrm{Cl}(U,c,\phi)$. As $\phi$ is trivial on $T_0^{(c)}$, $F$ is invariant under right translation by $I_1(c)$, and so, by Lemma~\ref{sha-hecke-ideal-tw-augmented-hida-inputs}(3), it lies in the sum of the automorphic parts $Y$ of weight $\xi$. Each $Y$ is stable under the groups $GL_n(L_{w_j})\subset G_w$, through $\pi_\xi$, and $GL_n(L_{\tilde v})$, $v\mid\ell$, through right translation, and each element of $Y$ is fixed by an open subgroup of each of them, as $Y$ is a sum of images of smooth representations; an open subgroup meets a compact subgroup in a subgroup of finite index. The actions of these groups commute with each other. By the previous step, with the groups $K_j\subset GL_n(L_{w_j})$ and $GL_n(\mathcal O_{L_{\tilde v}})\subset GL_n(L_{\tilde v})$, $F\in\sum_Y(Y\cap E)$, where $E$ is the intersection of the $E_j(K_j,F)$ and the $E_v(GL_n(\mathcal O_{L_{\tilde v}}),F)$. Let $y\in Y\cap E$. It satisfies $y(\gamma g)=\xi(\gamma_\ell)y(g)$ and is fixed by the elements of $\mathcal K_0K_w$ whose components at the $v_j$ are $1$, as all elements of $Y$ are. It is fixed by the $K_j$, as $F$ is. So it is fixed by $U$: every $u\in U$ is $(uy_0^{-1})y_0$, where $y_0$ is the product of the components $u_{w_j}\in K_j$, and $uy_0^{-1}$ has components $1$ at the $v_j$. The groups $N_0$ and $\bar N_c$ are products of their components in the $GL_n(\mathcal O_{L_{\tilde v}})$, which fix $F$, and $T_0$ is generated by the $e_{v,i}(u)$, which act on $F$ by $\phi(e_{v,i}(u))^{-1}$; so $y$ is fixed by $N_0$ and $\bar N_c$, and $T_0$ acts on it by $\phi^{-1}$. Hence $y\in\mathrm{Cl}(U,c,\phi)$, and $\mathrm{Cl}(U,c,\phi)=\sum_Y\big(Y\cap\mathrm{Cl}(U,c,\phi)\big)$.

\emph{Ordinary part.} Let $W=\mathrm{Cl}(U,c,\phi)$, finite-dimensional and stable under $U'_z$ by Lemma~\ref{sha-hecke-ideal-tw-augmented-hida-classical}(1), and let $W_\nu$ be the generalized $\nu$-eigenspace of $U'_z$ in $W$. Each $Y$ is stable under $U'_z$, a combination of right translations by elements of $G_\ell$. So $Y\cap W$ is the sum of the $Y\cap W_\nu$, and $W=\sum_Y\sum_\nu(Y\cap W_\nu)$. As the $W_\nu$ for distinct $\nu$ are independent, $W_\mu=\sum_Y(Y\cap W_\mu)$ for every $\mu$. Summing over the $\mu$ in the definition of $\mathrm{Cl}^{\mathrm{ord}}(U,c,\phi)$ gives the claim for $\mathrm{Cl}^{\mathrm{ord}}(U,c,\phi)$.
% clParts_alt (TClParts.lean, 350 lines): h1 is classical_in_iwahori_parts, with classical_fixed_by_iwahori (the Lean reads the triviality of phi on Tc c0 from IsSpec.triv) and the hypothesis AutomorphicHeckeSet; cl_avg (the subgroup S of K, the character chi, cl_proj for the averaging P), cl_multi (induction over a Finset), cl_iso (smoothness of the elements of an IsoSum), cl_QQ, cl_VV, cl_QV for the commutation, cl_back (mem_fixd_lev for the level, cl_KO for N, Nm e and T inside KO). The ordinary part is cl_gen (Module.End.iSup_maxGenEigenspace_eq_top, independent_maxGenEigenspace) with cl_stab and classical_parts_commute. The Lean's ClAt is the book's Cl(lev K, e, phi) by spec_cl, and its OrdAt sums over the mu with |iota^-1 mu| = 1, the integral units (spec_unit).
\end{proof}

\begin{lemma}\label{sha-hecke-ideal-tw-augmented-hida-eigen}
Let $A$ be a commutative ring, $\kappa$ a field with a ring map $\chi:A\to\kappa$, $M$ an $A$-module and $W$ a $\kappa$-vector space, regarded as an $A$-module through $\chi$. Let $(t_i)_{i\in I}$ be $A$-linear endomorphisms of $M$ and $(t'_i)_{i\in I}$ $\kappa$-linear endomorphisms of $W$. Let $Y\subset Y_0$ be $A$-submodules of $M$, with $Y_0$ free and $Y$ finite, such that the $t_i$ map $Y_0$ and $Y$ into themselves and commute on $Y$, let $e$ be an idempotent of the $A$-algebra of endomorphisms of $Y_0$ generated by the $t_i$, with $Y=eY_0$, and let $\mathbf T_Y$ be the $A$-algebra of endomorphisms of $Y$ generated by the $t_i$. Let $\operatorname{sp}:M\to W$ be additive on $Y_0$, with $\operatorname{sp}(bf)=\chi(b)\operatorname{sp}(f)$ and $\operatorname{sp}(t_if)=t'_i\operatorname{sp}(f)$ for $b\in A$ and $f\in Y_0$, and mapping every basis of $Y_0$ over $A$ to vectors linearly independent over $\kappa$. Let $O$ and $E$ be the $\kappa$-spans of $\operatorname{sp}(Y_0)$ and $\operatorname{sp}(Y)$, and $(Z_j)$ subspaces of $W$ stable under the $t'_i$ with $O=\sum_j(Z_j\cap O)$. Then:
\begin{enumerate}
\item $E\subset O$ and $E=\sum_j(Z_j\cap E)$;
\item there is a ring map $\rho:\mathbf T_Y\to\operatorname{End}_\kappa(E)$ with $\rho(t)\operatorname{sp}(f)=\operatorname{sp}(tf)$ for $t\in\mathbf T_Y$ and $f\in Y$, so that $\rho(t_i)=t'_i$ on $E$ and $\rho(b)=\chi(b)$ for $b\in A$;
\item if $Z'\subset W$ is a subspace with $Z'\cap E\ne0$ such that every $t'_i$ acts on $Z'\cap O$ by a scalar, then there is a ring map $z:\mathbf T_Y\to\kappa$, equal to $\chi$ on $A$, such that every $\rho(t)$ acts on $Z'\cap E$ by $z(t)$.
\end{enumerate}
\end{lemma}
\leanref{Development/AutomorphicArithmetic/AutomorphicVarieties/ModularityArgument/ThresholdArguments/DenseThresholdCases/HeckeOperatorPreservesPart.lean}{pa_fin}
\begin{proof}
\emph{Base change.} Let $(b_k)$ be a basis of $Y_0$ over $A$. The $A$-linear map $f\mapsto\operatorname{sp}(f)$ from $Y_0$ to $W$ extends to a $\kappa$-linear map $\kappa\otimes_AY_0\to W$, which is injective, as it maps the basis $(1\otimes b_k)$ to linearly independent vectors, and whose image is $O$. So an $A$-linear endomorphism $\varphi$ of $Y_0$ induces the endomorphism $\rho_0(\varphi)=1\otimes\varphi$ of $O$, with $\rho_0(\varphi)\operatorname{sp}(f)=\operatorname{sp}(\varphi f)$, and $\rho_0$ is a ring map with $\rho_0(b)=\chi(b)$ for $b\in A$. As $Y=eY_0$ is a direct summand of $Y_0$, $\kappa\otimes_AY\to\kappa\otimes_AY_0$ is injective, so $E\cong\kappa\otimes_AY$ in the same way, and an endomorphism $t$ of $Y$ induces $\rho(t)$ on $E$, with $\rho(t)\operatorname{sp}(f)=\operatorname{sp}(tf)$ for $f\in Y$.

\emph{Generators.} If $t$ is the restriction of $t_i$ to $Y$, then $\rho(t)=t'_i$ on $E$: both are $\kappa$-linear and agree on $\operatorname{sp}(Y)$, which spans $E$. Hence, by induction on the generation of the algebra, for every $t\in\mathbf T_Y$ and every subspace $Z\subset W$ stable under the $t'_i$, $\rho(t)$ maps $Z\cap E$ into $Z$; and in the same way $\rho_0(\varphi)$ maps $Z\cap O$ into $Z$ for $\varphi$ in the algebra generated by the $t_i$ on $Y_0$.

(1) As $\operatorname{sp}(Y)\subset\operatorname{sp}(Y_0)$, $E\subset O$. The endomorphism $\rho_0(e)$ maps $O$ into $E$, as $\operatorname{sp}(ef)\in\operatorname{sp}(Y)$, and fixes $E$, as $\operatorname{sp}(ef)=\operatorname{sp}(f)$ for $f\in Y$. Let $F\in E$, and write $F=\sum_jF_j$ with $F_j\in Z_j\cap O$. Then $F=\rho_0(e)F=\sum_j\rho_0(e)F_j$, with $\rho_0(e)F_j\in Z_j\cap E$.

(2) This is $\rho$, restricted to $\mathbf T_Y$.

(3) By induction on the generation of $\mathbf T_Y$, every $\rho(t)$ acts on $Z'\cap E$ by a scalar $z(t)$: the generators act on $Z'\cap E\subset Z'\cap O$ by scalars, the elements of $A$ by $\chi$, and sums and products of operators acting by scalars act by scalars. Comparing on a nonzero $v\in Z'\cap E$, $t\mapsto z(t)$ is a ring map, and $z=\chi$ on $A$.
% pa_fin, from pa_main and pa_inst: pa_rho (Psi : K tensor N = E through LinearMap.liftBaseChange, injective by hind on Module.Free.chooseBasis, and rho = conjugation of baseChangeHom), pa_op, pa_stab (Algebra.adjoin_induction), hfix, hin and hsum for (1), pa_chr for (3), and pa_pt to pass from z to an O-algebra map y with iota y = z. The Lean's fourth conclusion, a nonzero eigenvector in some Z_j meet O for every character z of T_Y equal to chi on A (pa_eig, AnnihilatorQuotientFinite.pointEig), is not stated.
\end{proof}

\begin{lemma}\label{sha-hecke-ideal-tw-augmented-hida}
In Situation~\ref{sha-hecke-ideal-tw-augmented-setup}:
\begin{enumerate}
\item every level $\mathcal K_0K$ is neat and of finite level;
\item for every level $U=\mathcal K_0K$, $\mathcal F^{\mathrm{ord}}(U)$ and $M_K$ are finite and free over $\Lambda$, and the operators $T_{w,j}$ and $U_{v,j}$ map them into themselves and commute on them;
\item for every level $U=\mathcal K_0K$ there is a $\Lambda$-linear map $e_U:\mathcal F(U)\to\mathcal F^{\mathrm{ord}}(U)$, equal to the identity on $\mathcal F^{\mathrm{ord}}(U)$, such that for every $f\in\mathcal F(U)$ and $k\ge0$ there is $N'$ such that $U_z^{N!}f-e_Uf$ has values in $\mathfrak m_\Lambda^k\Lambda^m$ for all $N\ge N'$;
\item for every weight $\lambda$ and every arithmetic point $x$ of $\Lambda$ of weight $\lambda$ there are a representation $\xi$ of $G_\ell$ on some $\mathbf C^d$, algebraic of weight $\lambda$ through $\iota$, and a map $\operatorname{sp}:\mathcal F\to V^\xi$ (Definition~\ref{sha-hecke-ideal-tw-augmented-parts}) such that for every level $U=\mathcal K_0K$: $\operatorname{sp}$ is additive on $\mathcal F^{\mathrm{ord}}(U)$ and $\operatorname{sp}(bf)=\iota x(b)\operatorname{sp}(f)$ for $b\in\Lambda$ and $f\in\mathcal F^{\mathrm{ord}}(U)$; the images under $\operatorname{sp}$ of a basis of $\mathcal F^{\mathrm{ord}}(U)$ over $\Lambda$ are linearly independent over $\mathbf C$; for each operator $t$ among the $T_{w,j}$ and $U_{v,j}$ there is an endomorphism $t^\xi$ of $V^\xi$, with $T_{w,j}^\xi$ as in Definition~\ref{sha-hecke-ideal-tw-augmented-parts}, such that $\operatorname{sp}(tf)=t^\xi\operatorname{sp}(f)$ for $f\in\mathcal F^{\mathrm{ord}}(U)$; $\operatorname{sp}(\rho(g)f)=\pi_\xi(g)\operatorname{sp}(f)$ whenever $g$ lies in one factor $GL_n(L_{w_j})$ of $G_w$, $f\in\mathcal F^{\mathrm{ord}}(U)$ and $\rho(g)f\in\mathcal F^{\mathrm{ord}}(U')$ for a level $U'$; and the span of $\operatorname{sp}(\mathcal F^{\mathrm{ord}}(U))$ is the sum of its intersections with the automorphic parts of weight $\xi$;
\item in (4), for every level $\mathcal K_0K$ with $M_K\ne0$, let $E_K$ be the span of $\operatorname{sp}(M_K)$; there is a ring map $\mathbf T_K\to\operatorname{End}_{\mathbf C}(E_K)$, $t\mapsto t^x$, with $t^x\operatorname{sp}(f)=\operatorname{sp}(tf)$ for $f\in M_K$; and if $x$ is of distinct type, then $E_K$ is the sum of its intersections with the automorphic parts of weight $\xi$, and for every automorphic part $Y$ of weight $\xi$ with $Y\cap E_K\ne0$ there is an $\mathcal O$-algebra map $y:\mathbf T_K\to\bar{\mathbf Q}_\ell$, equal to $x$ on $\Lambda$, such that every $t^x$ acts on $Y\cap E_K$ by multiplication by $\iota y(t)$.
\end{enumerate}
\end{lemma}
\leanref{Development/AutomorphicArithmetic/AutomorphicVarieties/HeckeAlgebras/OrdinaryHeckeCharacters/NtExpand.lean}{UpperTheory.V.Hi.neatU}
\leanref{Development/AutomorphicArithmetic/AutomorphicVarieties/HeckeAlgebras/CoefficientRepTrivialOnUnramified.lean}{UpperTheory.V.Hi.HSet.Ax}
\leanref{Development/AutomorphicArithmetic/AutomorphicVarieties/HeckeAlgebras/OrdinaryHeckeCharacters/FuCmpMem.lean}{UpperTheory.V.Hi.finU}
\leanref{Development/AutomorphicArithmetic/AutomorphicVarieties/HeckeAlgebras/OrdinaryHeckeCharacters/HiComm.lean}{UpperTheory.V.Hi.hiComm}
\leanref{Development/AutomorphicArithmetic/AutomorphicVarieties/HeckeAlgebras/OrdinaryHeckeCharacters/OrdFree.lean}{UpperTheory.V.Hi.ordStruct}
\leanref{Development/AutomorphicArithmetic/AutomorphicVarieties/HeckeAlgebras/OrdinaryHeckeCharacters/SpecIntIff.lean}{UpperTheory.V.Hi.spec}
\leanref{Development/AutomorphicArithmetic/AutomorphicVarieties/HeckeAlgebras/CoefficientRepTrivialOnUnramified.lean}{UpperTheory.V.Hi.HSet.IsSpec}
\leanref{Development/AutomorphicArithmetic/AutomorphicVarieties/ModularityArgument/ThresholdArguments/DenseThresholdCases/UnitaryPointProperties.lean}{UpperTheory.V.Hi.HSet.AutomorphicHeckeSet}
\leanref{Development/AutomorphicArithmetic/AutomorphicVarieties/ModularityArgument/ThresholdArguments/DenseThresholdCases/TClParts.lean}{UpperTheory.V.Hi.clParts_alt}
\leanref{Development/AutomorphicArithmetic/AutomorphicVarieties/ModularityArgument/ThresholdArguments/DenseThresholdCases/HeckeOperatorPreservesPart.lean}{UpperTheory.V.Hi.partAct0_alt}
\begin{proof}
(1) \emph{Neatness.} Let $U=\mathcal K_0K$, and let $\gamma g=g(u,b)$ with $\gamma\in U(n)(L^+)$, $g\in G^{(\ell)}\times G_\ell$, $u\in U$ and $b\in I_\ell$. View $\gamma$ as a matrix $M\in GL_n(L)$ with ${}^tc(M)M=1$, and for a place $w$ of $L$ write $g_w$, $u_w$ and $b_w$ for the components in $GL_n(L_w)$, through $U(n)(L^+_v)\subset GL_n(L_w)$ for the place $v$ below $w$. Comparing components, $g_w^{-1}M^kg_w=u_w^k$ for every place $w$ of $L$ not above $\ell$, and $g_{\tilde v}^{-1}M^kg_{\tilde v}=b_{\tilde v}^k$ for $v\mid\ell$, for all $k\ge0$. For each $j$, $u_{w_j}$ lies in the compact group $K_j$, in which $K_j\cap GL_n(\mathcal O_{L_{w_j}})$ is open, hence of finite index; so $u_{w_j}^{e_j}\in GL_n(\mathcal O_{L_{w_j}})$ for some $e_j\ge1$. Let $N=\prod_je_j$. For every finite place $v$ of $L^+$ there is a place $w$ of $L$ above $v$ at which $M^N$ is conjugate into $GL_n(\mathcal O_{L_w})$: for $v\mid\ell$ the place $\tilde v$, as $b\in I_\ell$ is integral; for $v\in S_{\mathrm a}$ and $v=v_0$ the place $\tilde v$, as $u$ is integral there; for $v=v_j$ the place $w_j$; and for $v\notin S'$ any place above $v$, as the components of $u$ there are integral. By Lemma~\ref{sha-hecke-ideal-tw-augmented-finite-order}, applied to $M^N$, $M^N$ has finite order, so $M$ has, and so has $u_{\tilde v_0}=g_{\tilde v_0}^{-1}Mg_{\tilde v_0}$. As $u_{\tilde v_0}$ lies in the principal congruence subgroup of level $\tilde v_0$, and $p$ is odd and unramified in $L$, $u_{\tilde v_0}=1$ by Lemma~\ref{sha-hecke-ideal-tw-augmented-torsion}. Hence $M=1$, so $\gamma=1$, and $g=g(u,b)$ gives $u=1$ and $b=1$.

\emph{Finite level.} Let $U=\mathcal K_0K$. Let $C\subset U(n)(\mathbf A^\infty_{L^+})$ be the subgroup of the elements whose components at all finite places of $L$ are integral, whose component at $\tilde v_0$ lies in the principal congruence subgroup of level $\tilde v_0$, whose component at each $w_j$ lies in $K_j$, and whose components at the places $\tilde v$, $v\mid\ell$, lie in the subgroup $I_1(0)$ of Lemma~\ref{sha-hecke-ideal-tw-augmented-uop}, formed for the factors $GL_n(L_{\tilde v})$. It is open: the integral elements form an open subgroup, and the other conditions are open conditions at finitely many places. Let $\Sigma$ be a finite set as in Lemma~\ref{sha-hecke-ideal-tw-augmented-hida-classes} for $C$.

Let $\Phi:U(n)(\mathbf A_{L^+})\to G^{(\ell)}\times G_\ell$ take an element to its components at the finite places away from $\ell$ and at the places $\tilde v$, $v\mid\ell$. It restricts to the embedding of $U(n)(L^+)$, and it is trivial on $U(n)(L^+\otimes\mathbf R)$. It is surjective. Indeed, let $(a,b)\in G^{(\ell)}\times G_\ell$, and let $h$ have components $a$ away from $\ell$. Each $v\mid\ell$ is split in $L$, so $U(n)(L^+_v)=GL_n(L_{\tilde v})$, and there is $y_v\in U(n)(\mathbf A_{L^+})$ with component $h_{\tilde v}^{-1}b_v$ at $\tilde v$ and $1$ at the places of $L$ not above $v$. With $y=\prod_{v\mid\ell}y_v$, $\Phi(hy)=(a,b)$.

Let $g\in G^{(\ell)}\times G_\ell$. Write $g=\Phi(h)$, then $h=\gamma sxk$ with $\gamma\in U(n)(L^+)$, $s\in\Sigma$, $x\in U(n)(L^+\otimes\mathbf R)$ and $k\in C$, and $\Phi(k)=(u,b)$. Then $g=\gamma\,\Phi(s)\,(u,b)$. The element $b$ lies in $I_1(0)$, which is generated by $N_0$, $T_0^{(0)}$ and $\bar N_0$, so $b\in I_\ell$. The element $u$ lies in $G^{(\ell)}_0$, as $k$ is integral at the places of $S_{\mathrm a}$ and lies in the principal congruence subgroup at $\tilde v_0$. It lies in $U$. Indeed, for each $j$ let $y_j\in G^{(\ell)}_0$ have component $k_{w_j}\in K_j$ at $w_j$ and $1$ at the places of $L$ not above $v_j$, and let $y=y_1\cdots y_q$, an element of $U$. The components of $uy^{-1}$ are $1$ at the $w_j$ and integral at the places outside $S'$, so $uy^{-1}\in U$, and $u=(uy^{-1})y\in U$. Hence the finitely many $\Phi(s)$, $s\in\Sigma$, represent all the double cosets of $U(n)(L^+)\backslash(G^{(\ell)}\times G_\ell)/(U\times I_\ell)$.

(2) \emph{Finiteness and freeness.} Let $U=\mathcal K_0K$. By (1), $U$ is of finite level, and it satisfies the freeness condition of Lemma~\ref{sha-hecke-ideal-tw-augmented-sections}, since by neatness $\gamma g=g(u,b)$ forces $u=1$ and $b=1$. By Lemma~\ref{sha-hecke-ideal-tw-augmented-ordinary}(1), $\mathcal F^{\mathrm{ord}}(U)$ is finite and free over $\Lambda$.

\emph{The operators.} Let $w$ be a split place of $L$ above a place $v\notin S'$ of $L^+$, so that $GL_n(L_w)=U(n)(L^+_v)\subset G^{(\ell)}_0$, and let $K_w=GL_n(\mathcal O_{L_w})$. The component of $U$ at $v$ is $K_w$, and $U=K_wU^v$, where $U^v$ is the subgroup of the elements of $U$ with component $1$ at $v$, which commute with $GL_n(L_w)$. For $f\in\mathcal F^{\mathrm{ord}}(U)$, $T_{w,j}f=\sum_hh\cdot f$, the sum over representatives $h$ of $K_w\alpha_jK_w/K_w$, which is finite by Lemma~\ref{sha-hecke-ideal-tw-augmented-spherical}. It is fixed by $K_w$, as a double coset operator, and by $U^v$, whose elements commute with the $h$ and fix $f$. Each $h\cdot f$ is an ordinary $\Lambda$-adic form, as in the proof of Lemma~\ref{sha-hecke-ideal-tw-augmented-projective}; so $T_{w,j}f\in\mathcal F^{\mathrm{ord}}(U)$. The $U_{v,j}=U_{\alpha_j}$, with $\alpha_j$ dominant, map $\mathcal F^{\mathrm{ord}}(U)$ into itself by Lemma~\ref{sha-hecke-ideal-tw-augmented-forms}(1).

\emph{Commutation.} Let $w$ and $w'$ be split places of $L$ above places $v,v'\notin S'$. If $v\ne v'$, the subgroups $GL_n(L_w)$ and $GL_n(L_{w'})$ of $G^{(\ell)}_0$ commute, so $T_{w,j}$ and $T_{w',j'}$, sums of translations by their elements, commute. If $v=v'$, both are $U(n)(L^+_v)$, through an isomorphism $GL_n(L_{w'})\cong GL_n(L_w)$ that carries $K_{w'}$ onto $K_w$; so $T_{w',j'}$ is the double coset operator $[K_w\beta K_w]$ for the image $\beta$ of $\alpha_{j'}$, and it commutes with $T_{w,j}$ on functions fixed by $K_w$, by Lemma~\ref{sha-hecke-ideal-tw-augmented-spherical}. The operators $U_{v'',j}$, $v''\mid\ell$, commute with the $T_{w,j}$, which are sums of maps $f\mapsto h\cdot f$, by Lemma~\ref{sha-hecke-ideal-tw-augmented-forms}(1); and they commute with each other, as $U_\alpha U_\beta=U_{\alpha\beta}=U_\beta U_\alpha$ on functions fixed by $N_0$, by Lemma~\ref{sha-hecke-ideal-tw-augmented-uop}(1).

\emph{The localizations.} Let $Y=\mathcal F^{\mathrm{ord}}(U)$ and let $e$ be an idempotent as in Lemma~\ref{sha-hecke-ideal-tw-augmented-idempotent}(1) for $Y$ and the operators. By (3) of that lemma, $M_K=eY$ and $e$ commutes with the operators. So $Y=eY\oplus(1-e)Y$, the operators map $M_K$ into itself and commute on it, and $M_K$, a direct summand of the finite free module $Y$, is finite and projective, hence free over the local ring $\Lambda$.

(3) This is Lemma~\ref{sha-hecke-ideal-tw-augmented-ordinary}(3), which applies to $U=\mathcal K_0K$ as in (2).

(4) Let $\lambda$ be a weight, $x$ an arithmetic point of $\Lambda$ of weight $\lambda$, and $e_0$ as in Definition~\ref{sha-hecke-ideal-tw-augmented-parts}(1).

\emph{The ring of values.} By Lemma~\ref{sha-automorphy-unramified-near-tangent-points}, $x(\Lambda)$ is a finitely generated $\mathcal O$-module and $|x(b)|<1$ for $b\in\mathfrak m_\Lambda$. So the values of $x$ lie in the ring of integers $\mathcal O_1$ of a finite extension of $\mathbf Q_\ell$. Let $\mathcal O'\supset\mathcal O_1$ and $(\xi,\chi,c_0)$ be as in Lemma~\ref{sha-hecke-ideal-tw-augmented-hida-inputs}(1) for $\lambda$ and $\mathcal O_1$, and let $c=\max(e_0,c_0)$. Then $x$ is an $\mathcal O$-algebra map $\Lambda\to\mathcal O'$ with $x(\mathfrak m_\Lambda)\subset\mathfrak m'$, as the elements of $\mathcal O'$ of absolute value less than $1$ lie in $\mathfrak m'$.

\emph{The character $\phi$.} For $t\in T_0$, $\xi(t)$ and $\xi(t^{-1})$ have entries in $\mathcal O'$, and their entries $(0,0)$ are $\chi(t)$ and $\chi(t)^{-1}$; so $\chi(t)\in\mathcal O'^\times$. Let $\phi(t)=x(\theta(t))\chi(t)^{-1}$, a character of $T_0$ with values in $\mathcal O'^\times$, so that $x(\theta(t))=\phi(t)\chi(t)$. For $v\mid\ell$, $1\le i\le n$ and $u\in\mathcal O_{L_{\tilde v}}^\times$, the formula for $\chi$ in Lemma~\ref{sha-hecke-ideal-tw-augmented-hida-inputs}(1) gives
\[
\phi(e_{v,i}(u))=x(\theta(e_{v,i}(u)))\prod_\tau\tau(u)^{-\lambda_{\tau,n+1-i}}=\chi_{x,v,i}(u),
\]
the character of Definition~\ref{sha-hecke-ideal-tw-augmented-parts}(1), which is $1$ for $u\equiv1$ modulo $\varpi^{c+1}$, as $c\ge e_0$. As $T_0^{(c)}$ is generated by these $e_{v,i}(u)$, $\phi$ is trivial on $T_0^{(c)}$.

\emph{The map $\operatorname{sp}$.} The levels are of finite level by (1), and they satisfy the freeness condition by neatness, as in (2); any two levels $\mathcal K_0K$ and $\mathcal K_0K'$ contain $\mathcal K_0(K\cap K')$. Let $\operatorname{sp}$ be the map of Lemma~\ref{sha-hecke-ideal-tw-augmented-hida-control}(2) for the set of all levels, $x$, $(\xi,\chi,c_0)$, $c$ and $\phi$, and let $U=\mathcal K_0K$ be a level. The additivity, the compatibility with $\Lambda$ and the linear independence are in Lemma~\ref{sha-hecke-ideal-tw-augmented-hida-control}(2). For $U_{v,j}$, let $U^\xi_{v,j}=U'_{\alpha_j}$, with $\alpha_j$ in the factor at $v$; then $\operatorname{sp}(U_{v,j}f)=U^\xi_{v,j}\operatorname{sp}(f)$ by Lemma~\ref{sha-hecke-ideal-tw-augmented-hida-control}(2). For $T_{w,j}$, $T_{w,j}f=\sum_hh\cdot f$ and $T^\xi_{w,j}F=\sum_hh\cdot F$, the sums over representatives $h\in GL_n(L_w)\subset G^{(\ell)}_0$ of $GL_n(\mathcal O_{L_w})\alpha_jGL_n(\mathcal O_{L_w})/GL_n(\mathcal O_{L_w})$. The argument of \emph{The operators} in (2) shows that $\sum_hh\cdot$ maps the elements of $V^\xi$ fixed by $U$ to elements fixed by $U$; so $\operatorname{sp}(T_{w,j}f)=T^\xi_{w,j}\operatorname{sp}(f)$ by Lemma~\ref{sha-hecke-ideal-tw-augmented-hida-control}(2). Let $g\in GL_n(L_{w_j})$, and let $K'$ be the family $K$ with $K_j$ replaced by $gK_jg^{-1}$. For $F\in V^\xi$ fixed by $U$, $g\cdot F$ is fixed by the elements of $\mathcal K_0K'$ with components $1$ at the $v_{j'}$ and by the $K_{j'}$ with $j'\ne j$, which commute with $g$, and by $gK_jg^{-1}$. These generate $\mathcal K_0K'$, as in the proof of Lemma~\ref{sha-hecke-ideal-tw-augmented-hida-decomposition}. So $\operatorname{sp}(\rho(g)f)=\pi_\xi(g)\operatorname{sp}(f)$ by Lemma~\ref{sha-hecke-ideal-tw-augmented-hida-control}(2), with the single element $g$.

\emph{The decomposition.} By Lemma~\ref{sha-hecke-ideal-tw-augmented-hida-control}(2), the span of $\operatorname{sp}(\mathcal F^{\mathrm{ord}}(U))$ is $\mathrm{Cl}^{\mathrm{ord}}(U,c,\phi)$, which is the sum of its intersections with the automorphic parts of weight $\xi$ by Lemma~\ref{sha-hecke-ideal-tw-augmented-hida-decomposition}.

(5) Let $x$, $\xi$, $c$, $\phi$ and $\operatorname{sp}$ be as in (4), let $U=\mathcal K_0K$ be a level with $M_K\ne0$, and let $Y_0=\mathcal F^{\mathrm{ord}}(U)$ and $S=\mathrm{Cl}^{\mathrm{ord}}(U,c,\phi)$.

\emph{The Hecke action.} By (2), $Y_0$ and $M_K$ are finite and free over $\Lambda$, and the operators map them into themselves and commute on them. Let $e$ be the idempotent of Lemma~\ref{sha-hecke-ideal-tw-augmented-idempotent}(1) for $Y_0$ and the operators, so that $M_K=eY_0$ by (3) of that lemma. Each automorphic part $Y$ of weight $\xi$ is stable under the $T^\xi_{w,j}$, which act on it by the scalars $a_{Y,w,j}$, and under the $U^\xi_{v,j}$: it is stable under right translation by each $GL_n(L_{\tilde v})$, as a sum of images of equivariant maps, and $U^\xi_{v,j}$ is a combination of such translations. By (4), the span of $\operatorname{sp}(Y_0)$ is $S$, the sum of its intersections with the automorphic parts. So Lemma~\ref{sha-hecke-ideal-tw-augmented-hida-eigen} applies to $A=\Lambda$, $\kappa=\mathbf C$, the map $\iota x$, $M=\mathcal F$, $W=V^\xi$, the operators $T_{w,j}$ and $U_{v,j}$ and their counterparts $T^\xi_{w,j}$ and $U^\xi_{v,j}$, $Y_0$, $M_K$, $e$, $\operatorname{sp}$ and the automorphic parts of weight $\xi$. It gives the ring map $t\mapsto t^x$, and $E_K=\sum_Y(Y\cap E_K)$, with $E_K\subset S$.

\emph{The $U_{v,j}$ are onto.} For every dominant $\alpha$, $U_\alpha$ maps $Y_0$ bijectively onto itself. Indeed, write the component of $\alpha$ at each $v\mid\ell$ as $\operatorname{diag}(\varpi^{a_1},\dots,\varpi^{a_n})$ with $a_1\ge\dots\ge a_n$, and let $k\ge0$ be at least every difference $a_i-a_{i+1}$. The components of $z^k$ are $\operatorname{diag}(\varpi^{k(n-1)},\dots,\varpi^k,1)$, so $\beta=z^k\alpha^{-1}$ is dominant, and $\alpha\beta=\beta\alpha=z^k$, as diagonal matrices commute. The operators $U_\alpha$, $U_\beta$ and $U_z$ map $Y_0$ into itself by Lemma~\ref{sha-hecke-ideal-tw-augmented-forms}(1), and the elements of $Y_0$ are fixed by $N_0$; so $U_z^k=U_{z^k}=U_\beta U_\alpha=U_\alpha U_\beta$ on $Y_0$, by Lemma~\ref{sha-hecke-ideal-tw-augmented-uop}(1). By Lemma~\ref{sha-hecke-ideal-tw-augmented-ordinary}(1), $U_z$, and so $U_z^k$, is bijective on $Y_0$. If $U_\alpha f=U_\alpha f'$ with $f,f'\in Y_0$, then $U_z^kf=U_\beta U_\alpha f=U_\beta U_\alpha f'=U_z^kf'$, so $f=f'$. Every $g\in Y_0$ is $U_z^kh$ for some $h\in Y_0$, so $g=U_\alpha(U_\beta h)$ with $U_\beta h\in Y_0$. In particular each $U_{v,j}=U_{\alpha_j}$, with $\alpha_j$ in the factor at $v$ and $1$ in the others, maps $Y_0$ onto itself. As $\operatorname{sp}(U_{v,j}f)=U^\xi_{v,j}\operatorname{sp}(f)$ and $\operatorname{sp}(Y_0)$ spans $S$, $U^\xi_{v,j}$ maps $S$ onto $S$, hence bijectively, as $S$ is finite-dimensional by Lemma~\ref{sha-hecke-ideal-tw-augmented-hida-control}(1). For an automorphic part $Y$, $U^\xi_{v,j}$ maps $Y\cap S$ into itself, injectively, hence onto.

\emph{The characters.} Let $x$ be of distinct type, and $Y$ an automorphic part of weight $\xi$ with $Y\cap E_K\ne0$. Let $v\mid\ell$. Apply Lemma~\ref{sha-hecke-ideal-tw-augmented-hida-inputs}(4) to $L_{\tilde v}$, right translation by $GL_n(L_{\tilde v})$ on $V^\xi$, $Y$, $\sigma_{Y,v}$, the subspace $Y\cap S$, which is nonzero as it contains $Y\cap E_K$, $c$, and the maps $f_i(u)=\phi(e_{v,i}(u))=\chi_{x,v,i}(u)$ on $\mathcal O_{L_{\tilde v}}^\times$. Its hypotheses hold: $Y\cap S$ is fixed by $N_0$ and $\bar N_c$, and $e_{v,i}(u)$ acts on it by $f_i(u)^{-1}$, as $S\subset\mathrm{Cl}(U,c,\phi)$; $f_i(u)=1$ for $u\equiv1$ modulo $\varpi^{c+1}$, as $\phi$ is trivial on $T_0^{(c)}$; the $f_i$ are pairwise distinct, as $x$ is of distinct type; and the $U_{\alpha_j}=\chi(\alpha_j)U^\xi_{v,j}$ map $Y\cap S$ onto itself. So every $U^\xi_{v,j}$ acts on $Y\cap S$ by a scalar. As the $T^\xi_{w,j}$ act on $Y$ by scalars, Lemma~\ref{sha-hecke-ideal-tw-augmented-hida-eigen}(3), with $Z'=Y$, gives a ring map $z:\mathbf T_K\to\mathbf C$, equal to $\iota x$ on $\Lambda$, such that every $t^x$ acts on $Y\cap E_K$ by $z(t)$. Then $y=\iota^{-1}\circ z$ is an $\mathcal O$-algebra map, equal to $x$ on $\Lambda$, and $t^x$ acts on $Y\cap E_K$ by $\iota y(t)$.
% (1) is neatU (nt_rat, nt_pow, nt_Iw, nt_out, nt_Q, nt_ord applied to M ^ (prod e q), then nt_free at s.w s.v0 with hp0 and hl0) and the field fin of HSet.Ax, supplied by finU in axU_s (FuCmpMem.lean, 195 lines), which the book follows: Phi = (away, cLoc at the s.w v), surjective by hsurj (fu_patch over the v in Vl with the sections of secU, trivial away from v), and trivial on infU by hinf; K0 (KO at every place, Kc 1 at s.w v0, K q at the w q, I1 0 at the s.w v) open by ho (fu_open, fu_Kc, fu_iInf); classSet_finite (the book's Lemma sha-hecke-ideal-tw-augmented-hida-classes) for K0; the representatives Phi of S; u in Ub from hk1 and hk2 (fu_cmp_mem), the l-part in Iw 0 by fu_Iw, and u in lev K, a supremum UQ sup the sec (K q), by fu_patch over the q with secU: y in the supremum with the components of k at the w q, and u y^-1 in UQ. The book's level is componentwise (levelUR2LevOut and mem_heckeLevel_of_local_components_mem identify the two), so the book's y and u y^-1 lie in U by their components. (2) is hiComm: ordFree (from ordStruct) with the freeness condition free_of_neat from neatU (as in LevelUFastA, where IsWt is built), hiC_main (hiC_Tfix with W_le, W_dec and sectionsCommuteAtSplit; hiC_dcOrd and hiC_ord; hiC_TT with Wm_mul_Wm_comm, hiC_two, hiC_equiv and sphComm; hUT through spStab; uOpMul and Iwa.comm for the U operators), then AnnihilatorQuotientFinite.locFree (locIdem, the complement ker e, Module.free_of_flat_of_isLocalRing). (3) is ordStruct.
% (4) is spec (SpecIntIff.lean, 201 lines), followed step by step: GradeZeroInterpolatingElement.coPt (the book cites Lemma sha-automorphy-unramified-near-tangent-points, which is Reformulated.B.CoPoint.coPt with the same content), wtRep (Lemma sha-hecke-ideal-tw-augmented-hida-inputs (1)), the local map y (IsLocalHom through spec_intIff and the bound of coPt), spec_chi for chi on T0, phi = y theta / psi with hphichi and hphie (Sp.chi_eV), htriv and hphi1 (Sp.gen_Tc), then Sp.absSpec (Lemma sha-hecke-ideal-tw-augmented-hida-control) with H.lev, Sp.lev_dir, Ax.fin and IsWt.free; spec_cl identifies ClAt with Cl(lev K, e, phi), hU and hz (Sp.chi_alV, weightChar_central_eq_prod_uniformizerDiag) identify opC and uzC with the UF, spec_unit identifies OrdAt with ClOrd; the T operators through Sp.lev_sph (W_le, W_dec, sectionsCommuteAtSplit) and Sp.sum_cl, act through Sp.lev_act (Function.update of K at q by the conjugate of K q, with Q_UQ and Q_Q; its hypothesis that rho(g) f lies in some Ord is not used). The decomposition is the second conjunct of clParts_alt (Lemma sha-hecke-ideal-tw-augmented-hida-decomposition), under AutomorphicHeckeSet (Lemma sha-hecke-ideal-tw-augmented-hida-inputs (3)). (5) is partAct0_alt (HeckeOperatorPreservesPart.lean, 336 lines): hiComm, locIdem, pa_fin (Lemma sha-hecke-ideal-tw-augmented-hida-eigen) with Z the parts satisfying P, hecke_operator_preserves_part (pa_pi, pa_iso) and clParts_alt for hOZ; it gives E = sum of the Y meet E at every arithmetic x. Under HSet.Good, goodScalB_alt (DenseThresholdCases/CharOfScalar.lean, 307 lines): ordFree for the bijectivity of U_a on Ord (lev K), a in I.P, which the book follows (OrdFree.lean: Iwa.exists_pow_mul_inv_mem for k with b = z^k a^-1 in P, Iwa.comm, hpow and hmul through uOpMul, hbij of ordStruct iterated k times; then injectivity through U_b and surjectivity through U_a (U_b h)); hSmap; onto_inf for honto; local_step (Lemma sha-hecke-ideal-tw-augmented-hida-inputs (4)) at each v; unit_eig, which makes the eigenvalue an l-adic unit, is not needed here. Then h5 of pa_fin and pa_pt give y.
\end{proof}

\subsection{The Hecke ideal}\label{subsection-sha-hecke-ideal}

\begin{lemma}\label{sha-hecke-ideal-tw-augmented-projective}
In Situation~\ref{sha-hecke-ideal-tw-augmented-setup}, let $V\subset U$ be subgroups of $G^{(\ell)}_0$ with $V$ normal in $U$ of finite index, and $U$ neat and of finite level. For $q\in U/V$ let $\tilde q\in U$ be a representative. Then $U$ maps $\mathcal F^{\mathrm{ord}}(V)$ into itself, and:
\begin{enumerate}
\item $\mathcal F^{\mathrm{ord}}(V)$, with this action of $U/V$, is projective over the group ring $\Lambda[U/V]$;
\item the trace $\operatorname{tr}f=\sum_{q\in U/V}\tilde q\cdot f$ maps $\mathcal F^{\mathrm{ord}}(V)$ onto $\mathcal F^{\mathrm{ord}}(U)$;
\item the elements $f\in\mathcal F^{\mathrm{ord}}(V)$ with $\operatorname{tr}f=0$ form the $\Lambda$-span of the $u\cdot f'-f'$ with $u\in U$ and $f'\in\mathcal F^{\mathrm{ord}}(V)$.
\end{enumerate}
\end{lemma}
\leanref{Development/AutomorphicArithmetic/AutomorphicVarieties/HeckeAlgebras/OrdinaryHeckeCharacters/OrdProj.lean}{UpperTheory.V.Hi.ordProj}
\begin{proof}
\emph{The level $V$.} If $\gamma g=g(u,b)$ with $\gamma\in U(n)(L^+)$, $g\in G^{(\ell)}\times G_\ell$, $u\in U$ and $b\in I_\ell$, then $u=1$ and $b=1$, as $U$ is neat; so $U$, and its subgroup $V$, satisfy the freeness condition of Lemma~\ref{sha-hecke-ideal-tw-augmented-sections}. The subgroup $V$ is of finite level: if $s$ is a finite set of representatives of $U(n)(L^+)\backslash(G^{(\ell)}\times G_\ell)/(U\times I_\ell)$, every $g$ is $\gamma x(u,b)$ with $x\in s$, $u\in U$ and $b\in I_\ell$, and writing $u=\tilde qv$ with $q=uV$ and $v\in V$ gives $g=\gamma\,x(\tilde q,1)(v,b)$; so the finitely many $x(\tilde q,1)$ are representatives for $V$. By Lemma~\ref{sha-hecke-ideal-tw-augmented-ordinary}, $\mathcal F^{\mathrm{ord}}(V)$ is finite and free over $\Lambda$, and there is a $\Lambda$-linear map $e_V:\mathcal F(V)\to\mathcal F^{\mathrm{ord}}(V)$, equal to the identity on $\mathcal F^{\mathrm{ord}}(V)$, such that for every $f\in\mathcal F(V)$ and $k\ge0$ there is $N'$ such that $U_z^{N!}f-e_Vf$ has values in $\mathfrak m_\Lambda^k\Lambda^m$ for all $N\ge N'$. Also $\mathcal F^{\mathrm{ord}}(U)$ is the set of elements of $\mathcal F^{\mathrm{ord}}(V)$ fixed by $U$, as $V\subset U$.

\emph{The action.} For $u\in G^{(\ell)}_0$ the map $f\mapsto u\cdot f$ preserves $\Lambda$-adic forms: it translates on the right in the $G^{(\ell)}$-component and multiplies the values by the matrix $\tau_0(u)$ over $\mathcal O$, while the conditions of Situation~\ref{sha-hecke-ideal-tw-augmented-setup}(2) involve left multiplication by $U(n)(L^+)$, right multiplication in $G_\ell$ and congruences modulo $\mathfrak m_\Lambda^k$. For the same reason it commutes with $U_z$ and maps functions with values in $\mathfrak m_\Lambda^k\Lambda^m$ to such functions; so it preserves ordinary forms. For $u\in U$ and $f\in\mathcal F^{\mathrm{ord}}(V)$, $u\cdot f$ is fixed by $V$, as $V$ is normal in $U$. So $U/V$ acts on $\mathcal F^{\mathrm{ord}}(V)$.

\emph{A fundamental domain.} Call $x,y\in G^{(\ell)}\times G_\ell$ equivalent if $y=\gamma x(u,b)$ with $\gamma\in U(n)(L^+)$, $u\in U$ and $b\in I_\ell$, and choose a representative $R(x)$ of the class of each $x$. By neatness, if $\gamma R(x)(u,b)=\gamma'R(x)(u',b')$ with $u,u'\in U$ and $b,b'\in I_\ell$, then $u=u'$. Let $\chi(x)=1$ if $x=\gamma R(x)(u,b)$ with $u\in V$, and $\chi(x)=0$ otherwise. Then $\chi(\gamma x(u,b))=\chi(x)$ for $\gamma\in U(n)(L^+)$, $u\in V$ and $b\in I_\ell$. For every $x$, $\sum_{q\in U/V}\chi(x\tilde q)=1$: write $x=\gamma_0R(x)(u_0,b_0)$ with $u_0\in U$; by the uniqueness of $u$, $\chi(x\tilde q)=1$ if and only if $u_0\tilde q\in V$, which holds for the single class $q=u_0^{-1}V$.

\emph{The map $\varphi$.} Since $\chi$ is invariant under $U(n)(L^+)$ on the left and under $V$ and $I_\ell\supset N_0T_0\bar N_c$ on the right, multiplication by $\chi$ maps $\mathcal F(V)$ into itself. Let $\varphi=e_V\circ\chi:\mathcal F^{\mathrm{ord}}(V)\to\mathcal F^{\mathrm{ord}}(V)$. For $f\in\mathcal F^{\mathrm{ord}}(V)$ and $q\in U/V$, $(\tilde q\cdot(\chi\,(\tilde q^{-1}\cdot f)))(x)=\chi(x\tilde q)f(x)$, so $\sum_q\tilde q\cdot(\chi\,(\tilde q^{-1}\cdot f))=f$. Fix $k$ and choose $N$ so large that, with $n=N!$, $U_z^nh-e_Vh$ has values in $\mathfrak m_\Lambda^k\Lambda^m$ for $h=f$ and for the $h=\chi\,(\tilde q^{-1}\cdot f)$; here $e_Vf=f$. Since the action of $U$ commutes with $U_z$ and preserves values in $\mathfrak m_\Lambda^k\Lambda^m$,
\[
\sum_q\tilde q\cdot\varphi(\tilde q^{-1}\cdot f)-f\equiv\sum_q\tilde q\cdot U_z^n\big(\chi\,(\tilde q^{-1}\cdot f)\big)-U_z^nf=U_z^n\Big(\sum_q\tilde q\cdot\big(\chi\,(\tilde q^{-1}\cdot f)\big)\Big)-U_z^nf=0
\]
modulo functions with values in $\mathfrak m_\Lambda^k\Lambda^m$. As $\bigcap_k\mathfrak m_\Lambda^k=0$, $\sum_qq\cdot\varphi(q^{-1}\cdot f)=f$.

\emph{Projectivity.} Let $F=\mathcal F^{\mathrm{ord}}(V)$. The map $s:F\to\Lambda[U/V]\otimes_\Lambda F$, $s(f)=\sum_qq\otimes\varphi(q^{-1}\cdot f)$, is $\Lambda[U/V]$-linear, since replacing $f$ by $q'\cdot f$ permutes the terms, and it is a section of the multiplication map $\Lambda[U/V]\otimes_\Lambda F\to F$ by the last identity. The module $\Lambda[U/V]\otimes_\Lambda F$ is projective over $\Lambda[U/V]$, since $F$ is free over $\Lambda$. So $F$, a direct summand of it, is projective over $\Lambda[U/V]$. This proves (1).

\emph{The trace.} For $f\in F$ the sum $\operatorname{tr}f$ does not depend on the representatives, as $f$ is fixed by $V$, and $q'\cdot\operatorname{tr}f=\operatorname{tr}f=\operatorname{tr}(q'\cdot f)$ for $q'\in U/V$, since multiplication by $q'$ on either side permutes $U/V$. So $\operatorname{tr}f$ is an element of $F$ fixed by $U$, that is, of $\mathcal F^{\mathrm{ord}}(U)$. Conversely, for $f\in\mathcal F^{\mathrm{ord}}(U)$, $q^{-1}\cdot f=f$ for every $q$, so the identity $\sum_qq\cdot\varphi(q^{-1}\cdot f)=f$ reads $\operatorname{tr}\varphi(f)=f$. This proves (2). For (3), $\operatorname{tr}(u\cdot f'-f')=0$. Conversely, let $\operatorname{tr}f=0$. Then $\sum_q\varphi(q^{-1}\cdot f)=\varphi(\operatorname{tr}f)=0$, as $q\mapsto q^{-1}$ permutes $U/V$, so
\[
f=\sum_q\big(\tilde q\cdot\varphi(q^{-1}\cdot f)-\varphi(q^{-1}\cdot f)\big),
\]
a sum of elements $u\cdot f'-f'$.
% ordProj: the level V is ordLevel (the representatives x (q.out, 1), the freeness condition inherited from U, Ord U = Ord V meet the U-fixed vectors) after hFU (the freeness condition from Neat U), with ordFree for the freeness of Ord V; ordProj_sp and ordProj_ord (spStab, ordProj_pow, ordProj_valra) for the action; ordProj_chi (the setoid ordProj_rel, R = Quotient.out, uniq from Neat); ordProj_mul; ordProj_phi (e from ordStruct, the vanishing by Ideal.iInf_pow_eq_bot_of_isLocalRing); ordProj_higman (Higman's criterion); ordProj_dc for the trace as the double coset operator dc (ra tau) U V 1 on V-fixed vectors, then htrinv, htrinvL and hsumw for (2) and (3).
\end{proof}

\begin{lemma}\label{sha-hecke-ideal-tw-augmented-sum}
In Situation~\ref{sha-hecke-ideal-tw-augmented-setup}:
\begin{enumerate}
\item if $K'\subset K$, then $M_K\subset M_{K'}$, and every element of $P$ lies in some $M_K$;
\item $\rho(g)$ maps $P$ into itself for every $g\in G_w$, so $\rho$ is a representation of $G_w$ on $P$;
\item for every $K$, the vectors of $P$ fixed by $\prod_jK_j$ are the elements of $M_K$;
\item if $K'\subset K$ and $\prod_jK'_j$ is normal in $\prod_jK_j$, then $M_{K'}$, with the action of $\prod_jK_j/\prod_jK'_j$ through $\rho$, is projective over $\Lambda[\prod_jK_j/\prod_jK'_j]$.
\end{enumerate}
\end{lemma}
\leanref{Development/AutomorphicArithmetic/AutomorphicVarieties/HeckeAlgebras/OrdinaryHeckeCharacters/PiBasic.lean}{UpperTheory.V.Hi.piBasic}
\begin{proof}
Write $K^\Pi=\prod_jK_j\subset G_w$, and $\mathcal K_0^w$ for the subgroup of elements of $\mathcal K_0K_w$ whose components at the $v_j$ are $1$; so $\mathcal K_0K=\mathcal K_0^wK^\Pi$, and the elements of $\mathcal K_0^w$ commute with $G_w$. For $g\in G_w$, $\rho(g)$ commutes with the operators: with the $T_{w,j}$, which are sums of translations by elements of $GL_n(L_w)$, commuting with $G_w$, and with the $U_{v,j}$, which act in the variable in $G_\ell$. Also, $u\cdot f$ is an ordinary $\Lambda$-adic form for every $u\in G^{(\ell)}_0$ and every ordinary $\Lambda$-adic form $f$, as in the proof of Lemma~\ref{sha-hecke-ideal-tw-augmented-projective}; so if $f\in\mathcal F^{\mathrm{ord}}(U)$ and $u\cdot f$ is fixed by a subgroup $U'$, then $u\cdot f\in\mathcal F^{\mathrm{ord}}(U')$.

(1) If $K'\subset K$, then $\mathcal K_0K'\subset\mathcal K_0K$, so $\mathcal F^{\mathrm{ord}}(\mathcal K_0K)\subset\mathcal F^{\mathrm{ord}}(\mathcal K_0K')$, and $\mathfrak m_\Lambda^k\mathcal F^{\mathrm{ord}}(\mathcal K_0K)\subset\mathfrak m_\Lambda^k\mathcal F^{\mathrm{ord}}(\mathcal K_0K')$; hence $M_K\subset M_{K'}$. As $K\cap K'$ is contained in $K$ and in $K'$, the $M_K$ form a directed family, and their sum $P$ is their union.

(2) Let $g\in G_w$, $f\in M_K$, and let $K'$ be the family $(g_jK_jg_j^{-1})$, so that $g^{-1}K'^\Pi g=K^\Pi$. The form $\rho(g)f$ is fixed by $\mathcal K_0^w$, which commutes with $g$ and fixes $f$, and by every $k\in K'^\Pi$, since $\rho(k)\rho(g)f=\rho(g)\rho(g^{-1}kg)f=\rho(g)f$. So $\rho(g)f\in\mathcal F^{\mathrm{ord}}(\mathcal K_0K')$. Since $\rho(g)$ commutes with the operators, $(t-a_t)^N\rho(g)f=\rho(g)(t-a_t)^Nf$ for an operator $t$ with eigenvalue $a_t$, and this lies in $\rho(g)\mathfrak m_\Lambda^k\mathcal F^{\mathrm{ord}}(\mathcal K_0K)\subset\mathfrak m_\Lambda^k\mathcal F^{\mathrm{ord}}(\mathcal K_0K')$ for a suitable $N$; so $\rho(g)f\in M_{K'}\subset P$.

(3) The elements of $M_K$ are fixed by $\mathcal K_0K\supset K^\Pi$. Conversely, let $f\in P$ be fixed by $K^\Pi$. By (1), $f\in M_{K'}$ for some $K'$, and $f\in M_{K''}$ for $K''=K\cap K'$. So $f\in\mathcal F^{\mathrm{ord}}(\mathcal K_0K'')$, and $f$ is fixed by $\mathcal K_0^w\subset\mathcal K_0K''$ and by $K^\Pi$, hence by $\mathcal K_0K$; so $f\in Y=\mathcal F^{\mathrm{ord}}(\mathcal K_0K)$. Let $Y''=\mathcal F^{\mathrm{ord}}(\mathcal K_0K'')\supset Y$, which is finite over $\Lambda$ (Lemma~\ref{sha-hecke-ideal-tw-augmented-hida}(2)); the operators map $Y$ and $Y''$ into themselves. By the Artin-Rees lemma there is $c$ with $\mathfrak m_\Lambda^{k+c}Y''\cap Y\subset\mathfrak m_\Lambda^kY$ for all $k$. For an operator $t$ and $k\ge0$ there is $N$ with $(t-a_t)^Nf\in\mathfrak m_\Lambda^{k+c}Y''$, as $f\in M_{K''}$; this element also lies in $Y$, hence in $\mathfrak m_\Lambda^kY$. So $f\in M_K$.

(4) Let $U=\mathcal K_0K$ and $V=\mathcal K_0K'$. Then $V$ is normal in $U=\mathcal K_0^wK^\Pi$, as $K'^\Pi$ is normal in $K^\Pi$ and $\mathcal K_0^w$ commutes with $G_w$. The map $K^\Pi/K'^\Pi\to U/V$ is surjective, since $U=\mathcal K_0^wK^\Pi$, and injective, since an element of $K^\Pi$ that lies in $V$ has its components in the $K'_j$. The group $K^\Pi$ is compact and $K'^\Pi$ open, so these quotients are finite. By Lemma~\ref{sha-hecke-ideal-tw-augmented-hida}(1), $U$ is neat and of finite level. By Lemma~\ref{sha-hecke-ideal-tw-augmented-projective}(1), $\mathcal F^{\mathrm{ord}}(V)$ is projective over $\Lambda[U/V]=\Lambda[K^\Pi/K'^\Pi]$, the action being $\rho$. Let $e$ be an idempotent as in Lemma~\ref{sha-hecke-ideal-tw-augmented-idempotent}(1) for $\mathcal F^{\mathrm{ord}}(V)$ and the operators. By (3) of that lemma, $e\mathcal F^{\mathrm{ord}}(V)=M_{K'}$, and $e$ commutes with the $\rho(k)$, $k\in K^\Pi$, which commute with the operators. So the inclusion $M_{K'}\to\mathcal F^{\mathrm{ord}}(V)$ and $e:\mathcal F^{\mathrm{ord}}(V)\to M_{K'}$ are maps of $\Lambda[K^\Pi/K'^\Pi]$-modules whose composite is the identity of $M_{K'}$, and $M_{K'}$ is projective.
% piBasic: pb_J (J the product of the sections), hlev (lev K = UQ sup J(pi K)), hop (pb_dc with sectionsCommuteAtSplit for the T_{w,j}, pb_commU for the U_{v,j}), then the private pb_main: hMle (pb_locMap with phi = 1), hdir, hact (pb_ordMap with spStab, pb_mem), hst (pb_conj, pb_locMap), h1 (pb_locInf: Artin-Rees as Ideal.exists_pow_inf_eq_pow_smul), h2, and for the projectivity pb_normal, pb_equiv (from the faithfulness hypothesis, supplied by levelUR2Faith in s3), pb_idx', ordProj, locIdem and pb_proj. The neatness hypothesis of piBasic is neatU, supplied in umodAbs_alt.
\end{proof}

\begin{lemma}\label{sha-hecke-ideal-tw-augmented-unique}
In Situation~\ref{sha-hecke-situation}, let $\theta$ be as in Lemma~\ref{sha-hecke-data}, let $S\subset S_1\subset S_2$ be finite sets of places of $L^+$, $\Gamma_{i,0}$ the subgroup fixing $L$ of the Galois group of the maximal extension of $L^+$ unramified outside $S_i$ and the infinite places, and $p:\Gamma_{2,0}\to\Gamma_{1,0}$ the projection. For $i=1,2$ let $T_i$ be a complete local noetherian $\mathcal O$-algebra with residue field $k$ and an $\mathcal O$-algebra map $\Lambda\to T_i$, with elements $t^{(i)}_{w,j}$ and $u^{(i)}_{v,j}$, and $(D_i,\psi_i)$ a pair over $T_i$ for $\Gamma_{i,0}$ attached to them (Definition~\ref{sha-hecke-ideal-tw-augmented-parts}(3)). Let $\phi:T_2\to T_1$ be a map of $\Lambda$-algebras with $\phi(t^{(2)}_{w,j})=t^{(1)}_{w,j}$ for the split places $w$ above places outside $S_2$ and $\phi(u^{(2)}_{v,j})=u^{(1)}_{v,j}$. Then $\phi(D_2)=D_1\circ p$ and $\phi\circ\psi_{2,v,i}=\psi_{1,v,i}$ for all $v\mid\ell$ and $i$.
\end{lemma}
\leanref{Development/AutomorphicArithmetic/AutomorphicVarieties/ModularityArgument/ThresholdArguments/DenseThresholdCases/LevelUR2BaseEq.lean}{UpperTheory.V.Hi.levelUR2BaseEq_alt}
\leanref{Development/AutomorphicArithmetic/AutomorphicVarieties/ModularityArgument/ThresholdArguments/DenseThresholdCases/TCoreUniq.lean}{UpperTheory.V.Hi.core_unique_of_axioms_weight_char}
\begin{proof}
\emph{The determinants.} The map $\phi$ is a local homomorphism, as both residue fields are $k$ and $\phi$ is $\mathcal O$-linear, so it is continuous for the adic topologies, and the coefficients of the characteristic polynomials of $\phi(D_2)$ and of $D_1\circ p$ are continuous functions on $\Gamma_{2,0}$ with values in the Hausdorff ring $T_1$. Let $w$ be a split place of $L$ above a place outside $S_2$ and $\sigma$ the inverse of a Frobenius element at $w$ in $\Gamma_{2,0}$; then $p(\sigma)$ is such an element in $\Gamma_{1,0}$. By (1) of the attachment, the characteristic polynomial of $\phi(D_2)$ at $\sigma$ is the image under $\phi$ of the Hecke polynomial with coefficients $t^{(2)}_{w,j}$, which is the Hecke polynomial with coefficients $t^{(1)}_{w,j}$, the characteristic polynomial of $D_1$ at $p(\sigma)$. These $\sigma$ are dense in $\Gamma_{2,0}$ by the Chebotarev density theorem. So the characteristic polynomials of $\phi(D_2)$ and $D_1\circ p$ agree at every element of $\Gamma_{2,0}$, and the two determinants are equal by Amitsur's formula \cite{Che14}, as in the proof of Lemma~\ref{sha-hecke-glue}.

\emph{The characters.} Fix $v\mid\ell$, let $W_v\subset D_v$ be the Weil group and $a:W_v\to L_{\tilde v}^\times$ a local Artin map, which is surjective, and let $s_0\in W_v$ with $a(s_0)$ a uniformizer. Put $\psi'_i=\phi\circ\psi_{2,v,i}$. For $s\in W_v$ with $a(s)$ a unit, (2) of the attachment, transported by $\phi$, which is the identity on the images of $\Lambda$ and of $\mathcal O$, gives $\psi'_i(s)\varepsilon(s)^{i-1}\theta(e_{v,i}(a(s)))=1$, and the same relation holds for $\psi_{1,v,i}$; so $\psi'_i(s)=\psi_{1,v,i}(s)$. At $s_0$, the products $\prod_{i\le j}\psi'_i\varepsilon^{i-1}$ and $\prod_{i\le j}\psi_{1,v,i}\varepsilon^{i-1}$ are $\phi(u^{(2)}_{v,j})=u^{(1)}_{v,j}$ for every $j$, so by induction on $j$, $\psi'_j(s_0)=\psi_{1,v,j}(s_0)$. Every $s\in W_v$ is $s's_0^k$ with $a(s')$ a unit, so $\psi'_i=\psi_{1,v,i}$ on $W_v$. Both are continuous and $W_v$ is dense in $D_v$, so they agree on $D_v$.
% levelUR2BaseEq_alt: the private abs_det (cu_ext for the density, push_cont, CA.d_gext for the determinant from the coefficients at group elements; the density is the field dense of HSet.Ax, with p the level-change map p0 h) and abs_chr (cu_ext with weil_group_dense_in_galois_group, cu_gen with the surjectivity of the Artin map from HSet.Ax.art_spec, cu_prod for the induction on j). core_unique_of_axioms_weight_char is the case S_1 = S_2 with phi the identity, by the same cu_ext, cu_gen and cu_prod.
\end{proof}

\begin{lemma}\label{sha-hecke-ideal-tw-augmented-compatible}
In Situation~\ref{sha-hecke-ideal-tw-augmented-setup}, let $B'$ be a universal ring of ordinary pairs for $\Gamma'$, with universal pair $(D',\psi')$. Let $K$ and $K'$ be families with $M_K\ne0$ and $M_{K'}\ne0$, and let $f_K:B'\to\mathbf T_K$ and $f_{K'}:B'\to\mathbf T_{K'}$ be $\mathcal O$-algebra maps such that $(f_K(D'),f_K\circ\psi')$ is attached to the operators $T_{w,j}$ and $U_{v,j}$ of $\mathbf T_K$, and similarly for $K'$. Let $\varphi:M_K\to M_{K'}$ be an injective $\Lambda$-linear map with $\varphi(ty)=t\varphi(y)$ for every operator $t$ and $y\in M_K$. Then $\varphi(f_K(b)y)=f_{K'}(b)\varphi(y)$ for $b\in B'$ and $y\in M_K$.
\end{lemma}
\leanref{Development/AutomorphicArithmetic/AutomorphicVarieties/ModularityArgument/ThresholdArguments/DenseThresholdCases/LevelUR2Compat.lean}{UpperTheory.V.Hi.level_change_compat_rank_two_of_dense_weight}
\leanref{Development/AutomorphicArithmetic/AutomorphicVarieties/HeckeAlgebras/OrdinaryHeckeCharacters/LevelUR2HeckeMap.lean}{UpperTheory.V.Hi.levelUR2HeckeMap}
\begin{proof}
\emph{The restriction map.} For every $s\in\mathbf T_{K'}$ there is a unique $\Phi(s)\in\mathbf T_K$ with $\varphi(\Phi(s)y)=s\varphi(y)$ for all $y\in M_K$. Uniqueness holds because $\varphi$ is injective. Existence holds for the operators, with $\Phi(t)=t$, and for the elements of $\Lambda$, and it is preserved by sums and products; since $\mathbf T_{K'}$ is generated over $\Lambda$ by the operators, it holds for every $s$. So $\Phi:\mathbf T_{K'}\to\mathbf T_K$ is a map of $\Lambda$-algebras with $\Phi(t)=t$ for the operators.

\emph{The two maps from $B'$.} The rings $\mathbf T_K$ and $\mathbf T_{K'}$ are complete local noetherian $\mathcal O$-algebras with residue field $k$ (Lemma~\ref{sha-hecke-algebra}, for $S'$ and these levels). By Lemma~\ref{sha-hecke-ideal-tw-augmented-unique}, with $S_1=S_2=S'$, $p$ the identity and $\phi=\Phi$, $\Phi(f_{K'}(D'))=f_K(D')$ and $\Phi\circ f_{K'}\circ\psi'=f_K\circ\psi'$. So $\Phi\circ f_{K'}$ and $f_K$ carry $(D',\psi')$ to the same pair, and by the uniqueness in the universal property of $B'$, $f_K=\Phi\circ f_{K'}$. Hence $\varphi(f_K(b)y)=\varphi(\Phi(f_{K'}(b))y)=f_{K'}(b)\varphi(y)$.
% level_change_compat_rank_two_of_dense_weight: levelUR2HeckeMap (with j the identity), the private lc_univ (transport of IsUniv along Phi), core_unique_of_axioms_weight_char, and the private lc_comp with the uniqueness hU of maps out of B' (base_object_hom_ext B').
\end{proof}

\begin{lemma}\label{sha-hecke-ideal-tw-augmented-action}
In Situation~\ref{sha-hecke-ideal-tw-augmented-setup}, let $B'$ be a universal ring of ordinary pairs for $\Gamma'$, with universal pair $(D',\psi')$, and $\lambda':\Lambda\to B'$ the $\mathcal O$-algebra map that carries the universal characters $\kappa_{v,i}$ of the $I_v$ to the restrictions of the $\psi'_{v,i}$. Then there is a ring map $\alpha:B'\to\operatorname{End}_\Lambda(P)$ such that $\alpha(\lambda'(b))$ is multiplication by $b$ for $b\in\Lambda$, $\alpha(t)$ commutes with $\rho(g)$ for $t\in B'$ and $g\in G_w$, and for every family $K$ with $M_K\ne0$ there is an $\mathcal O$-algebra map $f_K:B'\to\mathbf T_K$ such that $(f_K(D'),f_K\circ\psi')$ is attached to the operators of $\mathbf T_K$ and $\alpha(t)y=f_K(t)y$ for $t\in B'$ and $y\in M_K$.
\end{lemma}
\leanref{Development/AutomorphicArithmetic/AutomorphicVarieties/ModularityArgument/ThresholdArguments/DenseThresholdCases/LevelUR2ActEx.lean}{UpperTheory.V.Hi.levelUR2ActEx_alt}
\leanref{Development/AutomorphicArithmetic/AutomorphicVarieties/HeckeAlgebras/OrdinaryHeckeCharacters/LevelUR2Glue.lean}{UpperTheory.V.Hi.levelUR2Glue}
\begin{proof}
\emph{One level.} The data $\tau_0$ and $\theta$, and the eigenvalues $a_{w,j}$ at the split places $w$ above places outside $S'$ and $a_{v,j}$, are as in Lemma~\ref{sha-hecke-data} for $S'$ in place of $S$. Fix a family $K$. If $M_K=0$, let $\mu_K=0$. Otherwise the proof of Proposition~\ref{sha-hecke-ideal}, for $S'$, $\Gamma'$, $B'$ and the level $\mathcal K_0K$, gives an $\mathcal O$-algebra map $f_K:B'\to\mathbf T_K$, compatible with the maps from $\Lambda$, that carries $(D',\psi')$ to a pair attached to the operators of $\mathbf T_K$ (Lemma~\ref{sha-hecke-determinant}); let $\mu_K(t)$ be multiplication by $f_K(t)$ on $M_K$. In both cases $\mu_K:B'\to\operatorname{End}_\Lambda(M_K)$ is a ring map with $\mu_K(\lambda'(b))=b$.

\emph{Compatibility.} Let $\varphi:M_K\to M_{K'}$ be an injective $\Lambda$-linear map that commutes with the operators. Then $\varphi\circ\mu_K(t)=\mu_{K'}(t)\circ\varphi$ for every $t\in B'$: if $M_K=0$ both sides vanish; if $M_K\ne0$, then $M_{K'}\ne0$, and this is Lemma~\ref{sha-hecke-ideal-tw-augmented-compatible}.

\emph{Gluing.} For $K'\subset K$, $M_K\subset M_{K'}$ (Lemma~\ref{sha-hecke-ideal-tw-augmented-sum}(1)), and the inclusion commutes with the operators, so $\mu_{K'}(t)$ restricts to $\mu_K(t)$ on $M_K$. The $M_K$ form a directed family with union $P$. So there is a unique map $\alpha(t):P\to P$ that agrees with $\mu_K(t)$ on every $M_K$: on an element of $M_K\cap M_{K'}$ the two definitions agree with $\mu_{K\cap K'}(t)$, as $M_{K\cap K'}$ contains $M_K$ and $M_{K'}$. Each $\alpha(t)$ is $\Lambda$-linear, and $t\mapsto\alpha(t)$ is a ring map with $\alpha(\lambda'(b))=b$, since any finitely many elements of $P$ lie in one $M_K$.

\emph{Commutation with $G_w$.} Let $g\in G_w$, $z\in M_K$ and $K'=(g_jK_jg_j^{-1})$, as in the proof of Lemma~\ref{sha-hecke-ideal-tw-augmented-sum}(2). Then $\rho(g)$ maps $M_K$ injectively into $M_{K'}$ and commutes with the operators: for $g$ in one factor $GL_n(L_{w_j})$, because the translations commute with the $T_{w,j}$ and $U_{v,j}$, and for a product of such elements by composition. By the compatibility, $\alpha(t)\rho(g)z=\mu_{K'}(t)\rho(g)z=\rho(g)\mu_K(t)z=\rho(g)\alpha(t)z$.
% levelUR2ActEx_alt: hJ (for every K, from Hi.exists_hecke_ideal at the level V, supplied by LvW.step1_alt with theta' = theta, a' = a o opUp, lam = LamObj.toPair, type_for_r_up_of_type_for and res_up), the private ae_ex (mu from IsJ: ae_bot or ae_mu), ae_cmp (ae_zero, level_change_compat_rank_two_of_dense_weight), ae_abs (levelUR2Glue with ae_inf; the commutation with ae_conj and ae_gen, ae_op), ae_out.
\end{proof}

\begin{lemma}\label{sha-hecke-ideal-tw-augmented-eigensystems}
In Situation~\ref{sha-hecke-ideal-tw-augmented-setup}, let $\lambda$ be a weight, $x$ an arithmetic point of $\Lambda$ of weight $\lambda$, and $\xi$ and $\operatorname{sp}$ as in Lemma~\ref{sha-hecke-ideal-tw-augmented-hida}(4). Let $U=\mathcal K_0K$ be a level, $\mathbf T(U)$ the $\Lambda$-algebra of endomorphisms of $\mathcal F^{\mathrm{ord}}(U)$ generated by the operators $T_{w,j}$ and $U_{v,j}$, and $z:\mathbf T(U)\to\bar{\mathbf Q}_\ell$ a ring map equal to $x$ on $\Lambda$. Then there is an automorphic part $Y$ of weight $\xi$ with $a_{Y,w,j}=\iota z(T_{w,j})$ for all split places $w$ above places outside $S'$ and $1\le j\le n$.
\end{lemma}
\leanref{Development/AutomorphicArithmetic/AutomorphicVarieties/ModularityArgument/ThresholdArguments/DenseThresholdCases/LevelUR2OrdPt.lean}{UpperTheory.V.Hi.levelUR2OrdPt_alt}
\leanref{Development/AutomorphicArithmetic/AutomorphicVarieties/HeckeAlgebras/OrdinaryHeckeCharacters/LevelUR2EigPart.lean}{UpperTheory.V.Hi.levelUR2EigPart}
\begin{proof}
\emph{An eigenvector.} Let $F=\mathcal F^{\mathrm{ord}}(U)$, finite and free over $\Lambda$ with commuting operators (Lemma~\ref{sha-hecke-ideal-tw-augmented-hida}(2)), and view $\mathbf C$ as a $\Lambda$-algebra through $\iota x$. The map $\mathbf C\otimes_\Lambda F\to V^\xi$, $c\otimes f\mapsto c\operatorname{sp}(f)$, is well defined and injective, since $\operatorname{sp}$ is additive and $\Lambda$-semilinear on $F$ and the images of a basis are linearly independent. The module $F$ is faithful and finite over $\mathbf T(U)$, so $\mathbf C\otimes_{\mathbf T(U)}F\ne0$ for the map $\iota z:\mathbf T(U)\to\mathbf C$, by Nakayama's lemma. The quotient map $q:\mathbf C\otimes_\Lambda F\to\mathbf C\otimes_{\mathbf T(U)}F$ satisfies $q((1\otimes t)v)=\iota z(t)q(v)$. On a generalized eigenspace of the commuting operators $1\otimes t$ for a system of eigenvalues other than $(\iota z(t))_t$, some $1\otimes t-\iota z(t)$ is invertible, so $q$ vanishes there; hence the generalized eigenspace for $(\iota z(t))_t$ is nonzero, and it contains a common eigenvector $v\ne0$. Its image $G\ne0$ lies in the span of $\operatorname{sp}(F)$, and $t^\xi G=\iota z(t)G$ for every operator $t$.

\emph{The part.} By Lemma~\ref{sha-hecke-ideal-tw-augmented-hida}(4), $G$ lies in the sum of the automorphic parts $Y$ of weight $\xi$, on each of which $T^\xi_{w,j}$ acts by $a_{Y,w,j}$. Write $G$ as a finite sum $\sum_{s=1}^rG_s$ with $G_s\in Y_s$, and argue by induction on $r$. If $a_{Y_1,w,j}=\iota z(T_{w,j})$ for all $w,j$, take $Y=Y_1$. Otherwise choose $(w,j)$ with $c=a_{Y_1,w,j}-\iota z(T_{w,j})\ne0$; then $cG_1=\iota z(T_{w,j})G'-T^\xi_{w,j}G'$ with $G'=\sum_{s\ge2}G_s$, which lies in $Y_2+\dots+Y_r$, as $T^\xi_{w,j}$ preserves each $Y_s$. So $G=G_1+G'$ is a nonzero eigenvector lying in $Y_2+\dots+Y_r$, and the induction applies.
% levelUR2OrdPt_alt: the private op_eig (liftBaseChange injective from IsSpec.indep; AnnihilatorQuotientFinite.pointEig, which is fibreNe and eigOfQuot), IsSpec.span and clParts_alt for the decomposition of OrdAt, and levelUR2EigPart.
\end{proof}

\begin{lemma}\label{sha-hecke-ideal-tw-augmented-base}
In Situation~\ref{sha-hecke-ideal-tw-augmented-setup}, let $B'$, $(D',\psi')$, $\alpha$ and the $f_K$ be as in Lemma~\ref{sha-hecke-ideal-tw-augmented-action}, $\beta:B'\to B$ the $\mathcal O$-algebra map carrying $(D',\psi')$ to $(D^{\mathrm u}\circ p,\psi^{\mathrm u})$, and $M$ the module of Lemma~\ref{sha-hecke-algebra}. Let $\mu:B\to\operatorname{End}_\Lambda(M)$ be $0$ if $M=0$, and otherwise the composite of $\mathbf T\subset\operatorname{End}_\Lambda(M)$ with an $\mathcal O$-algebra map $f:B\to\mathbf T$ carrying $(D^{\mathrm u},\psi^{\mathrm u})$ to the pair $(D_{\mathbf T},\psi_{\mathbf T})$ of Lemma~\ref{sha-hecke-determinant}. Then:
\begin{enumerate}
\item $M_{K_w}=M$;
\item $\alpha(t)y=\mu(\beta(t))y$ for $t\in B'$ and $y\in M$.
\end{enumerate}
\end{lemma}
\leanref{Development/AutomorphicArithmetic/AutomorphicVarieties/ModularityArgument/ThresholdArguments/DenseThresholdCases/UmodAbs.lean}{UpperTheory.V.Hi.LvW.step2_alt}
\leanref{Development/AutomorphicArithmetic/AutomorphicVarieties/ModularityArgument/ThresholdArguments/DenseThresholdCases/LevelUR2LocEq.lean}{UpperTheory.V.Hi.levelUR2LocEq_alt}
\leanref{Development/AutomorphicArithmetic/AutomorphicVarieties/HeckeAlgebras/OrdinaryHeckeCharacters/LevelUR2LevEq.lean}{UpperTheory.V.Hi.levelUR2LevEq}
\leanref{Development/AutomorphicArithmetic/AutomorphicVarieties/ModularityArgument/ThresholdArguments/DenseThresholdCases/ModuleStructureOfCommuting.lean}{UpperTheory.V.Hi.levelUR2H3Abs_alt}
\leanref{Development/AutomorphicArithmetic/AutomorphicVarieties/HeckeAlgebras/OrdinaryHeckeCharacters/GalResContCoeff.lean}{UpperTheory.V.Hi.levelUR2GalRes}
\leanref{Development/AutomorphicArithmetic/AutomorphicVarieties/ModularityArgument/ThresholdArguments/DenseThresholdCases/LevelUR2SmulKey.lean}{UpperTheory.V.Hi.levelUR2SmulKey_alt}
\begin{proof}
\emph{The levels and the operators.} The level $\mathcal K_0K_w$ is $\mathcal K_0$: both consist of the elements of $G^{(\ell)}_0$ whose components at all places outside $S$ are integral. So $M_{K_w}$ and $M$ are localizations of the same module $Y=\mathcal F^{\mathrm{ord}}(\mathcal K_0)$: $M$ at the operators $T_{w,j}$ for the split places $w$ above places outside $S$ and the $U_{v,j}$, and $M_{K_w}$ at those of them with $w$ above a place outside $S'$, with the same eigenvalues.

(1) Apply Lemma~\ref{sha-hecke-ideal-tw-augmented-fewer} to $A=\Lambda$, which is a domain (Lemma~\ref{sha-hecke-algebra}), to $Y$, which is finite and free with commuting operators (Lemma~\ref{sha-hecke-ideal-tw-augmented-hida}(2), for $S$ in place of $S'$), to the operators at level $S$ and the subset of those at level $S'$, and to a point $x$ of $\Lambda$ of distinct type of some weight $\lambda$; such a point exists, since by the proof of Lemma~\ref{sha-hecke-algebra} an element of $\Lambda$ that vanishes at all points of distinct type of a weight is $0$, and $1\ne0$. Let $z:\mathbf T(\mathcal K_0)\to\bar{\mathbf Q}_\ell$ be a ring map, equal to $x$ on $\Lambda$, with $|z(t)-a_t|<1$ for the operators $t$ at level $S'$ and their eigenvalues $a_t$. By Lemma~\ref{sha-hecke-ideal-tw-augmented-eigensystems}, for $S$ in place of $S'$, there is an automorphic part $Y_1$ with $a_{Y_1,w,j}=\iota z(T_{w,j})$ for all split $w$ above places outside $S$. The coefficient of $X^{n-j}$ in the Hecke polynomial at $w$ with coefficients $b_1,\dots,b_n$ is $(-1)^j\operatorname{N}w^{j(j-1)/2}b_j$, and $\operatorname{N}w$ is an $\ell$-adic unit. So, by Lemma~\ref{sha-hecke-data}(3), at the inverse $\sigma$ of a Frobenius element at a split $w$ above a place outside $S'$, the coefficients of the characteristic polynomial of $\rho_{Y_1}(\sigma)$ are integral and reduce to those of $\bar D$ at $\sigma$. These $\sigma$ are dense in $\Gamma_0$, as the image of the dense set of such elements of $\Gamma'_0$. The characteristic polynomial of $\bar D$ is a locally constant function on $\Gamma_0$, and the coefficients of the characteristic polynomial of $\rho_{Y_1}(g)$ are continuous in $g$; so for every $g\in\Gamma_0$ there is such a $\sigma$ at which $\bar D$ has the same characteristic polynomial as at $g$ and the coefficients for $\rho_{Y_1}$ are within distance less than $1$ of those at $g$. Hence the coefficients of the characteristic polynomial of $\rho_{Y_1}(g)$ are integral and reduce to those of $\bar D$ at $g$, for every $g$. Taking for $g$ the inverse of a Frobenius element at any split $w$ above a place outside $S$, and using the coefficient formula again, $|z(T_{w,j})-a_{w,j}|<1$. The $U_{v,j}$ are operators at level $S'$. So Lemma~\ref{sha-hecke-ideal-tw-augmented-fewer} applies, and $M_{K_w}=M$.

(2) If $M=0$, both sides vanish. Otherwise $\mathbf T_{K_w}$ and $\mathbf T$ are the algebras of endomorphisms of $M$ generated by the operators at levels $S'$ and $S$, so $\mathbf T_{K_w}\subset\mathbf T$, and the inclusion $\phi$ fixes the operators at level $S'$. By Lemma~\ref{sha-hecke-ideal-tw-augmented-unique}, with $S_1=S$, $S_2=S'$ and the projection $p:\Gamma'_0\to\Gamma_0$, $\phi(f_{K_w}(D'))=D_{\mathbf T}\circ p$ and $\phi\circ f_{K_w}\circ\psi'=\psi_{\mathbf T}$. Also $f(\beta(D'))=f(D^{\mathrm u})\circ p=D_{\mathbf T}\circ p$ and $f\circ\beta\circ\psi'=f\circ\psi^{\mathrm u}=\psi_{\mathbf T}$. By the uniqueness in the universal property of $B'$, $\phi\circ f_{K_w}=f\circ\beta$. Hence, for $y\in M=M_{K_w}$, $\alpha(t)y=f_{K_w}(t)y=f(\beta(t))y=\mu(\beta(t))y$.
% (1) is LvW.step2_alt = levelUR2LocEq_alt: hop (the operators at V are those at V0 through opUp), levelUR2LevEq (lev K = lev K0 when K is KO at every w_q, through levelUR2LevOut and mem_heckeLevel_of_local_components_mem), and levelUR2H3Abs_alt: hiComm, a good point from DenseWt (isDomain_and_denseWt_of_isTheta), spec, levelUR2LocSub, levelUR2OrdPt_alt for the HSet at V0, levelUR2GalRes with phi = lq_lc r (the residual characteristic polynomial; IsRes.frob for its Frobenius values) and lq_dense, and the private hecke_part_frobenius (hecke_poly_coeff_formula). (2) is levelUR2SmulKey_alt: sk_zero, sk_hecke (levelUR2HeckeMap for the inclusion), levelUR2BaseEq_alt, sk_comp with base_object_hom_ext B'.
\end{proof}

\begin{lemma}\label{sha-hecke-ideal-tw-augmented-points}
In Situation~\ref{sha-hecke-ideal-tw-augmented-setup}, let $B'$, $(D',\psi')$ and $\alpha$ be as in Lemma~\ref{sha-hecke-ideal-tw-augmented-action}. Then for every nonzero $z\in P$ there are a ring map $\theta_z:B'\to\bar{\mathbf Q}_\ell$, a continuous homomorphism $\rho_z:\Gamma'_0\to GL_n(\bar{\mathbf Q}_\ell)$, a representation $\pi_z$ of $G_w$ on a complex vector space $V_z$ with an additive map $\operatorname{sp}_z:P\to V_z$, and irreducible smooth representations $\sigma_{z,j}$ of the $GL_n(L_{w_j})$, which satisfy the conditions on a point in Definition~\ref{sha-hecke-ideal-tw-level-module}(2) for $R=B'$, the pair $(D',\psi')$, $P$ with the action of $B'$ through $\alpha$, $\rho$ and $\iota$, and such that $\operatorname{sp}_z(z)\ne0$.
\end{lemma}
\leanref{Development/AutomorphicArithmetic/AutomorphicVarieties/ModularityArgument/ThresholdArguments/DenseThresholdCases/VLPiPts.lean}{UpperTheory.V.Hi.piPts_n_alt}
\begin{proof}
\emph{A point of $\Lambda$.} By Lemma~\ref{sha-hecke-ideal-tw-augmented-sum}(1), $z\in M_{K_1}$ for some family $K_1$, so $z\in\mathcal F^{\mathrm{ord}}(\mathcal K_0K_1)$, which is free over $\Lambda$ (Lemma~\ref{sha-hecke-ideal-tw-augmented-hida}(2)). Choose a basis; some coordinate $c\in\Lambda$ of $z$ is nonzero. By the proof of Lemma~\ref{sha-hecke-algebra}, for a weight $\lambda$ there is a point $x$ of $\Lambda$ of distinct type of weight $\lambda$ with $x(c)\ne0$. Let $\xi$ and $\operatorname{sp}$ be as in Lemma~\ref{sha-hecke-ideal-tw-augmented-hida}(4) for $x$. Then $\operatorname{sp}(z)\ne0$: otherwise the linear independence of the images of the basis gives $\iota x(c)=0$.

\emph{The map $s$.} Any finitely many elements of $P$ lie in one $M_K$ (Lemma~\ref{sha-hecke-ideal-tw-augmented-sum}(1)), so the restriction $s$ of $\operatorname{sp}$ to $P$ is additive. Let $\pi=\pi_\xi$. Then $s(\rho(g)f)=\pi(g)s(f)$ for $g\in G_w$ and $f\in P$: for $g$ in one factor, by Lemma~\ref{sha-hecke-ideal-tw-augmented-hida}(4) applied to a level $\mathcal K_0K$ with $f,\rho(g)f\in M_K$, and in general by composition.

\emph{The levels containing $z$.} Let $K$ be a family with $z\in M_K$, so $M_K\ne0$. Let $f_K$ be as in Lemma~\ref{sha-hecke-ideal-tw-augmented-action}, and $E_K$ and $t\mapsto t^x$ as in Lemma~\ref{sha-hecke-ideal-tw-augmented-hida}(5). For $t\in B'$ put $\mathrm{act}_K(t)=f_K(t)^x$; then $\mathrm{act}_K(t)s(f)=s(\alpha(t)f)$ for $f\in M_K$. Let $Y$ be an automorphic part with $Y\cap E_K\ne0$, and $y:\mathbf T_K\to\bar{\mathbf Q}_\ell$ as in Lemma~\ref{sha-hecke-ideal-tw-augmented-hida}(5); put $x_Y=y\circ f_K$. Then $\mathrm{act}_K(t)$ acts on $Y\cap E_K$ by $\iota x_Y(t)$. Moreover $x_Y(D')$ is the determinant of $\rho_Y$. Indeed, for a nonzero $F\in Y\cap E_K$, $T_{w,j}^xF=T^\xi_{w,j}F=a_{Y,w,j}F$, since $t^x$ and $t^\xi$ agree on $\operatorname{sp}(M_K)$, so $y(T_{w,j})=\iota^{-1}a_{Y,w,j}$. As $(f_K(D'),f_K\circ\psi')$ is attached to the operators, the characteristic polynomial of $x_Y(D')$ at the inverse of a Frobenius element at a split $w$ above a place outside $S'$ is the Hecke polynomial with coefficients $y(T_{w,j})$, which is that of $\rho_Y$ there. The map $y$ carries the maximal ideal of $\mathbf T_K$ into the maximal ideal of the ring of integers of $\bar{\mathbf Q}_\ell$, so the coefficients of the characteristic polynomials of $x_Y(D')$ are continuous; by the density of these Frobenius elements and Amitsur's formula \cite{Che14}, as in the proof of Lemma~\ref{sha-hecke-ideal-tw-augmented-unique}, $x_Y(D')$ is the determinant of $\rho_Y$.

\emph{The operators on the span.} Let $V_0$ be the span of $s(P)$. For $t\in B'$ there is a linear map $T_t:V^\xi\to V^\xi$ with $T_ts(f)=s(\alpha(t)f)$ for $f\in P$: every linear relation $\sum_rc_rs(f_r)=0$ gives $\sum_rc_rs(\alpha(t)f_r)=\mathrm{act}_K(t)\sum_rc_rs(f_r)=0$, for a family $K$ with $z$ and the $f_r$ in $M_K$; so $T_t$ is well defined on $V_0$, and it extends to $V^\xi$. The subspace $V_0$ is $\pi$-stable, since $s(\rho(g)f)=\pi(g)s(f)$, and $T_t\pi(g)=\pi(g)T_t$ on $V_0$, since $\alpha(t)$ commutes with $\rho(g)$ (Lemma~\ref{sha-hecke-ideal-tw-augmented-action}). Let $J$ be the set of triples $(Y,x',W)$ of an automorphic part $Y$, a ring map $x':B'\to\bar{\mathbf Q}_\ell$ with $x'(D')$ the determinant of $\rho_Y$, and a subspace $W\subset Y\cap V_0$ on which every $T_t$ acts by $\iota x'(t)$. Then $V_0=\sum_{(Y,x',W)\in J}W$: each $s(f)$ lies in $E_K$ for a family $K$ with $f,z\in M_K$, $E_K$ is the sum of the $Y\cap E_K$, and $(Y,x_Y,Y\cap E_K)\in J$, since $T_t$ agrees with $\mathrm{act}_K(t)$ on $E_K$.

\emph{The projection.} Lemma~\ref{sha-hecke-ideal-tw-augmented-projection}, applied to $V^\xi$, the $T_t$, the $\pi(g)$, $E=V_0$, the subspaces $W$ for $(Y,x',W)\in J$ with $\chi=\iota x'$, and $F_0=s(z)$, gives $(Y_0,x_0,W_0)\in J$, a $\pi$-stable subspace $Y^\circ$ of the sum of the $W$ for the triples $(Y,x',W)$ with $x'=x_0$, and a linear map $p:V^\xi\to Y^\circ$ with $p(s(z))\ne0$, $p(T_tF)=\iota x_0(t)p(F)$ and $p(\pi(g)F)=\pi(g)p(F)$ for $F\in V_0$.

\emph{Local components.} Fix $j$. Each automorphic part $Y$ with a triple $(Y,x_0,W)\in J$ is the sum of the images of the equivariant maps from $\sigma_{Y,j}$ under $GL_n(L_{w_j})$, so their sum is the sum of the images of the equivariant maps from these $\sigma_{Y,j}$, and by Lemma~\ref{sha-hecke-ideal-tw-augmented-isotypic}(1) so is its subrepresentation $Y^\circ$. By Lemma~\ref{sha-hecke-ideal-tw-augmented-isotypic}(2), applied to $Y^\circ$ with these $\sigma_{Y,j}$ and the vector $p(s(z))$, there are a $G_w$-equivariant surjection $p':Y^\circ\to V_z$ with $p'(p(s(z)))\ne0$ and, for each $j$, an automorphic part $Y_j$ with a triple $(Y_j,x_0,W)\in J$, such that $V_z$, as a representation of $GL_n(L_{w_j})$, is the sum of the images of the equivariant maps from $\sigma_{Y_j,j}$.

\emph{The point.} Put $\theta_z=x_0$, $\rho_z=\rho_{Y_0}$, $\pi_z$ the action of $G_w$ on $V_z$, $\operatorname{sp}_z=p'\circ p\circ s$ and $\sigma_{z,j}=\sigma_{Y_j,j}$. Then $\rho_z$ is continuous, $\theta_z(D')$ is its determinant, $\operatorname{sp}_z(z)\ne0$, $\operatorname{sp}_z(\alpha(t)f)=p'(p(T_ts(f)))=\iota\theta_z(t)\operatorname{sp}_z(f)$, and $\operatorname{sp}_z(\rho(g)f)=\pi_z(g)\operatorname{sp}_z(f)$. For each $j$, $\rho_{Y_0}$ and $\rho_{Y_j}$ are semisimple with the same determinant $x_0(D')$, hence the same characteristic polynomials, so they are conjugate by the Brauer-Nesbitt theorem \cite[(30.16)]{CR62}. Condition (d) of Definition~\ref{sha-hecke-ideal-tw-augmented-parts}(2) for $Y_j$ at $w_j$ therefore holds for $\rho_{Y_0}$ and $\sigma_{z,j}$, since conjugate representations have isomorphic Weil-Deligne representations. This is the condition on a point at $w_j$.
% piPts_n_alt: piPts_dir (finite sets in one MK, from piBasic), hiComm and Module.Free.chooseBasis, hden for x0, spec, piPts_coord; lev: hiAlg, partAct0_alt (act, a1, a2, the decomposition hE, a5), piPts_hecke, core_eigen_of_weight_char_hecke_level (first conjunct: eig_det with the density of HSet.Ax and eig_pt for the continuity) with piPts_mm; hsrho (IsSpec.act and Submonoid.pi_mem_of_mulSingle_mem); the private piPts_abs (PiLin.exists_desc for T, then PiLin.exists_proj); isoRefine with isoPick and isoSub for V_z; piPts_conj (AlgebraTools.Det.bnRep, woConj'_alt).
\end{proof}

\begin{lemma}\label{sha-hecke-ideal-tw-augmented}
Assume that Galois representations exist, in the sense of Definition~\ref{gal-galois-representations}, relative to the maps satisfying the laws of Definition~\ref{gal-llc-laws}. In Situation~\ref{sha-hecke-situation}, let $\iota:\bar{\mathbf Q}_\ell\cong\mathbf C$, let $\tau_0$, $\theta$ and $a$ be data as in Lemma~\ref{sha-hecke-data}, let $\mathcal K_0$ be a level and $M$ the module of Lemma~\ref{sha-hecke-algebra}, and let $\mu:B\to\operatorname{End}_\Lambda(M)$ be an $\mathcal O$-algebra map that is $0$ if $M=0$, and otherwise the composite of the inclusion $\mathbf T\subset\operatorname{End}_\Lambda(M)$ with an $\mathcal O$-algebra map $B\to\mathbf T$ carrying $(D^{\mathrm u},\psi^{\mathrm u})$ to the pair $(D_{\mathbf T},\psi_{\mathbf T})$ of Lemma~\ref{sha-hecke-determinant}. Let $v_1,\dots,v_q$ be pairwise distinct places of $L^+$ outside $S$ and split in $L$, $w_j$ a place of $L$ above $v_j$, $S'=S\cup\{v_1,\dots,v_q\}$, and $\Gamma'$, $p$, $\Gamma'_0$, $G_w$ and $K_w$ as in Definition~\ref{sha-hecke-ideal-tw-level-module}. Let $B'$ be a universal ring of ordinary pairs for $\Gamma'$, with universal pair $(D',\psi')$, and $\beta:B'\to B$ the $\mathcal O$-algebra map carrying $(D',\psi')$ to $(D^{\mathrm u}\circ p,\psi^{\mathrm u})$. Then there are a $\Lambda$-module $P$ with a $\Lambda$-linear representation $\rho$ of $G_w$, a ring map $\alpha:B'\to\operatorname{End}_\Lambda(P)$ whose values commute with $\rho$, and an $\mathcal O$-algebra map $\lambda':\Lambda\to B'$ that carries the universal characters $\kappa_{v,i}$ of the $I_v$ to the restrictions of the $\psi'_{v,i}$ and such that $\alpha(\lambda'(b))$ is multiplication by $b$, with the following properties.
\begin{enumerate}
\item For every compact open subgroup $K\subset G_w$ that is a product of subgroups of the factors, $P^K$ is finite over $\Lambda$.
\item For two such subgroups $K'\subset K$ with $K'$ normal in $K$, $P^{K'}$ is projective over the group ring $\Lambda[K/K']$.
\item There is an additive isomorphism $e_0:P^{K_w}\to M$ with $e_0(\alpha(t)m)=\mu(\beta(t))e_0(m)$ for $t\in B'$ and $m\in P^{K_w}$.
\item For every nonzero $z\in P$ there are a ring map $\theta_z:B'\to\bar{\mathbf Q}_\ell$, a continuous homomorphism $\rho_z:\Gamma'_0\to GL_n(\bar{\mathbf Q}_\ell)$, a representation $\pi_z$ of $G_w$ on a complex vector space $V_z$ with an additive map $\operatorname{sp}_z:P\to V_z$, and irreducible smooth representations $\sigma_{z,j}$ of the $GL_n(L_{w_j})$, which satisfy the conditions on a point in Definition~\ref{sha-hecke-ideal-tw-level-module}(2) for $R=B'$, the pair $(D',\psi')$, $P$ with the action of $B'$ through $\alpha$, and $\iota$, and such that $\operatorname{sp}_z(z)\ne0$.
\end{enumerate}
\end{lemma}
\leanref{Development/AutomorphicArithmetic/AutomorphicVarieties/ModularityArgument/ThresholdArguments/DenseThresholdCases/UmodAbs.lean}{UpperTheory.V.Hi.LvW.step1_alt}
\leanref{Development/AutomorphicArithmetic/AutomorphicVarieties/ModularityArgument/ThresholdArguments/DenseThresholdCases/UmodAbs.lean}{UpperTheory.V.Hi.LvW.step2_alt}
\leanref{Development/AutomorphicArithmetic/AutomorphicVarieties/ModularityArgument/ThresholdArguments/DenseThresholdCases/UmodAbs.lean}{UpperTheory.V.Hi.LvW.umodAbs_alt}
\leanref{Development/AutomorphicArithmetic/AutomorphicVarieties/ModularityArgument/ThresholdArguments/DenseThresholdCases/LevelUR2SmulKey.lean}{UpperTheory.V.Hi.levelUR2SmulKey_alt}
\begin{proof}
Apply Situation~\ref{sha-hecke-ideal-tw-augmented-setup} to $\iota$, the data $\tau_0$, $\theta$, $a$ and $\mathcal K_0$, $S'$ and the places $w_j$. A compact open subgroup $K\subset G_w$ that is a product of subgroups of the factors is the product $\prod_jK_j$ of a family $K=(K_j)$, and for it $M_K$ is the module of Lemma~\ref{sha-hecke-algebra} for $S'$ in place of $S$, the data $\tau_0$ and $\theta$, the eigenvalues $a_{w,j}$ at the split places $w$ above places outside $S'$ and the $a_{v,j}$ at the places above $\ell$, and the level $\mathcal K_0K$; these data are as in Lemma~\ref{sha-hecke-data} for $S'$. Let $P=\sum_KM_K$ with the representation $\rho$ of $G_w$ by right translation, and $\lambda':\Lambda\to B'$ the map given by the universal property of $\Lambda$ for the restrictions of the $\psi'_{v,i}$ to the $I_v$.

The proof rests on four facts about the $M_K$:
(a) $\rho$ is a representation of $G_w$ on $P$, with $P^K=M_K$, and for $K'\subset K$ with $K'$ normal in $K$, $P^{K'}$ is projective over $\Lambda[K/K']$ (Lemma~\ref{sha-hecke-ideal-tw-augmented-sum});
(b) there is a ring map $\alpha:B'\to\operatorname{End}_\Lambda(P)$ commuting with $\rho$, with $\alpha(\lambda'(b))=b$, that on each nonzero $M_K$ is the composite of the inclusion of its Hecke algebra $\mathbf T_K$ with a map $f_K:B'\to\mathbf T_K$ as in the proof of Proposition~\ref{sha-hecke-ideal}, for $S'$, $\Gamma'$, $B'$ and the level $\mathcal K_0K$ (Lemma~\ref{sha-hecke-ideal-tw-augmented-action});
(c) $M_{K_w}=M$, and $\alpha(t)=\mu(\beta(t))$ on it (Lemma~\ref{sha-hecke-ideal-tw-augmented-base});
(d) every nonzero $z\in P$ has a point as in (4) (Lemma~\ref{sha-hecke-ideal-tw-augmented-points}).
Then (1) follows from (a) and the finiteness of the $M_K$ over $\Lambda$ (Lemma~\ref{sha-hecke-ideal-tw-augmented-hida}(2)), (2) is part of (a), (3) is (c) with $e_0$ the identification of $P^{K_w}=M_{K_w}$ with $M$, and (4) is (d).
% The private s3 of BaseObjectHomExt.lean: step1_alt (theta' = theta, a' = a o opUp, lam = LamObj.toPair for lamUp h, IsWt through cont_up and free_of_neat with neatU, DenseWt, and Hi.exists_hecke_ideal at the level V for every K, which feeds (b)); step2_alt for (c)(1); umodAbs_alt (piBasic for rho, P the supremum of the HSet.MK, h2 for P^K = MK; hiComm for the finiteness; h4 with neatU and levelUR2Faith for the projectivity; levelUR2ActEx_alt for alpha; piPts_n_alt for the points); levelUR2SmulKey_alt for alpha = mu o beta on P^{K_w}; eq_sub for e_0. The four facts are followed in Lemmas sha-hecke-ideal-tw-augmented-sum, -action, -base and -points; the remaining open step, the inputs from Hida theory, is in Lemma sha-hecke-ideal-tw-augmented-hida.
\end{proof}

\begin{proposition}\label{sha-hecke-ideal}
Assume that Galois representations exist, in the sense of Definition~\ref{gal-galois-representations}, relative to the maps satisfying the laws of Definition~\ref{gal-llc-laws}. In Situation~\ref{sha-hecke-situation}, let $\iota:\bar{\mathbf Q}_\ell\cong\mathbf C$. Then there is an ideal $J\subset B$, the \emph{Hecke ideal}, with the following properties; write $D_J$ and $\psi_J$ for the universal pair modulo $J$.
\begin{enumerate}
\item (Finiteness.) $B/J$ is finite over $\Lambda$; if $\lambda b\in J$ with $\lambda\in\Lambda$ nonzero, then $b\in J$; $D_J$ is polarized as in (1) of Situation~\ref{sha-hecke-situation}; the characteristic polynomial of $D_J$ at $g\in D_v$ is $\prod_j(X-\psi_{J,v,j}(g))$; $D_J$ kills the elements of (2) and (3) of Situation~\ref{sha-hecke-situation} and their images under $g\mapsto g^*$; and on $I_u$, $u\in S_{\mathrm a}$, $D_J$ is the determinant of $\bigoplus_it_i$ for every family of Teichm\"uller lifts $t_i$ of the $\bar\chi_i$. If the $\bar\chi_i$ are trivial on $I_{v_0}$ and agree on $D_{v_0}$, and $\varepsilon\bmod\ell$ is nontrivial on $D_{v_0}$, then $D_J$ kills $g-1$ for $g\in I_{v_0}$.
\item (Density.) For every increasing weight $h$, an element of $B$ that vanishes at every point of distinct type of weight $h$ on the zero locus of $J$ lies in $J$.
\item (Automorphic points.) Let $2\le n\le4$ and let $x$ be a point, arithmetic of some weight, with $x(D^{\mathrm u})$ the determinant of a continuous $\rho:\operatorname{Gal}(\bar L/L)\to GL_n(\bar{\mathbf Q}_\ell)$ that is attached through $\iota$ to a regular algebraic cuspidal automorphic representation $\pi$ of $GL_n(\mathbf A_L)$ with $\pi^\vee\cong\pi^c$, windowed (Theorem~\ref{sha-residually-reducible-automorphy}) and unramified at $\tilde v_0$, such that for every $v\mid\ell$, $\rho|_{D_v}$ is triangular with diagonal $\psi_{x,v,1},\dots,\psi_{x,v,n}$ and $\psi_{x,v,i}\ne\psi_{x,v,j}\varepsilon$ for all $i,j$. Then $J\subset\ker x$.
\item (Purity.) Every point $x$ on the zero locus of $J$ that is arithmetic of some weight has $x(D^{\mathrm u})$ the determinant of a continuous semisimple $\rho$, and there are continuous $\rho'_1,\dots,\rho'_a$ and integers $m_1,\dots,m_a\ge1$ with $\det(X-\rho(g))=\prod_i\prod_{0\le j<m_i}\det(X-\rho'_i(g)\varepsilon(g)^{-j})$; if all $m_i=1$, then for all but finitely many places $w$ of $L$, any two Frobenius eigenvalues $z,z'$ of $\rho$ at $w$ satisfy $|\iota z|<\operatorname{N}w\,|\iota z'|$.
\item (Classicality.) If $x$ is as in (4) and $\rho$ is irreducible, then $\rho$ is attached through $\iota$ to a regular algebraic cuspidal automorphic representation $\pi$ of $GL_n(\mathbf A_L)$ with $\pi^\vee\cong\pi^c$, and it is windowed at the places not above $v_0$.
\item (Twists.) Let $x$ be a point of distinct type on the zero locus of $J$ and $\beta$ a character of $\Gamma_0$ of $\ell$-power order with open kernel, trivial on $I_{v_0}$ and on the $I_u$, with $\beta(cgc^{-1})=\beta(g)^{-1}$. If $x'$ is a point with $x'(D^{\mathrm u})=x(D^{\mathrm u})\otimes\beta$ and $\psi_{x',v,j}=\psi_{x,v,j}\beta$, then $J\subset\ker x'$.
\item (Taylor-Wiles data.) Let $D=n[L^+:\mathbf Q]$. There are a map $\lambda:\mathcal O[[X_1,\dots,X_D]]\to B$ sending each $X_k$ to $c_k\psi^{\mathrm u}_{v,i}(\gamma)-1$ for some $c_k\in\mathcal O^\times$ and $\gamma\in D_v$, and a $B$-module $M_0$, finite and free over $\mathcal O[[X_1,\dots,X_D]]$ through $\lambda$, whose support in $\operatorname{Spec}B$ is the zero locus of $J$. For every family of places $w_{N,j}$, with inverses $\phi_{N,j}$ of Frobenius elements and tame generators $\tau_{N,j}$, as in Lemma~\ref{sha-patching-places}, there is an integer $e$, and for every $N$ there are: an ordinary pair for the group $\Gamma_N$ of the level augmented by the $w_{N,j}$ (Definition~\ref{sha-hecke-ideal-tw-level-module}), over a complete local noetherian $\mathcal O$-algebra $R_N$ with residue field $k$; elements $U_{j,i}\in R_N$, $1\le j\le q$, $1\le i\le n$, lifting the eigenvalues of $\bar D$ at the $\phi_{N,j}$; a map $s_N$ to $R_N$ from the power series ring $S$ over $\mathcal O[[X_1,\dots,X_D]]$ in variables $y_{j,i}$, compatible with $\lambda$ on the weight variables; and an $R_N$-module $M_N$ isomorphic, as an $S$-module, to a free module over $S/\big((1+y_{j,i})^{\ell^{a_j}}-1\big)$ for some $a_j\ge N$. Let $B[u]$ be the quotient of $B[[y_{j,i}]]$ that makes the $u_{j,i}=\tilde u_{j,i}+y_{j,i}$, for lifts $\tilde u_{j,i}\in\mathcal O$ of the residual eigenvalues, the roots of the characteristic polynomials of $D^{\mathrm u}$ at the $\phi_{N,j}$. There are a map $R_N\to B[u]$ carrying the pair of $R_N$ to $D^{\mathrm u}$ and $\psi^{\mathrm u}$, each $U_{j,i}$ to $u_{j,i}$ and each $y_{j,i}$ to $0$, and maps between $M_N/(y_{j,i})M_N$ and $B[u]\otimes_BM_0$, compatible with it, whose two composites are multiplication by $f_N=\prod_j\prod_i\prod_{i'\ne i}(U_{j,i}-U_{j,i'})(U_{j,i}-\operatorname{N}w_{N,j}\,U_{j,i'})$. Finally, the $e$-th powers of $f_N\operatorname{tr}\big(h\prod_i(\phi_{N,j}-U_{j,i})\big)$ and of $f_N\operatorname{tr}\big(h\prod_{k\ne i}(\phi_{N,j}-U_{j,k})\,(\tau_{N,j}-1-y_{j,i})\big)$, for $h$ in the subgroup $\Gamma_{N,0}$ of $\Gamma_N$ fixing $L$, kill $M_N$.
\end{enumerate}
\end{proposition}
\leanref{Development/AutomorphicArithmetic/AutomorphicVarieties/ModularityArgument/ThresholdArguments/DenseThresholdCases/HidaU.lean}{UpperTheory.V.W.HidaUAlt}
\leanref{Development/AutomorphicArithmetic/AutomorphicVarieties/ModularityArgument/ThresholdArguments/DenseThresholdCases/THidaU.lean}{UpperTheory.V.Hi.hida_unitary_of_printed_hida}
\leanref{Development/AutomorphicArithmetic/AutomorphicVarieties/ModularityArgument/ThresholdArguments/DenseThresholdCases/PointTheoremBaseObject.lean}{UpperTheory.V.W.UniversalOrdinaryDeterminantRing.HidaPropsAlt}
\leanref{Development/AutomorphicArithmetic/AutomorphicVarieties/ModularityArgument/ThresholdArguments/DenseThresholdCases/ExistsHeckeIdeal.lean}{UpperTheory.V.Hi.exists_hecke_ideal}
\leanref{Development/AutomorphicArithmetic/AutomorphicVarieties/PatchingModules/ThetaCoefficientMapEq.lean}{UpperTheory.V.Pa.TWAut}
\begin{proof}
Take the data of Lemma~\ref{sha-hecke-data} and the level $\mathcal K_0$ of Lemma~\ref{sha-hecke-algebra}. If the localized module $M$ of Lemma~\ref{sha-hecke-algebra} is $0$, let $J=B$; all the properties hold trivially. Otherwise, let $\mathbf T$, $D_{\mathbf T}$ and $\psi_{\mathbf T}$ be as in Lemmas~\ref{sha-hecke-algebra} and~\ref{sha-hecke-determinant}.

\emph{The map from $B$.} The pair $(D_{\mathbf T},\psi_{\mathbf T})$ is an ordinary pair over $\mathbf T$. It reduces to $\bar D$ and the $\bar\psi_{v,i}$: by Lemma~\ref{sha-hecke-determinant}(1) and Lemma~\ref{sha-hecke-data}(3) the reductions agree at the Frobenius elements at split places outside $S$, which are dense. At $v\mid\ell$, on $I_v$, Lemma~\ref{sha-hecke-determinant}(2) and Lemma~\ref{sha-hecke-data}(2) make $\psi_{\mathbf T,v,i}$ the image of $\kappa_{v,i}$, which lifts $\bar\psi_{v,i}$; at the image of a uniformizer, Lemma~\ref{sha-hecke-determinant}(2) and Lemma~\ref{sha-hecke-data}(3), with $U_{v,j}\equiv a_{v,j}$, give the products over $i\le j$ of $\psi_{\mathbf T,v,i}\varepsilon^{i-1}$ modulo the maximal ideal, and induction on $j$ gives the factors. Each relation (1) to (3) of Situation~\ref{sha-hecke-situation} says that a continuous function on $\Gamma_0$ with values in $\mathbf T$ vanishes; it vanishes at every point of distinct type of a fixed weight, by Lemma~\ref{sha-hecke-determinant}(3), so it vanishes by Lemma~\ref{sha-hecke-algebra}. The universal property gives $f:B\to\mathbf T$ compatible with the maps from $\Lambda$. It is surjective, because $\mathbf T$ is generated over $\Lambda$ by the $T_{w,j}$, which are coefficients of characteristic polynomials of $D_{\mathbf T}$, and by the $U_{v,j}$, which are products of values of the $\psi_{\mathbf T,v,i}$. Let $J=\ker f$, so $B/J\cong\mathbf T$.

\emph{The properties.} (1) follows from Lemma~\ref{sha-hecke-algebra} for finiteness and torsion freeness; the relations of $D_J$ hold at the points of distinct type by Lemma~\ref{sha-hecke-determinant}(3), hence on $\mathbf T$ by density; the statement at $I_u$ holds because the representations of type $\tau_0$ have inertial type $\bigoplus_it_i$ there; and the statement at $I_{v_0}$ holds at these points by a local argument: under the residual conditions at $v_0$, a nonzero monodromy operator of $\rho$ at $\tilde v_0$ would give two Frobenius eigenvalues with ratio $p$ that are congruent modulo the maximal ideal, which is impossible as $p\not\equiv1\pmod\ell$, so the image of inertia is finite of order prime to $\ell$ and residually trivial, hence trivial. (2) is the density of Lemma~\ref{sha-hecke-algebra} transported by $f$.

\emph{Contributing representations.} Properties (3) and (6) rest on the following statement. Let $x$ be a point, arithmetic of weight $h$, let $\sigma$ be an automorphic representation of $U(n)$, and let $\rho:\operatorname{Gal}(\bar L/L)\to GL_n(\bar{\mathbf Q}_\ell)$ be continuous and semisimple, with $x(D^{\mathrm u})$ the determinant of $\rho|_{\Gamma_0}$. Suppose:
(a) at every place $w$ of $L$ split over $L^+$ and not above $\ell$ at which $\sigma_w$ has Satake parameters $B$, $\rho$ is unramified and the eigenvalues of $\rho$ at the inverse of a Frobenius element are, through $\iota$, the numbers $\operatorname{N}w^{(n-1)/2}\alpha$ for $\alpha\in B$;
(b) $\sigma$ has weight $h$ at infinity;
(c) $\sigma$ is unramified at the places of $L$ above the places of $L^+$ outside $S$;
(d) the local component of $\sigma$ at $\tilde v_0$ has nonzero vectors fixed by the principal congruence subgroup of level $\tilde v_0$;
(e) for $v\mid\ell$, the local component $\sigma_v$ is $\mathrm{PS}(\omega_v)$ for smooth characters $\omega_{v,i}$ of $L_{\tilde v}^\times$ with $x(\psi^{\mathrm u}_{v,i}(s))=\iota^{-1}(\omega_{v,i}(a(s)))\prod_\tau\tau(a(s))^{-h_{v,i,\tau}}$ for $s\in W_v$, in the notation of the flags at $\ell$ in the proof of Lemma~\ref{sha-hecke-determinant};
(f) at $u\in S_{\mathrm a}$, $\mathrm{WD}_\iota(\rho|_{W_{\tilde u}})=\operatorname{rc}(\sigma_u\nu^{(1-n)/2})$ for every $\operatorname{rc}$ that satisfies the laws.
Then $J\subset\ker x$.

\emph{The forms.} By (f) and Lemma~\ref{sha-hecke-ideal-inertia}(2), for every family of Teichm\"uller lifts $t_i$ of the $\bar\chi_i$, every $u\in S_{\mathrm a}$ and every $\operatorname{rc}$ that satisfies the laws, $\operatorname{rc}(\sigma_u\nu^{(1-n)/2})$ has monodromy $0$ and restriction $\bigoplus_i\iota t_i$ to inertia. With (b), (c) and (d), Lemma~\ref{sha-hecke-ideal-forms} shows that the space $Y$ of classical forms of weight $h$, type $\tau_0$ and level $\mathcal K_0$ in $\sigma$ is nonzero. For a split place $w\notin S$, $\sigma_w$ is spherical by (c), so each $T_{w,j}$ acts on $Y$ by a scalar $\mu_{w,j}$.

\emph{The representation $\rho_\sigma$.} Let $\rho_\sigma$ be the representation of Lemma~\ref{sha-hecke-ideal-unitary}. By (c) and (1) of that lemma it is unramified at the places above the places of $L^+$ outside $S$, so it factors through $\Gamma_0$. At a split place $w\notin S$, Lemma~\ref{sha-hecke-ideal-unitary}(3) and the Satake isomorphism make the characteristic polynomial of $\rho_\sigma$ at the inverse of a Frobenius element at $w$, carried by $\iota$, the Hecke polynomial with coefficients $\mu_{w,j}$; by (a), so is that of $\rho$. These inverse Frobenius elements are dense in $\Gamma_0$, and the coefficients of the characteristic polynomials of both determinants are continuous, so $x(D^{\mathrm u})$ is the determinant of $\rho_\sigma$.

\emph{The eigenvalues away from $\ell$.} Let $w\notin S$ be split and $1\le j\le n$. The characteristic polynomial of $\rho$ at the inverse of a Frobenius element at $w$ is the image under $x$ of that of $D^{\mathrm u}$, which reduces to that of $\bar D$, the Hecke polynomial with coefficients $a_{w,j}\bmod\mathfrak m_{\mathcal O}$ (Lemma~\ref{sha-hecke-data}(3)). The coefficient of $X^{n-j}$ in the Hecke polynomial with coefficients $b_1,\dots,b_n$ is $(-1)^j\operatorname{N}w^{j(j-1)/2}b_j$, and $\operatorname{N}w$ is an $\ell$-adic unit since $w\nmid\ell$. So $\iota^{-1}\mu_{w,j}-a_{w,j}$ lies in the maximal ideal of the ring of integers of $\bar{\mathbf Q}_\ell$.

\emph{The eigenvalues at $\ell$.} Let $\chi_{v,i}$ and $c_{v,j}$ be defined as in the flags at $\ell$ in the proof of Lemma~\ref{sha-hecke-determinant}, with $x$ in place of $y$. Let $v\mid\ell$, with uniformizer $\varpi$ of $L_{\tilde v}$, residue field of $q_v$ elements and local Artin map $a$. Let $z\in\mathcal O_{L_{\tilde v}}^\times$ and $s\in W_v$ with $a(s)=z$. On $I_v$, $\psi^{\mathrm u}_{v,i}$ is the image of $\kappa_{v,i}$, so by Lemma~\ref{sha-hecke-data}(2), $x(\psi^{\mathrm u}_{v,i}(s))\varepsilon(s)^{i-1}x(\theta(e_{v,i}(z)))=1$; also $\varepsilon(s)=\prod_\tau\tau(z)$, as in the proof of Lemma~\ref{sha-hecke-local-characters}. With (e) this gives $\omega_{v,i}(z)\chi_{v,i}(z)=1$. Now let $s\in W_v$ with $a(s)=\varpi$, so that $\varepsilon(s)=q_v^{-1}\prod_\tau\tau(\varpi)$. By (e),
\[
c_{v,j}^{-1}E_j(\omega_v)=\iota x\Big(\prod_{i\le j}\psi^{\mathrm u}_{v,i}(s)\varepsilon(s)^{i-1}\Big).
\]
This is $\iota$ of the image under $x$ of a unit of $B$, so its image under $\iota^{-1}$ is a unit. It is congruent to $a_{v,j}$, because $\psi^{\mathrm u}_{v,i}$ lifts $\bar\psi_{v,i}$ and by Lemma~\ref{sha-hecke-data}(3).

\emph{The point of $\mathbf T$.} By (e) and the last two steps, Lemma~\ref{sha-hecke-ideal-eigenform} applies to $\sigma$, $Y$, the $\mu_{w,j}$ and $\mu_{v,j}=c_{v,j}^{-1}E_j(\omega_v)$. So $M\ne0$, and there is an $\mathcal O$-algebra map $y:\mathbf T\to\bar{\mathbf Q}_\ell$, equal to $x$ on $\Lambda$, with $\iota y(T_{w,j})=\mu_{w,j}$ and $\iota y(U_{v,j})=c_{v,j}^{-1}E_j(\omega_v)$. By Lemma~\ref{sha-hecke-determinant}(1) and Lemma~\ref{sha-hecke-local-characters}, $y(D_{\mathbf T})$ is the determinant of $\rho_\sigma$, and $y\circ\psi_{\mathbf T,v,i}$ satisfies on $W_v$ the relation of (e); a continuous character of $D_v$ is determined by its values on the dense subgroup $W_v$. So $y\circ f$ and $x$ agree on the coefficients of the characteristic polynomials of $D^{\mathrm u}$ and on the values of the $\psi^{\mathrm u}_{v,i}$, which topologically generate $B$ over $\mathcal O$; hence $x=y\circ f$ and $J=\ker f\subset\ker x$.

\emph{Properties (3) to (7).} For (3), let $x$, $\rho$ and $\pi$ be as there, and let $h$ be the weight of $x$, with $\psi_{x,v,i}\phi_{v,i}=\prod_\tau\tau(\operatorname{Art}^{-1}\cdot)^{-h_{v,i,\tau}}$ on $I_v$ as in Definition~\ref{sha-hecke-arcs}(1).

\emph{Compatibility.} The pair formed by $\pi$ and the $n$-th power of the quadratic Hecke character of $L^+$ attached to $L$ is polarized, so Lemma~\ref{sha-hecke-ideal-attached} applies to $\rho$ and $\pi$, with multisets $\mathrm{HT}_\tau$. At $v\mid\ell$, $\rho|_{D_v}$ is triangular with diagonal $\psi_{x,v,i}$ and $\psi_{x,v,i}\ne\psi_{x,v,j}\varepsilon$, so ($\alpha$) of the windowed condition gives its second condition at $v$. So by ($\beta$), at every place $w\nmid\ell$ and every $v\mid\ell$, the supercuspidal support of the local component of $\pi$ consists of characters $\xi$ with $q_w^{s-1/2}<|\xi(\varpi_w)|<q_w^{s+1/2}$. No two of them differ by $\nu^{\pm1}$, so the local component is plain, and $\operatorname{rc}(\pi_w\nu^{(1-n)/2})$ is abelian and monodromy-free for every $\operatorname{rc}$ that satisfies the laws, as in the proof of Lemma~\ref{sha-hecke-ideal-attached}(2). Hence $\mathrm{WD}_\iota(\rho|_{W_w})=\operatorname{rc}(\pi_w\nu^{(1-n)/2})$ at every $w\nmid\ell$.

\emph{Spherical components.} An element $g$ of an inertia group at a place above a place of $L^+$ outside $S$ maps to $1$ in $\Gamma_0$, so $\operatorname{tr}\rho(hg)=\operatorname{tr}\rho(h)$ for all $h$, since $x(D^{\mathrm u})$ is the determinant of $\rho$; as in the proof of Lemma~\ref{sha-hecke-ideal-inertia}(1), $\rho(g)=1$. So $\rho$ is unramified at these places. At every place $w\nmid\ell$ at which $\rho$ is unramified, $\pi_w$ is spherical by Lemma~\ref{sha-hecke-ideal-spherical}.

\emph{The descent.} So $\pi$ is spherical at every place above a place of $L^+$ outside $S$. Since $2\le n\le4$, the descent of conjugate self-dual cuspidal representations of $GL_n(\mathbf A_L)$ to $U(n)$, one of the published results used \cite{Lab11}, gives an automorphic representation $\sigma$ of $U(n)$, unramified at the places not split over $L^+$, with the Satake parameters of $\pi$ at almost all split places. By base change (Theorem~\ref{gal-base-change-definite}) and the strong multiplicity one theorem \cite{JS81b}, the base change of $\sigma$ is $\pi$; so $\sigma_w\cong\pi_w$ at every split place $w$, and the infinitesimal characters of $\sigma$ at infinity are those of $\pi$.

\emph{The places above $\ell$.} Let $v\mid\ell$. The local component $\pi_v$ is plain, so for a map $\operatorname{rc}$ that satisfies the laws (Theorem~\ref{gal-llc-exists}), Lemma~\ref{sha-hecke-ideal-attached} makes $\rho|_{D_v}$ de Rham with Hodge-Tate multisets $\mathrm{HT}_\tau$, and potentially semistable with $\mathrm{WD}_\iota(\rho|_{D_v})=\operatorname{rc}(\pi_v\nu^{(1-n)/2})$. Since $h$ is increasing, Lemma~\ref{sha-hecke-ideal-above}, applied to $\rho|_{D_v}$, the $h_{v,i,\tau}$, the $\psi_{x,v,i}$ and $\pi_v$, gives smooth characters $\omega_{v,i}$ with $\pi_v\cong\mathrm{PS}(\omega_v)$ and the relation of (e), and $\mathrm{HT}_\tau=\{h_{v,1,\tau},\dots,h_{v,n,\tau}\}$ for the embeddings $\tau$ above $\tilde v$.

\emph{Conditions (a) to (f).} (a) At a split place $w\notin S$ at which $\sigma_w$ has Satake parameters $B$, $\pi_w\cong\sigma_w$, $\rho$ is unramified, and $\mathrm{WD}_\iota(\rho|_{W_w})=\operatorname{rc}(\pi_w\nu^{(1-n)/2})$ gives the Frobenius eigenvalues, as in the proof of Lemma~\ref{sha-hecke-ideal-unitary}(3). (b) For $v\mid\ell$ and every $\tau$ above $\tilde v$, $\{h_{v,i,\tau}\}=\mathrm{HT}_\tau=\{\frac{n-1}2-z:z\in Z_{\iota\circ\tau}(\pi)\}$, so the infinitesimal character of $\pi$ at $\iota\circ\tau$ is that of the weight $h$, and $\sigma$, which has the infinitesimal characters of $\pi$, has weight $h$. (c) At a split place outside $S$, $\rho$ is unramified, so $\pi_w$ is spherical and so is $\sigma_w\cong\pi_w$; at the other places outside $S$ this is the descent. (d) Since $\operatorname{N}v_0=p\ne\ell$, $v_0\nmid\ell$; $\rho$ is unramified at $\tilde v_0$, so $\pi_{v_0}$ is spherical, and $\sigma_{v_0}\cong\pi_{v_0}$ has nonzero vectors fixed by $GL_n(\mathcal O_{L_{\tilde v_0}})$, hence by the principal congruence subgroup. (e) is the last step. (f) At $u\in S_{\mathrm a}$, $\sigma_u\cong\pi_u$ and $\mathrm{WD}_\iota(\rho|_{W_{\tilde u}})=\operatorname{rc}(\pi_u\nu^{(1-n)/2})$. So (a) to (f) hold, and $J\subset\ker x$.

For (4), let $x$ be a point on the zero locus of $J$, arithmetic of weight $h$. Then $J\ne B$, and $x=y\circ f$ for an $\mathcal O$-algebra map $y:\mathbf T\to\bar{\mathbf Q}_\ell$ that is arithmetic of weight $h$. Lemma~\ref{sha-hecke-determinant}(3) gives a continuous semisimple $\rho$ with $x(D^{\mathrm u})=y(D_{\mathbf T})$ the determinant of $\rho$, together with an automorphic representation $\sigma$ of $U(n)$ with which $\rho$ is compatible at the split places in the sense of Lemma~\ref{sha-hecke-ideal-speh}. That lemma gives the $\rho'_i$, the $m_i$ and the bound on Frobenius eigenvalues.

\emph{Classicality.} For (5), let $x$ and $\rho$ be as in (4), with $\rho$ irreducible, and let $h$ be the weight of $x$. Then $J\ne B$, and $x=y\circ f$ for an $\mathcal O$-algebra map $y:\mathbf T\to\bar{\mathbf Q}_\ell$ that is arithmetic of weight $h$. By the first paragraph of the proof of Lemma~\ref{sha-hecke-determinant}, there is an automorphic representation $\sigma$ of $U(n)$ whose space of classical forms of weight $h$, type $\tau_0$ and level $\mathcal K_0$ is nonzero, and whose representation $\rho_\sigma$ of Lemma~\ref{sha-hecke-ideal-unitary} factors through $\Gamma_0$, with determinant $x(D^{\mathrm u})$. The representations $\rho$ and $\rho_\sigma$ are semisimple with the same characteristic polynomials, so they are conjugate by the Brauer-Nesbitt theorem \cite[(30.16)]{CR62}, and $\rho_\sigma$ is irreducible. By Lemma~\ref{sha-hecke-ideal-cuspidal}, applied to $\sigma$ and to $\rho_\sigma$, whose Frobenius eigenvalues at the split places are given by Lemma~\ref{sha-hecke-ideal-unitary}(3), $\rho_\sigma$, and so $\rho$, is weakly attached through $\iota$ to a regular algebraic cuspidal $\pi$ with $\pi^\vee\cong\pi^c$, which has Satake parameters at the places where $\sigma$ is unramified and whose local components at the split places are those of $\sigma$.

\emph{Plain local components.} Let $w\nmid\ell$ be a finite place of $L$ not above $v_0$. Then $\pi$ is plain at $w$ in the sense of Lemma~\ref{sha-hecke-ideal-plain}. If $w$ does not lie above a place of $S$, $\sigma$ is unramified at $w$, so $\pi$ has Satake parameters at $w$, and (2) of that lemma applies. Otherwise $w$ lies above a place $u\in S_{\mathrm a}$, and by (3) of that lemma it suffices to treat $w=\tilde u$. A local component $\tau$ of $\pi$ at $\tilde u$ is a local component of $\sigma$. For a map $\operatorname{rc}$ that satisfies the laws (Theorem~\ref{gal-llc-exists}), Lemma~\ref{sha-hecke-data}(1) makes $\operatorname{rc}(\tau\nu^{(1-n)/2})$ of monodromy $0$, with restriction $\bigoplus_i\iota t_i$ to inertia. So $\tau$ is plain, by (5) of that lemma.

\emph{The windowed condition away from $\ell$.} Let $\pi'$ be a regular algebraic cuspidal automorphic representation of $GL_n(\mathbf A_L)$ to which $\rho$ is attached through $\iota$, and let $t$ be as in Lemma~\ref{sha-hecke-ideal-plain}(1) for $\pi'$. At a place $w\nmid\ell$ not above $v_0$, $\pi$ is plain, so $\pi'$ is plain, by (4) of that lemma, and the bound of ($\beta$) holds at $w$ with $s=t$, by (1).

\emph{A second cuspidal representation.} Lemma~\ref{sha-hecke-ideal-ordinary-point}, applied to $y$, gives an automorphic representation $\sigma_0$ of $U(n)$ and, for each $v\mid\ell$, smooth characters $\omega_{v,i}$ with (1) to (3) of that lemma; there $y\circ\psi_{\mathbf T,v,i}=\psi_{x,v,i}$. The representations $\rho$ and $\rho_{\sigma_0}$ are semisimple with the same characteristic polynomials, so they are conjugate by the Brauer-Nesbitt theorem, and $\rho_{\sigma_0}$ is irreducible. By Lemma~\ref{sha-hecke-ideal-cuspidal}, applied to $\sigma_0$ and to $\rho_{\sigma_0}$, whose Frobenius eigenvalues are given by Lemma~\ref{sha-hecke-ideal-unitary}(3), $\rho_{\sigma_0}$, and so $\rho$, is weakly attached through $\iota$ to a regular algebraic cuspidal $\pi_0$ with $\pi_0^\vee\cong\pi_0^c$, whose local components at the split places are those of $\sigma_0$. So for $v\mid\ell$ every local component of $\pi_0$ at $\tilde v$ is isomorphic to the local component of $\sigma_0$ at $\tilde v$, which is a subquotient of $\mathrm{PS}(\omega_v)$.

\emph{The windowed condition above $\ell$.} Let $w$ be a place of $L$ above $\ell$ and $v$ the place of $L^+$ under it, and suppose that at every valuation ring $A$ of $\bar L$ over $w$, $\rho|_{D_A}$ is conjugate to an upper triangular representation with diagonal characters $\kappa_1,\dots,\kappa_n$ such that $\kappa_i\ne\kappa_j\varepsilon$ for all $i,j$. Property (1) holds for $J$, so by Lemma~\ref{sha-hecke-ideal-flag}, $\rho|_{D_v}$ is conjugate to an upper triangular representation with diagonal $\psi_{x,v,1},\dots,\psi_{x,v,n}$. The places $\tilde v$ and $w$ lie above $v$. By Lemma~\ref{sha-hecke-ideal-conjugate}(1), applied to $\pi_0$, $\rho$, $\tilde v$ and $w$, $\psi_{x,v,i}\ne\psi_{x,v,j}\varepsilon$ for all $i,j$. Since $h$ is increasing, Lemma~\ref{sha-hecke-ideal-exponents}(1), applied to the local components of $\pi_0$ at $\tilde v$ and the characters $\omega_{v,i}$, shows that $\pi_0$ is plain at $\tilde v$. Let $t_0$ be as in Lemma~\ref{sha-hecke-ideal-plain}(1) for $\pi_0$. The supercuspidal supports of the local components of $\pi_0$ at $\tilde v$ satisfy the bound of that lemma with $t_0$, so by Lemma~\ref{sha-hecke-ideal-exponents}(2), $\psi_{x,v,i}^m\ne(\psi_{x,v,j}\varepsilon)^m$ for all $i,j$ and $m\ge1$. By Lemma~\ref{sha-hecke-ideal-conjugate}(2), at every valuation ring $A$ over $w$, $\rho|_{D_A}$ is conjugate to an upper triangular representation with diagonal characters $\kappa_i$ such that $\kappa_i^m\ne(\kappa_j\varepsilon)^m$ for all $i,j$ and $m\ge1$. This is condition ($\alpha$) at $w$. Moreover, $\pi'\cong\pi_0$ by Theorem~\ref{lit-jacquet-shalika}(1), since at almost all places both have the Satake parameters determined through $\iota$ by the Frobenius eigenvalues of $\rho$. As $\pi_0$ is plain at $\tilde v$ and $\pi_0^\vee\cong\pi_0^c$, $\pi'$ is plain at $w$, by Lemma~\ref{sha-hecke-ideal-plain}(3).

Finally, let $w$ be a place of $L$ above $\ell$ that satisfies the second condition of ($\alpha$). With $m=1$ it satisfies the first, so $\pi'$ is plain at $w$ by the last step, and the bound of ($\beta$) holds at $w$ with $s=t$, by Lemma~\ref{sha-hecke-ideal-plain}(1). So $\rho$ is windowed at the places not above $v_0$.
% (5) is irrArithPointsAutomorphicAway_of_isJ. The cuspidal representation is the private coreA (core_hecke_of_automorphic_hecke_dense_weight with autAway_alt, then absPt), iso_of_det (Brauer-Nesbitt) and AwayOk; autAway_alt uses PartLinkAlt = projected_forms_link_of_printed_hida (partU_alt, then link2_alt: Lemma sha-hecke-ideal-cuspidal), sphUnl off S, and at the places above S_a hC (ConjTor), plainIff, partUR2Type_alt and typePlain. The windowed condition is WitOff: away from l, winOf and tu_same; above l, BornJ = born_j_of_printed_polarization_window, with ENDO = endOrd_of_printsH (the private core of IsoOfDet.lean, with core_ordinary_of_literal_automorphy_n = Lemma sha-hecke-ideal-ordinary-point and autUO_of_printsH, whose part2 is Lemma sha-hecke-ideal-cuspidal; LocUniq.isLocComp_unique; ExpOf with h = htW lam), flagJ_alt (Lemma sha-hecke-ideal-flag), conj_unlinked_of_printed_results (Lemma sha-hecke-ideal-conjugate, with w = s.w v and w' = w), exp_torsion_of_literal_automorphy_n (Lemma sha-hecke-ideal-exponents), winOf, strongMultOneC and ConjTor. In irrArithPointsAutomorphicAway_of_isJ, (alpha) is (born v hv hu).1, and above l the bound of (beta) at a place with PotUnl comes from Unl1 (m = 1), (born v hl _).2 and winOf for pi'.

\emph{Twists.} For (6), let $x$, $\beta$ and $x'$ be as there, with $x$ of distinct type of weight $h$, and let $\phi_{v,i}$ be characters as in Definition~\ref{sha-hecke-arcs}(1), pairwise distinct for each $v$. Then $J\ne B$, and $x=y\circ f$ for an $\mathcal O$-algebra map $y:\mathbf T\to\bar{\mathbf Q}_\ell$ of distinct type of weight $h$. Let $\sigma$ and the $\omega_{v,i}$ be given for $y$ by Lemma~\ref{sha-hecke-ideal-ordinary-point}, and let $\rho=\rho_\sigma$. Then $\rho$ factors through $\Gamma_0$, $x(D^{\mathrm u})$ is its determinant, $\rho$ is compatible with $\sigma$ at the split places in the sense of Lemma~\ref{sha-hecke-ideal-speh} (Lemma~\ref{sha-hecke-ideal-unitary}(2)), and the space of classical forms of weight $h$, type $\tau_0$ and level $\mathcal K_0$ in $\sigma$ is nonzero. So $\sigma$ has weight $h$ at infinity; it is unramified at the places above the places of $L^+$ outside $S$; its local component at $\tilde v_0$ has nonzero vectors fixed by the principal congruence subgroup; and, by Lemma~\ref{sha-hecke-data}(1), for every $u\in S_{\mathrm a}$, every local component $\tau$ of $\sigma$ at $\tilde u$ and every $\operatorname{rc}$ that satisfies the laws, $\operatorname{rc}(\tau\nu^{(1-n)/2})$ has monodromy $0$ and restriction $\bigoplus_i\iota t_i$ to inertia.

For $v\mid\ell$ the local component $\sigma_v$ is a subquotient of $\mathrm{PS}(\omega_v)$, where the $\omega_{v,i}$ satisfy the relation of (e) for $x$, and it is isomorphic to $\mathrm{PS}(\omega_v)$. For $s\in I_v$, comparing the relation of (e) with that of Definition~\ref{sha-hecke-arcs}(1) gives $\iota^{-1}(\omega_{v,i}(a(s)))=\phi_{v,i}(s)^{-1}$. The $\phi_{v,i}$ are pairwise distinct and $a(I_v)=\mathcal O_{L_{\tilde v}}^\times$, so the $\omega_{v,i}$ are pairwise distinct on $\mathcal O_{L_{\tilde v}}^\times$, on which $|\cdot|$ is trivial. So no two of the $\omega_{v,i}$ are equal or differ by $|\cdot|^{\pm1}$, $\mathrm{PS}(\omega_v)$ is irreducible, and it is its only irreducible subquotient \cite{BZ77,Zel80}. Hence $\sigma_v\cong\mathrm{PS}(\omega_v)$.

The order of $\beta$ is a power of the odd prime $\ell$. View $\beta$ and $\rho$ as representations of $\operatorname{Gal}(\bar L/L)$ through $\Gamma_0$; then $\beta(cgc^{-1})=\beta(g)^{-1}$. Let $\kappa$ and $\sigma'$ be as in Lemma~\ref{sha-hecke-ideal-twist-global}, for $\sigma$ and $\beta$, and put $\rho'=\rho\otimes\beta$. It is continuous and semisimple, and $x'(D^{\mathrm u})=x(D^{\mathrm u})\otimes\beta$ is its determinant. The point $x'$ is arithmetic of weight $h$: on $I_v$, $\psi_{x',v,i}\cdot\phi_{v,i}\beta^{-1}=\psi_{x,v,i}\phi_{v,i}$, and $\phi_{v,i}\beta^{-1}$ has open kernel. Conditions (a) to (f) of the contributing representations hold for $x'$, $\sigma'$ and $\rho'$.
(a) Let $w\nmid\ell$ be a place split over $L^+$, and $\tau'$ a local component of $\sigma'$ at $w$ with Satake parameters $B$. By Lemma~\ref{sha-hecke-ideal-twist-global}(4), $\tau'\cong\tau\otimes(\kappa_w\circ\det)$ for a local component $\tau$ of $\sigma$ at $w$. By Lemma~\ref{sha-hecke-ideal-twist-local}(2), $\tau$ is good and $\operatorname{rc}(\tau\nu^{(1-n)/2})$ is abelian and monodromy-free for every $\operatorname{rc}$ that satisfies the laws; so $\mathrm{WD}_\iota(\rho|_{W_w})=\operatorname{rc}(\tau\nu^{(1-n)/2})$ by the compatibility of $\rho$ with $\sigma$, and $\mathrm{WD}_\iota(\rho'|_{W_w})=\operatorname{rc}(\tau'\nu^{(1-n)/2})$ by Lemma~\ref{sha-hecke-ideal-twist-local}(3). As in the proof of Lemma~\ref{sha-hecke-ideal-unitary}(3), $\rho'$ is unramified at $w$, and its eigenvalues at the inverse of a Frobenius element are, through $\iota$, the numbers $\operatorname{N}w^{(n-1)/2}\alpha$ for $\alpha\in B$.
(b) By Lemma~\ref{sha-hecke-ideal-twist-global}(1), $\sigma'$ has weight $h$ at infinity.
(c) Let $w$ be a place of $L$ above a place of $L^+$ outside $S$. Since $\beta$ factors through $\Gamma_0$, it is trivial on the inertia groups at $w$; and $\sigma$ is unramified at $w$. So $\sigma'$ is unramified at $w$, by Lemma~\ref{sha-hecke-ideal-twist-global}(2).
(d) Since $\beta$ is trivial on $I_{v_0}$, Lemma~\ref{sha-hecke-ideal-twist-global}(3) gives a local component of $\sigma'$ at $\tilde v_0$ with nonzero vectors fixed by the principal congruence subgroup.
(e) Let $v\mid\ell$. By Lemma~\ref{sha-hecke-ideal-twist-global}(3), $\sigma'$ has a local component $\tau'\cong\mathrm{PS}(\omega_v)\otimes(\kappa_{\tilde v}\circ\det)$ at $\tilde v$. By Lemma~\ref{sha-hecke-ideal-twist-local}(4), applied with $\psi_i=\psi_{x,v,i}$ and $\mu=\beta$, $\tau'\cong\mathrm{PS}(\omega'_v)$ for smooth characters $\omega'_{v,i}$ that satisfy the relation of (e) with the characters $\psi_{x,v,i}\beta=\psi_{x',v,i}$.
(f) Let $u\in S_{\mathrm a}$ and $\tau'$ a local component of $\sigma'$ at $\tilde u$. By Lemma~\ref{sha-hecke-ideal-twist-global}(4), $\tau'\cong\tau\otimes(\kappa_{\tilde u}\circ\det)$ for a local component $\tau$ of $\sigma$ at $\tilde u$. Let $\operatorname{rc}$ satisfy the laws. Then $\operatorname{rc}(\tau\nu^{(1-n)/2})$ has monodromy $0$ and restriction $\bigoplus_i\iota t_i$ to inertia, so by Lemma~\ref{sha-hecke-ideal-twist-local}(1), $\tau$ is good and $\operatorname{rc}(\tau\nu^{(1-n)/2})$ is abelian and monodromy-free. As in (a), $\mathrm{WD}_\iota(\rho|_{W_{\tilde u}})=\operatorname{rc}(\tau\nu^{(1-n)/2})$, and $\mathrm{WD}_\iota(\rho'|_{W_{\tilde u}})=\operatorname{rc}(\tau'\nu^{(1-n)/2})$ by Lemma~\ref{sha-hecke-ideal-twist-local}(3), since $\tilde u\nmid\ell$.
So $J\subset\ker x'$, by the contributing representations.
% (6) is Z.Desc.W.twistOk_of_isJ = jTwist_of_alt with typeGd, twist_away_statement, sphUntw_alt, ordTw and principal_series_full_statement. The representation of x is the private pt_part (Z.Ord.core_ordinary_of_literal_automorphy_n: ordinary_scalar_of_literal_automorphy_n gives sigma_v as a subquotient of PS(omega_v) with OrdChar, and PsFullAlt, from distinct type through Loc.ordCharIff and prin_series_irreducible_and_unique_subquotient, makes it an isomorphism: Lemma sha-hecke-ideal-ordinary-point and the paragraph after the properties of sigma). The properties of sigma are partUR2Wt, partUR2Unr, Hi.partUR2V0 and Hi.partUR2Type_alt (with R.type, the type of Lemma sha-hecke-data(1)); the compatibility is s1 with lgAway' and woGd'. The twist is the private tw_aut (Lemma sha-hecke-ideal-twist-global) with beta^(2N+1) = 1 from the odd order l^k; rho' is continuous and semisimple by galoisRep_irreducibility_criterion; det_tw and part_tw give the determinant and the arithmetic weight of x'. Then upU_w_alt: (a) t5, hSp, hTw, satFrobT_alt; (b) t2 with partUR2Wt; (c) t3 with pr_ine and partUR2Unr; (d) t4 with Up.comp_ex and partUR2V0; (e) t4, hOr (with hPs); (f) t5, partUR2Type_alt, hTy, hTw.
\emph{The base of the Taylor-Wiles data.} For (7), let $D=n[L^+:\mathbf Q]$. By local class field theory, as in the proof of Lemma~\ref{sha-hecke-algebra}, there are places $v_k\mid\ell$, indices $i_k$, elements $\gamma_k\in I_{v_k}\subset D_{v_k}$ and $c_k\in\mathcal O$, for $1\le k\le D$, with $\bar c_k\bar\psi_{v_k,i_k}(\gamma_k)=1$ in $k$ for the residue $\bar c_k$ of $c_k$, and an isomorphism of $\mathcal O$-algebras from $\mathcal O[[X_1,\dots,X_D]]$ onto $\Lambda$ sending each $X_k$ to $c_k\kappa_{v_k,i_k}(\gamma_k)-1$. In particular $c_k\in\mathcal O^\times$. Let $\lambda$ be its composite with the map $\Lambda\to B$; it sends $X_k$ to $c_k\psi^{\mathrm u}_{v_k,i_k}(\gamma_k)-1$. If $M=0$, let $M_0=0$; otherwise let $M_0=M$, a $B$-module through $f:B\to\mathbf T\subset\operatorname{End}_\Lambda(M)$. By Hida's theorem, as in the proof of Lemma~\ref{sha-hecke-algebra}, $M$ is finite and free over $\Lambda$, so $M_0$ is finite and free over $\mathcal O[[X_1,\dots,X_D]]$ through $\lambda$. The annihilator of $M_0$ in $B$ is $\ker f=J$, or $B=J$ if $M=0$, so the support of $M_0$ in $\operatorname{Spec}B$ is the zero locus of $J$.

So $(\lambda,M_0)$ is a Taylor-Wiles base for $J$ (Definition~\ref{sha-hecke-ideal-tw-base}).

\emph{The data at the augmented levels.} Let $\sigma_1,\dots,\sigma_q$, $H$, the places $w_{N,j}$ and the elements $\phi_{N,j}$ and $\tau_{N,j}$ be as in Lemma~\ref{sha-patching-places}, let $v_{N,j}$ be the place of $L^+$ below $w_{N,j}$, and let $e=n!$. Choose $\tilde u_{j,i}\in\mathcal O$ lifting $\bar\chi_i(\sigma_j)$; these are the eigenvalues of $\bar D$ at the $\phi_{N,j}$, since the image of $\phi_{N,j}$ in $\Gamma$ lies in $\sigma_jH$ and $H$ kills the $\bar\chi_i$. Fix $N$, and let $\Gamma_N$ be the Galois group of the level augmented by the $v_{N,j}$, with $p_N:\Gamma_N\to\Gamma$.

(i) There is a module $(R,P)$ at level $\Gamma_N$ for the places $w_{N,j}$, with the base $(\lambda,M_0)$, that is separated by its points for $\iota$ (Definition~\ref{sha-hecke-ideal-tw-level-module}). Indeed, the conditions of Situation~\ref{sha-hecke-situation} on $(D^{\mathrm u},\psi^{\mathrm u})$ pull back along $p_N$, so $(D^{\mathrm u}\circ p_N,\psi^{\mathrm u})$ is an ordinary pair for $\Gamma_N$ over $B$. Hence there is a universal ring $B'$ of ordinary pairs for $\Gamma_N$, with universal pair $(D',\psi')$, and an $\mathcal O$-algebra map $\beta:B'\to B$ carrying $(D',\psi')$ to $(D^{\mathrm u}\circ p_N,\psi^{\mathrm u})$. Let $P$, $\rho$, $\alpha$, $\lambda'$ and $e_0$ be as in Lemma~\ref{sha-hecke-ideal-tw-augmented}, for $\iota$, the data $\tau_0$, $\theta$, $a$, $\mathcal K_0$ and $M$ of the base step, $\mu$ the action of $B$ on $M_0$, $B'$, $\beta$ and the places $v_{N,j}$ and $w_{N,j}$. Take $R=B'$ with $(D',\psi')$; the map $\Lambda_D\to B'$ that is $\lambda'$ composed with the isomorphism $\Lambda_D\to\Lambda$ of the base step, which sends $X_k$ to $c_k\psi'_{v_k,i_k}(\gamma_k)-1$ since $\gamma_k\in I_{v_k}$; the map $\beta$; the $B'$-module $P$, through $\alpha$, with $\rho$, which is $B'$-linear; and $\epsilon=e_0$, for which $\epsilon(tm)=\beta(t)\epsilon(m)$ by (3) of that lemma. The groups $\mathrm{Iw}_{\ell',w}\subset\mathrm{Iw}_w$ are compact open products of subgroups of the factors, and $\mathrm{Iw}_{\ell',w}$ is normal in $\mathrm{Iw}_w$ (Lemma~\ref{sha-hecke-ideal-tw-diamond}(1)); so $P^{\mathrm{Iw}_{\ell',w}}$ is finite over $\Lambda_D$ and projective over $\Lambda_D[\mathrm{Iw}_w/\mathrm{Iw}_{\ell',w}]$, by (1) and (2) of that lemma. By (4) of that lemma, the module is separated by its points.

(ii) The places $w_{N,j}$ do not lie above $\ell$, as the $v_{N,j}$ are outside $S$. By Lemma~\ref{sha-patching-places}, $\phi_{N,j}$ and $\tau_{N,j}$ are the images of elements $\tilde\phi_j$ and $\tilde\tau_j$ of the decomposition group $D_{A_j}\subset\operatorname{Gal}(\bar L/L)$ of a valuation ring $A_j$ over $w_{N,j}$, with $\tilde\phi_j^{-1}$ a Frobenius element and $\tilde\tau_j$ a topological generator of the $\ell$-part of tame inertia. Lemma~\ref{sha-hecke-ideal-tw-local}, applied to $L$, $w_{N,j}$, $A_j$, $\tilde\phi_j$ and $\tilde\tau_j$, gives $x_j$ and $\delta_j$, with $q_{w_{N,j}}=\operatorname{N}w_{N,j}$. For a point $y$ of $(R,P)$, the representation $\rho_y$, viewed on $\operatorname{Gal}(\bar L/L)$, and $\sigma_{y,j}$ satisfy the hypothesis of that lemma, by the definition of a point, and $\iota\rho_y(\tilde\phi_j)=\iota\rho_y(\phi_{N,j})$, $\iota\rho_y(\tilde\tau_j)=\iota\rho_y(\tau_{N,j})$. So $(R,P)$ satisfies the local conditions for the $x_j$, $\delta_j$, $\phi_{N,j}$, $\tau_{N,j}$ and $\iota$.

(iii) Since $\operatorname{N}w_{N,j}\equiv1\pmod{\ell^N}$, Lemma~\ref{sha-hecke-ideal-tw-operators} gives an operator module of level $N$ for $\Gamma_N$ that satisfies the relations of exponent $n!$ for $b_j=\operatorname{N}w_{N,j}$ and the $\phi_{N,j}$ and $\tau_{N,j}$.

(iv) Lemma~\ref{sha-hecke-ideal-tw-patching}, applied to this operator module and the $\tilde u_{j,i}$, gives $R_N=R'$ with its ordinary pair for $\Gamma_N$, the $U_{j,i}=U'_{j,i}$, $s_N=s'$, $M_N=M'$ with $a_j\ge N$, the map $\theta:R_N\to B[u]$, which kills the $s'(y_{j,i})$, the maps between $M_N/(y_{j,i})M_N$ and $B[u]\otimes_BM_0$ with composites $f_N=f'$, and the relations with $e=n!$. These are the data of (7) at level $N$.
% Base: the private s4 of BaseObjectHomExt.lean: Hi.levelUR2WtVars (W.locCFT from global Artin reciprocity, then the private weight_variables_exist with W.lam_ex_free), Hi.levelUR2BaseAct (mu = the Hecke action through f, ker mu = J, or J = B and mu = 0 when HSet.MK is zero), hiComm (MK finite and free over Lambda) and Hi.levelUR2BasePack (b : TWBase, with supp from Module.annihilator = ker mu = J). The augmented levels follow Pa.has_automorphy_of_statement for each N: (i) is the LevMod of s4 (Definition sha-hecke-ideal-tw-level-module) with ClassicalAlt, for V = Q.V N, h = homComp (Q.e N) (T.h N); (ii) is Z.Win.locDataT_alt (Lemma sha-hecke-ideal-tw-local) for Q.fr N j and Q.τ N j, giving LocalLevelConditions with the ι of ClassicalAlt; (iii) is Pa.Lev.opOf_alt (Lemma sha-hecke-ideal-tw-operators); then Pa.Lev.opTransport moves the OpMod and OpRel from globOf r (Q.V N) to T.GN N along the injective, hence bijective, Q.e N, which the book treats as the identification of Gamma_N with the Galois group of the augmented level; (iv) is Pa.Lev.opToTW (Lemma sha-hecke-ideal-tw-patching), with ut chosen by residue_surjective as lifts of chib i (T.sg j). Step (i) is the LevMod of s4 at the augmented level: LvR.pairUp (W.exists_determinantPair_pullback) for the pair over B, W.b0Ex for the universal ring B' of the augmented level (it assumes W.phiAll from the published CohP), beta from B'.univ, the private s3 (Lemma sha-hecke-ideal-tw-augmented, whose remaining open part is there), levelUR2IwLev, and levelUR2ModPack_alt (module_structure_of_commuting and levelUR2ProjTr), with ClassicalAlt from the points of s3; eM is e0 followed by the eb of levelUR2BasePack. The relation (a) of step (iii) is Pa.Lev.twoOf_alt (Lemma sha-hecke-ideal-tw-comparison).
% Contributing representations: Z.Desc.W.upU_w_alt (SphEq is (a), HasWeightAt (b), IsUnramU (c), Kc n 1 at v0 (d), PS and OrdChar (e), StrongEquivalenceHida (f)). The forms: the private Up.types (Lemma sha-hecke-ideal-inertia) and Hi.nonempty_unitary_hecke_parts (Lemma sha-hecke-ideal-forms). rho_sigma: Z.Desc.P.partU_alt = partX_alt (Lemma sha-hecke-ideal-unitary; Y.frob from satake_local_of_printed_results, satT and satFrobT_alt), then hfr and AlgebraTools.Det.extDense. Away from l: Up.hecke_cong (hsph, with chi_cong and IsRes.frob). At l: Up.ordChar_unit_mul_finite_component_eq_one (hfin, with R.theta, LamObj.toPair_spec and ord_val) and Up.norm_lowVal_inv_mul_psEig_eq_one_and_norm_sub_lt_one (pt_unit, pt_res, IsRes.unif). The point: Hi.coreUp_alt (Lemma sha-hecke-ideal-eigenform), then core_eigen_of_weight_char_hecke_level, ord_uniq, pt_ext with CA.exists_adjoin_coeff_character_sub_mem_maximalIdeal_pow.
% (3): Z.Desc.W.twisted_point_property_of_isJ_arith_rep_autCsd = point_property_of_isJ_arith_rep_autCsd = jPT_of_alt with recTExists, compatTOf (Lemma sha-hecke-ideal-attached), plainRec, guard_of and unlOf (the plain local components), sphT_alt (Lemma sha-hecke-ideal-spherical), the private ker_le (rho unramified outside S), DescWeak (the descent), same_of_prints_of_rT (Same: sigma_w = pi_w at split w and the infinitesimal characters; the Lean derives it from the Speh decomposition of the base change and Galois representations, the book cites strong multiplicity one), the private above (Lemma sha-hecke-ideal-above), dual_of_weights for (b), unram_of_fixed for (c), Kc n 1 inside KO for (d), then upU_w_alt.
% (4): Hi.arithmeticPointDecomposition_of_isJ, through hecke_ideal_end (absPt with core_hecke_of_automorphic_hecke_dense_weight; hP.aut) and Z.Win.Gu.endShape'_alt (Lemma sha-hecke-ideal-speh).
% (7): Pa.has_automorphy_of_statement, from hLev = Pa.LevelAlt (Hi.level_of_is_j_unitary_hecke_data = level_of_isJ_unitaryHeckeData through the private s4: the base b : TWBase and, at each augmented level, LevMod with ClassicalAlt), Z.Win.locDataT_alt (the relations LocRelT at the w_{N,j}), Pa.Lev.opOf_alt, opTransport and opToTW (TWRing and TWMaps with ex = n!). Property (1) at v0 is Hi.point_rep_unramified_at_auxiliary_place.
\end{proof}

\subsection{Normalized arcs and Taylor-Wiles places}\label{subsection-sha-normalized-arcs}

\begin{lemma}\label{sha-patching-blocks}
Let $\kappa\subset\kappa'$ be finite fields of characteristic different from $2$ such that every element of $\kappa$ is a square in $\kappa'$, let $\mathcal K$ and $\mathcal K'$ be the fraction fields of $\kappa[[T]]$ and $\kappa'[[T]]$, and let $j:\kappa[[T]]\to\kappa'[[T]]$ be the homomorphism that is the inclusion on $\kappa$ and sends $T$ to $T^2$, extended to $j:\mathcal K\to\mathcal K'$. Let $G$ be a group, $\sigma:G\to G$ a surjective map, $m\ge0$ an integer, $r:G\to GL_m(\mathcal K)$ a homomorphism, $H$ a symmetric invertible $m\times m$ matrix over $\mathcal K$ and $u:G\to\kappa[[T]]$ a function, such that ${}^tr(g)\,H\,r(\sigma(g))=u(g)H$ for all $g\in G$, and such that there is a nonzero $d\in\kappa[[T]]$ with $d\,r(g)\in M_m(\kappa[[T]])$ for all $g\in G$. Then there are $Q\in GL_m(\mathcal K')$ and a homomorphism $r':G\to GL_m(\kappa'[[T]])$ with ${}^tQ\,j(H)\,Q=1$ and $j(r(g))\,Q=Q\,r'(g)$ for all $g\in G$.
\end{lemma}
\leanref{Development/AutomorphicArithmetic/AutomorphicVarieties/PatchingModules/DeterminantPadicFree/OneBlock.lean}{UpperTheory.V.Pa.Det.oneBlock}
\begin{proof}
Write $A=\kappa[[T]]$. A \emph{lattice} is a finitely generated $A$-submodule of $\mathcal K^m$ that spans $\mathcal K^m$; it is \emph{stable} if $r(g)Y\subset Y$ for all $g$; and for a symmetric invertible matrix $H_1$ its \emph{dual} is $Y^\vee=\{x\in\mathcal K^m:{}^txH_1y\in A\text{ for all }y\in Y\}$, a lattice with $Y^{\vee\vee}=Y$.

\emph{A stable lattice.} The $A$-span $Y_0$ of the vectors $r(g)e_i$, for $g\in G$ and the standard basis vectors $e_i$, is stable, contains $A^m$ and is contained in $d^{-1}A^m$, so it is a lattice. Let $Q_0$ be the matrix whose columns are a basis of $Y_0$, and put $r_0(g)=Q_0^{-1}r(g)Q_0$ and $H_1={}^tQ_0HQ_0$. Then $r_0(g)\in GL_m(A)$, $H_1$ is symmetric and invertible, and ${}^tr_0(g)\,H_1\,r_0(\sigma(g))=u(g)H_1$.

\emph{Duals of stable lattices.} Let $Y$ be a lattice stable under $r_0$. Then $Y^\vee$ is stable: for $x\in Y^\vee$, $g\in G$ and $y\in Y$, the vector $y'=r_0(\sigma(g))^{-1}y$ lies in $Y$, and ${}^t(r_0(g)x)\,H_1\,y={}^t(r_0(g)x)\,H_1\,r_0(\sigma(g))\,y'=u(g)\,{}^txH_1y'\in A$.

\emph{An almost self-dual lattice.} The lattice $Y_1=A^m\cap(A^m)^\vee$ is stable and $Y_1\subset Y_1^\vee$; since both are lattices, $T^kY_1^\vee\subset Y_1$ for some $k\ge0$. Suppose that $Y$ is a stable lattice with $Y\subset Y^\vee$ and $T^{k+1}Y^\vee\subset Y$ for some $k\ge1$, and let $Y'=Y+T^kY^\vee$. Then $Y'$ is a stable lattice. It satisfies $Y'\subset Y'^\vee$: the pairings of $Y$ with $Y$ and with $T^kY^\vee$ lie in $A$, and $T^{2k}Y^\vee\subset T^{k+1}Y^\vee\subset Y=Y^{\vee\vee}$, so the pairings of $T^kY^\vee$ with itself lie in $A$. And $T^kY'^\vee\subset T^kY^\vee\subset Y'$, since $Y\subset Y'$. By induction on $k$ there is a stable lattice $Y$ with $TY^\vee\subset Y\subset Y^\vee$.

\emph{An adapted basis.} By the theory of elementary divisors over the discrete valuation ring $A$, applied to $TY^\vee\subset Y\subset Y^\vee$, there are a basis $b_1,\dots,b_m$ of $Y^\vee$ and a set $S$ of indices such that $Y$ has the basis $c_1,\dots,c_m$ with $c_i=Tb_i$ for $i\in S$ and $c_i=b_i$ for $i\notin S$. The matrix $M=({}^tb_iH_1c_{i'})_{i,i'}$ has entries in $A$, since $b_i\in Y^\vee$, and is invertible over $A$, since $Y^\vee$ is the dual of $Y$.

\emph{The ramified extension.} Put $s_i=T$ for $i\in S$ and $s_i=1$ for $i\notin S$, and let $P\in GL_m(\mathcal K')$ be the matrix with columns $s_ij(b_i)$. Since $j(T)=T^2$, the $(i,i')$ entry of $H'={}^tP\,j(H_1)\,P$ is $s_is_{i'}\,j({}^tb_iH_1b_{i'})$, which is $s_i\,j(M_{ii'})$ if $i'\notin S$, is $T\,j(M_{i'i})$ if $i\notin S$ and $i'\in S$ (by symmetry of $H_1$), and is $T^2j(T^{-1}M_{ii'})=j(M_{ii'})$ if $i,i'\in S$. So $H'$ is symmetric with entries in $\kappa[[T]]\subset\kappa'[[T]]$, and its determinant is $j(\det M)$, a unit. Since $r_0(g)$ preserves $Y^\vee$ and $Y$, there are matrices $R,R'$ over $A$ with $r_0(g)b_{i'}=\sum_iR_{ii'}b_i$ and $r_0(g)c_{i'}=\sum_iR'_{ii'}c_i$. The $(i,i')$ entry of $P^{-1}j(r_0(g))P$ is $s_i^{-1}j(R_{ii'})s_{i'}$. It lies in $\kappa'[[T]]$: this is clear unless $i\in S$ and $i'\notin S$, and in that case $R_{ii'}=TR'_{ii'}$, so the entry is $T\,j(R'_{ii'})$. So $r''(g)=P^{-1}j(r_0(g))P$ lies in $GL_m(\kappa'[[T]])$.

\emph{An orthonormal basis.} Since $2$ is a unit in the local ring $\kappa[[T]]$ and $H'$ is a symmetric matrix over it with unit determinant, there are $P_1\in GL_m(\kappa[[T]])$ and units $u_1,\dots,u_m$ of $\kappa[[T]]$ with ${}^tP_1H'P_1=\operatorname{diag}(u_1,\dots,u_m)$, by diagonalizing the form one anisotropic unit vector at a time. The residue of each $u_i$ is a square in $\kappa'$, so by Hensel's lemma, as $2$ is a unit, $u_i=a_i^2$ with $a_i\in\kappa'[[T]]^\times$. Let $P_2=P_1\operatorname{diag}(a_1^{-1},\dots,a_m^{-1})\in GL_m(\kappa'[[T]])$. Then $Q=j(Q_0)PP_2$ and $r'(g)=P_2^{-1}r''(g)P_2$ satisfy ${}^tQ\,j(H)\,Q={}^tP_2H'P_2=1$ and $j(r(g))Q=Q\,r'(g)$.
% Stable lattice: Det.latInt (exists_stable_lattice). Duals: map_mem_dualSubmodule_of_bilinForm_similitude. Almost self-dual: AlgebraTools.almost_selfdual (the induction key). Adapted basis: AlgebraTools.smith_zero_one, and M with its inverse N in selfdual_ramified_mat. Ramified extension: AlgebraTools.selfdual_ramified (j_1 with s = T, s^2 = j_1(T)). Orthonormal basis: AlgebraTools.orthonormal_local, from diag_local and sqrt_of_residue over the Henselian ring K'[[T]].
\end{proof}

\begin{lemma}\label{sha-patching-normalization}
In Situation~\ref{sha-hecke-situation}, let $(\rho,\psi)$ be a nice arc of type I or E over $\kappa[[T]]$, and for $u\in S_{\mathrm a}$ let $I'_u\subset I_u\cap\bigcap_i\ker\bar\chi_i$ be subgroups. Then there are a finite field $\kappa'\supset\kappa$, an injective local finite homomorphism $j:\kappa[[T]]\to\kappa'[[T]]$, an integer $1\le a\le n$, with $a=n$ for type I, and an arc $(\rho',\psi')$ over $\kappa'[[T]]$ with the same determinant and characters after $j$, still satisfying (N1), (N3) and (N4), which is \emph{normalized}: ${}^t\rho'(g)\rho'(cgc^{-1})=\varepsilon(g)^{1-n}$; $\rho'$ takes values in the block diagonal matrices with blocks of sizes $a$ and $n-a$, and for every block diagonal matrix $Z$ there is $s\ge0$ such that $T^sZ$ lies in the $\kappa'[[T]]$-span of $\rho'(\Gamma')$; $\rho'$ is trivial on the $I'_u$; each $\psi'_{v,i}$ is trivial on $\mu_v$; and for each $v\mid\ell$ there is $w_v\in GL_n(\kappa'[[T]])$, each of whose columns lies in one of the two blocks, with $w_v^{-1}\rho'(g)w_v$ upper triangular with diagonal $\psi'_{v,1}(g),\dots,\psi'_{v,n}(g)$ for $g\in D_v$.
\end{lemma}
\leanref{Development/AutomorphicArithmetic/AutomorphicVarieties/PatchingModules/TaylorWilesObstructionBound.lean}{UpperTheory.V.Pa.Norm}
\leanref{Development/AutomorphicArithmetic/AutomorphicVarieties/PatchingModules/AutomorphicDetection/ConjugationDescendsToSubring.lean}{UpperTheory.V.Pa.norm}
\leanref{Development/AutomorphicArithmetic/AutomorphicVarieties/PatchingModules/PairedDualFirstOrderDualMatrixSpace.lean}{UpperTheory.V.Pa.NP}
\begin{proof}
\emph{Block form.} By (N2), there are $1\le a\le n$, with $a=n$ for type I, and $Q\in GL_n(\mathcal K)$ such that $Q\rho Q^{-1}$ takes values in the block diagonal matrices with blocks of sizes $a$ and $n-a$, its $\mathcal K$-span on $\Gamma'$ is the whole algebra of block diagonal matrices, and the symmetric matrix $H={}^tQ^{-1}PQ^{-1}$, for which ${}^t(Q\rho Q^{-1})(g)\,H\,(Q\rho Q^{-1})(cgc^{-1})=\varepsilon(g)^{1-n}H$, is block diagonal.

\emph{The extension.} Let $j_1:\kappa[[T]]\to\kappa[[T]]$ be the homomorphism $T\mapsto T^2$, let $\kappa'$ be a finite extension of $\kappa$ in which every element of $\kappa$ is a square, with $j_2:\kappa[[T]]\to\kappa'[[T]]$ the inclusion, and put $j=j_2\circ j_1$, which is injective, local and finite. Let $\mathcal K'$ be the fraction field of $\kappa'[[T]]$.

\emph{The form.} There are a block diagonal $Q'\in GL_n(\mathcal K')$ with ${}^tQ'\,j(H)\,Q'=1$ and a representation $\rho':\Gamma_0\to GL_n(\kappa'[[T]])$, block diagonal, with $j(Q\rho Q^{-1})Q'=Q'\rho'$. Indeed, $\ell$ is odd, and there is a nonzero $d\in\kappa[[T]]$ such that $d\,QZQ^{-1}$ has entries in $\kappa[[T]]$ for every matrix $Z$ over $\kappa[[T]]$, obtained by clearing the denominators of $Q$ and $Q^{-1}$. Since $Q\rho Q^{-1}$ and $H$ are block diagonal, Lemma~\ref{sha-patching-blocks}, applied to each block of $Q\rho Q^{-1}$ and of $H$ with $\sigma(g)=cgc^{-1}$ and $u(g)=\varepsilon(g)^{1-n}$, gives the blocks of $Q'$ and $\rho'$. So $\rho'$ is conjugate over $\mathcal K'$ to $j\circ\rho$, by $T'=j(Q)^{-1}Q'$, and ${}^t\rho'(g)\rho'(cgc^{-1})=\varepsilon(g)^{1-n}$. Conjugating by $Q'$, which is block diagonal, shows that the $\mathcal K'$-span of $\rho'(\Gamma')$ is the algebra of block diagonal matrices; clearing denominators gives, for each block diagonal $Z$ over $\kappa'[[T]]$, an $s$ with $T^sZ$ in the $\kappa'[[T]]$-span.

\emph{The other properties.} The entries of $\rho'$ are continuous, since $\rho'$ is a conjugate of $j\circ\rho$. The characteristic polynomials of $\rho'$ are the images under $j$ of those of $\rho$, so (N3) holds and the determinants agree after $j$. If $X$ over $\kappa'[[T]]$ satisfies $\varepsilon(g)\rho'(g)X=X\rho'(g)$ for $g\in D_{v_0}$, then $T'XT'^{-1}$ satisfies the same relation for $j\circ\rho$ over $\mathcal K'$; this relation is a system of linear equations with coefficients in the image of $\mathcal K$, whose only solution over $\mathcal K$ is $0$ by (N4) for $\rho$, so its only solution over $\mathcal K'$ is $0$, and (N4) holds. $\rho'$ is trivial on the $I'_u$ because $\rho$ is. Put $\psi'_{v,i}=j\circ\psi_{v,i}$; they satisfy (N1) because $j$ is injective. Each $\psi_{v,i}$ is trivial on $\mu_v$: it is continuous, since on $D_v$ it is a combination of entries of $\rho$ with coefficients from the triangularizing matrix, so it kills the closed commutator subgroup of $D_v$; an element of $\mu_v$ has an $\ell$-power in that subgroup; and $\kappa[[T]]^\times$ has no nontrivial elements of $\ell$-power order, since $(x-1)^{\ell^k}=x^{\ell^k}-1$ in characteristic $\ell$. So the $\psi'_{v,i}$ are trivial on $\mu_v$. Finally, for $v\mid\ell$, $\rho'$ restricted to $D_v$ is conjugate over $\mathcal K'$ to an upper triangular representation with diagonal $\psi'_{v,1},\dots,\psi'_{v,n}$, by the conjugate by $T'$ of the triangular form of $\rho$; since $\rho'$ is block diagonal and these characters are pairwise distinct by (N1), the flag has a basis in $\kappa'[[T]]^n$ each of whose vectors lies in one block, and these vectors are the columns of $w_v$.
% Block form: Det.blockForm. The extension: Det.sqExt (j1) and Det.resSq (j2). The form: Det.twoBlock, from Det.oneBlock on each block, with the common denominator from bdd_den. The span over the ring: Det.spanInt. (N4): conjugation_descends_to_subring with Det.solBase. Trivial on mu_v: Det.psiMu. The flag: Det.flagBasis. All in Pa.norm (ConjugationDescendsToSubring.lean).
\end{proof}

\begin{definition}\label{sha-patching-selmer}
In Situation~\ref{sha-hecke-situation}, let $x$ be a normalized arc over $\kappa[[T]]$ as in Lemma~\ref{sha-patching-normalization}, with representation $\rho$, blocks of sizes $a$ and $n-a$, subgroups $I'_u$ and matrices $w_v$. An $n\times n$ matrix is \emph{even} if it is block diagonal with blocks of these sizes, and \emph{odd} if its two diagonal blocks vanish. For a matrix $Y$ write $\delta Y(g)=\rho(g)Y\rho(g)^{-1}-Y$.
\begin{enumerate}
\item The \emph{Selmer group} of $x$ is the $\kappa[[T]]$-module of maps $z:\Gamma_0\to M_n(\kappa[[T]])$ with odd values, continuous for the $T$-adic topology, such that $z(gh)=z(g)+\rho(g)z(h)\rho(g)^{-1}$ and $z(cgc^{-1})=-{}^tz(g)$ for $g,h\in\Gamma_0$, $z$ vanishes on each $I'_u$, and for each $v\mid\ell$ there is an odd matrix $Y_v$ such that $w_v^{-1}(z(g)-\delta Y_v(g))w_v$ is upper triangular for $g\in D_v$, with zero diagonal for $g\in\mu_v$. Its \emph{coboundaries} are the maps $\delta Y$ with $Y$ odd and ${}^tY=-Y$.
\item A \emph{Selmer basis} of $x$ consists of elements $\kappa_1,\dots,\kappa_g$ of the Selmer group, with matrices $Y_{j,v}$ as in (1), such that $\sum_jc_j\kappa_j$, for $c_j\in\kappa[[T]]$, is a coboundary only if all $c_j$ vanish, and such that for every $z$ in the Selmer group there are $s\ge0$ and $c_1,\dots,c_g\in\kappa[[T]]$ with $T^sz-\sum_jc_j\kappa_j$ a coboundary.
\item Given a Selmer basis, let $R_g$ be the $\kappa[[T]]$-algebra that is free as a module on $1$, $\epsilon_1,\dots,\epsilon_g$, $\epsilon'_1,\dots,\epsilon'_g$ and $\eta$, with $\epsilon_j\epsilon'_j=\eta$ for each $j$ and all other products of two of $\epsilon_1,\dots,\epsilon_g,\epsilon'_1,\dots,\epsilon'_g,\eta$ equal to $0$. The \emph{first-order deformation} is the homomorphism $\rho_1:\Gamma_0\to GL_n(R_g)$, $\rho_1(g)=(1+\sum_j\epsilon_j\kappa_j(g))\rho(g)$. Let $V^*$ be the module of $n\times n$ matrices over $R_g$ whose entries lie in $\kappa[[T]]\eta$ in the diagonal blocks and in $\sum_j\kappa[[T]]\epsilon'_j$ in the off-diagonal blocks, and view $\varepsilon$ as a character with values in $\mathbf F_\ell^\times\subset\kappa$. A \emph{dual cocycle} is a map $z:\Gamma_0\to V^*$, continuous for the $T$-adic topology, with $z(gh)=z(g)+\varepsilon(g)\rho_1(g)z(h)\rho_1(g)^{-1}$ and $z(cgc^{-1})=-\varepsilon(c)\,{}^tz(g)$; the \emph{dual coboundaries} are the maps $g\mapsto\varepsilon(g)\rho_1(g)Z\rho_1(g)^{-1}-Z$ with $Z\in V^*$ and $-\varepsilon(c)\,{}^tZ=Z$. For $\sigma\in\Gamma_0$, the \emph{pairing} $\langle z,\sigma\rangle$ is the coefficient of $\eta$ in $\operatorname{tr}(z(\sigma)\rho_1(\sigma))$.
\item A subset $\Sigma\subset\Gamma_0$ \emph{detects the dual Selmer group} of $x$, for a given Selmer basis, if for every dual cocycle $z$ with $\langle z,\sigma\rangle=0$ for all $\sigma\in\Sigma$ there is $s\ge0$ such that $T^sz$ is a dual coboundary.
\end{enumerate}
\end{definition}
\leanref{Development/AutomorphicArithmetic/AutomorphicVarieties/PatchingModules/PairedDualFirstOrderDualMatrixSpace.lean}{UpperTheory.V.Pa.SelZ}
\leanref{Development/AutomorphicArithmetic/AutomorphicVarieties/PatchingModules/PairedDualFirstOrderDualMatrixSpace.lean}{UpperTheory.V.Pa.Kap}
\leanref{Development/AutomorphicArithmetic/AutomorphicVarieties/PatchingModules/PairedDualFirstOrderDualMatrixSpace.lean}{UpperTheory.V.Pa.DD}
\leanref{Development/AutomorphicArithmetic/AutomorphicVarieties/PatchingModules/PairedDualFirstOrderDualMatrixSpace.lean}{UpperTheory.V.Pa.Kap.ρ1}
\leanref{Development/AutomorphicArithmetic/AutomorphicVarieties/PatchingModules/PairedDualFirstOrderDualMatrixSpace.lean}{UpperTheory.V.Pa.pairedDualFirstOrderDualMatrixSpace}
\leanref{Development/AutomorphicArithmetic/AutomorphicVarieties/PatchingModules/ThetaCoefficientMapEq.lean}{UpperTheory.V.Pa.TwZ}
\leanref{Development/AutomorphicArithmetic/AutomorphicVarieties/PatchingModules/ThetaCoefficientMapEq.lean}{UpperTheory.V.Pa.TwB}
\leanref{Development/AutomorphicArithmetic/AutomorphicVarieties/PatchingModules/ThetaCoefficientMapEq.lean}{UpperTheory.V.Pa.ell}
\leanref{Development/AutomorphicArithmetic/AutomorphicVarieties/PatchingModules/ThetaCoefficientMapEq.lean}{UpperTheory.V.Pa.Detects}

\begin{lemma}\label{sha-patching-twisted}
Let $F$ be a field of characteristic different from $2$, $G$ a group, $G'\subset G$ a subgroup containing every commutator, $\chi:G\to F^\times$ a nontrivial character that is trivial on $G'$, and $\sigma$ an automorphism of $G$ with $\sigma^2=1$ and $\chi\circ\sigma=\chi$. Let $1\le a\le n$; call an $n\times n$ matrix over $F$ \emph{even} if it is block diagonal with blocks of sizes $a$ and $n-a$, and \emph{odd} if its two diagonal blocks vanish. Let $r:G\to GL_n(F)$ be a homomorphism with even values such that the $F$-span of $r(G')$ is the space of even matrices and ${}^tr(g)\,r(\sigma(g))$ is a scalar matrix for every $g$. For a function $\theta:G\to F$ and a matrix $Y$ write $\delta_\theta Y(g)=\theta(g)r(g)Yr(g)^{-1}-Y$, and call $z:G\to M_n(F)$ a \emph{$\theta$-cocycle} if $z(gh)=z(g)+\theta(g)r(g)z(h)r(g)^{-1}$ for all $g,h$. Let:
\begin{enumerate}
\item $\kappa_1,\dots,\kappa_m$ be $1$-cocycles with odd values and $\kappa_k(\sigma(g))=-{}^t\kappa_k(g)$, such that $\sum_kc_k\kappa_k=\delta_1Y$, with $c_k\in F$ and $Y$ odd with ${}^tY=-Y$, only if all $c_k$ vanish;
\item $s\in F$ with $s^2=1$, and $\lambda_1,\dots,\lambda_m$ $\chi$-cocycles with odd values and $\lambda_k(\sigma(g))=s\,{}^t\lambda_k(g)$;
\item $\varphi:G\to M_n(F)$ with even values, $\varphi(\sigma(g))=s\,{}^t\varphi(g)$, and
\[\varphi(gh)=\varphi(g)+\chi(g)\Big(r(g)\varphi(h)r(g)^{-1}+\sum_k\big[\kappa_k(g),\,r(g)\lambda_k(h)r(g)^{-1}\big]\Big)\]
for all $g,h$, such that $\operatorname{tr}(\varphi(g)r(g))+\sum_k\operatorname{tr}(\lambda_k(g)\kappa_k(g)r(g))=0$ for $g\in G'$.
\end{enumerate}
Then there are an even matrix $Y_0$ and odd matrices $Y_1,\dots,Y_m$, with $s\,{}^tY_k=Y_k$ for $0\le k\le m$, such that $\lambda_k=\delta_\chi Y_k$ for $1\le k\le m$ and $\varphi(g)=\delta_\chi Y_0(g)+\chi(g)\sum_k[\kappa_k(g),r(g)Y_kr(g)^{-1}]$ for all $g$.
\end{lemma}
\leanref{Development/AutomorphicArithmetic/AutomorphicVarieties/PatchingModules/DeterminantPadicFree/DetAlg.lean}{UpperTheory.V.Pa.Det.detAlg}
\begin{proof}
Fix $g_0\in G$ with $\chi(g_0)\ne1$. Let $S$ be the diagonal matrix with entries $1$ in the first block and $-1$ in the second; a matrix $X$ is even exactly when $SXS=X$ and odd exactly when $SXS=-X$, and $\delta_\theta(SYS)=S(\delta_\theta Y)S$ since $r$ has even values. For a finitely supported function $c:G'\to F$ and a function $z$ on $G'$ write $r(c)=\sum_gc(g)r(g)$ and $z\cdot c=\sum_gc(g)z(g)r(g)$; call $c$ a \emph{relation} if $r(c)=0$, and say that $z$ \emph{kills relations at} $(i,j)$ if $(z\cdot c)_{ij}=0$ for every relation $c$. Two facts are used throughout. (i) If a matrix $X$ commutes with $r(h)$ for all $h\in G'$, it commutes with all even matrices; if moreover $X$ is odd, then $X=0$, as comparing the entries of $Xe_{ii}$ and $e_{ii}X$ for the diagonal unit matrices $e_{ii}$ shows. (ii) Since ${}^tr(g)r(\sigma(g))$ is scalar, $r(\sigma(g))Yr(\sigma(g))^{-1}={}^tr(g)^{-1}Y\,{}^tr(g)$, so $\delta_\theta Y(\sigma(g))={}^t\big(\delta_\theta({}^tY)(g)\big)$ when $\theta\circ\sigma=\theta$.

\emph{Step 1: derivations.} Let $z$ be a $1$-cocycle on $G'$ that kills relations at every $(i,j)$. Then $r(c)\mapsto z\cdot c$ is a well defined linear map $d$ on the span of $r(G')$, the even matrices. The cocycle relation gives $z(gh)r(gh)=z(g)r(g)\,r(h)+r(g)\,z(h)r(h)$, so $d$ is a derivation from the algebra of even matrices to $M_n(F)$. It is inner: with $\beta(k)$ the first index of the block of $k$, the matrix $Y=\sum_kd(e_{k\beta(k)})e_{\beta(k)k}$ satisfies $d(X)=YX-XY$ for even $X$. Hence $z(g)r(g)=Yr(g)-r(g)Y$, that is $z=\delta_1(-Y)$ on $G'$.

\emph{Step 2: moving relations.} Let $z$ be a $1$-cocycle on $G'$ that kills relations at $(i,j)$. For $u,v\in G'$ and a relation $c$, the function $g\mapsto c(u^{-1}gv^{-1})$ is again a relation, and the cocycle relation turns this into $(r(u)(z\cdot c)r(v))_{ij}=0$. By linearity $(A(z\cdot c)B)_{ij}=0$ for all even $A,B$, and with $A=e_{ii'}$ and $B=e_{j'j}$ this shows that $z$ kills relations at every $(i',j')$ with $i'$ in the block of $i$ and $j'$ in the block of $j$.

\emph{Step 3: separating the twist.} Let $\kappa$ be a $1$-cocycle and $\lambda$ a $\chi$-cocycle on $G$ with $\kappa+\lambda=\delta_1Y$ on $G'$. Put $c_0=\chi(g_0)$ and $Y_1=(c_0-1)^{-1}\big(c_0\,r(g_0)Yr(g_0)^{-1}-c_0\kappa(g_0)-\lambda(g_0)-Y\big)$. For $g\in G'$ the element $g'=g_0^{-1}gg_0$ lies in $G'$ and $gg_0=g_0g'$; expanding $\kappa$ and $\lambda$ at both sides of this equation, and using $\kappa+\lambda=\delta_1Y$ at $g$ and at $g'$, gives $(c_0-1)\kappa(g)=(c_0-1)\delta_1Y_1(g)$. So $\kappa=\delta_1Y_1$ and $\lambda=\delta_1(Y-Y_1)=\delta_\chi(Y-Y_1)$ on $G'$.

\emph{Step 4: inflation.} Let $\theta$ be $1$ or $\chi$, and let $\lambda$ be a $\theta$-cocycle on $G$ with odd values and $\lambda=\delta_\theta Y$ on $G'$. Then $\lambda=\delta_\theta Y'$ on $G$ for the odd matrix $Y'=\frac12(Y-SYS)$. Indeed $-\lambda=S\lambda S=\delta_\theta(SYS)$ on $G'$, so $\lambda=\delta_\theta Y'$ on $G'$. The difference $w=\lambda-\delta_\theta Y'$ is a $\theta$-cocycle with odd values that vanishes on $G'$. For $g\in G$ and $h\in G'$ we have $g^{-1}hg\in G'$, and expanding $w(hg)=w(g\cdot g^{-1}hg)$ gives $r(h)w(g)r(h)^{-1}=w(g)$. By (i), $w(g)=0$.

\emph{Step 5: polarization.} Let $P(Z)$ be the upper right corner of a matrix $Z$ (the entries $(i,j)$ with $i\le a<j$, the others set to $0$) and $P'(Z)$ its lower left corner; then ${}^tP(Z)=P'({}^tZ)$, and since $r$ has even values, $P$ and $P'$ carry $\theta$-cocycles to $\theta$-cocycles and commute with $\delta_\theta$. Let $\theta\circ\sigma=\theta$, let $t\in F$ with $t^2=1$, and let $z$ be a $\theta$-cocycle with odd values, $z(\sigma(g))=t\,{}^tz(g)$ and $P(z)=\delta_\theta Y$. Then $Y''=P(Y)+t\,{}^tP(Y)$ is odd, $t\,{}^tY''=Y''$, and $z=\delta_\theta Y''$: indeed $P(z)=P(\delta_\theta Y)=\delta_\theta P(Y)$, and by (ii) $P'(z(g))={}^tP({}^tz(g))=t\,{}^tP(z(\sigma(g)))=t\,{}^t\big(\delta_\theta P(Y)(\sigma(g))\big)=t\,\delta_\theta({}^tP(Y))(g)$, while $z=P(z)+P'(z)$ as $z$ is odd. The same holds with $P$ and $P'$ exchanged.

\emph{Step 6: corners.} Let $\alpha,\beta\in F^m$, put $K=\sum_k\alpha_k\kappa_k$ and $L=\sum_k\beta_k\lambda_k$, and suppose that $K+L$ kills relations at some $(i,j)$ with $i$ and $j$ in different blocks. Then $\alpha=0$, and $L=\delta_\chi Y$ for an odd $Y$ with $s\,{}^tY=Y$. Indeed, since $\chi$ is trivial on $G'$, $K+L$ is a $1$-cocycle on $G'$, so by Step 2 it kills relations at every entry of the corner containing $(i,j)$, say the upper right one. Because $r$ has even values, $(P(z)\cdot c)_{i'j'}$ is $(z\cdot c)_{i'j'}$ in that corner and $0$ outside it, so $P(K)+P(L)$ kills relations at every entry. By Step 1, $P(K)+P(L)=\delta_1Y$ on $G'$; by Step 3, $P(K)=\delta_1Y_1$ and $P(L)=\delta_\chi Y_2$ on $G'$; by Step 4 these hold on $G$ after replacing $Y_1,Y_2$ by odd matrices. By Step 5, with $t=-1$ for $K$ and $t=s$ for $L$, $K=\delta_1Y_K$ with $Y_K$ odd and ${}^tY_K=-Y_K$, and $L=\delta_\chi Y$ with $Y$ odd and $s\,{}^tY=Y$. By the hypothesis on the $\kappa_k$, $\alpha=0$.

\emph{Step 7: a relation between relations.} Let $c,c'$ be relations and let $i,j$ lie in different blocks. Then
\[\sum_k\Big((\kappa_k\cdot c)_{ij}(\lambda_k\cdot c')_{ji}+(\lambda_k\cdot c)_{ij}(\kappa_k\cdot c')_{ji}\Big)=0.\]
Indeed, let $\psi(g)=\varphi(g)r(g)+\sum_k\lambda_k(g)\kappa_k(g)r(g)$. For finitely supported $c_1,c_2$ on $G'$ let $c_1*c_2$ be their convolution. The cocycle relations of $\kappa_k$, $\lambda_k$ and $\varphi$ on $G'$, where $\chi=1$, give $r(c_1*c_2)=r(c_1)r(c_2)$, $\kappa_k\cdot(c_1*c_2)=(\kappa_k\cdot c_1)r(c_2)+r(c_1)(\kappa_k\cdot c_2)$, the same for $\lambda_k$, and
\[\psi\cdot(c_1*c_2)=(\psi\cdot c_1)r(c_2)+r(c_1)(\psi\cdot c_2)+\sum_k\big((\kappa_k\cdot c_1)(\lambda_k\cdot c_2)+(\lambda_k\cdot c_1)(\kappa_k\cdot c_2)\big),\]
where $\psi\cdot c=\sum_gc(g)\psi(g)$. The trace condition of (3) says that $\operatorname{tr}(\psi\cdot c_1)=0$ for every $c_1$. The diagonal unit matrices are even, so there are $d_i,d_j$ on $G'$ with $r(d_i)=e_{ii}$ and $r(d_j)=e_{jj}$. Apply the trace condition to $(c*d_j)*(c'*d_i)$: all terms containing $r(c)$ or $r(c')$ vanish, and $\operatorname{tr}(Ae_{jj}Be_{ii})=A_{ij}B_{ji}$, which gives the identity.

\emph{Step 8: the odd part.} We show that $\lambda_k=\delta_\chi Y_k$ with $Y_k$ odd and $s\,{}^tY_k=Y_k$, for each $k$. If $a=n$, the odd matrices are $0$, so $\lambda_k=0$ and $Y_k=0$. Otherwise let $i=1$ and $j=a+1$. For a relation $c'$, the combination $\sum_q\big((\lambda_q\cdot c')_{ji}\kappa_q+(\kappa_q\cdot c')_{ji}\lambda_q\big)$ kills relations at $(i,j)$ by Step 7, so Step 6 gives $(\lambda_q\cdot c')_{ji}=0$ for all $q$. Hence $\lambda_k$ kills relations at $(j,i)$, and Step 6 with $\alpha=0$ and $\beta$ the $k$-th unit vector gives $Y_k$.

\emph{Step 9: the even part.} Let $\psi_0(g)=\varphi(g)-\chi(g)\sum_k[\kappa_k(g),r(g)Y_kr(g)^{-1}]$. Using $\lambda_k=\delta_\chi Y_k$, a direct computation shows that $\psi_0$ is a $\chi$-cocycle with even values; using $\kappa_k(\sigma(g))=-{}^t\kappa_k(g)$, $s\,{}^tY_k=Y_k$ and (ii), that $\psi_0(\sigma(g))=s\,{}^t\psi_0(g)$; and using the trace condition and $\lambda_k=\delta_1Y_k$ on $G'$, that $\operatorname{tr}(\psi_0(g)r(g))=0$ for $g\in G'$.
On $G'$, $\psi_0$ is a $1$-cocycle that kills relations at every entry: for a relation $c$ the matrix $M=\psi_0\cdot c$ is even, and for $h\in G'$ the trace condition at $gh$ and the cocycle relation give $\operatorname{tr}(Mr(h))=-\operatorname{tr}(r(c)\psi_0(h)r(h))=0$; so $\operatorname{tr}(MX)=0$ for every even $X$, and $M=0$. By Step 1, $\psi_0=\delta_1Y'_0$ on $G'$, and also $\psi_0=\delta_1Y'_1$ for the even matrix $Y'_1=\frac12(Y'_0+SY'_0S)$, since $\psi_0$ has even values.
The difference $w=\psi_0-\delta_\chi Y'_1$ is a $\chi$-cocycle with even values that vanishes on $G'$. As in Step 4, each $w(k)$ commutes with every $r(h)$, $h\in G'$, hence with every $r(g)$ by (i); so $w(gk)=w(g)+\chi(g)w(k)$. Since $g_0gg_0^{-1}g^{-1}\in G'$, $w(g_0g)=w(gg_0)$, which gives $(\chi(g_0)-1)w(g)=(\chi(g)-1)w(g_0)$. So $w=\delta_\chi Z$ for the even matrix $Z=(\chi(g_0)-1)^{-1}w(g_0)$, and $\psi_0=\delta_\chi Y$ with $Y=Y'_1+Z$ even. Finally $Y_0=\frac12(Y+s\,{}^tY)$ is even with $s\,{}^tY_0=Y_0$, and $\psi_0=\delta_\chi Y_0$, by $\psi_0(\sigma(g))=s\,{}^t\psi_0(g)$ and (ii). This is the formula for $\varphi$.
% Step 1: Det.der with Det.derEv. Step 2: Det.relMove (move_one). Step 3: Det.eigen. Step 4: Det.infres, with (i) Det.central. Step 5: Det.shapiro with Det.cornerCoc, and (ii) Det.dSigma. Step 6: Det.ann (kills_cut, block_princ). Step 7: Det.relPair. Step 8: Det.oddPart. Step 9: Det.detAlg (the cocycle psi, hcoc, hpol, htr) and Det.evenPart (evP_kills, evP_even, evP_sah, then the symmetrization).
\end{proof}

\begin{lemma}\label{sha-patching-detection}
In Situation~\ref{sha-hecke-situation}, suppose that $\varepsilon\bmod\ell$ is nontrivial on $\Gamma_0$. Let $x$ be a normalized arc with a Selmer basis, as in Definition~\ref{sha-patching-selmer}. Then $\Gamma'$ detects the dual Selmer group of $x$.
\end{lemma}
\leanref{Development/AutomorphicArithmetic/AutomorphicVarieties/PatchingModules/DeterminantPadicFree/DetAsm.lean}{UpperTheory.V.Pa.Det.detAsm}
\begin{proof}
Let $F$ be the fraction field of $\kappa[[T]]$. The data $G=\Gamma_0$, $G'=\Gamma'$, $\chi=\varepsilon$ with values in $\mathbf F_\ell^\times\subset F$, $\sigma(g)=cgc^{-1}$ and $r=\rho$ satisfy the hypotheses of Lemma~\ref{sha-patching-twisted}: $2\ne0$ in $F$ as $\ell$ is odd; $\Gamma'=\Gamma_0\cap\ker\varepsilon$ contains the commutators of $\Gamma_0$; $\chi$ is trivial on $\Gamma'$ and nontrivial, since $\varepsilon\bmod\ell$ is nontrivial on $\Gamma_0$; $\varepsilon(cgc^{-1})=\varepsilon(g)$; ${}^t\rho(g)\rho(cgc^{-1})=\varepsilon(g)^{1-n}$; and the $F$-span of $\rho(\Gamma')$ is the space of even matrices, since for every even $Z$ over $\kappa[[T]]$ some $T^sZ$ lies in the $\kappa[[T]]$-span of $\rho(\Gamma')$. The Selmer basis $\kappa_1,\dots,\kappa_g$ satisfies (1) of that lemma: if $\sum_kc_k\kappa_k=\delta_1Y$ over $F$ with $Y$ odd and antisymmetric, multiplying by a power of $T$ makes the $c_k$ and $Y$ integral, and then the independence of the Selmer basis modulo coboundaries gives $c=0$.

Let $z$ be a dual cocycle with $\langle z,\sigma\rangle=0$ for $\sigma\in\Gamma'$. Write $z(g)=\eta\,\varphi(g)+\sum_k\epsilon'_k\lambda_k(g)$, with $\varphi(g)$ even and $\lambda_k(g)$ odd. Since $\rho_1(g)^{-1}=\rho(g)^{-1}(1-\sum_k\epsilon_k\kappa_k(g))$, for $Z=\eta Z_0+\sum_k\epsilon'_kZ_k$ in $V^*$ the matrix $\rho_1(g)Z\rho_1(g)^{-1}$ has $\epsilon'_k$-component $\rho(g)Z_k\rho(g)^{-1}$ and $\eta$-component $\rho(g)Z_0\rho(g)^{-1}+\sum_k[\kappa_k(g),\rho(g)Z_k\rho(g)^{-1}]$, and the coefficient of $\eta$ in $\operatorname{tr}(Z\rho_1(g))$ is $\operatorname{tr}(Z_0\rho(g))+\sum_k\operatorname{tr}(Z_k\kappa_k(g)\rho(g))$. So the cocycle relation and the symmetry of $z$, with $s=-\varepsilon(c)$, which satisfies $s^2=\varepsilon(c^2)=1$, and the vanishing of the pairing on $\Gamma'$ are the conditions (2) and (3) of Lemma~\ref{sha-patching-twisted}. That lemma gives $Y_0,\dots,Y_g$ over $F$. Choose $s'$ such that $T^{s'}Y_k$ has entries in $\kappa[[T]]$ for all $k$, and let $Z\in V^*$ have $\eta$-component $T^{s'}Y_0$ and $\epsilon'_k$-components $T^{s'}Y_k$. Then $-\varepsilon(c)\,{}^tZ=Z$, and by the same formulas $T^{s'}z(g)=\varepsilon(g)\rho_1(g)Z\rho_1(g)^{-1}-Z$, so $T^{s'}z$ is a dual coboundary.
% Det.detAsm: Det.detFr (the frame over FractionRing x.A), the private isKap (independence over the fraction field), Det.w1sAct (the components of the conjugation and of the trace), Det.detAlg, the private clr (clearing denominators) and Det.w1sMem (assembling Z).
\end{proof}

\begin{lemma}\label{sha-patching-frobenius-cocycles}
Let $\ell$ be a prime, $\kappa$ a finite field of characteristic $\ell$, and $G$ a compact totally disconnected topological group such that every open subgroup of $G$ has finitely many continuous homomorphisms to $\mathbf Z/\ell$. Let $R$ be a commutative $\kappa[[T]]$-algebra that is free of finite rank as a $\kappa[[T]]$-module, and call a map $f:G\to R$ \emph{continuous} if $g\mapsto f(g)\bmod T^NR$ is locally constant for every $N\ge0$. Let $\chi:G\to R$ be a continuous multiplicative map and $r:G\to M_n(R)$ a multiplicative map with continuous entries. Then every $\kappa[[T]]$-module of maps $z:G\to M_n(R)$ with continuous entries such that $z(gh)=z(g)+\chi(g)\,r(g)z(h)r(g^{-1})$ for all $g,h\in G$ is finitely generated.
\end{lemma}
\leanref{Development/AutomorphicArithmetic/AutomorphicVarieties/PatchingModules/DeterminantPadicFree/AcConst.lean}{UpperTheory.V.Pa.Det.matFg}
\leanref{Development/AutomorphicArithmetic/AutomorphicVarieties/PatchingModules/DeterminantPadicFree/MemPowSmul.lean}{UpperTheory.V.Pa.Det.z1Fg}
\leanref{Development/AutomorphicArithmetic/AutomorphicVarieties/PatchingModules/DeterminantPadicFree/LocallyConstantCocyclesFiniteOfPhi.lean}{UpperTheory.V.Pa.Det.locallyConstant_cocycles_finite_of_phi}
\begin{proof}
Let $M=M_n(R)$, a free $\kappa[[T]]$-module of finite rank, on which $g\in G$ acts by $g\cdot m=\chi(g)r(g)mr(g^{-1})$, and let $C$ be the module of all maps $z:G\to M$ with continuous entries such that $z(gh)=z(g)+g\cdot z(h)$. Since $\kappa[[T]]$ is noetherian, it suffices to show that $C$ is finitely generated.

\emph{Separation.} $\bigcap_kT^kC=0$, since the values of an element of $T^kC$ lie in $T^kM$ and $\bigcap_kT^kM=0$.

\emph{Reduction modulo $T$.} Let $\Phi$ send $z\in C$ to the map $g\mapsto z(g)\bmod TM$. Its kernel is $TC$. Indeed, if every $z(g)$ lies in $TM$, write $z(g)=Tw(g)$; then $w(g)$ is unique, since $M$ is torsion free, so $w$ satisfies the cocycle relation, and $w$ has continuous entries, since $T(w(g)-w(g'))\in T^{N+1}M$ gives $w(g)-w(g')\in T^NM$.

\emph{Finiteness modulo $T$.} The image of $\Phi$ consists of locally constant maps $\bar z:G\to\bar M=M/TM$ satisfying the cocycle relation for the action of $G$ on $\bar M$, and $g\mapsto g\cdot m$ is locally constant for each $m\in\bar M$. The group $\bar M$ is finite and killed by $\ell$. There are finitely many such $\bar z$. Indeed, the kernel $K$ of the action on $\bar M$ is the intersection of the finitely many open sets $\{g:g\cdot m=m\}$, so it is an open subgroup, of finite index since $G$ is compact. The cocycle relation gives $\bar z(1)=0$ and makes the restriction of $\bar z$ to $K$ a homomorphism. For each of the finitely many homomorphisms $f:\bar M\to\mathbf Z/\ell$, $f\circ\bar z$ is a continuous homomorphism $K\to\mathbf Z/\ell$, of which there are finitely many; these composites determine $\bar z$ on $K$, since an element of the $\mathbf F_\ell$-vector space $\bar M$ killed by every $f$ is $0$. Finally, $\bar z$ is determined by its restriction to $K$ and its values at a set of representatives $c$ of $G/K$: for $g=ch$ with $h\in K$, $\bar z(g)=\bar z(c)+c\cdot\bar z(h)$.

\emph{Nakayama.} So $C/TC$ is finite, hence finitely generated over $\kappa[[T]]$. Since $\kappa[[T]]$ is complete for the $T$-adic topology and $C$ is separated for it, $C$ is finitely generated \cite[Theorem~8.4]{Mat86}.
% Separation: hH. Reduction: hker (dividing by X). Finiteness: Det.locallyConstant_cocycles_finite_of_phi (the open kernel of the action, the maps phi to the homomorphisms of the kernel, hinj). Nakayama: UpperTheory.V.abelian₂.nakayama, then isNoetherian for the submodule. Det.matFg applies Det.z1Fg to M_n(R) with the twisted conjugation action.
\end{proof}

\begin{lemma}\label{sha-patching-frobenius-lattice}
Let $A$ be a discrete valuation ring with uniformizer $t$, $M$ an $A$-module, $S\subset M$ a finitely generated submodule and $B\subset M$ a submodule. Then there are $g\ge0$ and $\kappa_1,\dots,\kappa_g\in S$ such that $\sum_jc_j\kappa_j\in B$, with $c_j\in A$, only if all $c_j$ vanish, and such that for every $z\in S$ there are $s\ge0$ and $c_1,\dots,c_g\in A$ with $t^sz-\sum_jc_j\kappa_j\in B$.
\end{lemma}
\leanref{Development/AutomorphicArithmetic/AutomorphicVarieties/PatchingModules/DeterminantPadicFree/MaxIndep.lean}{UpperTheory.V.Pa.Det.maxIndep}
\begin{proof}
The image $\bar S$ of $S$ in $M/B$ is a finitely generated module over the principal ideal domain $A$, and every irreducible element of $A$ is associate to $t$, so $\bar S\cong A^g\oplus\bigoplus_iA/t^{e_i}A$ for some $g$ and exponents $e_i$. Let $\kappa_j\in S$ lift the image in $\bar S$ of the $j$-th basis vector of $A^g$. Their images in $M/B$ are linearly independent, which is the first property. For $z\in S$, let $K=\sum_ie_i$; then $t^K$ kills the torsion summands, so the image of $t^Kz$ lies in the image of $A^g$, which is spanned by the images of the $\kappa_j$.
\end{proof}

\begin{lemma}\label{sha-patching-frobenius-selmer}
In Situation~\ref{sha-hecke-situation}, suppose that every open subgroup of $\Gamma$ has finitely many continuous homomorphisms to $\mathbf Z/\ell$. Then every normalized arc has a Selmer basis, in the sense of Definition~\ref{sha-patching-selmer}(2).
\end{lemma}
\leanref{Development/AutomorphicArithmetic/AutomorphicVarieties/PatchingModules/DeterminantPadicFree/CharPOfPadic.lean}{UpperTheory.V.Pa.Det.kapEx}
\begin{proof}
Let $x$ be a normalized arc over $\kappa[[T]]$, with representation $\rho$. The subgroup $\Gamma_0$ is open in the profinite group $\Gamma$, so it is compact and totally disconnected, and its open subgroups are open subgroups of $\Gamma$; so every open subgroup of $\Gamma_0$ has finitely many continuous homomorphisms to $\mathbf Z/\ell$. The Selmer group and the coboundaries are $\kappa[[T]]$-submodules of the module of maps $\Gamma_0\to M_n(\kappa[[T]])$, since the conditions of Definition~\ref{sha-patching-selmer}(1) are linear in $z$ and in the matrices $Y_v$. The elements of the Selmer group have continuous entries and satisfy $z(gh)=z(g)+\rho(g)z(h)\rho(g)^{-1}$, so by Lemma~\ref{sha-patching-frobenius-cocycles}, with $G=\Gamma_0$, $R=\kappa[[T]]$, $\chi=1$ and $r=\rho$, the Selmer group is finitely generated. Lemma~\ref{sha-patching-frobenius-lattice}, applied to the Selmer group and the coboundaries, gives $\kappa_1,\dots,\kappa_g$, and matrices $Y_{j,v}$ for them exist by the definition of the Selmer group.
% Det.kapEx: Det.phiSub (PhiCond for G_0), Det.selSub (the Selmer group SelZ and the coboundaries Bd are submodules), Det.matFg, Det.maxIndep with B = Bd.
\end{proof}

\begin{lemma}\label{sha-patching-frobenius-forms}
Let $A$ be a discrete valuation ring with uniformizer $t$, $M$ an $A$-module, $Z\subset M$ a finitely generated submodule, $B\subset M$ a submodule, and $(f_i)_{i\in I}$ a family of $A$-linear maps $M\to A$ that vanish on $B$. Suppose that for every $z\in Z$ with $f_i(z)=0$ for all $i\in I$ there is $s\ge0$ with $t^sz\in B$. Then there is a finite subset $I_0\subset I$ such that for every $z\in Z$ with $f_i(z)=0$ for all $i\in I_0$ there is $s\ge0$ with $t^sz\in B$.
\end{lemma}
\leanref{Development/AutomorphicArithmetic/AutomorphicVarieties/PatchingModules/DeterminantPadicFree/FinSub.lean}{UpperTheory.V.Pa.Det.finSub}
\begin{proof}
Let $z_1,\dots,z_k$ generate $Z$, and for $i\in I$ let $h_i=(f_i(z_1),\dots,f_i(z_k))\in A^k$. Since $A^k$ is a noetherian module, among the submodules spanned by the sets $\{h_i:i\in I'\}$ for finite $I'\subset I$ there is a maximal one; let $I_0$ be such an $I'$. For $j\in I$, the span for $I_0\cup\{j\}$ contains that for $I_0$, so by maximality it is equal to it, and $h_j$ lies in the span of the $h_i$, $i\in I_0$. Let $z=\sum_ac_az_a\in Z$ with $f_i(z)=0$ for $i\in I_0$, and let $E:A^k\to A$ be the linear map $(y_a)\mapsto\sum_ac_ay_a$. Then $f_j(z)=E(h_j)$ for every $j\in I$. Since $E$ vanishes on the $h_i$, $i\in I_0$, it vanishes on their span, so $f_j(z)=0$ for all $j\in I$, and the hypothesis gives $s$ with $t^sz\in B$.
\end{proof}

\begin{lemma}\label{sha-patching-frobenius-finite}
In Situation~\ref{sha-hecke-situation}, suppose that every open subgroup of $\Gamma$ has finitely many continuous homomorphisms to $\mathbf Z/\ell$. Let $x$ be a normalized arc with a Selmer basis, and let $\Sigma\subset\Gamma_0$ be a subset that detects its dual Selmer group, in the sense of Definition~\ref{sha-patching-selmer}(4). Then some finite subset of $\Sigma$ detects it.
\end{lemma}
\leanref{Development/AutomorphicArithmetic/AutomorphicVarieties/PatchingModules/DeterminantPadicFree/DetFin.lean}{UpperTheory.V.Pa.Det.detFin}
\begin{proof}
Let $R_g$ and $\rho_1$ be as in Definition~\ref{sha-patching-selmer}(3). The ring $R_g$ is free of rank $2g+2$ over $\kappa[[T]]$. View $\varepsilon$ as a map to $R_g$ through the map $\mathbf Z_\ell\to\kappa[[T]]$. It is continuous for the $T$-adic topology: modulo $T^N$ it factors through $\varepsilon\bmod\ell$, since $\ell=0$ in $\kappa$, and $\varepsilon\bmod\ell$ is locally constant. The entries of $\rho_1$ are continuous, since those of $\rho$ and of the $\kappa_j$ are. As in the proof of Lemma~\ref{sha-patching-frobenius-selmer}, every open subgroup of $\Gamma_0$ has finitely many continuous homomorphisms to $\mathbf Z/\ell$. By Lemma~\ref{sha-patching-frobenius-cocycles}, with $G=\Gamma_0$, $R=R_g$, $\chi=\varepsilon$ and $r=\rho_1$, the module $Z$ of maps $z:\Gamma_0\to M_n(R_g)$ with continuous entries such that $z(gh)=z(g)+\varepsilon(g)\rho_1(g)z(h)\rho_1(g)^{-1}$ is finitely generated, and it contains the dual cocycles. For $\sigma\in\Sigma$ let $f_\sigma$ send a map $z:\Gamma_0\to M_n(R_g)$ to the coefficient of $\eta$ in $\operatorname{tr}(z(\sigma)\rho_1(\sigma))$, a $\kappa[[T]]$-linear map with $f_\sigma(z)=\langle z,\sigma\rangle$ for dual cocycles, and let $B$ be the intersection of the kernels of the $f_\sigma$. Lemma~\ref{sha-patching-frobenius-forms}, applied to $Z$, $B$ and the $f_\sigma$, whose hypothesis holds with $s=0$, gives a finite subset $\Sigma_0\subset\Sigma$. Let $z$ be a dual cocycle with $\langle z,\sigma\rangle=0$ for $\sigma\in\Sigma_0$. Then $T^sz\in B$ for some $s$, so $T^s\langle z,\sigma\rangle=0$, and hence $\langle z,\sigma\rangle=0$, for every $\sigma\in\Sigma$. Since $\Sigma$ detects the dual Selmer group, some $T^{s'}z$ is a dual coboundary.
% Det.detFin: Zb (continuous twisted cocycles with values in M_n(DD)), Det.matFg with chi = eps (hA, from lc_zp) and r = rho1, Det.finSub with f sigma = etaL o trace o (right multiplication by rho1(sigma)) and B the intersection of their kernels.
\end{proof}

\begin{lemma}\label{sha-patching-frobenius-extension}
In Situation~\ref{sha-hecke-situation}, let $x$ be a normalized arc over $\kappa[[T]]$, with representation $\rho$, block sizes $a$ and $n-a$, characters $\psi_{v,i}$ and a Selmer basis. Let $\sigma_1,\dots,\sigma_q\in\Gamma_0$ be such that the characteristic polynomial of each $\rho(\sigma_b)$ is separable over the fraction field of $\kappa[[T]]$ and $\{\sigma_1,\dots,\sigma_q\}$ detects the dual Selmer group of $x$. Then there is an integer $d\ge1$ such that for every integer $D\ge1$ divisible by $d$ there are a finite field $\kappa'$ with $\ell^D$ elements, an injective local homomorphism $j:\kappa[[T]]\to\kappa'[[T]]$ for which $\kappa'[[T]]$ is a finite $\kappa[[T]]$-module, and a normalized arc $x'$ over $\kappa'[[T]]$ with representation $j\circ\rho$, block sizes $a$ and $n-a$ and characters $j\circ\psi_{v,i}$, such that $x'$ has a Selmer basis for which $\{\sigma_1,\dots,\sigma_q\}$ detects its dual Selmer group, and the characteristic polynomial of $j(\rho(\sigma_b))$ is $\prod_{m=1}^n(X-\alpha_{b,m})$ with pairwise distinct $\alpha_{b,1},\dots,\alpha_{b,n}\in\kappa'[[T]]$, for each $b$.
\end{lemma}
\leanref{Development/AutomorphicArithmetic/AutomorphicVarieties/PatchingModules/DeterminantPadicFree/CharPOfPadicAlt.lean}{UpperTheory.V.Pa.Det.onePt}
\leanref{Development/AutomorphicArithmetic/AutomorphicVarieties/PatchingModules/DeterminantPadicFree/KapPush.lean}{UpperTheory.V.Pa.Det.kapPush}
\leanref{Development/AutomorphicArithmetic/AutomorphicVarieties/PatchingModules/DeterminantPadicFree/DetPush.lean}{UpperTheory.V.Pa.Det.detPush}
\begin{proof}
\emph{Splitting.} Let $P_b$ be the characteristic polynomial of $\rho(\sigma_b)$, and let $E$ be a splitting field of $\prod_bP_b$ over the fraction field $F$ of $\kappa[[T]]$. Every root of $\prod_bP_b$ is a root of a separable $P_b$, so $E$ is separable over $F$. Hence the integral closure $A_1$ of $\kappa[[T]]$ in $E$ is a complete discrete valuation ring, finite and free over $\kappa[[T]]$, with finite residue field, and the map $\kappa[[T]]\to A_1$ is local \cite[Chapter~II, \S2]{Ser79}. Each $P_b$ is monic with coefficients in $\kappa[[T]]$, so its $n$ roots in $E$ lie in $A_1$; they are pairwise distinct since $P_b$ is separable, and $P_b$ is the product of the corresponding linear factors over $A_1$.

\emph{Power series.} The ring $A_1$ has characteristic $\ell$ and a finite residue field $\kappa''$, so it is isomorphic to $\kappa''[[T]]$, with $T$ going to a uniformizer \cite[Chapter~II, \S4]{Ser79}. This gives an injective local homomorphism $j_1:\kappa[[T]]\to\kappa''[[T]]$, for which $\kappa''[[T]]$ is finite over $\kappa[[T]]$. Let $\ell^d$ be the number of elements of $\kappa''$.

\emph{The residue field.} Let $D\ge1$ be divisible by $d$, and let $\kappa'$ be the field with $\ell^D$ elements; it contains a subfield isomorphic to $\kappa''$. Let $j_2:\kappa''[[T]]\to\kappa'[[T]]$ apply such an embedding to the coefficients; it is injective and local, and $\kappa'[[T]]$ is finite over $\kappa''[[T]]$. Put $j=j_2\circ j_1$.

\emph{The arc.} Let $x'$ be obtained by applying $j$ to $\rho$, the $\psi_{v,i}$ and the matrices $w_v$, with the same block sizes. It is a normalized arc. The conditions of Lemma~\ref{sha-patching-normalization} are identities between matrices and characters and the inequalities (N1), which an injective ring homomorphism preserves, using for the symmetric form that $j$ carries the map $\mathbf Z_\ell\to\kappa[[T]]$ to the map $\mathbf Z_\ell\to\kappa'[[T]]$, both being the composite of reduction modulo $\ell$ with the inclusion of $\mathbf F_\ell$. The continuity and the condition on the span of $\rho(\Gamma')$ are preserved since $j(T)\in T\kappa'[[T]]$. The characteristic polynomial of $j(\rho(\sigma_b))$ is the image of $P_b$, which is $\prod_m(X-\alpha_{b,m})$ for the images $\alpha_{b,m}$ in $\kappa'[[T]]$ of its roots in $A_1$, pairwise distinct since $j_2$ and the isomorphism $A_1\cong\kappa''[[T]]$ are injective.

\emph{Coordinates.} The ring $\kappa'[[T]]$ is a finitely generated torsion free, hence free, module over $\kappa[[T]]$ through $j$. Let $b_1,\dots,b_r$ be a basis and $p_1,\dots,p_r:\kappa'[[T]]\to\kappa[[T]]$ the coordinates, so that $y=\sum_ij(p_i(y))b_i$ and $p_i(j(c)y)=c\,p_i(y)$. Write $j(T)=uT^e$ with $u$ a unit and $e\ge1$. Then $p_i(T^{eN}\kappa'[[T]])\subset T^N\kappa[[T]]$, since $T^{eN}y=j(T^N)u^{-N}y$. Apply $j$ and the $p_i$ to matrices and maps entrywise. For matrices $X,X'$ over $\kappa[[T]]$ and $Y$ over $\kappa'[[T]]$, $p_i(j(X)\,Y\,j(X'))=X\,p_i(Y)\,X'$. Two facts follow from the definition of the Selmer group. (i) If $z$ lies in the Selmer group of $x$, with matrices $Y_v$, then $j\circ z$ lies in that of $x'$, with matrices $j(Y_v)$, and $j$ carries coboundaries of $x$ to coboundaries of $x'$. (ii) If $z$ lies in the Selmer group of $x'$, with matrices $Y_v$, then $p_i\circ z$ lies in that of $x$, with matrices $p_i(Y_v)$, and $p_i$ carries coboundaries of $x'$ to coboundaries of $x$: the conditions involve $z$ through products as above, and continuity follows from $p_i(T^{eN}\kappa'[[T]])\subset T^N\kappa[[T]]$.

\emph{The Selmer basis.} Let $\kappa_1,\dots,\kappa_g$ be the Selmer basis of $x$, and put $\kappa'_k=j\circ\kappa_k$, with matrices $j(Y_{k,v})$; these lie in the Selmer group of $x'$ by (i). If $\sum_kc_k\kappa'_k$ is a coboundary, with $c_k\in\kappa'[[T]]$, then by (ii) so is $p_i\circ\sum_kc_k\kappa'_k=\sum_kp_i(c_k)\kappa_k$, so $p_i(c_k)=0$ for all $i$ and $k$, and $c_k=\sum_ij(p_i(c_k))b_i=0$. Let $z$ lie in the Selmer group of $x'$. By (ii) and the Selmer basis of $x$, there are $s_i\ge0$ and $c_{i,k}\in\kappa[[T]]$ with $T^{s_i}p_i\circ z-\sum_kc_{i,k}\kappa_k$ a coboundary of $x$. Let $N=\sum_is_i$ and $d_i=u^{-N}b_i\,j(T^{N-s_i})$. Then $\sum_id_i\,j(T^{s_i}p_i(y))=T^{eN}y$ for every $y\in\kappa'[[T]]$, since $j(T^{N-s_i})j(T^{s_i})=j(T)^N=u^NT^{eN}$. Hence
\[T^{eN}z-\sum_k\Big(\sum_id_i\,j(c_{i,k})\Big)\kappa'_k=\sum_id_i\,j\circ\Big(T^{s_i}p_i\circ z-\sum_kc_{i,k}\kappa_k\Big),\]
which is a coboundary of $x'$ by (i). So the $\kappa'_k$ form a Selmer basis of $x'$.

\emph{Detection.} Apply $j$ and the $p_i$ to the coefficients of elements of $R_g$ in the basis $1,\epsilon_k,\epsilon'_k,\eta$. Then $j$ carries the first-order deformation $\rho_1$ of $x$ to that of $x'$, and $\varepsilon$ to $\varepsilon$. Let $z$ be a dual cocycle of $x'$ with $\langle z,\sigma_b\rangle=0$ for all $b$. As in (ii), each $p_i\circ z$ is a dual cocycle of $x$, and $\langle p_i\circ z,\sigma_b\rangle=p_i(\langle z,\sigma_b\rangle)=0$, since the trace and the coefficient of $\eta$ commute with $p_i$. By the detection for $x$ there are $s_i\ge0$ and $Z_i\in V^*$ with $-\varepsilon(c)\,{}^tZ_i=Z_i$ such that $T^{s_i}p_i\circ z$ is the dual coboundary of $Z_i$. With $N$ and $d_i$ as above, $T^{eN}z=\sum_id_i\,j\circ(T^{s_i}p_i\circ z)$, which is the dual coboundary of $\sum_id_i\,j(Z_i)$.
% Det.onePt: Det.splitEx (splitting field, AlgebraTools.integralClosure_complete_dvr, exists_injective_roots_in_integral_closure), Det.extEx (AlgebraTools.powerSeries_equiv_of_complete_dvr and the private pushNP, with padic_hom_unique for eps), Det.resExt (j2 from GaloisField l d into GaloisField l D), Det.extTrans, Det.kapPush (Det.extFree for the basis b and coordinates p, Det.pushSel (i), Det.compSel (ii), with mapR and mapL), Det.detPush (exJP for J and P on DD).
\end{proof}

\begin{lemma}\label{sha-patching-frobenius}
In Situation~\ref{sha-hecke-situation}, suppose that every open subgroup of $\Gamma$ has finitely many continuous homomorphisms to $\mathbf Z/\ell$ and that $\varepsilon\bmod\ell$ is nontrivial on $\Gamma_0$. Let $x_0,\dots,x_h$ be normalized arcs satisfying (N3). Then there are $q\ge0$ and $\sigma_1,\dots,\sigma_q\in\Gamma'$, and for each $i$ a normalized arc $x'_i$ obtained from $x_i$ by a finite extension of $\kappa_i[[T]]$, such that the residue fields of the $x'_i$ have the same number of elements and:
\begin{enumerate}
\item the characteristic polynomial of $\rho'_i(\sigma_j)$ is $\prod_{m=1}^n(X-\alpha_{i,j,m})$ with distinct $\alpha_{i,j,m}\in\kappa'_i[[T]]$;
\item $x'_i$ has a Selmer basis $\kappa_1,\dots,\kappa_g$, in the sense of Definition~\ref{sha-patching-selmer}(2); its classes have values in the matrices whose diagonal blocks vanish;
\item the set $\{\sigma_1,\dots,\sigma_q\}$ detects the dual Selmer group of $x'_i$ for this basis, in the sense of Definition~\ref{sha-patching-selmer}(4), with the first-order deformation $\rho_{i,1}=(1+\sum_m\epsilon_m\kappa_m)\rho'_i$.
\end{enumerate}
\end{lemma}
\leanref{Development/AutomorphicArithmetic/AutomorphicVarieties/PatchingModules/PatchedSupportCompletions.lean}{UpperTheory.V.Pa.Choice}
\leanref{Development/AutomorphicArithmetic/AutomorphicVarieties/PatchingModules/AutomorphicDetection/Choice.lean}{UpperTheory.V.Pa.choice}
\leanref{Development/AutomorphicArithmetic/AutomorphicVarieties/PatchingModules/AutomorphicDetection/Detection.lean}{UpperTheory.V.Pa.detection}
\begin{proof}
Let $\Sigma$ be the set of $\sigma\in\Gamma'$ at which the characteristic polynomials of all the $\rho_i(\sigma)$ are separable over the fraction fields; by (N3) it is dense in $\Gamma'$. By Lemma~\ref{sha-patching-frobenius-selmer}, each $x_i$ has a Selmer basis. By Lemma~\ref{sha-patching-detection}, $\Gamma'$ detects the dual Selmer group of each $x_i$. If a dual cocycle $z$ has $\langle z,\sigma\rangle=0$ for all $\sigma\in\Sigma$, then the set of $\sigma$ with $\langle z,\sigma\rangle=0$, which is closed since $z$ and $\rho_{i,1}$ are continuous, contains the closure of $\Sigma$ and so $\Gamma'$. Hence $\Sigma$ detects the dual Selmer group of each $x_i$. By Lemma~\ref{sha-patching-frobenius-finite}, a finite subset $\Sigma_i\subset\Sigma$ detects the dual Selmer group of $x_i$. Let $\sigma_1,\dots,\sigma_q$ enumerate the union of the $\Sigma_i$; it detects the dual Selmer group of each $x_i$, since a larger set imposes more conditions. Lemma~\ref{sha-patching-frobenius-extension}, applied to each $x_i$ and $\sigma_1,\dots,\sigma_q$, gives integers $d_i\ge1$. With $D=\prod_id_i$, which every $d_i$ divides, it gives normalized arcs $x'_i$ obtained from the $x_i$ by finite extensions, whose residue fields all have $\ell^D$ elements, with Selmer bases for which $\{\sigma_1,\dots,\sigma_q\}$ detects the dual Selmer group, and with characteristic polynomials of the $\rho'_i(\sigma_j)$ as in (1).
% Pa.choice: Det.omDense (density), Det.kapEx (Lemma sha-patching-frobenius-selmer), Pa.detection (with Det.ellClosed), Det.detFin (Lemma sha-patching-frobenius-finite), detects_mono, Det.onePt (Lemma sha-patching-frobenius-extension) with D the product of the d_i.
\end{proof}

\begin{lemma}\label{sha-patching-places}
In Situation~\ref{sha-hecke-situation}, every open subgroup of $\Gamma$ has finitely many continuous homomorphisms to $\mathbf Z/\ell$. Let $\sigma_1,\dots,\sigma_q\in\Gamma'$, and let $H$ be the intersection of $\Gamma_0$, the kernel of $\varepsilon\bmod\ell$ and the kernels of the $\bar\chi_i$, an open normal subgroup of $\Gamma$. Then for each $N\ge1$ there are places $v_{N,1},\dots,v_{N,q}$ of $L^+$, pairwise distinct, outside $S$ and split in $L$, places $w_{N,j}$ of $L$ above them with $\operatorname{N}w_{N,j}\equiv1\pmod{\ell^N}$, and, in the Galois group $\Gamma_N$ of the level augmented by the $v_{N,j}$, inverses $\phi_{N,j}$ of Frobenius elements and topological generators $\tau_{N,j}$ of the $\ell$-part of tame inertia at places above the $w_{N,j}$, such that, writing $H_N$ for the inverse image of $H$ in $\Gamma_N$ and $P(H_N)$ for the kernel of its maximal pro-$\ell$ quotient:
\begin{enumerate}
\item the image of $\phi_{N,j}$ in $\Gamma$ lies in $\sigma_jH$, and for every open normal subgroup $V$ of $\Gamma$ containing $P(H)$ it lies in $\sigma_jV$ for all large $N$;
\item $\tau_{N,j}$ maps to $1$ in $\Gamma$, and $\tau_{N,j}^{-\operatorname{N}w_{N,j}}\phi_{N,j}^{-1}\tau_{N,j}\phi_{N,j}\in P(H_N)$;
\item the kernel of $\Gamma_N\to\Gamma$ lies in the closed subgroup generated by the conjugates of the $\tau_{N,j}$ and by $P(H_N)$;
\item $H_N$ has at most $h\,\ell^{q[\Gamma:H]}$ continuous homomorphisms to $\mathbf Z/\ell$, where $h$ is the finite number of continuous homomorphisms $H\to\mathbf Z/\ell$;
\item every open subgroup of the groups $D_v$, $v\mid\ell$, and of the local Galois groups at the places $\tilde u$, $u\in S_{\mathrm a}$, has finitely many continuous homomorphisms to $\mathbf Z/\ell$.
\end{enumerate}
\end{lemma}
\leanref{Development/AutomorphicArithmetic/AutomorphicVarieties/PatchingModules/TaylorWilesObstructionBound.lean}{UpperTheory.V.Pa.ArPl}
\leanref{Development/AutomorphicArithmetic/AutomorphicVarieties/PatchingModules/AutomorphicDetection/ArPl.lean}{UpperTheory.V.Pa.arPl}
\leanref{Development/AutomorphicArithmetic/AutomorphicVarieties/PatchingModules/ThetaCoefficientMapEq.lean}{UpperTheory.V.Pa.TWGrp}
\leanref{Development/AutomorphicArithmetic/AutomorphicVarieties/PatchingModules/TaylorWilesObstructionBound.lean}{UpperTheory.V.Pa.TWPl}
\begin{proof}
\emph{Finiteness.} Let $U$ be an open subgroup of $\Gamma$. A continuous homomorphism $U\to\mathbf Z/\ell$ gives a continuous cocycle of $\Gamma$ with values in the finite module of functions $\Gamma/U\to\mathbf Z/\ell$, and distinct homomorphisms give distinct classes (Shapiro's lemma, written out with cocycles). That cohomology group is finite, by the finiteness of the cohomology with finite coefficients of the Galois group of the maximal extension of a number field unramified outside a finite set of places \cite[Chapter~VIII, \S3]{NSW08}. The same argument with the local finiteness theorem \cite[Chapter~VII, \S2]{NSW08} gives (5).

\emph{Chebotarev.} Iterating the passage from an open normal subgroup to the intersection of the kernels of its continuous homomorphisms to $\mathbf Z/\ell$, which is open by the finiteness, gives open normal subgroups $W_N\subset H$ of $\Gamma$ such that every open normal subgroup containing $P(H)$ contains $W_N$ for all large $N$. Lift each $\sigma_j$ to $g_j\in\operatorname{Gal}(\bar L/L)$, and let $U_N$ be the open subgroup of $\operatorname{Gal}(\bar L/L)$ of the elements that map into $W_N$ and on which $\varepsilon\equiv1\pmod{\ell^N}$. By the Chebotarev density theorem, there are places $v_{N,j}$ of $L^+$, pairwise distinct, outside $S$ and split in $L$, places $w_{N,j}$ of $L$ above them, and Frobenius elements $\sigma_{N,j}$ at places above the $w_{N,j}$ with $\sigma_{N,j}g_j\in U_N$. Let $\phi_{N,j}$ be the image of $\sigma_{N,j}^{-1}$. Its image in $\Gamma$ lies in $\sigma_jW_N\subset\sigma_jH$, since $W_N$ is normal, which gives (1). Since $\sigma_j\in\Gamma'$, $\varepsilon(g_j)=1$, so $\varepsilon(\sigma_{N,j})\equiv1\pmod{\ell^N}$, that is, $\operatorname{N}w_{N,j}\equiv1\pmod{\ell^N}$.

\emph{Tame inertia.} Let $\tau_{N,j}$ be a topological generator of the $\ell$-part of the tame inertia at the place of $\sigma_{N,j}$. The inertia groups at places above $v_{N,j}\notin S$ map to $1$ in $\Gamma$, hence into $H$; so the kernel of the map from such an inertia group to its tame $\ell$-quotient, which has no nontrivial pro-$\ell$ quotient, maps into $P(H_N)$. In the decomposition group, $\sigma_{N,j}\tau_{N,j}\sigma_{N,j}^{-1}$ and $\tau_{N,j}^{\operatorname{N}w_{N,j}}$ agree modulo that kernel, by the structure of the tame quotient of a local Galois group at a place prime to $\ell$; this gives (2). The kernel of $\Gamma_N\to\Gamma$ is the closed normal subgroup generated by the inertia groups at the places above the $v_{N,j}$; each is conjugate to the inertia group at the place of $\sigma_{N,j}$, which is topologically generated by $\tau_{N,j}$ and the kernel of the tame $\ell$-quotient; this gives (3).

\emph{The bound (4).} Choose representatives $c$ of $\Gamma_N/H_N$, which has $[\Gamma:H]$ elements, and for a continuous $\chi:H_N\to\mathbf Z/\ell$ record the $q[\Gamma:H]$ values $\chi(c\tau_{N,j}c^{-1})$; the elements $c\tau_{N,j}c^{-1}$ lie in $H_N$ by (2). Fix, for each possible tuple of values, one $\chi_0$ that takes them. If $\chi$ and $\chi_0$ take the same values, then $\chi-\chi_0$ vanishes on every conjugate of every $\tau_{N,j}$ (conjugation by $H_N$ does not change its values) and on $P(H_N)$, so by (3) on the kernel of $H_N\to H$, and $\chi-\chi_0$ is the composite of a continuous homomorphism $H\to\mathbf Z/\ell$ with $H_N\to H$. So $\chi$ is determined by the tuple of values and a homomorphism $H\to\mathbf Z/\ell$, which gives (4).
% Pa.arPl: phiGS (GlobalFiniteness), frat (W_N), chebFam, the tame quotient hC.tame, levUp, kerLe (h1, h2 with prokerPL, and the kernel through conj_ine), homBound (twisted_kernel_forms_nonvanishing, closure_normal_frattini_le_kernel); (5) is phiDec from the local finiteness hypothesis.
\end{proof}

\begin{definition}\label{sha-patching-tangent-modules}
In Situation~\ref{sha-hecke-situation}, let $x$ be a normalized arc over $\kappa[[T]]$, with representation $\rho$, blocks of sizes $a$ and $n-a$, matrices $w_v$ and a Selmer basis $\kappa_1,\dots,\kappa_g$ with matrices $Y_{k,v}$, and let $R_g$, $\rho_1$ and $V^*$ be as in Definition~\ref{sha-patching-selmer}. Let $\sigma_1,\dots,\sigma_q$ and the places $w_{N,j}$ be as in Lemma~\ref{sha-patching-places}, let $\Gamma_N$ be the Galois group of the level augmented by the $w_{N,j}$ and $\Gamma_{N,0}$ its subgroup fixing $L$, and view $\rho_1$ and $\varepsilon$ as functions on $\Gamma_{N,0}$ through the map $\Gamma_{N,0}\to\Gamma_0$. For $m\ge0$ let $A_m=\kappa[[T]]/T^m$, reduce $R_g$, $\rho_1$, the $w_v$ and the $Y_{k,v}$ modulo $T^m$, and put $\nu=\binom a2+\binom{n-a}2+g\,a(n-a)$.
\begin{enumerate}
\item $W_m$ is the $A_m$-module of $n\times n$ matrices over $R_g/T^m$ whose entries lie in $A_m$ in the two diagonal blocks and in $\sum_kA_m\epsilon_k$ in the two off-diagonal blocks, and $W^*_m$ is the module $V^*$ reduced modulo $T^m$. The group $\Gamma_N$ acts on $W_m$ by $g\cdot Z=\rho_1(g)Z\rho_1(g)^{-1}$ for $g\in\Gamma_{N,0}$ and $c\cdot Z=-{}^tZ$, and on $W^*_m$ by $g\cdot Z=\varepsilon(g)\rho_1(g)Z\rho_1(g)^{-1}$ for $g\in\Gamma_{N,0}$ and $c\cdot Z=-\varepsilon(c)\,{}^tZ$.
\item For $v\mid\ell$ put $w'_v=(1-\sum_k\epsilon_kY_{k,v})w_v$. An \emph{ordinary cocycle} at $v$ is a locally constant map $\beta:D_v\to W_m$ with $\beta(gh)=\beta(g)+\rho_1(g)\beta(h)\rho_1(g)^{-1}$ such that ${w'_v}^{-1}\beta(g)w'_v$ is upper triangular for $g\in D_v$, with zero diagonal for $g\in\mu_v$.
\item The \emph{Selmer group of $\rho_1$ modulo $T^m$ at level $N$} is the set of locally constant maps $z:\Gamma_{N,0}\to W_m$ with $z(gh)=z(g)+\rho_1(g)z(h)\rho_1(g)^{-1}$ and $z(cgc^{-1})=-{}^tz(g)$ for all $g,h$, vanishing on the images of the $I'_u$, such that for each $v\mid\ell$ there is $Y_v\in W_m$ for which $g\mapsto z(g)-\rho_1(g)Y_v\rho_1(g)^{-1}+Y_v$ is an ordinary cocycle at $v$. Its \emph{coboundaries} are the maps $g\mapsto\rho_1(g)Y\rho_1(g)^{-1}-Y$ with $Y\in W_m$ and ${}^tY=-Y$.
\end{enumerate}
\end{definition}
\leanref{Development/AutomorphicArithmetic/AutomorphicVarieties/PatchingModules/PairedDualFirstOrderDualMatrixSpace.lean}{UpperTheory.V.Pa.W1}
\leanref{Development/AutomorphicArithmetic/AutomorphicVarieties/PatchingModules/PairedDualFirstOrderDualMatrixSpace.lean}{UpperTheory.V.Pa.Sel1}
\leanref{Development/AutomorphicArithmetic/AutomorphicVarieties/PatchingModules/PairedDualFirstOrderDualMatrixSpace.lean}{UpperTheory.V.Pa.Bd1}
\leanref{Development/AutomorphicArithmetic/AutomorphicVarieties/PatchingModules/EpsilonEtaMatrices/EpsilonEtaTensorDual.lean}{UpperTheory.V.Pa.Kap.Zb}
\leanref{Development/AutomorphicArithmetic/AutomorphicVarieties/PatchingModules/DegreeCountCohomology/IsDualPair.lean}{UpperTheory.V.Pa.IsAd}
\leanref{Development/AutomorphicArithmetic/AutomorphicVarieties/PatchingModules/DegreeCountCohomology/IsDualPair.lean}{UpperTheory.V.Pa.nu}

\begin{lemma}\label{sha-patching-tangent-counts}
In Situation~\ref{sha-hecke-situation}, let $\Gamma_N$ be the Galois group of a level augmented by places $w_{N,1},\dots,w_{N,q}$ as in Lemma~\ref{sha-patching-places}, and let $P_N$ be the set of places consisting of the $v\mid\ell$, the $u\in S_{\mathrm a}$, $v_0$ and the $w_{N,j}$. For $w\in P_N$ let $D_w\to\Gamma_N$ be the decomposition group at a place above $w$, with $J_u=I'_u$ for $u\in S_{\mathrm a}$, let $d=[L^+:\mathbf Q]$, and let $c_1,\dots,c_d\in\Gamma_N$ be complex conjugations at the infinite places of $L^+$, so that $\eta(c_i)=\varepsilon(c_i)=-1$. Let $M$ and $M'$ be finite modules over a ring, with actions of $\Gamma_N$ for which $g\mapsto g\cdot y$ is locally constant, and let $b:M\times M'\to\mathbf Z/\ell$ be biadditive with trivial left and right kernels and $b(g\cdot y,g\cdot y')=\bar\varepsilon(g)b(y,y')$, where $\bar\varepsilon=\varepsilon\bmod\ell$. Write $Z^1$ and $B^1$ for locally constant cocycles and coboundaries.
\begin{enumerate}
\item (Greenberg-Wiles, Lemma~\ref{sha-descent-link-greenberg-wiles}, \cite[Proposition~1.6]{Wil95}, \cite[Chapter~VIII, \S7]{NSW08}, written with cocycles and with the factor $|H^0(\Gamma_N,M')|\ge1$ omitted.) Let $B^1(D_w,M)\subset\mathcal L_w\subset Z^1(D_w,M)$ be subgroups for $w\in P_N$, let $Z^1_{\mathcal L}$ be the set of $z\in Z^1(\Gamma_N,M)$ with $z|_{D_w}\in\mathcal L_w$ for all $w$, and let $Z^1_{\mathcal L^\perp}$ be the set of $y\in Z^1(\Gamma_N,M')$ with $y|_{D_w}\in B^1(D_w,M')$ for every $w$ with $\mathcal L_w=Z^1(D_w,M)$. Then $Z^1(\Gamma_N,M)$ and $Z^1(\Gamma_N,M')$ are finite, and
\[\#Z^1_{\mathcal L}\cdot\#B^1(\Gamma_N,M')\cdot|M|^{|P_N|}\cdot\prod_{i=1}^d|M^{c_i}|\le|M|\cdot\#Z^1_{\mathcal L^\perp}\cdot\prod_{w\in P_N}|\mathcal L_w|.\]
\item (Tate's local Euler characteristic formula and local duality, Lemma~\ref{sha-descent-link-local-tate} with $r=1$, \cite[Chapter~VII, \S\S2--3]{NSW08}.) For $w\in P_N$, and more generally for $M$, $M'$ and $b$ as above with actions of $D_w$ only and $b(g\cdot y,g\cdot y')=\bar\varepsilon(g)b(y,y')$ for $g\in D_w$, $Z^1(D_w,M)$ is finite and $|Z^1(D_w,M)|=|M|^{1+d_w}\,|H^0(D_w,M')|$, with $d_w=[L^+_w:\mathbf Q_\ell]$ for $w\mid\ell$ and $d_w=0$ otherwise.
\item (Inflation-restriction, since $I_u/J_u$ has order prime to $\ell$, and the cohomology of the procyclic group $D_u/I_u$ \cite[Chapter~I, \S6]{NSW08}.) If $M$ is killed by $\ell$ and $J_u$ acts trivially on $M$, then the cocycles of $D_u$ that vanish on $J_u$ form a finite set with $|M|$ elements.
\item (Local class field theory \cite[Chapter~XIV]{Ser79}, with $d'_v=[L^+_v:\mathbf Q_\ell]$.) There are integers $d'_v$, for $v\mid\ell$, with $\sum_vd'_v=d$, and continuous surjections $D_v\to\mathbf Z_\ell^{d'_v}\times\mathbf Z_\ell$ carrying $I_v$ onto $\mathbf Z_\ell^{d'_v}\times0$, through which every continuous homomorphism from $D_v$ to a finite abelian $\ell$-group that is trivial on $\mu_v$ factors.
\end{enumerate}
\end{lemma}
\leanref{Development/AutomorphicArithmetic/AutomorphicVarieties/PatchingModules/DegreeCountCohomology/IsDualPair.lean}{UpperTheory.V.Pa.GWBare}
\leanref{Development/AutomorphicArithmetic/AutomorphicVarieties/PatchingModules/DegreeCountCohomology/IsDualPair.lean}{UpperTheory.V.Pa.LocCount}
\leanref{Development/AutomorphicArithmetic/AutomorphicVarieties/PatchingModules/DegreeCountCohomology/IsDualPair.lean}{UpperTheory.V.Pa.UnrCount}
\leanref{Development/AutomorphicArithmetic/AutomorphicVarieties/PatchingModules/DegreeCountCohomology/ArSys.lean}{UpperTheory.V.Pa.arSys}
\leanref{Development/AutomorphicArithmetic/AutomorphicVarieties/PatchingModules/TaylorWilesObstructionBound.lean}{UpperTheory.V.Pa.LocCFT}

\begin{lemma}\label{sha-patching-tangent-involution}
Let $\Gamma$ be a group, $c\in\Gamma$ with $c^2=1$, and $M$ a finite module over a ring in which $2$ is a unit, with an action of $\Gamma$. Let $\mathcal Z$ be a finite subgroup of the cocycles $\Gamma\to M$ containing the coboundaries $B$. Then $\#\mathcal Z\cdot|M^c|=\#\{z\in\mathcal Z:z(c)=0\}\cdot|M|$, $\#B\cdot|M^\Gamma|=|M|$, and $\#\{z\in B:z(c)=0\}\cdot|M^\Gamma|=|M^c|$.
\end{lemma}
\leanref{Development/AutomorphicArithmetic/AutomorphicVarieties/PatchingModules/DegreeCountCohomology/CCount.lean}{UpperTheory.V.Pa.cCount}

\begin{lemma}\label{sha-patching-tangent-symmetry}
In the setting of Definition~\ref{sha-patching-tangent-modules}, let $M$ be $W_m$ or $W^*_m$, with $\chi=1$ or $\chi=\varepsilon$ respectively. A locally constant map $z:\Gamma_{N,0}\to M$ with $z(gh)=z(g)+\chi(g)\rho_1(g)z(h)\rho_1(g)^{-1}$ and $z(cgc^{-1})=-\chi(c)\,{}^tz(g)$ extends to a cocycle $y$ of $\Gamma_N$ with $y(c)=0$. Conversely a cocycle $y$ of $\Gamma_N$ with $y(c)=0$ satisfies $y(cgc^{-1})=-\chi(c)\,{}^ty(g)$ for $g\in\Gamma_{N,0}$, and it vanishes if its restriction to $\Gamma_{N,0}$ does.
\end{lemma}
\leanref{Development/AutomorphicArithmetic/AutomorphicVarieties/PatchingModules/DegreeCountCohomology/PolZCases.lean}{UpperTheory.V.Pa.polZ}

\begin{lemma}\label{sha-patching-tangent-selmer}
In the setting of Definition~\ref{sha-patching-tangent-modules}, let $\mathcal L_v$ be the group generated by $B^1(D_v,W_m)$ and the ordinary cocycles at $v$, for $v\mid\ell$; $\mathcal L_u$ the cocycles of $D_u$ vanishing on $I'_u$, for $u\in S_{\mathrm a}$; and $\mathcal L_w=Z^1(D_w,W_m)$ for $w=v_0$ and $w=w_{N,j}$. Then the Selmer group of $\rho_1$ modulo $T^m$ at level $N$ is finite, and, with $Z^1_{\mathcal L}$ as in Lemma~\ref{sha-patching-tangent-counts}(1) for $M=W_m$,
\[\#\mathrm{Sel}\cdot|W_m|\le\#Z^1_{\mathcal L}\cdot|W_m^c|,\qquad\#\mathrm{Cob}\cdot|W_m^{\Gamma_N}|=|W_m^c|,\]
where $\mathrm{Sel}$ is the Selmer group and $\mathrm{Cob}$ its group of coboundaries.
\end{lemma}
\leanref{Development/AutomorphicArithmetic/AutomorphicVarieties/PatchingModules/EpsilonEtaMatrices/BdCard.lean}{UpperTheory.V.Pa.selN}
\begin{proof}
The group $\mathcal Z=Z^1_{\mathcal L}$ is finite by Lemma~\ref{sha-patching-tangent-counts}(1), and it contains the coboundaries of $\Gamma_N$: a coboundary $g\mapsto g\cdot Y-Y$ restricts into $\mathcal L_v$, and it vanishes on $I'_u$ since $\rho_1$ is trivial there. Since $\ell$ is odd, Lemma~\ref{sha-patching-tangent-involution} gives $\#\mathcal Z\cdot|W_m^c|=\#\{y\in\mathcal Z:y(c)=0\}\cdot|W_m|$. Let $z$ be in the Selmer group. By Lemma~\ref{sha-patching-tangent-symmetry} it extends to a cocycle $y$ of $\Gamma_N$ with $y(c)=0$. This $y$ lies in $\mathcal Z$: at $v\mid\ell$, $y|_{D_v}$ is the sum of the coboundary of $Y_v$ and the ordinary cocycle $g\mapsto z(g)-\rho_1(g)Y_v\rho_1(g)^{-1}+Y_v$; at $u$ it vanishes on $I'_u$; and the other conditions are empty. Since $y$ is determined by $z$, the Selmer group injects into $\{y\in\mathcal Z:y(c)=0\}$, which gives the first inequality and the finiteness. A coboundary $g\mapsto\rho_1(g)Y\rho_1(g)^{-1}-Y$ with ${}^tY=-Y$ is the restriction of the coboundary of $Y$ on $\Gamma_N$, which vanishes at $c$ since $c\cdot Y=Y$; conversely, by Lemma~\ref{sha-patching-tangent-symmetry}, the coboundaries $y$ of $\Gamma_N$ with $y(c)=0$ restrict injectively to coboundaries of the Selmer group. So $\#\mathrm{Cob}=\#\{y\in B^1(\Gamma_N,W_m):y(c)=0\}$, and Lemma~\ref{sha-patching-tangent-involution} gives the identity.
% selN: cCount (c1 for Z, c3 with bd_card for Bd1), polZ (hP) for the extension y with y(c) = 0, the local conditions from s6 and s5.
\end{proof}

\begin{lemma}\label{sha-patching-tangent-fixed}
Let $A$ be a discrete valuation ring with maximal ideal $\mathfrak m$ and finite residue field $\kappa$, let $E,E'$ be finitely generated $A$-modules and $(f_s)_{s\in S}$ a family of $A$-linear maps $E\to E'$, and let $r$ be the rank of $K=\bigcap_s\ker f_s$. Then there is $e\ge0$ such that for every $j\ge0$ the image in $E/\mathfrak m^jE$ of $\bigcap_sf_s^{-1}(\mathfrak m^jE')$ has length at most $rj+e$, and so at most $|\kappa|^{rj+e}$ elements.
\end{lemma}
\leanref{Development/AutomorphicArithmetic/AutomorphicVarieties/AlgebraTools/LemmaB.lean}{UpperTheory.V.AlgebraTools.lemma_B}
\leanref{Development/AutomorphicArithmetic/AutomorphicVarieties/PatchingModules/EpsilonEtaMatrices/FixBdQZero.lean}{UpperTheory.V.Pa.fixBd}
\begin{proof}
Since $\operatorname{Hom}_A(E,E')$ is a noetherian $A$-module, among the spans of the $f_s$ for $s$ in finite subsets of $S$ there is a maximal one, for a finite $S_0\subset S$, and it contains every $f_s$; so $K=\bigcap_{s\in S_0}\ker f_s$. Let $F=(f_s)_{s\in S_0}:E\to E'^{S_0}$, with kernel $K$. By the Artin-Rees lemma \cite[Theorem~8.5]{Mat86}, applied to the image of $F$, there is $c$ such that $F^{-1}(\mathfrak m^{i+c}E'^{S_0})\subset K+\mathfrak m^iE$ for all $i$. For a finitely generated $A$-module $X$ of rank $r_X$ there is $c_X$ with $\operatorname{length}(X/\mathfrak m^jX)\le r_Xj+c_X$ for all $j$. If $j<c$, the image has length at most $\operatorname{length}(E/\mathfrak m^jE)\le r_Ej+c_E$. If $j=i+c$, the image is contained in the sum of the images of $K$ and of $\mathfrak m^iE$; the first is a quotient of $K/\mathfrak m^jK$, of length at most $rj+c_K$, and the second is the kernel of $E/\mathfrak m^jE\to E/\mathfrak m^iE$, of length at most $r_Ec+c_E$. So $e=c_K+r_Ec+c_E$ works. A module of finite length $\lambda$ over $A$ has $|\kappa|^\lambda$ elements.
% lemma_B: lemma_B_le (the finite S_0 by set_has_maximal_iff_noetherian, then comap_pow_le, which is Artin-Rees), len_fg_quot, lemma_B_sup; card_eq_pow_length. Pa.fixBd applies it (through the private fixBd_abs) to the reductions modulo T^m of sets of fixed matrices.
\end{proof}

\begin{lemma}\label{sha-patching-tangent-invariants}
In the setting of Definition~\ref{sha-patching-tangent-modules}, with $\ell$ odd:
\begin{enumerate}
\item there is $c_1$ such that, for every $m$, at most $|\kappa|^{c_1}$ matrices $Y\in W_m$ satisfy ${}^tY=-Y$ and $\rho_1(g)Y\rho_1(g)^{-1}=Y$ for all $g\in\Gamma_0$;
\item if $x$ satisfies (N4), there is $c_0$ such that, for every $m$, at most $|\kappa|^{c_0}$ matrices $Y\in W^*_m$ satisfy $\varepsilon(g)\rho_1(g)Y\rho_1(g)^{-1}=Y$ for all $g\in D_{v_0}$;
\item if $\sigma\in\Gamma_0$ and the characteristic polynomial of $\rho(\sigma)$ is $\prod_i(X-\alpha_i)$ with pairwise distinct $\alpha_i\in\kappa[[T]]$, there is $C_\sigma$ such that, for every $m$, at most $|\kappa|^{mn+C_\sigma}$ matrices $Y\in W^*_m$ commute with $\rho_1(\sigma)$.
\end{enumerate}
\end{lemma}
\leanref{Development/AutomorphicArithmetic/AutomorphicVarieties/PatchingModules/EpsilonEtaMatrices/CommutingMatrixOffdiagZero.lean}{UpperTheory.V.Pa.h0G}
\leanref{Development/AutomorphicArithmetic/AutomorphicVarieties/PatchingModules/EpsilonEtaMatrices/Rho1MapReductionLevel.lean}{UpperTheory.V.Pa.h0v0}
\leanref{Development/AutomorphicArithmetic/AutomorphicVarieties/PatchingModules/EpsilonEtaMatrices/FinrankCentralizerW1Le.lean}{UpperTheory.V.Pa.h0q}
\begin{proof}
In each case the matrices counted are the reductions modulo $T^m$ of matrices $Y$ over $R_g$, in the module $W$ (that is $W_\infty$) or $W^*=V^*$, with $f_s(Y)\in T^m(\cdot)$ for the linear maps $f_s$ given by $Y\mapsto{}^tY+Y$ and $Y\mapsto\rho_1(g)Y\rho_1(g)^{-1}-Y$, respectively $Y\mapsto\varepsilon(g)\rho_1(g)Y\rho_1(g)^{-1}-Y$ and $Y\mapsto\rho_1(\sigma)Y\rho_1(\sigma)^{-1}-Y$. By Lemma~\ref{sha-patching-tangent-fixed}, with $A=\kappa[[T]]$ and $j=m$, it suffices to show that the module of $Y$ over $\kappa[[T]]$ with $f_s(Y)=0$ for all $s$ has rank $0$, $0$ and at most $n$ respectively.

(1) First, an antisymmetric matrix $Z$ over $\kappa[[T]]$ that commutes with $\rho(g)$ for every $g\in\Gamma_0$ is $0$. Let $i\ne j$. The diagonal unit matrix $e_{jj}$ is even, so by Lemma~\ref{sha-patching-normalization} some $T^se_{jj}$ lies in the $\kappa[[T]]$-span of $\rho(\Gamma')$, and $Z$ commutes with $T^se_{jj}$. The $(i,j)$ entries of $ZT^se_{jj}$ and $T^se_{jj}Z$ are $T^sZ_{ij}$ and $0$, so $Z_{ij}=0$. And $2Z_{ii}=0$ since ${}^tZ=-Z$, so $Z_{ii}=0$, as $2$ is a unit in $\kappa[[T]]$ for $\ell$ odd.

Now let $Y\in W$ be antisymmetric and invariant. Its component on $1$ is an antisymmetric matrix over $\kappa[[T]]$ commuting with $\rho(\Gamma_0)$, so it is $0$. Then the component of the invariance on $\epsilon_k$ says that the $\epsilon_k$-component of $Y$ commutes with $\rho(\Gamma_0)$, and it is antisymmetric, so it is $0$. The components on $\epsilon'_k$ and $\eta$ vanish by the definition of $W$.

(2) Let $Y\in V^*$ be invariant under $D_{v_0}$. Its components on $1$ and the $\epsilon_k$ vanish. For $\phi$ the coefficient of $\epsilon'_k$ or of $\eta$, the products with $\rho_1(g)$ only involve the components on $1$ of the factors, so the matrix $\phi(Y)$ satisfies $\varepsilon(g)\rho(g)\phi(Y)\rho(g)^{-1}=\phi(Y)$ for $g\in D_{v_0}$; by (N4) it is $0$, first for the $\epsilon'_k$ and then for $\eta$.

(3) Let $P=\rho(\sigma)$, an even matrix. Over the fraction field of $\kappa[[T]]$, $P$ has $n$ distinct eigenvalues, so it is diagonalizable; the matrices commuting with $P$ form a space of dimension $n$, the polynomials in $P$, and any two of them commute. Let $Y\in V^*$ commute with $\rho_1(\sigma)$, and write $Y=\eta Y_0+\sum_k\epsilon'_kY_k$ with $Y_0$ even and the $Y_k$ odd. Since the components of $Y$ on $1$ and on the $\epsilon_k$ vanish, the $\epsilon'_k$-component of $\rho_1(\sigma)Y=Y\rho_1(\sigma)$ is $PY_k=Y_kP$. The diagonal matrix $E$ with entries $1$ in the first block and $0$ in the second commutes with $P$, since $P$ is even, so $EY_k=Y_kE$; comparing the entries in the two off-diagonal blocks gives $Y_k=0$. Then the $\eta$-component of $\rho_1(\sigma)Y=Y\rho_1(\sigma)$ is $PY_0=Y_0P$. So $Y\mapsto Y_0$ is an injective linear map from the module of these $Y$ to the centralizer of $P$ in $M_n(\kappa[[T]])$, which has rank $n$; hence that module has rank at most $n$.
% (1): skew_invariant_M1_zero, from skew_commuting_matrix_zero (commuting_matrix_offdiag_zero with x.span for the off-diagonal entries, 2 a unit for the diagonal) and map_ep_mul_eq_leibniz for the epsilon_k components. (2): site_vector_first with sandwich_map_smul. (3): finrank_centralizer_W1_le, with centralizer_rank_le from AlgebraTools.centraliser_rank (over the fraction field: dimension n, and commuting matrices are polynomials in P); the epsilon'_k components vanish through E = diagonal (1 on the first block), then the eta components lie in the centralizer. All three go through Pa.fixBd.
\end{proof}

\begin{lemma}\label{sha-patching-tangent-trace}
Let $\ell$ be a prime, $\kappa$ a finite field of characteristic $\ell$, $m\ge0$ an integer and $A_m=\kappa[[T]]/T^m$. Let $\tau:A_m\to\mathbf Z/\ell$ be $0$ if $m=0$, and otherwise the map sending $y$ to the trace from $\kappa$ to $\mathbf F_\ell$ of the coefficient of $T^{m-1}$ in $y$. Then $\tau$ is additive, $\tau(ay)=\bar a\,\tau(y)$ for $a\in\mathbf Z_\ell$ with residue $\bar a\in\mathbf F_\ell$, and if $y\in A_m$ satisfies $\tau(yz)=0$ for all $z\in A_m$, then $y=0$.
\end{lemma}
\leanref{Development/AutomorphicArithmetic/AutomorphicVarieties/PatchingModules/DegreeCountCohomology/AuxiliaryVectorCutByHermitianLinearForm.lean}{UpperTheory.V.Pa.pair1}
\leanref{Development/AutomorphicArithmetic/AutomorphicVarieties/AlgebraTools/FrobFormTrunc.lean}{UpperTheory.V.AlgebraTools.frob_form_trunc}
\begin{proof}
The map $\tau$ is additive, and $a\in\mathbf Z_\ell$ acts on $A_m$, a ring of characteristic $\ell$, as $\bar a$. Let $y\ne0$, so that $m\ge1$, and let $j<m$ be such that the coefficient $y_j$ of $T^j$ in $y$ is nonzero. The trace form of $\kappa$ over $\mathbf F_\ell$ is nondegenerate, so there is $b\in\kappa$ with $\operatorname{Tr}(y_jb)\ne0$. The coefficient of $T^{m-1}$ in $y\cdot bT^{m-1-j}$ is $y_jb$, so $\tau(y\cdot bT^{m-1-j})\ne0$.
\end{proof}

\begin{lemma}\label{sha-patching-tangent-local}
In the setting of Definition~\ref{sha-patching-tangent-modules}, let $v\mid\ell$, let $d_v=[L^+_v:\mathbf Q_\ell]$, and let $d'_v$ be as in Lemma~\ref{sha-patching-tangent-counts}(4). Then there is $C_v$ such that, for every $m$, the ordinary cocycles at $v$ form a finite group with at most $|\kappa|^{m(n(d'_v+1)+\nu(1+d_v))+C_v}$ elements.
\end{lemma}
\leanref{Development/AutomorphicArithmetic/AutomorphicVarieties/PatchingModules/EpsilonEtaMatrices/EntryZero.lean}{UpperTheory.V.Pa.locV}
\leanref{Development/AutomorphicArithmetic/AutomorphicVarieties/PatchingModules/DegreeCountCohomology/HomCount.lean}{UpperTheory.V.Pa.homCount}
\begin{proof}
Let $\beta$ be an ordinary cocycle and $B(g)={w'_v}^{-1}\beta(g)w'_v$. Each column of $w_v$ lies in one of the two blocks, so $B(g)$ has entries in $A_m$ at the positions $(i,j)$ whose columns $i$ and $j$ of $w_v$ lie in the same block, and in $\sum_kA_m\epsilon_k$ at the other positions; and $B(g)$ is upper triangular. Its coordinates are therefore indexed by the diagonal positions, the positions $i<j$ of the same block, and the pairs of a position $i<j$ of different blocks with an index $k$; there are $n$, $\binom a2+\binom{n-a}2$ and $g\,a(n-a)$ of them, $n+\nu$ in all. Order them by the height $j-i$ of their position. By the triangular form of $\rho_1$ on $D_v$, if all coordinates of smaller height vanish, the coordinate at a diagonal position is a homomorphism $D_v\to A_m$ trivial on $\mu_v$, by the diagonal condition, and the coordinate at a position $i<j$ is a locally constant cocycle for the character $\psi_{v,i}\psi_{v,j}^{-1}$ with values in $A_m$. A homomorphism $D_v\to A_m$ trivial on $\mu_v$ factors through $\mathbf Z_\ell^{d'_v+1}$ by Lemma~\ref{sha-patching-tangent-counts}(4), since $A_m$ is an abelian $\ell$-group, and so through $(\mathbf Z/\ell)^{d'_v+1}$; there are at most $|A_m|^{d'_v+1}=|\kappa|^{m(d'_v+1)}$ of them. The cocycles for $\chi=\psi_{v,i}\psi_{v,j}^{-1}$, with $i<j$, number at most $|\kappa|^{m(1+d_v)+C_{ij}}$, for a constant $C_{ij}$. Indeed, $\chi$ is continuous, and $\chi\ne\varepsilon$ on $D_v$ by (N1), so there is $g_0\in D_v$ with $\chi(g_0)-\varepsilon(g_0)=uT^{C_{ij}}$ for a unit $u$ of $\kappa[[T]]$ and an integer $C_{ij}\ge0$. Let $M=A_m$, on which $g\in D_v$ acts by multiplication by $\chi(g)$, let $M'=A_m$, on which $g$ acts by $\chi'(g)=\varepsilon(g)\chi(g)^{-1}$, and let $b(y,y')=\tau(yy')$, with $\tau$ as in Lemma~\ref{sha-patching-tangent-trace}. These actions are locally constant, $b$ has trivial left and right kernels by that lemma, and $b(\chi(g)y,\chi'(g)y')=\tau(\varepsilon(g)yy')=\bar\varepsilon(g)b(y,y')$. By Lemma~\ref{sha-patching-tangent-counts}(2), the cocycles number $|A_m|^{1+d_v}\,|H^0(D_v,M')|$. If $y'\in M'$ is invariant, then $(\chi(g_0)-\varepsilon(g_0))y'=\chi(g_0)\big(y'-\chi'(g_0)y'\big)=0$, so $T^{C_{ij}}y'=0$, and at most $|\kappa|^{C_{ij}}$ elements of $A_m$ are killed by $T^{C_{ij}}$. A group with coordinates that vanish together only at $0$, such that each coordinate lies in a set $P_t$ whenever the coordinates of smaller height vanish, has at most $\prod_t|P_t|$ elements. The product of the bounds is $|\kappa|^{nm(d'_v+1)+\nu m(1+d_v)+\sum C_{ij}}$.
% locV: entry_zero and tri (the type pattern TyPat of w1^-1 Z w1), crd and crd_coc (the coordinate cocycles), homCount (LocAb), pair_count (from (x.N1 v).2 and psi_cont, then oneCount: the dual pair chRep chi, chRep (eps chi^-1) with b = tau(y z) from pair1, LocCount at D_v, and ker_le for the invariants killed by T^C), filt_ht (the filtration count), tyCard (the count nu).
\end{proof}

\begin{lemma}\label{sha-patching-tangent-skew}
In the setting of Definition~\ref{sha-patching-tangent-modules}, for every complex conjugation $c_i\in\Gamma_N$ with $\eta(c_i)=\varepsilon(c_i)=-1$, at least $|\kappa|^{m\nu}$ elements of $W_m$ are fixed by $c_i$.
\end{lemma}
\leanref{Development/AutomorphicArithmetic/AutomorphicVarieties/PatchingModules/EpsilonEtaMatrices/OneLevel.lean}{UpperTheory.V.Pa.oneLevel}
\begin{proof}
Let $h=c_ic\in\Gamma_{N,0}$, so that $c_i=hc$ and $chc^{-1}=c\,c_i=h^{-1}$, as $c^2=c_i^2=1$. Then $\varepsilon(h)=\varepsilon(c_i)\varepsilon(c)=1$. The relation ${}^t\rho_1(g)\rho_1(cgc^{-1})=\varepsilon(g)^{1-n}$ holds for $\rho_1$, since it holds for $\rho$, $\kappa_k(cgc^{-1})=-{}^t\kappa_k(g)$ and $\epsilon_k\epsilon_{k'}=0$; at $h$ it gives ${}^t\rho_1(h)\rho_1(h)^{-1}=1$, so $\rho_1(h)$ is symmetric. For $Y\in W_m$ antisymmetric, $\rho_1(h)Y\in W_m$, and $c_i\cdot(\rho_1(h)Y)=\rho_1(h)\big(-{}^t(\rho_1(h)Y)\big)\rho_1(h)^{-1}=\rho_1(h)Y\rho_1(h)\rho_1(h)^{-1}=\rho_1(h)Y$. So $Y\mapsto\rho_1(h)Y$ injects the antisymmetric elements of $W_m$ into the fixed points of $c_i$. The antisymmetric elements of $W_m$ are determined by their entries above the diagonal, which range over $A_m$ at the $\binom a2+\binom{n-a}2$ positions in the diagonal blocks and over $\sum_kA_m\epsilon_k$ at the $a(n-a)$ positions in the off-diagonal block; so there are $|A_m|^\nu=|\kappa|^{m\nu}$ of them.
% hfix in oneLevel, with card_pow_nu_le_card_skew_w1 (a private lemma of OneLevel.lean) and rho1_mul_mem_m1_of_mem_m1.
\end{proof}

\begin{lemma}\label{sha-patching-tangent-solutions}
Let $A$ be a complete discrete valuation ring with uniformizer $t$ and finite residue field $\kappa$. Let $C$ and $C'$ be $A$-modules on which multiplication by $t$ is injective, with $C$ complete and separated and $C'$ separated for the $t$-adic topology, let $\Phi:C\to C'$ be an $A$-linear map, and let $B\subset\ker\Phi$ be a submodule such that every $z\in\ker\Phi$ has a multiple $t^sz$, $s\ge0$, in $B$. For $m\ge0$ let $S_m=\Phi^{-1}(t^mC')$, and let $Q_m$ and $B_m$ be the images of $S_m$ and of $B$ in $C/t^mC$. Suppose that $Q_1$ is finite. Then $\ker\Phi$ is finitely generated, and there is $e\ge0$ such that for every $m$, $Q_m$ is finite and $|Q_m|\le|B_m|\cdot|\kappa|^e$.
\end{lemma}
\leanref{Development/AutomorphicArithmetic/AutomorphicVarieties/PatchingModules/DegreeCountCohomology/SolCnt.lean}{UpperTheory.V.Pa.solCnt}
\begin{proof}
\emph{Division.} If $c\in C$ and $t^jc\in S_{m+j}$, then $c\in S_m$: from $t^j\Phi(c)=t^{m+j}y$ it follows that $\Phi(c)=t^my$, as $t$ is injective on $C'$. Also $S_{m'}\subset S_m$ for $m\le m'$.

\emph{Finiteness of $Q_m$.} We argue by induction on $m$; $Q_0=0$. Let $F\subset S_{m+1}$ be a finite set with the same image in $C/t^mC$ as $S_{m+1}$, which is possible since that image lies in $Q_m$, and let $G\subset S_1$ be a finite set with image $Q_1$. For $x\in S_{m+1}$ there are $f\in F$ and $c$ with $x-f=t^mc$. By the division $c\in S_1$, so $c-g\in tC$ for some $g\in G$, and $x-f-t^mg\in t^{m+1}C$. So $Q_{m+1}$ is finite.

\emph{The kernel.} Since $\ker\Phi\subset S_1$, there is a finite $F_1\subset\ker\Phi$ such that every $x\in\ker\Phi$ is $f+tc$ for some $f\in F_1$ and $c\in C$; then $t\Phi(c)=0$, so $c\in\ker\Phi$. Hence $\ker\Phi/t\ker\Phi$ is finite. The module $\ker\Phi$ is separated, as $C$ is, and $A$ is complete, so $\ker\Phi$ is finitely generated \cite[Theorem~8.4]{Mat86}. Choosing for each of finitely many generators an exponent that puts its multiple in $B$, and letting $s$ be their sum, gives $t^s\ker\Phi\subset B$. The module $\ker\Phi/t^s\ker\Phi$ is finite, so there is a finite $G_0\subset\ker\Phi$ with $\ker\Phi\subset G_0+t^s\ker\Phi\subset G_0+B$.

\emph{Stabilization.} For $M\ge1$ the images $I_M$ of the $S_M$ in $C/tC$ lie in $Q_1$ and decrease with $M$. Choose $m_0\ge1$ with $|I_{m_0}|$ least. Then $I_M=I_{m_0}$ for $M\ge m_0$, so for every $x\in S_{m_0}$ and every $M$ there is $s\in S_M$ with $x-s\in tC$.

\emph{Approximation.} For $j\ge0$, $x\in S_{m_0+j}$ and every $M$, there is $s\in S_M$ with $x-s\in t^jC$. We argue by induction on $j$; for $j=0$ take $s=0$. Let $x\in S_{m_0+j+1}$ and $M'=\max(M,m_0+j+1)$. The stabilization gives $s_1\in S_{M'}$ and $c_1$ with $x-s_1=tc_1$. By the division, $c_1\in S_{m_0+j}$, and the induction hypothesis gives $s_2\in S_M$ with $c_1-s_2\in t^jC$. Then $s_1+ts_2\in S_M$ and $x-s_1-ts_2\in t^{j+1}C$.

\emph{A kernel element.} Let $x\in S_{m_0+j}$. Then there are $k\in\ker\Phi$ and $c\in C$ with $x-k=t^jc$. Indeed, the approximation, applied repeatedly, gives $r_0=x$ and $r_{i+1}\in S_{m_0+j+i+1}$ with $r_i-r_{i+1}=t^{j+i}d_i$ for some $d_i\in C$. The series $c=\sum_it^id_i$ converges in $C$, and $x-t^jc=r_i-t^j\sum_{i'\ge i}t^{i'}d_{i'}$ for every $i$. So $\Phi(x-t^jc)\in t^iC'$ for every $i$, and $k=x-t^jc$ lies in $\ker\Phi$, as $C'$ is separated.

\emph{The count.} Let $P=\sum_{i<m_0}|Q_i|$, let $F_0\subset S_{m_0}$ be a finite set with image $Q_{m_0}$, and let $N=|G_0|\cdot|F_0|$. If $m<m_0$, then $|Q_m|\le P$. Let $m=m_0+j$ and $x\in S_m$, and take $k$ and $c$ as above. By the division, $c\in S_{m_0}$, so $c-f\in t^{m_0}C$ for some $f\in F_0$; and $k-g\in B$ for some $g\in G_0$. Then $x-(g+t^jf)-(k-g)=t^j(c-f)\in t^mC$. So every element of $Q_m$ is the image of some $g+t^jf$ plus an element of $B_m$, and $|Q_m|\le N\cdot|B_m|$. Since $|B_m|\ge1$ and $N+P\le|\kappa|^{N+P}$, the exponent $e=N+P$ works.
% solCnt: sc_kerfin (exists_finset_ker_cover_mod_span_smul_module, then UpperTheory.V.abelian₂.nakayama), sc_unif (the exponent s), sc_cover with scS_fin for the zero map on ker Phi (the finite set G), scS_stab (m0), scS_step (approximation), scS_lim (the kernel element, from IsPrecomplete and IsHausdorff), sc_count; scS_fin is the finiteness of Q_m and scS_div the division.
\end{proof}

\begin{lemma}\label{sha-patching-tangent-detection}
In Situation~\ref{sha-hecke-situation}, suppose that every open subgroup of $\Gamma$ has finitely many continuous homomorphisms to $\mathbf Z/\ell$. Let $x$ be a normalized arc over $\kappa[[T]]$ with a Selmer basis, and let $R_g$, $\rho_1$ and $V^*$ be as in Definition~\ref{sha-patching-selmer}. For $m\ge0$, reduce $\rho_1$ and $\varepsilon$ modulo $T^m$. A \emph{dual cocycle modulo $T^m$} is a locally constant map $z:\Gamma_0\to V^*/T^mV^*$ with $z(gh)=z(g)+\varepsilon(g)\rho_1(g)z(h)\rho_1(g)^{-1}$ and $z(cgc^{-1})=-\varepsilon(c)\,{}^tz(g)$ for all $g,h$; the \emph{dual coboundaries modulo $T^m$} are the maps $g\mapsto\varepsilon(g)\rho_1(g)Z\rho_1(g)^{-1}-Z$ with $Z\in V^*/T^mV^*$ and $-\varepsilon(c)\,{}^tZ=Z$; and for $\sigma\in\Gamma_0$, $\langle z,\sigma\rangle\in\kappa[[T]]/T^m$ is the coefficient of $\eta$ in $\operatorname{tr}(z(\sigma)\rho_1(\sigma))$. Let $\sigma_1,\dots,\sigma_q\in\Gamma'$ be such that $\{\sigma_1,\dots,\sigma_q\}$ detects the dual Selmer group of $x$. Then there is $c_2$ such that, for every $m$, the dual cocycles $z$ modulo $T^m$ with $\langle z,\sigma_j\rangle=0$ for all $j$ are finite in number, and number at most $|\kappa|^{c_2}$ times the dual coboundaries modulo $T^m$.
\end{lemma}
\leanref{Development/AutomorphicArithmetic/AutomorphicVarieties/PatchingModules/EpsilonEtaMatrices/ExistsNcardAnnihilatedCocyclesLeCoboundariesMul.lean}{UpperTheory.V.Pa.exists_ncard_annihilated_cocycles_le_coboundaries_mul}
\leanref{Development/AutomorphicArithmetic/AutomorphicVarieties/PatchingModules/EpsilonEtaMatrices/FinTwLc.lean}{UpperTheory.V.Pa.finTw}
\begin{proof}
Apply Lemma~\ref{sha-patching-tangent-solutions} with $A=\kappa[[T]]$ and $t=T$. Let $C$ be the module of maps $z:\Gamma_0\to M_n(R_g)$ that are continuous for the $T$-adic topology, and $C'$ the module of families of matrices over $R_g$ indexed by the disjoint union of $\Gamma_0$, $\Gamma_0\times\Gamma_0$, $\Gamma_0$ and $\{1,\dots,q\}$. Let $\Phi(z)$ be the family whose entry at $g$ in the first copy of $\Gamma_0$ is the matrix of the coordinates of $z(g)$ that vanish on $V^*$; whose entry at $(g,h)$ is $z(gh)-z(g)-\varepsilon(g)\rho_1(g)z(h)\rho_1(g)^{-1}$; whose entry at $g$ in the second copy is $z(cgc^{-1})+\varepsilon(c)\,{}^tz(g)$; and whose entry at $j$ is the scalar matrix with entries the coefficient of $\eta$ in $\operatorname{tr}(z(\sigma_j)\rho_1(\sigma_j))$. Let $B$ be the module of dual coboundaries over $\kappa[[T]]$. Then $\ker\Phi$ is the module of dual cocycles $z$ with $\langle z,\sigma_j\rangle=0$ for all $j$.

The hypotheses of that lemma hold. Multiplication by $T$ is injective on $C$ and $C'$, since $R_g$ is free over $\kappa[[T]]$, and $C'$ is separated. The module $C$ is separated, and it is complete: if $z_m\equiv z_{m'}$ modulo $T^m$ for $m\le m'$, the entrywise limit $z$ is continuous, since $z\equiv z_m$ modulo $T^m$. A map in $C$ whose values lie in $T^mM_n(R_g)$ is $T^m$ times a map in $C$, as in the step \emph{Reduction modulo $T$} of the proof of Lemma~\ref{sha-patching-frobenius-cocycles}. A dual coboundary is a dual cocycle, and its pairing with $\sigma_j$ vanishes, since $\varepsilon(\sigma_j)=1$ for $\sigma_j\in\Gamma'$ and $\operatorname{tr}\big((\rho_1(\sigma_j)Z\rho_1(\sigma_j)^{-1}-Z)\rho_1(\sigma_j)\big)=0$; so $B\subset\ker\Phi$. Every $z\in\ker\Phi$ has a multiple $T^sz$ in $B$, since $\{\sigma_1,\dots,\sigma_q\}$ detects the dual Selmer group.

Reduction modulo $T^m$ identifies $Q_m$ with the set of dual cocycles modulo $T^m$ whose pairings with the $\sigma_j$ vanish. Indeed, $\Phi(z)\in T^mC'$ says that the reduction of $z$ is such a cocycle; conversely, such a cocycle takes finitely many values and is locally constant, so choosing a lift of each value gives an element of $C$; and the kernel of the reduction on $C$ is $T^mC$. In the same way $B_m$ is identified with the set of reductions of dual coboundaries over $\kappa[[T]]$, which are dual coboundaries modulo $T^m$. The dual cocycles modulo $T$ are finite in number, since they are locally constant cocycles with values in a finite module killed by $\ell$ (the step \emph{Finiteness modulo $T$} of the proof of Lemma~\ref{sha-patching-frobenius-cocycles}), so $Q_1$ is finite. Lemma~\ref{sha-patching-tangent-solutions} gives $c_2=e$.
% exists_ncard_annihilated_cocycles_le_coboundaries_mul: dg_abs (C the continuous maps, with mC, hH, hP for separation and completeness, hS for the identification of Q_m, dg_equiv) applied with dgL = dgPhi (dgCut for V^*, the cocycle relation, the symmetry, the scalar ell(sigma_j)), dgB = TwB over x.A, hdet from Detects; Z 1 finite by finTw (which extends to cocycles of Gamma by polZ and counts them with finZ under PhiCond).
\end{proof}

\begin{lemma}\label{sha-patching-tangent-descent}
In the setting of Definition~\ref{sha-patching-tangent-modules}, suppose that $\rho$ is trivial on the kernel $P(H)$ of the maximal pro-$\ell$ quotient of $H$. Let $y:\Gamma_N\to W^*_m$ be a locally constant cocycle with $y(c)=0$ and $y(\tau_{N,j})=0$ for all $j$. Then there is a dual cocycle $z$ modulo $T^m$, in the sense of Lemma~\ref{sha-patching-tangent-detection}, such that $y(g)=z(\bar g)$ for $g\in\Gamma_{N,0}$ with image $\bar g$ in $\Gamma_0$.
\end{lemma}
\leanref{Development/AutomorphicArithmetic/AutomorphicVarieties/PatchingModules/EpsilonEtaMatrices/Descend.lean}{UpperTheory.V.Pa.descend}
\begin{proof}
Let $U$ be the set of $g\in\Gamma_N$ that act trivially on $W^*_m$ and have $y(g)=0$. By the cocycle relation, $g\mapsto(Y\mapsto y(g)+g\cdot Y)$ is an action of $\Gamma_N$ on the finite set $W^*_m$, and $U$ is its kernel; since $y$ and the action are locally constant, $U$ is an open, hence closed, normal subgroup. We show that $U$ contains the kernel of $\Gamma_N\to\Gamma$.

An element $g\in H_N$ whose image in $\Gamma_0$ has $\rho_1(g)\equiv1$ modulo $T^m$ acts trivially on $W^*_m$: the image of $g$ lies in $H$, so $\varepsilon(g)\equiv1\pmod\ell$, and $\varepsilon$ acts on $W^*_m$ through its residue modulo $\ell$, as $\ell=0$ in $\kappa$. In particular each $\tau_{N,j}$, which maps to $1$ in $\Gamma$, acts trivially, and $y(\tau_{N,j})=0$, so $\tau_{N,j}\in U$.

For an open normal subgroup $H'$ of a profinite group $\Gamma'$, let $\mathcal F(H')$ be the family of open normal subgroups $U'$ of $\Gamma'$ such that every $h\in H'$ has a power $h^{\ell^e}$ in $U'$. These are the open normal subgroups of $\Gamma'$ that contain $P(H')$, since $H'/(H'\cap U')$ is finite; and $P(H')$ is their intersection, since an open normal subgroup $V'$ of $H'$ with $H'/V'$ an $\ell$-group contains the intersection of its finitely many conjugates under $\Gamma'$, which lies in $\mathcal F(H')$. The family $\mathcal F(H')$ is closed under finite intersections: if $h^{\ell^e}\in U'$ and $h^{\ell^{e'}}\in U''$, then $h^{\ell^{e+e'}}$ lies in both.

Let $h\in H_N$. The set $V$ of $g\in\Gamma_0$ with $\rho(g)\equiv1$ modulo $T^m$ is open in $\Gamma$, as $\rho$ is continuous and $\Gamma_0$ is open, and it contains $P(H)$, on which $\rho$ is trivial. Its complement is closed in the compact group $\Gamma$, hence compact, and it does not meet the intersection $P(H)$ of the closed subgroups $U'\in\mathcal F(H)$. Since $\mathcal F(H)$ is closed under finite intersections, some $U'\in\mathcal F(H)$ is contained in $V$. So the image of $h^{\ell^e}$ in $\Gamma$ lies in $V$ for some $e$, that is, $\rho(h^{\ell^e})\equiv1$ modulo $T^m$. Then $\rho_1(h^{\ell^e})=1+E$ modulo $T^m$ with $E=\sum_k\epsilon_k\kappa_k(h^{\ell^e})$ and $E^2=0$, so $\rho_1(h^{\ell^{e+1}})=1+\ell E=1$ modulo $T^m$. So $h_1=h^{\ell^{e+1}}$ acts trivially on $W^*_m$, and $y(h_1^\ell)=\ell\,y(h_1)=0$, as $W^*_m$ is killed by $\ell$; hence $h_1^\ell\in U$. So $U\in\mathcal F(H_N)$, and $U$ contains $P(H_N)$. By Lemma~\ref{sha-patching-places}(3), $U$ contains the kernel of $\Gamma_N\to\Gamma$, as $U$ is closed and normal.

For $g\in\Gamma_N$ and $k$ in that kernel, $y(gk)=y(g)+g\cdot y(k)=y(g)$. The map $\Gamma_N\to\Gamma$ is a continuous surjection of compact groups, hence a quotient map, so $y$ is the composite of it with a locally constant map $Y$ on $\Gamma$. Let $z$ be the restriction of $Y$ to $\Gamma_0$, each of whose elements is the image of an element of $\Gamma_{N,0}$. The cocycle relation of $y$ on $\Gamma_{N,0}$ gives that of $z$, and since $y(c)=0$, Lemma~\ref{sha-patching-tangent-symmetry} gives $z(cgc^{-1})=-\varepsilon(c)\,{}^tz(g)$.
% descend: affine_ker (the open normal subgroup U), htriv with epsD_one and T.H_eps, rho1_pow_eq_one_pro_l (pow_mem_of_PL: PL l H is the infimum over the family F(H), and elim_directed_family_closed on the compact complement of V; then one_add_pow with epM_mul), cocycle_pow, T.ker_le (Lemma sha-patching-places (3)) with normalClosure_le_normal and the infimum defining PL, factor_lc (the quotient map), and polZ for the symmetry.
\end{proof}

\begin{lemma}\label{sha-patching-tangent-dual}
In the setting of Definition~\ref{sha-patching-tangent-modules}, suppose that $\ell$ is odd, that every open subgroup of $\Gamma$ has finitely many continuous homomorphisms to $\mathbf Z/\ell$, that $\rho$ is trivial on the kernel of the maximal pro-$\ell$ quotient of $H$, and that $\{\sigma_1,\dots,\sigma_q\}$ detects the dual Selmer group of $x$. Then there is $c_2$ such that for every $m$ and all large $N$, the cocycles $y\in Z^1(\Gamma_N,W^*_m)$ for which, for each $j$, some $Y\in W^*_m$ has $y(\phi_{N,j})=\phi_{N,j}\cdot Y-Y$ and $y(\tau_{N,j})=\tau_{N,j}\cdot Y-Y$ number at most $\#B^1(\Gamma_N,W^*_m)\cdot|\kappa|^{c_2}$.
\end{lemma}
\leanref{Development/AutomorphicArithmetic/AutomorphicVarieties/PatchingModules/EpsilonEtaMatrices/DualN.lean}{UpperTheory.V.Pa.dualN}
\leanref{Development/AutomorphicArithmetic/AutomorphicVarieties/PatchingModules/EpsilonEtaMatrices/TwistedCocycleApplyOne.lean}{UpperTheory.V.Pa.phEq}
\begin{proof}
\emph{Detection modulo $T^m$.} The $\sigma_j$ lie in $\Gamma'$ by Lemma~\ref{sha-patching-places}. Let $c_2$ be the constant of Lemma~\ref{sha-patching-tangent-detection}: for every $m$, the dual cocycles modulo $T^m$, in the sense of that lemma, whose pairings with all $\sigma_j$ vanish number at most $|\kappa|^{c_2}$ times the dual coboundaries modulo $T^m$.

\emph{Large $N$.} For all large $N$, $\rho_1(\phi_{N,j})=\rho_1(\sigma_j)$ and $\varepsilon(\phi_{N,j})=1$ modulo $T^m$, and every dual cocycle $z$ modulo $T^m$ has $z(\phi_{N,j})=z(\sigma_j)$, for all $j$. Indeed, there are finitely many dual cocycles modulo $T^m$, since they are locally constant cocycles with values in a finite module killed by $\ell$ (the step \emph{Finiteness modulo $T$} of the proof of Lemma~\ref{sha-patching-frobenius-cocycles}), so it suffices to treat one $z$ and one $j$. Since $\rho_1$, $\varepsilon$ and $z$ modulo $T^m$ are locally constant, the $u\in\Gamma_0$ with $\rho_1(u)=1$, $\varepsilon(u)=1$ and $z(u)=0$ form an open subgroup $U$, on which $z$ is a homomorphism by the cocycle relation. Every $h\in H$ has a power $h^{\ell^e}$ in $U$: $\varepsilon(h)=1$ modulo $T^m$, since $\varepsilon\bmod\ell$ is trivial on $H$; $\rho_1(h)$ modulo $T^m$ has $\ell$-power order, since $\rho$ is trivial on $P(H)$ and $\rho_1$ is residually trivial on $H$; and once $\rho_1(h^{\ell^{e}})=1$, $z(h^{\ell^{e+1}})=\ell z(h^{\ell^e})=0$. Hence the normal core of $U\cap H$ in $H$ is an open normal subgroup with quotient an $\ell$-group, so it contains $P(H)$, and there is an open normal subgroup $V\subset U$ of $\Gamma$ containing $P(H)$. By Lemma~\ref{sha-patching-places}(1), $\phi_{N,j}\in V\sigma_j$ for all large $N$, and $\phi_{N,j}=u\sigma_j$ with $u\in U$ gives the three identities, since $\varepsilon(\sigma_j)=1$.

\emph{Descent.} Let $\mathcal Z$ be the set of cocycles of the statement, a group containing $B^1(\Gamma_N,W^*_m)$. Since $\tau_{N,j}$ maps to $1$ in $\Gamma$, it acts trivially on $W^*_m$, so $y(\tau_{N,j})=0$ for $y\in\mathcal Z$. By Lemma~\ref{sha-patching-tangent-descent}, the restriction to $\Gamma_{N,0}$ of a cocycle $y\in\mathcal Z$ with $y(c)=0$ is the composite with $\Gamma_{N,0}\to\Gamma_0$ of a dual cocycle $\Psi(y)$ modulo $T^m$. The pairing of $\Psi(y)$ with $\sigma_j$ vanishes: $\Psi(y)(\sigma_j)=y(\phi_{N,j})=\rho_1(\sigma_j)Y\rho_1(\sigma_j)^{-1}-Y$, and $\operatorname{tr}\big((\rho_1(\sigma_j)Y\rho_1(\sigma_j)^{-1}-Y)\rho_1(\sigma_j)\big)=0$. By Lemma~\ref{sha-patching-tangent-symmetry}, $\Psi$ is injective on $\{y\in\mathcal Z:y(c)=0\}$, and every dual coboundary modulo $T^m$, of a $Z$ with $-\varepsilon(c)\,{}^tZ=Z$, is $\Psi$ of the coboundary of $Z$ on $\Gamma_N$, which vanishes at $c$.

\emph{Count.} So $\#\{y\in\mathcal Z:y(c)=0\}\le|\kappa|^{c_2}\#\{y\in B^1(\Gamma_N,W^*_m):y(c)=0\}$. By Lemma~\ref{sha-patching-tangent-involution}, $\#\mathcal Z\cdot\#\{y\in B^1:y(c)=0\}=\#\{y\in\mathcal Z:y(c)=0\}\cdot\#B^1$, so $\#\mathcal Z\le|\kappa|^{c_2}\#B^1(\Gamma_N,W^*_m)$.
% dualN: exists_ncard_annihilated_cocycles_le_coboundaries_mul (Lemma sha-patching-tangent-detection), phEq (with phEq_one and finTw), descend (Lemma sha-patching-tangent-descent), polZ, cCount.
\end{proof}

\begin{lemma}\label{sha-patching-tangent}
In Situation~\ref{sha-hecke-situation}, suppose that $\varepsilon\bmod\ell$ is nontrivial on $\Gamma_0$ (this holds in the situation, since $L$ contains no primitive $\ell$-th root of unity, and the proof does not use it). Let $\sigma_1,\dots,\sigma_q$ and the places $w_{N,j}$ be as in Lemma~\ref{sha-patching-places}, let $x$ be a normalized arc trivial on the kernel of the maximal pro-$\ell$ quotient of $H$, satisfying (N4), with a Selmer basis detected by the $\sigma_j$ as in Lemma~\ref{sha-patching-frobenius}, and with distinct eigenvalues at each $\rho(\sigma_j)$. Then there is $C$ such that for every $m$ and all large $N$, the Selmer group of $\rho_1$ modulo $T^m$ at the level augmented by the $w_{N,j}$ is finite and has at most $|\kappa|^{m(n[L^+:\mathbf Q]+nq)+C}$ times as many elements as its group of coboundaries.
\end{lemma}
\leanref{Development/AutomorphicArithmetic/AutomorphicVarieties/PatchingModules/TaylorWilesObstructionBound.lean}{UpperTheory.V.Pa.FirstOrderDualArithmeticBound}
\leanref{Development/AutomorphicArithmetic/AutomorphicVarieties/PatchingModules/AutomorphicDetection/BoundW1Arithmetic.lean}{UpperTheory.V.Pa.b1Ar}
\begin{proof}
Use the notation of Definition~\ref{sha-patching-tangent-modules}, with $d=[L^+:\mathbf Q]$. Let $c_1$, $c_0$ and $C_j=C_{\sigma_j}$ be the constants of Lemma~\ref{sha-patching-tangent-invariants}, $C_v$ those of Lemma~\ref{sha-patching-tangent-local}, $c_2$ that of Lemma~\ref{sha-patching-tangent-dual}, and $d'_v$ as in Lemma~\ref{sha-patching-tangent-counts}(4); the hypotheses of these lemmas hold by Lemma~\ref{sha-patching-places}. Fix $m$, and let $N$ be large enough for Lemma~\ref{sha-patching-tangent-dual} and for the identities $\rho_1(\phi_{N,j})=\rho_1(\sigma_j)$ and $\varepsilon(\phi_{N,j})=1$ modulo $T^m$ of its proof. Write $W=W_m$ and $W^*=W^*_m$.

\emph{The pairing.} There is a biadditive pairing $b:W\times W^*\to\mathbf Z/\ell$ with trivial left and right kernels and $b(g\cdot Y,g\cdot Y')=\bar\varepsilon(g)b(Y,Y')$ for $g\in\Gamma_N$. Let $\tau:A_m\to\mathbf Z/\ell$ be as in Lemma~\ref{sha-patching-tangent-trace}, and let $b(Y,Y')$ be $\tau$ of the coefficient of $\eta$ in $\operatorname{tr}(YY')$. Let $Y\in W$ with $b(Y,Y')=0$ for all $Y'$, let $e_{ji}$ be the unit matrices, and let $z\in A_m$. If $i$ and $j$ lie in the same block, $Y_{ij}\in A_m$, and $Y'=\eta ze_{ji}\in W^*$ gives $\tau(Y_{ij}z)=0$; so $Y_{ij}=0$ by Lemma~\ref{sha-patching-tangent-trace}. If they lie in different blocks, $Y_{ij}=\sum_k\epsilon_ka_k$ with $a_k\in A_m$, and $Y'=\epsilon'_kze_{ji}\in W^*$ gives $\tau(a_kz)=0$, as $\epsilon_{k'}\epsilon'_k$ is $\eta$ for $k'=k$ and $0$ otherwise; so $a_k=0$. So the left kernel is trivial, and the right kernel is trivial in the same way, with $Y=ze_{ji}$ and $Y=\epsilon_kze_{ji}$. For $g\in\Gamma_{N,0}$, $\operatorname{tr}\big(\rho_1(g)Y\rho_1(g)^{-1}\cdot\varepsilon(g)\rho_1(g)Y'\rho_1(g)^{-1}\big)=\varepsilon(g)\operatorname{tr}(YY')$, and $\tau(\varepsilon(g)u)=\bar\varepsilon(g)\tau(u)$ for $u\in A_m$; for $c$, $\operatorname{tr}\big({}^tY\cdot\varepsilon(c)\,{}^tY'\big)=\varepsilon(c)\operatorname{tr}\big({}^t(Y'Y)\big)=\varepsilon(c)\operatorname{tr}(YY')$. Every element of $\Gamma_N$ outside $\Gamma_{N,0}$ is $gc$ with $g\in\Gamma_{N,0}$, and $c^2=1$, so $b(g\cdot Y,g\cdot Y')=\bar\varepsilon(g)b(Y,Y')$ for all $g\in\Gamma_N$.

\emph{The Greenberg-Wiles count.} Apply Lemma~\ref{sha-patching-tangent-counts}(1) to $M=W$, $M'=W^*$, $b$, and the local conditions $\mathcal L$ of Lemma~\ref{sha-patching-tangent-selmer}. The terms are bounded as follows.
\begin{enumerate}
\item At $v\mid\ell$. Let $b_v$ be the set of $Y\in W$ with ${w'_v}^{-1}Yw'_v$ upper triangular. It has $|\kappa|^{m(n+\nu)}$ elements, with one coordinate in $A_m$ at each diagonal position and each of the $\nu$ coordinates above the diagonal described in the proof of Lemma~\ref{sha-patching-tangent-local}, and for $Y\in b_v$ the coboundary $g\mapsto\rho_1(g)Y\rho_1(g)^{-1}-Y$ is an ordinary cocycle, since ${w'_v}^{-1}\rho_1(g)w'_v$ is upper triangular for $g\in D_v$. The map $(Y,\beta)\mapsto(g\mapsto g\cdot Y-Y)+\beta$ from $W$ times the ordinary cocycles $O_v$ onto $\mathcal L_v$ has a kernel containing the pairs $(Y,-\delta Y)$, $Y\in b_v$; so $|\mathcal L_v|\cdot|\kappa|^{m(n+\nu)}\le|W|\cdot|O_v|\le|W|\cdot|\kappa|^{e_v}$, with $e_v=m(n(d'_v+1)+\nu(1+d_v))+C_v$, by Lemma~\ref{sha-patching-tangent-local}.
\item At $u\in S_{\mathrm a}$. The group $I'_u$ acts trivially on $W$, since $\rho$ and the $\kappa_k$ are trivial on it, and $W$ is killed by $\ell$; so $|\mathcal L_u|=|W|$ by Lemma~\ref{sha-patching-tangent-counts}(3).
\item At $w=v_0$ and $w=w_{N,j}$. Here $\mathcal L_w=Z^1(D_w,W)$, so $|\mathcal L_w|=|W|\cdot|H^0(D_w,W^*)|$ by Lemma~\ref{sha-patching-tangent-counts}(2). The invariants of $D_{v_0}$ in $W^*$ number at most $|\kappa|^{c_0}$, by Lemma~\ref{sha-patching-tangent-invariants}(2). The decomposition group at $w_{N,j}$ contains $\phi_{N,j}$, which acts on $W^*$ by $Y\mapsto\rho_1(\sigma_j)Y\rho_1(\sigma_j)^{-1}$ by the choice of $N$, so its invariants number at most $|\kappa|^{mn+C_j}$, by Lemma~\ref{sha-patching-tangent-invariants}(3) and the hypothesis on the eigenvalues.
\item At the infinite places. $|W^{c_i}|\ge|\kappa|^{m\nu}$ for each $i$, by Lemma~\ref{sha-patching-tangent-skew}.
\item The dual term. The places $w$ with $\mathcal L_w=Z^1(D_w,W)$ include the $w_{N,j}$, and $\phi_{N,j}$ and $\tau_{N,j}$ lie in the decomposition group at $w_{N,j}$. So a $y\in Z^1_{\mathcal L^\perp}$ satisfies the condition of Lemma~\ref{sha-patching-tangent-dual}, and $\#Z^1_{\mathcal L^\perp}\le\#B^1(\Gamma_N,W^*)\cdot|\kappa|^{c_2}$.
\end{enumerate}
The set $P_N$ has $|V_\ell|+|S_{\mathrm a}|+1+q$ elements, where $V_\ell$ is the set of places above $\ell$. Putting (1) to (5) into Lemma~\ref{sha-patching-tangent-counts}(1) and dividing by $\#B^1(\Gamma_N,W^*)\cdot|W|^{|P_N|}$ gives
\[\#Z^1_{\mathcal L}\cdot|\kappa|^{m(n+\nu)|V_\ell|+m\nu d}\le|W|\cdot|\kappa|^{c_2+\sum_ve_v+c_0+\sum_j(mn+C_j)}.\]

\emph{The Selmer group.} By Lemma~\ref{sha-patching-tangent-selmer}, the Selmer group is finite, $\#\mathrm{Sel}\cdot|W|\le\#Z^1_{\mathcal L}\cdot|W^c|$ and $|W^c|=\#\mathrm{Cob}\cdot|W^{\Gamma_N}|$. An element of $W^{\Gamma_N}$ is fixed by $c$, so it is antisymmetric, and it is fixed by $\Gamma_{N,0}$, which acts through $\Gamma_0$; so $|W^{\Gamma_N}|\le|\kappa|^{c_1}$ by Lemma~\ref{sha-patching-tangent-invariants}(1). Hence
\[\#\mathrm{Sel}\cdot|\kappa|^{m(n+\nu)|V_\ell|+m\nu d}\le\#\mathrm{Cob}\cdot|\kappa|^{c_1+c_2+c_0+\sum_ve_v+\sum_j(mn+C_j)}.\]

\emph{The exponent.} Since $\sum_vd'_v=d$ by Lemma~\ref{sha-patching-tangent-counts}(4) and $\sum_vd_v=d$, the exponent on the right is $m(nd+nq)+C+m(n+\nu)|V_\ell|+m\nu d$, with $C=c_1+c_2+c_0+\sum_vC_v+\sum_jC_j$ independent of $m$ and $N$. Cancelling $|\kappa|^{m(n+\nu)|V_\ell|+m\nu d}$ gives $\#\mathrm{Sel}\le\#\mathrm{Cob}\cdot|\kappa|^{m(nd+nq)+C}$.
% The pairing: Pa.pairEx (pair1 for tau, blk_nd with nd1 to nd4 for the kernels, e0 for G_0, ec for c, and gamma = (gamma c) c). The count: oneLevel (sumCount, bCard and bMem for (1); ArData.Counts.unr for (2); ArData.Counts.locQ with h0v0 and h0q for (3); hfix for (4); dualN for (5); selN and h0G), then sum_exp and cancel_pow in W1_bound_arithmetic_rank_two. ArSys (arSys) supplies the counts from CohP, with ph and ta placing phi_{N,j} and tau_{N,j} in the decomposition groups and d0 identifying D_{v_0}.
% Pa.b1Ar applies Pa.b1abs, that is W1_bound_arithmetic_rank_two, with the arithmetic system arSys (from CohP), LocCFT and the PhiCond of arPl. The terms are h0G (1), dualN (2), h0v0 (3), h0q with phEq (4) and locV (5), combined by oneLevel; the local exponent at v is n(d'_v+1)+nu(1+dv_v) with sum_v d'_v = sum_v dv_v = [L^+:Q], and cancel_pow removes m((n+nu)#{v | l}+nu[L^+:Q]). The hypothesis that eps mod l is nontrivial on G_0 is not used (it is the unnamed argument in b1Ar).
\end{proof}

\begin{lemma}\label{sha-patching-modules}
Let $A$ be a discrete valuation ring with uniformizer $t$ and finite residue field $\kappa$, let $\mathcal H$ be a finite $A$-module and let $g,C,m\ge0$ be integers such that $t^m\mathcal H=0$ and $|\mathcal H/t^e\mathcal H|\le|\kappa|^{ge+C}$ for every $e\le m$. Then for all $\xi_0,\dots,\xi_g\in\mathcal H$ there are $c_0,\dots,c_g\in A$, at least one of them a unit, with $t^C\sum_ic_i\xi_i=0$.
\end{lemma}
\leanref{Development/AutomorphicArithmetic/AutomorphicVarieties/PatchingModules/ArchimedeanH0Cards/GenTrick.lean}{UpperTheory.V.Pa.B2.genTrick}
\begin{proof}
If $m\le C$, take all $c_i=1$, since $t^C$ kills $\mathcal H$. Otherwise suppose that no such $c_i$ exist, and put $e=C+1\le m$. Let $\phi:A^{g+1}\to\mathcal H$ be the map $c\mapsto\sum_ic_i\xi_i$, with image $N$.

First, if $d,d_0\in A^{g+1}$ satisfy $\phi(d)=t^e\phi(d_0)$, then $t^e$ divides every $d_i$. We show by induction on $k\le e$ that $t^k$ divides every $d_i$. If $d=t^kd'$ with $k<e$ and some $d'_{i_0}$ is not divisible by $t$, then $t^k\phi(d'-t^{e-k}d_0)=\phi(d)-t^e\phi(d_0)=0$, so $t^C\phi(d'-t^{e-k}d_0)=0$ as $k\le C$; and the coordinate $d'_{i_0}-t^{e-k}d_{0,i_0}$ is a unit, as $e-k\ge1$. This contradicts the assumption, so $t$ divides every $d'_i$.

Hence the map $(A/t^e)^{g+1}\to N/t^eN$ induced by $\phi$ is injective, and $|\kappa|^{e(g+1)}\le|N/t^eN|$. For a finite $A$-module $X$, multiplication by $t^e$ has kernel and cokernel of the same size. So $|N/t^eN|$ is the number of elements of $N$ killed by $t^e$, which is at most the number of elements of $\mathcal H$ killed by $t^e$, which is $|\mathcal H/t^e\mathcal H|\le|\kappa|^{ge+C}$. Since $e(g+1)=ge+C+1$, this is a contradiction.
\end{proof}

\begin{lemma}\label{sha-patching-obstruction-count}
Let $\Gamma$ be a topological group and $P$ a finite set, and for $w\in P$ let $D_w$ be a topological group with a continuous homomorphism $\varphi_w:D_w\to\Gamma$ and a subgroup $J_w\subset D_w$. Let $M$ be a finite module over a ring $A$ with a locally constant $A$-linear action of $\Gamma$, and for $w\in P$ let $Q_w$ be a finite $A$-module with a locally constant $A$-linear action of $D_w$ and a surjective $A$-linear map $\pi_w:M\to Q_w$ with $\pi_w(\varphi_w(g)y)=g\,\pi_w(y)$. Cohomology is that of locally constant cochains. Let $H^1_J(D_w,Q_w)$ be the kernel of the restriction map $H^1(D_w,Q_w)\to H^1(J_w,Q_w)$; let $\mathrm{Sel}$ be the group of classes in $H^1(\Gamma,M)$ whose image in $H^1(D_w,Q_w)$ lies in $H^1_J(D_w,Q_w)$ for every $w$; and let $\mathrm{Ck}$ be the cokernel of the map $H^2(\Gamma,M)\to\bigoplus_wH^2(D_w,Q_w)$. Let $H^2_{\mathrm{loc}}$ be the quotient of the group of families $(z,(y_w)_w)$, of a locally constant $2$-cocycle $z$ of $\Gamma$ with values in $M$ and locally constant maps $y_w:D_w\to Q_w$ whose coboundaries are the maps $(g,h)\mapsto\pi_w(z(\varphi_wg,\varphi_wh))$, by the families $(dx,(\pi_w\circ x\circ\varphi_w+a_w)_w)$, where $x:\Gamma\to M$ is locally constant and $a_w$ is a cocycle of $D_w$ with values in $Q_w$ whose restriction to $J_w$ is a coboundary. Suppose that $H^1$ and $H^2$ of $\Gamma$ with values in $M$, and of each $D_w$ with values in $Q_w$, are finite, and that there are integers $s\ge1$, $d_w\ge0$ and $r_w\ge0$ with
\[|H^2(\Gamma,M)|\,s^{\sum_wd_w}\le|H^1(\Gamma,M)|\quad\text{and}\quad|H^1(D_w,Q_w)|\le|H^1_J(D_w,Q_w)|\,|H^2(D_w,Q_w)|\,r_w\,s^{d_w}\text{ for every }w.\]
Then $H^2_{\mathrm{loc}}$ is finite and $|H^2_{\mathrm{loc}}|\le|\mathrm{Sel}|\cdot\prod_wr_w\cdot|\mathrm{Ck}|$.
\end{lemma}
\leanref{Development/AutomorphicArithmetic/AutomorphicVarieties/PatchingModules/ArchimedeanH0Cards/CardKerRange.lean}{UpperTheory.V.Pa.B2.count}
\begin{proof}
Let $\mathcal T=\bigoplus_wH^1(D_w,Q_w)/H^1_J(D_w,Q_w)$, and let $a:H^1(\Gamma,M)\to\mathcal T$ be the sum of the localization maps; its kernel is $\mathrm{Sel}$. Let $c:H^2_{\mathrm{loc}}\to H^2(\Gamma,M)$ send the class of $(z,(y_w))$ to the class of $z$, and let $\lambda:H^2(\Gamma,M)\to\bigoplus_wH^2(D_w,Q_w)$ be the localization map. The image of $c$ lies in $\ker\lambda$, since the image of $z$ in $H^2(D_w,Q_w)$ is the class of the coboundary of $y_w$.

For a family $(y_w)$ of cocycles $y_w$ of $D_w$ with values in $Q_w$, let $b((y_w))$ be the class of $(0,(y_w))$. If the class of $(z,(y_w))$ lies in $\ker c$, then $z=dx$ for some $x$, and $(z,(y_w))$ has the same class as $(0,(y_w-\pi_w\circ x\circ\varphi_w))$, whose components are cocycles; so $\ker c$ lies in the image of $b$. If the image of $(y_w)$ in $\mathcal T$ is $a$ of the class of a cocycle $x$ of $\Gamma$, then each $y_w-\pi_w\circ x\circ\varphi_w$ restricts to a coboundary on $J_w$, and $(0,(y_w))=(dx,(\pi_w\circ x\circ\varphi_w+(y_w-\pi_w\circ x\circ\varphi_w)))$, so $b((y_w))=0$. Choose for each element of $\ker c$ a family $(y_w)$ that $b$ sends to it, and send the element to the image of $(y_w)$ in $\mathcal T/a(H^1(\Gamma,M))$. This map is injective, since two families with the same image differ by a family that $b$ sends to $0$. Hence
\[|H^2_{\mathrm{loc}}|=|\ker c|\cdot|\operatorname{im}c|\le\frac{|\mathcal T|}{|a(H^1(\Gamma,M))|}\cdot\frac{|H^2(\Gamma,M)|}{|\operatorname{im}\lambda|}.\]
By the second hypothesis, $|\mathcal T|\le\prod_w|H^2(D_w,Q_w)|\,r_w\,s^{d_w}=|\operatorname{im}\lambda|\cdot|\mathrm{Ck}|\cdot\prod_wr_w\cdot s^{\sum_wd_w}$. By the first, $|H^2(\Gamma,M)|\,s^{\sum_wd_w}\le|H^1(\Gamma,M)|=|\mathrm{Sel}|\cdot|a(H^1(\Gamma,M))|$. Multiplying these bounds gives $|H^2_{\mathrm{loc}}|\le|\mathrm{Sel}|\cdot\prod_wr_w\cdot|\mathrm{Ck}|$.
% count: aMap (ker_aMap), bMap with ker_qMap and ker_cMap, inj_of (the injection of ker cMap into Tq / range aMap), range_cMap (into ker loc2), card_ker_range, and the arithmetic key.
\end{proof}

\begin{lemma}\label{sha-patching-obstruction}
In Situation~\ref{sha-hecke-situation}, suppose that $\varepsilon\bmod\ell$ is nontrivial on $\Gamma_0$ (as in Lemma~\ref{sha-patching-tangent}, the proof does not use it). Let $\sigma_1,\dots,\sigma_q$ and the places $w_{N,j}$ be as in Lemma~\ref{sha-patching-places}, and let $x$ be a normalized arc, with representation $\rho$, trivial on the kernel of the maximal pro-$\ell$ quotient of $H$, with distinct eigenvalues at each $\rho(\sigma_j)$, and with a Selmer basis of size $g$ as in Lemma~\ref{sha-patching-frobenius}(2), not necessarily detected by the $\sigma_j$; condition (N4) is not assumed. Then there is $C$ such that for every $m$ and all large $N$: for any $g+1$ families $(z_{i,1},\dots,z_{i,g})$, $0\le i\le g$, of Selmer classes of $\rho$ modulo $T^m$ at the augmented level, there are $c_0,\dots,c_g\in\kappa[[T]]$, one of them a unit, and a second-order Selmer class of $\rho_1$ whose first-order components are $T^C\sum_ic_iz_{i,1},\dots,T^C\sum_ic_iz_{i,g}$.
\end{lemma}
\leanref{Development/AutomorphicArithmetic/AutomorphicVarieties/PatchingModules/TaylorWilesObstructionBound.lean}{UpperTheory.V.Pa.TaylorWilesObstructionBound}
\leanref{Development/AutomorphicArithmetic/AutomorphicVarieties/PatchingModules/AutomorphicDetection/TaylorWilesObstructionBound.lean}{UpperTheory.V.Pa.b2Ar}
\begin{proof}
Fix $m$, let $N$ be large, and let $\mathcal H$ be the second cohomology group, with local conditions, of the adjoint of $\rho$ modulo $T^m$ at the augmented level; it is finite and killed by $T^m$.
\begin{enumerate}
\item \emph{Obstructions.} The families $(z_{i,1},\dots,z_{i,g})$ give obstruction classes $\xi_0,\dots,\xi_g\in\mathcal H$ such that, whenever a combination $\sum_ic_i\xi_i$ vanishes, there is a second-order Selmer class of $\rho_1$ whose first-order components are $\sum_ic_iz_{i,1},\dots,\sum_ic_iz_{i,g}$.
\item \emph{Reduction to $T^e$.} For $e\le m$, the exact sequences of coefficients $0\to\kappa[[T]]/T^{m-e}\to\kappa[[T]]/T^m\to\kappa[[T]]/T^e\to0$, given by multiplication by $T^e$ and reduction, and $0\to\kappa[[T]]/T^e\to\kappa[[T]]/T^m\to\kappa[[T]]/T^{m-e}\to0$, together with the surjectivity of global onto local second cohomology, bound $|\mathcal H/T^e\mathcal H|$ by the number of elements of the second cohomology group with local conditions modulo $T^e$ times that of a cokernel of the same level.
\item \emph{Counting modulo $T^e$.} For $e\ge0$ let $M_e$ be the module of matrices over $\kappa[[T]]/T^e$ whose two diagonal blocks vanish, on which $\Gamma_N$ acts by $g\cdot Y=\rho(g)Y\rho(g)^{-1}$ for $g\in\Gamma_{N,0}$ and $c\cdot Y=-{}^tY$. For $v\mid\ell$ let $Q_v$ be the quotient of $M_e$ by the $Y$ with $w_v^{-1}Yw_v$ upper triangular, a representation of $D_v$; for $u\in S_{\mathrm a}$ let $Q_u=M_e$. Let $\mathrm{Sel}_e(N)$ be the classes in $H^1(\Gamma_N,M_e)$ whose image in $H^1(D_v,Q_v)$ is $0$ for $v\mid\ell$ and whose restriction to $I'_u$ is $0$ for $u\in S_{\mathrm a}$; let $H^2_e(N)$ be the second cohomology group with these local conditions, the group of the step (2) at level $e$, that is the group $H^2_{\mathrm{loc}}$ of Lemma~\ref{sha-patching-obstruction-count} for $\Gamma_N$, $M_e$ and the $Q_w$, with $J_v=D_v$ for $v\mid\ell$ and $J_u=I'_u$ for $u\in S_{\mathrm a}$; and let $\mathrm{Ck}_e(N)$ be the cokernel of $H^2(\Gamma_N,M_e)\to\bigoplus_wH^2(D_w,Q_w)$. Write $\mathrm{Sel}_e(0)$ for the Selmer group at the base level, with $\Gamma$ in place of $\Gamma_N$, and $\mathrm{Hom}_\varepsilon(M_e)$ for the homomorphisms $f:M_e\to\mathbf Z/\ell$ with $f(g\cdot Y)=\bar\varepsilon(g)f(Y)$ for $g\in\Gamma_N$. There are constants $C_0$, $C_n$, $C_h$ and $C_d$ such that for every $m$, all large $N$ and all $e\le m$:
\begin{enumerate}
\item[(a)] $|H^2_e(N)|\le\#\mathrm{Sel}_e(N)\cdot\prod_{v\mid\ell}|H^0(D_v,Q_v)|\cdot\#\mathrm{Hom}_\varepsilon(M_e)$, by the global Euler characteristic formula and Tate's local formulas \cite[Chapter~VII, \S3 and Chapter~VIII, \S7]{NSW08};
\item[(b)] $\#\mathrm{Ck}_e(N)\le\#\mathrm{Hom}_\varepsilon(M_e)$, by Poitou-Tate duality \cite[Chapter~VIII, \S6]{NSW08};
\item[(c)] $\#\mathrm{Sel}_e(0)\le|\kappa|^{ge+C_0}$, since the Selmer basis has $g$ elements;
\item[(d)] $\#\mathrm{Sel}_e(N)\le\#\mathrm{Sel}_e(0)\cdot|\kappa|^{C_n}$, since $\rho$ is trivial on $P(H)$ and the eigenvalues of each $\rho(\sigma_j)$ are distinct;
\item[(e)] $|H^0(D_v,Q_v)|\le|\kappa|^{C_h}$ for $v\mid\ell$;
\item[(f)] $\#\mathrm{Hom}_\varepsilon(M_e)\le|\kappa|^{C_d}$.
\end{enumerate}
Let $A_e=\kappa[[T]]/T^e$, let $a$ and $n-a$ be the block sizes of $x$, and let $N_e=|A_e|^{a(n-a)}$.

\emph{Proof of (a).} An odd matrix is the sum of its upper right and lower left corners, each with $a(n-a)$ entries, so $|M_e|=N_e^2$. Let $\gamma\in\Gamma_N$ with $\eta(\gamma)=-1$ and $\gamma^2=1$, and write $\gamma=hc$ with $h\in\Gamma_{N,0}$. Then $\gamma$ acts by the involution $Y\mapsto\rho(h)(-{}^tY)\rho(h)^{-1}$, which carries upper right corners to lower left ones and conversely, as $\rho(h)$ is even. So sending $Y$ to its lower left corner is a bijection from $M_e^\gamma$ onto the lower left corners, with inverse $Y\mapsto Y+\gamma\cdot Y$, and $|M_e^\gamma|=N_e$. For $v\mid\ell$, since each column of $w_v$ lies in one block, $Y\mapsto w_v^{-1}Yw_v$ is a bijection from the odd matrices onto the matrices whose entries vanish at the positions $(i,j)$ for which the columns $i$ and $j$ of $w_v$ lie in the same block. So $Q_v$ is identified with the entries of $w_v^{-1}Yw_v$ at the positions $(i,j)$ with $i>j$ and with the columns $i$ and $j$ in different blocks, of which there are $a(n-a)$, and $|Q_v|=N_e$.

Apply Lemma~\ref{sha-patching-obstruction-count} to $\Gamma_N$, $M_e$ and the $Q_w$, with $J_v=D_v$ and $J_u=I'_u$; then $\mathrm{Sel}=\mathrm{Sel}_e(N)$, $H^2_{\mathrm{loc}}=H^2_e(N)$ and $\mathrm{Ck}=\mathrm{Ck}_e(N)$. As $M_e$ is killed by $\ell$, the cohomology groups are finite by Lemma~\ref{sha-descent-link-global-tate}(1) with $r=1$ and Lemma~\ref{sha-descent-link-local}. By the global Euler characteristic formula, Lemma~\ref{sha-descent-link-global-tate}(2) with $r=1$ for $\Gamma_N$, with $d=[L^+:\mathbf Q]$, no complex places, and the complex conjugations $c_1,\dots,c_d$,
\[|H^0(\Gamma_N,M_e)|\cdot|H^2(\Gamma_N,M_e)|\cdot|M_e|^d=|H^1(\Gamma_N,M_e)|\cdot\prod_{i=1}^d|M_e^{c_i}|,\]
so $|H^2(\Gamma_N,M_e)|\,N_e^d\le|H^1(\Gamma_N,M_e)|$. For $v\mid\ell$, Tate's local formula, Lemma~\ref{sha-descent-link-local}(2), gives $|H^1(D_v,Q_v)|=|H^0(D_v,Q_v)|\cdot|H^2(D_v,Q_v)|\cdot N_e^{d_v}$, with $d_v=[L^+_v:\mathbf Q_\ell]$ and $\sum_vd_v=d$. For $u\in S_{\mathrm a}$ it gives $|H^1(D_u,Q_u)|=|H^0(D_u,Q_u)|\cdot|H^2(D_u,Q_u)|$; and $|H^0(D_u,Q_u)|\le|H^1_J(D_u,Q_u)|$. Indeed, $I'_u$ acts trivially on $Q_u=M_e$, so the set $Z_u$ of cocycles of $D_u$ with values in $M_e$ that vanish on $I'_u$ has $|M_e|$ elements, by Lemma~\ref{sha-patching-tangent-counts}(3). Since $I'_u$ acts trivially, every coboundary vanishes on $I'_u$, so $Y\mapsto(g\mapsto g\cdot Y-Y)$ maps $M_e$ into $Z_u$, with kernel $H^0(D_u,Q_u)$. For the same reason a cocycle whose restriction to $I'_u$ is a coboundary vanishes on $I'_u$, so sending a cocycle to its class maps $Z_u$ onto $H^1_J(D_u,Q_u)$, with kernel the image of $M_e$. Hence $|M_e|=|Z_u|=|M_e|\,|H^0(D_u,Q_u)|^{-1}\,|H^1_J(D_u,Q_u)|$. So the hypotheses of Lemma~\ref{sha-patching-obstruction-count} hold with $s=N_e$, $r_v=|H^0(D_v,Q_v)|$, $d_u=0$ and $r_u=1$, and it gives $|H^2_e(N)|\le\#\mathrm{Sel}_e(N)\cdot\prod_v|H^0(D_v,Q_v)|\cdot\#\mathrm{Ck}_e(N)$. With (b), this is (a).

\emph{Proof of (c).} Let $\mathcal Z_e$ be the Selmer group of $x$ modulo $T^e$, the locally constant maps $z:\Gamma_0\to M_n(A_e)$ satisfying the conditions of Definition~\ref{sha-patching-selmer}(1) with odd matrices $Y_v$ over $A_e$, and let $\mathcal B_e$ be the set of maps $\delta Y$ with $Y$ odd over $A_e$ and ${}^tY=-Y$.

First, $\#\mathrm{Sel}_e(0)\cdot\#\mathcal B_e\le\#\mathcal Z_e$. Each class in $H^1(\Gamma,M_e)$ has a representative $y$ with $y(c)=0$: from $y(c^2)=0$ it follows that $c\cdot y(c)=-y(c)$, so $y+d\big(\tfrac12y(c)\big)$ vanishes at $c$, as $2$ is a unit. Choose such a representative $y_s$ of each $s\in\mathrm{Sel}_e(0)$. An element $\delta Y$ of $\mathcal B_e$ is the restriction to $\Gamma_0$ of the coboundary of $Y\in M_e$, and $c\cdot Y=Y$. The restriction $z$ of $y_s+dY$ to $\Gamma_0$ lies in $\mathcal Z_e$. Indeed, $y_s+dY$ vanishes at $c$, so $z(cgc^{-1})=c\cdot z(g)=-{}^tz(g)$; its restriction to $I'_u$ is a coboundary of the trivial module $Q_u=M_e$, hence $0$; and at $v\mid\ell$ its class in $H^1(D_v,Q_v)$ vanishes, so there is $Y_v\in M_e$ with $w_v^{-1}(z(g)-\delta Y_v(g))w_v$ upper triangular for $g\in D_v$, and its diagonal vanishes, as $w_v^{-1}Ww_v$ has zero diagonal for every odd $W$. The map $(s,\delta Y)\mapsto z$ is injective: a cocycle $y$ of $\Gamma$ vanishing at $c$ is determined by its restriction to $\Gamma_0$, since $y(gc)=y(g)$; so $z$ determines $y_s+dY$, hence $s$, and then $dY$.

Second, there is $C_0$ with $\#\mathcal Z_e\le\#\mathcal B_e\cdot|\kappa|^{ge+C_0}$ for every $e$. Apply Lemma~\ref{sha-patching-tangent-solutions} with $A=\kappa[[T]]$ and $t=T$. Let $C$ be the module of maps $f$ from the disjoint union of $\Gamma_0$ and the set $V_\ell$ of places above $\ell$ to $M_n(\kappa[[T]])$ that are continuous on $\Gamma_0$, let $C'$ be the module of families of elements of $\kappa[[T]]$ indexed by the following conditions, and let $\Phi(f)$ be the family of: the entries of each $f(y)$ in the two diagonal blocks; the entries of $f(gh)-f(g)-\rho(g)f(h)\rho(g)^{-1}$ and of $f(cgc^{-1})+{}^tf(g)$; the entries of $f(g)$ for $g\in I'_u$; and, for $v\mid\ell$ and $g\in D_v$, the entries below the diagonal of $w_v^{-1}(f(g)-\delta f(v)(g))w_v$, and its diagonal entries when $g\in\mu_v$. So the restriction to $\Gamma_0$ of an $f\in\ker\Phi$ lies in the Selmer group of $x$, with $Y_v=f(v)$. Let $B$ be the module of $f\in\ker\Phi$ whose restriction to $\Gamma_0$ is a combination of the Selmer basis plus a coboundary. The conditions on $C$ and $C'$ hold as in the proof of Lemma~\ref{sha-patching-tangent-detection}. Every $f\in\ker\Phi$ has some $T^sf$ in $B$, by the definition of a Selmer basis. And $Q_1$ is finite, as $\mathcal Z_1$ is: its elements extend to cocycles of $\Gamma$ vanishing at $c$, as in Lemma~\ref{sha-patching-tangent-symmetry}, and there are finitely many of these since $H^1(\Gamma,M_1)$ is finite. The lemma gives $e_0$ with $\#Q_e\le\#B_e\cdot|\kappa|^{e_0}$ for every $e$. Restriction to $\Gamma_0$ maps $Q_e$ onto $\mathcal Z_e$, since reduction modulo $T^e$ identifies $Q_e$ with the families of an element of $\mathcal Z_e$ and matrices $Y_v$ for it; and the index of the image of $B_e$ in $\mathcal Z_e$ is at most the index of $B_e$ in $Q_e$. The image of $B_e$ consists of reductions of maps $\sum_jc_j\kappa_j+b$ with $b$ a coboundary, so it has at most $|A_e|^g\cdot\#\mathcal B_e=|\kappa|^{ge}\cdot\#\mathcal B_e$ elements. Hence $\#\mathcal Z_e\le\#\mathcal B_e\cdot|\kappa|^{ge+e_0}$; let $C_0=e_0$.

Since $\mathcal B_e$ is not empty, the two bounds give $\#\mathrm{Sel}_e(0)\le|\kappa|^{ge+C_0}$.

\emph{Proof of (d).} Fix $j$ and let $S=\rho(\sigma_j)$, an even matrix with distinct eigenvalues. For a matrix $Y$ write $Y_{\mathrm{ev}}$ and $Y_{\mathrm{od}}$ for its even and odd parts. The $\kappa[[T]]$-linear map $\Phi(Y)=SY_{\mathrm{od}}-Y_{\mathrm{od}}S+Y_{\mathrm{ev}}$ on $M_n(\kappa[[T]])$ is injective: if $\Phi(Y)=0$, then $Y_{\mathrm{ev}}=0$, as $SY_{\mathrm{od}}-Y_{\mathrm{od}}S$ is odd; then $Y_{\mathrm{od}}$ commutes with $S$, hence, as in the proof of Lemma~\ref{sha-patching-tangent-invariants}(3), with the diagonal matrix $E$ of the first block, and so $Y_{\mathrm{od}}=0$. Let $d_j=\det\Phi\ne0$, and let $\Psi$ be the adjugate of $\Phi$, so that $\Psi\Phi=d_j$. If $Z$ is an odd matrix over $A_e$ commuting with $S$, lift it to a matrix $Y$ over $\kappa[[T]]$ with $Y_{\mathrm{ev}}=0$; then $\Phi(Y)\equiv0$ modulo $T^e$, so $d_jY=\Psi(\Phi(Y))\equiv0$, and $d_jZ=0$. Writing $d_j=uT^{s_j}$ with $u$ a unit, the entries of $Z$ are killed by $T^{s_j}$, so there are at most $|\kappa|^{C_j}$ such $Z$, with $C_j=s_jn^2$. Let $C_n=\sum_jC_j$.

Let $N\ge1$ be large enough that $\rho(\phi_{N,j})\equiv\rho(\sigma_j)$ modulo $T^m$ for every $j$, as in the step \emph{Large $N$} of the proof of Lemma~\ref{sha-patching-tangent-dual}. Let $z$ be a cocycle of $\Gamma_N$ with values in $M_e$. As in the proof of Lemma~\ref{sha-patching-tangent-descent}, the set $U_z$ of elements that act trivially on $M_e$ and at which $z$ vanishes is an open normal subgroup containing $P(H_N)$: since $\rho$ is trivial on $P(H)$, every $h\in H_N$ has a power $h^{\ell^k}$ acting trivially, and $z(h^{\ell^{k+1}})=\ell\,z(h^{\ell^k})=0$. Let $\tau=\tau_{N,j}$, $\phi=\phi_{N,j}$ and $q=\operatorname{N}w_{N,j}$; then $\tau$ acts trivially on $M_e$, and $q\equiv1\pmod\ell$. By Lemma~\ref{sha-patching-places}(2), $\tau^{-q}\phi^{-1}\tau\phi\in P(H_N)\subset U_z$, so $z(\phi^{-1}\tau\phi)=z(\tau^q)+\tau^q\cdot z(\tau^{-q}\phi^{-1}\tau\phi)=q\,z(\tau)=z(\tau)$. Since $\tau$ acts trivially, $z(\phi^{-1}\tau\phi)=\phi^{-1}\cdot z(\tau)$. So $z(\tau)$ commutes with $\rho(\phi)$, hence with $\rho(\sigma_j)$ modulo $T^e$.

The map from $\mathrm{Sel}_e(N)$ to $\prod_jM_e$ sending the class of $z$ to $(z(\tau_{N,j}))_j$ is well defined, as coboundaries vanish at the $\tau_{N,j}$, and its values lie in the product of the sets of odd matrices commuting with the $\rho(\sigma_j)$. If $z$ represents a class in its kernel, then $U_z$ contains the $\tau_{N,j}$ and $P(H_N)$, hence the kernel of $\Gamma_N\to\Gamma$ by Lemma~\ref{sha-patching-places}(3). Then, as in the proof of Lemma~\ref{sha-patching-tangent-descent}, $z$ is the inflation of a cocycle $z_0$ of $\Gamma$, whose class lies in $\mathrm{Sel}_e(0)$, since the local conditions at the $v\mid\ell$ and the $u\in S_{\mathrm a}$ are the same at both levels. So the kernel has at most $\#\mathrm{Sel}_e(0)$ elements, and $\#\mathrm{Sel}_e(N)\le\#\mathrm{Sel}_e(0)\cdot|\kappa|^{C_n}$.

\emph{Proof of (e).} By (N1), for each $v\mid\ell$ and $i\ne j$ there is $g_{v,i,j}\in D_v$ with $\psi_{v,i}(g_{v,i,j})\ne\psi_{v,j}(g_{v,i,j})$. Let $d_0\ne0$ be the product over all $v$ and all $i\ne j$ of $\psi_{v,i}(g_{v,i,j})\psi_{v,j}(g_{v,i,j})^{-1}-1$. Let $Y\in M_e$ represent an element of $H^0(D_v,Q_v)$, so that $w_v^{-1}(\rho(g)Y\rho(g)^{-1}-Y)w_v$ is upper triangular for $g\in D_v$, and let $Z=w_v^{-1}Yw_v$. We show by induction on $k$ that $d_0^kZ_{ij}=0$ whenever $i>j$ and $n+j\le k+i$; for $k=n$ this gives $d_0^nZ_{ij}=0$ for all $i>j$. Let $i>j$ with $n+j\le k+i$, let $g=g_{v,i,j}$, and let $B=w_v^{-1}\rho(g)w_v$ and $B'=w_v^{-1}\rho(g^{-1})w_v$, upper triangular with diagonal entries $\psi_{v,p}(g)$ and $\psi_{v,p}(g)^{-1}$. The $(i,j)$ entry of $BZB'$ is $Z_{ij}$, since $BZB'-Z$ is upper triangular. It is also $B_{ii}Z_{ij}B'_{jj}$ plus a combination of entries $Z_{pq}$ with $p\ge i$, $q\le j$ and $(p,q)\ne(i,j)$, each of which satisfies $n+q\le(k-1)+p$ and so is killed by $d_0^{k-1}$. So $d_0^{k-1}(B_{ii}B'_{jj}-1)Z_{ij}=0$, and $B_{ii}B'_{jj}-1=\psi_{v,i}(g)\psi_{v,j}(g)^{-1}-1$ divides $d_0$. Hence the class of $Y$ is determined by the entries $Z_{ij}$, $i>j$, which lie in the group of elements of $A_e$ killed by $d_0^n$, of order at most $|\kappa[[T]]/d_0^n|$. So $|H^0(D_v,Q_v)|\le|\kappa[[T]]/d_0^n|^{n^2}\le|\kappa|^{C_h}$ with $C_h=|\kappa[[T]]/d_0^n|^{n^2}$, independent of $e$ and $N$.

\emph{Proof of (f).} Let $E$ be the diagonal matrix with entries $1$ in the first block and $0$ in the second. It is even, so by Lemma~\ref{sha-patching-normalization} there is $s\ge0$ such that $T^sE$ lies in the $\kappa[[T]]$-span of $\rho(\Gamma')$. Let $f\in\mathrm{Hom}_\varepsilon(M_e)$. For $g\in\Gamma'$ and $Y\in M_e$, $f(\rho(g)Y)=f(Y\rho(g))$, since $\rho(g)Y$ and $Y\rho(g)$ are odd, $\rho(g)Y=g\cdot(Y\rho(g))$ for any lift of $g$ to $\Gamma_N$, and $\bar\varepsilon(g)=1$. By additivity, $f(RY)=f(YR)$ for every $R$ in the $\kappa[[T]]$-span of $\rho(\Gamma')$, in particular for $R=T^sE$. An odd $Y$ equals $EY+YE$, and $EYE=0$. So $f(T^sEY)=f(EY\cdot T^sE)=0$ and $f(T^sYE)=f(T^sE\cdot YE)=0$, and $f(T^sY)=0$ for every $Y\in M_e$. Every $Y\in M_e$ is $T^sY'$, with $Y'\in M_e$, plus an odd matrix whose entries are polynomials in $T$ of degree less than $s$ with coefficients in $\kappa$; so $f$ is determined by its values at the matrices $bT^ke_{ij}$ with $b\in\kappa$, $k<s$ and $i$ and $j$ in different blocks. There are at most $n^2s|\kappa|$ of these, and $\mathbf Z/\ell$ has at most $|\kappa|$ elements. So $\#\mathrm{Hom}_\varepsilon(M_e)\le|\kappa|^{C_d}$ with $C_d=|\kappa|n^2s$.

By (2), $|\mathcal H/T^e\mathcal H|\le|H^2_e(N)|\cdot\#\mathrm{Ck}_e(N)$. By (a) to (f) this is at most $|\kappa|^{ge+C}$ with $C=C_0+C_n+|V_\ell|C_h+2C_d$, where $V_\ell$ is the set of places above $\ell$, for every $e\le m$; and $C$ is independent of $m$ and $N$.
\item \emph{A relation.} By (3) and Lemma~\ref{sha-patching-modules}, applied to $\mathcal H$ over $\kappa[[T]]$, there are $c_0,\dots,c_g\in\kappa[[T]]$, one of them a unit, with $T^C\sum_ic_i\xi_i=0$. By (1), applied to the coefficients $T^Cc_i$, the families $T^C\sum_ic_iz_{i,1},\dots,T^C\sum_ic_iz_{i,g}$ are the first-order components of a second-order Selmer class.
\end{enumerate}
% (a) B2.odCount: odHalf (swap_count with cut for |M_e| and M_e^gamma, card_of_ker for Q_v), odKill, CohAx.ge, leV, leU and the first half of ur (B2.arUr, which is B2.urOf applied to UnrCount: the map d0 into Zu with kernel H0, and the surjection of Zu onto H1ur with kernel the range of d0), then B2.count (Lemma sha-patching-obstruction-count) and ck. (b) the field ck of CohAx. (c) B2.sel0: card_selmer_mul_card_boundary_le_card_selmerMatrix (selCard_norm, selCard_mem with selCard_diag and selCard_loc, selCard_inj), and sel_cnt (eqnL, Jx, eqn_iff, finite_truncated_locallyConstant_matrix_solutions, which applies solCnt and passes to the restriction by fc_idx, and img_card); Sol 1 finite from finite_selmerModule_of_finite_cohomology. (d) B2.selN: fix_count (odd_kill with quasi_inv and odd_comm, tors_card), ph_eq, coc_ker with PL, T.tame and coc_conj for z(tau) commuting with rho(phi), T.ker_le, descend and th_sel, sel_count. (e) B2.H0_local_card_bounded_of_cls_ker_le_ker: low_ann (d0 from (x.N1 v).1, kill_low), ann_card. (f) B2.dualBound: x.span for e1, f_symm_eq_of_left_right_mul_mat, comm_span, kill, od_dec, card_zmod_nsmul_eq_zero_le; the combination is the calc of hbound in obstructionBound_of_duality, with the private step for (2). CohAx at each level is B2.cohAxAr, transported to the augmented levels by cohAxIso. M_e is Od over x.Am e with IsOd; Q_v is D.Q (inl v), with kernel given by kerV, and J = the whole group at v; Sel_e is Dat.Sel, H^2_e is Dat.H2LocalConditions and Ck_e is Dat.Ck2. Closed: genTrick (4), Lemma sha-patching-modules.
% Pa.b2Ar applies B2.obstructionBound_of_duality with the cohomology axioms B2.cohAxAr (from CohP) at each level. Its steps: transl (1); amMaps and the private lemma step, with onto and cokerBound (2); odCount, sel0 (base level), selN (augmented level, from the triviality on PL H and Eig), H0_local_card_bounded_of_cls_ker_le_ker and dualBound (3); genTrick (4). The two unnamed hypotheses of b2Ar, the level V0.S' = s.S and the nontriviality of eps mod l, are not used.
\end{proof}

\subsection{The patching method}\label{subsection-sha-patching-method}

\begin{lemma}\label{sha-patching-chain-patched}
(Patching: Taylor and Wiles \cite{TW95}, in the form of Kisin \cite{Kis09} and Calegari and Geraghty \cite{CG18}.) Let $\Lambda$ be a complete noetherian local ring with finite residue field, $\Sigma$ a finite set, $S=\Lambda[[y_a:a\in\Sigma]]$, and for $k,j\ge0$ let $\mathfrak a_{k,j}=(y_a)_a^k+\mathfrak m_\Lambda^jS$. Let $R$ be a noetherian local ring with a homomorphism $s:S\to R$, let $f\in R$, $\mathfrak i\subset R$ an ideal, $\mathcal U$ a nonprincipal ultrafilter on $\mathbf N$, $(M_N)_{N\ge1}$ a sequence of $R$-modules, and $X$ an $R$-module that is finite and free over $\Lambda$, with $a\cdot x=s(a)x$ for $a\in\Lambda$ and $s(y_a)X=0$. Suppose that:
\begin{enumerate}
\item for all $k,j$ and $\mathcal U$-almost all $N$, there is an additive isomorphism $M_N/\mathfrak a_{k,j}M_N\cong(S/\mathfrak a_{k,j})^r$, for some $r$, compatible with the action of $S$;
\item for all $j$ and $\mathcal U$-almost all $N$, there are $R$-linear maps $\varphi:X/\mathfrak a_{1,j}X\to M_N/\mathfrak a_{1,j}M_N$ and $\psi$ in the other direction with $\varphi\psi$ and $\psi\varphi$ equal to multiplication by $f$;
\item there are $e\ge0$ and generators $G$ of $\mathfrak i$ such that for $r\in G$, all $j$ and $\mathcal U$-almost all $N$ there is $a\in R$ with $(fr)^e-a\in\mathfrak m_R^j$ and $aM_N=0$.
\end{enumerate}
Let $\mathfrak q_0,\dots,\mathfrak q_h$ be primes of $R$ containing $\mathfrak i$ and the $s(y_a)$ and not containing $f$, such that the completion of the localization of $\Lambda$ at $s^{-1}(\mathfrak q_i)\cap\Lambda$ is regular of dimension $d$. Then for each $i$ there are a completion $B_i$ of the localization of $R/\mathfrak i$ at $\mathfrak q_i$ and an ideal $\mathfrak S_i\subset B_i$ such that:
\begin{enumerate}
\item[(a)] $\dim B_i/P\ge d+|\Sigma|$ for every minimal prime $P$ of $\mathfrak S_i$;
\item[(b)] a prime $Q$ of $B_i$ containing the images of the $s(y_a)$ contains $\mathfrak S_i$ if and only if its inverse image in $R$ contains the annihilator of $X$;
\item[(c)] if $Z\subset Q$ are primes of $B_i$ with $\mathfrak S_i\subset Z$, $Q$ contains the images of the $s(y_a)$, and the inverse image of $Q$ in $R$ is contained in $\mathfrak q_{i'}$, then there is a prime $Z'\supset\mathfrak S_{i'}$ of $B_{i'}$ with the same inverse image in $R/\mathfrak i$ as $Z$.
\end{enumerate}
\end{lemma}
\leanref{Development/AutomorphicArithmetic/AutomorphicVarieties/PatchingModules/PatchedSupportCompletions.lean}{UpperTheory.V.Pa.S5}
\leanref{Development/AutomorphicArithmetic/AutomorphicVarieties/PatchingModules/PatchedSupportCompletions.lean}{UpperTheory.V.Pa.PatchedSupportCompletions}
\leanref{Development/AutomorphicArithmetic/AutomorphicVarieties/PatchingModules/AutomorphicDetection/AnnBc.lean}{UpperTheory.V.Pa.s5}

\begin{lemma}\label{sha-patching-chain-limit}
In Situation~\ref{sha-hecke-situation}, let $\sigma_1,\dots,\sigma_q$, $H$, the places $w_{N,j}$, the groups $\Gamma_N$ with their maps $p_N:\Gamma_N\to\Gamma$, and the elements $\phi_{N,j}$ and $\tau_{N,j}$ be as in Lemma~\ref{sha-patching-places}, and suppose that $P(H)=1$ and $P(H_N)=1$ for all $N$. Then there are a nonprincipal ultrafilter $\mathcal U$ on $\mathbf N$; a compact totally disconnected group $\tilde\Gamma$, topologically generated by finitely many elements, with an element $\tilde c$ and continuous surjections $\pi_N:\tilde\Gamma\to\Gamma_N$ with $\pi_N(\tilde c)=c$; and the quotient $\pi_\infty:\tilde\Gamma\to\Gamma_\infty$ of $\tilde\Gamma$ by a closed normal subgroup, with a continuous surjection $p_\infty:\Gamma_\infty\to\Gamma$ that carries $\pi_\infty(\tilde c)$ to $c$, such that the following hold. Write $\Gamma_{\infty,0}=p_\infty^{-1}(\Gamma_0)$ and $\tilde\Gamma_0=\pi_\infty^{-1}(\Gamma_{\infty,0})$, and say that $a\in\Gamma_N$ and $b\in\Gamma_\infty$ are \emph{close modulo} a subgroup $V$ of $\tilde\Gamma$ if $a=\pi_N(g)$ and $b=\pi_\infty(g')$ for some $g,g'\in\tilde\Gamma$ with $g^{-1}g'\in V$.
\begin{enumerate}
\item For all $N$ and $g\in\tilde\Gamma$, $p_N(\pi_N(g))\,p_\infty(\pi_\infty(g))^{-1}\in H$; and for every open normal subgroup $V$ of $\Gamma$ and $\mathcal U$-almost all $N$, $p_N(\pi_N(g))\,p_\infty(\pi_\infty(g))^{-1}\in V$ for all $g$.
\item For every open normal subgroup $V$ of $\tilde\Gamma$ and $\mathcal U$-almost all $N$, $(\ker\pi_N)V=(\ker\pi_\infty)V$.
\item There are continuous homomorphisms to $\Gamma_{\infty,0}$ from the groups $D_v$, $v\mid\ell$, and from the decomposition groups at the $\tilde u$, $u\in S_{\mathrm a}$, whose composites with $p_\infty$ are their maps to $\Gamma_0$, such that for every open normal subgroup $V$ of $\tilde\Gamma$ and $\mathcal U$-almost all $N$, the images of every element in $\Gamma_N$ and in $\Gamma_\infty$ are close modulo $V$.
\item There are $\phi_j,\tau_j\in\Gamma_{\infty,0}$, for $1\le j\le q$, with $\phi_j\tau_j=\tau_j\phi_j$, $p_\infty(\phi_j)=\sigma_j$ and $p_\infty(\tau_j)=1$, such that for every open normal subgroup $V$ of $\tilde\Gamma$ and $\mathcal U$-almost all $N$, $\phi_{N,j}$ and $\phi_j$ are close modulo $V$, and so are $\tau_{N,j}$ and $\tau_j$, for all $j$.
\end{enumerate}
\end{lemma}
\leanref{Development/AutomorphicArithmetic/AutomorphicVarieties/PatchingModules/ThetaCoefficientMapEq.lean}{UpperTheory.V.Pa.Lim}
\leanref{Development/AutomorphicArithmetic/AutomorphicVarieties/PatchingModules/PatchedSupportCompletions.lean}{UpperTheory.V.Pa.S4}
\leanref{Development/AutomorphicArithmetic/AutomorphicVarieties/PatchingModules/SeparableMatrixStability/SelmerCocyclesLimitAndDescentProperties.lean}{UpperTheory.V.Pa.selmer_cocycles_limit_and_descent_properties}
\begin{proof}
\emph{A covering group.} Let $\gamma_1,\dots,\gamma_m\in\Gamma$ represent the cosets of $H$. By Lemma~\ref{sha-patching-places}(4), each $H_N$ has at most $b=h\,\ell^{q[\Gamma:H]}$ continuous homomorphisms to $\mathbf Z/\ell$; since $P(H_N)=1$, $H_N$ is a pro-$\ell$ group, so it is topologically generated by $b$ elements. Let $\tilde\Gamma$ be the free profinite group on generators $x_0,x_1,\dots,x_m,x'_1,\dots,x'_b$, and $\tilde c=x_0$. For each $N$, let $\pi_N$ send $x_0$ to $c$, each $x_k$ to an element of $\Gamma_N$ over $\gamma_k$, and the $x'_k$ to topological generators of $H_N$. Its image contains $H_N$ and meets every coset of $H_N$, so $\pi_N$ is surjective; and the composite of $p_N\circ\pi_N$ with $\Gamma\to\Gamma/H$ sends $x_0$ to $cH$, $x_k$ to $\gamma_kH$ and $x'_k$ to $H$, for every $N$. So $p_N(\pi_N(g))$ and $p_{N'}(\pi_{N'}(g))$ agree modulo $H$ for all $N,N'$ and $g$.

\emph{The limit.} Let $\mathcal U$ be a nonprincipal ultrafilter on $\mathbf N$, and $K$ the set of $g\in\tilde\Gamma$ such that, for every open normal subgroup $V$ of $\tilde\Gamma$, $g\in(\ker\pi_N)V$ for $\mathcal U$-almost all $N$. It is a closed normal subgroup, and for every $V$, $(\ker\pi_N)V=KV$ for $\mathcal U$-almost all $N$. Let $\Gamma_\infty=\tilde\Gamma/K$; this gives (2). Since $\tilde\Gamma$ is topologically finitely generated, the maps $p_N\circ\pi_N$ have a limit along $\mathcal U$: a continuous surjection $F:\tilde\Gamma\to\Gamma$ that kills $K$, such that for every open normal subgroup $V$ of $\Gamma$ and $\mathcal U$-almost all $N$, $F(g)^{-1}p_N(\pi_N(g))\in V$ for all $g$. Let $p_\infty$ be the map induced by $F$; this gives the second part of (1). For the first part, given $g$ and $N$, choose $N'$ with $p_{N'}(\pi_{N'}(g))\,p_\infty(\pi_\infty(g))^{-1}\in H$, and use the previous step. As $p_N(\pi_N(\tilde c))=c$ for all $N$, the same limit gives $p_\infty(\pi_\infty(\tilde c))=c$.

\emph{Local maps and Frobenius elements.} Let $Y$ be one of the local groups of (3), with its maps $f_N:Y\to\Gamma_N$ and $f:Y\to\Gamma$, so that $p_N\circ f_N=f$. By Lemma~\ref{sha-patching-places}(5) there is a closed normal subgroup $P_Y$ of $Y$, with $Y/P_Y$ topologically finitely generated, that lies in the kernel of every continuous homomorphism from $Y$ to a profinite group carrying $f^{-1}(H)$ into a subgroup $H'$ with $P(H')=1$; in particular in the kernel of each $f_N$, as $P(H_N)=1$. As $Y/P_Y$ is topologically finitely generated, the maps $f_N$ have a limit: a continuous homomorphism $Y\to\Gamma_\infty$ such that for every $V$ and $\mathcal U$-almost all $N$ the images of every $y\in Y$ are close modulo $V$; its composite with $p_\infty$ is $f$ by (1), as $p_N\circ f_N=f$ for all $N$, so it lands in $\Gamma_{\infty,0}$. In the same way, any sequence of elements of the $\Gamma_N$ has a limit in $\Gamma_\infty$; let $\phi_j$ and $\tau_j$ be elements that are close modulo $V$ to $\phi_{N,j}$ and $\tau_{N,j}$, for every $V$ and $\mathcal U$-almost all $N$. By (1) and Lemma~\ref{sha-patching-places}(1), with $P(H)=1$, $p_\infty(\phi_j)=\sigma_j$; and by Lemma~\ref{sha-patching-places}(2), $p_\infty(\tau_j)=1$. So $\phi_j,\tau_j\in\Gamma_{\infty,0}$.

\emph{Commutation.} Let $W$ be an open normal subgroup of $\Gamma_\infty$ and $V=\pi_\infty^{-1}(W)$. Choose $N\ge[\Gamma_\infty:W]$ such that (2) holds for $V$ and $\phi_{N,j}$ and $\tau_{N,j}$ are close modulo $V$ to $\phi_j$ and $\tau_j$. Then $\ker\pi_N\subset V$. Since $P(H_N)=1$, Lemma~\ref{sha-patching-places}(2) gives $\phi_{N,j}^{-1}\tau_{N,j}\phi_{N,j}=\tau_{N,j}^{\operatorname{N}w_{N,j}}$, so $\phi_j^{-1}\tau_j\phi_j=\tau_j^{\operatorname{N}w_{N,j}}$ modulo $W$. As $\pi_N$ is a closed map and $\ker\pi_N\subset V$, some open normal subgroup $W'$ of $\Gamma_N$ has $\pi_N^{-1}(W')\subset V$; as $\tau_{N,j}$ lies in the pro-$\ell$ group $H_N$, some $\tau_{N,j}^{\ell^a}$ lies in $W'$, and so $\tau_j^{\ell^a}\in W$. Hence the image of $\tau_j$ in $\Gamma_\infty/W$ has order $\ell^{a'}$ with $a'<\ell^{a'}\le[\Gamma_\infty:W]\le N$. As $\operatorname{N}w_{N,j}\equiv1\pmod{\ell^N}$, $\tau_j^{\operatorname{N}w_{N,j}}=\tau_j$ modulo $W$, and $\phi_j$ and $\tau_j$ commute modulo $W$. As $W$ is arbitrary and $\Gamma_\infty$ is profinite, they commute.
% selmer_cocycles_limit_and_descent_properties: sys_cover (exists_free_profinite on n+1+hb generators, pl_gen for the generators of H_N from PL = bot and finN, exists_surj_over_finite), the hyperfilter, limKer and limKer_sup for K and (2), exists_lim_hom_out for F0 and hp (with sep for hct), unif and modH for (1); loc_fin and in_lim for the local maps (lift), with close_map; elt_lim and close_map for ph and ta (T.ph_lim with PL l H = bot, T.ta_ker); comm_lim (the tame relation T.tame with PL = bot, T.Nq_mod, pl_props for the l-power orders) for comm.
\end{proof}

\begin{lemma}\label{sha-patching-chain-ring}
In the setting of Lemma~\ref{sha-patching-chain-limit}, let $\tilde u_{j,i}\in\mathcal O$ for $1\le j\le q$ and $1\le i\le n$. For a complete local noetherian $\mathcal O$-algebra $R'$ with residue field $k$, a \emph{patching tuple} over $R'$ consists of: a continuous determinant $D$ of $\tilde\Gamma_0$ of dimension $n$ with values in $R'$, whose reduction modulo $\mathfrak m_{R'}$ is the composite of $\bar D$ with $p_\infty\circ\pi_\infty$; continuous characters $\psi_{v,i}:D_v\to R'^\times$, for $v\mid\ell$ and $1\le i\le n$, lifting $\bar\psi_{v,i}$ and equal to $\theta_{v,i}$ on $\mu_v$, where $D_v$ maps to $\Gamma_{\infty,0}$ as in Lemma~\ref{sha-patching-chain-limit}(3); elements $z_0$ and $z_{j,i}$ of $\mathfrak m_{R'}$; and elements $U_{j,i}\in R'$ with $U_{j,i}\equiv\tilde u_{j,i}$ modulo $\mathfrak m_{R'}$. Then there are a complete local noetherian $\mathcal O$-algebra $R$ with residue field $k$ and a patching tuple $(D_R,\psi_R,z_0,U,z)$ over $R$ from which every patching tuple over every $R'$ is obtained by a unique $\mathcal O$-algebra map $R\to R'$.
\end{lemma}
\leanref{Development/AutomorphicArithmetic/AutomorphicVarieties/PatchingModules/AnnihilatorBaseChange/SpecializationMapCompUniversalSection.lean}{UpperTheory.V.Pa.Core.Tup}
\leanref{Development/AutomorphicArithmetic/AutomorphicVarieties/PatchingModules/AnnihilatorBaseChange/SpecializationMapCompUniversalSection.lean}{UpperTheory.V.Pa.Core.RRObj}
\leanref{Development/AutomorphicArithmetic/AutomorphicVarieties/PatchingModules/AnnihilatorBaseChange/FiniteOpenHomsToFiniteGroup.lean}{UpperTheory.V.Pa.Core.rrEx}
\begin{proof}
The group $\tilde\Gamma_0$ is open in $\tilde\Gamma$, so it is topologically finitely generated, and each of its open subgroups has finitely many continuous homomorphisms to $\mathbf Z/\ell$. Since $p_\infty\circ\pi_\infty$ is surjective, the characters $\bar\chi_i\circ p_\infty\circ\pi_\infty$ are pairwise distinct, with open kernels, so the determinant of their sum has a universal deformation ring $R_0$, a complete local noetherian $\mathcal O$-algebra with residue field $k$ \cite{Che14}. By Lemma~\ref{sha-patching-places}(5), each $\bar\psi_{v,i}$ has a universal ring $R_{v,i}$ of continuous lifts equal to $\theta_{v,i}$ on $\mu_v$. The completed tensor product $R_1$ over $\mathcal O$ of $R_0$ and the $R_{v,i}$ carries a pair $(D,\psi)$ from which every pair over every $R'$ comes by a unique map $R_1\to R'$. Let $R=R_1[[Z_0,Z'_{j,i},Z_{j,i}]]$, with $z_0=Z_0$, $U_{j,i}=\tilde u_{j,i}+Z'_{j,i}$ and $z_{j,i}=Z_{j,i}$. A patching tuple over $R'$ is a pair over $R'$ together with the elements $z_0$, $U_{j,i}-\tilde u_{j,i}$ and $z_{j,i}$ of $\mathfrak m_{R'}$; since $R'$ is complete, the map $R_1\to R'$ of the pair extends uniquely to $R$ with these images of the variables.
% rrEx: pre_univ (univDet_mf, universalDeformationRing_localCharacters with T.phiV, coprod_opt from co_coprod), phi_of_fg for Ph0, then MvPowerSeries over R1 in the variables Unit + (q x n) + (q x n) with co_ev for the universal property.
\end{proof}

\begin{lemma}\label{sha-patching-chain-maps}
In the setting of Lemma~\ref{sha-patching-chain-ring}, suppose that $H$ kills the $\bar\chi_i$ and that each $\tilde u_{j,i}$ lifts $\bar\chi_i(\sigma_j)$. Let $J$, $\lambda$, $M_0$, $e$ and, for each $N$, $R_N$ with its ordinary pair $(D_N,\psi_N)$, the $U_{N,j,i}$, $s_N$, $M_N$ and the map $\theta_N:R_N\to B[u]_N$ be Taylor-Wiles data as in Proposition~\ref{sha-hecke-ideal}(7) for the elements $\phi_{N,j}$ and $\tau_{N,j}$, except that the ring of weight variables is $\Lambda=\mathcal O[[X_0,X_1,\dots,X_D]]$, with one further variable $X_0$; here $B[u]_N$ is the ring called $B[u]$ in (7), the residual datum of $R_N$ is that of $B$ composed with $p_N$, and $U_{N,j,i}\equiv\tilde u_{j,i}$ modulo $\mathfrak m_{R_N}$. Let $S=\Lambda[[y_{j,i}]]$, let $u_{j,i}=\tilde u_{j,i}+y_{j,i}$, and let $B[u]$ be the quotient of $B[[y_{j,i}]]$ by the coefficients of the polynomials $\prod_i(X-u_{j,i})-\det(X-D^{\mathrm u}(\sigma_j))$, $1\le j\le q$, where $\det(X-D^{\mathrm u}(\sigma_j))$ is the characteristic polynomial of $D^{\mathrm u}$ at $\sigma_j$. Then:
\begin{enumerate}
\item $B[u]$ is a complete local noetherian $\mathcal O$-algebra with residue field $k$, finite and free as a $B$-module, and generated as a $B$-algebra by the $u_{j,i}$; and $\det(X-D^{\mathrm u}(\sigma_j))=\prod_i(X-u_{j,i})$ in $B[u][X]$.
\item There is a unique $\mathcal O$-algebra map $s:S\to R$ with $s(X_0)=z_0$, $s(X_k)=c_k\psi_{R,v,i}(\gamma)-1$ for the $c_k$, $v$, $i$ and $\gamma$ that (7) attaches to $X_k$, and $s(y_{j,i})=z_{j,i}$.
\item For each $N$ there is an $\mathcal O$-algebra map $\rho_N:R\to R_N$ carrying the patching tuple of $R$ to $(D_N\circ\pi_N,\psi_N,s_N(X_0),(U_{N,j,i}),(s_N(y_{j,i})))$, and $\rho_N\circ s=s_N$.
\item There is an $\mathcal O$-algebra map $\pi:R\to B[u]$ carrying the patching tuple of $R$ to $(D^{\mathrm u}\circ p_\infty\circ\pi_\infty,\psi^{\mathrm u},\lambda(X_0),(u_{j,i}),0)$, and $\pi(s(a))=\lambda(a_0)$ for $a\in S$ with constant term $a_0\in\Lambda$ in the $y_{j,i}$.
\end{enumerate}
\end{lemma}
\leanref{Development/AutomorphicArithmetic/AutomorphicVarieties/PatchingModules/AnnihilatorBaseChange/SpecializationMapCompUniversalSection.lean}{UpperTheory.V.Pa.Core.RMaps}
\leanref{Development/AutomorphicArithmetic/AutomorphicVarieties/PatchingModules/AnnihilatorBaseChange/RmapsEx.lean}{UpperTheory.V.Pa.Core.rmapsEx}
\leanref{Development/AutomorphicArithmetic/AutomorphicVarieties/PatchingModules/AnnihilatorBaseChange/SpoCo.lean}{UpperTheory.V.Pa.Core.spoCo}
\leanref{Development/AutomorphicArithmetic/AutomorphicVarieties/PatchingModules/AnnihilatorBaseChange/SplittingAlgebraNontrivialFreeGenerated.lean}{UpperTheory.V.Pa.Core.spCo}
\begin{proof}
(1) The polynomial $\det(X-D^{\mathrm u}(\sigma_j))$ is monic of degree $n$, and its reduction is $\prod_i(X-\bar\chi_i(\sigma_j))=\prod_i(X-\bar u_{j,i})$. The universal splitting algebra over $B$ of these $q$ polynomials is finite and free over $B$, with roots $r_{j,i}$, and it has a maximal ideal $\mathfrak n$ at which $r_{j,i}$ reduces to $\bar u_{j,i}$. As $B$ is complete, its localization at $\mathfrak n$ is a direct factor, so it is finite and free over $B$, and complete local with residue field $k$. Sending $y_{j,i}$ to $r_{j,i}-\tilde u_{j,i}$, and $r_{j,i}$ to $u_{j,i}$, gives mutually inverse maps between $B[u]$ and this localization; the second uses that $\prod_i(X-u_{j,i})$ is the characteristic polynomial of $D^{\mathrm u}$ at $\sigma_j$ in $B[u]$, by the definition of $B[u]$, and that the maximal ideal of $B[u]$ lies over $\mathfrak n$. Under this identification $u_{j,i}$ is $r_{j,i}$, which gives the rest of (1).

(2) $S$ is a ring of power series over $\mathcal O$ in the $X_k$ and the $y_{j,i}$, and the prescribed images lie in $\mathfrak m_R$: $c_k\psi_{R,v,i}(\gamma)-1$ reduces to $\bar c_k\bar\psi_{v,i}(\gamma)-1=0$, by the choice of $c_k$ in (7).

(3) The tuple over $R_N$ is a patching tuple. Indeed, $\pi_N$ maps $\tilde\Gamma_0$ into $\Gamma_{N,0}$ by Lemma~\ref{sha-patching-chain-limit}(1), as $H\subset\Gamma_0$; the reduction of $D_N\circ\pi_N$ is the determinant of $\bigoplus_i\bar\chi_i\circ p_N\circ\pi_N$; and $\bar\chi_i(p_N(\pi_N(g)))=\bar\chi_i(p_\infty(\pi_\infty(g)))$ by Lemma~\ref{sha-patching-chain-limit}(1), since $H$ kills $\bar\chi_i$. The $\psi_{N,v,i}$ lift the $\bar\psi_{v,i}$ and equal $\theta_{v,i}$ on $\mu_v$, and $s_N(X_0)$ and $s_N(y_{j,i})$ lie in $\mathfrak m_{R_N}$. So $\rho_N$ exists by Lemma~\ref{sha-patching-chain-ring}. The maps $\rho_N\circ s$ and $s_N$ agree on $X_0$ and the $y_{j,i}$, and on $X_k$ since $\rho_N(c_k\psi_{R,v,i}(\gamma)-1)=c_k\psi_{N,v,i}(\gamma)-1=s_N(X_k)$ by (7); so they are equal, by uniqueness as in (2).

(4) In the same way, the tuple over $B[u]$ is a patching tuple, as $\lambda(X_0)\in\mathfrak m_B$ and $u_{j,i}-\tilde u_{j,i}=y_{j,i}$; so $\pi$ exists. The maps $\pi\circ s$ and $a\mapsto\lambda(a_0)$ agree on $X_0$, on the $X_k$ by (7), and on the $y_{j,i}$, which both send to $0$; so they are equal.
% rmapsEx: sinfExt for s, tupN (with Adapted and modH for the residual determinant, compat for psi) and W.univ for rho_N, tupPi and W.univ for pi, rm_ext for rho_N o s = s_N and pi o s. spoCo: spCo with chiRes and A.ut_res; spCo: splitting_algebra_nontrivial_free_generated (mSplitAlg_basis, local_factor, spCo_desc) and BC.sp_fin; spRoots for the factorization.
\end{proof}

\begin{lemma}\label{sha-patching-chain-relations}
In the setting of Lemma~\ref{sha-patching-chain-maps}, let $\operatorname{tr}$ be the trace of $D_R$, extended linearly to $R[\tilde\Gamma_0]$, let $\varepsilon$ act on $\tilde\Gamma_0$ through $p_\infty\circ\pi_\infty$, and for $g\in\tilde\Gamma_0$ let $g^*=\varepsilon(g)^{1-n}(\tilde cg\tilde c^{-1})^{-1}$. Let $\mathfrak i\subset R$ be the ideal generated by the following elements, for $h\in\tilde\Gamma_0$:
\begin{enumerate}
\item[(i)] $\operatorname{tr}(h(\delta-1))$, for $\delta\in\tilde\Gamma_0$ with $\pi_\infty(\delta)=1$;
\item[(ii)] the differences of the coefficients of the characteristic polynomials of $D_R$ at $g^*$ and at $g$, for $g\in\tilde\Gamma_0$;
\item[(iii)] $\operatorname{tr}\big(h\prod_i(g_i-\psi_{R,v,i}(g'_i))\big)$, for $v\mid\ell$, $g'_1,\dots,g'_n\in D_v$ and $g_i\in\tilde\Gamma_0$ with $\pi_\infty(g_i)$ the image of $g'_i$;
\item[(iv)] $\operatorname{tr}(h(g-1))$, for $u\in S_{\mathrm a}$ and $g\in\tilde\Gamma_0$ with $\pi_\infty(g)$ the image of an element of $I_u\cap\bigcap_i\ker\bar\chi_i$;
\item[(v)] $\operatorname{tr}\big(h\prod_i(\phi-U_{j,i})\big)$, for $\phi\in\tilde\Gamma_0$ with $\pi_\infty(\phi)=\phi_j$;
\item[(vi)] $\operatorname{tr}\big(h\prod_{k\ne i}(\phi-U_{j,k})\,(\tau-1-z_{j,i})\big)$, for $\phi,\tau\in\tilde\Gamma_0$ with $\pi_\infty(\phi)=\phi_j$ and $\pi_\infty(\tau)=\tau_j$.
\end{enumerate}
Let $f=\prod_j\prod_i\prod_{i'\ne i}(U_{j,i}-U_{j,i'})^2\in R$. Then:
\begin{enumerate}
\item $\pi(\mathfrak i)=0$; and if $B$ is topologically generated over $\mathcal O$ by the coefficients of $D^{\mathrm u}$, the values of $\psi^{\mathrm u}$ and $\lambda(X_0)$, then $\pi$ is surjective.
\item There is an ordinary pair $(D_\infty,\psi_\infty)$ over $R/\mathfrak i$ for $\Gamma_\infty$, with the residual datum of $B$ composed with $p_\infty$, such that $D_\infty\circ\pi_\infty$ and $\psi_\infty$ are $D_R$ and $\psi_R$ modulo $\mathfrak i$; and the map $R/\mathfrak i\to B[u]$ induced by $\pi$ carries it to $(D^{\mathrm u}\circ p_\infty,\psi^{\mathrm u})$.
\end{enumerate}
\end{lemma}
\leanref{Development/AutomorphicArithmetic/AutomorphicVarieties/PatchingModules/AnnihilatorBaseChange/SpecializationMapCompUniversalSection.lean}{UpperTheory.V.Pa.Core.Rel}
\leanref{Development/AutomorphicArithmetic/AutomorphicVarieties/PatchingModules/AnnihilatorBaseChange/SpecializationMapCompUniversalSection.lean}{UpperTheory.V.Pa.Core.ii}
\leanref{Development/AutomorphicArithmetic/AutomorphicVarieties/PatchingModules/AnnihilatorBaseChange/SpecializationMapCompUniversalSection.lean}{UpperTheory.V.Pa.Core.fo}
\leanref{Development/AutomorphicArithmetic/AutomorphicVarieties/PatchingModules/AnnihilatorBaseChange/CharpolyElementEqAevalGroupElement.lean}{UpperTheory.V.Pa.Core.piKer}
\leanref{Development/AutomorphicArithmetic/AutomorphicVarieties/PatchingModules/AnnihilatorBaseChange/PiSurj.lean}{UpperTheory.V.Pa.Core.piSurj}
\leanref{Development/AutomorphicArithmetic/AutomorphicVarieties/PatchingModules/AnnihilatorBaseChange/OoPair.lean}{UpperTheory.V.Pa.Core.ooPair}
\leanref{Development/AutomorphicArithmetic/AutomorphicVarieties/PatchingModules/AnnihilatorBaseChange/OoMaps.lean}{UpperTheory.V.Pa.Core.ooMaps}
\begin{proof}
(1) By Lemma~\ref{sha-patching-chain-maps}(4), $\pi(\operatorname{tr}(x))$ is the trace of $D^{\mathrm u}$ at the image of $x$ in $B[u][\Gamma_0]$ under $p_\infty\circ\pi_\infty$, and $\pi$ carries $\psi_R$, $U_{j,i}$ and $z_{j,i}$ to $\psi^{\mathrm u}$, $u_{j,i}$ and $0$. So $\pi$ carries the element (i) to a trace of $h'(1-1)=0$; (ii) to the differences of the coefficients of the characteristic polynomials of $D^{\mathrm u}$ at $g'^*$ and $g'$, $g'=p_\infty(\pi_\infty(g))$, which vanish by (1) of Situation~\ref{sha-hecke-situation}, as $p_\infty(\pi_\infty(\tilde c))=c$; (iii) and (iv) to the relations (2) and (3) of Situation~\ref{sha-hecke-situation}, which vanish; and (vi) to $0$, since its last factor maps to $1-1-0$, as $p_\infty(\tau_j)=1$. For (v), $p_\infty(\phi_j)=\sigma_j$, so the image is the trace of $h'\prod_i(\sigma_j-u_{j,i})$, which by Lemma~\ref{sha-patching-chain-maps}(1) is the characteristic polynomial of $D^{\mathrm u}$ at $\sigma_j$ evaluated at $\sigma_j$. This element lies in the kernel of $D^{\mathrm u}$ by the Cayley-Hamilton theorem for determinants, so the trace of its product with $h'$ vanishes \cite{Che14}.

For the surjectivity, the image of $\pi$ contains the coefficients of $D^{\mathrm u}$ at the elements of $\Gamma_0$, as $p_\infty\circ\pi_\infty$ is surjective, the values of $\psi^{\mathrm u}$, $\lambda(X_0)=\pi(s(X_0))$ and the $u_{j,i}=\pi(U_{j,i})$. By the generation hypothesis, every element of $B$ is congruent modulo any power of $\mathfrak m_B$ to an element of the $\mathcal O$-algebra generated by the first three, hence modulo the same power of $\mathfrak m_{B[u]}$ to an element of the image; and the elements of $B[u]$ with this property form a $B$-subalgebra, which contains the $u_{j,i}$ and so is all of $B[u]$. In particular $\mathfrak m_{B[u]}\subset\pi(\mathfrak m_R)+\mathfrak m_{B[u]}^2$, and as $R$ and $B[u]$ are complete, $\pi$ is surjective.

(2) The ring $R/\mathfrak i$ is complete local noetherian with residue field $k$. Let $\delta\in\tilde\Gamma_0$ with $\pi_\infty(\delta)=1$. By (i), the trace of $D_R\bmod\mathfrak i$ at $g(\delta-1)$ vanishes for all $g\in\tilde\Gamma_0$. The reduction of $D_R$ is the determinant of a sum of pairwise distinct characters, so it is multiplicity free, and an element whose products with all $g$ have trace zero lies in the kernel of $D_R\bmod\mathfrak i$ \cite{Che14}. Hence $D_R\bmod\mathfrak i$ is the composite with $\pi_\infty$ of a determinant $D_\infty$ of $\Gamma_{\infty,0}$, which is continuous since $\pi_\infty$ is a closed map. By (ii), (iii) and (iv), $D_\infty$ and $\psi_R\bmod\mathfrak i$ satisfy (1), (2) and (3) of an ordinary pair, for the residual datum of $B$ composed with $p_\infty$. Its image under $R/\mathfrak i\to B[u]$ is $(D^{\mathrm u}\circ p_\infty,\psi^{\mathrm u})$ by Lemma~\ref{sha-patching-chain-maps}(4), as $\pi_\infty$ is surjective.
% piKer: lam_push with pi_maps, then the six cases (ker; pol with u.P and Lambda.ct_o; ord with u.Ord; ine with u.I and Adapted; q1 with charpoly_element_eq_aeval_group_element, spRoots, ch_ker and tr_ker; q2 with p_ta). piSurj: surj_of_gen (the subalgebra P of elements approximable modulo every power) and co_surj. ooPair: ooDet (oo_kills from g_kills_tr for multiplicity-free residual determinants, gdet_desc, continuity through the closed map), resPull, and the relations pol, ord, ine; ooMaps from pi_maps and comap_inj.
\end{proof}

\begin{lemma}\label{sha-patching-chain-hypotheses}
In the setting of Lemma~\ref{sha-patching-chain-relations}, let $R$ act on each $M_N$ through $\rho_N$, and on $X=B[u]\otimes_BM_0$ through $\pi$, and let $\mathfrak a_{k,j}=(y_{j,i})^k+\mathfrak m_\Lambda^jS$. Then $X$ is finite and free over $\Lambda$, with $a\cdot x=s(a)x$ for $a\in\Lambda$ and $s(y_{j,i})X=0$, and:
\begin{enumerate}
\item for all $k,j$ and $\mathcal U$-almost all $N$, $M_N/\mathfrak a_{k,j}M_N\cong(S/\mathfrak a_{k,j})^r$ for some $r$, compatibly with the action of $S$ through $s$;
\item for all $j$ and $\mathcal U$-almost all $N$, there are $R$-linear maps $\varphi:X/\mathfrak a_{1,j}X\to M_N/\mathfrak a_{1,j}M_N$ and $\psi$ in the other direction with $\varphi\psi$ and $\psi\varphi$ equal to multiplication by $f$;
\item for each generator $r$ of $\mathfrak i$ listed in Lemma~\ref{sha-patching-chain-relations}, all $j$ and $\mathcal U$-almost all $N$, there is $a\in R$ with $(fr)^{e+1}-a\in\mathfrak m_R^j$ and $aM_N=0$.
\end{enumerate}
\end{lemma}
\leanref{Development/AutomorphicArithmetic/AutomorphicVarieties/PatchingModules/AnnihilatorBaseChange/PatchedModuleFreeFiniteOverIwasawaAlgebra.lean}{UpperTheory.V.Pa.Core.patchedModule_free_finite_over_iwasawaAlgebra}
\leanref{Development/AutomorphicArithmetic/AutomorphicVarieties/PatchingModules/AnnihilatorBaseChange/FreeNQuot.lean}{UpperTheory.V.Pa.Core.freeN}
\leanref{Development/AutomorphicArithmetic/AutomorphicVarieties/PatchingModules/AnnihilatorBaseChange/EventuallyExistsPatchedModuleQuotientMaps.lean}{UpperTheory.V.Pa.Core.eventually_exists_patchedModule_quotient_maps}
\leanref{Development/AutomorphicArithmetic/AutomorphicVarieties/PatchingModules/AnnihilatorBaseChange/T3Of.lean}{UpperTheory.V.Pa.Core.t3Of}
\begin{proof}
\emph{The module $X$.} By Lemma~\ref{sha-patching-chain-maps}(1), $B[u]\cong B^m$ as a $B$-module, so $X\cong M_0^m$, which is finite and free over $\Lambda$, with $a\in\Lambda$ acting by $\lambda(a)$. By Lemma~\ref{sha-patching-chain-maps}(4), $\lambda(a)=\pi(s(a))$ and $\pi(s(y_{j,i}))=0$.

\emph{(1).} Let $N\ge j+k$. By Proposition~\ref{sha-hecke-ideal}(7), $M_N\cong(S/I_N)^r$ compatibly with $S$ acting through $s_N=\rho_N\circ s$, where $I_N$ is generated by the $(1+y_{j',i})^{\ell^{a_{j'}}}-1$ with $a_{j'}\ge N$. Expanding by the binomial theorem, the terms of $(1+y)^{\ell^a}-1$ of degree at least $k$ lie in $(y)^k$, and those of degree $d<k\le\ell^{a-j}$ have coefficient $\binom{\ell^a}d$, which is divisible by $\ell^j$. As $\ell\in\mathfrak m_\Lambda$, $I_N\subset\mathfrak a_{k,j}$, and $M_N/\mathfrak a_{k,j}M_N\cong(S/\mathfrak a_{k,j})^r$.

\emph{(2).} Fix $j$. Since $M_0$ is finite over $\Lambda$ and the residue field of $\Lambda$ is finite, $M_0/\mathfrak m_\Lambda^jM_0$ is finite, so $\mathfrak m_B^bM_0\subset\lambda(\mathfrak m_\Lambda^j)M_0$ for some $b$. By the continuity of $D^{\mathrm u}$, there is an open normal subgroup $V$ of $\Gamma$ such that the coefficients of the characteristic polynomials of $D^{\mathrm u}$ at elements of $\Gamma_0$ that are congruent modulo $V$ agree modulo $\mathfrak m_B^b$. By Lemma~\ref{sha-patching-places}(1), with $P(H)=1$, and Lemma~\ref{sha-patching-chain-limit}(1), for $\mathcal U$-almost all $N$ we have $N\ge j$, $p_N(\phi_{N,j'})\in\sigma_{j'}V$ for all $j'$, and $p_N(\pi_N(g))\equiv p_\infty(\pi_\infty(g))$ modulo $V$ for all $g$. Fix such an $N$. Then the characteristic polynomials of $D^{\mathrm u}$ at $\sigma_{j'}$ and at $p_N(\phi_{N,j'})$ agree modulo $\mathfrak m_B^b$, so $B[u]$ and $B[u]_N$ have the same quotient $C$ by $\mathfrak m_B^b$, both being quotients of $B[[y_{j,i}]]$. The two maps $R\to C$ obtained from $\pi$ and from $\theta_N\circ\rho_N$ carry the patching tuple of $R$ to the same tuple over $C$: the determinants are $D^{\mathrm u}\circ p_\infty\circ\pi_\infty$ and $D^{\mathrm u}\circ p_N\circ\pi_N$, which agree modulo $\mathfrak m_B^b$ by the choice of $V$; the characters are $\psi^{\mathrm u}$; and the elements are $\lambda(X_0)$, $u_{j,i}$ and $0$ in both cases, by Lemma~\ref{sha-patching-chain-maps} and (7). By the uniqueness in Lemma~\ref{sha-patching-chain-ring}, the two maps agree. Next, let $f'_N=\prod_j\prod_i\prod_{i'\ne i}(U_{j,i}-U_{j,i'})(U_{j,i}-\operatorname{N}w_{N,j}\,U_{j,i'})\in R$, so that $\rho_N(f'_N)$ is the element $f_N$ of (7). Since $\operatorname{N}w_{N,j}\equiv1\pmod{\ell^N}$ and $N\ge j$, $f-f'_N\in\ell^jR$, and $\ell^j\in s(\mathfrak a_{1,j})$. Now $X/\mathfrak a_{1,j}X$ and $(B[u]_N\otimes_BM_0)/\mathfrak a_{1,j}$ are both $C\otimes_BM_0$ modulo the image of $\mathfrak a_{1,j}$, compatibly with the actions of $R$ through $\pi$ and through $\theta_N\circ\rho_N$; and $(y_{j,i})\subset\mathfrak a_{1,j}$. So the maps of (7) between $M_N/(y_{j,i})M_N$ and $B[u]_N\otimes_BM_0$, which are compatible with $\theta_N$ and have composites $f_N$, induce $R$-linear maps $\varphi$ and $\psi$ between $X/\mathfrak a_{1,j}X$ and $M_N/\mathfrak a_{1,j}M_N$ whose composites are multiplication by $f'_N$, that is by $f$.

\emph{(3).} We first show: for each generator $r$ and each $j$, for $\mathcal U$-almost all $N$ there is $r'\in R$ with $r-r'\in\mathfrak m_R^j$ and either $\rho_N(r')=0$ or $(f_N\rho_N(r'))^eM_N=0$. By the continuity of $D_R$, there is an open normal subgroup $V\subset\tilde\Gamma_0$ of $\tilde\Gamma$ such that the coefficients of the characteristic polynomials of $D_R$ at elements of $R[\tilde\Gamma_0]$ with the same image in $R[\tilde\Gamma_0/V]$ agree modulo $\mathfrak m_R^j$. Take $N$ such that (2), (3) and (4) of Lemma~\ref{sha-patching-chain-limit} hold for $V$, and (1) holds for the kernel of $\varepsilon$ modulo $\mathfrak m_R^j$, an open normal subgroup of $\Gamma$. By (2) of that lemma, if $\gamma_N\in\Gamma_N$ and $\gamma_\infty\in\Gamma_\infty$ are close modulo $V$ and $\pi_\infty(g)=\gamma_\infty$, then there is $g'$ with $\pi_N(g')=\gamma_N$ and $g^{-1}g'\in V$; and every $\delta$ with $\pi_\infty(\delta)=1$ is congruent modulo $V$ to some $\delta'$ with $\pi_N(\delta')=1$. Replace accordingly each group element in $r$ whose image under $\pi_\infty$ is prescribed: $\delta$ in (i) by $\delta'$ with $\pi_N(\delta')=1$; $g_i$ in (iii) and $g$ in (iv) by elements whose images under $\pi_N$ are the images of the same local elements in $\Gamma_N$; $\phi$ and $\tau$ in (v) and (vi) by elements over $\phi_{N,j}$ and $\tau_{N,j}$; and in (ii) replace $\varepsilon(g)$, computed through $p_\infty\circ\pi_\infty$, by $\varepsilon(p_N(\pi_N(g)))$, which agrees with it modulo $\mathfrak m_R^j$. Let $r'$ be the result; then $r-r'\in\mathfrak m_R^j$. By Lemma~\ref{sha-patching-chain-maps}(3), $\rho_N(r')$ is the same expression for the pair $(D_N,\psi_N)$ over $R_N$, with $\phi_{N,j}$, $\tau_{N,j}$, $U_{N,j,i}$ and $s_N(y_{j,i})$. For (i) it is a trace of $h'(1-1)$, which is $0$. For (ii), (iii) and (iv) it is one of the relations (1), (2) and (3) of the ordinary pair over $R_N$, so it vanishes; in (iv) the element lies in $I_u\cap\bigcap_i\ker(\bar\chi_i\circ p_N)$, as the residual datum of $R_N$ is that of $B$ composed with $p_N$. For (v) and (vi), $(f_N\rho_N(r'))^e$ kills $M_N$ by (7).

Now let $N$ also satisfy $N\ge j$, and let $f'_N$ be as in (2), so that $f-f'_N\in\mathfrak m_R^j$ and $\rho_N(f'_N)=f_N$. Then $fr-f'_Nr'=(f-f'_N)r+f'_N(r-r')\in\mathfrak m_R^j$, so $a=(f'_Nr')^{e+1}$ has $(fr)^{e+1}-a\in\mathfrak m_R^j$; and $a$ acts on $M_N$ as $(f_N\rho_N(r'))^{e+1}$, which is $0$.
% Module X: tens_free with spoCo. (1): freeN with freeN_le (one_add_pow_sub_one_mem, prime_pow_dvd_choose) and freeN_quot. (2): power_maximal_ideal_le_map_power_smul_top, piCong (sees_glob, T.ph_lim with PL l H = bot, p_lim, coeff_cong, sp_le, sp_lift, rrExt), spCong, foCong, quotient_equiv_of_compatible_ring_homs with the maps phi and psi of TWMaps. (3): t3Of with foCong, and the six cases patched_trace_single_mul_sub_one_approx_annihilated (ker), patched_lam_conj_twist_approx_annihilated (pol, ker_open for eps), patched_trace_local_charPoly_approx_annihilated (ord), eventually_trace_inertia_congruence_of_adapted (ine), t3Q1 and t3Q2 (rel1 and rel2 of TWMaps); each uses phSees (dSees) and closeLift.
\end{proof}

\begin{lemma}\label{sha-patching-chain-points}
In the setting of Lemma~\ref{sha-patching-chain-relations}, let $x_0,\dots,x_h$ be arcs, with representations $\rho_t$ over $\kappa_t[[T]]$, attached to $\mathcal O$-algebra maps $f_t:B\to\kappa_t[[T]]$, such that for each $t$ and $j$ the characteristic polynomial of $\rho_t(\sigma_j)$ is $\prod_i(X-\alpha_{t,j,i})$ with pairwise distinct $\alpha_{t,j,1},\dots,\alpha_{t,j,n}\in\kappa_t[[T]]$; and let $e_1,\dots,e_h$ be primes of $B$ with $e_t\subset\ker f_{t-1}\cap\ker f_t$. Then:
\begin{enumerate}
\item there are primes $P_0,\dots,P_h$ and $E_1,\dots,E_h$ of $B[u]$ with $P_t\cap B=\ker f_t$, $E_t\cap B=e_t$ and $E_t\subset P_{t-1}\cap P_t$; and for every prime $P$ of $B[u]$ and prime $\mathfrak q\subset P\cap B$ of $B$ there is a prime $Q\subset P$ of $B[u]$ with $Q\cap B=\mathfrak q$;
\item each $f_t$ extends to an $\mathcal O$-algebra map $f'_t:B[u]\to\kappa_t[[T]]$ with kernel $P_t$;
\item $f'_t(\pi(f))\ne0$, so the prime $\mathfrak q_t=\pi^{-1}(P_t)$ of $R$ contains $\mathfrak i$ and the $s(y_{j,i})$ and does not contain $f$.
\end{enumerate}
\end{lemma}
\leanref{Development/AutomorphicArithmetic/AutomorphicVarieties/PatchingModules/PatchingBaseChange/ChainLift.lean}{UpperTheory.V.Pa.BC.chain_lift}
\leanref{Development/AutomorphicArithmetic/AutomorphicVarieties/AlgebraTools/CaOverLift.lean}{UpperTheory.V.AlgebraTools.exists_lift_ker_eq}
\leanref{Development/AutomorphicArithmetic/AutomorphicVarieties/PatchingModules/AnnihilatorBaseChange/InjOfProd.lean}{UpperTheory.V.Pa.Core.ooPt}
\begin{proof}
(1) By Lemma~\ref{sha-patching-chain-maps}(1), $B[u]$ is finite and free over $B$, so it is integral and flat over $B$, and $B\to B[u]$ is injective. By lying over, there is a prime $P_0$ over $\ker f_0$. By going down for the flat extension, there is $E_1\subset P_0$ over $e_1$, and by going up for the integral extension, there is $P_1\supset E_1$ over $\ker f_1$; continuing gives the chain \cite[Theorems~9.4 and~9.5]{Mat86}. The last statement is going down.

(2) As $P_t\cap B=\ker f_t$, the fraction field $K_0$ of $B/\ker f_t$ embeds into the fraction field $K_P$ of $B[u]/P_t$, and through $f_t$ into the fraction field of $\kappa_t[[T]]$. The field $K_P$ is generated over $K_0$ by the images of the $u_{j,i}$, which are roots of the polynomials $\det(X-D^{\mathrm u}(\sigma_j))$ by Lemma~\ref{sha-patching-chain-maps}(1); their images under $f_t$ are the characteristic polynomials of the $\rho_t(\sigma_j)$, which split over $\kappa_t[[T]]$. So the embedding of $K_0$ extends to $K_P$. The composite of $B[u]\to K_P$ with this extension has kernel $P_t$ and extends $f_t$, and it carries each $u_{j,i}$ to a root of a polynomial that splits over $\kappa_t[[T]]$, which lies in $\kappa_t[[T]]$. As the $u_{j,i}$ generate $B[u]$ over $B$, it takes values in $\kappa_t[[T]]$; this is $f'_t$.

(3) In $\kappa_t[[T]][X]$, $\prod_i(X-f'_t(u_{j,i}))$ is the image of $\det(X-D^{\mathrm u}(\sigma_j))$, by Lemma~\ref{sha-patching-chain-maps}(1), that is $\prod_i(X-\alpha_{t,j,i})$. As $\kappa_t[[T]]$ is a domain, the $f'_t(u_{j,i})$ are the $\alpha_{t,j,i}$ in some order, so they are pairwise distinct, and $f'_t(\pi(f))=\prod_j\prod_i\prod_{i'\ne i}(f'_t(u_{j,i})-f'_t(u_{j,i'}))^2\ne0$. By Lemma~\ref{sha-patching-chain-relations}(1), $\pi(\mathfrak i)=0$, and $\pi(s(y_{j,i}))=0$ by Lemma~\ref{sha-patching-chain-maps}(4).
% BC.chain_lift (exists_ideal_over_prime_of_isIntegral, chain_lift_down by exists_ideal_le_liesOver_of_le for flat extensions, prime_chain_lift_extension). lift_pt applies exists_lift_ker_eq (exists_field_lift into the fraction field of A) with spRoots, chiRes and x.eig. ooPt: heig and inj_of_prod for f'_t(pi f) != 0; it also builds the OoPt used in Step 4.
\end{proof}

\begin{lemma}\label{sha-patching-chain-completion}
Let $\mathcal O$ be a complete discrete valuation ring with uniformizer $\varpi$ and residue field $\kappa$, let $A$ be a complete local noetherian $\mathcal O$-algebra with residue field $\kappa$, let $f:A\to\kappa[[T]]$ be a homomorphism of $\mathcal O$-algebras, and let $z_0\in A$ with $f(z_0)=T$. Let $\hat A$ be the completion of the localization of $A$ at the prime $\mathfrak p=\ker f$, and let $\mathcal O'$ be the completion of the localization of $\mathcal O[[X]]$ at the prime $\varpi\mathcal O[[X]]$. Then:
\begin{enumerate}
\item there is a unique homomorphism $\hat f:\hat A\to\kappa((T))$ whose composite with $A\to\hat A$ is $f$, and $\ker\hat f=\mathfrak m_{\hat A}$;
\item there is a unique continuous homomorphism of $\mathcal O$-algebras $F:\mathcal O[[X]]\to A$ with $F(X)=z_0$, and $f\circ F$ is the surjection $\mathcal O[[X]]\to\kappa[[T]]$ that reduces the coefficients modulo $\varpi$ and sends $X$ to $T$;
\item $\mathcal O'$ is a complete discrete valuation ring with uniformizer $\varpi$, and there is a homomorphism $\mathcal O'\to\hat A$, whose composite with $\mathcal O[[X]]\to\mathcal O'$ is the composite of $F$ with $A\to\hat A$, that makes $\hat A$ a complete local noetherian $\mathcal O'$-algebra with the residue field of $\mathcal O'$;
\item the composite of $\mathcal O'\to\hat A$ with $\hat f$ is surjective.
\end{enumerate}
\end{lemma}
\leanref{Development/AutomorphicArithmetic/AutomorphicVarieties/PatchingModules/SelmerScaling/RfEx.lean}{UpperTheory.V.Pa.rfEx}
\leanref{Development/AutomorphicArithmetic/AutomorphicVarieties/PatchingModules/SelmerScaling/OoExKer.lean}{UpperTheory.V.Pa.ooEx}
\begin{proof}
Two properties of $\hat A$ are used \cite[\S8]{Mat86}: its maximal ideal is $\mathfrak p\hat A$; and a homomorphism $\varphi$ from $A$ to a complete local ring $C$ with $\varphi^{-1}(\mathfrak m_C)=\mathfrak p$ extends uniquely to $\hat A$, since it carries $A\setminus\mathfrak p$ to units and $\mathfrak p^k$ into $\mathfrak m_C^k$. The same holds for $\mathcal O'$ and the prime $\varpi\mathcal O[[X]]$.

(1) The field $\kappa((T))$ is a complete local ring with maximal ideal $0$, and the inverse image of $0$ under $A\to\kappa[[T]]\to\kappa((T))$ is $\mathfrak p$. So $f$ extends uniquely to $\hat f$. As $\hat f$ carries units to nonzero elements, $\ker\hat f\subset\mathfrak m_{\hat A}$; and $\mathfrak m_{\hat A}=\mathfrak p\hat A$ lies in $\ker\hat f$, as $\hat f$ kills $\mathfrak p$.

(2) Since $T$ is not a unit, $z_0\in\mathfrak m_A$, so $F$ exists and is unique, as $A$ is complete. The map $f\circ F$ and the reduction map are continuous homomorphisms of $\mathcal O$-algebras to $\kappa[[T]]$, a complete local $\mathcal O$-algebra, and both send $X$ to $T$; so they agree. The reduction map is surjective, and its kernel is $\varpi\mathcal O[[X]]$, since a power series reduces to $0$ if and only if all its coefficients lie in $\varpi\mathcal O$.

(3) The ring $\mathcal O[[X]]$ is a noetherian domain and $\varpi$ is a prime element of it. So the maximal ideal of $\mathcal O'$ is generated by $\varpi$, and no power $\varpi^k$ vanishes in $\mathcal O'$, as $s\varpi^k\notin\varpi^{k+1}\mathcal O[[X]]$ for $s\notin\varpi\mathcal O[[X]]$. By the Krull intersection theorem every nonzero element of $\mathcal O'$ is a power of $\varpi$ times a unit, so $\mathcal O'$ is a discrete valuation ring with uniformizer $\varpi$. By (2), the inverse image of $\mathfrak m_{\hat A}$ under the composite of $F$ with $A\to\hat A$ is $F^{-1}(\mathfrak p)=\varpi\mathcal O[[X]]$, so this composite extends to a local homomorphism $\mathcal O'\to\hat A$. Every element of $\hat A$ is congruent modulo $\mathfrak m_{\hat A}$ to $a/s$ for some $a\in A$ and $s\in A\setminus\mathfrak p$. By (2) there are $c_a,c_s\in\mathcal O[[X]]$ with $a-F(c_a)$ and $s-F(c_s)$ in $\mathfrak p$; then $c_s\notin\varpi\mathcal O[[X]]$, so $c_s$ is a unit of $\mathcal O'$, and $a/s$ is congruent to the image of $c_a/c_s$. So $\mathcal O'$ maps onto the residue field of $\hat A$.

(4) Let $u=x/y\in\kappa((T))$ with $x,y\in\kappa[[T]]$ and $y\ne0$. By (2) there are $c_x,c_y\in\mathcal O[[X]]$ with $f(F(c_x))=x$ and $f(F(c_y))=y$; then $c_y\notin\varpi\mathcal O[[X]]$, so $c_y$ is a unit of $\mathcal O'$, and the image of $c_x/c_y$ under $\hat f$ is $u$.
% rfEx: IsCplAt.existsUnique_lift into FractionRing x.A (maximal ideal bot) and IsCplAt.map_pow_eq with n = 1 for the kernel. ooEx: ooEx_ev (co_ev for F, and the uniqueness in co_ev for f o F), ooEx_ker, IsCplAt.dvr (CaCplDvr: maximal ideal generated by pi, hnil from IsCplAt.strict, Krull intersection) on CplAt (PowerSeries O) p, IsCplAt.existsUnique_lift for O' -> Bh, IsCplAt.isCO (CaCplCo, the residue argument through IsCplAt.dense), and the surjectivity from IsFractionRing.div_surjective.
\end{proof}

\begin{lemma}\label{sha-patching-chain-chart}
Let $P$ be a regular local ring, $g\ge1$, $z=(z_{ij})$ a $g\times g$ matrix over $P$ and $m=(m_j)$ a column vector of length $g$ over $P$, such that the $g^2+g$ elements $z_{ij}$ and $m_j$ lie in $\mathfrak m_P$ and their images in $\mathfrak m_P/\mathfrak m_P^2$ are linearly independent over $P/\mathfrak m_P$. Let $I_z\subset P$ be the ideal generated by the entries of $zm$, and $I'_z=I_z+(\det z)$. Then there is a prime $K\supset I'_z$ of $P$ that is contained in every prime of $P$ containing $I'_z$.
\end{lemma}
\leanref{Development/AutomorphicArithmetic/AutomorphicVarieties/AlgebraTools/CaDetMin.lean}{UpperTheory.V.AlgebraTools.exists_least_prime_detIdeal}
\leanref{Development/AutomorphicArithmetic/AutomorphicVarieties/AlgebraTools/CaDetKerPrime.lean}{UpperTheory.V.AlgebraTools.genKer_isPrime}
\leanref{Development/AutomorphicArithmetic/AutomorphicVarieties/AlgebraTools/CaDetKerLe.lean}{UpperTheory.V.AlgebraTools.genKer_le_ker}
\leanref{Development/AutomorphicArithmetic/AutomorphicVarieties/AlgebraTools/CaDetKerIndep.lean}{UpperTheory.V.AlgebraTools.genKer_eq}
\begin{proof}
Let $P[\lambda]=P[\lambda_1,\dots,\lambda_g]$ be a polynomial ring, and $M=\mathfrak m_PP[\lambda]$, a prime ideal, since $P[\lambda]/M$ is a polynomial ring over the field $P/\mathfrak m_P$. For $1\le c\le g$ let $\lambda^{(c)}$ be the column vector with entries $\lambda^{(c)}_j=\lambda_j$ for $j\ne c$ and $\lambda^{(c)}_c=1$; let $I_c\subset P[\lambda]$ be the ideal generated by the entries of $z\lambda^{(c)}$ and the elements $m_j-\lambda^{(c)}_jm_c$, $j\ne c$; and let $K_c$ be the set of $x\in P$ with $sx\in I_c$ for some $s\in P[\lambda]\setminus M$. The generators of $I_c$ lie in $M$.

\emph{(i) For every prime $\mathfrak Q\supset M$ of $P[\lambda]$, the ideal $I_cP[\lambda]_{\mathfrak Q}$ is prime.} The local ring $P[\lambda]_{\mathfrak Q}$ is flat over $P$, and its fibre over the closed point is a localization of a polynomial ring over a field, so it is regular. So $P[\lambda]_{\mathfrak Q}$ is regular, and a regular system of parameters of $P$ is part of one of $P[\lambda]_{\mathfrak Q}$ \cite[Theorem~23.7]{Mat86}; so the $z_{ij}$ and $m_j$ remain linearly independent in its cotangent space. Then so are the generators of $I_c$. Indeed, write $\mathfrak n$ for the maximal ideal of $P[\lambda]_{\mathfrak Q}$, and suppose that $\sum_ia_i(z\lambda^{(c)})_i+\sum_{j\ne c}b_j(m_j-\lambda_jm_c)\in\mathfrak n^2$. The left side is $\sum_{i,j}a_i\lambda^{(c)}_jz_{ij}+\sum_{j\ne c}b_jm_j-\big(\sum_{j\ne c}b_j\lambda_j\big)m_c$, so all $b_j$ and all $a_i\lambda^{(c)}_j$ lie in $\mathfrak n$, in particular $a_i=a_i\lambda^{(c)}_c$. The quotient of a regular local ring by part of a regular system of parameters is regular, hence a domain \cite[\S14]{Mat86}.

\emph{(ii) $K_c$ is a prime containing $I'_z$.} It is the inverse image of $I_cP[\lambda]_M$ under $P\to P[\lambda]_M$, which is prime by (i). The $c$-th entry of $\operatorname{adj}(z)\,z\lambda^{(c)}=\det(z)\lambda^{(c)}$ is $\det z$, so $\det z\in I_c$. And $(zm)_i=(z\lambda^{(c)})_im_c+\sum_jz_{ij}(m_j-\lambda^{(c)}_jm_c)$, whose term at $j=c$ vanishes, lies in $I_c$.

\emph{(iii) Let $D'$ be a local ring, $\varphi:P\to D'$ a local homomorphism and $l\in D'^g$ with $l_c=1$, $\varphi(z)l=0$ and $\varphi(m_j)=l_j\varphi(m_c)$ for all $j$. Then $K_c\subset\ker\varphi$.} Let $\psi:P[\lambda]\to D'$ extend $\varphi$ by $\lambda_j\mapsto l_j$ for $j\ne c$, so that $\psi(\lambda^{(c)})=l$, and let $\mathfrak Q=\psi^{-1}(\mathfrak m_{D'})$, a prime containing $M$. Then $\psi$ kills $I_c$. Let $x\in K_c$, with $sx\in I_c$ and $s\notin M$. Then $sx$ lies in the ideal $I_cP[\lambda]_{\mathfrak Q}$, which is prime by (i). It does not contain $s$: otherwise $ts\in I_c\subset M$ for some $t\notin\mathfrak Q$, and $t\notin M$, so $s\in M$ as $M$ is prime. So $x\in I_cP[\lambda]_{\mathfrak Q}$, that is $t'x\in I_c$ for some $t'\notin\mathfrak Q$. Then $\psi(t')\varphi(x)=0$ and $\psi(t')$ is a unit, so $\varphi(x)=0$.

\emph{(iv) $K_c\subset K_d$ for all $c,d$.} Let $D'$ be the quotient of $P[\lambda]_M$ by the prime $I_dP[\lambda]_M$, a local ring, and $\varphi:P\to D'$ the map; it is local, since $\mathfrak m_P$ maps into $MP[\lambda]_M$. Let $l_j$ be the image of $\lambda^{(d)}_j$. Then $\varphi(z)l=0$ and $\varphi(m_j)=l_j\varphi(m_d)$ for all $j$. As $\lambda^{(d)}_c\notin M$, $l_c$ is a unit, and $l'=l_c^{-1}l$ satisfies $l'_c=1$, $\varphi(z)l'=0$ and $\varphi(m_j)=l'_j\varphi(m_c)$, since $\varphi(m_c)=l_c\varphi(m_d)$. By (iii), $K_c\subset\ker\varphi$, and $\ker\varphi=K_d$ by the definition of the localization.

\emph{(v) The prime $K=K_1$.} It contains $I'_z$ by (ii). Let $Q\supset I'_z$ be a prime, $L$ the fraction field of $P/Q$, and $V$ a valuation ring of $L$ that dominates $P/Q$ \cite[Theorem~10.2]{Mat86}; the map $\varphi:P\to V$ is local, with kernel $Q$. Write $\bar z$ and $\bar m$ for the images over $L$; then $\bar z\bar m=0$ and $\det\bar z=0$. If $\bar m\ne0$, choose $c$ with $\bar m_j/\bar m_c\in V$ for all $j$, and put $l_j=\bar m_j/\bar m_c$. If $\bar m=0$, choose $w\ne0$ with $\bar zw=0$, $c$ with $w_j/w_c\in V$ for all $j$, and put $l_j=w_j/w_c$. In both cases $l\in V^g$, $l_c=1$, $\bar zl=0$ and $\bar m_j=l_j\bar m_c$. By (iv) and (iii), $K=K_1\subset K_c\subset\ker\varphi=Q$.
% genKer (DetLocus.lean) is K_c, with lam P c j = 1 at c and X j elsewhere and chartIdeal I_c. (i) chartIdeal_localization_isPrime: regular_polynomial_localization (fibre_reg and regular_of_flat_fibre), chartGen_cotInd, regular_quotient_cotInd. (ii) genKer_isPrime. (iii) genKer_le_ker. (iv) generic_kernel_monotone and genKer_eq. (v) exists_least_prime_detIdeal with IsLocalRing.exists_factor_valuationRing and exists_kernel_vector (CaDetField.lean).
\end{proof}

\begin{lemma}\label{sha-patching-chain-incidence}
Let $P$, $g$, $z$, $m$, $I_z$ and $I'_z$ be as in Lemma~\ref{sha-patching-chain-chart}.
\begin{enumerate}
\item $\dim P/I_z+g\le\dim P$.
\item Let $Z\supset I_z$ be a prime of $P$ with $m_j\notin Z$ for some $j$ and $\dim P\le\dim P/Z+g$. Then $Z\subset Q$ for every prime $Q\supset I'_z$ of $P$.
\item Let $B'\subset A'$ be rings with $A'$ integral over $B'$, let $\pi:P\to A'$ be a surjective homomorphism with $\pi(I_z)=0$, and let $Z$ be a prime of $B'$ with $\dim P\le\dim B'/Z+g$ such that every prime of $A'$ over $Z$ omits some $\pi(m_j)$. Then $Z\subset Q'\cap B'$ for every prime $Q'$ of $A'$ with $\pi(\det z)\in Q'$.
\end{enumerate}
\end{lemma}
\leanref{Development/AutomorphicArithmetic/AutomorphicVarieties/AlgebraTools/CaDetDim.lean}{UpperTheory.V.AlgebraTools.dim_relIdeal_le}
\leanref{Development/AutomorphicArithmetic/AutomorphicVarieties/AlgebraTools/CaDetLocus.lean}{UpperTheory.V.AlgebraTools.le_of_detIdeal_le}
\leanref{Development/AutomorphicArithmetic/AutomorphicVarieties/AlgebraTools/CaDetCore.lean}{UpperTheory.V.AlgebraTools.under_le_of_det_mem}
\begin{proof}
(1) Let $K_0$ be the ideal generated by all the $z_{ij}$ and $m_j$, and $G$ the ideal generated by the $g^2$ elements $z_{ij}-\delta_{ij}m_i$. A prime containing $I_z+G$ contains $m_i^2=(zm)_i-\sum_j(z_{ij}-\delta_{ij}m_i)m_j$, hence $m_i$, hence every $z_{ij}$; so it contains $K_0$, and $\dim P/(I_z+G)\le\dim P/K_0$. As $P/K_0$ is the quotient of a regular local ring by part of a regular system of parameters, $\dim P/K_0=\dim P-g^2-g$; and as $G$ is generated by $g^2$ elements, $\dim P/I_z\le\dim P/(I_z+G)+g^2$ \cite[\S14]{Mat86}. Together these give (1).

(2) The $j$-th entry of $\operatorname{adj}(z)\,zm=\det(z)m$ gives $\det(z)m_j\in Z$, so $\det z\in Z$ and $I'_z\subset Z$. Let $K$ be as in Lemma~\ref{sha-patching-chain-chart}; then $K\subset Z$ and $K\subset Q$. If $Z\ne K$, then, as $P/K$ is a domain, $\dim P/Z+1\le\dim P/K$; and $\dim P/K\le\dim P/I_z\le\dim P-g\le\dim P/Z$ by (1) and the hypothesis, which is impossible since these dimensions are finite. So $Z=K\subset Q$.

(3) Let $Z'$ be a prime of $A'$ over $Z$ \cite[Theorem~9.3]{Mat86}, and choose $j$ with $\pi(m_j)\notin Z'$. Since $\pi$ is surjective, $P/\pi^{-1}(Z')\cong A'/Z'$, and $\dim A'/Z'=\dim B'/Z$ by going up and incomparability for the integral extension $B'/Z\subset A'/Z'$ \cite[Theorem~9.3]{Mat86}. So (2) applies to $\pi^{-1}(Z')$, which contains $I_z$ and omits $m_j$, and to $\pi^{-1}(Q')$, which contains $I'_z$. It gives $\pi^{-1}(Z')\subset\pi^{-1}(Q')$, so $Z'\subset Q'$ as $\pi$ is surjective, and $Z=Z'\cap B'\subset Q'\cap B'$.
% (1) dim_relIdeal_le: regular_quotient_cotInd for K_0, dim_quot_le_of_primes, height_le_ringKrullDim_quotient_add_spanFinrank for G. (2) le_of_detIdeal_le: det_mem_of_rel, exists_least_prime_detIdeal, ringKrullDim_succ_le_of_surjective; the case of an empty index set is excluded here by g >= 1. (3) under_le_of_det_mem, with ringKrullDim_quotient_eq_quotient_comap_of_isIntegral.
\end{proof}

\begin{lemma}\label{sha-patching-chain-square-zero}
Let $\mathcal O$ be a complete discrete valuation ring, $B'$ a complete local noetherian $\mathcal O$-algebra such that $\mathcal O$ maps onto its residue field, $A'$ a commutative $B'$-algebra, finite as a $B'$-module, into which $B'$ injects, and $N\subset A'$ a $B'$-submodule with $A'=B'\cdot1\oplus N$ and $xy\in\mathfrak m_{B'}\cdot1$ for $x,y\in N$. Let $I_N\subset B'$ be the ideal generated by the $d\in B'$ with $d\cdot1=xy$ for some $x,y\in N$. Then $A'$ is a complete local noetherian $\mathcal O$-algebra such that $\mathcal O$ maps onto its residue field, $\mathfrak m_{A'}=\mathfrak m_{B'}A'+NA'$, and $I_NA'\subset\mathfrak m_{A'}^2$.
\end{lemma}
\leanref{Development/AutomorphicArithmetic/AutomorphicVarieties/PatchingModules/SelmerScaling/Ired.lean}{UpperTheory.V.Pa.aaCO}
\begin{proof}
Every element of $A'$ is $b+y$ with $b\in B'$ and $y\in N$. Let $S$ be the set of such elements with $b\in\mathfrak m_{B'}$. It is closed under addition, and under multiplication by $A'$: $(b'+y')(b+y)=(b'b+d)+(b'y+by')$, where $y'y=d\cdot1$ with $d\in\mathfrak m_{B'}$. It does not contain $1$: if $1=b+y$, then $y=(1-b)\cdot1$ lies in $B'\cdot1\cap N=0$, so $b=1$, which is not in $\mathfrak m_{B'}$. An element $b+y\notin S$ is a unit: $b\notin\mathfrak m_{B'}$, and $(b+y)(b-y)=(b^2-d)\cdot1$ with $y^2=d\cdot1$ and $d\in\mathfrak m_{B'}$, where $b^2-d$ is a unit of $B'$. So $A'$ is local with maximal ideal $S=\mathfrak m_{B'}A'+NA'$. As $A'$ is finite over $B'$, it is noetherian and complete \cite[\S8]{Mat86}. If $b\equiv o$ modulo $\mathfrak m_{B'}$ with $o\in\mathcal O$, then $b+y\equiv o$ modulo $S$, so $\mathcal O$ maps onto the residue field of $A'$. Finally, a generator $d$ of $I_N$ is $xy$ with $x,y\in N\subset\mathfrak m_{A'}$, so it lies in $\mathfrak m_{A'}^2$.
% aaCO: hdec, the set S with hSmul, hSadd, h1, hnon, hunit, IsLocalRing.of_nonunits_add, co_finite for IsCO O A, and hM.
\end{proof}

\begin{lemma}\label{sha-patching-chain-cotangent}
Let $\mathcal O$ be a complete discrete valuation ring with residue field $\kappa$, let $B'$, $A'$, $N$ and $I_N$ be as in Lemma~\ref{sha-patching-chain-square-zero}, and let $J\subset B'$ be the ideal generated by the image of $\mathfrak m_{\mathcal O}$, $\mathfrak m_{B'}^2$ and $I_N$. Let $D_0,g\ge0$, let $X_1,\dots,X_{D_0}\in\mathfrak m_{B'}$ with $\mathfrak m_{B'}\subset(X_1,\dots,X_{D_0})+J$, let $m_1,\dots,m_g$ generate $N$ as a $B'$-module, and let $c_{i,i',k}\in\mathcal O$, for $1\le i,i'\le g$ and $1\le k\le D_0$, be such that the $g^2$ vectors $(\bar c_{i,i',k})_k$ of residues in $\kappa^{D_0}$ are linearly independent and $\sum_{i'}\big(\sum_kc_{i,i',k}X_k\big)m_{i'}\in JN$ for each $i$. Let $Z$ and $Q$ be primes of $B'$ with $\dim B'/Z\ge1+D_0$ and $I_N\not\subset Z$, such that $I_N\not\subset Q$ or $Q$ contains the annihilator of $N$. Then $Z\subset Q$.
\end{lemma}
\leanref{Development/AutomorphicArithmetic/AutomorphicVarieties/PatchingModules/SelmerScaling/CotangentIndependentOfResidues.lean}{UpperTheory.V.Pa.presOut1}
\begin{proof}
As $I_N\not\subset Z$, $N\ne0$ and $g\ge1$. By Lemma~\ref{sha-patching-chain-square-zero}, $A'$ is a complete local noetherian $\mathcal O$-algebra with residue field $\kappa$, $\mathfrak m_{A'}=\mathfrak m_{B'}A'+NA'$ and $I_NA'\subset\mathfrak m_{A'}^2$. Let $\varpi$ be a uniformizer of $\mathcal O$, and $P=\mathcal O[[x_1,\dots,x_{D_0},y_1,\dots,y_g]]$, a regular local ring of dimension $D_0+g+1$ whose maximal ideal has the regular system of parameters $\varpi$, $x_k$, $y_j$ \cite[\S\S14 and 19]{Mat86}. Let $F:P\to A'$ be the continuous homomorphism of $\mathcal O$-algebras with $F(x_k)=X_k$ and $F(y_j)=m_j$, which lie in $\mathfrak m_{A'}$.

\emph{$F$ is surjective.} The ideal $JA'$ lies in $F(\varpi P)A'+\mathfrak m_{A'}^2$, since $\mathfrak m_{\mathcal O}=\varpi\mathcal O$, $\mathfrak m_{B'}^2A'\subset\mathfrak m_{A'}^2$ and $I_NA'\subset\mathfrak m_{A'}^2$. As $\mathfrak m_{B'}\subset(X_1,\dots,X_{D_0})+J$, the $X_k=F(x_k)$ lie in $F(\mathfrak m_P)$, and $N$ lies in the ideal generated by the $m_j=F(y_j)$, it follows that $\mathfrak m_{A'}\subset F(\mathfrak m_P)A'+\mathfrak m_{A'}^2$. As $P$ and $A'$ are complete with residue field $\kappa$, $F$ is surjective \cite[Theorem~8.4]{Mat86}; and $F(\mathfrak m_P)=\mathfrak m_{A'}$ by Nakayama's lemma.

\emph{The relations.} As the $m_{i'}$ generate $N$, $JN=\sum_{i'}Jm_{i'}$, so for each $i$ there are $w_{i,i'}\in J$ with $\sum_{i'}w_{i,i'}m_{i'}=\sum_{i'}\big(\sum_kc_{i,i',k}X_k\big)m_{i'}$. The images of the $w_{i,i'}$ in $A'$ lie in $JA'\subset F(\varpi P+\mathfrak m_P^2)$, so there are $v_{i,i'}\in\varpi P+\mathfrak m_P^2$ with $F(v_{i,i'})=w_{i,i'}$. Let $\zeta$ be the $g\times g$ matrix over $B'$ with entries $\zeta_{i,i'}=\sum_kc_{i,i',k}X_k-w_{i,i'}$, and $z$ the matrix over $P$ with entries $z_{i,i'}=\sum_kc_{i,i',k}x_k-v_{i,i'}$. Then $F(z_{i,i'})=\zeta_{i,i'}$ and $\sum_{i'}\zeta_{i,i'}m_{i'}=0$, so $F$ kills the entries of $zy$, where $y=(y_j)$.

\emph{The cotangent condition.} The $z_{i,i'}$ and $y_j$ lie in $\mathfrak m_P$, and their images in $\mathfrak m_P/\mathfrak m_P^2$ are linearly independent. Indeed, $\mathfrak m_P/\mathfrak m_P^2$ has the basis $\varpi$, $x_k$, $y_j$ over $\kappa$, and $z_{i,i'}$ is congruent to $\sum_kc_{i,i',k}x_k$ plus a multiple of $\varpi$. If $\sum b_{i,i'}z_{i,i'}+\sum_jb_jy_j\in\mathfrak m_P^2$, the coefficients of the $y_j$ give $\bar b_j=0$, and those of the $x_k$ give $\sum_{i,i'}\bar b_{i,i'}\bar c_{i,i',k}=0$ for every $k$, so all $\bar b_{i,i'}=0$ by the independence of the vectors $(\bar c_{i,i',k})_k$.

\emph{$\det\zeta\in Q$.} Multiplying the relations $\sum_{i'}\zeta_{i,i'}m_{i'}=0$ by the adjugate of $\zeta$ gives $\det(\zeta)m_{i'}=0$ for all $i'$, so $\det\zeta$ annihilates $N$. If $Q$ contains the annihilator of $N$, then $\det\zeta\in Q$. Otherwise $I_N\not\subset Q$. For a generator $d\cdot1=xy$ of $I_N$, $\det(\zeta)d\cdot1=(\det(\zeta)x)y=0$, so $\det(\zeta)d=0\in Q$, as $B'\to A'$ is injective. If $\det\zeta\notin Q$, every such $d$ lies in $Q$, and $I_N\subset Q$, a contradiction.

\emph{Conclusion.} Let $Q'$ be a prime of $A'$ over $Q$ \cite[Theorem~9.3]{Mat86}, so $F(\det z)=\det\zeta\in Q'$. Apply Lemma~\ref{sha-patching-chain-incidence}(3) to $P$, $z$, $y$, $B'\subset A'$, $\pi=F$ and $Z$. Here $\dim P=D_0+g+1\le\dim B'/Z+g$. Every prime $Z'$ of $A'$ over $Z$ omits some $m_j=F(y_j)$: otherwise $N\subset Z'$, so $I_N\subset Z'\cap B'=Z$, which is excluded. So $Z\subset Q'\cap B'=Q$.
% presOut1: aaCO (Lemma sha-patching-chain-square-zero), mvPowerSeries_cotInd for P = MvPowerSeries (Fin D0 + Fin g) O, co_ev for F, hJ and hcot, co_surj and hmP for the surjectivity, hw and hw2 for w and v, zeta and z with hzeta and hzzeta, cotangent_independent_of_residues, det_mem_annihilator_of_relations and hdetQ, Ideal.exists_ideal_over_prime_of_isIntegral for Q', then under_le_of_det_mem (Lemma sha-patching-chain-incidence(3)).
\end{proof}

\begin{lemma}\label{sha-patching-chain-stable}
Let $G$ be a group, $n\ge1$ and $1\le a\le n$. Let $B'$ be a local ring, $A'$ a commutative $B'$-algebra, finite as a $B'$-module, into which $B'$ injects, and $N\subset A'$ a $B'$-submodule with $A'=B'\cdot1\oplus N$ and $xy\in\mathfrak m_{B'}\cdot1$ for $x,y\in N$. Let $D$ be a determinant of $G$ of dimension $n$ with values in $B'$, and $r:G\to M_n(A')$ a homomorphism of monoids such that $D\otimes A'$ is the determinant of $r$ and the $B'$-span of $r(G)$ is the module of matrices whose entries lie in $B'\cdot1$ in the two diagonal blocks, of sizes $a$ and $n-a$, and in $N$ in the two off-diagonal blocks. Let $\mathfrak p$ be a prime of $B'$. Then there are $g_1,\dots,g_{n^2},h_1,\dots,h_{n^2}\in G$ with $\det\big(\operatorname{tr}_D(g_sh_{s'})\big)_{s,s'}\notin\mathfrak p$ if and only if $a=n$ or there are $y,y'\in N$ and $d\in B'\setminus\mathfrak p$ with $yy'=d\cdot1$.
\end{lemma}
\leanref{Development/AutomorphicArithmetic/AutomorphicVarieties/PatchingModules/ClosureConjugation/UrepStable.lean}{UpperTheory.V.Pa.urep_stable}
\begin{proof}
As $A'$ is integral over $B'$, there is a prime $\mathfrak P$ of $A'$ over $\mathfrak p$ \cite[Theorem~9.3]{Mat86}. Let $k'$ be its residue field, $\phi:A'\to k'$ the quotient map and $\rho=\phi\circ r$; for $b\in B'$, $\phi(b)=0$ if and only if $b\in\mathfrak p$. The determinant $D\otimes k'$ is the determinant of $\rho$, so $\phi(\operatorname{tr}_D(g))=\operatorname{tr}\rho(g)$ for $g\in G$ \cite{Che14}, and $\phi$ carries $\det(\operatorname{tr}_D(g_sh_{s'}))$ to the determinant of the matrix $(\operatorname{tr}(\rho(g_s)\rho(h_{s'})))_{s,s'}$.

\emph{The condition holds if and only if the $k'$-span of $\rho(G)$ is $M_n(k')$.} If that determinant is nonzero, the $n^2$ matrices $\rho(g_s)$ are linearly independent, since a relation $\sum_sc_s\rho(g_s)=0$ gives $\sum_sc_s\operatorname{tr}(\rho(g_s)\rho(h_{s'}))=0$ for all $s'$; so they span $M_n(k')$. Conversely, if $\rho(G)$ spans $M_n(k')$, it contains a basis $x_1,\dots,x_{n^2}$, and the matrix $(\operatorname{tr}(x_sx_{s'}))$ is invertible, as the trace form on $M_n(k')$ is nondegenerate; take $g_s=h_s$ with $\rho(g_s)=x_s$.

Since the $B'$-span of $r(G)$ is the module of block matrices described above, the $k'$-span of $\rho(G)$ contains the image under $\phi$ of every such block matrix. Write $e_{ij}$ for the unit matrices.

Suppose that $\rho(G)$ spans $M_n(k')$ and $a<n$. If $\phi(y)=0$ for all $y\in N$, then every matrix in the span of $\rho(G)$ has entry $0$ at $(1,a+1)$, since that entry of each $r(g)$ lies in $N$; this is false for $e_{1,a+1}$. So $\phi(y)\ne0$ for some $y\in N$. Write $y^2=d\cdot1$ with $d\in\mathfrak m_{B'}$; then $\phi(d)=\phi(y)^2\ne0$, so $d\notin\mathfrak p$, and $y'=y$ and $d$ will do.

Conversely, it suffices that the span of $\rho(G)$ contains a nonzero multiple of each $e_{ij}$. If $i$ and $j$ lie in the same block, $e_{ij}$ is a block matrix of the module above, and its image is $e_{ij}$. Otherwise $a<n$, so there are $y,y'\in N$ and $d\notin\mathfrak p$ with $yy'=d\cdot1$; then $\phi(y)\phi(y')=\phi(d)\ne0$, so $\phi(y)\ne0$, and $ye_{ij}$ is a block matrix of the module, with image $\phi(y)e_{ij}$.
% urep_stable: Ideal.exists_ideal_over_prime_of_isIntegral for the prime over Q, tr_of_map (d_gmap, d_toDet), discSt_iff with span_of_gram and CA.g_gram, single_mem_blk and top_of_single; V.mul_mem gives d. discImp turns it into: a = n or ired N is not contained in Q.
\end{proof}

\subsection{Arc data and Cayley-Hamilton algebras}\label{subsection-sha-arc-data}

\begin{definition}\label{sha-patching-chain-datum}
In Situation~\ref{sha-hecke-situation}, let $\Gamma_\infty$ be a profinite group with a continuous surjection $p:\Gamma_\infty\to\Gamma$, an element $c_\infty$ with $c_\infty^2=1$ and $p(c_\infty)=c$, and continuous homomorphisms to $\Gamma_{\infty,0}=p^{-1}(\Gamma_0)$ from the groups $D_v$, $v\mid\ell$, and from the decomposition groups at the places $\tilde u$, $u\in S_{\mathrm a}$, whose composites with $p$ are their maps to $\Gamma_0$. View $\varepsilon$ as a character of $\Gamma_\infty$ through $p$, and for $g\in\Gamma_{\infty,0}$ put $g^*=\varepsilon(g)^{1-n}(c_\infty gc_\infty^{-1})^{-1}$. Let $x$ be a normalized arc over $A=\kappa[[T]]$ satisfying (N1), as in Lemma~\ref{sha-patching-normalization}, with representation $\rho$ and with $I'_u=I_u\cap\bigcap_i\ker\bar\chi_i$ for $u\in S_{\mathrm a}$. Let $\mathcal R$ be a complete local noetherian $\mathcal O$-algebra with residue field $k$, and $(D,\psi)$ an ordinary pair over $\mathcal R$ for $\Gamma_\infty$ and the residual datum of Situation~\ref{sha-hecke-situation} composed with $p$.
\begin{enumerate}
\item An \emph{arc datum} for $x$ over $(\mathcal R,D,\psi)$ consists of an integer $q\ge0$, elements $\phi_j,\tau_j\in\Gamma_{\infty,0}$ for $1\le j\le q$, elements $U_{j,i}$, $y_{j,i}$ and $s$ of $\mathcal R$ for $1\le j\le q$ and $1\le i\le n$, and a homomorphism of $\mathcal O$-algebras $f:\mathcal R\to A$, such that:
\begin{enumerate}
\item $\Gamma_\infty$ is topologically finitely generated, and $\mathcal O$ maps onto $\kappa$;
\item $\phi_j\tau_j=\tau_j\phi_j$ and $p(\tau_j)=1$ for all $j$;
\item $\operatorname{tr}_D\big(g\prod_i(\phi_j-U_{j,i})\big)=0$ and $\operatorname{tr}_D\big(g\prod_{k\ne i}(\phi_j-U_{j,k})\,(\tau_j-1-y_{j,i})\big)=0$ for all $g\in\Gamma_{\infty,0}$, $i$ and $j$;
\item $\mathcal R$ is topologically generated over $\mathcal O$ by the coefficients of $D$, the values of $\psi$, the $U_{j,i}$, the $y_{j,i}$ and $s$;
\item $f$ carries $(D,\psi)$ to the determinant of $\rho\circ p$ and the characters of $x$; $f(s)=T$ and $f(y_{j,i})=0$; and for each $j$ the $f(U_{j,i})$ are pairwise distinct and $\prod_i(X-f(U_{j,i}))$ is the characteristic polynomial of $\rho(p(\phi_j))$.
\end{enumerate}
\item The \emph{relation ideal} of an arc datum is the intersection, over the integers $N\ge1$ and the open normal subgroups $V$ of $\Gamma_{\infty,0}$, of the two-sided ideals of $\mathcal R[\Gamma_{\infty,0}]$ generated by $\mathfrak m_{\mathcal R}^N$, the elements $h-1$ with $h\in V$, the Cayley-Hamilton ideal of $D$ \cite{Che14}, the elements $\prod_i(g_i-\psi_{v,i}(g_i))$ with $v\mid\ell$ and $g_i\in D_v$, the elements $g-1$ with $g$ in the image of some $I'_u$, the elements $\prod_i(\phi_j-U_{j,i})$ and $\prod_{k\ne i}(\phi_j-U_{j,k})(\tau_j-1-y_{j,i})$, and the images of all these elements under the $\mathcal R$-linear anti-automorphism of $\mathcal R[\Gamma_{\infty,0}]$ that sends $g$ to $g^*$.
\end{enumerate}
\end{definition}
\leanref{Development/AutomorphicArithmetic/AutomorphicVarieties/PatchingModules/PatchedSupportCompletions.lean}{UpperTheory.V.Pa.OoPt}
\leanref{Development/AutomorphicArithmetic/AutomorphicVarieties/PatchingModules/PatchedSupportCompletions.lean}{UpperTheory.V.Pa.OoPt.ker}
\leanref{Development/AutomorphicArithmetic/AutomorphicVarieties/PatchingModules/ThetaCoefficientMapEq.lean}{UpperTheory.V.Pa.qGens}
\leanref{Development/AutomorphicArithmetic/AutomorphicVarieties/WeilGroupCharacters/RepAlg.lean}{UpperTheory.V.W.closureIdeal}

\begin{lemma}\label{sha-patching-chain-cayley-factor}
Let $x$ be a normalized arc over $A=\kappa[[T]]$, with representation $\rho$, and let $(q,\phi_j,\tau_j,U_{j,i},y_{j,i},s,f)$ be an arc datum for $x$ over $(\mathcal R,D,\psi)$, with relation ideal $I$, in the sense of Definition~\ref{sha-patching-chain-datum}. Then the homomorphism $\mathcal R[\Gamma_{\infty,0}]\to M_n(A)$ that is $f$ on $\mathcal R$ and sends $g$ to $\rho(p(g))$ kills $I$.
\end{lemma}
\leanref{Development/AutomorphicArithmetic/AutomorphicVarieties/PatchingModules/ClosureConjugation/RepQuot.lean}{UpperTheory.V.Pa.rep_quot}
\leanref{Development/AutomorphicArithmetic/AutomorphicVarieties/PatchingModules/ClosureConjugation/RepGens.lean}{UpperTheory.V.Pa.rep_gens}
\leanref{Development/AutomorphicArithmetic/AutomorphicVarieties/CommutativeAlgebraHeight/LiftEq.lean}{UpperTheory.V.CA.f_rep}
\begin{proof}
Call the homomorphism $R$. As in the proof of Lemma~\ref{sha-patching-chain-cayley}(3), $R(g^*)={}^t\rho(p(g))$ for $g\in\Gamma_{\infty,0}$, since $x$ is normalized and $f$ carries the image of $\varepsilon(g)$ in $\mathcal R$ to its image in $A$; so $R(y^*)={}^tR(y)$, and $R$ kills the images under $*$ of the elements it kills. It kills the other generators of the ideals in Definition~\ref{sha-patching-chain-datum}(2), other than $\mathfrak m_{\mathcal R}^N$ and the $h-1$:
\begin{enumerate}
\item[(i)] the Cayley-Hamilton ideal of $D$, as $f$ carries $D$ to the determinant of $\rho\circ p$ \cite{Che14};
\item[(ii)] $\prod_i(g_i-\psi_{v,i}(g_i))$ for $g_i\in D_v$: $f$ carries $\psi_{v,i}$ to the character of $x$, and after conjugation by $w_v$ the factor $\rho(p(g_i))-f(\psi_{v,i}(g_i))$ is upper triangular with $(i,i)$ entry $0$, so the product vanishes, as in step (i) of the proof of Lemma~\ref{sha-patching-chain-cocycles-ordinary};
\item[(iii)] $g-1$ for $g$ in the image of $I'_u$, as $\rho$ is trivial on $I'_u$;
\item[(iv)] $\prod_i(\phi_j-U_{j,i})$, which maps to $\prod_i(\rho(p(\phi_j))-f(U_{j,i}))=0$ by the Cayley-Hamilton theorem, as $\prod_i(X-f(U_{j,i}))$ is the characteristic polynomial of $\rho(p(\phi_j))$;
\item[(v)] $\prod_{k\ne i}(\phi_j-U_{j,k})(\tau_j-1-y_{j,i})$, whose last factor maps to $\rho(p(\tau_j))-1-f(y_{j,i})=1-1-0=0$.
\end{enumerate}
Let $N\ge1$. As $\rho$ and $p$ are continuous, there is an open normal subgroup $V$ of $\Gamma_{\infty,0}$ with $\rho(p(h))\equiv1$ modulo $T^N$ for $h\in V$; and $f(\mathfrak m_{\mathcal R}^N)\subset T^NA$, as $f$ is local. So $R$ maps the two-sided ideal generated by $\mathfrak m_{\mathcal R}^N$, the $h-1$ with $h\in V$, and the elements above into $M_n(T^NA)$. Every element of $I$ lies in this ideal, so $R(I)\subset\bigcap_NM_n(T^NA)=0$.
% rep_quot: f_rep (key: the closure through the open normal subgroup H and lift_CH for the Cayley-Hamilton ideal, then RingCon.lift) with rep_gens (rep_star for *, conj_prod_zero with x.tri and P.maps for the ordinary relations, x.ine, ch_prod_zero with P.eig, P.p_ta and P.f_y).
\end{proof}

\begin{definition}\label{sha-patching-chain-cayley-closure}
Let $B$ be a commutative ring and $E$ a $B$-algebra with a determinant $D$ of dimension $n$ \cite{Che14}. An element $y\in E$ is \emph{killed} by $D$ if $D(1+hy)=1$ for every commutative $B$-algebra $B'$ and every $h\in B'\otimes_BE$; these elements form the kernel $\ker D$, a two-sided ideal of $E$ \cite[Lemma~1.19]{Che14}. Now let $B$ be local and $G$ a topological group. The \emph{closure} of a subset $S\subset B[G]$ is the intersection, over the integers $N\ge1$ and the open normal subgroups $V$ of $G$, of the two-sided ideals of $B[G]$ generated by $S$, $\mathfrak m_B^N$ and the elements $v-1$ with $v\in V$. For a determinant $D$ of $G$ with values in $B$ and a subset $Z\subset B[G]$, let $E_{D,Z}$ be the quotient of $B[G]$ by the closure of the union of $Z$ and the Cayley-Hamilton ideal $\operatorname{CH}(D)$ of $D$.
\end{definition}
\leanref{Development/AutomorphicArithmetic/AutomorphicVarieties/WeilGroupCharacters/RepAlg.lean}{UpperTheory.V.W.GDet.Kills}
\leanref{Development/AutomorphicArithmetic/AutomorphicVarieties/WeilGroupCharacters/RepAlg.lean}{UpperTheory.V.W.closureIdeal}
\leanref{Development/AutomorphicArithmetic/AutomorphicVarieties/CommutativeAlgebraHeight/OfDetToDet.lean}{UpperTheory.V.W.GDet.chBar}
\leanref{Development/AutomorphicArithmetic/AutomorphicVarieties/CommutativeAlgebraHeight/OfDetToDet.lean}{UpperTheory.V.W.GDet.CHQuot}

\begin{definition}\label{sha-patching-chain-cayley-extension-bound}
Let $G$ be a topological group, $k$ a field and $\chi,\chi':G\to k^\times$ characters. Let $Z^1(\chi,\chi')$ be the $k$-vector space of locally constant maps $a:G\to k$ with $a(gh)=\chi(g)a(h)+a(g)\chi'(h)$ for all $g,h\in G$. Let $h\ge0$ be an integer. A family of characters $\chi_1,\dots,\chi_n:G\to k^\times$ has \emph{extension bound} $h$ if the $\chi_i$ are pairwise distinct with open kernels and, for $i\ne j$, $Z^1(\chi_i,\chi_j)$ is finite-dimensional of dimension at most $h$. Put $c(n,h)=n+\sum_{p=0}^{n-1}(n^2h)^p$.
\end{definition}
\leanref{Development/AutomorphicArithmetic/AutomorphicVarieties/CommutativeAlgebraHeight/FinDefs.lean}{UpperTheory.V.CA.Z1}
\leanref{Development/AutomorphicArithmetic/AutomorphicVarieties/CommutativeAlgebraHeight/FinDefs.lean}{UpperTheory.V.CA.MultFree}
\leanref{Development/AutomorphicArithmetic/AutomorphicVarieties/CommutativeAlgebraHeight/FinDefs.lean}{UpperTheory.V.CA.finBound}

\begin{lemma}\label{sha-patching-chain-cayley-bound}
Let $\ell$ be a prime and $G$ a profinite group.
\begin{enumerate}
\item If $G$ is topologically finitely generated, then every open subgroup of $G$ has finitely many continuous homomorphisms to $\mathbf Z/\ell$.
\item Suppose that every open subgroup of $G$ has finitely many continuous homomorphisms to $\mathbf Z/\ell$. Let $k$ be a finite field of characteristic $\ell$ and $\chi_1,\dots,\chi_n:G\to k^\times$ pairwise distinct characters with open kernels. Then the family $(\chi_i)$ has an extension bound.
\end{enumerate}
\end{lemma}
\leanref{Development/AutomorphicArithmetic/AutomorphicVarieties/PatchingModules/ClosureConjugation/PhiFg.lean}{UpperTheory.V.Pa.phi_fg}
\leanref{Development/AutomorphicArithmetic/AutomorphicVarieties/AlgebraTools/SubgroupQuot.lean}{UpperTheory.V.AlgebraTools.TopFG.subgroup_quot}
\leanref{Development/AutomorphicArithmetic/AutomorphicVarieties/CommutativeAlgebraHeight/GammaOneGammaZeroHomogeneousCocyclesFinite.lean}{UpperTheory.V.CA.exists_multFree_of_phiCond_of_injective_chars}
\begin{proof}
(1) Let $\Delta$ be a dense subgroup of $G$ generated by a finite set, and $H\subset G$ an open subgroup. As $H$ has finite index in $G$, $H\cap\Delta$ has finite index in $\Delta$, so it is finitely generated by Schreier's lemma; and it is dense in $H$, as $H$ is open. Let $t$ be a finite set generating $H\cap\Delta$. Two continuous homomorphisms $\phi,\psi:H\to\mathbf Z/\ell$ that agree on $t$ are equal: they agree on the open subgroup $\ker\phi\cap\ker\psi$, so the subgroup on which they agree is open, hence closed, and it contains $t$. So there are at most $\ell^{|t|}$ of them.

(2) Let $i,j$ be indices and $N=\ker\chi_i\cap\ker\chi_j$, an open normal subgroup of finite index. For $a\in Z^1(\chi_i,\chi_j)$ and $x,y\in N$, $a(xy)=a(x)+a(y)$. So, for a basis $b_1,\dots,b_m$ of $k$ over $\mathbf F_\ell$, the coordinates of the restriction of $a$ to $N$ are homomorphisms $N\to\mathbf Z/\ell$, continuous as $a$ is locally constant. For $g\in G$ write $g=rx$ with $r$ in a fixed set of representatives of $G/N$ and $x\in N$; then $a(g)=\chi_i(r)a(x)+a(r)$. So $a$ is determined by its values at the representatives and by $m$ continuous homomorphisms $N\to\mathbf Z/\ell$, of which there are finitely many. Hence $Z^1(\chi_i,\chi_j)$ is finite, and $h=\sum_{i,j}\dim_kZ^1(\chi_i,\chi_j)$ is an extension bound.
% phi_fg: TopFG.subgroup_quot (Group.closure_finite_fg, the finite index subgroupOf, density through hH.isOpenMap_subtype_val) and finite_open_homs_to_finite_group (the eqLocus is open, hence closed). exists_multFree_of_phiCond_of_injective_chars: module_finite_z1_padicInt (hadd, the coordinates phi in Module.finBasis (ZMod l) k, the injective map F to the values at the representatives and the open homomorphisms).
\end{proof}

\begin{lemma}\label{sha-patching-chain-cayley-lift}
Let $A$ be a complete local noetherian ring, $E$ an $A$-algebra that is integral over $A$, $F$ an $A$-algebra and $\pi:E\to F$ a surjective homomorphism of $A$-algebras whose kernel lies in the Jacobson radical of $E$. Let $f_1,\dots,f_m\in F$ be orthogonal idempotents with $\sum_if_i=1$. Then there are orthogonal idempotents $e_1,\dots,e_m\in E$ with $\sum_ie_i=1$ and $\pi(e_i)=f_i$ for all $i$.
\end{lemma}
\leanref{Development/AutomorphicArithmetic/AutomorphicVarieties/CommutativeAlgebraHeight/DLiftIdem.lean}{UpperTheory.V.CA.d_lift_idem}
\begin{proof}
\emph{An idempotent $u$ in the Jacobson radical is $0$.} Indeed $1-u$ has a left inverse $z$, and $u=z(1-u)u=0$.

\emph{One idempotent.} Let $x\in E$ with $\pi(x)$ idempotent. Then the $A$-subalgebra $A[x]$ contains an idempotent $e$ with $\pi(e)=\pi(x)$. Indeed, as $x$ is integral over $A$, the commutative ring $A[x]$ is finite over $A$, so it is complete and separated for the ideal $K=\mathfrak m_AA[x]$ \cite[\S8]{Mat86}. The element $y=x-x^2$ lies in every maximal ideal $M$ of $A[x]$. To see this, note that $\pi(y)=0$, so $y$ lies in the Jacobson radical of $E$. If $y\notin M$, there are $b\in A[x]$ and $m\in M$ with $by=1-m$, and $m=1-by$ has a left inverse $z$ in $E$. Let $z^d+a_{d-1}z^{d-1}+\dots+a_0=0$ with $a_i\in A$. Multiplying on the right by $m^d$ and using $z^im^i=1$ gives $1+a_{d-1}m+\dots+a_0m^d=0$, so $1\in M$, which is absurd. As $A[x]$ is integral over $A$, the primes of $A[x]$ that contain $K$ are maximal \cite[\S9]{Mat86}; so $y$ lies in all of them, and $y^r\in K$ for some $r\ge1$. Let $a=1-(1-x^r)^r$. The polynomial $1-(1-X^r)^r$ is divisible by $X^r$, and $1-X^r$ by $1-X$; so $(1-x^r)^ra$ is a multiple of $(x-x^2)^r$, and $a^2-a=-(1-x^r)^ra\in K$. As $\pi(x)$ and $1-\pi(x)$ are idempotent, $\pi(a)=1-(1-\pi(x))=\pi(x)$. The derivative $2a-1$ of $X^2-X$ at $a$ is a unit modulo $K$, as $(2a-1)^2=1+4(a^2-a)$; so Hensel's lemma for $X^2-X$ in the $K$-adically complete ring $A[x]$ gives an idempotent $e\in A[x]$ with $w=e-a\in K$, as in Lemma~\ref{sha-patching-chain-idempotents-lift}. Since $e^2=e$, $w(1-w^2)=-(3e-a-1)(a^2-a)$. Here $1-w^2$ is a unit of $A[x]$, as $w\in K$ and $K$ lies in every maximal ideal of $A[x]$, and $\pi(a^2-a)=0$; so $\pi(w)=0$ and $\pi(e)=\pi(a)=\pi(x)$.

\emph{Orthogonal lifts.} Let $f\in F$ be idempotent and $e'\in E$ an idempotent with $f\pi(e')=\pi(e')f=0$. Then there is an idempotent $e\in E$ with $\pi(e)=f$ and $ee'=e'e=0$. Indeed, choose $y\in E$ with $\pi(y)=f$, and let $x=(1-e')y(1-e')$, so that $\pi(x)=f$. Let $e\in A[x]$ be as in the previous step. As $e'x=xe'=0$, $e'$ commutes with every element of $A[x]$, in particular with $e$. So $e'e$ is idempotent, and $\pi(e'e)=\pi(e')f=0$; by the first step $e'e=0$, and $ee'=e'e=0$.

\emph{Conclusion.} By induction on $m$, there are orthogonal idempotents $e_1,\dots,e_m$ with $\pi(e_i)=f_i$: given such $e_2,\dots,e_m$ for $f_2,\dots,f_m$, apply the previous step to $f_1$ and $e'=e_2+\dots+e_m$; then $e_1e_i=e_1e'e_i=0$ and $e_ie_1=e_ie'e_1=0$ for $i\ge2$. Finally $1-\sum_ie_i$ is an idempotent in the kernel of $\pi$, so it is $0$ by the first step.
% d_lift_idem: dli_fin (induction with dli_orth and OrthogonalIdempotents.option), then dli_idem_zero for the sum. dli_orth: dli_one for (1 - e2) y (1 - e2), Algebra.commute_of_mem_adjoin_singleton_of_commute and dli_idem_zero. dli_one: Algebra.adjoin A {x} is finite, dli_mem_max (the integral equation of the left inverse z), then dli_comm: isAdicComplete_map_of_finite, the radical of K through isMaximal_of_isIntegral_of_isMaximal_comap, isIdempotentElem_one_sub_one_sub_pow_pow, HenselianRing.is_henselian for X^2 - X, and the identity (e - a)((e - a)(-(e - a)) + 1) = -(3e - a - 1)(a^2 - a).
\end{proof}

\begin{lemma}\label{sha-patching-chain-cayley-gma}
Let $B$ be a complete local noetherian ring, $E$ a $B$-algebra with a Cayley-Hamilton determinant $D$ of dimension $n$ \cite{Che14}, $k$ a field with a local homomorphism $B\to k$, and $\pi:E\to k^n$ a surjective homomorphism of $B$-algebras such that, for every $x\in E$, the image in $k[X]$ of the characteristic polynomial of $D$ at $x$ is $\prod_{i=1}^n(X-\pi(x)_i)$. Then the kernel of $\pi$ is the Jacobson radical of $E$, and there are orthogonal idempotents $e_1,\dots,e_n\in E$ with $\sum_ie_i=1$ such that $\pi(e_i)$ is the $i$-th unit vector $\delta_i$ of $k^n$ and $D(1+Xe_i)=1+X$ in $B[X]$, for all $i$.
\end{lemma}
\leanref{Development/AutomorphicArithmetic/AutomorphicVarieties/CommutativeAlgebraHeight/FgChi.lean}{UpperTheory.V.CA.f_gma}
\leanref{Development/AutomorphicArithmetic/AutomorphicVarieties/CommutativeAlgebraHeight/DUnit.lean}{UpperTheory.V.CA.d_unit}
\leanref{Development/AutomorphicArithmetic/AutomorphicVarieties/CommutativeAlgebraHeight/DIdemRank.lean}{UpperTheory.V.CA.d_idem_rank}
\begin{proof}
For $x\in E$ write $\chi_x=D(X-x)\in B[X]$ for the characteristic polynomial of $D$ at $x$, which is monic of degree $n$.

\emph{Units.} As $D$ is Cayley-Hamilton, $\chi_x(x)=0$. So $E$ is integral over $B$, and $x$ is a unit whenever $D(x)=(-1)^n\chi_x(0)$ is a unit of $B$ \cite[proof of Lemma~2.7(i)]{Che14}. If $\pi(x)=1$, the image of $\chi_x(0)$ in $k$ is $\prod_i(0-1)=(-1)^n\ne0$; as $B\to k$ is local, $\chi_x(0)$ is a unit, and so is $x$.

\emph{The radical.} If $\pi(x)=0$, then $\pi(yx+1)=1$ for every $y\in E$, so $yx+1$ is a unit, and $x$ lies in the Jacobson radical. Conversely, let $x$ lie in the Jacobson radical, and suppose that $\pi(x)_i\ne0$ for some $i$. Choose $y\in E$ with $\pi(y)=-\pi(x)_i^{-1}(1,\dots,1)$, and $z\in E$ with $z(yx+1)=1$. The $i$-th coordinate of the image of this equation under $\pi$ is $\pi(z)_i(-1+1)=1$, which is absurd.

\emph{Idempotents.} Lemma~\ref{sha-patching-chain-cayley-lift}, applied to $\pi$ and to $\delta_1,\dots,\delta_n$, gives orthogonal idempotents $e_i$ with $\sum_ie_i=1$ and $\pi(e_i)=\delta_i$. Fix $i$ and write $e=e_i$. By \cite[Lemma~2.4(1)]{Che14} there are $r,s\le n$ with $D(1+Xe)=(1+X)^r$ and $D(1+X(1-e))=(1+X)^s$. Substituting $X-2$ and $X-1$ for $X$, and using $(1+(X-2)e)(1+(X-1)(1-e))=X-e$, gives $\chi_e=(X-1)^rX^s$. Its image in $k[X]$ is $\prod_j(X-\pi(e)_j)=(X-1)X^{n-1}$. If $r=0$, evaluating at $1$ gives $1=0$; if $r\ge2$, cancelling $X-1$ and evaluating at $1$ gives $0=1$. So $r=1$.
% f_gma: d_unit (hu, hint) and hunit, hk for ker pi in the radical, the converse by Ideal.mem_jacobson_iff; d_lift_idem for the unit vectors; d_idem_rank for e and 1 - e, fg_ev, fg_chi and fg_one for r = 1.
\end{proof}

\begin{lemma}\label{sha-patching-chain-cayley-trace}
Let $A$ be a commutative ring, $E$ an $A$-algebra with a Cayley-Hamilton determinant $D$ of dimension $n\ge1$ \cite{Che14}, and $e_1,\dots,e_n\in E$ orthogonal idempotents with $\sum_ie_i=1$ and $D(1+Xe_i)=1+X$ in $A[X]$ for all $i$. Then $y\in E$ is killed by $D$ if and only if $\operatorname{tr}_D(xy)=0$ for all $x\in E$.
\end{lemma}
\leanref{Development/AutomorphicArithmetic/AutomorphicVarieties/CommutativeAlgebraHeight/GmaKills.lean}{UpperTheory.V.CA.gma_kills}
\begin{proof}
If $D$ kills $y$, then $D(1+Xxy)=1$ in $A[X]$, so its coefficient $\operatorname{tr}_D(xy)$ of $X$ vanishes. Conversely, suppose that $\operatorname{tr}_D(xy)=0$ for all $x$. For each $j$ and $x\in E$, Lemma~\ref{sha-patching-chain-idempotents-corner} gives $e_jxye_j=\operatorname{tr}_D(e_jxy)e_j=0$. Let $B'$ be a commutative $A$-algebra and $h\in B'\otimes_AE$; by linearity, $e_jhye_j=0$. So $D(1+hye_j)=D(1+(hye_j)e_j)=D(1+e_jhye_j)=1$ by \cite[Lemma~1.12(i)]{Che14}, and $D$ kills $ye_j$. Hence $D$ kills $y=\sum_jye_j$, as $\ker D$ is an ideal.
% gma_kills: d_ker (the trace of h x for x in the kernel), d_corner (Lemma sha-patching-chain-idempotents-corner), h0 by induction on the tensor, d_comm for D(1 + ab) = D(1 + ba).
\end{proof}

\begin{lemma}\label{sha-patching-chain-cayley-chars}
Let $k$ be a field, $G$ a group, $\chi_1,\dots,\chi_n:G\to k^\times$ pairwise distinct characters, and $D_\chi$ the determinant of $\chi_1\oplus\dots\oplus\chi_n$. Let $\pi_\chi:k[G]\to k^n$ be the homomorphism of $k$-algebras with $\pi_\chi(g)=(\chi_1(g),\dots,\chi_n(g))$. Then $\pi_\chi$ is surjective; $D_\chi$ kills $y\in k[G]$ if and only if $\pi_\chi(y)=0$; and the characteristic polynomial of $D_\chi$ at $y$ is $\prod_i(X-\pi_\chi(y)_i)$.
\end{lemma}
\leanref{Development/AutomorphicArithmetic/AutomorphicVarieties/CommutativeAlgebraHeight/FChars.lean}{UpperTheory.V.CA.f_chars}
\begin{proof}
For a commutative $k$-algebra $S$, let $\pi_S:S[G]\to S^n$ be the $S$-linear extension of $\pi_\chi$; then $D_\chi(y)=\prod_i\pi_S(y)_i$ for $y\in S[G]$, as $\chi_1\oplus\dots\oplus\chi_n$ is diagonal. This gives the characteristic polynomial. The $\chi_i$ are linearly independent as functions on $G$, by the linear independence of characters, so the vectors $\pi_\chi(g)$ span $k^n$. If $\pi_\chi(y)=0$, then $\pi_S(hy)=0$ and $D_\chi(1+hy)=1$ for all $h$. If $\pi_\chi(y)_i\ne0$, choose $h\in k[G]$ with $\pi_\chi(h)=-\pi_\chi(y)_i^{-1}(1,\dots,1)$; then the $i$-th factor of $D_\chi(1+hy)$ is $0$, so $D_\chi$ does not kill $y$.
% f_chars: piB, D_eq, linearIndependent_monoidHom with span_flip_eq_top_iff_linearIndependent for the surjectivity.
\end{proof}

\begin{lemma}\label{sha-patching-chain-cayley-dimension}
Let $k$ be a field, $G$ a topological group, $\chi_1,\dots,\chi_n:G\to k^\times$ a family of characters with extension bound $h$, and $E$ a finite-dimensional $k$-algebra. Let $\rho:G\to E$ be a homomorphism of monoids with open kernel whose image spans $E$, $\pi:E\to k^n$ a homomorphism of $k$-algebras with $\pi(\rho(g))=(\chi_1(g),\dots,\chi_n(g))$ for $g\in G$, and $e_1,\dots,e_n\in E$ orthogonal idempotents with $\sum_ie_i=1$, $\pi(e_i)=\delta_i$ and $e_iEe_i=ke_i$ for all $i$. Then $\dim_kE\le c(n,h)$.
\end{lemma}
\leanref{Development/AutomorphicArithmetic/AutomorphicVarieties/CommutativeAlgebraHeight/NatCastPowLePowOfExponentLe.lean}{UpperTheory.V.CA.f_cot}
\leanref{Development/AutomorphicArithmetic/AutomorphicVarieties/CommutativeAlgebraHeight/FNil.lean}{UpperTheory.V.CA.f_nil}
\begin{proof}
If $n=0$, then $1=0$ in $E$ and $E=0$. Let $n\ge1$ and $J=\ker\pi$. For $x\in E$, $e_ixe_i=ae_i$ for some $a\in k$, and applying the $i$-th coordinate of $\pi$ gives $a=\pi(x)_i$; so $e_ixe_i=\pi(x)_ie_i$, which vanishes for $x\in J$. Also $\pi(e_ixe_j)=0$ for $i\ne j$.

\emph{$J^n=0$.} For $T\subset\{1,\dots,n\}$ let $e_T=\sum_{i\in T}e_i$. We show by induction on $|T|$ that $e_Ty_1e_Ty_2\cdots e_Ty_me_T=0$ for $y_1,\dots,y_m\in J$ with $m\ge|T|$; for $T=\{1,\dots,n\}$ this gives $y_1\cdots y_n=0$. For $T=\emptyset$ it is clear. Otherwise write the first factor $e_T$ as $\sum_{i\in T}e_i$, fix $i\in T$ and let $T'=T\setminus\{i\}$, so that $e_T=e_i+e_{T'}$. For $a,z_1,\dots,z_r\in J$, $e_iae_Tz_1e_T\cdots z_re_T=e_iae_{T'}z_1e_{T'}\cdots z_re_{T'}$, by induction on $r$: as $e_iae_i=0$, $e_iae_T=e_iae_{T'}$, and $ae_{T'}z_1\in J$. With $a=y_1$ and $z_s=y_{s+1}$, the term of $i$ is $e_iy_1\cdot e_{T'}y_2e_{T'}\cdots y_me_{T'}$, which vanishes by induction, as $m-1\ge|T'|$.

\emph{$\dim J\le\dim J^2+n^2h$.} Let $\Lambda$ be the space of linear forms on $E$ that vanish on $J^2$. For $f\in\Lambda$ and $i\ne j$, the map $a_{f,ij}(g)=f(e_i\rho(g)e_j)$ lies in $Z^1(\chi_i,\chi_j)$. It is locally constant, as $\rho$ has open kernel. And $e_i\rho(gg')e_j=\sum_le_i\rho(g)e_l\cdot e_l\rho(g')e_j$: for $l\ne i,j$ both factors lie in $J$, so $f$ vanishes on the product; for $l=i$ the term is $\chi_i(g)e_i\rho(g')e_j$, and for $l=j$ it is $e_i\rho(g)e_j\chi_j(g')$. Let $\Psi:\Lambda\to\bigoplus_{i\ne j}Z^1(\chi_i,\chi_j)$ be the linear map $f\mapsto(a_{f,ij})$. If $\Psi(f)=0$, then $f$ vanishes on $J$: for $x\in J$, $x=\sum_{i,j}e_ixe_j$, the terms with $i=j$ vanish, and $f(e_ixe_j)=0$ for $i\ne j$, as $x$ is a linear combination of elements $\rho(g)$. So $\dim\Lambda\le\dim J^\perp+n^2h$, that is, $\dim E-\dim J^2\le\dim E-\dim J+n^2h$.

\emph{Conclusion.} Let $V$ be a complement of $J^2$ in $J$, so that $\dim V\le n^2h$ and $J\subset V+J^2$. By induction on $p$, $J^p\subset V^p+J^{p+1}$, where $V^p$ is spanned by products of $p$ elements of a basis of $V$; so $\dim J^p\le(n^2h)^p+\dim J^{p+1}$. Summing over $1\le p\le n-1$ and using $J^n=0$ gives $\dim J\le\sum_{p=1}^{n-1}(n^2h)^p$. As $\pi(E)\subset k^n$, $\dim E\le n+\dim J\le c(n,h)$.
% f_cot: f_nil (idempotent_sandwich_list_prod_eq_of_kernel and the induction D on the Finset T), finrank_ker_le_finrank_ker_sq_add (dual_offdiag_corner_cocycle with idempotent_corner_eq_proj_smul and proj_offdiag_corner_eq_zero, dual_vanishes_on_ker_of_offdiag_vanishing, the dual annihilators), Submodule.exists_isCompl for V, natCast_pow_le_pow_of_exponent_le, fr_pow and fr_nil.
\end{proof}

\begin{lemma}\label{sha-patching-chain-cayley-field}
Let $k$ be a field, $G$ a profinite group, $\chi_1,\dots,\chi_n:G\to k^\times$ a family of characters with extension bound $h$, $D_\chi$ the determinant of $\chi_1\oplus\dots\oplus\chi_n$, and $U$ an open normal subgroup of $G$. Then there is a finite set $s\subset G$ with at most $c(n,h)$ elements such that every element of $k[G]$ is congruent to a $k$-linear combination of elements of $s$ modulo the two-sided ideal generated by $\operatorname{CH}(D_\chi)$ and the elements $u-1$ with $u\in U$.
\end{lemma}
\leanref{Development/AutomorphicArithmetic/AutomorphicVarieties/CommutativeAlgebraHeight/FLevel.lean}{UpperTheory.V.CA.f_level}
\begin{proof}
As $G$ is profinite and the $\chi_i$ have open kernels, there is an open normal subgroup $H\subset U\cap\bigcap_i\ker\chi_i$. Let $J$ be the two-sided ideal of $k[G]$ generated by $\operatorname{CH}(D_\chi)$ and the $h'-1$ with $h'\in H$, $E=k[G]/J$ and $q:k[G]\to E$ the quotient map. Let $\pi_\chi$ be as in Lemma~\ref{sha-patching-chain-cayley-chars}.

\emph{The determinant of $E$.} $D_\chi$ kills $J$: it kills $\operatorname{CH}(D_\chi)$ \cite[Lemma~1.21]{Che14}, and the $h'-1$ by Lemma~\ref{sha-patching-chain-cayley-chars}, as $\pi_\chi(h'-1)=0$. So $D_\chi=D_E\circ q$ for a determinant $D_E$ of $E$ \cite[Lemma~1.19(ii)]{Che14}. It is Cayley-Hamilton: the Cayley-Hamilton ideal is generated by the coefficients of the elements $\chi_y(y)$, $y=t_1r_1+\dots+t_mr_m$, where $\chi_y$ is the characteristic polynomial at $y$ \cite[\S1.17]{Che14}, so as $q$ is surjective the Cayley-Hamilton ideal of $D_E$ is the image of $\operatorname{CH}(D_\chi)\subset J$. By Lemma~\ref{sha-patching-chain-cayley-chars}, $J$ lies in the kernel of $\pi_\chi$, so $\pi_\chi=\pi\circ q$ for a surjective homomorphism $\pi:E\to k^n$, and the characteristic polynomial of $D_E$ at $q(y)$ is $\prod_i(X-\pi(q(y))_i)$.

\emph{$E$ is finite-dimensional.} Let $\rho:G\to E$ be $g\mapsto q(g)$. Its kernel contains $H$, so it is open, and $\rho(G)$ is finite, as $G/H$ is finite; and $\rho(G)$ spans $E$.

\emph{The bound.} Lemma~\ref{sha-patching-chain-cayley-gma}, for the complete local noetherian ring $k$, gives orthogonal idempotents $e_1,\dots,e_n$ of $E$ with $\sum_ie_i=1$, $\pi(e_i)=\delta_i$ and $D_E(1+Xe_i)=1+X$; and $e_iEe_i=ke_i$ by Lemma~\ref{sha-patching-chain-idempotents-corner}, as $n\ge i\ge1$. By Lemma~\ref{sha-patching-chain-cayley-dimension}, $\dim_kE\le c(n,h)$. As $\rho(G)$ spans $E$, it contains a basis $\rho(g_1),\dots,\rho(g_d)$ with $d\le c(n,h)$; let $s=\{g_1,\dots,g_d\}$. Every $y\in k[G]$ is congruent modulo $J$ to a combination of elements of $s$, and $J$ lies in the two-sided ideal generated by $\operatorname{CH}(D_\chi)$ and the $u-1$, $u\in U$, as $H\subset U$.
% f_level, through the private level_field: ProfiniteGrp.exist_openNormalSubgroup_sub_open_nhds_of_one for H, f_chars, d_ch_ker with d_toDet_CH and hk for hJ, d_factor for DE, d_pull for hCH, RingCon.lift for pi', Subgroup.quotient_finite_of_isOpen for the finiteness, f_gma with IsLocalRing.maximalIdeal_eq_bot, d_corner and f_cot, Submodule.exists_fun_fin_finrank_span_eq for the basis.
\end{proof}

\begin{lemma}\label{sha-patching-chain-cayley-level}
Let $B$ be a complete local noetherian ring with finite residue field $k$, $G$ a profinite group and $D$ a continuous determinant of $G$ of dimension $n$ with values in $B$ \cite{Che14} whose reduction $D\otimes_Bk$ is the determinant of $\chi_1\oplus\dots\oplus\chi_n$, for a family of characters $\chi_i:G\to k^\times$ with extension bound $h$. Let $N\ge1$ and let $U$ be an open normal subgroup of $G$. Then there is a finite set $s\subset G$ with at most $c(n,h)$ elements such that every element of $B[G]$ is congruent to a $B$-linear combination of elements of $s$ modulo the two-sided ideal $Q$ generated by $\operatorname{CH}(D)$, $\mathfrak m_B^N$ and the elements $u-1$ with $u\in U$.
\end{lemma}
\leanref{Development/AutomorphicArithmetic/AutomorphicVarieties/CommutativeAlgebraHeight/FLevel.lean}{UpperTheory.V.CA.f_level}
\begin{proof}
Let $s$ be given by Lemma~\ref{sha-patching-chain-cayley-field} for the $\chi_i$ and $U$, and let $\varphi:B[G]\to k[G]$ be the reduction, which is surjective. The two-sided ideal generated by $\operatorname{CH}(D\otimes_Bk)$ is generated by $\varphi(\operatorname{CH}(D))$, since the formation of the Cayley-Hamilton quotient commutes with base change \cite[(1.8)]{Che14}. As $\varphi$ is surjective, every element of the two-sided ideal of $k[G]$ generated by $\varphi(\operatorname{CH}(D))$ and the $u-1$ is the image of an element of $Q$. So for $y\in B[G]$ there are $z$ in the $B$-span of $s$ and $x\in Q$ with $\varphi(y-z-x)=0$, that is, $y-z-x\in\mathfrak m_BB[G]$. Hence $B[G]$ is the sum of the span of $s$, $Q$ and $\mathfrak m_BB[G]$. Multiplying by $\mathfrak m_B^j$ and using $\mathfrak m_B^jQ\subset Q$, by induction on $j$ it is also the sum of the span of $s$, $Q$ and $\mathfrak m_B^jB[G]$ for every $j$. Finally $\mathfrak m_B^NB[G]\subset Q$.
% f_level: level_field, g_chmap for the Cayley-Hamilton ideal of the reduction, the lift of the two-sided span along phi, step, iter and last.
\end{proof}

\begin{lemma}\label{sha-patching-chain-cayley-limit}
Let $B$ be a complete local noetherian ring with finite residue field, $G$ a profinite group, $S\subset B[G]$ a subset and $c\ge0$. Suppose that for every open normal subgroup $U$ of $G$ there is a finite set $s_U\subset G$ with at most $c$ elements such that every element of $B[G]$ is congruent to a $B$-linear combination of elements of $s_U$ modulo the two-sided ideal generated by $S$, $\mathfrak m_B$ and the elements $u-1$ with $u\in U$. Then there is a finite set $s\subset G$ with at most $c$ elements such that:
\begin{enumerate}
\item every element of $B[G]$ is congruent to a $B$-linear combination of elements of $s$ modulo the closure $\bar S$ of $S$;
\item for every $N\ge1$ there is an open normal subgroup $U$ of $G$ such that $u-1\in\mathfrak m_B^NB[G]+\bar S$ for all $u\in U$.
\end{enumerate}
\end{lemma}
\leanref{Development/AutomorphicArithmetic/AutomorphicVarieties/CommutativeAlgebraHeight/FLimit.lean}{UpperTheory.V.CA.f_limit}
\begin{proof}
Write $M=B[G]$, viewed as a $B$-module, $\mathfrak m=\mathfrak m_B$, and for an open normal subgroup $U$ let $K_U$ be the two-sided ideal generated by $S$ and the $u-1$ with $u\in U$. The two-sided ideal generated by $S$, $\mathfrak m^n$ and the $u-1$ is $K_U+\mathfrak m^nM$; so $\bar S$ is the intersection of the $K_U+\mathfrak m^nM$. For $U$ and $n\ge0$ write $\langle s_U\rangle$ for the $B$-span of $s_U$ and $L^n_U=K_U+\mathfrak m^nM$.

\emph{Bounded indices.} If $P$ is a submodule with $P+\mathfrak mM=M$, then $P+\mathfrak m^nM=M$ for all $n$, by induction, as $\mathfrak m^nM=\mathfrak m^n(P+\mathfrak mM)\subset P+\mathfrak m^{n+1}M$. So $\langle s_U\rangle+L^n_U=M$, and $M/L^n_U$ is a quotient of $(B/\mathfrak m^n)^{s_U}$. As $B$ is noetherian with finite residue field, $B/\mathfrak m^n$ is finite, and $M/L^n_U$ has at most $|B/\mathfrak m^n|^c$ elements.

\emph{Stabilization.} Fix $n$. The submodules $L^n_U$ shrink as $U$ shrinks, and their indices are bounded; let $U_n$ be one of maximal index. For every $U$, $L^n_{U_n\cap U}\subset L^n_{U_n}$ with the same index, so they are equal, and $L^n_{U_n}=L^n_{U_n\cap U}\subset L^n_U$.

\emph{Approximation.} Let $s=s_{U_1}$ and $L^n=L^n_{U_n}$. Then $L^n\subset L^n_{U_m}\subset L^m$ for $m\le n$, and $L^n\subset L^{n+1}+\mathfrak m^nM$, as $L^n\subset L^n_{U_{n+1}}$. Also $\langle s\rangle+L^n=M$ for every $n$: by stabilization $K_{U_1}+\mathfrak mM\subset K_{U_n}+\mathfrak mM$, so $\langle s\rangle+K_{U_n}+\mathfrak mM=M$, and the first step applies. Hence $\mathfrak m^nM\subset\mathfrak m^n\langle s\rangle+L^{n+1}$ and $L^n\subset L^{n+1}+\mathfrak m^n\langle s\rangle$. Let $x\in L^N$. Choose recursively coefficients $a^{(n)}_t\in B$, for $t\in s$, with $a^{(n)}=0$ for $n\le N$, $a^{(n+1)}-a^{(n)}\in\mathfrak m^n$ and $x-\sum_ta^{(n)}_tt\in L^n$: if $x-\sum_ta^{(n)}_tt\in L^{n+1}$, keep the coefficients, and otherwise use $L^n\subset L^{n+1}+\mathfrak m^n\langle s\rangle$. As $B$ is complete, $a^{(n)}_t$ converges to some $a_t\in\mathfrak m^N$ with $a^{(n)}_t-a_t\in\mathfrak m^n$, and $x-\sum_ta_tt\in L^n+\mathfrak m^nM=L^n$ for every $n$.

\emph{Conclusion.} Every $y\in M$ lies in $L^0=M$, so this gives $z$ in the span of $s$ with $y-z\in L^n\subset L^n_U$ for all $n$ and $U$, that is, $y-z\in\bar S$; this is (1). For (2), let $N\ge1$ and $u\in U_N$. Then $u-1\in K_{U_N}\subset L^N$, so there is $z=\sum_ta_tt$ with $a_t\in\mathfrak m^N$ and $u-1-z\in\bar S$.
% f_limit: fl_abs with K U = the two-sided span of S and the u - 1 (restricted to B), fl_mem for the identification with the span including m^n, fl_nak, fl_card, fl_stab (Nat.sSup_mem and relIndex), fl_approx (step, the recursion f, IsPrecomplete.prec), hcl for the closure.
\end{proof}

\begin{lemma}\label{sha-patching-chain-cayley-finite}
Let $B$ be a complete local noetherian ring with finite residue field $k$, $G$ a profinite group and $D$ a continuous determinant of $G$ of dimension $n$ with values in $B$ \cite{Che14} whose reduction $D\otimes_Bk$ is the determinant of $\chi_1\oplus\dots\oplus\chi_n$, for a family of characters $\chi_i:G\to k^\times$ with extension bound $h$. Let $Z\subset B[G]$ be a subset and $E_{D,Z}$ as in Definition~\ref{sha-patching-chain-cayley-closure}. Then:
\begin{enumerate}
\item there is a finite set $s\subset G$ with at most $c(n,h)$ elements whose images span $E_{D,Z}$ as a $B$-module; in particular $E_{D,Z}$ is finite over $B$;
\item for every $N\ge1$ there is an open normal subgroup $V$ of $G$ such that the image of $v-1$ in $E_{D,Z}$ lies in $\mathfrak m_B^NE_{D,Z}$ for all $v\in V$.
\end{enumerate}
\end{lemma}
\leanref{Development/AutomorphicArithmetic/AutomorphicVarieties/CommutativeAlgebraHeight/FinE.lean}{UpperTheory.V.CA.finE}
\begin{proof}
Let $S=\operatorname{CH}(D)\cup Z$. For every open normal subgroup $U$, Lemma~\ref{sha-patching-chain-cayley-level} with $N=1$ gives a set $s_U$ with at most $c(n,h)$ elements spanning $B[G]$ modulo the two-sided ideal generated by $\operatorname{CH}(D)$, $\mathfrak m_B$ and the $u-1$, which lies in the two-sided ideal generated by $S$, $\mathfrak m_B$ and the $u-1$. Lemma~\ref{sha-patching-chain-cayley-limit}, with $c=c(n,h)$, gives $s$. By its part (1) the images of $s$ span $E_{D,Z}=B[G]/\bar S$. By its part (2), for $v$ in the open normal subgroup $V=U$ given there, $v-1$ is congruent modulo $\bar S$ to an element of $\mathfrak m_B^NB[G]$, so its image lies in $\mathfrak m_B^NE_{D,Z}$.
% finE: f_level for every U (the inclusion of spans), f_limit for D.CH union Z with c = finBound n h, then key (equality in CHQuot from the closure), hspan, Module.Finite from the image of s, and hle for the open subgroups.
\end{proof}

\begin{lemma}\label{sha-patching-chain-cayley-descent}
Let $B$ be a noetherian local ring, $G$ a profinite group, $D$ a continuous determinant of $G$ of dimension $n$ with values in $B$ \cite{Che14}, and $Z\subset B[G]$ a set of elements killed by $D$. Then:
\begin{enumerate}
\item $D$ kills every element of the closure of $Z$;
\item there is a Cayley-Hamilton determinant $D_Z$ of $E_{D,Z}$ whose composite with the quotient map $B[G]\to E_{D,Z}$ is $D$.
\end{enumerate}
\end{lemma}
\leanref{Development/AutomorphicArithmetic/AutomorphicVarieties/CommutativeAlgebraHeight/DClosed.lean}{UpperTheory.V.CA.d_closed}
\leanref{Development/AutomorphicArithmetic/AutomorphicVarieties/CommutativeAlgebraHeight/DQuotDet.lean}{UpperTheory.V.CA.d_quotDet}
\begin{proof}
(1) Let $y$ lie in the closure of $Z$, and let $N\ge1$. The determinant $D_N=D\otimes_BB/\mathfrak m_B^N$ has values in a discrete ring and is continuous, so it kills $v-1$ for all $v$ in some open normal subgroup $V$ of $G$ \cite[Lemma~2.33]{Che14}. It kills the images of the elements of $Z$, and the images of the elements of $\mathfrak m_B^N$, which are $0$. As its kernel is a two-sided ideal, it kills the image of the two-sided ideal generated by $Z$, $\mathfrak m_B^N$ and the $v-1$, which contains $y$. Now let $S=B[t_1,\dots,t_m]$ and $h\in S[G]$. The image of $D(1+hy)-1$ in $(B/\mathfrak m_B^N)[t_1,\dots,t_m]$ is $D_N(1+\bar h\bar y)-1=0$, for every $N$. So the coefficients of $D(1+hy)-1$ lie in $\bigcap_N\mathfrak m_B^N=0$ \cite[Theorem~8.10]{Mat86}. As it suffices to test polynomial rings over $B$ \cite[\S1.17, after Lemma~1.19]{Che14}, $D$ kills $y$.

(2) $D$ kills $\operatorname{CH}(D)$ \cite[Lemma~1.21]{Che14} and $Z$, so by (1) it kills the closure of their union, which is the kernel of the quotient map $q:B[G]\to E_{D,Z}$. So $D=D_Z\circ q$ for a determinant $D_Z$ of $E_{D,Z}$ \cite[Lemma~1.19(ii)]{Che14}. As $q$ is surjective, the Cayley-Hamilton ideal of $D_Z$ is the image of $\operatorname{CH}(D)$, as in the proof of Lemma~\ref{sha-patching-chain-cayley-field}, which is $0$.
% d_closed: exists_openNormalSubgroup_kills_of_groupDetLaw_isCont (the continuity of the lam_i modulo I^N), kills_map, the span hle, key2 with map_nat and IsHausdorff.haus coefficient by coefficient, d_kills_poly. d_quotDet: hle and hsub, d_ch_ker for the Cayley-Hamilton ideal, d_closed, d_factor, d_pull (the Cayley-Hamilton ideal of DE is the image of that of D).
\end{proof}

\begin{lemma}\label{sha-patching-chain-cayley-kills-trace}
Let $B$ be a complete local noetherian ring with finite residue field $k$, $G$ a profinite group and $D$ a continuous determinant of $G$ of dimension $n\ge1$ with values in $B$ \cite{Che14} whose reduction $D\otimes_Bk$ is the determinant of $\chi_1\oplus\dots\oplus\chi_n$, for a family of characters $\chi_i:G\to k^\times$ with extension bound $h$. Then $D$ kills $y\in B[G]$ if and only if $\operatorname{tr}_D(gy)=0$ for all $g\in G$.
\end{lemma}
\leanref{Development/AutomorphicArithmetic/AutomorphicVarieties/CommutativeAlgebraHeight/GKillsTr.lean}{UpperTheory.V.CA.g_kills_tr}
\leanref{Development/AutomorphicArithmetic/AutomorphicVarieties/CommutativeAlgebraHeight/FinEGma.lean}{UpperTheory.V.CA.finE_gma}
\begin{proof}
Let $E=E_{D,\emptyset}$ and $q:B[G]\to E$ the quotient map. By Lemma~\ref{sha-patching-chain-cayley-descent}(2), $D=D_E\circ q$ for a Cayley-Hamilton determinant $D_E$ of $E$. Let $\pi_\chi$ be as in Lemma~\ref{sha-patching-chain-cayley-chars} and $\rho:B[G]\to k^n$ its composite with the reduction $B[G]\to k[G]$, which is surjective. For $x\in B[G]$, the image in $k[X]$ of the characteristic polynomial of $D_E$ at $q(x)$ is that of $D\otimes_Bk$ at the reduction of $x$, which is $\prod_i(X-\rho(x)_i)$ by Lemma~\ref{sha-patching-chain-cayley-chars}. If $q(x)=0$, this polynomial is $X^n$, and evaluating at $\rho(x)_i$ gives $\rho(x)_i=0$. So $\rho=\pi\circ q$ for a surjective homomorphism $\pi:E\to k^n$, to which Lemma~\ref{sha-patching-chain-cayley-gma} applies; it gives idempotents $e_1,\dots,e_n$ of $E$ as in Lemma~\ref{sha-patching-chain-cayley-trace}. By that lemma, $D_E$ kills $q(y)$ if and only if $\operatorname{tr}_D(wy)=\operatorname{tr}_{D_E}(q(w)q(y))=0$ for all $w\in B[G]$, and as $\operatorname{tr}_D$ is $B$-linear this holds if and only if $\operatorname{tr}_D(gy)=0$ for all $g\in G$. Finally $D$ kills $y$ if and only if $D_E$ kills $q(y)$, as $D(1+hy)=D_E(1+q(h)q(y))$ and $q$ stays surjective after base change.
% g_kills_tr: finE_gma with Z empty (d_quotDet; f_chars composed with the reduction for rho; lift_surj for pi; f_gma), d_tr for the linearity of the trace, gma_kills, and d_pull for the kernel along the surjection toQuot.
\end{proof}

\begin{lemma}\label{sha-patching-chain-cayley-star}
Let $\ell$ be a prime, $\mathcal O$ a commutative $\mathbf Z_\ell$-algebra, $n\ge0$, and $\Gamma'$ a profinite group with a continuous homomorphism $\varepsilon:\Gamma'\to\mathbf Z_\ell^\times$, an element $c'$ with $c'^2=1$, and an open normal subgroup $\Gamma'_0$. For a commutative $\mathcal O$-algebra $R$ and $g\in\Gamma'_0$ let $g^*=\varepsilon(g)^{1-n}(c'gc'^{-1})^{-1}\in R[\Gamma'_0]$, and extend $g\mapsto g^*$ to an $R$-linear map $y\mapsto y^*$ of $R[\Gamma'_0]$.
\begin{enumerate}
\item The map $y\mapsto y^*$ is an anti-automorphism of $R$-algebras with $y^{**}=y$, and it commutes with the maps $R[\Gamma'_0]\to R'[\Gamma'_0]$ induced by homomorphisms of $\mathcal O$-algebras $R\to R'$.
\item Let $D$ be a determinant of $\Gamma'_0$ of dimension $n$ with values in $R$ \cite{Che14} such that the characteristic polynomials of $D$ at $g$ and at $g^*$ agree for all $g\in\Gamma'_0$. Then $D(y^*)=D(y)$ for every commutative $R$-algebra $R'$ and $y\in R'[\Gamma'_0]$; in particular the characteristic polynomials of $D$ at $y$ and at $y^*$ agree. And $y\mapsto y^*$ maps the kernel and the Cayley-Hamilton ideal of $D$ into themselves.
\item Let $R$ be local with $\ell\in\mathfrak m_R$, and $S\subset R[\Gamma'_0]$ a subset whose image under $*$ lies in the two-sided ideal generated by $S$. Then $*$ maps the closure of $S$ into itself.
\end{enumerate}
\end{lemma}
\leanref{Development/AutomorphicArithmetic/AutomorphicVarieties/PatchingModules/ClosureConjugation/StarBasic.lean}{UpperTheory.V.Pa.star_basic}
\leanref{Development/AutomorphicArithmetic/AutomorphicVarieties/PatchingModules/PatchingBaseChange/StarMap.lean}{UpperTheory.V.Pa.BC.star_map}
\leanref{Development/AutomorphicArithmetic/AutomorphicVarieties/CommutativeAlgebraHeight/DStar.lean}{UpperTheory.V.CA.d_star}
\leanref{Development/AutomorphicArithmetic/AutomorphicVarieties/PatchingModules/ClosureConjugation/StarClosure.lean}{UpperTheory.V.Pa.star_closure}
\begin{proof}
(1) For $g,g'\in\Gamma'_0$, $(gg')^*=\varepsilon(g')^{1-n}(c'g'c'^{-1})^{-1}\varepsilon(g)^{1-n}(c'gc'^{-1})^{-1}=g'^*g^*$, so $*$ is an anti-automorphism. As $\varepsilon(c'gc'^{-1})=\varepsilon(g)$ and $c'^2=1$, $(g^*)^*=\varepsilon(g)^{1-n}\varepsilon(g)^{n-1}g=g$. The compatibility with $R\to R'$ holds on group elements, where it is the compatibility of the images of $\varepsilon(g)$.

(2) Let $D^{\mathrm{op}}$ be $D$ viewed as a determinant of the opposite algebra of $R[\Gamma'_0]$. It has the same characteristic polynomials, the same kernel, by the two descriptions of the kernel in \cite[Lemma~1.19(i)]{Che14}, and the same Cayley-Hamilton ideal, as the powers of an element are the same in both algebras. The map $\sigma:y\mapsto y^*$ is an isomorphism of $R$-algebras from $R[\Gamma'_0]$ to its opposite algebra, and $D^{\mathrm{op}}\circ\sigma$ is a determinant of $R[\Gamma'_0]$ whose characteristic polynomial at each $g\in\Gamma'_0$ is that of $D$ at $g^*$, hence that of $D$ at $g$. The group elements form a multiplicative monoid that spans $R[\Gamma'_0]$, so $D^{\mathrm{op}}\circ\sigma=D$ by Amitsur's formula \cite[Lemma~1.12(ii)]{Che14}. This is the first statement. As $\sigma$ is an isomorphism, it carries the kernel and the Cayley-Hamilton ideal of $D=D^{\mathrm{op}}\circ\sigma$ to those of $D^{\mathrm{op}}$, which are those of $D$.

(3) Let $y$ lie in the closure of $S$, $N\ge1$ and $V$ an open normal subgroup of $\Gamma'_0$. As conjugation by $c'$ and $\varepsilon$ are continuous, there is an open normal subgroup $V'\subset V$ with $(c'vc'^{-1})^{-1}\in V$ and $\varepsilon(v)\equiv1\pmod{\ell^N}$ for $v\in V'$. Let $T$ be the two-sided ideal generated by $S$, $\mathfrak m_R^N$ and the $v-1$ with $v\in V$, and $T'$ the one generated by $S$, $\mathfrak m_R^N$ and the $v-1$ with $v\in V'$, so that $y\in T'$. The image under $*$ of each generator of $T'$ lies in $T$: for $s\in S$ by hypothesis; for $a\in\mathfrak m_R^N$, as $a^*=a$; and for $v\in V'$, $v^*-1=w(u-1)+(w-1)$ with $w=(c'vc'^{-1})^{-1}\in V$ and $u=\varepsilon(v)^{1-n}$, where $u-1\in\ell^NR\subset\mathfrak m_R^N$. As $*$ is an anti-automorphism, it maps $T'$ into $T$, so $y^*\in T$.
% star_basic (star_star_monoidAlgebra, starOp a monoid hom to the opposite), BC.star_map (pushCh of eps and the formula on single). pair_kills: d_star with hlam from Pair.P (d_ext, Amitsur's formula through d_poly_ext; d_op for chi, ker and CH of the opposite; d_pull), pk_surj, pk_bc for IsPolar; trinv's hinv for the characteristic polynomials. star_closure: exists_unit_valuation_le_subset_nhds for H', heps, hsplit, TwoSidedIdeal.span_induction.
\end{proof}

\begin{lemma}\label{sha-patching-chain-cayley-kills}
Let $x$ be a normalized arc over $A=\kappa[[T]]$, with blocks of sizes $a$ and $n-a$, and let $(q,\phi_j,\tau_j,U_{j,i},y_{j,i},s,f)$ be an arc datum for $x$ over $(\mathcal R,D,\psi)$, with relation ideal $I$, in the sense of Definition~\ref{sha-patching-chain-datum}. Let $S\subset\mathcal R[\Gamma_{\infty,0}]$ be the set of the elements $\prod_i(g_i-\psi_{v,i}(g_i))$ with $v\mid\ell$ and $g_i\in D_v$, the elements $g-1$ with $g$ in the image of some $I'_u$, the elements $\prod_i(\phi_j-U_{j,i})$ and $\prod_{k\ne i}(\phi_j-U_{j,k})(\tau_j-1-y_{j,i})$, and their images under $*$. Then:
\begin{enumerate}
\item the characters $\bar\chi_i\circ p$ of $\Gamma_{\infty,0}$ have an extension bound (Definition~\ref{sha-patching-chain-cayley-extension-bound}), and $D$ kills $y\in\mathcal R[\Gamma_{\infty,0}]$ if and only if $\operatorname{tr}_D(gy)=0$ for all $g\in\Gamma_{\infty,0}$;
\item $S^*=S$, and $D$ kills every element of $S$;
\item $I$ is the closure of the union of $S$ and $\operatorname{CH}(D)$, so that $\mathcal R[\Gamma_{\infty,0}]/I=E_{D,S}$ (Definition~\ref{sha-patching-chain-cayley-closure}).
\end{enumerate}
\end{lemma}
\leanref{Development/AutomorphicArithmetic/AutomorphicVarieties/PatchingModules/ClosureConjugation/GensKill.lean}{UpperTheory.V.Pa.gens_kill}
\leanref{Development/AutomorphicArithmetic/AutomorphicVarieties/CommutativeAlgebraHeight/PairKills.lean}{UpperTheory.V.CA.pair_kills}
\leanref{Development/AutomorphicArithmetic/AutomorphicVarieties/PatchingModules/ClosureConjugation/PairMf.lean}{UpperTheory.V.Pa.pair_mf}
\begin{proof}
(1) The group $\Gamma_\infty$ is topologically finitely generated, so by Lemma~\ref{sha-patching-chain-cayley-bound}(1) every open subgroup of $\Gamma_\infty$, and so every open subgroup of $\Gamma_{\infty,0}$, has finitely many continuous homomorphisms to $\mathbf Z/\ell$. As $p$ maps $\Gamma_{\infty,0}$ onto $\Gamma_0$, the $\bar\chi_i\circ p$ are pairwise distinct, with open kernels; and $k$ is finite of characteristic $\ell$. So Lemma~\ref{sha-patching-chain-cayley-bound}(2) gives an extension bound. The ring $\mathcal R$ is complete local noetherian with finite residue field $k$, $D$ is continuous with reduction the determinant of $\bigoplus_i\bar\chi_i\circ p$, and $n\ge a\ge1$; so the second statement is Lemma~\ref{sha-patching-chain-cayley-kills-trace}.

(2) As $y^{**}=y$ by Lemma~\ref{sha-patching-chain-cayley-star}(1), $S^*=S$. By (1), $D$ kills $\prod_i(g_i-\psi_{v,i}(g_i))$ and $g-1$, for $g$ in the image of $I'_u$, by the conditions (2) and (3) of an ordinary pair, and it kills $\prod_i(\phi_j-U_{j,i})$ and $\prod_{k\ne i}(\phi_j-U_{j,k})(\tau_j-1-y_{j,i})$ by condition (c) of an arc datum. By condition (1) of an ordinary pair, Lemma~\ref{sha-patching-chain-cayley-star}(2) applies to $D$, with $\Gamma'=\Gamma_\infty$, $c'=c_\infty$ and $\Gamma'_0=\Gamma_{\infty,0}$; so the kernel of $D$ is stable under $*$, and $D$ kills the images of these elements under $*$.

(3) Fix $N$ and $V$. Every generator of the ideal of Definition~\ref{sha-patching-chain-datum}(2) for $N$ and $V$ that is not in $\operatorname{CH}(D)\cup S$, $\mathfrak m_{\mathcal R}^N$ or $\{v-1:v\in V\}$ is the image under $*$ of an element of $\operatorname{CH}(D)$, of $\mathfrak m_{\mathcal R}^N$ or of that set. The first lies in $\operatorname{CH}(D)$, by Lemma~\ref{sha-patching-chain-cayley-star}(2), and the second is fixed by $*$. By the proof of Lemma~\ref{sha-patching-chain-cayley-star}(3), there is an open normal subgroup $V'\subset V$ such that $(v-1)^*$ lies in the two-sided ideal generated by $\mathfrak m_{\mathcal R}^N$ and the $w-1$, $w\in V$, for $v\in V'$. So the ideal of Definition~\ref{sha-patching-chain-datum}(2) for $N$ and $V'$ lies in the two-sided ideal generated by $\operatorname{CH}(D)\cup S$, $\mathfrak m_{\mathcal R}^N$ and the $v-1$ with $v\in V$, which lies in the ideal of the definition for $N$ and $V$. So the two intersections agree.
% gens_kill: phi_fg from P.fg, pair_kills (hks for the images under star, hoEl from Pair.Ord, hIp from Pair.I, both through g_kills_tr), pair_mf (exists_multFree_of_phiCond_of_injective_chars), g_kills_tr for P.q1 and P.q2, and star_basic (hss) for the stability of P.gens = pairGens union qGens. The Lean's relation ideal OoPt.ker is chBar of P.gens, which is (3).
\end{proof}

\begin{lemma}\label{sha-patching-chain-cayley}
Let $x$ be a normalized arc over $A=\kappa[[T]]$, with representation $\rho$ and blocks of sizes $a$ and $n-a$, and let $(q,\phi_j,\tau_j,U_{j,i},y_{j,i},s,f)$ be an arc datum for $x$ over $(\mathcal R,D,\psi)$, with relation ideal $I$, in the sense of Definition~\ref{sha-patching-chain-datum}. Let $\mathcal K$ be the fraction field of $A$, $\hat{\mathcal R}$ a local $\mathcal R$-algebra and $\hat f:\hat{\mathcal R}\to\mathcal K$ a homomorphism whose composite with $\mathcal R\to\hat{\mathcal R}$ is $f$; view $\mathcal K$ as an $\hat{\mathcal R}$-algebra through $\hat f$. Then:
\begin{enumerate}
\item $f$ is surjective; and if $\ker\hat f=\mathfrak m_{\hat{\mathcal R}}$, then $\hat f$ is surjective.
\item There are an $\hat{\mathcal R}$-algebra $E$, finite as an $\hat{\mathcal R}$-module; a surjective homomorphism of $\hat{\mathcal R}$-algebras $\iota:\hat{\mathcal R}[\Gamma_{\infty,0}]\to E$ whose kernel is the two-sided ideal generated by the image of $I$; a Cayley-Hamilton determinant $D_E$ of $E$ of dimension $n$ with $D_E\circ\iota=D\otimes\hat{\mathcal R}$ \cite{Che14}; an $\hat{\mathcal R}$-linear anti-automorphism $y\mapsto y^*$ of $E$ with $y^{**}=y$, $\operatorname{tr}_{D_E}(y^*)=\operatorname{tr}_{D_E}(y)$ for $y\in E$ and $\iota(g)^*=\iota(g^*)$ for $g\in\Gamma_{\infty,0}$; a homomorphism of $\hat{\mathcal R}$-algebras $r:E\to M_n(\mathcal K)$ with $r(\iota(g))=\rho(p(g))$; and a finitely generated $\mathcal R$-submodule $P$ of $E$ containing the $\iota(g)$ such that for every $j\ge0$ there is an open normal subgroup $V$ of $\Gamma_{\infty,0}$ with $\iota(gv)-\iota(g)\in\mathfrak m_{\mathcal R}^jP$ for all $g\in\Gamma_{\infty,0}$ and $v\in V$.
\item For $E$, $\iota$, $D_E$, $*$ and $r$ as in (2) and $y\in E$: $r(y^*)={}^tr(y)$; and the image under $\hat f$ of the characteristic polynomial of $D_E$ at $y$ is the characteristic polynomial of $r(y)$. If moreover $\hat f$ is surjective, then $r(E)$ is the set of block diagonal matrices over $\mathcal K$ with blocks of sizes $a$ and $n-a$.
\end{enumerate}
\end{lemma}
\leanref{Development/AutomorphicArithmetic/AutomorphicVarieties/PatchingModules/ClosureConjugation/GrAlg.lean}{UpperTheory.V.Pa.Hat}
\leanref{Development/AutomorphicArithmetic/AutomorphicVarieties/PatchingModules/ClosureConjugation/HatEx.lean}{UpperTheory.V.Pa.hat_ex}
\leanref{Development/AutomorphicArithmetic/AutomorphicVarieties/PatchingModules/ClosureConjugation/FSurj.lean}{UpperTheory.V.Pa.f_surj}
\leanref{Development/AutomorphicArithmetic/AutomorphicVarieties/PatchingModules/ClosureConjugation/RFSurj.lean}{UpperTheory.V.Pa.rF_surj}
\leanref{Development/AutomorphicArithmetic/AutomorphicVarieties/PatchingModules/ClosureConjugation/HatSt.lean}{UpperTheory.V.Pa.hat_st}
\leanref{Development/AutomorphicArithmetic/AutomorphicVarieties/PatchingModules/ClosureConjugation/HatChi.lean}{UpperTheory.V.Pa.hat_chi}
\leanref{Development/AutomorphicArithmetic/AutomorphicVarieties/PatchingModules/ClosureConjugation/HatRange.lean}{UpperTheory.V.Pa.hat_range}
\begin{proof}
(1) Since $\mathcal O$ maps onto the residue fields of $\mathcal R$ and of $A$ and $f$ is a homomorphism of $\mathcal O$-algebras, an element of $\mathcal R$ and its image under $f$ have the same residue in $\kappa$. As $f(s)=T$ has residue $0$, $s\in\mathfrak m_{\mathcal R}$. So $f(\mathfrak m_{\mathcal R})$ contains $T$, which generates $\mathfrak m_A$, and $f$ is surjective, since $\mathcal R$ and $A$ are complete \cite[Theorem~8.4]{Mat86}. Now suppose $\ker\hat f=\mathfrak m_{\hat{\mathcal R}}$, and let $z=a/b\in\mathcal K$ with $a,b\in A$ and $b\ne0$. Choose $a',b'\in\mathcal R$ with $f(a')=a$ and $f(b')=b$. The image of $b'$ in $\hat{\mathcal R}$ is a unit, since $\hat f$ carries it to $b\ne0$. So $z$ is the image under $\hat f$ of $a'b'^{-1}$.

(2) Let $E_0$ be the quotient of $\mathcal R[\Gamma_{\infty,0}]$ by $I$, and $S$ the set of Lemma~\ref{sha-patching-chain-cayley-kills}. By that lemma, $D$ kills every element of $S$, $S^*=S$, and $E_0=E_{D,S}$. So $D$ induces a Cayley-Hamilton determinant $D_0$ of $E_0$ whose composite with the quotient map is $D$, by Lemma~\ref{sha-patching-chain-cayley-descent}(2). By (1) of the definition of an ordinary pair, Lemma~\ref{sha-patching-chain-cayley-star}(2) applies to $D$, with $\Gamma'=\Gamma_\infty$, $c'=c_\infty$ and $\Gamma'_0=\Gamma_{\infty,0}$: the characteristic polynomials of $D$ at $y$ and at $y^*$ agree, and $*$ maps $\operatorname{CH}(D)$ into itself. As $S^*=S$ and $\ell\in\mathfrak m_{\mathcal R}$, Lemma~\ref{sha-patching-chain-cayley-star}(3) shows that $*$ maps $I$ into itself. So $*$ induces an $\mathcal R$-linear anti-automorphism of $E_0$, with $y^{**}=y$ by Lemma~\ref{sha-patching-chain-cayley-star}(1), and $\operatorname{tr}_{D_0}(y^*)=\operatorname{tr}_{D_0}(y)$. By Lemma~\ref{sha-patching-chain-cayley-kills}(1), the characters $\bar\chi_i\circ p$, whose sum has determinant the reduction of $D$, have an extension bound. As $D$ is continuous and $\mathcal R$ is complete local noetherian with finite residue field $k$, Lemma~\ref{sha-patching-chain-cayley-finite} shows that $E_0$ is finite over $\mathcal R$, and that for every $j$ there is an open normal subgroup $V$ of $\Gamma_{\infty,0}$ such that the image of $v-1$ in $E_0$ lies in $\mathfrak m_{\mathcal R}^jE_0$ for $v\in V$. The homomorphism $\mathcal R[\Gamma_{\infty,0}]\to M_n(A)$ that sends $g$ to $\rho(p(g))$, with $\mathcal R$ acting through $f$, kills $I$ by Lemma~\ref{sha-patching-chain-cayley-factor}, so it factors through $E_0$.

Let $E=\hat{\mathcal R}\otimes_{\mathcal R}E_0$, with $\iota$ the map induced by the quotient map, $D_E$ and $*$ the base changes of $D_0$ and of the anti-automorphism of $E_0$, $r$ the $\hat{\mathcal R}$-linear extension of the map $E_0\to M_n(A)\subset M_n(\mathcal K)$, and $P$ the image of $E_0$ under $y\mapsto1\otimes y$. Then $E$ is finite over $\hat{\mathcal R}$, $\iota$ is surjective, its kernel is the two-sided ideal generated by the image of $I$, $D_E$ is a Cayley-Hamilton determinant with $D_E\circ\iota=D\otimes\hat{\mathcal R}$, and $\operatorname{tr}_{D_E}$ is the base change of $\operatorname{tr}_{D_0}$. For $g\in\Gamma_{\infty,0}$ the element $g^*$ is $\varepsilon(g)^{1-n}$ times a group element, so $\iota(g)^*=\iota(g^*)$; and $r(\iota(g))=\rho(p(g))$. Finally, $\iota(gv)-\iota(g)$ is the image under $y\mapsto1\otimes y$ of the image of $g(v-1)$ in $E_0$, which lies in $\mathfrak m_{\mathcal R}^jE_0$ for $v\in V$; so it lies in $\mathfrak m_{\mathcal R}^jP$.
% hat_ex builds E = Bh (x)_B u.D.CHQuot P.gens. gens_kill is Lemma sha-patching-chain-cayley-kills; trinv (d_quotDet, d_star with hop and hCHs, star_closure, star_basic for st_st, d_tr and hinv for the trace) is Lemmas sha-patching-chain-cayley-descent and sha-patching-chain-cayley-star; pair_mf (exists_multFree_of_phiCond_of_injective_chars under PhiCond from phi_fg) with finE (Module.Finite and the open subgroups) is Lemmas sha-patching-chain-cayley-bound and sha-patching-chain-cayley-finite. The factorization rep_quot is Lemma sha-patching-chain-cayley-factor. The base change is tens (the kernel of iota) with trinv_baseChange, d_chS for IsCH, lift_ex for r and tr_bc for the trace; the continuity is the last bullet of hat_ex.

(3) Let $g\in\Gamma_{\infty,0}$. The image of $\varepsilon(g)\in\mathbf Z_\ell^\times$ in $\mathcal R$ is carried by $f$ to its image in $A$, since $A$ has characteristic $\ell$ and both images are those of an integer congruent to $\varepsilon(g)$ modulo $\ell$; and $p(c_\infty gc_\infty^{-1})=c\,p(g)\,c^{-1}$. As $x$ is normalized, ${}^t\rho(p(g))\,\rho(c\,p(g)\,c^{-1})=\varepsilon(g)^{1-n}$, so ${}^t\rho(p(g))=\varepsilon(g)^{1-n}\rho\big(p((c_\infty gc_\infty^{-1})^{-1})\big)=r(\iota(g^*))=r(\iota(g)^*)$. Both sides of $r(y^*)={}^tr(y)$ are $\hat{\mathcal R}$-linear in $y$ and $\iota$ is surjective, so this identity holds on $E$. For $y=\iota(z)$ with $z\in\hat{\mathcal R}[\Gamma_{\infty,0}]$, the characteristic polynomial of $D_E$ at $y$ is that of $D\otimes\hat{\mathcal R}$ at $z$. Its image under $\hat f$ is the characteristic polynomial, at the image of $z$ in $\mathcal K[\Gamma_{\infty,0}]$, of $D$ mapped to $\mathcal K$ through $f$. As $f$ carries $D$ to the determinant of $\rho\circ p$, this is the characteristic polynomial of $r(y)$.

Every $r(\iota(g))=\rho(p(g))$ is block diagonal, as $x$ is normalized; so $r(E)$, which is spanned over $\hat{\mathcal R}$ by these matrices, consists of block diagonal matrices. Conversely, suppose that $\hat f$ is surjective. Then $r(E)$ is a $\mathcal K$-subspace of $M_n(\mathcal K)$, and it contains $\rho(\Gamma')$, as $p$ maps $\Gamma_{\infty,0}$ onto $\Gamma_0$. For a matrix unit $e_{ij}$ with $i$ and $j$ in the same block, some $T^se_{ij}$ lies in the $A$-span of $\rho(\Gamma')$, as $x$ is normalized; since $T$ is invertible in $\mathcal K$, $e_{ij}\in r(E)$. So $r(E)$ contains every block diagonal matrix.
% (1) f_surj (co_surj, with s in the maximal ideal from res_map) and rF_surj (IsLocalization.mk'_surjective). (3) hat_st with eps_eq and x.form; hat_chi (characteristic_polynomial_compatibility, from H.pull and P.maps); hat_range (hr_span from x.span, and x.blk for the converse).
\end{proof}

\begin{lemma}\label{sha-patching-chain-idempotents-lift}
Let $B$ be a commutative ring, $I\subset B$ an ideal, and $E$ a $B$-algebra that is $I$-adically complete and separated as a $B$-module. Let $z\in E$ and $M\ge0$ with $(z^2-z)^M\in IE$. Then there is an idempotent $f\in E$ such that for every $N\ge0$ there is a polynomial $q\in B[X]$ with $f-\big(z+(z^2-z)q(z)\big)\in I^NE$.
\end{lemma}
\leanref{Development/AutomorphicArithmetic/AutomorphicVarieties/AlgebraTools/IdemLiftComplete.lean}{UpperTheory.V.AlgebraTools.idem_lift_complete}
\begin{proof}
Let $p_0=X$ and $p_{k+1}=3p_k^2-2p_k^3$ in $B[X]$, and write $\delta=X^2-X$. For any polynomial $P$, $(3P^2-2P^3)^2-(3P^2-2P^3)=(P^2-P)^2\big(4(P^2-P)-3\big)$ and $(3P^2-2P^3)-P=-(P^2-P)(2P-1)$. By induction on $k$, $p_k\in X+\delta B[X]$, $p_k^2-p_k\in\delta^{2^k}B[X]$ and $p_{k+1}-p_k\in\delta^{2^k}B[X]$; so $p_{k+j}-p_k\in\delta^{2^k}B[X]$ for all $j$. As $IE$ is a two-sided ideal and $(z^2-z)^M\in IE$, $(z^2-z)^cy\in I^mE$ for $y\in E$ whenever $c\ge Mm$; and $2^{Mm}\ge Mm$. Put $f_m=p_{Mm}(z)$. Then $f_m-f_{m'}\in I^mE$ for $m\le m'$, and $f_m^2-f_m\in I^mE$. As $E$ is complete, the $f_m$ converge to some $f$ with $f-f_m\in I^mE$ for all $m$. Then $f^2-f$ is congruent to $f_m^2-f_m$ modulo $I^mE$, so it lies in every $I^mE$, and $f^2=f$ as $E$ is separated. Finally $p_{MN}=X+\delta q$ for some $q$, so $f-(z+(z^2-z)q(z))=f-f_N\in I^NE$.
% idem_lift_complete: ilcP for p_k, PoincareLelongAlt (p_k in X + delta B[X]), ProjectiveLineTwo (p_k^2 - p_k), PL3 and PL4 (the differences), Mem, hcauchy, IsPrecomplete.prec for the limit L, hsq and hidem (IsHausdorff).
\end{proof}

\begin{lemma}\label{sha-patching-chain-idempotents-orthogonal}
Let $B$ be a commutative ring in which $2$ is a unit, $I\subset B$ an ideal, and $E$ a $B$-algebra that is $I$-adically complete and separated as a $B$-module, with a $B$-linear anti-automorphism $y\mapsto y^*$ such that $y^{**}=y$. Let $F$ be a commutative $B$-algebra in which $I$ maps to $0$, let $m\ge1$, and let $r:E\to M_m(F)$ be a homomorphism of $B$-algebras such that $r(y^*)={}^tr(y)$ for $y\in E$, every $y$ with $r(y)=0$ has a power in $IE$, and every matrix unit $e_{ii}$ lies in $r(E)$. Then there are orthogonal idempotents $e_1,\dots,e_m$ of $E$ with $\sum_ie_i=1$, $e_i^*=e_i$ and $r(e_i)=e_{ii}$.
\end{lemma}
\leanref{Development/AutomorphicArithmetic/AutomorphicVarieties/AlgebraTools/StCoiLift.lean}{UpperTheory.V.AlgebraTools.st_coi_lift}
\begin{proof}
As $IE$ is a two-sided ideal and $I$ maps to $0$ in $F$, $r$ kills $IE$; and $1^*=1$, as $1=1^{**}=(1^*\cdot1)^*=1^*\cdot1^{**}=1^*$.

\emph{One lift.} Let $e\in E$ be an idempotent with $e^*=e$, and $y\in E$ with $r(y)$ idempotent, ${}^tr(y)=r(y)$ and $r(e)r(y)r(e)=r(y)$. Then there is an idempotent $f$ with $f^*=f$, $ef=f=fe$ and $r(f)=r(y)$. Indeed, let $u=\frac12(y+y^*)$ and $z=eue$. Then $u^*=u$, so $z^*=e^*u^*e^*=z$; $r(u)=\frac12(r(y)+{}^tr(y))=r(y)$, so $r(z)=r(e)r(y)r(e)=r(y)$; and $ez=z=ze$. As $r(z^2-z)=r(y)^2-r(y)=0$, some power of $z^2-z$ lies in $IE$, and Lemma~\ref{sha-patching-chain-idempotents-lift} gives an idempotent $f$ and polynomials $q_N$ with $f-w_N\in I^NE$, where $w_N=z+(z^2-z)q_N(z)$. Each $w_N$ is a polynomial in $z$ with coefficients in $B$, so $w_N^*=w_N$, as $z^*=z$ and $*$ reverses products; and $w_N=z(1+(z-1)q_N(z))=(1+(z-1)q_N(z))z$, so $ew_N=w_N=w_Ne$. For each $N$, $f^*-f=(f-w_N)^*-(f-w_N)$, $ef-f=e(f-w_N)-(f-w_N)$ and $fe-f=(f-w_N)e-(f-w_N)$ lie in $I^NE$, since $*$ is $B$-linear; so they vanish, as $E$ is separated. Finally $r(f)=r(w_1)=r(z)+r(z^2-z)r(q_1(z))=r(y)$.

\emph{All the lifts.} Suppose that orthogonal idempotents $e_i$ with $e_i^*=e_i$ and $r(e_i)=e_{ii}$ are chosen for $i$ in a set $S$, and let $a\notin S$. Then $e=1-\sum_{i\in S}e_i$ is an idempotent with $e^*=e$, $ee_i=0=e_ie$ for $i\in S$, and $r(e)=1-\sum_{i\in S}e_{ii}$, so $r(e)e_{aa}r(e)=e_{aa}$. Choose $y$ with $r(y)=e_{aa}$, which is idempotent and symmetric, and let $f$ be as in the first step; put $e_a=f$. Then $e_ae_i=fe\,e_i=0$ and $e_ie_a=e_ie\,f=0$ for $i\in S$. Continuing gives $e_1,\dots,e_m$. Then $g=1-\sum_ie_i$ is an idempotent with $r(g)=1-\sum_ie_{ii}=0$, so $g^M\in IE$ for some $M$, that is $g\in IE$. Then $g=g^N\in I^NE$ for all $N$, so $g=0$.
% st_coi_lift: scl_step for one lift (u, z with hsz, hrz, hez, hze; idem_lift_complete; hfix through scl_fix for st, mulLeft e and mulRight e; hr0 for r f), the induction key over Finset.insert, and scl_idem_zero for 1 - sum. scl_st_one, scl_st_pow and scl_st_aeval for the stability of polynomials in z under st.
\end{proof}

\begin{lemma}\label{sha-patching-chain-idempotents-corner}
Let $B$ be a commutative ring, $E$ a $B$-algebra with a Cayley-Hamilton determinant $D$ of dimension $n\ge1$ \cite{Che14}, and $e\in E$ an idempotent with $D(1+Xe)=1+X$ in $B[X]$. Then $\operatorname{tr}_D(e)=1$ and $eye=\operatorname{tr}_D(ey)e$ for every $y\in E$.
\end{lemma}
\leanref{Development/AutomorphicArithmetic/AutomorphicVarieties/CommutativeAlgebraHeight/DCorner.lean}{UpperTheory.V.CA.d_corner}
\begin{proof}
For a commutative $B$-algebra $S$, write $D_S$ for the multiplicative map $S\otimes_BE\to S$ of $D$, and $e$, $y$ for $1\otimes e$, $1\otimes y$; recall that $D_{S[X]}(1+Xu)=1+\operatorname{tr}_D(u)X+\cdots$ for $u\in E$ \cite{Che14}.

\emph{(i)} For $b\in S$, $D_S(1+be)=1+b$, by specializing $X$ to $b$.

\emph{(ii) The corner functional.} For $w\in S\otimes_BE$ put $c_S(w)=D_S(1-e+ewe)$; it is compatible with base change in $S$. For $b\in S$, $1-e+e(bw)e=(1+(b-1)e)(1-e+ewe)$, so $c_S(bw)=bc_S(w)$ by (i); and $c_S(be)=b\,D_S(1)=b$. It is additive. Indeed, let $S'=S\oplus S\varepsilon_1\oplus S\varepsilon_2$ and $S''=S\oplus S\varepsilon$ with all products of the $\varepsilon$'s zero. For $\lambda_1,\lambda_2\in S$, the map $S'\to S''$ with $\varepsilon_k\mapsto\lambda_k\varepsilon$ carries $\varepsilon_1w+\varepsilon_2w'$ to $\varepsilon(\lambda_1w+\lambda_2w')$, and $c_{S''}(\varepsilon v)=\varepsilon c_S(v)$ for $v\in S\otimes_BE$ by homogeneity and base change. Write $c_{S'}(\varepsilon_1w+\varepsilon_2w')=s_0+s_1\varepsilon_1+s_2\varepsilon_2$; then $\lambda_1s_1+\lambda_2s_2=c_S(\lambda_1w+\lambda_2w')$. Taking $(\lambda_1,\lambda_2)$ to be $(1,0)$, $(0,1)$ and $(1,1)$ gives $c_S(w+w')=c_S(w)+c_S(w')$.

\emph{(iii) Traces of corner elements.} Let $u\in E$ with $eue=u$. Then $1+Xu=1-e+e(e+Xu)e$, so $D_{B[X]}(1+Xu)=c_{B[X]}(e+Xu)=1+Xc_B(u)$ by (ii), and $\operatorname{tr}_D(u)=c_B(u)$. With $u=e$, $\operatorname{tr}_D(e)=c_B(e)=1$.

\emph{(iv) A characteristic polynomial.} Since $(X-Xe+e)(1+(X-1)e)=X$, (i) and homogeneity give $D_{S[X]}(X-Xe+e)\cdot X=X^n$, so $D_{S[X]}(X-Xe+e)=X^{n-1}$. For $u$ as in (iii) and $s\in S$, $X-e-su=(X-Xe+e)\big(1-e+e((X-1)e-su)e\big)$, so by (ii) and (iii) the characteristic polynomial of $D$ at $e+su$ is $X^{n-1}\big(X-1-s\operatorname{tr}_D(u)\big)$.

\emph{(v) The corner is a line.} Take $S=B[\varepsilon]$ with $\varepsilon^2=0$, $s=\varepsilon$ and $y'=e+\varepsilon u$. By the Cayley-Hamilton property $y'^{n-1}(y'-1-\varepsilon\operatorname{tr}_D(u))=0$. Now $(y'-1-\varepsilon\operatorname{tr}_D(u))e=\varepsilon(u-\operatorname{tr}_D(u)e)$, as $ue=u$, and $y'\varepsilon(u-\operatorname{tr}_D(u)e)=\varepsilon(u-\operatorname{tr}_D(u)e)$, as $eu=u$ and $\varepsilon^2=0$. Multiplying the relation on the right by $e$ gives $\varepsilon(u-\operatorname{tr}_D(u)e)=0$ in $B[\varepsilon]\otimes_BE=E\oplus\varepsilon E$, so $u=\operatorname{tr}_D(u)e$. For $y\in E$, $u=eye$ satisfies $eue=u$, and $\operatorname{tr}_D(eye)=\operatorname{tr}_D(e\cdot ey)=\operatorname{tr}_D(ey)$; so $eye=\operatorname{tr}_D(ey)e$.
% d_corner: one_add_e (i); cornerF, cornerF_nat, cornerF_hom, cornerF_e and cornerF_add (ii), with add_of_homog (the dual numbers TrivSqZeroExt B (B x B) and TrivSqZeroExt B B); cornerF_tr (iii, through d_chi); compl_e and chiS_corner (iv); corner_eq (v, with d_chS for IsCH, h2 and h3); d_tr for tr(exe) = tr(ex).
\end{proof}

\begin{lemma}\label{sha-patching-chain-idempotents}
Let $B$ be a complete local noetherian ring in which $2$ is a unit, and $E$ a $B$-algebra, finite as a $B$-module, with a Cayley-Hamilton determinant $D$ of dimension $n$ \cite{Che14} and a $B$-linear anti-automorphism $y\mapsto y^*$ with $y^{**}=y$ and $\operatorname{tr}_D(y^*)=\operatorname{tr}_D(y)$. Let $F$ be a field with a homomorphism $B\to F$ whose kernel is $\mathfrak m_B$, let $1\le a\le n$, and let $r:E\to M_n(F)$ be a homomorphism of $B$-algebras such that, for $y\in E$, $r(y^*)={}^tr(y)$ and the characteristic polynomial of $r(y)$ is the image of that of $D$ at $y$, and such that $r(E)$ contains the block diagonal matrices with blocks of sizes $a$ and $n-a$. For $1\le i\le n$ let $o(i)=1$ if $i\le a$ and $o(i)=n$ if $i>a$. Then there are $e_1,\dots,e_n$ and $x_1,\dots,x_n$ in $E$ such that:
\begin{enumerate}
\item the $e_i$ are orthogonal idempotents with $\sum_ie_i=1$, $e_i^*=e_i$, $\operatorname{tr}_D(e_i)=1$, $e_iye_i\in Be_i$ for $y\in E$, and $D(1+Xe_i)=1+X$ in $B[X]$;
\item $e_{o(i)}x_ie_i=x_i$, $x_{o(i)}=e_{o(i)}$, $x_ix_i^*=e_{o(i)}$ and $x_i^*x_i=e_i$;
\item $r(e_i)=e_{ii}$ and $r(x_i)=e_{o(i)i}$, where the $e_{ij}$ are the matrix units.
\end{enumerate}
\end{lemma}
\leanref{Development/AutomorphicArithmetic/AutomorphicVarieties/PatchingModules/ClosureConjugation/OBLt.lean}{UpperTheory.V.Pa.idem}
\leanref{Development/AutomorphicArithmetic/AutomorphicVarieties/PatchingModules/ClosureConjugation/IdemNil.lean}{UpperTheory.V.Pa.idem_nil}
\leanref{Development/AutomorphicArithmetic/AutomorphicVarieties/AlgebraTools/InvAlg.lean}{UpperTheory.V.AlgebraTools.TrInv.Units}
\leanref{Development/AutomorphicArithmetic/AutomorphicVarieties/AlgebraTools/UnitsExist.lean}{UpperTheory.V.AlgebraTools.units_exist}
\begin{proof}
\emph{The kernel of $r$.} Let $y\in E$ with $r(y)=0$. The characteristic polynomial $\chi$ of $D$ at $y$ reduces to that of the zero matrix, so the coefficients of $X^n-\chi$ lie in $\mathfrak m_B$. As $D$ is Cayley-Hamilton, $\chi(y)=0$, so $y^n=y^n-\chi(y)$ lies in $\mathfrak m_BE$.

\emph{Rank-one idempotents.} Let $e\in E$ be an idempotent with $r(e)=e_{ii}$. Then $D(1+Xe)=(1+X)^k$ for some $k$. Its image in $F[X]$ is computed from the characteristic polynomial $X^n-X^{n-1}$ of $e_{ii}$, and is $1+X$; comparing degrees, $k=1$.

\emph{Idempotents.} As $E$ is finite over the complete noetherian ring $B$, it is $\mathfrak m_B$-adically complete and separated \cite[\S8]{Mat86}. By the first step every element of the kernel of $r$ has a power in $\mathfrak m_BE$, and the matrix units $e_{ii}$ lie in $r(E)$, as they are block diagonal. So Lemma~\ref{sha-patching-chain-idempotents-orthogonal}, with $I=\mathfrak m_B$, gives orthogonal idempotents $e_1,\dots,e_n$ with $\sum_ie_i=1$, $e_i^*=e_i$ and $r(e_i)=e_{ii}$. By the second step $D(1+Xe_i)=1+X$; and as $D$ is Cayley-Hamilton and $n\ge1$, Lemma~\ref{sha-patching-chain-idempotents-corner} gives $e_iye_i=\operatorname{tr}_D(e_iy)e_i$ for $y\in E$ and $\operatorname{tr}_D(e_i)=1$.

\emph{The elements $x_i$.} Choose $y_i\in E$ with $r(y_i)=e_{o(i)i}$, and put $z_i=e_{o(i)}y_ie_i$; then $r(z_i)=e_{o(i)i}$. By (1), $z_i^*z_i=b_ie_i$ and $z_iz_i^*=b'_ie_{o(i)}$ for some $b_i,b'_i\in B$. Taking traces, $b'_i=\operatorname{tr}_D(z_iz_i^*)=\operatorname{tr}_D(z_i^*z_i)=b_i$. Applying $r$, $e_{ii}=b_ie_{ii}$ in $M_n(F)$, so $b_i$ maps to $1$ in $F$. As $2$ is a unit and $B$ is complete, $b_i=\alpha_i^2$ for a unit $\alpha_i\in B$ with $\alpha_i\equiv1$ modulo $\mathfrak m_B$. Put $x_i=e_i$ if $o(i)=i$ and $x_i=\alpha_i^{-1}z_i$ otherwise. Then $e_{o(i)}x_ie_i=x_i$, $x_ix_i^*=\alpha_i^{-2}b_ie_{o(i)}=e_{o(i)}$, $x_i^*x_i=e_i$, and $r(x_i)=e_{o(i)i}$, as $\alpha_i$ maps to $1$ in $F$; and $o(o(i))=o(i)$, so $x_{o(i)}=e_{o(i)}$.
% Kernel: idem_nil (first part, from d_chS). Rank one: idem_nil (second part, d_idem_rank and the charpoly of single i i 1). Idempotents: AlgebraTools.st_coi_lift (Lemma sha-patching-chain-idempotents-orthogonal, with GlobalTheory.isAdicComplete_of_finite) and CA.d_corner (Lemma sha-patching-chain-idempotents-corner). The x_i: AlgebraTools.units_exist (b and b', hbb from the trace, hbF from r, AlgebraTools.sqrt_one_add for alpha, with HenselianRing B m). The Lean's idem states r(x_i); r(e_i) = e_ii is the hrep of st_coi_lift.
\end{proof}

\begin{lemma}\label{sha-patching-chain-gma}
Let $B$ be a local ring, $E$ a $B$-algebra, finite as a $B$-module, with a $B$-linear anti-automorphism $y\mapsto y^*$ such that $y^{**}=y$, and a $B$-linear map $\operatorname{tr}:E\to B$ with $\operatorname{tr}(yz)=\operatorname{tr}(zy)$ and $\operatorname{tr}(y^*)=\operatorname{tr}(y)$. Let $F$ be a field with a homomorphism $B\to F$ whose kernel is $\mathfrak m_B$, let $1\le a\le n$, let $o(i)=1$ for $i\le a$ and $o(i)=n$ for $i>a$, and let $r:E\to M_n(F)$ be a homomorphism of $B$-algebras whose values are block diagonal with blocks of sizes $a$ and $n-a$, with $r(y^*)={}^tr(y)$. Let $e_1,\dots,e_n$ be orthogonal idempotents of $E$ with $\sum_ie_i=1$, $e_i^*=e_i$, $\operatorname{tr}(e_i)=1$ and $e_iye_i\in Be_i$ for $y\in E$, and $x_1,\dots,x_n\in E$ with $e_{o(i)}x_ie_i=x_i$, $x_{o(i)}=e_{o(i)}$, $x_ix_i^*=e_{o(i)}$, $x_i^*x_i=e_i$ and $r(x_i)=e_{o(i)i}$. Then there are a commutative $B$-algebra $\mathcal A$, finite as a $B$-module, into which $B$ injects, a $B$-submodule $\mathcal N$ of $\mathcal A$ with $\mathcal A=B\cdot1\oplus\mathcal N$, $\mathcal N=0$ if $a=n$, and $yz\in\mathfrak m_B\cdot1$ for $y,z\in\mathcal N$, and an injective homomorphism of $B$-algebras $\varphi:E\to M_n(\mathcal A)$, such that:
\begin{enumerate}
\item the image of $\varphi$ is the module of matrices whose entries lie in $B\cdot1$ in the two diagonal blocks and in $\mathcal N$ in the two off-diagonal blocks;
\item $\varphi(y^*)={}^t\varphi(y)$ and $\varphi(e_i)=e_{ii}$;
\item if $y\in E$ and $b\in B$ satisfy $\varphi(y)_{ij}-b\in\mathcal N$, then the image of $b$ in $F$ is $r(y)_{ij}$.
\end{enumerate}
\end{lemma}
\leanref{Development/AutomorphicArithmetic/AutomorphicVarieties/PatchingModules/ClosureConjugation/GrAlg.lean}{UpperTheory.V.Pa.GrAlg}
\leanref{Development/AutomorphicArithmetic/AutomorphicVarieties/PatchingModules/ClosureConjugation/ComplOddPart.lean}{UpperTheory.V.Pa.gr_E}
\leanref{Development/AutomorphicArithmetic/AutomorphicVarieties/PatchingModules/ClosureConjugation/GrI.lean}{UpperTheory.V.Pa.gr_I}
\leanref{Development/AutomorphicArithmetic/AutomorphicVarieties/AlgebraTools/InvAlg.lean}{UpperTheory.V.AlgebraTools.TrInv.Units.vr}
\leanref{Development/AutomorphicArithmetic/AutomorphicVarieties/AlgebraTools/ThetaDiag.lean}{UpperTheory.V.AlgebraTools.vr_hom}
\leanref{Development/AutomorphicArithmetic/AutomorphicVarieties/AlgebraTools/VrRange.lean}{UpperTheory.V.AlgebraTools.vr_range}
\leanref{Development/AutomorphicArithmetic/AutomorphicVarieties/AlgebraTools/VrReduction.lean}{UpperTheory.V.AlgebraTools.vr_reduction}
\leanref{Development/AutomorphicArithmetic/AutomorphicVarieties/AlgebraTools/UnitsOneBlock.lean}{UpperTheory.V.AlgebraTools.units_one_block}
\begin{proof}
\emph{The case $a<n$.} Let $\mathcal N_0=e_1Ee_n$, and for $y,z\in\mathcal N_0$ let $\beta(y,z)=\operatorname{tr}(yz^*)$. Then $\beta$ is a symmetric $B$-bilinear form with $\beta(y,z)w=\beta(z,w)y$ for $y,z,w\in\mathcal N_0$, by the rank-one property of $e_1Ee_1$ and $e_nEe_n$ and $\operatorname{tr}(e_i)=1$. Let $\mathcal A=B\oplus\mathcal N_0$ with the product $(b,y)(b',z)=(bb'+\beta(y,z),bz+b'y)$, a commutative $B$-algebra, finite over $B$ as $E$ is, and let $\mathcal N=0\oplus\mathcal N_0$. For $y\in E$ let $z_{ij}=x_iyx_j^*$ and
\[\varphi(y)_{ij}=\big(\operatorname{tr}(z_{ij}),\ e_1z_{ij}e_n+e_1z_{ij}^*e_n\big).\]
The following are direct computations with the identities $x_ix_j^*=0$ for $i\ne j$, $x_ix_i^*=e_{o(i)}$, $\sum_ix_i^*x_i=1$ and the rank-one property: $\varphi$ is an injective homomorphism of $B$-algebras with $\varphi(y^*)={}^t\varphi(y)$; a matrix $Y$ over $\mathcal A$ is in its image if and only if $Y_{ij}\in B\cdot1$ when $o(i)=o(j)$ and $Y_{ij}\in\mathcal N$ when $o(i)\ne o(j)$, which is (1); the $B$-component of $\varphi(y)_{ij}$ maps to $r(y)_{ij}$ in $F$; and every $\beta(y,z)$ maps to $0$ in $F$. The last two use $r(y^*)={}^tr(y)$, $r(x_i)=e_{o(i)i}$ and that the values of $r$ are block diagonal.

From these: $B\to\mathcal A$, $b\mapsto(b,0)$, is injective, and $\mathcal A=B\cdot1\oplus\mathcal N$. For $y,z\in\mathcal N$, $yz=(\beta(y,z),0)$, and $\beta(y,z)\in\mathfrak m_B$, the kernel of $B\to F$. The entries of $\varphi(e_i)$ are computed from $x_ke_i=\delta_{ki}x_k$ and $e_ix_l^*=\delta_{il}x_l^*$, which give $\varphi(e_i)=e_{ii}$. If $\varphi(y)_{ij}-b\in\mathcal N$, then $b$ is the $B$-component of $\varphi(y)_{ij}$, which maps to $r(y)_{ij}$.

\emph{The case $a=n$.} Here $o(i)=1$ for all $i$. Let $\mathcal A=B$ and $\mathcal N=0$. The map with $\varphi(y)_{ij}=\operatorname{tr}(x_iyx_j^*)$ is an isomorphism of $B$-algebras $E\to M_n(B)$ with $\varphi(y^*)={}^t\varphi(y)$, by the same computations, and $\varphi(e_i)=e_{ii}$. As $x_iyx_j^*\in e_1Ee_1$, it equals $\varphi(y)_{ij}e_1$. Applying $r$, with $r(e_1)=e_{11}$, $r(x_i)=e_{1i}$ and $r(x_j^*)=e_{j1}$, gives $r(y)_{ij}e_{11}$ on the left and the image of $\varphi(y)_{ij}$ times $e_{11}$ on the right, which is (3).
% gr_E: p = 0 and q = n - 1 in Fin n (the book's 1 and n), U.oddForm p q on corner B (ep p) (ep q), whose Ext is the book's A; vr_hom (theta_mul, theta_ep, theta_st, theta_diag from units_basic), vr_range (ent_spec, rec_spec), vr_reduction (key, and red <= ker(B -> F)), with oddPart and compl_oddPart for N and mem_one_iff for B.1. gr_I: units_one_block with units_basic. The computations of the first paragraph are those of ThetaDiag.lean, VrRange.lean and VrReduction.lean.
\end{proof}

\begin{lemma}\label{sha-patching-chain-semisimple}
Let $B$ be a local ring, $E$ a $B$-algebra, $1\le a\le n$, $\mathcal A$ a commutative $B$-algebra with a $B$-submodule $\mathcal N$ such that $\mathcal A=B\cdot1\oplus\mathcal N$, and $\varphi:E\to M_n(\mathcal A)$ an injective homomorphism of $B$-algebras whose image is the module of matrices with entries in $B\cdot1$ in the two diagonal blocks, of sizes $a$ and $n-a$, and in $\mathcal N$ in the two off-diagonal blocks. Let $\kappa'$ be a field with a homomorphism $B\to\kappa'$, and $t\in B$ an element with nonzero image in $\kappa'$ that annihilates $\mathcal N$. Then $\kappa'\otimes_BE$ is semisimple.
\end{lemma}
\leanref{Development/AutomorphicArithmetic/AutomorphicVarieties/PatchingModules/ClosureConjugation/GrSs.lean}{UpperTheory.V.Pa.gr_ss}
\begin{proof}
For $i$ and $k$ in the same block let $\epsilon_{ik}\in E$ be the element with $\varphi(\epsilon_{ik})=e_{ik}$, which exists by the hypothesis on the image. As $\varphi$ is injective, $\epsilon_{ik}\epsilon_{kl}=\epsilon_{il}$, $\epsilon_{ik}\epsilon_{lm}=0$ for $k\ne l$, and $\sum_i\epsilon_{ii}=1$. Let $y\in E$. The entries of $\varphi(y)$ in the diagonal blocks lie in $B\cdot1$, and those in the off-diagonal blocks lie in $\mathcal N$, which $t$ kills; so $\varphi(ty)$ is a $B$-linear combination of the $e_{ik}$ with $i$ and $k$ in the same block, and $ty$ is the same combination of the $\epsilon_{ik}$. The elements $1\otimes\epsilon_{ik}$ define a homomorphism from the product of the matrix algebras over $\kappa'$ of the sizes of the blocks to $\kappa'\otimes_BE$. Its image contains each $1\otimes ty$, hence each $1\otimes y$, as $t$ is invertible in $\kappa'$; so it is surjective. A quotient of a semisimple ring is semisimple.
% gr_ss: hM (the elements e i k), hmul, hzero, hcross, hsum, hdec (t . y as a combination), alghom_of_units, RingHom.isSemisimpleRing_of_surjective.
\end{proof}

\begin{lemma}\label{sha-patching-chain-representation-determinant}
Let $B$ be a commutative ring, $E$ a $B$-algebra with a Cayley-Hamilton determinant $D$ of dimension $n$ \cite{Che14}, and $e_1,\dots,e_n$ orthogonal idempotents of $E$ with $\sum_ie_i=1$ and $D(1+Xe_i)=1+X$ in $B[X]$ for each $i$. Let $S$ be a commutative $B$-algebra and $\varphi:E\to M_n(S)$ a homomorphism of $B$-algebras with $\varphi(e_i)=e_{ii}$ for each $i$. Then $D\otimes S$ is the determinant of $\varphi$: for every commutative $S$-algebra $T$, all $t_1,\dots,t_k\in T$ and all $y_1,\dots,y_k\in E$, $D_T(\sum_at_a\otimes y_a)=\det\big(\sum_at_a\varphi(y_a)\big)$.
\end{lemma}
\leanref{Development/AutomorphicArithmetic/AutomorphicVarieties/CommutativeAlgebraHeight/DGmaDet.lean}{UpperTheory.V.CA.d_gma_det}
\begin{proof}
Let $\varphi_S:S\otimes_BE\to M_n(S)$ be the $S$-linear extension of $\varphi$. Both $D\otimes S$ and $\det\circ\varphi_S$ are determinants of dimension $n$ of the $S$-algebra $S\otimes_BE$, and the claim is that they are equal. By Amitsur's formula, two determinants of the same dimension of an algebra that have the same characteristic polynomial at every element of a multiplicative monoid spanning the algebra are equal \cite{Che14}. Let $M$ be the set of the $1\otimes y$ with $y=1$ or $e_iye_j=y$ for some $i$ and $j$. It is a monoid, since $e_iye_j=y$ and $e_ky'e_l=y'$ give $e_iyy'e_l=yy'$; and it spans $S\otimes_BE$, as $y=\sum_{i,j}e_iye_j$ for $y\in E$. At $1$ both characteristic polynomials are $(X-1)^n$. Now let $e_iye_j=y$.

If $i=j$, then $y=ae_i$ with $a=\operatorname{tr}_D(e_iy)$, by Lemma~\ref{sha-patching-chain-idempotents-corner}, which applies as $n\ge1$. For a commutative $S$-algebra $T'$ and $b\in T'$, $D_{T'}(1+b\,ae_i)=1+ab$, by specializing $X$ to $ab$ in $D(1+Xe_i)=1+X$, and $\det(1+ab\,e_{ii})=1+ab$. The characteristic polynomial $D(X-y)$ is $X^nD(1-X^{-1}y)$, computed over $S[X,X^{-1}]$, so it is determined by these values, and the two characteristic polynomials at $1\otimes y$ agree.

If $i\ne j$, then $y=(e_iy)e_j$. For any determinant the characteristic polynomials at $zz'$ and at $z'z$ agree \cite{Che14}, so at $y$ they are those at $e_je_iy=0$, which is $X^n$ for both.
% d_gma_det: main is d_ext (Amitsur's formula: d_amitsur, d_poly_ext and d_ext_poly) for D.baseChange S and Det.ofRep (rS r), with the private submonoid adapted; chi_one, chi_congr with ev_one_add and det_one_add_single and d_corner for i = j, d_tr (lam of a product) and chi_zero for i != j; then d_bc for the base change to T.
\end{proof}

\begin{lemma}\label{sha-patching-chain-representation}
Let $x$ be a normalized arc over $A=\kappa[[T]]$, with representation $\rho$ and blocks of sizes $a$ and $n-a$, and let $(q,\phi_j,\tau_j,U_{j,i},y_{j,i},s,f)$ be an arc datum for $x$ over $(\mathcal R,D,\psi)$, with relation ideal $I$, in the sense of Definition~\ref{sha-patching-chain-datum}. Let $\mathcal K$ be the fraction field of $A$, $\hat{\mathcal R}$ the completion of the localization of $\mathcal R$ at $\ker f$, and $\hat f:\hat{\mathcal R}\to\mathcal K$ a homomorphism with kernel $\mathfrak m_{\hat{\mathcal R}}$ whose composite with $\mathcal R\to\hat{\mathcal R}$ is $f$. Then there are a commutative $\hat{\mathcal R}$-algebra $\mathcal A$, finite as an $\hat{\mathcal R}$-module, into which $\hat{\mathcal R}$ injects, an $\hat{\mathcal R}$-submodule $\mathcal N$ of $\mathcal A$ with $\mathcal A=\hat{\mathcal R}\cdot1\oplus\mathcal N$, $\mathcal N=0$ if $a=n$, and $yz\in\mathfrak m_{\hat{\mathcal R}}\cdot1$ for $y,z\in\mathcal N$, and a multiplicative map $r:\Gamma_{\infty,0}\to M_n(\mathcal A)$, such that:
\begin{enumerate}
\item the $\hat{\mathcal R}$-span of $r(\Gamma_{\infty,0})$ is the module of matrices whose entries lie in $\hat{\mathcal R}\cdot1$ in the two diagonal blocks and in $\mathcal N$ in the two off-diagonal blocks;
\item $D\otimes\mathcal A$ is the determinant of $r$, and ${}^tr(g)\,r(c_\infty gc_\infty^{-1})=\varepsilon(g)^{1-n}$ for $g\in\Gamma_{\infty,0}$;
\item if $g\in\Gamma_{\infty,0}$ and $b\in\hat{\mathcal R}$ satisfy $r(g)_{ij}-b\in\mathcal N$, then $\hat f(b)=\rho(p(g))_{ij}$;
\item the homomorphism $\mathcal R[\Gamma_{\infty,0}]\to M_n(\mathcal A)$ induced by $r$ kills $I$;
\item a prime $Q$ of $\hat{\mathcal R}$ contains the annihilator of $\mathcal N$ if, for some field $\kappa'$ and homomorphism $\hat{\mathcal R}\to\kappa'$ with kernel $Q$, the quotient of $\kappa'[\Gamma_{\infty,0}]$ by the two-sided ideal generated by the image of $I$ is not semisimple;
\item there is a finitely generated $\mathcal R$-submodule $P$ of $M_n(\mathcal A)$ containing $r(\Gamma_{\infty,0})$ such that for every $j\ge0$ there is an open normal subgroup $V$ of $\Gamma_{\infty,0}$ with $r(gv)-r(g)\in\mathfrak m_{\mathcal R}^jP$ for $g\in\Gamma_{\infty,0}$ and $v\in V$.
\end{enumerate}
\end{lemma}
\leanref{Development/AutomorphicArithmetic/AutomorphicVarieties/PatchingModules/PatchedSupportCompletions.lean}{UpperTheory.V.Pa.URep}
\leanref{Development/AutomorphicArithmetic/AutomorphicVarieties/PatchingModules/PatchedSupportCompletions.lean}{UpperTheory.V.Pa.S6}
\leanref{Development/AutomorphicArithmetic/AutomorphicVarieties/PatchingModules/AutomorphicDetection/S6.lean}{UpperTheory.V.Pa.s6}
\leanref{Development/AutomorphicArithmetic/AutomorphicVarieties/PatchingModules/ClosureConjugation/UrepOf.lean}{UpperTheory.V.Pa.urep_of}
\leanref{Development/AutomorphicArithmetic/AutomorphicVarieties/PatchingModules/ClosureConjugation/GrDet.lean}{UpperTheory.V.Pa.gr_det}
\leanref{Development/AutomorphicArithmetic/AutomorphicVarieties/PatchingModules/ClosureConjugation/GrThick.lean}{UpperTheory.V.Pa.gr_thick}
\begin{proof}
The ring $\hat{\mathcal R}$ is complete local noetherian \cite[\S8]{Mat86}. The element $2$ is a unit of $\hat{\mathcal R}$: otherwise it lies in $\mathfrak m_{\hat{\mathcal R}}=\ker\hat f$, so $2=0$ in $\kappa$, while $\kappa$ has characteristic $\ell$, which is odd. By Lemma~\ref{sha-patching-chain-cayley}(1), $\hat f$ is surjective. Let $E$, $\iota$, $D_E$, $*$ and $P_E$ be as in Lemma~\ref{sha-patching-chain-cayley}(2), and $r_E$ the map called $r$ there. By Lemma~\ref{sha-patching-chain-cayley}(3), Lemma~\ref{sha-patching-chain-idempotents} applies to $\hat{\mathcal R}$, $E$, $D_E$, $*$, $F=\mathcal K$ and $r_E$, and gives elements $e_i$ and $x_i$; and the values of $r_E$ are block diagonal. So Lemma~\ref{sha-patching-chain-gma} applies, with $\operatorname{tr}=\operatorname{tr}_{D_E}$, and gives $\mathcal A$, $\mathcal N$ and $\varphi$. Put $r(g)=\varphi(\iota(g))$.

(1) As $\iota$ is surjective, the $\hat{\mathcal R}$-span of $r(\Gamma_{\infty,0})$ is $\varphi(E)$, which is the module of (1) by Lemma~\ref{sha-patching-chain-gma}(1).

(2) As $D_E$ is Cayley-Hamilton, the $e_i$ are orthogonal idempotents with sum $1$, $\varphi(e_i)=e_{ii}$ and $D_E(1+Xe_i)=1+X$ for all $i$, Lemma~\ref{sha-patching-chain-representation-determinant}, with $S=\mathcal A$, shows that $D_E\otimes\mathcal A$ is the determinant of $\varphi$. With $D_E\circ\iota=D\otimes\hat{\mathcal R}$, this gives that $D\otimes\mathcal A$ is the determinant of $r=\varphi\circ\iota$ on $\Gamma_{\infty,0}$. By Lemma~\ref{sha-patching-chain-gma}(2), ${}^tr(g)=\varphi(\iota(g)^*)=\varphi(\iota(g^*))=\varepsilon(g)^{1-n}r\big((c_\infty gc_\infty^{-1})^{-1}\big)$, so ${}^tr(g)\,r(c_\infty gc_\infty^{-1})=\varepsilon(g)^{1-n}$.

(3) By Lemma~\ref{sha-patching-chain-gma}(3), $\hat f(b)$ is the entry $(i,j)$ of $r_E(\iota(g))=\rho(p(g))$.

(4) The homomorphism induced by $r$ is $\varphi\circ\iota$ on $\hat{\mathcal R}[\Gamma_{\infty,0}]$, composed with $\mathcal R[\Gamma_{\infty,0}]\to\hat{\mathcal R}[\Gamma_{\infty,0}]$, and $\iota$ kills the image of $I$.

(5) Suppose that $Q$ does not contain the annihilator of $\mathcal N$, and let $t$ be an element of the annihilator outside $Q$. Its image in $\kappa'$ is nonzero, so $\kappa'\otimes_{\hat{\mathcal R}}E$ is semisimple by Lemma~\ref{sha-patching-chain-semisimple}. Since the kernel of $\iota$ is the two-sided ideal generated by the image of $I$, the quotient of $\kappa'[\Gamma_{\infty,0}]$ by the two-sided ideal generated by the image of $I$ is a quotient of $\kappa'\otimes_{\hat{\mathcal R}}E$, through $1\otimes\iota(g)\mapsto g$. So it is semisimple.

(6) Take $P=\varphi(P_E)$, which is finitely generated over $\mathcal R$ as $\varphi$ is $\mathcal R$-linear; then $r(gv)-r(g)=\varphi(\iota(gv)-\iota(g))$.
% s6: rF_surj, hat_ex, hat_st, hat_chi, hat_range, idem, then gr_E (a < n) or gr_I (a = n), then urep_of. urep_of: span from g.range, det = gr_det (CA.d_gma_det, Lemma sha-patching-chain-representation-determinant, through the private gdet_eq_ofRep and d_gmap), form from g.st and H.st, red from g.red and H.r_iota, kills from H.ker, thick = gr_thick (quot_ss and gr_ss), P = map of H.P with ker_le_span_sub_sup_map (P.FG), vr_mem and cont.
\end{proof}

\subsection{Selmer groups and tangent spaces at arcs}\label{subsection-sha-arc-tangent-spaces}

\begin{definition}\label{sha-patching-chain-selmer}
In Situation~\ref{sha-hecke-situation}, let $\Gamma_\infty$, $p$ and $c_\infty$ be as in Definition~\ref{sha-patching-chain-datum}. Let $x$ be a normalized arc over $A=\kappa[[T]]$, with representation $\rho$, blocks of sizes $a$ and $n-a$, subgroups $I'_u$, matrices $w_v$ and a Selmer basis $\kappa_1,\dots,\kappa_g$ with matrices $Y_{j,v}$; let $R_g$ and $\rho_1$ be as in Definition~\ref{sha-patching-selmer}(3), and $w'_v=(1-\sum_j\epsilon_jY_{j,v})w_v$. Let $X$ be $A$ or the fraction field $\mathcal K$ of $A$, and $R_{g,X}=R_g\otimes_AX$.
\begin{enumerate}
\item Let $X'$ be $X$ or $R_{g,X}$, $\rho':\Gamma_{\infty,0}\to GL_n(X')$ a homomorphism, $w''_v\in GL_n(X')$ for $v\mid\ell$, and $V\subset M_n(X')$ an $A$-submodule. Then $\mathrm{Sel}(\rho',w'',V)$ is the set of maps $z:\Gamma_{\infty,0}\to V$ such that: the reduction of $z$ modulo $T^mX'$ is locally constant for every $m$; $z(gh)=z(g)+\rho'(g)z(h)\rho'(g)^{-1}$ and $z(c_\infty gc_\infty^{-1})=-{}^tz(g)$ for $g,h\in\Gamma_{\infty,0}$; $z$ vanishes on the image of each $I'_u$; and for each $v\mid\ell$ there is $Y\in V$ such that ${w''_v}^{-1}\big(z(g)-\rho'(g)Y\rho'(g)^{-1}+Y\big)w''_v$ is upper triangular for $g\in D_v$, with zero diagonal for $g\in\mu_v$. Its \emph{coboundaries} are the maps $g\mapsto\rho'(g)Y\rho'(g)^{-1}-Y$ with $Y\in V$ and ${}^tY=-Y$.
\item $\mathrm{Sel}^{\mathrm{ev}}(X)$ and $\mathrm{Sel}^{\mathrm{od}}(X)$ are $\mathrm{Sel}(\rho\circ p,w,V)$ with $V$ the even, respectively the odd, matrices over $X$, in the sense of Definition~\ref{sha-patching-selmer}; $\mathrm B^{\mathrm{od}}(X)$ is the set of coboundaries for the odd matrices.
\item $\mathrm{Sel}^1(X)$ and $\mathrm{Sel}^2(X)$ are $\mathrm{Sel}(\rho_1\circ p,w',V)$ with $V$ the matrices over $R_{g,X}$ whose entries lie in $X$ in the two diagonal blocks and in $\sum_jX\epsilon_j$ in the two off-diagonal blocks, respectively in $\sum_jX\epsilon'_j$ in the two diagonal blocks and in $X\eta$ in the two off-diagonal blocks; $\mathrm B^1(X)$ is the set of coboundaries for the first $V$. For $\Psi\in\mathrm{Sel}^2(X)$ and $1\le j\le g$, let $\Psi'_j$ be the map whose entries are the coefficients of $\epsilon'_j$ in the entries of $\Psi$.
\item A map $\Gamma_{\infty,0}\to\mathcal K$ has \emph{bounded denominators} if its product with some power of $T$ is a map to $A$ whose reductions modulo the powers of $T$ are locally constant. A map to $M_n(\mathcal K)$, or to $M_n(R_{g,\mathcal K})$, has bounded denominators if every entry has, respectively if every coefficient of every entry in the basis $1$, $\epsilon_j$, $\epsilon'_j$, $\eta$ has.
\item A family $(z_1,\dots,z_g)$ of maps $\Gamma_{\infty,0}\to M_n(\mathcal K)$ \emph{lifts} if there is $\Psi\in\mathrm{Sel}^2(\mathcal K)$ with bounded denominators and $\Psi'_j=z_j$ for all $j$.
\end{enumerate}
\end{definition}
\leanref{Development/AutomorphicArithmetic/AutomorphicVarieties/PatchingModules/PairedDualFirstOrderDualMatrixSpace.lean}{UpperTheory.V.Pa.SelZ}
\leanref{Development/AutomorphicArithmetic/AutomorphicVarieties/PatchingModules/PairedDualFirstOrderDualMatrixSpace.lean}{UpperTheory.V.Pa.SelP}
\leanref{Development/AutomorphicArithmetic/AutomorphicVarieties/PatchingModules/PairedDualFirstOrderDualMatrixSpace.lean}{UpperTheory.V.Pa.SelM}
\leanref{Development/AutomorphicArithmetic/AutomorphicVarieties/PatchingModules/PairedDualFirstOrderDualMatrixSpace.lean}{UpperTheory.V.Pa.BdM}
\leanref{Development/AutomorphicArithmetic/AutomorphicVarieties/PatchingModules/PairedDualFirstOrderDualMatrixSpace.lean}{UpperTheory.V.Pa.Sel1}
\leanref{Development/AutomorphicArithmetic/AutomorphicVarieties/PatchingModules/PairedDualFirstOrderDualMatrixSpace.lean}{UpperTheory.V.Pa.Bd1}
\leanref{Development/AutomorphicArithmetic/AutomorphicVarieties/PatchingModules/PairedDualFirstOrderDualMatrixSpace.lean}{UpperTheory.V.Pa.Sel2}
\leanref{Development/AutomorphicArithmetic/AutomorphicVarieties/PatchingModules/SelmerScaling/SelmerSpannedByKappaLaurent.lean}{UpperTheory.V.Pa.BndD}
\leanref{Development/AutomorphicArithmetic/AutomorphicVarieties/PatchingModules/SelmerScaling/SelmerSpannedByKappaLaurent.lean}{UpperTheory.V.Pa.Lifts}

\begin{lemma}\label{sha-patching-chain-limit-bounds-kernel}
In the setting of Lemma~\ref{sha-patching-chain-limit}, the kernel of $p_\infty:\Gamma_\infty\to\Gamma$ is contained in the closed normal subgroup of $\Gamma_\infty$ generated by $\tau_1,\dots,\tau_q$.
\end{lemma}
\leanref{Development/AutomorphicArithmetic/AutomorphicVarieties/PatchingModules/SeparableMatrixStability/KerOo.lean}{UpperTheory.V.Pa.ker_oo}
\begin{proof}
Let $K$ be that closed normal subgroup. As $\Gamma_\infty/K$ is profinite, $K$ is the intersection of the subgroups $KO$ with $O$ open normal in $\Gamma_\infty$; so it suffices to show that $\ker p_\infty\subset W=KO$ for each such $O$. Let $V=\pi_\infty^{-1}(W)$, an open normal subgroup of $\tilde\Gamma$ containing $\ker\pi_\infty$. The group $\Gamma$ is topologically finitely generated, as a continuous image of $\tilde\Gamma$, so it has finitely many open subgroups of index $[\tilde\Gamma:V]$. By Lemma~\ref{sha-patching-chain-limit}(1), (2) and (4), there is $N$ such that $(\ker\pi_N)V=(\ker\pi_\infty)V$; $\tau_{N,j}$ and $\tau_j$ are close modulo $V$ for all $j$; and $p_N(\pi_N(g))\,p_\infty(\pi_\infty(g))^{-1}\in A$ for all $g\in\tilde\Gamma$ and every open normal subgroup $A$ of $\Gamma$ of index $[\tilde\Gamma:V]$.

Then $\ker\pi_N\subset V$, so $\pi_N(V)$ is an open normal subgroup of $\Gamma_N$, as $\pi_N$ is a closed surjection. It contains every $\tau_{N,j}$: if $\tau_{N,j}=\pi_N(g)$ and $\tau_j=\pi_\infty(g')$ with $g^{-1}g'\in V$, then $g'\in V$, as $\tau_j\in K$, so $g\in V$. As $P(H_N)=1$, Lemma~\ref{sha-patching-places}(3) puts the kernel of $p_N$ inside the closed normal subgroup generated by the $\tau_{N,j}$, hence inside $\pi_N(V)$. Let $f=p_N\circ\pi_N:\tilde\Gamma\to\Gamma$, a continuous surjection. Then $\ker f\subset\pi_N^{-1}(\pi_N(V))=V$, so $f(V)$ is an open normal subgroup of $\Gamma$ of index $[\tilde\Gamma:V]$. Let $y\in\ker p_\infty$, and $y=\pi_\infty(g)$. By the choice of $N$, $f(g)=p_N(\pi_N(g))\,p_\infty(\pi_\infty(g))^{-1}\in f(V)$, so $g\in V$ as $\ker f\subset V$, and $y\in\pi_\infty(V)=W$.
% ker_oo: totallyDisconnected_quot (AlgebraTools, QuotProf.lean), W = K sup O, V = comap of W, finite_hom_of_topFG for the open subgroups of the index of V in G0.Gamma (with TopFG.subgroup_quot), Lambda.ker_lim, Lambda.ta_lim and Lambda.p_lim for N, T.ker_le with T.ProL for the kernel of p_N, the private isOpen_map_of_ker_le.
\end{proof}

\begin{lemma}\label{sha-patching-chain-limit-bounds-transfer}
In the setting of Lemma~\ref{sha-patching-chain-limit}, let $R$ be a finite commutative $\kappa[[T]]$-algebra in which $T$ is nilpotent, $\rho:\Gamma_0\to M_n(R)$ a locally constant map with $\rho(1)=1$ and $\rho(gh)=\rho(g)\rho(h)$, $w_v$ and $w^\flat_v$ matrices over $R$ for $v\mid\ell$, and $V\subset M_n(R)$ a $\kappa[[T]]$-submodule. Let $\Gamma^\dagger$ be $\Gamma_\infty$, with $p_\infty$, $c_\infty=\pi_\infty(\tilde c)$ and the local maps of Lemma~\ref{sha-patching-chain-limit}(3), or one of the $\Gamma_N$, with $p_N$, $c$ and its local maps, and write $\Gamma^\dagger_0$ for the inverse image of $\Gamma_0$. Let $\mathrm{Sel}(\Gamma^\dagger)$ be the set of maps $z:\Gamma^\dagger_0\to V$ that satisfy the conditions of Definition~\ref{sha-patching-chain-selmer}(1) for $\rho$ composed with the map to $\Gamma$, with $w_v$ and $w^\flat_v$ in place of $w''_v$ and ${w''_v}^{-1}$. Then for $\mathcal U$-almost all $N$:
\begin{enumerate}
\item $\rho(p_N(\pi_N(g)))=\rho(p_\infty(\pi_\infty(g)))$ for $g\in\tilde\Gamma_0$;
\item for every $z\in\mathrm{Sel}(\Gamma_\infty)$ there is $z_N\in\mathrm{Sel}(\Gamma_N)$ with $z_N(\pi_N(g))=z(\pi_\infty(g))$ for $g\in\tilde\Gamma_0$, and for every $z_N\in\mathrm{Sel}(\Gamma_N)$ there is such a $z\in\mathrm{Sel}(\Gamma_\infty)$.
\end{enumerate}
\end{lemma}
\leanref{Development/AutomorphicArithmetic/AutomorphicVarieties/PatchingModules/SeparableMatrixStability/LcOfAdic.lean}{UpperTheory.V.Pa.t1}
\leanref{Development/AutomorphicArithmetic/AutomorphicVarieties/PatchingModules/SeparableMatrixStability/CocFin.lean}{UpperTheory.V.Pa.coc_fin}
\leanref{Development/AutomorphicArithmetic/AutomorphicVarieties/PatchingModules/SeparableMatrixStability/EventuallyFactorThroughTowerIffLimit.lean}{UpperTheory.V.Pa.eventually_factor_through_tower_iff_limit}
\leanref{Development/AutomorphicArithmetic/AutomorphicVarieties/PatchingModules/SeparableMatrixStability/EventuallyFactorizationCompatibleWithLocalMaps.lean}{UpperTheory.V.Pa.eventually_factorization_compatible_with_local_maps}
\leanref{Development/AutomorphicArithmetic/AutomorphicVarieties/PatchingModules/SeparableMatrixStability/EventuallyLocallyConstantCompTowerEq.lean}{UpperTheory.V.Pa.eventually_locallyConstant_comp_tower_eq}
\begin{proof}
\emph{(1)} As $\rho$ is locally constant on the compact group $\Gamma_0$, it takes finitely many values, each on an open and compact set; so there is an open normal subgroup $V'$ of $\Gamma$ with $\rho(ab)=\rho(b)$ for $a\in V'$ and $b\in\Gamma_0$. By Lemma~\ref{sha-patching-chain-limit}(1), for $\mathcal U$-almost all $N$, $p_N(\pi_N(g))\,p_\infty(\pi_\infty(g))^{-1}\in V'$ for all $g$, which gives (1).

\emph{The groups $\tilde\Gamma_0$.} By Lemma~\ref{sha-patching-chain-limit}(1), $p_N(\pi_N(g))$ and $p_\infty(\pi_\infty(g))$ agree modulo $H\subset\Gamma_0$; so $\pi_N$ maps $\tilde\Gamma_0$ onto $\Gamma_{N,0}$, and $\ker\pi_N\subset\tilde\Gamma_0$. The open subgroup $\tilde\Gamma_0$ has finite index in $\tilde\Gamma$, so it is topologically finitely generated. Let $\rho'(g)=\rho(p_\infty(\pi_\infty(g)))$ for $g\in\tilde\Gamma_0$, and let $\mathcal C$ be the set of locally constant maps $c:\tilde\Gamma_0\to M_n(R)$ with $c(gh)=c(g)+\rho'(g)c(h)\rho'(g)^{-1}$. It is finite: for two such maps, the set where they agree is a closed subgroup, so each $c$ is determined by its values on finitely many topological generators, and $R$ is finite.

\emph{Factorization.} Let $c\in\mathcal C$, and let $S_c$ be the set of $k\in\tilde\Gamma_0$ with $c(gk)=c(g)$ for all $g$. It is a subgroup, and it is open, as $c$ is locally constant on the compact group $\tilde\Gamma_0$. For a surjective homomorphism $\pi$ from $\tilde\Gamma_0$, $c=z\circ\pi$ for some $z$ if and only if $\ker\pi\subset S_c$. Let $V_c$ be the intersection of the conjugates of $S_c$ in $\tilde\Gamma$, an open normal subgroup of $\tilde\Gamma$. As $\ker\pi_N$ and $\ker\pi_\infty$ are normal in $\tilde\Gamma$ and contained in $\tilde\Gamma_0$, $c$ factors through $\pi_N$, respectively $\pi_\infty$, if and only if $\ker\pi_N\subset V_c$, respectively $\ker\pi_\infty\subset V_c$. By Lemma~\ref{sha-patching-chain-limit}(2), for $\mathcal U$-almost all $N$, $(\ker\pi_N)V_c=(\ker\pi_\infty)V_c$; so $c$ factors through $\pi_N$ if and only if it factors through $\pi_\infty$.

\emph{Local maps.} Let $c\in\mathcal C$. There is an open normal subgroup $V'_c$ of $\tilde\Gamma$ with $c(g)=c(g')$ whenever $g^{-1}g'\in V'_c$. By Lemma~\ref{sha-patching-chain-limit}(3), for $\mathcal U$-almost all $N$, the images in $\Gamma_N$ and in $\Gamma_\infty$ of every element of a local group are close modulo $V'_c$. If moreover $c=z_N\circ\pi_N=z\circ\pi_\infty$, then $z_N$ and $z$ agree at the images of each element of the local groups: if these images are $\pi_N(g)$ and $\pi_\infty(g')$ with $g^{-1}g'\in V'_c$, then $g,g'\in\tilde\Gamma_0$ and $z_N(\pi_N(g))=c(g)=c(g')=z(\pi_\infty(g'))$.

\emph{(2)} As $\mathcal C$ is finite, $\mathcal U$-almost all $N$ satisfy (1) and the conclusions of the last two paragraphs for every $c\in\mathcal C$. Fix such an $N$, and let $z\in\mathrm{Sel}(\Gamma_\infty)$. As $T$ is nilpotent in $R$, $z$ is locally constant, so $c=z\circ\pi_\infty$ lies in $\mathcal C$; let $z_N$ be the map with $z_N\circ\pi_N=c$. It lies in $\mathrm{Sel}(\Gamma_N)$. Its values are values of $z$, in $V$. It is locally constant, since $\pi_N$ is a continuous surjection of compact groups, hence a quotient map. The cocycle relation for $z_N$ follows from that for $z$ and (1), as $\pi_N$ is surjective onto $\Gamma_{N,0}$; the relation with $c$ follows from that with $c_\infty$, as $\pi_N(\tilde c)=c$, $\pi_\infty(\tilde c)=c_\infty$ and conjugation by $\tilde c$ preserves $\tilde\Gamma_0$. The vanishing on the images of the $I'_u$ and the local conditions at $v\mid\ell$, with the matrices $Y$ of $z$, hold for $z_N$, as $z_N$ and $z$ agree at the images of the local groups and $\rho\circ p_N$ and $\rho\circ p_\infty$ agree there with $\rho$ composed with the maps to $\Gamma_0$. The converse is proved in the same way, with $\Gamma_N$ and $\Gamma_\infty$ exchanged.
% t1: eventually_locallyConstant_comp_tower_eq for (1) (Lim.p_lim with an open normal subgroup inside a separating neighbourhood, compact_open_separated_mul_left); coc_fin for the finite set C (Ph0 is TopFG by the private topFG_subgroup_of_finiteIndex_of_isOpen); eventually_factor_through_tower_iff_limit (the private stabR, stabR_open, fac_iff and ker_iff, with Lim.ker_lim); eventually_factorization_compatible_with_local_maps (the private exists_open_normal with Lim.phi_lim, phiu_lim, ph_lim and ta_lim); the private t1_lemma for the conditions of SelZ, with isLocallyConstant_matrix_of_adicCont_of_isNilpotent and isLocallyConstant_of_comp_surjective_of_compactSpace.
\end{proof}

\begin{lemma}\label{sha-patching-chain-limit-bounds-artin-rees}
Let $A$ be a noetherian ring, $I\subset A$ an ideal, $M$ and $M'$ finite $A$-modules, and $(f_s)_{s\in S}$ a family of $A$-linear maps $M\to M'$. Then there is $c\ge0$ such that for every $m\ge0$ and every $y\in M$ with $f_s(y)\in I^{m+c}M'$ for all $s$, and every $A$-linear $\varphi:M\to A$ that vanishes on $\bigcap_s\ker f_s$, $\varphi(y)\in I^m$.
\end{lemma}
\leanref{Development/AutomorphicArithmetic/AutomorphicVarieties/PatchingModules/SeparableMatrixStability/CoefficientsBoundedOfSelmerBoundary.lean}{ar_family}
\begin{proof}
The span of the $f_s$ in $\operatorname{Hom}_A(M,M')$ is finitely generated, as $\operatorname{Hom}_A(M,M')$ is a noetherian $A$-module; let $v_1,\dots,v_k$ generate it, and $F=(v_1,\dots,v_k):M\to M'^k$. Then $\ker F=\bigcap_s\ker f_s$, and $F(y)\in I^{m+c}M'^k$ if every $f_s(y)$ lies in $I^{m+c}M'$. By the Artin-Rees lemma \cite[Theorem~8.5]{Mat86}, applied to the image of $F$, there is $c$ with $F^{-1}(I^{m+c}M'^k)\subset\ker F+I^mM$ for all $m$. Then $\varphi(y)\in\varphi(\ker F)+I^m\varphi(M)\subset I^m$.
% ar_family: IsNoetherian.noetherian for the span, Submodule.fg_iff_exists_fin_generating_family, the map Phi = LinearMap.pi v, comap_pow_le (Artin-Rees, Ideal.exists_pow_inf_eq_pow_smul) and the bound hle.
\end{proof}

\begin{lemma}\label{sha-patching-chain-limit-bounds-closed}
In the setting of Definition~\ref{sha-patching-chain-selmer}, write $A_m=A/T^mA$, and $\mathrm{Sel}^2_m$ for the set defined as $\mathrm{Sel}^2(X)$ in Definition~\ref{sha-patching-chain-selmer}(3), with $A_m$ in place of $X$ and $R_g/T^mR_g$ in place of $R_{g,X}$. Let $\Psi:\Gamma_{\infty,0}\to M_n(R_g)$ be a map whose reduction modulo $T^m$ lies in $\mathrm{Sel}^2_m$ for every $m\ge1$. Then $\Psi\in\mathrm{Sel}^2(A)$.
\end{lemma}
\leanref{Development/AutomorphicArithmetic/AutomorphicVarieties/PatchingModules/SeparableMatrixStability/SelmerSecondClosedUnderLimits.lean}{UpperTheory.V.Pa.selmer_second_closed_under_limits}
\begin{proof}
As $R_g$ is free over $A$ and $\bigcap_mT^mA=0$, a matrix over $R_g$ is determined by its reductions modulo the $T^m$, and it lies in the second module $V$ of Definition~\ref{sha-patching-chain-selmer}(3) if its reductions do. So the values of $\Psi$ lie in $V$, and the cocycle relation, the relation with $c_\infty$ and the vanishing on the images of the $I'_u$ hold for $\Psi$, as they hold modulo every $T^m$. The reductions of $\Psi$ modulo the $T^m$ are locally constant, by hypothesis.

Let $v\mid\ell$. For each $m$ choose a matrix $Y_m$ over $R_g$ whose reduction modulo $T^m$ is a matrix for the local condition of $\Psi\bmod T^m$ at $v$, with values in $V\bmod T^m$. Let $\mathcal U'$ be a nonprincipal ultrafilter on $\mathbf N$. As $\kappa$ is finite, for each $c$ the coefficient of $T^c$ in each coordinate, in the basis $1,\epsilon_k,\epsilon'_k,\eta$, of each entry of $Y_m$ is constant for $\mathcal U'$-almost all $m$; let $Y$ be the matrix with these coefficients. Then for every $j$ there is $N\ge j$ with $Y\equiv Y_N$ modulo $T^j$. So $Y$ lies in $V$, as its reductions do. Let $g\in D_v$ and $j\ge1$, and choose such an $N$. The entries of ${w'_v}^{-1}(\Psi(g)-\rho_1(p(g))Y\rho_1(p(g))^{-1}+Y)w'_v$ below the diagonal, and on the diagonal if $g\in\mu_v$, are congruent modulo $T^j$ to those of the same matrix with $Y_N$ in place of $Y$, which vanish modulo $T^N$, hence modulo $T^j$. As $j$ is arbitrary, they vanish.
% selmer_second_closed_under_limits: mem_second_order_of_reductions for the values, red_ker and pow_bot for the continuity, mat_sep (red_sep) for the relations, and for LocAt the choice Yt with red_surj, the private ulim_D (ulim_A over the finite field, Filter.hyperfilter), hev, hE and hent (red_mono).
\end{proof}

\begin{lemma}\label{sha-patching-chain-limit-bounds-ultrafilter}
In the setting of Lemma~\ref{sha-patching-chain-limit-bounds-closed}, suppose that $\Gamma_\infty$ is topologically finitely generated. Let $\mathcal U'$ be an ultrafilter on $\mathbf N$ that contains the complements of the finite sets, and for each $m\ge1$ let $Z_m:\Gamma_{\infty,0}\to M_n(R_g)$ be a map whose reduction modulo $T^m$ lies in $\mathrm{Sel}^2_m$. Then there is $\Psi\in\mathrm{Sel}^2(A)$ such that for every $j\ge1$ and $\mathcal U'$-almost all $m$, $\Psi\equiv Z_m$ modulo $T^j$.
\end{lemma}
\leanref{Development/AutomorphicArithmetic/AutomorphicVarieties/PatchingModules/SeparableMatrixStability/SelmerSecondExistsUltrafilterLimit.lean}{UpperTheory.V.Pa.selmer_second_exists_ultrafilter_limit}
\begin{proof}
For $m\ge j$, the reduction modulo $T^j$ of an element of $\mathrm{Sel}^2_m$ lies in $\mathrm{Sel}^2_j$, as the conditions of Definition~\ref{sha-patching-chain-selmer}(1) are preserved by homomorphisms of $A$-algebras. The set $\mathrm{Sel}^2_j$ is finite: $\Gamma_{\infty,0}$ is open of finite index in $\Gamma_\infty$, hence topologically finitely generated; $R_g/T^jR_g$ is finite and $T$ is nilpotent in it, so the elements of $\mathrm{Sel}^2_j$ are locally constant cocycles, each determined by its values on finitely many topological generators. So for each $j$, the reductions modulo $T^j$ of the $Z_m$, for $\mathcal U'$-almost all $m$, lie in a finite set, and are therefore equal for $\mathcal U'$-almost all $m$; let $\mu(j)$ be one such $m$. Let $\Psi$ be the map whose coefficient of $T^c$, in each coordinate of each entry in the basis $1,\epsilon_k,\epsilon'_k,\eta$, is that of $Z_{\mu(c+1)}$. For each $j$ and $\mathcal U'$-almost all $m$, $Z_m\equiv Z_{\mu(c+1)}$ modulo $T^{c+1}$ for all $c<j$, so $\Psi\equiv Z_m$ modulo $T^j$. In particular every reduction of $\Psi$ modulo $T^j$ is that of some $Z_m$ with $m\ge j$, which lies in $\mathrm{Sel}^2_j$; so $\Psi\in\mathrm{Sel}^2(A)$ by Lemma~\ref{sha-patching-chain-limit-bounds-closed}.
% selmer_second_exists_ultrafilter_limit: hred (sel_red with Ideal.Quotient.factor), the private selmer_second_finite_modulo_power (coc_fin, TopFG.subgroup_quot, lc_of_nil, lc_mat), ultra_fin for mu, the diagonal Psi with diag_mk, then selmer_second_closed_under_limits (Lemma sha-patching-chain-limit-bounds-closed).
\end{proof}

\begin{lemma}\label{sha-patching-chain-limit-bounds}
In Situation~\ref{sha-hecke-situation}, let $\sigma_1,\dots,\sigma_q$, $H$ and the places $w_{N,j}$ be as in Lemma~\ref{sha-patching-places}, with $P(H)=1$ and $P(H_N)=1$ for all $N$, and let $\Gamma_\infty$ and $p_\infty$ be as in Lemma~\ref{sha-patching-chain-limit}. Let $x$ be a normalized arc over $A=\kappa[[T]]$, with representation $\rho$ and a Selmer basis $\kappa_1,\dots,\kappa_g$, such that for each $j$ the characteristic polynomial of $\rho(\sigma_j)$ is $\prod_i(X-\alpha_{j,i})$ with pairwise distinct $\alpha_{j,i}\in A$. Let $D_0\ge0$, and suppose that the conclusion of Lemma~\ref{sha-patching-tangent} holds for $x$ with $n[L^+:\mathbf Q]+nq$ replaced by $D_0$, and that the conclusion of Lemma~\ref{sha-patching-obstruction} holds for $x$. Then, in the notation of Definition~\ref{sha-patching-chain-selmer} with $p=p_\infty$:
\begin{enumerate}
\item if $z_1,\dots,z_e\in\mathrm{Sel}^1(A)$ are such that $\sum_ic_iz_i\in\mathrm B^1(A)$, with $c_i\in A$, only if all $c_i$ vanish, then $e\le D_0$;
\item for any $z_{i,j}\in\mathrm{Sel}^{\mathrm{ev}}(A)$, $0\le i\le g$, $1\le j\le g$, there are $c_0,\dots,c_g\in A$, not all zero, and $\Psi\in\mathrm{Sel}^2(A)$ with $\Psi'_j=\sum_ic_iz_{i,j}$ for all $j$;
\item for every $z\in\mathrm{Sel}^{\mathrm{od}}(A)$ there are $s\ge0$ and $c_1,\dots,c_g\in A$ with $T^sz-\sum_jc_j\,\kappa_j\circ p_\infty\in\mathrm B^{\mathrm{od}}(A)$.
\end{enumerate}
\end{lemma}
\leanref{Development/AutomorphicArithmetic/AutomorphicVarieties/PatchingModules/PatchedSupportCompletions.lean}{UpperTheory.V.Pa.CohLim}
\leanref{Development/AutomorphicArithmetic/AutomorphicVarieties/PatchingModules/AutomorphicDetection/CohLim.lean}{UpperTheory.V.Pa.cohLim}
\leanref{Development/AutomorphicArithmetic/AutomorphicVarieties/PatchingModules/ThetaCoefficientMapEq.lean}{UpperTheory.V.Pa.K1}
\leanref{Development/AutomorphicArithmetic/AutomorphicVarieties/PatchingModules/ThetaCoefficientMapEq.lean}{UpperTheory.V.Pa.K2}
\leanref{Development/AutomorphicArithmetic/AutomorphicVarieties/PatchingModules/ThetaCoefficientMapEq.lean}{UpperTheory.V.Pa.K3}
\begin{proof}
Write $p=p_\infty$ and $\rho'=\rho\circ p$.

(3) Let $z\in\mathrm{Sel}^{\mathrm{od}}(A)$.

\emph{$z$ vanishes at the $\tau_j$.} As $\phi_j\tau_j=\tau_j\phi_j$, $p(\tau_j)=1$ and $p(\phi_j)=\sigma_j$, the cocycle relation gives $z(\phi_j)+\rho(\sigma_j)z(\tau_j)\rho(\sigma_j)^{-1}=z(\phi_j\tau_j)=z(\tau_j\phi_j)=z(\tau_j)+z(\phi_j)$, so $\rho(\sigma_j)$ commutes with $Z=z(\tau_j)$. Let $\Delta$ be the diagonal matrix with entries $1$ in the first block and $0$ in the second; it commutes with the even matrix $\rho(\sigma_j)$. Over an algebraic closure of $\mathcal K$, $\rho(\sigma_j)$ has the pairwise distinct eigenvalues $\alpha_{j,i}$, so the matrices commuting with it commute with each other; so $Z\Delta=\Delta Z$, and $Z$ is even. As $Z$ is odd, $Z=0$.

\emph{$z$ factors through $p$.} Let $H_z$ be the set of $k\in\ker p$ with $z(k)=0$. As $\rho'$ is trivial on $\ker p$, the cocycle relation shows that $H_z$ is a subgroup, that $z(gk)=z(g)$ for $k\in H_z$, and that $z(gkg^{-1})=\rho'(g)z(k)\rho'(g)^{-1}$ for $k\in\ker p$ and $g\in\Gamma_{\infty,0}$; with $z(c_\infty kc_\infty^{-1})=-{}^tz(k)$, $H_z$ is normal in $\Gamma_\infty$. It is closed, since each entry of $z$ has locally constant reductions modulo every power of $T$, so its zero set is closed. It contains the $\tau_j$, so it contains $\ker p$ by Lemma~\ref{sha-patching-chain-limit-bounds-kernel}. So $z(g)=z(g')$ whenever $p(g)=p(g')$, and $z=y\circ p$ for a map $y:\Gamma_0\to M_n(A)$, as $p$ maps $\Gamma_{\infty,0}$ onto $\Gamma_0$.

\emph{Conclusion.} The conditions defining the Selmer group of $x$ for $y$ are those for $z=y\circ p$, as $p$ is a continuous surjection onto $\Gamma_0$ compatible with $c_\infty$, the local maps and the $I'_u$; so $y$ lies in the Selmer group of $x$. By Definition~\ref{sha-patching-selmer}(2) there are $s\ge0$, $c_1,\dots,c_g\in A$ and an odd $Y$ with ${}^tY=-Y$ and $T^sy-\sum_jc_j\kappa_j=\delta Y$. Composing with $p$ gives $T^sz-\sum_jc_j\,\kappa_j\circ p\in\mathrm B^{\mathrm{od}}(A)$.

\emph{Notation for (1) and (2).} For $m\ge1$ write $A_m=A/T^mA$, and $\mathrm{Sel}^{\mathrm{ev}}_m$, $\mathrm{Sel}^1_m$, $\mathrm B^1_m$ and $\mathrm{Sel}^2_m$ for the sets defined as in Definition~\ref{sha-patching-chain-selmer}(2) and (3), with $A_m$ in place of $X$ and $R_g/T^mR_g$ in place of $R_{g,X}$; they are defined for $\Gamma_\infty$ with $p_\infty$, and in the same way for each $\Gamma_N$ with $p_N$. They are $A$-modules, since their conditions are $A$-linear; $\mathrm B^1_m\subset\mathrm{Sel}^1_m$, as the local condition for $\delta_1Y$ holds with the matrix $Y$ and the relation with $c_\infty$ follows from ${}^tY=-Y$ and ${}^t\rho_1(g)\rho_1(c_\infty gc_\infty^{-1})=\varepsilon(g)^{1-n}$; and homomorphisms of $A$-algebras, such as $A\to A_m$ and $A_m\to A_j$ for $j\le m$, carry each of these sets over one ring into the same set over the other, as the conditions are polynomial identities with coefficients from $A$. The group $\Gamma_\infty$ is topologically finitely generated, as a quotient of $\tilde\Gamma$.

(1) \emph{The bound modulo $T^m$.} Let $C$ be the constant of the hypothesis on Lemma~\ref{sha-patching-tangent}, and fix $m\ge1$. Lemma~\ref{sha-patching-chain-limit-bounds-transfer}, applied to $R=R_g/T^mR_g$, $\rho_1\bmod T^m$, the $w'_v$ and their inverses modulo $T^m$, and the first module $V$ over $A_m$, holds for $\mathcal U$-almost all $N$, and the hypothesis holds for all large $N$, a set in the nonprincipal ultrafilter $\mathcal U$; fix an $N$ in both sets. Composition with $\pi_\infty$ and with $\pi_N$ is injective on maps from $\Gamma_{\infty,0}$ and from $\Gamma_{N,0}$, as these maps are surjective on $\tilde\Gamma_0$. By (2) of that lemma, composition with $\pi_\infty$ maps $\mathrm{Sel}^1_m(\Gamma_\infty)$ into the image of $\mathrm{Sel}^1_m(\Gamma_N)$ under composition with $\pi_N$; by (1) of that lemma, composition with $\pi_N$ maps $\mathrm B^1_m(\Gamma_N)$ into the image of $\mathrm B^1_m(\Gamma_\infty)$, a finite set. So $\mathrm{Sel}^1_m(\Gamma_\infty)$ is finite, and
\[\#\mathrm{Sel}^1_m(\Gamma_\infty)\le\#\mathrm{Sel}^1_m(\Gamma_N)\le\#\mathrm B^1_m(\Gamma_N)\,|\kappa|^{mD_0+C}\le\#\mathrm B^1_m(\Gamma_\infty)\,|\kappa|^{mD_0+C}.\]

\emph{Independence modulo powers of $T$.} Let $z_1,\dots,z_e$ be as in (1). There is $C_1$ such that for all $m\ge C_1$ and $c\in A^e$, if the reduction of $\sum_ic_iz_i$ modulo $T^m$ lies in $\mathrm B^1_m$, then every $c_i$ lies in $T^{m-C_1}A$. Indeed, let $V_0$ be the $A$-module of matrices $Y$ in the first module over $A$ with ${}^tY=-Y$, and for $g\in\Gamma_{\infty,0}$ let $f_g:A^e\times V_0\to M_n(R_g)$ be the $A$-linear map $(c,Y)\mapsto\sum_ic_iz_i(g)-\delta_1Y(g)$, with $\delta_1Y(g)=\rho_1(p(g))Y\rho_1(p(g))^{-1}-Y$. Both modules are finite over $A$. If $f_g(c,Y)=0$ for all $g$, then $\sum_ic_iz_i\in\mathrm B^1(A)$, so $c=0$ by hypothesis; so the coordinates of $c$ vanish on $\bigcap_g\ker f_g$. Let $C_1$ be the constant of Lemma~\ref{sha-patching-chain-limit-bounds-artin-rees} for $I=TA$. Suppose that the reduction of $\sum_ic_iz_i$ modulo $T^m$ is $\delta_1\bar Y$, with $\bar Y$ in the first module over $A_m$ and ${}^t\bar Y=-\bar Y$. Choose an additive section of $A\to A_m$, and apply it to the coefficients of the entries of $\bar Y$; this gives $Y\in V_0$ lifting $\bar Y$, antisymmetric because the section is additive. Then $f_g(c,Y)$ reduces to $0$ modulo $T^m$, so it lies in $T^mM_n(R_g)$, for every $g$. Lemma~\ref{sha-patching-chain-limit-bounds-artin-rees}, with $m-C_1$ in place of $m$ and $\varphi$ the coordinates of $c$, gives $c_i\in T^{m-C_1}A$.

\emph{Counting.} Let $m\ge C_1$, and choose a section $\sigma$ of $A\to A/T^{m-C_1}A$. The map $\mathrm B^1_m\times(A/T^{m-C_1}A)^e\to\mathrm{Sel}^1_m$ that sends $(b,\bar c)$ to $b$ plus the reduction of $\sum_i\sigma(\bar c_i)z_i$ is injective: if $(b,\bar c)$ and $(b',\bar c')$ have the same image, the reduction of $\sum_i(\sigma(\bar c_i)-\sigma(\bar c'_i))z_i$ is $b'-b\in\mathrm B^1_m$, so $\sigma(\bar c_i)-\sigma(\bar c'_i)\in T^{m-C_1}A$ by the previous paragraph, and $\bar c=\bar c'$ and $b=b'$. As $A/T^{m-C_1}A$ has $|\kappa|^{m-C_1}$ elements and $\#\mathrm B^1_m\ge1$, the bound modulo $T^m$ gives $(m-C_1)e\le mD_0+C$ for every $m\ge C_1$. Suppose $e\ge D_0+1$, and let $K=C_1D_0+C+1$ and $m=C_1+K$. Then $(m-C_1)e\ge KD_0+K$, while $mD_0+C=KD_0+K-1$, a contradiction. So $e\le D_0$.

(2) \emph{The bound modulo $T^m$.} Let $C$ be the constant of the hypothesis on Lemma~\ref{sha-patching-obstruction}. For every $m\ge1$ and all $z_{i,j}\in\mathrm{Sel}^{\mathrm{ev}}_m(\Gamma_\infty)$, $0\le i\le g$, $1\le j\le g$, there are $c_0,\dots,c_g\in A$, one of them a unit, and $\Psi\in\mathrm{Sel}^2_m(\Gamma_\infty)$ whose matrix of coefficients of $\epsilon'_j$ is $T^C\sum_ic_iz_{i,j}$ for each $j$. Indeed, Lemma~\ref{sha-patching-chain-limit-bounds-transfer}, applied to $\rho\bmod T^m$ with the $w_v$ and the even matrices over $A_m$, and to $\rho_1\bmod T^m$ with the $w'_v$ and the second module over $A_m$, gives an $N$ at which the hypothesis and the conclusions of that lemma hold. It gives $z_{N,i,j}\in\mathrm{Sel}^{\mathrm{ev}}_m(\Gamma_N)$ with $z_{N,i,j}\circ\pi_N=z_{i,j}\circ\pi_\infty$ on $\tilde\Gamma_0$; the hypothesis gives $c$ and $\Psi_N\in\mathrm{Sel}^2_m(\Gamma_N)$ for these; and the lemma gives $\Psi\in\mathrm{Sel}^2_m(\Gamma_\infty)$ with $\Psi\circ\pi_\infty=\Psi_N\circ\pi_N$. As $\pi_\infty$ maps $\tilde\Gamma_0$ onto $\Gamma_{\infty,0}$, the coefficients of $\epsilon'_j$ of $\Psi$ are $T^C\sum_ic_iz_{i,j}$.

\emph{The limit.} Let $z_{i,j}\in\mathrm{Sel}^{\mathrm{ev}}(A)$. Their reductions modulo $T^m$ lie in $\mathrm{Sel}^{\mathrm{ev}}_m$; so for each $m$ the previous step gives $c^{(m)}\in A^{g+1}$, with a unit coordinate $c^{(m)}_{i(m)}$, and an element of $\mathrm{Sel}^2_m$ with coefficients of $\epsilon'_j$ equal to $T^C\sum_ic^{(m)}_iz_{i,j}$ modulo $T^m$; lift it, coefficient by coefficient, to a map $Z_m:\Gamma_{\infty,0}\to M_n(R_g)$. Let $\mathcal U'$ be an ultrafilter on $\mathbf N$ that contains the complements of the finite sets. Lemma~\ref{sha-patching-chain-limit-bounds-ultrafilter} gives $\Psi\in\mathrm{Sel}^2(A)$ with $\Psi\equiv Z_m$ modulo $T^J$ for $\mathcal U'$-almost all $m$, for each $J$. As $\kappa$ is finite, for each $i$ and $k$ the coefficient of $T^k$ in $c^{(m)}_i$ is constant for $\mathcal U'$-almost all $m$; let $c'_i\in A$ have these coefficients, so that $c^{(m)}_i\equiv c'_i$ modulo $T^J$ for $\mathcal U'$-almost all $m$, for each $J$. The index $i(m)$ takes finitely many values, so it equals some $i_0$ for $\mathcal U'$-almost all $m$. For such an $m$ with $c^{(m)}_{i_0}\equiv c'_{i_0}$ modulo $T$, $c'_{i_0}$ is a unit; so the elements $c_i=T^Cc'_i$ are not all zero. Let $J\ge1$, and choose $m\ge J$ in the $\mathcal U'$-almost sure sets above for $J$. Modulo $T^J$, the coefficients of $\epsilon'_j$ of $\Psi$ are those of $Z_m$, which are $T^C\sum_ic^{(m)}_iz_{i,j}\equiv\sum_ic_iz_{i,j}$. As $J$ is arbitrary, $\Psi'_j=\sum_ic_iz_{i,j}$ for all $j$.
% (1) and (2): the notation paragraph is sel_lin, coboundary_first_subset_selmer_first, sel_red, am_fin and rho_hom. The bounds modulo T^m are taylorWiles_first_condition_at_limit_of_group_condition (t1, Lemma sha-patching-chain-limit-bounds-transfer, with the private card_le_of_rel and Lambda.U_le) and taylorWiles_second_condition_at_limit_of_group_condition, giving FirstSelmerCardBoundAtAllLevels and SecondSelmerSpanBoundAtAllLevels for Lambda.h; independence modulo powers of T is coefficients_bounded_of_selmer_boundary (ar_family, Lemma sha-patching-chain-limit-bounds-artin-rees, and the private exists_section and exists_antisymmetric_lift_of_section); counting is firstCohomologyBound_of_boundaryBound (the private card_submodule_mul_quotient_le); the limit in (2) is secondCohomologyBound_of_boundaryBound (Ultrafilter.of atTop, selmer_second_exists_ultrafilter_limit, the private power_series_exists_ultrafilter_limit and power_series_eq_of_dvd_all_powers). (3) is koszulThree_of_det: od_comm (with the private mat_comm_of_sep), coc_ker (H with hHconj, hHc and the private closed_zero), ker_oo (Lemma sha-patching-chain-limit-bounds-kernel), Glob.Hom.sel_iff, k.gen and Glob.Hom.bd_iff.
% cohLim: Lambda.G is topologically finitely generated (TopFG.subgroup_quot from Lambda.fg); (1) firstCohomologyBound_of_boundaryBound applied to taylorWiles_first_condition_at_limit_of_group_condition (from T.B1); (2) secondCohomologyBound_of_boundaryBound applied to taylorWiles_second_condition_at_limit_of_group_condition (from T.B2); (3) koszulThree_of_det (from T.Eig and ProL). All in SeparableMatrixStability.
\end{proof}

\begin{lemma}\label{sha-patching-chain-laurent-scaling}
Let $\Gamma_\infty$, $p$, $x$, $\kappa_1,\dots,\kappa_g$, $A$, $\mathcal K$, $R_g$ and $R_{g,\mathcal K}$ be as in Definition~\ref{sha-patching-chain-selmer}.
\begin{enumerate}
\item Every constant map $\Gamma_{\infty,0}\to\mathcal K$ has bounded denominators. If $\phi:\Gamma_{\infty,0}\to\mathcal K$ has bounded denominators, there is $s_0\ge0$ such that for every $s\ge s_0$, $T^s\phi$ is a map to $A$ whose reductions modulo the powers of $T$ are locally constant.
\item Let $z$ be an element of $\mathrm{Sel}^{\mathrm{ev}}(\mathcal K)$, of $\mathrm{Sel}^{\mathrm{od}}(\mathcal K)$ or of $\mathrm{Sel}^1(\mathcal K)$ with bounded denominators. Then there is $s_0\ge0$ such that for every $s\ge s_0$, $T^sz$ is the image of an element of $\mathrm{Sel}^{\mathrm{ev}}(A)$, of $\mathrm{Sel}^{\mathrm{od}}(A)$, respectively of $\mathrm{Sel}^1(A)$.
\item The image in $M_n(R_{g,\mathcal K})$ of an element of $\mathrm{Sel}^2(A)$ lies in $\mathrm{Sel}^2(\mathcal K)$ and has bounded denominators.
\end{enumerate}
\end{lemma}
\leanref{Development/AutomorphicArithmetic/AutomorphicVarieties/PatchingModules/SelmerScaling/BndOps.lean}{UpperTheory.V.Pa.bndOps}
\leanref{Development/AutomorphicArithmetic/AutomorphicVarieties/PatchingModules/SelmerScaling/SelScale.lean}{UpperTheory.V.Pa.selScale}
\leanref{Development/AutomorphicArithmetic/AutomorphicVarieties/PatchingModules/SelmerScaling/SelScale1.lean}{UpperTheory.V.Pa.selScale1}
\leanref{Development/AutomorphicArithmetic/AutomorphicVarieties/PatchingModules/SelmerScaling/AlgebraMapMemSecondOrder.lean}{UpperTheory.V.Pa.selMap2}
\begin{proof}
(1) Let $c=a/b$ with $a,b\in A$ and $b\ne0$. Then $b=T^kv$ with $k\ge0$ and $v$ a unit of $A$, so $T^kc=av^{-1}$ lies in $A$, and a constant map has locally constant reductions. Now let $T^{s_1}\phi=\phi'$, with $\phi'$ a map to $A$ whose reductions modulo the powers of $T$ are locally constant. For $s\ge s_1$, $T^s\phi=T^{s-s_1}\phi'$, and multiplying by $T^{s-s_1}$ keeps the reductions locally constant.

(2) In the even and odd cases let $A'=A$, $X'=\mathcal K$, $\rho'=\rho\circ p$ and $w''=w$; in the case of $\mathrm{Sel}^1$ let $A'=R_g$, $X'=R_{g,\mathcal K}$, $\rho'=\rho_1\circ p$ and $w''=w'$. In each case $\rho'$ and the $w''_v$ come from $A'$, and the map $A'\to X'$ is injective, as it is so on the coefficients in the basis $1,\epsilon_j,\epsilon'_j,\eta$. For each $v\mid\ell$ choose a matrix $Y_v$ over $X'$, in the module $V$ of $z$, as in the local condition at $v$. The coefficients of the entries of $z$ have bounded denominators, and so have those of the entries of the $Y_v$, viewed as constant maps, by (1). By (1), applied to these finitely many maps, there is $s_0$ such that for every $s\ge s_0$ there are a map $z'$ from $\Gamma_{\infty,0}$ to the matrices over $A'$ and matrices $Y'_v$ over $A'$ whose images are $T^sz$ and $T^sY_v$, where the coefficients of the entries of $z'$ have locally constant reductions modulo every power of $T$. So the reductions of $z'$ modulo the powers of $T$ are locally constant. An entry, or a coefficient of an entry, that the module $V$ requires to vanish maps to $T^s$ times the corresponding entry or coefficient of $z$ or $Y_v$, which is $0$; so it vanishes, by injectivity, and the values of $z'$ and the $Y'_v$ lie in $V$. The cocycle relation, the relation with $c_\infty$, the vanishing on the $I'_u$, and the triangularity and the vanishing of the diagonal in the local condition are linear in $z$ and $Y_v$, with coefficients $\rho'(g)$ and $w''_v$, ${w''_v}^{-1}$ that come from $A'$. Mapped to $X'$, those for $z'$ and $Y'_v$ become $T^s$ times those for $z$ and $Y_v$, which hold. So they hold for $z'$ and $Y'_v$ by injectivity, and $z'$ lies in the Selmer group over $A$.

(3) The map $R_g\to R_{g,\mathcal K}$ is a ring homomorphism that carries $\rho_1\circ p$ and the $w'_v$ over $R_g$ to those over $R_{g,\mathcal K}$ and the second module $V$ of Definition~\ref{sha-patching-chain-selmer}(3) over $A$ into that over $\mathcal K$. It carries $T^mR_g$ into $T^mR_{g,\mathcal K}$, so reductions modulo $T^m$ stay locally constant. Applying it to an element $\Psi$ of $\mathrm{Sel}^2(A)$ and to its matrices $Y_v$ gives an element of $\mathrm{Sel}^2(\mathcal K)$. The coefficients of its entries are the images of those of $\Psi$, which are maps to $A$ with locally constant reductions modulo every power of $T$, since $T^mR_g$ consists of the elements whose coefficients lie in $T^mA$; so they have bounded denominators, with no power of $T$ needed.
% (1) bndOps: hconst (isUnit_divided_by_X_pow_order) and the last conjunct. (2) selScale (even and odd, the Bool b, with liftF) and selScale1 (Sel1, with liftM, adicDD and mapInj); both conclude by selLift, which checks each condition of SelZ through the injective map. (3) selMap2, with selMap2_lemma, dd_components_mem_span_pow_of_mem_span_algebraMap_pow and the maps kap_rho1/w1/wi1_map_dd_map_algebraMap. About 500 lines.
\end{proof}

\begin{lemma}\label{sha-patching-chain-laurent}
Let $\Gamma_\infty$, $p$, $x$, $\kappa_1,\dots,\kappa_g$ and $\mathcal K$ be as in Definition~\ref{sha-patching-chain-selmer}, and let $D_0\ge0$.
\begin{enumerate}
\item Suppose that (1) of Lemma~\ref{sha-patching-chain-limit-bounds} holds for $p$ and $D_0$. If $z_1,\dots,z_e\in\mathrm{Sel}^1(\mathcal K)$ have bounded denominators and $\sum_ic_iz_i\in\mathrm B^1(\mathcal K)$, with $c_i\in\mathcal K$, only if all $c_i$ vanish, then $e\le D_0$.
\item Suppose that (2) of Lemma~\ref{sha-patching-chain-limit-bounds} holds for $p$. If $z_{i,j}\in\mathrm{Sel}^{\mathrm{ev}}(\mathcal K)$, for $0\le i\le g$ and $1\le j\le g$, have bounded denominators, then there are $c_0,\dots,c_g\in\mathcal K$, not all zero, such that the family $(\sum_ic_iz_{i,j})_{1\le j\le g}$ lifts.
\item Suppose that (3) of Lemma~\ref{sha-patching-chain-limit-bounds} holds for $p$. If $z\in\mathrm{Sel}^{\mathrm{od}}(\mathcal K)$ has bounded denominators, then there are $c_1,\dots,c_g\in\mathcal K$ with $z-\sum_jc_j\,\kappa_j\circ p\in\mathrm B^{\mathrm{od}}(\mathcal K)$.
\end{enumerate}
\end{lemma}
\leanref{Development/AutomorphicArithmetic/AutomorphicVarieties/PatchingModules/SelmerScaling/SelmerRankBoundLaurent.lean}{UpperTheory.V.Pa.k1F}
\leanref{Development/AutomorphicArithmetic/AutomorphicVarieties/PatchingModules/SelmerScaling/EvenCocyclesLiftLaurent.lean}{UpperTheory.V.Pa.k2F}
\leanref{Development/AutomorphicArithmetic/AutomorphicVarieties/PatchingModules/SelmerScaling/SelmerSpannedByKappaLaurentAlt.lean}{UpperTheory.V.Pa.k3F}
\leanref{Development/AutomorphicArithmetic/AutomorphicVarieties/PatchingModules/SelmerScaling/SelmerSpannedByKappaLaurent.lean}{UpperTheory.V.Pa.SelmerRankBoundLaurent}
\leanref{Development/AutomorphicArithmetic/AutomorphicVarieties/PatchingModules/SelmerScaling/SelmerSpannedByKappaLaurent.lean}{UpperTheory.V.Pa.EvenCocyclesLiftLaurent}
\leanref{Development/AutomorphicArithmetic/AutomorphicVarieties/PatchingModules/SelmerScaling/SelmerSpannedByKappaLaurent.lean}{UpperTheory.V.Pa.SelmerSpannedByKappaLaurent}
\begin{proof}
The map $A\to\mathcal K$ is injective and $T\ne0$ in $\mathcal K$.

(1) By Lemma~\ref{sha-patching-chain-laurent-scaling}(2), applied to each $z_i$, there are $s\ge0$ and $z'_1,\dots,z'_e\in\mathrm{Sel}^1(A)$ whose images are $T^sz_1,\dots,T^sz_e$; take $s$ at least each of the bounds $s_0$. Let $c_1,\dots,c_e\in A$ with $\sum_ic_iz'_i=\delta Y$, where $Y$ lies in the first module $V$ of Definition~\ref{sha-patching-chain-selmer}(3) over $A$, ${}^tY=-Y$, and $\delta Y(g)=\rho_1(p(g))Y\rho_1(p(g))^{-1}-Y$. The image of $Y$ in $M_n(R_{g,\mathcal K})$ lies in that module over $\mathcal K$ and is antisymmetric, and mapping the relation gives $\sum_i(c_iT^s)z_i\in\mathrm B^1(\mathcal K)$. By hypothesis all $c_iT^s$ vanish in $\mathcal K$, so all $c_i$ vanish. So $z'_1,\dots,z'_e$ satisfy the hypothesis of Lemma~\ref{sha-patching-chain-limit-bounds}(1), and $e\le D_0$.

(2) By Lemma~\ref{sha-patching-chain-laurent-scaling}(2), there are $S\ge0$ and $z'_{i,j}\in\mathrm{Sel}^{\mathrm{ev}}(A)$ whose images are the $T^Sz_{i,j}$. By Lemma~\ref{sha-patching-chain-limit-bounds}(2) there are $c_0,\dots,c_g\in A$, not all zero, and $\Psi\in\mathrm{Sel}^2(A)$ with $\Psi'_j=\sum_ic_iz'_{i,j}$ for all $j$. By Lemma~\ref{sha-patching-chain-laurent-scaling}(3), the image of $\Psi$ lies in $\mathrm{Sel}^2(\mathcal K)$ and has bounded denominators; the coefficient of $\epsilon'_j$ in its entries is the image of $\Psi'_j$, which is $\sum_i(c_iT^S)z_{i,j}$. So the family $(\sum_ic''_iz_{i,j})_j$ lifts, with $c''_i=c_iT^S\in\mathcal K$, which are not all zero.

(3) By Lemma~\ref{sha-patching-chain-laurent-scaling}(2), there are $s_0\ge0$ and $z'\in\mathrm{Sel}^{\mathrm{od}}(A)$ whose image is $T^{s_0}z$. By Lemma~\ref{sha-patching-chain-limit-bounds}(3) there are $s'\ge0$, $c_1,\dots,c_g\in A$ and an odd matrix $Y$ over $A$ with ${}^tY=-Y$ and $T^{s'}z'(g)-\sum_jc_j\kappa_j(p(g))=\rho(p(g))Y\rho(p(g))^{-1}-Y$ for all $g$. Mapping to $\mathcal K$ and multiplying by $u=T^{-s'-s_0}$ gives $z-\sum_j(uc_j)\,\kappa_j\circ p=\delta(uY)$, where $uY$ is odd and antisymmetric. So $z-\sum_j(uc_j)\,\kappa_j\circ p\in\mathrm B^{\mathrm{od}}(\mathcal K)$.
% k1F: selScale1 with s = sup of the s0, then K1 (the independence transported by key and hc'); dual_numbers_map_smul, map_mem_dual_numbers_part and map_first_order_lift_residue_field for the image of Y. k2F: selScale (b = true) with S = sup, K2, then selMap2, with c i * t^S. k3F: selScale (b = false), K3, key2, and u = (t^(s' + s0))^(-1).
\end{proof}

\begin{definition}\label{sha-patching-chain-tangent}
Let $x$ be a normalized arc over $A=\kappa[[T]]$, with representation $\rho$, blocks of sizes $a$ and $n-a$ and a Selmer basis $\kappa_1,\dots,\kappa_g$, let $(q,\phi_j,\tau_j,U_{j,i},y_{j,i},s,f)$ be an arc datum for $x$ over $(\mathcal R,D,\psi)$, in the sense of Definition~\ref{sha-patching-chain-datum}, and let $\mathcal K$, $\hat{\mathcal R}$ and $\hat f$ be as in Lemma~\ref{sha-patching-chain-representation}. Let $\mathcal A=\hat{\mathcal R}\cdot1\oplus\mathcal N$ and $r$ satisfy the conclusions of that lemma, and write $y\mapsto(y_0,y_1)$ for the projections of $\mathcal A$ onto $\hat{\mathcal R}$ and $\mathcal N$. Let $\mathcal O'$ be a complete discrete valuation ring with a homomorphism $\mathcal O'\to\hat{\mathcal R}$ that makes $\hat{\mathcal R}$ a complete local noetherian $\mathcal O'$-algebra whose residue field is reached from $\mathcal O'$, such that the images in $\hat{\mathcal R}$ of $\mathcal O$ and of $s$ come from $\mathcal O'$, every element of $A$ is $f(b)$ for some $b\in\mathcal R$ whose image in $\hat{\mathcal R}$ comes from $\mathcal O'$, and $\hat f$ maps $\mathcal O'$ onto $\mathcal K$.
\begin{enumerate}
\item $I_{\mathcal N}\subset\hat{\mathcal R}$ is the ideal generated by the $d$ with $d\cdot1=yz$ for some $y,z\in\mathcal N$, and $J=\mathfrak m_{\mathcal O'}\hat{\mathcal R}+\mathfrak m_{\hat{\mathcal R}}^2+I_{\mathcal N}$.
\item The \emph{tangent space} $\mathcal T$ is the $\mathcal K$-vector space of additive maps $\alpha:\hat{\mathcal R}\to\mathcal K$ with $\alpha(bb')=\hat f(b)\alpha(b')+\hat f(b')\alpha(b)$, vanishing on the image of $\mathcal O'$ and on $I_{\mathcal N}$.
\item $\Lambda$ is the $\mathcal K$-vector space of additive maps $\mu:\mathcal N\to\mathcal K$ with $\mu(by)=\hat f(b)\mu(y)$ for $b\in\hat{\mathcal R}$.
\item $\Pi$ is the $\mathcal K$-vector space of additive maps $\Phi:\mathcal N\to\mathcal K$ with $\Phi(oy)=\hat f(o)\Phi(y)$ for $o$ in the image of $\mathcal O'$ and $\Phi(by)=0$ for $b\in J$.
\item For additive maps $\alpha:\hat{\mathcal R}\to\mathcal K$ and $\mu:\mathcal N\to\mathcal K$, let $\alpha_C$ and $\mu_C$ be the maps $\Gamma_{\infty,0}\to M_n(\mathcal K)$ with
\[\alpha_C(g)=\big(\alpha(r(g)_{ij,0})\big)_{ij}\,\rho(p(g))^{-1},\qquad\mu_C(g)=\big(\mu(r(g)_{ij,1})\big)_{ij}\,\rho(p(g))^{-1}.\]
\item $\mathcal V\subset\mathcal T^g$ is the set of families $(\alpha_1,\dots,\alpha_g)$ such that $(\alpha_{1,C},\dots,\alpha_{g,C})$ lifts, in the sense of Definition~\ref{sha-patching-chain-selmer}(5).
\item Elements $m_1,\dots,m_g$ of $\mathcal N$ are \emph{dual to the Selmer basis} if $\mu_C-\sum_j\mu(m_j)\,\kappa_j\circ p\in\mathrm B^{\mathrm{od}}(\mathcal K)$ for every $\mu\in\Lambda$.
\end{enumerate}
\end{definition}
\leanref{Development/AutomorphicArithmetic/AutomorphicVarieties/PatchingModules/SelmerScaling/SelmerSpannedByKappaLaurent.lean}{UpperTheory.V.Pa.Ctx}
\leanref{Development/AutomorphicArithmetic/AutomorphicVarieties/PatchingModules/SelmerScaling/SelmerSpannedByKappaLaurent.lean}{UpperTheory.V.Pa.ired}
\leanref{Development/AutomorphicArithmetic/AutomorphicVarieties/PatchingModules/SelmerScaling/SelmerSpannedByKappaLaurent.lean}{UpperTheory.V.Pa.Ctx.J2}
\leanref{Development/AutomorphicArithmetic/AutomorphicVarieties/PatchingModules/SelmerScaling/SelmerSpannedByKappaLaurent.lean}{UpperTheory.V.Pa.Ctx.Tp}
\leanref{Development/AutomorphicArithmetic/AutomorphicVarieties/PatchingModules/SelmerScaling/SelmerSpannedByKappaLaurent.lean}{UpperTheory.V.Pa.Ctx.La}
\leanref{Development/AutomorphicArithmetic/AutomorphicVarieties/PatchingModules/SelmerScaling/SelmerSpannedByKappaLaurent.lean}{UpperTheory.V.Pa.Ctx.Ph}
\leanref{Development/AutomorphicArithmetic/AutomorphicVarieties/PatchingModules/SelmerScaling/SelmerSpannedByKappaLaurent.lean}{UpperTheory.V.Pa.Ctx.alC}
\leanref{Development/AutomorphicArithmetic/AutomorphicVarieties/PatchingModules/SelmerScaling/SelmerSpannedByKappaLaurent.lean}{UpperTheory.V.Pa.Ctx.laC}
\leanref{Development/AutomorphicArithmetic/AutomorphicVarieties/PatchingModules/SelmerScaling/SelmerSpannedByKappaLaurent.lean}{UpperTheory.V.Pa.Ctx.DualM}

\begin{lemma}\label{sha-patching-chain-duality}
In the setting of Definition~\ref{sha-patching-chain-tangent}:
\begin{enumerate}
\item $I_{\mathcal N}\subset\mathfrak m_{\hat{\mathcal R}}$ and $J\subset\mathfrak m_{\hat{\mathcal R}}$;
\item $\mathcal T$ is finite-dimensional; an element $b\in\mathfrak m_{\hat{\mathcal R}}$ with $\alpha(b)=0$ for all $\alpha\in\mathcal T$ lies in $J$; and for every linear form $\xi$ on $\mathcal T$ there is $b\in\mathfrak m_{\hat{\mathcal R}}$ with $\alpha(b)=\xi(\alpha)$ for all $\alpha\in\mathcal T$;
\item $\Lambda$ is finite-dimensional; an element $y\in\mathcal N$ with $\mu(y)=0$ for all $\mu\in\Lambda$ lies in $\mathfrak m_{\hat{\mathcal R}}\mathcal N$; and for every linear form $\xi$ on $\Lambda$ there is $y\in\mathcal N$ with $\mu(y)=\xi(\mu)$ for all $\mu\in\Lambda$;
\item an element $y\in\mathcal N$ with $\Phi(y)=0$ for all $\Phi\in\Pi$ lies in $J\mathcal N$;
\item if $\mu\in\Lambda$ and $\mu_C\in\mathrm B^{\mathrm{od}}(\mathcal K)$, then $\mu=0$;
\item if $c_1,\dots,c_g\in\mathcal K$ and $\sum_jc_j\,\kappa_j\circ p\in\mathrm B^{\mathrm{od}}(\mathcal K)$, then all $c_j$ vanish.
\end{enumerate}
\end{lemma}
\leanref{Development/AutomorphicArithmetic/AutomorphicVarieties/PatchingModules/SelmerScaling/PrAdd.lean}{UpperTheory.V.Pa.prAdd}
\leanref{Development/AutomorphicArithmetic/AutomorphicVarieties/PatchingModules/SelmerScaling/TpDual.lean}{UpperTheory.V.Pa.tpDual}
\leanref{Development/AutomorphicArithmetic/AutomorphicVarieties/PatchingModules/SelmerScaling/LaDual.lean}{UpperTheory.V.Pa.laDual}
\leanref{Development/AutomorphicArithmetic/AutomorphicVarieties/PatchingModules/SelmerScaling/ExistsSemilinearFunctionalNonzero.lean}{UpperTheory.V.Pa.phSep}
\leanref{Development/AutomorphicArithmetic/AutomorphicVarieties/PatchingModules/SelmerScaling/LaInj.lean}{UpperTheory.V.Pa.laInj}
\leanref{Development/AutomorphicArithmetic/AutomorphicVarieties/PatchingModules/SelmerScaling/KapF.lean}{UpperTheory.V.Pa.indF}
\begin{proof}
The map $\hat f$ is surjective with kernel $\mathfrak m_{\hat{\mathcal R}}$, since it maps $\mathcal O'$ onto $\mathcal K$; and its restriction to $\mathcal O'$ has kernel $\mathfrak m_{\mathcal O'}$, a maximal ideal of the local ring $\mathcal O'$. As $\mathcal A$ is finite over $\hat{\mathcal R}$ and $\mathcal N$ is a direct summand, $\mathcal N$ is a finite $\hat{\mathcal R}$-module.

(1) If $d\cdot1=yz$ with $y,z\in\mathcal N$, then $yz=c\cdot1$ with $c\in\mathfrak m_{\hat{\mathcal R}}$, and $d=c$ as $\hat{\mathcal R}$ injects into $\mathcal A$. So $I_{\mathcal N}\subset\mathfrak m_{\hat{\mathcal R}}$. The map $\mathcal O'\to\hat{\mathcal R}$ is local, so $J\subset\mathfrak m_{\hat{\mathcal R}}$.

(2) Every $b\in\hat{\mathcal R}$ is $o+y$ with $o$ in the image of $\mathcal O'$ and $y\in\mathfrak m_{\hat{\mathcal R}}$. Let $W$ be the quotient of $\hat{\mathcal R}$ by the sum of $J$ and the image of $\mathcal O'$, and $\pi:\hat{\mathcal R}\to W$ the quotient map. The image of $\mathfrak m_{\mathcal O'}$ acts by $0$ on $W$, so $W$ is a vector space over $\mathcal O'/\mathfrak m_{\mathcal O'}\cong\mathcal K$, with $\mathcal K$ acting through $\hat f$. Writing $b=o+y$ and $b'=o'+y'$ as above, with $yy'\in\mathfrak m_{\hat{\mathcal R}}^2\subset J$, gives $\pi(bb')=\hat f(b)\pi(b')+\hat f(b')\pi(b)$. The space $W$ is spanned by the images of generators of $\mathfrak m_{\hat{\mathcal R}}$, by this product rule; so it is finite-dimensional, as $\hat{\mathcal R}$ is noetherian. For a linear form $\lambda$ on $W$, $\lambda\circ\pi$ lies in $\mathcal T$, as $\pi$ kills $J\supset I_{\mathcal N}$ and the image of $\mathcal O'$. Conversely, let $\alpha\in\mathcal T$. It kills $\mathfrak m_{\mathcal O'}\hat{\mathcal R}$, since $\alpha(ob)=\hat f(o)\alpha(b)+\hat f(b)\alpha(o)=0$ for $o$ in the image of $\mathfrak m_{\mathcal O'}$; it kills $\mathfrak m_{\hat{\mathcal R}}^2$, since $\hat f$ kills $\mathfrak m_{\hat{\mathcal R}}$; and it kills $I_{\mathcal N}$ and the image of $\mathcal O'$. So $\alpha=\lambda\circ\pi$ for a unique additive $\lambda:W\to\mathcal K$, which is $\mathcal K$-linear by the product rule. Hence $\mathcal T$ is the dual of $W$, and it is finite-dimensional. An element $b\in\mathfrak m_{\hat{\mathcal R}}$ killed by $\mathcal T$ has $\pi(b)=0$, so $b=j+o$ with $j\in J$ and $o$ in the image of $\mathcal O'$; then $o=b-j\in\mathfrak m_{\hat{\mathcal R}}$ comes from $\mathfrak m_{\mathcal O'}$, as $\mathcal O'\to\hat{\mathcal R}$ is local, so $o\in J$ and $b\in J$. A linear form on $\mathcal T$ is the evaluation at some $w\in W$, and $w=\pi(b)$ for some $b\in\mathfrak m_{\hat{\mathcal R}}$.

(3) An element of $\Lambda$ is determined by its values on generators of the finite $\hat{\mathcal R}$-module $\mathcal N$, so $\Lambda$ is finite-dimensional. The module $\mathcal N/\mathfrak m_{\hat{\mathcal R}}\mathcal N$ is a vector space over $\hat{\mathcal R}/\mathfrak m_{\hat{\mathcal R}}\cong\mathcal K$, and the composites of its linear forms with $\mathcal N\to\mathcal N/\mathfrak m_{\hat{\mathcal R}}\mathcal N$ lie in $\Lambda$. If $y\notin\mathfrak m_{\hat{\mathcal R}}\mathcal N$, some linear form is nonzero at the image of $y$, so some $\mu\in\Lambda$ has $\mu(y)\ne0$. The linear forms on $\Lambda$ of the form $\mu\mapsto\mu(y)$, $y\in\mathcal N$, make up a subspace $W'$ of the dual of $\Lambda$, closed under multiplication by $\hat f(b)$ through $y\mapsto by$. If $\xi\notin W'$, there is a linear form on the dual of $\Lambda$ that vanishes on $W'$ and not at $\xi$; it is the evaluation at some $\mu\in\Lambda$, as $\Lambda$ is finite-dimensional, and then $\mu(y)=0$ for all $y$, so $\mu=0$, which is impossible.

(4) Suppose $y\notin J\mathcal N$. The module $\mathcal N/J\mathcal N$ is killed by the image of $\mathfrak m_{\mathcal O'}$, so it is a vector space over $\mathcal O'/\mathfrak m_{\mathcal O'}\cong\mathcal K$. A linear form on it that is nonzero at the image of $y$, composed with $\mathcal N\to\mathcal N/J\mathcal N$, lies in $\Pi$ and is nonzero at $y$.

(5) Let $\mu_C(g)=\rho(p(g))Y\rho(p(g))^{-1}-Y$. By Lemma~\ref{sha-patching-chain-representation}(3), $\hat f(r(g)_{ij,0})=\rho(p(g))_{ij}$. For $Z$ in the $\hat{\mathcal R}$-span of $r(\Gamma_{\infty,0})$, write $Z_0=(\hat f(Z_{ij,0}))_{ij}$ and $Z_1=(\mu(Z_{ij,1}))_{ij}$; both are additive in $Z$ and are multiplied by $\hat f(b)$ when $Z$ is multiplied by $b\in\hat{\mathcal R}$. For $Z=r(g)$, $Z_1=\mu_C(g)\rho(p(g))=Z_0Y-YZ_0$; so $Z_1=Z_0Y-YZ_0$ on the whole span. If $a=n$, then $\mathcal N=0$ and $\mu=0$. Otherwise let $y\in\mathcal N$, and let $Z$ be the matrix with entry $y$ at $(1,a+1)$ and $0$ elsewhere, which lies in the span by Lemma~\ref{sha-patching-chain-representation}(1). Then $Z_0=0$, so $\mu(y)=(Z_1)_{1,a+1}=0$.

(6) Let $\sum_jc_j\,\kappa_j\circ p$ be the map $g\mapsto\rho(p(g))Y\rho(p(g))^{-1}-Y$ with $Y$ odd and ${}^tY=-Y$. Choose a nonzero $d\in A$ with $dc_j\in A$ for all $j$ and with $dY$ over $A$. As $p$ maps $\Gamma_{\infty,0}$ onto $\Gamma_0$, $\sum_j(dc_j)\kappa_j$ is the coboundary of $dY$ over $A$. By the independence in Definition~\ref{sha-patching-selmer}(2), $dc_j=0$, so $c_j=0$.
% (1) prAdd (ired <= maximal ideal from U.mul_mem and U.inj; J2 <= maximal ideal from IsCO OO Bh). (2) tpDual = tp_abs with dual_of_univ (N = J2 + O', the derivation hder, finiteness from IsNoetherian). (3) laDual = dual_functionals_finite_separating_surjective (hzero, hfin, hkill, and W with Submodule.exists_dual_map_eq_bot_of_notMem). (4) phSep with exists_semilinear_functional_nonzero. (5) laInj (key by span induction, prVr for the reduction, the matrix single 0 a). (6) indF (IsLocalization.exist_integer_multiples_of_finite, k.indep).
\end{proof}

\begin{lemma}\label{sha-patching-chain-cocycles-ordinary}
Let $F$ be an infinite field, $R$ a commutative $F$-algebra, $e:R\to F$ a homomorphism of $F$-algebras whose kernel consists of nilpotent elements, and $I\subset R$ an ideal with $I^2=0$. Let $H$ be a group, $n\ge0$, and $E$ an $F$-subalgebra of $M_n(R)$. Let $\rho_0,\rho_1:H\to GL_n(R)$ be homomorphisms, $\psi_{0,1},\dots,\psi_{0,n}$ and $\psi_{1,1},\dots,\psi_{1,n}$ homomorphisms $H\to R^\times$, and $w\in GL_n(R)$.
\begin{enumerate}
\item Suppose that:
\begin{enumerate}
\item[(a)] $w^{-1}\rho_0(g)w$ is upper triangular with diagonal $\psi_{0,1}(g),\dots,\psi_{0,n}(g)$ for $g\in H$;
\item[(b)] the entries of $\rho_1(g)-\rho_0(g)$ and the elements $\psi_{1,i}(g)-\psi_{0,i}(g)$ lie in $I$, for $g\in H$ and all $i$;
\item[(c)] the homomorphisms $e\circ\psi_{0,1},\dots,e\circ\psi_{0,n}:H\to F^\times$ are pairwise distinct;
\item[(d)] $\prod_{i=1}^n\big(\rho_1(g_i)-\psi_{1,i}(g_i)\big)=0$ for all $g_1,\dots,g_n\in H$, the product taken in increasing order of $i$;
\item[(e)] $E$ contains the matrices $\rho_0(g)$ and $\rho_1(g)$ and the scalar matrices $\psi_{0,i}(g)$ and $\psi_{1,i}(g)$, for $g\in H$ and all $i$.
\end{enumerate}
Then there is $Y\in E$ with entries in $I$ such that, for every $g\in H$, the matrix
\[w^{-1}\Big(\big(\rho_1(g)\rho_0(g)^{-1}-1\big)-\big(\rho_0(g)Y\rho_0(g)^{-1}-Y\big)\Big)w\]
is upper triangular with $(i,i)$ entry $\big(\psi_{1,i}(g)-\psi_{0,i}(g)\big)\psi_{0,i}(g)^{-1}$ for every $i$.
\item In particular, let $F[\varepsilon]$ be the ring of dual numbers, $E$ an $F$-subalgebra of $M_n(F[\varepsilon])$, $\rho:H\to GL_n(F)$ and $U:H\to GL_n(F[\varepsilon])$ homomorphisms, and $\psi_1,\dots,\psi_n:H\to F^\times$ and $\chi_1,\dots,\chi_n:H\to F[\varepsilon]^\times$ homomorphisms, with $U(g)=\rho(g)+\varepsilon U'(g)$ and $\chi_i(g)=\psi_i(g)+\varepsilon\chi'_i(g)$. Suppose that the $\psi_i$ are pairwise distinct; that there is $w\in GL_n(F)$ such that $w^{-1}\rho(g)w$ is upper triangular with diagonal $\psi_1(g),\dots,\psi_n(g)$ for $g\in H$; that $\prod_{i=1}^n\big(U(g_i)-\chi_i(g_i)\big)=0$ for all $g_1,\dots,g_n\in H$, the product taken in increasing order of $i$; and that $E$ contains the matrices $\rho(g)$ and $U(g)$ and the scalar matrices $\psi_i(g)$ and $\chi_i(g)$, for $g\in H$. Then there is $Y\in M_n(F)$ with $\varepsilon Y\in E$ such that, for every $g\in H$, $w^{-1}\big(U'(g)\rho(g)^{-1}-\rho(g)Y\rho(g)^{-1}+Y\big)w$ is upper triangular with diagonal $\chi'_1(g)\psi_1(g)^{-1},\dots,\chi'_n(g)\psi_n(g)^{-1}$.
\end{enumerate}
\end{lemma}
\leanref{Development/AutomorphicArithmetic/AutomorphicVarieties/PatchingModules/SelmerScaling/LocLemma.lean}{UpperTheory.V.Pa.locLemma}
\leanref{Development/AutomorphicArithmetic/AutomorphicVarieties/PatchingModules/SelmerScaling/DefUSel.lean}{locDual}
\begin{proof}
(1) As $I^2=0$, $e$ kills $I$: for $c\in I$, $e(c)^2=e(c^2)=0$. Let $\pi:M_n(R)\to M_n(R/I)$ be the reduction; the product of two matrices with entries in $I$ is $0$. Extend $\rho_0$, $\rho_1$ and the $\psi_{0,i}$, $\psi_{1,i}$ linearly to the group algebra $R[H]$. As in the proof of Lemma~\ref{sha-hecke-ideal-triangular}, the relation (d) then holds for all $y_1,\dots,y_n\in R[H]$, since it is multilinear.

\emph{(i) The relation for $\rho_0$.} After conjugation by $w$, the factor $\rho_0(g_k)-\psi_{0,k}(g_k)$ is upper triangular with $(k,k)$ entry $0$, by (a), so it maps the span of the first $k$ standard basis vectors into that of the first $k-1$. So $\prod_k(\rho_0(g_k)-\psi_{0,k}(g_k))=0$, and this extends to $R[H]$ as well.

\emph{(ii) A separating element.} The $e\circ\psi_{0,i}$ are pairwise distinct homomorphisms $H\to F^\times$, by (c), hence linearly independent over $F$; so $y\mapsto(e(\psi_{0,i}(y)))_i$ maps the $F$-span of $H$ in $R[H]$ onto $F^n$. As $F$ is infinite, there is $z$ in that span with $e(\psi_{0,1}(z)),\dots,e(\psi_{0,n}(z))$ pairwise distinct. As (a), (b) and (e) are $F$-linear in $g$, they hold at $z$: $w^{-1}\rho_0(z)w$ is upper triangular with diagonal $\psi_{0,i}(z)$; $\rho_1(z)-\rho_0(z)$ and $\psi_{1,i}(z)-\psi_{0,i}(z)$ have entries in $I$; and $E$ contains $\rho_0(z)$, $\rho_1(z)$ and the scalars $\psi_{0,i}(z)$, $\psi_{1,i}(z)$.

\emph{(iii) Units.} Let $S$ be an $F$-subalgebra of $R$ and $u\in S$ with $e(u)\ne0$. Then $u-e(u)$ lies in $S$ and in the kernel of $e$, so it is nilpotent, and $u=e(u)+(u-e(u))$ is a unit of $S$; so $u$ is a unit of $R$ whose inverse lies in $S$.

\emph{(iv) Lagrange idempotents.} Put $a_i=\psi_{0,i}(z)$ and $a'_i=\psi_{1,i}(z)$. Then $e(a'_i)=e(a_i)$, as $e$ kills $I$, so $e(a_i-a_j)$ and $e(a'_i-a'_j)$ are nonzero for $i\ne j$, and $a_i-a_j$ and $a'_i-a'_j$ are units by (iii). Let $f_a=\prod_{b\ne a}(a_a-a_b)^{-1}(\rho_0(z)-a_b)$ and $f'_a=\prod_{b\ne a}(a'_a-a'_b)^{-1}(\rho_1(z)-a'_b)$. Steps (1) and (2) of the proof of Lemma~\ref{sha-rigidity-lagrange}, which use only the relation at elements of the group algebra and that the differences are units, show that the $f'_a$ are orthogonal idempotents with sum $1$ such that $f'_a\rho_1(y)f'_b=0$ for $b<a$ and $f'_a\rho_1(y)f'_a=\psi_{1,a}(y)f'_a$, for $y\in R[H]$; by (i) the same holds for the $f_a$, $\rho_0$ and the $\psi_{0,a}$. The $f_a$ and $f'_a$ lie in $E$: the products of the $\rho_0(z)-a_b$, respectively of the $\rho_1(z)-a'_b$, lie in $E$, and the scalars $\prod_{b\ne a}(a_a-a_b)^{-1}$ and $\prod_{b\ne a}(a'_a-a'_b)^{-1}$ lie in $E$ by (iii), applied to the $F$-subalgebra of the $c\in R$ whose scalar matrix lies in $E$.

\emph{(v) Moving the idempotents.} The reduction $\pi$ carries $f'_a$ to $f_a$, as it carries $\rho_1(z)$ and the $a'_b$ to the images of $\rho_0(z)$ and the $a_b$. So $Y=1-\sum_af'_af_a$, which lies in $E$, has $\pi(Y)=1-\sum_a\pi(f_a)=0$, that is, its entries lie in $I$, and $Y^2=0$. By orthogonality $(1-Y)f_a=f'_af_a=f'_a(1-Y)$, and $(1+Y)(1-Y)=1=(1-Y)(1+Y)$, so $(1+Y)f'_a(1-Y)=f_a$. Let $g\in H$, $r_0=\rho_0(g)$ and $r_1=\rho_1(g)$. As $\pi(r_1)=\pi(r_0)$ and $\pi(Y)=0$, $r_1Y=r_0Y$, $Yr_1=Yr_0$ and $Yr_1Y=0$; so $M_g=r_1-(r_0Y-Yr_0)=(1+Y)r_1(1-Y)$. Then $f_aM_gf_b=(1+Y)f'_ar_1f'_b(1-Y)$, which is $0$ for $b<a$ and $\psi_{1,a}(g)f_a$ for $b=a$.

\emph{(vi) Triangular form.} After conjugation by $w$, $\rho_0(z)$ is upper triangular with diagonal $a_1,\dots,a_n$, so $w^{-1}f_aw$, its Lagrange idempotent, is upper triangular with $(a,a)$ entry $1$. Let $V$ be the matrix whose $a$-th column is that of $w^{-1}f_aw$. It is upper triangular with diagonal entries $1$, so $V$ and $V^{-1}$ are upper triangular with diagonal entries $1$; and $w^{-1}f_awV=Ve_{aa}$, as the $f_a$ are orthogonal idempotents. So the $(a,b)$ entry of $V^{-1}w^{-1}M_gwV$ is that of $V^{-1}w^{-1}f_aM_gf_bwV$, and $V^{-1}w^{-1}M_gwV$ is upper triangular with diagonal $\psi_{1,1}(g),\dots,\psi_{1,n}(g)$; so is $w^{-1}M_gw$.

\emph{(vii) Conclusion.} The matrix of the statement is $\big(w^{-1}M_gw-w^{-1}\rho_0(g)w\big)\,w^{-1}\rho_0(g^{-1})w$. The first factor is upper triangular with diagonal $\psi_{1,i}(g)-\psi_{0,i}(g)$, by (vi) and (a), and the second is upper triangular with diagonal $\psi_{0,i}(g)^{-1}$, by (a) for $g^{-1}$. So their product is upper triangular with diagonal $(\psi_{1,i}(g)-\psi_{0,i}(g))\psi_{0,i}(g)^{-1}$.

(2) Apply (1) with $R=F[\varepsilon]$; $e$ the homomorphism $c_0+c_1\varepsilon\mapsto c_0$, whose kernel $\varepsilon F$ consists of elements of square $0$; $I=\varepsilon F$, with $I^2=0$; $\rho_0=\rho$, viewed over $F[\varepsilon]$, and $\rho_1=U$; $\psi_{0,i}=\psi_i$ and $\psi_{1,i}=\chi_i$; and $w$. Then (a), (d) and (e) are hypotheses, (b) holds as $U(g)-\rho(g)=\varepsilon U'(g)$ and $\chi_i(g)-\psi_i(g)=\varepsilon\chi'_i(g)$, and (c) holds as $e\circ\psi_i=\psi_i$. So (1) gives $Y_1\in E$ with entries in $\varepsilon F$; write $Y_1=\varepsilon Y$ with $Y\in M_n(F)$. As $U(g)\rho(g)^{-1}-1=\varepsilon U'(g)\rho(g)^{-1}$, the matrix of (1) is $\varepsilon\,w^{-1}\big(U'(g)\rho(g)^{-1}-\rho(g)Y\rho(g)^{-1}+Y\big)w$, upper triangular with diagonal $\varepsilon\chi'_i(g)\psi_i(g)^{-1}$. Comparing coefficients of $\varepsilon$ gives (2).
% (1) is locLemma: heI (e kills I), hsq (products of matrices over I vanish), tri_prod_zero for (i) (hord0), sep_good for (ii) (with the F-linear predicate G of (a), (b) and (e) at z), unit_sub for (iii), lag_triangular with ord_ext and lag_mem for (iv), lag_map and move_idem for (v), lag_tri, lag_map (conjugation by w) and tri_of_idem for (vi), hbeta and tri_diag for (vii). About 400 lines with its private lemmas.
% (2) is the private locDual of DefUSel.lean (toDual, evE, scal_mem; R = DualNumber F and I = ker e), which states only the vanishing of the diagonal where chi_i(g) = 1, the case used.
\end{proof}

\begin{lemma}\label{sha-patching-chain-cocycles-selmer}
In the setting of Definition~\ref{sha-patching-chain-selmer}, let $e\ge0$, let $R_{e,\mathcal K}$ be the $\mathcal K$-algebra defined as $R_g\otimes_A\mathcal K$ in Definition~\ref{sha-patching-selmer}(3) with $e$ in place of $g$, with basis $1$, $\epsilon_k$, $\epsilon'_k$, $\eta$, and write $y_0$ for the coefficient of $1$ in $y\in R_{e,\mathcal K}$. Let $u:\Gamma_{\infty,0}\to M_n(R_{e,\mathcal K})$ be a homomorphism of monoids, and $\psi'_{v,i}:D_v\to R_{e,\mathcal K}^\times$, for $v\mid\ell$ and $1\le i\le n$, homomorphisms, such that:
\begin{enumerate}
\item[(a)] the matrix of the coefficients of $1$ in $u(g)$ is $\rho(p(g))$;
\item[(b)] the entries of $u(g)$ in the two diagonal blocks lie in the span of $1$ and the $\epsilon'_k$, and those in the off-diagonal blocks in the span of the $\epsilon_k$ and $\eta$;
\item[(c)] ${}^tu(g)\,u(c_\infty gc_\infty^{-1})=\varepsilon(g)^{1-n}$ for $g\in\Gamma_{\infty,0}$;
\item[(d)] $u(g)=1$ for $g$ in the image of each $I'_u$;
\item[(e)] each $\psi'_{v,i}(g)$ lies in the span of $1$ and the $\epsilon'_k$, its coefficient of $1$ is $\psi_{v,i}(g)$, the character of $x$, and $\psi'_{v,i}=1$ on $\mu_v$;
\item[(f)] $\prod_{i=1}^n\big(u(g_i)-\psi'_{v,i}(g_i)\big)=0$ for $v\mid\ell$ and $g_1,\dots,g_n\in D_v$, with $D_v$ mapped to $\Gamma_{\infty,0}$;
\item[(g)] $u$ has bounded denominators.
\end{enumerate}
For $1\le k\le e$ let $u_k(g)$ and $u'_k(g)$ be the matrices of the coefficients of $\epsilon_k$ and of $\epsilon'_k$ in the entries of $u(g)$. Then $g\mapsto u'_k(g)\rho(p(g))^{-1}$ lies in $\mathrm{Sel}^{\mathrm{ev}}(\mathcal K)$, $g\mapsto u_k(g)\rho(p(g))^{-1}$ lies in $\mathrm{Sel}^{\mathrm{od}}(\mathcal K)$, and both have bounded denominators.
\end{lemma}
\leanref{Development/AutomorphicArithmetic/AutomorphicVarieties/PatchingModules/SelmerScaling/EpC.lean}{UpperTheory.V.Pa.DefU}
\leanref{Development/AutomorphicArithmetic/AutomorphicVarieties/PatchingModules/SelmerScaling/DefUSel.lean}{UpperTheory.V.Pa.defUSel}
\begin{proof}
Fix $k$, and let $d(y)$ be the coefficient of $\epsilon'_k$ in $y\in R_{e,\mathcal K}$ in the even case, and that of $\epsilon_k$ in the odd case. Products of two elements of the span of the $\epsilon_j$, $\epsilon'_j$ and $\eta$ lie in $\mathcal K\eta$, so $d(yy')=y_0d(y')+d(y)y'_0$; and $d$ vanishes on $\mathcal K$. Write $d(u(g))$ for the matrix of the $d$ of the entries, and put $z(g)=d(u(g))\rho(p(g))^{-1}$ and $\rho'=\rho\circ p$.

\emph{Values.} By (b), $d(u(g))$ vanishes in the off-diagonal blocks in the even case and in the diagonal blocks in the odd case. As $\rho'(g)^{-1}$ is block diagonal, $z(g)$ is even, respectively odd. The reductions of $z$ modulo $T^m\mathcal K=\mathcal K$ are constant.

\emph{Cocycle relations.} Applying $d$ to $u(gh)=u(g)u(h)$ and using (a) gives $d(u(gh))=\rho'(g)d(u(h))+d(u(g))\rho'(h)$, so $z(gh)=z(g)+\rho'(g)z(h)\rho'(g)^{-1}$. Applying $d$ to (c), whose right side lies in $\mathcal K$, gives ${}^t\rho'(g)\,d(u(g^c))+{}^td(u(g))\,\rho'(g^c)=0$ with $g^c=c_\infty gc_\infty^{-1}$; multiplying by ${}^t\rho'(g)^{-1}$ on the left and by $\rho'(g^c)^{-1}$ on the right gives $z(g^c)=-{}^tz(g)$. By (d), $z$ vanishes on the images of the $I'_u$, as $d(1)=0$.

\emph{The local condition.} Let $v\mid\ell$. The map $y\mapsto y_0+d(y)\varepsilon$ is a ring homomorphism $R_{e,\mathcal K}\to\mathcal K[\varepsilon]$, by the product rule for $d$. Apply Lemma~\ref{sha-patching-chain-cocycles-ordinary}(2) with $F=\mathcal K$, which is infinite; $H=D_v$; $\rho'$ restricted to $D_v$; $U$ and $\chi_i$ the images of $u$ on $D_v$ and of $\psi'_{v,i}$, so that $U'=d(u)$ and $\chi'_i=d(\psi'_{v,i})$ by (a) and (e); $\psi_i=\psi_{v,i}$, which are pairwise distinct by (N1); $w=w_v$; and $E$ the set of matrices over $\mathcal K[\varepsilon]$ whose coefficient of $1$ is even and whose coefficient of $\varepsilon$ is even, respectively odd, which is a subalgebra. The relation holds by (f). The $U(g)$ lie in $E$ by (a) and (b). The scalars $\chi_i(g)$ lie in $E$: in the odd case their coefficient of $\varepsilon$ is the coefficient of $\epsilon_k$ in $\psi'_{v,i}(g)$, which vanishes by (e). The lemma gives $Y$, even, respectively odd, such that $w_v^{-1}(z(g)-\rho'(g)Y\rho'(g)^{-1}+Y)w_v$ is upper triangular with diagonal $d(\psi'_{v,i}(g))\psi_{v,i}(g)^{-1}$, for $g\in D_v$. On $\mu_v$, $\psi'_{v,i}=1$ by (e), so this diagonal vanishes.

\emph{Bounded denominators.} The entries of $z(g)$ are sums of products of the entries of $d(u(g))$, which have bounded denominators by (g), and of $\rho'(g^{-1})$, which are maps to $A$ with locally constant reductions modulo every $T^m$, as $\rho$, $p$ and inversion are continuous. Sums and products of maps with bounded denominators have bounded denominators: if $T^s\phi$ and $T^{s'}\phi'$ are maps to $A$ with locally constant reductions, so are $T^{s+s'}(\phi+\phi')$ and $T^{s+s'}\phi\phi'$.
% defUSel applies the private selGen with d = ep' k (b = False, Ev) and d = ep k (b = True, Od): mapd_mul for the product rule, the values from hdU, the cocycle relation from hu.mul, the form from hu.form and hsc, ine, the local condition through the private locDual (toDual, evE, scal_mem) and locLemma (Lemma sha-patching-chain-cocycles-ordinary) with oddDiag and homBasic, and bndOps for the bounded denominators.
\end{proof}

\begin{lemma}\label{sha-patching-chain-cocycles-bounded}
In the setting of Definition~\ref{sha-patching-chain-tangent}, let $R$ be a commutative $\mathcal K$-algebra, finite-dimensional over $\mathcal K$, with a homomorphism of $\mathcal K$-algebras $e_R:R\to\mathcal K$ whose kernel is nilpotent, and let $\vartheta:\mathcal A\to R$ be a ring homomorphism with $e_R(\vartheta(y))=\hat f(y_0)$ for $y\in\mathcal A$ and $\vartheta(o)=\hat f(o)$ for $o$ in the image of $\mathcal O'$. Then for every $\mathcal K$-linear map $\lambda:R\to\mathcal K$ and all $i$ and $j$, the map $g\mapsto\lambda\big(\vartheta(r(g)_{ij})\big)$ has bounded denominators.
\end{lemma}
\leanref{Development/AutomorphicArithmetic/AutomorphicVarieties/PatchingModules/SelmerScaling/TrLatRange.lean}{UpperTheory.V.Pa.trLat}
\leanref{Development/AutomorphicArithmetic/AutomorphicVarieties/PatchingModules/SelmerScaling/TrBndDen.lean}{UpperTheory.V.Pa.trBnd}
\begin{proof}
View $R$ as an $A$-algebra through $A\subset\mathcal K$. Let $J=\ker e_R$, with $J^N=0$, and let $\varphi:\mathcal R\to R$ be the composite of $\mathcal R\to\hat{\mathcal R}\to\mathcal A$ with $\vartheta$. For $b\in\mathcal R$, $e_R(\varphi(b))=\hat f(b)=f(b)$, so $\varphi(b)-f(b)\in J$. Every $a\in A$ is $f(b)$ for some $b\in\mathcal R$ whose image in $\hat{\mathcal R}$ comes from $\mathcal O'$, and then $\varphi(b)=\hat f(b)=a$.

\emph{(i) A finite subalgebra.} Let $g_1,\dots,g_m$ generate $\ker f$, which is finitely generated as $\mathcal R$ is noetherian; then $\varphi(g_i)\in J$. By induction on $k$, every $\varphi(b)$ lies in $A[\varphi(g_1),\dots,\varphi(g_m)]+J^k$: choose $b'$ with $\varphi(b')=f(b')=f(b)$, write $b-b'=\sum_ic_ig_i$, and apply the induction hypothesis to the $c_i$. With $k=N$, $\varphi(\mathcal R)$ lies in $M_0=A[\varphi(g_1),\dots,\varphi(g_m)]$, which is finite over $A$, as the $\varphi(g_i)$ are nilpotent.

\emph{(ii) A lattice.} Let $P_r$ be the finitely generated $\mathcal R$-module of Lemma~\ref{sha-patching-chain-representation}(6), generated by $p_1,\dots,p_t$, and let $E'$ be the $A$-span of the entries of the $\vartheta(p_q)$. For $p\in P_r$ the entries of $\vartheta(p)$ lie in $M_0E'$. Let $P_0=(M_0+M_0E')\cap(A+J)$, a finitely generated $A$-module, each of whose elements is an element of $A$ plus a nilpotent one. So the $A$-algebra $M$ generated by $P_0$ is finite over $A$, and $M$ contains $\varphi(\mathcal R)$ and the $\vartheta(r(g)_{ij})$, which lie in $M_0E'$ and differ from $\rho(p(g))_{ij}$ by elements of $J$. Let $P'=ME'$, a finitely generated $A$-module with $MP'\subset P'$ that contains the entries of $\vartheta(p)$ for $p\in P_r$.

\emph{(iii) Powers of the maximal ideal.} As $f$ is surjective, $f(\mathfrak m_{\mathcal R})\subset TA$, so $\varphi(\mathfrak m_{\mathcal R})\subset TM+(J\cap M)$, and $\varphi(\mathfrak m_{\mathcal R}^{j+N})\subset T^jM$, as $J^N=0$.

\emph{(iv) Conclusion.} The image $\lambda(P')$ is a finitely generated $A$-submodule of $\mathcal K$, so there is $s$ with $T^s\lambda(P')\subset A$. Let $\phi(g)=T^s\lambda(\vartheta(r(g)_{ij}))\in A$. Given $N'$, Lemma~\ref{sha-patching-chain-representation}(6) gives an open normal subgroup $V$ with $r(gv)-r(g)\in\mathfrak m_{\mathcal R}^{N'+N}P_r$ for $v\in V$. By (iii) the entries of $\vartheta$ of such an element lie in $T^{N'}MP'\subset T^{N'}P'$, so $\phi(gv)-\phi(g)\in T^{N'}A$. So $\phi$ has locally constant reductions modulo every power of $T$, and $g\mapsto\lambda(\vartheta(r(g)_{ij}))$ has bounded denominators.
% trLat: the private trLat_range for (i) (h1 from oo_A and Test.oo, h2 from Test.res), E from ker_le_span_sub_sup_map for (ii), with adjoin_nil_finite for M and trLat_mul, and the private trLat_pow for (iii). trBnd: the private trBnd_den (exist_integer_multiples_of_finset) and cont for (iv). About 250 lines.
\end{proof}

\begin{lemma}\label{sha-patching-chain-cocycles-test}
In the setting of Definition~\ref{sha-patching-chain-tangent}, let $e\ge0$, $\alpha_1,\dots,\alpha_e\in\mathcal T$, $\mu_1,\dots,\mu_e\in\Lambda$, and $\Phi:\mathcal N\to\mathcal K$ an additive map with $\Phi(by)=\hat f(b)\Phi(y)+\sum_k\alpha_k(b)\mu_k(y)$ for $b\in\hat{\mathcal R}$ and $y\in\mathcal N$. Let $R_{e,\mathcal K}$ be as in Lemma~\ref{sha-patching-chain-cocycles-selmer}, and let $\vartheta:\mathcal A\to R_{e,\mathcal K}$ be the map
\[\vartheta(b\cdot1+y)=\hat f(b)+\sum_k\mu_k(y)\epsilon_k+\sum_k\alpha_k(b)\epsilon'_k+\Phi(y)\eta\qquad(b\in\hat{\mathcal R},\ y\in\mathcal N).\]
Then $\vartheta$ is a ring homomorphism. The map $u=\vartheta\circ r$ and the characters $\psi'_{v,i}=\vartheta\circ\psi_{v,i}$, with $\psi_{v,i}$ the characters of the ordinary pair mapped to $\mathcal A$, satisfy (a) to (g) of Lemma~\ref{sha-patching-chain-cocycles-selmer}, and in its notation $u_k(g)\rho(p(g))^{-1}=\mu_{k,C}(g)$ and $u'_k(g)\rho(p(g))^{-1}=\alpha_{k,C}(g)$.
\end{lemma}
\leanref{Development/AutomorphicArithmetic/AutomorphicVarieties/PatchingModules/SelmerScaling/TestEx.lean}{UpperTheory.V.Pa.testEx}
\leanref{Development/AutomorphicArithmetic/AutomorphicVarieties/PatchingModules/SelmerScaling/TestVal.lean}{UpperTheory.V.Pa.testVal}
\leanref{Development/AutomorphicArithmetic/AutomorphicVarieties/PatchingModules/SelmerScaling/DefUTest.lean}{UpperTheory.V.Pa.defUTest}
\begin{proof}
\emph{$\vartheta$ is a ring homomorphism.} It is additive. Let $a=b+y$ and $a'=b'+y'$ with $b,b'\in\hat{\mathcal R}$ and $y,y'\in\mathcal N$, and write $yy'=c\cdot1$; then $c\in I_{\mathcal N}$, and $aa'=(bb'+c)+(by'+b'y)$. By Lemma~\ref{sha-patching-chain-duality}(1), $\hat f(c)=0$, and $\alpha_k(c)=0$ as $\alpha_k$ vanishes on $I_{\mathcal N}$. So the coefficient of $1$ in $\vartheta(aa')$ is $\hat f(b)\hat f(b')$, that of $\epsilon_k$ is $\hat f(b)\mu_k(y')+\hat f(b')\mu_k(y)$, that of $\epsilon'_k$ is $\hat f(b)\alpha_k(b')+\hat f(b')\alpha_k(b)$, and that of $\eta$ is $\hat f(b)\Phi(y')+\hat f(b')\Phi(y)+\sum_k(\alpha_k(b)\mu_k(y')+\alpha_k(b')\mu_k(y))$; these are the coefficients of $\vartheta(a)\vartheta(a')$, as $\epsilon_k\epsilon'_k=\eta$. Finally $\vartheta(1)=1$, as $\alpha_k$ vanishes on the image of $\mathcal O'$, which contains $1$.

\emph{Elements of $\mathcal O$.} For $o$ in the image of $\mathcal O'$, $\vartheta(o)=\hat f(o)$. The images of $\mathcal O$ in $\hat{\mathcal R}$ come from $\mathcal O'$, so $\vartheta$ carries the image of $c\in\mathcal O$ to $\hat f(c)=f(c)$, the image of $c$ in $A$. In particular $\vartheta(\varepsilon(g)^{1-n})$ is the image of $\varepsilon(g)^{1-n}$ in $\kappa$, as $\mathbf Z_\ell$ has only one homomorphism to the finite field $\kappa$.

\emph{The conditions.} (a) The coefficient of $1$ in $\vartheta(r(g)_{ij})$ is $\hat f(r(g)_{ij,0})=\rho(p(g))_{ij}$ by Lemma~\ref{sha-patching-chain-representation}(3). (b) By Lemma~\ref{sha-patching-chain-representation}(1), the entries of $r(g)$ lie in $\hat{\mathcal R}\cdot1$ in the diagonal blocks and in $\mathcal N$ in the off-diagonal blocks, and $\vartheta$ carries these into the spans of $1,\epsilon'_k$ and of $\epsilon_k,\eta$. (c) follows from Lemma~\ref{sha-patching-chain-representation}(2) and the previous paragraph. (d) and (f): $r$ kills the relation ideal by Lemma~\ref{sha-patching-chain-representation}(4), which contains $g-1$ for $g$ in the images of the $I'_u$ and the products $\prod_i(g_i-\psi_{v,i}(g_i))$; so $r(g)=1$ on the images of the $I'_u$ and $\prod_i(r(g_i)-\psi_{v,i}(g_i))=0$, and applying $\vartheta$ gives (d) and (f). (e) $\psi'_{v,i}(g)=\hat f(\psi_{v,i}(g))+\sum_k\alpha_k(\psi_{v,i}(g))\epsilon'_k$, and $\hat f(\psi_{v,i}(g))=f(\psi_{v,i}(g))$ is the character of $x$ by (e) of Definition~\ref{sha-patching-chain-datum}(1). On $\mu_v$, $\psi_{v,i}=\theta_{v,i}$ takes values in the image of $\mathcal O$, so $\psi'_{v,i}(g)$ is the image of $\theta_{v,i}(g)$ in $A$, which is $f(\psi_{v,i}(g))$, the character of the normalized arc $x$, and is $1$. (g) $R_{e,\mathcal K}$ is finite-dimensional, the kernel of $y\mapsto y_0$ has zero cube, $\vartheta(y)_0=\hat f(y_0)$ and $\vartheta(o)=\hat f(o)$ on the image of $\mathcal O'$; so Lemma~\ref{sha-patching-chain-cocycles-bounded}, with $\lambda$ the coefficient maps, gives (g).

\emph{The coefficients.} The coefficients of $\epsilon_k$ and $\epsilon'_k$ in $\vartheta(r(g)_{ij})$ are $\mu_k(r(g)_{ij,1})$ and $\alpha_k(r(g)_{ij,0})$, which give $\mu_{k,C}$ and $\alpha_{k,C}$ after multiplying by $\rho(p(g))^{-1}$.
% testEx (prAdd: the product rule p5 with the element of ired, p6, the private dd_fin and dd_nil for the Test fields); testVal (oo_O and Test.oo for O, padicInt_ringHom_finite_field_eq for eps, the mu clause from the pair's psi.mu); defUTest assembles DefU: red from prVr, even from prVr, form from U.form, ine and ord from vrRel (U.kills of the pair and q generators), psi_mu from testVal, bnd from trBnd (Lemma sha-patching-chain-cocycles-bounded).
\end{proof}

\begin{lemma}\label{sha-patching-chain-cocycles}
In the setting of Definition~\ref{sha-patching-chain-tangent}, the maps $\alpha\mapsto\alpha_C$ and $\mu\mapsto\mu_C$ are $\mathcal K$-linear. For $\alpha\in\mathcal T$, $\alpha_C$ lies in $\mathrm{Sel}^{\mathrm{ev}}(\mathcal K)$ and has bounded denominators; for $\mu\in\Lambda$, $\mu_C$ lies in $\mathrm{Sel}^{\mathrm{od}}(\mathcal K)$ and has bounded denominators.
\end{lemma}
\leanref{Development/AutomorphicArithmetic/AutomorphicVarieties/PatchingModules/SelmerScaling/AlSel.lean}{UpperTheory.V.Pa.alSel}
\begin{proof}
The linearity is clear from the definitions, as $\rho(p(g))^{-1}$ does not depend on $\alpha$ or $\mu$. Let $\alpha\in\mathcal T$. Lemma~\ref{sha-patching-chain-cocycles-test}, with $e=1$, $\alpha_1=\alpha$, $\mu_1=0$ and $\Phi=0$, gives $u$ and $\psi'_{v,i}$ satisfying (a) to (g) of Lemma~\ref{sha-patching-chain-cocycles-selmer}, with $u'_1(g)\rho(p(g))^{-1}=\alpha_C(g)$; so $\alpha_C\in\mathrm{Sel}^{\mathrm{ev}}(\mathcal K)$ has bounded denominators by that lemma. In the same way, with $\alpha_1=0$, $\mu_1=\mu$ and $\Phi=0$, $\mu_C=u_1\rho(p(\cdot))^{-1}$ lies in $\mathrm{Sel}^{\mathrm{od}}(\mathcal K)$ and has bounded denominators.
% alSel: defUTest (Lemma sha-patching-chain-cocycles-test) with e = 1 and defUSel (Lemma sha-patching-chain-cocycles-selmer).
\end{proof}

\begin{lemma}\label{sha-patching-chain-linear}
Let $F$ be a field, $\mathcal T$ a finite-dimensional $F$-vector space, $g,D_0\ge0$ integers and $\mathcal V\subset\mathcal T^g$ a subspace. Suppose that:
\begin{enumerate}
\item for all $z_0,\dots,z_g\in\mathcal T^g$ there are $c_0,\dots,c_g\in F$, not all zero, with $\sum_ic_iz_i\in\mathcal V$;
\item if $\alpha_1,\dots,\alpha_e\in\mathcal T$ are linearly independent and, for all $i$ and $j$, the family with $\alpha_i$ in position $j$ and $0$ elsewhere lies in $\mathcal V$, then $e+g^2\le D_0$.
\end{enumerate}
Then there are linear forms $X_1,\dots,X_{D_0}$ spanning the dual of $\mathcal T$, and $c_{i,i',k}\in F$, for $1\le i,i'\le g$ and $1\le k\le D_0$, such that the $g^2$ vectors $(c_{i,i',k})_k\in F^{D_0}$ are linearly independent and $\sum_{i'}\sum_kc_{i,i',k}X_k(\alpha_{i'})=0$ for every $(\alpha_1,\dots,\alpha_g)\in\mathcal V$ and every $i$.
\end{lemma}
\leanref{Development/AutomorphicArithmetic/AutomorphicVarieties/PatchingModules/SelmerScaling/LinAlg.lean}{UpperTheory.V.Pa.linAlg}
\begin{proof}
By (1), any $g+1$ elements of $\mathcal T^g/\mathcal V$ are linearly dependent, so $\dim\mathcal T^g/\mathcal V\le g$, and the annihilator $\mathcal V^\perp$ of $\mathcal V$ in the dual of $\mathcal T^g$ has dimension at most $g$. Let $\Phi_1,\dots,\Phi_g$ span $\mathcal V^\perp$, and let $\varphi_{i,i'}$ be the linear form on $\mathcal T$ that sends $t$ to $\Phi_i$ of the family with $t$ in position $i'$ and $0$ elsewhere. Then $\sum_{i'}\varphi_{i,i'}(\alpha_{i'})=\Phi_i(\alpha)=0$ for $\alpha\in\mathcal V$.

Let $\mathcal W$ be the subspace of $t\in\mathcal T$ such that every family with $t$ in one position and $0$ elsewhere lies in $\mathcal V$. By (2), applied to a basis of $\mathcal W$, $\dim\mathcal W+g^2\le D_0$. Let $S$ be the span of the $\varphi_{i,i'}$. If $t\in\mathcal T$ is killed by every $\varphi_{i,i'}$, then each family with $t$ in position $i'$ and $0$ elsewhere is killed by $\mathcal V^\perp$, so it lies in $\mathcal V$; hence $t\in\mathcal W$. So $\dim\mathcal T-\dim S\le\dim\mathcal W$, and a complement $C$ of $S$ in the dual of $\mathcal T$ has dimension at most $\dim\mathcal W\le D_0-g^2$.

Let $X_1,\dots,X_{g^2}$ be the $\varphi_{i,i'}$, with $\varphi_{i,i'}=X_{k(i,i')}$ for a bijection $k$ from the pairs to $\{1,\dots,g^2\}$, and let $X_{g^2+1},\dots,X_{D_0}$ span $C$, some of them $0$ if needed. They span the dual of $\mathcal T$. Let $c_{i,i',k}$ be $1$ for $k=k(i,i')$ and $0$ otherwise; these $g^2$ vectors are distinct standard basis vectors, and $\sum_{i'}\sum_kc_{i,i',k}X_k(\alpha_{i'})=\sum_{i'}\varphi_{i,i'}(\alpha_{i'})=0$ for $\alpha\in\mathcal V$.
% linAlg: finrank_quotient_le_of_forall_dependent_mod, Subspace.finrank_add_finrank_dualAnnihilator_eq, exists_spanning_family_of_finrank_le, hrow, W as an iInf of comaps, hSco and hSrank, Submodule.exists_isCompl, the equivalence e of Fin D0 with Fin (g*g) + Fin (D0 - g*g).
\end{proof}

\begin{lemma}\label{sha-patching-chain-rank-lift-triangular}
In the setting of Definition~\ref{sha-patching-chain-selmer} with $X=\mathcal K$, let $v\mid\ell$, write $\psi_{v,1},\dots,\psi_{v,n}$ for the characters of $x$ at $v$, and put $E_v=\sum_k\epsilon_kY_{k,v}$. Then:
\begin{enumerate}
\item $w_v^{-1}Zw_v$ has zero diagonal for every odd matrix $Z$ over $\mathcal K$;
\item $w'_v$ is invertible, with inverse $w_v^{-1}(1+E_v)$;
\item for $g\in D_v$, ${w'_v}^{-1}\rho_1(p(g))w'_v$ is upper triangular with diagonal $\psi_{v,1}(g),\dots,\psi_{v,n}(g)$;
\item $g\mapsto\rho_1(p(g))$ is a homomorphism from $\Gamma_{\infty,0}$ with values in the first module $V$ of Definition~\ref{sha-patching-chain-selmer}(3), trivial on the images of the $I'_u$, and ${}^t\rho_1(p(g))\,\rho_1(p(c_\infty gc_\infty^{-1}))=\varepsilon(g)^{1-n}$ for $g\in\Gamma_{\infty,0}$.
\end{enumerate}
\end{lemma}
\leanref{Development/AutomorphicArithmetic/AutomorphicVarieties/PatchingModules/SelmerScaling/InverseWeylLiftEqPair.lean}{UpperTheory.V.Pa.rho1Tri}
\leanref{Development/AutomorphicArithmetic/AutomorphicVarieties/PatchingModules/SelmerScaling/OddDiag.lean}{UpperTheory.V.Pa.oddDiag}
\begin{proof}
Write $w=w_v$ and $E=E_v$. (1) Each column of $w$ lies in one of the two blocks, so $w$ is block diagonal after a permutation of its columns, and each row of $w^{-1}$ is supported on the block in which the column of $w$ with the same index lies. For an odd $Z$, $Z$ maps the $i$-th column of $w$ into the other block, so the $(i,i)$ entry of $w^{-1}Zw$ vanishes.

(2) As $Y_{k,v}$ is odd, $E$ has entries in $\sum_k\mathcal K\epsilon_k$, and $E^2=0$. So $(1-E)w\cdot w^{-1}(1+E)=1-E^2=1$, and $w^{-1}(1+E)(1-E)w=1$.

(3) Let $g\in D_v$, and write $r=\rho(p(g))$ and $\kappa'_k=\kappa_k(p(g))$; let $\delta Y_{k,v}=rY_{k,v}r^{-1}-Y_{k,v}$. As products of two of the $\epsilon_k$ vanish, ${w'_v}^{-1}\rho_1(p(g))w'_v=w^{-1}(1+E)(1+\sum_k\epsilon_k\kappa'_k)r(1-E)w$ is
\[w^{-1}rw+\sum_k\epsilon_k\,w^{-1}\big(\kappa'_kr+Y_{k,v}r-rY_{k,v}\big)w=w^{-1}rw+\sum_k\epsilon_k\,\big(w^{-1}(\kappa'_k-\delta Y_{k,v})w\big)\big(w^{-1}rw\big).\]
Here $w^{-1}rw$ is upper triangular with diagonal $\psi_{v,1}(g),\dots,\psi_{v,n}(g)$, as $x$ is normalized, and $w^{-1}(\kappa'_k-\delta Y_{k,v})w$ is upper triangular by the local condition of $\kappa_k$ at $v$ in Definition~\ref{sha-patching-selmer}(1). So the matrix is upper triangular. The matrix $\kappa'_kr+Y_{k,v}r-rY_{k,v}$ is odd, as $\kappa'_k$ and $Y_{k,v}$ are odd and $r$ is even; so by (1) the coefficients of $\epsilon_k$ on the diagonal vanish, and the diagonal is that of $w^{-1}rw$.

(4) As $p$ is compatible with $c_\infty$, $c$, the local maps and the $I'_u$, it suffices to prove the statements for $\rho_1$ on $\Gamma_0$. For $h\in\Gamma_0$, $\rho_1(h)=\rho(h)+\sum_j\epsilon_j\kappa_j(h)\rho(h)$, with $\rho(h)$ even and $\kappa_j(h)\rho(h)$ odd; so $\rho_1(h)$ lies in the first module. As products of two of the $\epsilon_j$ vanish, $\rho_1(h)\rho_1(h')=\rho(hh')+\sum_j\epsilon_j\big(\kappa_j(h)+\rho(h)\kappa_j(h')\rho(h)^{-1}\big)\rho(hh')=\rho_1(hh')$, by the cocycle relation of $\kappa_j$. On $I'_u$, $\rho=1$ and $\kappa_j=0$. Finally, with $h^c=chc^{-1}$, ${}^t\rho_1(h)\rho_1(h^c)={}^t\rho(h)\rho(h^c)+\sum_j\epsilon_j\,{}^t\rho(h)\big({}^t\kappa_j(h)+\kappa_j(h^c)\big)\rho(h^c)=\varepsilon(h)^{1-n}$, as $x$ is normalized and $\kappa_j(h^c)=-{}^t\kappa_j(h)$.
% (4) is the first five conjuncts of rho1Tri (rho1_eq with congruenceCondition_mul and k.mem, pair_matrix_transpose and x.form, x.ine, od_ev), composed with p by homBasic.
% rho1Tri, last conjunct: weyl_lift_eq_pair, inverse_weyl_lift_eq_pair and rho1_eq (the private PP with congruenceCondition_mul), key, tri_mul with k.loc and oddDiag (h4), and h3 of oddDiag (conjugate_diagonal_zero_of_parity with x.w_blk) with od_ev and ev_od for the diagonal; the second and third conjuncts of rho1Tri's last part are (2), from weyl_lift_eq_pair and inverse_weyl_lift_eq_pair.
\end{proof}

\begin{lemma}\label{sha-patching-chain-rank-lift-invariants}
In the setting of Definition~\ref{sha-patching-chain-selmer} with $X=\mathcal K$, let $v\mid\ell$ and $Q\in M_n(\mathcal K)$, and suppose that $w_v^{-1}\big(\rho(p(g))Q\rho(p(g))^{-1}-Q\big)w_v$ is upper triangular for every $g\in D_v$. Then $w_v^{-1}Qw_v$ is upper triangular.
\end{lemma}
\leanref{Development/AutomorphicArithmetic/AutomorphicVarieties/PatchingModules/SelmerScaling/LowerEntriesZeroOfCommutingTriangular.lean}{UpperTheory.V.Pa.localInvariants_eq_zero}
\begin{proof}
For $g\in D_v$ let $T_g=w_v^{-1}\rho(p(g))w_v$, upper triangular with diagonal $\psi_{v,1}(g),\dots,\psi_{v,n}(g)$, the characters of $x$ at $v$, and let $M=w_v^{-1}Qw_v$. Then $w_v^{-1}\big(\rho(p(g))Q\rho(p(g))^{-1}-Q\big)w_v\,T_g=T_gM-MT_g$, a product of two upper triangular matrices; so $T_gM-MT_g$ is upper triangular. We show $M_{ij}=0$ for $i>j$, by induction on $(n-i)+j$. Since $T_g$ is upper triangular and, by induction, $M_{kj}=0$ for $k>i$ and $M_{ik}=0$ for $k<j$, the $(i,j)$ entries of $T_gM$ and $MT_g$ are $\psi_{v,i}(g)M_{ij}$ and $M_{ij}\psi_{v,j}(g)$. So $(\psi_{v,i}(g)-\psi_{v,j}(g))M_{ij}=0$ for every $g\in D_v$. By (N1), $\psi_{v,i}\ne\psi_{v,j}$, so $M_{ij}=0$.
% localInvariants_eq_zero: oddDiag (h1, h4, h5, h6), the private conjugate_commutator_identity and upper_triangular_mul, and lower_entries_zero_of_commuting_triangular (the induction on (n - i) + j, with x.N1 v for the distinct diagonals).
\end{proof}

\begin{lemma}\label{sha-patching-chain-rank-lift}
In the setting of Definition~\ref{sha-patching-chain-selmer} with $X=\mathcal K$, let $z\in\mathrm{Sel}^{\mathrm{ev}}(\mathcal K)$ have bounded denominators, and suppose that for each $j$ the family with $z$ in position $j$ and $0$ elsewhere lifts. Then there is $y\in\mathrm{Sel}^1(\mathcal K)$ with bounded denominators whose matrix of coefficients of $1$ is $z(g)$ at every $g$.
\end{lemma}
\leanref{Development/AutomorphicArithmetic/AutomorphicVarieties/PatchingModules/SelmerScaling/FirstOrderMatrixMemSquareZero.lean}{UpperTheory.V.Pa.lift1}
\begin{proof}
For each $j$ let $\Psi_j\in\mathrm{Sel}^2(\mathcal K)$ be a lift of the family with $z$ in position $j$, and let $y(g)=z(g)+\sum_j\epsilon_j\,\eta_j(g)$, where $\eta_j(g)$ is the matrix of the coefficients of $\eta$ in $\Psi_j(g)$. Then $\epsilon'_jy(g)=\Psi_j(g)$, as the coefficients of $\epsilon'_i$ in $\Psi_j$ are $\delta_{ij}z$ and $\Psi_j$ has no other coefficients besides those of $\eta$. A matrix with entries in the span of $1$ and the $\epsilon_m$ is determined by its coefficients of $1$ and its products with the $\epsilon'_j$. Applied to both sides of the cocycle relation, the relation with $c_\infty$ and the vanishing on the $I'_u$, this reduces them for $y$ to those for $z$ and the $\Psi_j$, as $\epsilon'_j$ is central. The values of $y$ lie in the first module $V$, as $z$ is even and the $\eta_j$ are odd; and $y$ has bounded denominators, as $z$ and the $\Psi_j$ have.

\emph{The local condition.} Let $v\mid\ell$, and write $\rho'=\rho\circ p$, $\rho'_1=\rho_1\circ p$, $w=w_v$, $w'=w'_v$ and $E=\sum_k\epsilon_kY_{k,v}$, so that $w'=(1-E)w$. Call a matrix over $R_{g,\mathcal K}$ \emph{first-order} if its entries lie in the span of $1$ and the $\epsilon_k$; these matrices form a ring, and $\rho'_1(g)$, $\rho'_1(g)^{-1}=\rho'_1(g^{-1})$, $w'$ and ${w'}^{-1}$ are first-order, with matrices of coefficients of $1$ equal to $\rho'(g)$, $\rho'(g)^{-1}$, $w$ and $w^{-1}$ by Lemma~\ref{sha-patching-chain-rank-lift-triangular}(2). For first-order $P$ and $P'$ and any $M$, the matrix of coefficients of $\epsilon'_i$ in $PMP'$ is $P_0M'_iP'_0$, where $P_0$ and $P'_0$ are the matrices of coefficients of $1$ and $M'_i$ that of the coefficients of $\epsilon'_i$ in $M$. Let $Y$, even, be a matrix for the local condition of $z$ at $v$, and $Y_j$, in the second module $V$, one for that of $\Psi_j$. Let $Y'_{j,i}$ be the matrix of coefficients of $\epsilon'_i$ in $Y_j$, and $Q_{j,i}=Y'_{j,i}-\delta_{ij}Y$, an even matrix.
\begin{enumerate}
\item[(i)] \emph{$w^{-1}Q_{j,i}w$ is upper triangular.} Let $g\in D_v$. The coefficients of $\epsilon'_i$ in the upper triangular matrix ${w'}^{-1}(\Psi_j(g)-\rho'_1(g)Y_j\rho'_1(g)^{-1}+Y_j)w'$ form the matrix $w^{-1}(\delta_{ij}z(g)-\rho'(g)Y'_{j,i}\rho'(g)^{-1}+Y'_{j,i})w$, which is therefore upper triangular. Subtracting it from $\delta_{ij}$ times the upper triangular matrix $w^{-1}(z(g)-\rho'(g)Y\rho'(g)^{-1}+Y)w$ shows that $w^{-1}(\rho'(g)Q_{j,i}\rho'(g)^{-1}-Q_{j,i})w$ is upper triangular. By Lemma~\ref{sha-patching-chain-rank-lift-invariants}, so is $w^{-1}Q_{j,i}w$.
\item[(ii)] \emph{A correction.} Let $C_j=(1-E)\big(\sum_i\epsilon'_iQ_{j,i}\big)(1+E)$. Its matrix of coefficients of $\epsilon'_i$ is $Q_{j,i}$. It lies in the second module $V$: $1\pm E$ lies in the first module, $\sum_i\epsilon'_iQ_{j,i}$ in the second, as the $Q_{j,i}$ are even, and the products of an element of the first module and one of the second, in either order, lie in the second. As $E^2=0$, ${w'}^{-1}C_jw'=w^{-1}(1+E)(1-E)\big(\sum_i\epsilon'_iQ_{j,i}\big)(1+E)(1-E)w=\sum_i\epsilon'_i\,w^{-1}Q_{j,i}w$, which is upper triangular by (i).
\item[(iii)] \emph{The matrix for $y$.} The matrix $Y_j-C_j$ lies in the second module, and its coefficients of $\epsilon'_i$ are $Y'_{j,i}-Q_{j,i}=\delta_{ij}Y$. Let $H_j$ be its matrix of coefficients of $\eta$, which is odd, and put $Y^\sharp=Y+\sum_j\epsilon_jH_j$. Then $Y^\sharp$ lies in the first module $V$, and $\epsilon'_jY^\sharp=\epsilon'_jY+\eta H_j=Y_j-C_j$.
\item[(iv)] \emph{The condition.} Let $g\in D_v$, and $M={w'}^{-1}(y(g)-\rho'_1(g)Y^\sharp\rho'_1(g)^{-1}+Y^\sharp)w'$, a first-order matrix. An entry $c+\sum_kc_k\epsilon_k$ of a first-order matrix vanishes if $c=0$ and $\epsilon'_j(c+\sum_kc_k\epsilon_k)=c\epsilon'_j+c_j\eta$ vanishes for all $j$. The matrix of coefficients of $1$ of $M$ is $w^{-1}(z(g)-\rho'(g)Y\rho'(g)^{-1}+Y)w$, which is upper triangular, with zero diagonal if $g\in\mu_v$. As $\epsilon'_j$ is central, $\epsilon'_jy(g)=\Psi_j(g)$ and $\epsilon'_jY^\sharp=Y_j-C_j$,
\[\epsilon'_jM={w'}^{-1}\big(\Psi_j(g)-\rho'_1(g)Y_j\rho'_1(g)^{-1}+Y_j\big)w'+\big(BC'_jB'-C'_j\big),\]
with $B={w'}^{-1}\rho'_1(g)w'$, $B'={w'}^{-1}\rho'_1(g^{-1})w'$ and $C'_j={w'}^{-1}C_jw'$. The first term is upper triangular, with zero diagonal if $g\in\mu_v$. By Lemma~\ref{sha-patching-chain-rank-lift-triangular}(3), $B$ and $B'$ are upper triangular with diagonals $\psi_{v,i}(g)$ and $\psi_{v,i}(g)^{-1}$, the characters of $x$; with (ii), $BC'_jB'-C'_j$ is upper triangular with diagonal entries $\psi_{v,i}(g)(C'_j)_{ii}\psi_{v,i}(g)^{-1}-(C'_j)_{ii}=0$. So $\epsilon'_jM$ is upper triangular for every $j$, with zero diagonal if $g\in\mu_v$, and hence so is $M$.
\end{enumerate}
So $Y^\sharp$ satisfies the local condition for $y$ at $v$.
% lift1, about 725 lines: mk1 for y, toE for e_j . y = Psi_j, eq_of and first_order_matrix_mem_square_zero for the relations; the local condition: SM (the first-order matrices) with pM_conj, Q (the difference of the pM i (Yj j) and Pi.single j Y) with hQe, hQt from localInvariants_eq_zero (Lemma sha-patching-chain-rank-lift-invariants), corr with weyl_correction_mem_second_order (blk_mul, twoSubspaceIntersection, second_order_mul_first_order_mem), pM_corr and corr_conj for (ii), Y' = mk1 Y (eta of Yj - corr) with toE for (iii), zero_of, hA, hB and hC (tri_conj with rho1Tri, Lemma sha-patching-chain-rank-lift-triangular) for (iv).
\end{proof}

\begin{lemma}\label{sha-patching-chain-rank-basis}
In the setting of Definition~\ref{sha-patching-chain-selmer} with $X=\mathcal K$, let $1\le i,j\le g$. Then the map $g'\mapsto\epsilon_j\,\kappa_i(p(g'))$ lies in $\mathrm{Sel}^1(\mathcal K)$ and has bounded denominators.
\end{lemma}
\leanref{Development/AutomorphicArithmetic/AutomorphicVarieties/PatchingModules/SelmerScaling/EOneTranspose.lean}{UpperTheory.V.Pa.kapSel1}
\begin{proof}
Write $\rho'=\rho\circ p$ and $\kappa'=\kappa_i\circ p$. Products of two matrices with entries in $\sum_m\mathcal K\epsilon_m$ vanish, so for matrices $P,P'$ over $\mathcal K$ and $u,u'$ with entries in that span, $(1+u)P\cdot\epsilon_jZ\cdot(1+u')P'=\epsilon_jPZP'$ and $P(1+u)\cdot\epsilon_jZ\cdot(1-u')P'=\epsilon_jPZP'$. With the first identity, $\rho_1(p(g))\,\epsilon_j\kappa'(h)\,\rho_1(p(g))^{-1}=\epsilon_j\rho'(g)\kappa'(h)\rho'(g)^{-1}$, so the cocycle relation of $\epsilon_j\kappa'$ is $\epsilon_j$ times that of $\kappa'$; the relation with $c_\infty$ and the vanishing on the images of the $I'_u$ follow from those of $\kappa_i$, as $p$ is compatible with $c_\infty$ and the local maps. The values of $\kappa_i$ are odd, so those of $\epsilon_j\kappa'$ lie in the first module $V$ of Definition~\ref{sha-patching-chain-selmer}(3). At $v\mid\ell$, let $Y_v$ be the odd matrix of the local condition for $\kappa_i$, and take $\epsilon_jY_v$. By the second identity, ${w'_v}^{-1}\big(\epsilon_j\kappa'(g)-\rho_1(p(g))\epsilon_jY_v\rho_1(p(g))^{-1}+\epsilon_jY_v\big)w'_v=\epsilon_jw_v^{-1}\big(\kappa_i(p(g))-\rho(p(g))Y_v\rho(p(g))^{-1}+Y_v\big)w_v$ for $g\in D_v$, which is upper triangular, with zero diagonal on $\mu_v$. Finally, the only nonzero coefficients of the entries of $\epsilon_j\kappa'$ are the entries of $\kappa'$, which are maps to $A$ with locally constant reductions modulo every power of $T$, as $p$ is continuous; so $\epsilon_j\kappa'$ has bounded denominators, and its reductions modulo $T^m R_{g,\mathcal K}=R_{g,\mathcal K}$ are constant.
% kapSel1: epM_single_map_mem_W1_of_mem_Od for the values, adicCont_epM_single_map_apply, the private one_add_epM_mul_map_C_mul_epM_mul_one_add_epM_mul_map_C and map_C_mul_one_add_epM_mul_epM_mul_one_sub_epM_mul_map_C for the two identities, epM_neg_transpose for c, and bndD_epM_single_map_of_adicCont.
\end{proof}

\begin{lemma}\label{sha-patching-chain-rank-invariant}
In the setting of Definition~\ref{sha-patching-chain-selmer}, let $Y$ be an even matrix over $\mathcal K$ with ${}^tY=-Y$ and $\rho(p(g))Y\rho(p(g))^{-1}=Y$ for all $g\in\Gamma_{\infty,0}$. Then $Y=0$.
\end{lemma}
\leanref{Development/AutomorphicArithmetic/AutomorphicVarieties/PatchingModules/SelmerScaling/InvEv.lean}{UpperTheory.V.Pa.invEv}
\begin{proof}
As $p$ maps $\Gamma_{\infty,0}$ onto $\Gamma_0$, $Y$ commutes with $\rho(g)$ for $g\in\Gamma_0$, hence with every matrix in the $A$-span of $\rho(\Gamma')$. Let $i\ne j$. The matrix unit $e_{ii}$ is even, so, as $x$ is normalized, $T^se_{ii}$ lies in that span for some $s$, and $Y$ commutes with $e_{ii}$. The entry $(i,j)$ of $e_{ii}Y=Ye_{ii}$ gives $Y_{ij}=0$. Finally ${}^tY=-Y$ gives $2Y_{ii}=0$, so $Y_{ii}=0$, as $\ell$ is odd.
% invEv: hc (Y commutes with rho through the surjectivity of h.p0), hs (span induction), x.span for single i i 1, and the unit 2 in Z_l.
\end{proof}

\begin{lemma}\label{sha-patching-chain-rank-injective-generators}
In the setting of Definition~\ref{sha-patching-chain-tangent}, let $R$, $e_R$ and $\vartheta:\mathcal A\to R$ be as in Lemma~\ref{sha-patching-chain-cocycles-bounded}, view $A\subset\mathcal K$ inside $R$, and map $\mathcal R$ to $\mathcal A$ through $\hat{\mathcal R}$. Suppose that $\vartheta(b)=f(b)$ for every $b\in\mathcal R$ that is a coefficient of the characteristic polynomial of $D$ at an element of $\Gamma_{\infty,0}$, a value of $\psi$, one of the $U_{j,i}$ or $y_{j,i}$, or $s$. Then $\vartheta(b\cdot1)=\hat f(b)$ for every $b\in\hat{\mathcal R}$.
\end{lemma}
\leanref{Development/AutomorphicArithmetic/AutomorphicVarieties/PatchingModules/SelmerScaling/TrConst.lean}{UpperTheory.V.Pa.trConst}
\begin{proof}
Let $\varphi:\mathcal R\to R$ be the composite of $\mathcal R\to\hat{\mathcal R}\to\mathcal A$ with $\vartheta$. Both $\varphi$ and $f$, viewed as a map to $R$, are ring homomorphisms.

\emph{(i) $\varphi=f$ on $\mathcal R$.} For $c\in\mathcal O$, the image of $c$ in $\hat{\mathcal R}$ is that of some $o\in\mathcal O'$, so $\varphi(c)=\vartheta(o)=\hat f(o)=f(c)$. So $\varphi$ and $f$ agree on the $\mathcal O$-subalgebra $\mathcal R_0$ of $\mathcal R$ generated by the elements of the hypothesis. Steps (i) to (iii) of the proof of Lemma~\ref{sha-patching-chain-cocycles-bounded} give an $A$-subalgebra $M$ of $R$, finite as an $A$-module, containing $\varphi(\mathcal R)$, and for every $j$ an integer $j'$ with $\varphi(\mathfrak m_{\mathcal R}^{j'})\subset T^jM$. As $\mathcal O$ maps onto $\kappa$, the composite of $f$ with $A\to\kappa$ is surjective, so its kernel is $\mathfrak m_{\mathcal R}$, and $f(\mathfrak m_{\mathcal R}^j)\subset T^jA$. Let $y\in\mathcal R$ and $j\ge0$. By (d) of Definition~\ref{sha-patching-chain-datum}(1), there is $a\in\mathcal R_0$ with $y-a\in\mathfrak m_{\mathcal R}^{\max(j,j')}$. Then $\varphi(y)-f(y)=\varphi(y-a)-f(y-a)$ lies in $T^jM$. As $M$ is a finite module over the local noetherian ring $A$, $\bigcap_jT^jM=0$ by Krull's intersection theorem \cite[Theorem~8.10]{Mat86}; so $\varphi(y)=f(y)$.

\emph{(ii) $\vartheta$ kills $\mathfrak m_{\hat{\mathcal R}}\cdot1$.} Let $N$ be such that $(\ker e_R)^N=0$, and $m\in\mathfrak m_{\hat{\mathcal R}}$. Then $e_R(\vartheta(m))=\hat f(m)=0$, so $\vartheta$ kills $\mathfrak m_{\hat{\mathcal R}}^N$. As $\hat{\mathcal R}/\mathfrak m_{\hat{\mathcal R}}^N$ is the quotient of the localization of $\mathcal R$ at $\ker f$ by the $N$-th power of its maximal ideal, there are $a,t\in\mathcal R$ with $t\notin\ker f$ and $tm-a\in\mathfrak m_{\hat{\mathcal R}}^N$. Applying $\vartheta$ and (i) gives $f(t)\vartheta(m)=f(a)$, with $f(t)\ne0$. So $\vartheta(m)=f(t)^{-1}f(a)$ lies in $\mathcal K$, and it equals $e_R(\vartheta(m))=0$.

\emph{(iii) Conclusion.} Let $b\in\hat{\mathcal R}$. As $\hat f$ maps $\mathcal O'$ onto $\mathcal K$, there is $o$ in the image of $\mathcal O'$ with $\hat f(o)=\hat f(b)$. Then $b-o\in\ker\hat f=\mathfrak m_{\hat{\mathcal R}}$, and $\vartheta(b)=\vartheta(o)=\hat f(o)=\hat f(b)$ by (ii).
% trConst: trLat for M (Lemma sha-patching-chain-cocycles-bounded, steps (i) to (iii)), hO from oo_O and Test.oo, hadj (adjoin induction), hB with S.P.gen (GenBy), the private f_local and krull_sub (Ideal.iInf_pow_smul_eq_bot_of_isLocalRing); hker from Test.res and prAdd, hkill from Test.nil, hmax from IsCplAt.dense; S.oo_surj with Test.oo for (iii).
\end{proof}

\begin{lemma}\label{sha-patching-chain-rank-injective-constant}
In the setting of Definition~\ref{sha-patching-chain-tangent}, let $R$, $e_R$ and $\vartheta:\mathcal A\to R$ be as in Lemma~\ref{sha-patching-chain-cocycles-bounded}, view $\mathcal K$ inside $R$, and write $\vartheta(Z)$ for the matrix of the $\vartheta(Z_{ij})$, for $Z\in M_n(\mathcal A)$. Suppose that there is $C\in GL_n(R)$ with $C\vartheta(r(g))C^{-1}=\rho(p(g))$ for all $g\in\Gamma_{\infty,0}$. Then $\vartheta(b\cdot1)=\hat f(b)$ for every $b\in\hat{\mathcal R}$.
\end{lemma}
\leanref{Development/AutomorphicArithmetic/AutomorphicVarieties/PatchingModules/SelmerScaling/ConjConst.lean}{conj_const}
\begin{proof}
Map $\mathcal R$ to $\mathcal A$ through $\hat{\mathcal R}$. By Lemma~\ref{sha-patching-chain-rank-injective-generators}, it suffices that $\vartheta(b)=f(b)$ for the elements $b$ of its hypothesis. For $b\in\mathcal R$, $e_R(\vartheta(b))=\hat f(b)=f(b)$; so it suffices that $\vartheta(b)$ lies in $\mathcal K$. An element $u$ of $R$ with $e_R(u)\ne0$ is a unit, since $u-e_R(u)$ is nilpotent. Let $\Phi:M_n(\mathcal A)\to M_n(R)$ be the ring homomorphism $Z\mapsto C\vartheta(Z)C^{-1}$. Then $\Phi(r(g))=\rho(p(g))$, and $\Phi$ maps the scalar matrix $c$ to the scalar matrix $\vartheta(c)$. By Lemma~\ref{sha-patching-chain-representation}(4), $r$ kills the relation ideal, and so the generators listed in Definition~\ref{sha-patching-chain-datum}(2).

\emph{(a) Coefficients of characteristic polynomials.} By Lemma~\ref{sha-patching-chain-representation}(2), the characteristic polynomial of $D$ at $g\in\Gamma_{\infty,0}$ maps to that of $r(g)$, and $\vartheta$ maps this to the characteristic polynomial of $\vartheta(r(g))$, which is that of its conjugate $\rho(p(g))$. Its coefficients lie in $\mathcal K$.

\emph{(b) Values of $\psi$.} Let $v\mid\ell$, and write $\psi^x_{v,i}$ for the characters of $x$, so that $f(\psi_{v,i}(g))=\psi^x_{v,i}(g)$ by (e) of Definition~\ref{sha-patching-chain-datum}(1). For $g\in D_v$ let $M(g)=w_v^{-1}\rho(p(g))w_v$, which is upper triangular over $\mathcal K$ with diagonal $\psi^x_{v,1}(g),\dots,\psi^x_{v,n}(g)$, and $c_i(g)=\vartheta(\psi_{v,i}(g))$, so that $e_R(c_i(g))=\psi^x_{v,i}(g)$. Applying $\Phi$, and then conjugation by $w_v$, to $\prod_i(r(g_i)-\psi_{v,i}(g_i))=0$ gives $\prod_i(M(g_i)-c_i(g_i))=0$ for $g_1,\dots,g_n\in D_v$. Fix $i$ and $g\in D_v$. By (N1), for each $j\ne i$ there is $s_j\in D_v$ with $\psi^x_{v,i}(s_j)\ne\psi^x_{v,j}(s_j)$; put $g_j=s_j$ for $j\ne i$ and $g_i=g$. The $(i,i)$ entry of a product of upper triangular matrices is the product of their $(i,i)$ entries, so $\prod_j(\psi^x_{v,i}(g_j)-c_j(g_j))=0$. For $j\ne i$ the factor has $e_R$-value $\psi^x_{v,i}(s_j)-\psi^x_{v,j}(s_j)\ne0$, so it is a unit. Hence $c_i(g)=\psi^x_{v,i}(g)$, which lies in $\mathcal K$.

\emph{(c) The $U_{j,i}$.} Fix $j$, let $P=\rho(p(\phi_j))\in M_n(\mathcal K)$ and $c_k=\vartheta(U_{j,k})$, so that $e_R(c_k)=f(U_{j,k})$; by (e) of Definition~\ref{sha-patching-chain-datum}(1), these are pairwise distinct and are the roots of the characteristic polynomial of $P$. Applying $\Phi$ to $\prod_k(r(\phi_j)-U_{j,k})=0$ gives $\prod_k(P-c_k)=0$. Fix $i$, and let $\xi\in\mathcal K^n$ be nonzero with $P\xi=f(U_{j,i})\xi$. Applying the product to $\xi$ gives $\prod_k(f(U_{j,i})-c_k)\,\xi=0$. For $k\ne i$ the factor has $e_R$-value $f(U_{j,i})-f(U_{j,k})\ne0$, so it is a unit; and some entry of $\xi$ is a nonzero element of $\mathcal K$, hence a unit of $R$. So $c_i=f(U_{j,i})$, which lies in $\mathcal K$.

\emph{(d) The $y_{j,i}$.} With $P$, $c_k$ and $\xi$ as in (c), apply $\Phi$ to $\prod_{k\ne i}(r(\phi_j)-U_{j,k})\,(r(\tau_j)-1-y_{j,i})=0$. As $p(\tau_j)=1$, $\Phi(r(\tau_j))=1$, so $\prod_{k\ne i}(P-c_k)\cdot\vartheta(y_{j,i})=0$. Applied to $\xi$ this gives $\prod_{k\ne i}(f(U_{j,i})-c_k)\,\vartheta(y_{j,i})\,\xi=0$, so $\vartheta(y_{j,i})=0$, by the units of (c).

\emph{(e) The element $s$.} The image of $s$ in $\hat{\mathcal R}$ comes from $\mathcal O'$, so $\vartheta(s)=\hat f(s)$, which lies in $\mathcal K$.
% conj_const: he (Test.res with prAdd) and key for the reduction to K, the private conj_ex for Phi, vrRel (r1, r3, r4, r5: the relations killed by U.vr, from U.kills), testVal (t3 for s, t4 for the characters), Matrix.charpoly_mul_comm for (a), the private psi_lem (tri_prod, unit_of, one_zero) with oddDiag and x.N1 for (b), the private u_lem (eig_ex, eig_prod, eig_map, eig_zero) with P.eig and P.f_U for (c) and (d); then trConst (Lemma sha-patching-chain-rank-injective-generators).
\end{proof}

\begin{lemma}\label{sha-patching-chain-rank-injective}
In the setting of Definition~\ref{sha-patching-chain-tangent}, let $\alpha\in\mathcal T$, and suppose that there is an even matrix $Y$ over $\mathcal K$ with ${}^tY=-Y$ and $\alpha_C(g)=\rho(p(g))Y\rho(p(g))^{-1}-Y$ for all $g\in\Gamma_{\infty,0}$. Then $\alpha=0$.
\end{lemma}
\leanref{Development/AutomorphicArithmetic/AutomorphicVarieties/PatchingModules/SelmerScaling/ConjConst.lean}{UpperTheory.V.Pa.alInj}
\begin{proof}
Let $\vartheta:\mathcal A\to R_{1,\mathcal K}$ be the homomorphism of Lemma~\ref{sha-patching-chain-cocycles-test} with $e=1$, $\alpha_1=\alpha$, $\mu_1=0$ and $\Phi=0$, so that $\vartheta(b\cdot1+y)=\hat f(b)+\alpha(b)\epsilon'_1$. By the hypothesis, $\vartheta(r(g))=\rho(p(g))+\epsilon'_1(\rho(p(g))Y-Y\rho(p(g)))=(1-\epsilon'_1Y)\rho(p(g))(1+\epsilon'_1Y)$ for all $g$, as $\epsilon_1'^2=0$. So $C\vartheta(r(g))C^{-1}=\rho(p(g))$ with $C=1+\epsilon'_1Y$, whose inverse is $1-\epsilon'_1Y$. The algebra $R_{1,\mathcal K}$ is finite-dimensional over $\mathcal K$, the kernel of $y\mapsto y_0$ has zero cube, $\vartheta(y)_0=\hat f(y_0)$, and $\vartheta(o)=\hat f(o)$ on the image of $\mathcal O'$, as in (g) of the proof of Lemma~\ref{sha-patching-chain-cocycles-test}. So Lemma~\ref{sha-patching-chain-rank-injective-constant} gives $\vartheta(b\cdot1)=\hat f(b)$ for $b\in\hat{\mathcal R}$. The coefficient of $\epsilon'_1$ in $\vartheta(b\cdot1)$ is $\alpha(b)$, so $\alpha=0$.
% alInj: testEx with e = 1, hA and hu for the conjugation by 1 + E Y, then the private conj_const of ConjConst.lean (Lemma sha-patching-chain-rank-injective-constant), with Test from testEx; the coefficient of ep' 0 through prAdd.
\end{proof}

\begin{lemma}\label{sha-patching-chain-rank}
In the setting of Definition~\ref{sha-patching-chain-tangent}, let $D_0\ge0$, and suppose that the conclusion of Lemma~\ref{sha-patching-chain-laurent}(1) holds for $D_0$. Let $\alpha_1,\dots,\alpha_e\in\mathcal T$ be linearly independent, such that for all $i$ and $j$ the family with $\alpha_{i,C}$ in position $j$ and $0$ elsewhere lifts. Then $e+g^2\le D_0$.
\end{lemma}
\leanref{Development/AutomorphicArithmetic/AutomorphicVarieties/PatchingModules/SelmerScaling/FirstOrderLiftConstantPart.lean}{UpperTheory.V.Pa.claim1}
\begin{proof}
For each $i$, $\alpha_{i,C}\in\mathrm{Sel}^{\mathrm{ev}}(\mathcal K)$ has bounded denominators by Lemma~\ref{sha-patching-chain-cocycles}, so Lemma~\ref{sha-patching-chain-rank-lift} gives $y_i\in\mathrm{Sel}^1(\mathcal K)$ with bounded denominators and with coefficients of $1$ given by $\alpha_{i,C}$. For $1\le i',j\le g$ let $w_{i',j}(g)=\epsilon_j\,\kappa_{i'}(p(g))$, which lies in $\mathrm{Sel}^1(\mathcal K)$ and has bounded denominators by Lemma~\ref{sha-patching-chain-rank-basis}. It suffices to show that the $e+g^2$ maps $y_i$ and $w_{i',j}$ satisfy the hypothesis of Lemma~\ref{sha-patching-chain-laurent}(1); then $e+g^2\le D_0$.

Let $c_i,c_{i',j}\in\mathcal K$ with $\sum_ic_iy_i+\sum_{i',j}c_{i',j}w_{i',j}=\delta_1Y$, where $Y$ lies in the first module $V$ of Definition~\ref{sha-patching-chain-selmer}(3) over $\mathcal K$, ${}^tY=-Y$ and $\delta_1Y(g)=\rho_1(p(g))Y\rho_1(p(g))^{-1}-Y$. Write $Y=Y_0+\sum_m\epsilon_mY_m$ with $Y_0$ even and the $Y_m$ odd, as $Y\in V$; then ${}^tY_0=-Y_0$ and ${}^tY_m=-Y_m$.

\emph{The coefficients of $1$.} The coefficient of $1$ in $\delta_1Y(g)$ is $\rho(p(g))Y_0\rho(p(g))^{-1}-Y_0$, that of $y_i$ is $\alpha_{i,C}$ and that of $w_{i',j}$ is $0$. So $(\sum_ic_i\alpha_i)_C(g)=\rho(p(g))Y_0\rho(p(g))^{-1}-Y_0$, by the linearity in Lemma~\ref{sha-patching-chain-cocycles}, and $\sum_ic_i\alpha_i=0$ by Lemma~\ref{sha-patching-chain-rank-injective}. As the $\alpha_i$ are linearly independent, all $c_i$ vanish. Then $\rho(p(g))Y_0\rho(p(g))^{-1}=Y_0$ for all $g$, so $Y_0=0$ by Lemma~\ref{sha-patching-chain-rank-invariant}.

\emph{The coefficients of $\epsilon_m$.} As $Y_0=0$, the coefficient of $\epsilon_m$ in $\delta_1Y(g)$ is $\rho(p(g))Y_m\rho(p(g))^{-1}-Y_m$, since products of two of the $\epsilon$'s vanish. As the $c_i$ vanish, the coefficient of $\epsilon_m$ on the left side is $\sum_{i'}c_{i',m}\kappa_{i'}(p(g))$. So $\sum_{i'}c_{i',m}\,\kappa_{i'}\circ p\in\mathrm B^{\mathrm{od}}(\mathcal K)$, and all $c_{i',m}$ vanish by Lemma~\ref{sha-patching-chain-duality}(6).
% claim1: lift1 (FirstOrderMatrixMemSquareZero) gives y_i in Sel1 with constant part alC (alpha i); w (i', j) = DD.epM (Pi.single j (kapF i')); kapSel1 (EOneTranspose), alInj (ConjConst), invEv, indF; then the hypothesis SelmerRankBoundLaurent on the family Sum.elim y w.
\end{proof}

\subsection{The patching bound}\label{subsection-sha-patching-bound}

\begin{lemma}\label{sha-patching-chain-functionals-decomposition}
In the setting of Definition~\ref{sha-patching-chain-tangent}, let $\Phi\in\Pi$. For $y\in\mathcal N$ the map $D_y:b\mapsto\Phi(by)-\hat f(b)\Phi(y)$ lies in $\mathcal T$, and $D_{cy}=\hat f(c)D_y$ for $c\in\hat{\mathcal R}$. There are $e\ge0$, $\alpha'_1,\dots,\alpha'_e\in\mathcal T$ and $\mu_1,\dots,\mu_e\in\Lambda$ with $\Phi(by)=\hat f(b)\Phi(y)+\sum_k\alpha'_k(b)\mu_k(y)$ for $b\in\hat{\mathcal R}$ and $y\in\mathcal N$.
\end{lemma}
\leanref{Development/AutomorphicArithmetic/AutomorphicVarieties/PatchingModules/SelmerScaling/PhDecDApply.lean}{UpperTheory.V.Pa.phDec}
\begin{proof}
Every $b\in\hat{\mathcal R}$ is $o+q$ with $o$ in the image of $\mathcal O'$ and $q\in\mathfrak m_{\hat{\mathcal R}}$, as $\hat f$ maps $\mathcal O'$ onto $\mathcal K$ with kernel $\mathfrak m_{\hat{\mathcal R}}$. Then $D_y(b)=\Phi(qy)$, as $\Phi(oy)=\hat f(o)\Phi(y)$ and $\hat f(b)=\hat f(o)$. For $b'=o'+q'$, $(bb')y=o(o'y)+o(q'y)+o'(qy)+(qq')y$ with $qq'\in\mathfrak m_{\hat{\mathcal R}}^2\subset J$, so $\Phi((qq')y)=0$ and $D_y(bb')=\hat f(o)\Phi(q'y)+\hat f(o')\Phi(qy)=\hat f(b)D_y(b')+\hat f(b')D_y(b)$. $D_y$ vanishes on the image of $\mathcal O'$, and on $I_{\mathcal N}\subset J\cap\mathfrak m_{\hat{\mathcal R}}$ (Lemma~\ref{sha-patching-chain-duality}(1)), as $\Phi$ kills $J\mathcal N$. So $D_y\in\mathcal T$. Comparing $D_y(bc)=\Phi(bcy)-\hat f(bc)\Phi(y)$ with the product rule and $\Phi(cy)=\hat f(c)\Phi(y)+D_y(c)$ gives $D_{cy}(b)=\hat f(c)D_y(b)$.

By Lemma~\ref{sha-patching-chain-duality}(2), $\mathcal T$ has a finite basis $\alpha'_1,\dots,\alpha'_e$. Let $\mu_k(y)$ be the $k$-th coordinate of $D_y$ in this basis. Then $\mu_k$ is additive, $\mu_k(cy)=\hat f(c)\mu_k(y)$ by the previous paragraph, so $\mu_k\in\Lambda$; and $\Phi(by)-\hat f(b)\Phi(y)=D_y(b)=\sum_k\alpha'_k(b)\mu_k(y)$.
% phDec: the private phDecD with phDecD_mem (dec, hqq, h1, h2) and hs, tpDual for the finiteness, Module.finBasis and the private phDec_expand.
\end{proof}

\begin{lemma}\label{sha-patching-chain-functionals-lift-conjugate}
In the setting of Definition~\ref{sha-patching-chain-selmer} with $X=\mathcal K$, let $e\ge0$, let $u$ and $\psi'_{v,i}$ satisfy (a) to (g) of Lemma~\ref{sha-patching-chain-cocycles-selmer}, with $u_k$ and $u'_k$ the matrices of the coefficients of $\epsilon_k$ and $\epsilon'_k$, and let $z_1,\dots,z_e:\Gamma_{\infty,0}\to M_n(\mathcal K)$ be maps with $u_k\rho(p(\cdot))^{-1}-z_k\in\mathrm B^{\mathrm{od}}(\mathcal K)$ for each $k$. Then there is $u^\natural$ such that $u^\natural$ and the same $\psi'_{v,i}$ satisfy (a) to (g), and $u^\natural_k\rho(p(\cdot))^{-1}=z_k$ and $u'^\natural_k=u'_k$ for each $k$.
\end{lemma}
\leanref{Development/AutomorphicArithmetic/AutomorphicVarieties/PatchingModules/SelmerScaling/DefUConj.lean}{UpperTheory.V.Pa.defUConj}
\begin{proof}
For each $k$ let $Y_k$ be odd with ${}^tY_k=-Y_k$ and $u_k(g)\rho(p(g))^{-1}-z_k(g)=\rho(p(g))Y_k\rho(p(g))^{-1}-Y_k$ for all $g$. Let $P=1+\sum_k\epsilon_kY_k$ and $Q=1-\sum_k\epsilon_kY_k$. As products of two of the $\epsilon_k$ vanish, $PQ=QP=1$; and ${}^tP=Q$, ${}^tQ=P$, as the $Y_k$ are antisymmetric. Put $u^\natural(g)=Pu(g)Q$.

\emph{The conditions.} $u^\natural$ is a homomorphism of monoids, as $QP=1$. (a) The coefficients of $1$ in $P$ and $Q$ form the identity matrix, and $P$ and $Q$ have no coefficients of $\epsilon'_k$ or $\eta$. (b) The matrices whose entries lie in the span of $1$ and the $\epsilon'_k$ in the two diagonal blocks and in the span of the $\epsilon_k$ and $\eta$ in the two off-diagonal blocks form a ring, as products of two of the $\epsilon_k$, $\epsilon'_k$, $\eta$ lie in $\mathcal K\eta$ and vanish unless one is an $\epsilon_k$ and the other $\epsilon'_k$; $P$ and $Q$ belong to it, as the $Y_k$ are odd. (c) ${}^tu^\natural(g)\,u^\natural(c_\infty gc_\infty^{-1})=P\,{}^tu(g)\,QP\,u(c_\infty gc_\infty^{-1})Q=\varepsilon(g)^{1-n}PQ=\varepsilon(g)^{1-n}$. (d) $u^\natural(g)=PQ=1$ on the images of the $I'_u$. (e) holds, as the characters are unchanged. (f) $\prod_i\big(Pu(g_i)Q-\psi'_{v,i}(g_i)\big)=P\prod_i\big(u(g_i)-\psi'_{v,i}(g_i)\big)Q=0$, as the scalars commute with $P$ and $PQ=1$. (g) The coefficients of the entries of $u^\natural$ are $\mathcal K$-linear combinations of those of $u$.

\emph{The coefficients.} The coefficient of $\epsilon_k$ in $Pu(g)Q$ is $u_k(g)+Y_k\rho(p(g))-\rho(p(g))Y_k$, as the coefficients of $1$ of $u(g)$ are $\rho(p(g))$; so $u^\natural_k(g)\rho(p(g))^{-1}=u_k(g)\rho(p(g))^{-1}-\big(\rho(p(g))Y_k\rho(p(g))^{-1}-Y_k\big)=z_k(g)$. The coefficient of $\epsilon'_k$ in $Pu(g)Q$ is $u'_k(g)$, as $P$ and $Q$ have no coefficients of $\epsilon'_k$ and products of the $\epsilon_j$ with anything other than $1$ and the $\epsilon'_j$ vanish.
% defUConj: P and Q with twoPow_mul_twoPow (epM_mul), P_a, P_ep, P_ep', P_tr, even_P and even_mul for (b), hc and conj_prod for (f), bnd_mul and bnd_const for (g), mmul_ep and mmul_ep' for the coefficients.
\end{proof}

\begin{lemma}\label{sha-patching-chain-functionals-lift-push}
In the setting of Definition~\ref{sha-patching-chain-selmer} with $X=\mathcal K$, let $e\ge0$, let $u$ and $\psi'_{v,i}$ satisfy (a) to (g) of Lemma~\ref{sha-patching-chain-cocycles-selmer}, with $u_k$ and $u'_k$ as there, and let $c_{k,j}\in\mathcal K$, for $1\le k\le e$ and $1\le j\le g$, be such that $u_k\rho(p(\cdot))^{-1}=\sum_jc_{k,j}\,\kappa_j\circ p$ for each $k$. Then there are $u^+:\Gamma_{\infty,0}\to M_n(R_{g,\mathcal K})$ and characters $\psi^+_{v,i}$ that satisfy (a) to (g) of Lemma~\ref{sha-patching-chain-cocycles-selmer} with $g$ in place of $e$, such that $u^+_j\rho(p(\cdot))^{-1}=\kappa_j\circ p$ and $u'^+_j=\sum_kc_{k,j}u'_k$ for each $j$.
\end{lemma}
\leanref{Development/AutomorphicArithmetic/AutomorphicVarieties/PatchingModules/SelmerScaling/DefUPush.lean}{UpperTheory.V.Pa.defUPush}
\leanref{Development/AutomorphicArithmetic/AutomorphicVarieties/PatchingModules/SelmerScaling/RelSub.lean}{UpperTheory.V.Pa.relSub}
\begin{proof}
Write $y_0$, $y_{\epsilon_k}$, $y_{\epsilon'_k}$ and $y_\eta$ for the coefficients of an element $y$ of $R_{e,\mathcal K}$ or $R_{g,\mathcal K}$. For $y\in R_{e,\mathcal K}$ and $y'\in R_{g,\mathcal K}$ write $y\sim y'$ if $y_0=y'_0$, $y_\eta=y'_\eta$, $y_{\epsilon_k}=\sum_jc_{k,j}y'_{\epsilon_j}$ for all $k$, and $y'_{\epsilon'_j}=\sum_kc_{k,j}y_{\epsilon'_k}$ for all $j$.

\emph{The relation.} (i) $r\sim r$ for $r\in\mathcal K$. (ii) If $y\sim y'$ and $z\sim z'$, then $y\pm z\sim y'\pm z'$ and $yz\sim y'z'$. The first is clear. For the product, the coefficients of $1$, $\epsilon_k$ and $\epsilon'_j$ are $y_0z_0$, $y_0z_{\epsilon_k}+y_{\epsilon_k}z_0$ and $y'_0z'_{\epsilon'_j}+y'_{\epsilon'_j}z'_0$, which are related as required; for $\eta$, $(yz)_\eta=y_0z_\eta+y_\eta z_0+\sum_k(y_{\epsilon_k}z_{\epsilon'_k}+y_{\epsilon'_k}z_{\epsilon_k})$, and
\[\sum_jy'_{\epsilon_j}z'_{\epsilon'_j}=\sum_jy'_{\epsilon_j}\sum_kc_{k,j}z_{\epsilon'_k}=\sum_k\Big(\sum_jc_{k,j}y'_{\epsilon_j}\Big)z_{\epsilon'_k}=\sum_ky_{\epsilon_k}z_{\epsilon'_k},\]
and in the same way $\sum_jy'_{\epsilon'_j}z'_{\epsilon_j}=\sum_ky_{\epsilon'_k}z_{\epsilon_k}$. So the relation passes to entrywise related matrices and their sums and products. (iii) If $y\sim y'_1$, $y\sim y'_2$ and $y'_1$ and $y'_2$ have the same coefficients of $1$ and the $\epsilon_j$, then $y'_1=y'_2$, as the coefficients of $\epsilon'_j$ and $\eta$ are determined by $y$.

\emph{The construction.} Let $u^+(\gamma)$ be the matrix whose entries have the coefficients of $1$ and $\eta$ of those of $u(\gamma)$, the coefficients of $\epsilon_j$ of those of $\rho_1(p(\gamma))$, which form the matrix $\kappa_j(p(\gamma))\rho(p(\gamma))$, and the coefficients of $\epsilon'_j$ given by $\sum_kc_{k,j}u'_k(\gamma)$. Then $u(\gamma)\sim u^+(\gamma)$ entrywise, as $u_k(\gamma)=\sum_jc_{k,j}\kappa_j(p(\gamma))\rho(p(\gamma))$. Let $\psi^+_{v,i}(g)$ have the coefficient of $1$ of $\psi'_{v,i}(g)$, the coefficients of $\epsilon'_j$ given by $\sum_kc_{k,j}\psi'_{v,i}(g)_{\epsilon'_k}$, and no others. The map $y\mapsto y^+$ so defined, from the span of $1$ and the $\epsilon'_k$ to the span of $1$ and the $\epsilon'_j$, is a ring homomorphism, as products of the $\epsilon'$ vanish; so $\psi^+_{v,i}$ is a homomorphism, and $\psi'_{v,i}(g)\sim\psi^+_{v,i}(g)$, since $\psi'_{v,i}(g)$ has no coefficients of $\epsilon_k$ or $\eta$ by (e). Let $\tau:R_{g,\mathcal K}\to R_{g,\mathcal K}$ be the map that keeps the coefficients of $1$ and the $\epsilon_j$ and kills the others; it is a ring homomorphism, as products of two of the $\epsilon_j$ vanish. Then $\tau(u^+(\gamma))=\rho_1(p(\gamma))$, by (a), and $\tau(\psi^+_{v,i}(g))=\psi_{v,i}(g)$, the character of $x$, by (e).

\emph{The conditions.} The identities $u^+(1)=1$ and $u^+(\gamma\gamma')=u^+(\gamma)u^+(\gamma')$, and those of (c), (d) and (f), are proved by (ii) and (iii): the two sides of each are related to the two sides of the same identity for $u$ and the $\psi'_{v,i}$, which are equal, and their images under $\tau$ agree. For all but (f), the images under $\tau$ are the two sides of the same identity for $\rho_1\circ p$, which hold by Lemma~\ref{sha-patching-chain-rank-lift-triangular}(4). For (f), $\tau$ carries $\prod_i(u^+(g_i)-\psi^+_{v,i}(g_i))$ to $\prod_i(\rho_1(p(g_i))-\psi_{v,i}(g_i))$, which vanishes: after conjugation by $w'_v$ the $i$-th factor is upper triangular with $(i,i)$ entry $0$, by Lemma~\ref{sha-patching-chain-rank-lift-triangular}(3), as in step (i) of the proof of Lemma~\ref{sha-patching-chain-cocycles-ordinary}. (a) and (e) hold by construction. (b) holds, as the coefficients of $\epsilon_j$ of $\rho_1(p(\gamma))$ are odd and those of $u(\gamma)$ satisfy (b). (g): the coefficients of $1$, $\eta$ and $\epsilon'_j$ are those of $u$ or $\mathcal K$-linear combinations of them, and those of $\epsilon_j$ are the entries of $\kappa_j(p(\gamma))\rho(p(\gamma))$, sums of products of maps to $A$ with locally constant reductions modulo every power of $T$.

\emph{The coefficients.} $u^+_j(\gamma)\rho(p(\gamma))^{-1}=\kappa_j(p(\gamma))$ and $u'^+_j=\sum_kc_{k,j}u'_k$ by construction.
% defUPush: the private PP (the book's u^+ entrywise) and pushU (the book's psi^+), TT (the book's tau), RM with relSub (Lemma: rel_C, rel_add, rel_mul, the uniqueness, and the products of lists) for RM_uniq, RM_mul, RM_one, RM_smul, RM_prod; rho1Tri for the identities of rho_1, conj_prod with tri_prod_zero for (f), bndRho and bndOps for (g).
\end{proof}

\begin{lemma}\label{sha-patching-chain-functionals-lift-second}
In the setting of Definition~\ref{sha-patching-chain-selmer} with $X=\mathcal K$, let $u:\Gamma_{\infty,0}\to M_n(R_{g,\mathcal K})$ and $\psi'_{v,i}$ satisfy (a) to (g) of Lemma~\ref{sha-patching-chain-cocycles-selmer} with $e=g$, and suppose that $u_j\rho(p(\cdot))^{-1}=\kappa_j\circ p$ for each $j$. Then the map $\Psi(\gamma)=u(\gamma)\rho_1(p(\gamma))^{-1}-1$ lies in $\mathrm{Sel}^2(\mathcal K)$ and has bounded denominators, and $\Psi'_j=u'_j\rho(p(\cdot))^{-1}$ for each $j$. In particular, the family $(u'_j\rho(p(\cdot))^{-1})_{1\le j\le g}$ lifts.
\end{lemma}
\leanref{Development/AutomorphicArithmetic/AutomorphicVarieties/PatchingModules/SelmerScaling/DefULift.lean}{UpperTheory.V.Pa.defULift}
\begin{proof}
Let $\mathcal I\subset R_{g,\mathcal K}$ be the ideal of the elements whose coefficients of $1$ and of the $\epsilon_j$ vanish, the span of the $\epsilon'_j$ and $\eta$; then $\mathcal I^2=0$. Write $\rho'_1=\rho_1\circ p$, a homomorphism by Lemma~\ref{sha-patching-chain-rank-lift-triangular}(4).

\emph{First properties.} The matrix $u(\gamma)-\rho'_1(\gamma)$ has entries in $\mathcal I$: the coefficients of $1$ agree by (a), and those of $\epsilon_j$ are $u_j(\gamma)=\kappa_j(p(\gamma))\rho(p(\gamma))$ on both sides. So $\Psi(\gamma)=(u(\gamma)-\rho'_1(\gamma))\rho'_1(\gamma)^{-1}$ has entries in $\mathcal I$, and $\Psi(\gamma)\Psi(\gamma')=0$. Let $E$ be the set of matrices whose entries lie in the span of $1$ and the $\epsilon'_j$ in the two diagonal blocks and in the span of the $\epsilon_j$ and $\eta$ in the two off-diagonal blocks; as in the proof of Lemma~\ref{sha-patching-chain-functionals-lift-conjugate}, it is a $\mathcal K$-subalgebra of $M_n(R_{g,\mathcal K})$. It contains $u(\gamma)$, by (b), and $\rho'_1(\gamma)^{\pm1}$, which lie in the first module. A matrix in $E$ with entries in $\mathcal I$ lies in the second module $V$ of Definition~\ref{sha-patching-chain-selmer}(3); so $\Psi(\gamma)\in V$.

\emph{The relations.} The reductions modulo $T^m\mathcal K=\mathcal K$ are constant. As $(u(g)-\rho'_1(g))\Psi(g')=0$,
\[\Psi(gg')-\Psi(g)-\rho'_1(g)\Psi(g')\rho'_1(g)^{-1}=(u(g)-\rho'_1(g))\Psi(g')\rho'_1(g)^{-1}=0.\]
For the relation with $c_\infty$, let $g^c=c_\infty gc_\infty^{-1}$. As $\Psi(g)^2=0$, $(1-\Psi(g))(1+\Psi(g))=1$, and $1+\Psi(g)=u(g)\rho'_1(g)^{-1}$, so $\rho'_1(g)u(g^{-1})=1-\Psi(g)$. By (c) and Lemma~\ref{sha-patching-chain-rank-lift-triangular}(4), ${}^tu(g)u(g^c)={}^t\rho'_1(g)\rho'_1(g^c)$, so
\[\Psi(g^c)={}^tu(g^{-1})\,{}^tu(g)\,u(g^c)\,\rho'_1(g^c)^{-1}-1={}^tu(g^{-1})\,{}^t\rho'_1(g)-1={}^t\big(1-\Psi(g)\big)-1=-{}^t\Psi(g).\]
On the images of the $I'_u$, $u=1$ by (d) and $\rho'_1=1$, so $\Psi=0$.

\emph{The local condition.} Let $v\mid\ell$. Apply Lemma~\ref{sha-patching-chain-cocycles-ordinary}(1) with $F=\mathcal K$, which is infinite; $R=R_{g,\mathcal K}$, with $e$ the coefficient of $1$, whose kernel consists of elements of zero cube; $I=\mathcal I$; $H=D_v$; $\rho_0=\rho'_1$ and $\rho_1=u$ on $D_v$; $\psi_{0,i}$ the characters $\psi_{v,i}$ of $x$ and $\psi_{1,i}=\psi'_{v,i}$; $w=w'_v$, with inverse $w_v^{-1}(1+\sum_k\epsilon_kY_{k,v})$; and the subalgebra $E$. Its hypotheses hold: (a) by Lemma~\ref{sha-patching-chain-rank-lift-triangular}(3); (b) by the first paragraph, and by (e), which gives $\psi'_{v,i}(g)-\psi_{v,i}(g)\in\mathcal I$; (c) by (N1); (d) by (f); and (e), as $\rho'_1$, $u$ and the scalars, which lie in the span of $1$ and the $\epsilon'_j$ by (e), lie in $E$. It gives $Y\in E$ with entries in $\mathcal I$, so $Y\in V$, such that for $g\in D_v$
\[{w'_v}^{-1}\Big(\Psi(g)-\big(\rho'_1(g)Y\rho'_1(g)^{-1}-Y\big)\Big)w'_v\]
is upper triangular with diagonal $(\psi'_{v,i}(g)-\psi_{v,i}(g))\psi_{v,i}(g)^{-1}$; for $g\in\mu_v$ this is $0$, as $\psi'_{v,i}=1$ there by (e) and $\psi_{v,i}=1$ there, $x$ being normalized. So $\Psi\in\mathrm{Sel}^2(\mathcal K)$.

\emph{Bounded denominators and the coefficients.} The coefficients of the entries of $\Psi(\gamma)$ are sums of products of those of $u(\gamma)$, which have bounded denominators by (g), and of $\rho'_1(\gamma^{-1})$, which are entries of $\rho(p(\gamma^{-1}))$ and of $\kappa_j(p(\gamma^{-1}))\rho(p(\gamma^{-1}))$, maps to $A$ with locally constant reductions. As $\rho'_1(\gamma)^{-1}$ has no coefficients of $\epsilon'_j$, and a product has a coefficient of $\epsilon'_j$ only from the coefficient of $1$ of one factor and that of $\epsilon'_j$ of the other, $\Psi'_j(\gamma)=u'_j(\gamma)\rho(p(\gamma))^{-1}$.
% defULift: the private II (the book's ideal I, II_mul, II_mat, II_right), aug_nil, EE (the book's E, with EE_scal, EE_of_even, mem_even_entries_of_first_order), mem_second_order_of_reductions for the values; hD, hZ and hsq; the cocycle relation by noncomm_ring with II_mat; the form through uinv, e2, e3 and e4; the local condition through locLemma (Lemma sha-patching-chain-cocycles-ordinary(1)) with R = DD x.Fr k.g, I = II and w = k.w1, from rho1Tri; bnd_mul with bndRho and bndOps; first_order_lift_apply and sum_cmp for the coefficients.
\end{proof}

\begin{lemma}\label{sha-patching-chain-functionals-lift}
In the setting of Definition~\ref{sha-patching-chain-selmer} with $X=\mathcal K$, let $e\ge0$, and let $u$ and $\psi'_{v,i}$ satisfy (a) to (g) of Lemma~\ref{sha-patching-chain-cocycles-selmer}, in whose notation $u_k$ and $u'_k$ are the matrices of the coefficients of $\epsilon_k$ and $\epsilon'_k$. Let $c_{k,j}\in\mathcal K$, for $1\le k\le e$ and $1\le j\le g$, be such that $u_k\rho(p(\cdot))^{-1}-\sum_jc_{k,j}\,\kappa_j\circ p\in\mathrm B^{\mathrm{od}}(\mathcal K)$ for each $k$. Then the family $\big(\sum_kc_{k,j}\,u'_k\rho(p(\cdot))^{-1}\big)_{1\le j\le g}$ lifts.
\end{lemma}
\leanref{Development/AutomorphicArithmetic/AutomorphicVarieties/PatchingModules/SelmerScaling/DefUConj.lean}{UpperTheory.V.Pa.defUConj}
\leanref{Development/AutomorphicArithmetic/AutomorphicVarieties/PatchingModules/SelmerScaling/DefUPush.lean}{UpperTheory.V.Pa.defUPush}
\leanref{Development/AutomorphicArithmetic/AutomorphicVarieties/PatchingModules/SelmerScaling/DefULift.lean}{UpperTheory.V.Pa.defULift}
\begin{proof}
Lemma~\ref{sha-patching-chain-functionals-lift-conjugate}, with $z_k=\sum_jc_{k,j}\,\kappa_j\circ p$, gives $u^\natural$, with the same characters, satisfying (a) to (g), with $u^\natural_k\rho(p(\cdot))^{-1}=\sum_jc_{k,j}\,\kappa_j\circ p$ and $u'^\natural_k=u'_k$. Lemma~\ref{sha-patching-chain-functionals-lift-push}, applied to $u^\natural$, gives $u^+$ over $R_{g,\mathcal K}$ and characters $\psi^+_{v,i}$ satisfying (a) to (g), with $u^+_j\rho(p(\cdot))^{-1}=\kappa_j\circ p$ and $u'^+_j=\sum_kc_{k,j}u'_k$. By Lemma~\ref{sha-patching-chain-functionals-lift-second}, the family $(u'^+_j\rho(p(\cdot))^{-1})_j=\big(\sum_kc_{k,j}u'_k\rho(p(\cdot))^{-1}\big)_j$ lifts.
% claim2 composes defUConj, defUPush and defULift (Lemmas sha-patching-chain-functionals-lift-conjugate, -push and -second).
\end{proof}

\begin{lemma}\label{sha-patching-chain-functionals}
In the setting of Definition~\ref{sha-patching-chain-tangent}, let $m_1,\dots,m_g\in\mathcal N$ be dual to the Selmer basis and $\Phi\in\Pi$. For $1\le j\le g$ let $\alpha_j(b)=\Phi(bm_j)-\hat f(b)\Phi(m_j)$ for $b\in\hat{\mathcal R}$. Then $\alpha_j\in\mathcal T$ for all $j$, and $(\alpha_1,\dots,\alpha_g)\in\mathcal V$.
\end{lemma}
\leanref{Development/AutomorphicArithmetic/AutomorphicVarieties/PatchingModules/SelmerScaling/Claim2.lean}{UpperTheory.V.Pa.claim2}
\begin{proof}
Lemma~\ref{sha-patching-chain-functionals-decomposition} gives $\alpha'_1,\dots,\alpha'_e\in\mathcal T$ and $\mu_1,\dots,\mu_e\in\Lambda$ with $\Phi(by)=\hat f(b)\Phi(y)+\sum_k\alpha'_k(b)\mu_k(y)$. Put $c_{k,j}=\mu_k(m_j)$. Then $\alpha_j(b)=\sum_kc_{k,j}\alpha'_k(b)$, so $\alpha_j=\sum_kc_{k,j}\alpha'_k\in\mathcal T$. Lemma~\ref{sha-patching-chain-cocycles-test}, for $e$, the $\alpha'_k$, the $\mu_k$ and $\Phi$, gives $u$ and $\psi'_{v,i}$ satisfying (a) to (g) of Lemma~\ref{sha-patching-chain-cocycles-selmer}, with $u_k\rho(p(\cdot))^{-1}=\mu_{k,C}$ and $u'_k\rho(p(\cdot))^{-1}=\alpha'_{k,C}$. As the $m_j$ are dual to the Selmer basis, $\mu_{k,C}-\sum_jc_{k,j}\,\kappa_j\circ p\in\mathrm B^{\mathrm{od}}(\mathcal K)$. So by Lemma~\ref{sha-patching-chain-functionals-lift} the family $(\sum_kc_{k,j}\alpha'_{k,C})_j$ lifts, and this family is $(\alpha_{j,C})_j$ by the linearity in Lemma~\ref{sha-patching-chain-cocycles}. So $(\alpha_1,\dots,\alpha_g)\in\mathcal V$.
% claim2: phDec (PhDecDApply: Phi(b y) = rF b Phi y + sum_k alpha_k b mu_k y), defUTest, defUConj (with the duality of m), defUPush with c k j = mu_k (m j), defULift; alpha_j = sum_k c k j alpha_k, and alSel for the linearity.
\end{proof}

\begin{lemma}\label{sha-patching-chain-dual-basis}
In the setting of Definition~\ref{sha-patching-chain-tangent}, suppose that the conclusion of Lemma~\ref{sha-patching-chain-laurent}(3) holds. Then there are $m_1,\dots,m_g\in\mathcal N$, dual to the Selmer basis, that generate $\mathcal N$ as an $\hat{\mathcal R}$-module.
\end{lemma}
\leanref{Development/AutomorphicArithmetic/AutomorphicVarieties/PatchingModules/SelmerScaling/DualEx.lean}{UpperTheory.V.Pa.dualEx}
\begin{proof}
The set $\mathrm B^{\mathrm{od}}(\mathcal K)$ is a $\mathcal K$-subspace. Let $\mu\in\Lambda$. By Lemma~\ref{sha-patching-chain-cocycles}, $\mu_C\in\mathrm{Sel}^{\mathrm{od}}(\mathcal K)$ has bounded denominators, so there is $c(\mu)\in\mathcal K^g$ with $\mu_C-\sum_jc_j(\mu)\,\kappa_j\circ p\in\mathrm B^{\mathrm{od}}(\mathcal K)$. It is unique, by Lemma~\ref{sha-patching-chain-duality}(6); so $\mu\mapsto c(\mu)$ is $\mathcal K$-linear, as $\mu\mapsto\mu_C$ is. It is injective: if $c(\mu)=0$, then $\mu_C\in\mathrm B^{\mathrm{od}}(\mathcal K)$ and $\mu=0$ by Lemma~\ref{sha-patching-chain-duality}(5). By Lemma~\ref{sha-patching-chain-duality}(3), each linear form $\mu\mapsto c_j(\mu)$ is $\mu\mapsto\mu(m_j)$ for some $m_j\in\mathcal N$; then the $m_j$ are dual to the Selmer basis.

As $\mathcal N$ is finite over $\hat{\mathcal R}$, by Nakayama's lemma it suffices that every $y\in\mathcal N$ lies in the span of the $m_j$ plus $\mathfrak m_{\hat{\mathcal R}}\mathcal N$. Since $\mu\mapsto c(\mu)$ is injective, the linear form $\mu\mapsto\mu(y)$ on $\Lambda$ is $\mu\mapsto\sum_jc_j(\mu)a_j$ for some $a_1,\dots,a_g\in\mathcal K$. Choose $b_j\in\hat{\mathcal R}$ with $\hat f(b_j)=a_j$. Then $\mu(y-\sum_jb_jm_j)=\mu(y)-\sum_ja_jc_j(\mu)=0$ for every $\mu\in\Lambda$, so $y-\sum_jb_jm_j\in\mathfrak m_{\hat{\mathcal R}}\mathcal N$ by Lemma~\ref{sha-patching-chain-duality}(3).
% dualEx: the submodule Bs of BdM, T = laC (linear by alSel), laDual, indF, laInj, then exists_dual_functional_complement (key for uniqueness, the linear map C, hCinj, hCj for m, and Submodule.le_of_le_smul_of_le_jacobson_bot with LinearMap.dualMap_surjective_of_injective for the span).
\end{proof}

\begin{lemma}\label{sha-patching-chain-presentation}
In the setting of Definition~\ref{sha-patching-chain-tangent}, let $D_0\ge0$, and suppose that the conclusions of Lemma~\ref{sha-patching-chain-laurent}(1), for $D_0$, (2) and (3) hold. Then there are $X_1,\dots,X_{D_0}\in\mathfrak m_{\hat{\mathcal R}}$ with $\mathfrak m_{\hat{\mathcal R}}\subset(X_1,\dots,X_{D_0})+J$, elements $m_1,\dots,m_g$ generating $\mathcal N$ as an $\hat{\mathcal R}$-module, and $c_{i,i',k}\in\mathcal O'$, for $1\le i,i'\le g$ and $1\le k\le D_0$, such that the $g^2$ vectors $(\bar c_{i,i',k})_k$ of residues are linearly independent and $\sum_{i'}\big(\sum_kc_{i,i',k}X_k\big)m_{i'}\in J\mathcal N$ for each $i$.
\end{lemma}
\leanref{Development/AutomorphicArithmetic/AutomorphicVarieties/PatchingModules/SelmerScaling/SelmerSpannedByKappaLaurent.lean}{UpperTheory.V.Pa.Ctx.Pres}
\leanref{Development/AutomorphicArithmetic/AutomorphicVarieties/PatchingModules/SelmerScaling/AlCSmul.lean}{UpperTheory.V.Pa.pres}
\begin{proof}
\emph{The subspace $\mathcal V$.} The map $\alpha\mapsto\alpha_C$ is $\mathcal K$-linear by Lemma~\ref{sha-patching-chain-cocycles}. The families that lift form a $\mathcal K$-subspace: $\mathrm{Sel}^2(\mathcal K)$ is a $\mathcal K$-subspace, and sums and $\mathcal K$-multiples of maps with bounded denominators have bounded denominators, as every element of $\mathcal K$ is a power of $T$ times an element of $A$. So $\mathcal V$ is a subspace of $\mathcal T^g$, and $\mathcal T$ is finite-dimensional by Lemma~\ref{sha-patching-chain-duality}(2).

\emph{The linear algebra.} Hypothesis (1) of Lemma~\ref{sha-patching-chain-linear} holds for $\mathcal V$: given $z_0,\dots,z_g\in\mathcal T^g$, with components $z_{i,j}\in\mathcal T$, the maps $(z_{i,j})_C$ lie in $\mathrm{Sel}^{\mathrm{ev}}(\mathcal K)$ and have bounded denominators, by Lemma~\ref{sha-patching-chain-cocycles}; Lemma~\ref{sha-patching-chain-laurent}(2) gives $c_0,\dots,c_g$, not all zero, such that $(\sum_ic_i(z_{i,j})_C)_j$ lifts, and this family is $\big((\sum_ic_iz_{i,j})_C\big)_j$. Hypothesis (2) is Lemma~\ref{sha-patching-chain-rank}. So Lemma~\ref{sha-patching-chain-linear} gives linear forms $X'_1,\dots,X'_{D_0}$ spanning the dual of $\mathcal T$ and $c'_{i,i',k}\in\mathcal K$. By Lemma~\ref{sha-patching-chain-duality}(2), there are $X_k\in\mathfrak m_{\hat{\mathcal R}}$ with $\alpha(X_k)=X'_k(\alpha)$ for $\alpha\in\mathcal T$. Choose $c_{i,i',k}\in\mathcal O'$ with $\hat f(c_{i,i',k})=c'_{i,i',k}$.

\emph{The residues.} The restriction of $\hat f$ to $\mathcal O'$ has kernel $\mathfrak m_{\mathcal O'}$, so it identifies the residue field of $\mathcal O'$ with $\mathcal K$, and it carries $\bar c_{i,i',k}$ to $c'_{i,i',k}$. So the vectors of residues are linearly independent.

\emph{A product rule.} For $\alpha\in\mathcal T$, $o$ in the image of $\mathcal O'$ and $b\in\mathfrak m_{\hat{\mathcal R}}$, $\alpha(ob)=\hat f(o)\alpha(b)$, since $\alpha(o)=0$ and $\hat f(b)=0$.

\emph{The maximal ideal.} Let $b\in\mathfrak m_{\hat{\mathcal R}}$. The linear form $\alpha\mapsto\alpha(b)$ on $\mathcal T$ is $\sum_k\lambda_kX'_k$ for some $\lambda_k\in\mathcal K$; choose $o_k$ in the image of $\mathcal O'$ with $\hat f(o_k)=\lambda_k$. Then $b-\sum_ko_kX_k\in\mathfrak m_{\hat{\mathcal R}}$ and $\alpha(b-\sum_ko_kX_k)=\alpha(b)-\sum_k\lambda_kX'_k(\alpha)=0$ for all $\alpha\in\mathcal T$, by the product rule. So $b-\sum_ko_kX_k\in J$ by Lemma~\ref{sha-patching-chain-duality}(2), and $b\in(X_1,\dots,X_{D_0})+J$.

\emph{The relations.} Let $m_1,\dots,m_g$ be as in Lemma~\ref{sha-patching-chain-dual-basis}. Fix $i$, put $b_{i'}=\sum_kc_{i,i',k}X_k\in\mathfrak m_{\hat{\mathcal R}}$ and $y=\sum_{i'}b_{i'}m_{i'}$. By Lemma~\ref{sha-patching-chain-duality}(4) it suffices that $\Phi(y)=0$ for every $\Phi\in\Pi$. Let $\alpha_{i'}$ be as in Lemma~\ref{sha-patching-chain-functionals}; then $(\alpha_1,\dots,\alpha_g)\in\mathcal V$. As $\hat f(b_{i'})=0$, $\Phi(b_{i'}m_{i'})=\alpha_{i'}(b_{i'})=\sum_kc'_{i,i',k}X'_k(\alpha_{i'})$, by the product rule. Summing over $i'$, $\Phi(y)=\sum_{i'}\sum_kc'_{i,i',k}X'_k(\alpha_{i'})=0$ by Lemma~\ref{sha-patching-chain-linear}.
% pres: tpDual, alSel, liftsSub (Lifts is a subspace: L0, Ladd, Lsmul, through bndOps), hA from EvenCocyclesLiftLaurent, hB from claim1 (SelmerRankBoundLaurent), linAlg, hlift for Xb, oo_surj for cO, dualEx (SelmerSpannedByKappaLaurent), hmulT for the product rule, hker for the residues, hsep for the maximal ideal, phSep with claim2 and hrel for the relations.
\end{proof}

\begin{lemma}\label{sha-patching-chain-thick}
Let $R$ be a commutative ring, $G$ a group and $I\subset R[G]$ a subset, and suppose that there is a finite set $s\subset G$ such that every element of $R[G]$ is congruent to an $R$-linear combination of elements of $s$ modulo the two-sided ideal generated by $I$. For a homomorphism $\varphi$ from $R$ to a field $\kappa$, let $E_\varphi$ be the quotient of $\kappa[G]$ by the two-sided ideal generated by the image of $I$.
\begin{enumerate}
\item Let $\varphi:R\to\kappa$ and $\varphi':R\to\kappa'$ be homomorphisms to fields with the same kernel, and suppose that $E_\varphi$ is not semisimple. Then there is a field extension $K$ of $\kappa'$ such that $E_{\varphi''}$ is not semisimple for the composite $\varphi'':R\to\kappa'\to K$.
\item Let $R'$ be a commutative ring, $p:G'\to G$ a surjective homomorphism of groups, $I'\subset R'[G']$ a subset, and $\varphi:R\to\kappa$ and $\varphi':R'\to\kappa$ homomorphisms such that the image of $I'$ in $\kappa[G]$, under $\varphi'$ and $p$, lies in the two-sided ideal generated by the image of $I$. If $E_\varphi$ is not semisimple, then neither is the quotient of $\kappa[G']$ by the two-sided ideal generated by the image of $I'$.
\end{enumerate}
\end{lemma}
\leanref{Development/AutomorphicArithmetic/AutomorphicVarieties/PatchingModules/PatchingBaseChange/NonSSField.lean}{UpperTheory.V.Pa.BC.nonSS_field}
\leanref{Development/AutomorphicArithmetic/AutomorphicVarieties/PatchingModules/PatchingBaseChange/NonSSMono.lean}{UpperTheory.V.Pa.BC.nonSS_mono}
\begin{proof}
(1) Let $\kappa_0$ be the fraction field of $R/\ker\varphi$. Both $\varphi$ and $\varphi'$ factor through $\kappa_0$, and there is a field $K$ with embeddings of $\kappa$ and $\kappa'$ over $\kappa_0$, for instance a quotient of $\kappa\otimes_{\kappa_0}\kappa'$ by a maximal ideal. Then $\varphi''$ is also the composite of $\varphi$ with $\kappa\to K$, so $E_{\varphi''}=K\otimes_\kappa E_\varphi$. The algebra $E_\varphi$ is finite-dimensional over $\kappa$, spanned by the images of $s$, so its Jacobson radical $\mathfrak r$ is nilpotent and nonzero, as $E_\varphi$ is not semisimple. Then $K\otimes_\kappa\mathfrak r$ is a nonzero nilpotent two-sided ideal of $E_{\varphi''}$, which is therefore not semisimple.

(2) The homomorphism $\kappa[G']\to\kappa[G]$ induced by $p$ is surjective and carries the ideal generated by the image of $I'$ into that generated by the image of $I$. So $E_\varphi$ is a quotient of the quotient of $\kappa[G']$, and a quotient of a semisimple ring is semisimple.
% nonSS_field: common_field for K, the private fin_quot for the finiteness, grp_quot_bc for the base change of the quotient, ss_descend (the radical argument) for the descent of semisimplicity. nonSS_mono: RingHom.isSemisimpleRing_of_surjective.
\end{proof}

\begin{lemma}\label{sha-patching-chain-ideals}
Let $g:B\to B'$ be a local homomorphism of local rings, $p:G\to G'$ a continuous homomorphism of topological groups, $D$ a determinant of $G$ of dimension $n$ with values in $B$ and $D'$ a determinant of $G'$ of dimension $n$ with values in $B'$ \cite{Che14}, such that $D\otimes_BB'=D'\circ p$. Let $\varphi:B[G]\to B'[G']$ be the ring homomorphism induced by $g$ and $p$, and $Z\subset B[G]$ a subset.
\begin{enumerate}
\item $\varphi$ maps the closure of $\operatorname{CH}(D)\cup Z$ into the closure of $\operatorname{CH}(D')\cup\varphi(Z)$, in the sense of Definition~\ref{sha-patching-chain-cayley-closure}.
\item Suppose that $G'=G$ is profinite, $p$ is the identity and $D'=D\otimes_BB'$; that $B$ is complete local noetherian with finite residue field $k$ and $B'$ is noetherian; and that $D$ is continuous, with reduction $D\otimes_Bk$ the determinant of $\chi_1\oplus\dots\oplus\chi_n$ for a family of characters $\chi_i:G\to k^\times$ with an extension bound (Definition~\ref{sha-patching-chain-cayley-extension-bound}). Let $J$ be the two-sided ideal of $B'[G]$ generated by the image under $\varphi$ of the closure of $\operatorname{CH}(D)\cup Z$, and let $Z'\subset J$. Then the closure of $\operatorname{CH}(D')\cup Z'$ lies in $J$.
\end{enumerate}
\end{lemma}
\leanref{Development/AutomorphicArithmetic/AutomorphicVarieties/PatchingModules/AnnihilatorBaseChange/ChBarPush.lean}{UpperTheory.V.Pa.Core.chBarPush}
\leanref{Development/AutomorphicArithmetic/AutomorphicVarieties/PatchingModules/PatchingBaseChange/ClosureBc.lean}{UpperTheory.V.Pa.BC.closure_bc}
\begin{proof}
(1) Let $N\ge1$ and $V'$ an open normal subgroup of $G'$; then $V=p^{-1}(V')$ is an open normal subgroup of $G$. The closure of $\operatorname{CH}(D)\cup Z$ lies in the two-sided ideal generated by $\operatorname{CH}(D)$, $Z$, $\mathfrak m_B^N$ and the $v-1$ with $v\in V$, and $\varphi$ maps it into the two-sided ideal generated by the images of these elements. These images lie in the two-sided ideal generated by $\operatorname{CH}(D')$, $\varphi(Z)$, $\mathfrak m_{B'}^N$ and the $v'-1$ with $v'\in V'$. Indeed, $g$ maps $\operatorname{CH}(D)$ into $\operatorname{CH}(D\otimes_BB')$, and $p$ maps $\operatorname{CH}(D'\circ p)$ into $\operatorname{CH}(D')$, as both ideals are generated by the coefficients of the elements $\chi_y(y)$, $y=t_1r_1+\dots+t_mr_m$, where $\chi_y$ is the characteristic polynomial at $y$ \cite[\S1.17]{Che14}; $g(\mathfrak m_B^N)\subset\mathfrak m_{B'}^N$, as $g$ is local; and $\varphi(v-1)=p(v)-1$ with $p(v)\in V'$.

(2) By Lemma~\ref{sha-patching-chain-cayley-finite}, there is a finite set $s\subset G$ whose images span $E=E_{D,Z}$ over $B$, and for every $N$ an open normal subgroup $V_N$ of $G$ with the image of $v-1$ in $\mathfrak m_B^NE$ for $v\in V_N$. Let $Q=B'[G]/J$. The composite of $\varphi$ with $B'[G]\to Q$ vanishes on the closure of $\operatorname{CH}(D)\cup Z$, so it induces an additive map $\psi:E\to Q$ with $\psi(bx)=g(b)\psi(x)$. The $B'$-module $Q$ is spanned by the images of the group elements, which are the $\psi$-images of their images in $E$; so it is spanned by the images of the elements of $s$, and it is finite over $B'$. By Krull's intersection theorem, $\bigcap_N\mathfrak m_{B'}^NQ=0$ \cite[Theorem~8.10]{Mat86}. So it suffices that the image in $Q$ of every $z$ in the closure of $\operatorname{CH}(D')\cup Z'$ lies in $\mathfrak m_{B'}^NQ$, for every $N$. Fix $N$. Then $z$ lies in the two-sided ideal generated by $\operatorname{CH}(D')$, $Z'$, $\mathfrak m_{B'}^N$ and the $v-1$ with $v\in V_N$, and the elements of $B'[G]$ whose image in $Q$ lies in $\mathfrak m_{B'}^NQ$ form a two-sided ideal. It contains these generators. The two-sided ideal generated by $\operatorname{CH}(D')$ is generated by $\varphi(\operatorname{CH}(D))$, as the formation of the Cayley-Hamilton quotient commutes with base change \cite[(1.8)]{Che14}, and so it lies in $J$; $Z'\subset J$ by hypothesis; and for $v\in V_N$ the image of $v-1$ is $\psi$ of an element of $\mathfrak m_B^NE$, so it lies in $g(\mathfrak m_B^N)Q\subset\mathfrak m_{B'}^NQ$.
% chBarPush: the open subgroup comap p, gDet_mapRingHom_mem_ch_map_of_mem_ch and gDet_mapDomainRingHom_mem_ch_of_mem_ch_comap for the Cayley-Hamilton ideals, Ideal.map_pow with hm. closure_bc: finE (hs, hU), cbc_exists for Q, cbc_semi for psi, hsp and hfin (Module.Finite), cbc_krull, cbc_span for the two-sided ideal, g_chmap with cbc_sub for the Cayley-Hamilton generators.
\end{proof}

\begin{lemma}\label{sha-patching-chain}
In Situation~\ref{sha-hecke-situation}, let $J\subset B$ be an ideal with Taylor-Wiles data as in Proposition~\ref{sha-hecke-ideal}(7) for a family of places as in Lemma~\ref{sha-patching-places}, for which $H$ kills the $\bar\chi_i$. Suppose that $B$ is topologically generated over $\mathcal O$ by the coefficients of $D^{\mathrm u}$ and the values of $\psi^{\mathrm u}$. Let $x_0,\dots,x_h$ be normalized arcs with residue fields of the same size, each with a Selmer basis, distinct eigenvalues at the $\sigma_j$, the bounds of Lemmas~\ref{sha-patching-tangent} and~\ref{sha-patching-obstruction}, and an $\mathcal O$-algebra map $f_i:B\to\kappa_i[[T]]$ to which it is attached. Let $e_1,\dots,e_h$ be thick primes with $e_i\subset\ker f_{i-1}\cap\ker f_i$, and suppose that some stable prime $\mathfrak q_1$ with $J\subset\mathfrak q_1\subset\ker f_0$ exists. Then $J$ is contained in every stable or thick prime of $B$ contained in $\ker f_h$.
\end{lemma}
\leanref{Development/AutomorphicArithmetic/AutomorphicVarieties/PatchingModules/PatchedSupportCompletions.lean}{UpperTheory.V.Pa.Raw}
\leanref{Development/AutomorphicArithmetic/AutomorphicVarieties/PatchingModules/AutomorphicDetection/Raw.lean}{UpperTheory.V.Pa.raw}
\leanref{Development/AutomorphicArithmetic/AutomorphicVarieties/PatchingModules/SeparableMatrixStability/RawMain.lean}{UpperTheory.V.Pa.raw_main}
\leanref{Development/AutomorphicArithmetic/AutomorphicVarieties/PatchingModules/AutomorphicDetection/CoreBC.lean}{UpperTheory.V.Pa.coreBC}
\leanref{Development/AutomorphicArithmetic/AutomorphicVarieties/PatchingModules/AutomorphicDetection/LiftPt.lean}{UpperTheory.V.Pa.core}
\begin{proof}
\emph{Step 1: a pro-$\ell$ quotient.} Replace $\Gamma$ by its quotient by the kernel of the maximal pro-$\ell$ quotient of $H$. The residual datum, the universal pair, the Taylor-Wiles data and the arcs descend, since $H$ acts on them through pro-$\ell$ groups, and stability and thickness of primes are unchanged.

\emph{Step 2: base change.} Replace $\mathcal O$ by a finite extension $\mathcal O'$ whose residue field has the size of the residue fields of the arcs, and $B$ by its completed base change to $\mathcal O'$, with one further weight variable, which each arc sends to $T$. The arcs and the maps $f_i$ lift, and the primes of the chain lift to a chain of the same kind.

\emph{Step 3: patching.} After Step 1, $P(H)=1$ and $P(H_N)=1$ for all $N$; after Step 2, $B$ is topologically generated over $\mathcal O$ by the coefficients of $D^{\mathrm u}$, the values of $\psi^{\mathrm u}$ and the image $\lambda(X_0)$ of the further weight variable. Let $\Lambda=\mathcal O[[X_0,X_1,\dots,X_D]]$ be the ring of weight variables after Step 2, and $S=\Lambda[[y_{j,i}]]$, with $nq$ variables $y_{j,i}$. Let $\mathcal U$, $\tilde\Gamma$, the $\pi_N$, $\Gamma_\infty$ and the $\phi_j,\tau_j$ be as in Lemma~\ref{sha-patching-chain-limit}. The lifts $\tilde u_{j,i}\in\mathcal O$ of the residual eigenvalues in (7) lift the $\bar\chi_i(\sigma_j)$, since $p_N(\phi_{N,j})\in\sigma_jH$ and $H$ kills the $\bar\chi_i$. Let $R$, $s$, the $\rho_N$, $B[u]$, $\pi$, $\mathfrak i$ and $f$ be as in Lemmas~\ref{sha-patching-chain-ring}, \ref{sha-patching-chain-maps} and~\ref{sha-patching-chain-relations}. By Lemma~\ref{sha-patching-chain-relations}(1), $\pi$ is surjective and $\mathfrak i\subset\ker\pi$. Let $R$ act on the $M_N$ through the $\rho_N$, and on $X=B[u]\otimes_BM_0$ through $\pi$. By Lemma~\ref{sha-patching-chain-hypotheses}, $X$ is finite and free over $\Lambda$, with $\Lambda$ acting through $s$ and the $s(y_{j,i})$ acting by $0$, and hypotheses (1), (2) and (3) of Lemma~\ref{sha-patching-chain-patched} hold, with $e+1$ in place of $e$ and the generators of $\mathfrak i$ of Lemma~\ref{sha-patching-chain-relations}.

Write $x_t$, $0\le t\le h$, for the arcs, with maps $f_t:B\to\kappa_t[[T]]$. By Lemma~\ref{sha-patching-chain-points}, there are primes $P_t$ and $E_t$ of $B[u]$ over the $\ker f_t$ and the $e_t$, with $E_t\subset P_{t-1}\cap P_t$, and maps $f'_t:B[u]\to\kappa_t[[T]]$ extending $f_t$, with kernels $P_t$; the primes $\mathfrak q_t=\pi^{-1}(P_t)$ of $R$ contain $\mathfrak i$ and the $s(y_{j,i})$, and $f\notin\mathfrak q_t$. The completion of $\Lambda$ at the inverse image of $\mathfrak q_t$ is regular of dimension $1+D$. Indeed, $\pi(s(a))=\lambda(a)$ for $a\in\Lambda$ by Lemma~\ref{sha-patching-chain-maps}(4), so the inverse image of $\mathfrak q_t$ in $\Lambda$ is the kernel $\mathfrak p_t$ of the composite $g_t=f_t\circ\lambda:\Lambda\to\kappa_t[[T]]$. It is surjective, as a local map of complete local rings that maps $\mathcal O$ onto the residue field $\kappa_t$ and $X_0$ to $T$, which generates the maximal ideal of $\kappa_t[[T]]$. So $\Lambda/\mathfrak p_t\cong\kappa_t[[T]]$ is regular of dimension $1$. Since $\Lambda$ is regular of dimension $D+2$, the localization $\Lambda_{\mathfrak p_t}$ is regular of dimension $D+2-1=1+D$, and so is its completion \cite[\S\S14, 17 and 19]{Mat86}. Lemma~\ref{sha-patching-chain-patched}, with $d=1+D$ and $|\Sigma|=nq$, gives completions $B_t$ of $R/\mathfrak i$ at the $\mathfrak q_t$ and ideals $\mathfrak S_t$ whose minimal primes have dimension at least $1+D+nq$. Finally, let $P$ be a prime of $B[u]$. Since $\pi$ is surjective, the annihilator of $X$ in $R$ is contained in $\pi^{-1}(P)$ if and only if the annihilator of $X$ in $B[u]$ is contained in $P$. As $M_0$ is finite over $B$, the support of $B[u]\otimes_BM_0$ is the inverse image of the support of $M_0$, which is the zero locus of $J$. So the annihilator of $X$ in $R$ is contained in $\pi^{-1}(P)$ if and only if $J\subset P\cap B$.

\emph{Step 4: the bound at an arc.} Fix $t$, and let $Z$ and $Q$ be primes of $B_t$, with $Z$ lying over a stable prime and $\dim B_t/Z\ge1+D+nq$, and $Q$ lying over a stable prime or over a thick one, for the pair $(D_\infty,\psi_\infty)$ of Lemma~\ref{sha-patching-chain-relations}(2). Then $Z\subset Q$. Write $D_0=D+nq$, and let $a$ and $n-a$ be the block sizes of the normalized arc $x_t$. The argument has four parts.
\begin{enumerate}
\item[(a)] Let $\varpi$ be a uniformizer of $\mathcal O$. Lemma~\ref{sha-patching-chain-completion} applies to $\mathcal O$, whose residue field is $\kappa_t$ after Step 2; to $A=R/\mathfrak i$, which is complete local noetherian with that residue field by Lemma~\ref{sha-patching-chain-relations}(2); to the map $R/\mathfrak i\to\kappa_t[[T]]$ induced by $f'_t\circ\pi$, which has kernel $\mathfrak q_t/\mathfrak i$; and to the image of $z_0$, which this map sends to $f_t(\lambda(X_0))=T$ by Lemma~\ref{sha-patching-chain-maps}(4) and Step 2. So $f'_t\circ\pi$ extends to a map $\hat f_t$ from $B_t$ to the fraction field $\kappa_t((T))$ of $\kappa_t[[T]]$ with kernel $\mathfrak m_{B_t}$; and $B_t$ is a complete local noetherian algebra, with the same residue field, over the complete discrete valuation ring $\mathcal O_t$, the completion of the localization of $\mathcal O[[X]]$ at $\varpi\mathcal O[[X]]$, whose map to $\kappa_t((T))$ is surjective. The map $\mathcal O[[X]]\to\mathcal O_t\to B_t$ is the map to $R/\mathfrak i$ that sends $X$ to $z_0$, followed by $R/\mathfrak i\to B_t$.
\item[(b)] The elements $\phi_j,\tau_j$ of Lemma~\ref{sha-patching-chain-limit}(4), the images in $R/\mathfrak i$ of the $U_{j,i}$, the $z_{j,i}$ and $z_0$, and the map $f'_t\circ\pi$ form an arc datum for $x_t$ over $R/\mathfrak i$ and $(D_\infty,\psi_\infty)$, in the sense of Definition~\ref{sha-patching-chain-datum}, with $c_\infty=\pi_\infty(\tilde c)$ and the local maps of Lemma~\ref{sha-patching-chain-limit}(3). Indeed: $c_\infty^2=1$, since $\tilde c^2$ lies in every $\ker\pi_N$, as $\pi_N(\tilde c)=c$, so in $(\ker\pi_\infty)V$ for every open normal subgroup $V$ of $\tilde\Gamma$ by Lemma~\ref{sha-patching-chain-limit}(2), and $\ker\pi_\infty$ is closed; $\Gamma_\infty$ is a quotient of the topologically finitely generated group $\tilde\Gamma$, and $\mathcal O$ maps onto $\kappa_t$ after Step 2; $\phi_j$ and $\tau_j$ commute and $p_\infty(\tau_j)=1$ by Lemma~\ref{sha-patching-chain-limit}(4); the two trace relations are the images of the generators (v) and (vi) of $\mathfrak i$ of Lemma~\ref{sha-patching-chain-relations}, as $\pi_\infty$ maps $\tilde\Gamma_0$ onto $\Gamma_{\infty,0}$ and $D_R$ is $D_\infty\circ\pi_\infty$ modulo $\mathfrak i$; $R$ is topologically generated over $\mathcal O$ by the coefficients of $D_R$, the values of $\psi_R$, $z_0$, the $U_{j,i}$ and the $z_{j,i}$, by the uniqueness in Lemma~\ref{sha-patching-chain-ring}, and their images in $R/\mathfrak i$ are the generators required; $f'_t\circ\pi$ carries $(D_\infty,\psi_\infty)$ to the determinant of $\rho_t\circ p_\infty$ and the characters of $x_t$, by Lemma~\ref{sha-patching-chain-relations}(2), as $f'_t$ extends $f_t$, to which $x_t$ is attached; it sends $z_0$ to $T$, as in (a), and the $z_{j,i}$ to $0$, by Lemma~\ref{sha-patching-chain-maps}(4); and it sends $U_{j,i}$ to $f'_t(u_{j,i})=\alpha_{t,j,i}$, which by the proof of Lemma~\ref{sha-patching-chain-points}(3) are pairwise distinct roots of the characteristic polynomial of $\rho_t(\sigma_j)=\rho_t(p_\infty(\phi_j))$. Lemma~\ref{sha-patching-chain-representation}, for this arc datum, $B_t$ and $\hat f_t$, gives a commutative $B_t$-algebra $A_t$, finite as a $B_t$-module, into which $B_t$ injects, and a $B_t$-submodule $N_t$ with $A_t=B_t\cdot1\oplus N_t$, $xy\in\mathfrak m_{B_t}\cdot1$ for $x,y\in N_t$, and $N_t=0$ if $a=n$; and a continuous representation $r_t$ of $\Gamma_{\infty,0}$ over $A_t$, with ${}^tr_t(g)\,r_t(c_\infty gc_\infty^{-1})=\varepsilon(g)^{1-n}$, such that $D_\infty\otimes A_t$ is the determinant of $r_t$, the $B_t$-span of the $r_t(g)$ is the module of matrices with entries in $B_t\cdot1$ in the two diagonal blocks, of sizes $a$ and $n-a$, and in $N_t$ in the two off-diagonal blocks, and the entries of $r_t$ modulo $N_t$, mapped by $\hat f_t$, are those of $\rho_t\circ p_\infty$. Every prime of $B_t$ lying over a thick prime satisfies the hypothesis of (5) of that lemma, so it contains the annihilator of $N_t$. Indeed, let $Q$ be a prime of $B_t$ and $P$ a prime of $B[u]$ with the same inverse image in $R/\mathfrak i$, such that $P\cap B$ is thick. Let $I_B\subset B[\Gamma_0]$ be the closed two-sided ideal that defines $E_0$ in Definition~\ref{sha-hecke-arcs}(4), and $I$ the relation ideal of the arc datum. The quotient of $B[\Gamma_0]$ by $I_B$ is spanned by the images of finitely many group elements. Indeed, after Step 1, $H$ is a pro-$\ell$ group with finitely many continuous homomorphisms to $\mathbf Z/\ell$, by Lemma~\ref{sha-patching-places}, so it is topologically finitely generated, and so is $\Gamma$, in which $H$ is open; by Lemma~\ref{sha-patching-chain-cayley-bound}, the $\bar\chi_i$ have an extension bound. Let $Z_B$ be the set of the elements $\prod_i(g_i-\psi^{\mathrm u}_{v,i}(g_i))$ and $g-1$ of (2) and (3) of Situation~\ref{sha-hecke-situation} and of their images under $*$. By Lemma~\ref{sha-patching-chain-cayley-star}(2), for $\Gamma'=\Gamma$, $c'=c$ and $\Gamma'_0=\Gamma_0$, the image under $*$ of $\operatorname{CH}(D^{\mathrm u})$ lies in $\operatorname{CH}(D^{\mathrm u})$; so $I_B$ is the closure of $\operatorname{CH}(D^{\mathrm u})\cup Z_B$, and $B[\Gamma_0]/I_B=E_{D^{\mathrm u},Z_B}$. As $B$ is complete local noetherian with finite residue field $k$ and $D^{\mathrm u}$ is continuous with reduction the determinant of $\bigoplus_i\bar\chi_i$, Lemma~\ref{sha-patching-chain-cayley-finite}(1) applies. In the same way, the quotient of $(R/\mathfrak i)[\Gamma_{\infty,0}]$ by $I$ is $E_{D_\infty,S}$ for the set $S$ of Lemma~\ref{sha-patching-chain-cayley-kills}, whose (1) gives an extension bound; and $R/\mathfrak i$ is complete local noetherian with residue field $k$ by Lemma~\ref{sha-patching-chain-relations}(2), so Lemma~\ref{sha-patching-chain-cayley-finite}(1) applies again. Next, the image of $I$ in $B[u][\Gamma_0]$, under $\pi$ and $p_\infty$, lies in the two-sided ideal $J_B$ generated by the image of $I_B$. The ring $B[u]$ is complete local noetherian and finite over $B$, by Lemma~\ref{sha-patching-chain-maps}(1); the map $\pi:R/\mathfrak i\to B[u]$ is surjective, hence local, and it carries $D_\infty$ to $D^{\mathrm u}\circ p_\infty$, by Lemma~\ref{sha-patching-chain-relations}(2). Let $S'$ be the image of $S$ in $B[u][\Gamma_0]$. By Lemma~\ref{sha-patching-chain-ideals}(1), for $\pi$ and $p_\infty$, the image of $I$ lies in the closure of $\operatorname{CH}(D^{\mathrm u}\otimes_BB[u])\cup S'$; by Lemma~\ref{sha-patching-chain-ideals}(2), for $B\to B[u]$ and $Z=Z_B$, it suffices that $S'\subset J_B$. The map induced by $\pi$ and $p_\infty$ commutes with $*$, by Lemma~\ref{sha-patching-chain-cayley-star}(1), as $p_\infty(c_\infty)=c$ and $\varepsilon$ is viewed on $\Gamma_\infty$ through $p_\infty$. It carries $\prod_i(g_i-\psi_{\infty,v,i}(g_i))$, for $g_i\in D_v$, to the image of $\prod_i(g_i-\psi^{\mathrm u}_{v,i}(g_i))\in Z_B$, by Lemma~\ref{sha-patching-chain-relations}(2); and $g-1$, for $g$ in the image of $I'_u$, to the image of $p_\infty(g)-1\in Z_B$, as $p_\infty(g)$ lies in the image of $I'_u$ in $\Gamma_0$. It carries $\prod_i(\phi_j-U_{j,i})$ to $\prod_i(\sigma_j-u_{j,i})$, the image of the characteristic polynomial of $D^{\mathrm u}$ at $\sigma_j$ evaluated at $\sigma_j$, by Lemma~\ref{sha-patching-chain-maps}(1); this element of $B[\Gamma_0]$ lies in $\operatorname{CH}(D^{\mathrm u})$, and so does its image under $*$. And it carries $\prod_{k\ne i}(\phi_j-U_{j,k})(\tau_j-1-z_{j,i})$ to $0$, as $p_\infty(\tau_j)=1$ and $\pi(z_{j,i})=0$. With the images under $*$, this gives $S'\subset J_B$. As $P\cap B$ is thick, the quotient of $\bar\kappa[\Gamma_0]$ by the ideal generated by the image of $I_B$, for an algebraic closure $\bar\kappa$ of the residue field of $P\cap B$, is not semisimple. Lemma~\ref{sha-patching-chain-thick}(1), for $B\to\bar\kappa$ and $B\to B[u]\to\kappa(P)$, gives a field extension $K$ of $\kappa(P)$ over which the same holds. Lemma~\ref{sha-patching-chain-thick}(2), for $p_\infty$ and $R/\mathfrak i\to B[u]\to K$, shows that the quotient of $K[\Gamma_{\infty,0}]$ by the ideal generated by the image of $I$ is not semisimple. Finally Lemma~\ref{sha-patching-chain-thick}(1), for $R/\mathfrak i\to K$ and $R/\mathfrak i\to B_t\to\kappa(Q)$, whose kernels agree, gives a field extension of $\kappa(Q)$, with its map from $B_t$ of kernel $Q$, over which the quotient by the image of $I$ is not semisimple; this is the hypothesis of (5).
% (b): Core.ooPt builds the OoPt (fg from Lim.fg, Ip from Adapted, q1 and q2 from Rel.q1 and Rel.q2 with lam_lift, gen from rrGen with genBy_img, maps from ooMaps, f_s, f_y, f_U by inj_of_prod, eig from spRoots), then s6 (Lemma sha-patching-chain-representation) in galOut. The hypothesis of (5) at primes over thick ones is Core.nonssT, used in chainAbs: BC.nonSS_field and BC.nonSS_mono are Lemma sha-patching-chain-thick; the spanning sets hsp and hsp' come from BC.pair_mf (exists_multFree_of_phiCond_of_injective_chars and the first conclusion of CA.finE, Lemmas sha-patching-chain-cayley-bound and sha-patching-chain-cayley-finite), with PhiCond from BC.phi0 (isProL_of_PL, topGen_of_hom with T.fin0, fg_of_open, phi_fg) for Gamma_0 and from phi_fg with fg_image for Gamma_infty; the containment h2 is chBarPush with ooMaps and BC.closure_bc (Lemma sha-patching-chain-ideals), applied to gensImg (pair_kills for the star of the Cayley-Hamilton ideal, BC.star_map and star_dom, oEl_dom, spRoots with charpoly_element_eq_aeval_group_element and k_mem for q1El, Lambda.p_ta and m.pi_maps for q2El).
\item[(c)] Let $\kappa_1,\dots,\kappa_g$ be the Selmer basis of $x_t$. After Step 1, $P(H)=1$ and $P(H_N)=1$ for all $N$; so Lemma~\ref{sha-patching-chain-limit-bounds}, applied to $x_t$ with the bounds of Lemmas~\ref{sha-patching-tangent} and~\ref{sha-patching-obstruction}, whose exponent $n[L^+:\mathbf Q]+nq$ is $D_0$, and Lemma~\ref{sha-patching-chain-laurent} give the three statements of Lemma~\ref{sha-patching-chain-laurent} for $p_\infty$. The ring $\mathcal O_t$ of (a) satisfies the conditions on $\mathcal O'$ in Definition~\ref{sha-patching-chain-tangent}, for the arc datum of (b), by (a) and Lemma~\ref{sha-patching-chain-completion}(2): the images of $\mathcal O$ and of $z_0$ come from $\mathcal O_t$, and every element of $\kappa_t[[T]]$ is the image under $f'_t\circ\pi$ of some $F(c)$ with $c\in\mathcal O[[X]]$. Lemma~\ref{sha-patching-chain-presentation}, with $\mathcal O'=\mathcal O_t$ and $A_t$, $N_t$ and $r_t$ of (b), gives elements $X_1,\dots,X_{D_0}\in\mathfrak m_{B_t}$ with $\mathfrak m_{B_t}\subset(X_1,\dots,X_{D_0})+J_t$, where $J_t$ is the ideal generated by $\mathfrak m_{\mathcal O_t}$, $\mathfrak m_{B_t}^2$ and the $b\in B_t$ with $b=xy$ for some $x,y\in N_t$; and elements $m_1,\dots,m_g$ generating $N_t$, and $c_{i,i',k}\in\mathcal O_t$, for $1\le i,i'\le g$ and $1\le k\le D_0$, such that the $g^2$ vectors $(\bar c_{i,i',k})_k$ of residues are linearly independent and $\sum_{i'}\big(\sum_kc_{i,i',k}X_k\big)m_{i'}\in J_tN_t$ for each $i$.
\item[(d)] If $a=n$, then $N_t=0$ and $J_t=\mathfrak m_{\mathcal O_t}B_t+\mathfrak m_{B_t}^2$. Then the map $\mathcal O_t[[X_1,\dots,X_{D_0}]]\to B_t$ sending each variable to the element of the same name is surjective, since the image of its maximal ideal together with $\mathfrak m_{B_t}^2$ contains $\mathfrak m_{B_t}$ and $B_t$ is complete. Its source is a domain of dimension $1+D_0$. If a nonzero $r$ maps into $Z$, then $B_t/Z$ is a quotient of $\mathcal O_t[[X_1,\dots,X_{D_0}]]$ by an ideal containing the nonzerodivisor $r$, so $\dim B_t/Z\le D_0$, a contradiction. So only $0$ maps into $Z$, and as the map is surjective, $Z=0\subset Q$. Let now $a<n$, and let $I_t\subset B_t$ be the ideal generated by the $b\in B_t$ with $b\cdot1=xy$ for some $x,y\in N_t$. As $Z$ lies over a stable prime, $\det\big(\operatorname{tr}_{D_\infty}(g_sh_{s'})\big)_{s,s'}\notin Z$ for some $g_1,\dots,g_{n^2},h_1,\dots,h_{n^2}\in\Gamma_{\infty,0}$, by Definition~\ref{sha-hecke-arcs}(4). Lemma~\ref{sha-patching-chain-stable}, applied to $B_t$, $A_t$, $N_t$, $D_\infty\otimes B_t$ and $r_t$ of (b), with $a<n$, gives $y,y'\in N_t$ and $d\notin Z$ with $yy'=d\cdot1$; so $I_t\not\subset Z$. In the same way $I_t\not\subset Q$ if $Q$ lies over a stable prime; and if $Q$ lies over a thick prime, it contains the annihilator of $N_t$ by (b). By (a), $B_t$ is a complete local noetherian $\mathcal O_t$-algebra such that $\mathcal O_t$ maps onto its residue field, and $J_t$ is the ideal $J$ of Lemma~\ref{sha-patching-chain-cotangent} for $\mathcal O_t$, $B_t$, $A_t$ and $N_t$, with $I_N=I_t$. So that lemma, with the elements of (c) and $\dim B_t/Z\ge1+D_0$, gives $Z\subset Q$.
\end{enumerate}
By (a) of Lemma~\ref{sha-patching-chain-patched}, Step 4 applies to every minimal prime $Z$ of $\mathfrak S_t$ over a stable prime.

\emph{Step 5: the chain.} The support of $M_0$ is the zero locus of $J$. Starting from $\mathfrak q_1$, Step 4 at $x_0$ puts every stable or thick prime contained in $\ker f_0$ in the support, in particular $e_1$; Step 4 at $x_1$, through $e_1$, does the same for $\ker f_1$; and so on up to $\ker f_h$.
% Pa.core. Step 3: s4 (Lemma sha-patching-chain-limit), rrEx (sha-patching-chain-ring), rmapsEx and spoCo (sha-patching-chain-maps), piKer, piSurj, ooPair and ooMaps (sha-patching-chain-relations), patchedModule_free_finite_over_iwasawaAlgebra, freeN, eventually_exists_patchedModule_quotient_maps and t3Of (sha-patching-chain-hypotheses), BC.chain_lift, lift_pt and ooPt (sha-patching-chain-points), laDim (dimension 1 + D; Lambda = abelian₂ O Unit D, with laDim_surj, mvPowerSeries_cotInd, regular_localization_of_regular_quotient and IsCplAt.regular), s5 (Lemma sha-patching-chain-patched, cited), and patchedModule_annihilator_le_comap_iff_ideal_le (ann_baseChange_le_iff from mem_support_baseChange, with A.b.supp). The generation hypothesis after Step 2 is BC.gen_bc in coreBC.
% Step 4 is galOut with D0 = D + nq: rfEx and ooEx for (a) (Lemma sha-patching-chain-completion; the fields rF, rF_f, rF_ker, OO, oo_A, oo_O, ordinaryPointSection and oo_surj of Ctx; the base ring O surjects onto x.K by (hx i).1 and P.f s = x.t by ooPt); s6 (URep, from ClosureConjugation: hat_ex, rF_surj, hat_st, hat_range, idem, gr_E, gr_I, urep_of) for (b), in Lemmas sha-patching-chain-cayley to sha-patching-chain-representation, with the fields fin, inj, N, compl, mul_mem, typeI, vr, span (blkMat n x.a 1 N), det, form, red, kills (the set P.ker), thick (NonSS P.ker), and P, vr_mem, cont for the continuity; (d) for a < n uses fin, inj, compl, mul_mem, span, det and thick; cohLim (K1, K2, K3: firstCohomologyBound_of_boundaryBound, secondCohomologyBound_of_boundaryBound, koszulThree_of_det, from T.B1, T.B2 and T.Eig at the limit level) and pres (Ctx.Pres, with J2 and ired; from tpDual, alSel, liftsSub, claim1, linAlg and dualEx) for (c), in Lemmas sha-patching-chain-limit-bounds to sha-patching-chain-presentation (k1F, k2F and k3F pass from K1, K2 and K3 to the Laurent forms); ctx_out splits on x.a = n: presOut0 (written in (d), with co_ev, co_surj, ringKrullDim_mvPowerSeries and ringKrullDim_succ_le_of_surjective) with U.typeI, and presOut1 (Lemma sha-patching-chain-cotangent) with discImp (urep_stable, Lemma sha-patching-chain-stable; Z and Q enter as DiscSt of u.D mapped to Bh, which is the stability of Definition sha-hecke-arcs(4)) and U.thick for a < n. Step 5 is chainAbs, with hdown (going down in B[u]), hSt (DiscSt in R/ii and in B through ooMaps and BC.discSt_map) and nonssT for the thick primes.
\end{proof}

\begin{remark}\label{sha-patching-chain-burnside}
In Situation~\ref{sha-hecke-situation}, let $\mathfrak q$ be a prime of $B$ and $\bar\kappa$ an algebraic closure of its residue field. Then $\mathfrak q$ is stable, in the sense of Definition~\ref{sha-hecke-arcs}(4), if and only if $D^{\mathrm u}\bmod\mathfrak q$, over $\bar\kappa$, is the determinant of an irreducible representation $\rho:\Gamma_0\to GL_n(\bar\kappa)$. For such a $\rho$ the traces of $D^{\mathrm u}$ map to those of $\rho$. If $\rho$ is irreducible, its images span $M_n(\bar\kappa)$ by Burnside's theorem, so they contain a basis $\rho(g_1),\dots,\rho(g_{n^2})$, and $\det\big(\operatorname{tr}(\rho(g_s)\rho(g_{s'}))\big)\ne0$ as the trace form on matrices is nondegenerate; so $\det\big(\operatorname{tr}_{D^{\mathrm u}}(g_sg_{s'})\big)\notin\mathfrak q$. Conversely, $D^{\mathrm u}\bmod\mathfrak q$ over $\bar\kappa$ is the determinant of a semisimple representation $\rho$ \cite{Che14}. If $\det\big(\operatorname{tr}_{D^{\mathrm u}}(g_sh_{s'})\big)\notin\mathfrak q$, a relation $\sum_sc_s\rho(g_s)=0$ gives $\sum_sc_s\operatorname{tr}(\rho(g_s)\rho(h_{s'}))=0$ for all $s'$, so all $c_s$ vanish; then the $n^2$ matrices $\rho(g_s)$ span $M_n(\bar\kappa)$, and $\rho$ is irreducible.
% Det.stableIff (B.IsStable q iff DiscSt B.u.D q) from Det.discStIff: groupDet_exists_semisimple_representation for the semisimple rho, discSt_iff with span_of_gram and CA.g_gram, irr_iff_span (AlgebraTools.irred_iff_span, Burnside). IsStable is the Lean's name for the irreducible form, used in the statements P1core and P3core of Proposition sha-patching-bound.
\end{remark}
\leanref{Development/AutomorphicArithmetic/AutomorphicVarieties/WeilGroupCharacters/UniversalOrdinaryDeterminantRing.lean}{UpperTheory.V.W.UniversalOrdinaryDeterminantRing.IsStable}
\leanref{Development/AutomorphicArithmetic/AutomorphicVarieties/PatchingModules/DeterminantPadicFree/StableIff.lean}{UpperTheory.V.Pa.Det.stableIff}
\leanref{Development/AutomorphicArithmetic/AutomorphicVarieties/PatchingModules/DeterminantPadicFree/DiscStIff.lean}{UpperTheory.V.Pa.Det.discStIff}

\begin{proposition}\label{sha-patching-bound}
In Situation~\ref{sha-hecke-situation}, let $J\subset B$ be an ideal with Taylor-Wiles data as in Proposition~\ref{sha-hecke-ideal}(7).
\begin{enumerate}
\item Let $f:B\to\kappa[[T]]$ be attached to a nice arc of type I, with $J\subset\ker f$. Then $J$ is contained in every prime of $B$ contained in $\ker f$.
\item Let $f_0,\dots,f_h$ be attached to nice arcs of type E, and $e_1,\dots,e_h$ thick primes with $e_i\subset\ker f_{i-1}\cap\ker f_i$. If some stable prime $\mathfrak q_1$ with $J\subset\mathfrak q_1\subset\ker f_0$ exists, then $J$ is contained in every stable or thick prime of $B$ contained in $\ker f_h$.
\end{enumerate}
\end{proposition}
\leanref{Development/AutomorphicArithmetic/AutomorphicVarieties/PatchingModules/TaylorWilesObstructionBound.lean}{UpperTheory.V.Pa.PfromAut}
\leanref{Development/AutomorphicArithmetic/AutomorphicVarieties/PatchingModules/AutomorphicDetection/PfromAut.lean}{UpperTheory.V.Pa.pFromAut}
\leanref{Development/AutomorphicArithmetic/AutomorphicVarieties/WeilGroupCharacters/UniversalOrdinaryDeterminantRing.lean}{UpperTheory.V.W.UniversalOrdinaryDeterminantRing.P1core}
\leanref{Development/AutomorphicArithmetic/AutomorphicVarieties/WeilGroupCharacters/UniversalOrdinaryDeterminantRing.lean}{UpperTheory.V.W.UniversalOrdinaryDeterminantRing.P3core}
\begin{proof}
Both statements follow from one about chains: let $f_0,\dots,f_h$ be attached to nice arcs, of type I or E, and $e_1,\dots,e_h$ and $\mathfrak q_1$ as in (2); then $J$ is contained in every stable or thick prime contained in $\ker f_h$, and if $f_i$ is attached to an arc of type I, every prime contained in $\ker f_i$ is stable. Statement (2) is the case of type E. For (1), take $h=0$: $\ker f$ is stable, so $\mathfrak q_1=\ker f$ serves, and every prime contained in $\ker f$ is stable.

For the chain statement: let $f$ be attached to an arc of type I, and let $\rho'$ over $\kappa'[[T]]$ be its normalization by Lemma~\ref{sha-patching-normalization}, with $a=n$, attached to $j\circ f$ for the injective map $j$ there. For every matrix $Z$, some $T^sZ$ lies in the $\kappa'[[T]]$-span of $\rho'(\Gamma')$; so $\rho'(\Gamma')$ spans the matrices over the fraction field of $\kappa'[[T]]$. As the trace form on matrices is nondegenerate, there are $g_1,\dots,g_{n^2}\in\Gamma'$ with $\det\big(\operatorname{tr}\rho'(g_sg_{s'})\big)\ne0$. This determinant is the image under $j\circ f$ of $\det\big(\operatorname{tr}_{D^{\mathrm u}}(g_sg_{s'})\big)$, which therefore lies in no prime $\mathfrak q\subset\ker f$; so every such prime is stable. Next, normalize each arc by Lemma~\ref{sha-patching-normalization}, with $I'_u$ the subgroups of $I_u$ on which the $\bar\chi_i$ are trivial; the kernels of the $f_i$ do not change, since the maps $j$ are injective. Since $L$ has no primitive $\ell$-th root of unity, $\varepsilon\bmod\ell$ is nontrivial on $\Gamma_0$. Lemma~\ref{sha-patching-places} gives the finiteness hypothesis of Lemma~\ref{sha-patching-frobenius}, which gives $\sigma_1,\dots,\sigma_q$, the extensions $x'_i$, their Selmer bases and eigenvalues; Lemma~\ref{sha-patching-places} then gives the places $w_{N,j}$ and the group $H$. Each $\rho'_i$ is trivial on the kernel of the maximal pro-$\ell$ quotient of $H$, because $\rho'_i$ is residually trivial on $H$, so its image of $H$ is a pro-$\ell$ group. Lemmas~\ref{sha-patching-tangent} and~\ref{sha-patching-obstruction} give the bounds at each $x'_i$, and the hypothesis on $J$ gives its Taylor-Wiles data for these places. The generation hypothesis holds, since $B$ is the universal ring of ordinary pairs. Lemma~\ref{sha-patching-chain} gives the chain statement, after identifying its thick primes with those of Definition~\ref{sha-hecke-arcs}.
\end{proof}

\subsection{Descent data and links}\label{subsection-sha-descent-data}

\begin{definition}\label{sha-descent-datum}
Let $\ell$ be an odd prime, let $\mathcal O$ be a local ring whose residue field $\mathbf F$ has characteristic $\ell$, and write $G_{\mathbf Q}=\operatorname{Gal}(\bar{\mathbf Q}/\mathbf Q)$.
For a place $A$ of $\bar{\mathbf Q}$ let $D_A\supset I_A$ be its decomposition and inertia groups in $G_{\mathbf Q}$, and for $n\ge1$ let $\mu_n\subset\bar{\mathbf Q}$ be the group of $n$-th roots of unity.
Let $\eta:G_{\mathbf Q}\to\{\pm1\}$ and $\Psi:G_{\mathbf Q}\to\mathcal O^\times$ be homomorphisms, let $\bar\Psi:G_{\mathbf Q}\to\mathbf F^\times$ be the reduction of $\Psi$ and $\omega:G_{\mathbf Q}\to\mathbf F^\times$ the reduction of $\chi_{\mathrm{cyc}}$ modulo $\ell$, let $c\in G_{\mathbf Q}$ be a complex conjugation, $\mathfrak B$ a place of $\bar{\mathbf Q}$ above $\ell$, $M\ge1$ an integer and $H_b\subset G_{\mathbf Q}$ a subgroup.
The tuple $(\eta,\Psi,c,\mathfrak B,M,H_b)$ is a \emph{descent datum} if:
\begin{enumerate}
\item $\mathbf F$ is finite, and $2\ell$ divides $M$;
\item there are an integer $D\ge1$ whose prime factors divide $M$ and $\delta\in\bar{\mathbf Q}$ with $\delta^2=-D$ and $g(\delta)=\eta(g)\delta$ for all $g\in G_{\mathbf Q}$;
\item there is $a\ge0$ such that $\eta$ and $\bar\Psi$ are trivial on the elements of $G_{\mathbf Q}$ that fix $\mu_{M^a}$;
\item $\eta$ is trivial on $I_A$ for every place $A$ above $\ell$, $\bar\Psi(c)=1$, and there is $j\ge0$ with $\bar\Psi=\omega^j$ on these inertia groups;
\item $\ker\omega\cap\ker\bar\Psi$ is not contained in $\ker\eta$;
\item there is $s\in D_{\mathfrak B}$ with $\eta(s)=1$ and $\bar\Psi(s)^2\ne1$;
\item $H_b$ is open, contains the commutator subgroup of $G_{\mathbf Q}$, has index a power of $\ell$ and contains $I_A$ for every place $A$ above $\ell$; and $\Psi(s)=1$ for every $s\in H_b$ with $\bar\Psi(s)=1$ that lies in $I_A$ for a place $A$ not above $\ell$.
\end{enumerate}
\end{definition}
\leanref{Development/AutomorphicArithmetic/AutomorphicVarieties/FrobeniusSmoothVectors/FrobeniusSplitPrimeDatum.lean}{UpperTheory.V.Fr.Gen}

\begin{lemma}\label{sha-descent-auxiliary}
Let $(\eta,\Psi,c,\mathfrak B,M,H_b)$ be a descent datum, with the notation of Definition~\ref{sha-descent-datum}.
\begin{enumerate}
\item There is an open subgroup $H_0\subset G_{\mathbf Q}$ of index $2$ that contains every complex conjugation, such that $D_A\cap H_0=D_A\cap\ker\eta$ for every place $A$ above a prime dividing $M$, and $I_A\not\subset H_0$ for some place $A$ above a prime not dividing $M$.
\item Let $w$ be a prime with $w\equiv1\pmod\ell$, let $\chi_w:G_{\mathbf Q}\to(\mathbf Z/w)^\times$ give the action on $\mu_w$, and let $H_{\mathrm{th}}(w)$ be the inverse image under $\chi_w$ of the group of $\ell$-th powers. Then $H_{\mathrm{th}}(w)$ is open, contains the commutator subgroup and has index $\ell$; it contains $I_A$ for every place $A$ not above $w$, and does not contain $I_A$ for the places $A$ above $w$.
\item For $m\ge0$ let $H_{\mathrm{lay}}(m)$ be the subgroup of the $g\in G_{\mathbf Q}$ with $\chi_{\mathrm{cyc}}(g)^{\ell-1}\equiv1\pmod{\ell^{m+1}}$. It is open, contains the commutator subgroup and has index $\ell^m$; its intersection with $\ker\omega$ is the kernel of $\chi_{\mathrm{cyc}}\bmod\ell^{m+1}$; it contains $I_A$ for every place $A$ not above $\ell$; and for every place $A$ above $\ell$, $I_A\cap H_{\mathrm{lay}}(m)$ has index $\ell^m$ in $I_A$, and $\chi_{\mathrm{cyc}}\bmod\ell^r$ maps $I_A$ onto $(\mathbf Z/\ell^r)^\times$ for every $r\ge0$.
\item For every $Q\ge1$ there is a prime $w\equiv1\pmod Q$ that does not divide $M$, with $w\equiv1\pmod\ell$ and $D_A\subset\ker\eta\cap\ker\bar\Psi\cap\ker\omega$ for every place $A$ above $w$. Such a prime is an \emph{auxiliary prime} of the datum.
\item For an open subgroup $V\subset G_{\mathbf Q}$ let $r_\ell(V)$ be the dimension over $\mathbf F_\ell$ of the space of continuous homomorphisms $V\to\mathbf Z/\ell$ that vanish on $V\cap I_A$ for every place $A$. Let $H_0$ be as in (1), let $w$ be an auxiliary prime, and put $U_m=H_0\cap H_{\mathrm{th}}(w)\cap H_b\cap H_{\mathrm{lay}}(m)$. Then there is $C_0\ge0$ with $r_\ell(U_m\cap\ker\eta\cap\ker\bar\Psi\cap\ker\omega)\le C_0$ for every $m\ge0$.
\end{enumerate}
\end{lemma}
\leanref{Development/AutomorphicArithmetic/AutomorphicVarieties/FrobeniusSmoothVectors/EtaleSimplicialCochainSchemeSectionExtensionNatu.lean}{UpperTheory.V.Fr.exists_twistCharacter_continuousCohomology}
\leanref{Development/AutomorphicArithmetic/AutomorphicVarieties/FrobeniusSmoothVectors/FrobeniusSplitPrimeDatum.lean}{UpperTheory.V.Fr.FrobeniusSplitPrimeDatum}
\leanref{Development/AutomorphicArithmetic/AutomorphicVarieties/FrobeniusSmoothVectors/FrobeniusSplitPrimeDatum.lean}{UpperTheory.V.Fr.Hth}
\leanref{Development/AutomorphicArithmetic/AutomorphicVarieties/FrobeniusSmoothVectors/HthPlace.lean}{UpperTheory.V.Fr.hth}
\leanref{Development/AutomorphicArithmetic/AutomorphicVarieties/FrobeniusSmoothVectors/FrobeniusSplitPrimeDatum.lean}{UpperTheory.V.Fr.Hlay}
\leanref{Development/AutomorphicArithmetic/AutomorphicVarieties/FrobeniusSmoothVectors/HlayPlace.lean}{UpperTheory.V.Fr.hlay}
\leanref{Development/AutomorphicArithmetic/AutomorphicVarieties/FrobeniusSmoothVectors/AuxwEx.lean}{UpperTheory.V.Fr.auxw_ex}
\leanref{Development/AutomorphicArithmetic/AutomorphicVarieties/FrobeniusSmoothVectors/FrobeniusSplitPrimeDatum.lean}{UpperTheory.V.Fr.FrobeniusWitnessPrime}
\leanref{Development/AutomorphicArithmetic/AutomorphicVarieties/FrobeniusSmoothVectors/RankLeOfCardLe.lean}{UpperTheory.V.Fr.b0}
\leanref{Development/AutomorphicArithmetic/AutomorphicVarieties/GaloisCochains/GLinkQuot.lean}{UpperTheory.V.GC.clRk}
\begin{proof}
(1) Let $D$ and $\delta$ be as in (b).
By Dirichlet's theorem there is a prime $p>4D+M$ with $p\equiv-1\pmod{4M}$; write $-p=1+4Mt$.
Since $p\equiv3\pmod4$, a quadratic Gauss sum gives $x\in\bar{\mathbf Q}$ with $x^2=-p$ and an element $s$ of $I_A$, for a place $A$ above $p$, with $s(x)=-x$.
Since $4D$ is prime to $p$, $s(\delta)=\delta$.
Put $z=\delta x$, so that $z^2=Dp>0$ and $s(z)=-z\ne z$.
Let $H_0$ be the stabilizer of $z$: it is open, and it has index $2$ because every $g$ maps $z$ to $\pm z$ and $s\notin H_0$.
A complex conjugation fixes $z$, which is real because $z^2>0$.
Let $A'$ be a place above a prime dividing $M$ and $g\in D_{A'}$.
Then $x^2=1+4k$ with $k=Mt$ in the maximal ideal of $A'$, and $g(x)=x$: otherwise $y=(1+x)/2$ has $g(y)=1-y$, so $y$ and $1-y$ have the same valuation, and their product $-k$ has valuation less than $1$; so both have valuation less than $1$, which contradicts $y+(1-y)=1$.
Hence $g(z)=\eta(g)z$, and $g\in H_0$ if and only if $\eta(g)=1$.
Finally $p$ does not divide $M$, since $p>M$.

(2) An element $g$ lies in $H_{\mathrm{th}}(w)$ if and only if $\chi_w(g)$ is an $\ell$-th power.
The map $\chi_w$ is surjective, already on $I_A$ for a place $A$ above $w$, and $(\mathbf Z/w)^\times$ is cyclic of order $w-1$, which $\ell$ divides; so the $\ell$-th powers have index $\ell$, and $I_A\not\subset H_{\mathrm{th}}(w)$ for $A$ above $w$.
$H_{\mathrm{th}}(w)$ contains $\ker\chi_w$, which is open and contains the commutator subgroup.
For $A$ not above $w$, $I_A$ acts trivially on $\mu_w$.

(3) For $A$ above $\ell$, $\chi_{\mathrm{cyc}}$ maps $I_A$ onto $\mathbf Z_\ell^\times$, so $\chi_{\mathrm{cyc}}\bmod\ell^r$ maps $I_A$ onto $(\mathbf Z/\ell^r)^\times$.
$H_{\mathrm{lay}}(m)$ is the inverse image under $\chi_{\mathrm{cyc}}\bmod\ell^{m+1}$ of the kernel of $x\mapsto x^{\ell-1}$ on $(\mathbf Z/\ell^{m+1})^\times\cong\mathbf Z/(\ell-1)\times\mathbf Z/\ell^m$, whose image has order $\ell^m$; this gives both indices.
It contains the kernel of $\chi_{\mathrm{cyc}}\bmod\ell^{m+1}$, which is open and contains the commutator subgroup.
If $x\equiv1\pmod\ell$ and $x^{\ell-1}\equiv1\pmod{\ell^{m+1}}$, then $x\equiv1\pmod{\ell^{m+1}}$, since $(1+\ell\mathbf Z)/(1+\ell^{m+1}\mathbf Z)$ is an $\ell$-group; this gives the intersection with $\ker\omega$.
For $A$ not above $\ell$, $\chi_{\mathrm{cyc}}$ is trivial on $I_A$.

(4) Let $a$ be as in (c).
By Dirichlet's theorem there is a prime $w>M^{a+1}Q$ with $w\equiv1\pmod{M^{a+1}Q}$.
Then $w$ does not divide $M$, and $w\equiv1$ modulo $\ell$ and modulo $Q$, as $\ell$ divides $M$.
Let $A$ lie above $w$ and $s\in D_A$.
For every $m$ dividing $M^{a+1}Q$, $m$ is prime to $w$, so $s$ acts on $\mu_m$ by $\zeta\mapsto\zeta^{w^i}$ for some $i$, which is the identity since $w\equiv1\pmod m$.
So $s$ fixes $\mu_{M^a}$ and $\mu_\ell$, and $\eta(s)=\bar\Psi(s)=\omega(s)=1$ by (c).

(5) Let $a$ be as in (c), and $V=H_0\cap H_{\mathrm{th}}(w)\cap H_b\cap\ker\eta\cap\ker\bar\Psi\cap\ker\omega$.
By (c), $\ker\eta$ and $\ker\bar\Psi$ contain the open subgroup fixing $\mu_{M^a}$, so $V$ is open.
It contains the commutator subgroup: $H_0$ has index $2$, $H_{\mathrm{th}}(w)$ by (2), $H_b$ by (g), and the kernels of characters do.
So its fixed field $E$ is an abelian number field, which contains $\mu_\ell$ because $V\subset\ker\omega$.
By Theorem~\ref{sha-ferrero-washington} there is $r$ with $|\operatorname{Cl}(K')/\operatorname{Cl}(K')^\ell|\le\ell^r$ for every $n$ and $K'=E(\mu_{\ell^n})$.
The fixed field of $V\cap\ker(\chi_{\mathrm{cyc}}\bmod\ell^n)$ is $E(\mu_{\ell^n})$, and by class field theory the continuous homomorphisms $\operatorname{Gal}(\bar{\mathbf Q}/K')\to\mathbf Z/\ell$ that vanish on all inertia groups factor injectively through $\operatorname{Cl}(K')/\operatorname{Cl}(K')^\ell$.
So $r_\ell(V\cap\ker(\chi_{\mathrm{cyc}}\bmod\ell^n))\le r$ for every $n$.
By (3), $U_m\cap\ker\eta\cap\ker\bar\Psi\cap\ker\omega=V\cap\ker(\chi_{\mathrm{cyc}}\bmod\ell^{m+1})$, and $C_0=r$ will do.
\end{proof}

\begin{lemma}\label{sha-descent-subgroup}
Let $(\eta,\Psi,c,\mathfrak B,M,H_b)$ be a descent datum, with the notation of Definition~\ref{sha-descent-datum}.
Let $H_0$, and $H_{\mathrm{th}}(w)$ for an auxiliary prime $w$ and $H_{\mathrm{lay}}(m)$ for an integer $m\ge0$, be as in Lemma~\ref{sha-descent-auxiliary}, and put $U=H_0\cap H_{\mathrm{th}}(w)\cap H_b\cap H_{\mathrm{lay}}(m)$. Then:
\begin{enumerate}
\item $U$ is open, contains the commutator subgroup of $G_{\mathbf Q}$, and contains every complex conjugation;
\item $I_A\cap U\subset\ker\eta$ for every place $A$, and $D_A\cap U\subset\ker\eta$ for every place $A$ above a prime dividing $M$;
\item $\Psi(s)=1$ for every $s\in U$ with $\bar\Psi(s)=1$ that lies in $I_A$ for a place $A$ not above $\ell$;
\item for every place $A$ above $\ell$, some $s\in I_A\cap U$ has $\omega(s)\ne1$, the index of $D_A\cap U$ in $D_A$ is at least $\ell^m$, and some $s\in D_A\cap U$ has $\bar\Psi(s)^2\ne1$;
\item for $j\in\{\pm1,\pm2\}$ the character $\omega\bar\Psi^j$ is nontrivial on $U\cap\ker\eta$, and there is no integer $j$ with $\omega=\bar\Psi^j$ on $U\cap\ker\eta$;
\item the fixed field $L$ of $U\cap\ker\eta$ is an abelian CM field, and the fixed field of $U$ is its maximal totally real subfield $L^+$; so $\operatorname{Gal}(\bar{\mathbf Q}/L^+)=U$ and $\operatorname{Gal}(\bar{\mathbf Q}/L)=U\cap\ker\eta$.
\end{enumerate}
\end{lemma}
\leanref{Development/AutomorphicArithmetic/AutomorphicVarieties/FrobeniusSmoothVectors/Ufacts.lean}{UpperTheory.V.Fr.ufacts}
\leanref{Development/AutomorphicArithmetic/AutomorphicVarieties/FrobeniusSmoothVectors/FrobeniusSplitPrimeDatum.lean}{UpperTheory.V.Fr.UFacts}
\leanref{Development/AutomorphicArithmetic/AutomorphicVarieties/FrobeniusSmoothVectors/FrobeniusSplitPrimeDatum.lean}{UpperTheory.V.Fr.Upl}
\leanref{Development/AutomorphicArithmetic/AutomorphicVarieties/FrobeniusSmoothVectors/RootOfUnityEqOneOfCloseResidue.lean}{UpperTheory.V.Fr.l15}
\leanref{Development/AutomorphicArithmetic/AutomorphicVarieties/FrobeniusSmoothVectors/FieldEx.lean}{UpperTheory.V.Fr.field_ex}
\begin{proof}
Let $a$ be as in Definition~\ref{sha-descent-datum}(c), and let $\ell^b$ be the index of $H_b$.

(1) Each of the four subgroups is open and contains the commutator subgroup, by Lemma~\ref{sha-descent-auxiliary} and (g).
Let $\tau$ be a complex conjugation; $\tau\in H_0$ by Lemma~\ref{sha-descent-auxiliary}(1).
If $H$ contains the commutator subgroup and has index $\ell^e$, then $\tau^{\ell^e}\in H$, and $\tau^{\ell^e}=\tau$ because $\tau^2=1$ and $\ell$ is odd; this applies to $H_{\mathrm{th}}(w)$, $H_b$ and $H_{\mathrm{lay}}(m)$.

(2) If $A$ lies above a prime dividing $M$, then $D_A\cap U\subset D_A\cap H_0=D_A\cap\ker\eta$ by Lemma~\ref{sha-descent-auxiliary}(1).
Otherwise $I_A$ acts trivially on $\mu_{M^a}$, and $\eta$ is trivial on $I_A$ by (c).

(3) This is (g), as $U\subset H_b$.

(4) Let $A$ lie above $\ell$.
By Lemma~\ref{sha-descent-auxiliary}(3) there is $s\in I_A$ with $\chi_{\mathrm{cyc}}(s)\equiv-1\pmod{\ell^{m+1}}$.
Then $s\in H_{\mathrm{lay}}(m)$ because $\ell-1$ is even; $\eta(s)=1$ by (d), so $s\in H_0$ because $\ell$ divides $M$; $s\in H_{\mathrm{th}}(w)$ because $A$ does not lie above $w\ne\ell$; and $s\in H_b$ by (g).
If $\omega(s)=1$, then $s$ would lie in $H_{\mathrm{lay}}(m)\cap\ker\omega$, the kernel of $\chi_{\mathrm{cyc}}\bmod\ell^{m+1}$, which is impossible since $\ell^{m+1}>2$.
Next, the index of $D_A\cap U$ in $D_A$ is at least that of $I_A\cap H_{\mathrm{lay}}(m)$ in $I_A$, which is $\ell^m$.
Finally let $s\in D_{\mathfrak B}$ be as in (f), and $g\in G_{\mathbf Q}$ with $g\mathfrak B=A$.
Then $t=gsg^{-1}\in D_A$ has $\eta(t)=1$ and $\bar\Psi(t)=\bar\Psi(s)$, so $t\in H_0$.
For $N=b+m+1$, $t^{\ell^N}$ lies in $H_{\mathrm{th}}(w)$, $H_b$ and $H_{\mathrm{lay}}(m)$, which contain the commutator subgroup and have index dividing $\ell^N$; so $t^{\ell^N}\in D_A\cap U$.
And $\bar\Psi(t^{\ell^N})^2=(\bar\Psi(s)^2)^{\ell^N}\ne1$, because $x\mapsto x^\ell$ is injective on $\mathbf F^\times$.

(5) There is $s'$ with $\eta(s')=1$, $\omega(s')=-1$ and $\bar\Psi(s')=1$: if $\eta(c)=1$, take $s'=c$, as $\omega(c)=-1$ and $\bar\Psi(c)=1$ by (d); otherwise take $t\in\ker\omega\cap\ker\bar\Psi$ with $\eta(t)=-1$, by (e), and $s'=tc$.
Next, let $\kappa=\omega^i\bar\Psi^j$ for integers $i,j$, and suppose that $\kappa$ is trivial on $U\cap\ker\eta$.
Then $\kappa$ is trivial on $\ker\eta$.
Indeed, Lemma~\ref{sha-descent-auxiliary}(1) gives a place $A$ above a prime not dividing $M$ and $t\in I_A$ with $t\notin H_0$.
Then $t$ fixes $\mu_{M^a}$, so $\eta(t)=\bar\Psi(t)=1$ by (c), and $\omega(t)=1$ since $A$ does not lie above $\ell$; so $\kappa(t)=1$.
For $s\in H_0\cap\ker\eta$ and $N=1+b+m$, $s^{\ell^N}\in U\cap\ker\eta$, as in (4); so $\kappa(s)^{\ell^N}=1$ and $\kappa(s)=1$.
For $s\in\ker\eta$ outside $H_0$, $st\in H_0\cap\ker\eta$, so $\kappa(s)=\kappa(st)=1$.
Applied to $\omega\bar\Psi^j$, and to $\omega\bar\Psi^{-j}$ if $\omega=\bar\Psi^j$ on $U\cap\ker\eta$, this contradicts $\omega(s')\bar\Psi(s')^{\pm j}=-1\ne1$.

(6) By (b), $\eta(g)=1$ if and only if $g(\delta)=\delta$; so $\ker\eta$ is open and $\delta\in L$.
Since $U\cap\ker\eta$ is open and contains the commutator subgroup, $L$ is an abelian number field.
It has no real embedding, since $\delta^2=-D<0$; so it is a totally complex abelian field, hence a CM field.
A complex conjugation $c$ lies in $U$ by (1), and $\eta(c)=-1$ since $\delta$ is not real.
So $U$ is generated by $U\cap\ker\eta$ and $c$, and its fixed field is the subfield of $L$ fixed by complex conjugation, which is $L^+$.
\end{proof}

\begin{definition}\label{sha-descent-link}
Let $(\eta,\Psi,c,\mathfrak B,M,H_b)$ be a descent datum, with the notation of Definition~\ref{sha-descent-datum}, let $g_1\in G_{\mathbf Q}$, and let $L\subset\bar{\mathbf Q}$ be an abelian CM field with maximal totally real subfield $L^+$. View $\operatorname{Gal}(\bar{\mathbf Q}/L^+)$ as a subgroup of $G_{\mathbf Q}$.
\begin{enumerate}
\item A \emph{link} of the datum to $L$ with base element $g_1$ consists of the following data.
(i) A profinite group $\Gamma$ with a continuous surjective homomorphism $p:\operatorname{Gal}(\bar{\mathbf Q}/L^+)\to\Gamma$ whose kernel is the closed normal subgroup generated by the elements $s$ of $\operatorname{Gal}(\bar{\mathbf Q}/L^+)$ that lie in $I_A$ for a place $A$ not above $\ell$ and have $\eta(s)=\bar\Psi(s)=1$.
(ii) Homomorphisms $\delta_\Gamma:\Gamma\to\{\pm1\}$, $\varepsilon_\Gamma:\Gamma\to\mathbf Z_\ell^\times$ and $\bar\Psi_\Gamma:\Gamma\to\mathbf F^\times$ with $\delta_\Gamma\circ p=\eta$, $\varepsilon_\Gamma\circ p=\chi_{\mathrm{cyc}}$ and $\bar\Psi_\Gamma\circ p=\bar\Psi$ on $\operatorname{Gal}(\bar{\mathbf Q}/L^+)$; let $\omega_\Gamma$ be the reduction of $\varepsilon_\Gamma$ in $\mathbf F^\times$, and $\Gamma_0=\ker\delta_\Gamma$.
(iii) The elements $c_\Gamma=p(c)$ and $g_0=p(g_1)$, where $c$ and $g_1$ lie in $\operatorname{Gal}(\bar{\mathbf Q}/L^+)$ and $g_0\in\Gamma_0$.
(iv) For each place $v$ of $L^+$ above $\ell$, a place $A_v=\gamma_v\mathfrak B$ of $\bar{\mathbf Q}$ above $v$ with $\eta(\gamma_v)=1$, such that every place of $\bar{\mathbf Q}$ above $\ell$ is $uA_v$ for a unique $v$ and some $u\in\operatorname{Gal}(\bar{\mathbf Q}/L^+)$; the decomposition group $G_v=D_{A_v}\cap\operatorname{Gal}(\bar{\mathbf Q}/L^+)$, its inertia group $I_v$, and the restriction $\mathrm{loc}_v:G_v\to\Gamma_0$ of $p$; and the index $d_v$ of $G_v$ in $D_{A_v}$, which is the local degree of $L^+$ at $v$. The character $\omega$ is nontrivial on each $G_v$.
(v) A finite family $(c_w)$ of elements of $\Gamma$.
Moreover $\operatorname{Gal}(\bar{\mathbf Q}/L)=\operatorname{Gal}(\bar{\mathbf Q}/L^+)\cap\ker\eta$; $I_A\cap\operatorname{Gal}(\bar{\mathbf Q}/L^+)\subset\ker\eta$ for every place $A$, and $D_A\cap\operatorname{Gal}(\bar{\mathbf Q}/L^+)\subset\ker\eta$ for every place $A$ above a prime dividing $M$; and $\Psi(s)=1$ for every $s\in\operatorname{Gal}(\bar{\mathbf Q}/L^+)$ with $\bar\Psi(s)=1$ that lies in $I_A$ for a place $A$ not above $\ell$.
The link is \emph{realized} if $(c_w)$ is indexed by the infinite places $w$ of $L^+$, with $c_w=p(\tau_w)$ for a complex conjugation $\tau_w\in\operatorname{Gal}(\bar{\mathbf Q}/L^+)$ at $w$; the subgroup $H_{\eta,\Psi}=\operatorname{Gal}(\bar{\mathbf Q}/L^+)\cap\ker\eta\cap\ker\bar\Psi$ is open and normal in $\operatorname{Gal}(\bar{\mathbf Q}/L^+)$, of index prime to $\ell$; $\Gamma_0$ is open; and $\delta_\Gamma(c_\Gamma)=-1$.
\item Let $n\in\{2,3\}$, let $(e_{n,1},\dots,e_{n,n})$ be $(1,0,-1)$ if $n=3$ and $(1,-1)$ if $n=2$, and let $J_n$ be the antidiagonal $n\times n$ matrix with entries $1$. Let $A$ be a local ring with maximal ideal $\mathfrak m_A$ and a local homomorphism $\mathcal O\to A$. A \emph{linked point} of rank $n$ with values in $A$ consists of a homomorphism of monoids $\rho:\Gamma_0\to M_n(A)$ and homomorphisms $\psi_{v,i}:G_v\to A^\times$, for the places $v$ of (iv) and $1\le i\le n$, such that:
(a) $\rho(c_\Gamma g^{-1}c_\Gamma^{-1})=J_n\,{}^t\rho(g)J_n$ for $g\in\Gamma_0$;
(b) the entries off the diagonal of $\rho(g_0)-\rho(c_\Gamma g_0^{-1}c_\Gamma^{-1})$ vanish;
(c) $(\rho(\mathrm{loc}_vh_1)-\psi_{v,1}(h_1))\cdots(\rho(\mathrm{loc}_vh_n)-\psi_{v,n}(h_n))=0$ for every $v$ and $h_1,\dots,h_n\in G_v$;
(d) the entries of $\rho$ and the $\psi_{v,i}$ are continuous for the $\mathfrak m_A$-adic topology;
(e) modulo $\mathfrak m_A$, $\rho(g)$ is the diagonal matrix with entries $\bar\Psi_\Gamma(g)^{e_{n,i}}$, and $\psi_{v,i}(h)$ is $\bar\Psi_\Gamma(\mathrm{loc}_vh)^{e_{n,i}}$.
\end{enumerate}
\end{definition}
\leanref{Development/AutomorphicArithmetic/AutomorphicVarieties/CurveDVRSpecialization/HostFrameEx.lean}{UpperTheory.V.Pr.GalLink}
\leanref{Development/AutomorphicArithmetic/AutomorphicVarieties/FrobeniusSmoothVectors/ContinuousCoboundaries2.lean}{UpperTheory.V.Fr.Frame}
\leanref{Development/AutomorphicArithmetic/AutomorphicVarieties/FrobeniusSmoothVectors/FrobeniusSplitPrimeDatum.lean}{UpperTheory.V.Fr.Linked}
\leanref{Development/AutomorphicArithmetic/AutomorphicVarieties/GaloisCochains/GLinkQuot.lean}{UpperTheory.V.GC.FLink}
\leanref{Development/AutomorphicArithmetic/AutomorphicVarieties/FrobeniusSmoothVectors/LeAdicClosure.lean}{UpperTheory.V.Fr.NRep}
\leanref{Development/AutomorphicArithmetic/AutomorphicVarieties/CurveDVRSpecialization/PropagationTheoremGL2.lean}{UpperTheory.V.Pr.IsLiftA}

\begin{lemma}\label{sha-descent-link-exists}
Let $(\eta,\Psi,c,\mathfrak B,M,H_b)$ be a descent datum, let $U$ and $L$ be as in Lemma~\ref{sha-descent-subgroup}, and let $g_1\in\operatorname{Gal}(\bar{\mathbf Q}/L)$. Then there is a realized link of the datum to $L$ with base element $g_1$ (Definition~\ref{sha-descent-link}).
\end{lemma}
\leanref{Development/AutomorphicArithmetic/AutomorphicVarieties/FrobeniusSmoothVectors/LinkEx.lean}{UpperTheory.V.Fr.link_ex}
\begin{proof}
By Lemma~\ref{sha-descent-subgroup}(6), $\operatorname{Gal}(\bar{\mathbf Q}/L^+)=U$ and $\operatorname{Gal}(\bar{\mathbf Q}/L)=U\cap\ker\eta$.
Let $H_{\eta,\Psi}=U\cap\ker\eta\cap\ker\bar\Psi$.
It is open by Definition~\ref{sha-descent-datum}(c), and normal in $U$; and $U/H_{\eta,\Psi}$ embeds in $\{\pm1\}\times\mathbf F^\times$, whose order $2(|\mathbf F|-1)$ is prime to $\ell$.
Let $N$ be the closed normal subgroup of $U$ generated by the groups $I_A\cap H_{\eta,\Psi}$ for the places $A$ not above $\ell$, let $\Gamma=U/N$, a profinite group, and let $p$ be the quotient map.
Then $N\subset H_{\eta,\Psi}$, and $N\subset\ker\chi_{\mathrm{cyc}}$ because $\chi_{\mathrm{cyc}}$ is trivial on these inertia groups; so $\eta$, $\chi_{\mathrm{cyc}}$ and $\bar\Psi$ factor through $\Gamma$.
The complex conjugation $c$ lies in $U$ by Lemma~\ref{sha-descent-subgroup}(1), $\eta(c)=-1$ as in the proof of its part (6), and $g_1\in U\cap\ker\eta$.
Let $v$ be a place of $L^+$ above $\ell$ and $A$ a place of $\bar{\mathbf Q}$ above $v$.
The places of $\bar{\mathbf Q}$ above $\ell$ are conjugate, so $A=\gamma\mathfrak B$ for some $\gamma$; if $\eta(\gamma)=-1$, replace $A$ by $cA$, which lies above $v$ since $c\in U$, and $\gamma$ by $c\gamma$.
This gives $A_v$; two places of $\bar{\mathbf Q}$ above the same place of $L^+$ are conjugate under $U$.
The group $G_v=D_{A_v}\cap U$ lies in $\ker\eta$ by Lemma~\ref{sha-descent-subgroup}(2), since $\ell$ divides $M$, so $\mathrm{loc}_v$ takes values in $\Gamma_0$; its index in $D_{A_v}$ is the local degree, and $\omega$ is nontrivial on it by Lemma~\ref{sha-descent-subgroup}(4).
For each infinite place $w$ of the totally real field $L^+$, let $\tau_w\in U$ be a complex conjugation at $w$ and $c_w=p(\tau_w)$.
The group $\Gamma_0$ is open because its inverse image $U\cap\ker\eta$ is, and $\delta_\Gamma(c_\Gamma)=\eta(c)=-1$.
The remaining conditions are parts (2), (3) and (6) of Lemma~\ref{sha-descent-subgroup}.
\end{proof}

\begin{definition}\label{sha-descent-residual}
Let a link of a descent datum to an abelian CM field $L$ be given, with the notation of Definitions~\ref{sha-descent-datum} and~\ref{sha-descent-link}; let $V$ be the set of places of $L^+$ above $\ell$ and $d=\sum_{v\in V}d_v$.
Let $\mathbf F'$ be a finite field of characteristic $\ell$, $W$ a finite-dimensional $\mathbf F'$-vector space with a linear action of $\Gamma$ with open kernel, and $b=(b_v)_{v\in V}$ subspaces of $W$ with $b_v$ stable under $\mathrm{loc}_v(G_v)$.
\begin{enumerate}
\item Let $H^0_b(\Gamma,W)$ be the space of $X\in W^\Gamma$ with $X\in b_v$ for every $v$.
\item Let $H^1_b(\Gamma,W)$ be the quotient of the space of pairs $(x,(X_v))$, of a continuous cocycle $x:\Gamma\to W$ and elements $X_v\in W$ with $x(\mathrm{loc}_vh)-(\mathrm{loc}_v(h)X_v-X_v)\in b_v$ for all $v$ and $h\in G_v$, by the pairs $(g\mapsto gX-X,(X_v))$ with $X\in W$ and $X_v-X\in b_v$.
\item Let $H^2_b(\Gamma,W)$ be the quotient of the space of pairs $(z,(y_v))$, of a continuous $2$-cocycle $z$ of $\Gamma$ with values in $W$ and continuous maps $y_v:G_v\to W$ such that $z(\mathrm{loc}_vh,\mathrm{loc}_vh')$ and the coboundary of $y_v$ at $(h,h')$ differ by an element of $b_v$, by the pairs $(d\beta,(y_v))$ with $\beta:\Gamma\to W$ continuous and $y_v(h)-\beta(\mathrm{loc}_vh)+(\mathrm{loc}_v(h)Y_v-Y_v)\in b_v$ for some $Y_v\in W$.
\end{enumerate}
Write $H^i_\emptyset$ when every $b_v$ is $0$ and $H^i_{\mathrm{all}}$ when every $b_v$ is $W$, and $\mathbf F(\chi)$ for $\mathbf F$ with $\Gamma$ acting through a character $\chi$. Define $H^1_\emptyset(\Gamma_0,\mathbf F(\chi))$ for a character $\chi$ of $\Gamma_0$ in the same way, with the $2|V|$ homomorphisms $\mathrm{loc}_v$ and $h\mapsto c_\Gamma\mathrm{loc}_v(h)c_\Gamma^{-1}$ from $G_v$ to $\Gamma_0$. Let $I'_v$ be the set of $h\in I_v$ with $\varepsilon_\Gamma(\mathrm{loc}_vh)=1$.
For an integer $C_0\ge0$, the link satisfies the \emph{residual conditions} with constant $C_0$ if:
\begin{enumerate}
\setcounter{enumi}{3}
\item $\mathbf F$ is finite; $\Gamma_0$ and $\ker\bar\Psi_\Gamma$ are open; $c_\Gamma^2=1$, $\delta_\Gamma(c_\Gamma)=-1$ and $\bar\Psi_\Gamma(c_\Gamma)=1$; $\bar\Psi_\Gamma(g_0)^2\ne1$; and $\varepsilon_\Gamma$ and the maps $\mathrm{loc}_v$ are continuous;
\item every open subgroup of $\Gamma$ has only finitely many continuous homomorphisms to $\mathbf Z/\ell$;
\item the family $(c_w)$ has $d$ members, and $c_w^2=1$ and $\delta_\Gamma(c_w)=-1$ for each of them;
\item $\omega_\Gamma$, $\omega_\Gamma\bar\Psi_\Gamma$ and $\omega_\Gamma\bar\Psi_\Gamma^{-1}$ are nontrivial on $\Gamma_0$; $\omega_\Gamma\bar\Psi_\Gamma^{2}\delta_\Gamma$ and $\omega_\Gamma\bar\Psi_\Gamma^{-2}\delta_\Gamma$ are nontrivial; there is no integer $j$ with $\omega_\Gamma=\bar\Psi_\Gamma^j$ on $\Gamma_0$; and for every $v$, $\bar\Psi_\Gamma^2$ is nontrivial on $\mathrm{loc}_v(G_v)$ and $\bar\Psi_\Gamma$ is trivial on $\mathrm{loc}_v(I'_v)$;
\item $H^1_\emptyset(\Gamma,\mathbf F(\bar\Psi_\Gamma^{-1}\delta_\Gamma))$ and $H^1_\emptyset(\Gamma,\mathbf F(\bar\Psi_\Gamma^{-2}\delta_\Gamma))$ are nonzero;
\item for every $v$ there are $u_1,\dots,u_{d_v}\in I_v$ with $\bar\Psi_\Gamma(\mathrm{loc}_vu_i)=1$, all but one of them in $I'_v$, such that every continuous homomorphism from $G_v$ to a finite abelian $\ell$-group that kills $u_1,\dots,u_{d_v}$ kills $I_v$;
\item there is $h\ge1$ such that $\chi^h=1$ for every continuous homomorphism $\chi$ from $\Gamma_0$ to a finite abelian $\ell$-group that kills $\mathrm{loc}_v(I_v)$ and $c_\Gamma\mathrm{loc}_v(I_v)c_\Gamma^{-1}$ for every $v$;
\item for every $\mathbf F'$, $W$ and $b$ as above such that no nonzero $\Gamma$-equivariant linear map $W\to\mathbf F'(\omega_\Gamma)$ exists, $H^2_b(\Gamma,W)$ is finite and
\[\dim H^2_b(\Gamma,W)+\dim H^0_b(\Gamma,W)+\sum_vd_v\dim b_v\le\dim H^1_b(\Gamma,W)+\sum_w\dim\ker(c_w-1\mid W);\]
\item $\dim H^1_\emptyset(\Gamma,\mathbf F(\bar\Psi_\Gamma^{-2}\delta_\Gamma))\le C_0$, $\dim H^1_\emptyset(\Gamma_0,\mathbf F(\bar\Psi_\Gamma^{-1}))\le C_0$ and $\dim H^1_{\mathrm{all}}(\Gamma,\mathbf F)\le1+C_0$.
\end{enumerate}
The link is \emph{large} for $C_0$ if $d>3C_0+7$ and $d_v>5$ for every $v$.
\end{definition}
\leanref{Development/AutomorphicArithmetic/AutomorphicVarieties/FrobeniusSmoothVectors/ContinuousCoboundaries2.lean}{UpperTheory.V.Fr.Frame.Ok}
\leanref{Development/AutomorphicArithmetic/AutomorphicVarieties/FrobeniusSmoothVectors/ContinuousCoboundaries2.lean}{UpperTheory.V.Fr.Frame.Big}
\leanref{Development/AutomorphicArithmetic/AutomorphicVarieties/FrobeniusSmoothVectors/ContinuousCoboundaries2.lean}{UpperTheory.V.Fr.Frame.EC}
\leanref{Development/AutomorphicArithmetic/AutomorphicVarieties/FrobeniusSmoothVectors/ContinuousCoboundaries2.lean}{UpperTheory.V.Fr.H1C}
\leanref{Development/AutomorphicArithmetic/AutomorphicVarieties/FrobeniusSmoothVectors/ContinuousCoboundaries2.lean}{UpperTheory.V.Fr.coneH2}

\begin{lemma}\label{sha-descent-strict}
Let $(\eta,\Psi,c,\mathfrak B,M,H_b)$ be a descent datum, with the notation of Definition~\ref{sha-descent-datum}, let $U$ be as in Lemma~\ref{sha-descent-subgroup}, for an auxiliary prime $w$ and an integer $m\ge0$, let $a\in\{1,2\}$, and let $\kappa=\bar\Psi^{-a}\eta:G_{\mathbf Q}\to\mathbf F^\times$. Then there is a continuous map $y:G_{\mathbf Q}\to\mathbf F$ with $y(st)=y(s)+\kappa(s)y(t)$ for all $s,t$, such that:
\begin{enumerate}
\item there is no $X\in\mathbf F$ with $y(s)=\kappa(s)X-X$ for all $s\in U$;
\item for every place $A$ above $\ell$ there is $X\in\mathbf F$ with $y(s)=\kappa(s)X-X$ for all $s\in D_A$;
\item $y(s)=0$ for every $s\in U$ with $\eta(s)=\bar\Psi(s)=1$ that lies in $I_A$ for a place $A$ not above $\ell$.
\end{enumerate}
\end{lemma}
\leanref{Development/AutomorphicArithmetic/AutomorphicVarieties/FrobeniusSmoothVectors/Strict.lean}{UpperTheory.V.Fr.strict}
\leanref{Development/AutomorphicArithmetic/AutomorphicVarieties/FrobeniusSmoothVectors/FrobeniusSplitPrimeDatum.lean}{UpperTheory.V.Fr.StrictCl}
\leanref{Development/AutomorphicArithmetic/AutomorphicVarieties/FrobeniusSmoothVectors/StSel.lean}{UpperTheory.V.Fr.st_sel}
\leanref{Development/AutomorphicArithmetic/AutomorphicVarieties/FrobeniusSmoothVectors/StCount.lean}{UpperTheory.V.Fr.st_count}
\leanref{Development/AutomorphicArithmetic/AutomorphicVarieties/FrobeniusSmoothVectors/StTame.lean}{UpperTheory.V.Fr.st_tame}
\leanref{Development/AutomorphicArithmetic/AutomorphicVarieties/FrobeniusSmoothVectors/StRes.lean}{UpperTheory.V.Fr.st_res}
\begin{proof}
By Definition~\ref{sha-descent-datum}(c), there is $b\ge0$ such that $\kappa$ is trivial on the elements fixing $\mu_{M^b}$; so $\ker\kappa$ is open and $\kappa$ is trivial on $I_A$ for the places $A$ above primes not dividing $M$.
Since $\delta$ is not real, $\eta(c)=-1$, and $\kappa(c)=\bar\Psi(c)^{-a}\eta(c)=-1\ne1$.
By Lemma~\ref{sha-descent-subgroup}(5) there is $s_0\in U\cap\ker\eta$ with $\omega(s_0)\bar\Psi(s_0)^a\ne1$, that is, $\kappa(s_0)\ne\omega(s_0)$.
For a place $A$ above $w$, $D_A\subset\ker\kappa\cap\ker\omega$, since $w$ is an auxiliary prime.
For a place $A$ above $\ell$, Lemma~\ref{sha-descent-subgroup}(4) gives $s\in D_A\cap U$ with $\bar\Psi(s)^2\ne1$; then $\eta(s)=1$ by its part (2), so $\kappa(s)=\bar\Psi(s)^{-a}\ne1$.

Let $S$ be the finite set consisting of $\ell$, $w$ and the primes $q$ with $I_A\not\subset\ker\kappa$ for the places $A$ above $q$, which divide $M$; let $G_S$ be the Galois group of the maximal extension of $\mathbf Q$ unramified outside $S$ and the infinite place.
The character $\kappa$ factors through $G_S$.
Consider the Selmer group of $\mathbf F(\kappa)$ with the local condition everything at $w$ and $0$ at the other primes of $S$.
Here $H^0(G_S,\mathbf F(\kappa))=0$ and $H^0$ at the infinite place vanishes, because $\kappa(c)\ne1$; $H^0(G_S,\mathbf F(\kappa^{-1}\omega))=0$ because $\kappa(s_0)\ne\omega(s_0)$; $H^0(D_A,\mathbf F(\kappa))=0$ for the primes of $S$ other than $w$, because $\kappa$ is nontrivial on $D_A$; and since $\kappa$ and $\omega$ are trivial on $D_A$ for $A$ above $w\ne\ell$, the local Euler characteristic formula and local duality \cite[Chapter~VII, \S\S2--3]{NSW08} give $|H^0(D_A,\mathbf F)|<|H^1(D_A,\mathbf F)|$.
The Greenberg-Wiles formula \cite[Proposition~1.6]{Wil95} then shows that the Selmer group has more than one element.
Let $y$ be the inflation to $G_{\mathbf Q}$ of a cocycle representing a nonzero element of it.
It is not a coboundary on $G_{\mathbf Q}$; it vanishes on $I_A$ for the places $A$ above primes outside $S$; and it is a coboundary on $D_A$ for one place $A$ above each prime of $S$ other than $w$, hence, by conjugation, for every such place.
This gives (2).

For (3), let $A$ be a place not above $\ell$ and $s\in I_A\cap U$ with $\eta(s)=\bar\Psi(s)=1$, so that $\kappa(s)=1$.
If $A$ lies above a prime outside $S$, then $y(s)=0$.
If $A$ lies above a prime of $S$ other than $w$, then $y(s)=\kappa(s)X-X=0$.
If $A$ lies above $w$, then $y$ is additive on $I_A$, because $\kappa$ is trivial on $D_A$, and $s\in H_{\mathrm{th}}(w)$.
Let $\phi:G_{\mathbf Q}\to\mathbf Z/\ell$ be a character with kernel $H_{\mathrm{th}}(w)$; it is nonzero on $I_A$ by Lemma~\ref{sha-descent-auxiliary}(2).
The maximal pro-$\ell$ quotient of $I_A$ is procyclic \cite[Chapter~VII]{NSW08}, so there is $\gamma\in I_A$ such that every continuous homomorphism $I_A\to\mathbf Z/\ell$ vanishing at $\gamma$ is $0$; then $\phi(\gamma)\ne0$.
For an $\mathbf F_\ell$-linear form $\lambda$ on $\mathbf F$, the homomorphism $\lambda\circ y-(\lambda(y(\gamma))/\phi(\gamma))\phi$ on $I_A$ vanishes at $\gamma$, hence everywhere; at $s$ it gives $\lambda(y(s))=0$.
So $y(s)=0$.

For (1), suppose that $y(s)=\kappa(s)X-X$ for $s\in U$.
Some $u\in U$ has $\kappa(u)\ne1$, as shown above for a place above $\ell$.
For $g\in G_{\mathbf Q}$, $gug^{-1}\in U$, because $U$ contains the commutator subgroup, and $\kappa(gug^{-1})=\kappa(u)$.
Comparing $y(gu)=y(g)+\kappa(g)y(u)$ with $y(gu)=y(gug^{-1})+\kappa(u)y(g)$ gives $(\kappa(u)-1)(y(g)-\kappa(g)X+X)=0$.
So $y$ would be a coboundary on $G_{\mathbf Q}$, which it is not.
\end{proof}

\begin{lemma}[Tate \cite{NSW08}]\label{sha-descent-link-local-tate}
Let $\ell$ be a prime, $F\subset\bar{\mathbf Q}$ a number field, $v$ a finite place of $F$, $A$ a place of $\bar{\mathbf Q}$ above $v$, and $D$ the decomposition group of $A$ in $\operatorname{Gal}(\bar{\mathbf Q}/F)$. Let $r\ge1$, let $\chi_r:D\to(\mathbf Z/\ell^r)^\times$ be the reduction of $\chi_{\mathrm{cyc}}$ modulo $\ell^r$, and write $\mathbf Z/\ell^r(\chi_r)$ for $\mathbf Z/\ell^r$ with $D$ acting through $\chi_r$.
Let $M$ be a finite $\mathbf Z/\ell^r$-module with a $\mathbf Z/\ell^r$-linear action of $D$ that is continuous for the discrete topology on $M$. Write $H^i(D,M)$ for the cohomology of continuous cochains, and $M^*(1)$ for the $\mathbf Z/\ell^r$-module of $\mathbf Z/\ell^r$-linear maps $\lambda:M\to\mathbf Z/\ell^r$, on which $g\in D$ acts by $(g\lambda)(m)=\chi_r(g)\lambda(g^{-1}m)$.
\begin{enumerate}
\item $H^0(D,M)$, $H^1(D,M)$ and $H^2(D,M)$ are finite, and $H^n(D,M)=0$ for $n\ge3$.
\item If $v$ lies above $\ell$, then $|H^1(D,M)|=|H^0(D,M)|\,|H^2(D,M)|\,|M|^{[F_v:\mathbf Q_\ell]}$. If $v$ does not lie above $\ell$, then $|H^1(D,M)|=|H^0(D,M)|\,|H^2(D,M)|$.
\item $H^2(D,\mathbf Z/\ell^r(\chi_r))$ is isomorphic to $\mathbf Z/\ell^r$. For the cup products $H^i(D,M)\times H^{2-i}(D,M^*(1))\to H^2(D,\mathbf Z/\ell^r(\chi_r))$ given by the pairing $(m,\lambda)\mapsto\lambda(m)$, the map sending $x\in H^i(D,M)$ to the homomorphism $y\mapsto x\cup y$ on $H^{2-i}(D,M^*(1))$ is bijective for $i=0,1,2$, and so is the map sending $y\in H^1(D,M^*(1))$ to the homomorphism $x\mapsto x\cup y$ on $H^1(D,M)$.
\end{enumerate}
\end{lemma}
\leanref{ReciprocityTheory/Hypotheses/GalCoh.lean}{ReciprocityTheory.Gal.LocalFiniteness}
\leanref{ReciprocityTheory/Hypotheses/GalCoh.lean}{ReciprocityTheory.Gal.LocalTateDuality}
\leanref{ReciprocityTheory/Hypotheses/GalCoh.lean}{ReciprocityTheory.Gal.LocalEulerChar}

\begin{lemma}\label{sha-descent-link-local}
Let $\ell$ be a prime, $F\subset\bar{\mathbf Q}$ a number field and $\omega:\operatorname{Gal}(\bar{\mathbf Q}/F)\to\mathbf F_\ell^\times$ the reduction of $\chi_{\mathrm{cyc}}$ modulo $\ell$.
Let $v$ be a finite place of $F$, $A$ a place of $\bar{\mathbf Q}$ above $v$, and $D$ the decomposition group of $A$ in $\operatorname{Gal}(\bar{\mathbf Q}/F)$.
Let $k$ be a finite field of characteristic $\ell$ and $M$ a finite-dimensional $k$-vector space with a continuous linear action of $D$; write $H^i(D,M)$ for the cohomology of continuous cochains, and $\operatorname{Hom}_D(M,k(\omega))$ for the space of linear forms $\lambda$ on $M$ with $\lambda(gm)=\omega(g)\lambda(m)$ for $g\in D$.
\begin{enumerate}
\item $H^1(D,M)$ and $H^2(D,M)$ are finite-dimensional, and $\dim H^2(D,M)=\dim\operatorname{Hom}_D(M,k(\omega))$.
\item If $v$ lies above $\ell$, then $\dim H^1(D,M)=\dim H^0(D,M)+\dim H^2(D,M)+[F_v:\mathbf Q_\ell]\dim M$. If $v$ does not lie above $\ell$, then $\dim H^1(D,M)=\dim H^0(D,M)+\dim H^2(D,M)$.
\item For every surjective $D$-equivariant linear map $M\to Q$ onto a $k$-vector space $Q$ with a linear action of $D$, the induced map $H^2(D,M)\to H^2(D,Q)$ is surjective.
\end{enumerate}
\end{lemma}
\leanref{Development/AutomorphicArithmetic/AutomorphicVarieties/GaloisCochains/LocEuler.lean}{UpperTheory.V.GC.locEuler}
\leanref{Development/AutomorphicArithmetic/AutomorphicVarieties/GaloisCochains/H2RightExactModL.lean}{UpperTheory.V.GC.local_second_cohomology_surjective_at_place}
\begin{proof}
View $M$ as a $\mathbf Z/\ell$-module, and apply Lemma~\ref{sha-descent-link-local-tate} with $r=1$, so that $\chi_1=\omega$. The continuous cochains of $D$ with values in $M$ do not depend on the ring of scalars, so neither do the groups $H^i(D,M)$.
Each $H^i(D,M)$ is a $k$-vector space, finite by Lemma~\ref{sha-descent-link-local-tate}(1); so it is finite-dimensional, with $|k|^{\dim H^i(D,M)}$ elements, and $|M|=|k|^{\dim M}$. So (2) is Lemma~\ref{sha-descent-link-local-tate}(2), written with exponents of $|k|$.

For (1), Lemma~\ref{sha-descent-link-local-tate}(3) with $i=2$ gives a bijection from $H^2(D,M)$ onto the homomorphisms from the finite $\mathbf F_\ell$-vector space $H^0(D,M^*(1))$ to $H^2(D,\mathbf F_\ell(\omega))\cong\mathbf F_\ell$; so $|H^2(D,M)|=|H^0(D,M^*(1))|$. The group $H^0(D,M^*(1))$ is the set of $\mathbf F_\ell$-linear forms $\mu$ on $M$ with $\mu(gm)=\omega(g)\mu(m)$ for $g\in D$.
Let $T(\lambda)=\operatorname{Tr}_{k/\mathbf F_\ell}\circ\lambda$ for a $k$-linear form $\lambda$ on $M$. If $T(\lambda)=0$, then $\operatorname{Tr}_{k/\mathbf F_\ell}(a\lambda(m))=T(\lambda)(am)=0$ for all $a\in k$, so $\lambda(m)=0$, as the trace form of $k$ over $\mathbf F_\ell$ is nondegenerate. So $T$ is injective, and it is bijective, as both spaces of forms have $|M|$ elements.
As $\omega$ takes values in $\mathbf F_\ell^\times$, $T$ maps $\operatorname{Hom}_D(M,k(\omega))$ into $H^0(D,M^*(1))$. Conversely, if $T(\lambda)\in H^0(D,M^*(1))$, then for all $a\in k$, $\operatorname{Tr}_{k/\mathbf F_\ell}\bigl(a\lambda(gm)\bigr)=T(\lambda)(g(am))=\omega(g)T(\lambda)(am)=\operatorname{Tr}_{k/\mathbf F_\ell}\bigl(a\,\omega(g)\lambda(m)\bigr)$, so $\lambda(gm)=\omega(g)\lambda(m)$.
Hence $|H^2(D,M)|=|\operatorname{Hom}_D(M,k(\omega))|$, which is (1).

For (3), let $N$ be the kernel of $q:M\to Q$, a subspace stable under $D$, and $z$ a continuous $2$-cocycle of $D$ with values in $Q$. Choose a map $s:Q\to M$ with $q\circ s=\mathrm{id}$, and let $\Phi=s\circ z$, which is continuous as $M$ and $Q$ are discrete. Its coboundary $d\Phi$ is a continuous $3$-cochain with values in $q^{-1}(0)=N$, and a $3$-cocycle. As $N$ is a finite $\mathbf Z/\ell$-module, $H^3(D,N)=0$ by Lemma~\ref{sha-descent-link-local-tate}(1), so $d\Phi=db$ for a continuous $2$-cochain $b$ with values in $N$. Then $\Phi-b$ is a $2$-cocycle with values in $M$, and $q\circ(\Phi-b)=z$.
\end{proof}
% locEuler (LocEuler.lean) for (1) and (2): Loc.resCoef to Z/l^1, the equivalences e0, e1, e2 of H^i, the fields of LocalFiniteness and LocalEulerChar at r = 1, and, for (1), cup20 of LocalTateDuality with card_H0dual (the trace map T, HcLocEuler.lean) and card_twDual. local_second_cohomology_surjective_at_place (H2RightExactModL.lean) for (3): the private H2_map_surjective_of_H3_subsingleton (sec, Phi, c and b as in the text), with cohDimTwoLocal_of_localFiniteness and CohDimTwoLocal.of_nsmul for H^3 of the kernel.

\begin{lemma}\label{sha-descent-link-unramified}
Let $\ell$ be a prime, $F\subset\bar{\mathbf Q}$ a number field, $v$ a finite place of $F$, $A$ a place of $\bar{\mathbf Q}$ above $v$, and $D\supset I$ the decomposition and inertia groups of $A$ in $\operatorname{Gal}(\bar{\mathbf Q}/F)$.
\begin{enumerate}
\item There is $\varphi\in D$ such that $D$ is the closure of the subgroup generated by $I$ and $\varphi$, and for every $n\ge1$ there is a continuous homomorphism $D\to\mathbf Z/n$ that vanishes on $I$ and takes $\varphi$ to $1$.
\item Let $N\subset I$ be a closed subgroup, normal in $D$, of finite index prime to $\ell$ in $I$. Let $k$ be a finite field of characteristic $\ell$ and $M$ a finite-dimensional $k$-vector space with a continuous linear action of $D$ that is trivial on $N$. Then every continuous $2$-cocycle $z$ of $D$ with values in $M$ such that $z(sn,tn')=z(s,t)$ for all $s,t\in D$ and $n,n'\in N$ is the coboundary of a continuous map $b:D\to M$ with $b(sn)=b(s)$ for all $s\in D$ and $n\in N$.
\end{enumerate}
\end{lemma}
\leanref{Development/AutomorphicArithmetic/AutomorphicVarieties/GaloisCochains/DecZhat.lean}{UpperTheory.V.GC.decZhat}
\leanref{Development/AutomorphicArithmetic/AutomorphicVarieties/GaloisCochains/IsClosedOfClosedFiniteIndex.lean}{UpperTheory.V.GC.H2_unramified_inflation}
\begin{proof}
(1) The quotient $D/I$ is the Galois group of the extension of residue fields at $A$, which is isomorphic to $\hat{\mathbf Z}$ and topologically generated by the Frobenius \cite[Chapter~II, \S9]{Neu99}.
Take for $\varphi$ an element of $D$ that maps to the Frobenius, and for the homomorphism to $\mathbf Z/n$ the composite of $D\to D/I\cong\hat{\mathbf Z}$ with the reduction modulo $n$.

(2) The group $D/N$ is an extension of $D/I\cong\hat{\mathbf Z}$, of $\ell$-cohomological dimension $1$, by the finite group $I/N$ of order prime to $\ell$, of $\ell$-cohomological dimension $0$.
So $D/N$ has $\ell$-cohomological dimension at most $1$ \cite[Chapter~III, \S3]{NSW08}, and $H^2(D/N,M)=0$.
The cocycles $z$ of the statement are the $2$-cocycles of $D/N$, and the maps $b$ the cochains of $D/N$.
\end{proof}

\begin{lemma}[Tate, Poitou \cite{NSW08}]\label{sha-descent-link-global-tate}
Let $\ell$ be a prime, $F\subset\bar{\mathbf Q}$ a number field with $r_2$ complex places, and $r\ge1$.
Let $S'$ be a finite set of finite places of $F$ containing the places above $\ell$, and $G_{S'}$ the Galois group over $F$ of the maximal extension of $F$ in $\bar{\mathbf Q}$ unramified at every finite place outside $S'$; the reduction of $\chi_{\mathrm{cyc}}$ modulo $\ell^r$ factors through a homomorphism $\chi_r:G_{S'}\to(\mathbf Z/\ell^r)^\times$.
For each $v\in S'$ let $A_v$ be a place of $\bar{\mathbf Q}$ above $v$ and $D_v$ its decomposition group in $\operatorname{Gal}(\bar{\mathbf Q}/F)$, which maps to $G_{S'}$; for each real place $w$ of $F$ let $c_w\in\operatorname{Gal}(\bar{\mathbf Q}/F)$ be a complex conjugation at $w$.
Let $M$ be a finite $\mathbf Z/\ell^r$-module with a $\mathbf Z/\ell^r$-linear action of $G_{S'}$ that is continuous for the discrete topology on $M$, and $M^*(1)$ the module of $\mathbf Z/\ell^r$-linear maps $\lambda:M\to\mathbf Z/\ell^r$, on which $g$ acts by $(g\lambda)(m)=\chi_r(g)\lambda(g^{-1}m)$.
Write $H^i$ for the cohomology of continuous cochains, $X^\vee$ for the module of $\mathbf Z/\ell^r$-linear maps $X\to\mathbf Z/\ell^r$, and, for $i=1,2$ and $N=M$ or $N=M^*(1)$, $\text{Ш}^i_{S'}(N)$ for the kernel of the restriction map $H^i(G_{S'},N)\to\prod_{v\in S'}H^i(D_v,N)$.
\begin{enumerate}
\item $H^0(G_{S'},M)$, $H^1(G_{S'},M)$ and $H^2(G_{S'},M)$ are finite; if $\ell$ is odd, $H^n(G_{S'},M)=0$ for $n\ge3$.
\item $|H^0(G_{S'},M)|\,|H^2(G_{S'},M)|\,|M|^{[F:\mathbf Q]}=|H^1(G_{S'},M)|\prod_{w\text{ real}}|M^{c_w}|\,|M|^{r_2}$.
\item Suppose that $\ell$ is odd or $F$ is totally complex, and let $\zeta\in\bar{\mathbf Q}$ be a primitive $\ell^r$-th root of unity. For each place $A$ of $\bar{\mathbf Q}$ above a finite place $v$ of $F$, with decomposition group $D_A$ in $\operatorname{Gal}(\bar{\mathbf Q}/F)$, write $\mathbf Z/\ell^r(\chi_r)$ for $\mathbf Z/\ell^r$ with $D_A$ acting through $\chi_r$. There are bijective $\mathbf Z/\ell^r$-linear maps $\operatorname{inv}_A:H^2(D_A,\mathbf Z/\ell^r(\chi_r))\to\mathbf Z/\ell^r$, one for each such $A$ and depending only on $F$, $\ell$, $r$ and $\zeta$, with the following properties.
\begin{enumerate}
\item[(a)] $\operatorname{inv}_A$ takes the value $1$ at the class of every $2$-cocycle $(g,h)\mapsto\kappa(g)\phi(h)$ in which $\kappa:D_A\to\mathbf Z/\ell^r$ satisfies $g(y)=\zeta^{\kappa(g)}y$ for some $y\in\bar{\mathbf Q}$ with $y^{\ell^r}=\varpi$, for some $\varpi\in\mathcal O_F$ that lies in $v$ but not in $v^2$, and $\phi:D_A\to\mathbf Z/\ell^r$ is a homomorphism that vanishes on the inertia group and takes the value $1$ at every Frobenius element.
\item[(b)] For $S'$, the $A_v$ and $M$ as above, write $\operatorname{inv}_v=\operatorname{inv}_{A_v}$. For some homomorphisms $H^2(G_{S'},M^*(1))^\vee\to H^1(G_{S'},M)$ and $H^1(G_{S'},M^*(1))^\vee\to H^2(G_{S'},M)$, the sequence
\[
\begin{aligned}
0&\to H^0(G_{S'},M)\to\textstyle\bigoplus_{v\in S'}H^0(D_v,M)\to H^2(G_{S'},M^*(1))^\vee\to H^1(G_{S'},M)\to\textstyle\bigoplus_{v\in S'}H^1(D_v,M)\\
&\to H^1(G_{S'},M^*(1))^\vee\to H^2(G_{S'},M)\to\textstyle\bigoplus_{v\in S'}H^2(D_v,M)\to H^0(G_{S'},M^*(1))^\vee\to0
\end{aligned}
\]
is exact. Here the maps $H^i(G_{S'},M)\to\bigoplus_vH^i(D_v,M)$ are the restrictions, and the maps $\bigoplus_vH^i(D_v,M)\to H^{2-i}(G_{S'},M^*(1))^\vee$ send $(x_v)$ to $y\mapsto\sum_v\operatorname{inv}_v(x_v\cup y|_{D_v})$, for the cup products given by the pairing $(m,\lambda)\mapsto\lambda(m)$.
\item[(c)] $\text{Ш}^1_{S'}(M)$, $\text{Ш}^2_{S'}(M)$, $\text{Ш}^1_{S'}(M^*(1))$ and $\text{Ш}^2_{S'}(M^*(1))$ are finite, $\text{Ш}^2_{S'}(M)\cong\text{Ш}^1_{S'}(M^*(1))^\vee$ and $\text{Ш}^2_{S'}(M^*(1))\cong\text{Ш}^1_{S'}(M)^\vee$, and $H^n(G_{S'},M)=0$ for $n\ge3$.
\end{enumerate}
\end{enumerate}
\end{lemma}
\leanref{ReciprocityTheory/Hypotheses/GalCoh.lean}{ReciprocityTheory.Gal.GlobalFiniteness}
\leanref{ReciprocityTheory/Hypotheses/GalCoh.lean}{ReciprocityTheory.Gal.GlobalEulerChar}
\leanref{Development/AbelianVarietyTheory/PublishedTheorems/GaloisCohomology₂/PoitouTateNormalised.lean}{GlobalTheory.Published.PoitouTateNormalised}

\begin{lemma}\label{sha-descent-link-global}
Let $\ell$ be an odd prime, $F\subset\bar{\mathbf Q}$ a number field and $\omega$ the reduction of $\chi_{\mathrm{cyc}}$ modulo $\ell$.
Let $S'$ be a finite set of finite places of $F$ containing the places above $\ell$, and $G_{S'}$ the Galois group over $F$ of the maximal extension of $F$ in $\bar{\mathbf Q}$ unramified at every finite place outside $S'$.
For each $v\in S'$ let $A_v$ be a place of $\bar{\mathbf Q}$ above $v$ and $D_v$ its decomposition group in $\operatorname{Gal}(\bar{\mathbf Q}/F)$, which maps to $G_{S'}$; for each infinite place $w$ of $F$ let $c_w\in\operatorname{Gal}(\bar{\mathbf Q}/F)$ be a complex conjugation at $w$.
Let $k$ be a finite field of characteristic $\ell$ and $M$ a finite-dimensional $k$-vector space with a continuous linear action of $G_{S'}$.
Let $\text{Ш}^2_{S'}(M)$ be the kernel of the restriction map $H^2(G_{S'},M)\to\prod_{v\in S'}H^2(D_v,M)$, and $\operatorname{Hom}_{G_{S'}}(M,k(\omega))$ the space of linear forms $\lambda$ on $M$ with $\lambda(gm)=\omega(g)\lambda(m)$ for $g\in G_{S'}$.
Then $H^1(G_{S'},M)$ and $H^2(G_{S'},M)$ are finite-dimensional, and:
\begin{enumerate}
\item if $F$ is totally real, $\dim H^1(G_{S'},M)+\sum_w\dim M^{c_w}=\dim H^0(G_{S'},M)+\dim H^2(G_{S'},M)+[F:\mathbf Q]\dim M$;
\item $\dim\text{Ш}^2_{S'}(M)+\sum_{v\in S'}\dim H^2(D_v,M)=\dim H^2(G_{S'},M)+\dim\operatorname{Hom}_{G_{S'}}(M,k(\omega))$.
\end{enumerate}
\end{lemma}
\leanref{Development/AutomorphicArithmetic/AutomorphicVarieties/GaloisCochains/GlobEuler.lean}{UpperTheory.V.GC.globEuler}
\leanref{Development/AutomorphicArithmetic/AutomorphicVarieties/GaloisCochains/PtTail.lean}{UpperTheory.V.GC.ptTail}
\begin{proof}
View $M$ as a $\mathbf Z/\ell$-module, and apply Lemma~\ref{sha-descent-link-global-tate} with $r=1$, so that $\chi_1=\omega$. The continuous cochains of $G_{S'}$ and of the $D_v$ with values in $M$ do not depend on the ring of scalars, so neither do the groups $H^i(G_{S'},M)$, $H^i(D_v,M)$ and $\text{Ш}^2_{S'}(M)$.
Each of them is a $k$-vector space. $H^1(G_{S'},M)$ and $H^2(G_{S'},M)$ are finite by Lemma~\ref{sha-descent-link-global-tate}(1), and each $H^2(D_v,M)$ by Lemma~\ref{sha-descent-link-local-tate}(1). So they are finite-dimensional, a finite-dimensional $k$-vector space $X$ has $|k|^{\dim X}$ elements, and $|M|=|k|^{\dim M}$.

For (1), $F$ has no complex places, so every infinite place $w$ is real, and each $M^{c_w}$ is a subspace of $M$. So Lemma~\ref{sha-descent-link-global-tate}(2) reads $|k|^{\dim H^0(G_{S'},M)+\dim H^2(G_{S'},M)+[F:\mathbf Q]\dim M}=|k|^{\dim H^1(G_{S'},M)+\sum_w\dim M^{c_w}}$, which is (1).

For (2), let $\zeta$ be a primitive $\ell$-th root of unity and $\operatorname{inv}_v$ the maps of Lemma~\ref{sha-descent-link-global-tate}(3), which applies as $\ell$ is odd. By the exactness of its sequence at the last two terms, the restriction map $H^2(G_{S'},M)\to\bigoplus_vH^2(D_v,M)$, whose kernel is $\text{Ш}^2_{S'}(M)$, has as image the kernel of a surjection onto $H^0(G_{S'},M^*(1))^\vee$. Hence
\[
|\text{Ш}^2_{S'}(M)|\prod_{v\in S'}|H^2(D_v,M)|=|H^2(G_{S'},M)|\,|H^0(G_{S'},M^*(1))^\vee|.
\]
The group $H^0(G_{S'},M^*(1))$ is a finite $\mathbf F_\ell$-vector space, so it has as many elements as its dual. Its elements are the $\mathbf F_\ell$-linear forms $\mu$ on $M$ with $\mu(gm)=\omega(g)\mu(m)$ for $g\in G_{S'}$, since $g\mu=\mu$ means $\mu(g^{-1}m)=\omega(g)^{-1}\mu(m)$ for every $m$. As in the proof of Lemma~\ref{sha-descent-link-local}, $\lambda\mapsto\operatorname{Tr}_{k/\mathbf F_\ell}\circ\lambda$ is a bijection from $\operatorname{Hom}_{G_{S'}}(M,k(\omega))$ onto these forms. Writing the cardinalities as powers of $|k|$ gives (2).
\end{proof}
% globEuler (GlobEuler.lean, 57 lines) for the finiteness and (1): M as a Z/l^1-module through ZMod.castHom, the equivalences e0, e1, e2 (subquotEquiv) of the H^i, ec for the fixed points of the c w, the fields of GlobalFiniteness and GlobalEulerChar at r = 1, and archH0Card with nrComplexPlaces = 0. ptTail (PtTail.lean, 176 lines) for (2): resCoef, exists_primitiveRoot_pow, exists_lift_cycMod for chi_1, the fields hex and hsurj of NineTermExact with card_exact, and hd from card_twDual and card_H0dual (the trace map).

\begin{lemma}[Greenberg, Wiles \cite{Wil95}]\label{sha-descent-link-greenberg-wiles}
Let $\ell$ be an odd prime, $F\subset\bar{\mathbf Q}$ a number field with $r_2$ complex places, and $r\ge1$.
Let $S'$ be a finite set of finite places of $F$ containing the places above $\ell$, and $G_{S'}$ the Galois group over $F$ of the maximal extension of $F$ in $\bar{\mathbf Q}$ unramified at every finite place outside $S'$; the reduction of $\chi_{\mathrm{cyc}}$ modulo $\ell^r$ factors through a homomorphism $\chi_r:G_{S'}\to(\mathbf Z/\ell^r)^\times$.
For each $v\in S'$ let $D_v$ be the decomposition group in $\operatorname{Gal}(\bar{\mathbf Q}/F)$ of a place of $\bar{\mathbf Q}$ above $v$, and for each real place $w$ of $F$ let $c_w\in\operatorname{Gal}(\bar{\mathbf Q}/F)$ be a complex conjugation at $w$.
Let $M$ be a finite $\mathbf Z/\ell^r$-module with a $\mathbf Z/\ell^r$-linear action of $G_{S'}$ that is continuous for the discrete topology on $M$, and $M^*(1)$ the module of $\mathbf Z/\ell^r$-linear maps $\lambda:M\to\mathbf Z/\ell^r$, on which $g$ acts by $(g\lambda)(m)=\chi_r(g)\lambda(g^{-1}m)$.
For each $v\in S'$ let $L_v\subset H^1(D_v,M)$ be a submodule, and $L_v^\perp\subset H^1(D_v,M^*(1))$ the submodule of classes whose cup product with every class of $L_v$, given by the pairing $(m,\lambda)\mapsto\lambda(m)$, vanishes in $H^2(D_v,\mathbf Z/\ell^r(\chi_r))$.
Let $\operatorname{Sel}_L(M)$ be the submodule of classes in $H^1(G_{S'},M)$ whose restriction to $D_v$ lies in $L_v$ for every $v\in S'$, and define $\operatorname{Sel}_{L^\perp}(M^*(1))\subset H^1(G_{S'},M^*(1))$ in the same way.
Then $H^1(G_{S'},M)$ and $H^1(G_{S'},M^*(1))$ are finite, and
\[
|\operatorname{Sel}_L(M)|\,|H^0(G_{S'},M^*(1))|\prod_{v\in S'}|H^0(D_v,M)|\prod_{w\text{ real}}|M^{c_w}|\,|M|^{r_2}=|\operatorname{Sel}_{L^\perp}(M^*(1))|\,|H^0(G_{S'},M)|\prod_{v\in S'}|L_v|.
\]
\end{lemma}
\leanref{ReciprocityTheory/Hypotheses/GalCoh.lean}{ReciprocityTheory.Gal.GreenbergWiles}
\leanref{ReciprocityTheory/Hypotheses/GalCoh.lean}{ReciprocityTheory.Gal.GreenbergWilesAt}

\begin{lemma}[Hermite \cite{Neu99}]\label{sha-descent-link-hermite}
Let $n,N\ge1$ be integers and $E\subset\bar{\mathbf Q}$ a number field that contains a primitive $n$-th root of unity.
Then there are only finitely many continuous homomorphisms $\operatorname{Gal}(\bar{\mathbf Q}/E)\to\mathbf Z/n$ that vanish on the inertia group of every place of $\bar{\mathbf Q}$ not above a prime dividing $N$.
\end{lemma}
\leanref{Development/AutomorphicArithmetic/AutomorphicVarieties/GaloisLibrary/KummerTheory/FinChar.lean}{UpperTheory.V.Kum.finChar}
\begin{proof}
Such a homomorphism cuts out a cyclic extension of $E$ of degree dividing $n$, unramified at the places not above the primes dividing $N$.
The discriminants over $\mathbf Q$ of these extensions are bounded, so there are only finitely many of them by Hermite's theorem \cite[Chapter~III, \S2]{Neu99}, and each has finitely many homomorphisms of its Galois group to $\mathbf Z/n$.
\end{proof}

\begin{lemma}[Hilbert class field \cite{Neu99}]\label{sha-descent-link-hilbert}
Let $E\subset\bar{\mathbf Q}$ be a number field with class group $\operatorname{Cl}(E)$ and class number $h_E$, and let $\chi:\operatorname{Gal}(\bar{\mathbf Q}/E)\to C$ be a homomorphism to a commutative group, with open kernel.
Suppose that $\chi$ is trivial on the inertia group of every place of $\bar{\mathbf Q}$, and on every element of $\operatorname{Gal}(\bar{\mathbf Q}/E)$ that acts as complex conjugation for some embedding $\bar{\mathbf Q}\to\mathbf C$.
Then $\chi$ is trivial on the kernel of the Artin map $\operatorname{Gal}(\bar{\mathbf Q}/E)\to\operatorname{Cl}(E)$ of the Hilbert class field of $E$; in particular $\chi^{h_E}=1$.
\end{lemma}
\leanref{Development/AutomorphicArithmetic/AutomorphicVarieties/GaloisLibrary/KummerTheory/ArtKer.lean}{UpperTheory.V.Kum.artKer}
\leanref{Development/AutomorphicArithmetic/AutomorphicVarieties/GaloisLibrary/KummerTheory/ClPullApply.lean}{UpperTheory.V.Kum.artCl}

\begin{lemma}\label{sha-descent-link-cone}
Let $\Gamma$ be a profinite group, $V$ a finite set, and for $v\in V$ let $G_v$ be a profinite group with a continuous homomorphism $\mathrm{loc}_v:G_v\to\Gamma$.
Let $k$ be a finite field of characteristic $\ell$, $W$ a finite-dimensional $k$-vector space with a continuous linear action of $\Gamma$, and $b_v\subset W$ a subspace stable under $\mathrm{loc}_v(G_v)$ for each $v$; let $Q_v=W/b_v$, on which $G_v$ acts through $\mathrm{loc}_v$.
Define $H^0_b(\Gamma,W)$, $H^1_b(\Gamma,W)$ and $H^2_b(\Gamma,W)$ by (1) to (3) of Definition~\ref{sha-descent-residual}, for these maps $\mathrm{loc}_v$.
Suppose that $H^1(\Gamma,W)$, $H^2(\Gamma,W)$ and the $H^1(G_v,Q_v)$ are finite-dimensional, and let $X\subset H^2(\Gamma,W)$ be the subspace of classes whose image in $H^2(G_v,Q_v)$ vanishes for every $v$.
Then $H^1_b(\Gamma,W)$ and $H^2_b(\Gamma,W)$ are finite-dimensional, and
\[
\dim H^2_b(\Gamma,W)+\dim H^0_b(\Gamma,W)+\dim H^1(\Gamma,W)+\sum_v\dim H^0(G_v,Q_v)=\dim H^1_b(\Gamma,W)+\dim H^0(\Gamma,W)+\sum_v\dim H^1(G_v,Q_v)+\dim X.
\]
\end{lemma}
\leanref{Development/AutomorphicArithmetic/AutomorphicVarieties/GaloisCochains/ConeCount.lean}{UpperTheory.V.GC.coneCount}
\begin{proof}
The sequence
\[
0\to H^0_b\to H^0(\Gamma,W)\to\bigoplus_vH^0(G_v,Q_v)\to H^1_b\to H^1(\Gamma,W)\to\bigoplus_vH^1(G_v,Q_v)\to H^2_b\to H^2(\Gamma,W)
\]
is exact, and the image of its last map is $X$.
Its maps send $X\in W^\Gamma$ to $(X\bmod b_v)_v$; $(Y_v)$ to the class of $(0,(\tilde Y_v))$, for lifts $\tilde Y_v\in W$ of the $Y_v$; the class of $(x,(X_v))$ to the class of $x$; the class of $x$ to the classes of the cocycles $x\circ\mathrm{loc}_v\bmod b_v$; the classes of cocycles $a_v$ to the class of $(0,(\tilde a_v))$, for lifts $\tilde a_v:G_v\to W$ of the $a_v$; and the class of $(z,(y_v))$ to the class of $z$.
Each exactness, and the description of the image, is a direct check on cochains.
The dimension formula follows by adding the dimensions along the sequence.
\end{proof}

\begin{lemma}\label{sha-descent-link-inflation}
Let $p:\Gamma'\to\Gamma$ be a continuous surjective homomorphism of profinite groups, $T$ a finite set, and $D_q$ and $J_q$, for $q\in T$, subgroups of $\Gamma'$ such that $\ker p$ is the closed normal subgroup generated by the $J_q$.
Let $k$ be a finite field of characteristic $\ell$ and $M$ a finite-dimensional $k$-vector space with a continuous linear action of $\Gamma$, on which $\Gamma'$ acts through $p$.
For $q\in T$ let $\operatorname{Hom}_{D_q}(J_q,M)$ be the space of continuous homomorphisms $f:J_q\to M$ with $f(djd^{-1})=df(j)$ whenever $d\in D_q$ and $j,djd^{-1}\in J_q$.
Suppose that $H^1(\Gamma',M)$ and the spaces $\operatorname{Hom}_{D_q}(J_q,M)$ are finite-dimensional.
Then $H^1(\Gamma,M)$ and the kernel $\ker(\mathrm{inf})$ of the inflation map $H^2(\Gamma,M)\to H^2(\Gamma',M)$ are finite-dimensional, and
\[
\dim\ker(\mathrm{inf})+\dim H^1(\Gamma',M)\le\dim H^1(\Gamma,M)+\sum_{q\in T}\dim\operatorname{Hom}_{D_q}(J_q,M).
\]
\end{lemma}
\leanref{Development/AutomorphicArithmetic/AutomorphicVarieties/GaloisCochains/InflKer.lean}{UpperTheory.V.GC.inflKer}
\begin{proof}
Let $N=\ker p$, which acts trivially on $M$, and let $\operatorname{Hom}_{\Gamma'}(N,M)$ be the space of continuous homomorphisms $f:N\to M$ with $f(gng^{-1})=gf(n)$ for $g\in\Gamma'$ and $n\in N$.
The five-term exact sequence \cite[Chapter~I, \S6]{NSW08}
\[
0\to H^1(\Gamma,M)\to H^1(\Gamma',M)\to\operatorname{Hom}_{\Gamma'}(N,M)\to H^2(\Gamma,M)\to H^2(\Gamma',M)
\]
has inflation as its first and last maps, and restriction to $N$ as its second.
Restriction to the $J_q$ maps $\operatorname{Hom}_{\Gamma'}(N,M)$ injectively to $\bigoplus_q\operatorname{Hom}_{D_q}(J_q,M)$.
Indeed, the zero set of such an $f$ is a closed normal subgroup of $\Gamma'$; if it contains every $J_q$, it is $N$.
So $H^1(\Gamma,M)$, $\operatorname{Hom}_{\Gamma'}(N,M)$ and $\ker(\mathrm{inf})$ are finite-dimensional, and $\dim\ker(\mathrm{inf})+\dim H^1(\Gamma',M)=\dim H^1(\Gamma,M)+\dim\operatorname{Hom}_{\Gamma'}(N,M)$, which gives the inequality.
\end{proof}

\begin{lemma}\label{sha-descent-link-tame}
Let $\ell$ be a prime, $F\subset\bar{\mathbf Q}$ a number field and $G_F=\operatorname{Gal}(\bar{\mathbf Q}/F)$.
Let $q$ be a finite place of $F$ not above $\ell$, $A$ a place of $\bar{\mathbf Q}$ above $q$, and $D\supset I$ the decomposition and inertia groups of $A$ in $G_F$.
Let $K\subset G_F$ be an open normal subgroup of index prime to $\ell$, and $J=I\cap K$.
Let $k$ be a finite field of characteristic $\ell$ and $M$ a finite-dimensional $k$-vector space with a continuous linear action of $G_F$ that is trivial on $J$.
Let $\operatorname{Hom}_D(J,M)$ be the space of continuous homomorphisms $f:J\to M$ with $f(djd^{-1})=df(j)$ for $d\in D$ and $j\in J$.
Then $\operatorname{Hom}_D(J,M)$ is finite-dimensional, and $\dim\operatorname{Hom}_D(J,M)=\dim H^2(D,M)$.
\end{lemma}
\leanref{Development/AutomorphicArithmetic/AutomorphicVarieties/GaloisCochains/TameEq.lean}{UpperTheory.V.GC.tameEq}
\begin{proof}
The index of $J$ in $I$ divides the index of $K$, so it is prime to $\ell$, and $J$ is closed and normal in $D$.
Let $z:D\to M$ be a continuous cocycle.
Its restriction to $J$ is a homomorphism, since $J$ acts trivially on $M$.
For $d\in D$ and $j\in J$, comparing $z(djd^{-1}\cdot d)=z(djd^{-1})+z(d)$ with $z(dj)=z(d)+dz(j)$ gives $z(djd^{-1})=dz(j)$.
Coboundaries vanish on $J$.
So restriction to $J$ is a linear map $\mathrm{res}:H^1(D,M)\to\operatorname{Hom}_D(J,M)$.

The kernel of $\mathrm{res}$ consists of the classes whose cocycles vanish on $J$.
These are the classes whose restriction to $I$ is a coboundary.
Indeed, a cocycle $z$ that vanishes on $J$ factors on $I$ through $I/J$, whose order is prime to $\ell$, and averaging over $I/J$ writes it as a coboundary on $I$; conversely, coboundaries vanish on $J$.
By Lemma~\ref{sha-descent-link-unramified}(1), $D/I$ is procyclic, topologically generated by the image of $\varphi$, and maps onto every $\mathbf Z/n$, so it is isomorphic to $\hat{\mathbf Z}$.
So these classes form the image of the injective inflation map from $H^1(D/I,M^I)$, which has dimension $\dim M^I/(\varphi-1)M^I=\dim\ker(\varphi-1\mid M^I)=\dim H^0(D,M)$.

The map $\mathrm{res}$ is surjective.
Let $f\in\operatorname{Hom}_D(J,M)$.
Choose an open normal subgroup $U$ of $D$ that acts trivially on $M$ and with $f(U\cap J)=0$, and a set $R$ of representatives of $D/JU$, and put $c(rju)=rf(j)$ for $r\in R$, $j\in J$ and $u\in U$.
This defines a continuous map $c:D\to M$ with $c(gj)=c(g)+gf(j)$ for $g\in D$ and $j\in J$.
Its coboundary $dc(s,t)=sc(t)-c(st)+c(s)$ satisfies $dc(sj,tj')=dc(s,t)$ for $j,j'\in J$, by the equivariance of $f$.
By Lemma~\ref{sha-descent-link-unramified}(2) for $N=J$, $dc$ is the coboundary of a continuous map $b:D\to M$ with $b(sj)=b(s)$.
Then $x=c-b$ is a cocycle, and $x(j)=f(j)$ for $j\in J$, because $x(1)=0$, $c(j)=c(1)+f(j)$ and $b(j)=b(1)$.

Hence $\dim\operatorname{Hom}_D(J,M)=\dim H^1(D,M)-\dim H^0(D,M)$, which is $\dim H^2(D,M)$ by Lemma~\ref{sha-descent-link-local}(2), since $q$ does not lie above $\ell$.
\end{proof}

\begin{lemma}\label{sha-descent-link-generators}
Let $\ell$ be an odd prime, $F\subset\bar{\mathbf Q}$ a number field and $G_F=\operatorname{Gal}(\bar{\mathbf Q}/F)$.
Let $v$ be a place of $F$ above $\ell$, $d_v=[F_v:\mathbf Q_\ell]$, $A$ a place of $\bar{\mathbf Q}$ above $v$, and $D\supset I$ the decomposition and inertia groups of $A$ in $G_F$; suppose that the reduction $\omega$ of $\chi_{\mathrm{cyc}}$ modulo $\ell$ is nontrivial on $D$.
Let $K\subset G_F$ be an open normal subgroup of index prime to $\ell$.
Then there are $u_1,\dots,u_{d_v}\in I\cap K$, with $\chi_{\mathrm{cyc}}(u_i)=1$ for $i\ge2$, such that every continuous homomorphism from $D$ to a finite abelian $\ell$-group that kills $u_1,\dots,u_{d_v}$ kills $I$.
\end{lemma}
\leanref{Development/AutomorphicArithmetic/AutomorphicVarieties/GaloisCochains/GenDec.lean}{UpperTheory.V.GC.genDec}
\begin{proof}
Let $\varphi\in D$, and the homomorphisms $d_n:D\to\mathbf Z/n$ for $n\ge1$, be as in Lemma~\ref{sha-descent-link-unramified}(1).
Let $X$ be the $\mathbf F_\ell$-vector space of continuous homomorphisms $D\to\mathbf F_\ell$, which is $H^1(D,\mathbf F_\ell)$ for the trivial action.
A linear form $\lambda$ on $\mathbf F_\ell$ with $\lambda=\omega(g)\lambda$ for all $g\in D$ vanishes, since $\omega$ is nontrivial on $D$.
So Lemma~\ref{sha-descent-link-local} gives $H^2(D,\mathbf F_\ell)=0$ and $\dim X=\dim H^0(D,\mathbf F_\ell)+d_v=1+d_v$.

Let $J=I\cap K$, a closed subgroup of $D$, and $P$ the set of $g\in J$ with $\chi_{\mathrm{cyc}}(g)=1$, which contains $1$.
Let $Y\subset X$ be the subspace of the $f$ with $f(\varphi)=0$, and $Y_P\subset Y$ the subspace of the $f$ that vanish on $P$.
Since $d_\ell(\varphi)=1$, $Y\ne X$, so $\dim Y\le d_v$.
The values at the points of $P$ separate the elements of $Y/Y_P$, so there are $\dim Y-\dim Y_P$ points of $P$ such that every $f\in Y$ that vanishes at them lies in $Y_P$; adding the point $1$, any larger number of points will also do.

If $Y_P=0$, this gives $u_1,\dots,u_{d_v}\in P$ such that every $f\in Y$ that vanishes at all $u_i$ is $0$.

If $Y_P\ne0$, then $\dim Y-\dim Y_P\le d_v-1$, and this gives $u_2,\dots,u_{d_v}\in P$ such that every $f\in Y$ that vanishes at them lies in $Y_P$.
Let $N_P$ be the set of common zeros of the elements of the finite space $Y_P$; it is open, and it contains $P$.
There is $u_1\in J$ such that for every $g\in J$ some $m\ge0$ has $g^{-1}u_1^m\in N_P$.
Indeed, the closed sets $C_n$ of the $g\in J$ outside $N_P$ with $\chi_{\mathrm{cyc}}(g)\equiv1\pmod{\ell^n}$ decrease with $n$, and their intersection is empty, as $P\subset N_P$.
Since $D$ is compact, some $C_n$ is empty.
The group $(\mathbf Z/\ell^n)^\times$ is cyclic, because $\ell$ is odd, so the image of $J$ in it is generated by the image of some $u_1\in J$.
For $g\in J$ choose $m\ge0$ with $\chi_{\mathrm{cyc}}(g)\equiv\chi_{\mathrm{cyc}}(u_1)^m\pmod{\ell^n}$; then $g^{-1}u_1^m\in J$ does not lie in $C_n$, so it lies in $N_P$.
Now let $f\in Y$ vanish at $u_1,\dots,u_{d_v}$.
Then $f\in Y_P$, and for $g\in J$, with $m$ as above, $f(g)=mf(u_1)-f(g^{-1}u_1^m)=0$.
For $g\in I$ and $r$ the index of $K$, $g^r\in J$ because $K$ is normal, so $rf(g)=0$, and $f(g)=0$ because $\ell\nmid r$.
So $f$ vanishes on $I$ and at $\varphi$, and its kernel is closed, so $f=0$.

In both cases $u_1,\dots,u_{d_v}$ lie in $I\cap K$ and $\chi_{\mathrm{cyc}}(u_i)=1$ for $i\ge2$, and every continuous homomorphism $f:D\to\mathbf F_\ell$ with $f(\varphi)=0$ and $f(u_i)=0$ for all $i$ is $0$.

Let $\chi:D\to B$ be a continuous homomorphism to a finite abelian $\ell$-group of order $\ell^e$ with $\chi(u_i)=1$ for all $i$, and let $\theta=(\chi,d_{\ell^e}):D\to B\times\mathbf Z/\ell^e$.
Then $\theta(u_i)=(1,0)$, as $u_i\in I$.
The finite abelian $\ell$-group $\theta(D)$ is generated by $\theta(\varphi)$.
Otherwise there is a nonzero homomorphism $\theta(D)\to\mathbf F_\ell$ that kills $\theta(\varphi)$, and its composite with $\theta$ is a nonzero continuous homomorphism $D\to\mathbf F_\ell$ that vanishes at $\varphi$ and at the $u_i$, which is impossible.
Let $g\in I$, and write $\theta(g)=\theta(\varphi)^m$.
The second components give $0=d_{\ell^e}(g)=m$ in $\mathbf Z/\ell^e$, so $\ell^e$ divides $m$, and $\chi(g)=\chi(\varphi)^m=1$.
\end{proof}

\begin{lemma}\label{sha-descent-link-strict}
Let $\ell$ be a prime, $F\subset\bar{\mathbf Q}$ a number field, $G_F=\operatorname{Gal}(\bar{\mathbf Q}/F)$, and for a place $A$ of $\bar{\mathbf Q}$ let $I_A$ be its inertia group in $G_F$.
Let $k$ be a field of characteristic $\ell$.
Let $U\subset G_F$ be a subgroup and $K_1\subset U$ an open subgroup, normal in $G_F$, of index prime to $\ell$ in $U$.
Let $r_\ell(K_1)$ be the dimension over $\mathbf F_\ell$ of the space of continuous homomorphisms $K_1\to\mathbf Z/\ell$ that vanish on $K_1\cap I_A$ for every place $A$ of $\bar{\mathbf Q}$.
Let $\Gamma$ be a topological group, $\pi:U\to\Gamma$ a continuous surjective homomorphism, $V$ a finite set, $G_v$ for $v\in V$ groups with homomorphisms $\mathrm{loc}_v:G_v\to\Gamma$, and $\kappa:\Gamma\to k^\times$ a homomorphism with $\kappa(\pi(K_1))=1$.
Suppose that for every place $A$ and every $u\in I_A\cap K_1$, either $\pi(u)=1$ or $\pi(u)=g\,\mathrm{loc}_v(s)g^{-1}$ for some $v\in V$, $s\in G_v$ and $g\in\Gamma$.
Let $H^1_\emptyset(\Gamma,k(\kappa))$ be the quotient of the space of pairs $(x,(X_v))$, of a locally constant map $x:\Gamma\to k$ with $x(st)=x(s)+\kappa(s)x(t)$ and elements $X_v\in k$ with $x(\mathrm{loc}_vh)=\kappa(\mathrm{loc}_vh)X_v-X_v$ for all $v$ and $h\in G_v$, by the pairs $(s\mapsto\kappa(s)X-X,(X)_v)$ with $X\in k$.
Then
\[
\dim H^1_\emptyset(\Gamma,k(\kappa))\le r_\ell(K_1)+\#\{v\in V:\kappa\circ\mathrm{loc}_v=1\}.
\]
\end{lemma}
\leanref{Development/AutomorphicArithmetic/AutomorphicVarieties/GaloisCochains/StrGen.lean}{UpperTheory.V.GC.strGen}
\begin{proof}
Let $(x,(X_v))$ be a pair as in the statement.
Since $\kappa(\pi(K_1))=1$, the map $u\mapsto x(\pi(u))$ is a continuous homomorphism $K_1\to k$.
It vanishes on $K_1\cap I_A$ for every place $A$.
Indeed, let $u\in I_A\cap K_1$; if $\pi(u)=1$, then $x(\pi(u))=x(1)=0$.
Otherwise $\pi(u)=g\,\mathrm{loc}_v(s)g^{-1}$, and $y=\mathrm{loc}_v(s)$ has $\kappa(y)=\kappa(\pi(u))=1$, so $x(y)=\kappa(y)X_v-X_v=0$.
Then $x(gyg^{-1})=x(g)+\kappa(g)x(y)+\kappa(gy)x(g^{-1})=x(g)+\kappa(g)x(g^{-1})=x(1)=0$.
The pairs of the form $(s\mapsto\kappa(s)X-X,(X)_v)$ give $0$, since $\kappa(\pi(K_1))=1$.
So $u\mapsto x(\pi(u))$ defines a linear map $\Phi$ from $H^1_\emptyset(\Gamma,k(\kappa))$ to the $k$-vector space of continuous homomorphisms $K_1\to k$ that vanish on the groups $K_1\cap I_A$.

This space has dimension at most $r_\ell(K_1)$ over $k$.
Indeed, $K_1$ is compact, so such a homomorphism takes finitely many values, which span a finite-dimensional $\mathbf F_\ell$-subspace of $k$; its coordinates in a basis of that subspace are continuous homomorphisms $K_1\to\mathbf F_\ell$ that vanish on the groups $K_1\cap I_A$.
So the space is spanned over $k$ by homomorphisms with values in $\mathbf F_\ell$.

Let $(x,(X_v))$ represent a class in the kernel of $\Phi$, so that $x$ vanishes on $N=\pi(K_1)$.
The subgroup $N$ is normal in $\Gamma$, of finite index prime to $\ell$, since its index divides that of $K_1$ in $U$.
Then $x(sn)=x(s)$ for $n\in N$, and summing $\kappa(s)x(q)=x(sq)-x(s)$ over a set of representatives $q$ of $\Gamma/N$ shows that $x(s)=\kappa(s)X-X$ with $X=-[\Gamma:N]^{-1}\sum_qx(q)$.
So the class is that of $(0,(X_v-X))$.
For $v$ with $\kappa\circ\mathrm{loc}_v\ne1$, choose $h\in G_v$ with $\kappa(\mathrm{loc}_vh)\ne1$; then $(\kappa(\mathrm{loc}_vh)-1)X_v$ and $(\kappa(\mathrm{loc}_vh)-1)X$ both equal $x(\mathrm{loc}_vh)$, so $X_v=X$.
Hence the kernel of $\Phi$ is spanned by the classes of the pairs $(0,(\delta_{vv'})_{v'})$ for the $v$ with $\kappa\circ\mathrm{loc}_v=1$, and the bound follows.
\end{proof}

\begin{situation}\label{sha-descent-link-setting}
Let $\ell$ be an odd prime, $F\subset\bar{\mathbf Q}$ a number field and $G_F=\operatorname{Gal}(\bar{\mathbf Q}/F)$; for a place $A$ of $\bar{\mathbf Q}$ let $D_A\supset I_A$ be its decomposition and inertia groups in $G_F$, and let $\omega:G_F\to\mathbf F_\ell^\times$ be the reduction of $\chi_{\mathrm{cyc}}$ modulo $\ell$.
Let $S$ be a finite set of finite places of $F$ that contains the places above $\ell$, and $K\subset G_F$ an open normal subgroup of index prime to $\ell$.
Let $\Gamma$ be a profinite group and $\pi:G_F\to\Gamma$ a continuous surjective homomorphism whose kernel is the closed normal subgroup generated by the groups $I_A\cap K$ for the places $A$ of $\bar{\mathbf Q}$ above finite places of $F$ outside $S$.
When $\Gamma$ acts on a vector space, $G_F$ and its subgroups act on it through $\pi$.
\end{situation}
\leanref{Development/AutomorphicArithmetic/AutomorphicVarieties/GaloisCochains/GLinkQuot.lean}{UpperTheory.V.GC.GLink}
\leanref{Development/AutomorphicArithmetic/AutomorphicVarieties/GaloisCochains/GLinkQuot.lean}{UpperTheory.V.GC.tameN}

\begin{lemma}\label{sha-descent-link-finite}
In Situation~\ref{sha-descent-link-setting}, every open subgroup of $\Gamma$ has only finitely many continuous homomorphisms to $\mathbf Z/\ell$.
\end{lemma}
\leanref{Development/AutomorphicArithmetic/AutomorphicVarieties/GaloisCochains/PhiGam.lean}{UpperTheory.V.GC.phiGam}
\begin{proof}
Let $H\subset\Gamma$ be an open subgroup, $\zeta\in\bar{\mathbf Q}$ a primitive $\ell$-th root of unity, and $N\ge1$ the product of the norms of the places of $S$, so that every place of $S$ lies above a prime dividing $N$.
The subgroup $V=\pi^{-1}(H)\cap K\cap\{g\in G_F:g\zeta=\zeta\}$ is open in $G_F$, so its fixed field $E$ is a number field, which contains $\zeta$, and $\operatorname{Gal}(\bar{\mathbf Q}/E)=V$.
Let $e:V\to H$ be the restriction of $\pi$.
Let $A$ be a place of $\bar{\mathbf Q}$ not above a prime dividing $N$.
Then $A$ lies above a place of $F$ outside $S$, so $I_A\cap V\subset I_A\cap K\subset\ker\pi$.
So for every continuous homomorphism $\phi:H\to\mathbf Z/\ell$, the composite $\phi\circ e$ vanishes on the inertia groups of $V$ at the places not above primes dividing $N$, and it is one of the finitely many homomorphisms of Lemma~\ref{sha-descent-link-hermite}.
The image $e(V)=\pi(V)$ has finite index in $H$, because $V$ has finite index in $G_F$ and $\pi$ is surjective.
A homomorphism $\phi$ is determined by $\phi\circ e$ and by its values on a set of representatives of $H/e(V)$, so there are only finitely many.
\end{proof}

\begin{lemma}\label{sha-descent-link-level}
In Situation~\ref{sha-descent-link-setting}, there are a finite set $S'\supset S$ of finite places of $F$ and a continuous surjective homomorphism $p:G_{S'}\to\Gamma$, where $G_{S'}$ is the Galois group over $F$ of the maximal extension of $F$ in $\bar{\mathbf Q}$ unramified at every finite place outside $S'$, with the following properties: $\pi$ is the composite of the quotient map $G_F\to G_{S'}$ with $p$; and for every choice of places $A_q$ of $\bar{\mathbf Q}$ above the places $q\in S'\setminus S$, the kernel of $p$ is the closed normal subgroup generated by the images of the groups $I_{A_q}\cap K$.
\end{lemma}
\leanref{Development/AutomorphicArithmetic/AutomorphicVarieties/GaloisCochains/LevelS.lean}{UpperTheory.V.GC.levelS}
\begin{proof}
The fixed field of $K$ is a finite extension of $F$, so there is a finite set $T$ of finite places of $F$ such that $I_A\subset K$ for every place $A$ above a place outside $T$.
Let $S'=S\cup T$.
The kernel of $G_F\to G_{S'}$ is the closed normal subgroup generated by the groups $I_A$ for the places $A$ above finite places outside $S'$, and for these $I_A=I_A\cap K\subset\ker\pi$.
So $\pi$ is the composite of $G_F\to G_{S'}$ with a continuous surjection $p$.
The images of the groups $I_{A_q}\cap K$ lie in $\ker p$, since $q\notin S$.
Conversely, let $N'$ be the closed normal subgroup generated by these images.
Every element of $\ker p$ is the image of an element of $\ker\pi$, so it suffices that the inverse image of $N'$ in $G_F$, a closed normal subgroup, contains the groups $I_A\cap K$ for the places $A$ above finite places $q\notin S$.
If $q\notin S'$, then $I_A$ maps to $1$ in $G_{S'}$.
If $q\in S'\setminus S$, then $A=uA_q$ for some $u\in G_F$, and $I_A\cap K=u(I_{A_q}\cap K)u^{-1}$ because $K$ is normal.
\end{proof}

\begin{lemma}\label{sha-descent-link-h2}
In Situation~\ref{sha-descent-link-setting}, suppose that $F$ is totally real.
For each $v\in S$ let $A_v$ be a place of $\bar{\mathbf Q}$ above $v$, and for each infinite place $w$ of $F$ let $c_w\in G_F$ be a complex conjugation at $w$.
Let $k$ be a finite field of characteristic $\ell$ and $M$ a finite-dimensional $k$-vector space with a continuous linear action of $\Gamma$.
Let $\text{Ш}^2_S(\Gamma,M)$ be the kernel of the restriction maps $H^2(\Gamma,M)\to H^2(D_{A_v},M)$ along $\pi$, for $v\in S$, and $\operatorname{Hom}_{G_F}(M,k(\omega))$ the space of linear forms $\lambda$ on $M$ with $\lambda(gm)=\omega(g)\lambda(m)$ for $g\in G_F$.
Then $H^1(\Gamma,M)$ and $H^2(\Gamma,M)$ are finite-dimensional, and
\[
\dim\text{Ш}^2_S(\Gamma,M)+\sum_{v\in S}\dim H^2(D_{A_v},M)+\dim H^0(\Gamma,M)+[F:\mathbf Q]\dim M\le\dim H^1(\Gamma,M)+\sum_w\dim M^{c_w}+\dim\operatorname{Hom}_{G_F}(M,k(\omega)).
\]
\end{lemma}
\leanref{Development/AutomorphicArithmetic/AutomorphicVarieties/GaloisCochains/H2GaloisQuotientEq.lean}{UpperTheory.V.GC.H2_galois_quotient_eq}
\begin{proof}
Let $S'$ and $p:G_{S'}\to\Gamma$ be as in Lemma~\ref{sha-descent-link-level}, choose places $A_q$ above the places $q\in S'\setminus S$, and for $v\in S'$ let $n_v=\dim H^2(D_{A_v},M)$, which is finite by Lemma~\ref{sha-descent-link-local}.

Let $q\in S'\setminus S$; it does not lie above $\ell$, and $I_{A_q}\cap K\subset\ker\pi$ acts trivially on $M$.
Let $D_q$ and $J_q$ be the images of $D_{A_q}$ and $I_{A_q}\cap K$ in $G_{S'}$.
Composing with $G_F\to G_{S'}$ embeds $\operatorname{Hom}_{D_q}(J_q,M)$, in the notation of Lemma~\ref{sha-descent-link-inflation}, in the space $\operatorname{Hom}_{D_{A_q}}(I_{A_q}\cap K,M)$, whose dimension is $n_q$ by Lemma~\ref{sha-descent-link-tame}.
By Lemma~\ref{sha-descent-link-global}, $H^1(G_{S'},M)$ and $H^2(G_{S'},M)$ are finite-dimensional.
Lemma~\ref{sha-descent-link-inflation}, applied to $p$ and the groups $D_q$ and $J_q$, shows that $H^1(\Gamma,M)$ and the kernel $\ker(\mathrm{inf})$ of the inflation $H^2(\Gamma,M)\to H^2(G_{S'},M)$ are finite-dimensional, and that
\[
\dim\ker(\mathrm{inf})+\dim H^1(G_{S'},M)\le\dim H^1(\Gamma,M)+\sum_{q\in S'\setminus S}n_q.\tag{a}
\]
Inflation maps $\text{Ш}^2_S(\Gamma,M)$ into the group $\text{Ш}^2_{S'}(M)$ of Lemma~\ref{sha-descent-link-global}.
Indeed, the restriction to $D_{A_v}$ of the inflation of a class $x$ is the restriction of $x$ along $\pi$.
For $v\in S$ it vanishes when $x\in\text{Ш}^2_S(\Gamma,M)$.
For $q\in S'\setminus S$, the map $D_{A_q}\to\Gamma$ is trivial on $I_{A_q}\cap K$, so the restriction of any class is represented by a cocycle inflated from $D_{A_q}/(I_{A_q}\cap K)$, and it vanishes by Lemma~\ref{sha-descent-link-unramified}(2).
Since $\ker(\mathrm{inf})$ and $H^2(G_{S'},M)$ are finite-dimensional, so is $H^2(\Gamma,M)$, and
\[
\dim\text{Ш}^2_S(\Gamma,M)\le\dim\ker(\mathrm{inf})+\dim\text{Ш}^2_{S'}(M).\tag{b}
\]
Since $p$ is surjective, $G_{S'}$ and $\Gamma$ have the same invariants in $M$, and the linear forms of Lemma~\ref{sha-descent-link-global} for $G_{S'}$ are those of the statement.
So Lemma~\ref{sha-descent-link-global} gives
\[
\dim\text{Ш}^2_{S'}(M)+\sum_{v\in S'}n_v=\dim H^2(G_{S'},M)+\dim\operatorname{Hom}_{G_F}(M,k(\omega)),\tag{c}
\]
\[
\dim H^1(G_{S'},M)+\sum_w\dim M^{c_w}=\dim H^0(\Gamma,M)+\dim H^2(G_{S'},M)+[F:\mathbf Q]\dim M.\tag{d}
\]
Adding (a), (b), (c) and (d) gives the inequality.
\end{proof}

\begin{lemma}\label{sha-descent-link-character}
In Situation~\ref{sha-descent-link-setting}, suppose that $F$ is totally real.
For each $v\in S$ let $A_v$ be a place of $\bar{\mathbf Q}$ above $v$, and for each infinite place $w$ of $F$ let $c_w\in G_F$ be a complex conjugation at $w$.
Let $k$ be a finite field of characteristic $\ell$, and $\kappa:\Gamma\to k^\times$ a homomorphism with $\kappa(\pi(g))=1$ for $g\in K\cap\ker\omega$.
Let $C_0\ge r_\ell(K\cap\ker\omega)$, with $r_\ell$ as in Lemma~\ref{sha-descent-link-strict}.
Let $\epsilon$ be $1$ if $\kappa=1$ and $0$ otherwise, and $S_\kappa$ the set of $v\in S$ with $\kappa(\pi(g))=\omega(g)$ for all $g\in D_{A_v}$.
Then
\[
\dim H^1(\Gamma,k(\kappa))+\#\{w:\kappa(\pi(c_w))=1\}\le\epsilon+C_0+[F:\mathbf Q]+\#S_\kappa.
\]
\end{lemma}
\leanref{Development/AutomorphicArithmetic/AutomorphicVarieties/GaloisCochains/CharacterH1GaloisQuotient.lean}{UpperTheory.V.GC.character_H1_galois_quotient}
\leanref{Development/AutomorphicArithmetic/AutomorphicVarieties/GaloisCochains/GwTame.lean}{UpperTheory.V.GC.gwTame}
\begin{proof}
Let $K'=K\cap\ker\omega$, an open normal subgroup of $G_F$ whose index is prime to $\ell$, since the index of $\ker\omega$ divides $\ell-1$.
Let $S'$ and $p:G_{S'}\to\Gamma$ be as in Lemma~\ref{sha-descent-link-level}, and choose places $A_v$ above the places $v\in S'\setminus S$.
The character $\omega$ factors through $G_{S'}$, as $S'$ contains the places above $\ell$.
Let $\kappa'=(\kappa\circ p)^{-1}\omega:G_{S'}\to k^\times$, which is trivial on the image of $K'$, and let $M$ be $k$ with $G_{S'}$ acting by $\kappa'$, viewed as an $\mathbf F_\ell$-vector space.
The trace form identifies $k$ with $M^*(1)$, in the notation of Lemma~\ref{sha-descent-link-greenberg-wiles}, by $a\mapsto(m\mapsto\operatorname{Tr}_{k/\mathbf F_\ell}(am))$, and $g$ acts on the image of $a$ as on the image of $\kappa(p(g))a$; so $M^*(1)$ is $k(\kappa\circ p)$.
For $v\in S$ let $L_v=0$, and for $v\in S'\setminus S$ let $L_v\subset H^1(D_{A_v},M)$ be the space of classes whose cocycles vanish on $I_{A_v}\cap K'$.

(i) Inflation along $p$, followed by this identification, is an injective map $H^1(\Gamma,k(\kappa))\to H^1(G_{S'},M^*(1))$, since $p$ is surjective.
Its image lies in $\operatorname{Sel}_{L^\perp}(M^*(1))$.
For $v\in S$ this is because $L_v=0$.
Let $v\in S'\setminus S$, let $y$ be a cocycle in the image, and $x$ a cocycle representing a class of $L_v$.
Then $y$ vanishes on $I_{A_v}\cap K'$, because $\pi$ kills $I_{A_v}\cap K$, and so does $x$.
So the cup product of $y$ and $x$ is a $2$-cocycle with values in $\mathbf F_\ell(\omega)$ inflated from $D_{A_v}/(I_{A_v}\cap K')$, on which $I_{A_v}\cap K'$ acts trivially, and it is a coboundary by Lemma~\ref{sha-descent-link-unramified}(2).

(ii) For $v\in S'\setminus S$, $|H^0(D_{A_v},M)|\le|L_v|$.
Indeed, $L_v$ is the kernel of restriction from $H^1(D_{A_v},M)$ to $\operatorname{Hom}_{D_{A_v}}(I_{A_v}\cap K',M)$, as $I_{A_v}\cap K'$ acts trivially on $M$.
By Lemma~\ref{sha-descent-link-tame} and Lemma~\ref{sha-descent-link-local}(2), at the place $v$ not above $\ell$, $\dim H^1(D_{A_v},M)\le\dim L_v+\dim H^2(D_{A_v},M)=\dim L_v+\dim H^1(D_{A_v},M)-\dim H^0(D_{A_v},M)$.

(iii) Every class of $\operatorname{Sel}_L(M)$ is represented by a cocycle on $G_{S'}$ that vanishes on the images of the groups $I_{A_v}\cap K'$ for all $v\in S'$: for $v\in S$ because its restriction to $D_{A_v}$ is a coboundary and $K'$ acts trivially on $M$, and for $v\in S'\setminus S$ by the definition of $L_v$.
There are at most $|k|^{C_0}$ such classes.
Indeed, let $z$ be such a cocycle, viewed on $G_F$.
The set of $g\in K'$ with $z(g)=0$ is a subgroup, normal in $G_F$ because $\kappa'$ is trivial on $K'$.
So $z$ vanishes on $I_B\cap K'$ for every place $B$ of $\bar{\mathbf Q}$: if $B$ lies above $v\in S'$, then $B=uA_v$ with $u\in G_F$ and $I_B\cap K'=u(I_{A_v}\cap K')u^{-1}$, and if $B$ lies above a finite place outside $S'$, then $I_B$ maps to $1$ in $G_{S'}$.
Let $N$ be the normal subgroup of $G_F$ generated by the groups $I_B\cap K'$, for all places $B$.
The cocycles $z$ factor through $G_F/N$, and classes that differ on $G_{S'}$ still differ on $G_F/N$.
Lemma~\ref{sha-descent-link-strict}, applied with $U=G_F$, $K_1=K'$, the quotient $G_F/N$, no maps $\mathrm{loc}_v$, and the character of $G_F/N$ induced by $\kappa'$, bounds the dimension of the space of these classes on $G_F/N$ by $r_\ell(K')\le C_0$.

(iv) $|H^0(G_{S'},M^*(1))|\le|k|^\epsilon$, since a nonzero invariant of $k(\kappa\circ p)$ forces $\kappa\circ p=1$, and so $\kappa=1$.

(v) For $v\in S$, $|H^0(D_{A_v},M)|$ is $|k|$ if $\kappa'$ is trivial on $D_{A_v}$, that is, if $v\in S_\kappa$, and $1$ otherwise.

(vi) For each $w$, $\kappa'(c_w)=-\kappa(\pi(c_w))^{-1}$, as $\omega(c_w)=-1$.
If $\kappa(\pi(c_w))=1$, then $\kappa'(c_w)=-1\ne1$, as $\ell$ is odd, and $M^{c_w}=0$; otherwise $|M^{c_w}|\le|k|$.

By (i) and Lemma~\ref{sha-descent-link-greenberg-wiles} for $M$ and the $L_v$, with $|H^0(G_{S'},M)|\ge1$, $r_2=0$ and (ii),
\[
|H^1(\Gamma,k(\kappa))|\le|\operatorname{Sel}_{L^\perp}(M^*(1))|\le|\operatorname{Sel}_L(M)|\,|H^0(G_{S'},M^*(1))|\prod_{v\in S}|H^0(D_{A_v},M)|\prod_w|M^{c_w}|.
\]
By (iii) to (vi) the right side is at most $|k|^{C_0+\epsilon+\#S_\kappa+n}$, where $n$ is the number of $w$ with $\kappa(\pi(c_w))\ne1$.
Since $F$ is totally real, $n+\#\{w:\kappa(\pi(c_w))=1\}=[F:\mathbf Q]$, and the bound follows.
\end{proof}

\begin{lemma}\label{sha-descent-link-cohomology}
Let a realized link of a descent datum to an abelian CM field $L$ be given, with the notation of Definitions~\ref{sha-descent-datum}, \ref{sha-descent-link} and~\ref{sha-descent-residual}, and let $C_0\ge r_\ell(H_{\eta,\Psi}\cap\ker\omega)$, with $r_\ell$ as in Lemma~\ref{sha-descent-auxiliary}(5).
Then the link satisfies conditions (5), (9), (10) and (11) of Definition~\ref{sha-descent-residual}.
If $\mathbf F$ is finite, then $\dim H^1_{\mathrm{all}}(\Gamma,\mathbf F)\le1+C_0$.
If $\bar\Psi_\Gamma^2$ is nontrivial on $\mathrm{loc}_v(G_v)$ for every $v$, then $\dim H^1_\emptyset(\Gamma,\mathbf F(\bar\Psi_\Gamma^{-2}\delta_\Gamma))\le C_0$ and $\dim H^1_\emptyset(\Gamma_0,\mathbf F(\bar\Psi_\Gamma^{-1}))\le C_0$.
\end{lemma}
\leanref{Development/AutomorphicArithmetic/AutomorphicVarieties/GaloisCochains/Imports.lean}{UpperTheory.V.GC.phi}
\leanref{Development/AutomorphicArithmetic/AutomorphicVarieties/GaloisCochains/ToZModEqOne.lean}{UpperTheory.V.GC.gen}
\leanref{Development/AutomorphicArithmetic/AutomorphicVarieties/GaloisCochains/EqOneOfPow.lean}{UpperTheory.V.GC.cl}
\leanref{Development/AutomorphicArithmetic/AutomorphicVarieties/GaloisCochains/ResH1SurjectiveOfIsomorphism.lean}{UpperTheory.V.GC.ec}
\leanref{Development/AutomorphicArithmetic/AutomorphicVarieties/GaloisCochains/H2GaloisQuotientSurjective.lean}{UpperTheory.V.GC.H2_galois_quotient_surjective}
\leanref{Development/AutomorphicArithmetic/AutomorphicVarieties/GaloisCochains/B1p.lean}{UpperTheory.V.GC.b1p}
\leanref{Development/AutomorphicArithmetic/AutomorphicVarieties/GaloisCochains/PsiLocalPairEqLocal.lean}{UpperTheory.V.GC.b1}
\leanref{Development/AutomorphicArithmetic/AutomorphicVarieties/GaloisCochains/CastToZModPowOne.lean}{UpperTheory.V.GC.h1}
\begin{proof}
% (5) is GC.phi (phiGam), (9) GC.gen (genDec), (10) GC.cl, (11) GC.ec (H2_galois_quotient_surjective, coneCount, locEuler, sum_deg), (12) GC.h1 (character_H1_galois_quotient) and GC.b1p, GC.b1 (strGen); each is stated for a link GC.FLink.
Let $F=L^+$, which is totally real, $G_F=\operatorname{Gal}(\bar{\mathbf Q}/L^+)$, $S=V$ the set of places of $L^+$ above $\ell$, and $K=H_{\eta,\Psi}$.
Since the link is realized, $K$ is open and normal in $G_F$ of index prime to $\ell$, and by (i) of Definition~\ref{sha-descent-link} the kernel of $p$ is the closed normal subgroup generated by the groups $I_A\cap K$ for the places $A$ not above $\ell$.
So $\ell$, $F$, $S$, $K$, $\Gamma$ and $\pi=p$ are as in Situation~\ref{sha-descent-link-setting}.
For $v\in V$, $G_v$ and $I_v$ are the decomposition and inertia groups of $A_v$ in $G_F$, and $\mathrm{loc}_v$ is the restriction of $p$; the family $(c_w)$ is $(p(\tau_w))$, for complex conjugations $\tau_w\in G_F$ at the infinite places $w$ of $L^+$.
Since $\delta_\Gamma\circ p=\eta$, $\bar\Psi_\Gamma\circ p=\bar\Psi$ and $\varepsilon_\Gamma\circ p=\chi_{\mathrm{cyc}}$, the group $p(K)$ lies in $\Gamma_0\cap\ker\bar\Psi_\Gamma$, $\omega_\Gamma\circ p=\omega$, and $\operatorname{Gal}(\bar{\mathbf Q}/L)=p^{-1}(\Gamma_0)$.
By (iv) of Definition~\ref{sha-descent-link}, every place $A$ of $\bar{\mathbf Q}$ above $\ell$ is $\sigma A_v$ for some $v$ and some $\sigma\in G_F$, and then $\sigma^{-1}I_A\sigma\cap G_F=I_v$ and $\sigma^{-1}D_A\sigma\cap G_F=G_v$.

Condition (5) is Lemma~\ref{sha-descent-link-finite}.

Condition (9). Let $v\in V$.
The character $\omega$ is nontrivial on $G_v$ by (iv) of Definition~\ref{sha-descent-link}, and $d_v$ is the local degree of $L^+$ at $v$.
Lemma~\ref{sha-descent-link-generators}, for $A_v$ and $K$, gives $u_1,\dots,u_{d_v}\in I_v\cap K$ with $\chi_{\mathrm{cyc}}(u_i)=1$ for $i\ge2$, such that every continuous homomorphism from $G_v$ to a finite abelian $\ell$-group that kills the $u_i$ kills $I_v$.
Then $\bar\Psi_\Gamma(\mathrm{loc}_vu_i)=\bar\Psi(u_i)=1$, and $\varepsilon_\Gamma(\mathrm{loc}_vu_i)=\chi_{\mathrm{cyc}}(u_i)=1$ for $i\ge2$, so $u_i\in I'_v$ for $i\ge2$.

Condition (10). Let $h_L$ be the class number of $L$, and $h=2h_L$.
Let $\chi:\Gamma_0\to B$ be a continuous homomorphism to a finite abelian $\ell$-group that kills $\mathrm{loc}_v(I_v)$ and $c_\Gamma\mathrm{loc}_v(I_v)c_\Gamma^{-1}$ for every $v$.
First, $\chi(p(\tau))=1$ for every place $A$ of $\bar{\mathbf Q}$ and every $\tau\in I_A\cap\operatorname{Gal}(\bar{\mathbf Q}/L)$.
If $A$ lies above $\ell$, write $A=\sigma A_v$ as above; then $s=\sigma^{-1}\tau\sigma\in I_v$ and $p(\tau)=p(\sigma)\mathrm{loc}_v(s)p(\sigma)^{-1}$.
If $\delta_\Gamma(p(\sigma))=1$, then $\chi(p(\tau))=\chi(\mathrm{loc}_v(s))=1$, since $B$ is commutative.
Otherwise $g=c_\Gamma^{-1}p(\sigma)$ lies in $\Gamma_0$, and $p(\tau)$ is the conjugate of $c_\Gamma\mathrm{loc}_v(s)c_\Gamma^{-1}$ by $c_\Gamma gc_\Gamma^{-1}\in\Gamma_0$, so again $\chi(p(\tau))=1$.
If $A$ does not lie above $\ell$, let $n$ be the index of $K$ in $G_F$, which is prime to $\ell$; then $\tau^n\in I_A\cap K\subset\ker p$, so $\chi(p(\tau))^n=1$, and $\chi(p(\tau))=1$ because $B$ is an $\ell$-group.
Now let $\psi:\operatorname{Gal}(\bar{\mathbf Q}/L)\to B$ be $\tau\mapsto\chi(p(\tau))^2$.
It has open kernel and is trivial on every inertia group, by the above.
It is trivial on every element $\sigma$ that acts as complex conjugation for some embedding $\bar{\mathbf Q}\to\mathbf C$, since $\sigma^2=1$.
By Lemma~\ref{sha-descent-link-hilbert}, $\psi^{h_L}=1$.
Every $g\in\Gamma_0$ is $p(\tau)$ for some $\tau\in\operatorname{Gal}(\bar{\mathbf Q}/L)$, so $\chi(g)^h=\psi(\tau)^{h_L}=1$.

Condition (11). Let $\mathbf F'$, $W$ and $b$ be as in (11), and $Q_v=W/b_v$, on which $G_v$ acts through $\mathrm{loc}_v$.
A linear form $\lambda$ on $W$ with $\lambda(p(g)m)=\omega(g)\lambda(m)$ for $g\in G_F$ is a $\Gamma$-equivariant map $W\to\mathbf F'(\omega_\Gamma)$, as $p$ is surjective, so it is $0$.
Lemma~\ref{sha-descent-link-h2}, for the places $A_v$ and the $\tau_w$, therefore shows that $H^1(\Gamma,W)$ and $H^2(\Gamma,W)$ are finite-dimensional and that
\[
\dim\text{Ш}^2_V(\Gamma,W)+\sum_v\dim H^2(G_v,W)+\dim H^0(\Gamma,W)+[L^+:\mathbf Q]\dim W\le\dim H^1(\Gamma,W)+\sum_w\dim\ker(c_w-1\mid W).
\]
Let $X\subset H^2(\Gamma,W)$ be the subspace of classes whose image in $H^2(G_v,Q_v)$ vanishes for every $v$.
Restriction maps $X$ to $\bigoplus_v\ker\big(H^2(G_v,W)\to H^2(G_v,Q_v)\big)$, with kernel $\text{Ш}^2_V(\Gamma,W)$, and the maps $H^2(G_v,W)\to H^2(G_v,Q_v)$ are surjective by Lemma~\ref{sha-descent-link-local}(3).
So
\[
\dim X+\sum_v\dim H^2(G_v,Q_v)\le\dim\text{Ш}^2_V(\Gamma,W)+\sum_v\dim H^2(G_v,W).
\]
By Lemma~\ref{sha-descent-link-local}, at the places $v$ above $\ell$, $H^1(G_v,Q_v)$ is finite-dimensional and
\[
\dim H^1(G_v,Q_v)=\dim H^0(G_v,Q_v)+\dim H^2(G_v,Q_v)+d_v(\dim W-\dim b_v).
\]
By Lemma~\ref{sha-descent-link-cone}, $H^1_b(\Gamma,W)$ and $H^2_b(\Gamma,W)$ are finite-dimensional, and
\[
\dim H^2_b+\dim H^0_b+\dim H^1(\Gamma,W)+\sum_v\dim H^0(G_v,Q_v)=\dim H^1_b+\dim H^0(\Gamma,W)+\sum_v\dim H^1(G_v,Q_v)+\dim X.
\]
Finally $\sum_vd_v\le[L^+:\mathbf Q]$, as the $d_v$ are the local degrees at distinct places above $\ell$.
Combining the last five relations gives
\[
\dim H^2_b+\dim H^0_b+\sum_vd_v\dim b_v\le\dim H^1_b+\sum_w\dim\ker(c_w-1\mid W)+\Big(\sum_vd_v-[L^+:\mathbf Q]\Big)\dim W,
\]
and the last term is at most $0$.

Condition (12). Let $K_1=K\cap\ker\omega$.
It is open and normal in $G_F$, of index prime to $\ell$, since the index of $\ker\omega$ divides $\ell-1$.
The number $r_\ell(K_1)$ of Lemma~\ref{sha-descent-link-strict} is that of Lemma~\ref{sha-descent-auxiliary}(5), since the inertia groups of $G_F$ are those of $G_{\mathbf Q}$ intersected with $G_F$; so $r_\ell(K_1)\le C_0$.
Let $A$ be a place of $\bar{\mathbf Q}$ and $u\in I_A\cap K_1$.
If $A$ lies above $\ell$, write $A=\sigma A_v$ as above; then $s=\sigma^{-1}u\sigma\in I_v$ and $p(u)=p(\sigma)\mathrm{loc}_v(s)p(\sigma)^{-1}$.
If $A$ does not lie above $\ell$, then $u\in I_A\cap K\subset\ker p$.

Suppose that $\mathbf F$ is finite.
The map $(x,(X_v))\mapsto x$ identifies $H^1_{\mathrm{all}}(\Gamma,\mathbf F)$ with $H^1(\Gamma,\mathbf F)$.
Lemma~\ref{sha-descent-link-character}, for the trivial character, the places $A_v$, the $\tau_w$ and $C_0$, gives $\dim H^1(\Gamma,\mathbf F)+[L^+:\mathbf Q]\le1+C_0+[L^+:\mathbf Q]$; the set $S_\kappa$ there is empty, because $\omega$ is nontrivial on each $G_v$.
So $\dim H^1_{\mathrm{all}}(\Gamma,\mathbf F)\le1+C_0$.

Suppose that $\bar\Psi_\Gamma^2$ is nontrivial on $\mathrm{loc}_v(G_v)$ for every $v$.
Apply Lemma~\ref{sha-descent-link-strict} with $U=G_F$, $K_1$, $\pi=p$, the maps $\mathrm{loc}_v$ and $\kappa=\bar\Psi_\Gamma^{-2}\delta_\Gamma$, which is trivial on $p(K_1)$; its hypothesis on inertia was shown above.
Since $\mathrm{loc}_v(G_v)\subset\Gamma_0$, where $\delta_\Gamma=1$, $\kappa\circ\mathrm{loc}_v=\bar\Psi_\Gamma^{-2}\circ\mathrm{loc}_v\ne1$ for every $v$.
So $\dim H^1_\emptyset(\Gamma,\mathbf F(\bar\Psi_\Gamma^{-2}\delta_\Gamma))\le r_\ell(K_1)\le C_0$.
Next apply Lemma~\ref{sha-descent-link-strict} with $U=\operatorname{Gal}(\bar{\mathbf Q}/L)=p^{-1}(\Gamma_0)$, which contains $K_1$ with index dividing that of $K_1$ in $G_F$, the restriction $U\to\Gamma_0$ of $p$, the $2|V|$ maps $\mathrm{loc}_v$ and $h\mapsto c_\Gamma\mathrm{loc}_v(h)c_\Gamma^{-1}$ from $G_v$ to $\Gamma_0$, and $\kappa=\bar\Psi_\Gamma^{-1}$ on $\Gamma_0$.
Its hypothesis on inertia holds: for $u\in I_A\cap K_1$ with $A$ above $\ell$, $p(u)=p(\sigma)\mathrm{loc}_v(s)p(\sigma)^{-1}$; if $\delta_\Gamma(p(\sigma))=1$, this is a conjugate of $\mathrm{loc}_v(s)$ in $\Gamma_0$, and otherwise $g=p(\sigma)c_\Gamma^{-1}$ lies in $\Gamma_0$ and $p(u)=g\,(c_\Gamma\mathrm{loc}_v(s)c_\Gamma^{-1})\,g^{-1}$.
Since $\bar\Psi_\Gamma(c_\Gamma xc_\Gamma^{-1})=\bar\Psi_\Gamma(x)$, $\kappa$ is nontrivial on the images of both maps for every $v$.
So $\dim H^1_\emptyset(\Gamma_0,\mathbf F(\bar\Psi_\Gamma^{-1}))\le r_\ell(K_1)\le C_0$.
\end{proof}

\begin{lemma}\label{sha-descent-residual-holds}
Let $(\eta,\Psi,c,\mathfrak B,M,H_b)$ be a descent datum, with the notation of Definition~\ref{sha-descent-datum}. Let $H_0$ and an auxiliary prime $w$ be as in Lemma~\ref{sha-descent-auxiliary}, let $C_0$ be as in its part (5), let $m\ge0$ with $3C_0+7<\ell^m$, let $U$ and $L$ be as in Lemma~\ref{sha-descent-subgroup} for $w$ and $m$, and let $g_1\in\operatorname{Gal}(\bar{\mathbf Q}/L)$ with $\bar\Psi(g_1)^2\ne1$. Then every realized link of the datum to $L$ with base element $g_1$ satisfies the residual conditions with constant $C_0$ and is large for $C_0$ (Definitions~\ref{sha-descent-link} and~\ref{sha-descent-residual}).
\end{lemma}
\leanref{Development/AutomorphicArithmetic/AutomorphicVarieties/FrobeniusSmoothVectors/OkBasic.lean}{UpperTheory.V.Fr.ok_basic}
\leanref{Development/AutomorphicArithmetic/AutomorphicVarieties/FrobeniusSmoothVectors/FrobeniusGaloisFieldFourCompatibility.lean}{UpperTheory.V.Fr.frobenius_galoisField_four_compatibility}
\begin{proof}
By Lemma~\ref{sha-descent-subgroup}(6), $\operatorname{Gal}(\bar{\mathbf Q}/L^+)=U$ and $p(U\cap\ker\eta)=\Gamma_0$; and $\omega_\Gamma\circ p=\omega$.

Condition (4). $\mathbf F$ is finite by Definition~\ref{sha-descent-datum}(a), and $\Gamma_0$ is open since the link is realized.
$c_\Gamma^2=p(c^2)=1$, $\delta_\Gamma(c_\Gamma)=-1$, and $\bar\Psi_\Gamma(c_\Gamma)=\bar\Psi(c)=1$ by (d); $\bar\Psi_\Gamma(g_0)^2=\bar\Psi(g_1)^2\ne1$.
The maps $\mathrm{loc}_v$ are continuous because $p$ is, and $\varepsilon_\Gamma$ and $\bar\Psi_\Gamma$ are continuous, with open kernel for $\bar\Psi_\Gamma$, because $p$ is a quotient map and $H_{\eta,\Psi}$ is open.

Condition (6). The number of infinite places of the totally real field $L^+$ is $[L^+:\mathbf Q]=\sum_vd_v$.
Each $\tau_w$ is a complex conjugation, so $c_w^2=1$, and $\delta_\Gamma(c_w)=\eta(\tau_w)=-1$ since $\tau_w(\delta)=-\delta$.

Condition (7). The statements on $\Gamma_0$ are Lemma~\ref{sha-descent-subgroup}(5), and $\omega_\Gamma$ is nontrivial on $\Gamma_0$ because otherwise $\omega_\Gamma=\bar\Psi_\Gamma^0$ there; the characters $\omega_\Gamma\bar\Psi_\Gamma^{\pm2}\delta_\Gamma$ are nontrivial on $\Gamma_0$, where $\delta_\Gamma=1$.
For every $v$, Lemma~\ref{sha-descent-subgroup}(4) at $A_v$ gives $h\in G_v=D_{A_v}\cap U$ with $\bar\Psi(h)^2\ne1$.
For $h\in I'_v$, $\omega(h)=1$, and $\bar\Psi(h)=\omega(h)^j=1$ by Definition~\ref{sha-descent-datum}(d).

Condition (8).
Let $a\in\{1,2\}$, let $y$ be as in Lemma~\ref{sha-descent-strict} for $\kappa=\bar\Psi^{-a}\eta$, and let $T$ be the set of $t\in\ker p$ with $y(t)=0$.
Since $\kappa$ is trivial on $\ker p$, $T$ is a subgroup; it is closed, normal in $U$, and contains the generators of $\ker p$ by Lemma~\ref{sha-descent-strict}(3); so $T=\ker p$.
Hence $y(st)=y(s)$ for $t\in\ker p$, and $y$ descends to a continuous cocycle $\phi$ of $\Gamma$ with values in $\mathbf F(\bar\Psi_\Gamma^{-a}\delta_\Gamma)$.
For every $v$, $A_v$ lies above $\ell$, so Lemma~\ref{sha-descent-strict}(2) gives $X_v$ with $\phi(\mathrm{loc}_vh)=\kappa(h)X_v-X_v$ for $h\in G_v$.
So $(\phi,(X_v))$ defines an element of $H^1_\emptyset(\Gamma,\mathbf F(\bar\Psi_\Gamma^{-a}\delta_\Gamma))$.
If it were $0$, then $\phi$ would be a coboundary, and so would $y$ on $U$, contrary to Lemma~\ref{sha-descent-strict}(1).

Conditions (5) and (9) to (12). The group $H_{\eta,\Psi}\cap\ker\omega$ is $U\cap\ker\eta\cap\ker\bar\Psi\cap\ker\omega$, so $r_\ell(H_{\eta,\Psi}\cap\ker\omega)\le C_0$ by Lemma~\ref{sha-descent-auxiliary}(5); $\mathbf F$ is finite and $\bar\Psi_\Gamma^2$ is nontrivial on each $\mathrm{loc}_v(G_v)$, as shown above. Lemma~\ref{sha-descent-link-cohomology} gives these conditions.

Largeness. By Lemma~\ref{sha-descent-subgroup}(4), $d_v$, the index of $D_{A_v}\cap U$ in $D_{A_v}$, is at least $\ell^m>3C_0+7\ge7$.
The place of $L^+$ below $\mathfrak B$ shows that $V$ is not empty, so $d>3C_0+7$.
\end{proof}

\begin{definition}\label{sha-descent-seed}
Let $(\eta,\Psi,c,\mathfrak B,M,H_b)$ be a descent datum, with the notation of Definition~\ref{sha-descent-datum}, where $\mathcal O\subset\bar{\mathbf Q}_\ell$ is the ring of integers of a finite extension of $\mathbf Q_\ell$, and let $L\subset\bar{\mathbf Q}$ be a number field. View $\eta$ as a character with values in $\bar{\mathbf Q}_\ell^\times$, and call $x,y\in\bar{\mathbf Q}_\ell$ \emph{congruent} if $|x-y|<1$. A \emph{seed datum} for $L$ consists of an imaginary quadratic field $K$, integers $N_{\mathrm s}\ge1$ and $j\ge0$, a prime $w$, and homomorphisms $\chi',\theta:G_{\mathbf Q}\to\bar{\mathbf Q}_\ell^\times$, such that:
\begin{enumerate}
\item $\operatorname{Gal}(\bar{\mathbf Q}/K)=\ker\eta$, $\ell$ does not divide the discriminant $d_K$, and $L$ is abelian over $\mathbf Q$ and contains $K$;
\item $\ell$ and $w$ do not divide $N_{\mathrm s}$; $\theta\ne1$ and $\theta^\ell=1$; and $\theta$ is trivial on the elements fixing $\mu_w$ and on $\operatorname{Gal}(\bar{\mathbf Q}/L)$;
\item $\eta$ and $\chi'$ are trivial on the elements of $G_{\mathbf Q}$ that fix $\mu_{N_{\mathrm s}}$, and on those that act on $\mu_{N_{\mathrm s}}$ by $\zeta\mapsto\zeta^w$;
\item $\chi'\ne\eta$, $\ell$ does not divide the order of $\chi'$, and $\chi'(c)=(-1)^j$;
\item $(\chi_{\mathrm{cyc}}^j\chi')(s)$ is congruent to $\Psi(s)$ for every $s\in G_{\mathbf Q}$, and $(\chi_{\mathrm{cyc}}^j\chi')(s)^2$ is not congruent to $1$ for some $s\in\operatorname{Gal}(\bar{\mathbf Q}/L)$.
\end{enumerate}
\end{definition}
\leanref{Development/AutomorphicArithmetic/AutomorphicVarieties/CurveDVRSpecialization/HostFrameEx.lean}{UpperTheory.V.Pr.SeedClause}
\leanref{Development/AutomorphicArithmetic/AutomorphicVarieties/IdeleGenerators/SlopeTestRemainingChoices.lean}{UpperTheory.V.Sd.FieldPairCharacterTripleConditions}
\leanref{Development/AutomorphicArithmetic/AutomorphicVarieties/IdeleGenerators/SlopeTestRemainingChoices.lean}{UpperTheory.V.Sd.SlopeTestRemainingChoices}
\leanref{Development/AutomorphicArithmetic/AutomorphicVarieties/IdeleGenerators/SlopeTestRemainingChoices.lean}{UpperTheory.V.Sd.SplitAt}

\begin{lemma}\label{sha-descent-seed-exists}
Let $(\eta,\Psi,c,\mathfrak B,M,H_b)$ be a descent datum, with the notation of Definition~\ref{sha-descent-datum}, where $\mathcal O\subset\bar{\mathbf Q}_\ell$ is the ring of integers of a finite extension of $\mathbf Q_\ell$.
Then there is $Q\ge1$ with the following property.
Let $L\subset\bar{\mathbf Q}$ be an abelian CM field with maximal totally real subfield $L^+$ and $\operatorname{Gal}(\bar{\mathbf Q}/L)=\operatorname{Gal}(\bar{\mathbf Q}/L^+)\cap\ker\eta$, let $w\equiv1\pmod Q$ be an auxiliary prime (Lemma~\ref{sha-descent-auxiliary}(4)) with $\operatorname{Gal}(\bar{\mathbf Q}/L^+)\subset H_{\mathrm{th}}(w)$, and suppose that $\bar\Psi(s)^2\ne1$ for some $s\in\operatorname{Gal}(\bar{\mathbf Q}/L)$.
Then there is a seed datum for $L$ (Definition~\ref{sha-descent-seed}).
\end{lemma}
\leanref{Development/AutomorphicArithmetic/AutomorphicVarieties/FrobeniusSmoothVectors/SeedCl.lean}{UpperTheory.V.Fr.seed_cl}
\leanref{Development/AutomorphicArithmetic/AutomorphicVarieties/FrobeniusSmoothVectors/ScQuad.lean}{UpperTheory.V.Fr.sc_quad}
\leanref{Development/AutomorphicArithmetic/AutomorphicVarieties/FrobeniusSmoothVectors/ScTeich.lean}{UpperTheory.V.Fr.sc_teich}
\leanref{Development/AutomorphicArithmetic/AutomorphicVarieties/FrobeniusSmoothVectors/ScStrip.lean}{UpperTheory.V.Fr.sc_strip}
\leanref{Development/AutomorphicArithmetic/AutomorphicVarieties/FrobeniusSmoothVectors/ScTheta.lean}{UpperTheory.V.Fr.sc_theta}
\leanref{Development/AutomorphicArithmetic/AutomorphicVarieties/FrobeniusSmoothVectors/ScDisc.lean}{UpperTheory.V.Fr.sc_disc}
\leanref{Development/AutomorphicArithmetic/AutomorphicVarieties/FrobeniusSmoothVectors/ScCong.lean}{UpperTheory.V.Fr.sc_cong}
\begin{proof}
Let $a$, $D$, $\delta$ and $j$ be as in Definition~\ref{sha-descent-datum}, and write $M^a\ell=\ell^rN_{\mathrm s}$ with $\ell\nmid N_{\mathrm s}$; put $Q=N_{\mathrm s}$.
Let $K=\mathbf Q(\delta)$, an imaginary quadratic field because $x^2+D$ is irreducible and $D>0$; by (b), $\operatorname{Gal}(\bar{\mathbf Q}/K)=\ker\eta$.
Let $T:\mathbf F^\times\to\mathcal O^\times\subset\bar{\mathbf Q}_\ell^\times$ be the Teichm\"uller lift, which is injective and reduces to the identity.
Let $\chi'=T\circ(\bar\Psi\omega^{-j})$.

The characters $\chi'$ and $\eta$ are trivial on the elements fixing $\mu_{M^a\ell}$, by (c), and on $I_{\mathfrak B}$, by (d).
They are then trivial on the elements fixing $\mu_{N_{\mathrm s}}$.
Indeed, let $s$ fix $\mu_{N_{\mathrm s}}$, and let $\sigma\in I_{\mathfrak B}$ with $\chi_{\mathrm{cyc}}(\sigma)=\chi_{\mathrm{cyc}}(s)$, which exists because $\chi_{\mathrm{cyc}}$ maps $I_{\mathfrak B}$ onto $\mathbf Z_\ell^\times$.
Then $\sigma$ fixes $\mu_{N_{\mathrm s}}$, as $\ell\nmid N_{\mathrm s}$, and $\sigma^{-1}s$ fixes $\mu_{\ell^r}$ and $\mu_{N_{\mathrm s}}$, hence $\mu_{\ell^rN_{\mathrm s}}$; so the character takes the value $1$ at $\sigma^{-1}s$ and at $\sigma$, hence at $s$.
Since $w\equiv1\pmod{N_{\mathrm s}}$, an element acting on $\mu_{N_{\mathrm s}}$ by $\zeta\mapsto\zeta^w$ fixes $\mu_{N_{\mathrm s}}$; this gives condition (3).

The prime $w$ does not divide $N_{\mathrm s}$: it does not divide $M$, and $w\ne\ell$ because $w\equiv1\pmod\ell$.
By Lemma~\ref{sha-descent-auxiliary}(2), there is a character $\theta:G_{\mathbf Q}\to\bar{\mathbf Q}_\ell^\times$ of order $\ell$ with kernel $H_{\mathrm{th}}(w)$; it is trivial on the elements fixing $\mu_w$, and on $\operatorname{Gal}(\bar{\mathbf Q}/L)\subset\operatorname{Gal}(\bar{\mathbf Q}/L^+)\subset H_{\mathrm{th}}(w)$.
This gives condition (2).

For condition (1): $\operatorname{Gal}(\bar{\mathbf Q}/L)\subset\ker\eta=\operatorname{Gal}(\bar{\mathbf Q}/K)$, so $K\subset L$; and $\ell\nmid d_K$ because $I_A\subset\ker\eta$ for the places $A$ above $\ell$, by (d), so that $K$ is unramified at $\ell$.

For condition (4): $\chi'^{|\mathbf F^\times|}=1$ and $|\mathbf F^\times|=|\mathbf F|-1$ is prime to $\ell$.
Since $T$ is injective and $T(-1)^2=1$, $T(-1)=-1$; so $\chi'(c)=T(\bar\Psi(c)\omega(c)^{-j})=T((-1)^j)=(-1)^j$, by (d) and $\omega(c)=-1$.
If $\chi'=\eta$, then every $s\in\ker\omega\cap\ker\bar\Psi$ has $\eta(s)=\chi'(s)=T(1)=1$, contrary to (e).

For condition (5): $(\chi_{\mathrm{cyc}}^j\chi')(s)=\chi_{\mathrm{cyc}}(s)^jT(\bar\Psi(s)\omega(s)^{-j})$ lies in $\mathcal O^\times$ and has reduction $\omega(s)^j\bar\Psi(s)\omega(s)^{-j}=\bar\Psi(s)$, so it is congruent to $\Psi(s)$.
If $(\chi_{\mathrm{cyc}}^j\chi')(s)^2$ were congruent to $1$ for the element $s\in\operatorname{Gal}(\bar{\mathbf Q}/L)$ of the statement, then $\Psi(s)^2$ would be congruent to $1$, and $\bar\Psi(s)^2=1$.
\end{proof}

\begin{proposition}\label{sha-descent-construction}
Let $(\eta,\Psi,c,\mathfrak B,M,H_b)$ be a descent datum, with the notation of Definition~\ref{sha-descent-datum}, where $\mathcal O\subset\bar{\mathbf Q}_\ell$ is the ring of integers of a finite extension of $\mathbf Q_\ell$.
Then there are an abelian CM field $L\subset\bar{\mathbf Q}$, with maximal totally real subfield $L^+$ and $\operatorname{Gal}(\bar{\mathbf Q}/L)=\operatorname{Gal}(\bar{\mathbf Q}/L^+)\cap\ker\eta$, an element $g_1\in\operatorname{Gal}(\bar{\mathbf Q}/L)$ with $\bar\Psi(g_1)^2\ne1$, a link of the datum to $L$ with base element $g_1$, an integer $C_0\ge0$ such that the link satisfies the residual conditions with constant $C_0$ and is large for $C_0$, and a seed datum for $L$ (Definitions~\ref{sha-descent-link}, \ref{sha-descent-residual} and~\ref{sha-descent-seed}).
\end{proposition}
\leanref{Development/AutomorphicArithmetic/AutomorphicVarieties/FrobeniusSmoothVectors/FrameEx.lean}{UpperTheory.V.Fr.frame_ex}
\leanref{Development/AutomorphicArithmetic/AutomorphicVarieties/FrobeniusSmoothVectors/FrobeniusSplitPrimeDatum.lean}{UpperTheory.V.Fr.HasFrameW}
\leanref{Development/AutomorphicArithmetic/AutomorphicVarieties/CurveDVRSpecialization/HostFrameEx.lean}{UpperTheory.V.Pr.HasFrame}
\leanref{Development/AutomorphicArithmetic/AutomorphicVarieties/FrobeniusSmoothVectors/GdEx.lean}{UpperTheory.V.Fr.gd_ex}
\leanref{Development/AutomorphicArithmetic/AutomorphicVarieties/CurveDVRSpecialization/HostFrameEx.lean}{UpperTheory.V.Pr.GdFrameEx}
\begin{proof}
Let $Q$ be as in Lemma~\ref{sha-descent-seed-exists}.
Let $H_0$ be as in Lemma~\ref{sha-descent-auxiliary}(1), $w\equiv1\pmod Q$ an auxiliary prime, given by its part (4), and $C_0$ as in its part (5).
Let $m=3C_0+7$, so that $3C_0+7<\ell^m$, and let $U$ and $L$ be as in Lemma~\ref{sha-descent-subgroup} for $w$ and $m$.
By Lemma~\ref{sha-descent-subgroup}(4) at $\mathfrak B$, there is $g_1\in D_{\mathfrak B}\cap U$ with $\bar\Psi(g_1)^2\ne1$; by its part (2), $\eta(g_1)=1$, since $\ell$ divides $M$; so $g_1\in U\cap\ker\eta=\operatorname{Gal}(\bar{\mathbf Q}/L)$.
Lemma~\ref{sha-descent-link-exists} gives a realized link of the datum to $L$ with base element $g_1$, and Lemma~\ref{sha-descent-residual-holds} shows that it satisfies the residual conditions with constant $C_0$ and is large for $C_0$.
Finally $\operatorname{Gal}(\bar{\mathbf Q}/L^+)=U\subset H_{\mathrm{th}}(w)$, and Lemma~\ref{sha-descent-seed-exists}, with $s=g_1$, gives a seed datum for $L$.
\end{proof}

\begin{lemma}\label{sha-descent-eisenstein}
Let $(\ell,k,\psi,K)$ be an Eisenstein datum, with $\psi$ a Dirichlet character modulo $N$, let $(\eta,c,\Psi)$ be deformation data for it (Definition~\ref{sha-ordinary-deformation}), and let $\mathfrak B$ be a place of $\bar{\mathbf Q}$ above $\ell$.
Then $(\eta,\Psi,c,\mathfrak B,\ell N|d_K|,\operatorname{Gal}(\bar{\mathbf Q}/\mathbf Q))$ is a descent datum (Definition~\ref{sha-descent-datum}).
\end{lemma}
\leanref{Development/AutomorphicArithmetic/AutomorphicVarieties/FrobeniusSmoothVectors/SqrtDiscr.lean}{UpperTheory.V.Fr.gen_gd}
\begin{proof}
Let $M=\ell N|d_K|$, let $\omega$ and $\bar\Psi$ be as in Definition~\ref{sha-descent-datum}, and let $t:G_{\mathbf Q}\to\mathcal O^\times$ be the Teichm\"uller lift of $\omega$.
On the inertia groups above $\ell$, $\psi$ acts through its component $\psi_\ell$, which has order prime to $\ell$; so $\psi=t^{b}$ there for some $b\ge0$, and $\chi'_\ell(\psi)=\psi t^{-b}$.

(a) The residue field of $\mathcal O$ is finite; $\ell$ divides $M$, and so does $2$, as $N$ is even.

(b) Take $D=|d_K|$ and $\delta$ a square root of $d_K$ in $K$; then $\delta^2=-D$, since $K$ is imaginary, and $g(\delta)=\eta(g)\delta$ because $\ker\eta=\operatorname{Gal}(\bar{\mathbf Q}/K)$.

(c) Take $a=1$. Since $K\subset\mathbf Q(\mu_{|d_K|})$, $\eta$ is trivial on the elements fixing $\mu_M$; $\psi$ and $\omega$ are trivial on them too, so $\bar\Psi=\omega^k\bar\psi$ is.

(d) $\eta$ is unramified at $\ell$, since $\ell\nmid d_K$; $\Psi(c)=(-1)^k\psi(-1)=1$; and on the inertia groups above $\ell$, $\bar\Psi=\omega^{k+b}$.

(e) Suppose that $\ker\omega\cap\ker\bar\Psi\subset\ker\eta$, and let $s$ satisfy $\omega(s)=1$ and $\chi'_\ell(\psi)(s)=1$.
Then $t(s)=1$, so $\psi(s)=1$ and $\bar\Psi(s)=1$; hence $s\in\ker\eta=\operatorname{Gal}(\bar{\mathbf Q}/K)$.
So $K$ would be contained in the field cut out by $\chi_{\mathrm{cyc}}\bmod\ell$ and $\chi'_\ell(\psi)$, which Definition~\ref{sha-ordinary-deformation}(a) excludes.

(f) The inertia group $I_{\mathfrak B}$ lies in $\ker\eta$ by (d).
On $I_{\mathfrak B}$, $\Psi^2=\chi_{\mathrm{cyc}}^{2k}\psi_\ell^2=\chi_{\mathrm{cyc}}^{2k}$, and $\chi_{\mathrm{cyc}}$ maps $I_{\mathfrak B}$ onto $\mathbf Z_\ell^\times$; since $\ell-1\nmid2k$, some $s\in I_{\mathfrak B}$ has $\chi_{\mathrm{cyc}}(s)^{2k}\not\equiv1\pmod\ell$, so $\Psi(s)^2-1\in\mathcal O^\times$.

(g) The whole group $G_{\mathbf Q}$ is open of index $\ell^0$.
Let $s$ lie in an inertia group at a place not above $\ell$, with $\bar\Psi(s)=1$.
Then $\chi_{\mathrm{cyc}}(s)=1$, so $\Psi(s)=\psi(s)$ is a root of unity of order prime to $\ell$, because $\ell\nmid\operatorname{ord}\psi$, and it is congruent to $1$ modulo $\mathfrak m_{\mathcal O}$; so $\Psi(s)=1$.
\end{proof}

\begin{lemma}\label{sha-descent-points}
Let $(\eta,\Psi,c,\mathfrak B,M,H_b)$ be a descent datum, with the notation of Definition~\ref{sha-descent-datum}, let $L$ be an abelian CM field, let $g_1\in G_{\mathbf Q}$, and let there be a link of the datum to $L$ with base element $g_1$ that satisfies the residual conditions with some constant (Definitions~\ref{sha-descent-link} and~\ref{sha-descent-residual}).
Then $\eta(c)=-1$, $\eta(g_1)=1$ and $\Psi(g_1)^2-1\in\mathcal O^\times$.
Let $A$ be a local ring that is a noetherian domain, with a local homomorphism $\mathcal O\to A$, and let $r:G_{\mathbf Q}\to\mathcal G_3(A)$ be a homomorphism with the properties that define an ordinary point with values in $A$ relative to $g_1$ and $\mathfrak B$ in Definition~\ref{sha-ordinary-deformation}(c), for the characters $\eta$ and $\Psi$ and the complex conjugation $c$; these properties make sense without completeness of $A$.
\begin{enumerate}
\item There is a linked point $(\rho,(\psi_{v,i}))$ of rank three with values in $A$ such that $\rho(p(s))$ is the matrix part of $r(s)$ for $s\in\operatorname{Gal}(\bar{\mathbf Q}/L)$, and $\psi_{v,i}(h)=\Psi(h)^{e_{3,i}}$ for every $v$ and $h\in I_v$.
\item If moreover the entries $(1,2)$, $(2,1)$, $(2,3)$ and $(3,2)$ of the matrix part of $r$ vanish on $\ker\eta$, then the matrices formed by the rows and columns $1$ and $3$ of $\rho$, with the characters $\psi_{v,1}$ and $\psi_{v,3}$, form a linked point of rank two with values in $A$.
\end{enumerate}
\end{lemma}
\leanref{Development/AutomorphicArithmetic/AutomorphicVarieties/CurveDVRSpecialization/LkMarks.lean}{UpperTheory.V.Pr.marks_of_link}
\leanref{Development/AutomorphicArithmetic/AutomorphicVarieties/CurveDVRSpecialization/ExistsNrep3.lean}{UpperTheory.V.Pr.exists_nrep3}
\leanref{Development/AutomorphicArithmetic/AutomorphicVarieties/CurveDVRSpecialization/ExistsNrep2.lean}{UpperTheory.V.Pr.exists_nrep2}
\begin{proof}
Since $c_\Gamma=p(c)$ and $\delta_\Gamma(c_\Gamma)=-1$, $\eta(c)=-1$; since $g_0=p(g_1)\in\Gamma_0$, $\eta(g_1)=1$; and $\bar\Psi(g_1)^2=\bar\Psi_\Gamma(g_0)^2\ne1$, so $\Psi(g_1)^2-1$ is a unit.

(1) First, $r$ is trivial on $\ker p$.
The generators $s$ of $\ker p$ lie in an inertia group at a place not above $\ell$ and have $\eta(s)=\bar\Psi(s)=1$; so $\Psi(s)=1$ by the definition of a link, and $r(s)=1$ by the definition of an ordinary point.
The set of $s$ at which the matrix part of $r$ is $1$ is closed, because $r$ is continuous for the $\mathfrak m_A$-adic topology and $\bigcap_n\mathfrak m_A^n=0$; so the matrix part of $r$ is $1$ on $\ker p$, and so is the multiplier $\eta=\delta_\Gamma\circ p$.
Since $p$ is a continuous surjection from a compact group to a Hausdorff group, $r$ factors as $r'\circ p$ with $r'$ continuous.
Let $\rho(g)$ be the matrix part of $r'(g)$ for $g\in\Gamma_0$, where the multiplier is $1$; $\rho$ is multiplicative.
The matrix part of $r'(c_\Gamma)$ is $1$ and its multiplier is $-1$, so the matrix part of $r'(c_\Gamma g^{-1}c_\Gamma^{-1})$ is $J_3\,{}^t\rho(g)J_3$; this is (a).
For (b), let $X=\rho(g_0)$, so that $\rho(c_\Gamma g_0^{-1}c_\Gamma^{-1})=J_3\,{}^tXJ_3$, whose entry $(i,i')$ is $X_{4-i',4-i}$.
The entries $(1,2)$ and $(2,1)$ of $X-J_3\,{}^tXJ_3$ vanish because $r$ is normalized at $g_1$; the entries $(1,3)$ and $(3,1)$ vanish identically; and the entries $(2,3)$ and $(3,2)$ are the negatives of the entries $(1,2)$ and $(2,1)$.

For (c), let $v$ be a place of $L^+$ above $\ell$, with $A_v=\gamma_v\mathfrak B$ and $\eta(\gamma_v)=1$.
The ordinary condition at $\mathfrak B$ gives $h_0\in GL_3(A)$, congruent to $1$ modulo $\mathfrak m_A$, such that $h_0^{-1}r(g)h_0$ is upper triangular for $g\in D_{\mathfrak B}\cap\ker\eta$, with diagonal $\Psi(g),1,\Psi(g)^{-1}$ for $g\in I_{\mathfrak B}$.
For $h\in G_v\subset D_{A_v}\cap\ker\eta$, $\gamma_v^{-1}h\gamma_v\in D_{\mathfrak B}\cap\ker\eta$; let $T(h)=h_0^{-1}r(\gamma_v^{-1}h\gamma_v)h_0$, which is upper triangular, and $\psi_{v,i}(h)=T(h)_{ii}$, which are characters.
Then $r(h)=PT(h)P^{-1}$ with $P=r(\gamma_v)h_0$.
For upper triangular $3\times3$ matrices $T_1,T_2,T_3$, $(T_1-(T_1)_{11})(T_2-(T_2)_{22})(T_3-(T_3)_{33})=0$; conjugating by $P$ gives (c).
The entries of $T$ are those of $P^{-1}r(h)P$, which gives (d) for $\psi_{v,i}$, and (d) for $\rho$ follows from the continuity of $r$.
Modulo $\mathfrak m_A$, $r$ is $r_0(\eta,\Psi)$ and $h_0$ is $1$, so $\rho(g)$ and $T(h)$ reduce to the diagonal matrices with entries $\bar\Psi(g)^{e_{3,i}}$ and $\bar\Psi(\gamma_v^{-1}h\gamma_v)^{e_{3,i}}=\bar\Psi(h)^{e_{3,i}}$; this is (e).
For $h\in I_v$, $\gamma_v^{-1}h\gamma_v\in I_{\mathfrak B}$, so $\psi_{v,i}(h)=\Psi(h)^{e_{3,i}}$.
Finally, for $s\in\operatorname{Gal}(\bar{\mathbf Q}/L)$, $\rho(p(s))$ is the matrix part of $r(s)$.

(2) Every $g\in\Gamma_0$ is $p(s)$ for some $s\in\operatorname{Gal}(\bar{\mathbf Q}/L)\subset\ker\eta$, so the entries $(1,2)$, $(2,1)$, $(2,3)$ and $(3,2)$ of $\rho(g)$ vanish.
For such matrices, taking the rows and columns $1$ and $3$ is multiplicative and carries $J_3\,{}^tXJ_3$ to $J_2\,{}^tX'J_2$; this gives a homomorphism $\rho'$ with (a), (b), (d) and (e), for the exponents $(1,-1)$.
For (c), let $h_1,h_3\in G_v$, and let $t\in G_v$ with $\bar\Psi_\Gamma(\mathrm{loc}_vt)^2\ne1$, given by the residual conditions.
Write $S=\rho'(\mathrm{loc}_vt)-\psi_{v,2}(t)$, $X(h)=\rho'(\mathrm{loc}_vh)-\psi_{v,1}(h)$ for $h\in G_v$, and $Z=\rho'(\mathrm{loc}_vh_3)-\psi_{v,3}(h_3)$.
The relation (c) for $\rho$ at $(h,t,h_3)$, restricted to the rows and columns $1$ and $3$, gives $X(h)SZ=0$ for every $h\in G_v$.
The determinant of $S$ is a unit, since modulo $\mathfrak m_A$ it is $(\bar\Psi(t)-1)(\bar\Psi(t)^{-1}-1)$.
By the Cayley-Hamilton theorem, $(\operatorname{tr}S-S)S=\det(S)$.
A direct expansion gives $X(h_1)(\operatorname{tr}S-S)=(\operatorname{tr}S+\psi_{v,2}(t))X(h_1)-X(h_1t)+\psi_{v,1}(h_1)X(t)$.
Multiplying on the right by $SZ$ gives $\det(S)X(h_1)Z=0$, so $X(h_1)Z=0$, which is (c) for $\rho'$.
\end{proof}

\subsection{Eigenforms congruent to Eisenstein series}\label{subsection-sha-eisenstein-eigenforms}

\begin{lemma}[Eisenstein series of level $\Gamma_1(N)$ \cite{DS05}]\label{sha-automorphy-eisenstein-series}
Let $N\ge1$ and $k\ge3$, and let $\zeta=e^{2\pi i/N}$.
For $a\in\mathbf Z/N$ let
\[
E_{k,a}(z)=\frac{N^k(k-1)!}{(-2\pi i)^k}\sum_{(c,d)}(cz+d)^{-k}\qquad(z\in\mathfrak H),
\]
the sum over the pairs $(c,d)\ne(0,0)$ of integers with $c\equiv0$ and $d\equiv a\pmod N$.
Let $A_0=0$, $A_1=X$ and $A_{j+1}=X\bigl((1-X)A_j'+jA_j\bigr)$ for $j\ge1$, in $\mathbf Z[X]$, and for $x\ne1$ in a field let $L_j(x)=A_j(x)/(1-x)^j$; so $L_{j+1}(x)=xL_j'(x)$, and $L_j(x)=\sum_{e\ge1}e^{j-1}x^e$ for $x\in\mathbf C$ with $|x|<1$.
\begin{enumerate}
\item $E_{k,a}$ is a modular form of weight $k$ for $\Gamma_1(N)$, and $E_{k,a}|_k\gamma=E_{k,ad}$ for every $\gamma=\begin{pmatrix}*&*\\ *&d\end{pmatrix}\in\Gamma_0(N)$.
\item If $N>1$ and $a$ is a unit, then $E_{k,a}|_k\gamma$ vanishes at $\infty$ for every $\gamma\in SL_2(\mathbf Z)$ whose lower left entry is not divisible by $N$.
\item If $a\ne0$, the $q$-expansion of $E_{k,a}$ has constant term $L_k(\zeta^a)$ and, for $n\ge1$, $n$-th coefficient $\sum_{e\mid n}e^{k-1}\bigl(\zeta^{ae}+(-1)^k\zeta^{-ae}\bigr)$.
\item For every even $j\ge0$ with $j\ne2$ there is a modular form of weight $j$ for $SL_2(\mathbf Z)$ with constant term $1$ and integer $q$-expansion coefficients.
\end{enumerate}
\end{lemma}
\leanref{Development/AutomorphicArithmetic/AutomorphicVarieties/IdeleGenerators/HolomorphicIrreducibility/TsumNormEisensteinTranslate.lean}{UpperTheory.V.Sd.eisMF}
\leanref{Development/AutomorphicArithmetic/AutomorphicVarieties/IdeleGenerators/HolomorphicIrreducibility/EisSlash.lean}{UpperTheory.V.Sd.eisSlash}
\leanref{Development/AutomorphicArithmetic/AutomorphicVarieties/IdeleGenerators/HolomorphicIrreducibility/EisVan.lean}{UpperTheory.V.Sd.eisVan}
\leanref{Development/AutomorphicArithmetic/AutomorphicVarieties/IdeleGenerators/HolomorphicIrreducibility/EisCoef.lean}{UpperTheory.V.Sd.eisCoef}
\leanref{Development/AutomorphicArithmetic/AutomorphicVarieties/IdeleGenerators/HolomorphicIrreducibility/LevelOne.lean}{UpperTheory.V.Sd.levelOne}

\begin{lemma}\label{sha-automorphy-eulerian}
Let $\ell$ be an odd prime, let $N\ge1$ with $\ell\nmid N$, and let $L_j$ be as in Lemma~\ref{sha-automorphy-eisenstein-series}. Let $x,y\in\bar{\mathbf Q}_\ell$ with $x^N=y^N=1$.
\begin{enumerate}
\item If $x\ne1$, then $|L_j(x)|\le1$ for every $j\ge1$.
\item If $x^2\ne1$, then $|L_3(x)|=1$.
\item If $x^2\ne1$, $y^2\ne1$ and $y\notin\{x,x^{-1}\}$, then one of $L_4(x)-L_4(y)$, $L_3(x)^2-L_3(y)^2$ and $L_3(x)L_5(x)-L_3(y)L_5(y)$ has absolute value $1$.
\end{enumerate}
\end{lemma}
\leanref{Development/AutomorphicArithmetic/AutomorphicVarieties/IdeleGenerators/HolomorphicIrreducibility/EulerianNumeratorThree.lean}{UpperTheory.V.Sd.lkSep}
\begin{proof}
For a root of unity $z\ne1$ of order dividing $N$, $1-z$ is a unit, since $\ell\nmid N$.
So (1) holds, because $A_j$ has integer coefficients.
For (2), $L_3(x)=x(1+x)/(1-x)^3$, and $1+x$ is a unit because $1-x^2=(1-x)(1+x)$ is.
For (3), let $u=x/(1-x)^2$ and $v=y/(1-y)^2$. They are units, and so is
\[
u-v=\frac{x(1-y/x)(1-xy)}{(1-x)^2(1-y)^2},
\]
since $y\ne x$ and $xy\ne1$.
From $A_4=X(1+4X+X^2)$ and $A_5=X(1+X)(1+10X+X^2)$ one gets $L_4(x)=u+6u^2$ and $L_3(x)L_5(x)=L_3(x)^2(1+12u)$, so
\[
L_4(x)-L_4(y)=(u-v)\bigl(1+6(u+v)\bigr),\quad
L_3(x)L_5(x)-L_3(y)L_5(y)=12L_3(x)^2(u-v)+\bigl(L_3(x)^2-L_3(y)^2\bigr)(1+12v).
\]
If $\ell=3$, the first difference is a unit.
Otherwise $12$ is a unit, and if $L_3(x)^2-L_3(y)^2$ is not a unit, the second identity and (2) make the third difference a unit.
\end{proof}

\begin{lemma}\label{sha-automorphy-unit-form}
Let $\ell$ be an odd prime, let $\iota:\bar{\mathbf Q}_\ell\cong\mathbf C$, and let $N\ge3$ with $\ell\nmid N$.
Then there is $k_1$ such that for every $k\ge k_1$ and every Dirichlet character $\chi$ modulo $N$ with $\chi(-1)=(-1)^k$ there is a modular form $G$ of weight $k$ for $\Gamma_1(N)$ with the following properties: $G|_k\gamma=\chi(d)G$ for every $\gamma=\begin{pmatrix}*&*\\ *&d\end{pmatrix}\in\Gamma_0(N)$; the $q$-expansion coefficients $b_n$ of $G$ satisfy $|\iota^{-1}(b_n)|\le1$; $G|_k\gamma$ vanishes at $\infty$ for every $\gamma\in SL_2(\mathbf Z)$ whose lower left entry is not divisible by $N$; and $G$ has constant term $1$.
\end{lemma}
\leanref{Development/AutomorphicArithmetic/AutomorphicVarieties/IdeleGenerators/HolomorphicIrreducibility/UnitForm.lean}{UpperTheory.V.Sd.unitForm}
\leanref{Development/AutomorphicArithmetic/AutomorphicVarieties/IdeleGenerators/HolomorphicIrreducibility/UnitFam.lean}{UpperTheory.V.Sd.unitFam}
\begin{proof}
Write $|x|$ for $|\iota^{-1}(x)|$, $c_0(f)$ for the constant term of a modular form $f$, and $\zeta=e^{2\pi i/N}$.
Consider families $f=(f_a)_{a\in(\mathbf Z/N)^\times}$ of modular forms of one weight $m$ for $\Gamma_1(N)$, with coefficients of absolute value at most $1$, such that $f_a|_m\gamma=f_{ad}$ for every $\gamma\in\Gamma_0(N)$ with lower right entry $d$.
Termwise products of such families are again such families, with $c_0(f_ag_a)=c_0(f_a)c_0(g_a)$; so is $f+bg$ for two families of the same weight and $|b|\le1$.
\begin{enumerate}
\item For $m\ge3$, the family $E_m=(E_{m,a})_a$ of Lemma~\ref{sha-automorphy-eisenstein-series} is such a family, by parts (1) and (3) of that lemma and Lemma~\ref{sha-automorphy-eulerian}(1) applied to $\zeta^a\ne1$; and $c_0(E_{m,a})=L_m(\zeta^a)$. For even $m\ne2$, a form of level one as in Lemma~\ref{sha-automorphy-eisenstein-series}(4) gives a constant family $F_m$ with constant terms $1$.
\item Let $d\in(\mathbf Z/N)^\times$ with $d\ne\pm1$. Lemma~\ref{sha-automorphy-eulerian}(3), applied to $x=\zeta$ and $y=\zeta^d$, shows that one of the families $H=E_4$, $E_3E_3$ and $E_3E_5$, of weight $m=4$, $6$ or $8$, has $|c_0(H_1)-c_0(H_d)|=1$; here $x^2\ne1$ because $N\ge3$, and $y\notin\{x,x^{-1}\}$ because $d\ne\pm1$. Then $h^{(d)}=H-c_0(H_d)F_m$ has $c_0(h^{(d)}_d)=0$ and $|c_0(h^{(d)}_1)|=1$.
\item The product $h$ of the families $h^{(d)}$ over all $d\ne\pm1$ has some even weight $m_0$, with $c_0(h_d)=0$ for $d\ne\pm1$ and $|c_0(h_1)|=1$.
\item For $j\ge10$, let $f=E_3E_3F_{j-6}$ if $j$ is even and $f=E_3F_{j-3}$ if $j$ is odd. For every $\gamma\in SL_2(\mathbf Z)$ whose lower left entry is not divisible by $N$, each $f_a|_j\gamma$ vanishes at $\infty$, because $E_{3,a}|_3\gamma$ does by Lemma~\ref{sha-automorphy-eisenstein-series}(2) and the other factor is bounded there. Also $|c_0(f_1)|=1$, since $|L_3(\zeta)|=1$ by Lemma~\ref{sha-automorphy-eulerian}(2).
\item Let $k_1=10+m_0$. For $k\ge k_1$, let $F=fh$, with $f$ as in (4) of weight $k-m_0$. Then $F$ is such a family of weight $k$, each $F_a$ vanishes at the cusps of (4), $|c_0(F_1)|=1$, and $c_0(F_a)=0$ for $a\ne\pm1$.
\item Let $G=(2c_0(F_1))^{-1}\sum_a\chi(a^{-1})F_a$. For $\gamma\in\Gamma_0(N)$ with lower right entry $d$, $G|_k\gamma=(2c_0(F_1))^{-1}\sum_a\chi(a^{-1})F_{ad}=\chi(d)G$. The coefficients of $G$ have absolute value at most $1$, since $2$ and $c_0(F_1)$ are units and the values of $\chi$ are roots of unity; and $G$ vanishes at the cusps of (4). Since $-1\in\Gamma_0(N)$, $F_{-1}=F_1|_k(-1)=(-1)^kF_1$, so
\[
c_0(G)=(2c_0(F_1))^{-1}\bigl(c_0(F_1)+\chi(-1)(-1)^kc_0(F_1)\bigr)=1.
\]
\end{enumerate}
\end{proof}

\begin{lemma}\label{sha-automorphy-constant-term}
Let $\ell$ be an odd prime and $\iota:\bar{\mathbf Q}_\ell\cong\mathbf C$.
Let $M>1$ with $\ell\nmid M$, and let $\varphi$ be a primitive Dirichlet character modulo $M$.
Let $w$ be a prime with $w\nmid M$, $\ell\mid w-1$ and $\varphi(w)=1$, let $\theta\ne1$ be a Dirichlet character modulo $w$ with $\theta^\ell=1$, and let $\chi$ be the Dirichlet character modulo $Mw$ with $\chi(n)=\varphi(n)\theta(n)$.
Let $\zeta=e^{2\pi i/Mw}$, let $\tau=\sum_{a\in\mathbf Z/Mw}\chi^{-1}(a)\zeta^a$, let $k\ge1$, and let $L_k$ be as in Lemma~\ref{sha-automorphy-eisenstein-series}.
If $c\in\mathbf C$ satisfies
\[
2\tau c=\sum_{a\in(\mathbf Z/Mw)^\times}\chi^{-1}(a)L_k(\zeta^a),
\]
then $|\iota^{-1}(c)|<1$.
\end{lemma}
\leanref{Development/AutomorphicArithmetic/AutomorphicVarieties/IdeleGenerators/HolomorphicIrreducibility/ConstSmall.lean}{UpperTheory.V.Sd.constSmall}
\leanref{Development/AutomorphicArithmetic/AutomorphicVarieties/IdeleGenerators/HolomorphicIrreducibility/EPrim.lean}{UpperTheory.V.Sd.ePrim}
\begin{proof}
Write $|x|$ for $|\iota^{-1}(x)|$, and $x\equiv y$ for $|x-y|<1$.
The character $\chi$ is primitive, since $\varphi$ is primitive modulo $M$ and $\theta$ is nontrivial modulo the prime $w\nmid M$.
So $\tau\cdot\bigl(\sum_a\chi(a)\zeta^a\bigr)\cdot\chi(-1)=Mw$, which is prime to $\ell$ because $\ell\nmid M$ and $\ell\mid w-1$; as both Gauss sums have absolute value at most $1$, $|\tau|=1$.
Since also $|2|=1$, it suffices that the right-hand side $S$ satisfies $S\equiv0$.
Each $L_k(\zeta^n)$ with $\zeta^n\ne1$ has absolute value at most $1$, by Lemma~\ref{sha-automorphy-eulerian}(1).
Write $S=\sum_{n=0}^{Mw-1}\varphi^{-1}(n)\theta^{-1}(n)L_k(\zeta^n)$, where $\varphi^{-1}(n)=0$ for $n$ not prime to $M$, and similarly for $\theta$.
\begin{enumerate}
\item For $w\nmid n$, $\theta^{-1}(n)$ is an $\ell$-th root of unity, so $\theta^{-1}(n)\equiv1$. Hence $S\equiv S_1-S_w$, where $S_1=\sum_n\varphi^{-1}(n)L_k(\zeta^n)$ and $S_w$ is the same sum over the $n$ divisible by $w$.
\item For $x$ with $x^w\ne1$ and a primitive $w$-th root of unity $\xi$, $\sum_{j=0}^{w-1}L_k(x\xi^j)=w^kL_k(x^w)$. Write $n=b+Mj$ with $0\le b<M$ and $0\le j<w$. Applying this identity to $x=\zeta^b$ and $\xi=\zeta^M$ for $b\ne0$ gives $S_1=w^kT$, with $T=\sum_{0<b<M}\varphi^{-1}(b)L_k(\zeta^{wb})$.
\item Writing $n=wb$ with $0\le b<M$ and using $\varphi(w)=1$ gives $S_w=T$.
\end{enumerate}
So $S\equiv(w^k-1)T$. Since $\ell\mid w-1\mid w^k-1$ and $|T|\le1$, $S\equiv0$.
\end{proof}
% The distribution identity in (2) is Sd.lkDist; the evaluation of the Gauss sums is gauss_mul_gauss_inv and gauss_eq in ConstSmall.lean.

\begin{lemma}[Deligne-Serre \cite{DS74}]\label{sha-automorphy-deligne-serre}
Let $\ell$ be a prime, let $\iota:\bar{\mathbf Q}_\ell\cong\mathbf C$, let $N,k\ge1$, and let $\chi$ be a Dirichlet character modulo $N$.
Let $f$ be a cusp form of weight $k$ for $\Gamma_1(N)$ with $f|_k\gamma=\chi(d)f$ for every $\gamma=\begin{pmatrix}*&*\\ *&d\end{pmatrix}\in\Gamma_0(N)$, whose $q$-expansion coefficients $a_n(f)$ satisfy $|\iota^{-1}(a_n(f))|\le1$ and $|\iota^{-1}(a_1(f))|=1$.
For each prime $p$ let $a_p\in\mathbf C$ with $|\iota^{-1}(a_p)|\le1$, and suppose that for every prime $p$ and every $n\ge1$
\[
\bigl|\iota^{-1}\bigl(a_{pn}(f)+\chi(p)p^{k-1}a_{n/p}(f)-a_pa_n(f)\bigr)\bigr|<1,
\]
where $a_{n/p}(f)=0$ if $p\nmid n$.
Then there is a \emph{normalized eigenform} $g$ of weight $k$ for $\Gamma_1(N)$ with character $\chi$, that is, a cusp form with $g|_k\gamma=\chi(d)g$ as above, $a_1(g)=1$ and $a_{pn}(g)+\chi(p)p^{k-1}a_{n/p}(g)=a_p(g)a_n(g)$ for every prime $p$ and every $n\ge1$, such that $|\iota^{-1}(a_p(g)-a_p)|<1$ for every prime $p$.
\end{lemma}
\leanref{Development/AutomorphicArithmetic/AutomorphicVarieties/IdeleGenerators/HolomorphicIrreducibility/EigLift.lean}{UpperTheory.V.Sd.eigLift}

\begin{proposition}\label{sha-automorphy-eisenstein-eigenform}
Let $\ell$ be an odd prime and $\iota:\bar{\mathbf Q}_\ell\cong\mathbf C$.
Let $M>1$ with $\ell\nmid M$, and let $\varphi$ be a primitive Dirichlet character modulo $M$.
Let $w$ be a prime with $w\nmid M$, $\ell\mid w-1$ and $\varphi(w)=1$, let $\theta\ne1$ be a Dirichlet character modulo $w$ with $\theta^\ell=1$, and let $\chi$ be the Dirichlet character modulo $Mw$ with $\chi(n)=\varphi(n)\theta(n)$.
Then there is $k_0$ such that for every $k\ge k_0$ with $\varphi(-1)=(-1)^k$ there is a normalized eigenform $g$ of weight $k$ for $\Gamma_1(Mw)$ with character $\chi$ (Lemma~\ref{sha-automorphy-deligne-serre}) such that $\bigl|\iota^{-1}\bigl(a_p(g)-1-\chi(p)p^{k-1}\bigr)\bigr|<1$ for every prime $p$.
\end{proposition}
\leanref{Development/AutomorphicArithmetic/AutomorphicVarieties/IdeleGenerators/ConstSmall.lean}{UpperTheory.V.Sd.PropE}
\leanref{Development/AutomorphicArithmetic/AutomorphicVarieties/IdeleGenerators/HolomorphicIrreducibility/SumDivHecke.lean}{UpperTheory.V.Sd.propE}
\leanref{Development/AutomorphicArithmetic/AutomorphicVarieties/IdeleGenerators/HolomorphicIrreducibility/EisChi.lean}{UpperTheory.V.Sd.eisChi}
\leanref{Development/AutomorphicArithmetic/AutomorphicVarieties/IdeleGenerators/HolomorphicIrreducibility/CuspOf.lean}{UpperTheory.V.Sd.cuspOf}
\begin{proof}
Write $|x|$ for $|\iota^{-1}(x)|$, $N=Mw$ and $\zeta=e^{2\pi i/N}$.
The character $\chi$ is primitive, as in the proof of Lemma~\ref{sha-automorphy-constant-term}, and $\chi(-1)=\varphi(-1)$, since $\theta$ has odd order $\ell$.
Also $N\ge3$ and $\ell\nmid N$.
Let $k_1$ be as in Lemma~\ref{sha-automorphy-unit-form} for $N$, let $k_0=\max(k_1,3)$, and let $k\ge k_0$ with $\varphi(-1)=(-1)^k$, so that $\chi(-1)=(-1)^k$.
\begin{enumerate}
\item Lemma~\ref{sha-automorphy-unit-form} gives a form $G$ of weight $k$ for $\Gamma_1(N)$ with character $\chi$, coefficients of absolute value at most $1$, constant term $1$, and vanishing at $\infty$ after $|_k\gamma$ for every $\gamma\in SL_2(\mathbf Z)$ whose lower left entry is not divisible by $N$.
\item Let $\tau=\sum_a\chi^{-1}(a)\zeta^a$, which is nonzero since $\chi$ is primitive, and let $E=(2\tau)^{-1}\sum_{a\in\mathbf Z/N}\chi^{-1}(a)E_{k,a}$, with $E_{k,a}$ as in Lemma~\ref{sha-automorphy-eisenstein-series}. By part (1) of that lemma, $E|_k\gamma=\chi(d)E$ for $\gamma\in\Gamma_0(N)$ with lower right entry $d$; by part (2), $E$ vanishes at the same cusps as $G$, since $\chi^{-1}(a)=0$ for $a$ not a unit. By part (3) and the identity $\sum_a\chi^{-1}(a)\zeta^{ae}=\chi(e)\tau$ for primitive $\chi$, together with $\chi(-1)=(-1)^k$, the coefficients of $E$ are $a_n(E)=\sigma(n)$ for $n\ge1$, where $\sigma(n)=\sum_{e\mid n}\chi(e)e^{k-1}$, and its constant term $c$ satisfies $2\tau c=\sum_a\chi^{-1}(a)L_k(\zeta^a)$.
\item By Lemma~\ref{sha-automorphy-constant-term}, $|c|<1$.
\item The form $F=E-cG$ has character $\chi$, constant term $0$, and vanishes at the cusps of (1). Every cusp is $\gamma\infty$ for some $\gamma\in SL_2(\mathbf Z)$; if $N$ divides the lower left entry of $\gamma$, then $\gamma\in\Gamma_0(N)$ and $F|_k\gamma=\chi(d)F$ vanishes at $\infty$. So $F$ is a cusp form.
\item For $n\ge1$, $a_n(F)=\sigma(n)-c\,a_n(G)$. These have absolute value at most $1$, and $a_1(F)=1-c\,a_1(G)$ is a unit.
\item For every prime $p$ and every $n\ge1$, $\sigma(pn)+\chi(p)p^{k-1}\sigma(n/p)=(1+\chi(p)p^{k-1})\sigma(n)$, with $\sigma(n/p)=0$ if $p\nmid n$, because $e\mapsto\chi(e)e^{k-1}$ is completely multiplicative. With $a_p=1+\chi(p)p^{k-1}$, which has $|a_p|\le1$, it follows that
\[
a_{pn}(F)+\chi(p)p^{k-1}a_{n/p}(F)-a_pa_n(F)=-c\bigl(a_{pn}(G)+\chi(p)p^{k-1}a_{n/p}(G)-a_pa_n(G)\bigr),
\]
which has absolute value less than $1$.
\item Lemma~\ref{sha-automorphy-deligne-serre}, applied to $F$ and the $a_p$, gives $g$.
\end{enumerate}
\end{proof}
% (2) is Sd.eisChi, with eisChi_gs for the Gauss sum identity; (6) is sig_hecke (sum_div_hecke).

\begin{lemma}\label{sha-automorphy-ordinary-forms}
Assume that Galois representations exist, in the sense of Definition~\ref{gal-galois-representations}, relative to the maps satisfying the laws of Definition~\ref{gal-llc-laws}.
Let $\ell$ be an odd prime, let $\iota:\bar{\mathbf Q}_\ell\cong\mathbf C$, and let $k\ge3$.
Let $\xi:\operatorname{Gal}(\bar{\mathbf Q}/\mathbf Q)\to\bar{\mathbf Q}_\ell^\times$ be a character with open kernel, and let $c$ be a function from the primes to $\mathbf Z_{\ge0}$, zero at all but finitely many primes, with $c(\ell)=1$, such that for every prime $p\ne\ell$ the inertia groups at $p$ lie in $\ker\xi$ if and only if $c(p)=0$.
Let $\pi$ be a polarized regular algebraic cuspidal automorphic representation of $GL_2(\mathbf A_{\mathbf Q})$, and let $r:\operatorname{Gal}(\bar{\mathbf Q}/\mathbf Q)\to GL_2(\bar{\mathbf Q}_\ell)$ be continuous and attached to $\pi$ through $\iota$, with $\det r=\chi_{\mathrm{cyc}}^{1-k}\xi$.
Suppose that:
$\pi_p$ is unramified for every $p$ with $c(p)=0$;
for every $p\ne\ell$ with $c(p)\ne0$, $\pi_p$ has a nonzero vector fixed by $K_1(p^{c(p)})$ on which $U_p$ acts by some $\lambda$ with $|\iota^{-1}(\lambda)|=1$;
$\pi_\ell$ has a nonzero vector fixed by $K_1(\ell)$ on which $U_\ell$ acts by some $\lambda$ with $|\iota^{-1}(\lambda)-1|<1$;
and $\pi_\infty$ has the infinitesimal character of the weights $(0,1-k)$.
Then:
\begin{enumerate}
\item $r$ is irreducible;
\item for every prime $p\ne\ell$, $r$ is diagonalizable on the inertia groups at $p$, with diagonal characters $1$ and $\xi$;
\item at every place $A$ above $\ell$, $r$ is upper triangular on the decomposition group $D_A$ in a suitable basis, with diagonal characters $\kappa_0,\kappa_1$ such that $\kappa_0=1$ and $\kappa_1=\chi_{\mathrm{cyc}}^{1-k}$ on the inertia group, and $|\kappa_0(\sigma^{-1})-1|<1$ for every Frobenius element $\sigma\in D_A$.
\end{enumerate}
\end{lemma}
\leanref{Development/AutomorphicArithmetic/AutomorphicVarieties/IdeleGenerators/WeightChange.lean}{UpperTheory.V.Sd.OrdIrr}
\leanref{Development/AutomorphicArithmetic/AutomorphicVarieties/IdeleGenerators/WeightChange.lean}{UpperTheory.V.Sd.OrdAway}
\leanref{Development/AutomorphicArithmetic/AutomorphicVarieties/IdeleGenerators/WeightChange.lean}{UpperTheory.V.Sd.OrdAbove}
\leanref{Development/AutomorphicArithmetic/AutomorphicVarieties/ModularityArgument/ThresholdArguments/DenseThresholdCases/SVROrdIrr.lean}{UpperTheory.V.Sd.ordinary_irreducible_of_printed_hida}
\leanref{Development/AutomorphicArithmetic/AutomorphicVarieties/ModularityArgument/ThresholdArguments/DenseThresholdCases/AutNewform2A.lean}{UpperTheory.V.Z.Win.RR.ordAway_a_alt}
\leanref{Development/AutomorphicArithmetic/AutomorphicVarieties/ModularityArgument/ThresholdArguments/DenseThresholdCases/AutNewform2A.lean}{UpperTheory.V.Z.Win.RR.ordAbove_of_printsH}
\begin{proof}
Parts (2) and (3) are local-global compatibility at the primes away from $\ell$ and at $\ell$ for the representation attached to $\pi$ by Definition~\ref{gal-galois-representations}; for ordinary modular forms they are due to Carayol \cite{Car86} and Wiles \cite[Theorem~2.1.4]{Wil88}.

For (1), suppose that $r$ is reducible, with constituents the continuous characters $c_1$ and $c_2$.
Let $A$ be a place above $\ell$, with $\kappa_0,\kappa_1$ as in (3).
On the inertia group $I_A$, the characteristic polynomials give $\{c_1,c_2\}=\{\kappa_0,\kappa_1\}=\{1,\chi_{\mathrm{cyc}}^{1-k}\}$ pointwise, and since a group is not the union of two proper subgroups, $c_1=\kappa_0$ on all of $I_A$ or $c_1=\kappa_1$ on all of $I_A$.
So $c_1/c_2=\chi_{\mathrm{cyc}}^{d}$ on $I_A$ with $d=\pm(k-1)\ne0$, at every place above $\ell$.
Hence $c_1/c_2$ is the $\ell$-adic character of an algebraic Hecke character $\chi$ of $\mathbf Q$ \cite{Ser68}, of some weight $n$.
At a prime $p\ne\ell$ at which $\pi$ and $\chi$ are unramified, the roots $b,b'$ of the characteristic polynomial of $r$ at the inverse of a Frobenius element are, through $\iota$, the Satake parameters of $\pi_p$, and they satisfy $|b|<p|b'|$ and $|b'|<p|b|$ by the bound of Jacquet and Shalika \cite{JS81a}, applied to a unitary twist of $\pi$.
Since $\iota(c_1/c_2)$ of that element is $b/b'$ or $b'/b$, the value of $\chi$ at $p$ lies strictly between $p^{-1}$ and $p$ in absolute value, which forces $n=0$.
Then $\chi$ has finite order on the local units at $\ell$, so some power $(c_1/c_2)^{e}$ with $e\ge1$ is trivial on $I_A$; but $\chi_{\mathrm{cyc}}^{de}$ is not.
\end{proof}
% (1): chars_of_red, hom_or, Af.algOfCyc_s, GalN.V.ratio_roots with Gal.SV.ratioW_of_hasSat_alt, type_zero_window, Corev.UnrPow.exists_pow_unrAt and cyc_ine. (2) and (3) are RR.ordAway_of_alt and RR.ordAbove_of_alt from Gal.SV.exRS_alt.

\begin{lemma}\label{sha-automorphy-eigenform-representation}
Assume that Galois representations exist, in the sense of Definition~\ref{gal-galois-representations}, relative to the maps satisfying the laws of Definition~\ref{gal-llc-laws}.
Let $\ell$ be an odd prime, let $\iota:\bar{\mathbf Q}_\ell\cong\mathbf C$, let $k\ge3$, let $N\ge1$ with $\ell\nmid N$, and let $\chi$ be a primitive Dirichlet character modulo $N$.
Let $g$ be a normalized eigenform of weight $k$ for $\Gamma_1(N)$ with character $\chi$ (Lemma~\ref{sha-automorphy-deligne-serre}) such that $\bigl|\iota^{-1}\bigl(a_p(g)-1-\chi(p)p^{k-1}\bigr)\bigr|<1$ for every prime $p$.
Let $\psi:\operatorname{Gal}(\bar{\mathbf Q}/\mathbf Q)\to\bar{\mathbf Q}_\ell^\times$ be the character with $\psi(s)=\iota^{-1}\chi(a)$ when $s$ acts on the $N$-th roots of unity by $\zeta\mapsto\zeta^a$.
Then there are $\pi$ and $r$ satisfying the hypotheses of Lemma~\ref{sha-automorphy-ordinary-forms} for this $k$, for $\xi=\psi^{-1}$, and for $c(p)$ the exponent of $p$ in $N$ for $p\ne\ell$, such that
\[
\bigl|\operatorname{tr}r(s^{-1})-1-\chi_{\mathrm{cyc}}(s)^{k-1}\psi(s)\bigr|<1\qquad\text{for every }s\in\operatorname{Gal}(\bar{\mathbf Q}/\mathbf Q).
\]
\end{lemma}
\leanref{Development/AutomorphicArithmetic/AutomorphicVarieties/ModularityArgument/ThresholdArguments/DenseThresholdCases/SVREGOrd.lean}{UpperTheory.V.Sd.EG.exists_ordinary_form_trace_near_one_add_cyc}
\begin{proof}
Let $D=\chi_{\mathrm{cyc}}^{k-1}\psi$.
\begin{enumerate}
\item The cuspidal automorphic representation attached to $g$ \cite[\S5]{Gel75} is a polarized regular algebraic cuspidal representation $\pi$ of $GL_2(\mathbf A_{\mathbf Q})$, with $\pi_\infty$ of the weights $(0,1-k)$, such that for every prime $p\mid N$, $\pi_p$ has a nonzero vector fixed by $K_1(p^{e_p})$, with $e_p$ the exponent of $p$ in $N$, on which $U_p$ acts by $a_p(g)$, and such that for $p\nmid N$, $\pi_p$ is unramified with Satake parameters the roots of $X^2-a_p(g)X+\chi(p)p^{k-1}$.
\item By Definition~\ref{gal-galois-representations}, there is a continuous $r$ attached to $\pi$ through $\iota$.
\item For a prime $p\nmid\ell N$ and a Frobenius element $\sigma$ at $p$, the roots of the characteristic polynomial of $r(\sigma^{-1})$ are, through $\iota$, those of $X^2-a_p(g)X+\chi(p)p^{k-1}$. Since $\iota^{-1}(\chi(p)p^{k-1})=D(\sigma)$, this gives $\operatorname{tr}r(\sigma^{-1})=\iota^{-1}a_p(g)$, so $|\operatorname{tr}r(\sigma^{-1})-1-D(\sigma)|<1$, and $\det r(\sigma^{-1})=D(\sigma)$. These Frobenius elements are dense (Theorem~\ref{lit-chebotarev-degree-one}), and both conditions are closed, so they hold for every $s$; in particular $\det r=D^{-1}=\chi_{\mathrm{cyc}}^{1-k}\psi^{-1}$.
\item For $p\ne\ell$ with $p\mid N$, $\chi(p)=0$, so $a_p(g)\equiv1$ is a unit. For $p\ne\ell$, the inertia groups at $p$ lie in $\ker\psi$ if and only if $p\nmid N$, because $\chi$ is primitive; and $\ker\psi$ is open.
\item At $\ell$, $\pi_\ell$ is unramified, and $a_\ell(g)\equiv1$ while $|\iota^{-1}(\chi(\ell)\ell^{k-1})|<|\ell|$. So $X^2-a_\ell(g)X+\chi(\ell)\ell^{k-1}$ has a root $\lambda$ with $|\iota^{-1}(\lambda)-1|<1$, and the $\ell$-stabilization of the spherical vector of $\pi_\ell$ is fixed by $K_1(\ell)$, with $U_\ell$-eigenvalue $\lambda$.
\end{enumerate}
\end{proof}
% (1) is Sd.adelise with Sd.holRep_alt; (2) is Gal.SV.exRS_alt; (3) uses Published.ChebotarevDegreeOne.dense; (5) is the private lemma stab.

\begin{proposition}\label{sha-automorphy-eisenstein-galois}
Assume that Galois representations exist, in the sense of Definition~\ref{gal-galois-representations}, relative to the maps satisfying the laws of Definition~\ref{gal-llc-laws}.
Let $\ell$ be an odd prime, let $\iota:\bar{\mathbf Q}_\ell\cong\mathbf C$, and let $c\in\operatorname{Gal}(\bar{\mathbf Q}/\mathbf Q)$ be a complex conjugation.
Let $N\ge1$ with $\ell\nmid N$ and let $w$ be a prime with $w\nmid N$.
Let $\varphi,\theta:\operatorname{Gal}(\bar{\mathbf Q}/\mathbf Q)\to\bar{\mathbf Q}_\ell^\times$ be characters such that: $\varphi\ne1$ is trivial on $\operatorname{Gal}(\bar{\mathbf Q}/\mathbf Q(\mu_N))$ and on every element that acts on $\mu_N$ by $\zeta\mapsto\zeta^w$; and $\theta\ne1$ is trivial on $\operatorname{Gal}(\bar{\mathbf Q}/\mathbf Q(\mu_w))$, with $\theta^\ell=1$.
Then there is $k_0$ such that for every $k\ge k_0$ with $\varphi(c)=(-1)^k$ there is a continuous irreducible $\rho:\operatorname{Gal}(\bar{\mathbf Q}/\mathbf Q)\to GL_2(\bar{\mathbf Q}_\ell)$ with the following properties.
\begin{enumerate}
\item $\rho$ is attached through $\iota$ to a regular algebraic cuspidal automorphic representation $\pi$ of $GL_2(\mathbf A_{\mathbf Q})$ such that, at every prime $p$, $\pi_p$ is unramified or all its local components are plain (Definition~\ref{gal-plain}).
\item $\det\rho=\chi_{\mathrm{cyc}}^{k-1}\varphi\theta$.
\item For every $s$, the coefficients of the characteristic polynomial of $\rho(s)$ and of $(X-1)(X-\chi_{\mathrm{cyc}}(s)^{k-1}\varphi(s))$ differ by elements of absolute value less than $1$.
\item For every prime $p\ne\ell$, $\rho$ is diagonalizable on the inertia groups at $p$, with diagonal characters $1$ and $\varphi\theta$.
\item At every place $A$ above $\ell$, $\rho$ is upper triangular on $D_A$ in a suitable basis, with diagonal characters $\kappa_0,\kappa_1$ such that $\kappa_0=\chi_{\mathrm{cyc}}^{k-1}$ and $\kappa_1=1$ on the inertia group, and $|\kappa_0-\chi_{\mathrm{cyc}}^{k-1}\varphi|<1$ and $|\kappa_1-1|<1$ on $D_A$.
\end{enumerate}
\end{proposition}
\leanref{Development/AutomorphicArithmetic/AutomorphicVarieties/ModularityArgument/AutomorphyLifting/PrimeIndexConstruction/MixSeedPAlt.lean}{UpperTheory.V.Z.Win.Birth.PropEGP}
\leanref{Development/AutomorphicArithmetic/AutomorphicVarieties/ModularityArgument/ThresholdArguments/DenseThresholdCases/BirthEG.lean}{UpperTheory.V.Z.Win.Birth.propEGP_of_printsH}
\leanref{Development/AutomorphicArithmetic/AutomorphicVarieties/IdeleGenerators/SlopeTestRemainingChoices.lean}{UpperTheory.V.Sd.EisRep}
\begin{proof}
\begin{enumerate}
\item The character $\varphi$ is the Galois character, through $\iota$, of a primitive Dirichlet character $\varphi_0$ modulo some $M\mid N$, and $\theta$ that of a Dirichlet character $\theta_0$ modulo $w$; here $\varphi_0(a)$ is the value of $\varphi$ at the elements acting on $\mu_M$ by $\zeta\mapsto\zeta^a$.
\item Then $M>1$ since $\varphi\ne1$, $\ell\nmid M$ and $w\nmid M$; $\theta_0\ne1$ and $\theta_0^\ell=1$, so $\theta_0$ has order $\ell$, which divides $w-1$; and $\varphi_0(w)=1$ by the hypothesis on $\varphi$. Let $\chi$ be the Dirichlet character modulo $Mw$ with $\chi(n)=\varphi_0(n)\theta_0(n)$; its Galois character $\psi$ is $\varphi\theta$, and $\ell\nmid Mw$.
\item Let $k_0'$ be as in Proposition~\ref{sha-automorphy-eisenstein-eigenform} for $\varphi_0,\theta_0,\chi$, and let $k_0=\max(k_0',3)$. Let $k\ge k_0$ with $\varphi(c)=(-1)^k$. Since $c$ acts on $\mu_M$ by $\zeta\mapsto\zeta^{-1}$, $\varphi_0(-1)=(-1)^k$, and Proposition~\ref{sha-automorphy-eisenstein-eigenform} gives an eigenform $g$. The character $\chi$ is primitive (proof of Lemma~\ref{sha-automorphy-constant-term}).
\item Lemma~\ref{sha-automorphy-eigenform-representation} gives $\pi$ and $r$, for $\xi=\psi^{-1}$. Let $D=\chi_{\mathrm{cyc}}^{k-1}\psi$, $\Phi=\chi_{\mathrm{cyc}}^{k-1}\varphi$, and $\rho=D\otimes r$. Since $\theta$ takes values in $\ell$-th roots of unity, $|\theta-1|<1$ and $|D-\Phi|<1$.
\item $\rho$ is continuous, irreducible by Lemma~\ref{sha-automorphy-ordinary-forms}(1), and $\det\rho=D^2\det r=D$, which is (2).
\item Since $\det r(s)=D(s)^{-1}$, $\operatorname{tr}\rho(s)=D(s)\operatorname{tr}r(s)=\operatorname{tr}r(s^{-1})$, which is congruent to $1+D(s)$ and so to $1+\Phi(s)$; with $\det\rho(s)=D(s)$ this gives (3).
\item For $p\ne\ell$, $\chi_{\mathrm{cyc}}$ is trivial on the inertia groups at $p$, so $D=\psi$ there, and Lemma~\ref{sha-automorphy-ordinary-forms}(2) gives the diagonal characters $D=\varphi\theta$ and $D\psi^{-1}=1$, which is (4).
\item At a place $A$ above $\ell$, with $\kappa_0,\kappa_1$ as in Lemma~\ref{sha-automorphy-ordinary-forms}(3), $\rho$ is triangular on $D_A$ with diagonal $\kappa_0D$ and $\kappa_1D$. On the inertia group, $\psi$ is trivial because $\ell\nmid Mw$, so $\kappa_0D=\chi_{\mathrm{cyc}}^{k-1}$ and $\kappa_1D=1$. The set of $s\in D_A$ with $|\kappa_0(s)-1|<1$ is a closed subgroup containing the inertia group and a Frobenius element, so it is $D_A$. Since $\kappa_0\kappa_1=\det r=D^{-1}$ on $D_A$, $\kappa_1D=\kappa_0^{-1}$; so $|\kappa_1D-1|<1$, and $\kappa_0D=\kappa_0\theta\Phi$ is congruent to $\Phi$. This is (5).
\item For (1), $D$ is the $\ell$-adic character of an algebraic Hecke character, being a power of $\chi_{\mathrm{cyc}}$ times a character of finite order, so $\rho$ is attached through $\iota$ to the twist of $\pi$ by it, which is regular algebraic and cuspidal. At a prime $p\mid\ell Mw$, the nonzero eigenvalue of $U_p$ on a vector of level $K_1(p^{c(p)})$ given by Lemma~\ref{sha-automorphy-eigenform-representation}, with the condition on the weights at $\ell$, makes every local component of $\pi$ at $p$ a constituent of a principal series that is plain; at the other primes $\pi_p$ is unramified. Twisting by a Hecke character preserves both alternatives.
\end{enumerate}
\end{proof}
% (1) is Sd.dirOf; (4)-(9) are the private lemma core in BirthEG.lean, with above_of and dec_le for (8), Sd.twist_real_selfDual, Z.Win.Birth.prOrd (ordPlain: uLoc, psPlain) and prTw for (9).

\begin{definition}\label{sha-automorphy-has-conjugates}
Let $N,k\ge1$ and let $\chi$ be a Dirichlet character modulo $N$. For a field automorphism $\sigma$ of $\mathbf C$ let $\chi^\sigma=\sigma\circ\chi$, again a Dirichlet character modulo $N$.
A cusp form $f$ of weight $k$ for $\Gamma_1(N)$ with character $\chi$, that is, with $f|_k\gamma=\chi(d)f$ for every $\gamma=\begin{pmatrix}*&*\\ *&d\end{pmatrix}\in\Gamma_0(N)$, \emph{has conjugate forms} if for every field automorphism $\sigma$ of $\mathbf C$ there is a cusp form $f^\sigma$ of weight $k$ for $\Gamma_1(N)$ with character $\chi^\sigma$ whose $q$-expansion coefficients are $a_n(f^\sigma)=\sigma(a_n(f))$ for $n\ge1$.
\end{definition}
\leanref{Development/AutomorphicArithmetic/AutomorphicVarieties/ModularityArgument/SymmetricPowerTransfer/UnlinkedHalfIntegral/UnitFormC.lean}{UpperTheory.V.Z.Sym.Unl.Cj}

\begin{lemma}\label{sha-automorphy-conjugation}
Let $N\ge1$, let $\zeta=e^{2\pi i/N}$, let $\sigma$ be a field automorphism of $\mathbf C$, and for a Dirichlet character $\chi$ modulo $N$ let $\chi^\sigma=\sigma\circ\chi$.
\begin{enumerate}
\item If $\chi$ is primitive, then so is $\chi^\sigma$.
\item If $N>1$, there is $s\in(\mathbf Z/N)^\times$ with $\sigma(\zeta^b)=\zeta^{sb}$ for every $b\in\mathbf Z/N$, and for every $k\ge3$ and $a\in(\mathbf Z/N)^\times$ the $q$-expansion coefficients of the Eisenstein series $E_{k,sa}$ of Lemma~\ref{sha-automorphy-eisenstein-series} are the images under $\sigma$ of those of $E_{k,a}$.
\item Let $N>1$, let $\chi$ be primitive, let $k\ge1$, let $L_k$ be as in Lemma~\ref{sha-automorphy-eisenstein-series}, and let $\tau=\sum_a\chi^{-1}(a)\zeta^a$ and $\tau'=\sum_a(\chi^\sigma)^{-1}(a)\zeta^a$, the sums over $a\in\mathbf Z/N$. If $c,c'\in\mathbf C$ satisfy $2\tau c=\sum_a\chi^{-1}(a)L_k(\zeta^a)$ and $2\tau'c'=\sum_a(\chi^\sigma)^{-1}(a)L_k(\zeta^a)$, then $c'=\sigma(c)$.
\item Let $\ell$ be an odd prime, let $\iota:\bar{\mathbf Q}_\ell\cong\mathbf C$, and suppose that $N\ge3$ and $\ell\nmid N$. Then there is $k_1$ such that for every $k\ge k_1$ and every Dirichlet character $\chi$ modulo $N$ with $\chi(-1)=(-1)^k$ there is a modular form $G$ of weight $k$ for $\Gamma_1(N)$ with the following properties: $G|_k\gamma=\chi(d)G$ for every $\gamma=\begin{pmatrix}*&*\\ *&d\end{pmatrix}\in\Gamma_0(N)$; its $q$-expansion coefficients $b_n$ satisfy $|\iota^{-1}(b_n)|\le1$; $G|_k\gamma$ vanishes at $\infty$ for every $\gamma\in SL_2(\mathbf Z)$ whose lower left entry is not divisible by $N$; $G$ has constant term $1$; and for every field automorphism $\sigma$ of $\mathbf C$ there is a modular form $G^\sigma$ of weight $k$ for $\Gamma_1(N)$ with $G^\sigma|_k\gamma=\chi^\sigma(d)G^\sigma$ for the same $\gamma$, which vanishes at $\infty$ after $|_k\gamma$ for the same $\gamma\in SL_2(\mathbf Z)$, and whose $q$-expansion coefficients are $\sigma(b_n)$.
\end{enumerate}
\end{lemma}
\leanref{Development/AutomorphicArithmetic/AutomorphicVarieties/ModularityArgument/SymmetricPowerTransfer/UnlinkedHalfIntegral/UnlinkedHeckeRoots/PrimCj.lean}{UpperTheory.V.Z.Sym.Unl.primCj}
\leanref{Development/AutomorphicArithmetic/AutomorphicVarieties/ModularityArgument/SymmetricPowerTransfer/UnlinkedHalfIntegral/UnlinkedHeckeRoots/MapLk.lean}{UpperTheory.V.Z.Sym.Unl.eisCj}
\leanref{Development/AutomorphicArithmetic/AutomorphicVarieties/ModularityArgument/SymmetricPowerTransfer/UnlinkedHalfIntegral/UnlinkedHeckeRoots/ConstCj.lean}{UpperTheory.V.Z.Sym.Unl.constCj}
\leanref{Development/AutomorphicArithmetic/AutomorphicVarieties/ModularityArgument/SymmetricPowerTransfer/UnlinkedHalfIntegral/UnlinkedHeckeRoots/UnitFormC.lean}{UpperTheory.V.Z.Sym.Unl.unitFormC}
\leanref{Development/AutomorphicArithmetic/AutomorphicVarieties/ModularityArgument/SymmetricPowerTransfer/UnlinkedHalfIntegral/UnlinkedHeckeRoots/UnitFamC.lean}{UpperTheory.V.Z.Sym.Unl.unitFamC}
\begin{proof}
(1) Since $\sigma$ is injective, $\chi$ and $\chi^\sigma$ have the same kernel on $(\mathbf Z/N)^\times$. So they factor through the same quotients $(\mathbf Z/N)^\times\to(\mathbf Z/d)^\times$, and they have the same conductor.

(2) Since $\sigma(\zeta)$ is a primitive $N$-th root of unity, $\sigma(\zeta)=\zeta^s$ for some $s$ prime to $N$, and then $\sigma(\zeta^b)=\zeta^{sb}$.
The polynomials $A_j$ of Lemma~\ref{sha-automorphy-eisenstein-series} have integer coefficients, so $\sigma(L_k(x))=L_k(\sigma(x))$.
By part (3) of that lemma, applied to $a$ and to $sa$, which are nonzero, $\sigma$ carries the constant term $L_k(\zeta^a)$ of $E_{k,a}$ to $L_k(\zeta^{sa})$, and its $n$-th coefficient $\sum_{e\mid n}e^{k-1}(\zeta^{ae}+(-1)^k\zeta^{-ae})$ to the same sum with $sa$ in place of $a$.

(3) Let $s$ be as in (2), and let $F$ be the identity or $L_k$, so that $\sigma(F(x))=F(\sigma(x))$. Substituting $b=sa$,
\[
\sigma\Bigl(\sum_a\chi^{-1}(a)F(\zeta^a)\Bigr)=\sum_a(\chi^\sigma)^{-1}(a)F(\zeta^{sa})=(\chi^\sigma)^{-1}(s^{-1})\sum_b(\chi^\sigma)^{-1}(b)F(\zeta^b).
\]
For $F$ the identity this gives $\sigma(\tau)=(\chi^\sigma)^{-1}(s^{-1})\tau'$, and for $F=L_k$ it gives $\sigma(2\tau c)=(\chi^\sigma)^{-1}(s^{-1})\,2\tau'c'$.
Since $\chi$ is primitive, $\tau\cdot\sum_a\chi(a)\zeta^a\cdot\chi(-1)=N$, so $\tau\ne0$ and $\sigma(\tau)\ne0$.
Dividing the second identity by the first gives $\sigma(c)=c'$.

(4) Follow the proof of Lemma~\ref{sha-automorphy-unit-form}, and call a family $f=(f_a)_a$ as there \emph{conjugation-stable} if for every $\sigma$ there is a family $f^\sigma=(f^\sigma_a)_a$ of modular forms of the same weight $m$ for $\Gamma_1(N)$, with $f^\sigma_a|_m\gamma=f^\sigma_{ad}$ for every $\gamma\in\Gamma_0(N)$ with lower right entry $d$, whose coefficients are the images under $\sigma$ of those of $f_a$.
By (2), $E_m$ is conjugation-stable, with $(E_m)^\sigma_a=E_{m,sa}$. The constant families $F_m$ have integer coefficients, so $F_m^\sigma=F_m$.
Termwise products and the combinations $f+bg$ of conjugation-stable families are conjugation-stable, with $(fg)^\sigma=f^\sigma g^\sigma$ and $(f+bg)^\sigma=f^\sigma+\sigma(b)g^\sigma$.
So the families of steps (1) to (5) of that proof are conjugation-stable.
For the family $F$ of step (5), each $F^\sigma_a$ vanishes at the cusps of step (4), since it is the product of $E_{3,sa}$, which vanishes there by Lemma~\ref{sha-automorphy-eisenstein-series}(2), with a form that is bounded there.
With $k_1$ and $G$ as in step (6), let
\[
G^\sigma=\sigma\bigl(2c_0(F_1)\bigr)^{-1}\sum_a\chi^\sigma(a^{-1})F^\sigma_a.
\]
The computation of step (6), with $\chi^\sigma$ in place of $\chi$, gives $G^\sigma|_k\gamma=\chi^\sigma(d)G^\sigma$; $G^\sigma$ vanishes at the cusps of step (4); and its coefficients are the images under $\sigma$ of those of $G$.
\end{proof}
% (4): unitFormC from unitFamC, whose conjugate families are the field cg of the private structure Good and the predicate VC (vc_eis, VC.mul, base, killer); (2) is exists_unit and eisCj; (3) is sum_conj and gauss_inv_ne in ConstCj.lean.

\begin{lemma}\label{sha-automorphy-deligne-serre-conjugates}
Let $\ell$ be a prime, let $\iota:\bar{\mathbf Q}_\ell\cong\mathbf C$, let $N,k\ge1$, and let $\chi$ be a Dirichlet character modulo $N$.
Let $f$ be a cusp form of weight $k$ for $\Gamma_1(N)$ with character $\chi$ that has conjugate forms (Definition~\ref{sha-automorphy-has-conjugates}), and whose $q$-expansion coefficients satisfy $|\iota^{-1}(a_n(f))|\le1$ and $|\iota^{-1}(a_1(f))|=1$.
For each prime $p$ let $a_p\in\mathbf C$ with $|\iota^{-1}(a_p)|\le1$, and suppose that for every prime $p$ and every $n\ge1$
\[
\bigl|\iota^{-1}\bigl(a_{pn}(f)+\chi(p)p^{k-1}a_{n/p}(f)-a_pa_n(f)\bigr)\bigr|<1,
\]
where $a_{n/p}(f)=0$ if $p\nmid n$.
Then there is a normalized eigenform $g$ of weight $k$ for $\Gamma_1(N)$ with character $\chi$ (Lemma~\ref{sha-automorphy-deligne-serre}) that has conjugate forms and satisfies $|\iota^{-1}(a_p(g)-a_p)|<1$ for every prime $p$.
\end{lemma}
\leanref{Development/AutomorphicArithmetic/AutomorphicVarieties/ModularityArgument/SymmetricPowerTransfer/UnlinkedHalfIntegral/UnlinkedHeckeRoots/EigLiftC.lean}{UpperTheory.V.Z.Sym.Unl.eigLiftC}
\begin{proof}
Let $W$ be the set of cusp forms of weight $k$ for $\Gamma_1(N)$ with character $\chi$ that have conjugate forms. It is a complex vector space, since $h^\sigma+\sigma(b)h'^\sigma$ is a conjugate form of $h+bh'$.
For a prime $p$ and $h\in W$, the cusp form $T_ph$ with character $\chi$ and coefficients $a_{pn}(h)+\chi(p)p^{k-1}a_{n/p}(h)$ for $n\ge1$ lies in $W$: the same operator for $\chi^\sigma$, applied to $h^\sigma$, gives a cusp form with character $\chi^\sigma$ whose coefficients are the images under $\sigma$ of those of $T_ph$, because $\sigma(\chi(p)p^{k-1})=\chi^\sigma(p)p^{k-1}$.
Let $V$ be the space of the sequences $(\iota^{-1}(a_n(h)))_{n\ge1}$ for $h\in W$. It is finite dimensional over $\bar{\mathbf Q}_\ell$, because the cusp forms of weight $k$ for $\Gamma_1(N)$ form a finite-dimensional space.
On sequences let $T_p$ act by $(T_pv)_n=v_{pn}+\iota^{-1}(\chi(p)p^{k-1})v_{n/p}$, with $v_{n/p}=0$ if $p\nmid n$. These operators preserve $V$ and commute, and since $|\iota^{-1}(\chi(p)p^{k-1})|\le1$ they preserve the sequences whose terms all have absolute value at most $1$, and those whose terms all have absolute value less than $1$.
The sequence $v$ of $f$ has terms of absolute value at most $1$, $|v_1|=1$, and every term of $T_pv-\iota^{-1}(a_p)v$ has absolute value less than $1$.
The lemma of Deligne and Serre \cite[Lemme~6.11]{DS74} gives a nonzero $h\in W$ whose sequence is an eigenvector of every $T_p$, with eigenvalue $b_p$ satisfying $|b_p-\iota^{-1}(a_p)|<1$.
Then $a_{pn}(h)+\chi(p)p^{k-1}a_{n/p}(h)=\iota(b_p)a_n(h)$ for all $p$ and $n\ge1$. If $a_1(h)$ were $0$, these relations would give $a_n(h)=0$ for every $n$, by induction on $n$; so $a_1(h)\ne0$, and $g=a_1(h)^{-1}h$ is a normalized eigenform in $W$ with $a_p(g)=\iota(b_p)$.
\end{proof}
% The subspace W is the private cjSub, heckeC is the stability under T_p, delSerre the lemma of Deligne and Serre, and eq_zero_of_eig the induction.

\begin{proposition}\label{sha-automorphy-eisenstein-conjugates}
Let $\ell$ be an odd prime and $\iota:\bar{\mathbf Q}_\ell\cong\mathbf C$.
Let $M>1$ with $\ell\nmid M$, and let $\varphi$ be a primitive Dirichlet character modulo $M$.
Let $w$ be a prime with $w\nmid M$, $\ell\mid w-1$ and $\varphi(w)=1$, let $\theta\ne1$ be a Dirichlet character modulo $w$ with $\theta^\ell=1$, and let $\chi$ be the Dirichlet character modulo $Mw$ with $\chi(n)=\varphi(n)\theta(n)$.
Then there is $k_0$ such that for every $k\ge k_0$ with $\varphi(-1)=(-1)^k$ there is a normalized eigenform $g$ of weight $k$ for $\Gamma_1(Mw)$ with character $\chi$ (Lemma~\ref{sha-automorphy-deligne-serre}) that has conjugate forms (Definition~\ref{sha-automorphy-has-conjugates}) and satisfies $\bigl|\iota^{-1}\bigl(a_p(g)-1-\chi(p)p^{k-1}\bigr)\bigr|<1$ for every prime $p$.
\end{proposition}
\leanref{Development/AutomorphicArithmetic/AutomorphicVarieties/ModularityArgument/SymmetricPowerTransfer/UnlinkedHalfIntegral/UnitFormC.lean}{UpperTheory.V.Z.Sym.Unl.PropEC}
\leanref{Development/AutomorphicArithmetic/AutomorphicVarieties/ModularityArgument/SymmetricPowerTransfer/UnlinkedHalfIntegral/UnlinkedHeckeRoots/PropEC.lean}{UpperTheory.V.Z.Sym.Unl.propEC}
\begin{proof}
Write $N=Mw$ and $\zeta=e^{2\pi i/N}$.
Follow the proof of Proposition~\ref{sha-automorphy-eisenstein-eigenform}, with $k_1$ from Lemma~\ref{sha-automorphy-conjugation}(4) and $k_0=\max(k_1,3)$, so that the form $G$ of its step (1) also has the forms $G^\sigma$ of Lemma~\ref{sha-automorphy-conjugation}(4).
Let $F=E-cG$ be the cusp form of its step (4). For $n\ge1$, $a_n(F)=S_\chi(n)-c\,b_n$, where $S_\chi(n)=\sum_{e\mid n}\chi(e)e^{k-1}$ is the function written $\sigma(n)$ in that proof, and $b_n$ are the coefficients of $G$.
It suffices that $F$ has conjugate forms: then Lemma~\ref{sha-automorphy-deligne-serre-conjugates}, in place of Lemma~\ref{sha-automorphy-deligne-serre} in step (7), gives $g$.

Let $\sigma$ be a field automorphism of $\mathbf C$. By Lemma~\ref{sha-automorphy-conjugation}(1), $\chi^\sigma$ is primitive, and $\chi^\sigma(-1)=\sigma(\chi(-1))=(-1)^k$.
Step (2) of that proof, for $\chi^\sigma$, gives a modular form $E'$ with character $\chi^\sigma$ that vanishes at the cusps of step (1), with $a_n(E')=S_{\chi^\sigma}(n)=\sigma(S_\chi(n))$ for $n\ge1$, and with constant term $c'$ satisfying $2\tau'c'=\sum_a(\chi^\sigma)^{-1}(a)L_k(\zeta^a)$, where $\tau'=\sum_a(\chi^\sigma)^{-1}(a)\zeta^a$.
By Lemma~\ref{sha-automorphy-conjugation}(3), $c'=\sigma(c)$.
As in step (4), $F^\sigma=E'-c'G^\sigma$ is a cusp form with character $\chi^\sigma$, since $G^\sigma$ has constant term $\sigma(1)=1$; and for $n\ge1$
\[
a_n(F^\sigma)=\sigma(S_\chi(n))-\sigma(c)\sigma(b_n)=\sigma(a_n(F)).
\]
\end{proof}
% F is cusp_comb (with Sd.cuspOf); the conjugate forms are hCj in propEC, with primCj, Sd.eisChi for chi^sigma, constCj and sig_cj.

\begin{lemma}[\cite{Miy89}]\label{sha-automorphy-petersson}
Let $k\ge3$ and $N\ge1$, let $\chi$ be a primitive Dirichlet character modulo $N$, let $g$ be a normalized eigenform of weight $k$ for $\Gamma_1(N)$ with character $\chi$ (Lemma~\ref{sha-automorphy-deligne-serre}), and let $p$ be a prime.
\begin{enumerate}
\item If $p\nmid N$, then $\chi(p)\overline{a_p(g)}=a_p(g)$.
\item If $p\mid N$, then $|a_p(g)|^2=p^{k-1}$.
\end{enumerate}
\end{lemma}
\leanref{Development/AutomorphicArithmetic/AutomorphicVarieties/ModularityArgument/SymmetricPowerTransfer/UnlinkedHalfIntegral/UnlinkedHeckeRoots/SelfAdj.lean}{UpperTheory.V.Z.Sym.Unl.selfAdj}
\leanref{Development/AutomorphicArithmetic/AutomorphicVarieties/ModularityArgument/SymmetricPowerTransfer/UnlinkedHalfIntegral/UnlinkedHeckeRoots/BadAbs.lean}{UpperTheory.V.Z.Sym.Unl.badAbs}

\begin{lemma}\label{sha-automorphy-half-step}
Let $p\ge2$ and $m\ge0$ be integers, let $a,x,\alpha,\beta\in\mathbf C$, and let $\sigma$ be a field automorphism of $\mathbf C$ with $\sigma(\sqrt p)=-\sqrt p$.
Suppose that $|x|=|\sigma(x)|=1$, $\bar ax=a$, $\overline{\sigma(a)}\sigma(x)=\sigma(a)$, $\alpha+\beta=a$ and $\alpha\beta=xp^m$.
Then $|\alpha|^2\ne p|\beta|^2$.
\end{lemma}
\leanref{Development/AutomorphicArithmetic/AutomorphicVarieties/ModularityArgument/SymmetricPowerTransfer/UnlinkedHalfIntegral/UnlinkedHeckeRoots/HalfAlg.lean}{UpperTheory.V.Z.Sym.Unl.halfAlg}
\begin{proof}
Suppose that $|\alpha|^2=p|\beta|^2$. Since $\alpha\beta=xp^m\ne0$, $\alpha$ and $\beta$ are nonzero; let $z=\alpha/\beta$, so that $z\bar z=p$.
From $\bar ax=a$ follows $a^2=a\bar ax=|a|^2x$, so
\[
z+2+z^{-1}=\frac{(\alpha+\beta)^2}{\alpha\beta}=\frac{|a|^2x}{xp^m}=\frac{|a|^2}{p^m},
\]
a real number $t\ge0$.
Since $z^{-1}=\bar z/p$, the imaginary part of the left side is $(1-1/p)\operatorname{Im}z$, so $z$ is real, and $z=\pm\sqrt p$.
If $z=-\sqrt p$, then $t=2-\sqrt p-1/\sqrt p=-(\sqrt p-1)^2/\sqrt p<0$, since $p\ge2$.
So $\alpha=\sqrt p\,\beta$ and $a=(1+\sqrt p)\beta$, and therefore $\sqrt p\,a^2=(1+\sqrt p)^2\sqrt p\,\beta^2=(1+\sqrt p)^2xp^m$.
Applying $\sigma$ gives $-\sqrt p\,\sigma(a)^2=(1-\sqrt p)^2\sigma(x)p^m$.
From $\overline{\sigma(a)}\sigma(x)=\sigma(a)$ follows $\sigma(a)^2=|\sigma(a)|^2\sigma(x)$, and $\sigma(x)\ne0$, so $-\sqrt p\,|\sigma(a)|^2=(1-\sqrt p)^2p^m$.
The left side is at most $0$ and the right side is positive, a contradiction.
\end{proof}
% The private lemmas sq_of_conj, ratio_eq (z is sqrt p) and abs_ratio_ne_of_sqrt_flip.

\begin{lemma}\label{sha-automorphy-hecke-roots}
Let $k\ge3$ and $N\ge1$, let $\chi$ be a primitive Dirichlet character modulo $N$, and let $g$ be a normalized eigenform of weight $k$ for $\Gamma_1(N)$ with character $\chi$ (Lemma~\ref{sha-automorphy-deligne-serre}) that has conjugate forms (Definition~\ref{sha-automorphy-has-conjugates}).
Let $p$ be a prime with $p\nmid N$, and let $\alpha,\beta\in\mathbf C$ with $\alpha+\beta=a_p(g)$ and $\alpha\beta=\chi(p)p^{k-1}$.
Then $|\alpha|^2\ne p|\beta|^2$.
\end{lemma}
\leanref{Development/AutomorphicArithmetic/AutomorphicVarieties/ModularityArgument/SymmetricPowerTransfer/UnlinkedHalfIntegral/UnlinkedHeckeRoots/GoodHalf.lean}{UpperTheory.V.Z.Sym.Unl.goodHalf}
\leanref{Development/AutomorphicArithmetic/AutomorphicVarieties/ModularityArgument/SymmetricPowerTransfer/UnlinkedHalfIntegral/UnlinkedHeckeRoots/ConjEig.lean}{UpperTheory.V.Z.Sym.Unl.conjEig}
\leanref{Development/AutomorphicArithmetic/AutomorphicVarieties/ModularityArgument/SymmetricPowerTransfer/UnlinkedHalfIntegral/UnlinkedHeckeRoots/SqrtFlip.lean}{UpperTheory.V.Z.Sym.Unl.sqrtFlip}
\begin{proof}
Since $\sqrt p$ is irrational, $\mathbf Q(\sqrt p)$ has an automorphism with $\sqrt p\mapsto-\sqrt p$, and it extends to a field automorphism $\sigma$ of $\mathbf C$.
Let $g^\sigma$ be the conjugate form of $g$ for $\sigma$. Then $a_1(g^\sigma)=\sigma(1)=1$, and applying $\sigma$ to the relations $a_{qn}(g)+\chi(q)q^{k-1}a_{n/q}(g)=a_q(g)a_n(g)$, for primes $q$ and $n\ge1$, gives the same relations for $g^\sigma$ with $\chi^\sigma$ in place of $\chi$.
So $g^\sigma$ is a normalized eigenform with character $\chi^\sigma$, which is primitive by Lemma~\ref{sha-automorphy-conjugation}(1).
Write $a=a_p(g)$ and $x=\chi(p)$, so that $a_p(g^\sigma)=\sigma(a)$ and $\chi^\sigma(p)=\sigma(x)$.
Lemma~\ref{sha-automorphy-petersson}(1), applied to $g$ and to $g^\sigma$, gives $\bar ax=a$ and $\overline{\sigma(a)}\sigma(x)=\sigma(a)$.
Since $p\nmid N$, $x$ and $\sigma(x)$ are roots of unity.
Lemma~\ref{sha-automorphy-half-step}, with $m=k-1$, gives the claim.
\end{proof}

\begin{lemma}\label{sha-automorphy-frobenius-ratios}
Assume that Galois representations exist, in the sense of Definition~\ref{gal-galois-representations}, relative to the maps satisfying the laws of Definition~\ref{gal-llc-laws}.
Let $\ell$ be an odd prime, let $\iota:\bar{\mathbf Q}_\ell\cong\mathbf C$, let $k\ge3$, let $N\ge1$ with $\ell\nmid N$, and let $\chi$ be a primitive Dirichlet character modulo $N$.
Let $g$ be a normalized eigenform of weight $k$ for $\Gamma_1(N)$ with character $\chi$ (Lemma~\ref{sha-automorphy-deligne-serre}) that has conjugate forms (Definition~\ref{sha-automorphy-has-conjugates}) and satisfies $\bigl|\iota^{-1}\bigl(a_p(g)-1-\chi(p)p^{k-1}\bigr)\bigr|<1$ for every prime $p$.
Let $\psi:\operatorname{Gal}(\bar{\mathbf Q}/\mathbf Q)\to\bar{\mathbf Q}_\ell^\times$ be the character with $\psi(s)=\iota^{-1}\chi(a)$ when $s$ acts on the $N$-th roots of unity by $\zeta\mapsto\zeta^a$.
Then there are $\pi$ and $r$ satisfying the hypotheses of Lemma~\ref{sha-automorphy-ordinary-forms} for this $k$, for $\xi=\psi^{-1}$, and for $c(p)$ the exponent of $p$ in $N$ for $p\ne\ell$, with
\[
\bigl|\operatorname{tr}r(s^{-1})-1-\chi_{\mathrm{cyc}}(s)^{k-1}\psi(s)\bigr|<1\qquad\text{for every }s\in\operatorname{Gal}(\bar{\mathbf Q}/\mathbf Q),
\]
such that in addition, for every prime $p\ne\ell$, every place $A$ of $\bar{\mathbf Q}$ above $p$, every $s$ in the decomposition group $D_A$ that acts on the residue field of $A$ by $x\mapsto x^p$, and every two roots $x,y$ of the characteristic polynomial of $r(s)$,
\[
|\iota(x)|<p\,|\iota(y)|\qquad\text{and}\qquad|\iota(x)|^2\ne p\,|\iota(y)|^2.
\]
\end{lemma}
\leanref{Development/AutomorphicArithmetic/AutomorphicVarieties/ModularityArgument/ThresholdArguments/DenseThresholdCases/SVOrdExT.lean}{UpperTheory.V.Sd.EG.ordExT_sv_alt}
\leanref{Development/AutomorphicArithmetic/AutomorphicVarieties/ModularityArgument/ThresholdArguments/DenseThresholdCases/OrdWinA.lean}{UpperTheory.V.Z.Sym.W.ordWin_a_alt}
\leanref{Development/AutomorphicArithmetic/AutomorphicVarieties/ModularityArgument/ThresholdArguments/DenseThresholdCases/THalfOf.lean}{UpperTheory.V.Z.Sym.Unl.half_of_printed_hida}
\leanref{Development/AutomorphicArithmetic/AutomorphicVarieties/ModularityArgument/ThresholdArguments/DenseThresholdCases/OrdWinA.lean}{UpperTheory.V.Z.Sym.W.tieS_a_alt}
\leanref{Development/AutomorphicArithmetic/AutomorphicVarieties/ModularityArgument/ThresholdArguments/DenseThresholdCases/OrdWinA.lean}{UpperTheory.V.Z.Sym.W.tieU_a_alt}
\begin{proof}
Let $\pi$ and $r$ be those of the proof of Lemma~\ref{sha-automorphy-eigenform-representation}. By its step (1), for every prime $p\mid N$, $\pi_p$ has a nonzero vector fixed by $K_1(p^{e_p})$, with $e_p\ge1$ the exponent of $p$ in $N$, on which $U_p$ acts by $a_p(g)$; and for $p\nmid N$, $\pi_p$ is unramified with Satake parameters the roots of $X^2-a_p(g)X+\chi(p)p^{k-1}$.
Let $p$, $A$, $s$, $x$ and $y$ be as in the statement. The roots of the characteristic polynomial of $r(s^{-1})$ are $x^{-1}$ and $y^{-1}$.
Since $\det r=\chi_{\mathrm{cyc}}^{1-k}\psi^{-1}$, $\chi_{\mathrm{cyc}}(s)=p$ and $\psi$ has finite order, $|\iota(\det r(s^{-1}))|=p^{k-1}$.

For the first inequality: every local component of $\pi$ is unramified or plain, by step (9) of the proof of Proposition~\ref{sha-automorphy-eisenstein-galois}. At such a prime $p\ne\ell$, local-global compatibility \cite{Car86} identifies the roots of the characteristic polynomial of $r(s)$, through $\iota$, with the values at a uniformizer of the characters of the supercuspidal support of $\pi_p$, up to a common factor; the bound of Jacquet and Shalika \cite{JS81a}, applied to a unitary twist of $\pi$, puts the ratio of these values strictly between $p^{-1}$ and $p$ in absolute value.

For the second, suppose that $|\iota(x)|^2=p|\iota(y)|^2$.
If $p\mid N$, then $|a_p(g)|^2=p^{k-1}$ by Lemma~\ref{sha-automorphy-petersson}(2), so $a_p(g)\ne0$. By local-global compatibility at $p$ \cite{Car86}, the vector of $\pi_p$ fixed by $K_1(p^{e_p})$ with the nonzero $U_p$-eigenvalue $a_p(g)$ gives a root $u_0$ of the characteristic polynomial of $r(s^{-1})$ with $|\iota(u_0)|=|a_p(g)|$. If $u\ne u_0$ is the other root, then $|\iota(u_0)||\iota(u)|=p^{k-1}=|a_p(g)|^2$, so $|\iota(u)|=|a_p(g)|$ as well. Hence $|\iota(x)|=|\iota(y)|\ne0$, and $|\iota(x)|^2=p|\iota(y)|^2$ gives $p=1$, which is absurd.
If $p\nmid N$, the roots of the characteristic polynomial of $r(s^{-1})$ are, through $\iota$, those of $X^2-a_p(g)X+\chi(p)p^{k-1}$ (step (3) of the proof of Lemma~\ref{sha-automorphy-eigenform-representation}). If $x=y$, the same contradiction arises. Otherwise $\alpha=\iota(y^{-1})$ and $\beta=\iota(x^{-1})$ are the two roots of this polynomial, so $\alpha+\beta=a_p(g)$, $\alpha\beta=\chi(p)p^{k-1}$ and $|\alpha|^2=p|\beta|^2$, which contradicts Lemma~\ref{sha-automorphy-hecke-roots}.
\end{proof}
% The first inequality is GalN.V.RW, from ordWin_a_alt: GalN.V.L.rwwd and the field ratio of Gal.SV.IsGalRS (from exRS_alt) at the places where pi is plain (plainAll, from Z.Win.Birth.ordPlain and the field plain). The second is Unl.NoHalf, from half_of_printed_hida: the case p | N uses badAbs, tieU_a_alt (jacqOfU_alt, exists_ps, eigW_alt) and det_inv; the case p not dividing N uses tieS_a_alt (the field frob), quad_pair and goodHalf. The U_p eigenvalues and the Satake parameters are the conjuncts hU and hS of ordExT_sv_alt.

\begin{definition}\label{sha-automorphy-tate-free}
Let $F$ be a number field, let $\ell$ be a prime, let $n\ge1$, and let $r:\operatorname{Gal}(\bar F/F)\to GL_n(\bar{\mathbf Q}_\ell)$ be a homomorphism. For a finite place $w$ of $F$ let $q_w$ be the order of its residue field.
Then $r$ is \emph{Tate-free} if for every finite place $w$ of $F$ not above $\ell$ and every place $A$ of $\bar F$ above $w$ such that $r$ maps the inertia group $I_A$ to a finite set, there are $j\ge1$ and an element $s$ of the decomposition group $D_A$ that acts on the residue field of $A$ by $x\mapsto x^{q_w^j}$ such that $\alpha^m\ne(q_w^j\beta)^m$ for all roots $\alpha,\beta$ of the characteristic polynomial of $r(s)$ and all $m\ge1$.
\end{definition}
\leanref{Development/AutomorphicArithmetic/Foundations/GaloisRepresentations/AdmissibleResidues/IsAutomorphicOverImagQuad.lean}{UpperTheory.GalN.V.Unlinked}

\begin{lemma}\label{sha-automorphy-adjoint-tate-free}
Let $F$ be a number field, let $\ell$ be a prime, let $\iota:\bar{\mathbf Q}_\ell\cong\mathbf C$, let $r:\operatorname{Gal}(\bar F/F)\to GL_2(\bar{\mathbf Q}_\ell)$ be a homomorphism, and let $\operatorname{ad}^0r=\operatorname{Sym}^2r\otimes(\det r)^{-1}$.
Suppose that for every finite place $w$ of $F$ not above $\ell$, with residue field of order $q_w$, every place $A$ of $\bar F$ above $w$, every $s\in D_A$ that acts on the residue field of $A$ by $x\mapsto x^{q_w}$, and all roots $\alpha,\beta$ of the characteristic polynomial of $r(s)$,
\[
|\iota(\alpha)|<q_w|\iota(\beta)|\qquad\text{and}\qquad|\iota(\alpha)|^2\ne q_w|\iota(\beta)|^2.
\]
Then $\operatorname{ad}^0r$ is Tate-free (Definition~\ref{sha-automorphy-tate-free}).
\end{lemma}
\leanref{Development/AutomorphicArithmetic/AutomorphicVarieties/IdeleGenerators/SlopeTestRemainingChoices.lean}{UpperTheory.V.Sd.ad0}
\leanref{Development/AutomorphicArithmetic/AutomorphicVarieties/ModularityArgument/SymmetricPowerTransfer/UnlinkedHalfIntegral/UnlinkedHeckeRoots/SplitTwo.lean}{UpperTheory.V.Z.Sym.Unl.adUnl}
\begin{proof}
Let $w$ and $A$ be as in Definition~\ref{sha-automorphy-tate-free}; the finiteness of $r(I_A)$ is not used. Take $j=1$ and $s\in D_A$ acting on the residue field by $x\mapsto x^{q}$, where $q=q_w$.
Let $a,b$ be the roots of the characteristic polynomial of $r(s)$, which are nonzero, and let $A_0=|\iota(a)|$ and $B_0=|\iota(b)|$.
The hypothesis gives $A_0<qB_0$, $B_0<qA_0$, $A_0^2\ne qB_0^2$ and $B_0^2\ne qA_0^2$, and, for $\alpha=\beta=a$, $q>1$.
The characteristic polynomial of $\operatorname{ad}^0r(s)$ is $(X-a/b)(X-1)(X-b/a)$, so the absolute values through $\iota$ of its roots lie in $\{A_0/B_0,1,B_0/A_0\}$.
If $\alpha^m=(q\beta)^m$ for roots $\alpha,\beta$ and some $m\ge1$, then $|\iota(\alpha)|=q|\iota(\beta)|$.
Among the nine pairs of values, $x=qx$ contradicts $q>1$; $A_0/B_0=q$ and $1=qB_0/A_0$ contradict $A_0<qB_0$; $B_0/A_0=q$ and $1=qA_0/B_0$ contradict $B_0<qA_0$; $A_0/B_0=qB_0/A_0$ contradicts $A_0^2\ne qB_0^2$; and $B_0/A_0=qA_0/B_0$ contradicts $B_0^2\ne qA_0^2$.
\end{proof}
% The characteristic polynomial of ad0 is the fifth part of Sd.ad0Facts, with split_two; the nine cases are the private ratio_ne.

\begin{proposition}\label{sha-automorphy-eisenstein-tate-free}
Assume that Galois representations exist, in the sense of Definition~\ref{gal-galois-representations}, relative to the maps satisfying the laws of Definition~\ref{gal-llc-laws}.
Let $\ell$ be an odd prime, let $\iota:\bar{\mathbf Q}_\ell\cong\mathbf C$, and let $c\in\operatorname{Gal}(\bar{\mathbf Q}/\mathbf Q)$ be a complex conjugation.
Let $N\ge1$ with $\ell\nmid N$ and let $w$ be a prime with $w\nmid N$.
Let $\varphi,\theta:\operatorname{Gal}(\bar{\mathbf Q}/\mathbf Q)\to\bar{\mathbf Q}_\ell^\times$ be characters such that: $\varphi\ne1$ is trivial on $\operatorname{Gal}(\bar{\mathbf Q}/\mathbf Q(\mu_N))$ and on every element that acts on $\mu_N$ by $\zeta\mapsto\zeta^w$; and $\theta\ne1$ is trivial on $\operatorname{Gal}(\bar{\mathbf Q}/\mathbf Q(\mu_w))$, with $\theta^\ell=1$.
Then there is $k_0$ such that for every $k\ge k_0$ with $\varphi(c)=(-1)^k$ there is a continuous irreducible $\rho:\operatorname{Gal}(\bar{\mathbf Q}/\mathbf Q)\to GL_2(\bar{\mathbf Q}_\ell)$ with the following properties.
\begin{enumerate}
\item $\rho$ is attached through $\iota$ to a regular algebraic cuspidal automorphic representation $\pi$ of $GL_2(\mathbf A_{\mathbf Q})$ such that, at every prime $p$, $\pi_p$ is unramified or all its local components are plain (Definition~\ref{gal-plain}).
\item $\det\rho=\chi_{\mathrm{cyc}}^{k-1}\varphi\theta$.
\item For every $s$, the coefficients of the characteristic polynomial of $\rho(s)$ and of $(X-1)(X-\chi_{\mathrm{cyc}}(s)^{k-1}\varphi(s))$ differ by elements of absolute value less than $1$.
\item For every prime $p\ne\ell$, $\rho$ is diagonalizable on the inertia groups at $p$, with diagonal characters $1$ and $\varphi\theta$.
\item At every place $A$ above $\ell$, $\rho$ is upper triangular on $D_A$ in a suitable basis, with diagonal characters $\kappa_0,\kappa_1$ such that $\kappa_0=\chi_{\mathrm{cyc}}^{k-1}$ and $\kappa_1=1$ on the inertia group, and $|\kappa_0-\chi_{\mathrm{cyc}}^{k-1}\varphi|<1$ and $|\kappa_1-1|<1$ on $D_A$.
\item For every prime $p\ne\ell$, every place $A$ above $p$, every $s\in D_A$ that acts on the residue field of $A$ by $x\mapsto x^p$, and every two roots $x,y$ of the characteristic polynomial of $\rho(s)$, $|\iota(x)|<p\,|\iota(y)|$ and $|\iota(x)|^2\ne p\,|\iota(y)|^2$.
\item $\operatorname{ad}^0\rho=\operatorname{Sym}^2\rho\otimes(\det\rho)^{-1}$ is Tate-free (Definition~\ref{sha-automorphy-tate-free}).
\end{enumerate}
\end{proposition}
\leanref{Development/AutomorphicArithmetic/AutomorphicVarieties/ModularityArgument/SymmetricPowerTransfer/WeilDeligneUnramified/EGCore.lean}{UpperTheory.V.Z.Sym.W.PropEGc}
\leanref{Development/AutomorphicArithmetic/AutomorphicVarieties/ModularityArgument/ThresholdArguments/DenseThresholdCases/WEGCore.lean}{UpperTheory.V.Z.Sym.W.prop_eg_complex_of_printed_hida}
\leanref{Development/AutomorphicArithmetic/AutomorphicVarieties/ModularityArgument/SymmetricPowerTransfer/WeilDeligneUnramified/UnitaryTwoForm.lean}{UpperTheory.V.Z.Sym.W.PropEGW}
\leanref{Development/AutomorphicArithmetic/AutomorphicVarieties/ModularityArgument/ThresholdArguments/DenseThresholdCases/TPropEGW.lean}{UpperTheory.V.Z.Sym.W.prop_eg_w_of_printed_hida}
\leanref{Development/AutomorphicArithmetic/AutomorphicVarieties/ModularityArgument/SymmetricPowerTransfer/UnlinkedHalfIntegral/UnlinkedHeckeRoots/RootOfSmul.lean}{UpperTheory.V.Z.Sym.Unl.twOk}
\begin{proof}
Follow the proof of Proposition~\ref{sha-automorphy-eisenstein-galois}, with two changes.
In its step (3), Proposition~\ref{sha-automorphy-eisenstein-conjugates} replaces Proposition~\ref{sha-automorphy-eisenstein-eigenform}, so that the eigenform $g$ also has conjugate forms.
In its step (4), Lemma~\ref{sha-automorphy-frobenius-ratios} replaces Lemma~\ref{sha-automorphy-eigenform-representation}, so that $r$ also satisfies the two inequalities of that lemma.
Steps (5) to (9) then give (1) to (5).
For (6): $\rho=D\otimes r$, so the roots of the characteristic polynomial of $\rho(s)$ are $D(s)$ times those of $r(s)$, and multiplying both roots by the same nonzero number preserves the two inequalities.
For (7): a finite place of $\mathbf Q$ not above $\ell$ is a prime $p\ne\ell$, with $q_p=p$, so (6) is the hypothesis of Lemma~\ref{sha-automorphy-adjoint-tate-free} for $\rho$.
\end{proof}
% prop_eg_complex_of_printed_hida repeats propEGP_of_printsH with the private core of WEGCore.lean, which uses propEC (through hcj) and ordExT_sv_alt; (6) is twOk applied to ordWin_a_alt and half_of_printed_hida; (7) is prop_eg_w_of_printed_hida with Unl.adUnl.

\subsection{Twists and automorphy of the adjoint}\label{subsection-sha-adjoint-automorphy}

\begin{lemma}\label{sha-automorphy-twist-glue}
Let $I,J,K,B$ be abelian groups, written multiplicatively, with $B$ divisible. Let $i:J\to I$, $d:K\to I$, $c:I\to I$, $\sigma:K\to K$, $N:J\to B$ and $g:I\to B$ be homomorphisms, let $U,V\subset I$ be subgroups, and let $n\ge0$ and $M\ge1$ be integers. Suppose that:
\begin{enumerate}
\item $c(i(a))=i(a)$ for $a\in J$, $c(d(x))=d(\sigma(x))$ for $x\in K$, and $c(U)\subset U$;
\item $g(i(a))=N(a)$ for every $a\in J$ with $i(a)\in U$;
\item for every $x\in K$ with $\sigma(x)=x$ there is $a\in J$ with $i(a)=d(x)$ and $N(a)=1$;
\item for $x\in K$, if $d(\sigma(x)/x)\in U$ then $(\sigma(x)/x)^n=1$, and if $d(\sigma(x)/x)\in V$ then $\sigma(x)=x$;
\item $u^M\in V$ for every $u\in U$.
\end{enumerate}
Then there are homomorphisms $\psi,\theta:I\to B$ and an integer $k\ge1$ such that $\psi\circ i=N$, $\psi\circ d=1$, $\psi(u)=g(u)\,\theta(u/c(u))$ for $u\in U$, $\theta=1$ on $V$, and $\theta^k=1$ on $U$. If $\ell$ is a prime not dividing $n$, $k$ can be taken prime to $\ell$.
\end{lemma}
\leanref{Development/AutomorphicArithmetic/AutomorphicVarieties/AutomorphicForms/TwistedAutomorphicForms/GlueExt.lean}{UpperTheory.V.Af.Tw.glue}
\begin{proof}
A homomorphism $f:G\to B$ from an abelian group that is trivial on the kernel of a homomorphism $p:G\to H$ of abelian groups is $F\circ p$ for a homomorphism $F:H\to B$: it factors through $p(G)$, and extends from $p(G)$ to $H$ because $B$ is divisible.
Choose integers $e,k\ge1$ with $e\equiv1\pmod n$ and $M\mid ek$. If $\ell\nmid n$, write $M=\ell^aM'$ with $\ell\nmid M'$, take $e$ with $e\equiv1\pmod n$ and $e\equiv0\pmod{\ell^a}$, and $k=M'$; otherwise take $e=1$ and $k=M$.
Let $p_0:J\times K\times U\to I$ be $p_0(a,x,u)=i(a)d(x)u$, and let $Z=\ker p_0$.
If $(a,x,u)\in Z$, then $u=(i(a)d(x))^{-1}$ and $c(u)=(i(a)d(\sigma(x)))^{-1}$ by (1), so $u/c(u)=d(\sigma(x)/x)$.
Define $p_1:Z\times V\to I$ and $f_1:Z\times V\to B$ by $p_1((a,x,u),v)=d(\sigma(x)/x)\,v$ and $f_1((a,x,u),v)=(N(a)g(u))^{-1}$.
Then $f_1$ is trivial on $\ker p_1$. Indeed, if $d(\sigma(x)/x)v=1$, then $d(\sigma(x)/x)=v^{-1}\in V$, so $\sigma(x)=x$ by (4); (3) gives $a'$ with $i(a')=d(x)$ and $N(a')=1$; then $i(aa')=i(a)d(x)=u^{-1}\in U$, and (2) gives $N(a)=N(aa')=g(u)^{-1}$.
So $f_1=\theta_1\circ p_1$ for a homomorphism $\theta_1:I\to B$. Let $\theta=\theta_1^e$.
Evaluating at $(1,v)$ shows that $\theta_1=1$ on $V$.
For $(a,x,u)\in Z$, evaluating at $((a,x,u),1)$ gives $\theta_1(u/c(u))=\theta_1(d(\sigma(x)/x))=(N(a)g(u))^{-1}$. Since $u/c(u)\in U$ by (1), (4) gives $(\sigma(x)/x)^n=1$, so $(N(a)g(u))^{-n}=\theta_1(d((\sigma(x)/x)^n))=1$; as $e\equiv1\pmod n$, $\theta(u/c(u))=(N(a)g(u))^{-1}$.
So the homomorphism $f:J\times K\times U\to B$, $f(a,x,u)=N(a)g(u)\theta(u/c(u))$, is trivial on $Z$, and $f=\psi\circ p_0$ for a homomorphism $\psi:I\to B$.
Evaluating at $(a,1,1)$, $(1,x,1)$ and $(1,1,u)$ gives $\psi(i(a))=N(a)$, $\psi(d(x))=1$ and $\psi(u)=g(u)\theta(u/c(u))$.
Finally, $\theta=\theta_1^e=1$ on $V$, and for $u\in U$, $\theta(u)^k=\theta_1(u^{ek})=1$, because $M\mid ek$ and $u^M\in V$.
\end{proof}
% The extension through the divisible group is the private glue_ext (KStar.Choice.IsDiv.exists_extension); the choice of e and k is the private glue_exp.

\begin{lemma}[Global class field theory \cite{Tat67}]\label{sha-automorphy-twist-artin}
Let $E$ be a number field, let $\bar E$ be an algebraic closure of $E$, let $G_E=\operatorname{Gal}(\bar E/E)$, and let $B$ be a Hausdorff abelian topological group. For a finite place $v$ of $E$, an idele is \emph{placed at $v$} if its components at the other places are $1$, and $q_v$ is the order of the residue field of $v$.
There is a homomorphism $\operatorname{art}_E:\mathbf A_E^\times\to G_E^{\mathrm{ab}}$, a continuous open surjection whose kernel is the closure of $E^\times D_E$, where $D_E$ is the connected component of $1$ in $(E\otimes\mathbf R)^\times$, such that for every finite place $v$ of $E$ and every place $A$ of $\bar E$ above $v$, $\operatorname{art}_E$ maps the units of $E_v$ placed at $v$ onto the image of the inertia group $I_A$, and the uniformizers of $E_v$ placed at $v$ onto the images of the inverses of the Frobenius elements at $A$, the elements of the decomposition group $D_A$ that act on the residue field of $A$ by $x\mapsto x^{q_v}$. Moreover:
\begin{enumerate}
\item if $\nu:G_E\to B$ is a continuous homomorphism trivial on the complex conjugations, then $\nu\circ\operatorname{art}_E$ is trivial on the ideles whose finite components are $1$;
\item if $\psi:\mathbf A_E^\times\to B$ is a continuous homomorphism trivial on $E^\times$ and on the ideles whose finite components are $1$, there is a continuous homomorphism $\kappa:G_E\to B$ with $\kappa(\sigma)=\psi(x)$ whenever $\operatorname{art}_E(x)$ is the image of $\sigma$;
\item two continuous homomorphisms $G_E\to B$ that agree on the Frobenius elements at the places above all but finitely many finite places of $E$ are equal.
\end{enumerate}
\end{lemma}
\leanref{Development/AbelianVarietyTheory/PublishedTheorems/AutomorphicForms₂/ExistsUniqueHcharOfFiniteOrder.lean}{GlobalTheory.Published.GlobalArtinReciprocity}
\leanref{Development/AutomorphicArithmetic/AutomorphicVarieties/AutomorphicForms/TwistedAutomorphicForms/InfTriv.lean}{UpperTheory.V.Af.Tw.infTriv}
\leanref{Development/AutomorphicArithmetic/AutomorphicVarieties/AutomorphicForms/TwistedAutomorphicForms/Descend.lean}{UpperTheory.V.Af.Tw.descend}
\leanref{Development/AbelianVarietyTheory/Proofs/IwasawaDecomposition/InertialOrderFifteenCharAuxiliaryPrimeSpec.lean}{GlobalTheory.GalChar.eq_of_frob}
\begin{proof}
The first statement is global class field theory \cite{Tat67}, \cite[Chapter~VI]{Neu99}.
(1) $\operatorname{art}_E$ is trivial on $D_E$ and maps $-1$ at a real place to the image of a complex conjugation, and the ideles whose finite components are $1$ are generated by $D_E$ and these elements.
(2) $\psi$ is trivial on $E^\times D_E$, hence on its closure, which is the kernel of the open quotient map $\operatorname{art}_E$.
(3) If $\phi$ is the quotient of the two homomorphisms, $\phi\circ\operatorname{art}_E$ is trivial on $E^\times$ and, for $v$ outside a finite set $S$, on the uniformizers of $E_v$ placed at $v$, which generate $E_v^\times$. By weak approximation, $E^\times$ and the groups $E_v^\times$ with $v\notin S$ generate a dense subgroup of $\mathbf A_E^\times$, so $\phi\circ\operatorname{art}_E=1$, and $\phi=1$ because $\operatorname{art}_E$ is surjective.
\end{proof}
% (1) uses HilbertFromArtin.art_sign_eq and exists_conj; (3) is FrobDense.eq_one_of_frob, through eq_one_of_away and eq_one_of_unif.

\begin{lemma}\label{sha-automorphy-twist-units}
Let $K$ be a number field, and for $N\ge1$ let $V_N$ be the group of ideles $u$ of $K$ with $u_w\in1+N\mathcal O_w$ at every finite place $w$.
\begin{enumerate}
\item For every $N\ge1$ there is $e\ge1$ with $u^e\in V_N$ for every idele $u$ of $K$ that is a unit at every finite place.
\item There is $n\ge1$ such that $y^n=1$ for every $y\in K^\times$ that is a unit at every finite place and satisfies $|\phi(y)|=1$ for every embedding $\phi:K\to\mathbf C$.
\item If $N\ge3$, every $y\in K^\times$ with $y\in V_N$ and $|\phi(y)|=1$ for every embedding $\phi:K\to\mathbf C$ is $1$.
\item Let $S$ be a set of finite places of $K$, let $B$ be a Hausdorff abelian topological group, and let $\chi:\mathbf A_K^\times\to B$ be a continuous homomorphism trivial on the ideles whose finite components are $1$ and on the units of $K_v$ placed at $v$ for every finite place $v\notin S$. Then $\chi(b)=1$ for every idele $b$ that is a unit at every finite place outside $S$ and has component $1$ at every place of $S$.
\end{enumerate}
\end{lemma}
\leanref{Development/AutomorphicArithmetic/AutomorphicVarieties/AutomorphicForms/TwistedAutomorphicForms/LevPow.lean}{UpperTheory.V.Af.Tw.levPow}
\leanref{Development/AutomorphicArithmetic/AutomorphicVarieties/AutomorphicForms/TwistedAutomorphicForms/Kron.lean}{UpperTheory.V.Af.Tw.kron}
\leanref{Development/AutomorphicArithmetic/AutomorphicVarieties/AutomorphicForms/TwistedAutomorphicForms/EqOneOfMemLev.lean}{UpperTheory.V.Af.Tw.eq_one_of_mem_lev}
\leanref{Development/AutomorphicArithmetic/AutomorphicVarieties/AutomorphicForms/TwistedAutomorphicForms/AwayTriv.lean}{UpperTheory.V.Af.Tw.awayTriv}
\begin{proof}
(1) Only the finitely many places $w$ dividing $N$ impose a condition, and there $u_w^e\in1+N\mathcal O_w$ for $e$ divisible by the order of $(\mathcal O_w/N\mathcal O_w)^\times$.
(2) Such a $y$ is a root of unity by Kronecker's theorem; take $n$ the number of roots of unity in $K$.
(3) $(y-1)/N$ is integral at every finite place, so it is an algebraic integer, and $|\phi((y-1)/N)|\le2/N<1$ for every $\phi$. An algebraic integer all of whose conjugates have absolute value at most $1$ is $0$ or a root of unity, and a root of unity has absolute value $1$; so $y=1$.
(4) Write $b=b_\infty b_f$, with $b_\infty$ having the components of $b$ at the infinite places and $1$ elsewhere. Then $\chi(b_\infty)=1$, and $b_f$ is a limit of finite products of its components placed at their places, each of which $\chi$ kills; as $\chi$ is continuous, $\chi(b_f)=1$.
\end{proof}

\begin{lemma}\label{sha-automorphy-twist-places}
Let $F$ be a number field, let $L\supset F$ be an extension of degree $2$, let $\Omega$ be an algebraic closure of $L$, and let $G_L=\operatorname{Gal}(\Omega/L)\subset G_F=\operatorname{Gal}(\Omega/F)$. For a finite place $v$ let $q_v$ be the order of its residue field, and for a place $A$ of $\Omega$ above $v$ call an element of the decomposition group $D_A$ that acts on the residue field of $A$ by $x\mapsto x^{q_v}$ a Frobenius element at $A$ for $v$.
\begin{enumerate}
\item $G_L$ is an open subgroup of $G_F$ of index $2$.
\item If $w\ne w'$ are finite places of $L$ over the same place $v$ of $F$, the map $F_v\to L_w$ is an isomorphism of topological fields.
\item All but finitely many finite places of $F$ are unramified in $L$. For a finite place $v$ of $F$ and $x\in F_v^\times$, the image in $\mathbf A_L^\times$ of $x$ placed at $v$ is the product, over the places $w$ of $L$ over $v$, of elements $y_w\in L_w^\times$ placed at $w$; if $v$ is unramified in $L$ and $x$ is a uniformizer, each $y_w$ is a uniformizer.
\item Let $v$ be a finite place of $F$ unramified in $L$ and $A$ a place of $\Omega$ above $v$. If $s\in G_L$ is a Frobenius element at $A$ for $v$, then $v$ has two places $w\ne w'$ in $L$, $A$ lies over $w$, and $q_w=q_{w'}=q_v$. If $\sigma\in G_F\setminus G_L$ is a Frobenius element at $A$ for $v$, then $v$ has one place $w$ in $L$, $A$ lies over $w$, and $\sigma^2$ is a Frobenius element at $A$ for $w$.
\item Let $g\in G_F\setminus G_L$, let $w\ne w'$ be places of $L$ over the same place of $F$, and let $A$ be a place of $\Omega$ above $w$. Then there is a place $A'$ of $\Omega$ above $w'$ such that for all $s,s'\in G_L$ with $s'=gsg^{-1}$ and every $q$: if $s\in I_A$ then $s'\in I_{A'}$, and if $s\in D_A$ acts on the residue field of $A$ by $x\mapsto x^q$, then $s'\in D_{A'}$ acts on that of $A'$ by $x\mapsto x^q$.
\end{enumerate}
\end{lemma}
\leanref{Development/AutomorphicArithmetic/AutomorphicVarieties/AutomorphicForms/TwistedAutomorphicForms/ResOpen.lean}{UpperTheory.V.Af.Tw.resOpen}
\leanref{Development/AutomorphicArithmetic/AutomorphicVarieties/AutomorphicForms/TwistedAutomorphicForms/SplitIso.lean}{UpperTheory.V.Af.Tw.splitIso}
\leanref{Development/AutomorphicArithmetic/AutomorphicVarieties/AutomorphicForms/TwistedAutomorphicForms/UnrAE.lean}{UpperTheory.V.Af.Tw.unrAE}
\leanref{Development/AutomorphicArithmetic/AutomorphicVarieties/AutomorphicForms/TwistedAutomorphicForms/LocProd.lean}{UpperTheory.V.Af.Tw.locProd}
\leanref{Development/AutomorphicArithmetic/AutomorphicVarieties/AutomorphicForms/TwistedAutomorphicForms/FrobSplit.lean}{UpperTheory.V.Af.Tw.frobSplit}
\leanref{Development/AutomorphicArithmetic/AutomorphicVarieties/AutomorphicForms/TwistedAutomorphicForms/FrobInert.lean}{UpperTheory.V.Af.Tw.frobInert}
\leanref{Development/AutomorphicArithmetic/AutomorphicVarieties/AutomorphicForms/TwistedAutomorphicForms/ConjPlace.lean}{UpperTheory.V.Af.Tw.conjPlace}
\begin{proof}
These are standard facts on decomposition and inertia groups in a quadratic extension \cite[Chapter~I, \S\S8--9]{Neu99}. In (2), $v$ has two places in $L$, so $e_wf_w=1$. In (5), $A'=g(A)$.
\end{proof}

\begin{lemma}\label{sha-automorphy-twist-transfer}
Let $\Gamma$ be a topological group, let $H\subset\Gamma$ be an open subgroup of index $2$, let $B$ be an abelian topological group, and let $\kappa:H\to B$ be a continuous homomorphism. Then there is a continuous homomorphism $\nu:\Gamma\to B$ such that for every $g\in\Gamma\setminus H$, $\kappa(s)\kappa(gsg^{-1})=\nu(s)$ for $s\in H$ and $\kappa(g^2)=\nu(g)$.
\end{lemma}
\leanref{Development/AutomorphicArithmetic/AutomorphicVarieties/AutomorphicForms/TwistedAutomorphicForms/VerGroup.lean}{UpperTheory.V.Af.Tw.verGroup}
\begin{proof}
$H$ is normal. Fix $g_0\notin H$. For $g\notin H$ and $x\in H$, $g_0^{-1}g\in H$, so $\kappa(gxg^{-1})=\kappa(g_0xg_0^{-1})$, as $B$ is abelian.
Let $\nu(\gamma)=\kappa(\gamma)\kappa(g_0\gamma g_0^{-1})$ for $\gamma\in H$ and $\nu(\gamma)=\kappa(\gamma^2)$ for $\gamma\notin H$; this is $\kappa$ composed with the transfer from $\Gamma$ to $H$.
It is multiplicative, by the four cases for $x,y$ in or out of $H$ and the identities $xy\,xy=x(yxy^{-1})y^2$, $xy\,xy=(xyx^{-1})x^2y$ and $xy\,yx=x\,y^2\,x^{-1}\,x^2$, with the remark above. It is continuous at $1$, because $H$ is open and $s\mapsto\kappa(s)\kappa(g_0sg_0^{-1})$ is continuous on $H$. The two identities hold by the remark above.
\end{proof}

\begin{lemma}\label{sha-automorphy-twist-character}
Let $F$ be a totally real number field, let $L\supset F$ be a CM field with $[L:F]=2$, let $\Omega$ be an algebraic closure of $L$, and let $G_L=\operatorname{Gal}(\Omega/L)\subset G_F=\operatorname{Gal}(\Omega/F)$. Let $\ell$ be a prime and let $M\ge3$ be an integer.
Let $\Phi$ be a set of finite places of $L$ such that for every finite place $w$ of $L$ dividing $M$ there is a place $w'\ne w$ of $L$ over the same place of $F$ with $w\in\Phi$ if and only if $w'\notin\Phi$.
Let $\nu:G_F\to\bar{\mathbf Q}_\ell^\times$ be a continuous homomorphism, trivial on the complex conjugations and on the inertia groups at the places of $\Omega$ over the finite places of $F$ not dividing $M$.
Then there are a continuous homomorphism $\kappa:G_L\to\bar{\mathbf Q}_\ell^\times$ and an integer $k\ge1$ such that:
\begin{enumerate}
\item for every $g\in G_F\setminus G_L$, $\kappa(s)\kappa(gsg^{-1})=\nu(s)$ for $s\in G_L$, and $\kappa(g^2)=\nu(g)$;
\item $\kappa$ is trivial on the inertia groups at the places of $\Omega$ over the finite places of $L$ not dividing $M$;
\item $\kappa^k$ is trivial on the inertia groups at the places of $\Omega$ over the finite places of $L$ not in $\Phi$.
\end{enumerate}
\end{lemma}
\leanref{Development/AutomorphicArithmetic/AutomorphicVarieties/ModularityArgument/SymmetricPowerTransfer/GeneralizedEigenspaces/PlacesOverFinite/LinkedEndEOf.lean}{UpperTheory.V.Z.Sym.Pl.LemTWM}
\leanref{Development/AutomorphicArithmetic/AutomorphicVarieties/ModularityArgument/SymmetricPowerTransfer/GeneralizedEigenspaces/PlacesOverFinite/LemTWM.lean}{UpperTheory.V.Z.Sym.Pl.lemTWM}
\leanref{Development/AutomorphicArithmetic/AutomorphicVarieties/AutomorphicForms/TwistedAutomorphicForms/PhiPart.lean}{UpperTheory.V.Af.Tw.phiPart}
\begin{proof}
Let $\operatorname{art}_F$ and $\operatorname{art}_L$ be the Artin maps of Lemma~\ref{sha-automorphy-twist-artin}, and use its terms ``placed at'' and ``Frobenius element''. Let $i:\mathbf A_F^\times\to\mathbf A_L^\times$ be the inclusion and $c$ the nontrivial automorphism of $L$ over $F$, acting on $L^\times$ and on $\mathbf A_L^\times$. Let $U$ be the group of ideles of $L$ that are units at every finite place, and $V_M\subset U$ the group of Lemma~\ref{sha-automorphy-twist-units} for $L$ and $M$.
\begin{enumerate}
\item \emph{The character of $F$.} Let $N_F=\nu\circ\operatorname{art}_F$, with $\nu$ viewed on $G_F^{\mathrm{ab}}$. It is continuous and trivial on $F^\times$; it is trivial on the ideles whose finite components are $1$, by Lemma~\ref{sha-automorphy-twist-artin}(1); and for a finite place $v\nmid M$ it is trivial on the units of $F_v$ placed at $v$, which $\operatorname{art}_F$ maps into the image of an inertia group at $v$.
\item \emph{A projection.} Let $v$ be a finite place of $F$ dividing $M$, and $w$ a place of $L$ over $v$. The hypothesis on $\Phi$ gives a place $w'\ne w$ over $v$ with exactly one of $w,w'$ in $\Phi$; let $w_v$ be that one. By Lemma~\ref{sha-automorphy-twist-places}(2), $F_v\to L_{w_v}$ is an isomorphism; let $r_v$ be its inverse. Let $p:\mathbf A_L^\times\to\mathbf A_F^\times$ be the product over the places $v\mid M$ of $x\mapsto r_v(x_{w_v})$ placed at $v$. Then $p$ is continuous; $p(i(a))$ has the component $a_v$ at each $v\mid M$ and $1$ elsewhere; $p(x)$ has component $1$ at every $v\nmid M$; and $p(x)=1$ if $x_w=1$ at every $w\in\Phi$ dividing $M$.
\item \emph{Gluing.} Apply Lemma~\ref{sha-automorphy-twist-glue} with $I=\mathbf A_L^\times$, $J=\mathbf A_F^\times$, $K=L^\times$, $B=\bar{\mathbf Q}_\ell^\times$ (which is divisible), the maps $i$, the diagonal map $d:L^\times\to\mathbf A_L^\times$, $c$ and $\sigma=c|_{L^\times}$, the groups $U$ and $V=V_M$, the characters $N_F$ and $g=N_F\circ p$, the integer $n$ of Lemma~\ref{sha-automorphy-twist-units}(2) for $L$, and in place of $M$ the integer $e$ of Lemma~\ref{sha-automorphy-twist-units}(1) for $L$ and $M$. Its hypotheses hold:
(1) is clear;
for (2), if $i(a)\in U$, then $a$ is a unit at every finite place, and $b=a/p(i(a))$ has component $a_v$ at $v\nmid M$ and $1$ at $v\mid M$, so $N_F(b)=1$ by step (1) and Lemma~\ref{sha-automorphy-twist-units}(4) for the places dividing $M$, which is $g(i(a))=N_F(a)$;
for (3), an $x\in L^\times$ with $c(x)=x$ lies in $F^\times$, and $N_F(x)=1$;
for (4), $y=c(x)/x$ satisfies $|\phi(y)|=1$ for every embedding $\phi:L\to\mathbf C$, since $L$ is CM, so Lemma~\ref{sha-automorphy-twist-units}(2) and (3), with $M\ge3$, give the two conditions;
and (5) is Lemma~\ref{sha-automorphy-twist-units}(1).
This gives $\psi,\theta:\mathbf A_L^\times\to\bar{\mathbf Q}_\ell^\times$ and $k\ge1$ with $\psi\circ i=N_F$, $\psi$ trivial on $L^\times$, $\psi(u)=N_F(p(u))\theta(u/c(u))$ for $u\in U$, $\theta=1$ on $V_M$ and $\theta^k=1$ on $U$.
\item \emph{Properties of $\psi$.} For $u\in V_M$, also $c(u)\in V_M$, so $\psi(u)=N_F(p(u))$. Since $V_M$ is a neighbourhood of $1$, $\psi$ is continuous at $1$, hence continuous. An idele $z$ with finite components $1$ lies in $V_M$ and has $p(z)=1$, so $\psi(z)=1$. Let $w$ be a finite place of $L$ and $x$ a unit of $L_w$ placed at $w$. If $w\nmid M$, then $x\in V_M$ and $p(x)=1$, so $\psi(x)=1$. If $w\notin\Phi$, then $p(x)=1$, so $\psi(x)=\theta(x/c(x))$ and $\psi(x)^k=1$.
\item \emph{The character $\kappa$.} Lemma~\ref{sha-automorphy-twist-artin}(2) gives a continuous $\kappa:G_L\to\bar{\mathbf Q}_\ell^\times$ with $\kappa(\tau)=\psi(x)$ whenever $\operatorname{art}_L(x)$ is the image of $\tau$. Every element of an inertia group $I_A$ at a place $A$ over a finite place $w$ of $L$ has image $\operatorname{art}_L(x)$ for a unit $x$ of $L_w$ placed at $w$, so step (4) gives (2) and (3).
Moreover, if $\kappa$ is trivial on $I_A$, $s$ is a Frobenius element at $A$ for $w$ and $y$ is a uniformizer of $L_w$ placed at $w$, then $\psi(y)=\kappa(s)^{-1}$: $\operatorname{art}_L(y)$ differs from the image of $s^{-1}$ by the image of an element of $I_A$.
\item \emph{The transfer.} By Lemma~\ref{sha-automorphy-twist-places}(1) and Lemma~\ref{sha-automorphy-twist-transfer}, there is a continuous $\nu':G_F\to\bar{\mathbf Q}_\ell^\times$ such that (1) holds with $\nu'$ in place of $\nu$. By Lemma~\ref{sha-automorphy-twist-artin}(3) it suffices to show $\nu'(\sigma)=\nu(\sigma)$ for every Frobenius element $\sigma$ at a place $A$ above a finite place $v$ of $F$ that does not divide $M$ and is unramified in $L$; by Lemma~\ref{sha-automorphy-twist-places}(3) this excludes finitely many $v$, and by (2) $\kappa$ is unramified above such $v$.
Let $x$ be a uniformizer of $F_v$ placed at $v$ with $\operatorname{art}_F(x)$ the image of $\sigma^{-1}$. Then $\psi(i(x))=N_F(x)=\nu(\sigma)^{-1}$, and by Lemma~\ref{sha-automorphy-twist-places}(3), $i(x)=\prod_{w\mid v}y_w$ with uniformizers $y_w$ of $L_w$ placed at $w$.
If $\sigma\in G_L$, Lemma~\ref{sha-automorphy-twist-places}(4) gives $w\ne w'$ over $v$ with $A$ over $w$ and $q_w=q_{w'}=q_v$, so $\sigma$ is a Frobenius element at $A$ for $w$. Take $g\in G_F\setminus G_L$ and $s'=g\sigma g^{-1}\in G_L$; by Lemma~\ref{sha-automorphy-twist-places}(5), $s'$ is a Frobenius element at a place $A'$ over $w'$. By step (5), $\nu(\sigma)^{-1}=\psi(y_w)\psi(y_{w'})=(\kappa(\sigma)\kappa(s'))^{-1}=\nu'(\sigma)^{-1}$.
If $\sigma\notin G_L$, Lemma~\ref{sha-automorphy-twist-places}(4) gives the one place $w$ over $v$, with $A$ over $w$ and $\sigma^2$ a Frobenius element at $A$ for $w$, and $\nu(\sigma)^{-1}=\psi(y_w)=\kappa(\sigma^2)^{-1}=\nu'(\sigma)^{-1}$.
So $\nu'=\nu$, which gives (1).
\end{enumerate}
\end{proof}
% Steps (1) to (4) are the private ideleExt_of (with phiPart, levPow, kron, awayTriv and eq_one_of_mem_lev); step (5) is descend and h7 (GalChar.exists_unit, and frob_val for the Frobenius elements); step (6) is the private verArt_of (resOpen, verGroup, conjPlace, unrAE, locProd, frobSplit, frobInert) with GalChar.eq_of_frob and GalChar.away_ae. The uniformizer x is ha.frob.

\begin{lemma}\label{sha-automorphy-algebraic-twist}
Let $L\subset\bar{\mathbf Q}$ be a CM field with maximal totally real subfield $L^+$, let $\ell$ be an odd prime, let $\iota:\bar{\mathbf Q}_\ell\cong\mathbf C$, and let $c\in\operatorname{Gal}(\bar{\mathbf Q}/\mathbf Q)$ be a complex conjugation.
Let $M\ge1$ with $\ell\mid M$, and suppose that $L$ is unramified over $L^+$ at every finite place and that every place of $L^+$ dividing $M$ splits in $L$.
Let $\chi:\operatorname{Gal}(\bar{\mathbf Q}/\mathbf Q)\to\bar{\mathbf Q}_\ell^\times$ be a character with open kernel and let $a\in\mathbf Z$ with $\chi(c)=(-1)^a$, and suppose that $\chi$ is trivial on the inertia groups of $\operatorname{Gal}(\bar{\mathbf Q}/L)$ at the places not dividing $M$.
Then there is a character $\kappa:\operatorname{Gal}(\bar{\mathbf Q}/L)\to\bar{\mathbf Q}_\ell^\times$, the $\ell$-adic character through $\iota$ of an algebraic Hecke character of $L$, such that:
\begin{enumerate}
\item $\kappa(s)\kappa(csc^{-1})=\chi_{\mathrm{cyc}}(s)^a\chi(s)$ for every $s\in\operatorname{Gal}(\bar{\mathbf Q}/L)$;
\item $\kappa$ is trivial on the inertia groups at the places not dividing $M$;
\item for every place $A$ of $\bar{\mathbf Q}$ above $\ell$ there are $b\in\mathbf Z$ and $e\ge1$ with $\kappa^e=\chi_{\mathrm{cyc}}^{be}$ on the inertia group of $\operatorname{Gal}(\bar{\mathbf Q}/L)$ at $A$.
\end{enumerate}
\end{lemma}
\leanref{Development/AutomorphicArithmetic/AutomorphicVarieties/ModularityArgument/SymmetricPowerTransfer/UnrSplit.lean}{UpperTheory.V.Z.Sym.TwEx}
\leanref{Development/AutomorphicArithmetic/AutomorphicVarieties/ModularityArgument/SymmetricPowerTransfer/GeneralizedEigenspaces/PlacesOverFinite/ExistsType.lean}{UpperTheory.V.Z.Sym.Pl.twEx}
\begin{proof}
Let $\nu=\chi_{\mathrm{cyc}}^a\chi$ on $\operatorname{Gal}(\bar{\mathbf Q}/L^+)$. It is continuous. It is trivial on every complex conjugation $c'$ in $\operatorname{Gal}(\bar{\mathbf Q}/L^+)$, since $\chi(c')=\chi(c)=(-1)^a$ and $\chi_{\mathrm{cyc}}(c')=-1$. It is trivial on the inertia groups at the places of $L^+$ not dividing $M$: such a place is prime to $\ell$, and since $L/L^+$ is unramified, its inertia groups lie in $\operatorname{Gal}(\bar{\mathbf Q}/L)$.
Since every place of $L^+$ dividing $M$ splits in $L$, there is a set $\Phi$ of places of $L$ dividing $M$ that contains exactly one of the two places above each place of $L^+$ dividing $M$.
Since $\ell$ is odd and $\ell\mid M$, $M\ge3$. Lemma~\ref{sha-automorphy-twist-character}, applied to $F=L^+$, $\nu$, $M$ and $\Phi$, gives a continuous character $\kappa$ of $\operatorname{Gal}(\bar{\mathbf Q}/L)$ and $k\ge1$ such that $\kappa(s)\kappa(gsg^{-1})=\nu(s)$ for $s\in\operatorname{Gal}(\bar{\mathbf Q}/L)$ and $g\in\operatorname{Gal}(\bar{\mathbf Q}/L^+)$ outside $\operatorname{Gal}(\bar{\mathbf Q}/L)$, $\kappa$ is trivial on the inertia groups at the places not dividing $M$, and $\kappa^k$ is trivial on the inertia groups at the places not in $\Phi$.
Taking $g=c$ gives (1), and (2) holds.
Let $n\ge1$ with $\chi^n=1$, and let $A$ be a place above $\ell$, over the place $w$ of $L$, which divides $M$. If $w\notin\Phi$, then $\kappa^{kn}=1$ on the inertia group at $A$. If $w\in\Phi$, let $w'$ be the other place of $L$ over the same place of $L^+$; conjugation by $c$ carries the inertia group at $A$ to an inertia group above $w'$, so for $s$ in the inertia group at $A$, $\kappa(s)=\chi_{\mathrm{cyc}}(s)^a\chi(s)\kappa(csc^{-1})^{-1}$ with $\kappa(csc^{-1})^k=1$, and $\kappa(s)^{kn}=\chi_{\mathrm{cyc}}(s)^{akn}$. This gives (3), with $e=kn$ and $b=0$ or $b=a$.
Finally, a continuous character of $\operatorname{Gal}(\bar{\mathbf Q}/L)$ that agrees on an inertia group above each place over $\ell$ with a power of $\chi_{\mathrm{cyc}}$, up to a character of finite order, is the $\ell$-adic character of an algebraic Hecke character \cite{Ser68}; so is $\kappa$.
\end{proof}
% nu and its properties are hνc, hνconj and hνunr in Pl.twEx; Phi is exists_type (Af.Tw.IsType); the character kappa is Pl.LemTWM (VerEq), proved as Pl.lemTWM; the computation at the places of Phi is hloc, with Tw.conjPlace; algebraicity is the first part of Af.AlgOfCyc (algOfCyc_s).

\begin{definition}\label{sha-automorphy-adjoint-input}
Let $L$ be a CM field with maximal totally real subfield $L^+$, let $\ell$ be a prime, let $\iota:\bar{\mathbf Q}_\ell\cong\mathbf C$, and let $M\ge1$ and $m\ge0$ be integers. Write $G_L=\operatorname{Gal}(\bar L/L)$.
\begin{enumerate}
\item A continuous $r:G_L\to GL_2(\bar{\mathbf Q}_\ell)$ is an \emph{adjoint input of level $M$ and weight $m$} if: (a) $r$ is attached through $\iota$ to a regular algebraic cuspidal automorphic representation $\Pi$ of $GL_2(\mathbf A_L)$ with $\Pi^\vee\cong\Pi^c$, and $r$ is windowed in the sense of Theorem~\ref{sha-residually-reducible-automorphy}; (b) $\operatorname{ad}^0r$ is Tate-free (Definition~\ref{sha-automorphy-tate-free}); (c) for every complex conjugation $c'\in\operatorname{Gal}(\bar L/L^+)$ there is an invertible $P$ with ${}^tr(s)Pr(c'sc'^{-1})=\chi_{\mathrm{cyc}}(s)^{-1}P$ for $s\in G_L$; (d) $r$ is irreducible on every open subgroup; (e) $m\ge2$; (f) $r$ is trivial on the inertia groups at the places not dividing $M$; (g) there are continuous characters $\tau_0,\tau_1$ of $G_L$ such that $r$ is diagonalizable on every inertia group at a place not above $\ell$, with diagonal characters $\tau_0,\tau_1$; (h) for every place $A$ of $\bar L$ above $\ell$ there are $b\in\mathbf Z$, $k\ge1$ and characters $\kappa_0,\kappa_1$ of $D_A$ such that $r|_{D_A}$ is conjugate to an upper triangular representation with diagonal $\kappa_0,\kappa_1$, and $\kappa_0^k=\chi_{\mathrm{cyc}}^{(b+m)k}$ and $\kappa_1^k=\chi_{\mathrm{cyc}}^{bk}$ on $I_A$.
\item For a finite place $w$ of $L$ and smooth characters $\psi_1,\dots,\psi_n$ of $L_w^\times$, $I(\psi_1,\dots,\psi_n)$ is the normalized parabolic induction to $GL_n(L_w)$, from the upper triangular Borel subgroup, of $\operatorname{diag}(t_1,\dots,t_n)\mapsto\prod_i\psi_i(t_i)$; it is $\mathrm{PS}(\omega)$ of Definition~\ref{sha-hecke-local-ordinary} for $\omega_i=\psi_i|\cdot|^{(1-n)/2}$.
\item An element $s\in G_L$ is \emph{exceptional} for $r$ if, for the eigenvalues $\alpha,\beta$ of $r(s)$, $(\alpha/\beta)^2+(\beta/\alpha)^2=\chi_{\mathrm{cyc}}(s)^j+\chi_{\mathrm{cyc}}(s)^{-j}$ for some integer $j$ with $|j|\le2$.
\end{enumerate}
\end{definition}
\leanref{Development/AutomorphicArithmetic/AutomorphicVarieties/ModularityArgument/SymmetricPowerTransfer/WeilDeligneUnramified/UnitaryTwoForm.lean}{UpperTheory.V.Z.Sym.W.AdInW}
\leanref{Development/AutomorphicArithmetic/AutomorphicVarieties/ModularityArgument/SymmetricPowerTransfer/GeneralizedEigenspaces/WeightKernels/GoodSet.lean}{UpperTheory.V.Z.Sym.Ga.Bad}

\begin{lemma}\label{sha-automorphy-unitary-comparison}
Assume that Galois representations exist, in the sense of Definition~\ref{gal-galois-representations}, relative to the maps satisfying the laws of Definition~\ref{gal-llc-laws}.
Let $E$ be a CM field with maximal totally real subfield $E^+$ such that $E/E^+$ is unramified at every finite place, let $\ell$ be a prime, let $\iota:\bar{\mathbf Q}_\ell\cong\mathbf C$, and let $n\ge1$. Let $\Pi$ be a cuspidal automorphic representation of $GL_n(\mathbf A_E)$ and $r:\operatorname{Gal}(\bar E/E)\to GL_n(\bar{\mathbf Q}_\ell)$ a continuous representation attached to $\Pi$ through $\iota$. Let $\sigma$ be an automorphic representation of $U_n$ (Definition~\ref{gal-definite-unitary}) such that at all but finitely many finite places $w$ of $E$ split over $E^+$, $\sigma$ has a spherical local component whose Satake parameters are Satake parameters of $\Pi$ at $w$. Then:
\begin{enumerate}
\item at every finite place $w$ split over $E^+$, every local component of $\Pi$ at $w$ is a local component of $\sigma$ at $w$;
\item at every embedding $u:E\to\mathbf C$, $\sigma$ has a weight $\lambda$ with $\{\lambda_k+\frac{n-1}2-k:0\le k<n\}=Z_u(\Pi)$ (Definition~\ref{gal-polarized}).
\end{enumerate}
\end{lemma}
\leanref{Development/AutomorphicArithmetic/AutomorphicVarieties/ModularityArgument/CuspidalDescent/SphericalTwists/RowProps.lean}{UpperTheory.V.Z.Desc.W.Same}
\leanref{Development/AutomorphicArithmetic/AutomorphicVarieties/ModularityArgument/ThresholdArguments/DenseThresholdCases/TSameOf.lean}{UpperTheory.V.Z.Desc.W.same_of_statement}
\leanref{Development/AutomorphicArithmetic/AutomorphicVarieties/ModularityArgument/ThresholdArguments/DenseThresholdCases/JPTRow.lean}{UpperTheory.V.Z.Desc.W.same_of_prints_of_rT}
\begin{proof}
\begin{enumerate}
\item \emph{The Speh decomposition.} Theorem~\ref{gal-base-change-definite}, applied to $\sigma$, together with Bernstein's theorem that a product of irreducible unitary representations of general linear groups is irreducible \cite{Ber84}, gives $c\ge1$, integers $a_i,d_i\ge1$ with $\sum_id_ia_i=n$, cuspidal automorphic representations $P_i$ of $GL_{a_i}(\mathbf A_E)$ with $P_i^c\cong P_i^\vee$, and at each finite place $w$ local components $\rho_{i,w}$ of $P_i$ and the Speh representations $\sigma_{i,w}$ of $GL_{d_ia_i}(E_w)$ attached to $\rho_{i,w}$ and $d_i$ \cite{MW89}, such that: at every split place $w$ the product $\sigma_{1,w}\times\dots\times\sigma_{c,w}$ is a local component of $\sigma$; and if $\sigma$ has weight $\lambda$ at $u$, then $\{\lambda_k+\frac{n-1}2-k\}$ is the sum over $i$ and $0\le k<d_i$ of the multisets $Z_u(P_i)+\frac{d_i-1}2-k$.
\item \emph{The Galois side of the parts.} Since $\sigma$ has an algebraic weight at each $u$, by (1) the elements of each $Z_u(P_i)$ are distinct, and lie in a coset of $\mathbf Z$ fixed by $n$, $d_i$ and $k$. So for each $i$ and $0\le k<d_i$ there is a real number $s_{i,k}$ such that $P_i|\det|^{s_{i,k}}$ is regular algebraic, and it is polarized since $P_i^c\cong P_i^\vee$. Let $r_{i,k}$ be its Galois representation (Definition~\ref{gal-galois-representations}), and $f(g)=\prod_{i,k}\det\big(X-r_{i,k}(g^{-1})\big)$. The coefficients of $f$ are continuous, and by local-global compatibility at the places $w\nmid\ell$ where every $P_i$ is unramified, $\iota f$ at a Frobenius element at $w$ is the polynomial whose roots are $q_w^{(n-1)/2}$ times the Satake parameters of $\sigma_{1,w}\times\dots\times\sigma_{c,w}$, namely the union over $i,k$ of $q_w^{(d_i-1)/2-k}$ times those of $P_i$.
\item \emph{Comparison with $r$.} Let $w$ be one of the almost all split places at which $\sigma$ has a spherical local component $\tau$ with Satake parameters $B$ that are Satake parameters of $\Pi$, with $w\nmid\ell$ and every $P_i$ unramified. Since the local component of $\sigma$ at a split place is unique up to isomorphism \cite{Fla79}, $\tau\cong\sigma_{1,w}\times\dots\times\sigma_{c,w}$, so $B$ is the Satake parameter of the Speh product computed in (2) \cite{Zel80}. As $r$ is attached to $\Pi$, the characteristic polynomial of $r(g^{-1})$ at a Frobenius element $g$ at $w$ is then $f(g)$. The Frobenius elements at these places, which include almost all places of degree one, are dense in $\operatorname{Gal}(\bar E/E)$ by the Chebotarev density theorem, since a place of degree one is split over $E^+$; by continuity, $\det(X-r(g^{-1}))=f(g)$ for every $g$.
\item \emph{One part.} So at almost all places $w$ of $E$, $\Pi$ has as Satake parameters the union over $i,k$ of $q_w^{(d_i-1)/2-k}$ times those of $P_i$, which are the Satake parameters of the twists $P_i|\det|^{k-(d_i-1)/2}$. By the theorem of Jacquet and Shalika on isobaric representations \cite{JS81b}, the cuspidal representation $\Pi$ and the family of these $\sum_id_i$ twists then have the same members; so $\sum_id_i=1$, that is $c=1$ and $d_1=1$, and $a_1=n$.
\item \emph{Strong multiplicity one.} Now $P_1$ and $\Pi$ have the same Satake parameters at almost all places, so $\Pi=P_1$ \cite{JS81b}.
\item \emph{Conclusion.} (1) At a split place $w$, $\sigma_{1,w}=\rho_{1,w}$ is a local component of $\Pi$ and of $\sigma$; every local component of $\Pi$ at $w$ is isomorphic to it \cite{Fla79}, so it is a local component of $\sigma$. (2) The automorphic representation $\sigma$ has a weight $\lambda$ at $u$, and by (1) of this proof, with $c=d_1=1$, $\{\lambda_k+\frac{n-1}2-k\}=Z_u(P_1)=Z_u(\Pi)$.
\end{enumerate}
\end{proof}
% same_of_prints_of_rT = same_of_statement (TSameOf.lean) with: (1) bcSpeh_alt (bridge of litN.aut.bc and litN.loc.bern), whose Parts give hP.weight, hP.inf, hP.csd and hP.loc, and hsplit; (2) speh_parameter_strings for the integrality, sumAEOf with galAEOf_alt (RowT for the Galois representations, rec_sat for the local-global compatibility) and satTw; (3) building_component_uniqueness, splitSat_alt and eq_of_split (Published.chebotarevDegreeOne', isSplitAt_of_unramified_of_inertiaDeg_eq_one); (4) onePart from Published.JacquetShalikaRigidity; (5) strongMultOneC; (6) Z.Desc.irrComp and spehInf_one', with exists_algebraicWeight_of_automorphic_definiteUnitary.

\begin{lemma}\label{sha-automorphy-unitary-two}
Assume that Galois representations exist, in the sense of Definition~\ref{gal-galois-representations}, relative to the maps satisfying the laws of Definition~\ref{gal-llc-laws}.
Let $L$ be a CM field such that $L/L^+$ is unramified at every finite place, let $\ell$ be an odd prime, let $\iota:\bar{\mathbf Q}_\ell\cong\mathbf C$, let $M\ge1$ with $\ell\mid M$ be such that every place of $L^+$ dividing $M$ splits in $L$, let $m\ge0$, and let $r$ be an adjoint input of level $M$ and weight $m$ (Definition~\ref{sha-automorphy-adjoint-input}).
Then there are an automorphic representation $\sigma$ of $U_2$ (Definition~\ref{gal-definite-unitary}, for $E=L$) and multisets $\mathrm{HT}_\tau$ of integers, for the embeddings $\tau:L\to\bar{\mathbf Q}_\ell$, such that:
\begin{enumerate}
\item at every finite place $w$ of $L$ split over $L^+$, $\sigma$ has a local component; every local component $\sigma_w$ of $\sigma$ at $w$ is plain (Definition~\ref{gal-plain}), $r$ is compatible with $\sigma_w$ in the sense of Lemma~\ref{sha-hecke-ideal-attached}(1), and if $w\mid\ell$, $r$ is de Rham at the places above $w$ with Hodge-Tate multisets $\mathrm{HT}_\tau$ at the embeddings $\tau$ that extend continuously to $L_w$;
\item $\sigma$ is unramified at every finite place not dividing $M$;
\item at every infinite place $v$, $\sigma$ has a weight $\lambda$ at an embedding inducing $v$ with $\lambda_0-\lambda_1=m-1$.
\end{enumerate}
\end{lemma}
\leanref{Development/AutomorphicArithmetic/AutomorphicVarieties/ModularityArgument/ThresholdArguments/DenseThresholdCases/WStmtxx.lean}{UpperTheory.V.Z.Sym.W.AboveWeightAlt}
\leanref{Development/AutomorphicArithmetic/AutomorphicVarieties/ModularityArgument/ThresholdArguments/DenseThresholdCases/AdjointLiftTransferToUnitaryGroupTwo.lean}{UpperTheory.V.Z.Sym.W.aboveWeight}
\begin{proof}
By Definition~\ref{sha-automorphy-adjoint-input}(1)(a), $r$ is windowed and attached through $\iota$ to a regular algebraic cuspidal $\Pi$ of $GL_2(\mathbf A_L)$ with $\Pi^\vee\cong\Pi^c$.

\emph{Compatibility.} Let $\chi$ be the quadratic Hecke character of $L^+$ attached to $L$. At every real place $v$, $\chi^2_v(-1)=1=(-1)^2$, and $\chi^2$ composed with the norm from $L$ is trivial; so $(\Pi,\chi^2)$ is polarized (Definition~\ref{gal-polarized}), since $\Pi^\vee\cong\Pi^c$. Lemma~\ref{sha-hecke-ideal-attached} applies to $r$ and $\Pi$, with the multisets $\mathrm{HT}_\tau=\{\tfrac12-z:z\in Z_{\iota\circ\tau}(\Pi)\}$.

\emph{Plain local components.} Let $w$ be a place of $L$ above $\ell$ and $A$ a place above $w$, with $b$, $k$, $\kappa_0$ and $\kappa_1$ as in Definition~\ref{sha-automorphy-adjoint-input}(1)(h), and let $(c_0,c_1)=(b+m,b)$. As in the proof of Lemma~\ref{sha-automorphy-local-ratios}(3), some $s\in I_A$ has $\chi_{\mathrm{cyc}}(s)$ of infinite order. If $\kappa_i^e=(\kappa_j\chi_{\mathrm{cyc}})^e$ for some $e\ge1$, then raising to the power $k$ at $s$ gives $\chi_{\mathrm{cyc}}(s)^{(c_i-c_j-1)ke}=1$, while $c_i\ne c_j+1$ because $m\ge2$. So $r$ satisfies the second condition of ($\alpha$) of the windowed condition at $w$.
By ($\beta$) there is $t\in\mathbf R$ such that at every finite place $w$ the local components of $\Pi$ at $w$ have supercuspidal support consisting of characters $\xi$ with $q_w^{t-1/2}<|\xi(\varpi_w)|<q_w^{t+1/2}$. No two of them differ by $\nu^{\pm1}$, so every local component $\Pi_w$ is plain, and $\operatorname{rc}(\Pi_w\nu^{-1/2})$ is abelian and monodromy-free for every $\operatorname{rc}$ that satisfies the laws, as in the proof of Lemma~\ref{sha-hecke-ideal-attached}(2). By Lemma~\ref{sha-hecke-ideal-attached}, $r$ is compatible with $\Pi_w$, and if $w\mid\ell$, $r$ is de Rham at the places above $w$ with the multisets $\mathrm{HT}_\tau$. Every finite place has a local component of $\Pi$ \cite{Fla79}.

\emph{Hodge-Tate weights.} Let $\phi:L\to\bar{\mathbf Q}_\ell$ be an embedding, $w$ the place above $\ell$ to whose completion it extends, $A$ a place above $w$ and $\Pi_w$ a local component at $w$. Lemma~\ref{sha-automorphy-local-ratios}(3), with (e) and (h) of Definition~\ref{sha-automorphy-adjoint-input}(1), gives $b_\phi\in\mathbf Z$ with $\mathrm{HT}_\phi=\{b_\phi,b_\phi+m\}$.
So for an embedding $u:L\to\mathbf C$ and $\phi=\iota^{-1}\circ u$, $Z_u(\Pi)=\{\tfrac12-b_\phi,\tfrac12-b_\phi-m\}$. This is the infinitesimal character $\{\lambda_0+\tfrac12,\lambda_1-\tfrac12\}$ of the weight $\lambda=(-b_\phi,1-b_\phi-m)$, and $\lambda_0>\lambda_1$ because $m\ge2$.

\emph{Spherical local components.} Let $w$ be a finite place not dividing $M$. Then $w\nmid\ell$, since $\ell\mid M$, and $r$ is trivial on the inertia groups at $w$ by (f). Every local component $\Pi_w$ is plain, and $\mathrm{WD}_\iota(r|_{W_A})=\operatorname{rc}(\Pi_w\nu^{-1/2})$ by the compatibility; so $\Pi_w$ has nonzero vectors fixed by $GL_2(\mathcal O_w)$ by Lemma~\ref{sha-hecke-ideal-spherical}, and $\Pi$ has Satake parameters at $w$.

\emph{The descent.} Let $S$ be the set of places of $L^+$ dividing $M$. It is finite, its places split in $L$, and every place of $L$ above $\ell$ lies above $S$. Whether a place of $L$ divides $M$ depends only on the place of $L^+$ below it, so $\Pi$ has Satake parameters at every place above a place outside $S$. Since $\Pi$ is regular algebraic with $\Pi^\vee\cong\Pi^c$ and of strictly decreasing weights at infinity, the descent of conjugate self-dual cuspidal representations of $GL_2(\mathbf A_L)$ to $U_2$ \cite{Lab11} gives an automorphic representation $\sigma$ of $U_2$, unramified at the places not split over $L^+$, which at almost all split places has a spherical local component with Satake parameters of $\Pi$. Lemma~\ref{sha-automorphy-unitary-comparison}, applied to $\Pi$, $r$ and $\sigma$, which compares $\sigma$ with $\Pi$ through the Speh decomposition of its base change, shows that every local component of $\Pi$ at a split place $w$ is a local component of $\sigma$ at $w$, and that at every embedding $u:L\to\mathbf C$, $\sigma$ has a weight $\lambda$ whose infinitesimal character is $Z_u(\Pi)$.

\emph{Conclusion.} (1) Let $w$ be split over $L^+$, and $\Pi_w$ a local component of $\Pi$ at $w$; it is a local component of $\sigma$. Every local component of $\sigma$ at $w$ is isomorphic to $\Pi_w$, since the local component of the irreducible representation $\sigma$ at a split place is unique up to isomorphism \cite{Fla79}. So (1) follows from the properties of $\Pi_w$.
(2) Let $w\nmid M$. If $w$ is split, the image in $\sigma$ of a nonzero vector of $\Pi_w$ fixed by $GL_2(\mathcal O_w)$ is fixed by $c_w^{-1}(GL_2(\mathcal O_w))$, so $\sigma$ is unramified at $w$. Otherwise $\sigma$ is unramified at $w$ by the descent.
(3) Let $v$ be an infinite place, $u$ an embedding inducing it, and $\lambda$ a weight of $\sigma$ at $u$ with infinitesimal character $Z_u(\Pi)$. Then $\{\lambda_0+\tfrac12,\lambda_1-\tfrac12\}=\{\tfrac12-b_\phi,\tfrac12-b_\phi-m\}$ for $\phi=\iota^{-1}\circ u$. Since $\lambda_0\ge\lambda_1$ and $m\ge0$, $\lambda_0+\tfrac12=\tfrac12-b_\phi$ and $\lambda_1-\tfrac12=\tfrac12-b_\phi-m$, so $\lambda_0-\lambda_1=m-1$.
% aboveWeight: the polarization of Pi by quadChar^2; P.compat (Lemma sha-hecke-ideal-attached) for HT and the compatibility; guard_of with P.plain and pot_of_above (flag_exp, cyc_inf_over) for the plain local components; above_w_statement for HT = {b, b+m}; Z.Desc.RegWt; sphT_alt (Lemma sha-hecke-ideal-spherical) at the places not dividing M; the descent P.desc (Z.Desc.DescUse) and P.same (Z.Desc.W.Same) for sigma; building_component_uniqueness; Ga.wtOf for the weights. Local components exist by Af.local_comp_complex_of_printed_results; unr_of_comp gives (2) at the split places.
\end{proof}

\begin{lemma}\label{sha-automorphy-local-ratios}
Assume that Galois representations exist, in the sense of Definition~\ref{gal-galois-representations}, relative to the maps satisfying the laws of Definition~\ref{gal-llc-laws}.
Let $L$ be a number field, let $\ell$ be a prime, let $\iota:\bar{\mathbf Q}_\ell\cong\mathbf C$, and let $r:G_L\to GL_2(\bar{\mathbf Q}_\ell)$ be continuous. Let $w$ be a finite place of $L$, with residue field of $q_w$ elements, and let $\tau$ be an irreducible smooth representation of $GL_2(L_w)$ that is plain (Definition~\ref{gal-plain}) and such that $r$ is compatible with $\tau$ in the sense of Lemma~\ref{sha-hecke-ideal-attached}(1).
\begin{enumerate}
\item Suppose that $w\nmid\ell$, that $r$ is trivial on the inertia groups at $w$, and that some Frobenius element at a place above $w$ is not exceptional for $r$ (Definition~\ref{sha-automorphy-adjoint-input}). Then $\tau\cong I(\psi_0,\psi_1)$ for unramified characters $\psi_0,\psi_1$ of $L_w^\times$ with $(\psi_0/\psi_1)^2\ne|\cdot|^j$ for every integer $j$ with $|j|\le2$.
\item Suppose that $w\nmid\ell$ and that $\operatorname{ad}^0r$ is Tate-free. Then $\tau\cong I(\psi_0,\psi_1)$ for smooth characters with $(\psi_0/\psi_1)^j\ne|\cdot|^{\pm1}$ for $j=1,2$.
\item Suppose that $\ell$ is odd and $w\mid\ell$, let $m\ge2$, let $\mathrm{HT}_\phi$ be multisets, indexed by the embeddings $\phi:L\to\bar{\mathbf Q}_\ell$, such that $r$ is de Rham at the places above $w$ with Hodge-Tate multisets $\mathrm{HT}_\phi$, and let $A$ be a place above $w$ with $b$, $k$, $\kappa_0$ and $\kappa_1$ as in Definition~\ref{sha-automorphy-adjoint-input}(1)(h). Then $\tau\cong I(\psi_0,\psi_1)$ for smooth characters with $(\psi_0/\psi_1)^j\ne|\cdot|^{\pm1}$ for $j=1,2$, and there is $b'\in\mathbf Z$ with $\mathrm{HT}_\phi=\{b',b'+m\}$ for every embedding $\phi$ that extends continuously to $L_w$.
\end{enumerate}
\end{lemma}
\leanref{Development/AutomorphicArithmetic/AutomorphicVarieties/ModularityArgument/ThresholdArguments/DenseThresholdCases/TGoodW.lean}{UpperTheory.V.Z.Sym.W.good_w_statement}
\leanref{Development/AutomorphicArithmetic/AutomorphicVarieties/ModularityArgument/ThresholdArguments/DenseThresholdCases/TAwayW.lean}{UpperTheory.V.Z.Sym.W.away_w_statement}
\leanref{Development/AutomorphicArithmetic/AutomorphicVarieties/ModularityArgument/ThresholdArguments/DenseThresholdCases/TAboveW.lean}{UpperTheory.V.Z.Sym.W.above_w_statement}
\begin{proof}
Let $A$ be a place above $w$, in (1) one at which some Frobenius element is not exceptional, let $a$ be a local Artin map for $A$, and let $\operatorname{rc}$ satisfy the laws. Since $\tau$ is plain, its supercuspidal support consists of characters $\psi_0,\psi_1$ with $\psi_0\ne\psi_1|\cdot|^{\pm1}$, so $I(\psi_0,\psi_1)$ is irreducible and $\tau\cong I(\psi_0,\psi_1)$ \cite{BZ77,Zel80}; and $y=\operatorname{rc}(\tau\nu^{-1/2})$ is abelian and monodromy-free (proof of Lemma~\ref{sha-hecke-ideal-attached}(2)).
If $w\nmid\ell$, the compatibility gives $\mathrm{WD}_\iota(r|_{W_A})=y$, so by Lemma~\ref{sha-hecke-local-rec}(1), for $s\in W_A$ the eigenvalues of $\iota r(s)$ are $\psi_i(a(s))|a(s)|^{-1/2}$, $i=0,1$; in particular $\iota$ carries the ratio $\alpha/\beta$ of the eigenvalues of $r(s)$ to $(\psi_0/\psi_1)(a(s))$.

(1) An element $t\in I_A$ has $r(t)=1$, so $\psi_i(a(t))=1$; since $a(I_A)=\mathcal O_w^\times$, the $\psi_i$ are unramified. Let $s$ be a Frobenius element at $A$, of degree one, that is not exceptional, so that $\iota(\chi_{\mathrm{cyc}}(s))=q_w=|a(s)|$. If $(\psi_0/\psi_1)^2=|\cdot|^j$ with $|j|\le2$, then $(\alpha/\beta)^2=\chi_{\mathrm{cyc}}(s)^j$, so $(\alpha/\beta)^2+(\beta/\alpha)^2=\chi_{\mathrm{cyc}}(s)^j+\chi_{\mathrm{cyc}}(s)^{-j}$, and $s$ is exceptional.

(2) Since $y$ is monodromy-free, $r$ is trivial on an open subgroup of $I_A$, by Grothendieck's monodromy theorem \cite[Appendix]{ST68}, so $\operatorname{ad}^0r$ maps $I_A$ to a finite set. As $\operatorname{ad}^0r$ is Tate-free, there are $j'\ge1$ and $s\in D_A$ acting on the residue field by $x\mapsto x^{q_w^{j'}}$, an element of $W_A$ of degree $j'$ with $|a(s)|=q_w^{j'}$, such that no two eigenvalues $x,y$ of $\operatorname{ad}^0r(s)$ satisfy $x=q_w^{j'}y$. These eigenvalues are $c,1,c^{-1}$ with $c=\alpha/\beta$. If $(\psi_0/\psi_1)^j=|\cdot|^{\pm1}$ with $j\in\{1,2\}$, then $c^j=q_w^{j'}$ or $c^jq_w^{j'}=1$; so $c=q_w^{j'}\cdot1$, $1=q_w^{j'}c$, $c=q_w^{j'}c^{-1}$ or $c^{-1}=q_w^{j'}c$, a contradiction.

(3) Let $c_0=b+m$ and $c_1=b$. The character $\chi_{\mathrm{cyc}}$ has infinite order on $I_A$, since by local class field theory some $s\in I_A$ has $\chi_{\mathrm{cyc}}(s)$ a power of $1+\ell$. So $\kappa_i\ne\kappa_j\chi_{\mathrm{cyc}}$ for all $i,j$: otherwise $\chi_{\mathrm{cyc}}^{(c_i-c_j-1)k}=1$ on $I_A$, while $c_i\ne c_j+1$ because $m\ge2$. On $I_A$, $\kappa_i\chi_{\mathrm{cyc}}^{-c_i}$ has finite order, so, by the description of $\chi_{\mathrm{cyc}}$ on $I_A$ through $a$ (proof of Lemma~\ref{sha-hecke-local-characters}), $\kappa_i$ satisfies the hypothesis of Lemma~\ref{sha-hecke-ideal-above} with $h_{i,\phi}=-c_i$ for every continuous embedding $\phi$ of $L_w$, which increases with $i$. That lemma, with $n=2$ and the representation $\tau$ in place of its $\sigma$, gives smooth characters $\omega_0,\omega_1$ with $\kappa_i(s)=\iota^{-1}(\omega_i(a(s)))\prod_\phi\phi(a(s))^{c_i}$ on $W_A$ and $\tau\cong\mathrm{PS}(\omega)=I(\psi_0,\psi_1)$, $\psi_i=\omega_i|\cdot|^{1/2}$, and $\mathrm{HT}_\phi=\{-b-m,-b\}$; take $b'=-b-m$.
Let $u=\omega_0/\omega_1=\psi_0/\psi_1$. The characters $\kappa_0/\kappa_1$, $1$ and $\kappa_1/\kappa_0$ of $D_A$ correspond in the same way to $u$, $1$ and $u^{-1}$, with the exponents $h=(-m,0,m)$, which increase, and no two of them differ by $\chi_{\mathrm{cyc}}$, by the argument above with $(c_0,c_1,c_2)=(m,0,-m)$. The argument ``No two $\omega_j$ are linked'' in the proof of Lemma~\ref{sha-hecke-ideal-above} then shows that no two of $u,1,u^{-1}$ differ by $|\cdot|^{\pm1}$. Hence $u\ne|\cdot|^{\pm1}$, and $u^2\ne|\cdot|^{\pm1}$, since $u^2=|\cdot|^{\pm1}$ would give $u=u^{-1}|\cdot|^{\pm1}$.
\end{proof}
% The plain tau is a principal series by principal_series_of_statement_w (PsOf); the eigenvalues are eigW_alt (EigWAlt). (1) is good_w_statement (Ga.ufFacts for (alpha/beta)^2+(beta/alpha)^2, Ga.cyc_char_eq_abs_pow_and_norm_pow_deg_away_l). (2) is away_w_statement (the private adjoint_inertia_image_finite, ad0Facts for the eigenvalues c,1,c^-1, link_false, SatChar.absPow_art). (3) is above_w_statement: not_link with the private cyc_inf_over, locAlg, the private above' (a copy of the private above of DualOfWeights.lean, Lemma sha-hecke-ideal-above) and the private noLk_above (not_linked for the triple, noLk_of).

\begin{definition}\label{sha-automorphy-lift-forms}
Let $L$ be a CM field, let $n\ge1$, and let $U_n$, $U_{n,w}$ and $c_w$ be as in Definition~\ref{gal-definite-unitary}.
\begin{enumerate}
\item Let $w$ be a finite place of $L$ split over $L^+$, so that $c_w:U_{n,w}\to GL_n(L_w)$ is an isomorphism, and let $N_0$ and $\bar N_c$ be as in Definition~\ref{sha-hecke-local-ordinary}. For $c\ge0$, $J_w(c)\subset GL_n(L_w)$ is the subgroup generated by $N_0$, $\bar N_c$ and the diagonal matrices with entries in $1+\varpi_w^{c+1}\mathcal O_w$; the \emph{Iwahori subgroup} $J_w$ is generated by $N_0$, $\bar N_0$ and the diagonal matrices with entries in $\mathcal O_w^\times$.
\item Let $K_0$ be the group of elements of the finite part of $U_n$ whose image under every $c_w$ lies in $GL_n(\mathcal O_w)$. A \emph{level} is an open subgroup $K\subset K_0$; it is \emph{maximal} at $w$ if it contains $K_0\cap U_{n,w}$. For a finite set $S$ of places split over $L^+$ and groups $J'_w\in\{J_w,J_w(c)\}$ for $w\in S$, $K(S)$ is the level of the $k\in K$ with $c_w(k)\in J'_w$ for $w\in S$; it is written $K(S_0,c)$ when $J'_w=J_w(c)$ for all $w\in S_0$, and $K(S_1)$ when $J'_w=J_w$ for all $w\in S_1$.
\item Let $V$ be a finite-dimensional complex representation of $\prod_{v\mid\infty}GL_n(\mathbf C)$, whose subgroup $\prod_vU(n)$ is the part of $U_n$ at infinity. A \emph{form of weight $V$ and level $K$} is a function $\phi:U_n\to V$ with $\phi(\gamma g)=\phi(g)$ for $\gamma\in U_n(L^+)$, $\phi(gx)=x^{-1}\phi(g)$ for $x$ in the part at infinity, and $\phi(gk)=\phi(g)$ for $k\in K$.
\item Let $\phi$ be such a form with $c_w(K)\subset J'_w$ for a split place $w$. For $d\in(L_w^\times)^n$ let $T_d\phi=\sum_y\phi(\,\cdot\,c_w^{-1}(y))$, the sum over $y\in J'_w\operatorname{diag}(d)J'_w/J'_w$. For smooth characters $\psi=(\psi_1,\dots,\psi_n)$ of $L_w^\times$, $\phi$ is a \emph{$(J'_w,\psi)$-eigenform} if $T_d\phi=\prod_i\psi_i(d_i)\,\delta_B(\operatorname{diag}(d))^{-1/2}\phi$ whenever $|d_1|<\dots<|d_n|$, where $\delta_B$ is the modulus character of the upper triangular Borel subgroup. For $n=2$ and a character $b$ of $L_w^\times$, $\phi$ is a \emph{$(J'_w,b)$-eigenform for the $U_t$} if $T_{(t,t^{-1})}\phi=b(t)\delta_B(\operatorname{diag}(t,t^{-1}))^{-1/2}\phi$ for every $t$ with $|t|<1$.
\item For $j\ge0$ and $e:\{v\mid\infty\}\to\mathbf Z$, $V_j\otimes{\det}^e$ is the space of polynomials in variables $x_{v,0},x_{v,1}$, one pair for each infinite place $v$, homogeneous of degree $j$ in each pair, with $GL_2(\mathbf C)$ at $v$ acting by substitution in the pair at $v$, twisted by $\prod_v\det(g_v)^{e(v)}$. For $j\ge1$, $W_j$ is the representation of $\prod_vGL_3(\mathbf C)$ that is the sum, over the sets $T$ of infinite places, of the spaces of polynomials in two triples of variables at each $v$, homogeneous of degree $j-1$ if $v\in T$ and $j$ if $v\notin T$ in each triple, with $GL_3(\mathbf C)$ at $v$ acting through the triples at $v$.
\end{enumerate}
\end{definition}
\leanref{Development/AutomorphicArithmetic/AutomorphicVarieties/ModularityArgument/SymmetricPowerTransfer/SignatureIntegrality/UnitaryDoubleCosetIndicator.lean}{UpperTheory.V.Z.Sym.Sig.Lev}
\leanref{Development/AutomorphicArithmetic/AutomorphicVarieties/ModularityArgument/SymmetricPowerTransfer/SignatureIntegrality/UnitaryDoubleCosetIndicator.lean}{UpperTheory.V.Z.Sym.Sig.Lev.J}
\leanref{Development/AutomorphicArithmetic/AutomorphicVarieties/ModularityArgument/SymmetricPowerTransfer/GeneralizedEigenspaces/IsLevel.lean}{UpperTheory.V.Z.Sym.Dict.lev}
\leanref{Development/AutomorphicArithmetic/AutomorphicVarieties/ModularityArgument/SymmetricPowerTransfer/GeneralizedEigenspaces/IsLevel.lean}{UpperTheory.V.Z.Sym.Dict.IsAF}
\leanref{Development/AutomorphicArithmetic/AutomorphicVarieties/ModularityArgument/SymmetricPowerTransfer/GeneralizedEigenspaces/IsLevel.lean}{UpperTheory.V.Z.Sym.Dict.EigT}
\leanref{Development/AutomorphicArithmetic/AutomorphicVarieties/ModularityArgument/SymmetricPowerTransfer/GeneralizedEigenspaces/IsLevel.lean}{UpperTheory.V.Z.Sym.Dict.EigS}
\leanref{Development/AutomorphicArithmetic/AutomorphicVarieties/ModularityArgument/SymmetricPowerTransfer/WeightTraceProducts/RepGApply.lean}{UpperTheory.V.Z.Sym.Wt.repG}

\begin{proposition}[Alpöge's comparison of forms]\label{sha-automorphy-lift-comparison}
Let $L$ be a CM field such that $L/L^+$ is unramified at every finite place, let $S_0$ be a finite set of finite places of $L$, split over $L^+$, no two over the same place of $L^+$, such that every place of $L$ dividing $2$ lies over the same place of $L^+$ as a place of $S_0$, and let $j\ge1$. Then there is $c_0$ such that for every $c\ge c_0$ there is a level $K_2$ of $U_2$ with the following property, in the notation of Definition~\ref{sha-automorphy-lift-forms}.
Let $S_1$ be a finite set of places split over $L^+$, disjoint from $S_0$, such that no two places of $S_0\cup S_1$ lie over the same place of $L^+$ and $K_2$ is maximal at the places of $S_1$. Let $b_w$ be characters of $L_w^\times$, let $e:\{v\mid\infty\}\to\mathbf Z$, and let $\phi\ne0$ be a form on $U_2$ of weight $V_j\otimes{\det}^e$ and level $K_0(S_0,c)(S_1)$ that is a $(J_w(c),b_w)$-eigenform for the $U_t$ at every $w\in S_0$ and a $(J_w,b_w)$-eigenform for the $U_t$ at every $w\in S_1$. Then:
\begin{enumerate}
\item either there is a nonzero form on $U_3$ of weight $W_j$ and level $K_0(S_0,c)(S_1)$ that is a $(J_w(c),(b_w,1,b_w^{-1}))$-eigenform at every $w\in S_0$ and a $(J_w,(b_w,1,b_w^{-1}))$-eigenform at every $w\in S_1$;
\item or there is a nonzero form on $U_2$ of weight $V_{2j+1}$ and level $K_2(S_1)$ that is a $(J_w,b_w^2)$-eigenform for the $U_t$ at every $w\in S_1$.
\end{enumerate}
\end{proposition}
\leanref{Development/AutomorphicArithmetic/AutomorphicVarieties/ModularityArgument/SymmetricPowerTransfer/WeilDeligneUnramified/LevUnion.lean}{UpperTheory.V.Z.Sym.W.SpecOutS}
\leanref{Development/AutomorphicArithmetic/AutomorphicVarieties/ModularityArgument/SymmetricPowerTransfer/WeilDeligneUnramified/LevUnion.lean}{UpperTheory.V.Z.Sym.W.specOutS}
\begin{proof}
Take $c_0=0$, and let $c\ge0$. For a totally negative $\Delta\in L^+$ with $L=L^+(\sqrt\Delta)$, let $D$ be the $L^+$-algebra of the matrices $\begin{pmatrix}a&b\\-\bar b&\bar a\end{pmatrix}$ with $a,b\in L$, a definite quaternion algebra over $L^+$, and $D^1$ its group of elements of determinant $1$. A matrix $g$ of this shape has ${}^t\bar gg=\det(g)\cdot1$; conversely, if $g\in U_2(L^+)$ has determinant $1$, then ${}^t\bar g=g^{-1}$ is the adjugate of $g$, which forces this shape. So $D^1$ is the group of the elements of $U_2(L^+)$ of determinant $1$ (Definition~\ref{gal-definite-unitary}).
\begin{enumerate}
\item[(1)] \emph{Quaternionic data.} As $L/L^+$ is unramified at every finite place, there is such a $\Delta$ for which the identification of $D^1$ with the elements of determinant $1$ of $U_2(L^+)$ extends to the finite adeles, and the unitary group in three variables over $L^+(\sqrt\Delta)$ is identified with $U_3$, compatibly with the splittings at the places of $L^+$ split in $L$ and with the maximal compact subgroups. For $S_0$ and the groups $J_w(c)$, $w\in S_0$, there are local data for $D$ at the places of $L^+$ below $S_0$ that match these groups, a finite set $Y$ of representatives of the double cosets of $D^\times(\mathbf A_{L^+,f})$ modulo $D^\times(L^+)$ on the left and $D^1(\mathbf A_{L^+,f})$ and the level of these data on the right, and an integer $m\ge3$. Let $K_2$ be the principal congruence subgroup of level $m$ of $U_2$.
\item[(2)] \emph{The enlarged set.} Let $S_1$, the $b_w$, $e$ and $\phi$ be as in the statement. Put $S=S_0\cup S_1$, with $J'_w=J_w(c)$ for $w\in S_0$ and $J'_w=J_w$ for $w\in S_1$, so that $K_0(S_0,c)(S_1)=K_0(S)$. The data of (1) extend to local data at the places below $S$ that match the $J'_w$, and there are characters $\beta_v$ of $L_v^{+\times}$, for the places $v$ of $L^+$ below $S$, that are the $b_w$ through the splittings; then the $\beta_v^2$ are the $b_w^2$.
\item[(3)] \emph{The dichotomy.} A comparison of the traces of the Hecke operators on the forms on $D^1$, on $D^\times$ and on $U_3$ shows that either there is a form on $U_3$ as in (1) of the statement, or there is a nonzero automorphic form $\phi'$ on $D^\times(\mathbf A_{L^+,f})$ of weight $V_{2j+1}$ and of the level of the data of (2) that is an eigenform of the Hecke operators at the places below $S$ with the characters $\beta_v^2$.
\item[(4)] \emph{The return.} In the second case, $\phi'$ does not vanish at some $tx$ with $t\in D^1(\mathbf A_{L^+,f})$ and $x\in Y$; the form $g\mapsto\phi'(gx)$ on $D^1(\mathbf A_{L^+,f})$ is then nonzero, it is an eigenform with the characters $\beta_v^2$, and its level contains the image of $K_2$ at the places of $S_1$. Through the identification of (1), it gives a nonzero form on $U_2$ of weight $V_{2j+1}$ and level $K_2(S_1)$ that is a $(J_w,b_w^2)$-eigenform for the $U_t$ at every $w\in S_1$. This is (2) of the statement.
\end{enumerate}
\openstep{Steps (1) to (4) are not yet carried out in detail: the choice of $\Delta$ with the identifications and local data of (1) and (2), the comparison of traces in (3), and the passage from $D^\times$ back to $U_2$ in (4).}
% Size: about 24,000 lines of Lean (see the next comment). By parts: (1) and (2) about 3,000 (bridgeEx alone 2,950), (3) about 21,800 (specF with inS, baseG, baseS, outG), (4) about 7,000 (kNon, kRes, levExt, outS), with large overlaps. Each step should become a lemma once its data are stated in full.
% (1): bridgeEx gives the Bridge B (Dict.Br.Bridge: RingBr with Delta, e : Lsq(L^+, Delta) = L, and the isomorphisms bG : GA Delta = finU L 3, bS : SA Delta = finU L 2 cap SUA L, rG : U3 Delta = ratU L 3, rS : SU Delta = ratU L 2 cap SUA L) with B.Ok (split, chi, finG, finS, KG, oneG, splG, KS, oneS, splS, conj, cong); DX Delta is the group of g in GL_2(L) with conj(g) = wJ^-1 g wJ, which is D^x, and SU Delta its determinant-one subgroup; T : Spl B chooses the splittings. fitEx for S_0 and deep c gives the datum A_0 (So.SD: Sg, s, V, Ry, chi) with Fit; iotaEx gives iota; repsEx gives Y with SD.Reps; congLev gives m with 3 <= m and the conditions hm.2.1, hm.2.2; congEx gives U = K_2 with IsCong m U and IsLevel. c_0 = 0 in the Lean.
% (2): lev_union; fitEx for S and V (deep c on S_0, iw on S_1) gives A with Fit and b with FitCh; fitExt gives A.Ext A_0; fitSq gives FitCh for a^2 and b^2.

% specOutS (LevUnion.lean): bridgeEx (the quaternionic Bridge B of L), fitEx, iotaEx, repsEx, congLev, congEx, fitExt and lev_union; the private core (inS, baseG, baseS, specF, outG; baseK' with Wt.cK and Wt.ch) and second (kNon, kRes, levExt, outS), with fitSq and offAll. Its dependencies lie in SymmetricPowerTransfer/GeneralizedEigenspaces/CharacterShapes (about 8,100 lines), SymmetricPowerTransfer/GlobalUnitaryForms (about 13,500) and SymmetricPowerTransfer/GeneralizedEigenspaces/GlobalRationalGL (about 2,600).
\end{proof}

\begin{lemma}\label{sha-automorphy-lift-forms-facts-reps-projections}
Let $L$ be a CM field, let $n\ge1$, and use the notation of Definition~\ref{gal-definite-unitary}. Call a continuous function on $U_n$ an \emph{automorphic form} if it is invariant on the left under $U_n(L^+)$, invariant on the right under an open subgroup of the finite part, and its right translates by the part at infinity span a finite-dimensional space. For every finite-dimensional space $M$ of automorphic forms there are finitely many automorphic representations $\sigma_1,\dots,\sigma_r$ of $U_n$ and linear maps $P_1,\dots,P_r$ from the space of automorphic forms to itself, commuting with right translation by $U_n$, such that $P_k(\varphi)\in\sigma_k$ for every $\varphi$ and $\sum_kP_k(\varphi)=\varphi$ for $\varphi\in M$.
\end{lemma}
\leanref{Development/AutomorphicArithmetic/AutomorphicVarieties/ModularityArgument/SymmetricPowerTransfer/GeneralizedEigenspaces/SmoothDepth/OpenLevelSubgroups/AutProj.lean}{UpperTheory.V.Z.Sym.Al.autProj}
\begin{proof}
\begin{enumerate}
\item \emph{A compact quotient.} By Lemma~\ref{sha-hecke-ideal-tw-augmented-hida-classes}, applied to a compact open subgroup $C$ of the finite part, every element of $U_n$ is $\gamma sxk$ with $\gamma\in U_n(L^+)$, $s$ in a finite set $\Sigma$, $x$ in the part at infinity, which is compact, and $k\in C$. So $U_n=U_n(L^+)\cdot\Sigma U_\infty C$ with $\Sigma U_\infty C$ compact, and integration over a fundamental domain of finite Haar measure gives a Hermitian form on the continuous functions invariant on the left under $U_n(L^+)$, invariant under right translation by $U_n$, with $\langle\varphi,\varphi\rangle\ne0$ for $\varphi\ne0$.
\item \emph{Automorphic representations.} Every automorphic form lies in a finite sum of automorphic representations: this is the discreteness of the space of automorphic forms on a group whose automorphic quotient is compact \cite{GGPS69}, which Alpöge derives from the invariant form of (1), applied to the finite-dimensional, hence completely reducible, span of the translates of the form by the part at infinity.
\item \emph{Projections.} By (2), the space of automorphic forms is a sum of irreducible representations of $U_n$ under right translation, so it is semisimple as a module over the group algebra of $U_n$: it is a direct sum of automorphic representations. A basis of $M$, and so $M$, lies in the sum $\sigma_1\oplus\dots\oplus\sigma_r$ of finitely many of the summands. The projections $P_k$ onto $\sigma_k$, along the sum of the other summands, have the stated properties.
\end{enumerate}
\end{proof}
% autProj (AutProj.lean): exists_compact_fundamental_set_unitary (definiteUnitary_classSet_finite) and exists_nondegenerate_right_invariant_form_of_compact_quotient for (1); mem_iSup_automorphicReps_of_invariant_form (mem_iSup_irreducible_of_admissible_unitary and admissible_equivariantForms_representation) for (2); exists_equivariant_projections_of_sum_irreducible (IsSemisimpleModule.of_sSup_simples_eq_top) for (3).

\begin{lemma}\label{sha-automorphy-lift-forms-facts-reps-subquotient-form}
Let $G$ be a commutative topological group whose topology is separable, and let $\rho$ be a smooth representation of $G$ of finite length on a complex vector space $V$: every vector has an open stabilizer, and every nonempty set of subrepresentations has a maximal and a minimal element. Let $\chi:G\to\mathbf C^\times$ be a homomorphism and $x\in V$ a nonzero vector such that, for every $g\in G$, some power of $\rho(g)-\chi(g)$ kills $x$. Then $\chi$ has open kernel, and there is a nonzero linear form $l$ on $V$ with $l(\rho(g)y)=\chi(g)\,l(y)$ for all $g\in G$ and $y\in V$.
\end{lemma}
\leanref{Development/AutomorphicArithmetic/AutomorphicVarieties/AlgebraTools/IwJqEmbed.lean}{exists_form}
\begin{proof}
Subrepresentations are the subspaces stable under $\rho(G)$.
\begin{enumerate}
\item \emph{A subrepresentation that misses $x$.} Let $N$ be maximal among the subrepresentations not containing $x$; the subrepresentation $0$ is one. Let $g\in G$ and $f=\rho(g)-\chi(g)$. As $G$ is commutative, $f$ commutes with $\rho(G)$, so the subspaces $f^{-m}(N)=\{y:f^my\in N\}$ are subrepresentations; they increase with $m$, so they are equal for $m\ge a$, for some $a$. Let $b$ be such that $f^bx=0$. The subrepresentation $R=N+f^{a+b}(V)$ does not contain $x$: if $x=y+f^{a+b}z$ with $y\in N$, then $f^{a+2b}z=f^b(x-y)=-f^by\in N$, so $z\in f^{-(a+2b)}(N)=f^{-(a+b)}(N)$ and $x\in N$, which is false. By maximality $R=N$, that is, $f^{a+b}$ maps $V$ into $N$.
\item \emph{A character.} As $x\notin N$, $N\ne V$, so there is a subrepresentation $M\supseteq N$ maximal among the proper ones. The quotient $V/M$ is irreducible and smooth, so by Schur's lemma for the commutative group $G$, whose topology is separable \cite{BZ76}, each $\rho(g)$ acts on it by a scalar $\mu(g)$, and $\mu:G\to\mathbf C^\times$ is a homomorphism. Let $l\ne0$ be a linear form on $V$ that vanishes on $M\ne V$. As $\rho(g)y-\mu(g)y\in M$, $l(\rho(g)y)=\mu(g)\,l(y)$.
\item \emph{$\mu=\chi$.} Choose $y$ with $l(y)\ne0$. For $g\in G$, with $f$, $a$ and $b$ as in (1), $l(f^{a+b}y)=(\mu(g)-\chi(g))^{a+b}l(y)$, and $f^{a+b}y\in N\subseteq M$, so this is $0$ and $\mu(g)=\chi(g)$.
\item \emph{The kernel of $\chi$.} The stabilizer of $y$ is open, and for $g$ in it $l(y)=l(\rho(g)y)=\chi(g)\,l(y)$, so $\chi(g)=1$.
\end{enumerate}
\end{proof}
% exists_form (IwJqEmbed.lean, private; a copy in TLocSubD.lean): (1) is the private exists_nil (N maximal among the subrepresentations missing x, the chain K m = comap f^m, R = N + range f^(a+b)); (2) eq_top_or_exists_le_coatom, isIrreducible_subQuot and GlobalTheory.TorusEigenform.exists_char_of_isIrreducible (LocGL.exists_eq_smul_of_isSmooth, Schur's lemma), Submodule.exists_le_ker_of_lt_top; (3) hmuchi; (4) Subgroup.isOpen_of_mem_nhds.

\begin{lemma}\label{sha-automorphy-lift-forms-facts-reps-subquotient-theta}
Let $F$ be a nonarchimedean local field with uniformizer $\varpi$ and residue field $k_F$ of $q$ elements, let $n\ge1$, let $C\subset k_F^\times$ be a subgroup, and let $\pi$ be a representation of $GL_n(F)$ on a complex vector space $V$. Let $\Theta$ be the homomorphism of Lemma~\ref{sha-hecke-ideal-tw-torus}, for the single group $GL_n(F)$, the subgroup $\mathrm{Iw}_C$ and $A=\mathbf C$, from $T(F)$ to the automorphisms of $V^{\mathrm{Iw}_C}$ (Definition~\ref{sha-hecke-ideal-tw-iwahori}). Let $N\subset GL_n(F)$ be the group of upper triangular unipotent matrices, $V(N)$ the span of the vectors $\pi(u)y-y$ with $u\in N$ and $y\in V$, $p:V\to V_N=V/V(N)$ the projection to the Jacquet module, and $\delta$ the modulus character of the upper triangular Borel subgroup, $\delta(\operatorname{diag}(t_1,\dots,t_n))=\prod_{i<j}|t_i/t_j|$ with $|\varpi|=q^{-1}$. Then
\[p(\Theta(t)w)=\delta(t)^{-1}\,p(\pi(t)w)\]
for every $t\in T(F)$ and every $w\in V^{\mathrm{Iw}_C}$.
\end{lemma}
\leanref{Development/AutomorphicArithmetic/AutomorphicVarieties/AlgebraTools/IwJqTheta.lean}{UpperTheory.V.AlgebraTools.jqTheta}
\begin{proof}
As $T(F)$ normalizes $N$, $\pi(t)$ maps $V(N)$ into itself for $t\in T(F)$, and $\sigma(t)p(y)=p(\pi(t)y)$ defines an action $\sigma$ of $T(F)$ on $V_N$. Let $S$ be the set of the $t\in T(F)$ with $p(\Theta(t)w)=\delta(t)^{-1}\sigma(t)p(w)$ for every $w\in V^{\mathrm{Iw}_C}$.
\begin{enumerate}
\item \emph{$S$ is a subgroup.} It contains $1$. Let $s,t\in S$ and $w\in V^{\mathrm{Iw}_C}$. As $\Theta(t)w\in V^{\mathrm{Iw}_C}$, $p(\Theta(st)w)=p(\Theta(s)\Theta(t)w)=\delta(s)^{-1}\sigma(s)p(\Theta(t)w)=\delta(st)^{-1}\sigma(st)p(w)$. The identity for $t$ at $\Theta(t^{-1})w$ reads $p(w)=\delta(t)^{-1}\sigma(t)p(\Theta(t^{-1})w)$, so $p(\Theta(t^{-1})w)=\delta(t^{-1})^{-1}\sigma(t^{-1})p(w)$.
\item \emph{Dominant powers of $\varpi$.} Let $a_1\ge\dots\ge a_n$ be integers and $t=\varpi^a=\operatorname{diag}(\varpi^{a_1},\dots,\varpi^{a_n})$, which is dominant, so that $\Theta(t)$ is the double coset operator $[\mathrm{Iw}_Ct\,\mathrm{Iw}_C]$. As in the proof of Proposition~\ref{gal-iwahori-operator-invertible}(3), there is a finite set $s$ of upper triangular unipotent matrices with entries in $\mathcal O_F$ such that the $ut$, $u\in s$, represent the cosets of $\mathrm{Iw}_Ct\,\mathrm{Iw}_C/\mathrm{Iw}_C$ once each. So $\Theta(t)w=\sum_{u\in s}\pi(u)\pi(t)w$, and, as $p(\pi(u)y)=p(y)$ for $u\in N$, $p(\Theta(t)w)=|s|\,\sigma(t)p(w)$. The number $|s|$ of these cosets is $q^{\sum_{i<j}(a_i-a_j)}$. Indeed, on the trivial representation of $GL_n(F)$ on $\mathbf Z$ the operator of a dominant $t'$ is multiplication by the number of cosets of $\mathrm{Iw}_Ct'\mathrm{Iw}_C/\mathrm{Iw}_C$, and these operators multiply along products of dominant matrices (Lemma~\ref{gal-iwahori-operator-dominant}). For $0\le m\le n$, let $\varpi_m$, $w_m$ and $z=P_{w_m}\varpi_m$ be as in Lemma~\ref{gal-iwahori-operator-fundamental}; then $\varpi_m=P_{w_m^{-1}}z$ with $z$ normalizing $\mathrm{Iw}_C$, so $\varpi_m$ has as many cosets as $P_{w_m^{-1}}$, which is $q^{m(n-m)}$ by Lemma~\ref{gal-iwahori-operator-permutation}(2), as $w_m^{-1}$ has $m(n-m)$ inversions. Now argue by induction on $N\ge0$, for $a$ with entries in $[c,c+N]$, as in the proof of Proposition~\ref{gal-iwahori-operator-invertible}(1). If $N=0$, $\varpi^a=\varpi^c$ is central and has one coset. If $N\ge1$, let $m$ be the number of indices $j$ with $a_j>c$, which are $1,\dots,m$, and $a'$ the vector $a$ with $1$ subtracted from its first $m$ entries; then $\varpi^a=\varpi^{a'}\varpi_m$ with both factors dominant, and $\sum_{i<j}(a_i-a_j)=\sum_{i<j}(a'_i-a'_j)+m(n-m)$. Finally $\delta(t)=\prod_{i<j}|\varpi|^{a_i-a_j}=q^{-\sum_{i<j}(a_i-a_j)}$, so $t\in S$.
\item \emph{Units.} Let $t\in T(F)$ have unit entries. It lies in the Iwahori subgroup, which normalizes $\mathrm{Iw}_C$, so $\mathrm{Iw}_Ct\,\mathrm{Iw}_C=t\,\mathrm{Iw}_C$ and $\Theta(t)w=\pi(t)w$ (Lemma~\ref{gal-iwahori-operator-normalizer}); and $\delta(t)=1$. So $t\in S$.
\item \emph{All of $T(F)$.} Let $t\in T(F)$, and write $t_i=u_i\varpi^{c_i}$ with $u_i\in\mathcal O_F^\times$ and $c_i\in\mathbf Z$. Let $M=\sum_i|c_i|$ and $b_i=-2M(i-1)$. Then $b$ and $c+b$ are decreasing: for $i<j$, $(c_i+b_i)-(c_j+b_j)=c_i-c_j+2M(j-i)\ge0$. So $t=\varpi^{c+b}(\varpi^b)^{-1}\operatorname{diag}(u)$ lies in $S$ by (1), (2) and (3).
\end{enumerate}
\end{proof}
% jqTheta (IwJqTheta.lean): (1) is the private jq_gen; (2) is hA, with jq_dom (domSum for the representatives s and jq_card for their number), domCard (IwDomCard.lean: hk_triv, hfund through exists_rot, rot_norm, hkNorm, permCard, invLen_inv and invLen_eq; the induction key with exists_cut and domMul), jq_sum and jq_delta; (3) is hB with unit_norm (memIwC) and hdelta1; (4) is unif_dec and anti_diff.

\begin{lemma}[Borel {\cite{Bor76}}]\label{sha-automorphy-lift-forms-facts-reps-subquotient-injective}
In Lemma~\ref{sha-automorphy-lift-forms-facts-reps-subquotient-theta}, suppose that $\pi$ is smooth. Then $p$ is injective on $V^{\mathrm{Iw}_C}$.
\end{lemma}
\leanref{Development/AutomorphicArithmetic/AutomorphicVarieties/AlgebraTools/IwJqInj.lean}{UpperTheory.V.AlgebraTools.jqInj}
\begin{proof}
For a subgroup $H$ of $GL_n(F)$ let $V(H)$ be the span of the vectors $\pi(h)y-y$ with $h\in H$ and $y\in V$. Let $w\in V^{\mathrm{Iw}_C}$ with $p(w)=0$, so that $w$ is a finite sum of vectors $\pi(u)y-y$ with $u\in N$. For $m\ge0$ let $a_m=\operatorname{diag}(1,\varpi^{-m},\varpi^{-2m},\dots,\varpi^{-(n-1)m})$, which is dominant. For $u\in N$, the entry $(i,j)$ of $a_mua_m^{-1}$, for $i<j$, is $\varpi^{m(j-i)}u_{ij}$, so $a_mua_m^{-1}$ has entries in $\mathcal O_F$, and lies in $\mathrm{Iw}_C$, for $m$ large. As $\pi(a_m)(\pi(u)y-y)=\pi(a_mua_m^{-1})y'-y'$ with $y'=\pi(a_m)y$, there is $m$ with $\pi(a_m)w\in V(\mathrm{Iw}_C)$. As in Lemma~\ref{sha-automorphy-lift-forms-facts-reps-subquotient-theta}(2), $[\mathrm{Iw}_Ca_m\mathrm{Iw}_C]w=\sum_{u\in s}\pi(u)\pi(a_m)w$ with $s\subset N\cap\mathrm{Iw}_C$, and $\pi(u)z=(\pi(u)z-z)+z\in V(\mathrm{Iw}_C)$ for $z\in V(\mathrm{Iw}_C)$ and $u\in\mathrm{Iw}_C$. So $w'=[\mathrm{Iw}_Ca_m\mathrm{Iw}_C]w$ is fixed by $\mathrm{Iw}_C$ and lies in $V(\mathrm{Iw}_C)$. Write $w'$ as a finite sum of vectors $\pi(h)y-y$ with $h\in\mathrm{Iw}_C$. As $\pi$ is smooth and $\mathrm{Iw}_C$ is compact, there is an open normal subgroup $L$ of finite index in $\mathrm{Iw}_C$ that fixes all these $y$. The sum of the $\pi(g)$ over representatives $g$ of $\mathrm{Iw}_C/L$ kills each $\pi(h)y-y$, as the $gh$ are again such representatives, and it multiplies $w'$ by $[\mathrm{Iw}_C:L]$. So $w'=0$. As $[\mathrm{Iw}_Ca_m\mathrm{Iw}_C]$ is invertible (Proposition~\ref{gal-iwahori-operator-invertible}(1)), $w=0$.
\end{proof}
% jqInj (IwJqInj.lean): jqInj_ev (a_m u a_m^-1 in unipO eventually), jqInj_dom, domSum, domMul (the unit), jqInj_compact and GlobalTheory.IwahoriInj.eq_zero_of_fixed (TwistedKernelFixed.eq_zero_of_fixed: sum_shift over H / L for the normal core L of the stabilizers). The published statement is Borel's: the vectors fixed by an Iwahori subgroup map isomorphically to the invariants of the integral torus in the Jacquet module [Bor76]; the Lean proves the injectivity its own way.

\begin{lemma}\label{sha-automorphy-lift-forms-facts-reps-subquotient-jacquet}
Let $F$, $q$, $n$, $C$, $V$, $\Theta$, $N$, $p$ and $\delta$ be as in Lemma~\ref{sha-automorphy-lift-forms-facts-reps-subquotient-theta}, for an irreducible smooth representation $\pi$ of $GL_n(F)$ on $V$. Let $\chi:T(F)\to\mathbf C^\times$ be a homomorphism, and suppose that some nonzero $w\in V^{\mathrm{Iw}_C}$ lies, for every $t\in T(F)$, in the generalized eigenspace of $\Theta(t)$ for $\chi(t)$. Then there are smooth characters $\psi_1,\dots,\psi_n$ of $F^\times$, trivial on the units whose residues lie in $C$, with $\prod_i\psi_i(t_i)=\chi(t)\delta(t)^{1/2}$ for $t=\operatorname{diag}(t_1,\dots,t_n)\in T(F)$, such that $\pi$ is a subquotient of the normalized parabolic induction $I(\psi)$, from the upper triangular Borel subgroup, of $\operatorname{diag}(t_1,\dots,t_n)\mapsto\prod_i\psi_i(t_i)$.
\end{lemma}
\leanref{Development/AutomorphicArithmetic/AutomorphicVarieties/AlgebraTools/IwJqEmbed.lean}{UpperTheory.V.AlgebraTools.jqEmbed}
\begin{proof}
Let $T(F)$ act on $V_N$ by $\sigma(t)p(y)=\delta(t)^{-1/2}p(\pi(t)y)$. This is a smooth representation of finite length, as $\pi$ is irreducible and smooth \cite{BZ77}. By Lemma~\ref{sha-automorphy-lift-forms-facts-reps-subquotient-theta}, $\sigma(t)p(y)=\delta(t)^{1/2}p(\Theta(t)y)$ for $y\in V^{\mathrm{Iw}_C}$.
\begin{enumerate}
\item \emph{A generalized eigenvector.} By Lemma~\ref{sha-automorphy-lift-forms-facts-reps-subquotient-injective}, $p(w)\ne0$. Let $t\in T(F)$ and $m$ with $(\Theta(t)-\chi(t))^mw=0$. Then $(\sigma(t)-\chi(t)\delta(t)^{1/2})^mp(w)=\delta(t)^{m/2}p\big((\Theta(t)-\chi(t))^mw\big)=0$.
\item \emph{A form.} The group $T(F)$ is commutative and its topology is separable. Lemma~\ref{sha-automorphy-lift-forms-facts-reps-subquotient-form}, for $\sigma$, $x=p(w)$ and the character $\chi\delta^{1/2}$, shows that $\chi\delta^{1/2}$ has open kernel and gives a nonzero linear form $l$ on $V_N$ with $l(\sigma(t)y)=\chi(t)\delta(t)^{1/2}l(y)$.
\item \emph{The characters.} For $a\in F^\times$ let $e_i(a)\in T(F)$ have $a$ in place $i$ and $1$ elsewhere, and $\psi_i(a)=\chi(e_i(a))\delta(e_i(a))^{1/2}$. Each $\psi_i$ has open kernel, and $\prod_i\psi_i(t_i)=\chi(t)\delta(t)^{1/2}$, as $t=\prod_ie_i(t_i)$. Let $a$ be a unit with residue in $C$. Then $e_i(a)\in\mathrm{Iw}_C$ fixes $w$, and $\delta(e_i(a))=1$, so $\sigma(e_i(a))p(w)=p(w)$, and (1) gives $(1-\psi_i(a))^mp(w)=0$ for some $m$; hence $\psi_i(a)=1$.
\item \emph{The principal series.} By Frobenius reciprocity \cite{BZ77}, the form $l$ of (2) makes $\pi$ a subquotient of $I(\psi)$.
\end{enumerate}
\end{proof}
% jqEmbed (IwJqEmbed.lean): hl from the hypothesis BernsteinZelevinskyFiniteness (JacquetFiniteLength); hq from jqTheta; hx from jqInj; hxc for (1); the private exists_form with isSmooth_jacB for (2); htc and the unitsC argument (hfix with UnramFixed.deltaHalf_torusO) for (3); GlobalTheory.SatakeExists.isSubquotient_of_eigenform (Model.evalOne_bijective, Frobenius reciprocity) for (4).

\begin{lemma}\label{sha-automorphy-lift-forms-facts-reps-subquotient}
Let $K$, $\varpi$, $n$, $N_0$, $\bar N_c$ and $T_0$ be as in Definition~\ref{sha-hecke-local-ordinary}. Let $J$ be the subgroup of $GL_n(K)$ generated by $N_0$, $\bar N_0$ and $T_0$, for $c\ge0$ let $J(c)$ be the subgroup generated by $N_0$, $\bar N_c$ and the diagonal matrices with entries in $1+\varpi^{c+1}\mathcal O_K$, and let $J'\in\{J,J(c)\}$. For a representation $\tau$ of $GL_n(K)$ and $g\in GL_n(K)$, let $[J'gJ']$ be the operator $v\mapsto\sum_y\tau(y)v$ on the vectors fixed by $\tau(J')$, the sum over representatives $y$ of $J'gJ'/J'$. Let $\psi_1,\dots,\psi_n$ be smooth characters of $K^\times$, let $I(\psi)$ be the normalized parabolic induction, from the upper triangular Borel subgroup $B$, of $\operatorname{diag}(t_1,\dots,t_n)\mapsto\prod_i\psi_i(t_i)$, and let $\delta_B$ be the modulus character of $B$. Let $(\tau,T)$ be an irreducible smooth representation of $GL_n(K)$ with a nonzero $v$ fixed by $\tau(J')$ such that $[J'\operatorname{diag}(d)J']v=\prod_i\psi_i(d_i)\,\delta_B(\operatorname{diag}(d))^{-1/2}\,v$ for every $d$ with $|d_1|<\dots<|d_n|$. Then $\tau$ is a subquotient of $I(\psi)$.
\end{lemma}
\leanref{Development/AutomorphicArithmetic/AutomorphicVarieties/ModularityArgument/ThresholdArguments/DenseThresholdCases/TLocSub.lean}{UpperTheory.V.Z.Sym.Al.subquotient_prin_series_of_iwahori_eigenvector}
\leanref{Development/AutomorphicArithmetic/AutomorphicVarieties/ModularityArgument/ThresholdArguments/DenseThresholdCases/TLocSubD.lean}{UpperTheory.V.Z.Sym.W.local_sub_d_statement}
\begin{proof}
For $d=(d_1,\dots,d_n)\in(K^\times)^n$ write $\operatorname{diag}(d)$ for the diagonal matrix with entries $d_i$; call $d$ \emph{dominant} if $|d_1|\le\dots\le|d_n|$ and \emph{strictly dominant} if $|d_1|<\dots<|d_n|$. Let $\chi(d)=\prod_i\psi_i(d_i)\,\delta_B(\operatorname{diag}(d))^{-1/2}$, a homomorphism from $(K^\times)^n$ to $\mathbf C^\times$, so that $[J'\operatorname{diag}(d)J']v=\chi(d)v$ for strictly dominant $d$. Let $k$ be the residue field of $K$, and $N$ the group of upper triangular unipotent matrices.
\begin{enumerate}
\item[(1)] \emph{Strictly dominant multiples.} For every $t\in(K^\times)^n$ there is a strictly dominant $d$ with $td$ strictly dominant. Indeed, choose a pair $(i_0,j_0)$ at which $|t_{j_0}/t_{i_0}|$ is smallest, and let $c_0=t_{j_0}/t_{i_0}$. Then $|c_0|\le1$, by the pair $(i_0,i_0)$, and $|t_i||c_0|\le|t_j|$ for all $i,j$. Let $c'=c_0\varpi$, so that $|c'|<1$, and $d_i=c'^{\,n-i}$. Then $|d_i|=|c'|^{n-i}$ increases strictly with $i$. For $i<j$, $|t_i||c'|<|t_j|$, so $|t_id_i|=|t_i||c'|\cdot|c'|^{n-j}\cdot|c'|^{j-i-1}\le|t_i||c'|\,|c'|^{n-j}<|t_j||c'|^{n-j}=|t_jd_j|$.
\end{enumerate}
\emph{The case $J'=J$.}
\begin{enumerate}
\item[(2)] \emph{The torus on $v$.} The group $J$ is the Iwahori subgroup $\mathrm{Iw}$ of Definition~\ref{sha-hecke-ideal-tw-iwahori}. Let $\Theta$ be the homomorphism of Lemma~\ref{sha-hecke-ideal-tw-torus}, for the single group $GL_n(K)$, the subgroup $\mathrm{Iw}$ and $A=\mathbf C$, from the diagonal torus to the automorphisms of the space of vectors fixed by $\tau(J)$. For strictly dominant $d$, $\operatorname{diag}(d)$ is dominant, so $\Theta(\operatorname{diag}(d))v=[J\operatorname{diag}(d)J]v=\chi(d)v$. Let $t\in(K^\times)^n$, and let $d$ be as in (1). Then $\chi(td)v=\Theta(\operatorname{diag}(t))\Theta(\operatorname{diag}(d))v=\chi(d)\,\Theta(\operatorname{diag}(t))v$, so $\Theta(\operatorname{diag}(t))v=\chi(t)v$.
\item[(3)] \emph{The principal series.} Lemma~\ref{sha-automorphy-lift-forms-facts-reps-subquotient-jacquet}, for $F=K$, $C=k^\times$, $w=v$ and the character $\operatorname{diag}(t)\mapsto\chi(t)$, gives smooth characters $\psi'_1,\dots,\psi'_n$ of $K^\times$ with $\prod_i\psi'_i(t_i)=\chi(t)\delta_B(\operatorname{diag}(t))^{1/2}=\prod_i\psi_i(t_i)$ for every $t$, such that $\tau$ is a subquotient of $I(\psi')$. Taking for $t$ the element with $a\in K^\times$ in place $i$ and $1$ elsewhere gives $\psi'_i=\psi_i$.
\end{enumerate}
\emph{The case $J'=J(c)$.} For a subgroup $H$ of $GL_n(K)$ let $\Delta(H)$ be the span of the vectors $\tau(h)y-y$ with $h\in H$ and $y\in T$, let $p:T\to T_N=T/\Delta(N)$ be the projection to the Jacquet module, and let $(K^\times)^n$ act on $T_N$ by $\sigma(t)p(y)=\delta_B(\operatorname{diag}(t))^{-1/2}p(\tau(\operatorname{diag}(t))y)$.
\begin{enumerate}
\item[(4)] \emph{Double cosets.} Let $d$ be dominant, and let $s_d\subset N_0$ be a set of representatives of $N_0/\operatorname{diag}(d)N_0\operatorname{diag}(d)^{-1}$. Step (4) of the proof of Lemma~\ref{sha-automorphy-lift-forms-facts-span-local} uses of $\varpi^a$ only that it is dominant: conjugation by $\operatorname{diag}(d)$ maps $N_0$ into itself, and conjugation by its inverse maps $\bar N_c$ into itself. So for $J''\in\{J,J(c)\}$, every representation $\rho$ of $GL_n(K)$ and every $y$ fixed by $\rho(J'')$, $[J''\operatorname{diag}(d)J'']y=\sum_{x\in s_d}\rho(x)\rho(\operatorname{diag}(d))y$. For the trivial representation of $GL_n(K)$ on $\mathbf C$, whose Jacquet module is $\mathbf C$, and $J''=J=\mathrm{Iw}$, Lemma~\ref{sha-automorphy-lift-forms-facts-reps-subquotient-theta} with $C=k^\times$ and $w=1$ gives $|s_d|=\delta_B(\operatorname{diag}(d))^{-1}$.
\item[(5)] \emph{$p(v)\ne0$.} Suppose $p(v)=0$. For $m\ge1$ let $a_m=(1,\varpi^{-m},\varpi^{-2m},\dots,\varpi^{-(n-1)m})$, which is strictly dominant. As in the proof of Lemma~\ref{sha-automorphy-lift-forms-facts-reps-subquotient-injective}, there is $m\ge1$ with $\tau(\operatorname{diag}(a_m))v\in\Delta(N_0)\subseteq\Delta(J(c))$. By (4), $[J(c)\operatorname{diag}(a_m)J(c)]v=\sum_{x\in s_{a_m}}\tau(x)\tau(\operatorname{diag}(a_m))v$, which lies in $\Delta(J(c))$, as $N_0\subset J(c)$, and is fixed by $\tau(J(c))$. As $\tau$ is smooth and $J(c)$ is compact, the averaging argument of that proof shows that this vector is $0$. But it is $\chi(a_m)v$, with $\chi(a_m)\ne0$ and $v\ne0$.
\item[(6)] \emph{The torus on $p(v)$.} Let $d$ be strictly dominant. By (4), and as $p(\tau(x)y)=p(y)$ for $x\in N$,
\[\chi(d)\,p(v)=p\big([J(c)\operatorname{diag}(d)J(c)]v\big)=|s_d|\,p(\tau(\operatorname{diag}(d))v)=\delta_B(\operatorname{diag}(d))^{-1}p(\tau(\operatorname{diag}(d))v),\]
so $\sigma(d)p(v)=\delta_B(\operatorname{diag}(d))^{1/2}\chi(d)p(v)=\prod_i\psi_i(d_i)\,p(v)$. For $t\in(K^\times)^n$ and $d$ as in (1), $\prod_i\psi_i(t_id_i)\,p(v)=\sigma(td)p(v)=\prod_i\psi_i(d_i)\,\sigma(t)p(v)$, so $\sigma(t)p(v)=\prod_i\psi_i(t_i)\,p(v)$.
\item[(7)] \emph{The principal series.} The representation $\sigma$ is smooth and of finite length \cite{BZ77}, and $(K^\times)^n$ is commutative with a separable topology. Lemma~\ref{sha-automorphy-lift-forms-facts-reps-subquotient-form}, for $\sigma$, the vector $p(v)\ne0$ of (5) and the character $t\mapsto\prod_i\psi_i(t_i)$, gives a nonzero linear form $l$ on $T_N$ with $l(\sigma(t)y)=\prod_i\psi_i(t_i)\,l(y)$. By Frobenius reciprocity \cite{BZ77}, $\tau$ is a subquotient of $I(\psi)$.
\end{enumerate}
% Case J, Al.subquotient_prin_series_of_iwahori_eigenvector (TLocSub.lean): (1) is the private exists_sdom (with sdom_of_step); (2) is Sig.link (J = iwC top), thetaEx, hsd (hkApply) and ht; (3) is jqEmbed with H.loc.bz and the private torusChar_injective.
% Case J(c), W.local_sub_d_statement (TLocSubD.lean), stated there for every level Sig.Lev: (4) is Sig.hkMul (hk_eq_uOp) with the private reps and idx (thetaEx and jqTheta on the trivial representation); (5) is mk_ne (jqInj_ev, jqInj_dom, Sig.finFacts for the compactness, IwahoriInj.eq_zero_of_fixed); (6) is jac_eig and ht with exists_sdom; (7) is H.loc.bz (BernsteinZelevinskyFiniteness), the private exists_form and SatakeExists.isSubquotient_of_eigenform.
\end{proof}

\begin{lemma}\label{sha-automorphy-lift-forms-facts-reps}
Assume that Galois representations exist, in the sense of Definition~\ref{gal-galois-representations}, relative to the maps satisfying the laws of Definition~\ref{gal-llc-laws}.
Let $L$ be a CM field such that $L/L^+$ is unramified at every finite place, use the notation of Definition~\ref{sha-automorphy-lift-forms}, let $n\ge1$ and let $V$ be a finite-dimensional representation of $\prod_vGL_n(\mathbf C)$ with continuous matrix coefficients. Say that a function $\phi$ on $U_n$ with values in $V$ has \emph{coefficients in} an automorphic representation $\pi$ of $U_n$ if $f\circ\phi\in\pi$ for every linear form $f$ on $V$. Then:
\begin{enumerate}
\item[(a)] for every level $K$ there is a finite set $A$ of automorphic representations of $U_n$, each with a nonzero vector fixed by $K$, such that for every automorphic representation $\pi$ of $U_n$ with a nonzero form of weight $V$ and level $K$ with coefficients in $\pi$, some $\sigma\in A$ has every local component of $\pi$, at every finite place, as a local component;
\item[(b)] for every level $K$ and every nonzero form $\phi$ of weight $V$ and level $K$ there are an automorphic representation $\pi$ of $U_n$ and a nonzero form $\phi'$ of weight $V$ and level $K$ with coefficients in $\pi$ such that, for every split place $w$, every $J'_w\in\{J_w,J_w(c)\}$, every $d\in(L_w^\times)^n$ and every $\mu\in\mathbf C$, $T_d\phi'=\mu\phi'$ whenever $T_d\phi=\mu\phi$, with $T_d$ formed for $J'_w$;
\item[(c)] let $\pi$ be an automorphic representation of $U_n$, $w$ a split place, $J'_w\in\{J_w,J_w(c)\}$, $\psi$ smooth characters of $L_w^\times$, and $\phi\ne0$ a function on $U_n$ with values in $V$ and coefficients in $\pi$ such that $\phi(gx)=\phi(g)$ for $x\in U_{n,w}$ with $c_w(x)\in J'_w$, and that is a $(J'_w,\psi)$-eigenform; then $\pi$ has at $w$ a local component that is a subquotient of $I(\psi)$;
\item[(d)] if an automorphic representation $\pi$ of $U_n$ has at a split place $w$ a local component that is a subquotient of $I(\psi)$ for separated $\psi$ (Lemma~\ref{sha-automorphy-lift-forms-facts}), then $\pi$ is unramified at $w$;
\item[(e)] let $K$ be a level, $S$ a finite set of split places, no two over the same place of $L^+$, at which $K$ is maximal, $\pi$ an automorphic representation of $U_n$ unramified at the places of $S$, and $J'_w\in\{J_w,J_w(c_w)\}$ for $w\in S$; if there is a nonzero form of weight $V$ and level $K(S)$ with coefficients in $\pi$, then there is a nonzero form of weight $V$ and level $K$ with coefficients in $\pi$.
\end{enumerate}
\end{lemma}
\leanref{Development/AutomorphicArithmetic/AutomorphicVarieties/ModularityArgument/SymmetricPowerTransfer/GeneralizedEigenspaces/SmoothDepth/OpenLevelSubgroups/FinFam.lean}{UpperTheory.V.Z.Sym.Al.finFam}
\leanref{Development/AutomorphicArithmetic/AutomorphicVarieties/ModularityArgument/SymmetricPowerTransfer/GeneralizedEigenspaces/ProjectionComponents/CompOfProj.lean}{UpperTheory.V.Z.Sym.U.finFamU}
\leanref{Development/AutomorphicArithmetic/AutomorphicVarieties/ModularityArgument/SymmetricPowerTransfer/GeneralizedEigenspaces/SmoothDepth/OpenLevelSubgroups/Imports.lean}{UpperTheory.V.Z.Sym.Al.proj}
\leanref{Development/AutomorphicArithmetic/AutomorphicVarieties/ModularityArgument/ThresholdArguments/DenseThresholdCases/TCompOfEig.lean}{UpperTheory.V.Z.Sym.Al.comp_sub_of_eigen_vector_fixed_by_j}
\leanref{Development/AutomorphicArithmetic/AutomorphicVarieties/ModularityArgument/ThresholdArguments/DenseThresholdCases/UnramifiedOfSplitDistinctCharacters.lean}{UpperTheory.V.Z.Sym.Al.unramified_of_split_distinct_characters}
\leanref{Development/AutomorphicArithmetic/AutomorphicVarieties/ModularityArgument/SymmetricPowerTransfer/GeneralizedEigenspaces/SmoothDepth/OpenLevelSubgroups/ImportsAlt.lean}{UpperTheory.V.Z.Sym.Al.raise}
\begin{proof}
Automorphic forms are as in Lemma~\ref{sha-automorphy-lift-forms-facts-reps-projections}, and for a function $\phi$ with values in $V$ and a linear form $f$ on $V$, $f\circ\phi$ is a \emph{coordinate} of $\phi$.
\begin{enumerate}
\item[(i)] \emph{Coordinates of forms.} Let $\phi$ be a form of weight $V$ and level $K$, and $f$ a linear form on $V$. Then $f\circ\phi$ is an automorphic form. It is invariant on the left under $U_n(L^+)$ and on the right under $K$. It is continuous: near $g_0$, $g=g_0x\kappa$ with $x$ near $1$ in the part at infinity and $\kappa\in K$, so $f(\phi(g))=f(x^{-1}\phi(g_0))$, a continuous function of $x$ as the matrix coefficients of $V$ are continuous. Its right translates by $x$ in the part at infinity are $(f\circ x^{-1})\circ\phi$, which lie in the finite-dimensional space of the coordinates of $\phi$.
\item[(a)] Let $F$ be the space of forms of weight $V$ and level $K$. It is finite-dimensional, by the argument of Lemma~\ref{sha-automorphy-lift-forms-facts-operators}(1), which does not depend on $n$. Let $M$ be the span of the coordinates of the elements of $F$, a finite-dimensional space of automorphic forms by (i), and let $\sigma_k$ and $P_k$ be as in Lemma~\ref{sha-automorphy-lift-forms-facts-reps-projections} for $M$. Let $A$ be the set of the $\sigma_k$ that have a nonzero vector fixed by $K$. Let $\pi$ and $\phi\in F$, $\phi\ne0$, be as in (a). Some coordinate $f\circ\phi$ is nonzero; it lies in $M\cap\pi$, and $\sum_kP_k(f\circ\phi)=f\circ\phi$, so $P_k(f\circ\phi)\ne0$ for some $k$. This vector of $\sigma_k$ is fixed by $K$, as $P_k$ commutes with right translation, so $\sigma_k\in A$. The restriction of $P_k$ to $\pi$ commutes with right translation and is nonzero, so it is injective, as $\pi$ is irreducible. If $\tau$ is a local component of $\pi$ at $w$, composing a nonzero intertwining map $\tau\to\pi$ with $P_k$ shows that $\tau$ is a local component of $\sigma_k$.
\item[(b)] Let $M$ be the space of the coordinates of $\phi$, finite-dimensional by (i), and let $\sigma_k$ and $P_k$ be as in Lemma~\ref{sha-automorphy-lift-forms-facts-reps-projections} for $M$. For each $k$ let $\phi_k$ be the function with values in $V$ whose coordinates are $f\circ\phi_k=P_k(f\circ\phi)$; it exists and is unique, as $f\mapsto P_k(f\circ\phi)$ is linear. Since $\sum_kf\circ\phi_k=f\circ\phi$ for every $f$, some $\phi_k$ is nonzero; let $\pi=\sigma_k$ and $\phi'=\phi_k$. Its coordinates lie in $\sigma_k$, and they are invariant on the left under $U_n(L^+)$. For $x$ in the part at infinity, $f\circ\phi'(\,\cdot\,x)=P_k\big((f\circ x^{-1})\circ\phi\big)=(f\circ x^{-1})\circ\phi'$, so $\phi'(gx)=x^{-1}\phi'(g)$; and in the same way $\phi'(gk)=\phi'(g)$ for $k\in K$. If $T_d\phi=\mu\phi$, then $f\circ T_d\phi'=P_k(f\circ T_d\phi)=\mu\,f\circ\phi'$ for every $f$, as $T_d$ is a sum of right translations, so $T_d\phi'=\mu\phi'$.
\item[(c)] By Theorem~\ref{gal-base-change-definite}(5), as $L/L^+$ is unramified, $\pi$ has a local component $\tau$ at $w$. Choose a linear form $f$ with $f\circ\phi\ne0$; then $f\circ\phi\in\pi$, and Lemma~\ref{sha-automorphy-lift-forms-facts-span-isotypic}(3) gives a linear map $p:\pi\to\tau$, compatible with the actions, with $v=p(f\circ\phi)\ne0$. As $f\circ\phi$ is invariant under right translation by the $x\in U_{n,w}$ with $c_w(x)\in J'_w$, $v$ is fixed by $\tau(J'_w)$; and as $T_d$ is a sum of right translations, $[J'_w\operatorname{diag}(d)J'_w]v=p(f\circ T_d\phi)=\prod_i\psi_i(d_i)\delta_B(\operatorname{diag}(d))^{-1/2}v$ whenever $|d_1|<\dots<|d_n|$. By Lemma~\ref{sha-automorphy-lift-forms-facts-reps-subquotient}, $\tau$ is a subquotient of $I(\psi)$.
\item[(d)] Let $\tau$ be the local component, a subquotient of $I(\psi)$. As $\psi$ is separated, no two of the characters $\psi_i$ differ by $|\cdot|^{\pm1}$, so $I(\psi)$ is irreducible and is its only irreducible subquotient \cite{BZ77,Zel80}; hence $\tau\cong I(\psi)$. As the $\psi_i$ are unramified, $I(\psi)$ has a nonzero vector fixed by $GL_n(\mathcal O_w)$: by the Iwasawa decomposition, the function $bk\mapsto\prod_i\psi_i(b_{ii})\delta_B(b)^{1/2}$, for $b$ upper triangular and $k\in GL_n(\mathcal O_w)$, is well defined and lies in $I(\psi)$. An intertwining map $\tau\to\pi$ is injective, as $\tau$ is irreducible, and carries this vector to a nonzero vector of $\pi$ fixed by $c_w^{-1}(GL_n(\mathcal O_w))$. So $\pi$ is unramified at $w$.
\item[(e)] By induction on the number of places of $S$; for $S=\emptyset$, $K(S)=K$. Let $w\in S$, $S'=S\setminus\{w\}$, and let $\phi\ne0$ be a form of weight $V$ and level $K(S)$ with coefficients in $\pi$. By Lemma~\ref{sha-automorphy-lift-forms-facts-span}(c), as $\pi$ is unramified at $w$, the span of $\phi$ at $w$ contains a nonzero $\phi'$ invariant under right translation by $U_{n,w}\cap K_0$. By Lemma~\ref{sha-automorphy-lift-forms-facts-span}(d), $\phi'$ has coefficients in $\pi$, is invariant on the left under $U_n(L^+)$, is equivariant under the part at infinity, and is invariant under right translation by the elements of $K(S)$ whose components over the place $v$ of $L^+$ under $w$ are $1$. Let $u\in K(S')$, and write $u=xy$ with $x\in U_{n,w}$ its component at $v$ and $y$ with component $1$ at $v$. Then $x\in U_{n,w}\cap K_0$; it lies in $K$, which is maximal at $w$, and in $K(S)$, as its components at the places of $S'$, which lie over other places of $L^+$, are $1$; so $y=x^{-1}u\in K(S)$. Hence $\phi'(gu)=\phi'(gxy)=\phi'(gx)=\phi'(g)$, and $\phi'$ is a nonzero form of weight $V$ and level $K(S')$ with coefficients in $\pi$. The induction hypothesis for $S'$ gives the conclusion.
\end{enumerate}
% (a) is finFam (FinFam.lean) and finFamU (CompOfProj.lean): afFin, coefAut ((i)), autProj, and the private comp_of_proj; (b) is proj (Imports.lean: autProj, F k with hF, htr for the translations, dcMap); (c) is comp_sub_of_eigen_vector_fixed_by_j for J_w and the private compD of MemLevelOfLocalDepth.lean for J_w(c): exists_local_component_at_split_of_unitary_aut (BaseChangeIdentities.splitComp), the first part of isoCopy with isoTypic, dcMap, and the local statements subquotient_prin_series_of_iwahori_eigenvector and W.local_sub_d_statement; (d) is unramified_of_split_distinct_characters (Hi.prin_series_irreducible_and_unique_subquotient, locSph); (e) is raise (spanSph, spanKeep, factor, by induction over S).
\end{proof}

\begin{lemma}\label{sha-automorphy-lift-forms-facts-operators-product}
Let $G$ be a group and $J\subset G$ a subgroup such that $JgJ/J$ is finite for every $g\in G$. For a representation $\rho$ of $G$ on a module $M$ over a commutative ring and $g\in G$, let $[JgJ]$ be the endomorphism $v\mapsto\sum_y\rho(y)v$ of the module $M^J$ of vectors fixed by $\rho(J)$, the sum over representatives $y$ of $JgJ/J$. Let $U$, $Z$ and $\bar U$ be subgroups of $G$ such that $Z$ normalizes $U$ and $\bar U$; such that $u=1$ and $z=1$ whenever $uz\bar u=1$ with $u\in U$, $z\in Z$ and $\bar u\in\bar U$; and such that every element of $J$ is a product $uz\bar u$ with $u\in U\cap J$, $z\in Z\cap J$ and $\bar u\in\bar U\cap J$. Call $m\in Z$ \emph{positive} if $mum^{-1}\in J$ for $u\in U\cap J$ and $m^{-1}\bar um\in J$ for $\bar u\in\bar U\cap J$. Let $m,m'\in Z$ be positive, and suppose that one of them, $m''$, satisfies $m''zm''^{-1}\in J$ and $m''^{-1}zm''\in J$ for $z\in Z\cap J$. Then $mm'$ is positive, and $[JmJ][Jm'J]=[Jmm'J]$ on $M^J$.
\end{lemma}
\leanref{Development/AutomorphicArithmetic/AutomorphicVarieties/AlgebraTools/IwLeviMul.lean}{UpperTheory.V.AlgebraTools.leviMul}
\leanref{Development/AutomorphicArithmetic/AutomorphicVarieties/AlgebraTools/IwHkMul.lean}{UpperTheory.V.AlgebraTools.hkMul}
\begin{proof}
For $x\in G$ let $J_x=J\cap xJx^{-1}$, the stabilizer of $xJ$ in $J$, so that $JxJ/J$ has $[J:J_x]$ elements.
\begin{enumerate}
\item \emph{Positivity of $mm'$.} For $u\in U\cap J$, $m'um'^{-1}$ lies in $U$, as $Z$ normalizes $U$, and in $J$, so $(mm')u(mm')^{-1}=m(m'um'^{-1})m^{-1}\in J$. In the same way $(mm')^{-1}\bar u(mm')\in J$ for $\bar u\in\bar U\cap J$.
\item \emph{Splitting $J$.} For $b\in J$ there is $b'\in J$ with $mb'm^{-1}\in J$ and $m'^{-1}b'^{-1}bm'\in J$. Write $b=uz\bar u$ as in the statement. If $m''=m$, take $b'=uz$: then $mb'm^{-1}=(mum^{-1})(mzm^{-1})\in J$ and $m'^{-1}b'^{-1}bm'=m'^{-1}\bar um'\in J$. If $m''=m'$, take $b'=u$: then $mum^{-1}\in J$ and $m'^{-1}z\bar um'=(m'^{-1}zm')(m'^{-1}\bar um')\in J$. In both cases $mbm'=(mb'm^{-1})\,mm'\,(m'^{-1}b'^{-1}bm')\in Jmm'J$.
\item \emph{$J_{mm'}\subseteq J_m$.} Let $g=uz\bar u\in J$ with $(mm')^{-1}g(mm')\in J$. Put $u_1=(mm')^{-1}u(mm')\in U$ and $z_1=(mm')^{-1}z(mm')\in Z$. By (1), $(mm')^{-1}\bar u(mm')\in J$, so $u_1z_1\in J$. Write $u_1z_1=u'z'\bar u'$ as in the statement; the hypothesis on products, applied to $(z'^{-1}u'^{-1}u_1z')(z'^{-1}z_1)\bar u'^{-1}=1$, gives $u_1=u'\in J$ and $z_1=z'\in J$. Then $m^{-1}um=m'u_1m'^{-1}\in J$, as $m'$ is positive; $m^{-1}zm\in J$, directly if $m''=m$ and as $m'z_1m'^{-1}$ if $m''=m'$; and $m^{-1}\bar um\in J$, as $m$ is positive. So $m^{-1}gm\in J$, that is, $g\in mJm^{-1}$.
\item \emph{Counting.} By (2), the orbit of $mm'J$ under $J_m$ is $m$ applied to the orbit $Jm'J/J$ of $m'J$ under $J$: for $b\in J$, $mbm'J=(mb'm^{-1})mm'J$ with $mb'm^{-1}\in J_m$, and conversely $jmm'J=m(m^{-1}jm)m'J$ with $m^{-1}jm\in J$ for $j\in J_m$. Since the stabilizer of $mm'J$ in $J_m$ is $J_{mm'}$ by (3), $[J_m:J_{mm'}]$ is the number of elements of $Jm'J/J$, and $[J:J_{mm'}]=[J:J_m][J_m:J_{mm'}]$ is the product of the numbers of elements of $JmJ/J$ and $Jm'J/J$.
\item \emph{The product.} Let $y$ and $y'$ run over representatives of $JmJ/J$ and $Jm'J/J$. For $v\in M^J$, $[JmJ][Jm'J]v=\sum_{y,y'}\rho(yy')v$, and $\rho(yy')v$ depends only on $yy'J$. Write $y=amj$ with $a,j\in J$ and $y'J=bm'J$ with $b\in J$; then $yy'J=am(jb)m'J$, which lies in $Jmm'J/J$ by (2). The map $(y,y')\mapsto yy'J$ to $Jmm'J/J$ is onto: for $a\in J$, write $am=yi$ with $i\in J$ and $im'=y'j$ with $j\in J$; then $amm'J=yy'J$. By (4) it is a bijection, so $[JmJ][Jm'J]v=[Jmm'J]v$.
\end{enumerate}
\end{proof}
% leviMul (IwLeviMul.lean): hP is (1), hB is (2), hA is (3) through the private split_mem, the private count is (4) (orbits and stabilizers), and hkMul (IwHkMul.lean) is (5) (hmem, hsurj and the cardinality).

\begin{lemma}\label{sha-automorphy-lift-forms-facts-span-local}
Let $K$, $\varpi$, $q$, $n$, $N_0$, $\bar N_c$, $T_0$ and the operators $U_\alpha$ be as in Definition~\ref{sha-hecke-local-ordinary}. Let $J$ be the subgroup of $GL_n(K)$ generated by $N_0$, $\bar N_0$ and $T_0$, and for $c\ge0$ let $J(c)$ be the subgroup generated by $N_0$, $\bar N_c$ and the diagonal matrices with entries in $1+\varpi^{c+1}\mathcal O_K$. Let $J'\in\{J,J(c)\}$, let $\psi_1,\dots,\psi_n$ be smooth characters of $K^\times$ that are trivial on the diagonal entries of the diagonal matrices in $J'$, and let $(\tau,T)$ be a representation of $GL_n(K)$ isomorphic to the normalized parabolic induction $I(\psi)$, from the upper triangular Borel subgroup $B$, of $\operatorname{diag}(t_1,\dots,t_n)\mapsto\prod_i\psi_i(t_i)$. For $g\in GL_n(K)$ let $[J'gJ']$ be the operator $v\mapsto\sum_y\tau(y)v$ on the vectors fixed by $\tau(J')$, the sum over representatives $y$ of $J'gJ'/J'$. Then there is a nonzero $v\in T$ fixed by $\tau(J')$ with
\[[J'\operatorname{diag}(d)J']\,v=\prod_i\psi_i(d_i)\,\delta_B(\operatorname{diag}(d))^{-1/2}\,v\]
for every $d\in(K^\times)^n$ with $|d_1|<\dots<|d_n|$, where $\delta_B$ is the modulus character of $B$.
\end{lemma}
\leanref{Development/AutomorphicArithmetic/AutomorphicVarieties/ModularityArgument/ThresholdArguments/DenseThresholdCases/TLocEig.lean}{UpperTheory.V.Z.Sym.Al.exists_fixd_dc_eq_expo_smul_of_equiv_prinSeries}
\begin{proof}
For $d\in(K^\times)^n$ let $\mu(d)=\prod_i\psi_i(d_i)\,\delta_B(\operatorname{diag}(d))^{-1/2}$, a character of $(K^\times)^n$, equal to $\prod_i\psi_i(d_i)$ when the $d_i$ are units. Let $c'=0$ if $J'=J$ and $c'=c$ if $J'=J(c)$.
\begin{enumerate}
\item \emph{An ordinary vector.} Let $\omega_i=\psi_i|\cdot|^{(1-n)/2}$, so that $\tau\cong\mathrm{PS}(\omega)$. The $\omega_i$ agree with the $\psi_i$ on units, so they are trivial on $1+\varpi^{c'+1}\mathcal O_K$. Lemma~\ref{sha-hecke-ideal-ordinary-vector}, for $c'$ and the character $\phi$ of $T_0$ with $\phi(e_i(u))=\omega_i(u)^{-1}$, gives a nonzero $v$ fixed by $N_0$ and $\bar N_{c'}$, on which $t\in T_0$ acts by $\prod_i\psi_i(t_i)$, with $U_{\alpha_j}v=E_j(\omega)v$ for $0\le j\le n$.
\item \emph{$v$ is fixed by $J'$.} The generators $N_0$ and $\bar N_{c'}$ of $J'$ fix $v$, and a diagonal matrix of $J'$ acts on $v$ by $\prod_i\psi_i(t_i)=1$.
\item \emph{The operators $U_\alpha$.} We have $E_j(\omega)=\prod_{i\le j}\psi_i(\varpi)q^{(n+1-2i)/2}=\mu(\alpha_j)$, as $\delta_B(\alpha_j)=q^{-j(n-j)}$ and $\sum_{i\le j}(n+1-2i)=j(n-j)$. For integers $a_1\ge\dots\ge a_n$ let $\varpi^a=\operatorname{diag}(\varpi^{a_1},\dots,\varpi^{a_n})$. We show $U_{\varpi^a}v=\mu(\varpi^a)v$. The $a$ for which this holds are closed under addition, by Lemma~\ref{sha-hecke-ideal-tw-augmented-uop}(1), and contain the $a$ of the $\alpha_j$, $1\le j\le n$. They contain $(-1,\dots,-1)$: $U_{\alpha_n^{-1}}U_{\alpha_n}v=U_1v=v$ by the same lemma, so $U_{\alpha_n^{-1}}v=\mu(\alpha_n)^{-1}v$. Every $a$ with $a_1\ge\dots\ge a_n$ is a sum of these, as $a=a_n(1,\dots,1)+\sum_{j<n}(a_j-a_{j+1})\,a(\alpha_j)$, where $a(\alpha_j)$ has its first $j$ entries $1$ and the others $0$.
\item \emph{Double cosets of dominant elements.} Let $Z'$ be the group of diagonal matrices in $J'$. Every element of $J'$ is $xz\bar y$ with $x\in N_0$, $z\in Z'$ and $\bar y\in\bar N_{c'}$, and this decomposition is unique (the Iwahori factorization). For $a_1\ge\dots\ge a_n$, $\varpi^{-a}z\bar y\varpi^a\in J'$, so $J'\varpi^aJ'=N_0\varpi^aJ'$; and for $x,x'\in N_0$, $x\varpi^aJ'=x'\varpi^aJ'$ if and only if $\varpi^{-a}x^{-1}x'\varpi^a$ lies in $J'$, hence in $N_0$ by the uniqueness, that is, if and only if $x^{-1}x'\in\varpi^aN_0\varpi^{-a}$. So $[J'\varpi^aJ']=U_{\varpi^a}$ on the vectors fixed by $\tau(J')$. For $t\in T_0$, which normalizes $J'$, $J'tJ'=tJ'$ and $[J'tJ']=\tau(t)$.
\item \emph{Conclusion.} Let $|d_1|<\dots<|d_n|$, and write $\operatorname{diag}(d)=t\varpi^a$ with $t\in T_0$ and $a_i=\operatorname{ord}(d_i)$, so that $a_1>\dots>a_n$. Lemma~\ref{sha-automorphy-lift-forms-facts-operators-product} applies to $G=GL_n(K)$, $J'$, the upper triangular unipotent, diagonal and lower triangular unipotent subgroups, and the positive elements $t$ and $\varpi^a$, which commute with the diagonal matrices: conjugation by them maps $N_0$ into $N_0$, and conjugation by their inverses maps $\bar N_{c'}$ into $\bar N_{c'}$. So $[J'\operatorname{diag}(d)J']v=[J'tJ'][J'\varpi^aJ']v=\tau(t)U_{\varpi^a}v=\mu(\varpi^a)\mu(t)v=\mu(d)v$, by (1), (3) and (4).
\end{enumerate}
\end{proof}
% exists_fixd_dc_eq_expo_smul_of_equiv_prinSeries (TLocEig.lean): hE (om_shift), Hi.exists_ne_zero_mem_eig_uOp_eq_psEig_smul_of_equiv_ps (Lemma sha-hecke-ideal-ordinary-vector) with phi and phi_inv for (1); hT and hvJ for (2); mu, expo_eq, mu_fund, uOp_upow with antitone_mem (the monoid generated by fund n j and -1) and uOpMul for (3); Sig.hkMul (hk_eq_uOp, sum_reps, and leviMul) and u_operator_torus_eigen for (4); the factorization d = u p for (5). The Lean states it for every level Sig.Lev (the subgroups J'); the book states the two it uses.

\begin{lemma}\label{sha-automorphy-lift-forms-facts-span-isotypic}
Let $L$ be a CM field, let $n\ge1$, use the notation of Definition~\ref{gal-definite-unitary} and Definition~\ref{sha-automorphy-lift-forms}, and let $\pi$ be an automorphic representation of $U_n$, $w$ a split place and $\tau$ a local component of $\pi$ at $w$. Say that a linear map from a representation of $GL_n(L_w)$ to a space of functions on $U_n$ \emph{intertwines} if it carries the action of $c_w(x)$ to right translation by $x$, for $x\in U_{n,w}$.
\begin{enumerate}
\item $\pi$ is the sum of the images of the intertwining maps from $\tau$ to $\pi$.
\item Let $N\ne0$ be a complex representation of $GL_n(L_w)$, and let $q_1,\dots,q_r:N\to\pi$ be linear maps compatible with the actions of $GL_n(L_w)$, through $c_w$ on $\pi$, such that $x=0$ whenever $q_i(x)=0$ for all $i$. Then there is an injective linear map from $\tau$ to $N$ compatible with the actions.
\item For every nonzero $x\in\pi$ there is a linear map $p:\pi\to\tau$, compatible with the actions, with $p(x)\ne0$.
\end{enumerate}
\end{lemma}
\leanref{Development/AutomorphicArithmetic/AutomorphicVarieties/ModularityArgument/SymmetricPowerTransfer/GeneralizedEigenspaces/SmoothDepth/OpenLevelSubgroups/IsoTypic.lean}{UpperTheory.V.Z.Sym.Al.isoTypic}
\leanref{Development/AutomorphicArithmetic/AutomorphicVarieties/ModularityArgument/SymmetricPowerTransfer/GeneralizedEigenspaces/SmoothDepth/OpenLevelSubgroups/IsoCopy.lean}{UpperTheory.V.Z.Sym.Al.isoCopy}
\begin{proof}
(1) Let $P\subseteq\pi$ be the sum of these images; it is nonzero, as $\tau$ is a local component. Let $g\in U_n$, and write $g=xy$ with $x\in U_{n,w}$ its component at the place of $L^+$ under $w$, and $y$ with component $1$ there, so that $y$ commutes with $U_{n,w}$. For an intertwining map $p$, right translation by $y$ composed with $p$ is again intertwining, and right translation by $x$ maps $p(t)$ to $p(\tau(c_w(x))t)$. So $P$ is stable under right translation by $U_n$, and $P=\pi$, as $\pi$ is irreducible.

(2) Consider representations of $GL_n(L_w)$ as modules over its group algebra. The image of an intertwining map from the irreducible $\tau$ is $0$ or isomorphic to $\tau$, so by (1) $\pi$ is a sum of simple modules isomorphic to $\tau$; hence it is semisimple, and so is $\pi^r$, whose simple submodules are all isomorphic to $\tau$. The map $x\mapsto(q_i(x))_i$ identifies $N$ with a nonzero submodule of $\pi^r$, which is semisimple and so contains a simple submodule isomorphic to $\tau$. Such a map from $\tau$ is injective.

(3) As in (2), $\pi$ is semisimple, a direct sum of simple modules isomorphic to $\tau$. A nonzero $x$ has a nonzero component in one of them, and the projection onto it, followed by an isomorphism with $\tau$, is the map $p$.
\end{proof}
% isoTypic (IsoTypic.lean): factor for g = xy, the maps repOn y o p, irreducibility of pi. isoCopy (IsoCopy.lean): the private isoCopy_struct (the sum is semisimple and isotypic), then linearEquiv_finsupp and the injective LinearMap.pi Q.

\begin{lemma}\label{sha-automorphy-lift-forms-facts-span}
Assume that Galois representations exist, in the sense of Definition~\ref{gal-galois-representations}, relative to the maps satisfying the laws of Definition~\ref{gal-llc-laws}.
Let $L$ be a CM field such that $L/L^+$ is unramified at every finite place, let $n\ge1$, and use the notation of Definition~\ref{sha-automorphy-lift-forms} and Lemma~\ref{sha-automorphy-lift-forms-facts-reps}. For a split place $w$ and a function $\phi$ on $U_n$ with values in a vector space, the \emph{span of $\phi$ at $w$} is the span of the right translates $g\mapsto\phi(gx)$, $x\in U_{n,w}$. Let $\pi$ be an automorphic representation of $U_n$.
\begin{enumerate}
\item[(a)] Let $n=2$. If $j\ge0$ and $\pi$ has at every infinite place a weight $\lambda$ with $\lambda_0-\lambda_1=j$, then for some $e$ there is a nonzero function $\phi$ on $U_2$ with values in $V_j\otimes{\det}^e$ and coefficients in $\pi$ such that $\phi(\gamma g)=\phi(g)$ for $\gamma\in U_2(L^+)$ and $\phi(gx)=x^{-1}\phi(g)$ for $x$ in the part at infinity.
\item[(b)] Let $w$ be a split place, $J'_w\in\{J_w,J_w(c)\}$ and $\psi$ smooth characters with $I(\psi)$ a local component of $\pi$ at $w$, trivial on the diagonal entries of $J'_w$. For every nonzero $\phi$ with values in a finite-dimensional space $V$ and coefficients in $\pi$, the span of $\phi$ at $w$ contains a nonzero $\phi'$ with $\phi'(gx)=\phi'(g)$ for $x\in U_{n,w}$ with $c_w(x)\in J'_w$ that is a $(J'_w,\psi)$-eigenform.
\item[(c)] Let $w$ be a split place at which $\pi$ is unramified. For every nonzero $\phi$ with values in a finite-dimensional space and coefficients in $\pi$, the span of $\phi$ at $w$ contains a nonzero $\phi'$ with $\phi'(gx)=\phi'(g)$ for $x\in U_{n,w}\cap K_0$.
\item[(d)] Let $w$ be a split place and $\phi'$ in the span of $\phi$ at $w$. If $\phi$ has coefficients in $\pi$, so has $\phi'$. If $\phi(\gamma g)=\phi(g)$ for $\gamma\in U_n(L^+)$, $\phi(gx)=x^{-1}\phi(g)$ for $x$ in the part at infinity, and $\phi(gu)=\phi(g)$ for $u$ in a subgroup of the elements whose components at the places over the place of $L^+$ under $w$ are $1$, then $\phi'$ has the same properties. If $w'$ is a split place over another place of $L^+$ and $T_d\phi=\mu\phi$ for an operator $T_d$ at $w'$, then $T_d\phi'=\mu\phi'$.
\end{enumerate}
\end{lemma}
\leanref{Development/AutomorphicArithmetic/AutomorphicVarieties/ModularityArgument/SymmetricPowerTransfer/GeneralizedEigenspaces/SmoothDepth/OpenLevelSubgroups/InfVec.lean}{UpperTheory.V.Z.Sym.Al.infVec}
\leanref{Development/AutomorphicArithmetic/AutomorphicVarieties/ModularityArgument/ThresholdArguments/DenseThresholdCases/TSpanEig.lean}{UpperTheory.V.Z.Sym.Al.exists_eigenvector_in_span_of_comp_eq_trivial}
\leanref{Development/AutomorphicArithmetic/AutomorphicVarieties/ModularityArgument/SymmetricPowerTransfer/GeneralizedEigenspaces/SmoothDepth/OpenLevelSubgroups/SpanSph.lean}{UpperTheory.V.Z.Sym.Al.spanSph}
\leanref{Development/AutomorphicArithmetic/AutomorphicVarieties/ModularityArgument/SymmetricPowerTransfer/GeneralizedEigenspaces/SmoothDepth/OpenLevelSubgroups/SpanKeep.lean}{UpperTheory.V.Z.Sym.Al.spanKeep}
\begin{proof}
(d) Each property is linear in $\phi$, so it suffices to check it for $\phi'=\phi(\,\cdot\,x)$ with $x\in U_{n,w}$. The representation $\pi$ is stable under right translation, which gives the first. Right translation by $x$ commutes with left translation by $U_n(L^+)$, and with right translation by the part at infinity and by the elements whose components over the place of $L^+$ under $w$ are $1$, which lie at other places; this gives the second. The operator $T_d$ at $w'$ is a sum of right translations by elements of $U_{n,w'}$, which commute with $x$; this gives the third.
% (d) is spanKeep (SpanKeep.lean, span_induction with hpi.stable, and the commutation of the translations at w with away n w and with the T operators at w').

(a) For each infinite place $v$ let $\lambda_v$ be a weight of $\pi$ at $v$ with $\lambda_{v,0}-\lambda_{v,1}=j$, and let $H_v\cong U(2)$ be the factor of the part at infinity at $v$, which acts on $\nabla(\lambda_v)$ through the embedding $\tau_v$ attached to $v$.
\begin{enumerate}
\item \emph{Isotypic parts.} For each $v$, $\pi$ is the sum of the images of the linear maps $\nabla(\lambda_v)\to\pi$ that carry the action of $H_v$ to right translation. This sum is nonzero, as $\pi$ has weight $\lambda_v$ at $v$; it is stable under right translation by $H_v$, and by the rest of $U_2$, which commutes with $H_v$; so it is $\pi$, which is irreducible.
\item \emph{A map from the tensor product.} As $\nabla(\lambda_v)$ is irreducible and finite-dimensional, the density theorem and (1) extend the action of $H_v$ on $\pi$ to an algebra map $\Phi_v$ from the endomorphisms of $\nabla(\lambda_v)$ to those of $\pi$, whose image commutes with every endomorphism commuting with $H_v$; so the images of $\Phi_v$ and $\Phi_{v'}$ commute for $v\ne v'$. Choose $w_v\ne0$ in $\nabla(\lambda_v)$ and a linear form $\ell_v$ with $\ell_v(w_v)=1$, and for $x\in\nabla(\lambda_v)$ write $x\ell_v$ for the endomorphism $y\mapsto\ell_v(y)x$. A nonzero subspace of $\pi$ stable under the image of $\Phi_v$ is a module over the simple algebra of endomorphisms of $\nabla(\lambda_v)$, so it is not killed by the idempotent $\Phi_v(w_v\ell_v)$; applying these commuting idempotents one place at a time gives $y\ne0$ in $\pi$ with $\Phi_v(w_v\ell_v)y=y$ for every $v$. The map $F(x)=\big(\prod_v\Phi_v(x_v\ell_v)\big)y$ is multilinear in $x=(x_v)_v$, has $F(w)=y\ne0$, and satisfies $F(\dots,hx_v,\dots)=hF(x)$ for $h\in H_v$. So it induces a nonzero linear map $T:\bigotimes_v\nabla(\lambda_v)\to\pi$ that carries the action of the part at infinity to right translation.
\item \emph{A polynomial model.} Realize $\nabla(\lambda_v)$, as Alpöge does, on the polynomials of degree $j$ in two variables $Y_0,Y_1$, with $\lambda_{v,1}$ as the power of the determinant, so that the vectors $D_v(a,b)=(aY_1-bY_0)^j$ satisfy $h\,D_v(a,b)=\det(h)^{\lambda_{v,1}+j}D_v\big((a,b)h^{-1}\big)$ for $h\in GL_2(\mathbf C)$, with $(a,b)h^{-1}$ the product of a row vector and a matrix. For $x=(x_{v,0},x_{v,1})_v$ put $c(x)=\bigotimes_vD_v(x_{v,0},x_{v,1})$. Expanding the binomials, $c(x)=\sum_\kappa a_\kappa(x)c_\kappa$, where $\kappa$ runs over the families $(\kappa_v)$ with $0\le\kappa_v\le j$, $c_\kappa=\bigotimes_vY_1^{\kappa_v}Y_0^{j-\kappa_v}$, and $a_\kappa=\prod_v\binom{j}{\kappa_v}x_{v,0}^{\kappa_v}(-x_{v,1})^{j-\kappa_v}$ lies in $V_j\otimes{\det}^e$.
\item \emph{The function.} Let $e(v)=-(\lambda_{v,1}+j)$ and $\phi(g)=\sum_\kappa T(c_\kappa)(g)\,a_\kappa$, so that $\phi(g)$ is the polynomial $x\mapsto T(c(x))(g)$. For a linear form $f$ on $V_j\otimes{\det}^e$, $f\circ\phi=\sum_\kappa f(a_\kappa)T(c_\kappa)\in\pi$, so $\phi$ has coefficients in $\pi$, and $\phi(\gamma g)=\phi(g)$ for $\gamma\in U_2(L^+)$, as the functions in $\pi$ have this property. For $u$ in the part at infinity, $T(c(x))(gu)=T(u\,c(x))(g)=\prod_v\det(u_v)^{\lambda_{v,1}+j}\,T(c(xu^{-1}))(g)$ by (3), where $xu^{-1}$ multiplies each pair by $u_v^{-1}$; this says $\phi(gu)=u^{-1}\phi(g)$ in $V_j\otimes{\det}^e$.
\item \emph{$\phi\ne0$.} If $\phi=0$, then $T(c(x))=0$ for all $x$. By (3), the span of the $c(x)$ is stable under the part at infinity, and it is nonzero, as $c(x)\ne0$ for $x_{v,0}=1$ and $x_{v,1}=0$. As $\bigotimes_v\nabla(\lambda_v)$ is irreducible under the part at infinity, the $c(x)$ span it, and $T=0$, a contradiction.
\end{enumerate}
% infVec (InfVec.lean): (1) and (2) are GlobalTheory.T.exists_ne_zero_weightTensor_intertwining_autRep (exists_ne_zero_wsp_intertwining: iSup_range_weightMaps_eq_top_definiteUnitary, then exists_multilinearMap_of_isotypic_commuting with exists_algHom_endomorphism_of_span_intertwiners, exists_ne_zero_fixed_by_commuting_rankOne_images and exists_multilinearMap_of_commuting_endAlgHom); (3) is the private D, sq_D, cv, cv_eq, A and A_mem; (4) is key and wInf_cv with Wt.eval_mRep; (5) is cv_ne with isIrreducible_wInfU.

(b) Let $N$ be the span of $\phi$ at $w$, a representation of $GL_n(L_w)$ by right translation through $c_w^{-1}$; it contains $\phi\ne0$. By (d), its elements have coefficients in $\pi$. For a basis of $V$, the coordinate maps $q_i:N\to\pi$ are compatible with the actions, and an element of $N$ whose coordinates all vanish is $0$. By Lemma~\ref{sha-automorphy-lift-forms-facts-span-isotypic}(2), applied to the local component $\tau\cong I(\psi)$, there is an injective map $f:\tau\to N$ compatible with the actions. Lemma~\ref{sha-automorphy-lift-forms-facts-span-local}, for $K=L_w$ and $J'=J'_w$, gives $v\ne0$ in $\tau$, fixed by $J'_w$, with $[J'_w\operatorname{diag}(d)J'_w]v=\prod_i\psi_i(d_i)\delta_B(\operatorname{diag}(d))^{-1/2}v$ whenever $|d_1|<\dots<|d_n|$. Let $\phi'=f(v)\ne0$, in $N$. It is invariant under right translation by the $x\in U_{n,w}$ with $c_w(x)\in J'_w$, and, as $f$ is compatible with the actions and $T_d$ is a sum of right translations, $T_d\phi'=f([J'_w\operatorname{diag}(d)J'_w]v)$. So $\phi'$ is a $(J'_w,\psi)$-eigenform.
% exists_eigenvector_in_span_of_comp_eq_trivial (TSpanEig.lean): N, rho and hst; hco (spanKeep); the coordinates q with Module.finBasis; isoCopy with isoTypic; hinj; exists_fixd_dc_eq_expo_smul_of_equiv_prinSeries; dcMap for the operators.

(c) Let $K_w=U_{n,w}\cap K_0$, the group of the $x\in U_{n,w}$ with $c_w(x)\in GL_n(\mathcal O_w)$, and let $\varphi_0\ne0$ in $\pi$ be fixed by right translation by $K_w$, which exists as $\pi$ is unramified at $w$. Choose a linear form $f$ on the space of values with $f\circ\phi\ne0$. Then $f\circ\phi\in\pi$, and as $\pi$ is irreducible, the right translates of $f\circ\phi$ span $\pi$: $\varphi_0=\sum_gc_g\,(f\circ\phi)(\,\cdot\,g)$, a finite sum. Write each $g$ with $c_g\ne0$ as $x_gy_g$, with $x_g\in U_{n,w}$ and $y_g$ with component $1$ at the place of $L^+$ under $w$, so that $y_g$ commutes with $U_{n,w}$. The function $\phi_g(t)=\phi(tx_g)$ has coefficients in $\pi$, so it is invariant under right translation by an open subgroup $C_g$ of the finite part, as its coordinates are; let $C=\bigcap_gC_g$. The group $K_w$ is compact, so $K_w\cap C$ has finite index in it; let $r$ run over representatives of $K_w/(K_w\cap C)$, and put $A_g(t)=\sum_r\phi_g(tr)$. For every $t$,
\[[K_w:K_w\cap C]\,\varphi_0(t)=\sum_r\varphi_0(tr)=\sum_gc_g\sum_rf(\phi(trx_gy_g))=\sum_gc_g\,f\big(A_g(ty_g)\big),\]
as $rx_gy_g=y_grx_g$. Since $\varphi_0\ne0$, some $A_g$ is nonzero. It lies in the span of $\phi$ at $w$, being a sum of right translates of $\phi$ by the elements $rx_g$ of $U_{n,w}$. It is invariant under right translation by $K_w$: for $k\in K_w$, $kr$ runs over representatives of the same cosets, and $\phi_g$ is invariant under $K_w\cap C$. So $\phi'=A_g$ has the required properties.
% spanSph (SpanSph.lean): the private finite_quot (K_w compact, through intPts), openLev for the open subgroups U g, factor for g = x y, key, and sum_translate for the invariance.
\end{proof}

\begin{lemma}\label{sha-automorphy-lift-forms-facts-operators}
Let $L$ be a CM field, use the notation of Definition~\ref{sha-automorphy-lift-forms} with $n=2$, let $V$ be a finite-dimensional representation of $\prod_vGL_2(\mathbf C)$ with continuous matrix coefficients, let $K$ be a level, and let $S$ be a finite set of places split over $L^+$, no two over the same place of $L^+$, at which $K$ is maximal. Let $F$ be the space of forms of weight $V$ and level $K(S)$, with $J'_w=J_w$ for $w\in S$. Then $F$ is finite-dimensional, and for $w\in S$ and $d,d'\in(L_w^\times)^2$ with $|d_1|\le|d_2|$ and $|d'_1|\le|d'_2|$: $T_d$ maps $F$ into itself; $T_{dd'}=T_dT_{d'}$ on $F$; $T_d=1$ on $F$ if $d_1$ and $d_2$ are units; and $T_d$ commutes on $F$ with the operators $T_{d''}$, for $d''$ with $|d''_1|\le|d''_2|$, at the other places of $S$.
\end{lemma}
\leanref{Development/AutomorphicArithmetic/AutomorphicVarieties/ModularityArgument/SymmetricPowerTransfer/GeneralizedEigenspaces/SmoothDepth/OpenLevelSubgroups/AfFin.lean}{UpperTheory.V.Z.Sym.Al.afFin}
\leanref{Development/AutomorphicArithmetic/AutomorphicVarieties/ModularityArgument/SymmetricPowerTransfer/GeneralizedEigenspaces/SmoothDepth/OpenLevelSubgroups/SecMemLev.lean}{UpperTheory.V.Z.Sym.Al.afOps}
\begin{proof}
\begin{enumerate}
\item \emph{Finiteness.} The level $K(S)$ is an open subgroup of the finite part of $U_2$, so by Lemma~\ref{sha-hecke-ideal-tw-augmented-hida-classes} there is a finite set $\Sigma\subset U_2$ such that every $g\in U_2$ is $\gamma sxk$ with $\gamma\in U_2(L^+)$, $s\in\Sigma$, $x$ in the part at infinity and $k\in K(S)$. For $\phi\in F$, $\phi(\gamma sxk)=\phi(sx)=x^{-1}\phi(s)$. So $\phi\mapsto(\phi(s))_{s\in\Sigma}$ is an injective linear map from $F$ to $V^\Sigma$, and $F$ is finite-dimensional.
\item \emph{The setting at $w$.} Let $w\in S$, let $v$ be the place of $L^+$ under $w$, so that $U_{2,w}$ is the part of $U_2$ at $v$, and let $\rho(g)$, for $g\in GL_2(L_w)$, be right translation by $c_w^{-1}(g)$ on functions on $U_2$ with values in $V$. An element $x\in U_{2,w}$ with $c_w(x)\in J_w$ lies in $K(S)$: it lies in $K_0$, as $J_w\subset GL_2(\mathcal O_w)$, hence in $K$, which is maximal at $w$, and its components at the other places of $S$, which lie over other places of $L^+$, are $1$. So every $\phi\in F$ is fixed by $\rho(J_w)$, and on such functions $T_d=[J_w\operatorname{diag}(d)J_w]$ in the notation of Lemma~\ref{sha-automorphy-lift-forms-facts-operators-product}; the sets $J_w\operatorname{diag}(d)J_w/J_w$ are finite, as $J_w$ is compact and open.
\item \emph{$T_d$ maps $F$ into itself.} Let $\phi\in F$. Then $T_d\phi$ is a finite sum of right translates of $\phi$ by elements of $U_{2,w}$. Each of them is invariant on the left under $U_2(L^+)$, equivariant under the part at infinity, and invariant under right translation by the elements of $K(S)$ whose component at $v$ is $1$, since these commute with $U_{2,w}$. Moreover $T_d\phi$ is fixed by $\rho(J_w)$, as left multiplication by $J_w$ permutes $J_w\operatorname{diag}(d)J_w/J_w$. Every $k\in K(S)$ is $xy$ with $x\in U_{2,w}$ its component at $v$ and $y$ with component $1$ at $v$; then $c_w(x)\in J_w$, so $x\in K(S)$ by (2), and $y=x^{-1}k\in K(S)$. Hence $T_d\phi(gk)=T_d\phi(gxy)=T_d\phi(gx)=T_d\phi(g)$, and $T_d\phi\in F$.
\item \emph{Products.} Apply Lemma~\ref{sha-automorphy-lift-forms-facts-operators-product} to $G=GL_2(L_w)$, $J=J_w$, the groups $U$, $Z$ and $\bar U$ of upper triangular unipotent, diagonal and lower triangular unipotent matrices, and $\rho$. Every element of $J_w$ is $uz\bar u$ with $u\in U$ integral, $z$ diagonal with unit entries and $\bar u\in\bar U$ with entry in $\mathfrak m_w$ (the Iwahori factorization); and $uz\bar u=1$ forces $\bar u=(uz)^{-1}$ to be upper triangular, so $\bar u=1$ and $u=z=1$. For $m=\operatorname{diag}(d)$ with $|d_1|\le|d_2|$, $mum^{-1}$ has entry $d_1u_{12}d_2^{-1}$ and $m^{-1}\bar um$ has entry $d_2^{-1}\bar u_{21}d_1$, of absolute values at most $|u_{12}|$ and $|\bar u_{21}|$; so $m$ is positive, and it commutes with $Z$. So $T_dT_{d'}=T_{dd'}$ on $F$.
\item \emph{Units.} If $d_1$ and $d_2$ are units, $\operatorname{diag}(d)\in J_w$, so $J_w\operatorname{diag}(d)J_w/J_w$ has the single element $J_w$, and $T_d\phi=\phi$ for $\phi\in F$.
\item \emph{Commutation.} For $w'\in S$ with $w'\ne w$, $w'$ lies over a place of $L^+$ other than $v$. So $T_{d''}$ at $w'$ is a sum of right translations by elements of $U_2$ at that place, which commute with the right translations by elements of $U_{2,w}$ that make up $T_d$.
\end{enumerate}
% afFin (AfFin.lean): GlobalTheory.T.definiteUnitary_classSet_finite and the injective map T to S -> V. afOps (SecMemLev.lean): (2) is the private sec_mem_lev and fix_of_AF, with hk_val for the operator; (3) is op_keep (spanKeep for the translates, factor for k = xy); (4) is Sig.hkMul (HkEqUOp.lean: leviMul with the private lf, pos and normJ, Lemma sha-automorphy-lift-forms-facts-operators-product); (5) is op_unit (hkApply); (6) is op_comm (factor). The Lean states afOps for every n.
\end{proof}

\begin{lemma}\label{sha-automorphy-lift-forms-facts}
Assume that Galois representations exist, in the sense of Definition~\ref{gal-galois-representations}, relative to the maps satisfying the laws of Definition~\ref{gal-llc-laws}.
Let $L$ be a CM field such that $L/L^+$ is unramified at every finite place, and use the notation of Definition~\ref{sha-automorphy-lift-forms}. Call characters $\psi_1,\dots,\psi_n$ of $L_w^\times$ \emph{separated} if they are unramified and $\psi_i\ne\psi_{i'}$ and $\psi_i\ne\psi_{i'}|\cdot|$ for all $i\ne i'$. Let $n\ge1$ and let $V$ be a finite-dimensional representation of $\prod_vGL_n(\mathbf C)$ with continuous matrix coefficients.
\begin{enumerate}
\item For every level $K$ there is a finite set $A$ of automorphic representations of $U_n$ with the following property. Let $S$ be a finite set of places split over $L^+$, no two over the same place of $L^+$, at which $K$ is maximal, let $\psi_w$ be separated for $w\in S$, and let $\phi\ne0$ be a form of weight $V$ and level $K(S)$ that is a $(J_w,\psi_w)$-eigenform at every $w\in S$. Then some $\sigma\in A$ has, at every $w\in S$, a local component that is a subquotient of $I(\psi_w)$.
\item For a finite set $S_0$ of places split over $L^+$, no two over the same place of $L^+$, and $c\ge0$, there is a finite set $A$ of automorphic representations of $U_n$, each with a nonzero vector fixed by $K_0(S_0,c)$, with the property of (1) for $K=K_0(S_0,c)$, the sets $S$ disjoint from $S_0$ with no two places of $S_0\cup S$ over the same place of $L^+$, and forms that are in addition $(J_w(c),\psi_w)$-eigenforms for smooth $\psi_w$ at every $w\in S_0$: some $\sigma\in A$ has, at every $w\in S_0\cup S$, a local component that is a subquotient of $I(\psi_w)$.
\item Let $\pi$ be an automorphic representation of $U_2$ and $j\ge0$ such that at every infinite place $\pi$ has a weight $\lambda$ with $\lambda_0-\lambda_1=j$. Let $S$ be a finite set of places split over $L^+$, no two over the same place of $L^+$, such that $\pi$ is unramified at every finite place $w$ over whose place of $L^+$ no place of $S$ lies. For $w\in S$ let $J'_w\in\{J_w,J_w(c_w)\}$, and let $\psi_w$ be characters with $I(\psi_w)$ a local component of $\pi$ at $w$ that are trivial on the diagonal entries of $J'_w$. Then for some $e$ there is a nonzero form of weight $V_j\otimes{\det}^e$ and level $K_0(S)$ that is a $(J'_w,\psi_w)$-eigenform at every $w\in S$.
\item Let $n=2$, let $K$ be a level, let $S$ be a finite set of places split over $L^+$, no two over the same place of $L^+$, at which $K$ is maximal, and let $\phi\ne0$ be a form of weight $V$ and level $K(S)$ that is a $(J_w,b_w)$-eigenform for the $U_t$ at every $w\in S$. Then there are a nonzero form $\phi'$ of weight $V$ and level $K(S)$ and unramified characters $\psi_{w,0},\psi_{w,1}$ with $\psi_{w,0}/\psi_{w,1}=b_w$ such that $\phi'$ is a $(J_w,\psi_w)$-eigenform at every $w\in S$.
\item For $n=2$, a $(J'_w,\psi)$-eigenform is a $(J'_w,\psi_0/\psi_1)$-eigenform for the $U_t$.
\end{enumerate}
\end{lemma}
\leanref{Development/AutomorphicArithmetic/AutomorphicVarieties/ModularityArgument/ThresholdArguments/DenseThresholdCases/TOutRep.lean}{UpperTheory.V.Z.Sym.Al.outRep_alt}
\leanref{Development/AutomorphicArithmetic/AutomorphicVarieties/ModularityArgument/ThresholdArguments/DenseThresholdCases/TInVec.lean}{UpperTheory.V.Z.Sym.Al.inVec_alt}
\leanref{Development/AutomorphicArithmetic/AutomorphicVarieties/ModularityArgument/SymmetricPowerTransfer/GeneralizedEigenspaces/SmoothDepth/OpenLevelSubgroups/Refine.lean}{UpperTheory.V.Z.Sym.Al.refine}
\leanref{Development/AutomorphicArithmetic/AutomorphicVarieties/ModularityArgument/SymmetricPowerTransfer/GeneralizedEigenspaces/SmoothDepth/OpenLevelSubgroups/EigTS.lean}{UpperTheory.V.Z.Sym.Al.eigTS}
\begin{proof}
(5) For $d=(t,t^{-1})$, $\psi_0(t)\psi_1(t^{-1})=(\psi_0/\psi_1)(t)$.

(1) Let $A$ be the finite set of Lemma~\ref{sha-automorphy-lift-forms-facts-reps}(a) for $K$, and let $S$, $\psi_w$ and $\phi$ be as in (1). The group $K(S)$ is a level. By Lemma~\ref{sha-automorphy-lift-forms-facts-reps}(b) there are an automorphic representation $\pi$ and a nonzero form $\phi_1$ of weight $V$ and level $K(S)$ with coefficients in $\pi$ that is a $(J_w,\psi_w)$-eigenform at every $w\in S$. Let $w\in S$. An element $x\in U_{n,w}$ with $c_w(x)\in J_w$ lies in $K(S)$: $J_w\subset GL_n(\mathcal O_w)$, so $x\in K_0\cap U_{n,w}\subset K$, as $K$ is maximal at $w$, and its components at the other places of $S$, which lie over other places of $L^+$, are $1$. So $\phi_1(gx)=\phi_1(g)$, and, as the $\psi_{w,i}$ are unramified, hence smooth, Lemma~\ref{sha-automorphy-lift-forms-facts-reps}(c) gives a local component of $\pi$ at $w$ that is a subquotient of $I(\psi_w)$. By Lemma~\ref{sha-automorphy-lift-forms-facts-reps}(d), $\pi$ is unramified at the places of $S$, and by Lemma~\ref{sha-automorphy-lift-forms-facts-reps}(e) it has a nonzero form of weight $V$ and level $K$ with coefficients in $\pi$. So some $\sigma\in A$ has every local component of $\pi$ as a local component, in particular, at every $w\in S$, one that is a subquotient of $I(\psi_w)$.

(2) Let $K=K_0(S_0,c)$, a level, and let $A$ be the finite set of Lemma~\ref{sha-automorphy-lift-forms-facts-reps}(a) for $K$; its elements have nonzero vectors fixed by $K$. Let $S$, $\psi_w$ and $\phi$ be as in (2). As in (1), Lemma~\ref{sha-automorphy-lift-forms-facts-reps}(b) gives $\pi$ and $\phi_1$, of level $K(S)$ and with the eigenform properties of $\phi$ at the places of $S_0\cup S$, and at every $w\in S$ a local component of $\pi$ that is a subquotient of $I(\psi_w)$. Let $w\in S_0$. An element $x\in U_{n,w}$ with $c_w(x)\in J_w(c)$ lies in $K(S)$: it lies in $K_0$, as $J_w(c)\subset GL_n(\mathcal O_w)$, its component at $w$ lies in $J_w(c)$, and its components at the other places of $S_0\cup S$, which lie over other places of $L^+$, are $1$. So Lemma~\ref{sha-automorphy-lift-forms-facts-reps}(c), with $J'_w=J_w(c)$ and the smooth $\psi_w$, gives a local component of $\pi$ at $w$ that is a subquotient of $I(\psi_w)$. As in (1), $\pi$ is unramified at the places of $S$ and has a nonzero form of weight $V$ and level $K$ with coefficients in $\pi$, and some $\sigma\in A$ has the local components of $\pi$ at every $w\in S_0\cup S$.

(3) \emph{A first vector.} Lemma~\ref{sha-automorphy-lift-forms-facts-span}(a) gives $e$ and a nonzero $\phi_0$ with values in $V_j\otimes{\det}^e$ and coefficients in $\pi$, invariant on the left under $U_2(L^+)$ and equivariant under the part at infinity. Its values are finite sums of functions in $\pi$, which are invariant under open subgroups, so $\phi_0$ is invariant under right translation by an open subgroup of the finite part of $U_2$; this subgroup contains the elements of $K_0$ whose components at the places of a finite set $T$ of places of $L$ are $1$.
\emph{The places.} Choose a finite set $T'\supseteq S$ of split places, no two over the same place of $L^+$, such that each place of $T$ lies over the same place of $L^+$ as a place of $T'$ and the places of $T'\setminus S$ lie over places of $L^+$ different from those under $S$: add to $S$, one at a time, a place of $T$ over a place of $L^+$ under no place chosen so far. For $D\subseteq T'$ let $K_D$ be the group of the $k\in K_0(S)$ whose components at the places over the places of $L^+$ under $T'\setminus D$ are $1$.
\emph{The induction.} We show, by induction on $D$, that there is a nonzero $\phi_D$ with coefficients in $\pi$, invariant on the left under $U_2(L^+)$, equivariant under the part at infinity, invariant under right translation by $K_D$, and a $(J'_w,\psi_w)$-eigenform at the places $w\in S\cap D$. For $D=\emptyset$, $\phi_\emptyset=\phi_0$ is invariant under $K_\emptyset$, which lies in $K_0$ and has components $1$ over the places under $T'$, hence over those under $T$. Let $D=D_0\cup\{w\}$ with $w\notin D_0$. If $w\in S$, Lemma~\ref{sha-automorphy-lift-forms-facts-span}(b), applied to $\phi_{D_0}$, gives $\phi_D$ in its span at $w$, invariant under the $x\in U_{2,w}$ with $c_w(x)\in J'_w$ and a $(J'_w,\psi_w)$-eigenform. If $w\notin S$, no place of $S$ lies over the place of $L^+$ under $w$, so $\pi$ is unramified at $w$, and Lemma~\ref{sha-automorphy-lift-forms-facts-span}(c) gives $\phi_D$ in the span of $\phi_{D_0}$ at $w$, invariant under $U_{2,w}\cap K_0$. By Lemma~\ref{sha-automorphy-lift-forms-facts-span}(d), $\phi_D$ has coefficients in $\pi$, keeps the properties at infinity and on the left, is invariant under the elements of $K_{D_0}$ whose components over the place of $L^+$ under $w$ are $1$, and keeps the eigenform properties at the places of $S\cap D_0$, which lie over other places of $L^+$. Every $u\in K_D$ is $xy$ with $x\in U_{2,w}$, its component at $w$, and $y$ with component $1$ at $w$; then $x\in K_0$, $c_w(x)\in J'_w$ if $w\in S$, and $y\in K_{D_0}$, as its components at the places of $S$ other than $w$ are those of $u$. So $\phi_D$ is invariant under $K_D$.
\emph{Conclusion.} For $D=T'$, $K_{T'}=K_0(S)$, and $\phi_{T'}$ is a nonzero form of weight $V_j\otimes{\det}^e$ and level $K_0(S)$ that is a $(J'_w,\psi_w)$-eigenform at every $w\in S$.

(4) Let $F$ be the finite-dimensional space of forms of weight $V$ and level $K(S)$ (Lemma~\ref{sha-automorphy-lift-forms-facts-operators}), and call $d\in(L_w^\times)^2$ \emph{dominant} if $|d_1|\le|d_2|$. By that lemma, the operators $T_d$ on $F$, for $w\in S$ and dominant $d$, commute with each other: at the same place $T_dT_{d'}=T_{dd'}=T_{d'}T_d$. Let $E\subset F$ be the space of the $\phi''$ with $T_{(t,t^{-1})}\phi''=b_w(t)\delta_B(\operatorname{diag}(t,t^{-1}))^{-1/2}\phi''$ for $w\in S$ and $|t|<1$. It contains $\phi\ne0$, and every $T_d$ maps it into itself, as $T_d$ commutes with the $T_{(t,t^{-1})}$. Commuting endomorphisms of a nonzero finite-dimensional complex vector space have a common eigenvector; let $\phi'\in E$ be one, with $T_d\phi'=\chi_w(d)\phi'$ for $w\in S$ and dominant $d$. Fix $w\in S$ and write $\chi=\chi_w$.
\begin{enumerate}
\item[(i)] By Lemma~\ref{sha-automorphy-lift-forms-facts-operators}, $\chi(dd')=\chi(d)\chi(d')$ for dominant $d,d'$, and $\chi(d)=1$ if $d_1,d_2$ are units.
\item[(ii)] $\chi(d)\ne0$ if $|d_1|<|d_2|$. Indeed, let $z=(d_1d_2,d_1d_2)$ and $t=d_1/d_2$, so that $|t|<1$ and $d\cdot d=z\cdot(t,t^{-1})$. As $z$ and $z^{-1}$ are dominant, $\chi(z)\chi(z^{-1})=\chi(1)=1$, so $\chi(z)\ne0$, and $\chi(d)^2=\chi(z)\chi(t,t^{-1})=\chi(z)b_w(t)\delta_B(\operatorname{diag}(t,t^{-1}))^{-1/2}\ne0$.
\item[(iii)] Every element of $(L_w^\times)^2$ is a quotient of two elements $d$ with $|d_1|<|d_2|$, on which $\chi$ is multiplicative, nonzero and unchanged under multiplication by units. So $\chi$ extends to a character of $(L_w^\times)^2$ trivial on $(\mathcal O_w^\times)^2$, and there are unramified characters $\psi_{w,0},\psi_{w,1}$ of $L_w^\times$ with $\chi(d)=\psi_{w,0}(d_1)\psi_{w,1}(d_2)\delta_B(\operatorname{diag}(d))^{-1/2}$ whenever $|d_1|<|d_2|$.
\item[(iv)] For $|t|<1$, (iii) and the definition of $E$ give $(\psi_{w,0}/\psi_{w,1})(t)=b_w(t)$; every element of $L_w^\times$ is a quotient of two such $t$, so $\psi_{w,0}/\psi_{w,1}=b_w$.
\end{enumerate}
So $\phi'$ is a nonzero form of weight $V$ and level $K(S)$ that is a $(J_w,\psi_w)$-eigenform at every $w\in S$.
% (1) is Al.outRep_alt: U.finFam (finFamU, through autProj), proj, comp_sub_of_eigen_vector_fixed_by_j (Al.subquotient_prin_series_of_iwahori_eigenvector, isoCopy, isoTypic, dcMap), unramified_of_split_distinct_characters and raise. (2) is the private outW of MemLevelOfLocalDepth.lean, the same with the private compD (W.local_sub_d_statement) at S_0 and mem_level_of_local_depth. (3) is Al.inVec_alt: infVec, openLev, cof, thinOut and the induction over the places with the private steps of TInVec.lean. (4) is Al.refine: afOps, afFin and the common eigenvector of the commuting operators u w d. The files are SymmetricPowerTransfer/GeneralizedEigenspaces/SmoothDepth and ProjectionComponents (about 3,200 lines) and TOutRep, TInVec, TCompOfEig, TLocSubD, TCompEx and UnramifiedOfSplitDistinctCharacters (about 760 lines).
\end{proof}

\begin{lemma}[Alpöge's adjoint lift]\label{sha-automorphy-adjoint-lift}
Assume that Galois representations exist, in the sense of Definition~\ref{gal-galois-representations}, relative to the maps satisfying the laws of Definition~\ref{gal-llc-laws}.
Let $L$ be a CM field such that $L/L^+$ is unramified at every finite place. Let $S_0$ be a finite set and $P$ a set of finite places of $L$, disjoint, all split over $L^+$, with no two places of $S_0\cup P$ over the same place of $L^+$, and such that every place of $L$ dividing $2$ lies over the same place of $L^+$ as a place of $S_0$. Let $\pi$ be an automorphic representation of $U_2$ (Definition~\ref{gal-definite-unitary}) and $j\ge1$ such that at every infinite place $\pi$ has a weight $\lambda$ with $\lambda_0-\lambda_1=j$, and $\pi$ is unramified at every finite place $w$ such that no place of $S_0$ lies over the same place of $L^+$ as $w$. For $w\in S_0\cup P$ let $\psi_{w,0},\psi_{w,1}$ be smooth characters of $L_w^\times$ with $I(\psi_{w,0},\psi_{w,1})$ a local component of $\pi$ at $w$, and let $a_w=\psi_{w,0}/\psi_{w,1}$. Suppose that for $w\in P$ the $\psi_{w,i}$ are unramified and $a_w^2\ne|\cdot|^k$ for every integer $k$ with $|k|\le2$. Then one of the following holds:
\begin{enumerate}
\item there is an automorphic representation $\sigma$ of $U_3$, unramified at every finite place $w$ such that no place of $S_0$ lies over the same place of $L^+$ as $w$, which at every $w\in S_0$, and at all but finitely many $w\in P$, has a local component that is a subquotient of $I(a_w,1,a_w^{-1})$;
\item there is an automorphic representation $\sigma'$ of $U_2$ which, at all but finitely many $w\in P$, has a local component that is a subquotient of $I(\psi')$ for characters $\psi'_0,\psi'_1$ with $\psi'_0/\psi'_1=a_w^2$.
\end{enumerate}
\end{lemma}
\leanref{Development/AutomorphicArithmetic/AutomorphicVarieties/ModularityArgument/SymmetricPowerTransfer/WeilDeligneUnramified/UnitaryTwoForm.lean}{UpperTheory.V.Z.Sym.W.AdLiftW}
\leanref{Development/AutomorphicArithmetic/AutomorphicVarieties/ModularityArgument/ThresholdArguments/DenseThresholdCases/MemLevelOfLocalDepth.lean}{UpperTheory.V.Z.Sym.W.adjoint_lift_w_of_literal_automorphy_spec_out}
\begin{proof}
We use the notation of Definition~\ref{sha-automorphy-lift-forms}. Since $a_w$ is unramified and $a_w^2\ne|\cdot|^k$ for $|k|\le2$, for $w\in P$ the characters $(a_w,1,a_w^{-1})$ are separated in the sense of Lemma~\ref{sha-automorphy-lift-forms-facts}, and so is every pair $(\psi'_0,\psi'_1)$ of unramified characters with $\psi'_0/\psi'_1=a_w^2$.
\begin{enumerate}
\item \emph{The level.} Let $c_0$ be given by Proposition~\ref{sha-automorphy-lift-comparison} for $S_0$ and $j$. Each $\psi_{w,i}$, $w\in S_0$, is smooth, so it is trivial on $1+\varpi_w^{c+1}\mathcal O_w$ for all large $c$. Fix $c\ge c_0$ with this property for all $w\in S_0$ and $i$, and let $K_2$ be the level of $U_2$ given by the proposition.
\item \emph{Finitely many candidates.} Let $A_3$ be the set of the representations of Lemma~\ref{sha-automorphy-lift-forms-facts}(2), for $n=3$, the weight $W_j$, $S_0$ and $c$, that have at every $w\in S_0$ a local component that is a subquotient of $I(a_w,1,a_w^{-1})$; and let $A_2$ be the set of Lemma~\ref{sha-automorphy-lift-forms-facts}(1) for $n=2$, the weight $V_{2j+1}$ and the level $K_2$. Both are finite. Let $P'$ be the set of $w\in P$ at which both $K_2$ and $K_0(S_0,c)\subset U_3$ are maximal; it contains all but finitely many places of $P$.
\item \emph{Finite subsets of $P'$.} Let $S_1\subset P'$ be finite. Then some $\sigma\in A_3$ has at every $w\in S_1$ a local component that is a subquotient of $I(a_w,1,a_w^{-1})$, or some $\sigma'\in A_2$ has at every $w\in S_1$ a local component that is a subquotient of $I(\psi')$ for characters with $\psi'_0/\psi'_1=a_w^2$.
Indeed, $S_1$ is disjoint from $S_0$, and no two places of $S_0\cup S_1$ lie over the same place of $L^+$. Lemma~\ref{sha-automorphy-lift-forms-facts}(3), applied to $\pi$, $S_0\cup S_1$, the groups $J_w(c)$ for $w\in S_0$ and $J_w$ for $w\in S_1$, and the characters $\psi_w$, gives $e$ and a nonzero form $\phi$ on $U_2$ of weight $V_j\otimes{\det}^e$ and level $K_0(S_0,c)(S_1)$, which by part (5) of that lemma is a $(J_w(c),a_w)$-eigenform for the $U_t$ at $w\in S_0$ and a $(J_w,a_w)$-eigenform for the $U_t$ at $w\in S_1$. Proposition~\ref{sha-automorphy-lift-comparison} applies.
In its case (1), Lemma~\ref{sha-automorphy-lift-forms-facts}(2) gives a representation of its set $A$ with local components that are subquotients of $I(a_w,1,a_w^{-1})$ at every $w\in S_0\cup S_1$; it lies in $A_3$.
In its case (2), Lemma~\ref{sha-automorphy-lift-forms-facts}(4) gives a nonzero form of weight $V_{2j+1}$ and level $K_2(S_1)$ that is a $(J_w,\psi'_w)$-eigenform at every $w\in S_1$, for unramified $\psi'_w$ with $\psi'_{w,0}/\psi'_{w,1}=a_w^2$; and Lemma~\ref{sha-automorphy-lift-forms-facts}(1) gives $\sigma'\in A_2$.
\item \emph{All of $P'$.} Some element of the finite set $A_3\sqcup A_2$ has the property of (3) at every $w\in P'$: otherwise, choosing for each element a place of $P'$ at which it fails, the finite set $S_1$ of these places contradicts (3).
\item If it is $\sigma\in A_3$: $\sigma$ has a nonzero vector fixed by $K_0(S_0,c)$, and this group contains $K_0\cap U_{3,w}$ for every $w$ over whose place of $L^+$ no place of $S_0$ lies, so $\sigma$ is unramified at such $w$. With its local components at $S_0$ and at the places of $P'$, this is alternative (1). If it is $\sigma'\in A_2$, this is alternative (2).
\end{enumerate}
\end{proof}
% adjoint_lift_w_of_literal_automorphy_spec_out: smoothDepth for c_1, SpecOutS (hS) for c_0 and U_2, contE for the continuity of W_j and V_{2j+1}, the private outW for A_3 (with the condition at S_0) and Al.outRep_alt for A_2, Gd from cof, then the private pigeon over the finite sets S of P' with Al.inVec_alt (levels W = deep c on S_0 and iw on S_1), eigTS, SpecOutS, and in the second case Al.refine; regUnl for the separated characters; the private unram_of for the unramified places.

\begin{lemma}[\cite{CH13,BLGGT14}]\label{sha-automorphy-unitary-galois}
Assume that Galois representations exist, in the sense of Definition~\ref{gal-galois-representations}, relative to the maps satisfying the laws of Definition~\ref{gal-llc-laws}.
Let $L$ be a CM field such that $L/L^+$ is unramified at every finite place, let $\ell$ be an odd prime, let $\iota:\bar{\mathbf Q}_\ell\cong\mathbf C$, let $n\ge1$, and let $\sigma$ be an automorphic representation of $U_n$ (Definition~\ref{gal-definite-unitary}). Then there is a continuous $r:\operatorname{Gal}(\bar L/L)\to GL_n(\bar{\mathbf Q}_\ell)$ such that:
\begin{enumerate}
\item $\sigma$ has a local component at every finite place split over $L^+$, and $r$ is compatible with each of them in the sense of Lemma~\ref{sha-hecke-ideal-attached}(1);
\item for every complex conjugation $c\in\operatorname{Gal}(\bar L/L^+)$ there is an invertible $P$ with ${}^tr(s)Pr(csc^{-1})=\chi_{\mathrm{cyc}}(s)^{1-n}P$ for $s\in\operatorname{Gal}(\bar L/L)$;
\item if $r$ is irreducible, it is attached through $\iota$ to a regular algebraic cuspidal automorphic representation $\pi$ of $GL_n(\mathbf A_L)$ with $\pi^\vee\cong\pi^c$, which has Satake parameters at every finite place at which $\sigma$ is unramified and whose local components at the places split over $L^+$ are local components of $\sigma$.
\end{enumerate}
\end{lemma}
\leanref{Development/AutomorphicArithmetic/AutomorphicVarieties/ModularityArgument/ThresholdArguments/DenseThresholdCases/WStmtxx.lean}{UpperTheory.V.Z.Sym.W.GalExWAlt}
\leanref{Development/AutomorphicArithmetic/AutomorphicVarieties/ModularityArgument/ThresholdArguments/DenseThresholdCases/TGalExW.lean}{UpperTheory.V.Z.Sym.W.galois_exists_w_statement}

\begin{lemma}\label{sha-automorphy-spread}
Let $L$ be a CM field such that $L/L^+$ is unramified at every finite place, and let $G_L=\operatorname{Gal}(\bar L/L)$. A \emph{Frobenius element} at a finite place $w$ of $L$ is an element of the decomposition group of a place above $w$ that acts on the residue field by $x\mapsto x^{q_w}$. Let $E,B\subset G_L$ be closed, and suppose that $s\in E$ whenever $s,s'\in G_L$ and a complex conjugation $c\in\operatorname{Gal}(\bar L/L^+)$ satisfy $s'=csc^{-1}$ and $s'\in E$. Let $M\ge1$ and let $P$ be a set of finite places of $L$ such that:
\begin{enumerate}
\item every place $w$ split over $L^+$ and not dividing $M$, at which some Frobenius element lies outside $B$, lies over the same place of $L^+$ as a place of $P$;
\item for all but finitely many $w\in P$, every Frobenius element at $w$ lies in $E$.
\end{enumerate}
Then $G_L\setminus B\subset E$.
\end{lemma}
\leanref{Development/AutomorphicArithmetic/AutomorphicVarieties/ModularityArgument/SymmetricPowerTransfer/GeneralizedEigenspaces/WeightKernels/TwistRatios/Spread.lean}{UpperTheory.V.Z.Sym.Ga.spread}
\begin{proof}
Let $T$ be the set of places of $L$ over the same place of $L^+$ as one of the finitely many $w\in P$ at which some Frobenius element is not in $E$, and let $S$ be the union of $T$ with the set of places dividing $M$; $S$ is finite.
By Theorem~\ref{lit-chebotarev-degree-one}, applied to the finite Galois extensions of $L$ inside $\bar L$, every nonempty open subset of $G_L$ contains a Frobenius element at a place $v\notin S$ of absolute residue degree one. Since $E$ is closed and $G_L\setminus B$ is open, it suffices to show that such an element $s$, at such a place $v$, lies in $E$ if $s\notin B$.
Since $L/L^+$ is unramified and $v$ has residue degree one, $v$ is split over $L^+$; and $v\nmid M$. By (1) there is $w'\in P$ over the same place of $L^+$ as $v$, and $w'\notin T$ because $v\notin T$, so every Frobenius element at $w'$ lies in $E$. Either $w'=v$, and $s\in E$; or $w'$ is the conjugate of $v$, and for a complex conjugation $c$, $s'=csc^{-1}$ is a Frobenius element at $w'$. Then $s'\in E$, so $s\in E$.
\end{proof}
% Ga.spread: ChebotarevDegreeOne.dense (open_subset_closure_inter), isSplitAt_of_unramified_of_inertiaDeg_eq_one, Ga.conjEx and Ga.conjFrob (with cf_primes: w' = v or w' is the conjugate of v).

\begin{lemma}\label{sha-automorphy-regular-elements}
Let $L$ be a CM field, let $\ell$ be a prime, let $\iota:\bar{\mathbf Q}_\ell\cong\mathbf C$, let $M\ge1$ and $m\ge0$, and let $r$ be an adjoint input of level $M$ and weight $m$ (Definition~\ref{sha-automorphy-adjoint-input}). Then every open subgroup $H$ of $\operatorname{Gal}(\bar L/L)$ contains an element $s$ that is not exceptional for $r$ and such that $r(s)$ has eigenvalues $\alpha,\beta$ with $\alpha/\beta$ not of finite order.
\end{lemma}
\leanref{Development/AutomorphicArithmetic/AutomorphicVarieties/ModularityArgument/ThresholdArguments/DenseThresholdCases/TRegElt.lean}{UpperTheory.V.Z.Sym.Ga.exists_regular_elt_in_open_subgroup_infinite_order}
\begin{proof}
Let $A$ be a place of $\bar L$ above $\ell$, with $b$, $k$, $\kappa_0$ and $\kappa_1$ as in Definition~\ref{sha-automorphy-adjoint-input}(1)(h). Since $H\cap I_A$ is open in $I_A$, it contains an element $s$ with $\chi_{\mathrm{cyc}}(s)$ of infinite order, as in the proof of Lemma~\ref{sha-automorphy-local-ratios}(3). The eigenvalues of $r(s)$ are $\alpha=\kappa_0(s)$ and $\beta=\kappa_1(s)$, and $(\alpha/\beta)^k=\chi_{\mathrm{cyc}}(s)^{mk}$. Since $m\ge2$, $\alpha/\beta$ is not of finite order.
If $s$ were exceptional, then $(\alpha/\beta)^2+(\beta/\alpha)^2=\chi_{\mathrm{cyc}}(s)^j+\chi_{\mathrm{cyc}}(s)^{-j}$ for some $|j|\le2$, so $(\alpha/\beta)^2=\chi_{\mathrm{cyc}}(s)^{\pm j}$, the two sides being roots of the same quadratic polynomial; then $\chi_{\mathrm{cyc}}(s)^{(2m\mp j)k}=1$ with $2m\mp j>0$, which is impossible.
\end{proof}
% Ga.ineInf for s, then the private regElt_key (regElt_zpow_ne); the place A is the private regElt_place.

\begin{lemma}\label{sha-automorphy-traces}
Let $\Gamma$ be a profinite group, let $\ell$ be a prime and let $k=\bar{\mathbf Q}_\ell$. For $g\in GL_2(k)$ let $u(g)=\operatorname{tr}(g)^2/\det(g)-2$, so that $u(g)=\alpha/\beta+\beta/\alpha$ for the eigenvalues $\alpha,\beta$ of $g$. For a homomorphism $\rho:\Gamma\to GL_2(k)$ let $\operatorname{Sym}^4_0\rho=\operatorname{Sym}^4\rho\otimes(\det\rho)^{-2}$.
\begin{enumerate}
\item Let $\rho$ be continuous and irreducible on every open subgroup, and suppose that every open subgroup contains an element $s$ such that $\rho(s)$ has eigenvalues $\alpha,\beta$ with $\alpha/\beta$ not of finite order. Then $\operatorname{Sym}^4_0\rho$ is irreducible on every open subgroup.
\item Let $\rho,\rho':\Gamma\to GL_2(k)$ be homomorphisms with $\operatorname{Sym}^4_0\rho$ irreducible on every open subgroup. Then there is no nonempty open $V\subset\Gamma$ with $u(\rho'(g))=u(\rho(g))^2-2$ for all $g\in V$.
\item Let $R,R':\Gamma\to GL_n(k)$ be homomorphisms with $R$ irreducible on every open subgroup, and suppose that $\operatorname{tr}R(g)=\operatorname{tr}R'(g)$ for all $g$ in a nonempty open $V\subset\Gamma$. Then $R'\cong R\otimes\xi$ for a character $\xi$ of $\Gamma$ with open kernel.
\item Let $\rho:\Gamma\to GL_2(k)$ be continuous with $\operatorname{Sym}^4_0\rho$ irreducible on every open subgroup, let $\varepsilon:\Gamma\to k^\times$ be a continuous character, and call $g$ \emph{exceptional} if $u(\rho(g))^2-2=\varepsilon(g)^j+\varepsilon(g)^{-j}$ for some $|j|\le2$. Let $\xi$ be a character of $\Gamma$ with open kernel such that $\xi(g)(u(\rho(g))+1)=u(\rho(g))+1$ for every $g$ that is not exceptional. Then $\xi=1$.
\end{enumerate}
\end{lemma}
\leanref{Development/AutomorphicArithmetic/AutomorphicVarieties/ModularityArgument/SymmetricPowerTransfer/GeneralizedEigenspaces/WeightKernels/TwistRatios/ExistsNonzeroLowerEntryOfIrreducible.lean}{UpperTheory.V.Z.Sym.Ga.sym4Irr}
\leanref{Development/AutomorphicArithmetic/AutomorphicVarieties/ModularityArgument/SymmetricPowerTransfer/GeneralizedEigenspaces/WeightKernels/TwistRatios/NoKappa.lean}{UpperTheory.V.Z.Sym.Ga.noKappa}
\leanref{Development/AutomorphicArithmetic/AutomorphicVarieties/ModularityArgument/SymmetricPowerTransfer/GeneralizedEigenspaces/WeightKernels/TwistRatios/TwistOf.lean}{UpperTheory.V.Z.Sym.Ga.twistOf}
\leanref{Development/AutomorphicArithmetic/AutomorphicVarieties/ModularityArgument/SymmetricPowerTransfer/GeneralizedEigenspaces/ProjectionComponents/UfDg.lean}{UpperTheory.V.Z.Sym.U.xiOne}
\leanref{Development/AutomorphicArithmetic/AutomorphicVarieties/ModularityArgument/SymmetricPowerTransfer/GeneralizedEigenspaces/WeightKernels/TwistRatios/KerLe.lean}{UpperTheory.V.Z.Sym.Ga.kerLe}
\begin{proof}
(1) Let $H$ be open, $s\in H$ as in the hypothesis, and conjugate so that $\rho(s)=\operatorname{diag}(\alpha,\beta)$; the twist by $(\det\rho)^{-2}$ does not affect irreducibility. The diagonal entries $\alpha^4,\alpha^3\beta,\dots,\beta^4$ of $\operatorname{Sym}^4\rho(s)$ are distinct, since $\alpha^m\ne\beta^m$ for $m\ge1$, so a subspace stable under $\operatorname{Sym}^4\rho(H)$ is spanned by basis vectors. The irreducibility of $\rho$ on the open subgroups of $H$ gives elements of every open subgroup with nonzero lower left entry and elements with nonzero upper right entry, and from these, by continuity, an element $g\in H$ such that all four entries of $\rho(g)$ are nonzero. Then the first row and the first column of $\operatorname{Sym}^4\rho(g)$ have no zero entry, so a nonzero stable subspace contains the first basis vector, and then all of them.

\emph{A dimension bound.} Let $H$ be a group, $R:H\to GL_n(k)$ irreducible, $x_0\in GL_n(k)$, $\varphi:H\to\mathcal B$ a homomorphism to the units of a finite-dimensional $k$-algebra, and $f$ a linear form on $\mathcal B$ with $\operatorname{tr}(x_0R(h))=f(\varphi(h))$ for all $h$. If $\sum_hc_h\varphi(h)=0$ in $\mathcal B$, then $\sum_hc_h\varphi(h'hh'')=0$, so $\operatorname{tr}(x_0R(h')(\sum_hc_hR(h))R(h''))=0$ for all $h',h''$; since the $R(h)$ span $M_n(k)$ (Burnside's theorem) and $x_0$ is invertible, $\sum_hc_hR(h)=0$. So $R$ factors through the subalgebra generated by $\varphi(H)$, which has dimension at least $n^2$.

(2) Suppose $V$ exists, let $g\in V$, and let $H$ be an open normal subgroup with $gH\subset V$. For $2\times2$ matrices, $\operatorname{tr}\operatorname{Sym}^4_0(x)=u(x)^2+u(x)-1$ and $\operatorname{tr}\operatorname{ad}^0(x)=u(x)+1$. So for $h\in H$,
$\operatorname{tr}\bigl(\operatorname{Sym}^4_0\rho(g)\operatorname{Sym}^4_0\rho(h)\bigr)=u(\rho(gh))+u(\rho'(gh))+1=\operatorname{tr}\bigl(\operatorname{ad}^0\rho(g)\operatorname{ad}^0\rho(h)\bigr)+\operatorname{tr}\bigl(\operatorname{ad}^0\rho'(g)\operatorname{ad}^0\rho'(h)\bigr)-1$,
a linear form in $\varphi(h)=(\operatorname{ad}^0\rho(h),\operatorname{ad}^0\rho'(h),1)\in M_3(k)\times M_3(k)\times k$. The dimension bound, for $R=\operatorname{Sym}^4_0\rho|_H$ and $n=5$, gives $25\le9+9+1$, a contradiction.

(3) Let $g_0\in V$ and $H$ an open normal subgroup with $g_0H\subset V$. Then $\operatorname{tr}(R(g_0)R(h))=\operatorname{tr}(R'(g_0)R'(h))$ for $h\in H$. By the dimension bound for $R|_H$ and $x_0=R(g_0)$, the linear extension of $R'|_H$ to the group algebra kills the kernel of that of $R|_H$, whose image is $M_n(k)$; so $R'|_H=\Theta\circ R|_H$ for an algebra endomorphism $\Theta$ of $M_n(k)$, which is inner. Hence $R'|_H\cong R|_H$, and since $R|_H$ is irreducible and $H$ is normal, Schur's lemma gives a character $\xi$, trivial on $H$, with $R'\cong R\otimes\xi$.

(4) Suppose $\xi(g_0)\ne1$, and let $U=g_0\ker\xi$, an open set on which $\xi=\xi(g_0)$. If some $g\in U$ is not exceptional, the non-exceptional elements of $U$ form a nonempty open set, since the exceptional elements form a closed set; on it $u(\rho(g))=-1$, so $u(\rho(g))=u(\rho(g))^2-2$, contradicting (2) with $\rho'=\rho$. Otherwise $U$ is contained in the union of the closed sets $B_j=\{g:u(\rho(g))^2-2=\varepsilon(g)^j+\varepsilon(g)^{-j}\}$, $|j|\le2$, so one of them contains a nonempty open set $V'$. With $\rho'=\operatorname{diag}(\varepsilon^j,1)$, $u(\rho'(g))=\varepsilon(g)^j+\varepsilon(g)^{-j}=u(\rho(g))^2-2$ on $V'$, contradicting (2).
\end{proof}
% (1) is Ga.sym4Irr (S4.eig_basis, S4.exists_nonzero_lower_entry_of_irreducible and its upper analogue, S4.exists_full, S4.core). The dimension bound is Ga.kerLe (Reformulated.B.Qe26.kerLe). (2) is Ga.noKappa, with Ga.coset, Ga.trSym4 and Ga.ufFacts, and the private phiB and fB. (3) is Ga.twistOf, with Ga.coset, Ga.conjOn and the private twistOf_ratio. (4) is U.xiOne, with Ga.badClosed and the private dg and uf_dg.

\begin{lemma}\label{sha-automorphy-adjoint-branches}
Assume that Galois representations exist, in the sense of Definition~\ref{gal-galois-representations}, relative to the maps satisfying the laws of Definition~\ref{gal-llc-laws}.
Let $L$ be a CM field such that $L/L^+$ is unramified at every finite place, let $\ell$ be an odd prime, let $\iota:\bar{\mathbf Q}_\ell\cong\mathbf C$, let $M\ge1$ and $m\ge0$, and let $r$ be an adjoint input of level $M$ and weight $m$ (Definition~\ref{sha-automorphy-adjoint-input}). Let $\sigma$ be an automorphic representation of $U_2$ that has a local component at every finite place split over $L^+$, with which $r$ is compatible in the sense of Lemma~\ref{sha-hecke-ideal-attached}(1). Let $P$ be a set of finite places of $L$ and, for $w\in P$, let $\psi_{w,0},\psi_{w,1}$ be smooth characters of $L_w^\times$ and $a_w=\psi_{w,0}/\psi_{w,1}$, such that: every $w\in P$ is split over $L^+$ and prime to $\ell$, $I(\psi_{w,0},\psi_{w,1})$ is a local component of $\sigma$ at $w$, $a_w$ is unramified, and $a_w^2\ne|\cdot|^k$ for $|k|\le2$; and every place $w$ split over $L^+$, not dividing $M$, at which some Frobenius element is not exceptional for $r$, lies over the same place of $L^+$ as a place of $P$.
\begin{enumerate}
\item Let $\sigma_3$ be an automorphic representation of $U_3$ that, at all but finitely many $w\in P$, has a local component that is a subquotient of $I(a_w,1,a_w^{-1})$. Then there is a regular algebraic cuspidal automorphic representation $\pi$ of $GL_3(\mathbf A_L)$ with $\pi^\vee\cong\pi^c$, attached through $\iota$ to $\operatorname{ad}^0r\otimes\chi_{\mathrm{cyc}}^{-1}$, that has Satake parameters at every finite place at which $\sigma_3$ is unramified, and whose local components at every place split over $L^+$ are local components of $\sigma_3$.
\item There is no automorphic representation $\sigma'$ of $U_2$ which, at all but finitely many $w\in P$, has a local component that is a subquotient of $I(\psi')$ with $\psi'_0/\psi'_1=a_w^2$.
\end{enumerate}
\end{lemma}
\leanref{Development/AutomorphicArithmetic/AutomorphicVarieties/ModularityArgument/ThresholdArguments/DenseThresholdCases/ResidualCharpolyRankFour.lean}{UpperTheory.V.Z.Sym.W.branch1W_alt}
\leanref{Development/AutomorphicArithmetic/AutomorphicVarieties/ModularityArgument/ThresholdArguments/DenseThresholdCases/UfOfCharpoly.lean}{UpperTheory.V.Z.Sym.W.branch_two_w_statement}
\begin{proof}
Let $G_L=\operatorname{Gal}(\bar L/L)$, $\varepsilon=\chi_{\mathrm{cyc}}$, and let $u$ be as in Lemma~\ref{sha-automorphy-traces}. The following are direct computations with matrices: $u(g^{-1})=u(g)$, $\operatorname{tr}\operatorname{ad}^0(g)=u(g)+1$, and $u(g')=u(g)$ if ${}^tgPg'=\mu P$ for an invertible $P$ and a scalar $\mu$; and for $g,g'\in GL_n$ with ${}^tgPg'=\mu P$, $\operatorname{tr}g'=\mu\operatorname{tr}g^{-1}$. The set $B$ of elements of $G_L$ that are exceptional for $r$ is closed, being defined by finitely many equations between continuous functions; by Lemma~\ref{sha-automorphy-regular-elements} its complement is not empty.

\emph{Eigenvalues at $P$.} Let $w\in P$, $A$ a place above $w$, $a$ a local Artin map for $A$ and $s\in W_A$. Since $|\cdot|\ne1$, $a_w\ne|\cdot|^{\pm1}$ and $a_w^2\ne|\cdot|^{\pm1}$, each of the families $(\psi_{w,0},\psi_{w,1})$, $(a_w,1,a_w^{-1})$ and $(\psi'_0,\psi'_1)$ with $\psi'_0/\psi'_1=a_w^2$ is unlinked, so a local component that is a subquotient of the corresponding principal series is plain. If $r'$ is compatible with such a local component $\tau$ of rank $n$, then, as in the proof of Lemma~\ref{sha-automorphy-local-ratios}, Lemma~\ref{sha-hecke-local-rec}(1) shows that the eigenvalues of $\iota r'(s)$ are the values $\psi_i(a(s))|a(s)|^{(1-n)/2}$ for the characters $\psi_i$ of the family; and $\iota\varepsilon(s)=|a(s)|$. With $b=a_w(a(s))$ this gives $\iota u(r(s))=b+b^{-1}$.

(1) Let $r_3$ be the representation of Lemma~\ref{sha-automorphy-unitary-galois} for $\sigma_3$, and let $E$ be the set of $g\in G_L$ with $\operatorname{tr}r_3(g)\varepsilon(g)=u(r(g))+1$ and $\operatorname{tr}r_3(g^{-1})\varepsilon(g^{-1})=u(r(g^{-1}))+1$. It is closed.
\begin{enumerate}
\item \emph{Stability.} Let $c$ be a complex conjugation and $s'=csc^{-1}$ with $s'\in E$. The pairings of Lemma~\ref{sha-automorphy-unitary-galois}(2) for $r_3$ and of Definition~\ref{sha-automorphy-adjoint-input}(1)(c) for $r$ give $\operatorname{tr}r_3(s')=\varepsilon(s)^{-2}\operatorname{tr}r_3(s^{-1})$, $\operatorname{tr}r_3(s'^{-1})=\varepsilon(s)^2\operatorname{tr}r_3(s)$ and $u(r(s'))=u(r(s))$; and $\varepsilon(s')=\varepsilon(s)$. So $\operatorname{tr}r_3(s)\varepsilon(s)=\operatorname{tr}r_3(s'^{-1})\varepsilon(s'^{-1})=u(r(s'))+1=u(r(s))+1$, and in the same way for $s^{-1}$; so $s\in E$.
\item \emph{Frobenius elements.} For all but finitely many $w\in P$, $\sigma_3$ has a local component at $w$ that is a subquotient of $I(a_w,1,a_w^{-1})$, and $r_3$ is compatible with it. For a Frobenius element $s$ at $w$ and $t\in\{s,s^{-1}\}$, both in $W_A$, the eigenvalues computed above give $\iota\operatorname{tr}r_3(t)=(b+1+b^{-1})|a(t)|^{-1}$ with $b=a_w(a(t))$, so $\operatorname{tr}r_3(t)\varepsilon(t)=u(r(t))+1$. So $s\in E$.
\item By Lemma~\ref{sha-automorphy-spread}, with the hypothesis on $P$, $G_L\setminus B\subset E$.
\item \emph{A twist.} Let $R=\operatorname{ad}^0r\otimes\varepsilon^{-1}$. It is irreducible on every open subgroup, since $r$ is (Definition~\ref{sha-automorphy-adjoint-input}(1)(d)). For $g\notin B$, $\operatorname{tr}r_3(g)=\varepsilon(g)^{-1}(u(r(g))+1)=\operatorname{tr}R(g)$. Lemma~\ref{sha-automorphy-traces}(3), with $V=G_L\setminus B$, gives a character $\xi$ with open kernel and $r_3\cong R\otimes\xi$.
\item By Lemmas~\ref{sha-automorphy-regular-elements} and~\ref{sha-automorphy-traces}(1), $\operatorname{Sym}^4_0r$ is irreducible on every open subgroup. For $g\notin B$, $\operatorname{tr}r_3(g)=\xi(g)\operatorname{tr}R(g)=\operatorname{tr}R(g)$, so $\xi(g)(u(r(g))+1)=u(r(g))+1$. Lemma~\ref{sha-automorphy-traces}(4), whose exceptional elements are those of $B$, gives $\xi=1$, so $r_3\cong R$.
\item So $r_3$ is irreducible, and Lemma~\ref{sha-automorphy-unitary-galois}(3) gives $\pi$, attached to $r_3$ and hence to $R$, with the stated properties.
\end{enumerate}

(2) Suppose that $\sigma'$ exists, and let $r'$ be the representation of Lemma~\ref{sha-automorphy-unitary-galois} for $\sigma'$. Let $E$ be the set of $g\in G_L$ with $u(r'(g))=u(r(g))^2-2$; it is closed. If $s'=csc^{-1}$ for a complex conjugation $c$, the pairings give $u(r(s'))=u(r(s))$ and $u(r'(s'))=u(r'(s))$, so $s'\in E$ implies $s\in E$. For all but finitely many $w\in P$, $\sigma'$ has a local component at $w$ that is a subquotient of $I(\psi')$ with $\psi'_0/\psi'_1=a_w^2$; for a Frobenius element $s$ at $w$, the eigenvalues computed above give $\iota u(r'(s))=b^2+b^{-2}=(b+b^{-1})^2-2$, with $b=a_w(a(s))$, so $s\in E$. By Lemma~\ref{sha-automorphy-spread}, $G_L\setminus B\subset E$, a nonempty open set on which $u(r'(g))=u(r(g))^2-2$. But $\operatorname{Sym}^4_0r$ is irreducible on every open subgroup, by Lemmas~\ref{sha-automorphy-regular-elements} and~\ref{sha-automorphy-traces}(1), which contradicts Lemma~\ref{sha-automorphy-traces}(2).
\end{proof}
% (1), branch1W_alt: galois_exists_w_statement for r3; Ga.badClosed; E and hstab (Ga.trPair, Ga.ufFacts); hfrob with the private local_trace_compatible (plain_of_statement with unl_two and unl_three, eigW_alt, Ga.cyc_char_eq_abs_pow_and_norm_pow_deg_away_l, trace_det_of_charpoly_two, AuxBlock12); Ga.spread; Ga.exists_regular_elt_in_open_subgroup_infinite_order; Sd.ad0Irr and Sd.shapeFacts for R; Ga.twistOf; Ga.sym4Irr; U.xiOne; Ga.isAutC_isAutCsdA_isIrr_irrOpen_symm_tw_of_repIso and att_iso.
% (2), branch_two_w_statement: the Galois representation r' of sigma' (galois_exists_w_statement); the closed set E of g with uf(r' g) = uf(r g)^2 - 2, stable under the complex conjugations by the pairings (Ga.ufFacts); E contains the Frobenius elements at almost all w in P, by plain_of_statement (reg_one, reg_two), eigW_alt and uf_of_pol; Ga.spread (with Ga.badClosed and the condition on P) puts the complement of the exceptional elements in E; Ga.exists_regular_elt_in_open_subgroup_infinite_order, Ga.sym4Irr and Ga.noKappa give the contradiction.

\begin{lemma}\label{sha-automorphy-plain-at}
Let $L$ be a CM field, let $\sigma$ be an automorphic representation of $U_3$ (Definition~\ref{gal-definite-unitary}), and let $\pi$ be a cuspidal automorphic representation of $GL_3(\mathbf A_L)$ whose local components at every place split over $L^+$ are local components of $\sigma$. Let $w$ be a place split over $L^+$ and $a$ a smooth character of $L_w^\times$ with $a^j\ne|\cdot|^{\pm1}$ for $j=1,2$, and suppose that $\sigma$ has a local component at $w$ that is a subquotient of $I(a,1,a^{-1})$ (Definition~\ref{sha-automorphy-adjoint-input}). Then every local component of $\pi$ at $w$ is plain (Definition~\ref{gal-plain}).
\end{lemma}
\leanref{Development/AutomorphicArithmetic/AutomorphicVarieties/ModularityArgument/ThresholdArguments/DenseThresholdCases/TPlainAtW.lean}{UpperTheory.V.Z.Sym.W.plain_at_w_statement}
\begin{proof}
A local component $\tau$ of $\pi$ at $w$ is a local component of $\sigma$ at $w$. The local component of the irreducible representation $\sigma$ at a split place is unique up to isomorphism, by its decomposition as a restricted tensor product \cite{Fla79}; so $\tau$ is a subquotient of $I(a,1,a^{-1})$, and its supercuspidal support is $\{a,1,a^{-1}\}$. No two of $a,1,a^{-1}$ differ by $|\cdot|$: that would give $a=|\cdot|^{\pm1}$, $a^2=|\cdot|^{\pm1}$ or $|\cdot|=1$, and the last is false. So the support consists of characters and is unlinked, and $\tau$ is plain.
\end{proof}
% The private ratios and absOne_ne; UpperTheory.building_component_uniqueness; plain_of_statement (PlainOf).

\begin{lemma}\label{sha-automorphy-windowed-glue}
Let $L$ be a CM field, let $\ell$ be a prime, let $\iota:\bar{\mathbf Q}_\ell\cong\mathbf C$, and let $R:\operatorname{Gal}(\bar L/L)\to GL_3(\bar{\mathbf Q}_\ell)$ be a homomorphism. Suppose that for every finite place $w$ of $L$ there is a cuspidal automorphic representation $\pi_w$ of $GL_3(\mathbf A_L)$, attached through $\iota$ to $R$ (Definition~\ref{gal-weakly-attached}), that has Satake parameters at $w$ or all of whose local components at $w$ are plain. Then every regular algebraic cuspidal automorphic representation $\pi$ of $GL_3(\mathbf A_L)$ attached through $\iota$ to $R$ satisfies the windowed condition at every finite place: there is $t\in\mathbf R$ such that at every finite place $w$, the local components of $\pi$ at $w$ have supercuspidal support consisting of characters $\xi$ with $q_w^{t-1/2}<|\xi(\varpi_w)|<q_w^{t+1/2}$.
\end{lemma}
\leanref{Development/AutomorphicArithmetic/AutomorphicVarieties/ModularityArgument/SymmetricPowerTransfer/WeilDeligneUnramified/WeilCharpolySupport/GlueW.lean}{UpperTheory.V.Z.Sym.W.glueW}
\begin{proof}
For a cuspidal $\pi$ there is $t\in\mathbf R$ such that the bound holds at every finite place $w$ at which $\pi$ has Satake parameters or plain local components. Indeed, $\pi$ is the twist of a unitary cuspidal $\pi_0$ by $|\det|^{-t}$. At such a place the supercuspidal support of the local components consists of characters and is unlinked, for a spherical component because it is an irreducible unramified principal series; the local component of $\pi_0$ is unitary and generic, so the characters of its support have absolute values at $\varpi_w$ strictly between $q_w^{-1/2}$ and $q_w^{1/2}$ \cite{JS81a}.
Let $\pi$ be regular algebraic cuspidal and attached to $R$, with such a $t$, and let $w$ be a finite place. At all but finitely many places, $\pi_w$ and $\pi$ both have Satake parameters, determined by the eigenvalues of $R$ at a Frobenius element; so $\pi_w=\pi$ by strong multiplicity one \cite{JS81b}. Hence $\pi$ has Satake parameters at $w$ or plain local components at $w$, and the bound holds at $w$.
\end{proof}
% ImaginaryQuadraticEtaCharacter.winOfP_of (WinOf, PlainIff, SphUnl) for t, with winOf through CG.exists_isUnitary_twist_absPow, ResW.isUnitarisable_of_isLocComp_cusp, Generic.isGeneric_of_isLocComp and exponent_of_unl; Af.strongMultOneC and hecke_eigenvalue_basic_properties for pi_w = pi.

\begin{proposition}\label{sha-automorphy-adjoint-automorphy}
Assume that Galois representations exist, in the sense of Definition~\ref{gal-galois-representations}, relative to the maps satisfying the laws of Definition~\ref{gal-llc-laws}.
Let $L$ be a CM field, let $\ell$ be an odd prime, let $\iota:\bar{\mathbf Q}_\ell\cong\mathbf C$, and let $M\ge1$ with $2\ell\mid M$ be such that $L/L^+$ is unramified at every finite place and every place of $L^+$ dividing $M$ splits in $L$. Let $m\ge0$ and let $r$ be an adjoint input of level $M$ and weight $m$ (Definition~\ref{sha-automorphy-adjoint-input}). Then $\operatorname{ad}^0r\otimes\chi_{\mathrm{cyc}}^{-1}$ is attached through $\iota$ to a regular algebraic cuspidal automorphic representation $\pi$ of $GL_3(\mathbf A_L)$ with $\pi^\vee\cong\pi^c$, and every such $\pi$ satisfies the windowed condition of Lemma~\ref{sha-automorphy-windowed-glue} at every finite place.
\end{proposition}
\leanref{Development/AutomorphicArithmetic/AutomorphicVarieties/ModularityArgument/SymmetricPowerTransfer/WeilDeligneUnramified/UnitaryTwoForm.lean}{UpperTheory.V.Z.Sym.W.AdAutW}
\leanref{Development/AutomorphicArithmetic/AutomorphicVarieties/ModularityArgument/ThresholdArguments/DenseThresholdCases/MemOfUnder.lean}{UpperTheory.V.Z.Sym.W.adjoint_automorphy_w_of_pure_weight_adjoint_lift}
\begin{proof}
\emph{Step 1.} Let $w_0$ be a place of $L$ dividing $M$. There is a regular algebraic cuspidal $\pi$ of $GL_3(\mathbf A_L)$ with $\pi^\vee\cong\pi^c$, attached through $\iota$ to $\operatorname{ad}^0r\otimes\chi_{\mathrm{cyc}}^{-1}$, with Satake parameters at every finite place not dividing $M$, and with plain local components at $w_0$.
Whether a place of $L$ divides $M$ depends only on the place of $L^+$ below it.
\begin{enumerate}
\item Lemma~\ref{sha-automorphy-unitary-two} gives $\sigma$ and $\mathrm{HT}$.
\item Choose a set $S$ of places of $L$ dividing $M$ that contains $w_0$ and one place over each place of $L^+$ below a place dividing $M$, and no two over the same place of $L^+$; it is finite, and its places split over $L^+$. Let $X$ be the set of places $w$ of $L$ split over $L^+$, not dividing $M$, at which some Frobenius element is not exceptional for $r$, and choose $Q\subset X$ containing one place over each place of $L^+$ below a place of $X$. Then $S$ and $Q$ are disjoint, and no two places of $S\cup Q$ lie over the same place of $L^+$. Every place dividing $2$ divides $M$, so it lies over the same place of $L^+$ as a place of $S$.
\item For $w\in Q$, Lemma~\ref{sha-automorphy-local-ratios}(1), applied to a local component of $\sigma$ at $w$, which exists, is plain and is compatible with $r$ by step (1), gives unramified $\psi_{w,0},\psi_{w,1}$ with $I(\psi_{w,0},\psi_{w,1})$ a local component of $\sigma$ and $a_w^2\ne|\cdot|^k$ for $|k|\le2$, where $a_w=\psi_{w,0}/\psi_{w,1}$; here $w\nmid\ell$ since $\ell\mid M$, $r$ is trivial on the inertia groups at $w$ by Definition~\ref{sha-automorphy-adjoint-input}(1)(f), and some Frobenius element at $w$ is not exceptional since $w\in X$.
For $w\in S$, Lemma~\ref{sha-automorphy-local-ratios}(3), if $w\mid\ell$, with Definition~\ref{sha-automorphy-adjoint-input}(1)(e) and (h), and Lemma~\ref{sha-automorphy-local-ratios}(2), if $w\nmid\ell$, with Definition~\ref{sha-automorphy-adjoint-input}(1)(b), give $\psi_{w,0},\psi_{w,1}$ with $I(\psi_{w,0},\psi_{w,1})$ a local component of $\sigma$ and $a_w^j\ne|\cdot|^{\pm1}$ for $j=1,2$.
\item By step (2), $\sigma$ is unramified at every place $w$ such that no place of $S$ lies over the same place of $L^+$ as $w$, since such a $w$ does not divide $M$. So Lemma~\ref{sha-automorphy-adjoint-lift}, with $S_0=S$, $P=Q$, $\pi=\sigma$ and $j=m-1\ge1$, applies.
\item By the choice of $Q$, $P=Q$ satisfies the hypotheses of Lemma~\ref{sha-automorphy-adjoint-branches}, so by its part (2) the second alternative of Lemma~\ref{sha-automorphy-adjoint-lift} does not occur. Let $\sigma_3$ be as in its first alternative.
\item Lemma~\ref{sha-automorphy-adjoint-branches}(1) gives $\pi$, regular algebraic cuspidal with $\pi^\vee\cong\pi^c$ and attached to $\operatorname{ad}^0r\otimes\chi_{\mathrm{cyc}}^{-1}$. At a place $w\nmid M$, no place of $S$ lies over the same place of $L^+$, so $\sigma_3$ is unramified at $w$ and $\pi$ has Satake parameters at $w$. At $w_0$, $\sigma_3$ has a local component that is a subquotient of $I(a_{w_0},1,a_{w_0}^{-1})$, with $a_{w_0}^j\ne|\cdot|^{\pm1}$ for $j=1,2$ by step (3), so the local components of $\pi$ at $w_0$ are plain by Lemma~\ref{sha-automorphy-plain-at}.
\end{enumerate}

\emph{Step 2.} Since $\ell\mid M$, some place $w_1$ of $L$ divides $M$, and Step 1 for $w_0=w_1$ gives $\pi$, which proves the first statement. For every finite place $w$ there is a cuspidal representation attached to $\operatorname{ad}^0r\otimes\chi_{\mathrm{cyc}}^{-1}$ with Satake parameters or plain local components at $w$: the one from Step 1 for $w_0=w$ if $w\mid M$, and $\pi$ if $w\nmid M$. Lemma~\ref{sha-automorphy-windowed-glue} gives the second statement.
\end{proof}
% The private core of MemOfUnder.lean is Step 1: aboveWeight, selectW (S, with hfin), Ga.select (Q, from X with Ga.FrobAt and Ga.Bad), good_w_statement, above_w_statement and away_w_statement (key), the GoodSet hgood, AdLiftW with S_0 = S and P = Q, branch1W_alt and branch_two_w_statement, and plain_at_w_statement at w_0. Step 2 is adjoint_automorphy_w_of_pure_weight_adjoint_lift, with GlobalTheory.EvExtend.placeOf for w_1 and glueW. The windowed condition of the statement is Birth.AllW (ImaginaryQuadraticEtaCharacter.AllWin).

\subsection{Realizations of linked points}\label{subsection-sha-realizations}

\begin{situation}\label{sha-automorphy-link-situation}
Assume that Galois representations exist, in the sense of Definition~\ref{gal-galois-representations}, relative to the maps satisfying the laws of Definition~\ref{gal-llc-laws}.
Let $(\ell,k,\psi,K)$ be an Eisenstein datum with deformation data $(\eta,c,\Psi)$ (Definition~\ref{sha-ordinary-deformation}), let $\iota:\bar{\mathbf Q}_\ell\cong\mathbf C$ and let $\mathfrak B$ be a place above $\ell$. Links, linked points, the residual conditions and seed data (Definitions~\ref{sha-descent-link}, \ref{sha-descent-residual} and~\ref{sha-descent-seed}) are those of the descent datum $(\eta,\Psi,c,\mathfrak B,M,\operatorname{Gal}(\bar{\mathbf Q}/\mathbf Q))$ of Lemma~\ref{sha-descent-eisenstein}, with $M=\ell N|d_K|$ for the modulus $N$ of $\psi$. Fix a link to an abelian CM field $L$ that satisfies the residual conditions with a constant $C_0$ for which it is large, and let $d=[L^+:\mathbf Q]$. For $n\in\{2,3\}$ let $w_n=(1,0,-1)$ or $(1,-1)$; the \emph{prescribed inertial characters} of a linked point of rank $n$ are $\psi_{v,i}=\Psi^{w_{n,i}}$ on $I_v$.
\end{situation}

\begin{definition}\label{sha-automorphy-realization}
In Situation~\ref{sha-automorphy-link-situation}, let $n\in\{2,3\}$, let $\mathcal O_1\supset\mathcal O$ be the ring of integers of a finite extension of $\mathbf Q_\ell$, and let $\kappa:\operatorname{Gal}(\bar{\mathbf Q}/L)\to\bar{\mathbf Q}_\ell^\times$ be a character.
\begin{enumerate}
\item $\kappa$ is an \emph{admissible twist of rank $n$} if it is the $\ell$-adic character through $\iota$ of an algebraic Hecke character, $\kappa(s)\kappa(s')=\chi_{\mathrm{cyc}}(s)^{1-n}$ whenever $s'=csc^{-1}$, $\kappa$ is trivial on the inertia groups at the places not dividing $M$, and $\kappa^m$ is trivial on the inertia groups at the places not above $\ell$ for some $m$ prime to $\ell$.
\item A \emph{realization of rank $n$} of the link, \emph{twisted by $\kappa$}, with coefficients $\mathcal O_1$, consists of: a constant for which the link, with its residual characters composed with the map of residue fields, satisfies the residual conditions and is large; a residual datum of rank $n$ over $\mathcal O_1$ in Situation~\ref{sha-hecke-situation}, for $L$, the places above $\ell$ of the link and an auxiliary place $v_0$ not dividing $M$, at which the residual characters $\bar\chi_i$ are trivial on $I_{v_0}$ and agree on $D_{v_0}$ and $\varepsilon\bmod\ell$ is nontrivial on $D_{v_0}$, with its universal ordinary determinant ring $B_n$ and Iwasawa algebra $\Lambda$; a continuous surjection $\pi$ from the Galois group of the residual datum onto the group $\Gamma_0$ of the link, compatible with the conjugations, the cyclotomic characters and the decomposition groups above $\ell$, through which the map $p$ of the link factors on $\operatorname{Gal}(\bar{\mathbf Q}/L)$, and a character $\kappa_0$ of that Galois group with values in $\mathcal O_1^\times$ that equals $\kappa$ on $\operatorname{Gal}(\bar{\mathbf Q}/L)$, such that the residual characters of the residual datum are those of the link composed with $\pi$ and twisted by $\kappa_0$, and such that $\kappa_0$ is continuous for the adic topology of $\mathcal O_1$, $\kappa_0(g)\kappa_0(cgc^{-1})=\varepsilon(g)^{1-n}$ for every $g$, $\kappa_0$ is trivial on $I_{v_0}$, and for every $u\in S_{\mathrm a}$ and every $g\in I_u$ there is an integer $m$ prime to $\ell$ with $\pi(g)^m=1$ and $\kappa_0(g)^m=1$; such that, for every $v\mid\ell$, $\mu_v$ lies in the closure of the commutator subgroup of $D_v$; and such that the kernel of $\pi$ lies in every closed normal subgroup $N$ of the Galois group of the residual datum that is stable under $g\mapsto cgc^{-1}$, contains $I_{v_0}$, and contains every $g\in I_u$, for $u\in S_{\mathrm a}$, with $\bar\Psi_\Gamma(\pi(g))=1$; and the universal ring $A_n$ of the linked points of rank $n$ over $\mathcal O_1$ lifting the residual one, with its universal linked point $(\rho^{\mathrm u},\psi^{\mathrm u})$. Let $B\subseteq A_n$ be the closure, for the $\mathfrak m_{A_n}$-adic topology, of the $\mathcal O_1$-subalgebra generated by the values of the characters $\psi^{\mathrm u}_{v,i}$ and by the products of entries of $\rho^{\mathrm u}$ of total weight $0$, the entry $(i,j)$ having weight $w_{n,i}-w_{n,j}$; and let $A_n^\circ\subseteq A_n$ be the $\mathcal O_1$-subalgebra generated by $B$ and the entries of $\rho^{\mathrm u}$.
\item A linked point $(\rho,\psi)$ of rank $n$ gives, through $\pi$ and the twist by $\kappa$, a determinant and characters as in Situation~\ref{sha-hecke-situation}. A map $\phi:B_n\to B$ is \emph{compatible} if, followed by the inclusion of $B$ in $A_n$, it carries the universal determinant and characters of $B_n$ to those given by the universal linked point.
\item A prime $\mathfrak p$ of $A_n^\circ$ is \emph{off-triangular} if, for the universal linked point $\rho^{\mathrm u}$, some entry $\rho^{\mathrm u}(g)_{ij}$ with $i>j$ and some entry $\rho^{\mathrm u}(g')_{ij}$ with $i<j$ lie outside $\mathfrak p$, and, if $n=3$, some entry $\rho^{\mathrm u}(g'')_{ij}$ with $w_{3,i}-w_{3,j}$ odd lies outside $\mathfrak p$.
\end{enumerate}
\end{definition}
\leanref{Development/AutomorphicArithmetic/AutomorphicVarieties/CurveDVRSpecialization/EveryPointSatisfiesFrameCondition.lean}{UpperTheory.V.Pr.IsKa}
\leanref{Development/AutomorphicArithmetic/AutomorphicVarieties/CurveDVRSpecialization/EveryPointSatisfiesFrameCondition.lean}{UpperTheory.V.Pr.CReal}
\leanref{Development/AutomorphicArithmetic/AutomorphicVarieties/CurveDVRSpecialization/PropagationTheoremGL2.lean}{UpperTheory.V.Pr.TwLink}
\leanref{Development/AutomorphicArithmetic/AutomorphicVarieties/FrobeniusSmoothVectors/LeAdicClosure.lean}{UpperTheory.V.Fr.Space}
\leanref{Development/AutomorphicArithmetic/AutomorphicVarieties/FrobeniusSmoothVectors/LeAdicClosure.lean}{UpperTheory.V.Fr.Space.B}
\leanref{Development/AutomorphicArithmetic/AutomorphicVarieties/FrobeniusSmoothVectors/LeAdicClosure.lean}{UpperTheory.V.Fr.Space.A}
\leanref{Development/AutomorphicArithmetic/AutomorphicVarieties/FrobeniusSmoothVectors/LeAdicClosure.lean}{UpperTheory.V.Fr.NRep.TI}
% Rule 13: the place above l is \mathfrak B in Definition sha-ordinary-deformation, the descent section (sha-descent-*), the Situation and Theorem sha-residually-reducible-automorphy. From the Situation to the proof of the Theorem, B is the ring of (2) (the Lean's Space.B), except in the statements set in Situation sha-hecke-situation (Lemma sha-automorphy-realization-exists-determinant, Definition sha-automorphy-unramified-multiplicity(3), Lemma sha-automorphy-unramified-tangent), where B is that Situation's universal ordinary determinant ring, and in the lemmas of commutative algebra, where B is any ring; the universal ordinary determinant ring of a realization is B_n. The later sections on rigidity, admissibility, lifts and the local model (sha-rigidity-*, sha-admissibility-stable, sha-lifts-*, sha-model-*) still write B for a place above l, with about 350 symbols (D_B, I_B, R_B, u_B, x_B, a_B, ...).

\begin{lemma}\label{sha-automorphy-realization-twists}
In Situation~\ref{sha-automorphy-link-situation}, $\chi_{\mathrm{cyc}}^{-1}$ is an admissible twist of rank three (Definition~\ref{sha-automorphy-realization}). If $\kappa$ is a character of $\operatorname{Gal}(\bar{\mathbf Q}/K)$, the $\ell$-adic character through $\iota$ of an algebraic Hecke character, with $\kappa(s)\kappa(csc^{-1})=\chi_{\mathrm{cyc}}(s)^{-1}$, trivial on the inertia groups at the places dividing neither $\ell$ nor $d_K$, and with $\kappa^m$ trivial on the inertia groups at the places not above $\ell$ for some $m$ prime to $\ell$, then its restriction to $\operatorname{Gal}(\bar{\mathbf Q}/L)$ is an admissible twist of rank two.
\end{lemma}
\leanref{Development/AutomorphicArithmetic/AutomorphicVarieties/CurveDVRSpecialization/ValIntLe.lean}{UpperTheory.V.Pr.kaOf}
\begin{proof}
The character $\chi_{\mathrm{cyc}}^{-1}$ is that of an algebraic Hecke character; $\chi_{\mathrm{cyc}}(csc^{-1})=\chi_{\mathrm{cyc}}(s)$, so $\chi_{\mathrm{cyc}}^{-1}(s)\chi_{\mathrm{cyc}}^{-1}(csc^{-1})=\chi_{\mathrm{cyc}}(s)^{1-3}$; and $\chi_{\mathrm{cyc}}$ is trivial on the inertia groups at the places not above $\ell$, which include those not dividing $M$ since $\ell\mid M$; take $m=1$.
For $\kappa$: $L$ contains $K$, so the restriction is algebraic, and it satisfies the relation with $c$ for $n=2$. A place not dividing $M$ divides neither $\ell$ nor $d_K$, so $\kappa$ is trivial on its inertia groups, and the condition on $\kappa^m$ restricts.
\end{proof}
% kaOf: Sd.AlgCyc for chi^-1 and Sd.AlgRes for the restriction (with exists_alg, algClo and IsAbelianGalois.tower_bot), the field conj of Sd.IsTwist, discr_over with LinkHyp.ram for the unramified places, the field ord.

\begin{lemma}\label{sha-automorphy-realization-exists-determinant}
In Situation~\ref{sha-hecke-situation}, suppose that each $\bar\psi_{v,i}$ is the restriction of $\bar\chi_i$ to $D_v$. Then the universal ordinary determinant ring exists: there are a complete local noetherian $\mathcal O$-algebra $B$ with residue field $k$ and an ordinary pair $(D^{\mathrm u},\psi^{\mathrm u})$ over $B$ such that every ordinary pair over every complete local noetherian $\mathcal O$-algebra $R$ with residue field $k$ is obtained from it by a unique $\mathcal O$-algebra map $B\to R$.
\end{lemma}
\leanref{Development/AutomorphicArithmetic/AutomorphicVarieties/WeilGroupCharacters/WeilCharactersModular/ProModularRingExists.lean}{UpperTheory.V.W.b0Ex}
\leanref{Development/AutomorphicArithmetic/AutomorphicVarieties/CommutativeAlgebraHeight/UniversalDeformationRingNonemptyOfNonemptyPair.lean}{UpperTheory.V.CA.universalDeformationRing_nonempty_of_nonempty_pair}
\leanref{Development/AutomorphicArithmetic/AutomorphicVarieties/WeilGroupCharacters/ResidualDiagonalDeterminantIsPolar.lean}{UpperTheory.V.W.residual_diagonal_determinant_isPolar}
\leanref{Development/AutomorphicArithmetic/AutomorphicVarieties/CurveDVRSpecialization/ResidualOkOfLocalCharsRestrict.lean}{UpperTheory.residual_ok_of_local_chars_restrict}
\begin{proof}
\begin{enumerate}
\item \emph{A pair over $k$.} The residual determinant $\bar D$, with the characters $\bar\psi_{v,i}$, is an ordinary pair over $k$. It is continuous, as the $\bar\chi_i$ and the $\bar\psi_{v,i}$ have open kernels, and it reduces to itself. On $\mu_v$, $\bar\psi_{v,i}$ equals the reduction of $\theta_{v,i}$, both being trivial there: for $h\in\mu_v$ some $h^{\ell^m}$ lies in the closed commutator subgroup of $D_v$, on which $\bar\psi_{v,i}$, whose kernel is closed, and $\theta_{v,i}$ are trivial, and the only element of $\ell$-power order of $k^\times$ is $1$. Condition (1) is the condition on $\bar D$ in the residual datum. For (2), $\bar D$ is the determinant of $\bar\chi_1\oplus\dots\oplus\bar\chi_n$, and the product $\prod_i\big(\operatorname{diag}(\bar\chi_1(g_i),\dots,\bar\chi_n(g_i))-\bar\psi_{v,i}(g_i)\big)$ is a diagonal matrix whose $j$-th entry has the factor $\bar\chi_j(g_j)-\bar\psi_{v,j}(g_j)=0$; so it is $0$, and so is the trace of $h$ times it. For (3), the representation is trivial at $g\in I_u\cap\bigcap_i\ker\bar\chi_i$.
\item \emph{Universal rings for the pieces.} By Lemma~\ref{sha-patching-places}, every open subgroup of $\Gamma$, hence of its open subgroup $\Gamma_0$, and every open subgroup of each $D_v$ has finitely many continuous homomorphisms to $\mathbf Z/\ell$. So, as in the proof of Lemma~\ref{sha-patching-chain-ring}: the determinant $\bar D$ of the pairwise distinct characters $\bar\chi_i$, which have open kernels, has a universal deformation ring $R_0$, a complete local noetherian $\mathcal O$-algebra with residue field $k$ carrying a continuous determinant $D_0$ that reduces to $\bar D$ \cite{Che14}; each $\bar\psi_{v,i}$ has a universal ring $R_{v,i}$ of continuous lifts equal to $\theta_{v,i}$ on $\mu_v$; and the completed tensor product $R_1$ over $\mathcal O$ of $R_0$ and the $R_{v,i}$ carries a continuous determinant $D_1$ reducing to $\bar D$ and continuous characters $\psi_{1,v,i}$ lifting $\bar\psi_{v,i}$ and equal to $\theta_{v,i}$ on $\mu_v$, from which every such family over every $R$ is obtained by a unique $\mathcal O$-algebra map $R_1\to R$.
\item \emph{The relations.} Let $I$ be the ideal of $R_1$ generated by: the differences of the coefficients of the characteristic polynomials of $D_1$ at $g^*$ and at $g$, for $g\in\Gamma_0$; the elements $\operatorname{tr}_{D_1}\big(h\prod_i(g_i-\psi_{1,v,i}(g_i))\big)$ for $v\mid\ell$, $g_1,\dots,g_n\in D_v$ and $h\in\Gamma_0$; and the elements $\operatorname{tr}_{D_1}\big(h(g-1)\big)$ for $u\in S_{\mathrm a}$, $g\in I_u\cap\bigcap_i\ker\bar\chi_i$ and $h\in\Gamma_0$. These coefficients and traces are compatible with $\mathcal O$-algebra maps, so an $\mathcal O$-algebra map $f:R_1\to R$ carries $(D_1,\psi_1)$ to an ordinary pair if and only if $f$ kills $I$.
\item \emph{The quotient.} Let $B=R_1/I$. By (1) and (3), the map $R_1\to k$ of the residual pair kills $I$, so $B$ is not zero, and it is a complete local noetherian $\mathcal O$-algebra with residue field $k$. By (3), the image $(D^{\mathrm u},\psi^{\mathrm u})$ of $(D_1,\psi_1)$ in $B$ is an ordinary pair. For an ordinary pair over $R$, the map $R_1\to R$ of (2) kills $I$ by (3), so it factors through $B$; and the factorization is unique, as $R_1\to B$ is surjective.
\end{enumerate}
\end{proof}
% b0Ex (ProModularRingExists.lean) from PhiAll (W.phiAll from Pa.CohP: phiGS from GlobalFiniteness, Det.phiSub, phiDec from LocalFiniteness), through universalDeformationRing_nonempty_of_nonempty_pair: pre_univ (univDet_mf, universalDeformationRing_localCharacters, coprod_opt from co_coprod) for (2), rels and rels_iff for (3), the quotient by Ideal.span rels with co_quot and the universal property through Ideal.Quotient.liftₐ for (4). The pair over the residue field of (1) is residual_diagonal_determinant_isPolar, from residual_ok_of_local_chars_restrict (Resid.Ok, with residualDatum_inertia_characters_trivial on mu_v); cRealAv uses it to apply ProModularRingExists.

\begin{lemma}\label{sha-automorphy-realization-exists-linked}
In Situation~\ref{sha-automorphy-link-situation}, let $n\in\{2,3\}$, and let $\mathcal O_1\supseteq\mathcal O$ be the ring of integers of a finite extension of $\mathbf Q_\ell$, with residue field $k_1$, such that the link, with its residual characters composed with the map of residue fields, satisfies the residual conditions with some constant (Definition~\ref{sha-descent-residual}). Then the universal ring $A_n$ of linked points of rank $n$ over $\mathcal O_1$ lifting the residual one exists: there are a complete local noetherian $\mathcal O_1$-algebra $A_n$ with residue field $k_1$ and a linked point $(\rho^{\mathrm u},\psi^{\mathrm u})$ of rank $n$ with values in $A_n$, lifting the residual one, such that every linked point of rank $n$ lifting the residual one, with values in a complete local noetherian $\mathcal O_1$-algebra $C$ with residue field $k_1$, is the image of $(\rho^{\mathrm u},\psi^{\mathrm u})$ under a unique $\mathcal O_1$-algebra map $A_n\to C$.
\end{lemma}
\leanref{Development/AutomorphicArithmetic/AutomorphicVarieties/FrobeniusSmoothVectors/SpaceEx.lean}{UpperTheory.V.Fr.space_ex}
\leanref{Development/AutomorphicArithmetic/AutomorphicVarieties/FrobeniusSmoothVectors/RbarProd.lean}{UpperTheory.V.Fr.rbar}
\leanref{Development/AutomorphicArithmetic/AutomorphicVarieties/CommutativeAlgebra/UnivLift.lean}{UpperTheory.V.Alg.univ_lift}
\leanref{Development/AutomorphicArithmetic/AutomorphicVarieties/AlgebraTools/MonomialRepresentations/HenselDiag.lean}{UpperTheory.V.AlgebraTools.henselDiag}
\begin{proof}
View $\bar\Psi_\Gamma$ with values in $k_1$, and write $g^*=c_\Gamma g^{-1}c_\Gamma^{-1}$ and $w_i=w_{n,i}$.
\begin{enumerate}
\item[(i)] \emph{The residual point.} Let $\rho_0(g)=\operatorname{diag}\big(\bar\Psi_\Gamma(g)^{w_1},\dots,\bar\Psi_\Gamma(g)^{w_n}\big)$ for $g\in\Gamma_0$ and $\bar\psi_{v,i}=\bar\Psi_\Gamma^{w_i}\circ\mathrm{loc}_v$. This is a linked point with values in $k_1$. For (a): $\bar\Psi_\Gamma(g^*)=\bar\Psi_\Gamma(g)^{-1}$, since $\bar\Psi_\Gamma(c_\Gamma)=1$ by residual condition (4), and $w_{n+1-i}=-w_i$. For (b): $\rho_0(g_0)-\rho_0(g_0^*)$ is diagonal. For (c): the product is a diagonal matrix whose $j$-th entry has the factor $\bar\Psi_\Gamma(\mathrm{loc}_vh_j)^{w_j}-\bar\psi_{v,j}(h_j)=0$. The map $\rho_0$ is locally constant, as $\ker\bar\Psi_\Gamma$ is open by residual condition (4).
\item[(ii)] \emph{Lifts.} The group $\Gamma_0$ is open in the profinite group $\Gamma$, so it is compact, and by residual condition (5) each of its open subgroups has finitely many continuous homomorphisms to $\mathbf Z/\ell$. So the continuous homomorphisms of monoids $\Gamma_0\to M_n(R)$ lifting $\rho_0$, for complete local noetherian $\mathcal O_1$-algebras $R$ with residue field $k_1$, are represented by such an algebra $R$ with a universal lift $u$ \cite{Maz89}.
\item[(iii)] \emph{Diagonalization at $v$.} By residual condition (7), for each $v$ there is $g_v\in G_v$ such that $a_v=\bar\Psi_\Gamma(\mathrm{loc}_vg_v)$ has $a_v^2\ne1$. The residues $a_v^{w_i}$ are pairwise distinct, since $w_i-w_j$ divides $2$ for $i\ne j$. So $u(\mathrm{loc}_vg_v)$ has its entries off the diagonal in $\mathfrak m_R$ and diagonal entries with pairwise distinct residues, and by Hensel's lemma there are matrices $P_v$ and $Q_v$ over $R$, congruent to $1$, with $P_vQ_v=1$ and $Q_vu(\mathrm{loc}_vg_v)P_v$ diagonal. Let $T_v(s)=Q_vu(\mathrm{loc}_vs)P_v$ for $s\in G_v$; it reduces to $\rho_0(\mathrm{loc}_vs)$.
\item[(iv)] \emph{The quotient.} Let $I$ be the ideal of $R$ generated by the entries of $u(g^*)-J_n\,{}^tu(g)J_n$ for $g\in\Gamma_0$, the entries off the diagonal of $u(g_0)-u(g_0^*)$, and the entries $(i,j)$ with $j<i$ of the matrices $T_v(s)$. The map $R\to k_1$ of $\rho_0$ kills $I$, by (i) and since each $T_v(s)$ reduces to a diagonal matrix; so $A_n=R/I$ is not zero, and it is a complete local noetherian $\mathcal O_1$-algebra with residue field $k_1$. Let $\rho^{\mathrm u}$ be the image of $u$ in $A_n$. Over $A_n$ each $T_v(s)$ is upper triangular, and $T_v(ss')=T_v(s)T_v(s')$; so $\psi^{\mathrm u}_{v,i}(s)=T_v(s)_{ii}$ defines characters of $G_v$, and the products of (c) for $\rho^{\mathrm u}$ vanish, since after conjugation by $Q_v$ the $i$-th factor maps the span of the first $i$ basis vectors into that of the first $i-1$. The $\psi^{\mathrm u}_{v,i}$ are continuous, being entries of a conjugate of the continuous $u\circ\mathrm{loc}_v$, and they reduce to the $\bar\psi_{v,i}$. So $(\rho^{\mathrm u},\psi^{\mathrm u})$ is a linked point lifting the residual one.
\item[(v)] \emph{Universality.} Let $(\rho',\psi')$ be a linked point with values in $C$ lifting the residual one. As $\rho'$ is a continuous lift of $\rho_0$, $\rho'=f\circ u$ for a unique $\mathcal O_1$-algebra map $f:R\to C$, and $f$ kills the generators of $I$ of the first two kinds by conditions (a) and (b) for $\rho'$. Fix $v$, and write $T'(s)=f(T_v(s))$ and $T'(g_v)=\operatorname{diag}(d_1,\dots,d_n)$. Then $d_j$ reduces to $a_v^{w_j}$ and $\psi'_{v,i}(g_v)$ to $a_v^{w_i}$, so the diagonal matrix $D_i=T'(g_v)-\psi'_{v,i}(g_v)$ has a unit as its $j$-th entry for every $j\ne i$. Conjugating condition (c) for $\rho'$ by $f(Q_v)$ gives $\prod_i\big(T'(h_i)-\psi'_{v,i}(h_i)\big)=0$ for $h_1,\dots,h_n\in G_v$. Take $h_i=g_v$ for $i\ne a$ and $h_a=s$. For $b\le a$, the entry $(a,b)$ of the product is the product of the entry $(a,a)$ of $D_1\cdots D_{a-1}$, the entry $(a,b)$ of $T'(s)-\psi'_{v,a}(s)$ and the entry $(b,b)$ of $D_{a+1}\cdots D_n$; the first and the last are units, as $i\ne a$ for $i<a$ and $i\ne b$ for $i>a\ge b$. So $T'(s)$ is upper triangular with diagonal entries $\psi'_{v,a}(s)$. Hence $f$ kills $I$ and factors through $A_n$ by a map that carries $(\rho^{\mathrm u},\psi^{\mathrm u})$ to $(\rho',\psi')$; it is unique, as $f$ is unique and $R\to A_n$ is surjective.
\end{enumerate}
\end{proof}
% space_ex (SpaceEx.lean): rbar (RbarProd.lean) for (i), with hc for local constancy; univ_lift with Det.phiSub from Ok.del_open and Ok.phi for (ii); Ok.f3, w_inj and diag_ex (henselDiag, the Henselian property of a complete local ring) for (iii); the set S, co_quot, ord_of_tri (with ord_of_flag) and cont_conj for (iv); tri_diag (by cases on n = 2, 3, which the text does for all n at once), ord_conj and Ideal.Quotient.liftₐ for (v).

\begin{lemma}\label{sha-automorphy-realization-exists}
In Situation~\ref{sha-automorphy-link-situation}:
\begin{enumerate}
\item for every ring of integers $\mathcal O'\supseteq\mathcal O$ of a finite extension of $\mathbf Q_\ell$ in $\bar{\mathbf Q}_\ell$ and all continuous characters $\kappa,\kappa'$ of $\operatorname{Gal}(\bar{\mathbf Q}/L)$ with values in $\bar{\mathbf Q}_\ell^\times$, there is the ring of integers $\mathcal O_1\supseteq\mathcal O'$ of a finite extension, large enough (it contains the images of all embeddings of $\mathcal O_L$), that contains the values of $\kappa$ and $\kappa'$; the inclusion $\mathcal O'\subseteq\mathcal O_1$ induces a map of residue fields;
\item let $n\in\{2,3\}$, let $\mathcal O_1\supseteq\mathcal O$ be the ring of integers of a finite extension of $\mathbf Q_\ell$, with the map of residue fields induced by the inclusion, and let $\kappa$ be an admissible twist of rank $n$ with values in $\mathcal O_1$; then for every integer $p_0\ne0$ there is a realization of rank $n$ of the link, twisted by $\kappa$, with coefficients $\mathcal O_1$, whose place $v_0$ has residue characteristic prime to $p_0$ (Definition~\ref{sha-automorphy-realization}).
\end{enumerate}
\end{lemma}
\leanref{Development/AutomorphicArithmetic/AutomorphicVarieties/CurveDVRSpecialization/CoefEx.lean}{UpperTheory.V.Pr.coefEx}
\leanref{Development/AutomorphicArithmetic/Foundations/GaloisRepresentations/PlainTorsionPairs/CRealAv.lean}{UpperTheory.Gal.projectedFormsSystemPart.cRealAv}
\leanref{Development/AutomorphicArithmetic/AutomorphicVarieties/CurveDVRSpecialization/EveryPointSatisfiesFrameCondition.lean}{UpperTheory.V.Pr.CRealEx}
\begin{proof}
(1) The group $\operatorname{Gal}(\bar{\mathbf Q}/L)$ is compact, so by the Baire category theorem the values of $\kappa$ lie in a finite extension of $\mathbf Q_\ell$; they form a compact, hence bounded, subgroup of $\bar{\mathbf Q}_\ell^\times$, so they have absolute value at most $1$. The same holds for $\kappa'$. The ring $\mathcal O'$ lies in a finite extension, and so do the images of $\mathcal O_L$ under its finitely many embeddings. Let $\mathcal O_1$ be the ring of integers of the compositum of these four finite extensions; it contains $\mathcal O'$, the images of $\mathcal O_L$ and the values of $\kappa$ and $\kappa'$, since these are integral over $\mathbf Z_\ell$. The inclusion $\mathcal O'\to\mathcal O_1$ is injective and integral, hence local, and so induces a map of residue fields.
(2) Write $G_{L^+}=\operatorname{Gal}(\bar{\mathbf Q}/L^+)$ and $G_L=\operatorname{Gal}(\bar{\mathbf Q}/L)$, and let $\Gamma$, $p$, $\Gamma_0$, $c_\Gamma$, $g_0$, $A_v$ and $G_v$ be those of the link (Definition~\ref{sha-descent-link}(1)), so that $p(G_L)=\Gamma_0$.
\begin{enumerate}
\item[(a)] \emph{An element of $G_L$.} By the residual conditions (Definition~\ref{sha-descent-residual}(7)), $\omega_\Gamma$ is not a power of $\bar\Psi_\Gamma$ on $\Gamma_0$. By the duality of finite abelian groups, applied to the image of $\Gamma_0$ under $(\bar\Psi_\Gamma,\omega_\Gamma)$, some $g\in\Gamma_0$ has $\bar\Psi_\Gamma(g)=1$ and $\omega_\Gamma(g)\ne1$. Let $g'\in G_L$ with $p(g')=g$; then $\bar\Psi(g')=1$ and $\chi_{\mathrm{cyc}}(g')\not\equiv1\pmod\ell$.
\item[(b)] \emph{The place $v_0$.} Let $U\subseteq G_{L^+}$ be the open normal subgroup of the elements of $G_L$ at which $\bar\Psi$ and $\chi_{\mathrm{cyc}}\bmod\ell$ are trivial. By the Chebotarev density theorem (Theorem~\ref{lit-chebotarev-degree-one}), applied to the fixed field of $U$ over $L^+$, there are a place $v_0$ of $L^+$ of degree one over a prime $q$ not dividing $2Mp_0|d_L|$, a place $A_0$ of $\bar{\mathbf Q}$ above $v_0$ with $I_{A_0}\cap G_{L^+}\subseteq U$, and a Frobenius element $\sigma\in D_{A_0}\cap G_{L^+}$ with $\sigma\in g'U$. The group $D_{A_0}\cap G_{L^+}$ is the closure of the subgroup generated by $\sigma$ and its inertia group, so it lies in the closed subgroup of the elements of $G_L$ at which $\bar\Psi$ is trivial. And $\chi_{\mathrm{cyc}}(\sigma)\equiv\chi_{\mathrm{cyc}}(g')\not\equiv1\pmod\ell$. The prime $q$ is odd, different from $\ell$ since $\ell\mid M$, unramified in $L$, and prime to $p_0$.
\item[(c)] \emph{The data of Situation~\ref{sha-hecke-situation}.} Take for $S$ the places of $L^+$ above $\ell$, the set $S_{\mathrm a}$ of places dividing $M$ but not $\ell$, and $v_0$; for $v\mid\ell$ let $\tilde v$ be the place of $L$ below $A_v$, which gives a bijection with the places of $L^+$ above $\ell$ by Definition~\ref{sha-descent-link}(1)(iv), and let $\tilde v_0$ be the place below $A_0$; let $c$ be the complex conjugation of the link, and $\mathcal O_1$ the coefficient ring. These are data of that Situation: $\mathcal O_1$ contains the images of the embeddings of $\mathcal O_L$; $L/L^+$ is unramified at every finite place, as $I_A\cap G_{L^+}\subseteq\ker\eta$; the places of $S$ split in $L$, as $D_A\cap G_{L^+}\subseteq\ker\eta$ for $A$ above a prime dividing $M$, and by (b) for $v_0$; and $L$ contains no primitive $\ell$-th root of unity, since $\omega$ is nontrivial on each $G_v\subseteq G_L$. Let $G_0$ be the group of that Situation that fixes $L$, and $\mathrm{pr}:G_L\to G_0$ the quotient map; its kernel is the closed normal subgroup generated by the inertia groups at the places of $L$ above places outside $S$.
\item[(d)] \emph{The map $\pi$.} Let $A$ be a place of $\bar{\mathbf Q}$ above a place of $L^+$ outside $S$. It is not above $\ell$ and does not divide $M$; so $\eta$ is trivial on $I_A\cap G_{L^+}$, and so is $\Psi$, which is unramified outside $\ell N$. By Definition~\ref{sha-descent-link}(1)(i), $I_A\cap G_{L^+}$ lies in the kernel of $p$. So $p$ factors through the Galois group of the maximal extension of $L^+$ unramified outside $S$ and the infinite places, and gives a continuous $\pi:G_0\to\Gamma_0$ with $\pi\circ\mathrm{pr}=p$ on $G_L$. It is surjective, as $p(G_L)=\Gamma_0$; it carries conjugation by $c$ to conjugation by $c_\Gamma=p(c)$; and $\varepsilon_\Gamma\circ\pi=\varepsilon$, since $\varepsilon_\Gamma\circ p=\chi_{\mathrm{cyc}}$.
\item[(e)] \emph{The character $\kappa_0$.} By Definition~\ref{sha-automorphy-realization}(1), $\kappa$ is trivial on the inertia groups at the places not dividing $M$, hence on the kernel of $\mathrm{pr}$ by (c), as the places outside $S$ do not divide $M$. So $\kappa=\kappa_0\circ\mathrm{pr}$ for a character $\kappa_0:G_0\to\mathcal O_1^\times$, which is continuous since $\mathrm{pr}$ is a quotient map, and $\kappa_0(g)\kappa_0(cgc^{-1})=\varepsilon(g)^{1-n}$ by the identity for $\kappa$. It is trivial on $I_{v_0}$, as $v_0$ does not divide $M$. Let $m'$ be prime to $\ell$ with $\kappa^{m'}$ trivial on the inertia groups at the places not above $\ell$, and let $m=(|\mathbf F|-1)m'$, which is prime to $\ell$. For $u\in S_{\mathrm a}$ and $g\in I_u$, write $g=\mathrm{pr}(s)$ with $s$ in the inertia group of $G_L$ at a place $A$ above $\tilde u$, which is not above $\ell$. Then $s^{|\mathbf F|-1}$ lies in $I_A$, with $\eta(s^{|\mathbf F|-1})=1$ and $\bar\Psi(s^{|\mathbf F|-1})=1$, so it lies in the kernel of $p$ by Definition~\ref{sha-descent-link}(1)(i); hence $\pi(g)^m=p(s)^m=1$, and $\kappa_0(g)^m=\kappa(s)^m=1$.
\item[(f)] \emph{The residual datum.} Let $\bar\chi_i(g)=\bar\Psi_\Gamma(\pi(g))^{e_{n,i}}\bar\kappa_0(g)$, with $\bar\Psi_\Gamma$ composed with the map of residue fields and $\bar\kappa_0$ the reduction of $\kappa_0$. These characters have open kernels. They are pairwise distinct: at $g$ with $\pi(g)=g_0$, $\bar\chi_i=\bar\chi_j$ would give $\bar\Psi_\Gamma(g_0)^{e_{n,i}-e_{n,j}}=1$, while $e_{n,i}-e_{n,j}\in\{\pm1,\pm2\}$ and $\bar\Psi_\Gamma(g_0)^2\ne1$. Since $e_{n,n+1-i}=-e_{n,i}$, (d) and (e) give $\bar\varepsilon(g)^{1-n}\bar\chi_i(cgc^{-1})^{-1}=\bar\chi_{n+1-i}(g)$; so $g\mapsto g^*$ permutes the factors of $\bar D=\prod_i\bar\chi_i$, and $\bar D$ takes the same value at $y$ and $y^*$. With $\bar\psi_{v,i}$ the restriction of $\bar\chi_i$ to $D_v$ and $\theta_{v,i}=1$, this is a residual datum of rank $n$, whose residual characters are those of the link composed with $\pi$ and twisted by $\kappa_0$.
\item[(g)] \emph{Decomposition groups.} For $v\mid\ell$, the decomposition group $D_v$ of $G_L$ at $A_v$ equals $G_v=D_{A_v}\cap G_{L^+}$, which lies in $G_L$ since $\ell\mid M$. This identifies $D_v$ with $G_v$, compatibly with $\mathrm{loc}_v$ and $\pi$ and with the inertia groups.
\item[(h)] \emph{The place $v_0$.} The characters $\bar\chi_i$ are trivial on $I_{v_0}$, since $\bar\Psi_\Gamma\circ\pi$ is trivial on $D_{v_0}$ by (b) and $\kappa$ is trivial on $I_{v_0}$, as $v_0\nmid M$. They agree on $D_{v_0}$, where each equals $\bar\kappa_0$. And $\varepsilon\bmod\ell$ is nontrivial on $D_{v_0}$, at the image of $\sigma$.
\item[(i)] \emph{The constant.} The link, with its residual characters composed with the injective map of residue fields, satisfies the residual conditions with the constant of the Situation, and is large for it.
\item[(j)] \emph{The groups $\mu_v$.} Let $v\mid\ell$. By (g), $D_v=G_v$, the decomposition group of $G_L$ at $A_v$, which is the absolute Galois group of the completion $L_{\tilde v}$; and $\omega$ is nontrivial on it (Definition~\ref{sha-descent-link}(1)(iv)), so $L_{\tilde v}$ contains no primitive $\ell$-th root of unity. By local class field theory \cite[Chapter~VII]{NSW08}, the quotient of $D_v$ by the closure of its commutator subgroup is the profinite completion of $L_{\tilde v}^\times$, whose elements of $\ell$-power order are the $\ell$-power roots of unity of $L_{\tilde v}$. So $1$ is its only element of $\ell$-power order, and $\mu_v$ lies in the closure of the commutator subgroup of $D_v$.
\item[(k)] \emph{The kernel of $\pi$.} Let $G_{L^+,S}$ be the Galois group of the maximal extension of $L^+$ unramified outside $S$ and the infinite places, and $q:G_{L^+}\to G_{L^+,S}$ the quotient map; $G_0$ is the subgroup of index two that fixes $L$, $q=\mathrm{pr}$ on $G_L$, and $G_{L^+,S}=G_0\cup q(c)G_0$. Let $N$ be a closed normal subgroup of $G_0$, stable under $g\mapsto cgc^{-1}$, that contains $I_{v_0}$ and every $g\in I_u$, for $u\in S_{\mathrm a}$, with $\bar\Psi_\Gamma(\pi(g))=1$. Then $N$ is normal in $G_{L^+,S}$, and closed, as $G_0$ is open.
Let $A$ be a place of $\bar{\mathbf Q}$ not above $\ell$, $w$ the place of $L^+$ below it, and $s\in I_A\cap G_{L^+}$ with $\bar\Psi(s)=1$; we show $q(s)\in N$. If $w\notin S$, then $q(s)=1$. Otherwise $w=v_0$ or $w\in S_{\mathrm a}$, as $w\nmid\ell$. Let $I_w\subseteq G_0$ be the inertia group of the Situation at $w$, the image under $\mathrm{pr}$ of $I_{A_w}\cap G_L$ for a place $A_w$ of $\bar{\mathbf Q}$ above $\tilde w$, and let $h\in G_{L^+}$ with $hA_w=A$. Then $t=h^{-1}sh$ lies in $I_{A_w}\cap G_{L^+}$, which lies in $G_L$ by Definition~\ref{sha-descent-link}(1); so $\mathrm{pr}(t)\in I_w$. If $w=v_0$, then $\mathrm{pr}(t)\in N$. If $w\in S_{\mathrm a}$, then $\bar\Psi_\Gamma(\pi(\mathrm{pr}(t)))=\bar\Psi_\Gamma(p(t))=\bar\Psi(t)=\bar\Psi(s)=1$ by (d), so again $\mathrm{pr}(t)\in N$. As $N$ is normal in $G_{L^+,S}$, $q(s)=q(h)\,\mathrm{pr}(t)\,q(h)^{-1}\in N$.
So the closed normal subgroup $q^{-1}(N)$ of $G_{L^+}$ contains the generators of $\ker p$ in Definition~\ref{sha-descent-link}(1)(i), hence $\ker p$. Let $g\in\ker\pi$, and $s\in G_L$ with $\mathrm{pr}(s)=g$. Then $p(s)=\pi(g)=1$ by (d), so $s\in\ker p$ and $g=q(s)\in N$.
\end{enumerate}
These are the data of a realization of rank $n$, twisted by $\kappa$, with coefficients $\mathcal O_1$, other than its universal rings $B_n$ and $A_n$; its place $v_0$ has residue characteristic $q$, prime to $p_0$.
The universal ordinary determinant ring $B_n$ of the residual datum of (f) exists by Lemma~\ref{sha-automorphy-realization-exists-determinant}, since $\bar\psi_{v,i}$ is the restriction of $\bar\chi_i$ to $D_v$. The universal ring $A_n$ of linked points of rank $n$ over $\mathcal O_1$ lifting the residual one exists by Lemma~\ref{sha-automorphy-realization-exists-linked}, by (i).
% (2) is cRealAv (CRealAv.lean): (a) is the private exists_kernel_element_nontrivial_cyclotomic of CRealEx.lean (Loc.v0e and Ok.f5); (b) is frame_criterion_decomposition_groups (Loc.v00, Kum.decGen) with M' = 2 M p_0 |disc L|; (c) is cr_set, with frame_criterion_galois_equivariance for the bijection e, cr_split, and the level V from the hypothesis W.LevEx; (d) is frame_criterion_projection; (e) is frame_criterion_level_kernel (ker_pr_le) for kappa_0 with kappa_cont and kappa_conj, then hram (hka.unr at v0M) for the TwLink field ramκ and htame for its field tame (m = card (F^x) * m_1 with hcard and hka.ord; cr_tame puts s^card(F^x) in ker p through k.p_ker); (f) is frame_criterion_residual; (g) is frame_criterion_lambda_characters; (h) is hv0 (Resid.V0ok); (i) is Fr.ok_ext. (j) is the TwLink field mu, cr_mu of CRealEx.lean (Loc.muNone, MuNone.lean, proved in the Lean from local class field theory; the element with cyclotomic character not 1 mod l is hmu, from the field mu of the link). (k) is the TwLink field ker, frame_criterion_kernel (FrameCriterionKernel.lean: normal_map for the normality in V.frame.Gamma, conj_ine and s.unr for t, ine_NS for w outside S, k.p_ker and pr_surj). The universal rings: B_0 from the hypothesis W.ProModularRingExists (W.b0Ex from W.phiAll, Lemma sha-automorphy-realization-exists-determinant), applied to the pair of residual_diagonal_determinant_isPolar and residual_ok_of_local_chars_restrict, and S from Fr.space_ex (Lemma sha-automorphy-realization-exists-linked).
% (1) is Pr.coefEx (CoefEx.lean): ce_char (CompactRepField.exists_finiteDimensional_entries_mem and boundedness), ce_coef, ce_large, X.So.coefOf and ResidueField.map. (2) is cRealAv (from W.LevEx and W.ProModularRingExists; cr_set chooses the setting with v0 prime to p0), and Pr.cRealEx without the condition on p0. Size: CRealAv.lean (about 330 lines) and Foundations/GaloisRepresentations/PlainTorsionPairs (about 970 lines in all).
\end{proof}

\begin{lemma}\label{sha-automorphy-realization-seed-normal}
Let $K$ be an algebraically closed field in which $2\ne0$, let $n\in\{2,3\}$, and let $(e_{n,i})$ and $J_n$ be as in Definition~\ref{sha-descent-link}(2). Let $Q\in GL_n(K)$ be symmetric, let $D\in M_n(K)$ satisfy ${}^tDQ=-QD$, and let $\zeta\in K$ be nonzero with $\operatorname{tr}(D^2)=2\zeta^2$. Then there is $P\in GL_n(K)$ with ${}^tPQP=J_n$ and $DP=P\operatorname{diag}(e_{n,1}\zeta,\dots,e_{n,n}\zeta)$.
\end{lemma}
\leanref{Development/AutomorphicArithmetic/AutomorphicVarieties/FrobeniusSmoothVectors/OfForm.lean}{skew_eig}
\leanref{Development/AutomorphicArithmetic/AutomorphicVarieties/FrobeniusSmoothVectors/OfForm.lean}{skew_norm}
\begin{proof}
\emph{Eigenvalues.} From ${}^tD=-QDQ^{-1}$ follows $\operatorname{tr}D=-\operatorname{tr}D$, so $\operatorname{tr}D=0$; and if $n=3$, $\det D=-\det D$, so $\det D=0$. Expressing the characteristic polynomial through traces, $2\det(X-D)$ is $2X^2-\operatorname{tr}(D^2)=2(X^2-\zeta^2)$ if $n=2$, and $2X^3-\operatorname{tr}(D^2)X=2X(X^2-\zeta^2)$ if $n=3$. So $\mu_i=e_{n,i}\zeta$ is an eigenvalue of $D$ for each $i$.

\emph{Eigenvectors.} The $\mu_i$ are pairwise distinct, as $\zeta\ne0$ and $2\ne0$. Let $V$ be the matrix whose $i$-th column is an eigenvector of $D$ for $\mu_i$; it is invertible, and $DV=V\operatorname{diag}(\mu)$. The matrix $M={}^tVQV$ is symmetric and invertible, and $\operatorname{diag}(\mu)M={}^t(DV)QV={}^tV\,{}^tDQV=-{}^tVQDV=-M\operatorname{diag}(\mu)$. So $(\mu_i+\mu_j)M_{ij}=0$, and $\mu_i+\mu_j=0$ only for $j=n+1-i$. Hence $M$ is antidiagonal, with nonzero antidiagonal entries since it is invertible.

\emph{Scaling.} Let $c=(1,M_{12}^{-1})$ if $n=2$, and $c=(1,y,M_{13}^{-1})$ with $y^2=M_{22}^{-1}$ if $n=3$. By the symmetry of $M$, $c_ic_{n+1-i}M_{i,n+1-i}=1$ for every $i$. Let $P=V\operatorname{diag}(c)$. Then ${}^tPQP=\operatorname{diag}(c)M\operatorname{diag}(c)=J_n$ and $DP=V\operatorname{diag}(\mu)\operatorname{diag}(c)=P\operatorname{diag}(\mu)$.
\end{proof}
% skew_eig (the characteristic polynomials through det_fin_two and det_fin_three) and skew_norm (eigenvectors_linearIndependent', the antidiagonal M and the scaling c) in OfForm.lean.

\begin{lemma}\label{sha-automorphy-realization-seed-point}
In Situation~\ref{sha-automorphy-link-situation}, let $n\in\{2,3\}$ and let a realization of rank $n$ with coefficients $\mathcal O_1$ be given (Definition~\ref{sha-automorphy-realization}). Let $G_0$ be the Galois group of its residual datum, with its decomposition groups $D_v$ above $\ell$, which map to $G_0$ and are identified with the groups $G_v$ of the link compatibly with $\mathrm{loc}_v$; let $\pi:G_0\to\Gamma_0$ and $\kappa_0:G_0\to\mathcal O_1^\times$ be its map and its character, and $(D^{\mathrm u},\psi^{\mathrm u})$ the universal ordinary pair over $B_n$. Let $K$ be an algebraically closed field that is an $\mathcal O_1$-algebra, $x:B_n\to K$ an $\mathcal O_1$-algebra map, $r:\Gamma_0\to GL_n(K)$ a homomorphism such that $0$ and $K^n$ are the only $r(\Gamma_0)$-stable subspaces of $K^n$, and $P\in GL_n(K)$ a symmetric matrix with ${}^tr(g)Pr(c_\Gamma gc_\Gamma^{-1})=P$ for $g\in\Gamma_0$. Let $\rho(g)=\kappa_0(g)r(\pi(g))$ for $g\in G_0$, and let $\psi_{v,i}:D_v\to K^\times$ be characters such that $x$ carries $D^{\mathrm u}$ to the determinant of $\rho$ and each $\psi^{\mathrm u}_{v,i}$ to $\psi_{v,i}$, and such that $\rho$ is conjugate on each $D_v$ to an upper triangular representation with diagonal characters $\psi_{v,1},\dots,\psi_{v,n}$. Then there are a homomorphism $\rho':\Gamma_0\to M_n(K)$ and characters $\psi'_{v,i}:G_v\to K^\times$ satisfying conditions (a), (b) and (c) of Definition~\ref{sha-descent-link}(2), such that $0$ and $K^n$ are the only $\rho'(\Gamma_0)$-stable subspaces of $K^n$ and $x(D^{\mathrm u})$ is the determinant of $g\mapsto\kappa_0(g)\rho'(\pi(g))$.
\end{lemma}
\leanref{Development/AutomorphicArithmetic/AutomorphicVarieties/CurveDVRSpecialization/UniversalMapsPointOfPolarizedFlaggedRep.lean}{UpperTheory.universalMaps_point_of_polarized_flagged_rep}
\leanref{Development/AutomorphicArithmetic/AutomorphicVarieties/FrobeniusSmoothVectors/OfForm.lean}{UpperTheory.V.Fr.of_form}
\begin{proof}
Let $C=x(B_n)\subseteq K$, a quotient of the complete noetherian local $\mathcal O_1$-algebra $B_n$, hence one too, with the same residue field; the universal pair gives an ordinary pair $(D_C,\psi_C)$ over $C$. For $g\in\Gamma_0$ write $g^*=c_\Gamma g^{-1}c_\Gamma^{-1}$. By the residual conditions (Definition~\ref{sha-descent-residual}(4)), $c_\Gamma^2=1$ and $\bar\Psi_\Gamma(c_\Gamma)=1$; so $(g^*)^*=g$, $(gh)^*=h^*g^*$ and $\bar\Psi_\Gamma(g^*)=\bar\Psi_\Gamma(g)^{-1}$.
\begin{enumerate}
\item \emph{Traces.} For $g\in G_0$ let $\tau(g)=\kappa_0(g)^{-1}\operatorname{tr}_{D_C}(g)\in C$. In $K$, $\tau(g)=\operatorname{tr}r(\pi(g))$, since $x$ carries $D^{\mathrm u}$ to the determinant of $\rho$. As $\pi$ is surjective and $C\to K$ is injective, there is $t:\Gamma_0\to C$ with $t(\pi(g))=\tau(g)$, and $t(h)=\operatorname{tr}r(h)$ in $K$. Modulo $\mathfrak m_C$, $D_C$ reduces to the determinant of the residual characters, which are the characters $(\bar\Psi_\Gamma\circ\pi)^{e_{n,i}}$ of the link twisted by the reduction of $\kappa_0$; so $t(h)\equiv\sum_i\bar\Psi_\Gamma(h)^{e_{n,i}}$.
\item \emph{Ordered products.} Let $\psi'_{v,i}=\psi_{v,i}\kappa_0^{-1}$ on $D_v=G_v$. In a basis in which $\rho$ is upper triangular on $D_v$, the $i$-th factor of $(\rho(h_1)-\psi_{v,1}(h_1))\cdots(\rho(h_n)-\psi_{v,n}(h_n))$, for $h_1,\dots,h_n\in D_v$, maps the span of the first $i$ basis vectors into that of the first $i-1$; so the product is $0$. Since $\rho(h)=\kappa_0(h)r(\mathrm{loc}_vh)$ and $\psi_{v,i}(h)=\kappa_0(h)\psi'_{v,i}(h)$, dividing by the units $\kappa_0(h_i)$ gives $(r(\mathrm{loc}_vh_1)-\psi'_{v,1}(h_1))\cdots(r(\mathrm{loc}_vh_n)-\psi'_{v,n}(h_n))=0$.
\item \emph{A skew matrix.} The identity for $P$ at $g^{-1}$ gives $Pr(g^*)={}^tr(g)P$ for $g\in\Gamma_0$; so $\operatorname{tr}r(g^*)=\operatorname{tr}r(g)$. Let $A=r(g_0)$, $A^*=r(g_0^*)$ and $D=A-A^*$. Then $PA^*={}^tAP$ and $PA={}^tA^*P$, so ${}^tDP=-PD$. Also $\operatorname{tr}(A^{*2})=\operatorname{tr}r((g_0^2)^*)=\operatorname{tr}(A^2)$, so $\operatorname{tr}(D^2)=2\bigl(\operatorname{tr}r(g_0^2)-\operatorname{tr}r(g_0g_0^*)\bigr)$.
\item \emph{A square root.} Let $a=\bar\Psi_\Gamma(g_0)$ and $z=a-a^{-1}$, which is nonzero since $a^2\ne1$ by the residual conditions. As $\bar\Psi_\Gamma(g_0g_0^*)=1$, (1) shows that $X=t(g_0^2)-t(g_0g_0^*)\in C$ is congruent to $\sum_ia^{2e_{n,i}}-n=(a-a^{-1})^2=z^2$. Since $\ell$ is odd, $2$ is a unit of the complete local ring $C$, so by Hensel's lemma $X=\zeta^2$ for some $\zeta\in C$ congruent to $z$. Then $\zeta\ne0$, and $\operatorname{tr}(D^2)=2\zeta^2$ in $K$ by (1) and (3).
\item \emph{The linked point.} Lemma~\ref{sha-automorphy-realization-seed-normal}, with $Q=P$, gives $P'\in GL_n(K)$ with ${}^tP'PP'=J_n$ and $DP'=P'\operatorname{diag}(e_{n,i}\zeta)$. Let $\rho'(g)=P'^{-1}r(g)P'$. From $PP'={}^tP'^{-1}J_n$ and $J_n^2=1$ follows $\rho'(g^*)=P'^{-1}P^{-1}\,{}^tr(g)PP'=J_n\,{}^t\rho'(g)J_n$, which is (a). The matrix $\rho'(g_0)-\rho'(g_0^*)=P'^{-1}DP'$ is diagonal, which is (b). The products of (c) for $(\rho',\psi')$ are the conjugates by $P'$ of those of (2), so they vanish.
\item \emph{Conclusion.} A $\rho'(\Gamma_0)$-stable subspace $W$ gives the $r(\Gamma_0)$-stable subspace $P'W$, so $W$ is $0$ or $K^n$. The representation $g\mapsto\kappa_0(g)\rho'(\pi(g))$ is conjugate to $\rho$ by $P'$, so it has the same determinant, which is $x(D^{\mathrm u})$.
\end{enumerate}
\end{proof}
% universalMaps_point_of_polarized_flagged_rep: C = x.range (co_quot) with the pair of exists_determinantPair_map_of_algHom; then the private main: tt, tt_map and tt_res for (1); chi_ex, flagRel and prod_smul for (2); of_form for (3) to (5) (hst, hskew, hDD; hresX, co_sqrt; skew_eig, skew_norm, hpol, hdm, ordProd_conj); irr_conj and repApi for (6). The data Integral that of_form also gives (continuity of t and of the characters over C, and zeta) are not used here.

\begin{lemma}\label{sha-automorphy-realization-seed}
In Situation~\ref{sha-automorphy-link-situation}, let $(j,\chi',\theta)$ be a seed datum for $L$, and let a realization of rank $n$, twisted by an admissible $\kappa$, be given. Let $m\ge2$, and let $\psi_i$ and $\tau_i$ be the restrictions of $(\chi_{\mathrm{cyc}}^j\chi')^{w_{n,i}}$ and $\chi'^{w_{n,i}}$ to $\operatorname{Gal}(\bar{\mathbf Q}/L)$. Let $r:\operatorname{Gal}(\bar{\mathbf Q}/L)\to GL_n(\bar{\mathbf Q}_\ell)$ be a \emph{seed} for $\kappa$ with characters $\psi_i$, inertial type $(\tau_i)$ and weights $(mw_{n,i})$: a continuous irreducible representation such that (i) there is a symmetric invertible $P$ with ${}^tr(s)Pr(csc^{-1})=P$ for all $s$; (ii) the coefficients of the characteristic polynomial of each $r(s)$ are congruent to those of $\prod_i(X-\psi_i(s))$; (iii) at every place $A$ not above $\ell$, $r$ is conjugate on $I_A$ to the diagonal representation with entries $\tau_i$; (iv) at every place $A$ above $\ell$, $r$ is conjugate on $D_A$ to an upper triangular representation with diagonal characters $\kappa_{A,i}$ congruent to $\psi_i$ and equal to $\chi_{\mathrm{cyc}}^{mw_{n,i}}$ on $I_A$; and (v) $r\otimes\kappa$ is attached through $\iota$ to a regular algebraic cuspidal automorphic representation $\pi$ of $GL_n(\mathbf A_L)$ with $\pi^\vee\cong\pi^c$ and is windowed. Such seeds are given by (a) for $n=2$ and by (b) for $n=3$ in the proof of Theorem~\ref{sha-residually-reducible-automorphy}. Let $\pi$ and $\kappa_0$ be the map and the character of the realization (Definition~\ref{sha-automorphy-realization}(2)), and $D^{\mathrm u}$ the universal determinant over $B_n$. Then for every ideal $J$ of $B_n$ with the property of automorphic points of Proposition~\ref{sha-hecke-ideal}, there are an algebraically closed field $F$ that is an $\mathcal O_1$-algebra, an $\mathcal O_1$-algebra map $x:B_n\to F$ that kills $J$, a homomorphism of monoids $\rho':\Gamma_0\to M_n(F)$ and homomorphisms $\psi'_{v,i}:G_v\to F^\times$ that satisfy conditions (a), (b) and (c) of Definition~\ref{sha-descent-link}(2), such that $0$ and $F^n$ are the only subspaces of $F^n$ stable under every $\rho'(g)$, and such that $x(D^{\mathrm u})$ is the determinant of $g\mapsto\kappa_0(g)\rho'(\pi(g))$.
\end{lemma}
\leanref{Development/AutomorphicArithmetic/AutomorphicVarieties/ModularityArgument/ThresholdArguments/DenseThresholdCases/PrSeedPtOf.lean}{UpperTheory.V.Pr.seedPtOf_alt}
\leanref{Development/AutomorphicArithmetic/AutomorphicVarieties/CurveDVRSpecialization/SecondThirdPrimeConditionTransferAlt.lean}{UpperTheory.V.Pr.Setup.IdealSpecializationPoint}
\begin{proof}
Write $G=\operatorname{Gal}(\bar{\mathbf Q}/L)$, let $G_0$ be the Galois group of the residual datum of the realization, $\mathrm{pr}:G\to G_0$ its map, $\pi:G_0\to\Gamma_0$ the map of the realization and $\kappa_0$ its character, so that $\pi\circ\mathrm{pr}=p$ and $\kappa_0\circ\mathrm{pr}=\kappa$ on $G$. By the definition of a seed datum, $\chi'$ has order prime to $\ell$ and $\chi_{\mathrm{cyc}}^j\chi'$ is congruent to $\Psi$. Let $J$ have the property of automorphic points.
\begin{enumerate}
\item \emph{Inertia away from $\ell$.} Let $A$ be a place not above $\ell$ and $s\in I_A$. Then $\chi_{\mathrm{cyc}}(s)=1$, so $\chi'(s)$ is congruent to $\Psi(s)$. If $\Psi(s)=1$, then $\chi'(s)$ is a root of unity of order prime to $\ell$ congruent to $1$, so $\chi'(s)=1$; then every $\tau_i(s)=1$, and $r(s)=1$ by (iii).
\item \emph{Factorization through the link.} By (1) and Definition~\ref{sha-descent-link}(1), $r$ is trivial on the kernel of $p$, so $r=r'\circ p$ for a continuous $r':\Gamma_0\to GL_n(\bar{\mathbf Q}_\ell)$, and (i) gives ${}^tr'(g)Pr'(c_\Gamma gc_\Gamma^{-1})=P$ for $g\in\Gamma_0$.
\item \emph{The representation of $G_0$.} Let $\rho=(r'\circ\pi)\otimes\kappa_0$. Then $\rho\circ\mathrm{pr}=r\otimes\kappa$, and $\rho$ is continuous, since $\mathrm{pr}$ is a continuous surjection of compact groups and $r\otimes\kappa$ is continuous.
\item \emph{Characters above $\ell$.} For $v\mid\ell$ let $\psi_{v,i}=\kappa_{A_v,i}\kappa$ on $D_v$, with $\kappa_{A_v,i}$ as in (iv). Then $\rho$ is triangular on $D_v$ with diagonal characters $\psi_{v,i}$.
\item \emph{Tame inertia.} Let $A$ be a place not above $\ell$ and $s\in I_A$ with $\bar\chi_i(\mathrm{pr}(s))=1$ for all $i$. The residual characters are the reductions of $\beta_i=\Psi^{w_{n,i}}\kappa$, so each $\tau_i(s)\kappa(s)$ is congruent to $\beta_i(s)$, hence to $1$. Its order divides the product of the order of $\chi'$ and an integer $o$ prime to $\ell$ with $\kappa^o$ trivial on $I_A$ (Definition~\ref{sha-automorphy-realization}(1)), so it is $1$. By (iii), $(r\otimes\kappa)(s)=1$.
\item \emph{A point of $B_n$.} The pair $(\rho,(\psi_{v,i}))$ satisfies the conditions that define ordinary pairs, with values in $\bar{\mathbf Q}_\ell$: $\rho$ is polarized by (2), since $\kappa_0(g)\kappa_0(cgc^{-1})=\varepsilon(g)^{1-n}$; the characteristic polynomial of $\rho(\mathrm{pr}(s))$ is congruent, by (ii), to $\prod_i(X-\psi_i(s)\kappa(s))$, hence to $\prod_i(X-\beta_i(s))$, whose reduction is that of $\bar D$; $\rho$ is trivial on $I_u\cap\bigcap_i\ker\bar\chi_i$ for $u\in S_{\mathrm a}$ by (5); $\psi_{v,i}$ is congruent to $\Psi^{w_{n,i}}\kappa$ and so lifts $\bar\psi_{v,i}$; $\psi_{v,i}$ and $\theta_{v,i}$ are trivial on $\mu_v$, which lies in the closure of the commutator subgroup of $D_v$ by Definition~\ref{sha-automorphy-realization}(2), as $\psi_{v,i}$ is continuous and $\theta_{v,i}$ is trivial on that closure; and the ordinary relations hold by (4). The values lie in the ring of integers of a finite extension of $\mathbf Q_\ell$, so by the universal property of $B_n$ there is an $\mathcal O_1$-algebra map $x:B_n\to\bar{\mathbf Q}_\ell$ that carries $(D^{\mathrm u},\psi^{\mathrm u})$ to the determinant of $\rho$ and the $\psi_{v,i}$.
\item \emph{$J$ lies in the kernel of $x$.} On $I_v$, $\psi_{v,i}=\chi_{\mathrm{cyc}}^{mw_{n,i}}\kappa$, and $\kappa$ is the character of an algebraic Hecke character. So, by global class field theory as in step (4) of the proof of Lemma~\ref{sha-automorphy-realization-away}, $x$ is arithmetic of a weight, which is increasing since $m\ge2$ and $w_{n,1}>\dots>w_{n,n}$; and $\psi_{v,i}\ne\psi_{v,j}\varepsilon$ for all $i,j$, since $m(w_{n,i}-w_{n,j})\ne1$ for $m\ge2$ and $\chi_{\mathrm{cyc}}$ has infinite order on $I_v$. The determinant $x(D^{\mathrm u})$ is that of $r\otimes\kappa$, which is continuous, attached to a conjugate self-dual cuspidal representation and windowed by (v), and trivial on $I_{v_0}$ by (5), as $v_0\nmid M$ and the $\bar\chi_i$ are trivial on $I_{v_0}$. With (4), the property of automorphic points gives $J\subset\ker x$.
\item \emph{The linked point.} A subspace of $\bar{\mathbf Q}_\ell^n$ stable under $r'(\Gamma_0)$ is stable under $r$, as $r(s)=r'(p(s))$; so $r'$ is irreducible. Lemma~\ref{sha-automorphy-realization-seed-point}, applied to $K=\bar{\mathbf Q}_\ell$, to $x$ of (6), to $r'$ and $P$ of (2) and to the characters $\psi_{v,i}$ of (4), gives a homomorphism $\rho':\Gamma_0\to M_n(\bar{\mathbf Q}_\ell)$ and characters $\psi'_{v,i}:G_v\to\bar{\mathbf Q}_\ell^\times$ satisfying conditions (a), (b) and (c) of Definition~\ref{sha-descent-link}(2), such that $0$ and $\bar{\mathbf Q}_\ell^n$ are the only $\rho'(\Gamma_0)$-stable subspaces of $\bar{\mathbf Q}_\ell^n$ and $x(D^{\mathrm u})$ is the determinant of $g\mapsto\kappa_0(g)\rho'(\pi(g))$. With $F=\bar{\mathbf Q}_\ell$ this proves the lemma.
\end{enumerate}
\end{proof}
% seedPtOf_alt = main (PrSeedPtOf.lean), from GlobalArtinReciprocity: (1) hε, hχo, eq_one_of_near and one_of_diag; (2) exists_frame_factorization_of_trivial_on_inertia; (3) hρr, htw, hρc; (4) hr.above and hfl, hfl'; (5) htame with hka.ord; (6) pt_ex = W.exists_determinantPair_over_algebraicIntegers, with beta_ex and Sd.shapeFacts for the congruences, comm_one and K.mu for mu_v; (7) diagonal_characters_regular_algebraic_of_inertia_weights and hJ = PointTheoremBaseObject, with R.v0M and R.v0.1; (8) universalMaps_point_of_polarized_flagged_rep (Lemma sha-automorphy-realization-seed-point), with the irreducibility of r' from hr.irr and k.p0.

\begin{lemma}\label{sha-automorphy-realization-dimension}
Let $\mathcal O$ be a discrete valuation ring, let $m\ge0$, let $\Lambda$ be an $\mathcal O$-algebra isomorphic to $\mathcal O[[X_1,\dots,X_m]]$, let $B$ be a noetherian ring with a ring homomorphism $\Lambda\to B$, and let $J$ be an ideal of $B$ such that $B/J$ is finite over $\Lambda$ and such that $\lambda b\in J$ with $\lambda\in\Lambda$ nonzero implies $b\in J$. Then $B/\mathfrak p$ has Krull dimension $1+m$ for every minimal prime $\mathfrak p$ over $J$.
\end{lemma}
\leanref{Development/AutomorphicArithmetic/AutomorphicVarieties/CurveDVRSpecialization/DimOf.lean}{UpperTheory.V.Pr.dimOf}
\leanref{Development/AutomorphicArithmetic/AutomorphicVarieties/AlgebraTools/HidaCompleteRings/Component.lean}{UpperTheory.V.AlgebraTools.AnnihilatorQuotientFinite.component}
\begin{proof}
Let $T=B/J$ and $P=\mathfrak p/J$, a minimal prime of $T$.
The nonzero elements of $\Lambda$ are nonzerodivisors on $T$, and a minimal prime consists of zerodivisors, so $P\cap\Lambda=0$.
So $\Lambda\to T/P$ is injective, and it is integral since $T$ is finite over $\Lambda$; hence $\dim T/P=\dim\Lambda$ \cite[\S9]{Mat86}.
Finally $T/P=B/\mathfrak p$, and $\dim\mathcal O[[X_1,\dots,X_m]]=1+m$ \cite[\S15]{Mat86}.
\end{proof}

\begin{lemma}\label{sha-automorphy-realization-arcs}
In Situation~\ref{sha-hecke-situation}, suppose that the residual characters $\bar\chi_1,\dots,\bar\chi_n$ agree on $D_{v_0}$ and that $\varepsilon\bmod\ell$ is nontrivial on $D_{v_0}$. Then every arc (Definition~\ref{sha-hecke-arcs}) satisfies condition (N4) there.
\end{lemma}
\leanref{Development/AutomorphicArithmetic/AutomorphicVarieties/CurveDVRSpecialization/EveryPointSatisfiesFrameConditionAlt.lean}{UpperTheory.V.Pr.n4Of}
\begin{proof}
Let $(\rho,\psi)$ be an arc over $\kappa[[T]]$, with fraction field $\mathcal K$, and let $X$ be a matrix over $\mathcal K$ with $\varepsilon(g)\rho(g)X=X\rho(g)$ for $g\in D_{v_0}$.
Take $g\in D_{v_0}$ with $\varepsilon(g)\not\equiv1\pmod\ell$, and let $A=\rho(g)$, $z=\varepsilon(g)$, $P$ the characteristic polynomial of $A$ and $c=\bar\chi_1(g)$.
Since the determinant of $\rho$ reduces to $\bar D$ and the $\bar\chi_i(g)$ are all equal to $c$, $P$ reduces modulo $T$ to $(X-c)^n$.
From $zAX=XA$ follows $P(zA)X=XP(A)=0$, by the Cayley-Hamilton theorem.
Modulo $T$, $A$ reduces to $\bar A$ with $N=\bar A-c$ nilpotent, and $P(zA)$ reduces to $(\bar z\bar A-c)^n=\big((\bar z-1)c+\bar zN\big)^n$, which is invertible because $(\bar z-1)c\ne0$ and $\bar zN$ is nilpotent.
So $\det P(zA)$ is a unit of $\kappa[[T]]$, and $X=0$.
\end{proof}

\begin{lemma}\label{sha-automorphy-realization-ideal}
In Situation~\ref{sha-automorphy-link-situation}, let a realization of rank $n$ with coefficients $\mathcal O_1$, large enough, satisfy the conclusion of Lemma~\ref{sha-automorphy-realization-seed} for every ideal with the property of automorphic points. Then there is an ideal $J_n$ of $B_n$ such that: $B_n/J_n$ is finite over $\Lambda$; every minimal prime over $J_n$ has dimension $1+nd$; the determinant over $B_n/J_n$ is trivial on the inertia group at $v_0$; $J_n$ satisfies the bounds of Proposition~\ref{sha-patching-bound} at the nice arcs of type I and along the chains of type E; the conclusion of Lemma~\ref{sha-automorphy-realization-seed} holds for $J_n$; and the irreducible arithmetic points of $B_n/J_n$ are classical, as in Proposition~\ref{sha-hecke-ideal}(5). Moreover every arc of the residual datum of the realization satisfies (N4) of Definition~\ref{sha-hecke-arcs}.
\end{lemma}
\leanref{Development/AutomorphicArithmetic/AutomorphicVarieties/ModularityArgument/ThresholdArguments/DenseThresholdCases/PrLinkedEnd.lean}{UpperTheory.V.Pr.modularity_linked_ends}
\leanref{Development/AutomorphicArithmetic/AutomorphicVarieties/CurveDVRSpecialization/SecondThirdPrimeConditionTransferAlt.lean}{UpperTheory.V.Pr.Setup.ModIn}
\begin{proof}
Let $B_n$, $\Lambda$ and $v_0$ be those of the realization.
\begin{enumerate}
\item \emph{The field and the degree.} Each group $G_v$ of the link lies in $\operatorname{Gal}(\bar{\mathbf Q}/L)$, since $\ell\mid M$, and $\omega$ is nontrivial on it (Definition~\ref{sha-descent-link}(1)); so $L$ contains no primitive $\ell$-th root of unity, as Situation~\ref{sha-hecke-situation} requires. The places $A_v$ represent the places of $L^+$ above $\ell$, and $d_v$ is the local degree of $L^+$ at $v$, so $\sum_vd_v=[L^+:\mathbf Q]=d$.
\item \emph{The Iwasawa algebra.} By local class field theory, $\Lambda$ is a power series ring over $\mathcal O_1$ in variables indexed by $v\mid\ell$, $1\le i\le n$ and a basis of the free part of $\mathcal O_{L_{\tilde v}}^\times\otimes\mathbf Z_\ell$ (proof of Lemma~\ref{sha-hecke-algebra}). That free part has rank $[L_{\tilde v}:\mathbf Q_\ell]=d_v$, since $v$ splits in $L$; so $\Lambda\cong\mathcal O_1[[X_1,\dots,X_{nd}]]$ by (1).
\item \emph{The ideal.} Proposition~\ref{sha-hecke-ideal} gives the Hecke ideal $J_n\subset B_n$, with Taylor-Wiles data, and Proposition~\ref{sha-patching-bound} gives the bounds for it. By Proposition~\ref{sha-hecke-ideal}(1), $B_n/J_n$ is finite over $\Lambda$, and $\lambda b\in J_n$ with $\lambda\in\Lambda$ nonzero implies $b\in J_n$; since the residual characters of the realization are trivial on $I_{v_0}$ and agree on $D_{v_0}$, and $\varepsilon\bmod\ell$ is nontrivial on $D_{v_0}$, the determinant over $B_n/J_n$ kills $g-1$ for $g\in I_{v_0}$.
\item \emph{The dimension.} The ring $B_n$ is noetherian, so Lemma~\ref{sha-automorphy-realization-dimension}, with (2) and (3), gives $\dim B_n/\mathfrak p=1+nd$ for every minimal prime $\mathfrak p$ over $J_n$.
\item \emph{The seed and classicality.} By Proposition~\ref{sha-hecke-ideal}(3), $J_n$ has the property of automorphic points, so by hypothesis the conclusion of Lemma~\ref{sha-automorphy-realization-seed} holds for it. Classicality is Proposition~\ref{sha-hecke-ideal}(5).
\item \emph{(N4).} This is Lemma~\ref{sha-automorphy-realization-arcs}, by the conditions at $v_0$.
\end{enumerate}
\end{proof}
% modularity_linked_ends (PrLinkedEnd.lean): (1) is degEq (orbit_count; k.mu and cycMod_one for the roots of unity); (2) is hLam (W.lamEx from LocCFTAll); (3) is hMod = W.ModUAlt (W.modU_of_printsH from HidaUAlt and PfromAut), with hJ.ring.fin, tf and unr0 R.v0; (4) is dimOf; (5) is hs J hJ.pt and hJ.end2; (6) is n4Of. The fields of Setup.ModIn are fin, dim, unr (Unr0), n4, p1, p3 and seed.

\subsection{Components of the ring of linked points}\label{subsection-sha-linked-components}

\begin{definition}\label{sha-automorphy-realization-components-transfer-primes}
In Situation~\ref{sha-automorphy-link-situation}, let $n\in\{2,3\}$, let $\mathcal O_1\supseteq\mathcal O$ be the ring of integers of a finite extension of $\mathbf Q_\ell$, and let $A_n$, its universal linked point $(\rho^{\mathrm u},\psi^{\mathrm u})$, $B$ and $A_n^\circ$ be as in Definition~\ref{sha-automorphy-realization}(2).
\begin{enumerate}
\item Let $\Lambda_B$ be the ideal of $B$ generated by $\mathfrak m_{\mathcal O_1}$ and the $\psi^{\mathrm u}_{v,i}(g)-1$ for $v\mid\ell$, $1\le i\le n$ and $g\in I_v$ with $\bar\Psi(g)=1$. For a prime $\mathfrak p$ of $A_n^\circ$ write $\mathfrak p_B=\mathfrak p\cap B$.
\item For a prime $x$ of $A_n^\circ$, a \emph{completion at $x$} consists of a complete discrete valuation ring $\mathcal O'$ finite over $\mathcal O_1$ with $\mathfrak m_{\mathcal O_1}\mathcal O'=\mathfrak m_{\mathcal O'}$, a complete local noetherian $\mathcal O'$-algebra $R$ with the residue field of $\mathcal O'$, and an $\mathcal O_1$-algebra map $\iota:A_n^\circ\to R$ with $\iota^{-1}(\mathfrak m_R)=x$, such that every $\mathcal O_1$-algebra map $f$ from $A_n^\circ$ to a complete local noetherian $\mathcal O'$-algebra $C$ with the residue field of $\mathcal O'$, inducing the same map as $\iota$ on residue fields, is $g\circ\iota$ for a unique $\mathcal O'$-algebra map $g:R\to C$. Write $(\rho_R,\psi_R)$ for the image of the universal linked point under $\iota$, and $P_B$ for the inverse image in $B$ of a prime $P$ of $R$.
\item Let $(\mathcal O',R,\iota)$ be a completion at a prime of $A_n^\circ$ and $P$ a prime of $R$. Call $P$ \emph{off-triangular} if it satisfies the condition of Definition~\ref{sha-automorphy-realization}(4) for $\rho_R$, and, for $n=3$, \emph{endoscopic} if some entry of $\rho_R$ below the diagonal and some entry above it lie outside $P$ while all entries of odd weight lie in $P$. Call $P$ \emph{nice} if $\dim R/P=1$, $\ell\in P$, some entry of $\rho_R$ below the diagonal and some entry above it lie outside $P$, and for every $v$ and $i\ne j$ some $g\in I_v$ with $\chi_{\mathrm{cyc}}(g)=1$ has $\psi_{R,v,i}(g)-\psi_{R,v,j}(g)\notin P$. For $n=3$, call a prime $E$ of $R$ \emph{thick} if it is endoscopic and there is a multiplicative map $r$ from $\Gamma_0$ to the $3\times3$ matrices over $R/E$, with entries continuous for the topology of the maximal ideal, which agrees with $\rho_R$ modulo $E$ on the entries of even weight, has $r(g)_{21}=r(g)_{23}=0$, satisfies (c) of Definition~\ref{sha-descent-link}(2) with the characters $\psi_{R,v,i}$ modulo $E$, and also with $r(c_\Gamma(\mathrm{loc}_vh)^{-1}c_\Gamma^{-1})$ in place of $r(\mathrm{loc}_vh)$ and the characters in the reverse order, and for which no field $K'$ containing $R/E$ has a vector $y\in K'^3$ with $y_2=1$ and $r(g)y=r(g)_{22}y$ for all $g$. Call two nice endoscopic primes \emph{thickly linked} if they contain a common thick prime.
\end{enumerate}
\end{definition}
\leanref{Development/AutomorphicArithmetic/AutomorphicVarieties/FrobeniusSmoothVectors/LeAdicClosure.lean}{UpperTheory.V.Fr.Space.lamId}
\leanref{Development/AutomorphicArithmetic/AutomorphicVarieties/FrobeniusSmoothVectors/LeAdicClosure.lean}{UpperTheory.V.Fr.Space.Loc}
\leanref{Development/AutomorphicArithmetic/AutomorphicVarieties/FrobeniusSmoothVectors/LeAdicClosure.lean}{UpperTheory.V.Fr.Thick}
\leanref{Development/AutomorphicArithmetic/AutomorphicVarieties/CurveDVRSpecialization/PropagationTheoremGL2.lean}{UpperTheory.V.Pr.Nice}
\leanref{Development/AutomorphicArithmetic/AutomorphicVarieties/CurveDVRSpecialization/PropagationTheoremGL2.lean}{UpperTheory.V.Pr.ThickLink}

\begin{lemma}\label{sha-automorphy-realization-components-transfer-twist}
In Situation~\ref{sha-automorphy-link-situation}, let $n\in\{2,3\}$ (a restriction of Alpöge's statement, which holds, with the same proof, for every rank $n\ge1$, for linked points and the properties below defined by the same conditions with the weights $w_{n,i}=n+1-2i$ for $n\ne3$; this book defines linked points only in ranks two and three), let $\mathcal O_1\supseteq\mathcal O$ be the ring of integers of a finite extension of $\mathbf Q_\ell$, with residue field $k_1$, and let a residual datum of rank $n$ over $\mathcal O_1$ (Situation~\ref{sha-hecke-situation}), a map $\pi:G_0\to\Gamma_0$ and a character $\kappa_0:G_0\to\mathcal O_1^\times$ be given with the properties that Definition~\ref{sha-automorphy-realization}(2) requires of them, as for a realization of rank $n$ with coefficients $\mathcal O_1$; here $G_0$ is the Galois group of the residual datum, whose decomposition groups $D_v$ above $\ell$ are identified with the groups $G_v$ of the link. Let $C$ be a complete local noetherian $\mathcal O_1$-algebra with residue field $k_1$, and $(\rho,\psi)$ a linked point of rank $n$ with values in $C$ lifting the residual one (Definition~\ref{sha-descent-link}(2)). Then the determinant $D$ of the representation $\tilde\rho:g\mapsto\kappa_0(g)\rho(\pi(g))$ of $G_0$, with the characters $\tilde\psi_{v,i}:h\mapsto\psi_{v,i}(h)\kappa_0(h)$ of $D_v$, is an ordinary pair over $C$ for the residual datum of the realization (Situation~\ref{sha-hecke-situation}).
\end{lemma}
\leanref{Development/AutomorphicArithmetic/AutomorphicVarieties/CurveDVRSpecialization/PxIsCont.lean}{px_pair}
\leanref{Development/AutomorphicArithmetic/AutomorphicVarieties/CurveDVRSpecialization/SelfTwistEigenvalues/UnipTorsion.lean}{UpperTheory.V.Pr.T.unipTorsion}
\begin{proof}
For $y=\sum_ga_gg$ in $C[G_0]$, the characteristic polynomial of $D$ at $y$ is that of the matrix $\tilde\rho(y)=\sum_ga_g\tilde\rho(g)$, and $\operatorname{tr}_D(y)$ is its trace.
\begin{enumerate}
\item \emph{Continuity and reduction.} The entries of $\tilde\rho$ are continuous for the $\mathfrak m_C$-adic topology, as $\kappa_0$, $\pi$ and the entries of $\rho$ are; so the coefficients of the characteristic polynomials of $D$ at the elements of $G_0$ are continuous. Modulo $\mathfrak m_C$, $\rho(\pi(g))$ is the diagonal matrix with entries $\bar\Psi_\Gamma(\pi(g))^{w_{n,i}}$, so $\tilde\rho(g)$ reduces to the diagonal matrix with entries $\bar\chi_i(g)$, the residual characters of the realization. Hence $D$ reduces to $\bar D$.
\item \emph{The characters.} Each $\tilde\psi_{v,i}$ is continuous, and it reduces to $\bar\Psi_\Gamma(\mathrm{loc}_vh)^{w_{n,i}}\bar\kappa_0(h)=\bar\chi_i(h)=\bar\psi_{v,i}(h)$. A continuous character $\chi$ of $D_v$ with values in $C^\times$ is trivial on the closure of the commutator subgroup of $D_v$: for each $N$, the $h$ with $\chi(h)\equiv1$ modulo $\mathfrak m_C^N$ form a closed set containing the commutators, and $\bigcap_N\mathfrak m_C^N=0$. So $\tilde\psi_{v,i}$ is trivial on $\mu_v$, which lies in that closure by Definition~\ref{sha-automorphy-realization}(2), and so is $\theta_{v,i}$, which is trivial on that closure (Situation~\ref{sha-hecke-situation}).
\item \emph{Condition (1) of Situation~\ref{sha-hecke-situation}.} For $g\in G_0$, $\kappa_0(g)\kappa_0(cgc^{-1})=\varepsilon(g)^{1-n}$ and $\pi(cgc^{-1})=c_\Gamma\pi(g)c_\Gamma^{-1}$. By condition (a) of Definition~\ref{sha-descent-link}(2), $g^*=\varepsilon(g)^{1-n}(cgc^{-1})^{-1}$ maps to $\varepsilon(g)^{1-n}\kappa_0(cgc^{-1})^{-1}J_n\,{}^t\rho(\pi(g))J_n=J_n\,{}^t\tilde\rho(g)J_n$. As $J_n^2=1$, this matrix has the characteristic polynomial of $\tilde\rho(g)$.
\item \emph{Condition (2).} For $h_1,\dots,h_n\in D_v$, the element $\prod_i(h_i-\tilde\psi_{v,i}(h_i))$ maps to $\prod_i\kappa_0(h_i)\big(\rho(\mathrm{loc}_vh_i)-\psi_{v,i}(h_i)\big)$, which is $0$ by condition (c). So $\operatorname{tr}_D$ vanishes on its products with the elements of $G_0$.
\item \emph{Condition (3).} Let $u\in S_{\mathrm a}$, $g\in I_u\cap\bigcap_i\ker\bar\chi_i$ and $\gamma=\tilde\rho(g)$. By the definition of a realization there is $m$ prime to $\ell$ with $\pi(g)^m=1$ and $\kappa_0(g)^m=1$, so $\gamma^m=1$. By (1), $\gamma$ reduces to the identity, so its characteristic polynomial $\chi$ is monic and reduces to $(X-1)^n$. Let $s=\sum_{0\le k<m}X^k$. The ring $C[X]/(\chi)$ is finite over the local ring $C$, so each of its maximal ideals $\mathfrak M$ lies over $\mathfrak m_C$; as $\chi$ reduces to $(X-1)^n$, $X$ maps to $1$ in the residue field of $\mathfrak M$, and $s$ maps to $m\ne0$. So $s$ lies in no maximal ideal and is a unit of $C[X]/(\chi)$. By the Cayley-Hamilton theorem, $X\mapsto\gamma$ defines a ring map from $C[X]/(\chi)$ to the matrices over $C$; it carries $s$ to an invertible matrix $s(\gamma)$, and $(1-\gamma)s(\gamma)=1-\gamma^m=0$. So $\gamma=1$, and $\operatorname{tr}_D\big(h(g-1)\big)=\operatorname{tr}\big(\tilde\rho(h)(\gamma-1)\big)=0$ for $h\in G_0$.
\end{enumerate}
\end{proof}
% px_pair (PxIsCont.lean): cont is px_isCont with px_mul and px_cont; red is px_red; psi is px_psi (px_comm for mu_v, with the field mu of TwLink and r.th_comm); P is px_star with Matrix.charpoly_transpose; Ord is px_ord with px_prod; I is px_ine (the field tame of TwLink, then Pr.T.unipTorsion, whose unit is unipTorsion_unit).

\begin{lemma}\label{sha-automorphy-realization-components-transfer-dense}
\begin{enumerate}
\item Let $\mathcal O$ be a local ring and $R$ a complete local noetherian $\mathcal O$-algebra, with $\mathcal O\to R$ local, such that every element of $R$ is congruent modulo $\mathfrak m_R$ to an element of $\mathcal O$. Let $T\subseteq R$ be a set such that any two $\mathcal O$-algebra maps from $R$ to an $\mathcal O$-algebra $R'$ with the same properties as $R$ that agree on $T$ are equal. Then for every $x\in R$ and $N\ge0$ there is $y$ in the $\mathcal O$-subalgebra of $R$ generated by $T$ with $x-y\in\mathfrak m_R^N$.
\item In Situation~\ref{sha-hecke-situation}, let $(D^{\mathrm u},\psi^{\mathrm u})$ be the universal ordinary pair over $B$. Then for every $x\in B$ and $N\ge0$ there is $y$ in the $\mathcal O$-subalgebra of $B$ generated by the coefficients of the characteristic polynomials of $D^{\mathrm u}$ at the elements of $\Gamma_0$ and the values of the characters $\psi^{\mathrm u}_{v,i}$, with $x-y\in\mathfrak m_B^N$.
\end{enumerate}
\end{lemma}
\leanref{Development/AutomorphicArithmetic/AutomorphicVarieties/CommutativeAlgebraHeight/UGenApprox.lean}{UpperTheory.V.CA.u_gen}
\leanref{Development/AutomorphicArithmetic/AutomorphicVarieties/CommutativeAlgebraHeight/ExistsAdjoinCoeffCharacterSubMemMaximalIdealPow.lean}{UpperTheory.V.CA.exists_adjoin_coeff_character_sub_mem_maximalIdeal_pow}
\begin{proof}
(1) Let $S$ be the $\mathcal O$-subalgebra generated by $T$ and $k$ the residue field of $\mathcal O$. For $t\in T$ choose $a_t\in\mathcal O$ with $t-a_t\in\mathfrak m_R$, and let $I$ be the ideal of $R$ generated by the $t-a_t$ and the image of $\mathfrak m_{\mathcal O}$; its generators lie in $S\cap\mathfrak m_R$.
\begin{enumerate}
\item[(i)] \emph{$\mathfrak m_R\subseteq I+\mathfrak m_R^2$.} Suppose not, and let $x\in\mathfrak m_R$ lie outside $I+\mathfrak m_R^2$. Let $N$ be the $\mathcal O$-submodule $I+\mathfrak m_R^2+\mathcal O\cdot1$ of $R$. Then $x\notin N$: if $x=y+a$ with $y\in I+\mathfrak m_R^2$ and $a\in\mathcal O$, then $a\in\mathfrak m_R$, so $a\in\mathfrak m_{\mathcal O}$ as $\mathcal O\to R$ is local, and $x\in I+\mathfrak m_R^2$. The module $R/N$ is killed by $\mathfrak m_{\mathcal O}$, so it is a $k$-vector space, and some $k$-linear map on it is nonzero at the class of $x$; let $\delta:R\to k$ be its composite with the projection. For $y,z\in R$ with $y\equiv a$ and $z\equiv b$ modulo $\mathfrak m_R$, $a,b\in\mathcal O$, the element $(y-a)(z-b)$ lies in $\mathfrak m_R^2$, which gives $\delta(yz)=\bar y\delta(z)+\delta(y)\bar z$. So $y\mapsto\bar y+\delta(y)\epsilon$ is an $\mathcal O$-algebra map from $R$ to $k[\epsilon]/(\epsilon^2)$, which has the properties of $R$. It agrees with $y\mapsto\bar y$ on $T$, as $\delta$ kills $I$ and $\mathcal O\cdot1$, but not at $x$; this contradicts the hypothesis.
\item[(ii)] \emph{$R=S+\mathfrak m_R^N$ for every $N$.} For $N=1$ this is the hypothesis on residues. By (i), $\mathfrak m_R^N\subseteq(I+\mathfrak m_R^2)^N\subseteq I^N+\mathfrak m_R^{N+1}$. An element of $I^N$ is a sum of products $rs$ with $s$ a product of $N$ generators of $I$, hence $s\in S\cap\mathfrak m_R^N$, and $r\in R=\mathcal O+\mathfrak m_R$; so $rs\in S+\mathfrak m_R^{N+1}$. Hence $\mathfrak m_R^N\subseteq S+\mathfrak m_R^{N+1}$, and $R=S+\mathfrak m_R^N$ gives $R=S+\mathfrak m_R^{N+1}$.
\end{enumerate}
(2) Two $\mathcal O$-algebra maps $f,f'$ from $B$ to a complete local noetherian $\mathcal O$-algebra $R'$ with residue field $k$ that agree on these elements carry $(D^{\mathrm u},\psi^{\mathrm u})$ to ordinary pairs with the same characters and the same characteristic polynomials at the elements of $\Gamma_0$. A determinant of a group is determined by its characteristic polynomials at the elements of the group, by Amitsur's formula \cite[(1.5) and Corollary~1.14]{Che14}; so the two pairs are equal, and $f=f'$ by the universal property of $B$. Now apply (1).
\end{proof}
% u_gen (UGenApprox.lean): (i) is the private u_gen_span (u_gen_dual for the functional, the derivation D and the map f to DualNumber); (ii) is u_gen_approx, which is Reformulated.B.LayerApproximation.co_ext_approx. (2) is exists_adjoin_coeff_character_sub_mem_maximalIdeal_pow: pair_map, B.univ and d_gext (the determinant from the coefficients at the group elements).

\begin{lemma}\label{sha-automorphy-realization-components-transfer-coefficients}
In Situation~\ref{sha-automorphy-link-situation}, let $n\in\{2,3\}$, let $\mathcal O_1\supseteq\mathcal O$ be the ring of integers of a finite extension of $\mathbf Q_\ell$ such that the link, with its residual characters composed with the map of residue fields, satisfies the residual conditions with some constant (Definition~\ref{sha-descent-residual}; neither this hypothesis nor $n\in\{2,3\}$ is used in the proof), and let $A_n$, its universal linked point $(\rho^{\mathrm u},\psi^{\mathrm u})$ and $B$ be as in Definition~\ref{sha-automorphy-realization}(2). Let $g_1,\dots,g_m\in\Gamma_0$. Then the coefficients of the polynomial $\det\big(\sum_aX_a\rho^{\mathrm u}(g_a)\big)$ in $A_n[X_1,\dots,X_m]$ lie in $B$. In particular, for $c\in B$ and $g\in\Gamma_0$, the coefficients of the characteristic polynomial of $c\rho^{\mathrm u}(g)$ lie in $B$.
\end{lemma}
\leanref{Development/AutomorphicArithmetic/AutomorphicVarieties/FrobeniusSmoothVectors/DetMem.lean}{UpperTheory.V.Fr.det_mem}
\leanref{Development/AutomorphicArithmetic/AutomorphicVarieties/CurveDVRSpecialization/PxIsCont.lean}{px_coef}
\begin{proof}
Expanding the determinant over the permutations $\sigma$ of $\{1,\dots,n\}$ and the maps $f:\{1,\dots,n\}\to\{1,\dots,m\}$,
\[\det\Big(\sum_aX_a\rho^{\mathrm u}(g_a)\Big)=\sum_\sigma\operatorname{sgn}(\sigma)\sum_f\prod_i\rho^{\mathrm u}(g_{f(i)})_{\sigma(i),i}\prod_iX_{f(i)}.\]
So each coefficient is a sum of products $\pm\prod_i\rho^{\mathrm u}(g_{f(i)})_{\sigma(i),i}$, whose total weight is $\sum_i(w_{n,\sigma(i)}-w_{n,i})=0$; these lie in $B$. The characteristic polynomial of $c\rho^{\mathrm u}(g)$ is the image of $\det\big(X_1\cdot1+X_2\rho^{\mathrm u}(g)\big)$ under $X_1\mapsto X$, $X_2\mapsto-c$, so its coefficients are polynomials in $c$ with coefficients in $B$.
\end{proof}

\begin{lemma}\label{sha-automorphy-realization-components-transfer-onto}
In Situation~\ref{sha-automorphy-link-situation}, let $n\in\{2,3\}$, let $\mathcal O_1\supseteq\mathcal O$ be the ring of integers of a finite extension of $\mathbf Q_\ell$, with residue field $k_1$, such that the link, with its residual characters composed with the map of residue fields, satisfies the residual conditions with some constant, and let $A_n$, its universal linked point $(\rho^{\mathrm u},\psi^{\mathrm u})$ and $B$ be as in Definition~\ref{sha-automorphy-realization}(2). Let $C$ be a complete local noetherian $\mathcal O_1$-algebra with residue field $k_1$ and $\phi:C\to A_n$ an $\mathcal O_1$-algebra map whose image contains the traces of the $\rho^{\mathrm u}(g)$, $g\in\Gamma_0$, and the values of the $\psi^{\mathrm u}_{v,i}$. Then $B\subseteq\phi(C)$.
\end{lemma}
\leanref{Development/AutomorphicArithmetic/AutomorphicVarieties/FrobeniusSmoothVectors/TraceMem.lean}{UpperTheory.V.Fr.B_gen}
\leanref{Development/AutomorphicArithmetic/AutomorphicVarieties/FrobeniusSmoothVectors/AlgClosed.lean}{UpperTheory.V.Fr.alg_closed}
\begin{proof}
Write $\rho=\rho^{\mathrm u}$, $g^*=c_\Gamma g^{-1}c_\Gamma^{-1}$ and $w_i=w_{n,i}$. Let $A'=\phi(C)$, a quotient of $C$, hence a complete local noetherian $\mathcal O_1$-algebra with residue field $k_1$, whose maximal ideal is $A'\cap\mathfrak m_{A_n}$. Let $\Sigma$ be the $A'$-span of the $\rho(g)$, $g\in\Gamma_0$, an $A'$-algebra of matrices; the traces of its elements lie in $A'$, and so does the $(i,i)$ entry of $M$ whenever $M$ and $E_{ii}M$ lie in $\Sigma$, where $E_{ij}$ are the matrix units.
\begin{enumerate}
\item \emph{A diagonal matrix.} By conditions (a) and (b) of Definition~\ref{sha-descent-link}(2), $z=\rho(g_0)-\rho(g_0^*)$ is $\zeta\operatorname{diag}(w_1,\dots,w_n)$ with $\zeta=\rho(g_0)_{11}-\rho(g_0)_{nn}$, as $\rho(g^*)_{ii}=\rho(g)_{n+1-i,n+1-i}$ and $w_{n+1-i}=-w_i$. By condition (e), $\zeta$ reduces to $a-a^{-1}$ with $a=\bar\Psi_\Gamma(g_0)$, which is nonzero since $a^2\ne1$ by the residual conditions; so $\zeta$ is a unit of $A_n$.
\item \emph{$\zeta\in A'$.} As $z^2\in\Sigma$, $\operatorname{tr}(z^2)=\zeta^2\sum_iw_i^2=2\zeta^2$ lies in $A'$, and so does $\zeta^2$, as $\ell$ is odd. Choose $b\in\mathcal O_1$ with $\zeta\equiv b$ modulo $\mathfrak m_{A_n}$; it is a unit. By Hensel's lemma in $A'$, applied to $X^2-\zeta^2$ at $b$, whose derivative $2b$ is a unit, there is $d\in A'$ with $d^2=\zeta^2$ and $d\equiv b\equiv\zeta$. Then $(d+\zeta)(d-\zeta)=0$ with $d+\zeta\equiv2\zeta$ a unit, so $\zeta=d\in A'$, and $\zeta$ is a unit of $A'$.
\item \emph{Matrix units.} So $Y=\operatorname{diag}(w_1,\dots,w_n)=\zeta^{-1}z$ lies in $\Sigma$, and so do the $E_{ii}$: if $n=2$, $2E_{11}=1+Y$ and $2E_{22}=1-Y$; if $n=3$, $2E_{11}=Y^2+Y$, $E_{22}=1-Y^2$ and $2E_{33}=Y^2-Y$.
\item \emph{Products of weight $0$.} For $i,j$ let $E(i,j)$ be the $A'$-span of the entries $\rho(g)_{ij}$, $g\in\Gamma_0$. As $E_{ii}\rho(g)E_{jj}\rho(h)E_{kk}\in\Sigma$ has $(i,k)$ entry $\rho(g)_{ij}\rho(h)_{jk}$, $E(i,j)E(j,k)\subseteq E(i,k)$, and $E(i,i)\subseteq A'$. By condition (a), $E(i,j)=E(n+1-j,n+1-i)$. Let $x$ be a product of entries of $\rho$ of total weight $0$; we show $x\in A'$ by induction on the number of factors, allowing factors from $A'$ and from the $E(i,j)$, of weight $w_i-w_j$. A factor of weight $0$ lies in some $E(i,i)$, as the $w_i$ are distinct, hence in $A'$. Otherwise there are factors $y\in E(i,j)$ of positive weight and $y'\in E(k,m)$ of negative weight. As the $w_i$ lie in $\{1,0,-1\}$ and $w_{n+1-i}=-w_i$, one of $j=k$, $m=i$, $n+1-i=k$ or $m=n+1-j$ holds; after replacing $y$ by its copy in $E(n+1-j,n+1-i)$ in the last two cases, $yy'$ lies in a single $E(a,b)$, of weight the sum of the two. This lowers the number of factors.
\item \emph{Closure.} By (4) and the hypothesis on the $\psi^{\mathrm u}_{v,i}$, the $\mathcal O_1$-subalgebra $D$ whose closure is $B$ lies in $A'$. Let $x\in B$, and choose $y_N\in D$ with $x-y_N\in\mathfrak m_{A_n}^{N+1}$. The ideals $I_N=A'\cap\mathfrak m_{A_n}^N$ of $A'$ decrease and meet in $0$, so by Lemma~\ref{sha-automorphy-unramified-near-determined-cofinal} each $\mathfrak m_{A'}^k$ contains some $I_N$; as $y_{N+1}-y_N\in I_{N+1}$, the $y_N$ converge in the complete ring $A'$ to some $y$ with $y-y_N\in I_{N+1}$. Then $x-y$ lies in every $\mathfrak m_{A_n}^N$, so $x=y\in A'$.
\end{enumerate}
\end{proof}
% B_gen (TraceMem.lean): co_quot for A' and mem_max; z_diag and hzeta for (1); sqrt_mem (HenselianRing.is_henselian) and trace_mem for (2); units_mem for (3); mon_mem (ent_mul, single_mem, hrev, hdiag and the combinatorial lemma cycles with w_facts) for (4); closed (alg_closed, from GlobalTheory.exists_le_maximalIdeal_pow_of_iInf_eq_bot, which is Chevalley's lemma) for (5).

\begin{lemma}\label{sha-automorphy-realization-components-transfer-arc-curve}
Let $p$ be a prime and $R$ a complete local noetherian domain of Krull dimension one and characteristic $p$, with finite residue field. Then there are a finite field $\kappa$ and an injective local homomorphism $R\to\kappa[[T]]$ that makes $\kappa[[T]]$ a finite $R$-module.
\end{lemma}
\leanref{Development/AutomorphicArithmetic/AutomorphicVarieties/CurveDVRSpecialization/SelfTwistEigenvalues/CurveNorm.lean}{UpperTheory.V.Pr.T.curveNorm}
\leanref{Development/AutomorphicArithmetic/AutomorphicVarieties/CurveDVRSpecialization/SelfTwistEigenvalues/CurveOverSeries.lean}{UpperTheory.V.Pr.T.curveOverSeries}
\leanref{Development/AutomorphicArithmetic/AutomorphicVarieties/CurveDVRSpecialization/SelfTwistEigenvalues/CurveDVR.lean}{UpperTheory.V.Pr.T.curveDVR}
\begin{proof}
\begin{enumerate}
\item \emph{A power series ring.} As $\dim R=1$, $\mathfrak m_R\ne0$; choose $\pi\ne0$ in $\mathfrak m_R$. As $R$ is complete, there is a ring homomorphism $f:\mathbf F_p[[X]]\to R$ with $f(X)=\pi$. It is injective: a nonzero series is $X^mu$ with $u$ a unit, and $f(X^mu)=\pi^mf(u)\ne0$, as $R$ is a domain. It is local: a series that is not a unit is $Xh$, and $f(Xh)=\pi f(h)\in\mathfrak m_R$. It makes $R$ a finite $\mathbf F_p[[X]]$-module. Indeed, a nonzero prime of the one-dimensional local domain $R$ is maximal, so $\mathfrak m_R$ is the only prime containing $\pi$, and $\mathfrak m_R^e\subseteq\pi R$ for some $e$, as $\mathfrak m_R$ is finitely generated. So $R/\pi R$, a quotient of $R/\mathfrak m_R^e$, is finite, as $R/\mathfrak m_R$ is. As $R$ is $\pi$-adically separated, being $\mathfrak m_R$-adically separated, and $\mathbf F_p[[X]]$ is $X$-adically complete, $R$ is a finite $\mathbf F_p[[X]]$-module \cite[Theorem~8.4]{Mat86}.
\item \emph{The normalization.} Let $A=\mathbf F_p[[X]]$, with fraction field $F$, let $L$ be the fraction field of $R$, which is finite over $F$ by (1), and let $A'$ be the integral closure of $A$ in $L$. As $F$ is complete for the valuation of $A$, $A'$ is a discrete valuation ring, finite over $A$, and $L$ is complete for its valuation \cite[Chapter~II, \S2, Proposition~3]{Ser79}; so $A'$ is complete. As $R$ is integral over $A$, $R\subseteq A'$, and $A'$ is finite over $R$. The inclusion is local: some prime of $A'$ lies over $\mathfrak m_R$ \cite[Theorem~9.3(i)]{Mat86}, and it lies in $\mathfrak m_{A'}$, so $\mathfrak m_{A'}\cap R=\mathfrak m_R$. The residue field $\kappa$ of $A'$ is finite, being a quotient of $A'/XA'$, which is finite over $\mathbf F_p$.
\item \emph{Power series.} The complete discrete valuation ring $A'$ has characteristic $p$, and its residue field $\kappa$ is perfect, being finite. So $A'\cong\kappa[[T]]$ \cite[Chapter~II, \S4, Theorem~2]{Ser79}, and the inclusion $R\subseteq A'$ gives the homomorphism.
\end{enumerate}
\end{proof}
% curveNorm (CurveNorm.lean) is curveDVR (CurveDVR.lean) followed by (3). curveDVR: curveOverSeries for (1) (GlobalTheory.FormalIsogeny.finite_of_finite_quotient_map for the finiteness) and the private exists_dvr_quotient_of_complete_local for (2) (completeDVRExt, from GlobalTheory.FormalIsogeny.finite_integralClosure and isDiscreteValuationRing_integralClosure; curveDVR_complete for the completeness of A', finite over R; isLocalHom of an integral injective map). (3): the private curveNorm_coeffField (a coefficient field kappa -> A', by Algebra.FormallySmooth of the finite field kappa over F_p) and GlobalTheory.ArcAvoid.exists_ringEquiv_powerSeries_of_complete (T to a uniformizer).

\begin{lemma}\label{sha-automorphy-realization-components-transfer-arc-triangular}
Let $C$ be a local ring, $H$ a group, $n\ge1$, $r:H\to M_n(C)$ a homomorphism of monoids, $\chi_1,\dots,\chi_n:H\to C^\times$ homomorphisms and $g_0\in H$. Suppose that $\chi_i(g_0)-\chi_j(g_0)$ is a unit for $i\ne j$; that the characteristic polynomial of $r(g_0)$ reduces modulo $\mathfrak m_C$ to $\prod_i(X-\bar\chi_i(g_0))$, where the bar denotes reduction; and that $(r(h_1)-\chi_1(h_1))\cdots(r(h_n)-\chi_n(h_n))=0$ for all $h_1,\dots,h_n\in H$. Then there is an invertible matrix $P$ over $C$ such that, for every $h\in H$, $Pr(h)P^{-1}$ is upper triangular with diagonal entries $\chi_1(h),\dots,\chi_n(h)$.
\end{lemma}
\leanref{Development/AutomorphicArithmetic/AutomorphicVarieties/FrobeniusSmoothVectors/FlagLoc.lean}{UpperTheory.V.Fr.flag_loc}
\begin{proof}
Let $A=r(g_0)$ and $a_i=\chi_i(g_0)$, and let $e_i=\prod_{j\ne i}(a_i-a_j)^{-1}(A-a_j)$.
\begin{enumerate}
\item \emph{Idempotents.} Taking every $h_i$ equal to $g_0$ gives $\prod_j(A-a_j)=0$. So $(A-a_i)e_i=0$, that is, $Ae_i=e_iA=a_ie_i$; and $e_ie_j=0$ for $i\ne j$, as $e_ie_j$ is a multiple of $\prod_k(A-a_k)$. The polynomial $\sum_i\prod_{j\ne i}(X-a_j)(a_i-a_j)^{-1}-1$ has degree less than $n$ and vanishes at $a_1,\dots,a_n$, whose differences are units, so it is $0$, as the Vandermonde determinant of the $a_k$ is a unit. Hence $\sum_ie_i=1$, and $e_i=e_i\sum_je_j=e_i^2$.
\item \emph{Triangularity.} Let $h\in H$ and $c\ge i\ge d$. Take $h_j=g_0$ for $j\ne i$ and $h_i=h$, and multiply the product of the hypothesis by $e_c$ on the left and by $e_d$ on the right. For $j<i$, $e_c(A-a_j)=(a_c-a_j)e_c$, with $a_c-a_j$ a unit as $j\ne c$; for $j>i$, $(A-a_j)e_d=(a_d-a_j)e_d$, with $a_d-a_j$ a unit. So $e_c(r(h)-\chi_i(h))e_d=0$. With $c=i>d$ this gives $e_cr(h)e_d=\chi_c(h)e_ce_d=0$, and with $c=d=i$ it gives $e_ir(h)e_i=\chi_i(h)e_i$.
\item \emph{The $e_i$ are nonzero.} Modulo $\mathfrak m_C$, $\bar a_i$ is a root of the characteristic polynomial of $\bar A$, so $\bar Av=\bar a_iv$ for a nonzero vector $v$ over the residue field, and then $\bar e_iv=v$. So $e_i\ne0$.
\item \emph{Conjugation.} A nonzero idempotent matrix $e$ over $C$ has an entry that is a unit: otherwise $\det(1-e)\equiv1$ modulo $\mathfrak m_C$, so $1-e$ is invertible, and $e(1-e)=0$ gives $e=0$. Choose $(\alpha_i,\beta_i)$ such that $u_i=(e_i)_{\alpha_i\beta_i}$ is a unit. Let $P'$ be the matrix whose $k$-th column is the $\beta_k$-th column of $e_k$, and $P$ the matrix whose $j$-th row is $u_j^{-1}$ times the $\alpha_j$-th row of $e_j$. For every matrix $M$, $(PMP')_{jk}=u_j^{-1}(e_jMe_k)_{\alpha_j\beta_k}$. With $M=1$, (1) gives $PP'=1$, so $P'=P^{-1}$. With $M=r(h)$, (2) gives $(Pr(h)P^{-1})_{jk}=0$ for $j>k$ and $(Pr(h)P^{-1})_{jj}=u_j^{-1}\chi_j(h)u_j=\chi_j(h)$.
\end{enumerate}
\end{proof}
% flag_loc (FlagLoc.lean): ord_ext (the relation on the group algebra; only group elements are used here) and lag_triangular (lag_coi, mid_zero) for (1) and (2), lag_ne_zero (through lag_ne_zero_field) for (3), coi_conj_single (exists_unit_entry) and entry_of_conj for (4).

\begin{lemma}\label{sha-automorphy-realization-components-transfer-arc-flag}
In Situation~\ref{sha-automorphy-link-situation}, let $n\in\{2,3\}$, let $\mathcal O_1\supseteq\mathcal O$ be the ring of integers of a finite extension of $\mathbf Q_\ell$ such that the link, with its residual characters composed with the map of residue fields, satisfies the residual conditions with some constant, and let $A_n$, its universal linked point $(\rho^{\mathrm u},\psi^{\mathrm u})$, $B$ and $A_n^\circ$ be as in Definition~\ref{sha-automorphy-realization}(2). Let $C$ be a local ring, $f:A_n^\circ\to C$ a ring homomorphism and $v$ a place of the link above $\ell$. Then there is an invertible matrix $P$ over $C$ such that, for every $h\in G_v$, $Pf(\rho^{\mathrm u}(\mathrm{loc}_vh))P^{-1}$ is upper triangular with diagonal entries $f(\psi^{\mathrm u}_{v,1}(h)),\dots,f(\psi^{\mathrm u}_{v,n}(h))$.
\end{lemma}
\leanref{Development/AutomorphicArithmetic/AutomorphicVarieties/FrobeniusSmoothVectors/ZpowWNe.lean}{UpperTheory.V.Fr.flag}
\begin{proof}
We apply Lemma~\ref{sha-automorphy-realization-components-transfer-arc-triangular} to $H=G_v$, $r(h)=f(\rho^{\mathrm u}(\mathrm{loc}_vh))$ and $\chi_i=f\circ\psi^{\mathrm u}_{v,i}$. Its last hypothesis is condition (c) of Definition~\ref{sha-descent-link}(2) for the universal linked point, carried by $f$. By residual condition (7) of Definition~\ref{sha-descent-residual} there is $g_v\in G_v$ such that $a=\bar\Psi_\Gamma(\mathrm{loc}_vg_v)$ has $a^2\ne1$; let $g_0=g_v$.
\begin{enumerate}
\item \emph{Units.} For $i\ne j$, $w_{n,i}-w_{n,j}$ is $\pm1$ or $\pm2$, so $a^{w_{n,i}}\ne a^{w_{n,j}}$. The difference $\psi^{\mathrm u}_{v,i}(g_v)-\psi^{\mathrm u}_{v,j}(g_v)$ lies in $B$, and by condition (e) its residue is $a^{w_{n,i}}-a^{w_{n,j}}\ne0$. As $B$ is local with maximal ideal $B\cap\mathfrak m_{A_n}$ (Lemma~\ref{sha-automorphy-realization-presentation}(3)), it is a unit of $B$, hence of $A_n^\circ$, and its image under $f$ is a unit of $C$.
\item \emph{The characteristic polynomial.} By Lemma~\ref{sha-automorphy-realization-components-propagation-diagonal}(1), the characteristic polynomial of $\rho^{\mathrm u}(\mathrm{loc}_vg_v)$ is $\prod_i\big(X-\psi^{\mathrm u}_{v,i}(g_v)\big)$. Applying $f$ and reducing modulo $\mathfrak m_C$ gives the remaining hypothesis.
\end{enumerate}
\end{proof}
% flag (ZpowWNe.lean): hF.f3 for g_v; co_B and zpow_w_ne for (1); o1 for (2); then flag_loc with the field ord of the universal point.

\begin{lemma}\label{sha-automorphy-realization-components-transfer-arc}
In Situation~\ref{sha-automorphy-link-situation}, let $n\in\{2,3\}$, let a realization of rank $n$ with coefficients $\mathcal O_1$ be given, with $G_0$ the Galois group of its residual datum, $\pi:G_0\to\Gamma_0$ and $\kappa_0$ its map and its character, universal ordinary determinant ring $B_n$, and $A_n$, $B$ and $A_n^\circ$ as in Definition~\ref{sha-automorphy-realization}, and let $\phi:B_n\to B$ be a compatible map. Let $(\mathcal O',R,\iota)$ be a completion at a prime of $A_n^\circ$ (Definition~\ref{sha-automorphy-realization-components-transfer-primes}), and $\mathfrak q$ a prime of $R$ with $\dim R/\mathfrak q=1$ and $\ell\in\mathfrak q$. Then there are an arc $(\rho_y,\psi_y)$ over $\kappa[[T]]$ for the residual datum of the realization (Definition~\ref{sha-hecke-arcs}(2)) and a local $\mathcal O_1$-algebra map $j:R\to\kappa[[T]]$ with kernel $\mathfrak q$, such that $\rho_y(g)=\kappa_0(g)j(\rho_R(\pi(g)))$ for $g\in G_0$ and $\psi_{y,v,i}(h)=j(\psi_{R,v,i}(h))\kappa_0(h)$ for $h\in D_v$, and such that the arc is attached to $j\circ\iota\circ\phi$.
\end{lemma}
\leanref{Development/AutomorphicArithmetic/AutomorphicVarieties/CurveDVRSpecialization/PointEx.lean}{UpperTheory.V.Pr.pointEx}
\leanref{Development/AutomorphicArithmetic/AutomorphicVarieties/CurveDVRSpecialization/PointEx.lean}{UpperTheory.V.Pr.pointEx_of_le}
\begin{proof}
\begin{enumerate}
\item \emph{The curve.} As $\ell\in\mathfrak q$ and $\ell\ne0$ in $\mathcal O_1$, the inverse image of $\mathfrak q$ in the discrete valuation ring $\mathcal O_1$ is nonzero, hence maximal; so $\mathfrak q$ contains $\mathfrak m_{\mathcal O_1}$. The ring $R/\mathfrak q$ is a complete local noetherian domain of dimension one and characteristic $\ell$, with finite residue field. By Lemma~\ref{sha-automorphy-realization-components-transfer-arc-curve} there are a finite field $\kappa$ and an injective local homomorphism $j_0:R/\mathfrak q\to\kappa[[T]]$; let $j$ be its composite with $R\to R/\mathfrak q$, a local homomorphism with kernel $\mathfrak q$. The map $\mathcal O_1\to R/\mathfrak q$ kills $\mathfrak m_{\mathcal O_1}$, so composing it with $j_0$ gives a ring homomorphism from $k_1$ to $\kappa[[T]]$. Its values $y$ are constant: with $q$ the number of elements of $k_1$, $y^q=y$, so $y(0)^q=y(0)$; if $y-y(0)$ had order $m\ge1$, then $(y-y(0))^q=y-y(0)$ would have order $qm>m$. So it is a homomorphism $k_1\to\kappa$ followed by the inclusion of the constants; this makes $\kappa$ a field containing $k_1$, and $j$ an $\mathcal O_1$-algebra map.
\item \emph{The arc.} Let $\rho_y$ and $\psi_y$ be as in the statement. The pair $(j(\rho_R),j(\psi_R))$ is a linked point with values in the local ring $\kappa[[T]]$, continuous as $j$ is local, so by Lemma~\ref{sha-automorphy-realization-components-transfer-twist} and its proof, applied to the twist, $\rho_y$ has continuous entries, is trivial on $I_u\cap\bigcap_i\ker\bar\chi_i$ for $u\in S_{\mathrm a}$, and the $\psi_{y,v,i}$ are continuous. With $P=J_n$, condition (a) of Definition~\ref{sha-descent-link}(2) and $\kappa_0(g)\kappa_0(cgc^{-1})=\varepsilon(g)^{1-n}$ give ${}^t\rho_y(g)P\rho_y(cgc^{-1})=\varepsilon(g)^{1-n}P$. As $\phi$ is compatible, $j\circ\iota\circ\phi$ carries $(D^{\mathrm u},\psi^{\mathrm u})$ to the determinant of $\rho_y$ and the $\psi_{y,v,i}$; reducing modulo $T$, the determinant of $\rho_y$ reduces to $\bar D$ and the $\psi_{y,v,i}$ lift the $\bar\psi_{v,i}$, as $(D^{\mathrm u},\psi^{\mathrm u})$ does.
\item \emph{Flags.} Let $v$ be a place above $\ell$. By Lemma~\ref{sha-automorphy-realization-components-transfer-arc-flag}, applied to the local ring $\kappa[[T]]$ and $f=j\circ\iota$, there is an invertible matrix $P$ over $\kappa[[T]]$ such that $Pj(\rho_R(\mathrm{loc}_vh))P^{-1}$ is upper triangular with diagonal entries $j(\psi_{R,v,i}(h))$ for every $h\in G_v=D_v$. Multiplying by $\kappa_0(h)$, $P\rho_y(h)P^{-1}$ is upper triangular with diagonal entries $\psi_{y,v,i}(h)$. This holds over the fraction field of $\kappa[[T]]$ too, so $(\rho_y,\psi_y)$ is an arc.
\end{enumerate}
\end{proof}
% pointEx (PointEx.lean): T.lne and the maximality of the inverse image of p in O for (1), then pointEx_of_le: pe_curve (Pr.T.curveNorm for R/p, then pe_const for the constants) for (1); pe_maps_of, pe_maps for the attachment; pe_tw_cont, pe_form (the form Jgl), hred (pe_res), pe_ine (as px_ine) and the reduction of psi for (2); Fr.flag with f = j o L.iota and pe_flag for (3).

\begin{lemma}\label{sha-automorphy-realization-components-transfer-nice-separate}
In Situation~\ref{sha-automorphy-link-situation}, let $n\in\{2,3\}$, let a realization of rank $n$ with coefficients $\mathcal O_1$ be given, with $G_0$ the Galois group of its residual datum, $\pi:G_0\to\Gamma_0$ and $\kappa_0$ its map and its character, let $(\mathcal O',R,\iota)$ be a completion at a prime of $A_n^\circ$, and let $(\rho_y,\psi_y)$ be an arc over $\kappa[[T]]$ for the residual datum of the realization and $j:R\to\kappa[[T]]$ a local $\mathcal O_1$-algebra map with $\rho_y(g)=\kappa_0(g)j(\rho_R(\pi(g)))$ and $\psi_{y,v,i}(h)=j(\psi_{R,v,i}(h))\kappa_0(h)$, whose kernel $\mathfrak q$ is nice (Definition~\ref{sha-automorphy-realization-components-transfer-primes}); the relation between $\rho_y$ and $\rho_R$, and the locality of $j$, are not used. Let $U$ be an open subgroup of $G_0$ and $v$ a place above $\ell$. Then there is $g\in D_v$ whose image in $G_0$ lies in $U$ and in the subgroup $\Gamma'$ of Situation~\ref{sha-hecke-situation}, such that $\psi_{y,v,i}(g)^m\ne\psi_{y,v,i'}(g)^m$ for all $i\ne i'$ and $m\ge1$.
\end{lemma}
\leanref{Development/AutomorphicArithmetic/AutomorphicVarieties/CurveDVRSpecialization/NiceSep.lean}{UpperTheory.V.Pr.niceSep}
\leanref{Development/AutomorphicArithmetic/AutomorphicVarieties/CurveDVRSpecialization/SelfTwistEigenvalues/OneUnitTF.lean}{UpperTheory.V.Pr.T.oneUnitTF}
\leanref{Development/AutomorphicArithmetic/AutomorphicVarieties/CurveDVRSpecialization/CommonFree.lean}{UpperTheory.V.Pr.commonFree}
\begin{proof}
Identify $D_v$ with $G_v$. Let $N\subseteq D_v$ be the subgroup of the $h\in I_v$ with $\varepsilon(h)=1$, and $V\subseteq D_v$ the inverse image of $U$, which has finite index, as $U$ is open and $D_v$ is compact. For $i\ne i'$ let $\chi_{ii'}=\psi_{y,v,i}/\psi_{y,v,i'}$, a homomorphism from $D_v$ to $\kappa[[T]]^\times$.
\begin{enumerate}
\item \emph{Equal residues.} By the definition of a realization, $\psi_{y,v,i}(h)$ reduces to $\bar\Psi_\Gamma(\mathrm{loc}_vh)^{w_{n,i}}\bar\kappa_0(h)$. For $h\in N$, $\bar\Psi_\Gamma(\mathrm{loc}_vh)=1$ by residual condition (7) of Definition~\ref{sha-descent-residual}, so $\chi_{ii'}(h)\equiv1$ modulo $T$.
\item \emph{No torsion.} Let $x\in\kappa[[T]]$ with $x\equiv1$ modulo $T$ and $x^m=1$ for some $m\ge1$. Write $m=\ell^km'$ with $\ell\nmid m'$, and let $z=x^{\ell^k}$, so that $z-1=(x-1)^{\ell^k}\in T\kappa[[T]]$ and $z^{m'}=1$. Then $1+z+\dots+z^{m'-1}\equiv m'$ modulo $T$ is a unit, and $(z-1)(1+z+\dots+z^{m'-1})=0$, so $z=1$; hence $(x-1)^{\ell^k}=0$ and $x=1$. By (1), this applies to $x=\chi_{ii'}(h)$ for $h\in N$.
\item \emph{Nontrivial ratios.} Let $i\ne i'$. As $\mathfrak q$ is nice, some $h\in N$ has $\psi_{R,v,i}(h)-\psi_{R,v,i'}(h)\notin\mathfrak q$, so $\chi_{ii'}(h)\ne1$. As $V$ has finite index, $h^k\in V$ for some $k\ge1$, and $\chi_{ii'}(h^k)=\chi_{ii'}(h)^k\ne1$ by (2). So $\chi_{ii'}$ is nontrivial on $N\cap V$.
\item \emph{A common element.} If every $h\in N\cap V$ had $\chi_{ii'}(h)=1$ for some $i\ne i'$, then $N\cap V$ would be the union of the finitely many subgroups $N\cap V\cap\ker\chi_{ii'}$, so one of them would have finite index $d$ in $N\cap V$ \cite[Lemma~4.1]{Neu54}; then $\chi_{ii'}(h)^d=1$, hence $\chi_{ii'}(h)=1$ by (2), for every $h\in N\cap V$, which contradicts (3). So some $g\in N\cap V$ has $\chi_{ii'}(g)\ne1$, hence $\chi_{ii'}(g)^m\ne1$ for every $m\ge1$ by (2), for all $i\ne i'$.
\end{enumerate}
As $g\in V$, its image lies in $U$; as $g\in N$, $\varepsilon(g)=1$, so its image lies in $\Gamma'$.
\end{proof}
% niceSep (NiceSep.lean): N = (F.Nv (e v)).comap lam, V = U.comap phi; hred with T.ok.Psi_N and psib, chib of TwLink for (1); oneUnitTF (the private tf) for (2); hnice.sep and Subgroup.exists_pow_mem_of_index_ne_zero for (3); the private exists_element_separating_powers with commonFree (Subgroup.exists_finiteIndex_of_leftCoset_cover, Neumann's lemma) for (4); FreeRatios over y.Fr by injectivity.

\begin{lemma}\label{sha-automorphy-realization-components-transfer-nice-irreducible}
In Situation~\ref{sha-automorphy-link-situation}, let $n\in\{2,3\}$, let a realization of rank $n$ with coefficients $\mathcal O_1$ be given, with $G_0$ the Galois group of its residual datum, $\pi:G_0\to\Gamma_0$ and $\kappa_0$ its map and its character, let $(\mathcal O',R,\iota)$ be a completion at a prime of $A_n^\circ$, and let $(\rho_y,\psi_y)$ be an arc over $\kappa[[T]]$ for the residual datum of the realization and $j:R\to\kappa[[T]]$ a local $\mathcal O_1$-algebra map with $\rho_y(g)=\kappa_0(g)j(\rho_R(\pi(g)))$ and $\psi_{y,v,i}(h)=j(\psi_{R,v,i}(h))\kappa_0(h)$, whose kernel $\mathfrak q$ is nice (Definition~\ref{sha-automorphy-realization-components-transfer-primes}). Let $\bar{\mathcal K}$ be an algebraic closure of the fraction field of $\kappa[[T]]$, $\Gamma'$ the subgroup of Situation~\ref{sha-hecke-situation}, and $U$ an open subgroup of $G_0$.
\begin{enumerate}
\item If $\mathfrak q$ is off-triangular, then $0$ and $\bar{\mathcal K}^n$ are the only subspaces of $\bar{\mathcal K}^n$ stable under $\rho_y(g)$ for every $g\in U\cap\Gamma'$.
\item If $\mathfrak q$ is endoscopic, then $0$ and $\bar{\mathcal K}^2$ are the only subspaces of $\bar{\mathcal K}^2$ stable under the matrices $\begin{pmatrix}\rho_y(g)_{11}&\rho_y(g)_{13}\\\rho_y(g)_{31}&\rho_y(g)_{33}\end{pmatrix}$ for every $g\in U\cap\Gamma'$.
\end{enumerate}
\end{lemma}
\leanref{Development/AutomorphicArithmetic/AutomorphicVarieties/CurveDVRSpecialization/NiceIrr.lean}{UpperTheory.V.Pr.niceIrr}
\leanref{Development/AutomorphicArithmetic/AutomorphicVarieties/CurveDVRSpecialization/NiceIrrE.lean}{UpperTheory.V.Pr.niceIrrE}
\begin{proof}
\openstep{The irreducibility on the open subgroups of $\Gamma'$, at an off-triangular nice prime and for the corner at an endoscopic nice prime, is not yet written out.}
% niceIrr (NiceIrr.lean, with the open subgroups of MatIrrOpenGd.lean) for (1); niceIrrE (NiceIrrE.lean) for (2). Size: 547 lines by tools/lean_size.py for the two together.
\end{proof}

\begin{lemma}\label{sha-automorphy-realization-components-transfer-nice-type-e}
In Situation~\ref{sha-automorphy-link-situation}, let a realization of rank three with coefficients $\mathcal O_1$ be given, with $G_0$ the Galois group of its residual datum, $\pi:G_0\to\Gamma_0$ and $\kappa_0$ its map and its character, let $(\mathcal O',R,\iota)$ be a completion at a prime of $A_3^\circ$, and let $(\rho_y,\psi_y)$ be an arc over $\kappa[[T]]$ for the residual datum of the realization and $j:R\to\kappa[[T]]$ a local $\mathcal O_1$-algebra map with $\rho_y(g)=\kappa_0(g)j(\rho_R(\pi(g)))$ and $\psi_{y,v,i}(h)=j(\psi_{R,v,i}(h))\kappa_0(h)$, whose kernel $\mathfrak q$ is nice and endoscopic (Definition~\ref{sha-automorphy-realization-components-transfer-primes}). Then the arc satisfies (N2) for type E (Definition~\ref{sha-hecke-arcs}(3)).
\end{lemma}
\leanref{Development/AutomorphicArithmetic/AutomorphicVarieties/CurveDVRSpecialization/NiceTypeE.lean}{UpperTheory.V.Pr.niceTypeE}
\begin{proof}
\openstep{The decomposition of the arc at a nice endoscopic prime into the corner block and the middle character, with their forms, their irreducibility on $\Gamma'$ and the absence of equivariant maps between them, is not yet written out.}
% niceTypeE (NiceTypeE.lean): the corner block and the middle entry of rho_y, whose odd entries vanish modulo q; niceIrrE with U the whole group for the irreducibility of the corner. Size: 333 lines by tools/lean_size.py.
\end{proof}

\begin{lemma}\label{sha-automorphy-realization-components-transfer-nice}
In Situation~\ref{sha-automorphy-link-situation}, let $n\in\{2,3\}$, let a realization of rank $n$ with coefficients $\mathcal O_1$ be given, with $G_0$ the Galois group of its residual datum, $\pi:G_0\to\Gamma_0$ and $\kappa_0$ its map and its character, let $(\mathcal O',R,\iota)$ be a completion at a prime of $A_n^\circ$, and let $(\rho_y,\psi_y)$ be an arc over $\kappa[[T]]$ for the residual datum of the realization and $j:R\to\kappa[[T]]$ a local $\mathcal O_1$-algebra map with $\rho_y(g)=\kappa_0(g)j(\rho_R(\pi(g)))$ and $\psi_{y,v,i}(h)=j(\psi_{R,v,i}(h))\kappa_0(h)$, whose kernel $\mathfrak q$ is nice (Definition~\ref{sha-automorphy-realization-components-transfer-primes}). Then the arc satisfies (N1) and (N3) of Definition~\ref{sha-hecke-arcs}(3); it satisfies (N2) for type I if $\mathfrak q$ is off-triangular, and (N2) for type E if $\mathfrak q$ is endoscopic.
\end{lemma}
\leanref{Development/AutomorphicArithmetic/AutomorphicVarieties/CurveDVRSpecialization/NiceOf.lean}{UpperTheory.V.Pr.niceOf}
\begin{proof}
Let $\bar{\mathcal K}$ be an algebraic closure of the fraction field of $\kappa[[T]]$, and $\Gamma'$ the subgroup of Situation~\ref{sha-hecke-situation}. Identify $D_v$ with $G_v$.
\begin{enumerate}
\item \emph{(N1).} Let $v\mid\ell$. Lemma~\ref{sha-automorphy-realization-components-transfer-nice-separate}, with $U=G_0$, gives $g\in D_v$ with $\varepsilon(g)=1$ and $\psi_{y,v,i}(g)\ne\psi_{y,v,i'}(g)$ for $i\ne i'$. So $\psi_{y,v,i}\ne\psi_{y,v,i'}$, and, for $i<i'$, $\psi_{y,v,i}\ne\psi_{y,v,i'}\varepsilon$, as the two differ at $g$.
\item \emph{(N2) for type I.} If $\mathfrak q$ is off-triangular, this is Lemma~\ref{sha-automorphy-realization-components-transfer-nice-irreducible}(1) with $U=G_0$.
\item \emph{(N2) for type E.} If $\mathfrak q$ is endoscopic, this is Lemma~\ref{sha-automorphy-realization-components-transfer-nice-type-e}.
\item \emph{(N3).} Let $U\subseteq G_0$ be open with $s\in U\cap\Gamma'$, and $V$ an open subgroup of $G_0$ with $sV\subseteq U$. For every open subgroup $V'$ of $G_0$ and every $v\mid\ell$, Lemma~\ref{sha-automorphy-realization-components-transfer-nice-separate} gives $u\in D_v$ with image in $V'\cap\Gamma'$ at which no positive powers of the $\psi_{y,v,i}(u)$ agree; by the flag of the arc at $v$, the characteristic polynomial of $\rho_y(u)$ is $\prod_i(X-\psi_{y,v,i}(u))$. By Lemma~\ref{sha-automorphy-realization-components-transfer-nice-irreducible}, $\rho_y$ is irreducible over $\bar{\mathcal K}$ on every open subgroup of $\Gamma'$ if $\mathfrak q$ is off-triangular; if $\mathfrak q$ is endoscopic, the entries of $\rho_y$ of odd weight vanish, so $\rho_y$ is the sum of its corner block, irreducible on every open subgroup of $\Gamma'$, and its middle entry. \openstep{That these give an element of the coset $s(V\cap\Gamma')\subseteq U\cap\Gamma'$ at which the characteristic polynomial of $\rho_y$ is separable is not yet written out.}
% The private cosetI and cosetE of NiceSpecializationTwistSatisfiesThirdCondition.lean, with coset_in, strong_of (strongIrrOfOpen), linClosedOfAdic (the linear conditions on the entries of rho_y over the algebraic closure are closed), charpoly_flag, cornerHom, midHom and charpoly_block, and regCoset or regCosetE (RegCoset.lean, RegCosetE.lean). Size: 196 lines by tools/lean_size.py for regCoset, regCosetE, strongIrrOfOpen and linClosedOfAdic together.
\end{enumerate}
\end{proof}
% niceOf (NiceOf.lean): nice_specialization_twist_satisfies_first_condition (niceSep) for (1), niceIrr with U the whole group for (2), niceTypeE for (3), nice_specialization_twist_satisfies_third_condition for (4).

\begin{lemma}\label{sha-automorphy-realization-components-transfer-stable}
In Situation~\ref{sha-automorphy-link-situation}, let $n\in\{2,3\}$, let a realization of rank $n$ with coefficients $\mathcal O_1$ be given, with universal ordinary determinant ring $B_n$ and with $A_n$, $B$ and $A_n^\circ$ as in Definition~\ref{sha-automorphy-realization}, let $\phi:B_n\to B$ be a compatible map, $(\mathcal O',R,\iota)$ a completion at a prime of $A_n^\circ$ and $P$ an off-triangular prime of $R$ (Definition~\ref{sha-automorphy-realization-components-transfer-primes}). Then the prime $(\iota\circ\phi)^{-1}(P)$ of $B_n$ is stable (Definition~\ref{sha-hecke-arcs}(4)).
\end{lemma}
\leanref{Development/AutomorphicArithmetic/AutomorphicVarieties/CurveDVRSpecialization/StableOf.lean}{UpperTheory.V.Pr.stableOf}
\leanref{Development/AutomorphicArithmetic/AutomorphicVarieties/CommutativeAlgebraHeight/StableField.lean}{UpperTheory.V.CA.stable_field}
\begin{proof}
Let $G_0$, $\pi$ and $\kappa_0$ be as in Lemma~\ref{sha-automorphy-realization-components-transfer-twist}, $K$ an algebraic closure of the fraction field of $R/P$, and $f:A_n^\circ\to K$ the composite of $\iota$ with $R\to K$; its kernel is $\iota^{-1}(P)$.
\begin{enumerate}
\item \emph{The kernel of $f$ is off-triangular.} The entries of $\rho_R$ are the images under $\iota$ of those of $\rho^{\mathrm u}$. So if all the entries of $\rho^{\mathrm u}$ of one of the kinds in Definition~\ref{sha-automorphy-realization}(4) lay in $\iota^{-1}(P)$, all the entries of $\rho_R$ of that kind would lie in $P$.
\item \emph{Irreducibility.} By Lemma~\ref{sha-automorphy-unramified-realization-irreducible}, $0$ and $K^n$ are the only subspaces of $K^n$ stable under every $f(\rho^{\mathrm u}(h))$. As $\pi$ is onto and the $\kappa_0(g)$ are nonzero scalars, the same holds for the representation $\sigma:g\mapsto\kappa_0(g)f(\rho^{\mathrm u}(\pi(g)))$ of $G_0$.
\item \emph{Stability.} Since $\phi$ is compatible, $f\circ\phi$ carries $D^{\mathrm u}$ to the determinant of $\sigma$, and its kernel is $\mathfrak q=(\iota\circ\phi)^{-1}(P)$. By Burnside's theorem the $\sigma(g)$ span the $n\times n$ matrices over $K$, so some $\sigma(g_1),\dots,\sigma(g_{n^2})$ form a basis, and $\det\big(\operatorname{tr}\sigma(g_sg_{s'})\big)\ne0$, as the trace form on matrices is nondegenerate. This determinant is the image under $f\circ\phi$ of $\det\big(\operatorname{tr}_{D^{\mathrm u}}(g_sg_{s'})\big)$, which therefore lies outside $\mathfrak q$. So $\mathfrak q$ is stable.
\end{enumerate}
\end{proof}
% stableOf (StableOf.lean): exAlgCl for K; hTIS from loc_rho (the direction from A_n^circ to R) for (1); types (the direction of Lemma sha-automorphy-unramified-realization-irreducible) and TwLink.tw_irr for (2); TwLink.IsPhi.maps and CA.stable_field (through Det.irrOpenDet, the Gram determinants of Remark sha-patching-chain-burnside) for (3).

\begin{lemma}\label{sha-automorphy-realization-components-transfer-thick-irreducible}
Let $G$ be a group, $\kappa$ an algebraically closed field, $j:\kappa\to K$ a homomorphism to a field, $\sigma:G\to M_3(K)$ a homomorphism of monoids with $\sigma(g)_{21}=\sigma(g)_{23}=0$ for every $g$, and $\rho:G\to M_3(\kappa)$ a homomorphism of monoids such that $j$ carries the determinant of $\rho$ to that of $\sigma$. Then some subspace of $\kappa^3$ other than $0$ and $\kappa^3$ is stable under every $\rho(g)$.
\end{lemma}
\leanref{Development/AutomorphicArithmetic/AutomorphicVarieties/CurveDVRSpecialization/ThickAlg.lean}{UpperTheory.V.Pr.Thick.not_irr}
\begin{proof}
Suppose not. Extend $\rho$ and $\sigma$ to ring homomorphisms $\rho(\sum_ga_gg)=\sum_ga_g\rho(g)$ and $\sigma_j(\sum_ga_gg)=\sum_gj(a_g)\sigma(g)$ from $\kappa[G]$. As $j$ carries the determinant of $\rho$ to that of $\sigma$, the characteristic polynomial of $\sigma_j(y)$ is the image under $j$ of that of $\rho(y)$, for $y\in\kappa[G]$.
\begin{enumerate}
\item \emph{$\rho$ is onto.} By Burnside's theorem, as $\kappa$ is algebraically closed and only $0$ and $\kappa^3$ are stable, $\rho(\kappa[G])$ is the algebra of all $3\times3$ matrices over $\kappa$.
\item \emph{A character.} Every $\sigma_j(y)$ has $\sigma_j(y)_{21}=\sigma_j(y)_{23}=0$, so the $(2,2)$ entry of a product of two of them is the product of their $(2,2)$ entries, and $\psi(y)=\sigma_j(y)_{22}$ is a ring homomorphism $\kappa[G]\to K$. Expanding $\det(X-\sigma_j(y))$ along the second row shows that $\psi(y)$ is a root of the characteristic polynomial of $\sigma_j(y)$.
\item \emph{$\psi$ factors through $\rho$.} If $\rho(y)=0$, the characteristic polynomial of $\sigma_j(y)$ is $X^3$, so $\psi(y)^3=0$ and $\psi(y)=0$. By (1), $\psi=\psi'\circ\rho$ for a ring homomorphism $\psi'$ from the $3\times3$ matrices over $\kappa$ to $K$.
\item \emph{Contradiction.} Let $E_{ab}$ be the matrix units. For $a\ne b$, $\psi'(E_{ab})^2=\psi'(E_{ab}^2)=0$, so $\psi'(E_{aa})=\psi'(E_{ab})\psi'(E_{ba})=0$. Then $1=\psi'(E_{11}+E_{22}+E_{33})=0$.
\end{enumerate}
\end{proof}
% not_irr (ThickAlg.lean): repAlg_surjective_of_matIrr_fin_three (Schur's lemma and the Jacobson density theorem) for (1); th_row, mid and root_row for (2); charpoly_th and root_row for (3); the matrix units e a b with hsq and hd for (4).

\begin{lemma}\label{sha-automorphy-realization-components-transfer-thick-split}
Let $G$ be a group, $\kappa$ an algebraically closed field, $j:\kappa\to K$ a homomorphism to a field, $\sigma:G\to M_3(K)$ a homomorphism of monoids with $\sigma(g)_{21}=\sigma(g)_{23}=0$ for every $g$, and $D$ a determinant of $G$ of dimension three with values in $\kappa$ that $j$ carries to the determinant of $\sigma$. Let $I$ be a two-sided ideal of $\kappa[G]$ such that $\kappa[G]/I$ is a semisimple ring and $\sum_gj(a_g)\sigma(g)=0$ for every $\sum_ga_gg\in I$. Then there is $y\in K^3$ with $y_2=1$ and $\sigma(g)y=\sigma(g)_{22}y$ for every $g\in G$.
\end{lemma}
\leanref{Development/AutomorphicArithmetic/AutomorphicVarieties/CurveDVRSpecialization/ThickAlg.lean}{UpperTheory.V.Pr.Thick.split_of_ss}
\begin{proof}
Let $\sigma_j(\sum_ga_gg)=\sum_gj(a_g)\sigma(g)$, a ring homomorphism from $\kappa[G]$.
\begin{enumerate}
\item \emph{Roots.} Let $g\in G$. Expanding along the second row, $\sigma(g)_{22}$ is a root of the characteristic polynomial of $\sigma(g)$, which is the image under $j$ of the characteristic polynomial of $D$ at $g$. That polynomial splits over $\kappa$, so $\sigma(g)_{22}=j(a)$ for some $a\in\kappa$.
\item \emph{A complement.} As $\sigma_j$ kills $I$, $K^3$ is a module over $E=\kappa[G]/I$ through $\sigma_j$. The subspace $W$ of vectors $v$ with $v_2=0$ is a submodule, as the second row of each $\sigma_j(y)$ is zero off the diagonal. As $E$ is semisimple, $W$ has a complementary submodule $C$. Write the second basis vector as $w+c$ with $w\in W$ and $c\in C$; then $c_2=1$.
\item \emph{An eigenvector.} Let $g\in G$, and $a$ as in (1). The vector $\sigma(g)c-j(a)c=\sigma_j(g-a)c$ lies in $C$, and its second coordinate is $\sigma(g)_{22}c_2-j(a)c_2=0$, so it lies in $W\cap C=0$. Hence $\sigma(g)c=\sigma(g)_{22}c$, and $y=c$ will do.
\end{enumerate}
\end{proof}
% split_of_ss (ThickAlg.lean): charpoly_th, root_row and IsAlgClosed.splits for hrow (1); th_row, the submodule W and exists_isCompl for (2); the element single g 1 - single 1 a for (3).

\begin{lemma}\label{sha-automorphy-realization-components-transfer-thick}
In Situation~\ref{sha-automorphy-link-situation}, let a realization of rank three with coefficients $\mathcal O_1$ be given, with universal ordinary determinant ring $B_3$ and with $A_3$, $B$ and $A_3^\circ$ as in Definition~\ref{sha-automorphy-realization}, let $\phi:B_3\to B$ be a compatible map, $(\mathcal O',R,\iota)$ a completion at a prime of $A_3^\circ$ and $E$ a thick prime of $R$ (Definition~\ref{sha-automorphy-realization-components-transfer-primes}). Then the prime $(\iota\circ\phi)^{-1}(E)$ of $B_3$ is thick (Definition~\ref{sha-hecke-arcs}(4)).
\end{lemma}
\leanref{Development/AutomorphicArithmetic/AutomorphicVarieties/CurveDVRSpecialization/ThickOf.lean}{UpperTheory.V.Pr.thickOf}
\begin{proof}
Let $G_0$, $\pi$ and $\kappa_0$ be as in Lemma~\ref{sha-automorphy-realization-components-transfer-twist}, $r$ a map for $E$ as in Definition~\ref{sha-automorphy-realization-components-transfer-primes}(3), $\mathfrak q=(\iota\circ\phi)^{-1}(E)$, $f:B_3\to R/E$ the composite of $\iota\circ\phi$ with the projection, $\sigma(g)=\kappa_0(g)r(\pi(g))$ for $g\in G_0$, and $K'$ an algebraic closure of the fraction field of $R/E$.
\begin{enumerate}
\item \emph{The determinant.} As $\phi$ is compatible, $f$ carries $D^{\mathrm u}$ to the determinant of $g\mapsto\kappa_0(g)\rho_R(\pi(g))$ modulo $E$. As $E$ is endoscopic, the entries of $\rho_R$ of odd weight lie in $E$, so modulo $E$ this representation preserves the span of the first and third basis vectors and the span of the second. The matrices $r(g)$ agree with it on the entries of even weight and have $r(g)_{21}=r(g)_{23}=0$, so they preserve the span of the first and third basis vectors, with the same blocks on it and on the quotient. So $\sigma$ has the same characteristic polynomials, and $f$ carries $D^{\mathrm u}$ to the determinant of $\sigma$.
\item \emph{The algebra $E_0$.} Let $Q=R/E$, a complete local noetherian ring with finite residue field, and $\theta:B_3[G_0]\to M_3(Q)$ the ring homomorphism with $\theta(bg)=f(b)\sigma(g)$. We show that $\theta$ kills the closure of the two-sided ideal that defines $E_0$ in Definition~\ref{sha-hecke-arcs}(4). First, $f$ carries $\mathfrak m_{B_3}$ into $\mathfrak m_Q$: if $a\in\mathfrak m_{B_3}$ and $f(a)$ were a unit, then, with $q$ the number of elements of the residue field of $Q$, $f(1-a^{q-1})$ would have residue $0$, while $1-a^{q-1}$ is a unit. Recall that $f(\psi^{\mathrm u}_{v,i}(h))=\kappa_0(h)\psi_{R,v,i}(h)$ modulo $E$, as $\phi$ is compatible, and $\sigma(h)=\kappa_0(h)r(\mathrm{loc}_vh)$ for $h\in D_v=G_v$. For $g\in G_0$, $\theta(g^*)=\varepsilon(g)^{-2}\kappa_0(cg^{-1}c^{-1})r(c_\Gamma\pi(g)^{-1}c_\Gamma^{-1})=\kappa_0(g)r(c_\Gamma\pi(g)^{-1}c_\Gamma^{-1})$, as $\kappa_0(g)\kappa_0(cgc^{-1})=\varepsilon(g)^{-2}$; and $y\mapsto y^*$ reverses products and fixes $B_3$.
\begin{enumerate}
\item[(i)] \emph{The Cayley-Hamilton ideal.} It is generated by the coefficients of the elements $\chi_y(y)$, $y=t_1y_1+\dots+t_my_m$ with $y_i\in B_3[G_0]$, where $\chi_y$ is the characteristic polynomial of $D^{\mathrm u}$ at $y$ \cite[\S1.17]{Che14}. As $f$ carries $D^{\mathrm u}$ to the determinant of $\sigma$ by (1), $\theta$ carries $\chi_y(y)$ to $\chi(\theta(y))$, where $\chi$ is the characteristic polynomial of the matrix $\theta(y)$ over $Q[t_1,\dots,t_m]$; this is $0$ by the Cayley-Hamilton theorem.
\item[(ii)] \emph{The ordinary relations.} For $v\mid\ell$ and $h_1,h_2,h_3\in D_v$, $\theta$ carries $\prod_i(h_i-\psi^{\mathrm u}_{v,i}(h_i))$ to $\prod_i\kappa_0(h_i)\big(r(\mathrm{loc}_vh_i)-\psi_{R,v,i}(h_i)\big)$ modulo $E$, which is $0$ by condition (c) for $r$ in the definition of a thick prime. Its image under $y\mapsto y^*$ is the product $\prod_{i=3}^1(h_i^*-\psi^{\mathrm u}_{v,i}(h_i))$ in the reverse order, which $\theta$ carries to $\prod_{i=3}^1\kappa_0(h_i)\big(r(c_\Gamma(\mathrm{loc}_vh_i)^{-1}c_\Gamma^{-1})-\psi_{R,v,i}(h_i)\big)$ modulo $E$; this is $0$ by the second form of (c) in that definition, with the characters in the reverse order.
\item[(iii)] \emph{Inertia.} Let $u\in S_{\mathrm a}$ and $g\in I_u\cap\bigcap_i\ker\bar\chi_i$, and let $\gamma$ be $\theta(g)=\sigma(g)$ or $\theta(g^*)$. By the definition of a realization there is $m$ prime to $\ell$ with $\pi(g)^m=1$ and $\kappa_0(g)^m=1$, so $\gamma^m=1$. The characteristic polynomial of $\sigma(g)$ is the image under $f$ of that of $D^{\mathrm u}$ at $g$, by (1); by condition (1) of Situation~\ref{sha-hecke-situation} that of $\theta(g^*)$ is the same. As $f$ is local and $D^{\mathrm u}$ reduces to $\bar D$, with $\bar\chi_i(g)=1$, it reduces to $(X-1)^3$. So $\gamma=1$, by the argument of step (3) in the proof of Lemma~\ref{sha-automorphy-realization-components-transfer-twist}, and $\theta$ kills $g-1$ and $g^*-1$.
\item[(iv)] \emph{The closure.} Let $N\ge1$. The entries of $\sigma$ are continuous, as those of $r$ are and $\kappa_0$ and $\pi$ are continuous, so the $g$ with $\sigma(g)\equiv1$ modulo $\mathfrak m_Q^N$ form an open normal subgroup $H_N$ of $G_0$. The closure lies in the two-sided ideal generated by the elements of (i) to (iii), their images under $y\mapsto y^*$, $\mathfrak m_{B_3}^N$ and the $h-1$ with $h\in H_N$. The map $\theta$ followed by reduction modulo $\mathfrak m_Q^N$ kills all of these, by (i) to (iii) and as $f(\mathfrak m_{B_3})\subseteq\mathfrak m_Q$. So $\theta$ carries the closure into the matrices with entries in $\bigcap_N\mathfrak m_Q^N=0$ \cite[Theorem~8.10]{Mat86}.
\end{enumerate}
\item \emph{Not stable.} Choose an embedding of an algebraic closure $\bar\kappa(\mathfrak q)$ of the residue field of $\mathfrak q$ into $K'$ that extends the map induced by $f$. Call this embedding $j$, and view $\sigma$ with values in $K'$. By (1), $j$ carries $D^{\mathrm u}$ modulo $\mathfrak q$ to the determinant of $\sigma$, and $\sigma(g)_{21}=\sigma(g)_{23}=0$ for all $g$. If $D^{\mathrm u}$ modulo $\mathfrak q$ were, over $\bar\kappa(\mathfrak q)$, the determinant of a representation with no stable subspace other than $0$ and $\bar\kappa(\mathfrak q)^3$, this would contradict Lemma~\ref{sha-automorphy-realization-components-transfer-thick-irreducible}.
\item \emph{Not semisimple.} The ring $E_0\otimes_{B_3}\bar\kappa(\mathfrak q)$ is the quotient of $\bar\kappa(\mathfrak q)[G_0]$ by the two-sided ideal $I$ generated by the image of the ideal that defines $E_0$. For $z$ in that ideal and $\bar z$ its image, $\sum_gj(\bar a_g)\sigma(g)$, with $\bar z=\sum_g\bar a_gg$, is the image of $\theta(z)$ in $K'$, which is $0$ by (2). So the map $\sum_g\bar a_gg\mapsto\sum_gj(\bar a_g)\sigma(g)$ kills $I$. If $E_0\otimes_{B_3}\bar\kappa(\mathfrak q)$ were semisimple, Lemma~\ref{sha-automorphy-realization-components-transfer-thick-split}, with $D=D^{\mathrm u}$ modulo $\mathfrak q$, would give $y\in K'^3$ with $y_2=1$ and $\sigma(g)y=\sigma(g)_{22}y$ for all $g\in G_0$. Dividing by the nonzero scalars $\kappa_0(g)$, with $\pi$ onto, this would give $r(h)y=r(h)_{22}y$ for all $h\in\Gamma_0$, which the definition of a thick prime excludes.
\end{enumerate}
So $\mathfrak q$ is not stable and $E_0\otimes_{B_3}\bar\kappa(\mathfrak q)$ is not semisimple: $\mathfrak q$ is thick.
\end{proof}
% thickOf (ThickOf.lean): TwLink.IsPhi.maps and ofRep_twr for (1); the private kills for (2), with local_of_finite (f local), th_CH (i), the fields ord and ord_star of IsThickMod with th_star (ii), the field tame of TwLink with th_Ip, chi_res and charpoly_star (iii), th_closure (iv); exists_emb and not_irr for (3); th_comp, th_map and split_of_ss, then hs.nonsplit, for (4).

\begin{lemma}\label{sha-automorphy-realization-components-transfer-field-point}
In Situation~\ref{sha-automorphy-link-situation}, let $n\in\{2,3\}$, let $\mathcal O_1\supseteq\mathcal O$ be the ring of integers of a finite extension of $\mathbf Q_\ell$, with residue field $k_1$, such that the link, with its residual characters composed with the map of residue fields, satisfies the residual conditions with some constant, and let $A_n$, its universal linked point $(\rho^{\mathrm u},\psi^{\mathrm u})$, $B$ and $A_n^\circ$ be as in Definition~\ref{sha-automorphy-realization}(2). Let $C$ be a complete local noetherian $\mathcal O_1$-algebra with residue field $k_1$, $K$ a field with an injective $\mathcal O_1$-algebra map $C\to K$, and $(\rho,\psi)$ a homomorphism of monoids $\rho:\Gamma_0\to M_n(K)$ with characters $\psi_{v,i}:G_v\to K^\times$ satisfying (a), (b) and (c) of Definition~\ref{sha-descent-link}(2). Suppose that $(\rho,\psi)$ is \emph{integral over $C$}: the traces of the $\rho(g)$ are the images of the values of a function $t:\Gamma_0\to C$, continuous for the $\mathfrak m_C$-adic topology, with $t(g)\equiv\sum_i\bar\Psi_\Gamma(g)^{w_{n,i}}$; the $\psi_{v,i}$ are the images of continuous functions $G_v\to C$ with residues $\bar\Psi_\Gamma(\mathrm{loc}_vh)^{w_{n,i}}$; and there is $\zeta\in C$ with residue $\bar\Psi_\Gamma(g_0)-\bar\Psi_\Gamma(g_0)^{-1}$ such that the $(i,i)$ entry of $\rho(g_0)-\rho(c_\Gamma g_0^{-1}c_\Gamma^{-1})$ is $w_{n,i}\zeta$ for every $i$. Suppose also that some entry of $\rho$ below the diagonal and some entry above it are nonzero. Then there are $\mathcal O_1$-algebra maps $\xi:A_n^\circ\to K$ and $\beta:B\to C$ such that $\xi$ carries $(\rho^{\mathrm u},\psi^{\mathrm u})$ to $(\rho,\psi)$ and $\xi(b)$ is the image of $\beta(b)$ for $b\in B$.
\end{lemma}
\leanref{Development/AutomorphicArithmetic/AutomorphicVarieties/FrobeniusSmoothVectors/PointField.lean}{UpperTheory.V.Fr.point_field}
\begin{proof}
We check the hypotheses (i) to (iv) of Lemma~\ref{sha-automorphy-unramified-near-completion-points-cone}, for $C'=C$ and the field $K$, which is a $C$-algebra with $C\to K$ injective; that lemma then gives $\xi$ and $\beta$, which is the statement. Hypothesis (iv) is part of the integrality of $(\rho,\psi)$: the $\psi_{v,i}$ are the images of maps $G_v\to C$ that are continuous for the $\mathfrak m_C$-adic topology, that is, locally constant modulo every power of $\mathfrak m_C$, with residues $\bar\Psi_\Gamma(\mathrm{loc}_vh)^{w_{n,i}}$.
\openstep{That an integral point with an entry below and an entry above the diagonal that are nonzero satisfies hypotheses (i) to (iii) of Lemma~\ref{sha-automorphy-unramified-near-completion-points-cone} (diagonal entries from $C$ with the residual values, products of weight $0$ with a factor off the diagonal from $\mathfrak m_C$, and finitely generated spans of the entries, on which they are locally constant) is not yet written out.}
% point_field (PointField.lean) is alg (Lemma sha-automorphy-unramified-near-completion-points-cone) applied to PF.cone_of_integral; the Cone field psi is the field psi of Integral. Open: the Cone fields dg, cyc, fg and cont, from PF.res_of_integral (idem for the matrix units in the C-span of the rho(g), residues with the trace t and zeta of Integral) for dg and the off-diagonal products, conn (from OffTri bot) with fin for cyc, and trace_cont and artin_rees for fg and cont. Size: 598 lines by tools/lean_size.py for PF.cone_of_integral, of which 277 for PF.res_of_integral.
\end{proof}

\begin{lemma}\label{sha-automorphy-realization-components-transfer-untwist-trace}
In Situation~\ref{sha-automorphy-link-situation}, let $n\in\{2,3\}$, and let $\mathcal O_1\supseteq\mathcal O$ be the ring of integers of a finite extension of $\mathbf Q_\ell$, with residue field $k_1$, such that the link, with its residual characters composed with the map of residue fields, satisfies the residual conditions with some constant. Let $C$ be a complete local noetherian $\mathcal O_1$-algebra with residue field $k_1$, $K$ an algebraically closed field with an injective ring homomorphism $C\to K$, $b:\Gamma_0\to M_n(K)$ a homomorphism of monoids, $t:\Gamma_0\to C$ a function and $\chi_{v,i}:G_v\to C^\times$ homomorphisms, for the places $v$ above $\ell$ and $1\le i\le n$, such that:
\begin{enumerate}
\item $t$ is continuous for the $\mathfrak m_C$-adic topology, $t(x)$ has residue $\sum_i\bar\Psi_\Gamma(x)^{w_{n,i}}$, its image in $K$ is $\operatorname{tr}b(x)$, and $t(c_\Gamma x^{-1}c_\Gamma^{-1})=t(x)$, for $x\in\Gamma_0$;
\item each $\chi_{v,i}$ is continuous for the $\mathfrak m_C$-adic topology, and $\chi_{v,i}(h)$ has residue $\bar\Psi_\Gamma(\mathrm{loc}_vh)^{w_{n,i}}$;
\item $\operatorname{tr}\big(b(x)\prod_i(b(\mathrm{loc}_vh_i)-\chi_{v,i}(h_i))\big)=0$ for every $v$, $x\in\Gamma_0$ and $h_1,\dots,h_n\in G_v$, the $\chi_{v,i}(h_i)$ taken in $K$.
\end{enumerate}
Let $K'$ be an algebraically closed field, $\varphi:C\to K'$ a ring homomorphism, and $\rho':\Gamma_0\to M_n(K')$ a homomorphism of monoids with characters $\psi'_{v,i}:G_v\to K'^\times$ satisfying (a), (b) and (c) of Definition~\ref{sha-descent-link}(2), such that $0$ and $K'^n$ are the only $\rho'(\Gamma_0)$-stable subspaces of $K'^n$ and $\varphi(t(x))=\operatorname{tr}\rho'(x)$ for $x\in\Gamma_0$. Then there are a homomorphism of monoids $\rho:\Gamma_0\to M_n(K)$ with characters $\psi_{v,i}:G_v\to K^\times$ satisfying (a), (b) and (c), integral over $C$ in the sense of Lemma~\ref{sha-automorphy-realization-components-transfer-field-point}, such that $0$ and $K^n$ are the only $\rho(\Gamma_0)$-stable subspaces of $K^n$, and an invertible matrix $h$ over $K$, such that $\rho(x)=hb(x)h^{-1}$ for $x\in\Gamma_0$ and $\psi_{v,i}$ is the image of $\chi_{v,i}$ in $K$.
\end{lemma}
\leanref{Development/AutomorphicArithmetic/AutomorphicVarieties/FrobeniusSmoothVectors/OfTrace.lean}{UpperTheory.V.Fr.of_trace}
\begin{proof}
\openstep{The conjugation of a representation with integral traces and ordinary characters to a linked point, integral and irreducible, is not yet written out.}
% of_trace (OfTrace.lean). Size: 530 lines by tools/lean_size.py.
\end{proof}

\begin{lemma}\label{sha-automorphy-realization-components-transfer-untwist}
In Situation~\ref{sha-automorphy-link-situation}, let $n\in\{2,3\}$, let a realization of rank $n$ with coefficients $\mathcal O_1$ be given, with residue field $k_1$ of $\mathcal O_1$, with $G_0$ the Galois group of its residual datum, and with $\pi:G_0\to\Gamma_0$ and $\kappa_0$ its map and its character. Let $C$ be a complete local noetherian $\mathcal O_1$-algebra with residue field $k_1$ and $(D,\psi)$ an ordinary pair over $C$ for the residual datum of the realization, such that $g-1$ lies in the kernel of $D$ for $g\in I_{v_0}$. Let $K$ be an algebraically closed field with an injective $\mathcal O_1$-algebra map $C\to K$. Let $K'$ be an algebraically closed field, $\varphi:C\to K'$ an $\mathcal O_1$-algebra map, and $\rho':\Gamma_0\to M_n(K')$ a homomorphism of monoids with characters $\psi'_{v,i}:G_v\to K'^\times$ satisfying (a), (b) and (c) of Definition~\ref{sha-descent-link}(2), such that $0$ and $K'^n$ are the only $\rho'(\Gamma_0)$-stable subspaces of $K'^n$ and $\varphi(D)$ is the determinant of $g\mapsto\kappa_0(g)\rho'(\pi(g))$. Then there are a homomorphism of monoids $\rho:\Gamma_0\to M_n(K)$ and characters $\psi_{v,i}:G_v\to K^\times$ satisfying (a), (b) and (c), integral over $C$ in the sense of Lemma~\ref{sha-automorphy-realization-components-transfer-field-point}, such that $0$ and $K^n$ are the only $\rho(\Gamma_0)$-stable subspaces of $K^n$, $D$ is, over $K$, the determinant of $g\mapsto\kappa_0(g)\rho(\pi(g))$, and the characters of the pair are $h\mapsto\psi_{v,i}(h)\kappa_0(h)$.
\end{lemma}
\leanref{Development/AutomorphicArithmetic/AutomorphicVarieties/CurveDVRSpecialization/SeedOf.lean}{untwist}
\begin{proof}
Let $j:C\to K$ be the given map, $\operatorname{tr}_D$ the trace of $D$, and identify $D_v$ with $G_v$. Recall that the trace form $(X,Y)\mapsto\operatorname{tr}(XY)$ on the $n\times n$ matrices over a field is nondegenerate.
\begin{enumerate}
\item \emph{A semisimple representation.} As $K$ is algebraically closed, there is a semisimple representation $r_0:G_0\to M_n(K)$ whose determinant is the image of $D$ under $j$, and whose kernel on $K[G_0]$ is that of this determinant \cite[Theorem~2.12]{Che14}. So $\operatorname{tr}r_0(y)=j(\operatorname{tr}_D(y))$ for $y\in C[G_0]$, and $r_0(y)=0$ if $y$ lies in the kernel of $D$.
\item \emph{Irreducibility.} As $K'$ is algebraically closed and only $0$ and $K'^n$ are stable, the $\rho'(x)$, $x\in\Gamma_0$, span the matrices over $K'$ by Burnside's theorem, and so do the $\sigma'(g)=\kappa_0(g)\rho'(\pi(g))$, $g\in G_0$, as $\pi$ is onto. By the nondegeneracy of the trace form, some $g_1,\dots,g_{n^2}\in G_0$ have $\det\big(\operatorname{tr}(\sigma'(g_a)\sigma'(g_b))\big)\ne0$. This determinant is the image under $\varphi$ of $\det\big(\operatorname{tr}_D(g_ag_b)\big)$, as $\varphi$ carries $D$ to the determinant of $\sigma'$. So $\det\big(\operatorname{tr}_D(g_ag_b)\big)\ne0$, its image $\det\big(\operatorname{tr}(r_0(g_a)r_0(g_b))\big)$ under the injective $j$ is nonzero, and the $r_0(g)$ span the matrices over $K$. In particular a matrix $X$ over $K$ with $\operatorname{tr}(r_0(h)X)=0$ for all $h\in G_0$ is $0$.
\item \emph{Inertia at $v_0$.} For $g\in I_{v_0}$, $g-1$ lies in the kernel of $D$, so $r_0(g)=1$ by (1).
\item \emph{Untwisting.} Let $u(g)=\kappa_0(g)^{-1}r_0(g)$, a homomorphism of monoids, and $t'(g)=\kappa_0(g)^{-1}\operatorname{tr}_D(g)\in C$, so that $j(t'(g))=\operatorname{tr}u(g)$. By (2), $g$ lies in the kernel $N$ of $u$ if and only if $t'(hg)=t'(h)$ for every $h\in G_0$. So:
\begin{enumerate}
\item[(i)] $N$ is a closed normal subgroup, as $t'$ is continuous, being the product of $\kappa_0^{-1}$ and a coefficient of the continuous determinant $D$, and $C$ is separated;
\item[(ii)] $N$ is stable under $g\mapsto cgc^{-1}$: by condition (1) of Situation~\ref{sha-hecke-situation}, $\operatorname{tr}_D(g)=\varepsilon(g)^{1-n}\operatorname{tr}_D((cgc^{-1})^{-1})$, and with $\kappa_0(g)\kappa_0(cgc^{-1})=\varepsilon(g)^{1-n}$ this gives $t'((cgc^{-1})^{-1})=t'(g)$ for every $g$; for $g\in N$ it follows that $t'(h\,cgc^{-1})=t'\big(g^{-1}(chc^{-1})^{-1}\big)=t'\big((chc^{-1})^{-1}\big)=t'(h)$;
\item[(iii)] $I_{v_0}\subseteq N$, by (3), as $\kappa_0$ is trivial on $I_{v_0}$;
\item[(iv)] $N$ contains every $g\in I_u$, $u\in S_{\mathrm a}$, with $\bar\Psi_\Gamma(\pi(g))=1$. Indeed, by the definition of a realization there is $m$ prime to $\ell$ with $\pi(g)^m=1$ and $\kappa_0(g)^m=1$. Then $\bar\chi_i(g^m)=1$ for every $i$, so $g^m-1$ is killed by $\operatorname{tr}_D(h\,\cdot)$ for all $h$, by condition (3) of Situation~\ref{sha-hecke-situation}; by (1) and (2), $r_0(g^m)=1$, and $u(g)^m=1$. The characteristic polynomial $P$ of $D$ at $\kappa_0(g)^{-1}g$ is monic, $P(u(g))=0$ by (1) and the Cayley-Hamilton theorem, and $P$ reduces to $(X-1)^n$, as $\bar\chi_i(g)=\bar\kappa_0(g)$. As in step (3) of the proof of Lemma~\ref{sha-automorphy-realization-components-transfer-twist}, $1+X+\dots+X^{m-1}$ is a unit of $C[X]/(P)$, and $(X-1)(1+X+\dots+X^{m-1})=X^m-1$; mapping $X$ to $u(g)$ gives $u(g)=1$.
\end{enumerate}
By (i) to (iv), $N$ is one of the subgroups of the last condition on $\pi$ in Definition~\ref{sha-automorphy-realization}(2), so the kernel of $\pi$ lies in $N$.
% The field ker of TwLink (PropagationTheoremGL2.lean), applied in ker_le: ker pi lies in every closed normal subgroup N of G_0 stable under conj, containing the g in I_u with Psi(pi g) = 1 and containing I_0. It is the last condition on pi in Definition sha-automorphy-realization(2), proved for the realization of -exists(2) in its step (k).
Then $u$ is constant on the fibres of $\pi$, so $u=b\circ\pi$ for a homomorphism of monoids $b:\Gamma_0\to M_n(K)$: for $g,g'$ with $\pi(g)=\pi(g')$, $gg'^{-1}$ lies in the kernel of $\pi$, so $u(g)=u(g')$.
\item \emph{Characters.} Let $\chi_{v,i}(h)=\kappa_0(h)^{-1}\psi_{v,i}(h)$ for $h\in D_v$. These are continuous homomorphisms $G_v\to C^\times$, and $\chi_{v,i}(h)$ reduces to $\bar\Psi_\Gamma(\mathrm{loc}_vh)^{w_{n,i}}$, as $\psi_{v,i}$ lifts $\bar\psi_{v,i}=\bar\Psi_\Gamma(\mathrm{loc}_v\,\cdot)^{w_{n,i}}\bar\kappa_0$.
\item \emph{Traces.} Let $t(\pi(g))=t'(g)$; this is well defined, as $j(t'(g))=\operatorname{tr}b(\pi(g))$ and $j$ is injective. The function $t$ is continuous, as $t'$ is and $\pi$ is a quotient map, $G_0$ being compact; $t'(g)$ reduces to $\bar\kappa_0(g)^{-1}\sum_i\bar\chi_i(g)=\sum_i\bar\Psi_\Gamma(\pi(g))^{w_{n,i}}$; and $t(c_\Gamma x^{-1}c_\Gamma^{-1})=t(x)$ by (4)(ii). For $v$, $h\in G_0$ and $g_1,\dots,g_n\in D_v$, condition (2) of Situation~\ref{sha-hecke-situation} with (1) gives $\operatorname{tr}\big(r_0(h)\prod_i(r_0(g_i)-\psi_{v,i}(g_i))\big)=0$, so $\prod_i(r_0(g_i)-\psi_{v,i}(g_i))=0$ by (2); dividing by the units $\kappa_0(g_i)$, $\prod_i\big(b(\mathrm{loc}_vg_i)-\chi_{v,i}(g_i)\big)=0$, and so $\operatorname{tr}\big(b(x)\prod_i(b(\mathrm{loc}_vg_i)-\chi_{v,i}(g_i))\big)=0$ for $x\in\Gamma_0$. Finally $\varphi(t(\pi(g)))=\kappa_0(g)^{-1}\operatorname{tr}\sigma'(g)=\operatorname{tr}\rho'(\pi(g))$.
\item \emph{The point.} By (4) to (6), Lemma~\ref{sha-automorphy-realization-components-transfer-untwist-trace} applies to $b$, $t$, the $\chi_{v,i}$, $\varphi$ and $\rho'$, and gives $(\rho,\psi)$ satisfying (a), (b) and (c), integral over $C$, with only $0$ and $K^n$ stable, and an invertible matrix $h$ over $K$ with $\rho(x)=hb(x)h^{-1}$ and $\psi_{v,i}$ the image of $\chi_{v,i}$. Then $\kappa_0(g)\rho(\pi(g))=hr_0(g)h^{-1}$, so over $K$, $D$ is the determinant of $g\mapsto\kappa_0(g)\rho(\pi(g))$; and $\psi_{v,i}(h')\kappa_0(h')$ is the image of the character of the pair at $h'\in D_v$.
\end{enumerate}
\end{proof}
% untwist (SeedOf.lean, private): groupDet_exists_semisimple_representation (Chenevier's Theorem 2.12) for (1); hs' and span_of_spec (g_gram, gram_of_map) for (2); hI0 (kills_map) for (3); ker_le (mem_ker_iff with tt_map and eq_zero_of_tr; ac_closed with tt_cont; tt_star; ramkappa; tame_ker with T.unipTorsion) and desc for (4), through the field ker of TwLink; chi_ex for (5); tt, ht0, tt_res, tt_star, the ordinary relation through eq_zero_of_tr and prod_smul, and tt_map with phi for (6); Fr.of_trace and ofRep_conj for (7).

\begin{lemma}\label{sha-automorphy-realization-components-transfer-local-point}
In Situation~\ref{sha-automorphy-link-situation}, let $n\in\{2,3\}$, let $\mathcal O_1\supseteq\mathcal O$ be the ring of integers of a finite extension of $\mathbf Q_\ell$, with residue field $k_1$, such that the link, with its residual characters composed with the map of residue fields, satisfies the residual conditions with some constant, and let $A_n$, its universal linked point $(\rho^{\mathrm u},\psi^{\mathrm u})$, $B$ and $A_n^\circ$ be as in Definition~\ref{sha-automorphy-realization}(2). Let $C$ be a complete local noetherian ring with a local map $\mathcal O_1\to C$, whose residue field is finite over $k_1$, and $(\rho,\psi)$ a homomorphism of monoids $\rho:\Gamma_0\to M_n(C)$ with characters $\psi_{v,i}:G_v\to C^\times$ satisfying (a), (b) and (c) of Definition~\ref{sha-descent-link}(2), continuous for the $\mathfrak m_C$-adic topology, such that modulo $\mathfrak m_C$ the diagonal entries of $\rho(g)$ are $\bar\Psi_\Gamma(g)^{w_{n,i}}$ and the $\psi_{v,i}(h)$ are $\bar\Psi_\Gamma(\mathrm{loc}_vh)^{w_{n,i}}$. Then there is an $\mathcal O_1$-algebra map $\xi:A_n^\circ\to C$ that carries $(\rho^{\mathrm u},\psi^{\mathrm u})$ to $(\rho,\psi)$ and the maximal ideal of $B$ into $\mathfrak m_C$.
\end{lemma}
\leanref{Development/AutomorphicArithmetic/AutomorphicVarieties/FrobeniusSmoothVectors/PointLoc.lean}{UpperTheory.V.Fr.point_loc}
\begin{proof}
This is Lemma~\ref{sha-automorphy-unramified-near-completion-points}, proved below without using the present lemma: its hypotheses are among those here, and its conclusion is the statement.
% point_loc (PointLoc.lean) is also the Lean declaration of Lemma sha-automorphy-unramified-near-completion-points, written out there (CA.cl_sub, the Cone conditions with cyc_mem and cont_sub, alg, co_B).
\end{proof}

\begin{lemma}\label{sha-automorphy-realization-components-transfer}
In Situation~\ref{sha-automorphy-link-situation}, let $n\in\{2,3\}$, let a realization of rank $n$ with coefficients $\mathcal O_1$ be given, with $B_n$, $A_n$, $B$ and $A_n^\circ$ as in Definition~\ref{sha-automorphy-realization}, and let $J_n$ be an ideal of $B_n$ with the properties of Lemma~\ref{sha-automorphy-realization-ideal}. Let $\Lambda_B$, $\mathfrak p_B$, the completions $(\mathcal O',R,\iota)$ at the primes of $A_n^\circ$ with $(\rho_R,\psi_R)$ and $P_B$, and the off-triangular, endoscopic, nice and thick primes of $R$ be as in Definition~\ref{sha-automorphy-realization-components-transfer-primes}. Then:
\begin{enumerate}
\item[(a)] there is a compatible surjection $\phi:B_n\to B$;
\end{enumerate}
and for every compatible surjection $\phi:B_n\to B$:
\begin{enumerate}
\item[(b)] the maximal ideal of $B$ lies in the radical of $\phi(J_n)+\Lambda_B$;
\item[(c)] for every prime $x$ of $A_n^\circ$, every completion $(\mathcal O',R,\iota)$ at $x$ and every nice off-triangular prime $\mathfrak q$ of $R$ with $\phi(J_n)\subseteq\mathfrak q_B$, $\phi(J_n)\subseteq P_B$ for every prime $P\subseteq\mathfrak q$ of $R$;
\item[(d)] if $n=3$: for every prime $x$ of $A_3^\circ$, every completion $(\mathcal O',R,\iota)$ at $x$, and nice endoscopic primes $\mathfrak q,\mathfrak q'$ of $R$ joined by a chain of thickly linked nice endoscopic primes, if $\phi(J_3)\subseteq P_B$ for some off-triangular prime $P\subseteq\mathfrak q$, then $\phi(J_3)\subseteq P'_B$ for every off-triangular prime $P'\subseteq\mathfrak q'$;
\item[(e)] there are an off-triangular prime $\xi$ of $A_n^\circ$ and a prime $Z$ of $B$ with $\phi(J_n)\subseteq Z\subseteq\xi_B$ and $\dim B/Z=1+nd$;
\item[(f)] if $\phi(J_n)\subseteq\mathfrak p_B$ for every off-triangular prime $\mathfrak p$ of $A_n^\circ$, then $J_n$ has the property of Lemma~\ref{sha-automorphy-realization-components}(1).
\end{enumerate}
\end{lemma}
\leanref{Development/AutomorphicArithmetic/AutomorphicVarieties/CurveDVRSpecialization/PxIsCont.lean}{UpperTheory.V.Pr.phiEx}
\leanref{Development/AutomorphicArithmetic/AutomorphicVarieties/CurveDVRSpecialization/FinLamOf.lean}{UpperTheory.V.Pr.finLamOf}
\leanref{Development/AutomorphicArithmetic/AutomorphicVarieties/CurveDVRSpecialization/FirstPrimeConditionTransfer.lean}{UpperTheory.V.Pr.p1Of}
\leanref{Development/AutomorphicArithmetic/AutomorphicVarieties/CurveDVRSpecialization/SecondThirdPrimeConditionTransfer.lean}{UpperTheory.V.Pr.rankThree_deformationSetup_splitting_condition}
\leanref{Development/AutomorphicArithmetic/AutomorphicVarieties/CurveDVRSpecialization/SeedOf.lean}{UpperTheory.V.Pr.seedOf}
\leanref{Development/AutomorphicArithmetic/AutomorphicVarieties/CurveDVRSpecialization/GoodPtOf.lean}{UpperTheory.V.Pr.goodPtOf}
\begin{proof}
(a) Let $G_0$, $\pi$ and $\kappa_0$ be as in Lemma~\ref{sha-automorphy-realization-components-transfer-twist}. By that lemma, applied to $C=A_n$ and the universal linked point, the determinant of $g\mapsto\kappa_0(g)\rho^{\mathrm u}(\pi(g))$, with the characters $h\mapsto\psi^{\mathrm u}_{v,i}(h)\kappa_0(h)$, is an ordinary pair over $A_n$. The universal property of $B_n$ gives an $\mathcal O_1$-algebra map $f:B_n\to A_n$ that carries $(D^{\mathrm u},\psi^{\mathrm u})$ to it.
\begin{enumerate}
\item[(i)] \emph{$f(B_n)\subseteq B$.} For $g\in G_0$, the coefficients of the characteristic polynomial of $f(D^{\mathrm u})$ at $g$ are those of $\kappa_0(g)\rho^{\mathrm u}(\pi(g))$, which lie in $B$ by Lemma~\ref{sha-automorphy-realization-components-transfer-coefficients}, as $\kappa_0(g)\in\mathcal O_1$; and $f(\psi^{\mathrm u}_{v,i}(h))=\psi^{\mathrm u}_{v,i}(h)\kappa_0(h)\in B$. Let $x\in B_n$ and $N\ge0$. By Lemma~\ref{sha-automorphy-realization-components-transfer-dense}(2), $x-y\in\mathfrak m_{B_n}^N$ for some $y$ in the $\mathcal O_1$-subalgebra generated by these coefficients and values, so $f(y)\in B$; as $f$ is local, $f(x)-f(y)\in\mathfrak m_{A_n}^N$. As $B$ is the closure of an $\mathcal O_1$-subalgebra $D$, there is $z\in D$ with $f(y)-z\in\mathfrak m_{A_n}^N$, and then $f(x)-z\in\mathfrak m_{A_n}^N$. As $N$ is arbitrary, $f(x)\in B$. Let $\phi:B_n\to B$ be $f$ with values in $B$; it is compatible by construction.
\item[(ii)] \emph{$\phi$ is onto.} Let $h\in\Gamma_0$ and $g\in G_0$ with $\pi(g)=h$. Then $f(\operatorname{tr}_{D^{\mathrm u}}(g))=\kappa_0(g)\operatorname{tr}\rho^{\mathrm u}(h)$, so $\operatorname{tr}\rho^{\mathrm u}(h)=f\big(\kappa_0(g)^{-1}\operatorname{tr}_{D^{\mathrm u}}(g)\big)$ lies in $f(B_n)$, as $\kappa_0(g)^{-1}\in\mathcal O_1$; likewise $\psi^{\mathrm u}_{v,i}(h)=f\big(\kappa_0(h)^{-1}\psi^{\mathrm u}_{v,i}(h)\big)$. By Lemma~\ref{sha-automorphy-realization-components-transfer-onto}, applied to $C=B_n$ and $f$, $B\subseteq f(B_n)$.
\end{enumerate}
% phiEx (PxIsCont.lean): B.univ with px_pair; hmem (CA.exists_adjoin_coeff_character_sub_mem_maximalIdeal_pow, hlam with px_chi, px_coef and Fr.det_mem, hpsi with Space.psi_mem) for (i); hun and Fr.B_gen (with px_tr) for (ii).

(b) Let $Q$ be a prime of $B$ that contains $\phi(J_n)+\Lambda_B$; we show that $Q$ is the maximal ideal $\mathfrak m_B$ of the local ring $B$. Then $\mathfrak m_B$ is the only prime containing $\phi(J_n)+\Lambda_B$, and it lies in the radical. Let $P=\phi^{-1}(Q)$, a prime of $B_n$ that contains $J_n$ and the image of $\mathfrak m_{\mathcal O_1}$, and let $q$ be the number of elements of the residue field $k_1$ of $\mathcal O_1$.
\begin{enumerate}
\item[(i)] \emph{Powers of the characters.} Let $v\mid\ell$, $1\le i\le n$ and $g\in I_v$. As $\phi$ is compatible, $\phi(\psi^{\mathrm u}_{v,i}(g))=xy$, where $x$ is the value of the character $\psi^{\mathrm u}_{v,i}$ of the universal linked point at the image $h$ of $g$ under $\pi$, and $y\in\mathcal O_1^\times$ the value of $\kappa_0$ at $g$. Then $x^{q-1}$ is the value at $h^{q-1}\in I_v$, at which $\bar\Psi$ is $\bar\Psi(h)^{q-1}=1$, so $x^{q-1}-1\in\Lambda_B\subseteq Q$; and $y^{q-1}-1\in\mathfrak m_{\mathcal O_1}\subseteq Q$. As $x^{q-1}y^{q-1}-1=(x^{q-1}-1)y^{q-1}+(y^{q-1}-1)$, $\psi^{\mathrm u}_{v,i}(g)^{q-1}-1\in P$.
\item[(ii)] \emph{The Iwasawa algebra.} Consider two $\mathcal O_1$-algebra maps from $\Lambda$ to $B_n/P$: the map $\Lambda\to B_n$ followed by the projection, and the reduction $\Lambda\to k_1$ followed by the map $k_1\to B_n/P$ induced by $\mathcal O_1\to B_n/P$, which kills $\mathfrak m_{\mathcal O_1}$. They carry the universal family of characters of the $I_v$ to the same family. Indeed, at $g\in I_v$ the first gives the class of the unit $y=\psi^{\mathrm u}_{v,i}(g)$, and the second the class of $a\in\mathcal O_1$ with the residue of $y$. Let $s=\sum_{0\le k<q-1}y^ka^{q-2-k}$, so that $(y-a)s=y^{q-1}-a^{q-1}$; the residue of $s$ is $(q-1)\bar y^{q-2}=-\bar y^{q-2}\ne0$, so $s$ is a unit, and $y^{q-1}-a^{q-1}=(y^{q-1}-1)-(a^{q-1}-1)\in P$ by (i), as $a^{q-1}\equiv1$ modulo $\mathfrak m_{\mathcal O_1}$. So $y-a\in P$. By the universal property of $\Lambda$ (Situation~\ref{sha-hecke-situation}), the two maps are equal, so the first one kills $\mathfrak m_\Lambda$.
\item[(iii)] \emph{$P$ is maximal.} As $J_n\subseteq P$, $B_n/P$ is finite over $\Lambda$, and by (ii) the kernel of $\Lambda\to B_n/P$ is $\mathfrak m_\Lambda$, as it is proper. So the domain $B_n/P$ is integral over the field $\Lambda/\mathfrak m_\Lambda$, hence a field, and $P$ is maximal.
\item[(iv)] \emph{$Q=\mathfrak m_B$.} As $Q\subseteq\mathfrak m_B$, $P\subseteq\phi^{-1}(\mathfrak m_B)$, which is proper, so $P=\phi^{-1}(\mathfrak m_B)$ by (iii). As $\phi$ is onto, $Q=\phi(P)=\mathfrak m_B$.
\end{enumerate}
% (b) is finLamOf (FinLamOf.lean): hJQ, hLQ, hOQ; (i) is hpsi (hu, the private flo for kappa through hB); (ii) is the private flo_lam (key with flo_sub_mem, hkappa' with flo_cont, the uniqueness of LamObj.univ, toPair_spec); (iii) is the private flo_max (Ideal.isMaximal_of_isIntegral_of_isMaximal_comap'); (iv) is hM and hQmax. The Lean states (b) for every Lambda (LamObj) with B_n / J finite over it; Lambda_B is Space.lamId.

(c) Let $x$, $(\mathcal O',R,\iota)$, $\mathfrak q$ and $P$ be as in (c).
\begin{enumerate}
\item[(i)] As $\mathfrak q$ is nice, $\dim R/\mathfrak q=1$ and $\ell\in\mathfrak q$, so Lemma~\ref{sha-automorphy-realization-components-transfer-arc} gives an arc and a local map $j$ with kernel $\mathfrak q$, the arc being attached to $f=j\circ\iota\circ\phi$.
\item[(ii)] By Lemma~\ref{sha-automorphy-realization-components-transfer-nice}, the arc satisfies (N1), (N3) and (N2) for type I, as $\mathfrak q$ is off-triangular; it satisfies (N4) by Lemma~\ref{sha-automorphy-realization-ideal}. So it is nice of type I.
\item[(iii)] As $\phi(J_n)\subseteq\mathfrak q_B$, $\iota(\phi(J_n))\subseteq\mathfrak q=\ker j$, so $J_n\subseteq\ker f$. By Proposition~\ref{sha-patching-bound}(1), which applies to $J_n$ by Lemma~\ref{sha-automorphy-realization-ideal}, $J_n$ lies in every prime of $B_n$ contained in $\ker f=(\iota\circ\phi)^{-1}(\mathfrak q)$, in particular in $(\iota\circ\phi)^{-1}(P)$. So $\phi(J_n)\subseteq P_B$.
\end{enumerate}
% p1Of (FirstPrimeConditionTransfer.lean): pointEx for (i), niceOf with the field n4 of ModIn for (ii), the field p1 (P1core) and the equivalence key for (iii).

(d) Let $\mathfrak q=\mathfrak c_0,\mathfrak c_1,\dots,\mathfrak c_h=\mathfrak q'$ be nice endoscopic primes with $\mathfrak c_{i-1}$ and $\mathfrak c_i$ thickly linked, and thick primes $E_i\subseteq\mathfrak c_{i-1}\cap\mathfrak c_i$, for $1\le i\le h$.
\begin{enumerate}
\item[(i)] For each $i$, Lemma~\ref{sha-automorphy-realization-components-transfer-arc} gives an arc and a local map $j_i$ with kernel $\mathfrak c_i$, attached to $f_i=j_i\circ\iota\circ\phi$; by Lemma~\ref{sha-automorphy-realization-components-transfer-nice} and (N4), as in (c), it is nice of type E, as $\mathfrak c_i$ is endoscopic.
\item[(ii)] The primes $e_i=(\iota\circ\phi)^{-1}(E_i)$ of $B_3$ are thick by Lemma~\ref{sha-automorphy-realization-components-transfer-thick}, and $e_i\subseteq\ker f_{i-1}\cap\ker f_i$.
\item[(iii)] The prime $\mathfrak q_1=(\iota\circ\phi)^{-1}(P)$ is stable by Lemma~\ref{sha-automorphy-realization-components-transfer-stable}; it contains $J_3$, as $\phi(J_3)\subseteq P_B$, and it lies in $\ker f_0$, as $P\subseteq\mathfrak q$.
\item[(iv)] By Proposition~\ref{sha-patching-bound}(2), which applies to $J_3$ by Lemma~\ref{sha-automorphy-realization-ideal}, $J_3$ lies in every stable or thick prime of $B_3$ contained in $\ker f_h$. For an off-triangular prime $P'\subseteq\mathfrak q'$, $(\iota\circ\phi)^{-1}(P')$ is stable by Lemma~\ref{sha-automorphy-realization-components-transfer-stable} and lies in $\ker f_h$. So $\phi(J_3)\subseteq P'_B$.
\end{enumerate}
% rankThree_deformationSetup_splitting_condition (SecondThirdPrimeConditionTransfer.lean): chain_nat for the chain; pointEx (hpt) and niceOf with the hypothesis on N4 (hNE) for (i); thickOf and back_le for (ii); stableOf for (iii); the field p3 (P3core) and stableOf for (iv).

(e) By Lemma~\ref{sha-automorphy-realization-ideal}, the conclusion of Lemma~\ref{sha-automorphy-realization-seed} holds for $J_n$: there are an algebraically closed field $K'$ that is an $\mathcal O_1$-algebra, an $\mathcal O_1$-algebra map $x:B_n\to K'$ that kills $J_n$, and $(\rho',\psi')$ satisfying (a), (b) and (c) of Definition~\ref{sha-descent-link}(2), with $0$ and $K'^n$ the only $\rho'(\Gamma_0)$-stable subspaces, such that $x(D^{\mathrm u})$ is the determinant of $g\mapsto\kappa_0(g)\rho'(\pi(g))$.
\begin{enumerate}
\item[(i)] \emph{A component.} Let $Z_0\subseteq\ker x$ be a minimal prime over $J_n$, $C=B_n/Z_0$ with the projection $q:B_n\to C$, and $K$ an algebraic closure of the fraction field of $C$. The domain $C$ is a complete local noetherian $\mathcal O_1$-algebra with residue field $k_1$, and the image $(D,\psi)$ of $(D^{\mathrm u},\psi^{\mathrm u})$ is an ordinary pair over $C$. For $g\in I_{v_0}$, $g-1$ lies in the kernel of $D$, as it lies in that of the determinant over $B_n/J_n$ by Lemma~\ref{sha-automorphy-realization-ideal}, and $J_n\subseteq Z_0$. Also $x=x'\circ q$ for some $x':C\to K'$.
\item[(ii)] \emph{The untwisted point.} Lemma~\ref{sha-automorphy-realization-components-transfer-untwist}, with $\varphi=x'$, gives $(\rho,\psi)$ with values in $K$, satisfying (a), (b) and (c), integral over $C$ and irreducible, such that $D$ is the determinant of $g\mapsto\kappa_0(g)\rho(\pi(g))$ over $K$ and the characters of the pair are $h\mapsto\psi_{v,i}(h)\kappa_0(h)$.
\item[(iii)] \emph{Entries off the diagonal.} As $\rho$ is irreducible over the algebraically closed field $K$, the $\rho(g)$ span the matrices over $K$, by Burnside's theorem. If $\rho(g)_{ij}=0$ for all $g$, then $\operatorname{tr}(\rho(g)E_{ji})=0$ for all $g$, which the nondegeneracy of the trace form excludes. So some entry of $\rho$ below the diagonal and some entry above it are nonzero.
\item[(iv)] \emph{The point of $A_n^\circ$.} Lemma~\ref{sha-automorphy-realization-components-transfer-field-point} gives $\xi:A_n^\circ\to K$ and $\beta:B\to C$ with $\xi(\rho^{\mathrm u},\psi^{\mathrm u})=(\rho,\psi)$ and $\xi(b)$ the image of $\beta(b)$ for $b\in B$.
\item[(v)] \emph{$\beta\circ\phi=q$.} As $\phi$ is compatible, $\xi\circ\phi$ carries $D^{\mathrm u}$ to the determinant of $g\mapsto\kappa_0(g)\xi(\rho^{\mathrm u}(\pi(g)))$, which is $D$ over $K$, and $\psi^{\mathrm u}$ to the characters of the pair. So $\beta\circ\phi$ and $q$, followed by the injective map $C\to K$, carry $(D^{\mathrm u},\psi^{\mathrm u})$ to the same pair. As a determinant over $C$ is determined by the coefficients of its characteristic polynomials \cite[(1.5) and Corollary~1.14]{Che14}, which are determined by their images in $K$, $\beta\circ\phi$ and $q$ carry $(D^{\mathrm u},\psi^{\mathrm u})$ to the same ordinary pair over $C$, and $\beta\circ\phi=q$ by the universal property of $B_n$. In particular $\beta$ is onto.
\item[(vi)] \emph{An off-triangular prime.} The kernel $\xi_0$ of $\xi$ is off-triangular. Indeed, if all entries of $\rho^{\mathrm u}$ below the diagonal lay in $\xi_0$, the first basis vector would span a $\rho(\Gamma_0)$-stable line; if all entries above it did, the last basis vector would; and if $n=3$ and all entries of odd weight did, the second basis vector would; each contradicts the irreducibility of $\rho$.
\item[(vii)] \emph{The prime $Z$.} Let $Z=\ker\beta$, a prime of $B$, as $C$ is a domain. For $a\in J_n\subseteq Z_0$, $\beta(\phi(a))=q(a)=0$, so $\phi(J_n)\subseteq Z$. For $b\in Z$, $\xi(b)$ is the image of $\beta(b)=0$, so $Z\subseteq(\xi_0)_B$. As $\beta$ is onto, $B/Z\cong C=B_n/Z_0$, whose dimension is $1+nd$ by Lemma~\ref{sha-automorphy-realization-ideal}.
\end{enumerate}
% seedOf (SeedOf.lean): Ideal.exists_minimalPrimes_le, co_quot and exists_determinantPair_map_of_algHom for (i), with kills_map and the field unr (Unr0); untwist for (ii); offTri_of_irr (eq_zero_of_tr, irr_iff_span) for (iii); Fr.point_field for (iv); maps_A, gdet_inj and the uniqueness in T.B.univ for (v); Fr.types (the direction from MatIrr to TI) for (vi); hD (the field dim, DimT) with Ideal.quotientKerAlgEquivOfSurjective for (vii).

(f) Let $A$ be a complete discrete valuation ring with finite residue field and a local map $\mathcal O_1\to A$, $(\rho,\psi)$ a linked point of rank $n$ with values in $A$ lifting the residual one, and $j:A\to K'$ an injective map to a field such that $0$ and $K'^n$ are the only subspaces of $K'^n$ stable under every $j(\rho(g))$.
\begin{enumerate}
\item[(i)] Lemma~\ref{sha-automorphy-realization-components-transfer-local-point}, applied to $C=A$, whose residue field is finite, gives $\xi:A_n^\circ\to A$ that carries $(\rho^{\mathrm u},\psi^{\mathrm u})$ to $(\rho,\psi)$.
\item[(ii)] The kernel of $j\circ\xi$ is off-triangular, by the argument of (e)(vi). So $\phi(J_n)$ lies in it, by hypothesis, and $\xi\circ\phi$ kills $J_n$, as $j$ is injective.
\item[(iii)] As $\phi$ is compatible and $A_n^\circ\subseteq A_n$, $\phi$ carries $(D^{\mathrm u},\psi^{\mathrm u})$ to the determinant of $g\mapsto\kappa_0(g)\rho^{\mathrm u}(\pi(g))$ over $A_n^\circ$ and the characters $h\mapsto\psi^{\mathrm u}_{v,i}(h)\kappa_0(h)$. Applying $\xi$, $\xi\circ\phi$ carries them to the determinant of $g\mapsto\kappa_0(g)\rho(\pi(g))$ and the characters $h\mapsto\psi_{v,i}(h)\kappa_0(h)$.
\end{enumerate}
So $\xi\circ\phi:B_n\to A$ is a map as in Lemma~\ref{sha-automorphy-realization-components}(1).
% goodPtOf (GoodPtOf.lean): Fr.point_loc for (i); Fr.types for (ii); maps_of_inj (Det.injTest) and maps_comp for (iii).
% Open steps remain in Lemmas sha-automorphy-realization-components-transfer-nice (N3), -nice-irreducible, -nice-type-e, -field-point, -untwist (the kernel of pi) and -untwist-trace; -local-point cites Lemma sha-automorphy-unramified-near-completion-points.
\end{proof}

\begin{lemma}\label{sha-automorphy-realization-components-propagation-global}
Let $X$ be a noetherian topological space and $T$, $G$, $C$, $C_{\mathrm{cl}}$ and $E$ subsets of $X$. Call an irreducible component of $X$ that meets $T$ a \emph{$T$-component}. Suppose that:
\begin{enumerate}
\item $C_{\mathrm{cl}}\subseteq C$, and $C$ is preconnected;
\item every closed subset of $X$ that meets $C$ meets $C_{\mathrm{cl}}$;
\item $C$ lies in the closure of $C_{\mathrm{cl}}\setminus E$;
\item every $T$-component meets $C$;
\item for $x\in C_{\mathrm{cl}}$ and $T$-components $Y,Y'$ containing $x$, if $Y\subseteq G$ then $Y'\subseteq G$;
\item every point of $C_{\mathrm{cl}}\setminus E$ lies in a $T$-component;
\item some $T$-component lies in $G$.
\end{enumerate}
Then $T\subseteq G$.
\end{lemma}
\leanref{Development/AutomorphicArithmetic/AutomorphicVarieties/CurveDVRSpecialization/GlobProp.lean}{UpperTheory.V.Pr.globProp}
\leanref{Development/AutomorphicArithmetic/AutomorphicVarieties/CurveDVRSpecialization/FlipTwoThree.lean}{UpperTheory.V.Pr.GlobData.Ok}
\begin{proof}
\begin{enumerate}
\item[(i)] \emph{Every $T$-component lies in $G$.} Suppose not. Let $G'$ be the union of the $T$-components contained in $G$, and $B'$ the union of the others; as $X$ has finitely many irreducible components, each closed, $G'$ and $B'$ are closed. By (6), $C_{\mathrm{cl}}\setminus E\subseteq G'\cup B'$, so $C\subseteq G'\cup B'$ by (3). By (7) and (4), $C$ meets $G'$, and by the assumption and (4), $C$ meets $B'$. As $C$ is preconnected, $C\cap G'\cap B'\ne\emptyset$. By (2), the closed set $G'\cap B'$ contains a point $z\in C_{\mathrm{cl}}$. Then $z$ lies in a $T$-component $Y\subseteq G$ and in a $T$-component $Y'\not\subseteq G$, which contradicts (5).
\item[(ii)] \emph{Conclusion.} A point $x\in T$ lies in an irreducible component of $X$, which is a $T$-component, so $x\in G$ by (i).
\end{enumerate}
\end{proof}

\begin{lemma}\label{sha-automorphy-realization-components-propagation-local}
Let $R$ be a noetherian ring, $D_1\ge0$ an integer, $I$ an ideal of $R$, and $\mathcal T$, $\mathcal E$, $\mathcal G$, $\mathcal N$ and $\mathcal H$ sets of primes of $R$. Call the minimal primes over $I$ \emph{components}; call two components $P,Q$ \emph{adjacent} if $\dim R/(P+Q)\ge D_1$; and call two primes \emph{linked} if both lie in $\mathcal N\cap\mathcal E$ and both contain a common component that lies in $\mathcal H$. Suppose that:
\begin{enumerate}
\item a prime that contains a prime in $\mathcal G$ lies in $\mathcal G$, and a prime contained in a prime in $\mathcal T$ lies in $\mathcal T$;
\item $\mathcal N\subseteq\mathcal T\cup\mathcal E$;
\item every prime $P\in\mathcal G$ with $\dim R/P\ge D_1$ lies in some prime in $\mathcal N$;
\item every prime $P$ with $\dim R/P\ge D_1$ that contains a component in $\mathcal H$ lies in some prime in $\mathcal N\cap\mathcal E$;
\item every minimal prime of $R$ that lies in $\mathcal T$ is a component; any two components are joined by a chain of adjacent components; and every component lies in $\mathcal T$ or in $\mathcal H$;
\item for $\mathfrak p\in\mathcal N\cap\mathcal T\cap\mathcal G$, every prime contained in $\mathfrak p$ lies in $\mathcal G$;
\item for $\mathfrak p\in\mathcal N\cap\mathcal E$ and $\mathfrak q$ joined to $\mathfrak p$ by a chain of linked primes, if some prime contained in $\mathfrak p$ lies in $\mathcal T\cap\mathcal G$, then every prime in $\mathcal T$ contained in $\mathfrak q$ lies in $\mathcal G$.
\end{enumerate}
If some prime $P_1\in\mathcal G\cap\mathcal T$ has $\dim R/P_1\ge D_1$, then every minimal prime of $R$ that lies in $\mathcal T$ lies in $\mathcal G$.
\end{lemma}
\leanref{Development/AutomorphicArithmetic/AutomorphicVarieties/CurveDVRSpecialization/LocProp.lean}{UpperTheory.V.Pr.locProp}
\leanref{Development/AutomorphicArithmetic/AutomorphicVarieties/CurveDVRSpecialization/FlipTwoThree.lean}{UpperTheory.V.Pr.LocData.Ok}
\begin{proof}
\begin{enumerate}
\item[(i)] \emph{Nice primes.} Let $\mathfrak p\in\mathcal N$ contain a prime $P\in\mathcal T\cap\mathcal G$. Then every prime in $\mathcal T$ contained in $\mathfrak p$ lies in $\mathcal G$: if $\mathfrak p\in\mathcal T$, then $\mathfrak p\in\mathcal G$ by (1) and (6) applies; otherwise $\mathfrak p\in\mathcal E$ by (2), and (7) applies with $\mathfrak q=\mathfrak p$.
\item[(ii)] \emph{An invariant.} Say that a component $Q$ is \emph{good} if either $Q\in\mathcal T\cap\mathcal G$, or $Q\in\mathcal H\setminus\mathcal T$ and there are $\mathfrak p_0,\mathfrak p\in\mathcal N\cap\mathcal E$ such that some prime contained in $\mathfrak p_0$ lies in $\mathcal T\cap\mathcal G$, $Q\subseteq\mathfrak p$, and $\mathfrak p$ is joined to $\mathfrak p_0$ by a chain of linked primes.
\item[(iii)] \emph{Above a good component.} Let $Q$ be a good component and $M\supseteq Q$ a prime with $\dim R/M\ge D_1$. Then $M$ lies in some $\mathfrak p\in\mathcal N$ such that every prime in $\mathcal T$ contained in $\mathfrak p$ lies in $\mathcal G$. If $Q\in\mathcal T\cap\mathcal G$, then $M\in\mathcal G$ by (1), (3) gives $\mathfrak p\supseteq M$ in $\mathcal N$, and (i) applies with $P=Q$. Otherwise, with $\mathfrak p_0$ and $\mathfrak p$ as in (ii), (4) gives $\mathfrak p'\supseteq M$ in $\mathcal N\cap\mathcal E$; $\mathfrak p$ and $\mathfrak p'$ are linked through $Q$, so $\mathfrak p'$ is joined to $\mathfrak p_0$ by a chain of linked primes, and (7) applies.
\item[(iv)] \emph{Adjacent components.} Let $Q,Q'$ be adjacent components with $Q$ good. As $\dim R/(Q+Q')\ge D_1$, some prime $M\supseteq Q+Q'$ has $\dim R/M\ge D_1$: the first prime of a chain of length at least $D_1$ of primes containing $Q+Q'$. Let $\mathfrak p\supseteq M$ be as in (iii). If $Q'\in\mathcal T$, then $Q'\subseteq\mathfrak p$ gives $Q'\in\mathcal G$, and $Q'$ is good. Otherwise $Q'\in\mathcal H$ by (5), and (4) gives $\mathfrak p'\supseteq M$ in $\mathcal N\cap\mathcal E$, which contains $Q'$. If $Q\in\mathcal T\cap\mathcal G$, take $\mathfrak p_0=\mathfrak p'$ and the trivial chain. If not, with $\mathfrak p_0$ and $\mathfrak p_1$ as in (ii) for $Q$, the primes $\mathfrak p_1$ and $\mathfrak p'$ are linked through $Q$, and the chain from $\mathfrak p_0$ extends to $\mathfrak p'$. In both cases $Q'$ is good.
\item[(v)] \emph{Conclusion.} Let $P_0\subseteq P_1$ be a minimal prime of $R$. Then $P_0\in\mathcal T$ by (1), so $P_0$ is a component by (5). By (3), $P_1$ lies in some $\mathfrak p\in\mathcal N$, and by (i), with $P=P_1$, $P_0\in\mathcal G$; so $P_0$ is good. A minimal prime $P\in\mathcal T$ of $R$ is a component joined to $P_0$ by a chain of adjacent components, by (5), so it is good by (iv); as $P\in\mathcal T$, $P\in\mathcal G$.
\end{enumerate}
\end{proof}
% locProp (LocProp.lean): hW for (i); Inv for (ii); K for (iii); step with exists_prime_qdim for (iv); the final induction on Relation.ReflTransGen for (v). The hypotheses are the fields good_up, ti_down, nice_type, nice_good, nice_thick, comp, chain, dich, p1 and p23 of LocData.Ok, with sharp = I, D1, TI, TE, Good, Nice and Thick; the conclusion is LocData.Concl.

\begin{lemma}\label{sha-automorphy-realization-components-propagation-charts}
Let $X$ be a topological space. Let $C$ be a ring, $j:\operatorname{Spec}C\to X$ an open embedding, $R$ a local ring with a flat ring map $\iota:C\to R$, $f=j\circ\operatorname{Spec}(\iota):\operatorname{Spec}R\to X$ and $x=f(\mathfrak m_R)$. Then $f$ is continuous; for every minimal prime $P$ of $R$ the closure of $\{f(P)\}$ is an irreducible component of $X$; and every irreducible component of $X$ that contains $x$ is such a closure.
Now let $T\subseteq X$ be open, $G\subseteq X$ closed, $x\in X$, $R$ a noetherian local ring and $f:\operatorname{Spec}R\to X$ a map with the properties just stated: $f$ is continuous, $f(\mathfrak m_R)=x$, the closure of $\{f(P)\}$ is an irreducible component of $X$ for every minimal prime $P$ of $R$, and every irreducible component of $X$ that contains $x$ is such a closure. Let $T_R=f^{-1}(T)$, $G_R=f^{-1}(G)$ and $D_1\ge0$ an integer, and suppose that $(\ast)$: if some prime $P\in G_R\cap T_R$ has $\dim R/P\ge D_1$, then every minimal prime of $R$ that lies in $T_R$ lies in $G_R$. Then:
\begin{enumerate}
\item if $\dim R/P\ge D_1$ for every minimal prime $P$ of $R$, then for irreducible components $Y,Y'$ of $X$ that meet $T$ and contain $x$, $Y\subseteq G$ implies $Y'\subseteq G$;
\item if every minimal prime of $R$ lies in $T_R$, then $x$ lies in an irreducible component of $X$ that meets $T$;
\item if some prime $P\in G_R\cap T_R$ has $\dim R/P\ge D_1$, then some irreducible component of $X$ that meets $T$ lies in $G$.
\end{enumerate}
\end{lemma}
\leanref{Development/AutomorphicArithmetic/AutomorphicVarieties/CurveDVRSpecialization/ChartOfFlat.lean}{UpperTheory.V.Pr.chartOfFlat}
\leanref{Development/AutomorphicArithmetic/AutomorphicVarieties/CurveDVRSpecialization/GlobOfCharts.lean}{UpperTheory.V.Pr.globOfCharts}
\begin{proof}
\begin{enumerate}
\item[(i)] \emph{The first statement.} The map $f$ is continuous, as $\operatorname{Spec}(\iota)$ and $j$ are. As $\iota$ is flat, it has the going-down property \cite[Theorem~9.5]{Mat86}: for a prime $P$ of $R$ and a prime $q\subseteq\iota^{-1}(P)$ of $C$ there is a prime $P'\subseteq P$ of $R$ with $\iota^{-1}(P')=q$.
First, let $Z\subseteq X$ be irreducible and $c$ a prime of $C$ with $j(c)\in Z$. Then there is a prime $q\subseteq c$ of $C$ with $Z\subseteq\overline{\{j(q)\}}$. Indeed, $Z\cap j(\operatorname{Spec}C)$ is a nonempty open subset of $Z$, so it is irreducible and dense in $Z$; as $j$ is an open embedding, $j^{-1}(Z)$ is irreducible, and the closure of $j^{-1}(Z)$ in $\operatorname{Spec}C$ is an irreducible closed set, with a generic point $q$. Then $c\in\overline{\{q\}}$, that is $q\subseteq c$, and $Z\cap j(\operatorname{Spec}C)\subseteq j(\overline{\{q\}})\subseteq\overline{\{j(q)\}}$, as $j$ is continuous; so $Z\subseteq\overline{\{j(q)\}}$.
Let $P$ be a minimal prime of $R$, and $Z\supseteq\overline{\{f(P)\}}$ an irreducible subset of $X$. With $c=\iota^{-1}(P)$, $j(c)=f(P)\in Z$, so there is $q\subseteq\iota^{-1}(P)$ with $Z\subseteq\overline{\{j(q)\}}$; going down gives $P'\subseteq P$ with $\iota^{-1}(P')=q$, and $P'=P$ as $P$ is minimal. So $Z\subseteq\overline{\{f(P)\}}$, and the irreducible closed set $\overline{\{f(P)\}}$ is an irreducible component of $X$.
Conversely, let $Y$ be an irreducible component of $X$ that contains $x=j(\iota^{-1}(\mathfrak m_R))$. There is $q\subseteq\iota^{-1}(\mathfrak m_R)$ with $Y\subseteq\overline{\{j(q)\}}$; going down gives a prime $P'$ of $R$ with $f(P')=j(q)$; and a minimal prime $P\subseteq P'$ has $f(P')\in\overline{\{f(P)\}}$, as $f$ is continuous. So $Y\subseteq\overline{\{f(P)\}}$, which is irreducible, and $Y=\overline{\{f(P)\}}$ as $Y$ is an irreducible component.
\item[(ii)] \emph{Openness.} In the second statement, if $\overline{\{f(P)\}}$ meets the open set $T$, then $f(P)\in T$. And $T_R$ is open, as $f$ is continuous.
\item[(iii)] \emph{Proof of (1).} By the last property of $f$, $Y=\overline{\{f(P)\}}$ and $Y'=\overline{\{f(P')\}}$ for minimal primes $P,P'$ of $R$, which lie in $T_R$ by (ii). As $f(P)\in Y\subseteq G$, $P\in G_R$, and $\dim R/P\ge D_1$; so $P'\in G_R$ by $(\ast)$, and $Y'\subseteq G$, as $G$ is closed.
\item[(iv)] \emph{Proof of (2).} Some irreducible component $Y$ contains $x$; by the last property of $f$, $Y=\overline{\{f(P)\}}$ with $P$ minimal, and $f(P)\in T$.
\item[(v)] \emph{Proof of (3).} Let $P_0\subseteq P$ be a minimal prime of $R$. As $P$ lies in the closure of $P_0$ and in the open set $T_R$, $P_0\in T_R$, and $P_0\in G_R$ by $(\ast)$. By the third property of $f$, $\overline{\{f(P_0)\}}$ is an irreducible component; it meets $T$ at $f(P_0)$ and lies in the closed set $G$.
\end{enumerate}
\end{proof}
% chartOfFlat (ChartOfFlat.lean): the first statement, Chart (j (comap (closedPoint R))) (j o comap), for an open embedding j of Spec C into X; Algebra.HasGoingDown (iff_generalizingMap_primeSpectrumComap) and the private exp_product_chart_local (the generic point q, by QuasiSober, and subset_closure_inter_of_isPreirreducible_of_isOpen) for (i). The propagation uses it with C = A_n^circ and j the identity. globOfCharts (GlobOfCharts.lean): the second statement, for any Chart f at x, with T = M.TI open and G = M.Good closed; hmeet for (ii), then the three parts.

\begin{lemma}\label{sha-automorphy-realization-components-propagation-geometry-lift}
In Situation~\ref{sha-automorphy-link-situation}, let $n\in\{2,3\}$, let $\mathcal O_1\supseteq\mathcal O$ be the ring of integers of a finite extension of $\mathbf Q_\ell$ such that the link, with its residual characters composed with the map of residue fields, satisfies the residual conditions with some constant, and let $A_n$, its universal linked point, $B$ and $A_n^\circ$ be as in Definition~\ref{sha-automorphy-realization}(2), with the parts $(A_n)_j$ of Lemma~\ref{sha-automorphy-realization-presentation}. Let $Z$ be a prime of $B$ and $\xi$ an off-triangular prime of $A_n^\circ$ (Definition~\ref{sha-automorphy-realization}(4)) with $Z\subseteq\xi\cap B$. Then there is an off-triangular prime $\eta$ of $A_n^\circ$, generated by elements that lie in the parts $(A_n)_j$, with $\eta\cap B=Z$, such that the fraction field of $A_n^\circ/\eta$ has transcendence degree $1$ over that of $B/Z$.
\end{lemma}
\leanref{Development/AutomorphicArithmetic/AutomorphicVarieties/FrobeniusSmoothVectors/FrobeniusCornerLiftTwoToThree.lean}{UpperTheory.V.Fr.frobenius_corner_lift_two_to_three}
\begin{proof}
\openstep{The construction of a homogeneous off-triangular prime above $Z$ with transcendence degree one is not yet written out.}
% frobenius_corner_lift_two_to_three (FrobeniusCornerLiftTwoToThree.lean). Size: 202 lines by tools/lean_size.py.
\end{proof}

\begin{lemma}\label{sha-automorphy-realization-components-propagation-geometry-dimension}
In Situation~\ref{sha-automorphy-link-situation}, let $n\in\{2,3\}$, let $\mathcal O_1\supseteq\mathcal O$ be the ring of integers of a finite extension of $\mathbf Q_\ell$ such that the link, with its residual characters composed with the map of residue fields, satisfies the residual conditions with some constant, and let $A_n$, its universal linked point, $B$ and $A_n^\circ$ be as in Definition~\ref{sha-automorphy-realization}(2). For a prime $\eta$ of $A_n^\circ$ let $t(\eta)$ be the transcendence degree of the fraction field of $A_n^\circ/\eta$ over that of $B/(\eta\cap B)$. Let $x$ be a prime of $A_n^\circ$ and $(\mathcal O',R,\iota)$ a completion at $x$ (Definition~\ref{sha-automorphy-realization-components-transfer-primes}(2)). Then:
\begin{enumerate}
\item for every prime $P$ of $R$ and every integer $t\ge t(\iota^{-1}(P))$, $\dim R/P\le\dim B/P_B+t$;
\item for every prime $\eta\subseteq x$ of $A_n^\circ$ with $t(\eta)$ finite, there is a prime $P$ of $R$ with $\iota^{-1}(P)=\eta$ and $\dim R/P=\dim B/(\eta\cap B)+t(\eta)$.
\end{enumerate}
\end{lemma}
\leanref{Development/AutomorphicArithmetic/AutomorphicVarieties/FrobeniusSmoothVectors/LocDim.lean}{UpperTheory.V.Fr.loc_dim}
\begin{proof}
\openstep{The comparison of the dimensions of the primes of a completion with those of $B$ and the transcendence degrees is not yet written out.}
% loc_dim (LocDim.lean). Its first part is used, for the completion R_x, in Lemma sha-automorphy-unramified-near-completion(2). Size: 2,794 lines by tools/lean_size.py.
\end{proof}

\begin{lemma}\label{sha-automorphy-realization-components-propagation-geometry}
In Situation~\ref{sha-automorphy-link-situation}, let $n\in\{2,3\}$, let $\mathcal O_1\supseteq\mathcal O$ be the ring of integers of a finite extension of $\mathbf Q_\ell$ such that the link, with its residual characters composed with the map of residue fields, satisfies the residual conditions with some constant, and let $A_n$, its universal linked point $(\rho^{\mathrm u},\psi^{\mathrm u})$, $B$ and $A_n^\circ$ be as in Definition~\ref{sha-automorphy-realization}(2), with the notation of Definition~\ref{sha-automorphy-realization-components-transfer-primes}. Let $X=\operatorname{Spec}A_n^\circ$ and $X_{\mathrm c}\subseteq X$ the set of primes $\mathfrak p$ with $\mathfrak m_B\subseteq\mathfrak p_B$ outside which lies some entry of $\rho^{\mathrm u}$ below the diagonal. Then:
\begin{enumerate}
\item if $n=2$, $X_{\mathrm c}$ is preconnected; if $n=3$, $X_{\mathrm c}$ is nonempty and preconnected, and it lies in the closure of the set of its primes outside which lies some entry of $\rho^{\mathrm u}$ of odd weight;
\item for every closed $C$ and open $U$ of $X$ such that $C\cap U\cap X_{\mathrm c}$ is nonempty, some maximal ideal of $A_n^\circ$ lies in $C\cap U\cap X_{\mathrm c}$;
\item every minimal prime of $A_n^\circ$ is generated by its elements that lie in the parts $(A_n)_j$ of Lemma~\ref{sha-automorphy-realization-presentation}; and every prime of $A_n^\circ$ generated by such elements, outside which lies some entry of $\rho^{\mathrm u}$ below the diagonal, lies in a maximal ideal of $A_n^\circ$ that belongs to $X_{\mathrm c}$;
\item for every prime $Z$ of $B$ and every off-triangular prime $\xi$ of $A_n^\circ$ with $Z\subseteq\xi_B$, there is a maximal ideal $x\in X_{\mathrm c}$ of $A_n^\circ$ such that every completion $(\mathcal O',R,\iota)$ at $x$ has an off-triangular prime $P$ with $P_B=Z$ and $\dim R/P=\dim B/Z+1$.
\end{enumerate}
\end{lemma}
\leanref{Development/AutomorphicArithmetic/AutomorphicVarieties/FrobeniusSmoothVectors/PreconnectedRankTwo.lean}{UpperTheory.V.Fr.preconnected_rank_two}
\leanref{Development/AutomorphicArithmetic/AutomorphicVarieties/FrobeniusSmoothVectors/ExistsNonzeroLowCocycle.lean}{UpperTheory.V.Fr.preconnected_rank_three}
\leanref{Development/AutomorphicArithmetic/AutomorphicVarieties/FrobeniusSmoothVectors/FibreClosed.lean}{UpperTheory.V.Fr.jac}
\leanref{Development/AutomorphicArithmetic/AutomorphicVarieties/FrobeniusSmoothVectors/MinHom.lean}{UpperTheory.V.Fr.min_hom}
\leanref{Development/AutomorphicArithmetic/AutomorphicVarieties/FrobeniusSmoothVectors/MonMemPow.lean}{UpperTheory.V.Fr.l41}
\leanref{Development/AutomorphicArithmetic/AutomorphicVarieties/FrobeniusSmoothVectors/Seedpt.lean}{UpperTheory.V.Fr.seedpt}
\begin{proof}
Write $\rho=\rho^{\mathrm u}$ and $A=A_n^\circ$, and call an element of $A$ \emph{homogeneous} if it lies in some part $(A_n)_j$ of Lemma~\ref{sha-automorphy-realization-presentation}. By that lemma, $B=(A_n)_0$ is local with maximal ideal $\mathfrak m_B=B\cap\mathfrak m_{A_n}$, $A$ is a finitely generated $B$-algebra, hence noetherian, and the direct sum of the $(A_n)_j$, each a finitely generated $B$-module, with $(A_n)_j(A_n)_{j'}\subseteq(A_n)_{j+j'}$; and $\rho(g)_{ik}\in(A_n)_{w_{n,i}-w_{n,k}}$. An ideal generated by homogeneous elements contains the homogeneous components of each of its elements.
\begin{enumerate}
\item \emph{Connectedness.} \openstep{The connectedness of $X_{\mathrm c}$, and, in rank three, its nonemptiness and the density in it of the primes outside which lies some entry of odd weight, are not yet written out.}
% preconnected_rank_two (PreconnectedRankTwo.lean) and preconnected_rank_three (ExistsNonzeroLowCocycle.lean). Size: 1,150 lines by tools/lean_size.py for the two together.
\item \emph{Maximal ideals.} Let $\mathfrak p\in C\cap U\cap X_{\mathrm c}$. As $\mathfrak p\cap B$ is a proper ideal containing $\mathfrak m_B$, it is $\mathfrak m_B$. The set $U'$ of the primes in $U$ outside which lies some entry of $\rho$ below the diagonal is open and contains $\mathfrak p$, so there is $f\in A$ with $\mathfrak p\in D(f)\subseteq U'$. The domain $A/\mathfrak p$ is a finitely generated algebra over the field $B/\mathfrak m_B$, so its zero ideal is the intersection of its maximal ideals \cite[Theorem~5.5]{Mat86}. As $f\notin\mathfrak p$, some maximal ideal of $A/\mathfrak p$ does not contain the image of $f$; its inverse image $x$ is a maximal ideal of $A$ with $\mathfrak p\subseteq x$ and $f\notin x$. Then $x\in C$, as $C$ is closed; $x\in U'$; and $x\cap B\supseteq\mathfrak m_B$. So $x\in C\cap U\cap X_{\mathrm c}$.
\item \emph{Minimal primes.} Let $\eta$ be a minimal prime of $A$ and $\eta^*$ the ideal generated by the homogeneous elements of $\eta$. Then $\eta^*$ is prime. Indeed, let $a,b\notin\eta^*$ with $ab\in\eta^*$. Removing from $a$ and $b$ their homogeneous components that lie in $\eta$, which lie in $\eta^*$, we may assume that every homogeneous component of $a$ and of $b$ lies outside $\eta$. If $a_s$ and $b_t$ are their components of least degree, the component of degree $s+t$ of $ab$ is $a_sb_t$, which lies in $\eta^*\subseteq\eta$; this contradicts the primality of $\eta$. As $\eta^*\subseteq\eta$ and $\eta$ is minimal, $\eta=\eta^*$.

Now let $\eta$ be a prime generated by homogeneous elements, outside which lies some entry of $\rho$ below the diagonal, and let $\mathfrak a=\eta+\mathfrak m_BA$.
\begin{enumerate}
\item[(i)] \emph{Graded Nakayama.} If $(A_n)_j\subseteq\mathfrak a$, then $(A_n)_j\subseteq\eta$. Indeed, taking components of degree $j$, which is possible as $\eta$ and $\mathfrak m_BA$ are generated by homogeneous elements and $\mathfrak m_B\subseteq(A_n)_0$, gives $(A_n)_j\subseteq(\eta\cap(A_n)_j)+\mathfrak m_B(A_n)_j$. So the finitely generated $B$-module $M=(A_n)_j/(\eta\cap(A_n)_j)$ has $M=\mathfrak m_BM$, and $M=0$ by Nakayama's lemma.
\item[(ii)] \emph{The entries below the diagonal.} Let $\mathfrak t$ be the ideal generated by the entries of $\rho$ below the diagonal. Suppose $\mathfrak t$ lay in the radical of $\mathfrak a$. As $\mathfrak t$ is finitely generated, $\mathfrak t^M\subseteq\mathfrak a$ for some $M$. An entry $r=\rho(g)_{ik}$ with $i>k$ has weight $w_{n,i}-w_{n,k}\le-1$, so $r^{2M}\in(A_n)_j$ with $j=2M(w_{n,i}-w_{n,k})\le-2M$. By Lemma~\ref{sha-automorphy-realization-presentation}(2), $(A_n)_j$ is spanned over $B$ by products of entries of total weight $j$. The weights of the entries are at least $-2$, so such a product has at least $M$ factors of negative weight, which are entries below the diagonal; so it lies in $\mathfrak t^M$. Hence $(A_n)_j\subseteq\mathfrak a$, so $r^{2M}\in\eta$ by (i), and $r\in\eta$. This contradicts the choice of $\eta$.
\item[(iii)] \emph{Conclusion.} By (ii), some prime $\mathfrak p\supseteq\mathfrak a$ does not contain $\mathfrak t$, so $\mathfrak p$ lies in $X_{\mathrm c}$ and in the closed set of primes containing $\eta$. By (2), with this closed set and $U=X$, a maximal ideal of $A$ that contains $\eta$ lies in $X_{\mathrm c}$.
\end{enumerate}
\item \emph{The seed.} Let $Z$ and $\xi$ be as in (4). Lemma~\ref{sha-automorphy-realization-components-propagation-geometry-lift} gives an off-triangular prime $\eta$ of $A$, generated by homogeneous elements, with $\eta\cap B=Z$, such that the fraction field of $A/\eta$ has transcendence degree $1$ over that of $B/Z$. By (3) there is a maximal ideal $x\in X_{\mathrm c}$ of $A$ with $\eta\subseteq x$. Let $(\mathcal O',R,\iota)$ be a completion at $x$. By Lemma~\ref{sha-automorphy-realization-components-propagation-geometry-dimension}(2), there is a prime $P$ of $R$ with $\iota^{-1}(P)=\eta$ and $\dim R/P=\dim B/Z+1$. Then $P_B=\eta\cap B=Z$, and $P$ is off-triangular by Lemma~\ref{sha-automorphy-realization-components-propagation-completions}(2), as $\eta$ is.
\end{enumerate}
\end{proof}
% (2) is jac (FibreClosed.lean): gr_A, co_B and the private fibre_closed (isJacobsonRing_of_finiteType, PrimeSpectrum.isBasis_basic_opens). (3) is min_hom (MinHom.lean: the private aux, with Ideal.IsPrime.homogeneousCore) and l41 (MonMemPow.lean: the private gr_nak for (i), mon_mem_pow and w_lt for (ii), jac for (iii)). (4) is seedpt (Seedpt.lean): frobenius_corner_lift_two_to_three, l41, loc_dim and loc_rho.

\begin{lemma}\label{sha-automorphy-realization-components-propagation-completions}
In Situation~\ref{sha-automorphy-link-situation}, let $n\in\{2,3\}$, let $\mathcal O_1\supseteq\mathcal O$ be the ring of integers of a finite extension of $\mathbf Q_\ell$ such that the link, with its residual characters composed with the map of residue fields, satisfies the residual conditions with some constant $C_0$, and let $A_n$, its universal linked point $(\rho^{\mathrm u},\psi^{\mathrm u})$, $B$ and $A_n^\circ$ be as in Definition~\ref{sha-automorphy-realization}(2), with the notation of Definition~\ref{sha-automorphy-realization-components-transfer-primes}. Let $x$ be a prime of $A_n^\circ$. Then:
\begin{enumerate}
\item if $x$ is maximal and $\mathfrak m_B\subseteq x_B$, there is a completion at $x$;
\item for every completion $(\mathcal O',R,\iota)$ at $x$: $\iota$ is flat, $x$ is maximal and $\mathfrak m_B\subseteq x_B$; $(\rho_R,\psi_R)$ is continuous, and its diagonal entries and characters reduce to those of the residual linked point; a prime $P$ of $R$ contains all entries of $\rho_R$ below the diagonal, all entries above it, or all entries of odd weight, if and only if $\iota^{-1}(P)$ contains the corresponding entries of $\rho^{\mathrm u}$; and if some entry of $\rho^{\mathrm u}$ below the diagonal lies outside $x$, then the entries of $\rho_R$ below the diagonal generate $R$ and those above it lie in $\mathfrak m_R$.
\end{enumerate}
Suppose moreover that the link is large for $C_0$ and that some entry of $\rho^{\mathrm u}$ below the diagonal lies outside $x$, and let $(\mathcal O',R,\iota)$ be a completion at $x$. Then:
\begin{enumerate}
\setcounter{enumi}{2}
\item if $n=2$, every minimal prime $P$ of $R$ is off-triangular with $\dim R/P\ge2d+2$, and any two minimal primes are joined by a chain of minimal primes in which consecutive primes $P',Q'$ have $\dim R/(P'+Q')\ge2d+1$;
\item if $n=3$, there is an ideal $I$ of $R$ such that: $I=0$ if some entry of $\rho^{\mathrm u}$ of odd weight lies outside $x$; $I\subseteq P$ for every prime $P$ of $R$ outside which lies some entry of $\rho_R$ of odd weight; every minimal prime $P$ over $I$ has $\dim R/P\ge3d+2$, and any two are joined by a chain of minimal primes over $I$ in which consecutive primes $P',Q'$ have $\dim R/(P'+Q')\ge3d+1$; and every minimal prime over $I$ is off-triangular or thick;
\item for every ideal $I'$ of $B$ such that $\mathfrak m_B$ lies in the radical of $I'+\Lambda_B$, and every prime $P$ of $R$ with $I'\subseteq P_B$ and $\dim R/P\ge nd+1$, there is a nice prime of $R$ containing $P$.
\end{enumerate}
\end{lemma}
\leanref{Development/AutomorphicArithmetic/AutomorphicVarieties/FrobeniusSmoothVectors/ResEqOf.lean}{UpperTheory.V.Fr.loc_ex}
\leanref{Development/AutomorphicArithmetic/AutomorphicVarieties/FrobeniusSmoothVectors/LocFlat.lean}{UpperTheory.V.Fr.loc_flat}
\leanref{Development/AutomorphicArithmetic/AutomorphicVarieties/FrobeniusSmoothVectors/LocRho.lean}{UpperTheory.V.Fr.loc_rho}
\leanref{Development/AutomorphicArithmetic/AutomorphicVarieties/FrobeniusSmoothVectors/ChainedBot.lean}{UpperTheory.V.Fr.sharp2}
\leanref{Development/AutomorphicArithmetic/AutomorphicVarieties/FrobeniusSmoothVectors/ChainedOfComap.lean}{UpperTheory.V.Fr.sharp3}
\leanref{Development/AutomorphicArithmetic/AutomorphicVarieties/CurveDVRSpecialization/ExistsNicePrimeAboveOfLargeDim.lean}{UpperTheory.V.Pr.exists_nice_prime_above_of_large_dim}
\begin{proof}
\openstep{The facts about the completions of $A_n^\circ$ are not yet written out: their existence and flatness, the comparison of the entries, the chains of minimal primes and their types in ranks two and three, and the nice primes above primes of large dimension.}
% loc_ex (ResEqOf.lean) for (1); loc_flat (LocFlat.lean) and loc_rho (LocRho.lean) for (2); sharp2 (ChainedBot.lean: Chained bot (2d+1)) for (3); sharp3 (ChainedOfComap.lean: Loc.Sharp, with the fields triv, off, chained (Chained I (3d+1)) and dich) for (4); exists_nice_prime_above_of_large_dim (ExistsNicePrimeAboveOfLargeDim.lean, through NicePrimeEx) for (5). Size: 11,860 lines by tools/lean_size.py for the six declarations together (loc_ex 762, loc_flat 751, loc_rho 168, sharp2 5,590, sharp3 10,660, exists_nice_prime_above_of_large_dim 2,049).
\end{proof}

\begin{lemma}\label{sha-automorphy-realization-components-propagation-endoscopic-corner}
In Situation~\ref{sha-automorphy-link-situation}, let $\mathcal O_1\supseteq\mathcal O$ be the ring of integers of a finite extension of $\mathbf Q_\ell$ such that the link, with its residual characters composed with the map of residue fields, satisfies the residual conditions with some constant, and for $n\in\{2,3\}$ let $A_n$, its universal linked point $(\rho_n,\psi_n)$, $B^{(n)}=B$ and $A_n^\circ$ be as in Definition~\ref{sha-automorphy-realization}(2). Let $\mathfrak e$ be the ideal of $A_3^\circ$ generated by the entries of $\rho_3$ of odd weight, and suppose $\mathfrak e\ne A_3^\circ$. Let $K=A_3^\circ/\mathfrak e$ and $C=B^{(3)}/(\mathfrak e\cap B^{(3)})$, with the projection $q:A_3^\circ\to K$ and the injective map $j:C\to K$ induced by the inclusion, and let $c(1)=1$ and $c(2)=3$. Then:
\begin{enumerate}
\item the matrices $\sigma(g)=\big(q(\rho_3(g)_{c(a)c(b)})\big)_{1\le a,b\le2}$, for $g\in\Gamma_0$, and the characters $\sigma_{v,a}=q\circ\psi_{3,v,c(a)}$ of $G_v$ form a homomorphism of monoids $\sigma:\Gamma_0\to M_2(K)$ with characters satisfying (a), (b) and (c) of Definition~\ref{sha-descent-link}(2); and $q(\rho_3(\mathrm{loc}_vh)_{22})=q(\psi_{3,v,2}(h))$ for every $v$ and $h\in G_v$;
\item there are $\mathcal O_1$-algebra maps $\xi:A_2^\circ\to K$ and $\beta:B^{(2)}\to C$ such that $\xi$ carries $(\rho_2,\psi_2)$ to $(\sigma,(\sigma_{v,a}))$ and $\xi(b)=j(\beta(b))$ for $b\in B^{(2)}$.
\end{enumerate}
\end{lemma}
\leanref{Development/AutomorphicArithmetic/AutomorphicVarieties/FrobeniusSmoothVectors/ExistsFrobeniusCorner.lean}{UpperTheory.V.Fr.exists_frobenius_corner}
\leanref{Development/AutomorphicArithmetic/AutomorphicVarieties/FrobeniusSmoothVectors/AlgVring.lean}{UpperTheory.V.Fr.alg}
\begin{proof}
\openstep{The corner of the universal linked point of rank three modulo its entries of odd weight is not yet shown to be a linked point of rank two with a cone over $C$, nor is the map from $A_2^\circ$ that it gives constructed.}
% The first half of the private finite_lambda_two_three_surjection_compat (QuotientHomNonvanishingBound.lean): cone_u and frobenius_corner_cone_map (the image of the universal point in K has a Cone over C), hodd, frobenius_corner_units and exists_frobenius_corner (IsCorner: odd, r, psi, mid) for (1); frobenius_corner_mem_cone and alg (AlgVring.lean, with K a ring and C -> K injective) for (2). Size: 1,266 lines by tools/lean_size.py for the six declarations together (alg alone 767).
\end{proof}

\begin{lemma}\label{sha-automorphy-realization-components-propagation-endoscopic-lambda}
In Situation~\ref{sha-automorphy-link-situation}, let $\mathcal O_1\supseteq\mathcal O$ be the ring of integers of a finite extension of $\mathbf Q_\ell$ such that the link, with its residual characters composed with the map of residue fields, satisfies the residual conditions with some constant, and for $n\in\{2,3\}$ let $A_n$, its universal linked point $(\rho_n,\psi_n)$, $B^{(n)}=B$ and $A_n^\circ$ be as in Definition~\ref{sha-automorphy-realization}(2), with $\Lambda_{B^{(n)}}$ as in Definition~\ref{sha-automorphy-realization-components-transfer-primes}; let $c(1)=1$ and $c(2)=3$. Let $C$ and $K$ be $\mathcal O_1$-algebras, $p:B^{(3)}\to C$ a surjective and $q:A_3^\circ\to K$ an $\mathcal O_1$-algebra map, and $j:C\to K$ an injective $\mathcal O_1$-algebra map with $q(b)=j(p(b))$ for $b\in B^{(3)}$. Suppose that $q$ kills the entries of $\rho_3$ of odd weight and carries $\rho_3(\mathrm{loc}_vh)_{22}$ to $q(\psi_{3,v,2}(h))$ for every $v$ and $h\in G_v$, and let $\xi:A_2^\circ\to K$ and $\beta:B^{(2)}\to C$ be $\mathcal O_1$-algebra maps with $\xi(\rho_2(g)_{aa'})=q(\rho_3(g)_{c(a)c(a')})$, $\xi(\psi_{2,v,a}(h))=q(\psi_{3,v,c(a)}(h))$ and $\xi(b)=j(\beta(b))$ for all $g\in\Gamma_0$, $a,a'\in\{1,2\}$, places $v$, $h\in G_v$ and $b\in B^{(2)}$. Let $I_2$ be an ideal of $B^{(2)}$ such that the maximal ideal of $B^{(2)}$ lies in the radical of $I_2+\Lambda_{B^{(2)}}$. Then the maximal ideal of $B^{(3)}$ lies in the radical of $p^{-1}(\beta(I_2)C)+\Lambda_{B^{(3)}}$.
\end{lemma}
\leanref{Development/AutomorphicArithmetic/AutomorphicVarieties/FrobeniusSmoothVectors/QuotientHomNonvanishingBound.lean}{finite_lambda_surjective_injective_lift}
\begin{proof}
\openstep{The comparison of the radicals, through the middle entry and the monomials of weight zero in the corner, is not yet written out.}
% The private finite_lambda_surjective_injective_lift (QuotientHomNonvanishingBound.lean): for a prime P containing the ideal, hQ (the maximal ideal of B^(2) lies in the inverse image of p(P) under beta), hlift, frobenius_corner_middle_mem (the middle entries are 1 modulo P), frobenius_corner_monomial for the monomials of weight zero, and frobenius_corner_maximalIdeal_le. Size: 689 lines by tools/lean_size.py beyond the declarations of Lemma sha-automorphy-realization-components-propagation-endoscopic-corner (789 alone).
\end{proof}

\begin{lemma}\label{sha-automorphy-realization-components-propagation-endoscopic}
In Situation~\ref{sha-automorphy-link-situation}, let $\mathcal O_1\supseteq\mathcal O$ be the ring of integers of a finite extension of $\mathbf Q_\ell$ such that the link, with its residual characters composed with the map of residue fields, satisfies the residual conditions with some constant, and for $n\in\{2,3\}$ let $A_n$, $B^{(n)}=B$ and $A_n^\circ$ be as in Definition~\ref{sha-automorphy-realization}(2), with $\Lambda_{B^{(n)}}$ as in Definition~\ref{sha-automorphy-realization-components-transfer-primes}. Let $I_2$ be an ideal of $B^{(2)}$ such that the maximal ideal of $B^{(2)}$ lies in the radical of $I_2+\Lambda_{B^{(2)}}$ and $I_2\subseteq\xi_{B^{(2)}}$ for every off-triangular prime $\xi$ of $A_2^\circ$. Then there is an ideal $I_3$ of $B^{(3)}$ such that the maximal ideal of $B^{(3)}$ lies in the radical of $I_3+\Lambda_{B^{(3)}}$ and $I_3\subseteq\xi_{B^{(3)}}$ for every prime $\xi$ of $A_3^\circ$ outside which lie some entry of the universal linked point below the diagonal and some entry above it, while all its entries of odd weight lie in $\xi$.
\end{lemma}
\leanref{Development/AutomorphicArithmetic/AutomorphicVarieties/FrobeniusSmoothVectors/QuotientHomNonvanishingBound.lean}{UpperTheory.V.Fr.l53}
\begin{proof}
Write $(\rho_n,\psi_n)$ for the universal linked point of rank $n$, and let $\mathfrak e$ be the ideal of $A_3^\circ$ generated by the entries of $\rho_3$ of odd weight.
\begin{enumerate}
\item[(i)] \emph{If $\mathfrak e=A_3^\circ$.} No prime of $A_3^\circ$ contains all the entries of odd weight, and $I_3=B^{(3)}$ will do.
\item[(ii)] \emph{The ideal.} Otherwise let $K=A_3^\circ/\mathfrak e$ and $C=B^{(3)}/(\mathfrak e\cap B^{(3)})$, with the projections $q:A_3^\circ\to K$ and $p:B^{(3)}\to C$ and the injective map $j:C\to K$ induced by the inclusion, so that $q(b)=j(p(b))$ for $b\in B^{(3)}$. Lemma~\ref{sha-automorphy-realization-components-propagation-endoscopic-corner} gives $\xi:A_2^\circ\to K$ and $\beta:B^{(2)}\to C$. Let $I_3$ be the inverse image under $p$ of the ideal $\beta(I_2)C$. By Lemma~\ref{sha-automorphy-realization-components-propagation-endoscopic-lambda}, whose hypotheses on $p$, $q$, $j$, $\xi$ and $\beta$ are given by Lemma~\ref{sha-automorphy-realization-components-propagation-endoscopic-corner}, the maximal ideal of $B^{(3)}$ lies in the radical of $I_3+\Lambda_{B^{(3)}}$.
\item[(iii)] \emph{The endoscopic primes.} Let $\mathfrak P$ be a prime of $A_3^\circ$ outside which lie some entry of $\rho_3$ below the diagonal and some entry above it, while all its entries of odd weight lie in $\mathfrak P$. Then $\mathfrak e\subseteq\mathfrak P$, so $\mathfrak P=q^{-1}(\bar{\mathfrak P})$ for the prime $\bar{\mathfrak P}=q(\mathfrak P)$ of $K$. Let $\mathfrak P_2=\xi^{-1}(\bar{\mathfrak P})$, a prime of $A_2^\circ$. The entries of $\rho_3$ below the diagonal at $(2,1)$ and $(3,2)$ have odd weight, so some $\rho_3(g)_{31}$ lies outside $\mathfrak P$; then $\xi(\rho_2(g)_{21})=q(\rho_3(g)_{31})\notin\bar{\mathfrak P}$, so $\rho_2(g)_{21}\notin\mathfrak P_2$. In the same way some $\rho_2(g')_{12}$ lies outside $\mathfrak P_2$, so $\mathfrak P_2$ is off-triangular. By hypothesis $I_2\subseteq\mathfrak P_2$, so $j(\beta(b))=\xi(b)\in\bar{\mathfrak P}$ for $b\in I_2$. For $x\in I_3$, $p(x)\in\beta(I_2)C$, so $q(x)=j(p(x))\in\bar{\mathfrak P}$ and $x\in\mathfrak P$. So $I_3\subseteq\mathfrak P\cap B^{(3)}$.
\end{enumerate}
\end{proof}
% l53 (QuotientHomNonvanishingBound.lean): the case ee = top for (i); the private finite_lambda_two_three_surjection_compat for (ii) and (iii), with the private finite_lambda_surjective_injective_lift for the radical and the case analysis hTI (with hxir, the entries of xi through the corner) for (iii).

\begin{lemma}\label{sha-automorphy-realization-components-propagation-diagonal}
In Situation~\ref{sha-automorphy-link-situation}, let $\mathcal O_1\supseteq\mathcal O$ be the ring of integers of a finite extension of $\mathbf Q_\ell$, with residue field $k_1$, such that the link, with its residual characters composed with the map of residue fields, satisfies the residual conditions with some constant.
\begin{enumerate}
\item Let $n\in\{2,3\}$ and $A_n$, its universal linked point $(\rho^{\mathrm u},\psi^{\mathrm u})$ and $A_n^\circ$ be as in Definition~\ref{sha-automorphy-realization}(2). For every $v$ and $h\in G_v$, the characteristic polynomial of $\rho^{\mathrm u}(\mathrm{loc}_vh)$ is $\prod_i\big(X-\psi^{\mathrm u}_{v,i}(h)\big)$.
\item Let $A$ be a noetherian local domain with finite residue field, $j:k_1\to A$ a ring map, and $\rho:\Gamma_0\to M_3(A)$ a homomorphism of monoids with characters $\psi_{v,i}$ satisfying (a), (b) and (c) of Definition~\ref{sha-descent-link}(2), continuous for the $\mathfrak m_A$-adic topology, whose diagonal entries are multiplicative in $g$ and congruent to $j(\bar\Psi_\Gamma(g)^{w_{3,i}})$ modulo $\mathfrak m_A$, with $\rho(\mathrm{loc}_vh)_{ii}=1$ for $h\in I_v$ with $\bar\Psi_\Gamma(\mathrm{loc}_vh)=1$. Then $\rho(g)_{ii}=j(\bar\Psi_\Gamma(g)^{w_{3,i}})$ for every $g$ and $i$.
\end{enumerate}
\end{lemma}
\leanref{Development/AutomorphicArithmetic/AutomorphicVarieties/FrobeniusSmoothVectors/CharpolyFlag.lean}{UpperTheory.V.Fr.o1}
\leanref{Development/AutomorphicArithmetic/AutomorphicVarieties/CurveDVRSpecialization/FinPtOf.lean}{diag_eq}
\begin{proof}
(1) By residual condition (7) of Definition~\ref{sha-descent-residual} there is $g_v\in G_v$ such that $a=\bar\Psi_\Gamma(\mathrm{loc}_vg_v)$ has $a^2\ne1$. We apply Lemma~\ref{sha-automorphy-realization-components-transfer-arc-triangular} to the local ring $A_n$, $H=G_v$, $r(h)=\rho^{\mathrm u}(\mathrm{loc}_vh)$, $\chi_i=\psi^{\mathrm u}_{v,i}$ and $g_0=g_v$. By condition (e) of Definition~\ref{sha-descent-link}(2), $\psi^{\mathrm u}_{v,i}(g_v)$ has residue $a^{w_{n,i}}$, and these residues are pairwise distinct, as $w_{n,i}-w_{n,j}$ is $\pm1$ or $\pm2$ for $i\ne j$; so the differences $\psi^{\mathrm u}_{v,i}(g_v)-\psi^{\mathrm u}_{v,j}(g_v)$ are units of the local ring $A_n$. By (e) again, $\rho^{\mathrm u}(\mathrm{loc}_vg_v)$ reduces to the diagonal matrix with entries $a^{w_{n,i}}$, whose characteristic polynomial is $\prod_i(X-a^{w_{n,i}})$. The product condition is (c). So there is an invertible $P$ with $P\rho^{\mathrm u}(\mathrm{loc}_vh)P^{-1}$ upper triangular with diagonal entries $\psi^{\mathrm u}_{v,1}(h),\dots,\psi^{\mathrm u}_{v,n}(h)$ for every $h\in G_v$, and $\rho^{\mathrm u}(\mathrm{loc}_vh)$ has the characteristic polynomial of that matrix.

(2) Fix $i$, and write $w=w_{3,i}$ and $\mathfrak m=\mathfrak m_A$. As $j$ is injective, $A$ has characteristic $\ell$. Let $h\ge1$ be the integer of residual condition (10).
\begin{enumerate}
\item[(i)] \emph{A character.} Let $\alpha(g)=\rho(g)_{ii}$. It is multiplicative, and $\alpha(1)=1$, so $\alpha(g)\alpha(g^{-1})=1$ and $\alpha$ is a homomorphism $\Gamma_0\to A^\times$. Let $\beta(g)=j(\bar\Psi_\Gamma(g)^{w})$, a homomorphism whose kernel contains the open subgroup $\ker\bar\Psi_\Gamma\cap\Gamma_0$ (residual condition (4)). Let $\chi=\alpha\beta^{-1}$. Then $\chi(g)\equiv1$ modulo $\mathfrak m$ for every $g$.
\item[(ii)] \emph{$\chi$ kills the generators of inertia.} Let $v$ be a place of the link and $u\in I_v$ with $\bar\Psi_\Gamma(\mathrm{loc}_vu)=1$, and let $x=\mathrm{loc}_vu$. Then $\alpha(x)=1$ by hypothesis, and $\beta(x)=1$, so $\chi(x)=1$. By condition (a), the $(i,i)$ entry of $\rho(c_\Gamma y^{-1}c_\Gamma^{-1})$ is the $(4-i,4-i)$ entry of $\rho(y)$; with $y=x^{-1}$, $\alpha(c_\Gamma xc_\Gamma^{-1})=\rho(x^{-1})_{4-i,4-i}=\rho(x)_{4-i,4-i}^{-1}=1$. As $\bar\Psi_\Gamma$ takes values in an abelian group, $\beta(c_\Gamma xc_\Gamma^{-1})=\beta(x)=1$. So $\chi(c_\Gamma xc_\Gamma^{-1})=1$.
\item[(iii)] \emph{Finite quotients.} Let $N\ge1$. The ring $A/\mathfrak m^N$ is finite, as $A$ is noetherian with finite residue field. For $g\in\Gamma_0$ write $\chi(g)=1+t$ with $t\in\mathfrak m$; then $\chi(g)^{\ell^N}=1+t^{\ell^N}\equiv1$ modulo $\mathfrak m^N$, as $A$ has characteristic $\ell$. So the reduction $\chi_N$ of $\chi$ modulo $\mathfrak m^N$ takes values in the subgroup of the elements $y$ of $(A/\mathfrak m^N)^\times$ with $y^{\ell^N}=1$, a finite abelian $\ell$-group. Its kernel is open, as it contains the intersection of $\ker\beta$ with the set of $g$ at which $\alpha(g)\equiv\alpha(1)$ modulo $\mathfrak m^N$, an open neighbourhood of $1$ by the continuity of $\rho$.
\item[(iv)] \emph{$\chi^h=1$.} Let $v$ be a place of the link and $u_1,\dots,u_{d_v}\in I_v$ as in residual condition (9). The homomorphisms $\chi_N\circ\mathrm{loc}_v$ and $h'\mapsto\chi_N(c_\Gamma\mathrm{loc}_v(h')c_\Gamma^{-1})$ from $G_v$ to a finite abelian $\ell$-group are continuous, as $\mathrm{loc}_v$ and conjugation are, and kill $u_1,\dots,u_{d_v}$ by (ii); by (9) they kill $I_v$. So $\chi_N$ kills $\mathrm{loc}_v(I_v)$ and $c_\Gamma\mathrm{loc}_v(I_v)c_\Gamma^{-1}$ for every $v$, and $\chi_N^h=1$ by residual condition (10). That is, $\chi(g)^h-1\in\mathfrak m^N$ for every $N$, and $\chi(g)^h=1$ by Krull's intersection theorem \cite[Theorem~8.10]{Mat86}.
\item[(v)] \emph{$\chi=1$.} Let $g\in\Gamma_0$, write $h=\ell^km$ with $\ell\nmid m$, and let $y=\chi(g)^{\ell^k}$. Then $y-1=(\chi(g)-1)^{\ell^k}\in\mathfrak m$ and $y^m=1$. So $(y-1)(1+y+\dots+y^{m-1})=0$, and the second factor, congruent to $m$ modulo $\mathfrak m$, is a unit, since $m$ is prime to the characteristic $\ell$ of the residue field. Hence $(\chi(g)-1)^{\ell^k}=y-1=0$, and $\chi(g)=1$ as $A$ is a domain.
\end{enumerate}
So $\rho(g)_{ii}=\alpha(g)=\beta(g)=j(\bar\Psi_\Gamma(g)^{w_{3,i}})$.
\end{proof}
% o1 (CharpolyFlag.lean): hF.f3 for g_v, zpow_w_ne' for the units, hd (the diagonal reduction from S.lift) and Matrix.charpoly_diagonal for the residual characteristic polynomial, flag_loc (Lemma sha-automorphy-realization-components-transfer-arc-triangular) and charpoly_flag' for the conclusion; the Lean works over S.R and maps back to S.A by injectivity. diag_eq (FinPtOf.lean, private): a and b for (i); hT, with hinv1 and the star facts of Fr.star, for (ii); hpowN, Ideal.finite_quotient_pow and hopen for (iii); hF.gen and hF.cl through hcl', and Ideal.iInf_pow_eq_bot_of_isLocalRing for (iv); Pr.T.oneUnitTF for (v).

\begin{lemma}\label{sha-automorphy-realization-components-propagation}
In Situation~\ref{sha-automorphy-link-situation}, let $\mathcal O_1\supseteq\mathcal O$ be the ring of integers of a finite extension of $\mathbf Q_\ell$, with residue field $k_1$, such that the link, with its residual characters composed with the map of residue fields, satisfies the residual conditions with a constant $C_0$ and is large for $C_0$ (Definition~\ref{sha-descent-residual}). For $n\in\{2,3\}$ let $A_n$, $B$ and $A_n^\circ$ be as in Definition~\ref{sha-automorphy-realization}(2), use the notation of Definition~\ref{sha-automorphy-realization-components-transfer-primes}, and let $M_n$ be an ideal of the subalgebra $B$ of $A_n$ with the properties (b), (c) and (e) of Lemma~\ref{sha-automorphy-realization-components-transfer} in place of $\phi(J_n)$, and for $n=3$ also (d). Then:
\begin{enumerate}
\item[(a)] $M_2\subseteq\mathfrak p_B$ for every off-triangular prime $\mathfrak p$ of $A_2^\circ$;
\item[(b)] $M_3\subseteq\mathfrak p_B$ for every off-triangular prime $\mathfrak p$ of $A_3^\circ$, and there is an ideal $I$ of $B\subseteq A_3$ such that the maximal ideal of $B$ lies in the radical of $I+\Lambda_B$ and $I\subseteq\mathfrak p_B$ for every prime $\mathfrak p$ of $A_3^\circ$ outside which lie some entry of the universal linked point below the diagonal and some entry above it;
\item[(c)] let $A$ be a complete discrete valuation ring with finite residue field, with a local map $\mathcal O_1\to A$ through which $A$ is a $k_1$-algebra, and $(\rho,\psi)$ a linked point of rank three with values in $A$ lifting the residual one, such that $\psi_{v,i}(h)=1$ for $h\in I_v$ with $\bar\Psi_\Gamma(\mathrm{loc}_vh)=1$; then $\rho(g)$ has the characteristic polynomial $\prod_i(X-\bar\Psi_\Gamma(g)^{w_{3,i}})$ for every $g\in\Gamma_0$.
\end{enumerate}
\end{lemma}
\leanref{Development/AutomorphicArithmetic/AutomorphicVarieties/CurveDVRSpecialization/ComponentsGoodRankTwo.lean}{UpperTheory.V.Pr.components_good_rank_two}
\leanref{Development/AutomorphicArithmetic/AutomorphicVarieties/CurveDVRSpecialization/PropagationTheoremGL2.lean}{UpperTheory.V.Pr.PropagationTheoremGL2}
\leanref{Development/AutomorphicArithmetic/AutomorphicVarieties/CurveDVRSpecialization/PropagationTheoremGL3.lean}{UpperTheory.V.Pr.components_good_rank_three}
\leanref{Development/AutomorphicArithmetic/AutomorphicVarieties/CurveDVRSpecialization/PropagationTheoremGL2.lean}{UpperTheory.V.Pr.PropagationTheoremGL3}
\leanref{Development/AutomorphicArithmetic/AutomorphicVarieties/CurveDVRSpecialization/FinPtOf.lean}{UpperTheory.V.Pr.finPtOf}
\begin{proof}
For $n\in\{2,3\}$, the ring $A_n^\circ$ is noetherian by Lemma~\ref{sha-automorphy-realization-presentation}(4), so $X_n=\operatorname{Spec}A_n^\circ$ is a noetherian space. Let $T_n\subseteq X_n$ be the set of off-triangular primes, which is open, being cut out by the conditions that the ideals generated by the entries below the diagonal, above it, and, if $n=3$, of odd weight are not contained in the prime; and let $G_n\subseteq X_n$ be the set of primes $\mathfrak p$ with $M_n\subseteq\mathfrak p_B$, the zero locus of the ideal generated by $M_n$, which is closed. Let $X_{\mathrm c}\subseteq X_n$ be as in Lemma~\ref{sha-automorphy-realization-components-propagation-geometry}, and $X_{\mathrm{cl}}$ the set of its maximal ideals. In both (a) and (b) we apply Lemma~\ref{sha-automorphy-realization-components-propagation-global} to $X_n$, $T_n$, $G_n$, $C=X_{\mathrm c}$, $C_{\mathrm{cl}}=X_{\mathrm{cl}}$ and a set $E$; its conclusion $T_n\subseteq G_n$ is the first statement of (a) and of (b).

(a) Let $n=2$ and $E=\emptyset$.
\begin{enumerate}
\item[(i)] \emph{Hypotheses (1) to (4) of Lemma~\ref{sha-automorphy-realization-components-propagation-global}.} (1) is Lemma~\ref{sha-automorphy-realization-components-propagation-geometry}(1), and (2) is its part (2) with $U=X_2$. For (3), let $\mathfrak p\in X_{\mathrm c}$ and $U\ni\mathfrak p$ open; part (2), with the closed set $X_2$, gives a maximal ideal in $U\cap X_{\mathrm c}$, so $\mathfrak p$ lies in the closure of $X_{\mathrm{cl}}$. For (4), a $T_2$-component is the zero locus of a minimal prime $\eta$ of $A_2^\circ$; as it contains an off-triangular prime, some entry below the diagonal lies outside $\eta$, and part (3) gives a maximal ideal in $X_{\mathrm c}$ containing $\eta$.
\item[(ii)] \emph{The local statement.} Let $x\in X_{\mathrm{cl}}$. By Lemma~\ref{sha-automorphy-realization-components-propagation-completions}(1) there is a completion $(\mathcal O',R,\iota)$ at $x$, and $\iota$ is flat by part (2). Apply Lemma~\ref{sha-automorphy-realization-components-propagation-local} to $R$, $D_1=2d+1$, $I=0$, the off-triangular primes $\mathcal T$, $\mathcal E=\emptyset$, the set $\mathcal G$ of primes $P$ with $M_2\subseteq P_B$, the nice primes $\mathcal N$ and $\mathcal H=\emptyset$. Its hypotheses hold: (1) is clear; nice primes are off-triangular, as $n=2$, which is (2); (3) is Lemma~\ref{sha-automorphy-realization-components-propagation-completions}(5) with $I'=M_2$, by property (b) of $M_2$; (4) and (7) are empty; (5) is Lemma~\ref{sha-automorphy-realization-components-propagation-completions}(3), the components being the minimal primes of $R$; and (6) is property (c) of $M_2$. By Lemma~\ref{sha-automorphy-realization-components-propagation-completions}(2), $\mathcal T$ and $\mathcal G$ are the inverse images of $T_2$ and $G_2$ under the map $f:\operatorname{Spec}R\to X_2$ induced by $\iota$. By the first statement of Lemma~\ref{sha-automorphy-realization-components-propagation-charts}, applied to $C=A_2^\circ$, $j$ the identity of $X_2$ and the flat map $\iota$, $f$ has the properties that its second statement requires at $x=\iota^{-1}(\mathfrak m_R)$; and condition $(\ast)$ there holds, as just shown.
\item[(iii)] \emph{Hypotheses (5) and (6).} By Lemma~\ref{sha-automorphy-realization-components-propagation-completions}(3), every minimal prime $P$ of $R$ is off-triangular with $\dim R/P\ge2d+2$. So Lemma~\ref{sha-automorphy-realization-components-propagation-charts}(1) and (2) give hypotheses (5) and (6) at $x$.
\item[(iv)] \emph{Hypothesis (7).} By property (e) of $M_2$ there are an off-triangular prime $\xi$ of $A_2^\circ$ and a prime $Z$ of $B$ with $M_2\subseteq Z\subseteq\xi_B$ and $\dim B/Z=1+2d$. Lemma~\ref{sha-automorphy-realization-components-propagation-geometry}(4) gives $x\in X_{\mathrm{cl}}$, and, for a completion at $x$, an off-triangular prime $P$ with $P_B=Z$, so $M_2\subseteq P_B$, and $\dim R/P=2+2d\ge2d+1$. By (ii) and Lemma~\ref{sha-automorphy-realization-components-propagation-charts}(3), some $T_2$-component lies in $G_2$.
\end{enumerate}
% (a) is components_good_rank_two (ComponentsGoodRankTwo.lean): gr_A; the fields of GlobData.Ok for gd S M (TI, Good = M <= pi p, Core = S.BB, Cl, EE = empty): preconnected_rank_two, jac twice, comp_min with min_hom and l41 for (i); at_pt with loc_ex, ld_ok (sharp2, exists_nice_prime_above_of_large_dim, P1'), locProp, chartOfFlat with loc_flat, loc_rho and globOfCharts for (ii) and (iii); seedpt and at_pt for (iv); then globProp.

(b) Let $n=3$, and let $I_3$ be the ideal of $B\subseteq A_3$ that Lemma~\ref{sha-automorphy-realization-components-propagation-endoscopic} gives for $I_2=M_2$, which applies by property (b) of $M_2$ and by (a). Let $E\subseteq X_3$ be the set of primes that contain all the entries of $\rho^{\mathrm u}$ of odd weight.
\begin{enumerate}
\item[(i)] \emph{Hypotheses (1) to (4).} (1), (2) and (4) are as in (a). For (3): by Lemma~\ref{sha-automorphy-realization-components-propagation-geometry}(1), $X_{\mathrm c}$ lies in the closure of $X_{\mathrm c}\setminus E$; and for $\mathfrak p\in X_{\mathrm c}\setminus E$ and $U\ni\mathfrak p$ open, part (2), with the closed set $X_3$ and the open set $U\setminus E$, gives a maximal ideal in $(U\cap X_{\mathrm c})\setminus E$. So $X_{\mathrm c}$ lies in the closure of $X_{\mathrm{cl}}\setminus E$.
\item[(ii)] \emph{The local statement.} Let $x\in X_{\mathrm{cl}}$, $(\mathcal O',R,\iota)$ a completion at $x$, and $I$ an ideal of $R$ as in Lemma~\ref{sha-automorphy-realization-components-propagation-completions}(4). Apply Lemma~\ref{sha-automorphy-realization-components-propagation-local} to $R$, $D_1=3d+1$, $I$, the off-triangular primes $\mathcal T$, the endoscopic primes $\mathcal E$, the set $\mathcal G$ of primes $P$ with $M_3\subseteq P_B$, the nice primes $\mathcal N$ and the thick primes $\mathcal H$. Its hypotheses hold. (1) is clear. (2): a nice prime is off-triangular if some entry of odd weight lies outside it, and endoscopic otherwise. (3) is Lemma~\ref{sha-automorphy-realization-components-propagation-completions}(5) with $I'=M_3$. (4): let $P$ contain a thick component $Q$, with $\dim R/P\ge3d+1$. As $Q$ is endoscopic, so is $\iota^{-1}(Q)$, by Lemma~\ref{sha-automorphy-realization-components-propagation-completions}(2), so $I_3\subseteq\iota^{-1}(Q)_B\subseteq P_B$; part (5) of that lemma, with $I'=I_3$, gives a nice prime containing $P$, which contains the entries of odd weight, as $Q$ does, and so is endoscopic. (5): a minimal prime of $R$ in $\mathcal T$ contains $I$, by Lemma~\ref{sha-automorphy-realization-components-propagation-completions}(4), so it is minimal over $I$; the chains and the dichotomy are part (4) of that lemma. (6) and (7) are properties (c) and (d) of $M_3$, as two linked primes are thickly linked. As in (a), $(\ast)$ of Lemma~\ref{sha-automorphy-realization-components-propagation-charts} holds at $x$.
\item[(iii)] \emph{Hypothesis (5).} Let $Y,Y'$ be $T_3$-components containing $x$, with $Y\subseteq G_3$. By Lemma~\ref{sha-automorphy-realization-components-propagation-charts}, $Y$ and $Y'$ are the closures of the images of minimal primes $P,P'$ of $R$, which lie in $\mathcal T$, as $T_3$ is open. Then $P$ is minimal over $I$, so $\dim R/P\ge3d+2$ by Lemma~\ref{sha-automorphy-realization-components-propagation-completions}(4), and $P\in\mathcal G$. By $(\ast)$, $P'\in\mathcal G$, and $Y'\subseteq G_3$, as $G_3$ is closed.
\item[(iv)] \emph{Hypothesis (6).} Let $x\in X_{\mathrm{cl}}\setminus E$. Then $I=0$, and every minimal prime $P$ of $R$ is off-triangular or thick. A thick prime is endoscopic, so if $P$ were thick, $\iota^{-1}(P)\subseteq x$ would contain all the entries of odd weight, by Lemma~\ref{sha-automorphy-realization-components-propagation-completions}(2), and $x\in E$. So every minimal prime of $R$ lies in $\mathcal T$, and Lemma~\ref{sha-automorphy-realization-components-propagation-charts}(2) gives (6).
\item[(v)] \emph{Hypothesis (7).} As in (a)(iv), with $\dim R/P=2+3d\ge3d+1$.
\item[(vi)] \emph{The ideal $I$.} Let $I=M_3\cap I_3$. A prime of $B$ that contains $I+\Lambda_B$ contains $M_3$ or $I_3$, and so contains $\mathfrak m_B$, by property (b) of $M_3$ and by Lemma~\ref{sha-automorphy-realization-components-propagation-endoscopic}; so $\mathfrak m_B$ lies in the radical of $I+\Lambda_B$. Let $\mathfrak p$ be a prime of $A_3^\circ$ outside which lie some entry below the diagonal and some entry above it. If $\mathfrak p$ contains all the entries of odd weight, then $I\subseteq I_3\subseteq\mathfrak p_B$; otherwise $\mathfrak p$ is off-triangular, and $I\subseteq M_3\subseteq\mathfrak p_B$ by the first statement.
\end{enumerate}
% (b) is components_good_rank_three (PropagationTheoremGL3.lean): l53 for I_3; gd S M with EE = {ee <= p}; preconnected_rank_three and jac for (i); loc_ex, sharp3, ld_ok and locProp for (ii); chart_ex (chartOfFlat with loc_flat), ld_ti (loc_rho) and comp_of_ti for (iii); Sharp.triv, Sharp.dich and globOfCharts for (iv); seedpt and globOfCharts for (v); FinLam of M inf I_3 and the case split on ee <= p for (vi).

(c) Let $A$ and $(\rho,\psi)$ be as in (c), and write $\rho_{ij}$ for the entries of $\rho(g)$. Lemma~\ref{sha-automorphy-realization-components-transfer-local-point}, applied to $C=A$, whose residue field is finite, gives $\xi:A_3^\circ\to A$ that carries $(\rho^{\mathrm u},\psi^{\mathrm u})$ to $(\rho,\psi)$. It suffices to show, for every $g$, that $\rho_{12}\rho_{21}$, $\rho_{13}\rho_{31}$, $\rho_{23}\rho_{32}$, $\rho_{12}\rho_{23}\rho_{31}$ and $\rho_{13}\rho_{32}\rho_{21}$ vanish and that $\rho_{ii}$ is the image of $\bar\Psi_\Gamma(g)^{w_{3,i}}$ in $A$: then, expanding the determinant, the characteristic polynomial of $\rho(g)$ is $\prod_i(X-\rho_{ii})$.
\begin{enumerate}
\item[(i)] \emph{If the kernel of $\xi$ is off-triangular.} Let $I$ be as in (b), so $I$ lies in the prime $\ker(\xi)\cap B$. So does $\Lambda_B$: $\xi$ kills $\mathfrak m_{\mathcal O_1}$, as $\mathcal O_1\to A$ factors through $k_1$, and $\xi(\psi^{\mathrm u}_{v,i}(h)-1)=\psi_{v,i}(h)-1=0$ for $h\in I_v$ with $\bar\Psi_\Gamma(\mathrm{loc}_vh)=1$. By (b), $\mathfrak m_B\subseteq\ker\xi$, so $\xi(b)$ is the image in $A$ of the residue of $b$ in $k_1$, for $b\in B$. The five products above and the $\rho^{\mathrm u}(g)_{ii}$ are products of entries of $\rho^{\mathrm u}$ of total weight $0$, so they lie in $B$; the residual point is diagonal, so the residues of the five products are $0$, and that of $\rho^{\mathrm u}(g)_{ii}$ is $\bar\Psi_\Gamma(g)^{w_{3,i}}$.
\item[(ii)] \emph{If not.} All entries of $\rho^{\mathrm u}$ below the diagonal, or all above it, lie in $\ker\xi$, so every $\rho(g)$ is triangular of the same kind. Then the five products vanish, and the $\rho_{ii}$ are multiplicative in $g$. For $h\in I_v$ with $\bar\Psi_\Gamma(\mathrm{loc}_vh)=1$, Lemma~\ref{sha-automorphy-realization-components-propagation-diagonal}(1), carried by $\xi$, shows that $\rho(\mathrm{loc}_vh)$ has the characteristic polynomial $\prod_i(X-\psi_{v,i}(h))=(X-1)^3$; as $\rho(\mathrm{loc}_vh)$ is triangular, its diagonal entries are roots of $(X-1)^3$ in the domain $A$, hence equal to $1$. As $(\rho,\psi)$ is continuous and lifts the residual point, Lemma~\ref{sha-automorphy-realization-components-propagation-diagonal}(2) gives $\rho_{ii}=\bar\Psi_\Gamma(g)^{w_{3,i}}$.
\end{enumerate}
% (c) is finPtOf (FinPtOf.lean): point_loc; charpoly_of_cyc for the reduction to the five products and the diagonal; the case OffTri (ker xi) with OffFin (the ideal I of (b)), hlam, hmB, hval and hprod for (i); htri, hcyc, hmul, hbar, h1 (Fr.o1) and diag_eq for (ii).
% The hypotheses F.Ok C_0 and F.Big C_0 of PropagationTheoremGL2 and GL3 are the residual conditions and the largeness; FinPtOf needs only F.Ok C_0 and OffFin. Open steps remain in Lemmas sha-automorphy-realization-components-propagation-geometry ((1)), -geometry-lift, -geometry-dimension, -completions, -endoscopic-corner and -endoscopic-lambda.
\end{proof}

\begin{lemma}\label{sha-automorphy-realization-components}
In Situation~\ref{sha-automorphy-link-situation}, let realizations of ranks two and three with coefficients $\mathcal O_1$ be given, with ideals $J_2$ and $J_3$ as in Lemma~\ref{sha-automorphy-realization-ideal}. Then:
\begin{enumerate}
\setcounter{enumi}{-1}
\item there is a compatible map $\phi$ from $B_3$ onto the subalgebra $B$ of $A_3$, the closure of the $\mathcal O_1$-subalgebra generated by the values of the characters $\psi^{\mathrm u}_{v,i}$ and by the products of entries of the universal linked point of total weight $0$, such that every off-triangular prime of $A_3^\circ$ contains $\phi(J_3)$ (Definition~\ref{sha-automorphy-realization});
\item for $n\in\{2,3\}$, let $A$ be a complete discrete valuation ring with finite residue field and a local map $\mathcal O_1\to A$, and $(\rho,\psi)$ a linked point of rank $n$ with values in $A$ lifting the residual one, such that for some field $K'$ and injective ring map $j:A\to K'$, $0$ and $K'^n$ are the only subspaces of $K'^n$ stable under every $j(\rho(g))$; then there is an $\mathcal O_1$-algebra map $B_n\to A$ that kills $J_n$ and carries the universal determinant and characters to the determinant of $g\mapsto\kappa_0(g)\rho(\pi(g))$ and the characters $h\mapsto\psi_{v,i}(h)\kappa_0(h)$, where $\pi$ and $\kappa_0$ are the map and the character of the realization of rank $n$;
\item let $A$ be a complete discrete valuation ring with finite residue field and a local map $\mathcal O_1\to A$ that factors through the residue field $k_1$ of $\mathcal O_1$, and $(\rho,\psi)$ a linked point of rank three with values in $A$ lifting the residual one, such that $\psi_{v,i}(h)=1$ for every $v$, every $i$ and every $h\in I_v$ with $\bar\Psi_\Gamma(\mathrm{loc}_vh)=1$; then for every $g\in\Gamma_0$, $\rho(g)$ has the characteristic polynomial $\prod_i(X-\bar\Psi_\Gamma(g)^{w_{3,i}})$, its coefficients taken in $A$ through $k_1\to A$.
\end{enumerate}
\end{lemma}
\leanref{Development/AutomorphicArithmetic/AutomorphicVarieties/ModularityArgument/LinkedUnramifiedChains/DerivativeDivisibility/AbsGood.lean}{UpperTheory.V.J.absGood}
\leanref{Development/AutomorphicArithmetic/AutomorphicVarieties/CurveDVRSpecialization/PropagationTheoremGL2.lean}{UpperTheory.V.Pr.AllGood}
\leanref{Development/AutomorphicArithmetic/AutomorphicVarieties/CurveDVRSpecialization/AbsMain.lean}{UpperTheory.V.Pr.absMain}
\leanref{Development/AutomorphicArithmetic/AutomorphicVarieties/CurveDVRSpecialization/SecondThirdPrimeConditionTransferAlt.lean}{UpperTheory.V.Pr.Setup.GoodPt}
\leanref{Development/AutomorphicArithmetic/AutomorphicVarieties/CurveDVRSpecialization/PropagationTheoremGL2.lean}{UpperTheory.V.Pr.FinPt}
\begin{proof}
For $n\in\{2,3\}$, Lemma~\ref{sha-automorphy-realization-components-transfer}(a) gives a compatible surjection $\phi_n$ from $B_n$ onto the subalgebra $B$ of $A_n$, and by (b), (c) and (e) of that lemma, and (d) for $n=3$, the ideals $M_n=\phi_n(J_n)$ satisfy the hypotheses of Lemma~\ref{sha-automorphy-realization-components-propagation}. By (a) of that lemma, $\phi_2(J_2)\subseteq\mathfrak p_B$ for every off-triangular prime $\mathfrak p$ of $A_2^\circ$, and by its (b), $\phi_3(J_3)\subseteq\mathfrak p_B$ for every off-triangular prime $\mathfrak p$ of $A_3^\circ$. This is (0), with $\phi=\phi_3$. By Lemma~\ref{sha-automorphy-realization-components-transfer}(f), for $n=2$ and $n=3$, these inclusions give (1). Lemma~\ref{sha-automorphy-realization-components-propagation}(c) is (2).
% (0) is J.absGood: phiEx, then components_good_rank_two and components_good_rank_three as in absMain, keeping the conclusion AllGood (the off-triangular primes are those with S.u.TI). absMain: phiEx, finLamOf, p1Of, rankThree_deformationSetup_splitting_condition, seedOf, components_good_rank_two (PropagationTheoremGL2: preconnected_rank_two, sharp2, exists_nice_prime_above_of_large_dim, P1', SeedIn, globProp) and components_good_rank_three (PropagationTheoremGL3: P23', l53, OffFin), goodPtOf and finPtOf. Size: CurveDVRSpecialization, about 14,000 lines.
\end{proof}

\begin{lemma}\label{sha-automorphy-realization-away}
In Situation~\ref{sha-automorphy-link-situation}, let a realization of rank $n$, twisted by an admissible $\kappa$, with coefficients $\mathcal O_1$, be given with an ideal $J_n$ that satisfies Lemma~\ref{sha-automorphy-realization-components}(1) and the classicality of Lemma~\ref{sha-automorphy-realization-ideal}. Let $A\subset\bar{\mathbf Q}_\ell$ be a discrete valuation ring finite over $\mathcal O$ and $(\rho,\psi)$ a linked point of rank $n$ with values in $A$, lifting the residual one, with the prescribed inertial characters, such that $r=\rho\circ p$ is irreducible on $\operatorname{Gal}(\bar{\mathbf Q}/L)$. Then $r\otimes\kappa$ is attached through $\iota$ to a regular algebraic cuspidal automorphic representation $\pi$ of $GL_n(\mathbf A_L)$ with $\pi^\vee\cong\pi^c$, and every such $\pi$ satisfies the windowed condition at the places not above the place $v_0$ of the realization.
\end{lemma}
\leanref{Development/AutomorphicArithmetic/AutomorphicVarieties/ModularityArgument/ThresholdArguments/DenseThresholdCases/TwoRuns.lean}{UpperTheory.Gal.projectedFormsSystemPart.autPtO_alt}
\leanref{Development/AutomorphicArithmetic/Foundations/GaloisRepresentations/PlainTorsionPairs/IrrArithPointsAutomorphicAway.lean}{UpperTheory.Gal.projectedFormsSystemPart.AutOffN}
\begin{proof}
\begin{enumerate}
\item \emph{Coefficients.} The finite extension of $\mathbf Q_\ell$ generated by the fraction fields of $\mathcal O_1$ and $A$ has a ring of integers $\mathcal O_2$ that contains $\mathcal O_1$ and $A$, since their elements are integral over $\mathbf Z_\ell$. The inclusions $A\to\mathcal O_2$ and $\mathcal O_1\to\mathcal O_2$ are injective and integral, hence local, and $\mathcal O_2$ is a complete discrete valuation ring finite over $\mathcal O_1$.
\item \emph{The point of $B_n$.} Through $A\subset\mathcal O_2$, $(\rho,\psi)$ is a linked point with values in $\mathcal O_2$: its entries stay continuous, and it lifts the residual one, with the map of residue fields of the realization. With $j$ the inclusion of $\mathcal O_2$ in $\bar{\mathbf Q}_\ell$, $0$ and $\bar{\mathbf Q}_\ell^n$ are the only subspaces of $\bar{\mathbf Q}_\ell^n$ stable under $\rho(\Gamma_0)$, because such a subspace is stable under $r$, as $r(s)=\rho(p(s))$. The ring $\mathcal O_2$ is a complete discrete valuation ring with finite residue field, and $\mathcal O_1\to\mathcal O_2$ is local, by (1). So Lemma~\ref{sha-automorphy-realization-components}(1) gives $F:B_n\to\mathcal O_2$ that kills $J_n$ and carries $(D^{\mathrm u},\psi^{\mathrm u})$ to the determinant of $\rho$ and the characters $\psi_{v,i}$, composed with the map $\pi$ of the realization and twisted by its character. Let $x:B_n\to\bar{\mathbf Q}_\ell$ be $F$ followed by the inclusion.
\item \emph{The determinant.} Restricted along the map from $\operatorname{Gal}(\bar{\mathbf Q}/L)$ to the Galois group of the residual datum, $\pi$ becomes $p$ and the twisting character becomes $\kappa$, so $x(D^{\mathrm u})$ is the determinant of $r\otimes\kappa$.
\item \emph{The weight.} Let $v\mid\ell$. Since $\kappa$ is the $\ell$-adic character of an algebraic Hecke character, there are integers $a_{v,\tau}$, for the embeddings $\tau$ of $L_{\tilde v}$, such that every character of $D_v$ equal to $\chi_{\mathrm{cyc}}^m\chi\kappa$ on $I_v$, for an integer $m$ and a character $\chi$ with open kernel, satisfies the condition of Definition~\ref{sha-hecke-arcs}(1) with the exponents $a_{v,\tau}-m$; this follows from global class field theory \cite{Ser68}. By (2), $\psi_{x,v,i}=\Psi^{w_{n,i}}\kappa$ on $I_v$, and $\Psi^{w_{n,i}}=\chi_{\mathrm{cyc}}^{kw_{n,i}}\psi^{w_{n,i}}$ with $\psi$ of finite order. So $x$ is arithmetic of the weight $h_{v,i,\tau}=a_{v,\tau}-kw_{n,i}$, which is increasing because $k\ge1$ and $w_{n,1}>\dots>w_{n,n}$.
\item \emph{Classicality.} The point $x$ lies on the zero locus of $J_n$ and is arithmetic, and by (3) its determinant is that of $r\otimes\kappa$, which is irreducible since $r$ is. The classicality of Lemma~\ref{sha-automorphy-realization-ideal}, which is Proposition~\ref{sha-hecke-ideal}(5), gives the conclusion.
\end{enumerate}
\end{proof}
% autPtO_alt (TwoRuns.lean) with the private helpers of AutPtOf.lean: (1) ap_coef (ap_field, ap_norm, X.So.coefOf) and ap_local; (2) apExt, ap_lift (ap_adic for continuity) and hJ = GoodPt; (3) ap_det with kappa_pr and pi_pr; (4) ap_locAlg (GlobalArtinReciprocity, Loc.laHecke, Loc.laCyc, Loc.laMul) and ap_psi, with H.e_pos; (5) end2oN, which is arithmeticPointClassicalityDense, with isIrr_twist_and_repIso_of_twist_repIso.

\subsection{The multiplicity bound}\label{subsection-sha-multiplicity-bound}

\begin{definition}\label{sha-automorphy-unramified-multiplicity}
Let $\mathcal O$ be a ring, let $\Lambda$ and $B$ be $\mathcal O$-algebras with an $\mathcal O$-algebra map $\lambda:\Lambda\to B$, and let $J$ be an ideal of $B$.
\begin{enumerate}
\item A \emph{Hecke module} for $(\lambda,J)$ is a $B$-module $M$, on which $\Lambda$ acts through $\lambda$, that is finite and free over $\Lambda$, such that: $J$ is radical; a prime of $B$ contains the annihilator of $M$ if and only if it contains $J$; and $\lambda(c)a\in J$ with $c\in\Lambda$ nonzero implies $a\in J$.
\item Let $M$ be a Hecke module and $e$ a $\Lambda$-linear endomorphism of $M$ with $e^2=e$ that commutes with the action of $B$. For a field $K$ and an $\mathcal O$-algebra map $y:B\to K$, view $K$ as a $\Lambda$-algebra through $y\circ\lambda$. The \emph{multiplicity} $\mu_e(y)$ is the dimension over $K$ of the space of vectors of $M\otimes_\Lambda K$ fixed by $e$ that lie, for every $a\in B$, in the generalized eigenspace of $a$ for the eigenvalue $y(a)$. Write $\mu(y)$ for $\mu_e(y)$ with $e=1$.
\item In Situation~\ref{sha-hecke-situation}, call a point $x$ of $B$ \emph{irreducible} if $x(D^{\mathrm u})$ is the determinant of a continuous irreducible representation of $\operatorname{Gal}(\bar L/L)$. An ideal $J$ of the universal ordinary determinant ring $B$ satisfies the \emph{multiplicity bound} if there are a Hecke module $M$ for the map $\Lambda\to B$ and $J$, and an idempotent $e$ as in (2), such that $1\le\mu_e(x)\le\mu_e(y)$ for all irreducible points $x,y$ on the zero locus of $J$ and every increasing weight $h$ such that $x$ is arithmetic of weight $h$ and $y$ is of distinct type of weight $h$ (Definition~\ref{sha-hecke-arcs}).
\end{enumerate}
\end{definition}
\leanref{Development/AutomorphicArithmetic/AutomorphicVarieties/ModularityArgument/LinkedUnramifiedChains/CmpPrime.lean}{UpperTheory.V.J.Dat}
\leanref{Development/AutomorphicArithmetic/AutomorphicVarieties/ModularityArgument/CuspidalDescent/GLDenseRanges/MuRange.lean}{UpperTheory.V.Z.Desc.V0.muE}
\leanref{Development/AutomorphicArithmetic/AutomorphicVarieties/ModularityArgument/CuspidalDescent/GLDenseRanges/MuRange.lean}{UpperTheory.V.Z.Desc.V0.Idem}
\leanref{Development/AutomorphicArithmetic/AutomorphicVarieties/ModularityArgument/LinkedUnramifiedChains/AbsGood.lean}{UpperTheory.V.J.IrrPt}
\leanref{Development/AutomorphicArithmetic/AutomorphicVarieties/ModularityArgument/LinkedUnramifiedChains/AbsGood.lean}{UpperTheory.V.J.Num}
\leanref{Development/AutomorphicArithmetic/AutomorphicVarieties/ModularityArgument/CuspidalDescent/GLDenseRanges/EigBoxE.lean}{UpperTheory.V.Z.Desc.V0.NumE}

\begin{definition}\label{sha-automorphy-unramified-hecke-local-kirillov-forms}
Let $F$ be a nonarchimedean local field with ring of integers $\mathcal O_F$, maximal ideal $\mathfrak m_F$ and residue field $k_F=\mathcal O_F/\mathfrak m_F$ of $q$ elements, and let $n\ge1$. Let $K_0=GL_n(\mathcal O_F)$, let $K_1$ be the kernel of the reduction $K_0\to GL_n(k_F)$, let $N\subset GL_n(F)$ be the group of upper triangular unipotent matrices, and let $N_0=N\cap K_0$.
\begin{enumerate}
\item A \emph{character} of $F$ is a continuous homomorphism $\psi$ from the additive group of $F$ to $\mathbf C^\times$. It is \emph{of level one} if it is trivial on $\mathfrak m_F$ and nontrivial on $\mathcal O_F$. For $t\in F^n$ and $u\in N$ put $\psi_t(u)=\psi\big(\sum_{i<n}t_iu_{i,i+1}\big)$; this is a character of $N$, since $(uu')_{i,i+1}=u_{i,i+1}+u'_{i,i+1}$ for $u,u'\in N$.
\item Let $(\rho,W)$ be a complex representation of $GL_n(F)$. A \emph{Kirillov form} for $(\psi,t)$ is a linear form $l$ on $W$ such that $l(\rho(u)v)=\psi_t(u)\,l(v)$ for $u\in N$ and $v\in W$, and such that $v=0$ whenever $l(\rho(p)v)=0$ for every $p\in GL_n(F)$ with last row $(0,\dots,0,1)$.
\item $W(\psi,t)$ is the space of the $v\in W$ fixed by $\rho(K_1)$ with $\rho(u)v=\psi_t(u)v$ for every $u\in N_0$.
\item For a representation $\rho$ of $GL_n(F)$ on a module over a commutative ring and a finitely supported function $a:GL_n(F)\to\mathbf Z$, $E_a=\sum_ga(g)\rho(g)$.
\end{enumerate}
\end{definition}
\leanref{Development/AutomorphicArithmetic/AutomorphicVarieties/ModularityArgument/CuspidalDescent/GLDenseRanges/KirCusp.lean}{UpperTheory.V.Z.Desc.V0.chr}
\leanref{Development/AutomorphicArithmetic/AutomorphicVarieties/ModularityArgument/CuspidalDescent/GLDenseRanges/KirCusp.lean}{UpperTheory.V.Z.Desc.V0.Lev1}
\leanref{Development/AutomorphicArithmetic/AutomorphicVarieties/ModularityArgument/CuspidalDescent/GLDenseRanges/KirCusp.lean}{UpperTheory.V.Z.Desc.V0.KirF}
\leanref{Development/AutomorphicArithmetic/AutomorphicVarieties/ModularityArgument/CuspidalDescent/GLDenseRanges/KirCusp.lean}{UpperTheory.V.Z.Desc.V0.isoPart}
\leanref{Development/AutomorphicArithmetic/AutomorphicVarieties/ModularityArgument/CuspidalDescent/GLDenseRanges/KirCusp.lean}{UpperTheory.V.Z.Desc.V0.opE}

\begin{lemma}\label{sha-automorphy-unramified-hecke-local-sums}
Let $G$ and $A$ be groups, $K\subset G$ a subgroup, $Q$ a finite group, and $\iota:A\to G$ and $\pi:A\to Q$ homomorphisms such that, for $y\in A$, $\pi(y)=1$ if and only if $\iota(y)\in K$, and $\iota(y)^{-1}K\iota(y)\subseteq K$. Let $s:Q\to A$ be a map with $\pi(s(x))=x$ for $x\in Q$, let $(\chi_f)_{f\in I}$ be a finite family of pairwise distinct homomorphisms $Q\to\mathbf C^\times$, and let $c:Q\to\mathbf Z$ satisfy $c(x)=\sum_{f\in I}\chi_f(x)^{-1}$ for $x\in Q$. For a representation $\rho$ of $G$ let $E=\sum_{x\in Q}c(x)\rho(\iota(s(x)))$.
\begin{enumerate}
\item For every representation $\rho$ of $G$ on a module over a commutative ring and every $v$ fixed by $\rho(K)$, $E(Ev)=|Q|\,Ev$.
\item For every complex representation $(\rho,W)$ of $G$, let $W^K$ be the space of the vectors fixed by $\rho(K)$, and for $f\in I$ let $W_f$ be the space of the $v\in W^K$ with $\rho(\iota(y))v=\chi_f(\pi(y))v$ for all $y\in A$. Then the space of the $v\in W^K$ with $Ev=|Q|v$ has dimension $\sum_{f\in I}\dim W_f$.
\end{enumerate}
\end{lemma}
\leanref{Development/AutomorphicArithmetic/AutomorphicVarieties/BuildingsResidueChars/MemFixedOfCharacterSumProjection.lean}{UpperTheory.mem_fixed_of_character_sum_projection}
\leanref{Development/AutomorphicArithmetic/AutomorphicVarieties/BuildingsResidueChars/RankFixedInfCharacterEigenspace.lean}{UpperTheory.rank_fixed_inf_character_eigenspace}
\begin{proof}
Let $v$ be fixed by $\rho(K)$, and let $x,y\in Q$. The element $s(yx)^{-1}s(y)s(x)$ of $A$ lies in the kernel of $\pi$, so its image under $\iota$ lies in $K$, and $\rho(\iota s(y))\rho(\iota s(x))v=\rho(\iota s(yx))v$.
For $f,f'\in I$, $\sum_{y\in Q}\chi_f(y)^{-1}\chi_{f'}(y)$ is $|Q|$ if $f=f'$ and $0$ otherwise: if $f\ne f'$, choose $z$ with $\chi_f(z)\ne\chi_{f'}(z)$; replacing $y$ by $zy$ multiplies the sum by $\chi_f(z)^{-1}\chi_{f'}(z)\ne1$.
\begin{enumerate}
\item For $z\in Q$, $c(y^{-1}z)=\sum_{f'}\chi_{f'}(y)\chi_{f'}(z)^{-1}$, so
\[\sum_{y\in Q}c(y)c(y^{-1}z)=\sum_{f,f'}\chi_{f'}(z)^{-1}\sum_y\chi_f(y)^{-1}\chi_{f'}(y)=|Q|\,c(z).\]
By the first paragraph, $E(Ev)=\sum_{x,y}c(y)c(x)\rho(\iota s(yx))v=\sum_z\big(\sum_yc(y)c(y^{-1}z)\big)\rho(\iota s(z))v=|Q|\,Ev$.
\item For $y\in A$, $\rho(\iota(y))$ maps $W^K$ into itself, as $\iota(y)^{-1}K\iota(y)\subseteq K$, and it agrees there with $\rho(\iota s(\pi(y)))$, since $\iota(s(\pi(y))^{-1}y)\in K$. For $f\in I$ let $e_f=|Q|^{-1}\sum_x\chi_f(x)^{-1}\rho(\iota s(x))$ on $W^K$. By the first paragraph and the orthogonality above, $e_f$ maps $W^K$ into $W_f$, it is the identity on $W_f$, and $e_{f'}$ vanishes on $W_f$ for $f'\ne f$. On $W^K$, $E=|Q|\sum_fe_f$, so $Ev=|Q|v$ if and only if $v=\sum_fe_fv$. Hence $(w_f)_f\mapsto\sum_fw_f$ is an isomorphism from $\bigoplus_fW_f$ onto the space of these $v$, with inverse $v\mapsto(e_fv)_f$.
\end{enumerate}
\end{proof}
% mem_fixed_of_character_sum_projection: hT (the first paragraph), chi_orth and c_conv (the convolution identity), opE_sum. rank_fixed_inf_character_eigenspace: hfix, hV, hmul, hR, then the private rank_isotypic_character_projection (the projections e, heV, hproj, hP, hker and the sum map, hsum, hsumV).

\begin{lemma}\label{sha-automorphy-unramified-hecke-local-projector}
In Definition~\ref{sha-automorphy-unramified-hecke-local-kirillov-forms}, there are a character $\psi$ of $F$ of level one, a finitely supported function $a:GL_n(F)\to\mathbf Z$ with support in $N_0$, an integer $j\ge0$ and a nonempty finite set $T\subset(\mathcal O_F^\times)^n$ such that:
\begin{enumerate}
\item $g^{-1}K_1g=K_1$ for $g\in K_0$;
\item $E_a(E_av)=q^jE_av$ for every representation $\rho$ of $GL_n(F)$ on a module over a commutative ring and every $v$ fixed by $\rho(K_1)$;
\item for every complex representation $(\rho,W)$ of $GL_n(F)$, the space of the $v\in W$ fixed by $\rho(K_1)$ with $E_av=q^jv$ has dimension $\sum_{t\in T}\dim W(\psi,t)$.
\end{enumerate}
\end{lemma}
\leanref{Development/AutomorphicArithmetic/AutomorphicVarieties/ModularityArgument/CuspidalDescent/GLDenseRanges/InertiaPartMaps/MapInj.lean}{UpperTheory.V.Z.Desc.V0.eOp}
\leanref{Development/AutomorphicArithmetic/AutomorphicVarieties/ModularityArgument/CuspidalDescent/GLDenseRanges/KirCusp.lean}{UpperTheory.V.Z.Desc.V0.EOp}
\begin{proof}
\begin{enumerate}
\item[(i)] \emph{The character.} The subgroup $\mathfrak m_F$ of $F$ is open and does not contain $1$. By Pontryagin duality for the discrete group $F/\mathfrak m_F$ there is a homomorphism $\psi:F\to\mathbf C^\times$, with values of absolute value $1$, trivial on $\mathfrak m_F$ and with $\psi(1)\ne1$. It is continuous, being trivial on an open subgroup, so it is a character of level one. On $\mathcal O_F$, $\psi(y)=\phi(\bar y)$ for a character $\phi$ of the additive group of $k_F$, nontrivial as $\phi(1)=\psi(1)$.
\item[(ii)] \emph{The residual characters.} Let $Q$ be the group of upper triangular unipotent matrices over $k_F$. Its entries above the diagonal are arbitrary, so $|Q|=q^j$ with $j=n(n-1)/2$. For $f=(f_1,\dots,f_{n-1})\in(k_F^\times)^{n-1}$ let $\chi_f(x)=\phi\big(\sum_{i<n}f_ix_{i,i+1}\big)$, a homomorphism $Q\to\mathbf C^\times$ as in Definition~\ref{sha-automorphy-unramified-hecke-local-kirillov-forms}(1). These are pairwise distinct: if $f_i\ne f'_i$ and $b\in k_F$, the matrix $x$ equal to $1$ except for the entry $(f_i-f'_i)^{-1}b$ at $(i,i+1)$ has $\chi_f(x)\chi_{f'}(x)^{-1}=\phi(b)$, which is $\ne1$ for some $b$. As $\sum_{b\in k_F^\times}\phi(-by)$ is $q-1$ for $y=0$ and $-1$ otherwise,
\[\sum_f\chi_f(x)^{-1}=\prod_{i<n}\sum_{b\in k_F^\times}\phi(-bx_{i,i+1})=c(x),\]
where $c(x)=\prod_{i<n}c_i(x)$, with $c_i(x)=q-1$ if $x_{i,i+1}=0$ and $c_i(x)=-1$ otherwise. So $c$ takes values in $\mathbf Z$.
\item[(iii)] \emph{The operator.} Reduction modulo $\mathfrak m_F$ is a homomorphism $\pi:N_0\to Q$, onto since a matrix over $k_F$ lifts entry by entry, with kernel $N_0\cap K_1$. Choose $s:Q\to N_0$ with $\pi\circ s=\mathrm{id}$, and let $a$ be the function with value $c(x)$ at $s(x)$ for $x\in Q$ and $0$ elsewhere, so that $E_a=\sum_{x\in Q}c(x)\rho(s(x))$ and $a$ has support in $N_0$.
\item[(iv)] \emph{(1) to (3).} The group $K_1$ is the kernel of a homomorphism on $K_0$, which gives (1). Lemma~\ref{sha-automorphy-unramified-hecke-local-sums} applies to $G=GL_n(F)$, $K=K_1$, $A=N_0$, $\iota$ the inclusion, $\pi$, $s$, the $\chi_f$ and $c$: an element $y\in N_0$ has $\pi(y)=1$ if and only if $y\in K_1$, and $N_0\subset K_0$ normalizes $K_1$ by (1). Its part (1), with $|Q|=q^j$, gives (2). Its part (2) gives that the space of (3) has dimension $\sum_f\dim W_f$, where $W_f$ is the space of the $v$ fixed by $\rho(K_1)$ with $\rho(y)v=\chi_f(\pi(y))v$ for $y\in N_0$.
\item[(v)] \emph{The set $T$.} For $f$ as in (ii), let $t(f)\in(\mathcal O_F^\times)^n$ have $t(f)_n=1$ and $t(f)_i$ a lift of $f_i$ for $i<n$, a unit as $f_i\ne0$. For $y\in N_0$ the element $\sum_it(f)_iy_{i,i+1}$ lies in $\mathcal O_F$ and reduces to $\sum_if_i\bar y_{i,i+1}$, so $\psi_{t(f)}(y)=\chi_f(\pi(y))$ by (i), and $W_f=W(\psi,t(f))$. The map $f\mapsto t(f)$ is injective, as $f_i$ is the reduction of $t(f)_i$. Its image $T$ is finite and nonempty, and $\sum_f\dim W_f=\sum_{t\in T}\dim W(\psi,t)$, which gives (3).
\end{enumerate}
\end{proof}
% eOp (MapInj.lean): (i) is UpperTheory.exists_level_one_addChar_lifting_residue (the private sep_char and cont_char); (ii) is UpperTheory.exists_injective_generic_unipotent_characters (chi, chi_inj, sum_units_psi, chi_inv) and UpperTheory.card_unipotent_upper_eq_pow; (iii) is pA, pA_surj (unipotent_upper_map_surjective), pA_ker and s = surjInv; (iv) is Sig.Kc_conj, mem_fixed_of_character_sum_projection and rank_fixed_inf_character_eigenspace; (v) is lift, hlp, hinj, hunit and key.

\begin{lemma}\label{sha-automorphy-unramified-hecke-local-move}
In Definition~\ref{sha-automorphy-unramified-hecke-local-kirillov-forms}, let $(\rho,W)$ be a complex representation of $GL_n(F)$ with a Kirillov form for $(\psi_0,(1,\dots,1))$, for a nontrivial character $\psi_0$ of $F$. Then $(\rho,W)$ has a Kirillov form for $(\psi,t)$ for every nontrivial character $\psi$ of $F$ and every $t\in F^n$ with nonzero entries.
\end{lemma}
\leanref{Development/AutomorphicArithmetic/AutomorphicVarieties/ModularityArgument/CuspidalDescent/GLDenseRanges/InertiaPartMaps/KirMove.lean}{UpperTheory.V.Z.Desc.V0.kirMove}
\leanref{Development/AutomorphicArithmetic/AutomorphicVarieties/ModularityArgument/CuspidalDescent/GLDenseRanges/KirCusp.lean}{UpperTheory.V.Z.Desc.V0.KirMove}
\begin{proof}
Let $l_0$ be the Kirillov form for $(\psi_0,(1,\dots,1))$. There is $b\in F^\times$ with $\psi(x)=\psi_0(bx)$ for all $x\in F$ \cite[\S2.2]{Tat67}. Put $c_i=bt_i\ne0$ and $e_i=\prod_{i\le i'<n}c_{i'}$, so that $e_n=1$ and $e_i=c_ie_{i+1}$ for $i<n$. Let $g=\operatorname{diag}(e_1,\dots,e_n)$ and $l=l_0\circ\rho(g)$.
For $u\in N$, $gug^{-1}\in N$ and $(gug^{-1})_{i,i+1}=e_iu_{i,i+1}e_{i+1}^{-1}=bt_iu_{i,i+1}$, so
\[l(\rho(u)v)=l_0\big(\rho(gug^{-1})\rho(g)v\big)=\psi_0\Big(\sum_{i<n}bt_iu_{i,i+1}\Big)l_0(\rho(g)v)=\psi_t(u)\,l(v).\]
Suppose that $l(\rho(p)v)=0$ for every $p$ with last row $(0,\dots,0,1)$, and let $p'$ be such a matrix. Since $e_n=1$, $g^{-1}p'$ has the same last row, so $l_0(\rho(p')v)=l(\rho(g^{-1}p')v)=0$. As $l_0$ is a Kirillov form, $v=0$.
\end{proof}
% kirMove (KirMove.lean): GlobalTheory.LocGL.exists_unit_eq_mulShift for b; c, ee (e_i), ee_step, ee_last and g = dg (ee c); conj_apply for the conjugate and inv_mul_apply for the last row.

\begin{lemma}\label{sha-automorphy-unramified-hecke-local-nonzero}
In Definition~\ref{sha-automorphy-unramified-hecke-local-kirillov-forms}, let $\psi$ be a character of $F$ of level one, let $t\in(\mathcal O_F^\times)^n$, and let $(\rho,W)$ be a smooth complex representation of $GL_n(F)$, that is, one in which every vector has an open stabilizer (the proof does not use this hypothesis), with a Kirillov form $l$ for $(\psi,t)$ and a nonzero vector fixed by $\rho(K_1)$. Then $W(\psi,t)\ne0$.
\end{lemma}
\leanref{Development/AutomorphicArithmetic/AutomorphicVarieties/ModularityArgument/CuspidalDescent/GLDenseRanges/InertiaPartMaps/ValOrd.lean}{UpperTheory.V.Z.Desc.V0.isoNe}
\leanref{Development/AutomorphicArithmetic/AutomorphicVarieties/ModularityArgument/CuspidalDescent/GLDenseRanges/KirCusp.lean}{UpperTheory.V.Z.Desc.V0.IsoNe}
\begin{proof}
Let $\varpi$ be a uniformizer and $\operatorname{ord}$ the valuation of $F$ with $\operatorname{ord}(\varpi)=1$. For $x\in N\cap K_1$, $\psi_t(x)=1$: the entries $x_{i,i+1}$ lie in $\mathfrak m_F$, hence so does $\sum_it_ix_{i,i+1}$. As $K_1$ is normal in $K_0$, $\rho(g)$ maps vectors fixed by $\rho(K_1)$ to vectors fixed by $\rho(K_1)$ for $g\in K_0$.
\begin{enumerate}
\item \emph{A diagonal matrix.} Let $w_0\ne0$ be fixed by $\rho(K_1)$. As $l$ is a Kirillov form, $l(\rho(p)w_0)\ne0$ for some $p\in GL_n(F)$. By the Iwasawa decomposition, $p=uDk$ with $u\in N$, $D=\operatorname{diag}(d_1,\dots,d_n)$ and $k\in K_0$. Then $w_1=\rho(k)w_0$ is fixed by $\rho(K_1)$, and $l(\rho(p)w_0)=\psi_t(u)\,l(\rho(D)w_1)$, so $l(\rho(D)w_1)\ne0$.
\item \emph{Dominance.} We show $\operatorname{ord}(d_1)\ge\dots\ge\operatorname{ord}(d_n)$. For $x\in N\cap K_1$, $DxD^{-1}\in N$ and $\rho(DxD^{-1})\rho(D)w_1=\rho(D)\rho(x)w_1=\rho(D)w_1$, so $l(\rho(D)w_1)=\psi_t(DxD^{-1})\,l(\rho(D)w_1)$ and $\psi_t(DxD^{-1})=1$. Suppose $\operatorname{ord}(d_{i+1})>\operatorname{ord}(d_i)$ for some $i<n$. Let $x_0\in\mathcal O_F$ with $\psi(x_0)\ne1$, let $y=x_0d_{i+1}/(t_id_i)$, and let $x\in N$ be equal to $1$ except for the entry $y$ at $(i,i+1)$. Then $\operatorname{ord}(y)\ge\operatorname{ord}(d_{i+1})-\operatorname{ord}(d_i)\ge1$, so $x\in N\cap K_1$, while $\psi_t(DxD^{-1})=\psi(t_id_iyd_{i+1}^{-1})=\psi(x_0)\ne1$, a contradiction.
\item \emph{A subgroup fixing $\rho(D)w_1$.} Let $B=DK_1D^{-1}$, which fixes $\rho(D)w_1$, and let $m=\sum_i|\operatorname{ord}(d_i)|$. An element $g\in K_1$ congruent to $1$ modulo $\varpi^{2m+1}$ lies in $B$: the entries of $D^{-1}gD-1$ are $d_i^{-1}(g-1)_{ii'}d_{i'}$, of order at least $2m+1+\operatorname{ord}(d_{i'})-\operatorname{ord}(d_i)\ge1$, and likewise for $g^{-1}$. So $B\cap K_1$ contains the principal congruence subgroup of level $\varpi^{2m+1}$, and it has finite index in $K_1$.
\item \emph{$K_1$ keeps the value of $l$.} Let $g\in K_1$. By the Iwahori factorization, $g=xsy$ with $x\in N\cap K_1$, $s$ diagonal in $K_1$ and $y$ lower triangular unipotent in $K_1$. The entries of $D^{-1}yD$ below the diagonal are $d_i^{-1}y_{ii'}d_{i'}$ with $i>i'$, of order at least $\operatorname{ord}(y_{ii'})\ge1$ by (2); so $D^{-1}yD\in K_1$, $y\in B$, and $\rho(y)\rho(D)w_1=\rho(D)w_1$. The matrix $s$ commutes with $D$ and fixes $w_1$, so $\rho(s)\rho(D)w_1=\rho(D)w_1$. Hence $l(\rho(g)\rho(D)w_1)=\psi_t(x)\,l(\rho(D)w_1)=l(\rho(D)w_1)$.
\item \emph{Averaging over $K_1$.} Let $w_3=\sum_g\rho(g)\rho(D)w_1$, the sum over representatives $g$ of $K_1/(B\cap K_1)$; its terms do not depend on the representatives, by (3). It is fixed by $\rho(K_1)$, which permutes the cosets, and $l(w_3)=[K_1:B\cap K_1]\,l(\rho(D)w_1)\ne0$ by (4).
\item \emph{Averaging over $N_0$.} The subgroup $N_0\cap K_1$ has finite index in $N_0$, and $\psi_t$ is trivial on it. Let $w=\sum_u\psi_t(u)^{-1}\rho(u)w_3$, the sum over representatives $u$ of $N_0/(N_0\cap K_1)$; its terms do not depend on the representatives, as $\rho(N_0\cap K_1)$ fixes $w_3$. Each term is fixed by $\rho(K_1)$, since $u\in K_0$, so $w$ is. For $u'\in N_0$, $\rho(u')w=\psi_t(u')w$, as $u'$ permutes the cosets. Since $l(\rho(u)w_3)=\psi_t(u)\,l(w_3)$, $l(w)=[N_0:N_0\cap K_1]\,l(w_3)\ne0$. So $w$ is a nonzero element of $W(\psi,t)$.
\end{enumerate}
\end{proof}
% isoNe (ValOrd.lean): chr_one for psi_t on N cap K_1; hconj (Kc_conj'); (1) iwasawa (exists_unipB_mul_diagT_mul_KO); (2) dom (el, el_unip, el_Kc, chr_el); (3) Bg, hle with isoNe_lemma_2 and Kc_index; (4) hl1 with iwahori (GlAx.udl) and conj_Nm; (5) and (6) are the private avg, with chi = 1 on K_1 and chi = chr psi t on unipB cap KO (isoNe_lemma for multiplicativity). The smoothness hypothesis of IsoNe is not used.

\begin{lemma}\label{sha-automorphy-unramified-hecke-local-injective}
In Definition~\ref{sha-automorphy-unramified-hecke-local-kirillov-forms}, let $\psi$ be a character of $F$ of level one, let $t\in(\mathcal O_F^\times)^n$, and let $(\rho,W)$ be a complex representation of $GL_n(F)$ with a Kirillov form $l$ for $(\psi,t)$. Then the only $v\in W(\psi,t)$ with $l(v)=0$ is $v=0$.
\end{lemma}
\leanref{Development/AutomorphicArithmetic/AutomorphicVarieties/ModularityArgument/CuspidalDescent/GLDenseRanges/InertiaPartMaps/TvApply.lean}{UpperTheory.V.Z.Desc.V0.isoLe}
\leanref{Development/AutomorphicArithmetic/AutomorphicVarieties/ModularityArgument/CuspidalDescent/GLDenseRanges/KirCusp.lean}{UpperTheory.V.Z.Desc.V0.IsoLe}
\begin{proof}
Let $v\in W(\psi,t)$ with $l(v)=0$. Call row $a$ of a matrix $p$ \emph{unitriangular} if $p_{aa'}=0$ for $a'<a$ and $p_{aa}=1$; write $e_1,\dots,e_n$ for the standard basis of column vectors. We show, by induction on $r$ with $0\le r\le n-1$, that $l(\rho(p)v)=0$ for every $p\in GL_n(F)$ whose rows $a>r$ are unitriangular. For $r=n-1$ these are the $p$ with last row $(0,\dots,0,1)$, so then $v=0$, as $l$ is a Kirillov form.
For $r=0$, $p\in N$ and $l(\rho(p)v)=\psi_t(p)\,l(v)=0$. Let $1\le i\le n-1$, assume the claim for $r=i-1$, and let $p$ have unitriangular rows $a>i$. Suppose that $l(\rho(p)v)\ne0$.
\begin{enumerate}
\item \emph{Row $i$ modulo $\mathfrak m_F$.} Let $j\le i$ and $b\in\mathcal O_F$, and let $u=1+be_je_{i+1}^{\mathsf T}$, equal to $1$ except for the entry $b$ at $(j,i+1)$. Then $u\in N_0$, and $\psi_t(u)=\psi(\delta_{ij}t_ib)$, so $\rho(u)v=\psi(\delta_{ij}t_ib)v$. The rows $a>i$ of $p^{-1}$ are unitriangular, as those of $p$ are, so row $i+1$ of $p^{-1}$ is a row vector $w$ with $w_{a'}=0$ for $a'\le i$ and $w_{i+1}=1$. Then $pu=u'p$ with $u'=1+yw$ and $y=bpe_j$. As $j\le i$, $y_a=bp_{aj}=0$ for $a>i$; so $u'\in N$, and of its entries at the places $(a,a+1)$ only the one at $(i,i+1)$, which is $y_iw_{i+1}=bp_{ij}$, can be nonzero. Hence
\[\psi(\delta_{ij}t_ib)\,l(\rho(p)v)=l(\rho(pu)v)=l(\rho(u'p)v)=\psi(t_ibp_{ij})\,l(\rho(p)v),\]
and $\psi(t_i(p_{ij}-\delta_{ij})b)=1$ for all $b\in\mathcal O_F$. If $t_i(p_{ij}-\delta_{ij})$ were not in $\mathfrak m_F$, $\psi$ would be trivial on $t_i(p_{ij}-\delta_{ij})\mathcal O_F\supseteq\mathcal O_F$. So $p_{ij}\equiv\delta_{ij}\bmod\mathfrak m_F$ for $j\le i$, as $t_i$ is a unit.
\item \emph{Clearing row $i$.} Let $d\in F^n$ be the row vector with $d_j=p_{ij}-\delta_{ij}$ for $j\le i$ and $d_j=0$ for $j>i$; its entries lie in $\mathfrak m_F$. The matrix $M=1+e_id$, equal to $1$ except in row $i$, which is $e_i^{\mathsf T}+d$, lies in $K_1$, as $1+d_i$ is a unit. Let $p'=pM^{-1}$, so that $p_{aj}=p'_{aj}+p'_{ai}d_j$ for all $a,j$. For $a=i$ and $j=i$ this gives $p'_{ii}(1+d_i)=p_{ii}=1+d_i$, so $p'_{ii}=1$, and then $p'_{ij}=p_{ij}-d_j=0$ for $j<i$. For $a>i$ it gives $p'_{ai}(1+d_i)=p_{ai}=0$, so $p'_{ai}=0$ and row $a$ of $p'$ is that of $p$. So the rows $a\ge i$ of $p'$ are unitriangular, and $l(\rho(p')v)=0$ by the induction hypothesis. But $M\in K_1$ fixes $v$, so $l(\rho(p)v)=l(\rho(p')\rho(M)v)=0$, a contradiction.
\end{enumerate}
\end{proof}
% isoLe (TvApply.lean): descent (the induction on r), rowStep (tv x e, unip_of, chr_of, conj_eq with inv_rows), lev_lt for the level one character, and exists_mem_Kc_rows_unitriangular (M = tv y0 d0, in Kc 1).

\begin{lemma}\label{sha-automorphy-unramified-hecke-local}
Let $F$ be a nonarchimedean local field with ring of integers $\mathcal O_F$ and residue field of $q$ elements, let $n\ge1$, let $K_0=GL_n(\mathcal O_F)$, and let $K_1$ be the kernel of the reduction $K_0\to GL_n(\mathcal O_F/\mathfrak m_F)$. For a representation $\rho$ of $GL_n(F)$ on a module over a commutative ring and a finitely supported function $a:GL_n(F)\to\mathbf Z$, write $E_a=\sum_ga(g)\rho(g)$. Say that a complex representation $(\rho,W)$ of $GL_n(F)$ has a \emph{Kirillov model} if there are a nontrivial continuous character $\psi$ of $F$ and a linear form $l$ on $W$ such that $l(\rho(u)v)=\psi(u_{12}+u_{23}+\dots+u_{n-1,n})\,l(v)$ for every upper triangular unipotent $u\in GL_n(F)$ and every $v\in W$, and such that $v=0$ whenever $l(\rho(p)v)=0$ for every $p\in GL_n(F)$ with last row $(0,\dots,0,1)$. Then there are a finitely supported $a:GL_n(F)\to\mathbf Z$ with support in $K_0$ and integers $j\ge0$ and $d_0\ge1$ such that:
\begin{enumerate}
\item $g^{-1}K_1g=K_1$ for $g\in K_0$, and $E_a(E_av)=q^jE_av$ for every representation $\rho$ of $GL_n(F)$ on a module over a commutative ring and every $v$ fixed by $\rho(K_1)$;
\item for every smooth complex representation $(\rho,W)$ of $GL_n(F)$ that has a Kirillov model and a nonzero vector fixed by $\rho(K_1)$, the space of the $v\in W$ fixed by $\rho(K_1)$ with $E_av=q^jv$ has dimension $d_0$.
\end{enumerate}
\end{lemma}
\leanref{Development/AutomorphicArithmetic/AutomorphicVarieties/ModularityArgument/CuspidalDescent/GLDenseRanges/InertiaPartMaps/RankEqOneOfInj.lean}{UpperTheory.V.Z.Desc.V0.locE}
\leanref{Development/AutomorphicArithmetic/AutomorphicVarieties/ModularityArgument/CuspidalDescent/GLDenseRanges/KirCusp.lean}{UpperTheory.V.Z.Desc.V0.LocE}
\leanref{Development/AutomorphicArithmetic/AutomorphicVarieties/ModularityArgument/CuspidalDescent/GLDenseRanges/KirCusp.lean}{UpperTheory.V.Z.Desc.V0.Kir}
\begin{proof}
Use the notation of Definition~\ref{sha-automorphy-unramified-hecke-local-kirillov-forms}; a complex representation has a Kirillov model if and only if it has a Kirillov form for $(\psi,(1,\dots,1))$ for some nontrivial character $\psi$ of $F$. Take $\psi$, $a$, $j$ and $T$ as in Lemma~\ref{sha-automorphy-unramified-hecke-local-projector}, and $d_0=|T|\ge1$. The support of $a$ lies in $N_0\subset K_0$, and (1) is (1) and (2) of that lemma. Let $(\rho,W)$ be as in (2), and let $t\in T$. The entries of $t$ are units, hence nonzero, and $\psi$ is nontrivial, so Lemma~\ref{sha-automorphy-unramified-hecke-local-move} gives a Kirillov form $l$ for $(\psi,t)$. By Lemma~\ref{sha-automorphy-unramified-hecke-local-nonzero}, $W(\psi,t)\ne0$, and by Lemma~\ref{sha-automorphy-unramified-hecke-local-injective}, $l$ is injective on $W(\psi,t)$. So $\dim W(\psi,t)=1$, and by Lemma~\ref{sha-automorphy-unramified-hecke-local-projector}(3) the space of (2) has dimension $|T|=d_0$.
% EOp is eOp (Lemma sha-automorphy-unramified-hecke-local-projector), KirMove is kirMove (-move), IsoNe is isoNe (-nonzero) and IsoLe is isoLe (-injective); W(psi,t) is isoPart and psi_t is chr psi t.
\end{proof}
% locE (RankEqOneOfInj.lean): EOp gives psi, a, j and T; hpsine; rank_eq_one_of_inj with KirMove, IsoNe and IsoLe for each t in T; d_0 = T.card.

\begin{lemma}[{\cite{Sha74,GK75}}]\label{sha-automorphy-unramified-hecke-generic}
Let $E$ be a totally complex number field, let $n\ge1$, let $\Pi$ be a cuspidal automorphic representation of $GL_n(\mathbf A_E)$, and let $w$ be a finite place of $E$. Then every local component of $\Pi$ at $w$ has a Kirillov model, in the sense of Lemma~\ref{sha-automorphy-unramified-hecke-local}.
\end{lemma}
\leanref{Development/AutomorphicArithmetic/AutomorphicVarieties/ModularityArgument/CuspidalDescent/GLDenseRanges/InertiaPartMaps/KirCusp.lean}{UpperTheory.V.Z.Desc.V0.kirCusp}
\leanref{Development/AutomorphicArithmetic/AutomorphicVarieties/ModularityArgument/CuspidalDescent/GLDenseRanges/KirCusp.lean}{UpperTheory.V.Z.Desc.V0.KirCusp}

\begin{lemma}\label{sha-automorphy-unramified-hecke-projector}
In Situation~\ref{sha-hecke-ideal-tw-augmented-setup}, suppose that $S'=S$, and let $U=\mathcal K_0$ be the level, so that $M=M_{()}$. Let $K_0=GL_n(\mathcal O_{L_{\tilde v_0}})$, let $K_1$ be the kernel of its reduction modulo $\tilde v_0$, let $q$ be the number of elements of the residue field of $L_{\tilde v_0}$, and let $a$ and $j$ be as in Lemma~\ref{sha-automorphy-unramified-hecke-local} for $F=L_{\tilde v_0}$. For $g\in GL_n(L_{\tilde v_0})$, viewed in $G^{(\ell)}$ through $\tilde v_0$, let $\rho_0(g)$ be right translation by $g$ on functions on $G^{(\ell)}\times G_\ell$, and let $E=\sum_ga(g)\rho_0(g)$. Then $E$ maps $M$ into itself, $q^j$ is a unit of $\Lambda$, and $e=q^{-j}E$ is a $\Lambda$-linear endomorphism of $M$ with $e^2=e$ that commutes with the operators $T_{w,i}$ and $U_{v,i}$ on $M$, and so with $\mathbf T$ if $M\ne0$.
\end{lemma}
\leanref{Development/AutomorphicArithmetic/AutomorphicVarieties/ModularityArgument/CuspidalDescent/GLDenseRanges/InertiaPartMaps/CommFinsum.lean}{UpperTheory.V.Z.Desc.V0.eEx}
\leanref{Development/AutomorphicArithmetic/AutomorphicVarieties/ModularityArgument/CuspidalDescent/GLDenseRanges/InertiaPartMaps/CmpOverL.lean}{UpperTheory.V.Z.Desc.V0.normS}
\begin{proof}
\begin{enumerate}
\item \emph{Normalization.} Let $g\in K_0$ and $u\in U$. The element $g^{-1}ug$ has the components of $u$ away from $\tilde v_0$, and its component at $\tilde v_0$ lies in $K_1$, as $K_1$ is normal in $K_0$ by Lemma~\ref{sha-automorphy-unramified-hecke-local}(1). So $g^{-1}ug\in U$, and $\tau_0(g^{-1}ug)=\tau_0(u)$, as $\tau_0$ depends only on the components at the places of $S_{\mathrm a}$.
\item \emph{Translation by $K_0$.} Let $g\in K_0$. Right translation by $g$ commutes with left translation by $U(n)(L^+)$, with right translation by $G_\ell$, and with right translation by $GL_n(L_w)$ for every split place $w$ of $L$ above a place outside $S$, as these act at other places. So $\rho_0(g)$ maps $\Lambda$-adic forms to $\Lambda$-adic forms; it commutes with $U_z$, so it maps ordinary forms to ordinary forms; and it commutes with the operators $T_{w,i}$ and $U_{v,i}$, which are sums of such translations. If $f$ is fixed by $U$, then for $u\in U$ and $h\in G^{(\ell)}\times G_\ell$, by (1),
\[(u\cdot\rho_0(g)f)(h)=\tau_0(u)f(hug)=\tau_0(g^{-1}ug)f\big(hg\cdot g^{-1}ug\big)=f(hg),\]
so $\rho_0(g)f$ is fixed by $U$. Hence $\rho_0(g)$ is a $\Lambda$-linear endomorphism of $\mathcal F^{\mathrm{ord}}(U)$ that commutes with the operators. If $f\in M$, then for every operator $t$, with eigenvalue $a_t$, and every $k\ge0$, some power of $t-a_t$ maps $f$ into $\mathfrak m_\Lambda^k\mathcal F^{\mathrm{ord}}(U)$; as $\rho_0(g)$ commutes with $t$ and maps $\mathfrak m_\Lambda^k\mathcal F^{\mathrm{ord}}(U)$ into itself, the same power maps $\rho_0(g)f$ into it. So $\rho_0(g)$ maps the localization $M$ (Lemma~\ref{sha-hecke-ideal-tw-augmented-idempotent}) into itself, and so does $E$, as $a$ has support in $K_0$.
\item \emph{$K_1$ acts trivially.} For $g\in K_1$, the element of $G^{(\ell)}_0$ with component $g$ at $\tilde v_0$ and $1$ elsewhere lies in $U$, and $\tau_0$ is $1$ on it. So $\rho_0(g)f=f$ for $f\in\mathcal F(U)$.
\item \emph{The idempotent.} By (3) and Lemma~\ref{sha-automorphy-unramified-hecke-local}(1), $E(Ef)=q^jEf$ for $f\in M$. The number $q$ is a power of $p$, and $p\ne\ell$, so $q^j$ is a unit of $\mathbf Z_\ell$, hence of $\Lambda$. Let $e=q^{-j}E$ on $M$. Then $q^je=E$, $e^2=q^{-2j}E^2=q^{-j}E=e$ on $M$, and $e$ commutes with the operators by (2). If $M\ne0$, $\mathbf T$ is generated over $\Lambda$ by the operators, so $e$ commutes with $\mathbf T$.
\end{enumerate}
\end{proof}
% eEx (CommFinsum.lean): hrhoM from tr_MK (tr_Ord with ord_pres, amb_pres, fix_pres and valIn_pres, then iSup_pres for the localization), hfixM from tr_fix, hsq from IsE.sq, hcomm from tr_op (comm_ra with Norm.sph for the T operators, comm_U for the U operators), and e = u^-1 E with u the unit N; the commutation with TK through Algebra.commute_of_mem_adjoin_of_forall_mem_commute. normS (CmpOverL.lean): secSp and secS_spec; the conjugate of an element of UQ by sec0 g (huq with Kc_conj and h0, the components at the places of S_a being 1, and htau); sph through comm_of and sec_one; triv from hdesc. The number of elements of the residue field is prime to l by the private card_not_dvd of TZEqDimE.lean.

\begin{lemma}\label{sha-automorphy-unramified-hecke-weights}
In Situation~\ref{sha-hecke-situation}, let $\theta$ be as in Lemma~\ref{sha-hecke-data}(2), and let $h$ be an increasing weight (Definition~\ref{sha-hecke-arcs}(1)). For an embedding $\tau:L\to\bar{\mathbf Q}_\ell$ that extends continuously to $L_{\tilde v}$ for some $v\mid\ell$ (the place $v$ is then unique), put $\lambda_{\tau,n+1-i}=h_{v,i,\tau}-(i-1)$ for $1\le i\le n$, with $h_{v,i,\tau}$ taken for the extension of $\tau$; for the other embeddings put $\lambda_\tau=(0,\dots,0)$. Then $\lambda$ is a weight in the sense of Definition~\ref{sha-hecke-ideal-tw-augmented-parts}(1). A point $x$ of $B$ is arithmetic of weight $h$ if and only if its restriction $x|_\Lambda$ to $\Lambda$ is an arithmetic point of $\Lambda$ of weight $\lambda$, and it is of distinct type of weight $h$ if and only if $x|_\Lambda$ is a point of $\Lambda$ of distinct type of weight $\lambda$.
\end{lemma}
\leanref{Development/AutomorphicArithmetic/AutomorphicVarieties/ModularityArgument/ThresholdArguments/DenseThresholdCases/TWtBridge.lean}{UpperTheory.V.Hi.arith_good_iff_hset_of_theta_incr_weight}
\begin{proof}
As $h$ is increasing, $h_{v,i+1,\tau}-i\ge h_{v,i,\tau}-(i-1)$, so $\lambda_{\tau,1}\ge\dots\ge\lambda_{\tau,n}$, and $\lambda$ is a weight.
Let $v\mid\ell$, $1\le i\le n$ and $u\in\mathcal O_{L_{\tilde v}}^\times$, and let $\sigma\in I_v$ with $\operatorname{Art}^{-1}\sigma=u$. The point $x$ carries the universal character $\kappa_{v,i}$ of $I_v$ to $\psi_{x,v,i}|_{I_v}$, so Lemma~\ref{sha-hecke-data}(2) gives $x(\theta(e_{v,i}(u)))=\psi_{x,v,i}(\sigma)^{-1}\varepsilon(\sigma)^{1-i}$, and $\varepsilon(\sigma)=\prod_\tau\tau(u)$, the product over the embeddings $\tau:L_{\tilde v}\to\bar{\mathbf Q}_\ell$ (proof of Lemma~\ref{sha-hecke-local-characters}). So the character $\chi_{x,v,i}$ of Definition~\ref{sha-hecke-ideal-tw-augmented-parts}(1) is
\[\chi_{x,v,i}(u)=\psi_{x,v,i}(\sigma)^{-1}\prod_\tau\tau(u)^{1-i-\lambda_{\tau,n+1-i}}=\Big(\psi_{x,v,i}(\sigma)\prod_\tau\tau(u)^{h_{v,i,\tau}}\Big)^{-1}.\]
Hence the only characters $\phi_{v,i}$ of $I_v$ that can satisfy the condition of Definition~\ref{sha-hecke-arcs}(1) are $\phi_{v,i}(\sigma)=\chi_{x,v,i}(\operatorname{Art}^{-1}\sigma)$. By local class field theory, $\operatorname{Art}^{-1}$ is a continuous map from $I_v$ onto $\mathcal O_{L_{\tilde v}}^\times$, so $\phi_{v,i}$ has open kernel if and only if $\chi_{x,v,i}$ is trivial on the units congruent to $1$ modulo $\varpi^e$ for some $e$; a common $e$ for the finitely many pairs $(v,i)$ is the largest of them. This gives the first equivalence. Since $\operatorname{Art}^{-1}$ is onto, the $\phi_{v,i}$, $1\le i\le n$, are pairwise distinct if and only if the $\chi_{x,v,i}$ are, which gives the second.
\end{proof}
% arith_good_iff_hset_of_theta_incr_weight (TWtBridge.lean): wt_ex for lambda (mu v tau i = h v i.rev tau - i.rev, with emb_place and emb_inj for the uniqueness of v); hf from IsTheta (h1), toPair_spec (h2) and Loc.laCyc (h3, h4: epsV is the cyclotomic character and equals the product of the tau at art sigma), h5; locAlg_iff for each (v, j) and the bound E of the finitely many e; part_val and art_onto for the distinct characters. The hypothesis LiteratureRankNHypothesis supplies the local Artin maps (laArtinLit').

\begin{lemma}\label{sha-automorphy-unramified-hecke-eigenspace-multiplicity}
Let $\mathcal O$ be a commutative ring, $\Lambda$ and $B$ commutative $\mathcal O$-algebras, $\lambda:\Lambda\to B$ an $\mathcal O$-algebra map, and $M$ a $B$-module that is finite and free over $\Lambda$ through $\lambda$. Let $K$ be an algebraically closed field with an $\mathcal O$-algebra structure, $y:B\to K$ an $\mathcal O$-algebra map that vanishes on every element of $B$ that kills $M$, and $(g_i)_{i\in I}$ elements of $B$ such that every element of $B$ acts on $M$ by an element of the $\Lambda$-algebra of endomorphisms generated by the actions of the $g_i$. View $K$ as a $\Lambda$-algebra through $y\circ\lambda$. Let $\iota:K\to C$ be an isomorphism of fields, $W$ a $C$-vector space, $\operatorname{sp}:M\to W$ an additive map with $\operatorname{sp}(am)=\iota y(\lambda(a))\operatorname{sp}(m)$ for $a\in\Lambda$ and $m\in M$, such that the images under $\operatorname{sp}$ of every basis of $M$ over $\Lambda$ are linearly independent over $C$, and $\gamma_i$, for $i\in I$, $C$-linear endomorphisms of $W$ with $\operatorname{sp}(g_im)=\gamma_i\operatorname{sp}(m)$. Let $Z$ be the space of the vectors of the span of $\operatorname{sp}(M)$ that lie, for every $i$, in the generalized eigenspace of $\gamma_i$ for $\iota y(g_i)$, and let $P\subset M\otimes_\Lambda K$ be the space of the vectors that lie, for every $b\in B$, in the generalized eigenspace of $b$ for $y(b)$.
\begin{enumerate}
\item $Z\ne0$, and $\dim_KP=\dim_CZ$.
\item Let $e$ be a $\Lambda$-linear endomorphism of $M$, $N\ge1$ an integer that is nonzero in $C$, and $E$ a $C$-linear endomorphism of $W$ with $N\operatorname{sp}(em)=E\operatorname{sp}(m)$ for $m\in M$. Then the space of the vectors of $P$ fixed by $e$ has dimension $\dim_C\big(Z\cap\ker(E-N)\big)$ over $K$.
\end{enumerate}
\end{lemma}
\leanref{Development/AutomorphicArithmetic/AutomorphicVarieties/ModularityArgument/LinkedUnramifiedChains/TeichmullerParts/SubalgebraIntegrality/FieldNontrivial.lean}{UpperTheory.V.Z.J.D.mu_eq}
\leanref{Development/AutomorphicArithmetic/AutomorphicVarieties/ModularityArgument/CuspidalDescent/GLDenseRanges/InertiaPartMaps/MuHeckeE.lean}{UpperTheory.V.Z.Desc.V0.mu_eqE}
\begin{proof}
\begin{enumerate}
\item[(i)] \emph{A model of the span.} Let $(b_r)$ be a basis of $M$ over $\Lambda$. Then the $1\otimes b_r$ form a basis of $M\otimes_\Lambda K$ over $K$, and the additive map $\Psi:M\otimes_\Lambda K\to W$, $m\otimes k\mapsto\iota(k)\operatorname{sp}(m)$, sends $\sum_rk_r(1\otimes b_r)$ to $\sum_r\iota(k_r)\operatorname{sp}(b_r)$. As the $\operatorname{sp}(b_r)$ are linearly independent over $C$ and $\iota$ is bijective, $\Psi$ is injective, with $\Psi(kv)=\iota(k)\Psi(v)$; its image is the span of the $\operatorname{sp}(b_r)$, which is the span of $\operatorname{sp}(M)$, as $\operatorname{sp}(\sum_ra_rb_r)=\sum_r\iota y(\lambda(a_r))\operatorname{sp}(b_r)$.
\item[(ii)] \emph{Generators suffice.} The space $P$ is the space of the vectors that lie, for every $i$, in the generalized eigenspace of $g_i$ for $y(g_i)$. Indeed, let $v$ be such a vector. The endomorphisms $\phi$ of $M$ for which some $b\in B$ acts by $\phi$ with $v$ in the generalized eigenspace of $b\otimes K$ for $y(b)$ include the actions of the $g_i$ and of the $\lambda(a)$, $a\in\Lambda$, which act on $M\otimes_\Lambda K$ by the scalar $y(\lambda(a))$; and they are closed under sums and products, as the actions of elements of $B$ commute, so that generalized eigenvectors for commuting operators are generalized eigenvectors of their sums and products, for the sums and products of the eigenvalues. So they include the action of every $b\in B$, by some $b'\in B$ with $v$ in the generalized eigenspace of $b'$ for $y(b')$. As $b-b'$ kills $M$, $y(b)=y(b')$ and $b$ acts as $b'$.
\item[(iii)] \emph{Transport.} For $v\in M\otimes_\Lambda K$, $\Psi(g_iv)=\gamma_i\Psi(v)$, so $\Psi((g_i-y(g_i))^kv)=(\gamma_i-\iota y(g_i))^k\Psi(v)$. By (i) and (ii), $\Psi$ maps $P$ bijectively onto $Z$, and $\dim_KP=\dim_CZ$, as $\Psi$ is semilinear for $\iota$. In (2), $N\Psi(ev)=E\Psi(v)$, so $ev=v$ if and only if $E\Psi(v)=N\Psi(v)$, as $\Psi$ is injective and $N\ne0$ in $C$; hence $\Psi$ maps the vectors of $P$ fixed by $e$ bijectively onto $Z\cap\ker(E-N)$, which gives (2).
\item[(iv)] \emph{$Z\ne0$.} As $M$ is finite over $\Lambda$, it is finite over $B$. The prime $\ker(y)$ of $B$ contains the annihilator of $M$, so it lies in the support of $M$, and $K\otimes_BM\ne0$. The surjection $\pi:M\otimes_\Lambda K\to K\otimes_BM$ satisfies $\pi(bv)=y(b)\pi(v)$ for $b\in B$. The space $M\otimes_\Lambda K$ is finite-dimensional, and the commuting operators $b$ make it the sum of their common generalized eigenspaces, as $K$ is algebraically closed. On the common generalized eigenspace for a family of eigenvalues $\chi\ne y$, some $(b-\chi(b))^k$ with $\chi(b)\ne y(b)$ kills $v$, so $(y(b)-\chi(b))^k\pi(v)=0$ and $\pi(v)=0$. As $\pi\ne0$, $P\ne0$, and $Z\ne0$ by (iii).
\end{enumerate}
\end{proof}
% mu_eq (FieldNontrivial.lean, private): (i) is exists_semi (Reformulated.B.Ra9.exists_semi); (ii) is joint_gen (ge_add, ge_mul, ge_scal, act_lam); (iii) is transport; (iv) is joint_ne (fibre_nontrivial through Module.mem_support_iff_nontrivial_residueField_tensorProduct, and ne_bot_of_quot with iSup_iInf_maxGenEigenspace_eq_top_of_commute). mu_eqE (MuHeckeE.lean, private): exists_semi, joint_gen, hPsi, hPsie, q1, q2 and rank_semi; the Lean does not assume K algebraically closed or M finite there.

\begin{lemma}\label{sha-automorphy-unramified-hecke-eigenspace-specialization}
In Situation~\ref{sha-hecke-ideal-tw-augmented-setup}, let $\lambda$ be a weight and $x$ an arithmetic point of $\Lambda$ of weight $\lambda$. For $g\in G^{(\ell)}$ let $r_g$ be right translation by $g$ on functions on $G^{(\ell)}\times G_\ell$. Then there are $\xi$ and $\operatorname{sp}$ as in Lemma~\ref{sha-hecke-ideal-tw-augmented-hida}(4) with the following property. Let $g\in G^{(\ell)}$ be such that: for every level $U=\mathcal K_0K$ and every $u\in U$, $g^{-1}ug\in U$ and $\tau_0(g^{-1}ug)=\tau_0(u)$; the same holds for every $u$ in the group of the elements of $G^{(\ell)}_0$ whose components at the places of $L$ above places outside $S'$ are integral and whose components at $w_1,\dots,w_q$ are $1$; and $g$ commutes with $GL_n(L_w)$ for every split place $w$ of $L$ above a place outside $S'$ (only the first condition is used in the proof). Then $\operatorname{sp}(r_gf)=r_g\operatorname{sp}(f)$ for every level $U$ and every $f\in\mathcal F^{\mathrm{ord}}(U)$.
\end{lemma}
\leanref{Development/AutomorphicArithmetic/AutomorphicVarieties/ModularityArgument/CuspidalDescent/GLDenseRanges/EigBoxE.lean}{UpperTheory.V.Z.Desc.V0.Spec2}
\leanref{Development/AutomorphicArithmetic/AutomorphicVarieties/ModularityArgument/CuspidalDescent/GLDenseRanges/InertiaPartMaps/AbsSpec2.lean}{UpperTheory.V.Z.Desc.V0.spec2}
\begin{proof}
Let $(\xi,\chi,c_0)$, $c$, $\phi$ and $\operatorname{sp}$ be as constructed in the proof of Lemma~\ref{sha-hecke-ideal-tw-augmented-hida}(4), which have the properties of (4). By its construction in the proof of Lemma~\ref{sha-hecke-ideal-tw-augmented-hida-control}(2), for a level $U$ and $f\in\mathcal F^{\mathrm{ord}}(U)$, $\operatorname{sp}(f)$ lies in $\mathrm{Cl}^{\mathrm{ord}}(U,c,\phi)$ and has control map $\iota x\circ f$, and every $F\in\mathrm{Cl}^{\mathrm{ord}}(U,c,\phi)$ with control map $\iota x\circ r_gf$ is $\operatorname{sp}(r_gf)$. Let $g$ be as in the statement, $U$ a level, $f\in\mathcal F^{\mathrm{ord}}(U)$ and $F=\operatorname{sp}(f)$.
\begin{enumerate}
\item \emph{$r_gF$ is classical.} Right translation by $g$ commutes with left translation by $U(n)(L^+)$ and with right translation by $G_\ell$, so $r_gF$ satisfies $r_gF(\gamma h)=\xi(\gamma_\ell)r_gF(h)$ and the conditions at $N_0$, $\bar N_c$ and $T_0$. For $u\in U$, with $u'=g^{-1}ug\in U$ and $\tau_0(u')=\tau_0(u)$, $(u\cdot r_gF)(h)=\iota\tau_0(u)F(hug)=\iota\tau_0(u')F(hgu')=F(hg)$, as $F$ is fixed by $U$. So $r_gF\in\mathrm{Cl}(U,c,\phi)$. As $r_g$ commutes with $U'_z$, a combination of right translations by $G_\ell$, it maps each generalized eigenspace of $U'_z$ into itself, and $r_gF\in\mathrm{Cl}^{\mathrm{ord}}(U,c,\phi)$.
\item \emph{The control map.} The component in $G_\ell$ of $hg$ is that of $h$, so $\operatorname{ctrl}(r_gF)=r_g\operatorname{ctrl}(F)=r_g(\iota x\circ f)=\iota x\circ r_gf$ (Definition~\ref{sha-hecke-ideal-tw-augmented-hida-weight}(3)).
\end{enumerate}
By (1) and (2), $r_gF=\operatorname{sp}(r_gf)$.
\end{proof}
% spec2 (AbsSpec2.lean) repeats the steps of spec (Lemma sha-hecke-ideal-tw-augmented-hida (4)) with the private Sp.absSpec2 in place of Sp.absSpec; absSpec2 adds to absSpec the properties hgood (sp f in ClOrd with ctrl (sp f) = mapFn f) and huniq (the classical specialization is unique), which are the construction of the proof of Lemma sha-hecke-ideal-tw-augmented-hida-control (2). Tr: hU from hN.lev, (1) is Sp.tr_clOrd (tr_cl and tr_uOp), (2) is Sp.ctrl_tr, then huniq. The hypothesis Norm H g also has the fields uq (for UQ) and sph (for the sections at Wp), which spec2 does not use.

\begin{lemma}\label{sha-automorphy-unramified-hecke-eigenspace-localization-eigenvector}
Let $\ell$ be a prime, $|\cdot|$ the absolute value of $\bar{\mathbf Q}_\ell$ and $\iota:\bar{\mathbf Q}_\ell\to\mathbf C$ an isomorphism of fields. Let $A$ be a complete local noetherian ring with finite residue field, $M$ an $A$-module, $Y\subseteq M$ a submodule that is finite and free over $A$, and $(t_i)_{i\in I}$ $A$-linear endomorphisms of $M$ that map $Y$ into itself and commute on $Y$. Let $(a_i)_{i\in I}$ be elements of $A$ and $Y_{(a)}$ the localization of $Y$ at the $t_i$ for the eigenvalues $a_i$ (Lemma~\ref{sha-hecke-ideal-tw-augmented-idempotent}), and suppose that the $t_i$ map $Y_{(a)}$ into itself and commute on it; this also follows from Lemma~\ref{sha-hecke-ideal-tw-augmented-idempotent}(3). Let $x:A\to\bar{\mathbf Q}_\ell$ be a ring map with $|x(b)|\le1$ for $b\in A$ and $|x(b)|<1$ for $b\in\mathfrak m_A$. Let $E$ be a complex vector space, $\operatorname{sp}:M\to E$ a map with $\operatorname{sp}(f+f')=\operatorname{sp}(f)+\operatorname{sp}(f')$ and $\operatorname{sp}(bf)=\iota x(b)\operatorname{sp}(f)$ for $f,f'\in Y$ and $b\in A$, such that the images under $\operatorname{sp}$ of every basis of $Y$ over $A$ are linearly independent over $\mathbf C$, and let $\gamma_i$, for $i\in I$, be $\mathbf C$-linear endomorphisms of $E$ with $\operatorname{sp}(t_if)=\gamma_i\operatorname{sp}(f)$ for $f\in Y$. Let $F\ne0$ be a vector of the span of $\operatorname{sp}(Y)$, and $\mu_i\in\mathbf C$ with $\gamma_iF=\mu_iF$ and $|\iota^{-1}(\mu_i)-x(a_i)|<1$ for every $i$. Then $Y_{(a)}\ne0$, and there is a ring map $z$ from the $A$-subalgebra of $\operatorname{End}_A(Y_{(a)})$ generated by the restrictions of the $t_i$ to $\bar{\mathbf Q}_\ell$ that equals $x$ on $A$ and has $\iota z(t_i)=\mu_i$ for every $i$.
\end{lemma}
\leanref{Development/AutomorphicArithmetic/AutomorphicVarieties/HeckeAlgebras/UpOperatorEigenvalues/EigChar.lean}{UpperTheory.V.Hi.Up.eigPt}
\leanref{Development/AutomorphicArithmetic/AutomorphicVarieties/HeckeAlgebras/UpOperatorEigenvalues/EigChar.lean}{UpperTheory.V.Hi.Up.eigChar}
\leanref{Development/AutomorphicArithmetic/AutomorphicVarieties/AlgebraTools/HidaCompleteRings/HeckeModel.lean}{UpperTheory.V.AlgebraTools.AnnihilatorQuotientFinite.heckeModel}
\leanref{Development/AutomorphicArithmetic/AutomorphicVarieties/AlgebraTools/HidaCompleteRings/HeckeRestr.lean}{UpperTheory.V.AlgebraTools.AnnihilatorQuotientFinite.heckeRestr}
\begin{proof}
Let $T\subseteq\operatorname{End}_A(Y)$ be the $A$-subalgebra generated by the restrictions of the $t_i$ to $Y$, which we also write $t_i$; it is commutative, as the $t_i$ commute on $Y$.
\begin{enumerate}
\item \emph{A character of $T$.} Let $(b_k)_{k\in K}$ be a basis of $Y$ over $A$, with $K$ finite. The linear map $L:\mathbf C^K\to E$, $w\mapsto\sum_kw_k\operatorname{sp}(b_k)$, is injective by hypothesis. For $f=\sum_kc_kb_k\in Y$, $\operatorname{sp}(f)=L\big((\iota x(c_k))_k\big)$, by the two properties of $\operatorname{sp}$. So the image of $L$ contains the span of $\operatorname{sp}(Y)$, and $F=L(w)$ for some $w\ne0$. For $u\in T$ let $\rho(u)$ be the matrix of $u$ in the basis $(b_k)$, with $\iota x$ applied to its entries; $\rho$ is a ring map from $T$ to the $K\times K$ matrices over $\mathbf C$, and $\rho(b)$ is the scalar $\iota x(b)$ for $b\in A$. For the $k$-th standard basis vector $\epsilon_k$, $L(\rho(t_i)\epsilon_k)=\operatorname{sp}(t_ib_k)=\gamma_iL(\epsilon_k)$; so $L(\rho(t_i)w)=\gamma_iF=\mu_iF=L(\mu_iw)$, and $\rho(t_i)w=\mu_iw$. The $u\in T$ for which $w$ is an eigenvector of $\rho(u)$ include the elements of $A$ and the $t_i$, and they are closed under sums and products; so they are all of $T$. Writing $\rho(u)w=\chi(u)w$ defines a ring map $\chi:T\to\mathbf C$ with $\chi(b)=\iota x(b)$ for $b\in A$ and $\chi(t_i)=\mu_i$.
\item \emph{Its values are integral.} Let $\chi'=\iota^{-1}\circ\chi$. Then $|\chi'(b)|=|x(b)|\le1$ for $b\in A$, and $\chi'(t_i)=\big(\iota^{-1}(\mu_i)-x(a_i)\big)+x(a_i)$ has $|\chi'(t_i)|\le1$. By the ultrametric inequality, the $u\in T$ with $|\chi'(u)|\le1$ are closed under sums and products; so $\chi'$ takes values in the ring of integers of $\bar{\mathbf Q}_\ell$, a local ring whose maximal ideal consists of the elements of absolute value less than $1$.
\item \emph{$Y_{(a)}\ne0$.} By Lemma~\ref{sha-hecke-ideal-tw-augmented-idempotent}(1) and (3), there is an idempotent $e\in T$ with $1-e\in\mathfrak m_AT+\sum_i(t_i-a_i)T$ and $Y_{(a)}=eY$. The map $\chi'$ carries $\mathfrak m_AT$ and the $t_i-a_i$ into the maximal ideal of the ring of integers, as $|x(b)|<1$ for $b\in\mathfrak m_A$ and $|\chi'(t_i-a_i)|=|\iota^{-1}(\mu_i)-x(a_i)|<1$. So $1-\chi'(e)$ lies in that maximal ideal, $\chi'(e)$ is a unit, and $\chi'(e)=1$, as $\chi'(e)^2=\chi'(e)$. Hence $e\ne0$, and $Y_{(a)}=eY\ne0$.
\item \emph{The map $z$.} Each $s\in T$ commutes with $e$, so it maps $Y_{(a)}=eY$ into itself, and restriction is an $A$-algebra map $r$ from $T$ onto the $A$-subalgebra $T'$ of $\operatorname{End}_A(Y_{(a)})$ generated by the restrictions of the $t_i$. If $r(s)=0$, then $se=0$, as $e$ maps $Y$ into $Y_{(a)}$, and $\chi'(s)=\chi'(s)\chi'(e)=\chi'(se)=0$. So $\chi'=z\circ r$ for a ring map $z:T'\to\bar{\mathbf Q}_\ell$, which equals $x$ on $A$ and has $\iota z(t_i)=\iota\chi'(t_i)=\mu_i$.
\end{enumerate}
\end{proof}
% eigPt (EigChar.lean): heckeModel for T (the commutative heckeAlg Y op, with t its generators); (1) is eigChar (the injective L = Fintype.linearCombination of sp (b i), hs, the matrix map rho through LinearMap.toMatrixAlgEquiv, hrhot, hrhoa and hall by Algebra.adjoin_induction); (2) is hmem; (3) and (4) are heckeRestr (locAbs for e, locChar for Y_(a) = eY, hxe for chi'(e) = 1, hb for e /= 0, hker and RingHom.liftOfSurjective for z), applied to the codomain restriction of chi' to the valuation ring.

\begin{lemma}\label{sha-automorphy-unramified-hecke-eigenspace-localization}
In Situation~\ref{sha-hecke-ideal-tw-augmented-setup}, let $\lambda$ be a weight, $x$ an arithmetic point of $\Lambda$ of weight $\lambda$, and $\xi$ and $\operatorname{sp}$ as in Lemma~\ref{sha-hecke-ideal-tw-augmented-hida}(4). Let $U=\mathcal K_0K$ be a level with $M_K\ne0$, and $y:\mathbf T_K\to\bar{\mathbf Q}_\ell$ an $\mathcal O$-algebra map equal to $x$ on $\Lambda$. Then:
\begin{enumerate}
\item the images under $\operatorname{sp}$ of every basis of $M_K$ over $\Lambda$ are linearly independent over $\mathbf C$;
\item a vector of the span of $\operatorname{sp}(\mathcal F^{\mathrm{ord}}(U))$ lies, for every operator $t$ among the $T_{w,j}$ and $U_{v,j}$, in the generalized eigenspace of $t^\xi$ for $\iota y(t)$, if and only if it lies in the span of $\operatorname{sp}(M_K)$ and has the same property.
\end{enumerate}
\end{lemma}
\leanref{Development/AutomorphicArithmetic/AutomorphicVarieties/ModularityArgument/LinkedUnramifiedChains/TeichmullerParts/UnramifiedTrivialInertialType.lean}{UpperTheory.V.Z.J.D.LocSp}
\leanref{Development/AutomorphicArithmetic/AutomorphicVarieties/ModularityArgument/LinkedUnramifiedChains/TeichmullerParts/SubalgebraIntegrality/ExistsJointEigenvectorBis.lean}{UpperTheory.V.Z.J.D.locSp}
\begin{proof}
Let $Y_0=\mathcal F^{\mathrm{ord}}(U)$, and for each operator $t$ let $a_t\in\mathcal O$ be its eigenvalue (Situation~\ref{sha-hecke-ideal-tw-augmented-setup}(4)). By Lemma~\ref{sha-hecke-ideal-tw-augmented-hida}(2), $Y_0$ and $M_K$ are finite and free over $\Lambda$, and the operators map them into themselves and commute on them. Let $e$ be the idempotent of Lemma~\ref{sha-hecke-ideal-tw-augmented-idempotent}(1) for $Y_0$, the operators and the $a_t$, so that $M_K=eY_0$, and let $Z'=(1-e)Y_0$.
\begin{enumerate}
\item[(i)] \emph{Congruences.} We show $|y(t)-x(a_t)|<1$ for every operator $t$. Let $D$ be the action of $t-a_t$ on $M_K$. For $m\in M_K$, some power of $t-a_t$ maps $m$ into $\mathfrak m_\Lambda Y_0$, by Lemma~\ref{sha-hecke-ideal-tw-augmented-idempotent}(2), and then, applying $e$, into $\mathfrak m_\Lambda M_K$. As $M_K$ is finite over $\Lambda$, a single power $D^N$ maps $M_K$ into $\mathfrak m_\Lambda M_K$. By the Cayley-Hamilton theorem \cite[Theorem~2.1]{Mat86}, $D^N$ is a root of a monic polynomial $p$ of some degree $d$ whose coefficient of $X^k$ lies in $\mathfrak m_\Lambda^{d-k}$. Applying $y$, $u=(y(t)-x(a_t))^N$ is a root of the polynomial with coefficients $x(p_k)$, which have absolute value less than $1$ for $k<d$ (Lemma~\ref{sha-automorphy-unramified-near-tangent-points}). If $|u|\ge1$, then $|\sum_{k<d}x(p_k)u^k|<|u|^d$, which is impossible. So $|u|<1$.
\item[(ii)] \emph{A complement.} As $e$ lies in the algebra generated by the operators, $Y_0=M_K\oplus Z'$ with both summands stable under the operators. A direct summand of the free module $Y_0$ over the local ring $\Lambda$ is projective, hence free \cite[Theorem~2.5]{Mat86}. A basis of $M_K$ and a basis of $Z'$ together form a basis of $Y_0$, so by Lemma~\ref{sha-hecke-ideal-tw-augmented-hida}(4) their images under $\operatorname{sp}$ are linearly independent over $\mathbf C$. This gives (1). It also shows that the span $E_1$ of $\operatorname{sp}(M_K)$ and the span $E_2$ of $\operatorname{sp}(Z')$, which are the spans of the images of these bases, meet in $0$, and the span of $\operatorname{sp}(Y_0)$ is $E_1+E_2$. Each $t^\xi$ maps $E_1$ and $E_2$ into themselves, as $\operatorname{sp}(tm)=t^\xi\operatorname{sp}(m)$, and the $t^\xi$ commute on $E_2$.
\item[(iii)] \emph{Components.} Let $F=F_1+F_2$, with $F_1\in E_1$ and $F_2\in E_2$, lie in the generalized eigenspace of every $t^\xi$ for $\iota y(t)$. The component in $E_2$ of $(t^\xi-\iota y(t))^kF$ is $(t^\xi-\iota y(t))^kF_2$, so $F_2$ lies in the same generalized eigenspaces.
\item[(iv)] \emph{$F_2=0$.} Suppose $F_2\ne0$. The commuting operators $t^\xi$ on the nonzero finite-dimensional space of the vectors of $E_2$ that lie in these generalized eigenspaces have a common eigenvector $G\ne0$, with $t^\xi G=\iota y(t)G$, by induction on the dimension. The eigenvalues $\iota y(t)$ satisfy the congruences of (i).
So Lemma~\ref{sha-automorphy-unramified-hecke-eigenspace-localization-eigenvector}, applied to $A=\Lambda$, $Y=Z'$, the operators $t$ and the $a_t$, $x$, $\operatorname{sp}$, the $t^\xi$, $F=G$ and $\mu_t=\iota y(t)$, shows that the localization of $Z'$ at the operators for the eigenvalues $a_t$ is nonzero. Its hypotheses hold: $Z'$ is finite and free over $\Lambda$, stable under the operators, which commute on it, by (ii); its localization is stable under them by Lemma~\ref{sha-hecke-ideal-tw-augmented-idempotent}(3); $|x|\le1$ on $\Lambda$ and $|x|<1$ on $\mathfrak m_\Lambda$ by Lemma~\ref{sha-automorphy-unramified-near-tangent-points}; $\operatorname{sp}$ is additive on $Z'$, with $\operatorname{sp}(bm)=\iota x(b)\operatorname{sp}(m)$ for $b\in\Lambda$ and $\operatorname{sp}(tm)=t^\xi\operatorname{sp}(m)$, by Lemma~\ref{sha-hecke-ideal-tw-augmented-hida}(4); a basis of $Z'$ and a basis of $M_K$ together form a basis of $Y_0$, so the images of every basis of $Z'$ are linearly independent over $\mathbf C$, as in (ii); and $|y(t)-x(a_t)|<1$ by (i).
The localization of $Z'$ at the operators for the eigenvalues $a_t$ lies in $Z'$ and in the localization $M_K$ of $Y_0$, which meet in $0$. So $F_2=0$, and $F=F_1\in E_1$.
\end{enumerate}
The span of $\operatorname{sp}(M_K)$ lies in that of $\operatorname{sp}(Y_0)$, which gives the converse in (2).
\end{proof}
% locSp (ExistsJointEigenvectorBis.lean): Up.normLt and hx1 for the absolute values of x; (i) is the private abs_cong (abs_nil, LinearMap.exists_monic_and_coeff_mem_pow_and_aeval_eq_zero_of_range_le_smul and root_lt); (ii) to (iv) are the private abs_main: AnnihilatorQuotientFinite.locFree for the complement Z (the book uses the idempotent of Lemma sha-hecke-ideal-tw-augmented-idempotent for it), free_of_flat_of_isLocalRing, hboth, sp_span, hdisE, hst and hcm; joint_vec with exists_joint_eigenvector_bis for G; Up.eigPt and locPart_sub for the contradiction. Finally hsp.span identifies the span of sp(Ord) with OrdAt.

\begin{lemma}\label{sha-automorphy-unramified-hecke-eigenspace-galois}
In Situation~\ref{sha-hecke-ideal-tw-augmented-setup}, let $\lambda$ be a weight, $x$ an arithmetic point of $\Lambda$ of weight $\lambda$, and $\xi$ and $\operatorname{sp}$ as in Lemma~\ref{sha-hecke-ideal-tw-augmented-hida}(4). Let $U=\mathcal K_0K$ be a level with $M_K\ne0$, $y:\mathbf T_K\to\bar{\mathbf Q}_\ell$ an $\mathcal O$-algebra map equal to $x$ on $\Lambda$, and $Z$ the space of the vectors of the span of $\operatorname{sp}(\mathcal F^{\mathrm{ord}}(U))$ that lie, for every operator $t$ among the $T_{w,j}$ and $U_{v,j}$, in the generalized eigenspace of $t^\xi$ for $\iota y(t)$. Let $(D,\psi)$ be attached to the elements $T_{w,j}$ and $U_{v,j}$ of $\mathbf T_K$, for $S''=S'$ (Definition~\ref{sha-hecke-ideal-tw-augmented-parts}(3)), and let $\rho:\Gamma'_0\to GL_n(\bar{\mathbf Q}_\ell)$ be a continuous irreducible representation whose determinant is $y\circ D$. Then for every automorphic part $Y$ of weight $\xi$ with $Y\cap Z\ne0$ there is $P\in GL_n(\bar{\mathbf Q}_\ell)$ with $\rho_Y(h)=P\rho(h)P^{-1}$ for all $h\in\Gamma'_0$.
\end{lemma}
\leanref{Development/AutomorphicArithmetic/AutomorphicVarieties/ModularityArgument/ThresholdArguments/DenseThresholdCases/PartRep.lean}{UpperTheory.V.Z.J.D.PartRepAlt}
\leanref{Development/AutomorphicArithmetic/AutomorphicVarieties/ModularityArgument/ThresholdArguments/DenseThresholdCases/TPartRep.lean}{UpperTheory.V.Z.J.D.partRep_s}
\begin{proof}
Let $F\ne0$ lie in $Y\cap Z$, and let $w$ be a split place of $L$ above a place outside $S'$.
\begin{enumerate}
\item \emph{Eigenvalues.} For $1\le j\le n$, $T^\xi_{w,j}$ acts on $Y$ by $a_{Y,w,j}$, and $F$ lies in its generalized eigenspace for $\iota y(T_{w,j})$. So $(a_{Y,w,j}-\iota y(T_{w,j}))^kF=(T^\xi_{w,j}-\iota y(T_{w,j}))^kF=0$ for some $k$, and $a_{Y,w,j}=\iota y(T_{w,j})$; for $j=0$ both are $1$.
\item \emph{Determinants.} By (1), the characteristic polynomial of $\rho_Y$ at the inverse of a Frobenius element at $w$, carried by $\iota$, is the Hecke polynomial with coefficients $\iota y(T_{w,j})$. As $(D,\psi)$ is attached, this is also the characteristic polynomial of $y\circ D$, that is, of $\rho$, at the same element. These elements are dense in $\Gamma'_0$, by the Chebotarev density theorem (Theorem~\ref{lit-chebotarev-degree-one}), and the coefficients of the characteristic polynomials of $\rho_Y$ and $\rho$ are continuous; so $\rho_Y$ and $\rho$ have the same characteristic polynomials everywhere.
\item \emph{Conjugacy.} As $\rho_Y$ is semisimple and $\rho$ is irreducible, they are conjugate, by the Brauer-Nesbitt theorem \cite[(30.16)]{CR62}.
\end{enumerate}
\end{proof}
% partRep_s (TPartRep.lean): hgen; hval with Y.sph and val_eq for (1), genW and Y.a_zero for j = 0; AlgebraTools.Det.extDense with hH.dense (the field dense of HSet.Ax), Det.repCont, hU.frob, Y.frob and pol_map for (2); AlgebraTools.Det.bnRep with Y.ss and ss_irr for (3).

\begin{lemma}\label{sha-automorphy-unramified-hecke-eigenspace}
In Situation~\ref{sha-hecke-ideal-tw-augmented-setup}, suppose that $S'=S$ and $M\ne0$, let $U=\mathcal K_0$, let $f:B\to\mathbf T$ be the surjection of the proof of Proposition~\ref{sha-hecke-ideal}, through which $B$ acts on $M$, and let $K_0$, $\rho_0$, $q$, $a$, $j$ and $e$ be as in Lemma~\ref{sha-automorphy-unramified-hecke-projector}. Let $y:\mathbf T\to\bar{\mathbf Q}_\ell$ be an $\mathcal O$-algebra map such that $x=y\circ f$ restricts to an arithmetic point of $\Lambda$ of a weight $\lambda$ (Definition~\ref{sha-hecke-ideal-tw-augmented-parts}(1)) and $x(D^{\mathrm u})$ is the determinant of a continuous irreducible representation of $\operatorname{Gal}(\bar L/L)$. Write $\mu_e(x)$ for the dimension over $\bar{\mathbf Q}_\ell$ of the space of vectors of $M\otimes_\Lambda\bar{\mathbf Q}_\ell$ fixed by $e$ that lie, for every $b\in B$, in the generalized eigenspace of $b$ for $x(b)$, as in Definition~\ref{sha-automorphy-unramified-multiplicity}(2). Then:
\begin{enumerate}
\item $|y(U_{v,i})|=1$ for all $v\mid\ell$ and $1\le i\le n$;
\item there are $\xi$ and $\operatorname{sp}$ as in Lemma~\ref{sha-hecke-ideal-tw-augmented-hida}(4) for $x|_\Lambda$, with moreover $\operatorname{sp}(\rho_0(g)f')=\rho_0(g)\operatorname{sp}(f')$ for $g\in K_0$ and $f'\in\mathcal F^{\mathrm{ord}}(U)$, with the following properties. Let $Z$ be the space of the vectors of the span of $\operatorname{sp}(\mathcal F^{\mathrm{ord}}(U))$ that lie in the generalized eigenspace of $t^\xi$ for $\iota y(t)$, for every operator $t$ among the $T_{w,i}$ and $U_{v,i}$, and let $E^\xi=\sum_ga(g)\rho_0(g)$ on $V^\xi$. Then $Z\ne0$; $\rho_Y$ is irreducible for every automorphic part $Y$ of weight $\xi$ with $Y\cap Z\ne0$; and $\mu_e(x)=\dim_{\mathbf C}\big(Z\cap\ker(E^\xi-q^j)\big)$.
\end{enumerate}
\end{lemma}
\leanref{Development/AutomorphicArithmetic/AutomorphicVarieties/ModularityArgument/LinkedUnramifiedChains/TeichmullerParts/SubalgebraIntegrality/UnitEv.lean}{UpperTheory.V.Z.J.D.unitEv}
\leanref{Development/AutomorphicArithmetic/AutomorphicVarieties/ModularityArgument/ThresholdArguments/DenseThresholdCases/ZEqDimE.lean}{UpperTheory.V.Z.Desc.V0.IsDimEAlt}
\leanref{Development/AutomorphicArithmetic/AutomorphicVarieties/ModularityArgument/CuspidalDescent/GLDenseRanges/EigBoxE.lean}{UpperTheory.V.Z.Desc.V0.Spec2}
\leanref{Development/AutomorphicArithmetic/AutomorphicVarieties/ModularityArgument/LinkedUnramifiedChains/TeichmullerParts/HodgeTateFromSandwich.lean}{UpperTheory.V.Z.J.D.MuZ}
\leanref{Development/AutomorphicArithmetic/AutomorphicVarieties/ModularityArgument/CuspidalDescent/GLDenseRanges/EigBoxE.lean}{UpperTheory.V.Z.Desc.V0.MuZE}
\leanref{Development/AutomorphicArithmetic/AutomorphicVarieties/ModularityArgument/ThresholdArguments/DenseThresholdCases/PartRep.lean}{UpperTheory.V.Z.J.D.PartRepAlt}
\begin{proof}
(1) The operator $U_{v,i}$ is a unit of $\mathbf T$ (Lemma~\ref{sha-hecke-algebra}), so it suffices to show $|y(t)|\le1$ for every $t\in\mathbf T$: then $|y(U_{v,i})|\le1$ and $|y(U_{v,i})^{-1}|\le1$. By Lemma~\ref{sha-automorphy-unramified-near-tangent-points}, applied to $\Lambda$, $|y(b)|\le1$ for $b\in\Lambda$. Every $t\in\mathbf T$ is integral over $\Lambda$, as $\mathbf T$ is finite over $\Lambda$; so $y(t)$ is a root of a monic polynomial whose coefficients have absolute value at most $1$, and $|y(t)|\le1$, as the absolute value is nonarchimedean.

(2) Let $h$ be the weight with $h_{v,i,\tau}=\lambda_{\tau,n+1-i}+(i-1)$, for the embeddings $\tau$ of $L_{\tilde v}$ and their restrictions to $L$; it is increasing, as $\lambda_{\tau,1}\ge\dots\ge\lambda_{\tau,n}$. The weight that Lemma~\ref{sha-automorphy-unramified-hecke-weights} attaches to $h$ agrees with $\lambda$ at the embeddings of $L$ that extend to some $L_{\tilde v}$, and only these enter Definition~\ref{sha-hecke-ideal-tw-augmented-parts}(1); so by that lemma $x$ is arithmetic of weight $h$. The map $f$ carries $\psi^{\mathrm u}$ to the characters $\psi_{\mathbf T}$, so $y$ is arithmetic of weight $h$ for the characters $y\circ\psi_{\mathbf T,v,i}$. Let $\rho$ be the continuous semisimple representation of $\Gamma_0$ given by Lemma~\ref{sha-hecke-determinant}(3) at $y$, so that $x(D^{\mathrm u})=y(D_{\mathbf T})$ is its determinant. Then $\rho$ is irreducible. Indeed, let $\rho'$ be a continuous irreducible representation of $\operatorname{Gal}(\bar L/L)$ with determinant $x(D^{\mathrm u})$. The composite of $\rho$ with the surjection $\operatorname{Gal}(\bar L/L)\to\Gamma_0$ is semisimple with the same determinant as $\rho'$, so it is isomorphic to $\rho'$, as semisimple representations over an algebraically closed field are determined by their determinants \cite{Che14}; hence it is irreducible, and so is $\rho$.
The rest follows from four facts about Hida's forms, for $x$, $y$ and $\rho$. Here $M=M_{()}$ and $\mathbf T=\mathbf T_{()}$, and the only level is $U=\mathcal K_0$.
\begin{enumerate}
\item[(a)] \emph{The specialization.} Let $\xi$ and $\operatorname{sp}$ be as in Lemma~\ref{sha-automorphy-unramified-hecke-eigenspace-specialization} for $x|_\Lambda$; they are as in Lemma~\ref{sha-hecke-ideal-tw-augmented-hida}(4). Every $g\in K_0$, viewed in $G^{(\ell)}$ through $\tilde v_0$, satisfies the conditions of that lemma: $g^{-1}ug\in U$ and $\tau_0(g^{-1}ug)=\tau_0(u)$ for $u\in U$, by step (1) of the proof of Lemma~\ref{sha-automorphy-unramified-hecke-projector}; for $S'=S$ the group of the second condition is $U$ itself; and $g$ commutes with $GL_n(L_w)$ for the split places $w$ above places outside $S$, which act at other places. So $\operatorname{sp}(\rho_0(g)f')=\rho_0(g)\operatorname{sp}(f')$ for $g\in K_0$ and $f'\in\mathcal F^{\mathrm{ord}}(U)$.
\item[(b)] \emph{$Z\ne0$.} For each operator $t$ among the $T_{w,i}$ and $U_{v,i}$ choose $g_t\in B$ with $f(g_t)=t$, as $f$ is onto. Apply Lemma~\ref{sha-automorphy-unramified-hecke-eigenspace-multiplicity}(1) to $\Lambda\to B$, $M$, $K=\bar{\mathbf Q}_\ell$, the map $x=y\circ f$, the elements $g_t$, $\iota$, $W=V^\xi$, $\operatorname{sp}$ on $M$ and the endomorphisms $t^\xi$. Its hypotheses hold: an element $b$ of $B$ that kills $M$ has $f(b)=0$, as $\mathbf T$ acts faithfully on $M$, so $x(b)=0$; every $b\in B$ acts by $f(b)$, which lies in the $\Lambda$-algebra generated by the operators $t=f(g_t)$; $\operatorname{sp}$ is additive on $M$, with $\operatorname{sp}(am)=\iota x(a)\operatorname{sp}(m)$ for $a\in\Lambda$ and $\operatorname{sp}(tm)=t^\xi\operatorname{sp}(m)$, by Lemma~\ref{sha-hecke-ideal-tw-augmented-hida}(4); and the images of every basis of $M$ are linearly independent over $\mathbf C$, by Lemma~\ref{sha-automorphy-unramified-hecke-eigenspace-localization}(1). By Lemma~\ref{sha-automorphy-unramified-hecke-eigenspace-localization}(2), the space $Z$ of that lemma, formed with the span of $\operatorname{sp}(M)$, is the space $Z$ of the statement. So $Z\ne0$.
\item[(c)] \emph{The multiplicity.} Apply Lemma~\ref{sha-automorphy-unramified-hecke-eigenspace-multiplicity}(2), with the data of (b), $e$, $N=q^j$ and $E=E^\xi$. For $m\in M$, $q^jem=Em=\sum_ga(g)\rho_0(g)m$, and the $\rho_0(g)m$ lie in $\mathcal F^{\mathrm{ord}}(U)$, by step (2) of the proof of Lemma~\ref{sha-automorphy-unramified-hecke-projector}. As $\operatorname{sp}$ is additive on $\mathcal F^{\mathrm{ord}}(U)$ and $\operatorname{sp}(bf')=\iota x(b)\operatorname{sp}(f')$ for $b\in\mathbf Z\subset\Lambda$, (a) gives $q^j\operatorname{sp}(em)=\operatorname{sp}(Em)=\sum_ga(g)\rho_0(g)\operatorname{sp}(m)=E^\xi\operatorname{sp}(m)$. So $\mu_e(x)=\dim_{\mathbf C}\big(Z\cap\ker(E^\xi-q^j)\big)$.
\item[(d)] \emph{Galois representations.} By Lemma~\ref{sha-hecke-determinant}(1) and (2), $(D_{\mathbf T},\psi_{\mathbf T})$ is attached to the operators $T_{w,i}$ and $U_{v,i}$ of $\mathbf T$, in the sense of Definition~\ref{sha-hecke-ideal-tw-augmented-parts}(3) for $S''=S$, and $y\circ D_{\mathbf T}=x(D^{\mathrm u})$ is the determinant of the irreducible $\rho$. By Lemma~\ref{sha-automorphy-unramified-hecke-eigenspace-galois}, every automorphic part $Y$ of weight $\xi$ with $Y\cap Z\ne0$ has $\rho_Y$ conjugate to $\rho$, hence irreducible.
\end{enumerate}
Together, (a) to (d) give (2).
% (a) is Spec2 (spec2, AbsSpec2.lean; IsSpec and Tr) with normS for the translations at v0; (b) is the first half of MuZ (muZ locSp: the private mu_hecke, which is mu_eq for Y = MK); (c) is MuZE (muZE locSp: sp_opE, ord_tr and Tr for hE, then the private mu_heckeE, which is mu_eqE); (d) is PartRepAlt (partRep_s). Up.eigPt, inside Lemma sha-automorphy-unramified-hecke-eigenspace-localization, is Lemma sha-automorphy-unramified-hecke-eigenspace-localization-eigenvector.
\end{proof}
% (1) is unitEv (UnitEv.lean): hU.psi_unif makes U_{v,i} the value of a unit, norm_unit, and norm_le_of_isIntegral with isIntegral_of_subalgebra and, on Lambda, res_surj, Up.normLt and Up.intIff. (2) is the private absDimE of TNumJE.lean: factor for y, out.pt (core_hecke_of_automorphic_hecke_dense_weight, Lemma sha-hecke-determinant (3)) for rho, the private irrOf and iso_of_det (AlgebraTools.Det.bnRep) for its irreducibility, then Spec2, MuZ, MuZE, UnitEv and PartRepAlt with irr_iso; its conclusion is IsDimEAlt, with ev the values iota y(t).

\begin{lemma}\label{sha-automorphy-unramified-hecke-count-inputs}
Assume the Galois representations as in Theorem~\ref{sha-residually-reducible-automorphy}. In Situation~\ref{sha-hecke-ideal-tw-augmented-setup}, suppose that $S'=S$ and $n\le4$, let $U=\mathcal K_0$, and let $K_0$, $\rho_0$ and $E^\xi_b$ be as in Lemma~\ref{sha-automorphy-unramified-hecke-count}. Let $\lambda$ be a weight, $x$ an arithmetic point of $\Lambda$ of weight $\lambda$, $\xi$ and $\operatorname{sp}$ as in Lemma~\ref{sha-hecke-ideal-tw-augmented-hida}(4) for $x$, and $(\epsilon_t)$ and $Z$ as in Lemma~\ref{sha-automorphy-unramified-hecke-count}. For an automorphic representation $\sigma$ of $U(n)$ let $Y(\sigma)$ be the space of the functions $F\in V^\xi$ with $F(\gamma g)=\xi(\gamma_\ell)F(g)$ for $\gamma\in U(n)(L^+)$ that is spanned by those whose coordinates, made scalar through the identification of $\mathbf C^d$ with $\bigotimes\nabla(\lambda_\tau)$, are forms in $\sigma$. Let $Z_\ell$ be the space of the vectors of $V^\xi$ that satisfy the conditions defining $Z$ at the places above $\ell$ only: invariance under the level at the $\tilde v$, $v\mid\ell$, and the generalized eigenspaces of the $U^\xi_{v,i}$ for the $\epsilon_{U_{v,i}}$.
\begin{enumerate}
\item There is a family $\lambda'$ of dominant weights at the infinite places of $L$ such that an automorphic representation of $U(n)$ has weight $\lambda'_v$ at every infinite place $v$ if and only if it has, at the embedding $\iota\circ\tau$ for each embedding $\tau$ of $L$ that extends to some $L_{\tilde v}$ with $v\mid\ell$, the weight dual to $\lambda_\tau$; call such representations \emph{of weight $\lambda'$}.
\item There is a finite family $(\sigma_i)_{i\in J}$ of automorphic representations of $U(n)$ of weight $\lambda'$, unramified at the places of $L$ above places outside $S$, such that the $Y(\sigma_i)$ are independent and every function in $V^\xi$ of the kind defining $Y(\sigma)$ that is fixed by the level of $\operatorname{sp}$ lies in their sum.
\item For such a family, every finitely supported $b:K_0\to\mathbf Z$ and every $N\ge0$, $\dim\big(Z\cap\ker(E^\xi_b-N)\big)=\sum_i\dim\big(Y(\sigma_i)\cap Z\cap\ker(E^\xi_b-N)\big)$.
\item If $Y(\sigma)\ne0$, then $Y(\sigma)$ is an automorphic part of weight $\xi$ (Definition~\ref{sha-hecke-ideal-tw-augmented-parts}(2)), whose Galois representation is compatible, at almost all split places $w$, with the local components of $\sigma$ at $w$.
\item If $Y(\sigma)\ne0$ and the Galois representation of that automorphic part is irreducible, there is a regular algebraic cuspidal automorphic representation $\Pi$ of $GL_n(\mathbf A_L)$ with $\Pi^\vee\cong\Pi^c$ that is a weak base change of $\sigma$, with the infinitesimal characters given by $\lambda'$ and Satake parameters at the places where $\sigma$ is unramified.
\item Let $\sigma$ and $\sigma'$ be automorphic representations of $U(n)$ with $Y(\sigma)\cap Z\ne0$. (a) If $Y(\sigma')\cap Z\ne0$, every weak base change of $\sigma$ is one of $\sigma'$. (b) If $\sigma$ and $\sigma'$ have the same local components at the split places, then $Y(\sigma')\cap Z=Y(\sigma')\cap Z_\ell$. (c) If $\sigma'$ is of weight $\lambda'$, unramified at the places of $L$ above places outside $S$, and has local components $\tau_w$ at the split places $w$, then $\dim\big(Y(\sigma')\cap Z_\ell\cap\ker(E^\xi_b-N)\big)=\dim\mathcal E\cdot\dim I(b,N)$, where $\mathcal E$ is the space of the vectors of $\bigotimes_{v\mid\ell}\tau_{\tilde v}$ cut out by the conditions defining $Z_\ell$, and $I(b,N)$ is formed as in Lemma~\ref{sha-automorphy-unramified-hecke-count} with $\sigma_u=\tau_{\tilde u}$ and $\sigma_0=\tau_{\tilde v_0}$. (d) In (c), $\dim\mathcal E\le1$.
\item If $\Pi$ is a regular algebraic cuspidal automorphic representation of $GL_n(\mathbf A_L)$ with $\Pi^\vee\cong\Pi^c$ that is a weak base change of an automorphic representation $\sigma$ of $U(n)$ that is unramified at the places that do not split in $L$, then at every split place $w$ the local components of $\Pi$ at $w$ are those of $\sigma$.
\item For every weight $\lambda'$ of $U(n)$ there is an integer $a_0\ge1$ such that, for every $\Pi$ as in (7) with the infinitesimal characters given by $\lambda'$ and Satake parameters at the places that do not split, and every family $(\sigma_i)$ as in (2), the number of the $i$ such that $\Pi$ is a weak base change of $\sigma_i$ is $a_0$.
\end{enumerate}
\end{lemma}
\leanref{Development/AutomorphicArithmetic/AutomorphicVarieties/ModularityArgument/LinkedUnramifiedChains/TeichmullerParts/UnramifiedTrivialInertialType.lean}{UpperTheory.V.Z.J.D.WtEx}
\leanref{Development/AutomorphicArithmetic/AutomorphicVarieties/ModularityArgument/LinkedUnramifiedChains/TeichmullerParts/UnramifiedTrivialInertialType.lean}{UpperTheory.V.Z.J.D.Decomp}
\leanref{Development/AutomorphicArithmetic/AutomorphicVarieties/ModularityArgument/CuspidalDescent/GLDenseRanges/InertiaPartMaps/ZeInfRep.lean}{UpperTheory.V.Z.Desc.V0.zDecE}
\leanref{Development/AutomorphicArithmetic/AutomorphicVarieties/ModularityArgument/ThresholdArguments/DenseThresholdCases/PartU.lean}{UpperTheory.V.Z.Desc.P.partU_alt}
\leanref{Development/AutomorphicArithmetic/AutomorphicVarieties/ModularityArgument/ThresholdArguments/DenseThresholdCases/BcOf.lean}{UpperTheory.V.Z.Desc.P.dense_threshold_bc_of_prop}
\leanref{Development/AutomorphicArithmetic/AutomorphicVarieties/ModularityArgument/ThresholdArguments/DenseThresholdCases/TAtSame.lean}{UpperTheory.V.Z.J.D.atSame_alt}
\leanref{Development/AutomorphicArithmetic/AutomorphicVarieties/ModularityArgument/ThresholdArguments/DenseThresholdCases/TSphSame.lean}{UpperTheory.V.Z.J.D.spherical_same_of_z_eq}
\leanref{Development/AutomorphicArithmetic/AutomorphicVarieties/ModularityArgument/CuspidalDescent/GLDenseRanges/EigBoxE.lean}{UpperTheory.V.Z.Desc.V0.YDimE}
\leanref{Development/AutomorphicArithmetic/AutomorphicVarieties/ModularityArgument/LinkedUnramifiedChains/TeichmullerParts/UnramifiedTrivialInertialType.lean}{UpperTheory.V.Z.J.D.EigPi}
\leanref{Development/AutomorphicArithmetic/AutomorphicVarieties/ModularityArgument/MultiplicityOne/CompOfBis.lean}{UpperTheory.V.Z.Mult.CompOf}
\leanref{Development/AutomorphicArithmetic/AutomorphicVarieties/ModularityArgument/MultiplicityOne/CompOfBis.lean}{UpperTheory.V.Z.Mult.Mult}
\begin{proof}
\openstep{The proofs of (1) to (8) are not yet written out: the weight of $U(n)$; the decomposition of the space of forms of type $\tau_0$ and weight $\xi$ into automorphic representations of $U(n)$ and of the dimensions of the eigenspaces along it; the automorphic part and the base change of a representation that meets $Z$; the comparison of two representations that meet $Z$; the dimension of the part of a representation in $Z_\ell$ as a product of a local eigenspace above $\ell$, of dimension at most one, and the space $I(b,N)$; and the local components and the count of the representations with a given base change.}
% (1) is WtEx, proved as V.J.wtEx (ScalarIrreducibility/PlaceFree.lean); (2) is Decomp, proved as decomp indepS (GlobalRationalGL/IntegralConjugation/SplitCore.lean); (3) is ZDecE, proved as zDecE zEq normS (InertiaPartMaps/ZeInfRep.lean, 212 lines); (4) is Z.Desc.P.partU_alt (PartU.lean, 156 lines, through partX_alt); (5) is BcOfAlt, proved as Z.Desc.P.dense_threshold_bc_of_prop (BcOf.lean) from HypothesesH'; (6)(a) is AtSameAlt (atSame_alt, TAtSame.lean), (6)(b) is SphSameAlt (spherical_same_of_z_eq zEq, TSphSame.lean), (6)(c) is YDimE (yDimE flathE eigBoxE secSp, InertiaPartMaps/EPartPi.lean and EigBoxE.lean, 214 lines), (6)(d) is Slope, Ord1Alt (dense_threshold_order_one, ExistsEig.lean) and EigPi (eigPi, RankLeOneOfSlices.lean); (7) is Z.Mult.CompOf and (8) is Z.Mult.Mult, hypotheses of countE'_alt that the Lean proves as compOf_free and mult_free (MultiplicityOne/CompOfFree.lean), by comparisons of trace formulas of several hundred thousand lines. Together, with the Hida infrastructure shared with Lemma sha-automorphy-unramified-hecke-eigenspace, about 5,700 lines besides (7) and (8).
% The statements are given in words; their exact forms are: Y(sigma) is Hi.CaseU.YOf (Big meet the span of the F whose toForm coordinates lie in sigma), Z_l is zL, the family of (2) is IsDec S lam' K' sigma with a section S and a level K' equal to KL outside S, iSupIndep of the YOf and the inclusion of Big meet the I1-fixed vectors in their sum; (6)(c) is the rank of YOf meet eigL meet eC equal to rank eigL of C.repL times rank cInv; weak base change is At.
\end{proof}

\begin{lemma}\label{sha-automorphy-unramified-hecke-count}
In Situation~\ref{sha-hecke-ideal-tw-augmented-setup}, suppose that $S'=S$ and $n\le4$, let $U=\mathcal K_0$, let $K_0$, $K_1$ and $\rho_0$ be as in Lemma~\ref{sha-automorphy-unramified-hecke-projector}, and for a finitely supported $b:K_0\to\mathbf Z$ let $E^\xi_b=\sum_gb(g)\rho_0(g)$ on $V^\xi$. Let $\lambda$ be a weight. Then there is an integer $a_0\ge0$ with the following property. Let $x$ be an arithmetic point of $\Lambda$ of weight $\lambda$, let $\xi$ and $\operatorname{sp}$ be as in Lemma~\ref{sha-hecke-ideal-tw-augmented-hida}(4) for $x$, and let $(\epsilon_t)$ be complex numbers, one for each operator $t$ among the $T_{w,i}$ and $U_{v,i}$, with $|\iota^{-1}(\epsilon_{U_{v,i}})|=1$. Let $Z$ be the space of the vectors of the span of $\operatorname{sp}(\mathcal F^{\mathrm{ord}}(U))$ that lie in the generalized eigenspace of $t^\xi$ for $\epsilon_t$, for every $t$. Suppose that $Z\ne0$ and that $\rho_Y$ is irreducible for every automorphic part $Y$ of weight $\xi$ with $Y\cap Z\ne0$. Then there are a cuspidal automorphic representation $\Pi$ of $GL_n(\mathbf A_L)$ and an automorphic representation $\sigma$ of $U(n)$ that contains a nonzero form of type $\tau_0$ and weight $\xi$, with local components $\sigma_u$ at the places $\tilde u$, $u\in S_{\mathrm a}$, and $\sigma_0$ at $\tilde v_0$, on spaces $W_u$ and $W_0$, which are local components of $\Pi$, such that for every finitely supported $b:K_0\to\mathbf Z$ and every integer $N\ge0$,
\[\dim_{\mathbf C}\big(Z\cap\ker(E^\xi_b-N)\big)=a_0\dim_{\mathbf C}I(b,N),\]
where $I(b,N)$ is the space of the $F=(F_1,\dots,F_m)\in\big(\bigotimes_uW_u\otimes W_0\big)^m$ such that $\sum_k\iota\tau_0(g)_{ik}\big(\bigotimes_u\sigma_u(g_u)\otimes1\big)F_k=F_i$ for $g=(g_u)\in K_{\mathrm a}$ and $1\le i\le m$, and such that each $F_i$ is fixed by $1\otimes\sigma_0(K_1)$ and satisfies $\big(1\otimes\sum_gb(g)\sigma_0(g)\big)F_i=NF_i$.
\end{lemma}
\leanref{Development/AutomorphicArithmetic/AutomorphicVarieties/ModularityArgument/ThresholdArguments/DenseThresholdCases/ZEqDimE.lean}{UpperTheory.V.Z.Desc.V0.CountEAlt}
\leanref{Development/AutomorphicArithmetic/AutomorphicVarieties/ModularityArgument/ThresholdArguments/DenseThresholdCases/CountE.lean}{UpperTheory.V.Z.Desc.V0.countE'_alt}
\begin{proof}
Let $\lambda'$ be the family of weights given by Lemma~\ref{sha-automorphy-unramified-hecke-count-inputs}(1), and $a_0\ge1$ the integer given by part (8) of that lemma for $\lambda'$; it depends only on $\lambda'$, hence only on $\lambda$. Let $x$, $\xi$, $\operatorname{sp}$, $(\epsilon_t)$ and $Z$ be as in the statement, and for an automorphic representation $\sigma$ of $U(n)$ let $Y(\sigma)\subseteq V^\xi$ be as in Lemma~\ref{sha-automorphy-unramified-hecke-count-inputs}. Let $b$ and $N$ be as in the statement, and write $E=\ker(E^\xi_b-N)$.
\begin{enumerate}
\item[(1)] \emph{Decomposition.} Lemma~\ref{sha-automorphy-unramified-hecke-count-inputs}(2) gives a finite family $(\sigma_i)_{i\in J}$ of automorphic representations of $U(n)$, and by part (3) of that lemma, $\dim(Z\cap E)=\sum_{i\in J}\dim\big(Y(\sigma_i)\cap Z\cap E\big)$. For $b=0$ and $N=0$, $E=V^\xi$; as $Z\ne0$, some $i_0\in J$ has $Y(\sigma_{i_0})\cap Z\ne0$.
\item[(2)] \emph{The base change.} By part (4) of that lemma, $Y(\sigma_{i_0})$ is an automorphic part $Y$ of weight $\xi$; by hypothesis $\rho_Y$ is irreducible, as $Y\cap Z\ne0$. By part (5), there is a regular algebraic cuspidal automorphic representation $\Pi$ of $GL_n(\mathbf A_L)$ with $\Pi^\vee\cong\Pi^c$ that is a weak base change of $\sigma_{i_0}$, has the infinitesimal characters given by $\lambda'$, and has Satake parameters at the places where $\sigma_{i_0}$ is unramified. For every finite place $w$ of $L$ let $\tau_w$ be a local component of $\Pi$ at $w$.
\item[(3)] \emph{Local components.} By part (7), for every $i\in J$ of which $\Pi$ is a weak base change, the local component of $\sigma_i$ at every split place $w$ is $\tau_w$.
\item[(4)] \emph{Which parts meet $Z$.} By part (6)(a), every $i$ with $Y(\sigma_i)\cap Z\ne0$ has $\Pi$ as a weak base change. Let $i$ have $\Pi$ as a weak base change. By (3), $\sigma_i$ and $\sigma_{i_0}$ have the same local components at the split places, so by part (6)(b), $Y(\sigma_i)\cap Z=Y(\sigma_i)\cap Z_\ell$; and by (3) and part (6)(c), $\dim\big(Y(\sigma_i)\cap Z_\ell\cap E\big)=a\cdot\dim I(b,N)$, where $a$ is the dimension of the space $\mathcal E$ of part (6)(c) for the local components $\tau_{\tilde v}$, $v\mid\ell$, and $I(b,N)$ is formed with $\sigma_u=\tau_{\tilde u}$ and $\sigma_0=\tau_{\tilde v_0}$.
\item[(5)] \emph{$a=1$.} By part (6)(d), $a\le1$. By (4) for $i=i_0$, with $b=0$ and $N=0$, $a\cdot\dim I(0,0)=\dim\big(Y(\sigma_{i_0})\cap Z\big)\ne0$; so $a=1$.
\item[(6)] \emph{The count.} By (1), (4) and (5), $\dim(Z\cap E)$ is $\dim I(b,N)$ times the number of $i\in J$ of which $\Pi$ is a weak base change. By part (8), this number is $a_0$.
\end{enumerate}
So $\Pi$ and $\sigma=\sigma_{i_0}$, which contains a nonzero form of type $\tau_0$ and weight $\xi$ as $Y(\sigma_{i_0})\ne0$, with $\sigma_u=\tau_{\tilde u}$ and $\sigma_0=\tau_{\tilde v_0}$, which are local components of $\sigma_{i_0}$ by (3) as the places of $S$ split in $L$, have the stated property.
% The proof follows count_e_of_sph_same_at_same_bc_dims (TCountE.lean, 193 lines): lambda' is wtEx (hWt), a_0 is Mult (hM); (1) is Decomp (hDec), ZDecE (hsumE, hZD), eigenpart_zero_eq_top, ce_top and ce_ex; (2) is partU_alt, hP, BcOfAlt (hBc, with ce_irr) and local_comp_complex_of_printed_results for tau; (3) is CompOf (hCO, with hns for the unramified places); (4) is AtSameAlt (h1), SphSameAlt (h2) and YDimE (h3); (5) is ce_one (Slope, Ord1Alt, EigPi) and cardinal_factor_eq_one; (6) is ce_cnt with Mult (hcard). The output is ce_mk with hS (the places of S split, ce_S).
% countE'_alt (CountE.lean) = count_e_of_sph_same_at_same_bc_dims with wtEx, decomp indepS, zDecE (ZDecE: the dimension of Z meet eC is the sum over the parts YOf sigma_i), zEq, spherical_same_of_z_eq, atSame_alt, dense_threshold_bc_of_prop, yDimE (YDimE with flathE, eigBoxE, secSp), compsEx, slope, dense_threshold_order_one and eigPi. The count is the hypothesis Z.Mult.Mult and the local components the hypothesis Z.Mult.CompOf (MultiplicityOne/CompOfBis.lean); the Lean proves them as Z.Mult.mult_free and Z.Mult.compOf_free (CompOfFree.lean), by a comparison of trace formulas of several hundred thousand lines. Size of the rest: about 5,800 lines of Lean, with the Hida infrastructure it shares with Lemma sha-automorphy-unramified-hecke-eigenspace.
\end{proof}

\begin{lemma}\label{sha-automorphy-unramified-hecke-invariants}
Let $G$ and $H$ be groups with complex representations $\rho$ on $W$ and $\kappa$ on $W'$, let $K\subset G$ and $K'\subset H$ be subgroups, $\tau:K\to GL_m(\mathbf C)$ a homomorphism, $b:H\to\mathbf Z$ finitely supported and $N\ge0$ an integer, and write $E_b=\sum_hb(h)\kappa(h)$. Let $I_\rho$ be the space of the $F\in W^m$ with $\sum_k\tau(g)_{ik}\rho(g)F_k=F_i$ for $g\in K$ and $1\le i\le m$; let $W'_{b,N}$ be the space of the $w\in W'$ fixed by $\kappa(K')$ with $E_bw=Nw$; and let $I$ be the space of the $F\in(W\otimes W')^m$ with $\sum_k\tau(g)_{ik}(\rho(g)\otimes1)F_k=F_i$ for $g\in K$ and $1\le i\le m$, such that each $F_i$ is fixed by $1\otimes\kappa(K')$ and satisfies $(1\otimes E_b)F_i=NF_i$. Then $I\cong I_\rho\otimes W'_{b,N}$, so $\dim I=\dim I_\rho\cdot\dim W'_{b,N}$, as cardinals.
\end{lemma}
\leanref{Development/AutomorphicArithmetic/AutomorphicVarieties/ModularityArgument/CuspidalDescent/GLDenseRanges/InertiaPartMaps/InvBox.lean}{UpperTheory.V.Z.Desc.V0.invBox}
\leanref{Development/AutomorphicArithmetic/AutomorphicVarieties/ModularityArgument/CuspidalDescent/GLDenseRanges/EigBoxE.lean}{UpperTheory.V.Z.Desc.V0.InvBox}
\begin{proof}
Identify $(W\otimes W')^m$ with $W^m\otimes W'$, so that $F\otimes w$ corresponds to $(F_i\otimes w)_i$. For $g\in K$ let $T_g$ be the endomorphism $F\mapsto(\sum_k\tau(g)_{ik}\rho(g)F_k)_i$ of $W^m$; the first condition defining $I$ becomes $(T_g\otimes1)z=z$, and the others become $(1\otimes\kappa(h'))z=z$ for $h'\in K'$ and $(1\otimes(E_b-N))z=0$.
The map $\iota:I_\rho\otimes W'_{b,N}\to W^m\otimes W'$ induced by the inclusions is injective, as tensor products over a field preserve injections, and its image lies in $I$, since the conditions hold on each $F\otimes w$ with $F\in I_\rho$ and $w\in W'_{b,N}$.
Conversely, let $z\in I$. Write $z=\sum_rc_r\otimes w_r$ with $(c_r)$ part of a basis of $W^m$. The conditions in the second factor give $\kappa(h')w_r=w_r$ for $h'\in K'$ and $E_bw_r=Nw_r$, by the linear independence of the $c_r$; so $z$ lies in the image of $W^m\otimes W'_{b,N}$. Write $z=\sum_sF_s\otimes w'_s$ with $(w'_s)$ part of a basis of $W'_{b,N}$. Then $(T_g\otimes1)z=z$ gives $T_gF_s=F_s$ for every $s$ and $g\in K$, by the linear independence of the $w'_s$; so $F_s\in I_\rho$, and $z=\iota(\sum_sF_s\otimes w'_s)$. Hence $\iota$ is an isomorphism onto $I$.
\end{proof}
% invBox (InvBox.lean): Psi = TensorProduct.piLeft, T k, hV1 (inv1 through mem_invT_iff_sum_smul_inl_eq_and_inr_fixed), hV2, the injective iota (TensorProduct.map_injective_of_flat_flat), and hrange: memL_lemma_4 for the second factor, then mem_range_rTensor_subtype_of_forall_sl_coord_mem for the first; the rank is that of a tensor product.

\begin{lemma}\label{sha-automorphy-unramified-hecke-types}
Assume that Galois representations exist, in the sense of Definition~\ref{gal-galois-representations}, relative to the maps satisfying the laws of Definition~\ref{gal-llc-laws}. Let $L$ be a CM field, $n\ge1$, $\Pi$ and $\Pi'$ cuspidal automorphic representations of $GL_n(\mathbf A_L)$, $w$ a finite place of $L$ and $A$ a place of $\bar L$ above $w$, with inertia group $I_A$ and Weil group $W_A$, and let $\tau$ and $\tau'$ be local components of $\Pi$ and $\Pi'$ at $w$. Let $\chi_1,\dots,\chi_n$ be characters of $I_A$ that are restrictions of characters of $W_A$, and suppose that for every $\operatorname{rc}$ that satisfies the laws, the Weil-Deligne representations $\operatorname{rc}(\tau\nu^{(1-n)/2})$ and $\operatorname{rc}(\tau'\nu^{(1-n)/2})$ have monodromy $0$ and restrict to $\bigoplus_i\chi_i$ on $I_A$. Then the restrictions of $\tau$ and $\tau'$ to $GL_n(\mathcal O_{L_w})$ are isomorphic.
\end{lemma}
\leanref{Development/AutomorphicArithmetic/AutomorphicVarieties/ModularityArgument/ThresholdArguments/DenseThresholdCases/TypePS2.lean}{UpperTheory.V.Z.J.D.ResidualEqualityRankTwo}
\leanref{Development/AutomorphicArithmetic/AutomorphicVarieties/ModularityArgument/ThresholdArguments/DenseThresholdCases/ResEq2.lean}{UpperTheory.V.Z.J.D.resEq2_alt}
\leanref{Development/AutomorphicArithmetic/AutomorphicVarieties/ModularityArgument/ThresholdArguments/DenseThresholdCases/TTypeV.lean}{UpperTheory.V.Z.J.D.type_v_statement}
\leanref{Development/AutomorphicArithmetic/AutomorphicVarieties/ModularityArgument/LinkedUnramifiedChains/TeichmullerParts/SubalgebraIntegrality/EquivOfChar.lean}{UpperTheory.V.Z.J.D.psRes}
\begin{proof}
Let $a$ be a local Artin map for $A$ (Definition~\ref{gal-local-artin}) and fix $\operatorname{rc}$ that satisfies the laws.
\begin{enumerate}
\item \emph{Principal series.} Write $\operatorname{rc}(\tau\nu^{(1-n)/2})=(r,0)$, with $r(s)$ the diagonal matrix with entries $\chi_i(s)$ for $s\in I_A$. For $s,t\in W_A$ the commutator $[s,t]$ has degree $0$, so it lies in $I_A$, and $\chi_i([s,t])=1$ since $\chi_i$ is the restriction of a character of $W_A$; so $r([s,t])=1$, and the $r(s)$ commute. As $r$ is semisimple, it is a sum of characters $\eta_1,\dots,\eta_n$ of $W_A$, trivial on an open subgroup of $I_A$. Comparing the traces on $I_A$, which determine the characters of $I_A$ with their multiplicities, the $\eta_i$ can be numbered so that $\eta_i=\chi_i$ on $I_A$. By local class field theory, $\eta_i=\omega_i\circ a$ for smooth characters $\omega_i$ of $L_w^\times$. As $\tau$ is a local component of a cuspidal representation, it is generic \cite{Sha74}. A generic irreducible representation whose Weil-Deligne representation under $\operatorname{rc}$, after the twist by $\nu^{(1-n)/2}$, has monodromy $0$ and characteristic polynomials $\prod_i(X-\omega_i(a(s)))$ is isomorphic to $\mathrm{PS}(\omega)$, by the laws (Definition~\ref{gal-llc-laws}) and the classification of generic representations by segments \cite{BZ77,Zel80}: its segments all have length $1$, as the monodromy is $0$, and their supports are the $\omega_i\nu^{(n-1)/2}$. So $\tau\cong\mathrm{PS}(\omega)$, with $\omega_i(a(s))=\chi_i(s)$ for $s\in I_A$.
\item \emph{The units.} In the same way $\tau'\cong\mathrm{PS}(\omega')$ with $\omega'_i(a(s))=\chi_i(s)$ for $s\in I_A$. As $a$ maps $I_A$ onto $\mathcal O_{L_w}^\times$, $\omega_i=\omega'_i$ on $\mathcal O_{L_w}^\times$.
\item \emph{Restriction to $GL_n(\mathcal O_{L_w})$.} Let $\beta_i=\omega'_i\omega_i^{-1}$, an unramified character, and let $\phi$ be the function on $GL_n(L_w)$ with $\phi(pk)=\prod_i\beta_i(p_{ii})$ for $p$ upper triangular and $k\in GL_n(\mathcal O_{L_w})$; it is well defined by the Iwasawa decomposition, as the $\beta_i$ are trivial on $\mathcal O_{L_w}^\times$. Multiplication by $\phi$ carries the functions of $\mathrm{PS}(\omega)$ to those of $\mathrm{PS}(\omega')$, bijectively, as $\omega'_i=\beta_i\omega_i$, and commutes with right translation by $GL_n(\mathcal O_{L_w})$, as $\phi(gk)=\phi(g)$. So the restrictions of $\mathrm{PS}(\omega)$ and $\mathrm{PS}(\omega')$ to $GL_n(\mathcal O_{L_w})$ are isomorphic, and so are those of $\tau$ and $\tau'$.
\end{enumerate}
\end{proof}
% resEq2_alt (ResEq2.lean): TypePS2Alt (typePST'_alt = type_pst_of_type_v type_v_statement) for tau and tau', ha.surj and ha.ord_eq for the units, and psRes. type_v_statement (TTypeV.lean): hart (the Artin map is artOf h), hcomm (commutators of degree 0 in the inertia group), exists_chars for eta, sum_of_prod and exists_perm for the numbering, ha.factor for omega, then TypeVAux.nonempty_equiv_principal_series_of_char_poly_of_generic with the genericity exists_isGeneric_of_isLocComp. psRes (EquivOfChar.lean): equiv_of_char with parab, KO, exists_pw (Iwasawa), sB and the character torusChar rho, trivial on the integral points of the Borel subgroup.

\begin{lemma}\label{sha-automorphy-unramified-hecke-dimension}
In Situation~\ref{sha-hecke-ideal-tw-augmented-setup}, suppose that $S'=S$ and $n\le4$, let $U=\mathcal K_0$, let $K_0$, $K_1$, $q$, $\rho_0$, $a$ and $j$ be as in Lemma~\ref{sha-automorphy-unramified-hecke-projector}, and for $\xi$ let $E^\xi=\sum_ga(g)\rho_0(g)$ on $V^\xi$. Let $\lambda$ be a weight. For $i=1,2$ let $x_i$, $\xi_i$, $\operatorname{sp}_i$, $(\epsilon^{(i)}_t)$ and $Z_i$ be as in Lemma~\ref{sha-automorphy-unramified-hecke-count}, with $x_i$ an arithmetic point of $\Lambda$ of weight $\lambda$, such that $Z_i\ne0$ and $\rho_Y$ is irreducible for every automorphic part $Y$ of weight $\xi_i$ with $Y\cap Z_i\ne0$. Then
\[1\le\dim_{\mathbf C}\big(Z_1\cap\ker(E^{\xi_1}-q^j)\big)=\dim_{\mathbf C}\big(Z_2\cap\ker(E^{\xi_2}-q^j)\big).\]
\end{lemma}
\leanref{Development/AutomorphicArithmetic/AutomorphicVarieties/ModularityArgument/ThresholdArguments/DenseThresholdCases/ZEqDimE.lean}{UpperTheory.V.Z.Desc.V0.ZEqDimEAlt}
\leanref{Development/AutomorphicArithmetic/AutomorphicVarieties/ModularityArgument/ThresholdArguments/DenseThresholdCases/TZEqDimE.lean}{UpperTheory.V.Z.Desc.V0.z_eq_dim_e_of_count_e_res_eq_inv_loc_kir}
\begin{proof}
Let $d_0\ge1$ be as in Lemma~\ref{sha-automorphy-unramified-hecke-local} for $F=L_{\tilde v_0}$, and $a_0$ as in Lemma~\ref{sha-automorphy-unramified-hecke-count} for $\lambda$. For $i=1,2$, that lemma gives $\Pi_i$ and $\sigma_i$, with local components $\sigma_{i,u}$ on $W_{i,u}$ and $\sigma_{i,0}$ on $W_{i,0}$, such that $\dim(Z_i\cap\ker(E^{\xi_i}_b-N))=a_0\dim I_i(b,N)$ for all $b$ and $N$.
\begin{enumerate}
\item \emph{Factorization.} Let $\delta_i$ be the dimension of the space of the $F\in(\bigotimes_uW_{i,u})^m$ with $\sum_k\iota\tau_0(g)_{lk}(\bigotimes_u\sigma_{i,u}(g_u))F_k=F_l$ for $g\in K_{\mathrm a}$ and $1\le l\le m$. By Lemma~\ref{sha-automorphy-unramified-hecke-invariants}, applied to $\prod_uGL_n(L_{\tilde u})$ acting on $\bigotimes_uW_{i,u}$, $GL_n(L_{\tilde v_0})$ acting on $W_{i,0}$, $K_{\mathrm a}$, $K_1$ and $\iota\tau_0$, $\dim I_i(b,N)=\delta_i\dim W_{i,0}(b,N)$, where $W_{i,0}(b,N)$ is the space of the vectors of $W_{i,0}$ fixed by $\sigma_{i,0}(K_1)$ with $\sum_gb(g)\sigma_{i,0}(g)w=Nw$.
\item \emph{Nonvanishing.} For $b=0$ and $N=0$, $\ker(E^{\xi_i}_0-0)$ is all of $V^{\xi_i}$, so $\dim Z_i=a_0\delta_i\dim W_{i,0}^{K_1}$. The space $Z_i$ lies in the span of $\operatorname{sp}_i(\mathcal F^{\mathrm{ord}}(U))$, which is finite-dimensional (Lemma~\ref{sha-hecke-ideal-tw-augmented-hida-control}(1)), and it is nonzero. So $a_0\ne0$, $\delta_i\ne0$ and $W_{i,0}^{K_1}\ne0$.
\item \emph{The factor at $\tilde v_0$.} The representation $\sigma_{i,0}$ is a local component of the cuspidal representation $\Pi_i$, so it has a Kirillov model by Lemma~\ref{sha-automorphy-unramified-hecke-generic}; it is smooth, and it has nonzero vectors fixed by $K_1$ by (2). By Lemma~\ref{sha-automorphy-unramified-hecke-local}(2), $\dim W_{i,0}(a,q^j)=d_0$. So $\dim(Z_i\cap\ker(E^{\xi_i}-q^j))=a_0\delta_id_0\ge1$.
\item \emph{The factors at $S_{\mathrm a}$.} Let $(t_k)$ be a family of Teichm\"uller lifts of the $\bar\chi_k$, and $u\in S_{\mathrm a}$. The restrictions of the characters $\iota t_k$ to the Weil group at a place above $\tilde u$ extend their restrictions to the inertia group. As $\sigma_i$ contains a nonzero form of type $\tau_0$, the representation $\bigotimes_u\sigma_{i,u}$ contains $\iota\tau_0$, so by Lemma~\ref{sha-hecke-data}(1), for every $\operatorname{rc}$ that satisfies the laws, $\operatorname{rc}(\sigma_{i,u}\nu^{(1-n)/2})$ has monodromy $0$ and restricts to $\bigoplus_k\iota t_k$ on the inertia group. The $\sigma_{i,u}$ are local components of the cuspidal representations $\Pi_i$, so by Lemma~\ref{sha-automorphy-unramified-hecke-types}, $\sigma_{1,u}$ and $\sigma_{2,u}$ have isomorphic restrictions to $GL_n(\mathcal O_{L_{\tilde u}})$. The tensor product of these isomorphisms carries the space defining $\delta_1$ onto the space defining $\delta_2$, so $\delta_1=\delta_2$.
\end{enumerate}
By (3) and (4), the two dimensions are $a_0\delta_1d_0=a_0\delta_2d_0\ge1$.
\end{proof}
% z_eq_dim_e_of_count_e_res_eq_inv_loc_kir (TZEqDimE.lean): hLoc (Lemma sha-automorphy-unramified-hecke-local) for a, j, D_0, with card_not_dvd; hCount (Lemma sha-automorphy-unramified-hecke-count) for a_0; au_isTeich and hext for the Teichmuller characters and their extensions, htau (the type, Lemma sha-hecke-data (1)) for chi; key: InvBox (hb) for (1), h5 0 0 with ePart_zero for (2), hKir (Lemma sha-automorphy-unramified-hecke-generic) and hdim for (3), partUR2Type_alt for the types; then InvEq1 (invEq1, rank_isotypic_invariants_eq) with ResidualEqualityRankTwo (Lemma sha-automorphy-unramified-hecke-types) for (4).

\begin{lemma}\label{sha-automorphy-unramified-hecke-bound}
Assume that Galois representations exist, in the sense of Definition~\ref{gal-galois-representations}, relative to the maps satisfying the laws of Definition~\ref{gal-llc-laws}. In Situation~\ref{sha-hecke-situation}, let $n\le4$ and $\iota:\bar{\mathbf Q}_\ell\cong\mathbf C$, and let $J$ be the Hecke ideal of Proposition~\ref{sha-hecke-ideal}, built from data as in Lemma~\ref{sha-hecke-data} and the level of Lemma~\ref{sha-hecke-algebra}. If $J$ is radical, then $J$ satisfies the multiplicity bound (Definition~\ref{sha-automorphy-unramified-multiplicity}(3)).
\end{lemma}
\leanref{Development/AutomorphicArithmetic/AutomorphicVarieties/ModularityArgument/ThresholdArguments/DenseThresholdCases/ZEqDimE.lean}{UpperTheory.V.Z.Desc.V0.NumJEAlt}
\leanref{Development/AutomorphicArithmetic/AutomorphicVarieties/ModularityArgument/ThresholdArguments/DenseThresholdCases/TNumJE.lean}{UpperTheory.V.Z.Desc.V0.numJE_alt}
\leanref{Development/AutomorphicArithmetic/AutomorphicVarieties/ModularityArgument/ThresholdArguments/DenseThresholdCases/CountE.lean}{UpperTheory.V.Z.Desc.V0.numJEDenseThresholdCases}
\begin{proof}
Use Situation~\ref{sha-hecke-ideal-tw-augmented-setup} with $S'=S$, so that the module $M$ of Lemma~\ref{sha-hecke-algebra} is $M_{()}$, of level $U=\mathcal K_0$. If $M=0$, then $J=B$ (proof of Proposition~\ref{sha-hecke-ideal}); the zero module is a Hecke module for $J=B$, with $e=0$, and no point lies on the zero locus of $B$, so the bound holds. Otherwise $J$ is the kernel of the surjection $f:B\to\mathbf T$ of the proof of Proposition~\ref{sha-hecke-ideal}, which is compatible with the maps from $\Lambda$.
\begin{enumerate}
\item \emph{The Hecke module.} Let $B$ act on $M$ through $f$. Then $\Lambda$ acts through $\lambda$, and $M$ is finite and free over $\Lambda$ (Lemma~\ref{sha-hecke-algebra}); $J$ is radical by hypothesis. As $\mathbf T\subset\operatorname{End}_\Lambda(M)$, the annihilator of $M$ in $B$ is $\ker f=J$, so a prime of $B$ contains it if and only if it contains $J$. If $\lambda(c)b\in J$ with $c\in\Lambda$ nonzero, then $c\,f(b)m=0$ for every $m\in M$; the coordinates of $f(b)m$ in a basis of $M$ over $\Lambda$ are killed by $c$, so they vanish, as $\Lambda$ is a domain (Lemma~\ref{sha-hecke-algebra}); so $f(b)=0$ and $b\in J$. Hence $M$ is a Hecke module for $(\lambda,J)$.
\item \emph{The idempotent.} Let $a$ and $j$ be as in Lemma~\ref{sha-automorphy-unramified-hecke-local} for $F=L_{\tilde v_0}$, and $e$ as in Lemma~\ref{sha-automorphy-unramified-hecke-projector}. It is $\Lambda$-linear with $e^2=e$, and it commutes with every $f(b)$, so $e(bm)=be(m)$ for $b\in B$ and $m\in M$.
\item \emph{The points.} Let $x$ and $x'$ be irreducible points on the zero locus of $J$ and $h$ an increasing weight such that $x$ is arithmetic of weight $h$ and $x'$ is of distinct type of weight $h$. Let $\lambda$ be the weight of Lemma~\ref{sha-automorphy-unramified-hecke-weights} for $h$; by that lemma, $x|_\Lambda$ and $x'|_\Lambda$ are arithmetic points of $\Lambda$ of weight $\lambda$. As $J=\ker f$ lies in their kernels, $x=y\circ f$ and $x'=y'\circ f$ for $\mathcal O$-algebra maps $y,y':\mathbf T\to\bar{\mathbf Q}_\ell$.
\item \emph{The multiplicities.} Lemma~\ref{sha-automorphy-unramified-hecke-eigenspace}, applied to $y$, gives $\xi$, $\operatorname{sp}$ and $Z$, with $|\iota^{-1}(\iota y(U_{v,i}))|=1$, $Z\ne0$, $\rho_Y$ irreducible for every automorphic part $Y$ of weight $\xi$ with $Y\cap Z\ne0$, and $\mu_e(x)=\dim(Z\cap\ker(E^\xi-q^j))$; and similarly for $y'$, with $\xi'$, $\operatorname{sp}'$ and $Z'$. These are data as in Lemma~\ref{sha-automorphy-unramified-hecke-count}, with $\epsilon_t=\iota y(t)$ and $\epsilon'_t=\iota y'(t)$. By Lemma~\ref{sha-automorphy-unramified-hecke-dimension}, $1\le\mu_e(x)=\mu_e(x')$.
\end{enumerate}
So $1\le\mu_e(x)\le\mu_e(x')$, which is the multiplicity bound.
\end{proof}
% numJE_alt (TNumJE.lean): the case J = top (datTop: PUnit) and the case of a surjection f; hiComm for freeness and finiteness; the private datHecke and datEnd for (1) (the annihilator is J, and tf through Module.Free.chooseBasis); hZ (ZEqDimE, Lemma sha-automorphy-unramified-hecke-dimension) for a and N = q^j with not l | N, normS and eEx for (2), with unitN; arith_good_iff_hset_of_theta_incr_weight (Lemma sha-automorphy-unramified-hecke-weights) for (3), with isDomain_and_denseWt_of_isTheta; core_hecke_of_automorphic_hecke_dense_weight for out; the private absDimE (Lemma sha-automorphy-unramified-hecke-eigenspace) twice, and hz. numJEDenseThresholdCases (CountE.lean) supplies the inputs muZ, muZE, eEx, spec2, normS, the proof of ZEqDimE from countE'_alt, resEq2_alt, invEq1, invBox, locE and kirCusp, unitEv and partRep_s. The Lean states NumJEAlt for any Taylor-Wiles data Q and level K; the book needs only the empty data.

\begin{lemma}\label{sha-automorphy-unramified-hecke}
Assume that Galois representations exist, in the sense of Definition~\ref{gal-galois-representations}, relative to the maps satisfying the laws of Definition~\ref{gal-llc-laws}. In Situation~\ref{sha-hecke-situation}, let $n=3$ and $\iota:\bar{\mathbf Q}_\ell\cong\mathbf C$, and suppose that the residual characters $\bar\chi_i$ are trivial on $I_{v_0}$ and agree on $D_{v_0}$ and that $\varepsilon\bmod\ell$ is nontrivial on $D_{v_0}$. Then there is an ideal $J$ of $B$ with the properties (1) to (6) of Proposition~\ref{sha-hecke-ideal} and the bounds of Proposition~\ref{sha-patching-bound} that satisfies the multiplicity bound (Definition~\ref{sha-automorphy-unramified-multiplicity}).
\end{lemma}
\leanref{Development/AutomorphicArithmetic/AutomorphicVarieties/ModularityArgument/ThresholdArguments/DenseThresholdCases/LevelRelationTwoX.lean}{UpperTheory.V.J.ModNAlt}
\leanref{Development/AutomorphicArithmetic/AutomorphicVarieties/ModularityArgument/ThresholdArguments/DenseThresholdCases/RadOfDense.lean}{UpperTheory.V.Z.J.D.modularity_n_of_num_j_printed_mult_comp}
\begin{proof}
Let $J$ be the Hecke ideal of Proposition~\ref{sha-hecke-ideal}, built from the data of Lemma~\ref{sha-hecke-data}; it has the properties (1) to (6), and Proposition~\ref{sha-patching-bound} gives the bounds. By the density (2), $J$ is radical: if $a^k\in J$, then $a$ vanishes at every point on the zero locus of $J$, so $a\in J$.
As $n=3\le4$ and $J$ is radical, $J$ satisfies the multiplicity bound by Lemma~\ref{sha-automorphy-unramified-hecke-bound}.
% ModNAlt is modularity_n_of_printed_mult_comp = modularity_n_of_num_j_printed_mult_comp numJ'_alt (RadOfDense.lean): Hi.nonempty_unitaryHeckeData and Hi.exists_hecke_ideal with no Taylor-Wiles places (QData.empty), the properties of HidaPropsAlt, Pa.pFromAut, and NumJAlt (numJ'_alt = numJEDenseThresholdCases, Lemma sha-automorphy-unramified-hecke-bound) with J radical (rad_of_dense) and Lambda a domain (isDomain_and_denseWt_of_isTheta). The hypotheses Z.Mult.Mult and Z.Mult.CompOf (MultiplicityOne/CompOfBis.lean) are used in Lemma sha-automorphy-unramified-hecke-count; the Lean proves them as Z.Mult.mult_free and Z.Mult.compOf_free.
\end{proof}

\subsection{Tangent spaces at linked points}\label{subsection-sha-linked-tangents}

\begin{lemma}\label{sha-automorphy-realization-presentation}
In Situation~\ref{sha-automorphy-link-situation}, let $n\in\{2,3\}$, let $\mathcal O_1\supseteq\mathcal O$ be the ring of integers of a finite extension of $\mathbf Q_\ell$, with residue field $k_1$, and let $A_n$, its universal linked point $(\rho^{\mathrm u},\psi^{\mathrm u})$, $B$ and $A_n^\circ$ be as in Definition~\ref{sha-automorphy-realization}(2). Then $A_n$ has a weighted presentation $\pi:\mathcal O_1[[X_s:s\in\Sigma]]\to A_n$ with weights in $\mathbf Z$ (Definition~\ref{sha-counts-vd-model}(1), with $\mathcal O_1$ and $k_1$ in place of $\mathcal O$ and $\mathbf F$), whose variables map into $\mathfrak m_{A_n}$, such that, with $(A_n)_j$ its parts:
\begin{enumerate}
\item each entry $\rho^{\mathrm u}(g)_{ij}$ lies in $(A_n)_{w_{n,i}-w_{n,j}}$, and each value of each $\psi^{\mathrm u}_{v,i}$ lies in $(A_n)_0$;
\item $B=(A_n)_0$; for every $j\in\mathbf Z$, $(A_n)_j$ is the $B$-submodule of $A_n$ spanned by the products of entries of $\rho^{\mathrm u}$ of total weight $j$; and $A_n^\circ$ is the sum of the $(A_n)_j$, which is direct;
\item $B$ is a complete local noetherian $\mathcal O_1$-algebra with residue field $k_1$, whose maximal ideal is $B\cap\mathfrak m_{A_n}$;
\item each $(A_n)_j$ is a finitely generated $B$-module, and $A_n^\circ$ is a finitely generated $B$-algebra, hence a noetherian ring.
\end{enumerate}
\end{lemma}
\leanref{Development/AutomorphicArithmetic/AutomorphicVarieties/FrobeniusSmoothVectors/ContMap.lean}{UpperTheory.V.Fr.wpres}
\leanref{Development/AutomorphicArithmetic/AutomorphicVarieties/FrobeniusSmoothVectors/CoB.lean}{UpperTheory.V.Fr.co_B}
\leanref{Development/AutomorphicArithmetic/AutomorphicVarieties/FrobeniusSmoothVectors/IsInternalCongr.lean}{UpperTheory.V.Fr.gr_A}
\begin{proof}
Let $\mathcal C$ be the class of complete local noetherian $\mathcal O_1$-algebras with residue field $k_1$. For an $\mathcal O_1$-algebra map $f:A_n\to R'$ with $R'$ in $\mathcal C$, the image of $(\rho^{\mathrm u},\psi^{\mathrm u})$ under $f$ is a linked point lifting the residual one: conditions (a) to (c) of Definition~\ref{sha-descent-link}(2) are identities, $f$ is local and so keeps the continuity (d), and $f$ is compatible with the maps to $k_1$, which keeps the reductions (e).
\begin{enumerate}
\item[(i)] \emph{Coordinates.} For $g\in\Gamma_0$ and $1\le i,j\le n$ let $e_{g,i,j}=\rho^{\mathrm u}(g)_{ij}$ if $i\ne j$, and $e_{g,i,i}=\rho^{\mathrm u}(g)_{ii}-a_{g,i}$ with $a_{g,i}\in\mathcal O_1$ of the same residue; give $e_{g,i,j}$ the weight $w_{n,i}-w_{n,j}$. For $v$, $i$ and $h\in G_v$ let $e_{v,i,h}=\psi^{\mathrm u}_{v,i}(h)-a_{v,i,h}$ with $a_{v,i,h}\in\mathcal O_1$ of the same residue, of weight $0$. All these elements lie in $\mathfrak m_{A_n}$, those off the diagonal by condition (e).
\item[(ii)] \emph{Uniqueness.} Two $\mathcal O_1$-algebra maps $f,f':A_n\to R'$ with $R'$ in $\mathcal C$ that agree at every coordinate agree on the entries of $\rho^{\mathrm u}$ and on the values of the $\psi^{\mathrm u}_{v,i}$. So they carry $(\rho^{\mathrm u},\psi^{\mathrm u})$ to the same linked point, and $f=f'$ by the universal property of $A_n$.
\item[(iii)] \emph{Twists.} Let $f:A_n\to R'$ be an $\mathcal O_1$-algebra map with $R'$ in $\mathcal C$, let $u\in R'^\times$, and let $d=\operatorname{diag}(u^{w_{n,1}},\dots,u^{w_{n,n}})$. Then $g\mapsto d\,f(\rho^{\mathrm u}(g))\,d^{-1}$, with the characters $f\circ\psi^{\mathrm u}_{v,i}$, is a linked point lifting the residual one. Indeed $w_{n,n+1-i}=-w_{n,i}$, so $J_ndJ_n=d^{-1}$, which gives (a); conjugation by a diagonal matrix keeps the vanishing of the entries in (b) and of the products in (c); and it multiplies the entry $(i,j)$ by $u^{w_{n,i}-w_{n,j}}$, which is $1$ on the diagonal, so it keeps (d) and (e). The universal property gives $f':A_n\to R'$ with $f'(\rho^{\mathrm u}(g)_{ij})=u^{w_{n,i}-w_{n,j}}f(\rho^{\mathrm u}(g)_{ij})$ and $f'\circ\psi^{\mathrm u}_{v,i}=f\circ\psi^{\mathrm u}_{v,i}$. So $f'(e)=u^{w}f(e)$ for every coordinate $e$ of weight $w$. A torus element $t$ of $R'$ for the group $\mathbf Z$ (Lemma~\ref{sha-counts-vd-model-torus-grading}) is $j\mapsto u^j$ with $u=t(1)$.
\item[(iv)] \emph{The presentation.} By (i) to (iii), Lemma~\ref{sha-counts-vd-model-torus-grading}(1), with $\mathcal O_1$, $k_1$ and the group $\mathbf Z$, gives a weighted presentation $\pi$ whose variables map to coordinates, such that every coordinate of weight $w$ lies in $(A_n)_w$. The constants $a_{g,i}$ and $a_{v,i,h}$ lie in $(A_n)_0$, which contains $1$; this gives (1).
\item[(v)] \emph{$B\subseteq(A_n)_0$.} Let $D$ be the $\mathcal O_1$-subalgebra generated by the values of the $\psi^{\mathrm u}_{v,i}$ and the products of entries of total weight $0$, so that $B$ is the closure of $D$. Since $(A_n)_j(A_n)_{j'}\subseteq(A_n)_{j+j'}$, (1) gives $D\subseteq(A_n)_0$. An element of $B$ is congruent modulo every $\mathfrak m_{A_n}^N$ to an element of $D$, so it lies in $(A_n)_0$ by Lemma~\ref{sha-counts-vd-presentations}(e)(i).
\item[(vi)] \emph{Monomials.} Each variable maps either to an entry $\rho^{\mathrm u}(g)_{ij}$ with $i\ne j$, of weight $w_{n,i}-w_{n,j}$, or to an element of $D$ of weight $0$, a diagonal entry or a value of a character minus a constant. So for every monomial $X^m$ of weight $j$, $\pi(X^m)=dx$ with $d\in D$ and $x$ a product of entries of total weight $j$; in particular $\pi(X^m)\in D$ if $j=0$.
\item[(vii)] \emph{$(A_n)_0\subseteq B$.} Let $y=\pi(f)$ with $f$ of weight $0$, and let $N\ge1$. The sum $f_{\le N}$ of the terms of $f$ whose exponents are all at most $N$ is a polynomial of weight $0$, so $\pi(f_{\le N})\in D$ by (vi). By Lemma~\ref{sha-counts-vd-presentations}(b), $f-f_{\le N}$ is a finite sum $\sum_mb_mX^m$ in which each $X^m$ has an exponent greater than $N$; as the variables map into $\mathfrak m_{A_n}$, $\pi(f-f_{\le N})\in\mathfrak m_{A_n}^N$. So $y$ lies in the closure of $D$, which is $B$.
\item[(viii)] \emph{The parts and $A_n^\circ$.} By (1), the products of entries of total weight $j$ lie in $(A_n)_j$, which is a module over $(A_n)_0=B$. Conversely, by Lemma~\ref{sha-counts-vd-presentations}(b), every element of $(A_n)_j$ is a finite sum $\sum_m\pi(b_m)\pi(X^m)$ with $b_m$ of weight $0$ and $X^m$ of weight $j$, and by (vi) each term is $\pi(b_m)d\,x$ with $\pi(b_m)d\in B$ and $x$ a product of entries of total weight $j$. This is the second statement of (2). By Lemma~\ref{sha-counts-vd-presentations}(b), the sum of the $(A_n)_j$ is direct, and it is an $\mathcal O_1$-subalgebra; it contains $B$ and the entries, and it lies in $A_n^\circ$ by the second statement. So it is $A_n^\circ$.
\item[(ix)] \emph{(3) and (4).} By Lemma~\ref{sha-counts-vd-presentations}(c), $(A_n)_0=B$ is a complete local noetherian $\mathcal O_1$-algebra with residue field $k_1$ and maximal ideal $B\cap\mathfrak m_{A_n}$. By Lemma~\ref{sha-counts-vd-presentations}(b), each $(A_n)_j$ is $\sum_{m\in\mu}B\,\pi(X^m)$ for a finite set $\mu$, and the sum of the parts, which is $A_n^\circ$, is generated as a $B$-algebra by the images of the finitely many variables. As $B$ is noetherian, so is $A_n^\circ$, by Hilbert's basis theorem.
\end{enumerate}
\end{proof}
% wpres (ContMap.lean): coord (the coordinates e, he, huniq from uniq with twist_lemma, htw from twist with alg_twist, then coord_grading), mono_lemma_2 and mono for (vi), dense for (vii) (the truncation low over Finset.Iic c and part_span), part_closed for (v), dec for (viii). co_B: wPres_gr_zero_isCO_of_mem_maximalIdeal and co_quot for (3). gr_A: gr_transfer with gr_noeth and gr_internal for (4) and the grading of A by S.gr. The hypotheses F.Ok C_0 and n in {2,3} of wpres are not used by coord.

\begin{lemma}\label{sha-automorphy-unramified-realization-irreducible}
In Situation~\ref{sha-automorphy-link-situation}, let $n\in\{2,3\}$, let $\mathcal O_1\supseteq\mathcal O$ be the ring of integers of a finite extension of $\mathbf Q_\ell$, and let $A_n$, its universal linked point $(\rho^{\mathrm u},\psi^{\mathrm u})$, $B$ and $A_n^\circ$ be as in Definition~\ref{sha-automorphy-realization}(2). Let $K$ be a field and $f:A_n^\circ\to K$ a ring homomorphism whose kernel is off-triangular (Definition~\ref{sha-automorphy-realization}(4)). Then $0$ and $K^n$ are the only subspaces of $K^n$ stable under every $f(\rho^{\mathrm u}(g))$, $g\in\Gamma_0$.
\end{lemma}
\leanref{Development/AutomorphicArithmetic/AutomorphicVarieties/FrobeniusSmoothVectors/IrrOfConn.lean}{UpperTheory.V.Fr.types}
\leanref{Development/AutomorphicArithmetic/AutomorphicVarieties/FrobeniusSmoothVectors/AddMulSelfEqOfOrthogonalIdempotents.lean}{UpperTheory.V.Fr.idem}
\leanref{Development/AutomorphicArithmetic/AutomorphicVarieties/ModularityArgument/LinkedUnramifiedChains/ScalarIrreducibility/IrrTI.lean}{UpperTheory.V.J.irrTI}
\begin{proof}
Write $\rho=\rho^{\mathrm u}$, $g^*=c_\Gamma g^{-1}c_\Gamma^{-1}$, and $E_{ii}$ for the diagonal matrix units.
\begin{enumerate}
\item \emph{$B$ is local.} By Lemma~\ref{sha-automorphy-realization-presentation}(3), which rests on the weighted presentation of $A_n$, $B$ is a complete local noetherian ring with maximal ideal $B\cap\mathfrak m_{A_n}$.
\item \emph{A diagonal matrix.} By condition (a) of Definition~\ref{sha-descent-link}(2), $\rho(g^*)_{ij}=\rho(g)_{n+1-j,n+1-i}$, so the diagonal entries of $\rho(g_0)-\rho(g_0^*)$ are $e_{n,i}\zeta$ with $\zeta=\rho(g_0)_{11}-\rho(g_0)_{nn}$, and by condition (b) its other entries vanish. The diagonal entries of $\rho$ have weight $0$, so $\zeta\in B$; by condition (e), $\zeta$ reduces to $a-a^{-1}$ with $a=\bar\Psi_\Gamma(g_0)$, which is nonzero since $a^2\ne1$. By (1), $\zeta$ is a unit of $B$, so $Y=\operatorname{diag}(e_{n,1},\dots,e_{n,n})=\zeta^{-1}\bigl(\rho(g_0)-\rho(g_0^*)\bigr)$ lies in the $B$-span $\Sigma$ of the matrices $\rho(g)$. This span is a $B$-algebra, since $\rho$ is multiplicative and $\rho(1)=1$.
\item \emph{The idempotents.} The element $2$ is a unit of $B$, as $\ell$ is odd. If $n=2$, then $2E_{11}=1+Y$ and $2E_{22}=1-Y$; if $n=3$, then $2E_{11}=Y^2+Y$, $E_{22}=1-Y^2$ and $2E_{33}=Y^2-Y$. So every $E_{ii}$ lies in $\Sigma$.
\item \emph{Stable subspaces.} Let $W\subseteq K^n$ be stable under every $f(\rho(g))$. Since $f(b)$ acts on $K^n$ as a scalar for $b\in B$, $W$ is stable under the image of $\Sigma$, hence under every $f(E_{ii})$. So if $w\in W$ has $w_i\ne0$, the basis vector $\epsilon_i$ lies in $W$. If $\epsilon_j\in W$ and $f(\rho(g)_{ij})\ne0$ for some $g$, then $\epsilon_i\in W$, since $f(\rho(g))\epsilon_j$ has $i$-th coordinate $f(\rho(g)_{ij})$.
\item \emph{Connectivity.} Write $Z(i,j)$ if $f(\rho(g)_{ij})=0$ for every $g$; by (2), $Z(i,j)$ implies $Z(n+1-j,n+1-i)$. Since the kernel of $f$ is off-triangular, $Z(i,j)$ fails for some $i>j$ and for some $i<j$. If $n=2$, $Z(2,1)$ and $Z(1,2)$ both fail. If $n=3$, $Z(2,1)$ and $Z(3,2)$ are equivalent, as are $Z(1,2)$ and $Z(2,3)$, so $Z(2,1)$ and $Z(3,1)$ do not both hold, $Z(1,2)$ and $Z(1,3)$ do not both hold, and, by the condition on the entries of odd weight, $Z(1,2)$ and $Z(2,1)$ do not both hold. If neither holds, then none of $Z(1,2)$, $Z(2,1)$, $Z(2,3)$, $Z(3,2)$ holds. If $Z(2,1)$ holds, then $Z(3,1)$, $Z(1,2)$ and $Z(2,3)$ fail. If $Z(1,2)$ holds, then $Z(1,3)$, $Z(2,1)$ and $Z(3,2)$ fail. In each case, for all $i,j$ there is a chain $j=i_0,i_1,\dots,i_m=i$ along which $Z(i_{k+1},i_k)$ fails for every $k$.
\item \emph{Conclusion.} If $W\ne0$, then by (4) some $\epsilon_j$ lies in $W$, and by (4) and (5) every $\epsilon_i$ does; so $W=K^n$.
\end{enumerate}
\end{proof}
% Fr.types (IrrOfConn.lean), the direction from TI to MatIrr: irr_of_conn with idem (AddMulSelfEqOfOrthogonalIdempotents.lean: Y from Integral.ze, then add_mul_self_eq_of_orthogonal_idempotents) for (2) and (3), the hypothesis Integral O S.B S.u from cone_u (ArtinRees.lean, the field ze: zeta = aB g0 0 - aB g0 (n-1)), and conn2, conn3 (with jtj for the symmetry) for (5). Step (1) is the instance IsCO O S.B of co_B (Lemma sha-automorphy-realization-presentation(3)), which cone_u takes as a hypothesis. irrTI does not use the integrality of the point.

\begin{lemma}\label{sha-automorphy-unramified-realization}
In Situation~\ref{sha-automorphy-link-situation}, suppose that $N$ is even and let $(j,\chi',\theta)$ be a seed datum for $L$. For every subring $A$ of $\bar{\mathbf Q}_\ell$ finite over $\mathcal O$ there are $\mathcal O_1\supset A$, large enough, a realization of rank three twisted by $\chi_{\mathrm{cyc}}^{-1}$ with coefficients $\mathcal O_1$, an isomorphism of $\Lambda$ with a power series ring over $\mathcal O_1$ in finitely many variables, an ideal $J_3$ of $B_3$ and a compatible map $\phi$ from $B_3$ onto the subalgebra $B$ of $A_3$ (Definition~\ref{sha-automorphy-realization}(2)), the closure of the $\mathcal O_1$-subalgebra generated by the values of the characters $\psi^{\mathrm u}_{v,i}$ and by the products of entries of the universal linked point of total weight $0$, such that: $J_3$ has the properties (1) to (6) of Proposition~\ref{sha-hecke-ideal} and the bounds of Proposition~\ref{sha-patching-bound}, and every minimal prime over $J_3$ has dimension $1+3d$; $J_3$ satisfies the multiplicity bound (Definition~\ref{sha-automorphy-unramified-multiplicity}); every off-triangular prime of $A_3^\circ$ (Definition~\ref{sha-automorphy-realization}) contains $\phi(J_3)$; and for every $\mathcal O_1$-algebra map $\xi:A_3^\circ\to\bar{\mathbf Q}_\ell$ with values of absolute value at most $1$, with $|\xi(b)|<1$ for $b$ in the maximal ideal of $B$ and $|\xi(a)|<1$ for the entries $a$ off the diagonal of the universal linked point, whose kernel is off-triangular, and such that $\xi\circ\phi$ is arithmetic of some weight, the representation $r_\xi\otimes\chi_{\mathrm{cyc}}^{-1}$, where $r_\xi$ is the universal linked point composed with $\xi$ and $p$, is attached through $\iota$ to a regular algebraic cuspidal automorphic representation $\pi$ of $GL_3(\mathbf A_L)$ with $\pi^\vee\cong\pi^c$ and is windowed.
\end{lemma}
\leanref{Development/AutomorphicArithmetic/AutomorphicVarieties/ModularityArgument/ThresholdArguments/DenseThresholdCases/UnitaryRankTwoFlags.lean}{UpperTheory.Gal.projectedFormsSystemPart.OpenEx2EAlt}
\leanref{Development/AutomorphicArithmetic/AutomorphicVarieties/ModularityArgument/ThresholdArguments/DenseThresholdCases/UnitaryRankTwoFlags.lean}{UpperTheory.Gal.projectedFormsSystemPart.openEx2E_alt}
\begin{proof}
Let $K$ be the quadratic field of the seed datum and $M$ the integer of the link.
\begin{enumerate}
\item \emph{Twists.} As in step (1) of the proof of (c) in Theorem~\ref{sha-residually-reducible-automorphy}, there is an algebraic Hecke character $\kappa$ of $K$ of infinity type $(-1,0)$ with $\kappa\kappa^c=\chi_{\mathrm{cyc}}^{-1}$. By Lemma~\ref{sha-automorphy-realization-twists}, $\kappa_3=\chi_{\mathrm{cyc}}^{-1}$ and the restriction $\kappa_2$ of $\kappa$ to $\operatorname{Gal}(\bar{\mathbf Q}/L)$ are admissible twists of ranks three and two.
\item \emph{Coefficients.} The rings $A$ and $\mathcal O$ lie in the ring of integers $\mathcal O'$ of a finite extension of $\mathbf Q_\ell$. Lemma~\ref{sha-automorphy-realization-exists}(1), applied to $\mathcal O'$, $\kappa_2$ and $\kappa_3$, gives $\mathcal O_1\supseteq\mathcal O'$, large enough, that contains the values of the twists.
\item \emph{Realizations.} Lemma~\ref{sha-automorphy-realization-exists}(2) gives realizations $R_2$ and $R_3$ of ranks two and three, twisted by $\kappa_2$ and $\kappa_3$, with coefficients $\mathcal O_1$, and a realization $R_0$ of rank three twisted by $\kappa_3$ whose place $v_0$ has residue characteristic prime to that of the place $v_0$ of $R_3$. The universal ring $A_3$ of linked points of rank three over $\mathcal O_1$ depends only on the link, so $R_0$ is taken with the same $A_3$ and the same $B\subset A_3$ as $R_3$. The realization of the statement is $R_3$.
\item \emph{Seeds.} Since $N$ is even and $\ell N$ divides $M$, $2\ell$ divides $M$, and $M$ satisfies the other hypotheses of (b) in the proof of Theorem~\ref{sha-residually-reducible-automorphy}. So (a) and (b) there, with $m\ge2$, give seeds of ranks two and three, and by Lemma~\ref{sha-automorphy-realization-seed}, $R_2$, $R_3$ and $R_0$ satisfy the conclusion of that lemma for every ideal with the property of automorphic points.
\item \emph{The Iwasawa algebras and the ideals.} By steps (1) and (2) of the proof of Lemma~\ref{sha-automorphy-realization-ideal}, $L$ contains no primitive $\ell$-th root of unity, and the Iwasawa algebras of $R_2$, $R_3$ and $R_0$ are power series rings over $\mathcal O_1$ in $2d$, $3d$ and $3d$ variables. Propositions~\ref{sha-hecke-ideal} and~\ref{sha-patching-bound} give ideals $J_2$ and $J_0$ for $R_2$ and $R_0$, and Lemma~\ref{sha-automorphy-unramified-hecke}, by the conditions at $v_0$ of $R_3$, gives an ideal $J_3$ for $R_3$ with the same properties and the multiplicity bound. Steps (3) to (6) of the proof of Lemma~\ref{sha-automorphy-realization-ideal}, with (4) above, show that $J_2$, $J_3$ and $J_0$ have the properties of that lemma; in particular every minimal prime over $J_3$ has dimension $1+3d$.
\item \emph{The surjections.} Lemma~\ref{sha-automorphy-realization-components}(0), applied to $R_2$ and $R_3$, gives the compatible surjection $\phi$ onto $B$, with $\phi(J_3)$ in every off-triangular prime of $A_3^\circ$. Applied to $R_2$ and $R_0$, it gives a compatible map $\phi_0$ from the universal ordinary determinant ring of $R_0$ to $B$, with $\phi_0(J_0)$ in every off-triangular prime of $A_3^\circ$.
\item \emph{Automorphy at the points of $B$.} Let $\xi$ be as in the statement, with $x=\xi\circ\phi$ arithmetic of weight $h$, and let $\rho_\xi$ be the linked point given by $\xi$. The kernel of $\xi$ is an off-triangular prime, so it contains $\phi(J_3)$, and $x$ lies on the zero locus of $J_3$. The linked point $\rho_\xi$ is absolutely irreducible, by Lemma~\ref{sha-automorphy-unramified-realization-irreducible} applied to $\xi$ and $K=\bar{\mathbf Q}_\ell$, since its kernel is off-triangular.
By compatibility, $x(D^{\mathrm u})$ is the determinant of $r_\xi\otimes\chi_{\mathrm{cyc}}^{-1}$, which is irreducible. The classicality of $J_3$ (Proposition~\ref{sha-hecke-ideal}(5)) shows that $r_\xi\otimes\chi_{\mathrm{cyc}}^{-1}$ is attached through $\iota$ to a regular algebraic cuspidal $\pi$ with $\pi^\vee\cong\pi^c$ and is windowed at the places not above the place $v_0$ of $R_3$.
The point $x_0=\xi\circ\phi_0$ is also arithmetic of some weight: the places above $\ell$ of $R_0$ and $R_3$ correspond to the same places $A_v$ of the link, both realizations twist by $\chi_{\mathrm{cyc}}^{-1}$, and so the characters of $x_0$ on the decomposition groups above $\ell$ are those of $x$. The same argument for $R_0$, $J_0$ and $\phi_0$ shows that $r_\xi\otimes\chi_{\mathrm{cyc}}^{-1}$ is windowed at the places not above the place $v_0$ of $R_0$. As the two places $v_0$ have different residue characteristics, it is windowed at every place, as in step (7) of the proof of (c) in Theorem~\ref{sha-residually-reducible-automorphy}.
\end{enumerate}
\end{proof}
% openEx2E_alt (UnitaryRankTwoFlags.lean): (1) Sd.Twist and Pr.kaOf; (2) oe_big and Pr.coefEx, oe_res; (3) Pr.cRealEx for R_2 and R_3, cRealAv with p_0 = R_3.s.p for R_0, and setS R_0 R_3.S; (4) hS3 (with H.3, the divisibility 2 l | M from Even N and H.2, Z.Sym.Pl.linkUnr and Z.Sym.Pl.chiUnr) and hS2, each with bound 2, through seedPtOf_alt; (5) degEq, hLam, hMod for J_2 and J_0, hN = J.ModNAlt for J_3 with Num, and J.mod_l_open_existence (COpenEx.lean, the analogue of modularity_linked_ends); (6) J.absGood twice; (7) endPt_of_allGood_of_arithmeticPointClassicality (EndPtOf.lean): off_self for R_3 and for setS R_0 R_3.S (AllGood, IsPhi.maps, end2oN, irrTI = Fr.types in IrrOfConn.lean), la_move with psi_eq, kap and psi_j for the weight of x_0, and glueN with v0_ne.

\begin{lemma}\label{sha-automorphy-unramified-point}
In the situation of Lemma~\ref{sha-automorphy-unramified-realization}, let $(\rho,\psi)$ be a linked point of rank three with values in a discrete valuation ring $A\subset\mathcal O_1$ finite over $\mathcal O$, lifting the residual one, and let $r=\rho\circ p$.
\begin{enumerate}
\item There is an $\mathcal O_1$-algebra map $\xi:A_3\to\mathcal O_1$ that carries the universal linked point to $(\rho,\psi)$.
\item If $r$ is irreducible, then for each $j$ there are $g$ and $i\ne j$ such that $\xi$ does not vanish on the entry $(i,j)$ of the universal linked point at $g$; and the restriction of $\xi$ to $A_3^\circ$, with values in $\bar{\mathbf Q}_\ell$, has values of absolute value at most $1$, with $|\xi(b)|<1$ for $b$ in the maximal ideal of $B$.
\item If $r$ is irreducible and $(\rho,\psi)$ has the prescribed inertial characters, the point $\xi\circ\phi$ of $B_3$ is arithmetic of some weight and irreducible.
\end{enumerate}
\end{lemma}
\leanref{Development/AutomorphicArithmetic/AutomorphicVarieties/ModularityArgument/LinkedUnramifiedChains/DerivativeDivisibility/XiAdic.lean}{UpperTheory.V.J.xiEx}
\leanref{Development/AutomorphicArithmetic/AutomorphicVarieties/ModularityArgument/LinkedUnramifiedChains/DerivativeDivisibility/NotIrrOfLine.lean}{UpperTheory.V.J.xiTI}
\leanref{Development/AutomorphicArithmetic/AutomorphicVarieties/ModularityArgument/LinkedUnramifiedChains/DerivativeDivisibility/XiArith.lean}{UpperTheory.V.J.xiArith}
\begin{proof}
(1) Through $A\subset\mathcal O_1$, $(\rho,\psi)$ is a linked point over $\mathcal O_1$ lifting the residual one, with the residue map of the realization; the universal property of $A_3$ gives $\xi$.
(2) If for some $j$ all entries $(i,j)$ with $i\ne j$ vanished under $\xi$, the line spanned by the $j$-th basis vector would be stable under $r$, which is irreducible. The values of $\xi$ lie in $\mathcal O_1$, and $\xi$ is local, so it maps the maximal ideal $B\cap\mathfrak m_{A_3}$ of $B$ (Lemma~\ref{sha-automorphy-realization-presentation}(3)) into $\mathfrak m_{\mathcal O_1}$.
(3) Let $F=\xi\circ\phi:B_3\to\mathcal O_1$ and let $x$ be $F$ followed by the inclusion in $\bar{\mathbf Q}_\ell$. Since $\phi$ is compatible and $\xi$ carries the universal linked point to $(\rho,\psi)$, $F$ carries $(D^{\mathrm u},\psi^{\mathrm u})$ to the determinant of $\rho$ and the characters $\psi_{v,i}$, composed with the map $\pi$ of the realization and twisted by its character. The twist $\chi_{\mathrm{cyc}}^{-1}$ is admissible (Lemma~\ref{sha-automorphy-realization-twists}). Steps (3) and (4) of the proof of Lemma~\ref{sha-automorphy-realization-away}, with $\kappa=\chi_{\mathrm{cyc}}^{-1}$ and $n=3$, show that $x(D^{\mathrm u})$ is the determinant of $r\otimes\chi_{\mathrm{cyc}}^{-1}$ and that $x$ is arithmetic of an increasing weight, as $(\rho,\psi)$ has the prescribed inertial characters. The representation $r\otimes\chi_{\mathrm{cyc}}^{-1}$ is continuous: the entries of $r$ are those of $\rho\circ p$, which are continuous for the $\mathfrak m_{\mathcal O_1}$-adic topology, and that topology is the $\ell$-adic one on $\mathcal O_1$. It is irreducible because $r$ is. So $x$ is irreducible in the sense of Definition~\ref{sha-automorphy-unramified-multiplicity}(3).
% xiEx: xi_loc, xi_lift and the field univ of Fr.Space. xiTI: not_irr_of_line and ti_of (the condition TI on ker xi), norm_le_one and map_mem_max for IntPt. xiArith (XiArith.lean), from GlobalArtinReciprocity: hD and hP from IsPhi; the weight by the private ap_locAlg and ap_psi (copies of those of AutPtOf.lean) with H.e_pos; continuity by xa_gl, xa_cont (xa_pow, xa_key, xa_lnorm) and xa_adic; the determinant by ap_det; irreducibility by isIrr_twist_and_repIso_of_twist_repIso.
\end{proof}

\begin{lemma}\label{sha-automorphy-unramified-algebra-finite}
Let $\mathcal O$ be a ring, let $\Lambda$ and $B$ be $\mathcal O$-algebras with an $\mathcal O$-algebra map $\lambda:\Lambda\to B$, let $J$ be an ideal of $B$, and let $M$ be a Hecke module for $(\lambda,J)$ (Definition~\ref{sha-automorphy-unramified-multiplicity}). For $a\in B$ let $\chi_a\in\Lambda[X]$ be the characteristic polynomial of the endomorphism $m\mapsto am$ of the finite free $\Lambda$-module $M$, and for $q\in\Lambda[X]$ write $q(a)$ for its value at $a$ through $\lambda$.
\begin{enumerate}
\item Every element of $B$ that kills $M$ lies in $J$, and every element of $J$ has a power that kills $M$.
\item For every $a\in B$, $\chi_a(a)\in J$.
\item If $\Lambda$ is noetherian, there are $b_1,\dots,b_r\in B$ such that every $a\in B$ is congruent modulo $J$ to $\sum_i\lambda(c_i)b_i$ for some $c_1,\dots,c_r\in\Lambda$.
\item Let $K$ be a field that is an $\mathcal O$-algebra, $y:B\to K$ an $\mathcal O$-algebra map and $t\in B$, and let $\chi_t^y\in K[X]$ be the image of $\chi_t$ under $y\circ\lambda$. Then the multiplicity $\mu(y)$ is at most the multiplicity of $y(t)$ as a root of $\chi_t^y$.
\end{enumerate}
\end{lemma}
\leanref{Development/AutomorphicArithmetic/AutomorphicVarieties/ModularityArgument/LinkedUnramifiedChains/DerivativeDivisibility/IntMod.lean}{UpperTheory.V.J.intMod}
\leanref{Development/AutomorphicArithmetic/AutomorphicVarieties/ModularityArgument/LinkedUnramifiedChains/DerivativeDivisibility/MuLe.lean}{UpperTheory.V.J.muLe}
\begin{proof}
(1) A prime of $B$ contains $J$ if and only if it contains the annihilator of $M$. So an element that kills $M$ lies in every prime that contains $J$, hence in $J$, which is radical. An element of $J$ lies in every prime that contains the annihilator, hence in its radical.
(2) By the Cayley-Hamilton theorem, $\chi_a$ kills the endomorphism $m\mapsto am$, and $\chi_a(a)$ acts on $M$ by its value there, since $\Lambda$ acts through $\lambda$. So $\chi_a(a)$ kills $M$, and it lies in $J$ by (1).
(3) The endomorphisms $m\mapsto bm$, for $b\in B$, form a $\Lambda$-submodule of $\operatorname{End}_\Lambda(M)$, since $\lambda(c)b$ acts as $c$ times $b$. As $\operatorname{End}_\Lambda(M)$ is a finite module over the noetherian ring $\Lambda$, this submodule is generated by the endomorphisms of some $b_1,\dots,b_r$. If $a$ acts as $\sum_ic_ib_i$, then $a-\sum_i\lambda(c_i)b_i$ kills $M$ and lies in $J$ by (1).
(4) The space that defines $\mu(y)$ lies in the generalized eigenspace of $t$ on $M\otimes_\Lambda K$ for $y(t)$. Its dimension is the multiplicity of $y(t)$ as a root of the characteristic polynomial of $t$ on $M\otimes_\Lambda K$, which is $\chi_t^y$.
\end{proof}
% (1) and (3) are the private finite_span_mod_and_integral_of_j_dat of PrimEx.lean (the image of B in End M is a finitely generated Lambda-module); (2) is intMod (act_eval and aeval_self_charpoly); (4) is muLe (finrank_maxGenEigenspace_eq and charpoly_baseChange).

\begin{lemma}\label{sha-automorphy-unramified-algebra-primitive}
Let $\mathcal O$ be a ring, let $\Lambda$ be a noetherian domain of characteristic $0$ and $B$ a ring, both $\mathcal O$-algebras, with an $\mathcal O$-algebra map $\lambda:\Lambda\to B$, let $J$ be an ideal of $B$ and $M$ a Hecke module for $(\lambda,J)$ (Definition~\ref{sha-automorphy-unramified-multiplicity}). Let $K$ be a field of characteristic $0$ that is an $\mathcal O$-algebra, and let $x:B\to K$ be an $\mathcal O$-algebra map that vanishes on $J$. Then there is $t\in B$ such that:
\begin{enumerate}
\item for every $a\in B$ there are a nonzero $c\in\Lambda$ and $P\in\Lambda[X]$ with $\lambda(c)a-P(t)\in J$, where $P(t)$ is the value at $t$ through $\lambda$;
\item every $\mathcal O$-algebra map $x':B\to K$ that vanishes on $J$, with $x'\circ\lambda=x\circ\lambda$ and $x'(t)=x(t)$, equals $x$.
\end{enumerate}
\end{lemma}
\leanref{Development/AutomorphicArithmetic/AutomorphicVarieties/ModularityArgument/LinkedUnramifiedChains/DerivativeDivisibility/PrimEx.lean}{UpperTheory.V.J.primEx}
\begin{proof}
Let $b_1,\dots,b_r$ be as in Lemma~\ref{sha-automorphy-unramified-algebra-finite}(3), and let $f_i=\chi_{b_i}$, which is monic, with $f_i(b_i)\in J$ by part (2) of that lemma. For an integer $s\ge0$ let $t_s=\sum_{i=1}^rs^{i-1}b_i$.

\emph{Finite sets of tuples.} Let $E$ be a field of characteristic $0$ and $\rho:\Lambda\to E$ a ring map. The set $T$ of tuples $v\in E^r$ such that each $v_i$ is a root of the image of $f_i$ under $\rho$ is finite, since the $f_i$ are monic. For distinct $v,w\in T$ the polynomial $\sum_i(v_i-w_i)X^{i-1}$ is nonzero and has finitely many roots, and the integers are distinct in $E$. So for all but finitely many $s$, the map $v\mapsto\sum_iv_is^{i-1}$ is injective on $T$.

\emph{The quotient.} Let $R=B/J$, $F$ the fraction field of $\Lambda$, $\Omega$ an algebraic closure of $F$ and $A=R\otimes_\Lambda F$. The ring $R$ is reduced, as $J$ is radical; it is a finite $\Lambda$-module, generated by the images of the $b_i$; and nonzero elements of $\Lambda$ are nonzerodivisors on $R$, by the last condition on a Hecke module. So $R\to A$ is injective, and $A$ is a finite reduced $F$-algebra.

\emph{A primitive element of $A$.} Apply the finiteness above to $\rho:\Lambda\to\Omega$, with its set $T$. Every $F$-algebra map $\chi:A\to\Omega$ gives the tuple $(\chi(b_i))_i$ in $T$, since $f_i(b_i)\in J$. Let $s$ be such that $v\mapsto\sum_iv_is^{i-1}$ is injective on $T$. If two maps $\chi,\chi'$ agree at $t_s$, then they agree at every $b_i$, hence on $R$, which the $b_i$ generate as a $\Lambda$-module, and hence on $A$. The ring $A$ is a finite product of finite field extensions of $F$ \cite[Chapter~8]{AM69}, each separable since $F$ has characteristic $0$; so the number of $F$-algebra maps $A\to\Omega$ is $\dim_FA$ \cite[Chapter~V, \S4]{Lan02}. These maps restrict to distinct maps on the subalgebra $F[t_s]$, of which there are at most $\dim_FF[t_s]$. So $A=F[t_s]$.

\emph{Part (1) for such $s$.} Let $a\in B$. The image of $a$ in $A$ is $Q(t_s)$ for some $Q\in F[X]$. Choose a nonzero $c\in\Lambda$ with $P=cQ\in\Lambda[X]$. Then $\lambda(c)a-P(t_s)$ maps to $0$ in $A$, so its image in $R$ is killed by a nonzero element of $\Lambda$, and is $0$. That is, $\lambda(c)a-P(t_s)\in J$.

\emph{Part (2).} Apply the finiteness to $\rho=x\circ\lambda:\Lambda\to K$, with its set $T_K$, and choose $s$ such that $v\mapsto\sum_iv_is^{i-1}$ is injective on $T$ and on $T_K$; all but finitely many $s$ qualify. Let $x'$ be as in (2). Since $x$ and $x'$ vanish on $J$ and agree on $\lambda(\Lambda)$, the tuples $(x(b_i))_i$ and $(x'(b_i))_i$ lie in $T_K$, and they have the same image $x(t_s)=x'(t_s)$; so they are equal. Every $a\in B$ is congruent modulo $J$ to some $\sum_i\lambda(c_i)b_i$, so $x'(a)=\sum_ix(\lambda(c_i))x(b_i)=x(a)$. Take $t=t_s$.
\end{proof}
% primEx (PrimEx.lean): finite_span_mod_and_integral_of_j_dat; finite_setOf_roots_of_monic_family and finite_setOf_nat_polynomial_collision for the finite sets; finite_setOf_nat_not_primitive_element (with adjoin_singleton_eq_whole_of_separates_homs: IsArtinianRing.equivPi and AlgHom.natCard_of_splits) for (1); the tuples in K for (2).

\begin{lemma}\label{sha-automorphy-unramified-algebra-factors}
Let $\mathcal O$ be a ring, let $\Lambda$ be an integrally closed domain with fraction field $F$ and $B$ a ring, both $\mathcal O$-algebras, with an $\mathcal O$-algebra map $\lambda:\Lambda\to B$, let $J$ be an ideal of $B$ and $M$ a Hecke module for $(\lambda,J)$ (Definition~\ref{sha-automorphy-unramified-multiplicity}). Let $t\in B$ be such that for every $a\in B$ there are a nonzero $c\in\Lambda$ and $P\in\Lambda[X]$ with $\lambda(c)a-P(t)\in J$, and let $\chi_t$ be the characteristic polynomial of $t$ on $M$ (Lemma~\ref{sha-automorphy-unramified-algebra-finite}).
\begin{enumerate}
\item There are pairwise distinct monic $g_1,\dots,g_s\in\Lambda[X]$, irreducible in $F[X]$, and integers $m_i\ge1$ with $\chi_t=\prod_ig_i^{m_i}$.
\item Let $g=\prod_ig_i$. Then $g(t)\in J$, and every $q\in\Lambda[X]$ with $q(t)\in J$ is divisible by each $g_i$.
\item For each $i$ let $e_i=\prod_{j\ne i}g_j(t)$ and $Q_i=\{a\in B:ae_i\in J\}$. Then $Q_i$ is a prime ideal that contains $J$; $q(t)\in Q_i$ if and only if $g_i$ divides $q$; and $\lambda(c)\notin Q_i$ for every nonzero $c\in\Lambda$.
\item Let $K$ be an algebraically closed field that is an $\mathcal O$-algebra, $p:\Lambda\to K$ an $\mathcal O$-algebra map, and $z\in K$ a root of the image of $g_i$ under $p$. Then there is an $\mathcal O$-algebra map $x':B\to K$ that vanishes on $Q_i$, with $x'\circ\lambda=p$ and $x'(t)=z$.
\end{enumerate}
\end{lemma}
\leanref{Development/AutomorphicArithmetic/AutomorphicVarieties/ModularityArgument/LinkedUnramifiedChains/DerivativeDivisibility/FactEx.lean}{UpperTheory.V.J.factEx}
\leanref{Development/AutomorphicArithmetic/AutomorphicVarieties/ModularityArgument/LinkedUnramifiedChains/DerivativeDivisibility/MinPol.lean}{UpperTheory.V.J.minPol}
\leanref{Development/AutomorphicArithmetic/AutomorphicVarieties/ModularityArgument/LinkedUnramifiedChains/DerivativeDivisibility/CmpPrime.lean}{UpperTheory.V.J.cmpPrime}
\leanref{Development/AutomorphicArithmetic/AutomorphicVarieties/ModularityArgument/LinkedUnramifiedChains/DerivativeDivisibility/LyOver.lean}{UpperTheory.V.J.lyOver}
\begin{proof}
(1) Factor $\chi_t$ in $F[X]$ into monic irreducible factors. Each factor is a monic polynomial over $F$ dividing the monic $\chi_t\in\Lambda[X]$, so its coefficients lie in $\Lambda$, which is integrally closed \cite[Proposition~5.15]{AM69}. Group equal factors.

(2) The polynomial $\chi_t$ divides $g^m$ with $m=\sum_im_i$, and $\chi_t(t)\in J$ by Lemma~\ref{sha-automorphy-unramified-algebra-finite}(2), so $g(t)^m\in J$; as $J$ is radical, $g(t)\in J$. Let $q(t)\in J$. By Lemma~\ref{sha-automorphy-unramified-algebra-finite}(1), some $q(t)^N$ kills $M$, so $q(T)^N=0$, where $T$ is the endomorphism $m\mapsto tm$ of $M$. Let $z$ be a root of $g_i$ in an algebraic closure $\bar F$ of $F$. As $g_i$ divides $\chi_t$, $z$ is an eigenvalue of $T$ on $M\otimes_\Lambda\bar F$; for an eigenvector $v$, $q(z)^Nv=q(T)^Nv=0$, so $q(z)=0$. The minimal polynomial of $z$ over $F$ is the monic irreducible $g_i$, so $g_i$ divides $q$ in $F[X]$, and in $\Lambda[X]$ since $g_i$ is monic.

(3) Clearly $J\subseteq Q_i$. If $q(t)\in Q_i$, then $g_i$ divides $q\prod_{j\ne i}g_j$ by (2). In $F[X]$, $g_i$ is prime and divides no $g_j$ with $j\ne i$, these being distinct monic irreducible polynomials; so $g_i$ divides $q$ in $F[X]$, and in $\Lambda[X]$. Conversely, if $g_i$ divides $q$, then $q\prod_{j\ne i}g_j$ is a multiple of $g$, and $q(t)e_i\in J$ by (2). A nonzero constant $c$ is not divisible by $g_i$, which has positive degree; so $\lambda(c)\notin Q_i$, and in particular $1\notin Q_i$.
Membership in $Q_i$ depends only on the class modulo $J$, and for $c\ne0$, $\lambda(c)a\in Q_i$ if and only if $a\in Q_i$, by the last condition on a Hecke module applied to $\lambda(c)(ae_i)$. Let $ab\in Q_i$. By the hypothesis on $t$ there are nonzero $c,d\in\Lambda$ and $P_a,P_b\in\Lambda[X]$ with $\lambda(c)a\equiv P_a(t)$ and $\lambda(d)b\equiv P_b(t)$ modulo $J$. Then $(P_aP_b)(t)\equiv\lambda(cd)ab\in Q_i$, so $g_i$ divides $P_aP_b$. As $g_i$ is prime in $F[X]$, it divides $P_a$ or $P_b$ there, hence in $\Lambda[X]$; so $P_a(t)$ or $P_b(t)$ lies in $Q_i$, and so does $a$ or $b$. So $Q_i$ is prime.

(4) Let $f:\Lambda[X]\to B/Q_i$ be the map $q\mapsto q(t)\bmod Q_i$. By (3), its kernel consists of the multiples of $g_i$. The ring $B/Q_i$ is integral over the image of $f$, since each $a\in B$ is a root of the monic $\chi_a$ modulo $J\subseteq Q_i$ (Lemma~\ref{sha-automorphy-unramified-algebra-finite}(2)). The map $\Lambda[X]\to K$ that applies $p$ to the coefficients and evaluates at $z$ kills $g_i$, hence the kernel of $f$, so it factors through the image of $f$. Since $K$ is algebraically closed, it extends to $B/Q_i$ \cite[Chapter~5, Exercise~2]{AM69}. The composite of $B\to B/Q_i$ with this extension is $x'$.
\end{proof}
% (1) is factEx (normalizedFactors and IsIntegrallyClosed.eq_map_mul_C_of_dvd); (2) is minPol (dvd_of_aeval through an eigenvector over AlgebraicClosure and minpoly); (3) is cmpPrime (cmpPrime_core); (4) is lyOver (exists_ideal_over_prime_of_isIntegral and IsAlgClosed.lift in ext_integral and ext_injective).

\begin{lemma}\label{sha-automorphy-unramified-algebra-separating}
Let $\mathcal O$ be a ring, let $\Lambda$ and $B$ be $\mathcal O$-algebras with an $\mathcal O$-algebra map $\lambda:\Lambda\to B$, with $\Lambda$ nonzero, let $J$ be an ideal of $B$ and $M$ a Hecke module for $(\lambda,J)$ (Definition~\ref{sha-automorphy-unramified-multiplicity}). Let $K$ be an algebraically closed field that is an $\mathcal O$-algebra, $x:B\to K$ an $\mathcal O$-algebra map, and $t\in B$ such that every $\mathcal O$-algebra map $x':B\to K$ that vanishes on $J$, with $x'\circ\lambda=x\circ\lambda$ and $x'(t)=x(t)$, equals $x$. Then the multiplicity of $x(t)$ as a root of $\chi_t^x$ (Lemma~\ref{sha-automorphy-unramified-algebra-finite}(4)) is at most $\mu(x)$.
\end{lemma}
\leanref{Development/AutomorphicArithmetic/AutomorphicVarieties/ModularityArgument/LinkedUnramifiedChains/DerivativeDivisibility/LeJoint.lean}{UpperTheory.V.J.muGe}
\begin{proof}
Let $V=M\otimes_\Lambda K$, with $K$ a $\Lambda$-algebra through $x\circ\lambda$, and for $a\in B$ let $F_a$ be the endomorphism of $V$ given by $a$. The $F_a$ commute, $F_{\lambda(c)}$ is multiplication by $x(\lambda(c))$, and $F_a=0$ if $a$ kills $M$. The multiplicity of $x(t)$ as a root of $\chi_t^x$ is the dimension of the generalized eigenspace $V_t$ of $F_t$ for $x(t)$. So it suffices that every vector of $V_t$ lies in the generalized eigenspace of every $F_a$ for $x(a)$.

The space $V_t$ is stable under the commuting family $(F_a)$, and since $K$ is algebraically closed it is the sum of its joint generalized eigenspaces, one for each function $\chi:B\to K$. Let $v\ne0$ lie in the joint generalized eigenspace of $\chi$ in $V_t$. A nonzero vector lies in the generalized eigenspace of an endomorphism for at most one eigenvalue. If $v$ lies in those of $F_a$ for $\chi(a)$ and of $F_b$ for $\chi(b)$, then it lies in those of $F_{a+b}=F_a+F_b$ for $\chi(a)+\chi(b)$ and of $F_{ab}=F_aF_b$ for $\chi(a)\chi(b)$, since $F_a$ and $F_b$ commute. So $\chi$ is additive and multiplicative, and $\chi\circ\lambda=x\circ\lambda$; so $\chi$ is an $\mathcal O$-algebra map. For $a\in J$, some $a^n$ kills $M$ (Lemma~\ref{sha-automorphy-unramified-algebra-finite}(1)), so $F_{a^n}=0$ and $\chi(a)^n=0$. Also $\chi(t)=x(t)$, as $v\in V_t$. By the hypothesis, $\chi=x$. So $V_t$ is its own joint generalized eigenspace for $x$, which proves the claim.
\end{proof}
% muGe (LeJoint.lean): le_joint (iSup_iInf_maxGenEigenspace_eq_top_of_forall_mapsTo for the restrictions to the generalized eigenspace of t), ev_uniq, ev_add, ev_mul and ev_scal for the map x', and Jrad.

\begin{lemma}\label{sha-automorphy-unramified-algebra-conductor}
Let $\mathcal O$, $\Lambda$, $F$, $B$, $\lambda$, $J$, $M$, $t$, the $g_i$ and the $Q_i$ be as in Lemma~\ref{sha-automorphy-unramified-algebra-factors}, with $\Lambda$ of characteristic $0$, and let $g_i'$ be the derivative of $g_i$. Then for each $i$ and every $a\in B$ there is $W\in\Lambda[X]$ with $g_i'(t)a-W(t)\in Q_i$.
\end{lemma}
\leanref{Development/AutomorphicArithmetic/AutomorphicVarieties/ModularityArgument/LinkedUnramifiedChains/DerivativeDivisibility/CondField.lean}{UpperTheory.V.J.cond}
\begin{proof}
Let $E=F[X]/(g_i)$, a field finite over $F$ and separable since $F$ has characteristic $0$, and let $r$ be the class of $X$, so that $g_i'(r)\ne0$. Let $k$ be the fraction field of the domain $B/Q_i$, and write $\bar a$ for the class of $a\in B$. By Lemma~\ref{sha-automorphy-unramified-algebra-factors}(3), $\lambda(c)\notin Q_i$ for $c\ne0$, so $\Lambda\to B/Q_i$ is injective and extends to $F\to k$; as $g_i(t)\in Q_i$, there is a field homomorphism $\theta:E\to k$ over $\Lambda$ with $\theta(r)=\bar t$.
By the hypothesis on $t$ there are a nonzero $c\in\Lambda$ and $P\in\Lambda[X]$ with $\lambda(c)a-P(t)\in J\subseteq Q_i$. Let $\alpha=P(r)/c\in E$, so that $\theta(\alpha)=\bar a$. The element $\bar a$ is a root of the monic $\chi_a$ (Lemma~\ref{sha-automorphy-unramified-algebra-finite}(2)), and $\theta$ is injective, so $\alpha$ is integral over $\Lambda$.
For $y\in\Lambda[r]$, $\alpha y$ is integral over $\Lambda$, so its trace to $F$, a sum of conjugates of $\alpha y$, is integral over $\Lambda$ and lies in $\Lambda$, which is integrally closed. By Euler's formula \cite[Chapter~III, \S6]{Ser79}, the elements of $E$ whose traces against $\Lambda[r]$ all lie in $\Lambda$ form $g_i'(r)^{-1}\Lambda[r]$. So $g_i'(r)\alpha=W(r)$ for some $W\in\Lambda[X]$. Applying $\theta$ gives $g_i'(t)a-W(t)\in Q_i$, since $B/Q_i\to k$ is injective.
\end{proof}
% cond (CondField.lean) from cmpPrime: cond_core (AdjoinRoot.lift into FractionRing (B/Q)) and cond_field (traceForm_dualSubmodule_adjoin, Algebra.isIntegral_trace).

\begin{lemma}\label{sha-automorphy-unramified-algebra-derivation}
Let $\mathcal O$, $\Lambda$, $B$, $\lambda$, $J$, $M$, $t$, the $g_i$, $g$, the $e_i$ and the $Q_i$ be as in Lemma~\ref{sha-automorphy-unramified-algebra-factors}, with $\Lambda$ of characteristic $0$. Let $K$ be a field that is an $\mathcal O$-algebra and $x:B\to K$ an $\mathcal O$-algebra map that vanishes on $Q_i$ for some $i$, such that $x(t)$ is a simple root of the image $g^x$ of $g$ under $x\circ\lambda$. Then every $\mathcal O$-linear $\delta:B\to K$ with $\delta(ab)=x(a)\delta(b)+x(b)\delta(a)$ that vanishes on $J$ and on $\lambda(\Lambda)$ is $0$.
\end{lemma}
\leanref{Development/AutomorphicArithmetic/AutomorphicVarieties/ModularityArgument/LinkedUnramifiedChains/DerivativeDivisibility/DerZero.lean}{UpperTheory.V.J.derZero}
\begin{proof}
For $q\in\Lambda[X]$, the Leibniz rule and $\delta\circ\lambda=0$ give $\delta(q(t))=x(q'(t))\delta(t)$.
Since $g(t)\in J$, $0=\delta(g(t))=x(g'(t))\delta(t)$, and $x(g'(t))$ is the value of the derivative of $g^x$ at $x(t)$, which is nonzero as $x(t)$ is a simple root. So $\delta(t)=0$, and $\delta(q(t))=0$ for every $q\in\Lambda[X]$.
Next, $g'=g_i'\prod_{j\ne i}g_j+g_i\bigl(\prod_{j\ne i}g_j\bigr)'$, and $x(g_i(t))=0$ since $g_i(t)\in Q_i$ (Lemma~\ref{sha-automorphy-unramified-algebra-factors}(3)). So $x(g_i'(t))x(e_i)=x(g'(t))\ne0$.
For $b\in Q_i$, $be_i\in J$ and $x(b)=0$, so $0=\delta(be_i)=x(e_i)\delta(b)$, and $\delta(b)=0$.
Finally let $a\in B$, and let $W$ be given by Lemma~\ref{sha-automorphy-unramified-algebra-conductor}. Then $\delta(g_i'(t)a)=\delta(W(t))+\delta(g_i'(t)a-W(t))=0$, while $\delta(g_i'(t)a)=x(g_i'(t))\delta(a)+x(a)\delta(g_i'(t))=x(g_i'(t))\delta(a)$. So $\delta(a)=0$.
\end{proof}
% derZero (DerZero.lean): derivation_eval2_eq_eval2_derivative_mul, eval_derivative_ne_zero_of_rootMultiplicity_eq_one, and cond for the last step.

\begin{lemma}\label{sha-automorphy-unramified-algebra-rigid}
Let $\mathcal O$ be a ring, let $\Lambda$ be a noetherian integrally closed domain of characteristic $0$ and $B$ a ring, both $\mathcal O$-algebras, with an $\mathcal O$-algebra map $\lambda:\Lambda\to B$, let $J$ be an ideal of $B$ and $M$ a Hecke module for $(\lambda,J)$, with multiplicities $\mu$ (Definition~\ref{sha-automorphy-unramified-multiplicity}). Let $K$ be an algebraically closed field of characteristic $0$ that is an $\mathcal O$-algebra, let $x:B\to K$ be an $\mathcal O$-algebra map that vanishes on $J$, let $\mathcal G$ be a set of $\mathcal O$-algebra maps $B\to K$, and let $b_0\in B$ with $x(b_0)\ne0$. Suppose that $ab_0\in J$ whenever $a\in B$ vanishes at every $y\in\mathcal G$, and that every $y\in\mathcal G$ vanishes on $J$ and has $\mu(x)\le\mu(y)$. Then every $\mathcal O$-linear $\delta:B\to K$ with $\delta(ab)=x(a)\delta(b)+x(b)\delta(a)$ that vanishes on $J$ and on $\lambda(\Lambda)$ is $0$.
\end{lemma}
\leanref{Development/AutomorphicArithmetic/AutomorphicVarieties/ModularityArgument/CuspidalDescent/GLDenseRanges/InertiaPartMaps/UnrAlgW.lean}{UpperTheory.V.Z.Desc.V0.unrAlgW}
\begin{proof}
Lemma~\ref{sha-automorphy-unramified-algebra-primitive} gives $t\in B$ with its properties (1) and (2) for $x$. Let the $g_i$, $m_i$, $g$, $e_i$ and $Q_i$ be given by Lemma~\ref{sha-automorphy-unramified-algebra-factors} for $t$. For $y:B\to K$ and $q\in\Lambda[X]$ let $q^y$ be the image of $q$ under $y\circ\lambda$, and let $k_i$ be the multiplicity of $x(t)$ as a root of $g_i^x$.
\begin{enumerate}
\item \emph{The multiplicity at $x$.} The multiplicity of $x(t)$ as a root of $\chi_t^x=\prod_i(g_i^x)^{m_i}$ is $\sum_im_ik_i$, so $\sum_im_ik_i\le\mu(x)$ by Lemma~\ref{sha-automorphy-unramified-algebra-separating}.
\item \emph{Some $k_i$ is positive.} Since $g(t)\in J$ and $x$ vanishes on $J$, $x(t)$ is a root of $g^x=\prod_ig_i^x$.
\item \emph{The components through $x$.} Let $k_i>0$. Lemma~\ref{sha-automorphy-unramified-algebra-factors}(4), with $p=x\circ\lambda$ and $z=x(t)$, gives $x'$ that vanishes on $Q_i\supseteq J$, with $x'\circ\lambda=x\circ\lambda$ and $x'(t)=x(t)$. By (2) of Lemma~\ref{sha-automorphy-unramified-algebra-primitive}, $x'=x$; so $x$ vanishes on $Q_i$.
\item \emph{The bound $\mu(x)\le m_i$ when $k_i>0$.} The polynomial $g_i$ is not constant, since $g_i^x$ has a root. Suppose that every $y\in\mathcal G$ vanishes at $g_i'(t)$ or at $e_i$. Then $g_i'(t)e_i$ vanishes at every $y\in\mathcal G$, so $g_i'(t)e_ib_0\in J$, that is, $g_i'(t)b_0\in Q_i$. As $Q_i$ is prime, either $g_i'(t)\in Q_i$, and then $g_i$ divides $g_i'$ by Lemma~\ref{sha-automorphy-unramified-algebra-factors}(3), which is impossible as $g_i'$ is nonzero of smaller degree, $\Lambda$ having characteristic $0$; or $b_0\in Q_i$, and then $x(b_0)=0$ by (3), which is false. So some $y\in\mathcal G$ has $y(g_i'(t))\ne0$ and $y(e_i)\ne0$. Since $y$ vanishes on $J$, $y(g(t))=0$, so $y(t)$ is a root of $g_i^y$. It is a simple root, since $y(g_i'(t))\ne0$, and not a root of $g_j^y$ for $j\ne i$, since $y(e_i)\ne0$. So $y(t)$ is a root of $\chi_t^y$ of multiplicity $m_i$, and $\mu(x)\le\mu(y)\le m_i$ by Lemma~\ref{sha-automorphy-unramified-algebra-finite}(4).
\item \emph{A simple root.} Let $N=\mu(x)$. By (1), (2) and (4), $N\ge m_ik_i\ge1$ for some $i$, and $N\sum_ik_i\le\sum_im_ik_i\le N$. So $\sum_ik_i=1$, and $x(t)$ is a simple root of $g^x$.
\item \emph{Conclusion.} Take $i$ with $k_i>0$. By (3) and (5), Lemma~\ref{sha-automorphy-unramified-algebra-derivation} gives $\delta=0$.
\end{enumerate}
\end{proof}
% unrAlgW (UnrAlgW.lean): primEx, factEx, minPol and the structure Cmp (with intMod); (1) is muGe with rm_map_prod; (2) is hex; (3) is comp_bound (lyOver and Sep); (4) is hle (cmpPrime, not_dvd_der, good_point and muLe); (5) is num_core; (6) is derZero.

\begin{lemma}\label{sha-automorphy-unramified-algebra-image}
Let $\mathcal O$ be a ring, let $\Lambda$ be a local domain and $B$ a ring, both $\mathcal O$-algebras, with an $\mathcal O$-algebra map $\lambda:\Lambda\to B$, let $J$ be an ideal of $B$, $M$ a Hecke module for $(\lambda,J)$ and $e$ an idempotent as in Definition~\ref{sha-automorphy-unramified-multiplicity}(2), and let $K$ be a field that is an $\mathcal O$-algebra. Then there are an ideal $J'\supseteq J$ of $B$ and a Hecke module $M'$ for $(\lambda,J')$ such that for every $\mathcal O$-algebra map $y:B\to K$, the multiplicity $\mu'(y)$ of $y$ on $M'$ is $\mu_e(y)$, and $\mu_e(y)\ge1$ implies that $y$ vanishes on $J'$.
\end{lemma}
\leanref{Development/AutomorphicArithmetic/AutomorphicVarieties/ModularityArgument/CuspidalDescent/GLDenseRanges/InertiaPartMaps/MuRange.lean}{UpperTheory.V.Z.Desc.V0.muRange}
\begin{proof}
\emph{The module.} Let $M'=e(M)$, a $B$-submodule since $e$ commutes with $B$. Let $i:M'\to M$ be the inclusion and $s:M\to M'$ the map $m\mapsto e(m)$; then $s\circ i$ is the identity, as $e^2=e$, and $i\circ s=e$. So $M'$ is a direct summand of $M$ over $\Lambda$, hence finite and projective, hence free since $\Lambda$ is local \cite[Theorem~2.5]{Mat86}.

\emph{The ideal.} Let $J'$ be the radical of the annihilator of $M'$ in $B$. It is radical, and a prime contains the annihilator of $M'$ if and only if it contains $J'$. If $\lambda(c)a\in J'$ with $c\ne0$, then some $(\lambda(c)a)^n$ kills $M'$, so $c^n(a^nm)=0$ for $m\in M'$; as $M'$ is free over the domain $\Lambda$, $a^nm=0$, and $a\in J'$. So $M'$ is a Hecke module for $(\lambda,J')$. By Lemma~\ref{sha-automorphy-unramified-algebra-finite}(1), $J$ lies in the radical of the annihilator of $M$, which lies in $J'$ since $M'\subseteq M$.

\emph{Multiplicities.} Let $y:B\to K$, and view $K$ as a $\Lambda$-algebra through $y\circ\lambda$. The maps $1\otimes i:M'\otimes_\Lambda K\to M\otimes_\Lambda K$ and $1\otimes s$ satisfy $(1\otimes s)(1\otimes i)=1$, so $1\otimes i$ is injective; it commutes with the action of every $a\in B$, and $(1\otimes e)(1\otimes i)=1\otimes i$, $(1\otimes i)(1\otimes s)=1\otimes e$. So $1\otimes i$ maps the joint generalized eigenspace of $M'\otimes_\Lambda K$ for $y$ into the space of vectors of $M\otimes_\Lambda K$ fixed by $1\otimes e$ in the joint generalized eigenspace for $y$. It maps onto that space: such a vector $v$ is $(1\otimes i)(w)$ with $w=(1\otimes s)(v)$, and $(1\otimes i)\bigl((a-y(a))^kw\bigr)=(a-y(a))^kv$, so $w$ lies in the joint generalized eigenspace by injectivity. Hence $\mu'(y)=\mu_e(y)$.

\emph{Vanishing on $J'$.} Let $\mu_e(y)\ge1$, and take $w\ne0$ in the joint generalized eigenspace of $M'\otimes_\Lambda K$ for $y$. For $a\in J'$, some $a^n$ kills $M'$, so $a$ acts nilpotently on $M'\otimes_\Lambda K$. If $y(a)\ne0$, then $a-y(a)=y(a)\bigl(y(a)^{-1}a-1\bigr)$ is invertible there, against $(a-y(a))^kw=0$. So $y(a)=0$.
\end{proof}
% muRange (MuRange.lean): P = range of e with the splitting i, s (Module.Projective.of_split, Module.free_of_flat_of_isLocalRing), J' = radical of the annihilator of P, the equality key of the two eigenspaces, and the nilpotence of G a for the vanishing.

\begin{lemma}\label{sha-automorphy-unramified-algebra}
Let $\mathcal O$ be a ring, let $\Lambda$ be a noetherian integrally closed local domain of characteristic $0$ and $B$ a ring, both $\mathcal O$-algebras, with an $\mathcal O$-algebra map $\lambda:\Lambda\to B$. Let $J$ be an ideal of $B$, $M$ a Hecke module for $(\lambda,J)$ and $e$ an idempotent as in Definition~\ref{sha-automorphy-unramified-multiplicity}(2). Let $K$ be an algebraically closed $\mathcal O$-algebra of characteristic $0$, let $x:B\to K$ be an $\mathcal O$-algebra map, let $\mathcal G$ be a set of $\mathcal O$-algebra maps $B\to K$, and let $b_0\in B$ with $x(b_0)\ne0$. Suppose that every $a\in B$ that vanishes at every $y\in\mathcal G$ lies in $J$, that $\mu_e(x)\ge1$, and that $\mu_e(x)\le\mu_e(y)$ for every $y\in\mathcal G$ with $y(b_0)\ne0$. Then every tangent vector $\delta$ of $B$ at $x$ with values in $K$, that is, every $\mathcal O$-linear $\delta:B\to K$ with $\delta(ab)=x(a)\delta(b)+x(b)\delta(a)$, that vanishes on $J$ and on $\lambda(\Lambda)$ is $0$.
\end{lemma}
\leanref{Development/AutomorphicArithmetic/AutomorphicVarieties/ModularityArgument/CuspidalDescent/GLDenseRanges/MuRange.lean}{UpperTheory.V.Z.Desc.V0.UnrAlgE}
\leanref{Development/AutomorphicArithmetic/AutomorphicVarieties/ModularityArgument/CuspidalDescent/GLDenseRanges/InertiaPartMaps/UnrAlgE.lean}{UpperTheory.V.Z.Desc.V0.unrAlgE}
\begin{proof}
Let $J'\supseteq J$ and the Hecke module $M'$ for $(\lambda,J')$ be given by Lemma~\ref{sha-automorphy-unramified-algebra-image}, with multiplicities $\mu'$ on $M'$.
If $a\in B$ vanishes at every $y\in\mathcal G$ with $y(b_0)\ne0$, then $ab_0$ vanishes on all of $\mathcal G$, so $ab_0\in J$.
Since $\mu_e(x)\ge1$, $x$ vanishes on $J'$; and every $y\in\mathcal G$ with $y(b_0)\ne0$ has $\mu_e(y)\ge\mu_e(x)\ge1$, so it vanishes on $J'$, and $\mu'(x)=\mu_e(x)\le\mu_e(y)=\mu'(y)$.
So every $a\in J'$ vanishes at these $y$, and $ab_0\in J$; then $0=\delta(ab_0)=x(a)\delta(b_0)+x(b_0)\delta(a)=x(b_0)\delta(a)$, so $\delta$ vanishes on $J'$.
Now Lemma~\ref{sha-automorphy-unramified-algebra-rigid}, applied to $J'$, $M'$, $x$, the set $\mathcal G'$ of $y\in\mathcal G$ with $y(b_0)\ne0$, and $b_0$, gives $\delta=0$: an $a$ that vanishes on $\mathcal G'$ has $ab_0\in J\subseteq J'$.
\end{proof}
% unrAlgE (UnrAlgE.lean) from muRange (Lemma sha-automorphy-unramified-algebra-image) and unrAlgW (Lemma sha-automorphy-unramified-algebra-rigid).

\begin{lemma}\label{sha-automorphy-unramified-tangent}
Let a residual datum of rank $n\ge1$ be given in Situation~\ref{sha-hecke-situation}, with its universal ordinary determinant ring $B$, and suppose that its Iwasawa algebra $\Lambda$ is a power series ring over $\mathcal O$ in finitely many variables. Let $\iota:\bar{\mathbf Q}_\ell\cong\mathbf C$, and let $J$ be an ideal of $B$ with the density (2) and the purity (4) of Proposition~\ref{sha-hecke-ideal} that satisfies the multiplicity bound (Definition~\ref{sha-automorphy-unramified-multiplicity}). Then at every irreducible arithmetic point $x$ of $B$ on the zero locus of $J$, every tangent vector of $B$ at $x$ with values in $\bar{\mathbf Q}_\ell$ that kills $J$ and the universal characters on the inertia groups above $\ell$ is $0$.
\end{lemma}
\leanref{Development/AutomorphicArithmetic/AutomorphicVarieties/ModularityArgument/ThresholdArguments/DenseThresholdCases/LevelRelationTwoX.lean}{UpperTheory.V.J.UnramifiedEndsStatement}
\leanref{Development/AutomorphicArithmetic/AutomorphicVarieties/ModularityArgument/ThresholdArguments/DenseThresholdCases/ChainRankFourGood.lean}{UpperTheory.V.J.unramified_algebraic_weights}
\begin{proof}
Let $x$ be an irreducible point on the zero locus of $J$, arithmetic of an increasing weight $h$, and $\delta$ a tangent vector of $B$ at $x$ with values in $\bar{\mathbf Q}_\ell$ that kills $J$ and the $\psi^{\mathrm u}_{v,i}(g)$ for $g\in I_v$. Let $\lambda:\Lambda\to B$ be the structure map.
\begin{enumerate}
\item \emph{Irreducibility is open.} Let $\rho$ be a continuous irreducible representation with determinant $x(D^{\mathrm u})$. By Burnside's theorem the $\rho(g)$ span the matrices, so the trace form gives $g_1,\dots,g_{n^2}$ with $\det\big(\operatorname{tr}\rho(g_ag_b)\big)\ne0$. Let $b_0=\det\big(\operatorname{tr}_{D^{\mathrm u}}(g_ag_b)\big)\in B$; then $x(b_0)\ne0$. If $y$ is a point of distinct type of weight $h$ on the zero locus of $J$ with $y(b_0)\ne0$, then $y$ is arithmetic, so by purity $y(D^{\mathrm u})$ is the determinant of a continuous $\rho'$, and the same determinant for $\rho'$ is $y(b_0)\ne0$; so $\rho'$ is irreducible, and $y$ is irreducible.
\item \emph{$\delta$ vanishes on $\lambda(\Lambda)$.} The maps $x\circ\lambda$ and $x\circ\lambda+\epsilon\,\delta\circ\lambda$ from $\Lambda$ to a ring of dual numbers $T[\epsilon]$, over a complete local subring $T$ of $\bar{\mathbf Q}_\ell$ that contains their values, carry the universal characters of $\Lambda$ to the same characters, since $\delta$ kills the $\psi^{\mathrm u}_{v,i}$ on inertia. By the universal property of $\Lambda$ they are equal, so $\delta\circ\lambda=0$.
\item \emph{A point of distinct type.} Let $\mathcal G$ be the set of points of distinct type of weight $h$ on the zero locus of $J$. By the density (2), an element of $B$ that vanishes at every $y\in\mathcal G$ lies in $J$; as $b_0\notin J\subset\ker x$, some $y_0\in\mathcal G$ has $y_0(b_0)\ne0$, and $y_0$ is irreducible by (1).
\item \emph{Multiplicities.} Let $M$ and $e$ be given by the multiplicity bound. Applied to $x$ and $y_0$ it gives $\mu_e(x)\ge1$, and applied to $x$ and each $y\in\mathcal G$ with $y(b_0)\ne0$, which is irreducible by (1), it gives $\mu_e(x)\le\mu_e(y)$.
\item \emph{Conclusion.} The ring $\Lambda$, a power series ring over the discrete valuation ring $\mathcal O$ of characteristic $0$, is a regular local ring of characteristic $0$, hence a noetherian integrally closed local domain \cite[\S19]{Mat86}. Lemma~\ref{sha-automorphy-unramified-algebra}, with $K=\bar{\mathbf Q}_\ell$, $\mathcal G$ and $b_0$, gives $\delta=0$.
\end{enumerate}
\end{proof}
% unramified_algebraic_weights = Z.Desc.V0.unramified_at_of_dense_purity_num_e_power_series (UnrE.lean): (1) J.irrOpenW_alt (g_gram, gram_at, GalN.V.end1W' for purity); (2) J.lamDer (exists_pair_of_der and the field univ of LamObj); (3) the hypothesis B.Dense; (4) the hypothesis Num (NumE); (5) isDomain_mvPowerSeries, isIntegrallyClosed_mvPowerSeries and unrAlgE (from unrAlgW and muRange). The conclusion is J.UnrAt.

\begin{lemma}\label{sha-automorphy-unramified-near-dense}
Let $\mathcal O\subset\bar{\mathbf Q}_\ell$ be the ring of integers of a finite extension of $\mathbf Q_\ell$, let $B_0$ be a noetherian $\mathcal O$-algebra, $B$ an $\mathcal O$-algebra and $\phi:B_0\to B$ a surjective $\mathcal O$-algebra map. Let $J$ be an ideal of $B_0$, $D\ge0$ an integer and $\mathcal G$ a set of $\mathcal O$-algebra maps $B_0\to\bar{\mathbf Q}_\ell$. Suppose that $B_0/Z_0$ has Krull dimension $D$ for every minimal prime $Z_0$ over $J$, and that every $b\in B_0$ that vanishes at every $x\in\mathcal G$ vanishing on $J$ lies in $J$. Then every prime $Z$ of $B$ that contains $\phi(J)$ has $\dim B/Z\le D$; and if $\dim B/Z=D$, then for every $b\in B$ outside $Z$ there is an $\mathcal O$-algebra map $y:B\to\bar{\mathbf Q}_\ell$ that vanishes on $Z$, with $y\circ\phi\in\mathcal G$ and $y(b)\ne0$.
\end{lemma}
\leanref{Development/AutomorphicArithmetic/AutomorphicVarieties/ModularityArgument/LinkedUnramifiedChains/ScalarIrreducibility/DenseOf.lean}{UpperTheory.V.J.denseOf}
\begin{proof}
Let $Z'=\phi^{-1}(Z)$, a prime of $B_0$ that contains $J$; since $\phi$ is surjective, $B/Z\cong B_0/Z'$. Let $Z_0\subseteq Z'$ be a minimal prime over $J$. Then $\dim B_0/Z'\le\dim B_0/Z_0=D$.
Suppose that $\dim B/Z=D$. Then $Z_0=Z'$: otherwise some $r\in Z'$ has nonzero image in the domain $B_0/Z_0$, a chain of primes of $B_0/Z_0$ containing that image extends by the prime $0$, and $B_0/Z'$ is a quotient of $B_0/(Z_0+rB_0)$, so $\dim B_0/Z'+1\le D$.
Let $b=\phi(b')$; then $b'\notin Z_0$. The minimal primes over $J$ are finite in number, as $B_0$ is noetherian, so by prime avoidance some $e\notin Z_0$ lies in every other minimal prime over $J$. Then $b'e\notin Z_0$, so $b'e\notin J$, and by the hypothesis some $x\in\mathcal G$ vanishing on $J$ has $x(b'e)\ne0$. The prime $\ker x$ contains a minimal prime $Z_1$ over $J$; since $x(e)\ne0$, $Z_1=Z_0$. So $\ker\phi\subseteq Z'=Z_0\subseteq\ker x$, and $x=y\circ\phi$ for an $\mathcal O$-algebra map $y:B\to\bar{\mathbf Q}_\ell$, which vanishes on $Z$ and has $y(b)=x(b')\ne0$.
\end{proof}
% denseOf (DenseOf.lean): dim_succ_le (ringKrullDim_quotient_succ_le_of_nonZeroDivisor) and exists_sep.

\begin{lemma}\label{sha-automorphy-unramified-near-additive}
Let $\ell$ be a prime, $H$ a compact topological group, $u_1,\dots,u_m$ and $g$ elements of $H$ such that every continuous homomorphism from $H$ to a finite abelian $\ell$-group that kills $u_1,\dots,u_m$ kills $g$, and $f:H\to\bar{\mathbf Q}_\ell$ a continuous map with $f(hh')=f(h)+f(h')$ for all $h,h'$ and $f(u_a)=0$ for all $a$. Then $f(g)=0$.
\end{lemma}
\begin{proof}
Let $\varepsilon>0$. As the absolute value of $\bar{\mathbf Q}_\ell$ is ultrametric, the elements of absolute value less than $\varepsilon$ form a subgroup $W$ of $\bar{\mathbf Q}_\ell$. Let $\chi:H\to f(H)/(f(H)\cap W)$ be the composite of $f$ with the quotient map, a surjective homomorphism. Its kernel, the set of $h$ with $|f(h)|<\varepsilon$, is open since $f$ is continuous, so it has finite index in the compact group $H$, and the image of $\chi$ is finite. It is an $\ell$-group: for $h\in H$ and $k$ large, $|f(h^{\ell^k})|=|\ell^kf(h)|=\ell^{-k}|f(h)|<\varepsilon$. So $\chi$ is a continuous homomorphism to a finite abelian $\ell$-group that kills the $u_a$, hence $g$; that is, $|f(g)|<\varepsilon$. As $\varepsilon$ is arbitrary, $f(g)=0$.
\end{proof}
% The private add_vanish of TanBound.lean.

\begin{lemma}\label{sha-automorphy-unramified-near-tangent-points}
Let $\mathcal O\subset\bar{\mathbf Q}_\ell$ be the ring of integers of a finite extension of $\mathbf Q_\ell$, let $R$ be a complete local noetherian $\mathcal O$-algebra whose residue field is that of $\mathcal O$, and let $y:R\to\bar{\mathbf Q}_\ell$ be an $\mathcal O$-algebra map. Then $y(R)$ is a finitely generated $\mathcal O$-module, $|y(a)|\le1$ for every $a\in R$, and $|y(a)|<1$ for every $a\in\mathfrak m_R$.
\end{lemma}
\leanref{Development/AutomorphicArithmetic/IntegralCasts/CapelliIdentities/CoPoint.lean}{UpperTheory.Reformulated.B.CoPoint.coPt}
\leanref{Development/AutomorphicArithmetic/IntegralCasts/CapelliIdentities/ResidueHomNorm.lean}{UpperTheory.residue_hom_norm}
\begin{proof}
\emph{Absolute values.} The elements of $\mathcal O$ have absolute value at most $1$, and an element of $\bar{\mathbf Q}_\ell$ of absolute value at most $1$ is integral over $\mathcal O$. Let $m\in\mathfrak m_R$, and suppose that $|y(m)|\ge1$. Then $y(m)^{-1}$ is a root of a monic $p\in\mathcal O[X]$ of some degree $e$, and $q(X)=X^ep(X^{-1})$ is a polynomial with constant term $1$ and $q(y(m))=0$. As $q(m)\equiv1$ modulo $m$, $q(m)$ is a unit of $R$; but $y(q(m))=q(y(m))=0$. So $|y(m)|<1$. Every $a\in R$ is $c+m$ with $c\in\mathcal O$ and $m\in\mathfrak m_R$, since the residue fields agree; so $|y(a)|\le\max(|c|,|y(m)|)\le1$.

\emph{Finiteness.} Let $\varpi$ be a uniformizer of $\mathcal O$ and $I$ the ideal generated by $\ker y$ and $\varpi$. Each $y(a)$ has absolute value at most $1$, so it is a root of a monic $p_a\in\mathcal O[X]$, and $p_a(a)\in\ker y$. So $R/I$ is integral over the field $\mathcal O/\varpi$, every prime of the noetherian local ring $R/I$ is maximal, and $R/I$ is artinian with finite residue field, hence finite. Let $E\subset R$ be a finite set that maps onto $R/I$. Then every $s\in y(R)$ is $y(e)+\varpi s'$ with $e\in E$ and $s'\in y(R)$. Repeating this, $s-\sum_{k<K}\varpi^ky(e_k)\in\varpi^Ky(R)$ for every $K$, with $e_k\in E$. The elements of $\varpi^Ky(R)$ have absolute value at most $|\varpi|^K$, which tends to $0$, and the partial sums lie in the $\mathcal O$-module spanned by $y(E)$, which is finitely generated over the complete ring $\mathcal O$ and so closed in $\bar{\mathbf Q}_\ell$. So $s$ lies in it, and $y(R)$ is spanned by $y(E)$.
\end{proof}
% coPt (CoPoint.lean): integral_dvr_norm, residue_hom_norm (the reverse polynomial), monic_ideal_completion (R/(ker y + varpi) is finite) and digit_expansion_lattice.

\begin{lemma}\label{sha-automorphy-unramified-near-tangent-small}
Let $\mathcal O$, $R$ and $y$ be as in Lemma~\ref{sha-automorphy-unramified-near-tangent-points}, and let $D:R\to\bar{\mathbf Q}_\ell$ be an $\mathcal O$-linear map with $D(ab)=y(a)D(b)+y(b)D(a)$ for all $a,b\in R$. Then for every $\varepsilon>0$ there is $k\ge0$ such that $|D(a)|\le\varepsilon$ and $|y(a)|\le\varepsilon$ for every $a\in\mathfrak m_R^k$.
\end{lemma}
\leanref{Development/AutomorphicArithmetic/AutomorphicVarieties/ModularityArgument/LinkedUnramifiedChains/ScalarIrreducibility/DerCont.lean}{UpperTheory.V.J.derCont}
\leanref{Development/AutomorphicArithmetic/IntegralCasts/CapelliIdentities/ExistsSpanModSq.lean}{UpperTheory.Reformulated.B.Ra7.exists_span_mod_sq}
\begin{proof}
Let $K=\ker y$ and let $\varpi$ be a uniformizer of $\mathcal O$.
\begin{enumerate}
\item \emph{Vanishing on $K^2$.} For $f,g\in K$, $D(fg)=y(f)D(g)+y(g)D(f)=0$ and $y(fg)=0$. So $D$ and $y$ vanish on $K^2$.
\item \emph{A finite set.} By Lemma~\ref{sha-automorphy-unramified-near-tangent-points}, $y(R)$ is spanned over $\mathcal O$ by some $y(r_1),\dots,y(r_p)$, and $K$ is generated by some $k_1,\dots,k_m$, as $R$ is noetherian. Let $s$ be the finite set of the $r_t$ and the $r_tk_j$. For $a\in R$ there are $c_t\in\mathcal O$ with $a-\sum_tc_tr_t\in K$, say $a-\sum_tc_tr_t=\sum_jb_jk_j$; and there are $c_{jt}\in\mathcal O$ with $b_j-\sum_tc_{jt}r_t\in K$, so that $b_jk_j\equiv\sum_tc_{jt}r_tk_j$ modulo $K^2$. So every element of $R$ is congruent modulo $K^2$ to an $\mathcal O$-linear combination of $s$.
\item \emph{A bound.} Let $C=\max_{e\in s}|D(e)|$. By (1) and (2), as the elements of $\mathcal O$ have absolute value at most $1$ and the absolute value is ultrametric, $|D(a)|\le C$ for every $a\in R$.
\item \emph{A power of $\mathfrak m_R$.} By (2), the local ring $R/(\varpi R+K^2)$ is spanned over the finite field $\mathcal O/\varpi$ by the images of $s$, so it is finite, hence artinian, and its maximal ideal is nilpotent. So $\mathfrak m_R^j\subseteq\varpi R+K^2$ for some $j$.
\item \emph{Conclusion.} For $n\ge1$, $\mathfrak m_R^{jn}\subseteq(\varpi R+K^2)^n\subseteq\varpi^nR+K^2$. For $a=\varpi^nb+c$ with $c\in K^2$, (1) gives $D(a)=\varpi^nD(b)$ and $y(a)=\varpi^ny(b)$, so $|D(a)|\le|\varpi|^nC$ and $|y(a)|\le|\varpi|^n$ by (3) and Lemma~\ref{sha-automorphy-unramified-near-tangent-points}. As $|\varpi|<1$, some $n$ has $|\varpi|^n\max(C,1)\le\varepsilon$; take $k=jn$.
\end{enumerate}
\end{proof}
% derCont = derCont_abs (DerCont.lean): coPt; hDq and hyq for (1); exists_span_mod_sq for (2); hbd for (3); exists_pow_le (R/(varpi R + K^2) is finite and artinian) for (4); sup_pow_le and the choice of n for (5).

\begin{lemma}\label{sha-automorphy-unramified-near-tangent-filtration}
Let $X$ be a finitely generated free abelian group, and let $R$ be a complete local noetherian $\mathcal O$-algebra with residue field $\mathbf F$, with a weighted presentation $\pi$ by variables $X_s$ of weights $w_s\in X$ (Definition~\ref{sha-counts-vd-model}(1)) such that $z_s\in\mathfrak m_R$ for all $s$. Let $\mathfrak m_0=R_0\cap\mathfrak m_R$, the maximal ideal of $R_0$ (Lemma~\ref{sha-counts-vd-presentations}(c)). Then for every $w\in X$ and every $N\ge0$ there is $M$ such that $R_w\cap\mathfrak m_R^M\subseteq\mathfrak m_0^NR_w$.
\end{lemma}
\leanref{Development/AutomorphicArithmetic/AutomorphicVarieties/CommutativeAlgebra/PartFilt.lean}{UpperTheory.V.Alg.part_filt}
\begin{proof}
For an exponent vector $n$ write $|n|$ for its degree, $X^n=\prod_sX_s^{n_s}$ and $z^n=\prod_sz_s^{n_s}$.
\begin{enumerate}
\item \emph{Splitting monomials.} By Gordan's lemma the monoid of exponent vectors of weight $0$ is generated by finitely many nonzero vectors, and by Dickson's lemma the exponent vectors of weight $w$ that are minimal for the componentwise order form a finite set, as in the proof of Lemma~\ref{sha-counts-vd-presentations}(b) and (c). Let $D$ be the largest degree of these finitely many vectors. An exponent vector $n$ of weight $w$ with $|n|>D$ is $m+r$ with $m$ minimal of weight $w$ and $r\ne0$ of weight $0$, and $r$ is at least one of the generators $g$; so $n=g+n'$ with $g\ne0$ of weight $0$, $|g|\le D$ and $n'$ of weight $w$.
\item \emph{Monomials of large degree.} For every $k$ there is $N_k$ such that $z^n\in\mathfrak m_0^kR_w$ for every exponent vector $n$ of weight $w$ with $|n|\ge N_k$. For $k=0$ take $N_0=0$. Given $N_k$, let $N_{k+1}=N_k+D+1$ and write $n=g+n'$ as in (1); then $z^g\in\mathfrak m_0$, since $g\ne0$ has weight $0$ and the $z_s$ lie in $\mathfrak m_R$, and $|n'|\ge N_k$, so $z^n=z^gz^{n'}\in\mathfrak m_0^{k+1}R_w$.
\item \emph{The bound.} Let $M=N_N+N$ and $x\in R_w\cap\mathfrak m_R^M$. Write $x=\pi(f)$ with $f$ of weight $w$. As $\pi$ is surjective and local, $\mathfrak m_R^M$ is the image of the $M$-th power of the maximal ideal of the power series ring, so $x=\pi(g)$ for a series $g$ in that power, whose coefficient at $X^n$ lies in $\mathfrak m_{\mathcal O}^{M-|n|}$. Let $h$ be the sum of the terms of $g$ of weight $w$. The terms of weight $w$ of $f-g\in\ker\pi$ form $f-h$, which lies in $\ker\pi$; so $x=\pi(h)$. Let $h_<$ be the sum of the terms of $h$ of degree less than $N_N$. Each of them is $cX^n$ with $c\in\mathfrak m_{\mathcal O}^N$ and $X^n$ of weight $w$, and maps to $cz^n\in\mathfrak m_0^NR_w$. All monomials of $h-h_<$ have weight $w$ and degree at least $N_N$, so by Lemma~\ref{sha-counts-vd-presentations}(b) $h-h_<=\sum_{n\in\mu}b_nX^n$ with $b_n$ of weight $0$ and $\mu$ a finite set of such monomials; and $\pi(h-h_<)=\sum_n\pi(b_n)z^n\in R_0\mathfrak m_0^NR_w$ by (2). So $x\in\mathfrak m_0^NR_w$.
\end{enumerate}
\end{proof}
% part_filt (PartFilt.lean): part_filt_split (from gordan) for (1), hmono for (2), hconst for the constants, part_filt_coeff, W.homog with wcomp, the truncation low and part_span for (3). The Lean's bound is Q^N * R_0 * R_w with Q = R_0 meet m_R, which is the same ideal.

\begin{lemma}\label{sha-automorphy-unramified-near-tangent-adic}
In Situation~\ref{sha-automorphy-link-situation}, let $n\in\{2,3\}$, let $\mathcal O_1\supseteq\mathcal O$ be the ring of integers of a finite extension of $\mathbf Q_\ell$, and let $A_n$, its universal linked point $(\rho^{\mathrm u},\psi^{\mathrm u})$ and $B$ be as in Definition~\ref{sha-automorphy-realization}(2), with $\mathfrak m_B$ the maximal ideal of $B$. Then for every $N$ there is $M$ with $B\cap\mathfrak m_{A_n}^M\subseteq\mathfrak m_B^N$. In particular, for every place $v$, every $i$ and every $N$, the map $h\mapsto\psi^{\mathrm u}_{v,i}(h)\bmod\mathfrak m_B^N$ from $G_v$ to $B/\mathfrak m_B^N$ is locally constant.
\end{lemma}
\leanref{Development/AutomorphicArithmetic/AutomorphicVarieties/FrobeniusSmoothVectors/MemPowSmulLe.lean}{UpperTheory.V.Fr.cont_u}
\begin{proof}
By Lemma~\ref{sha-automorphy-realization-presentation}, $A_n$ has a weighted presentation with weights in $\mathbf Z$ whose variables map into $\mathfrak m_{A_n}$, with $(A_n)_0=B$, whose maximal ideal is $\mathfrak m_B=B\cap\mathfrak m_{A_n}$, and the values of the $\psi^{\mathrm u}_{v,i}$ lie in $B$. Lemma~\ref{sha-automorphy-unramified-near-tangent-filtration}, with $\mathcal O_1$ and $k_1$ in place of $\mathcal O$ and $\mathbf F$ and with $w=0$, gives $M$ with $B\cap\mathfrak m_{A_n}^M\subseteq\mathfrak m_B^NB=\mathfrak m_B^N$. By condition (d) of Definition~\ref{sha-descent-link}(2), $h\mapsto\psi^{\mathrm u}_{v,i}(h)\bmod\mathfrak m_{A_n}^M$ is locally constant; and if $\psi^{\mathrm u}_{v,i}(h)-\psi^{\mathrm u}_{v,i}(h')\in\mathfrak m_{A_n}^M$, this difference lies in $B\cap\mathfrak m_{A_n}^M\subseteq\mathfrak m_B^N$.
\end{proof}
% cont_u (MemPowSmulLe.lean): wpres, part_filt with j = 0, hmulB and hP, and hcB for the field AdicCont S.mB (S.psiB v i); cone_u (ArtinRees.lean) passes it on as the field psi of Cone.

\begin{lemma}\label{sha-automorphy-unramified-near-tangent}
In the situation of Lemma~\ref{sha-automorphy-unramified-realization}, let $\xi':A_3^\circ\to\bar{\mathbf Q}_\ell$ be an $\mathcal O_1$-algebra map with values of absolute value at most $1$ and $|\xi'(b)|<1$ for $b$ in the maximal ideal of $B$. Let $T$ be the $\bar{\mathbf Q}_\ell$-vector space of tangent vectors of $A_3^\circ$ at $\xi'$ with values in $\bar{\mathbf Q}_\ell$, that is, of $\mathcal O_1$-linear maps $D:A_3^\circ\to\bar{\mathbf Q}_\ell$ with $D(ab)=\xi'(a)D(b)+\xi'(b)D(a)$, and $T_0\subseteq T$ the subspace of those that vanish at $\psi^{\mathrm u}_{v,i}(h)$ for all $v$, $i$ and $h\in I_v$. Then $\dim T\le3d+\dim T_0$.
\end{lemma}
\leanref{Development/AutomorphicArithmetic/AutomorphicVarieties/ModularityArgument/LinkedUnramifiedChains/ScalarIrreducibility/TanBound.lean}{UpperTheory.V.J.tanBound}
\leanref{Development/AutomorphicArithmetic/AutomorphicVarieties/ModularityArgument/LinkedUnramifiedChains/RedDefs.lean}{UpperTheory.V.J.FixTan}
\leanref{Development/AutomorphicArithmetic/AutomorphicVarieties/ModularityArgument/LinkedUnramifiedChains/ScalarIrreducibility/DerRep.lean}{UpperTheory.V.J.derRep}
\begin{proof}
For each place $v$ of $L^+$ above $\ell$, residual condition (8) of Definition~\ref{sha-descent-residual} gives $u_{v,1},\dots,u_{v,d_v}\in I_v$ such that every continuous homomorphism from $G_v$ to a finite abelian $\ell$-group that kills them kills $I_v$.
\begin{enumerate}
\item \emph{Continuity.} For $D\in T$ and all $v$ and $i$, the maps $h\mapsto\xi'(\psi^{\mathrm u}_{v,i}(h))$ and $h\mapsto D(\psi^{\mathrm u}_{v,i}(h))$ on $G_v$ are continuous. Indeed, the restrictions of $\xi'$ and $D$ to $B$ are an $\mathcal O_1$-algebra map $y:B\to\bar{\mathbf Q}_\ell$ and an $\mathcal O_1$-linear map with $D(ab)=y(a)D(b)+y(b)D(a)$ for $a,b\in B$, and $B$ is a complete local noetherian $\mathcal O_1$-algebra with the residue field of $\mathcal O_1$ (Lemma~\ref{sha-automorphy-realization-presentation}(3)). Let $\varepsilon>0$ and $h\in G_v$. Lemma~\ref{sha-automorphy-unramified-near-tangent-small} gives $k$ with $|D(b)|\le\varepsilon$ and $|\xi'(b)|\le\varepsilon$ for $b\in\mathfrak m_B^k$, and by Lemma~\ref{sha-automorphy-unramified-near-tangent-adic} the $h'\in G_v$ with $\psi^{\mathrm u}_{v,i}(h')-\psi^{\mathrm u}_{v,i}(h)\in\mathfrak m_B^k$ form a neighbourhood of $h$. On it, $D(\psi^{\mathrm u}_{v,i}(h'))$ and $\xi'(\psi^{\mathrm u}_{v,i}(h'))$ differ from their values at $h$ by at most $\varepsilon$.
% derRep (DerRep.lean), its last part: der_small from derCont (Lemma sha-automorphy-unramified-near-tangent-small) and char_cont' with the field psi of cone_u, which is cont_u (Lemma sha-automorphy-unramified-near-tangent-adic). The bound der_bdd and entry_cont, for the entries of the universal linked point, are not used here.
\item \emph{Logarithmic derivatives.} Let $D\in T$, and fix $v$ and $i$. Each $\psi^{\mathrm u}_{v,i}(h)$ is a unit, so $\xi'(\psi^{\mathrm u}_{v,i}(h))\ne0$; let $f(h)=D(\psi^{\mathrm u}_{v,i}(h))/\xi'(\psi^{\mathrm u}_{v,i}(h))$ for $h\in G_v$. By (1), $f$ is continuous, and by the Leibniz rule $f(hh')=f(h)+f(h')$.
\item \emph{The count.} Consider the linear map $T\to\bar{\mathbf Q}_\ell^{3d}$, $D\mapsto\big(D(\psi^{\mathrm u}_{v,i}(u_{v,a}))\big)_{v,i,a}$, whose $3\sum_vd_v=3d$ coordinates are indexed by $v$, $1\le i\le3$ and $1\le a\le d_v$. Let $D$ be in its kernel, and fix $v$ and $i$. The function $f$ of (2) vanishes at the $u_{v,a}$, and $G_v$ is compact; so, by the choice of the $u_{v,a}$ and Lemma~\ref{sha-automorphy-unramified-near-additive}, $f$ vanishes on $I_v$, and $D(\psi^{\mathrm u}_{v,i}(g))=0$ for $g\in I_v$. So the kernel lies in $T_0$, and $\dim T\le3d+\dim T_0$.
\end{enumerate}
\end{proof}
% tanBound (TanBound.lean): hF.gen for the u_{v,a} (residual condition (8)), rank_le_card_add_rank_of_kernel for (3), derRep for (1), add_vanish (Lemma sha-automorphy-unramified-near-additive).

\subsection{The ideal near small points}\label{subsection-sha-near-small-points}

\begin{lemma}\label{sha-automorphy-unramified-near-scale}
In Situation~\ref{sha-automorphy-link-situation}, let $n\in\{2,3\}$, let $\mathcal O_1\supseteq\mathcal O$ be the ring of integers of a finite extension of $\mathbf Q_\ell$, and let $A_n$, its universal linked point $(\rho^{\mathrm u},\psi^{\mathrm u})$, $B$ and $A_n^\circ$ be as in Definition~\ref{sha-automorphy-realization}(2), with the parts $(A_n)_j$ of Lemma~\ref{sha-automorphy-realization-presentation}. Let $\xi:A_n^\circ\to\bar{\mathbf Q}_\ell$ be an $\mathcal O_1$-algebra map. Then there are a nonzero $u\in\bar{\mathbf Q}_\ell$ and an $\mathcal O_1$-algebra map $\xi':A_n^\circ\to\bar{\mathbf Q}_\ell$ with $\xi'(a)=u^j\xi(a)$ for every $j$ and every $a\in(A_n)_j$, such that $|\xi'(a)|\le1$ for every $a\in A_n^\circ$, $|\xi'(b)|<1$ for $b$ in the maximal ideal of $B$, and $|\xi'(\rho^{\mathrm u}(g)_{ij})|<1$ for all $g\in\Gamma_0$ and $i\ne j$.
\end{lemma}
\leanref{Development/AutomorphicArithmetic/AutomorphicVarieties/ModularityArgument/LinkedUnramifiedChains/ScalarIrreducibility/ExistsBetweenFinset.lean}{UpperTheory.V.J.scale}
\begin{proof}
Write $w_i=w_{n,i}$.
\begin{enumerate}
\item \emph{Values on $B$.} By Lemma~\ref{sha-automorphy-realization-presentation}(3) and Lemma~\ref{sha-automorphy-unramified-near-tangent-points}, applied to the restriction of $\xi$ to $B$, $|\xi(b)|\le1$ for $b\in B$ and $|\xi(b)|<1$ for $b\in\mathfrak m_B$.
\item \emph{Products of weight $0$.} Let $x$ be a product of entries $\rho^{\mathrm u}(g_t)_{i_tj_t}$ with $\sum_t(w_{i_t}-w_{j_t})=0$ and $i_t\ne j_t$ for some $t$. Then $x\in B$ by Definition~\ref{sha-automorphy-realization}(2), and $x\in\mathfrak m_{A_n}$, since the entries off the diagonal lie in $\mathfrak m_{A_n}$ by condition (e) of Definition~\ref{sha-descent-link}(2). So $x\in\mathfrak m_B=B\cap\mathfrak m_{A_n}$, and $|\xi(x)|<1$ by (1).
\item \emph{Largest entries.} Fix $i\ne j$. The $B$-submodule spanned by the entries $\rho^{\mathrm u}(g)_{ij}$ lies in $(A_n)_{w_i-w_j}$, which is a finitely generated module over the noetherian ring $B$ (Lemma~\ref{sha-automorphy-realization-presentation}(1), (3) and (4)); so it is spanned by finitely many of these entries. Among them choose one, $\rho^{\mathrm u}(g_{ij})_{ij}$, at which $|\xi|$ is largest, and let $a_{ij}$ be that value. As $|\xi|\le1$ on $B$ and the absolute value is ultrametric, $|\xi(\rho^{\mathrm u}(g)_{ij})|\le a_{ij}$ for every $g$.
\item \emph{The size of $u$.} Let $p=w_i-w_j>0$ and $q=w_k-w_m<0$. The product of $-q$ copies of $\rho^{\mathrm u}(g_{ij})_{ij}$ and $p$ copies of $\rho^{\mathrm u}(g_{km})_{km}$ has total weight $0$, so (2) gives $a_{ij}^{-q}a_{km}^p<1$. If $a_{ij}$ and $a_{km}$ are nonzero, this says $-\log a_{km}/q<-\log a_{ij}/p$. So there are real numbers $\alpha<\beta$ with $-\log a_{km}/(w_k-w_m)\le\alpha$ whenever $w_k<w_m$ and $a_{km}\ne0$, and $\beta\le-\log a_{ij}/(w_i-w_j)$ whenever $w_i>w_j$ and $a_{ij}\ne0$. A root $u$ of $X^e-\ell^{-c}$, for integers $c$ and $e\ge1$, has $\log|u|=(c/e)\log\ell$; these values are dense in $\mathbf R$, so some such $u$ has $\alpha<\log|u|<\beta$. Then $|u|^{w_i-w_j}a_{ij}<1$ for all $i\ne j$: this is clear if $a_{ij}=0$, and otherwise it says $(w_i-w_j)\log|u|<-\log a_{ij}$.
\item \emph{The twist.} By Lemma~\ref{sha-automorphy-realization-presentation}(2), $A_n^\circ$ is the direct sum of the $(A_n)_j$, with $(A_n)_j(A_n)_{j'}\subseteq(A_n)_{j+j'}$ and $\mathcal O_1\subseteq B=(A_n)_0$. So the map $\xi'$ equal to $a\mapsto u^j\xi(a)$ on $(A_n)_j$ is an $\mathcal O_1$-algebra map. It agrees with $\xi$ on $B$, and for $i\ne j$, $|\xi'(\rho^{\mathrm u}(g)_{ij})|=|u|^{w_i-w_j}|\xi(\rho^{\mathrm u}(g)_{ij})|\le|u|^{w_i-w_j}a_{ij}<1$ by (3) and (4).
\item \emph{Absolute values.} The ring $A_n^\circ$ is generated as an $\mathcal O_1$-algebra by $B$ and the entries of $\rho^{\mathrm u}$, and the diagonal entries lie in $B$. By (1) and (5), $|\xi'|\le1$ at these generators, hence on $A_n^\circ$, the absolute value being ultrametric. And $|\xi'(b)|=|\xi(b)|<1$ for $b\in\mathfrak m_B$.
\end{enumerate}
\end{proof}
% scale (ExistsBetweenFinset.lean): coPt for (1), Cone.cyc of cone_u for (2), exists_max_entry with the field fg of cone_u for (3), exists_scale (exists_between_finset, exists_log_norm with x = 1/l) for (4), twist through the decomposition of S.A by S.gr (gr_A) for (5), Algebra.adjoin_induction for (6). The Lean's hypothesis that the kernel of xi is off-triangular (OffTri) is not used.

\begin{lemma}\label{sha-automorphy-unramified-near-generic}
In the situation of Lemma~\ref{sha-automorphy-unramified-near-scale}, let $P$ be a minimal prime of $A_n^\circ$ that is off-triangular (Definition~\ref{sha-automorphy-realization}(4)), and let $\mathcal G$ be a set of $\mathcal O_1$-algebra maps $B\to\bar{\mathbf Q}_\ell$ such that for every $b\in B$ outside $P\cap B$ some $y\in\mathcal G$ vanishes on $P\cap B$ and has $y(b)\ne0$. Then there is an $\mathcal O_1$-algebra map $\xi':A_n^\circ\to\bar{\mathbf Q}_\ell$ whose kernel contains $P$ and is off-triangular, whose restriction to $B$ lies in $\mathcal G$, and such that $|\xi'(a)|\le1$ for every $a\in A_n^\circ$, $|\xi'(b)|<1$ for $b$ in the maximal ideal of $B$, and $|\xi'(\rho^{\mathrm u}(g)_{ij})|<1$ for all $g\in\Gamma_0$ and $i\ne j$.
\end{lemma}
\leanref{Development/AutomorphicArithmetic/AutomorphicVarieties/ModularityArgument/LinkedUnramifiedChains/ScalarIrreducibility/ChevPrime.lean}{UpperTheory.V.J.genPt}
\leanref{Development/AutomorphicArithmetic/AutomorphicVarieties/ModularityArgument/LinkedUnramifiedChains/ScalarIrreducibility/ExtIntegral.lean}{UpperTheory.V.J.chev}
\leanref{Development/AutomorphicArithmetic/AutomorphicVarieties/FrobeniusSmoothVectors/MinHom.lean}{UpperTheory.V.Fr.min_hom}
\begin{proof}
\begin{enumerate}
\item \emph{An element outside $P$.} As $P$ is off-triangular, some entry $x_1=\rho^{\mathrm u}(g_1)_{i_1j_1}$ with $i_1>j_1$ and some entry $x_2=\rho^{\mathrm u}(g_2)_{i_2j_2}$ with $i_2<j_2$ lie outside $P$; if $n=3$, some entry $x_3$ of odd weight lies outside $P$, and if $n=2$ let $x_3=1$. Each $x_t$ lies in one of the parts $(A_n)_j$ (Lemma~\ref{sha-automorphy-realization-presentation}(1)), and $a=x_1x_2x_3\notin P$, as $P$ is prime.
\item \emph{Extension of points.} The ring $A_n^\circ$ is a finitely generated $B$-algebra (Lemma~\ref{sha-automorphy-realization-presentation}(4)), so the domain $A_n^\circ/P$ is finitely generated over its subring $B/(P\cap B)$. By \cite[Proposition~5.23]{AM69} there is $b\in B$ outside $P\cap B$ such that every ring map from $B/(P\cap B)$ to an algebraically closed field that does not vanish at $b$ extends to $A_n^\circ/P$ without vanishing at $a$.
\item \emph{A point.} Let $y\in\mathcal G$ vanish on $P\cap B$ with $y(b)\ne0$. By (2) there is a ring map $\xi_0:A_n^\circ\to\bar{\mathbf Q}_\ell$ that vanishes on $P$, extends $y$ and has $\xi_0(a)\ne0$; it is an $\mathcal O_1$-algebra map, as $y$ is one and $\mathcal O_1\subseteq B$.
\item \emph{Rescaling.} Lemma~\ref{sha-automorphy-unramified-near-scale}, applied to $\xi_0$, gives $u\ne0$ and $\xi'$ with the bounds of the statement and $\xi'=u^j\xi_0$ on $(A_n)_j$. So $\xi'$ agrees with $\xi_0$ on $B=(A_n)_0$, and its restriction to $B$ is $y\in\mathcal G$; and an element of a part vanishes under $\xi'$ if and only if it vanishes under $\xi_0$. In particular $\xi'(x_t)\ne0$ for $t=1,2,3$, so the kernel of $\xi'$ is off-triangular.
\item \emph{$P$ lies in the kernel.} As $A_n^\circ$ is the direct sum of its parts (Lemma~\ref{sha-automorphy-realization-presentation}(2)), the ideal $P^*$ generated by the elements of $P$ that lie in some part is prime \cite[Lemma~1.5.6]{BH93}. It lies in $P$, so $P^*=P$ by the minimality of $P$. An element of $P$ that lies in a part vanishes under $\xi_0$, hence under $\xi'$. So $P\subseteq\ker\xi'$.
\end{enumerate}
\end{proof}
% genPt (ChevPrime.lean): gen_out for (1); chev_prime from chev (ChevalleyRestriction.main_fin, which is [AM69, Proposition 5.23]) with gr_A for (2) and (3); scale (Lemma sha-automorphy-unramified-near-scale) for (4); min_hom (MinHom.lean: the homogeneous core of a minimal prime is a prime contained in it) for (5).

\begin{lemma}\label{sha-automorphy-unramified-near-fibre}
In the situation of Lemma~\ref{sha-automorphy-unramified-near-scale}, let $\eta$ be a prime of $A_n^\circ$ such that some entry $\rho^{\mathrm u}(g)_{ij}$ with $i>j$ and some entry $\rho^{\mathrm u}(g')_{ij}$ with $i<j$ lie outside $\eta$. Then the fraction field of $A_n^\circ/\eta$ has transcendence degree at most $1$ over that of $B/(\eta\cap B)$.
\end{lemma}
\leanref{Development/AutomorphicArithmetic/AutomorphicVarieties/FrobeniusSmoothVectors/TrdegQuotLeOne.lean}{UpperTheory.V.Fr.fibre}
\begin{proof}
Let $D=A_n^\circ/\eta$ and $K$ the image of $B$ in $D$; write $\bar x$ for the image of $x$. The domain $D$ is generated over $K$ by the images of the entries of $\rho^{\mathrm u}$, and the entry $(i,j)$ has weight $w_{n,i}-w_{n,j}$, positive for $i<j$. So there are entries $a$ of weight $p>0$ and $b$ of weight $-q<0$ with $\bar a,\bar b\ne0$. A product of entries of total weight $0$ lies in $B$, so its image lies in $K$. Let $x$ be an entry of weight $-m\le0$. Then $\bar x^p\bar a^m\in K$, and $\bar a^m\ne0$, so $\bar x$ is algebraic over $K(\bar a)$. Let $x$ be an entry of weight $m>0$. Then $\bar x^q\bar b^m$ and $c=\bar b^p\bar a^q$ lie in $K$, $c\ne0$, and $\bar x^{pq}c^m=(\bar x^q\bar b^m)^p\bar a^{qm}$; so $\bar x$ is algebraic over $K(\bar a)$. Hence the fraction field of $D$ is algebraic over $K(\bar a)$.
\end{proof}
% fibre (TrdegQuotLeOne.lean): trdeg_quot_le_one and trdeg_le_one, with alg_of_rel for the two cases. The hypothesis is S.u.OffTri eta.

\begin{lemma}\label{sha-automorphy-unramified-near-completion-flag}
In the situation of Lemma~\ref{sha-automorphy-unramified-near-scale}, let $v$ be a place of $L^+$ above $\ell$. Then the characteristic polynomial of $\rho^{\mathrm u}(\mathrm{loc}_vh)$ is $\prod_i(X-\psi^{\mathrm u}_{v,i}(h))$ for every $h\in G_v$. Moreover, for every local ring $C$ and every ring homomorphism $f:A_n^\circ\to C$ there is $g_v\in GL_n(C)$ such that, for every $h\in G_v$, $g_vf(\rho^{\mathrm u}(\mathrm{loc}_vh))g_v^{-1}$ is upper triangular with diagonal entries $f(\psi^{\mathrm u}_{v,1}(h)),\dots,f(\psi^{\mathrm u}_{v,n}(h))$.
\end{lemma}
\leanref{Development/AutomorphicArithmetic/AutomorphicVarieties/FrobeniusSmoothVectors/ZpowWNe.lean}{UpperTheory.V.Fr.flag}
\leanref{Development/AutomorphicArithmetic/AutomorphicVarieties/FrobeniusSmoothVectors/CharpolyFlag.lean}{UpperTheory.V.Fr.o1}
\begin{proof}
Write $w_i=w_{n,i}$.
\begin{enumerate}
\item \emph{An element of $G_v$.} By residual condition (7) of Definition~\ref{sha-descent-residual} there is $h_0\in G_v$ such that $a=\bar\Psi_\Gamma(\mathrm{loc}_vh_0)$ has $a^2\ne1$. For $i\ne j$, $w_i-w_j$ divides $2$, so $a^{w_i}\ne a^{w_j}$, and the element $\psi^{\mathrm u}_{v,i}(h_0)-\psi^{\mathrm u}_{v,j}(h_0)$ of the local ring $B$, whose residue is $a^{w_i}-a^{w_j}$ by condition (e) of Definition~\ref{sha-descent-link}(2), is a unit of $B$ (Lemma~\ref{sha-automorphy-realization-presentation}(3)).
\item \emph{The characteristic polynomial.} By condition (e), $\rho^{\mathrm u}(\mathrm{loc}_vh_0)$ reduces to the diagonal matrix with entries $a^{w_i}$, so its characteristic polynomial reduces to $\prod_i(X-a^{w_i})$, the reduction of $\prod_i(X-\psi^{\mathrm u}_{v,i}(h_0))$. With (1) and condition (c), Lemma~\ref{sha-rigidity-lagrange}, applied to the local ring $A_n$, $H=G_v$, $r=\rho^{\mathrm u}\circ\mathrm{loc}_v$, the characters $\psi^{\mathrm u}_{v,i}$ and $g_0=h_0$, shows that $\rho^{\mathrm u}\circ\mathrm{loc}_v$ is conjugate over $A_n$ to an upper triangular representation with diagonal entries $\psi^{\mathrm u}_{v,i}$. So the characteristic polynomial of $\rho^{\mathrm u}(\mathrm{loc}_vh)$ is $\prod_i(X-\psi^{\mathrm u}_{v,i}(h))$ in $A_n$, hence in its subring $A_n^\circ$.
\item \emph{Over $C$.} The values $f(\psi^{\mathrm u}_{v,i}(h))$ are units, and by (1) the differences $f(\psi^{\mathrm u}_{v,i}(h_0))-f(\psi^{\mathrm u}_{v,j}(h_0))$, $i\ne j$, are units of $C$. By (2) the characteristic polynomial of $f(\rho^{\mathrm u}(\mathrm{loc}_vh_0))$ is $\prod_i(X-f(\psi^{\mathrm u}_{v,i}(h_0)))$, so its reduction modulo the maximal ideal of $C$ is $\prod_i(X-\bar c_i)$, with $\bar c_i$ the residue of $f(\psi^{\mathrm u}_{v,i}(h_0))$. Condition (c), applied through $f$, gives the vanishing of the ordered products for $f\circ\rho^{\mathrm u}\circ\mathrm{loc}_v$ and the characters $f\circ\psi^{\mathrm u}_{v,i}$. So Lemma~\ref{sha-rigidity-lagrange}, applied to $C$ and $g_0=h_0$, gives $g_v$.
\end{enumerate}
\end{proof}
% Fr.flag (ZpowWNe.lean): hF.f3 v for h_0 (residual condition (7)), zpow_w_ne for (1), co_B for the units of B, o1 (CharpolyFlag.lean: flag_loc over S.R with the diagonal reduction, then charpoly_flag') for (2), and flag_loc (FlagLoc.lean, Lemma sha-rigidity-lagrange) with the map of the charpoly through residue C . f for (3).

\begin{lemma}\label{sha-automorphy-unramified-near-completion-points-subring}
Let $\mathcal O_1$ be a complete noetherian local ring with residue field $k_1$, and let $C$ be a complete local noetherian ring with an $\mathcal O_1$-algebra structure for which $\mathcal O_1\to C$ is local and the residue field $k_C$ of $C$ is finite over $k_1$. Let $C_0$ be the set of $c\in C$ with $c-a\in\mathfrak m_C$ for some $a\in\mathcal O_1$. Then $C_0$ is an $\mathcal O_1$-subalgebra of $C$ that contains $\mathfrak m_C$; it is a complete local noetherian ring with maximal ideal $\mathfrak m_C$, $\mathcal O_1\to C_0$ is local and induces $k_1\cong C_0/\mathfrak m_C$; an element of $C_0$ lies in $\mathfrak m_{C_0}^N$ if and only if it lies in $\mathfrak m_C^N$; and $C$ is a finite $C_0$-module.
\end{lemma}
\leanref{Development/AutomorphicArithmetic/AutomorphicVarieties/CommutativeAlgebraHeight/ClSub.lean}{UpperTheory.V.CA.cl_sub}
\begin{proof}
\begin{enumerate}
\item \emph{A subalgebra.} If $c-a$ and $c'-a'$ lie in $\mathfrak m_C$, so do $(c+c')-(a+a')$ and $cc'-aa'=(c-a)c'+a(c'-a')$. So $C_0$ is a subalgebra; it contains $\mathcal O_1$, and $\mathfrak m_C$ (take $a=0$).
\item \emph{Units.} An element $c$ of $C_0$ is a unit of $C_0$ if and only if $c\notin\mathfrak m_C$. If $c\notin\mathfrak m_C$, write $c=a+m$ with $a\in\mathcal O_1$ and $m\in\mathfrak m_C$; then $a$ is a unit of $\mathcal O_1$, as $\mathcal O_1\to C$ is local, and $c^{-1}-a^{-1}=c^{-1}a^{-1}(a-c)\in\mathfrak m_C$, so $c^{-1}\in C_0$. So the non-units of $C_0$ form the ideal $\mathfrak m_C$: $C_0$ is local with maximal ideal $\mathfrak m_C$, every element is congruent to an element of $\mathcal O_1$ modulo it, and $\mathcal O_1\to C_0$ is local.
\item \emph{Powers.} Clearly $\mathfrak m_{C_0}^N\subseteq\mathfrak m_C^N$. Conversely, by induction on $N$, an element of $\mathfrak m_C^{N+1}$ is a sum of products $ij$ with $i\in\mathfrak m_C$ and $j\in\mathfrak m_C^N$, both in $C_0$, and $j\in\mathfrak m_{C_0}^N$ by induction if $N\ge1$; so it lies in $\mathfrak m_{C_0}^{N+1}$.
\item \emph{Finiteness.} Let $c_1,\dots,c_r\in C$ have residues that span $k_C$ over $k_1$. For $x\in C$ there are $a_t\in\mathcal O_1$ with $y=\sum_ta_tc_t$ of the same residue as $x$; then $x-y\in\mathfrak m_C\subseteq C_0$, and $x=(x-y)\cdot1+y$. So $1,c_1,\dots,c_r$ generate $C$ as a $C_0$-module.
\item \emph{Noetherian.} Let $N$ be a $C_0$-submodule of $C$, and $J=CN$, an ideal of $C$, finitely generated over $C$ and so, by (4), over $C_0$; so is $\mathfrak m_CJ$. As $\mathfrak m_C\subseteq C_0$, $\mathfrak m_CJ\subseteq N\subseteq J$. The module $J/\mathfrak m_CJ$ is finitely generated over $C_0$ and killed by $\mathfrak m_{C_0}=\mathfrak m_C$, so it is a finite-dimensional $k_1$-vector space, and its submodule $N/\mathfrak m_CJ$ is finitely generated. So $N$ is finitely generated: $C$ is a noetherian $C_0$-module, and so is its submodule $C_0$.
\item \emph{Completeness.} By (3), a sequence in $C_0$ is Cauchy for the $\mathfrak m_{C_0}$-adic topology if and only if it is so for the $\mathfrak m_C$-adic topology; its limit $L$ in $C$ is congruent modulo $\mathfrak m_C$ to its terms, so $L\in C_0$, and by (3) again it is the limit in $C_0$. Separatedness follows from that of $C$ in the same way.
\end{enumerate}
\end{proof}
% cl_sub (ClSub.lean): exists_subalgebra_of_residue_constants for (1); residue_constant_subalgebra_isUnit_iff and _isLocalRing for (2); _mem_maximalIdeal_pow_iff for (3); _module_finite for (4); _isNoetherian with isNoetherian_of_finite_torsion_by_maximalIdeal for (5); isAdicComplete_residue_constant_subalgebra for (6). The Lean's hypothesis on C is CA.IsCL O C.

\begin{lemma}\label{sha-automorphy-unramified-near-completion-points-triangular}
In the situation of Lemma~\ref{sha-automorphy-unramified-near-scale}, write $w_i=w_{n,i}$.
\begin{enumerate}
\item Let $K$ be a field, $f:k_1\to K$ a ring homomorphism, and $(\rho,\psi)$, with $\rho:\Gamma_0\to M_n(K)$ a homomorphism of monoids and $\psi_{v,i}:G_v\to K^\times$ homomorphisms, satisfy conditions (a) to (c) of Definition~\ref{sha-descent-link}(2), with $\rho(g)_{ii}=f(\bar\Psi_\Gamma(g)^{w_i})$ for all $g$ and $i$. Then $\rho(g)_{ij}=0$ for all $g$ and all $i>j$, or $\rho(g)_{ij}=0$ for all $g$ and all $i<j$.
\item Let $C$ be a local ring, $f:k_1\to k_C$ a ring homomorphism to its residue field, and $(\rho,\psi)$ with values in $C$ satisfy conditions (a) to (c), with $\rho(g)_{ii}$ of residue $f(\bar\Psi_\Gamma(g)^{w_i})$ for all $g$ and $i$. Then every product $\rho(g_1)_{i_1j_1}\cdots\rho(g_m)_{i_mj_m}$ with $\sum_t(w_{i_t}-w_{j_t})=0$ and $i_t\ne j_t$ for some $t$ lies in $\mathfrak m_C$.
\end{enumerate}
\end{lemma}
\begin{proof}
\begin{enumerate}
\item Write $r_{ij}(g)=\rho(g)_{ij}$, with $1\le i,j\le n$. The diagonal entries are multiplicative, $r_{ii}(gh)=r_{ii}(g)r_{ii}(h)$, as $f\circ\bar\Psi_\Gamma^{w_i}$ is. If $n=2$, $r_{11}(gh)=r_{11}(g)r_{11}(h)+r_{12}(g)r_{21}(h)$ gives $r_{12}(g)r_{21}(h)=0$ for all $g,h$; as $K$ is a field, $r_{12}$ or $r_{21}$ vanishes identically. Let $n=3$.
\begin{enumerate}
\item[(a)] \emph{A diagonal matrix.} Let $X=\rho(g_0)-\rho(g_0^*)$, with $g^*=c_\Gamma g^{-1}c_\Gamma^{-1}$; it is diagonal by condition (b). With $p=\bar\Psi_\Gamma(g_0)$ and $\bar\Psi_\Gamma(g_0^*)=p^{-1}$ (residual condition (4) of Definition~\ref{sha-descent-residual}), $X_{11}=f(p)-f(p^{-1})$, $X_{22}=0$ and $X_{33}=-X_{11}$. As $p^2\ne1$ and $f$ is injective, $X_{11}\ne0$; as $\ell$ is odd, $X_{11}-X_{33}=2X_{11}\ne0$. So $X_{22}\ne X_{33}$ and $X_{11}\ne X_{33}$.
\item[(b)] \emph{Products.} For all $g,h$ and $i$, $(\rho(g)\rho(h))_{ii}=r_{ii}(gh)=r_{ii}(g)r_{ii}(h)$, and $(\rho(g)X\rho(h))_{ii}=r_{ii}(gg_0h)-r_{ii}(gg_0^*h)=r_{ii}(g)X_{ii}r_{ii}(h)$. For $i=1$ these give $r_{12}(g)r_{21}(h)+r_{13}(g)r_{31}(h)=0$ and $X_{22}r_{12}(g)r_{21}(h)+X_{33}r_{13}(g)r_{31}(h)=0$, so both products vanish, by (a). For $i=2$ they give in the same way $r_{21}(g)r_{12}(h)=0$ and $r_{23}(g)r_{32}(h)=0$.
\item[(c)] \emph{Symmetry.} By condition (a), $r_{ij}(g^*)=r_{4-j,4-i}(g)$; so $r_{12}$ vanishes identically if and only if $r_{23}$ does, and $r_{21}$ if and only if $r_{32}$ does.
\item[(d)] \emph{Cases.} By (b), $r_{12}$ or $r_{21}$ vanishes identically. Suppose $r_{12}$ does, hence $r_{23}$. By (b), $r_{13}$ or $r_{31}$ vanishes identically; in the first case all entries above the diagonal vanish. In the second, $r_{12}(gh)=r_{13}(g)r_{32}(h)$, as the other terms contain $r_{12}$ or $r_{31}$; so $r_{13}$ vanishes identically, and all entries above the diagonal vanish, or $r_{32}$ does, hence $r_{21}$, and all entries below the diagonal vanish. If $r_{21}$ vanishes identically, the same argument with the transposed indices, using $r_{32}(gh)=r_{31}(g)r_{12}(h)$, applies.
\end{enumerate}
\item The reduction of $(\rho,\psi)$ modulo $\mathfrak m_C$ satisfies the hypotheses of (1) with $K=k_C$; suppose that its entries below the diagonal vanish, the other case being symmetric. If some $i_t>j_t$, the factor $\rho(g_t)_{i_tj_t}$ lies in $\mathfrak m_C$. Otherwise $i_t\le j_t$ for every $t$, so $w_{i_t}-w_{j_t}\ge0$, as $w$ decreases; since these numbers have sum $0$, all vanish, and $i_t=j_t$ for every $t$, against the hypothesis.
\end{enumerate}
\end{proof}
% (1) is the private tri_of_bar of PointLoc.lean (n = 2: or_zero; n = 3: hX0, hX1, hX2, hq and h2 for (a), hA and hB with two_zero for (b), s01 and s10 from Fr.star for (c), the cases with hC for (d)); (2) is the private cyc_mem, with wsum. The Lean states (1) for (rho.map (residue C)).IsBar f.

\begin{lemma}\label{sha-automorphy-unramified-near-completion-points-laurent}
Let $C'$ be a local ring, $K$ a $C'$-algebra and $(P_j)_{j\in\mathbf Z}$ $C'$-submodules of $K$ with $1\in P_0$ and $P_jP_{j'}\subseteq P_{j+j'}$, such that every element of $P_0$ is the image of an element of $C'$ and, for $j\ne0$, every product $xy$ with $x\in P_j$ and $y\in P_{-j}$ is the image of an element of $\mathfrak m_{C'}$. Let $R$ be the $C'$-subalgebra of $K[T,T^{-1}]$ of the Laurent polynomials $\sum_jx_jT^j$ with $x_j\in P_j$ for every $j$, and $\mathfrak n\subseteq R$ the set of those whose constant coefficient $x_0$ is the image of an element of $\mathfrak m_{C'}$. Then:
\begin{enumerate}
\item $\mathfrak n$ is an ideal of $R$, and every element of $R$ is congruent modulo $\mathfrak n$ to the image of an element of $C'$;
\item if $x\in\mathfrak n^N$, then $x_jT^j\in\mathfrak n^N$ for every $j$;
\item if $C'\to K$ is injective and $R$ is noetherian, then $C'\to R$ is injective, the inverse image of $\mathfrak n$ in $C'$ is $\mathfrak m_{C'}$, and $\bigcap_N\mathfrak n^N=0$.
\end{enumerate}
\end{lemma}
\begin{proof}
\begin{enumerate}
\item The set $R$ is a subalgebra, as $P_jP_{j'}\subseteq P_{j+j'}$ and $1\in P_0$. Sums of elements of $\mathfrak n$ lie in $\mathfrak n$. For $q\in R$ and $x\in\mathfrak n$, the constant coefficient of $qx$ is $\sum_aq_ax_{-a}$; the term $a=0$ is $q_0x_0$, the image of an element of $\mathfrak m_{C'}$ as $q_0\in P_0$ is the image of an element of $C'$, and the other terms are images of elements of $\mathfrak m_{C'}$ by hypothesis. For $x\in R$, $x_0$ is the image of some $c\in C'$, and $x-c\in\mathfrak n$.
\item If $x\in\mathfrak n$, then $x_jT^j\in\mathfrak n$, as its constant coefficient is $x_0$ or $0$. Suppose (2) holds for $N$. An element of $\mathfrak n^{N+1}$ is a sum of products $xy$ with $x\in\mathfrak n^N$ and $y\in\mathfrak n$, and $(xy)_jT^j=\sum_a(x_aT^a)(y_{j-a}T^{j-a})$, a sum of products of an element of $\mathfrak n^N$ and one of $\mathfrak n$.
\item The constant coefficient of the image of $c\in C'$ is the image of $c$ in $K$; so $C'\to R$ is injective, and $c$ maps into $\mathfrak n$ if and only if $c\in\mathfrak m_{C'}$. Let $x\in\bigcap_N\mathfrak n^N$ and $j\in\mathbf Z$. By (2), $q=x_jT^j\in\bigcap_N\mathfrak n^N$, so by Krull's intersection theorem \cite[Theorem~8.9]{Mat86} there is $r\in\mathfrak n$ with $rq=q$. Comparing the coefficients of $T^j$, $r_0x_j=x_j$, where $r_0$ is the image of some $c\in\mathfrak m_{C'}$. As $1-c$ is a unit, $x_j=0$. So $x=0$.
\end{enumerate}
\end{proof}
% vertex and mem_vertex (AlgDefs.lean, with HasVertex: the fields zero and cross) for (1); alg_homog (AlgHomog.lean, single_coeff_mul) for (2); the private alg_vring of AlgVring.lean (inj, comap, surj) and alg_sep (AlgSep.lean, Ideal.mem_iInf_smul_pow_eq_bot_iff) for (3). The ring R is lauSub P, an IsVertexRing C (vertex h) in the Lean's terms.

\begin{lemma}\label{sha-automorphy-unramified-near-completion-points-cone}
In the situation of Lemma~\ref{sha-automorphy-unramified-near-scale}, write $w_i=w_{n,i}$, let $C'$ be a complete local noetherian $\mathcal O_1$-algebra with residue field $k_1$, $K$ a $C'$-algebra with $C'\to K$ injective, and $(\rho,\psi)$, with $\rho:\Gamma_0\to M_n(K)$ a homomorphism of monoids and $\psi_{v,i}:G_v\to K^\times$ homomorphisms, satisfy conditions (a) to (c) of Definition~\ref{sha-descent-link}(2). Suppose:
\begin{enumerate}
\item[(i)] each $\rho(g)_{ii}$ is the image of an element of $C'$ of residue $\bar\Psi_\Gamma(g)^{w_i}$;
\item[(ii)] each product $\rho(g_1)_{i_1j_1}\cdots\rho(g_m)_{i_mj_m}$ with $\sum_t(w_{i_t}-w_{j_t})=0$ and $i_t\ne j_t$ for some $t$ is the image of an element of $\mathfrak m_{C'}$;
\item[(iii)] for all $i,j$, the $C'$-submodule $M_{ij}$ of $K$ spanned by the $\rho(g)_{ij}$, $g\in\Gamma_0$, is finitely generated, and for every $N$ the map $g\mapsto\rho(g)_{ij}\bmod\mathfrak m_{C'}^NM_{ij}$ is locally constant;
\item[(iv)] there are maps $\chi_{v,i}:G_v\to C'$, locally constant modulo every power of $\mathfrak m_{C'}$, of residues $\bar\Psi_\Gamma(\mathrm{loc}_vh)^{w_i}$, whose images in $K$ are the $\psi_{v,i}$.
\end{enumerate}
Then there are $\mathcal O_1$-algebra maps $\xi:A_n^\circ\to K$ and $\beta:B\to C'$ such that $\xi$ carries $(\rho^{\mathrm u},\psi^{\mathrm u})$ to $(\rho,\psi)$ and $\xi(b)$ is the image of $\beta(b)$ for every $b\in B$.
\end{lemma}
\leanref{Development/AutomorphicArithmetic/AutomorphicVarieties/FrobeniusSmoothVectors/AlgVring.lean}{UpperTheory.V.Fr.alg}
\leanref{Development/AutomorphicArithmetic/AutomorphicVarieties/FrobeniusSmoothVectors/AlgLimit.lean}{UpperTheory.V.Fr.alg_limit}
\leanref{Development/AutomorphicArithmetic/AutomorphicVarieties/FrobeniusSmoothVectors/AlgFamily.lean}{UpperTheory.V.Fr.alg_family}
\begin{proof}
\begin{enumerate}
\item \emph{A graded model.} For $j\in\mathbf Z$ let $P_j$ be the $C'$-span in $K$ of the products of entries of $\rho$ of total weight $j$, the entry $(i,i')$ having weight $w_i-w_{i'}$. Then $1\in P_0$ and $P_jP_{j'}\subseteq P_{j+j'}$. A product of total weight $0$ is the image of an element of $C'$, by (ii) if it has a factor off the diagonal and by (i) otherwise; and for $j\ne0$ a product of weight $j$ has a factor off the diagonal, so the product of one of weight $j$ and one of weight $-j$ is the image of an element of $\mathfrak m_{C'}$, by (ii). So Lemma~\ref{sha-automorphy-unramified-near-completion-points-laurent} applies to the $P_j$; let $R$ and $\mathfrak n$ be as there.
\item \emph{$R$ is noetherian.} By (iii), each $M_{ij}$ is spanned by finitely many entries; the $C'$-subalgebra of $K[T,T^{-1}]$ generated by the $xT^{w_i-w_j}$, for these finitely many $x$ and all $i,j$, contains every $\rho(g)_{ij}T^{w_i-w_j}$, hence every $yT^j$ with $y$ a product of total weight $j$, hence every $P_jT^j$, and so it is $R$. As $C'$ is noetherian, so is $R$, by Hilbert's basis theorem. By Lemma~\ref{sha-automorphy-unramified-near-completion-points-laurent}(3), $C'\subseteq R$, $\mathfrak n\cap C'=\mathfrak m_{C'}$ and $\bigcap_N\mathfrak n^N=0$; by its part (1), $\mathfrak n$ is a maximal ideal with $R/\mathfrak n=k_1$.
\item \emph{A point over $R$.} Let $\rho_D(g)$ be the matrix with entries $\rho(g)_{ij}T^{w_i-w_j}$, that is, $d\rho(g)d^{-1}$ with $d=\operatorname{diag}(T^{w_1},\dots,T^{w_n})$, and $\psi_D=\psi$. Then $(\rho_D,\psi_D)$ takes values in $R$ and satisfies conditions (a) to (c): $J_ndJ_n=d^{-1}$ as $w_{n+1-i}=-w_i$, which gives (a), and conjugation by a diagonal matrix keeps the vanishing in (b) and (c). Its entries off the diagonal lie in $\mathfrak n$, their constant coefficients being $0$; its diagonal entries and the values of $\psi_D$ lie in $C'$, with residues $\bar\Psi_\Gamma(g)^{w_i}$ and $\bar\Psi_\Gamma(\mathrm{loc}_vh)^{w_i}$, by (i) and (iv); and products of its entries of total weight $0$ lie in $C'$, by (1).
\item \emph{Continuity.} For $N\ge0$, the maps $g\mapsto\rho_D(g)_{ij}$ and $h\mapsto\psi_{D,v,i}(h)$ are locally constant modulo $\mathfrak n^N$. Indeed, by (iii), near any $g$ the difference $\rho(g')_{ij}-\rho(g)_{ij}$ lies in $\mathfrak m_{C'}^NM_{ij}$, and $\mathfrak m_{C'}^NM_{ij}T^{w_i-w_j}\subseteq\mathfrak n^N$, as $\mathfrak m_{C'}\subseteq\mathfrak n$ and $M_{ij}T^{w_i-w_j}\subseteq R$; and by (iv), near any $h$ the difference $\chi_{v,i}(h')-\chi_{v,i}(h)$ lies in $\mathfrak m_{C'}^N\subseteq\mathfrak n^N$.
\item \emph{Finite levels.} Let $N\ge0$. The ideal $\mathfrak n$ is maximal, so $R/\mathfrak n^{N+1}$ is a local ring whose maximal ideal $\mathfrak n/\mathfrak n^{N+1}$ is nilpotent; it is noetherian, complete, and an $\mathcal O_1$-algebra with residue field $k_1$, with $\mathcal O_1\to R/\mathfrak n^{N+1}$ local, by (2). By (3) and (4), the image of $(\rho_D,\psi_D)$ in it satisfies conditions (a) to (d), and its reduction is the residual linked point. By the universal property of $A_n$ (Lemma~\ref{sha-automorphy-realization-exists-linked}) there is a unique $\mathcal O_1$-algebra map $f_N:A_n\to R/\mathfrak n^{N+1}$ that carries $(\rho^{\mathrm u},\psi^{\mathrm u})$ to it. By uniqueness, $f_{N+1}$ followed by the projection is $f_N$. And $f_N$ kills $\mathfrak m_{A_n}^{N+1}$: it maps $\mathfrak m_{A_n}$ into the maximal ideal $\mathfrak n/\mathfrak n^{N+1}$, whose $(N+1)$-st power is $0$: for $y\in\mathfrak m_{A_n}$ and a unit $a$ of $\mathcal O_1$, $y-a$ is a unit, so $f_N(y)-a$ is one, and the residue of $f_N(y)$, which lies in $k_1$, is not that of $a$; so it is $0$.
\item \emph{Limits in $C'$.} Let $T'$ be the set of $a\in A_n$ for which some $c\in C'$ has $c\equiv f_N(a)$ modulo $\mathfrak n^{N+1}$ for every $N$; it is an $\mathcal O_1$-subalgebra. It contains the values of the $\psi^{\mathrm u}_{v,i}$ and the products of entries of $\rho^{\mathrm u}$ of total weight $0$, whose images under every $f_N$ are those of the corresponding elements of $C'$ given by (3). So it contains the $\mathcal O_1$-subalgebra $B_0$ that they generate. Let $b\in B$. As $B$ is the closure of $B_0$, there are $y_M\in B_0$ with $b-y_M\in\mathfrak m_{A_n}^M$, and $c_M\in C'$ with $c_M\equiv f_{N}(y_{M+1})$ modulo $\mathfrak n^{N+1}$ for all $N$. By (5), $f_N(b)=f_N(y_M)$ for $M\ge N+1$; so $c_{M+1}-c_M\in J_{M+1}$, where $J_N=C'\cap\mathfrak n^N$. The ideals $J_N$ decrease and have intersection $0$ by (2), so Lemma~\ref{sha-automorphy-unramified-near-determined-cofinal} gives $s_k$ with $J_{s_k}\subseteq\mathfrak m_{C'}^k$; a subsequence of the $c_M$ is then Cauchy for the $\mathfrak m_{C'}$-adic topology, and its limit $z$ in the complete ring $C'$ has $z-c_M\in J_{M+1}$ for all $M$, as $\mathfrak m_{C'}^{M+1}\subseteq J_{M+1}$. So $z\equiv f_N(b)$ modulo $\mathfrak n^{N+1}$ for every $N$, and $B\subseteq T'$.
\item \emph{The map.} The set of $a\in A_n$ for which some $d\in R$ has $d\equiv f_N(a)$ modulo $\mathfrak n^{N+1}$ for every $N$ is an $\mathcal O_1$-subalgebra; it contains $B$ by (6) and the entries of $\rho^{\mathrm u}$, with $d$ the entries of $\rho_D$, so it contains $A_n^\circ$. As $\bigcap_N\mathfrak n^N=0$, such a $d$ is unique, and $a\mapsto d$ is an $\mathcal O_1$-algebra map $\xi_R:A_n^\circ\to R$. By uniqueness, $\xi_R$ carries $(\rho^{\mathrm u},\psi^{\mathrm u})$ to $(\rho_D,\psi_D)$, and by (6) it maps $B$ into $C'$; let $\beta:B\to C'$ be the resulting map.
\item \emph{Evaluation.} Let $e:K[T,T^{-1}]\to K$ be the $C'$-algebra map with $e(T)=1$, and $\xi=e\circ\xi_R$. Then $\xi(\rho^{\mathrm u}(g)_{ij})=e(\rho(g)_{ij}T^{w_i-w_j})=\rho(g)_{ij}$, $\xi\circ\psi^{\mathrm u}_{v,i}=\psi_{v,i}$, and $\xi(b)$ is the image of $\beta(b)$.
\end{enumerate}
\end{proof}
% alg (AlgVring.lean): alg_vertex (AlgVertex.lean) for (1); alg_ft (SingleMemOfSpan.lean, single_mem_of_span) and the private alg_vring for (2); alg_lau (AlgLau.lean, with alg_twist of AlgTwist.lean) and alg_model (AlgModel.lean) for (3), alg_cont (AlgCont.lean) for (4); alg_limit (AlgLimit.lean) for (5) to (7): alg_family (AlgFamily.lean) with alg_quot (AlgQuot.lean), alg_lift (AlgLift.lean) and S.univ for (5), the subalgebra T' with alg_closed (AlgClosed.lean, through GlobalTheory.exists_le_maximalIdeal_pow_of_iInf_eq_bot, the statement of Lemma sha-automorphy-unramified-near-determined-cofinal) for (6), and alg_glue (AlgGlue.lean) for (7); the private alg_ev for (8). The Lean's hypotheses (i) to (iv) are the fields dg, cyc, fg with cont, and psi of Fr.Cone O C rho.

\begin{lemma}\label{sha-automorphy-unramified-near-completion-points}
In the situation of Lemma~\ref{sha-automorphy-unramified-near-scale}, let $C$ be a complete local noetherian ring with an $\mathcal O_1$-algebra structure for which $\mathcal O_1\to C$ is local and the residue field $k_C$ of $C$ is finite over $k_1$. Let $\rho:\Gamma_0\to M_n(C)$ and $\psi_{v,i}:G_v\to C^\times$ satisfy conditions (a) to (d) of Definition~\ref{sha-descent-link}(2), and suppose that the diagonal entries of each $\rho(g)$ and the values $\psi_{v,i}(h)$ reduce to $\bar\Psi_\Gamma(g)^{w_{n,i}}$ and $\bar\Psi_\Gamma(\mathrm{loc}_vh)^{w_{n,i}}$, viewed in $k_C$. Then there is an $\mathcal O_1$-algebra map $\xi:A_n^\circ\to C$ that carries $(\rho^{\mathrm u},\psi^{\mathrm u})$ to $(\rho,\psi)$ and the maximal ideal of $B$ into $\mathfrak m_C$.
\end{lemma}
\leanref{Development/AutomorphicArithmetic/AutomorphicVarieties/FrobeniusSmoothVectors/PointLoc.lean}{UpperTheory.V.Fr.point_loc}
\leanref{Development/AutomorphicArithmetic/AutomorphicVarieties/FrobeniusSmoothVectors/AlgVring.lean}{UpperTheory.V.Fr.alg}
\begin{proof}
Let $f:k_1\to k_C$ be the map of residue fields, and write $w_i=w_{n,i}$.
\begin{enumerate}
\item \emph{A subring.} Let $C_0\subseteq C$ be the subring of Lemma~\ref{sha-automorphy-unramified-near-completion-points-subring}: a complete local noetherian $\mathcal O_1$-algebra with residue field $k_1$ and maximal ideal $\mathfrak m_C$, over which $C$ is finite, with $\mathfrak m_{C_0}^N=C_0\cap\mathfrak m_C^N$.
\item \emph{Values in $C_0$.} An element $c\in C$ whose residue is $f(t)$ for some $t\in k_1$ lies in $C_0$, with residue $t$ there: $t$ is the residue of some $a\in\mathcal O_1$, and $c-a\in\mathfrak m_C$. By hypothesis this applies to the diagonal entries $\rho(g)_{ii}$, with $t=\bar\Psi_\Gamma(g)^{w_i}$, and to the values $\psi_{v,i}(h)$, with $t=\bar\Psi_\Gamma(\mathrm{loc}_vh)^{w_i}$.
\item \emph{The hypotheses (i) to (iv) of Lemma~\ref{sha-automorphy-unramified-near-completion-points-cone}, for $C'=C_0$ and $K=C$.} (i) is (2). (ii) is Lemma~\ref{sha-automorphy-unramified-near-completion-points-triangular}(2), as $\mathfrak m_C=\mathfrak m_{C_0}$. For (iii): $C$ is a finite module over the noetherian ring $C_0$, so the submodule $M_{ij}$ spanned by the $\rho(g)_{ij}$ is finitely generated. Let $N\ge0$. By the Artin-Rees lemma \cite[Theorem~8.5]{Mat86} there is $k$ with $\mathfrak m_{C_0}^{N+k}C\cap M_{ij}\subseteq\mathfrak m_{C_0}^NM_{ij}$; and $\mathfrak m_C^{N+k}\subseteq\mathfrak m_{C_0}^{N+k}C$, as $\mathfrak m_C=\mathfrak m_{C_0}$. By condition (d), near any $g$ the difference $\rho(g')_{ij}-\rho(g)_{ij}$ lies in $\mathfrak m_C^{N+k}\cap M_{ij}\subseteq\mathfrak m_{C_0}^NM_{ij}$. For (iv), take $\chi_{v,i}=\psi_{v,i}$, with values in $C_0$ and the residues of (2); it is locally constant modulo $\mathfrak m_{C_0}^N=C_0\cap\mathfrak m_C^N$ by condition (d).
\item \emph{The map.} Lemma~\ref{sha-automorphy-unramified-near-completion-points-cone} gives $\mathcal O_1$-algebra maps $\xi:A_n^\circ\to C$, carrying $(\rho^{\mathrm u},\psi^{\mathrm u})$ to $(\rho,\psi)$, and $\beta:B\to C_0$ with $\xi(b)=\beta(b)$ for $b\in B$.
\item \emph{The maximal ideal of $B$.} Let $b\in\mathfrak m_B$. For a unit $a$ of $\mathcal O_1$, $b-a$ is a unit of $B$, so $\beta(b)-a$ is a unit of $C_0$; so the residue of $\beta(b)$ in $k_1$ is $0$, and $\xi(b)=\beta(b)\in\mathfrak m_{C_0}=\mathfrak m_C$.
\end{enumerate}
\end{proof}
% point_loc (PointLoc.lean): CA.cl_sub for (1); hin and hres for (2); the Cone over S' for (3): its field dg from (2), cyc from cyc_mem, fg from IsNoetherian S' C, cont from the private cont_sub (Ideal.exists_pow_inf_eq_pow_smul, the Artin-Rees lemma, with hJN), psi with lc_lemma and hpow; alg (Lemma sha-automorphy-unramified-near-completion-points-cone) for (4); res_map and co_B (hmB) for (5). The Lean's C is a CA.IsCL O C with a residue map f satisfying IsResMap, and the conclusion S.x0 <= comap xi (maximalIdeal C) is (5).

\begin{lemma}\label{sha-automorphy-unramified-near-completion-unique}
In the situation of Lemma~\ref{sha-automorphy-unramified-near-scale}, let $C$ be a noetherian local ring that is an $\mathcal O_1$-algebra, and let $\xi,\xi':A_n^\circ\to C$ be $\mathcal O_1$-algebra maps that carry $(\rho^{\mathrm u},\psi^{\mathrm u})$ to the same pair and the maximal ideal $\mathfrak m_B$ of $B$ into $\mathfrak m_C$. Then $\xi=\xi'$.
\end{lemma}
\leanref{Development/AutomorphicArithmetic/AutomorphicVarieties/FrobeniusSmoothVectors/ExtA.lean}{UpperTheory.V.Fr.ext_A}
\begin{proof}
Let $E\subseteq A_n^\circ$ be the $\mathcal O_1$-subalgebra on which $\xi$ and $\xi'$ agree. It contains the entries of $\rho^{\mathrm u}$ and the values of the $\psi^{\mathrm u}_{v,i}$, so it contains the $\mathcal O_1$-subalgebra $B_0$ generated by these values and by the products of entries of total weight $0$, whose closure for the $\mathfrak m_{A_n}$-adic topology is $B$ (Definition~\ref{sha-automorphy-realization}(2)). Let $b\in B$ and $N\ge0$. By Lemma~\ref{sha-automorphy-unramified-near-tangent-adic} there is $M$ with $B\cap\mathfrak m_{A_n}^M\subseteq\mathfrak m_B^N$, and there is $b_0\in B_0$ with $b-b_0\in\mathfrak m_{A_n}^M$, hence $b-b_0\in\mathfrak m_B^N$. So $\xi(b-b_0)$ and $\xi'(b-b_0)$ lie in $\mathfrak m_C^N$, and $\xi(b)-\xi'(b)=\xi(b-b_0)-\xi'(b-b_0)\in\mathfrak m_C^N$. By Krull's intersection theorem \cite[Theorem~8.10]{Mat86}, $\xi(b)=\xi'(b)$. So $B\subseteq E$, and as $A_n^\circ$ is generated by $B$ and the entries of $\rho^{\mathrm u}$, $E=A_n^\circ$.
\end{proof}
% ext_A (ExtA.lean): the equalizer E, hr and hpsi for the generators, h0 for B_0 (the adjoin of S.rho.mon 0 and the characters), hpow and the field hc 0 N of cont_u (Lemma sha-automorphy-unramified-near-tangent-adic) for B, Ideal.iInf_pow_eq_bot_of_isLocalRing, and Algebra.adjoin_le for A_n^circ.

\begin{definition}\label{sha-automorphy-unramified-near-completion-tangent-triangulation}
In the situation of Lemma~\ref{sha-automorphy-unramified-near-scale}, let $C$ be a commutative ring. Let $G_n(C)=GL_n(C)\rtimes\{\pm1\}$, where $-1$ acts by $\theta(g)=J_n{}^tg^{-1}J_n$; for $x=(g,\epsilon)\in G_n(C)$ write $\mathrm{mat}(x)=g$, and let $x$ act on $M_n(C)$ by $x\cdot X=gXg^{-1}$ if $\epsilon=1$ and $x\cdot X=-gJ_n{}^tXJ_ng^{-1}$ if $\epsilon=-1$.
\begin{enumerate}
\item A \emph{triangulation} over $C$ is a pair $(r,(g_v))$ of a homomorphism $r:\Gamma\to G_n(C)$ whose sign is $\delta_\Gamma$ and matrices $g_v\in GL_n(C)$, for the places $v$ of $L^+$ above $\ell$, such that $g_v\,\mathrm{mat}(r(\mathrm{loc}_vh))\,g_v^{-1}$ is upper triangular for every $h\in G_v$. If $C$ is local, the triangulation is \emph{continuous} if the entries of $\mathrm{mat}\circ r$ are locally constant modulo every power of $\mathfrak m_C$.
\item It is a triangulation \emph{of} a pair $(\rho,\psi)$ with values in $C$ that satisfies conditions (a) to (c) of Definition~\ref{sha-descent-link}(2) if $\mathrm{mat}(r(c_\Gamma))=1$, $\mathrm{mat}(r(g))=\rho(g)$ for $g\in\Gamma_0$, and the diagonal entries of $g_v\rho(\mathrm{loc}_vh)g_v^{-1}$ are $\psi_{v,1}(h),\dots,\psi_{v,n}(h)$ for every $v$ and $h\in G_v$.
\item If $C$ is local with residue field $k'$, then $\Gamma$ acts on $W=M_n(k')$ through the reduction $\bar r$ of $r$, by $s\cdot X=\bar r(s)\cdot X$; let $b_v\subseteq W$ be the space of the $X$ with $\bar g_vX\bar g_v^{-1}$ upper triangular.
\end{enumerate}
\end{definition}
\leanref{Development/AutomorphicArithmetic/AutomorphicVarieties/FrobeniusSmoothVectors/PresDefs.lean}{UpperTheory.V.Fr.Box}
\leanref{Development/AutomorphicArithmetic/AutomorphicVarieties/FrobeniusSmoothVectors/PresDefs.lean}{UpperTheory.V.Fr.Box.Of}
\leanref{Development/AutomorphicArithmetic/AutomorphicVarieties/FrobeniusSmoothVectors/PresDefs.lean}{UpperTheory.V.Fr.Box.ad}
\leanref{Development/AutomorphicArithmetic/AutomorphicVarieties/UnitaryThreeGroups/UnitaryThreeReflectionRight.lean}{UpperTheory.V.U3.GG.ad}

\begin{lemma}\label{sha-automorphy-unramified-near-completion-tangent-correct}
In the situation of Lemma~\ref{sha-automorphy-unramified-near-scale}, let $\mathcal O'$ be a local ring with residue field $k'$, let $p:C'\to C$ be a map of complete local noetherian $\mathcal O'$-algebras with residue field $k'$, and let $\mu:k'\to C'$ be an additive map with $\mu(k')\subseteq\mathfrak m_{C'}$, $p\circ\mu=0$ and $\mu(\bar ac)=a\mu(c)$ for $a\in C'$ and $c\in k'$, where $\bar a$ is the residue of $a$; for a matrix $X$ over $k'$ write $\mu(X)$ for the matrix of the $\mu(X_{ij})$. Let $(r,(g_v))$ be a triangulation over $C$, with $W$ and the $b_v$ of Definition~\ref{sha-automorphy-unramified-near-completion-tangent-triangulation}(3). Let $\tilde r:\Gamma\to GL_n(C')$, $\tilde g_v\in GL_n(C')$ and $\tilde t_v:G_v\to GL_n(C')$ lift $\mathrm{mat}\circ r$, the $g_v$ and $h\mapsto g_v\,\mathrm{mat}(r(\mathrm{loc}_vh))\,g_v^{-1}$, with $\tilde r$ and the $\tilde t_v$ locally constant modulo every power of $\mathfrak m_{C'}$ and every $\tilde t_v(h)$ upper triangular. Let $\beta:\Gamma\to W$ be locally constant, $Y_v\in W$ and $a_v:G_v\to W$, and suppose, with $d\beta(s,t)=\beta(s)+s\cdot\beta(t)-\beta(st)$ and $\theta^1=\mathrm{id}$, $\theta^{-1}=\theta$:
\begin{enumerate}
\item[(i)] $\tilde r(s)\,\theta^{\delta_\Gamma(s)}(\tilde r(t))\,\tilde r(st)^{-1}=1+\mu(d\beta(s,t))$ for all $s,t\in\Gamma$;
\item[(ii)] $\tilde r(\mathrm{loc}_vh)\bigl(\tilde g_v^{-1}\tilde t_v(h)\tilde g_v\bigr)^{-1}=1+\mu(a_v(h))$ for all $v$ and $h\in G_v$;
\item[(iii)] $a_v(h)-\bigl(\beta(\mathrm{loc}_vh)-(\mathrm{loc}_v(h)\cdot Y_v-Y_v)\bigr)\in b_v$ for all $v$ and $h\in G_v$.
\end{enumerate}
Then there is a continuous triangulation $(r',(g'_v))$ over $C'$ whose image under $p$ is $(r,(g_v))$, with $\mathrm{mat}(r'(s))=(1-\mu(\beta(s)))\tilde r(s)$ for every $s\in\Gamma$.
\end{lemma}
\leanref{Development/AutomorphicArithmetic/AutomorphicVarieties/FrobeniusSmoothVectors/PrCorrect.lean}{UpperTheory.V.Fr.pr_correct}
\leanref{Development/AutomorphicArithmetic/AutomorphicVarieties/FrobeniusSmoothVectors/PrSqz.lean}{UpperTheory.V.Fr.pr_sqz}
\begin{proof}
\begin{enumerate}
\item \emph{Square zero.} For $c,d\in k'$, $\mu(c)\mu(d)=\mu(\overline{\mu(c)}d)=0$, as $\mu(c)\in\mathfrak m_{C'}$; and for a matrix $A$ over $C'$, $A\mu(X)=\mu(\bar AX)$ and $\mu(X)A=\mu(X\bar A)$. So $e(X)=1+\mu(X)$ satisfies $e(X)e(Y)=e(X+Y)$, it is invertible with $e(X)^{-1}=e(-X)$, $p(e(X))=1$, and $\theta(e(Y))=e(-J_n{}^tYJ_n)$. Hence, for $x\in G_n(C')$ with reduction $\bar x$, $x\,(e(X),1)\,x^{-1}=(e(\bar x\cdot X),1)$ in $G_n(C')$.
\item \emph{The homomorphism.} Let $x(s)=(\tilde r(s),\delta_\Gamma(s))\in G_n(C')$, whose reduction is $\bar r(s)$. By (i), $x(s)x(t)=(e(d\beta(s,t)),1)\,x(st)$. Let $r'(s)=(e(-\beta(s)),1)\,x(s)$. By (1), $r'(s)r'(t)=(e(-\beta(s)-s\cdot\beta(t)),1)\,x(s)x(t)=(e(-\beta(st)),1)\,x(st)=r'(st)$. So $r'$ is a homomorphism with sign $\delta_\Gamma$ and $\mathrm{mat}(r'(s))=(1-\mu(\beta(s)))\tilde r(s)$; it lifts $r$, as $p(e(X))=1$; and it is continuous, as $\beta$ is locally constant.
\item \emph{Triangular forms.} Let $g'_v=\tilde g_ve(-Y_v)$, a lift of $g_v$. Fix $v$ and $h\in G_v$, and write $\tilde r=\tilde r(\mathrm{loc}_vh)$, $\beta=\beta(\mathrm{loc}_vh)$, $a=a_v(h)$, $\tilde t=\tilde t_v(h)$. By (ii), $\tilde r\tilde g_v^{-1}=e(a)\tilde g_v^{-1}\tilde t$, and by (1), $\tilde re(Y_v)\tilde r^{-1}=e(\mathrm{loc}_v(h)\cdot Y_v)$. So
\[g'_v\,\mathrm{mat}(r'(\mathrm{loc}_vh))\,g_v'^{-1}=\tilde g_ve(-Y_v)e(-\beta)e(\mathrm{loc}_v(h)\cdot Y_v)e(a)\tilde g_v^{-1}\tilde t=e(\bar g_vZ\bar g_v^{-1})\tilde t\]
with $Z=a-(\beta-(\mathrm{loc}_v(h)\cdot Y_v-Y_v))$, by (1). As $Z\in b_v$ by (iii), $e(\bar g_vZ\bar g_v^{-1})$ is upper triangular, and so is the product.
\end{enumerate}
\end{proof}
% pr_correct (PrCorrect.lean): pr_sqz (PrSqz.lean: hsq, hL, hR, hu, e with hemul and heinv, hconj, htheta) for (1); x, hxmul, hcomm and hmul for (2), with the field r_cont of SetLift for the continuity; hr, h1, hX, key, hW and upT_mul for (3). The Lean packages the lifts as a SetLift p B (fields r, g, t with r_map, g_map, t_map, r_cont, t_cont, t_tri) and (i), (ii) as D.om x = (dOne ad beta x).map mu and D.a v s = (a v s).map mu.

\begin{lemma}\label{sha-automorphy-unramified-near-completion-tangent-normal}
In the situation of Lemma~\ref{sha-automorphy-unramified-near-scale}, write $w_i=w_{n,i}$ and $g^*=c_\Gamma g^{-1}c_\Gamma^{-1}$, and let $\mathcal O'$ be a local ring with residue field $k'$.
\begin{enumerate}
\item Let $K$ be a field, $f:k_1\to K$ a ring homomorphism, and $(\rho,\psi)$ with values in $K$ satisfy conditions (a) to (c) of Definition~\ref{sha-descent-link}(2), with $\rho(g)_{ii}=f(\bar\Psi_\Gamma(g)^{w_i})$ and $\psi_{v,i}(h)=f(\bar\Psi_\Gamma(\mathrm{loc}_vh)^{w_i})$. Then $2\ne0$ in $K$, the diagonal entries of $\rho(g_0)-\rho(g_0^*)$ are pairwise distinct, for every $v$ some $h\in G_v$ has the $\psi_{v,i}(h)$ pairwise distinct, and $T=\operatorname{diag}(w_1,\dots,w_n)$ is nonzero with $-J_n{}^tTJ_n=T$.
\item Let $C$ be a local ring and $(r,(g_v))$ a continuous triangulation over $C$ with $\mathrm{mat}(r(c_\Gamma))=1$ and with the entries off the diagonal of $\mathrm{mat}(r(g_0))-\mathrm{mat}(r(g_0^*))$ equal to $0$. Let $\rho(g)=\mathrm{mat}(r(g))$ for $g\in\Gamma_0$, and $\psi_{v,i}(h)$ the diagonal entries of $g_v\rho(\mathrm{loc}_vh)g_v^{-1}$. Then $(\rho,\psi)$ satisfies conditions (a) to (d), and $(r,(g_v))$ is a triangulation of it.
\item Let $p:C'\to C$ be a surjective map of complete local noetherian $\mathcal O'$-algebras with residue field $k'$, with $\mathfrak m_{C'}\ker p=0$. Let $(\rho,\psi)$ with values in $C$ satisfy conditions (a) to (c), with $2\ne0$ in $k'$ and the diagonal entries of $\bar\rho(g_0)-\bar\rho(g_0^*)$ pairwise distinct, where $\bar\rho$ is the reduction of $\rho$. Let $(r',(g'_v))$ be a continuous triangulation over $C'$ whose image under $p$ is a triangulation of $(\rho,\psi)$. Then there are $h\in GL_n(C')$ with $p(h)=1$ and $(\rho',\psi')$ with values in $C'$ that satisfies conditions (a) to (d), with image $(\rho,\psi)$ under $p$, such that $x=(h,1)$ has $\mathrm{mat}(x\,r'(c_\Gamma)\,x^{-1})=1$ and $\mathrm{mat}(x\,r'(g)\,x^{-1})=\rho'(g)$ for every $g\in\Gamma_0$.
\end{enumerate}
\end{lemma}
\leanref{Development/AutomorphicArithmetic/AutomorphicVarieties/FrobeniusSmoothVectors/PrGood.lean}{UpperTheory.V.Fr.pr_good}
\leanref{Development/AutomorphicArithmetic/AutomorphicVarieties/FrobeniusSmoothVectors/PrOfbox.lean}{UpperTheory.V.Fr.pr_ofbox}
\leanref{Development/AutomorphicArithmetic/AutomorphicVarieties/FrobeniusSmoothVectors/PrNorm.lean}{UpperTheory.V.Fr.pr_norm}
\leanref{Development/AutomorphicArithmetic/AutomorphicVarieties/FrobeniusSmoothVectors/PrNorm1.lean}{UpperTheory.V.Fr.pr_norm1}
\leanref{Development/AutomorphicArithmetic/AutomorphicVarieties/FrobeniusSmoothVectors/JconjApply.lean}{UpperTheory.V.Fr.pr_norm2}
\begin{proof}
\begin{enumerate}
\item As $\ell$ is odd, $2\ne0$ in $k_1$, hence in $K$. Let $a=\bar\Psi_\Gamma(g_0)$; then $\bar\Psi_\Gamma(g_0^*)=a^{-1}$, as $\bar\Psi_\Gamma(c_\Gamma)=1$ (residual condition (4) of Definition~\ref{sha-descent-residual}), so the diagonal entries of $\rho(g_0)-\rho(g_0^*)$ are $f(a^{w_i}-a^{-w_i})$: for $n=3$ these are $f(a-a^{-1})$, $0$ and $-f(a-a^{-1})$, and for $n=2$ the first and last; they are pairwise distinct as $a^2\ne1$, $2\ne0$ and $f$ is injective. By residual condition (7) some $h\in G_v$ has $\bar\Psi_\Gamma(\mathrm{loc}_vh)^2\ne1$, and then the $\psi_{v,i}(h)$ are pairwise distinct, as $w_i-w_j$ divides $2$ for $i\ne j$. Finally $w_1=1$, and $w_{n+1-i}=-w_i$ gives $-J_n{}^tTJ_n=T$.
\item As $\delta_\Gamma$ is trivial on $\Gamma_0$, $\rho$ is multiplicative. For $g\in\Gamma_0$, $r(g^*)=r(c_\Gamma)r(g)^{-1}r(c_\Gamma)^{-1}=(1,-1)(\rho(g),1)^{-1}(1,-1)=(\theta(\rho(g)^{-1}),1)$, so $\rho(g^*)=J_n{}^t\rho(g)J_n$, which is condition (a); condition (b) is a hypothesis. The matrices $g_v\rho(\mathrm{loc}_vh)g_v^{-1}$ are upper triangular, so their diagonal entries $\psi_{v,i}$ are multiplicative, and the $i$-th factor of the product in (c), conjugated by $g_v$, maps the span of the first $i$ basis vectors into that of the first $i-1$; so the product vanishes. The entries of $\rho$ and of the $\psi_{v,i}$, which are entries of conjugates of $\mathrm{mat}\circ r$ by fixed matrices, are locally constant modulo every power of $\mathfrak m_C$, which is (d).
\item Products of two elements of $\ker p$ vanish, as $\ker p\subseteq\mathfrak m_{C'}$.
\begin{enumerate}
\item[(a)] \emph{The conjugation.} Let $r'(c_\Gamma)=(X,-1)$. As $c_\Gamma^2=1$, $X\theta(X)=1$, that is, $X=J_n{}^tXJ_n$; and $p(X)=1$, so $X=1+N$ with entries of $N$ in $\ker p$, $N^2=0$ and $J_n{}^tNJ_n=N$. As $2$ is a unit of $C'$, let $Y=N/2$ and $h_1=1-Y$, with inverse $1+Y$ and $J_n{}^th_1J_n=h_1$. Then $\mathrm{mat}((h_1,1)(X,-1)(h_1,1)^{-1})=h_1X\theta(h_1^{-1})=(1-Y)(1+N)(1-Y)=1$. Let $r_1=(h_1,1)\,r'\,(h_1,1)^{-1}$, with $\mathrm{mat}(r_1(c_\Gamma))=1$; as in (2), $\mathrm{mat}(r_1(g^*))=J_n{}^t\mathrm{mat}(r_1(g))J_n$ for $g\in\Gamma_0$.
\item[(b)] \emph{The diagonal matrix.} Let $Z=\mathrm{mat}(r_1(g_0))-\mathrm{mat}(r_1(g_0^*))$. Then $-J_n{}^tZJ_n=Z$, so $Z_{n+1-j,n+1-i}=-Z_{ij}$; the image of $Z$ under $p$ is $\rho(g_0)-\rho(g_0^*)$, whose entries off the diagonal vanish by condition (b), so these entries of $Z$ lie in $\ker p$; and the $Z_{ii}-Z_{jj}$, $i\ne j$, are units, by hypothesis. Let $A_{ij}=Z_{ij}(Z_{ii}-Z_{jj})^{-1}$ for $i\ne j$ and $A_{ii}=0$. Then the entries of $A$ lie in $\ker p$, $A^2=0$, and $J_n{}^tAJ_n=-A$, so $h_2=1+A$ has inverse $1-A$ and $\theta(h_2)=h_2$. The entry $(i,j)$, $i\ne j$, of $h_2Zh_2^{-1}=Z+AZ-ZA-AZA$ is $Z_{ij}+A_{ij}Z_{jj}-Z_{ii}A_{ij}=0$, the other terms being sums of products of two elements of $\ker p$.
\item[(c)] \emph{Conclusion.} Let $h=h_2h_1$ and $r_2=(h,1)\,r'\,(h,1)^{-1}$, with $g''_v=g'_vh^{-1}$. Then $(r_2,(g''_v))$ is a continuous triangulation over $C'$; $\mathrm{mat}(r_2(c_\Gamma))=h_2\theta(h_2)^{-1}=1$ by (a) and $\theta(h_2)=h_2$; and the entries off the diagonal of $\mathrm{mat}(r_2(g_0))-\mathrm{mat}(r_2(g_0^*))=h_2Zh_2^{-1}$ vanish by (b). By (2), $\rho'=\mathrm{mat}\circ r_2$ on $\Gamma_0$, with the $\psi'$ of (2), satisfies conditions (a) to (d). As $p(h)=1$, the image of $(r_2,(g''_v))$ under $p$ is a triangulation of $(\rho,\psi)$, so $(\rho',\psi')$ maps to $(\rho,\psi)$.
\end{enumerate}
\end{enumerate}
\end{proof}
% (1) is pr_good (PrGood.lean: two_ne, diff_inj with hs, the field f3 of Frame.Ok and pow_inj, then lieT and thetaLie case by case). (2) is pr_ofbox (PrOfbox.lean: glPart, diagHom, conjHom, hpol, ord_of_flag, cont_conj). (3) is pr_norm (PrNorm.lean): pr_norm1 (PrNorm1.lean) for (a), with cj, cj_map and cj_cont; star_mat, hθ, hZp, hoff and hdist, then pr_norm2 (JconjApply.lean: A, hAA, hAr, the entries of (1 + A) Z (1 - A)) for (b); hB2c, hB2diag, pr_ofbox and of_unique for (c). The Lean's hypothesis is (rho.bar O).Good, of which only the fields two and zd are used here.

\begin{lemma}\label{sha-automorphy-unramified-near-completion-tangent}
In the situation of Lemma~\ref{sha-automorphy-unramified-near-scale}, let $\mathcal O'\supseteq\mathcal O_1$ be the ring of integers of a finite unramified extension of the fraction field of $\mathcal O_1$, with residue field $k'$, and let $R$ be a complete local noetherian $\mathcal O'$-algebra with residue field $k'$. Let $(\rho_R,\psi_R)$, with values in $R$, satisfy conditions (a) to (d) of Definition~\ref{sha-descent-link}(2), and suppose that it is universal: for every complete local noetherian $\mathcal O'$-algebra $C$ with residue field $k'$, every $(\rho',\psi')$ with values in $C$ that satisfies (a) to (d) and has the same reduction $(\bar\rho_R,\bar\psi_R)$ is the image of $(\rho_R,\psi_R)$ under a unique $\mathcal O'$-algebra map $R\to C$. Suppose that the diagonal entries of $\bar\rho_R(g)$ and the values $\bar\psi_{R,v,i}(h)$ are $\bar\Psi_\Gamma(g)^{w_{n,i}}$ and $\bar\Psi_\Gamma(\mathrm{loc}_vh)^{w_{n,i}}$, viewed in $k'$, and let $g_v\in GL_n(R)$ be such that $g_v\rho_R(\mathrm{loc}_vh)g_v^{-1}$ is upper triangular with diagonal entries $\psi_{R,v,1}(h),\dots,\psi_{R,v,n}(h)$ for $h\in G_v$.
Let $r_R:\Gamma\to GL_n(R)\rtimes\{\pm1\}$, where $-1$ acts by $g\mapsto J_n{}^tg^{-1}J_n$, be the homomorphism with $r_R(g)=(\rho_R(g),1)$ for $g\in\Gamma_0$ and $r_R(c_\Gamma)=(1,-1)$; it is one by condition (a) and $c_\Gamma^2=1$. Let $M(s)$ be the matrix part of $r_R(s)$ and $\bar M(s)$ its reduction, let $\Gamma$ act on $W=M_n(k')$ by $(g,1)\cdot X=gXg^{-1}$ and $(g,-1)\cdot X=-gJ_n{}^tXJ_ng^{-1}$ through the reduction $\bar r$ of $r_R$, and let $b_v$ be the space of $X\in W$ with $\bar g_vX\bar g_v^{-1}$ upper triangular. Let $t_R$ be the $k'$-vector space of $\mathcal O'$-linear maps $\delta:R\to k'$ with $\delta(ab)=\bar a\delta(b)+\bar b\delta(a)$, and for $\delta\in t_R$ let $z_\delta(s)=\delta(M(s))\bar M(s)^{-1}$, with $\delta$ applied to each entry. Then:
\begin{enumerate}
\item for every continuous cocycle $z:\Gamma\to W$ and elements $X_v\in W$ with $z(\mathrm{loc}_vh)-(\mathrm{loc}_v(h)X_v-X_v)\in b_v$ for all $v$ and $h\in G_v$, there are $\delta\in t_R$ and $X\in W$ with $z_\delta(s)=z(s)+sX-X$ for all $s\in\Gamma$;
\item there is $\delta_0\in t_R$ with $z_{\delta_0}(s)=sT-T$ for all $s\in\Gamma$, where $T$ is the diagonal matrix with entries $w_{n,1},\dots,w_{n,n}$.
\end{enumerate}
\end{lemma}
\leanref{Development/AutomorphicArithmetic/AutomorphicVarieties/FrobeniusSmoothVectors/PrTsurj.lean}{UpperTheory.V.Fr.pr_tsurj}
\leanref{Development/AutomorphicArithmetic/AutomorphicVarieties/FrobeniusSmoothVectors/GgMapOfGL.lean}{UpperTheory.V.Fr.pr_tor}
\begin{proof}
Write $w_i=w_{n,i}$, $T=\operatorname{diag}(w_1,\dots,w_n)$, $\bar r$ for the reduction of $r_R$ and $k'[\epsilon]$ for the dual numbers over $k'$, a complete local noetherian $\mathcal O'$-algebra with residue field $k'$. With $G_n(R)=GL_n(R)\rtimes\{\pm1\}$ as in Definition~\ref{sha-automorphy-unramified-near-completion-tangent-triangulation}, $(r_R,(g_v))$ is a triangulation of $(\rho_R,\psi_R)$, continuous by condition (d), as $\mathrm{mat}(r_R(gc_\Gamma))=\rho_R(g)$ for $g\in\Gamma_0$; $W$, its action and the $b_v$ are those of that Definition, and $\bar r$ has locally constant entries.
\begin{enumerate}
\item \emph{The reduction.} By Lemma~\ref{sha-automorphy-unramified-near-completion-tangent-normal}(1), applied to $(\bar\rho_R,\bar\psi_R)$ and the map $k_1\to k'$: $2\ne0$ in $k'$, the diagonal entries of $\bar\rho_R(g_0)-\bar\rho_R(g_0^*)$ are pairwise distinct, and $-J_n{}^tTJ_n=T$.
\item \emph{Points over the dual numbers.} Let $(\rho',\psi')$ with values in $k'[\epsilon]$ satisfy conditions (a) to (d) and reduce to $(\bar\rho_R,\bar\psi_R)$, and let $r':\Gamma\to G_n(k'[\epsilon])$ be a homomorphism with sign $\delta_\Gamma$, $\mathrm{mat}(r'(c_\Gamma))=1$ and $\mathrm{mat}(r'(g))=\rho'(g)$ for $g\in\Gamma_0$. Then there is $\delta\in t_R$ such that the coefficient of $\epsilon$ in $\mathrm{mat}(r'(s))$ is $z_\delta(s)\bar M(s)$ for every $s\in\Gamma$. Indeed, the universal property gives an $\mathcal O'$-algebra map $f:R\to k'[\epsilon]$ that carries $(\rho_R,\psi_R)$ to $(\rho',\psi')$. Write $f(a)=\bar a+\epsilon\delta(a)$; then $\delta$ is $\mathcal O'$-linear and $\delta(ab)=\bar a\delta(b)+\bar b\delta(a)$, as $f$ is multiplicative, so $\delta\in t_R$. The homomorphisms $f\circ r_R$ and $r'$ have the same sign and agree at $c_\Gamma$ and on $\Gamma_0$, so they are equal, as $\Gamma=\Gamma_0\cup\Gamma_0c_\Gamma$. Hence $\mathrm{mat}(r'(s))=f(M(s))=(1+\epsilon z_\delta(s))\bar M(s)$.
\item \emph{Proof of (1).} Let $(z,(X_v))$ be as in (1). The map $\mu(c)=\epsilon c$ from $k'$ to $k'[\epsilon]$ satisfies the hypotheses of Lemma~\ref{sha-automorphy-unramified-near-completion-tangent-correct} for the reduction map $p:k'[\epsilon]\to k'$. Apply that lemma to the triangulation $(\bar r,(\bar g_v))$ over $k'$, with the lifts $\tilde r(s)=\bar M(s)$, $\tilde g_v=\bar g_v$ and $\tilde t_v(h)=\bar g_v\bar M(\mathrm{loc}_vh)\bar g_v^{-1}$, read in $k'\subseteq k'[\epsilon]$, and with $\beta=-z$, $Y_v=-X_v$ and $a_v=0$: condition (i) holds as $\bar r$ is a homomorphism and $dz=0$, $z$ being a cocycle; (ii) is clear; and (iii) is the condition on $(z,(X_v))$. It gives a continuous triangulation $(r',(g'_v))$ over $k'[\epsilon]$ that reduces to $(\bar r,(\bar g_v))$, with $\mathrm{mat}(r'(s))=(1+\epsilon z(s))\bar M(s)$. By (1), Lemma~\ref{sha-automorphy-unramified-near-completion-tangent-normal}(3), applied to $p$, to the pair $(\bar\rho_R,\bar\psi_R)$, of which $(\bar r,(\bar g_v))$ is a triangulation, and to $(r',(g'_v))$, gives $h=1+\epsilon Y$, with $Y\in W$, and $(\rho',\psi')$ over $k'[\epsilon]$ that satisfies conditions (a) to (d) and reduces to $(\bar\rho_R,\bar\psi_R)$, such that $x\,r'\,x^{-1}$, with $x=(h,1)$, has matrix $1$ at $c_\Gamma$ and $\rho'$ on $\Gamma_0$. By step (1) of the proof of Lemma~\ref{sha-automorphy-unramified-near-completion-tangent-correct}, $x\,r'(s)\,x^{-1}=(1+\epsilon(Y-s\cdot Y),1)\,r'(s)$, whose matrix is $(1+\epsilon(z(s)+Y-s\cdot Y))\bar M(s)$. So (2) gives $\delta\in t_R$ with $z_\delta(s)=z(s)+s\cdot X-X$ for $X=-Y$.
\item \emph{Proof of (2).} Read $\bar r$ and the $\bar g_v$ over $k'[\epsilon]$, let $x=(1-\epsilon T,1)$, $r'(s)=x\,\bar r(s)\,x^{-1}$ and $g'_v=\bar g_v(1+\epsilon T)$. Then $(r',(g'_v))$ is a triangulation, as $g'_v\,\mathrm{mat}(r'(\mathrm{loc}_vh))\,g_v'^{-1}=\bar g_v\bar M(\mathrm{loc}_vh)\bar g_v^{-1}$; it is continuous, its entries being locally constant; and by step (1) of the proof of Lemma~\ref{sha-automorphy-unramified-near-completion-tangent-correct}, $\mathrm{mat}(r'(s))=(1+\epsilon(s\cdot T-T))\bar M(s)$. As $c_\Gamma\cdot T=-J_n{}^tTJ_n=T$ by (1) and $\bar M(c_\Gamma)=1$, $\mathrm{mat}(r'(c_\Gamma))=1$. For $g\in\Gamma_0$, $\mathrm{mat}(r'(g))=(1-\epsilon T)\bar\rho_R(g)(1+\epsilon T)$, conjugate by a diagonal matrix; so the entries off the diagonal of $\mathrm{mat}(r'(g_0))-\mathrm{mat}(r'(g_0^*))$ vanish, by condition (b). By Lemma~\ref{sha-automorphy-unramified-near-completion-tangent-normal}(2), $\mathrm{mat}\circ r'$ on $\Gamma_0$, with the characters given by the $g'_v$, satisfies conditions (a) to (d); it reduces to $(\bar\rho_R,\bar\psi_R)$, as $x$ reduces to $1$. So (2) gives $\delta_0\in t_R$ with $z_{\delta_0}(s)=s\cdot T-T$.
\end{enumerate}
\end{proof}
% pr_tsurj (PrTsurj.lean): pr_good for (1) (the hypothesis (rho.bar O').Good and thetaLie T = T); pr_dof (ResDual.lean, with res_dual, pr_ext and hU.univ) for (2); for (3), pr_open (local constancy of the reduced Box), pr_tbox (PrTbox.lean: pr_correct for the trivial SetLift of B.bar with beta = -z, Y = -X_v, a = 0, and mu = inr), pr_norm, pr_dof, pr_sqz for h = e Y and hconj, and X = -Y. pr_tor (GgMapOfGL.lean) for (4): h = e(-T), B' = conjBy h B_0 with g v * h^-1, hform, hc (with thetaLie T = T), hdiagE and hd, hcont (lc_of_ker with pr_open), pr_ofbox, of_unique and pr_dof. The Lean's hypotheses are rho.Univ O' (the universality), a Box B with B.Of rho and B.Cont (r_R and the g_v), and (rho.bar O').Good.

\begin{definition}\label{sha-automorphy-unramified-near-completion-obstructions}
Let $\mathcal O'$ be a complete local noetherian ring with residue field $k'$, $R$ a complete local noetherian $\mathcal O'$-algebra with residue field $k'$, and $H$ a $k'$-vector space. Then $H$ is an \emph{obstruction space} for $R$ if for every complete local noetherian $\mathcal O'$-algebra $C'$ with residue field $k'$ and every surjective $\mathcal O'$-algebra map $p:C'\to R$ with $\mathfrak m_{C'}\ker p=0$, there is a $k'$-linear map $\mathrm{ob}$ from the space of $\mathcal O'$-linear maps $\ker p\to k'$ to $H$ with the following property. If $\mathrm{ob}(u)=0$, $q:C'\to C''$ is a surjective $\mathcal O'$-algebra map onto a complete local noetherian $\mathcal O'$-algebra with residue field $k'$ that kills the $y\in\ker p$ with $u(y)=0$, and $p'':C''\to R$ is an $\mathcal O'$-algebra map with $p''\circ q=p$, then there is an $\mathcal O'$-algebra map $s:R\to C''$ with $p''\circ s=\mathrm{id}_R$.
\end{definition}
\leanref{Development/AutomorphicArithmetic/AutomorphicVarieties/FrobeniusSmoothVectors/PresDefs.lean}{UpperTheory.V.Fr.HasObs}

\begin{lemma}\label{sha-automorphy-unramified-near-completion-obstruction}
Let $\mathcal O'$, $k'$, $R$, $(\rho_R,\psi_R)$, $g_v$, $\Gamma$ acting on $W$, and the $b_v$ be as in Lemma~\ref{sha-automorphy-unramified-near-completion-tangent}. Then $H^2_b(\Gamma,W)$, in the notation of Definition~\ref{sha-descent-residual}, is an obstruction space for $R$ (Definition~\ref{sha-automorphy-unramified-near-completion-obstructions}).
\end{lemma}
\leanref{Development/AutomorphicArithmetic/AutomorphicVarieties/FrobeniusSmoothVectors/PrSetlift.lean}{UpperTheory.V.Fr.pr_setlift}
\leanref{Development/AutomorphicArithmetic/AutomorphicVarieties/FrobeniusSmoothVectors/ValSubOneMap.lean}{UpperTheory.V.Fr.pr_obs}
\begin{proof}
Let $r_R$ and $M(s)$ be as in Lemma~\ref{sha-automorphy-unramified-near-completion-tangent}; as in its proof, $(r_R,(g_v))$ is a continuous triangulation of $(\rho_R,\psi_R)$. Let $C'$ be a complete local noetherian $\mathcal O'$-algebra with residue field $k'$ and $p:C'\to R$ a surjective $\mathcal O'$-algebra map with $\mathfrak m_{C'}K=0$, where $K=\ker p$; then $K\subseteq\mathfrak m_{C'}$, and $K$ is a $k'$-vector space.
\begin{enumerate}
\item \emph{Continuous lifts.} By the Artin-Rees lemma \cite[Theorem~8.5]{Mat86} there is $N_0$ with $K\cap\mathfrak m_{C'}^{N_0}\subseteq\mathfrak m_{C'}K=0$. As $p$ maps $\mathfrak m_{C'}^N$ onto $\mathfrak m_R^N$, there are $\sigma_0:R\to C'$, with $\sigma_0(c)$ depending only on $c$ modulo $\mathfrak m_R^{N_0}$ and $c-p(\sigma_0(c))\in\mathfrak m_R^{N_0}$, and for each $c$ an element $w(c)\in\mathfrak m_{C'}^{N_0}$ with $p(w(c))=c-p(\sigma_0(c))$; let $\sigma=\sigma_0+w$. Then $p\circ\sigma=\mathrm{id}$. If $c_1-c_2\in\mathfrak m_R^N$ with $N\ge N_0$, then $c_1-c_2=p(z)$ with $z\in\mathfrak m_{C'}^N$, and $\sigma(c_1)-\sigma(c_2)-z=w(c_1)-w(c_2)-z$ lies in $K\cap\mathfrak m_{C'}^{N_0}=0$; so $\sigma(c_1)-\sigma(c_2)\in\mathfrak m_{C'}^N$. Applying $\sigma$ to the entries, an entry that vanishes identically being lifted by $0$, gives $\tilde r:\Gamma\to GL_n(C')$, $\tilde g_v$ and $\tilde t_v:G_v\to GL_n(C')$ lifting $\mathrm{mat}\circ r_R$, the $g_v$ and $h\mapsto g_vM(\mathrm{loc}_vh)g_v^{-1}$, locally constant modulo every power of $\mathfrak m_{C'}$, with $\tilde t_v(h)$ upper triangular; these matrices are invertible, as their determinants lift units and $C'$ is local.
\item \emph{The cochains.} Let $\mathrm{om}(s,t)=\tilde r(s)\theta^{\delta_\Gamma(s)}(\tilde r(t))\tilde r(st)^{-1}-1$ and $a_v(h)=\tilde r(\mathrm{loc}_vh)\bigl(\tilde g_v^{-1}\tilde t_v(h)\tilde g_v\bigr)^{-1}-1$, with $\theta^{\pm1}$ as in Lemma~\ref{sha-automorphy-unramified-near-completion-tangent-correct}. Their entries lie in $K$, as $r_R$ is a homomorphism and $\tilde t_v$ lifts $g_vM(\mathrm{loc}_v\cdot)g_v^{-1}$. They are locally constant: modulo $\mathfrak m_{C'}^{N_0}$ they are, as $\tilde r$, $\tilde t_v$ and $\delta_\Gamma$ are, and two of their values that agree modulo $\mathfrak m_{C'}^{N_0}$ differ by a matrix over $K\cap\mathfrak m_{C'}^{N_0}=0$.
\item \emph{Matrices over $K$.} Let $u:K\to k'$ be $\mathcal O'$-linear, applied to matrices entrywise; then $u(az)=\bar au(z)$ for $a\in C'$, as $\mathfrak m_{C'}K=0$. For matrices $Z,Z'$ over $K$ and $A$ over $C'$: $ZZ'=0$; $AZ$ and $ZA$ are over $K$, with $u(AZ)=\bar Au(Z)$ and $u(ZA)=u(Z)\bar A$; and for $x\in G_n(C')$, $x\cdot Z$ is over $K$, $u(x\cdot Z)=\bar x\cdot u(Z)$ and $x(1+Z,1)x^{-1}=(1+x\cdot Z,1)$.
\item \emph{The cocycle.} For $u$ as in (3) let $\omega_u=u(\mathrm{om})$ and $\alpha_{u,v}=u(a_v)$, locally constant by (2); $u\mapsto(\omega_u,(\alpha_{u,v}))$ is $k'$-linear. With $x(s)=(\tilde r(s),\delta_\Gamma(s))$, $x(s)x(t)x(st)^{-1}=(1+\mathrm{om}(s,t),1)$, and the identity
\[x(s)x(t)x(st)^{-1}\cdot x(st)x(t')x(stt')^{-1}=x(s)\bigl(x(t)x(t')x(tt')^{-1}\bigr)x(s)^{-1}\cdot x(s)x(tt')x(stt')^{-1}\]
gives, by (3), $\mathrm{om}(s,t)+\mathrm{om}(st,t')=x(s)\cdot\mathrm{om}(t,t')+\mathrm{om}(s,tt')$; applying $u$, $s\cdot\omega_u(t,t')+\omega_u(s,tt')=\omega_u(st,t')+\omega_u(s,t)$.
\item \emph{The local condition.} Fix $v$ and $s,t\in G_v$, and let $T=\tilde t_v(s)\tilde t_v(t)\tilde t_v(st)^{-1}-1$, upper triangular with entries in $K$. With $\tilde r=\tilde r(\mathrm{loc}_vs)$, the definitions give
\[(1+\mathrm{om}(\mathrm{loc}_vs,\mathrm{loc}_vt))(1+a_v(st))=(1+\tilde ra_v(t)\tilde r^{-1})(1+a_v(s))(1+\tilde g_v^{-1}T\tilde g_v),\]
both sides being $\tilde r\,\tilde r(\mathrm{loc}_vt)\,\tilde g_v^{-1}\tilde t_v(st)^{-1}\tilde g_v$. By (3), $\mathrm{om}(\mathrm{loc}_vs,\mathrm{loc}_vt)-(\tilde ra_v(t)\tilde r^{-1}-a_v(st)+a_v(s))=\tilde g_v^{-1}T\tilde g_v$; applying $u$, $\omega_u(\mathrm{loc}_vs,\mathrm{loc}_vt)-d\alpha_{u,v}(s,t)=\bar g_v^{-1}u(T)\bar g_v$, with $d\alpha(s,t)=\alpha(s)+s\cdot\alpha(t)-\alpha(st)$. As $u(T)$ is upper triangular, this lies in $b_v$. By (4) and (5), $(\omega_u,(\alpha_{u,v}))$ defines a class $\mathrm{ob}(u)\in H^2_b(\Gamma,W)$, and $\mathrm{ob}$ is $k'$-linear.
\item \emph{A square-zero map.} Suppose $\mathrm{ob}(u)=0$: there are $\beta:\Gamma\to W$ locally constant and $Y_v\in W$ with $\omega_u=d\beta$ and $\alpha_{u,v}(h)-(\beta(\mathrm{loc}_vh)-(\mathrm{loc}_v(h)\cdot Y_v-Y_v))\in b_v$. Let $q:C'\to C''$ and $p'':C''\to R$ be as in Definition~\ref{sha-automorphy-unramified-near-completion-obstructions}. There is an additive $\mu:k'\to C''$ with $\mu(u(y))=q(y)$ for $y\in K$: if $u=0$, $q$ kills $K$ and $\mu=0$; otherwise $u$ maps onto $k'$, by (3), and $\mu(c)=q(y)$ for any $y\in K$ with $u(y)=c$ is well defined, as $q$ kills the $y\in K$ with $u(y)=0$. Then $\mu(k')=q(K)\subseteq\mathfrak m_{C''}$, $p''\circ\mu=0$, and for $a=q(b)$, $\mu(\bar ac)=q(by)=a\mu(c)$ when $u(y)=c$.
\item \emph{The section.} The lifts $q\circ\tilde r$, $q(\tilde g_v)$ and $q\circ\tilde t_v$ of $(r_R,(g_v))$ along $p''$ are locally constant modulo the powers of $\mathfrak m_{C''}$, and by (6) they satisfy conditions (i) to (iii) of Lemma~\ref{sha-automorphy-unramified-near-completion-tangent-correct} with $\beta$, the $Y_v$ and $a_v=\alpha_{u,v}$: $q(1+\mathrm{om})=1+\mu(\omega_u)=1+\mu(d\beta)$ and $q(1+a_v)=1+\mu(\alpha_{u,v})$. That lemma gives a continuous triangulation over $C''$ whose image under $p''$ is $(r_R,(g_v))$. As $q$ is surjective, $\ker p''=q(K)$ and $\mathfrak m_{C''}\ker p''=q(\mathfrak m_{C'}K)=0$. By step (1) of the proof of Lemma~\ref{sha-automorphy-unramified-near-completion-tangent}, Lemma~\ref{sha-automorphy-unramified-near-completion-tangent-normal}(3) applies and gives $(\rho'',\psi'')$ with values in $C''$ that satisfies conditions (a) to (d) and maps to $(\rho_R,\psi_R)$ under $p''$; so it has the reduction of $(\rho_R,\psi_R)$. The universal property gives $s:R\to C''$ carrying $(\rho_R,\psi_R)$ to $(\rho'',\psi'')$. Then $p''\circ s$ carries $(\rho_R,\psi_R)$ to itself, so it is the identity, by the uniqueness in the universal property.
\end{enumerate}
\end{proof}
% The hobs part of pres (HasObs.lean). (1) is pr_setlift (PrSetlift.lean: gl_lift, through pr_sec of InfPowBot.lean with inf_pow_bot, the Artin-Rees lemma Ideal.exists_pow_inf_eq_pow_smul); (2) is pr_omk (LcMatrix.lean: mem_ker, lc_ker, lc_gl, lc_hom); (3) is pr_evk (MapThetaLie.lean); (4) is borelProjection_isTwoCocycle (with grp_id) and (5) is pr_cocv (PrCocv.lean: upT_gl, ring_id), which with pr_omk put SetLift.co u in continuousCocycles2, so that ob is the composite of SetLift.co with the quotient coneH2; (6) is pr_mu (PrMu.lean) and (7) is pr_obs (ValSubOneMap.lean: push, push_om, push_a, pr_correct, pr_norm, hU.univ twice). The Lean takes O' local and R with IsCO O' R, rho.Univ, a Box with Of and Cont, and (rho.bar O').Good.

\begin{lemma}[Mazur {\cite{Maz89}}]\label{sha-automorphy-unramified-near-completion-mazur}
Let $\mathcal O'$ be a complete discrete valuation ring with finite residue field $k'$, $R$ a complete local noetherian $\mathcal O'$-algebra with residue field $k'$, $t_R$ the $k'$-vector space of $\mathcal O'$-linear maps $\delta:R\to k'$ with $\delta(ab)=\bar a\delta(b)+\bar b\delta(a)$, and $H$ a finite-dimensional obstruction space for $R$ (Definition~\ref{sha-automorphy-unramified-near-completion-obstructions}). Then $R$ is the quotient of the power series ring over $\mathcal O'$ in $\dim t_R$ variables by an ideal generated by $\dim H$ elements.
\end{lemma}
\leanref{Development/AutomorphicArithmetic/AutomorphicVarieties/FrobeniusSmoothVectors/PrMazur.lean}{UpperTheory.V.Fr.pr_mazur}
\leanref{Development/AutomorphicArithmetic/AutomorphicVarieties/CommutativeAlgebra/ObsPres.lean}{UpperTheory.V.Alg.obs_pres}
% Mazur proves this for universal deformation rings, with the obstruction space H^2 of the adjoint representation; the statement here, for any obstruction space, is the same argument. The Lean proves it in its own way: pr_mazur (PrMazur.lean) from obs_pres (CommutativeAlgebra/ObsPres.lean), pr_dualb (PrDualb.lean) and pr_tan (TanFinite.lean), which identifies t_R with the O'-algebra maps from R to the dual numbers over k'.

\begin{lemma}\label{sha-automorphy-unramified-near-completion-deformation}
In the situation of Lemma~\ref{sha-automorphy-unramified-near-scale}, let $x$ be a maximal ideal of $A_n^\circ$ that contains the maximal ideal of $B$, let $\mathcal O'$ be the ring of integers of a finite unramified extension of the fraction field of $\mathcal O_1$, with residue field $k'$, and let $e:A_n^\circ/x\to k'$ be an $\mathcal O_1$-algebra map. Let $y$ be the kernel of the map $D=\mathcal O'\otimes_{\mathcal O_1}A_n^\circ\to k'$, $c\otimes a\mapsto\bar c\,e(a\bmod x)$, let $R_x$ be the completion of $D$ at the maximal ideal $y$ and $\iota:A_n^\circ\to R_x$ the map $a\mapsto1\otimes a$, and let $\rho_x=\iota\circ\rho^{\mathrm u}$ and $\psi_{x,v,i}=\iota\circ\psi^{\mathrm u}_{v,i}$, with reductions $\bar\rho_x$ and $\bar\psi_{x,v,i}$ modulo $\mathfrak m_{R_x}$. Let $t_x$ be the $k'$-vector space of $\mathcal O'$-linear maps $\delta:R_x\to k'$ with $\delta(ab)=\bar a\delta(b)+\bar b\delta(a)$, where $\bar a$ is the residue of $a$. Then:
\begin{enumerate}
\item $(\rho_x,\psi_x)$ satisfies conditions (a) to (d) of Definition~\ref{sha-descent-link}(2); and for every complete local noetherian $\mathcal O'$-algebra $C$ with residue field $k'$, every $(\rho',\psi')$ with values in $C$ that satisfies (a) to (d) and reduces to $(\bar\rho_x,\bar\psi_x)$ is the image of $(\rho_x,\psi_x)$ under a unique $\mathcal O'$-algebra map $R_x\to C$;
\item there are $g_v\in GL_n(R_x)$, for the places $v$ of $L^+$ above $\ell$, such that $g_v\rho_x(\mathrm{loc}_vh)g_v^{-1}$ is upper triangular with diagonal entries $\psi_{x,v,1}(h),\dots,\psi_{x,v,n}(h)$ for every $h\in G_v$.
\end{enumerate}
Let $\bar r:\Gamma\to GL_n(k')\rtimes\{\pm1\}$, where $-1$ acts by $g\mapsto J_n{}^tg^{-1}J_n$, be the homomorphism with $\bar r(g)=(\bar\rho_x(g),1)$ for $g\in\Gamma_0$ and $\bar r(c_\Gamma)=(1,-1)$. Let $\Gamma$ act on $W=M_n(k')$ through $\bar r$, by $(g,1)\cdot X=gXg^{-1}$ and $(g,-1)\cdot X=-gJ_n{}^tXJ_ng^{-1}$, and let $b_v$ be the space of $X\in W$ with $\bar g_vX\bar g_v^{-1}$ upper triangular. Then:
\begin{enumerate}
\setcounter{enumi}{2}
\item the kernel of the action of $\Gamma$ on $W$ is open, each $b_v$ is stable under $\mathrm{loc}_v(G_v)$, and there is no nonzero $\Gamma$-equivariant linear map from $W$ to $k'(\omega_\Gamma)$;
\item $t_x$ is finite-dimensional, and $\dim H^1_b(\Gamma,W)+1\le\dim t_x+\dim H^0_b(\Gamma,W)$, in the notation of Definition~\ref{sha-descent-residual};
\item if $H^2_b(\Gamma,W)$ is finite-dimensional, then $R_x$ is the quotient of the power series ring over $\mathcal O'$ in $\dim t_x$ variables by an ideal generated by $\dim H^2_b(\Gamma,W)$ elements.
\end{enumerate}
\end{lemma}
\leanref{Development/AutomorphicArithmetic/AutomorphicVarieties/FrobeniusSmoothVectors/LocRho.lean}{UpperTheory.V.Fr.loc_rho}
\leanref{Development/AutomorphicArithmetic/AutomorphicVarieties/FrobeniusSmoothVectors/DefIso.lean}{UpperTheory.V.Fr.def_iso}
\leanref{Development/AutomorphicArithmetic/AutomorphicVarieties/FrobeniusSmoothVectors/ConjVal.lean}{UpperTheory.V.Fr.pr_box}
\leanref{Development/AutomorphicArithmetic/AutomorphicVarieties/FrobeniusSmoothVectors/PrDual.lean}{UpperTheory.V.Fr.pr_dual}
\leanref{Development/AutomorphicArithmetic/AutomorphicVarieties/FrobeniusSmoothVectors/PrCount.lean}{UpperTheory.V.Fr.pr_count}
\leanref{Development/AutomorphicArithmetic/AutomorphicVarieties/FrobeniusSmoothVectors/PrMazur.lean}{UpperTheory.V.Fr.pr_mazur}
\begin{proof}
Write $w_i=w_{n,i}$, and $\bar\Psi_\Gamma$ for its composite with the map of residue fields $k_1\to k'$.
\begin{enumerate}
\item \emph{The pair $(\rho_x,\psi_x)$.} It is the image of the universal linked point under the ring map $\iota$, so it satisfies conditions (a) to (c) of Definition~\ref{sha-descent-link}(2). By step (5) of the proof of Lemma~\ref{sha-automorphy-unramified-near-small-continuity}, $g\mapsto\rho^{\mathrm u}(g)_{ij}$ is locally constant modulo $\mathfrak m_B^NA_n^\circ$ for every $N$, and by Lemma~\ref{sha-automorphy-unramified-near-tangent-adic} so is $h\mapsto\psi^{\mathrm u}_{v,i}(h)$ modulo $\mathfrak m_B^N$. As $x$ contains $\mathfrak m_B$ and $\iota^{-1}(\mathfrak m_{R_x})=x$, $\iota$ maps $\mathfrak m_B$ into $\mathfrak m_{R_x}$, and condition (d) follows.
\item \emph{Its reduction.} For $b\in B$ there is $a\in\mathcal O_1$ with $b-a\in\mathfrak m_B\subseteq x$, so $\iota(b)-a\in\mathfrak m_{R_x}$: the residue of $\iota(b)$ is the image in $k'$ of the residue of $b$. As the diagonal entries of $\rho^{\mathrm u}$ and the values of the $\psi^{\mathrm u}_{v,i}$ lie in $B$, condition (e) shows that the diagonal entries of $\bar\rho_x(g)$ and the values $\bar\psi_{x,v,i}(h)$ are $\bar\Psi_\Gamma(g)^{w_i}$ and $\bar\Psi_\Gamma(\mathrm{loc}_vh)^{w_i}$.
\item \emph{Proof of (1).} Let $C$ and $(\rho',\psi')$ be as in (1). By (2), Lemma~\ref{sha-automorphy-unramified-near-completion-points}, applied to $C$ as an $\mathcal O_1$-algebra, gives an $\mathcal O_1$-algebra map $\xi:A_n^\circ\to C$ that carries $(\rho^{\mathrm u},\psi^{\mathrm u})$ to $(\rho',\psi')$ and $\mathfrak m_B$ into $\mathfrak m_C$. The reductions of $\xi$ and $\iota$, as maps $A_n^\circ\to k'$, carry $(\rho^{\mathrm u},\psi^{\mathrm u})$ to the same pair and $\mathfrak m_B$ to $0$, so they are equal by Lemma~\ref{sha-automorphy-unramified-near-completion-unique}. Hence the $\mathcal O'$-algebra map $D\to C$, $c\otimes a\mapsto c\,\xi(a)$, carries $y$ into $\mathfrak m_C$; as $C$ is complete, it extends to an $\mathcal O'$-algebra map $g:R_x\to C$ \cite[\S8]{Mat86}, and $g\circ\iota=\xi$, so $g$ carries $(\rho_x,\psi_x)$ to $(\rho',\psi')$. If $g'$ is another such map, $g'\circ\iota$ and $\xi$ carry $(\rho^{\mathrm u},\psi^{\mathrm u})$ to the same pair and $\mathfrak m_B$ into $\mathfrak m_C$, as $g'$ is local; so $g'\circ\iota=\xi$ by Lemma~\ref{sha-automorphy-unramified-near-completion-unique}, $g'$ and $g$ agree on $D$, and $g'=g$ by the uniqueness of the extension to the completion.
\item \emph{Proof of (2).} This is Lemma~\ref{sha-automorphy-unramified-near-completion-flag}, applied to the local ring $R_x$ and $f=\iota$.
\item \emph{Proof of (3): openness and stability.} By (d), the entries of $\bar\rho_x$ are locally constant on $\Gamma_0$, which is open in $\Gamma$ (residual condition (4) of Definition~\ref{sha-descent-residual}). So the subgroup of the $g\in\Gamma_0$ with $\bar\rho_x(g)=1$, which acts trivially on $W$, is open. For $h\in G_v$, $\bar r(\mathrm{loc}_vh)=(\bar\rho_x(\mathrm{loc}_vh),1)$, and $U=\bar g_v\bar\rho_x(\mathrm{loc}_vh)\bar g_v^{-1}$ is upper triangular by (2); if $\bar g_vX\bar g_v^{-1}$ is upper triangular, so is $\bar g_v\bigl(\bar\rho_x(\mathrm{loc}_vh)X\bar\rho_x(\mathrm{loc}_vh)^{-1}\bigr)\bar g_v^{-1}=U\bar g_vX\bar g_v^{-1}U^{-1}$. So $b_v$ is stable under $\mathrm{loc}_v(G_v)$.
\item \emph{A triangular reduction.} Either every entry of $\rho^{\mathrm u}$ below the diagonal lies in $x$, or every entry above it does. Otherwise there are an entry $e$ below the diagonal, of weight $-p<0$, and an entry $e'$ above it, of weight $q>0$, outside $x$; the product $e^qe'^p$ has weight $0$ and an entry off the diagonal, so it lies in $\mathfrak m_B\subseteq x$ (step (2) of the proof of Lemma~\ref{sha-automorphy-unramified-near-scale}), against the primality of $x$. As $\iota^{-1}(\mathfrak m_{R_x})=x$, $\bar\rho_x(g)$ is upper triangular for every $g$, or lower triangular for every $g$.
\item \emph{Proof of (3): duality.} Let $f:W\to k'(\omega_\Gamma)$ be $\Gamma$-equivariant, and $\Phi_{ij}=f(E_{ij})$ for the matrix units $E_{ij}$. For $g\in\Gamma_0$ and $X\in W$, $f(\bar\rho_x(g)X)=\omega_\Gamma(g)f(X\bar\rho_x(g))$; with $X=E_{ij}$ this reads
\[\textstyle\sum_p\bar\rho_x(g)_{pi}\Phi_{pj}=\omega_\Gamma(g)\sum_q\Phi_{iq}\bar\rho_x(g)_{jq}.\]
\emph{Diagonal case.} Suppose that $\bar\rho_x(g)_{ii}\Phi_{ij}=\omega_\Gamma(g)\Phi_{ij}\bar\rho_x(g)_{jj}$ for all $g\in\Gamma_0$, and $\Phi_{ij}\ne0$. By (2), $\omega_\Gamma\bar\Psi_\Gamma^{w_j-w_i}=1$ on $\Gamma_0$. If $w_j-w_i\in\{0,\pm1\}$, this contradicts residual condition (7). Otherwise $w_j-w_i=\pm2$ and $j=n+1-i$. Then $c_\Gamma$ acts on $W$ by $X\mapsto-J_n{}^tXJ_n$, which sends $E_{ij}$ to $-E_{ij}$, so $-\Phi_{ij}=\omega_\Gamma(c_\Gamma)\Phi_{ij}$ and $\omega_\Gamma(c_\Gamma)=-1$. For $g\notin\Gamma_0$, $gc_\Gamma\in\Gamma_0$, $\bar\Psi_\Gamma(c_\Gamma)=1$ and $\delta_\Gamma(g)=-1$, so $\omega_\Gamma(g)\bar\Psi_\Gamma(g)^{w_j-w_i}\delta_\Gamma(g)=1$; that is, $\omega_\Gamma\bar\Psi_\Gamma^{\pm2}\delta_\Gamma$ is trivial on $\Gamma$, which again contradicts residual condition (7). So $\Phi_{ij}=0$.
\emph{Induction.} Suppose that every $\bar\rho_x(g)$ is upper triangular (step (6)). By induction on $i$ upwards, and for each $i$ on $j$ downwards, the terms with $p<i$ and with $q>j$ in the identity above vanish, and those with $p>i$ and $q<j$ vanish by triangularity; so the diagonal case gives $\Phi_{ij}=0$. If every $\bar\rho_x(g)$ is lower triangular, the same induction runs on $i$ downwards and on $j$ upwards. So $\Phi=0$, and $f=0$.
\item \emph{Proof of (4): the tangent space.} Let $s_1,\dots,s_m$ generate the ideal $\mathfrak m_{R_x}$, as $R_x$ is noetherian. An element $\delta\in t_x$ vanishes on $\mathcal O'$, since $\delta(1)=2\delta(1)$, and $\delta(cz)=\bar c\delta(z)$ for $c\in R_x$ and $z\in\mathfrak m_{R_x}$. So $\delta$ is determined on $\mathfrak m_{R_x}$ by the $\delta(s_k)$, and on $R_x=\mathcal O'+\mathfrak m_{R_x}$. Hence $\delta\mapsto(\delta(s_k))_k$ is an injective linear map $t_x\to k'^m$, and $t_x$ is finite-dimensional.
\item \emph{Rigidity of the $b_v$.} Let $Y\in W$ with $\mathrm{loc}_v(h)\cdot Y-Y\in b_v$ for all $h\in G_v$; then $Y\in b_v$. Indeed, by residual condition (7) there is $h\in G_v$ with $\bar\Psi_\Gamma(\mathrm{loc}_vh)^2\ne1$, and then the diagonal entries $\bar\psi_{x,v,i}(h)=\bar\Psi_\Gamma(\mathrm{loc}_vh)^{w_i}$ of $U=\bar g_v\bar\rho_x(\mathrm{loc}_vh)\bar g_v^{-1}$ are pairwise distinct, as $w_i-w_j$ divides $2$. With $Z=\bar g_vY\bar g_v^{-1}$, the matrix $UZU^{-1}-Z$ is upper triangular, and so is $UZ-ZU=(UZU^{-1}-Z)U$. For $j<i$, by induction on $n-i+j$, the entry $(i,j)$ of $UZ-ZU$ is $(U_{ii}-U_{jj})Z_{ij}$ plus terms $U_{ic}Z_{cj}$ with $c>i$ and $Z_{ic}U_{cj}$ with $c<j$, which vanish by induction; so $Z_{ij}=0$, and $Y\in b_v$.
\item \emph{Proof of (4): the count.} By (1), (2) and step (2), Lemma~\ref{sha-automorphy-unramified-near-completion-tangent} applies to $R=R_x$, $(\rho_x,\psi_x)$ and the $g_v$; let $\delta\mapsto z_\delta$ and $\delta_0$ be as there, with $\bar r$ the reduction of $r_R$. Let $\mathcal C$ be the space of maps $s\mapsto sX-X$ from $\Gamma$ to $W$. A pair $(z,(X_v))$ that defines a class of $H^1_b(\Gamma,W)$ with $z\in\mathcal C$, say $z(s)=sX-X$, has $\mathrm{loc}_v(h)(X-X_v)-(X-X_v)\in b_v$ for all $h$, so $X-X_v\in b_v$ by (9), and the class is $0$. So $(z,(X_v))\mapsto z+\mathcal C$ induces an injective linear map from $H^1_b(\Gamma,W)$ to the space of maps modulo $\mathcal C$, whose image, by Lemma~\ref{sha-automorphy-unramified-near-completion-tangent}(1), lies in the image of $t_x$ under $\delta\mapsto z_\delta+\mathcal C$. With $K$ the kernel of this map on $t_x$, $\dim H^1_b(\Gamma,W)\le\dim t_x-\dim K$. Let $T=\operatorname{diag}(w_1,\dots,w_n)\ne0$. If $sT=T$ for every $s\in\Gamma$, then $T\in b_v$ by (9), so $T\in H^0_b(\Gamma,W)$ and $\dim H^0_b(\Gamma,W)\ge1$. Otherwise some $s$ has $z_{\delta_0}(s)=sT-T\ne0$, so $\delta_0$ is a nonzero element of $K$ by Lemma~\ref{sha-automorphy-unramified-near-completion-tangent}(2), and $\dim K\ge1$. In both cases $\dim H^1_b(\Gamma,W)+1\le\dim t_x+\dim H^0_b(\Gamma,W)$.
\item \emph{Proof of (5).} By Lemma~\ref{sha-automorphy-unramified-near-completion-obstruction}, applied as in (10), $H^2_b(\Gamma,W)$ is an obstruction space for $R_x$. If it is finite-dimensional, Lemma~\ref{sha-automorphy-unramified-near-completion-mazur} shows that $R_x$ is the quotient of the power series ring over $\mathcal O'$ in $\dim t_x$ variables by an ideal generated by $\dim H^2_b(\Gamma,W)$ elements.
\end{enumerate}
% pres (HasObs.lean) with def_iso. (1) and (2) of the text are loc_rho (LocRho.lean: cont_map with the fields of cont_u for (d), res_b for the reductions); (3) is def_iso (DefIso.lean): point_loc (Lemma sha-automorphy-unramified-near-completion-points), ext_A (Lemma sha-automorphy-unramified-near-completion-unique) twice and L.univ (cpl_univ, the universal property of the completion). (4) is the flag hypothesis of pr_box (Fr.flag, Lemma sha-automorphy-unramified-near-completion-flag); pr_box (ConjVal.lean) builds the Box (r_R and the g_v) from it. (5) is pr_open and pr_bstab (upT_mul); (6) is pr_res with the second part of loc_rho and Fr.types (IrrOfConn.lean, hcone.cyc); (7) is pr_dual (tri_zero, w_cases, thetaLie_single, ok.f2a, f2b, f2c). (8) is pr_tan (TanFinite.lean); (9) is pr_brig (upT_of_comm) with pr_good (PrGood.lean) for the distinct diagonal entries; (10) is pr_count (PrCount.lean) with pr_tsurj and pr_tor (Lemma sha-automorphy-unramified-near-completion-tangent); (11) is the HasObs part of pres (Lemma sha-automorphy-unramified-near-completion-obstruction) and pr_mazur (Lemma sha-automorphy-unramified-near-completion-mazur). Lean names: rho_x is L.rho (an NRep over L.R), the g_v and r-bar are the Box B of pres and Box.bar, W with its action and the b_v are (Box.bar L.O' B).ad and .b, t_x is Tan L.O' L.R. The inputs point_loc, pr_tsurj with pr_tor, and the HasObs part of pres are written out in Lemmas sha-automorphy-unramified-near-completion-points, -tangent and -obstruction.
\end{proof}

\begin{lemma}\label{sha-automorphy-unramified-near-completion}
In the situation of Lemma~\ref{sha-automorphy-unramified-near-scale}, let $d=[L^+:\mathbf Q]$, and for a prime $\eta$ of $A_n^\circ$ let $t(\eta)$ be the transcendence degree of the fraction field of $A_n^\circ/\eta$ over that of $B/(\eta\cap B)$. Let $x$ be a maximal ideal of $A_n^\circ$ that contains the maximal ideal of $B$. Then there are the ring of integers $\mathcal O'$ of a finite unramified extension of the fraction field of $\mathcal O_1$, a complete local noetherian $\mathcal O'$-algebra $R_x$ with the residue field of $\mathcal O'$, and an $\mathcal O_1$-algebra map $\iota:A_n^\circ\to R_x$ such that:
\begin{enumerate}
\item $\iota$ is flat, and $\iota^{-1}(\mathfrak m_{R_x})=x$;
\item $\dim R_x/P\le\dim B/(\iota^{-1}(P)\cap B)+t(\iota^{-1}(P))$ for every prime $P$ of $R_x$;
\item for some $N\ge nd+1$, $R_x$ is the quotient of the power series ring over $\mathcal O'$ in $N$ variables by an ideal generated by $N-nd-1$ elements.
\end{enumerate}
\end{lemma}
\leanref{Development/AutomorphicArithmetic/AutomorphicVarieties/FrobeniusSmoothVectors/ResEqOf.lean}{UpperTheory.V.Fr.loc_ex}
\leanref{Development/AutomorphicArithmetic/AutomorphicVarieties/FrobeniusSmoothVectors/LocFlat.lean}{UpperTheory.V.Fr.loc_flat}
\leanref{Development/AutomorphicArithmetic/AutomorphicVarieties/FrobeniusSmoothVectors/LocDim.lean}{UpperTheory.V.Fr.loc_dim}
\leanref{Development/AutomorphicArithmetic/AutomorphicVarieties/FrobeniusSmoothVectors/HasObs.lean}{UpperTheory.V.Fr.pres}
\begin{proof}
Write $A=A_n^\circ$. By Lemma~\ref{sha-automorphy-realization-presentation}(3) and (4), $B$ is a complete local noetherian $\mathcal O_1$-algebra with the residue field $k_1$ of $\mathcal O_1$, which is finite, and $A$ is a finitely generated $B$-algebra, hence noetherian. As $x$ is a proper ideal that contains $\mathfrak m_B$, $x\cap B=\mathfrak m_B$.
\begin{enumerate}
\item \emph{The residue field at $x$.} The field $A/x$ is a finitely generated algebra over the finite field $B/\mathfrak m_B=k_1$, so it is finite over $k_1$ \cite[Corollary~5.24]{AM69}.
\item \emph{The coefficient ring.} Let $g$ be the minimal polynomial over $k_1$ of a primitive element of $A/x$, and $f\in\mathcal O_1[X]$ a monic lift of $g$. Then $\mathcal O'=\mathcal O_1[X]/(f)$ is the ring of integers of a finite unramified extension of the fraction field of $\mathcal O_1$, with maximal ideal $\mathfrak m_{\mathcal O_1}\mathcal O'$ and residue field $k'=k_1[X]/(g)$ \cite[Chapter~III, \S5]{Ser79}; so there is an $\mathcal O_1$-algebra map $e:A/x\to k'$.
\item \emph{The ring $R_x$.} Let $D=\mathcal O'\otimes_{\mathcal O_1}A$, finite over $A$ and so noetherian, and $y$ the kernel of the surjective map $D\to k'$, $c\otimes a\mapsto\bar c\,e(a\bmod x)$; it is a maximal ideal, and its inverse image in $A$ is $x$, as $e$ is injective. Let $R_x$ be the completion of $D$ at $y$, that is, the $y$-adic completion of the localization $D_y$, and $\iota(a)=1\otimes a$. Then $R_x$ is a complete local noetherian $\mathcal O'$-algebra whose maximal ideal meets $D$ in $y$, with residue field $D/y\cong k'$ \cite[\S8]{Mat86}; so $\iota^{-1}(\mathfrak m_{R_x})=x$.
\item \emph{Flatness.} The ring $\mathcal O'$ is torsion-free over the discrete valuation ring $\mathcal O_1$, hence flat over it, so $A\to D$ is flat. The localization $D\to D_y$ is flat, and so is the completion $D_y\to R_x$ of a noetherian local ring \cite[Theorem~8.8]{Mat86}. So $\iota$ is flat, and with (3) this proves (1).
\item \emph{Dimensions.} Let $P$ be a prime of $R_x$, $p=P\cap D$ and $\eta=p\cap A=\iota^{-1}(P)$. As $D$ is finite over $A$, $D/p$ is integral over $A/\eta$, so $p\cap B=\eta\cap B$ and the transcendence degree of $D/p$ over $B/(p\cap B)$ is $t(\eta)$. The ring $R_x/P$ is a quotient of $R_x/pR_x$, the completion of the local ring $(D/p)_y$, which has the same dimension \cite[\S13]{Mat86}; so $\dim R_x/P\le\dim(D/p)_y$. The ring $D$ is a finitely generated $B$-algebra, $y$ is a maximal ideal with $y\cap B=\mathfrak m_B$, and $D/y$ is finite over $k_1$. The complete local ring $B$ is a quotient of a regular local ring \cite[Theorem~29.4]{Mat86}, hence universally catenary \cite[Theorem~17.9]{Mat86}, so the dimension formula \cite[Theorem~15.6]{Mat86} for the domain $D/p$, finitely generated over $B/(p\cap B)$, at the prime $y/p$ over $\mathfrak m_B/(p\cap B)$ gives $\dim(D/p)_y=\dim B/(p\cap B)+t(\eta)$. This proves (2).
\item \emph{The presentation.} Let $\bar r$, $W$, the $b_v$ and $t_x$ be as in Lemma~\ref{sha-automorphy-unramified-near-completion-deformation}, for $x$, $\mathcal O'$ and $e$. By its part (3), residual condition (10) of Definition~\ref{sha-descent-residual} applies to $k'$, $W$ and $b$: $H^2_b(\Gamma,W)$ is finite-dimensional, and with $h^i=\dim H^i_b(\Gamma,W)$,
\[h^2+h^0+\sum_vd_v\dim b_v\le h^1+\sum_w\dim\ker(c_w-1\mid W).\]
Each $b_v$ is a conjugate of the space of upper triangular matrices, of dimension $n(n+1)/2$, and $\sum_vd_v=d$. By residual condition (6) there are $d$ elements $c_w$, each with $c_w^2=1$ and $\delta_\Gamma(c_w)=-1$; so $\bar r(c_w)=(g,-1)$ with $(g,-1)^2=(gJ_n{}^tg^{-1}J_n,1)=1$, that is, $S=gJ_n$ is symmetric, and $c_w$ acts on $W$ by $X\mapsto-S\,{}^tX\,S^{-1}$. So $X$ is fixed by $c_w$ if and only if $XS$ is antisymmetric, and as $2\ne0$ in $k'$ the fixed space has dimension $n(n-1)/2$. Hence $h^2+h^0+nd\le h^1$, and by Lemma~\ref{sha-automorphy-unramified-near-completion-deformation}(4), $h^2+nd+1\le N$ with $N=\dim t_x$. By its part (5), $R_x$ is the quotient of the power series ring over $\mathcal O'$ in $N$ variables by an ideal generated by $h^2$ elements; adding $N-nd-1-h^2$ generators equal to $0$ gives (3).
\end{enumerate}
\end{proof}
% loc_ex (ResEqOf.lean): res_fin for (1), unr_ex (O' = AdjoinRoot f, with KStar.Binom.adjoinRoot_unramified) for (2), cpl_ex (cpl_is and cpl_univ, the completion of O' (x) A at ker phi with its universal property) for (3). loc_flat (LocFlat.lean): flat_dvr, loc_cpl and V.CA.c3a (localization and completion are flat), max_of_fin, for (4). loc_dim (LocDim.lean), the first part: dim_gen with td_fin, V.CA.c3c (dim of quotients of the completion) and V.CA.c3b (the dimension formula, which the Lean proves through cl_cmquot: a complete local ring is a quotient of a Cohen-Macaulay ring), for (5). pres (HasObs.lean) for (6): hF.ec with pr_bdim (dim b_v), pr_cfix (the fixed space of c_w), hsum1, hsum2 and hT, then HasObs padded by Fin (N - (n d + 1) - h^2) -> k' and pr_mazur. Lean names: R_x is L.R of a Space.Loc x, O' is L.O' (with L.unr), iota is L.iota with L.lies; (3) is Loc.Pres with N = N' + (n d + 1).

\begin{lemma}\label{sha-automorphy-unramified-near-heights}
In the situation of Lemma~\ref{sha-automorphy-unramified-near-scale}, let $d=[L^+:\mathbf Q]$, and let $M$ be an ideal of $B$ such that every off-triangular prime $\eta$ of $A_n^\circ$ has $M\subseteq\eta\cap B$, and $\dim B/Z\le1+nd$ for every prime $Z$ of $B$ that contains $M$. Let $\xi:A_n^\circ\to\bar{\mathbf Q}_\ell$ be an $\mathcal O_1$-algebra map with $|\xi(a)|\le1$ for $a\in A_n^\circ$ and $|\xi(b)|<1$ for $b$ in the maximal ideal of $B$, and let $\mathfrak q=\ker\xi$. Then:
\begin{enumerate}
\item $\mathfrak q$ has height at least $1+nd$;
\item if $\mathfrak q$ is off-triangular, then every associated prime of $A_n^\circ$ contained in $\mathfrak q$ is minimal, and $\dim B/(P\cap B)\ge1+nd$ for every minimal prime $P$ of $A_n^\circ$ contained in $\mathfrak q$.
\end{enumerate}
\end{lemma}
\leanref{Development/AutomorphicArithmetic/AutomorphicVarieties/ModularityArgument/LinkedUnramifiedChains/ScalarIrreducibility/IntOfNorm.lean}{UpperTheory.V.J.locAt}
\leanref{Development/AutomorphicArithmetic/AutomorphicVarieties/ModularityArgument/LinkedUnramifiedChains/ScalarIrreducibility/HeightLeOfDim.lean}{UpperTheory.V.J.htUp}
\leanref{Development/AutomorphicArithmetic/AutomorphicVarieties/ModularityArgument/LinkedUnramifiedChains/ScalarIrreducibility/HtOfSum.lean}{UpperTheory.V.J.htLow}
\leanref{Development/AutomorphicArithmetic/AutomorphicVarieties/ModularityArgument/LinkedUnramifiedChains/ScalarIrreducibility/CompDim.lean}{UpperTheory.V.J.compDim}
\begin{proof}
Write $A=A_n^\circ$, a noetherian ring by Lemma~\ref{sha-automorphy-realization-presentation}(4), and let $t$ be as in Lemma~\ref{sha-automorphy-unramified-near-completion}. A prime contained in an off-triangular prime is off-triangular, since the entries outside the larger prime lie outside the smaller one.
\begin{enumerate}
\item \emph{A maximal ideal.} Let $x$ be the set of $a\in A$ with $|\xi(a)|<1$. As $|\xi|\le1$ and the absolute value is ultrametric and multiplicative, $x$ is a prime ideal; it contains $\mathfrak q$ and the maximal ideal of $B$. Each $\xi(a)$ has absolute value at most $1$, so it is a root of a monic polynomial $f_a$ over $\mathbf Z_\ell$, and $f_a(a)\in\mathfrak q$. So $A/\mathfrak q$ is integral over the image of $\mathbf Z_\ell$, and the domain $A/x$ is integral over $\mathbf F_\ell$, as $\ell\in x$; hence $A/x$ is a field, and $x$ is maximal.
\item \emph{A prime over $\mathfrak q$.} Let $\mathcal O'$, $R=R_x$ and $\iota$ be given by Lemma~\ref{sha-automorphy-unramified-near-completion}. As $\iota$ is flat, it has the going-down property \cite[\S9]{Mat86}, so some prime of $R$ contained in $\mathfrak m_R$ lies over $\mathfrak q\subseteq x$; a prime $Q$ minimal over $\mathfrak qR$ inside it then has $\iota^{-1}(Q)=\mathfrak q$.
\item \emph{$\dim R/Q=1$.} By (1), $t(\mathfrak q)=0$. The ring $B/(\mathfrak q\cap B)$ is isomorphic to $\xi(B)$, a finitely generated $\mathcal O_1$-module by Lemma~\ref{sha-automorphy-unramified-near-tangent-points}, so its dimension is at most $1$. By Lemma~\ref{sha-automorphy-unramified-near-completion}(2), $\dim R/Q\le1$. And $\ell\in\mathfrak m_R$ while $\ell\notin Q$, as $\xi(\ell)\ne0$; so $Q\ne\mathfrak m_R$ and $\dim R/Q\ge1$.
\item \emph{Heights in $R$.} By Lemma~\ref{sha-automorphy-unramified-near-completion}(3), $R=S/(f_1,\dots,f_{N'})$ with $S$ the power series ring over $\mathcal O'$ in $N$ variables and $N'=N-nd-1$. For a prime $P$ of $R$ with inverse image $\tilde P$ in $S$: $S$ is a regular local ring of dimension $N+1$, so $\operatorname{ht}\tilde P+\dim R/P=N+1$ \cite[\S17]{Mat86}; and $\operatorname{ht}\tilde P\le\operatorname{ht}P+N'$, by Krull's height theorem in the localization of $S$ at $\tilde P$ \cite[\S13]{Mat86}. For $P=Q$ this gives $\operatorname{ht}\tilde Q=N$ by (3), and $\operatorname{ht}Q\ge nd+1$.
\item \emph{Proof of (1).} The map $A_{\mathfrak q}\to R_Q$ is a flat local homomorphism, and $R_Q/\mathfrak qR_Q$ has dimension $0$, as $Q$ is minimal over $\mathfrak qR$. So $\operatorname{ht}\mathfrak q=\operatorname{ht}Q\ge1+nd$ \cite[\S15]{Mat86}.
\item \emph{$\operatorname{ht}Q\le nd+1$ if $\mathfrak q$ is off-triangular.} Let $P\subseteq Q$ be a prime of $R$. Then $\iota^{-1}(P)\subseteq\mathfrak q$ is off-triangular, so $t(\iota^{-1}(P))\le1$ by Lemma~\ref{sha-automorphy-unramified-near-fibre}, and $\iota^{-1}(P)\cap B$ contains $M$, so its quotient has dimension at most $1+nd$. By Lemma~\ref{sha-automorphy-unramified-near-completion}(2), $\dim R/P\le nd+2$. A chain of primes $P_0\subsetneq\dots\subsetneq P_k=Q$ extends by $\mathfrak m_R\ne Q$ to a chain of length $k+1$ in $R/P_0$; so $k\le nd+1$.
\item \emph{Associated primes.} Let $\mathfrak q$ be off-triangular. By (4) and (6), $\dim R_Q=\operatorname{ht}Q\le nd+1=\dim S_{\tilde Q}-N'$. So the images of $f_1,\dots,f_{N'}$ in the regular local ring $S_{\tilde Q}$ are part of a system of parameters, hence a regular sequence, and $R_Q=S_{\tilde Q}/(f_1,\dots,f_{N'})$ is Cohen-Macaulay; so the associated primes of $R_Q$ are minimal \cite[\S17]{Mat86}, that is, every associated prime of $R$ contained in $Q$ is minimal. Let $P$ be an associated prime of $A$ contained in $\mathfrak q$. As $A/P$ embeds in $A$ and $\iota$ is flat, $R/PR$ embeds in $R$, and the primes minimal over $PR$ are associated primes of $R/PR$, hence of $R$ \cite[\S6]{Mat86}. One of them, $P'$, lies in $Q$, since $PR\subseteq Q$; so $P'$ is a minimal prime of $R$, and $P\subseteq\iota^{-1}(P')$. Let $P_0\subseteq P$ be a prime of $A$. Going down gives a prime of $R$ inside $P'$ that lies over $P_0$; it is $P'$, as $P'$ is minimal, so $P_0=\iota^{-1}(P')\supseteq P$. So $P$ is minimal.
\item \emph{Dimension of the components.} Let $\mathfrak q$ be off-triangular and $P\subseteq\mathfrak q$ a minimal prime of $A$. By going down some prime of $R$ inside $Q$ lies over $P$, and a minimal prime $P'$ of $R$ inside it has $\iota^{-1}(P')\subseteq P$, hence $\iota^{-1}(P')=P$. As $P$ is off-triangular, $t(P)\le1$ by Lemma~\ref{sha-automorphy-unramified-near-fibre}, and Lemma~\ref{sha-automorphy-unramified-near-completion}(2) gives $\dim R/P'\le\dim B/(P\cap B)+1$. By (4), with $\operatorname{ht}P'=0$, $\dim R/P'=N+1-\operatorname{ht}\tilde P'\ge N+1-N'=nd+2$. So $\dim B/(P\cap B)\ge nd+1$.
\end{enumerate}
\end{proof}
% locAt (IntOfNorm.lean): nx, nx_prime and nx_max (int_of_norm, int_pt) for (1); Fr.loc_ex, Fr.loc_flat and down (going down) for (2); dim_up (td_pt, Fr.loc_dim, dim_pt from coPt) and dim_ge with l_mem and l_not for (3); its conclusion is Over. htLow (HtOfSum.lean): Fr.pres, htPres (HtPres.lean) for (4), flatHt (FlatHt.lean) for (5), unmix (CmOfRegularLocal.lean) and flatAss (FlatAss.lean) for (7). htUp (HeightLeOfDim.lean): height_le_of_dim, Fr.fibre and the first part of Fr.loc_dim, with the hypotheses AllGood and the bound of DenseB, for (6). compDim (CompDim.lean): min_below, Fr.fibre, Fr.loc_dim, Fr.pres, htPres and dim_arith for (8).

\begin{lemma}[Chevalley]\label{sha-automorphy-unramified-near-determined-cofinal}
Let $B$ be a complete local noetherian ring with maximal ideal $\mathfrak m$, and let $I_0\supseteq I_1\supseteq\cdots$ be ideals of $B$ with $\bigcap_NI_N=0$. Then for every $k\ge0$ there is $N$ with $I_N\subseteq\mathfrak m^k$.
\end{lemma}
\begin{proof}
\begin{enumerate}
\item \emph{Stabilization.} For each $j$ the ring $B/\mathfrak m^j$ is artinian, as its only prime is $\mathfrak m/\mathfrak m^j$. So the decreasing images of the $I_N$ in it are eventually constant: there is $e_j$ such that for every $N\ge e_j$, every $a\in I_{e_j}$ is congruent modulo $\mathfrak m^j$ to some element of $I_N$.
\item \emph{A sequence.} Suppose that no $I_N$ lies in $\mathfrak m^k$. Let $e'_j$ be the $e_{k+j}$ of (1), and $E_j=\max(e'_0,\dots,e'_j)+j$, which increases with $j$ and has $E_j\ge e'_j$ and $E_j\ge j$. Choose $a_0\in I_{E_0}$ outside $\mathfrak m^k$. Given $a_j\in I_{E_j}\subseteq I_{e'_j}$ outside $\mathfrak m^k$, (1) gives $a_{j+1}\in I_{E_{j+1}}$ with $a_j-a_{j+1}\in\mathfrak m^{k+j}$; then $a_{j+1}\notin\mathfrak m^k$.
\item \emph{The limit.} So $a_m-a_{m'}\in\mathfrak m^{k+m}$ for $m\le m'$, and as $B$ is complete the $a_j$ converge to some $a\in B$ with $a_j-a\in\mathfrak m^j$ for all $j$. Fix $N$. For $j\ge N$, $a_j\in I_{E_j}\subseteq I_N$, so $a\in I_N+\mathfrak m^j$ for all $j$; by Krull's intersection theorem \cite[Theorem~8.10]{Mat86} for the finite $B$-module $B/I_N$, $a\in I_N$. Hence $a\in\bigcap_NI_N=0$. But $a_k-a\in\mathfrak m^k$ and $a_k\notin\mathfrak m^k$, a contradiction.
\end{enumerate}
\end{proof}
% The private cofinal of DerExt.lean (no leanref, as it is private): artinian_quot and stab for (1), the sequence a with E for (2), IsPrecomplete.prec and Ideal.iInf_pow_smul_eq_bot_of_isLocalRing for (3).

\begin{lemma}\label{sha-automorphy-unramified-near-determined}
In the situation of Lemma~\ref{sha-automorphy-unramified-near-scale}, let $\xi:A_n^\circ\to\bar{\mathbf Q}_\ell$ be an $\mathcal O_1$-algebra map, and let $D:A_n^\circ\to\bar{\mathbf Q}_\ell$ be an $\mathcal O_1$-linear map with $D(ab)=\xi(a)D(b)+\xi(b)D(a)$ that vanishes at every entry of every $\rho^{\mathrm u}(g)$ and at every value of every $\psi^{\mathrm u}_{v,i}$. Then $D=0$.
\end{lemma}
\leanref{Development/AutomorphicArithmetic/AutomorphicVarieties/ModularityArgument/LinkedUnramifiedChains/ScalarIrreducibility/DerExt.lean}{UpperTheory.V.J.derExt}
\begin{proof}
\begin{enumerate}
\item \emph{A dense subalgebra.} $D(1)=D(1\cdot1)=2D(1)$, so $D(1)=0$, and $D$ vanishes on $\mathcal O_1$. By the Leibniz rule, $D$ vanishes on the $\mathcal O_1$-subalgebra generated by the entries of $\rho^{\mathrm u}$ and the values of the $\psi^{\mathrm u}_{v,i}$, and so on the subalgebra $D_0$ generated by these values and the products of entries of total weight $0$, whose closure is $B$ (Definition~\ref{sha-automorphy-realization}(2)).
\item \emph{$D$ vanishes on $B$.} The restriction of $D$ to $B$ satisfies the Leibniz rule at the restriction of $\xi$, and $B$ is a complete local noetherian $\mathcal O_1$-algebra with the residue field of $\mathcal O_1$ (Lemma~\ref{sha-automorphy-realization-presentation}(3)). Let $b\in B$ and $\varepsilon>0$. Lemma~\ref{sha-automorphy-unramified-near-tangent-small} gives $k$ with $|D(b')|\le\varepsilon$ for $b'\in\mathfrak m_B^k$. The ideals $B\cap\mathfrak m_{A_n}^N$ of $B$ decrease, and their intersection is $0$, since $\bigcap_N\mathfrak m_{A_n}^N=0$ by Krull's intersection theorem \cite[Theorem~8.10]{Mat86}; so Lemma~\ref{sha-automorphy-unramified-near-determined-cofinal} gives $M$ with $B\cap\mathfrak m_{A_n}^M\subseteq\mathfrak m_B^k$. As $b$ lies in the closure of $D_0$, there is $b_0\in D_0$ with $b-b_0\in\mathfrak m_{A_n}^M$, hence in $\mathfrak m_B^k$. So $|D(b)|=|D(b-b_0)|\le\varepsilon$ by (1), and $D(b)=0$.
\item \emph{Conclusion.} $A_n^\circ$ is generated as an $\mathcal O_1$-algebra by $B$ and the entries of $\rho^{\mathrm u}$, so $D=0$ by (2) and the Leibniz rule.
\end{enumerate}
\end{proof}
% derExt (DerExt.lean): kill for (1) and (3), close_zero with derCont (Lemma sha-automorphy-unramified-near-tangent-small) for (2); close_zero finds M by cofinal (Lemma sha-automorphy-unramified-near-determined-cofinal) applied to I_N = B meet m_R^N, with Ideal.iInf_pow_eq_bot_of_isLocalRing for the intersection.

\begin{lemma}\label{sha-automorphy-unramified-near-torus}
In the situation of Lemma~\ref{sha-automorphy-unramified-near-scale}, let $\xi:A_n^\circ\to\bar{\mathbf Q}_\ell$ be an $\mathcal O_1$-algebra map with off-triangular kernel, with $|\xi(a)|\le1$ for $a\in A_n^\circ$ and $|\xi(b)|<1$ for $b$ in the maximal ideal of $B$, and let $\rho_\xi(g)$ be the matrix $\xi(\rho^{\mathrm u}(g))$. Let $T_0$ be the $\bar{\mathbf Q}_\ell$-vector space of $\mathcal O_1$-linear maps $D:A_n^\circ\to\bar{\mathbf Q}_\ell$ with $D(ab)=\xi(a)D(b)+\xi(b)D(a)$ that vanish at $\psi^{\mathrm u}_{v,i}(h)$ for all $v$, $i$ and $h\in I_v$. Suppose that for every $D\in T_0$ there is $X\in M_n(\bar{\mathbf Q}_\ell)$ with $D(\rho^{\mathrm u}(g)_{ij})=(X\rho_\xi(g)-\rho_\xi(g)X)_{ij}$ for all $g\in\Gamma_0$ and all $i,j$. Then $\dim T_0\le1$.
\end{lemma}
\leanref{Development/AutomorphicArithmetic/AutomorphicVarieties/ModularityArgument/LinkedUnramifiedChains/ScalarIrreducibility/FixLe.lean}{UpperTheory.V.J.fixLe}
\leanref{Development/AutomorphicArithmetic/AutomorphicVarieties/ModularityArgument/LinkedUnramifiedChains/ScalarIrreducibility/FreeEtaleSheafZmodExistsHomSendingGenerator.lean}{UpperTheory.V.J.torus}
\begin{proof}
Write $w_i=w_{n,i}$, $\bar\imath=n+1-i$ and $g^*=c_\Gamma g^{-1}c_\Gamma^{-1}$; by the residual conditions (Definition~\ref{sha-descent-residual}(4)), $(g^*)^*=g$.
\begin{enumerate}
\item \emph{Residues.} If $b\in B$ has nonzero residue, then $\xi(b)\ne0$: $b=c+m$ with $c$ a unit of $\mathcal O_1$ and $m\in\mathfrak m_B$, and $|\xi(c)|=1>|\xi(m)|$.
\item \emph{Irreducibility.} As the kernel of $\xi$ is off-triangular, $0$ and $\bar{\mathbf Q}_\ell^n$ are the only subspaces stable under every $\rho_\xi(g)$ (Lemma~\ref{sha-automorphy-unramified-realization-irreducible}).
\item \emph{$X$ is diagonal.} Let $D\in T_0$, with its matrix $X$. By condition (b) of Definition~\ref{sha-descent-link}(2), $Z=\rho_\xi(g_0)-\rho_\xi(g_0^*)$ is diagonal. Its diagonal entries are the images of elements of $B$ whose residues are $a^{w_i}-a^{-w_i}$, with $a=\bar\Psi_\Gamma(g_0)$; for $i\ne j$ these residues differ, since $a^2\ne1$ (residual condition (4)) and $2\ne0$ in the residue field. So the diagonal entries of $Z$ are pairwise distinct, by (1). The off-diagonal entries of $\rho^{\mathrm u}(g_0)-\rho^{\mathrm u}(g_0^*)$ vanish, so for $i\ne j$, $0=(XZ-ZX)_{ij}=X_{ij}(Z_{jj}-Z_{ii})$, and $X_{ij}=0$.
\item \emph{The entries.} By (3), $D(\rho^{\mathrm u}(g)_{ij})=(X_{ii}-X_{jj})\xi(\rho^{\mathrm u}(g)_{ij})$. By condition (a) of Definition~\ref{sha-descent-link}(2), $\rho^{\mathrm u}(g)_{ij}=\rho^{\mathrm u}(g^*)_{\bar\jmath\bar\imath}$; applying $D$ to both sides gives $(y_i-y_j)\xi(\rho^{\mathrm u}(g)_{ij})=0$ with $y_i=X_{ii}+X_{\bar\imath\bar\imath}$. So the diagonal matrix with entries $y_i$ commutes with every $\rho_\xi(g)$, and the kernel of its difference with $y_1$ is a nonzero subspace stable under every $\rho_\xi(g)$; by (2), all $y_i$ are equal. For $n=2$ let $t=(X_{11}-X_{22})/2$; for $n=3$, $X_{11}+X_{33}=2X_{22}$, and let $t=X_{11}-X_{22}$. Then $X_{ii}-X_{jj}=t(w_i-w_j)$, and $D(\rho^{\mathrm u}(g)_{ij})=t(w_i-w_j)\xi(\rho^{\mathrm u}(g)_{ij})$.
\item \emph{The characters.} Fix $v$, and for $h\in G_v$ let $R(h)=\rho_\xi(\mathrm{loc}_vh)$, $p_i(h)=\xi(\psi^{\mathrm u}_{v,i}(h))$, $d_i(h)=D(\psi^{\mathrm u}_{v,i}(h))$ and $M_i(h)=R(h)-p_i(h)$. By condition (c) of Definition~\ref{sha-descent-link}(2), $\prod_i\big(\rho^{\mathrm u}(\mathrm{loc}_vh_i)-\psi^{\mathrm u}_{v,i}(h_i)\big)=0$ for $h_1,\dots,h_n\in G_v$. Applying $\xi$ gives $M_1(h_1)\cdots M_n(h_n)=0$. Applying $D$ to the entries, with the Leibniz rule and $D(\rho^{\mathrm u}(\mathrm{loc}_vh))=XR(h)-R(h)X$, gives $\sum_iM_1(h_1)\cdots\big(XM_i(h_i)-M_i(h_i)X-d_i(h_i)\big)\cdots M_n(h_n)=0$; the terms with $X$ telescope to $XM_1\cdots M_n-M_1\cdots M_nX=0$, so $\sum_id_i(h_i)W_i=0$, where $W_i$ is the product of the $M_k(h_k)$, $k\ne i$, in order.
By residual condition (7) there is $s\in G_v$ with $\bar\Psi_\Gamma(\mathrm{loc}_vs)^2\ne1$. The residue of $\psi^{\mathrm u}_{v,i}(s)$ is $\bar\Psi_\Gamma(\mathrm{loc}_vs)^{w_i}$, and that of the diagonal entry $\rho^{\mathrm u}(\mathrm{loc}_vs)_{ii}$ is the same; these are pairwise distinct for distinct $i$, as $w_i-w_j$ divides $2$. By (1), the $a_i=p_i(s)$ are pairwise distinct. With all $h_k=s$, let $R=R(s)$ and $W_i^s=\prod_{k\ne i}(R-a_k)$. Its entry $(i,i)$ is the image of an element of $B$ with residue $\prod_{k\ne i}(\bar\Psi_\Gamma(\mathrm{loc}_vs)^{w_i}-\bar\Psi_\Gamma(\mathrm{loc}_vs)^{w_k})\ne0$, since the diagonal entries of $\rho^{\mathrm u}(\mathrm{loc}_vs)$ and of $\rho^{\mathrm u}(\mathrm{loc}_vs^2)$ lie in $B$; so $W_i^s\ne0$ by (1). As $\prod_k(R-a_k)=0$, $W_k^sW_i^s=0$ for $k\ne i$ and $W_i^sW_i^s=\prod_{k\ne i}(a_i-a_k)\,W_i^s$. Multiplying $\sum_kd_k(s)W_k^s=0$ by $W_i^s$ gives $d_i(s)\prod_{k\ne i}(a_i-a_k)W_i^s=0$, so $d_i(s)=0$ for every $i$. Now let $h\in G_v$, fix $i$, and take $h_i=h$ and $h_k=s$ for $k\ne i$: the relation becomes $d_i(h)W_i^s=0$, so $d_i(h)=0$. So $D$ vanishes at every value of every $\psi^{\mathrm u}_{v,i}$.
\item \emph{The bound.} For $D\in T_0$ let $t_D$ be the $t$ of (4). If $T_0\ne0$, choose $D_1\ne0$ in $T_0$. Then $t_{D_1}\ne0$: otherwise $D_1$ vanishes at the entries by (4) and at the values of the characters by (5), and $D_1=0$ by Lemma~\ref{sha-automorphy-unramified-near-determined}. For $D\in T_0$, $t_{D_1}D-t_DD_1$ vanishes at the entries and at the values of the characters, so it is $0$ by the same lemma. So $T_0$ is spanned by $D_1$.
\end{enumerate}
\end{proof}
% fixLe (FixLe.lean) for (6), with irrTI (Lemma sha-automorphy-unramified-realization-irreducible) for (2) and derExt (Lemma sha-automorphy-unramified-near-determined). torus (FreeEtaleSheafZmodExistsHomSendingGenerator.lean): ne_of_res for (1); offDiag_zero and z_ne for (3); comm_entry, hpol, scalar_of_irr and free_etale_sheaf_zmod_exists_hom_sending_generator for (4); the field f3 of Frame.Ok, der_shift, ordDer2 and ordDer3, kill2 and kill3, finish2 and finish3, with ne_lin and ne_quad for the nonvanishing of W_i^s, for (5), case by case for n = 2 and n = 3, which the text does for both at once. torus does not use that D vanishes on inertia.

\begin{lemma}\label{sha-automorphy-unramified-near-weight}
In the situation of Lemma~\ref{sha-automorphy-unramified-realization}, let $B_3$ be the universal ordinary determinant ring of the realization, and for each place $v$ of $L^+$ above $\ell$ fix an embedding $\tau_v$ of $L_{\tilde v}$ in $\bar{\mathbf Q}_\ell$. Let $h_{v,j,\tau}=2(j-1)$ if $\tau=\tau_v$ and $h_{v,j,\tau}=j-1$ otherwise. Then $h$ is an increasing weight (Definition~\ref{sha-hecke-arcs}), and for every point $x$ of $B_3$ of distinct type of weight $h$ and every $v$ there is $g\in D_v$ with $\varepsilon(g)=1$ such that $\psi_{x,v,i}(g)^m\ne\psi_{x,v,j}(g)^m$ for all $i\ne j$ and $m\ge1$.
\end{lemma}
\leanref{Development/AutomorphicArithmetic/AutomorphicVarieties/ModularityArgument/LinkedUnramifiedChains/ScalarIrreducibility/PlaceFree.lean}{UpperTheory.V.J.wtEx}
\leanref{Development/AutomorphicArithmetic/AutomorphicVarieties/ModularityArgument/LinkedUnramifiedChains/ScalarIrreducibility/WtUnit.lean}{UpperTheory.V.J.Wu.exists_unit}
\begin{proof}
The weight is increasing, as $j\mapsto2(j-1)$ and $j\mapsto j-1$ are. Fix $v$, let $F=L_{\tilde v}$, let $T$ be the set of embeddings of $F$ in $\bar{\mathbf Q}_\ell$, and $e=|T|$.
\begin{enumerate}
\item \emph{A unit.} The link is large, so the local degree $d_v=[F:\mathbf Q_\ell]$ exceeds $5$; so there are $b\in\mathcal O_L$ and $\tau_1,\tau_2\in T$ with $\tau_1(b)\ne\tau_2(b)$. Let $z=1+\ell^2b$, $c=\prod_{\tau\in T}\tau(z)\in\mathbf Q_\ell$ and $y=z^e/c$. Each $|\tau(z)-1|\le\ell^{-2}$, so $|\tau(z)|=1$, $|c|=1$ and $y$ is a unit of $\mathcal O_F$, with $\prod_\tau\tau(y)=c^e/c^e=1$. No $\tau(y)$ is a root of unity: if $\tau(y)^m=1$ with $m\ge1$, then $y^m=1$, so $\tau'(z)^{em}=c^m$ for every $\tau'\in T$, and $u=\tau_1(z)/\tau_2(z)$ satisfies $u^{em}=1$ and $|u-1|\le\ell^{-2}$; the only such root of unity is $1$, so $\tau_1(z)=\tau_2(z)$ and $\tau_1(b)=\tau_2(b)$, which is false.
\item \emph{An element of inertia.} By local class field theory there is $\sigma\in I_v$ with $\operatorname{Art}^{-1}\sigma=y$, and $\varepsilon(\sigma)=\prod_\tau\tau(y)^{-1}=1$.
\item \emph{The ratios.} Let $x$ be of distinct type of weight $h$, with characters $\phi_{v,i}$ of $I_v$ with open kernels as in Definition~\ref{sha-hecke-arcs}(1). Then $\psi_{x,v,i}(\sigma)\phi_{v,i}(\sigma)=\prod_\tau\tau(y)^{-h_{v,i,\tau}}=\tau_v(y)^{-(i-1)}\prod_\tau\tau(y)^{-(i-1)}=\tau_v(y)^{-(i-1)}$. Each $\phi_{v,i}(\sigma)$ is a root of unity, of order $N_i$ say, as $\phi_{v,i}$ has open kernel in the compact group $I_v$. If $\psi_{x,v,i}(\sigma)^m=\psi_{x,v,j}(\sigma)^m$ with $i\ne j$ and $m\ge1$, raising to the power $N_iN_j$ gives $\tau_v(y)^{(j-i)mN_iN_j}=1$, against (1). Take $g=\sigma$.
\end{enumerate}
\end{proof}
% wtEx (PlaceFree.lean): the weight 2j at tau_0 = Classical.choice and j elsewhere (Fin 3, j = 0, 1, 2), increasing; then place_free: R.big.2 (d_v > 5) for an element of Dec not in Gal(Qbar/L^+) and so b with tau_1(b) != tau_2(b); Wu.exists_unit (WtUnit.lean: z = 1 + l^2 b, y = z^d / N(z), eq_one_of_pow_eq_one) for (1); art_elt (Loc.laArtin, ineIncl_onto_units, Loc.laCyc) for (2); prod_wt and free_of for (3); mem_G' for epsilon(g) = 1. The Lean asks GlobalArtinReciprocity. It uses only that x is arithmetic of weight h.

\begin{definition}\label{sha-automorphy-unramified-near-locally-generic}
In Situation~\ref{sha-hecke-situation}, let $\iota:\bar{\mathbf Q}_\ell\cong\mathbf C$.
\begin{enumerate}
\item A point $x$ of $B$ and a continuous homomorphism $\rho:\Gamma_0\to GL_n(\bar{\mathbf Q}_\ell)$ form an \emph{automorphic locally generic point} if: $x(D^{\mathrm u})$ is the determinant of $\rho$, and $x$ is arithmetic of some weight; $\rho$, as a representation of $\operatorname{Gal}(\bar L/L)$, is attached through $\iota$ to a regular algebraic cuspidal automorphic representation $\pi$ of $GL_n(\mathbf A_L)$ with $\pi^\vee\cong\pi^c$, and is windowed (Theorem~\ref{sha-residually-reducible-automorphy}); $\rho$ is trivial on $I_{v_0}$; for each $v\mid\ell$, $\rho$ is conjugate on $D_v$ to an upper triangular representation with diagonal characters $\psi_{x,v,1},\dots,\psi_{x,v,n}$; for each $v\mid\ell$ there is $g\in D_v$ with $\varepsilon(g)=1$ and $\psi_{x,v,i}(g)^m\ne\psi_{x,v,j}(g)^m$ for all $i\ne j$ and $m\ge1$; $\rho$ is irreducible on $\Gamma'$; and every open subset of $\Gamma_0$ that meets $\Gamma'$ contains some $\sigma\in\Gamma'$ at which the characteristic polynomial of $\rho(\sigma)$ is separable.
\item Let $\rho:\Gamma_0\to GL_n(\bar{\mathbf Q}_\ell)$ be a homomorphism, $\psi_{v,i}:D_v\to\bar{\mathbf Q}_\ell^\times$ characters, and $P\in GL_n(\bar{\mathbf Q}_\ell)$ with ${}^t\rho(g)P\rho(cgc^{-1})=\varepsilon(g)^{1-n}P$ for $g\in\Gamma_0$. An \emph{adjoint extension} for $P$ is a representation $R$ of $\Gamma$ on $M_n(\bar{\mathbf Q}_\ell)$ with $R(g)X=\rho(g)X\rho(g)^{-1}$ for $g\in\Gamma_0$ and $R(c)X=-P^{-1}\,{}^tX\,P$. A continuous map $z:\Gamma\to M_n(\bar{\mathbf Q}_\ell)$ with $z(gg')=z(g)+R(g)z(g')$ is an \emph{ordinary cocycle} if $z$ vanishes on $I_u\cap\bigcap_i\ker\bar\chi_i$ for every $u\in S_{\mathrm a}$, and if for every $v\mid\ell$ there are $a\in GL_n(\bar{\mathbf Q}_\ell)$ and $X\in M_n(\bar{\mathbf Q}_\ell)$ such that for every $g\in D_v$: $a\rho(g)a^{-1}$ is upper triangular with diagonal entries $\psi_{v,1}(g),\dots,\psi_{v,n}(g)$; $a\big(z(g)-\rho(g)X\rho(g)^{-1}+X\big)a^{-1}$ is upper triangular; and its diagonal entries vanish if $g\in I_v$.
\end{enumerate}
\end{definition}
\leanref{Development/AutomorphicArithmetic/AutomorphicVarieties/ModularityArgument/ThresholdArguments/DenseThresholdCases/AutomorphicPointGalois.lean}{UpperTheory.V.Th.AutomorphicPointLocallyGeneric}
\leanref{Development/AutomorphicArithmetic/AutomorphicVarieties/ModularityArgument/ThresholdArguments/DenseThresholdCases/AutomorphicPointGalois.lean}{UpperTheory.V.Th.AutomorphicPointGalois}
\leanref{Development/AutomorphicArithmetic/AutomorphicVarieties/TheoremsMuLinked/OrdZero.lean}{UpperTheory.V.Th.OrdCoc}
\leanref{Development/AutomorphicArithmetic/AutomorphicVarieties/TheoremsMuLinked/OrdZero.lean}{UpperTheory.V.Th.OrdAt}
\leanref{Development/AutomorphicArithmetic/AutomorphicVarieties/TheoremsMuLinked/OrdZero.lean}{UpperTheory.V.Th.OrdZero}
\leanref{Development/AutomorphicArithmetic/AutomorphicVarieties/WeakShimuraTheory/EigenbasisLifts/IsTanZero.lean}{UpperTheory.V.Ws.Eb.IsAdExt}

\begin{theorem}\label{sha-automorphy-unramified-near-vanishing}
Assume that Galois representations exist, in the sense of Definition~\ref{gal-galois-representations}, relative to the maps satisfying the laws of Definition~\ref{gal-llc-laws}. In Situation~\ref{sha-hecke-situation}, let $2\le n\le4$ and $\iota:\bar{\mathbf Q}_\ell\cong\mathbf C$, and suppose that the residual characters $\bar\chi_i$ are trivial on $I_{v_0}$ and agree on $D_{v_0}$ and that $\varepsilon\bmod\ell$ is nontrivial on $D_{v_0}$. Let a point $x$ of $B$ and $\rho:\Gamma_0\to GL_n(\bar{\mathbf Q}_\ell)$ form an automorphic locally generic point (Definition~\ref{sha-automorphy-unramified-near-locally-generic}), and let $\psi_{v,i}=\psi_{x,v,i}$. Then for every $P$ as in Definition~\ref{sha-automorphy-unramified-near-locally-generic}(2) and every adjoint extension $R$ for $P$, every ordinary cocycle $z$ is a coboundary: there is $X$ with $z(g)=R(g)X-X$ for all $g\in\Gamma$.
\end{theorem}
\leanref{Development/AutomorphicArithmetic/AutomorphicVarieties/ModularityArgument/ThresholdArguments/DenseThresholdCases/AutomorphicPointGalois.lean}{UpperTheory.V.Th.DenseThresholdMainCase}
\leanref{Development/AutomorphicArithmetic/AutomorphicVarieties/ModularityArgument/ThresholdArguments/DenseThresholdCases/ThFinal.lean}{UpperTheory.V.Th.core_alt_bis}
\begin{proof}
\openstep{The proof is not yet written out. It passes through the Hida family at the point $x$, given by the Hecke ideal of Proposition~\ref{sha-hecke-ideal}, and through the published inputs on local and global Galois cohomology.}
% DenseThresholdMainCase is proved by core_alt_bis (ThFinal.lean) from Af.HypothesesH' and Pa.CohP: core_of_ord_vanishing_two_phi_all (TCoreOf.lean, 74 lines: Det.redB, Ws.Eb.relKill, Det.polNorm, regDet, and OrdVan2Alt applied to z minus a coboundary) with W.phiAll, and ordinary_vanishing_two_of_printed_hida (TOrdVan.lean, 193 lines: nmAbs, Th.Nm.regular_arithmetic_hida_of_literal_automorphy, b0Ex, lamEx, hida_x_of_printed_hida (THidaX.lean) and core_o_of_literal_automorphy_n (TCoreO.lean)). These draw on TheoremsMuLinked (about 18,700 lines) and on the Hida inputs of Proposition sha-hecke-ideal.
\end{proof}

\begin{lemma}\label{sha-automorphy-unramified-near-small-continuity}
In the situation of Lemma~\ref{sha-automorphy-unramified-near-scale}, let $\xi:A_n^\circ\to\bar{\mathbf Q}_\ell$ be an $\mathcal O_1$-algebra map with $|\xi(a)|\le1$ for $a\in A_n^\circ$ and $|\xi(b)|<1$ for $b$ in the maximal ideal $\mathfrak m_B$ of $B$, and let $D:A_n^\circ\to\bar{\mathbf Q}_\ell$ be an $\mathcal O_1$-linear map with $D(ab)=\xi(a)D(b)+\xi(b)D(a)$ for all $a,b$. Then for all $i,j$ the maps $g\mapsto\xi(\rho^{\mathrm u}(g)_{ij})$ and $g\mapsto D(\rho^{\mathrm u}(g)_{ij})$ from $\Gamma_0$ to $\bar{\mathbf Q}_\ell$ are continuous.
\end{lemma}
\leanref{Development/AutomorphicArithmetic/AutomorphicVarieties/ModularityArgument/LinkedUnramifiedChains/ScalarIrreducibility/DerRep.lean}{UpperTheory.V.J.derRep}
\leanref{Development/AutomorphicArithmetic/AutomorphicVarieties/FrobeniusSmoothVectors/ArtinRees.lean}{UpperTheory.V.Fr.cone_u}
\begin{proof}
Fix $i,j$, let $w=w_{n,i}-w_{n,j}$, and let $E$ be the $B$-submodule of $A_n$ spanned by the entries $\rho^{\mathrm u}(g)_{ij}$, $g\in\Gamma_0$.
\begin{enumerate}
\item \emph{Small on powers of $\mathfrak m_B$.} The restrictions of $\xi$ and $D$ to $B$ are an $\mathcal O_1$-algebra map and an $\mathcal O_1$-linear map with the Leibniz rule, and $B$ is a complete local noetherian $\mathcal O_1$-algebra with the residue field $k_1$ of $\mathcal O_1$ (Lemma~\ref{sha-automorphy-realization-presentation}(3)). So by Lemma~\ref{sha-automorphy-unramified-near-tangent-small}, for every $\delta>0$ there is $k$ with $|\xi(b)|\le\delta$ and $|D(b)|\le\delta$ for $b\in\mathfrak m_B^k$.
\item \emph{A bound on $B$.} Take $k_0$ for $\delta=1$ in (1). The ring $B/\mathfrak m_B^{k_0}$ is finite, as $B$ is noetherian with finite residue field. Choosing representatives $b_1,\dots,b_q$ of its elements, every $b\in B$ is some $b_t$ plus an element of $\mathfrak m_B^{k_0}$; as the absolute value is ultrametric, $|D(b)|\le C_0=\max(1,|D(b_1)|,\dots,|D(b_q)|)$.
\item \emph{A bound on $E$.} By Lemma~\ref{sha-automorphy-realization-presentation}(1) and (4), $E$ lies in $(A_n)_w$, a finitely generated module over the noetherian ring $B$; so $E$ is spanned by finitely many of the entries, $e_1,\dots,e_r$. For $b\in B$ and $e\in E$, $D(be)=\xi(b)D(e)+\xi(e)D(b)$, so $|D(be)|\le\max(|D(e)|,C_0)$. Hence $|D(e)|\le C=\max(C_0,|D(e_1)|,\dots,|D(e_r)|)$ for every $e\in E$.
\item \emph{Small on $\mathfrak m_B^kE$.} Let $\varepsilon>0$, $\delta=\varepsilon/C$, and $k$ as in (1) for $\delta$. For $b\in\mathfrak m_B^k$ and $e\in E$, $|\xi(be)|\le\delta\le\varepsilon$, and $|D(be)|\le\max(\delta C,\delta)=\varepsilon$ by (2) and (3), since $|\xi(e)|\le1$. Sums of such elements satisfy the same bounds, so $|\xi|\le\varepsilon$ and $|D|\le\varepsilon$ on $\mathfrak m_B^kE$.
\item \emph{Local constancy.} Let $N\ge0$. By the Artin-Rees lemma \cite[Theorem~8.5]{Mat86}, applied to the submodule $E$ of the finitely generated $B$-module $(A_n)_w$, there is $K$ with $E\cap\mathfrak m_B^K(A_n)_w\subseteq\mathfrak m_B^NE$. By Lemma~\ref{sha-automorphy-unramified-near-tangent-filtration}, applied to the presentation of Lemma~\ref{sha-automorphy-realization-presentation} with $\mathcal O_1$ and $k_1$ in place of $\mathcal O$ and $\mathbf F$, there is $M$ with $(A_n)_w\cap\mathfrak m_{A_n}^M\subseteq\mathfrak m_B^K(A_n)_w$. By condition (d) of Definition~\ref{sha-descent-link}(2), for $g\in\Gamma_0$ the $g'$ with $\rho^{\mathrm u}(g')_{ij}-\rho^{\mathrm u}(g)_{ij}\in\mathfrak m_{A_n}^M$ form a neighbourhood of $g$; for them this difference lies in $E\cap(A_n)_w\cap\mathfrak m_{A_n}^M\subseteq\mathfrak m_B^NE$.
\item \emph{Conclusion.} Let $g\in\Gamma_0$ and $\varepsilon>0$. With $k$ as in (4) and $N=k$ in (5), on a neighbourhood of $g$ the values $\xi(\rho^{\mathrm u}(g')_{ij})$ and $D(\rho^{\mathrm u}(g')_{ij})$ differ from those at $g$ by at most $\varepsilon$.
\end{enumerate}
\end{proof}
% derRep (DerRep.lean), its first two parts: der_small (derCont, Lemma sha-automorphy-unramified-near-tangent-small) for (1), der_bdd (finite_quot, bdd_of_small) for (2), small_smul for (3) and (4), entry_cont' and cont_of_small for (6). The local constancy (5) is the field cont of cone_u (ArtinRees.lean): artin_rees for the span of the entries inside S.gr w, hAR from cont_u (part_filt, Lemma sha-automorphy-unramified-near-tangent-filtration) and S.lift.1.1; the finiteness of E is the field fg of cone_u (gr_A, isNoetherian_of_le).

\begin{lemma}\label{sha-automorphy-unramified-near-small}
Assume that Galois representations exist, in the sense of Definition~\ref{gal-galois-representations}, relative to the maps satisfying the laws of Definition~\ref{gal-llc-laws}. In the situation of Lemma~\ref{sha-automorphy-unramified-realization}, let $G_0$ be the Galois group of the residual datum of the realization, $\pi:G_0\to\Gamma_0$ its map and $\kappa_0$ its character, and $B_3$ its universal ordinary determinant ring. Let $\xi':A_3^\circ\to\bar{\mathbf Q}_\ell$ be an $\mathcal O_1$-algebra map with off-triangular kernel, with values of absolute value at most $1$, with $|\xi'(b)|<1$ for $b$ in the maximal ideal of $B$ and $|\xi'(e)|<1$ for the entries $e$ off the diagonal of the universal linked point. Let $\rho_{\xi'}(g)=\xi'(\rho^{\mathrm u}(g))$ for $g\in\Gamma_0$, and $\rho(g)=\kappa_0(g)\rho_{\xi'}(\pi(g))$ for $g\in G_0$. Let $h$ be an increasing weight, and suppose that $x=\xi'\circ\phi$ is of distinct type of weight $h$ and that for every $v$ some $g\in D_v$ with $\varepsilon(g)=1$ has $\psi_{x,v,i}(g)^m\ne\psi_{x,v,j}(g)^m$ for all $i\ne j$ and $m\ge1$. Then $x$ and $\rho$ form an automorphic locally generic point (Definition~\ref{sha-automorphy-unramified-near-locally-generic}), and $\psi_{x,v,i}(g)=\kappa_0(g)\xi'(\psi^{\mathrm u}_{v,i}(g))$ for $g\in D_v$.
\end{lemma}
\leanref{Development/AutomorphicArithmetic/AutomorphicVarieties/ModularityArgument/ThresholdArguments/DenseThresholdCases/EndPtOf.lean}{UpperTheory.Gal.projectedFormsSystemPart.niceAt2_alt}
\leanref{Development/AutomorphicArithmetic/AutomorphicVarieties/ModularityArgument/ThresholdArguments/DenseThresholdCases/TNiceSmall.lean}{UpperTheory.V.Th.nice_zero_of_rank_two_or_three_irreducible_free}
\begin{proof}
\begin{enumerate}
\item \emph{Determinant and characters.} Since $\phi$ is compatible (Definition~\ref{sha-automorphy-realization}(3)), $x$ carries $D^{\mathrm u}$ to the determinant of $\rho$ and $\psi^{\mathrm u}_{v,i}$ to $g\mapsto\kappa_0(g)\xi'(\psi^{\mathrm u}_{v,i}(g))$.
\item \emph{Arithmetic.} A point of distinct type of weight $h$ is arithmetic of weight $h$.
\item \emph{Automorphy.} As $x=\xi'\circ\phi$ is arithmetic and $\xi'$ satisfies the bounds of Lemma~\ref{sha-automorphy-unramified-realization}, that lemma shows that $\rho$, as a representation of $\operatorname{Gal}(\bar{\mathbf Q}/L)$, which is $r_{\xi'}\otimes\chi_{\mathrm{cyc}}^{-1}$, is attached through $\iota$ to a regular algebraic cuspidal $\pi$ with $\pi^\vee\cong\pi^c$, and is windowed.
\item \emph{Irreducibility.} By Lemma~\ref{sha-automorphy-unramified-realization-irreducible}, $\rho_{\xi'}$ is irreducible, and so is $\rho$, as $\pi$ is surjective.
\item \emph{Local genericity.} This is the hypothesis on $x$, by (1).
\item \emph{Continuity.} The character $\kappa_0$ is continuous for the adic topology of $\mathcal O_1$ (Definition~\ref{sha-automorphy-realization}(2)), hence as a map to $\bar{\mathbf Q}_\ell$, since the elements of $\varpi^N\mathcal O_1$, for a uniformizer $\varpi$, have absolute value at most $|\varpi|^N$. The map $\pi$ is continuous, and so are the entries $g\mapsto\xi'(\rho^{\mathrm u}(g)_{ij})$ of $\rho_{\xi'}$, by Lemma~\ref{sha-automorphy-unramified-near-small-continuity} with $D=0$. So $\rho$ is continuous.
\item \emph{Inertia at $v_0$.} Let $\mathrm{pr}$ be the map from $\operatorname{Gal}(\bar{\mathbf Q}/L)$ onto $G_0$, and let $g\in I_{v_0}$, the image under $\mathrm{pr}$ of an element $s$ of the inertia group of $\operatorname{Gal}(\bar{\mathbf Q}/L)$ at a place $A$ above $\tilde v_0$. As $v_0$ does not divide $M$ and $\ell N$ divides $M$, $A$ is not above $\ell$ and $\Psi$ is unramified at $A$; so $\Psi(s)=1$, and $\eta(s)=1$ as $s\in\operatorname{Gal}(\bar{\mathbf Q}/L)$. By Definition~\ref{sha-descent-link}(1)(i), $s$ lies in the kernel of $p$, so $\pi(g)=p(s)=1$. And $\kappa_0(g)=1$, as $\kappa_0$ is trivial on $I_{v_0}$ (Definition~\ref{sha-automorphy-realization}(2)). So $\rho(g)=1$.
\item \emph{Triangular forms.} Fix $v\mid\ell$, and identify $D_v$ with the group $G_v$ of the link, so that $\pi$ restricts to $\mathrm{loc}_v$ (Definition~\ref{sha-automorphy-realization}(2)). Let $C\subset\bar{\mathbf Q}_\ell$ be the local ring of the elements of absolute value at most $1$; the values of $\xi'$ lie in $C$, and an element of $B$ with nonzero residue is a unit of $\mathcal O_1$ plus an element of $\mathfrak m_B$, so its image under $\xi'$ is a unit of $C$. By residual condition (7) of Definition~\ref{sha-descent-residual} there is $h_0\in G_v$ such that $a=\bar\Psi_\Gamma(\mathrm{loc}_vh_0)$ has $a^2\ne1$.
For $i\ne j$, the element $\psi^{\mathrm u}_{v,i}(h_0)-\psi^{\mathrm u}_{v,j}(h_0)$ of $B$ has residue $a^{w_{3,i}}-a^{w_{3,j}}\ne0$, since $w_{3,i}-w_{3,j}$ divides $2$; so its image under $\xi'$ is a unit of $C$. The diagonal entries of $\rho^{\mathrm u}(\mathrm{loc}_vh_0)$ have weight $0$, so each $\rho^{\mathrm u}(\mathrm{loc}_vh_0)_{ii}-\psi^{\mathrm u}_{v,i}(h_0)$ lies in $B$, and in $\mathfrak m_B$ by condition (e) of Definition~\ref{sha-descent-link}(2); its image under $\xi'$ has absolute value less than $1$, as have the images of the entries off the diagonal. So modulo the maximal ideal of $C$, $\rho_{\xi'}(\mathrm{loc}_vh_0)$ is the diagonal matrix with the residues of the $\xi'(\psi^{\mathrm u}_{v,i}(h_0))$, and its characteristic polynomial reduces to $\prod_i(X-\xi'(\psi^{\mathrm u}_{v,i}(h_0)))$.
Condition (c) of Definition~\ref{sha-descent-link}(2), applied through $\xi'$, gives the vanishing of the ordered products for $\rho_{\xi'}\circ\mathrm{loc}_v$ and the characters $\xi'\circ\psi^{\mathrm u}_{v,i}$, which take values in $C^\times$. So Lemma~\ref{sha-rigidity-lagrange}, with $C$, $H=G_v$, $r=\rho_{\xi'}\circ\mathrm{loc}_v$, these characters and $g_0=h_0$, gives $g\in GL_3(C)$ such that $g\rho_{\xi'}(\mathrm{loc}_vh)g^{-1}$ is upper triangular with diagonal entries $\xi'(\psi^{\mathrm u}_{v,i}(h))$ for every $h\in G_v$. Multiplying by $\kappa_0(h)$, $g\rho(h)g^{-1}$ is upper triangular with diagonal entries $\psi_{x,v,1}(h),\psi_{x,v,2}(h),\psi_{x,v,3}(h)$, by (1).
\item \emph{Irreducibility on $\Gamma'$ and separability.} Let $G'_0=G_0\cap\ker\varepsilon$, the group $\Gamma'$ of Situation~\ref{sha-hecke-situation} for the residual datum of the realization; it is normal in the profinite group $G_0$. Let $A_x$ be the quotient of $B_3$ by the kernel of $x$, a local ring with the residue field $k_1$; let $j_x:A_x\to\bar{\mathbf Q}_\ell$ be the injective map induced by $x$, and $\phi_x:A_x\to k_1$ the residue map. Let $\bar\chi_1,\bar\chi_2,\bar\chi_3$ be the characters of the residual datum, which are the characters $(\bar\Psi_\Gamma\circ\pi)^{w_{3,i}}$ twisted by the reduction of $\kappa_0$, and $\bar\psi_{v,i}$, the restriction of $\bar\chi_i$ to $D_v$, its characters at $v\mid\ell$ (Definition~\ref{sha-automorphy-realization}(2)).
Since $x$ carries $D^{\mathrm u}$ to the determinant of $\rho$, and $D^{\mathrm u}$ reduces to the determinant of $\bar\chi_1\oplus\bar\chi_2\oplus\bar\chi_3$, the characteristic polynomial of $\rho(y)$, for $y\in G_0$, is the image under $j_x$ of that of $D^{\mathrm u}$ at $y$ taken modulo $\ker x$, whose image under $\phi_x$ is $\prod_i(X-\bar\chi_i(y))$. For $v\mid\ell$ and $d\in D_v$, the class $c$ of $\psi^{\mathrm u}_{v,i}(d)$ has $j_x(c)=\psi_{x,v,i}(d)$ and $\phi_x(c)=\bar\psi_{v,i}(d)$, and by (8) the characteristic polynomial of $\rho(d)$ is $\prod_i(X-\psi_{x,v,i}(d))$. By hypothesis some $d_0\in D_v$ with $\varepsilon(d_0)=1$, so with image in $G'_0$, has $\psi_{x,v,i}(d_0)^m\ne\psi_{x,v,j}(d_0)^m$ for $i\ne j$ and $m\ge1$. At the element $h_0$ of (8), $\bar\psi_{v,1}(h_0)\bar\psi_{v,2}(h_0)^{-1}=a\ne1$.
So Lemma~\ref{sha-rigidity-small}, with the compact group $G_0$, its normal subgroup $G'_0$, $n=3$, $r=\rho$, which is irreducible by (4), $A_x$, $j_x$, $\phi_x$, the $\bar\chi_i$, $D=D_v$ with its map to $G_0$, and the characters $\psi_{x,v,i}$ and $\bar\psi_{v,i}$, shows that $\rho$ is irreducible on $U\cap G'_0$ for every open subgroup $U$ of $G_0$, and in particular on $G'_0$. With the image of $d_0$ in $G'_0$ and the continuity (6), Lemma~\ref{sha-rigidity-separable}, for the profinite group $G_0$, shows that every open subset of $G_0$ that meets $G'_0$ contains an element $\sigma$ of $G'_0$ at which the characteristic polynomial of $\rho(\sigma)$ is separable.
\end{enumerate}
By (1) to (9), $x$ and $\rho$ satisfy every condition of Definition~\ref{sha-automorphy-unramified-near-locally-generic}(1).
% niceAt2_alt (EndPtOf.lean): hm = TwLink.IsPhi.maps for (1); harith for (2); the hypothesis hJ : EndPtAlt, which is the last part of Lemma sha-automorphy-unramified-realization (endPt_of_allGood_of_arithmeticPointClassicality), for (3); R.K.tw_irr with irrTI for (4); hcont for (6), from derRep (Lemma sha-automorphy-unramified-near-small-continuity) with D = 0 and the field kappa_cont of TwLink through X.So.adicCont (adic continuity into the DVR O_1 is continuity); unr0 for (7), from the private unramified_at_base_place of NiceAt.lean (cycZ and the Dirichlet class trivial on that inertia group, k.p_ker) and R.K.ramκ; flag for (8), from the private pt_flag of NiceAt.lean (flag_int, which is Fr.flag_loc over the valuation ring nb l of Qbar_l, with big_of_res, small_of_res and w_ne) and tw_flag; (9) is nice_zero_of_rank_two_or_three_irreducible_free (TNiceSmall.lean), with Th.Det.redB for A_x, j_x and phi_x, Th.Det.irrSmall (Lemma sha-rigidity-small) and Th.Det.n3small (Lemma sha-rigidity-separable), the element of (8) from R.ok.f3 and w_ne, and R.K.psib, R.K.chib for the residual characters.
\end{proof}

\begin{lemma}\label{sha-automorphy-unramified-near-matrix-adjoint}
Let $G$ be a group, $G_0\subset G$ a subgroup of index two, $c\in G\setminus G_0$ with $c^2=1$, and write $g^c=cgc^{-1}$. Let $K$ be a field, $n\ge1$, $\varepsilon:G_0\to K^\times$ a character with $\varepsilon(g^c)=\varepsilon(g)$, and $\rho:G_0\to GL_n(K)$ a homomorphism whose values span $M_n(K)$. Let $P\in GL_n(K)$ satisfy ${}^t\rho(g)P\rho(g^c)=\varepsilon(g)^{1-n}P$ for $g\in G_0$. Then there is a representation $R$ of $G$ on $M_n(K)$ with $R(g)X=\rho(g)X\rho(g)^{-1}$ for $g\in G_0$ and $R(c)X=-P^{-1}\,{}^tXP$.
\end{lemma}
\leanref{Development/AutomorphicArithmetic/AutomorphicVarieties/WeakShimuraTheory/EigenbasisLifts/AeComm.lean}{UpperTheory.V.Ws.Eb.adExt}
\begin{proof}
Let $\theta(X)=-P^{-1}\,{}^tXP$ and $\mathrm{Ad}_g(X)=\rho(g)X\rho(g)^{-1}$, and write $e(g)=\varepsilon(g)^{1-n}$.
\begin{enumerate}
\item \emph{$\theta\,\mathrm{Ad}_g=\mathrm{Ad}_{g^c}\theta$ for $g\in G_0$.} The relation for $P$ gives ${}^t\rho(g)P=e(g)P\rho(g^c)^{-1}$ and $P^{-1}\,{}^t\rho(g)^{-1}=e(g)^{-1}\rho(g^c)P^{-1}$. So $\theta(\rho(g)X\rho(g)^{-1})=-P^{-1}\,{}^t\rho(g)^{-1}\,{}^tX\,{}^t\rho(g)P=\rho(g^c)\theta(X)\rho(g^c)^{-1}$.
\item \emph{$\theta^2=1$.} Let $h\in G_0$ and $A={}^t\rho((h^{-1})^c)$. As $c^2=1$ and $e((h^{-1})^c)=e(h^{-1})$, the relation at $(h^{-1})^c$ gives $AP\rho(h^{-1})=e(h^{-1})P$, and the transpose of the relation at $h^{-1}$ gives $A\,{}^tP\rho(h^{-1})=e(h^{-1})\,{}^tP$. Hence $A=e(h^{-1})P\rho(h)P^{-1}=e(h^{-1})\,{}^tP\rho(h)\,{}^tP^{-1}$, that is, $P\rho(h)P^{-1}={}^tP\rho(h)\,{}^tP^{-1}$. As the $\rho(h)$ span $M_n(K)$, $PXP^{-1}={}^tPX\,{}^tP^{-1}$ for every $X$. Then $\theta(\theta(X))=P^{-1}\,{}^tPX\,{}^tP^{-1}P=P^{-1}(PXP^{-1})P=X$.
\item \emph{The representation.} Let $R(g)=\mathrm{Ad}_g$ for $g\in G_0$ and $R(g)=\mathrm{Ad}_{gc}\theta$ for $g\notin G_0$; then $R(c)=\theta$. For $g,h\in G_0$, $R(gh)=R(g)R(h)$. For $g\in G_0$ and $h\notin G_0$, $R(gh)=\mathrm{Ad}_{ghc}\theta=\mathrm{Ad}_g\mathrm{Ad}_{hc}\theta$. For $g\notin G_0$ and $h\in G_0$, $R(g)R(h)=\mathrm{Ad}_{gc}\theta\mathrm{Ad}_h=\mathrm{Ad}_{gc}\mathrm{Ad}_{h^c}\theta=\mathrm{Ad}_{ghc}\theta=R(gh)$ by (1), as $gc\,h^c=ghc$. For $g,h\notin G_0$, $gh\in G_0$, and $R(g)R(h)=\mathrm{Ad}_{gc}\theta\,\mathrm{Ad}_{hc}\theta=\mathrm{Ad}_{gc}\mathrm{Ad}_{(hc)^c}\theta^2=\mathrm{Ad}_{gh}$ by (1) and (2), as $gc\,(hc)^c=gh$. So $R$ is a representation.
\end{enumerate}
\end{proof}
% adExt (AeComm.lean): hkey (exists_rep_conj_chi_smul_and_q_twisted) for (1); hS (u_mul_ri_mul_v_eq_transpose_of_eigen, by induction on the span) and hthetatheta (neg_conj_neg_transpose_eq_self_of_symm) for (2); the map f with the four cases of map_mul for (3). The Lean's hypothesis Full is the spanning of M_n(K).

\begin{lemma}\label{sha-automorphy-unramified-near-matrix-unipotent}
Let $A$ be a local ring, $E$ an $A$-algebra and $\gamma\in E$. Let $\chi\in A[X]$ be a monic polynomial with $\chi(\gamma)=0$ whose reduction modulo the maximal ideal of $A$ is $(X-1)^n$ for some $n$, and let $m\ge1$ be an integer that is a unit of $A$, with $\gamma^m=1$. Then $\gamma=1$.
\end{lemma}
\leanref{Development/AutomorphicArithmetic/AutomorphicVarieties/CurveDVRSpecialization/SelfTwistEigenvalues/UnipTorsion.lean}{UpperTheory.V.Pr.T.unipTorsion}
\begin{proof}
Let $s=1+X+\dots+X^{m-1}$. Its image in $A[X]/(\chi)$ is a unit. Otherwise it lies in a maximal ideal $\mathfrak M$ of $A[X]/(\chi)$, which is finite over $A$, so $\mathfrak M\cap A$ is the maximal ideal of $A$; in the field $(A[X]/(\chi))/\mathfrak M$ the image $x$ of $X$ is a root of $(X-1)^n$, so $x=1$, and the image of $s$ is $m\ne0$, a contradiction. So there is $p\in A[X]$ with $s(\gamma)p(\gamma)=1$ in $E$. As $(1-\gamma)s(\gamma)=1-\gamma^m=0$, $1-\gamma=(1-\gamma)s(\gamma)p(\gamma)=0$.
\end{proof}
% unipTorsion with the private unipTorsion_unit (UnipTorsion.lean): Ideal.exists_le_maximal, isMaximal_comap_of_isIntegral_of_isMaximal and hroot for the unit; mul_neg_geom_sum for 1 - gamma^m.

\begin{lemma}\label{sha-automorphy-unramified-near-matrix}
Assume that Galois representations exist, in the sense of Definition~\ref{gal-galois-representations}, relative to the maps satisfying the laws of Definition~\ref{gal-llc-laws}. In the situation of Lemma~\ref{sha-automorphy-unramified-near-small}, let $\xi':A_3^\circ\to\bar{\mathbf Q}_\ell$ be an $\mathcal O_1$-algebra map with values of absolute value at most $1$ and $|\xi'(b)|<1$ for $b$ in the maximal ideal of $B$, with $\rho_{\xi'}$ and $\rho$ as there, and let $x$ be a point of $B_3$ such that $x$ and $\rho$ form an automorphic locally generic point and $\psi_{x,v,i}(g)=\kappa_0(g)\xi'(\psi^{\mathrm u}_{v,i}(g))$ for $g\in D_v$. Let $D:A_3^\circ\to\bar{\mathbf Q}_\ell$ be an $\mathcal O_1$-linear map with $D(ab)=\xi'(a)D(b)+\xi'(b)D(a)$ that vanishes at $\psi^{\mathrm u}_{v,i}(h)$ for all $v$, $i$ and $h\in I_v$. Then there is $X\in M_3(\bar{\mathbf Q}_\ell)$ with $D(\rho^{\mathrm u}(g)_{ij})=(X\rho_{\xi'}(g)-\rho_{\xi'}(g)X)_{ij}$ for all $g\in\Gamma_0$ and all $i,j$.
\end{lemma}
\leanref{Development/AutomorphicArithmetic/AutomorphicVarieties/ModularityArgument/ThresholdArguments/DenseThresholdCases/BFixOf.lean}{UpperTheory.V.J.fixOf_alt}
\begin{proof}
\begin{enumerate}
\item \emph{Vanishing.} The residual datum of the realization has rank $3$, its characters are trivial on $I_{v_0}$ and agree on $D_{v_0}$, and $\varepsilon\bmod\ell$ is nontrivial on $D_{v_0}$ (Definition~\ref{sha-automorphy-realization}(2)). So Theorem~\ref{sha-automorphy-unramified-near-vanishing} applies to $x$ and $\rho$: every ordinary cocycle for $\rho$ and the characters $\psi_{x,v,i}$ is a coboundary.
\item \emph{A first-order lift.} The map $a\mapsto\xi'(a)+\epsilon D(a)$ is an $\mathcal O_1$-algebra map from $A_3^\circ$ to $\bar{\mathbf Q}_\ell[\epsilon]/(\epsilon^2)$, by the Leibniz rule and $D(1)=0$. Applied to $\rho^{\mathrm u}$ it gives $\rho_{\xi'}+\epsilon\rho'$ with $\rho'(g)=D(\rho^{\mathrm u}(g))$, entrywise. Multiplicativity gives $\rho'(gg')=\rho_{\xi'}(g)\rho'(g')+\rho'(g)\rho_{\xi'}(g')$, so $y(g)=\rho'(\pi(g))\rho_{\xi'}(\pi(g))^{-1}$ satisfies $y(gg')=y(g)+\rho(g)y(g')\rho(g)^{-1}$ for $g,g'\in G_0$.
\item \emph{An ordinary cocycle.} There are $P$, an adjoint extension $R$ for $P$ and an ordinary cocycle $z$ for $\rho$ and the $\psi_{x,v,i}$ that restricts to $y$ on $G_0$. Let $G$ be the group $\Gamma$ of Situation~\ref{sha-hecke-situation} for the residual datum of the realization, which contains $G_0$ with index two, and $c\in G\setminus G_0$ its complex conjugation, with $c^2=1$; write $g^c=cgc^{-1}$. Let $\rho_D=\rho_{\xi'}+\epsilon\rho'$ be the image of $\rho^{\mathrm u}$ under the map of (2), and $\tilde\rho(g)=\kappa_0(g)\rho_D(\pi(g))$ for $g\in G_0$, so that $\tilde\rho(g)=(1+\epsilon y(g))\rho(g)$.
\begin{enumerate}
\item[(a)] \emph{Continuity.} By Lemma~\ref{sha-automorphy-unramified-near-small-continuity}, applied to $\xi'$ and $D$, the entries of $\rho_{\xi'}$ and of $\rho'$ are continuous on $\Gamma_0$. As $\pi$ is continuous and $\rho_{\xi'}(\pi(g))^{-1}=\rho_{\xi'}(\pi(g^{-1}))$, $y$ is continuous.
\item[(b)] \emph{The matrix $P$.} Condition (a) of Definition~\ref{sha-descent-link}(2), applied through the map of (2), gives $\rho_D(c_\Gamma h^{-1}c_\Gamma^{-1})=J_3\,{}^t\rho_D(h)J_3$, so ${}^t\rho_D(h)J_3\rho_D(c_\Gamma hc_\Gamma^{-1})=J_3$ for $h\in\Gamma_0$, as $J_3^2=1$. The map $\pi$ carries $g^c$ to $c_\Gamma\pi(g)c_\Gamma^{-1}$, and $\kappa_0(g)\kappa_0(g^c)=\varepsilon(g)^{-2}$ by Definition~\ref{sha-automorphy-realization}(2). So ${}^t\tilde\rho(g)J_3\tilde\rho(g^c)=\varepsilon(g)^{-2}J_3$ for $g\in G_0$. The part of this identity without $\epsilon$ says that $P=J_3$ satisfies ${}^t\rho(g)P\rho(g^c)=\varepsilon(g)^{1-3}P$, as in Definition~\ref{sha-automorphy-unramified-near-locally-generic}(2). Its coefficient of $\epsilon$ is ${}^t\rho(g)\bigl({}^ty(g)J_3+J_3y(g^c)\bigr)\rho(g^c)=0$, so $y(g^c)=-J_3^{-1}\,{}^ty(g)J_3$.
\item[(c)] \emph{The extension $R$.} By Definition~\ref{sha-automorphy-unramified-near-locally-generic}(1), $\rho$ is irreducible on the subgroup $\Gamma'$ of $G_0$, hence on $G_0$, so by Burnside's theorem its values span $M_3(\bar{\mathbf Q}_\ell)$. The cyclotomic character satisfies $\varepsilon(g^c)=\varepsilon(g)$. So Lemma~\ref{sha-automorphy-unramified-near-matrix-adjoint}, with $K=\bar{\mathbf Q}_\ell$, $n=3$ and $P=J_3$, gives an adjoint extension $R$ for $P$, with $R(c)X=-J_3^{-1}\,{}^tXJ_3$. By (b), $y(g^c)=R(c)y(g)$ for $g\in G_0$.
\item[(d)] \emph{The cocycle $z$.} Let $z(g)=y(g)$ for $g\in G_0$ and $z(g)=y(gc)$ for $g\in G\setminus G_0$. As $G_0$ and $G_0c$ are open and $y$ is continuous, $z$ is continuous. By (2), $z(gh)=z(g)+R(g)z(h)$ for $g,h\in G_0$. For $g\in G_0$ and $h\notin G_0$, $z(gh)=y(g\cdot hc)=z(g)+R(g)z(h)$. For $g\notin G_0$ and $h\in G_0$, $ghc=gc\cdot h^c$, so $z(gh)=y(gc)+R(gc)R(c)y(h)=z(g)+R(g)z(h)$ by (c). For $g,h\notin G_0$, $gh=gc\cdot(hc)^c$, so $z(gh)=y(gc)+R(gc)R(c)y(hc)=z(g)+R(g)z(h)$. So $z$ is a continuous cocycle of $G$ for $R$ that restricts to $y$ on $G_0$.
\item[(e)] \emph{Ordinarity away from $\ell$.} Let $u\in S_{\mathrm a}$ and $g\in I_u\cap\bigcap_i\ker\bar\chi_i$. By Definition~\ref{sha-automorphy-realization}(2) there is an integer $m$ prime to $\ell$ with $\pi(g)^m=1$ and $\kappa_0(g)^m=1$. The matrix $M=\kappa_0(g)\rho^{\mathrm u}(\pi(g))$ over $A_3$ has $M^m=1$. By condition (e), its residue is the diagonal matrix with entries $\bar\Psi_\Gamma(\pi(g))^{w_{3,i}}\bar\kappa_0(g)=\bar\chi_i(g)=1$, as the residual characters of the realization are those of the link composed with $\pi$ and twisted by $\kappa_0$ (Definition~\ref{sha-automorphy-realization}(2)); so its characteristic polynomial reduces to $(X-1)^3$. As $m$ is prime to the residue characteristic $\ell$, it is a unit of $A_3$, and the Cayley-Hamilton theorem and Lemma~\ref{sha-automorphy-unramified-near-matrix-unipotent} give $M=1$. Applying the map of (2) to its entries gives $\tilde\rho(g)=1$, so $\rho(g)=1$ and $y(g)=0$; that is, $z(g)=0$.
\item[(f)] \emph{Ordinarity above $\ell$.} Fix $v\mid\ell$, and identify $D_v$ with the group $G_v$ of the link, so that $\pi$ restricts to $\mathrm{loc}_v$. The ring of dual numbers $\bar{\mathbf Q}_\ell[\epsilon]/(\epsilon^2)$ is local, so Lemma~\ref{sha-automorphy-unramified-near-completion-flag}, applied to the map of (2), gives an invertible matrix $G_v^\epsilon$ over it such that $G_v^\epsilon\rho_D(\mathrm{loc}_vh)(G_v^\epsilon)^{-1}$ is upper triangular with diagonal entries $\xi'(\psi^{\mathrm u}_{v,i}(h))+\epsilon D(\psi^{\mathrm u}_{v,i}(h))$. Multiplying by $\kappa_0(h)$, $E(h)=G_v^\epsilon\tilde\rho(h)(G_v^\epsilon)^{-1}$ is upper triangular with diagonal entries $\psi_{x,v,i}(h)+\epsilon\kappa_0(h)D(\psi^{\mathrm u}_{v,i}(h))$, by the hypothesis on the $\psi_{x,v,i}$.
Write $G_v^\epsilon=a+\epsilon A_1$ with $a\in GL_3(\bar{\mathbf Q}_\ell)$, so that $(G_v^\epsilon)^{-1}=a^{-1}-\epsilon a^{-1}A_1a^{-1}$; let $X=a^{-1}A_1$, and let $E_1(h)$ be the coefficient of $\epsilon$ in $E(h)$. The part of $E(h)$ without $\epsilon$ is $a\rho(h)a^{-1}$, upper triangular with diagonal entries $\psi_{x,v,i}(h)$. Expanding $E(h)=(a+\epsilon A_1)(1+\epsilon y(h))\rho(h)(a^{-1}-\epsilon a^{-1}A_1a^{-1})$ gives
\[a\bigl(y(h)-\rho(h)X\rho(h)^{-1}+X\bigr)a^{-1}=E_1(h)\cdot a\rho(h)^{-1}a^{-1}.\]
The right side is a product of upper triangular matrices, as $a\rho(h)^{-1}a^{-1}=a\rho(h^{-1})a^{-1}$; so it is upper triangular, with diagonal entries $\kappa_0(h)D(\psi^{\mathrm u}_{v,i}(h))\psi_{x,v,i}(h)^{-1}$, which vanish for $h\in I_v$ by the hypothesis on $D$. As $z$ equals $y$ on the image of $D_v$ in $G_0$, $a$ and $X$ satisfy the condition at $v$ of Definition~\ref{sha-automorphy-unramified-near-locally-generic}(2).
\end{enumerate}
By (b) to (f), $R$ is an adjoint extension for $P=J_3$, and $z$ is an ordinary cocycle for $\rho$ and the $\psi_{x,v,i}$ that restricts to $y$ on $G_0$.
\item \emph{Conclusion.} By (1) there is $X'$ with $z(g)=R(g)X'-X'$, so $\rho'(\pi(g))\rho_{\xi'}(\pi(g))^{-1}=\rho(g)X'\rho(g)^{-1}-X'$ for $g\in G_0$. As $\rho(g)$ is a scalar multiple of $\rho_{\xi'}(\pi(g))$ and $\pi$ is surjective, $\rho'(h)=\rho_{\xi'}(h)X'-X'\rho_{\xi'}(h)$ for $h\in\Gamma_0$; take $X=-X'$.
\end{enumerate}
% fixOf_alt (BFixOf.lean, 343 lines in all): hC (DenseThresholdMainCase, Theorem sha-automorphy-unramified-near-vanishing) with degEq, R.r and R.v0 for (1); fxD for (2), with fx_tan and fx_coc; the private fx_core for (3) and (4). In (3): derRep (Lemma sha-automorphy-unramified-near-small-continuity) for (a); fx_form and Ws.Eb.cocRel (CocRel.lean) for (b), whose matrix Z = P^-1 snd(P') is 0 here, as P' = Jgl has no first-order part; Det.irrIffSpan (Burnside) and Ws.Eb.adExt (Lemma sha-automorphy-unramified-near-matrix-adjoint) for (c); Ws.Eb.cocExt (CocExt.lean) for (d), whose w = (R(c)Z - Z)/2 is 0, so that z(g) = y(gc) off G_0; fx_ine and fx_ineS for (e), with K.tame (the m prime to l, a condition of Definition sha-automorphy-realization(2), proved in Lemma sha-automorphy-realization-exists by cr_tame), K.chib for the residue of K.tw rho g = kappa(g) rho(pi g), Pr.T.unipTorsion (Lemma sha-automorphy-unramified-near-matrix-unipotent) and Matrix.aeval_self_charpoly; fx_ord (with fx_alg for the displayed identity and tri_mul), fx_twflag, Fr.flag (Lemma sha-automorphy-unramified-near-completion-flag) and fx_snd for (f). Then -X in (4).
\end{proof}

\begin{lemma}\label{sha-automorphy-unramified-near}
In the situation of Lemma~\ref{sha-automorphy-unramified-realization}, let $\xi:A_3^\circ\to\bar{\mathbf Q}_\ell$ be an $\mathcal O_1$-algebra map with values of absolute value at most $1$, with $|\xi(b)|<1$ for $b$ in the maximal ideal of $B$, that satisfies the condition on entries of Lemma~\ref{sha-automorphy-unramified-point}(2). Then for each $a\in J_3$ there is $s\in A_3^\circ$ with $\xi(s)\ne0$ and $s\phi(a)=0$.
\end{lemma}
\leanref{Development/AutomorphicArithmetic/AutomorphicVarieties/ModularityArgument/ThresholdArguments/DenseThresholdCases/EndPtOf.lean}{UpperTheory.Gal.projectedFormsSystemPart.Red2Alt}
\leanref{Development/AutomorphicArithmetic/AutomorphicVarieties/ModularityArgument/ThresholdArguments/DenseThresholdCases/EndPtOf.lean}{UpperTheory.Gal.projectedFormsSystemPart.red2_alt}
\begin{proof}
Call an $\mathcal O_1$-algebra map $\xi':A_3^\circ\to\bar{\mathbf Q}_\ell$ a \emph{small point} if it has values of absolute value at most $1$, $|\xi'(b)|<1$ for $b$ in the maximal ideal of $B$, and $|\xi'(e)|<1$ for the entries $e$ off the diagonal of the universal linked point. The proof uses the following facts about the realization of Lemma~\ref{sha-automorphy-unramified-realization} and its ring $A_3^\circ$, which is noetherian by Lemma~\ref{sha-automorphy-realization-presentation}(4).
\begin{enumerate}
\item[(A)] There is an increasing weight $h$ such that at every small point $\xi'$ with off-triangular kernel containing $\phi(J_3)$ and with $\xi'\circ\phi$ of distinct type of weight $h$, the tangent vectors of $A_3^\circ$ at $\xi'$ with values in $\bar{\mathbf Q}_\ell$ that kill the values of the $\psi^{\mathrm u}_{v,i}$ on the $I_v$ form a space of dimension at most $1$. By Lemma~\ref{sha-automorphy-unramified-near-torus}, it suffices that each such tangent vector $D$ have a matrix $X$ with $D(\rho^{\mathrm u}(g)_{ij})=(X\rho_{\xi'}(g)-\rho_{\xi'}(g)X)_{ij}$ for all $g$, $i$ and $j$, where $\rho_{\xi'}(g)=\xi'(\rho^{\mathrm u}(g))$.
\item[(B)] Every prime $Z$ of $B$ that contains $\phi(J_3)$ has $\dim B/Z\le1+3d$; if equality holds, then for every $b\notin Z$ there is $y:B\to\bar{\mathbf Q}_\ell$ vanishing on $Z$, with $y\circ\phi$ of distinct type of weight $h$ and $y(b)\ne0$. This is Lemma~\ref{sha-automorphy-unramified-near-dense}, applied to the noetherian ring $B_3$, the surjection $\phi$, $J_3$, $D=1+3d$ and the set of points of distinct type of weight $h$: every minimal prime over $J_3$ has dimension $1+3d$ by Lemma~\ref{sha-automorphy-unramified-realization}, and the density (2) of Proposition~\ref{sha-hecke-ideal} for the increasing weight $h$ is the second hypothesis.
\item[(C)] For a prime $\mathfrak q$ of $A_3^\circ$ that is the kernel of an $\mathcal O_1$-algebra map to $\bar{\mathbf Q}_\ell$ with values of absolute value at most $1$ and $|\cdot|<1$ on the maximal ideal of $B$: $\mathfrak q$ has height at least $1+3d$; if $\mathfrak q$ is off-triangular, then, by the bound of (B) and Lemma~\ref{sha-automorphy-unramified-realization}, every associated prime of $A_3^\circ$ contained in $\mathfrak q$ is minimal, and $\dim B/(P\cap B)\ge1+3d$ for every minimal prime $P\subseteq\mathfrak q$. This is Lemma~\ref{sha-automorphy-unramified-near-heights}, with $n=3$ and $M=\phi(J_3)$: every off-triangular prime of $A_3^\circ$ contains $\phi(J_3)$ by Lemma~\ref{sha-automorphy-unramified-realization}, and the bound on $\dim B/Z$ is the first part of (B).
\item[(D)] If $P$ is a minimal prime of $A_3^\circ$ that is off-triangular, and the points $y$ of $B$ of (B) vanishing on $P\cap B$ are dense in the sense of (B), then there is a small point $\xi'$ with off-triangular kernel containing $P$ and with $\xi'\circ\phi$ of distinct type of weight $h$. This is Lemma~\ref{sha-automorphy-unramified-near-generic}, with $\mathcal G$ the set of $y$ with $y\circ\phi$ of distinct type of weight $h$.
\item[(E)] At an $\mathcal O_1$-algebra map $\xi'$ to $\bar{\mathbf Q}_\ell$ with values of absolute value at most $1$ and $|\cdot|<1$ on the maximal ideal of $B$, the dimension of the tangent space of $A_3^\circ$ with values in $\bar{\mathbf Q}_\ell$ is at most $3d$ plus that of the subspace of tangent vectors that kill the values of the $\psi^{\mathrm u}_{v,i}$ on the $I_v$. This is Lemma~\ref{sha-automorphy-unramified-near-tangent}.
\end{enumerate}
\emph{Proof of (A).} Let $h$ be the increasing weight of Lemma~\ref{sha-automorphy-unramified-near-weight}. Let $\xi'$ be a small point with off-triangular kernel such that $x=\xi'\circ\phi$ is of distinct type of weight $h$, and let $\rho$ be as in Lemma~\ref{sha-automorphy-unramified-near-small}. By Lemma~\ref{sha-automorphy-unramified-near-weight}, for every $v$ some $g\in D_v$ with $\varepsilon(g)=1$ has $\psi_{x,v,i}(g)^m\ne\psi_{x,v,j}(g)^m$ for $i\ne j$ and $m\ge1$. So Lemma~\ref{sha-automorphy-unramified-near-small} shows that $x$ and $\rho$ form an automorphic locally generic point, with $\psi_{x,v,i}=\kappa_0\cdot\xi'\circ\psi^{\mathrm u}_{v,i}$. By Lemma~\ref{sha-automorphy-unramified-near-matrix}, every tangent vector $D$ of $A_3^\circ$ at $\xi'$ that kills the values of the $\psi^{\mathrm u}_{v,i}$ on the $I_v$ has a matrix $X$ as in (A), and Lemma~\ref{sha-automorphy-unramified-near-torus} bounds the dimension by $1$.
% (A): wtEx (Lemma sha-automorphy-unramified-near-weight), niceAt2_alt with EndPtAlt (Lemma sha-automorphy-unramified-near-small), fixOf_alt with Th.DenseThresholdMainCase (Lemmas sha-automorphy-unramified-near-matrix and sha-automorphy-unramified-near-vanishing). The bound is fixLe (Lemma sha-automorphy-unramified-near-torus, with torus and derExt). (C) and (D) are written out above: (C) is locAt, htUp, htLow and compDim (Lemma sha-automorphy-unramified-near-heights), (D) is genPt (Lemma sha-automorphy-unramified-near-generic).
Let $\mathfrak q=\ker\xi$, a prime of the noetherian ring $A_3^\circ$, let $a\in J_3$ and $b=\phi(a)$.
\begin{enumerate}
\item \emph{$\mathfrak q$ is off-triangular.} For the first column, the last column and, as $n=3$, the middle column, the condition on entries gives a nonzero entry below the diagonal, one above it, and one in position $(1,2)$ or $(3,2)$, where $w_{3,i}-w_{3,j}$ is odd. So every prime contained in $\mathfrak q$ is off-triangular, and contains $b$ by Lemma~\ref{sha-automorphy-unramified-realization}.
\item \emph{Minimal primes.} Let $P\subseteq\mathfrak q$ be a minimal prime of $A_3^\circ$. By (C) and (B), $\dim B/(P\cap B)=1+3d$, so (D) gives a small point $\xi'$ with $P\subseteq\ker\xi'$. By (A) and (E), the tangent space of $A_3^\circ$ at $\xi'$ has dimension at most $3d+1$, and by (C) the height of $\ker\xi'$ is at least $3d+1$. So the localization of $A_3^\circ$ at $\ker\xi'$ is a regular local ring, hence a domain \cite[\S19]{Mat86}, and $P$ is its only minimal prime; so every element of $P$ is killed by some element outside $P$.
\item \emph{Conclusion.} Suppose that every $s\in A_3^\circ$ with $sb=0$ lies in $\mathfrak q$. Then the annihilator $I=\{s\in A_3^\circ:sb=0\}$ of the ideal $bA_3^\circ$ lies in $\mathfrak q$, so some prime $P\subseteq\mathfrak q$ is minimal over $I$. Over a noetherian ring, a prime minimal over the annihilator of a finitely generated module is an associated prime of that module \cite[\S6]{Mat86}; so $P$ is an associated prime of $bA_3^\circ$, hence of $A_3^\circ$, which contains $bA_3^\circ$. By (C), $P$ is a minimal prime of $A_3^\circ$. By (1), $b\in P$, so by (2) some $s\notin P$ has $sb=0$. Then $s\in I\subseteq P$, a contradiction. So $sb=0$ for some $s\notin\mathfrak q$, that is, with $\xi(s)\ne0$.
\end{enumerate}
\end{proof}
% red2_alt (EndPtOf.lean): (A) is wtEx with niceAt2_alt (from EndPtAlt), J.fixOf_alt (from Th.DenseThresholdMainCase: OrdZero at the locally generic automorphic points) and fixLe; (B) is denseOf (DenseB); the rest is redAbs (RedAbs.lean): (C) is locAt, htUp and htLow (through a localization L.R with a prime Q over q) and compDim; (D) is genPt; (E) is tanBound (Lemma sha-automorphy-unramified-near-tangent); regDom and domMin give (2); redCrit gives (3). Step (1) is ti_of. The noetherian ring S.A of redAbs is Fr.gr_A. DenseThresholdMainCase is proved separately; the other inputs are in LinkedUnramifiedChains/ScalarIrreducibility and ThresholdArguments/DenseThresholdCases.

\subsection{Residually reducible automorphy}\label{subsection-sha-residually-reducible}

\begin{lemma}\label{sha-automorphy-unramified-tangents-rank}
Let $\ell$ be a prime, $\mathcal O$ a complete discrete valuation ring with finite residue field of characteristic $\ell$ in which $\ell\ne0$, $R$ a noetherian $\mathcal O$-algebra, $\xi:R\to\mathcal O$ an $\mathcal O$-algebra map, $T\subseteq R$ a subset and $C\ge0$ an integer. For $k\ge1$ let $X_k$ be the set of $\mathcal O$-linear maps $\delta:R\to\mathcal O/\ell^k$ with $\delta(ab)=\xi(a)\delta(b)+\xi(b)\delta(a)$ for all $a,b$ and $\delta(T)=0$. Suppose that $|\mathcal O/\ell^k|^2\le C|X_k|$ for every $k\ge1$ for which $X_k$ is finite. Then there are two $\mathcal O$-linearly independent $\mathcal O$-linear maps $\delta,\delta':R\to\mathcal O$ that satisfy the same rule at $\xi$ and vanish on $T$.
\end{lemma}
\leanref{Development/AutomorphicArithmetic/AutomorphicVarieties/ModularityArgument/LinkedUnramifiedChains/CrystallineCardinalities/RankTwo.lean}{UpperTheory.V.Z.J.rankTwo}
\begin{proof}
\begin{enumerate}
\item \emph{The module of relations.} Let $N$ be the $\mathcal O$-submodule of $R$ spanned by $T$ and the elements $ab-\xi(a)b-\xi(b)a$ for $a,b\in R$, and $Q=R/N$. An $\mathcal O$-linear map from $R$ to an $\mathcal O$-module satisfies the rule at $\xi$ and vanishes on $T$ if and only if it vanishes on $N$; so these maps are the composites of the $\mathcal O$-linear maps from $Q$ with the quotient map, and they correspond bijectively to them.
\item \emph{$Q$ is finite over $\mathcal O$.} The ideal $\ker\xi$ is generated by finitely many $s_1,\dots,s_m$, as $R$ is noetherian. The element $1\cdot1-\xi(1)1-\xi(1)1=-1$ lies in $N$, and so does $ab$ whenever $\xi(a)=\xi(b)=0$. For $r\in R$, $r-\xi(r)\cdot1$ lies in $\ker\xi$, so it is a sum of terms $r's_j$, and $r's_j=\xi(r')s_j+(r'-\xi(r')\cdot1)s_j$ with the last term in $N$. So $R$ is spanned over $\mathcal O$ by $1$, the $s_j$ and $N$, and $Q$ is spanned by the images of the $s_j$.
\item \emph{Two cases.} By the structure theorem for finitely generated modules over the principal ideal domain $\mathcal O$, $Q$ is isomorphic to $\mathcal O^m\oplus\bigoplus_i\mathcal O/(p_i^{e_i})$ with $p_i$ prime and $e_i\ge1$. If $m\ge2$, the first two coordinate functionals $Q\to\mathcal O$, composed with $R\to Q$, are two maps as in the statement, linearly independent because they take the values $(1,0)$ and $(0,1)$ at the images of the first two basis vectors. If $m\le1$, then, as an $\mathcal O$-linear map from $\mathcal O/(p_i^{e_i})$ to $\mathcal O/\ell^k$ is determined by the image of $1$, an element killed by $p_i^{e_i}$, and the kernel of multiplication by $p_i^{e_i}$ on $\mathcal O/\ell^k$ has as many elements as its cokernel, a quotient of $\mathcal O/(p_i^{e_i})$, there are at most $|\mathcal O/\ell^k|\cdot C'$ linear maps $Q\to\mathcal O/\ell^k$, with $C'=\prod_i|\mathcal O/(p_i^{e_i})|$, for every $k$.
\item \emph{The second case does not occur.} By (1), $|X_k|$ is the number of linear maps $Q\to\mathcal O/\ell^k$, which is finite since $Q$ is finitely generated and $\mathcal O/\ell^k$ is finite. In the second case the hypothesis gives $|\mathcal O/\ell^k|^2\le CC'|\mathcal O/\ell^k|$, so $|\mathcal O/\ell^k|\le CC'$ for every $k$. But $|\mathcal O/\ell^k|\ge2^k$, since $\ell\ne0$ in $\mathcal O$ is not a unit.
\end{enumerate}
\end{proof}
% rankTwo (RankTwo.lean): rel and rel_le for (1), rel_fin for (2), hom_alt (Module.equiv_free_prod_directSum, tors_card) for (3), and the final count with k = C C' + 1 for (4).

\begin{lemma}[Bloch-Kato {\cite[Proposition~3.8]{BK90}}]\label{sha-automorphy-unramified-tangents-unit-duality}
Let $K$ be a finite extension of $\mathbf Q_\ell$ in $\bar{\mathbf Q}_\ell$, $G_K=\operatorname{Gal}(\bar{\mathbf Q}_\ell/K)$, $\chi$ the $\ell$-adic cyclotomic character of $G_K$, and $B_{\mathrm{cris}}$ Fontaine's ring of crystalline periods. For a finite-dimensional $\mathbf Q_\ell$-vector space $V$ with a continuous linear action of $G_K$, call a continuous map $a:G_K\to V$ with $a(ss')=a(s)+sa(s')$ a \emph{crystalline cocycle} if there is $x\in V\otimes_{\mathbf Q_\ell}B_{\mathrm{cris}}$ with $a(s)\otimes1=sx-x$ for all $s\in G_K$; these are the cocycles whose classes lie in $H^1_f(K,V)$ \cite[(3.7.2)]{BK90}. Let $V$ and $V'$ be two such spaces, with $V$ de Rham, and let $P:V\times V'\to\mathbf Q_\ell$ be a perfect bilinear pairing with $P(sx,sy)=\chi(s)P(x,y)$. Then:
\begin{enumerate}
\item a continuous cocycle $a$ of $V$ is crystalline if and only if for every crystalline cocycle $b$ of $V'$ there is a continuous $f:G_K\to\mathbf Q_\ell$ with $P(a(s),sb(s'))=\chi(s)f(s')-f(ss')+f(s)$ for all $s,s'\in G_K$;
\item a continuous cocycle $b$ of $V'$ is crystalline if and only if the same holds for every crystalline cocycle $a$ of $V$.
\end{enumerate}
\end{lemma}
\leanref{Development/AbelianVarietyTheory/PublishedTheorems/GaloisCohomology₂/RationalDualityOfAllCoefficients.lean}{GlobalTheory.Published.BlochKato.RationalDuality}
\leanref{Development/AbelianVarietyTheory/PublishedTheorems/GaloisCohomology₂/RationalDualityOfAllCoefficients.lean}{GlobalTheory.Published.BlochKatoLocal}
% The Lean assumes this as the field RationalDuality of GlobalTheory.Published.BlochKatoLocal (with DimFormula and IntegralDuality), through the hypothesis LocalPublishedTheorems (field bk); it does not prove it. The locator was checked in the published paper: (3.7.2) defines H^1_f(K,V) as the kernel of H^1(K,V) -> H^1(K,B_crys (x) V), and Proposition 3.8 states that H^1_f(K,V) and H^1_f(K,V*(1)) are exact annihilators of each other under the cup-product pairing, for V de Rham. The Lean's form, with V' and a perfect pairing P to Q_l(1) in place of V*(1), is the same statement.

\begin{lemma}[Fontaine {\cite[\S\S4.4 and 9.1]{BC09}}]\label{sha-automorphy-unramified-tangents-unit-crystalline-periods}
Let $K$ be a finite extension of $\mathbf Q_\ell$ in $\bar{\mathbf Q}_\ell$, $G_K=\operatorname{Gal}(\bar{\mathbf Q}_\ell/K)$, $I_K\subseteq G_K$ the inertia subgroup and $\chi$ the $\ell$-adic cyclotomic character of $G_K$. Let $\mathbf Q_\ell^{\mathrm{ur}}\subseteq\bar{\mathbf Q}_\ell$ be the maximal unramified extension of $\mathbf Q_\ell$ and $K_0=K\cap\mathbf Q_\ell^{\mathrm{ur}}$. Let $C$ be the completion of $\bar{\mathbf Q}_\ell$, $\mathcal O_C$ its ring of integers, $\mathcal O_C^\flat=\varprojlim_{x\mapsto x^\ell}\mathcal O_C/\ell\mathcal O_C$ its tilt, a perfect domain of characteristic $\ell$ on which $G_K$ acts by ring automorphisms through its action on $\mathcal O_C$, and $A_{\mathrm{inf}}=W(\mathcal O_C^\flat)$ its ring of Witt vectors, on which $G_K$ acts by functoriality. Let $B_{\mathrm{cris}}\subseteq B_{\mathrm{dR}}$ be Fontaine's rings of crystalline and de Rham periods, with their actions of $G_K$. Then:
\begin{enumerate}
\item $B_{\mathrm{dR}}$ is a field containing $\bar{\mathbf Q}_\ell$, $G_K$ acts on it by field automorphisms that extend its action on $\bar{\mathbf Q}_\ell$, and $B_{\mathrm{dR}}^{G_K}=K$;
\item $B_{\mathrm{cris}}$ is a $G_K$-stable subring of $B_{\mathrm{dR}}$ that contains $K_0$, and $B_{\mathrm{cris}}^{G_K}=K_0$;
\item for every finite extension $L$ of $\mathbf Q_\ell$ in $\bar{\mathbf Q}_\ell$, with $L_0=L\cap\mathbf Q_\ell^{\mathrm{ur}}$, the map $L\otimes_{L_0}B_{\mathrm{cris}}\to B_{\mathrm{dR}}$, $x\otimes b\mapsto xb$, is injective;
\item every nonzero $b\in B_{\mathrm{cris}}$ such that for every $g\in G_K$ there is $x\in\mathbf Q_\ell$ with $g(b)=xb$ is a unit of $B_{\mathrm{cris}}$;
\item there is an injective $G_K$-equivariant ring map $A_{\mathrm{inf}}\to B_{\mathrm{cris}}$, whose restriction to $\mathbf Z_\ell=W(\mathbf F_\ell)\subseteq A_{\mathrm{inf}}$ is the inclusion $\mathbf Z_\ell\subseteq\mathbf Q_\ell\subseteq B_{\mathrm{cris}}$;
\item there is a unit $t$ of $B_{\mathrm{cris}}$ with $g(t)=\chi(g)t$ for every $g\in G_K$.
\end{enumerate}
\end{lemma}
\leanref{Development/AbelianVarietyTheory/Proofs/IwasawaDecomposition/ExistsLocalPrimeAinf.lean}{GlobalTheory.MemLoc.exists_locPrime_tied}
\leanref{Formalism/WeaklyAdmissibleFiltrations/Statement.lean}{Formalism.WeaklyAdmissibleFiltrations.PeriodData}
\leanref{Formalism/SharpRestriction/LevelLetterCocycles/DataNLoc.lean}{Formalism.BasicNotions.LocAinf}
\leanref{Development/AbelianVarietyTheory/Proofs/InertiaEigenvalues/LinearIndepOnOfInvariant.lean}{GlobalTheory.Per.FontaineBound.isField_of_tied}
\leanref{Development/AbelianVarietyTheory/Proofs/InertiaEigenvalues/AlgClosureTensorBcrisToBdRInjectiveStatement.lean}{GlobalTheory.T.algClosure_tensor_bcris_to_bdR_injective}
\leanref{Development/AbelianVarietyTheory/Proofs/InertiaEigenvalues/BcrisRegularOfTensorInjStatement.lean}{GlobalTheory.T.bcris_regular_of_tensorInj}
% Locators checked in the preliminary version (June 2009) of [BC09]: (1) is Definition 4.4.7, Lemma 4.4.10 and the text after it, and Theorem 4.4.13; (2) is Definition 9.1.4 and the claim after it; (3) is Theorem 9.1.5 (Brinon and Conrad note that B_cris does not depend on the base field); (4) is Proposition 9.1.6; (5) is (4.4.1) with Definition 9.1.4; (6) is the computation g(t) = chi(g) t in Section 4.4 with Proposition 9.1.3 and Definition 9.1.4. The Lean constructs these rings (Per.Ainf, Per.Acris, Per.Bcris, Per.BdR; PP l, AIP l and tiedP from exists_locPrime_tied, with the period data pd C K of exists_realPD) and proves (1) to (6) for them: (1) is isField_of_tied (from BdRField.bdR_isField), the field invdR of PeriodData and the map C.oD; (2) is the fields actB and invB of PeriodData; (3) is used in the Lean's form Core.TensorInj, that elements of Qbar_l linearly independent over Q_l^ur stay linearly independent over B_cris in B_dR, proved as algClosure_tensor_bcris_to_bdR_injective (TensorInj.tensorInj), which follows from (3) for a finite extension L containing the elements; (4) is Core.Regular, proved as bcris_regular_of_tensorInj; (5) is the fields iA, ιA_inj and ιA_act of LocAinf with ιA_zlW; (6) is the fields t, t_unit and t_act of PeriodData. In the Lean (1) to (6) are theorems, not hypotheses.

\begin{lemma}\label{sha-automorphy-unramified-tangents-unit-crystalline-tilt}
Let $K$, $G_K$, $I_K$ and $\mathcal O_C^\flat$ be as in Lemma~\ref{sha-automorphy-unramified-tangents-unit-crystalline-periods}, and for $m\ge1$ let $F_m$ be the set of $y\in\mathcal O_C^\flat$ with $y^{\ell^m}=y$.
\begin{enumerate}
\item For every $m\ge1$ and $y\in F_m$, the elements $g\in G_K$ with $g(y)=y$ form an open subgroup of $G_K$.
\item If $g\in G_K$ fixes every element of $F_m$ for every $m\ge1$, then $g\in I_K$.
\end{enumerate}
\end{lemma}
\leanref{Development/AbelianVarietyTheory/Proofs/InertiaEigenvalues/NormSubLtOneOfFix.lean}{GlobalTheory.UnramifiedCrys.eventually_actT_eq}
\leanref{Development/AbelianVarietyTheory/Proofs/InertiaEigenvalues/NormSubLtOneOfFix.lean}{GlobalTheory.UnramifiedCrys.norm_sub_lt_one_of_fix}
\leanref{Development/AutomorphicArithmetic/AutomorphicVarieties/PadicHodgeTheory/MemIK.lean}{UpperTheory.V.Loc.memIK}
\begin{proof}
Write $y\in\mathcal O_C^\flat$ as a sequence $(y^{(n)})_{n\ge0}$ of elements of $\mathcal O_C/\ell\mathcal O_C$ with $(y^{(n+1)})^\ell=y^{(n)}$; products are taken coordinate by coordinate, and $g(y)^{(n)}=g(y^{(n)})$ for $g\in G_K$.
\begin{enumerate}
\item[(a)] \emph{An element of $F_m$ is determined by $y^{(0)}$.} Let $y\in F_m$, $n\ge0$ and $k\ge1$ with $mk\ge n$. As $y=y^{\ell^{mk}}$, $y^{(n)}=(y^{(n)})^{\ell^{mk}}=\bigl((y^{(n)})^{\ell^n}\bigr)^{\ell^{mk-n}}=(y^{(0)})^{\ell^{mk-n}}$.
\item[(1)] \emph{Proof of (1).} Let $y\in F_m$. As $\bar{\mathbf Q}_\ell$ is dense in $C$, there is an integer $x$ of $\bar{\mathbf Q}_\ell$ whose class modulo $\ell$ is $y^{(0)}$. The stabilizer $U$ of $x$ in $G_K$ is open, as $x$ is algebraic over $\mathbf Q_\ell$. For $g\in U$, $g(y)$ lies in $F_m$ and $g(y)^{(0)}=g(y^{(0)})=y^{(0)}$, so $g(y)=y$ by (a). So the stabilizer of $y$ contains the open subgroup $U$.
\item[(2)] \emph{Proof of (2).} Let $x$ be an integer of $\bar{\mathbf Q}_\ell$. If $|x|<1$, then $|g(x)-x|<1$, as $|g(x)|=|x|$. Let $|x|=1$. As the residue field of $\bar{\mathbf Q}_\ell$ is algebraic over $\mathbf F_\ell$, $|x^n-1|<1$ for some $n\ge1$; write $n=\ell^kn'$ with $\ell\nmid n'$. Then $|x^{n'}-1|<1$, since in a field of characteristic $\ell$ an element $a$ with $a^{\ell^k}=1$ is $1$, and by Hensel's lemma there is $\zeta\in\bar{\mathbf Q}_\ell$ with $\zeta^{n'}=1$ and $|x-\zeta|<1$. Let $s=\varphi(n')$, so that $\zeta^{\ell^s}=\zeta$. The classes modulo $\ell$ of the $\zeta^{\ell^{(s-1)j}}$, $j\ge0$, form an element $y$ of $\mathcal O_C^\flat$ with $y^{\ell^s}=y$ and $y^{(0)}$ the class of $\zeta$. As $g(y)=y$, $g(\zeta)$ and $\zeta$ have the same class modulo $\ell$, so $|g(\zeta)-\zeta|<1$, and $|g(x)-x|\le\max(|g(x-\zeta)|,|g(\zeta)-\zeta|,|\zeta-x|)<1$. So $g$ acts trivially on the residue field of $\bar{\mathbf Q}_\ell$, that is, $g\in I_K$.
\end{enumerate}
\end{proof}
% eventually_actT_eq (with eq_of_coeff_zero for (a) and stabilizer_isOpen_of_isIntegral) for (1); norm_sub_lt_one_of_fix (ResidueRoot.exists_pow_alg, sub_one_lt_of_pow_p_pow, exists_root, the totient for s) for (2), and memIK: the Lean's I_K is the subgroup generated by the g with |g(x) - x| < 1 for every integer x. The action of G_K on the tilt is the field actT of LocAinf, with actT_coeff for g(y)^(n) = g(y^(n)).

\begin{lemma}\label{sha-automorphy-unramified-tangents-unit-crystalline-lang}
Let $\ell$ be a prime, $R$ a perfect domain of characteristic $\ell$ and $W(R)$ its ring of Witt vectors. Let $G$ be a compact topological group acting on $R$ by ring automorphisms, and on $W(R)$ by functoriality, such that every $y\in R$ with $y^{\ell^m}=y$ for some $m\ge1$ is fixed by a neighbourhood of $1$ in $G$. Let $r\ge0$ and let $M:G\to GL_r(\mathbf Z_\ell)$ be a continuous homomorphism with $M(g)=1$ for every $g\in G$ that fixes every $y\in R$ with $y^{\ell^m}=y$ for some $m\ge1$. Then there is $W\in M_r(W(R))$ whose determinant is a unit of $W(R)$, with $M(g)\,g(W)=W$ for every $g\in G$; here $\mathbf Z_\ell=W(\mathbf F_\ell)\subseteq W(R)$, and $g$ acts on matrices entry by entry.
\end{lemma}
\leanref{Development/AbelianVarietyTheory/Proofs/InertiaEigenvalues/UnramifiedIsCrysAndJumpEqZero.lean}{GlobalTheory.UnramifiedRep.exists_witt_matrix}
\begin{proof}
\begin{enumerate}
\item[(a)] \emph{Finite levels.} For $m\ge1$ let $F_m$ be the set of $y\in R$ with $y^{\ell^m}=y$, and $N_m$ the subgroup of the $g\in G$ that fix $F_m$ pointwise. As $R$ is a domain, $F_m$ is finite, the set of roots of $X^{\ell^m}-X$; it is a subfield of $R$, closed under sums as $y\mapsto y^{\ell^m}$ is additive, and containing the inverse $y^{\ell^m-2}$ of each nonzero $y$. As $F_m$ is finite and stable under $G$, $N_m$ is an open normal subgroup, by the hypothesis on $G$, and has finite index, as $G$ is compact; the finite group $G/N_m$ acts faithfully on $F_m$. If $m$ divides $m'$, then $F_m\subseteq F_{m'}$, so $N_{m'}\subseteq N_m$.
\item[(b)] \emph{Congruences.} For every $k\ge1$ there is $m\ge1$ with $M(g)\equiv1$ modulo $\ell^k$ for every $g\in N_m$. Indeed, the set $U$ of the $g$ with $M(g)\equiv1$ modulo $\ell^k$ is open, as $M$ is continuous, and it contains the intersection of the $N_m$, by the hypothesis on $M$. The $N_m$ are closed, being open subgroups, and $N_{mm'}\subseteq N_m\cap N_{m'}$; as $G\smallsetminus U$ is compact, some $N_m$ is contained in $U$.
\item[(c)] \emph{A solution modulo $\ell$.} Let $m$ be as in (b) for $k=1$, and let $\bar M(g)\in GL_r(\mathbf F_\ell)\subseteq GL_r(F_m)$ be the reduction of $M(g)$. Its entries are fixed by $G$ and $\bar M(g)=1$ for $g\in N_m$, so $\bar M$ is a $1$-cocycle of $G/N_m=\operatorname{Gal}(F_m/F_m^{G})$ with values in $GL_r(F_m)$. By Hilbert's Theorem~90 for $GL_r$ \cite[Chapter~X, Proposition~3]{Ser79} there is $W_1\in GL_r(F_m)$ with $\bar M(g)\,g(W_1)=W_1$ for every $g\in G$.
\item[(d)] \emph{Elements of trace one.} For every $m\ge1$ there is $\tau\in W(R)$, fixed by $N_m$, with $\sum_{x\in G/N_m}x(\tau)=1$. Indeed, the elements of $G/N_m$ restrict to distinct multiplicative maps $F_m\to R$, which are linearly independent over $R$ by Dedekind's lemma, so some $\theta\in F_m$ has $t=\sum_xx(\theta)\ne0$. Then $t\in F_m$ is a unit of $R$ fixed by $G$, and $T=\sum_xx([\theta])$, where $[\theta]$ is the Teichm\"uller lift, is fixed by $G$ and has constant coefficient $t$. As $R$ is perfect, the elements of $W(R)$ with constant coefficient $0$ are the multiples of $\ell$, and $W(R)$ is $\ell$-adically complete; so $T[t^{-1}]\in1+\ell W(R)$ is a unit, and so is $T$. Put $\tau=[\theta]T^{-1}$; it is fixed by $N_m$, and $\sum_xx(\tau)=T^{-1}T=1$.
\item[(e)] \emph{Averaging.} Let $N=N_m$, $\tau$ as in (d), $c,d\in W(R)$ with $M(g)\equiv1$ modulo $d$ for $g\in N$, and $W\in M_r(W(R))$ fixed by $N$ with $M(g)g(W)\equiv W$ modulo $c$ for every $g\in G$. Choose representatives $x$ of $G/N$ and put $W'=\sum_xM(x)\,x(\tau W)$. Then $W-W'=\sum_xx(\tau)\bigl(W-M(x)x(W)\bigr)\equiv0$ modulo $c$. For $n\in N$ and each $x$, $nx=xn'$ with $n'\in N$, as $N$ is normal, so $n(x(\tau W))=x(\tau W)$ and $W'$ is fixed by $N$. For $g\in G$, write $gx=yn_x$ with $y$ the representative of $gxN$ and $n_x\in N$; as $x\mapsto y$ is a bijection, $M(g)g(W')-W'=\sum_xM(y)\bigl(M(n_x)-1\bigr)y(\tau W)\equiv0$ modulo $d$.
\item[(f)] \emph{Successive approximation.} Let $W_0\in M_r(W(R))$ have as entries the Teichm\"uller lifts of the entries of $W_1$. It is fixed by $N_m$, and $M(g)g(W_0)\equiv W_0$ modulo $\ell$, as the entries of the difference have constant coefficients those of $\bar M(g)g(W_1)-W_1=0$. Put $m_0=m$. Given $W_k$ fixed by $N_{m_k}$ with $M(g)g(W_k)\equiv W_k$ modulo $\ell^{k+1}$ for all $g$, let $m'$ be as in (b) for $\ell^{k+2}$ and $m_{k+1}=m_km'$, and apply (e) with $N=N_{m_{k+1}}\subseteq N_{m_k}$, $c=\ell^{k+1}$ and $d=\ell^{k+2}$. It gives $W_{k+1}\equiv W_k$ modulo $\ell^{k+1}$, fixed by $N_{m_{k+1}}$, with $M(g)g(W_{k+1})\equiv W_{k+1}$ modulo $\ell^{k+2}$. As $W(R)$ is $\ell$-adically complete and separated, the $W_k$ converge to some $W\equiv W_0$ modulo $\ell$, and $M(g)g(W)=W$ for every $g$.
\item[(g)] \emph{The determinant.} As $\det W\equiv\det W_0$ modulo $\ell$, the constant coefficient of $\det W$ is $\det W_1$, a nonzero element of the field $F_m$; so $\det W$ is a unit, as in (d).
\end{enumerate}
\end{proof}
% exists_witt_matrix (UnramifiedIsCrysAndJumpEqZero.lean, with its helpers about 400 lines): (a) is finite_fixedPow, isOpen_fixer, normal_fixer and pow_pow_mul (NormSubLtOneOfFix.lean); (b) is exists_fixer_dvd and exists_fixer_subset; (c) is GlobalTheory.T.exists_unit_matrix_coboundary_of_fixed_frobenius_pow, which the Lean proves by Serre's argument (UnramifiedHilbertNinety.exists_matrix_coboundary, with Dedekind's lemma linearIndependent_monoidHom; it gives W_1 with entries in F_m); (d) is exists_witt_trace_one; (e) is InvariantLift.average_step; (f) is InvariantLift.exists_invariant_lift, with S the set of the N_m; (g) is isUnit_of_constantCoeff and dvd_iff_constantCoeff.

\begin{lemma}\label{sha-automorphy-unramified-tangents-unit-crystalline}
Let $K$, $G_K$, $\chi$, $B_{\mathrm{cris}}$ and the crystalline cocycles be as in Lemma~\ref{sha-automorphy-unramified-tangents-unit-duality}, let $I_K\subseteq G_K$ be the inertia subgroup and $K_0\subseteq K$ the largest subfield unramified over $\mathbf Q_\ell$. For a finite-dimensional $\mathbf Q_\ell$-vector space $V$ with a continuous linear action of $G_K$, let $D_K(V)=(V\otimes_{\mathbf Q_\ell}B_{\mathrm{cris}})^{G_K}$, for the diagonal action, and call $V$ \emph{crystalline} if $\dim_{\mathbf Q_\ell}D_K(V)=[K_0:\mathbf Q_\ell]\dim_{\mathbf Q_\ell}V$; a representation on a vector space over a finite extension of $\mathbf Q_\ell$ is crystalline if it is so as a representation over $\mathbf Q_\ell$. Then:
\begin{enumerate}
\item $\dim_{\mathbf Q_\ell}D_K(V)\le[K_0:\mathbf Q_\ell]\dim_{\mathbf Q_\ell}V$ for every such $V$;
\item if $f:V\to V'$ is a $G_K$-equivariant $\mathbf Q_\ell$-linear map, then $V$ is crystalline if $f$ is injective and $V'$ is crystalline, and $V'$ is crystalline if $f$ is surjective and $V$ is crystalline;
\item if $V$ is crystalline, then so is its restriction to $G_{K'}$, for every finite extension $K'$ of $K$ in $\bar{\mathbf Q}_\ell$;
\item if $L_1\subseteq L_2$ are finite extensions of $\mathbf Q_\ell$ and $V$ is a finite-dimensional $L_1$-vector space with a continuous $L_1$-linear action of $G_K$, then $V$ is crystalline if and only if $V\otimes_{L_1}L_2$ is;
\item if $E$ is a finite extension of $\mathbf Q_\ell$ and $\kappa:G_K\to E^\times$ a continuous character such that $E(\kappa)$, the space $E$ on which $g$ acts by multiplication by $\kappa(g)$, is crystalline, then $E(\kappa^{-1})$ is crystalline, and for every crystalline finite-dimensional $E$-vector space $V$ with a continuous $E$-linear action of $G_K$, the twist $V(\kappa)$, on which $g$ acts by $x\mapsto\kappa(g)gx$, is crystalline;
\item if there is an integer $w$ with $gx=\chi(g)^wx$ for all $g\in I_K$ and $x\in V$, then $V$ is crystalline;
\item let $V$ be crystalline, $V'$ a finite-dimensional $\mathbf Q_\ell$-vector space with a continuous linear action of $G_K$, $i:V\to V'$ an injective and $p:V'\to\mathbf Q_\ell$ a surjective $G_K$-equivariant linear map, for the trivial action on $\mathbf Q_\ell$, with $i(V)=\ker p$, and let $s\in V'$ with $p(s)=1$ and $c:G_K\to V$ with $i(c(g))=gs-s$; then $V'$ is crystalline if and only if $c$ is a crystalline cocycle.
\end{enumerate}
\end{lemma}
\leanref{Development/AbelianVarietyTheory/Proofs/InertiaEigenvalues/LinearIndepOnOfInvariant.lean}{GlobalTheory.Per.FontaineBound.finite_finrank_le_base}
\leanref{Development/AutomorphicArithmetic/AutomorphicVarieties/PadicHodgeTheory/AdmSub.lean}{UpperTheory.V.Loc.isCrys_isDR_sub_quot}
\leanref{Development/AbelianVarietyTheory/Proofs/InertiaEigenvalues/RestrictCrystallineDeRhamFiniteExt.lean}{GlobalTheory.Per.TiedRestrict.restrict_crys_deRham}
\leanref{Development/AbelianVarietyTheory/Proofs/IwasawaDecomposition/FiniteDimensionalFilDR.lean}{GlobalTheory.Pst.isCrys_ext_iff}
\leanref{Development/AutomorphicArithmetic/AutomorphicVarieties/PadicHodgeTheory/IsCrysOfInertia.lean}{UpperTheory.V.Loc.crUnr}
\leanref{Development/AutomorphicArithmetic/AutomorphicVarieties/PadicHodgeTheory/MatrixExtendByZeroApply.lean}{UpperTheory.V.Loc.crExt}
\leanref{Development/AutomorphicArithmetic/AutomorphicVarieties/PadicHodgeTheory/IsCrIff.lean}{UpperTheory.V.Loc.isCrIff}
\begin{proof}
Write $B=B_{\mathrm{cris}}$ and $d=\dim_{\mathbf Q_\ell}V$. The proof uses these properties of $B$, from Lemma~\ref{sha-automorphy-unramified-tangents-unit-crystalline-periods}(1), (2) and~(4): $G_K$ acts on $B$ by automorphisms of $K_0$-algebras, $B^{G_K}=K_0$, and every nonzero $b\in B$ such that each $g(b)$, $g\in G_K$, lies in $\mathbf Q_\ell b$ is a unit of $B$. Moreover, $B$ is a domain, being a subring of the field $B_{\mathrm{dR}}$ (Lemma~\ref{sha-automorphy-unramified-tangents-unit-crystalline-periods}(1) and~(2)), and it is \emph{regular}: if $b,b'\in B$, $b\ne0$ and $b\,g(b')=g(b)\,b'$ for every $g\in G_K$, then $b'\in K_0b$. Indeed, $c=b'/b\in B_{\mathrm{dR}}$ is fixed by $G_K$, so $c\in K$ by Lemma~\ref{sha-automorphy-unramified-tangents-unit-crystalline-periods}(1). If $c$ does not lie in the maximal unramified extension $\mathbf Q_\ell^{\mathrm{ur}}$ of $\mathbf Q_\ell$ in $\bar{\mathbf Q}_\ell$, then $1$ and $c$ are linearly independent over $\mathbf Q_\ell^{\mathrm{ur}}$, hence over $K_0$, so $1\otimes b'-c\otimes b$ is a nonzero element of $K\otimes_{K_0}B$, as $b\ne0$, whose image $b'-cb$ in $B_{\mathrm{dR}}$ vanishes; this contradicts Lemma~\ref{sha-automorphy-unramified-tangents-unit-crystalline-periods}(3) for $L=K$. So $c\in K\cap\mathbf Q_\ell^{\mathrm{ur}}=K_0$, and $b'=cb$.
\begin{enumerate}
\item[(a)] \emph{Invariant vectors (Fontaine).} Elements $x_1,\dots,x_m$ of $D_K(V)$ that are linearly independent over $K_0$ are linearly independent over $B$ in the free $B$-module $V\otimes_{\mathbf Q_\ell}B$. Indeed, let $T\subseteq\{1,\dots,m\}$ be a set of least size carrying a relation $\sum_{j\in T}b_jx_j=0$ with some $b_{j_0}\ne0$. For $g\in G_K$, applying $g$, which is semilinear and fixes the $x_j$, gives $\sum_jg(b_j)x_j=0$; so $\sum_{j\in T}\bigl(b_{j_0}g(b_j)-g(b_{j_0})b_j\bigr)x_j=0$, a relation whose coefficient at $j_0$ vanishes. By the minimality of $T$, $b_{j_0}g(b_j)=g(b_{j_0})b_j$ for all $j\in T$ and $g$, and regularity gives $c_j\in K_0$ with $b_j=c_jb_{j_0}$. Then $b_{j_0}\sum_jc_jx_j=0$, so $\sum_jc_jx_j=0$, as $B$ is a domain and $V\otimes B$ is free; hence every $c_j$ vanishes, and $b_{j_0}=c_{j_0}b_{j_0}=0$, a contradiction.
\item[(1)] \emph{Proof of (1).} As $G_K$ acts $K_0$-linearly on $B$, $D_K(V)$ is a $K_0$-vector space. By (a), it has at most $d$ linearly independent elements over $K_0$, since $V\otimes B$ is free of rank $d$ over the domain $B$. So $\dim_{\mathbf Q_\ell}D_K(V)\le[K_0:\mathbf Q_\ell]d$.
\item[(b)] \emph{Exact sequences.} Let $0\to V_0\to V\to V_1\to0$ be an exact sequence of representations as in the statement, with $V$ crystalline. Tensoring over the field $\mathbf Q_\ell$ with $B$ keeps it exact, and taking invariants gives an exact sequence $0\to D_K(V_0)\to D_K(V)\to D_K(V_1)$. So $[K_0:\mathbf Q_\ell]\dim V=\dim D_K(V)\le\dim D_K(V_0)+\dim D_K(V_1)$, and by (1) both inequalities $\dim D_K(V_i)\le[K_0:\mathbf Q_\ell]\dim V_i$ are equalities: $V_0$ and $V_1$ are crystalline.
\item[(2)] \emph{Proof of (2).} If $f$ is injective, apply (b) to $0\to V\to V'\to V'/f(V)\to0$; if $f$ is surjective, to $0\to\ker f\to V\to V'\to0$.
\item[(3)] \emph{Proof of (3).} Let $K'_0$ be the largest subfield of $K'$ unramified over $\mathbf Q_\ell$. As $G_{K'}\subseteq G_K$, $D_K(V)\subseteq D_{K'}(V)$. As $V$ is crystalline, $D_K(V)$ has a basis over $K_0$ of $d$ elements; by (a) they are linearly independent over $B$, hence over its subfield $K'_0$. So $\dim_{K'_0}D_{K'}(V)\ge d$, and (1), for $K'$, gives equality.
\item[(4)] \emph{Proof of (4).} A basis of $L_2$ over $L_1$ identifies $V\otimes_{L_1}L_2$, as a representation over $\mathbf Q_\ell$, with the direct sum of $[L_2:L_1]$ copies of $V$, and $D_K$ commutes with finite direct sums. So both sides of the equality that defines crystalline representations are multiplied by $[L_2:L_1]$.
\item[(c)] \emph{A criterion.} $V$ is crystalline if and only if $D_K(V)$ spans $V\otimes_{\mathbf Q_\ell}B$ over $B$. Let $V$ be crystalline, so that $\dim_{K_0}D_K(V)=d$, and let $x_1,\dots,x_d$ be a basis of $D_K(V)$ over $K_0$; by (a) it is linearly independent over $B$. Write $x_j=\sum_kv_k\otimes a_{kj}$ for a basis $v_1,\dots,v_d$ of $V$ and $A=(a_{kj})\in M_d(B)$. Then $\det A\ne0$, since otherwise some nonzero $c\in B^d$ has $Ac=0$, as $B$ is a domain, and $\sum_jc_jx_j=0$. For $g\in G_K$ let $P(g)\in GL_d(\mathbf Q_\ell)$ be the matrix of $g$ on $V$; the invariance of the $x_j$ reads $P(g)g(A)=A$, so $g(\det A)=\det P(g)^{-1}\det A\in\mathbf Q_\ell\det A$. So $\det A$ is a unit of $B$, and the $x_j$ form a basis of $V\otimes B$ over $B$. Conversely, if $D_K(V)$ spans $V\otimes B$, the $B$-linear map $B\otimes_{K_0}D_K(V)\to V\otimes B$ is surjective; its source is free of rank $\dim_{K_0}D_K(V)$, finite by (1), and its target is free of rank $d$, so $d\le\dim_{K_0}D_K(V)$, and (1) gives equality.
\item[(d)] \emph{Bilinear maps.} Let $V'$ and $V''$ be crystalline representations as in the statement, and $\beta:V'\times V''\to V$ a $G_K$-equivariant $\mathbf Q_\ell$-bilinear map whose image spans $V$. Then $V$ is crystalline. Indeed, $\beta$ extends to the $B$-bilinear map $\beta_B:(V'\otimes B)\times(V''\otimes B)\to V\otimes B$ with $\beta_B(x\otimes b,y\otimes b')=\beta(x,y)\otimes bb'$, and $g\beta_B(u,u')=\beta_B(gu,gu')$; so $\beta_B$ maps $D_K(V')\times D_K(V'')$ into $D_K(V)$. By (c), $D_K(V')$ and $D_K(V'')$ span $V'\otimes B$ and $V''\otimes B$; so the $B$-span of $D_K(V)$ contains every $\beta_B(u,u')$, in particular every $\beta(x,y)\otimes b$, and these span $V\otimes B$. By (c), $V$ is crystalline. In particular, applied to $(x,y)\mapsto x\otimes y$, the tensor product of two crystalline representations is crystalline.
\item[(e)] \emph{Duals and linear maps.} If $V$ is crystalline, so is its dual $V^\vee=\operatorname{Hom}_{\mathbf Q_\ell}(V,\mathbf Q_\ell)$, on which $g$ acts by $f\mapsto f\circ g^{-1}$. Indeed, by (c) and its proof a basis $x_1,\dots,x_d$ of $D_K(V)$ over $K_0$ is a basis of $V\otimes B$ over $B$. Identify $V^\vee\otimes B$ with $\operatorname{Hom}_B(V\otimes B,B)$, on which $g$ acts by $f\mapsto g\circ f\circ g^{-1}$. The dual basis $(x_j^*)$ is invariant, as $g(x_j^*(g^{-1}x_k))=g(x_j^*(x_k))=x_j^*(x_k)$; so $D_K(V^\vee)$ spans $V^\vee\otimes B$, and $V^\vee$ is crystalline by (c). Hence, if $V'$ and $V''$ are crystalline, so is $\operatorname{Hom}_{\mathbf Q_\ell}(V',V'')\cong V'^\vee\otimes V''$, with $g$ acting by $f\mapsto g\circ f\circ g^{-1}$, by (d); and by (2) so is every representation with an injective equivariant map to it.
\item[(5)] \emph{Proof of (5).} The trivial representation $E$ is crystalline by (6). The map $E(\kappa^{-1})\to\operatorname{Hom}_{\mathbf Q_\ell}(E(\kappa),E)$, $a\mapsto(x\mapsto ax)$, is injective and $G_K$-equivariant, so $E(\kappa^{-1})$ is crystalline by (e). The $\mathbf Q_\ell$-bilinear map $E(\kappa)\times V\to V(\kappa)$, $(a,x)\mapsto ax$, is $G_K$-equivariant with image spanning $V(\kappa)$, so $V(\kappa)$ is crystalline by (d).
\item[(6)] \emph{Proof of (6).} Let $V'$ be the space $V$ with $g$ acting by $x\mapsto\chi(g)^{-w}gx$; it is trivial on $I_K$. As $G_K$ is compact, the coordinates of the vectors $ge$, for $g\in G_K$ and $e$ in a basis of $V'$, are bounded, so these vectors span a $G_K$-stable $\mathbf Z_\ell$-lattice. In a basis $v_1,\dots,v_d$ of $V$ that is a $\mathbf Z_\ell$-basis of this lattice, $g$ acts on $V'$ by a matrix $M(g)\in GL_d(\mathbf Z_\ell)$, and $M$ is a continuous homomorphism with $M(g)=1$ for $g\in I_K$. By Lemma~\ref{sha-automorphy-unramified-tangents-unit-crystalline-tilt}, Lemma~\ref{sha-automorphy-unramified-tangents-unit-crystalline-lang} applies to $R=\mathcal O_C^\flat$, $G=G_K$ and $M$, and gives $W\in M_d(A_{\mathrm{inf}})$ with $\det W$ a unit and $M(g)g(W)=W$ for every $g$. Let $W_B\in M_d(B)$ be its image under the map of Lemma~\ref{sha-automorphy-unramified-tangents-unit-crystalline-periods}(5), $t$ the unit of Lemma~\ref{sha-automorphy-unramified-tangents-unit-crystalline-periods}(6), and $X=t^{-w}W_B$. Then $\det X=t^{-wd}\det W_B$ is a unit of $B$, and the matrix $P(g)=\chi(g)^wM(g)$ of $g$ on $V$ satisfies $P(g)g(X)=\chi(g)^wM(g)\chi(g)^{-w}t^{-w}g(W_B)=X$. An element $\sum_kv_k\otimes y_k$ of $V\otimes B$, with $y\in B^d$, lies in $D_K(V)$ if and only if $P(g)g(y)=y$ for every $g\in G_K$. This holds for $y=Xc$ with $c\in K_0^d$. Conversely, if it holds, then $z=X^{-1}y$ satisfies $Xg(z)=P(g)g(X)g(z)=P(g)g(y)=y=Xz$, so $g(z)=z$ for every $g$, and $z\in K_0^d$ as $B^{G_K}=K_0$. So $c\mapsto\sum_kv_k\otimes(Xc)_k$ is a $K_0$-linear bijection from $K_0^d$ onto $D_K(V)$, injective as $X$ is invertible, and $\dim_{\mathbf Q_\ell}D_K(V)=[K_0:\mathbf Q_\ell]d$.
\item[(7)] \emph{Proof of (7).} The map $c$ is continuous, as $i$ has a linear left inverse, and $c(gh)=c(g)+gc(h)$, as $i(c(gh))=g(hs-s)+gs-s$. As in (b), tensoring with $B$ and taking invariants gives, with $D_K(\mathbf Q_\ell)=B^{G_K}=K_0$, an exact sequence $0\to D_K(V)\to D_K(V')\xrightarrow{q}K_0$, where $q$ is the restriction of $p\otimes1$. As $V$ is crystalline, $\dim_{\mathbf Q_\ell}D_K(V')=[K_0:\mathbf Q_\ell]\dim V+\dim_{\mathbf Q_\ell}q(D_K(V'))$; as $\dim V'=\dim V+1$, $V'$ is crystalline if and only if $q$ is surjective. The map $q$ is $K_0$-linear, $K_0$ acting through $B$, so it is surjective if and only if $q(y)=1$ for some $y\in D_K(V')$. Such a $y$ exists if and only if $c$ is a crystalline cocycle. If $c(g)\otimes1=gx-x$ for some $x\in V\otimes B$ and every $g$, then $y=s\otimes1-(i\otimes1)(x)$ has $(p\otimes1)(y)=1$ and $gy=s\otimes1+i(c(g))\otimes1-(i\otimes1)(x+c(g)\otimes1)=y$. Conversely, if $y\in D_K(V')$ and $q(y)=1$, then $s\otimes1-y$ lies in the kernel of $p\otimes1$, which is the image of $i\otimes1$, as in (b): $s\otimes1-y=(i\otimes1)(x)$ with $x\in V\otimes B$. Then $(i\otimes1)(gx-x)=g(s\otimes1)-s\otimes1=i(c(g))\otimes1$, and as $i\otimes1$ is injective, $c(g)\otimes1=gx-x$.
\end{enumerate}
% (0) The properties of B: the fields actB and invB of PeriodData and Core.Regular through RegularRing.lineUnit (Lemma sha-automorphy-unramified-tangents-unit-crystalline-periods); isDomain_cris (and RegularRing.isDomain_of_tied) for the domain; regularity is GlobalTheory.Per.FracInv.tied_frac_invariant and RegularRing.fracInv, two Lean proofs of the same fact: fracInv argues as the text, with the Lean's form TensorInj of (3) of Lemma -periods (c in Q_l^ur, then c in K_0 = K meet Q_l^ur); tied_frac_invariant gets c in Q_l^ur from FracAlg.mem_Kur_of_ratio instead, which shows that every nonzero element of Q_l(c) has absolute value in l^Z (norm_of_ratio, through the Gauss norm on A_cris) and applies UnramValue.le_Kur_of_norm.
% (a) is GlobalTheory.SimpStatement.main (ImportsExtra.lean, through FontaineBound.linearIndepOn_of_invariant), by strong induction on T. (1) is FontaineBound.finite_finrank_le_base, used in AdmSub.lean as bound, with regularity from GlobalTheory.Per.FracInv.tied_frac_invariant (from LocTied, tiedP) and isDomain_cris. (2) is isCrys_isDR_sub_quot (AdmSub.lean): of_exact (rTensor_exact, Module.Flat) for the exact sequences and sub_quot. (3) is isCrys_restrict (RestrictCrystallineDeRhamFiniteExt.lean: finite_le_0, indepB, dcris_le). (4) is isCrys_ext_iff (finrank_Dcris_ext through vecCoord), with isCrField and isCrIff for the passage between dimensions over the coefficient field and over Q_l (IsCrys versus IsCr, whose index (phi.range sup Iq).index is [K_0 : Q_l] by sup_Iq). (c) is GlobalTheory.RegularRing.isCrys_iff_span (span_eq_top_of_invariant for the determinant, with LineUnit) and, in AdmBilin.lean, the private admissible_iff_span_top. (d) is UpperTheory.V.Loc.admissible_of_bilinear_openSubgroup_galois (AdmBilin.lean: span_bilin and admissible_bilinear_image); the Lean proves the tensor product case separately as RegularRing.tensorCrys, by the same argument. (e) is UpperTheory.V.Loc.isCrys_isDR_internalHom (AdmHom.lean: exists_inj_dual_tprod for W into V'^dual (x) V'', and crys_hom with RegularRing.dualCrys, tensorCrys and subCrys). (5) is isCr_inv (isCrys_isDR_internalHom with b the multiplication of A) and isCr_twist (admissible_of_bilinear_openSubgroup_galois), private in MatrixExtendByZeroApply.lean, with isCr_trivial = crUnr at w = 0. (6) is crUnr through the private isCrys_of_inertia (IsCrysOfInertia.lean): exists_int_basis for the lattice, exists_matrixHom for M, exists_witt_matrix (Lemma sha-automorphy-unramified-tangents-unit-crystalline-lang) with eventually_actT_eq and norm_sub_lt_one_of_fix (Lemma sha-automorphy-unramified-tangents-unit-crystalline-tilt), ιA_zlW and ιA_act for W_B, the unit t of PD.t_unit, and GlobalTheory.T.isCrys_of_split_alt (SplitCrys.isCrys_of_split: finrank_coordInv and finrank_coordInv_of_split, which uses invB). (7) is the private line of MatrixExtendByZeroApply.lean: GlobalTheory.CrysExt.isCrys_iff_surjective (with dcris_trivial, from invB), lift_iff_surjective and coboundary_iff_lift (IsCrystallineIffMemH1f.lean); the continuity and the cocycle rule of c are hZ in crExt. The Lean's notion is IsCr phi rho for an open embedding phi of G into Gal(Qbar_l/Q_l), which openField identifies with G_K; CrOn phi rho U is IsCr on an open subgroup U.
\end{proof}

\begin{lemma}\label{sha-automorphy-unramified-tangents-unit-cup}
Let $K$, $G_K$, $\chi$ and the crystalline cocycles be as in Lemma~\ref{sha-automorphy-unramified-tangents-unit-duality}, let $I_K\subseteq G_K$ be the inertia subgroup, and let $E$ be a finite extension of $\mathbf Q_\ell$. Let $c:G_K\to E$ be a crystalline cocycle of $E(\chi)$, the space $E$ on which $s$ acts by multiplication by $\chi(s)$; that is, $c$ is continuous, $c(ss')=c(s)+\chi(s)c(s')$, and its class lies in $H^1_f(K,E(\chi))$. Then there are continuous maps $u,\beta:G_K\to E$ such that $u$ is a nonzero homomorphism to $(E,+)$ that vanishes on $I_K$, and $\beta(ss')=\beta(s)+\chi(s)\beta(s')-c(s)u(s')$ for all $s,s'\in G_K$.
\end{lemma}
\leanref{Development/AutomorphicArithmetic/AutomorphicVarieties/ModularityArgument/LinkedUnramifiedChains/CrystallineCardinalities/CupUnr.lean}{UpperTheory.V.Z.J.cupUnr}
\begin{proof}
\begin{enumerate}
\item \emph{An unramified homomorphism.} Let $\sigma\in\operatorname{Gal}(\bar{\mathbf Q}_\ell/\mathbf Q_\ell)$ lift the Frobenius of $\bar{\mathbf F}_\ell$, and let $d:\operatorname{Gal}(\bar{\mathbf Q}_\ell/\mathbf Q_\ell)\to\mathbf Z_\ell$ be the composite of the map onto $\operatorname{Gal}(\mathbf Q_\ell^{\mathrm{ur}}/\mathbf Q_\ell)\cong\hat{\mathbf Z}$, which sends $\sigma$ to $1$, with the projection $\hat{\mathbf Z}\to\mathbf Z_\ell$. It is a continuous homomorphism, trivial on inertia, with $d(\sigma^n)=n\ne0$ for $n\ge1$. Let $u(s)=d(s)\cdot1\in E$ for $s\in G_K$. As $G_K$ has finite index in $\operatorname{Gal}(\bar{\mathbf Q}_\ell/\mathbf Q_\ell)$, some $\sigma^n$ with $n\ge1$ lies in $G_K$, and $u(\sigma^n)\ne0$.
\item \emph{The pairing.} Let $b_1,\dots,b_e$ be a basis of $E$ over $\mathbf Q_\ell$, let $V=E(\chi)$, and let $V'=E$ with the trivial action of $G_K$. The representation $V$ is a sum of copies of $\mathbf Q_\ell(1)$, so it is de Rham, and the pairing $P(x,y)=\sum_kx_ky_k$, in the coordinates of the basis, is perfect with $P(\chi(s)x,y)=\chi(s)P(x,y)$.
\item \emph{Unramified cocycles.} For each $k$, $c_k(s)=d(s)b_k$ is a continuous homomorphism from $G_K$ to $V'$ that is trivial on $I_K$; it is a crystalline cocycle of $V'$. Indeed, $G_K$ acts on $V'\times\mathbf Q_\ell$ by $g(y,a)=(y+a\,c_k(g),a)$, trivially on $I_K$, so this representation and $V'$ are crystalline by Lemma~\ref{sha-automorphy-unramified-tangents-unit-crystalline}(6), with $w=0$; Lemma~\ref{sha-automorphy-unramified-tangents-unit-crystalline}(7), with $i(y)=(y,0)$, $p(y,a)=a$ and $s=(0,1)$, shows that $c_k$ is a crystalline cocycle.
\item \emph{The cup products.} As $c$ is a crystalline cocycle of $V$, Lemma~\ref{sha-automorphy-unramified-tangents-unit-duality}(1) gives for each $k$ a continuous $f_k:G_K\to\mathbf Q_\ell$ with $P(c(s),d(s')b_k)=\chi(s)f_k(s')-f_k(ss')+f_k(s)$, that is, $f_k(ss')=f_k(s)+\chi(s)f_k(s')-d(s')c_k(s)$, where $c_k(s)$ is the $k$-th coordinate of $c(s)$.
\item \emph{Conclusion.} Let $\beta=\sum_kf_kb_k$. It is continuous, and by (4), $\beta(ss')=\beta(s)+\chi(s)\beta(s')-d(s')c(s)=\beta(s)+\chi(s)\beta(s')-c(s)u(s')$.
\end{enumerate}
\end{proof}
% cupUnr (CupUnr.lean): openField identifies the open subgroup G with G_K; exists_deg (frobCyc, frobFree) for (1), with Subgroup.exists_pow_mem_of_index_ne_zero for the nonvanishing; the private gk_case for (2) to (5): charFil and isCrIff for the de Rham property of E(chi), the dual basis b.toDual with hdual for P, z1fUnr (ExistsIntBasis.lean) for (3), and (hP.bk ...).2.1, the field RationalDuality of BlochKatoLocal (Lemma sha-automorphy-unramified-tangents-unit-duality), for (4). Step (3) follows the direction of rat_case that the text uses: crExt with kappa = 1 and H the whole group (its extension criterion is Lemma sha-automorphy-unramified-tangents-unit-crystalline(7)) and crUnr at w = 0 for extRep (part (6)); the converse, through crWt0, is not needed.

\begin{lemma}\label{sha-automorphy-unramified-tangents-unit}
In the situation of Lemma~\ref{sha-automorphy-unramified-point}, suppose that $(\rho,\psi)$ has the prescribed inertial characters, that $r$ is potentially crystalline at the place $A_v$ above a place $v$ of $L^+$ above $\ell$, and that $\psi_{v,i}=\psi_{v,j}\chi_{\mathrm{cyc}}$ on all of $G_v$ for some $i,j$. Let $\rho_\xi$ be the linked point with values in $\mathcal O_1$ given by $\xi$, and $\varepsilon$ the cyclotomic character of the link. Then $j=i+1$, and with $i'=i$ there are $g\in GL_3(\mathcal O_1)$ and a homomorphism $t:G_v\to GL_3(\mathcal O_1)$ with $t(s)=g\rho_\xi(\mathrm{loc}_vs)g^{-1}$ upper triangular with diagonal $\psi_{v,1}(s),\psi_{v,2}(s),\psi_{v,3}(s)$ and $t(s)_{i'i'}=\varepsilon(s)t(s)_{i'+1,i'+1}$ for $s\in G_v$, together with continuous maps $u,\beta:G_v\to\mathcal O_1$, where $u$ is a nonzero homomorphism to $(\mathcal O_1,+)$ that vanishes on $I_v$, and $\beta(ss')=\beta(s)+\varepsilon(s)\beta(s')-t(s)_{i',i'+1}\bigl(t(s)^{-1}\bigr)_{i'+1,i'+1}u(s')$ for $s,s'\in G_v$.
\end{lemma}
\leanref{Development/AutomorphicArithmetic/AutomorphicVarieties/ModularityArgument/LinkedUnramifiedChains/CrystallineCardinalities/UnitOf.lean}{UpperTheory.V.Z.J.unitOf}
\leanref{Development/AutomorphicArithmetic/AutomorphicVarieties/ModularityArgument/LinkedUnramifiedChains/CrystallineCardinalities/UnitStep.lean}{UpperTheory.V.Z.J.UnitAt}
\leanref{Development/AutomorphicArithmetic/AutomorphicVarieties/ModularityArgument/LinkedUnramifiedChains/CrystallineCardinalities/UnitStep.lean}{UpperTheory.V.Z.J.UnitStep}
\begin{proof}
\emph{The indices.} By local class field theory, $\chi_{\mathrm{cyc}}$ maps $I_v$ onto an open subgroup of $\mathbf Z_\ell^\times$, so some $g_0\in I_v$ has $\chi_{\mathrm{cyc}}(g_0)$ of infinite order. With the prescribed inertial characters, $\psi_{v,i}=\Psi^{w_{3,i}}$ on $I_v$, and $\Psi=\chi_{\mathrm{cyc}}^k\psi$ with $\psi$ of finite order and $k\ge1$ (Definition~\ref{sha-ordinary-deformation}). So $\psi_{v,i}(g_0)=\psi_{v,j}(g_0)\chi_{\mathrm{cyc}}(g_0)$ gives $\chi_{\mathrm{cyc}}(g_0)^{k(w_{3,i}-w_{3,j})-1}=\psi(g_0)^{-(w_{3,i}-w_{3,j})}$, which has finite order. Hence $k(w_{3,i}-w_{3,j})=1$, so $k=1$ and $w_{3,i}-w_{3,j}=1$, that is, $j=i+1$, as $w_3=(1,0,-1)$.

\emph{The local group.} Let $E$ be the fraction field of $\mathcal O_1$. The entries of $\rho_\xi$ are continuous: by condition (d) of Definition~\ref{sha-descent-link}(2) the entries of $\rho^{\mathrm u}$ are continuous for the $\mathfrak m_{A_3}$-adic topology, and $\xi$ maps $\mathfrak m_{A_3}$ into the maximal ideal of $\mathcal O_1$ (Lemma~\ref{sha-automorphy-unramified-near-tangent-points}); so $r$ is continuous. As $r$ is potentially crystalline at $A_v$, there are an embedding of $\bar{\mathbf Q}$ in $\bar{\mathbf Q}_\ell$ inducing the place $A_v$, through which $G_v$ becomes the group $G_K=\operatorname{Gal}(\bar{\mathbf Q}_\ell/K)$ of a finite extension $K$ of $\mathbf Q_\ell$, with inertia subgroup $I_v$ and cyclotomic character $\varepsilon\circ\mathrm{loc}_v$, and a finite extension $K'$ of $K$ such that a model of $r$ over a finite extension $L_1$ of $\mathbf Q_\ell$ is crystalline on $G_{K'}$, in the sense of Lemma~\ref{sha-automorphy-unramified-tangents-unit-crystalline}. Let $U$ be the normal core of the open subgroup $G_{K'}$ of $G_v$; it is the Galois group of a finite extension $K''$ of $K'$. By Lemma~\ref{sha-automorphy-unramified-tangents-unit-crystalline}(3), the model is crystalline on $U$. The map $s\mapsto\rho_\xi(\mathrm{loc}_vs)$, with values in $GL_3(E)$, has the same matrices in $\bar{\mathbf Q}_\ell$ as $r$; so after extending scalars to a finite extension $L_2$ containing $L_1$ and $E$ the two become the same representation, and by Lemma~\ref{sha-automorphy-unramified-tangents-unit-crystalline}(4), applied twice, $s\mapsto\rho_\xi(\mathrm{loc}_vs)$ is crystalline on $U$.

\emph{A triangular form.} Lemma~\ref{sha-automorphy-unramified-near-completion-flag}, applied to the local ring $\mathcal O_1$ and the restriction of $\xi$ to $A_3^\circ$, gives $g\in GL_3(\mathcal O_1)$ such that $t(s)=g\rho_\xi(\mathrm{loc}_vs)g^{-1}$ is upper triangular with diagonal entries $\psi_{v,1}(s),\psi_{v,2}(s),\psi_{v,3}(s)$ for $s\in G_v$, since $\xi$ carries $\psi^{\mathrm u}_{v,k}$ to $\psi_{v,k}$; and $t$ is a homomorphism. As $j=i'+1$, $t(s)_{i'i'}=\psi_{v,i'}(s)=\varepsilon(s)\psi_{v,i'+1}(s)=\varepsilon(s)t(s)_{i'+1,i'+1}$. Being conjugate to $s\mapsto\rho_\xi(\mathrm{loc}_vs)$ by $g$, $t$ is crystalline on $U$, by Lemma~\ref{sha-automorphy-unramified-tangents-unit-crystalline}(2) applied to the isomorphism $x\mapsto gx$.

\emph{A crystalline class.} Let $\alpha(s)=t(s)_{i'i'}$, $\kappa(s)=t(s)_{i'+1,i'+1}$ and $b(s)=t(s)_{i',i'+1}$. Since $t$ is upper triangular, $\alpha$ and $\kappa$ are characters, $\alpha=\varepsilon\kappa$, and $b(ss')=\alpha(s)b(s')+b(s)\kappa(s')$. The representation of $G_v$ on $E^2$ by the matrices $\begin{pmatrix}\alpha(s)&b(s)\\0&\kappa(s)\end{pmatrix}$ is the restriction of $t$ to the span of the first two basis vectors if $i'=1$, and the quotient of $t$ by the span of the first basis vector if $i'=2$. By Lemma~\ref{sha-automorphy-unramified-tangents-unit-crystalline}(2), this representation is crystalline on $U$, and so are its subrepresentation $E(\alpha)$ on the first basis vector and its quotient $E(\kappa)$ by it.
Its twist by $\kappa^{-1}$, on which $s$ acts by $(x_1,x_2)\mapsto(\varepsilon(s)x_1+c(s)x_2,x_2)$ with $c(s)=b(s)\kappa(s)^{-1}=t(s)_{i',i'+1}\bigl(t(s)^{-1}\bigr)_{i'+1,i'+1}$, is an extension of the trivial representation $E$ by $E(\varepsilon)=E(\alpha)(\kappa^{-1})$; $c$ is continuous and $c(ss')=c(s)+\varepsilon(s)c(s')$. By Lemma~\ref{sha-automorphy-unramified-tangents-unit-crystalline}(5), applied on $U$ to $\kappa$ and then to the character $\kappa^{-1}$, this twist and $E(\varepsilon)$ are crystalline on $U$. The $\mathbf Q_\ell$-subspace of the twist spanned by $E\times0$ and $s_0=(0,1)$ is stable, and $ss_0-s_0=(c(s),0)$; by Lemma~\ref{sha-automorphy-unramified-tangents-unit-crystalline}(2) it is crystalline on $U$, and by Lemma~\ref{sha-automorphy-unramified-tangents-unit-crystalline}(7), with $i(x)=(x,0)$ and $p$ the second coordinate, the restriction of $c$ to $U$ is a crystalline cocycle: there is $x$ with $c(s)\otimes1=sx-x$ for $s\in U$. For $g\in G_v$ the element $F(g)=gx-c(g)\otimes1$ depends only on $gU$, since $c(gh)=c(g)+g\,c(h)$ for $h\in U$; the average $x'$ of $F$ over $G_v/U$ satisfies $sx'=x'+c(s)\otimes1$ for every $s\in G_v$. So $c$ is a crystalline cocycle of $E(\varepsilon)$ for $G_v=G_K$.

\emph{The unit step.} Lemma~\ref{sha-automorphy-unramified-tangents-unit-cup}, applied to $K$, $E$ and $c$, gives continuous $u_0,\beta_0:G_v\to E$, with $u_0$ a nonzero homomorphism that vanishes on $I_v$ and $\beta_0(ss')=\beta_0(s)+\varepsilon(s)\beta_0(s')-c(s)u_0(s')$.

\emph{Integral values.} The continuous maps $u_0$ and $\beta_0$ on the compact group $G_v$ are bounded, and the elements of $E$ of absolute value at most $1$ form $\mathcal O_1$; so for some $m\ge0$, $u=\ell^mu_0$ and $\beta=\ell^m\beta_0$ take values in $\mathcal O_1$. They are continuous, $u$ is a nonzero homomorphism that vanishes on $I_v$, and $\beta(ss')=\beta(s)+\varepsilon(s)\beta(s')-c(s)u(s')$, as this relation is linear in the pair $(u_0,\beta_0)$. With $g$ and $t$ these are the data of the statement.
% unitOf (UnitOf.lean, 305 lines): hcont and hrc (Hi.Up.pt_cont) for the continuity; crOfPot (CrOfPot.lean: crys_res, crys_model, isCr_of_crys, the normal core of GK K') and locEmb (Loc.decEmb, galEmbFam) for the local group; cycInertia for g_0 and idxStep for the indices; flagStep for the rest: Fr.flag (Lemma sha-automorphy-unramified-near-completion-flag) for g and t, crConj for the crystalline conjugate, fieldOf for E (the Q_l-span of O_1, with l^m x in O_1 for x in E), and stepCore: stepF (TwoApply.lean: the block two, crOn_map with isCrys_isDR_sub_quot for the subrepresentation or quotient, and core with crExt, which twists by kappa^-1, descends from U to G_v by z1fDesc, Z1fDescAvg.lean, and identifies crystalline extensions of the trivial representation with crystalline cocycles), cupUnr (Lemma sha-automorphy-unramified-tangents-unit-cup) and clearDen (AdicOfCont.lean: clear_one for the bound, adic_of_cont for the continuity into O_1). Of the published local theorems (LocalPublishedTheorems: LocalFiniteness, LocalTateDuality, LocalEulerChar, their K-versions, and BlochKatoLocal), the unit step uses only the rational duality of Bloch-Kato (Lemma sha-automorphy-unramified-tangents-unit-duality). The facts on crystalline representations are Lemma sha-automorphy-unramified-tangents-unit-crystalline: crOfPot uses crys_res (part (3)) and crys_model (part (4), twice, through the compositum supExt), then isCr_of_crys and cr_pull; flagStep uses crConj (part (2) for the isomorphism g); stepF uses crOn_map (part (2)) for the block two, and its core uses crExt: isCr_inv and isCr_twist (part (5)), the private triv (part (2) for the injection j of the Q_l-extension ext0, then line, part (7)) and z1fDesc for the descent from U.
\end{proof}

\begin{definition}\label{sha-automorphy-unramified-tangents-count-counts}
Let a link be given, with group $\Gamma$, homomorphisms $\delta_\Gamma$ and $\varepsilon_\Gamma$, set $V$ of places $v$ of $L^+$ above $\ell$, groups $G_v\supseteq I_v$, maps $\mathrm{loc}_v$ and local degrees $d_v$ (Definition~\ref{sha-descent-link}). For a profinite group $D$ acting on a finite abelian group $N$ by group automorphisms, let $Z^1(D,N)$ be the set of locally constant maps $z:D\to N$ with $z(st)=z(s)+sz(t)$ for all $s,t\in D$. For $k\ge1$, a homomorphism $\epsilon:D\to\mathbf Z_\ell^\times$ and $N$ killed by $\ell^k$, let $\operatorname{Hom}_\epsilon(N,\mathbf Z/\ell^k)$ be the set of homomorphisms $f:N\to\mathbf Z/\ell^k$ with $f(sx)=\epsilon(s)f(x)$ for all $s\in D$ and $x\in N$.
\begin{enumerate}
\item The link satisfies the \emph{local counts} if for every $v\in V$, every $k\ge1$ and every finite abelian group $N$ killed by $\ell^k$ with an action of $G_v$ whose kernel is open, $Z^1(G_v,N)$ is finite and $|Z^1(G_v,N)|=|N|^{1+d_v}\,|\operatorname{Hom}_{\varepsilon_\Gamma\circ\mathrm{loc}_v}(N,\mathbf Z/\ell^k)|$.
\item Let $(c'_w)_{w\in\Omega}$ be a family of elements of $\Gamma$ indexed by a finite set $\Omega$. The link satisfies the \emph{global count} for $(c'_w)$ if the following holds for every $k\ge1$ and every finite abelian group $N$ killed by $\ell^k$ with an action of $\Gamma$ whose kernel is open. Let $Z_v\subseteq Z^1(G_v,N)$, for $v\in V$ and $G_v$ acting through $\mathrm{loc}_v$, be subgroups that contain the maps $h\mapsto\mathrm{loc}_v(h)Y-Y$ for $Y\in N$, and let $Z_L$ be the set of $z\in Z^1(\Gamma,N)$ with $z\circ\mathrm{loc}_v\in Z_v$ for every $v$. Then $Z_L$ is finite and
\[
|N|\prod_{v\in V}|Z_v|\le|Z_L|\,|\operatorname{Hom}_{\varepsilon_\Gamma}(N,\mathbf Z/\ell^k)|\prod_{w\in\Omega}|N^{c'_w}|\,|N|^{|V|},
\]
where $N^{c'_w}$ is the set of elements of $N$ fixed by $c'_w$.
\end{enumerate}
\end{definition}
\leanref{Development/AutomorphicArithmetic/AutomorphicVarieties/ModularityArgument/LinkedUnramifiedChains/CrystallineCardinalities/UnitStep.lean}{UpperTheory.V.Z.J.Zc}
\leanref{Development/AutomorphicArithmetic/AutomorphicVarieties/ModularityArgument/LinkedUnramifiedChains/CrystallineCardinalities/UnitStep.lean}{UpperTheory.V.Z.J.HomMu}
\leanref{Development/AutomorphicArithmetic/AutomorphicVarieties/ModularityArgument/LinkedUnramifiedChains/CrystallineCardinalities/UnitStep.lean}{UpperTheory.V.Z.J.LE}
\leanref{Development/AutomorphicArithmetic/AutomorphicVarieties/ModularityArgument/LinkedUnramifiedChains/CrystallineCardinalities/UnitStep.lean}{UpperTheory.V.Z.J.LocE}
\leanref{Development/AutomorphicArithmetic/AutomorphicVarieties/ModularityArgument/LinkedUnramifiedChains/CrystallineCardinalities/UnitStep.lean}{UpperTheory.V.Z.J.ZL}
\leanref{Development/AutomorphicArithmetic/AutomorphicVarieties/ModularityArgument/LinkedUnramifiedChains/CrystallineCardinalities/UnitStep.lean}{UpperTheory.V.Z.J.GWL}

\begin{lemma}\label{sha-automorphy-unramified-tangents-count-inputs-unramified}
Let $F\subset\bar{\mathbf Q}$ be a number field, $v$ a finite place of $F$, $A$ a place of $\bar{\mathbf Q}$ above $v$, and $D\supseteq I$ the decomposition and inertia groups of $A$ in $\operatorname{Gal}(\bar{\mathbf Q}/F)$. Let $N$ be a finite abelian group with an action of $D$ whose kernel is open, and let $Z$ be the set of $z\in Z^1(D,N)$ (Definition~\ref{sha-automorphy-unramified-tangents-count-counts}) for which there is $Y\in N$ with $z(s)=sY-Y$ for every $s\in I$. Then $Z$ is a finite subgroup of $Z^1(D,N)$ that contains the maps $s\mapsto sY-Y$ for $Y\in N$, and $|N|\le|Z|$.
\end{lemma}
\leanref{Development/AutomorphicArithmetic/AutomorphicVarieties/ModularityArgument/LinkedUnramifiedChains/CrystallineCardinalities/UnrLow.lean}{UpperTheory.V.Z.J.unrLow}
\begin{proof}
By Lemma~\ref{sha-descent-link-unramified}(1) there is $\varphi\in D$ such that $D$ is the closure of the subgroup generated by $I$ and $\varphi$, and for every $n\ge1$ a continuous homomorphism $d_n:D\to\mathbf Z/n$ that vanishes on $I$ with $d_n(\varphi)=1$. The group $I$ is normal in $D$, and the maps $s\mapsto sY$ are locally constant, as the kernel of the action is open.
\begin{enumerate}
\item[(a)] \emph{Vanishing.} A cocycle $z\in Z^1(D,N)$ that vanishes on $I$ and at $\varphi$ is $0$. Indeed, the set of $s$ with $z(s)=0$ is a subgroup, as $z(st)=z(s)+sz(t)$ and $z(s^{-1})=-s^{-1}z(s)$; it is closed, as $z$ is locally constant; and it contains $I$ and $\varphi$.
\item[(b)] \emph{A finite group.} If $z,z'\in Z$, with elements $Y$ and $Y'$, then $z-z'$ is a cocycle equal to $s\mapsto s(Y-Y')-(Y-Y')$ on $I$; and the map $s\mapsto sY-Y$ is a locally constant cocycle that lies in $Z$. So $Z$ is a subgroup that contains these maps. Choose $Y_z$ for each $z\in Z$. The map $z\mapsto(Y_z,z(\varphi))$ to $N\times N$ is injective: if $z$ and $z'$ have the same image, then $z-z'$ vanishes on $I$ and at $\varphi$, so $z=z'$ by (a). So $Z$ is finite.
\item[(c)] \emph{Cocycles from invariants.} Let $N^I$ be the set of elements of $N$ fixed by $I$; it is stable under $D$, as $I$ is normal. Let $x\in N^I$. There is $z\in Z^1(D,N)$ that vanishes on $I$ with $z(\varphi)=x$. Indeed, let $T$ be the permutation $y\mapsto x+\varphi y$ of $N$, $n$ its order, and $\tau(s)=T^{d_n(s)}$, a locally constant homomorphism from $D$ to the permutations of $N$, trivial on $I$, with $\tau(\varphi)=T$. Let $S$ be the set of $s\in D$ for which there is $a\in N^I$ with $\tau(s)y=a+sy$ for every $y\in N^I$. If $s,t\in S$, with elements $a$ and $b$, then $\tau(st)y=\tau(s)(b+ty)=(a+sb)+sty$ for $y\in N^I$, with $a+sb\in N^I$; and $\tau(s^{-1})y=-s^{-1}a+s^{-1}y$. So $S$ is a subgroup. It is closed, as membership depends only on $\tau(s)$ and on the action of $s$, which are locally constant in $s$. It contains $I$, with $a=0$, and $\varphi$, with $a=x$. So $S=D$. Let $z(s)=\tau(s)(0)$. For $s,t\in D$, $\tau(t)(0)=b$ with $b$ the element of $t$, and $\tau(s)(b)=a+sb=\tau(s)(0)+s\tau(t)(0)$; so $z(st)=z(s)+sz(t)$. And $z$ is locally constant, vanishes on $I$ and has $z(\varphi)=T(0)=x$.
\item[(d)] \emph{The count.} Let $r$ send $z\in Z$ to its restriction to $I$, and $\pi_0$ send $Y\in N$ to the map $s\mapsto sY-Y$ on $I$. By the definition of $Z$ and (b), $r$ and $\pi_0$ have the same image. The kernel of $\pi_0$ is $N^I$, and by (c) there is an injective map from $N^I$ to the kernel of $r$, $x\mapsto z$ with $z(\varphi)=x$. So $|Z|=|\ker r|\,|\operatorname{im}r|\ge|N^I|\,|\operatorname{im}\pi_0|=|N|$.
\end{enumerate}
\end{proof}
% unrLow (UnrLow.lean, 144 lines): GC.decZhat (Lemma sha-descent-link-unramified(1)) for phi and d_n, then the private ul_abs: (a) is ul_zero (UpperTheory.Reformulated.B.Rb1.ul_zero) with ul_top; (b) is the subgroup Z (AddSubgroup.ofSub) with hcob and hfin; (c) is ul_ex, with T, tau (zmodMulEquivOfGenerator), the subgroup S, stab, and ul_lc for the local constancy; (d) is r, pi, hr and hk.

\begin{lemma}\label{sha-automorphy-unramified-tangents-count-inputs-level}
Let $F\subset\bar{\mathbf Q}$ be a number field, $G_F=\operatorname{Gal}(\bar{\mathbf Q}/F)$, $S$ a finite set of finite places of $F$, $H$ an open normal subgroup of $G_F$, and $p:G_F\to\Gamma$ a continuous surjective homomorphism onto a profinite group, whose kernel is the closed normal subgroup generated by the groups $I_B\cap H$ for the places $B$ of $\bar{\mathbf Q}$ above the finite places of $F$ outside $S$. Then there is a finite set $S'\supseteq S$ of finite places of $F$ such that $p$ factors through a continuous surjection $p':G_{S'}\to\Gamma$, with $G_{S'}$ as in Lemma~\ref{sha-descent-link-global}. For every choice of places $A_q$ of $\bar{\mathbf Q}$ above the places $q\in S'\setminus S$, the kernel of $p'$ is the closed normal subgroup generated by the images in $G_{S'}$ of the groups $I_{A_q}\cap H$.
\end{lemma}
\leanref{Development/AutomorphicArithmetic/AutomorphicVarieties/GaloisCochains/LevelS.lean}{UpperTheory.V.GC.levelS}
\begin{proof}
The fixed field $E$ of $H$ is a finite extension of $F$, so there is a finite set $T$ of finite places of $F$ such that $E/F$ is unramified outside $T$ \cite[Chapter~III, \S2]{Neu99}; then $I_B\subseteq H$ for every place $B$ above a place outside $T$. Let $S'=S\cup T$. For $B$ above a place outside $S'$, $I_B=I_B\cap H$ lies in $\ker p$. So the kernel of $G_F\to G_{S'}$, the closed normal subgroup generated by these $I_B$, lies in $\ker p$, and $p$ factors through a continuous surjection $p':G_{S'}\to\Gamma$.

Let $N'$ be the closed normal subgroup of $G_{S'}$ generated by the images of the $I_{A_q}\cap H$, $q\in S'\setminus S$. They lie in the image of $\ker p$, so $N'\subseteq\ker p'$. Conversely, the inverse image of $N'$ in $G_F$ is a closed normal subgroup that contains every group $I_B\cap H$ with $B$ above a place outside $S$. If $B$ lies above $q\in S'\setminus S$, then $B=uA_q$ for some $u\in G_F$, and $I_B\cap H=u(I_{A_q}\cap H)u^{-1}$, as $H$ is normal. If $B$ lies above a place outside $S'$, then $I_B$ maps to $1$ in $G_{S'}$. So $\ker p$ lies in the inverse image of $N'$, and $\ker p'\subseteq N'$.
\end{proof}
% levelS (LevelS.lean): the private ine_le_open (Base.ramFinite for the fixed field) for T, hgen and hNS for the factorization through QuotientGroup.lift, and the two inclusions, the first with exists_smul_eq_of_isPlaceOver and conj_mem_Ine'.

\begin{lemma}\label{sha-automorphy-unramified-tangents-count-inputs-global}
Let $\ell$ be an odd prime, $F\subset\bar{\mathbf Q}$ a totally real number field, $S'$ a finite set of finite places of $F$ that contains the places above $\ell$, and $G_{S'}$, the groups $D_q$ for $q\in S'$, and complex conjugations $c_w$ at the infinite places $w$ of $F$, as in Lemma~\ref{sha-descent-link-greenberg-wiles}; let $\epsilon:G_{S'}\to\mathbf Z_\ell^\times$ be the homomorphism through which $\chi_{\mathrm{cyc}}$ factors. Let $k\ge1$, let $N$ be a finite abelian group killed by $\ell^k$ with an action of $G_{S'}$ whose kernel is open, and for $q\in S'$ let $Z_q\subseteq Z^1(D_q,N)$, for $D_q$ acting through $G_{S'}$, be a subgroup that contains the maps $s\mapsto sY-Y$ for $Y\in N$ (Definition~\ref{sha-automorphy-unramified-tangents-count-counts}). Let $Z_L$ be the set of $z\in Z^1(G_{S'},N)$ whose restriction to $D_q$ lies in $Z_q$ for every $q$. Then $Z^1(G_{S'},N)$ is finite and
\[
|N|\prod_{q\in S'}|Z_q|\le|Z_L|\,|\operatorname{Hom}_\epsilon(N,\mathbf Z/\ell^k)|\prod_w|N^{c_w}|\,|N|^{|S'|},
\]
where $N^{c_w}$ is the set of elements of $N$ fixed by $c_w$.
\end{lemma}
\leanref{Development/AutomorphicArithmetic/AutomorphicVarieties/ModularityArgument/LinkedUnramifiedChains/CrystallineCardinalities/FiniteCocyclesOfFiniteH1.lean}{bare}
\begin{proof}
View $N$ as a $\mathbf Z/\ell^k$-module, with the discrete topology; the action is continuous, and the locally constant cocycles are the continuous cocycles. Let $\chi_k$ be the reduction of $\epsilon$ modulo $\ell^k$, and write $H^i$ and $B^1$ for continuous cohomology and coboundaries.
\begin{enumerate}
\item[(a)] \emph{Finiteness.} By Lemma~\ref{sha-descent-link-greenberg-wiles} with $r=k$, $H^1(G_{S'},N)$ and $H^1(G_{S'},N^*(1))$ are finite. As $N$ is finite, so is $B^1(G_{S'},N)$, and hence $Z^1(G_{S'},N)$.
\item[(b)] \emph{Classes.} For a profinite group $D'$ acting on $N$, the map $Y\mapsto(s\mapsto sY-Y)$ has kernel $H^0(D',N)$ and image $B^1(D',N)$, so $|N|=|H^0(D',N)|\,|B^1(D',N)|$. For $q\in S'$ let $L_q$ be the image of $Z_q$ in $H^1(D_q,N)$, a submodule. As $Z_q$ contains $B^1(D_q,N)$, it is the set of cocycles whose class lies in $L_q$, so $|Z_q|=|B^1(D_q,N)|\,|L_q|$. Likewise $Z_L$ is the set of cocycles whose class lies in $\operatorname{Sel}_L(N)$, so $|Z_L|=|B^1(G_{S'},N)|\,|\operatorname{Sel}_L(N)|$.
\item[(c)] \emph{The dual term.} An element $\lambda$ of $H^0(G_{S'},N^*(1))$ satisfies $\lambda(x)=\chi_k(s)\lambda(s^{-1}x)$, that is, $\lambda(sx)=\chi_k(s)\lambda(x)$, for all $s$ and $x$. So $\lambda\in\operatorname{Hom}_\epsilon(N,\mathbf Z/\ell^k)$, and $|H^0(G_{S'},N^*(1))|\le|\operatorname{Hom}_\epsilon(N,\mathbf Z/\ell^k)|$.
\item[(d)] \emph{Greenberg-Wiles.} As $F$ is totally real, it has no complex places and every infinite place is real. Lemma~\ref{sha-descent-link-greenberg-wiles} with $r=k$ and the $L_q$, with $|\operatorname{Sel}_{L^\perp}(N^*(1))|\ge1$ and (c), gives
\[|H^0(G_{S'},N)|\prod_q|L_q|\le|\operatorname{Sel}_L(N)|\,|\operatorname{Hom}_\epsilon(N,\mathbf Z/\ell^k)|\prod_q|H^0(D_q,N)|\prod_w|N^{c_w}|.\]
\item[(e)] \emph{Conclusion.} Multiplying by $|B^1(G_{S'},N)|\prod_q|B^1(D_q,N)|$ and using (b) gives the inequality of the statement.
\end{enumerate}
\end{proof}
% bare (FiniteCocyclesOfFiniteH1.lean, private): zrep for the Z/l^k-module, PoitouTate.exists_lift_cycMod for mu, the field gf of HcPublished (GlobalFiniteness) for (a), harch for the archimedean factor, hdual for (c), and the field gw (GreenbergWiles) through the private core, which uses ex_cls, cls and card_cls for (b), finite_z1_of_finite_h1 for (a), and card_coboundaries_mul_card_h0.

\begin{lemma}\label{sha-automorphy-unramified-tangents-count-inputs}
In the situation of Lemma~\ref{sha-automorphy-unramified-point}, the link satisfies the local counts, and there is a family $(c'_w)$ of elements of $\Gamma$, indexed by the infinite places $w$ of $L^+$, with $c_w'^2=1$ and $\delta_\Gamma(c'_w)=-1$ for every $w$, for which the link satisfies the global count (Definition~\ref{sha-automorphy-unramified-tangents-count-counts}).
\end{lemma}
\leanref{Development/AutomorphicArithmetic/AutomorphicVarieties/ModularityArgument/LinkedUnramifiedChains/CrystallineCardinalities/LinkOf.lean}{UpperTheory.V.Z.J.linkOf}
\leanref{Development/AutomorphicArithmetic/AutomorphicVarieties/ModularityArgument/LinkedUnramifiedChains/CrystallineCardinalities/FiniteCocyclesOfFiniteH1.lean}{UpperTheory.V.Z.J.gwlOf}
\leanref{Development/AutomorphicArithmetic/AutomorphicVarieties/ModularityArgument/LinkedUnramifiedChains/CrystallineCardinalities/LeOf2.lean}{UpperTheory.V.Z.J.leOf2}
\begin{proof}
Let $G_{L^+}=\operatorname{Gal}(\bar{\mathbf Q}/L^+)$, let $p:G_{L^+}\to\Gamma$ be the map of the link, and $V$ the set of places of $L^+$ above $\ell$ (Definition~\ref{sha-descent-link}(1)).
\begin{enumerate}
\item[(a)] \emph{The Galois data.} Let $H$ be the set of $s\in G_{L^+}$ with $\eta(s)=\bar\Psi(s)=1$, that is, with $\delta_\Gamma(p(s))=\bar\Psi_\Gamma(p(s))=1$. It is a normal subgroup of $G_{L^+}$; it is open, as $p$ is continuous and $\Gamma_0$ and $\ker\bar\Psi_\Gamma$ are open (residual condition (4) of Definition~\ref{sha-descent-residual}); and its index is prime to $\ell$, as $G_{L^+}/H$ embeds in $\{\pm1\}\times\mathbf F^\times$, of order $2(|\mathbf F|-1)$, with $\ell$ odd. By condition (i) of Definition~\ref{sha-descent-link}(1), the kernel of $p$ is the closed normal subgroup generated by the groups $I_B\cap H$ for the places $B$ of $\bar{\mathbf Q}$ above the places of $L^+$ outside $V$. By its condition (iv), for $v\in V$ the group $G_v$ is the decomposition group in $G_{L^+}$ of the place $A_v=\gamma_v\mathfrak B$ above $v$, with inertia group $I_v$, $\mathrm{loc}_v$ is the restriction of $p$, and $d_v=[L^+_v:\mathbf Q_\ell]$; by (ii), $\varepsilon_\Gamma\circ p=\chi_{\mathrm{cyc}}$. For each infinite place $w$ of the totally real field $L^+$, choose an embedding $j_w:\bar{\mathbf Q}\to\mathbf C$ inducing $w$, and let $\tau_w\in G_{L^+}$ be the element with $j_w\circ\tau_w=\overline{j_w}$, a complex conjugation at $w$ with $\tau_w^2=1$. Let $c'_w=p(\tau_w)$. Then $c_w'^2=1$, and $\delta_\Gamma(c'_w)=\eta(\tau_w)=-1$: otherwise $\tau_w$ would lie in $\operatorname{Gal}(\bar{\mathbf Q}/L)=G_{L^+}\cap\ker\eta$, and $j_w$ would map $L$ into $\mathbf R$, which is impossible for the CM field $L$.
\item[(b)] \emph{The local counts.} Let $v\in V$, $k\ge1$, and $N$ a finite abelian group killed by $\ell^k$ with an action of $G_v$ whose kernel is open. Then $N$ is a finite $\mathbf Z/\ell^k$-module, the action is continuous for the discrete topology, and $Z^1(G_v,N)$ is the group of continuous cocycles. On $G_v$, $\varepsilon_\Gamma\circ\mathrm{loc}_v=\chi_{\mathrm{cyc}}$, with reduction $\chi_k$ modulo $\ell^k$. By Lemma~\ref{sha-descent-link-local-tate}, for $L^+$, $v$, $A_v$ and $r=k$, the groups $H^i(G_v,N)$, $i\le2$, are finite and $|H^1(G_v,N)|=|H^0(G_v,N)|\,|H^2(G_v,N)|\,|N|^{d_v}$. By its part (3) with $i=2$, $H^2(G_v,N)$ is in bijection with the homomorphisms from $H^0(G_v,N^*(1))$ to $H^2(G_v,\mathbf Z/\ell^k(\chi_k))\cong\mathbf Z/\ell^k$, and there are $|H^0(G_v,N^*(1))|$ of these, as $H^0(G_v,N^*(1))$ is killed by $\ell^k$. The elements of $H^0(G_v,N^*(1))$ are the homomorphisms $\lambda:N\to\mathbf Z/\ell^k$ with $\lambda(sx)=\chi_k(s)\lambda(x)$, that is, the elements of $\operatorname{Hom}_{\varepsilon_\Gamma\circ\mathrm{loc}_v}(N,\mathbf Z/\ell^k)$. Finally, $Z^1(G_v,N)$ is finite, with $|Z^1(G_v,N)|=|B^1(G_v,N)|\,|H^1(G_v,N)|$, and $|N|=|H^0(G_v,N)|\,|B^1(G_v,N)|$, so
\[|Z^1(G_v,N)|\,|H^0(G_v,N)|=|H^1(G_v,N)|\,|N|=|H^0(G_v,N)|\,|\operatorname{Hom}_{\varepsilon_\Gamma\circ\mathrm{loc}_v}(N,\mathbf Z/\ell^k)|\,|N|^{1+d_v},\]
which is the local count at $v$.
\item[(c)] \emph{The level.} Lemma~\ref{sha-automorphy-unramified-tangents-count-inputs-level}, applied to $L^+$, $V$, $H$ and $p$, gives a finite set $S'\supseteq V$ of places of $L^+$ and a continuous surjection $p':G_{S'}\to\Gamma$ through which $p$ factors. For $q\in S'\setminus V$ choose a place $A_q$ above $q$, and let $D_q\supseteq I_q$ be its decomposition and inertia groups in $G_{L^+}$; for $v\in V$ let $D_v=G_v$. Then $\ker p'$ is the closed normal subgroup generated by the images of the $I_q\cap H$, $q\in S'\setminus V$.
\item[(d)] \emph{The global count.} Let $k\ge1$, $N$ a finite abelian group killed by $\ell^k$ with an action of $\Gamma$ whose kernel is open, and for $v\in V$ let $Z_v\subseteq Z^1(G_v,N)$ be subgroups that contain the maps $h\mapsto\mathrm{loc}_v(h)Y-Y$; let $Z_L$ be as in Definition~\ref{sha-automorphy-unramified-tangents-count-counts}(2). Let $G_{S'}$ act on $N$ through $p'$, with open kernel. For $q\in S'\setminus V$ let $Z_q\subseteq Z^1(D_q,N)$ be the group of Lemma~\ref{sha-automorphy-unramified-tangents-count-inputs-unramified}, with $|N|\le|Z_q|$. Lemma~\ref{sha-automorphy-unramified-tangents-count-inputs-global}, applied to $L^+$, $S'$, the $\tau_w$, $N$, the $Z_q$ for $q\in S'$, and $\epsilon=\varepsilon_\Gamma\circ p'$, gives that $Z^1(G_{S'},N)$ is finite and
\[|N|\prod_{q\in S'}|Z_q|\le|Z'_L|\,|\operatorname{Hom}_{\varepsilon_\Gamma\circ p'}(N,\mathbf Z/\ell^k)|\prod_w|N^{\tau_w}|\,|N|^{|S'|},\]
where $Z'_L$ is the set of $z\in Z^1(G_{S'},N)$ whose restriction to $D_q$ lies in $Z_q$ for every $q\in S'$. As $p'$ is surjective, $\operatorname{Hom}_{\varepsilon_\Gamma\circ p'}(N,\mathbf Z/\ell^k)=\operatorname{Hom}_{\varepsilon_\Gamma}(N,\mathbf Z/\ell^k)$, and $N^{\tau_w}=N^{c'_w}$. As $|N|\le|Z_q|$ for the $|S'|-|V|$ places $q\notin V$, dividing by $|N|^{|S'|-|V|}$ gives
\[|N|\prod_{v\in V}|Z_v|\le|Z'_L|\,|\operatorname{Hom}_{\varepsilon_\Gamma}(N,\mathbf Z/\ell^k)|\prod_w|N^{c'_w}|\,|N|^{|V|}.\]
\item[(e)] \emph{Descent to $\Gamma$.} The map $y\mapsto y\circ p'$ is injective on maps $\Gamma\to N$, as $p'$ is surjective, and it carries $Z_L$ into the finite group $Z^1(G_{S'},N)$; so $Z_L$ is finite. It remains to show that every $z\in Z'_L$ is $y\circ p'$ for some $y\in Z_L$, so that $|Z'_L|\le|Z_L|$. Let $q\in S'\setminus V$. There is $Y$ with $z(s)=sY-Y$ for $s\in I_q$; for $s\in I_q\cap H$, $p(s)=1$, so $s$ acts trivially and $z(s)=0$. The set of $g\in G_{S'}$ with $p'(g)=1$ and $z(g)=0$ is a subgroup, by the cocycle relation. It is normal: if $x$ lies in it and $g\in G_{S'}$, then $z(gx)=z(g)+g\,z(x)=z(g)$, and, as $x$ acts trivially, $z(gxg^{-1})=z(gx)+gx\,z(g^{-1})=z(g)+g\,z(g^{-1})=z(1)=0$. It is closed, as $p'$ is continuous and $z$ is locally constant. So it contains $\ker p'$, by (c). Hence $z(g)$ depends only on $p'(g)$, and $z=y\circ p'$ for a map $y:\Gamma\to N$. It is locally constant, as $p'$ is a continuous surjection of compact groups and so $\Gamma$ carries the quotient topology, and it is a cocycle. For $v\in V$, $y\circ\mathrm{loc}_v$ is the restriction of $z$ to $G_v=D_v$, which lies in $Z_v$; so $y\in Z_L$.
\end{enumerate}
By (b) the link satisfies the local counts, and by (d) and (e) it satisfies the global count for $(c'_w)$.
% linkOf (LinkOf.lean, 202 lines) for (a): K for H (with del_open, Psi_open and the private idx_prime), S the places above l, hker from k.p_ker with the private clos_comap, the private equiv_ex identifying G_{L^+} with its image upl in G_Q, Pa.ccEx and cc for tau_w (the Lean indexes the c'_w by Phi.Inf through a bijection ee with the infinite places, using card_inf and Pr.degEq), and k.j_range, k.j_ine, k.dv (with the private deg_eq) and k.eps for the local data; the Lean packages these as GC.FLink. leOf2 (LeOf2.lean) for (b): leOf (LcOrbit.lean, 138 lines: zRep, the fields lf, lt and le of HcPublished, cardZ for the cocycles, cardInv and Th.dualZ for the dual term), transported by leTrans (LeTrans.lean) along the isomorphism isoOf of G_v with Dec (A v). gwlOf (FiniteCocyclesOfFiniteH1.lean, 321 lines) for (c) to (e): GC.levelS for (c), tran for the places of V and unrLow for the others, bare for the first inequality of (d), numer for the division, and desc with hker for (e).
\end{proof}

\begin{definition}\label{sha-automorphy-unramified-tangents-count-cocycles}
In the situation of Lemma~\ref{sha-automorphy-unramified-point}, let $(r,(g_v))$ be a continuous triangulation over $\mathcal O_1$, for $n=3$ (Definition~\ref{sha-automorphy-unramified-near-completion-tangent-triangulation}), and let $k\ge1$. Write $\mathcal O_k=\mathcal O_1/\ell^k\mathcal O_1$, $q_k=|\mathcal O_k|$, $\bar x$ for the image of an element or a matrix over $\mathcal O_1$ modulo $\ell^k$, and $r_k$ for $r$ followed by the reduction $G_3(\mathcal O_1)\to G_3(\mathcal O_k)$. Let $\Gamma$ act on $W_k=M_3(\mathcal O_k)$ by $s\cdot X=r_k(s)\cdot X$, and for each place $v$ of $L^+$ above $\ell$ let $t_v:G_v\to GL_3(\mathcal O_1)$ be the homomorphism $t_v(h)=g_v\,\mathrm{mat}(r(\mathrm{loc}_vh))\,g_v^{-1}$, with upper triangular values.
\begin{enumerate}
\item Let $Z^{\mathrm{tri}}_{v,k}$ be the set of locally constant maps $z:G_v\to W_k$ with $z(hh')=z(h)+\bar t_v(h)z(h')\bar t_v(h)^{-1}$ for all $h,h'\in G_v$, with $z(h)$ upper triangular for every $h$, and with the diagonal entries of $z(h)$ equal to $0$ for $h\in I_v$.
\item Let $Z_{v,k}$ be the set of maps $G_v\to W_k$ of the form $h\mapsto\bar g_v^{-1}z(h)\bar g_v+(\mathrm{loc}_v(h)\cdot Y-Y)$ with $z\in Z^{\mathrm{tri}}_{v,k}$ and $Y\in W_k$.
\item Let $Z_k$ be the set of $z\in Z^1(\Gamma,W_k)$ (Definition~\ref{sha-automorphy-unramified-tangents-count-counts}) with $z\circ\mathrm{loc}_v\in Z_{v,k}$ for every $v$ and $z(c_\Gamma)=0$, such that the entries off the diagonal of $z(g_0)\,\mathrm{mat}(r_k(g_0))-z(g_0^*)\,\mathrm{mat}(r_k(g_0^*))$ vanish, where $g_0^*=c_\Gamma g_0^{-1}c_\Gamma^{-1}$.
\end{enumerate}
\end{definition}
\leanref{Development/AutomorphicArithmetic/AutomorphicVarieties/ModularityArgument/LinkedUnramifiedChains/CrystallineCardinalities/UnitStep.lean}{UpperTheory.V.Z.J.triOf}
\leanref{Development/AutomorphicArithmetic/AutomorphicVarieties/ModularityArgument/LinkedUnramifiedChains/CrystallineCardinalities/UnitStep.lean}{UpperTheory.V.Z.J.adK}
\leanref{Development/AutomorphicArithmetic/AutomorphicVarieties/ModularityArgument/LinkedUnramifiedChains/CrystallineCardinalities/UnitStep.lean}{UpperTheory.V.Z.J.ZF}
\leanref{Development/AutomorphicArithmetic/AutomorphicVarieties/ModularityArgument/LinkedUnramifiedChains/CrystallineCardinalities/UnitStep.lean}{UpperTheory.V.Z.J.Zv}
\leanref{Development/AutomorphicArithmetic/AutomorphicVarieties/ModularityArgument/LinkedUnramifiedChains/CrystallineCardinalities/UnitStep.lean}{UpperTheory.V.Z.J.ZN}

\begin{lemma}\label{sha-automorphy-unramified-tangents-count-triangulation}
In the situation of Lemma~\ref{sha-automorphy-unramified-point}, let $\rho_\xi$ be the linked point with values in $\mathcal O_1$ given by $\xi$, with characters $\psi_{v,i}$.
\begin{enumerate}
\item There is a continuous triangulation $(r,(g_v))$ of $\rho_\xi$ over $\mathcal O_1$ (Definition~\ref{sha-automorphy-unramified-near-completion-tangent-triangulation}).
\item Let $(r,(g_w))$ be such a triangulation, $v$ a place of $L^+$ above $\ell$, $g\in GL_3(\mathcal O_1)$, and $t:G_v\to GL_3(\mathcal O_1)$ a homomorphism such that $t(h)=g\rho_\xi(\mathrm{loc}_vh)g^{-1}$ is upper triangular with diagonal entries $\psi_{v,1}(h),\psi_{v,2}(h),\psi_{v,3}(h)$ for every $h\in G_v$. Then $(r,(g'_w))$, with $g'_v=g$ and $g'_w=g_w$ for $w\ne v$, is a continuous triangulation of $\rho_\xi$, and its homomorphism $t_v$ (Definition~\ref{sha-automorphy-unramified-tangents-count-cocycles}) is $t$.
\end{enumerate}
\end{lemma}
\leanref{Development/AutomorphicArithmetic/AutomorphicVarieties/FrobeniusSmoothVectors/ConjVal.lean}{UpperTheory.V.Fr.pr_box}
\begin{proof}
Write $\rho=\rho_\xi$, $c=c_\Gamma$ and $\theta(y)=J_3{}^ty^{-1}J_3$.
\begin{enumerate}
\item \emph{The homomorphism.} For $s\in\Gamma$ let $s'=s$ if $\delta_\Gamma(s)=1$ and $s'=sc$ otherwise; as $\delta_\Gamma(c)=-1$, $s'\in\Gamma_0$. Put $r(s)=(\rho(s'),\delta_\Gamma(s))\in G_3(\mathcal O_1)$. Condition (a) of Definition~\ref{sha-descent-link}(2), applied to $y^{-1}$, gives $\theta(\rho(y))=\rho(cyc^{-1})$ for $y\in\Gamma_0$. If $\delta_\Gamma(s)=1$, then $(st)'=st'$ and $r(s)r(t)=(\rho(s)\rho(t'),\delta_\Gamma(t))=r(st)$. If $\delta_\Gamma(s)=-1$, then $r(s)r(t)=(\rho(sc)\theta(\rho(t')),-\delta_\Gamma(t))=(\rho(sct'c^{-1}),\delta_\Gamma(st))$, and $sct'c^{-1}=st'c=(st)'$, as $c^2=1$ by residual condition (4) of Definition~\ref{sha-descent-residual}. So $r$ is a homomorphism with sign $\delta_\Gamma$, $\mathrm{mat}(r(c))=\rho(c^2)=1$ and $\mathrm{mat}(r(y))=\rho(y)$ for $y\in\Gamma_0$.
\item \emph{The triangulation.} Lemma~\ref{sha-automorphy-unramified-near-completion-flag}, applied to the local ring $\mathcal O_1$ and the restriction of $\xi$ to $A_3^\circ$, gives $g_v\in GL_3(\mathcal O_1)$ with $g_v\rho(\mathrm{loc}_vh)g_v^{-1}$ upper triangular with diagonal entries $\psi_{v,1}(h),\psi_{v,2}(h),\psi_{v,3}(h)$, since $\xi$ carries $\psi^{\mathrm u}_{v,i}$ to $\psi_{v,i}$. So $(r,(g_v))$ is a triangulation of $\rho$. It is continuous: $\Gamma_0$ is open (residual condition (4)), so $s\mapsto s'$ is continuous, and the entries of $\rho$ are continuous for the $\mathfrak m_{\mathcal O_1}$-adic topology by condition (d).
\item \emph{Proof of (2).} The conditions on a triangulation of $\rho$ concern each place separately, and at $v$ the matrix $g\,\mathrm{mat}(r(\mathrm{loc}_vh))\,g^{-1}=g\rho(\mathrm{loc}_vh)g^{-1}=t(h)$ is upper triangular with the diagonal entries $\psi_{v,i}(h)$. The homomorphism $r$ is unchanged, so the triangulation stays continuous.
\end{enumerate}
\end{proof}
% box_pt (BoxUnit.lean, private) for (1): pr_box (ConjVal.lean: rp with pr and ee, theta_r for condition (a), pr_cont from del_open) with the flags of Fr.flag (Lemma sha-automorphy-unramified-near-completion-flag) for xi composed with the inclusion of A_3^circ (S.A.val), and lift_pt for the continuity. box_unit (BoxUnit.lean, private) for (2): Function.update B_0.g v g, with triOf B v = t.

\begin{lemma}\label{sha-automorphy-unramified-tangents-count-local-conj}
Let $C$ be a commutative ring, $T\in GL_3(C)$ an upper triangular matrix and $B=T^{-1}$, and let $E_{13}$ be the matrix with entry $1$ at $(1,3)$ and $0$ elsewhere.
\begin{enumerate}
\item $B$ is upper triangular, and $T_{ii}B_{ii}=1$ for $i=1,2,3$.
\item For every upper triangular $X\in M_3(C)$, $TXB$ is upper triangular, $(TXB)_{ii}=X_{ii}$ for $i=1,2,3$, and $(TXB)_{i,i+1}=T_{ii}B_{i+1,i+1}X_{i,i+1}-T_{i,i+1}B_{i+1,i+1}(X_{ii}-X_{i+1,i+1})$ for $i=1,2$.
\item $T(xE_{13})B=T_{11}B_{33}\,xE_{13}$ for every $x\in C$.
\item If $Y\in M_3(C)$ is such that $TYB-Y$ is upper triangular, and $T_{ii}-T_{jj}$ is a unit for all $i\ne j$, then $Y$ is upper triangular.
\end{enumerate}
\end{lemma}
\leanref{Development/AutomorphicArithmetic/AutomorphicVarieties/ModularityArgument/LinkedUnramifiedChains/CrystallineCardinalities/CjFacts.lean}{UpperTheory.V.Z.J.cjFacts}
\begin{proof}
(1) The inverse of an invertible upper triangular matrix is upper triangular, and the diagonal entries of $TB=1$ are $T_{ii}B_{ii}=1$.

(2) As $T$, $X$ and $B$ are upper triangular, the entry $(i,j)$ of $TXB$ is the sum of the $T_{ip}X_{pq}B_{qj}$ over $i\le p\le q\le j$. It vanishes for $i>j$, and it is $T_{ii}X_{ii}B_{ii}=X_{ii}$ for $i=j$. For $j=i+1$ it is $T_{ii}X_{ii}B_{i,i+1}+T_{ii}X_{i,i+1}B_{i+1,i+1}+T_{i,i+1}X_{i+1,i+1}B_{i+1,i+1}$, and the entry $(i,i+1)$ of $TB=1$ gives $T_{ii}B_{i,i+1}=-T_{i,i+1}B_{i+1,i+1}$, which gives the formula.

(3) The entry $(i,j)$ of $TE_{13}B$ is $T_{i1}B_{3j}$, which vanishes unless $i=1$ and $j=3$.

(4) Let $Z=TYB-Y$. Multiplying on the right by $T$ gives $TY-YT=ZT$, which is upper triangular. Its entry $(3,1)$ gives $(T_{33}-T_{11})Y_{31}=0$, so $Y_{31}=0$. Then its entry $(2,1)$ gives $(T_{22}-T_{11})Y_{21}=-T_{23}Y_{31}=0$, and its entry $(3,2)$ gives $(T_{33}-T_{22})Y_{32}=Y_{31}T_{12}=0$; so $Y_{21}=Y_{32}=0$.
\end{proof}
% cjFacts (CjFacts.lean, 214 lines), for T the values of a homomorphism t with TriRep t, reduced modulo l^k: (1) is hB (Matrix.blockTriangular_inv_of_blockTriangular) and hdiag (inv_entries); (2) is conj_facts; (3) is conj_single; (4) is up_of_conj, applied at the element of TriRep.dist. Its first part, the open kernel of the action, is step (c) of the proof of Lemma sha-automorphy-unramified-tangents-count-local.

\begin{lemma}\label{sha-automorphy-unramified-tangents-count-local-lift}
Let $D$ be a profinite group, $\epsilon:D\to\mathbf Z_\ell^\times$ a homomorphism and $k\ge1$; let $Z^1(D,N)$ and $\operatorname{Hom}_\epsilon(N,\mathbf Z/\ell^k)$ be as in Definition~\ref{sha-automorphy-unramified-tangents-count-counts}.
\begin{enumerate}
\item Let $d\ge0$, and suppose that for every $k'\ge1$ and every finite abelian group $N$ killed by $\ell^{k'}$ with an action of $D$ whose kernel is open, $Z^1(D,N)$ is finite and $|Z^1(D,N)|=|N|^{1+d}\,|\operatorname{Hom}_\epsilon(N,\mathbf Z/\ell^{k'})|$. Let $N$ be such a group for $k'=k$, $N'\subseteq N$ a subgroup stable under $D$, and $C\ge0$ an integer with $|\operatorname{Hom}_\epsilon(N',\mathbf Z/\ell^k)|\le C$. Let $z:D\to N$ be a locally constant map with $z(st)-z(s)-sz(t)\in N'$ for all $s,t\in D$. Then there is $y\in Z^1(D,N)$ with $y(s)-C!\,z(s)\in N'$ for every $s\in D$.
\item Let $\mathcal O$ be a discrete valuation ring with finite residue field of characteristic $\ell$ and with $\ell\ne0$ in $\mathcal O$, together with a homomorphism $\mathbf Z_\ell\to\mathcal O$, and let $\mathcal O_k=\mathcal O/\ell^k\mathcal O$. Let $N$ be an abelian group with an action of $D$, $e:\mathcal O_k\to N$ an additive bijection and $\chi:D\to\mathcal O$ a map with $s\cdot e(x)=e(\chi(s)x)$ for $s\in D$ and $x\in\mathcal O_k$. If $c=\chi(s_0)-\epsilon(s_0)$ is nonzero for some $s_0\in D$, then $\mathcal O/c\mathcal O$ is finite and $|\operatorname{Hom}_\epsilon(N,\mathbf Z/\ell^k)|\le|\mathcal O/c\mathcal O|$.
\end{enumerate}
\end{lemma}
\leanref{Development/AutomorphicArithmetic/AutomorphicVarieties/ModularityArgument/LinkedUnramifiedChains/CrystallineCardinalities/HomMuLe.lean}{UpperTheory.V.Z.J.lift}
\leanref{Development/AutomorphicArithmetic/AutomorphicVarieties/ModularityArgument/LinkedUnramifiedChains/CrystallineCardinalities/HomBd.lean}{UpperTheory.V.Z.J.homBd}
\leanref{Development/AutomorphicArithmetic/AutomorphicVarieties/ModularityArgument/LinkedUnramifiedChains/CrystallineCardinalities/DualCard.lean}{UpperTheory.V.Z.J.dualCard}
\begin{proof}
A finite abelian group $A$ killed by $\ell^k$ has $|A|$ homomorphisms to $\mathbf Z/\ell^k$, as it is a product of cyclic groups of orders dividing $\ell^k$.

(1) Let $Q=N/N'$, with the induced action. The groups $N'$ and $Q$ are killed by $\ell^k$, and the kernels of the actions on them contain that on $N$, so they are open. Reduction modulo $N'$ is a homomorphism $\pi:Z^1(D,N)\to Z^1(D,Q)$, whose kernel consists of cocycles with values in $N'$, so $|\ker\pi|\le|Z^1(D,N')|$. Composition with $N\to Q$ embeds $\operatorname{Hom}_\epsilon(Q,\mathbf Z/\ell^k)$ in $\operatorname{Hom}_\epsilon(N,\mathbf Z/\ell^k)$, which contains $0$. Let $m$ be the index of the image of $\pi$ in the finite group $Z^1(D,Q)$. Then $m\,|Z^1(D,N)|=|Z^1(D,Q)|\,|\ker\pi|$, and by the hypothesis for $N$, $N'$ and $Q$,
\[m\,|N|^{1+d}\,|\operatorname{Hom}_\epsilon(N,\mathbf Z/\ell^k)|\le|Q|^{1+d}\,|\operatorname{Hom}_\epsilon(N,\mathbf Z/\ell^k)|\,|N'|^{1+d}\,C=|N|^{1+d}\,|\operatorname{Hom}_\epsilon(N,\mathbf Z/\ell^k)|\,C.\]
So $1\le m\le C$, and $m$ divides $C!$. The map $w:s\mapsto z(s)\bmod N'$ lies in $Z^1(D,Q)$, so $m\,w$, and hence $C!\,w$, lies in the image of $\pi$: there is $y\in Z^1(D,N)$ with $\pi(y)=C!\,w$, that is, $y(s)-C!\,z(s)\in N'$ for every $s$.

(2) As $c\ne0$ and the residue field is finite, $\mathcal O/c\mathcal O$ is finite. The element $\epsilon(s_0)\in\mathbf Z_\ell$ is congruent modulo $\ell^k$ to an integer $n\ge0$, so it acts on $\mathcal O_k$ and on $\mathbf Z/\ell^k$ as $n$. Let $f\in\operatorname{Hom}_\epsilon(N,\mathbf Z/\ell^k)$ and $a\in\mathcal O_k$. Then $f(e(ca))=f(s_0\cdot e(a))-f(e(na))=\epsilon(s_0)f(e(a))-nf(e(a))=0$. So $f\circ e$ factors through $\mathcal O_k/c\mathcal O_k$, and $f$ is determined by the induced homomorphism $\mathcal O_k/c\mathcal O_k\to\mathbf Z/\ell^k$. The group $\mathcal O_k/c\mathcal O_k$ is killed by $\ell^k$, so it has $|\mathcal O_k/c\mathcal O_k|$ such homomorphisms, and it is a quotient of $\mathcal O/c\mathcal O$. Hence $|\operatorname{Hom}_\epsilon(N,\mathbf Z/\ell^k)|\le|\mathcal O/c\mathcal O|$.
\end{proof}
% lift (HomMuLe.lean, 108 lines): ZcA, card_ZcA, the map ZcQ (pi), ker_le, homMu_le, idx_mul, then Nat.dvd_factorial and AddSubgroup.nsmul_index_mem. homBd (HomBd.lean, 75 lines): kill and wd, the injection Phi, dualCard (DualCard.lean) for O_k / c O_k, and the surjection psi from O / c O; the congruence of epsilon(s_0) is the field hsm of okFin. The Lean states lift for any topological group D with LE l D eps d.

\begin{lemma}\label{sha-automorphy-unramified-tangents-count-local}
In the situation of Definition~\ref{sha-automorphy-unramified-tangents-count-cocycles}, suppose that the link satisfies the local counts (Definition~\ref{sha-automorphy-unramified-tangents-count-counts}) and that $(r,(g_v))$ is a continuous triangulation of a linked point $(\rho,\psi)$ with values in $\mathcal O_1$ that lifts the residual one. Let $v$ be a place of $L^+$ above $\ell$.
\begin{enumerate}
\item For every $k\ge1$, $Z_{v,k}$ is a subgroup of $Z^1(G_v,W_k)$, for $G_v$ acting through $\mathrm{loc}_v$, that contains the maps $h\mapsto\mathrm{loc}_v(h)\cdot Y-Y$ for $Y\in W_k$, and $q_k^{9+3d_v}\le|Z_{v,k}|$.
\item Suppose that for some $i'\in\{1,2\}$, $t_v(h)_{i'i'}=\varepsilon_\Gamma(\mathrm{loc}_vh)\,t_v(h)_{i'+1,i'+1}$ for every $h\in G_v$, and that there are continuous maps $u,\beta:G_v\to\mathcal O_1$, with $u$ a nonzero homomorphism to $(\mathcal O_1,+)$ that vanishes on $I_v$, such that $\beta(hh')=\beta(h)+\varepsilon_\Gamma(\mathrm{loc}_vh)\beta(h')-t_v(h)_{i',i'+1}\bigl(t_v(h)^{-1}\bigr)_{i'+1,i'+1}u(h')$ for all $h,h'\in G_v$. Then there is an integer $C$ with $q_k^{10+3d_v}\le C|Z_{v,k}|$ for every $k\ge1$.
\end{enumerate}
\end{lemma}
\leanref{Development/AutomorphicArithmetic/AutomorphicVarieties/ModularityArgument/LinkedUnramifiedChains/CrystallineCardinalities/AdicContinuousSumThree.lean}{UpperTheory.V.Z.J.triOk}
\leanref{Development/AutomorphicArithmetic/AutomorphicVarieties/ModularityArgument/LinkedUnramifiedChains/CrystallineCardinalities/LocLow.lean}{UpperTheory.V.Z.J.locLow}
\leanref{Development/AutomorphicArithmetic/AutomorphicVarieties/ModularityArgument/LinkedUnramifiedChains/CrystallineCardinalities/LocHigh.lean}{UpperTheory.V.Z.J.locHigh}
\leanref{Development/AutomorphicArithmetic/AutomorphicVarieties/ModularityArgument/LinkedUnramifiedChains/CrystallineCardinalities/LcOfKer.lean}{UpperTheory.V.Z.J.zvCnt}
\begin{proof}
Let $k\ge1$ and write $q=q_k$ and $d=d_v$.
\begin{enumerate}
\item[(a)] \emph{The homomorphism $t_v$.} Its entries are continuous, being those of conjugates of $\mathrm{mat}\circ r\circ\mathrm{loc}_v$ by fixed matrices, as $\mathrm{loc}_v$ is continuous (residual condition (4) of Definition~\ref{sha-descent-residual}); its values are upper triangular; and its diagonal entries are the $\psi_{v,i}$. By residual condition (7) there is $h_0\in G_v$ with $a=\bar\Psi_\Gamma(\mathrm{loc}_vh_0)$ satisfying $a^2\ne1$. As $(\rho,\psi)$ lifts the residual one, the residue of $\psi_{v,i}(h_0)$ is $a^{w_{3,i}}$ (condition (e) of Definition~\ref{sha-descent-link}(2)), and $a^{w_{3,i}}\ne a^{w_{3,j}}$ for $i\ne j$, as $w_{3,i}-w_{3,j}$ divides $2$; so the differences $t_v(h_0)_{ii}-t_v(h_0)_{jj}$, $i\ne j$, are units. By residual condition (8) there are $u_1,\dots,u_d\in I_v$ such that every locally constant homomorphism from $G_v$ to a finite abelian group killed by $\ell^k$ that vanishes at $u_1,\dots,u_d$ vanishes on $I_v$.
\item[(b)] \emph{A lower bound for $Z^{\mathrm{tri}}_{v,k}$.} Let $T_k\subseteq W_k$ be the group of upper triangular matrices, with $|T_k|=q^6$. The group $G_v$ acts on $T_k$ by $X\mapsto\bar t_v(h)X\bar t_v(h)^{-1}$, with open kernel by (a), and this action does not change the diagonal entries, as $t_v(h)$ and $t_v(h)^{-1}$ are upper triangular with $t_v(h)_{ii}(t_v(h)^{-1})_{ii}=1$. By the local count, $Z^1(G_v,T_k)$ is finite with at least $|T_k|^{1+d}=q^{6+6d}$ elements, as $\operatorname{Hom}_{\varepsilon_\Gamma\circ\mathrm{loc}_v}(T_k,\mathbf Z/\ell^k)$ contains $0$. For $z\in Z^1(G_v,T_k)$, each diagonal entry $h\mapsto z(h)_{jj}$ is a locally constant homomorphism to $\mathcal O_k$. The homomorphism $z\mapsto(z(u_a)_{jj})_{a,j}$ from $Z^1(G_v,T_k)$ to $\mathcal O_k^{3d}$ has kernel $Z^{\mathrm{tri}}_{v,k}$, by (a). So $q^{6+6d}\le|Z^{\mathrm{tri}}_{v,k}|\,q^{3d}$, that is, $q^{6+3d}\le|Z^{\mathrm{tri}}_{v,k}|$; and $Z^{\mathrm{tri}}_{v,k}$ is finite.
\item[(c)] \emph{The action on $W_k$.} For $h\in G_v$ and $X\in W_k$ let $h\star X=\bar t_v(h)X\bar t_v(h)^{-1}$. This is an action of $G_v$ on $W_k$ by group automorphisms, and the elements of $Z^{\mathrm{tri}}_{v,k}$ are cocycles for it. Its kernel is open: $\ell^k\mathcal O_1$ contains a power $\mathfrak m^n$ of the maximal ideal $\mathfrak m$ of $\mathcal O_1$, and the $h$ with $t_v(h)\equiv1$ modulo $\mathfrak m^n$ form an open subgroup by (a), which acts trivially. For $h\in G_v$ and $j=1,2$ let $p_j(h)=t_v(h)_{jj}(t_v(h)^{-1})_{j+1,j+1}$ and $q_j(h)=t_v(h)_{j,j+1}(t_v(h)^{-1})_{j+1,j+1}$, and let $\chi(h)=t_v(h)_{11}(t_v(h)^{-1})_{33}$. Lemma~\ref{sha-automorphy-unramified-tangents-count-local-conj}, applied to $\bar t_v(h)$ over $\mathcal O_k$, gives: for $X$ upper triangular, $h\star X$ is upper triangular with the diagonal of $X$, and $(h\star X)_{j,j+1}=p_j(h)X_{j,j+1}-q_j(h)(X_{jj}-X_{j+1,j+1})$; $h\star(xE_{13})=\chi(h)xE_{13}$ for $x\in\mathcal O_k$; and, with $h_0$ of (a), if $h\star Y-Y$ is upper triangular for every $h$, then $Y$ is upper triangular. Write $\varepsilon(h)=\varepsilon_\Gamma(\mathrm{loc}_vh)$. If $t_v(h)_{i'i'}=\varepsilon(h)t_v(h)_{i'+1,i'+1}$ for every $h$, then $p_{i'}(h)=\varepsilon(h)$, as $t_v(h)_{i'+1,i'+1}(t_v(h)^{-1})_{i'+1,i'+1}=1$.
\item[(d)] \emph{A cocycle that moves the diagonal.} Suppose the hypotheses of (2), and let $t=t_v(h_0)$. Let $c_0=\chi(h_0)-\varepsilon(h_0)$. As $t_{33}(t^{-1})_{33}=1$, $c_0=\varepsilon(h_0)(t^{-1})_{33}(t_{22}-t_{33})$ if $i'=1$ and $c_0=(t^{-1})_{33}(t_{11}-t_{22})$ if $i'=2$; so $c_0\ne0$ by (a). Let $C_1=|\mathcal O_1/c_0\mathcal O_1|$, which is finite. Choose $h_1\in G_v$ with $u(h_1)\ne0$, and let $c=C_1!\,u(h_1)$. It is nonzero, as $\mathcal O_1$ has characteristic $0$: its residue field has characteristic $\ell$, and $\ell\ne0$ in $\mathcal O_1$. Neither $C_1$ nor $c$ depends on $k$.

The multiples $xE_{13}$, $x\in\mathcal O_k$, form a subgroup $N'$ of $W_k$ stable under $\star$, on which $h$ acts by $\chi(h)$, by (c). So Lemma~\ref{sha-automorphy-unramified-tangents-count-local-lift}(2), with $e(x)=xE_{13}$ and $s_0=h_0$, gives $|\operatorname{Hom}_{\varepsilon_\Gamma\circ\mathrm{loc}_v}(N',\mathbf Z/\ell^k)|\le C_1$.

For $x,b\in\mathcal O_k$ let $Z(x,b)=xE_{11}+bE_{12}$ if $i'=1$ and $Z(x,b)=bE_{23}-xE_{33}$ if $i'=2$, and let $z(h)=Z(\bar u(h),\bar\beta(h))$. The map $z$ is locally constant, as $u$ and $\beta$ are continuous. Let $h,h'\in G_v$ and $Y=h\star z(h')$. By (c), $Y$ is upper triangular with the diagonal of $z(h')$. Its entry $(i',i'+1)$ is $\varepsilon(h)\bar\beta(h')-q_{i'}(h)\bar u(h')$, as $Z(x,b)_{i'i'}-Z(x,b)_{i'+1,i'+1}=x$, and its other entry $(j,j+1)$ vanishes, as $Z(x,b)_{j,j+1}=0$ and $Z(x,b)_{jj}=Z(x,b)_{j+1,j+1}$ for $j\ne i'$. As $u$ is additive and $\beta(hh')=\beta(h)+\varepsilon(h)\beta(h')-q_{i'}(h)u(h')$, the matrix $z(hh')-z(h)-Y$ vanishes outside the entry $(1,3)$, that is, it lies in $N'$.

By the local counts, Lemma~\ref{sha-automorphy-unramified-tangents-count-local-lift}(1), for $D=G_v$, $\epsilon=\varepsilon_\Gamma\circ\mathrm{loc}_v$, $d$, $W_k$ with the action $\star$, $N'$ and $C_1$, gives a cocycle $y$ of $G_v$ for $\star$ with $y(h)-C_1!\,z(h)\in N'$ for every $h$. Outside the entry $(1,3)$, $y(h)$ agrees with $C_1!\,z(h)$. So $y(h)$ is upper triangular, its diagonal vanishes for $h\in I_v$, where $u$ vanishes, and $y\in Z^{\mathrm{tri}}_{v,k}$. Moreover $y(h_1)_{i'i'}-y(h_1)_{i'+1,i'+1}=C_1!\,\bar u(h_1)=\bar c$.
\item[(e)] \emph{Cocycles with equal diagonal entries.} Suppose that $t_v(h)_{i'i'}=\varepsilon(h)t_v(h)_{i'+1,i'+1}$ for every $h$, and let $Z'$ be the set of $z\in Z^{\mathrm{tri}}_{v,k}$ with $z(h)_{i'i'}=z(h)_{i'+1,i'+1}$ for every $h$. Let $B_0\subseteq W_k$ be the group of upper triangular $X$ with $X_{i'i'}=X_{i'+1,i'+1}$. It has $q^5$ elements, it is stable under $\star$ by (c), and the kernel of the action on it is open. The map $\lambda:B_0\to\mathcal O_k$, $X\mapsto X_{i',i'+1}$, is surjective, and $\lambda(h\star X)=\varepsilon(h)\lambda(X)$ by (c), as $p_{i'}(h)=\varepsilon(h)$ and $X_{i'i'}=X_{i'+1,i'+1}$. As $\mathcal O_k$ is killed by $\ell^k$, it has $q$ homomorphisms to $\mathbf Z/\ell^k$, and their composites with $\lambda$ are $q$ distinct elements of $\operatorname{Hom}_{\varepsilon_\Gamma\circ\mathrm{loc}_v}(B_0,\mathbf Z/\ell^k)$. So by the local count, $Z^1(G_v,B_0)$, for the action $\star$, is finite with at least $(q^5)^{1+d}q=q^{6+5d}$ elements. Let $P=\{2,3\}$ if $i'=1$ and $P=\{1,2\}$ if $i'=2$. For $z\in Z^1(G_v,B_0)$ and $j\in P$, $h\mapsto z(h)_{jj}$ is a locally constant homomorphism to $\mathcal O_k$, as $\star$ does not change diagonal entries. The homomorphism $z\mapsto(z(u_a)_{jj})_{a\le d,\,j\in P}$ from $Z^1(G_v,B_0)$ to $\mathcal O_k^{2d}$ has its kernel in the set of the $z$ with $z(h)_{jj}=0$ for all $j\in P$ and $h\in I_v$, by (a); so this set has at least $q^{6+3d}$ elements. Each of its elements, viewed with values in $W_k$, lies in $Z'$: its diagonal vanishes on $I_v$, as its entries $(j,j)$ with $j\in P$ vanish there and its entries $(i',i')$ and $(i'+1,i'+1)$ are equal, one of them in $P$. So $|Z'|\ge q^{6+3d}$.
\item[(f)] \emph{The unit step.} Suppose the hypotheses of (2), and let $c$, $y$ and $h_1$ be as in (d) and $Z'$ as in (e). Let $J=\bar c\,\mathcal O_k$, and for $b\in J$ choose $a_b\in\mathcal O_k$ with $a_b\bar c=b$. As $Z^{\mathrm{tri}}_{v,k}$ is stable under sums and multiplication by elements of $\mathcal O_k$, $(z,b)\mapsto z+a_by$ maps $Z'\times J$ to $Z^{\mathrm{tri}}_{v,k}$; it is injective, since the difference of the entries $(i',i')$ and $(i'+1,i'+1)$ of $z(h_1)+a_by(h_1)$ is $b$. As $\mathcal O_k/J$ is a quotient of $\mathcal O_1/c\mathcal O_1$, $q\le|\mathcal O_1/c\mathcal O_1|\,|J|$. So $q^{7+3d}\le C|Z^{\mathrm{tri}}_{v,k}|$ with $C=|\mathcal O_1/c\mathcal O_1|$, which does not depend on $k$.
\item[(g)] \emph{From $Z^{\mathrm{tri}}_{v,k}$ to $Z_{v,k}$.} Let $\iota(X)=\bar g_v^{-1}X\bar g_v$ for $X\in W_k$. For $h\in G_v$, $\mathrm{loc}_v(h)\in\Gamma_0$ acts on $W_k$ by conjugation by $\mathrm{mat}(r_k(\mathrm{loc}_vh))=\bar g_v^{-1}\bar t_v(h)\bar g_v$, so $\mathrm{loc}_v(h)\cdot\iota(X)=\iota(h\star X)$. The conditions defining $Z^{\mathrm{tri}}_{v,k}$ are additive, so it is a group, and $\psi(z,Y)=\bigl(h\mapsto\iota(z(h))+\mathrm{loc}_v(h)\cdot Y-Y\bigr)$ is a homomorphism from $Z^{\mathrm{tri}}_{v,k}\times W_k$ onto $Z_{v,k}$. So $Z_{v,k}$ is a group. Its elements are locally constant cocycles for the action through $\mathrm{loc}_v$, as $\iota$ carries cocycles for $\star$ to such cocycles and the kernel of $\star$ is open; and it contains $\psi(0,Y)$, the map $h\mapsto\mathrm{loc}_v(h)\cdot Y-Y$. If $\psi(z,Y)=0$, then $Y'=-\iota^{-1}(Y)$ satisfies $z(h)=h\star Y'-Y'$ for every $h$; so $h\star Y'-Y'$ is upper triangular for every $h$, and $Y'$ is upper triangular by (c). As $(z,Y)$ is determined by $Y'$, the kernel of $\psi$ has at most $q^6$ elements. As $Z^{\mathrm{tri}}_{v,k}$ is finite by (b), $|Z^{\mathrm{tri}}_{v,k}|\,q^9=|Z_{v,k}|\,|\ker\psi|\le|Z_{v,k}|\,q^6$.
\item[(h)] \emph{Conclusion.} By (b) and (g), $q^{6+3d}q^9\le|Z_{v,k}|q^6$; with the rest of (g), this is (1). Under the hypotheses of (2), (f) and (g) give $q^{7+3d}q^9\le C|Z_{v,k}|q^6$, which is (2).
\end{enumerate}
% (a) is triOk (AdicContinuousSumThree.lean: adic_continuous_sum_three for the continuity, hF.f3 and pow_w_ne for the units, genI_of with the field gen of Frame.Ok for u_1, ..., u_d; the Lean calls this GenI). (b) is locLow (LocLow.lean: local_lower_triangular_cohomology, with triG, triG_card, cjFacts for the action and its open kernel, and wtCut). (c) is cjFacts (Lemma sha-automorphy-unramified-tangents-count-local-conj; the open kernel is its first part, from okFin and TriRep.cont). (d) is yEx (MatrixEntryZeroTwoSubmodule.lean, 198 lines): hχne for c_0 (from TriRep.dist and the first field of UnitStep), natCast_ne for c, the submodule matrixEntryZeroTwoSubmodule (N') with e02 and homBd for the bound C_1, zM (Z(x, b)) with zM_coc for the cocycle relation modulo N', and lift (Lemma sha-automorphy-unramified-tangents-count-local-lift). (e) is low0 (LowerTriangularAbsorbsTorsion.lean, 182 lines): b0 (B_0) with card_borel_torsion_lower_bound, ent (lambda) with ent_surj, the positions q j (P), and lower_triangular_absorbs_torsion (dualCard for the q homomorphisms, then wtCut for the 2d values at u_1, ..., u_d). (f) is locHigh (LocHigh.lean): hmem, the injection hcount and hJ. (g) is zvCnt (LcOfKer.lean, 205 lines): ad_cj and cjE (iota), cnt (psi, with key for the kernel and hT from the last part of cjFacts), card_upT, lc_of_ker. (h) is the private sel_cut of SelLow.lean.
\end{proof}

\begin{lemma}\label{sha-automorphy-unramified-tangents-count-adjoint}
In the situation of Definition~\ref{sha-automorphy-unramified-tangents-count-cocycles}, let $(r,(g_v))$ be a continuous triangulation of a linked point $(\rho,\psi)$ with values in $\mathcal O_1$ that lifts the residual one, and let $k\ge1$.
\begin{enumerate}
\item The kernel of the action of $\Gamma$ on $W_k$ is open.
\item $\operatorname{Hom}_{\varepsilon_\Gamma}(W_k,\mathbf Z/\ell^k)=\{0\}$ (Definition~\ref{sha-automorphy-unramified-tangents-count-counts}).
\item If $x=(A,-1)\in G_3(\mathcal O_1)$ has $x^2=1$, then at most $q_k^3$ elements of $W_k$ are fixed by the action of the reduction $\bar x$ of Definition~\ref{sha-automorphy-unramified-near-completion-tangent-triangulation}.
\item If the set $Z_L$ of the $z\in Z^1(\Gamma,W_k)$ with $z\circ\mathrm{loc}_v\in Z_{v,k}$ for every $v$ is finite, then $|Z_L|\le|Z_k|\,q_k^8$.
\end{enumerate}
\end{lemma}
\leanref{Development/AutomorphicArithmetic/AutomorphicVarieties/ModularityArgument/LinkedUnramifiedChains/CrystallineCardinalities/AdOpen.lean}{UpperTheory.V.Z.J.adOpen}
\leanref{Development/AutomorphicArithmetic/AutomorphicVarieties/ModularityArgument/LinkedUnramifiedChains/CrystallineCardinalities/DualZero.lean}{UpperTheory.V.Z.J.dualZero}
\leanref{Development/AutomorphicArithmetic/AutomorphicVarieties/ModularityArgument/LinkedUnramifiedChains/CrystallineCardinalities/ArchUp.lean}{UpperTheory.V.Z.J.archUp}
\leanref{Development/AutomorphicArithmetic/AutomorphicVarieties/ModularityArgument/LinkedUnramifiedChains/CrystallineCardinalities/NormCnt.lean}{UpperTheory.V.Z.J.normCnt}
\begin{proof}
\begin{enumerate}
\item[(1)] \emph{Proof of (1).} As $\mathcal O_1$ is a discrete valuation ring with $\ell\ne0$, $\ell^k\mathcal O_1$ contains a power of $\mathfrak m_{\mathcal O_1}$, so the entries of $\mathrm{mat}\circ r_k$ are locally constant; and $\delta_\Gamma$ is locally constant, as $\Gamma_0$ is open (residual condition (4) of Definition~\ref{sha-descent-residual}). So the kernel of $r_k$ is open, and it acts trivially on $W_k$.
\item[(3)] \emph{Proof of (3).} Write $J=J_3$, so that $J^2=1$ and ${}^tJ=J$. As $x^2=1$, $A\theta(A)=1$ with $\theta(A)=J{}^tA^{-1}J$, so $AJ=J{}^tA$, that is, ${}^t(AJ)=AJ$. Let $X\in W_k$ be fixed by $\bar x$, that is, $X=-\bar AJ{}^tXJ\bar A^{-1}$, and let $Y=X\bar AJ$. Then $Y=-\bar AJ{}^tX=-{}^t(\bar AJ){}^tX=-{}^tY$. As $\ell$ is odd (Definition~\ref{sha-descent-datum}), $2$ is a unit of $\mathcal O_k$, so the diagonal entries of $Y$ vanish and $Y$ is determined by $Y_{12}$, $Y_{13}$ and $Y_{23}$. As $X\mapsto X\bar AJ$ is injective, at most $q_k^3$ matrices $X$ are fixed.
\end{enumerate}
For (2), let $f\in\operatorname{Hom}_{\varepsilon_\Gamma}(W_k,\mathbf Z/\ell^k)$. Write $\bar a$ for the image in $\mathcal O_k$ of $a\in\mathbf Z_\ell$, $k_1$ for the residue field of $\mathcal O_1$, $\phi:\mathcal O_k\to k_1$ for the reduction, $I=\ker\phi$, the image of $\mathfrak m_{\mathcal O_1}$, $E_{ab}$ for the matrix with entry $1$ at $(a,b)$ and $0$ elsewhere, and $e_{3,i}$ as in Definition~\ref{sha-descent-link}(2).
\begin{enumerate}
\item[(2a)] \emph{Linearity.} The action of $\Gamma$ on $W_k$ is $\mathcal O_k$-linear. An element $x$ of $\mathcal O_k$ with $\phi(x)\ne0$ is a unit, being the image of a unit of $\mathcal O_1$; so for $x\in I$ and $y\in\mathcal O_k$, $1+yx$ is a unit, and $I$ lies in the Jacobson radical of $\mathcal O_k$. For $a\in\mathbf Z_\ell$ there is an integer $n\ge0$ with $a-n\in\ell^k\mathbf Z_\ell$, so $\bar a=n$ in $\mathcal O_k$, and $f(\bar aX)=nf(X)=af(X)$ for $X\in W_k$, with $a$ acting on $\mathbf Z/\ell^k$ through its reduction.
\item[(2b)] \emph{A submodule.} Let $N$ be the set of $X\in W_k$ with $f(yX)=0$ for every $y\in\mathcal O_k$, an $\mathcal O_k$-submodule. For $s\in\Gamma$ and $X\in W_k$, $s\cdot X-\overline{\varepsilon_\Gamma(s)}X\in N$: for $y\in\mathcal O_k$, $f(y(s\cdot X))=f(s\cdot(yX))=\varepsilon_\Gamma(s)f(yX)=f(\overline{\varepsilon_\Gamma(s)}yX)$ by (2a).
\item[(2c)] \emph{A criterion.} Let $a,b\in\{1,2,3\}$ and $s\in\Gamma$ be such that $\phi\bigl((s\cdot E_{ab})_{ab}\bigr)\ne\omega_\Gamma(s)$ and $\phi\bigl((s\cdot E_{ab})_{ij}\bigr)=0$ for $(i,j)\ne(a,b)$. Then $E_{ab}\in N+IW_k$. Indeed, $Y=s\cdot E_{ab}-\overline{\varepsilon_\Gamma(s)}E_{ab}$ lies in $N$ by (2b). As $\phi(\overline{\varepsilon_\Gamma(s)})=\omega_\Gamma(s)$, the entry $u=Y_{ab}$ has $\phi(u)\ne0$, so it is a unit, and the other entries of $Y$ lie in $I$. So $E_{ab}=u^{-1}Y-u^{-1}(Y-uE_{ab})$ lies in $N+IW_k$.
\item[(2d)] \emph{The action modulo $\mathfrak m$.} For $h\in\Gamma_0$, $r(h)=(\rho(h),1)$, and $\rho(h)$ reduces to the diagonal matrix with entries $\bar\Psi_\Gamma(h)^{e_{3,i}}$ in $k_1$, as $(\rho,\psi)$ lifts the residual one (condition (e) of Definition~\ref{sha-descent-link}(2)), and its inverse reduces to the diagonal matrix with entries $\bar\Psi_\Gamma(h)^{-e_{3,i}}$. Hence $\phi\bigl((h\cdot E_{ab})_{ij}\bigr)$ is $\bar\Psi_\Gamma(h)^{e_{3,a}-e_{3,b}}$ if $(i,j)=(a,b)$ and $0$ otherwise. As $\mathrm{mat}(r(c_\Gamma))=1$ and $\delta_\Gamma(c_\Gamma)=-1$, $c_\Gamma$ acts by $X\mapsto-J_3{}^tXJ_3$, whose entry $(i,j)$ is $-X_{4-j,4-i}$.
\item[(2e)] \emph{The matrices $E_{ab}$.} Let $e=e_{3,a}-e_{3,b}$. If $e=0$, residual condition (7) of Definition~\ref{sha-descent-residual} gives $h\in\Gamma_0$ with $\omega_\Gamma(h)\ne1=\bar\Psi_\Gamma(h)^e$, and (2c) and (2d) with $s=h$ give $E_{ab}\in N+IW_k$. If $e=\pm1$, residual condition (7) gives $h\in\Gamma_0$ with $\omega_\Gamma(h)\bar\Psi_\Gamma(h)^{-e}\ne1$, and the same argument applies. If $e=\pm2$, then $\{a,b\}=\{1,3\}$ and $b=4-a$, and residual condition (7) gives $g\in\Gamma$ with $\omega_\Gamma(g)\bar\Psi_\Gamma(g)^{-e}\delta_\Gamma(g)\ne1$. If $\delta_\Gamma(g)=1$, then $g\in\Gamma_0$ and the same argument applies. If $\delta_\Gamma(g)=-1$, let $h=c_\Gamma^{-1}g\in\Gamma_0$; then $\bar\Psi_\Gamma(h)=\bar\Psi_\Gamma(g)$, as $\bar\Psi_\Gamma(c_\Gamma)=1$ (residual condition (4)). By (2d), the entry $(i,j)$ of $g\cdot E_{ab}=c_\Gamma\cdot(h\cdot E_{ab})$ is $-(h\cdot E_{ab})_{4-j,4-i}$; its residue is $-\bar\Psi_\Gamma(g)^e$ at $(a,b)=(4-b,4-a)$ and $0$ elsewhere. As $-\bar\Psi_\Gamma(g)^e\ne\omega_\Gamma(g)$, (2c) applies with $s=g$.
\item[(2f)] \emph{Conclusion.} By (2e), the finitely generated $\mathcal O_k$-module $W_k$ is $N+IW_k$, and $I$ lies in the Jacobson radical; so $W_k=N$ by Nakayama's lemma \cite[Theorem~2.2]{Mat86}. For $X\in W_k$ this gives $f(X)=f(1\cdot X)=0$.
\end{enumerate}
For (4), write $c=c_\Gamma$, $g=g_0$, $g^*=cg^{-1}c$, which is $c_\Gamma g_0^{-1}c_\Gamma^{-1}$ as $c^2=1$ (residual condition (4)), $J=J_3$, $A=\mathrm{mat}(r_k(g))$, and $D(z)=z(g)\,\mathrm{mat}(r_k(g))-z(g^*)\,\mathrm{mat}(r_k(g^*))$.
\begin{enumerate}
\item[(4a)] \emph{A group.} The conditions defining $Z^{\mathrm{tri}}_{v,k}$ are additive, so it is a group, and $Z_{v,k}$ is the image of the homomorphism $(z,Y)\mapsto\bigl(h\mapsto\bar g_v^{-1}z(h)\bar g_v+\mathrm{loc}_v(h)\cdot Y-Y\bigr)$ on $Z^{\mathrm{tri}}_{v,k}\times W_k$. So the $Z_{v,k}$ and $Z_L$ are groups.
\item[(4b)] \emph{A homomorphism.} Let $E:Z_L\to\mathcal O_k^8$ send $z$ to the entries $z(c)_{11}$, $z(c)_{12}$, $z(c)_{21}$, $z(c)_{13}$, $z(c)_{31}$, $z(c)_{22}$, $D(z)_{12}$ and $D(z)_{21}$. It is a homomorphism, so $|Z_L|\le|\ker E|\,q_k^8$, and it remains to show that $\ker E\subseteq Z_k$. Let $z\in\ker E$.
\item[(4c)] \emph{The value at $c$.} As $r_k(c)=(1,-1)$, $c$ acts on $W_k$ by $X\mapsto-J{}^tXJ$, and $(J{}^tXJ)_{ij}=X_{4-j,4-i}$. The cocycle relation gives $z(1)=z(1)+z(1)$, so $z(1)=0$, and $0=z(c^2)=z(c)-J{}^tz(c)J$. So $z(c)_{ij}=z(c)_{4-j,4-i}$, which gives $z(c)_{33}=z(c)_{11}$, $z(c)_{23}=z(c)_{12}$ and $z(c)_{32}=z(c)_{21}$. With the six entries of $z(c)$ that $E$ kills, $z(c)=0$.
\item[(4d)] \emph{The entries at $g_0$.} As $g\in\Gamma_0$, $r_k(g)=(A,1)$, and $0=z(g^{-1}g)$ gives $z(g^{-1})=-A^{-1}z(g)A$. As $z(c)=0$, $z(cg^{-1}c)=c\cdot z(g^{-1})=J\,{}^t(A^{-1}z(g)A)\,J$, and $\mathrm{mat}(r_k(cg^{-1}c))=\theta(A^{-1})=J{}^tAJ$. Since ${}^t(A^{-1}z(g)A)\,{}^tA={}^tA\,{}^tz(g)$ and $J^2=1$, $D(z)=P-J{}^tPJ$ with $P=z(g)A$, that is, $D(z)_{ij}=P_{ij}-P_{4-j,4-i}$. So $D(z)_{13}=D(z)_{31}=0$, $D(z)_{23}=-D(z)_{12}$ and $D(z)_{32}=-D(z)_{21}$, and as $E$ kills $D(z)_{12}$ and $D(z)_{21}$, the entries of $D(z)$ off the diagonal vanish.
\end{enumerate}
By (4c) and (4d), $z\in Z_k$; so $|Z_L|\le|Z_k|\,q_k^8$.
% (1) is adOpen (AdOpen.lean: r_ker_open with okFin for the power of the maximal ideal). (3) is archUp (ArchUp.lean: arch_gen, with h2 from l != 2, the field l_odd of Frame.Ok). (2) is dualZero (DualZero.lean, 328 lines): (2a) is hlin, hunit, hjac and hbr with the private exists_nat_sub_mem_span_pow_and_toZModPow_eq_and_toZMod_eq; (2b) is N and hP; (2c) is hkey; (2d) is hG0 (hB.res and the fields of IsLift), hom, htheta and hc (hB.atc); (2e) is L1 (fields f2a and f2b of Frame.Ok) and L2 (field f2c, with Psi_c and del_c); (2f) is Submodule.le_of_le_smul_of_le_jacobson_bot. (4) is normCnt (NormCnt.lean, 162 lines): (4a) is the hypothesis hZ, supplied in selLow by zvCnt, and zlGrp; (4b) is ent and count; (4c) and (4d) are core, with jtj, z1, hsym, zgi, zh, Mh and hD, and hs for g_0^*.
\end{proof}

\begin{lemma}\label{sha-automorphy-unramified-tangents-count-selmer}
In the situation of Definition~\ref{sha-automorphy-unramified-tangents-count-cocycles}, suppose that the link satisfies the local counts, and the global count for a family $(c'_w)$ indexed by the infinite places $w$ of $L^+$ with $c_w'^2=1$ and $\delta_\Gamma(c'_w)=-1$ (Definition~\ref{sha-automorphy-unramified-tangents-count-counts}). Let $(r,(g_v))$ be a continuous triangulation of a linked point $(\rho,\psi)$ with values in $\mathcal O_1$ that lifts the residual one, and suppose that for some place $v$ the homomorphism $t_v$ satisfies the hypotheses of Lemma~\ref{sha-automorphy-unramified-tangents-count-local}(2). Then there is an integer $C$ such that for every $k\ge1$ the set $Z_k$ is finite and $q_k^2\le C|Z_k|$.
\end{lemma}
\leanref{Development/AutomorphicArithmetic/AutomorphicVarieties/ModularityArgument/LinkedUnramifiedChains/CrystallineCardinalities/SelLow.lean}{UpperTheory.V.Z.J.selLow}
\begin{proof}
Let $C$ be the integer of Lemma~\ref{sha-automorphy-unramified-tangents-count-local}(2) at $v$, let $k\ge1$, $q=q_k$, $V$ the set of places of $L^+$ above $\ell$ and $d=\sum_{w\in V}d_w$.
\begin{enumerate}
\item \emph{The global count.} By Lemma~\ref{sha-automorphy-unramified-tangents-count-local}(1), the $Z_{w,k}$, $w\in V$, are subgroups of $Z^1(G_w,W_k)$ that contain the maps $h\mapsto\mathrm{loc}_w(h)\cdot Y-Y$. The group $W_k$ is killed by $\ell^k$, has $q^9$ elements, and the kernel of the action of $\Gamma$ on it is open, by Lemma~\ref{sha-automorphy-unramified-tangents-count-adjoint}(1). So the global count gives that the set $Z_L$ of Lemma~\ref{sha-automorphy-unramified-tangents-count-adjoint}(4) is finite, with
\[q^9\prod_{w\in V}|Z_{w,k}|\le|Z_L|\,|\operatorname{Hom}_{\varepsilon_\Gamma}(W_k,\mathbf Z/\ell^k)|\prod_{w'}|W_k^{c'_{w'}}|\,q^{9|V|}.\]
\item \emph{The factors.} By Lemma~\ref{sha-automorphy-unramified-tangents-count-adjoint}(2), $|\operatorname{Hom}_{\varepsilon_\Gamma}(W_k,\mathbf Z/\ell^k)|=1$. For each $w'$, $r(c'_{w'})$ has square $1$ and sign $\delta_\Gamma(c'_{w'})=-1$, so $|W_k^{c'_{w'}}|\le q^3$ by Lemma~\ref{sha-automorphy-unramified-tangents-count-adjoint}(3); and the totally real field $L^+$ has $[L^+:\mathbf Q]=d$ infinite places. By Lemma~\ref{sha-automorphy-unramified-tangents-count-adjoint}(4), $|Z_L|\le|Z_k|\,q^8$; and $Z_k\subseteq Z_L$ is finite.
\item \emph{The local factors.} By Lemma~\ref{sha-automorphy-unramified-tangents-count-local}, $q^{9+3d_w}\le|Z_{w,k}|$ for every $w$ and $q^{10+3d_v}\le C|Z_{v,k}|$. With $E=\prod_{w\in V}q^{9+3d_w}=q^{9|V|+3d}$, this gives $qE\le C\prod_w|Z_{w,k}|$.
\item \emph{Conclusion.} By (1) to (3), $q^{10}E\le Cq^9\prod_w|Z_{w,k}|\le C|Z_k|\,q^8q^{3d}q^{9|V|}=C|Z_k|\,q^8E$, so $q^2\le C|Z_k|$.
\end{enumerate}
\end{proof}
% selLow (SelLow.lean): triOk, locHigh, locLow and zvCnt for (1) and (3) (Lemma sha-automorphy-unramified-tangents-count-local, with the private sel_cut); the hypothesis GWL with adOpen for (1); dualZero, archUp (with hF.card_inf for the number of infinite places) and normCnt for (2) (Lemma sha-automorphy-unramified-tangents-count-adjoint); the private sel_arith for (4). The Lean states selLow for any coefficient ring O with IsCoef l O and (l : O) != 0 and any frame with Ok C_0.

\begin{lemma}\label{sha-automorphy-unramified-tangents-count-tangent}
In the situation of Lemma~\ref{sha-automorphy-unramified-point}, let $\rho_\xi$ be the linked point with values in $\mathcal O_1$ given by $\xi$, $(r,(g_v))$ a continuous triangulation of $\rho_\xi$ over $\mathcal O_1$, $k\ge1$ and $z\in Z_k$ (Definition~\ref{sha-automorphy-unramified-tangents-count-cocycles}). Then there is an $\mathcal O_1$-linear map $\delta:A_3\to\mathcal O_k$ with $\delta(ab)=\xi(a)\delta(b)+\xi(b)\delta(a)$ for all $a,b\in A_3$, with $\delta(\psi^{\mathrm u}_{w,i}(h))=0$ for every place $w$ of $L^+$ above $\ell$, every $i$ and every $h\in I_w$, and with $\delta(\rho^{\mathrm u}(s)_{ij})=\bigl(z(s)\,\mathrm{mat}(r_k(s))\bigr)_{ij}$ for all $s\in\Gamma_0$ and all $i,j$.
\end{lemma}
\leanref{Development/AutomorphicArithmetic/AutomorphicVarieties/ModularityArgument/LinkedUnramifiedChains/CrystallineCardinalities/LiftOfBar.lean}{UpperTheory.V.Z.J.derOf}
\begin{proof}
\begin{enumerate}
\item \emph{A ring of first order.} Let $C=\mathcal O_1\oplus\mathcal O_k\epsilon$, with $(a+x\epsilon)(a'+x'\epsilon)=aa'+(ax'+a'x)\epsilon$. It is a local $\mathcal O_1$-algebra, finite over $\mathcal O_1$, hence complete and noetherian, with residue field $k_1$; let $\varphi:C\to\mathcal O_1$ be the projection $a+x\epsilon\mapsto a$. For $X\in W_k$, $1+X\epsilon$ is invertible with inverse $1-X\epsilon$, and for $y\in G_3(\mathcal O_1)$, read in $G_3(C)$, $y(1+X\epsilon,1)y^{-1}=(1+(\bar y\cdot X)\epsilon,1)$, with the action of Definition~\ref{sha-automorphy-unramified-near-completion-tangent-triangulation}; for $y=(g,-1)$ this uses $\theta(1+X\epsilon)=1-J_3{}^tXJ_3\epsilon$.
\item \emph{A triangulation over $C$.} As $z\in Z_k$, for each place $w$ there are $z'_w\in Z^{\mathrm{tri}}_{w,k}$ and $Y_w\in W_k$ with $z(\mathrm{loc}_wh)=\bar g_w^{-1}z'_w(h)\bar g_w+(\mathrm{loc}_w(h)\cdot Y_w-Y_w)$. Let $r'(s)=(1+z(s)\epsilon,1)\,r(s)$ and $g'_w=g_w(1+Y_w\epsilon)$. By (1) and the cocycle rule, $r'(s)r'(t)=(1+(z(s)+s\cdot z(t))\epsilon,1)\,r(st)=r'(st)$, so $r'$ is a homomorphism to $G_3(C)$ with sign $\delta_\Gamma$. By (1), $\mathrm{mat}(r(\mathrm{loc}_wh))(1-Y_w\epsilon)=(1-(\mathrm{loc}_w(h)\cdot Y_w)\epsilon)\,\mathrm{mat}(r(\mathrm{loc}_wh))$, so
\[g'_w\,\mathrm{mat}(r'(\mathrm{loc}_wh))\,g_w'^{-1}=g_w\bigl(1+\bar g_w^{-1}z'_w(h)\bar g_w\epsilon\bigr)\mathrm{mat}(r(\mathrm{loc}_wh))\,g_w^{-1}=(1+z'_w(h)\epsilon)\,t_w(h),\]
which is upper triangular with diagonal entries $(1+z'_w(h)_{ii}\epsilon)\psi_{w,i}(h)$. So $(r',(g'_w))$ is a triangulation over $C$; it is continuous, as $z$ is locally constant and $r$ is continuous.
\item \emph{A linked point over $C$.} $\mathrm{mat}(r'(c_\Gamma))=(1+z(c_\Gamma)\epsilon)\mathrm{mat}(r(c_\Gamma))=1$, as $z(c_\Gamma)=0$. The matrix $\mathrm{mat}(r'(g_0))-\mathrm{mat}(r'(g_0^*))$ is $\rho_\xi(g_0)-\rho_\xi(g_0^*)+\bigl(z(g_0)\mathrm{mat}(r_k(g_0))-z(g_0^*)\mathrm{mat}(r_k(g_0^*))\bigr)\epsilon$, whose entries off the diagonal vanish by condition (b) of Definition~\ref{sha-descent-link}(2) and the definition of $Z_k$. By Lemma~\ref{sha-automorphy-unramified-near-completion-tangent-normal}(2), $\rho'=\mathrm{mat}\circ r'$ on $\Gamma_0$, with the characters $\psi'_{w,i}(h)=(1+z'_w(h)_{ii}\epsilon)\psi_{w,i}(h)$, satisfies conditions (a) to (d). Its image under $\varphi$ is $(\rho_\xi,\psi)$, which lifts the residual one; as $\ker\varphi=\mathcal O_k\epsilon$ lies in the maximal ideal of $C$, $(\rho',\psi')$ has the same reduction, so it is a linked point with values in $C$ lifting the residual one.
\item \emph{The derivation.} By Lemma~\ref{sha-automorphy-realization-exists-linked}, there is an $\mathcal O_1$-algebra map $f:A_3\to C$ that carries $(\rho^{\mathrm u},\psi^{\mathrm u})$ to $(\rho',\psi')$. Then $\varphi\circ f$ and $\xi$ both carry $(\rho^{\mathrm u},\psi^{\mathrm u})$ to $(\rho_\xi,\psi)$, so $\varphi\circ f=\xi$ by the uniqueness in that lemma. Write $f(a)=\xi(a)+\delta(a)\epsilon$. Then $\delta$ is $\mathcal O_1$-linear, and $\delta(ab)=\xi(a)\delta(b)+\xi(b)\delta(a)$ as $f$ is multiplicative. For $h\in I_w$, $z'_w(h)_{ii}=0$, so $f(\psi^{\mathrm u}_{w,i}(h))=\psi'_{w,i}(h)=\psi_{w,i}(h)$ and $\delta(\psi^{\mathrm u}_{w,i}(h))=0$. For $s\in\Gamma_0$, $f(\rho^{\mathrm u}(s)_{ij})=\rho'(s)_{ij}=\rho_\xi(s)_{ij}+\bigl(z(s)\mathrm{mat}(r_k(s))\bigr)_{ij}\epsilon$.
\end{enumerate}
\end{proof}
% derOf (LiftOfBar.lean, 281 lines): C = TrivSqZeroExt O Q with Q = Ok l O k (sq_facts, co_sqz for IsCO), pi and mu with sq_exp for (1); the private box_def for (2), with upT_mul; pr_ofbox (Lemma sha-automorphy-unramified-near-completion-tangent-normal(2)) with hB.atc, the field diag of the linked point and z4, z5 (the conditions of ZN) for (3), and lift_of_bar for the reduction; S.univ for C and for O, with h5, for (4).

\begin{lemma}\label{sha-automorphy-unramified-tangents-count}
In the situation of Lemma~\ref{sha-automorphy-unramified-point}, let $\rho_\xi$ be the linked point with values in $\mathcal O_1$ given by $\xi$, and suppose that for some place $v$ of $L^+$ above $\ell$ there are $i'$, $g$, $t$, $u$ and $\beta$ as in the conclusion of Lemma~\ref{sha-automorphy-unramified-tangents-unit}. Then there are an integer $C$ and, for every $k\ge1$, a finite set $Z_k$ of cocycles of the group $\Gamma$ of the link with values in $M_3(\mathcal O_1/\ell^k)$, with $|\mathcal O_1/\ell^k|^2\le C|Z_k|$, and an injective map from $Z_k$ to the set $X_k$ of $\mathcal O_1$-linear maps $\delta:A_3\to\mathcal O_1/\ell^k$ with $\delta(ab)=\xi(a)\delta(b)+\xi(b)\delta(a)$ that vanish at $\psi^{\mathrm u}_{w,i}(s)$ for every place $w$ of $L^+$ above $\ell$, every $i$ and every $s\in I_w$.
\end{lemma}
\leanref{Development/AutomorphicArithmetic/AutomorphicVarieties/ModularityArgument/LinkedUnramifiedChains/CrystallineCardinalities/BoxUnit.lean}{UpperTheory.V.Z.J.tan}
\begin{proof}
Let $\psi_{w,i}$ be the characters of $\rho_\xi$.
\begin{enumerate}
\item \emph{Inputs.} By Lemma~\ref{sha-automorphy-unramified-tangents-count-inputs}, the link satisfies the local counts, and the global count for a family $(c'_w)$ indexed by the infinite places of $L^+$ with $c_w'^2=1$ and $\delta_\Gamma(c'_w)=-1$.
\item \emph{A triangulation.} By Lemma~\ref{sha-automorphy-unramified-tangents-count-triangulation}, applied to $g$ and $t$, there is a continuous triangulation $(r,(g_w))$ of $\rho_\xi$ over $\mathcal O_1$ with $t_v=t$. The linked point $\rho_\xi$ lifts the residual one, being the image of the universal one (Lemma~\ref{sha-automorphy-unramified-point}(1)). By the conclusion of Lemma~\ref{sha-automorphy-unramified-tangents-unit}, $t_v$ satisfies the hypotheses of Lemma~\ref{sha-automorphy-unramified-tangents-count-local}(2) with $i'$, $u$ and $\beta$, the cyclotomic character of the link being $\varepsilon_\Gamma$.
\item \emph{The count.} Let $Z_k$ be the set of Definition~\ref{sha-automorphy-unramified-tangents-count-cocycles} for $(r,(g_w))$, a set of cocycles of $\Gamma$ with values in $M_3(\mathcal O_1/\ell^k)$. By (1) and (2), Lemma~\ref{sha-automorphy-unramified-tangents-count-selmer} gives an integer $C$ with $Z_k$ finite and $|\mathcal O_1/\ell^k|^2\le C|Z_k|$ for every $k\ge1$.
\item \emph{The map.} For $z\in Z_k$, Lemma~\ref{sha-automorphy-unramified-tangents-count-tangent} gives $\delta_z\in X_k$ with $\delta_z(\rho^{\mathrm u}(s)_{ij})=\bigl(z(s)\,\mathrm{mat}(r_k(s))\bigr)_{ij}$ for $s\in\Gamma_0$.
\item \emph{Injectivity.} Let $z,z'\in Z_k$ with $\delta_z=\delta_{z'}$. Then $z(s)\,\mathrm{mat}(r_k(s))=z'(s)\,\mathrm{mat}(r_k(s))$ for $s\in\Gamma_0$, and $\mathrm{mat}(r_k(s))$ is invertible, so $z=z'$ on $\Gamma_0$. For $s\notin\Gamma_0$, $sc_\Gamma\in\Gamma_0$, and $z(sc_\Gamma)=z(s)+s\cdot z(c_\Gamma)=z(s)$ as $z(c_\Gamma)=0$, and likewise for $z'$. So $z=z'$.
\end{enumerate}
% Z.J.tan (BoxUnit.lean) before rankTwo. (1) is the hypotheses GWL and LocE, from gwlOf and leOf2 through linkOf in tanA_of (Lemma sha-automorphy-unramified-tangents-count-inputs), with F.Ok C_0 (R.ok) and l != 0 (R.T.lne) from the realization. (2) is the private box_pt and box_unit (Lemma sha-automorphy-unramified-tangents-count-triangulation), with lift_pt for the lifting. (3) is selLow (Lemma sha-automorphy-unramified-tangents-count-selmer), with Z_k = ZN B k. (4) is derOf (Lemma sha-automorphy-unramified-tangents-count-tangent), and (5) is the private zn_ext with hF.del_c. The text's injective map is the Lean's f, chosen with derOf. LinkedUnramifiedChains/CrystallineCardinalities has about 4,700 lines, unitOf included; Lemmas sha-automorphy-unramified-tangents-count-inputs (about 900 lines), -count-local (about 800 lines) and -count-adjoint (about 490 lines) and the lemmas before them write out these parts.
\end{proof}

\begin{lemma}\label{sha-automorphy-unramified-tangents}
In the situation of Lemma~\ref{sha-automorphy-unramified-point}, suppose that $r$ is irreducible, that $(\rho,\psi)$ has the prescribed inertial characters, that $r$ is potentially crystalline at the place $A_v$ above a place $v$ of $L^+$ above $\ell$, and that $\psi_{v,i}=\psi_{v,j}\chi_{\mathrm{cyc}}$ on all of $G_v$ for some $i,j$. Then there are two $\mathcal O_1$-linearly independent tangent vectors of $A_3$ at $\xi$, with values in $\mathcal O_1$, that kill the universal characters on the inertia groups $I_v$.
\end{lemma}
\leanref{Development/AutomorphicArithmetic/AutomorphicVarieties/ModularityArgument/LinkedUnramifiedChains/LinkedUnlOf.lean}{UpperTheory.V.J.TanA}
\leanref{Development/AutomorphicArithmetic/AutomorphicVarieties/ModularityArgument/LinkedUnramifiedChains/CrystallineCardinalities/TanA.lean}{UpperTheory.V.J.tanA}
\begin{proof}
Let $\rho_\xi$ be the linked point with values in $\mathcal O_1$ given by $\xi$, and $\varepsilon$ the cyclotomic character of the link.
\begin{enumerate}
\item \emph{A unit step at $v$.} Lemma~\ref{sha-automorphy-unramified-tangents-unit} gives $i'$, $g$, $t$, $u$ and $\beta$ at $v$.
\item \emph{Cocycles modulo $\ell^k$.} By (1), Lemma~\ref{sha-automorphy-unramified-tangents-count} gives an integer $C$ and, for every $k\ge1$, a finite set $Z_k$ of cocycles with $|\mathcal O_1/\ell^k|^2\le C|Z_k|$ and an injective map from $Z_k$ to the set $X_k$ of $\mathcal O_1$-linear maps $\delta:A_3\to\mathcal O_1/\ell^k$ with $\delta(ab)=\xi(a)\delta(b)+\xi(b)\delta(a)$ that vanish at $\psi^{\mathrm u}_{w,i}(s)$ for every place $w$ of $L^+$ above $\ell$, every $i$ and every $s\in I_w$.
\item \emph{Two tangent vectors.} The ring $A_3$ is noetherian, $\ell\ne0$ in $\mathcal O_1$, and by (2) every finite $X_k$ has at least $|Z_k|$ elements, so $|\mathcal O_1/\ell^k|^2\le C|X_k|$. Lemma~\ref{sha-automorphy-unramified-tangents-rank}, applied to $A_3$, $\xi$, the set $T$ of the values $\psi^{\mathrm u}_{w,i}(s)$ with $s\in I_w$, and $C$, gives the two tangent vectors.
\end{enumerate}
% tanA_of (TanAOf.lean): (1) is the hypothesis UnitOfP, proved as Z.J.unitOf from Loc.LocalPublishedTheorems (crOfPot), Lemma sha-automorphy-unramified-tangents-unit; the unit step is Z.J.UnitAt (UnitStep.lean). (2) is Lemma sha-automorphy-unramified-tangents-count. (2) and (3) are Z.J.tan (BoxUnit.lean), with the hypotheses GWL from gwlOf (GC.Hypotheses₂) and LocE from leOf2, for the Galois data of linkOf and Pr.degEq: in (2), box_pt (pr_box, Fr.flag) and box_unit give the triangular forms B, selLow (triOk, locHigh, locLow, zvCnt) gives C and Z_k = ZN B k, and derOf (LiftOfBar.lean) with zn_ext gives the injection; (3) is rankTwo (Lemma sha-automorphy-unramified-tangents-rank) with T the values of the characters on inertia.
% tanA = tanA_of from GC.Hypotheses₂ (global_compat_printed_of_literal_automorphy_n: UnramifiedOrthogonality, GreenbergWiles) and Loc.LocalPublishedTheorems (locPrints_of_litN), through Pr.degEq, Z.J.unitOf, linkOf, Z.J.tan, gwlOf and leOf2. Size: LinkedUnramifiedChains/CrystallineCardinalities, about 4,700 lines.
\end{proof}

\begin{theorem}\label{sha-residually-reducible-automorphy}
Assume that Galois representations exist, in the sense of Definition~\ref{gal-galois-representations}, relative to the maps satisfying the laws of Definition~\ref{gal-llc-laws}.
Let $(\ell,k,\psi,K)$ be an Eisenstein datum, with deformation data $(\eta,c,\Psi)$ (Definition~\ref{sha-ordinary-deformation}), let $\iota:\bar{\mathbf Q}_\ell\cong\mathbf C$, and let $\mathfrak B$ be a place above $\ell$.
Say that a continuous representation $r:\operatorname{Gal}(\bar{\mathbf Q}/F)\to GL_n(\bar{\mathbf Q}_\ell)$, for a CM field $F$, is \emph{windowed} if:
($\alpha$) at every place $v$ of $F$ above $\ell$ at which $r$ is triangular on the decomposition groups with characters $\kappa_1,\dots,\kappa_n$ such that $\kappa_i\ne\kappa_j\chi_{\mathrm{cyc}}$ for all $i,j$, it is triangular with characters such that $\kappa_i^{m}\ne(\kappa_j\chi_{\mathrm{cyc}})^{m}$ for all $i,j$ and $m\ge1$; and
($\beta$) for every regular algebraic cuspidal automorphic representation $\pi$ of $GL_n(\mathbf A_F)$ attached to $r$ through $\iota$ there is $s\in\mathbf R$ such that at every place $v\nmid\ell$, and at every place $v\mid\ell$ satisfying the second condition of ($\alpha$), $\pi_v$ has supercuspidal support consisting of characters $\xi$ with $q_v^{s-1/2}<|\xi(\varpi_v)|<q_v^{s+1/2}$.
Then there is $g_1\in\operatorname{Gal}(\bar{\mathbf Q}/K)$ with $\Psi(g_1)^2-1\in\mathcal O^\times$ such that the following hold for the ordinary points relative to $g_1$ and $\mathfrak B$.
\begin{enumerate}
\item (Finiteness.) There is a subgroup $U$ of finite index such that for every complete discrete valuation ring $A$ over $\mathcal O$ with finite residue field in which $\mathfrak m_{\mathcal O}$ maps to $0$, and every ordinary point $r$ with values in $A$, the characteristic polynomial of $r(g)$ equals that of $r_0(\eta,\Psi)(g)$ for all $g\in U\cap\ker\eta$.
\item (Stable points.) Let $A\subset\bar{\mathbf Q}_\ell$ be a discrete valuation ring finite over $\mathcal O$ and $r$ an ordinary point with values in $A$ whose restriction to $\ker\eta$ is irreducible over $\bar{\mathbf Q}_\ell$. Let $\rho:\operatorname{Gal}(\bar{\mathbf Q}/K)\to GL_3(\bar{\mathbf Q}_\ell)$ be its restriction. Then $\rho\otimes\chi_{\mathrm{cyc}}^{-1}$ is attached through $\iota$ to a regular algebraic cuspidal automorphic representation $\pi$ of $GL_3(\mathbf A_K)$ with $\pi^\vee\cong\pi^{c}$, and it is windowed.
\item (Endoscopic points.) Let $A$ be as in (ii) and let $r$ be an ordinary point with values in $A$. Suppose that on $\ker\eta$ the entries $(1,2),(2,1),(2,3),(3,2)$ of $r$ vanish, the middle entry is $1$, and the corner block formed by rows and columns $1$ and $3$ is irreducible. Let $W:\operatorname{Gal}(\bar{\mathbf Q}/\mathbf Q)\to GL_2(\bar{\mathbf Q}_\ell)$ agree with the corner block on $\ker\eta$ and satisfy $\operatorname{tr}W(c)=0$. Then $W\otimes\chi_{\mathrm{cyc}}^{-k}$ is attached through $\iota$ to a newform of weight $2k+1$, and the restriction of $W$ to every imaginary quadratic field on which it is irreducible is windowed.
\end{enumerate}
\end{theorem}
\begin{proof}
The proof is assembled from four propositions about ordinary points over abelian CM fields and a descent step, all proved from the Galois representations of Definition~\ref{gal-galois-representations} and from published results.
The published results are of three kinds: the representation theory of $GL_n$ over local fields (the Bernstein-Zelevinsky and Langlands classifications, the local converse theorem, unrefined minimal types and depth zero); automorphic forms on $GL_n$ and on definite quaternion and unitary groups (rapid decay and finiteness of cusp forms, the cuspidal spectrum, the theorems of Jacquet-Shalika and Gelfand-Kazhdan, cyclic and quadratic base change and descent, base change between definite unitary groups and $GL_n$, local-global compatibility, Clozel's purity and cohomological criterion, strong multiplicity one, newforms, Galois representations of Hilbert modular forms); and Galois cohomology, $p$-adic Hodge theory and complex multiplication (global class field theory, Hecke characters, Berger's theorem, Poitou-Tate duality, local and global Euler characteristics, the Greenberg-Wiles formula, the Bloch-Kato local conditions, Deligne-Ribet and Siegel-Klingen).
None of them is an automorphy lifting theorem: the automorphy of ordinary points is proved inside the argument, by the Hecke ideal and the patching bound of Propositions~\ref{sha-hecke-ideal} and~\ref{sha-patching-bound}.

\emph{Links and seeds.} Links, linked points, the residual conditions and seed data are those of Definitions~\ref{sha-descent-link}, \ref{sha-descent-residual} and~\ref{sha-descent-seed}, for the descent datum $(\eta,\Psi,c,\mathfrak B,\ell N|d_K|,\operatorname{Gal}(\bar{\mathbf Q}/\mathbf Q))$ of Lemma~\ref{sha-descent-eisenstein}, where $N$ is the modulus of $\psi$.
They concern an abelian CM field $L\supset K$ with maximal totally real subfield $L^+$ and $\operatorname{Gal}(\bar{\mathbf Q}/L)=\ker\eta\cap\operatorname{Gal}(\bar{\mathbf Q}/L^+)$.
A linked point of rank $n$, composed with the map $p$ of its link, gives a representation of $\operatorname{Gal}(\bar{\mathbf Q}/L)$ in $M_n(A)$.
In (c) and (d), a link that satisfies the residual conditions does so with a constant $C_0$ for which it is also large.
The quadratic field of a seed datum is $K$, since its Galois group is $\ker\eta$, and a seed datum carries an integer $j$ and characters $\chi'$ and $\theta$.
Below, $m$ denotes an integer with $m\equiv j\pmod{\ell-1}$, which can be taken larger than any given bound.
\begin{enumerate}
\item \emph{Seeds of rank two.} For a seed datum and a suitable algebraic Hecke character $\kappa$ of $K$, of infinity type $(-1,0)$ and with $\kappa\kappa^{c}=\chi_{\mathrm{cyc}}^{-1}$, there is a continuous $\sigma:\operatorname{Gal}(\bar{\mathbf Q}/\mathbf Q)\to GL_2(\bar{\mathbf Q}_\ell)$ with $\det\sigma=\eta\theta$, whose reduction has constituents $\chi_{\mathrm{cyc}}^{j}\chi'\eta$ and $(\chi_{\mathrm{cyc}}^{j}\chi')^{-1}$, and whose restriction to $L$ is irreducible, polarized, of inertial type $(\chi',\chi'^{-1})$ away from $\ell$, ordinary at $\ell$ with weights $(m,-m)$, and, after the twist by the restriction of $\kappa$, attached to a conjugate self-dual cuspidal automorphic representation. It is constructed from modular forms of large weight congruent to Eisenstein series, followed by base change to $L$.
\item \emph{Seeds of rank three.} Let $M\ge1$ with $2\ell\mid M$ be such that $L$ is unramified over $L^+$ at every finite place, the places of $L^+$ dividing $M$ split in $L$, and $\chi'$ is trivial on the inertia groups of $\operatorname{Gal}(\bar{\mathbf Q}/L)$ at the places not dividing $M$; in (c) and (d), $M$ is the integer of the link. For a seed datum with $\chi'\ne\eta$, there is $r:\operatorname{Gal}(\bar{\mathbf Q}/L)\to GL_3(\bar{\mathbf Q}_\ell)$, irreducible on every open subgroup, with reduction of constituents $\chi_{\mathrm{cyc}}^{j}\chi'$, $1$, $(\chi_{\mathrm{cyc}}^{j}\chi')^{-1}$, polarized, of inertial type $(\chi',1,\chi'^{-1})$ away from $\ell$, ordinary at $\ell$ with weights $(m,0,-m)$, and such that $r\otimes\chi_{\mathrm{cyc}}^{-1}$ is attached to a conjugate self-dual cuspidal automorphic representation and is windowed. It is the restriction to $L$ of the adjoint of a modular form of large weight congruent to an Eisenstein series, made automorphic by an adjoint lift between definite unitary groups.
\item \emph{Automorphy of linked points.} Let a link to $L$ satisfy the residual conditions, and let a seed datum for $L$ exist. Then the linked points of rank three with values in complete discrete valuation rings of characteristic $\ell$ with finite residue field, whose characters $\psi_{v,i}$ are trivial on the elements of $I_v$ at which $\bar\Psi$ is trivial, have the characteristic polynomials of the residual representation; the linked points of rank three that are irreducible become automorphic after the twist by $\chi_{\mathrm{cyc}}^{-1}$; and there is an algebraic character by whose twist the irreducible linked points of rank two become automorphic. This follows from the automorphy of points described below, applied to the components of the deformation rings that contain the seeds of (a) and (b).
\item \emph{Linked unramified points.} Under the hypotheses of (c), at every place above $\ell$ at which it is potentially crystalline, the diagonal characters $\kappa_1,\kappa_2,\kappa_3$ of a linked point of rank three that is irreducible on every open subgroup satisfy $\kappa_i\ne\kappa_j\chi_{\mathrm{cyc}}$ on some element of the decomposition group.
\end{enumerate}

\emph{Automorphy of points.} The residual representation of a link of rank $n$ over a CM field $L$ without nontrivial $\ell$-th roots of unity, with the places of the link and a place $v_0$, is a residual datum in Situation~\ref{sha-hecke-situation}, with its universal ordinary determinant ring $B_n$ and Iwasawa algebra $\Lambda$. Proposition~\ref{sha-hecke-ideal} gives the Hecke ideal $J\subset B_n$, with Taylor-Wiles data, and Proposition~\ref{sha-patching-bound} shows that $J$ is contained in every prime of $B_n$ below a nice arc of type I on the zero locus of $J$, and in every stable or thick prime below the last of a chain of nice arcs of type E joined by thick primes, starting from a stable prime that contains $J$. So the irreducible components of $\operatorname{Spec}B_n$ through nice arcs of $\operatorname{Spec}B_n/J$ lie in $\operatorname{Spec}B_n/J$, whose irreducible arithmetic points are automorphic by Proposition~\ref{sha-hecke-ideal}(5).

\emph{Proof of (a).}
Let $(j,\chi',\theta)$ be a seed datum for $L$, let $\kappa$ be the Hecke character of (a), and let $\varphi=\chi'^2\eta$, so that $\varphi\ne1$ and $\varphi(c)=-1$.
\begin{enumerate}
\item \emph{Eisenstein congruences.} There is $k_0$ such that for every $k\ge k_0$ with $\varphi(c)=(-1)^k$ there is a continuous irreducible $\rho:\operatorname{Gal}(\bar{\mathbf Q}/\mathbf Q)\to GL_2(\bar{\mathbf Q}_\ell)$ attached through $\iota$ to a regular algebraic cuspidal representation of $GL_2(\mathbf A_{\mathbf Q})$ of weight $k$, with $\det\rho=\chi_{\mathrm{cyc}}^{k-1}\varphi\theta$, reduction with constituents $1$ and $\chi_{\mathrm{cyc}}^{k-1}\varphi$, inertial type $(1,\varphi\theta)$ away from $\ell$, and ordinary at $\ell$ with Hodge-Tate weights $(k-1,0)$ and diagonal characters $(\chi_{\mathrm{cyc}}^{k-1}\varphi,1)$. This is Proposition~\ref{sha-automorphy-eisenstein-galois}, applied to $\varphi$ and $\theta$ with the level $N$ and the prime $w$ of the splitting conditions of the seed datum: $\theta$ is a nontrivial character of $\operatorname{Gal}(\mathbf Q(\mu_w)/\mathbf Q)$ of order $\ell$, and $\chi'$ and $\eta$, hence $\varphi$, are trivial on $\operatorname{Gal}(\bar{\mathbf Q}/\mathbf Q(\mu_N))$ and on the elements acting on $\mu_N$ by $\zeta\mapsto\zeta^w$.
\item Given a bound $m_0$, take $t\ge m_0+k_0+1$, $m=j+(\ell-1)t\ge m_0$ and $k=2m+1$, so that $\varphi(c)=(-1)^k$, and let $\rho$ be as in (1). Let $\Xi=\chi_{\mathrm{cyc}}^m\chi'$ and $\sigma=\rho\otimes\Xi^{-1}$. Then $\det\sigma=\eta\theta$. Since $m\equiv j\pmod{\ell-1}$, $\Xi$ is congruent to $\chi_{\mathrm{cyc}}^j\chi'$, so the reduction of $\sigma$ has constituents $\chi_{\mathrm{cyc}}^j\chi'\eta$ and $(\chi_{\mathrm{cyc}}^j\chi')^{-1}$.
\item The restriction of $\rho$ to an open normal subgroup $H$ is irreducible unless $\chi_{\mathrm{cyc}}^{k-1}\varphi$ is congruent to $1$ on $H$: a reducible restriction of the irreducible $\rho$ is a sum of characters congruent to the residual constituents, because $H$ contains an element at which $\rho$ has the eigenvalue $1$ and an eigenvalue that is not a root of unity. On $\operatorname{Gal}(\bar{\mathbf Q}/L)$ the character $\chi_{\mathrm{cyc}}^{k-1}\varphi=\Xi^2\eta$ is congruent to $(\chi_{\mathrm{cyc}}^j\chi')^2$, which is not congruent to $1$ by the definition of a seed datum, so the restriction of $\sigma$ to $L$ is irreducible.
\item Since $\det\sigma$ is trivial on $\operatorname{Gal}(\bar{\mathbf Q}/L)$ and $\det\sigma(c)=\eta(c)\theta(c)=-1$ ($\theta$ has odd order), there is a symmetric invertible $P$ with ${}^t\sigma(s)P\sigma(csc)=P$ for $s\in\operatorname{Gal}(\bar{\mathbf Q}/L)$. Since $\eta$ and $\theta$ are trivial on $\operatorname{Gal}(\bar{\mathbf Q}/L)$ and $\chi_{\mathrm{cyc}}$ is trivial on inertia away from $\ell$, the restriction of $\sigma$ to $L$ has inertial type $(\chi',\chi'^{-1})$ away from $\ell$, and it is ordinary at $\ell$ with weights $(2m-m,-m)=(m,-m)$.
\item Since $\Xi$ is algebraic, $\sigma$ is attached to a regular algebraic cuspidal representation of $GL_2(\mathbf A_{\mathbf Q})$. Base change to $K$ and then along the abelian extension $L/K$ \cite{AC89}, with the twist by $\kappa$, makes the restriction of $\sigma$ to $L$, twisted by the restriction of $\kappa$, cuspidal automorphic, since it is irreducible by (3); it is conjugate self-dual by the pairing of (4). The windowed conditions follow along the same base changes, since $\sigma$ is ordinary at $\ell$ with distinct weights.
\end{enumerate}

\emph{Proof of (b).}
Let $(j,\chi',\theta)$ be a seed datum with $\chi'\ne\eta$, and let $\varphi=\chi'\eta\ne1$.
\begin{enumerate}
\item Proposition~\ref{sha-automorphy-eisenstein-tate-free}, applied to $\varphi$ and $\theta$ with the level and the prime $w$ of the splitting conditions of the seed datum, as in step (1) of (a), gives $k_0$ and, for each $k\ge k_0$ with $\varphi(c)=(-1)^k$, a representation $\rho$ with the properties of step (1) of (a) for $\varphi$, whose eigenvalues at the elements of the decomposition groups at primes $p\ne\ell$ acting on the residue field by $x\mapsto x^p$ satisfy in addition the two inequalities of (6) of that proposition, and such that $\operatorname{ad}^0\rho$ is Tate-free (Definition~\ref{sha-automorphy-tate-free}).
\item Given a bound $m_0$, take $t\ge m_0+k_0+2$, $m=j+(\ell-1)t\ge m_0$ and $k=m+1$; then $\varphi(c)=(-1)^{j+1}=(-1)^{m+1}$ because $\ell-1$ is even. Let $\Phi=\chi_{\mathrm{cyc}}^m\chi'\eta$, which is congruent to $\chi_{\mathrm{cyc}}^j\chi'$ on $\operatorname{Gal}(\bar{\mathbf Q}/L)$.
\item By the criterion of step (3) of (a), $\rho$ is irreducible on every open subgroup, and so is $\operatorname{ad}^0\rho$. Let $r$ be the restriction of $\operatorname{ad}^0\rho$ to $L$. Its reduction has constituents $\Phi,1,\Phi^{-1}$, congruent to $\chi_{\mathrm{cyc}}^j\chi',1,(\chi_{\mathrm{cyc}}^j\chi')^{-1}$; a symmetric pairing on $\operatorname{ad}^0$ polarizes it; it has inertial type $(\chi',1,\chi'^{-1})$ away from $\ell$ and is ordinary at $\ell$ with weights $(m,0,-m)$.
\item \emph{Twist.} Let $M$ be as in (b), and let $\Theta=\chi'\eta\theta$, a character with open kernel, so that $\det\rho=\chi_{\mathrm{cyc}}^m\Theta$ and $\Theta=\chi'$ on $\operatorname{Gal}(\bar{\mathbf Q}/L)$. Since $\theta$ has odd order and $c^2=1$, $\theta(c)=1$, so $\Theta(c)=\varphi(c)=(-1)^{m+1}$.
Lemma~\ref{sha-automorphy-algebraic-twist}, applied to $\Theta^{-1}$ with $a=-(m+1)$, gives an algebraic $\kappa$ with $\kappa(s)\kappa(csc^{-1})=\chi_{\mathrm{cyc}}(s)^{-m-1}\Theta(s)^{-1}$ for $s\in\operatorname{Gal}(\bar{\mathbf Q}/L)$, trivial on the inertia groups at the places not dividing $M$, and with $\kappa^e=\chi_{\mathrm{cyc}}^{be}$ on the inertia group at each place above $\ell$, for some $b\in\mathbf Z$ and $e\ge1$.
Let $r'=\rho|_L\otimes\kappa$. It has the following properties.
\begin{enumerate}
\item[(i)] $\det\rho(s)\kappa(s)\kappa(csc^{-1})=\chi_{\mathrm{cyc}}(s)^{-1}$ for $s\in\operatorname{Gal}(\bar{\mathbf Q}/L)$. With $J=\begin{pmatrix}0&1\\-1&0\end{pmatrix}$ and $P=J\rho(c)^{-1}$, the identity ${}^tXJX=\det(X)J$ for $2\times2$ matrices then gives ${}^tr'(s)Pr'(csc^{-1})=\chi_{\mathrm{cyc}}(s)^{-1}P$. The same holds, with $P=J\rho(c')^{-1}$, for every complex conjugation $c'$, since $c'c^{-1}\in\operatorname{Gal}(\bar{\mathbf Q}/L)$.
\item[(ii)] The restriction of $\rho$ to $K$ is automorphic by quadratic base change, and its restriction to $L$ by base change along the abelian extension $L/K$ \cite{AC89}, since $\rho$ is irreducible on every open subgroup. The twist by the algebraic character $\kappa$ is again automorphic, and with the pairing of (i), $r'$ is attached through $\iota$ to a regular algebraic cuspidal representation $\Pi$ of $GL_2(\mathbf A_L)$ with $\Pi^\vee\cong\Pi^c$.
\item[(iii)] $r'$ is windowed: the windowed condition holds for the restriction of $\rho$ to $K$, by the plain local components of step (1) of (a), and it passes along the base change to $L$ and the twist by $\kappa$; condition ($\alpha$) holds because $r'$ is ordinary at the places above $\ell$ with distinct weights.
\item[(iv)] $\operatorname{ad}^0r'=(\operatorname{ad}^0\rho)|_L$ is Tate-free. Indeed, if $s$ is an element of degree $j$ at a place of $\mathbf Q$ as in Definition~\ref{sha-automorphy-tate-free}, some power $s^{m'}$ lies in $\operatorname{Gal}(\bar{\mathbf Q}/L)$, and $s^{m'f}$, with $f$ the residue degree of the place of $L$, is an element of degree $jm'$ at that place; its eigenvalues are the $m'f$-th powers of those of $s$, so the condition for $s$, applied with $m'f$ times a given exponent, gives it for $s^{m'f}$.
\item[(v)] $r'$ is irreducible on every open subgroup, and it is trivial on the inertia groups at the places not dividing $M$: there, $\rho$ is diagonalizable with diagonal characters $1$ and $\Theta$ by step (1) of (a), and $\Theta=\chi'$ and $\kappa$ are trivial.
\item[(vi)] At a place $A$ above $\ell$, $\rho$ is triangular on $D_A$ with diagonal characters $\kappa_0,\kappa_1$ equal to $\chi_{\mathrm{cyc}}^m$ and $1$ on the inertia group, so $r'$ is triangular with diagonal characters $\kappa_0\kappa$ and $\kappa_1\kappa$, which satisfy $(\kappa_0\kappa)^e=\chi_{\mathrm{cyc}}^{(b+m)e}$ and $(\kappa_1\kappa)^e=\chi_{\mathrm{cyc}}^{be}$ on the inertia group.
\end{enumerate}
\item \emph{Adjoint lift.} By (i) to (vi), with $m\ge2$, $r'$ is an adjoint input of level $M$ and weight $m$ (Definition~\ref{sha-automorphy-adjoint-input}): (i) is (c) there, (ii) and (iii) are (a), (iv) is (b), (v) gives (d) and (f), (g) holds with the characters $\kappa$ and $\Theta\kappa$, since $\rho$ is diagonalizable with diagonal characters $1$ and $\Theta$ on the inertia groups away from $\ell$ by step (1) of (a), and (vi) is (h).
Proposition~\ref{sha-automorphy-adjoint-automorphy} then shows that $\operatorname{ad}^0r'\otimes\chi_{\mathrm{cyc}}^{-1}$ is attached through $\iota$ to a regular algebraic cuspidal representation $\pi$ of $GL_3(\mathbf A_L)$ with $\pi^\vee\cong\pi^c$, and that every such $\pi$ satisfies the bounds of condition ($\beta$) of the windowed condition at every finite place.
\item Since $\operatorname{ad}^0(\rho|_L\otimes\kappa)=(\operatorname{ad}^0\rho)|_L=r$, (5) makes $r\otimes\chi_{\mathrm{cyc}}^{-1}$ attached to a regular algebraic cuspidal representation $\pi$ of $GL_3(\mathbf A_L)$ with $\pi^\vee\cong\pi^c$, and windowed: condition ($\alpha$) holds because $r$ is ordinary at $\ell$ with the distinct weights $(m,0,-m)$.
\end{enumerate}
% (b)(4) is Z.Sym.W.seedInW_alt (SeedInW, whose conclusion is the structure AdInW): Z.Sym.TwEx (Pl.twEx) for kappa; pair_of with gram_two and cc_diff for (i); Sd.up_quad_of_printed_results, Z.Bc.upSol_of_alt, Sd.twist_conjugate_selfDual and GalN.V.aut_csd_zero_of_complex_conj_pair_automorphic for (ii); awReal, awUp, awTw and potOrd for (iii); adjoint_zero_twist and GalN.V.L.wup (exists_pow_lift, isFrob_pow, noLink_pow) for (iv); shapeFacts and resFacts for (v) and (vi).
% (b)(5) is AdAutW (Proposition sha-automorphy-adjoint-automorphy), proved by adjoint_automorphy_w_of_pure_weight_adjoint_lift; SeedInW concludes AdInW (Definition sha-automorphy-adjoint-input). (b)(6) is mkCsdW' with AllW.witA and potOrd in adjoint_lift_from_eisenstein_of_propagation.

\emph{Proof of (c).}
Let a link to $L$ satisfy the residual conditions with constant $C_0$, and let $(j,\chi',\theta)$ be a seed datum for $L$ with quadratic field $K$; we are in Situation~\ref{sha-automorphy-link-situation}.
Let $d=[L^+:\mathbf Q]$.
\begin{enumerate}
\item \emph{Twists.} In rank three the twist is $\kappa_3=\chi_{\mathrm{cyc}}^{-1}$. In rank two it is the restriction $\kappa_2$ to $\operatorname{Gal}(\bar{\mathbf Q}/L)$ of an algebraic Hecke character $\kappa$ of $K$ of infinity type $(-1,0)$ with $\kappa\kappa^c=\chi_{\mathrm{cyc}}^{-1}$, as in (a); it is algebraic. Both are admissible twists, by Lemma~\ref{sha-automorphy-realization-twists}.
\item \emph{Realizations.} After enlarging $\mathcal O$ to a finite extension $\mathcal O_1$, for $n\in\{2,3\}$ and every integer $p_0\ne0$ there is a realization of rank $n$ of the link, twisted by $\kappa_n$, whose place $v_0$ has residue characteristic prime to $p_0$ (Definition~\ref{sha-automorphy-realization} and Lemma~\ref{sha-automorphy-realization-exists}). It consists of the setting of the Hecke ideal with these places, a residual datum of rank $n$ obtained from the residual representation of the link by the twist, its universal ordinary determinant ring $B_n$, the universal ring $A_n$ of linked points of rank $n$, and the map from the Galois group of the residual datum to that of the link through which a linked point gives a point of $B_n$.
\item \emph{Seeds.} The seed of (a), respectively (b), with $m$ large, gives a point of $B_n$ with values in an algebraically closed field at which the representation is absolutely irreducible and which lies on the zero locus of every ideal with the property of automorphic points of the Hecke ideal (Lemma~\ref{sha-automorphy-realization-seed}): after the twist it is conjugate self-dual automorphic, unramified at $v_0$, and ordinary with characters of weights $(m,-m)$, respectively $(m,0,-m)$, which are distinct and not related by $\chi_{\mathrm{cyc}}$.
\item \emph{Automorphy of points.} For each realization it gives an ideal $J_n$ of $B_n$ with: $B_n/J_n$ finite over $\Lambda$; every irreducible component of $\operatorname{Spec}B_n/J_n$ of dimension $1+nd$; the determinant over $B_n/J_n$ trivial on the inertia group at $v_0$; the patching bound; the seed point of (3) on its zero locus; and the classicality of its irreducible arithmetic points (Lemma~\ref{sha-automorphy-realization-ideal}).
\item \emph{Components in rank two.} The following argument is Lemma~\ref{sha-automorphy-realization-components}. Every prime of $A_2^\circ$ at which the universal linked point is not triangular contains the image of $J_2$. The spectrum of $A_2^\circ$ is connected; through each prime at which the point is not triangular, its irreducible components are joined by chains in which consecutive components meet in dimension at least $2d+1$; every closed subset of large dimension contains a nice prime, of dimension one, containing $\ell$, at which the point is not triangular and its ordinary characters are distinct; the patching bound carries the zero locus of $J_2$ from a nice prime to every component through it; and the seed of (3) lies on a component of dimension $1+2d$ in that zero locus. These spread the zero locus of $J_2$ to every component.
\item \emph{Components in rank three.} The same argument applies to $A_3^\circ$, with the variant of the patching bound used to pass through primes at which the point is endoscopic, joined by thick primes, and with (5) for the components on which the point becomes reducible. It also gives an ideal $I$ of $B$, with quotient finite over $\Lambda$, contained in every prime of $A_3^\circ$ at which the point is not triangular; this gives the first statement of (c), for the linked points of rank three with values in complete discrete valuation rings of characteristic $\ell$ with finite residue field (Lemma~\ref{sha-automorphy-realization-components}(2), applies, since a local map from $\mathcal O_1$ to such a ring kills a uniformizer of $\mathcal O_1$, some power of which is a multiple of $\ell$, and so factors through $k_1$).
\item \emph{Automorphy.} An irreducible linked point of rank $n$ with values in a discrete valuation ring finite over $\mathcal O$ therefore lies on the zero locus of $J_n$, and it is arithmetic of the prescribed inertial characters. By classicality, its twist by $\kappa_n$ is conjugate self-dual automorphic and windowed at the places not above $v_0$ (Lemma~\ref{sha-automorphy-realization-away}). Two realizations whose places $v_0$ have different residue characteristics, given by Lemma~\ref{sha-automorphy-realization-exists} with $p_0=1$ and then with $p_0$ the residue characteristic of the first, make it windowed at every place.
% (7): twoRuns_alt (e_2, e_3 with p_0 = 1 and p_0 = R.s.p), autPtO_alt for each, glueN (witA_of with ImaginaryQuadraticEtaCharacter.winOf and winUnl), v0_ne; finPtDown descends FinPt from O_1 to O.
\end{enumerate}

\emph{Proof of (d).}
Let $\rho$ be a linked point of rank three with values in a discrete valuation ring finite over $\mathcal O$, irreducible on every open subgroup and potentially crystalline at a place $v$ above $\ell$, and suppose that $\kappa_i=\kappa_j\chi_{\mathrm{cyc}}$ on the whole decomposition group at $v$.
\begin{enumerate}
\item By Lemma~\ref{sha-automorphy-unramified-realization}, there is a realization of rank three twisted by $\chi_{\mathrm{cyc}}^{-1}$, with coefficients $\mathcal O_1$ containing those of $\rho$, as in (c), for which $\Lambda$ is a power series ring over $\mathcal O_1$, $J_3$ has the properties of the automorphy of points, $B_3$ maps onto the closure $B$ of the $\mathcal O_1$-subalgebra of $A_3$ generated by the values of the characters $\psi^{\mathrm u}_{v,i}$ and by the products of matrix entries of the universal linked point of total weight $0$, where the entry $(i,j)$ has weight $w_i-w_j$ with $w=(1,0,-1)$, every off-triangular prime of the $\mathcal O_1$-subalgebra $A_3^\circ$ generated by $B$ and the entries contains the image of $J_3$, and $J_3$ satisfies the multiplicity bound of Definition~\ref{sha-automorphy-unramified-multiplicity}: on a Hecke module with support in the zero locus of $J_3$ there is an idempotent whose multiplicity at each irreducible arithmetic point is at least $1$ and at most its multiplicity at the points of distinct type of the same weight.
\item The point $\rho$ gives an $\mathcal O_1$-point $\xi$ of $A_3$, irreducible, which is an arithmetic point of $B_3$ on the zero locus of $J_3$ (Lemma~\ref{sha-automorphy-unramified-point}).
\item At an irreducible arithmetic point of $B_3/J_3$, every tangent vector of $B_3$ that kills $J_3$ and the inertial characters is $0$; this comes from the density of arithmetic points, purity and the multiplicity bound (Lemma~\ref{sha-automorphy-unramified-tangent}).
\item The ideal $J_3$ vanishes near $\xi$: for each $a\in J_3$ there is $s\in A_3^\circ$ with $\xi(s)\ne0$ and $sa=0$ in $A_3^\circ$ (Lemma~\ref{sha-automorphy-unramified-near}, with Lemma~\ref{sha-automorphy-unramified-point}(2)). With (3), any two tangent vectors of $A_3$ at $\xi$ that kill the inertial characters are linearly dependent.
Indeed, let $E:A_3\to\mathcal O_1$ be such a tangent vector, an $\mathcal O_1$-linear map with $E(xy)=\xi(x)E(y)+\xi(y)E(x)$. For $a\in J_3$ and $s$ as above, $\xi(s)\xi(a)=0$ and $\xi(s)E(a)+\xi(a)E(s)=E(sa)=0$, so $\xi(a)=E(a)=0$. So the composite of $E$ with $B_3\to A_3$ is a tangent vector of $B_3$ at an irreducible arithmetic point on the zero locus of $J_3$ that kills $J_3$ and the inertial characters, which the map carries to those of $A_3$ times constants; by (3) it is $0$, and $E$ vanishes on $B$.
Since $\xi$ is irreducible, an entry of odd weight, a lower entry and an upper entry of the universal linked point are not in the kernel of $\xi$; combining them gives products $p$ and $q$ of entries, of weights $1$ and $-1$, with $\xi(p)\ne0$ and $\xi(q)\ne0$.
If $E$ vanishes on $B$ and at $p$, then $E=0$: since $pq$ has weight $0$, $E(q)=0$; for a product $x$ of weight $1$, $xq$ has weight $0$, so $E(x)=0$, and in the same way $E$ vanishes on the products of weights $-1$, $2$ and $-2$, hence on every entry and every inertial character; by the universal property of $A_3$, a tangent vector at $\xi$ is determined by its values on these.
For two such tangent vectors $D,D'$, the tangent vector $D'(p)D-D(p)D'$ vanishes on $B$ and at $p$, so it is $0$; if $D,D'$ were linearly independent, then $D(p)=D'(p)=0$ and $D=0$.
\item The hypotheses on $\rho$ at $v$ give two linearly independent such tangent vectors (Lemma~\ref{sha-automorphy-unramified-tangents}). This contradiction proves (d).
\end{enumerate}
% The steps of (b), (c) and (d) that are not yet written out are the open steps in the proofs of Proposition sha-automorphy-lift-comparison, Lemma sha-automorphy-lift-forms-facts and Lemmas sha-automorphy-realization-exists to sha-automorphy-unramified-tangents.
% In (b): W.aboveWeight is written out as the proof of Lemma sha-automorphy-unitary-two, branch1W_alt and branch_two_w_statement as that of Lemma sha-automorphy-adjoint-branches, and adjoint_lift_w_of_literal_automorphy_spec_out as that of Lemma sha-automorphy-adjoint-lift, from W.SpecOutS (Proposition sha-automorphy-lift-comparison) and outW, Al.outRep_alt, Al.inVec_alt, Al.refine (Lemma sha-automorphy-lift-forms-facts).
% (d)(4) after the first sentence is J.endJ: the derivation through phi, UnrAt, then J.tor (exists_mon from the TI condition, der_zero with mon_mul, and J.det from the universal property of the Space); the first sentence is Red2Alt (red2_alt).
% (d)(5) is J.TanA. It is a hypothesis of UnitaryRankTwoFlagsAlt and of linkedUnlE_alt, and linkedUnlE_holds_alt discharges it by J.tanA from GC.Hypotheses₂ (global_compat_printed_of_literal_automorphy_n: the published Galois cohomology inputs of the literature hypothesis, with UnramifiedOrthogonality and GreenbergWiles) and Loc.LocalPublishedTheorems (locPrints_of_litN), through J.tanA_of (Pr.degEq, Z.J.unitOf, linkOf, Z.J.tan, gwlOf, leOf2). So TanA is proved, not assumed: neither (d) nor the theorem has it as a hypothesis.
% (a) is Sd.Seed2Alt, proved by Z.Win.Birth.seed2_alt through Z.Bc.second_initial_of_prints (with awChain_of_printsH and propEGP_of_printsH).
% (b) is Z.Sym.UnitaryThreeAutomorphyDense, proved by W.unitary_three_automorphy_dense_of_prints through adjoint_lift_from_eisenstein_of_propagation, prop_eg_w_of_printed_hida, seedInW_alt and adjoint_automorphy_w_of_printed_hida_window.
% (c) is Z.Sym.Pl.LinkedEndEAlt, proved by Z.Sym.Top.linkedEnd_alt through Gal.projectedFormsSystemPart.linked_endomorphism_dense.
% (d) is Z.J.LinkedUnlE, proved by Z.J.linkedUnlE_holds_alt through Gal.projectedFormsSystemPart.linkedUnlE_alt.
% Automorphy of points: W.ModUAlt, from W.modU_of_printsH, combining Hi.hida_unitary_of_printed_hida (Proposition sha-hecke-ideal) and Pa.pFromAut (Proposition sha-patching-bound).
% (a) steps: (1) Z.Win.Birth.PropEGP, written out as Proposition sha-automorphy-eisenstein-galois; (3) Sd.sI with eisReg; (4) pair_two; (5) up, Sd.twist_real_selfDual, GalN.V.aut_csd_zero_of_complex_conj_pair_automorphic, and upW from AwChain (windowed).
% (b) steps: (1) Z.Sym.W.PropEGW (Unlinked of ad0), proved by prop_eg_w_of_printed_hida from prop_eg_complex_of_printed_hida (W.PropEGc, with RW and NoHalf) and Unl.adUnl; PropEGc follows the proof of Proposition sha-automorphy-eisenstein-galois with Unl.propEC (the eigenform also satisfies Cj, from unitFormC and cusp_comb) in place of Sd.propE; (3) Sd.sI, ad0Irr, exists_symmetric_pairing_adjoint_zero; (4) Z.Sym.W.SeedInW and AdAutW (adjoint_automorphy_w_of_pure_weight_adjoint_lift from PureWeightCompatibility and AdLiftW). A textual dependency search finds neither hida_unitary_of_printed_hida, exists_hecke_ideal, pFromAut nor modU_of_printsH below (b).
% (c) steps, in linked_endomorphism_dense and twoRuns_alt: (1) kaOf, Sd.Twist; (2) coefEx, Pr.CReal, cRealAv; (3) seedPtOf_alt (Setup.IdealSpecializationPoint); (4) Pr.modularity_linked_ends (Setup.ModIn, DimT, Unr0); (5) and (6) Pr.absMain through components_good_rank_two (PropagationTheoremGL2: preconnected_rank_two, sharp2, exists_nice_prime_above_of_large_dim, P1', SeedIn, globProp) and components_good_rank_three (PropagationTheoremGL3, P23', l53, OffFin), with finPtOf and finPtDown; (7) autPtO_alt (AutOffN) and glueN.
% (d) steps, in UnitaryRankTwoFlagsAlt: (1) OpenEx2EAlt (with J.Num); (2) J.XiEx, J.XiTI, J.XiArith; (3) J.UnramifiedEndsStatement (J.UnrAt); (4) Red2Alt and J.End; (5) J.TanA.

\emph{Irreducibility on open subgroups.} Let $g_1$ be as in the statement, let $r$ be an ordinary point relative to $g_1$ and $\mathfrak B$ with values in a discrete valuation ring $A\subset\bar{\mathbf Q}_\ell$ finite over $\mathcal O$, whose restriction to $\ker\eta$ is irreducible, and let $U\subset\ker\eta$ be open. Then $r$ is irreducible on $U$.
Indeed, replacing $U$ by its normal core, we may assume that $U$ is normal in $\operatorname{Gal}(\bar{\mathbf Q}/\mathbf Q)$.
The inertia group at $\mathfrak B$ lies in $\ker\eta$, and contains an element $x$ with $\Psi(x)^2-1\in\mathcal O^\times$, since $\ell-1\nmid2k$ and $\psi_\ell^2=1$, and an element $\sigma$ with $\chi_{\mathrm{cyc}}(\sigma)=1+\ell$.
Let $u=\sigma^m$, where $m$ is the index of $U$ times $\varphi(N)$; then $u\in U$, $\psi(u)=1$, and $\Psi(u)^2=(1+\ell)^{2km}\ne1$.
After conjugation, $r$ is upper triangular on the decomposition group at $\mathfrak B$ intersected with $\ker\eta$, with diagonal $\Psi,1,\Psi^{-1}$ on inertia; so $r(u)$ is triangular with three distinct eigenvalues, and after a further conjugation by an upper triangular matrix $r(u)$ is diagonal and $r(x)$ is still upper triangular.
Suppose that $r$ is reducible on $U$.
Every $U$-stable subspace is spanned by eigenvectors of $r(u)$, that is, by basis vectors, and $r(g)$ maps $U$-stable subspaces to $U$-stable subspaces for $g\in\ker\eta$, since $U$ is normal.
Let $W$ be a proper nonzero $U$-stable subspace; if it is a plane, it meets a plane $r(g)^{-1}W\ne W$, with $g\in\ker\eta$, in a line. So there is a $U$-stable basis line $L$, and the lines $r(g)L$ with $g\in\ker\eta$ are $U$-stable basis lines, and they span $\bar{\mathbf Q}_\ell^3$ by the irreducibility on $\ker\eta$.
So every $r(g)$ with $g\in\ker\eta$ permutes the three basis lines, by a permutation $\pi_g$; the group of the $\pi_g$ is transitive, so it contains a $3$-cycle $\pi_g$.
Since $r(x)$ is upper triangular, $\pi_x$ is the identity, so $r(gx^n)$ has zero diagonal and trace $0$ for every $n\ge0$.
Reducing modulo $\mathfrak m_A$, where $r$ becomes $r_0(\eta,\Psi)$, the residues $a=\bar\Psi(g)$ and $b=\bar\Psi(x)$ satisfy $ab^n+1+(ab^n)^{-1}=0$ for $n=0,1,2$.
So $a$, $ab$ and $ab^2$ are roots of $X^2+X+1$, two of them coincide, and $b^2=1$, which contradicts $\Psi(x)^2-1\in\mathcal O^\times$.
% This paragraph is Pr.matIrr_open_gd (frame, through core', exists_line, coord and roots, in MatIrrOpenGd.lean); x and the inclusion of inertia in ker eta come from U3.dat_facts.

\emph{Descent.} Let $(\ell,k,\psi,K)$ be an Eisenstein datum, with $\psi$ a Dirichlet character modulo $N$.
By Lemma~\ref{sha-descent-eisenstein}, $(\eta,\Psi,c,\mathfrak B,\ell N|d_K|,\operatorname{Gal}(\bar{\mathbf Q}/\mathbf Q))$ is a descent datum.
Proposition~\ref{sha-descent-construction} gives an abelian CM field $L\supset K$, an element $g_1\in\operatorname{Gal}(\bar{\mathbf Q}/L)$ with $\bar\Psi(g_1)^2\ne1$, a link to $L$ with base element $g_1$ that satisfies the residual conditions with a constant $C_0$ and is large for $C_0$, and a seed datum for $L$.
Since $\operatorname{Gal}(\bar{\mathbf Q}/L)\subset\ker\eta=\operatorname{Gal}(\bar{\mathbf Q}/K)$, $g_1\in\operatorname{Gal}(\bar{\mathbf Q}/K)$, and $\Psi(g_1)^2-1\in\mathcal O^\times$.
By Lemma~\ref{sha-descent-points}, an ordinary point relative to $g_1$ and $\mathfrak B$ with values in a complete local noetherian domain gives a linked point of rank three, and if it is block diagonal on $\ker\eta$ as in (iii), its corner block gives a linked point of rank two.
% The construction is Fr.gd_ex (GdFrameEx): Fr.seed_cl and Fr.frame_ex applied to Fr.gen_gd (Lemma sha-descent-eisenstein). The rest of the descent step is Z.Win.Dn.gdEnd_of_frame_linkedEnd_downCsdK_descent: marks_of_link gives the properties of g1, finEnd_of_finPt gives (i), and exists_nrep3 (inside autEnd3_alt) and exists_nrep2 (inside autEnd2_alt) give the linked points used for (ii) and (iii).
Statement (i) follows from the first part of (c), applied to the linked points of Lemma~\ref{sha-descent-points}(1) with $U=\operatorname{Gal}(\bar{\mathbf Q}/L)$: when $\mathfrak m_{\mathcal O}$ maps to $0$ in $A$, their characters $\psi_{v,i}$ are trivial on the $h\in I_v$ with $\bar\Psi(h)=1$.
For (ii), $r$ is irreducible on every open subgroup of $\ker\eta$, by the paragraph on irreducibility on open subgroups; then (c) makes the restriction of $\rho\otimes\chi_{\mathrm{cyc}}^{-1}$ to $L$ automorphic, and the descent of conjugate self-dual automorphy along the solvable extension $L/K$ gives the representation $\pi$ of $GL_3(\mathbf A_K)$; the windowed conditions follow from the ordinary local conditions at $\ell$ and from (d).
For (iii), (c) makes the restriction of $W$ to $L$ automorphic after an algebraic twist; the descent of automorphy in rank two from $L$ to $\mathbf Q$, and the attachment of a two-dimensional automorphic representation with Hodge-Tate weights $\pm k$ to a newform of weight $2k+1$, give the newform, and the windowed conditions again follow from the local conditions at $\ell$.
\end{proof}
\leanref{Development/AutomorphicArithmetic/AutomorphicVarieties/ModularityArgument/ThresholdArguments/DenseThresholdCases/AutEnd3AltAlt.lean}{UpperTheory.V.Pr.GdEndAlt}
\leanref{Development/AutomorphicArithmetic/AutomorphicVarieties/ModularityArgument/ThresholdArguments/DenseThresholdCases/AutEnd3AltAlt.lean}{UpperTheory.V.Pr.AutEnd3Alt}
\leanref{Development/AutomorphicArithmetic/AutomorphicVarieties/CurveDVRSpecialization/RatTorsionFree.lean}{UpperTheory.V.Pr.FinEnd}
\leanref{Development/AutomorphicArithmetic/AutomorphicVarieties/ModularityArgument/AutomorphyLifting/AutomorphyDescent/AutEnd2Ao.lean}{UpperTheory.V.Z.Win.Dn.AutEnd2Ao}
\leanref{Development/AutomorphicArithmetic/AutomorphicVarieties/CurveDVRSpecialization/MatIrrOpenGd.lean}{UpperTheory.V.Pr.matIrr_open_gd}
\leanref{Development/AutomorphicArithmetic/AutomorphicVarieties/ModularityArgument/ThresholdArguments/DenseThresholdCases/AutomorphyUnramifiedCoefficients.lean}{UpperTheory.V.Aut.AutBAlt}
\leanref{Development/AutomorphicArithmetic/AutomorphicVarieties/ModularityArgument/ThresholdArguments/DenseThresholdCases/DnEndsTop.lean}{UpperTheory.V.Z.Win.Dn.ends0_alt}

\subsection{Host data and Siegel data}\label{subsection-sha-siegel-data}

The plus part of a host field is split according to the characters of its Galois group, and these characters are sorted into cases according to the prime $\ell$ and the twist $k$.

\begin{definition}\label{sha-host-datum}
Let $\psi$ be a primitive Dirichlet character of conductor $N_\psi$ and order $o_\psi$.
\begin{enumerate}
\item Let $p_K(\psi)$ be the least prime $p\equiv3\pmod4$ with $p\ge7$ and $p\nmid N_\psi$; let $h^*(\psi)$ be the class number of $\mathbf Q(\sqrt{-p_K(\psi)})$; let $p_1(\psi)$ be the least prime $p\equiv1$ modulo the odd part of $o_\psi$ with $p\nmid2N_\psi p_K(\psi)$. The set of \emph{bad primes} of $\psi$ is
\[
 B_\psi=\{3\}\cup\{\ell\ \text{odd prime}:\ \ell\mid o_\psi N_\psi\,p_K(\psi)\,h^*(\psi)\,p_1(\psi)\}.
\]
\item A \emph{host datum} for $\psi$ is $h=(\psi,e,q,d)$ with $e\ge2$ and $(q,d)$ a prime pair for $M_\psi$. Put $N_h=\operatorname{lcm}(2^{e+2},N_\psi)\,qd$, and let $H_h\subset(\mathbf Z/N_h)^\times$ consist of the $a$ with $a\equiv\pm1\pmod{2^{e+2}}$, $\psi(a)=1$ and $a$ a square modulo $q$ and modulo $d$. For $t\ge0$, $h$ is \emph{large for $t$} if $2t+1<q$, $2t+1<d$, and $N_\psi,o_\psi,p_K(\psi),h^*(\psi),p_1(\psi)$ are all less than $q$ and $d$.
\item For a prime $\ell$, $\mathcal O$ as in Definition~\ref{sha-character-parts} and $k\ge0$, let $X_h(k)$ be the set of characters $\psi'$ of $(\mathbf Z/N_h)^\times$ with values in $\mathcal O$, trivial on $H_h$, with $\psi'(-1)=(-1)^k$.
\item Let $\ell$ be an odd prime, $k\ge1$ and $\psi'\in X_h(k)$, and let $\chi'_\ell(\psi')$ be as in Definition~\ref{sha-ordinary-deformation}. The triple $(\ell,k,\psi')$ is in
\begin{itemize}
\item \emph{case I} if $\ell\notin B_\psi$ and $\ell-1\nmid2k$;
\item \emph{case Z} if $\ell\notin B_\psi$, $\ell-1\mid2k$ and $\psi'=1$;
\item \emph{case II} if $k\ge2$, $\chi'_\ell(\psi')$ has even order, and $\ell\in B_\psi$ or $\ell-1\mid2k$;
\item \emph{case III} if $\ell\notin B_\psi$, $\ell-1\mid2k$, and $\psi'\ne1$ has odd order;
\item \emph{case IV} if $\ell\in B_\psi$, $k\ge2$ and $\chi'_\ell(\psi')$ has odd order;
\item \emph{case V} if $\ell\in B_\psi$ and $k=1$.
\end{itemize}
\end{enumerate}
\end{definition}
\leanref{Development/AutomorphicArithmetic/AutomorphicVarieties/StmtFields.lean}{UpperTheory.V.pK}
\leanref{Development/AutomorphicArithmetic/AutomorphicVarieties/StmtFields.lean}{UpperTheory.V.hStar}
\leanref{Development/AutomorphicArithmetic/AutomorphicVarieties/StmtFields.lean}{UpperTheory.V.pOne}
\leanref{Development/AutomorphicArithmetic/AutomorphicVarieties/StmtFields.lean}{UpperTheory.V.Bchi}
\leanref{Development/AutomorphicArithmetic/AutomorphicVarieties/HeightAndVanishing/DirichletPrimeToLPartMulLPowerPart.lean}{UpperTheory.V.HV.Host}
\leanref{Development/AutomorphicArithmetic/AutomorphicVarieties/HeightAndVanishing/DirichletPrimeToLPartMulLPowerPart.lean}{UpperTheory.V.HV.Host.H}
\leanref{Development/AutomorphicArithmetic/AutomorphicVarieties/HeightAndVanishing/DirichletPrimeToLPartMulLPowerPart.lean}{UpperTheory.V.HV.Host.Big}
\leanref{Development/AutomorphicArithmetic/AutomorphicVarieties/HeightAndVanishing/DirichletPrimeToLPartMulLPowerPart.lean}{UpperTheory.V.HV.Host.X}
\leanref{Development/AutomorphicArithmetic/AutomorphicVarieties/HeightAndVanishing/DirichletPrimeToLPartMulLPowerPart.lean}{UpperTheory.V.HV.Host.CaseI}

\begin{lemma}\label{sha-cases}
Let $h$ be a host datum for $\psi$ that is large for $t$, and let $E$ be a host field of $(\psi,e,q,d)$.
Then $N_E=N_h$, and $E$ is the subfield of $\mathbf Q(\zeta_{N_h})$ fixed by $H_h$.
Let $\ell$ be an odd prime, $0<k<t$ and $\psi'\in X_h(k)$.
If $\ell\notin B_\psi$, then $\ell\nmid\operatorname{ord}\psi'$ and $(\ell,k,\psi')$ is in case I, Z, II or III.
If $\ell\in B_\psi$, then $(\ell,k,\psi')$ is in case II, IV or V.
\end{lemma}
\begin{proof}
The conditions defining $H_h$ depend only on $a$ modulo $N_h$, so $E$ is the subfield of $\mathbf Q(\zeta_{N_h})$ fixed by $H_h$.
For every prime $p\mid N_h$ there is $a\equiv1$ modulo $N_h/p^{v_p(N_h)}$ with $a\notin H_h$, which shows that $E$ does not embed in $\mathbf Q(\zeta_{N'})$ for any $N'$ with $v_p(N')<v_p(N_h)$; so $N_E=N_h$.
Every $\psi'\in X_h(k)$ has order dividing $2^{e+1}o_\psi$, and its component at an odd prime $p\nmid N_\psi$ has order at most $2$.
Since $h$ is large for $t$, the primes of $B_\psi$ are less than $q$ and $d$; an odd prime $\ell\notin B_\psi$ satisfies $\ell\ge5$ and $\ell\nmid o_\psi$, and if $\ell-1\mid2k$ with $k<t$ then $k\ge2$, $\ell\le2t$ and $\ell\nmid N_h$.

Let $\ell\notin B_\psi$. Then $\ell\nmid\operatorname{ord}\psi'$, so $\psi'$ is its own part of order prime to $\ell$.
If $\ell-1\nmid2k$, the triple is in case I.
If $\ell-1\mid2k$, then $\ell\nmid N_h$, so the component of $\psi'$ at $\ell$ is trivial and $\chi'_\ell(\psi')=\psi'$; the triple is in case Z if $\psi'=1$, in case III if $\psi'\ne1$ has odd order, and in case II if $\psi'$ has even order, since $k\ge2$.
Let $\ell\in B_\psi$. If $k=1$ the triple is in case V; if $k\ge2$ it is in case II or IV according to the parity of the order of $\chi'_\ell(\psi')$.
\end{proof}
\leanref{Development/AutomorphicArithmetic/AutomorphicVarieties/HeightAndVanishing/HostCond.lean}{UpperTheory.V.HV.hostCond}
\leanref{Development/AutomorphicArithmetic/AutomorphicVarieties/HeightAndVanishing/HostLevel.lean}{UpperTheory.V.HV.hostLevel}
\leanref{Development/AutomorphicArithmetic/AutomorphicVarieties/HeightAndVanishing/HostComp.lean}{UpperTheory.V.HV.hostComp}
\leanref{Development/AutomorphicArithmetic/AutomorphicVarieties/HeightAndVanishing/BchiFacts.lean}{UpperTheory.V.HV.bchiFacts}
\leanref{Development/AutomorphicArithmetic/AutomorphicVarieties/HeightAndVanishing/Cases.lean}{UpperTheory.V.HV.cases}

\begin{definition}\label{sha-siegel-datum}
Let $\ell$ be an odd prime, let $\mathcal O$ be as in Definition~\ref{sha-character-parts}, with residue field $\mathbf F$, write $G=\operatorname{Gal}(\bar{\mathbf Q}/\mathbf Q)$, and write $\bar x$ for the image in $\mathbf F$ of $x\in\mathcal O$.
For a commutative ring $R$ and $X\in M_4(R)$ let $X^*$ be the matrix with entries $X^*_{ij}=\epsilon_i\epsilon_jX_{5-j,5-i}$, where $\epsilon=(1,1,-1,-1)$.
A matrix $g\in M_4(R)$ is \emph{symplectic with multiplier} $\nu\in R$ if $g^*g=\nu$, and $\mathfrak{sp}_4(R)$ is the set of $X$ with $X^*=-X$.
Give the entry $(i,j)$ of a $4\times4$ matrix the \emph{weight} $(\lambda_i-\lambda_j,\mu_i-\mu_j)\in\mathbf Z^2$, where $\lambda=(1,0,1,0)$ and $\mu=(1,1,0,0)$.
A matrix is \emph{block diagonal} if its entries $(i,j)$ with $\lambda_i\ne\lambda_j$ vanish; its \emph{first block} is the $2\times2$ matrix of its entries in rows and columns $1$ and $3$.
\begin{enumerate}
\item A \emph{Siegel datum} $D$ consists of a finite set $S$ of primes containing $2$ and $\ell$, an integer $k\ge2$, characters $\psi,\xi,\omega:G\to\mathcal O^\times$, a complex conjugation $c$, a place $v_p$ of $\bar{\mathbf Q}$ above each prime $p$, and elements $\phi$, $g_0$, $s_0$ and $\sigma_p$ (for the primes $p$ in the set $T$ below) of $G$, subject to the conditions below.
Put $\Psi=\chi_{\mathrm{cyc}}^{k}\psi$, $\alpha=\chi_{\mathrm{cyc}}\omega^{-1}\xi$, $\beta=\Psi\alpha^{-1}$, $\theta=\alpha\Psi$, $(\delta_1,\delta_2,\delta_3,\delta_4)=(\theta,\Psi,\alpha,1)$ and $r_0=\operatorname{diag}(\delta_1,\delta_2,\delta_3,\delta_4)$, which is symplectic with multiplier $\theta$.
Let $D_p$ and $I_p$ be the decomposition and inertia groups at $v_p$, let $S_b$ be the set of primes $p\in S$ with $p\ne\ell$ for which $\ell$ divides the order of $\psi(I_p)$, and let $T=S_b\cup\{\ell\}$.
The conditions are:
$\psi$ and $\xi$ have open kernels and are unramified outside $S$; $\Psi(c)=1$; $\omega^{\ell-1}=1$ and $\bar\omega=\bar\chi_{\mathrm{cyc}}$; $\xi$ has order prime to $\ell$, $\xi(c)=-1$, and $\xi$ is trivial on $I_p$ for $p\in T$;
for $p\in T$, $\sigma_p$ is a Frobenius element at $v_p$ and $\overline{\xi\Psi^{i}}(\sigma_p)\ne1$ for every integer $i$; $\chi_{\mathrm{cyc}}(\sigma_\ell)=1$, and $\phi$ is the limit of $\sigma_\ell^{\ell^{n!}}$ as $n\to\infty$;
$\bar\alpha(g_0)\ne1$ and $\bar\Psi(g_0)\ne1$; $s_0\in I_\ell$, $\psi(s_0)=\omega(s_0)=1$, and $\chi_{\mathrm{cyc}}(s_0)$ has infinite order;
and for each $\gamma\in\{\alpha^{\pm1},\beta^{\pm1},\theta^{\pm1}\}$, every locally constant cocycle $z:G\to\mathbf F$ for the action through $\bar\gamma$ that vanishes on the elements $g$ with $\bar r_0(g)=1$ of the inertia groups at places not above $T$, and is a coboundary on $D_p$ for each $p\in T$, is a coboundary.
\item Let $R$ be a complete local noetherian $\mathcal O$-algebra with residue field $\mathbf F$ and maximal ideal $\mathfrak m_R$. A \emph{Siegel point} of $D$ with values in $R$ is a homomorphism $r:G\to M_4(R)$, continuous for the $\mathfrak m_R$-adic topology, with $r\equiv r_0\pmod{\mathfrak m_R}$, such that:
$r(g)$ is symplectic with multiplier $\theta(g)$ for every $g$;
$r(g)=1$ for every $g$ with $\bar r_0(g)=1$ in an inertia group at a place not above $T$;
$r(\phi)$ is block diagonal, and the entries $(1,3)$ and $(3,1)$ of $r(g_0)+\theta(g_0)r(g_0^{-1})$ vanish;
the restriction of $r$ to $D_\ell$ is the image, under a homomorphism of $\mathcal O$-algebras $B\to R$, of a lift $\rho:D_\ell\to M_4(B)$ of $r_0\bmod\mathfrak m_{\mathcal O}$, where $B$ is a complete local noetherian $\mathcal O$-algebra with residue field $\mathbf F$ that is a domain and has a homomorphism of $\mathcal O$-algebras $B\to\mathcal O$ carrying $\rho$ to $r_0$, and $\rho$ is symplectic with multiplier $\theta$ and \emph{ordinary}: with $e_i=\prod_{j\ne i}(\rho(s_0)-\delta_j(s_0))$, the characteristic polynomial of $\rho(s_0)$ is $\prod_j(X-\delta_j(s_0))$, $e_j\rho(g)e_i=0$ for $i<j$ and $g\in D_\ell$, and $e_i(\rho(s)-\delta_i(s))e_i=0$ for $s\in I_\ell$;
and for each $p\in S_b$, the restriction of $r$ to $D_p$ is obtained in the same way from a lift over such a domain that is symplectic with multiplier $\theta$ and trivial on $I_p\cap\ker\psi$.
A \emph{$\bar{\mathbf Q}_\ell$-point} of $D$ is the image of a Siegel point under a homomorphism of $\mathcal O$-algebras $R\to\bar{\mathbf Q}_\ell$; points with values in $\bar{\mathbf Q}_\ell[\varepsilon]$, $\varepsilon^2=0$, are defined in the same way.
\item The datum $D$ is \emph{finite} if for every Siegel point $r$ with values in $R$, the $\mathcal O$-subalgebra of $R$ generated by the products $r(g_1)_{i_1j_1}\cdots r(g_m)_{i_mj_m}$ whose weights add up to $0$ is finite over $\mathcal O$.
A $\bar{\mathbf Q}_\ell$-point $r$ is \emph{endoscopic} if every $r(g)$ is block diagonal and the first blocks of the $r(g)$ span $M_2(\bar{\mathbf Q}_\ell)$.
The datum $D$ is \emph{rigid} at a $\bar{\mathbf Q}_\ell$-point $r$ if every point $r'$ with values in $\bar{\mathbf Q}_\ell[\varepsilon]$ with $r'\equiv r\pmod\varepsilon$ is $r+\varepsilon(Xr-rX)$ for some $X\in\mathfrak{sp}_4(\bar{\mathbf Q}_\ell)$.
For $\iota:\bar{\mathbf Q}_\ell\cong\mathbf C$, a continuous representation $r:G\to GL_4(\bar{\mathbf Q}_\ell)$ is \emph{automorphic through $\iota$} if, for every imaginary quadratic field $K$ such that $r$ is irreducible on $\operatorname{Gal}(\bar{\mathbf Q}/K)$, the restriction of $r$ to $\operatorname{Gal}(\bar{\mathbf Q}/K)$ is attached through $\iota$ to a regular algebraic cuspidal automorphic representation of $GL_4(\mathbf A_K)$ and is windowed (Theorem~\ref{sha-residually-reducible-automorphy}).
\item Let $h$ be a host datum for a primitive Dirichlet character $\psi_0$ (Definition~\ref{sha-host-datum}, with $\psi_0$ in place of $\psi$), let $k\ge1$ and $\psi\in X_h(k)$.
Let $2^{c_2}$ be the conductor of the $2$-component of $\psi$, let $\vartheta$ be the part of $\psi$ of order prime to $\ell$ with its $2$-component removed, let $f_0$ be the order of $\ell$ modulo $\operatorname{lcm}(4,\operatorname{ord}\vartheta)$, let $n_0=v_2(\ell^{f_0}-1)$, and put
\[
 n_{\mathrm{II}}=\big\lceil\log_2\max\big(\ell N_h4^{n_0},\,2^{c_2+6}o_{\psi_0}\ell^3N_{\psi_0}^2\big)\big\rceil,\qquad n_{\mathrm{III}}=\big\lceil\log_2\max\big(\ell N_{\psi_0}p_1(\psi_0)4^{n_0},\,8\ell^2\big)\big\rceil.
\]
A Siegel datum $D$ is \emph{attached} to $(h,k,\psi)$ if its integer is $k$ and its character $\psi$ is the Dirichlet character $\psi$ (viewed as in Definition~\ref{sha-character-parts}), and either $(\ell,k,\psi)$ is in case II, $S$ is the set of primes dividing $2\ell N_h$ and $\xi$ is a primitive odd character of conductor $2^{n_{\mathrm{II}}}$; or $(\ell,k,\psi)$ is in case III, $S$ is the set of primes dividing $2\ell N_hp_1(\psi_0)$, and $\xi=\xi_0\xi'$ with $\xi'$ primitive and odd of conductor $2^{n_{\mathrm{III}}}$ and $\xi_0$ a character modulo $p_1(\psi_0)$ of order $\operatorname{ord}\psi$.
\end{enumerate}
\end{definition}
\leanref{Development/AutomorphicArithmetic/AutomorphicVarieties/SymplecticGSp4/DecompSigmaOrd.lean}{UpperTheory.V.GradeZeroInterpolatingElement.Raw}
\leanref{Development/AutomorphicArithmetic/AutomorphicVarieties/SymplecticGSp4/DecompSigmaOrd.lean}{UpperTheory.V.GradeZeroInterpolatingElement.Raw.Good}
\leanref{Development/AutomorphicArithmetic/AutomorphicVarieties/SymplecticGSp4/DecompSigmaOrd.lean}{UpperTheory.V.GradeZeroInterpolatingElement.Raw.IsPt}
\leanref{Development/AutomorphicArithmetic/AutomorphicVarieties/SymplecticGSp4/DecompSigmaOrd.lean}{UpperTheory.V.GradeZeroInterpolatingElement.Raw.Fin0}
\leanref{Development/AutomorphicArithmetic/AutomorphicVarieties/SymplecticGSp4/DecompSigmaOrd.lean}{UpperTheory.V.GradeZeroInterpolatingElement.Raw.NoEnd}
\leanref{Development/AutomorphicArithmetic/AutomorphicVarieties/SymplecticGSp4/DecompSigmaOrd.lean}{UpperTheory.V.GradeZeroInterpolatingElement.Raw.RigidAt}
\leanref{Development/AutomorphicArithmetic/AutomorphicVarieties/ModularityArgument/AutomorphyLifting/IsAutCsdA.lean}{UpperTheory.V.Z.Win.IsAutQA}
\leanref{Development/AutomorphicArithmetic/AutomorphicVarieties/AutomorphicCases/SymplecticGSp4.lean}{UpperTheory.V.Aut.OfCase}
\leanref{Development/AutomorphicArithmetic/AutomorphicVarieties/HeightAndVanishing/DirichletPrimeToLPartMulLPowerPart.lean}{UpperTheory.V.HV.Host.nII}
\leanref{Development/AutomorphicArithmetic/AutomorphicVarieties/HeightAndVanishing/DirichletPrimeToLPartMulLPowerPart.lean}{UpperTheory.V.HV.Host.nIII}

\begin{lemma}\label{sha-flagged-residues}
Let $D$ be a Siegel datum, with the notation of Definition~\ref{sha-siegel-datum}.
\begin{enumerate}
\item For $p\in T$, the residues $\bar\delta_i(\sigma_p)$ and $\bar\delta_j(\sigma_p)$ are distinct whenever $\lambda_i\ne\lambda_j$.
\item The element $\phi$ lies in $D_\ell$, and $\bar r_0(\phi)=\bar r_0(\sigma_\ell)$.
\item $\alpha(c)=\beta(c)=\theta(c)=-1$ and $\Psi(c)=1$, so $r_0(c)=\operatorname{diag}(-1,1,-1,1)$.
\item Every $p\in S_b$ satisfies $p\equiv1\pmod\ell$, and $\bar\chi_{\mathrm{cyc}}$ is trivial on $D_p$.
\end{enumerate}
\end{lemma}
\begin{proof}
Since $\bar\omega=\bar\chi_{\mathrm{cyc}}$, the residue of $\alpha=\chi_{\mathrm{cyc}}\omega^{-1}\xi$ is $\bar\xi$.
The quotients $\delta_i/\delta_j$ with $\lambda_i\ne\lambda_j$ are $\alpha^{\pm1}$, $(\alpha\Psi)^{\pm1}$ and $(\alpha\Psi^{-1})^{\pm1}$, whose residues are $(\bar\xi\bar\Psi^{e})^{\pm1}$ with $e\in\{0,1,-1\}$; at $\sigma_p$ these differ from $1$ by the conditions on $\sigma_p$. This is (1).
For (2), $D_\ell$ is closed and contains the $\sigma_\ell^{\ell^{n!}}$. The $\bar\delta_i$ take values in $\mathbf F$, of order $\ell^f$ say, so $\bar\delta_i(\sigma_\ell^{\ell^{n!}})=\bar\delta_i(\sigma_\ell)$ once $f\mid n!$; since the $\bar\delta_i$ are locally constant, they take the same values at $\phi$.
For (3), $\chi_{\mathrm{cyc}}(c)=\xi(c)=-1$, and $\omega(c)=-1$, because $\omega(c)$ is an $(\ell-1)$-st root of unity with residue $\bar\chi_{\mathrm{cyc}}(c)=-1$ and distinct $(\ell-1)$-st roots of unity have distinct residues. So $\alpha(c)=-1$, $\theta(c)=\alpha(c)\Psi(c)=-1$ and $\beta(c)=\Psi(c)\alpha(c)^{-1}=-1$.
For (4), $\psi$ has open kernel, so by the Kronecker-Weber theorem $\psi(I_p)$ is a quotient of $\mathbf Z_p^\times$; its order is divisible by $\ell\ne p$, so $\ell\mid p-1$. The group $D_p$ is topologically generated by $I_p$ and a Frobenius element $\gamma$. The values of $\omega$ are $(\ell-1)$-st roots of unity, which are determined by their residues; the residue $\bar\chi_{\mathrm{cyc}}$ of $\omega$ is trivial on $I_p$ since $p\ne\ell$, and $\bar\chi_{\mathrm{cyc}}(\gamma)=p\equiv1$. So $\omega$, and with it $\bar\chi_{\mathrm{cyc}}=\bar\omega$, is trivial on $D_p$.
\end{proof}
\leanref{Development/AutomorphicArithmetic/AutomorphicVarieties/SymplecticGSp4/LMemMax.lean}{UpperTheory.V.GradeZeroInterpolatingElement.facts}

\begin{lemma}\label{sha-flagged-block-idempotent}
Let $R$ be a local ring with maximal ideal $\mathfrak m$, write $E_{ij}$ for the matrix units of $M_4(R)$ and $E=E_{11}+E_{33}$, and let $X\in M_4(R)$ be block diagonal (Definition~\ref{sha-siegel-datum}), with $X_{ij}\in\mathfrak m$ for $i\ne j$ and $X_{ii}-X_{jj}\in R^\times$ whenever $\lambda_i\ne\lambda_j$.
Then $E$ lies in every $R$-subalgebra of $M_4(R)$ that contains $X$.
\end{lemma}
\begin{proof}
Let $\Lambda$ be such a subalgebra, and let $t$ and $d$ be the trace and the determinant of the block of $X$ in rows and columns $2$ and $4$.
By the Cayley-Hamilton theorem for that block, $Y=X^2-tX+d$ lies in $\Lambda$ and is supported in rows and columns $1$ and $3$; let $Y'$ be its block there.
Modulo $\mathfrak m$, $Y'$ is diagonal with entries $(X_{11}-X_{22})(X_{11}-X_{44})$ and $(X_{33}-X_{22})(X_{33}-X_{44})$, so $\det Y'$ is a unit.
The Cayley-Hamilton theorem for $Y'$ gives $Y^2-(\operatorname{tr}Y')Y+(\det Y')E=0$, so $E=(\det Y')^{-1}\big((\operatorname{tr}Y')Y-Y^2\big)$ lies in $\Lambda$.
\end{proof}
\leanref{Development/AutomorphicArithmetic/AutomorphicVarieties/WeakShimuraTheory/CompactIntegralFields/BlockPart.lean}{ee_mem}
\leanref{Development/AutomorphicArithmetic/AutomorphicVarieties/SymplecticGSp4/EiiUnit.lean}{eii_step}

\begin{lemma}\label{sha-flagged-idempotents}
Let $D$ be a Siegel datum, let $R$ be a complete local noetherian $\mathcal O$-algebra with residue field $\mathbf F$, and let $r$ be a Siegel point of $D$ with values in $R$ (Definition~\ref{sha-siegel-datum}).
Then the $R$-submodule of $M_4(R)$ spanned by the $r(g)$, $g\in G$, contains the matrix units $E_{11}$, $E_{22}$, $E_{33}$ and $E_{44}$.
\end{lemma}
\begin{proof}
Let $M$ be that submodule. It contains $1=r(1)$ and is closed under multiplication, since $r(g)r(h)=r(gh)$. Since $r(g)^*r(g)=\theta(g)$, we have $r(g)^*=\theta(g)r(g^{-1})$, so $M$ is stable under $X\mapsto X^*$.
The matrix $r(\phi)$ is block diagonal, its off-diagonal entries lie in $\mathfrak m_R$, and its diagonal entries are congruent to the $\delta_i(\phi)$, whose residues are those of the $\delta_i(\sigma_\ell)$ and are distinct across the blocks, by Lemma~\ref{sha-flagged-residues}. By Lemma~\ref{sha-flagged-block-idempotent}, $E=E_{11}+E_{33}$ lies in $M$.
Let $H=r(g_0)+\theta(g_0)r(g_0^{-1})$, an element of $M$ whose entries $(1,3)$ and $(3,1)$ vanish. Then $EHE-H_{33}E=(H_{11}-H_{33})E_{11}$.
Modulo $\mathfrak m_R$, $H_{11}\equiv\theta(g_0)+1$ and $H_{33}\equiv\alpha(g_0)+\Psi(g_0)$, so $H_{11}-H_{33}\equiv(\alpha(g_0)-1)(\Psi(g_0)-1)$, which is a unit since $\bar\alpha(g_0)\ne1$ and $\bar\Psi(g_0)\ne1$.
Hence $E_{11}\in M$ and $E_{33}=E-E_{11}\in M$. Finally $E_{kk}^*=E_{5-k,5-k}$, so $E_{44}=E_{11}^*$ and $E_{22}=E_{33}^*$ lie in $M$.
\end{proof}
\leanref{Development/AutomorphicArithmetic/AutomorphicVarieties/SymplecticGSp4/EiiUnit.lean}{UpperTheory.V.GradeZeroInterpolatingElement.eii}

\begin{lemma}\label{sha-flagged-relation}
Let $K$ be a field, let $H$ be a group, and let $r:H\to M_4(K)$ be a homomorphism such that each $r(h)$ is symplectic with some multiplier (Definition~\ref{sha-siegel-datum}) and the $K$-span $\Lambda$ of the $r(h)$ contains the matrix units $E_{11},\dots,E_{44}$.
For $i,j\in\{1,2,3,4\}$ write $i\succeq j$ if $r(h)_{ij}\ne0$ for some $h\in H$.
\begin{enumerate}
\item $i\succeq j$ if and only if $E_{ij}\in\Lambda$.
\item The relation $\succeq$ is reflexive and transitive, and $i\succeq j$ implies $5-j\succeq5-i$.
\item $\Lambda=M_4(K)$ if and only if $i\succeq j$ for all $i$ and $j$.
\item If every $r(h)$ is block diagonal, $1\succeq3$ and $3\succeq1$, then the first blocks of the $r(h)$ span $M_2(K)$.
\end{enumerate}
\end{lemma}
\begin{proof}
The span $\Lambda$ is a subalgebra. If $r(h)_{ij}\ne0$, then $E_{ii}r(h)E_{jj}=r(h)_{ij}E_{ij}$ gives $E_{ij}\in\Lambda$; if $r(h)_{ij}=0$ for all $h$, every element of $\Lambda$ has entry $(i,j)$ equal to $0$. This is (1), and (3) follows.
Reflexivity is $E_{ii}\in\Lambda$, and transitivity follows from $E_{ij}E_{jk}=E_{ik}$.
Each $r(h)$ is invertible, with inverse $r(h^{-1})$, so its multiplier $\nu$ is nonzero and $r(h^{-1})=\nu^{-1}r(h)^*$; the entry $(5-j,5-i)$ of $r(h)^*$ is $\pm r(h)_{ij}$, which gives the rest of (2).
For (4), taking the first block is a $K$-linear map from $\Lambda$ to the span of the first blocks, and it carries $E_{11},E_{13},E_{31},E_{33}$, which lie in $\Lambda$ by (1), to the matrix units of $M_2(K)$.
\end{proof}
\leanref{Development/AutomorphicArithmetic/AutomorphicVarieties/SymplecticGSp4/GeAlgTop.lean}{UpperTheory.V.GradeZeroInterpolatingElement.geAlg}

\begin{lemma}\label{sha-flagged-cocycles}
Let $D$ be a Siegel datum, with the notation of Definition~\ref{sha-siegel-datum}, let $k$ be a field with a ring homomorphism $a:\mathcal O\to k$ that kills $\mathfrak m_{\mathcal O}$, let $\gamma\in\{\alpha^{\pm1},\beta^{\pm1},\theta^{\pm1}\}$, and write $\bar\gamma=a\circ\gamma$.
\begin{enumerate}
\item Let $z:G\to k$ be locally constant, with $z(gh)=z(g)+\bar\gamma(g)z(h)$ for all $g,h\in G$. Suppose that $z(g)=0$ for every $g$ with $\bar r_0(g)=1$ in an inertia group at a place not above $T$, and that for each $p\in T$ there is $m_p\in k$ with $z(g)=\bar\gamma(g)m_p-m_p$ for $g\in D_p$. Then there is $m\in k$ with $z(g)=\bar\gamma(g)m-m$ for all $g\in G$.
\item Let $p\in S_b$. Every locally constant $z:D_p\to k$ with $z(gh)=z(g)+\bar\gamma(g)z(h)$ for all $g,h\in D_p$ has the form $z(g)=\bar\gamma(g)m-m$ for some $m\in k$.
\end{enumerate}
\end{lemma}
\begin{proof}
Through $a$, the field $k$ is a vector space over $\mathbf F$.
For an $\mathbf F$-linear map $\mu:k\to\mathbf F$, the function $\mu\circ z$ is a locally constant cocycle with values in $\mathbf F$ for the action through the residue of $\gamma$, and it satisfies the hypotheses of the part in question, with $\mu(m_p)$ in place of $m_p$.
Suppose that each $\mu\circ z$ is a coboundary, $\mu(z(g))=\bar\gamma(g)m_\mu-m_\mu$.
If $\bar\gamma(g)=1$ for all $g$ in the domain of $z$, then $\mu(z(g))=0$ for every $\mu$, so $z=0$.
Otherwise choose $g_1$ with $\bar\gamma(g_1)\ne1$ and put $m=z(g_1)/(\bar\gamma(g_1)-1)$; then $\mu(m)=m_\mu$, so $\mu\big(z(g)-\bar\gamma(g)m+m\big)=0$ for every $\mu$, and $z(g)=\bar\gamma(g)m-m$.
So it suffices to treat $k=\mathbf F$.
In (1) this is the last condition of Definition~\ref{sha-siegel-datum}.
In (2), $\mathbf F(\bar\gamma)$ is a finite module for $D_p$ with $p\ne\ell$, and it suffices that $H^1(D_p,\mathbf F(\bar\gamma))=0$.
By the local Euler characteristic formula and local duality \cite[Chapter~VII]{NSW08}, its dimension is the sum of the dimensions of $H^0(D_p,\mathbf F(\bar\gamma))$ and $H^0(D_p,\mathbf F(\bar\gamma^{-1}\bar\chi_{\mathrm{cyc}}))$.
Write $\gamma=\delta_i\delta_j^{-1}$ with $\lambda_i\ne\lambda_j$. By Lemma~\ref{sha-flagged-residues}, $\bar\gamma(\sigma_p)\ne1$ and $\bar\chi_{\mathrm{cyc}}(\sigma_p)=1$, so both spaces of invariants vanish.
\end{proof}
\leanref{Development/AutomorphicArithmetic/AutomorphicVarieties/SymplecticGSp4/XvanExt.lean}{UpperTheory.V.GradeZeroInterpolatingElement.xvanExt}
\leanref{Development/AutomorphicArithmetic/AutomorphicVarieties/SymplecticGSp4/LocVanP.lean}{UpperTheory.V.GradeZeroInterpolatingElement.locVanP}

\begin{lemma}\label{sha-flagged-common-line}
Let $K$ be a field, let $H$ be a group, let $r:H\to M_4(K)$ be a homomorphism, and let $u\in GL_4(K)$ be such that $u^{-1}r(h)u$ is upper triangular for every $h\in H$.
Let $\Pi\subseteq\{1,2,3,4\}$ be such that $r(h)_{ab}=0$ whenever $b\in\Pi$ and $a\notin\Pi$, let $i\in\Pi$ and $j\ne i$, and suppose that $r(h)_{ik}=r(h)_{jk}=0$ for all $h$ and all $k\in\Pi$ other than $i$ and $j$.
Then there are $(x,y)\in K^2$, not both $0$, and a function $\varphi:H\to K$ with
\[
 r(h)_{ii}x+r(h)_{ij}y=\varphi(h)x,\qquad r(h)_{ji}x+r(h)_{jj}y=\varphi(h)y\qquad(h\in H).
\]
\end{lemma}
\begin{proof}
Let $W\subseteq K^4$ be spanned by the basis vectors $e_b$ with $b\in\Pi$; it is stable under the $r(h)$.
Let $\pi:W\to K^2$ be $w\mapsto(w_i,w_j)$, and let $B(h)$ be the matrix with rows $(r(h)_{ii},r(h)_{ij})$ and $(r(h)_{ji},r(h)_{jj})$.
For $w\in W$, $(r(h)w)_i=\sum_{k\in\Pi}r(h)_{ik}w_k=r(h)_{ii}w_i+r(h)_{ij}w_j$, since the other terms vanish and $w_j=0$ if $j\notin\Pi$; the same holds for the coordinate $j$. So $\pi(r(h)w)=B(h)\pi(w)$.
Let $F_n$ be the span of the first $n$ columns of $u$; it is stable under the $r(h)$, which act on $F_n/F_{n-1}$ by the $n$-th diagonal entry $t_n(h)$ of $u^{-1}r(h)u$.
Since $\pi(e_i)\ne0$, there is a least $n$ such that $F_n\cap W$ contains a vector $w$ with $\pi(w)\ne0$.
For $h\in H$, the vector $r(h)w-t_n(h)w$ lies in $F_{n-1}\cap W$, so $\pi$ kills it, and $B(h)\pi(w)=t_n(h)\pi(w)$. So $(x,y)=\pi(w)$ and $\varphi=t_n$ have the required property.
\end{proof}
\leanref{Development/AutomorphicArithmetic/AutomorphicVarieties/WeakShimuraTheory/CompactIntegralFields/CobOfCoord.lean}{line}

\begin{lemma}\label{sha-flagged-extension-cocycle}
Let $\ell$ be an odd prime, write $G=\operatorname{Gal}(\bar{\mathbf Q}/\mathbf Q)$, let $S$ be a finite set of primes containing $\ell$, let $D_\ell$ be a decomposition group at a place above $\ell$, let $c$ be a complex conjugation, and let $s\in D_\ell$ with $\chi_{\mathrm{cyc}}(s)=1$.
Let $A$ be a complete discrete valuation ring with finite residue field, with $\ell$ in its maximal ideal $\mathfrak m_A$, and let $T$ be a free $A$-module of finite rank with an $A$-linear action $\rho$ of $G$ such that $g\mapsto\rho(g)x\bmod\mathfrak m_A^nT$ is locally constant for all $x\in T$ and $n\ge1$.
Suppose that $\rho(c)=-1$, that $\rho(s)-1$ is bijective on $T$, and that $N\subseteq\ker\rho$ is a subgroup containing the inertia groups at the places not above $S$.
Let $P\subseteq T$ be a nonzero $A$-submodule, stable under $D_\ell$, such that $x\in P$ whenever $ax\in P$ with $a\in A$ nonzero.
Then there are a function $z:G\to T$, not identically $0$, and $y\in T$ such that: $z(gh)=z(g)+\rho(g)z(h)$ for all $g,h$; $z\bmod\mathfrak m_A^nT$ is locally constant for every $n$; $z(c)=0$; $z(g)=0$ for every $g\in N$ that lies in an inertia group at a place not above $\ell$; and $z(h)-(\rho(h)y-y)\in P$ for every $h\in D_\ell$.
\end{lemma}
\begin{proof}
For $n\ge1$ let $Z_n$ be the set of $z:G\to T/\mathfrak m_A^nT$ with the analogous properties, for some $y\in T/\mathfrak m_A^nT$, with the image $P_n$ of $P$ in place of $P$; it is an $A$-module. Let $\varpi$ be a uniformizer of $A$.
\begin{enumerate}
\item[(a)] Each $Z_n$ is finite: its elements are cocycles of $G_S$ with values in the finite module $T/\mathfrak m_A^nT$, and these form a finite set \cite[Chapter~VIII]{NSW08}.
\item[(b)] $|Z_n|\ge|P_n|\ge n$. For the first inequality, take the Selmer group of $T/\mathfrak m_A^nT$ whose local conditions are the unramified classes at the places of $S$ other than $\ell$ and, at $\ell$, a subgroup of $H^1(D_\ell,T/\mathfrak m_A^nT)$ with at least $|P_n|$ elements whose classes are represented by cocycles $h\mapsto\rho(h)y-y+p(h)$ with $p(h)\in P_n$; such a subgroup exists since $\rho(s)-1$ is bijective. The Greenberg-Wiles formula \cite[Proposition~1.6]{Wil95}, with Poitou-Tate duality and the local Euler characteristic formula \cite{NSW08}, shows that this Selmer group has at least $|P_n|$ elements; there the unramified classes have as many elements as the invariants, $c$ acts by $-1$, and $s$ satisfies $\chi_{\mathrm{cyc}}(s)=1$. Since $2$ is a unit and $c$ acts by $-1$, each class contains a unique cocycle with $z(c)=0$, and these cocycles lie in $Z_n$. For the second inequality, $P$ contains an element outside $\mathfrak m_AT$, since it is nonzero and saturated, and its multiples by $A$ have $|A/\mathfrak m_A|^n\ge n$ distinct images in $T/\mathfrak m_A^nT$.
\item[(c)] Call $z\in Z_n$ \emph{primitive} if some $z(g)$ lies outside $\mathfrak m_AT/\mathfrak m_A^nT$. Reduction modulo $\mathfrak m_A^n$ carries primitive elements of $Z_{n+1}$ to primitive elements of $Z_n$. Since $T$ is free and $P$ is saturated and stable under $\rho(s)$, every $z\in Z_{n+1}$ that is not primitive is $\varpi z'$ for some $z'\in Z_n$, which determines $z$. So if $Z_n$ had no primitive element, neither would any $Z_{n'}$ with $n'\ge n$, and then $|Z_{n'}|\le|Z_n|$ for all $n'\ge n$, contrary to (b). Hence every $Z_n$ has a primitive element.
\item[(d)] By (a), (c) and K\"onig's lemma, there are $(z_n,y_n)$ for $n\ge1$, compatible under reduction, with $z_n\in Z_n$ primitive. Since $A$ is complete, their limits $z$ and $y$ exist and have the required properties, and $z\ne0$ since $z_1\ne0$.
\end{enumerate}
\end{proof}
\leanref{Development/AutomorphicArithmetic/AutomorphicVarieties/WeakShimuraTheory/BlockMatrixPlaces/CocT.lean}{UpperTheory.V.Ws.Bl.cocT}
\leanref{Development/AutomorphicArithmetic/AutomorphicVarieties/WeakShimuraTheory/BlockMatrixPlaces/GwCount.lean}{UpperTheory.V.Ws.Bl.gwCount}
\leanref{Development/AutomorphicArithmetic/AutomorphicVarieties/WeakShimuraTheory/BlockMatrixPlaces/PrimEx.lean}{UpperTheory.V.Ws.Bl.primEx}
\leanref{Development/AutomorphicArithmetic/AutomorphicVarieties/WeakShimuraTheory/BlockMatrixPlaces/IsPrecompleteFinite.lean}{UpperTheory.V.Ws.Bl.limit}

\begin{lemma}\label{sha-counts-vd-model-torus-characters}
Let $\mathcal O$ be a local ring, let $X$ be a free abelian group with a basis $b_1,\dots,b_n$, and write $j=\sum_ij_ib_i$ for $j\in X$. For $j\in X$ let $t(j)=\prod_{i=1}^n(1+T_i)^{j_i}$, a unit of $\mathcal O[[T_1,\dots,T_n]]$ with constant term $1$, so that $t(j+j')=t(j)t(j')$. Let $R$ be an $\mathcal O$-algebra, let $s\subseteq X$ be a finite set, and let $x_j\in R$ for $j\in s$ be such that $\sum_{j\in s}t(j)x_j=0$ in $R[[T_1,\dots,T_n]]$.
Then $x_j=0$ for every $j\in s$.
\end{lemma}
\leanref{Development/AutomorphicArithmetic/AutomorphicVarieties/CommutativeAlgebra/TchPoly.lean}{UpperTheory.V.Alg.tchar}
\begin{proof}
Let $N=\sum_{j\in s}\sum_i|j_i|$, and for $j\in s$ let $e(j)=(j_1+N,\dots,j_n+N)$, a vector of nonnegative integers; the map $j\mapsto e(j)$ is injective. Multiplying the relation by the unit $\prod_i(1+T_i)^N$ gives $\sum_{j\in s}x_j\prod_i(1+T_i)^{e(j)_i}=0$, an identity in the polynomial ring $R[T_1,\dots,T_n]$, which embeds in $R[[T_1,\dots,T_n]]$. Substituting $T_i-1$ for $T_i$ gives $\sum_{j\in s}x_j\prod_iT_i^{e(j)_i}=0$, and the coefficient of $\prod_iT_i^{e(j)_i}$ is $x_j$.
\end{proof}

\begin{lemma}\label{sha-counts-vd-model-torus-action}
Let $\mathcal O$ be a local ring with residue field $\mathbf F$, let $X$ be an abelian group, and let $R$ and $R'$ be complete local noetherian $\mathcal O$-algebras with residue field $\mathbf F$. Let $\pi:\mathcal O[[X_1,\dots,X_m]]\to R$ be a surjective homomorphism of $\mathcal O$-algebras with $\pi(X_k)\in\mathfrak m_R$, where $X_k$ has a weight $w_k\in X$, such that for every $w\in X$ the sum of the terms of weight $w$ of every element of $\ker\pi$ lies in $\ker\pi$; let $R_w$ be the image of the series of weight $w$. Let $f:R\to R'$ be a homomorphism of $\mathcal O$-algebras and $t:X\to R'^\times$ a homomorphism.
Then there is a unique homomorphism of $\mathcal O$-algebras $f':R\to R'$ with $f'(y)=t(w)f(y)$ for every $w\in X$ and every $y\in R_w$.
\end{lemma}
\leanref{Development/AutomorphicArithmetic/AutomorphicVarieties/CommutativeAlgebra/TactVanish.lean}{UpperTheory.V.Alg.tact}
\begin{proof}
Let $z_k=\pi(X_k)$. Then $f(z_k)\in\mathfrak m_{R'}$, since $f$ is compatible with the maps to $\mathbf F$, and as $R'$ is complete there is a homomorphism of $\mathcal O$-algebras $\Psi:\mathcal O[[X_1,\dots,X_m]]\to R'$ with $\Psi(X_k)=t(w_k)f(z_k)$. For a series $G$ let $G^{<n}$ be the sum of its terms of degree less than $n$. Then $G-G^{<n}$ lies in the ideal generated by the monomials of degree $n$, so a homomorphism of $\mathcal O$-algebras from $\mathcal O[[X_1,\dots,X_m]]$ that maps the variables into $\mathfrak m_{R'}$, such as $\Psi$ or $f\circ\pi$, maps it into $\mathfrak m_{R'}^n$.
\begin{enumerate}
\item[(i)] If $G$ has weight $w$, then $\Psi(G)=t(w)f(\pi(G))$. Each monomial $X^d$ of $G^{<n}$ has weight $\sum_kd_kw_k=w$, so $\Psi(X^d)=\prod_k\big(t(w_k)f(z_k)\big)^{d_k}=t(w)f(\pi(X^d))$, and $\Psi(G^{<n})=t(w)f(\pi(G^{<n}))$. Hence the two sides are congruent modulo $\mathfrak m_{R'}^n$ for every $n$, and they are equal, since $\bigcap_n\mathfrak m_{R'}^n=0$.
\item[(ii)] $\Psi$ kills $\ker\pi$. Let $F\in\ker\pi$ and $n\ge1$, let $F_w$ be the sum of the terms of $F$ of weight $w$, and let $W_n$ be the finite set of the weights of the monomials of degree less than $n$ that occur in $F$. Then $F-\sum_{w\in W_n}F_w$ has no terms of degree less than $n$, so its image under $\Psi$ lies in $\mathfrak m_{R'}^n$. Each $F_w$ lies in $\ker\pi$, so $\Psi(F_w)=t(w)f(\pi(F_w))=0$ by (i). Hence $\Psi(F)\in\mathfrak m_{R'}^n$ for every $n$, and $\Psi(F)=0$.
\end{enumerate}
So $\Psi=f'\circ\pi$ for a homomorphism of $\mathcal O$-algebras $f':R\to R'$, which has the stated property by (i). If $f''$ also has it, then $f''(z_k)=t(w_k)f(z_k)$, as $z_k\in R_{w_k}$. So $f''\circ\pi$ and $\Psi$ agree on the variables, which they map into $\mathfrak m_{R'}$, and they are equal, as in step (ii) of the proof of Lemma~\ref{sha-counts-vd-model-coproduct}; since $\pi$ is surjective, $f''=f'$.
\end{proof}

\begin{lemma}\label{sha-counts-vd-model-torus-grading}
Let $\mathcal O$ be a complete local noetherian ring with residue field $\mathbf F$, let $X$ be a finitely generated free abelian group, and let $R$ be a complete local noetherian $\mathcal O$-algebra with residue field $\mathbf F$. Below, $R'$ is any complete local noetherian $\mathcal O$-algebra with residue field $\mathbf F$, and a \emph{torus element} of $R'$ is a homomorphism $t:X\to R'^\times$ with $t(j)\equiv1$ modulo $\mathfrak m_{R'}$ for every $j$.
\begin{enumerate}
\item Let $(e_s)_{s\in S}$ be a family of elements of $\mathfrak m_R$, indexed by a set $S$, with weights $\operatorname{wt}(s)\in X$. Suppose that two homomorphisms of $\mathcal O$-algebras $R\to R'$ that agree at every $e_s$ are equal, and that for every homomorphism of $\mathcal O$-algebras $f:R\to R'$ and every torus element $t$ of $R'$ there is a homomorphism of $\mathcal O$-algebras $f':R\to R'$ with $f'(e_s)=t(\operatorname{wt}(s))f(e_s)$ for every $s$. Then there are $s_1,\dots,s_m\in S$ and a surjective homomorphism of $\mathcal O$-algebras $\pi:\mathcal O[[X_1,\dots,X_m]]\to R$ with $\pi(X_k)=e_{s_k}$ such that, if $X_k$ is given the weight $\operatorname{wt}(s_k)$, then for every $w\in X$ the sum of the terms of weight $w$ of every element of $\ker\pi$ lies in $\ker\pi$; that is, $\pi$ is a weighted presentation (Definition~\ref{sha-counts-vd-model}(1)). Moreover $e_s$ lies in the image $R_{\operatorname{wt}(s)}$ of the series of weight $\operatorname{wt}(s)$, for every $s\in S$.
\item Let $\Gamma$ be a group, $\iota$ a finite set, $w:\iota\to X$ a function and $\rho:\Gamma\to M_\iota(R)$ a homomorphism, and let $c(g,a,b)\in\mathcal O$ for $g\in\Gamma$ and $a,b\in\iota$, with $\rho(g)_{ab}-c(g,a,b)\in\mathfrak m_R$, and $c(g,a,b)=0$ whenever $w_a\ne w_b$. Suppose that two homomorphisms of $\mathcal O$-algebras $R\to R'$ that agree on the entries of the $\rho(g)$ are equal, and that for every homomorphism of $\mathcal O$-algebras $f:R\to R'$ and every torus element $t$ of $R'$ there is a homomorphism of $\mathcal O$-algebras $f':R\to R'$ with $f'(\rho(g)_{ab})=t(w_a-w_b)f(\rho(g)_{ab})$ for all $g$, $a$ and $b$. Then $R$ has a weighted presentation $\pi$ as in (1), each of whose variables maps to some $\rho(g)_{ab}-c(g,a,b)$ and has the weight $w_a-w_b$; and $\rho(g)_{ab}$ lies in the part $R_{w_a-w_b}$ for all $g$, $a$ and $b$.
\end{enumerate}
\end{lemma}
\leanref{Development/AutomorphicArithmetic/AutomorphicVarieties/CommutativeAlgebra/CoordGrading.lean}{UpperTheory.V.Alg.coord_grading}
\leanref{Development/AutomorphicArithmetic/AutomorphicVarieties/CommutativeAlgebra/HomogOfStable.lean}{UpperTheory.V.Alg.homog_of_stable}
\leanref{Development/AutomorphicArithmetic/AutomorphicVarieties/CommutativeAlgebra/EigenMem.lean}{UpperTheory.V.Alg.eigen_mem}
\leanref{Development/AutomorphicArithmetic/AutomorphicVarieties/CommutativeAlgebra/LiftGrading.lean}{UpperTheory.V.Alg.lift_grading}
\begin{proof}
(1) For a series $h$ and $w\in X$ let $h_w$ be the sum of the terms of $h$ of weight $w$. Let $t$ be the homomorphism of Lemma~\ref{sha-counts-vd-model-torus-characters} for $\mathcal O$ and $X$, with values in the units of $\mathcal O[[T_1,\dots,T_n]]$; for an $\mathcal O$-algebra $A$, the ring $A[[T_1,\dots,T_n]]$ receives it coefficientwise, and if $A$ is a complete local noetherian $\mathcal O$-algebra with residue field $\mathbf F$, so is $A[[T_1,\dots,T_n]]$, and $t$ is a torus element of it, since $t(j)$ has constant term $1$.
\begin{enumerate}
\item[(i)] \emph{The presentation.} Two homomorphisms $R\to\mathbf F[\varepsilon]$ that agree at every $e_s$ are equal, $\mathbf F[\varepsilon]$ being a complete local noetherian $\mathcal O$-algebra with residue field $\mathbf F$. The argument of the first statement of Lemma~\ref{sha-counts-vd-presentations}(d), with $B=\mathcal O$, which does not use a presentation of $R$, shows that $\mathfrak m_R$ lies in the ideal generated by the $e_s$, $\mathfrak m_{\mathcal O}R$ and $\mathfrak m_R^2$. As $R$ is noetherian, the ideal generated by the $e_s$ is generated by finitely many of them, $e_{s_1},\dots,e_{s_m}$. As $R$ is complete, there is a homomorphism of $\mathcal O$-algebras $\pi:\mathcal O[[X_1,\dots,X_m]]\to R$ with $\pi(X_k)=e_{s_k}$, and the ideal of $R$ generated by the image of the maximal ideal contains the $e_s$ and $\mathfrak m_{\mathcal O}R$. So $\pi$ is surjective, as in step (i) of the proof of the second statement of Lemma~\ref{sha-counts-vd-presentations}(d).
\item[(ii)] \emph{The kernel.} Let $h\in\ker\pi$, $j\in X$ and $N\ge1$. Let $q:R\to\bar R=R/\mathfrak m_R^N$ be the quotient map and $p=q\circ\pi$; then $\bar R$ is a complete local noetherian $\mathcal O$-algebra with residue field $\mathbf F$ and $\mathfrak m_{\bar R}^N=0$. Let $R'=\bar R[[T_1,\dots,T_n]]$, $i:\bar R\to R'$ the inclusion of the constants, and $F':\mathcal O[[X_1,\dots,X_m]]\to R'$ the homomorphism with $F'(X_k)=t(\operatorname{wt}(s_k))\,p(X_k)$, which lies in the ideal $\mathfrak m_{\bar R}R'$. The hypothesis, applied to $i\circ q$ and $t$, gives $f'$ with $f'(e_s)=t(\operatorname{wt}(s))\,i(q(e_s))$ for every $s$. Then $f'\circ\pi$ and $F'$ agree on the variables, which they map into $\mathfrak m_{R'}$, so they are equal, as in step (ii) of the proof of Lemma~\ref{sha-counts-vd-model-coproduct}, and $F'(h)=f'(\pi(h))=0$.
The ideals $\mathfrak m_{\bar R}$ and $\mathfrak m_{\bar R}R'$ have zero $N$-th power, and $p$ and $F'$ map the variables into them. A series $g$ is congruent to the sum of its terms of degree less than $N$ modulo the ideal generated by the monomials of degree $N$, so $p(g)=\sum_dg_dp(X)^d$ and $F'(g)=\sum_dg_dF'(X)^d=\sum_dg_d\,t(\operatorname{wt}(d))\,p(X)^d$, where $d$ runs over the exponent vectors of degree less than $N$, $g_d$ is the coefficient of $X^d$ in $g$, and $\operatorname{wt}(d)=\sum_kd_k\operatorname{wt}(s_k)$. Let $s$ be the finite set of the weights of these exponent vectors. Grouping the terms by weight gives $0=F'(h)=\sum_{k\in s}t(k)\,p(h_k)$ in $\bar R[[T_1,\dots,T_n]]$, so $p(h_k)=0$ for $k\in s$, by Lemma~\ref{sha-counts-vd-model-torus-characters}; and if $j\notin s$, then $p(h_j)$ is an empty sum. So $\pi(h_j)\in\mathfrak m_R^N$ for every $N$, and $\pi(h_j)=0$ by Krull's intersection theorem \cite[Theorem~8.10]{Mat86}. Hence $\pi$ is a weighted presentation.
\item[(iii)] \emph{The parts.} Let $s\in S$, $x=e_s$ and $j=\operatorname{wt}(s)$. First, if $f,f':R\to R'$ and a torus element $t'$ of $R'$ satisfy $f'(e_{s_k})=t'(\operatorname{wt}(s_k))f(e_{s_k})$ for every $k$, then $f'(x)=t'(j)f(x)$. Indeed the hypothesis gives $g'$ with $g'(e_{s'})=t'(\operatorname{wt}(s'))f(e_{s'})$ for every $s'\in S$; then $g'\circ\pi$ and $f'\circ\pi$ agree on the variables, so $g'=f'$, and $f'(x)=g'(x)=t'(j)f(x)$.
Let $R'=R[[T_1,\dots,T_n]]$ and $f:R\to R'$ the inclusion of the constants. Lemma~\ref{sha-counts-vd-model-torus-action}, applied to $\pi$, which is a weighted presentation by (ii), gives $f':R\to R'$ with $f'(y)=t(w)f(y)$ for $y\in R_w$; in particular $f'(e_{s_k})=t(\operatorname{wt}(s_k))f(e_{s_k})$, so $f'(x)=t(j)f(x)$. Write $x=\pi(g)$, and let $N\ge1$. There is a finite set $s'\subseteq X$ containing $j$ such that $g-\sum_{k\in s'}g_k$ has no terms of degree less than $N$; as $f'\circ\pi$ maps the variables into the ideal $\mathfrak m_RR'$, it maps $g-\sum_{k\in s'}g_k$ into $(\mathfrak m_RR')^N$. Reducing the coefficients modulo $\mathfrak m_R^N$, with $q:R\to R/\mathfrak m_R^N$, gives $\sum_{k\in s'}t(k)\,q(\pi(g_k))-t(j)\,q(x)=0$ in $(R/\mathfrak m_R^N)[[T_1,\dots,T_n]]$, since $f'(\pi(g_k))=t(k)f(\pi(g_k))$ and $f'(\pi(g))=t(j)f(x)$. By Lemma~\ref{sha-counts-vd-model-torus-characters}, $q(\pi(g_j))=q(x)$. So $x-\pi(g_j)\in\mathfrak m_R^N$ for every $N$, and $x=\pi(g_j)\in R_j$ by Krull's intersection theorem.
\end{enumerate}
(2) Apply (1) to $S=\Gamma\times\iota\times\iota$ and $e_{(g,a,b)}=\rho(g)_{ab}-c(g,a,b)$, of weight $w_a-w_b$, which lies in $\mathfrak m_R$. Two homomorphisms of $\mathcal O$-algebras that agree at the $e_{(g,a,b)}$ agree on the entries of the $\rho(g)$, so they are equal. Given $f$ and $t$, the hypothesis gives $f'$ with $f'(\rho(g)_{ab})=t(w_a-w_b)f(\rho(g)_{ab})$; then $f'(e_{(g,a,b)})=t(w_a-w_b)f(e_{(g,a,b)})$, because $c(g,a,b)=0$ if $w_a\ne w_b$, and $t(0)=1$ if $w_a=w_b$. So (1) gives the presentation, and $e_{(g,a,b)}\in R_{w_a-w_b}$. Finally $\rho(g)_{ab}=e_{(g,a,b)}+c(g,a,b)$, and $c(g,a,b)$ is $0$ if $w_a\ne w_b$ and lies in $R_0$, which contains $1$, if $w_a=w_b$.
\end{proof}

\begin{lemma}\label{sha-flagged-universal}
Let $D$ be a Siegel datum that is not finite (Definition~\ref{sha-siegel-datum}).
Then there are a complete local noetherian $\mathcal O$-algebra $U$ with residue field $\mathbf F$, a Siegel point $\rho$ of $D$ with values in $U$, elements $g_1,\dots,g_m\in G$ with indices $a_s,b_s\in\{1,2,3,4\}$, and a surjective homomorphism of $\mathcal O$-algebras $\pi:\mathcal O[[Z_1,\dots,Z_m]]\to U$ such that:
\begin{enumerate}
\item $\pi(Z_s)=z_s$, where $z_s=\rho(g_s)_{a_sb_s}-r_0(g_s)_{a_sb_s}$ lies in $\mathfrak m_U$;
\item give $Z_s$ the weight $w_s\in\mathbf Z^2$ of the entry $(a_s,b_s)$; then for each $f\in\ker\pi$ and each $w\in\mathbf Z^2$, the sum of the terms of $f$ of weight $w$ lies in $\ker\pi$;
\item the $\mathcal O$-algebra $U_0$, the image under $\pi$ of the series all of whose monomials have weight $0$, is not finite over $\mathcal O$.
\end{enumerate}
\end{lemma}
\begin{proof}
The Siegel points of $D$ are the images of a universal one: the lifts of $\bar r_0$ are represented by a complete local noetherian $\mathcal O$-algebra \cite{Maz89}, and the conditions defining a Siegel point cut out a quotient $U$ of it, with universal point $\rho$.
By the universal property, two homomorphisms of $\mathcal O$-algebras from $U$ that agree on the entries of the $\rho(g)$ are equal.
Let $R'$ be a complete local noetherian $\mathcal O$-algebra with residue field $\mathbf F$, and $t_1,t_2\in R'^\times$ congruent to $1$ modulo $\mathfrak m_{R'}$. Conjugation by $\operatorname{diag}(t_1t_2,t_2,t_1,1)$ multiplies the entry $(i,j)$ by $t_1^{\lambda_i-\lambda_j}t_2^{\mu_i-\mu_j}$, and it carries Siegel points with values in $R'$ to Siegel points. So, by the universal property, for every $f:U\to R'$ there is $f':U\to R'$ with $f'(\rho(g)_{ij})=t_1^{\lambda_i-\lambda_j}t_2^{\mu_i-\mu_j}f(\rho(g)_{ij})$ for all $g$, $i$ and $j$.
Lemma~\ref{sha-counts-vd-model-torus-grading}(2), applied with $X=\mathbf Z^2$, $w_i=(\lambda_i,\mu_i)$ and $c(g,i,j)=r_0(g)_{ij}$, which vanishes when $(i,j)$ has a nonzero weight since $r_0$ is diagonal, gives the $g_s$, $a_s$, $b_s$ and a presentation $\pi$ with (1) and (2), and shows that the entry $(i,j)$ of every $\rho(g)$ lies in the image $U_w$ of the series of weight $w$, where $w$ is the weight of $(i,j)$.
For (3), let $r$ be a Siegel point with values in $R$, the image of $\rho$ under $f:U\to R$. Since $U_wU_{w'}\subseteq U_{w+w'}$, every product of entries of $\rho$ whose weights add up to $0$ lies in $U_0$, so the $\mathcal O$-algebra generated by the products of entries of $r$ of weight $0$ lies in $f(U_0)$. If $U_0$ were finite over $\mathcal O$, this algebra would be finite over the noetherian ring $\mathcal O$, for every $r$, and $D$ would be finite.
\end{proof}
\leanref{Development/AutomorphicArithmetic/AutomorphicVarieties/WeakShimuraTheory/PolynomialTori/UnivGr.lean}{UpperTheory.V.Ws.Po.univGr}

\begin{lemma}\label{sha-flagged-weight-zero}
Let $\mathcal O$ be as in Definition~\ref{sha-character-parts}, with residue field $\mathbf F$ of characteristic $\ell$.
Let $U$ be a complete local noetherian $\mathcal O$-algebra with residue field $\mathbf F$, and let $\pi:\mathcal O[[Z_1,\dots,Z_m]]\to U$ be a surjective homomorphism of $\mathcal O$-algebras with $z_s=\pi(Z_s)\in\mathfrak m_U$, whose kernel has the property (2) of Lemma~\ref{sha-flagged-universal} for weights $w_1,\dots,w_m\in\mathbf Z^2$ of the $Z_s$.
Let $U_0$ be the image under $\pi$ of the series all of whose monomials have weight $0$, and suppose that $U_0$ is not finite over $\mathcal O$.
Then there are an algebraically closed field $\Omega$ of characteristic $\ell$ containing $\bar{\mathbf F}_\ell$, whose valuation ring $V$ has residue field $\bar{\mathbf F}_\ell$, a complete discrete valuation ring $A'\subset V$ with finite residue field, a local homomorphism of $\mathcal O$-algebras $Y:U\to A'$, and a monomial $v=\prod_sz_s^{n_s}$ with $(n_s)\ne0$ and $\sum_sn_sw_s=0$, such that $Y(v)\ne0$.
\end{lemma}
\begin{proof}
Since $U_0$ is not finite over $\mathcal O$, there are a prime $\mathfrak p\supseteq\mathfrak m_{\mathcal O}U_0$ of $U_0$, an element $t\in\mathfrak m_{U_0}$ not in $\mathfrak p$, and $n\ge1$ with $\mathfrak m_{U_0}^n\subseteq\mathfrak p+tU_0$. Then $U_0/\mathfrak p$ is finite over the power series ring $\mathbf F[[\tau]]$, which maps injectively to it by $\tau\mapsto t$. Let $\Omega$ be an algebraically closed field of characteristic $\ell$, algebraic over the fraction field of $\mathbf F[[\tau]]\subset\Omega$, whose valuation ring $V$ contains $\mathbf F[[\tau]]$ as a local subring and has residue field $\bar{\mathbf F}_\ell$. Since $\Omega$ is algebraically closed, the inclusion of $\mathbf F[[\tau]]$ extends to $U_0/\mathfrak p$, which gives $y:U_0\to V$, with values in a finite extension of the fraction field of $\mathbf F[[\tau]]$. It kills $\mathfrak m_{\mathcal O}$, it is local since $y(\mathfrak m_{U_0})^n\subseteq\tau V$, and $y(t)=\tau\ne0$.
The maximal ideal of $U_0$ is generated by $\mathfrak m_{\mathcal O}$ and the monomials $v=\prod_sz_s^{n_s}$ with $(n_s)\ne0$ and $\sum_sn_sw_s=0$, so $y(v)\ne0$ for one of these.
The images $U_w$ under $\pi$ of the series of weight $w$ form a $\mathbf Z^2$-graded $\mathcal O$-algebra, generated by $U_0$ and the $z_s$, and $y$ extends to a homomorphism $\xi$ from it to $\Omega$. For $t\in(\Omega^\times)^2$, the map that multiplies $\xi$ on $U_w$ by $t^w$ is another such extension.
Every monomial in the $z_s$ of weight $0$ and positive degree lies in $\mathfrak m_{U_0}$, so its image under $\xi$ lies in $\mathfrak m_V$. By Fourier-Motzkin elimination, a theorem of the alternative for the linear inequalities $\operatorname{val}(t^{w_s})+\operatorname{val}\xi(z_s)>0$ in the valuations of $t_1,t_2$, whose values on $\Omega^\times$ are dense, there is $t$ with $t^{w_s}\xi(z_s)\in\mathfrak m_V$ for all $s$.
After this change, a complete discrete valuation ring $A'\subset V$ with finite residue field contains the values of $y$ and the $\xi(z_s)$, hence all values of $\xi$, and the $\xi(z_s)$ lie in $\mathfrak m_{A'}$. Since $A'$ is complete, $\xi$ extends by continuity to $Y:U\to A'$, which is local.
\end{proof}
\leanref{Development/AutomorphicArithmetic/AutomorphicVarieties/WeakShimuraTheory/PolynomialTori/AmbPt.lean}{UpperTheory.V.Ws.Po.ambPt}
\leanref{Development/AutomorphicArithmetic/AutomorphicVarieties/CommutativeAlgebra/WPresMaximalIdealGrZeroEqSupSpanMonomials.lean}{UpperTheory.V.Alg.wPres_maximalIdeal_gr_zero_eq_sup_span_monomials}
\leanref{Development/AutomorphicArithmetic/AutomorphicVarieties/WeakShimuraTheory/PolynomialTori/AmbScale.lean}{UpperTheory.V.Ws.Po.ambScale}
\leanref{Development/AutomorphicArithmetic/AutomorphicVarieties/WeakShimuraTheory/PolynomialTori/MonoMem.lean}{UpperTheory.V.Ws.Po.ptA}

\begin{lemma}\label{sha-flagged-relation-cases}
Let $\lambda=(1,0,1,0)$, as in Definition~\ref{sha-siegel-datum}, and let $\succeq$ be a relation on $\{1,2,3,4\}$ with the following properties:
\begin{enumerate}
\item it is reflexive and transitive;
\item $i\succeq j$ implies $5-j\succeq5-i$;
\item if $\lambda_i\ne\lambda_j$, $i\succeq j$ and $j\succeq i$, then $i\succeq k\succeq i$ for some $k\notin\{i,j\}$.
\end{enumerate}
Then one of the following holds: $i\succeq j$ for all $i$ and $j$; or, for all $i\ne j$, $i\succeq j$ implies $j\not\succeq i$; or $1\succeq3$, $3\succeq1$, and either $i\not\succeq j$ whenever $\lambda_i<\lambda_j$, or $i\not\succeq j$ whenever $\lambda_i>\lambda_j$.
\end{lemma}
\leanref{Development/AutomorphicArithmetic/AutomorphicVarieties/WeakShimuraTheory/CompactIntegralFields/CobOfCoord.lean}{tri_comb}
\begin{proof}
This is a finite check over the relations on a set of four elements: after recording the instances of (2) for the pairs $(i,j)$ equal to $(1,2)$, $(2,1)$, $(1,3)$ and $(3,1)$, and those of (3) for $(1,2)$, $(1,4)$, $(3,2)$ and $(3,4)$, a case analysis on the remaining entries of the relation gives the claim. Remark~\ref{sha-flagged-relation-cases-argument} gives an argument by hand.
\end{proof}

\begin{remark}\label{sha-flagged-relation-cases-argument}
Lemma~\ref{sha-flagged-relation-cases} also has a short argument by hand. Let $\sim$ be the equivalence relation given by $i\succeq j\succeq i$; by (2), $i\mapsto5-i$ permutes its classes. If every class has one element, the second alternative holds. If a class contains $i$ and $j$ with $\lambda_i\ne\lambda_j$, it has at least three elements by (3); if it had three, the remaining index would be a class by itself, fixed by $i\mapsto5-i$, which has no fixed points; so it is $\{1,2,3,4\}$, and the first alternative holds. Otherwise every class with more than one element is $\{1,3\}$ or $\{2,4\}$, and since $i\mapsto5-i$ exchanges these two sets, both are classes. Then $1\succeq3\succeq1$, and $i\succeq j$ cannot hold both for some $i\in\{1,3\}$ and $j\in\{2,4\}$ and for some $i\in\{2,4\}$ and $j\in\{1,3\}$, since the two classes would then merge; this is the third alternative.
\end{remark}

\subsection{Towers and flagged pairs}\label{subsection-sha-towers}

\begin{definition}\label{sha-siegel-tower-ambient}
Let $D$ be a Siegel datum, with the notation of Definition~\ref{sha-siegel-datum}, and let $\iota:\bar{\mathbf Q}_\ell\cong\mathbf C$.
\begin{enumerate}
\item An \emph{ambient ring} is the valuation ring $V$ of an algebraically closed field $\Omega_V$, with residue field $k_V$, given with local homomorphisms $\mathcal O'\to V$ for the rings $\mathcal O'\supseteq\mathcal O$ of Definition~\ref{sha-character-parts}, compatible with the inclusions, such that every element of $V$ is congruent modulo $\mathfrak m_V$ to the image of an element of some $\mathcal O'$. The valuation ring of $\bar{\mathbf Q}_\ell$, with the inclusions, is an ambient ring. A complete discrete valuation subring $A$ of $V$ with finite residue field is always taken to satisfy $\mathfrak m_A\subseteq\mathfrak m_V$, with $\Omega_V$ algebraic over the fraction field of $A$.
\item Let $V$ be an ambient ring. A \emph{flagged representation} over $V$ is a pair $(r,u)$ of a homomorphism $r:G\to M_4(V)$ and a matrix $u\in M_4(V)$, both with entries in a complete discrete valuation subring of $V$ with finite residue field through which $r$ factors continuously, such that: $r\equiv r_0\pmod{\mathfrak m_V}$; $r(g)$ is symplectic with multiplier $\theta(g)$; $r(g)=1$ for $g$ with $\bar r_0(g)=1$ in an inertia group at a place not above $T$, and for $g\in I_p\cap\ker\psi$ with $p\in S_b$; $u$ is symplectic with multiplier $1$; and for $s\in D_\ell$ the matrix $u^*r(s)u$ is upper triangular with diagonal entries congruent to the $\delta_i(s)$ modulo $\mathfrak m_V$, and equal to them when $s\in I_\ell$.
\item The homomorphism $r$ is \emph{of type $\star$} if it is irreducible over $\Omega_V$, or if it is a \emph{non-split block extension}: its entries $(i,j)$ with $i\in\{1,3\}$ and $j\in\{2,4\}$ vanish, the matrices of its entries in rows and columns $2$ and $4$ span $M_2(\Omega_V)$, and no conjugation by $1+Y$, with $Y$ supported on the entries $(i,j)$ with $i\in\{2,4\}$ and $j\in\{1,3\}$, makes it block diagonal.
\item A representation $r:G\to M_4(\bar{\mathbf Q}_\ell)$ is \emph{pure} through $\iota$ if there are an open subgroup $H$ of $G$ and an integer $M\ne0$ such that for every prime $p\nmid M$, every $f\ge1$, every $\sigma\in H$ that acts as $x\mapsto x^{p^f}$ on the residue field of a place above $p$, and all eigenvalues $x$ and $y\ne0$ of $r(\sigma)$, $|\iota(x)|<p^f|\iota(y)|$.
\end{enumerate}
\end{definition}
\leanref{Development/AutomorphicArithmetic/AutomorphicVarieties/WeakShimuraTheory/NearSplit.lean}{UpperTheory.V.Ws.Amb}
\leanref{Development/AutomorphicArithmetic/AutomorphicVarieties/WeakShimuraTheory/NearSplit.lean}{UpperTheory.V.Ws.KPair}
\leanref{Development/AutomorphicArithmetic/AutomorphicVarieties/WeakShimuraTheory/RedIdeal.lean}{UpperTheory.V.Ws.Star}
\leanref{Development/AutomorphicArithmetic/Foundations/GaloisRepresentations/AdmissibleResidues/ArithmeticPointDecomposition.lean}{UpperTheory.GalN.V.PureW}

\begin{definition}\label{sha-siegel-tower-towers}
Let $D$ be a Siegel datum (Definition~\ref{sha-siegel-datum}), with its notation.
For a subfield $E\subseteq\bar{\mathbf Q}$ write $G_E=\operatorname{Gal}(\bar{\mathbf Q}/E)\subseteq G$, and for a place $w$ of $\bar{\mathbf Q}$ write $D_w(E)$ and $I_w(E)$ for the decomposition and inertia groups of $G_E$ at $w$; bars denote reduction modulo $\mathfrak m_{\mathcal O}$.
\begin{enumerate}
\item The datum $D$ is \emph{of case II}, respectively \emph{of case III}, if it is attached to $(h,k,\psi)$ (Definition~\ref{sha-siegel-datum}) for a primitive Dirichlet character $\psi_0$, a host datum $h$ for $\psi_0$ that is large for some $t$, an integer $0<k<t$ and $\psi\in X_h(k)$ such that $(\ell,k,\psi)$ is in case II, respectively in case III (Definition~\ref{sha-host-datum}, with $\psi_0$ in place of $\psi$).
\item A \emph{relabelling} is an element $\epsilon\in\{0,1\}$. Its characters $(\chi_1,\chi_2,\chi_3,\chi_4)$ are $(\delta_1,\delta_2,\delta_3,\delta_4)$ if $\epsilon=0$ and $(\delta_3,\delta_4,\delta_1,\delta_2)$ if $\epsilon=1$. For $w=(a,b)\in\mathbf Z^2$ put $\chi^w=(\chi_1\chi_2^{-1})^a(\chi_1\chi_3^{-1})^b$; then $\chi_i\chi_j^{-1}=\chi^w$ for the weight $w$ of the entry $(i,j)$. The \emph{roots} are the weights $\pm(1,0)$, $\pm(0,1)$, $\pm(1,1)$, $\pm(-1,1)$ of the off-diagonal entries, and the \emph{odd roots} are the roots other than $\pm(0,1)$.
\item A \emph{tower} $T$ for $D$ consists of a relabelling $\epsilon$, an imaginary quadratic field $K$, a totally real number field $F$, Galois over $\mathbf Q$ with abelian Galois group, of degree $d_T$, the field $L=FK$, which is CM, an element $g_0^T\in G_F$, and a finite set $\Sigma_T$ of places of $\bar{\mathbf Q}$ above $\ell$, such that every place of $\bar{\mathbf Q}$ above $\ell$ is conjugate under $G_F$ to exactly one element of $\Sigma_T$, and every $v\in\Sigma_T$ induces on $K$ the place induced by $v_\ell$.
\item The tower $T$ is \emph{fit} if it satisfies the conditions (a) to (m) below, and either $D$ is of case II and $T$ satisfies the conditions (n) to (r) below, or $D$ is of case III, $\ell\ge5$, the set $S_b$ is empty and $\psi^m=1$ for some $m$ prime to $\ell$. Here $w$ denotes a place of $\bar{\mathbf Q}$ and $\mathbf F$ is the residue field of $\mathcal O$:
(a) $\epsilon=0$, or $\bar\Psi$ is trivial on $D_\ell$;
(b) $d_T$ is even;
(c) $I_w(F)\subseteq G_L$ for every $w$;
(d) $D_w(\mathbf Q)\subseteq G_K$ for every $w$ above a prime of $S$;
(e) for every $w'\in\mathbf Z^2$ some $g\in G_L$ has $\bar\chi^{w'}(g)\ne\bar\omega(g)$;
(f) for every root $w'$, $\bar\chi^{w'}$ is nontrivial on $G_L$;
(g) every homomorphism $\eta:G_L\to\mathbf F^\times$ with $\eta^2=1$ that is trivial on $G_L\cap\ker\omega$, and such that for each $i$ there is $j$ with $\eta\bar\chi_i=\bar\chi_j$ on $G_L$, is trivial;
(h) for some root $w'$, the restriction of $\bar\chi^{w'}$ to $G_L$ has order greater than $4$;
(i) $\psi(g)=1$ for every $g\in I_w(F)$ with $w$ not above $\ell$ and $\bar r_0(g)=1$;
(j) there is $m$ prime to $\ell$ with $\psi(g)^m=1$ for every $g\in G_F$ with $\chi_{\mathrm{cyc}}(g)=1$;
(k) $\bar r_0$ is trivial on $I_w(F)$ for every $w$ above $\ell$;
(l) the residues $\bar\chi_1(g_0^T),\dots,\bar\chi_4(g_0^T)$ are pairwise distinct;
(m) for every $w$ above $\ell$, either $\bar\omega$ is nontrivial on $I_w(F)$, or some $g\in D_w(F)$ has $\bar\omega(g)=1$, $\bar\chi^{(0,1)}(g)=1$ and $\bar\chi^{(1,0)}(g)\ne1$;
and the conditions (n) to (r) are:
(n) $(\bar\chi^{(0,1)})^2=1$ on $G_F$;
(o) $\bar\chi^{(0,1)}$ is trivial on $I_w(F)$ for every $w$, and on $D_w(F)$ for every $w$ above $\ell$;
(p) there is a real quadratic field $k_0$ in which $2$ and $\ell$ split, such that for $g\in G_F$, $\bar\chi^{(0,1)}(g)=1$ if and only if $g$ fixes $k_0$;
(q) $\xi$ has order a power of $2$, and is trivial on $I_w(\mathbf Q)$ for every $w$ not above $2$;
(r) for every $w$ above $\ell$ some $g\in D_w(F)$ has $\bar\chi^{(1,0)}(g)^4\ne1$, and for every $w$ above $2$ some $g\in I_w(F)$ has $\bar\chi^{(1,0)}(g)^4\ne1$.
\item The tower $T$ is \emph{deep} if $\psi^m$ is trivial on $G_F$ for some $m$ prime to $\ell$, and \emph{tame} if for every place $w$ above $\ell$ and every $g\in I_w(\mathbf Q)$ some power $g^m$ with $m$ prime to $\ell$ lies in $G_F$.
\item For towers $T$ and $T'$, with fields $F$ and $F'$, write $T\le T'$ if they have the same relabelling and the same field $K$, $F\subseteq F'$, there is $j\ge0$ with $g^{\ell^j}\in G_{F'}$ for all $g\in G_F$, and $g_0^{T'}=(g_0^T)^{\ell^{j'}}$ in $G$ for some $j'\ge0$. The tower $T'$ is a \emph{layer} of $T$, of index $m\ge0$, if $G_{F'}$ is the set of $g\in G_F$ with $\chi_{\mathrm{cyc}}(g)^{\ell-1}\equiv1\pmod{\ell^{m+1}}$.
\end{enumerate}
\end{definition}
\leanref{Development/AutomorphicArithmetic/AutomorphicVarieties/WeakShimuraTheory/NearSplit.lean}{UpperTheory.V.Ws.Ctx2}
\leanref{Development/AutomorphicArithmetic/AutomorphicVarieties/WeakShimuraTheory/NearSplit.lean}{UpperTheory.V.Ws.Tower}
\leanref{Development/AutomorphicArithmetic/AutomorphicVarieties/WeakShimuraTheory/NearSplit.lean}{UpperTheory.V.Ws.chi}
\leanref{Development/AutomorphicArithmetic/AutomorphicVarieties/WeakShimuraTheory/NearSplit.lean}{UpperTheory.V.Ws.Tower.Ok}
\leanref{Development/AutomorphicArithmetic/AutomorphicVarieties/WeakShimuraTheory/NearSplit.lean}{UpperTheory.V.Ws.Tower.Two}
\leanref{Development/AutomorphicArithmetic/AutomorphicVarieties/WeakShimuraTheory/NearSplit.lean}{UpperTheory.V.Ws.Tower.Fit}
\leanref{Development/AutomorphicArithmetic/AutomorphicVarieties/WeakShimuraTheory/NearSplit.lean}{UpperTheory.V.Ws.Tower.Deep}
\leanref{Development/AutomorphicArithmetic/AutomorphicVarieties/WeakShimuraTheory/NearSplit.lean}{UpperTheory.V.Ws.TameL}
\leanref{Development/AutomorphicArithmetic/AutomorphicVarieties/WeakShimuraTheory/NearSplit.lean}{UpperTheory.V.Ws.Tower.Le}
\leanref{Development/AutomorphicArithmetic/AutomorphicVarieties/WeakShimuraTheory/NearSplit.lean}{UpperTheory.V.Ws.IsLayer}

\begin{definition}\label{sha-siegel-tower-pairs}
Let $D$ be a Siegel datum, with the notation of Definitions~\ref{sha-siegel-datum} and~\ref{sha-siegel-tower-towers}.
A matrix is \emph{supported in} a set $W$ of weights if its off-diagonal entries of weight outside $W$ vanish; a homomorphism $r$ with values in $M_4(R)$ is supported in $W$ if every $r(g)$ is, and it \emph{has an entry of weight} $w$ if $r(g)_{ij}\ne0$ for some $g$ and some entry $(i,j)$ of weight $w$.
\begin{enumerate}
\item A \emph{deformation problem} $\Phi$ is one of the following two. The \emph{rational problem} of a relabelling $\epsilon$ and an element $g_0\in G$ has the field $E=\mathbf Q$; the one place $v_\ell$, with $D_{v_\ell}=D_\ell$ and $\mu_{v_\ell}=\emptyset$; the set $N_\Phi$ of the $g\in G$ that lie in an inertia group at a place not above $T$ and satisfy $\bar r_0(g)=1$, or lie in $I_p\cap\ker\psi$ for some $p\in S_b$; the characters $\chi_i$ of $\epsilon$; and $g_0$. The \emph{problem of a tower} $T$ has the field $E=F$; the places $v\in\Sigma_T$, with $D_v=D_v(F)$ and $\mu_v$ the set of $s\in D_v$ whose image in the quotient of $D_v$ by the closure of its commutator subgroup has order a power of $\ell$; the set $N_\Phi$ of the $g\in G_F$ that lie in $I_w(F)$ for a place $w$ not above $\ell$ and satisfy $\bar r_0(g)=1$; the restrictions to $G_F$ of the characters $\chi_i$ of the relabelling of $T$; and $g_0^T$.
\item Let $\Phi$ be a deformation problem with field $E$ and places $v$, let $R$ be a commutative ring and $a:\mathcal O\to R$ a ring homomorphism. A \emph{flagged pair} for $\Phi$ along $a$ is a pair $P=(r,(n_v)_v)$ of a homomorphism $r:G_E\to M_4(R)$ and matrices $n_v\in M_4(R)$ such that: $r(g)$ is symplectic with multiplier $a(\theta(g))$; $r(g)=1$ for $g\in N_\Phi$; $r(g_0)$ is diagonal; $n_v$ is symplectic with multiplier $1$ and lower triangular with diagonal entries $1$; for $s\in D_v$ the matrix $n_v^*r(s)n_v$ is upper triangular, with diagonal entries $\psi_{v,1}(s),\dots,\psi_{v,4}(s)$ say; and $\psi_{v,i}(s)=1$ for $s\in\mu_v$. The image of $P$ under a ring homomorphism $f:R\to R'$ is the flagged pair $(f\circ r,(f(n_v))_v)$ along $f\circ a$.
\item A flagged pair $P=(r,(n_v))$ for $\Phi$ along $a$ is \emph{residual} if the entries of $r$ are locally constant on $G_E$, $r(g)_{ii}=a(\chi_i(g))$ for $g\in G_E$, $\psi_{v,i}(s)=a(\chi_i(s))$ for $s\in D_v$, and either $r$ is supported in $\{(-1,1),(-1,0),(-1,-1),(0,-1)\}$ and each $n_v$ in $\{(-1,0),(0,-1),(-1,-1)\}$, or $r$ is supported in $\{(-1,0),(1,-1),(0,-1),(-1,-1)\}$. It is \emph{of type V} if $r$ is supported in $\{(-1,1),(-1,0),(-1,-1)\}$ and each $n_v$ is supported in $\{(-1,0),(0,-1),(-1,-1)\}$ and has a nonzero entry $(3,1)$. A residual flagged pair is \emph{ample} if one of the following holds: $r$ has an entry of weight $(0,-1)$ and an entry of weight $(-1,1)$ or $(-1,0)$; $r$ has entries of weights $(-1,0)$ and $(1,-1)$; $P$ is of type V and $r$ has an entry of weight $(-1,0)$; or $P$ is of type V, $r$ has entries of weights $(-1,1)$ and $(-1,-1)$, and for each $v$ some $s\in D_v$ has $a(\chi^{(1,0)}(s))\ne1$.
\item Residual flagged pairs $P$ and $P'$ for $\Phi$ along $a$ are \emph{linked} if there is a residual flagged pair for $\Phi$ along the composite of $a$ with $R\subset R[X]$ whose images under $X\mapsto0$ and $X\mapsto1$ are $P$ and $P'$. A \emph{chain} from $P$ to $P'$ is a sequence $P=P_0,P_1,\dots,P_m=P'$ of flagged pairs such that $P_0$ is ample, and each $P_{i}$ with $i\ge1$ is ample and linked to $P_{i-1}$.
\item Let $\mathcal O'\supseteq\mathcal O$ be one of the rings of Definition~\ref{sha-character-parts}, with residue field $k'$, and let $c$ be a residual flagged pair for $\Phi$ along $\mathcal O\to k'$. For a complete local noetherian $\mathcal O'$-algebra $R$ with residue field $k'$, a \emph{lift} of $c$ to $R$ is a flagged pair for $\Phi$ along $\mathcal O\to R$ whose entries are continuous for the $\mathfrak m_R$-adic topology and whose image in $k'$ is $c$. A lift $P$ to $R$ is \emph{universal} if for every lift $P'$ of $c$ to every such ring $R'$ there is a unique homomorphism of $\mathcal O'$-algebras $f:R\to R'$ that carries $P$ to $P'$.
\item Let $V$ be an ambient ring (Definition~\ref{sha-siegel-tower-ambient}), with residue field $k_V$, and let $\epsilon$ be a relabelling and $g_0\in G$. A \emph{rational flagged pair} for $(\epsilon,g_0)$ is a flagged pair $P=(r,n)$ for the rational problem of $(\epsilon,g_0)$ along $\mathcal O\to V$, with entries in a complete discrete valuation subring of $V$ with finite residue field through which $r$ factors continuously, such that $\psi_{v_\ell,i}(s)=\delta_i(s)$ for $s\in I_\ell$. It is \emph{near split} if $r\equiv\operatorname{diag}(\chi_1,\dots,\chi_4)$ and $\psi_{v_\ell,i}(s)\equiv\chi_i(s)$ for $s\in D_\ell$, modulo $\mathfrak m_V$, and the off-diagonal entries of $n$ of weight other than $(0,-1)$ lie in $\mathfrak m_V$.
\item Let $k$ be a field and $a:\mathcal O\to k$ a homomorphism. A flagged pair $c=(r,n)$ for the rational problem of $(\epsilon,g_0)$ along $a$ is \emph{of lattice shape} if it is residual, $r$ is supported in $\{(-1,1),(-1,0),(-1,-1),(0,-1)\}$, $n$ is supported in $\{(-1,0),(0,-1),(-1,-1)\}$, $r$ has an entry of weight $(-1,1)$, and either $r$ has an entry of weight $(0,-1)$, or: $r$ is supported in $\{(-1,1),(-1,0),(-1,-1)\}$, the entry $(3,1)$ of $n$ is nonzero, $\bar\Psi$ is trivial on $D_\ell$, and there is no $m\in k$ with $r(s)_{23}=m\big(a(\chi_3(s))-a(\chi_2(s))\big)$ for all $s\in D_\ell$.
A flagged pair $c'=(r',(n'_v))$ for the problem of a tower $T$ along $a$ is \emph{over} $c$ if $r'$ is the restriction of $r$ to $G_F$ and for each $v\in\Sigma_T$ there is $g\in G$ carrying $v_\ell$ to $v$ such that $n_v'^*\,r(g)\,n$ is upper triangular.
\end{enumerate}
\end{definition}
\leanref{Development/AutomorphicArithmetic/AutomorphicVarieties/WeakShimuraTheory/RedIdeal.lean}{UpperTheory.V.Ws.Fr}
\leanref{Development/AutomorphicArithmetic/AutomorphicVarieties/WeakShimuraTheory/NearSplit.lean}{UpperTheory.V.Ws.frQ}
\leanref{Development/AutomorphicArithmetic/AutomorphicVarieties/WeakShimuraTheory/NearSplit.lean}{UpperTheory.V.Ws.Tower.fr}
\leanref{Development/AutomorphicArithmetic/AutomorphicVarieties/WeakShimuraTheory/RedIdeal.lean}{UpperTheory.V.Ws.Fr.IsPair}
\leanref{Development/AutomorphicArithmetic/AutomorphicVarieties/WeakShimuraTheory/RedIdeal.lean}{UpperTheory.V.Ws.Fr.IsRes}
\leanref{Development/AutomorphicArithmetic/AutomorphicVarieties/WeakShimuraTheory/RedIdeal.lean}{UpperTheory.V.Ws.Fr.Res}
\leanref{Development/AutomorphicArithmetic/AutomorphicVarieties/WeakShimuraTheory/RedIdeal.lean}{UpperTheory.V.Ws.Fr.InV}
\leanref{Development/AutomorphicArithmetic/AutomorphicVarieties/WeakShimuraTheory/RedIdeal.lean}{UpperTheory.V.Ws.Fr.Good}
\leanref{Development/AutomorphicArithmetic/AutomorphicVarieties/WeakShimuraTheory/RedIdeal.lean}{UpperTheory.V.Ws.Fr.Linked}
\leanref{Development/AutomorphicArithmetic/AutomorphicVarieties/WeakShimuraTheory/RedIdeal.lean}{UpperTheory.V.Ws.Fr.Chain}
\leanref{Development/AutomorphicArithmetic/AutomorphicVarieties/WeakShimuraTheory/RedIdeal.lean}{UpperTheory.V.Ws.Fr.Lifts}
\leanref{Development/AutomorphicArithmetic/AutomorphicVarieties/WeakShimuraTheory/RedIdeal.lean}{UpperTheory.V.Ws.Fr.IsDef}
\leanref{Development/AutomorphicArithmetic/AutomorphicVarieties/WeakShimuraTheory/NearSplit.lean}{UpperTheory.V.Ws.QPair}
\leanref{Development/AutomorphicArithmetic/AutomorphicVarieties/WeakShimuraTheory/NearSplit.lean}{UpperTheory.V.Ws.NearSplit}
\leanref{Development/AutomorphicArithmetic/AutomorphicVarieties/WeakShimuraTheory/NearSplit.lean}{UpperTheory.V.Ws.LatShape}
\leanref{Development/AutomorphicArithmetic/AutomorphicVarieties/WeakShimuraTheory/NearSplit.lean}{UpperTheory.V.Ws.Over}

\begin{definition}\label{sha-siegel-tower-space}
Let $D$ be a Siegel datum, $T$ a tower for $D$, and $\mathcal O'\supseteq\mathcal O$ one of the rings of Definition~\ref{sha-character-parts}, with residue field $k'$; flagged pairs are for the problem of $T$ (Definition~\ref{sha-siegel-tower-pairs}).
\begin{enumerate}
\item A \emph{deformation space} $X$ for $T$ over $\mathcal O'$ consists of $\mathcal O'$-algebras $B$ and $A$ with an injective homomorphism of $\mathcal O'$-algebras $B\to A$, where $B$ is complete local noetherian with residue field $k'$ and $A$ is of finite type over $B$, and a flagged pair $P_X=(r_X,(n_{X,v}))$ along $\mathcal O\to A$ whose entries are continuous for the $\mathfrak m_BA$-adic topology, such that:
(1) $A$ is generated as a $B$-algebra by the entries of the $r_X(g)$ and of the $n_{X,v}$;
(2) for $g\in G_F$ there is $c_g\in B[X]$, with coefficients continuous in $g$, that maps to the characteristic polynomial of $r_X(g)$ and reduces modulo $\mathfrak m_B$ to $\prod_i(X-\bar\chi_i(g))$;
(3) for each $v$, $i$ and $s\in D_v$ there is $\psi^B_{v,i}(s)\in B$, continuous in $s$, that maps to the diagonal entry $\psi_{v,i}(s)$ of $P_X$ and has residue $\bar\chi_i(s)$;
(4) two homomorphisms of $\mathcal O'$-algebras from $B$ to a complete local noetherian $\mathcal O'$-algebra with residue field $k'$ that agree at the coefficients of the $c_g$ and at the $\psi^B_{v,i}(s)$ are equal;
(5) for every residual flagged pair $c$ along $\mathcal O\to k'$ and every lift $Q$ of $c$ to a complete local noetherian $\mathcal O'$-algebra $R$ with residue field $k'$, some homomorphism of $\mathcal O'$-algebras $A\to R$ carries $P_X$ to $Q$;
(6) every residual flagged pair along $\mathcal O\to k'\subset k'[X]$ is the image of $P_X$ under a homomorphism of $\mathcal O'$-algebras $A\to k'[X]$ that kills $\mathfrak m_B$.
A \emph{point} of $X$ is a homomorphism of $\mathcal O'$-algebras $x:A\to k'$, and its \emph{image} is the image of $P_X$ under $x$.
\item Let $J\subseteq B$ be an ideal and $\mathfrak p\subseteq A$ a prime. Then $J$ \emph{lies in some minimal prime below} $\mathfrak p$ if some minimal prime $\mathfrak q$ of $A$ with $\mathfrak q\subseteq\mathfrak p$ contains $JA$, and in \emph{all minimal primes below} $\mathfrak p$ if every such $\mathfrak q$ contains $JA$. A ring homomorphism $y$ from $A$ \emph{detects} $J$ if every prime of $B$ whose image in $A$ is killed by $y$ contains $J$.
\end{enumerate}
\end{definition}
\leanref{Development/AutomorphicArithmetic/AutomorphicVarieties/WeakShimuraTheory/RedIdeal.lean}{UpperTheory.V.Ws.Space}
\leanref{Development/AutomorphicArithmetic/AutomorphicVarieties/WeakShimuraTheory/RedIdeal.lean}{UpperTheory.V.Ws.Space.SomeIn}
\leanref{Development/AutomorphicArithmetic/AutomorphicVarieties/WeakShimuraTheory/RedIdeal.lean}{UpperTheory.V.Ws.Space.AllIn}
\leanref{Development/AutomorphicArithmetic/AutomorphicVarieties/WeakShimuraTheory/RedIdeal.lean}{UpperTheory.V.Ws.Space.LocZero}

\begin{definition}\label{sha-siegel-tower-ideal}
Let $D$ be a Siegel datum, $T$ a tower for $D$, with fields $K\subset L$ and $F\subset L$, and let $\iota:\bar{\mathbf Q}_\ell\cong\mathbf C$; flagged pairs are for the problem of $T$.
For $v\in\Sigma_T$ let $U_v$ be the set of $s\in I_v(F)$ with $\chi_{\mathrm{cyc}}(s)=1$ and $\bar r_0(s)=1$.
For a homomorphism $r:G_F\to M_4(\bar{\mathbf Q}_\ell)$ and matrices $n_v\in M_4(\bar{\mathbf Q}_\ell)$, $v\in\Sigma_T$, let $\psi_{v,i}(s)$ be the entry $(i,i)$ of $n_v^*r(s)n_v$; the pair $(r,(n_v))$ has \emph{inertial weights} $h\in\mathbf Z^4$ if there is an open subgroup $U$ of $G_F$ with $\psi_{v,i}(s)^2=\chi_{\mathrm{cyc}}(s)^{h_i}$ for all $v$, $i$ and $s\in U\cap I_v(F)$.
\begin{enumerate}
\item A \emph{seed} for $T$ is such a pair $(r,(n_v))$ such that the $r(g)$ with $g\in G_L\cap\ker\chi_{\mathrm{cyc}}$ span $M_4(\bar{\mathbf Q}_\ell)$, and there are: $h\in\mathbf Z^4$ with $h_i-h_{i+1}$ even and at least $4$, such that the pair has inertial weights $h$; a continuous character $\kappa$ of $G_L$, an integer $a$ with $a+h_1$ even, and $k\ge1$, with $(\kappa(s)^2\chi_{\mathrm{cyc}}(s)^{-a})^k=1$ for $s$ in the inertia groups of $G_L$ at the places above $\ell$; and a regular algebraic cuspidal automorphic representation $\pi$ of $GL_4(\mathbf A_L)$, for which there is $t\in\mathbf R$ such that at every finite place $u$ of $L$ the supercuspidal support of $\pi_u$ consists of characters $\xi$ of $L_u^\times$ with $q_u^{t-1/2}<|\xi(\varpi_u)|<q_u^{t+1/2}$, such that the twist of the restriction of $r$ to $G_L$ by $\kappa$ is attached to $\pi$ through $\iota$.
An \emph{integral seed} over a ring $\mathcal O_1\supseteq\mathcal O$ of Definition~\ref{sha-character-parts} is a flagged pair along $\mathcal O\subseteq\mathcal O_1$ with $\mathfrak m_{\mathcal O_1}$-adically continuous entries, whose image in the residue field $k_1$ of $\mathcal O_1$ satisfies the last three conditions of a residual flagged pair, and whose image in $\bar{\mathbf Q}_\ell$ is a seed.
\item Let $X=(A,B,P_X)$ be a deformation space for $T$ over $\mathcal O'$ (Definition~\ref{sha-siegel-tower-space}). A \emph{modularity ideal} for $X$ is an ideal $J\subseteq B$ with the following properties, where $m=(k+1,k,1,0)$ for the integer $k$ of $D$.
\begin{itemize}
\item (Finiteness.) Some power of $\mathfrak m_B$ lies in the ideal generated by $J$, $\mathfrak m_{\mathcal O'}$ and the elements $\psi^B_{v,i}(s)-1$ with $s\in I_v(F)$.
\item (Patching.) Let $\kappa$ be a finite field with an $\mathcal O'$-algebra structure and $y:A\to\kappa[[\tau]]$ a homomorphism of $\mathcal O'$-algebras that maps $\mathfrak m_B$ into $\tau\kappa[[\tau]]$ and kills $J$, such that for each $v$ and $i\ne j$ some $s\in U_v$ has $y(\psi_{v,i}(s))\ne y(\psi_{v,j}(s))$, and the $y(r_X(g))$ span $M_4$ over the fraction field of $\kappa[[\tau]]$. Then $y$ detects $J$.
\item (Seeds.) Every homomorphism of $\mathcal O'$-algebras $y:A\to\bar{\mathbf Q}_\ell$ that carries $P_X$ to a seed detects $J$.
\item (Purity and automorphy.) Let $y:A\to\bar{\mathbf Q}_\ell$ be a homomorphism of $\mathcal O'$-algebras that kills $J$ and carries $P_X$ to a flagged pair with inertial weights $2m$. Every homomorphism $r:G\to M_4(\bar{\mathbf Q}_\ell)$ whose restriction to $G_F$ is $y\circ r_X$ is pure through $\iota$ (Definition~\ref{sha-siegel-tower-ambient}); if moreover $r$ is continuous, $r(g)$ is symplectic with multiplier $\theta(g)$ for all $g$, and $r$ is irreducible on $G_L$, then $r$ is automorphic through $\iota$ (Definition~\ref{sha-siegel-datum}).
\end{itemize}
\item Let $k$ be a field and $a:\mathcal O\to k$ a homomorphism that kills $\mathfrak m_{\mathcal O}$. For $w\in\mathbf Z^2$ let $Z_w$ be the $k$-space of locally constant $z:G_F\to k$ with $z(gh)=z(g)+a(\chi^w(g))z(h)$, vanishing on the set $N_\Phi$ of the problem of $T$, and with $z(g_0^T)=0$ if $w\ne0$; let $S_w\subseteq Z_w$ consist of the $z$ for which each $v\in\Sigma_T$ has an $m_v\in k$ with $z(s)=m_v-a(\chi^w(s))m_v$ for $s\in D_v(F)$; and let $Z_{v,w}$ be the $k$-space of locally constant $z:D_v(F)\to k$ with $z(gh)=z(g)+a(\chi^w(g))z(h)$. For integers $d,C\ge0$, the tower $T$ satisfies the \emph{cohomology bounds} $(d,C)$ along $a$ if $\Sigma_T$ has at most $C$ elements, $\dim S_w\le C$ for every root $w$, $\dim Z_w\le C$ for $w\in\{0,(0,1),(0,-1)\}$, and $\dim Z_w\le d+C$ and $\sum_v\dim Z_{v,w}\le d+C$ for every odd root $w$.
\end{enumerate}
\end{definition}
\leanref{Development/AutomorphicArithmetic/AutomorphicVarieties/WeakShimuraTheory/NearSplit.lean}{UpperTheory.V.Ws.Tower.U1}
\leanref{Development/AutomorphicArithmetic/AutomorphicVarieties/WeakShimuraTheory/NearSplit.lean}{UpperTheory.V.Ws.Tower.Wts}
\leanref{Development/AutomorphicArithmetic/AutomorphicVarieties/ModularityArgument/ThresholdArguments/DenseThresholdCases/PotentialModularityTower.lean}{UpperTheory.V.Ws.Tower.IsSeedAlt}
\leanref{Development/AutomorphicArithmetic/AutomorphicVarieties/ModularityArgument/ThresholdArguments/DenseThresholdCases/PotentialModularityTower.lean}{UpperTheory.V.Ws.Tower.SeedPairAlt}
\leanref{Development/AutomorphicArithmetic/AutomorphicVarieties/ModularityArgument/ThresholdArguments/DenseThresholdCases/PotentialModularityTower.lean}{UpperTheory.V.Ws.Tower.PotentialModularityTower}
\leanref{Development/AutomorphicArithmetic/AutomorphicVarieties/WeakShimuraTheory/RedIdeal.lean}{UpperTheory.V.Ws.Fr.Bounds}

\begin{definition}\label{sha-siegel-step-towers-input}
Let $D$ be a Siegel datum, with the notation of Definitions~\ref{sha-siegel-datum} and~\ref{sha-siegel-tower-towers}; $\chi_{\mathrm{cyc}}$ denotes the cyclotomic character with values in $\mathbf Z_\ell^\times\subseteq\mathcal O^\times$.
\begin{enumerate}
\item Let $S'$ be a finite set of primes, $p\notin S'$ a prime and $s\in\{1,-1\}$. A \emph{quadratic character at $p$ of sign $s$ split over $S'$} is a character $\eta:G\to\{\pm1\}$ with open kernel of index $2$ such that: $\eta$ is trivial on $I_w(\mathbf Q)$ for every place $w$ of $\bar{\mathbf Q}$ not above $p$; for every place $w$ above $p$ some $g\in I_w(\mathbf Q)$ has $\eta(g)=-1$; $\eta(\sigma)=s$ for every complex conjugation $\sigma\in G$; and $\eta$ is trivial on $D_w(\mathbf Q)$ for every place $w$ above a prime of $S'$.
\item An \emph{auxiliary character} for $D$ is a character $\kappa:G\to\mathcal O^\times$ with open kernel and $\kappa^{\ell^j}=1$ for some $j\ge0$, trivial on $I_w(\mathbf Q)$ for every place $w$ above no prime of $S$ and for every place $w$ above $\ell$, such that $\psi(g)=1$ for every $g\in I_w(\mathbf Q)$ with $w$ not above $\ell$, $\bar\Psi(g)=1$ and $\kappa(g)=1$, and such that for some $m$ prime to $\ell$, $\psi(g)^m=1$ for every $g\in G$ with $\chi_{\mathrm{cyc}}(g)=1$ and $\kappa(g)=1$.
\item A \emph{tower input} $X$ for $D$ consists of an auxiliary character $\kappa$, primes $p_3,p,p'$ and characters $\varphi,\eta,\eta':G\to\{\pm1\}$, such that $\varphi$ is a quadratic character at $p_3$ of sign $1$ split over $S$, $\eta$ one at $p$ of sign $-1$ split over $S\cup\{p_3\}$, and $\eta'$ one at $p'$ of sign $-1$ split over $S\cup\{p_3,p\}$, and such that: $\Psi(\sigma)=1$ for every complex conjugation $\sigma$; $\xi$ has order a power of $2$ and is trivial on $I_w(\mathbf Q)$ for every place $w$ not above $2$; some $t\in G$ has $\omega(t)=1$, $\kappa(t)=1$ and $\bar\Psi(t)=-1$; for every place $w$ above $\ell$ some $g\in D_w(\mathbf Q)$ has $\bar\Psi(g)=1$, $\kappa(g)=1$ and $\bar\xi(g)^4\ne1$; and for every place $w$ above $2$ some $g\in I_w(\mathbf Q)$ has $\bar\Psi(g)=1$, $\kappa(g)=1$ and $\bar\xi(g)^4\ne1$.
Its \emph{subgroup} $H_X$ is the set of $g\in G$ with $\bar\Psi(g)=\varphi(g)$ (the image of $\pm1$ in $\mathbf F$), $\eta(g)=\eta'(g)$ and $\kappa(g)=1$.
\end{enumerate}
\end{definition}
\leanref{Development/AutomorphicArithmetic/AutomorphicVarieties/WeakShimuraTheory/FrobeniusFields/ExistsResQ.lean}{UpperTheory.V.Ws.QuadAt}
\leanref{Development/AutomorphicArithmetic/AutomorphicVarieties/WeakShimuraTheory/FrobeniusFields/ExistsResQ.lean}{UpperTheory.V.Ws.Kap}
\leanref{Development/AutomorphicArithmetic/AutomorphicVarieties/WeakShimuraTheory/FrobeniusFields/ExistsResQ.lean}{UpperTheory.V.Ws.Set2}
\leanref{Development/AutomorphicArithmetic/AutomorphicVarieties/WeakShimuraTheory/FrobeniusFields/ExistsResQ.lean}{UpperTheory.V.Ws.Set2.H}

\begin{lemma}\label{sha-siegel-step-towers-quadratic}
Let $S'$ be a finite set of primes and $s\in\{1,-1\}$. There are a prime $p\notin S'$ and a quadratic character at $p$ of sign $s$ split over $S'$ (Definition~\ref{sha-siegel-step-towers-input}).
\end{lemma}
\leanref{Development/AutomorphicArithmetic/AutomorphicVarieties/WeakShimuraTheory/FrobeniusFields/QuadEx.lean}{UpperTheory.V.Ws.quadEx}
\begin{proof}
For an integer $N\ge1$ let $c_N:G\to(\mathbf Z/N)^\times$ be the character by which $G$ acts on the $N$-th roots of unity. It has open kernel, it is trivial on $I_w(\mathbf Q)$ for every place $w$ above no prime dividing $N$, it maps $I_w(\mathbf Q)$ onto $(\mathbf Z/N)^\times$ for every place $w$ above a prime $N$, and $c_N(\sigma)=-1$ for every complex conjugation $\sigma$. For a prime $q\nmid N$ and a place $w$ above $q$, $D_w(\mathbf Q)$ is contained in the closed subgroup generated by $I_w(\mathbf Q)$ and an element $\gamma\in D_w(\mathbf Q)$ with $c_N(\gamma)=q$.
\begin{enumerate}
\item[(a)] \emph{The prime.} By Dirichlet's theorem there is a prime $p$ larger than every element of $S'$ with $p\equiv s$ modulo $8\prod_{q\in S'}q$. Then $p\notin S'$ and $p$ is odd. Let $(\frac{\cdot}{p})$ be the Legendre symbol. As $p\equiv1\pmod4$ if and only if $s=1$, $(\frac{-1}{p})=s$. Let $q\in S'$. If $q=2$, then $(\frac2p)=1$, as $p\equiv\pm1\pmod8$. If $q$ is odd, then $(\frac pq)=(\frac sq)$, and quadratic reciprocity gives $(\frac qp)=(\frac pq)$ if $s=1$, as then $p\equiv1\pmod4$, and $(\frac qp)=(\frac{-1}q)(\frac{-1}q)=1$ if $s=-1$, as then $p\equiv3\pmod4$ and $(\frac pq)=(\frac{-1}q)$. So $(\frac qp)=1$ for every $q\in S'$.
\item[(b)] \emph{The character.} Let $\eta(g)=(\frac{c_p(g)}p)$. Its kernel contains the open kernel of $c_p$, so it is open. Some $u\in(\mathbf Z/p)^\times$ is not a square, and for every place $w$ above $p$ some $g\in I_w(\mathbf Q)$ has $c_p(g)=u$, so $\eta(g)=-1$; in particular $\eta$ is onto $\{\pm1\}$ and its kernel has index $2$. It is trivial on $I_w(\mathbf Q)$ for $w$ not above $p$, as $c_p$ is. For a complex conjugation $\sigma$, $\eta(\sigma)=(\frac{-1}p)=s$ by (a).
\item[(c)] \emph{Splitting.} Let $q\in S'$ and $w$ above $q$. Then $q\ne p$, and with $\gamma$ as in the first paragraph, $\eta(\gamma)=(\frac qp)=1$ by (a), and $\eta$ is trivial on $I_w(\mathbf Q)$. As $\ker\eta$ is closed, it contains $D_w(\mathbf Q)$.
\end{enumerate}
\end{proof}
% quadEx (QuadEx.lean, 118 lines): (a) is the private quadEx_prime (Nat.forall_exists_prime_gt_and_zmodEq, Mathlib's Dirichlet theorem) with quadEx_jac (Mathlib's quadratic reciprocity for the Jacobi symbol); (b) and (c) use Chr.clsFacts, HV.compIne and Chr.decCyc, the Lean's facts on c_N, which the first paragraph states.

\begin{lemma}\label{sha-siegel-step-towers-kappa-dirichlet}
Let $N\ge1$, let $c_N:G\to(\mathbf Z/N)^\times$ be the character by which $G$ acts on the $N$-th roots of unity, and for a prime $r$ write $N=r^{e_r}N_r$ with $r\nmid N_r$. Let $\chi$ be a Dirichlet character modulo $N$ with values in a domain, viewed as the character $\chi\circ c_N$ of $G$, and for a prime $r$ let $\chi_r$ be its \emph{component at} $r$, $\chi_r(u)=\chi(u_r)$, where $u_r\equiv u$ modulo $r^{e_r}$ and $u_r\equiv1$ modulo $N_r$.
\begin{enumerate}
\item The map $\chi\mapsto\chi\circ c_N$ is an injective homomorphism, and $\chi\circ c_N$ has open kernel. The map $\chi\mapsto\chi_r$ is a homomorphism, $(\chi_r)_r=\chi_r$, $(\chi_{r'})_r=1$ for primes $r'\ne r$, and $\chi$ is the product of its components at the primes dividing $N$.
\item Let $w$ be a place of $\bar{\mathbf Q}$ above a prime $r$. Then $c_N$ maps $I_w(\mathbf Q)$ onto the subgroup of the $u\equiv1$ modulo $N_r$. So $\chi$ and $\chi_r$ agree on $I_w(\mathbf Q)$, and $\chi$ is trivial on $I_w(\mathbf Q)$ if and only if $\chi_r=1$. If $r\nmid N$ and $\sigma$ is a Frobenius element at $w$, then $\chi(\sigma)=\chi(r)$.
\item If $g\in G$ has $\chi_{\mathrm{cyc}}(g)=1$, then $\chi_\ell(g)=1$.
\item If the conductor of $\chi$ is $2^c$, then $\chi^{2^c}=1$.
\end{enumerate}
\end{lemma}
\leanref{Development/AutomorphicArithmetic/AutomorphicVarieties/CharacterLevelBounds/ClsFacts.lean}{UpperTheory.V.Chr.clsFacts}
\leanref{Development/AutomorphicArithmetic/AutomorphicVarieties/HeightAndVanishing/CompAlg.lean}{UpperTheory.V.HV.compAlg}
\leanref{Development/AutomorphicArithmetic/AutomorphicVarieties/HeightAndVanishing/ProjFacts.lean}{UpperTheory.V.HV.projFacts}
\leanref{Development/AutomorphicArithmetic/AutomorphicVarieties/HeightAndVanishing/CompIne.lean}{UpperTheory.V.HV.compIne}
\begin{proof}
(1) The map $c_N$ is onto, by the irreducibility of the cyclotomic polynomials, and its kernel, which fixes $\mu_N$, is open. By the Chinese remainder theorem, $u\mapsto u_r$ is the projection of $(\mathbf Z/N)^\times$ onto its factor $(\mathbf Z/r^{e_r})^\times$; it is an idempotent endomorphism, $(u_{r'})_r=1$ for $r'\ne r$, and $u$ is the product of the $u_r$. The claims follow.

(2) The image of $I_w(\mathbf Q)$ is the inertia group at $r$ of $\mathbf Q(\mu_N)$, which is $(\mathbf Z/r^{e_r})^\times$, the subgroup of the $u\equiv1$ modulo $N_r$ \cite[Chapter~2]{Was97}. On this subgroup $u=u_r$, which gives the second claim, and $u\mapsto u_r$ maps $(\mathbf Z/N)^\times$ onto it, which gives the third. If $r\nmid N$, then $N_r=N$, and $\sigma(\zeta)=\zeta^r$ for $\zeta\in\mu_N$, so $c_N(\sigma)=r$.

(3) Such $g$ fixes the $\ell$-power roots of unity, so $c_N(g)\equiv1$ modulo $\ell^{e_\ell}$, and $c_N(g)_\ell=1$.

(4) The character $\chi$ factors through $(\mathbf Z/2^c)^\times$, whose order divides $2^c$.
\end{proof}
% Chr.dirG is chi composed with Chr.cls N (c_N); Chr.clsFacts (open kernel, compatibility with N | M, the values on roots of unity); the component is the Lean's comp r chi = chi o proj r N, with compAlg and projFacts for (1); (2) is compIne (its first part, the image of inertia, is the private compIne_map through B1.cycIne, the cited inertia of Q(mu_N)) and the private dirG_frob of ExistsDecompositionElementAtLXiPowFourNeOne.lean; (3) is the private kx_cls and kx_kapL of KxKapL.lean; (4) is the private pow_cond of ExistsInertiaElementAtTwoXiPowFourNeOne.lean. psiBar gives the injectivity in (1) (dirHom_inj). About 400 lines with CompIne, CompAlg, ProjFacts and ClsFacts.

\begin{lemma}\label{sha-siegel-step-towers-kappa-units}
Let $D$ be a Siegel datum, with the notation of Definition~\ref{sha-siegel-datum}.
\begin{enumerate}
\item Let $m\ge1$ be prime to $\ell$, and $x,y\in\mathcal O^\times$ with $x^m=y^m=1$ and $\bar x=\bar y$. Then $x=y$.
\item A character $a:G\to\mathcal O^\times$ with $a^{\ell^s}=1$ for some $s$ has $\bar a=1$.
\item The kernel of $\omega$ is the kernel of $c_\ell$, in the notation of Lemma~\ref{sha-siegel-step-towers-kappa-dirichlet}; $\omega$ is trivial on $I_w(\mathbf Q)$ for every place $w$ not above $\ell$; and every character $\theta:G\to\mathcal O^\times$ that is trivial on $\ker\omega$ is $\omega^b$ for some $0\le b<\ell-1$.
\end{enumerate}
\end{lemma}
\leanref{Development/AutomorphicArithmetic/AutomorphicVarieties/HeightAndVanishing/CongrFacts.lean}{UpperTheory.V.HV.congrFacts}
\leanref{Development/AutomorphicArithmetic/AutomorphicVarieties/WeakShimuraTheory/DirichletGroupPowers/RedOne.lean}{UpperTheory.V.Ws.large₂.redOne}
\leanref{Development/AutomorphicArithmetic/AutomorphicVarieties/HeightAndVanishing/ExistsStrictHenselianAbsolutelyClosedLocalRing.lean}{UpperTheory.V.HV.omT}
\begin{proof}
(1) Let $z=xy^{-1}$, so that $z^m=1$ and $\bar z=1$. The element $\sum_{i<m}z^i$ has residue $m\ne0$, so it is a unit, and $(z-1)\sum_{i<m}z^i=z^m-1=0$; so $z=1$.

(2) The residue field has characteristic $\ell$, so $(\bar a(g)-1)^{\ell^s}=\bar a(g)^{\ell^s}-1=0$.

(3) As $\omega^{\ell-1}=1$ and $\bar\omega=\bar\chi_{\mathrm{cyc}}$, (1) shows that $\omega(g)=1$ if and only if $\chi_{\mathrm{cyc}}(g)\equiv1$ modulo $\ell$, that is $c_\ell(g)=1$. For $w$ not above $\ell$, $c_\ell$ is trivial on $I_w(\mathbf Q)$ by Lemma~\ref{sha-siegel-step-towers-kappa-dirichlet}(2). The group $G/\ker\omega\cong(\mathbf Z/\ell)^\times$ is cyclic; let $\sigma_0$ map to a generator. Then $\omega(\sigma_0)$ has order $\ell-1$, and $\theta(\sigma_0)^{\ell-1}=\theta(\sigma_0^{\ell-1})=1$. The $(\ell-1)$-st roots of unity of the domain $\mathcal O$ are the powers of $\omega(\sigma_0)$, so $\theta(\sigma_0)=\omega(\sigma_0)^b$ with $0\le b<\ell-1$, and $\theta=\omega^b$, as both factor through $G/\ker\omega$.
\end{proof}
% (1) is redOne (RedOne.lean) and the private eq_of_res_eq of CongrFacts.lean (congrFacts, fourth part); (2) is the private res_eq_one (congrFacts, third part); (3) is omT (its parts ker w = ker (Chr.cls l), the inertia away from l, and the last part, through gen_pow and clsOnto).

\begin{lemma}\label{sha-siegel-step-towers-kappa-dyadic}
Let $n\ge0$, let $\xi$ be a primitive Dirichlet character modulo $2^n$ with values in a domain, and let $v>1$ and $u$ be odd integers and $a\ge0$ with $\xi(v)^{2^au}=1$. Then $2^n<v^22^a$.
\end{lemma}
\leanref{Development/AutomorphicArithmetic/AutomorphicVarieties/SymplecticGSp4/TwoPowLtSqMulOfPrimitiveCharPowEqOne.lean}{UpperTheory.two_pow_lt_sq_mul_of_primitive_char_pow_eq_one}
\begin{proof}
If $n<2$, then $2^n\le2<9\le v^2$. Let $n=m+2$. As $v$ is odd, $v^2-1$ is divisible by $4$, so $v^2=1+2^{k+2}d$ with $k\ge0$ and $d$ odd. By induction, $(v^2)^u=1+2^{k+2}d'$ with $d'$ odd, as $u$ is odd, and $(v^2)^{u2^a}=1+2^{k+2+a}d''$ with $d''$ odd. Suppose that $m+2>k+2+a$. Raising to the power $2^{m+1-(k+2+a)}$ gives $(v^2)^{u2^a2^{m+1-(k+2+a)}}\equiv1+2^{m+1}$ modulo $2^{m+2}$, and $\xi$ of the left side is $1$, as $\xi(v)^{2^au}=1$. But $\xi(1+2^{m+1})\ne1$: otherwise $\xi$ would be trivial on the kernel of $(\mathbf Z/2^{m+2})^\times\to(\mathbf Z/2^{m+1})^\times$, which is $\{1,1+2^{m+1}\}$, and $\xi$ would not be primitive. So $m+2\le k+2+a$, and $2^n\le2^{k+2}2^a<v^22^a$.
\end{proof}
% two_pow_lt_sq_mul_of_primitive_char_pow_eq_one (120 lines): the private primitive_dyadic_character_ne_one, two_adic_unit_pow_odd and two_adic_unit_pow_two_pow.

\begin{lemma}\label{sha-siegel-step-towers-kappa-exponent}
Let $D$ be a Siegel datum of case II, attached to $(h,k,\psi)$ for a host datum $h$ for $\psi_0$ (Definitions~\ref{sha-siegel-tower-towers} and~\ref{sha-siegel-datum}), and let $2^{c_2}$ be the conductor of the component $\psi_2$ of $\psi$ at $2$. Then there are $a\ge0$ and an odd $u\ge1$ with $\psi^{2^au}=1$ and $\ell^22^{a+2}\le2^{n_{\mathrm{II}}}$.
\end{lemma}
\leanref{Development/AutomorphicArithmetic/AutomorphicVarieties/HeightAndVanishing/HostComp.lean}{UpperTheory.V.HV.hostComp}
\leanref{Development/AutomorphicArithmetic/AutomorphicVarieties/HeightAndVanishing/OrdProjNDvd.lean}{UpperTheory.V.HV.hostTh1}
\leanref{Development/AutomorphicArithmetic/AutomorphicVarieties/LevelCohomology/Th1Facts.lean}{UpperTheory.V.X.th1Facts}
\begin{proof}
Write $o=o_{\psi_0}$ and $N=N_h$, and use the components of Lemma~\ref{sha-siegel-step-towers-kappa-dirichlet}.
\begin{enumerate}
\item[(a)] \emph{Odd components.} Let $r$ be an odd prime and $u\in(\mathbf Z/N)^\times$. If $r\nmid N_{\psi_0}$, then $(u_r)^2$ is $1$ modulo $2^{e+2}$ and $N_{\psi_0}$ and a square modulo $q$ and $d$, so it lies in $H_h$ (Definition~\ref{sha-host-datum}), and $\psi_r^2=1$. If $r\mid N_{\psi_0}$, then $r\nmid qd$, and $(u_r)^o$ is $1$ modulo $2^{e+2}$, $q$ and $d$, and $\psi_0((u_r)^o)=\psi_0(u_r)^o=1$; so $(u_r)^o\in H_h$ and $\psi_r^o=1$. In both cases $\psi_r^{2o}=1$, as $\psi$ is trivial on $H_h$.
\item[(b)] \emph{Parts.} Let $\kappa_0$ be the part of $\psi$ of $\ell$-power order, $\ell^s$ say, and $\theta_0=\psi\kappa_0^{-1}$ the part of order prime to $\ell$; both are powers of $\psi$. The component $(\kappa_0)_2$ has order dividing $\ell^s$ and a power of $2$, so it is $1$, and $(\theta_0)_2=\psi_2$. Let $\theta_1=\theta_0\psi_2^{-1}$. Then $\psi=\psi_2\theta_1\kappa_0$, $(\theta_1)_2=1$, and for odd $r$, $(\theta_1)_r=(\theta_0)_r$ is a power of $\psi_r$, so $(\theta_1)_r^{2o}=1$ by (a). As $\theta_1$ is the product of its components, $\theta_1^{2o}=1$.
\item[(c)] \emph{The exponent.} By Lemma~\ref{sha-siegel-step-towers-kappa-dirichlet}(4), $\psi_2^{2^{c_2}}=1$. So $\psi^{2^{c_2}\cdot2o\cdot\ell^s}=1$. Write $o=2^{v}o'$ with $o'$ odd, and put $a=c_2+1+v$ and $u=o'\ell^s$, which is odd as $\ell$ is.
\item[(d)] \emph{The bound.} As $2^v\le o$, $\ell^22^{a+2}=2^{c_2+3}2^v\ell^2\le2^{c_2+6}o\,\ell^3N_{\psi_0}^2\le2^{n_{\mathrm{II}}}$, by the definition of $n_{\mathrm{II}}$ (Definition~\ref{sha-siegel-datum}).
\end{enumerate}
\end{proof}
% The private psi_exp of ExistsDecompositionElementAtLXiPowFourNeOne.lean: (a) is hostComp (HostComp.lean, 81 lines, with hostIndex for psi^{2^{e+1} o} = 1); (b) is th1Facts (Th1Facts.lean, 81 lines: psi = psi2 * th1 * kap) and hostTh1 (OrdProjNDvd.lean: h1 for comp 2 th0 = psi2, t3 for the order of th1 dividing 2 o); (c) is pow_cond with compCond (CompCond.lean, the conductor of psi2 is 2^{c2}); (d) is the final calc with Nat.le_pow_clog. About 400 lines.

\begin{lemma}\label{sha-siegel-step-towers-kappa}
Let $D$ be a Siegel datum of case II (Definition~\ref{sha-siegel-tower-towers}), with the notation of Definition~\ref{sha-siegel-datum}.
\begin{enumerate}
\item There is an auxiliary character for $D$ (Definition~\ref{sha-siegel-step-towers-input}).
\item The character $\xi$ has order a power of $2$ and is trivial on $I_w(\mathbf Q)$ for every place $w$ not above $2$, and $\Psi(\sigma)=1$ for every complex conjugation $\sigma\in G$.
\item Let $\kappa:G\to\mathcal O^\times$ be a character with $\kappa^{\ell^j}=1$ for some $j\ge0$. Then some $t\in G$ has $\omega(t)=1$, $\kappa(t)=1$ and $\bar\Psi(t)=-1$; for every place $w$ above $\ell$ some $g\in D_w(\mathbf Q)$ has $\omega(g)=1$, $\bar\Psi(g)=1$, $\kappa(g)=1$ and $\bar\xi(g)^4\ne1$; and for every place $w$ above $2$ some $g\in I_w(\mathbf Q)$ has $\bar\Psi(g)=1$, $\kappa(g)=1$ and $\bar\xi(g)^4\ne1$.
\end{enumerate}
\end{lemma}
\leanref{Development/AutomorphicArithmetic/AutomorphicVarieties/WeakShimuraTheory/FrobeniusFields/KxKapL.lean}{UpperTheory.V.Ws.kapEx}
\leanref{Development/AutomorphicArithmetic/AutomorphicVarieties/WeakShimuraTheory/FrobeniusFields/PsiCc.lean}{UpperTheory.V.Ws.xi2}
\leanref{Development/AutomorphicArithmetic/AutomorphicVarieties/WeakShimuraTheory/FrobeniusFields/NegVal.lean}{UpperTheory.V.Ws.ev2}
\leanref{Development/AutomorphicArithmetic/AutomorphicVarieties/ModularityArgument/UpperTriangularPlaces/BigOm.lean}{UpperTheory.V.Keep.bigOm}
\leanref{Development/AutomorphicArithmetic/AutomorphicVarieties/WeakShimuraTheory/FrobeniusFields/ExistsDecompositionElementAtLXiPowFourNeOne.lean}{UpperTheory.V.Ws.exists_decomposition_element_at_l_xi_pow_four_ne_one}
\leanref{Development/AutomorphicArithmetic/AutomorphicVarieties/WeakShimuraTheory/FrobeniusFields/ExistsInertiaElementAtTwoXiPowFourNeOne.lean}{UpperTheory.V.Ws.exists_inertia_element_at_two_xi_pow_four_ne_one}
\begin{proof}
(2) As $D$ is of case II, $\xi$ is the Galois character of a primitive Dirichlet character of conductor $2^{n}$, $n=n_{\mathrm{II}}$ (Definition~\ref{sha-siegel-datum}). It factors through $c_{2^n}$, in the notation of the proof of Lemma~\ref{sha-siegel-step-towers-quadratic}, so its order divides $|(\mathbf Z/2^n)^\times|$, a power of $2$, and it is trivial on $I_w(\mathbf Q)$ for $w$ not above $2$. Let $\sigma$ be a complex conjugation. Then $\chi_{\mathrm{cyc}}(\sigma)=-1$ and $\psi(\sigma)=\psi(-1)=(-1)^k$, as $\psi\in X_h(k)$ (Definition~\ref{sha-host-datum}); so $\Psi(\sigma)=(-1)^k(-1)^k=1$.

For (1) and (3), let $\psi$ also denote the Dirichlet character modulo $N=N_h$ of which it is the Galois character (Definition~\ref{sha-siegel-datum}), and use the components of Lemma~\ref{sha-siegel-step-towers-kappa-dirichlet}. Let $\kappa_0$ be the part of $\psi$ of $\ell$-power order, $\ell^s$ say, and $\theta_0=\psi\kappa_0^{-1}$ the part of order $m$ prime to $\ell$. Recall that $\Psi=\chi_{\mathrm{cyc}}^k\psi$, so $\bar\Psi=\bar\omega^k\bar\psi$, as $\bar\omega=\bar\chi_{\mathrm{cyc}}$.

(1) Let $\kappa_1=\kappa_0\,(\kappa_0)_\ell^{-1}$, and $\kappa$ its Galois character.
\begin{enumerate}
\item[(a)] \emph{Open kernel and order.} The kernel of $\kappa$ contains that of $c_N$, which is open. As $\kappa_0^{\ell^s}=1$, also $(\kappa_0)_\ell^{\ell^s}=1$, and $\kappa^{\ell^s}=1$.
\item[(b)] \emph{A useful identity.} Let $g\in G$ with $\chi_{\mathrm{cyc}}(g)=1$. Then $(\kappa_0)_\ell(g)=1$ by Lemma~\ref{sha-siegel-step-towers-kappa-dirichlet}(3), so $\kappa(g)=\kappa_0(g)$, and $\psi(g)=\theta_0(g)\kappa(g)$.
\item[(c)] \emph{Away from $S$.} Let $w$ be above no prime of $S$ and $g\in I_w(\mathbf Q)$. As $\ell\in S$, $\chi_{\mathrm{cyc}}(g)=1$; and $\psi(g)=1$, as $\psi$ is unramified outside $S$. As $\kappa_0$ is a power of $\psi$, $\kappa(g)=\kappa_0(g)=1$ by (b).
\item[(d)] \emph{Above $\ell$.} The component of $\kappa_1$ at $\ell$ is $(\kappa_0)_\ell\,((\kappa_0)_\ell)_\ell^{-1}=1$, so $\kappa$ is trivial on $I_w(\mathbf Q)$ for $w$ above $\ell$, by Lemma~\ref{sha-siegel-step-towers-kappa-dirichlet}(2).
\item[(e)] \emph{The last two conditions.} Let $g\in G$ with $\chi_{\mathrm{cyc}}(g)=1$ and $\kappa(g)=1$. By (b), $\psi(g)=\theta_0(g)$, so $\psi(g)^m=1$, with $m$ prime to $\ell$; this is the last condition. Let now $w$ be a place not above $\ell$ and $g\in I_w(\mathbf Q)$ with $\bar\Psi(g)=1$ and $\kappa(g)=1$. Then $\chi_{\mathrm{cyc}}(g)=1$, so $\psi(g)^m=1$ and $\bar\psi(g)=\bar\Psi(g)=1$; so $\psi(g)=1$ by Lemma~\ref{sha-siegel-step-towers-kappa-units}(1).
\end{enumerate}
So $\kappa$ is an auxiliary character for $D$.

(3) Let $\kappa$ and $j$ be as in (3); $\ell^j$ is odd. Let $n=n_{\mathrm{II}}$, so that $\xi$ is the Galois character of a primitive Dirichlet character $\xi'$ modulo $2^n$.
\begin{enumerate}
\item[(a)] \emph{The residue of $\Psi$.} Let $\chi'=\theta_0\,(\theta_0)_\ell^{-1}$, the character $\chi'_\ell(\psi)$ of Definition~\ref{sha-ordinary-deformation}. The component $(\theta_0)_\ell$ has order prime to $\ell$, so it is trivial on the kernel of $(\mathbf Z/\ell^{e_\ell})^\times\to(\mathbf Z/\ell)^\times$, which has $\ell$-power order: it factors through $(\mathbf Z/\ell)^\times$, and as a character of $G$ through $c_\ell$. As $\ker\omega=\ker c_\ell$, Lemma~\ref{sha-siegel-step-towers-kappa-units}(3) gives $(\theta_0)_\ell=\omega^b$ for some $b$. As $\psi=(\theta_0)_\ell\,\chi'\,\kappa_0$ and $\bar\kappa_0=1$ (Lemma~\ref{sha-siegel-step-towers-kappa-units}(2)), $\bar\Psi=\bar\omega^{k+b}\bar\chi'$.
\item[(b)] \emph{The element $t$.} As $D$ is of case II, $\chi'$ has even order $2k'$, and so has its Galois character, by Lemma~\ref{sha-siegel-step-towers-kappa-dirichlet}(1). So some $\sigma$ has $\chi'(\sigma)^{k'}\ne1$; as $(\chi'(\sigma)^{k'})^2=1$ in the domain $\mathcal O$, $\sigma_1=\sigma^{k'}$ has $\chi'(\sigma_1)=-1$. The component of $\chi'$ at $\ell$ is $1$, so $\chi'$ is trivial on $I_\ell$; and $c_\ell$ maps $I_\ell$ onto $(\mathbf Z/\ell)^\times$ (Lemma~\ref{sha-siegel-step-towers-kappa-dirichlet}(2)), so some $\tau\in I_\ell$ has $c_\ell(\tau)=c_\ell(\sigma_1)$. Let $t=(\sigma_1\tau^{-1})^{\ell^j}$. Then $c_\ell(t)=1$, so $\omega(t)=1$; $\kappa(t)=1$; and $\chi'(t)=(-1)^{\ell^j}=-1$. By (a), $\bar\Psi(t)=\bar\chi'(t)=-1$.
\item[(c)] \emph{The elements at $\ell$.} Let $\sigma_\ell$ be the Frobenius element at $v_\ell$ of $D$, with $\chi_{\mathrm{cyc}}(\sigma_\ell)=1$. Then $\bar\omega(\sigma_\ell)=1$, and as $\omega(\sigma_\ell)^{\ell-1}=1$, $\omega(\sigma_\ell)=1$ by Lemma~\ref{sha-siegel-step-towers-kappa-units}(1). Let $a$ and $u$ be as in Lemma~\ref{sha-siegel-step-towers-kappa-exponent}, and $g_0=\sigma_\ell^{\ell^j2^au}\in D_\ell$. Then $\omega(g_0)=1$, $\kappa(g_0)=1$ and $\psi(g_0)=1$, so $\bar\Psi(g_0)=1$. Suppose that $\bar\xi(g_0)^4=1$. As $\xi$ has order prime to $\ell$, $\xi(g_0)^4=1$ by Lemma~\ref{sha-siegel-step-towers-kappa-units}(1). As $\ell\nmid2^n$, $\xi(\sigma_\ell)=\xi'(\ell)$ by Lemma~\ref{sha-siegel-step-towers-kappa-dirichlet}(2), so $\xi'(\ell)^{2^{a+2}\ell^ju}=1$, and Lemma~\ref{sha-siegel-step-towers-kappa-dyadic} gives $2^n<\ell^22^{a+2}\le2^n$, a contradiction. So $\bar\xi(g_0)^4\ne1$. For $w$ above $\ell$, choose $h\in G$ with $hv_\ell=w$; then $g=hg_0h^{-1}\in D_w(\mathbf Q)$, and the characters $\omega$, $\bar\Psi$, $\kappa$ and $\bar\xi$ take the same values at $g$ and $g_0$.
\item[(d)] \emph{The elements at $2$.} Let $w$ be above $2$. By Lemma~\ref{sha-siegel-step-towers-kappa-dirichlet}(2), $c_{2^n}$ maps $I_w(\mathbf Q)$ onto $(\mathbf Z/2^n)^\times$, so some $u\in I_w(\mathbf Q)$ has $c_{2^n}(u)=3$, and $\xi(u)=\xi'(3)$. Let $2^{c_2}$ be the conductor of $\psi_2$. On $I_w(\mathbf Q)$, $\psi$ agrees with $\psi_2$, and $\psi_2^{2^{c_2}}=1$, by (2) and (4) of Lemma~\ref{sha-siegel-step-towers-kappa-dirichlet}; so $\psi(u)^{2^{c_2}}=1$. Let $g=u^{2^{c_2}\ell^j}\in I_w(\mathbf Q)$. Then $\omega(g)=1$, as $w$ is not above $\ell$ (Lemma~\ref{sha-siegel-step-towers-kappa-units}(3)), and $\psi(g)=1$, so $\bar\Psi(g)=1$; and $\kappa(g)=1$. Suppose that $\bar\xi(g)^4=1$. As $\xi(u)^{2^n}=1$ and $2^n$ is prime to $\ell$, $\xi(g)^4=1$ by Lemma~\ref{sha-siegel-step-towers-kappa-units}(1), that is $\xi'(3)^{2^{c_2+2}\ell^j}=1$. Lemma~\ref{sha-siegel-step-towers-kappa-dyadic} gives $2^n<9\cdot2^{c_2+2}<2^{c_2+6}$; but $2^{c_2+6}\le2^n$ by the definition of $n_{\mathrm{II}}$, a contradiction. So $\bar\xi(g)^4\ne1$.
\end{enumerate}
% (1) is kapEx (KxKapL.lean, 92 lines): kappa_0 is kap l psi (lPart), theta_0 is th0 l psi (lPrime), kappa_1 is kapB l psi = kap * (kapL)^{-1} with kapL = comp l kap; (a) is the opn and pow fields with kx_ord, (b) is kx_kapL, (c) is kx_unr with D.psi_unr, (d) is compIne with compAlg (comp l (kapB) = 1), (e) is kx_pow and congrFacts for kb and kc. (3)(a) is psiBar (PsiBar.lean, the private ker_le for (theta_0)_l = omega^b through compOrd, and partOrd); (3)(b) is ev2 (NegVal.lean, 59 lines, the private neg_val, compIne and projFacts for tau, omT for omega); (3)(c) is Keep.bigOm (BigOm.lean, 168 lines, with psi_exp, dirG_frob, redOne and Chr.placesQ; the Lean's Set2.bigL uses exists_decomposition_element_at_l_xi_pow_four_ne_one, 121 lines, which is the same argument without omega(g) = 1); (3)(d) is exists_inertia_element_at_two_xi_pow_four_ne_one (109 lines, with the private all_mem_range, tot_dvd and pow_cond). The four lemmas before this one carry the library of Dirichlet characters as Galois characters.
\end{proof}

\begin{lemma}\label{sha-siegel-step-towers-subgroup}
Let $D$ be a Siegel datum and $X=(\kappa,p_3,p,p',\varphi,\eta,\eta')$ a tower input for $D$, with subgroup $H_X$ (Definition~\ref{sha-siegel-step-towers-input}).
\begin{enumerate}
\item For $g\in G$ and $s_1,s_2,s_3\in\{\pm1\}$ there is $g'\in G$ with $\varphi(g')=s_1$, $\eta(g')=s_2$, $\eta'(g')=s_3$, $\Psi(g')=\Psi(g)$, $\xi(g')=\xi(g)$, $\omega(g')=\omega(g)$ and $\kappa(g')=\kappa(g)$.
\item The subgroup $H_X$ is open, contains every commutator and every complex conjugation, and has even index; and for every place $w$ above $\ell$ and $g\in I_w(\mathbf Q)$ some power $g^m$ with $m$ prime to $\ell$ lies in $H_X$.
\end{enumerate}
\end{lemma}
\leanref{Development/AutomorphicArithmetic/AutomorphicVarieties/WeakShimuraTheory/FrobeniusFields/Adj2.lean}{UpperTheory.V.Ws.adj2}
\leanref{Development/AutomorphicArithmetic/AutomorphicVarieties/WeakShimuraTheory/FrobeniusFields/CharacterSubgroupOpenEvenIndexTamePower.lean}{UpperTheory.V.Ws.characterSubgroup_open_even_index_tame_power}
\begin{proof}
First, the primes $p_3$, $p$ and $p'$ are distinct and not in $S$, and a place of $\bar{\mathbf Q}$ lies above at most one prime. Let $q$ be a prime not in $S$ and $w$ a place above $q$. Then $\Psi$, $\xi$, $\omega$ and $\kappa$ are trivial on $I_w(\mathbf Q)$: $\psi$ and $\xi$ are unramified outside $S$, $\chi_{\mathrm{cyc}}$ is unramified at $q\ne\ell$, $\kappa$ is trivial there as $w$ is above no prime of $S$, and for $g\in I_w(\mathbf Q)$, $\omega(g)$ is an $(\ell-1)$-st root of unity with residue $\bar\chi_{\mathrm{cyc}}(g)=1$, hence $\omega(g)=1$. Also, the residue of $\Psi$ is $\bar\Psi=\bar\omega^k\bar\psi$, as $\bar\omega=\bar\chi_{\mathrm{cyc}}$.

(1) Let $h_3\in I_{v_{p_3}}(\mathbf Q)$ with $\varphi(h_3)=-1$, $h\in I_{v_p}(\mathbf Q)$ with $\eta(h)=-1$ and $h'\in I_{v_{p'}}(\mathbf Q)$ with $\eta'(h')=-1$. By the first paragraph, $\eta$ and $\eta'$ are trivial at $h_3$, $\varphi$ and $\eta'$ at $h$, and $\varphi$ and $\eta$ at $h'$, as each of these characters is trivial on inertia at places not above its prime; and $\Psi$, $\xi$, $\omega$ and $\kappa$ are trivial at $h_3$, $h$ and $h'$. Multiplying $g$ on the right by $h_3$ if $\varphi(g)\ne s_1$, then by $h$ if $\eta$ of the product is not $s_2$, then by $h'$ if $\eta'$ of the product is not $s_3$, gives $g'$.

(2) Let $m+1$ be the number of elements of $\mathbf F$, which is a power of $\ell$; so $m$ is prime to $\ell$ and $\bar\Psi(g)^m=1$ for all $g$.
\begin{enumerate}
\item[(a)] \emph{Open.} The set $U$ of $g$ with $\varphi(g)=\eta(g)=\eta'(g)=\kappa(g)=1$ and $\bar\Psi(g)=1$ is contained in $H_X$. It is open, as the kernels of $\varphi$, $\eta$, $\eta'$ and $\kappa$ are, and the set of $g$ with $\bar\Psi(g)=1$ contains $\ker\omega\cap\ker\psi$, which is open. So the subgroup $H_X$, which contains the neighbourhood $U$ of $1$, is open.
\item[(b)] \emph{Commutators.} The characters $\varphi$, $\eta$, $\eta'$, $\kappa$ and $\bar\Psi$ take values in commutative groups, so they are trivial on commutators, which therefore lie in $U\subseteq H_X$.
\item[(c)] \emph{Complex conjugations.} Let $\sigma$ be a complex conjugation. Then $\sigma^2=1$ and $\kappa(\sigma)^{\ell^j}=1$ with $\ell$ odd, so $\kappa(\sigma)=1$. Also $\Psi(\sigma)=1=\varphi(\sigma)$ and $\eta(\sigma)=-1=\eta'(\sigma)$. So $\sigma\in H_X$.
\item[(d)] \emph{Even index.} The character $\eta\eta'$ is trivial on $H_X$, and it is onto $\{\pm1\}$ by (1). So its kernel has index $2$ and contains $H_X$, and $2$ divides the index of $H_X$.
\item[(e)] \emph{Tame powers.} Let $w$ be above $\ell$ and $g\in I_w(\mathbf Q)$. As $\ell\in S$, $w$ is not above $p_3$, $p$ or $p'$, so $\varphi$, $\eta$ and $\eta'$ are trivial at $g$, and $\kappa(g)=1$ as $\kappa$ is trivial on $I_w(\mathbf Q)$. With $m$ as above, $\bar\Psi(g^m)=1=\varphi(g^m)$, so $g^m\in H_X$.
\end{enumerate}
\end{proof}

\begin{lemma}\label{sha-siegel-step-towers-elements}
Let $D$ be a Siegel datum, with its complex conjugation $c$, and $X$ a tower input for $D$, with subgroup $H_X$ and characters $\kappa,\varphi,\eta,\eta'$ (Definition~\ref{sha-siegel-step-towers-input}).
There are $a,u,t\in H_X\cap\ker\eta$ such that $\bar\xi(a)=-1$, $\bar\Psi(a)=1$ and $\bar\omega(a)=-1$; $\bar\xi(u)^4\ne1$, $\bar\Psi(u)=1$ and $\omega(u)=1$; and $\bar\xi(t)^2\ne1$, $\bar\Psi(t)=-1$ and $\omega(t)=1$.
\end{lemma}
\leanref{Development/AutomorphicArithmetic/AutomorphicVarieties/WeakShimuraTheory/FrobeniusFields/Elts2.lean}{UpperTheory.V.Ws.elts2}
\begin{proof}
First, let $g\in G$ and $s\in\{\pm1\}$ with $\kappa(g)=1$ and $\bar\Psi(g)=s$. Lemma~\ref{sha-siegel-step-towers-subgroup}(1) gives $g'$ with $\varphi(g')=s$, $\eta(g')=\eta'(g')=1$, and with the values of $\Psi$, $\xi$, $\omega$ and $\kappa$ at $g$; then $g'\in H_X\cap\ker\eta$. Let $j$ be such that $\kappa^{\ell^j}=1$; $\ell^j$ is odd.
\begin{enumerate}
\item[(a)] Apply the first paragraph to $c^{\ell^j}$, with $s=1$: $\kappa(c^{\ell^j})=1$ and $\Psi(c)=1$. As $\xi(c)=-1$ and $\omega(c)=-1$, the latter because $\omega(c)$ is a root of unity of order prime to $\ell$ with residue $\bar\chi_{\mathrm{cyc}}(c)=-1$, the resulting element $a$ has $\bar\xi(a)=-1$, $\bar\Psi(a)=1$ and $\bar\omega(a)=-1$.
\item[(b)] Let $u_0\in I_{v_2}(\mathbf Q)$ have $\bar\Psi(u_0)=1$, $\kappa(u_0)=1$ and $\bar\xi(u_0)^4\ne1$, as in Definition~\ref{sha-siegel-step-towers-input}(3). Apply the first paragraph with $s=1$; the resulting $u$ has $\omega(u)=\omega(u_0)=1$, as $\omega$ is unramified at $2\ne\ell$.
\item[(c)] Let $t_0$ have $\omega(t_0)=1$, $\kappa(t_0)=1$ and $\bar\Psi(t_0)=-1$, and apply the first paragraph with $s=-1$, which gives $t_1$. If $\bar\xi(t_1)^2\ne1$, let $t=t_1$. Otherwise let $t=t_1u$; then $\bar\xi(t)^2=\bar\xi(u)^2\ne1$, since $\bar\xi(u)^4\ne1$, and $\bar\Psi(t)=-1$, $\omega(t)=1$, $t\in H_X\cap\ker\eta$.
\end{enumerate}
\end{proof}

\begin{lemma}\label{sha-siegel-step-towers-make}
Let $D$ be a Siegel datum, $\epsilon$ a relabelling, $K$ an imaginary quadratic field such that $D_w(\mathbf Q)\subseteq G_K$ for every place $w$ above $\ell$, $F$ a totally real number field, Galois over $\mathbf Q$ with abelian Galois group, and $g_0\in G_F$ (Definition~\ref{sha-siegel-tower-towers}).
Then there is a tower $T$ for $D$ with relabelling $\epsilon$, fields $K$ and $F$, and $g_0^T=g_0$.
\end{lemma}
\leanref{Development/AutomorphicArithmetic/AutomorphicVarieties/WeakShimuraTheory/CompactIntegralFields/MkTower.lean}{UpperTheory.V.Ws.mkTower}
\begin{proof}
Let $L=FK$.
\begin{enumerate}
\item[(a)] \emph{$L$ is CM.} It is totally complex, as it contains $K$, and $[L:F]\le[K:\mathbf Q]=2$. It is not $F$, which is totally real; so it is a totally complex quadratic extension of the totally real field $F$.
\item[(b)] \emph{An element of $G_F$ moving $K$.} If $G_F\subseteq G_K$, then $K\subseteq F$ would be totally real. So some $h\in G_F$ acts nontrivially on $K$.
\item[(c)] \emph{Moving places.} Let $P$ be the set of places of $\bar{\mathbf Q}$ above $\ell$ that induce on $K$ the place induced by $v_\ell$. Let $w$ be a place above $\ell$. As $G$ acts transitively on the places above $\ell$, some $\sigma\in G$ carries $v_\ell$ to $w$: $\sigma x$ is integral at $w$ if and only if $x$ is integral at $v_\ell$. As $\operatorname{Gal}(K/\mathbf Q)$ has order $2$ and $h$ acts nontrivially on $K$, one of $\sigma$ and $h^{-1}\sigma$ acts trivially on $K$; so there is $m\in\{1,h\}\subseteq G_F$ that agrees with $\sigma$ on $K$. The place $m^{-1}w$ lies above $\ell$, and for $x\in K$, $x$ is integral at $m^{-1}w$ if and only if $mx=\sigma x$ is integral at $w$, that is, if and only if $x$ is integral at $v_\ell$. So $m^{-1}w\in P$, and $w$ is conjugate under $G_F$ to an element of $P$.
\item[(d)] \emph{The places.} Conjugacy under $G_F$ is an equivalence relation on $P$. It has finitely many classes: each class contains $\tau v_\ell$ for some $\tau\in G$, and $\tau v_\ell$ and $\tau'v_\ell$ are conjugate under $G_F$ when $\tau$ and $\tau'$ have the same restriction to $F$, as then $\tau'\tau^{-1}\in G_F$ carries $\tau v_\ell$ to $\tau'v_\ell$. Let $\Sigma_T\subseteq P$ contain one element of each class. By (c), every place above $\ell$ is conjugate under $G_F$ to an element of $\Sigma_T$, and to only one, since two elements of $\Sigma_T$ conjugate to the same place are conjugate to each other.
\end{enumerate}
The relabelling $\epsilon$, the fields $K$ and $F$, the element $g_0$ and $\Sigma_T$ form a tower $T$.
\end{proof}
% mkTower (MkTower.lean, 238 lines): (a) is the private mt_cm, (b) mt_out, (c) exists_over_l_place_conj_agreeing_on_field with mt_conj and mem_or_mul_mem_fixing_of_quadratic, (d) the quotient by the setoid s, finite through restrictNormalHom. The Lean's mkTower assumes that D_w(Q) lies in G_K for w above l (hdec); its proof does not use this hypothesis.

\begin{lemma}\label{sha-siegel-step-towers-three-host}
Let $D$ be a Siegel datum of case III, attached to $(h,k,\psi)$ for a host datum $h=(\psi_0,e,q,d)$ that is large for $t$, with $0<k<t$ (Definitions~\ref{sha-siegel-tower-towers}, \ref{sha-siegel-datum} and~\ref{sha-host-datum}), and let $m$ be the order of $\psi$. Then $\ell\ge5$, $\ell-1\mid2k$, $\ell\nmid N_h$, $\ell\ne p_1(\psi_0)$, $m$ is odd and prime to $\ell$, the set $S_b$ is empty, and $\xi$ and $\psi$ are trivial on $I_w(\mathbf Q)$ for every place $w$ above $\ell$.
\end{lemma}
\leanref{Development/AutomorphicArithmetic/AutomorphicVarieties/WeakShimuraTheory/DirichletGroupPowers/DirGPowOne.lean}{UpperTheory.V.Ws.large₂.host3}
\leanref{Development/AutomorphicArithmetic/AutomorphicVarieties/HeightAndVanishing/BchiFacts.lean}{UpperTheory.V.HV.bchiFacts}
\leanref{Development/AutomorphicArithmetic/AutomorphicVarieties/HeightAndVanishing/HostOdd.lean}{UpperTheory.V.HV.hostOdd}
\leanref{Development/AutomorphicArithmetic/AutomorphicVarieties/HeightAndVanishing/HostIndex.lean}{UpperTheory.V.HV.hostIndex}
\begin{proof}
\begin{enumerate}
\item[(a)] \emph{The prime $\ell$.} As $(\ell,k,\psi)$ is in case III, $\ell\notin B_{\psi_0}$, $\ell-1\mid2k$, $\psi$ has odd order and $\psi\ne1$. As $3\in B_{\psi_0}$ and $\ell$ is odd, $\ell\ge5$; and $\ell$ divides none of $o_{\psi_0}$, $N_{\psi_0}$ and $p_1(\psi_0)$. As $\ell-1\mid2k$, $\ell\le2k+1<2t+1$, and $2t+1<q$ and $2t+1<d$ as $h$ is large for $t$; so $\ell\nmid qd$, and $\ell\nmid N_h=\operatorname{lcm}(2^{e+2},N_{\psi_0})\,qd$.
\item[(b)] \emph{The order of $\psi$.} For $u\in(\mathbf Z/N_h)^\times$, $u^{2^{e+1}o_{\psi_0}}$ is $1$ modulo $2^{e+2}$, lies in the kernel of $\psi_0$ and is a square modulo $q$ and $d$, so it lies in $H_h$. So $\psi^{2^{e+1}o_{\psi_0}}=1$, and as $m$ is odd, $m$ divides the odd part of $o_{\psi_0}$; so $\ell\nmid m$.
\item[(c)] \emph{$S_b$.} For $p\in S$, $\psi(I_p)$ consists of $m$-th roots of unity, so its order is prime to $\ell$; so $S_b$ is empty.
\item[(d)] \emph{Inertia at $\ell$.} The characters $\psi$, $\xi_0$ and $\xi'$ with $\xi=\xi_0\xi'$ (Definition~\ref{sha-siegel-datum}) have moduli $N_h$, $p_1(\psi_0)$ and a power of $2$, all prime to $\ell$ by (a). By Lemma~\ref{sha-siegel-step-towers-kappa-dirichlet}(2) they are trivial on $I_w(\mathbf Q)$ for $w$ above $\ell$.
\end{enumerate}
\end{proof}
% host3 (DirGPowOne.lean, 101 lines): (a) is bchiFacts (BchiFacts.lean); (b) is hostOdd (HostOdd.lean, through hostComp and hostIndex, the last part of which is the claim on u^(2^(e+1) o)); (c) is the private not_dvd_card; (d) is the private dirG_unr.

\begin{lemma}\label{sha-siegel-step-towers-three-fields}
Let $S'$ be a finite set of primes, and $H\subseteq G$ an open subgroup that contains the commutator subgroup of $G$ and a complex conjugation. Then there are an imaginary quadratic field $K$, a totally real number field $F$, Galois over $\mathbf Q$ with abelian Galois group, of even degree, and a subgroup $J\subseteq G$, such that: $D_w(\mathbf Q)\subseteq J$ for every place $w$ above a prime of $S'$; $g^2\in J$ for every $g\in G$; $J\subseteq G_K$; $G_F\subseteq H$; $G_F\cap G_K=H\cap J$; and $I_w(\mathbf Q)\cap G_F\subseteq G_K$ for every place $w$.
\end{lemma}
\leanref{Development/AutomorphicArithmetic/AutomorphicVarieties/WeakShimuraTheory/FrobeniusFields/FixAb.lean}{UpperTheory.V.Ws.fixAb}
\leanref{Development/AutomorphicArithmetic/AutomorphicVarieties/WeakShimuraTheory/FrobeniusFields/FixInf.lean}{UpperTheory.V.Ws.fixInf}
\begin{proof}
\begin{enumerate}
\item[(a)] \emph{Two characters.} By Lemma~\ref{sha-siegel-step-towers-quadratic} there are a prime $p\notin S'$ and a quadratic character $\eta$ at $p$ of sign $-1$ split over $S'$, and a prime $p'\notin S'\cup\{p\}$ and a quadratic character $\eta'$ at $p'$ of sign $-1$ split over $S'\cup\{p\}$ (Definition~\ref{sha-siegel-step-towers-input}). Let $J=\ker\eta\cap\ker\eta'$.
\item[(b)] \emph{The field $K$.} By infinite Galois theory, as in step (b) of the proof of Lemma~\ref{sha-siegel-step-towers}(1), $\ker\eta=G_K$ for a quadratic field $K$, which is imaginary as no complex conjugation lies in $\ker\eta$.
\item[(c)] \emph{The field $F$.} The subgroup $H\cap\ker(\eta\eta')$ is open and contains the commutator subgroup, so it is $G_F$ for a number field $F$, Galois over $\mathbf Q$ with abelian Galois group, whose degree is its index. Let $\sigma$ be a complex conjugation. It is conjugate to the complex conjugation in $H$, and $H$ is normal, as it contains the commutator subgroup; so $\sigma\in H$, and $\eta\eta'(\sigma)=(-1)(-1)=1$. So $F$ is totally real. Some $g\in I_{v_p}(\mathbf Q)$ has $\eta(g)=-1$, and $\eta'(g)=1$ as $\eta'$ is split over $p$; so $\eta\eta'$ is onto $\{\pm1\}$, its kernel has index $2$, and $2$ divides the degree of $F$.
\item[(d)] \emph{The properties.} $D_w(\mathbf Q)\subseteq J$ for $w$ above a prime of $S'$, as $\eta$ and $\eta'$ are split over $S'$. As $\eta$ and $\eta'$ have values $\pm1$, $g^2\in J$. Clearly $J\subseteq\ker\eta=G_K$ and $G_F\subseteq H$, and $G_F\cap G_K=H\cap\ker(\eta\eta')\cap\ker\eta=H\cap J$. Let $w$ be a place and $g\in I_w(\mathbf Q)\cap G_F$. If $w$ is above $p$, then $\eta'(g)=1$, as $\eta'$ is split over $p$, and $\eta\eta'(g)=1$, so $\eta(g)=1$; otherwise $\eta(g)=1$, as $\eta$ is trivial on inertia at places not above $p$. So $g\in G_K$.
\end{enumerate}
\end{proof}
% The private fieldsEx of NegSqPowNeOneOfOddOrder.lean: (a) is quadEx twice; (b) is fixAb and fixInf for K; (c) is fixAb and fixInf for F, with the private complex_conjugations_conjugate and the index of ker(eta eta'^{-1}); (d) is the final refine.

\begin{lemma}\label{sha-siegel-step-towers-three-elements}
Let $D$ be a Siegel datum of case III, attached to $(h,k,\psi)$, with $\xi=\xi_0\xi'$ as in Definition~\ref{sha-siegel-datum}, $\xi'$ primitive and odd modulo $2^n$ with $n=n_{\mathrm{III}}$, and let $m$ be the order of $\psi$.
\begin{enumerate}
\item Some $g_1\in I_2$ has $\xi(g_1)=-1$, $\psi(g_1)=1$ and $\omega(g_1)=1$.
\item Let $s\in\{1,-1\}$. Some $g_2\in G$ has $\omega(g_2)=1$, $\xi(g_2)^2=\psi(g_2)^s$, $\xi(g_2)^m=1$ and $\xi(g_2)\ne1$.
\end{enumerate}
\end{lemma}
\leanref{Development/AutomorphicArithmetic/AutomorphicVarieties/WeakShimuraTheory/DirichletGroupPowers/EltTwo.lean}{UpperTheory.V.Ws.large₂.eltTwo}
\leanref{Development/AutomorphicArithmetic/AutomorphicVarieties/WeakShimuraTheory/DirichletGroupPowers/EltOdd.lean}{UpperTheory.V.Ws.large₂.eltOdd}
\leanref{Development/AutomorphicArithmetic/AutomorphicVarieties/WeakShimuraTheory/DirichletGroupPowers/ClsIndep.lean}{UpperTheory.V.Ws.large₂.clsIndep}
\leanref{Development/AutomorphicArithmetic/AutomorphicVarieties/WeakShimuraTheory/DirichletGroupPowers/ImgRoots.lean}{UpperTheory.V.Ws.large₂.imgRoots}
\begin{proof}
Use the characters $c_M$ of Lemma~\ref{sha-siegel-step-towers-kappa-dirichlet}, and write $p_1=p_1(\psi_0)$, an odd prime that divides neither $N_h$ nor $\ell$ (Definition~\ref{sha-host-datum} and Lemma~\ref{sha-siegel-step-towers-three-host}). By Lemma~\ref{sha-siegel-step-towers-three-host}, $m$ is odd and $\ell\nmid N_h$.
\begin{enumerate}
\item[(a)] \emph{Independence.} Let $M,M'\ge1$ be coprime and $a\in(\mathbf Z/M)^\times$. By the Chinese remainder theorem there is $u\in(\mathbf Z/MM')^\times$ with $u\equiv a$ modulo $M$ and $u\equiv1$ modulo $M'$, and as $c_{MM'}$ is onto, some $g\in G$ has $c_M(g)=a$ and $c_{M''}(g)=1$ for every $M''\mid M'$.
\item[(b)] \emph{Values of $\xi_0$.} The character $\xi_0$ has order $m$, so its image consists of $m$-th roots of unity, and $\xi_0^{r}=1$ for the number $r$ of elements of its image, so $m\le r$. As the domain $\mathcal O$ has at most $m$ roots of unity of order dividing $m$, the image of $\xi_0$ is the set of all of them.
\item[(c)] \emph{Proof of (1).} By Lemma~\ref{sha-siegel-step-towers-kappa-dirichlet}(2), $c_{2^n}$ maps $I_2$ onto $(\mathbf Z/2^n)^\times$, and $c_{p_1}$ is trivial on $I_2$, as $p_1$ is odd. Let $g\in I_2$ with $c_{2^n}(g)=-1$; then $\xi(g)=\xi'(-1)=-1$, as $\xi'$ is odd. Let $g_1=g^m$. Then $\xi(g_1)=(-1)^m=-1$, $\psi(g_1)=1$, and $\omega(g_1)=1$, as $2$ is not $\ell$ (Lemma~\ref{sha-siegel-step-towers-kappa-units}(3)).
\item[(d)] \emph{Proof of (2).} As $\psi\ne1$, some $a\in(\mathbf Z/N_h)^\times$ has $\psi(a)\ne1$. By (a) there is $u$ with $c_{N_h}(u)=a$ and $c_\ell(u)=c_{p_1}(u)=1$. Let $y=\psi(a)^{2^ns}$. Then $y\ne1$, as $\psi(a)^m=1$ and $2^n$ is prime to $m$, and $y^m=1$. With $m=2i+1$, $x=y^{i+1}$ has $x^2=y$ and $x^m=1$. By (b), $x=\xi_0(b)$ for some $b\in(\mathbf Z/p_1)^\times$, and by (a) there is $v$ with $c_{p_1}(v)=b$ and $c_{N_h}(v)=c_\ell(v)=c_{2^n}(v)=1$. Let $g_2=u^{2^n}v$. Then $c_{p_1}(g_2)=b$ and $c_{2^n}(g_2)=c_{2^n}(u)^{2^n}=1$, as the exponent of $(\mathbf Z/2^n)^\times$ divides $2^n$; so $\xi(g_2)=\xi_0(b)=x$. Also $\psi(g_2)=\psi(a)^{2^n}$, and $c_\ell(g_2)=1$, so $\omega(g_2)=1$ by Lemma~\ref{sha-siegel-step-towers-kappa-units}(3). Hence $\xi(g_2)^2=y=\psi(g_2)^s$, $\xi(g_2)^m=1$, and $\xi(g_2)\ne1$ as $x^2=y\ne1$.
\end{enumerate}
\end{proof}
% (a) is clsIndep (ClsIndep.lean, 39 lines, the private crt_unit and Chr.clsOnto); (b) is imgRoots (ImgRoots.lean, 41 lines); (c) is eltTwo (EltTwo.lean, 55 lines, the private neg_mem_range and not_over_odd, omT); (d) is eltOdd (EltOdd.lean, 103 lines, the private half, pow_ord and pow_two_pow, with bchiFacts and pK_pOne_prime_conditions for the coprimality).

\begin{lemma}\label{sha-siegel-step-towers}
Let $D$ be a Siegel datum, with the notation of Definition~\ref{sha-siegel-tower-towers}.
\begin{enumerate}
\item Suppose that $D$ is of case II or of case III. For every relabelling $\epsilon$ satisfying (a) there is a fit tame tower for $D$ with relabelling $\epsilon$.
\item Let $T\le T'$ be towers for $D$. If $T$ satisfies (a) to (m), so does $T'$; if $T$ satisfies (n) to (r), so does $T'$; so if $T$ is fit, so is $T'$; and if $T$ is deep, so is $T'$. Moreover $T\le T$, and $T'\le T''$ implies $T\le T''$.
\item Let $T$ be a fit tame tower for $D$ and let $n\ge0$. There are a tower $T'\ge T$ and $m\ge0$ such that $T'$ is the layer of index $m$ of $T$, $T'$ is deep, and $d_{T'}\ge n$.
\end{enumerate}
\end{lemma}
\leanref{Development/AutomorphicArithmetic/AutomorphicVarieties/WeakShimuraTheory/CompactIntegralFields/ExistsTameTowerOkTwo.lean}{UpperTheory.V.Ws.exists_tame_tower_ok_two}
\leanref{Development/AutomorphicArithmetic/AutomorphicVarieties/WeakShimuraTheory/CompactIntegralFields/NegSqPowNeOneOfOddOrder.lean}{UpperTheory.V.Ws.exists_tame_tower_ok_three}
\leanref{Development/AutomorphicArithmetic/AutomorphicVarieties/WeakShimuraTheory/CompactIntegralFields/PowOneInj.lean}{UpperTheory.V.Ws.leFit}
\leanref{Development/AutomorphicArithmetic/AutomorphicVarieties/WeakShimuraTheory/CompactIntegralFields/LayerEx.lean}{UpperTheory.V.Ws.layerEx}
\begin{proof}
(1) Suppose first that $D$ is of case II, and let $\epsilon$ be a relabelling satisfying (a).
\begin{enumerate}
\item[(a)] \emph{A tower input.} By Lemma~\ref{sha-siegel-step-towers-kappa}(1) there is an auxiliary character $\kappa$ for $D$ (Definition~\ref{sha-siegel-step-towers-input}). Lemma~\ref{sha-siegel-step-towers-quadratic} gives a prime $p_3\notin S$ and a quadratic character $\varphi$ at $p_3$ of sign $1$ split over $S$; then a prime $p\notin S\cup\{p_3\}$ and a quadratic character $\eta$ at $p$ of sign $-1$ split over $S\cup\{p_3\}$; then a prime $p'\notin S\cup\{p_3,p\}$ and a quadratic character $\eta'$ at $p'$ of sign $-1$ split over $S\cup\{p_3,p\}$. By (2) and (3) of Lemma~\ref{sha-siegel-step-towers-kappa}, applied to $\kappa$, these data form a tower input $X$.
\item[(b)] \emph{The fields.} By Lemma~\ref{sha-siegel-step-towers-subgroup}(2), $H_X$ is an open subgroup of $G$ that contains every commutator. By infinite Galois theory, $H_X=G_F$ for a number field $F$, Galois over $\mathbf Q$ with abelian Galois group, whose degree is the index of $H_X$, which is even. The field $F$ is totally real: an embedding of $F$ into $\mathbf C$ extends to an embedding $j$ of $\bar{\mathbf Q}$, complex conjugation composed with $j$ is $j\circ\sigma$ for a complex conjugation $\sigma\in G$, and $\sigma\in H_X$ fixes $F$. In the same way $\ker\eta=G_K$ for a quadratic field $K$, which is imaginary since no complex conjugation lies in $\ker\eta$, and $\ker\varphi=G_{k_0}$ for a quadratic field $k_0$, which is real since every complex conjugation lies in $\ker\varphi$. As $\eta$ and $\varphi$ are split over $S$, $D_w(\mathbf Q)$ lies in $G_K$ for every place $w$ above a prime of $S$, and in $G_{k_0}$ for $w$ above $2$ or $\ell$.
\item[(c)] \emph{The tower.} Let $a,u,t$ be as in Lemma~\ref{sha-siegel-step-towers-elements}. Then $t\in H_X=G_F$, and Lemma~\ref{sha-siegel-step-towers-make}, which applies as $\ell\in S$, gives a tower $T$ with relabelling $\epsilon$, fields $K$ and $F$, and $g_0^T=t$. Put $L=FK$, so that $G_L=G_F\cap\ker\eta$ and $a,u,t\in G_L$; on $G_F$, $\bar\Psi=\varphi$, $\eta=\eta'$ and $\kappa=1$. For $g\in G$ and $w'=(w_1,w_2)\in\mathbf Z^2$, $\bar\chi^{w'}(g)=\bar\xi(g)^{w_1}\bar\Psi(g)^{\pm w_2}$, with the sign $+$ if $\epsilon=0$ and $-$ if $\epsilon=1$: indeed $\chi_1\chi_2^{-1}=\alpha$, whose residue is $\bar\xi$ as $\bar\omega=\bar\chi_{\mathrm{cyc}}$, and $\chi_1\chi_3^{-1}=\Psi^{\pm1}$. Likewise $\bar\theta=\bar\xi\bar\Psi$ and $\bar\alpha=\bar\xi$, so $\bar r_0(g)=1$ if and only if $\bar\xi(g)=1$ and $\bar\Psi(g)=1$. As $\xi$ has order a power of $2$, so has $\bar\xi(g)$; so if $\bar\xi(g)\ne1$, then $\bar\xi(g)^n\ne1$ for every odd $n$.
\item[(d)] \emph{Conditions (a) to (d), (i) to (k), and tameness.} (a) is the hypothesis on $\epsilon$, (b) holds as $d_T$ is the index of $H_X$, and (d) is in (b) above. For (c), let $g\in I_w(F)$. If $w$ is above $p$, it is not above $p'$, so $\eta(g)=\eta'(g)=1$; otherwise $\eta(g)=1$, as $\eta$ is trivial on inertia at places not above $p$. So $g\in G_F\cap\ker\eta=G_L$. For (i), let $g\in I_w(F)$ with $w$ not above $\ell$ and $\bar r_0(g)=1$; then $\bar\Psi(g)=1$ and $\kappa(g)=1$, so $\psi(g)=1$, as $\kappa$ is an auxiliary character. For (j), let $m$ be as in the definition of an auxiliary character: if $g\in G_F$ and $\chi_{\mathrm{cyc}}(g)=1$, then $\kappa(g)=1$, so $\psi(g)^m=1$. For (k), let $w$ be above $\ell$ and $g\in I_w(F)$. Then $\xi(g)=1$, as $\xi$ is trivial on the inertia groups above $\ell$, and $\varphi(g)=1$, as $w$ is not above $p_3$; so $\bar\Psi(g)=\varphi(g)=1$ and $\bar r_0(g)=1$. The tower is tame by the last assertion of Lemma~\ref{sha-siegel-step-towers-subgroup}(2).
\item[(e)] \emph{Conditions (e) to (h) and (l).} For (e), let $w'\in\mathbf Z^2$. If $w_1$ is even, $\bar\chi^{w'}(a)=(-1)^{w_1}=1\ne-1=\bar\omega(a)$. If $w_1$ is odd, $\bar\chi^{w'}(u)=\bar\xi(u)^{w_1}\ne1=\bar\omega(u)$, by (c), as $\bar\xi(u)\ne1$. For (f), a root $w'$ with $w_1$ odd has $\bar\chi^{w'}(u)\ne1$ in the same way, and the other roots are $\pm(0,1)$, with $\bar\chi^{w'}(t)=\bar\Psi(t)^{\pm1}=-1$. For (h), $\bar\chi^{(1,0)}(u)^n=\bar\xi(u)^n\ne1$ for $1\le n\le4$: for $n=1,2,4$ since $\bar\xi(u)^4\ne1$, and for $n=3$ by (c). For (g), let $\eta_0$ be as there, and choose $i$ with $\chi_i=\delta_4=1$; then $\eta_0=\bar\chi_j$ on $G_L$ for some $j$. As $\omega(u)=\omega(t)=1$, $\eta_0(u)=\eta_0(t)=1$. But $\bar\theta(u)=\bar\xi(u)\ne1$, $\bar\Psi(t)=-1$ and $\bar\alpha(u)=\bar\xi(u)\ne1$, so $\chi_j=\delta_4$ and $\eta_0=1$. For (l), put $x=\bar\xi(t)$; the residues of $\theta$, $\Psi$, $\alpha$ and $1$ at $g_0^T=t$ are $-x$, $-1$, $x$ and $1$, and the $\bar\chi_i(g_0^T)$ are these values in some order. They are pairwise distinct, since $x\ne0$, $x^2\ne1$ and $\mathbf F$ has odd characteristic.
\item[(f)] \emph{Condition (m).} Let $w$ be above $\ell$. Lemma~\ref{sha-siegel-step-towers-kappa}(3) gives $g\in D_w(\mathbf Q)$ with $\omega(g)=1$, $\bar\Psi(g)=1$, $\kappa(g)=1$ and $\bar\xi(g)^4\ne1$. Then $g^2\in H_X$, as $\bar\Psi(g^2)=1=\varphi(g)^2$, $\eta(g^2)=1=\eta'(g^2)$ and $\kappa(g^2)=1$. So $g^2\in D_w(F)$, $\bar\omega(g^2)=1$, $\bar\chi^{(0,1)}(g^2)=\bar\Psi(g)^{\pm2}=1$ and $\bar\chi^{(1,0)}(g^2)=\bar\xi(g)^2\ne1$.
\item[(g)] \emph{Conditions (n) to (r).} On $G_F$, $\bar\chi^{(0,1)}=\bar\Psi^{\pm1}=\varphi$, with values $\pm1$. This gives (n), and with (b) it gives (p) for the field $k_0$. For (o), let $g\in I_w(F)$: if $w$ is above $p_3$, then $\bar\Psi(g)=1$, as $p_3\notin S$ (first paragraph of the proof of Lemma~\ref{sha-siegel-step-towers-subgroup}); otherwise $\varphi(g)=1$. And for $w$ above $\ell$ and $g\in D_w(F)$, $\varphi(g)=1$, as $\varphi$ is split over $S$. Condition (q) is part of the tower input. For (r), let $w$ be above $\ell$ and $g\in D_w(\mathbf Q)$ with $\bar\Psi(g)=1$, $\kappa(g)=1$ and $\bar\xi(g)^4\ne1$, as in the tower input. As $\varphi$, $\eta$ and $\eta'$ are split over $S$, they are trivial at $g$, so $g\in H_X$, $g\in D_w(F)$ and $\bar\chi^{(1,0)}(g)^4=\bar\xi(g)^4\ne1$. For $w$ above $2$, the same argument applies to the element $g\in I_w(\mathbf Q)\subseteq D_w(\mathbf Q)$ of the tower input, and gives $g\in I_w(F)$.
\end{enumerate}
So $T$ is a fit tame tower with relabelling $\epsilon$.

Suppose now that $D$ is of case III, attached to $(h,k,\psi)$, and let $\epsilon$ be a relabelling satisfying (a). Let $m$ be the order of $\psi$, and let $s=1$ if $\epsilon=0$ and $s=-1$ if $\epsilon=1$, so that $\chi_1\chi_2^{-1}=\alpha$, $\chi_1\chi_3^{-1}=\Psi^s$ and $\chi_1\chi_4^{-1}=\alpha\Psi^s$. Recall that $\bar\alpha=\bar\xi$ and $\bar\Psi=\bar\omega^k\bar\psi$, as $\bar\omega=\bar\chi_{\mathrm{cyc}}$.
\begin{enumerate}
\item[(a)] \emph{The subgroup $H$.} By Lemma~\ref{sha-siegel-step-towers-three-host}, $\ell\ge5$, $\ell-1\mid2k$, $m$ is odd and prime to $\ell$, $S_b$ is empty, and $\xi$ and $\psi$ are trivial on the inertia groups above $\ell$. As $\omega^{\ell-1}=1$, $\omega^{2k}=1$. Let $H$ be the set of $g\in G$ with $\bar\xi(g)^2=\bar\Psi(g)^s$, the kernel of a character of $G$. It contains $\ker\xi\cap\ker\psi\cap\ker\omega$, so it is open; it contains the commutator subgroup; and it contains $c$, as $\xi(c)=-1$ and $\Psi(c)=1$. For $w$ above $\ell$ and $g\in I_w(\mathbf Q)$, $g^2\in H$, as $\xi(g^2)=\psi(g^2)=1$ and $\omega(g^2)^k=1$.
\item[(b)] \emph{The fields.} Lemma~\ref{sha-siegel-step-towers-three-fields}, with $S'=S$ and this $H$, gives $K$, $F$ and $J$. Put $L=FK$, so that $G_L=G_F\cap G_K=H\cap J$.
\item[(c)] \emph{Elements.} Lemma~\ref{sha-siegel-step-towers-three-elements} gives $g_1\in I_2$, and $g_2$ for $s$. Let $g_0=g_1g_2^2$ and $y=\xi(g_2)$. Then $\xi(g_0)=-y^2$, $\omega(g_0)=1$ and $\psi(g_0)=\psi(g_2)^2$, so $\xi(g_0)^2=y^4=\psi(g_0)^s$ and $\bar\Psi(g_0)=\bar\psi(g_0)$; so $g_0\in H$. As $2\in S$, $g_1\in D_{v_2}(\mathbf Q)\subseteq J$, and $g_2^2\in J$; so $g_0\in G_L$. For $0<i\le4$, $\bar\xi(g_0)^i\ne1$. Indeed, $y\ne1$ and $y^m=1$; if $i$ is even, $(-y^2)^i=y^{2i}$ with $2i\mid8$ prime to $m$, so $(-y^2)^i\ne1$; if $i$ is odd, $((-y^2)^i)^m=-1$. So $\xi(g_0)^i\ne1$, and as $(\xi(g_0)^i)^{2m}=1$ with $2m$ prime to $\ell$, $\bar\xi(g_0)^i\ne1$ by Lemma~\ref{sha-siegel-step-towers-kappa-units}(1). Finally, $c_\ell$ maps $I_\ell$ onto $(\mathbf Z/\ell)^\times$ (Lemma~\ref{sha-siegel-step-towers-kappa-dirichlet}(2)) and $\ker\omega=\ker c_\ell$, so $\omega(I_\ell)$ has $\ell-1>2$ elements, and some $g_3\in I_\ell$ has $\omega(g_3^2)\ne1$. Let $g_4=g_3^2$; it lies in $H$ by (a) and in $J$, so $g_4\in G_L$.
\item[(d)] \emph{The tower.} As $\ell\in S$, $D_w(\mathbf Q)\subseteq J\subseteq G_K$ for $w$ above $\ell$, and Lemma~\ref{sha-siegel-step-towers-make} gives a tower $T$ with relabelling $\epsilon$, fields $K$ and $F$, and $g_0^T=g_0$.
\item[(e)] \emph{Characters on $G_F$.} For $g\in G_F\subseteq H$, $\bar\chi^{(1,0)}(g)=\bar\alpha(g)=\bar\xi(g)$ and $\bar\chi^{(0,1)}(g)=\bar\Psi(g)^s=\bar\xi(g)^2$; so $\bar\chi^{w'}(g)=\bar\xi(g)^{w_1+2w_2}$ for $w'=(w_1,w_2)$, and $\bar\chi_i(g)=\bar\chi_1(g)\bar\xi(g)^{1-i}$. By (c), the $\bar\chi_i(g_0)$ are pairwise distinct; and for every root $w'$, $w_1+2w_2$ is nonzero of absolute value at most $3$, so $\bar\chi^{w'}(g_0)\ne1$.
\item[(f)] \emph{Conditions (a) to (m).} (a) is the hypothesis on $\epsilon$, and (b) holds as $F$ has even degree. (c): $I_w(F)\subseteq I_w(\mathbf Q)\cap G_F\subseteq G_K$. (d): $D_w(\mathbf Q)\subseteq J\subseteq G_K$ for $w$ above a prime of $S$. (e): $\bar\chi^{w'}(g_4)=\bar\xi(g_3)^{2(w_1+2w_2)}=1$, as $\xi$ is trivial on $I_\ell$, while $\bar\omega(g_4)\ne1$ by Lemma~\ref{sha-siegel-step-towers-kappa-units}(1), as $\omega(g_4)\ne1$ and $\omega(g_4)^{\ell-1}=1$. (f) and (h) follow from (e) and (c), with $g_0\in G_L$ and the root $(1,0)$. (g): let $\eta$ be as there, and $j$ with $\eta\bar\chi_1=\bar\chi_j$ on $G_L$; as $\omega(g_0)=1$, $\eta(g_0)=1$, so $\bar\chi_1(g_0)=\bar\chi_j(g_0)$, $j=1$ by (e), and $\eta=1$. (i): let $g\in I_w(F)$ with $w$ not above $\ell$ and $\bar r_0(g)=1$; then $\bar\Psi(g)=1$ and $\omega(g)=1$ (Lemma~\ref{sha-siegel-step-towers-kappa-units}(3)), so $\bar\psi(g)=1$, and as $\psi(g)^m=1$, $\psi(g)=1$ by Lemma~\ref{sha-siegel-step-towers-kappa-units}(1). (j): $\psi^m=1$. (k): let $w$ be above $\ell$ and $g\in I_w(F)$; then $\xi(g)=1$, so $\bar\alpha(g)=1$, and $g\in H$ gives $\bar\Psi(g)^s=\bar\xi(g)^2=1$; so $\bar r_0(g)=1$. (l) is in (e). (m): let $w$ be above $\ell$; if $\bar\omega$ were trivial on $I_w(F)$, then for $g\in I_w(\mathbf Q)$, $g^2\in H\cap J\subseteq G_F$ by (a), so $g^2\in I_w(F)$ and $\omega(g^2)=1$ by Lemma~\ref{sha-siegel-step-towers-kappa-units}(1); so $\omega^2$ would be trivial on $I_w(\mathbf Q)$, which is impossible as in (c). So $\bar\omega$ is nontrivial on $I_w(F)$.
\item[(g)] \emph{The rest.} The conditions on $D$ of a fit tower in case III are in (a). The tower is tame: for $w$ above $\ell$ and $g\in I_w(\mathbf Q)$, $g^2\in G_F$, as in (f), and $2$ is prime to $\ell$.
\end{enumerate}
So $T$ is a fit tame tower with relabelling $\epsilon$; moreover $\bar\chi^{(0,1)}=(\bar\chi^{(1,0)})^2$ on $G_F$, by (e).

(2) Let $F\subseteq F'$ be the fields of $T$ and $T'$, which have the same relabelling, the same field $K$ and hence the same characters $\chi_i$; put $L'=F'K$, so that $G_{L'}=G_{F'}\cap G_K\subseteq G_L$. Let $j$ be such that $g^{\ell^j}\in G_{F'}$ for all $g\in G_F$. The residue field $\mathbf F$ of $\mathcal O$ has characteristic $\ell$, so $x\mapsto x^{\ell^j}$ is an injective endomorphism of $\mathbf F$; in particular $x^{\ell^j}=1$ implies $x=1$. For $g\in G_F$ (respectively $g\in G_L$, $g\in D_w(F)$, $g\in I_w(F)$) the element $g'=g^{\ell^j}$ lies in $G_{F'}$ (respectively $G_{L'}$, $D_w(F')$, $I_w(F')$), and $\bar\gamma(g')=\bar\gamma(g)^{\ell^j}$ for every character $\gamma$ of $G$ with values in $\mathcal O^\times$. So every condition asking for an element $g$ of one of these groups at which finitely many such residues are $1$ or differ from $1$, or at which two residues differ, passes from $T$ to $T'$ with $g'$ in place of $g$; this gives (e), (f), (h), (m) and (r). Conditions that hold for all elements of a subgroup of $G_F$, or of $G_L$, or of an inertia or decomposition group of $G_F$, pass to the corresponding subgroups of $G_{F'}$, which are contained in them; this gives (c) (as $I_w(F')\subseteq I_w(F)\cap G_{F'}\subseteq G_L\cap G_{F'}=G_{L'}$), (i), (j), (k), (n), (o), (p) and depth. Conditions (a), (d) and (q) do not involve $F$, nor do the conditions on $D$ in the definition of a fit tower in case III, and (b) holds because $d_{T'}=d_T[F':F]$. For (l), $g_0^{T'}=(g_0^T)^{\ell^{j'}}$, so $\bar\chi_i(g_0^{T'})=\bar\chi_i(g_0^T)^{\ell^{j'}}$, and these are pairwise distinct by the injectivity of $x\mapsto x^{\ell^{j'}}$.
For (g), let $\eta:G_{L'}\to\mathbf F^\times$ be as there for $T'$. Choose $i_1$ with $\eta\bar\chi_1=\bar\chi_{i_1}$ on $G_{L'}$, and put $\eta_0=\bar\chi_{i_1}\bar\chi_1^{-1}$ on $G_L$, so that $\eta_0(g)^{\ell^j}=\eta(g^{\ell^j})$ for $g\in G_L$. Then $\eta_0$ satisfies the hypotheses of (g) for $T$: $(\eta_0(g)^2)^{\ell^j}=\eta(g^{\ell^j})^2=1$; if $\omega(g)=1$ then $\omega(g^{\ell^j})=1$, so $\eta_0(g)^{\ell^j}=1$; and if $\eta\bar\chi_i=\bar\chi_{i'}$ on $G_{L'}$, then $\big(\eta_0(g)\bar\chi_i(g)\bar\chi_{i'}(g)^{-1}\big)^{\ell^j}=1$ for $g\in G_L$. So $\eta_0=1$ by (g) for $T$, and $\eta=\bar\chi_{i_1}\bar\chi_1^{-1}=\eta_0=1$ on $G_{L'}\subseteq G_L$.
Finally $T\le T$ with $j=j'=0$, and if $T'\le T''$ with exponents $k$ and $k'$, then $T\le T''$ with exponents $j+k$ and $j'+k'$.

(3) Let $F$ be the field of $T$, and put $\chi'=\chi_{\mathrm{cyc}}^{\ell-1}$, a continuous character with values in $1+\ell\mathbf Z_\ell$.
\begin{enumerate}
\item[(a)] \emph{Depth.} By (j) there is $m_1$ prime to $\ell$ with $\psi(g)^{m_1}=1$ for $g\in G_F$ with $\chi_{\mathrm{cyc}}(g)=1$. Let $U$ be the kernel of $\psi^{(\ell-1)m_1}$, an open subgroup of $G$. If $g\in G_F$ and $\chi'(g)=1$, then $g^{\ell-1}\in G_F$ has $\chi_{\mathrm{cyc}}(g^{\ell-1})=1$, so $g\in U$. The sets of $g\in G_F\setminus U$ with $\chi'(g)\equiv1\pmod{\ell^{a}}$ are closed in the compact group $G$ and decrease with $a$, with empty intersection; so one of them is empty: there is $m_0$ such that every $g\in G_F$ with $\chi'(g)\equiv1\pmod{\ell^{m_0}}$ lies in $U$.
\item[(b)] \emph{Separation.} As $T$ is tame, some power $s=s_0^k$ with $k$ prime to $\ell$ lies in $G_F$; $\chi_{\mathrm{cyc}}(s)$ has infinite order, since $\chi_{\mathrm{cyc}}(s_0)$ has. So $\chi'(s^t)\ne1$ for $t>0$, and there is $m\ge m_0$ such that $\chi'(s^t)\not\equiv1\pmod{\ell^{m+1}}$ for $0<t<n$.
\item[(c)] \emph{The field.} Let $H$ be the set of $g\in G_F$ with $\chi'(g)\equiv1\pmod{\ell^{m+1}}$, an open subgroup of $G$. It contains every commutator of $G$, which lies in $G_F$ as $F$ is abelian over $\mathbf Q$ and is killed by the character $\chi'$; and it contains every complex conjugation $\sigma$, which fixes the totally real field $F$ and has $\chi'(\sigma)=(\pm1)^{\ell-1}=1$. So $H=G_{F'}$ for a totally real number field $F'\supseteq F$, Galois and abelian over $\mathbf Q$. For $g\in G_F$, $\chi'(g^{\ell^m})=\chi_{\mathrm{cyc}}(g)^{(\ell-1)\ell^m}\equiv1\pmod{\ell^{m+1}}$, as $(\ell-1)\ell^m$ is the order of $(\mathbf Z/\ell^{m+1})^\times$; so $g^{\ell^m}\in G_{F'}$.
\item[(d)] \emph{The tower.} By (d) of Definition~\ref{sha-siegel-tower-towers}, $\ell$ splits in $K$. Since $F'$ is totally real, some element of $G_{F'}$ acts nontrivially on $K$, and it exchanges the two places of $K$ above $\ell$; so each orbit of $G_{F'}$ on the places of $\bar{\mathbf Q}$ above $\ell$ contains a place that induces on $K$ the place induced by $v_\ell$. Choosing one in each orbit, the relabelling of $T$, the field $K$, the field $F'$, the element $(g_0^T)^{\ell^m}\in G_{F'}$ and these places form a tower $T'$. By (c), $T\le T'$ with $j=j'=m$, and $T'$ is the layer of index $m$ of $T$.
\item[(e)] \emph{$T'$ is deep.} For $g\in G_{F'}$, $\chi'(g)\equiv1\pmod{\ell^{m+1}}$ with $m+1\ge m_0$, so $g\in U$ by (a), and $\psi(g)^{(\ell-1)m_1}=1$, with $(\ell-1)m_1$ prime to $\ell$.
\item[(f)] \emph{The degree.} For $0\le a<b<n$, the element $s^{b-a}$ does not lie in $G_{F'}$, by (b), so $s^a$ and $s^b$ restrict to distinct embeddings of $F'$ in $\bar{\mathbf Q}$. Hence $d_{T'}=[F':\mathbf Q]\ge n$.
\end{enumerate}
% (1) in case II is exists_tame_tower_ok_two (ExistsTameTowerOkTwo.lean, 59 lines, from Ctx2 and LabOk): (a) is kapEx, quadEx three times, xi2, ev2, the two element lemmas and the structure Set2; (b) is fixAb and fixInf (FixAb.lean, FixInf.lean); (c) is elts2, exists_resQ and mkTower; (d) is ok2B (Ok2B.lean) with barFacts and away; (e) is ok2A (OddPowNe.lean, with the private odd_pow_ne, dist4, lab_surj); (f) is Keep.loc2 (BigOm.lean); (g) is two2 (Two2.lean). Lemma sha-siegel-step-towers-kappa is written out.
% (1) in case III is exists_tame_tower_ok_three (NegSqPowNeOneOfOddOrder.lean, 280 lines, from Ctx3), which gives Tower.Three and hence the conditions on D of a fit tower: (a) is large₂.host3 (Lemma sha-siegel-step-towers-three-host), h2e, hsq, hopen and hcc, with H the kernel of the private del D pr (del_iff); (b) is the private fieldsEx (Lemma sha-siegel-step-towers-three-fields) and hL; (c) is large₂.eltTwo and large₂.eltOdd (Lemma sha-siegel-step-towers-three-elements) with s = sgnP pr, hx, hw, the private neg_sq_pow_ne_one_of_odd_order and redOne for hord, and g3 from omT (o7); (d) is mkTower; (e) is the private residual_root_character_eq_pow, reduced_character_quotient_eq, hdist and hz; (f) is the refine of Tower.Ok field by field, with residual_rep_trivial_iff for (i) and (k) and Keep.loc3 (BigOm.lean) for (m); (g) is Three and TameL with 2 prime to l. With WeakShimuraTheory/DirichletGroupPowers, about 650 lines.
% (2) is leFit (PowOneInj.lean): the private up, pow_inj and pow_one_inj (the Frobenius of the residue field), ok_up for (a) to (m), with g3'_up for (g), two_up for (n) to (r) and three_up for case III. About 550 lines.
% (3) is layerEx (LayerEx.lean): (a) is Ok.kc and the private le_compact, (b) uses TameL with D.s0 and decompSigma_ord, (c) is the private le_field and le_pow, (d) is mkTower (MkTower.lean, with mt_out and mt_conj), (f) counts the embeddings s^t restricted to F'. About 530 lines. The lemma as a whole is about 3,100 lines of Lean.
\end{proof}

\begin{lemma}\label{sha-siegel-step-labels-flag}
Let $R$ be a local ring, let $X,u\in M_4(R)$ and $d,d'\in R^4$, with $\lambda=(1,0,1,0)$ as in Definition~\ref{sha-siegel-datum}. Suppose that $X\equiv\operatorname{diag}(d_1,\dots,d_4)$ modulo $\mathfrak m_R$, that $u$ is symplectic with multiplier $1$, that $u^*Xu$ is upper triangular with entries $(i,i)$ congruent to $d'_i$ modulo $\mathfrak m_R$, and that $d_i-d'_j\in R^\times$ whenever $\lambda_i\ne\lambda_j$.
Then $u_{21},u_{41}\in\mathfrak m_R$, one of $u_{11}$ and $u_{31}$ is a unit, and if $u_{11}$ is a unit, then so is $u_{11}u_{22}-u_{12}u_{21}$, and $u_{11}u_{32}-u_{12}u_{31}\in\mathfrak m_R$.
\end{lemma}
\leanref{Development/AutomorphicArithmetic/AutomorphicVarieties/WeakShimuraTheory/FlagsAndDiagonals/ResFlag.lean}{UpperTheory.V.Ws.Lab.resFlag}
\begin{proof}
Bars denote residues. As $uu^*=1$, $Xu=uT'$ with $T'=u^*Xu$; so $\operatorname{diag}(\bar d)\bar u=\bar uT$, where $T=\bar T'$ is upper triangular with $T_{jj}=\bar d'_j$. Comparing the first and second columns, $(\bar d_i-T_{11})\bar u_{i1}=0$ and $(\bar d_i-T_{22})\bar u_{i2}=\bar u_{i1}T_{12}$ for all $i$. As $\lambda_2,\lambda_4\ne\lambda_1$, the first gives $\bar u_{21}=\bar u_{41}=0$. The entry $(1,1)$ of $\bar u^*\bar u=1$ is then $\bar u_{44}\bar u_{11}-\bar u_{24}\bar u_{31}=1$, so $\bar u_{11}$ or $\bar u_{31}$ is nonzero. The identities for $i=1,3$ in both columns give $(\bar d_1-T_{22})(\bar d_3-T_{22})(\bar u_{11}\bar u_{32}-\bar u_{12}\bar u_{31})=0$, and $\bar d_1-T_{22}$ and $\bar d_3-T_{22}$ are nonzero, as $\lambda_1,\lambda_3\ne\lambda_2$. Suppose that $\bar u_{11}\ne0$ and $\bar u_{22}=0$. The entry $(4,2)$ of $\bar u^*\bar u=1$ gives $\bar u_{11}\bar u_{42}=0$, so $\bar u_{42}=0$, and then the entries $(2,1)$ and $(2,2)$ of $\bar u^*\bar u=1$, with $\bar u_{11}\bar u_{32}=\bar u_{12}\bar u_{31}$, give $\bar u_{11}=0$, a contradiction. So $\bar u_{11}\bar u_{22}-\bar u_{12}\bar u_{21}=\bar u_{11}\bar u_{22}\ne0$.
\end{proof}

\begin{lemma}\label{sha-siegel-step-labels-hensel}
Let $C$ be a henselian local ring, $I$ a finite set and $M\in M_I(C)$ a matrix with off-diagonal entries in $\mathfrak m_C$ and with $M_{ii}-M_{jj}\in C^\times$ for $i\ne j$.
Then there are $u\in M_I(C)$ congruent to $1$ modulo $\mathfrak m_C$ and $d\in C^I$ with $d_i\equiv M_{ii}$ modulo $\mathfrak m_C$ for every $i$, such that $Mu=u\operatorname{diag}(d)$.
\end{lemma}
\leanref{Development/AutomorphicArithmetic/AutomorphicVarieties/AlgebraTools/MonomialRepresentations/HenselDiag.lean}{UpperTheory.V.AlgebraTools.henselDiag}
\begin{proof}
Bars denote residues. The reduction of $M$ is $\operatorname{diag}(\bar M_{jj})$, so the characteristic polynomial $\chi(X)=\det(X-M)$ reduces to $\prod_j(X-\bar M_{jj})$. Fix $i\in I$.
\begin{enumerate}
\item[(1)] \emph{The eigenvalue.} The residue of $\chi(M_{ii})$ is $\prod_j(\bar M_{ii}-\bar M_{jj})=0$, and the residue of $\chi'(M_{ii})$ is $\prod_{j\ne i}(\bar M_{ii}-\bar M_{jj})$, which is nonzero. As $\chi$ is monic and $C$ is henselian, there is $d_i\in C$ with $\chi(d_i)=0$ and $d_i\equiv M_{ii}$.
\item[(2)] \emph{The eigenvector.} Let $A$ be the adjugate of $d_i-M$. Then $(d_i-M)A=\det(d_i-M)=\chi(d_i)=0$. The reduction of $d_i-M$ is $\operatorname{diag}(\bar M_{ii}-\bar M_{kk})_k$, so the reduction of $A$ is $\operatorname{diag}\big(\prod_{j\ne k}(\bar M_{ii}-\bar M_{jj})\big)_k$. Hence $A_{ii}$ is a unit and $A_{ki}\in\mathfrak m_C$ for $k\ne i$. The vector $v^{(i)}$ with entries $v^{(i)}_k=A_{ki}A_{ii}^{-1}$ satisfies $Mv^{(i)}=d_iv^{(i)}$, $v^{(i)}_i=1$ and $v^{(i)}_k\in\mathfrak m_C$ for $k\ne i$.
\item[(3)] \emph{The matrix.} Let $u$ be the matrix whose column $i$ is $v^{(i)}$. By (2), $u\equiv1$ and $Mu=u\operatorname{diag}(d)$.
\end{enumerate}
\end{proof}
% henselDiag (HenselDiag.lean, about 150 lines): (1) is the private exists_root (HenselianLocalRing.is_henselian on the characteristic polynomial), (2) the private exists_vec (the adjugate of scalar d - M), (3) the matrix of the columns. The Lean also proves uniqueness: any u' congruent to 1 with M u' = u' diag(d') has d' = d and u' = u diag(t); it is used only for the block diagonal conclusion of normal, which the text omits.

\begin{lemma}\label{sha-siegel-step-labels-diagonal}
Let $C$ be a complete local ring in which $2$ is a unit, and let $g\in M_4(C)$ be symplectic with a multiplier $\nu\in C^\times$, with off-diagonal entries in $\mathfrak m_C$ and with $g_{ii}-g_{jj}\in C^\times$ for $i\ne j$.
Then there is $v\in M_4(C)$, symplectic with multiplier $1$ and congruent to $1$ modulo $\mathfrak m_C$, such that $v^*gv$ is diagonal.
\end{lemma}
\leanref{Development/AutomorphicArithmetic/AutomorphicVarieties/WeakShimuraTheory/CompactIntegralFields/Normal.lean}{UpperTheory.V.Ws.normal}
\begin{proof}
The ring $C$ is henselian, being complete \cite[Tag 04GM]{Sta}. By Lemma~\ref{sha-siegel-step-labels-hensel} there are $u\equiv1$ modulo $\mathfrak m_C$ and $d_1,\dots,d_4$ with $d_i\equiv g_{ii}$ and $gu=u\operatorname{diag}(d_1,\dots,d_4)$. Let $W=u^*u$, which is congruent to $1$. From $g^*g=\nu$ we get $\nu W=u^*g^*gu=\operatorname{diag}(d)^*W\operatorname{diag}(d)$, that is $\nu W_{ij}=d_{5-i}W_{ij}d_j$. For $i=j$, $W_{ii}$ is a unit, so $d_{5-i}d_i=\nu$; in particular every $d_i$ is a unit. For $i\ne j$, $d_{5-i}d_j-\nu=d_{5-i}(d_j-d_i)$ is a unit, as $d_j-d_i\equiv g_{jj}-g_{ii}$, so $W_{ij}=0$. So $W=\operatorname{diag}(w_1,\dots,w_4)$ with $w_i\equiv1$, and $W^*=W$ gives $w_{5-i}=w_i$. Put $v=u\operatorname{diag}(1,1,w_2^{-1},w_1^{-1})$. Then $v\equiv1$, $v^*v=1$, as its entries $(i,i)$ are $t_{5-i}w_it_i=1$ for $t=(1,1,w_2^{-1},w_1^{-1})$, and $v^*gv=\operatorname{diag}(1,1,w_2^{-1},w_1^{-1})^*\,W\operatorname{diag}(d)\operatorname{diag}(1,1,w_2^{-1},w_1^{-1})$ is diagonal.
\end{proof}
% normal (Normal.lean, about 250 lines with Hensel's lemma in the form AlgebraTools.henselDiag): the Lean also shows that v is block diagonal when g is, which uses that 2 is a unit and is not used here. resFlag (ResFlag.lean, about 90 lines) reduces to the private resFlag_field.

\begin{lemma}\label{sha-siegel-step-labels}
Let $D$ be a Siegel datum, let $V$ be an ambient ring and let $(r,u)$ be a flagged representation over $V$ (Definition~\ref{sha-siegel-tower-ambient}).
There is a relabelling $\epsilon$ satisfying (a) of Definition~\ref{sha-siegel-tower-towers} with the following property. For every $g_0\in G$ at which $\bar\chi_1,\dots,\bar\chi_4$ take pairwise distinct values there is a near split rational flagged pair $(r_P,n)$ for $(\epsilon,g_0)$ (Definition~\ref{sha-siegel-tower-pairs}) such that $r_P$ is conjugate to $r$ over $\Omega_V$, and such that, if $r$ is of type $\star$, then $r_P$ has entries of weights $(0,1)$ and $(0,-1)$ and an entry of one of the weights $(-1,1)$, $(-1,0)$ and $(-1,-1)$.
\end{lemma}
\leanref{Development/AutomorphicArithmetic/AutomorphicVarieties/WeakShimuraTheory/CompactIntegralFields/Labels.lean}{UpperTheory.V.Ws.labels}
\begin{proof}
Let $A\subset V$ be the complete discrete valuation ring of Definition~\ref{sha-siegel-tower-ambient} for $(r,u)$; congruences are modulo $\mathfrak m_V$, and $\mathfrak m_A=A\cap\mathfrak m_V$. Let $w_0=1$, and let $w_1$ be the matrix with rows $(0,0,-1,0)$, $(0,0,0,-1)$, $(1,0,0,0)$ and $(0,1,0,0)$. Both are symplectic with multiplier $1$, they have vanishing entries $(i,j)$ with $\lambda_i\ne\lambda_j$, and $w_\epsilon^*\operatorname{diag}(\delta_1,\dots,\delta_4)w_\epsilon=\operatorname{diag}(\chi_1,\dots,\chi_4)$ for the characters $\chi_i$ of the relabelling $\epsilon$. As $\lambda_i=\lambda_{i'}$ when $\chi_i=\delta_{i'}$, Lemma~\ref{sha-flagged-residues}(1) shows that $\bar\chi_i(\sigma_\ell)\ne\bar\delta_j(\sigma_\ell)$ whenever $\lambda_i\ne\lambda_j$, for both relabellings.
\begin{enumerate}
\item[(1)] \emph{The relabelling.} As $r\equiv r_0$, $r(\sigma_\ell)\equiv\operatorname{diag}(\delta_i(\sigma_\ell))$, and $u^*r(\sigma_\ell)u$ is upper triangular with entries $(i,i)$ congruent to $\delta_i(\sigma_\ell)$. By Lemma~\ref{sha-siegel-step-labels-flag}, $u_{11}$ or $u_{31}$ is a unit. If $u_{11}$ is a unit, let $\epsilon=0$. Otherwise $u_{31}$ is a unit; let $\epsilon=1$. Then for $s\in D_\ell$, with $Y=u^*r(s)u$ upper triangular and $r(s)u=uY$, the entry $(3,1)$ gives $u_{31}Y_{11}=(r(s)u)_{31}\equiv\delta_3(s)u_{31}$, and $Y_{11}\equiv\delta_1(s)$; so $\bar\delta_3(s)=\bar\delta_1(s)$, that is $\bar\alpha(s)=\bar\theta(s)=\bar\alpha(s)\bar\Psi(s)$, and $\bar\Psi$ is trivial on $D_\ell$. So $\epsilon$ satisfies (a), and in both cases the entry $(1,1)$ of $u'=w_\epsilon^*u$ is a unit.
Let $w=w_\epsilon$. The matrix $w^*r(\sigma_\ell)w$ is congruent to $\operatorname{diag}(\chi_i(\sigma_\ell))$, and $u'^*(w^*r(\sigma_\ell)w)u'=u^*r(\sigma_\ell)u$. So Lemma~\ref{sha-siegel-step-labels-flag}, applied to $w^*r(\sigma_\ell)w$ and $u'$, shows that $u'_{11}u'_{22}-u'_{12}u'_{21}$ is a unit.
\item[(2)] \emph{The normal form.} Let $g_0\in G$ be such that the $\bar\chi_i(g_0)$ are pairwise distinct, and $X=w^*r(g_0)w\equiv\operatorname{diag}(\chi_i(g_0))$, a matrix over $A$ that is symplectic with multiplier $\theta(g_0)$, with off-diagonal entries in $\mathfrak m_A$ and with $X_{ii}-X_{jj}\in A^\times$ for $i\ne j$; $2$ is a unit of $A$, as $\ell$ is odd. Lemma~\ref{sha-siegel-step-labels-diagonal} gives $v\in M_4(A)$, symplectic with multiplier $1$, with $v\equiv1$ and $v^*Xv$ diagonal. The matrix $v^*u'$ is symplectic with multiplier $1$ and congruent to $u'$, so its entry $(1,1)$ and the determinant of its block in rows and columns $1$ and $2$ are units. By Lemma~\ref{sha-siegel-step-restriction-matrices}(2), with $\nu=1$ and $n_0=1$, there is a lower unipotent $n$ such that $b=n^*v^*u'$ is upper triangular; then $v^*u'=nb$, and $b$ is symplectic with multiplier $1$. Let $m=wv$, $r_P(g)=m^*r(g)m$ and $P=(r_P,n)$.
\item[(3)] \emph{A rational flagged pair.} Since $m^*m=1$, $r_P$ is a homomorphism conjugate to $r$ over $\Omega_V$; its entries lie in $A$ and it factors continuously through $A$. Each $r_P(g)$ is symplectic with multiplier $\theta(g)$; $r_P(g)=1$ whenever $r(g)=1$, in particular on the set $N_\Phi$ of the rational problem (Definition~\ref{sha-siegel-tower-ambient}(2)); and $r_P(g_0)=v^*Xv$ is diagonal. From $m^*u=nb$ and $b^*b=1$ we get $n=m^*ub^*$, so $n^*r_P(s)n=b\,(u^*r(s)u)\,b^*$. For $s\in D_\ell$ this is a product of upper triangular matrices, with entries $(i,i)$ equal to those of $u^*r(s)u$, since $b_{ii}(b^*)_{ii}=(bb^*)_{ii}=1$. So $\psi_{v_\ell,i}(s)=(u^*r(s)u)_{ii}$, which is $\delta_i(s)$ for $s\in I_\ell$. So $P$ is a rational flagged pair for $(\epsilon,g_0)$.
\item[(4)] \emph{Near split.} As $m\equiv w$ and $r(g)\equiv\operatorname{diag}(\delta_i(g))$, $r_P(g)\equiv w^*\operatorname{diag}(\delta_i(g))w=\operatorname{diag}(\chi_i(g))$. For $s\in D_\ell$, $\psi_{v_\ell,i}(s)\equiv\delta_i(s)\equiv\chi_i(s)$, the second congruence because $\bar\Psi$ is trivial on $D_\ell$ when $\epsilon=1$. Lemma~\ref{sha-siegel-step-labels-flag}, applied to $r_P(\sigma_\ell)\equiv\operatorname{diag}(\chi_i(\sigma_\ell))$ and $n$, whose entry $n_{11}=1$ is a unit, gives $n_{21},n_{41},n_{32}\in\mathfrak m_V$, and the entry $(2,1)$ of $n^*n=1$ gives $n_{43}=-n_{21}$. These are the entries of $n$ below the diagonal with weight other than $(0,-1)$.
\item[(5)] \emph{Type $\star$.} Suppose that $r$ is irreducible over $\Omega_V$. Its values span $M_4(\Omega_V)$, hence so do those of $r_P=m^{-1}rm$, and every entry of $r_P$ is somewhere nonzero, in particular the entries $(1,3)$, $(3,1)$ and $(2,3)$, of weights $(0,1)$, $(0,-1)$ and $(-1,1)$.
Suppose that $r$ is a non-split block extension, and call a matrix \emph{parabolic} if its entries $(i,j)$ with $\lambda_i>\lambda_j$ vanish. The matrix $X$ is parabolic, as $r(g_0)$ and $w$ are, and so is $v$: for $j\in\{2,4\}$, the entries $(1,j)$ and $(3,j)$ of $Xv=v\operatorname{diag}(v^*Xv)$ show that $(v_{1j},v_{3j})$ is killed by the block of $X-(v^*Xv)_{jj}$ in rows and columns $1$ and $3$, whose determinant is a unit, as it is congruent to $(X_{11}-X_{jj})(X_{33}-X_{jj})$. So $m$ and $m^*$ are parabolic, and the blocks of $r_P(g)$ in rows and columns $2$ and $4$ are those of $r(g)$ conjugated by the block of $m$ there; they span $M_2(\Omega_V)$, so the entries $(2,4)$ and $(4,2)$ of $r_P$, of weights $(0,1)$ and $(0,-1)$, are somewhere nonzero. The weights $(-1,1)$, $(-1,0)$ and $(-1,-1)$ are those of the entries $(i,j)$ with $\lambda_i<\lambda_j$. If all these entries of $r_P$ vanished, $r_P$ would be block diagonal. Let $B$ and $B'$ be the block diagonal parts of $m^*$ and $m$; then $mB+B'm^*=2$, as the entry $(i,j)$ of the left side is $(2+\lambda_i-\lambda_j)(mm^*)_{ij}$. With $Y=mB-1$, which is supported on the entries $(i,j)$ with $\lambda_i<\lambda_j$, we would get $(1-Y)r(g)(1+Y)=B'\,r_P(g)\,B$ block diagonal for all $g$, contrary to the hypothesis. So $r_P$ has an entry of one of these weights.
\end{enumerate}
% labels (Labels.lean): Lab.choice (the relabelling, with LabOk and a unit minor), sg_sep, Lab.resFlag, Lab.exNorm (the conjugating matrix wM pr v and the flag), Lab.pairOf (QPair and NearSplit), Lab.parEig and Lab.starConj (Star gives StarE), in WeakShimuraTheory/FlagsAndDiagonals. Size: about 1,600 lines of Lean with its dependencies there.
\end{proof}

\begin{definition}\label{sha-siegel-step-lattice-shapes}
Let $D$ be a Siegel datum, with the notation of Definitions~\ref{sha-siegel-datum} and~\ref{sha-siegel-tower-pairs}, and let $W_{\mathrm r}$, $W_{\mathrm n}$, $W_{\mathrm m}$ and $W_{\mathrm V}$ be as in Definition~\ref{sha-siegel-step-chains-notions}.
\begin{enumerate}
\item For a field $K$, $t=(t_1,t_2)\in K^\times\times K^\times$ and $X\in M_4(K)$, let $t\cdot X$ be the matrix with entries $t_1^{w_1}t_2^{w_2}X_{ij}$, where $(w_1,w_2)$ is the weight of $(i,j)$.
\item Let $k$ be a field and $c=(r,(n_v))$ a pair of a homomorphism $r$ with values in $M_4(k)$ and matrices $n_v\in M_4(k)$. Its \emph{tight set} is the set of roots $w$ such that $r$ has an entry of weight $w$, or some $n_v$ has a nonzero entry $(2,1)$, $(3,2)$, $(3,1)$ or $(4,1)$ if $w$ is $(-1,0)$, $(1,-1)$, $(0,-1)$ or $(-1,-1)$ respectively.
\item A set $S$ of roots has one of the \emph{six shapes} if one of the following holds: $(-1,1),(0,-1)\in S\subseteq W_{\mathrm r}$; $(-1,0),(0,-1)\in S\subseteq W_{\mathrm n}$; $(-1,-1),(0,-1)\in S\subseteq\{(-1,-1),(0,-1)\}$; $(-1,0),(1,-1)\in S\subseteq W_{\mathrm m}$; $(-1,-1),(1,-1)\in S\subseteq\{(-1,-1),(0,-1),(1,-1)\}$; or $(-1,-1),(1,0)\in S\subseteq\{(-1,-1),(0,-1),(1,-1),(1,0)\}$.
\end{enumerate}
\end{definition}
\leanref{Development/AutomorphicArithmetic/AutomorphicVarieties/WeakShimuraTheory/AdicLattices/CompanionTheory.lean}{UpperTheory.V.Ws.Lat.TConj}
\leanref{Development/AutomorphicArithmetic/AutomorphicVarieties/WeakShimuraTheory/AdicLattices/CompanionTheory.lean}{UpperTheory.V.Ws.Lat.Tight}
\leanref{Development/AutomorphicArithmetic/AutomorphicVarieties/WeakShimuraTheory/AdicLattices/CompanionTheory.lean}{UpperTheory.V.Ws.Lat.CompanionTheory}

\begin{lemma}\label{sha-siegel-step-lattice-vertex-program}
Let $W_{\mathrm V}=\{(-1,1),(-1,0),(-1,-1)\}$, and let $C$ be a set of pairs $(w,k)$ of a root $w$ (Definition~\ref{sha-siegel-tower-towers}) and an integer $k\ge0$ such that: $(w,0)\notin C$ for every root $w\ne(0,-1)$; $((0,-1),k)\in C$ for some $k$; and $(w,k)\in C$ for some $w\in W_{\mathrm V}$ and some $k$. For $(x,y)\in\mathbf Z^2$ and $(w,k)\in C$, with $w=(w_1,w_2)$, put $e_{x,y}(w,k)=w_1x+w_2y+2k$.
Then there is $(x,y)\in\mathbf Z^2$ such that $e_{x,y}(w,k)\ge0$ for all $(w,k)\in C$, and the set $S$ of the roots $w$ with $e_{x,y}(w,k)=0$ for some $(w,k)\in C$ has one of the six shapes (Definition~\ref{sha-siegel-step-lattice-shapes}).
\end{lemma}
\leanref{Development/AutomorphicArithmetic/AutomorphicVarieties/WeakShimuraTheory/CompactIntegralFields/RootSystemLinearProgramCases.lean}{UpperTheory.V.Ws.h7}
\begin{proof}
Let $F$ be the set of $(x,y)\in\mathbf Z^2$ with $e_{x,y}(w,k)\ge0$ for all $(w,k)\in C$. It contains $(0,0)$.
\begin{enumerate}
\item[(a)] \emph{The choice.} Let $((0,-1),k_0)\in C$. Every $(x,y)\in F$ has $-y+2k_0\ge0$, so $y\le2k_0$. Let $y$ be the largest integer with $(x',y)\in F$ for some $x'$. Let $(w,k_1)\in C$ with $w\in W_{\mathrm V}$; then $w_1=-1$, and every $(x',y)\in F$ has $x'\le w_2y+2k_1\le|y|+2k_1$. Let $x$ be the largest integer with $(x,y)\in F$, and write $e=e_{x,y}$. So $e(w,k)\ge0$ on $C$.
\item[(b)] \emph{Moves.} Let $(a,b)\in\mathbf Z^2$ with $(x+a,y+b)\notin F$. Then some $(w,k)\in C$ has $e(w,k)+w_1a+w_2b<0$, while $e(w,k)\ge0$. By the choice of $y$, $(x+a,y+1)\notin F$ for every $a$; by the choice of $x$, $(x+1,y)\notin F$.
\begin{itemize}
\item For $(a,b)=(1,0)$: $w_1<0$ and $e(w,k)<1$, so $e(w,k)=0$ and $w$ is $(-1,0)$, $(-1,1)$ or $(-1,-1)$. So one of these lies in $S$.
\item For $(a,b)=(0,1)$: likewise one of $(0,-1)$, $(1,-1)$ and $(-1,-1)$ lies in $S$.
\item \emph{Parity.} The two previous items give $-x+w_2y\equiv0$ and $w'_1x-y\equiv0$ modulo $2$ for some $w_2,w'_1\in\{-1,0,1\}$. If $w_2=0$, then $x$ and then $y$ are even; otherwise $x\equiv y$. So $x+y$ is even, and $e(w,k)$ is even when $w_1+w_2$ is.
\item For $(a,b)=(-1,1)$: $-w_1+w_2<0$, so $w$ is $(1,0)$ or $(0,-1)$, with $e(w,k)<1$, or $(1,-1)$, with $e(w,k)<2$ and $e(w,k)$ even. In each case $e(w,k)=0$. So one of $(0,-1)$, $(1,0)$ and $(1,-1)$ lies in $S$.
\item For $(a,b)=(1,1)$: $w_1+w_2<0$, and in the same way one of $(0,-1)$, $(-1,0)$ and $(-1,-1)$ lies in $S$.
\end{itemize}
\item[(c)] \emph{Signs.} Let $w\in S$, with $e(w,k)=0$ and $(w,k)\in C$. If $w\ne(0,-1)$, then $k\ge1$ by the first hypothesis, so $w_1x+w_2y=-2k<0$. If $w=(0,-1)$, then $y=2k\ge0$. So $(1,0)\in S$ implies $x<0$; $(-1,0)\in S$ implies $x>0$; $(-1,1)\in S$ implies $y<x$; $(1,-1)\in S$ implies $x<y$; $(1,1)\in S$ implies $x+y<0$; $(-1,-1)\in S$ implies $x+y>0$; $(0,1)\in S$ implies $y<0$; and $(0,-1)\in S$ implies $y\ge0$.
\item[(d)] \emph{The case $(0,-1)\in S$.} Then $y\ge0$, so $(0,1)\notin S$, and by (b) one of $(-1,1)$, $(-1,0)$ and $(-1,-1)$ lies in $S$. If $(-1,1)\in S$, then $x>y\ge0$, so $(1,0)$, $(1,-1)$ and $(1,1)$ are not in $S$, and $S$ has the first shape. If $(-1,1)\notin S$ and $(-1,0)\in S$, then $x>0$, so $(1,0)$ and $(1,1)$ are not in $S$; $S$ has the fourth shape if $(1,-1)\in S$, and the second otherwise. If $(-1,1)$ and $(-1,0)$ are not in $S$, then $(-1,-1)\in S$, so $x+y>0$ and $(1,1)\notin S$; $S$ has the sixth shape if $(1,0)\in S$, the fifth if $(1,0)\notin S$ and $(1,-1)\in S$, and the third otherwise.
\item[(e)] \emph{The case $(0,-1)\notin S$.} By (b), $S$ contains $(1,-1)$ or $(-1,-1)$, contains $(1,0)$ or $(1,-1)$, and contains $(-1,0)$ or $(-1,-1)$. Suppose that $(1,-1)\in S$, so $x<y$. If $(-1,0)\in S$, then $y>x>0$, so $(1,0)$, $(1,1)$, $(-1,1)$ and $(0,1)$ are not in $S$, and $S$ has the fourth shape. Otherwise $(-1,-1)\in S$, so $x+y>0$; then $(-1,1)$ and $(1,1)$ are not in $S$, and neither is $(0,1)$, as $y<0$ would give $x<y<0$. So $S$ has the sixth shape if $(1,0)\in S$ and the fifth otherwise. Suppose that $(1,-1)\notin S$. Then $(-1,-1)$ and $(1,0)$ lie in $S$, so $x<0<x+y$ and $y>0$; then $(-1,0)$, $(-1,1)$, $(1,1)$ and $(0,1)$ are not in $S$, and $S$ has the sixth shape.
\end{enumerate}
\end{proof}
% The private root_system_linear_program_cases (RootSystemLinearProgramCases.lean): (a) is the two applications of Int.exists_greatest_of_bdd, (b) is mv, up and m1 to m4 with par, (c) is sg and s0, and (d), (e) are the private sign_pattern_case_split (a case split checked by omega) and finite_sign_pattern_cases. h7 adds the choice of t from a square root of a uniformizer, which the text does in the proof of Lemma sha-siegel-step-lattice-vertex. The file has 180 lines.

\begin{lemma}\label{sha-siegel-step-lattice-vertex}
Let $D$ be a Siegel datum and $V$ an ambient ring with residue field $k_V$ (Definition~\ref{sha-siegel-tower-ambient}). Let $r:G\to M_4(V)$ be a homomorphism and $n\in M_4(V)$ a lower unipotent matrix (Lemma~\ref{sha-siegel-step-restriction-matrices}), with entries in a complete discrete valuation subring of $V$ with finite residue field through which $r$ factors continuously, such that the off-diagonal entries of $r$, and those of $n$ of weight other than $(0,-1)$, lie in $\mathfrak m_V$, and such that $r$ has entries of weights $(0,1)$ and $(0,-1)$ and an entry of one of the weights $(-1,1)$, $(-1,0)$ and $(-1,-1)$.
Then there is $t\in\Omega_V^\times\times\Omega_V^\times$ such that the matrices $t\cdot r(g)$ and $t\cdot n$ (Definition~\ref{sha-siegel-step-lattice-shapes}) have entries in $V$, and the tight set of their images over $k_V$ has one of the six shapes.
\end{lemma}
\leanref{Development/AutomorphicArithmetic/AutomorphicVarieties/WeakShimuraTheory/AdicLattices/Vertex.lean}{UpperTheory.V.Ws.Lat.vertex}
\begin{proof}
Let $A\subseteq V$ be the complete discrete valuation ring that contains the entries of $r$ and $n$. For a matrix $X$ and a root $w\in W_{\mathrm m}=\{(-1,0),(1,-1),(0,-1),(-1,-1)\}$, let $X_w$ be the entry that the tight set reads for $w$: $X_{21}$, $X_{32}$, $X_{31}$ or $X_{41}$ for $w=(-1,0)$, $(1,-1)$, $(0,-1)$ or $(-1,-1)$. For $t$ and an entry $(i,j)$ of weight $w=(w_1,w_2)$ write $t^w=t_1^{w_1}t_2^{w_2}$, so that $(t\cdot X)_{ij}=t^wX_{ij}$.
\begin{enumerate}
\item[(a)] \emph{Valuations.} Let $\pi$ be a uniformizer of $A$. As $\Omega_V$ is algebraically closed, there is $p\in\Omega_V$ with $p^2=\pi$. As $\pi$ is a nonzero element of $\mathfrak m_A\subseteq\mathfrak m_V$, $p$ is a nonzero element of $\mathfrak m_V$; so for an integer $m$, $p^m\in V$ if and only if $m\ge0$, and $p^m$ is a unit of $V$ if and only if $m=0$. Every nonzero $z\in A$ is $u\pi^k$ with $u\in A^\times$ and $k\ge0$, and $u$ is a unit of $V$. So for $(x,y)\in\mathbf Z^2$, $t=(p^x,p^y)$ and a weight $w$, $t^wz=up^{w_1x+w_2y+2k}$ lies in $V$ if and only if $w_1x+w_2y+2k\ge0$, and is a unit of $V$ if and only if $w_1x+w_2y+2k=0$. An element of $A$ that lies in $\mathfrak m_V$ lies in $\mathfrak m_A$, as units of $A$ are units of $V$.
\item[(b)] \emph{The program.} Let $C$ be the set of pairs $(w,k)$ such that either some off-diagonal entry $r(g)_{ij}$ of weight $w$ is $u\pi^k$ with $u\in A^\times$, or $w\in W_{\mathrm m}$ and $n_w=u\pi^k$ with $u\in A^\times$. By (a) and the hypotheses on the entries, $k\ge1$ for $(w,k)\in C$ with $w\ne(0,-1)$. As $r$ has an entry of weight $(0,-1)$ and one of weight in $W_{\mathrm V}$, $C$ satisfies the hypotheses of Lemma~\ref{sha-siegel-step-lattice-vertex-program}, which gives $(x,y)$ and a set $S$ of one of the six shapes. Let $t=(p^x,p^y)$.
\item[(c)] \emph{Integrality.} The diagonal entries of $t\cdot r(g)$ are those of $r(g)$, and its off-diagonal entries are $0$ or $t^wu\pi^k$ with $(w,k)\in C$; by (a) they lie in $V$. In the same way $t^wn_w\in V$ for $w\in W_{\mathrm m}$. The matrix $t\cdot n$ is lower triangular with diagonal entries $1$, and its entries below the diagonal other than $t^wn_w$ are those at $(4,3)$ and $(4,2)$. The entries $(4,3)$ and $(4,2)$ of $n^*n=1$ give $n_{43}=-n_{21}$ and $n_{42}=n_{31}-n_{21}n_{32}$. As the weight of $(4,3)$ is $(-1,0)$, and that of $(4,2)$ is $(0,-1)=(-1,0)+(1,-1)$,
\[
(t\cdot n)_{43}=-t^{(-1,0)}n_{21},\qquad(t\cdot n)_{42}=t^{(0,-1)}n_{31}-\big(t^{(-1,0)}n_{21}\big)\big(t^{(1,-1)}n_{32}\big),
\]
which lie in $V$.
\item[(d)] \emph{The tight set.} Let $w$ be a root. By (a), the image over $k_V$ of $t\cdot r$ has an entry of weight $w$ if and only if some entry $r(g)_{ij}=u\pi^k$ of weight $w$ has $w_1x+w_2y+2k=0$; and for $w\in W_{\mathrm m}$, the image of $t^wn_w$ is nonzero if and only if $n_w=u\pi^k$ with $w_1x+w_2y+2k=0$. So the tight set of the images of $t\cdot r$ and $t\cdot n$ is $S$, which has one of the six shapes.
\end{enumerate}
\end{proof}
% Lat.vertex (Vertex.lean, 111 lines, with tcEx). Its hypotheses are Amb.Small P, IsLow, the off-diagonal entries of P.r in m_V, the off-diagonal entries of P.n of weight other than -wPs in m_V, and StarE P.r. (a) and (b) are h7 (RootSystemLinearProgramCases.lean, 180 lines: the square root p of the uniformizer, val, le and eq), with the private vx_nCo_ideal for the entries of n; (c) is tcEx (LowerSymplecticEntriesEq.lean, 67 lines, with the private lower_symplectic_entries_eq and lowMem); (d) is hS, with vx_res, vx_nCo_map and vx_nCo_tconj. The hypothesis that r has an entry of weight (0,1) is not used. About 360 lines.

\begin{lemma}\label{sha-siegel-step-lattice-conjugate}
Let $D$ be a Siegel datum, $V$ an ambient ring, $\Phi$ a deformation problem with field $E$ and places $v$ (Definition~\ref{sha-siegel-tower-pairs}), and $a:\mathcal O\to V$ a ring homomorphism. Let $P=(r,(n_v))$ be a flagged pair for $\Phi$ along $a$, let $t\in\Omega_V^\times\times\Omega_V^\times$, and let $P'=(r',(n'_v))$ be a pair with entries in $V$ such that $r'(g)=t\cdot r(g)$ and $n'_v=t\cdot n_v$ (Definition~\ref{sha-siegel-step-lattice-shapes}).
\begin{enumerate}
\item $P'$ is a flagged pair for $\Phi$ along $a$, with $r'(g)_{ii}=r(g)_{ii}$, with the same diagonal entries $\psi_{v,i}(s)$ as $P$, and with $r'$ conjugate to $r$ over $\Omega_V$.
\item If the entries of $P$ lie in a complete discrete valuation subring of $V$ with finite residue field through which $r$ factors continuously, so do those of $P'$, for another such subring.
\end{enumerate}
\end{lemma}
\leanref{Development/AutomorphicArithmetic/AutomorphicVarieties/WeakShimuraTheory/AdicLattices/TcPair.lean}{UpperTheory.V.Ws.Lat.tcPair}
\leanref{Development/AutomorphicArithmetic/AutomorphicVarieties/WeakShimuraTheory/AdicLattices/TcSmall.lean}{UpperTheory.V.Ws.Lat.tcSmall}
\begin{proof}
(1) Let $d=\operatorname{diag}(t_1^{\lambda_i}t_2^{\mu_i})_i$, with $\lambda$ and $\mu$ as in Definition~\ref{sha-siegel-datum}. Then $t\cdot X=dXd^{-1}$, since the entry $(i,j)$ of $dXd^{-1}$ is $t_1^{\lambda_i-\lambda_j}t_2^{\mu_i-\mu_j}X_{ij}$. As $\lambda_i+\lambda_{5-i}=\mu_i+\mu_{5-i}=1$, $d^*d=t_1t_2$, so $d^*=t_1t_2d^{-1}$ and $(dXd^{-1})^*=(d^{-1})^*X^*d^*=dX^*d^{-1}$. So $X\mapsto t\cdot X$ is multiplicative, commutes with $X\mapsto X^*$, fixes scalar matrices and diagonal entries, and keeps vanishing entries zero. Hence: $r'(g)^*r'(g)=t\cdot(r(g)^*r(g))=a(\theta(g))$; $r'(g)=1$ when $r(g)=1$; $r'(g_0)$ is diagonal; each $n'_v$ is symplectic with multiplier $1$, lower triangular with diagonal entries $1$; and $n_v'^*r'(s)n'_v=t\cdot(n_v^*r(s)n_v)$ is upper triangular with the diagonal entries $\psi_{v,i}(s)$ of $n_v^*r(s)n_v$, which are $1$ for $s\in\mu_v$. Finally $r'(g)=d\,r(g)\,d^{-1}$ for all $g$.

(2) Let $A\subseteq V$ be such a subring for $P$, with fraction field $K_A$. For an entry $(i,j)$ of weight $(w_1,w_2)$ let $u_{ij}=t_1^{w_1}t_2^{w_2}\in\Omega_V^\times$, so that the entries $(i,j)$ of $r'(g)$ and $n'_v$ are $u_{ij}r(g)_{ij}$ and $u_{ij}(n_v)_{ij}$.
\begin{enumerate}
\item[(a)] \emph{The ring.} Let $L$ be the subfield of $\Omega_V$ generated by $K_A$ and the elements $u_{ij}$. As $\Omega_V$ is algebraic over $K_A$, $L$ is finite over $K_A$. Let $A'=V\cap L$, a valuation ring of $L$ containing $A$. An element of $A'$ that is a unit of $V$ is a unit of $A'$, as its inverse lies in $L$; so $\mathfrak m_{A'}=A'\cap\mathfrak m_V$, which contains $\mathfrak m_A$. As $A$ is complete, the valuation of $K_A$ extends uniquely to $L$, and $A'$ is the valuation ring of this extension; so $A'$ is a complete discrete valuation ring with finite residue field \cite[Chapter~II, \S2, Proposition~3]{Ser79}. And $\Omega_V$ is algebraic over $K_A$, which lies in the fraction field of $A'$.
\item[(b)] \emph{The entries.} The entries of $P'$ lie in $V$ by hypothesis, and in $L$, as products of an element $u_{ij}$ and an element of $A$; so they lie in $A'$.
\item[(c)] \emph{Continuity.} Fix an entry $(i,j)$. There are $c,d\in A'$ with $d\ne0$ and $du_{ij}=c$: if $u_{ij}\in V$, then $u_{ij}\in A'$, and $(c,d)=(u_{ij},1)$; otherwise $u_{ij}^{-1}\in V\cap L=A'$, and $(c,d)=(1,u_{ij}^{-1})$. Then $d\,r'(g)_{ij}=c\,r(g)_{ij}$ for $g\in G_E$. Let $\pi$ be a uniformizer of $A'$ and write $d=w\pi^e$ with $w\in A'^\times$ and $e\ge0$. Let $N\ge0$. As $g\mapsto r(g)_{ij}$ is continuous for the $\mathfrak m_A$-adic topology and $\mathfrak m_A\subseteq\pi A'$, the class of $r(g)_{ij}$ modulo $\pi^{N+e}A'$ is locally constant in $g$, and so is that of $d\,r'(g)_{ij}$. As $dx\in\pi^{N+e}A'$ implies $x\in\pi^NA'$, the class of $r'(g)_{ij}$ modulo $\pi^NA'$ is locally constant in $g$. So $r'$ factors continuously through $A'$.
\end{enumerate}
% (1) is Lat.tcPair (TcPair.lean, 87 lines: tc_diag, tc_zero, tc_mul, tc_adj, tc_scalar, tc_gsp, and the conjugating matrix from grp). (2) is Lat.tcSmall (TcSmall.lean, 128 lines): (a) is smallExt (CompactIntegralFields/SmallExt.lean, 230 lines with its private fin_quot, haus and key, which proves the cited facts for the valuation ring V cap L of rank at most [L:K_A] over A), applied to the u_ij and their inverses; (b) is mem_cl; (c) is frac and the private cont_div.
\end{proof}

\begin{lemma}\label{sha-siegel-step-lattice-vanishing}
Let $D$ be a Siegel datum, $\epsilon$ a relabelling and $g_0\in G$ an element at which the residues $\bar\chi_1,\dots,\bar\chi_4$ are pairwise distinct; let $k$ be a field and $a:\mathcal O\to k$ a homomorphism that kills $\mathfrak m_{\mathcal O}$. Let $c=(r,n)$ be a flagged pair for the rational problem of $(\epsilon,g_0)$ along $a$, with locally constant entries, $r(g)_{ii}=a(\chi_i(g))$ for all $g$ and $\psi_{v_\ell,i}(s)=a(\chi_i(s))$ for $s\in D_\ell$. Let $(i,j)$ be an entry whose weight $w$ is an odd root, such that $r(g)_{ix}r(h)_{xj}=0$ for all $g,h\in G$ and $x\ne i,j$, and let $m\in k$ with $r(s)_{ij}=m\big(a(\chi_j(s))-a(\chi_i(s))\big)$ for all $s\in D_\ell$. Then $r(g)_{ij}=0$ for all $g$, and $m=0$.
\end{lemma}
\leanref{Development/AutomorphicArithmetic/AutomorphicVarieties/WeakShimuraTheory/AdicLattices/EntVan.lean}{UpperTheory.V.Ws.Lat.entVan}
\leanref{Development/AutomorphicArithmetic/AutomorphicVarieties/WeakShimuraTheory/AdicLattices/OddVan.lean}{UpperTheory.V.Ws.Lat.oddVan}
\begin{proof}
Write $\bar\chi_i=a\circ\chi_i$ and $\bar\gamma=a\circ\chi^w$, so that $\chi^w\chi_j=\chi_i$, as $\chi_i\chi_{5-i}=\theta$ for all $i$.
\begin{enumerate}
\item[(a)] \emph{A cocycle.} For $g,h\in G$, $r(gh)_{ij}=\sum_xr(g)_{ix}r(h)_{xj}=\bar\chi_i(g)r(h)_{ij}+r(g)_{ij}\bar\chi_j(h)$, by the hypothesis on the products. So $z(g)=r(g)_{ij}\bar\chi_j(g)^{-1}$ is locally constant and satisfies $z(gh)=z(g)+\bar\gamma(g)z(h)$. It vanishes on the set $N_\Phi$ of the rational problem, where $r=1$, and at $g_0$, where $r$ is diagonal. For $s\in D_\ell$, $z(s)=m-\bar\gamma(s)m$.
\item[(b)] \emph{Vanishing.} The character $\chi^w$ of the odd root $w$ is one of $\alpha^{\pm1}$, $\beta^{\pm1}$, $\theta^{\pm1}$, and equals $\delta_{i'}\delta_{j'}^{-1}$ for some $i',j'$ with $\lambda_{i'}\ne\lambda_{j'}$. By (a), $z$ is a coboundary on $D_\ell$, and on $D_p$ for $p\in S_b$ by Lemma~\ref{sha-flagged-cocycles}(2); it vanishes on the elements of the inertia groups at places not above $T$ with $\bar r_0(g)=1$, which lie in $N_\Phi$. By Lemma~\ref{sha-flagged-cocycles}(1) there is $n_0\in k$ with $z(g)=\bar\gamma(g)n_0-n_0$ for all $g$. As the $\bar\chi_i(g_0)$ are pairwise distinct, $\bar\gamma(g_0)\ne1$, and $z(g_0)=0$ gives $n_0=0$; so $z=0$ and $r(g)_{ij}=0$ for all $g$.
\item[(c)] \emph{The constant $m$.} By Lemma~\ref{sha-flagged-residues}(1), $\bar\delta_{i'}(\sigma_\ell)\ne\bar\delta_{j'}(\sigma_\ell)$, so $\bar\gamma(\sigma_\ell)\ne1$. As $\sigma_\ell\in D_\ell$, $0=z(\sigma_\ell)=m-\bar\gamma(\sigma_\ell)m$, so $m=0$.
\end{enumerate}
\end{proof}
% (a) is Lat.entVan (EntVan.lean, 64 lines, with the private rt_wt and adic_bot); (b) and (c) are Lat.oddVan (OddVan.lean, 87 lines): ov_rt and ov_ij for the character, xvanExt (Lemma sha-flagged-cocycles (1)) and Lat.locVanK (Lemma sha-flagged-cocycles (2) over the field k), key for the elements at which the residue of the character is not 1, and facts for sigma_l.

\begin{lemma}\label{sha-siegel-step-lattice-exclusion}
Let $D$, $\epsilon$, $g_0$, $k$, $a$ and $c=(r,n)$ be as in Lemma~\ref{sha-siegel-step-lattice-vanishing}, and let $S$ be the tight set of $c$ (Definition~\ref{sha-siegel-step-lattice-shapes}). If $(-1,0)\notin S$, then $(-1,-1)\notin S$; and if $(-1,1)\notin S$ and $(0,1)\notin S$, then $(-1,0)\notin S$.
\end{lemma}
\leanref{Development/AutomorphicArithmetic/AutomorphicVarieties/WeakShimuraTheory/AdicLattices/NotTight.lean}{UpperTheory.V.Ws.Lat.noNeg}
\begin{proof}
Recall that $r(g)_{ii}=a(\chi_i(g))$ for $g\in G$, and that the entry $(1,1)$ of $n^*r(s)n$ is $\psi_{v_\ell,1}(s)=a(\chi_1(s))$ for $s\in D_\ell$.
\begin{enumerate}
\item[(a)] \emph{The first column.} As $n$ is symplectic with multiplier $1$, $n^*n=1$, so $nn^*=1$ and $r(s)n=n\,(n^*r(s)n)$. For $s\in D_\ell$ the matrix $n^*r(s)n$ is upper triangular, so the first column of $r(s)n$ is $a(\chi_1(s))$ times the first column of $n$. As $n$ is lower triangular with $n_{11}=1$, this reads
\[
r(s)_{i1}+r(s)_{i2}n_{21}+r(s)_{i3}n_{31}+r(s)_{i4}n_{41}=n_{i1}\,a(\chi_1(s))\qquad(1\le i\le 4,\ s\in D_\ell).
\]
\item[(b)] \emph{The weight $(-1,-1)$.} Suppose that $(-1,0)\notin S$. The entries of weight $(-1,0)$ are $(2,1)$ and $(4,3)$, so $r(g)_{21}=r(g)_{43}=0$ for all $g$, and $n_{21}=0$. The entry $(4,1)$ has weight $(-1,-1)$, an odd root. For $x\in\{2,3\}$ and $g,h\in G$, $r(g)_{4x}r(h)_{x1}=0$, as $r(h)_{21}=0$ and $r(g)_{43}=0$. The identity (a) for $i=4$, with $n_{21}=0$, $r(s)_{43}=0$ and $r(s)_{44}=a(\chi_4(s))$, gives $r(s)_{41}=n_{41}\big(a(\chi_1(s))-a(\chi_4(s))\big)$ for $s\in D_\ell$. By Lemma~\ref{sha-siegel-step-lattice-vanishing}, with $m=n_{41}$, $r(g)_{41}=0$ for all $g$ and $n_{41}=0$. As $(4,1)$ is the only entry of weight $(-1,-1)$, $(-1,-1)\notin S$.
\item[(c)] \emph{The weight $(-1,0)$.} Suppose that $(-1,1)\notin S$ and $(0,1)\notin S$. The entry of weight $(-1,1)$ is $(2,3)$, and the entries of weight $(0,1)$ are $(1,3)$ and $(2,4)$; so $r(g)_{23}=r(g)_{24}=0$ for all $g$. The entry $(2,1)$ has weight $(-1,0)$, an odd root. For $x\in\{3,4\}$ and $g,h\in G$, $r(g)_{2x}r(h)_{x1}=0$. The identity (a) for $i=2$ gives $r(s)_{21}=n_{21}\big(a(\chi_1(s))-a(\chi_2(s))\big)$ for $s\in D_\ell$. By Lemma~\ref{sha-siegel-step-lattice-vanishing}, with $m=n_{21}$, $r(g)_{21}=0$ for all $g$ and $n_{21}=0$. Let $g\in G$ and $\nu=a(\theta(g))$. From $r(g)^*r(g)=\nu$ and $r(g)r(g^{-1})=1$ we get $r(g)^*=\nu\,r(g^{-1})$. The entry $(2,1)$ of $r(g)^*$ is $r(g)_{43}$, so $r(g)_{43}=\nu\,r(g^{-1})_{21}=0$. So the entries of weight $(-1,0)$ of $r$ and the entry $(2,1)$ of $n$ vanish, and $(-1,0)\notin S$.
\end{enumerate}
\end{proof}
% Lat.noNeg (NotTight.lean, 101 lines, with SymplecticGSp4/SgnMulSelf), for any deformation problem with PreRes and EntVan; the text states it for the rational problem, where EntVan is Lemma sha-siegel-step-lattice-vanishing. (a) is the private col0 with PreRes.char; (b) and (c) apply EntVan at the Lean's entries (3,0) and (1,0), and (c) uses the inverse identity of grp for r(g)^*.

\begin{lemma}\label{sha-siegel-step-lattice-reading}
Let $D$, $\epsilon$, $g_0$, $k$, $a$ and $c=(r,n)$ be as in Lemma~\ref{sha-siegel-step-lattice-exclusion}, and suppose that the tight set $S$ of $c$ satisfies $(-1,1),(0,-1)\in S\subseteq W_{\mathrm r}$. Then $r$ is supported in $W_{\mathrm r}$, $n$ is supported in $W_{\mathrm n}$, $r$ has an entry of weight $(-1,1)$, and either $r$ has an entry of weight $(0,-1)$, or $r$ is supported in $W_{\mathrm V}$, $n_{31}\ne0$, $\bar\chi_1=\bar\chi_3$ on $D_\ell$, and there is no $m\in k$ with $r(s)_{23}=m\big(a(\chi_3(s))-a(\chi_2(s))\big)$ for all $s\in D_\ell$.
\end{lemma}
\leanref{Development/AutomorphicArithmetic/AutomorphicVarieties/WeakShimuraTheory/AdicLattices/LowCase.lean}{UpperTheory.V.Ws.Lat.shape}
\begin{proof}
The weight of every off-diagonal entry is a root. The only entry of weight $(-1,1)$ is $(2,3)$, and the entries of weight $(0,-1)$ are $(3,1)$ and $(4,2)$.
\begin{enumerate}
\item[(a)] \emph{The support of $r$.} If $r(g)_{ij}\ne0$ with $i\ne j$, the weight of $(i,j)$ lies in $S\subseteq W_{\mathrm r}$.
\item[(b)] \emph{The support of $n$.} The matrix $n$ is lower triangular. Below the diagonal, the entries $(2,1)$, $(3,1)$, $(4,1)$, $(4,2)$ and $(4,3)$ have weights $(-1,0)$, $(0,-1)$, $(-1,-1)$, $(0,-1)$ and $(-1,0)$, in $W_{\mathrm n}$, and $(3,2)$ has weight $(1,-1)$. If $n_{32}\ne0$, then $(1,-1)\in S$, which is not in $W_{\mathrm r}$. So $n_{32}=0$, and $n$ is supported in $W_{\mathrm n}$.
\item[(c)] \emph{The weight $(-1,1)$.} The tight set reads no entry of $n$ for the weight $(-1,1)$; so, as $(-1,1)\in S$, $r$ has an entry of weight $(-1,1)$, that is, $r(g)_{23}\ne0$ for some $g$.
\item[(d)] \emph{The second case.} Suppose that $r$ has no entry of weight $(0,-1)$. As $(0,-1)\in S$, $n_{31}\ne0$. A weight outside $W_{\mathrm V}$ is outside $W_{\mathrm r}$ or equal to $(0,-1)$, so by (a) $r$ is supported in $W_{\mathrm V}$.
\item[(e)] \emph{The characters on $D_\ell$.} In the case of (d), let $s\in D_\ell$ and $X=r(s)$. The entries $(1,2)$, $(1,3)$, $(1,4)$, $(3,1)$, $(3,2)$ and $(3,4)$ have weights $(1,0)$, $(0,1)$, $(1,1)$, $(0,-1)$, $(1,-1)$ and $(1,0)$, none of them in $W_{\mathrm V}$; so the first row of $X$ is $(X_{11},0,0,0)$ and the third is $(0,0,X_{33},0)$. The third row of $n^*$ is $(-n_{42},-n_{32},n_{22},n_{12})=(-n_{42},0,1,0)$, by (b). The entry $(3,1)$ of $n^*n=1$ gives $n_{31}-n_{42}=0$. The entry $(3,1)$ of $n^*Xn$, which vanishes as this matrix is upper triangular, is $-n_{42}(Xn)_{11}+(Xn)_{31}=-n_{42}X_{11}+X_{33}n_{31}=n_{31}(X_{33}-X_{11})$. As $n_{31}\ne0$, $a(\chi_1(s))=X_{11}=X_{33}=a(\chi_3(s))$. As $a$ kills $\mathfrak m_{\mathcal O}$ and $\mathbf F\to k$ is injective, $\bar\chi_1(s)=\bar\chi_3(s)$.
\item[(f)] \emph{The entry $(2,3)$.} In the case of (d), suppose that $m\in k$ has $r(s)_{23}=m\big(a(\chi_3(s))-a(\chi_2(s))\big)$ for all $s\in D_\ell$. The weight $(-1,1)$ of $(2,3)$ is an odd root. For $g,h\in G$, $r(g)_{21}r(h)_{13}=0$ and $r(g)_{24}r(h)_{43}=0$, since the entries $(1,3)$ and $(2,4)$ have weight $(0,1)$, outside $W_{\mathrm r}$. By Lemma~\ref{sha-siegel-step-lattice-vanishing}, $r(g)_{23}=0$ for all $g$, contrary to (c). So there is no such $m$.
\end{enumerate}
\end{proof}
% Lat.shape (LowCase.lean, 104 lines), for any deformation problem with PreRes and EntVan: (a) is hS, (b) is hN through the private low_case, (c) is hEbe, (d) is vV_case, (e) is the private entry_zero_zero_eq_two_two_of_isUp_conj, (f) is EntVan at the Lean's entry (1,2) with be_case.

\begin{lemma}\label{sha-siegel-step-lattice}
Let $D$ be a Siegel datum, let $V$ be an ambient ring with residue field $k_V$ (Definition~\ref{sha-siegel-tower-ambient}), let $\epsilon$ be a relabelling satisfying (a) of Definition~\ref{sha-siegel-tower-towers}, and let $g_0\in G$ be an element at which $\bar\chi_1,\dots,\bar\chi_4$ take pairwise distinct values.
Let $P=(r_P,n)$ be a near split rational flagged pair for $(\epsilon,g_0)$ such that $r_P$ has entries of weights $(0,1)$ and $(0,-1)$ and an entry of one of the weights $(-1,1)$, $(-1,0)$ and $(-1,-1)$ (Definition~\ref{sha-siegel-tower-pairs}).
Then there is a rational flagged pair $P'=(r_{P'},n')$ for $(\epsilon,g_0)$ such that $r_{P'}$ is conjugate to $r_P$ over $\Omega_V$ and the image of $P'$ over $k_V$ is of lattice shape.
\end{lemma}
\leanref{Development/AutomorphicArithmetic/AutomorphicVarieties/WeakShimuraTheory/CompactIntegralFields/LatSix.lean}{UpperTheory.V.Ws.lat}
\begin{proof}
Let $A\subset V$ be the complete discrete valuation ring that contains the entries of $P$ and through which $r_P$ factors continuously; bars denote images over $k_V$, and $a:\mathcal O\to k_V$ is the residue map, which kills $\mathfrak m_{\mathcal O}$.
\begin{enumerate}
\item[(1)] \emph{Conjugation.} As $P$ is near split, the off-diagonal entries of $r_P$, and those of $n$ of weight other than $(0,-1)$, lie in $\mathfrak m_V$, and $n$ is lower unipotent, $P$ being a flagged pair. Lemma~\ref{sha-siegel-step-lattice-vertex} gives $t$ such that $P'=(r_{P'},n')$, with $r_{P'}(g)=t\cdot r_P(g)$ and $n'=t\cdot n$, has entries in $V$ and the tight set $S$ of $c=\bar P'$ has one of the six shapes.
\item[(2)] \emph{A flagged pair.} By Lemma~\ref{sha-siegel-step-lattice-conjugate}, $P'$ is a flagged pair for the rational problem of $(\epsilon,g_0)$ along $\mathcal O\to V$, with the diagonal entries of $r_P$ and the same $\psi_{v_\ell,i}$; $r_{P'}$ is conjugate to $r_P$ over $\Omega_V$; and the entries of $P'$ lie in a complete discrete valuation subring of $V$ with finite residue field through which $r_{P'}$ factors continuously.
\item[(3)] \emph{The reduction.} As $P$ is near split, $\bar r_{P'}(g)_{ii}=\bar r_P(g)_{ii}=a(\chi_i(g))$ and $\bar\psi'_{v_\ell,i}(s)=\bar\psi_{v_\ell,i}(s)=a(\chi_i(s))$ for $s\in D_\ell$. The entries of $\bar r_{P'}$ are locally constant, since $r_{P'}$ factors continuously through a complete discrete valuation ring whose maximal ideal maps into $\mathfrak m_V$. So $c$ satisfies the hypotheses of Lemma~\ref{sha-siegel-step-lattice-vanishing}, as the $\bar\chi_i(g_0)$ are pairwise distinct.
\item[(4)] \emph{The shape of the tight set.} By Lemma~\ref{sha-siegel-step-lattice-exclusion}, if $(-1,0)\notin S$ then $(-1,-1)\notin S$, and if $(-1,1),(0,1)\notin S$ then $(-1,0)\notin S$. This excludes five of the six shapes. In the second and the fourth, $(-1,0)\in S$ while $(-1,1)$ and $(0,1)$ are not in $S$. In the third, the fifth and the sixth, $(-1,-1)\in S$ while $(-1,0)\notin S$. So $(-1,1),(0,-1)\in S\subseteq W_{\mathrm r}$.
\item[(5)] \emph{Lattice shape.} By Lemma~\ref{sha-siegel-step-lattice-reading}, $\bar r_{P'}$ is supported in $W_{\mathrm r}$ and $\bar n'$ in $W_{\mathrm n}$, so $c$ is residual; $\bar r_{P'}$ has an entry of weight $(-1,1)$; and either it has an entry of weight $(0,-1)$, or it is supported in $W_{\mathrm V}$, $\bar n'_{31}\ne0$, $\bar\chi_1=\bar\chi_3$ on $D_\ell$, and no $m\in k_V$ has $\bar r_{P'}(s)_{23}=m\big(a(\chi_3(s))-a(\chi_2(s))\big)$ for all $s\in D_\ell$. In the second case $\bar\Psi$ is trivial on $D_\ell$: $\{\chi_1,\chi_3\}=\{\theta,\alpha\}$ with $\theta=\alpha\Psi$, and $\bar\alpha(s)\ne0$. So $c$ is of lattice shape.
\item[(6)] \emph{Conclusion.} For $s\in I_\ell$, $\psi'_{v_\ell,i}(s)=\psi_{v_\ell,i}(s)=\delta_i(s)$, as $P$ is rational. With (2), $P'$ is a rational flagged pair for $(\epsilon,g_0)$, $r_{P'}$ is conjugate to $r_P$ over $\Omega_V$, and its image $c$ over $k_V$ is of lattice shape by (5).
\end{enumerate}
% lat (LatSix.lean, 61 lines): (1) is Lat.vertex (Lemma sha-siegel-step-lattice-vertex); (2) is tcPair and tcSmall; (3) is hpre (Lat.PreRes) with hN.red and hN.char; the vanishing is entVan with Lat.oddVan (hE); (4) is noNeg and the private lat_six over Lat.CompanionTheory; (5) is shape; (6) is the final refine, with the last case comparing chi 0 and chi 2 through mul_eq_left₀. Its lemmas are written out: Lemma sha-siegel-step-lattice-vertex (about 360 lines), (2) of Lemma sha-siegel-step-lattice-conjugate (128 lines), and Lemmas sha-siegel-step-lattice-exclusion (101 lines) and sha-siegel-step-lattice-reading (104 lines).
\end{proof}

\begin{lemma}\label{sha-siegel-step-restriction-matrices}
Let $R$ be a commutative ring. Call $n\in M_4(R)$ \emph{lower unipotent} if it is symplectic with multiplier $1$ (Definition~\ref{sha-siegel-datum}), lower triangular, and has diagonal entries $1$. Let $W_{\mathrm r}=\{(-1,1),(-1,0),(-1,-1),(0,-1)\}$ and $W_{\mathrm n}=\{(-1,0),(0,-1),(-1,-1)\}$.
\begin{enumerate}
\item Let $X\in M_4(R)$ be supported in $W_{\mathrm r}$ (Definition~\ref{sha-siegel-tower-pairs}), with diagonal entries in $R^\times$, and let $n_0$ be lower unipotent and supported in $W_{\mathrm n}$. Then the entry $(1,1)$ of $Xn_0$ and the determinant of its block in rows and columns $1$ and $2$ are units, and every lower unipotent $n$ with $n^*Xn_0$ upper triangular is supported in $W_{\mathrm n}$.
\item Let $\nu\in R^\times$, let $X$ be symplectic with multiplier $\nu$, and let $n_0$ be lower unipotent, such that the entry $(1,1)$ of $Xn_0$ and the determinant of its block in rows and columns $1$ and $2$ are units. Then there is a lower unipotent $n$ with $n^*Xn_0$ upper triangular.
\item Let $\nu\in R^\times$, let $X$ be symplectic with multiplier $\nu$, and let $n_0$ and $n$ be lower unipotent with $n^*Xn_0$ upper triangular. If $Y,Z\in M_4(R)$ satisfy $ZX=XY$ and $n_0^*Yn_0$ is upper triangular, then $n^*Zn$ is upper triangular, with the diagonal entries of $n_0^*Yn_0$.
\end{enumerate}
\end{lemma}
\leanref{Development/AutomorphicArithmetic/AutomorphicVarieties/WeakShimuraTheory/RootsOfUnityShapes/ShapeQ.lean}{UpperTheory.V.Ws.Rq.shapeQ}
\leanref{Development/AutomorphicArithmetic/AutomorphicVarieties/WeakShimuraTheory/RootsOfUnityShapes/DiagGsp.lean}{UpperTheory.V.Ws.Rq.luG}
\leanref{Development/AutomorphicArithmetic/AutomorphicVarieties/WeakShimuraTheory/RootsOfUnityShapes/ConjFlag.lean}{UpperTheory.V.Ws.Rq.conjFlag}
\begin{proof}
A lower unipotent $n$ has $n^*n=1$, so $nn^*=1$; a symplectic $X$ with multiplier $\nu\in R^\times$ has $XX^*=\nu$; and $X\mapsto X^*$ carries upper triangular matrices to upper triangular ones, with the diagonal reversed.

(1) The weights of the entries show that the only off-diagonal entries of $X$ that may be nonzero are $(2,1)$, $(2,3)$, $(3,1)$, $(4,1)$, $(4,2)$ and $(4,3)$, and that the entry $(3,2)$ of $n_0$ vanishes. So the first row of $Y=Xn_0$ is $(X_{11},0,0,0)$, $Y_{22}=X_{22}$ and $Y_{32}=0$, which gives the first claim. Let $n$ be lower unipotent with $U=n^*Xn_0$ upper triangular, so that $Y=nU$. Comparing the entries $(1,2)$, $(2,2)$ and $(3,2)$ gives $U_{12}=0$, $U_{22}=X_{22}$ and $n_{32}X_{22}=0$, so $n_{32}=0$; this is the only entry below the diagonal whose weight $(1,-1)$ lies outside $W_{\mathrm n}$.

(2) The matrix $Y=Xn_0\operatorname{diag}(1,1,\nu^{-1},\nu^{-1})$ is symplectic with multiplier $1$ and has the first two columns of $Xn_0$. A symplectic matrix with multiplier $1$ whose leading minors of sizes $1$ and $2$ are units is a product $nb$ with $n$ lower unipotent and $b$ upper triangular, by Gaussian elimination in the symplectic group. Then $n^*Xn_0=b\operatorname{diag}(1,1,\nu,\nu)$ is upper triangular.

(3) Let $U=n^*Xn_0$, which is upper triangular, and $U^*U=n_0^*X^*nn^*Xn_0=\nu$. Then $U(n_0^*Yn_0)U^*=n^*XYX^*n=n^*ZXX^*n=\nu\,n^*Zn$. The left side is a product of upper triangular matrices, so $n^*Zn$ is upper triangular, and its entry $(i,i)$ is $\nu^{-1}U_{ii}(n_0^*Yn_0)_{ii}(U^*)_{ii}=(n_0^*Yn_0)_{ii}$, since $(U^*)_{ii}U_{ii}=(U^*U)_{ii}=\nu$.
\end{proof}
% The elimination in (2) is Lab.luSp (FlagsAndDiagonals/LowAdj.lean, 56 lines), a computation; luG reduces to it with the private diag_gsp. ShapeQ.lean, DiagGsp.lean and ConjFlag.lean are about 190 lines.

\begin{lemma}\label{sha-siegel-step-restriction}
Let $D$ be a Siegel datum, let $V$ be an ambient ring with residue field $k_V$ (Definition~\ref{sha-siegel-tower-ambient}), and let $T_0\le T$ be towers for $D$ with $T_0$ fit (Definition~\ref{sha-siegel-tower-towers}), with relabelling $\epsilon$, and with $F$ the totally real field of $T$.
Let $P=(r_P,n)$ be a rational flagged pair for $(\epsilon,g_0^{T_0})$ whose image $c_0$ over $k_V$ is of lattice shape (Definition~\ref{sha-siegel-tower-pairs}).
Then there is a flagged pair $P'=(r',(n'_v))$ for the problem of $T$ along $\mathcal O\to V$, with entries in a complete discrete valuation subring of $V$ with finite residue field through which $r'$ factors continuously, such that $r'$ is the restriction of $r_P$ to $G_F$, the image of $P'$ over $k_V$ is residual and over $c_0$, and $\psi'_{v,i}(s)=\delta_i(s)$ for $v\in\Sigma_T$ and $s\in I_v(F)$.
\end{lemma}
\leanref{Development/AutomorphicArithmetic/AutomorphicVarieties/WeakShimuraTheory/CompactIntegralFields/RestrQ.lean}{UpperTheory.V.Ws.restrQ}
\begin{proof}
By Lemma~\ref{sha-siegel-step-towers}(2), $T$ satisfies (a) to (m). Let $A\subset V$ be the complete discrete valuation ring that contains the entries of $P$ and through which $r_P$ factors continuously, and let $W_{\mathrm r}$ and $W_{\mathrm n}$ be as in Lemma~\ref{sha-siegel-step-restriction-matrices}. Bars denote images over $k_V$. As $c_0=(\bar r_P,\bar n)$ is residual and of lattice shape, $\bar r_P$ is supported in $W_{\mathrm r}$ with diagonal entries the residues of the $\chi_i$, and $\bar n$ is supported in $W_{\mathrm n}$.
\begin{enumerate}
\item[(1)] \emph{Moving the places.} For $v\in\Sigma_T$ choose $g_v\in G$ with $g_vv_\ell=v$, which exists since $G$ acts transitively on the places above $\ell$. Then $g_v^{-1}sg_v\in D_\ell$ for $s\in D_v(F)$, and $g_v^{-1}sg_v\in I_\ell$ for $s\in I_v(F)$.
\item[(2)] \emph{The flags.} The matrix $r_P(g_v)$ is symplectic with multiplier $\theta(g_v)$, a unit. By Lemma~\ref{sha-siegel-step-restriction-matrices}(1) over $k_V$, the entry $(1,1)$ of $r_P(g_v)n$ and the determinant of its block in rows and columns $1$ and $2$ have nonzero images in $k_V$; as $\mathfrak m_A\subseteq\mathfrak m_V$, they are units of $A$. By (2) of that lemma over $A$, there is a lower unipotent $n'_v\in M_4(A)$ with $n_v'^*r_P(g_v)n$ upper triangular, and by its (1) over $k_V$, $\bar n'_v$ is supported in $W_{\mathrm n}$.
\item[(3)] \emph{The diagonal characters.} Let $r'$ be the restriction of $r_P$ to $G_F$ and $P'=(r',(n'_v))$. Let $s\in D_v(F)$ and $s'=g_v^{-1}sg_v\in D_\ell$. Then $r_P(s)r_P(g_v)=r_P(g_v)r_P(s')$, and $n^*r_P(s')n$ is upper triangular with diagonal entries $\psi_{v_\ell,i}(s')$, as $P$ is a flagged pair. By Lemma~\ref{sha-siegel-step-restriction-matrices}(3), $n_v'^*r'(s)n'_v$ is upper triangular with diagonal entries $\psi'_{v,i}(s)=\psi_{v_\ell,i}(s')$. If $s\in I_v(F)$, then $s'\in I_\ell$, so $\psi'_{v,i}(s)=\delta_i(s')=\delta_i(s)$, since $P$ is rational and $\delta_i$ is a character of $G$.
\item[(4)] \emph{A flagged pair.} Each $r'(g)$ is symplectic with multiplier $\theta(g)$, the $n'_v$ are lower unipotent, and (3) gives the flags. The matrix $r'(g_0^T)=r_P(g_0^{T_0})^{\ell^{j'}}$ is diagonal, where $T_0\le T$ with exponent $j'$. Let $g\in I_w(F)$ with $w$ not above $\ell$ and $\bar r_0(g)=1$. If $w$ is not above a prime of $S_b$, then $g$ lies in the set $N_\Phi$ of the rational problem, so $r'(g)=1$. If $w$ lies above $p\in S_b$, then $\psi(g)=1$ by (i), and for $h\in G$ with $hv_p=w$ the element $h^{-1}gh$ lies in $I_p\cap\ker\psi$; so $r_P(h^{-1}gh)=1$, and $r'(g)=1$. Finally let $s\in\mu_v$, so that $s^{\ell^a}$ lies in the closure of the commutator subgroup of $D_v(F)$ for some $a$. The quotient of $D_v(F)$ by $I_v(F)$ is abelian and torsion free, being isomorphic to the Galois group of an extension of finite fields of infinite degree; as $I_v(F)$ is closed, $s\in I_v(F)$, and $\psi'_{v,i}(s)=\delta_i(s)$ by (3). Every continuous character $\gamma$ of $G$ has $\gamma(s)^{\ell^a}=1$. So $\chi_{\mathrm{cyc}}(s)=1$, as $\mathbf Z_\ell^\times$ has no element of order $\ell$; $\psi(s)=1$, as $\psi(s)^m=1$ for the $m$ prime to $\ell$ of (j); $\xi(s)=1$, as $\xi$ has order prime to $\ell$; and $\omega(s)=1$, as $\omega(s)$ is an $(\ell-1)$-st root of unity. Hence $\delta_i(s)=1$. The entries of $P'$ lie in $A$, and $r'$ factors continuously through $A$.
\item[(5)] \emph{The image over $k_V$.} The homomorphism $\bar r'$ is the restriction of $\bar r_P$, so it has locally constant entries, diagonal entries the residues of the $\chi_i$, and support in $W_{\mathrm r}$; the $\bar n'_v$ are supported in $W_{\mathrm n}$ by (2); and for $s\in D_v(F)$, $\bar\psi'_{v,i}(s)=\bar\psi_{v_\ell,i}(g_v^{-1}sg_v)=\bar\chi_i(g_v^{-1}sg_v)=\bar\chi_i(s)$, by (3) and since $c_0$ is residual. So the image of $P'$ is residual. It is over $c_0$, since $g_v$ carries $v_\ell$ to $v$ and $\bar n_v'^*\bar r_P(g_v)\bar n$ is upper triangular. The identity $\psi'_{v,i}(s)=\delta_i(s)$ for $s\in I_v(F)$ is (3).
\end{enumerate}
% restrQ (RestrQ.lean): Rq.placeMove for g_v, the private low_desc (through Rq.shapeQ and Rq.luG, an LU decomposition whose pivots are units by the lattice shape), Rq.conjFlag for the flag at v, Rq.killQ for N_Phi, Rq.muPs and the private muIne (Rq.powIne, B1.decComm) for mu_v, leFit for Ok of T, in WeakShimuraTheory/RootsOfUnityShapes. Size: about 1,100 lines of Lean with its dependencies there.
\end{proof}

\subsection{Seeds and Hilbert Eisenstein series}\label{subsection-sha-hilbert-eisenstein}

\begin{lemma}\label{sha-siegel-step-seeds-final}
Let $D$ be a Siegel datum, $T$ a tower for $D$ with field $F$ and relabelling characters $\chi_i$ (Definition~\ref{sha-siegel-tower-towers}), and $\iota:\bar{\mathbf Q}_\ell\cong\mathbf C$; let $\mathfrak m$ be the maximal ideal of the valuation ring of $\bar{\mathbf Q}_\ell$, and call $x\in\bar{\mathbf Q}_\ell$ integral if $|x|\le1$. Let $r:G_F\to M_4(\bar{\mathbf Q}_\ell)$ be a homomorphism and $n_v\in M_4(\bar{\mathbf Q}_\ell)$, for $v\in\Sigma_T$, matrices, and for $s\in D_v(F)$ let $\psi_{v,i}(s)$ be the entry $(i,i)$ of $n_v^*r(s)n_v$. Suppose that:
\begin{enumerate}
\item[(a)] $r$ is continuous, $r(g)$ is symplectic with multiplier $\theta(g)$, $r(g)=1$ for $g$ in the set $N_\Phi$ of the problem of $T$ (Definition~\ref{sha-siegel-tower-pairs}), and $r(g_0^T)$ is diagonal; the entries of $r$ are integral, $r(g)_{ii}-\chi_i(g)\in\mathfrak m$, and the off-diagonal entries of $r(g)$ of weight outside $\{(-1,1),(-1,0),(-1,-1)\}$ lie in $\mathfrak m$; and either $r(g)_{21}$ is a unit for some $g$, or $r(g)_{23}$ and $r(g')_{41}$ are units for some $g$ and $g'$;
\item[(b)] each $n_v$ is lower unipotent (Lemma~\ref{sha-siegel-step-restriction-matrices}) with integral entries, $(n_v)_{32}\in\mathfrak m$ and $(n_v)_{31}$ a unit; $n_v^*r(s)n_v$ is upper triangular for $s\in D_v(F)$; $\psi_{v,i}(s)=1$ for $s\in\mu_v$; and $\psi_{v,i}(s)-\chi_i(s)\in\mathfrak m$ for $s\in D_v(F)$;
\item[(c)] $(r,(n_v))$ is a seed for $T$ (Definition~\ref{sha-siegel-tower-ideal}).
\end{enumerate}
Then there are a ring $\mathcal O_1\supseteq\mathcal O$ of Definition~\ref{sha-character-parts}, with residue field $k_1$, and an integral seed $Q$ for $T$ over $\mathcal O_1$ whose image $(\bar r,(\bar n_v))$ in $k_1$ is of type V, with $\bar r$ having an entry of weight $(-1,0)$ or entries of weights $(-1,1)$ and $(-1,-1)$.
\end{lemma}
\leanref{Development/AutomorphicArithmetic/AutomorphicVarieties/ModularityArgument/ThresholdArguments/DenseThresholdCases/InducedFinalStatement.lean}{UpperTheory.induced_final_statement_alt}
\leanref{Development/AutomorphicArithmetic/AutomorphicVarieties/ModularityArgument/ThresholdArguments/DenseThresholdCases/FinalS.lean}{UpperTheory.V.Sd.S2.Final_s}
\leanref{Development/AutomorphicArithmetic/AutomorphicVarieties/IdeleGenerators/QuadraticSplitting/QuadSplit.lean}{UpperTheory.V.Sd.S2.GlobOk}
\leanref{Development/AutomorphicArithmetic/AutomorphicVarieties/IdeleGenerators/QuadraticSplitting/QuadSplit.lean}{UpperTheory.V.Sd.S2.LocOk}
\leanref{Development/AutomorphicArithmetic/AutomorphicVarieties/BuildingsResidueChars/AutomorphicFormsDescentToSmallerLevel.lean}{UpperTheory.automorphic_forms_descent_to_smaller_level}
\begin{proof}
\begin{enumerate}
\item[(1)] \emph{Coefficients.} As $G_F$ is compact and $r$ is continuous, the entries of $r$ lie in a finite extension $E$ of $\mathbf Q_\ell$, by the Baire category theorem. Let $\mathcal O_1$ be the ring of integers of the compositum of $E$, the fraction field of $\mathcal O$ and the extension of $\mathbf Q_\ell$ generated by the entries of the $n_v$. It is a ring of Definition~\ref{sha-character-parts} containing $\mathcal O$, and it contains the entries of $r$ and of the $n_v$, which are integral. Let $Q=(r,(n_v))$, viewed over $\mathcal O_1$. Its entries are continuous for the $\mathfrak m_{\mathcal O_1}$-adic topology, which is the topology of $\bar{\mathbf Q}_\ell$ on $\mathcal O_1$.
\item[(2)] \emph{Residues.} An element $x$ of $\mathcal O_1$ lies in $\mathfrak m_{\mathcal O_1}$ if and only if $|x|<1$. So elements of $\mathcal O_1$ whose difference lies in $\mathfrak m$ have the same residue, and units of $\mathcal O_1$ have nonzero residues.
\item[(3)] \emph{A flagged pair.} The conditions on a flagged pair for the problem of $T$ along $\mathcal O\subseteq\mathcal O_1$ are identities between entries; they hold in $\bar{\mathbf Q}_\ell$ by (a) and (b), hence in $\mathcal O_1$.
\item[(4)] \emph{The image in $k_1$.} By (a), (b) and (2), the diagonal entries of $\bar r(g)$ are the $\bar\chi_i(g)$, and $\bar\psi_{v,i}(s)=\bar\chi_i(s)$ for $s\in D_v(F)$. The homomorphism $\bar r$ is supported in $\{(-1,1),(-1,0),(-1,-1)\}$, hence in $W_{\mathrm r}$ of Lemma~\ref{sha-siegel-step-restriction-matrices}. The only entry below the diagonal of weight outside $W_{\mathrm n}=\{(-1,0),(0,-1),(-1,-1)\}$ is $(3,2)$, of weight $(1,-1)$, so each $\bar n_v$ is supported in $W_{\mathrm n}$; and its entry $(3,1)$ is nonzero. So the image of $Q$ in $k_1$ is residual and of type V, and $Q$ is an integral seed, its image in $\bar{\mathbf Q}_\ell$ being the seed $(r,(n_v))$.
\item[(5)] \emph{Entries.} The entries $(2,1)$, $(2,3)$ and $(4,1)$ have weights $(-1,0)$, $(-1,1)$ and $(-1,-1)$, and the units among them have nonzero residues by (2).
\end{enumerate}
\end{proof}
% induced_final_statement_alt (InducedFinalStatement.lean, 93 lines): (1) is automorphic_forms_descent_to_smaller_level (37 lines: CompactRepField.exists_finiteDimensional_entries_mem, the Baire category argument, and integrality_of_norm_coefficients for O_1), (2) is hz, hne and heq through Fr.sc_norm, (3) is fin_gsp and fin_conj, (4) is fin_q and fin_v with fin_psi, (5) is the case split on GlobOk.ent. The hypotheses (a), (b) and (c) are GlobOk, LocOk (with LowQ for the n_v) and IsSeedAlt. About 300 lines with the descent and integrality lemmas.

\begin{definition}\label{sha-siegel-step-seeds-induced-datum}
Let $D$ be a Siegel datum (Definition~\ref{sha-siegel-datum}) and $T$ a deep tower for $D$ that satisfies the conditions (a) to (r) of Definition~\ref{sha-siegel-tower-towers}(4), with fields $F$, $K$ and $L$, element $g_0^T$ and relabelling characters $\chi_1,\dots,\chi_4$, and let $\iota:\bar{\mathbf Q}_\ell\cong\mathbf C$. Characters with values in $\mathcal O^\times$ are viewed with values in $\bar{\mathbf Q}_\ell^\times$. For $x$ and $y$ in the valuation ring of $\bar{\mathbf Q}_\ell$, write $x\equiv y$ if $|x-y|<1$. For a subfield $E\subseteq\bar{\mathbf Q}$ and a place $w$ of $\bar{\mathbf Q}$, $G_E$, $D_w(E)$ and $I_w(E)$ are as in Definition~\ref{sha-siegel-tower-towers}.
\begin{enumerate}
\item Let $E$ be a number field, $\rho:G_E\to GL_2(\bar{\mathbf Q}_\ell)$ a homomorphism, $w$ a place of $\bar{\mathbf Q}$ above $\ell$, $(m_1,m_2)\in\mathbf Z^2$ and $\psi_1,\psi_2$ characters of $G_E$. Then $\rho$ is \emph{ordinary at $w$ of weights $(m_1,m_2)$ with residues $(\psi_1,\psi_2)$} if there are $P\in GL_2(\bar{\mathbf Q}_\ell)$ and characters $\kappa_1,\kappa_2$ of $D_w(E)$ such that, for $s\in D_w(E)$, $P\rho(s)P^{-1}$ is upper triangular with diagonal entries $\kappa_1(s),\kappa_2(s)$, $\kappa_i(s)\equiv\psi_i(s)$, and $\kappa_i(s)=\chi_{\mathrm{cyc}}(s)^{m_i}$ if $s\in I_w(E)$.
\item For $\gamma\in\bar{\mathbf Q}_\ell^\times$ and $A,B\in M_2(\bar{\mathbf Q}_\ell)$ let
\[
M(\gamma,A,B)=\frac12\begin{pmatrix}
A_{22}+B_{22}&A_{21}/\gamma-B_{21}&B_{22}-A_{22}&A_{21}/\gamma+B_{21}\\
\gamma A_{12}-B_{12}&A_{11}+B_{11}&-\gamma A_{12}-B_{12}&A_{11}-B_{11}\\
B_{22}-A_{22}&-A_{21}/\gamma-B_{21}&A_{22}+B_{22}&-A_{21}/\gamma+B_{21}\\
\gamma A_{12}+B_{12}&A_{11}-B_{11}&-\gamma A_{12}+B_{12}&A_{11}+B_{11}
\end{pmatrix}.
\]
\item An \emph{induction datum} for $T$ and $\iota$ consists of a real quadratic field $k_0\subseteq\bar{\mathbf Q}$, a character $\varphi:G\to\bar{\mathbf Q}_\ell^\times$, integers $j$ and $t$, homomorphisms $\rho:G_{k_0}\to GL_2(\bar{\mathbf Q}_\ell)$ and $R:G\to GL_4(\bar{\mathbf Q}_\ell)$, elements $\alpha,\beta\in\bar{\mathbf Q}_\ell^\times$ and a character $\delta:G_F\to\bar{\mathbf Q}_\ell^\times$, such that the following hold, with $\mu=\chi_{\mathrm{cyc}}^j\xi\varphi$:
(a) $D_w(\mathbf Q)\subseteq G_{k_0}$ for every place $w$ above $2$ or $\ell$; for $g\in G_F$, $\bar\chi^{(0,1)}(g)=1$ if and only if $g\in G_{k_0}$; and $\varphi=1$ on $G_{k_0}$ and $\varphi=-1$ off $G_{k_0}$;
(b) $j\ge2$, $t\ge2$ and $\ell-1$ divides $j$; $\rho$ is continuous and irreducible on every open subgroup of $G_{k_0}$; $\rho$ is attached through $\iota$ to a cuspidal automorphic representation of $GL_2(\mathbf A_{k_0})$ that has the infinitesimal character of weight $(j+t,-t)$ at one infinite place of $k_0$ and of weight $(j,0)$ at the other; $\det\rho=\chi_{\mathrm{cyc}}^j\xi$ on $G_{k_0}$; for $s\in G_{k_0}$ the coefficients of the characteristic polynomial of $\rho(s)$ are congruent to those of $(X-1)(X-\xi(s))$; for every place $w$ not above $\ell$, $\rho$ is conjugate on $I_w(k_0)$ to the diagonal representation with entries $1$ and $\xi$; and for places $w,w'$ above $\ell$ with $w'=g^{-1}w$ for some $g\notin G_{k_0}$, either $\rho$ is ordinary at $w$ of weights $(j+t,-t)$ with residues $(\xi,1)$ and at $w'$ of weights $(j,0)$ with residues $(1,\xi)$, or the same holds with $w$ and $w'$ exchanged;
(c) for $g\in G_{k_0}\cap G_F$, $\rho(g)_{11}\equiv1$, $\rho(g)_{22}\equiv\xi(g)$, $|\rho(g)_{21}|<1$ and $|\rho(g)_{12}|\le1$; and $|\rho(g)_{12}|=1$ for some such $g$;
(d) $g_0^T\notin G_{k_0}$; $\rho((g_0^T)^2)=\operatorname{diag}(\alpha^2,\beta^2)$; $R(h)=M\big(\beta/\alpha,\rho(h),\rho((g_0^T)^{-1}hg_0^T)\big)$ for $h\in G_{k_0}$, and $R(g_0^T)=\operatorname{diag}(\beta,-\alpha,-\beta,\alpha)$; $R$ is continuous, and $R(g)$ is symplectic with multiplier $\mu(g)$ for $g\in G$;
(e) $\alpha\equiv\chi_4(g_0^T)$ and $\beta\equiv\chi_1(g_0^T)$;
(f) $\delta$ is continuous, and $\delta(g)^2\mu(g)=\theta(g)$ and $\delta(g)\equiv1$ for $g\in G_F$.
\item The \emph{twisted restriction} of an induction datum is the homomorphism $r_\delta:G_F\to M_4(\bar{\mathbf Q}_\ell)$, $r_\delta(g)=\delta(g)R(g)$.
\end{enumerate}
\end{definition}
\leanref{Development/AutomorphicArithmetic/AutomorphicVarieties/IdeleGenerators/QuadraticSplitting/QuadSplit.lean}{UpperTheory.V.Sd.S2.IndDat}
\leanref{Development/AutomorphicArithmetic/AutomorphicVarieties/IdeleGenerators/QuadraticSplitting/QuadSplit.lean}{UpperTheory.V.Sd.S2.Hyp}
\leanref{Development/AutomorphicArithmetic/AutomorphicVarieties/IdeleGenerators/QuadraticSplitting/QuadSplit.lean}{UpperTheory.V.Sd.S2.MixDat}
\leanref{Development/AutomorphicArithmetic/AutomorphicVarieties/IdeleGenerators/QuadraticSplitting/QuadSplit.lean}{UpperTheory.V.Sd.S2.Rib}
\leanref{Development/AutomorphicArithmetic/AutomorphicVarieties/IdeleGenerators/QuadraticSplitting/QuadSplit.lean}{UpperTheory.V.Sd.S2.IsMix}
\leanref{Development/AutomorphicArithmetic/AutomorphicVarieties/IdeleGenerators/QuadraticSplitting/QuadSplit.lean}{UpperTheory.V.Sd.S2.mix}
\leanref{Development/AutomorphicArithmetic/AutomorphicVarieties/IdeleGenerators/QuadraticSplitting/QuadSplit.lean}{UpperTheory.V.Sd.S2.twF}
\leanref{Development/AutomorphicArithmetic/AutomorphicVarieties/IdeleGenerators/SlopeTestRemainingChoices.lean}{UpperTheory.V.Sd.OrdAt}

\begin{lemma}\label{sha-siegel-step-seeds-induced-characters}
Let $D$, $T$ and $\iota$ be as in Definition~\ref{sha-siegel-step-seeds-induced-datum}, with its notation, and let $k_0$ and $\varphi$ satisfy (a) there. Then:
\begin{enumerate}
\item the characters of $G$ with values in $\mathcal O^\times$ take values of absolute value $1$, and $\chi_i\chi_{5-i}=\theta$ for each $i$;
\item for $g\in G_F\cap G_{k_0}$, $\chi_1(g)\equiv\chi_3(g)\equiv\xi(g)$ and $\chi_2(g)\equiv\chi_4(g)\equiv1$;
\item $g_0^T\notin G_{k_0}$, $\xi((g_0^T)^2)\not\equiv1$, $\chi_4(g_0^T)\equiv\pm1$, $\chi_2(g_0^T)\equiv-\chi_4(g_0^T)$ and $\chi_3(g_0^T)\equiv-\chi_1(g_0^T)$;
\item for every $j\ge0$, the character $\theta/(\chi_{\mathrm{cyc}}^j\xi\varphi)$ of $G$ is continuous, and if $\ell-1$ divides $j$, then $\theta(g)\equiv\chi_{\mathrm{cyc}}(g)^j\xi(g)\varphi(g)$ for $g\in G_F$;
\item for $g\in G_F$ and $n\ge0$, $\bar\chi^{(1,0)}(g)^n=1$ if and only if $\xi(g)^n\equiv1$;
\item for $g\in G$, $\xi(g)\equiv1$ implies $\xi(g)=1$, and $\omega(g)=1$ if and only if $\chi_{\mathrm{cyc}}(g)\equiv1$; the kernel of $\omega$ is open.
\end{enumerate}
\end{lemma}
\leanref{Development/AutomorphicArithmetic/AutomorphicVarieties/BuildingsResidueChars/GaloisCharacterResidualDecomposition.lean}{UpperTheory.galois_character_residual_decomposition}
\leanref{Development/AutomorphicArithmetic/AutomorphicVarieties/IdeleGenerators/QuadraticSplitting/QuadSplit.lean}{UpperTheory.V.Sd.S2.CharDat}
\begin{proof}
Bars denote residues. As $\bar\omega=\bar\chi_{\mathrm{cyc}}$, $\bar\alpha=\bar\xi$, and the residues of $(\chi_1,\dots,\chi_4)$ are $(\bar\xi\bar\Psi,\bar\Psi,\bar\xi,1)$ if $\epsilon=0$ and $(\bar\xi,1,\bar\xi\bar\Psi,\bar\Psi)$ if $\epsilon=1$. In both cases $\bar\chi^{(1,0)}=\bar\alpha=\bar\xi$, and $\bar\chi^{(0,1)}(g)^n=1$ if and only if $\bar\Psi(g)^n=1$. For $g\in G_F$, $\bar\Psi(g)^2=1$ by (n), and $\bar\Psi(g)=1$ if and only if $g\in G_{k_0}$, by (a); so $\bar\Psi(g)$ is the residue of $\varphi(g)$.

(1) The characters take values in $\mathcal O^\times$. The products $\chi_i\chi_{5-i}$ are $\delta_1\delta_4=\theta$ and $\delta_2\delta_3=\Psi\alpha=\theta$.

(2) For $g\in G_F\cap G_{k_0}$, $\bar\Psi(g)=1$; read off the residues.

(3) If $g_0^T\in G_{k_0}$, then $\bar\Psi(g_0^T)=1$ and $\bar\chi_2(g_0^T)=\bar\chi_4(g_0^T)$, against (l). So $\bar\Psi(g_0^T)=-1$, and with $x=\bar\xi(g_0^T)$ the residues of the $\chi_i(g_0^T)$ are $(-x,-1,x,1)$ or $(x,1,-x,-1)$. This gives the last three claims. If $\xi((g_0^T)^2)\equiv1$, then $x=\pm1$, and $\bar\chi_1(g_0^T)$ equals $\bar\chi_2(g_0^T)$ or $\bar\chi_4(g_0^T)$, against (l).

(4) The characters $\chi_{\mathrm{cyc}}$, $\omega$, $\xi$, $\psi$ and $\varphi$ are continuous, $\varphi$ because its kernel $G_{k_0}$ is open, and $\theta=\chi_{\mathrm{cyc}}\omega^{-1}\xi\chi_{\mathrm{cyc}}^k\psi$, with $k$ the integer of $D$. If $\ell-1$ divides $j$, then $\bar\chi_{\mathrm{cyc}}^j=\bar\omega^j=1$, as $\omega^{\ell-1}=1$; and for $g\in G_F$, $\bar\theta(g)=\bar\alpha(g)\bar\Psi(g)=\bar\xi(g)\bar\varphi(g)$.

(5) $\bar\chi^{(1,0)}=\bar\xi$.

(6) The values of $\xi$ are $m$-th roots of unity for some $m$ prime to $\ell$, and those of $\omega$ are $(\ell-1)$-st roots of unity; distinct such roots of unity have distinct residues. As $\bar\omega=\bar\chi_{\mathrm{cyc}}$, the kernel of $\omega$ is the set of $g$ with $\chi_{\mathrm{cyc}}(g)\equiv1$, which is open.
\end{proof}
% galois_character_residual_decomposition (GaloisCharacterResidualDecomposition.lean, about 115 lines): b_chi for the residues of the chi_i (chi_f and induced_characters_values), rb_ps and rb_al for the residues of chi^w, hin and hout for Psi bar on G_F (Hyp.ps and Tower.Two.sq), Tower.Ok.dist for (l), teichmuller_character_spec for (6), cont_of_ker for the continuity.

\begin{definition}\label{sha-siegel-step-seeds-induced-hilbert-form-ordinary}
Let $\ell$ be a prime, $\iota:\bar{\mathbf Q}_\ell\cong\mathbf C$ and $E\subseteq\bar{\mathbf Q}$ a totally real number field. Write $G_E=\operatorname{Gal}(\bar{\mathbf Q}/E)$, and $D_A(E)\supseteq I_A(E)$ for the decomposition and inertia groups of $G_E$ at a place $A$ of $\bar{\mathbf Q}$. For a prime $w$ of $E$, a uniformizer $\varpi_w$ of $E_w$ and $c\ge0$, let $K_1(w^c)$ be the group of $g\in GL_2(\mathcal O_{E_w})$ whose bottom row is congruent to $(0,1)$ modulo $w^c$, and let $U_w$ be the operator of the double coset $K_1(w^c)\operatorname{diag}(\varpi_w,1)K_1(w^c)$ on the vectors fixed by $K_1(w^c)$ in a smooth representation of $GL_2(E_w)$. Let $\xi:G_E\to\bar{\mathbf Q}_\ell^\times$ be a homomorphism and $s\in\mathbf Z$.
An \emph{ordinary form datum} for $\xi$ of weight $s$ consists of a function $c$ from the primes of $E$ to $\mathbf Z_{\ge0}$, a pair $m(w)=(m_0(w),m_1(w))\in\mathbf Z^2$ and an element $a(w)\in\bar{\mathbf Q}_\ell$ for each prime $w$ of $E$ above $\ell$, a cuspidal automorphic representation $\pi$ of $GL_2(\mathbf A_E)$ and a homomorphism $r:G_E\to GL_2(\bar{\mathbf Q}_\ell)$, such that:
\begin{enumerate}
\item[(a)] $\pi$ is regular algebraic and polarized: for some Hecke character $\chi$ of $E$ with the same sign at all real places, there is a nondegenerate pairing $B$ on $\pi$ with $B(g\phi,g\phi')=\chi(\det g)B(\phi,\phi')$ for $g$ in $GL_2(\mathbf A_{E,f})$ and in the maximal compact subgroup at infinity;
\item[(b)] $r$ is continuous and attached to $\pi$ through $\iota$, and $\det r=\chi_{\mathrm{cyc}}^s\xi$;
\item[(c)] $\xi$ has open kernel, $c(w)=0$ for all but finitely many $w$, and for every prime $w$ and place $A$ of $\bar{\mathbf Q}$ above $w$, $I_A(E)\subseteq\ker\xi$ if and only if $c(w)=0$ or $w$ lies above $\ell$;
\item[(d)] if $c(w)=0$, then $\pi_w$ has a nonzero vector fixed by $GL_2(\mathcal O_{E_w})$;
\item[(e)] if $c(w)\ne0$ and $w$ does not lie above $\ell$, then $\pi_w$ has a nonzero vector fixed by $K_1(w^{c(w)})$ on which $U_w$ acts by some $\lambda$ with $|\iota^{-1}(\lambda)|=1$;
\item[(f)] if $w$ lies above $\ell$, then $c(w)=1$, $m_1(w)+2\le m_0(w)$, $m_0(w)+m_1(w)=s$ and $|a(w)|=1$; $\pi_w$ has a nonzero vector fixed by $K_1(w)$ on which $U_w$ acts by some $\lambda$ with $|\iota^{-1}(\lambda)\operatorname{N}(w)^{m_0(w)}-a(w)|<1$; and for every embedding $\tau:E\to\bar{\mathbf Q}_\ell$ that induces $w$, $\pi$ has the infinitesimal character of weight $m(w)$ at the infinite place of $E$ induced by $\iota\circ\tau$.
\end{enumerate}
\end{definition}
\leanref{Development/AutomorphicArithmetic/AutomorphicVarieties/IdeleGenerators/IsCondAt.lean}{UpperTheory.V.Sd.OrdForm}
\leanref{Development/AutomorphicArithmetic/AutomorphicVarieties/IdeleGenerators/IsCondAt.lean}{UpperTheory.V.Sd.K1}
\leanref{Development/AutomorphicArithmetic/AutomorphicVarieties/IdeleGenerators/IsCondAt.lean}{UpperTheory.V.Sd.Uop}

\begin{definition}\label{sha-siegel-step-seeds-induced-hilbert-form-exists-holomorphic}
Let $E$ be a totally real number field, $\kappa$ a function from the infinite places of $E$ to $\mathbf Z_{\ge0}$, $c$ a function from the primes of $E$ to $\mathbf Z_{\ge0}$, and $\omega$ a Hecke character of $E$. For a prime $w$ of $E$, let $K_1(w^{c(w)})$ and $U_w$ be the group and the operator of Definition~\ref{sha-siegel-step-seeds-induced-hilbert-form-ordinary}, acting on functions on $GL_2(\mathbf A_E)$ by right translation at $w$. For an infinite place $v$ of $E$ and $a,b\in\{1,2\}$, let $X^v_{ab}$ be the derivative of right translation at $v$ in the direction of the matrix unit $e_{ab}$: $X^v_{ab}\phi(g)$ is the derivative at $x=0$ of $\phi(g\,(1+xe_{ab})_v)$, for $x\in\mathbf R$. For $s\in\mathbf C$, the \emph{Capelli operator} at $v$ is $C_v(s)=(s+X^v_{11})(s-1+X^v_{22})-X^v_{21}X^v_{12}$, acting on the functions on $GL_2(\mathbf A_E)$ that are smooth in the variable at infinity.
A \emph{holomorphic form} of weight $\kappa$, level $c$ and central character $\omega$ is a function $\phi$ on $GL_2(\mathbf A_E)$, smooth in the variable at infinity, such that:
\begin{enumerate}
\item[(a)] $\phi$ is an automorphic form: $\phi(\gamma g)=\phi(g)$ for $\gamma\in GL_2(E)$; $\phi$ is fixed by right translation by an open subgroup of $GL_2(\mathbf A_{E,f})$; its right translates by the maximal compact subgroup of $GL_2(E\otimes\mathbf R)$ span a finite-dimensional space, and so do its images under the algebra of operators generated by the Capelli operators $C_v(s)$, for the infinite places $v$ and $s\in\mathbf C$; and $\phi$ has moderate growth;
\item[(b)] $\phi$ is fixed by $K_1(w^{c(w)})$ for every prime $w$;
\item[(c)] $\phi(zg)=\omega(z)\phi(g)$ for $z\in\mathbf A_E^\times$ and $g\in GL_2(\mathbf A_E)$;
\item[(d)] for every infinite place $v$, $(X^v_{12}-X^v_{21})\phi=i\kappa(v)\phi$ and $(X^v_{11}-X^v_{22})\phi=i(X^v_{12}+X^v_{21})\phi$.
\end{enumerate}
It is a \emph{cusp form} if moreover $\int_{E\backslash\mathbf A_E}\phi\bigl(\begin{pmatrix}1&x\\0&1\end{pmatrix}g\bigr)\,dx=0$ for every $g\in GL_2(\mathbf A_E)$. For an integer $k\ge0$, a holomorphic form has \emph{parallel weight} $k$ if $\kappa(v)=k$ for every $v$. A \emph{holomorphic eigenform} of parallel weight $k$, level $c$ and central character $\omega$, with \emph{eigenvalues} $\lambda(w)\in\mathbf C$, is a nonzero holomorphic form $\phi$ of parallel weight $k$, level $c$ and central character $\omega$ with $\operatorname N(w)^{k/2-1}U_w\phi=\lambda(w)\phi$ for every prime $w$.
\end{definition}
\leanref{Development/AutomorphicArithmetic/AutomorphicVarieties/IdeleGenerators/IsCondAt.lean}{UpperTheory.V.Sd.IsHol}
\leanref{PositivityTheory/AutomorphicForms/PublishedInputs/JacquetLanglands/TorusKernels/IsHolVOf.lean}{UpperTheory.Esk.IsHolV}
\leanref{Development/AutomorphicArithmetic/AutomorphicVarieties/IdeleGenerators/IsCondAt.lean}{UpperTheory.V.Sd.IsEigHol}
\leanref{Development/AutomorphicArithmetic/AutomorphicVarieties/IdeleGenerators/IsCondAt.lean}{UpperTheory.V.Sd.hk}
\leanref{Development/AutomorphicArithmetic/AutomorphicVarieties/AutomorphicForms/RealCusp.lean}{UpperTheory.V.TR.IsCuspForm}

\begin{definition}\label{sha-siegel-step-seeds-induced-hilbert-form-exists-eisenstein-expansion}
Let $E$ be a totally real number field with different $\mathfrak d$, and fix a uniformizer $\varpi_w$ of $E_w$ at every prime $w$ of $E$. For a Hecke character $\eta$ of $E$ write $\eta(\varpi_w)$ for its value at the idele with component $\varpi_w$ at $w$ and $1$ elsewhere. Let $k\ge0$, let $c_0$ be a function on $\mathbf A_{E,f}^\times$ and $a$ a function on the nonzero integral ideals of $E$, and write $\mathcal O_E$ for the unit ideal.
\begin{enumerate}
\item For $y\in\mathbf A_{E,f}^\times$, let $a(y)=a(\mathfrak m)$ if $y$ generates the fractional ideal $\mathfrak m\mathfrak d^{-1}$ for a nonzero integral ideal $\mathfrak m$, that is, $\operatorname{ord}_w(y_w)=\operatorname{ord}_w(\mathfrak m)-\operatorname{ord}_w(\mathfrak d)$ at every prime $w$; let $a(y)=0$ if there is no such $\mathfrak m$.
\item A function $\phi$ on $GL_2(\mathbf A_E)$ has the \emph{Fourier expansion of weight $k$ with constant term $c_0$ and coefficients $a$} if for all $x\in E\otimes\mathbf R$, all totally positive $y\in E\otimes\mathbf R$ and all $y_f\in\mathbf A_{E,f}^\times$, the family of the numbers $a(\xi y_f)\exp\bigl(2\pi i\sum_vv(\xi)(x_v+iy_v)\bigr)$, indexed by the totally positive $\xi\in E$, with $v$ running over the infinite places, is summable with sum
\[
\phi\Bigl(\begin{pmatrix}y&x\\0&1\end{pmatrix}_\infty\begin{pmatrix}y_f&0\\0&1\end{pmatrix}_f\Bigr)\bigl(\operatorname N(y)\,|y_f|\bigr)^{-k/2}-c_0(y_f),
\]
where $\operatorname N(y)=\prod_vy_v$ and $|y_f|$ is the product of the normalized absolute values $|y_w|_w$ over the primes $w$.
\item Let $c$ be a function from the primes of $E$ to $\mathbf Z_{\ge0}$ and $\omega$ a Hecke character of $E$. Let $\chi_{c,\omega}$ be the multiplicative function on the nonzero integral ideals with $\chi_{c,\omega}(w)=\omega(\varpi_w)$ if $c(w)=0$ and $\chi_{c,\omega}(w)=0$ if $c(w)\ne0$. For a prime $w$, let $T^k_wa$ be the function $\mathfrak m\mapsto a(w\mathfrak m)+\chi_{c,\omega}(w)\operatorname N(w)^{k-1}a(\mathfrak mw^{-1})$ on the nonzero integral ideals, where $a(\mathfrak mw^{-1})=0$ if $w$ does not divide $\mathfrak m$.
\end{enumerate}
\end{definition}
\leanref{Development/AutomorphicArithmetic/AutomorphicVarieties/IdeleGenerators/IsCondAt.lean}{UpperTheory.V.Sd.HasExp}
\leanref{Development/AutomorphicArithmetic/AutomorphicVarieties/IdeleGenerators/IsCondAt.lean}{UpperTheory.V.Sd.coefAt}
\leanref{Development/AutomorphicArithmetic/AutomorphicVarieties/IdeleGenerators/IsCondAt.lean}{UpperTheory.V.Sd.GenD}
\leanref{Development/AutomorphicArithmetic/AutomorphicVarieties/IdeleGenerators/IsCondAt.lean}{UpperTheory.V.Sd.coefDiv}
\leanref{Development/AutomorphicArithmetic/AutomorphicVarieties/IdeleGenerators/IsCondAt.lean}{UpperTheory.V.Sd.chiI}
\leanref{Development/AutomorphicArithmetic/AutomorphicVarieties/IdeleGenerators/IsCondAt.lean}{UpperTheory.V.Sd.hkCoef}

\begin{lemma}\label{sha-siegel-step-seeds-induced-hilbert-form-exists-eisenstein-series}
Let $E$ be a totally real number field and $\omega$ a Hecke character of $E$ of finite order with $\omega_v(-1)=-1$ at every real place $v$. For each prime $w$ let $c(w)$ be the conductor exponent of $\omega$ at $w$, the least $n\ge0$ such that $\omega_w$ is trivial on the units of $E_w$ congruent to $1$ modulo $w^n$, and suppose that $c(w)\ne0$ for some $w$. Then there are a holomorphic form $\phi$ of parallel weight $1$, level $c$ and central character $\omega$ (Definition~\ref{sha-siegel-step-seeds-induced-hilbert-form-exists-holomorphic}) and a number $L\in\mathbf C$ such that $\phi$ has the Fourier expansion of weight $1$ with constant term the constant function $L$ and coefficients $\mathfrak m\mapsto\sum_{\mathfrak r\mid\mathfrak m}\chi_{c,\omega}(\mathfrak r)$ (Definition~\ref{sha-siegel-step-seeds-induced-hilbert-form-exists-eisenstein-expansion}).
\end{lemma}
\leanref{Development/AutomorphicArithmetic/AutomorphicVarieties/PublishedTheorems/Hilbert.lean}{UpperTheory.V.Published.HilbertEisOne}
\leanref{PositivityTheory/AutomorphicForms/PublishedInputs/HilbertEisenstein/HilbertEisensteinWeightOneHolds.lean}{UpperTheory.V.Published.hilbert_eisenstein_weight_one_holds}
\begin{proof}
\openstep{The construction of this Eisenstein series of parallel weight one and the computation of its Fourier expansion are not yet written out.}
% Published.HilbertEisOne is an input of Hypotheses4Z' that the Lean proves as hilbert_eisenstein_weight_one_holds (the directory PublishedInputs/HilbertEisenstein, 63 files): top applied to poisson_summation_holds, s_integrality, moderate_growth_holds, smoothness_holds, left_invariance_holds, level_holds, open_kernel_holds, central_character_holds, K_finiteness_holds, weight_holds, low_weight_holds, center_finiteness_holds, fourier_expansion_holds and scalar_action_holds. It is the Eisenstein series of parallel weight one attached to the pair of characters (1, omega), as in Shimura's work on Hilbert modular Eisenstein series; a citation with a checked locator would also reach the floor. Size: about 8,500 lines with its dependencies, net of the definitions.
\end{proof}

\begin{lemma}\label{sha-siegel-step-seeds-induced-hilbert-form-exists-eisenstein-congruent}
Let $E$ be a totally real number field, $\ell$ an odd prime and $\iota:\bar{\mathbf Q}_\ell\cong\mathbf C$. Then there is $N\ge0$ such that for every integer $k>0$ divisible by $(\ell-1)\ell^N$ there is a holomorphic form $\psi$ of parallel weight $k$, level $0$ and trivial central character (Definition~\ref{sha-siegel-step-seeds-induced-hilbert-form-exists-holomorphic}) with the Fourier expansion of weight $k$ with some constant term $d_0$ and coefficients $b$ (Definition~\ref{sha-siegel-step-seeds-induced-hilbert-form-exists-eisenstein-expansion}) such that $|\iota^{-1}(d_0(y)-1)|<1$ for every $y\in\mathbf A_{E,f}^\times$ and $|\iota^{-1}(b(\mathfrak m))|<1$ for every nonzero integral ideal $\mathfrak m$.
\end{lemma}
\leanref{Development/AutomorphicArithmetic/AutomorphicVarieties/PublishedTheorems/Hilbert.lean}{UpperTheory.V.Published.HilbertOne}
\leanref{Development/AutomorphicArithmetic/AutomorphicVarieties/GaussSumFiltrations/ExistsHilbertModularFormCongrOne.lean}{UpperTheory.V.Published.exists_hilbertModularForm_congr_one}
\begin{proof}
\openstep{The construction of a form of parallel weight $k$ whose Fourier expansion is congruent to $1$, for every $k$ divisible by $(\ell-1)\ell^N$, is not yet written out.}
% Published.HilbertOne is proved in the Lean as exists_hilbertModularForm_congr_one, that is GlobalTheory.HilbertThetaSeries.hilbertOne, from theta series. Lit.OneUse, the input of Hypotheses4Z', is Lit.oneUse = oneUse_of exists_hilbertModularForm_congr_one; oneUse_of takes k = (l - 1) l^N (B + 1), as in step (2) of the proof of Lemma sha-siegel-step-seeds-induced-hilbert-form-exists-eisenstein. Size: about 9,000 lines with its dependencies, net of the definitions.
\end{proof}

\begin{lemma}\label{sha-siegel-step-seeds-induced-hilbert-form-exists-eisenstein-raise}
Let $E$ be a totally real number field with uniformizers $\varpi_w$, $k\ge0$, $c$ a function from the primes of $E$ to $\mathbf Z_{\ge0}$, $\omega$ a Hecke character of $E$, and $\phi$ a holomorphic form of parallel weight $k$, level $c$ and central character $\omega$ with the Fourier expansion of weight $k$ with constant term $c_0$ and coefficients $a$ (Definition~\ref{sha-siegel-step-seeds-induced-hilbert-form-exists-eisenstein-expansion}). Let $w$ be a prime, $u\in\mathbf A_{E,f}^\times$ the idele with component $\varpi_w^{-1}$ at $w$ and $1$ elsewhere, and $c'$ the function with $c'(w)=c(w)+1$ and $c'=c$ at the other primes. Then $\psi(g)=|u|^{-k/2}\phi\bigl(g\operatorname{diag}(u,1)\bigr)$ is a holomorphic form of parallel weight $k$, level $c'$ and central character $\omega$ with the Fourier expansion of weight $k$ with constant term $y\mapsto c_0(yu)$ and coefficients $\mathfrak m\mapsto a(\mathfrak mw^{-1})$.
\end{lemma}
\leanref{Development/AutomorphicArithmetic/AutomorphicVarieties/IdeleGenerators/WeightChange.lean}{UpperTheory.V.Sd.HolB}
\leanref{Development/AutomorphicArithmetic/AutomorphicVarieties/IdeleGenerators/HolomorphicIrreducibility/MemMirabQuot.lean}{UpperTheory.V.Sd.holB}
\begin{proof}
Right translation by an element of $GL_2(\mathbf A_{E,f})$ maps automorphic forms to automorphic forms: it commutes with left translation by $GL_2(E)$, with right translation by $GL_2(E\otimes\mathbf R)$ and with the derivatives at infinity, it conjugates open subgroups of $GL_2(\mathbf A_{E,f})$ into open subgroups, and it preserves moderate growth. So $\psi$ is an automorphic form; it satisfies condition (c) of Definition~\ref{sha-siegel-step-seeds-induced-hilbert-form-exists-holomorphic}, as left multiplication by scalar matrices commutes with right translation, and condition (d), as the derivatives at infinity commute with right translation by $GL_2(\mathbf A_{E,f})$. For condition (b) at $w$, let $h\in K_1(w^{c(w)+1})$, with entries $h_{ij}$. Then $\operatorname{diag}(\varpi_w,1)\,h\,\operatorname{diag}(\varpi_w,1)^{-1}$ has entries $h_{11}$, $\varpi_wh_{12}$, $\varpi_w^{-1}h_{21}$ and $h_{22}$, so it lies in $K_1(w^{c(w)})$, and $\psi(gh)=\psi(g)$ as $\phi$ is fixed by $K_1(w^{c(w)})$; at the other primes $\operatorname{diag}(u,1)$ has component $1$. For the expansion, apply that of $\phi$ at $y_fu$: as $|y_fu|=|y_f||u|$, the left side becomes that of $\psi$ at $y_f$ minus $c_0(y_fu)$, and $\xi y_fu$ generates $\mathfrak m'\mathfrak d^{-1}$ with $\mathfrak m'$ integral if and only if $\xi y_f$ generates $\mathfrak m'w\mathfrak d^{-1}$; so $a(\xi y_fu)$ is $a(\mathfrak mw^{-1})$ when $\xi y_f$ generates $\mathfrak m\mathfrak d^{-1}$, and $0$ when $\xi y_f$ generates no such ideal.
\end{proof}
% holB (MemMirabQuot.lean, 50 lines with the private exists_idele, aut_fin, zalg_comm and lie_fin): u is exists_idele w varpi^-1, s = |u|^(-k/2); the level is K1_conj at w and locGL_comm elsewhere; cent, wt and low as in the text; the expansion is hE at yf * u with coefAt_shift and absF_mul. About 290 lines with its dependencies.

\begin{lemma}\label{sha-siegel-step-seeds-induced-hilbert-form-exists-eisenstein-stabilized}
Let $E$ be a totally real number field with uniformizers $\varpi_w$, $\ell$ a prime and $\iota:\bar{\mathbf Q}_\ell\cong\mathbf C$. Let $\omega$ be a Hecke character of $E$ of finite order with $\omega_v(-1)=-1$ at every real place $v$, and $c(w)$ its conductor exponent at each prime $w$; suppose that $c(w)\ne0$ for some $w$ and that $c(w)=0$ for every $w$ above $\ell$. Let $P$ be a set of primes of $E$ that contains a prime above $\ell$. Let $c_\ell(w)=1$ for the primes $w$ above $\ell$ and $c_\ell(w)=c(w)$ for the others, and let $e(w)$ be $\omega(\varpi_w)$ if $w$ lies above $\ell$ and $w\in P$, $1$ if $w$ lies above $\ell$ and $w\notin P$, $1+\omega(\varpi_w)$ if $w$ does not lie above $\ell$ and $c(w)=0$, and $1$ otherwise. Then there is a holomorphic form $\phi$ of parallel weight $1$, level $c_\ell$ and central character $\omega$ with the Fourier expansion of weight $1$ with constant term $0$ and coefficients $a$ (Definition~\ref{sha-siegel-step-seeds-induced-hilbert-form-exists-eisenstein-expansion}) such that $a(\mathcal O_E)=1$, and $|\iota^{-1}(a(\mathfrak m))|\le1$ and $T^1_wa(\mathfrak m)=e(w)a(\mathfrak m)$ for every prime $w$ and every nonzero integral ideal $\mathfrak m$, with $T^1_w$ for $c_\ell$ and $\omega$.
\end{lemma}
\leanref{Development/AutomorphicArithmetic/AutomorphicVarieties/IdeleGenerators/WeightChange.lean}{UpperTheory.V.Sd.Emix}
\leanref{Development/AutomorphicArithmetic/AutomorphicVarieties/IdeleGenerators/HolomorphicIrreducibility/NotDvdMul.lean}{UpperTheory.V.Sd.emix}
\begin{proof}
Let $\chi=\chi_{c,\omega}$. The values of $\omega$ are roots of unity, so $|\iota^{-1}(\omega(\varpi_w))|=1$ for every $w$. Sums and scalar multiples of holomorphic forms of parallel weight $1$, level $c'$ and central character $\omega$ are again such forms, the Fourier expansions add, and a form of level $c'$ is one of level $c''$ whenever $c'\le c''$. Let $T$ be the finite set of primes above $\ell$. For $U\subseteq T$ let $c_U(w)=1$ for $w\in U$ and $c_U(w)=c(w)$ for $w\notin U$. We show by induction on $|U|$ that there are a holomorphic form $\phi_U$ of parallel weight $1$, level $c_U$ and central character $\omega$, a number $L_U$ and a function $a_U$ such that:
\begin{enumerate}
\item[(i)] $\phi_U$ has the Fourier expansion of weight $1$ with constant term the constant $L_U$ and coefficients $a_U$, and $L_U=0$ if $U$ meets $P$;
\item[(ii)] $a_U(\mathcal O_E)=1$, and $|\iota^{-1}(a_U(\mathfrak m))|\le1$ for every $\mathfrak m$;
\item[(iii)] $a_U(w\mathfrak m)=e(w)a_U(\mathfrak m)$ for $w\in U$, and $a_U(w\mathfrak m)+\chi(w)a_U(\mathfrak mw^{-1})=(1+\chi(w))a_U(\mathfrak m)$ for $w\notin U$, for every $\mathfrak m$.
\end{enumerate}
For $U=\emptyset$, Lemma~\ref{sha-siegel-step-seeds-induced-hilbert-form-exists-eisenstein-series} gives $\phi_\emptyset$ and $L_\emptyset$, with $a_\emptyset(\mathfrak m)=\sum_{\mathfrak r\mid\mathfrak m}\chi(\mathfrak r)$. Then $a_\emptyset(\mathcal O_E)=\chi(\mathcal O_E)=1$, and $a_\emptyset$ is integral as the values of $\chi$ are. For a prime $w$, the divisors of $w\mathfrak m$ are the $w\mathfrak r$ with $\mathfrak r\mid\mathfrak m$ and the divisors of $\mathfrak m$ not divisible by $w$; and $\sum_{\mathfrak r\mid\mathfrak m,\,w\nmid\mathfrak r}\chi(\mathfrak r)=a_\emptyset(\mathfrak m)-\chi(w)a_\emptyset(\mathfrak mw^{-1})$. As $\chi$ is multiplicative, $a_\emptyset(w\mathfrak m)=\chi(w)a_\emptyset(\mathfrak m)+a_\emptyset(\mathfrak m)-\chi(w)a_\emptyset(\mathfrak mw^{-1})$, which is (iii).

Let $U\subseteq T$ satisfy (i) to (iii), $v\in T\setminus U$ and $U'=U\cup\{v\}$. As $c(v)=0$, $\chi(v)=\omega(\varpi_v)$. Let $y=1$ and $x=\omega(\varpi_v)$ if $v\in P$, and $y=\omega(\varpi_v)$ and $x=1$ if $v\notin P$; so $x=e(v)$, $x+y=1+\chi(v)$ and $xy=\chi(v)$. Let $\psi$ be the form of Lemma~\ref{sha-siegel-step-seeds-induced-hilbert-form-exists-eisenstein-raise} for $\phi_U$ and $v$, of level $c_{U'}$, and let $\phi_{U'}=\phi_U-y\psi$, of level $c_{U'}$. It has the Fourier expansion of weight $1$ with constant term $L_{U'}=L_U-yL_U$ and coefficients $a_{U'}(\mathfrak m)=a_U(\mathfrak m)-ya_U(\mathfrak mv^{-1})$.
\begin{enumerate}
\item[(i)] If $v\in P$, then $y=1$ and $L_{U'}=0$; if $U$ meets $P$, then $L_U=0$ and $L_{U'}=0$.
\item[(ii)] $a_{U'}(\mathcal O_E)=a_U(\mathcal O_E)=1$, as $v$ does not divide $\mathcal O_E$; and $a_{U'}$ is integral, as $y$ is.
\item[(iii)] At $v$, by (iii) for $U$ at $v\notin U$, $a_{U'}(v\mathfrak m)=a_U(v\mathfrak m)-ya_U(\mathfrak m)=(x+y)a_U(\mathfrak m)-xya_U(\mathfrak mv^{-1})-ya_U(\mathfrak m)=x\,a_{U'}(\mathfrak m)=e(v)a_{U'}(\mathfrak m)$. At a prime $w\ne v$, $w\mathfrak m$ is divisible by $v$ if and only if $\mathfrak m$ is, and then $(w\mathfrak m)v^{-1}=w(\mathfrak mv^{-1})$; likewise $\mathfrak mw^{-1}$, when integral, is divisible by $v$ if and only if $\mathfrak m$ is, and then $(\mathfrak mw^{-1})v^{-1}=(\mathfrak mv^{-1})w^{-1}$. So the relation of (iii) at $w$ for $a_U$ holds also for $\mathfrak m\mapsto a_U(\mathfrak mv^{-1})$, and hence, being linear, for $a_{U'}$.
\end{enumerate}
For $U=T$, $c_T=c_\ell$ and $L_T=0$, as $T$ meets $P$. Let $\phi=\phi_T$ and $a=a_T$, and let $\chi_\ell=\chi_{c_\ell,\omega}$. For $w\in T$, $c_\ell(w)=1$, so $\chi_\ell(w)=0$ and $T^1_wa(\mathfrak m)=a(w\mathfrak m)=e(w)a(\mathfrak m)$. For $w\notin T$, $c_\ell(w)=c(w)$ and $\chi_\ell(w)=\chi(w)$, so $T^1_wa(\mathfrak m)=(1+\chi(w))a(\mathfrak m)$; this is $e(w)a(\mathfrak m)$, as $1+\chi(w)=1+\omega(\varpi_w)$ if $c(w)=0$ and $1+\chi(w)=1$ otherwise.
\end{proof}
% emix (NotDvdMul.lean, 30 lines) applies the private step (the induction on the finset T of primes above l, the text's U, with Rel w e t a for a(wm) + e a(m/w) = t a(m)) to Published.HilbertEisOne and intS, the integral elements; the case U empty is rel_sum with sum_step, chi_mul and chi_top; the induction step is holB (Lemma sha-siegel-step-seeds-induced-hilbert-form-exists-eisenstein-raise), holLin for sums and levels, rel_same at v and rel_add with rel_div at the other primes. The text's uniformizers are the Lean's atFin w (unif w). About 600 lines with holB and holLin.

\begin{lemma}\label{sha-siegel-step-seeds-induced-hilbert-form-exists-eisenstein-product-holomorphic}
Let $E$ be a totally real number field, $k,k'\ge0$, $c$ a function from the primes of $E$ to $\mathbf Z_{\ge0}$ and $\omega$ a Hecke character of $E$ of finite order. Let $\phi$ be a holomorphic form of parallel weight $k$, level $c$ and central character $\omega$, and $\psi$ a holomorphic form of parallel weight $k'$, level $c$ and trivial central character (Definition~\ref{sha-siegel-step-seeds-induced-hilbert-form-exists-holomorphic}). Then the product $\phi\psi$ is a holomorphic form of parallel weight $k+k'$, level $c$ and central character $\omega$.
\end{lemma}
\leanref{Development/AutomorphicArithmetic/AutomorphicVarieties/IdeleGenerators/HolomorphicIrreducibility/HolMul.lean}{UpperTheory.V.Sd.hol_mul}
\leanref{Development/AutomorphicArithmetic/AutomorphicVarieties/IdeleGenerators/ScalarStabilizers/Mulsm.lean}{UpperTheory.V.Sd.Hm.mulsm}
\leanref{Development/AutomorphicArithmetic/AutomorphicVarieties/IdeleGenerators/ScalarStabilizers/Trzero.lean}{UpperTheory.V.Sd.Hm.trzero}
\leanref{Development/AutomorphicArithmetic/AutomorphicVarieties/IdeleGenerators/ScalarStabilizers/Zfinof.lean}{UpperTheory.V.Sd.Hm.zfinof}
\leanref{Development/AutomorphicArithmetic/AutomorphicVarieties/IdeleGenerators/ScalarStabilizers/Mulaut.lean}{UpperTheory.V.Sd.Hm.mulaut}
\begin{proof}
Let $\chi=\phi\psi$, and use the derivatives $X^v_{ab}$ and the Capelli operators $C_v(s)$ of Definition~\ref{sha-siegel-step-seeds-induced-hilbert-form-exists-holomorphic}.
\begin{enumerate}
\item[(1)] \emph{Derivatives.} The function $\chi$ is smooth in the variable at infinity, as $\phi$ and $\psi$ are, and every derivative of right translation at infinity satisfies the Leibniz rule $X(\phi\psi)=(X\phi)\psi+\phi(X\psi)$.
\item[(2)] \emph{The scalar directions.} By Lemma~\ref{sha-siegel-step-seeds-induced-hilbert-form-exists-representation-casimir}(1), applied to $\phi$ and to $\psi$, whose central characters $\omega$ and $1$ have finite order, $(X^v_{11}+X^v_{22})\phi=0$ and $(X^v_{11}+X^v_{22})\psi=0$; by (1), $(X^v_{11}+X^v_{22})\chi=0$, at every infinite place $v$.
\item[(3)] \emph{Condition (d).} By (1), $(X^v_{12}-X^v_{21})\chi=ik\phi\psi+ik'\phi\psi=i(k+k')\chi$, and $(X^v_{11}-X^v_{22})\chi=i\bigl((X^v_{12}+X^v_{21})\phi\bigr)\psi+i\phi(X^v_{12}+X^v_{21})\psi=i(X^v_{12}+X^v_{21})\chi$.
\item[(4)] \emph{The Capelli operators.} The computation of (2) in the proof of Lemma~\ref{sha-siegel-step-seeds-induced-hilbert-form-exists-representation-casimir} uses only the commutation relations of the $X^v_{ab}$, the vanishing of $(X^v_{11}+X^v_{22})$ on the function and condition (d). By (2) and (3) it applies to $\chi$ with weight $k+k'$: $C_v(u+\frac12)\chi=\bigl(u+\frac{k+k'-1}2\bigr)\bigl(u-\frac{k+k'-1}2\bigr)\chi$ for every infinite place $v$ and every $u\in\mathbf C$. So every $C_v(s)$ maps $\chi$ to a multiple of $\chi$. The operators that map $\chi$ to a multiple of $\chi$ form an algebra: it contains the scalars, and if $z\chi=a\chi$ and $z'\chi=b\chi$, then $(z+z')\chi=(a+b)\chi$ and $zz'\chi=ab\chi$. So every operator in the algebra generated by the $C_v(s)$ maps $\chi$ into $\mathbf C\chi$.
\item[(5)] \emph{Automorphy.} $\chi$ is invariant under left translation by $GL_2(E)$, as $\phi$ and $\psi$ are. Its stabilizer in $GL_2(\mathbf A_{E,f})$ contains the intersection of those of $\phi$ and $\psi$, so it is open. Its right translates by the maximal compact subgroup at infinity lie in the span of the products of the translates of $\phi$ and of $\psi$, which is finite-dimensional. Its images under the algebra generated by the Capelli operators lie in $\mathbf C\chi$ by (4). For moderate growth, let $g\in GL_2(\mathbf A_E)$, and let $|\phi(gh)|\le C\|h\|^N$ and $|\psi(gh)|\le C'\|h\|^{N'}$ for $h\in GL_2(E\otimes\mathbf R)$, with $\|h\|$ the height of the definition of moderate growth; then $|\chi(gh)|\le CC'\|h\|^{N+N'}$. So $\chi$ is an automorphic form.
\item[(6)] \emph{Level and central character.} $\chi$ is fixed by $K_1(w^{c(w)})$ for every prime $w$, as $\phi$ and $\psi$ are, and $\chi(zg)=\omega(z)\cdot1\cdot\chi(g)$ for $z\in\mathbf A_E^\times$.
\end{enumerate}
With (3), $\chi$ satisfies the conditions (a) to (d) of Definition~\ref{sha-siegel-step-seeds-induced-hilbert-form-exists-holomorphic} for parallel weight $k+k'$, level $c$ and central character $\omega$.
\end{proof}
% hol_mul (private, HolMul.lean, 35 lines): (1) is Hm.mulsm (Mulsm.lean, 27 lines); (2) is Hm.trzero (Trzero.lean, 68 lines: lie_scalar, as in Lemma sha-siegel-step-seeds-induced-hilbert-form-exists-representation-casimir(1)); (3) is h4 and h5; (4) is Hm.zfinof (Zfinof.lean, 65 lines, with its own copy casimir of the computation and qdet2 for the expansion of TR.cap v u = C_v(u + 1/2); the algebra is TR.Zalg, and key is the induction Algebra.adjoin_induction over its generators, scalars, sums and products); (5) is Hm.mulaut (Mulaut.lean, 46 lines: stab_open, fin_mul and the growth bounds); (6) is the last two fields.

\begin{lemma}\label{sha-siegel-step-seeds-induced-hilbert-form-exists-eisenstein-product-finite}
Let $E$ be a totally real number field with different $\mathfrak d$, $y\in\mathbf A_{E,f}^\times$ and $\xi\in E$. Then the set of totally positive $\eta\in E$ with $v(\eta)<v(\xi)$ at every infinite place $v$ and such that $\eta y$ generates $\mathfrak m\mathfrak d^{-1}$ for some nonzero integral ideal $\mathfrak m$ is finite.
\end{lemma}
\leanref{Development/AutomorphicArithmetic/AutomorphicVarieties/IdeleGenerators/ScalarStabilizers/Boxfin.lean}{UpperTheory.V.Sd.Hm.boxfin}
\begin{proof}
If no $\eta\in E$ has $\eta y$ generating such an ideal, the set is empty. Otherwise fix such an $\eta_0$, with $\eta_0y$ generating $\mathfrak m_0\mathfrak d^{-1}$, and a nonzero $\delta\in\mathfrak m_0$; then $\eta_0\ne0$. Let $\eta$ lie in the set, with $\eta y$ generating $\mathfrak m\mathfrak d^{-1}$. At every prime $w$, $\operatorname{ord}_w(\eta/\eta_0)=\operatorname{ord}_w(\mathfrak m)-\operatorname{ord}_w(\mathfrak m_0)$ and $\operatorname{ord}_w(\delta)\ge\operatorname{ord}_w(\mathfrak m_0)$, so $\operatorname{ord}_w(\delta\eta/\eta_0)\ge0$, and $\delta\eta/\eta_0$ is an algebraic integer. At every infinite place $v$, $|v(\delta\eta/\eta_0)|\le|v(\delta\xi/\eta_0)|$, as $0<v(\eta)<v(\xi)$. There are finitely many algebraic integers of $E$ whose absolute values at the infinite places are bounded by a given number, and $\eta\mapsto\delta\eta/\eta_0$ is injective.
\end{proof}
% Hm.boxfin (Boxfin.lean, 74 lines): eta_0, m_0 and d (the text's delta) with Submodule.exists_mem_ne_zero_of_ne_bot, mem_integers_of_valuation_le_one with val_le, and Embeddings.finite_of_norm_le with the bound the sum over v of v(d xi / eta_0).

\begin{lemma}\label{sha-siegel-step-seeds-induced-hilbert-form-exists-eisenstein-product}
Let $E$ be a totally real number field with uniformizers $\varpi_w$, $\ell$ a prime and $\iota:\bar{\mathbf Q}_\ell\cong\mathbf C$. Let $k,k'\ge0$, $c$ a function from the primes of $E$ to $\mathbf Z_{\ge0}$ and $\omega$ a Hecke character of $E$ of finite order. Let $\phi$ be a holomorphic form of parallel weight $k$, level $c$ and central character $\omega$ with the Fourier expansion of weight $k$ with constant term $0$ and coefficients $a$, with $|\iota^{-1}(a(\mathfrak m))|\le1$ for every $\mathfrak m$. Let $\psi$ be a holomorphic form of parallel weight $k'$, level $0$ and trivial central character with the Fourier expansion of weight $k'$ with constant term $d_0$ and coefficients $b$, with $|\iota^{-1}(d_0(y)-1)|<1$ for every $y$ and $|\iota^{-1}(b(\mathfrak m))|<1$ for every $\mathfrak m$ (Definitions~\ref{sha-siegel-step-seeds-induced-hilbert-form-exists-holomorphic} and~\ref{sha-siegel-step-seeds-induced-hilbert-form-exists-eisenstein-expansion}). Then there are a holomorphic form $\chi$ of parallel weight $k+k'$, level $c$ and central character $\omega$ and a function $a'$ on the nonzero integral ideals such that $\chi$ has the Fourier expansion of weight $k+k'$ with constant term $0$ and coefficients $a'$, and $|\iota^{-1}(a'(\mathfrak m)-a(\mathfrak m))|<1$ for every nonzero integral ideal $\mathfrak m$.
\end{lemma}
\leanref{Development/AutomorphicArithmetic/AutomorphicVarieties/IdeleGenerators/WeightChange.lean}{UpperTheory.V.Sd.HolMul}
\leanref{Development/AutomorphicArithmetic/AutomorphicVarieties/IdeleGenerators/HolomorphicIrreducibility/HolMul.lean}{UpperTheory.V.Sd.holMul}
\begin{proof}
Write $a(y,x)=\begin{pmatrix}y&x\\0&1\end{pmatrix}$ over any ring, $\mathfrak d$ for the different of $E$, and $a(z)$, $b(z)$ for the values at $z\in\mathbf A_{E,f}^\times$ given by Definition~\ref{sha-siegel-step-seeds-induced-hilbert-form-exists-eisenstein-expansion}(1). For $y\in\mathbf A_{E,f}^\times$ and $\eta\in E$ let $\alpha_y(\eta)=a(\eta y)$ and $\beta_y(\eta)=b(\eta y)$ if $\eta$ is totally positive, and $\alpha_y(\eta)=\beta_y(\eta)=0$ otherwise, and let
\[
\gamma_y(\xi)=\sum_{\eta\in E}\alpha_y(\eta)\beta_y(\xi-\eta)\qquad(\xi\in E).
\]
\begin{enumerate}
\item[(1)] \emph{The form.} As $K_1(w^{c(w)})\subseteq K_1(w^0)=GL_2(\mathcal O_{E_w})$, $\psi$ is also a holomorphic form of level $c$. By Lemma~\ref{sha-siegel-step-seeds-induced-hilbert-form-exists-eisenstein-product-holomorphic}, $\chi=\phi\psi$ is a holomorphic form of parallel weight $k+k'$, level $c$ and central character $\omega$.
\item[(2)] \emph{Units.} Let $u\in\mathbf A_{E,f}^\times$ have unit components at every prime. For $z\in\mathbf A_{E,f}^\times$, $zu$ generates the same fractional ideal as $z$, so $a(zu)=a(z)$ and $b(zu)=b(z)$; so $\alpha_{yu}=\alpha_y$, $\beta_{yu}=\beta_y$ and $\gamma_{yu}=\gamma_y$. Also $|u|=1$, and $\psi(g\,a(u,0)_f)=\psi(g)$: the stabilizer of $\psi$ contains the elements of $GL_2(\hat{\mathcal O}_E)$ whose components at the primes of some finite set $T$ are $1$, and $a(u_w,0)\in GL_2(\mathcal O_{E_w})$ fixes $\psi$ for $w\in T$, as in the proof of Lemma~\ref{sha-siegel-step-seeds-induced-hilbert-form-exists-eisenstein-lift-hecke-translation}(1). So the expansions of $\psi$ at $(0,1,y)$ and at $(0,1,yu)$ have the same terms and the same value of $\psi$, and $d_0(yu)=d_0(y)$.
\item[(3)] \emph{Totally positive elements.} Let $\xi\in E$ be totally positive and $y\in\mathbf A_{E,f}^\times$. Then $d_0(\xi y)=d_0(y)$. Indeed, as $\psi$ is invariant under left translation by $a(\xi,0)\in GL_2(E)$, $\psi\bigl(a(1,0)_\infty a(y,0)_f\bigr)=\psi\bigl(a(\xi,0)_\infty a(\xi y,0)_f\bigr)$. By the product formula, $\prod_vv(\xi)\cdot|\xi|=1$, so $\prod_vv(\xi)\cdot|\xi y|=|y|$. The terms of the expansion of $\psi$ at $(0,(v(\xi))_v,\xi y)$ are $b(\eta\xi y)\exp(-2\pi\sum_vv(\eta\xi))$, those of the expansion at $(0,1,y)$ reindexed by $\eta\mapsto\eta\xi$. So the two expansions have the same sum and the same value of $\psi$, and $d_0(\xi y)=d_0(y)$. Also $\alpha_{\xi y}(\eta)=\alpha_y(\eta\xi)$ and $\beta_{\xi y}(\eta)=\beta_y(\eta\xi)$, so, reindexing by $\eta\mapsto\eta\xi$, $\gamma_{\xi y}(1)=\gamma_y(\xi)$.
\item[(4)] \emph{Supports.} If $\gamma_y(\xi)\ne0$, then $\xi$ is totally positive and $\xi y$ generates $\mathfrak m\mathfrak d^{-1}$ for some nonzero integral ideal $\mathfrak m$. Indeed, some $\eta$ has $\alpha_y(\eta)\beta_y(\xi-\eta)\ne0$; then $\eta$ and $\xi-\eta$ are totally positive, and $\eta y$ and $(\xi-\eta)y$ generate $\mathfrak m_1\mathfrak d^{-1}$ and $\mathfrak m_2\mathfrak d^{-1}$ with $\mathfrak m_1,\mathfrak m_2$ integral. So $\xi$ is totally positive, hence nonzero, and $\operatorname{ord}_w(\xi y_w)\ge\min\bigl(\operatorname{ord}_w(\eta y_w),\operatorname{ord}_w((\xi-\eta)y_w)\bigr)\ge-\operatorname{ord}_w(\mathfrak d)$ at every prime $w$. So $\xi y$ generates $\mathfrak m\mathfrak d^{-1}$ with $\mathfrak m=\prod_ww^{\operatorname{ord}_w(\xi y_w)+\operatorname{ord}_w(\mathfrak d)}$ integral.
\item[(5)] \emph{The coefficients.} Let $A(y)=a(y)d_0(y)+\gamma_y(1)$ for $y\in\mathbf A_{E,f}^\times$. For each nonzero integral ideal $\mathfrak m$ choose $Y_\mathfrak m\in\mathbf A_{E,f}^\times$ that generates $\mathfrak m\mathfrak d^{-1}$, and let $a'(\mathfrak m)=A(Y_\mathfrak m)$. Then the value $a'(y)$ given by Definition~\ref{sha-siegel-step-seeds-induced-hilbert-form-exists-eisenstein-expansion}(1) is $A(y)$ for every $y\in\mathbf A_{E,f}^\times$. If $y$ generates $\mathfrak m\mathfrak d^{-1}$, then $u=y^{-1}Y_\mathfrak m$ has unit components, and $a'(y)=A(yu)=A(y)$ by (2). Otherwise $a'(y)=0$, $a(y)=0$, and $\gamma_y(1)=0$ by (4).
\item[(6)] \emph{The expansion.} Let $x\in E\otimes\mathbf R$, $t\in E\otimes\mathbf R$ totally positive and $y\in\mathbf A_{E,f}^\times$; let $g=a(t,x)_\infty a(y,0)_f$, $D=\prod_vt_v\cdot|y|$ and $Q(\xi)=\exp\bigl(2\pi i\sum_vv(\xi)(x_v+it_v)\bigr)$ for $\xi\in E$, so that $Q(\eta)Q(\xi-\eta)=Q(\xi)$. The expansions of $\phi$ and $\psi$ give
\[
\sum_{\xi\in E}\alpha_y(\xi)Q(\xi)=\phi(g)D^{-k/2},\qquad\sum_{\xi\in E}\beta_y(\xi)Q(\xi)=\psi(g)D^{-k'/2}-d_0(y),
\]
for summable, hence absolutely summable, families. The Cauchy product of two absolutely summable families on the additive group $E$, with $d_0(y)$ times the first family added, gives
\[
\sum_{\xi\in E}\bigl(\alpha_y(\xi)d_0(y)+\gamma_y(\xi)\bigr)Q(\xi)=\phi(g)D^{-k/2}d_0(y)+\phi(g)D^{-k/2}\bigl(\psi(g)D^{-k'/2}-d_0(y)\bigr)=\chi(g)D^{-(k+k')/2}.
\]
For $\xi$ totally positive, the coefficient is $a(\xi y)d_0(\xi y)+\gamma_{\xi y}(1)=A(\xi y)=a'(\xi y)$, by (3) and (5); otherwise it is $0$, by (4). So $\chi$ has the Fourier expansion of weight $k+k'$ with constant term $0$ and coefficients $a'$.
\item[(7)] \emph{The congruence.} Let $\mathfrak m$ be a nonzero integral ideal and $y=Y_\mathfrak m$. Then $a(y)=a(\mathfrak m)$ and
\[
a'(\mathfrak m)-a(\mathfrak m)=a(\mathfrak m)\bigl(d_0(y)-1\bigr)+\gamma_y(1).
\]
The first term has $|\iota^{-1}(\cdot)|<1$, as $|\iota^{-1}(a(\mathfrak m))|\le1$ and $|\iota^{-1}(d_0(y)-1)|<1$. The nonzero terms of $\gamma_y(1)$ have $\eta$ and $1-\eta$ totally positive and $\eta y$ generating some $\mathfrak m_1\mathfrak d^{-1}$; by Lemma~\ref{sha-siegel-step-seeds-induced-hilbert-form-exists-eisenstein-product-finite} there are finitely many of them, and each is $a(\mathfrak m_1)b(\mathfrak m_2)$ for nonzero integral ideals $\mathfrak m_1,\mathfrak m_2$, with $|\iota^{-1}(a(\mathfrak m_1)b(\mathfrak m_2))|<1$. By the ultrametric inequality, $|\iota^{-1}(a'(\mathfrak m)-a(\mathfrak m))|<1$.
\end{enumerate}
\end{proof}
% holMul (HolMul.lean, 62 lines, with the private lemmas of HolMul.lean, 290 lines): the text's alpha_y, beta_y and gamma_y are cf a y, cf b y and cv a b y. (1) is hol_mul with holLin (level 0 below c); (2) is coefAt_unit, cf_unit and const_unit with Hm.fixunit and Hm.absf; (3) is const_rat (rat_aff, Hm.absf for the product formula), cf_rat and cv_rat; (4) is cv_ne with genD_add and genD_of_le; (5) is exists_genD for Y, A and hA; (6) is hasSum_cf, conv (summable_mul_of_summable_norm and prod_fiberwise) and hE'; (7) is e and cv_norm with Hm.boxfin (Lemma sha-siegel-step-seeds-induced-hilbert-form-exists-eisenstein-product-finite). About 920 lines with the dependencies.

\begin{lemma}\label{sha-siegel-step-seeds-induced-hilbert-form-exists-eisenstein-lift-finite-rotation}
Let $E$ be a totally real number field with set of infinite places $\Sigma$. For $\theta\in E\otimes\mathbf R$ let $r(\theta)\in GL_2(E\otimes\mathbf R)$ be the rotation of Lemma~\ref{sha-siegel-step-seeds-induced-hilbert-form-exists-eisenstein-lift-unique-rotation}, with component $\begin{pmatrix}\cos\theta_v&\sin\theta_v\\-\sin\theta_v&\cos\theta_v\end{pmatrix}$ at each $v\in\Sigma$, and for a function $s:\Sigma\to\{\pm1\}$ let $\epsilon(s)\in GL_2(E\otimes\mathbf R)$ have component $\operatorname{diag}(1,s(v))$ at each $v$.
\begin{enumerate}
\item Let $z\in\mathbf C$ and $\phi$ a function on $GL_2(\mathbf A_E)$, smooth in the variable at infinity, with $(X^v_{12}-X^v_{21})\phi=z\phi$ at every infinite place $v$, for the derivatives $X^v_{ab}$ of Definition~\ref{sha-siegel-step-seeds-induced-hilbert-form-exists-holomorphic}. Then $\phi(g\,r(\theta)_\infty)=\exp\bigl(z\sum_v\theta_v\bigr)\phi(g)$ for every $g\in GL_2(\mathbf A_E)$ and $\theta\in E\otimes\mathbf R$.
\item Every element of $GL_2(E\otimes\mathbf R)$ whose components are orthogonal matrices is $\epsilon(s)r(\theta)$ for some $s:\Sigma\to\{\pm1\}$ and $\theta\in E\otimes\mathbf R$.
\end{enumerate}
\end{lemma}
\leanref{Development/AutomorphicArithmetic/AutomorphicVarieties/GaloisLibrary/CrystallineBaseChange/HfStmt.lean}{UpperTheory.V.Br.RotEig}
\leanref{Development/AutomorphicArithmetic/AutomorphicVarieties/GaloisLibrary/CrystallineBaseChange/RotEig.lean}{UpperTheory.V.Br.rotEig}
\leanref{Development/AutomorphicArithmetic/AutomorphicVarieties/GaloisLibrary/CrystallineBaseChange/CuspSet.lean}{UpperTheory.V.Br.CuspStabilizer.dec}
\begin{proof}
(1) Let $f(t)=\phi(g\,r(t\theta)_\infty)$ for $t\in\mathbf R$, and $\sigma=z\sum_v\theta_v$. As $r((t+s)\theta)=r(t\theta)r(s\theta)$, and the derivative of $s\mapsto r(s\theta)$ at $s=0$ has component $\theta_v(e_{12}-e_{21})$ at each $v$, the chain rule gives $f'(t)=\sum_v\theta_v\bigl((X^v_{12}-X^v_{21})\phi\bigr)(g\,r(t\theta)_\infty)=\sigma f(t)$. So $f(t)\exp(-\sigma t)$ has derivative $0$, and is constant. So $\phi(g\,r(\theta)_\infty)=f(1)=\exp(\sigma)f(0)=\exp(\sigma)\phi(g)$.

(2) Let $h$ be a real orthogonal $2\times2$ matrix with rows $(a,b)$ and $(c,d)$. As $a^2+b^2=1$, $(a,b)=(\cos t,\sin t)$ for some $t\in\mathbf R$. Let $s=ad-bc$. From $a^2+b^2=1$, $ac+bd=0$ and $c^2+d^2=1$ one gets $c=-sb$, $d=sa$ and $s^2=1$. So $h=\operatorname{diag}(1,s)\begin{pmatrix}\cos t&\sin t\\-\sin t&\cos t\end{pmatrix}$. Apply this at each infinite place.
\end{proof}
% rotEig (RotEig.lean, 76 lines): the private rot_eig, with X = toL [[0, theta], [-theta, 0]] = sum_v theta_v (E_01 - E_10 at v), hL for the derivative, TR.hasDerivAt_curve and is_const_of_deriv_eq_zero; it is the argument of norm_rot in Lemma sha-siegel-step-seeds-induced-hilbert-form-exists-eisenstein-lift-unique-rotation for any z. (2) is the private dec with real_dec, in CuspSet.lean and again in HolCuspFin.lean; the text's epsilon(s) and r(theta) are the private eps s and rot theta there.

\begin{lemma}\label{sha-siegel-step-seeds-induced-hilbert-form-exists-eisenstein-lift-finite-harish-chandra}
Let $E$ be a totally real number field and $a\ge0$ an integer. For an infinite place $v$ and $1\le i,j\le a$, let $X^v_{ij}$ be the derivative of right translation at $v$ in the direction of the matrix unit $e_{ij}$, acting on the functions on $GL_a(\mathbf A_E)$ that are smooth in the variable at infinity, as in Definition~\ref{sha-siegel-step-seeds-induced-hilbert-form-exists-holomorphic}. For $s\in\mathbf C$ let $T_{ij}(s)=s\delta_{ij}+X^v_{ij}$ and
\[
C^{(a)}_v(s)=\sum_{\sigma\in S_a}\operatorname{sgn}(\sigma)\,T_{\sigma(1)1}(s)\,T_{\sigma(2)2}(s-1)\cdots T_{\sigma(a)a}(s-a+1),
\]
so that $C^{(2)}_v(s)$ is the Capelli operator $C_v(s)$ of that definition. Call such a function $\phi$ an \emph{automorphic form} on $GL_a$ if it satisfies condition (a) of Definition~\ref{sha-siegel-step-seeds-induced-hilbert-form-exists-holomorphic} with $GL_a$ in place of $GL_2$ and the operators $C^{(a)}_v(s)$ in place of the $C_v(s)$, and a \emph{cusp form} on $GL_a$ if moreover $\int_D\phi(n(x)g)\,d\mu(x)=0$ for every $0<r<a$, every $g\in GL_a(\mathbf A_E)$, every Haar measure $\mu$ on the $r\times(a-r)$ matrices over $\mathbf A_E$ and every measurable fundamental domain $D$ for the translations by the matrices over $E$, where $n(x)=\begin{pmatrix}1_r&x\\0&1_{a-r}\end{pmatrix}$. Let $K_\infty\subseteq GL_a(E\otimes\mathbf R)$ be the group of the elements whose components are orthogonal matrices, $L$ an open subgroup of $GL_a(\mathbf A_{E,f})$, $\rho$ a representation of $K_\infty$ on a finite-dimensional complex vector space $W$ whose matrix coefficients $k\mapsto\lambda(\rho(k)w)$, for $\lambda\in W^*$ and $w\in W$, are continuous, and $\Lambda$ a complex function on the pairs of an infinite place and a complex number. Then the cusp forms $\phi$ on $GL_a$ such that $\phi$ is fixed by $L$ under right translation, for every $g\in GL_a(\mathbf A_E)$ the function $k\mapsto\phi(g\,k_\infty)$ on $K_\infty$ is a linear combination of matrix coefficients of $\rho$, and $C^{(a)}_v(u+\frac{a-1}2)\phi=\Lambda(v,u)\phi$ for every infinite place $v$ and every $u\in\mathbf C$, span a finite-dimensional space.
\end{lemma}
\leanref{Development/AutomorphicArithmetic/AutomorphicVarieties/AutomorphicForms/TemperedLocalComponents/PCuspFin.lean}{UpperTheory.V.Af.cuspFormsFiniteReal}
\begin{proof}
\openstep{The proof of this form of Harish-Chandra's finiteness theorem, through the reduction theory of $GL_a$ and the centre of the enveloping algebra, is not yet written out.}
% cuspFormsFiniteReal (PCuspFin.lean, 83 lines): with a basis b of the Lie algebra by matrix units (CuspFin.exists_stdBasis), the subalgebra S of operators with a common eigenvalue on the forms with Capelli eigenvalues c, the Casimir operators in S (CuspFin.casMem) and the trace X_11 + ... + X_aa in S (CuspRep.Central.trace_mem), a monic polynomial q in the Casimir of K_infinity killing the functions of type rho (CuspFin.typePoly), the Laplacian sum_l b_l^2 = sum_v Cas_v + 1/2 sum_l X_l^2 (sqId, hDelta) with aevalShift, and the finiteness CuspFin.eigenFinite_of with CuspFin.mem_cuspForms_of for cusp forms fixed by L that are killed by a monic polynomial in the Laplacian and are eigenvectors of the trace. Size: about 11,700 lines with its dependencies (tools/lean_size.py, 10,100 in theorems), mostly reduction theory (ReductionCore, ReductionFinite, ReductionSiegel, ReductionDirichlet, ReductionIwasawa) and the Harish-Chandra centre (HarishChandraCentre). A citation of Harish-Chandra's theorem with a checked locator, together with the statement that the Capelli operators generate the centre of the enveloping algebra, would also reach the floor.
\end{proof}

\begin{lemma}\label{sha-siegel-step-seeds-induced-hilbert-form-exists-eisenstein-lift-finite-cusp}
Let $E$ be a totally real number field, $k\ge0$, $c$ a function from the primes of $E$ to $\mathbf Z_{\ge0}$ that vanishes at all but finitely many primes, and $\omega$ a Hecke character of $E$ of finite order. For an integer $M\ge0$ let $K(M)$ be the group of $g\in GL_2(\hat{\mathcal O}_E)$ with $g\equiv1$ modulo $M$; so $K(0)=\{1\}$ and $K(1)=GL_2(\hat{\mathcal O}_E)$. Then:
\begin{enumerate}
\item there is $M\ge1$ such that every holomorphic form of parallel weight $k$, level $c$ and central character $\omega$ (Definition~\ref{sha-siegel-step-seeds-induced-hilbert-form-exists-holomorphic}) is fixed by $K(M)$ under right translation;
\item the holomorphic forms of parallel weight $k$, level $c$ and central character $\omega$ that are cusp forms span a finite-dimensional space.
\end{enumerate}
\end{lemma}
\leanref{Development/AutomorphicArithmetic/AutomorphicVarieties/GaloisLibrary/CrystallineBaseChange/HfStmt.lean}{UpperTheory.V.Br.HolCuspFin}
\leanref{Development/AutomorphicArithmetic/AutomorphicVarieties/GaloisLibrary/CrystallineBaseChange/HolCuspFin.lean}{UpperTheory.V.Br.holCuspFin}
\begin{proof}
(1) Let $M=\prod_w\operatorname N(w)^{c(w)}$, the product over the finitely many primes $w$ with $c(w)\ne0$; it is an integer $\ge1$. For every prime $w$, $\operatorname N(w)\in w$, so $M\in w^{c(w)}\mathcal O_{E_w}$. So for $u\in K(M)$ and every prime $w$, the component $u_w$ is congruent to $1$ modulo $w^{c(w)}$, and lies in $K_1(w^{c(w)})$. Let $\phi$ be a holomorphic form of parallel weight $k$, level $c$ and central character $\omega$, and $U$ its stabilizer in $GL_2(\mathbf A_{E,f})$, which is open. By condition (b) of Definition~\ref{sha-siegel-step-seeds-induced-hilbert-form-exists-holomorphic}, $U$ contains $(u_w)_w$, the element with component $u_w$ at $w$ and $1$ elsewhere, for every $u\in K(M)$ and every $w$. As $U$ is a neighbourhood of $1$, there is a finite set $S$ of primes such that $U$ contains every element of $K(M)$ whose components at the primes of $S$ are $1$. Let $u\in K(M)$. Multiplying $u$ on the right by $((u_w)_w)^{-1}$ for each $w\in S$ gives an element of $K(M)$ whose components at $S$ are $1$, which lies in $U$; so $u\in U$. So $\phi$ is fixed by $K(M)$.

(2) Let $\Sigma$ be the set of infinite places, $K_\infty\subseteq GL_2(E\otimes\mathbf R)$ the group of the elements whose components are orthogonal, and $\epsilon(s)$ and $r(\theta)$ as in Lemma~\ref{sha-siegel-step-seeds-induced-hilbert-form-exists-eisenstein-lift-finite-rotation}.
\begin{enumerate}
\item[(a)] \emph{A representation of $K_\infty$.} Let $V$ be the space of continuous functions $f$ on $K_\infty$ with $f(x\,r(\theta))=\exp(ik\sum_v\theta_v)f(x)$ for all $x\in K_\infty$ and $\theta\in E\otimes\mathbf R$. By Lemma~\ref{sha-siegel-step-seeds-induced-hilbert-form-exists-eisenstein-lift-finite-rotation}(2), every $x\in K_\infty$ is $\epsilon(s)r(\theta)$, so $f\in V$ is determined by its values at the finitely many $\epsilon(s)$, and $V$ is finite-dimensional. The group $K_\infty$ acts on $V$ by $(\lambda(k)f)(x)=f(k^{-1}x)$. Let $\rho$ be the dual representation on $W=V^*$. Its matrix coefficients are continuous: for $f\in V$ and $w\in W$, $k\mapsto w(\lambda(k^{-1})f)$ is continuous, as $w$ is a linear form on the finite-dimensional space $V$ and $k\mapsto(\lambda(k^{-1})f)(x)$ is continuous for every $x$. Every $f\in V$ is a matrix coefficient of $\rho$: with $e\in W$ the evaluation at $1$, $f(k)=(\lambda(k^{-1})f)(1)=(\rho(k)e)(f)$.
\item[(b)] \emph{The type.} Let $\phi$ be a holomorphic form of parallel weight $k$, level $c$ and central character $\omega$, and $g\in GL_2(\mathbf A_E)$. The function $x\mapsto\phi(g\,x_\infty)$ on $K_\infty$ is continuous, as $\phi$ is smooth in the variable at infinity, and lies in $V$, by Lemma~\ref{sha-siegel-step-seeds-induced-hilbert-form-exists-eisenstein-lift-finite-rotation}(1) with $z=ik$, applied at $g\,x_\infty$, and condition (d). By (a), it is a matrix coefficient of $\rho$.
\item[(c)] \emph{The Capelli operators.} By Lemma~\ref{sha-siegel-step-seeds-induced-hilbert-form-exists-representation-casimir}(2), $C_v(u+\frac12)\phi=\bigl(u+\frac{k-1}2\bigr)\bigl(u-\frac{k-1}2\bigr)\phi$ for every infinite place $v$ and $u\in\mathbf C$.
\item[(d)] \emph{Finiteness.} Let $M$ be as in (1). The group $K(M)$ is open, as $M\ge1$. A holomorphic form that is a cusp form is a cusp form on $GL_2$ in the sense of Lemma~\ref{sha-siegel-step-seeds-induced-hilbert-form-exists-eisenstein-lift-finite-harish-chandra}. By (1), (b) and (c), the holomorphic forms of parallel weight $k$, level $c$ and central character $\omega$ that are cusp forms lie in the set of that lemma for $a=2$, $L=K(M)$, $\rho$ and $\Lambda(v,u)=(u+\frac{k-1}2)(u-\frac{k-1}2)$. So they span a finite-dimensional space.
\end{enumerate}
\end{proof}
% holCuspFin (HolCuspFin.lean, 178 lines), followed step by step. (1) is exists_level (M the product of absNorm w ^ c w), local_mem_principal_congruence and fixed_KLev (exists_KAway_subset, mem_KAway and an induction over S, the text's product over S). (2): (a) is the private exists_rep with e = eps, r = rot, chi theta = exp(i k sum_v theta_v) and dec, rho the dual L.dual of the left translation on V; (b) is cont and rotEig (Lemma sha-siegel-step-seeds-induced-hilbert-form-exists-eisenstein-lift-finite-rotation(1)) with the field wt; (c) is holCap (HolCap.lean, 78 lines), the Lean's own copy for parallel weight of the computation of Lemma sha-siegel-step-seeds-induced-hilbert-form-exists-representation-casimir(2); (d) is cuspFormsFiniteReal (Lemma sha-siegel-step-seeds-induced-hilbert-form-exists-eisenstein-lift-finite-harish-chandra) with isOpen_KLev and Submodule.finiteDimensional_of_le. The open step left is in that lemma.

\begin{lemma}\label{sha-siegel-step-seeds-induced-hilbert-form-exists-eisenstein-lift-finite-iwasawa}
Let $E$ be a totally real number field, and write $a(y,x)=\begin{pmatrix}y&x\\0&1\end{pmatrix}$ over any ring. Then every $g\in GL_2(\mathbf A_E)$ is $z\,a(y,x)\,h_\infty u_f$ for some ideles $z$, viewed as a scalar matrix, and $y$, some adele $x$, some $h\in GL_2(E\otimes\mathbf R)$ whose components are orthogonal matrices, and some $u\in GL_2(\hat{\mathcal O}_E)$.
\end{lemma}
\leanref{Development/AutomorphicArithmetic/AutomorphicVarieties/GaloisLibrary/CrystallineBaseChange/HfStmt.lean}{UpperTheory.V.Br.IwasawaA}
\leanref{Development/AutomorphicArithmetic/AutomorphicVarieties/GaloisLibrary/CrystallineBaseChange/IwasawaA.lean}{UpperTheory.V.Br.iwasawaA}
\leanref{GroundTheory/AutomorphicForms/ReductionIwasawa.lean}{GroundTheory.Automorphic.Reduction.exists_adelicBorel_mul_maxCompact}
\begin{proof}
At an infinite place $v$, the Gram--Schmidt process applied to the rows of $g_v$, starting from the second, gives $g_v=b_vh_v$ with $b_v$ upper triangular and $h_v$ in the group of matrices that preserve the Euclidean norm of row vectors, that is, orthogonal. At a prime $w$, let $(c,d)$ be the second row of $g_w$; if $|c|\le|d|$, then $g_w\begin{pmatrix}1&0\\-c/d&1\end{pmatrix}$ is upper triangular, and otherwise $g_w\begin{pmatrix}0&1\\1&0\end{pmatrix}\begin{pmatrix}1&0\\-d/c&1\end{pmatrix}$ is. These matrices lie in $GL_2(\mathcal O_{E_w})$, so $g_w=b_wu_w$ with $b_w$ upper triangular and $u_w\in GL_2(\mathcal O_{E_w})$; take $b_w=1$ and $u_w=g_w$ when $g_w\in GL_2(\mathcal O_{E_w})$, which holds at all but finitely many $w$. So $g=bk$ with $b\in GL_2(\mathbf A_E)$ upper triangular, $k_\infty$ orthogonal at every infinite place and $k_f\in GL_2(\hat{\mathcal O}_E)$. The diagonal entries $b_{11}$ and $b_{22}$ are ideles, as their product $\det b$ is one, and $b=z\,a(y,x)$ with $z=b_{22}$, $y=b_{11}/b_{22}$ and $x=b_{12}/b_{22}$.
\end{proof}
% iwasawaA (IwasawaA.lean, 53 lines): Reduction.exists_adelicBorel_mul_maxCompact (ReductionIwasawa.lean, the decomposition g = b k with b in adelicBorel and k in maxCompact, from exists_upper_mul_archCompact' at the infinite places and Iwasawa.exists_borel_mul_integralGL at the primes, with p.1 = 1 where g_w is integral; the Lean proves both local decompositions for GL_n, at the primes through lattices), then the private upper_eq for b = z aff y x, orth (a matrix preserving the Euclidean norm is orthogonal) for k at infinity, and mem_integralFiniteGL_iff with mem_KLev for k_f in KLev F 2 1. The text gives the local decompositions for GL_2 by Gram-Schmidt and by column operations.

\begin{lemma}\label{sha-siegel-step-seeds-induced-hilbert-form-exists-eisenstein-lift-finite-points}
Let $E$ be a totally real number field. Write $a(y,x)=\begin{pmatrix}y&x\\0&1\end{pmatrix}$ and $n(x)=a(1,x)$ over any ring, $r(\theta)$ for the rotation of Lemma~\ref{sha-siegel-step-seeds-induced-hilbert-form-exists-eisenstein-lift-unique-rotation} and $K(M)$ as in Lemma~\ref{sha-siegel-step-seeds-induced-hilbert-form-exists-eisenstein-lift-finite-cusp}. For $M\ge0$, call a set $Z\subseteq GL_2(\mathbf A_E)$ \emph{closed under moves for $M$} if for every $g\in Z$ the following lie in $Z$: $n(x)g$ for $x\in\mathbf A_E$; $a(\alpha,0)g$ for $\alpha\in E^\times$; $a(y,0)_\infty g$ for totally positive $y\in E\otimes\mathbf R$; $zg$ for scalar ideles $z$; $g\,u_f$ for $u\in K(M)$; and $g\,r(\theta)_\infty$ for $\theta\in E\otimes\mathbf R$.
\begin{enumerate}
\item For every $M\ge1$ there is a finite set $S\subset GL_2(\mathbf A_E)$ such that every set containing $S$ and closed under moves for $M$ is $GL_2(\mathbf A_E)$.
\item Let $\mu$ be a Haar measure on $\mathbf A_E$ and $D$ a measurable fundamental domain for the translations by $E$. Let $k\ge0$, $c$ a function from the primes of $E$ to $\mathbf Z_{\ge0}$, $\omega$ a Hecke character of $E$ of finite order, $M\ge0$, and $\phi$ a holomorphic form of parallel weight $k$, level $c$ and central character $\omega$ that is fixed by $K(M)$. Then the set of $g$ with $\int_D\phi(n(x)g)\,d\mu(x)=0$ is closed under moves for $M$.
\item There are a Haar measure $\mu$ on $\mathbf A_E$ and a measurable fundamental domain $D$ for the translations by $E$ such that every function $\phi$ on $GL_2(\mathbf A_E)$ that is invariant under left translation by $GL_2(E)$ and has $\int_D\phi(n(x)g)\,d\mu(x)=0$ for every $g$ is cuspidal: $\int_{D'}\phi(n(x)g)\,d\mu'(x)=0$ for every Haar measure $\mu'$, fundamental domain $D'$ and $g$.
\end{enumerate}
\end{lemma}
\leanref{Development/AutomorphicArithmetic/AutomorphicVarieties/GaloisLibrary/CrystallineBaseChange/HfStmt.lean}{UpperTheory.V.Br.Moves}
\leanref{Development/AutomorphicArithmetic/AutomorphicVarieties/GaloisLibrary/CrystallineBaseChange/HfStmt.lean}{UpperTheory.V.Br.CuspSet}
\leanref{Development/AutomorphicArithmetic/AutomorphicVarieties/GaloisLibrary/CrystallineBaseChange/CuspSet.lean}{UpperTheory.V.Br.cuspSet}
\leanref{Development/AutomorphicArithmetic/AutomorphicVarieties/GaloisLibrary/CrystallineBaseChange/HfStmt.lean}{UpperTheory.V.Br.CtMoves}
\leanref{Development/AutomorphicArithmetic/AutomorphicVarieties/GaloisLibrary/CrystallineBaseChange/CtMoves.lean}{UpperTheory.V.Br.ctMoves}
\leanref{Development/AutomorphicArithmetic/AutomorphicVarieties/GaloisLibrary/CrystallineBaseChange/HfStmt.lean}{UpperTheory.V.Br.CtOne}
\leanref{Development/AutomorphicArithmetic/AutomorphicVarieties/GaloisLibrary/CrystallineBaseChange/CoRatAff.lean}{UpperTheory.V.Br.ctOne}
\begin{proof}
(3) Let $\nu$ be a Haar measure on $\mathbf A_E$ and $D_0$ a measurable fundamental domain for the translations by $E$; both exist. Suppose that $\phi$ is invariant under left translation by $GL_2(E)$ and $\int_{D_0}\phi(n(x)g)\,d\nu(x)=0$ for every $g$. Let $\mu'$, $D'$ and $g$ be given. For $\gamma\in E$, $n(\gamma+x)=n(\gamma)n(x)$ with $n(\gamma)\in GL_2(E)$, so $x\mapsto\phi(n(x)g)$ is invariant under translation by $E$. So its integrals against $\nu$ over the fundamental domains $D'$ and $D_0$ agree, and the first is $0$. As $\mu'$ is a positive multiple of $\nu$, $\int_{D'}\phi(n(x)g)\,d\mu'(x)=0$.

(1) Let $\Sigma$ be the set of infinite places, and $\epsilon(s)$, for $s:\Sigma\to\{\pm1\}$, as in Lemma~\ref{sha-siegel-step-seeds-induced-hilbert-form-exists-eisenstein-lift-finite-rotation}. As $a(t,0)a(t',0)=a(tt',0)$ for ideles $t,t'$, which commute,
\[
a(tu,0)g=a(u,0)\,a(t,0)g\qquad(t,u\in\mathbf A_E^\times,\ g\in GL_2(\mathbf A_E)).
\]
\begin{enumerate}
\item[(a)] \emph{Decomposition.} By Lemmas~\ref{sha-siegel-step-seeds-induced-hilbert-form-exists-eisenstein-lift-finite-iwasawa} and~\ref{sha-siegel-step-seeds-induced-hilbert-form-exists-eisenstein-lift-finite-rotation}(2), every $g\in GL_2(\mathbf A_E)$ is $z\,a(y,x)\,\epsilon(s)_\infty r(\theta)_\infty u_f$ for some ideles $z$ and $y$, an adele $x$, some $s:\Sigma\to\{\pm1\}$, some $\theta\in E\otimes\mathbf R$ and some $u\in K(1)$.
\item[(b)] \emph{Stabilizers.} For $g\in GL_2(\mathbf A_E)$, let $H_g$ be the set of ideles $u$ such that, for every set $Z$ closed under moves for $M$ and every idele $t$, $a(tu,0)g\in Z$ if and only if $a(t,0)g\in Z$. It is a subgroup of the ideles. By the display and the moves by $a(\alpha,0)$ and $a(y,0)_\infty$, applied to $\alpha$ and $\alpha^{-1}$ and to $y$ and $y^{-1}$, it contains $E^\times$ and the totally positive elements of $E\otimes\mathbf R$. Let $u$ be a finite idele with $g_f^{-1}a(u,0)g_f\in K(M)$. Then $a(u,0)g=g\,(g_f^{-1}a(u,0)g_f)_f$, so $a(tu,0)g=a(t,0)g\,(g_f^{-1}a(u,0)g_f)_f$, and likewise for $u^{-1}$, as $g_f^{-1}a(u^{-1},0)g_f=(g_f^{-1}a(u,0)g_f)^{-1}$; by the move by $K(M)$, $u\in H_g$. As $K(M)$ is open for $M\ge1$, these $u$ form a neighbourhood of $1$ in $\mathbf A_{E,f}^\times$. By step~(b) of the proof of Lemma~\ref{sha-siegel-step-seeds-induced-hilbert-form-exists-eisenstein-noncuspidal-characters}, $H_g$ has finite index in the ideles.
\item[(c)] \emph{The finite set.} The group $K(1)$ is compact and $K(M)$ is an open subgroup, so $K(1)/K(M)$ is finite; let $R\subseteq K(1)$ be a set of representatives. For $s:\Sigma\to\{\pm1\}$ and $\rho\in R$ let $b_{s,\rho}=\epsilon(s)_\infty\rho_f$, and let $T_{s,\rho}$ be a set of representatives of the ideles modulo $H_{b_{s,\rho}}$, finite by (b). As $\Sigma$ is finite, the set $S$ of the $a(t,0)b_{s,\rho}$ with $s:\Sigma\to\{\pm1\}$, $\rho\in R$ and $t\in T_{s,\rho}$ is finite.
\item[(d)] \emph{The moves reach everything.} Let $Z\supseteq S$ be closed under moves for $M$. For every $s$, $\rho$ and idele $t$, $a(t,0)b_{s,\rho}\in Z$: write $t=t'q$ with $t'\in T_{s,\rho}$ and $q\in H_{b_{s,\rho}}$; then $a(t',0)b_{s,\rho}\in S\subseteq Z$, and $a(t,0)b_{s,\rho}\in Z$ by the definition of $H_{b_{s,\rho}}$. Let $g\in GL_2(\mathbf A_E)$, written as in (a), and let $\rho\in R$ represent $uK(M)$, so $\rho=uh$ with $h\in K(M)$. Then $a(y,0)\epsilon(s)_\infty\rho_f\in Z$, and the moves by $h^{-1}\in K(M)$, by $r(\theta)$, by $n(x)$ and by $z$ give $z\,n(x)\,a(y,0)\epsilon(s)_\infty\rho_fh^{-1}_fr(\theta)_\infty\in Z$. As components at infinity commute with finite ones, $\rho h^{-1}=u$ and $n(x)a(y,0)=a(y,x)$, this element is $g$. So $Z=GL_2(\mathbf A_E)$.
\end{enumerate}

(2) Let $\psi(g)=\int_D\phi(n(x)g)\,d\mu(x)$, defined by Lemma~\ref{sha-siegel-step-seeds-induced-hilbert-form-exists-eisenstein-noncuspidal-constant-term}(1), and $Z$ the set of $g$ with $\psi(g)=0$. By (2) and (3) of that lemma, $\psi(n(x)g)=\psi(g)$, $\psi(a(\alpha,0)g)=\psi(g)$ and $\psi(a(y,0)_\infty g)=(\prod_vy_v)^{k/2}\psi(g)$; so $Z$ is closed under the moves by $n(x)$, $a(\alpha,0)$ and $a(y,0)_\infty$. Let $h\in GL_2(\mathbf A_E)$ and $e\in\mathbf C$ with $\phi(g'h)=e\phi(g')$ for every $g'$. Then $\psi(gh)=\int_D\phi(n(x)gh)\,d\mu(x)=e\psi(g)$, so $gh\in Z$ for $g\in Z$. This applies to a scalar idele $z$ with $e=\omega(z)$, as $\phi(g'z)=\phi(zg')=\omega(z)\phi(g')$ by condition (c) of Definition~\ref{sha-siegel-step-seeds-induced-hilbert-form-exists-holomorphic}; to $u_f$ with $u\in K(M)$ and $e=1$; and to $r(\theta)_\infty$ with $e=\exp(ik\sum_v\theta_v)$, by condition (d) and Lemma~\ref{sha-siegel-step-seeds-induced-hilbert-form-exists-eisenstein-lift-finite-rotation}(1). As $zg=gz$ for a scalar $z$, $Z$ is closed under the remaining moves.
% cuspSet (CuspSet.lean, 149 lines), followed step by step: the display is am; (a) is iwasawaA (Lemma sha-siegel-step-seeds-induced-hilbert-form-exists-eisenstein-lift-finite-iwasawa) and dec; (b) is the private stab M g, mem_stab and stab_fin, with h1 (rat), h2 (pos), h3 (lev) and ray_class_group_finite, the Lean's RayFin, which is step (b) of Lemma sha-siegel-step-seeds-induced-hilbert-form-exists-eisenstein-noncuspidal-characters; (c) is the finiteness of KLev F 2 1 / KLev F 2 M (isCompact_KLev, isOpen_KLev, Subgroup.quotient_finite_of_isOpen), b p = infGL (eps s) * finGL (out p.2) and the union over p of the ranges of q |-> aff q.out 0 * b p; (d) is h1, then h2 (lev), h3 (rot, unip, cent) and e1, with QuotientGroup.mk_out_eq_mul for rho = u h. ctMoves (CtMoves.lean, 87 lines) is (2): hN and hT from CtBasic (eisenstein_constantTerm_of_periodMean_of_upperBlock, with eisenstein_smooth_continuous), cm_arch with CtLeft (eisenstein_continuous_left_action) for the move by a(y, 0) at infinity, which is the argument of ct_arch behind Lemma sha-siegel-step-seeds-induced-hilbert-form-exists-eisenstein-noncuspidal-constant-term(3), cm_scal for the right translations (the text's e), cm_scalar_comm with the field cent, the hypothesis hM for K(M), and rotEig (Lemma sha-siegel-step-seeds-induced-hilbert-form-exists-eisenstein-lift-finite-rotation(1)) with the field wt for the rotations. (3) is ctOne (CoRatAff.lean, 45 lines): main for any Haar measure and fundamental domain, with hper (co_rat_aff), IsAddFundamentalDomain.setIntegral_eq and Measure.isAddLeftInvariant_eq_smul, applied to unipMeasure and unipFundamentalDomain.
\end{proof}

\begin{lemma}\label{sha-siegel-step-seeds-induced-hilbert-form-exists-eisenstein-lift-finite}
Let $E$ be a totally real number field, $k\ge0$, $c$ a function from the primes of $E$ to $\mathbf Z_{\ge0}$ that vanishes at all but finitely many primes, and $\omega$ a Hecke character of $E$ of finite order. Then the holomorphic forms of parallel weight $k$, level $c$ and central character $\omega$ (Definition~\ref{sha-siegel-step-seeds-induced-hilbert-form-exists-holomorphic}) span a finite-dimensional space.
\end{lemma}
\leanref{Development/AutomorphicArithmetic/AutomorphicVarieties/IdeleGenerators/WeightChange.lean}{UpperTheory.V.Sd.HolFin}
\leanref{Development/AutomorphicArithmetic/AutomorphicVarieties/ModularityArgument/ThresholdArguments/DenseThresholdCases/THolFin.lean}{UpperTheory.V.Sd.holomorphic_finite_of_printed_hida}
\begin{proof}
Let $X$ be the set of holomorphic forms of parallel weight $k$, level $c$ and central character $\omega$. It contains $0$, and $\phi+z\phi'\in X$ for $\phi,\phi'\in X$ and $z\in\mathbf C$, so it is a vector space. For a Haar measure $\mu$ on $\mathbf A_E$, a measurable fundamental domain $D$ for the translations by $E$ and $\phi\in X$, let $\psi_\phi(g)=\int_D\phi(n(x)g)\,d\mu(x)$, with $n(x)=\begin{pmatrix}1&x\\0&1\end{pmatrix}$; the integrals converge by Lemma~\ref{sha-siegel-step-seeds-induced-hilbert-form-exists-eisenstein-noncuspidal-constant-term}(1), and $\psi_\phi$ is linear in $\phi$.
\begin{enumerate}
\item[(1)] By Lemma~\ref{sha-siegel-step-seeds-induced-hilbert-form-exists-eisenstein-lift-finite-cusp}, there is $M\ge1$ such that every $\phi\in X$ is fixed by $K(M)$, and the cusp forms in $X$ span a finite-dimensional space $Y$.
\item[(2)] By Lemma~\ref{sha-siegel-step-seeds-induced-hilbert-form-exists-eisenstein-lift-finite-points}(3), choose $\mu$ and $D$ such that a function invariant under left translation by $GL_2(E)$ whose constant term $\int_D\phi(n(x)g)\,d\mu(x)$ vanishes for every $g$ is cuspidal; and by (1) of that lemma, a finite set $S\subset GL_2(\mathbf A_E)$ such that the only set containing $S$ and closed under moves for $M$ is $GL_2(\mathbf A_E)$.
\item[(3)] Let $L:X\to\mathbf C^S$ be the linear map $\phi\mapsto(\psi_\phi(s))_{s\in S}$. If $L(\phi)=0$, the set of $g$ with $\psi_\phi(g)=0$ contains $S$, and it is closed under moves for $M$ by Lemma~\ref{sha-siegel-step-seeds-induced-hilbert-form-exists-eisenstein-lift-finite-points}(2) and (1); so it is $GL_2(\mathbf A_E)$, and $\phi$ is a cusp form by (2). So the kernel of $L$ lies in $Y$.
\end{enumerate}
The kernel of $L$ and the image of $L$ are finite-dimensional, so $X$ is.
\end{proof}
% holomorphic_finite_of_printed_hida takes HypothesesH' but does not use it: it is UpperTheory.V.Br.holFin0 (FinOfKer.lean, 57 lines), followed step by step: X is P = {phi | IsHol k c omega phi} with holLin for h0 and hadd; (1) is holCuspFin (Lemma sha-siegel-step-seeds-induced-hilbert-form-exists-eisenstein-lift-finite-cusp); (2) is ctOne and cuspSet; (3) is ctMoves, with the integrability hint from eisenstein_constantTerm_of_periodMean_of_upperBlock; the conclusion is the private fin_of_ker. Size: about 1,300 lines with its dependencies, net of the definitions and of cuspFormsFiniteReal; the open part left is in Lemma sha-siegel-step-seeds-induced-hilbert-form-exists-eisenstein-lift-finite-harish-chandra.

\begin{lemma}\label{sha-siegel-step-seeds-induced-hilbert-form-exists-eisenstein-lift-hecke-cosets}
Let $E$ be a number field, $w$ a prime of $E$, $\mathcal O_w$ the ring of integers of $E_w$, $\varpi_w$ a uniformizer of $E_w$ and $c\ge0$, and let $K_1(w^c)$ be as in Definition~\ref{sha-siegel-step-seeds-induced-hilbert-form-ordinary}. For $y\in E_w^\times$ and $x\in E_w$ write $a(y,x)=\begin{pmatrix}y&x\\0&1\end{pmatrix}$. Let $S\subset\mathcal O_w$ be a finite set that maps bijectively onto $\mathcal O_w/\varpi_w\mathcal O_w$, $\rho$ a representation of $GL_2(E_w)$ on a complex vector space and $v$ a vector fixed by $K_1(w^c)$. Let $U_wv=\sum_g\rho(g)v$, the sum over representatives $g$ of the cosets $gK_1(w^c)$ contained in $K_1(w^c)\operatorname{diag}(\varpi_w,1)K_1(w^c)$. Then $U_wv$ is fixed by $K_1(w^c)$, and
\[
U_wv=\sum_{u\in S}\rho(a(\varpi_w,u))v+\begin{cases}\rho(\operatorname{diag}(1,\varpi_w))v&\text{if }c=0,\\0&\text{if }c\ne0.\end{cases}
\]
\end{lemma}
\leanref{Development/AutomorphicArithmetic/AutomorphicVarieties/IdeleGenerators/WeightChange.lean}{UpperTheory.V.Sd.UopSum}
\leanref{Development/AutomorphicArithmetic/AutomorphicVarieties/IdeleGenerators/HolomorphicIrreducibility/UopSum.lean}{UpperTheory.V.Sd.uopSum}
\begin{proof}
Let $K=K_1(w^c)$ and $d=a(\varpi_w,0)=\operatorname{diag}(\varpi_w,1)$. The cosets in $KdK$ are the $kdK$ with $k\in K$. Let $k=\begin{pmatrix}\alpha&\beta\\\gamma&\delta\end{pmatrix}\in K$, so that $\gamma$ and $\delta-1$ lie in $w^c\mathcal O_w$.
\begin{enumerate}
\item[(a)] If $\delta$ is a unit, let $u\in S$ with $u\equiv\beta\delta^{-1}$ modulo $\varpi_w$, and write $\beta=\varpi_wt\delta+u\delta$ with $t\in\mathcal O_w$. Then $kd=a(\varpi_w,u)\begin{pmatrix}\alpha-u\gamma&t\delta\\\varpi_w\gamma&\delta\end{pmatrix}$. The second matrix lies in $K$: its entries are integral, its determinant is $\alpha\delta-\beta\gamma$, a unit, and its bottom row is congruent to $(0,1)$ modulo $w^c$.
\item[(b)] If $\delta$ is not a unit, then $c=0$, as $\delta\equiv1$ modulo $w^c$; write $\delta=\varpi_wt$. Then $kd=\operatorname{diag}(1,\varpi_w)\begin{pmatrix}\varpi_w\alpha&\beta\\\gamma&t\end{pmatrix}$, and the second matrix lies in $GL_2(\mathcal O_w)=K$, as its determinant is $\alpha\delta-\beta\gamma$.
\end{enumerate}
Conversely, $a(\varpi_w,u)=a(1,u)d$ lies in $KdK$ for $u\in\mathcal O_w$, and if $c=0$, then $\operatorname{diag}(1,\varpi_w)=sds$ with $s=\begin{pmatrix}0&1\\1&0\end{pmatrix}\in K$. The cosets $a(\varpi_w,u)K$ with $u\in S$ are distinct: if $a(\varpi_w,u)m=a(\varpi_w,u')$ with $m\in K$, the second column gives $m_{22}=1$ and $u'=\varpi_wm_{12}+u$, so $u\equiv u'$ modulo $\varpi_w$. None of them is $\operatorname{diag}(1,\varpi_w)K$: if $a(\varpi_w,u)m=\operatorname{diag}(1,\varpi_w)$, the first column gives $m_{21}=0$ and $\varpi_wm_{11}=1$, which is impossible for integral $m_{11}$. So these are the cosets in $KdK$, which gives the formula. For $k\in K$, left multiplication by $k$ permutes the cosets $gK$ in $KdK$, and $\rho(g)v$ depends only on $gK$, as $v$ is fixed by $K$; so $\rho(k)U_wv=U_wv$.
\end{proof}
% uopSum (UopSum.lean, 258 lines) is the private gen for the subgroup KK of the mirabolic level: (a) and (b) are the two cases of cover, the converse is y_mem and z_mem, the distinctness is yy_res and yz_ne, and the formula and the fixedness come from Hi.dcSum (DcSumEqOfMk.lean) for the finite orbit hfin. The Lean's U_w is Hi.dc rho K K (aff varpi 0), the finite sum over the K-orbit of the coset of diag(varpi, 1) in GL_2(E_w)/K, and its diag(1, varpi) is the scalar varpi times aff(varpi, 0)^-1 (zz_val).

\begin{lemma}\label{sha-siegel-step-seeds-induced-hilbert-form-exists-eisenstein-lift-hecke-translation}
Let $E$ be a totally real number field, $k\ge0$, $c$ a function from the primes of $E$ to $\mathbf Z_{\ge0}$, $\omega$ a Hecke character of $E$, and $\phi$ a holomorphic form of parallel weight $k$, level $c$ and central character $\omega$ (Definition~\ref{sha-siegel-step-seeds-induced-hilbert-form-exists-holomorphic}). Write $a(y,x)=\begin{pmatrix}y&x\\0&1\end{pmatrix}$ over any ring, $\hat{\mathcal O}_E\subset\mathbf A_{E,f}$ for the finite adeles with integral components, and $x-\beta$ for $(x_v-v(\beta))_v$ when $x\in E\otimes\mathbf R$ and $\beta\in E$.
\begin{enumerate}
\item If $o\in\hat{\mathcal O}_E$, then $\phi\bigl(g\,a(1,o)_f\bigr)=\phi(g)$ for every $g\in GL_2(\mathbf A_E)$.
\item Let $y\in(E\otimes\mathbf R)^\times$, $x\in E\otimes\mathbf R$, $y_f\in\mathbf A_{E,f}^\times$, $t\in\mathbf A_{E,f}$ and $\beta\in E$ with $t-\beta\in y_f\hat{\mathcal O}_E$. Then $\phi\bigl(a(y,x)_\infty a(y_f,t)_f\bigr)=\phi\bigl(a(y,x-\beta)_\infty a(y_f,0)_f\bigr)$.
\end{enumerate}
\end{lemma}
\leanref{Development/AutomorphicArithmetic/AutomorphicVarieties/BuildingsResidueChars/HilbertHolomorphicInvariantIntegralTranslation.lean}{UpperTheory.hilbertHolomorphic_invariant_integralTranslation}
\leanref{Development/AutomorphicArithmetic/AutomorphicVarieties/IdeleGenerators/HolomorphicIrreducibility/ExistsIdele.lean}{UpperTheory.V.Sd.shift}
\begin{proof}
(1) The stabilizer of $\phi$ in $GL_2(\mathbf A_{E,f})$ is open, so there is a finite set $T$ of primes such that it contains every element of $GL_2(\hat{\mathcal O}_E)$ whose components at the primes of $T$ are $1$. Let $o'$ have components $0$ at the primes of $T$ and $o_v$ elsewhere. Then $a(1,o)=a(1,o')\prod_{w\in T}a(1,o_w)_w$. Each $a(1,o_w)$ lies in $K_1(w^{c(w)})$, which fixes $\phi$, and $a(1,o')$ lies in the stabilizer.

(2) Let $\gamma=a(1,\beta)\in GL_2(E)$. As $a(y,x)a(z,t')=a(yz,x+yt')$,
\[
a(y,x)_\infty a(y_f,t)_f=\gamma\,a(y,x-\beta)_\infty a(y_f,0)_f\,a\bigl(1,y_f^{-1}(t-\beta)\bigr)_f.
\]
The last factor has $y_f^{-1}(t-\beta)\in\hat{\mathcal O}_E$. So the claim follows from the invariance of $\phi$ under left translation by $GL_2(E)$ and from (1).
\end{proof}
% hilbertHolomorphic_invariant_integralTranslation (68 lines) is (1): exists_KAway_subset for T, and the induction key over T with the factorization e and the membership hK in K1 w (c w). shift (private, ExistsIdele.lean) is (2): h2 for the factorization, h.aut.left and (1).

\begin{lemma}\label{sha-siegel-step-seeds-induced-hilbert-form-exists-eisenstein-lift-hecke-characters}
Let $E$ be a totally real number field with ring of integers $\mathcal O_E$ and different $\mathfrak d$. For $x\in E$ write $\operatorname{Tr}(x)=\sum_vv(x)$, the sum over the infinite places, and $e(t)=\exp(2\pi it)$ for $t\in\mathbf R$. Let $w$ be a prime of $E$, $\varpi_w$ a uniformizer of $E_w$, and $s(r)\in\mathcal O_E$ a representative of each $r\in\mathcal O_E/w$.
\begin{enumerate}
\item For every $y\in\mathbf A_{E,f}^\times$ there is $\beta\in E$ with $\operatorname{ord}_w(\beta)=\operatorname{ord}_w(y_w)$ and $\operatorname{ord}_v(\beta)\ge\operatorname{ord}_v(y_v)$ at every prime $v$.
\item If $t$ is a unit of $\mathcal O_{E_w}$, the map $r\mapsto ts(r)$ induces a bijection from $\mathcal O_E/w$ onto $\mathcal O_{E_w}/\varpi_w\mathcal O_{E_w}$.
\item Let $\eta\in E$ with $\operatorname{ord}_v(\eta)\ge-\operatorname{ord}_v(\mathfrak d)$ at every prime $v$. Then $e(\operatorname{Tr}(\eta b))=1$ for every $b\in\mathcal O_E$.
\item Let $\eta\in E$ with $\operatorname{ord}_v(\eta)\ge-\operatorname{ord}_v(\mathfrak d)$ at every prime $v\ne w$ and $\operatorname{ord}_w(\eta)=-\operatorname{ord}_w(\mathfrak d)-1$. Then $\sum_{r\in\mathcal O_E/w}e(\operatorname{Tr}(\eta s(r)))=0$.
\end{enumerate}
\end{lemma}
\leanref{Development/AutomorphicArithmetic/AutomorphicVarieties/BuildingsResidueChars/ExistsGlobalElementValuationEqIdeleAtPlace.lean}{UpperTheory.exists_global_element_valuation_eq_idele_at_place}
\leanref{Development/AutomorphicArithmetic/AutomorphicVarieties/BuildingsResidueChars/BijectiveResidueUnitMulLift.lean}{UpperTheory.bijective_residue_unit_mul_lift}
\leanref{Development/AutomorphicArithmetic/AutomorphicVarieties/BuildingsResidueChars/AddCharInverseDifferentAndResidueSumEqZero.lean}{UpperTheory.addChar_inverseDifferent_and_residue_sum_eq_zero}
\begin{proof}
(1) Let $p\in w$ be nonzero. As $E$ is dense in $\mathbf A_{E,f}$, there is $\beta\in E$ with $\beta-y\in yp\hat{\mathcal O}_E$. Then $\operatorname{ord}_v(\beta-y_v)\ge\operatorname{ord}_v(y_v)+\operatorname{ord}_v(p)$ at every prime $v$. So $\operatorname{ord}_v(\beta)\ge\operatorname{ord}_v(y_v)$, and at $w$, where $\operatorname{ord}_w(p)\ge1$, $\operatorname{ord}_w(\beta)=\operatorname{ord}_w(y_w)$.

(2) The map $\mathcal O_E/w\to\mathcal O_{E_w}/\varpi_w\mathcal O_{E_w}$ induced by the inclusion is a homomorphism of fields, so injective; it is surjective, as $E$ is dense in $E_w$ and an element of $E$ close to an integer of $E_w$ is congruent modulo $w$ to an element of $\mathcal O_E$. Multiplication by the unit $t$ is a bijection of the residue field.

(3) The inverse different $\mathfrak d^{-1}$ is the set of $x\in E$ with $\operatorname{ord}_v(x)\ge-\operatorname{ord}_v(\mathfrak d)$ at every prime $v$, and also the set of $x$ with $\operatorname{Tr}_{E/\mathbf Q}(x\mathcal O_E)\subseteq\mathbf Z$. As $E$ is totally real, $\operatorname{Tr}_{E/\mathbf Q}$ is the sum of the real embeddings. So $\operatorname{Tr}(\eta b)\in\mathbf Z$, and $e(\operatorname{Tr}(\eta b))=1$.

(4) The function $\psi(b)=e(\operatorname{Tr}(\eta b))$ is a homomorphism from the additive group $\mathcal O_E$ to $\mathbf C^\times$. For $b\in w$, $\eta b$ satisfies the hypothesis of (3), as $\operatorname{ord}_w(b)\ge1$; so $\psi(b)=e(\operatorname{Tr}(\eta b\cdot1))=1$ by (3), and $\psi$ factors through $\mathcal O_E/w$. As $\operatorname{ord}_w(\eta)<-\operatorname{ord}_w(\mathfrak d)$, $\eta\notin\mathfrak d^{-1}$, so $\operatorname{Tr}(\eta b_0)\notin\mathbf Z$ and $\psi(b_0)\ne1$ for some $b_0\in\mathcal O_E$. Translating $r$ by the class of $b_0$ gives $\psi(b_0)\sum_r\psi(s(r))=\sum_r\psi(s(r))$, so the sum vanishes.
\end{proof}
% exists_global_element_valuation_eq_idele_at_place (ExistsGlobalElementValuationEqIdeleAtPlace.lean, 45 lines): denseRange_algebraMap_finiteAdeleRing with the open set of integral finite adeles, z = p^-1 y^-1 and key. bijective_residue_unit_mul_lift (60 lines): the lift g of the residue map through O_E/w, injective as a ring map of fields, surjective by denseRange_algebraMap and exists_valuation_sub_lt_of_integer. addChar_inverseDifferent_and_residue_sum_eq_zero (195 lines with its private lemmas): (3) is part1 (mem_dual_of_valuation_le_different, mem_dual_one_iff_trace_mul_int, sum_reAt_eq_trace' and cexp_two_pi_I_mul_eq_one_iff_int); (4) is hadd, hker, hcoset, b0 (false_of_mem_dual_of_valuation_eq_different_succ) and hsum. The Lean writes the conditions on valuations multiplicatively: v(eta) <= exp(ord_v d) is the text's ord_v(eta) >= -ord_v(d).

\begin{lemma}\label{sha-siegel-step-seeds-induced-hilbert-form-exists-eisenstein-lift-hecke-values}
Let $E$ be a totally real number field with uniformizers $\varpi_w$, $k\ge0$, $c$ a function from the primes of $E$ to $\mathbf Z_{\ge0}$, $\omega$ a Hecke character of $E$, and $\phi$ a holomorphic form of parallel weight $k$, level $c$ and central character $\omega$ (Definition~\ref{sha-siegel-step-seeds-induced-hilbert-form-exists-holomorphic}). Let $w$ be a prime of $E$, $h_w=\operatorname N(w)^{k/2-1}U_w$ with $U_w$ acting through right translation at $w$, and $s(r)\in\mathcal O_E$ a representative of each $r\in\mathcal O_E/w$. Use the notation of Lemma~\ref{sha-siegel-step-seeds-induced-hilbert-form-exists-eisenstein-lift-hecke-translation} and the function $\chi_{c,\omega}$ of Definition~\ref{sha-siegel-step-seeds-induced-hilbert-form-exists-eisenstein-expansion}. Then there is an idele $u\in\mathbf A_{E,f}^\times$ that generates $w^{-1}$, that is, with $\operatorname{ord}_w(u_w)=-1$ and $u_v$ a unit at every prime $v\ne w$, such that the following holds. Let $y\in(E\otimes\mathbf R)^\times$, $x\in E\otimes\mathbf R$, $y_f\in\mathbf A_{E,f}^\times$ and $\beta\in E$ with $\operatorname{ord}_w(\beta)=\operatorname{ord}_w(y_{f,w})$ and $\operatorname{ord}_v(\beta)\ge\operatorname{ord}_v(y_{f,v})$ at every prime $v$. Then
\[
(h_w\phi)\bigl(a(y,x)_\infty a(y_f,0)_f\bigr)=\operatorname N(w)^{k/2-1}\Bigl(\sum_{r\in\mathcal O_E/w}\phi\bigl(a(y,x-\beta s(r))_\infty a(y_fu^{-1},0)_f\bigr)+\chi_{c,\omega}(w)\,\phi\bigl(a(y,x)_\infty a(y_fu,0)_f\bigr)\Bigr).
\]
\end{lemma}
\leanref{Development/AutomorphicArithmetic/AutomorphicVarieties/IdeleGenerators/HolomorphicIrreducibility/ExistsIdele.lean}{UpperTheory.V.Sd.hk_val}
\begin{proof}
Let $p$ be the idele with component $\varpi_w$ at $w$ and $1$ elsewhere, and $u=p^{-1}$. Let $g=a(y,x)_\infty a(y_f,0)_f$ and $t=\beta/y_{f,w}$, a unit of $\mathcal O_{E_w}$. By Lemma~\ref{sha-siegel-step-seeds-induced-hilbert-form-exists-eisenstein-lift-hecke-characters}(2), the $ts(r)$ map bijectively onto $\mathcal O_{E_w}/\varpi_w\mathcal O_{E_w}$. As $\phi$ is fixed by $K_1(w^{c(w)})$, Lemma~\ref{sha-siegel-step-seeds-induced-hilbert-form-exists-eisenstein-lift-hecke-cosets}, applied to right translation at $w$ with these representatives, gives
\[
(h_w\phi)(g)=\operatorname N(w)^{k/2-1}\Bigl(\sum_r\phi\bigl(g\,a(\varpi_w,ts(r))_w\bigr)+\epsilon\,\phi\bigl(g\operatorname{diag}(1,\varpi_w)_w\bigr)\Bigr),
\]
with $\epsilon=1$ if $c(w)=0$ and $\epsilon=0$ otherwise.
\begin{enumerate}
\item[(a)] The term of $r$. In $GL_2(\mathbf A_{E,f})$, $a(y_f,0)\,a(\varpi_w,ts(r))_w=a(y_fp,\tau)$, where $\tau\in\mathbf A_{E,f}$ has component $y_{f,w}ts(r)=\beta s(r)$ at $w$ and $0$ elsewhere. The finite adele $\tau-\beta s(r)$ has component $0$ at $w$, and $-\beta s(r)$ at a prime $v\ne w$, where $\operatorname{ord}_v(\beta s(r))\ge\operatorname{ord}_v(\beta)\ge\operatorname{ord}_v(y_{f,v})=\operatorname{ord}_v((y_fp)_v)$. So $\tau-\beta s(r)\in y_fp\hat{\mathcal O}_E$, and by Lemma~\ref{sha-siegel-step-seeds-induced-hilbert-form-exists-eisenstein-lift-hecke-translation}(2) the term is $\phi\bigl(a(y,x-\beta s(r))_\infty a(y_fu^{-1},0)_f\bigr)$.
\item[(b)] The last term. If $c(w)\ne0$, then $\epsilon=0=\chi_{c,\omega}(w)$. If $c(w)=0$, then $\operatorname{diag}(1,\varpi_w)=\varpi_w\,a(\varpi_w^{-1},0)$, so $g\operatorname{diag}(1,\varpi_w)_w=p\,a(y,x)_\infty a(y_fp^{-1},0)_f$ with $p$ a scalar matrix. By (c) of Definition~\ref{sha-siegel-step-seeds-induced-hilbert-form-exists-holomorphic}, the term is $\omega(p)\phi\bigl(a(y,x)_\infty a(y_fu,0)_f\bigr)$, and $\omega(p)=\omega(\varpi_w)=\chi_{c,\omega}(w)$.
\end{enumerate}
\end{proof}
% hk_val (private, ExistsIdele.lean, 60 lines): p from exists_idele w (atFin w (unif w)) and u = p^-1; t0 is the text's t and hb, hS its representatives (bijective_residue_unit_mul_lift); the formula is uopSum with the image of f; (a) is hp3, aff_mul and shift; (b) is fin_scalar, scalar_comm, h.cent, aff_inv and chi_prime.

\begin{lemma}\label{sha-siegel-step-seeds-induced-hilbert-form-exists-eisenstein-lift-hecke-coefficients}
Let $E$ be a totally real number field with different $\mathfrak d$, $w$ a prime of $E$, $s(r)\in\mathcal O_E$ a representative of each $r\in\mathcal O_E/w$, and $u\in\mathbf A_{E,f}^\times$ an idele that generates $w^{-1}$. Use the notation $\operatorname{Tr}$ and $e$ of Lemma~\ref{sha-siegel-step-seeds-induced-hilbert-form-exists-eisenstein-lift-hecke-characters}. Let $a$ be a function on the nonzero integral ideals of $E$, extended to $\mathbf A_{E,f}$ by the rule of Definition~\ref{sha-siegel-step-seeds-induced-hilbert-form-exists-eisenstein-expansion}(1): $a(z)=a(\mathfrak m)$ if $z$ generates $\mathfrak m\mathfrak d^{-1}$ with $\mathfrak m$ a nonzero integral ideal, and $a(z)=0$ if there is no such $\mathfrak m$; let $a_w$ be the function $\mathfrak m\mapsto a(w\mathfrak m)$, extended in the same way. Let $z\in\mathbf A_{E,f}$ and $\eta\in E$ with $\operatorname{ord}_w(\eta)=\operatorname{ord}_w(z_w)$ and $\operatorname{ord}_v(\eta)\ge\operatorname{ord}_v(z_v)$ at every prime $v$. Then
\[
a(zu^{-1})\sum_{r\in\mathcal O_E/w}e\bigl(\operatorname{Tr}(\eta s(r))\bigr)=\operatorname N(w)\,a_w(z).
\]
\end{lemma}
\leanref{Development/AutomorphicArithmetic/AutomorphicVarieties/IdeleGenerators/HolomorphicIrreducibility/ExistsIdele.lean}{UpperTheory.V.Sd.char_sum}
\begin{proof}
For an integral ideal $\mathfrak m$, $z$ generates $\mathfrak m\mathfrak d^{-1}$ if and only if $zu^{-1}$ generates $w\mathfrak m\mathfrak d^{-1}$, as $u^{-1}$ generates $w$.
\begin{enumerate}
\item[(a)] Suppose that $z$ generates $\mathfrak m\mathfrak d^{-1}$ for a nonzero integral ideal $\mathfrak m$. Then $a(zu^{-1})=a(w\mathfrak m)=a_w(z)$. At every prime $v$, $\operatorname{ord}_v(\eta)\ge\operatorname{ord}_v(z_v)=\operatorname{ord}_v(\mathfrak m)-\operatorname{ord}_v(\mathfrak d)\ge-\operatorname{ord}_v(\mathfrak d)$. By Lemma~\ref{sha-siegel-step-seeds-induced-hilbert-form-exists-eisenstein-lift-hecke-characters}(3), every term of the sum is $1$, so the sum is the number $\operatorname N(w)$ of residues.
\item[(b)] Otherwise $a_w(z)=0$. If $zu^{-1}$ generates no $\mathfrak m\mathfrak d^{-1}$, then $a(zu^{-1})=0$. Suppose that $zu^{-1}$ generates $\mathfrak m\mathfrak d^{-1}$. Then $w$ does not divide $\mathfrak m$: otherwise $\mathfrak m=w\mathfrak m'$ and $z$ generates $\mathfrak m'\mathfrak d^{-1}$. So $\operatorname{ord}_v(z_v)=\operatorname{ord}_v(\mathfrak m)-\operatorname{ord}_v(\mathfrak d)-\operatorname{ord}_v(w)$ is at least $-\operatorname{ord}_v(\mathfrak d)$ at $v\ne w$, and is $-\operatorname{ord}_w(\mathfrak d)-1$ at $w$. As $\operatorname{ord}_v(\eta)\ge\operatorname{ord}_v(z_v)$ and $\operatorname{ord}_w(\eta)=\operatorname{ord}_w(z_w)$, Lemma~\ref{sha-siegel-step-seeds-induced-hilbert-form-exists-eisenstein-lift-hecke-characters}(4) shows that the sum vanishes.
\end{enumerate}
\end{proof}
% char_sum (private, ExistsIdele.lean, 33 lines): sh for the first sentence (genD_shift), the case h for (a) with c1, the cases h' and not h' for (b) with h0, hz and c2. The Lean's z is a finite adele and its u any unit with Valued.v (u_v) = exp(ord_v w), the text's u generating w^-1.

\begin{lemma}\label{sha-siegel-step-seeds-induced-hilbert-form-exists-eisenstein-lift-hecke}
Let $E$ be a totally real number field with uniformizers $\varpi_w$, $k\ge1$, $c$ a function from the primes of $E$ to $\mathbf Z_{\ge0}$, $\omega$ a Hecke character of $E$ of finite order (the finiteness of its order is not used), and $\phi$ a holomorphic form of parallel weight $k$, level $c$ and central character $\omega$. For a prime $w$ let $h_w=\operatorname N(w)^{k/2-1}U_w$, with $U_w$ acting on the vectors fixed by $K_1(w^{c(w)})$ through right translation at $w$. Then for all primes $w,w'$:
\begin{enumerate}
\item $h_w\phi$ is a holomorphic form of parallel weight $k$, level $c$ and central character $\omega$, and $h_wh_{w'}\phi=h_{w'}h_w\phi$;
\item if $\phi$ has the Fourier expansion of weight $k$ with constant term $c_0$ and coefficients $a$, then $h_w\phi$ has the Fourier expansion of weight $k$ with some constant term and coefficients $T^k_wa$, for $c$ and $\omega$ (Definition~\ref{sha-siegel-step-seeds-induced-hilbert-form-exists-eisenstein-expansion}).
\end{enumerate}
\end{lemma}
\leanref{Development/AutomorphicArithmetic/AutomorphicVarieties/IdeleGenerators/WeightChange.lean}{UpperTheory.V.Sd.HolHk}
\leanref{Development/AutomorphicArithmetic/AutomorphicVarieties/IdeleGenerators/HolomorphicIrreducibility/ExistsIdele.lean}{UpperTheory.V.Sd.holHk}
\begin{proof}
Write $R_w(g)$ for right translation by $g\in GL_2(E_w)$ at $w$, and $a(y,x)=\begin{pmatrix}y&x\\0&1\end{pmatrix}$.
\begin{enumerate}
\item[(1)] \emph{The operator.} Fix a prime $w$ and a finite set $S_w\subset\mathcal O_{E_w}$ that maps bijectively onto $\mathcal O_{E_w}/\varpi_w\mathcal O_{E_w}$, and let
\[
A_w=\operatorname N(w)^{k/2-1}\Bigl(\sum_{u\in S_w}R_w(a(\varpi_w,u))+\epsilon_wR_w(\operatorname{diag}(1,\varpi_w))\Bigr),
\]
with $\epsilon_w=1$ if $c(w)=0$ and $\epsilon_w=0$ otherwise. By Lemma~\ref{sha-siegel-step-seeds-induced-hilbert-form-exists-eisenstein-lift-hecke-cosets}, applied to right translation at $w$, if $\psi$ is fixed by $K_1(w^{c(w)})$, then $h_w\psi=A_w\psi$, and $h_w\psi$ is fixed by $K_1(w^{c(w)})$. So if $B$ is a linear operator that commutes with every $R_w(g)$, then $B\psi$ is also fixed by $K_1(w^{c(w)})$ and $Bh_w\psi=BA_w\psi=A_wB\psi=h_wB\psi$.
\item[(2)] \emph{Holomorphy.} By hypothesis $\phi$ is fixed by $K_1(w^{c(w)})$, so $h_w\phi=A_w\phi$. The conditions of Definition~\ref{sha-siegel-step-seeds-induced-hilbert-form-exists-holomorphic} hold for $h_w\phi$:
\begin{enumerate}
\item[(a)] Right translation by an element of $GL_2(\mathbf A_{E,f})$ maps automorphic forms to automorphic forms, as in the proof of Lemma~\ref{sha-siegel-step-seeds-induced-hilbert-form-exists-eisenstein-raise}, and so do linear combinations; so $h_w\phi$ is an automorphic form.
\item[(b)] At $w$, $h_w\phi$ is fixed by $K_1(w^{c(w)})$ by (1). At a prime $v\ne w$, right translations at $v$ commute with those at $w$, so by (1), $R_v(g)h_w\phi=h_wR_v(g)\phi=h_w\phi$ for $g\in K_1(v^{c(v)})$.
\item[(c)] Left multiplication by a scalar matrix commutes with right translation, so each $R_w(g)\phi$ satisfies condition (c), and so does $A_w\phi$.
\item[(d)] The derivatives $X^v_{ab}$ commute with right translation by elements of $GL_2(\mathbf A_{E,f})$, so by (1), $X^v_{ab}h_w\phi=h_wX^v_{ab}\phi$. Applying $h_w$ to the identities (d) for $\phi$ gives them for $h_w\phi$.
\end{enumerate}
\item[(3)] \emph{Commutation.} If $w=w'$ there is nothing to prove. Otherwise $\phi$ and $h_{w'}\phi$ are fixed by $K_1(w^{c(w)})$, by (2)(b) for $w'$, so $h_w=A_w$ on them. The operator $A_w$ is a combination of right translations at $w$, which commute with right translations at $w'$. By (1) for $w'$, $h_wh_{w'}\phi=A_wh_{w'}\phi=h_{w'}A_w\phi=h_{w'}h_w\phi$.
\item[(4)] \emph{Fourier coefficients.} Suppose that $\phi$ has the Fourier expansion of weight $k$ with constant term $c_0$ and coefficients $a$. Let $s(r)\in\mathcal O_E$ be a representative of each $r\in\mathcal O_E/w$, and $u$ the idele of Lemma~\ref{sha-siegel-step-seeds-induced-hilbert-form-exists-eisenstein-lift-hecke-values}, which generates $w^{-1}$; so $|u|=\operatorname N(w)$ and $|u^{-1}|=\operatorname N(w)^{-1}$. Write $\chi=\chi_{c,\omega}$, $q=\operatorname N(w)$ and
\[
c'_0(y_f)=c_0(y_fu^{-1})+\chi(w)q^{k-1}c_0(y_fu).
\]
Let $x\in E\otimes\mathbf R$, $y\in E\otimes\mathbf R$ totally positive and $y_f\in\mathbf A_{E,f}^\times$. By Lemma~\ref{sha-siegel-step-seeds-induced-hilbert-form-exists-eisenstein-lift-hecke-characters}(1) there is $\beta\in E$ with $\operatorname{ord}_w(\beta)=\operatorname{ord}_w(y_{f,w})$ and $\operatorname{ord}_v(\beta)\ge\operatorname{ord}_v(y_{f,v})$ at every prime $v$. Take the expansions of $\phi$ at $(x-\beta s(r),y,y_fu^{-1})$ for $r\in\mathcal O_E/w$, multiplied by $q^{-1}$, and at $(x,y,y_fu)$, multiplied by $\chi(w)q^{k-1}$, and add them.
\begin{enumerate}
\item[(a)] \emph{The terms.} For a totally positive $\xi\in E$ let $Q(\xi)=\exp\bigl(2\pi i\sum_vv(\xi)(x_v+iy_v)\bigr)$. Replacing $x$ by $x-\beta s(r)$ multiplies $Q(\xi)$ by $e(\operatorname{Tr}(-\xi\beta s(r)))$, in the notation of Lemma~\ref{sha-siegel-step-seeds-induced-hilbert-form-exists-eisenstein-lift-hecke-characters}. So the term of $\xi$ is $Q(\xi)$ times
\[
q^{-1}a(\xi y_fu^{-1})\sum_{r\in\mathcal O_E/w}e\bigl(\operatorname{Tr}(-\xi\beta s(r))\bigr)+\chi(w)q^{k-1}a(\xi y_fu).
\]
Let $z=\xi y_f$ and $\eta=-\xi\beta$. Then $\operatorname{ord}_w(\eta)=\operatorname{ord}_w(z_w)$ and $\operatorname{ord}_v(\eta)\ge\operatorname{ord}_v(z_v)$ at every prime $v$, so by Lemma~\ref{sha-siegel-step-seeds-induced-hilbert-form-exists-eisenstein-lift-hecke-coefficients} the first part is $a_w(z)$, with $a_w(\mathfrak m)=a(w\mathfrak m)$. As in the proof of Lemma~\ref{sha-siegel-step-seeds-induced-hilbert-form-exists-eisenstein-raise}, $zu$ generates $\mathfrak m'\mathfrak d^{-1}$ with $\mathfrak m'$ integral if and only if $z$ generates $\mathfrak m'w\mathfrak d^{-1}$; so $a(zu)$ is the value at $z$ of $\mathfrak m\mapsto a(\mathfrak mw^{-1})$. So the term of $\xi$ is $(T^k_wa)(\xi y_f)Q(\xi)$.
\item[(b)] \emph{The sum.} By the expansions of $\phi$, the family of these terms is summable with sum
\[
q^{-1}\sum_{r}\Bigl(\phi\bigl(a(y,x-\beta s(r))_\infty a(y_fu^{-1},0)_f\bigr)\bigl(\operatorname N(y)|y_f|q^{-1}\bigr)^{-k/2}-c_0(y_fu^{-1})\Bigr)+\chi(w)q^{k-1}\Bigl(\phi\bigl(a(y,x)_\infty a(y_fu,0)_f\bigr)\bigl(\operatorname N(y)|y_f|q\bigr)^{-k/2}-c_0(y_fu)\Bigr).
\]
There are $q$ residues $r$, and $q^{-1}q^{k/2}=q^{k-1}q^{-k/2}=q^{k/2-1}$, as $k\ge1$. So this sum is $\bigl(\operatorname N(y)|y_f|\bigr)^{-k/2}$ times $q^{k/2-1}\bigl(\sum_r\phi(\ldots)+\chi(w)\phi(\ldots)\bigr)$, minus $c'_0(y_f)$; by Lemma~\ref{sha-siegel-step-seeds-induced-hilbert-form-exists-eisenstein-lift-hecke-values}, it is $(h_w\phi)\bigl(a(y,x)_\infty a(y_f,0)_f\bigr)\bigl(\operatorname N(y)|y_f|\bigr)^{-k/2}-c'_0(y_f)$.
\end{enumerate}
So $h_w\phi$ has the Fourier expansion of weight $k$ with constant term $c'_0$ and coefficients $T^k_wa$.
\end{enumerate}
\end{proof}
% holHk (ExistsIdele.lean, 58 lines): (1) is hA and hk_comm with exists_reps and uopSum; (2) is hk_hol, with hecke_operators_commute_on_automorphic_forms for (a), locSm_comm for (b), the submodule W of functions with central character omega for (c) and eop_comm (TR.lie_rt_fin) for (d); (3) is hk_hk; (4) is the last part of holHk: hk_val with u = p^-1, exists_global_element_valuation_eq_idele_at_place for beta, H for the combined expansion, char_sum, coefAt_shift and coefAt_add for (a), and absF_of, absF_mul, q1 and q2 for (b), where q2 uses k >= 1. The hypothesis IsOfFinOrder omega is not used. About 750 lines with the lemmas before this one.

\begin{lemma}\label{sha-siegel-step-seeds-induced-hilbert-form-exists-eisenstein-lift-unique-exponentials}
Let $V$ be a finite set, $I$ a set, $\mu:I\to\mathbf R^V$ and $B:I\to\mathbf C$ functions, and $d\in\mathbf C$. For $m,x\in\mathbf R^V$ write $e(m\cdot x)=\exp\bigl(2\pi i\sum_{v\in V}m_vx_v\bigr)$.
\begin{enumerate}
\item Suppose that $\mu_i\ne0$ whenever $B_i\ne0$, and that for every $x\in\mathbf R^V$ the family $(B_ie(\mu_i\cdot x))_{i\in I}$ is summable with sum $d$. Then $d=0$.
\item Suppose that $\mu$ is injective, and that for every $x\in\mathbf R^V$ the family $(B_ie(\mu_i\cdot x))_{i\in I}$ is summable with sum $0$. Then $B_i=0$ for every $i\in I$.
\end{enumerate}
\end{lemma}
\leanref{Development/AutomorphicArithmetic/AutomorphicVarieties/IdeleGenerators/HolomorphicIrreducibility/ExpUnique.lean}{UpperTheory.V.Sd.coef_zero}
\leanref{Development/AutomorphicArithmetic/AutomorphicVarieties/IdeleGenerators/HolomorphicIrreducibility/ExpUnique.lean}{UpperTheory.V.Sd.kill}
\begin{proof}
(1) At $x=0$ the family $(B_i)$ is summable, so $\sum_i|B_i|$ is finite. We show by induction on a finite set $S\subseteq I$ that there is a function $C:I\to\mathbf C$ with $|C_i|\le|B_i|$ for every $i$, $C_i=0$ for $i\in S$, and $\sum_iC_ie(\mu_i\cdot x)=d$ for every $x$. For $S$ empty, take $C=B$. Let $C$ serve for $S$ and let $j\notin S$. If $B_j=0$, then $C_j=0$ and $C$ serves for $S\cup\{j\}$. Otherwise $\mu_j\ne0$; choose $v$ with $\mu_{j,v}\ne0$, and let $g\in\mathbf R^V$ have $g_v=1/(2\mu_{j,v})$ and $g_{v'}=0$ for $v'\ne v$, so that $e(\mu_j\cdot g)=-1$. Let $C'_i=C_i(1+e(\mu_i\cdot g))/2$. Then $|C'_i|\le|C_i|$, $C'_j=0$, and $C'_i=0$ for $i\in S$. As $e(\mu_i\cdot(x+g))=e(\mu_i\cdot x)e(\mu_i\cdot g)$, the average of the identities at $x$ and at $x+g$ is $\sum_iC'_ie(\mu_i\cdot x)=d$. So $C'$ serves for $S\cup\{j\}$. For each finite $S$, the identity at $x=0$ gives $d=\sum_{i\notin S}C_i$, so $|d|\le\sum_{i\notin S}|B_i|$, which tends to $0$ as $S$ grows. So $d=0$.

(2) Let $j\in I$, $\mu'_i=\mu_i-\mu_j$, $B'_i=B_i$ for $i\ne j$ and $B'_j=0$. Multiplying the identity by $e(-\mu_j\cdot x)$ and subtracting $B_j$ gives $\sum_iB'_ie(\mu'_i\cdot x)=-B_j$ for every $x$. As $\mu$ is injective, $\mu'_i\ne0$ for $i\ne j$. By (1), $B_j=0$.
\end{proof}
% kill, const_zero and coef_zero (private, ExpUnique.lean, about 100 lines with ee, exists_ee and norm_ee): (1) is const_zero with kill for the induction, tendsto_tsum_compl_atTop_zero for the tails and HasSum.norm_le_of_bounded; (2) is coef_zero.

\begin{lemma}\label{sha-siegel-step-seeds-induced-hilbert-form-exists-eisenstein-lift-unique-rotation}
Let $E$ be a totally real number field and $k\ge0$, and let $\phi$ be a function on $GL_2(\mathbf A_E)$, smooth in the variable at infinity, with $(X^v_{12}-X^v_{21})\phi=ik\phi$ at every infinite place $v$, for the derivatives $X^v_{ab}$ of Definition~\ref{sha-siegel-step-seeds-induced-hilbert-form-exists-holomorphic}. For $\theta\in E\otimes\mathbf R$ let $r(\theta)\in GL_2(E\otimes\mathbf R)$ have component $\begin{pmatrix}\cos\theta_v&\sin\theta_v\\-\sin\theta_v&\cos\theta_v\end{pmatrix}$ at each infinite place $v$. Then $|\phi(g\,r(\theta))|=|\phi(g)|$ for every $g\in GL_2(\mathbf A_E)$ and $\theta\in E\otimes\mathbf R$.
\end{lemma}
\leanref{Development/AutomorphicArithmetic/AutomorphicVarieties/IdeleGenerators/HolomorphicIrreducibility/ExpInj.lean}{UpperTheory.V.Sd.norm_rot}
\begin{proof}
Let $f(t)=\phi(g\,r(t\theta))$ for $t\in\mathbf R$, and $\sigma=ik\sum_v\theta_v$. As $r((t+s)\theta)=r(t\theta)r(s\theta)$, and the derivative of $s\mapsto r(s\theta)$ at $s=0$ is $\theta_v(e_{12}-e_{21})$ at each $v$, the chain rule gives $f'(t)=\sum_v\theta_v(X^v_{12}-X^v_{21})\phi(g\,r(t\theta))=\sigma f(t)$. So $f(t)\exp(-\sigma t)$ has derivative $0$, and is constant. So $\phi(g\,r(\theta))=f(1)=\exp(\sigma)f(0)=\exp(\sigma)\phi(g)$, and $|\exp(\sigma)|=1$ as $\sigma$ is imaginary.
\end{proof}
% norm_rot (private, ExpInj.lean, 27 lines): rotA for r, rotA_add, rotA_deriv, lie_rot for the derivative, TR.hasDerivAt_curve for the chain rule, and is_const_of_deriv_eq_zero.

\begin{lemma}\label{sha-siegel-step-seeds-induced-hilbert-form-exists-eisenstein-lift-unique-approximation}
Let $E$ be a totally real number field and $K$ an open subgroup of $GL_2(\mathbf A_{E,f})$. For $y\in\mathbf A_{E,f}^\times$ write $a(y,0)=\operatorname{diag}(y,1)$. Then every $g\in GL_2(\mathbf A_E)$ is $\gamma\,h_\infty\bigl(a(y,0)u\bigr)_f$ for some $\gamma\in GL_2(E)$, some $h\in GL_2(E\otimes\mathbf R)$ with totally positive determinant, some $y\in\mathbf A_{E,f}^\times$ and some $u\in K$.
\end{lemma}
\leanref{Development/AutomorphicArithmetic/AutomorphicVarieties/IdeleGenerators/WeightChange.lean}{UpperTheory.V.Sd.StrongApp}
\leanref{Development/AutomorphicArithmetic/AutomorphicVarieties/IdeleGenerators/HolomorphicIrreducibility/StrongApp.lean}{UpperTheory.V.Sd.strongApp}
\leanref{GroundTheory/AutomorphicForms/StrongApproximationSLDense.lean}{GroundTheory.Automorphic.AdelicGLn.finiteSL_le_topologicalClosure}
\begin{proof}
\begin{enumerate}
\item[(a)] \emph{Strong approximation.} $SL_2(E)$ is dense in $SL_2(\mathbf A_{E,f})$. Indeed, $E$ is dense in $\mathbf A_{E,f}$, so every elementary matrix $1+te_{ij}$ with $t\in\mathbf A_{E,f}$ and $i\ne j$ is a limit of elementary matrices with entries in $E$. An element $g'$ of $SL_2(\mathbf A_{E,f})$ is the limit of its truncations, which agree with $g'$ at the primes of a finite set $S$ and are $1$ elsewhere; each truncation is a product of elements of $SL_2(E_v)$ placed at single primes $v\in S$, and each of these is a product of elementary matrices over the field $E_v$, which are elementary matrices over $\mathbf A_{E,f}$.
\item[(b)] \emph{The infinite part.} As $E$ is dense in $E\otimes\mathbf R$, there is $\alpha\in E^\times$ with $\alpha\det g_\infty$ totally positive.
\item[(c)] \emph{The finite part.} Let $g'=a(\alpha,0)^{-1}g_f$, $y=\det g'$ and $s_0=g'a(y,0)^{-1}\in SL_2(\mathbf A_{E,f})$. The set of $s\in GL_2(\mathbf A_{E,f})$ with $a(y,0)^{-1}s^{-1}g'\in K$ is open and contains $s_0$, so by (a) it contains some $\sigma\in SL_2(E)$. With $u=a(y,0)^{-1}\sigma^{-1}g'\in K$, $g'=\sigma a(y,0)u$.
\item[(d)] \emph{The decomposition.} Let $\gamma=a(\alpha,0)\sigma\in GL_2(E)$, of determinant $\alpha$, and $h=\gamma^{-1}g_\infty$. Then $g_f=\gamma a(y,0)u$ and $g_\infty=\gamma h$, so $g=\gamma\,h_\infty(a(y,0)u)_f$; and $\det h=\alpha^{-1}\det g_\infty=\alpha^{-2}(\alpha\det g_\infty)$ is totally positive.
\end{enumerate}
\end{proof}
% strongApp (StrongApp.lean, 74 lines): (a) is GroundTheory.Automorphic.AdelicGLn.finiteSL_le_topologicalClosure (StrongApproximationSLDense.lean, 338 lines: denseRange_algebraMap_finiteAdeleRing, transvectionGL_mem_topologicalClosure, elementarySubgroup_le_topologicalClosure, truncate and mem_closure_elementarySubgroup), (b) is the private exists_sign, (c) the private exists_sl, and (d) the final refine.

\begin{lemma}\label{sha-siegel-step-seeds-induced-hilbert-form-exists-eisenstein-lift-unique}
Let $E$ be a totally real number field and $k\ge0$.
\begin{enumerate}
\item If a function $\phi$ on $GL_2(\mathbf A_E)$ has the Fourier expansions of weight $k$ with constant terms $c_0$ and $c'_0$ and coefficients $a$ and $a'$ (Definition~\ref{sha-siegel-step-seeds-induced-hilbert-form-exists-eisenstein-expansion}), then $c_0=c'_0$ and $a(\mathfrak m)=a'(\mathfrak m)$ for every nonzero integral ideal $\mathfrak m$.
\item Let $k\ge1$, $c$ a function from the primes of $E$ to $\mathbf Z_{\ge0}$, $\omega$ a Hecke character of $E$ of finite order, and $\phi$ a holomorphic form of parallel weight $k$, level $c$ and central character $\omega$ with the Fourier expansion of weight $k$ with constant term $c_0$ and coefficients $a$. If $a(\mathfrak m)=0$ for every $\mathfrak m$, then $\phi=0$.
\end{enumerate}
\end{lemma}
\leanref{Development/AutomorphicArithmetic/AutomorphicVarieties/IdeleGenerators/WeightChange.lean}{UpperTheory.V.Sd.ExpUnique}
\leanref{Development/AutomorphicArithmetic/AutomorphicVarieties/IdeleGenerators/HolomorphicIrreducibility/ExpUnique.lean}{UpperTheory.V.Sd.expUnique}
\leanref{Development/AutomorphicArithmetic/AutomorphicVarieties/IdeleGenerators/WeightChange.lean}{UpperTheory.V.Sd.ExpInj}
\leanref{Development/AutomorphicArithmetic/AutomorphicVarieties/IdeleGenerators/HolomorphicIrreducibility/ExpInj.lean}{UpperTheory.V.Sd.expInj}
\begin{proof}
Write $a(y,x)=\begin{pmatrix}y&x\\0&1\end{pmatrix}$ over any ring, $V$ for the set of infinite places of $E$, and $e(m\cdot x)$ as in Lemma~\ref{sha-siegel-step-seeds-induced-hilbert-form-exists-eisenstein-lift-unique-exponentials}. For a totally positive $\xi\in E$ let $\mu_\xi=(v(\xi))_{v\in V}\in\mathbf R^V$.

(1) Let $y_f\in\mathbf A_{E,f}^\times$, and take $y_v=1$ at every $v$. Subtracting the two expansions at $(x,y,y_f)$ shows that for every $x\in\mathbf R^V$ the family of the numbers $B_\xi e(\mu_\xi\cdot x)$, indexed by the totally positive $\xi$, with
\[
B_\xi=\bigl(a(\xi y_f)-a'(\xi y_f)\bigr)\exp\Bigl(-2\pi\sum_vv(\xi)\Bigr),
\]
is summable with sum $c'_0(y_f)-c_0(y_f)$. Each $\mu_\xi$ is nonzero, as $\xi$ is totally positive. By Lemma~\ref{sha-siegel-step-seeds-induced-hilbert-form-exists-eisenstein-lift-unique-exponentials}(1), $c'_0(y_f)=c_0(y_f)$. Let $\mathfrak m$ be a nonzero integral ideal. Choose $y_f\in\mathbf A_{E,f}^\times$ with $\operatorname{ord}_w(y_{f,w})=\operatorname{ord}_w(\mathfrak m)-\operatorname{ord}_w(\mathfrak d)$ at every prime $w$. It generates $\mathfrak m\mathfrak d^{-1}$ and no other $\mathfrak m'\mathfrak d^{-1}$, so $a(y_f)=a(\mathfrak m)$ and $a'(y_f)=a'(\mathfrak m)$. The map $\xi\mapsto\mu_\xi$ is injective, as each embedding of $E$ is. By Lemma~\ref{sha-siegel-step-seeds-induced-hilbert-form-exists-eisenstein-lift-unique-exponentials}(2), now with sum $0$, $B_1=0$, that is, $a(\mathfrak m)=a'(\mathfrak m)$.

(2) As every $a(\mathfrak m)$ is $0$, the function $y\mapsto a(y)$ on $\mathbf A_{E,f}^\times$ vanishes.
\begin{enumerate}
\item[(a)] \emph{Absolute values.} Let $h\in GL_2(E\otimes\mathbf R)$ have totally positive determinant, and $y_f\in\mathbf A_{E,f}^\times$. Then
\[
\bigl|\phi\bigl(h_\infty a(y_f,0)_f\bigr)\bigr|=\Bigl(\prod_vY_v\cdot|y_f|\Bigr)^{k/2}|c_0(y_f)|,\qquad Y_v=\frac{\det h_v}{h_{v,21}^2+h_{v,22}^2}.
\]
Indeed, let $z_v=h_{v,22}-ih_{v,21}$, which is nonzero, $\rho_v=|z_v|$ and $\theta_v=\arg z_v$. A direct computation gives $h=\rho\,a(Y,x)\,r(\theta)$, with $\rho$ the scalar matrix of $(\rho_v)$, $x_v=(h_{v,11}h_{v,21}+h_{v,12}h_{v,22})/(h_{v,21}^2+h_{v,22}^2)$ and $r(\theta)$ as in Lemma~\ref{sha-siegel-step-seeds-induced-hilbert-form-exists-eisenstein-lift-unique-rotation}; and $Y$ is totally positive. By (c) of Definition~\ref{sha-siegel-step-seeds-induced-hilbert-form-exists-holomorphic}, $\phi(h_\infty a(y_f,0)_f)=\omega(\rho)\phi\bigl(a(Y,x)_\infty a(y_f,0)_f\,r(\theta)_\infty\bigr)$, and $|\omega(\rho)|=1$ as $\omega$ has finite order. By Lemma~\ref{sha-siegel-step-seeds-induced-hilbert-form-exists-eisenstein-lift-unique-rotation} the factor $r(\theta)$ does not change the absolute value. The expansion of $\phi$ at $(x,Y,y_f)$ has all its terms $0$, so $\phi(a(Y,x)_\infty a(y_f,0)_f)=(\prod_vY_v\cdot|y_f|)^{k/2}c_0(y_f)$.
\item[(b)] \emph{The constant term.} Let $y_f\in\mathbf A_{E,f}^\times$ and $t=a(y_f,0)\in GL_2(\mathbf A_{E,f})$, and for an integer $N$ let $l(N)=\begin{pmatrix}1&0\\N&1\end{pmatrix}$. The integers $n!$ tend to $0$ in $\mathbf A_{E,f}$, and the stabilizer of $\phi$ in $GL_2(\mathbf A_{E,f})$ is open; so for some $N=n!\ge1$ it contains $t^{-1}l(N)t$. As $\phi$ is invariant under left translation by $l(N)\in GL_2(E)$,
\[
\phi(t_f)=\phi\bigl(l(N)_\infty l(N)_ft_f\bigr)=\phi\bigl(l(N)_\infty t_f(t^{-1}l(N)t)_f\bigr)=\phi\bigl(l(N)_\infty t_f\bigr).
\]
Apply (a) to $h=1$, with $Y_v=1$, and to $h=l(N)$, with $Y_v=1/(N^2+1)$: $|y_f|^{k/2}|c_0(y_f)|=\bigl((N^2+1)^{-[E:\mathbf Q]}|y_f|\bigr)^{k/2}|c_0(y_f)|$. As $N\ge1$ and $k\ge1$, $\bigl((N^2+1)^{-[E:\mathbf Q]}|y_f|\bigr)^{k/2}<|y_f|^{k/2}$, so $c_0(y_f)=0$.
\item[(c)] \emph{Every value.} Let $g\in GL_2(\mathbf A_E)$. By Lemma~\ref{sha-siegel-step-seeds-induced-hilbert-form-exists-eisenstein-lift-unique-approximation}, applied to the stabilizer $K$ of $\phi$, $g=\gamma\,h_\infty(a(y,0)u)_f$ with $\gamma\in GL_2(E)$, $h$ of totally positive determinant, $y\in\mathbf A_{E,f}^\times$ and $u\in K$. So $\phi(g)=\phi(h_\infty a(y,0)_f)$, and by (a) and (b) its absolute value is $0$.
\end{enumerate}
\end{proof}
% expUnique (ExpUnique.lean, 178 lines) is (1): key for the difference of the expansions at y = 1, const_zero for c0, exists_genD for y_f, hinj for the injectivity and coef_zero at the index 1, with coefAt_eq (genD_unique). expInj (ExpInj.lean, 252 lines) is (2): norm_inf for (a), with iwasawa, hφ.cent, the bound hn from IsOfFinOrder, norm_rot and the expansion; const_zero for (b), with tendsto_factorial_finiteAdele, lowA and e2; and the final strongApp (Lemma sha-siegel-step-seeds-induced-hilbert-form-exists-eisenstein-lift-unique-approximation) for (c).

\begin{lemma}\label{sha-siegel-step-seeds-induced-hilbert-form-exists-eisenstein-lift-deligne-serre}
Let $\ell$ be a prime, $I$ and $J$ sets, and $V$ a finite-dimensional $\bar{\mathbf Q}_\ell$-subspace of the functions $I\to\bar{\mathbf Q}_\ell$. Call $v\in V$ \emph{integral} if $|v(n)|\le1$ for every $n\in I$, and \emph{small} if $|v(n)|<1$ for every $n$. Let $(T_j)_{j\in J}$ be pairwise commuting linear operators on $V$ that map integral vectors to integral vectors and small vectors to small vectors. Let $f\in V$ be integral with $|f(n)|=1$ for some $n$, and $a_j\in\bar{\mathbf Q}_\ell$ with $|a_j|\le1$ and $T_jf-a_jf$ small for every $j$. Then there are an integral $h\in V$, $h\ne0$, and $b_j\in\bar{\mathbf Q}_\ell$ with $T_jh=b_jh$ and $|b_j-a_j|<1$ for every $j$.
\end{lemma}
\leanref{Development/AutomorphicArithmetic/AutomorphicVarieties/IdeleGenerators/ConstSmall.lean}{UpperTheory.V.Sd.DelSerre}
\leanref{Development/AutomorphicArithmetic/AutomorphicVarieties/IdeleGenerators/HolomorphicIrreducibility/DelSerre.lean}{UpperTheory.V.Sd.delSerre}
\begin{proof}
Let $B\subseteq V$ be the set of $v$ such that $cv$ is integral for some $c\ne0$; by the ultrametric inequality it is a subspace, and the $T_j$ map it into itself. We show, by induction on $d$, that for every $d$-dimensional subspace $U\subseteq B$ with $T_jU\subseteq U$ for every $j$, and every $f\in U$ with the properties of the statement, $h$ and $b_j$ exist. The statement is the case $U=B$, which contains $f$.

Suppose first that for every $j$, every eigenvalue $c$ of $T_j$ with an eigenvector in $U$ satisfies $|c-a_j|<1$. As $f\ne0$, $U\ne0$; as $\bar{\mathbf Q}_\ell$ is algebraically closed and the $T_j$ commute and map $U$ into itself, they have a common eigenvector $v\in U$, with eigenvalues $b_j$. Then $|b_j-a_j|<1$, and some nonzero multiple $h$ of $v$ is integral, as $v\in B$.

Otherwise there are $j$, $c\in\bar{\mathbf Q}_\ell$ and a nonzero $v\in U$ with $T_jv=cv$ and $|c-a_j|\ge1$. Replacing $v$ by a nonzero multiple, $v$ is integral. Then $|c|\le1$: otherwise, as the vectors $T_j^mv=c^mv$ are integral for all $m$, $|c|^m|v(n)|\le1$ for every $n$ and $m$, which fails for $m$ large and an $n$ with $v(n)\ne0$. So $|c-a_j|\le1$, and $|c-a_j|=1$. Let $S=T_j-c$; it commutes with the $T_i$, maps $U$ into itself and kills $v$, so $U'=S(U)\subseteq B$ is stable under the $T_i$ and has dimension less than $d$. The vector $Sf\in U'$ is integral, as $|c|\le1$. At an $n$ with $|f(n)|=1$, $(T_jf-a_jf)(n)-(Sf)(n)=(c-a_j)f(n)$ has absolute value $1$ and $(T_jf-a_jf)(n)$ is small, so $|(Sf)(n)|=1$. For every $i$, $T_iSf-a_iSf=S(T_if-a_if)$ is small, as $T_if-a_if$ is small, $T_j$ maps small vectors to small vectors and $|c|\le1$. By induction, $h$ and $b_j$ exist for $U'$ and $Sf$.
\end{proof}
% delSerre (DelSerre.lean, 165 lines): the private main by strong induction on finrank U, with bdd for B, UpperTheory.V.LinAlg.simEig for the common eigenvector, eig_le_one for |c| <= 1, finrank_map_lt for the dimension, and sub_le_one', sub_lt_one' for the ultrametric estimates. This is a form of the lemma of Deligne and Serre [DS74, Lemme 6.11], over the algebraically closed field Qbar l and with integrality measured coordinatewise; the Lean proves it.

\begin{lemma}\label{sha-siegel-step-seeds-induced-hilbert-form-exists-eisenstein-lift}
Assume the Galois representations as in Theorem~\ref{sha-residually-reducible-automorphy}. Let $E$ be a totally real number field with uniformizers $\varpi_w$, $\ell$ a prime and $\iota:\bar{\mathbf Q}_\ell\cong\mathbf C$, and write $x\equiv y$ for complex numbers with $|\iota^{-1}(x-y)|<1$. Let $k\ge1$, $c$ a function from the primes of $E$ to $\mathbf Z_{\ge0}$ that vanishes at all but finitely many primes, and $\omega$ a Hecke character of $E$ of finite order. Let $\phi$ be a holomorphic form of parallel weight $k$, level $c$ and central character $\omega$ with the Fourier expansion of weight $k$ with constant term $c_0$ and coefficients $a$ (Definition~\ref{sha-siegel-step-seeds-induced-hilbert-form-exists-eisenstein-expansion}), such that $|\iota^{-1}(a(\mathfrak m))|\le1$ for every $\mathfrak m$ and $|\iota^{-1}(a(\mathcal O_E))|=1$. Let $t(w)\in\mathbf C$ with $|\iota^{-1}(t(w))|\le1$ and $T^k_wa(\mathfrak m)\equiv t(w)a(\mathfrak m)$ for every prime $w$ and every $\mathfrak m$, with $T^k_w$ for $c$ and $\omega$. Then there is a holomorphic eigenform of parallel weight $k$, level $c$ and central character $\omega$ (Definition~\ref{sha-siegel-step-seeds-induced-hilbert-form-exists-holomorphic}) whose eigenvalues $\lambda(w)$ satisfy $\lambda(w)\equiv t(w)$ for every prime $w$.
\end{lemma}
\leanref{Development/AutomorphicArithmetic/AutomorphicVarieties/IdeleGenerators/WeightChange.lean}{UpperTheory.V.Sd.HolLift}
\leanref{Development/AutomorphicArithmetic/AutomorphicVarieties/ModularityArgument/ThresholdArguments/DenseThresholdCases/THolLift.lean}{UpperTheory.V.Sd.holLift_alt}
\begin{proof}
For a prime $w$ let $h_w=\operatorname N(w)^{k/2-1}U_w$, so that a holomorphic eigenform with eigenvalues $\lambda(w)$ is a nonzero form $g$ with $h_wg=\lambda(w)g$. Let $I$ be the set of nonzero integral ideals of $E$.
\begin{enumerate}
\item[(1)] \emph{The space.} Let $X$ be the set of holomorphic forms of parallel weight $k$, level $c$ and central character $\omega$ that have a Fourier expansion of weight $k$. Sums and scalar multiples of such forms are such forms, with the sums and multiples of their expansions; so $X$ is a vector space, finite-dimensional by Lemma~\ref{sha-siegel-step-seeds-induced-hilbert-form-exists-eisenstein-lift-finite}. It contains $\phi$.
\item[(2)] \emph{Coefficients.} By Lemma~\ref{sha-siegel-step-seeds-induced-hilbert-form-exists-eisenstein-lift-unique}(1), the coefficients $a_g(\mathfrak m)$, $\mathfrak m\in I$, of the expansion of $g\in X$ are determined by $g$. Let $\Phi(g)$ be the function $\mathfrak m\mapsto\iota^{-1}(a_g(\mathfrak m))$ on $I$. Then $\Phi(g+g')=\Phi(g)+\Phi(g')$ and $\Phi(zg)=\iota^{-1}(z)\Phi(g)$, and $\Phi$ is injective by Lemma~\ref{sha-siegel-step-seeds-induced-hilbert-form-exists-eisenstein-lift-unique}(2), as $k\ge1$. So $V=\Phi(X)$ is a finite-dimensional $\bar{\mathbf Q}_\ell$-subspace of the functions $I\to\bar{\mathbf Q}_\ell$.
\item[(3)] \emph{Hecke operators.} For a prime $w$ let $e(w)=\iota^{-1}\bigl(\chi_{c,\omega}(w)\operatorname N(w)^{k-1}\bigr)$, and let $F_w$ send a function $v$ on $I$ to $\mathfrak m\mapsto v(w\mathfrak m)+e(w)v(\mathfrak mw^{-1})$, where $v(\mathfrak mw^{-1})=0$ if $w$ does not divide $\mathfrak m$. By Lemma~\ref{sha-siegel-step-seeds-induced-hilbert-form-exists-eisenstein-lift-hecke}, $h_w$ maps $X$ into itself, the $h_w$ commute on $X$, and $\Phi(h_wg)=F_w\Phi(g)$ for $g\in X$, as the coefficients of $h_wg$ are $T^k_wa_g$. So $F_w$ maps $V$ into itself, and the restrictions of the $F_w$ to $V$ commute.
\item[(4)] \emph{Integrality.} The values of $\chi_{c,\omega}$ are $0$ or roots of unity and $\operatorname N(w)$ is an integer, so $|e(w)|\le1$. By the ultrametric inequality, each $F_w$ maps integral vectors of $V$ to integral vectors and small vectors to small vectors, in the sense of Lemma~\ref{sha-siegel-step-seeds-induced-hilbert-form-exists-eisenstein-lift-deligne-serre}.
\item[(5)] \emph{A common eigenvector.} The vector $f=\Phi(\phi)$ is integral and $|f(\mathcal O_E)|=1$, by the hypotheses on $a$. Let $a_w=\iota^{-1}(t(w))$, so $|a_w|\le1$, and $(F_wf-a_wf)(\mathfrak m)=\iota^{-1}\bigl(T^k_wa(\mathfrak m)-t(w)a(\mathfrak m)\bigr)$ is small. By (3), (4) and Lemma~\ref{sha-siegel-step-seeds-induced-hilbert-form-exists-eisenstein-lift-deligne-serre}, there are a nonzero $h\in V$ and $b_w\in\bar{\mathbf Q}_\ell$ with $F_wh=b_wh$ and $|b_w-a_w|<1$ for every prime $w$.
\item[(6)] \emph{The eigenform.} Let $h=\Phi(g)$ with $g\in X$; then $g\ne0$. By (3), $\Phi(h_wg)=F_wh=b_wh=\Phi(\iota(b_w)g)$, so $h_wg=\iota(b_w)g$, as $\Phi$ is injective. So $g$ is a holomorphic eigenform of parallel weight $k$, level $c$ and central character $\omega$ with eigenvalues $\lambda(w)=\iota(b_w)$, and $|\iota^{-1}(\lambda(w)-t(w))|=|b_w-a_w|<1$.
\end{enumerate}
\end{proof}
% holLift_alt (THolLift.lean, 81 lines, with the private SS, cf, Phi, Phi_eq, fT, fT_Phi, hkS, Phi_hkS, Phi_inj and norm_chiI), followed step by step: (1) is SS with holLin and holomorphic_finite_of_printed_hida (Lemma sha-siegel-step-seeds-induced-hilbert-form-exists-eisenstein-lift-finite); (2) is cf, cf_eq (expUnique), Phi (semilinear over iota^-1) and Phi_inj (expInj), with V the image of the top submodule; (3) is fT, e, hmem with hkS and Phi_hkS, and the commutation through holHk (Lemma sha-siegel-step-seeds-induced-hilbert-form-exists-eisenstein-lift-hecke); (4) is he with norm_chiI and the two bounds passed to delSerre; (5) is delSerre (Lemma sha-siegel-step-seeds-induced-hilbert-form-exists-eisenstein-lift-deligne-serre) with Phi_eq for f, the index top for |f(O_E)| = 1, and ht and hcong; (6) is the final refine with Phi_inj. The assumption on Galois representations is the HypothesesH' argument of holomorphic_finite_of_printed_hida, which does not use it. Size: about 3,100 lines with its dependencies, net of the definitions, of which about 800 are shared with Lemma sha-siegel-step-seeds-induced-hilbert-form-exists-eisenstein-noncuspidal; the open part left is in Lemma sha-siegel-step-seeds-induced-hilbert-form-exists-eisenstein-lift-finite-harish-chandra, on which Lemma sha-siegel-step-seeds-induced-hilbert-form-exists-eisenstein-lift-finite rests.

\begin{lemma}\label{sha-siegel-step-seeds-induced-hilbert-form-exists-eisenstein-noncuspidal-constant-term-periodic}
Let $A$ be an abelian topological group whose topology is second countable, $\Gamma\subseteq A$ a subgroup, and $C\subseteq A$ a compact set with $A=\Gamma+C$. Call a function $f$ on $A$ (or on $A\times\mathbf R$) \emph{$\Gamma$-periodic} if $f(\gamma+x)=f(x)$ (or $f(\gamma+x,t)=f(x,t)$) for all $\gamma\in\Gamma$ and all arguments.
\begin{enumerate}
\item Suppose moreover that $A$ is locally compact and Hausdorff and that $\Gamma$ is countable. Let $\mu$ be a Haar measure on $A$ and $D$ a measurable fundamental domain for the translations by $\Gamma$. Then $\mu(D)$ is finite, and for every $\Gamma$-periodic function $f:A\to\mathbf C$: if $f$ is continuous, it is integrable on $D$; $\int_Df(x+y)\,d\mu(x)=\int_Df(x)\,d\mu(x)$ for every $y\in A$; and $\int_Df(e(x))\,d\mu(x)=\int_Df(x)\,d\mu(x)$ for every automorphism $e$ of the topological group $A$ with $e(\Gamma)=\Gamma$.
\item Let $\mu$ be a Borel measure on $A$ and $D\subseteq A$ a set with $\mu(D)$ finite. Let $f,f':A\times\mathbf R\to\mathbf C$ be continuous and $\Gamma$-periodic, such that $t\mapsto f(x,t)$ has derivative $f'(x,t)$ at every $t\in\mathbf R$, for every $x\in A$. Then $t\mapsto\int_Df(x,t)\,d\mu(x)$ has derivative $\int_Df'(x,0)\,d\mu(x)$ at $t=0$.
\end{enumerate}
\end{lemma}
\leanref{Development/AutomorphicArithmetic/AutomorphicVarieties/IdeleGenerators/PerDeriv.lean}{UpperTheory.V.Sd.Eis.PerMean}
\leanref{Development/AutomorphicArithmetic/AutomorphicVarieties/BuildingsResidueChars/PeriodicIntegralFundamentalDomainInvariant.lean}{UpperTheory.periodic_integral_fundamentalDomain_invariant}
\leanref{Development/AutomorphicArithmetic/AutomorphicVarieties/IdeleGenerators/PerDeriv.lean}{UpperTheory.V.Sd.Eis.PerDeriv}
\leanref{Development/AutomorphicArithmetic/AutomorphicVarieties/BuildingsResidueChars/EisensteinPeriodDerivative.lean}{UpperTheory.eisenstein_period_derivative}
\begin{proof}
(1) As $D$ is a fundamental domain, $\mu(C)=\sum_{\gamma\in\Gamma}\mu\bigl((\gamma+C)\cap D\bigr)$, and the sets $(\gamma+C)\cap D$ cover $D$, as $A=\Gamma+C$. So $\mu(D)\le\mu(C)$, which is finite as $C$ is compact; and $\mu(D)\ne0$, as the translates of $D$ cover $A$ up to a null set. A continuous $\Gamma$-periodic $f$ is bounded on $C$, hence on $A$, and so integrable on $D$. For the invariances, let $e$ be a homeomorphism of $A$ with $e_*\mu=c\mu$ for some $c\ge0$, and $\sigma$ a bijection of $\Gamma$ with $e(\sigma(\gamma)+x)=\gamma+e(x)$. Then $e^{-1}(D)$ is again a fundamental domain, so $c\mu(D)=\mu(e^{-1}(D))=\mu(D)$, and $c=1$. The function $f\circ e$ is $\Gamma$-periodic, and the integrals of a $\Gamma$-periodic function over two fundamental domains agree; so $\int_Df(e(x))\,d\mu(x)=\int_{e^{-1}(D)}f(e(x))\,d\mu(x)=\int_Df(x)\,d\mu(x)$, the last as $e_*\mu=\mu$. For $e(x)=x+y$, take $c=1$ and $\sigma$ the identity. For an automorphism $e$ with $e(\Gamma)=\Gamma$, $e_*\mu$ is a Haar measure, hence a multiple of $\mu$; take $\sigma=e^{-1}$ on $\Gamma$.

(2) The functions $f'$ and $f$ are bounded on the compact set $C\times[-1,1]$, hence on $A\times[-1,1]$, as they are $\Gamma$-periodic; let $M$ bound $f'$. As $\mu(D)$ is finite, the constant $M$ is integrable on $D$, and so is $x\mapsto f(x,0)$. By the theorem on differentiation under the integral sign, with the bound $|f'(x,t)|\le M$ for $|t|\le1$, $t\mapsto\int_Df(x,t)\,d\mu(x)$ has derivative $\int_Df'(x,0)\,d\mu(x)$ at $0$.
\end{proof}
% periodic_integral_fundamentalDomain_invariant (PeriodicIntegralFundamentalDomainInvariant.lean, 70 lines): perMean_fin (IsAddFundamentalDomain.measure_eq_tsum), Measure.integrableOn_of_bounded, and perMean_key (IsAddFundamentalDomain.preimage_of_equiv, measure_set_eq, setIntegral_eq and setIntegral_map_equiv) applied to MeasurableEquiv.addRight y and to the automorphism e, with Measure.isAddLeftInvariant_eq_smul. eisenstein_period_derivative (EisensteinPeriodDerivative.lean, 39 lines): exists_bound_of_continuous_periodic_cocompact for both bounds and hasDerivAt_integral_of_dominated_loc_of_deriv_le on Icc (-1) 1. The Lean's measurable spaces are Borel; its Cocpt Gamma is the text's compact C.

\begin{lemma}\label{sha-siegel-step-seeds-induced-hilbert-form-exists-eisenstein-noncuspidal-constant-term-matrix}
Let $h\in GL_2(\mathbf R)$, $p=(p_{ab})$ a complex $2\times2$ matrix and $\kappa\in\mathbf C$, and write $e_{ab}$ for the matrix units. Suppose that $p_{11}+p_{22}=0$, $p_{12}-p_{21}=i\kappa$, $p_{11}-p_{22}=i(p_{12}+p_{21})$ and $\sum_{a,b}(h^{-1}e_{12}h)_{ab}\,p_{ab}=0$. Then $\sum_{a,b}(h^{-1}e_{11}h)_{ab}\,p_{ab}=\kappa/2$.
\end{lemma}
\leanref{Development/AutomorphicArithmetic/AutomorphicVarieties/IdeleGenerators/PerDeriv.lean}{UpperTheory.V.Sd.Eis.LieAlg}
\leanref{Development/AutomorphicArithmetic/AutomorphicVarieties/BuildingsResidueChars/EisensteinLieAlgebraAction.lean}{UpperTheory.eisenstein_lie_algebra_action}
\begin{proof}
Let $h$ have rows $(\alpha,\beta)$ and $(\gamma,\delta)$ and determinant $d$, and let $a=p_{11}$. The relations give $p_{22}=-a$, $p_{12}=i(\frac\kappa2-a)$ and $p_{21}=-i(\frac\kappa2+a)$. As $h^{-1}e_{1j}h$ is $d^{-1}$ times the product of the column $(\delta,-\gamma)$ with the $j$-th row of $h$, the hypothesis reads $2a\gamma\delta+i\delta^2(\frac\kappa2-a)+i\gamma^2(\frac\kappa2+a)=0$, that is, $-ia(\delta+i\gamma)^2+i\frac\kappa2(\delta+i\gamma)(\delta-i\gamma)=0$. As $\gamma$ and $\delta$ are real and not both $0$, $\delta+i\gamma\ne0$, so $a(\delta+i\gamma)=\frac\kappa2(\delta-i\gamma)$. The sum of the conclusion is
\[
d^{-1}\Bigl(a(\alpha\delta+\beta\gamma)+i\delta\beta\bigl(\tfrac\kappa2-a\bigr)+i\gamma\alpha\bigl(\tfrac\kappa2+a\bigr)\Bigr)=d^{-1}\Bigl(a(\alpha-i\beta)(\delta+i\gamma)+i\tfrac\kappa2(\beta\delta+\alpha\gamma)\Bigr)=d^{-1}\tfrac\kappa2(\alpha\delta-\beta\gamma)=\tfrac\kappa2.
\]
\end{proof}
% eisenstein_lie_algebra_action (EisensteinLieAlgebraAction.lean, 61 lines): the private complex_linear_combination_eq_half_of_relations, with the entries of h and h^-1 and the cancellation of c - i d (the text's -i(delta + i gamma)), which is hc.

\begin{lemma}\label{sha-siegel-step-seeds-induced-hilbert-form-exists-eisenstein-noncuspidal-constant-term}
Let $E$ be a totally real number field, $k\ge0$, $c$ a function from the primes of $E$ to $\mathbf Z_{\ge0}$, $\omega$ a Hecke character of $E$ of finite order, and $\phi$ a holomorphic form of parallel weight $k$, level $c$ and central character $\omega$ (Definition~\ref{sha-siegel-step-seeds-induced-hilbert-form-exists-holomorphic}). Let $\mu$ be a Haar measure on $\mathbf A_E$ and $D\subseteq\mathbf A_E$ a measurable fundamental domain for the translations by $E$. Write $a(y,x)=\begin{pmatrix}y&x\\0&1\end{pmatrix}$ and $n(x)=a(1,x)$ over any ring. Then:
\begin{enumerate}
\item for every $g\in GL_2(\mathbf A_E)$, the function $x\mapsto\phi(n(x)g)$ is integrable on $D$; let $\psi(g)=\int_D\phi(n(x)g)\,d\mu(x)$;
\item $\psi(n(x)g)=\psi(g)$ for $x\in\mathbf A_E$, and $\psi(a(\alpha,0)g)=\psi(g)$ for $\alpha\in E^\times$;
\item $\psi(a(y,0)_\infty g)=\bigl(\prod_vy_v\bigr)^{k/2}\psi(g)$ for every totally positive $y\in E\otimes\mathbf R$.
\end{enumerate}
\end{lemma}
\leanref{Development/AutomorphicArithmetic/AutomorphicVarieties/IdeleGenerators/PerDeriv.lean}{UpperTheory.V.Sd.Eis.CtBasic}
\leanref{Development/AutomorphicArithmetic/AutomorphicVarieties/BuildingsResidueChars/EisensteinConstantTermOfPeriodMeanOfUpperBlock.lean}{UpperTheory.eisenstein_constantTerm_of_periodMean_of_upperBlock}
\leanref{Development/AutomorphicArithmetic/AutomorphicVarieties/IdeleGenerators/PerDeriv.lean}{UpperTheory.V.Sd.Eis.CtLeft}
\leanref{Development/AutomorphicArithmetic/AutomorphicVarieties/BuildingsResidueChars/EisensteinContinuousLeftAction.lean}{UpperTheory.eisenstein_continuous_left_action}
\leanref{Development/AutomorphicArithmetic/AutomorphicVarieties/IdeleGenerators/HolomorphicIrreducibility/HolEis.lean}{UpperTheory.V.Sd.Eis.ct_arch}
\begin{proof}
For a function $\theta$ on $GL_2(\mathbf A_E)$ such that $x\mapsto\theta(n(x)g)$ is integrable on $D$ for every $g$, write $\psi_\theta(g)=\int_D\theta(n(x)g)\,d\mu(x)$; so $\psi=\psi_\phi$. For a $2\times2$ matrix $X$ with entries in $E\otimes\mathbf R$ and a function $\theta$ smooth in the variable at infinity, let $X\theta(g)$ be the derivative at $t=0$ of $\theta(g\,(1+tX)_\infty)$. It is real-linear in $X$; $X^v_{ab}\theta$ is $X\theta$ for $X$ the matrix unit $e_{ab}$ at $v$; and by the chain rule, $X\theta(g)$ is the derivative at $t=0$ of $\theta(g\,\tau(t)_\infty)$ for every differentiable $\tau:\mathbf R\to GL_2(E\otimes\mathbf R)$ with $\tau(0)=1$ and derivative $X$ at $0$. Call a continuous $\tau:\mathbf R\to GL_2(E\otimes\mathbf R)$ with $\tau(s+t)=\tau(s)\tau(t)$ for all $s,t$ and a derivative at $0$ a \emph{one-parameter subgroup}.
\begin{enumerate}
\item[(a)] \emph{The adeles.} The ring $\mathbf A_E$ is locally compact, second countable and Hausdorff, $E$ is countable, and $x\mapsto n(x)$ is continuous. There is a compact $C\subseteq\mathbf A_E$ with $\mathbf A_E=E+C$. Indeed, let $P$ be a compact fundamental parallelotope of the lattice $\mathcal O_E$ in $E\otimes\mathbf R$, and $x\in\mathbf A_E$. Choose an integer $N\ge1$ with $Nx_f\in\hat{\mathcal O}_E$. As $\mathcal O_E/N\mathcal O_E=\hat{\mathcal O}_E/N\hat{\mathcal O}_E$, there is $\beta\in\mathcal O_E$ with $x_f-\beta/N\in\hat{\mathcal O}_E$, and there is $\beta'\in\mathcal O_E$ with $x_\infty-\beta/N-\beta'\in P$. So $x-(\beta/N+\beta')\in P\times\hat{\mathcal O}_E$, and $C=P\times\hat{\mathcal O}_E$ will do.
\item[(b)] \emph{Continuity.} Let $U$ be the stabilizer of $\phi$ in $GL_2(\mathbf A_{E,f})$, which is open, and $g_0\in GL_2(\mathbf A_E)$ with finite part $k_0$. For $g$ with $g_f\in k_0U$, $\phi(g)=\phi(g_\infty(k_0)_f)$, a smooth function of $g_\infty$; so $\phi$ is continuous near $g_0$. Right translation by $GL_2(\mathbf A_{E,f})$ commutes with every $X$, so the stabilizer of $X\phi$ contains $U$, and $X\phi$ is continuous by the same argument. Both $\phi$ and $X\phi$ are invariant under left translation by $GL_2(E)$.
\item[(c)] \emph{Periodicity.} If $\theta$ is invariant under left translation by $GL_2(E)$, then $x\mapsto\theta(n(x)h)$ is $E$-periodic for every $h\in GL_2(\mathbf A_E)$, as $n(\gamma+x)=n(\gamma)n(x)$ with $n(\gamma)\in GL_2(E)$ for $\gamma\in E$.
\item[(d)] \emph{(1) and (2).} By (a), (b), (c) and Lemma~\ref{sha-siegel-step-seeds-induced-hilbert-form-exists-eisenstein-noncuspidal-constant-term-periodic}(1), $x\mapsto\theta(n(x)g)$ is integrable on $D$ for $\theta=\phi$ and for $\theta=X\phi$; this gives (1). For $y\in\mathbf A_E$, $n(x)n(y)=n(x+y)$, so $\psi(n(y)g)=\int_D\phi(n(x+y)g)\,d\mu(x)=\psi(g)$ by the same lemma. For $\alpha\in E^\times$, $n(x)a(\alpha,0)=a(\alpha,0)n(\alpha^{-1}x)$ with $a(\alpha,0)\in GL_2(E)$, so $\psi(a(\alpha,0)g)=\int_D\phi(n(\alpha^{-1}x)g)\,d\mu(x)$. Multiplication by $\alpha^{-1}$ is an automorphism of the topological group $\mathbf A_E$ that maps $E$ onto $E$, so this is $\psi(g)$ by the same lemma. This gives (2).
\item[(e)] \emph{Differentiating the constant term.} Let $\tau$ be a one-parameter subgroup with derivative $X$ at $0$, and $g\in GL_2(\mathbf A_E)$. Then $t\mapsto\psi(g\,\tau(t)_\infty)$ has derivative $\psi_{X\phi}(g)$ at $t=0$. Indeed, $f(x,t)=\phi(n(x)g\,\tau(t)_\infty)$ and $f'(x,t)=X\phi(n(x)g\,\tau(t)_\infty)$ are continuous by (a) and (b), and $E$-periodic in $x$ by (c). As $\tau(t+s)=\tau(t)\tau(s)$, $t\mapsto f(x,t)$ has derivative $f'(x,t)$ at every $t$. And $\mu(D)$ is finite by Lemma~\ref{sha-siegel-step-seeds-induced-hilbert-form-exists-eisenstein-noncuspidal-constant-term-periodic}(1). Part (2) of that lemma gives the claim.
\item[(f)] \emph{A linear form.} Fix $g$, let $h=g_\infty$, and let $L(X)=\psi_{X\phi}(g)$, which is real-linear in $X$ by (d). Let $\tau$ be a one-parameter subgroup with derivative $X$ at $0$. Then $\tau(t)_\infty g=g\,(h^{-1}\tau(t)h)_\infty$, and $t\mapsto h^{-1}\tau(t)h$ is a one-parameter subgroup with derivative $h^{-1}Xh$ at $0$. So, by (e), $t\mapsto\psi(\tau(t)_\infty g)$ has derivative $L(h^{-1}Xh)$ at $t=0$.
\item[(g)] \emph{The unipotent direction.} Let $x\in E\otimes\mathbf R$. The map $t\mapsto n(tx)$ is a one-parameter subgroup with derivative $xe_{12}$ at $0$, and $\psi(n(tx)_\infty g)=\psi(g)$ for every $t$ by (2). By (f), $L(h^{-1}xe_{12}h)=0$.
\item[(h)] \emph{At each infinite place.} For an infinite place $v$ let $h_v$ be the component of $h$ at $v$, and $p^v_{ab}=L(e_{ab}\text{ at }v)=\psi_{X^v_{ab}\phi}(g)$. For $x\in E\otimes\mathbf R$ and $i,j\in\{1,2\}$, $h^{-1}(xe_{ij})h$ has component $x_vh_v^{-1}e_{ij}h_v$ at $v$, so
\[
L(h^{-1}xe_{ij}h)=\sum_vx_v\sum_{a,b}(h_v^{-1}e_{ij}h_v)_{ab}\,p^v_{ab}.
\]
By Lemma~\ref{sha-siegel-step-seeds-induced-hilbert-form-exists-representation-casimir}(1), $(X^v_{11}+X^v_{22})\phi=0$, so $p^v_{11}+p^v_{22}=0$. By condition (d) of Definition~\ref{sha-siegel-step-seeds-induced-hilbert-form-exists-holomorphic} and the linearity of $\theta\mapsto\psi_\theta(g)$, $p^v_{12}-p^v_{21}=ik\psi(g)$ and $p^v_{11}-p^v_{22}=i(p^v_{12}+p^v_{21})$. By (g) and the display, with $x_v=1$ and $x_{v'}=0$ for $v'\ne v$, $\sum_{a,b}(h_v^{-1}e_{12}h_v)_{ab}\,p^v_{ab}=0$. By Lemma~\ref{sha-siegel-step-seeds-induced-hilbert-form-exists-eisenstein-noncuspidal-constant-term-matrix}, with $\kappa=k\psi(g)$, $\sum_{a,b}(h_v^{-1}e_{11}h_v)_{ab}\,p^v_{ab}=\frac k2\psi(g)$.
\item[(i)] \emph{The diagonal direction.} Let $u\in E\otimes\mathbf R$ and $y_t=(e^{tu_v})_v$ for $t\in\mathbf R$. The map $t\mapsto a(y_t,0)$ is a one-parameter subgroup with derivative $ue_{11}$ at $0$. By (f), (h) and the display, $t\mapsto\psi(a(y_t,0)_\infty g)$ has derivative $L(h^{-1}ue_{11}h)=\frac k2\bigl(\sum_vu_v\bigr)\psi(g)$ at $t=0$.
\end{enumerate}
(3) Let $u$ and $y_t$ be as in (i). As $a(y_{s+t},0)=a(y_s,0)a(y_t,0)$, (i) applied at $a(y_t,0)_\infty g$ shows that $f(t)=\psi(a(y_t,0)_\infty g)$ satisfies $f'(t)=\frac k2(\sum_vu_v)f(t)$ for every $t$, so $f(1)=\exp\bigl(\frac k2\sum_vu_v\bigr)f(0)$. For a totally positive $y$, take $u_v=\log y_v$.
\end{proof}
% The proof follows the Lean. (a) is BlockTop, proved by eisenstein_constantTerm_upperBlock (EisensteinConstantTermUpperBlock.lean, 36 lines, with GroundTheory.Global.exists_isCompact_forall_exists_sub_mem for C, through existsUnique_sub_mem_fundamentalSet). (b) is SmCont, proved by eisenstein_smooth_continuous (EisensteinSmoothContinuous.lean, 44 lines), and for X phi the private ct_level (TR.lie_rt_fin) of constant_term_automorphy. (c) is hper. (d) is CtBasic, proved by eisenstein_constantTerm_of_periodMean_of_upperBlock (65 lines: key from PerMean, the automorphism cb_scale u^-1 and cb_rat_aff); it holds for every continuous function invariant under GL_2(E) on the left. (e) is CtLie, proved by constant_term_automorphy (ConstantTermAutomorphy.lean, 78 lines) from PerDeriv, with TR.hasDerivAt_curve for the chain rule. (f) to (i) are eisenstein_continuous_left_action (EisensteinContinuousLeftAction.lean, 150 lines), which proves CtLeft: L, hsm, key (infGL_left, conj_deriv), hnil (nil_curve), hdec (conj_dec), hv (LieAlg, Lemma sha-siegel-step-seeds-induced-hilbert-form-exists-eisenstein-noncuspidal-constant-term-matrix, with CentLie = eisenstein_central_lie_action = Hm.trzero, and the fields wt and low) and tor_curve. (3) is the private ct_arch with ode (HolEis.lean, 35 lines).

\begin{lemma}\label{sha-siegel-step-seeds-induced-hilbert-form-exists-eisenstein-noncuspidal-characters}
Let $E$ be a totally real number field and $\nu$ a Hecke character of $E$. Write $a(y,x)=\begin{pmatrix}y&x\\0&1\end{pmatrix}$ over any ring and $R(h)$ for right translation by $h\in GL_2(\mathbf A_E)$. Let $\phi$ be a function on $GL_2(\mathbf A_E)$, smooth in the variable at infinity, whose stabilizer in $GL_2(\mathbf A_{E,f})$ under right translation is open. Let $\psi\ne0$ be a function on $GL_2(\mathbf A_E)$, and $\Lambda$ a linear map on the complex functions on $GL_2(\mathbf A_E)$ with $\Lambda(R(h)\phi)=R(h)\psi$ for every $h$. Suppose that $\psi(a(1,x)g)=\psi(g)$ for every adele $x$, $\psi(a(\alpha,0)g)=\psi(g)$ for every $\alpha\in E^\times$, and $\psi(a(y,0)g)=\nu(y)\psi(g)$ for every totally positive $y\in E\otimes\mathbf R$, viewed as an idele. Then there are a Hecke character $\eta$ of $E$ of finite order, a function $\Psi\ne0$ on $GL_2(\mathbf A_E)$ and a linear map $\Lambda'$ on the complex functions on $GL_2(\mathbf A_E)$ with $\Lambda'(R(h)\phi)=R(h)\Psi$ for every $h$, such that $\Psi(a(t,x)g)=(\eta\nu)(t)\Psi(g)$ for every idele $t$, adele $x$ and $g\in GL_2(\mathbf A_E)$.
\end{lemma}
\leanref{Development/AutomorphicArithmetic/AutomorphicVarieties/IdeleGenerators/HolomorphicIrreducibility/HolEis.lean}{UpperTheory.V.Sd.Eis.proj}
\leanref{Development/AutomorphicArithmetic/AutomorphicVarieties/IdeleGenerators/PerDeriv.lean}{UpperTheory.V.Sd.Eis.RayFin}
\leanref{Development/AutomorphicArithmetic/AutomorphicVarieties/BuildingsResidueChars/RayClassGroupFinite.lean}{UpperTheory.ray_class_group_finite}
\leanref{Development/AutomorphicArithmetic/AutomorphicVarieties/IdeleGenerators/PerDeriv.lean}{UpperTheory.V.Sd.Eis.CharPick}
\leanref{Development/AutomorphicArithmetic/AutomorphicVarieties/BuildingsResidueChars/ExistsCharacterSeparation.lean}{UpperTheory.exists_character_separation}
\begin{proof}
For $g\in GL_2(\mathbf A_E)$ let $f_g(t)=\nu(t)^{-1}\psi(a(t,0)g)$, a function on the ideles, and let $H_g$ be the subgroup of ideles $u$ with $f_g(tu)=f_g(t)$ for every $t$. As $a(tu,0)=a(u,0)a(t,0)$, $f_g(tu)=\nu(tu)^{-1}\psi(a(u,0)a(t,0)g)$.
\begin{enumerate}
\item[(a)] \emph{The subgroups $H_g$.} $H_g$ contains $E^\times$, as $\nu(\alpha)=1$ for $\alpha\in E^\times$, and the totally positive elements of $E\otimes\mathbf R$, by the hypotheses on $\psi$. Let $k$ be the finite part of $g$. The finite ideles $u$ with $\nu(u)=1$ and $k^{-1}a(u,0)k$ in the stabilizer of $\phi$ form a neighbourhood of $1$, as $\nu$ is continuous and the stabilizer is open. For such $u$, $a(u,0)g=g\,(k^{-1}a(u,0)k)_f$, and $\psi(g'(k^{-1}a(u,0)k)_f)=\psi(g')$ for every $g'$, by applying $\Lambda$ to $R(h)\phi=\phi$; so $u\in H_g$.
\item[(b)] \emph{Finite index.} A subgroup $H$ of the ideles with the three properties of (a) is open and of finite index. It contains a neighbourhood of $1$, the product of the totally positive elements at infinity and a neighbourhood of $1$ in the finite ideles, so it is open. For $\rho>0$, the square of the idele $z_\rho$ with component $\rho$ at every infinite place and $1$ at every prime lies in $H$, and has norm $\rho^{2[E:\mathbf Q]}$. The ideles of norm $1$ are $E^\times S$ for a compact set $S$, as the idele class group of norm $1$ is compact. So every idele $t$ is $z_\rho^2\alpha s$ with $\rho>0$, $\alpha\in E^\times$ and $s\in S$, and $t\in Hs$. As $H$ is open, the quotient by $H$ is discrete, and the image of the compact $S$ in it is finite. So $H$ has finite index.
\item[(c)] \emph{A character.} Let $T$ be the idele group, $U_0=H_{g_0}$ for some $g_0$ with $\psi(g_0)\ne0$, and $f_0=f_{g_0}$, which is $U_0$-invariant and nonzero, as $f_0(1)=\psi(g_0)$. By (a) and (b), $U_0$ has finite index. By Fourier analysis on the finite abelian group $T/U_0$, there is a character $\chi$ of $T$, trivial on $U_0$, with $\sum_{t\in T/U_0}\chi(t)^{-1}f_0(t)\ne0$. For a subgroup $V\subseteq U_0$ of finite index and a $V$-invariant function $f$, let $a_V(f)=\sum_{t\in T/V}\chi(t)^{-1}f(t)$; then $a_V(R(s)f)=\chi(s)a_V(f)$ for $s\in T$, with $R(s)f(t)=f(ts)$. Let $K$ be the span of the functions $R(s)f-\chi(s)f$, with $f$ invariant under some subgroup of finite index and $s\in T$. Every element of $K$ is invariant under some $V\subseteq U_0$ of finite index with $a_V=0$ on it, while $a_{U_0}(f_0)\ne0$ and $a_V(f_0)=[U_0:V]a_{U_0}(f_0)$; so $f_0\notin K$. Let $L$ be a linear form on the functions on $T$ that vanishes on $K$, with $L(f_0)\ne0$. Then $L(R(s)f)=\chi(s)L(f)$ for every $s$ and every $f$ invariant under a subgroup of finite index.
\item[(d)] \emph{The Hecke character.} $\chi$ is trivial on the open subgroup $U_0$, so continuous, and trivial on $E^\times\subseteq U_0$; so it is a Hecke character $\eta$, of finite order as $\chi^{[T:U_0]}=1$.
\item[(e)] \emph{The function.} Let $\Psi(g)=L(f_g)$; then $\Psi(g_0)\ne0$. Let $\Lambda'(\theta)$ be the function $g\mapsto L\bigl(t\mapsto\nu(t)^{-1}\Lambda(\theta)(a(t,0)g)\bigr)$; then $\Lambda'(R(h)\phi)=R(h)\Psi$. For an idele $t'$ and an adele $x$, $a(t,0)a(t',x)=n(tx)a(tt',0)$ with $n(y)=a(1,y)$, so $f_{a(t',x)g}(t)=\nu(t)^{-1}\psi(a(tt',0)g)=\nu(t')f_g(tt')$. As $f_g$ is invariant under $H_g$, of finite index by (a) and (b), $\Psi(a(t',x)g)=\nu(t')L(R(t')f_g)=\nu(t')\eta(t')\Psi(g)$.
\end{enumerate}
\end{proof}
% proj (private, HolEis.lean, 65 lines): f and H as in the text; (a) and (b) are hH with RayFin, which is ray_class_group_finite (RayClassGroupFinite.lean, 93 lines: hopen, hsq with posIdele, and exists_isCompact_mul_mem_of_ideleModule_mem_Icc, the compactness of the norm-one idele classes); (c) is CharPick, which is exists_character_separation (ExistsCharacterSeparation.lean, 185 lines: the Fourier coefficient exists_addChar_fourier_coeff_ne_zero_of_ne_zero, the averages aV, kill, the span Kset and Submodule.exists_dual_map_eq_bot_of_notMem); (d) is hc, hd and the order (H g0).index; (e) is M composed with Lambda and the last computation e. About 350 lines.

\begin{lemma}\label{sha-siegel-step-seeds-induced-hilbert-form-exists-eisenstein-noncuspidal-local}
Let $E$ be a number field, $w$ a prime of $E$, $\mathcal O_w$ the ring of integers of $E_w$, $\varpi_w$ a uniformizer and $q=\operatorname N(w)$. Let $\mu_1,\mu_2$ be homomorphisms from $E_w^\times$ to $\mathbf C^\times$, $V$ a complex vector space and $\Phi:GL_2(E_w)\to V$ a nonzero function with
\[
\Phi\Bigl(\begin{pmatrix}a&dx\\0&d\end{pmatrix}g\Bigr)=\mu_1(a)\mu_2(d)\Phi(g)\qquad(a,d\in E_w^\times,\ x\in E_w,\ g\in GL_2(E_w)).
\]
Let $c\ge0$ and $\lambda\in\mathbf C$, and suppose that $\Phi$ is fixed by $K_1(w^c)$ under right translation and $U_w\Phi=\lambda\Phi$, with $K_1(w^c)$ and $U_w$ as in Lemma~\ref{sha-siegel-step-seeds-induced-hilbert-form-exists-eisenstein-lift-hecke-cosets}.
\begin{enumerate}
\item If $c=0$, then $\mu_1$ and $\mu_2$ are trivial on $\mathcal O_w^\times$ and $\lambda=q\mu_1(\varpi_w)+\mu_2(\varpi_w)$.
\item If $c\ne0$, then $\lambda=0$, or $\mu_1$ is trivial on $\mathcal O_w^\times$ and $\lambda=q\mu_1(\varpi_w)$, or $\mu_2$ is trivial on $\mathcal O_w^\times$ and $\lambda=\mu_2(\varpi_w)$.
\end{enumerate}
\end{lemma}
\leanref{Development/AutomorphicArithmetic/AutomorphicVarieties/IdeleGenerators/WeightChange.lean}{UpperTheory.V.Sd.LocEig}
\leanref{Development/AutomorphicArithmetic/AutomorphicVarieties/IdeleGenerators/HolomorphicIrreducibility/LocEig.lean}{UpperTheory.V.Sd.locEig}
\begin{proof}
Write $a(y,x)=\begin{pmatrix}y&x\\0&1\end{pmatrix}$, $l(x)=\begin{pmatrix}1&0\\x&1\end{pmatrix}$, $s=\begin{pmatrix}0&1\\1&0\end{pmatrix}$, $\varpi=\varpi_w$, and call the matrices $\begin{pmatrix}a&dx\\0&d\end{pmatrix}$ \emph{Borel}. Let $S\subset\mathcal O_w$ be a set of representatives of $\mathcal O_w/\varpi\mathcal O_w$; it has $q$ elements. The matrices $l(x)$ with $x\in\varpi^c\mathcal O_w$ and $a(u,y)$ with $u\in\mathcal O_w^\times$ and $y\in\mathcal O_w$ lie in $K_1(w^c)$, and so does $s$ if $c=0$; so $\Phi(g\,l(x))=\Phi(g)$ and $\Phi(g\,a(u,y))=\Phi(g)$ for these, and $\Phi(gs)=\Phi(g)$ if $c=0$. By Lemma~\ref{sha-siegel-step-seeds-induced-hilbert-form-exists-eisenstein-lift-hecke-cosets},
\[
\sum_{u\in S}\Phi(g\,a(\varpi,u))+\epsilon\,\Phi(g\operatorname{diag}(1,\varpi))=\lambda\Phi(g)\qquad(g\in GL_2(E_w)),
\]
with $\epsilon=1$ if $c=0$ and $\epsilon=0$ otherwise.
\begin{enumerate}
\item[(a)] $\Phi(a(y,x))=\mu_1(y)\Phi(1)$ and $\Phi(\operatorname{diag}(1,d)g)=\mu_2(d)\Phi(g)$. So $\sum_{u\in S}\Phi(a(\varpi,u))=q\mu_1(\varpi)\Phi(1)$.
\item[(b)] If $\Phi(1)\ne0$, $\mu_1$ is trivial on $\mathcal O_w^\times$: for a unit $u$, $\Phi(1)=\Phi(a(u,0))=\mu_1(u)\Phi(1)$. If $\Phi(s)\ne0$, $\mu_2$ is trivial on $\mathcal O_w^\times$: $s\,a(u,0)=\operatorname{diag}(1,u)s$, so $\Phi(s)=\Phi(s\,a(u,0))=\mu_2(u)\Phi(s)$.
\item[(c)] If $g_{22}\ne0$, then $g$ is a Borel matrix times $l(g_{21}/g_{22})$; so $\Phi(g)=0$ if $\Phi(l(g_{21}/g_{22}))=0$. If $g_{21}\ne0$, then $g$ is a Borel matrix times $s\,a(1,g_{22}/g_{21})$.
\item[(d)] If $\Phi(l(x))=0$ for every $x\in\varpi\mathcal O_w$, then $\Phi(s)\ne0$. Otherwise $\Phi=0$: let $g\in GL_2(E_w)$. If $g_{21}=0$, then $g_{22}\ne0$ and $\Phi(g)=0$ by (c), with $x=0$. If $g_{21}\ne0$ and $y=g_{22}/g_{21}\in\mathcal O_w$, then by (c) $\Phi(g)$ is a multiple of $\Phi(s\,a(1,y))=\Phi(s)=0$. Otherwise $g_{21}/g_{22}\in\varpi\mathcal O_w$ and $\Phi(g)=0$ by (c).
\end{enumerate}
\emph{The case $c=0$.} Then $\Phi(s)=\Phi(1)$. If $\Phi(1)=0$, then $\Phi(l(x))=\Phi(1)=0$ for every $x\in\mathcal O_w$, so $\Phi(s)\ne0$ by (d), which is false. So $\Phi(1)=\Phi(s)\ne0$, and $\mu_1,\mu_2$ are trivial on $\mathcal O_w^\times$ by (b). At $g=1$, the identity for $U_w$ and (a) give $q\mu_1(\varpi)\Phi(1)+\mu_2(\varpi)\Phi(1)=\lambda\Phi(1)$.

\emph{The case $c\ne0$.} Suppose that $\lambda\ne0$. If $\Phi(1)\ne0$, then (a), (b) and the identity at $g=1$ give that $\mu_1$ is trivial on $\mathcal O_w^\times$ and $\lambda=q\mu_1(\varpi)$. Suppose that $\Phi(1)=0$. Then $\Phi(l(x))=\Phi(1)=0$ for $x\in\varpi^c\mathcal O_w$. If $\Phi(l(x'))=0$ for every $x'\in\varpi^{j+2}\mathcal O_w$, then $\Phi(l(x))=0$ for every $x\in\varpi^{j+1}\mathcal O_w$: for $u\in S$, $l(x)a(\varpi,u)=\begin{pmatrix}\varpi&u\\x\varpi&xu+1\end{pmatrix}$, and $xu+1$ is a unit, so by (c) $\Phi(l(x)a(\varpi,u))=0$, as $x\varpi/(xu+1)\in\varpi^{j+2}\mathcal O_w$; the identity at $g=l(x)$ gives $\lambda\Phi(l(x))=0$. By descending induction from $c\ge1$, $\Phi(l(x))=0$ for every $x\in\varpi\mathcal O_w$. By (d), $\Phi(s)\ne0$, and $\mu_2$ is trivial on $\mathcal O_w^\times$ by (b). Take $g=s$ in the identity. Let $u_0\in S$ be the representative of $0$, $u_0=\varpi y$; then $s\,a(\varpi,u_0)=\operatorname{diag}(1,\varpi)s\,a(1,y)$, so its term is $\mu_2(\varpi)\Phi(s\,a(1,y))=\mu_2(\varpi)\Phi(s)$. For a unit $u\in S$, $s\,a(\varpi,u)=\begin{pmatrix}0&1\\\varpi&u\end{pmatrix}$, with $\varpi/u\in\varpi\mathcal O_w$, so its term is $0$ by (c). So $\lambda\Phi(s)=\mu_2(\varpi)\Phi(s)$, and $\lambda=\mu_2(\varpi)$.
\end{proof}
% locEig (LocEig.lean, 321 lines) applies the private core to the valuation ring of E_w, the uniformizer atFin w (unif w) and the memberships of low, aff and sw in K1 w c (mem_parahoricTimesPrincipalCongruence_iff), with uopSum for the identity. In core: (a) is hA, hD and sum1; (b) is u1 and u2; (c) is z1 with dec1 and dec2; (d) is iw; the case c = 0 is the first constructor; the case c != 0 is step, down, hz and the computation at sw K. The Lean's Borel matrices are scalar d times aff (a/d) x (bor_val).

\begin{lemma}\label{sha-siegel-step-seeds-induced-hilbert-form-exists-eisenstein-noncuspidal}
Let $E$ be a totally real number field with uniformizers $\varpi_w$, $k\ge1$, $c$ a function from the primes of $E$ to $\mathbf Z_{\ge0}$ that vanishes at all but finitely many primes (this hypothesis is not used), $\omega$ a Hecke character of $E$ of finite order, and $\phi$ a holomorphic eigenform of parallel weight $k$, level $c$ and central character $\omega$ with eigenvalues $\lambda(w)$ (Definition~\ref{sha-siegel-step-seeds-induced-hilbert-form-exists-holomorphic}). Suppose that $\phi$ is not a cusp form. Then there are Hecke characters $\eta_1,\eta_2$ of $E$ of finite order with $\eta_1\eta_2=\omega$ such that, for every prime $w$:
\begin{enumerate}
\item if $c(w)=0$, then $\eta_1$ and $\eta_2$ are unramified at $w$ and $\lambda(w)=\eta_1(\varpi_w)+\eta_2(\varpi_w)\operatorname N(w)^{k-1}$;
\item if $c(w)\ne0$, then $\lambda(w)=0$, or $\eta_1$ is unramified at $w$ and $\lambda(w)=\eta_1(\varpi_w)$, or $\eta_2$ is unramified at $w$ and $\lambda(w)=\eta_2(\varpi_w)\operatorname N(w)^{k-1}$.
\end{enumerate}
\end{lemma}
\leanref{Development/AutomorphicArithmetic/AutomorphicVarieties/IdeleGenerators/WeightChange.lean}{UpperTheory.V.Sd.HolEis}
\leanref{Development/AutomorphicArithmetic/AutomorphicVarieties/IdeleGenerators/HolomorphicIrreducibility/HolEis.lean}{UpperTheory.V.Sd.holEis}
\begin{proof}
Write $a(y,x)=\begin{pmatrix}y&x\\0&1\end{pmatrix}$ and $n(x)=a(1,x)$ over any ring, $R(h)$ for right translation by $h$, and $q_w=\operatorname N(w)$. Say that a function $\psi$ on $GL_2(\mathbf A_E)$ is \emph{linked} to $\phi$ if there is a linear map $\Lambda$ on the complex functions on $GL_2(\mathbf A_E)$ with $\Lambda(R(h)\phi)=R(h)\psi$ for every $h\in GL_2(\mathbf A_E)$. If $\psi$ is linked to $\phi$, $h\in GL_2(\mathbf A_E)$, $z\in\mathbf C$ and $\phi(gh)=z\phi(g)$ for every $g$, then $\psi(gh)=z\psi(g)$ for every $g$, by applying $\Lambda$ to $R(h)\phi=z\phi$.
\begin{enumerate}
\item[(1)] \emph{The constant term.} As $\phi$ is not a cusp form, there are a Haar measure $\mu$ on $\mathbf A_E$, a measurable fundamental domain $D\subseteq\mathbf A_E$ for the translations by $E$ and $g_0\in GL_2(\mathbf A_E)$ with $\psi(g_0)\ne0$, where $\psi(g)=\int_D\phi(n(x)g)\,d\mu(x)$. By Lemma~\ref{sha-siegel-step-seeds-induced-hilbert-form-exists-eisenstein-noncuspidal-constant-term}, these integrals converge, $\psi(n(x)g)=\psi(g)$ for $x\in\mathbf A_E$, $\psi(a(\alpha,0)g)=\psi(g)$ for $\alpha\in E^\times$, and $\psi(a(y,0)_\infty g)=(\prod_vy_v)^{k/2}\psi(g)$ for totally positive $y\in E\otimes\mathbf R$.
\item[(2)] \emph{Linking.} The map $\theta\mapsto\bigl(g\mapsto\int_D\theta(n(x)g)\,d\mu(x)\bigr)$, on the space of the functions $\theta$ for which these integrals converge, is linear and commutes with right translations; this space contains every $R(h)\phi$. Extending it to a linear map on all functions shows that $\psi$ is linked to $\phi$.
\item[(3)] \emph{A power of the norm.} Let $\nu$ be the Hecke character $t\mapsto|t|^{k/2}$, with $|t|$ the norm of the idele $t$. Then $\nu(y)=(\prod_vy_v)^{k/2}$ for $y\in E\otimes\mathbf R$ totally positive, viewed as an idele, and at every prime $w$, $\nu$ is unramified and $\nu(\varpi_w)=q_w^{-k/2}$.
\item[(4)] \emph{A character.} The stabilizer of $\phi$ in $GL_2(\mathbf A_{E,f})$ is open. By (1), (2), (3) and Lemma~\ref{sha-siegel-step-seeds-induced-hilbert-form-exists-eisenstein-noncuspidal-characters}, there are a Hecke character $\eta$ of finite order and a nonzero function $\Psi$ linked to $\phi$ with $\Psi(a(t,x)g)=(\eta\nu)(t)\Psi(g)$ for all ideles $t$, adeles $x$ and $g\in GL_2(\mathbf A_E)$.
\item[(5)] \emph{The central character.} For a scalar matrix $z$, $\phi(gz)=\phi(zg)=\omega(z)\phi(g)$; as $\Psi$ is linked to $\phi$, $\Psi(zg)=\Psi(gz)=\omega(z)\Psi(g)$.
\item[(6)] \emph{Local functions.} Let $w$ be a prime and $\lambda'(w)=q_w^{1-k/2}\lambda(w)$, so that $U_w\phi=\lambda'(w)\phi$. Choose $g_1$ with $\Psi(g_1)\ne0$, let $k_1\in GL_2(E_w)$ be its component at $w$, and let $g'$ be $g_1$ with its component at $w$ replaced by $1$; so $g_1=(k_1)_wg'$, and $h_w$ commutes with $g'$ for $h\in GL_2(E_w)$. Let $\Phi(h)=\Psi(h_wg')$ for $h\in GL_2(E_w)$; then $\Phi(k_1)=\Psi(g_1)\ne0$. Let $\Lambda'$ link $\Psi$ to $\phi$, and $M(\theta)$ the function $h\mapsto\Lambda'(\theta)(h_wg')$. Then $M(R(h_w)\phi)=R(h)\Phi$ for $h\in GL_2(E_w)$, with $R(h)$ right translation on $GL_2(E_w)$, and $M(\phi)=\Phi$. So $\Phi$ is fixed by $K_1(w^{c(w)})$, and by Lemma~\ref{sha-siegel-step-seeds-induced-hilbert-form-exists-eisenstein-lift-hecke-cosets}, applied to $\phi$ and to $\Phi$ with the same representatives, $U_w\Phi=M(U_w\phi)=\lambda'(w)\Phi$. As $\begin{pmatrix}a&dx\\0&d\end{pmatrix}=d\,a(a/d,x)$, (4) and (5) give
\[
\Phi\Bigl(\begin{pmatrix}a&dx\\0&d\end{pmatrix}h\Bigr)=\omega(d)\,(\eta\nu)(a/d)\,\Phi(h)=\mu_1(a)\mu_2(d)\Phi(h),
\]
with the characters $\mu_1=(\eta\nu)_w$ and $\mu_2=(\omega\eta^{-1}\nu^{-1})_w$ of $E_w^\times$.
\item[(7)] \emph{The local eigenvalues.} By Lemma~\ref{sha-siegel-step-seeds-induced-hilbert-form-exists-eisenstein-noncuspidal-local}, if $c(w)=0$, then $\mu_1,\mu_2$ are trivial on $\mathcal O_{E_w}^\times$ and $\lambda'(w)=q_w\mu_1(\varpi_w)+\mu_2(\varpi_w)$; if $c(w)\ne0$, then $\lambda'(w)=0$, or $\mu_1$ is trivial on $\mathcal O_{E_w}^\times$ and $\lambda'(w)=q_w\mu_1(\varpi_w)$, or $\mu_2$ is trivial on $\mathcal O_{E_w}^\times$ and $\lambda'(w)=\mu_2(\varpi_w)$.
\item[(8)] \emph{Conclusion.} Let $\eta_1=\eta$ and $\eta_2=\omega\eta^{-1}$, Hecke characters of finite order with $\eta_1\eta_2=\omega$. As $\nu$ is unramified at $w$, $\mu_1$ and $\mu_2$ are trivial on $\mathcal O_{E_w}^\times$ if and only if $\eta_1$ and $\eta_2$ are unramified at $w$. As $\nu(\varpi_w)=q_w^{-k/2}$ and $k\ge1$, $q_w^{k/2-1}\cdot q_w\mu_1(\varpi_w)=\eta_1(\varpi_w)$ and $q_w^{k/2-1}\mu_2(\varpi_w)=\eta_2(\varpi_w)q_w^{k-1}$. Multiplying the alternatives of (7) by $q_w^{k/2-1}$ gives (1) and (2) of the statement for $\lambda(w)=q_w^{k/2-1}\lambda'(w)$.
\end{enumerate}
\end{proof}
% holEis = Sd.Eis.main (HolEis.lean, 303 lines with its private lemmas) applied to its six inputs, followed step by step: (1) is the negation of IsCuspidal (r = 1, mu, D, g0) with CtBasic (eisenstein_constantTerm_of_periodMean_of_upperBlock), SmCont (eisenstein_smooth_continuous) and ct_arch with CtLeft (eisenstein_continuous_left_action), Lemma sha-siegel-step-seeds-induced-hilbert-form-exists-eisenstein-noncuspidal-constant-term; (2) is ct_rel, the text's linking is the private Rel and its consequence Rel.scal; (3) is AbsPow (eisenstein_norm_eq_abs_pow, 34 lines: absChar F s); (4) is proj with RayFin (ray_class_group_finite) and CharPick (exists_character_separation), Lemma sha-siegel-step-seeds-induced-hilbert-form-exists-eisenstein-noncuspidal-characters; (5) is hZ; (6) is loc with uopSum; (7) is locEig (Lemma sha-siegel-step-seeds-induced-hilbert-form-exists-eisenstein-noncuspidal-local) inside key; (8) is num, u1, u2 and the final refine, with Nat.cast_sub hk. The hypothesis that c vanishes at all but finitely many primes is not used. About 1,200 lines with the dependencies, of which about 550 are in Lemma sha-siegel-step-seeds-induced-hilbert-form-exists-eisenstein-noncuspidal-constant-term and the two lemmas before it.

\begin{lemma}\label{sha-siegel-step-seeds-induced-hilbert-form-exists-eisenstein-chebotarev}
Let $E$ be a number field with uniformizers $\varpi_w$ at its primes, $\ell$ a prime and $\iota:\bar{\mathbf Q}_\ell\cong\mathbf C$, and write $x\equiv y$ for complex numbers with $|\iota^{-1}(x-y)|<1$. Let $\eta_1,\eta_2$ be Hecke characters of $E$ of finite order and $S$ a finite set of primes of $E$ such that, for every prime $w\notin S$, $\eta_1$ and $\eta_2$ are unramified at $w$ and $\eta_1(\varpi_w)\equiv1$ or $\eta_2(\varpi_w)\equiv1$ (Definition~\ref{sha-siegel-step-seeds-induced-hilbert-form-exists-eisenstein-expansion}). Then $\eta_1(x)\equiv1$ or $\eta_2(x)\equiv1$ for every idele $x$ of $E$.
\end{lemma}
\leanref{Development/AutomorphicArithmetic/AutomorphicVarieties/ModularityArgument/ThresholdArguments/DenseThresholdCases/MPropMHol.lean}{gal_of}
\begin{proof}
Let $G_E=\operatorname{Gal}(\bar E/E)$. As $\eta_i$ has finite order, global class field theory \cite[Chapter~VI]{Neu99} gives a continuous character $\delta_i:G_E\to\bar{\mathbf Q}_\ell^\times$ with $\delta_i(\sigma)=\iota^{-1}(\eta_i(x))$ whenever a reciprocity map sends the idele $x$ to the image of $\sigma$ in $G_E^{\mathrm{ab}}$; at a prime $w$ at which $\eta_i$ is unramified, $\delta_i(\sigma)=\iota^{-1}(\eta_i(\varpi_w))^{-1}$ for every Frobenius element $\sigma$ at a place above $w$. Let $C$ be the set of $\sigma\in G_E$ with $\delta_1(\sigma)\equiv1$ or $\delta_2(\sigma)\equiv1$. It is closed, as the $\delta_i$ are continuous and $\{z:|z-1|<1\}$ is closed in $\bar{\mathbf Q}_\ell$. Let $\sigma$ be a Frobenius element at a place above a prime $w\notin S$, and $i$ such that $\eta_i(\varpi_w)\equiv1$. Then $z=\iota^{-1}(\eta_i(\varpi_w))$ has $|z|=1$, and $z^{-1}-1=-z^{-1}(z-1)$, so $\delta_i(\sigma)=z^{-1}\equiv1$ and $\sigma\in C$. These Frobenius elements, at the primes of degree one outside $S$, are dense in $G_E$ (Theorem~\ref{lit-chebotarev-degree-one}), so $C=G_E$. For an idele $x$, let $\sigma\in G_E$ map to the image of $x$ under a reciprocity map. Then $\iota^{-1}(\eta_i(x))=\delta_i(\sigma)$ for $i=1,2$, and $\sigma\in C$.
\end{proof}
% gal_of (MPropMHol.lean, 25 lines, private): the Artin map art of hA (GlobalArtinReciprocity, from HypothesesH' through litN.aut.art), exists_isLAdicOf with isAlgPos_zero_of_orderDvd for delta_i, the closed set hcl, Published.ChebotarevDegreeOne.dense with chebotarevDegreeOne', IsLAdicOf.frob_unr and inv_cong at the Frobenius elements, and toAb_surjective with the field of IsLAdicOf at x.

\begin{lemma}\label{sha-siegel-step-seeds-induced-hilbert-form-exists-eisenstein}
Assume the Galois representations as in Theorem~\ref{sha-residually-reducible-automorphy}. Let $E$ be a totally real number field, $\ell$ an odd prime and $\iota:\bar{\mathbf Q}_\ell\cong\mathbf C$. Let $\omega$ be a Hecke character of $E$ of finite order with $\omega_v(-1)=-1$ at every real place $v$. For each prime $w$ of $E$ let $\varpi_w$ be a uniformizer of $E_w$, and let $c(w)$ be the conductor exponent of $\omega$ at $w$: the least $n\ge0$ such that $\omega_w$ is trivial on the units of $E_w$ congruent to $1$ modulo $w^n$. Suppose that $c(w)\ne0$ for some $w$, and that $c(w)=0$ for every $w$ above $\ell$. Let $P$ be a set of primes of $E$, and suppose that there are primes $w_1\in P$ and $w_2\notin P$ above $\ell$ with $|\iota^{-1}(\omega_{w_i}(\varpi_{w_i}))-1|\ge1$ for $i=1,2$. Let $c_\ell(w)=1$ for the primes $w$ above $\ell$ and $c_\ell(w)=c(w)$ for the others, and let
\[
e(w)=\begin{cases}\omega_w(\varpi_w)&\text{if $w$ lies above $\ell$ and $w\in P$,}\\1&\text{if $w$ lies above $\ell$ and $w\notin P$,}\\1+\omega_w(\varpi_w)&\text{if $w$ does not lie above $\ell$ and $c(w)=0$,}\\1&\text{if $w$ does not lie above $\ell$ and $c(w)\ne0$.}\end{cases}
\]
Then for every $N_0\ge0$ there are an integer $k\ge N_0$ divisible by $\ell-1$ and a holomorphic eigenform of parallel weight $k+1$, level $c_\ell$ and central character $\omega$ (Definition~\ref{sha-siegel-step-seeds-induced-hilbert-form-exists-holomorphic}) that is a cusp form, with eigenvalues $\lambda(w)$ such that $|\iota^{-1}(\lambda(w)-e(w))|<1$ for every prime $w$.
\end{lemma}
\leanref{Development/AutomorphicArithmetic/AutomorphicVarieties/IdeleGenerators/WeightChange.lean}{UpperTheory.V.Sd.PropMHol}
\leanref{Development/AutomorphicArithmetic/AutomorphicVarieties/ModularityArgument/ThresholdArguments/DenseThresholdCases/MPropMHol.lean}{UpperTheory.V.Z.Bc.propMHolZ_alt}
\leanref{Development/AutomorphicArithmetic/AutomorphicVarieties/IdeleGenerators/WeightChange.lean}{UpperTheory.V.Sd.levL}
\leanref{Development/AutomorphicArithmetic/AutomorphicVarieties/IdeleGenerators/WeightChange.lean}{UpperTheory.V.Sd.eisT}
\leanref{Development/AutomorphicArithmetic/AutomorphicVarieties/IdeleGenerators/IsCondAt.lean}{UpperTheory.V.Sd.IsCondAt}
\begin{proof}
Write $x\equiv y$ for complex numbers with $|\iota^{-1}(x-y)|<1$; call $x$ \emph{integral} if $|\iota^{-1}(x)|\le1$ and a \emph{unit} if $|\iota^{-1}(x)|=1$. The values of Hecke characters of finite order are roots of unity, hence units. Fourier expansions and the functions $\chi_{c,\omega}$ and $T^k_w$ are those of Definition~\ref{sha-siegel-step-seeds-induced-hilbert-form-exists-eisenstein-expansion}, for the uniformizers $\varpi_w$, and $T^k_w$ is taken for $c_\ell$ and $\omega$. Let $w_1\in P$ and $w_2\notin P$ be the primes above $\ell$ of the hypothesis.
\begin{enumerate}
\item[(1)] \emph{A form of weight one.} As $w_1\in P$ lies above $\ell$, Lemma~\ref{sha-siegel-step-seeds-induced-hilbert-form-exists-eisenstein-stabilized} gives a holomorphic form $\phi_1$ of parallel weight $1$, level $c_\ell$ and central character $\omega$ with the Fourier expansion of weight $1$ with constant term $0$ and integral coefficients $a$, with $a(\mathcal O_E)=1$ and $T^1_wa(\mathfrak m)=e(w)a(\mathfrak m)$ for every prime $w$ and every $\mathfrak m$.
\item[(2)] \emph{A form congruent to $1$.} Let $N$ be as in Lemma~\ref{sha-siegel-step-seeds-induced-hilbert-form-exists-eisenstein-congruent} and $k=(\ell-1)\ell^N(N_0+1)$. Then $k\ge N_0$, $k>0$ and $\ell-1$ divides $k$, and that lemma gives a holomorphic form $\psi$ of parallel weight $k$, level $0$ and trivial central character with the Fourier expansion of weight $k$ with a constant term $d_0\equiv1$ and coefficients $b\equiv0$.
\item[(3)] \emph{The product.} Lemma~\ref{sha-siegel-step-seeds-induced-hilbert-form-exists-eisenstein-product}, applied to $\phi_1$ and $\psi$, gives a holomorphic form $\phi_2$ of parallel weight $k+1$, level $c_\ell$ and central character $\omega$ with the Fourier expansion of weight $k+1$ with constant term $0$ and coefficients $a'$ with $a'(\mathfrak m)\equiv a(\mathfrak m)$ for every $\mathfrak m$.
\item[(4)] \emph{The level.} If $c_\ell(w)\ne0$, then $w$ lies above $\ell$ or $\omega$ is ramified at $w$; so $c_\ell$ vanishes at all but finitely many primes. If $c_\ell(w)=0$, then $w$ does not lie above $\ell$ and $c(w)=0$; as $\ell$ does not divide $\operatorname N(w)$ and $\ell-1$ divides $k$, $\operatorname N(w)^k\equiv1$ by Fermat's little theorem.
\item[(5)] \emph{Congruent coefficients.} The numbers $e(w)$, $\chi_{c_\ell,\omega}(w)$, $a(\mathfrak m)$ and $a'(\mathfrak m)$ are integral, the last by (3) and the ultrametric inequality, and $a'(\mathcal O_E)\equiv a(\mathcal O_E)=1$ is a unit. Let $w$ be a prime and $\mathfrak m$ a nonzero integral ideal. If $c_\ell(w)=0$, then by (3) and (4),
\[
T^{k+1}_wa'(\mathfrak m)=a'(w\mathfrak m)+\chi_{c_\ell,\omega}(w)\operatorname N(w)^ka'(\mathfrak mw^{-1})\equiv a(w\mathfrak m)+\chi_{c_\ell,\omega}(w)a(\mathfrak mw^{-1})=T^1_wa(\mathfrak m).
\]
If $c_\ell(w)\ne0$, then $\chi_{c_\ell,\omega}(w)=0$ and $T^{k+1}_wa'(\mathfrak m)=a'(w\mathfrak m)\equiv a(w\mathfrak m)=T^1_wa(\mathfrak m)$. In both cases, by (1), $T^{k+1}_wa'(\mathfrak m)\equiv e(w)a(\mathfrak m)\equiv e(w)a'(\mathfrak m)$.
\item[(6)] \emph{An eigenform.} By (3), (4) and (5), Lemma~\ref{sha-siegel-step-seeds-induced-hilbert-form-exists-eisenstein-lift} applies to $\phi_2$, of weight $k+1\ge1$, with $t=e$. It gives a holomorphic eigenform $g$ of parallel weight $k+1$, level $c_\ell$ and central character $\omega$ with eigenvalues $\lambda(w)\equiv e(w)$ for every prime $w$.
\item[(7)] \emph{Cuspidality.} Suppose that $g$ is not a cusp form. By (4), Lemma~\ref{sha-siegel-step-seeds-induced-hilbert-form-exists-eisenstein-noncuspidal} gives Hecke characters $\eta_1,\eta_2$ of $E$ of finite order with $\eta_1\eta_2=\omega$ and the properties (1) and (2) there, for $k+1$.
\begin{enumerate}
\item[(a)] Let $w$ be a prime with $c_\ell(w)=0$. Then $\eta_1$ and $\eta_2$ are unramified at $w$, and $\lambda(w)=\eta_1(\varpi_w)+\eta_2(\varpi_w)\operatorname N(w)^k\equiv\eta_1(\varpi_w)+\eta_2(\varpi_w)$ by (4). Also $\lambda(w)\equiv e(w)=1+\omega(\varpi_w)=1+\eta_1(\varpi_w)\eta_2(\varpi_w)$. So $(\eta_1(\varpi_w)-1)(\eta_2(\varpi_w)-1)\equiv0$, and as the absolute value is multiplicative, $\eta_1(\varpi_w)\equiv1$ or $\eta_2(\varpi_w)\equiv1$. By (4) and Lemma~\ref{sha-siegel-step-seeds-induced-hilbert-form-exists-eisenstein-chebotarev}, $\eta_1(x)\equiv1$ or $\eta_2(x)\equiv1$ for every idele $x$ of $E$.
\item[(b)] Let $w$ be a prime above $\ell$ with $e(w)$ a unit. Then $c_\ell(w)=1$, and $\lambda(w)\equiv e(w)$ is a unit, so $\lambda(w)\ne0$. As $\ell$ divides $\operatorname N(w)$ and $k>0$, $\eta_2(\varpi_w)\operatorname N(w)^k$ is not a unit, so $\lambda(w)$ is not this number. By (2) of Lemma~\ref{sha-siegel-step-seeds-induced-hilbert-form-exists-eisenstein-noncuspidal}, $\lambda(w)=\eta_1(\varpi_w)$, and so $\eta_1(\varpi_w)\equiv e(w)$.
\item[(c)] At $w_1$, $e(w_1)=\omega(\varpi_{w_1})=\eta_1(\varpi_{w_1})\eta_2(\varpi_{w_1})$ is a unit; by (b), $\eta_1(\varpi_{w_1})(1-\eta_2(\varpi_{w_1}))\equiv0$, so $\eta_2(\varpi_{w_1})\equiv1$, as $\eta_1(\varpi_{w_1})$ is a unit. At $w_2$, $e(w_2)=1$, so $\eta_1(\varpi_{w_2})\equiv1$ by (b).
\item[(d)] Apply (a) to the idele with components $\varpi_{w_1}$ at $w_1$, $\varpi_{w_2}$ at $w_2$ and $1$ elsewhere: $\eta_1(\varpi_{w_1})\eta_1(\varpi_{w_2})\equiv1$ or $\eta_2(\varpi_{w_1})\eta_2(\varpi_{w_2})\equiv1$. In the first case, $\eta_1(\varpi_{w_1})\equiv1$ by (c), and $\omega(\varpi_{w_1})=\eta_1(\varpi_{w_1})\eta_2(\varpi_{w_1})\equiv1$ by (c), against the hypothesis on $w_1$. In the second case, $\eta_2(\varpi_{w_2})\equiv1$ by (c), and $\omega(\varpi_{w_2})=\eta_1(\varpi_{w_2})\eta_2(\varpi_{w_2})\equiv1$ by (c), against the hypothesis on $w_2$.
\end{enumerate}
So $g$ is a cusp form. With (2) and (6), $k$ and $g$ have the properties of the statement.
\end{enumerate}
\end{proof}
% propMHolZ_alt (MPropMHol.lean, 117 lines, with the private sm_add, sm_sub, sm_mul, le_of_cong, eq_of_cong, inv_cong, norm_nat_lt, pow_cong, dvd_normV, norm_val and chiI_prime), followed step by step. Hypotheses4Z' is HypothesesH' with the published inputs Published.HilbertEisOne (PP.eis, Lemma sha-siegel-step-seeds-induced-hilbert-form-exists-eisenstein-series) and Lit.OneUse (PP.one, Lemma sha-siegel-step-seeds-induced-hilbert-form-exists-eisenstein-congruent with oneUse_of), which the Lean proves. (1) is emix PP.eis; (2) is PP.one; (3) is holMul; (4) is hfin (Ideal.finite_factors, UnrAe.unr_ae_and_level, IsCondAt) and hpow (pow_cong); (5) is the congruence passed to holLift_alt (hkCoef, chiI_prime, sm_mul, sm_add); (6) is Sd.holLift_alt PP.h; (7) is the by_contra: holEis for (7), key = gal_of PP.h.litN.aut.art for (a), hab for (b), hy for (c) and kk with the final rcases for (d). Size: about 3,400 lines with its dependencies, net of the definitions and of the two published inputs; its parts are the lemmas before this one.

\subsection{Ordinary Hilbert modular forms}\label{subsection-sha-ordinary-hilbert-forms}

\begin{definition}\label{sha-siegel-step-seeds-induced-hilbert-form-exists-weight-model}
Let $E$ be a totally real number field with set of infinite places $\Sigma$, and $c$ a function from the primes of $E$ to $\mathbf Z_{\ge0}$ that vanishes at all but finitely many primes. Let $K_c\subseteq GL_2(\mathbf A_{E,f})$ be the group of elements whose component at every prime $w$ lies in $K_1(w^{c(w)})$ (Definition~\ref{sha-siegel-step-seeds-induced-hilbert-form-ordinary}), and $H_c$ the subgroup generated by $K_c$ and the scalar matrices. Let $\mathfrak P=\mathbf C\setminus\mathbf R$, on which $GL_2(\mathbf R)$ acts by Möbius transformations, and for $n\in\mathbf Z_{\ge0}$, $g\in GL_2(\mathbf R)$ and $z\in\mathfrak P$ let $j_n(g,z)=(g_{21}z+g_{22})^n|\det g|^{-n/2}$. Let $\kappa$ be a function from $\Sigma$ to $\mathbf Z_{\ge0}$.
\begin{enumerate}
\item Let $\chi$ be a character of $H_c$. The space $M_\kappa(c,\chi)$ is the span of the functions $f$ on $\mathfrak P^\Sigma\times GL_2(\mathbf A_{E,f})$ such that: $z\mapsto f(z,x)$ is continuous and holomorphic in each variable, for every $x$; $f(\gamma z,\gamma x)=\prod_vj_{\kappa(v)}(v(\gamma),z_v)\,f(z,x)$ for $\gamma\in GL_2(E)$, which acts on $z$ through the real embeddings $v$; $f(z,xh)=\chi(h)f(z,x)$ for $h\in H_c$; and $|f(z,x)|\prod_v|\operatorname{Im}z_v|^{\kappa(v)/2}$ is bounded.
\item A function $f$ on $\mathfrak P^\Sigma\times GL_2(\mathbf A_{E,f})$ \emph{corresponds} in weight $\kappa$ to a function $\phi$ on $GL_2(\mathbf A_E)$ if $\phi(g_\infty x)\prod_vj_{\kappa(v)}(g_v,i)=f(g_\infty i,x)$ for all $g_\infty=(g_v)\in GL_2(E\otimes\mathbf R)$ and $x\in GL_2(\mathbf A_{E,f})$, where $g_\infty i$ has components $g_v\cdot i$.
\item Fix a uniformizer $\varpi_w$ of $E_w$ at every prime $w$. For a finite sequence $W=(w_1,\dots,w_m)$ of primes, \emph{representatives} for $W$ are elements $a_1,\dots,a_n$ of $GL_2(\mathbf A_{E,f})$ such that: for every representation $\rho$ of $GL_2(\mathbf A_{E,f})$ and every vector $x$ fixed by $K_c$, $\sum_i\rho(a_i)x$ is fixed by $K_c$ and equals $U_{w_1}\cdots U_{w_m}x$, where $U_w$ acts through $GL_2(E_w)$ on the vectors fixed by $K_1(w^{c(w)})$; $|\det a_i|=\prod_j\operatorname N(w_j)^{-1}$ for every $i$; and at every prime $w$ with $c(w)\ne0$, the component of every $a_i$ at $w$ is $\begin{pmatrix}\varpi_w^{m_w}&u\\0&1\end{pmatrix}$ for some $u\in\mathcal O_{E_w}$, where $m_w$ is the number of $j$ with $w_j=w$. For such $a$, $T_af(z,x)=\sum_if(z,xa_i)$.
\end{enumerate}
\end{definition}
\leanref{PositivityTheory/AutomorphicForms/PublishedInputs/JacquetLanglands/TorusKernels/MultivariableReproducingFormula.lean}{UpperTheory.Esk.Dat.IsForm}
\leanref{PositivityTheory/AutomorphicForms/PublishedInputs/JacquetLanglands/TorusKernels/MultivariableReproducingFormula.lean}{UpperTheory.Esk.Dat.M}
\leanref{PositivityTheory/AutomorphicForms/PublishedInputs/JacquetLanglands/TorusKernels/MultivariableReproducingFormula.lean}{UpperTheory.Esk.Dat.Op}
\leanref{PositivityTheory/AutomorphicForms/PublishedInputs/JacquetLanglands/TorusKernels/IsHolVOf.lean}{UpperTheory.Esk.mkDat}
\leanref{PositivityTheory/AutomorphicForms/PublishedInputs/JacquetLanglands/TorusKernels/IsHolVOf.lean}{UpperTheory.Esk.Rel}
\leanref{PositivityTheory/AutomorphicForms/PublishedInputs/JacquetLanglands/TorusKernels/IsHolVOf.lean}{UpperTheory.Esk.IsReps}
\leanref{Development/AutomorphicArithmetic/AutomorphicVarieties/IdeleGenerators/ShortLatticeUnitarize/CmpResplit.lean}{UpperTheory.V.Sd.Wc.Kc}
\leanref{Development/AutomorphicArithmetic/AutomorphicVarieties/IdeleGenerators/ShortLatticeUnitarize/CmpResplit.lean}{UpperTheory.V.Sd.Wc.Hc}

\begin{lemma}\label{sha-siegel-step-seeds-induced-hilbert-form-exists-weight-model-facts}
Let $E$ be a totally real number field with $[E:\mathbf Q]\ge2$, $c$ a function from the primes of $E$ to $\mathbf Z_{\ge0}$ that vanishes at all but finitely many primes, $\omega$ a Hecke character of $E$ of finite order, and $\kappa$ a function from the infinite places of $E$ to $\mathbf Z_{\ge1}$. Use the notation of Definition~\ref{sha-siegel-step-seeds-induced-hilbert-form-exists-weight-model}, and write $R(x)$ for right translation by $x\in GL_2(\mathbf A_{E,f})$.
\begin{enumerate}
\item If $\phi$ is a nonzero holomorphic form of weight $\kappa$, level $c$ and central character $\omega$ (Definition~\ref{sha-siegel-step-seeds-induced-hilbert-form-exists-holomorphic}), then $\omega_v(-1)=(-1)^{\kappa(v)}$ at every infinite place $v$, $\phi$ is fixed by $R(K_c)$, and there is a character $\chi$ of $H_c$ with $\chi(hz)=\omega(z)$ for $h\in K_c$ and scalar $z\in\mathbf A_{E,f}^\times$. Every such $\chi$ has finite order.
\item Every function on $\mathfrak P^\Sigma\times GL_2(\mathbf A_{E,f})$ corresponds in weight $\kappa$ to exactly one function on $GL_2(\mathbf A_E)$, and different functions correspond to different functions. The correspondence is linear, and if $f$ corresponds to $\phi$ and $a_1,\dots,a_n\in GL_2(\mathbf A_{E,f})$, then $\sum_if(\cdot,\cdot\,a_i)$ corresponds to $\sum_iR(a_i)\phi$.
\end{enumerate}
Let $\chi$ be a character of $H_c$ as in (1), and suppose that $\omega_v(-1)=(-1)^{\kappa(v)}$ at every infinite place $v$.
\begin{enumerate}
\item[(3)] $M_\kappa(c,\chi)$ is finite-dimensional and consists of the functions with the four properties of Definition~\ref{sha-siegel-step-seeds-induced-hilbert-form-exists-weight-model}(1). Every finite sequence $W$ of primes has representatives $a$, and $T_a$ maps $M_\kappa(c,\chi)$ into itself.
\item[(4)] A holomorphic form of weight $\kappa$, level $c$ and central character $\omega$ that is a cusp form corresponds in weight $\kappa$ to an element of $M_\kappa(c,\chi)$.
\item[(5)] The function that corresponds in weight $\kappa$ to an element of $M_\kappa(c,\chi)$ is a holomorphic form of weight $\kappa$, level $c$ and central character $\omega$ that is a cusp form.
\item[(6)] For every representation $\rho$ of $GL_2(\mathbf A_{E,f})$ and primes $w,w'$, the operators $U_w$ and $U_{w'}$ map the vectors fixed by $K_c$ to such vectors, and commute on them.
\end{enumerate}
\end{lemma}
\leanref{PositivityTheory/AutomorphicForms/PublishedInputs/JacquetLanglands/TorusKernels/IsHolVOf.lean}{UpperTheory.Esk.H0}
\leanref{PositivityTheory/AutomorphicForms/PublishedInputs/JacquetLanglands/TorusKernels/IsHolVOf.lean}{UpperTheory.Esk.H0b}
\leanref{PositivityTheory/AutomorphicForms/PublishedInputs/JacquetLanglands/TorusKernels/IsHolVOf.lean}{UpperTheory.Esk.HRel}
\leanref{PositivityTheory/AutomorphicForms/PublishedInputs/JacquetLanglands/TorusKernels/IsHolVOf.lean}{UpperTheory.Esk.HFin}
\leanref{PositivityTheory/AutomorphicForms/PublishedInputs/JacquetLanglands/TorusKernels/IsHolVOf.lean}{UpperTheory.Esk.H1}
\leanref{PositivityTheory/AutomorphicForms/PublishedInputs/JacquetLanglands/TorusKernels/IsHolVOf.lean}{UpperTheory.Esk.H2a}
\leanref{PositivityTheory/AutomorphicForms/PublishedInputs/JacquetLanglands/TorusKernels/IsHolVOf.lean}{UpperTheory.Esk.H2d}
\leanref{PositivityTheory/AutomorphicForms/PublishedInputs/JacquetLanglands/TorusKernels/IsHolVOf.lean}{UpperTheory.Esk.H2e}
\leanref{PositivityTheory/AutomorphicForms/PublishedInputs/JacquetLanglands/TorusKernels/IsHolVOf.lean}{UpperTheory.Esk.DAx}
\leanref{PositivityTheory/AutomorphicForms/PublishedInputs/JacquetLanglands/TorusKernels/IsHolVOf.lean}{UpperTheory.Esk.DReps}
\leanref{PositivityTheory/AutomorphicForms/PublishedInputs/JacquetLanglands/TorusKernels/IsHolVOf.lean}{UpperTheory.Esk.DGood}
\leanref{PositivityTheory/AutomorphicForms/PublishedInputs/JacquetLanglands/TorusKernels/MultivariableReproducingFormula.lean}{UpperTheory.Esk.FormsStableUnderOperators}
\leanref{Development/AutomorphicArithmetic/AutomorphicVarieties/IdeleGenerators/ShortLatticeUnitarize/Unitarize.lean}{UpperTheory.V.Sd.Wc.HkComm}
\begin{proof}
\openstep{The comparison between holomorphic forms on $GL_2(\mathbf A_E)$ and the space $M_\kappa(c,\chi)$ is not yet written out: the sign and the character $\chi$ of a holomorphic form, the correspondence and its compatibility with right translations, the finite dimensionality of $M_\kappa(c,\chi)$ and its stability under the operators $T_a$, the existence of representatives, the cuspidality of the forms obtained from $M_\kappa(c,\chi)$, and the commutation of the operators $U_w$.}
% The inputs of pf_HolWC supplied by holWC_of: hyperbolic_terms_setup_holds (H0), hyperbolic_orbit_integral_holds (H0b), s_HRel (HRel), s_HFin (HFin), hyperbolic_terms_integrable_holds (H1), hyperbolic_relation_holds (H2a), hyperbolic_tail_estimate_holds (H2d), hyperbolic_finite_sum_identity_holds (H2e), forms_stable_under_operators_holds, dAx (the axioms Dat.Ax of mkDat), dReps (DRepsP.lean), dGood' and UpperTheory.V.Sd.Wc.hkComm (hkReps uopSum, kcTop). Size: about 6,200 lines with their dependencies, net of the definitions and of Lemma sha-siegel-step-seeds-induced-hilbert-form-exists-weight-linear; the largest parts are s_HFin (about 2,400), dAx (about 1,800), hyperbolic_tail_estimate_holds (about 1,500) and dReps (about 1,100).
\end{proof}

\begin{lemma}\label{sha-siegel-step-seeds-induced-hilbert-form-exists-weight-traces}
Let $E$ be a totally real number field with $[E:\mathbf Q]\ge2$, $\ell$ an odd prime such that every prime $w$ of $E$ above $\ell$ has norm $\ell$ and $\ell\notin w^2$, and $\iota:\bar{\mathbf Q}_\ell\cong\mathbf C$; write $\nu(x)=|\iota^{-1}(x)|$ for $x\in\mathbf C$. For an infinite place $v$ of $E$, let $w_v$ be the prime of $E$ above $\ell$ induced by the embedding $\iota^{-1}\circ v$ of $E$ into $\bar{\mathbf Q}_\ell$. Let $k_0\ge4$, $c$ a function from the primes of $E$ to $\mathbf Z_{\ge0}$ that vanishes at all but finitely many primes, with $c(w_v)=1$ for every infinite place $v$ and $c(u)\ne0$ for some prime $u$ of norm $2$, $\omega$ a Hecke character of $E$ of finite order with $\omega_v(-1)=(-1)^{k_0}$ at every infinite place $v$, and $\chi$ a character of $H_c$ with $\chi(hz)=\omega(z)$ for $h\in K_c$ and scalar $z$ (Definition~\ref{sha-siegel-step-seeds-induced-hilbert-form-exists-weight-model}). Let $N\ge0$, and $t$ a function from the infinite places to $\mathbf Z_{\ge0}$ with $(\ell-1)\ell^N$ dividing every $t(v)$. For a function $s$ from the infinite places to $\mathbf Z_{\ge0}$ and a prime $w$, let $n_s(w)=\operatorname N(w)^{k_0/2-1}\prod_{v:\,w_v=w}\ell^{s(v)}$. Let $W$ be a finite sequence of primes and $r\ge1$ such that every $w_v$ occurs at least $r$ times in $W$, and $a$ representatives for $W$. Let
\[
A=\prod_{w\in W}n_0(w)\cdot\operatorname{tr}\bigl(T_a\mid M_{k_0}(c,\chi)\bigr),\qquad A'=\prod_{w\in W}n_t(w)\cdot\operatorname{tr}\bigl(T_a\mid M_{k_0+2t}(c,\chi)\bigr),
\]
the products running over the terms of $W$ with multiplicity, $M_{k_0}$ for the constant weight $k_0$ and $M_{k_0+2t}$ for $v\mapsto k_0+2t(v)$. Then $\nu(A)\le1$, $\nu(A')\le1$ and $\nu(A-A')\le\max(\nu(\ell)^N,\nu(\ell)^r)$.
\end{lemma}
\leanref{PositivityTheory/AutomorphicForms/PublishedInputs/JacquetLanglands/TorusKernels/IsHolVOf.lean}{UpperTheory.Esk.TrCong}
\leanref{PositivityTheory/AutomorphicForms/PublishedInputs/JacquetLanglands/TorusKernels/PfTrCong.lean}{UpperTheory.Esk.pf_TrCong}
\leanref{PositivityTheory/AutomorphicForms/PublishedInputs/JacquetLanglands/TorusKernels/IsHolVOf.lean}{UpperTheory.Esk.nz}
\begin{proof}
\openstep{The congruence between the traces is not yet written out: the trace formula expressing the trace of $T_a$ on $M_\kappa(c,\chi)$ as a finite sum of orbital terms of totally elliptic elements of $GL_2(E)$, the orders of their stabilizers, which are powers of $2$ by the prime of norm $2$ in the level, and the congruence of each orbital term in the weights $k_0$ and $k_0+2t$.}
% TrCong is proved by pf_TrCong (PfTrCong.lean, 136 lines) from s_TF (the trace formula TF: tr(Op a | M k) = geo D k a S, for a finite family S of elliptic classes, through pf_TF, kernel_sum_bounded_holds, s_Ecut, linear_trace_identity_holds, s_Fell and s_Fhyp), dAx, dTwo' (DTwo: the stabilizers have order a power of 2, through the prime u of norm 2), dGood', the elliptic term bound GTerm (pf_GTerm with s_GCh, s_GUnit and s_GMod: the Chebyshev polynomials ch of the orbital terms orb, units at the primes w_v, and absF times the l-adic part), H0b and HFin, with geo_bound, nu_card and nu_ee. Size: about 8,100 lines with its dependencies, net of the definitions and of Lemmas sha-siegel-step-seeds-induced-hilbert-form-exists-weight-model-facts and -linear, of which the trace formula s_TF is about 7,500.
\end{proof}

\begin{lemma}\label{sha-siegel-step-seeds-induced-hilbert-form-exists-weight-linear-systems}
Let $J$ be a set, $V$ a finite-dimensional complex vector space and $(T_j)_{j\in J}$ a family of pairwise commuting linear operators on $V$. Then there are a finite set $S$ of functions $J\to\mathbf C$ and integers $m(u)\ge1$ for $u\in S$ such that:
\begin{enumerate}
\item a function $u:J\to\mathbf C$ lies in $S$ if and only if there is a nonzero $f\in V$ with $T_jf=u(j)f$ for every $j\in J$;
\item $\sum_{u\in S}m(u)=\dim V$;
\item for every finite sequence $W$ in $J$, $\operatorname{tr}\prod_{j\in W}T_j=\sum_{u\in S}m(u)\prod_{j\in W}u(j)$, the products running over the terms of $W$.
\end{enumerate}
\end{lemma}
\leanref{PositivityTheory/AutomorphicForms/PublishedInputs/JacquetLanglands/TorusKernels/ComplexMinusRealLineMeasureCoordinates.lean}{UpperTheory.Esk.Sys}
\leanref{PositivityTheory/AutomorphicForms/PublishedInputs/JacquetLanglands/TorusKernels/RSubPowApply.lean}{UpperTheory.Esk.pf_Sys}
\begin{proof}
For $u:J\to\mathbf C$ let $E(u)$ be the space of $x\in V$ such that, for every $j$, some power of $T_j-u(j)$ kills $x$. As the $T_i$ commute, each $T_i$ maps each $E(u)$ into itself. For commuting operators on a finite-dimensional complex space the spaces $E(u)$ are independent, and they span $V$, as the generalized eigenspaces of each $T_j$ span $V$. So only finitely many $E(u)$ are nonzero. Let $S$ be the set of $u$ with $E(u)\ne0$, and $m(u)=\dim E(u)$.

For (3), the trace of $\prod_{j\in W}T_j$ is the sum over $u\in S$ of the traces of its restrictions to the $E(u)$. On $E(u)$ each $T_j-u(j)$ is nilpotent. So if $B$ is an operator on $E(u)$ that commutes with $T_j$, then $B(T_j-u(j))$ is nilpotent, and $\operatorname{tr}(BT_j)=u(j)\operatorname{tr}B$. By induction on the length of $W$, the trace of $\prod_{j\in W}T_j$ on $E(u)$ is $m(u)\prod_{j\in W}u(j)$. For $W$ empty, (3) gives (2).

For (1), a nonzero $f$ with $T_jf=u(j)f$ for every $j$ lies in $E(u)$, so $E(u)\ne0$. Conversely, let $E(u)\ne0$, let $U\subseteq E(u)$ be a nonzero subspace mapped into itself by every $T_i$, of least dimension, and let $x\in U$ be nonzero. For $j\in J$, let $k$ be least with $(T_j-u(j))^kx=0$. Then $k\ge1$, and $(T_j-u(j))^{k-1}x$ is a nonzero vector of $U$ killed by $T_j-u(j)$. So $U\cap\ker(T_j-u(j))$ is nonzero, and every $T_i$ maps it into itself, as $T_i$ commutes with $T_j$; by minimality it is $U$. So $T_jx=u(j)x$ for every $j$.
\end{proof}
% pf_Sys (RSubPowApply.lean, 160 lines with the lemmas of CommutingOperatorTraces): E is CommutingOperatorTraces.E, the infimum of the maxGenEigenspace of the T j; the independence and the spanning are Mathlib's Module.End.independent_iInf_maxGenEigenspace_of_forall_mapsTo and iSup_iInf_maxGenEigenspace_eq_top_of_iSup_maxGenEigenspace_eq_top_of_commute, the finiteness WellFoundedGT.finite_ne_bot_of_iSupIndep; (3) is htr with LinearMap.trace_eq_sum_trace_restrict', trace_restrict and trace_R_prod (LinearMap.trace_comp_eq_mul_of_commute_of_isNilpotent with R_nil); (1) is exists_eigvec (a minimal stable subspace, by well-foundedness) and eigenspace_le_maxGenEigenspace.

\begin{lemma}\label{sha-siegel-step-seeds-induced-hilbert-form-exists-weight-linear-powers}
Let $\nu$ be a multiplicative absolute value on $\mathbf C$ satisfying the ultrametric inequality, and $\ell$ a prime.
\begin{enumerate}
\item $\nu(n)\le1$ for every integer $n\ge0$.
\item If $y\in\mathbf C$ and $\nu(y)\le1$, then $\nu\bigl((1+y)^\ell-1\bigr)\le\max(\nu(\ell),\nu(y))\,\nu(y)$.
\item If $\nu(\ell)<1$, and $x\in\mathbf C$ with $\nu(x)<1$, then $\nu\bigl((1+x)^{\ell^a}-1\bigr)\le\max(\nu(\ell),\nu(x))^a\,\nu(x)$ for every integer $a\ge0$.
\item If $y_1,\dots,y_n\in\mathbf C$ and $\nu(y_i-1)<1$ for every $i$, then $\nu(y_1\cdots y_n-1)<1$.
\end{enumerate}
\end{lemma}
\leanref{PositivityTheory/AutomorphicForms/PublishedInputs/JacquetLanglands/TorusKernels/FiniteSystemProof.lean}{UpperTheory.Esk.Fin0Aux.nat_le_one}
\leanref{PositivityTheory/AutomorphicForms/PublishedInputs/JacquetLanglands/TorusKernels/FiniteSystemProof.lean}{UpperTheory.Esk.Fin0Aux.step}
\leanref{PositivityTheory/AutomorphicForms/PublishedInputs/JacquetLanglands/TorusKernels/FiniteSystemProof.lean}{UpperTheory.Esk.Fin0Aux.iter}
\leanref{PositivityTheory/AutomorphicForms/PublishedInputs/JacquetLanglands/TorusKernels/FiniteSystemProof.lean}{UpperTheory.Esk.Fin0Aux.prod_near_one}
\begin{proof}
(1) By induction on $n$: $\nu(n+1)\le\max(\nu(n),\nu(1))\le1$.

(2) By the binomial theorem, $(1+y)^\ell-1=\sum_{0<i<\ell}\binom\ell iy^i+y^\ell$. For $0<i<\ell$, $\ell$ divides $\binom\ell i$, so by (1) and as $\nu(y)\le1$, $\nu\bigl(\binom\ell iy^i\bigr)\le\nu(\ell)\nu(y)$. As $\ell\ge2$, $\nu(y^\ell)\le\nu(y)^2$. The ultrametric inequality gives the bound.

(3) By induction on $a$; for $a=0$ both sides are $\nu(x)$. Let $\theta=\max(\nu(\ell),\nu(x))$, so $\theta\le1$, and $y=(1+x)^{\ell^a}-1$, so that $(1+x)^{\ell^{a+1}}-1=(1+y)^\ell-1$. By induction, $\nu(y)\le\theta^a\nu(x)\le\nu(x)<1$. By (2), $\nu\bigl((1+y)^\ell-1\bigr)\le\max(\nu(\ell),\nu(y))\,\nu(y)\le\theta\cdot\theta^a\nu(x)$.

(4) By induction on $n$; for $n=0$ the product is $1$. Let $p=y_2\cdots y_n$, with $\nu(p-1)<1$ by induction; then $\nu(p)\le\max(\nu(p-1),1)=1$. As $y_1p-1=(y_1-1)p+(p-1)$, $\nu(y_1p-1)\le\max\bigl(\nu(y_1-1)\nu(p),\nu(p-1)\bigr)<1$.
\end{proof}

\begin{lemma}\label{sha-siegel-step-seeds-induced-hilbert-form-exists-weight-linear-factors}
Let $\nu$ be a multiplicative absolute value on $\mathbf C$ satisfying the ultrametric inequality, $J$ a set, and $\Lambda$ a linear form on the space of complex functions on the set of functions $J\to\mathbf C$. For a function $u:J\to\mathbf C$ and a finite sequence $W$ in $J$ write $u^W=\prod_{j\in W}u(j)$, the product running over the terms of $W$. Let $P$ be a property of finite sequences in $J$ such that, if $W$ has $P$, then so does $W$ followed by any $j\in J$, and let $B\ge0$ with $\nu\bigl(\Lambda(u\mapsto u^W)\bigr)\le B$ for every $W$ with $P$. Let $(\alpha_i,j_i,\beta_i)\in\mathbf C\times J\times\mathbf C$ for $1\le i\le n$, with $\nu(\alpha_i)\le1$ and $\nu(\beta_i)\le1$. Then for every $W$ with $P$,
\[
\nu\Bigl(\Lambda\Bigl(u\mapsto u^W\prod_{i=1}^n\bigl(\alpha_iu(j_i)+\beta_i\bigr)\Bigr)\Bigr)\le B.
\]
\end{lemma}
\leanref{PositivityTheory/AutomorphicForms/PublishedInputs/JacquetLanglands/TorusKernels/FiniteSystemProof.lean}{UpperTheory.Esk.Fin0Aux.factors}
\begin{proof}
By induction on $n$, for all $W$ with $P$ at once; for $n=0$ this is the hypothesis. For $n\ge1$, let $R(u)=\prod_{i=2}^n(\alpha_iu(j_i)+\beta_i)$ and $W'$ be $W$ followed by $j_1$. Then $u^W(\alpha_1u(j_1)+\beta_1)R(u)=\alpha_1u^{W'}R(u)+\beta_1u^WR(u)$, so the value of $\Lambda$ is $\alpha_1\Lambda(u\mapsto u^{W'}R(u))+\beta_1\Lambda(u\mapsto u^WR(u))$. Both $W'$ and $W$ have $P$, so by induction both values of $\Lambda$ have absolute value at most $B$. As $\nu(\alpha_1)\le1$ and $\nu(\beta_1)\le1$, the ultrametric inequality gives the bound.
\end{proof}

\begin{lemma}\label{sha-siegel-step-seeds-induced-hilbert-form-exists-weight-linear-integral}
Let $\nu$ be a multiplicative absolute value on $\mathbf C$ satisfying the ultrametric inequality, $J$ a set, $L\subseteq J$ finite, $S$ a finite set of functions $J\to\mathbf C$ and $a$ a complex function on $S$, such that $\nu\bigl(\sum_{u\in S}a(u)\prod_{j\in W}u(j)\bigr)\le1$ for every finite sequence $W$ in $J$ in which every $j\in L$ occurs. Let $t\in S$ with $a(t)\ne0$ and $t(j)\ne0$ for every $j\in L$. Then $\nu(t(j))\le1$ for every $j\in J$.
\end{lemma}
\leanref{PositivityTheory/AutomorphicForms/PublishedInputs/JacquetLanglands/TorusKernels/FiniteSystemProof.lean}{UpperTheory.Esk.Fin0Aux.integral}
\begin{proof}
For a function $g$ on the functions $J\to\mathbf C$ let $\Lambda(g)=\sum_{u\in S}a(u)g(u)$. For each $u\in S$ other than $t$ choose $j(u)\in J$ with $u(j(u))\ne t(j(u))$. For $x\in\mathbf C$ let $c(x)=1$ if $\nu(x)\le1$ and $c(x)=x^{-1}$ otherwise; then $c(x)\ne0$, $\nu(c(x))\le1$ and $\nu(c(x)x)\le1$. Apply Lemma~\ref{sha-siegel-step-seeds-induced-hilbert-form-exists-weight-linear-factors} to $\Lambda$, to the property that every $j\in L$ occurs in $W$, which is kept when a term is added, to $B=1$, and to the triples $\bigl(c(u(j(u))),\,j(u),\,-c(u(j(u)))u(j(u))\bigr)$ for $u\in S\setminus\{t\}$. Let $j_0\in J$, $n\ge0$, and $W$ the sequence of the elements of $L$ followed by $n$ copies of $j_0$. The factor of $u$ is $c(u(j(u)))\bigl(s(j(u))-u(j(u))\bigr)$ at a function $s$, which vanishes at $s=u$. So only the term of $t$ remains, and the lemma gives $\nu(K)\nu(t(j_0))^n\le1$, where
\[
K=a(t)\prod_{j\in L}t(j)\prod_{u\in S\setminus\{t\}}c(u(j(u)))\bigl(t(j(u))-u(j(u))\bigr).
\]
Every factor of $K$ is nonzero, so $\nu(K)>0$. As $\nu(K)\nu(t(j_0))^n\le1$ for every $n$, $\nu(t(j_0))\le1$.
\end{proof}

\begin{lemma}\label{sha-siegel-step-seeds-induced-hilbert-form-exists-weight-linear}
Let $\nu$ be a multiplicative absolute value on $\mathbf C$ satisfying the ultrametric inequality, $\ell$ a prime with $\nu(\ell)<1$, $J$ a set, and $V$ and $V'$ finite-dimensional complex vector spaces with families $(T_j)_{j\in J}$ and $(T'_j)_{j\in J}$ of pairwise commuting linear operators on them. Let $L\subseteq J$ be finite, $\varepsilon\ge0$ and $0\le\rho<1$ with $\varepsilon<\nu(d)$ for every integer $1\le d\le\dim V$. Suppose that for every finite sequence $W$ in $J$ and every $r\ge1$ such that every $j\in L$ occurs at least $r$ times in $W$, the traces $\tau=\operatorname{tr}\prod_{j\in W}T_j$ and $\tau'=\operatorname{tr}\prod_{j\in W}T'_j$ satisfy $\nu(\tau)\le1$, $\nu(\tau')\le1$ and $\nu(\tau-\tau')\le\max(\varepsilon,\rho^r)$. Let $f\in V$ be nonzero and $s_j\in\mathbf C$ with $T_jf=s_jf$ for every $j$ and $\nu(s_j)=1$ for $j\in L$. Then there are a nonzero $f'\in V'$ and $s'_j\in\mathbf C$ with $T'_jf'=s'_jf'$ and $\nu(s'_j-s_j)<1$ for every $j$.
\end{lemma}
\leanref{PositivityTheory/AutomorphicForms/PublishedInputs/JacquetLanglands/TorusKernels/ComplexMinusRealLineMeasureCoordinates.lean}{UpperTheory.Esk.Fin0}
\leanref{PositivityTheory/AutomorphicForms/PublishedInputs/JacquetLanglands/TorusKernels/FinalTF.lean}{UpperTheory.Esk.finite_system_holds}
\begin{proof}
Let $s$ be the function $j\mapsto s_j$. For a function $u:J\to\mathbf C$ and a finite sequence $W$ in $J$ write $u^W=\prod_{j\in W}u(j)$, the product running over the terms of $W$. Call $x\in\mathbf C$ \emph{integral} if $\nu(x)\le1$.
\begin{enumerate}
\item[(1)] \emph{Eigenvalue systems.} Lemma~\ref{sha-siegel-step-seeds-induced-hilbert-form-exists-weight-linear-systems} gives a finite set $S$ and integers $m(u)\ge1$ for the family $(T_j)$ on $V$, and $S'$ and $m'(u)\ge1$ for $(T'_j)$ on $V'$. As $f\ne0$, $s\in S$.
\item[(2)] \emph{Integrality.} Let $W$ be a finite sequence in which every $j\in L$ occurs. By the hypothesis with $r=1$ and (3) of Lemma~\ref{sha-siegel-step-seeds-induced-hilbert-form-exists-weight-linear-systems}, $\nu\bigl(\sum_{u\in S}m(u)u^W\bigr)\le1$ and $\nu\bigl(\sum_{u\in S'}m'(u)u^W\bigr)\le1$. As $m(u)\ne0$ and $m'(u)\ne0$, Lemma~\ref{sha-siegel-step-seeds-induced-hilbert-form-exists-weight-linear-integral} shows: if $u\in S\cup S'$ and $u(j)\ne0$ for every $j\in L$, then $u(j)$ is integral for every $j\in J$. As $\nu(s_j)=1$, $s_j\ne0$ for $j\in L$; so every $s_j$ is integral.
\item[(3)] \emph{The assumption.} Suppose that the conclusion fails. Every $u\in S'$ is the family of eigenvalues of a nonzero vector of $V'$, by (1) of Lemma~\ref{sha-siegel-step-seeds-induced-hilbert-form-exists-weight-linear-systems}; so there is $j$ with $\nu(u(j)-s_j)\ge1$. If $J$ is empty, then $S'$ is empty, and $\dim V\ge m(s)\ge1$ by (2) of that lemma. Choose $r\ge1$ with $\rho^r<\nu(\dim V)$, which is possible as $\nu(\dim V)>\varepsilon\ge0$ and $\rho<1$. For the empty sequence $W$, $\tau=\dim V$ and $\tau'=0$, and the hypothesis gives $\nu(\dim V)\le\max(\varepsilon,\rho^r)<\nu(\dim V)$, a contradiction. So $J$ is nonempty.
\item[(4)] \emph{Separating factors.} Let $N$ be the set of $u\in S\cup S'$ with $u(j)\ne0$ for every $j\in L$ and $\nu(u(j)-s_j)\ge1$ for some $j$, and choose such a $j(u)$ for each $u\in N$. By (2), the values of every $u\in N$ are integral, and $s_{j(u)}\ne u(j(u))$. For $v\in N$ let $\alpha_v=(s_{j(v)}-v(j(v)))^{-1}$ and $\beta_v=-v(j(v))\alpha_v$; they are integral, as $\nu(s_{j(v)}-v(j(v)))\ge1$ and $v(j(v))$ is integral. For a function $u:J\to\mathbf C$ let
\[
G(u)=\prod_{j\in L}u(j)\cdot\prod_{v\in N}\bigl(\alpha_vu(j(v))+\beta_v\bigr)=\prod_{j\in L}u(j)\cdot\prod_{v\in N}\frac{u(j(v))-v(j(v))}{s_{j(v)}-v(j(v))}.
\]
\item[(5)] \emph{Vanishing.} Let $u\in S\cup S'$ with $\nu(u(j)-s_j)\ge1$ for some $j$. Then $G(u)=0$: either $u(j)=0$ for some $j\in L$, or $u\in N$ and the factor of $v=u$ vanishes. By (3), $G(u)=0$ for every $u\in S'$.
\item[(6)] \emph{Near $1$.} Let $C=\prod_{j\in L}s_j^{-1}$; then $\nu(C)=1$. Let $S_0$ be the set of $u\in S$ with $\nu(u(j)-s_j)<1$ for every $j$. For $u\in S_0$, $CG(u)$ is the product of the numbers $u(j)/s_j$ for $j\in L$ and $(u(j(v))-v(j(v)))/(s_{j(v)}-v(j(v)))$ for $v\in N$. These numbers minus $1$ are $(u(j)-s_j)/s_j$ and $(u(j(v))-s_{j(v)})/(s_{j(v)}-v(j(v)))$, of absolute value at most $\nu(u(j)-s_j)$ and $\nu(u(j(v))-s_{j(v)})$, so less than $1$. By Lemma~\ref{sha-siegel-step-seeds-induced-hilbert-form-exists-weight-linear-powers}(4), $\nu(CG(u)-1)<1$.
\item[(7)] \emph{The exponent.} Let $d=\sum_{u\in S_0}m(u)$. As $s\in S_0$, $1\le m(s)\le d\le\dim V$ by (2) of Lemma~\ref{sha-siegel-step-seeds-induced-hilbert-form-exists-weight-linear-systems}, so $\nu(d)>\varepsilon\ge0$. Let $\delta$ be the largest of the numbers $\nu(CG(u)-1)$ with $u\in S_0$, and $\theta=\max(\nu(\ell),\delta,\rho)$; by (6), $\theta<1$. Choose $a\ge0$ with $\theta^a<\nu(d)$, and let $n=\ell^a\ge1$.
\item[(8)] \emph{The bound from the traces.} For a function $g$ on the functions $J\to\mathbf C$ let $\Lambda(g)=\sum_{u\in S}m(u)g(u)-\sum_{u\in S'}m'(u)g(u)$, a linear form. Let $W$ be a finite sequence in which every $j\in L$ occurs at least $n$ times. By (3) of Lemma~\ref{sha-siegel-step-seeds-induced-hilbert-form-exists-weight-linear-systems}, $\Lambda(u\mapsto u^W)=\tau-\tau'$, so $\nu(\Lambda(u\mapsto u^W))\le\max(\varepsilon,\rho^n)$ by the hypothesis with $r=n$. Apply Lemma~\ref{sha-siegel-step-seeds-induced-hilbert-form-exists-weight-linear-factors} to $\Lambda$, this property of $W$, which is kept when a term is added, $B=\max(\varepsilon,\rho^n)$, the triples $(\alpha_v,j(v),\beta_v)$ for $v\in N$, each taken $n$ times, and $W$ the sequence of the elements of $L$ repeated $n$ times. As the function there is $u\mapsto G(u)^n$, $\nu(\Lambda(G^n))\le\max(\varepsilon,\rho^n)$.
\item[(9)] \emph{The bound from the eigenvalues.} By (5), $G(u)=0$ for $u\in S'$ and for $u\in S\setminus S_0$. So
\[
C^n\Lambda(G^n)=\sum_{u\in S_0}m(u)\bigl(CG(u)\bigr)^n=d+\sum_{u\in S_0}m(u)\Bigl(\bigl(1+(CG(u)-1)\bigr)^{\ell^a}-1\Bigr).
\]
For $u\in S_0$, let $x=CG(u)-1$; then $\nu(x)\le\delta<1$, and by (1) and (3) of Lemma~\ref{sha-siegel-step-seeds-induced-hilbert-form-exists-weight-linear-powers}, the term of $u$ has absolute value at most $\max(\nu(\ell),\nu(x))^a\nu(x)\le\theta^a$.
\item[(10)] \emph{The contradiction.} As $\nu(C)=1$, (8), (9) and the ultrametric inequality give $\nu(d)\le\max(\varepsilon,\rho^n,\theta^a)$. But $\varepsilon<\nu(d)$ by (7), $\theta^a<\nu(d)$ by the choice of $a$, and $\rho^n\le\rho^a\le\theta^a$, as $n=\ell^a\ge a$. This contradicts the assumption of (3), so the conclusion holds.
\end{enumerate}
\end{proof}
% Fin0 is finite_system_holds = finite_system_proof s_Sys (FiniteSystemProof.lean, 409 lines with the lemmas of Fin0Aux; pf_Sys in RSubPowApply.lean, 160 lines). The Lean states it for a Valuation nu on C with values in the nonnegative reals. The proof follows finite_system_proof step by step: (1) is hSys for T and T' and hsS; (2) is intV, intV' (Fin0Aux.integral), hs0 and hsint; (3) is the by_contra, hS'nc, and the case isEmpty_or_nonempty J with hS'e, hd1 and r + 1 for the text's r; (4) is NC, jn, hNCint, hne, Fl, hFl and G; (5) is hG0; (6) is Cinv, hCinv, cong and hcong (prod_near_one); (7) is d, hscong, hd1, hd2, hεd, τ (the text's delta), θ, a and n; (8) is hbound (Fin0Aux.factors with lam, count_rep and prod_rep); (9) is hval and hsmall (iter, nat_le_one); (10) is hρn, hdeq and hfin.

\begin{lemma}\label{sha-siegel-step-seeds-induced-hilbert-form-exists-weight}
Let $E$ be a totally real number field with $[E:\mathbf Q]\ge2$, $\ell$ an odd prime such that every prime $w$ of $E$ above $\ell$ has norm $\ell$ and $\ell\notin w^2$, and $\iota:\bar{\mathbf Q}_\ell\cong\mathbf C$. For an infinite place $v$ of $E$, let $w_v$ be the prime of $E$ above $\ell$ induced by the embedding $\iota^{-1}\circ v$ of $E$ into $\bar{\mathbf Q}_\ell$. Let $k_0\ge4$, and let $c$ be a function from the primes of $E$ to $\mathbf Z_{\ge0}$ that vanishes at all but finitely many primes, with $c(w)=1$ for every $w$ above $\ell$ and $c(u)\ne0$ for some prime $u$ of norm $2$. Let $\omega$ be a Hecke character of $E$ of finite order, and $\phi$ a holomorphic eigenform of parallel weight $k_0$, level $c$ and central character $\omega$ that is a cusp form (Definition~\ref{sha-siegel-step-seeds-induced-hilbert-form-exists-holomorphic}), with eigenvalues $\lambda(w)$ such that $|\iota^{-1}(\lambda(w))|=1$ for every $w$ above $\ell$. Let $S$ be a set of infinite places of $E$ and $N_0\ge0$. Then there are an integer $t\ge N_0$ divisible by $\ell-1$, a nonzero holomorphic form $\phi'$ that is a cusp form, of weight $k_0+2t$ at the places in $S$ and $k_0$ at the other infinite places, of level $c$ and central character $\omega$, and complex numbers $\lambda'(w)$, such that $\operatorname N(w)^{k_0/2-1}U_w\phi'=\lambda'(w)\phi'$ for every prime $w$, $|\iota^{-1}(\ell^t\lambda'(w)-\lambda(w))|<1$ for the primes $w=w_v$ with $v\in S$, and $|\iota^{-1}(\lambda'(w)-\lambda(w))|<1$ for the other primes $w$.
\end{lemma}
\leanref{PositivityTheory/AutomorphicForms/PublishedInputs/JacquetLanglands/TorusKernels/IsHolVOf.lean}{UpperTheory.Esk.HolWC}
\leanref{PositivityTheory/AutomorphicForms/PublishedInputs/JacquetLanglands/TorusKernels/MixSeedOf.lean}{UpperTheory.Esk.holWC}
\leanref{Development/AutomorphicArithmetic/AutomorphicVarieties/IdeleGenerators/ShortLatticeUnitarize/CmpResplit.lean}{UpperTheory.V.Sd.Wc.LDat}
\begin{proof}
Write $\nu(x)=|\iota^{-1}(x)|$ for $x\in\mathbf C$, and use the notation of Definition~\ref{sha-siegel-step-seeds-induced-hilbert-form-exists-weight-model}, with the uniformizers $\varpi_w$ of the operators $U_w$; $M_\kappa$ stands for $M_\kappa(c,\chi)$, with $\chi$ chosen in (1).
\begin{enumerate}
\item[(1)] \emph{The character.} By Lemma~\ref{sha-siegel-step-seeds-induced-hilbert-form-exists-weight-model-facts}(1), applied to $\phi$ of constant weight $k_0$, $\omega_v(-1)=(-1)^{k_0}$ at every infinite place $v$, $\phi$ is fixed by right translation by $K_c$, and there is a character $\chi$ of $H_c$ with $\chi(hz)=\omega(z)$ for $h\in K_c$ and scalar $z$.
\item[(2)] \emph{The operators.} Let $\kappa$ be a function from the infinite places to $\mathbf Z_{\ge1}$ with $\omega_v(-1)=(-1)^{\kappa(v)}$ for every $v$. As $\chi(h)=\omega(1)=1$ for $h\in K_c$, the elements of $M_\kappa$ are fixed by right translation by $K_c$ in the second variable. For every prime $w$ fix representatives $a^w$ for the sequence $(w)$, by Lemma~\ref{sha-siegel-step-seeds-induced-hilbert-form-exists-weight-model-facts}(3), and let $T_w$ be the restriction of $T_{a^w}$ to $M_\kappa$. By the defining property of representatives, applied to right translation in the second variable, $T_af=U_{w_1}\cdots U_{w_m}f=T_{w_1}\cdots T_{w_m}f$ for $f\in M_\kappa$ and representatives $a$ for $W=(w_1,\dots,w_m)$. By Lemma~\ref{sha-siegel-step-seeds-induced-hilbert-form-exists-weight-model-facts}(6), the $T_w$ commute.
\item[(3)] \emph{Eigenvectors.} Let $\kappa$ be as in (2), $f\in M_\kappa$ correspond in weight $\kappa$ to a function $\psi$ fixed by right translation by $K_c$, $w$ a prime and $\alpha,\beta\in\mathbf C$. By Lemma~\ref{sha-siegel-step-seeds-induced-hilbert-form-exists-weight-model-facts}(2), $T_wf$ corresponds to $\sum_iR(a^w_i)\psi=U_w\psi$, and $\alpha T_wf-\beta f$ corresponds to $\alpha U_w\psi-\beta\psi$. As $0$ corresponds to $0$ and different functions correspond to different functions, $\alpha T_wf=\beta f$ if and only if $\alpha U_w\psi=\beta\psi$.
\item[(4)] \emph{The weight.} By Lemma~\ref{sha-siegel-step-seeds-induced-hilbert-form-exists-weight-model-facts}(3), $M_{k_0}$ is finite-dimensional; let $N=\dim M_{k_0}$ and $t=(\ell-1)\ell^N\max(N_0,1)$. Then $t\ge N_0$ and $\ell-1$ divides $t$. Let $s(v)=t$ for $v\in S$ and $s(v)=0$ otherwise, and $\kappa'=k_0+2s$. Then $(\ell-1)\ell^N$ divides every $s(v)$, $\omega_v(-1)=(-1)^{\kappa'(v)}$ for every $v$ by (1), and $M_{\kappa'}$ is finite-dimensional.
\item[(5)] \emph{Small integers.} As $\iota^{-1}(\ell)=\ell$, $\nu(\ell)<1$. Let $1\le d\le N$, and write $d=\ell^am$ with $\ell\nmid m$. As $\ell^a\le d\le N<\ell^N$, $a<N$. There are integers $x,y$ with $mx+\ell y=1$; as $\nu(\ell y)<1$, $\nu(mx)=1$, and as $\nu(m),\nu(x)\le1$, $\nu(m)=1$. So $\nu(d)=\nu(\ell)^a>\nu(\ell)^N$.
\item[(6)] \emph{The form in the model.} By Lemma~\ref{sha-siegel-step-seeds-induced-hilbert-form-exists-weight-model-facts}(4), $\phi$ corresponds in weight $k_0$ to some $f_0\in M_{k_0}$, and $f_0\ne0$, as $\phi\ne0$ and $0$ corresponds to $0$. For a function $s_1$ from the infinite places to $\mathbf Z_{\ge0}$, let $n_{s_1}(w)=\operatorname N(w)^{k_0/2-1}\prod_{v:\,w_v=w}\ell^{s_1(v)}$. As $\operatorname N(w)^{k_0/2-1}U_w\phi=\lambda(w)\phi$, (1) and (3) give $n_0(w)T_wf_0=\lambda(w)f_0$ for every prime $w$.
\item[(7)] \emph{Traces.} Let $W$ be a finite sequence of primes and $r\ge1$ such that every $w_v$ occurs at least $r$ times in $W$, and let $a$ be representatives for $W$. By (2), $\prod_{w\in W}n_0(w)T_w=\bigl(\prod_{w\in W}n_0(w)\bigr)T_a$ on $M_{k_0}$ and $\prod_{w\in W}n_s(w)T_w=\bigl(\prod_{w\in W}n_s(w)\bigr)T_a$ on $M_{\kappa'}$. By Lemma~\ref{sha-siegel-step-seeds-induced-hilbert-form-exists-weight-traces}, applied to $N$ and $s$, their traces $\tau$ and $\tau'$ satisfy $\nu(\tau)\le1$, $\nu(\tau')\le1$ and $\nu(\tau-\tau')\le\max(\nu(\ell)^N,\nu(\ell)^r)$.
\item[(8)] \emph{A congruent eigenvector.} Apply Lemma~\ref{sha-siegel-step-seeds-induced-hilbert-form-exists-weight-linear} to the families $(n_0(w)T_w)$ on $M_{k_0}$ and $(n_s(w)T_w)$ on $M_{\kappa'}$, indexed by the primes, which commute by (2); to $L=\{w_v\}$, $\varepsilon=\nu(\ell)^N$ and $\rho=\nu(\ell)$, which satisfy its hypotheses by (5) and (7); and to $f_0$ and $\lambda$, by (6) and as $\nu(\lambda(w_v))=1$ for every $v$. It gives a nonzero $f'\in M_{\kappa'}$ and numbers $s'(w)$ with $n_s(w)T_wf'=s'(w)f'$ and $\nu(s'(w)-\lambda(w))<1$ for every prime $w$.
\item[(9)] \emph{The new form.} Let $\phi'$ be the function to which $f'$ corresponds in weight $\kappa'$. By Lemma~\ref{sha-siegel-step-seeds-induced-hilbert-form-exists-weight-model-facts}(5), $\phi'$ is a holomorphic form of weight $\kappa'$, level $c$ and central character $\omega$ that is a cusp form, and $\phi'\ne0$, as $f'\ne0$. By Lemma~\ref{sha-siegel-step-seeds-induced-hilbert-form-exists-weight-model-facts}(1), $\phi'$ is fixed by right translation by $K_c$; so, by (3), $n_s(w)U_w\phi'=s'(w)\phi'$ for every prime $w$.
\item[(10)] \emph{The eigenvalues.} Two infinite places that induce the same prime $w$ above $\ell$ are equal, as $E_w=\mathbf Q_\ell$ has one embedding into $\bar{\mathbf Q}_\ell$. So $P(w)=\prod_{v:\,w_v=w}\ell^{s(v)}$ is $\ell^t$ if $w=w_v$ with $v\in S$, and $1$ otherwise, and $n_s(w)=\operatorname N(w)^{k_0/2-1}P(w)$. Let $\lambda'(w)=P(w)^{-1}s'(w)$. Then $\operatorname N(w)^{k_0/2-1}U_w\phi'=\lambda'(w)\phi'$ by (9), and $P(w)\lambda'(w)=s'(w)$, so $\nu(\ell^t\lambda'(w)-\lambda(w))<1$ for $w=w_v$ with $v\in S$ and $\nu(\lambda'(w)-\lambda(w))<1$ for the other primes $w$, by (8). With (4) and (9), $t$, $\phi'$ and $\lambda'$ have the properties of the statement.
\end{enumerate}
\end{proof}
% HolWC is proved by Esk.holWC (MixSeedOf.lean) = holWC_of dAx dTwo' dGood' dReps (pf_GTerm s_GCh s_GUnit s_GMod) (FinalH.lean), that is pf_HolWC (PfHolWC.lean, 304 lines) applied to finite_system_holds (Fin0, Lemma sha-siegel-step-seeds-induced-hilbert-form-exists-weight-linear), pf_TrCong (TrCong, Lemma sha-siegel-step-seeds-induced-hilbert-form-exists-weight-traces) and the inputs of Lemma sha-siegel-step-seeds-induced-hilbert-form-exists-weight-model-facts. The proof follows pf_HolWC step by step: (1) is hH0 and hH0b; (2) is hReps, hgood, hE, hfixM (fix_of_form), word, wordU and comm (hcomm = Wc.hkComm); (3) is eig_iff; (4) is hHFin, N, t, htB, htd, hdiv and hSgn'; (5) is hnul (from D.lt; the text uses iota^-1(l) = l) and hepsilon (nu_coprime, nu_nat_le); (6) is hH1, hf0ne, hnz0 and heig; (7) is htr; (8) is hFin0; (9) is lift (H2d, H2a, H2e, H0 and eig_iff); (10) is hP, smul_fix and hprod. The Lean states the lemma for any D : Wc.LDat F iota, whose map D.pl is v -> w_v (by the field embAt), injective by D.inj. Size: about 14,800 lines with its dependencies, net of the definitions: about 6,200 for Lemma sha-siegel-step-seeds-induced-hilbert-form-exists-weight-model-facts, 8,100 for Lemma sha-siegel-step-seeds-induced-hilbert-form-exists-weight-traces and 500 for Lemma sha-siegel-step-seeds-induced-hilbert-form-exists-weight-linear.

\begin{lemma}\label{sha-siegel-step-seeds-induced-hilbert-form-exists-representation-casimir}
Let $E$ be a totally real number field, $\kappa$ a function from the infinite places of $E$ to $\mathbf Z_{\ge0}$, $c$ a function from the primes of $E$ to $\mathbf Z_{\ge0}$, $\omega$ a Hecke character of $E$ of finite order, and $\phi$ a holomorphic form of weight $\kappa$, level $c$ and central character $\omega$, with the derivatives $X^v_{ab}$ and the Capelli operators $C_v(s)=(s+X^v_{11})(s-1+X^v_{22})-X^v_{21}X^v_{12}$ of Definition~\ref{sha-siegel-step-seeds-induced-hilbert-form-exists-holomorphic}. Let $X$ be the derivative of right translation by the scalar matrices at infinity: $X\phi(g)$ is the derivative at $x=0$ of $\phi(g\,e^x)$, where $e^x\in GL_2(E\otimes\mathbf R)$ is the scalar matrix $e^x$ at every infinite place. Then:
\begin{enumerate}
\item $X\phi=0$, and $(X^v_{11}+X^v_{22})\phi=0$ for every infinite place $v$;
\item $C_v(u+\frac12)\phi=\bigl(u+\frac{\kappa(v)-1}2\bigr)\bigl(u-\frac{\kappa(v)-1}2\bigr)\phi$ for every infinite place $v$ and every $u\in\mathbf C$;
\item if $\phi\ne0$, then $\omega_v(-1)=(-1)^{\kappa(v)}$ at every infinite place $v$.
\end{enumerate}
\end{lemma}
\leanref{PositivityTheory/AutomorphicForms/PublishedInputs/JacquetLanglands/TorusKernels/IsHolVOf.lean}{UpperTheory.Esk.HCap}
\leanref{PositivityTheory/AutomorphicForms/PublishedInputs/JacquetLanglands/TorusKernels/PfHCap.lean}{UpperTheory.Esk.pf_HCap}
\leanref{PositivityTheory/AutomorphicForms/PublishedInputs/JacquetLanglands/TorusKernels/PfHRep.lean}{UpperTheory.Esk.HeckeRepresentationSatake.sign_hol}
\begin{proof}
(1) Let $z=(z_v)$ be real numbers indexed by the infinite places, and for $t\in\mathbf R$ let $e^{tz}$ be the idele with component $e^{tz_v}$ at $v$ and $1$ at the primes, viewed as a scalar matrix. Let $N\ge1$ with $\omega^N=1$. Then $\omega(e^{tz})=\omega(e^{tz/N})^N=1$, so by condition (c), $\phi(g\,e^{tz})=\phi(e^{tz}g)=\phi(g)$ for all $g$ and $t$. The derivative at $t=0$ is the derivative of $\phi$ at $g$ in the direction of the scalar matrix $z$, which therefore vanishes. With $z_v=1$ for every $v$ this is $X\phi=0$; with $z_v=1$ at one place $v$ and $z_{v'}=0$ elsewhere it is $(X^v_{11}+X^v_{22})\phi=0$.

(2) Fix $v$ and write $X_{ab}=X^v_{ab}$, $H=X_{11}-X_{22}$, $Y=X_{12}+X_{21}$ and $W=X_{12}-X_{21}$. The derivatives satisfy $[X_{ab},X_{cd}]=\delta_{bc}X_{ad}-\delta_{da}X_{cb}$; so $X_{11}+X_{22}$ commutes with every $X_{ab}$, $[X_{12},X_{21}]=H$ and $[H,Y]=2W$. By (1), $X_{22}\phi=-X_{11}\phi$, $X_{11}\phi=\frac12H\phi$ and $X_{11}X_{11}\phi=\frac12X_{11}H\phi=\frac14H^2\phi$. Expanding $C_v(s)\phi$ and using $X_{21}X_{12}=\frac12(X_{12}X_{21}+X_{21}X_{12})-\frac12H$ gives
\[
C_v(s)\phi=s(s-1)\phi-\tfrac12H\phi-\tfrac14H^2\phi-X_{21}X_{12}\phi=s(s-1)\phi-\tfrac14(H^2+Y^2-W^2)\phi,
\]
as $Y^2-W^2=2(X_{12}X_{21}+X_{21}X_{12})$. Condition (d) says $W\phi=i\kappa(v)\phi$ and $(H-iY)\phi=0$. As $(H+iY)(H-iY)=H^2+Y^2-i[H,Y]=H^2+Y^2-2iW$,
\[
(H^2+Y^2-W^2)\phi=(H+iY)(H-iY)\phi+2iW\phi-W^2\phi=(\kappa(v)^2-2\kappa(v))\phi.
\]
At $s=u+\frac12$, $s(s-1)-\frac14(\kappa(v)^2-2\kappa(v))=u^2-\frac14(\kappa(v)-1)^2$, which is (2).

(3) As $\phi\ne0$, there is $g$ with $\phi(g)\ne0$. For $\theta\in\mathbf R$ let $r(\theta)\in GL_2(E\otimes\mathbf R)$ be $\begin{pmatrix}\cos\theta&\sin\theta\\-\sin\theta&\cos\theta\end{pmatrix}$ at $v$ and $1$ at the other infinite places, and $f(\theta)=\phi(g\,r(\theta))$. As $r(\theta+\theta')=r(\theta)r(\theta')$ and the derivative of $r$ at $0$ is $e_{12}-e_{21}$ at $v$, condition (d) gives $f'(\theta)=\bigl((X^v_{12}-X^v_{21})\phi\bigr)(g\,r(\theta))=i\kappa(v)f(\theta)$. So $f(\theta)e^{-i\kappa(v)\theta}$ is constant, and $f(\pi)=e^{i\pi\kappa(v)}\phi(g)$. The matrix $r(\pi)$ is the scalar idele with component $-1$ at $v$ and $1$ elsewhere, so $f(\pi)=\omega_v(-1)\phi(g)$ by condition (c). Hence $\omega_v(-1)=e^{i\pi\kappa(v)}=(-1)^{\kappa(v)}$.
\end{proof}
% HCap is pf_HCap (PfHCap.lean, 77 lines). (1) is lie_scalar with infGL_scalar and fin_expU (the field cent and the finite order of omega); (2) is the private casimir, that is UpperTheory.Reformulated.B.Qf10.casimir (CapelliIdentities/Qminor.lean): the expansion hq of the Capelli determinant qdet (Formalism.REAL.qdet, here TR.cap v u = qdet (Eop v) (u + 1/2)), the relations GlRel (TR.glRel_Eop) and a linear combination of h1 (part (1)), wt and low, which the text writes through H, Y and W. (3) is sign_hol (PfHRep.lean): rot, rot_add, TR.hasDerivAt_curve with the field wt, is_const_of_deriv_eq_zero, and cornerGL_scalar with the field cent at rot (pi e_v) = -1 at v. About 300 lines with these files.

\begin{lemma}\label{sha-siegel-step-seeds-induced-hilbert-form-exists-representation-constituent}
Let $E$ be a totally real number field. Call a space of functions on $GL_2(\mathbf A_E)$ that are smooth in the variable at infinity \emph{stable} if it is stable under right translation by $GL_2(\mathbf A_{E,f})$ and by the maximal compact subgroup of $GL_2(E\otimes\mathbf R)$, and under the derivatives of right translation in the directions of the Lie algebra of $GL_2(E\otimes\mathbf R)$. Let $\phi\ne0$ lie in the sum of the spaces of the cuspidal automorphic representations of $GL_2(\mathbf A_E)$. Then there are a cuspidal automorphic representation $\pi_0$ of $GL_2(\mathbf A_E)$ and a nonzero $\psi$ in the space of $\pi_0$ such that, for every linear operator $z$ on the functions smooth in the variable at infinity that maps every stable space into itself and every $a\in\mathbf C$, $z\phi=a\phi$ implies $z\psi=a\psi$.
\end{lemma}
\leanref{PositivityTheory/AutomorphicForms/PublishedInputs/JacquetLanglands/TorusKernels/PfHRep.lean}{UpperTheory.Esk.HeckeRepresentationSatake.proj}
\begin{proof}
Sums and intersections of stable spaces are stable. As $\phi$ lies in the sum, it lies in $\pi_1+\dots+\pi_n$ for finitely many cuspidal automorphic representations $\pi_i$, identified with their spaces. We argue by induction on $n$. If $n=0$, then $\phi=0$, which is false. Let $n\ge1$ and $S=\pi_2+\dots+\pi_n$, a stable space. The space $\pi_1\cap S$ is stable and contained in $\pi_1$, so it is $0$ or $\pi_1$, as $\pi_1$ is irreducible. If it is $\pi_1$, then $\phi\in S$, and the induction hypothesis applies. Otherwise write $\phi=\psi+\psi'$ with $\psi\in\pi_1$ and $\psi'\in S$. If $\psi=0$, then $\phi\in S$, and the induction hypothesis applies. If $\psi\ne0$, let $\pi_0=\pi_1$. For $z$ and $a$ as in the statement, $z\psi-a\psi=a\psi'-z\psi'$ lies in $\pi_1\cap S=0$, as $z$ maps $\pi_1$ and $S$ into themselves; so $z\psi=a\psi$.
\end{proof}
% proj (PfHRep.lean, 25 lines, with stable_iSup): Submodule.mem_iSup_iff_exists_finset and Finset.induction_on, with pi.irred applied to pi.V inf the supremum over S.

\begin{lemma}\label{sha-siegel-step-seeds-induced-hilbert-form-exists-representation}
Assume the Galois representations as in Theorem~\ref{sha-residually-reducible-automorphy}. Let $E$ be a totally real number field, $k\ge2$, $n$ a function from the infinite places of $E$ to $\mathbf Z_{\ge0}$, $c$ a function from the primes of $E$ to $\mathbf Z_{\ge0}$, and $\omega$ a Hecke character of $E$ of finite order. Let $\phi$ be a nonzero holomorphic form of weight $v\mapsto k+2n(v)$, level $c$ and central character $\omega$ that is a cusp form (Definition~\ref{sha-siegel-step-seeds-induced-hilbert-form-exists-holomorphic}), and let $\lambda(w)\in\mathbf C$ with $\operatorname N(w)^{k/2-1}U_w\phi=\lambda(w)\phi$ for every prime $w$. Then there is a cuspidal automorphic representation $\pi$ of $GL_2(\mathbf A_E)$, regular algebraic and polarized as in condition (a) of Definition~\ref{sha-siegel-step-seeds-induced-hilbert-form-ordinary}, such that:
\begin{enumerate}
\item $\pi$ has the infinitesimal character of weight $(n(v),1-k-n(v))$ at every infinite place $v$ of $E$;
\item for every prime $w$, the action $\pi_w$ of $GL_2(E_w)$ on the space of $\pi$ has a nonzero vector fixed by $K_1(w^{c(w)})$ on which $U_w$ acts by $\lambda(w)$;
\item for every prime $w$ with $c(w)=0$ and every uniformizer $\varpi_w$ of $E_w$, $\pi$ has Satake parameters $B$ at $w$ with $\prod_{b\in B}(X-\operatorname N(w)^{1/2}b)=X^2-\lambda(w)X+\omega_w(\varpi_w)\operatorname N(w)^{k-1}$.
\end{enumerate}
\end{lemma}
\leanref{PositivityTheory/AutomorphicForms/PublishedInputs/JacquetLanglands/TorusKernels/IsHolVOf.lean}{UpperTheory.Esk.HRep}
\leanref{Development/AutomorphicArithmetic/AutomorphicVarieties/ModularityArgument/ThresholdArguments/DenseThresholdCases/ExistsSat.lean}{UpperTheory.Esk.pf_HRep_alt}
\begin{proof}
Let $\kappa(v)=k+2n(v)$, and call a space of functions \emph{stable} as in Lemma~\ref{sha-siegel-step-seeds-induced-hilbert-form-exists-representation-constituent}. For a prime $w$ let $\varpi_w$ be a uniformizer of $E_w$ and $\mathcal O_w$ the ring of integers of $E_w$. Right translation by an element of $GL_2(E_w)$, viewed in $GL_2(\mathbf A_E)$, maps stable spaces into themselves; so do the operators $U_w$, which are finite sums of such translations, and the operators $C_v(s)$ of Lemma~\ref{sha-siegel-step-seeds-induced-hilbert-form-exists-representation-casimir}, which are polynomials in the derivatives $X^v_{ab}$.
\begin{enumerate}
\item[(1)] \emph{A cuspidal constituent.} By Lemma~\ref{sha-siegel-step-seeds-induced-hilbert-form-exists-representation-casimir}(1), $X\phi=0$. A cusp form on which $X$ acts by a scalar lies in the sum of the spaces of the cuspidal automorphic representations of $GL_2(\mathbf A_E)$, by the discreteness of the cuspidal spectrum \cite{GGPS69}. So Lemma~\ref{sha-siegel-step-seeds-induced-hilbert-form-exists-representation-constituent} gives a cuspidal automorphic representation $\pi_0$ and a nonzero $\psi$ in its space such that every linear operator $z$ that maps stable spaces into themselves and satisfies $z\phi=a\phi$ satisfies $z\psi=a\psi$.
\item[(2)] \emph{The twist.} Let $\nu=|\cdot|^{1-k/2}$, a Hecke character of $E$. It is trivial on the units of every $E_w$, $\nu_w(\varpi_w)=\operatorname N(w)^{k/2-1}$, $\nu_v(-1)=1$ at every infinite place $v$, and the derivative at $x=0$ of $\nu_v(e^x)$ is $1-\frac k2$. Let $\pi$ be the twist of $\pi_0$ by $\nu\circ\det$, the cuspidal automorphic representation on the functions $T\phi_0:g\mapsto\nu(\det g)\phi_0(g)$ with $\phi_0$ in the space of $\pi_0$; then $\pi(g)T\phi_0=\nu(\det g)\,T(\pi_0(g)\phi_0)$ for $g\in GL_2(\mathbf A_{E,f})$. Let $x=T\psi$; it is nonzero.
\item[(3)] \emph{Fixed vectors.} Let $w$ be a prime, $g\in GL_2(E_w)$ and $a\in\mathbf C$, and suppose that right translation by $g$ multiplies $\phi$ by $a$. By (1), $\pi_0(g)\psi=a\psi$. By condition (b) of Definition~\ref{sha-siegel-step-seeds-induced-hilbert-form-exists-holomorphic}, $\psi$ is therefore fixed by $K_1(w^{c(w)})$. The determinant of an element of $K_1(w^{c(w)})$ is a unit of $\mathcal O_w$, on which $\nu_w$ is trivial, so $x$ is fixed by $K_1(w^{c(w)})$ too.
\item[(4)] \emph{The operators $U_w$.} Let $w$ be a prime. As $\operatorname N(w)^{k/2-1}U_w\phi=\lambda(w)\phi$, (1) gives $\operatorname N(w)^{k/2-1}U_w\psi=\lambda(w)\psi$. On the vectors fixed by $K_1(w^c)$ of a representation $\rho$, $U_w=\sum_g\rho(g)$, where $g$ runs over the matrices $\begin{pmatrix}\varpi_w&u\\0&1\end{pmatrix}$ with $u$ in a set of representatives of $\mathcal O_w/\varpi_w\mathcal O_w$, together with $\varpi_w\operatorname{diag}(\varpi_w,1)^{-1}$ if $c=0$, by the computation of the cosets in the double coset, as in the proof of Lemma~\ref{sha-siegel-step-seeds-induced-hilbert-form-galois-jacquet}. Each of these matrices has determinant $\varpi_w$. So, by (2) and (3), $U_wx=\nu_w(\varpi_w)\,T(U_w\psi)=\operatorname N(w)^{k/2-1}T(U_w\psi)=\lambda(w)x$.
\item[(5)] \emph{The centre.} Let $\chi=\omega\nu^2$. By condition (c), right translation by the scalar matrix $\varpi_w$ at $w$ multiplies $\phi$ by $\omega_w(\varpi_w)$. By (3), $\pi_0(\varpi_w)\psi=\omega_w(\varpi_w)\psi$, so $\pi(\varpi_w)x=\nu_w(\varpi_w^2)\,\omega_w(\varpi_w)x=\chi_w(\varpi_w)x$.
\item[(6)] \emph{Satake parameters.} Let $w$ be a prime such that $x$ is fixed by $GL_2(\mathcal O_w)$. By Lemma~\ref{sha-siegel-step-seeds-induced-hilbert-form-galois-component}(1) and its proof, $\pi$ has Satake parameters $\{\alpha,\beta\}$ at $w$, and $\pi_w$ is a sum of images of a spherical local component with these Satake parameters. So, with the normalization of Definition~\ref{gal-definite-unitary}, the operators of the double cosets of $\operatorname{diag}(\varpi_w,1)$ and of $\operatorname{diag}(\varpi_w,\varpi_w)$ modulo $GL_2(\mathcal O_w)$ act on every vector of $\pi_w$ fixed by $GL_2(\mathcal O_w)$, in particular on $x$, by $\operatorname N(w)^{1/2}(\alpha+\beta)$ and by $\alpha\beta$. The first is $U_w$ for $c=0$, as $K_1(w^0)=GL_2(\mathcal O_w)$; the second is the scalar matrix $\varpi_w$, its own double coset. By (5), $\alpha\beta=\chi_w(\varpi_w)$.
\item[(7)] \emph{Infinitesimal characters.} Let $v$ be an infinite place and $u\in\mathbf C$. By Lemma~\ref{sha-siegel-step-seeds-induced-hilbert-form-exists-representation-casimir}(2) and (1), $C_v(u+\frac12)\psi=(u+\frac{\kappa(v)-1}2)(u-\frac{\kappa(v)-1}2)\psi$. The operator $C_v(s)$ commutes with right translation by $GL_2(\mathbf A_{E,f})$, which acts on another variable, and with the maximal compact subgroup and the Lie algebra at infinity, as $C_v(s)$ is central in the enveloping algebra of the Lie algebra of $GL_2(\mathbf R)$ at $v$ and invariant under conjugation by $GL_2(\mathbf R)$ \cite{HU91}. So the eigenspace of $C_v(u+\frac12)$ in the space of $\pi_0$ for this eigenvalue is stable. It contains $\psi\ne0$, and $\pi_0$ is irreducible, so $C_v(u+\frac12)$ acts on $\pi_0$ by this scalar: $\pi_0$ has the infinitesimal character $\{\frac{\kappa(v)-1}2,-\frac{\kappa(v)-1}2\}$ at $v$. The derivatives satisfy $X^v_{ab}T\phi_0=T(X^v_{ab}\phi_0)+\delta_{ab}(1-\frac k2)T\phi_0$ by (2), so $C_v(s)T\phi_0=T(C_v(s+1-\frac k2)\phi_0)$, and twisting adds $1-\frac k2$ to the infinitesimal character. So $\pi$ has the infinitesimal character $\{n(v)+\frac12,\frac32-k-n(v)\}$, that of weight $(n(v),1-k-n(v))$, at $v$. Its two elements lie in $\frac12+\mathbf Z$, and they are distinct, as $2n(v)=1-k$ is impossible for $k\ge2$. So $\pi$ is regular algebraic.
\item[(8)] \emph{The sign.} By Lemma~\ref{sha-siegel-step-seeds-induced-hilbert-form-exists-representation-casimir}(3), $\omega_v(-1)=(-1)^{k+2n(v)}=(-1)^k$, and $\nu_v(-1)=1$; so $\chi_v(-1)=(-1)^k$ at every infinite place $v$.
\item[(9)] \emph{The polarization.} Let $\pi^\vee$ be the cuspidal automorphic representation on the functions $g\mapsto\phi_1({}^tg^{-1})$ with $\phi_1$ in the space of $\pi$, and $\pi'$ its twist by $\chi\circ\det$. The stabilizer of $x$ in $GL_2(\mathbf A_{E,f})$ is open, so it contains $GL_2(\mathcal O_w)$ for all but finitely many primes $w$; and $\chi$ is unramified at all but finitely many primes. At a prime $w$ of both kinds, $\pi$ has Satake parameters $\{\alpha,\beta\}$ with $\alpha\beta=\chi_w(\varpi_w)$ by (6); so $\pi^\vee$ has the Satake parameters $\{\alpha^{-1},\beta^{-1}\}$ at $w$, and $\pi'$ has $\{\chi_w(\varpi_w)\alpha^{-1},\chi_w(\varpi_w)\beta^{-1}\}=\{\beta,\alpha\}$. By strong multiplicity one \cite{JS81b}, $\pi\cong\pi'$: there is a linear automorphism $S$ of the space of $\pi$ with $S(\pi(g)\phi_1)=\chi(\det g)\,\pi({}^tg^{-1})S\phi_1$ for $g\in GL_2(\mathbf A_{E,f})$, and with the corresponding relations for the maximal compact subgroup and the Lie algebra at infinity. By the theorem of Gelfand and Kazhdan \cite{GK75}, there is a nondegenerate bilinear form $B$ on the space of $\pi$ with $B(\pi(g)\phi_1,\pi({}^tg^{-1})\phi_2)=B(\phi_1,\phi_2)$ for $g\in GL_2(\mathbf A_{E,f})$, $B(\pi(h)\phi_1,\pi(h)\phi_2)=B(\phi_1,\phi_2)$ for $h$ in the maximal compact subgroup at infinity, and $B(Z\phi_1,\phi_2)+B(\phi_1,-{}^tZ\phi_2)=0$ for $Z$ in the Lie algebra. Then $B'(\phi_1,\phi_2)=B(\phi_1,S\phi_2)$ is nondegenerate, and $B'(\pi(g)\phi_1,\pi(g)\phi_2)=\chi(\det g)B(\pi(g)\phi_1,\pi({}^tg^{-1})S\phi_2)=\chi(\det g)B'(\phi_1,\phi_2)$; likewise for the maximal compact subgroup, which consists of orthogonal matrices, and for the Lie algebra. With (8), $\pi$ is polarized by $\chi$, as in condition (a) of Definition~\ref{sha-siegel-step-seeds-induced-hilbert-form-ordinary}.
\item[(10)] \emph{Conclusion.} Statement (1) is (7). For (2), $x$ is a nonzero vector fixed by $K_1(w^{c(w)})$ with $U_wx=\lambda(w)x$, by (3) and (4). For (3), let $c(w)=0$. Then $x$ is fixed by $K_1(w^0)=GL_2(\mathcal O_w)$, and by (6) $\pi$ has Satake parameters $B=\{\alpha,\beta\}$ at $w$ with $U_wx=\operatorname N(w)^{1/2}(\alpha+\beta)x$ and $\alpha\beta=\omega_w(\varpi_w)\operatorname N(w)^{k-2}$. By (4), $\lambda(w)=\operatorname N(w)^{1/2}(\alpha+\beta)$. So $\prod_{b\in B}(X-\operatorname N(w)^{1/2}b)=X^2-\operatorname N(w)^{1/2}(\alpha+\beta)X+\operatorname N(w)\alpha\beta=X^2-\lambda(w)X+\omega_w(\varpi_w)\operatorname N(w)^{k-1}$.
\end{enumerate}
\end{proof}
% pf_HRep_alt (ExistsSat.lean, 112 lines), followed step by step, with the copies in PfHRep.lean (359 lines) of the private lemmas of HolRep.lean. (1) is hC (Esk.HCap, supplied as s_HCap = pf_HCap, Lemma sha-siegel-step-seeds-induced-hilbert-form-exists-representation-casimir), P.ssR (Published.CuspSemisimpleReal, a field of HypothesesH') and proj (Lemma sha-siegel-step-seeds-induced-hilbert-form-exists-representation-constituent). (2) is nu = Model.absPow F (1 - k/2) with absPow_isAbsPow (hnu1, hnu2, hnu3), pi = pi_0.twist nu, T = pi_0.toTwist nu and hact (finPart_toTwist, chiDet_locGL). (3) is hloc, hfix_0 and hfix with unit_det. (4) is hU: hz for Sd.hk, uop_map (through Sd.uopSum) for the inclusion of pi_0 and for T. (5) is hZ with fin_scalar, CuspRep.scalar_comm and the field cent. (6) is hsat: exists_sat_alt (local_comp_real_of_printed_results, satake_local_of_printed_results and LocSph.exists_fixed_of_mem_iSup, the argument of Lemma sha-siegel-step-seeds-induced-hilbert-form-galois-component(1)) and sat_two (iwahori_level_zero_eq, hecke_Tj, the Hecke eigenvalues satEig of IsSatVecGL). (7) is hinf: cap_all (cap_act: TR.lie_rt_fin, GlobalTheory.CapAdInv.qdet_semiconj, qdet_ad, TR.cap_comm), hz with TR.Zalg_stable, HasInfChar.twist (cap_mulChi) and expo_abs; regularity is the Nodup and integrality bullets. (8) is HChar.locInf_mul, sign_abs, sign_hol and sign_eq. (9) is P.gkR (Published.GelfandKazhdanReal), TR.CuspSub.isSelfDualBy_of_iso and P.smoR (Published.StrongMultOneReal), with LocSph.eventually_mem, UnrAe.finite_ramified, HasSat.dual and HasSat.twist. (10) is the last bullet: iwahori_level_zero_eq, h6 to h9. About 700 lines with the dependencies.

\begin{lemma}\label{sha-siegel-step-seeds-induced-hilbert-form-exists}
Assume the Galois representations as in Theorem~\ref{sha-residually-reducible-automorphy}. Let $\ell$ be an odd prime, $\iota:\bar{\mathbf Q}_\ell\cong\mathbf C$, and write $x\equiv y$ for $|x-y|<1$ in $\bar{\mathbf Q}_\ell$. Let $E\subseteq\bar{\mathbf Q}$ be a totally real number field with $[E:\mathbf Q]\ge2$ such that every prime $w$ of $E$ above $\ell$ has norm $\ell$ and $\ell\notin w^2$; write $G_E=\operatorname{Gal}(\bar{\mathbf Q}/E)$, and $D_A(E)\supseteq I_A(E)$ for the decomposition and inertia groups of $G_E$ at a place $A$ of $\bar{\mathbf Q}$. Let $\xi:G_E\to\bar{\mathbf Q}_\ell^\times$ be a homomorphism with open kernel such that $\xi(c)=-1$ for every complex conjugation $c\in G_E$, $\xi$ is trivial on $I_A(E)$ for every place $A$ above $\ell$, and $\xi$ is nontrivial on $I_A(E)$ for some place $A$ above a prime of $E$ of norm $2$. Let $P$ be a set of primes of $E$, and suppose that there are a prime $w_1\in P$ and a prime $w_2\notin P$ above $\ell$, places $A_i$ above $w_i$ and elements $\sigma_i\in D_{A_i}(E)$ with $\xi(\sigma_i)\not\equiv1$, for $i=1,2$. Then there are integers $k\ge2$ and $t\ge2$, both divisible by $\ell-1$, and an ordinary form datum $(c,m,a,\pi,r)$ for $\xi$ of weight $-k$ (Definition~\ref{sha-siegel-step-seeds-induced-hilbert-form-ordinary}), such that:
\begin{enumerate}
\item $m(w)=(t,-k-t)$ for the primes $w\in P$ above $\ell$, and $m(w)=(0,-k)$ for the primes $w\notin P$ above $\ell$;
\item for $s\in G_E$ the coefficients of the characteristic polynomial of $r(s)$ are congruent to those of $(X-1)(X-\xi(s))$;
\item for every prime $w$ above $\ell$, place $A$ above $w$ and Frobenius element $\sigma\in D_A(E)$, $a(w)=\xi(\sigma)^{-1}$ if $w\in P$, and $a(w)=1$ if $w\notin P$.
\end{enumerate}
\end{lemma}
\leanref{PositivityTheory/AutomorphicForms/PublishedInputs/JacquetLanglands/TorusKernels/IsHolVOf.lean}{UpperTheory.Esk.PropMV}
\leanref{Development/AutomorphicArithmetic/AutomorphicVarieties/ModularityArgument/ThresholdArguments/DenseThresholdCases/SVRPropMV.lean}{UpperTheory.Esk.pf_PropMV_sv_alt}
\begin{proof}
Write $x\equiv y$ also for complex numbers $x,y$ with $\iota^{-1}(x)\equiv\iota^{-1}(y)$. For a prime $w$ of $E$, let $\mathcal O_w$ be the ring of integers of $E_w$ and $\varpi_w$ a uniformizer.
\begin{enumerate}
\item[(1)] \emph{The Hecke character of $\xi$.} As $\ker\xi$ is open, $G_E/\ker\xi$ is finite, of order $m$ say; so $\xi$ is continuous, $\xi^m=1$ and $|\xi(g)|=1$ for every $g\in G_E$. By global class field theory \cite[Chapter~VI]{Neu99} there is a Hecke character $\omega$ of $E$ with $\omega^m=1$ such that $\xi(g)=\iota^{-1}(\omega(x))$ whenever a reciprocity map sends the idele $x$ to the image of $g$ in $G_E^{\mathrm{ab}}$. Let $c(w)$ be the conductor exponent of $\omega$ at $w$, the least $n\ge0$ such that $\omega_w$ is trivial on the units congruent to $1$ modulo $w^n$. It vanishes exactly when $\omega$ is unramified at $w$, so at all but finitely many primes.
\item[(2)] \emph{Properties of $\omega$.}
\begin{enumerate}
\item[(a)] For a prime $w$ and a place $A$ above $w$, $\xi$ is trivial on $I_A(E)$ if and only if $c(w)=0$, and then $\iota^{-1}(\omega_w(\varpi_w))=\xi(\sigma^{-1})$ for every Frobenius element $\sigma\in D_A(E)$.
\item[(b)] By (a), $c(w)=0$ for every prime $w$ above $\ell$, as $\xi$ is trivial on the inertia groups above $\ell$; and $c(u)\ne0$ for the prime $u$ of norm $2$ of the hypothesis.
\item[(c)] $\omega_v(-1)=-1$ at every real place $v$: a reciprocity map sends the idele with component $-1$ at $v$ and $1$ elsewhere to the image of a complex conjugation $c_v\in G_E$ at $v$, and $\xi(c_v)=-1$.
\item[(d)] Let $w$ be a prime above $\ell$, $A$ a place above $w$ and $\sigma\in D_A(E)$ with $\xi(\sigma)\not\equiv1$. Then $\omega_w(\varpi_w)\not\equiv1$. Indeed, suppose that $\omega_w(\varpi_w)\equiv1$. As $\omega$ is unramified at $w$ by (b), every value of $\omega_w$ is a power of $\omega_w(\varpi_w)$, so $\xi$ takes values congruent to $1$ on the image of $E_w^\times$ under the reciprocity map. The image of $D_A(E)$ in $G_E^{\mathrm{ab}}$ is the closure of the image of $E_w^\times$, and the $x\in\bar{\mathbf Q}_\ell^\times$ with $x\equiv1$ form a closed subgroup; so $\xi(\sigma)\equiv1$, which is false.
\end{enumerate}
\item[(3)] \emph{An eigenform of parallel weight.} Let $c_\ell$ and $e$ be as in Lemma~\ref{sha-siegel-step-seeds-induced-hilbert-form-exists-eisenstein} for $\omega$, $c$ and $P$. By (1) and (2), with the primes $w_1\in P$ and $w_2\notin P$ of the hypothesis, that lemma applies, with $N_0=3$. It gives an integer $k\ge3$ divisible by $\ell-1$ and a holomorphic eigenform $\phi$ of parallel weight $k+1$, level $c_\ell$ and central character $\omega$ that is a cusp form, with eigenvalues $\lambda(w)\equiv e(w)$ for every prime $w$. Then $c_\ell$ vanishes at all but finitely many primes. For $w$ above $\ell$, $|\iota^{-1}(e(w))|=1$, as $e(w)$ is $1$ or, by (2)(a) and (b), $\omega_w(\varpi_w)=\iota(\xi(\sigma^{-1}))$ for a Frobenius element $\sigma$.
\item[(4)] \emph{The primes above $\ell$.} For an infinite place $v$ of $E$, let $w_v$ be the prime above $\ell$ induced by the embedding $\iota^{-1}\circ v$ of $E$ into $\bar{\mathbf Q}_\ell$. Every prime above $\ell$ is induced by an embedding of $E$ into $\bar{\mathbf Q}_\ell$, which is $\iota^{-1}\circ v$ for an infinite place $v$, as $E$ is totally real. As every prime $w$ above $\ell$ has norm $\ell$ and $\ell\notin w^2$, $E_w$ is $\mathbf Q_\ell$ and has one embedding into $\bar{\mathbf Q}_\ell$, so two infinite places that induce the same prime are equal. So $v\mapsto w_v$ is a bijection from the infinite places of $E$ to the primes above $\ell$. By (3), $|\iota^{-1}(\lambda(w_v))|=1$ for every $v$, as $\lambda(w_v)\equiv e(w_v)$. Also $c_\ell(w)=1$ for $w$ above $\ell$, and $c_\ell(u)\ne0$ for the prime $u$ of norm $2$: $c_\ell(u)=1$ if $u$ lies above $\ell$, and $c_\ell(u)=c(u)\ne0$ by (2)(b) otherwise.
\item[(5)] \emph{Raising the weight.} Let $S$ be the set of infinite places $v$ with $w_v\in P$. By (3) and (4), Lemma~\ref{sha-siegel-step-seeds-induced-hilbert-form-exists-weight} applies to $k_0=k+1\ge4$, $c_\ell$, $\omega$, $\phi$, $\lambda$, $S$ and $N_0=2$. It gives an integer $t\ge2$ divisible by $\ell-1$, a nonzero holomorphic form $\phi'$ that is a cusp form, of weight $k+1+2t$ at the places in $S$ and $k+1$ at the other infinite places, of level $c_\ell$ and central character $\omega$, and complex numbers $\lambda'(w)$ with $\operatorname N(w)^{(k+1)/2-1}U_w\phi'=\lambda'(w)\phi'$ for every prime $w$. By (4), the primes $w_v$ with $v\in S$ are the primes in $P$ above $\ell$. So, by the ultrametric inequality, $\ell^t\lambda'(w)\equiv\lambda(w)\equiv e(w)$ for the primes $w\in P$ above $\ell$, and $\lambda'(w)\equiv\lambda(w)\equiv e(w)$ for the other primes $w$.
\item[(6)] \emph{The automorphic representation.} Let $n(v)=t$ for $v\in S$ and $n(v)=0$ for the other infinite places $v$. Lemma~\ref{sha-siegel-step-seeds-induced-hilbert-form-exists-representation}, for the weight $k+1\ge2$, $n$, $c_\ell$, $\omega$, $\phi'$ and $\lambda'$, gives a cuspidal automorphic representation $\pi$ of $GL_2(\mathbf A_E)$, regular algebraic and polarized, with the infinitesimal character of weight $(n(v),-k-n(v))$ at every infinite place $v$, such that for every prime $w$, $\pi_w$ has a nonzero vector fixed by $K_1(w^{c_\ell(w)})$ on which $U_w$ acts by $\lambda'(w)$, and such that $\pi$ has Satake parameters $B_w$ with $\prod_{b\in B_w}(X-\operatorname N(w)^{1/2}b)=X^2-\lambda'(w)X+\omega_w(\varpi_w)\operatorname N(w)^k$ at every prime $w$ with $c_\ell(w)=0$.
\item[(7)] \emph{The Galois representation.} By the assumed Galois representations, there is a continuous homomorphism $r:G_E\to GL_2(\bar{\mathbf Q}_\ell)$ attached to $\pi$ through $\iota$.
\item[(8)] \emph{Determinant and residues.} Let $w$ be a prime not above $\ell$ with $c(w)=0$ such that $r$ has Frobenius eigenvalues $\operatorname N(w)^{1/2}B$ at $w$ for Satake parameters $B$ of $\pi$ at $w$; all but finitely many primes are such, and $B=B_w$ as Satake parameters are unique. Let $\sigma$ be a Frobenius element at a place above $w$. By (6), the characteristic polynomial of $\iota(r(\sigma^{-1}))$ is $X^2-\lambda'(w)X+\omega_w(\varpi_w)\operatorname N(w)^k$. As $\chi_{\mathrm{cyc}}(\sigma)=\operatorname N(w)$, (2)(a) gives $\det r(\sigma^{-1})=\xi(\sigma^{-1})\chi_{\mathrm{cyc}}(\sigma)^k$, and with (5) it gives $\operatorname{tr}r(\sigma^{-1})=\iota^{-1}(\lambda'(w))\equiv\iota^{-1}(e(w))=1+\xi(\sigma^{-1})$. These Frobenius elements are dense in $G_E$ (Theorem~\ref{lit-chebotarev-degree-one}); $r$, $\xi$ and $\chi_{\mathrm{cyc}}$ are continuous, and $\{x:|x|<1\}$ is closed in $\bar{\mathbf Q}_\ell$. So $\det r(g^{-1})=\xi(g^{-1})\chi_{\mathrm{cyc}}(g)^k$ and $\operatorname{tr}r(g^{-1})\equiv1+\xi(g^{-1})$ for every $g\in G_E$, and $\det r=\chi_{\mathrm{cyc}}^{-k}\xi$. As $\ell-1$ divides $k$, $\chi_{\mathrm{cyc}}(g)^{-k}\equiv1$; as $|\xi(g)|=1$, $\det r(g)\equiv\xi(g)$. So the coefficients of the characteristic polynomial of $r(g)$ are congruent to those of $(X-1)(X-\xi(g))$ for every $g$, which is (2).
\item[(9)] \emph{The ordinary form datum.} Let $m$ be as in (1) of the statement, and $a(w)=\iota^{-1}(e(w))$ for the primes $w$ above $\ell$. Then $(c_\ell,m,a,\pi,r)$ is an ordinary form datum for $\xi$ of weight $-k$. Condition (a) holds by (6), and (b) by (7) and (8). For (c), $\ker\xi$ is open, $c_\ell$ vanishes at all but finitely many primes, and, by (2)(a) and (b), $\xi$ is trivial on $I_A(E)$ if and only if $c(w)=0$, that is, if and only if $c_\ell(w)=0$ or $w$ lies above $\ell$. For (d), if $c_\ell(w)=0$, the vector of (6) is fixed by $K_1(w^0)=GL_2(\mathcal O_w)$. For (e), if $c_\ell(w)\ne0$ and $w$ does not lie above $\ell$, then $e(w)=1$, so $\lambda'(w)\equiv1$ by (5) and $|\iota^{-1}(\lambda'(w))|=1$. For (f), let $w$ lie above $\ell$. Then $c_\ell(w)=1$; $m_1(w)+2\le m_0(w)$, as $2\le k+2t$ and $2\le k$; $m_0(w)+m_1(w)=-k$; $|a(w)|=1$ by (3); as $\operatorname N(w)=\ell$, (5) gives $\iota^{-1}(\lambda'(w))\operatorname N(w)^{m_0(w)}\equiv a(w)$ for the eigenvalue $\lambda'(w)$ of $U_w$ of (6); and for an embedding $\tau:E\to\bar{\mathbf Q}_\ell$ that induces $w$, the infinite place $v$ induced by $\iota\circ\tau$ has $w_v=w$, so $v\in S$ if and only if $w\in P$, and by (6) $\pi$ has at $v$ the infinitesimal character of weight $(n(v),-k-n(v))=m(w)$.
\item[(10)] \emph{Conclusion.} The integers $k\ge3$ and $t\ge2$ are divisible by $\ell-1$; (1) holds by the choice of $m$, and (2) by (8). For (3), let $w$ lie above $\ell$, $A$ a place above $w$ and $\sigma\in D_A(E)$ a Frobenius element. If $w\in P$, then $a(w)=\iota^{-1}(\omega_w(\varpi_w))=\xi(\sigma)^{-1}$ by (2)(a) and (b); if $w\notin P$, then $a(w)=1$.
\end{enumerate}
% pf_PropMV_sv_alt (SVRPropMV.lean, 153 lines), followed step by step: (1) is hPH.litN.aut.art, hc, hξm, hn, HeckeGal.exists_unique_hchar_of_finite_order and PropMVAux.cond_ex; (2)(a) is hine and hfr (IsLAdicOf.unr_iff, IsLAdicOf.frob_unr), (b) is hcl and hex, (c) is hsg (HilbertFromArtin.exists_conj, art_sign_eq), (d) is hnc with PropMVAux.dec_cong (ha.dec, loc_zpow, closed_oneU); (3) is hMH (Sd.PropMHol) with B = 3, hfinL and heu; (4) is Wc.exL (the LDat D: pl, inj through exL_uniq, surj through exists_embedding_place, embAt), hu2 and hunit; (5) is hWC (Esk.HolWC) with S = {v | D.pl v in P} and B = 2, hlam' and haway (prod_at, prod_away, nu_lt, add_lt); (6) is hHR (Esk.HRep); (7) is Gal.SV.exRS_alt; (8) is PropMVAux.det_res (Published.ChebotarevDegreeOne.dense, cyc_mem) with HasSat.unique, tr_det and cyc_frob; (9) and (10) are the final refine. Of the three inputs, Lemma sha-siegel-step-seeds-induced-hilbert-form-exists-representation (Esk.HRep, by Esk.pf_HRep_alt, about 700 lines) is written out with its parts and has no open step. Lemmas sha-siegel-step-seeds-induced-hilbert-form-exists-eisenstein (Sd.PropMHol, by Z.Bc.propMHolZ_alt, about 3,400) and sha-siegel-step-seeds-induced-hilbert-form-exists-weight (Esk.HolWC, by Esk.holWC, about 14,700) are written out at their top level; the open steps left are in the lemmas they rest on, which their comments list. The whole is about 20,600 lines of Lean with the dependencies, net of the definitions, of which pf_PropMV_sv_alt with PfPropMV.lean is about 1,600; about 19,200 of them are not shared with Lemma sha-siegel-step-seeds-induced-hilbert-form-galois.
\end{proof}

\begin{lemma}\label{sha-siegel-step-seeds-induced-hilbert-form-galois-jacquet}
Let $E$ be a number field, $w$ a prime of $E$, $\mathcal O_w$ the ring of integers of $E_w$, $\varpi_w$ a uniformizer of $E_w$ and $q=\operatorname N(w)$. For $c\ge0$ let $K_1(w^c)$ be the group of $g\in GL_2(\mathcal O_w)$ whose bottom row is congruent to $(0,1)$ modulo $w^c$, and let $U_w$ be the operator of the double coset $K_1(w^c)\operatorname{diag}(\varpi_w,1)K_1(w^c)$ on the vectors fixed by $K_1(w^c)$ in a smooth representation of $GL_2(E_w)$. Let $\sigma$ be an irreducible smooth representation of $GL_2(E_w)$ on a space $V$, let $c\ge1$ and $\lambda\in\mathbf C^\times$, and suppose that $\sigma$ has a nonzero vector fixed by $K_1(w^c)$ on which $U_w$ acts by $\lambda$. Then there are smooth characters $\psi_0,\psi_1$ of $E_w^\times$ such that $\sigma$ is a subquotient of the normalized parabolic induction $\psi_0\times\psi_1$ (Definition~\ref{gal-irreducible-classes}), from the upper triangular Borel subgroup, of the character $\begin{pmatrix}t_0&x\\0&t_1\end{pmatrix}\mapsto\psi_0(t_0)\psi_1(t_1)$; $\psi_0$ is trivial on $\mathcal O_w^\times$; and $\lambda^2=q\,\psi_0(\varpi_w)^2$.
\end{lemma}
\leanref{Development/AutomorphicArithmetic/AutomorphicVarieties/ModularityArgument/AutomorphyLifting/PrimeIndexConstruction/LocTrans.lean}{UpperTheory.V.Z.Win.Birth.ULoc}
\leanref{Development/AutomorphicArithmetic/AutomorphicVarieties/ModularityArgument/AutomorphyLifting/PrimeIndexConstruction/PsOfForm.lean}{UpperTheory.V.Z.Win.Birth.uLoc}
\leanref{Development/AutomorphicArithmetic/AutomorphicVarieties/IdeleGenerators/WeightChange.lean}{UpperTheory.V.Sd.JacqOfU}
\leanref{Development/AutomorphicArithmetic/AutomorphicVarieties/IdeleGenerators/WeightChange.lean}{UpperTheory.V.Sd.UopSum}
\leanref{Development/AutomorphicArithmetic/AutomorphicVarieties/IdeleGenerators/HolomorphicIrreducibility/UopSum.lean}{UpperTheory.V.Sd.uopSum}
\leanref{Development/AbelianVarietyTheory/Proofs/SatakeParameters/SphericalRepGLExistsSatakeParameters.lean}{GlobalTheory.SatakeExists.isSubquotient_of_eigenform}
\begin{proof}
Let $v\in V$ be the vector of the hypothesis. For $y\in E_w^\times$ and $x\in E_w$ write $a(y,x)=\begin{pmatrix}y&x\\0&1\end{pmatrix}$, $n(x)=a(1,x)$ and $a(y)=a(y,0)$, and let $s=\begin{pmatrix}0&1\\1&0\end{pmatrix}$. Let $S\subset\mathcal O_w$ be a set of representatives of $\mathcal O_w/\varpi_w\mathcal O_w$; it has $q$ elements. The matrices $n(x)$ with $x\in\mathcal O_w$ and $a(u)$ with $u\in\mathcal O_w^\times$ lie in $K_1(w^c)$, so they fix $v$.
\begin{enumerate}
\item[(1)] \emph{The operator.} As $c\ge1$, a computation of the cosets in the double coset gives $U_wz=\sum_{x\in S}\sigma(a(\varpi_w,x))z$ for every vector $z$ fixed by $K_1(w^c)$. So $\lambda v=\sum_{x\in S}\sigma(a(\varpi_w,x))v$.
\item[(2)] \emph{The Jacquet module.} Let $V(N)$ be the span of the vectors $\sigma(n(x))z-z$ with $x\in E_w$ and $z\in V$, let $V_N=V/V(N)$, and let $p:V\to V_N$ be the projection. As $a(y)n(x)=n(yx)a(y)$, $V(N)$ is stable under the $\sigma(a(y))$, and $\tau(y)p(z)=p(\sigma(a(y))z)$ defines an action $\tau$ of $E_w^\times$ on $V_N$. As $a(y,x)=n(x)a(y)$, $p(\sigma(a(y,x))z)=\tau(y)p(z)$.
\item[(3)] \emph{The image of $v$.} Let $W\subseteq V(N)$ be the span of the vectors $\sigma(n(x))z-z$ with $x\in\mathcal O_w$ and $z\in V$.
\begin{enumerate}
\item[(a)] For $z\in V(N)$, $\sigma(a(\varpi_w^m))z\in W$ for all large $m$. Indeed, $\sigma(a(\varpi_w^m))(\sigma(n(x))z'-z')=\sigma(n(\varpi_w^mx))z''-z''$ with $z''=\sigma(a(\varpi_w^m))z'$, and $\varpi_w^mx\in\mathcal O_w$ for large $m$; for a finite sum of such vectors, take $m$ large for all of them.
\item[(b)] For $j\ge0$, if $\sigma(a(\varpi_w^{j+1}))v\in W$, then $\sigma(a(\varpi_w^j))v\in W$. Indeed, by (1), $\lambda\sigma(a(\varpi_w^j))v=\sum_{x\in S}\sigma(n(\varpi_w^jx))\sigma(a(\varpi_w^{j+1}))v$; each term lies in $W$, as $\sigma(n(y))z-z\in W$ for $y\in\mathcal O_w$; and $\lambda\ne0$.
\item[(c)] Every element of $W$ is killed by an element of the group algebra of $n(\mathcal O_w)$ with nonzero augmentation. Let $x\in\mathcal O_w$ and $z\in V$. The stabilizer of $z$ is open, so $\sigma(n(Mx))z=z$ for some power $M\ge1$ of the residue characteristic, and $\sum_{0\le i<M}n(x)^i$ kills $\sigma(n(x))z-z$ and has augmentation $M$. For a sum of such vectors, take the product of these elements of the commutative group algebra.
\end{enumerate}
Suppose that $p(v)=0$. By (a) and (b), $v\in W$, and by (c) some element of the group algebra with augmentation $e\ne0$ kills $v$. As $v$ is fixed by $n(\mathcal O_w)$, that element acts on $v$ by $e$, so $v=0$, which is false. So $p(v)\ne0$.
\item[(4)] \emph{A character.} By (1) and (2), $q\,\tau(\varpi_w)p(v)=\lambda p(v)$, and $\tau(u)p(v)=p(v)$ for $u\in\mathcal O_w^\times$. Every element of $E_w^\times$ is $u\varpi_w^n$ with $u\in\mathcal O_w^\times$ and $n\in\mathbf Z$, so there is a character $\mu_1$ of $E_w^\times$, trivial on $\mathcal O_w^\times$, with $\tau(y)p(v)=\mu_1(y)p(v)$ for every $y$ and $q\mu_1(\varpi_w)=\lambda$.
\item[(5)] \emph{Finite generation.} As $\sigma$ is irreducible and smooth, the centre of $GL_2(E_w)$ acts on it by a character $\omega_\sigma$. Every $g\in GL_2(E_w)$ is $z\,a(y,x)\,s\,n(u)$ or $z\,a(y,x)\,s\,n(u)\,s$ with $z$ scalar, $y\in E_w^\times$, $x\in E_w$ and $u\in\mathcal O_w$, according as $g_{21}\ne0$ and $g_{22}/g_{21}\in\mathcal O_w$ or not. As $\sigma$ is irreducible, the $\sigma(g)v$ span $V$; so, by (2), $V_N$ is spanned by the $\tau(y)p(\sigma(s)z)$ with $y\in E_w^\times$ and $z$ among the $\sigma(n(u))v$ and $\sigma(n(u))\sigma(s)v$ with $u\in\mathcal O_w$. There are finitely many such $z$, as $\mathcal O_w$ is compact and stabilizers are open. So $V_N$ is a finitely generated module over the group algebra $\mathbf C[E_w^\times]$, acting through $\tau$.
\item[(6)] \emph{A form.} Let $J$ be the kernel of the homomorphism $\mathbf C[E_w^\times]\to\mathbf C$ that extends $\mu_1$, and let $M'=JV_N$, the span of the $\tau(y)z-\mu_1(y)z$. If $M'=V_N$, Nakayama's lemma for the finitely generated module $V_N$ gives an element of $1+J$ that kills $V_N$; but it acts on $p(v)\ne0$ by $1$. So $M'\ne V_N$, and there is a nonzero linear form $l$ on $V_N$ that vanishes on $M'$, that is, with $l(\tau(y)z)=\mu_1(y)l(z)$.
\item[(7)] \emph{The Borel subgroup.} Let $L=l\circ p$, a nonzero linear form on $V$ that vanishes on $V(N)$, and let $\mu_2=\omega_\sigma\mu_1^{-1}$. For $t_0,t_1\in E_w^\times$ and $x\in E_w$, the matrix $\begin{pmatrix}t_0&t_1x\\0&t_1\end{pmatrix}$ is the scalar $t_1$ times $a(t_0/t_1,x)$, so by (2) and (6), $L(\sigma(b)z)=\omega_\sigma(t_1)\mu_1(t_0/t_1)L(z)=\mu_1(t_0)\mu_2(t_1)L(z)$ for every upper triangular $b$ with diagonal $(t_0,t_1)$.
\item[(8)] \emph{The principal series.} Let the diagonal torus act on $V_N$ by $\operatorname{diag}(t_0,t_1)p(z)=|t_0/t_1|^{-1/2}p(\sigma(\operatorname{diag}(t_0,t_1))z)$, the normalization by the square root of the modulus character of the Borel subgroup. By (7), $l$ transforms under this action by the character $\operatorname{diag}(t_0,t_1)\mapsto\psi_0(t_0)\psi_1(t_1)$, where $\psi_0=\mu_1|\cdot|^{-1/2}$ and $\psi_1=\mu_2|\cdot|^{1/2}$. This character is smooth, as the stabilizer of $p(z)$ in the torus is open for a $z$ with $L(z)\ne0$. By Frobenius reciprocity \cite{BZ77}, $l$ gives a nonzero $GL_2(E_w)$-equivariant map from $\sigma$ to $\psi_0\times\psi_1$, which is injective as $\sigma$ is irreducible. So $\sigma$ is a subquotient of $\psi_0\times\psi_1$.
\item[(9)] \emph{The eigenvalue.} The character $\psi_0$ is trivial on $\mathcal O_w^\times$, as $\mu_1$ and $|\cdot|$ are. As $|\varpi_w|=q^{-1}$, $\lambda=q\mu_1(\varpi_w)=q^{1/2}\psi_0(\varpi_w)$, so $\lambda^2=q\,\psi_0(\varpi_w)^2$.
\end{enumerate}
\end{proof}
% uLoc (PsOfForm.lean, 136 lines) is the private jacq, a copy of Sd.jacqOfU without its hypothesis bundle, followed by the private ps_of_form. jacq is the private gen of JacqOfU.lean (457 lines): (1) is uopSum (UopSum.lean, 258 lines) with reps, affine_mem_iwahori_level and unit_mem; (2) is N, hN, tau and hq; (3) is not_mem_ker, with h1 for (a), h2 for (b), fixed_coinv for (c) and tors for the orders; (4) is the subgroup H, hπH, hH and mu; (5) is cent (LocGL.exists_smooth_hasCentralChar), iwasawa, fin_orb and hgen; (6) is quot_of_sub (Submodule.exists_sub_one_mem_and_smul_eq_zero_of_fg_of_le_smul); (7) is the last field of gen. (8) is ps_of_form: hker, hle with deltaHalf, hs and GlobalTheory.SatakeExists.isSubquotient_of_eigenform (Frobenius reciprocity through Model.evalOne_bijective, then the kernel is a subrepresentation); (9) is h0 with nabs_unit and nabs_pi. About 900 lines with these files.

\begin{lemma}\label{sha-siegel-step-seeds-induced-hilbert-form-galois-component}
Let $E$ be a totally real number field, $\pi$ a cuspidal automorphic representation of $GL_2(\mathbf A_E)$ and $w$ a prime of $E$, with $\mathcal O_w$ the ring of integers of $E_w$; write $\pi_w$ for the action of $GL_2(E_w)$ on the space of $\pi$. For $c\ge0$ let $K_1(w^c)$ and $U_w$ be the group and the operator of Lemma~\ref{sha-siegel-step-seeds-induced-hilbert-form-galois-jacquet}.
\begin{enumerate}
\item If $\pi_w$ has a nonzero vector fixed by $GL_2(\mathcal O_w)$, then $\pi$ has Satake parameters at $w$.
\item Let $\sigma$ be a local component of $\pi$ at $w$, $c\ge0$ and $\lambda\in\mathbf C$. If $\pi_w$ has a nonzero vector fixed by $K_1(w^c)$ on which $U_w$ acts by $\lambda$, then so has $\sigma$.
\item Let $\omega$ be a Hecke character of $E$ with $\phi(zg)=\omega(z)\phi(g)$ for $z\in\mathbf A_E^\times$, $g\in GL_2(\mathbf A_E)$ and $\phi$ in the space of $\pi$. Then the centre $E_w^\times$ of $GL_2(E_w)$ acts on every local component of $\pi$ at $w$ by $\omega_w$.
\end{enumerate}
\end{lemma}
\leanref{Development/AutomorphicArithmetic/AutomorphicVarieties/ModularityArgument/AutomorphyLifting/PrimeIndexConstruction/LocTrans.lean}{UpperTheory.V.Z.Win.Birth.LocTrans}
\leanref{Development/AutomorphicArithmetic/AutomorphicVarieties/ModularityArgument/AutomorphyLifting/PrimeIndexConstruction/LocTransAlt.lean}{UpperTheory.V.Z.Win.Birth.locTrans}
\begin{proof}
Let $\sigma$ be a local component of $\pi$ at $w$. The sum of the images of all $GL_2(E_w)$-equivariant maps $\sigma\to\pi_w$ is stable under $GL_2(E_w)$ and under every linear operator on the space of $\pi$ that commutes with $GL_2(E_w)$. These operators include the elements of $GL_2(\mathbf A_{E,f})$ with component $1$ at $w$, which together with $GL_2(E_w)$ generate $GL_2(\mathbf A_{E,f})$, the maximal compact subgroup at infinity and the Lie algebra at infinity. As $\pi$ is irreducible and $\sigma$ is a local component, that sum is the whole space: $\pi_w$ is a sum of images of $\sigma$.

(1) By admissibility, $\pi$ has a local component $\sigma$ at $w$. Let $x\ne0$ be fixed by $GL_2(\mathcal O_w)$; it lies in the sum of finitely many images $f_1(\sigma),\dots,f_n(\sigma)$. Adding these images one at a time to the subspace $0$, as long as the sum does not contain $x$, gives a $GL_2(E_w)$-stable subspace $W$ with $x\notin W$ and an $i$ with $x=y+f_i(u)$ for some $y\in W$ and $u$ in the space of $\sigma$. As $f_i(u)\notin W$ and $\sigma$ is irreducible, $f_i^{-1}(W)=0$. For $k\in GL_2(\mathcal O_w)$, $f_i(\sigma(k)u-u)=y-\pi_w(k)y\in W$, so $\sigma(k)u=u$; and $u\ne0$. So $\sigma$ is spherical, with some Satake parameters $B$ \cite{Car79}, and as $\pi_w$ is a sum of images of $\sigma$, $\pi$ has Satake parameters $B$ at $w$.

(2) Let $x\ne0$ be such a vector. Let $M$ be a $GL_2(E_w)$-stable subspace maximal among those that do not contain $x$ (Zorn's lemma). Some image $f(\sigma)$ is not contained in $M$, and then $f^{-1}(M)=0$ as $\sigma$ is irreducible. By the maximality of $M$, $x$ and then every image of $\sigma$ lie in $M+f(\sigma)$, so the space of $\pi$ is $M\oplus f(\sigma)$. The projection $p$ onto $\sigma$ along $M$ is equivariant, with $p(x)\ne0$, so $p(x)$ is fixed by $K_1(w^c)$. On the vectors fixed by $K_1(w^c)$ of any smooth representation $\rho$, $U_w$ is $\sum_j\rho(g_j)$ for representatives $g_j$ of the finitely many cosets in the double coset, the same for every $\rho$. So $U_w$ commutes with $p$, and $U_wp(x)=\lambda p(x)$.

(3) Let $f$ be a nonzero equivariant map from $\sigma$ to $\pi_w$; it is injective, as $\sigma$ is irreducible. For $z\in E_w^\times$, the scalar matrix $z$ in $GL_2(E_w)$ acts on the space of $\pi$ by right translation by the scalar idele with component $z$ at $w$ and $1$ elsewhere, which is central, hence by $\omega_w(z)$. So $f(\sigma(z)u)=\omega_w(z)f(u)$, and $\sigma(z)u=\omega_w(z)u$ for every $u$.
\end{proof}
% locTrans (LocTransAlt.lean, 73 lines). The isotypy is the private RealLocalComponent.isotypic (LocCompR.lean), through bot_or_top and pi.irred. (1) is RealLocalComponent.exists_comp (from Published.AdmissibleReal), CuspRep.LocSph.exists_fixed_of_mem_iSup, level_zero_eq_maximal_compact and the private Nf.satLoc' (SatInd.exists_hasSatGL_of_isIrreducible, then isSatVecGL_of_mem_iSup). (2) is the private Wc.sep_of_sum (FinProj.lean, Zorn) and hasU_of_map with Wc.hkReps uopSum. (3) is scalar_fin and CuspRep.scalar_comm. About 450 lines with these files and the parts of LocCompR.lean they use.

\begin{lemma}\label{sha-siegel-step-seeds-induced-hilbert-form-galois-central}
Let $\ell$ be a prime, $\iota:\bar{\mathbf Q}_\ell\cong\mathbf C$, and $E\subseteq\bar{\mathbf Q}$ a totally real number field; write $G_E=\operatorname{Gal}(\bar{\mathbf Q}/E)$, and $I_A(E)$ for the inertia group of $G_E$ at a place $A$ of $\bar{\mathbf Q}$. Let $\xi:G_E\to\bar{\mathbf Q}_\ell^\times$ be a homomorphism with open kernel, $s\in\mathbf Z$, $\pi$ a cuspidal automorphic representation of $GL_2(\mathbf A_E)$ and $r:G_E\to GL_2(\bar{\mathbf Q}_\ell)$ a homomorphism attached to $\pi$ through $\iota$ with $\det r=\chi_{\mathrm{cyc}}^s\xi$. Let $\omega$ be a Hecke character of $E$ with $\phi(zg)=\omega(z)\phi(g)$ for $z\in\mathbf A_E^\times$, $g\in GL_2(\mathbf A_E)$ and $\phi$ in the space of $\pi$. Let $w$ be a prime of $E$, $\mathcal O_w$ the ring of integers of $E_w$ and $\varpi_w$ a uniformizer. Then:
\begin{enumerate}
\item $\bigl(\omega_w(\varpi_w)\operatorname N(w)^{s+1}\bigr)^N=1$ for some $N\ge1$;
\item if $\omega_w$ is trivial on $\mathcal O_w^\times$, then $\xi$ is trivial on $I_A(E)$ for every place $A$ of $\bar{\mathbf Q}$ above $w$.
\end{enumerate}
\end{lemma}
\leanref{Development/AutomorphicArithmetic/AutomorphicVarieties/ModularityArgument/AutomorphyLifting/PrimeIndexConstruction/LocTrans.lean}{UpperTheory.V.Z.Win.Birth.DetCent}
\leanref{Development/AutomorphicArithmetic/AutomorphicVarieties/ModularityArgument/AutomorphyLifting/PrimeIndexConstruction/DetCent.lean}{UpperTheory.V.Z.Win.Birth.detCent}
\begin{proof}
For a prime $v$ write $\mathcal O_v$ for the ring of integers of $E_v$ and $\varpi_v$ for a uniformizer.
\begin{enumerate}
\item[(a)] \emph{The Hecke character of $\xi$.} As $\ker\xi$ is open, $G_E/\ker\xi$ is finite, of order $m$ say; so $\xi$ is continuous and $\xi^m=1$. By global class field theory \cite[Chapter~VI]{Neu99} there is a Hecke character $\chi$ of $E$ with $\chi^m=1$ such that $\xi(g)=\iota^{-1}(\chi(x))$ whenever a reciprocity map sends the idele $x$ to the image of $g$ in $G_E^{\mathrm{ab}}$. So, for every prime $v$ and place $A$ above $v$, $\xi$ is trivial on $I_A(E)$ if and only if $\chi$ is unramified at $v$, and then $\iota(\xi(\sigma^{-1}))=\chi_v(\varpi_v)$ for every Frobenius element $\sigma$ at $A$.
\item[(b)] \emph{Almost all primes.} Let $P$ be the Hecke character $|\cdot|^{s+1}$; $P_v$ is trivial on $\mathcal O_v^\times$ and $P_v(\varpi_v)=\operatorname N(v)^{-s-1}$. Let $v$ be a prime not above $\ell$ at which $\chi$ is unramified and $\pi$ has Satake parameters $B$ such that $r$ has Frobenius eigenvalues $\operatorname N(v)^{1/2}B$ at $v$; all but finitely many primes are such. The scalar matrix $\varpi_v$ acts on the space of $\pi$ by $\omega_v(\varpi_v)$; it is its own double coset modulo $GL_2(\mathcal O_v)$, whose Hecke operator acts on the spherical vectors by $\prod_{b\in B}b$. So $\omega_v(\varpi_v)=\prod_{b\in B}b$. Let $\sigma$ be a Frobenius element at a place above $v$. Then $\iota(\det r(\sigma^{-1}))$ is the product $\operatorname N(v)\prod_{b\in B}b$ of the Frobenius eigenvalues; and $\det r(\sigma^{-1})=\chi_{\mathrm{cyc}}(\sigma)^{-s}\xi(\sigma^{-1})$, with $\chi_{\mathrm{cyc}}(\sigma)=\operatorname N(v)$ and $\iota(\xi(\sigma^{-1}))=\chi_v(\varpi_v)$ by (a). So $\omega_v(\varpi_v)=\operatorname N(v)^{-s-1}\chi_v(\varpi_v)=\chi_v(\varpi_v)P_v(\varpi_v)$.
\item[(c)] \emph{The central character.} Let $\nu=\omega\chi^{-1}P^{-1}$, and let $T$ be a finite set of primes outside which $\nu$ is unramified with $\nu_v(\varpi_v)=1$; it exists by (b). Then $\nu_v=1$ for $v\notin T$. For $k\in E^\times$, $\nu(k)=1$, so the continuous character $\nu_\infty\prod_{v\in T}\nu_v$ of $(E\otimes\mathbf R)^\times\times\prod_{v\in T}E_v^\times$ is trivial on the image of $E^\times$. That image is dense, by weak approximation \cite[Chapter~II, \S3]{Neu99}. So the character is trivial, and $\nu=1$, that is, $\omega=\chi P$.
\item[(d)] \emph{Conclusion.} As $\omega_w=\chi_wP_w$, $\omega_w(\varpi_w)\operatorname N(w)^{s+1}=\chi_w(\varpi_w)$, whose $m$-th power is $1$; this is (1). If $\omega_w$ is trivial on $\mathcal O_w^\times$, so is $\chi_w=\omega_wP_w^{-1}$; then $\chi$ is unramified at $w$, and (2) follows from (a).
\end{enumerate}
\end{proof}
% detCent (DetCent.lean, 184 lines): (a) is GlobalTheory.globalArtinReciprocity_holds, fin_order and HeckeGal.exists_unique_hchar_of_finite_order, with IsLAdicOf.frob_unr and IsLAdicOf.unr_iff; (b) is the filter_upwards of hnu: hatt.2, UnrAe.unr_ae_and_level, scalar_eq_prod with locRep_scalar for the central character, det_sat for the determinant, cycZ_of_isFrob and absVal; (c) is PositivityTheory.jumpWindows.FiniteInfinitePlaceProduct.eq_one (FiniteInfinitePlaceProduct.lean, 132 lines: loc_eq_one_of_val, CharModel.productRule and GroundTheory.Global.denseRange_units_prod_pi); (d) is the last two fields. About 600 lines with these files.

\begin{lemma}\label{sha-siegel-step-seeds-induced-hilbert-form-galois}
Assume the Galois representations as in Theorem~\ref{sha-residually-reducible-automorphy}. Let $\ell$ be an odd prime, $\iota:\bar{\mathbf Q}_\ell\cong\mathbf C$, and write $x\equiv y$ for $|x-y|<1$ in $\bar{\mathbf Q}_\ell$. Let $E\subseteq\bar{\mathbf Q}$ be a totally real number field such that every prime $w$ of $E$ above $\ell$ has norm $\ell$ and $\ell\notin w^2$; write $G_E=\operatorname{Gal}(\bar{\mathbf Q}/E)$, and $D_A(E)\supseteq I_A(E)$ for the decomposition and inertia groups of $G_E$ at a place $A$ of $\bar{\mathbf Q}$. Let $\xi:G_E\to\bar{\mathbf Q}_\ell^\times$ be a homomorphism, $s\in\mathbf Z$, and $(c,m,a,\pi,r)$ an ordinary form datum for $\xi$ of weight $s$ (Definition~\ref{sha-siegel-step-seeds-induced-hilbert-form-ordinary}). Then:
\begin{enumerate}
\item $r$ is irreducible;
\item for every place $A$ of $\bar{\mathbf Q}$ not above $\ell$, $r$ is conjugate on $I_A(E)$ to the diagonal representation with entries $1$ and $\xi$;
\item for every prime $w$ of $E$ above $\ell$ and place $A$ above $w$, there are characters $\kappa_0,\kappa_1$ of $D_A(E)$ such that $r$ is conjugate on $D_A(E)$ to an upper triangular representation with diagonal entries $\kappa_0,\kappa_1$, $\kappa_i=\chi_{\mathrm{cyc}}^{m_i(w)}$ on $I_A(E)$, and $\kappa_0(\sigma^{-1})\chi_{\mathrm{cyc}}(\sigma)^{m_0(w)}\equiv a(w)$ for every Frobenius element $\sigma\in D_A(E)$;
\item $r$ is attached through $\iota$ to a regular algebraic cuspidal automorphic representation of $GL_2(\mathbf A_E)$ that has, at every finite place of $E$, Satake parameters or only plain local components (Definition~\ref{gal-plain}).
\end{enumerate}
\end{lemma}
\leanref{Development/AutomorphicArithmetic/AutomorphicVarieties/IdeleGenerators/WeightChange.lean}{UpperTheory.V.Sd.OrdIrr}
\leanref{Development/AutomorphicArithmetic/AutomorphicVarieties/IdeleGenerators/WeightChange.lean}{UpperTheory.V.Sd.OrdAway}
\leanref{Development/AutomorphicArithmetic/AutomorphicVarieties/IdeleGenerators/WeightChange.lean}{UpperTheory.V.Sd.OrdAbove}
\leanref{Development/AutomorphicArithmetic/AutomorphicVarieties/ModularityArgument/ThresholdArguments/DenseThresholdCases/SVROrdIrr.lean}{UpperTheory.V.Sd.ordinary_irreducible_of_printed_hida}
\leanref{Development/AutomorphicArithmetic/AutomorphicVarieties/ModularityArgument/ThresholdArguments/DenseThresholdCases/AutNewform2A.lean}{UpperTheory.V.Z.Win.RR.ordAway_a_alt}
\leanref{Development/AutomorphicArithmetic/AutomorphicVarieties/ModularityArgument/ThresholdArguments/DenseThresholdCases/AutNewform2A.lean}{UpperTheory.V.Z.Win.RR.ordAbove_of_printsH}
\leanref{Development/AutomorphicArithmetic/AutomorphicVarieties/ModularityArgument/AutomorphyLifting/PrimeIndexConstruction/MixSeedPAlt.lean}{UpperTheory.V.Z.Win.Birth.PrOrd}
\leanref{Development/AutomorphicArithmetic/AutomorphicVarieties/ModularityArgument/AutomorphyLifting/PrimeIndexConstruction/PrOrd.lean}{UpperTheory.V.Z.Win.Birth.prOrd}
\leanref{Development/AutomorphicArithmetic/AutomorphicVarieties/ModularityArgument/AutomorphyLifting/PrimeIndexConstruction/OrdPlain.lean}{UpperTheory.V.Z.Win.Birth.ordPlain}
\begin{proof}
Parts (2) and (3) are local-global compatibility away from $\ell$ and at $\ell$ for the representation attached to $\pi$ by the assumed Galois representations; for Hilbert modular forms they are due to Carayol \cite{Car86} and Wiles \cite[Theorem~2.1.4]{Wil88}.

For (1), suppose that $r$ is reducible. As $r$ is continuous, a stable line gives continuous characters $c_1,c_2$ of $G_E$ such that the characteristic polynomial of $r(g)$ is $(X-c_1(g))(X-c_2(g))$ for every $g\in G_E$.
\begin{enumerate}
\item[(a)] \emph{Inertia above $\ell$.} Let $w$ be a prime above $\ell$ and $A$ a place above $w$, with $\kappa_0,\kappa_1$ as in (3). For $g\in I_A(E)$ the characteristic polynomial of $r(g)$ is also $(X-\chi_{\mathrm{cyc}}(g)^{m_0(w)})(X-\chi_{\mathrm{cyc}}(g)^{m_1(w)})$. So $c_1c_2=\chi_{\mathrm{cyc}}^{m_0(w)+m_1(w)}$ on $I_A(E)$, and each $c_1(g)$ is $\chi_{\mathrm{cyc}}(g)^{m_0(w)}$ or $\chi_{\mathrm{cyc}}(g)^{m_1(w)}$. The elements at which $c_1$ agrees with a given one of these characters form a subgroup, and a group is not the union of two proper subgroups; so $c_1=\chi_{\mathrm{cyc}}^{m_0(w)}$ on $I_A(E)$ or $c_1=\chi_{\mathrm{cyc}}^{m_1(w)}$ on $I_A(E)$. Hence $c_1/c_2=\chi_{\mathrm{cyc}}^d$ on $I_A(E)$ with $d=\pm(m_0(w)-m_1(w))\ne0$, as $m_1(w)+2\le m_0(w)$.
\item[(b)] \emph{A Hecke character.} By (a), $c_1/c_2$ is the $\ell$-adic character of an algebraic Hecke character $\chi$ of $E$ \cite[Chapter~III]{Ser68}. As $E$ is totally real, the infinity type of $\chi$ is parallel, given by an integer $n$, and $|\chi(\varpi_w)|=\operatorname{N}(w)^{-n}$ for a uniformizer $\varpi_w$ of $E_w$ at every prime $w$ at which $\chi$ is unramified.
\item[(c)] \emph{The weight.} Let $w$ be a prime not above $\ell$ at which $\chi$ is unramified and $\pi_w$ has Satake parameters that correspond through $\iota$ to the eigenvalues of $r$ at the Frobenius elements; all but finitely many primes are such. Let $A$ be a place above $w$ and $\phi\in D_A(E)$ a Frobenius element. The images $b,b'$ under $\iota$ of the roots of the characteristic polynomial of $r(\phi^{-1})$ satisfy $|b|<\operatorname{N}(w)|b'|$ and $|b'|<\operatorname{N}(w)|b|$, by the bound of Jacquet and Shalika \cite{JS81a}, applied to a unitary twist of $\pi$; the assumed Galois representations provide this bound at such primes. Since $\chi(\varpi_w)=\iota((c_1/c_2)(\phi^{-1}))$ is $b/b'$ or $b'/b$, $\operatorname{N}(w)^{-1}<|\chi(\varpi_w)|<\operatorname{N}(w)$, and so $n=0$.
\item[(d)] \emph{Finite order.} Some power $\chi^N$, with $N\ge1$, is unramified at every prime. Let $w$ and $A$ be as in (a). By class field theory, and as $n=0$, each $(c_1/c_2)(g)$ with $g\in I_A(E)$ is $\iota^{-1}$ of the value of $\chi$ at a unit of $E_w$; so $(c_1/c_2)^N=1$ on $I_A(E)$, that is, $\chi_{\mathrm{cyc}}^{dN}=1$ on $I_A(E)$. But $\chi_{\mathrm{cyc}}$ maps the inertia group of $A$ in $\operatorname{Gal}(\bar{\mathbf Q}/\mathbf Q)$ onto $\mathbf Z_\ell^\times$, as $\mathbf Q_\ell(\mu_{\ell^j})$ is totally ramified over $\mathbf Q_\ell$ with group $(\mathbf Z/\ell^j)^\times$ for every $j$ \cite[Chapter~I, \S10]{Neu99}. So for $u\in\mathbf Z_\ell^\times$ of infinite order there is $\sigma$ in that group with $\chi_{\mathrm{cyc}}(\sigma)=u$; some $\sigma^a$ with $a\ge1$ lies in $G_E$, hence in $I_A(E)$, and $\chi_{\mathrm{cyc}}(\sigma^a)^{dN}=u^{adN}\ne1$. This contradiction proves (1).
\end{enumerate}

For (4), we show that $\pi$ itself has the property. It is regular algebraic, and $r$ is attached to it through $\iota$, by conditions (a) and (b) of Definition~\ref{sha-siegel-step-seeds-induced-hilbert-form-ordinary}. Let $\omega$ be the central character of $\pi$, a Hecke character of $E$ with $\phi(zg)=\omega(z)\phi(g)$ for $z\in\mathbf A_E^\times$, $g\in GL_2(\mathbf A_E)$ and $\phi$ in the space of $\pi$. Let $w$ be a prime of $E$, $\mathcal O_w$ the ring of integers of $E_w$, $\varpi_w$ a uniformizer, $q=\operatorname N(w)$, and $\pi_w$ the action of $GL_2(E_w)$ on the space of $\pi$.
\begin{enumerate}
\item[(i)] \emph{Primes of level zero.} If $c(w)=0$, then $\pi_w$ has a nonzero vector fixed by $GL_2(\mathcal O_w)$ by condition (d), and $\pi$ has Satake parameters at $w$ by Lemma~\ref{sha-siegel-step-seeds-induced-hilbert-form-galois-component}(1).
\item[(ii)] \emph{The eigenvalue.} Suppose that $c(w)\ne0$. Then there is $\lambda\ne0$ such that $\pi_w$ has a nonzero vector fixed by $K_1(w^{c(w)})$ on which $U_w$ acts by $\lambda$, with $|\iota^{-1}(\lambda)q^{m_0(w)}|=1$ if $w$ lies above $\ell$. Indeed, if $w$ lies above $\ell$, condition (f) gives $c(w)=1$ and such a $\lambda$ with $|\iota^{-1}(\lambda)q^{m_0(w)}-a(w)|<1$; as $|a(w)|=1$, the ultrametric inequality gives $|\iota^{-1}(\lambda)q^{m_0(w)}|=1$, so $\lambda\ne0$. Otherwise condition (e) gives such a $\lambda$ with $|\iota^{-1}(\lambda)|=1$.
\item[(iii)] \emph{A principal series.} Let $c(w)\ne0$ and $\lambda$ be as in (ii), and let $\sigma$ be a local component of $\pi$ at $w$. By Lemma~\ref{sha-siegel-step-seeds-induced-hilbert-form-galois-component}(2), $\sigma$ has a nonzero vector fixed by $K_1(w^{c(w)})$ on which $U_w$ acts by $\lambda$. As $c(w)\ge1$ and $\lambda\ne0$, Lemma~\ref{sha-siegel-step-seeds-induced-hilbert-form-galois-jacquet} gives smooth characters $\psi_0,\psi_1$ of $E_w^\times$ such that $\sigma$ is a subquotient of $\psi_0\times\psi_1$, $\psi_0$ is trivial on $\mathcal O_w^\times$, and $\lambda^2=q\psi_0(\varpi_w)^2$. The centre $E_w^\times$ acts on $\psi_0\times\psi_1$ by $\psi_0\psi_1$, hence also on its subquotient $\sigma$; by Lemma~\ref{sha-siegel-step-seeds-induced-hilbert-form-galois-component}(3) it acts on $\sigma$ by $\omega_w$. So $\psi_0\psi_1=\omega_w$.
\item[(iv)] \emph{Linked characters.} The supercuspidal support of $\sigma$ is $\{\psi_0,\psi_1\}$ (Definition~\ref{gal-irreducible-classes}), which consists of characters. Suppose that it is not unlinked (Definition~\ref{gal-plain}): $\psi_i=\psi_j\nu$ for some $i,j$, that is, $\psi_i=\psi_j$ on $\mathcal O_w^\times$ and $q\psi_i(\varpi_w)=\psi_j(\varpi_w)$. If $i=j$, then $q=1$, which is false. So $\{i,j\}=\{0,1\}$, and $\psi_1$ is trivial on $\mathcal O_w^\times$, as $\psi_0$ is.
\item[(v)] \emph{Primes above $\ell$.} Suppose in (iv) that $w$ lies above $\ell$. Write $t=|q|$, $L=|\iota^{-1}(\lambda)|$, $A=|\iota^{-1}(\psi_0(\varpi_w))|$, $B=|\iota^{-1}(\psi_1(\varpi_w))|$ and $m_i=m_i(w)$. As $\ell$ divides $q>1$, $0<t<1$. By (ii), $Lt^{m_0}=1$, so $L^2t^{2m_0}=1$. By (iii), $L^2=tA^2$. By (iv), $At=B$ or $Bt=A$. By Lemma~\ref{sha-siegel-step-seeds-induced-hilbert-form-galois-central}(1), some power of $\omega_w(\varpi_w)q^{s+1}=\psi_0(\varpi_w)\psi_1(\varpi_w)q^{s+1}$ is $1$, so $ABt^{s+1}=1$. If $B=At$, then $L^2t^{s+1}=A^2t^{s+2}=ABt^{s+1}=1$, so $t^{2m_0}=t^{s+1}$ and $2m_0=s+1$. If $A=Bt$, then $L^2t^{s-1}=A^2t^s=ABt^{s+1}=1$, so $2m_0=s-1$. As $m_0+m_1=s$ by condition (f), $m_0=m_1\pm1$ in both cases, against $m_1+2\le m_0$.
\item[(vi)] \emph{Primes not above $\ell$.} Suppose in (iv) that $w$ does not lie above $\ell$. Then $\omega_w=\psi_0\psi_1$ is trivial on $\mathcal O_w^\times$, so, by Lemma~\ref{sha-siegel-step-seeds-induced-hilbert-form-galois-central}(2), $\xi$ is trivial on $I_A(E)$ for every place $A$ above $w$. By condition (c), $c(w)=0$ or $w$ lies above $\ell$, which is false.
\end{enumerate}
By (iv) to (vi), when $c(w)\ne0$ the supercuspidal support of every local component of $\pi$ at $w$ consists of characters and is unlinked, so these local components are plain. With (i), this proves (4).
% (4) is PrOrd, proved by prOrd (PrOrd.lean, 11 lines) through ordPlain (OrdPlain.lean, 151 lines), followed step by step; omega is Af.Rn.centChar. (i) is the case hc with h.sph and the first part of locTrans (Lemma sha-siegel-step-seeds-induced-hilbert-form-galois-component(1)). (ii) is hlam, with h.above and the private norm_of_near, and h.away. (iii) is the second and third parts of locTrans, uLoc and hprod with psCent (PsCent.lean, 38 lines: parInd_scalar and scalar_of_isSubquotient). (iv) is psPlain (PsPlain.lean, 57 lines: IwahoriSupport.hasSupport_chars for the support, and the translation of psi_i = psi_j nu through norm_nabs_valuation), the case i = j and hu01. (v) is the private real_core and norm_q, with hN of detCent, hsq and hq. (vi) is hI of detCent, hram and h.ram. About 2,100 lines with the dependencies, almost all shared with (2) and (3).
\end{proof}
% (2) and (3) are RR.ordAway_a_alt (RR.ordAway_of_alt) and RR.ordAbove_of_printsH (RR.ordAbove_of_alt), from Gal.SV.exRS_alt with local_comp_real_of_printed_results, jacqOfU_alt and isSubquotient_principalSeries_of_borel_functional; about 6,400 and 26,400 lines with their dependencies, overlapping. (1) is ordinary_irreducible_of_printed_hida (SVROrdIrr.lean): chars_of_red; (a) is key, hp and hom_or; (b) is Af.algOfCyc_s; (c) is Gal.SV.exRS_alt, Gal.SV.ratioW_of_hasSat_alt, GalN.V.ratio_roots and GalN.V.type_zero_window (HeckeNormPow.norm_val_eq_rpow, heckeChar_parallel_weight); (d) is Corev.UnrPow.exists_pow_unrAt, cyc_ine and the field ine_image of the l-adic character. About 900 lines.

\begin{lemma}\label{sha-siegel-step-seeds-induced-hilbert-form-twist}
Let $\ell$ be a prime, $\iota:\bar{\mathbf Q}_\ell\cong\mathbf C$, $E\subseteq\bar{\mathbf Q}$ a totally real number field with $G_E=\operatorname{Gal}(\bar{\mathbf Q}/E)$, $n\ge1$ and $k\ge0$. Let $r:G_E\to GL_n(\bar{\mathbf Q}_\ell)$ be a homomorphism and $r'(g)=\chi_{\mathrm{cyc}}(g)^kr(g)$.
\begin{enumerate}
\item If $r$ is attached through $\iota$ to a cuspidal automorphic representation $\pi$ of $GL_n(\mathbf A_E)$, then $r'$ is attached through $\iota$ to the twist $\pi'$ of $\pi$ by $|\det|^k$, and whenever $\pi$ has the infinitesimal character of weight $(m_1,\dots,m_n)$ at an infinite place of $E$, $\pi'$ has that of weight $(m_1+k,\dots,m_n+k)$ there.
\item Assume the Galois representations as in Theorem~\ref{sha-residually-reducible-automorphy}. If $r$ is attached through $\iota$ to a regular algebraic cuspidal automorphic representation of $GL_n(\mathbf A_E)$ that has, at every finite place of $E$, Satake parameters or only plain local components (Definition~\ref{gal-plain}), then so is $r'$.
\end{enumerate}
\end{lemma}
\leanref{Development/AutomorphicArithmetic/AutomorphicVarieties/ModularityArgument/BaseChangeUnits/BaseChangeNorms/MixSeed.lean}{attR_tw_cyc}
\leanref{Development/AutomorphicArithmetic/AutomorphicVarieties/ModularityArgument/AutomorphyLifting/PrimeIndexConstruction/MixSeedPAlt.lean}{UpperTheory.V.Z.Win.Birth.PrTw}
\leanref{Development/AutomorphicArithmetic/AutomorphicVarieties/ModularityArgument/AutomorphyLifting/PrimeIndexConstruction/PlainAtRTwist.lean}{UpperTheory.V.Z.Win.Birth.prTw}
\leanref{Development/AutomorphicArithmetic/AutomorphicVarieties/ModularityArgument/ThresholdArguments/DenseThresholdCases/TAlgCyc.lean}{UpperTheory.V.Sd.alg_cyclotomic_of_printed_results}
\begin{proof}
(1) As $r$ is semisimple, so is $r'$. Let $w$ be a prime not above $\ell$. The character $\chi_{\mathrm{cyc}}^k$ is trivial on the inertia groups at $w$, and its value at the inverse of a Frobenius element is $\operatorname{N}(w)^{-k}$, the value of $|\cdot|^k$ at a uniformizer of $E_w$. So if $\pi_w$ has Satake parameters that correspond through $\iota$ to the eigenvalues of $r$ at the Frobenius elements, then $\pi'_w$ has these Satake parameters multiplied by $\operatorname{N}(w)^{-k}$, and they correspond in the same way to the eigenvalues of $r'$. All but finitely many primes are such. Twisting by $|\det|^k$ adds $k$ to each entry of the infinitesimal character at every infinite place.

(2) The character $\chi_{\mathrm{cyc}}^k$ is the $\ell$-adic character of an algebraic Hecke character $\chi$ of $E$, through the Artin map of global class field theory. Let $\pi$ be a representation as in the hypothesis. Then $r'$ is attached through $\iota$ to $\pi\otimes(\chi\circ\det)$, which is regular algebraic, as its infinitesimal characters are those of $\pi$ shifted by the infinity type of $\chi$. Let $v$ be a finite place of $E$. The twist of a plain representation by a character is plain. Suppose that $\pi$ has Satake parameters at $v$. If $\chi$ is unramified at $v$, the twist has Satake parameters at $v$. If $\chi$ is ramified at $v$, the local components of $\pi$ at $v$ are generic and unramified, hence plain, and so are their twists.
\end{proof}
% (1) is the private attR_tw_cyc of MixSeed.lean (BaseChangeNorms): pi' = pi.twist (Model.absPow F k), the semisimplicity through galoisRep_irreducibility_criterion, the Frobenius eigenvalues through charEig_cyc_pow and frobenius_eigenvalues_twist, and the infinitesimal characters through HasInfChar.twist and expo_abs. (2) is prTw with Sd.alg_cyclotomic_of_printed_results for cycBar^k: expo_intR, sys_twist, plainAtR_twist and plainAtR_of_hasSat. About 400 lines with their dependencies.

\begin{lemma}\label{sha-siegel-step-seeds-induced-hilbert-form-ordinary-twist}
Let $\ell$ be a prime, write $x\equiv y$ for $|x-y|<1$ in $\bar{\mathbf Q}_\ell$, and let $E\subseteq\bar{\mathbf Q}$ be a number field; write $G_E=\operatorname{Gal}(\bar{\mathbf Q}/E)$, and $D_A(E)\supseteq I_A(E)$ for the decomposition and inertia groups of $G_E$ at a place $A$ of $\bar{\mathbf Q}$. Let $r:G_E\to GL_2(\bar{\mathbf Q}_\ell)$ be a continuous homomorphism, $s$ an integer divisible by $\ell-1$, and $\psi_0,\psi_1:G_E\to\bar{\mathbf Q}_\ell^\times$ homomorphisms with values of absolute value $1$, with $\psi_0$ continuous and $\det r=\chi_{\mathrm{cyc}}^s\psi_0\psi_1$. Let $A$ be a place of $\bar{\mathbf Q}$ above a prime $w$ of $E$ above $\ell$, with $\psi_0$ trivial on $I_A(E)$, and let $(m_0,m_1)\in\mathbf Z^2$ with $\ell-1$ dividing $m_0$. Let $\kappa_0,\kappa_1$ be characters of $D_A(E)$ such that $r$ is conjugate on $D_A(E)$ to an upper triangular representation with diagonal entries $\kappa_0,\kappa_1$, $\kappa_i=\chi_{\mathrm{cyc}}^{m_i}$ on $I_A(E)$, and $\kappa_0(\sigma^{-1})\chi_{\mathrm{cyc}}(\sigma)^{m_0}\equiv\psi_0(\sigma^{-1})$ for every Frobenius element $\sigma\in D_A(E)$. Let $k\ge0$ be divisible by $\ell-1$. Then $g\mapsto\chi_{\mathrm{cyc}}(g)^kr(g)$ is ordinary at $A$ of weights $(m_0+k,m_1+k)$ with residues $(\psi_0,\psi_1)$ (Definition~\ref{sha-siegel-step-seeds-induced-datum}(1)).
\end{lemma}
\leanref{Development/AutomorphicArithmetic/AutomorphicVarieties/ModularityArgument/BaseChangeUnits/BaseChangeNorms/MixSeed.lean}{ordAt_tw}
\begin{proof}
Let $Q\in GL_2(\bar{\mathbf Q}_\ell)$ conjugate $r$ on $D_A(E)$ to the upper triangular representation. Then $Q\chi_{\mathrm{cyc}}^krQ^{-1}$ is upper triangular on $D_A(E)$ with diagonal entries $\kappa_i\chi_{\mathrm{cyc}}^k$, equal to $\chi_{\mathrm{cyc}}^{m_i+k}$ on $I_A(E)$. It remains to show that $\kappa_0\chi_{\mathrm{cyc}}^k\equiv\psi_0$ and $\kappa_1\chi_{\mathrm{cyc}}^k\equiv\psi_1$ on $D_A(E)$.

By the ultrametric inequality, the elements $x\in\bar{\mathbf Q}_\ell^\times$ with $x\equiv1$ form a subgroup $U$, and for $|y|=1$, $x\equiv y$ if and only if $xy^{-1}\in U$. For $g\in G_E$ and $j$ divisible by $\ell-1$, $\chi_{\mathrm{cyc}}(g)^j\in U$, as $\chi_{\mathrm{cyc}}(g)\in\mathbf Z_\ell^\times$ and $u^{\ell-1}\equiv1\pmod\ell$ for $u\in\mathbf Z_\ell^\times$.

Let $e=\kappa_0\psi_0^{-1}$, a homomorphism on $D_A(E)$. It is continuous, as $\kappa_0$ is the entry $(1,1)$ of $QrQ^{-1}$ and $\psi_0$ is continuous. For $g\in I_A(E)$, $e(g)=\chi_{\mathrm{cyc}}(g)^{m_0}\in U$. For a Frobenius element $\sigma$, the hypothesis gives $\kappa_0(\sigma^{-1})\chi_{\mathrm{cyc}}(\sigma)^{m_0}\psi_0(\sigma)\in U$, and its inverse times $\chi_{\mathrm{cyc}}(\sigma)^{m_0}$ is $e(\sigma)\in U$. So $e^{-1}(U)$ is an open subgroup of $D_A(E)$, hence closed, which contains $I_A(E)$ and a Frobenius element. As $D_A(E)$ is the closure of the subgroup generated by $I_A(E)$ and any Frobenius element \cite[Chapter~II, \S9]{Neu99}, $e(D_A(E))\subseteq U$.

For $g\in D_A(E)$, $\kappa_0(g)\chi_{\mathrm{cyc}}(g)^k\psi_0(g)^{-1}=e(g)\chi_{\mathrm{cyc}}(g)^k\in U$, so $\kappa_0\chi_{\mathrm{cyc}}^k\equiv\psi_0$. As $\kappa_0\kappa_1=\det r=\chi_{\mathrm{cyc}}^s\psi_0\psi_1$ on $D_A(E)$, $\kappa_1(g)\chi_{\mathrm{cyc}}(g)^k\psi_1(g)^{-1}=e(g)^{-1}\chi_{\mathrm{cyc}}(g)^{s+k}\in U$, so $\kappa_1\chi_{\mathrm{cyc}}^k\equiv\psi_1$.
\end{proof}
% ordAt_tw (MixSeed.lean, 50 lines): shapeFacts (part 3) for the flag of the twist, U1, cong_iff and cyc_mem, e0 through dec_cong (GlobalTheory.TameQuot.closure_frob_sup_Ine and exists_isFrob), and ed for the determinant.

\begin{lemma}\label{sha-siegel-step-seeds-induced-hilbert-form-irreducible-normal}
Let $\ell$ be an odd prime and write $x\equiv y$ for $|x-y|<1$ in $\bar{\mathbf Q}_\ell$. Let $\Gamma$ be a group, $\rho:\Gamma\to GL_2(\bar{\mathbf Q}_\ell)$ an irreducible homomorphism and $\Phi:\Gamma\to\bar{\mathbf Q}_\ell^\times$ a homomorphism with $|\Phi(s)|=1$, such that for every $s\in\Gamma$ the coefficients of the characteristic polynomial of $\rho(s)$ are congruent to those of $(X-1)(X-\Phi(s))$. Let $H$ be a normal subgroup of $\Gamma$ that contains an element $s_0$ such that $\rho(s_0)$ has distinct eigenvalues, and suppose that $\rho$ is reducible on $H$. Then $\Phi(s)\equiv1$ for every $s\in H$, and $\Phi(s)^2\equiv1$ for every $s\in\Gamma$.
\end{lemma}
\leanref{Development/AutomorphicArithmetic/AutomorphicVarieties/IdeleGenerators/HolomorphicIrreducibility/LineOfNotIrr.lean}{UpperTheory.V.Sd.sIabs}
\leanref{Development/AutomorphicArithmetic/AutomorphicVarieties/IdeleGenerators/ShapeFacts.lean}{UpperTheory.V.Sd.SIabs}
\begin{proof}
As $\rho$ is reducible on $H$, there is a nonzero $v\in\bar{\mathbf Q}_\ell^2$ that is an eigenvector of $\rho(h)$ for every $h\in H$. Let $D$ be the set of $s\in\Gamma$ with $\rho(s)v$ a multiple of $v$; it contains $H$. As $\rho$ is irreducible, the line through $v$ is not stable under $\Gamma$, so there is $t_0\notin D$. If $s\in D$ and $t\notin D$, then $ts\notin D$, since $\rho(ts)v=c\rho(t)v$ with $\rho(s)v=cv$ and $c\ne0$.

Let $t\notin D$. Then $\operatorname{tr}\rho(t)=0$. Indeed, $t^{-1}s_0t$ and $t^{-2}s_0t^2$ lie in $H$, as $H$ is normal, so $\rho(t)v$ and $\rho(t)^2v$ are eigenvectors of $\rho(s_0)$, as is $v$. The vectors $v$ and $\rho(t)v$ are independent; as $\rho(s_0)$ has distinct eigenvalues, they span its two eigenlines, and the nonzero vector $\rho(t)^2v$ lies on one of them. If $\rho(t)^2v$ were a multiple of $\rho(t)v$, then $\rho(t)v$ would be a multiple of $v$. So $\rho(t)^2v=\mu v$ for some $\mu$, and in the basis $(v,\rho(t)v)$ the matrix of $\rho(t)$ is $\begin{pmatrix}0&\mu\\1&0\end{pmatrix}$, of trace $0$.

So for $t\notin D$, comparing the coefficients of $X$, $|1+\Phi(t)|<1$, that is $\Phi(t)\equiv-1$. For $s\in D$, $t_0s\notin D$, so $\Phi(t_0)\Phi(s)+1$ and $\Phi(t_0)+1$ have absolute value less than $1$; as $|\Phi(t_0)|=1$ and $\Phi(t_0)(\Phi(s)-1)$ is their difference, $\Phi(s)\equiv1$. Finally $\Phi(s)^2-1=(\Phi(s)+1)(\Phi(s)-1)$ for every $s\in\Gamma$; both factors have absolute value at most $1$, and one of them less than $1$, according as $s\notin D$ or $s\in D$.
\end{proof}
% sIabs (LineOfNotIrr.lean, 118 lines): the private core, with line_of_not_irr for v, moves_line for t_0, and trace_zero (Reformulated.B.Qf25.trace_zero) for the trace; then hneg, hpos and the final factorization.

\begin{lemma}\label{sha-siegel-step-seeds-induced-hilbert-form-irreducible}
Let $\ell$ be an odd prime and write $x\equiv y$ for $|x-y|<1$ in $\bar{\mathbf Q}_\ell$. Let $E\subseteq\bar{\mathbf Q}$ be a number field; write $G_E=\operatorname{Gal}(\bar{\mathbf Q}/E)$, and $D_A(E)\supseteq I_A(E)$ for the decomposition and inertia groups of $G_E$ at a place $A$ of $\bar{\mathbf Q}$. Let $r:G_E\to GL_2(\bar{\mathbf Q}_\ell)$ be an irreducible homomorphism and $\xi:G_E\to\bar{\mathbf Q}_\ell^\times$ a homomorphism with $|\xi(g)|=1$ for all $g$, such that for every $s\in G_E$ the coefficients of the characteristic polynomial of $r(s)$ are congruent to those of $(X-1)(X-\xi(s))$, and $\xi(g)^2\not\equiv1$ for some $g\in G_E$. Let $A$ be a place of $\bar{\mathbf Q}$ above $\ell$ and $k\ne0$ an integer, and suppose that $r$ is ordinary at $A$ of weights $(k,0)$ with some residues (Definition~\ref{sha-siegel-step-seeds-induced-datum}(1)). Then $r$ is irreducible on every open subgroup of $G_E$.
\end{lemma}
\leanref{Development/AutomorphicArithmetic/AutomorphicVarieties/ModularityArgument/BaseChangeUnits/BaseChangeNorms/MixSeed.lean}{irrOpen_of}
\begin{proof}
Let $H$ be an open subgroup of $G_E$. As $H$ has finite index in the compact group $G_E$, the intersection $N$ of its conjugates is an open normal subgroup of $G_E$ contained in $H$. It suffices to show that $r$ is irreducible on $N$.

There is $g\in N\cap I_A(E)$ with $\chi_{\mathrm{cyc}}(g)^k\ne1$. Indeed, let $u\in\mathbf Z_\ell^\times$ be of infinite order. As $\mathbf Q_\ell(\mu_{\ell^n})$ is totally ramified over $\mathbf Q_\ell$ with group $(\mathbf Z/\ell^n)^\times$ for every $n$ \cite[Chapter~I, \S10]{Neu99}, $\chi_{\mathrm{cyc}}$ maps the inertia group of $A$ in $\operatorname{Gal}(\bar{\mathbf Q}/\mathbf Q)$, which is compact, onto $\mathbf Z_\ell^\times$, so some $\sigma$ in it has $\chi_{\mathrm{cyc}}(\sigma)=u$. As $G_E$ and $N$ have finite index, there are $a,b\ge1$ with $\sigma^a\in G_E$, hence $\sigma^a\in I_A(E)$, and $\sigma^{ab}\in N$. Then $g=\sigma^{ab}$ has $\chi_{\mathrm{cyc}}(g)^k=u^{abk}\ne1$.

As $r$ is ordinary at $A$ of weights $(k,0)$, $r(g)$ is conjugate to an upper triangular matrix with diagonal entries $\chi_{\mathrm{cyc}}(g)^k$ and $1$, which are distinct; so $r(g)$ has distinct eigenvalues. If $r$ were reducible on $N$, Lemma~\ref{sha-siegel-step-seeds-induced-hilbert-form-irreducible-normal}, for $\Gamma=G_E$, $\rho=r$, $\Phi=\xi$, the normal subgroup $N$ and $s_0=g$, would give $\xi(g')^2\equiv1$ for every $g'\in G_E$, against the hypothesis.
\end{proof}
% irrOpen_of (MixSeed.lean, 27 lines): H.normalCore (open by Subgroup.isOpen_of_isClosed_of_finiteIndex), exists_ine (GlobalTheory.CycDec.exists_unit_infinite_order, GlobalTheory.CycInertia.cycZ_Ine_surjective and Subgroup.exists_pow_mem_of_index_ne_zero, 40 lines), the triangular form with Matrix.charpoly_of_isUpperTriangular for the separable characteristic polynomial, and sIabs.

\begin{lemma}\label{sha-siegel-step-seeds-induced-hilbert-form}
Assume the Galois representations as in Theorem~\ref{sha-residually-reducible-automorphy}. Let $\ell$ be an odd prime, $\iota:\bar{\mathbf Q}_\ell\cong\mathbf C$, and write $x\equiv y$ for $|x-y|<1$ in $\bar{\mathbf Q}_\ell$. Let $E\subseteq\bar{\mathbf Q}$ be a totally real number field of even degree such that every prime $w$ of $E$ above $\ell$ has norm $\ell$ and $\ell\notin w^2$. Write $G_E=\operatorname{Gal}(\bar{\mathbf Q}/E)$, and $D_w(E)$ and $I_w(E)$ for the decomposition and inertia groups of $G_E$ at a place $w$ of $\bar{\mathbf Q}$. Let $\xi:G_E\to\bar{\mathbf Q}_\ell^\times$ be a homomorphism with open kernel such that:
\begin{enumerate}
\item $\xi(c)=-1$ for every complex conjugation $c\in G_E$;
\item $\xi$ is trivial on $I_w(E)$ for every place $w$ of $\bar{\mathbf Q}$ above $\ell$;
\item for some prime $u$ of $E$ of norm $2$ and some place $w$ of $\bar{\mathbf Q}$ above $u$, $\xi$ is nontrivial on $I_w(E)$;
\item $\xi(g)^2\not\equiv1$ for some $g\in G_E$;
\item for every place $w$ of $\bar{\mathbf Q}$ above $\ell$, $\xi(\sigma)\not\equiv1$ for some $\sigma\in D_w(E)$.
\end{enumerate}
Let $P$ be a set of primes of $E$ that contains some prime above $\ell$ and does not contain some other prime above $\ell$. Then there are integers $k\ge2$ and $t\ge2$, both divisible by $\ell-1$, a cuspidal automorphic representation $\pi$ of $GL_2(\mathbf A_E)$, and a continuous homomorphism $r:G_E\to GL_2(\bar{\mathbf Q}_\ell)$, such that: $r$ is irreducible on every open subgroup of $G_E$; $r$ is attached to $\pi$ through $\iota$; $\det r=\chi_{\mathrm{cyc}}^k\xi$; for $s\in G_E$ the coefficients of the characteristic polynomial of $r(s)$ are congruent to those of $(X-1)(X-\xi(s))$; for every place $w$ of $\bar{\mathbf Q}$ not above $\ell$, $r$ is conjugate on $I_w(E)$ to the diagonal representation with entries $1$ and $\xi$; $r$ is attached through $\iota$ to a regular algebraic cuspidal automorphic representation of $GL_2(\mathbf A_E)$ that has, at every finite place of $E$, Satake parameters or only plain local components (Definition~\ref{gal-plain}); and for every prime $w_0$ of $E$ above $\ell$:
\begin{enumerate}
\item[(a)] if $w_0\in P$, then $r$ is ordinary of weights $(k+t,-t)$ with residues $(\xi,1)$ (Definition~\ref{sha-siegel-step-seeds-induced-datum}(1)) at every place of $\bar{\mathbf Q}$ above $w_0$, and $\pi$ has the infinitesimal character of weight $(k+t,-t)$ at the infinite place of $E$ induced by $\iota\circ\tau$, for every embedding $\tau:E\to\bar{\mathbf Q}_\ell$ that induces $w_0$;
\item[(b)] if $w_0\notin P$, the same holds with the weights $(k,0)$ and the residues $(1,\xi)$.
\end{enumerate}
\end{lemma}
\leanref{Development/AutomorphicArithmetic/AutomorphicVarieties/ModularityArgument/AutomorphyLifting/PrimeIndexConstruction/MixSeedPAlt.lean}{UpperTheory.V.Z.Win.Birth.MixSeedP}
\leanref{Development/AutomorphicArithmetic/AutomorphicVarieties/ModularityArgument/ThresholdArguments/DenseThresholdCases/MixSeedOf.lean}{UpperTheory.Esk.mixSeed_of_alt}
\leanref{Development/AutomorphicArithmetic/AutomorphicVarieties/ModularityArgument/ThresholdArguments/DenseThresholdCases/MPropMHol.lean}{UpperTheory.V.Z.Bc.propMHolZ_alt}
\leanref{Development/AutomorphicArithmetic/AutomorphicVarieties/ModularityArgument/ThresholdArguments/DenseThresholdCases/BirthMix.lean}{UpperTheory.V.Z.Win.Birth.mix_birth_of_prints}
\begin{proof}
Let $w_1\in P$ and $w_2\notin P$ be primes above $\ell$, and $A_1$ and $A_2$ places of $\bar{\mathbf Q}$ above them.
\begin{enumerate}
\item[(1)] \emph{Values of $\xi$.} As $\ker\xi$ is open, $G_E/\ker\xi$ is finite, so each $\xi(g)$ is a root of unity and $|\xi(g)|=1$. As $[E:\mathbf Q]$ is even and positive, $[E:\mathbf Q]\ge2$.
\item[(2)] \emph{The form.} By condition (5) on $\xi$, for $i=1,2$ there is $\sigma_i\in D_{A_i}(E)$ with $\xi(\sigma_i)\not\equiv1$. With conditions (1), (2) and (3) on $\xi$ and step (1), Lemma~\ref{sha-siegel-step-seeds-induced-hilbert-form-exists} applies to $E$, $\xi$ and $P$. It gives integers $k\ge2$ and $t\ge2$, both divisible by $\ell-1$, and an ordinary form datum $(c,m,a,\pi_1,r_1)$ for $\xi$ of weight $-k$ with the properties (1) to (3) of that lemma.
\item[(3)] \emph{The twist.} Let $r(g)=\chi_{\mathrm{cyc}}(g)^kr_1(g)$, continuous as $r_1$ and $\chi_{\mathrm{cyc}}$ are. By Lemma~\ref{sha-siegel-step-seeds-induced-hilbert-form-twist}(1) there is a cuspidal automorphic representation $\pi$ of $GL_2(\mathbf A_E)$ to which $r$ is attached through $\iota$, with the infinitesimal character of weight $(n_0+k,n_1+k)$ at every infinite place at which $\pi_1$ has that of weight $(n_0,n_1)$. As $\det r_1=\chi_{\mathrm{cyc}}^{-k}\xi$, $\det r=\chi_{\mathrm{cyc}}^{2k}\det r_1=\chi_{\mathrm{cyc}}^k\xi$.
\item[(4)] \emph{Residues.} Let $g\in G_E$ and $\delta=\chi_{\mathrm{cyc}}(g)^k$. As $\ell-1$ divides $k$, $\delta\equiv1$, as in the proof of Lemma~\ref{sha-siegel-step-seeds-induced-hilbert-form-ordinary-twist}; in particular $|\delta|=1$. The characteristic polynomial of $r(g)=\delta r_1(g)$ arises from that of $r_1(g)$ by multiplying the coefficient of $X^{2-i}$ by $\delta^i$, and $(X-\delta)(X-\delta\xi(g))$ arises from $(X-1)(X-\xi(g))$ in the same way. So, by Lemma~\ref{sha-siegel-step-seeds-induced-hilbert-form-exists}(2), the coefficients of the characteristic polynomial of $r(g)$ are congruent to those of $(X-\delta)(X-\delta\xi(g))$. As $\delta\equiv1$ and $|\xi(g)|=1$, these are congruent to those of $(X-1)(X-\xi(g))$.
\item[(5)] \emph{The primes above $\ell$.} Let $w$ be a prime above $\ell$ and $A$ a place above $w$. Lemma~\ref{sha-siegel-step-seeds-induced-hilbert-form-galois}(3) gives characters $\kappa_0,\kappa_1$ of $D_A(E)$ such that $r_1$ is conjugate on $D_A(E)$ to an upper triangular representation with diagonal entries $\kappa_0,\kappa_1$, $\kappa_i=\chi_{\mathrm{cyc}}^{m_i(w)}$ on $I_A(E)$, and $\kappa_0(\sigma^{-1})\chi_{\mathrm{cyc}}(\sigma)^{m_0(w)}\equiv a(w)$ for every Frobenius element $\sigma\in D_A(E)$. If $w\in P$, let $(\psi_0,\psi_1)=(\xi,1)$: then $\psi_0$ is trivial on $I_A(E)$ by condition (2) on $\xi$, $m_0(w)=t$, and $a(w)=\xi(\sigma)^{-1}=\psi_0(\sigma^{-1})$ by Lemma~\ref{sha-siegel-step-seeds-induced-hilbert-form-exists}(3). If $w\notin P$, let $(\psi_0,\psi_1)=(1,\xi)$: then $m_0(w)=0$ and $a(w)=1=\psi_0(\sigma^{-1})$. In both cases $\det r_1=\chi_{\mathrm{cyc}}^{-k}\psi_0\psi_1$, $\psi_0$ is continuous, as $\ker\xi$ is open, and $\psi_0$ and $\psi_1$ take values of absolute value $1$ by step (1). Lemma~\ref{sha-siegel-step-seeds-induced-hilbert-form-ordinary-twist}, with $s=-k$, shows that $r$ is ordinary at $A$ of weights $(t+k,-t)$ with residues $(\xi,1)$ if $w\in P$, and of weights $(k,0)$ with residues $(1,\xi)$ if $w\notin P$.
\item[(6)] \emph{Irreducibility.} By Lemma~\ref{sha-siegel-step-seeds-induced-hilbert-form-galois}(1), $r_1$ is irreducible, hence so is $r$, which has the same invariant subspaces. By step (5), $r$ is ordinary at $A_2$ of weights $(k,0)$, with $k\ne0$. With steps (1) and (4) and condition (4) on $\xi$, Lemma~\ref{sha-siegel-step-seeds-induced-hilbert-form-irreducible} shows that $r$ is irreducible on every open subgroup of $G_E$.
\item[(7)] \emph{Inertia away from $\ell$.} Let $A$ be a place of $\bar{\mathbf Q}$ not above $\ell$. By Lemma~\ref{sha-siegel-step-seeds-induced-hilbert-form-galois}(2), $r_1$ is conjugate on $I_A(E)$ to the diagonal representation with entries $1$ and $\xi$. As $\chi_{\mathrm{cyc}}$ is trivial on $I_A(E)$, $r=r_1$ on $I_A(E)$, and the same holds for $r$.
\item[(8)] \emph{Local components.} By Lemma~\ref{sha-siegel-step-seeds-induced-hilbert-form-galois}(4) and Lemma~\ref{sha-siegel-step-seeds-induced-hilbert-form-twist}(2), $r$ is attached through $\iota$ to a regular algebraic cuspidal automorphic representation of $GL_2(\mathbf A_E)$ that has, at every finite place, Satake parameters or only plain local components.
\item[(9)] \emph{Infinite places.} Let $w_0$ be a prime above $\ell$ and $\tau:E\to\bar{\mathbf Q}_\ell$ an embedding that induces $w_0$. By condition (f) of Definition~\ref{sha-siegel-step-seeds-induced-hilbert-form-ordinary}, $\pi_1$ has the infinitesimal character of weight $m(w_0)$ at the infinite place induced by $\iota\circ\tau$. By step (3), $\pi$ has there the infinitesimal character of weight $m(w_0)+(k,k)$, which is $(k+t,-t)$ if $w_0\in P$ and $(k,0)$ if $w_0\notin P$.
\end{enumerate}
By steps (3) to (9), $k$, $t$, $\pi$ and $r$ have the properties of the statement.
\end{proof}
% mix_birth_of_prints (BirthMix.lean, 106 lines), from HypothesesH' and Esk.PropMV: (1) is hn and hF2; (2) is hM (Lemma sha-siegel-step-seeds-induced-hilbert-form-exists) with hdec for the sigma_i; (3) is attR_tw_cyc, cont_units and the determinant computation; (4) is shapeFacts (its parts on ResSum) with cyc_mem; (5) is ordAbove_of_printsH (Lemma sha-siegel-step-seeds-induced-hilbert-form-galois(3)) and ordAt_tw, as hO₁ and hO₂; (6) is ordinary_irreducible_of_printed_hida, the part of shapeFacts on IsIrr, and irrOpen_of; (7) is ordAway_a_alt with the parts of shapeFacts on IsDiagOn and Base.places; (8) is prOrd, prTw and alg_cyclotomic_of_printed_results; (9) is the field above of OrdForm with attR_tw_cyc. About 570 lines with the private lemmas of MixSeed.lean, ShapeFacts.lean and LineOfNotIrr.lean, not counting the inputs of Lemmas sha-siegel-step-seeds-induced-hilbert-form-exists and sha-siegel-step-seeds-induced-hilbert-form-galois. MixSeedP is Esk.mixSeed_of_alt (MixSeedOf.lean), that is mixSeed_of_holWC_alt holWC, which applies mix_birth_of_prints to Esk.pf_PropMV_sv_alt; its input Sd.PropMHol is Z.Bc.propMHolZ_alt from Hypotheses4Z' (HypothesesH' with Published.hilbert_eisenstein_weight_one_holds and Lit.oneUse).

\subsection{Seeds induced from Hilbert modular forms}\label{subsection-sha-induced-seeds}

\begin{lemma}\label{sha-siegel-step-seeds-induced-hilbert}
Assume the Galois representations as in Theorem~\ref{sha-residually-reducible-automorphy}. Let $D$, $T$ and $\iota$ be as in Definition~\ref{sha-siegel-step-seeds-induced-datum}, and let $k_0$ and $\varphi$ satisfy (a) there. Then there are integers $j$ and $t$ and a homomorphism $\rho:G_{k_0}\to GL_2(\bar{\mathbf Q}_\ell)$ that satisfy (b) there, such that $\rho$ is attached through $\iota$ to a regular algebraic cuspidal automorphic representation of $GL_2(\mathbf A_{k_0})$ that has, at every finite place of $k_0$, Satake parameters or only plain local components (Definition~\ref{gal-plain}).
\end{lemma}
\leanref{Development/AutomorphicArithmetic/AutomorphicVarieties/ModularityArgument/AutomorphyLifting/PrimeIndexConstruction/MixSeedPAlt.lean}{UpperTheory.V.Z.Win.Birth.MixOfP}
\leanref{Development/AutomorphicArithmetic/AutomorphicVarieties/ModularityArgument/AutomorphyLifting/PrimeIndexConstruction/MixSeedPAlt.lean}{UpperTheory.V.Z.Win.Birth.PR}
\leanref{Development/AutomorphicArithmetic/AutomorphicVarieties/ModularityArgument/AutomorphyLifting/PrimeIndexConstruction/ContRes.lean}{UpperTheory.V.Z.Win.Birth.mixOfP}
\leanref{Development/AutomorphicArithmetic/AutomorphicVarieties/BuildingsResidueChars/QuadraticFieldSplitPrimeTwoPlaces.lean}{UpperTheory.quadratic_field_split_prime_two_places}
\leanref{Development/AutomorphicArithmetic/AutomorphicVarieties/BuildingsResidueChars/ExistsRingHomToLocalFieldInducingPlace.lean}{UpperTheory.exists_ringHom_to_localField_inducing_place}
\leanref{Development/AutomorphicArithmetic/AutomorphicVarieties/BuildingsResidueChars/ComplexConjugationsAreConjugate.lean}{UpperTheory.complex_conjugations_are_conjugate}
\begin{proof}
\begin{enumerate}
\item[(1)] \emph{Primes of $k_0$.} Let $p\in\{2,\ell\}$, $w$ a place of $\bar{\mathbf Q}$ above $p$ and $g\in G\setminus G_{k_0}$. The primes of $k_0$ below $w$ and below $g^{-1}w$ are distinct: otherwise $hg^{-1}w=w$ for some $h\in G_{k_0}$, which acts transitively on the places above a prime of $k_0$; then $hg^{-1}\in D_w(\mathbf Q)\subseteq G_{k_0}$ by (a), and $g\in G_{k_0}$. So $p$ has two distinct prime factors in $k_0$, which has degree $2$; comparing norms, $(p)$ is their product and each has norm $p$. For every such $w$ and $g$, one of them lies below $w$ and the other below $g^{-1}w$. Let $w_1\ne w_2$ be the primes of $k_0$ above $\ell$, and $u$ a prime of $k_0$ above $2$, with a place $w_u$ of $\bar{\mathbf Q}$ above it.
\item[(2)] \emph{The hypotheses on $\xi$.} Let $P=\{w_1\}$, and let $\xi_0$ be the restriction of $\xi$ to $G_{k_0}$; its kernel is open, as that of $\xi$ is. Every complex conjugation in $G_{k_0}$ is conjugate in $G$ to the complex conjugation $c$ of $D$, as any two complex conjugations of $G$ are; so $\xi_0$ takes the value $\xi(c)=-1$ at it. By condition (q), $\xi$ is trivial on $I_w(\mathbf Q)$ for $w$ above $\ell$, which is not above $2$. By condition (r), some $g\in I_{w_u}(F)$ has $\bar\chi^{(1,0)}(g)^4\ne1$, so $\xi(g)^4\not\equiv1$ by Lemma~\ref{sha-siegel-step-seeds-induced-characters}(5); by (a), $g\in D_{w_u}(\mathbf Q)\subseteq G_{k_0}$, so $g\in I_{w_u}(k_0)$, $\xi_0(g)\ne1$, and $\xi_0(g^2)^2\not\equiv1$. For a place $w$ above $\ell$, (r) gives $x\in D_w(F)$ with $\xi(x)^4\not\equiv1$, and $x\in G_{k_0}$ by (a); so $\sigma=x^4\in D_w(k_0)$ has $\xi_0(\sigma)\not\equiv1$. With (1), $k_0$, $\xi_0$ and $P$ satisfy the hypotheses of Lemma~\ref{sha-siegel-step-seeds-induced-hilbert-form}.
\item[(3)] \emph{The form.} That lemma gives integers $j$ (its $k$) and $t$, a cuspidal automorphic representation $\pi$ of $GL_2(\mathbf A_{k_0})$ and $\rho$ (its $r$). They have the properties asked in (b) other than the infinitesimal characters of $\pi$ and the condition at the places above $\ell$, and $\rho$ is attached through $\iota$ to a representation as in the claim.
\item[(4)] \emph{Infinite places.} For $i=1,2$ there is an embedding $\tau_i:k_0\to\bar{\mathbf Q}_\ell$ that induces $w_i$, as for every prime of a number field above $\ell$. Let $v_i$ be the infinite place of $k_0$ induced by $\iota\circ\tau_i$. As $k_0$ is real, $v_1=v_2$ would give $\iota\circ\tau_1=\iota\circ\tau_2$, so $\tau_1=\tau_2$ and $w_1=w_2$. So $v_1\ne v_2$, and by (3), $\pi$ has the infinitesimal character of weight $(j+t,-t)$ at $v_1$ and of weight $(j,0)$ at $v_2$.
\item[(5)] \emph{The places above $\ell$.} Let $w,w'$ be places above $\ell$ with $w'=g^{-1}w$ for some $g\notin G_{k_0}$. By (1), $w$ lies above $w_1$ and $w'$ above $w_2$, or the reverse. In the first case, by (3), $\rho$ is ordinary at $w$ of weights $(j+t,-t)$ with residues $(\xi,1)$ and at $w'$ of weights $(j,0)$ with residues $(1,\xi)$; in the second, the same holds with $w$ and $w'$ exchanged.
\end{enumerate}
So $j$, $t$ and $\rho$ satisfy (b).
\end{proof}
% mixOfP (ContRes.lean, 95 lines): (1) is quadratic_field_split_prime_two_places (Sd.S2.QuadSplit, 114 lines) for l and for 2, with Hyp.spl and Hyp.sp2; (2) is the list of arguments to MixSeedP: cont_res and Good.xi_open, complex_conjugations_are_conjugate (Sd.S2.ConjCc) with Good.isCC and Good.xi_odd, Two.xi with not_both, Two.big2 with chars.al and lift for u, the square s^2, Two.bigL with lift and y^4, and P = {w}; (4) is exists_ringHom_to_localField_inducing_place (Sd.S2.EmbEx, 115 lines); (5) is the case split on hpl. About 300 lines with these dependencies. MixSeedP is Lemma sha-siegel-step-seeds-induced-hilbert-form.

\begin{lemma}\label{sha-siegel-step-seeds-induced-ribet-lattice}
Let $\ell$ be a prime, and write $x\equiv y$ for $|x-y|<1$ in $\bar{\mathbf Q}_\ell$. Let $\Gamma$ be a compact Hausdorff topological group, $\rho:\Gamma\to GL_2(\bar{\mathbf Q}_\ell)$ a continuous irreducible homomorphism, and $a,b:\Gamma\to\bar{\mathbf Q}_\ell^\times$ homomorphisms with $|a(s)|=|b(s)|=1$ for all $s$, such that for every $s\in\Gamma$ the coefficients of the characteristic polynomial of $\rho(s)$ are congruent to those of $(X-a(s))(X-b(s))$. Let $s_0\in\Gamma$ with $a(s_0)\not\equiv b(s_0)$. Then there is $P\in GL_2(\bar{\mathbf Q}_\ell)$ such that $\rho'=P\rho P^{-1}$ has $\rho'(s_0)$ diagonal, $\rho'(s)_{11}\equiv a(s)$, $\rho'(s)_{22}\equiv b(s)$, $|\rho'(s)_{21}|<1$ and $|\rho'(s)_{12}|\le1$ for every $s\in\Gamma$, and $|\rho'(s)_{12}|=1$ for some $s\in\Gamma$.
\end{lemma}
\leanref{Development/AutomorphicArithmetic/AutomorphicVarieties/IdeleGenerators/HolomorphicIrreducibility/RibetLat.lean}{UpperTheory.V.Sd.ribetLat}
\leanref{Development/AutomorphicArithmetic/AutomorphicVarieties/IdeleGenerators/ShapeFacts.lean}{UpperTheory.V.Sd.RibetLat}
\leanref{Development/AutomorphicArithmetic/IntegralCasts/CapelliIdentities/IrreducibleInvariant.lean}{UpperTheory.Reformulated.B.IrreducibleInvariant.exists_ne}
\begin{proof}
Elements of absolute value at most $1$ form a ring, and $x\equiv y$ is a congruence modulo its maximal ideal.
\begin{enumerate}
\item[(1)] \emph{Traces.} Comparing the coefficients of $X$ and of $1$, $\operatorname{tr}\rho(s)\equiv a(s)+b(s)$ for every $s$, and $\det\rho(s_0)\equiv a(s_0)b(s_0)$.
\item[(2)] \emph{Eigenvalues at $s_0$.} Write $t_0$ and $d_0$ for the trace and the determinant of $\rho(s_0)$, and $a_0=a(s_0)$, $b_0=b(s_0)$. Let $x$ be a root of $X^2-t_0X+d_0$, so that $d_0=x(t_0-x)$. Then $(x-a_0)(t_0-x-a_0)=(d_0-a_0b_0)-a_0(t_0-a_0-b_0)$ has absolute value less than $1$, so $x\equiv a_0$ or $t_0-x\equiv a_0$. So the eigenvalues of $\rho(s_0)$ are $\lambda\equiv a_0$ and $\lambda'=t_0-\lambda$, and $\lambda'-b_0=(t_0-a_0-b_0)-(\lambda-a_0)$, so $\lambda'\equiv b_0$. As $a_0\not\equiv b_0$, $|\lambda-\lambda'|\ge1$; in particular $\lambda\ne\lambda'$, and there is $P_0\in GL_2(\bar{\mathbf Q}_\ell)$ with $P_0\rho(s_0)P_0^{-1}=\operatorname{diag}(\lambda,\lambda')$. Let $\sigma=P_0\rho P_0^{-1}$.
\item[(3)] \emph{Diagonal entries.} Let $s\in\Gamma$. As $\sigma(s_0)$ is diagonal, $\operatorname{tr}\sigma(s_0s)=\lambda\sigma(s)_{11}+\lambda'\sigma(s)_{22}$, and
\[
(\lambda-\lambda')\big(\sigma(s)_{11}-a(s)\big)=\big(\operatorname{tr}\sigma(s_0s)-a_0a(s)-b_0b(s)\big)-\lambda'\big(\operatorname{tr}\sigma(s)-a(s)-b(s)\big)-a(s)(\lambda-a_0)-b(s)(\lambda'-b_0).
\]
By (1) and (2), and as $a$ and $b$ are homomorphisms and $|\lambda'|\le1$, the right side has absolute value less than $1$. So $\sigma(s)_{11}\equiv a(s)$, as $|\lambda-\lambda'|\ge1$, and $\sigma(s)_{22}=\operatorname{tr}\sigma(s)-\sigma(s)_{11}\equiv b(s)$.
\item[(4)] \emph{Products.} For $s,s'\in\Gamma$,
\[
\sigma(s)_{12}\sigma(s')_{21}=\big(\sigma(ss')_{11}-a(s)a(s')\big)-\sigma(s)_{11}\big(\sigma(s')_{11}-a(s')\big)-a(s')\big(\sigma(s)_{11}-a(s)\big),
\]
which has absolute value less than $1$ by (3), as $|\sigma(s)_{11}|\le1$.
\item[(5)] \emph{Scaling.} As $\rho$ is irreducible, $\sigma(s)_{12}\ne0$ for some $s$: otherwise every $\sigma(s)$ is lower triangular, and the line spanned by $P_0^{-1}\binom01$ is stable under $\rho$. As $\Gamma$ is compact and $s\mapsto|\sigma(s)_{12}|$ is continuous, this function attains its maximum at some $s_1$, and $c=\sigma(s_1)_{12}\ne0$. Let $P=\operatorname{diag}(1,c)P_0$ and $\rho'=P\rho P^{-1}$, so that $\rho'(s)$ has entries $\sigma(s)_{11}$, $\sigma(s)_{12}/c$, $c\sigma(s)_{21}$ and $\sigma(s)_{22}$. Then $\rho'(s_0)$ is diagonal; the diagonal entries are as claimed by (3); $|\rho'(s)_{12}|\le1$ by the choice of $s_1$, with equality at $s=s_1$; and $|\rho'(s)_{21}|=|\sigma(s_1)_{12}\sigma(s)_{21}|<1$ by (4).
\end{enumerate}
\end{proof}
% ribetLat (RibetLat.lean, 214 lines): htr and hdet for (1); roots, hne and diag2 for (2); core for (3) to (5), with hA and hD for (3), hBC for (4), and IrreducibleInvariant.exists_ne (about 50 lines), isCompact_univ.exists_isMaxOn and the matrix g = diag(1, B) for (5).

\begin{lemma}\label{sha-siegel-step-seeds-induced-ribet}
Let $D$, $T$ and $\iota$ be as in Definition~\ref{sha-siegel-step-seeds-induced-datum}, let $k_0$ and $\varphi$ satisfy (a) there, and let $j$, $t$ and $\rho_0:G_{k_0}\to GL_2(\bar{\mathbf Q}_\ell)$ satisfy (b) there. Let $u\in G_{k_0}$ with $u=(g_0^T)^2$. Then there is $P\in GL_2(\bar{\mathbf Q}_\ell)$ such that $\rho=P\rho_0P^{-1}$ has $\rho(u)$ diagonal and satisfies (c) there.
\end{lemma}
\leanref{Development/AutomorphicArithmetic/AutomorphicVarieties/BuildingsResidueChars/ExistsRibetLattice.lean}{UpperTheory.exists_ribet_lattice}
\leanref{Development/AutomorphicArithmetic/AutomorphicVarieties/IdeleGenerators/QuadraticSplitting/QuadSplit.lean}{UpperTheory.V.Sd.S2.RibOf}
\begin{proof}
Let $H=G_{k_0}\cap G_F$, an open subgroup of $G_{k_0}$; it is closed, hence compact. As $g_0^T\in G_F$, $u\in H$. By (b), $\rho_0$ is continuous and irreducible on $H$, and for $s\in H$ the coefficients of the characteristic polynomial of $\rho_0(s)$ are congruent to those of $(X-1)(X-\xi(s))$. The characters $1$ and $\xi$ take values of absolute value $1$, by Lemma~\ref{sha-siegel-step-seeds-induced-characters}(1), and $1\not\equiv\xi(u)$, by part (3) of that lemma. Lemma~\ref{sha-siegel-step-seeds-induced-ribet-lattice}, for $\Gamma=H$, the restriction of $\rho_0$, $a=1$, $b=\xi$ and $s_0=u$, gives $P$ such that $\rho=P\rho_0P^{-1}$ has $\rho(u)$ diagonal and, for $s\in H$, $\rho(s)_{11}\equiv1$, $\rho(s)_{22}\equiv\xi(s)$, $|\rho(s)_{21}|<1$ and $|\rho(s)_{12}|\le1$, with equality for some $s\in H$. This is (c).
\end{proof}
% exists_ribet_lattice (ExistsRibetLattice.lean, 42 lines): H is the preimage of the fixing subgroup of F in G_{k0}, open by galEmbTop and compact; huH for u; ribetLat (Lemma sha-siegel-step-seeds-induced-ribet-lattice) with dat.cont, dat.irr, chars.unit, dat.res and chars.sq; and the conjugation by g.

\begin{lemma}\label{sha-siegel-step-seeds-induced-conjugate}
Let $D$, $T$, $\iota$, $k_0$ and $\varphi$ be as in Lemma~\ref{sha-siegel-step-seeds-induced-ribet}, let $j$, $t$ and $\rho_0$ satisfy (b) of Definition~\ref{sha-siegel-step-seeds-induced-datum}, and let $P\in GL_2(\bar{\mathbf Q}_\ell)$. Then $j$, $t$ and $P\rho_0P^{-1}$ satisfy (b). If $\rho_0$ is attached through $\iota$ to a regular algebraic cuspidal automorphic representation of $GL_2(\mathbf A_{k_0})$ with Satake parameters or only plain local components at every finite place, so is $P\rho_0P^{-1}$.
\end{lemma}
\leanref{Development/AutomorphicArithmetic/AutomorphicVarieties/BuildingsResidueChars/GaloisRepPropertiesInvariantUnderConjugation.lean}{UpperTheory.galoisRep_properties_invariant_under_conjugation}
\leanref{Development/AutomorphicArithmetic/AutomorphicVarieties/ModularityArgument/AutomorphyLifting/PrimeIndexConstruction/ConjEquivApply.lean}{UpperTheory.V.Z.Win.Birth.prIso}
% galoisRep_properties_invariant_under_conjugation (about 80 lines) and prIso (about 50 lines).

\begin{lemma}\label{sha-siegel-step-seeds-induced-mix}
Let $k$ be a field in which $2\ne0$, $\Gamma$ a group, $\Gamma_0\subset\Gamma$ a subgroup of index $2$, $\tau\in\Gamma\setminus\Gamma_0$, $\alpha,\beta\in k^\times$, and $\rho:\Gamma_0\to GL_2(k)$ and $\mu:\Gamma\to k^\times$ homomorphisms such that $\rho(\tau^2)=\operatorname{diag}(\alpha^2,\beta^2)$, $\det\rho(h)=\mu(h)$ for $h\in\Gamma_0$, and $\alpha\beta=\mu(\tau)$. Then there is a homomorphism $R:\Gamma\to GL_4(k)$ with $R(h)=M(\beta/\alpha,\rho(h),\rho(\tau^{-1}h\tau))$ for $h\in\Gamma_0$ (Definition~\ref{sha-siegel-step-seeds-induced-datum}(2), over $k$), $R(\tau)=\operatorname{diag}(\beta,-\alpha,-\beta,\alpha)$, and $R(g)$ symplectic with multiplier $\mu(g)$ for every $g\in\Gamma$. If moreover $k=\bar{\mathbf Q}_\ell$, $\Gamma$ is a topological group, $\Gamma_0$ is open and $\rho$ is continuous, then every such $R$ is continuous.
\end{lemma}
\leanref{Development/AutomorphicArithmetic/AutomorphicVarieties/BuildingsResidueChars/SymmetricSquareMixedExists.lean}{UpperTheory.symmetric_square_mixed_exists}
\leanref{Development/AutomorphicArithmetic/AutomorphicVarieties/BuildingsResidueChars/ContinuousMixFamily.lean}{UpperTheory.continuous_mix_family}
\leanref{Development/AutomorphicArithmetic/AutomorphicVarieties/IdeleGenerators/QuadraticSplitting/QuadSplit.lean}{UpperTheory.V.Sd.S2.MixEx}
\leanref{Development/AutomorphicArithmetic/AutomorphicVarieties/IdeleGenerators/QuadraticSplitting/QuadSplit.lean}{UpperTheory.V.Sd.S2.MixCont}
% symmetric_square_mixed_exists (SymmetricSquareMixedExists.lean, 127 lines): R is the representation induced from rho, in the basis in which it takes the form of M; a computation with 4 x 4 matrices. continuous_mix_family (ContinuousMixFamily.lean, 51 lines).

\begin{lemma}\label{sha-siegel-step-seeds-induced-square-root}
Let $\ell$ be an odd prime, and write $x\equiv y$ for $|x-y|<1$ in $\bar{\mathbf Q}_\ell$.
\begin{enumerate}
\item Every $a\in\bar{\mathbf Q}_\ell$ with $a\equiv1$ has exactly one square root $x$ with $x\equiv1$.
\item Let $\Gamma$ be a topological group and $\eta:\Gamma\to\bar{\mathbf Q}_\ell^\times$ a continuous homomorphism with $\eta(g)\equiv1$ for all $g$. Then there is a continuous homomorphism $\delta:\Gamma\to\bar{\mathbf Q}_\ell^\times$ with $\delta^2=\eta$ and $\delta(g)\equiv1$ for all $g$.
\end{enumerate}
\end{lemma}
\leanref{Development/AutomorphicArithmetic/AutomorphicVarieties/BuildingsResidueChars/ExistsSquareRootOfCharacterNearOne.lean}{UpperTheory.exists_square_root_of_character_near_one}
\leanref{Development/AutomorphicArithmetic/AutomorphicVarieties/IdeleGenerators/QuadraticSplitting/QuadSplit.lean}{UpperTheory.V.Sd.S2.SqrtChar}
\begin{proof}
As $\ell$ is odd, $|2|=1$. If $a\equiv1$ and $b\equiv1$, then $a+b=2+(a-1)+(b-1)$ has absolute value $1$, so $|a-b|=|a^2-b^2|$.

(1) Let $x^2=a$. If $x\not\equiv1$, then $|x-1|\ge1$, and $-x\equiv1$, as $(x-1)(-x-1)=1-a$ gives $|-x-1|=|1-a|/|x-1|<1$. Two square roots of $a$ that are $\equiv1$ are equal, by the first paragraph.

(2) Let $\delta(g)$ be the square root of $\eta(g)$ with $\delta(g)\equiv1$, by (1). For $g,h\in\Gamma$, $\delta(g)\delta(h)\equiv1$ is a square root of $\eta(gh)$, so it is $\delta(gh)$; and $\delta$ is continuous, as $|\delta(g)-\delta(h)|=|\eta(g)-\eta(h)|$.
\end{proof}

\begin{lemma}\label{sha-siegel-step-seeds-induced-exists}
Assume the Galois representations as in Theorem~\ref{sha-residually-reducible-automorphy}. Let $D$, $T$ and $\iota$ be as in Definition~\ref{sha-siegel-step-seeds-induced-datum}. Then there is an induction datum $(k_0,\varphi,j,t,\rho,\alpha,\beta,R,\delta)$ for $T$ and $\iota$ such that $\rho$ is attached through $\iota$ to a regular algebraic cuspidal automorphic representation of $GL_2(\mathbf A_{k_0})$ that has, at every finite place of $k_0$, Satake parameters or only plain local components (Definition~\ref{gal-plain}).
\end{lemma}
\leanref{Development/AutomorphicArithmetic/AutomorphicVarieties/IdeleGenerators/QuadraticSplitting/QuadSplit.lean}{UpperTheory.V.Sd.S2.IndEx}
\leanref{Development/AutomorphicArithmetic/AutomorphicVarieties/BuildingsResidueChars/ExistsQuadraticCharacterOfSubgroupIndexTwo.lean}{UpperTheory.exists_quadratic_character_of_subgroup_index_two}
\begin{proof}
\begin{enumerate}
\item[(1)] \emph{The field.} By (p) there is a real quadratic field $k_0$ such that $D_w(\mathbf Q)\subseteq G_{k_0}$ for every place $w$ above $2$ or $\ell$ and, for $g\in G_F$, $\bar\chi^{(0,1)}(g)=1$ if and only if $g\in G_{k_0}$. The subgroup $G_{k_0}$ of $G$ is open of index $2$, so there is a character $\varphi:G\to\{\pm1\}$ with kernel $G_{k_0}$. So $k_0$ and $\varphi$ satisfy (a), and Lemma~\ref{sha-siegel-step-seeds-induced-characters} applies to them.
\item[(2)] \emph{The form.} Lemma~\ref{sha-siegel-step-seeds-induced-hilbert} gives $j$, $t$ and $\rho_0$ that satisfy (b), with $\rho_0$ attached through $\iota$ to a representation as in the claim.
\item[(3)] \emph{The basis.} By Lemma~\ref{sha-siegel-step-seeds-induced-characters}(3), $g_0^T\notin G_{k_0}$, so $u=(g_0^T)^2$ lies in $G_{k_0}$, which has index $2$. Lemma~\ref{sha-siegel-step-seeds-induced-ribet} gives $\rho=P\rho_0P^{-1}$ with $\rho(u)$ diagonal that satisfies (c), and by Lemma~\ref{sha-siegel-step-seeds-induced-conjugate}, $\rho$ satisfies (b) and is attached through $\iota$ to a representation as in the claim.
\item[(4)] \emph{The determinant.} Let $\mu=\chi_{\mathrm{cyc}}^j\xi\varphi$. As $\varphi=1$ on $G_{k_0}$, (b) gives $\det\rho(h)=\mu(h)$ for $h\in G_{k_0}$.
\item[(5)] \emph{The scalars.} As $g_0^T\in G_F$, $u\in G_{k_0}\cap G_F$, and $\rho(u)_{11}\equiv1$ by (c). By Lemma~\ref{sha-siegel-step-seeds-induced-square-root}(1), $\rho(u)_{11}$ has a square root $x\equiv1$, and by Lemma~\ref{sha-siegel-step-seeds-induced-characters}(3), $\chi_4(g_0^T)\equiv1$ or $\chi_4(g_0^T)\equiv-1$. Let $\alpha=x$ in the first case and $\alpha=-x$ in the second, so that $\alpha^2=\rho(u)_{11}$ and $\alpha\equiv\chi_4(g_0^T)$, and let $\beta=\mu(g_0^T)/\alpha$. As $\rho(u)$ is diagonal, its entry $(2,2)$ is $\det\rho(u)/\alpha^2=\mu(g_0^T)^2/\alpha^2=\beta^2$, by (4). So $\rho(u)=\operatorname{diag}(\alpha^2,\beta^2)$ and $\alpha\beta=\mu(g_0^T)$.
\item[(6)] \emph{The induced representation.} Lemma~\ref{sha-siegel-step-seeds-induced-mix}, for $\Gamma=G$, $\Gamma_0=G_{k_0}$, $\tau=g_0^T$, $\rho$ and $\mu$, gives $R$ with the properties of (d); $R$ is continuous by its last claim, as $G_{k_0}$ is open and $\rho$ is continuous. With (3), (d) holds.
\item[(7)] \emph{Units.} As $\ell-1$ divides $j$, Lemma~\ref{sha-siegel-step-seeds-induced-characters}(4) gives $\mu(g)\equiv\theta(g)$ for $g\in G_F$, so $|\mu(g)|=1$ by part (1) of that lemma. Also $|\alpha|=1$, as $\alpha\equiv\chi_4(g_0^T)$.
\item[(8)] \emph{The congruences (e).} The first is (5). For the second, $\chi_1\chi_4=\theta$ by Lemma~\ref{sha-siegel-step-seeds-induced-characters}(1), so
\[
\beta-\chi_1(g_0^T)=\frac{\big(\mu(g_0^T)-\theta(g_0^T)\big)-\chi_1(g_0^T)\big(\alpha-\chi_4(g_0^T)\big)}{\alpha},
\]
whose numerator has absolute value less than $1$ by (5) and (7), as $g_0^T\in G_F$; and $|\alpha|=1$.
\item[(9)] \emph{The twist.} By Lemma~\ref{sha-siegel-step-seeds-induced-characters}(4), the restriction $\eta$ of $\theta/\mu$ to $G_F$ is continuous, and $\eta(g)\equiv1$ for $g\in G_F$, by (7). Lemma~\ref{sha-siegel-step-seeds-induced-square-root}(2) gives $\delta$ with $\delta^2=\eta$ and $\delta\equiv1$, which is (f).
\end{enumerate}
\end{proof}
% indEx (the private theorem of BirthSeedII.lean, about 70 lines with the private root_ex): (1) is exists_quadratic_character_of_subgroup_index_two (HypEx: Tower.Two.quad and exists_sign_character_index_two, about 40 lines), with galois_character_residual_decomposition for CharDat; (2) is MixOfP; (3) is exists_ribet_lattice, galoisRep_properties_invariant_under_conjugation and prIso; (4) is hdet; (5) is root_ex; (6) is symmetric_square_mixed_exists and continuous_mix_family; (7) is hmu, hmu1 and halpha1; (8) is the first bullet of the refine; (9) is exists_square_root_of_character_near_one and the last bullet.

\begin{lemma}\label{sha-siegel-step-seeds-induced-irreducible-twist}
Let $k$ be an algebraically closed field of characteristic $0$, $\Gamma$ a group, $H$ a normal subgroup of $\Gamma$ that contains $aba^{-1}b^{-1}$ for all $a,b\in\Gamma$, and $r:\Gamma\to GL_m(k)$ a homomorphism that is irreducible and whose restriction to $H$ is not irreducible. Then there are a homomorphism $\eta:\Gamma\to k^\times$ with $\eta\ne1$, $\eta(H)=1$ and $\eta^m=1$, and $Q\in GL_m(k)$ with $\eta(g)Qr(g)=r(g)Q$ for every $g\in\Gamma$.
\end{lemma}
\leanref{Development/AutomorphicArithmetic/AutomorphicVarieties/UnitaryThreeGroups/Clifford.lean}{UpperTheory.V.U3.exists_selfTwist}
\begin{proof}
Let $V=k^m$, which is nonzero as $r$ is irreducible, and say that a subspace is $H$-stable if it is stable under $r(H)$.
\begin{enumerate}
\item[(1)] \emph{A decomposition.} Some $H$-stable subspace is neither $0$ nor $V$. Let $W$ be a minimal nonzero $H$-stable subspace, and $U$ a maximal one among the $H$-stable subspaces other than $0$ and $V$. If $W\subseteq r(g)^{-1}U$ for every $g\in\Gamma$, then the intersection of the $r(g)^{-1}U$ is stable under $r(\Gamma)$ and nonzero, so it is $V$, and $U=V$, which is false. So $W\not\subseteq U'=r(g)^{-1}U$ for some $g$. As $H$ is normal, $U'$ is $H$-stable. Then $W\cap U'$ is $H$-stable and smaller than $W$, so it is $0$; and $r(g)(U'+W)$ is $H$-stable and contains $U$ properly, so it is $V$. Hence $V=U'\oplus W$, with $U'$ and $W$ nonzero and $H$-stable.
\item[(2)] \emph{A traceless endomorphism.} The projection $E$ of $V$ onto $U'$ along $W$ commutes with $r(H)$. So does $P=mE-\operatorname{tr}(E)$, which has trace $0$. It is nonzero: it acts on $W$ by $-\operatorname{tr}(E)$ and on $U'$ by $m-\operatorname{tr}(E)$, and $m\ne0$ in $k$.
\item[(3)] \emph{A common eigenvector.} Let $S$ be the space of endomorphisms of $V$ of trace $0$ that commute with $r(H)$; it contains $P\ne0$. For $g\in\Gamma$, $f_g(X)=r(g)Xr(g)^{-1}$ preserves $S$, as $H$ is normal, and $f_h$ is the identity on $S$ for $h\in H$. As $H$ contains the commutators, the $f_g$ commute on $S$, so they have a common eigenvector $X_0\in S$: $r(g)X_0r(g)^{-1}=c(g)X_0$ with $c(g)\in k$.
\item[(4)] \emph{The character.} As $X_0\ne0$, $c(g)\ne0$, and $c:\Gamma\to k^\times$ is a homomorphism with $c(H)=1$. If $X_0v=0$, then $X_0r(g)v=c(g)^{-1}r(g)X_0v=0$, so the kernel of $X_0$ is stable under $r(\Gamma)$; it is not $V$, so it is $0$, and $X_0$ is invertible. Let $\eta=c$ and $Q=X_0$, so that $\eta(g)Qr(g)=r(g)Q$; taking determinants, $\eta(g)^m=1$. If $\eta=1$, $Q$ commutes with $r(\Gamma)$, so an eigenspace of $Q$ is stable under $r(\Gamma)$, hence equal to $V$, and $Q$ is a scalar; as its trace is $0$ and $m\ne0$ in $k$, $Q=0$, which is false. So $\eta\ne1$.
\end{enumerate}
\end{proof}
% exists_selfTwist (Clifford.lean, 237 lines): exists_compl for (1), exists_traceless for (2), exists_twist with common_eigen for (3), and the end of exists_selfTwist (huniq, hker, hTu, hM and the determinants) for (4).

\begin{lemma}\label{sha-siegel-step-seeds-induced-irreducible-commutators}
Let $\ell$ be a prime, $\Gamma$ a topological group, and $r:\Gamma\to GL_m(\bar{\mathbf Q}_\ell)$ a continuous homomorphism that is irreducible on every open subgroup of $\Gamma$. Let $H$ be an open subgroup of $\Gamma$ and $H'\subseteq H$ a subgroup that contains $aba^{-1}b^{-1}$ for all $a,b\in H$. Then $r$ is irreducible on $H'$.
\end{lemma}
\leanref{Development/AutomorphicArithmetic/AutomorphicVarieties/IdeleGenerators/HolomorphicIrreducibility/IrrAb.lean}{UpperTheory.V.Sd.irrAb}
\leanref{Development/AutomorphicArithmetic/AutomorphicVarieties/IdeleGenerators/ShapeFacts.lean}{UpperTheory.V.Sd.IrrAb}
\begin{proof}
For $a\in H$ and $x\in H'$, $axa^{-1}=(axa^{-1}x^{-1})x\in H'$, so $H'$ is normal in $H$. Suppose that $r$ is not irreducible on $H'$. As $r$ is irreducible on $H$, Lemma~\ref{sha-siegel-step-seeds-induced-irreducible-twist}, for $H'\subseteq H$, gives $\eta:H\to\bar{\mathbf Q}_\ell^\times$ with $\eta\ne1$ and $\eta^m=1$, and $Q$ with $\eta(g)Qr(g)=r(g)Q$ for $g\in H$. As $Qr(g)\ne0$, $\eta(g)=1$ if and only if $Qr(g)=r(g)Q$. So the kernel of $\eta$ is closed in $H$, as $r$ is continuous, and of finite index, as $\eta$ takes values in the $m$-th roots of unity; so it is open in $H$, hence in $\Gamma$, and $r$ is irreducible on it. An eigenspace of $Q$ is stable under $r(\ker\eta)$, so it is everything, and $Q$ is a scalar. Then $\eta(g)Qr(g)=Qr(g)$ for $g\in H$, so $\eta=1$, which is false.
\end{proof}
% irrAb (IrrAb.lean, 78 lines): exists_selfTwist (Lemma sha-siegel-step-seeds-induced-irreducible-twist), hiff, hcl, hrg and Subgroup.isOpen_of_isClosed_of_finiteIndex for the kernel, and the eigenspace htop.

\begin{lemma}\label{sha-siegel-step-seeds-induced-irreducible-mackey}
Let $k$ be a field, $G$ a group, $V$ a $k$-vector space and $\pi:G\to GL(V)$ a homomorphism. Let $H\subseteq G$ be a subgroup of index $2$, $s\in G\setminus H$, and $V_1\subseteq V$ a subspace stable under $\pi(H)$ whose only subspaces stable under $\pi(H)$ are $0$ and $V_1$. Suppose that $V=V_1\oplus\pi(s)V_1$ and that there is no linear isomorphism $e:V_1\to\pi(s)V_1$ with $e(\pi(h)v)=\pi(h)e(v)$ for $h\in H$ and $v\in V_1$. Then the only subspaces of $V$ stable under $\pi(G)$ are $0$ and $V$.
\end{lemma}
\leanref{Development/AutomorphicArithmetic/AutomorphicVarieties/AlgebraTools/InducedAutomorphicForms/MackeyTwo.lean}{UpperTheory.V.AlgebraTools.mackeyTwo}
\begin{proof}
As $H$ has index $2$, $s^{-1}hs\in H$ for $h\in H$, and $s^2\in H$; so $\pi(s)V_1$ is stable under $\pi(H)$, and $\pi(s^2)V_1=V_1$. Let $W$ be stable under $\pi(G)$. Then $V_1\cap W$ is stable under $\pi(H)$, so it is $V_1$ or $0$. If it is $V_1$, then $W$ contains $V_1$ and $\pi(s)V_1$, so $W=V$. Otherwise $V_1\cap W=0$, and $\pi(s)V_1\cap W=\pi(s)(V_1\cap W)=0$. The subspace $V_1\cap(\pi(s)V_1+W)$ is stable under $\pi(H)$.
\begin{enumerate}
\item[(1)] If it is $0$, let $w\in W$ and write $w=a+b$ with $a\in V_1$ and $b\in\pi(s)V_1$. Then $a=w-b$ lies in $V_1\cap(\pi(s)V_1+W)=0$, so $w=b\in\pi(s)V_1\cap W=0$. So $W=0$.
\item[(2)] If it is $V_1$, then $V=\pi(s)V_1\oplus W$; and $\pi(s)V_1\subseteq\pi(s)(\pi(s)V_1+W)=V_1+W$, so $V=V_1\oplus W$. The projection along $W$ gives a linear isomorphism $e:V_1\to\pi(s)V_1$ with $v-e(v)\in W$ for $v\in V_1$. For $h\in H$, $e(\pi(h)v)-\pi(h)e(v)=\pi(h)\big(v-e(v)\big)-\big(\pi(h)v-e(\pi(h)v)\big)$ lies in $W$ and in $\pi(s)V_1$, so it is $0$. This is excluded.
\end{enumerate}
\end{proof}
% mackeyTwo (MackeyTwo.lean, 99 lines): hconj, hss and hst2 for the first sentence, the case split on V_1 inf W, d1 and d2, the case split on V_1 inf (V_1.map (rho s) sup W) for (1) and (2), c1, c2 and graph_equiv for e.

\begin{lemma}\label{sha-siegel-step-seeds-induced-irreducible-planes}
Let $k$ be a field in which $2\ne0$, $\Gamma$ a group, $\Gamma_0\subset\Gamma$ a subgroup of index $2$, $\tau\in\Gamma\setminus\Gamma_0$, $\alpha,\beta\in k^\times$, and $\rho:\Gamma_0\to GL_2(k)$ and $R:\Gamma\to GL_4(k)$ homomorphisms with $\rho(\tau^2)=\operatorname{diag}(\alpha^2,\beta^2)$, $R(h)=M(\beta/\alpha,\rho(h),\rho(\tau^{-1}h\tau))$ for $h\in\Gamma_0$ (Definition~\ref{sha-siegel-step-seeds-induced-datum}(2), over $k$) and $R(\tau)=\operatorname{diag}(\beta,-\alpha,-\beta,\alpha)$. Let $U\subseteq\Gamma$ be a subgroup with $U\not\subseteq\Gamma_0$, such that $\rho$ is irreducible on $U\cap\Gamma_0$ and there is no $Q\in GL_2(k)$ with $\rho(\tau^{-1}h\tau)=Q\rho(h)Q^{-1}$ for all $h\in U\cap\Gamma_0$. Then $R$ is irreducible on $U$.
\end{lemma}
\leanref{Development/AutomorphicArithmetic/AutomorphicVarieties/BuildingsResidueChars/MixInvariantPlaneDecomposition.lean}{UpperTheory.mix_invariant_plane_decomposition}
\leanref{Development/AutomorphicArithmetic/AutomorphicVarieties/IdeleGenerators/QuadraticSplitting/QuadSplit.lean}{UpperTheory.V.Sd.S2.MixIrr}
\begin{proof}
For $a\in k^2$ let $p_X(a)=(a_2/\beta,a_1/\alpha,-a_2/\beta,a_1/\alpha)$ and $p_Y(a)=(a_2,-a_1,a_2,a_1)$, as columns; these maps are injective, with images $X$ and $Y$. A computation gives, for $A,B\in M_2(k)$ and $\gamma\in k^\times$,
\[
M(\beta/\alpha,A,B)p_X(a)=p_X(Aa),\qquad M(\gamma,A,B)p_Y(a)=p_Y(Ba),\qquad\operatorname{diag}(\beta,-\alpha,-\beta,\alpha)p_X(a)=p_Y(a),
\]
and $k^4=X\oplus Y$, as $2\ne0$. So for $h\in\Gamma_0$, $R(h)p_X(a)=p_X(\rho(h)a)$ and $R(h)p_Y(a)=p_Y(\rho(\tau^{-1}h\tau)a)$. Let $s\in U\setminus\Gamma_0$. As $\Gamma_0$ has index $2$, $h_0=s\tau^{-1}\in\Gamma_0$, and $R(s)p_X(a)=R(h_0)R(\tau)p_X(a)=p_Y(\rho(\tau^{-1}h_0\tau)a)$; so $R(s)X=Y$. The subgroup $U\cap\Gamma_0$ has index $2$ in $U$, and $s\notin U\cap\Gamma_0$. Through $p_X$, $X$ with the action of $U\cap\Gamma_0$ is $k^2$ with the action $\rho$, so its only stable subspaces are $0$ and $X$. A linear isomorphism $e:X\to Y$ compatible with the action of $U\cap\Gamma_0$ would give $Q=p_Y^{-1}\circ e\circ p_X\in GL_2(k)$ with $Q\rho(h)=\rho(\tau^{-1}h\tau)Q$ for $h\in U\cap\Gamma_0$, which is excluded. Lemma~\ref{sha-siegel-step-seeds-induced-irreducible-mackey}, for $U\cap\Gamma_0\subset U$, $s$ and $V_1=X$, shows that $R$ is irreducible on $U$.
\end{proof}
% mix_invariant_plane_decomposition (MixInvariantPlaneDecomposition.lean, 175 lines): pX, pY, pX_inj, pY_inj, mix_pX, mix_pY, diag_pX and compl_XY for the computation (by field_simp and ring), hX and hY, hsX and hmap for R(s)X = Y, and mackeyTwo (Lemma sha-siegel-step-seeds-induced-irreducible-mackey). The hypothesis on rho(tau^2), a field of IsMix, is not used.

\begin{lemma}\label{sha-siegel-step-seeds-induced-irreducible-sign}
Let $k$ be an algebraically closed field, $\Gamma$ a group, $\rho,\rho':\Gamma\to GL_2(k)$ homomorphisms, $H'\subseteq\Gamma$ a subgroup on which $\rho$ is irreducible, and $Q\in GL_2(k)$ with $\rho'(h)=Q\rho(h)Q^{-1}$ for $h\in H'$. Let $s\in\Gamma$ with $shs^{-1}\in H'$ for $h\in H'$ and $\det\rho'(s)=\det\rho(s)$. Then $\rho'(s^2)=Q\rho(s^2)Q^{-1}$.
\end{lemma}
\leanref{Development/AutomorphicArithmetic/AutomorphicVarieties/BuildingsResidueChars/MixSignTwist.lean}{UpperTheory.mix_sign_twist}
\leanref{Development/AutomorphicArithmetic/AutomorphicVarieties/IdeleGenerators/QuadraticSplitting/QuadSplit.lean}{UpperTheory.V.Sd.S2.SignTw}
\begin{proof}
Let $N=\rho(s)^{-1}Q^{-1}\rho'(s)Q$. For $h\in H'$, as $shs^{-1}\in H'$,
\[
N\rho(h)=\rho(s)^{-1}Q^{-1}\rho'(s)\rho'(h)Q=\rho(s)^{-1}Q^{-1}\rho'(shs^{-1})\rho'(s)Q=\rho(s)^{-1}\rho(shs^{-1})Q^{-1}\rho'(s)Q=\rho(h)N.
\]
An eigenspace of $N$ is then stable under $\rho(H')$, so it is $k^2$, and $N=x$ is a scalar. As $\det N=1$, $x^2=1$, so $N^2=1$ and $N$ commutes with every matrix. Then $\rho'(s)=Q\rho(s)NQ^{-1}$, and $\rho'(s^2)=Q\rho(s)N\rho(s)NQ^{-1}=Q\rho(s)^2N^2Q^{-1}=Q\rho(s^2)Q^{-1}$.
\end{proof}
% mix_sign_twist (MixSignTwist.lean, 49 lines): hcomm, schur2, hd, hx2, hc, hNN and the final calc.

\begin{lemma}\label{sha-siegel-step-seeds-induced-irreducible-residue}
Let $D$ be a Siegel datum (Definition~\ref{sha-siegel-datum}) and $T$ a tower for $D$ that satisfies the conditions (a) to (r) of Definition~\ref{sha-siegel-tower-towers}(4), with fields $F$ and $L$ and relabelling characters $\chi_1,\dots,\chi_4$. Then some $g\in G_L\cap\ker\chi_{\mathrm{cyc}}$ has $\bar\chi^{(0,1)}(g)\ne1$.
\end{lemma}
\leanref{Development/AutomorphicArithmetic/AutomorphicVarieties/BuildingsResidueChars/ResidualCharactersPsiNe.lean}{UpperTheory.residual_characters_psi_ne}
\leanref{Development/AutomorphicArithmetic/AutomorphicVarieties/IdeleGenerators/QuadraticSplitting/QuadSplit.lean}{UpperTheory.V.Sd.S2.PsiNe}
\begin{proof}
Suppose that $\bar\chi^{(0,1)}(g)=1$ for every $g\in G_L\cap\ker\chi_{\mathrm{cyc}}$, and let $\eta$ be the restriction of $\bar\chi^{(0,1)}$ to $G_L$.
\begin{enumerate}
\item[(1)] \emph{Square.} By condition (n), $\eta^2=1$.
\item[(2)] \emph{Triviality on $\ker\omega$.} Let $k_0$ be the field of condition (p), so that for $g\in G_F$, $\bar\chi^{(0,1)}(g)=1$ if and only if $g\in G_{k_0}$. Let $g\in G_L$ with $\omega(g)=1$. As $\bar\omega=\bar\chi_{\mathrm{cyc}}$, $\chi_{\mathrm{cyc}}(g)\equiv1\pmod\ell$, so $\chi_{\mathrm{cyc}}(g^{\ell^n})=\chi_{\mathrm{cyc}}(g)^{\ell^n}\equiv1\pmod{\ell^{n+1}}$. As $\ell$ is odd and $\eta^2=1$, $\eta(g^{\ell^n})=\eta(g)$, so $g^{\ell^n}$ lies in the set $C$ of the $h\in G_L$ with $g^{-1}h\in G_{k_0}$, which is closed in the compact group $G_F$. So the sequence $(g^{\ell^n})_n$ has a cluster point $h\in C$, and $\chi_{\mathrm{cyc}}(h)=1$ by continuity. Then $\eta(h)=1$ by the assumption, and $\eta(g)=\eta(h)$, as $g^{-1}h\in G_F\cap G_{k_0}$. So $\eta(g)=1$.
\item[(3)] \emph{Permutation.} The entries $(1,3)$ and $(2,4)$ have weight $(0,1)$, so $\eta=\bar\chi_1\bar\chi_3^{-1}=\bar\chi_2\bar\chi_4^{-1}$ on $G_L$. With (1), $\eta\bar\chi_1=\bar\chi_3$, $\eta\bar\chi_2=\bar\chi_4$, $\eta\bar\chi_3=\bar\chi_1$ and $\eta\bar\chi_4=\bar\chi_2$ on $G_L$.
\item[(4)] \emph{Contradiction.} By (1) to (3) and condition (g), $\eta=1$. But by condition (f) for the root $(0,1)$, $\bar\chi^{(0,1)}$ is nontrivial on $G_L$.
\end{enumerate}
\end{proof}
% residual_characters_psi_ne (ResidualCharactersPsiNe.lean, 141 lines): eta, eta_sq (Two.sq) for (1); eta_om with Two.quad, dvd_of_om, lim_pt and closed_fix for (2); eta_mul, eta_mul' and etaL_perm for (3); Ok.g3' and Ok.g3 for wPs for (4).

\begin{lemma}\label{sha-siegel-step-seeds-induced-irreducible}
Let $D$, $T$ and $\iota$ be as in Definition~\ref{sha-siegel-step-seeds-induced-datum}, and let $(k_0,\varphi,j,t,\rho,\alpha,\beta,R,\delta)$ be an induction datum for $T$ and $\iota$. Then the restriction of $R$ to the subgroup $G_L\cap\ker\chi_{\mathrm{cyc}}$ of $G_F$ is irreducible.
\end{lemma}
\leanref{Development/AutomorphicArithmetic/AutomorphicVarieties/BuildingsResidueChars/InducedIrreducible.lean}{UpperTheory.induced_irreducible}
\leanref{Development/AutomorphicArithmetic/AutomorphicVarieties/IdeleGenerators/QuadraticSplitting/QuadSplit.lean}{UpperTheory.V.Sd.S2.IndIrr}
\begin{proof}
Let $\tau=g_0^T$ and $U=G_L\cap\ker\chi_{\mathrm{cyc}}$, a subgroup of $G$.
\begin{enumerate}
\item[(1)] \emph{Irreducibility of $\rho$.} The subgroup $H=G_L\cap G_{k_0}$ of $G_{k_0}$ is open, and $U\cap G_{k_0}\subseteq H$ contains $aba^{-1}b^{-1}$ for $a,b\in H$, as $\chi_{\mathrm{cyc}}(aba^{-1}b^{-1})=1$. By (b), $\rho$ is continuous and irreducible on every open subgroup of $G_{k_0}$. So $\rho$ is irreducible on $U\cap G_{k_0}$, by Lemma~\ref{sha-siegel-step-seeds-induced-irreducible-commutators}.
\item[(2)] \emph{$U\not\subseteq G_{k_0}$.} By Lemma~\ref{sha-siegel-step-seeds-induced-irreducible-residue}, some $g\in U$ has $\bar\chi^{(0,1)}(g)\ne1$, and $g\notin G_{k_0}$ by (a).
\item[(3)] \emph{$\rho$ and its conjugate.} Suppose that $Q\in GL_2(\bar{\mathbf Q}_\ell)$ has $\rho(\tau^{-1}h\tau)=Q\rho(h)Q^{-1}$ for $h\in U\cap G_{k_0}$. As $G_{k_0}$ is normal in $G$, $\rho'(h)=\rho(\tau^{-1}h\tau)$ is a homomorphism on $G_{k_0}$. The open subgroup $G_L\cap G_{k_0}$ has finite index in $G$, so $s=s_0^n$ lies in it for some $n\ge1$, where $s_0$ is the element of $D$. Conjugation by $s$ preserves $U\cap G_{k_0}$, and $\det\rho'(s)=\chi_{\mathrm{cyc}}(s)^j\xi(s)=\det\rho(s)$ by (b), as characters of $G$ take the same value at $s$ and at $\tau^{-1}s\tau$. By Lemma~\ref{sha-siegel-step-seeds-induced-irreducible-sign}, $\rho(\tau^{-1}s^2\tau)=Q\rho(s^2)Q^{-1}$, so $\rho(s^2)$ and $\rho(\tau^{-1}s^2\tau)$ have the same characteristic polynomial. As $s_0\in I_\ell$, $s^2\in I_{v_\ell}(k_0)$ and $\tau^{-1}s^2\tau\in I_{\tau^{-1}v_\ell}(k_0)$. If $\rho$ is ordinary at a place $w$ of weights $(m_1,m_2)$, then for $x\in I_w(k_0)$ the characteristic polynomial of $\rho(x)$ is $(X-\chi_{\mathrm{cyc}}(x)^{m_1})(X-\chi_{\mathrm{cyc}}(x)^{m_2})$. As $\tau\notin G_{k_0}$ (Lemma~\ref{sha-siegel-step-seeds-induced-characters}(3)), (b) applies to $v_\ell$ and $\tau^{-1}v_\ell$; so, with $e=\chi_{\mathrm{cyc}}(s^2)=\chi_{\mathrm{cyc}}(\tau^{-1}s^2\tau)$, one of the two characteristic polynomials is $(X-e^{j+t})(X-e^{-t})$ and the other is $(X-e^j)(X-1)$. As $\chi_{\mathrm{cyc}}(s_0)$ has infinite order, so has $e=\chi_{\mathrm{cyc}}(s_0)^{2n}$; as $j+t\ne0$ and $t\ne0$, the first polynomial does not vanish at $1$, and the second does. So there is no such $Q$.
\item[(4)] \emph{Conclusion.} By (d), Lemma~\ref{sha-siegel-step-seeds-induced-irreducible-planes} applies to $\Gamma=G$, $\Gamma_0=G_{k_0}$, $\tau$, $\alpha$, $\beta$, $\rho$, $R$ and $U$, and by (1) to (3) it shows that $R$ is irreducible on $U$, which is the subgroup $G_L\cap\ker\chi_{\mathrm{cyc}}$ of $G_F$.
\end{enumerate}
\end{proof}
% induced_irreducible (InducedIrreducible.lean, 141 lines): V = T.GalL and U = V inf ker cycZ, hU and hU' for the identification with T.G'; (1) is hirr through irrAb (Lemma sha-siegel-step-seeds-induced-irreducible-commutators); (2) is residual_characters_psi_ne with Hyp.ps; (3) is the second goal of key: conjHom and hchar for rho', exists_pow_mem_of_index_ne_zero with decompSigma_mem for s, mix_sign_twist with hdet and hcyc, cp_conj, cp_ord, ine_move, decompSigma_ord for he, dat.above and cp_ne; (4) is mix_invariant_plane_decomposition (Lemma sha-siegel-step-seeds-induced-irreducible-planes). With its dependencies (IrrAb, Clifford, MackeyTwo, MixInvariantPlaneDecomposition, MixSignTwist, ResidualCharactersPsiNe) about 920 lines.

\begin{lemma}\label{sha-siegel-step-seeds-induced-flags-line}
Let $\ell$ be a prime, and write $x\equiv y$ for $|x-y|<1$ in $\bar{\mathbf Q}_\ell$. Let $S$ be a set, and let $A:S\to GL_2(\bar{\mathbf Q}_\ell)$ and $x,\kappa_1,\kappa_2:S\to\bar{\mathbf Q}_\ell^\times$ be maps such that: $x(s_0)\not\equiv1$ for some $s_0\in S$; for every $s\in S$, $A(s)_{11}\equiv1$, $A(s)_{22}\equiv x(s)$, $|A(s)_{21}|<1$ and $|A(s)_{12}|\le1$; and there is $P\in GL_2(\bar{\mathbf Q}_\ell)$ such that, for every $s\in S$, $PA(s)P^{-1}$ is upper triangular with diagonal entries $\kappa_1(s),\kappa_2(s)$. Then $\det A(s)=\kappa_1(s)\kappa_2(s)$ for every $s$, and:
\begin{enumerate}
\item if $\kappa_1(s)\equiv x(s)$ for every $s$, there is $u\in\bar{\mathbf Q}_\ell$ with $|u|\le1$ and $A(s)\binom u1=\kappa_1(s)\binom u1$ for every $s$;
\item if $\kappa_1(s)\equiv1$ for every $s$, there is $y\in\bar{\mathbf Q}_\ell$ with $|y|<1$ and $A(s)\binom1y=\kappa_1(s)\binom1y$ for every $s$.
\end{enumerate}
\end{lemma}
\leanref{Development/AutomorphicArithmetic/AutomorphicVarieties/BuildingsResidueChars/MixStableLineLocal.lean}{UpperTheory.mix_stable_line_local}
\leanref{Development/AutomorphicArithmetic/AutomorphicVarieties/IdeleGenerators/QuadraticSplitting/QuadSplit.lean}{UpperTheory.V.Sd.S2.LocLine}
\begin{proof}
Let $\binom pq$ be the first column of $P^{-1}$, which is nonzero. Comparing the first columns of $A(s)P^{-1}=P^{-1}\cdot PA(s)P^{-1}$ gives $A(s)\binom pq=\kappa_1(s)\binom pq$, and $\det A(s)=\det(PA(s)P^{-1})=\kappa_1(s)\kappa_2(s)$. Write $A=A(s_0)$ and $\kappa=\kappa_1(s_0)$.

(1) If $\kappa\equiv A_{11}$, then $x(s_0)\equiv\kappa\equiv A_{11}\equiv1$, against the choice of $s_0$; so $|\kappa-A_{11}|\ge1$. The first row of $A\binom pq=\kappa\binom pq$ gives $(\kappa-A_{11})p=A_{12}q$, so $|p|\le|A_{12}||q|\le|q|$. Hence $q\ne0$, and $u=p/q$ has $|u|\le1$ and $\binom u1=q^{-1}\binom pq$.

(2) In the same way $|\kappa-A_{22}|\ge1$, as otherwise $1\equiv\kappa\equiv A_{22}\equiv x(s_0)$. The second row gives $(\kappa-A_{22})q=A_{21}p$, so $|q|\le|A_{21}||p|$. Hence $p\ne0$, and $y=q/p$ has $|y|\le|A_{21}|<1$ and $\binom1y=p^{-1}\binom pq$.
\end{proof}
% mix_stable_line_local (MixStableLineLocal.lean, 96 lines): col_eq for the eigenvector and the determinant, one_le_norm and norm_le_of for (1) and (2).

\begin{lemma}\label{sha-siegel-step-seeds-induced-flags-matrix}
Let $\ell$ be an odd prime and $\gamma,u,y\in\bar{\mathbf Q}_\ell$ with $|\gamma|=1$, $|u|\le1$ and $|y|<1$. For $p,q,r,s,t\in\bar{\mathbf Q}_\ell$ let
\[
n(p,q,r,s,t)=\begin{pmatrix}1&0&0&0\\p&1&0&0\\q&s&1&0\\r&t&-p&1\end{pmatrix},
\]
and let
\[
n_1=n\Big(\gamma u,-1,\gamma u,\frac{-2y}{1+\gamma uy},-\frac{1-\gamma uy}{1+\gamma uy}\Big),\qquad n_2=n\Big(-u,1,u,\frac{-2y}{\gamma+uy},\frac{\gamma-uy}{\gamma+uy}\Big).
\]
Then $n_1$ and $n_2$ are defined and lower unipotent (Lemma~\ref{sha-siegel-step-restriction-matrices}), their entries have absolute value at most $1$, their entries $(3,2)$ have absolute value less than $1$, and their entries $(3,1)$ have absolute value $1$. Let $A,B\in M_2(\bar{\mathbf Q}_\ell)$ and $a,a',b,b'\in\bar{\mathbf Q}_\ell$ with $a,b\ne0$, $\det A=aa'$ and $\det B=bb'$, and let $M=M(\gamma,A,B)$ (Definition~\ref{sha-siegel-step-seeds-induced-datum}(2)).
\begin{enumerate}
\item If $A\binom u1=a\binom u1$ and $B\binom1y=b\binom1y$, then $n_1^*Mn_1$ is upper triangular with diagonal entries $a,b,b',a'$.
\item If $A\binom1y=a\binom1y$ and $B\binom u1=b\binom u1$, then $n_2^*Mn_2$ is upper triangular with diagonal entries $b,a,a',b'$.
\end{enumerate}
\end{lemma}
\leanref{Development/AutomorphicArithmetic/AutomorphicVarieties/BuildingsResidueChars/MixFlagExists.lean}{UpperTheory.mix_flag_exists}
\leanref{Development/AutomorphicArithmetic/AutomorphicVarieties/IdeleGenerators/QuadraticSplitting/QuadSplit.lean}{UpperTheory.V.Sd.S2.MixFlag}
\begin{proof}
As $|\gamma uy|<1$ and $|uy|<1$, the elements $1\pm\gamma uy$ and $\gamma\pm uy$ have absolute value $1$. So $n_1$ and $n_2$ are defined, their entries have absolute value at most $1$, their entries $(3,2)$ have absolute value at most $|2y|<1$, and their entries $(3,1)$ are $\mp1$. A matrix $n(p,q,r,s,t)$ with $t+ps-q=0$ is lower unipotent, and this identity holds for $n_1$ and $n_2$. For (1), the eigenvectors and the determinants give $A=\left(\begin{smallmatrix}a'+uc&(a-a'-uc)u\\c&a-cu\end{smallmatrix}\right)$ with $c=A_{21}$ and $B=\left(\begin{smallmatrix}b-dy&d\\(b-b'-dy)y&b'+dy\end{smallmatrix}\right)$ with $d=B_{12}$, and the entries of $n_1^*Mn_1$ are computed from these. (2) is the same computation with the roles of $A$ and $B$ exchanged.
\end{proof}
% mix_flag_exists (MixFlagExists.lean, 144 lines): nn, adj_nn, low_nn and lowQ_nn for the matrices, par1 and par2 for the forms of A and B, case1 and case2 for the computation of the products.

\begin{lemma}\label{sha-siegel-step-seeds-induced-flags}
Let $D$, $T$ and $\iota$ be as in Definition~\ref{sha-siegel-step-seeds-induced-datum}, and let $(k_0,\varphi,j,t,\rho,\alpha,\beta,R,\delta)$ be an induction datum for $T$ and $\iota$. Then there are integers $m_1,\dots,m_4$ with $m_i\ge m_{i+1}+2$ such that, for every place $w$ of $\bar{\mathbf Q}$ above $\ell$, there is a lower unipotent $n\in M_4(\bar{\mathbf Q}_\ell)$ (Lemma~\ref{sha-siegel-step-restriction-matrices}) with integral entries, $|n_{32}|<1$ and $|n_{31}|=1$, such that for $s\in D_w(F)$ the matrix $n^*R(s)n$ is upper triangular, with diagonal entries $d_1(s),\dots,d_4(s)$ say, $(d_1(s),d_2(s),d_3(s),d_4(s))\equiv(\xi(s),1,\xi(s),1)$, and $d_i(s)=\chi_{\mathrm{cyc}}(s)^{m_i}$ if $s\in I_w(F)$.
\end{lemma}
\leanref{Development/AutomorphicArithmetic/AutomorphicVarieties/BuildingsResidueChars/InducedFlagCondition.lean}{UpperTheory.induced_flag_condition}
\leanref{Development/AutomorphicArithmetic/AutomorphicVarieties/IdeleGenerators/QuadraticSplitting/QuadSplit.lean}{UpperTheory.V.Sd.S2.IndFlag}
\begin{proof}
Let $m=(j+t,j,0,-t)$; then $m_i\ge m_{i+1}+2$, as $j\ge2$ and $t\ge2$ by (b). Let $w$ be a place above $\ell$, let $\tau=g_0^T$ and $w'=\tau^{-1}w$, which lies above $\ell$, and let $\gamma=\beta/\alpha$. For $s\in D_w(F)$ let $s'=\tau^{-1}s\tau$.
\begin{enumerate}
\item[(1)] \emph{The elements $s$ and $s'$.} Let $s\in D_w(F)$. By (a), $s\in D_w(\mathbf Q)\subseteq G_{k_0}$ and $s'\in D_{w'}(\mathbf Q)\subseteq G_{k_0}$; and $s'\in G_F$, as $\tau\in G_F$. If $s\in I_w(F)$, then $s\in I_w(k_0)$ and $s'\in I_{w'}(k_0)$. Characters of $G$ take the same value at $s$ and at $s'$. By (d), $R(s)=M(\gamma,\rho(s),\rho(s'))$, and by (c), $\rho(s)$ and $\rho(s')$ satisfy the conditions on entries of Lemma~\ref{sha-siegel-step-seeds-induced-flags-line}, with $x=\xi(s)$. By (e) and Lemma~\ref{sha-siegel-step-seeds-induced-characters}(1), $|\gamma|=1$.
\item[(2)] \emph{A residue.} By condition (r) of Definition~\ref{sha-siegel-tower-towers}(4), some $s_0\in D_w(F)$ has $\bar\chi^{(1,0)}(s_0)^4\ne1$. So $\bar\chi^{(1,0)}(s_0)\ne1$, and $\xi(s_0)\not\equiv1$ by Lemma~\ref{sha-siegel-step-seeds-induced-characters}(5).
\item[(3)] \emph{The flags of $\rho$.} By Lemma~\ref{sha-siegel-step-seeds-induced-characters}(3), $\tau\notin G_{k_0}$, so (b) applies to $w$ and $w'$. Suppose first that $\rho$ is ordinary at $w$ of weights $(j+t,-t)$ with residues $(\xi,1)$ and at $w'$ of weights $(j,0)$ with residues $(1,\xi)$. Evaluating these flags at $s$ and at $s'$ gives $P,P'\in GL_2(\bar{\mathbf Q}_\ell)$ and functions $\kappa_1,\kappa_2,\kappa'_1,\kappa'_2$ on $D_w(F)$ such that, for $s\in D_w(F)$, $P\rho(s)P^{-1}$ and $P'\rho(s')P'^{-1}$ are upper triangular with diagonal entries $\kappa_1(s),\kappa_2(s)$ and $\kappa'_1(s),\kappa'_2(s)$. By (1), $\kappa_1\equiv\xi$, $\kappa_2\equiv1$, $\kappa'_1\equiv1$ and $\kappa'_2\equiv\xi$, and for $s\in I_w(F)$, $(\kappa_1(s),\kappa_2(s),\kappa'_1(s),\kappa'_2(s))=(\chi_{\mathrm{cyc}}(s)^{j+t},\chi_{\mathrm{cyc}}(s)^{-t},\chi_{\mathrm{cyc}}(s)^j,1)$.
\item[(4)] \emph{Eigenvectors.} Lemma~\ref{sha-siegel-step-seeds-induced-flags-line}, with $s_0$ of (2), gives $\det\rho(s)=\kappa_1(s)\kappa_2(s)$ and $\det\rho(s')=\kappa'_1(s)\kappa'_2(s)$, an element $u$ with $|u|\le1$ and $\rho(s)\binom u1=\kappa_1(s)\binom u1$, and an element $y$ with $|y|<1$ and $\rho(s')\binom1y=\kappa'_1(s)\binom1y$, for all $s\in D_w(F)$.
\item[(5)] \emph{The flag of $R$.} Let $n=n_1$ be the matrix of Lemma~\ref{sha-siegel-step-seeds-induced-flags-matrix} for $\gamma$, $u$ and $y$. It is lower unipotent, with integral entries, $|n_{32}|<1$ and $|n_{31}|=1$, and by (1) of that lemma, for $s\in D_w(F)$ the matrix $n^*R(s)n$ is upper triangular with diagonal entries $(\kappa_1(s),\kappa'_1(s),\kappa'_2(s),\kappa_2(s))$. By (3), these are congruent to $(\xi(s),1,\xi(s),1)$, and they are the $\chi_{\mathrm{cyc}}(s)^{m_i}$ for $s\in I_w(F)$.
\item[(6)] \emph{The other case.} Otherwise, by (b), $\rho$ is ordinary at $w$ of weights $(j,0)$ with residues $(1,\xi)$ and at $w'$ of weights $(j+t,-t)$ with residues $(\xi,1)$. As in (3) and (4), with the roles of $w$ and $w'$ exchanged, now $\kappa_1\equiv1$, $\kappa_2\equiv\xi$, $\kappa'_1\equiv\xi$ and $\kappa'_2\equiv1$, with $(\kappa_1,\kappa_2,\kappa'_1,\kappa'_2)=(\chi_{\mathrm{cyc}}^j,1,\chi_{\mathrm{cyc}}^{j+t},\chi_{\mathrm{cyc}}^{-t})$ on $I_w(F)$, and there are $y$ with $|y|<1$ and $\rho(s)\binom1y=\kappa_1(s)\binom1y$, and $u$ with $|u|\le1$ and $\rho(s')\binom u1=\kappa'_1(s)\binom u1$. With $n=n_2$, (2) of Lemma~\ref{sha-siegel-step-seeds-induced-flags-matrix} makes $n^*R(s)n$ upper triangular with diagonal entries $(\kappa'_1(s),\kappa_1(s),\kappa_2(s),\kappa'_2(s))$, which are congruent to $(\xi(s),1,\xi(s),1)$ and are the $\chi_{\mathrm{cyc}}(s)^{m_i}$ for $s\in I_w(F)$.
\end{enumerate}
\end{proof}
% induced_flag_condition (InducedFlagCondition.lean, 138 lines): m = (k + t, k, 0, -t) in the Lean's names, with dat.hk and dat.ht; (1) is ex_lift, dec_conj, ine_conj, hcomm, hR, hR', hmix and hγ; (2) is hne, from Two.bigL and chars.al; (3) is the case split on dat.above (with chars.out) and pull; (4) is mix_stable_line_local (Lemma sha-siegel-step-seeds-induced-flags-line); (5) and (6) are the two parts of mix_flag_exists (Lemma sha-siegel-step-seeds-induced-flags-matrix).

\begin{lemma}\label{sha-siegel-step-seeds-induced-deletion}
Let $D$, $T$ and $\iota$ be as in Definition~\ref{sha-siegel-step-seeds-induced-datum}, and let $(k_0,\varphi,j,t,\rho,\alpha,\beta,R,\delta)$ be an induction datum for $T$ and $\iota$. Then there are $b\in\mathbf Z$, $N\ge1$ and an open subgroup $U$ of $G_F$ such that $(\delta(g)^2\chi_{\mathrm{cyc}}(g)^{-b})^N=1$ for $g\in G_F$, and such that for every place $w$ above $\ell$ and $s\in I_w(F)$: $\delta(s)^2=\chi_{\mathrm{cyc}}(s)^b$ if $s\in U$, and $\delta(s)=1$ if $\chi_{\mathrm{cyc}}(s)=1$.
\end{lemma}
\leanref{Development/AutomorphicArithmetic/AutomorphicVarieties/BuildingsResidueChars/InducedDeletionCondition.lean}{UpperTheory.induced_deletion_condition}
\leanref{Development/AutomorphicArithmetic/AutomorphicVarieties/IdeleGenerators/QuadraticSplitting/QuadSplit.lean}{UpperTheory.V.Sd.S2.IndDel}
\begin{proof}
Let $k$ be the integer of $D$ and $b=k+1-j$. As $T$ is deep, there is $m$ prime to $\ell$ with $\psi(g)^m=1$ for $g\in G_F$; and $m\ge1$, as $\ell>1$.
\begin{enumerate}
\item[(1)] \emph{The square of $\delta$.} As $\theta=\chi_{\mathrm{cyc}}\omega^{-1}\xi\chi_{\mathrm{cyc}}^k\psi$ and $\mu=\chi_{\mathrm{cyc}}^j\xi\varphi$, (f) gives $\delta(g)^2=\theta(g)/\mu(g)=\chi_{\mathrm{cyc}}(g)^b\omega(g)^{-1}\psi(g)\varphi(g)^{-1}$ for $g\in G_F$.
\item[(2)] \emph{The exponent.} Let $N=2(\ell-1)m\ge1$. As $\omega^{\ell-1}=1$, $\psi^m=1$ on $G_F$ and $\varphi^2=1$ by (a), (1) gives $(\delta(g)^2\chi_{\mathrm{cyc}}(g)^{-b})^N=(\omega(g)^{-1}\psi(g)\varphi(g)^{-1})^N=1$ for $g\in G_F$.
\item[(3)] \emph{The subgroup.} Let $U$ be the subgroup of the $g\in G_F$ with $\psi(g)=\omega(g)=1$. It is open, as $\psi$ has open kernel and so has $\omega$, by Lemma~\ref{sha-siegel-step-seeds-induced-characters}(6).
\item[(4)] \emph{Inertia.} Let $w$ be a place above $\ell$ and $s\in I_w(F)$. Then $s\in D_w(\mathbf Q)\subseteq G_{k_0}$ by (a), so $\varphi(s)=1$. If $s\in U$, (1) gives $\delta(s)^2=\chi_{\mathrm{cyc}}(s)^b$.
\item[(5)] \emph{Where $\chi_{\mathrm{cyc}}(s)=1$.} Let moreover $\chi_{\mathrm{cyc}}(s)=1$. Then $\omega(s)=1$ by Lemma~\ref{sha-siegel-step-seeds-induced-characters}(6). By condition (k) of Definition~\ref{sha-siegel-tower-towers}(4), $\bar r_0(s)=1$; its entry $(2,2)$ is $\bar\Psi(s)$, and $\Psi(s)=\chi_{\mathrm{cyc}}(s)^k\psi(s)=\psi(s)$, so $\bar\psi(s)=1$. Then $\psi(s)=1$: in $0=\psi(s)^m-1=(\psi(s)-1)(1+\psi(s)+\dots+\psi(s)^{m-1})$ the second factor has residue $m\ne0$, so it is a unit. So $s\in U$, and $\delta(s)^2=\chi_{\mathrm{cyc}}(s)^b=1$ by (4). As $\delta(s)\equiv1$ by (f), and $|-1-1|=|2|=1$ for $\ell$ odd, $\delta(s)\ne-1$; so $\delta(s)=1$.
\end{enumerate}
\end{proof}
% induced_deletion_condition (InducedDeletionCondition.lean, 106 lines): key for (1), with the Lean's integers D.e = k and k = j; h.hyp.deep and hm0 for m; b = D.e + 1 - k, N = (l - 1) m 2 with om_pow and hφ for (2); U the preimage of ker psi inf ker om, with psi_open and chars.om_open, for (3); hφI (Hyp.spl) and fin for (4); chars.om, Ok.ineL at the entry (2,2), teichmuller_character_spec (Sd.S3.Teich, 25 lines) and δ_one with norm_two for (5).

\begin{lemma}\label{sha-siegel-step-seeds-induced-local}
Let $D$, $T$ and $\iota$ be as in Definition~\ref{sha-siegel-step-seeds-induced-datum}, and let $(k_0,\varphi,j,t,\rho,\alpha,\beta,R,\delta)$ be an induction datum for $T$ and $\iota$, with twisted restriction $r_\delta$. Then there are matrices $n_v\in M_4(\bar{\mathbf Q}_\ell)$ for $v\in\Sigma_T$, integers $m_1,\dots,m_4$ with $m_i\ge m_{i+1}+2$, $b\in\mathbf Z$ and $N\ge1$ such that $(r_\delta,(n_v))$ satisfies (b) of Lemma~\ref{sha-siegel-step-seeds-final}, it has inertial weights $(2m_1+b,\dots,2m_4+b)$ (Definition~\ref{sha-siegel-tower-ideal}), and $(\delta(g)^2\chi_{\mathrm{cyc}}(g)^{-b})^N=1$ for $g\in G_F$.
\end{lemma}
\leanref{Development/AutomorphicArithmetic/Foundations/GaloisRepresentations/AdmissibleResidues/IndLoc0.lean}{UpperTheory.GalN.V.indLoc0}
\leanref{Development/AutomorphicArithmetic/AutomorphicVarieties/BuildingsResidueChars/InertiaActsTriviallyOnRootsOfUnity.lean}{UpperTheory.inertia_acts_trivially_on_roots_of_unity}
\begin{proof}
Let $m_1,\dots,m_4$ be as in Lemma~\ref{sha-siegel-step-seeds-induced-flags}, and $b$, $N$ and $U$ as in Lemma~\ref{sha-siegel-step-seeds-induced-deletion}. For $v\in\Sigma_T$, which lies above $\ell$, let $n_v$ be the matrix $n$ of Lemma~\ref{sha-siegel-step-seeds-induced-flags} at $w=v$, with diagonal entries $d_{v,i}(s)$ of $n_v^*R(s)n_v$.
\begin{enumerate}
\item[(1)] \emph{Flags.} As $n_v^*r_\delta(s)n_v=\delta(s)\,n_v^*R(s)n_v$, it is upper triangular for $s\in D_v(F)$, with diagonal entries $\psi_{v,i}(s)=\delta(s)d_{v,i}(s)$. The matrices $n_v$ have the properties asked in (b) by Lemma~\ref{sha-siegel-step-seeds-induced-flags}.
\item[(2)] \emph{The set $\mu_v$.} Let $s\in\mu_v$, so that $s^{\ell^a}$ lies in the closure of the commutator subgroup of $D_v(F)$ for some $a\ge0$. As in step (4) of the proof of Lemma~\ref{sha-siegel-step-restriction}, $s\in I_v(F)$ and $\chi_{\mathrm{cyc}}(s)^{\ell^a}=1$; so $\chi_{\mathrm{cyc}}(s)=1$, as $\mathbf Z_\ell^\times$ has no element of order $\ell$ for $\ell$ odd. So $d_{v,i}(s)=\chi_{\mathrm{cyc}}(s)^{m_i}=1$, and $\delta(s)=1$ by Lemma~\ref{sha-siegel-step-seeds-induced-deletion}; hence $\psi_{v,i}(s)=1$.
\item[(3)] \emph{Residues.} Let $s\in D_v(F)$. By (a) of Definition~\ref{sha-siegel-step-seeds-induced-datum}, $s\in G_{k_0}$, as $v$ lies above $\ell$. As $\delta(s)\equiv1$ and $|\xi(s)|=1$, $\psi_{v,i}(s)=\delta(s)d_{v,i}(s)$ is congruent to the $i$-th entry of $(\xi(s),1,\xi(s),1)$, hence to $\chi_i(s)$, by Lemma~\ref{sha-siegel-step-seeds-induced-characters}(2).
\item[(4)] \emph{Weights.} For $s\in U\cap I_v(F)$, $\psi_{v,i}(s)^2=\delta(s)^2\chi_{\mathrm{cyc}}(s)^{2m_i}=\chi_{\mathrm{cyc}}(s)^{2m_i+b}$, by Lemma~\ref{sha-siegel-step-seeds-induced-deletion}. So $(r_\delta,(n_v))$ has inertial weights $(2m_i+b)_i$.
\end{enumerate}
The last claim is Lemma~\ref{sha-siegel-step-seeds-induced-deletion}.
\end{proof}
% indLoc0 (IndLoc0.lean, about 50 lines): induced_flag_condition and induced_deletion_condition, key and key' for (1), inertia_acts_trivially_on_roots_of_unity (MuIne, about 150 lines with Rq.powIne and B1.decComm) for (2), Hyp.spl, near and CharDat.inn for (3), and the last bullet for (4).

\begin{lemma}[Base change and automorphic induction {\cite[Chapter~3, Theorems~4.2 and~5.1]{AC89}}]\label{sha-siegel-step-seeds-induced-automorphic-published}
Let $\ell$ be a prime and $\iota:\bar{\mathbf Q}_\ell\cong\mathbf C$. For a number field $E$ write $G_E=\operatorname{Gal}(\bar{\mathbf Q}/E)$.
\begin{enumerate}
\item Let $F$ be a totally real number field, $E\supset F$ a totally complex quadratic extension, and $r:G_F\to GL_2(\bar{\mathbf Q}_\ell)$ a continuous representation, irreducible on $G_E$ and attached through $\iota$ to a regular algebraic cuspidal automorphic representation $\pi$ of $GL_2(\mathbf A_F)$. Then the restriction of $r$ to $G_E$ is attached through $\iota$ to a regular algebraic cuspidal automorphic representation of $GL_2(\mathbf A_E)$ that has, at each embedding $\sigma:E\to\mathbf C$, the infinitesimal character of $\pi$ at the infinite place of $F$ induced by $\sigma$.
\item Let $F\subseteq E$ be totally complex number fields, with $E$ Galois over $F$ with solvable Galois group, and $r:G_F\to GL_2(\bar{\mathbf Q}_\ell)$ a continuous representation, irreducible on $G_E$ and attached through $\iota$ to a regular algebraic cuspidal automorphic representation $\pi$ of $GL_2(\mathbf A_F)$. Then the restriction of $r$ to $G_E$ is attached through $\iota$ to a regular algebraic cuspidal automorphic representation of $GL_2(\mathbf A_E)$ that has, at each embedding $\sigma:E\to\mathbf C$, the infinitesimal character of $\pi$ at the restriction of $\sigma$ to $F$.
\item Let $K\subset M$ be totally complex number fields with $[M:K]=2$, and $\pi$ a cuspidal automorphic representation of $GL_2(\mathbf A_M)$ such that:
\begin{enumerate}
\item[(i)] for every embedding $\phi:K\to\mathbf C$, with $\sigma\ne\sigma'$ the two embeddings of $M$ above $\phi$, $\pi$ has at $\sigma$ and $\sigma'$ the infinitesimal characters of weights $m$ and $m'$ in $\mathbf Z^2$ with $m_1,m_2,m'_1,m'_2$ pairwise distinct;
\item[(ii)] at every finite place $v$ of $K$ that does not split in $M$, $M$ is unramified over $K$ and $\pi$ has Satake parameters at the places above $v$.
\end{enumerate}
Then there is a cuspidal automorphic representation $\Pi$ of $GL_4(\mathbf A_K)$ with the following properties. For all but finitely many finite places $w$ of $M$, with $v$ the place of $K$ below $w$: if $w'\ne w$ is the other place above $v$, then $\pi$ has Satake parameters $B$ at $w$ and $B'$ at $w'$, and $\Pi$ has Satake parameters $B+B'$ at $v$; if $N(w)=N(v)^2$, then there is a multiset $B'$ such that $\pi$ has Satake parameters $\{b^2:b\in B'\}$ at $w$ and $\Pi$ has Satake parameters $B'+(-B')$ at $v$. For $\phi$, $\sigma$ and $\sigma'$ as in (i), the infinitesimal character of $\Pi$ at $\phi$ is the union of those of $\pi$ at $\sigma$ and at $\sigma'$. Hypothesis (ii) is not used: the published theorem \cite[Chapter~3, Theorem~4.2(e) and \S6]{AC89} gives $\Pi$ under (i) alone.
\end{enumerate}
\end{lemma}
\leanref{Development/AutomorphicArithmetic/AutomorphicVarieties/ModularityArgument/BaseChangeUnits/QuadraticInductionGL2Alt.lean}{UpperTheory.V.Z.Bc.UpQuadI}
\leanref{Development/AutomorphicArithmetic/AutomorphicVarieties/ModularityArgument/BaseChangeUnits/QuadraticInductionGL2Alt.lean}{UpperTheory.V.Z.Bc.UpSolI}
\leanref{Development/AutomorphicArithmetic/AutomorphicVarieties/ModularityArgument/InductionEstimates/IwahoriHeckeInjectivity/HeckeInj.lean}{UpperTheory.V.Z.InvariantSubmoduleMonicPolynomial.E.IndCs}
\leanref{Development/AutomorphicArithmetic/AutomorphicVarieties/ModularityArgument/BaseChangeUnits/QuadraticInductionGL2Alt.lean}{UpperTheory.V.Z.Bc.IndSys}
\leanref{Development/AutomorphicArithmetic/AutomorphicVarieties/ModularityArgument/BaseChangeUnits/QuadraticInductionGL2Alt.lean}{UpperTheory.V.Z.Bc.MixWtC}
\leanref{Development/AutomorphicArithmetic/AutomorphicVarieties/ModularityArgument/InductionEstimates/IwahoriHeckeInjectivity/HeckeInj.lean}{UpperTheory.V.Z.InvariantSubmoduleMonicPolynomial.E.SplitBad}

\begin{lemma}\label{sha-siegel-step-seeds-induced-automorphic-induction}
Let $\ell$ be a prime, $\iota:\bar{\mathbf Q}_\ell\cong\mathbf C$, $K\subset M$ totally complex number fields with $[M:K]=2$, and $\pi$ a cuspidal automorphic representation of $GL_2(\mathbf A_M)$ that satisfies (i) and (ii) of Lemma~\ref{sha-siegel-step-seeds-induced-automorphic-published}(3). Let $r:\operatorname{Gal}(\bar{\mathbf Q}/M)\to GL_2(\bar{\mathbf Q}_\ell)$ be attached through $\iota$ to $\pi$, and $R:\operatorname{Gal}(\bar{\mathbf Q}/K)\to GL_4(\bar{\mathbf Q}_\ell)$ a semisimple representation such that for every $s\in\operatorname{Gal}(\bar{\mathbf Q}/K)$ not in $\operatorname{Gal}(\bar{\mathbf Q}/M)$:
\begin{enumerate}
\item[(a)] for $g\in\operatorname{Gal}(\bar{\mathbf Q}/M)$, the characteristic polynomial of $R(g)$ is that of $r(g)$ times that of $r(s^{-1}gs)$;
\item[(b)] the characteristic polynomial of $R(s)$ is $f(X^2)$, where $f$ is the characteristic polynomial of $r(s^2)$.
\end{enumerate}
Then $R$ is attached through $\iota$ to a regular algebraic cuspidal automorphic representation of $GL_4(\mathbf A_K)$.
\end{lemma}
\leanref{Development/AutomorphicArithmetic/AutomorphicVarieties/ModularityArgument/ThresholdArguments/DenseThresholdCases/EGIndAutCs.lean}{UpperTheory.V.Z.InvariantSubmoduleMonicPolynomial.E.indAutCs_of_alt}
\leanref{Development/AutomorphicArithmetic/AutomorphicVarieties/ModularityArgument/BaseChangeUnits/IndGlue.lean}{UpperTheory.V.Z.Bc.attInd}
\begin{proof}
\begin{enumerate}
\item[(1)] \emph{The induced representation.} Lemma~\ref{sha-siegel-step-seeds-induced-automorphic-published}(3) gives a cuspidal automorphic representation $\Pi$ of $GL_4(\mathbf A_K)$ with the properties stated there.
\item[(2)] \emph{Regularity.} Let $\phi$, $\sigma$, $\sigma'$, $m$ and $m'$ be as in (i). The infinitesimal characters of $\pi$ at $\sigma$ and $\sigma'$ are $\{m_1+\frac12,m_2+\frac12\}$ and $\{m'_1+\frac12,m'_2+\frac12\}$, so that of $\Pi$ at $\phi$ is $\{n+\frac32:n=m_1-1,m_2-1,m'_1-1,m'_2-1\}$, with four distinct elements. So $\Pi$ is regular algebraic.
\item[(3)] \emph{The twist by $\chi_{\mathrm{cyc}}^{-1}$.} The representation $R\otimes\chi_{\mathrm{cyc}}^{-1}$ is semisimple, as $R$ is, and it is attached through $\iota$ to $\Pi$. This is a comparison of Frobenius eigenvalues at the places of $K$ not above $\ell$ where $r$, $R$ and $\Pi$ are unramified and the relations of (1) hold. At such a place $v$ that splits in $M$, into $w\ne w'$, a Frobenius element at a place above $w$ lies in $\operatorname{Gal}(\bar{\mathbf Q}/M)$, and its conjugate by an $s$ as in (a) is a Frobenius element above $w'$; by (a), as $N(w)=N(w')=N(v)$, the Frobenius eigenvalues of $R\otimes\chi_{\mathrm{cyc}}^{-1}$ at $v$ are $N(v)^{3/2}$ times the Satake parameters $B+B'$ of $\Pi$ at $v$. At such a place that does not split, $N(w)=N(v)^2$ for the place $w$ above it, a Frobenius element at $v$ lies outside $\operatorname{Gal}(\bar{\mathbf Q}/M)$ and its square is a Frobenius element at $w$; by (b), the Frobenius eigenvalues of $R$ at $v$ are the square roots of those of $r$ at $w$, which are $N(v)b^2$ for $b\in B'$, and those of $R\otimes\chi_{\mathrm{cyc}}^{-1}$ are $N(v)^{3/2}$ times $B'+(-B')$.
\item[(4)] \emph{Conclusion.} The character $\chi_{\mathrm{cyc}}$ is attached through $\iota$ to an algebraic Hecke character $\chi$ of $K$ \cite{Ser68}. So $R=(R\otimes\chi_{\mathrm{cyc}}^{-1})\otimes\chi_{\mathrm{cyc}}$ is attached through $\iota$ to $\Pi\otimes(\chi\circ\det)$, which is regular algebraic and cuspidal by (2).
\end{enumerate}
\end{proof}
% indAutCs_of_alt (EGIndAutCs.lean, about 30 lines), followed step by step: (1) is IndCs (h K M h2 pi hm hb); (2) is CG.isRegAlg_iff with infOf_add; (3) is Z.Bc.attInd (IndGlue.lean, 374 lines: exists_quadratic_sign_character, exists_moved, the private loc for the comparison at a place, algSat_four, frobenius_eigenvalues_twist and charEig_cyc_inv), applied to the Satake parameters of pi and Pi, which are unique; (4) is tw_back with Sd.twist_conjugate_selfDual and alg_cyclotomic_of_printed_results (the cyclotomic character is algebraic, from Hypotheses').

\begin{lemma}\label{sha-siegel-step-seeds-induced-automorphic-windows-plain}
Let $E$ be a totally complex number field, $n\ge1$, $\pi$ a cuspidal automorphic representation of $GL_n(\mathbf A_E)$, $w$ a finite place of $E$, with uniformizer $\varpi_w$ and residue field of order $q_w$, and $t\in\mathbf R$. Suppose that the supercuspidal support of every local component of $\pi$ at $w$ consists of characters $\xi$ of $E_w^\times$ with $q_w^{t-1/2}<|\xi(\varpi_w)|<q_w^{t+1/2}$.
\begin{enumerate}
\item Let $A$ be a valuation ring over $w$, as in Situation~\ref{gal-local-situation}, and $\operatorname{rc}$ a map that satisfies the laws (Definition~\ref{gal-llc-laws}). Then every local component $\sigma$ of $\pi$ at $w$ is plain (Definition~\ref{gal-plain}), and $\operatorname{rc}(\sigma\nu^{(1-n)/2})$ is abelian and monodromy-free (Definition~\ref{gal-wd-classes}).
\item If the characters in these supercuspidal supports are unramified, then $\pi$ has Satake parameters at $w$.
\end{enumerate}
\end{lemma}
\leanref{Development/AutomorphicArithmetic/AutomorphicVarieties/ModularityArgument/AutomorphyLifting/PolarizedLinkage/PlainRec.lean}{UpperTheory.V.Z.Win.plainRec}
\leanref{Development/AutomorphicArithmetic/AutomorphicVarieties/ModularityArgument/LeviSmoothLocalization/PrincipalSeriesMatching/WinUnl.lean}{UpperTheory.V.Z.ImaginaryQuadraticEtaCharacter.winUnl}
\leanref{Development/AutomorphicArithmetic/AutomorphicVarieties/ModularityArgument/LeviSmoothLocalization/PrincipalSeriesMatching/WinSph.lean}{UpperTheory.V.Z.ImaginaryQuadraticEtaCharacter.winSph}
\begin{proof}
Let $\sigma$ be a local component of $\pi$ at $w$, with supercuspidal support $\{\xi_1,\dots,\xi_n\}$.
\begin{enumerate}
\item[(a)] \emph{The support is unlinked.} If $\xi_i=\xi_j\nu$, then $|\xi_i(\varpi_w)|=q_w^{-1}|\xi_j(\varpi_w)|$. Both values lie in the open interval from $q_w^{t-1/2}$ to $q_w^{t+1/2}$, whose endpoints have ratio $q_w$, so this is impossible. So $\sigma$ is plain.
\item[(b)] \emph{The Weil-Deligne representation.} As no two of the $\xi_j$ differ by $\nu^{\pm1}$, $\sigma$ is isomorphic to the product $\xi_1\times\dots\times\xi_n$, which is irreducible (Theorem~\ref{lit-bundle-unitary}(2) and~(5)) and is the Langlands quotient of the segments $(\xi_j,1)$. By laws (3) and (2), $\operatorname{rc}(\sigma)=\bigoplus_j\xi_j\circ a$, for the local Artin map $a$ of the laws, and by law (5), $\operatorname{rc}(\sigma\nu^{(1-n)/2})=\bigoplus_j(\xi_j|\cdot|^{(1-n)/2})\circ a$. This is a sum of characters with monodromy $0$, which is (1).
\item[(c)] \emph{Satake parameters.} Suppose that the $\xi_j$ are unramified, and let $\alpha_j=\xi_j(\varpi_w)$. By (b), $\sigma$ is isomorphic to $\xi_1\times\dots\times\xi_n$, the representation induced from an unramified character of the upper triangular Borel subgroup. As in the proof of Lemma~\ref{sha-hecke-ideal-spherical}, it has nonzero vectors fixed by $GL_n(\mathcal O_w)$, by the Iwasawa decomposition, and its Satake parameters are $\{\alpha_1,\dots,\alpha_n\}$. This holds for every local component at $w$, which is (2).
\end{enumerate}
\end{proof}
% plainRec (PlainRec.lean, 60 lines): plainIff with winUnl for (a) (WinUnl.lean, quot_ne); plainTor_exists_chars, unlFull (the product of the characters is irreducible), isL_chars and the laws langlands, one and nu for (b), with tor_sum and monodromy_free_sum_list. winSph (WinSph.lean, 15 lines): suppFacts, exists_family (the unramified characters unrChar alpha_i), winUnl and unlFull, then hasSatGL_prinSeries for (c).

\begin{lemma}\label{sha-siegel-step-seeds-induced-automorphic-windows-twist}
Let $E$ be a totally complex number field. For a finite place $u$ of $E$ let $\varpi_u$ be a uniformizer and $q_u$ the order of the residue field. For $t\in\mathbf R$, a cuspidal automorphic representation $\pi$ of $GL_m(\mathbf A_E)$ satisfies the \emph{condition on supercuspidal supports with $t$} at a finite place $u$ if the supercuspidal support of every local component of $\pi$ at $u$ consists of characters $\xi$ of $E_u^\times$ with $q_u^{t-1/2}<|\xi(\varpi_u)|<q_u^{t+1/2}$; it satisfies the \emph{condition on supercuspidal supports} if for some $t$ it does so at every finite place.
\begin{enumerate}
\item Let $\chi$ be a Hecke character of $E$. There is $a\in\mathbf R$ such that for every $m\ge1$, cuspidal automorphic representation $\pi$ of $GL_m(\mathbf A_E)$, $t\in\mathbf R$ and finite place $u$: if $\pi$ satisfies the condition with $t$ at $u$, then $\pi\otimes(\chi\circ\det)$ satisfies it with $t+a$ at $u$.
\item Let $\ell$ be a prime, $\iota:\bar{\mathbf Q}_\ell\cong\mathbf C$, $n\ge2$ and $r:\operatorname{Gal}(\bar{\mathbf Q}/E)\to GL_n(\bar{\mathbf Q}_\ell)$ a homomorphism attached through $\iota$ to a cuspidal automorphic representation of $GL_n(\mathbf A_E)$ that satisfies the condition on supercuspidal supports. Then every regular algebraic cuspidal automorphic representation of $GL_n(\mathbf A_E)$ to which $r$ is attached through $\iota$ satisfies it.
\item Let $\ell$, $\iota$, $n\ge2$ and $r$ be as in (2), with $r$ attached through $\iota$ to a regular algebraic cuspidal automorphic representation of $GL_n(\mathbf A_E)$, and suppose that every such representation satisfies the condition on supercuspidal supports. Let $\kappa:\operatorname{Gal}(\bar{\mathbf Q}/E)\to\bar{\mathbf Q}_\ell^\times$ be a continuous character attached through $\iota$ to an algebraic Hecke character of $E$. Then every regular algebraic cuspidal automorphic representation of $GL_n(\mathbf A_E)$ to which $r\otimes\kappa$ is attached through $\iota$ satisfies the condition on supercuspidal supports.
\end{enumerate}
\end{lemma}
\leanref{Development/AutomorphicArithmetic/AutomorphicVarieties/ModularityArgument/LeviSmoothLocalization/PrincipalSeriesMatching/WindowMul.lean}{UpperTheory.V.Z.ImaginaryQuadraticEtaCharacter.winTw}
\leanref{Development/AutomorphicArithmetic/AutomorphicVarieties/ModularityArgument/AutomorphyLifting/PrimeIndexConstruction/AwOne.lean}{UpperTheory.V.Z.Win.Birth.awOne}
\leanref{Development/AutomorphicArithmetic/AutomorphicVarieties/ModularityArgument/AutomorphyLifting/PrimeIndexConstruction/AwTw.lean}{UpperTheory.V.Z.Win.Birth.awTw}
\begin{proof}
(1) The absolute value $|\chi|$ is a continuous homomorphism from the idele class group of $E$ to $\mathbf R_{>0}$. It is trivial on the idele classes of norm $1$, which form a compact group \cite[Chapter~VI, \S1]{Neu99}, so $|\chi|=\|\cdot\|^{-a}$ for some $a\in\mathbf R$, where $\|\cdot\|$ is the idelic norm; hence $|\chi_u(\varpi_u)|=q_u^{a}$ at every finite place $u$. The local components of $\pi\otimes(\chi\circ\det)$ at $u$ are the $\sigma\otimes(\chi_u\circ\det)$ for the local components $\sigma$ of $\pi$ at $u$, and the supercuspidal support of $\sigma\otimes(\chi_u\circ\det)$ consists of the $\xi\chi_u$ for $\xi$ in that of $\sigma$. For such a character $\xi$, $|\xi\chi_u(\varpi_u)|=q_u^a|\xi(\varpi_u)|$, which lies between $q_u^{t+a-1/2}$ and $q_u^{t+a+1/2}$.

(2) Let $\pi$ be the given representation, satisfying the condition with $t$, and $\pi'$ a regular algebraic cuspidal automorphic representation to which $r$ is attached. At all but finitely many finite places $u$, $\pi$ and $\pi'$ have Satake parameters $B$ and $B'$, and $r$ has Frobenius eigenvalues $q_u^{(n-1)/2}B$ and $q_u^{(n-1)/2}B'$ at $u$ (Definition~\ref{gal-weakly-attached}); so $B=B'$. By Theorem~\ref{lit-jacquet-shalika}(1), $\pi'\cong\pi$, so $\pi'$ has the local components of $\pi$ and satisfies the condition with $t$.

(3) Let $\pi$ be a regular algebraic cuspidal automorphic representation to which $r$ is attached, satisfying the condition with $t$, and $\chi$ the algebraic Hecke character to which $\kappa$ is attached. At all but finitely many places $u$, $\chi$ is unramified, $\pi$ has Satake parameters $B$, $r$ has Frobenius eigenvalues $q_u^{(n-1)/2}B$ and $\kappa$ has Frobenius eigenvalue $\chi_u(\varpi_u)$ at $u$; then $r\otimes\kappa$ has Frobenius eigenvalues $q_u^{(n-1)/2}\chi_u(\varpi_u)B$, and $\pi\otimes(\chi\circ\det)$ has Satake parameters $\chi_u(\varpi_u)B$ at $u$. As $r$ is semisimple, so is $r\otimes\kappa$. So $r\otimes\kappa$ is attached through $\iota$ to $\pi\otimes(\chi\circ\det)$, which satisfies the condition with $t+a$ by (1). By (2), applied to $r\otimes\kappa$, every regular algebraic cuspidal automorphic representation to which $r\otimes\kappa$ is attached satisfies the condition.
\end{proof}
% winTw (WindowMul.lean, 25 lines): charAbs (the exponent of the absolute value of chi, through HeckeNormPow.toQC), suppTw (the support of the twist) and window_mul. awOne (AwOne.lean, 10 lines): hecke_eigenvalue_basic_properties for B = B', and suppSmo (SuppSmo.lean: strongMultOneC, Theorem lit-jacquet-shalika (1)). awTw (AwTw.lean, 15 lines): attTwist (the twist is attached to pi twisted by chi, with sat_twist and CharEig of kappa), winTw and awOne. The Lean states (2) and (3) for every n >= 2 (AllC) and r without continuity.

\begin{lemma}\label{sha-siegel-step-seeds-induced-automorphic-windows-up}
Assume the Galois representations as in Theorem~\ref{sha-residually-reducible-automorphy}. Let $F\subseteq E$ be totally complex number fields, with $E$ Galois over $F$ with solvable Galois group, $\ell$ a prime, $\iota:\bar{\mathbf Q}_\ell\cong\mathbf C$, $n\ge2$, and $r:\operatorname{Gal}(\bar{\mathbf Q}/F)\to GL_n(\bar{\mathbf Q}_\ell)$ a continuous representation that is attached through $\iota$ to a regular algebraic cuspidal automorphic representation of $GL_n(\mathbf A_F)$ and irreducible on $\operatorname{Gal}(\bar{\mathbf Q}/E)$. Suppose that every regular algebraic cuspidal automorphic representation of $GL_n(\mathbf A_F)$ to which $r$ is attached through $\iota$ satisfies the condition on supercuspidal supports (Lemma~\ref{sha-siegel-step-seeds-induced-automorphic-windows-twist}). Then every regular algebraic cuspidal automorphic representation of $GL_n(\mathbf A_E)$ to which the restriction $r_E$ of $r$ to $\operatorname{Gal}(\bar{\mathbf Q}/E)$ is attached through $\iota$ satisfies it.
\end{lemma}
\leanref{Development/AutomorphicArithmetic/AutomorphicVarieties/ModularityArgument/AutomorphyLifting/PrimeIndexConstruction/ExistsPrimeIndex.lean}{UpperTheory.V.Z.Win.Birth.awUp}
\leanref{Development/AutomorphicArithmetic/AutomorphicVarieties/ModularityArgument/ThresholdArguments/DenseThresholdCases/BcUpCyc.lean}{UpperTheory.V.Z.Bc.upCyc_of_alt}
\begin{proof}
The proof is by induction on $d=[E:F]$, for all fields $F\subseteq E$ and all $n$ and $r$ at once. If $d=1$, then $E=F$ and there is nothing to prove. Let $d>1$.
\begin{enumerate}
\item[(a)] \emph{A field of prime degree.} The group $\operatorname{Gal}(E/F)$ is solvable and nontrivial, so its commutator subgroup is proper; a maximal proper subgroup $H$ containing it is normal, and $\operatorname{Gal}(E/F)/H$, an abelian group with no subgroups other than $1$ and itself, is cyclic of prime order $p$. Let $K$ be the fixed field of $H$. Then $K$ is totally complex and Galois over $F$ of degree $p$, $E$ is Galois over $K$ with the solvable Galois group $H$, and $[E:K]<d$. As $r$ is irreducible on $\operatorname{Gal}(\bar{\mathbf Q}/E)$, it is irreducible on $\operatorname{Gal}(\bar{\mathbf Q}/K)$.
\item[(b)] \emph{Automorphy over $K$.} Let $\pi$ be a regular algebraic cuspidal automorphic representation of $GL_n(\mathbf A_F)$ to which $r$ is attached, and $\eta$ a Hecke character of $F$ of order $p$ trivial on the norms from $K$ \cite[Chapter~VI]{Neu99}. Suppose that $\pi\cong\pi\otimes(\eta\circ\det)$. Let $\delta$ be the character of $\operatorname{Gal}(\bar{\mathbf Q}/F)$ of order $p$, trivial on $\operatorname{Gal}(\bar{\mathbf Q}/K)$, that corresponds to $\eta$. At all but finitely many places the Satake parameters of $\pi$ are stable under multiplication by the value of $\eta$ at a uniformizer, so $r$ and $r\otimes\delta$ have the same characteristic polynomials at Frobenius elements, hence everywhere by the Chebotarev density theorem (Theorem~\ref{lit-chebotarev-degree-one}), and $r\cong r\otimes\delta$ by the Brauer-Nesbitt theorem \cite[(30.16)]{CR62}. An isomorphism $Q$, with $Qr(g)=\delta(g)r(g)Q$ for all $g$, commutes with $r$ on $\operatorname{Gal}(\bar{\mathbf Q}/K)$, on which $r$ is irreducible; so $Q$ is a scalar, and $\delta=1$, which is false. So $\pi\not\cong\pi\otimes(\eta\circ\det)$, and by \cite[Chapter~3, Theorems~4.2(a) and~5.1]{AC89} $\pi$ has a cuspidal base change $P$ to $GL_n(\mathbf A_K)$: at a place $w$ of $K$ of residue degree $f$ over a place $v$ at which $\pi$ has Satake parameters $B$ and $K$ is unramified, $P$ has Satake parameters $\{b^f:b\in B\}$, and the infinitesimal character of $P$ at an embedding $\tau:K\to\mathbf C$ is that of $\pi$ at the restriction of $\tau$ to $F$. So $P$ is regular algebraic, and the restriction $r_K$ of $r$ to $\operatorname{Gal}(\bar{\mathbf Q}/K)$, which is irreducible by (a), is attached to $P$ through $\iota$, since a Frobenius element at $w$ acts as the $f$-th power of one at $v$.
\item[(c)] \emph{The condition over $K$.} Let $\pi$ be as in (b), satisfying the condition with $t$, and $P$ a regular algebraic cuspidal automorphic representation of $GL_n(\mathbf A_K)$ to which $r_K$ is attached. As in step (p4) of the proof of Lemma~\ref{sha-rigidity-windowed-up}, $P$ is a weak base change of $\pi$ (Definition~\ref{gal-depth-zero-lift-weak}). By Lemma~\ref{sha-rigidity-windowed-local}, at every finite place $w$ of $K$ the supercuspidal support of every local component of $P$ at $w$ consists of characters $\zeta$ with $q_w^{t-1/2}<|\zeta(\varpi_w)|<q_w^{t+1/2}$. By Lemma~\ref{sha-siegel-step-seeds-induced-automorphic-windows-twist}(2), every regular algebraic cuspidal automorphic representation to which $r_K$ is attached satisfies the condition on supercuspidal supports.
\item[(d)] \emph{Induction.} By (a) to (c), $K\subseteq E$ and $r_K$ satisfy the hypotheses, and $[E:K]<d$. The statement for $K\subseteq E$ gives the conclusion for $r_E$, the restriction of $r_K$.
\end{enumerate}
\end{proof}
% awUp (ExistsPrimeIndex.lean, 220 lines): a strong induction on [E:F]. Degree one is allW_one (win_one and DegOneLift.isLift: the Lean transports windows along the isomorphism of F with E). (a) is the private exists_prime_index (a coatom containing the commutator subgroup, with Group.IsSolvable.commutator_lt_top_of_nontrivial) and IsGalois.of_fixedField_normal_subgroup; (b) is the hypothesis UpCyc, discharged by Z.Bc.upCyc_of_alt (BcUpCyc.lean: Z.Bc.cycC_holds for the alternative of a self-twist or a weak base change with the infinitesimal characters, selfTw and schur_twist_holds against the self-twist, frobenius_power_compatibility for the attachment); (c) is allW_step (attBC for the weak base change, winUpDown, whose up half is Lemma sha-rigidity-windowed-local, and awOne); (d) is the induction hypothesis with resFacts.

\begin{lemma}\label{sha-siegel-step-seeds-induced-automorphic-windows-real}
Assume the Galois representations as in Theorem~\ref{sha-residually-reducible-automorphy}. Let $F$ be a totally real number field, $E\supset F$ a CM field with $[E:F]=2$, $\ell$ a prime, $\iota:\bar{\mathbf Q}_\ell\cong\mathbf C$, and $r:\operatorname{Gal}(\bar{\mathbf Q}/F)\to GL_2(\bar{\mathbf Q}_\ell)$ a continuous representation, irreducible on $\operatorname{Gal}(\bar{\mathbf Q}/E)$ and attached through $\iota$ to a regular algebraic cuspidal automorphic representation $\pi$ of $GL_2(\mathbf A_F)$ that has, at every finite place, Satake parameters or only plain local components (Definition~\ref{gal-plain}). Then every regular algebraic cuspidal automorphic representation of $GL_2(\mathbf A_E)$ to which the restriction of $r$ to $\operatorname{Gal}(\bar{\mathbf Q}/E)$ is attached through $\iota$ satisfies the condition on supercuspidal supports (Lemma~\ref{sha-siegel-step-seeds-induced-automorphic-windows-twist}).
\end{lemma}
\leanref{Development/AutomorphicArithmetic/AutomorphicVarieties/ModularityArgument/AutomorphyLifting/PrimeIndexConstruction/AwReal.lean}{UpperTheory.V.Z.Win.Birth.awReal}
\begin{proof}
Let $P$ be such a representation. By Lemma~\ref{sha-hecke-ideal-plain}(1), there is $t\in\mathbf R$ such that at every finite place $w$ of $E$ at which $P$ is plain, every character $\zeta$ in the supercuspidal support of a local component of $P$ at $w$ satisfies $q_w^{t-1/2}<|\zeta(\varpi_w)|<q_w^{t+1/2}$; and at such a place these supports consist of characters, by Definition~\ref{gal-plain}. So it suffices that $P$ is plain at every finite place.
\openstep{That $P$ is plain at every finite place $w$ of $E$ is not yet written out. With $v$ the place of $F$ below $w$, Alpöge chooses a CM quadratic extension $E'$ of $F$ in which $v$ splits completely, with $r$ irreducible on the Galois group of $M=EE'$; takes the quadratic base change $P'$ of $\pi$ to $E'$, whose local components at the places above $v$ are those of $\pi$, so that $P'$ is plain there; and passes the bound of Lemma~\ref{sha-hecke-ideal-plain}(1) for $P'$ up to a base change $Q$ of $P$ and $P'$ to $M$, with the characters coming through the norm from $F_v$, and down from $M$ to $E$ at $w$.}
% awReal (AwReal.lean, 237 lines): winPt with winOf (Lemma sha-hecke-ideal-plain (1)) is the reduction above; the open part is the private plainUp: splitCM (E' and M = EE', with the place split and r irreducible on G_M), ConvolutionTheory.R2.quadBC2 with noTw (baseChange_noSelfTwist_of_frobenius_sign_irreducible, through isoFrob, frobenius_sign_compatibility and schur_twist_holds), Bc.bcAtt, upCycle (Z.Bc.cycC_holds and selfTw) for Q, attBC, splitTr and plainOf for P' at the split places, winUpDown up from E' to M and thruTr for the characters through the norm, and winUpDown down from M to E (the down half is the open part of Lemma sha-siegel-step-ideal-descent-updown-local (3)), then plainOf. Size: lean_size of plainUp, stopping at the book's citations, 185,000 lines, of which about 10,400 in AutomorphicArithmetic/AutomorphicVarieties and 26,400 in AutomorphicArithmetic/Foundations; the rest is mostly the local and global theory of GL_n in GroundTheory/AutomorphicForms (38,700), AbelianVarietyTheory/Proofs (34,300) and AutomorphicSpectrum/Proofs (24,300).
\end{proof}

\begin{lemma}\label{sha-siegel-step-seeds-induced-automorphic-windows-induction-local}
Let $K\subset M$ be totally complex number fields with $[M:K]=2$, $Q$ a cuspidal automorphic representation of $GL_4(\mathbf A_K)$ and $\pi$ one of $GL_2(\mathbf A_M)$ such that for all but finitely many finite places $w$ of $M$, with $v$ the place of $K$ below $w$: if $w'\ne w$ is the other place above $v$, then $\pi$ has Satake parameters $B$ at $w$ and $B'$ at $w'$, and $Q$ has Satake parameters $B+B'$ at $v$; and if $N(w)=N(v)^2$, then there is a multiset $B'$ such that $\pi$ has Satake parameters $\{b^2:b\in B'\}$ at $w$ and $Q$ has Satake parameters $B'+(-B')$ at $v$. For a finite place $u$ let $\varpi_u$ be a uniformizer and $q_u$ the order of the residue field. Let $t\in\mathbf R$ and $v$ a finite place of $K$ such that at every place $w$ of $M$ above $v$ the supercuspidal support of every local component of $\pi$ at $w$ consists of characters $\zeta$ with $q_w^{t-1/2}<|\zeta(\varpi_w)|<q_w^{t+1/2}$, for each of which there is a smooth character $\chi'$ of $K_v^\times$ such that $\zeta\cdot(\chi'\circ N_{M_w/K_v})$ is unramified. Then the supercuspidal support of every local component of $Q$ at $v$ consists of characters $\xi$ with $q_v^{t-1/2}<|\xi(\varpi_v)|<q_v^{t+1/2}$.
\end{lemma}
\leanref{Development/AutomorphicArithmetic/AutomorphicVarieties/ModularityArgument/LeviSmoothLocalization/WinUpDown.lean}{UpperTheory.V.Z.ImaginaryQuadraticEtaCharacter.WinInd}
\leanref{Development/AutomorphicArithmetic/AutomorphicVarieties/ModularityArgument/LeviSmoothLocalization/PrincipalSeriesMatching/WinInd.lean}{UpperTheory.V.Z.ImaginaryQuadraticEtaCharacter.winInd}
\begin{proof}
\begin{enumerate}
\item[(a)] \emph{The local supports.} The supercuspidal support of every local component of $Q$ at $v$ consists of characters, and for each such character $\xi$ there are a place $w$ of $M$ above $v$ such that $\xi\circ N_{M_w/K_v}$ lies in the supercuspidal support of a local component of $\pi$ at $w$.
\openstep{Step (a) is not yet written out: that the supports of $Q$ at $v$ consist of characters, and that each comes through the norm from a character in a support of $\pi$ above $v$, which Alpöge proves by comparing gamma factors of the local components of $\pi$ and $Q$ twisted by characters.}
% (a) is indTor (TorAt for P from TorAt and ThruAt above v) and indX (the norm of each character of the support), through matchIndOf, the induced gamma factors and the local supports in GammaLaurentFactors and LeviSmoothLocalization. Size: lean_size of indTor and indX, stopping at the book's citations, 127,000 lines, of which about 8,300 in AutomorphicArithmetic/AutomorphicVarieties (the gamma factors); the rest is mostly the local theory of GL_n in GroundTheory/AutomorphicForms (51,900) and AbelianVarietyTheory/Proofs (36,900).
\item[(b)] \emph{The bound.} Let $\xi$ and $w$ be as in (a), and $f$ the residue degree of $w$ over $v$. Then $q_w=q_v^f$, and $N_{M_w/K_v}(\varpi_w)=u\varpi_v^f$ for a unit $u$ of $K_v$; the character $\xi$ is continuous on the compact group of units, so $|\xi(u)|=1$, and $|\xi\circ N_{M_w/K_v}(\varpi_w)|=|\xi(\varpi_v)|^f$. The bound for $\xi\circ N_{M_w/K_v}$ at $w$ is $q_v^{f(t-1/2)}<|\xi(\varpi_v)|^f<q_v^{f(t+1/2)}$, and taking $f$-th roots gives the bound for $\xi$ at $v$.
\end{enumerate}
\end{proof}
% winInd (WinInd.lean, 33 lines): indTor and indX for (a); normVal and the private rootBounds for (b).

\begin{lemma}\label{sha-siegel-step-seeds-induced-automorphic-windows}
Assume the Galois representations as in Theorem~\ref{sha-residually-reducible-automorphy}. Let $\ell$ be a prime and $\iota:\bar{\mathbf Q}_\ell\cong\mathbf C$. For a number field $E$ write $G_E=\operatorname{Gal}(\bar{\mathbf Q}/E)$. A cuspidal automorphic representation $\Pi$ of $GL_n(\mathbf A_E)$, for a totally complex field $E$, satisfies the \emph{condition on supercuspidal supports} if there is $t\in\mathbf R$ such that at every finite place $u$ of $E$ the supercuspidal support of $\Pi_u$ consists of characters $\xi$ of $E_u^\times$ with $q_u^{t-1/2}<|\xi(\varpi_u)|<q_u^{t+1/2}$, where $\varpi_u$ is a uniformizer and $q_u$ the order of the residue field.
\begin{enumerate}
\item Let $F$ be a totally real number field, $E\supset F$ a CM quadratic extension, and $M\supseteq E$ a totally complex field, Galois over $E$ with solvable Galois group. Let $r:G_F\to GL_2(\bar{\mathbf Q}_\ell)$ be a continuous representation, irreducible on $G_M$ and attached through $\iota$ to a regular algebraic cuspidal automorphic representation of $GL_2(\mathbf A_F)$ that has, at every finite place, Satake parameters or only plain local components (Definition~\ref{gal-plain}), and let $\kappa:G_E\to\bar{\mathbf Q}_\ell^\times$ be a continuous character attached through $\iota$ to an algebraic Hecke character of $E$. Then every regular algebraic cuspidal automorphic representation of $GL_2(\mathbf A_M)$ to which the restriction to $G_M$ of $r_E\otimes\kappa$, with $r_E$ the restriction of $r$ to $G_E$, is attached through $\iota$ satisfies the condition on supercuspidal supports.
\item Let $M$ be a CM field, $n\ge1$, and $r:G_M\to GL_n(\bar{\mathbf Q}_\ell)$ a continuous representation. Suppose that there are a totally real subfield $F_0$ of $M$, a character $\mu$ of $G_{F_0}$ attached through $\iota$ to an algebraic Hecke character of $F_0$ and taking the same value at all complex conjugations, a complex conjugation $c\in G_{F_0}$ and $P\in GL_n(\bar{\mathbf Q}_\ell)$ with ${}^tr(s)\,P\,r(csc^{-1})=\mu(s)P$ for $s\in G_M$. Let $\Pi$ be a regular algebraic cuspidal automorphic representation of $GL_n(\mathbf A_M)$ to which $r$ is attached through $\iota$, and suppose that every such representation satisfies the condition on supercuspidal supports. Let $w$ be a finite place of $M$ not above $\ell$ such that $r$ is trivial on the inertia group of $G_M$ at some place of $\bar{\mathbf Q}$ above $w$. Then $\Pi$ has Satake parameters at $w$.
\item Let $K\subset M$ be totally complex number fields with $[M:K]=2$, and $\pi$ a cuspidal automorphic representation of $GL_2(\mathbf A_M)$ that satisfies the condition on supercuspidal supports and (ii) of Lemma~\ref{sha-siegel-step-seeds-induced-automorphic-published}(3). Let $r:G_M\to GL_2(\bar{\mathbf Q}_\ell)$ be attached through $\iota$ to $\pi$, and $R:G_K\to GL_4(\bar{\mathbf Q}_\ell)$ a representation that satisfies (a) and (b) of Lemma~\ref{sha-siegel-step-seeds-induced-automorphic-induction}. Then every regular algebraic cuspidal automorphic representation of $GL_4(\mathbf A_K)$ to which $R$ is attached through $\iota$ satisfies the condition on supercuspidal supports.
\end{enumerate}
\end{lemma}
\leanref{Development/AutomorphicArithmetic/AutomorphicVarieties/ModularityArgument/AutomorphyLifting/PrimeIndexConstruction/RankFourAutomorphyOfTensor.lean}{UpperTheory.V.Z.Win.Birth.AwChain}
\leanref{Development/AutomorphicArithmetic/AutomorphicVarieties/ModularityArgument/ThresholdArguments/DenseThresholdCases/BirthAwChain.lean}{UpperTheory.V.Z.Win.Birth.awChain_of_printsH}
\leanref{Development/AutomorphicArithmetic/AutomorphicVarieties/ModularityArgument/AutomorphyLifting/PrimeIndexConstruction/RankFourAutomorphyOfTensor.lean}{UpperTheory.V.Z.Win.Birth.SphOfUnrW}
\leanref{Development/AutomorphicArithmetic/AutomorphicVarieties/ModularityArgument/ThresholdArguments/DenseThresholdCases/BirthSphOfUnrW.lean}{UpperTheory.V.Z.Win.Birth.sphOfUnrW_alt}
\leanref{Development/AutomorphicArithmetic/AutomorphicVarieties/ModularityArgument/AutomorphyLifting/PrimeIndexConstruction/ThruBad.lean}{UpperTheory.V.Z.Win.Birth.AwInd}
\leanref{Development/AutomorphicArithmetic/AutomorphicVarieties/ModularityArgument/AutomorphyLifting/PrimeIndexConstruction/ValAbsPowNegOne.lean}{UpperTheory.V.Z.Win.Birth.awInd}
\begin{proof}
(1) Let $r_M$ be the restriction of $r$ to $G_M$, and $\chi$ the algebraic Hecke character of $E$ to which $\kappa$ is attached.
\begin{enumerate}
\item[(a)] \emph{Over $E$.} As $r$ is irreducible on $G_M\subseteq G_E$, it is irreducible on $G_E$. By Lemma~\ref{sha-siegel-step-seeds-induced-automorphic-windows-real}, every regular algebraic cuspidal automorphic representation of $GL_2(\mathbf A_E)$ to which $r_E$ is attached through $\iota$ satisfies the condition on supercuspidal supports.
\item[(b)] \emph{Automorphy over $E$.} By Lemma~\ref{sha-siegel-step-seeds-induced-automorphic-published}(1), $r_E$ is attached through $\iota$ to a regular algebraic cuspidal automorphic representation $P_E$ of $GL_2(\mathbf A_E)$.
\item[(c)] \emph{The twist.} By (a), (b) and Lemma~\ref{sha-siegel-step-seeds-induced-automorphic-windows-twist}(3), applied to $r_E$ and $\kappa$, every regular algebraic cuspidal automorphic representation of $GL_2(\mathbf A_E)$ to which $r_E\otimes\kappa$ is attached through $\iota$ satisfies the condition. As in the proof of that lemma, $r_E\otimes\kappa$ is attached through $\iota$ to $P_E\otimes(\chi\circ\det)$, which is regular algebraic: as $\chi$ is algebraic, its infinitesimal character at each embedding is that of $P_E$ shifted by an integer.
\item[(d)] \emph{Up to $M$.} The restriction of $r_E\otimes\kappa$ to $G_M$ is the twist of $r_M$ by a character, so it is continuous and irreducible. By (c), Lemma~\ref{sha-siegel-step-seeds-induced-automorphic-windows-up}, applied to $E\subseteq M$, $n=2$ and $r_E\otimes\kappa$, gives (1).
\end{enumerate}

(2) Let $t$ be such that $\Pi$ satisfies the condition with $t$, and $A$ a valuation ring of $\bar{\mathbf Q}$ over $w$, as in Situation~\ref{gal-local-situation}, such that $r$ is trivial on $I_A$.
\begin{enumerate}
\item[(a)] \emph{Polarization.} By Lemma~\ref{sha-siegel-step-ideal-descent-updown-local}(1), applied to $r$, $F_0$, $\mu$, $c$, $P$ and $\Pi$, there is a Hecke character $\chi$ of the maximal totally real subfield of $M$ with $(\Pi,\chi)$ polarized (Definition~\ref{gal-polarized}).
\item[(b)] \emph{Local-global compatibility.} By Theorem~\ref{gal-llc-exists} there are a local Artin map $a$ and a map $\operatorname{rc}$ that satisfies the laws with respect to $a$. Let $\sigma$ be a local component of $\Pi$ at $w$; its supercuspidal support consists of characters $\xi_1,\dots,\xi_n$. By Lemma~\ref{sha-siegel-step-seeds-induced-automorphic-windows-plain}(1), $y=\operatorname{rc}(\sigma\nu^{(1-n)/2})$ is abelian and monodromy-free. As $\Pi$ is regular algebraic and polarized by (a), and $r$ is attached to it, Lemma~\ref{sha-hecke-ideal-attached}(1) gives $\mathrm{WD}_\iota(r|_{W_A})=y$, as $w\nmid\ell$.
\item[(c)] \emph{The characters are unramified.} The representation $\sigma$ is a subquotient of $\mathrm{PS}(\omega)$ for $\omega_j=\xi_j\nu^{(1-n)/2}$ (Definition~\ref{sha-hecke-local-ordinary}), so by Lemma~\ref{sha-hecke-local-rec}(1) the characteristic polynomial of $y(s)$ is $\prod_j(X-\xi_j(a(s))|a(s)|^{(1-n)/2})$ for $s\in W_A$. Since $r$ is trivial on $I_A$, $y=\mathrm{WD}_\iota(r|_{W_A})$ is trivial on $I_A$, so this polynomial is $(X-1)^n$ for $s\in I_A$, and $\xi_j(a(s))=1$ for every $j$, as $|a(s)|=1$. By (2) of Definition~\ref{gal-local-artin}, $a$ maps $I_A$, the elements of degree $0$, onto $\mathcal O_w^\times$, since $a$ is surjective. So every $\xi_j$ is unramified.
\item[(d)] \emph{Conclusion.} This holds for every local component of $\Pi$ at $w$, so by Lemma~\ref{sha-siegel-step-seeds-induced-automorphic-windows-plain}(2), $\Pi$ has Satake parameters at $w$.
\end{enumerate}

(3) Let $P$ be a regular algebraic cuspidal automorphic representation of $GL_4(\mathbf A_K)$ to which $R$ is attached through $\iota$, $t$ such that $\pi$ satisfies the condition with $t$, and $Q=P\otimes(\|\cdot\|^{-1}\circ\det)$, where $\|\cdot\|$ is the idelic norm of $K$.
\begin{enumerate}
\item[(a)] \emph{The Satake parameters of $Q$.} Let $v$ be a finite place of $K$, not above $\ell$ and outside a finite set, at which $P$ has Satake parameters $C$ and $R$ has Frobenius eigenvalues $N(v)^{3/2}C$, and at whose places $w$ above it $\pi$ has Satake parameters and $r$ has Frobenius eigenvalues $N(w)^{1/2}$ times them. Then $Q$ has Satake parameters $N(v)C$ at $v$. As in step (3) of the proof of Lemma~\ref{sha-siegel-step-seeds-induced-automorphic-induction}, which uses only (a) and (b) for $R$ and the attachment of $r$ to $\pi$: if $v$ splits in $M$ into $w\ne w'$, where $\pi$ has Satake parameters $B$ and $B'$, a Frobenius element $\phi$ at a place above $w$ lies in $G_M$, and its conjugate by an element of $G_K$ outside $G_M$ is one above $w'$, so by (a) the Frobenius eigenvalues of $R$ at $v$ are $N(v)^{1/2}(B+B')$; if $N(w)=N(v)^2$ for the place $w$ above $v$, where $\pi$ has Satake parameters $\{b^2:b\in B'\}$, a Frobenius element $\phi$ at $v$ lies outside $G_M$ and $\phi^2$ is one at $w$, so by (b) the Frobenius eigenvalues of $R$ at $v$ are the square roots of the $N(v)b^2$, that is, $N(v)^{1/2}(B'+(-B'))$. Comparing with $N(v)^{3/2}C$, $Q$ has Satake parameters $B+B'$, respectively $B'+(-B')$, at $v$: $Q$ and $\pi$ satisfy the relations of Lemma~\ref{sha-siegel-step-seeds-induced-automorphic-windows-induction-local}.
\item[(b)] \emph{The characters come through the norm.} Let $v$ be a finite place of $K$, $w$ a place of $M$ above $v$, and $\zeta$ a character in the supercuspidal support of a local component of $\pi$ at $w$. If $v$ splits in $M$, then $N_{M_w/K_v}$ is an isomorphism, and $\chi'=\zeta^{-1}\circ N_{M_w/K_v}^{-1}$ has $\zeta\cdot(\chi'\circ N_{M_w/K_v})=1$. Otherwise $\pi$ has Satake parameters at $w$ by (ii), so the supercuspidal supports of its local components at $w$ consist of unramified characters, as in the proof of Lemma~\ref{sha-hecke-ideal-plain}(2), and $\chi'=1$ will do.
\item[(c)] \emph{The condition for $Q$.} By (a), (b) and Lemma~\ref{sha-siegel-step-seeds-induced-automorphic-windows-induction-local}, applied to $K\subset M$, $Q$, $\pi$ and $t$ at every finite place $v$ of $K$, $Q$ satisfies the condition with $t$.
\item[(d)] \emph{The condition for $P$.} By Lemma~\ref{sha-siegel-step-seeds-induced-automorphic-windows-twist}(1), applied to the Hecke character $\|\cdot\|$ of $K$, there is $a$ such that $Q\otimes(\|\cdot\|\circ\det)=P$ satisfies the condition with $t+a$.
\end{enumerate}
% (1) is awChain_of_printsH (BirthAwChain.lean, 34 lines), followed step by step: hirrE and awReal (Lemma sha-siegel-step-seeds-induced-automorphic-windows-real) for (a); Sd.up_quad_of_printed_results, quadratic base change from HypothesesH', which the text cites as Lemma sha-siegel-step-seeds-induced-automorphic-published (1), for (b); awTw and Sd.twist_conjugate_selfDual (regAlg_twist, attTwist) for (c); shapeFacts and contTwist, then awUp with Z.Bc.upCyc_of_alt (Lemma sha-siegel-step-seeds-induced-automorphic-windows-up) for (d).
% (2) is sphOfUnrW_alt (BirthSphOfUnrW.lean, 97 lines), followed step by step: the private exists_polarized, a copy of polarized_pair_of_attached (Lemma sha-siegel-step-ideal-descent-updown-local (1)), for (a); recTExists (Theorem gal-llc-exists), plainRec (Lemma sha-siegel-step-seeds-induced-automorphic-windows-plain (1)) and compatTOf (Lemma sha-hecke-ideal-attached) for (b); plainTor_exists_chars, SupportUnique.hasSupport_unique, charPol_of_chars (the Lean's form of Lemma sha-hecke-local-rec (1) used here), charPol_wdOf and the private unram_of_charPol for (c); winSph (Lemma sha-siegel-step-seeds-induced-automorphic-windows-plain (2)) for (d).
% (3) is awInd (ValAbsPowNegOne.lean, 220 lines), followed step by step: indSys_of_att (attInd, the private indSys_indSat, indSat_unique and sq_surj, frobenius_eigenvalues_twist with charEig_cyc, val_absPow_neg_one and sat_tw) for (a); thruBad (ThruBadAlt.lean, 57 lines: thru_split with map_nm, and suppFacts at the places with Satake parameters) for (b); winInd (Lemma sha-siegel-step-seeds-induced-automorphic-windows-induction-local) for (c); winTw with absPow K 1 and suppSmo for (d), where the Lean identifies the twist back with P through strong multiplicity one.
\end{proof}

\begin{lemma}\label{sha-siegel-step-seeds-induced-automorphic}
Assume the Galois representations as in Theorem~\ref{sha-residually-reducible-automorphy}. Let $D$, $T$ and $\iota$ be as in Definition~\ref{sha-siegel-step-seeds-induced-datum}, and let $(k_0,\varphi,j,t,\rho,\alpha,\beta,R,\delta)$ be an induction datum for $T$ and $\iota$ such that $\rho$ is attached through $\iota$ to a regular algebraic cuspidal automorphic representation of $GL_2(\mathbf A_{k_0})$ that has, at every finite place of $k_0$, Satake parameters or only plain local components (Definition~\ref{gal-plain}). Then the restriction of $R$ to $G_L$ is attached through $\iota$ to a regular algebraic cuspidal automorphic representation of $GL_4(\mathbf A_L)$, and every regular algebraic cuspidal automorphic representation $\pi$ of $GL_4(\mathbf A_L)$ to which it is attached through $\iota$ satisfies the condition on supercuspidal supports of Definition~\ref{sha-siegel-tower-ideal}(1): there is $t'\in\mathbf R$ such that at every finite place $u$ of $L$ the supercuspidal support of $\pi_u$ consists of characters $\xi'$ of $L_u^\times$ with $q_u^{t'-1/2}<|\xi'(\varpi_u)|<q_u^{t'+1/2}$.
\end{lemma}
\leanref{Development/AutomorphicArithmetic/AutomorphicVarieties/ModularityArgument/AutomorphyLifting/PrimeIndexConstruction/RankFourAutomorphyOfTensor.lean}{UpperTheory.V.Z.Win.Birth.IndAutLP}
\leanref{Development/AutomorphicArithmetic/AutomorphicVarieties/ModularityArgument/ThresholdArguments/DenseThresholdCases/IndAutLOfSplit.lean}{UpperTheory.V.Z.InvariantSubmoduleMonicPolynomial.E.indAutLP_of_printsH_upQuad_upSol_indCs}
\begin{proof}
Let $\tau=g_0^T$, $E=k_0K$ and $M=k_0L$, and let $\pi_0$ be the cuspidal automorphic representation of $GL_2(\mathbf A_{k_0})$ of (b), to which $\rho$ is attached.
\begin{enumerate}
\item[(1)] \emph{Irreducibility.} By Lemma~\ref{sha-siegel-step-seeds-induced-irreducible}, $R$ is irreducible on $G_L\cap\ker\chi_{\mathrm{cyc}}$, hence on $G_L$.
\item[(2)] \emph{The fields.} By Lemma~\ref{sha-siegel-step-seeds-induced-irreducible-residue}, some $g\in G_L$ has $\bar\chi^{(0,1)}(g)\ne1$, so $g\notin G_{k_0}$ by (a), and $k_0\not\subseteq L$. As $K$ is imaginary and $k_0$ real, $K\not\subseteq k_0$. So $[E:k_0]=2$ and $[M:L]=2$. The fields $E$ and $M$ contain $K$, so they are totally complex; $L=FK$ is Galois over $\mathbf Q$ with abelian Galois group, as $F$ and $K$ are, and so are $M$ and $E\subseteq M$. So $E$ and $M$ are CM fields, $M$ is Galois over $E$ and over $L$ with abelian Galois groups, and $G_M=G_L\cap G_{k_0}$.
\item[(3)] \emph{Regularity.} By (b), $\pi_0$ has the infinitesimal characters of weights $(j+t,-t)$ and $(j,0)$ at the two infinite places of $k_0$; their entries are distinct, as $j,t\ge2$. So $\pi_0$ is regular algebraic.
\item[(4)] \emph{Base change.} By (b), $\rho$ is irreducible on $G_E$ and on $G_M$, which are open in $G_{k_0}$. By Lemma~\ref{sha-siegel-step-seeds-induced-automorphic-published}(1), applied to $k_0\subset E$, $\rho$ and $\pi_0$, the restriction of $\rho$ to $G_E$ is attached through $\iota$ to a regular algebraic cuspidal automorphic representation $P_1$ of $GL_2(\mathbf A_E)$ that has at each embedding of $E$ the infinitesimal character of $\pi_0$ at the infinite place of $k_0$ below it. By part (2) of that lemma, applied to $E\subseteq M$ and $P_1$, the restriction $\rho_M$ of $\rho$ to $G_M$ is attached through $\iota$ to a regular algebraic cuspidal automorphic representation $P$ of $GL_2(\mathbf A_M)$ that has at each embedding of $M$ the infinitesimal character of $\pi_0$ at the infinite place of $k_0$ below it.
\item[(5)] \emph{A polarization.} By (b) and condition (q) of Definition~\ref{sha-siegel-tower-towers}(4), $\det\rho=\chi_{\mathrm{cyc}}^j\xi$ with $\xi$ of finite order, so $\det\rho$ is attached through $\iota$ to an algebraic Hecke character of $k_0$ \cite{Ser68}; it takes the same value at all complex conjugations, as $\chi_{\mathrm{cyc}}$ takes the value $-1$ at each of them and they are conjugate in $G$. Let $c\in G_{k_0}$ be a complex conjugation, which exists as $k_0$ is real, $J=\big(\begin{smallmatrix}0&1\\-1&0\end{smallmatrix}\big)$ and $P'=J\rho(c)^{-1}$. As ${}^tAJA=(\det A)J$ for $A\in GL_2$, ${}^t\rho(s)P'\rho(csc^{-1})={}^t\rho(s)J\rho(s)\rho(c)^{-1}=\det\rho(s)P'$ for $s\in G_{k_0}$, and so for $s\in G_M$. So $\rho_M$ has the polarization of Lemma~\ref{sha-siegel-step-seeds-induced-automorphic-windows}(2), with the totally real field $k_0$, the character $\det\rho$, $c$ and $P'$.
\item[(6)] \emph{Supports over $M$.} By Lemma~\ref{sha-siegel-step-seeds-induced-automorphic-windows}(1), applied to $k_0\subset E\subseteq M$, $\rho$ and the trivial character $\kappa=1$, attached to the trivial Hecke character, every regular algebraic cuspidal automorphic representation of $GL_2(\mathbf A_M)$ to which $\rho_M$ is attached through $\iota$, in particular $P$, satisfies the condition on supercuspidal supports.
\item[(7)] \emph{The places of $L$ that do not split in $M$.} Let $v$ be a finite place of $L$ that does not split in $M$. If $v$ lay above $2$ or $\ell$, the decomposition groups of $G_L$ at the places above $v$ would fix $k_0$ by (a), and $v$ would split in $M=k_0L$; so $v$ lies above neither. By condition (o) of Definition~\ref{sha-siegel-tower-towers}(4) and (a), the inertia groups of $G_L$ at the places above $v$ fix $k_0$, so $M$ is unramified over $L$ at $v$. Let $w$ be a place of $M$ above $v$. By (b), $\rho$ is conjugate on the inertia groups of $G_{k_0}$ at $w$ to the diagonal representation with entries $1$ and $\xi$, and $\xi$ is trivial on them by condition (q), as $w$ is not above $2$; so $\rho_M$ is trivial on the inertia groups at $w$. By (4) to (6), Lemma~\ref{sha-siegel-step-seeds-induced-automorphic-windows}(2) applies to $M$, $\rho_M$, $P$ and $w$, which is not above $\ell$: $P$ has Satake parameters at $w$. So $P$ satisfies condition (ii) of Lemma~\ref{sha-siegel-step-seeds-induced-automorphic-published}(3) for $L\subset M$.
\item[(8)] \emph{Infinitesimal characters.} Let $\phi$ be an embedding of $L$ in $\mathbf C$, and $\sigma\ne\sigma'$ the two embeddings of $M$ above it. As $M=k_0L$, $\sigma$ and $\sigma'$ differ on $k_0$, so they induce the two infinite places of $k_0$. By (4) and (3), $P$ has at $\sigma$ and $\sigma'$ the infinitesimal characters of weights $(j+t,-t)$ and $(j,0)$, in some order, and the four integers $j+t$, $-t$, $j$, $0$ are pairwise distinct as $j,t\ge2$. So $P$ satisfies condition (i) of Lemma~\ref{sha-siegel-step-seeds-induced-automorphic-published}(3) for $L\subset M$.
\item[(9)] \emph{Characteristic polynomials.} Let $s\in G_L$, $s\notin G_M$; then $s\notin G_{k_0}$ by (2). For $h\in G_{k_0}$, $R(h)=M(\beta/\alpha,\rho(h),\rho(\tau^{-1}h\tau))$ by (d), and an invertible matrix that depends only on $\alpha$ and $\beta$ conjugates $M(\beta/\alpha,A,B)$ to a matrix with characteristic polynomial that of $A$ times that of $B$. For $g,g'\in G_M$ with $g'=s^{-1}gs$, let $h=\tau^{-1}s$, which lies in $G_{k_0}$ as $G_{k_0}$ has index $2$ in $G$ and $\tau,s\notin G_{k_0}$; then $\tau^{-1}g\tau=hg'h^{-1}$, and the characteristic polynomial of $R(g)$ is that of $\rho(g)$ times that of $\rho(g')$. If $g\in G_M$ and $g=s^2$, a computation with the same matrices, $R(\tau)$ of (d) and $s=(s\tau^{-1})\tau$, gives that the characteristic polynomial of $R(s)$ is $f(X^2)$, where $f$ is the characteristic polynomial of $\rho(g)$.
\end{enumerate}
By (1), the restriction of $R$ to $G_L$ is irreducible, hence semisimple. By (7) to (9), Lemma~\ref{sha-siegel-step-seeds-induced-automorphic-induction}, applied to $L\subset M$, $P$, $\rho_M$ and the restriction of $R$ to $G_L$, shows that this restriction is attached through $\iota$ to a regular algebraic cuspidal automorphic representation of $GL_4(\mathbf A_L)$. By (6), (7) and (9), Lemma~\ref{sha-siegel-step-seeds-induced-automorphic-windows}(3) shows that every regular algebraic cuspidal automorphic representation of $GL_4(\mathbf A_L)$ to which it is attached through $\iota$ satisfies the condition on supercuspidal supports.
% indAutLP_of_upQuad_upSol_indAutCs (BirthIndS.lean, 164 lines), followed step by step: (1) is hirr (induced_irreducible) and hL, with hfix for G_L; (2) is hkL (residual_characters_psi_ne with Hyp.ps), hKk, the private deg_two for dE and dM, and the instances for E and M (IsAbelianGalois.of_isCyclic, isAbelianGalois_sup, IsCMField.of_isAbelianGalois); (3) is hreg with two_places and inf_reg; (4) is hQ (UpQuadI) and hS (UpSolI), with resFacts for the irreducibility from MixDat.irr; (5) is hpol, polMix_alt (algOfCyc_s for det, the sign through complex_conjugations_are_conjugate and cyclotomic_char_conj_invariant, conj_el_ex) with solvable_polarized_restriction to E and M; (6) is hW, awChain_of_printsH with kappa = 1; (7) is hbad, with hsp (splitOfDec and SplitIn 2 and l), unrOfIne with Tower.Two.unr and Hyp.ps, and sphOfUnrW_alt with MixDat.away and Tower.Two.xi; (8) is the bullet with two_emb, pl, hom_ext_sup and hinj; (9) is symmetric_square_mixed_charpoly (P_mix, cp_conj, cp_blk) for the inclusion of G_{k_0}. The conclusion is IndAutCs (hI) and awInd, with the semisimplicity from galoisRep_irreducibility_criterion.
% IndAutLP is proved from HypothesesH' by indAutLP_of_printsH_upQuad_upSol_indCs, that is indAutLP_of_upQuad_upSol_indAutCs (BirthIndS.lean, 164 lines: induced_irreducible for the irreducibility on G_L, E = k_0 K and M = k_0 L, residual_characters_psi_ne for k_0 not in L, Z.Bc.UpQuadI and UpSolI for the base changes, IndAutCs for the induction, awInd for the supercuspidal supports), with Z.Bc.upQuadI_of_alt, Z.Bc.upSolI_of_alt and indAutCs_of_alt from Z.Bc.indCs. Size: the open parts below this lemma are Lemmas sha-siegel-step-seeds-induced-automorphic-windows-real (plainUp) and -automorphic-windows-induction-local (indTor and indX); lean_size of the three declarations together, stopping at the book's citations, is 237,000 lines, of which about 17,800 are in AutomorphicArithmetic/AutomorphicVarieties and 26,400 in AutomorphicArithmetic/Foundations, the rest mostly the local and global theory of GL_n (GroundTheory/AutomorphicForms 60,800, AbelianVarietyTheory/Proofs 50,900, AutomorphicSpectrum/Proofs 24,400). Automorphic induction itself, Z.Bc.indCs through InductionEstimates (about 54,000 lines), is cited as Lemma sha-siegel-step-seeds-induced-automorphic-published (3).
\end{proof}

\begin{lemma}\label{sha-siegel-step-seeds-induced-residues}
Let $\ell$ be an odd prime, $\gamma,x\in\bar{\mathbf Q}_\ell^\times$ with $|\gamma|=|x|=1$, and $A,B\in GL_2(\bar{\mathbf Q}_\ell)$ such that, for $C=A$ and for $C=B$, $|C_{11}-1|<1$, $|C_{22}-x|<1$, $|C_{21}|<1$ and $|C_{12}|\le1$. Let $M=M(\gamma,A,B)$ (Definition~\ref{sha-siegel-step-seeds-induced-datum}(2)). Then the entries of $M$ have absolute value at most $1$; $|M_{ii}-x|<1$ for $i=1,3$ and $|M_{ii}-1|<1$ for $i=2,4$; the off-diagonal entries of $M$ of weight outside $\{(-1,1),(-1,0),(-1,-1)\}$ (Definition~\ref{sha-siegel-datum}) have absolute value less than $1$; and if $|A_{12}|=1$, then $|M_{21}|=1$, or $|M_{23}|=|M_{41}|=1$.
\end{lemma}
\leanref{Development/AutomorphicArithmetic/AutomorphicVarieties/BuildingsResidueChars/MixMatrixResidueProperties.lean}{UpperTheory.mixMatrix_residue_properties}
\leanref{Development/AutomorphicArithmetic/AutomorphicVarieties/IdeleGenerators/QuadraticSplitting/QuadSplit.lean}{UpperTheory.V.Sd.S2.MixRes}
\begin{proof}
As $\ell$ is odd, $|2|=1$, so each entry of $M$ has the absolute value of the sum or difference of the two terms in its numerator. The entries of $A$ and $B$ have absolute value at most $1$, and so do those of $M$. The diagonal entries are $(A_{22}+B_{22})/2$ and $(A_{11}+B_{11})/2$, congruent to $x$ and to $1$. The off-diagonal entries of weight outside $\{(-1,1),(-1,0),(-1,-1)\}$ are those other than $(2,1)$, $(2,3)$, $(4,1)$ and $(4,3)$; each of them is half of $\pm A_{21}/\gamma\pm B_{21}$, of $\pm(A_{22}-B_{22})$ or of $\pm(A_{11}-B_{11})$, which have absolute value less than $1$. Finally $M_{21}+M_{41}=\gamma A_{12}$ and $M_{23}=-M_{41}$. If $|A_{12}|=1$ and $|M_{21}|<1$, then $|M_{41}|=1$, and $|M_{23}|=1$.
\end{proof}
% mixMatrix_residue_properties (MixMatrixResidueProperties.lean, 106 lines): the private key, entry by entry, with s and n for the last claim.

\begin{lemma}\label{sha-siegel-step-seeds-induced-kill}
Let $D$, $T$ and $\iota$ be as in Definition~\ref{sha-siegel-step-seeds-induced-datum}, and let $(k_0,\varphi,j,t,\rho,\alpha,\beta,R,\delta)$ be an induction datum for $T$ and $\iota$, with twisted restriction $r_\delta$. Then $r_\delta(g)=1$ for every $g$ in the set $N_\Phi$ of the problem of $T$ (Definition~\ref{sha-siegel-tower-pairs}(1)).
\end{lemma}
\leanref{Development/AutomorphicArithmetic/AutomorphicVarieties/BuildingsResidueChars/InducedRepresentationConstituentVanishing.lean}{UpperTheory.induced_representation_constituent_vanishing}
\leanref{Development/AutomorphicArithmetic/AutomorphicVarieties/IdeleGenerators/QuadraticSplitting/QuadSplit.lean}{UpperTheory.V.Sd.S2.IndKill}
\begin{proof}
Let $g\in N_\Phi$: $g\in I_w(F)$ for a place $w$ not above $\ell$, and $\bar r_0(g)=1$.
\begin{enumerate}
\item[(1)] \emph{Characters.} By condition (i) of Definition~\ref{sha-siegel-tower-towers}(4), $\psi(g)=1$. By condition (o), $\bar\chi^{(0,1)}(g)=1$, so $g\in G_{k_0}$ by (a) of Definition~\ref{sha-siegel-step-seeds-induced-datum}, and $\varphi(g)=1$. As $w$ is not above $\ell$, $g$ fixes the roots of unity of $\ell$-power order, so $\chi_{\mathrm{cyc}}(g)=1$, and $\omega(g)=1$ by Lemma~\ref{sha-siegel-step-seeds-induced-characters}(6). So $\alpha(g)=\chi_{\mathrm{cyc}}(g)\omega(g)^{-1}\xi(g)=\xi(g)$. The entry $(3,3)$ of $\bar r_0(g)=1$ is $\bar\alpha(g)$, so $\xi(g)\equiv1$, and $\xi(g)=1$ by Lemma~\ref{sha-siegel-step-seeds-induced-characters}(6). Hence $\theta(g)=\alpha(g)\chi_{\mathrm{cyc}}(g)^k\psi(g)=1$, with $k$ the integer of $D$, and $\mu(g)=1$.
\item[(2)] \emph{The twist.} By (f), $\delta(g)^2=\theta(g)/\mu(g)=1$. As $\delta(g)\equiv1$, $-1\not\equiv1$ and $\ell$ is odd, $\delta(g)=1$.
\item[(3)] \emph{The induced representation.} As $g\in I_w(k_0)$ and $\xi(g)=1$, $\rho(g)=1$ by the condition on inertia in (b). Let $g'=(g_0^T)^{-1}gg_0^T$. It lies in $G_{k_0}$, which is normal of index $2$ in $G$, and in the inertia group $I_{w'}(k_0)$ at $w'=(g_0^T)^{-1}w$, which is not above $\ell$; and $\xi(g')=\xi(g)=1$. So $\rho(g')=1$ as well, and by (d), $R(g)=M(\beta/\alpha,1,1)=1$.
\end{enumerate}
So $r_\delta(g)=\delta(g)R(g)=1$.
\end{proof}
% induced_representation_constituent_vanishing (InducedRepresentationConstituentVanishing.lean, about 160 lines with its dependencies): Tower.Ok.kb for psi, Tower.Two.unr with Hyp.ps for k_0, the private cyc_ine (through Chr.clsFacts), CharDat.om, Raw.rbar at (2,2) with Fr.sc_norm and CharDat.xi, the private sqrt_one, diag_one with MixDat.away at A and at the conjugate place (Two.ineConj), and mix_one.

\begin{lemma}\label{sha-siegel-step-seeds-induced-global}
Let $D$, $T$ and $\iota$ be as in Definition~\ref{sha-siegel-step-seeds-induced-datum}, and let $(k_0,\varphi,j,t,\rho,\alpha,\beta,R,\delta)$ be an induction datum for $T$ and $\iota$, with twisted restriction $r_\delta$. Then $r_\delta$ satisfies condition (a) of Lemma~\ref{sha-siegel-step-seeds-final}.
\end{lemma}
\leanref{Development/AutomorphicArithmetic/AutomorphicVarieties/IdeleGenerators/QuadraticSplitting/IndGlobOfKill.lean}{UpperTheory.V.Sd.S2.indGlob_of_kill}
\leanref{Development/AutomorphicArithmetic/AutomorphicVarieties/IdeleGenerators/QuadraticSplitting/QuadSplit.lean}{UpperTheory.V.Sd.S2.GlobOk}
\begin{proof}
Let $W=\{(-1,1),(-1,0),(-1,-1)\}$, and say that a matrix $X\in M_4(\bar{\mathbf Q}_\ell)$ is \emph{close to} $(x_1,\dots,x_4)$ if its entries have absolute value at most $1$, $X_{ii}\equiv x_i$, and its off-diagonal entries of weight outside $W$ have absolute value less than $1$. By (e) and Lemma~\ref{sha-siegel-step-seeds-induced-characters}(1), $|\alpha|=|\beta|=1$; let $\gamma=\beta/\alpha$.
\begin{enumerate}
\item[(1)] \emph{Elements of $G_{k_0}$.} Let $g\in G_F\cap G_{k_0}$ and $g'=(g_0^T)^{-1}gg_0^T$, which lies in $G_F\cap G_{k_0}$, as $g_0^T\in G_F$ and $G_{k_0}$ is normal in $G$; and $\xi(g')=\xi(g)$. By (d), $R(g)=M(\gamma,\rho(g),\rho(g'))$, and by (c), $\rho(g)$ and $\rho(g')$ satisfy the hypotheses of Lemma~\ref{sha-siegel-step-seeds-induced-residues} with $x=\xi(g)$. So $R(g)$ is close to $(\xi(g),1,\xi(g),1)$, hence to $(\chi_1(g),\dots,\chi_4(g))$ by Lemma~\ref{sha-siegel-step-seeds-induced-characters}(2); and if $|\rho(g)_{12}|=1$, then $|R(g)_{21}|=1$, or $|R(g)_{23}|=|R(g)_{41}|=1$.
\item[(2)] \emph{Elements outside $G_{k_0}$.} Let $g\in G_F$, $g\notin G_{k_0}$. Then $h=g(g_0^T)^{-1}\in G_F\cap G_{k_0}$, as $G_{k_0}$ has index $2$ and $g_0^T\notin G_{k_0}$, and $R(g)=R(h)\operatorname{diag}(\beta,-\alpha,-\beta,\alpha)$ by (d). By (e) and Lemma~\ref{sha-siegel-step-seeds-induced-characters}(3), $(\beta,-\alpha,-\beta,\alpha)\equiv(\chi_1(g_0^T),\chi_2(g_0^T),\chi_3(g_0^T),\chi_4(g_0^T))$, and these have absolute value $1$. Multiplying the columns of $R(h)$ by them keeps the absolute values of the entries, so $R(g)$ is close to $(\chi_i(h)\chi_i(g_0^T))_i=(\chi_i(g))_i$, by (1).
\item[(3)] \emph{The twist.} As $\delta(g)\equiv1$, $r_\delta(g)=\delta(g)R(g)$ has entries of the same absolute values as $R(g)$, and is close to $(\chi_i(g))_i$ for every $g\in G_F$.
\end{enumerate}
Now $r_\delta$ is continuous, as $\delta$ and $R$ are. By (d) and (f), $r_\delta(g)$ is symplectic with multiplier $\delta(g)^2\mu(g)=\theta(g)$. By Lemma~\ref{sha-siegel-step-seeds-induced-kill}, $r_\delta(g)=1$ for $g\in N_\Phi$. The matrix $r_\delta(g_0^T)=\delta(g_0^T)\operatorname{diag}(\beta,-\alpha,-\beta,\alpha)$ is diagonal. The conditions on the entries are (3). By (c) some $g\in G_F\cap G_{k_0}$ has $|\rho(g)_{12}|=1$, and by (1) and (3), $|r_\delta(g)_{21}|=1$, or $|r_\delta(g)_{23}|=|r_\delta(g)_{41}|=1$.
\end{proof}
% indGlob_of_kill (IndGlobOfKill.lean, about 80 lines; the property Sh is "close to"): inn is (1), through mixMatrix_residue_properties and CharDat.inn; all is (2), with Sh.mul_diag, CharDat.c1 and c2; all' is (3), with Sh.smul; then the fields of GlobOk: cont, mult (gsp_smul with delta_sq), kill (induced_representation_constituent_vanishing), nrm (mix.off), int, diag, supp and ent (from ns and inn).

\begin{lemma}\label{sha-siegel-step-seeds-induced}
Assume the Galois representations as in Theorem~\ref{sha-residually-reducible-automorphy}.
Let $D$ be a Siegel datum of case II, $T$ a fit deep tower for $D$ (Definition~\ref{sha-siegel-tower-towers}), and $\iota:\bar{\mathbf Q}_\ell\cong\mathbf C$.
Then there are a homomorphism $r:G_F\to M_4(\bar{\mathbf Q}_\ell)$ and matrices $n_v$, $v\in\Sigma_T$, that satisfy (a), (b) and (c) of Lemma~\ref{sha-siegel-step-seeds-final}.
\end{lemma}
\leanref{Development/AutomorphicArithmetic/AutomorphicVarieties/ModularityArgument/ThresholdArguments/DenseThresholdCases/BirthSeedII.lean}{UpperTheory.V.Z.Bc.seedII_of_p_alt}
\begin{proof}
As $D$ is of case II and $T$ is fit, $T$ satisfies the conditions (a) to (r), so $D$, $T$ and $\iota$ are as in Definition~\ref{sha-siegel-step-seeds-induced-datum}.
\begin{enumerate}
\item[(1)] \emph{The datum.} Lemma~\ref{sha-siegel-step-seeds-induced-exists} gives an induction datum $(k_0,\varphi,j,t,\rho,\alpha,\beta,R,\delta)$ for $T$ and $\iota$, with $\rho$ attached through $\iota$ to a regular algebraic cuspidal automorphic representation of $GL_2(\mathbf A_{k_0})$ that has, at every finite place, Satake parameters or only plain local components. Let $r=r_\delta$ be its twisted restriction, and let $n_v$, $m_1,\dots,m_4$, $b$ and $N$ be as in Lemma~\ref{sha-siegel-step-seeds-induced-local}.
\item[(2)] \emph{Conditions (a) and (b).} They hold by Lemmas~\ref{sha-siegel-step-seeds-induced-global} and~\ref{sha-siegel-step-seeds-induced-local}.
\item[(3)] \emph{Spanning.} By Lemma~\ref{sha-siegel-step-seeds-induced-irreducible}, $R$ is irreducible on $G'=G_L\cap\ker\chi_{\mathrm{cyc}}$. As $r(g)=\delta(g)R(g)$ with $\delta(g)$ a scalar, $r$ is irreducible on $G'$, and by Burnside's theorem the $r(g)$ with $g\in G'$ span $M_4(\bar{\mathbf Q}_\ell)$.
\item[(4)] \emph{Weights.} Let $h_i=2m_i+b$. Then $h_i-h_{i+1}=2(m_i-m_{i+1})$ is even and at least $4$, and $(r,(n_v))$ has inertial weights $h$ by Lemma~\ref{sha-siegel-step-seeds-induced-local}.
\item[(5)] \emph{The twisting character.} Let $\kappa$ be the restriction of $\delta^{-1}$ to $G_L$, which is continuous, $a=-b$ and $k=N\ge1$. Then $a+h_1=2m_1$ is even, and for $s$ in an inertia group of $G_L$ at a place above $\ell$, $(\kappa(s)^2\chi_{\mathrm{cyc}}(s)^{-a})^N=\big((\delta(s)^2\chi_{\mathrm{cyc}}(s)^{-b})^N\big)^{-1}=1$ by Lemma~\ref{sha-siegel-step-seeds-induced-local}.
\item[(6)] \emph{Automorphy.} The twist of the restriction of $r$ to $G_L$ by $\kappa$ is the restriction of $R$ to $G_L$. By Lemma~\ref{sha-siegel-step-seeds-induced-automorphic}, it is attached through $\iota$ to a regular algebraic cuspidal automorphic representation $\pi$ of $GL_4(\mathbf A_L)$, which satisfies the condition on supercuspidal supports of Definition~\ref{sha-siegel-tower-ideal}(1).
\end{enumerate}
By (3) to (6), $(r,(n_v))$ is a seed for $T$, which is (c).
% seedII_of_p_alt (BirthSeedII.lean, 30 lines), followed step by step: (1) is indEx and indLoc0; (2) is indGlob_of_kill with induced_representation_constituent_vanishing, and hloc; (3) is induced_irreducible with irrIff (Burnside: KStar.Schema.span_eq_top_of_irreducible) and shapeFacts for the scalar twist; (4) to (6) are the fields of IsSeedAlt.aut: h = 2m + b, kappa = (delta restricted to G_L)^-1, pi from IndAutLP, a = -b, k = N, the evenness through m 0, the inertia condition through hfin, mkReg' with the AllW of IndAutLP, and the attachment through tw_val and smul_smul. The data part, up to the call of induced_final_statement_alt, rests on Lemmas sha-siegel-step-seeds-induced-exists, -irreducible, -local, -automorphic and -global; its open parts are Lemma sha-siegel-step-seeds-induced-hilbert-form (about 13,500 lines when last measured, the Hilbert modular form) and, below Lemma -automorphic, whose proof is written, Lemmas -automorphic-windows-real and -automorphic-windows-induction-local (lean_size together 237,000 lines, of which about 17,800 in AutomorphicArithmetic/AutomorphicVarieties; the rest is mostly the published local and global theory of GL_n; see the comment at the end of the proof of -automorphic).
\end{proof}

\begin{lemma}\label{sha-siegel-step-seeds}
Assume the Galois representations as in Theorem~\ref{sha-residually-reducible-automorphy}.
Let $D$ be a Siegel datum of case II, let $T$ be a fit deep tower for $D$ (Definition~\ref{sha-siegel-tower-towers}), and let $\iota:\bar{\mathbf Q}_\ell\cong\mathbf C$.
There are a ring $\mathcal O_1\supseteq\mathcal O$ of Definition~\ref{sha-character-parts}, with residue field $k_1$, and an integral seed $Q$ for $T$ over $\mathcal O_1$ (Definition~\ref{sha-siegel-tower-ideal}) whose image $(\bar r,(\bar n_v))$ in $k_1$ is of type V, with $\bar r$ having an entry of weight $(-1,0)$ or entries of weights $(-1,1)$ and $(-1,-1)$.
\end{lemma}
\leanref{Development/AutomorphicArithmetic/AutomorphicVarieties/ModularityArgument/ThresholdArguments/DenseThresholdCases/PotentialModularityTower.lean}{UpperTheory.V.Ws.SeedIIAlt}
\leanref{Development/AutomorphicArithmetic/AutomorphicVarieties/ModularityArgument/ThresholdArguments/DenseThresholdCases/BcIndCs.lean}{UpperTheory.V.Z.Bc.seedIIs_alt}
\begin{proof}
Lemma~\ref{sha-siegel-step-seeds-induced} gives $r$ and matrices $n_v$ satisfying (a), (b) and (c) of Lemma~\ref{sha-siegel-step-seeds-final}, and that lemma gives $\mathcal O_1$ and $Q$.
% SeedIIAlt is proved from HypothesesH' by Z.Bc.seedIIs_alt through seedII_Z_alt and seedII_of_p_alt, whose last step is induced_final_statement_alt (Lemma sha-siegel-step-seeds-final) applied to the twist twF T R delta of the induced representation and the flags of indLoc0 (Lemma sha-siegel-step-seeds-induced, whose open parts are its Lemmas -hilbert-form, -automorphic-windows-real and -automorphic-windows-induction-local).
\end{proof}

\subsection{Chains of flagged pairs and cohomology bounds}\label{subsection-sha-chains}

\begin{definition}\label{sha-siegel-step-chains-notions}
Let $D$ be a Siegel datum and $T$ a tower for $D$ with field $F$, with the notation of Definitions~\ref{sha-siegel-tower-towers} and~\ref{sha-siegel-tower-pairs}. Put $W_{\mathrm r}=\{(-1,1),(-1,0),(-1,-1),(0,-1)\}$, $W_{\mathrm n}=\{(-1,0),(0,-1),(-1,-1)\}$, $W_{\mathrm m}=\{(-1,0),(1,-1),(0,-1),(-1,-1)\}$ and $W_{\mathrm V}=\{(-1,1),(-1,0),(-1,-1)\}$.
\begin{enumerate}
\item Let $W$ be a finite set of primes. The \emph{problem of $T$ with auxiliary primes $W$} is the problem of $T$, except that its set $N_\Phi$ consists of the $g\in G_F$ that lie in $I_w(F)$ for a place $w$ above no prime of $W\cup\{\ell\}$ and have $\bar r_0(g)=1$; for $W=\emptyset$ it is the problem of $T$. A prime $p$ \emph{splits} in $T$ if $D_w(\mathbf Q)\subseteq G_F$ and $\bar r_0$ and $\omega$ are trivial on $D_w(\mathbf Q)$ for every place $w$ above $p$. A tower $T'$ with field $F'$ \emph{kills} $p$ if every continuous homomorphism $G\to\mathbf Z/\ell$ that is trivial on $I_w(\mathbf Q)$ for every place $w$ not above $p$ is trivial on $G_{F'}$.
\item Let $T\le T'$ be towers with fields $F\subseteq F'$. A flagged pair $c'=(r',(n'_{v'}))$ for a problem of $T'$ is \emph{over} a flagged pair $c=(r,(n_v))$ for a problem of $T$, along the same homomorphism, if $r'$ is the restriction of $r$ to $G_{F'}$ and for every $v'\in\Sigma_{T'}$ there are $v\in\Sigma_T$ and $g\in G_F$ carrying $v$ to $v'$ such that $n_{v'}'^*\,r(g)\,n_v$ is upper triangular.
\item Let $c=(r,(n_v))$ and $c'=(r',(n'_v))$ be flagged pairs for the same problem. They have the \emph{same $(0,-1)$-part} if $r(g)_{31}=r'(g)_{31}$ for all $g$ and $(n_v)_{31}=(n'_v)_{31}$ for all $v$, and the \emph{same $(-1,0)$-part} if the same holds for the entries $(2,1)$. For a set $W'$ of weights, $c$ \emph{lies in} $W'$ if $r$ is supported in $W'$ and each $n_v$ is supported in $W_{\mathrm n}$.
\end{enumerate}
\end{definition}
\leanref{Development/AutomorphicArithmetic/AutomorphicVarieties/WeakShimuraTheory/NearSplit.lean}{UpperTheory.V.Ws.Tower.fr}
\leanref{Development/AutomorphicArithmetic/AutomorphicVarieties/WeakShimuraTheory/NearSplit.lean}{UpperTheory.V.Ws.Tower.Splits}
\leanref{Development/AutomorphicArithmetic/AutomorphicVarieties/WeakShimuraTheory/NearSplit.lean}{UpperTheory.V.Ws.Dies}
\leanref{Development/AutomorphicArithmetic/AutomorphicVarieties/WeakShimuraTheory/NearSplit.lean}{UpperTheory.V.Ws.Tower.OverT}
\leanref{Development/AutomorphicArithmetic/AutomorphicVarieties/WeakShimuraTheory/RedIdeal.lean}{UpperTheory.V.Ws.PD.SamePs}
\leanref{Development/AutomorphicArithmetic/AutomorphicVarieties/WeakShimuraTheory/RedIdeal.lean}{UpperTheory.V.Ws.PD.SameAl}
\leanref{Development/AutomorphicArithmetic/AutomorphicVarieties/WeakShimuraTheory/RedIdeal.lean}{UpperTheory.V.Ws.PD.In}

\begin{lemma}\label{sha-siegel-step-chains-lines-parabolic}
Let $R$ be a commutative ring, and let $\kappa$ be one of the vectors $\lambda=(1,0,1,0)$ and $\mu=(1,1,0,0)$ of Definition~\ref{sha-siegel-datum}, which define the weights. Call $x\in M_4(R)$ \emph{$\kappa$-parabolic} if $x_{ij}=0$ whenever $\kappa_i-\kappa_j=1$, and \emph{$\kappa$-radical} if $x_{ij}=0$ whenever $\kappa_i-\kappa_j\ne-1$. For $\kappa$-parabolic $x,y\in M_4(R)$ with $y-x$ $\kappa$-radical, put $L(x,y)=x+X(y-x)\in M_4(R[X])$.
\begin{enumerate}
\item Products of $\kappa$-parabolic matrices, and $x^*$ for $\kappa$-parabolic $x$, are $\kappa$-parabolic. The product of a $\kappa$-radical and a $\kappa$-parabolic matrix, in either order, and $u^*$ for $\kappa$-radical $u$, are $\kappa$-radical. The product of two $\kappa$-radical matrices is $0$.
\item Let $(x,y)$ and $(x',y')$ be pairs as above. Then so are $(xx',yy')$ and $(x^*,y^*)$, and $L(x,y)L(x',y')=L(xx',yy')$ and $L(x,y)^*=L(x^*,y^*)$.
\end{enumerate}
\end{lemma}
\begin{proof}
(1) The entries of $\kappa$ are $0$ and $1$. If $\kappa_i-\kappa_k=1$, then $\kappa_i=1$ and $\kappa_k=0$, so every $j$ has $\kappa_i-\kappa_j=1$ or $\kappa_j-\kappa_k=1$; hence $(xx')_{ik}=\sum_jx_{ij}x'_{jk}=0$ for $\kappa$-parabolic $x,x'$. Let $u$ be $\kappa$-radical and $x$ $\kappa$-parabolic. A nonzero term $u_{ij}x_{jk}$ needs $\kappa_i=0$, $\kappa_j=1$ and $\kappa_j-\kappa_k\ne1$, so $\kappa_k=1$ and $\kappa_i-\kappa_k=-1$; so $ux$ is $\kappa$-radical, and $xu$ likewise. A nonzero term $u_{ij}u'_{jk}$ would need $\kappa_j=1$ and $\kappa_j=0$. Finally $\kappa_{5-i}=1-\kappa_i$, so the entry $(i,j)$ of $x^*$, which is $\pm x_{5-j,5-i}$, comes from an entry with $\kappa_{5-j}-\kappa_{5-i}=\kappa_i-\kappa_j$.

(2) By (1), $xx'$, $yy'$, $x^*$ and $y^*$ are $\kappa$-parabolic, $y^*-x^*=(y-x)^*$ is $\kappa$-radical, and $(y-x)(y'-x')=0$, so $yy'-xx'=(y-x)x'+x(y'-x')$ is $\kappa$-radical. Then
\[
L(x,y)L(x',y')=xx'+X\big((y-x)x'+x(y'-x')\big)+X^2(y-x)(y'-x')=xx'+X(yy'-xx')=L(xx',yy'),
\]
and $L(x,y)^*=x^*+X(y-x)^*=L(x^*,y^*)$, as $z\mapsto z^*$ is linear.
\end{proof}
% For lambda these are the private Cong, Cong.adj, Cong.mul and Cong.aff_mul of AffSmulOne.lean (with aff x y = (1-X)x + Xy = L(x,y)); for mu the private GlobalTheory (radical), Par (parabolic), blk_mul, blk_par, par_blk, blk_sub_mul, comb_mul, adj_comb, blk_adj and par_adj of CompactIntegralFields/GlobalTheory.lean (with comb A B = L(A,B)).

\begin{lemma}\label{sha-siegel-step-chains-lines}
Let $D$ be a Siegel datum, $T$ a fit tower for $D$, $k$ a field and $a:\mathcal O\to k$ a homomorphism that kills $\mathfrak m_{\mathcal O}$. Let $c$ and $c'$ be residual flagged pairs for the problem of $T$ along $a$, with the notation of Definition~\ref{sha-siegel-step-chains-notions}.
\begin{enumerate}
\item If $c$ and $c'$ lie in $W_{\mathrm r}$ and have the same $(0,-1)$-part, they are linked.
\item If their homomorphisms are supported in $W_{\mathrm m}$ and they have the same $(-1,0)$-part, they are linked.
\end{enumerate}
\end{lemma}
\leanref{Development/AutomorphicArithmetic/AutomorphicVarieties/WeakShimuraTheory/CompactIntegralFields/AffSmulOne.lean}{UpperTheory.V.Ws.lineP}
\leanref{Development/AutomorphicArithmetic/AutomorphicVarieties/WeakShimuraTheory/CompactIntegralFields/GlobalTheory.lean}{UpperTheory.V.Ws.lineM}
\begin{proof}
Write $c=(r,(n_v))$ and $c'=(r',(n'_v))$, and use Lemma~\ref{sha-siegel-step-chains-lines-parabolic}, with $\kappa=\lambda$ in (1) and $\kappa=\mu$ in (2).
\begin{enumerate}
\item[(a)] \emph{Pairs in (1).} Let $x$ and $y$ be symplectic with multipliers $\nu$ and $\nu'$, supported in $W_{\mathrm r}$, with $x_{ii}=y_{ii}$ for all $i$, $x_{11}\ne0$ and $x_{31}=y_{31}$. The weights $(1,0)$, $(1,1)$ and $(1,-1)$ of the entries $(i,j)$ with $\lambda_i-\lambda_j=1$ are not in $W_{\mathrm r}$, so $x$ and $y$ are $\lambda$-parabolic. The entries $(i,j)$ with $\lambda_i=\lambda_j$ are the diagonal ones and $(1,3)$, $(2,4)$, $(3,1)$ and $(4,2)$. The first two have weight $(0,1)$, so they vanish in $x$ and $y$. The entry $(3,1)$ of $x^*x=\nu$ is $-x_{11}x_{42}-x_{21}x_{32}+x_{31}x_{22}+x_{41}x_{12}=0$, and $x_{32}=x_{12}=0$; so $x_{11}x_{42}=x_{22}x_{31}$, and likewise for $y$, which gives $x_{42}=y_{42}$. So $y-x$ is $\lambda$-radical. This applies to $x=r(g)$ and $y=r'(g)$ for $g\in G_F$, whose diagonal entries are $a(\chi_i(g))$, and to $x=n_v$ and $y=n'_v$, which are symplectic with multiplier $1$, supported in $W_{\mathrm n}\subseteq W_{\mathrm r}$ and have diagonal entries $1$; the entries $(3,1)$ agree as $c$ and $c'$ have the same $(0,-1)$-part.
\item[(b)] \emph{Pairs in (2).} Let $x$ and $y$ be lower triangular and symplectic with multipliers $\nu$ and $\nu'$, with $x_{ii}=y_{ii}$ for all $i$, $x_{11}\ne0$ and $x_{21}=y_{21}$. The entries $(i,j)$ with $\mu_i-\mu_j=1$ lie above the diagonal, so $x$ and $y$ are $\mu$-parabolic. The entries $(i,j)$ with $\mu_i=\mu_j$ are the diagonal ones and $(1,2)$, $(3,4)$, $(2,1)$ and $(4,3)$. The entry $(2,1)$ of $x^*x=\nu$ is $x_{43}x_{11}+x_{33}x_{21}-x_{23}x_{31}-x_{13}x_{41}=x_{43}x_{11}+x_{33}x_{21}=0$, and likewise for $y$, so $x_{43}=y_{43}$. So $y-x$ is $\mu$-radical. The weights of $W_{\mathrm m}$ are those of the entries below the diagonal, so this applies to $x=r(g)$ and $y=r'(g)$, and to $x=n_v$ and $y=n'_v$, as $c$ and $c'$ have the same $(-1,0)$-part.
\item[(c)] \emph{The linking pair.} Let $\rho(g)=L(r(g),r'(g))$ and $N_v=L(n_v,n'_v)$, and $P=(\rho,(N_v))$, along the composite of $a$ with $k\subset k[X]$. By Lemma~\ref{sha-siegel-step-chains-lines-parabolic}(2), $\rho(gh)=\rho(g)\rho(h)$, $\rho(g)^*\rho(g)=L(r(g)^*r(g),r'(g)^*r'(g))=a(\theta(g))$, $N_v^*N_v=1$, and $N_v^*\rho(s)N_v=L(n_v^*r(s)n_v,n_v'^*r'(s)n'_v)$. The entry $(i,j)$ of $L(x,y)$ is $x_{ij}+X(y_{ij}-x_{ij})$: it vanishes when $x_{ij}=y_{ij}=0$, and equals $x_{ij}$ when $x_{ij}=y_{ij}$. As $c$ and $c'$ are residual flagged pairs, it follows that $\rho(1)=1$, $\rho(g)=1$ for $g\in N_\Phi$, $\rho(g_0^T)$ is diagonal, $N_v$ is lower triangular with diagonal entries $1$, and for $s\in D_v(F)$ the matrix $N_v^*\rho(s)N_v$ is upper triangular with diagonal entries $a(\chi_i(s))$, which are $1$ for $s\in\mu_v$; that $\rho(g)_{ii}=a(\chi_i(g))$; and that $\rho$ is supported in $W_{\mathrm r}$ and the $N_v$ in $W_{\mathrm n}$ in (1), and $\rho$ in $W_{\mathrm m}$ in (2). The entries of $\rho$ are locally constant, as those of $r$ and $r'$ are. So $P$ is a residual flagged pair, and its images under $X\mapsto0$ and $X\mapsto1$ are $c$ and $c'$.
\end{enumerate}
\end{proof}
% lineP (AffSmulOne.lean, 215 lines) and lineM (CompactIntegralFields/GlobalTheory.lean, 198 lines), for any deformation problem with Fr.Ok along a, over a compact totally disconnected group; the problem of a fit tower has Fr.Ok along a by Lemma sha-siegel-step-coefficients-standard (1), which neither proof uses. (a) is the private entries_zero_and_relation_of_suppM_of_isGSp and cong_of, with supp_nP for the n_v; (b) is the private e32, blk_of and low_of; (c) is aff (comb), gsp (comb_gsp), conj and the final refine. About 410 lines.

\begin{lemma}\label{sha-siegel-step-chains-levi}
Let $D$, $T$, $k$ and $a$ be as in Lemma~\ref{sha-siegel-step-chains-lines}.
\begin{enumerate}
\item Let $c$ be a residual flagged pair for the problem of $T$ along $a$ whose homomorphism is supported in $W_{\mathrm m}$. There is a residual flagged pair $c'=(r',(n'_v))$ with the same $(-1,0)$-part as $c$, such that $r'$ and each $n'_v$ are supported in $\{(-1,0)\}$.
\item Suppose that $a(\chi^{(0,1)}(s))=1$ for every $v\in\Sigma_T$ and $s\in D_v(F)$. Let $c=(r,(n_v))$ be a residual flagged pair for the problem of $T$ along $a$ such that $r$ and each $n_v$ are supported in $\{(-1,0)\}$, and let $y_v\in k$ for $v\in\Sigma_T$. There is a residual flagged pair $(r,(n''_v))$ with each $n''_v$ supported in $W_{\mathrm n}$, $(n''_v)_{21}=(n_v)_{21}$ and $(n''_v)_{31}=y_v$.
\end{enumerate}
\end{lemma}
\leanref{Development/AutomorphicArithmetic/AutomorphicVarieties/WeakShimuraTheory/ResidualChains/LvSmulOne.lean}{UpperTheory.V.Ws.Chain.levi}
\leanref{Development/AutomorphicArithmetic/AutomorphicVarieties/WeakShimuraTheory/ResidualChains/MovePs.lean}{UpperTheory.V.Ws.Chain.movePs}
\begin{proof}
(1) Write $c=(r,(n_v))$. The weights in $W_{\mathrm m}$ are those of the entries below the diagonal, so each $r(g)$ is lower triangular, and so is each $n_v$. For a lower triangular $x$ let $\operatorname{lv}(x)$ be the matrix with the entries $x_{ij}$ with $\mu_i=\mu_j$, and $0$ elsewhere. In the notation of Lemma~\ref{sha-siegel-step-chains-lines-parabolic} with $\kappa=\mu$, a lower triangular $x$ is $\mu$-parabolic, $\operatorname{lv}(x)$ is $\mu$-parabolic and block diagonal for the blocks $\{1,2\}$ and $\{3,4\}$, and $x-\operatorname{lv}(x)$ is $\mu$-radical. For lower triangular $x,x'$, writing $x=l+u$ and $x'=l'+u'$ with $l=\operatorname{lv}(x)$ and $l'=\operatorname{lv}(x')$, the product $ll'$ is block diagonal and $lu'+ul'+uu'$ is $\mu$-radical by (1) of that lemma; so $\operatorname{lv}(xx')=\operatorname{lv}(x)\operatorname{lv}(x')$. As $z\mapsto z^*$ preserves block diagonal and $\mu$-radical matrices, $\operatorname{lv}(x^*)=\operatorname{lv}(x)^*$. The map $\operatorname{lv}$ keeps the diagonal entries and the entries $(2,1)$ and $(4,3)$, and kills the others.
Let $r'(g)=\operatorname{lv}(r(g))$, $n'_v=\operatorname{lv}(n_v)$ and $c'=(r',(n'_v))$. Then $r'$ is a homomorphism; $r'(g)^*r'(g)=\operatorname{lv}(r(g)^*r(g))=a(\theta(g))$; $r'(g)=1$ when $r(g)=1$; $r'(g_0^T)$ is diagonal; $n'_v$ is lower triangular with diagonal entries $1$ and $n_v'^*n'_v=\operatorname{lv}(n_v^*n_v)=1$; and for $s\in D_v(F)$, $n_v'^*r'(s)n'_v=\operatorname{lv}(n_v^*r(s)n_v)$, as $z\mapsto z^*$ keeps lower triangular matrices lower triangular; it is upper triangular, as $n_v^*r(s)n_v$ is, with the same diagonal entries. So $c'$ is a flagged pair, with the diagonal entries and the $\psi_{v,i}$ of $c$, and with locally constant entries; so it is residual. The off-diagonal entries of $r'(g)$ and $n'_v$ are at $(2,1)$ and $(4,3)$, of weight $(-1,0)$, and the entries $(2,1)$ are those of $c$.

(2) Write $c=(r,(n_v))$. For $y\in k$ let $u(y)=1+y(E_{31}+E_{42})$, where $E_{ij}$ are the matrix units. As $(E_{31}+E_{42})^*=-(E_{31}+E_{42})$ and $(E_{31}+E_{42})^2=0$, $u(y)^*u(y)=1$, and $u(y)$ is lower unipotent. Let $n''_v=n_vu(y_v)$. As $n_v=1+(n_v)_{21}E_{21}+(n_v)_{43}E_{43}$,
\[
n''_v=1+(n_v)_{21}E_{21}+y_vE_{31}+y_vE_{42}+(n_v)_{43}E_{43}+(n_v)_{43}y_vE_{41},
\]
which is lower unipotent, supported in $W_{\mathrm n}$, with $(n''_v)_{21}=(n_v)_{21}$ and $(n''_v)_{31}=y_v$. Let $s\in D_v(F)$ and $M=n_v^*r(s)n_v$. It is upper triangular, as $c$ is a flagged pair, and lower triangular, as $n_v^*$, $r(s)$ and $n_v$ are; so $M$ is diagonal, with diagonal entries $a(\chi_i(s))$, as $c$ is residual. By hypothesis $a(\chi_1(s))=a(\chi_3(s))$; and $\chi_1\chi_4=\chi_2\chi_3=\theta$ gives $a(\chi_2(s))=a(\chi_4(s))$. So $M$ commutes with $E_{31}+E_{42}$, hence with $u(y_v)$, and $n_v''^*r(s)n''_v=u(y_v)^*Mu(y_v)=M$. Hence $(r,(n''_v))$ is a flagged pair whose diagonal entries $\psi_{v,i}(s)$ are those of $c$, in particular $1$ on $\mu_v\subseteq D_v(F)$. Its homomorphism $r$ is supported in $\{(-1,0)\}\subseteq W_{\mathrm m}$, so it is residual.
% Chain.levi (LvSmulOne.lean, 187 lines) and Chain.movePs (MovePs.lean, 185 lines), for any deformation problem over a field with mu_v in D_v; movePs also uses chi_i chi_{5-i} = nu. In (1), lv is the private Levi.lv, with levi_lemma (lv commutes with adj) and levi_lv_mul_of_par; in (2), u(y) is the private Chain.uu, with adj_uu_mul_diag_mul_uu_eq_of_diag_pairs and mul_uu_eq_explicit_of_suppM_neg_wAl_unit_diag. About 370 lines.
\end{proof}

\begin{lemma}\label{sha-siegel-step-chains-rational-flags}
Let $k$ be a field, and let $W_{\mathrm r}$, $W_{\mathrm n}$ and $W_{\mathrm m}$ be as in Definition~\ref{sha-siegel-step-chains-notions}. Call $n\in M_4(k)$ \emph{lower unipotent} if it is symplectic with multiplier $1$, lower triangular, and has diagonal entries $1$.
\begin{enumerate}
\item Let $\nu\in k^\times$, let $X\in M_4(k)$ be symplectic with multiplier $\nu$, and let $n$ be lower unipotent. Suppose that $X$ is supported in $W_{\mathrm r}$ and $n$ in $W_{\mathrm n}$, or that $X$ is supported in $W_{\mathrm m}$. Then there is a lower unipotent $n'$ with $n'^*Xn$ upper triangular, which in the first case is supported in $W_{\mathrm n}$.
\item Let $X,X',Z\in M_4(k)$ with $XX'=1$, and let $n$ and $n'$ be lower unipotent, with $n'^*Xn$ and $n^*Zn$ upper triangular. Then $n'^*(XZX')n'$ is upper triangular, with the diagonal entries of $n^*Zn$.
\end{enumerate}
\end{lemma}
\leanref{Development/AutomorphicArithmetic/AutomorphicVarieties/WeakShimuraTheory/GradedFlagMoves/FlagNew.lean}{UpperTheory.V.Ws.Rt.flagNew}
\leanref{Development/AutomorphicArithmetic/AutomorphicVarieties/WeakShimuraTheory/GradedFlagMoves/FlagMove.lean}{UpperTheory.V.Ws.Rt.flagMove}
\begin{proof}
A lower unipotent $n$ has $n^*n=1$, so $nn^*=1$.

(1) Put $u=Xn$, which is symplectic with multiplier $\nu$.
\begin{enumerate}
\item[(a)] \emph{Entries of $X$.} The entries $(1,2)$, $(1,3)$, $(1,4)$, $(2,4)$ and $(3,4)$ have the weights $(1,0)$, $(0,1)$, $(1,1)$, $(0,1)$ and $(1,0)$, which lie in neither $W_{\mathrm r}$ nor $W_{\mathrm m}$; so they vanish in $X$. Moreover $X_{23}X_{32}=0$ and $X_{23}n_{32}=0$: in the first case $X_{32}=n_{32}=0$, as the weight $(1,-1)$ of $(3,2)$ is not in $W_{\mathrm r}$, and in the second $X_{23}=0$, as the weight $(-1,1)$ of $(2,3)$ is not in $W_{\mathrm m}$.
\item[(b)] \emph{Pivots.} By (a), the entries $(1,1)$ and $(2,2)$ of $X^*X=\nu$ are $X_{44}X_{11}=\nu$ and $X_{33}X_{22}-X_{23}X_{32}=X_{33}X_{22}=\nu$, so $X_{11}$ and $X_{22}$ are nonzero. As the first row of $X$ is $(X_{11},0,0,0)$, $n_{12}=0$ and $n_{22}=1$, this gives $u_{11}=X_{11}\ne0$, $u_{12}=0$ and $u_{22}=X_{22}+X_{23}n_{32}=X_{22}\ne0$.
\item[(c)] \emph{The flag.} As $k$ is a field, there are $x,t,y,z\in k$ with
\[
xu_{11}+u_{21}=0,\quad tu_{22}+u_{32}=0,\quad(y+tx)u_{11}+tu_{21}+u_{31}=0,\quad zu_{11}+yu_{21}-xu_{31}+u_{41}=0.
\]
Let $n'$ be the matrix with rows $(1,0,0,0)$, $(-x,1,0,0)$, $(-y,-t,1,0)$ and $(-z,-y-tx,x,1)$. A computation shows that $n'$ is lower unipotent and that $n'^*$ has rows $(1,0,0,0)$, $(x,1,0,0)$, $(y+tx,t,1,0)$ and $(z,y,-x,1)$.
\item[(d)] \emph{Upper triangular.} Let $w=n'^*u=n'^*Xn$, which is symplectic with multiplier $\nu$. By the choice of $x$, $y$ and $z$ its first column is $(u_{11},0,0,0)$, and $w_{32}=(y+tx)u_{12}+tu_{22}+u_{32}=0$. Then the entries $(4,2)$ and $(4,3)$ of $w^*w=\nu$ are $w_{11}w_{42}=0$ and $w_{11}w_{43}=0$, so $w_{42}=w_{43}=0$, and $w$ is upper triangular.
\item[(e)] \emph{Support.} In the first case, $u_{32}=X_{31}n_{12}+X_{32}+X_{33}n_{32}+X_{34}n_{42}=0$ by (a), so $t=0$ and $n'_{32}=0$. The other entries of $n'$ below the diagonal have weights in $W_{\mathrm n}$, so $n'$ is supported in $W_{\mathrm n}$.
\end{enumerate}

(2) Let $B=n'^*Xn$ and $C=n^*X'n'$. As $nn^*=1$, $XX'=1$ and $n'^*n'=1$, $BC=1$; so $C=B^{-1}$ is upper triangular, as the inverse of an invertible upper triangular matrix, and $B_{ii}C_{ii}=(BC)_{ii}=1$. As $nn^*=1$, $n'^*(XZX')n'=B\,(n^*Zn)\,C$, a product of upper triangular matrices, whose entry $(i,i)$ is $B_{ii}(n^*Zn)_{ii}C_{ii}=(n^*Zn)_{ii}$.
\end{proof}
% flagNew (FlagNew.lean, 139 lines): (a) is key, (b) is d0, d1, h00, h11 and h01, (c) is the private lowMk and adjMk with the solutions in core, (d) is the private upOf, (e) is the last clause of core (u_{32} = 0 gives t = 0) with e21. flagMove (FlagMove.lean, 47 lines): the private up_inv for C, up_mul for the product and its diagonal.

\begin{lemma}\label{sha-siegel-step-chains-rational}
Let $D$ be a Siegel datum, $T$ a fit tower for $D$ with relabelling $\epsilon$, $k$ a field and $a:\mathcal O\to k$ a homomorphism that kills $\mathfrak m_{\mathcal O}$; flagged pairs for rational problems and for problems of $T$ are along $a$.
\begin{enumerate}
\item Let $c_0$ be a residual flagged pair for the rational problem of $(\epsilon,g_0^T)$ that lies in $W_{\mathrm r}$. There is a residual flagged pair for the problem of $T$ that is over $c_0$.
\item Let $g_0\in G$ be such that $a(\chi^w(g_0))\ne1$ for every root $w$, let $c_0$ be a flagged pair of lattice shape for the rational problem of $(\epsilon,g_0)$, and let $s'$ be a residual flagged pair for the problem of $T$ that is over $c_0$. Then $s'$ is residual for the problem of $T$ with auxiliary primes $W$ for every finite set $W$ of primes, $s'$ lies in $W_{\mathrm r}$, and either $s'$ has entries of weights $(0,-1)$ and $(-1,1)$, or $s'$ is of type V and has an entry of weight $(-1,0)$.
\end{enumerate}
\end{lemma}
\leanref{Development/AutomorphicArithmetic/AutomorphicVarieties/WeakShimuraTheory/ResidualChains/ExistsResidualDatumOverRestrictionToRationals.lean}{UpperTheory.V.Ws.Chain.exists_residual_datum_over_restriction_to_rationals}
\leanref{Development/AutomorphicArithmetic/AutomorphicVarieties/WeakShimuraTheory/CompactIntegralFields/FlagFacts.lean}{UpperTheory.V.Ws.goodQ}
\begin{proof}
Let $F$ be the field of $T$. Since $a$ kills $\mathfrak m_{\mathcal O}$, the field $k$ has characteristic $\ell$, and $a(\chi_i(g))\ne0$ for all $g$.

(1) Write $c_0=(r_0,n_0)$, so that $r_0$ is supported in $W_{\mathrm r}$, $n_0$ is lower unipotent and supported in $W_{\mathrm n}$, and $r_0(g_0^T)$ is diagonal.
\begin{enumerate}
\item[(a)] \emph{Moving the places.} For $v\in\Sigma_T$ choose $g_v\in G$ with $g_vv_\ell=v$, as in step (1) of the proof of Lemma~\ref{sha-siegel-step-restriction}; then $g_v^{-1}sg_v\in D_\ell$ for $s\in D_v(F)$.
\item[(b)] \emph{The flags.} The matrix $r_0(g_v)$ is symplectic with multiplier $a(\theta(g_v))\ne0$ and supported in $W_{\mathrm r}$, and $n_0$ is supported in $W_{\mathrm n}$. By Lemma~\ref{sha-siegel-step-chains-rational-flags}(1), there is a lower unipotent $n_v$, supported in $W_{\mathrm n}$, with $n_v^*r_0(g_v)n_0$ upper triangular.
\item[(c)] \emph{The diagonal characters.} Let $s\in D_v(F)$ and $s'=g_v^{-1}sg_v\in D_\ell$. Then $r_0(s)=r_0(g_v)r_0(s')r_0(g_v^{-1})$ with $r_0(g_v)r_0(g_v^{-1})=1$, and $n_0^*r_0(s')n_0$ is upper triangular with diagonal entries $a(\chi_i(s'))=a(\chi_i(s))$, as $c_0$ is residual. By Lemma~\ref{sha-siegel-step-chains-rational-flags}(2), $n_v^*r_0(s)n_v$ is upper triangular with diagonal entries $a(\chi_i(s))$.
\item[(d)] \emph{A flagged pair.} Let $r$ be the restriction of $r_0$ to $G_F$ and $s_1=(r,(n_v))$. Each $r(g)$ is symplectic with multiplier $a(\theta(g))$, the $n_v$ are lower unipotent, (c) gives the flags, and $r(g_0^T)$ is diagonal. Let $g\in I_w(F)$ with $w$ not above $\ell$ and $\bar r_0(g)=1$. If $w$ lies above $p\in S_b$, then $\psi(g)=1$ by (i), and for $h\in G$ with $hv_p=w$, $h^{-1}gh\in I_p\cap\ker\psi$, so $r_0(h^{-1}gh)=1$ and $r(g)=1$. Otherwise $w$ lies above no prime of $S_b\cup\{\ell\}$, $g$ lies in the set $N_\Phi$ of the rational problem, and $r(g)=r_0(g)=1$. Let $s\in\mu_v$, so that $s^{\ell^m}$ lies in the closure of the commutator subgroup of $D_v(F)$ for some $m$. The homomorphism $g\mapsto a(\chi_i(g))$ from $G_F$ to $k^\times$ is locally constant, as $c_0$ is residual, so its kernel is closed and contains $s^{\ell^m}$. As $k$ has characteristic $\ell$, $(a(\chi_i(s))-1)^{\ell^m}=a(\chi_i(s))^{\ell^m}-1=0$, so $\psi_{v,i}(s)=a(\chi_i(s))=1$ by (c).
\item[(e)] \emph{Residual and over $c_0$.} The entries of $r$ are locally constant, $r(g)_{ii}=a(\chi_i(g))$, (c) gives $\psi_{v,i}(s)=a(\chi_i(s))$ for $s\in D_v(F)$, $r$ is supported in $W_{\mathrm r}$ and the $n_v$ in $W_{\mathrm n}$; so $s_1$ is residual. It is over $c_0$ by (b), as $g_v$ carries $v_\ell$ to $v$.
\end{enumerate}

(2) Write $c_0=(r_0,n_0)$ and $s'=(r',(n'_v))$, so that $r'$ is the restriction of $r_0$ to $G_F$.
\begin{enumerate}
\item[(a)] \emph{Auxiliary primes.} For a finite set $W$ of primes, the set $N_\Phi$ of the problem of $T$ with auxiliary primes $W$ is contained in that of the problem of $T$, and the other conditions are the same; so $s'$ is residual for it.
\item[(b)] \emph{The flags.} Let $x\in M_4(k)$ be supported in $W_{\mathrm r}$ with $x_{11},x_{22}\ne0$, let $n$ be lower unipotent and supported in $W_{\mathrm n}$, and let $n'$ be lower unipotent with $U=n'^*xn$ upper triangular. As $n'n'^*=1$, $xn=n'U$. The entries $(1,2)$, $(1,3)$, $(1,4)$, $(2,4)$, $(3,2)$ and $(3,4)$ of $x$ vanish, their weights being outside $W_{\mathrm r}$, and $n_{12}=n_{32}=0$. So the second column of $xn$ is $(0,x_{22},0,(xn)_{42})$. Comparing the first rows gives $U_{11}=x_{11}$ and $U_{12}=0$, so the second column of $n'U$ is $U_{22}$ times that of $n'$; hence $U_{22}=x_{22}\ne0$ and $n'_{32}=0$. Comparing the first columns,
\[
n'_{31}x_{11}=x_{31}+x_{33}n_{31},\qquad n'_{21}x_{11}=x_{21}+x_{22}n_{21}+x_{23}n_{31}.
\]
\item[(c)] \emph{$s'$ lies in $W_{\mathrm r}$.} The homomorphism $r'$ is supported in $W_{\mathrm r}$, as $r_0$ is. For $v\in\Sigma_T$, let $g\in G$ carry $v_\ell$ to $v$ with $n_v'^*r_0(g)n_0$ upper triangular. By (b), with $x=r_0(g)$, $n=n_0$ and $n'=n'_v$, $(n'_v)_{32}=0$; so $n'_v$ is supported in $W_{\mathrm n}$.
\item[(d)] \emph{Vanishing on $G$.} Let $i\ne j$ be such that $r_0(gh)_{ij}=r_0(g)_{ii}r_0(h)_{ij}+r_0(g)_{ij}r_0(h)_{jj}$ for all $g,h\in G$, and suppose that $r'(g)_{ij}=0$ for all $g\in G_F$. Put $f(g)=r_0(g)_{ij}$ and $h=g_0^T\in G_F$. As $F$ is Galois over $\mathbf Q$, $ghg^{-1}\in G_F$, so $f(ghg^{-1})=0$ for all $g$, and $f(h)=0$. So $f(gh)=f\big((ghg^{-1})g\big)=a(\chi_i(h))f(g)$ and $f(gh)=f(g)a(\chi_j(h))$. By (l), $a(\chi_i(h))\ne a(\chi_j(h))$. So $f(g)=0$ for all $g\in G$.
\item[(e)] \emph{The first case of the lattice shape.} Suppose that $r_0$ has an entry of weight $(0,-1)$. It also has one of weight $(-1,1)$. The entries of these weights are $(3,1)$, $(4,2)$ and $(2,3)$, and each satisfies the relation of (d): the terms $r_0(g)_{3x}r_0(h)_{x1}$ with $x\in\{2,4\}$, $r_0(g)_{4x}r_0(h)_{x2}$ with $x\in\{1,3\}$, and $r_0(g)_{2x}r_0(h)_{x3}$ with $x\in\{1,4\}$ vanish, as $r_0(g)_{32}$, $r_0(g)_{34}$, $r_0(h)_{12}$, $r_0(h)_{32}$, $r_0(h)_{13}$ and $r_0(g)_{24}$ do. By (d), $r'$ has entries of weights $(0,-1)$ and $(-1,1)$.
\item[(f)] \emph{The second case.} Otherwise $r_0$ is supported in $W_{\mathrm V}$, $(n_0)_{31}\ne0$, and no $m\in k$ has $r_0(s)_{23}=m\big(a(\chi_3(s))-a(\chi_2(s))\big)$ for all $s\in D_\ell$. Then $r'$ is supported in $W_{\mathrm V}$. For $v$ and $g$ as in (c), (b) gives $(n'_v)_{31}r_0(g)_{11}=r_0(g)_{33}(n_0)_{31}\ne0$, as $r_0(g)_{31}=0$; so $(n'_v)_{31}\ne0$, and $s'$ is of type V. Suppose that $r'$ has no entry of weight $(-1,0)$, so that $r'(g)_{21}=0$ for $g\in G_F$. The entry $(2,1)$ satisfies the relation of (d), as $r_0(h)_{31}=r_0(g)_{24}=0$; so $r_0(g)_{21}=0$ for all $g\in G$. For $s\in D_\ell$, (b) with $x=r_0(s)$ and $n=n'=n_0$ gives $(n_0)_{31}r_0(s)_{11}=r_0(s)_{33}(n_0)_{31}$, so $r_0(s)_{11}=r_0(s)_{33}$, and then $(n_0)_{21}r_0(s)_{33}=r_0(s)_{22}(n_0)_{21}+r_0(s)_{23}(n_0)_{31}$. So $r_0(s)_{23}=m\big(a(\chi_3(s))-a(\chi_2(s))\big)$ with $m=(n_0)_{21}/(n_0)_{31}$, a contradiction. So $r'$ has an entry of weight $(-1,0)$.
\end{enumerate}
% (1) is Chain.exists_residual_datum_over_restriction_to_rationals (139 lines, with Rq.placeMove, Rt.flagNew and Rt.flagMove): (a) is placeMove, (b) is Rt.flagNew (FlagNew.lean, 139 lines, Lemma sha-siegel-step-chains-rational-flags (1)), (c) is Rt.flagMove (FlagMove.lean, 47 lines, its (2)), (d) is the kill case through Tower.Ok.kb and the mu case through the closed kernel and the Frobenius of k. (2) is goodQ (FlagFacts.lean, 141 lines): (b) is the private flag_facts, (d) is the private coc_zero and conj_mem with Ok.dist for (l), (e) and (f) are the case split on LatShape.ps. The hypothesis of (2) on g_0 is not used by goodQ. About 300 lines with these helpers.
\end{proof}

\begin{definition}\label{sha-siegel-step-chains-gysin-cochains}
Let $D$ be a Siegel datum, $T$ a tower for $D$ with field $F$, $k$ a field and $a:\mathcal O\to k$ a homomorphism that kills $\mathfrak m_{\mathcal O}$, and let $\Phi$ be the problem of $T$ with auxiliary primes $W$ for a finite set $W$ of primes (Definition~\ref{sha-siegel-step-chains-notions}), with its set $N_\Phi$. Let $w,w'\in\mathbf Z^2$.
\begin{enumerate}
\item A \emph{cocycle of weight $w$} for $\Phi$ is a locally constant $x:G_F\to k$ with $x(gh)=x(g)+a(\chi^w(g))x(h)$ for all $g,h\in G_F$ and $x(g)=0$ for $g\in N_\Phi$. It is \emph{locally trivial with constants} $(m_v)_{v\in\Sigma_T}$, $m_v\in k$, if $x(s)=m_v-a(\chi^w(s))m_v$ for every $v\in\Sigma_T$ and $s\in D_v(F)$.
\item Let $x$ be a cocycle of weight $w$, locally trivial with constants $(m_v)$, and $y$ a cocycle of weight $w'$. A \emph{cup cochain} for $(x,y)$ is a locally constant $z:G_F\to k$, vanishing on $N_\Phi$, with
\[
z(gh)=z(g)+a(\chi^{w+w'}(g))z(h)+y(g)a(\chi^w(g))x(h)\qquad(g,h\in G_F),
\]
such that each $v\in\Sigma_T$ has $o_v\in k$ with $z(s)=o_v-a(\chi^{w+w'}(s))o_v-y(s)a(\chi^w(s))m_v$ for $s\in D_v(F)$; the $o_v$ are its \emph{constants}.
\end{enumerate}
\end{definition}
\leanref{Development/AutomorphicArithmetic/AutomorphicVarieties/WeakShimuraTheory/RamifiedLineCharacters/MemRamLine.lean}{UpperTheory.V.Ws.Fr.Coc}
\leanref{Development/AutomorphicArithmetic/AutomorphicVarieties/WeakShimuraTheory/RamifiedLineCharacters/MemRamLine.lean}{UpperTheory.V.Ws.Fr.Cup}
\leanref{Development/AutomorphicArithmetic/AutomorphicVarieties/WeakShimuraTheory/RedIdeal.lean}{UpperTheory.V.Ws.Fr.LocTriv}

\begin{lemma}\label{sha-siegel-step-chains-gysin-data}
Let $D$, $T$, $F$, $k$, $a$ and $\Phi$ be as in Definition~\ref{sha-siegel-step-chains-gysin-cochains}, and let $p=(r,(n_v))$ be a residual flagged pair for $\Phi$ along $a$ with $r$ supported in $W_{\mathrm r}$. Then $y(g)=r(g)_{31}a(\chi_1(g))^{-1}$ is a cocycle of weight $(0,-1)$ for $\Phi$, locally trivial with constants $((n_v)_{31})_v$, with $y(g_0^T)=0$.
\end{lemma}
\leanref{Development/AutomorphicArithmetic/AutomorphicVarieties/WeakShimuraTheory/RamifiedLineCharacters/DataPs.lean}{UpperTheory.V.Ws.Gys.dataPs}
\begin{proof}
The function $y$ is locally constant, as the entries of $r$ are. As $r(g)_{32}=r(g)_{34}=0$, $r(gh)_{31}=r(g)_{31}r(h)_{11}+r(g)_{33}r(h)_{31}$; dividing by $a(\chi_1(gh))=a(\chi_1(g))a(\chi_1(h))$ gives $y(gh)=y(g)+a(\chi^{(0,-1)}(g))y(h)$, as $\chi^{(0,-1)}=\chi_3\chi_1^{-1}$. For $g\in N_\Phi$, $r(g)=1$, so $y(g)=0$; and $r(g_0^T)$ is diagonal, so $y(g_0^T)=0$. For $s\in D_v(F)$, $r(s)n_v=n_vU$ with $U=n_v^*r(s)n_v$ upper triangular and $U_{11}=a(\chi_1(s))$; the entry $(3,1)$ gives $r(s)_{31}+r(s)_{33}(n_v)_{31}=(n_v)_{31}a(\chi_1(s))$, that is $y(s)=(n_v)_{31}-a(\chi^{(0,-1)}(s))(n_v)_{31}$.
\end{proof}
% dataPs (DataPs.lean, 66 lines), for any Fr with IsRes and Supp nP; y is the Lean's p.r g 2 0 * Ring.inverse (p.r g 0 0).

\begin{lemma}\label{sha-siegel-step-chains-gysin-raw-module}
Let $D$, $T$, $F$, $k$, $a$ and $\Phi$ be as in Definition~\ref{sha-siegel-step-chains-gysin-cochains}, let $w,w'\in\mathbf Z^2$, and let $y_j$, for $j$ in a finite set $J$, be cocycles of weight $w'$ for $\Phi$. Then there are a finite $\mathbf F_\ell$-vector space $M$ with a continuous linear action $\rho$ of $G_F$, additive maps $p,q_j:M\to k$ for $j\in J$, and $m_0\in M$, such that:
\begin{enumerate}
\item $\rho(g)=1$ for every $g\in G_F$ with $a(\chi_i(g))=1$ for all $i$ and $y_j(g)=0$ for all $j$;
\item $p(\rho(g)u)=a(\chi^w(g))p(u)$ and $q_j(\rho(g)u)=a(\chi^{w+w'}(g))q_j(u)+y_j(g)a(\chi^w(g))p(u)$ for $g\in G_F$ and $u\in M$;
\item $m_0\ne0$, and $p(u)\ne0$ for every nonzero $u\in\mathbf F_\ell m_0$.
\end{enumerate}
\end{lemma}
\leanref{Development/AutomorphicArithmetic/AutomorphicVarieties/WeakShimuraTheory/RamifiedLineCharacters/ExtRep.lean}{UpperTheory.V.Ws.Gys.extRep}
\begin{proof}
Write $\gamma=a\circ\chi^w$ and $\gamma'=a\circ\chi^{w+w'}$; then $\gamma'=\gamma\cdot(a\circ\chi^{w'})$, and $y_j(1)=0$ by the cocycle relation at $(1,1)$.
\begin{enumerate}
\item[(a)] \emph{The module.} The image $A=a(\mathcal O)$ is a finite subring of $k$, as $a$ kills $\mathfrak m_{\mathcal O}$. Each $y_j$ is locally constant on the compact group $G_F$, so it takes finitely many values; let $V\subseteq k$ be the $A$-submodule they span, which is finite. Let $M=A\times V^J$, an $\mathbf F_\ell$-vector space as $k$ has characteristic $\ell$, and for $g\in G_F$ let
\[
\rho(g)\big(u,(v_j)_j\big)=\big(\gamma(g)u,\ (\gamma'(g)v_j+\gamma(g)u\,y_j(g))_j\big).
\]
\item[(b)] \emph{An action.} As $y_j(1)=0$, $\rho(1)=1$. For $g,h\in G_F$ the second component of $\rho(g)\rho(h)(u,v)$ is $\gamma'(gh)v_j+\gamma(gh)u\big(y_j(g)+a(\chi^{w'}(g))y_j(h)\big)=\gamma'(gh)v_j+\gamma(gh)u\,y_j(gh)$, by the cocycle relation for $y_j$; so $\rho(gh)=\rho(g)\rho(h)$.
\item[(c)] \emph{Continuity and (1).} If $a(\chi_i(g))=1$ for all $i$, then $\gamma(g)=\gamma'(g)=1$, as $\chi^w$ and $\chi^{w+w'}$ are products of powers of the $\chi_i$; if also $y_j(g)=0$ for all $j$, then $\rho(g)=1$. The set $U$ of such $g$ is an open neighbourhood of $1$, as the $a\circ\chi_i$ and the $y_j$ are locally constant, and $\rho(gh)=\rho(g)$ for $h\in U$; so $g\mapsto\rho(g)u$ is locally constant for each $u$, and $\rho$ is continuous.
\item[(d)] \emph{The maps.} Let $p(u,v)=u$, $q_j(u,v)=v_j$ and $m_0=(1,0)$. Then (2) is the definition of $\rho$, and $p(cm_0)=c$ for $c\in\mathbf F_\ell$, which gives (3).
\end{enumerate}
\end{proof}
% extRep (ExtRep.lean, 192 lines): (a) is A = a.range, V the span of the ranges of the y j, M = A x (J -> V) with the private actM and mkRep; (b) is hm with hsplit; (c) is hker and the private continuous_of_right_invariant_on_nhds_one; (d) is p, q and m0 = (1, 0).

\begin{lemma}\label{sha-siegel-step-chains-gysin-raw-local}
Let $\ell$ be a prime, $F\subset\bar{\mathbf Q}$ a number field, $v$ a finite place of $F$ not above $\ell$, $A$ a place of $\bar{\mathbf Q}$ above $v$, and $D\supseteq I$ the decomposition and inertia groups of $A$ in $\operatorname{Gal}(\bar{\mathbf Q}/F)$. Suppose that $\chi_{\mathrm{cyc}}$ is trivial modulo $\ell$ on $D$. Let $M$ be a finite $\mathbf F_\ell$-vector space with the trivial action of $D$, let $m_0\in M$ be nonzero, and let $L\subseteq H^1(D,M)$ be the set of classes of continuous cocycles $f:D\to M$ with $f(g)\in\mathbf F_\ell m_0$ for every $g\in I$. Then $L$ has at least $2|M|$ elements.
\end{lemma}
\leanref{Development/AutomorphicArithmetic/AutomorphicVarieties/WeakShimuraTheory/RamifiedLineCharacters/CharRep.lean}{UpperTheory.V.Ws.Gys.locNe}
\leanref{Development/AutomorphicArithmetic/AutomorphicVarieties/WeakShimuraTheory/RamifiedLineCharacters/LocHoms.lean}{UpperTheory.V.Ws.Gys.locHoms}
\leanref{Development/AutomorphicArithmetic/AutomorphicVarieties/WeakShimuraTheory/RamifiedLineCharacters/RamCard.lean}{UpperTheory.V.Ws.Gys.ramCard}
\begin{proof}
For the trivial action of $D$ on a group, cocycles are continuous homomorphisms and coboundaries vanish, so distinct cocycles have distinct classes. As $\chi_{\mathrm{cyc}}$ is trivial modulo $\ell$ on $D$, the module $\mathbf F_\ell^*(1)$ of Lemma~\ref{sha-descent-link-greenberg-wiles}, the linear forms on $\mathbf F_\ell$ with $D$ acting through $\chi_{\mathrm{cyc}}$ modulo $\ell$, is $\mathbf F_\ell$ with the trivial action. Call a homomorphism $D\to\mathbf F_\ell$ unramified if it vanishes on $I$.
\begin{enumerate}
\item[(a)] \emph{A nonzero homomorphism.} By Lemma~\ref{sha-descent-link-local}(2), as $v$ is not above $\ell$, $|H^1(D,\mathbf F_\ell)|=|H^0(D,\mathbf F_\ell)|\,|H^2(D,\mathbf F_\ell)|\ge\ell$. So there is a nonzero continuous homomorphism $f:D\to\mathbf F_\ell$.
\item[(b)] \emph{Duality.} By local Tate duality the cup product pairing $H^1(D,\mathbf F_\ell^*(1))\times H^1(D,\mathbf F_\ell)\to\mathbf F_\ell$ is perfect, and, as $v$ is not above $\ell$, the unramified classes of the two groups are the exact annihilators of each other \cite[Chapter~VII, \S2]{NSW08}.
\item[(c)] \emph{An unramified homomorphism.} Suppose that every unramified homomorphism $D\to\mathbf F_\ell$ is $0$. Then the unramified subgroup of $H^1(D,\mathbf F_\ell)$ is $0$, its annihilator is all of $H^1(D,\mathbf F_\ell^*(1))$, and by (b) every class of $H^1(D,\mathbf F_\ell^*(1))$ is unramified. Applied to the class of $f$, this shows that $f$ vanishes on $I$, so $f=0$, a contradiction. So there is an unramified $\varphi_1\ne0$.
\item[(d)] \emph{A ramified homomorphism.} Suppose that every homomorphism $D\to\mathbf F_\ell$ vanishes on $I$. Then every class of $H^1(D,\mathbf F_\ell)$ and of $H^1(D,\mathbf F_\ell^*(1))$ is unramified, so by (b) the class of $f$ in $H^1(D,\mathbf F_\ell^*(1))$ is orthogonal to all of $H^1(D,\mathbf F_\ell)$, hence $0$ by the perfectness, and $f=0$, a contradiction. So there is a homomorphism $\varphi_2$ that does not vanish on $I$.
\item[(e)] \emph{The count.} Scaling, $\varphi_1(g_1)=1$ for some $g_1\in D$ and $\varphi_2(g_2)=1$ for some $g_2\in I$. For $t\in\{0,1\}$ and $m\in M$ the map $f_{t,m}(g)=\varphi_1(g)m+t\varphi_2(g)m_0$ is a cocycle, as $D$ acts trivially, with values in $\mathbf F_\ell m_0$ on $I$. If $f_{t,m}=f_{t',m'}$, then evaluating at $g_2$ gives $tm_0=t'm_0$, so $t=t'$, and evaluating at $g_1$ gives $m=m'$. So the classes of the $f_{t,m}$ are $2|M|$ distinct elements of $L$.
\end{enumerate}
\end{proof}
% (a) is locNe (CharRep.lean, 26 lines, LocalEulerChar at v not above l); (b) to (d) are locHoms (LocHoms.lean, 77 lines): (b) is LocalTateDuality (bijectivity of cupH1H1TwistedDual) and UnramifiedOrthogonality, (c) is ha with hU2, (d) is hb with hU1 and hinj, then rescale; (e) is ramCard (RamCard.lean, 46 lines), with the cocycles fn t m.

\begin{lemma}\label{sha-siegel-step-chains-gysin-raw-split}
Let $D$ be a Siegel datum, $T$ a tower for $D$ with field $F$, and $U\subseteq G_F$ an open subgroup (Definition~\ref{sha-siegel-step-chains-notions}). There are infinitely many finite places $v$ of $F$ such that $v$ lies above a prime that splits in $T$ and $D_A(F)\subseteq U$ for every place $A$ of $\bar{\mathbf Q}$ above $v$.
\end{lemma}
\leanref{Development/AutomorphicArithmetic/AutomorphicVarieties/WeakShimuraTheory/RamifiedLineCharacters/PlacesQ.lean}{UpperTheory.V.Ws.Gys.placesQ}
\leanref{Development/AutomorphicArithmetic/AutomorphicVarieties/WeakShimuraTheory/RamifiedLineCharacters/SplitPrimes.lean}{UpperTheory.V.Ws.Gys.splitPrimes}
\leanref{Development/AutomorphicArithmetic/AutomorphicVarieties/PadicHodgeTheory/V00.lean}{UpperTheory.V.Loc.v00}
\leanref{Development/AutomorphicArithmetic/AutomorphicVarieties/GaloisLibrary/KummerTheory/DecGen.lean}{UpperTheory.V.Kum.decGen}
\begin{proof}
\begin{enumerate}
\item[(a)] \emph{An open subgroup of $G$.} As $G_F$ is open in $G$, the image of $U$ is open in $G$. Let $H$ be its intersection with the kernels of $\bar r_0$ and $\omega$, an open subgroup of $G$.
\item[(b)] \emph{Split primes.} Infinitely many primes $p$ have $D_A(\mathbf Q)\subseteq H$ for every place $A$ above $p$. Indeed, let $H'\subseteq H$ be an open normal subgroup of $G$, and suppose that only finitely many primes have this property for $H'$; let $N$ be their product. By the Chebotarev density theorem (Theorem~\ref{lit-chebotarev-degree-one}), applied to the fixed field of $H'$, the identity element and the primes dividing $N$, there are a prime $p\nmid N$, a place $A$ above $p$ with $I_A(\mathbf Q)\subseteq H'$, and a Frobenius element $\sigma\in D_A(\mathbf Q)$ with $\sigma\in H'$. As $D_A(\mathbf Q)$ is the closure of the subgroup generated by $I_A(\mathbf Q)$ and $\sigma$ \cite[Chapter~II, \S9]{Neu99}, and $H'$ is closed, $D_A(\mathbf Q)\subseteq H'$. The places above $p$ are conjugate to $A$ under $G$, and $H'$ is normal, so $D_{A'}(\mathbf Q)\subseteq H'$ for every $A'$ above $p$. So $p$ divides $N$, a contradiction.
\item[(c)] \emph{Places of $F$.} A prime $p$ as in (b) splits in $T$, as $D_A(\mathbf Q)\subseteq H\subseteq G_F$ and $\bar r_0$ and $\omega$ are trivial on $H$. Choose a place $v_p$ of $F$ above each such $p$; distinct primes give distinct places. For $A$ above $v_p$, $D_A(F)\subseteq D_A(\mathbf Q)$, which lies in the image of $U$; so $D_A(F)\subseteq U$.
\end{enumerate}
\end{proof}
% placesQ (PlacesQ.lean, 94 lines): (a) is H with galEmbTop; (b) is splitPrimes (SplitPrimes.lean, 39 lines: an open normal N inside H, Loc.v00 with g = 1 from GlobalTheory.Published.cheb_holds, Kum.decGen for D_A(Q) in N, Chr.placesQ for the conjugate places); (c) is the private ex_place and prime_eq.

\begin{lemma}\label{sha-siegel-step-chains-gysin-raw-count}
Let $\ell$ be an odd prime, $F\subset\bar{\mathbf Q}$ a number field and $G_F=\operatorname{Gal}(\bar{\mathbf Q}/F)$, and for a finite place $v$ of $F$ and a place $A$ above it write $D_A(F)\supseteq I_A(F)$ for the decomposition and inertia groups. Let $M$ be a finite $\mathbf F_\ell$-vector space with a continuous linear action $\rho$ of $G_F$, and $S_0$ a finite set of finite places of $F$ that contains the places above $\ell$, such that $\rho$ is trivial on $I_A(F)$ for every place $A$ above a place not in $S_0$. Let $Q$ be an infinite set of finite places of $F$ such that $\rho$ and $\chi_{\mathrm{cyc}}$ modulo $\ell$ are trivial on $D_A(F)$ for every $v\in Q$ and $A$ above $v$. Let $m_0\in M$ be nonzero.
Then there are a finite set $S\subseteq S_0\cup Q$ and a continuous cocycle $c:G_F\to M$ such that: $c$ vanishes on $I_A(F)$ for every place $A$ above a place not in $S$; for every $v\in S_0$ and $A$ above $v$ there is $m\in M$ with $c(g)=\rho(g)m-m$ for $g\in D_A(F)$; and some $g\in G_F$ has $\rho(g)=1$, $c(g)\ne0$ and $c(g)\in\mathbf F_\ell m_0$.
\end{lemma}
\leanref{Development/AutomorphicArithmetic/AutomorphicVarieties/WeakShimuraTheory/RamifiedLineCharacters/GysCore.lean}{UpperTheory.V.Ws.Gys.gysCore}
\leanref{Development/AutomorphicArithmetic/AutomorphicVarieties/WeakShimuraTheory/RamifiedLineCharacters/ZunrFin.lean}{UpperTheory.V.Ws.Gys.zunrFin}
\leanref{Development/AutomorphicArithmetic/AutomorphicVarieties/WeakShimuraTheory/RamifiedLineCharacters/SelBig.lean}{UpperTheory.V.Ws.Gys.selBig}
\leanref{Development/AutomorphicArithmetic/AutomorphicVarieties/WeakShimuraTheory/RamifiedLineCharacters/GwLower.lean}{UpperTheory.V.Ws.Gys.gwLower}
\leanref{Development/AutomorphicArithmetic/AutomorphicVarieties/WeakShimuraTheory/RamifiedLineCharacters/SelMem.lean}{UpperTheory.V.Ws.Gys.selMem}
\leanref{Development/AutomorphicArithmetic/AutomorphicVarieties/WeakShimuraTheory/RamifiedLineCharacters/CocConj.lean}{UpperTheory.V.Ws.Gys.cocConj}
\begin{proof}
For a set $S$ of finite places of $F$, let $G_S$ be the quotient of $G_F$ by the closed normal subgroup generated by the $I_A(F)$ with $A$ above a place not in $S$, as in Lemma~\ref{sha-descent-link-global}.
\begin{enumerate}
\item[(a)] \emph{Moving places.} Let $c$ be a continuous cocycle, $v$ a place of $F$, and $A$, $A'=sA$ places above $v$, with $s\in G_F$. If $c(g)=\rho(g)m-m$ for $g\in D_A(F)$, then $c(g)=\rho(g)m'-m'$ for $g\in D_{A'}(F)=sD_A(F)s^{-1}$, with $m'=\rho(s)m-c(s)$; and if $\rho$ is trivial and $c$ vanishes on $I_A(F)$, then $c$ vanishes on $I_{A'}(F)$. Both follow from the cocycle relation applied to $g=s(s^{-1}gs)s^{-1}$.
\item[(b)] \emph{Finiteness.} Let $Z_0$ be the group of continuous cocycles $c:G_F\to M$ that vanish on $I_A(F)$ for every $A$ above a place not in $S_0$. For such $c$ the set of $g$ with $\rho(g)=1$ and $c(g)=0$ is a closed normal subgroup, so it contains the kernel of $G_F\to G_{S_0}$, and $\rho$ and $c$ factor through $G_{S_0}$. As $H^1(G_{S_0},M)$ is finite (Lemma~\ref{sha-descent-link-global-tate}(1) with $r=1$), and so is the group of coboundaries, $Z_0$ is finite. Let $N=|Z_0|$.
\item[(c)] \emph{The places.} Let $M^*$ be the dual of $M$ and $C=|M^*|\,|M|^{[F:\mathbf Q]+|S_0|}$. As $Q$ is infinite, there is a set $S_1\subseteq Q$ of $NC$ places not in $S_0$. Let $S=S_0\cup S_1$; then $\rho$ factors through $G_S$. For $v\in S$ choose a place $A_v$ above $v$, and write $D_v\supseteq I_v$ for $D_{A_v}(F)\supseteq I_{A_v}(F)$.
\item[(d)] \emph{The Selmer group.} For $v\in S_0$ let $L_v=0$, and for $v\in S_1$ let $L_v\subseteq H^1(D_v,M)$ be the set of classes of cocycles with values in $\mathbf F_\ell m_0$ on $I_v$. Let $\mathrm{Sel}\subseteq H^1(G_S,M)$ be the set of classes whose restriction to $D_v$ lies in $L_v$ for every $v\in S$.
\item[(e)] \emph{Its size.} By Lemma~\ref{sha-descent-link-greenberg-wiles}, $\prod_{v\in S}|L_v|$ divides $|\mathrm{Sel}|\,|H^0(G_S,M^*(1))|\prod_{v\in S}|H^0(D_v,M)|\prod_{w\text{ real}}|M^{c_w}|\,|M|^{r_2}$. Here $|H^0(G_S,M^*(1))|\le|M^*|$, $\prod_{v\in S}|H^0(D_v,M)|\le|M|^{|S|}$, and the archimedean factor is at most $|M|^{r_1+r_2}\le|M|^{[F:\mathbf Q]}$. For $v\in S_1$, $v$ is not above $\ell$, and $D_v$ acts trivially on $M$ and on $\mu_\ell$, so $|L_v|\ge2|M|$ by Lemma~\ref{sha-siegel-step-chains-gysin-raw-local}. Hence $(2|M|)^{NC}\le|\mathrm{Sel}|\,|M^*|\,|M|^{[F:\mathbf Q]}|M|^{|S_0|}|M|^{NC}$, that is $2^{NC}\le|\mathrm{Sel}|\,C$.
\item[(f)] \emph{Elements of $\mathrm{Sel}$.} Let $f$ be a cocycle of $G_S$ whose class lies in $\mathrm{Sel}$, and $c$ its composite with $G_F\to G_S$. Then $c$ is a continuous cocycle that vanishes on $I_A(F)$ for $A$ above a place not in $S$. For $v\in S_0$, $c$ is a coboundary on $D_{A_v}(F)$, hence on $D_A(F)$ for every $A$ above $v$, by (a). For $v\in S_1$ and $g\in I_v$, $c(g)\in\mathbf F_\ell m_0$: the restriction of $c$ to $D_v$ differs from a cocycle with values in $\mathbf F_\ell m_0$ on $I_v$ by a coboundary $\rho(g)m-m$, which vanishes as $D_v$ acts trivially.
\item[(g)] \emph{Conclusion.} Suppose that no $S$ and $c$ as in the statement exist. For each $x\in\mathrm{Sel}$ choose a cocycle $f_x$ of $G_S$ representing it, and let $c_x$ be its composite with $G_F\to G_S$. Let $v\in S_1$ and $g\in I_v$. Then $\rho(g)=1$ and $c_x(g)\in\mathbf F_\ell m_0$, so $c_x(g)=0$, since otherwise $S$ and $c_x$ would satisfy the conclusion, by (f). By (a), $c_x$ vanishes on $I_A(F)$ for every $A$ above $v$. So $c_x\in Z_0$, and $x\mapsto c_x$ is injective. Hence $|\mathrm{Sel}|\le N$, and by (e) $2^{NC}\le NC$, which is false.
\end{enumerate}
\end{proof}
% gysCore (GysCore.lean, 61 lines): (a) is cocConj (CocConj.lean, 57 lines); (b) is zunrFin (ZunrFin.lean, 86 lines, the private killSub and module_finite_z1_padicInt, GlobalFiniteness) with descend (Descend.lean, 37 lines); (c) is C, N, S_1 and S, with descend for rho_S; (d) is gysSel with gysL and ramLine (MemRamLine.lean); (e) is selBig (SelBig.lean, 42 lines) through gwLower (GwLower.lean, 51 lines, GreenbergWiles with exists_lift_cycMod and ccEx for the complex conjugations) and locHoms with ramCard (Lemma sha-siegel-step-chains-gysin-raw-local); (f) is selMem (SelMem.lean, 49 lines); (g) is hz, hinj and the final count with Nat.lt_two_pow_self.

\begin{lemma}\label{sha-siegel-step-chains-gysin-raw-cocycles}
Let $D$ be a Siegel datum, $T$ a tower for $D$ with field $F$, $W$ a finite set of primes that split in $T$, and $N_\Phi$ the set of the problem of $T$ with auxiliary primes $W$ (Definition~\ref{sha-siegel-step-chains-notions}). Let $M$ be a finite $\mathbf F_\ell$-vector space with a continuous linear action $\rho$ of $G_F$ that is trivial on $N_\Phi$, and let $m_0\in M$ be nonzero.
Then there are a finite set $W'\supseteq W$ of primes that split in $T$ and a continuous cocycle $c:G_F\to M$ such that: $c$ vanishes on the set $N_{\Phi'}$ of the problem of $T$ with auxiliary primes $W'$; for each $v\in\Sigma_T$ there is $m_v\in M$ with $c(s)=\rho(s)m_v-m_v$ for $s\in D_v(F)$; and some $g\in G_F$ has $\rho(g)=1$, $c(g)\ne0$ and $c(g)\in\mathbf F_\ell m_0$.
\end{lemma}
\leanref{Development/AutomorphicArithmetic/AutomorphicVarieties/WeakShimuraTheory/RamifiedLineCharacters/OpenKer.lean}{UpperTheory.V.Ws.Gys.gysT}
\leanref{Development/AutomorphicArithmetic/AutomorphicVarieties/WeakShimuraTheory/RamifiedLineCharacters/MapRep.lean}{UpperTheory.V.Ws.Gys.unrS}
\begin{proof}
\begin{enumerate}
\item[(a)] \emph{The places $S_0$.} Let $S_0$ be the set of finite places of $F$ above the primes of $S\cup W$, a finite set containing the places above $\ell$. Let $A$ be a place above a place of $F$ not in $S_0$, and $g\in I_A(F)$. Then $A$ lies above no prime of $S\cup W$, so $\bar r_0(g)=1$, as $\psi$, $\xi$, $\omega$ and $\chi_{\mathrm{cyc}}$ are unramified outside $S$ (Definition~\ref{sha-siegel-datum}); so $g\in N_\Phi$ and $\rho(g)=1$.
\item[(b)] \emph{The places $Q$.} Let $Q$ be the set of finite places $v$ of $F$ above a prime that splits in $T$ such that $\rho$ and $\chi_{\mathrm{cyc}}$ modulo $\ell$ are trivial on $D_A(F)$ for every $A$ above $v$. The kernels of $\rho$ and of $\chi_{\mathrm{cyc}}$ modulo $\ell$ are open, so $Q$ is infinite by Lemma~\ref{sha-siegel-step-chains-gysin-raw-split}.
\item[(c)] \emph{The cocycle.} As $\ell$ is odd, Lemma~\ref{sha-siegel-step-chains-gysin-raw-count} gives $S\subseteq S_0\cup Q$, a cocycle $c$ and an element $g$. Each $v\in S$ not in $S_0$ lies in $Q$, so above a prime $p_v$ that splits in $T$. Let $W'=W\cup\{p_v:v\in S\setminus S_0\}$.
\item[(d)] \emph{$c$ vanishes on $N_{\Phi'}$.} Let $g\in I_A(F)$ with $A$ above no prime of $W'\cup\{\ell\}$ and $\bar r_0(g)=1$. Then $g\in N_\Phi$, as $W\subseteq W'$, so $\rho(g)=1$. Let $v$ be the place of $F$ below $A$. If $v\in S_0$, then $c$ is a coboundary $\rho(h)m-m$ on $D_A(F)\supseteq I_A(F)$, so $c(g)=0$. If $v\in S\setminus S_0$, then $A$ lies above $p_v\in W'$, which is excluded. If $v\notin S$, then $c(g)=0$ by (c).
\item[(e)] \emph{The places of $\Sigma_T$.} For $v\in\Sigma_T$, the place of $F$ below $v$ lies above $\ell$, hence in $S_0$, so $c$ is a coboundary on $D_v(F)$.
\end{enumerate}
\end{proof}
% gysT (OpenKer.lean, 98 lines): (a) is S_0 from the private fin_over and hunr with Gys.unrS (MapRep.lean, from the facts of D.Good); (b) is Gys.placesQ applied to the private open_ker and open_cyc; (c) is gysCore and hnext (Tower.Next); (d) is the first bullet, with Tower.Next.sub and the case A = top (where Ine A is trivial); (e) is the second bullet.

\begin{lemma}\label{sha-siegel-step-chains-gysin-raw}
Let $D$ be a Siegel datum, $T$ a fit tower for $D$ with field $F$, $k$ a field and $a:\mathcal O\to k$ a homomorphism that kills $\mathfrak m_{\mathcal O}$, and $W$ a finite set of primes that split in $T$. Let $w,w'\in\mathbf Z^2$, and let $y_j$, for $j$ in a finite set $J$, be cocycles of weight $w'$ for the problem of $T$ with auxiliary primes $W$ (Definition~\ref{sha-siegel-step-chains-gysin-cochains}).
Then there are a finite set $W'\supseteq W$ of primes that split in $T$, a cocycle $x$ of weight $w$ for the problem of $T$ with auxiliary primes $W'$, locally trivial with constants $(m_v)$, such that $x(g)\ne0$ for some $g$ with $a(\chi^w(g))=1$, and for each $j$ a cup cochain for $(x,y_j)$ for the problem of $T$ with auxiliary primes $W'$.
\end{lemma}
\leanref{Development/AutomorphicArithmetic/AutomorphicVarieties/WeakShimuraTheory/RamifiedLineCharacters/CupRaw.lean}{UpperTheory.V.Ws.Gys.cupRaw}
\leanref{Development/AutomorphicArithmetic/AutomorphicVarieties/WeakShimuraTheory/RamifiedLineCharacters/OpenKer.lean}{UpperTheory.V.Ws.Gys.gysT}
\begin{proof}
Let $\Phi$ and $\Phi'$ be the problems of $T$ with auxiliary primes $W$ and $W'$, with sets $N_\Phi$ and $N_{\Phi'}$, once $W'$ is chosen.
\begin{enumerate}
\item[(a)] \emph{The module.} Lemma~\ref{sha-siegel-step-chains-gysin-raw-module} gives $M$, $\rho$, $p$, the $q_j$ and $m_0$. For $g\in N_\Phi$, $\bar r_0(g)=1$, so $a(\chi_i(g))=1$ for all $i$, and $y_j(g)=0$; so $\rho$ is trivial on $N_\Phi$.
\item[(b)] \emph{The cocycle.} Lemma~\ref{sha-siegel-step-chains-gysin-raw-cocycles} gives $W'\supseteq W$, a cocycle $c$, elements $m_v$ for $v\in\Sigma_T$, and $g$ with $\rho(g)=1$ and $c(g)$ a nonzero element of $\mathbf F_\ell m_0$.
\item[(c)] \emph{The cocycle $x$.} Let $x=p\circ c$. It is locally constant, as $c$ is continuous with values in the finite discrete $M$; $x(gh)=p(c(g)+\rho(g)c(h))=x(g)+a(\chi^w(g))x(h)$; and $x$ vanishes on $N_{\Phi'}$. For $s\in D_v(F)$, $x(s)=p(\rho(s)m_v-m_v)=m'_v-a(\chi^w(s))m'_v$ with $m'_v=-p(m_v)$. So $x$ is a cocycle of weight $w$ for $\Phi'$, locally trivial with constants $(m'_v)$.
\item[(d)] \emph{The element.} $x(g)=p(c(g))\ne0$ by (3) of Lemma~\ref{sha-siegel-step-chains-gysin-raw-module}, and $p(c(g))=p(\rho(g)c(g))=a(\chi^w(g))p(c(g))$, so $a(\chi^w(g))=1$.
\item[(e)] \emph{The cup cochains.} Let $z_j=q_j\circ c$. It is locally constant and vanishes on $N_{\Phi'}$, and by (2) of Lemma~\ref{sha-siegel-step-chains-gysin-raw-module},
\[
z_j(gh)=q_j(c(g))+q_j(\rho(g)c(h))=z_j(g)+a(\chi^{w+w'}(g))z_j(h)+y_j(g)a(\chi^w(g))x(h).
\]
For $s\in D_v(F)$, $z_j(s)=q_j(\rho(s)m_v-m_v)=o_v-a(\chi^{w+w'}(s))o_v-y_j(s)a(\chi^w(s))m'_v$ with $o_v=-q_j(m_v)$. As $N_{\Phi'}\subseteq N_\Phi$, $y_j$ is a cocycle for $\Phi'$, and $z_j$ is a cup cochain for $(x,y_j)$ for $\Phi'$.
\end{enumerate}
% cupRaw (CupRaw.lean, 50 lines), for any Fr with Ok along a and the property Fr.Gys l (T.Next W): (a) is Gys.extRep (Lemma sha-siegel-step-chains-gysin-raw-module) with Ok.kill; (b) is the property Fr.Gys, which is gysT (Lemma sha-siegel-step-chains-gysin-raw-cocycles, from CohomologyHyp and D.Good); (c) to (e) are the four bullets of cupRaw, with -p (mm v) and -q j (mm v) for the constants. With ExtRep, OpenKer, GysCore, ZunrFin, Descend, SelBig, GwLower, SelMem, CocConj, LocHoms, CharRep, RamCard, CcOne, PlacesQ, SplitPrimes and MapRep, about 1,300 lines in RamifiedLineCharacters, besides v00 and decGen.
\end{proof}

\begin{lemma}\label{sha-siegel-step-chains-gysin-cup}
Let $D$, $T$, $k$, $a$, $W$, $w$, $w'$, $J$ and the $y_j$ be as in Lemma~\ref{sha-siegel-step-chains-gysin-raw}, and suppose that $y_j(g_0^T)=0$ for every $j$, $a(\chi^w(g_0^T))\ne1$ and $a(\chi^{w+w'}(g_0^T))\ne1$.
Then there are a finite set $W'\supseteq W$ of primes that split in $T$, a cocycle $x$ of weight $w$ for the problem of $T$ with auxiliary primes $W'$, locally trivial with constants $(m_v)$, with $x(g_0^T)=0$ and $x(g)\ne0$ for some $g$, and for each $j$ a cup cochain $z_j$ for $(x,y_j)$ with $z_j(g_0^T)=0$, for the problem of $T$ with auxiliary primes $W'$; and the $y_j$ are cocycles for that problem.
\end{lemma}
\leanref{Development/AutomorphicArithmetic/AutomorphicVarieties/WeakShimuraTheory/RamifiedLineCharacters/GysinPsOf.lean}{UpperTheory.V.Ws.Gys.cupEx}
\leanref{Development/AutomorphicArithmetic/AutomorphicVarieties/WeakShimuraTheory/RamifiedLineCharacters/CupNorm.lean}{UpperTheory.V.Ws.Gys.cupNorm}
\begin{proof}
Lemma~\ref{sha-siegel-step-chains-gysin-raw} gives $W'$, a cocycle $x_0$ with constants $(m^0_v)$ and an element $g$ with $a(\chi^w(g))=1$ and $x_0(g)\ne0$, and cup cochains $z^0_j$ for $(x_0,y_j)$ with constants $(o^0_{j,v})$. The set $N_\Phi$ only shrinks when $W$ grows, so the $y_j$ are cocycles for the problem with auxiliary primes $W'$. Write $\gamma=a\circ\chi^w$ and $\gamma'=a\circ\chi^{w+w'}$, homomorphisms $G_F\to k^\times$ that are trivial on $N_\Phi$.
\begin{enumerate}
\item[(a)] \emph{Normalizing $x$.} As $\gamma(g_0^T)\ne1$, there is $m_1\in k$ with $x_0(g_0^T)=m_1-\gamma(g_0^T)m_1$. Let $x(h)=x_0(h)-m_1+\gamma(h)m_1$ and $m_v=m^0_v-m_1$. Then $x$ is a cocycle of weight $w$, vanishing on $N_\Phi$ and at $g_0^T$, locally trivial with constants $(m_v)$, and $x(g)=x_0(g)\ne0$.
\item[(b)] \emph{Normalizing $z_j$.} As $\gamma'(g_0^T)\ne1$, there is $o_1\in k$ with $z^0_j(g_0^T)=o_1-\gamma'(g_0^T)o_1$. Let $z_j(h)=z^0_j(h)-o_1+\gamma'(h)o_1+y_j(h)\gamma(h)m_1$. A direct computation with the relations for $x_0$, $y_j$ and $z^0_j$ shows that $z_j$ is a cup cochain for $(x,y_j)$, with constants $o^0_{j,v}-o_1$; it vanishes on $N_\Phi$, and $z_j(g_0^T)=0$ as $y_j(g_0^T)=0$.
\end{enumerate}
\end{proof}
% cupEx (GysinPsOf.lean): coefFacts, std for Ok, cupRaw with gysT (Lemma sha-siegel-step-chains-gysin-raw), and cupNorm (CupNorm.lean, 64 lines) for (a) and (b).

\begin{lemma}\label{sha-siegel-step-chains-gysin-matrix}
Let $D$, $T$, $F$, $k$, $a$ and $\Phi$ be as in Definition~\ref{sha-siegel-step-chains-gysin-cochains}, with $T$ fit, and suppose that $a(\chi_i(s))=1$ for every $v\in\Sigma_T$, $s\in\mu_v$ and $i$. For $x,y,t\in k$ let $u(x,y,t)$ be the matrix with rows $(1,0,0,0)$, $(x,1,0,0)$, $(y,0,1,0)$ and $(t,y,-x,1)$. Let $x$ be a cocycle of weight $(-1,0)$ for $\Phi$, locally trivial with constants $(m_v)$, $y$ a cocycle of weight $(0,-1)$, locally trivial with constants $(q_v)$, and $z$ a cup cochain for $(x,y)$ with constants $(o_v)$, such that $x$, $y$ and $z$ vanish at $g_0^T$. Put
\[
r(g)=u\big(x(g),y(g),2z(g)-x(g)y(g)\big)\operatorname{diag}\big(a(\chi_1(g)),\dots,a(\chi_4(g))\big),\qquad n_v=u(m_v,q_v,2o_v-m_vq_v).
\]
Then $(r,(n_v))$ is a residual flagged pair for $\Phi$ along $a$ that lies in $W_{\mathrm n}$, with $r(g)_{21}=x(g)a(\chi_1(g))$, $(n_v)_{21}=m_v$, $r(g)_{31}=y(g)a(\chi_1(g))$ and $(n_v)_{31}=q_v$.
\end{lemma}
\leanref{Development/AutomorphicArithmetic/AutomorphicVarieties/WeakShimuraTheory/RamifiedLineCharacters/SuppMMono.lean}{UpperTheory.V.Ws.Gys.ob}
\begin{proof}
\begin{enumerate}
\item[(a)] \emph{The matrices $u$.} A computation gives $u(x,y,t)u(x',y',t')=u(x+x',y+y',t+t'+yx'-xy')$ and $u(x,y,t)^*=u(-x,-y,-t)$, so $u(x,y,t)^*u(x,y,t)=u(0,0,0)=1$: each $u(x,y,t)$ is lower unipotent, and it is supported in $W_{\mathrm n}$. If $d\in k^4$ has $d_2=\rho d_1$, $d_3=\sigma d_1$ and $d_4=\rho\sigma d_1$, then $\operatorname{diag}(d)u(x,y,t)=u(\rho x,\sigma y,\rho\sigma t)\operatorname{diag}(d)$.
\item[(b)] \emph{A homomorphism.} Let $D_g=\operatorname{diag}(a(\chi_i(g)))$, $\rho_g=a(\chi^{(-1,0)}(g))$ and $\sigma_g=a(\chi^{(0,-1)}(g))$. As $\chi^{(-1,0)}=\chi_2\chi_1^{-1}$, $\chi^{(0,-1)}=\chi_3\chi_1^{-1}$ and $\chi_1\chi_4=\chi_2\chi_3$, (a) applies to $D_g$. Write $t(g)=2z(g)-x(g)y(g)$. By (a), $r(g)r(h)=u\big(x(g)+\rho_gx(h),\,y(g)+\sigma_gy(h),\,t(g)+\rho_g\sigma_gt(h)+\rho_gy(g)x(h)-\sigma_gx(g)y(h)\big)D_gD_h$. The relations for $x$, $y$ and $z$ give $x(gh)=x(g)+\rho_gx(h)$, $y(gh)=y(g)+\sigma_gy(h)$ and $t(gh)=t(g)+\rho_g\sigma_gt(h)+\rho_gy(g)x(h)-\sigma_gx(g)y(h)$, as $\chi^{(-1,-1)}=\chi^{(-1,0)}\chi^{(0,-1)}$. So $r(gh)=r(g)r(h)$; and $x(1)=y(1)=z(1)=0$, so $r(1)=1$.
\item[(c)] \emph{A flagged pair.} By (a), $r(g)^*r(g)=D_g^*D_g=a(\theta(g))$, as $\chi_i\chi_{5-i}=\theta$. For $g\in N_\Phi$, $x(g)=y(g)=z(g)=0$, and $a(\chi_i(g))=1$ as $\bar r_0(g)=1$; so $r(g)=1$. As $x$, $y$ and $z$ vanish at $g_0^T$, $r(g_0^T)=D_{g_0^T}$ is diagonal. Let $s\in D_v(F)$. Then $x(s)=m_v-\rho_sm_v$, $y(s)=q_v-\sigma_sq_v$ and $z(s)=o_v-\rho_s\sigma_so_v-\rho_sy(s)m_v$, and a computation with (a) gives $n_v^*r(s)n_v=u(-m_v,-q_v,-t_v)\,u(x(s),y(s),t(s))\,u(\rho_sm_v,\sigma_sq_v,\rho_s\sigma_st_v)\,D_s=D_s$, where $t_v=2o_v-m_vq_v$. So $n_v^*r(s)n_v$ is upper triangular with diagonal entries $a(\chi_i(s))$, which are $1$ for $s\in\mu_v$ by hypothesis.
\item[(d)] \emph{Residual.} The entries of $r$ are locally constant, as $x$, $y$, $z$ and the $a\circ\chi_i$ are. The diagonal entries of $r(g)$ are $a(\chi_i(g))$, and (c) gives $\psi_{v,i}(s)=a(\chi_i(s))$ on $D_v(F)$. By (a), $r(g)$ and $n_v$ are supported in $W_{\mathrm n}\subseteq W_{\mathrm r}$. So $(r,(n_v))$ is residual and lies in $W_{\mathrm n}$. The entries $(2,1)$ and $(3,1)$ of $r(g)$ are $x(g)a(\chi_1(g))$ and $y(g)a(\chi_1(g))$, and those of $n_v$ are $m_v$ and $q_v$.
\end{enumerate}
\end{proof}
% ob (SuppMMono.lean, 232 lines with MemRamLine.lean for uQ, uQ_mul, adj_uQ, isLow_uQ, suppM_uQ and diag_mul_uQ): (b) is the private exists_hom_M4_eq_uQ_mul_diagonal_of_cup with chi1, chi2 and chi3, (c) is gsp_uD and the private adj_uQ_conj_uQ_mul_diagonal_eq_diagonal, (d) is the final refine. The Lean's hypotheses are Fr.Ok along a and the condition on mu_v.

\begin{lemma}\label{sha-siegel-step-chains-gysin}
Let $D$, $T$, $k$ and $a$ be as in Lemma~\ref{sha-siegel-step-chains-lines}, and let $W$ be a finite set of primes that split in $T$. Let $p$ and $p'$ be residual flagged pairs for the problem of $T$ with auxiliary primes $W$ along $a$ that lie in $W_{\mathrm r}$.
Then there are a finite set $W'\supseteq W$ of primes that split in $T$ and residual flagged pairs $s$ and $s'$ for the problem of $T$ with auxiliary primes $W'$ along $a$, both lying in $W_{\mathrm n}$, such that $s$ has the same $(0,-1)$-part as $p$, $s'$ the same $(0,-1)$-part as $p'$, $s$ and $s'$ have the same $(-1,0)$-part, and $s$ has an entry of weight $(-1,0)$.
\end{lemma}
\leanref{Development/AutomorphicArithmetic/AutomorphicVarieties/WeakShimuraTheory/CompactIntegralFields/GysinAl.lean}{UpperTheory.V.Ws.gysinAl}
\begin{proof}
Let $F$ be the field of $T$; cocycles and cup cochains are as in Definition~\ref{sha-siegel-step-chains-gysin-cochains}. Write $p=(r_p,(n_{p,v}))$ and $p'=(r_{p'},(n_{p',v}))$.
\begin{enumerate}
\item[(a)] \emph{The base point.} By (l), the $a(\chi_i(g_0^T))$ are pairwise distinct; so $a(\chi^{(-1,0)}(g_0^T))=a(\chi_2(g_0^T))/a(\chi_1(g_0^T))\ne1$ and $a(\chi^{(-1,-1)}(g_0^T))=a(\chi_4(g_0^T))/a(\chi_1(g_0^T))\ne1$.
\item[(b)] \emph{The $(0,-1)$-parts.} By Lemma~\ref{sha-siegel-step-chains-gysin-data}, $y(g)=r_p(g)_{31}a(\chi_1(g))^{-1}$ and $y'(g)=r_{p'}(g)_{31}a(\chi_1(g))^{-1}$ are cocycles of weight $(0,-1)$ for the problem of $T$ with auxiliary primes $W$, vanishing at $g_0^T$ and locally trivial with constants $((n_{p,v})_{31})$ and $((n_{p',v})_{31})$.
\item[(c)] \emph{Auxiliary primes.} Lemma~\ref{sha-siegel-step-chains-gysin-cup}, with $w=(-1,0)$, $w'=(0,-1)$ and the cocycles $y$ and $y'$, which applies by (a), gives a finite set $W'\supseteq W$ of primes that split in $T$, a cocycle $x$ of weight $(-1,0)$ for the problem of $T$ with auxiliary primes $W'$, locally trivial with constants $(m_v)$, with $x(g_0^T)=0$ and $x(g)\ne0$ for some $g$, and cup cochains $z$ for $(x,y)$ and $z'$ for $(x,y')$ that vanish at $g_0^T$; and $y$ and $y'$ are cocycles for that problem.
\item[(d)] \emph{The pairs.} For $v\in\Sigma_T$ and $s\in\mu_v\subseteq D_v(F)$, $a(\chi_i(s))=1$: as $p$ is residual, the diagonal entries $\psi_{v,i}(s)$ of $n_{p,v}^*r_p(s)n_{p,v}$ are $a(\chi_i(s))$, and they are $1$ on $\mu_v$. So Lemma~\ref{sha-siegel-step-chains-gysin-matrix}, applied to $(x,y,z)$ and to $(x,y',z')$, gives residual flagged pairs $s$ and $s'$ for the problem of $T$ with auxiliary primes $W'$, lying in $W_{\mathrm n}$.
\item[(e)] \emph{The parts.} The entries $(3,1)$ of the homomorphism of $s$ are $y(g)a(\chi_1(g))=r_p(g)_{31}$, and those of its flags are $(n_{p,v})_{31}$; so $s$ has the same $(0,-1)$-part as $p$, and likewise $s'$ as $p'$. The entries $(2,1)$ of $s$ and of $s'$ are $x(g)a(\chi_1(g))$ and $m_v$ for both, so $s$ and $s'$ have the same $(-1,0)$-part. As $x(g)\ne0$ for some $g$, the homomorphism of $s$ has an entry of weight $(-1,0)$.
\end{enumerate}
% gysinAl (GysinAl.lean, 27 lines): (a) is Gys.rootCharacters_ne_one_at_basepoint (35 lines, from Ok.dist), (b) is Gys.dataPs (Lemma sha-siegel-step-chains-gysin-data), (c) is Gys.cupEx (Lemma sha-siegel-step-chains-gysin-cup, with its raw cocycles, Lemma sha-siegel-step-chains-gysin-raw), (d) is Gys.mu_of and Gys.ob (Lemma sha-siegel-step-chains-gysin-matrix), (e) is the final tuple. In WeakShimuraTheory/RamifiedLineCharacters, about 1,750 lines.
\end{proof}

\begin{lemma}\label{sha-siegel-step-chains-dies-cyclotomic}
Let $M\ge1$ be an integer, and let $V\subseteq G$ be an open subgroup that contains the commutator subgroup of $G$ and the inertia group $I_w(\mathbf Q)$ of every place $w$ of $\bar{\mathbf Q}$ above no prime dividing $M$. Then there is $n\ge0$ such that $V$ contains every $g\in G$ that fixes the $M^n$-th roots of unity.
\end{lemma}
\leanref{Development/AutomorphicArithmetic/AutomorphicVarieties/GaloisLibrary/ClassFieldRamification/KwAbsRam.lean}{UpperTheory.V.CF.kwAbsRam}
\leanref{Development/AutomorphicArithmetic/AutomorphicVarieties/GaloisLibrary/ClassFieldRamification/Kw.lean}{UpperTheory.V.CF.kw}
\leanref{Development/AutomorphicArithmetic/AutomorphicVarieties/GaloisLibrary/ClassFieldRamification/KwRamSplit.lean}{UpperTheory.V.CF.kwRam}
\leanref{Development/AutomorphicArithmetic/AutomorphicVarieties/GaloisLibrary/ClassFieldRamification/CycDescend.lean}{UpperTheory.V.CF.cycDescend}
\leanref{Development/AutomorphicArithmetic/AutomorphicVarieties/GaloisLibrary/ClassFieldRamification/UnitsGen.lean}{UpperTheory.V.CF.unitsGen}
\leanref{Development/AutomorphicArithmetic/CyclotomicBounds/ClassNumberBounds/InertiaRamification/CycInertia.lean}{KStar.Ram.map_inertia_eq}
\begin{proof}
\begin{enumerate}
\item[(a)] \emph{The field.} Let $L$ be the fixed field of $V$. As $V$ is open and contains the commutator subgroup, it is normal, and $L$ is a number field, Galois over $\mathbf Q$ with abelian Galois group $G/V$. Let $p$ be a prime not dividing $M$ and $w$ a place above $p$. Then $M$ is a unit at $w$, as $1=ap+bM$ for some integers $a,b$; so $w$ lies above no prime dividing $M$, and $I_w(\mathbf Q)\subseteq V$. The inertia groups of $L$ at the places above $p$ are the images of these $I_w(\mathbf Q)$ in $G/V$, so $L$ is unramified at $p$.
\item[(b)] \emph{Kronecker-Weber.} There is $N\ge1$ with $L\subseteq\mathbf Q(\mu_N)$ \cite[Theorem~14.1]{Was97}.
\item[(c)] \emph{Splitting $N$.} There are $a,b\ge1$ and $n\ge0$ with $N=ab$, $a\mid M^n$ and $b$ prime to $M$: by induction on $N$, if $N$ is not prime to $M$, apply the induction hypothesis to $N/q$ for a prime $q$ dividing $N$ and $M$, and replace the resulting $a$ by $qa$ and $n$ by $n+1$.
\item[(d)] \emph{Inertia in $\mathbf Q(\mu_N)$.} Let $C=\mathbf Q(\mu_N)$, and let $\chi:\operatorname{Gal}(C/\mathbf Q)\to(\mathbf Z/N)^\times$ be the isomorphism given by the action on $\mu_N$. For a prime $q$ with $q^e\,\|\,N$, $e\ge1$, the inertia groups of $C$ at the places above $q$ are mapped by $\chi$ onto the subgroup $U_q$ of the $u\equiv1$ modulo $N/q^e$ \cite[Chapter~2]{Was97}. For $q\mid b$, $q$ does not divide $M$ by (c), so $L$ is unramified at $q$ by (a), and these inertia groups fix $L$; so $\chi^{-1}(U_q)$ fixes $L$.
\item[(e)] \emph{Generation.} The subgroups $U_q$ with $q\mid b$ generate the kernel of $(\mathbf Z/N)^\times\to(\mathbf Z/a)^\times$. This is proved by induction on $b$. For $b=1$ the kernel is trivial. Otherwise let $q\mid b$, with $q^e\,\|\,b$, so that $q^e\,\|\,N$ as $a$ and $b$ are coprime, and put $a'=aq^e$ and $b'=b/q^e$. Let $u\equiv1$ modulo $a$. The greatest common divisor of $a'$ and $N/q^e$ is $a$, so by the Chinese remainder theorem there is $k$ with $k\equiv1$ modulo $a'$ and $k\equiv u$ modulo $N/q^e$; it is prime to $N$, as $N$ divides $a'\cdot N/q^e$. Then $u=k\cdot(k^{-1}u)$ with $k^{-1}u\in U_q$, and $k$ lies in the kernel of the map to $(\mathbf Z/a')^\times$, which is generated by the $U_{q'}$ with $q'\mid b'$ by induction.
\item[(f)] \emph{Descent.} Let $\sigma\in\operatorname{Gal}(C/\mathbf Q)$ fix $\mu_a$. Then $\chi(\sigma)\equiv1$ modulo $a$, so by (e) and (d), $\sigma$ fixes $L$. So $L\subseteq\mathbf Q(\mu_a)\subseteq\mathbf Q(\mu_{M^n})$.
\item[(g)] \emph{Conclusion.} If $g\in G$ fixes the $M^n$-th roots of unity, it fixes $L$ by (f), so $g\in V$.
\end{enumerate}
\end{proof}
% kwAbsRam (KwAbsRam.lean, 109 lines, the private abelian_ramification_cyclotomic_bound): (a) is hLV, hab and hunr (through unramIff and ineAbs); (b) to (f) are kwRamSub (KwRamSub.lean) through kwRam (KwRamSplit.lean): (b) is CF.kw (Kw.lean, through kwHom, from GroundTheory.Global.exists_le_rayClassField and kwRay), (c) is the private kwRam_split, (d) is KStar.Ram.map_inertia_eq (CycInertia.lean, 177 lines) with CF.ineFixQ, (e) is unitsGen (UnitsGen.lean, 119 lines, the private ker_le_sup and aux), (f) is cycDescend (CycDescend.lean, 79 lines); (g) is the stabilizer of the root of unity. About 600 lines.

\begin{lemma}\label{sha-siegel-step-chains-dies}
Let $D$ be a Siegel datum, $T$ a tame tower for $D$ satisfying (a) to (m) of Definition~\ref{sha-siegel-tower-towers}, and $W$ a finite set of primes that split in $T$ (Definition~\ref{sha-siegel-step-chains-notions}). There is a tame tower $T'\ge T$ that kills every prime of $W$.
\end{lemma}
\leanref{Development/AutomorphicArithmetic/AutomorphicVarieties/WeakShimuraTheory/CompactIntegralFields/DiesEx.lean}{UpperTheory.V.Ws.diesEx}
\begin{proof}
Let $F$ be the field of $T$.
\begin{enumerate}
\item[(a)] \emph{$\ell\notin W$.} Suppose that $\ell$ splits in $T$. The Frobenius element $\sigma_\ell$ lies in $D_{v_\ell}(\mathbf Q)$, so $\bar r_0(\sigma_\ell)=1$ and $\omega(\sigma_\ell)=1$. The entry $(3,3)$ gives $\bar\alpha(\sigma_\ell)=1$; as $\bar\alpha=\bar\chi_{\mathrm{cyc}}\bar\omega^{-1}\bar\xi$ and $\bar\chi_{\mathrm{cyc}}=\bar\omega$, this is $\bar\xi(\sigma_\ell)=1$, contrary to the condition $\overline{\xi\Psi^{0}}(\sigma_\ell)\ne1$ of Definition~\ref{sha-siegel-datum}. So every $p\in W$ differs from $\ell$.
\item[(b)] \emph{The subgroup $V$.} For $p\in W$ let $\Theta_p$ be the set of continuous homomorphisms $\theta:G\to\mathbf Z/\ell$ that are trivial on $I_w(\mathbf Q)$ for every place $w$ not above $p$, and let $V$ be the intersection of the kernels of the $\theta\in\Theta_p$, $p\in W$. Let $\theta\in\Theta_p$. Its kernel is open, contains the commutator subgroup of $G$ and the inertia groups $I_w(\mathbf Q)$ at the places $w$ not above $p$, so by Lemma~\ref{sha-siegel-step-chains-dies-cyclotomic} with $M=p$ it contains every element that fixes $\mu_{p^n}$, for some $n$. If $g\in G$ fixes $\mu_p$, then $g^{p^n}$ fixes $\mu_{p^{n+1}}$, by induction on $n$, so $p^n\theta(g)=0$; as $\ell\theta(g)=0$ and $p\ne\ell$ by (a), $\theta(g)=0$. So $V$ contains the open subgroup of the $g$ that fix $\mu_p$ for every $p\in W$, and $V$ is open. It contains $g^\ell$ for every $g\in G$, and the commutator subgroup of $G$.
\item[(c)] \emph{The field.} The subgroup $H=G_F\cap V$ is open and contains the commutator subgroup of $G$, as $F$ is abelian over $\mathbf Q$; so $H=G_{F'}$ for a number field $F'\supseteq F$, Galois and abelian over $\mathbf Q$. A complex conjugation $\sigma$ lies in $G_F$, as $F$ is totally real, and $\sigma=\sigma^\ell\in V$, as $\ell$ is odd; so $F'$ is totally real. For $g\in G_F$, $g^\ell\in G_{F'}$.
\item[(d)] \emph{The tower.} As in step (d) of the proof of Lemma~\ref{sha-siegel-step-towers}(3), with $F'$ and $(g_0^T)^\ell\in G_{F'}$, the relabelling of $T$, the field $K$, the field $F'$, the element $(g_0^T)^\ell$ and a choice of places above $\ell$ form a tower $T'$; by (c), $T\le T'$ with $j=j'=1$.
\item[(e)] \emph{$T'$ is tame.} Let $w$ be a place above $\ell$ and $g\in I_w(\mathbf Q)$. As $T$ is tame, $g^m\in G_F$ for some $m$ prime to $\ell$. For $p\in W$ and $\theta\in\Theta_p$, $\theta$ is trivial on $I_w(\mathbf Q)$, as $w$ is not above $p\ne\ell$; so $g\in V$, and $g^m\in G_F\cap V=G_{F'}$.
\item[(f)] \emph{$T'$ kills $W$.} For $p\in W$ and $\theta\in\Theta_p$, $G_{F'}\subseteq V\subseteq\ker\theta$.
\end{enumerate}
% diesEx (DiesEx.lean, 183 lines with its private helpers): (a) is hne with Raw.condD at i = 0; (b) is V with the private dx_open, dx_ker (through CF.kwAbsRam, Lemma sha-siegel-step-chains-dies-cyclotomic) and dx_mu; (c) is the private dx_field, with dx_conj for the complex conjugations; (d) is mkTower with Ok.split; (e) and (f) are the final refine.
\end{proof}

\begin{lemma}\label{sha-siegel-step-chains-restrict-kill-roots}
Let $w$ be a place of $\bar{\mathbf Q}$, with valuation ring $\mathcal O_w\subset\bar{\mathbf Q}$ and maximal ideal $\mathfrak m_w$, and let $m\ge1$ be an integer with $m\notin\mathfrak m_w$.
\begin{enumerate}
\item Let $\eta\in\bar{\mathbf Q}$ with $\eta^m=1$ and $\eta-1\in\mathfrak m_w$. Then $\eta=1$. In particular every element of $I_w(\mathbf Q)$ fixes the $m$-th roots of unity.
\item Let $\tau$ be an automorphism of $\bar{\mathbf Q}$ such that $\tau(y)/y-1\in\mathfrak m_w$ for every nonzero $y\in\bar{\mathbf Q}$. Then $\tau^i(y)/y-1\in\mathfrak m_w$ for every $i\ge0$ and every nonzero $y$, and every $x\in\bar{\mathbf Q}$ with $\tau^m(x)=x$ has $\tau(x)=x$.
\end{enumerate}
\end{lemma}
\leanref{Development/AbelianVarietyTheory/Proofs/FoundationProofs/FixedOfPowFixed.lean}{GlobalTheory.WildInertiaChar.wild_pow}
\leanref{Development/AbelianVarietyTheory/Proofs/FoundationProofs/FixedOfPowFixed.lean}{GlobalTheory.WildInertiaChar.fixed_of_pow_fixed}
\begin{proof}
(1) Suppose that $\eta\ne1$. Then $\sum_{i<m}\eta^i=0$, so $m=\sum_{i<m}(1-\eta^i)$. By induction on $i$, $1-\eta^i\in\mathfrak m_w$, as $1-\eta^{i+1}=(1-\eta^i)+\eta^i(1-\eta)$ and $\eta$ is a unit at $w$. So $m\in\mathfrak m_w$, a contradiction. Let $g\in I_w(\mathbf Q)$ and $\zeta^m=1$. Then $\zeta$ is a unit at $w$, so $g(\zeta)-\zeta\in\mathfrak m_w$, and $\eta=g(\zeta)/\zeta$ has $\eta^m=1$ and $\eta-1\in\mathfrak m_w$; so $g(\zeta)=\zeta$.

(2) For the first claim, by induction on $i$: with $z=\tau^i(y)$,
\[
\frac{\tau^{i+1}(y)}y-1=\Big(\frac{\tau(z)}z-1\Big)\Big(\frac{\tau^i(y)}y-1\Big)+\Big(\frac{\tau(z)}z-1\Big)+\Big(\frac{\tau^i(y)}y-1\Big).
\]
For the second, suppose that $\tau(x)\ne x$ and put $d=\tau(x)-x$. Then $\sum_{i<m}\tau^i(d)=\tau^m(x)-x=0$, so $m=-\sum_{i<m}\big(\tau^i(d)/d-1\big)$, which lies in $\mathfrak m_w$ by the first claim, a contradiction.
\end{proof}
% (1) is the private thetaEx_eta, thetaEx_pow and thetaEx_fix of Rt.thetaEx (ThetaEx.lean); (2) is GlobalTheory.WildInertiaChar.wild_pow and fixed_of_pow_fixed (FixedOfPowFixed.lean, 82 lines).

\begin{lemma}\label{sha-siegel-step-chains-restrict-kill-tame}
Let $E\subset\bar{\mathbf Q}$ be a number field and $w$ a place of $\bar{\mathbf Q}$, with valuation ring $\mathcal O_w$, maximal ideal $\mathfrak m_w$ and valuation $v$, written multiplicatively, and let $I=I_w(E)$. Call $\sigma\in I$ \emph{wild} if $\sigma(y)/y-1\in\mathfrak m_w$ for every nonzero $y\in\bar{\mathbf Q}$.
\begin{enumerate}
\item Let $\ell$ be a prime with $\ell\notin\mathfrak m_w$, and let $\psi:I\to N$ be a homomorphism to a group, with open kernel and with $\psi(\sigma)^\ell=1$ for all $\sigma\in I$. Then $\psi(\sigma)=1$ for every wild $\sigma$.
\item There is $\gamma\in I$ such that every closed subgroup of $I$ that contains $\gamma$ and every wild element is $I$.
\end{enumerate}
\end{lemma}
\leanref{Development/AbelianVarietyTheory/Proofs/FoundationProofs/ExistsTameGenAtOfValuationSubring.lean}{GlobalTheory.TameInertiaGen.exists_tameGenAt_of_valuationSubring}
\leanref{Development/AbelianVarietyTheory/Proofs/IwasawaDecomposition/ValuationApplyOfWild.lean}{GlobalTheory.TameCharUnique.mem_of_tameGen}
\begin{proof}
(1) Let $\sigma$ be wild. The kernel of $\psi$ is open in $I$, so it contains $I\cap G_L$ for a finite extension $L$ of $E$. As $L$ has finitely many embeddings into $\bar{\mathbf Q}$ over $E$, two powers $\sigma^i$ and $\sigma^j$ with $i<j$ agree on $L$, and $\sigma^n$ fixes $L$ for $n=j-i>0$. Write $n=\ell^am$ with $\ell\nmid m$. By Lemma~\ref{sha-siegel-step-chains-restrict-kill-roots}(2), $\sigma^m$ satisfies its hypothesis, and as $(\sigma^m)^{\ell^a}$ fixes $L$ and $\ell^a\notin\mathfrak m_w$, $\sigma^m$ fixes $L$. So $\psi(\sigma)^m=1$; with $\psi(\sigma)^\ell=1$ and $m$ prime to $\ell$, $\psi(\sigma)=1$.

(2) Let $U=1+\mathfrak m_w$, a subgroup of $\bar{\mathbf Q}^\times$.
\begin{enumerate}
\item[(a)] \emph{An integer.} Let $p$ be the prime with $v(p)<1$ if there is one, and $p=1$ otherwise. There is at most one such prime, and every nonzero integer $a$ has $v(a)=v(p)^k$ for some $k\ge0$: for a prime $q\ne p$, $1=xq+yp$ with integers $x,y$ shows that $v(q)=1$.
\item[(b)] \emph{Valuations.} Every nonzero $y\in\bar{\mathbf Q}$ has $v(y)^m=v(p)^k$ for some $m\ge1$ and $k\in\mathbf Z$. Indeed $\sum_ic_iy^i=0$ for some nonzero polynomial with integer coefficients; two of the nonzero terms $c_iy^i$ and $c_jy^j$, $i<j$, have the same largest valuation, and with (a) this gives $v(y)^{j-i}=v(c_i)/v(c_j)=v(p)^k$.
\item[(c)] \emph{Tame characters.} For $\sigma\in I$ and a unit $u$ of $\mathcal O_w$, $\sigma(u)/u\in U$, as $\sigma(u)-u\in\mathfrak m_w$. For $n\ge1$ choose $r_n\in\bar{\mathbf Q}$ with $r_n^{n!}=p$, and let $t_n(\sigma)$ be the class of $\sigma(r_n)/r_n$ in $\bar{\mathbf Q}^\times/U$. As $c=\tau(r_n)/r_n$ is an $n!$-th root of unity, hence a unit, $\sigma\tau(r_n)/r_n=\big(\sigma(c)/c\big)\,c\,\sigma(r_n)/r_n$ gives $t_n(\sigma\tau)=t_n(\sigma)t_n(\tau)$. So $t_n$ is a homomorphism on $I$, its image is cyclic, being contained in the image of the $n!$-th roots of unity, and its fibres are closed. If $t_{n+1}(\sigma)=1$, then $t_n(\sigma)=1$: as $v(r_n)=v(r_{n+1})^{n+1}$, $r_n=u\,r_{n+1}^{n+1}$ with $u$ a unit, and $\sigma(r_n)/r_n=(\sigma(u)/u)\big(\sigma(r_{n+1})/r_{n+1}\big)^{n+1}\in U$.
\item[(d)] \emph{A generator.} The group $I$ is compact. For each $n$ the set $S_n$ of $\gamma\in I$ such that $t_n(\gamma)$ generates $t_n(I)$ is closed and nonempty, and $S_{n+1}\subseteq S_n$ by (c). So some $\gamma$ lies in every $S_n$. Let $\sigma\in I$. For each $n$ the set of $h$ in the closure of the powers of $\gamma$ with $t_n(h)=t_n(\sigma)$ is closed and nonempty, as it contains a power of $\gamma$, and these sets decrease with $n$, by (c). So there is $h$ in the closure of the powers of $\gamma$ with $t_n(h^{-1}\sigma)=1$ for all $n$.
\item[(e)] \emph{The quotient is wild.} Let $\sigma'=h^{-1}\sigma$ and $y\ne0$. By (b), $v(y)^m=v(p)^k$ with $m\ge1$, so $v(y)^m=v(r_m)^{m!k}$, and as the value group is torsion free, $v(y)=v(r_m)^{(m-1)!k}$. So $y=u\,r_m^{(m-1)!k}$ with $u$ a unit, and $\sigma'(y)/y=(\sigma'(u)/u)\big(\sigma'(r_m)/r_m\big)^{(m-1)!k}\in U$, as $t_m(\sigma')=1$. So $\sigma'$ is wild.
\item[(f)] \emph{Conclusion.} Let $S$ be a closed subgroup of $I$ containing $\gamma$ and the wild elements, and $\sigma\in I$. With $h$ as in (d), $h\in S$ as $S$ is closed, and $h^{-1}\sigma\in S$ by (e); so $\sigma\in S$.
\end{enumerate}
\end{proof}
% (1) is the private Rt.wild_mem_ker of TameExp.lean (with krullTopology_mem_nhds_one_iff); its twin GlobalTheory.WildInertiaChar.trivial_on_wild is not used here. (2) is exists_tameGenAt_of_valuationSubring (ExistsTameGenAtOfValuationSubring.lean, 268 lines): (a) is exists_nat_gen, (b) exists_pow_eq, (c) cob, cob_mul, cob_mem_of_unit, cob_mem_of_zpow, tameChar, isCyclic_range and isClosed_preimage, (d) exists_topGen, (e) the last part of exists_tameGenAt_of_valuationSubring, which gives TameGenAt A gamma; (f) is TameCharUnique.mem_of_tameGen (ValuationApplyOfWild.lean). About 400 lines.

\begin{lemma}\label{sha-siegel-step-chains-restrict-kill}
Let $D$ be a Siegel datum, $T$ a tower for $D$ with field $F$, $k$ a field and $a:\mathcal O\to k$ a homomorphism that kills $\mathfrak m_{\mathcal O}$, and $W$ a finite set of primes. Let $s=(r,(n_v))$ be a residual flagged pair for the problem of $T$ with auxiliary primes $W$ along $a$ (Definition~\ref{sha-siegel-step-chains-notions}) such that $\ell\ge5$ or $r$ is supported in $W_{\mathrm n}$. Let $p$ be a prime that splits in $T$, let $F'$ be a subfield of $\bar{\mathbf Q}$ such that every continuous homomorphism $G\to\mathbf Z/\ell$ that is trivial on $I_w(\mathbf Q)$ for every place $w$ not above $p$ is trivial on $G_{F'}$, and let $w$ be a place above $p$ and not above $\ell$.
Then $r(g)=1$ for every $g\in I_w(F)\cap G_{F'}$.
\end{lemma}
\leanref{Development/AutomorphicArithmetic/AutomorphicVarieties/WeakShimuraTheory/GradedFlagMoves/SplitMod.lean}{UpperTheory.V.Ws.Rt.splitMod}
\leanref{Development/AutomorphicArithmetic/AutomorphicVarieties/WeakShimuraTheory/GradedFlagMoves/ThetaEx.lean}{UpperTheory.V.Ws.Rt.thetaEx}
\leanref{Development/AutomorphicArithmetic/AutomorphicVarieties/WeakShimuraTheory/GradedFlagMoves/UnipPow.lean}{UpperTheory.V.Ws.Rt.unipPow}
\leanref{Development/AutomorphicArithmetic/AutomorphicVarieties/WeakShimuraTheory/GradedFlagMoves/TameExp.lean}{UpperTheory.V.Ws.Rt.tameExp}
\begin{proof}
\begin{enumerate}
\item[(a)] \emph{$p\equiv1$ modulo $\ell$.} As $w$ is not above $\ell$, $p\ne\ell$, and $D_w(\mathbf Q)$ contains an element $\gamma$ with $\chi_{\mathrm{cyc}}(\gamma)=p$. As $p$ splits in $T$, $\omega(\gamma)=1$, so $p\equiv\bar\chi_{\mathrm{cyc}}(\gamma)=\bar\omega(\gamma)=1$ modulo $\ell$.
\item[(b)] \emph{A character of order $\ell$.} By (a), $\ell$ divides the order $p-1$ of the cyclic group $(\mathbf Z/p)^\times$, so there is a surjection $(\mathbf Z/p)^\times\to\mathbf Z/\ell$. Composed with the character $c_p:G\to(\mathbf Z/p)^\times$ by which $G$ acts on $\mu_p$, it gives a homomorphism $\theta_0:G\to\mathbf Z/\ell$ with open kernel. It is trivial on $I_{w'}(\mathbf Q)$ for every place $w'$ not above $p$, as $p$ is a unit at $w'$ and $I_{w'}(\mathbf Q)$ fixes $\mu_p$ by Lemma~\ref{sha-siegel-step-chains-restrict-kill-roots}(1). It is nontrivial on $I_w(\mathbf Q)$, as $c_p$ maps $I_w(\mathbf Q)$ onto $(\mathbf Z/p)^\times$, $\mathbf Q(\mu_p)$ being totally ramified at $p$ \cite[Chapter~2]{Was97}. By hypothesis, $\theta_0$ is trivial on $G_{F'}$.
\item[(c)] \emph{Exponent $\ell$.} Let $\sigma\in I_w(F)$. As $p$ splits in $T$, $\bar r_0$ is trivial on $D_w(\mathbf Q)$, so the diagonal entries $a(\chi_i(\sigma))$ of $r(\sigma)$ are $1$. Let $N=r(\sigma)-1$, and let $h$ be $(3,1,2,0)$ if $r$ is supported in $W_{\mathrm r}$, $(3,2,1,0)$ if in $W_{\mathrm m}$, and $(2,1,1,0)$ if in $W_{\mathrm n}$. In each case every entry $(i,j)$ with $i\ne j$ whose weight lies in the set has $h_i<h_j$, so $N_{ij}=0$ whenever $h_j\le h_i$, and by induction $(N^m)_{ij}=0$ whenever $h_j<h_i+m$. As $s$ is residual, $r$ is supported in $W_{\mathrm r}$ or in $W_{\mathrm m}$; if $\ell\ge5$ the entries of $h$ for these sets are less than $\ell$, and otherwise $r$ is supported in $W_{\mathrm n}$, whose $h$ has entries less than $3\le\ell$. So $N^\ell=0$, and $r(\sigma)^\ell=1+N^\ell=1$, as $k$ has characteristic $\ell$.
\item[(d)] \emph{Tame inertia.} Let $\varphi$ and $\theta$ be the restrictions of $r$ and $\theta_0$ to $I_w(F)$. They are homomorphisms with open kernels, $\varphi(\sigma)^\ell=1$ by (c), and $\theta$ is nontrivial: some $g'\in I_w(\mathbf Q)$ has $\theta_0(g')\ne0$, and $g'\in G_F$ as $D_w(\mathbf Q)\subseteq G_F$. As $w$ is not above $\ell$, $\ell$ is a unit at $w$, so by Lemma~\ref{sha-siegel-step-chains-restrict-kill-tame}(1), $\varphi$ and $\theta$ are trivial on the wild elements of $I_w(F)$. Let $\gamma\in I_w(F)$ be as in (2) of that lemma. Then $\theta(\gamma)\ne0$: otherwise the closed subgroup $\ker\theta$ would contain $\gamma$ and the wild elements, and be $I_w(F)$. The set $S$ of $\sigma\in I_w(F)$ such that $(\varphi(\sigma),\theta(\sigma))$ is a power of $(\varphi(\gamma),\theta(\gamma))$ is a subgroup containing $\ker\varphi\cap\ker\theta$, so it is open, hence closed; it contains $\gamma$ and the wild elements; so $S=I_w(F)$. So for each $\sigma\in I_w(F)$ some integer $m$ has $\varphi(\sigma)=\varphi(\gamma)^m$ and $\theta(\sigma)=m\theta(\gamma)$. As $\theta(\gamma)$ has order $\ell$, $\theta(\sigma)=0$ implies $\ell\mid m$ and $\varphi(\sigma)=1$.
\item[(e)] \emph{Conclusion.} For $g\in I_w(F)\cap G_{F'}$, $\theta(g)=0$ by (b), so $r(g)=\varphi(g)=1$ by (d).
\end{enumerate}
\end{proof}
% The private killW of RestrT.lean: (a) is Rt.splitMod (SplitMod.lean, 35 lines, through Chr.decCyc); (b) is Rt.thetaEx (ThetaEx.lean, 114 lines: thetaEx_hom for the surjection, thetaEx_fix for the places not above p, and B1.cycIne, the cited inertia of Q(mu_p)) and Dies; (c) is Rt.unipPow (UnipPow.lean, 50 lines, with the private pow_graded and pow_char_graded and the gradings ![3,1,2,0], ![3,2,1,0], ![2,1,1,0]); (d) is Rt.tameExp (TameExp.lean, 122 lines), through the private wild_mem_ker and GlobalTheory.TameInertiaGen (Lemma sha-siegel-step-chains-restrict-kill-tame), with S the comap of the zpowers of (phi gamma, theta gamma).

\begin{lemma}\label{sha-siegel-step-chains-restrict-flags}
Let $k$ be a field.
\begin{enumerate}
\item Let $B\in GL_4(k)$, and let $n,m\in M_4(k)$ be lower unipotent (Lemma~\ref{sha-siegel-step-restriction-matrices}) with $n^*B$ and $m^*B$ upper triangular. Then $n=m$.
\item Let $D$ be a Siegel datum, $a:\mathcal O\to k$ a homomorphism, and $c=(r,(n_v))$ a flagged pair along $a$ either for a problem of a tower $T$ with field $E$ and places $\Sigma_T$, or for a rational problem, with $E=\mathbf Q$ and the one place $v_\ell$ (Definitions~\ref{sha-siegel-tower-pairs} and~\ref{sha-siegel-step-chains-notions}). Let $T'$ be a tower, with $T\le T'$ in the first case, and $c'=(r',(n'_{v'}))$ a flagged pair along $a$ for a problem of $T'$ that is over $c$. If $v'\in\Sigma_{T'}$, $v$ is a place of the problem of $c$ and $g\in G_E$ carries $v$ to $v'$, then $n_{v'}'^*r(g)n_v$ is upper triangular.
\item In (2), every flagged pair along $a$ for a problem of $T'$ that is over $c$ is $c'$.
\end{enumerate}
\end{lemma}
\begin{proof}
(1) As $n^*n=1$, $(n^*)^{-1}=n$, so $m^*n=(m^*B)(n^*B)^{-1}$, which is upper triangular. It is also lower triangular with diagonal entries $1$, as $m^*$ and $n$ are. So $m^*n=1$, and $n=mm^*n=m$.

(2) As $c'$ is over $c$, there are a place $v_2$ of the problem of $c$ and $g_2\in G_E$ carrying $v_2$ to $v'$ with $n_{v'}'^*r(g_2)n_{v_2}$ upper triangular. In the first case $v$ and $v_2$ are conjugate under $G_E$, and $\Sigma_T$ meets each orbit once; so $v_2=v$, as in the second case. Then $g_2^{-1}g$ fixes $v$, so it lies in the decomposition group of $G_E$ at $v$, and $n_v^*r(g_2^{-1}g)n_v$ is upper triangular. As $n_vn_v^*=1$,
\[
n_{v'}'^*r(g)n_v=\big(n_{v'}'^*r(g_2)n_v\big)\big(n_v^*r(g_2^{-1}g)n_v\big),
\]
a product of upper triangular matrices.

(3) Let $c''=(r'',(n''_{v'}))$ be over $c$. Both $r''$ and $r'$ are the restriction of $r$ to $G_{F'}$. For $v'\in\Sigma_{T'}$ choose $v$ and $g$ as in the definition of $c''$ being over $c$; by (2), $n_{v'}'^*r(g)n_v$ is also upper triangular. The matrix $B=r(g)n_v$ is invertible, with inverse $n_v^*r(g^{-1})$, so $n''_{v'}=n'_{v'}$ by (1).
\end{proof}
% The private eq_of_isLow_of_isUp_adj_mul_of_isUnit_det (1), over_fact_of_isLow_of_flag with inv_mul_mem_dec_of_both_map_place (2), and over_pd_unique_of_isLow_of_flag (3) of OverFacts.lean, about 150 lines.

\begin{lemma}\label{sha-siegel-step-chains-restrict-powers}
Let $k$ be a field of characteristic $\ell$, $\Gamma$ a group, $r:\Gamma\to M_4(k)$ a homomorphism and $i\ne j$, such that $x\mapsto r(x)_{ii}$ and $x\mapsto r(x)_{jj}$ are multiplicative with nonzero values, and $r(xz)_{ij}=r(x)_{ij}r(z)_{jj}+r(x)_{ii}r(z)_{ij}$ for all $x,z\in\Gamma$. If $r(x)_{ij}\ne0$ for some $x$ and $r(y)_{ii}\ne r(y)_{jj}$ for some $y$, then there is $z\in\Gamma$ with $r(z^{\ell^e})_{ij}\ne0$ for every $e\ge0$.
\end{lemma}
\begin{proof}
Write $d_i(x)=r(x)_{ii}$. By induction on $m$, $r(z^m)_{ij}\big(d_i(z)-d_j(z)\big)=r(z)_{ij}\big(d_i(z)^m-d_j(z)^m\big)$. For $m=\ell^e$ the right side is $r(z)_{ij}\big(d_i(z)-d_j(z)\big)^{\ell^e}$. So any $z$ with $r(z)_{ij}\ne0$ and $d_i(z)\ne d_j(z)$ works. If $d_i(x)\ne d_j(x)$, take $z=x$. Otherwise, if $r(y)_{ij}\ne0$, take $z=y$; and if $r(y)_{ij}=0$, take $z=yx$, for which $r(yx)_{ij}=d_i(y)r(x)_{ij}\ne0$ and $d_i(yx)=d_i(y)d_i(x)\ne d_j(y)d_j(x)$, as $d_i(x)=d_j(x)\ne0$.
\end{proof}
% The private entry_pow_mul_sub_eq_and_diag_pow and exists_entry_pow_l_pow_ne_zero of OverFacts.lean.

\begin{lemma}\label{sha-siegel-step-chains-restrict}
Let $D$ be a Siegel datum, $T\le T'$ towers for $D$ with $T$ fit, $k$ a field and $a:\mathcal O\to k$ a homomorphism that kills $\mathfrak m_{\mathcal O}$; flagged pairs are along $a$. Let $W$ be a finite set of primes that split in $T$ and are killed by $T'$ (Definition~\ref{sha-siegel-step-chains-notions}).
\begin{enumerate}
\item Let $s$ be a residual flagged pair for the problem of $T$ with auxiliary primes $W$ such that $\ell\ge5$, or the homomorphism of $s$ is supported in $W_{\mathrm n}$, or $s$ is residual for the problem of $T$. There is a residual flagged pair for the problem of $T'$ that is over $s$.
\item Let $W_1$ and $W'_1$ be finite sets of primes, let $s$ and $t$ be residual flagged pairs for the problem of $T$ with auxiliary primes $W_1$, and let $s'$ and $t'$ be residual flagged pairs for the problem of $T'$ with auxiliary primes $W'_1$, over $s$ and $t$ respectively. Then:
(a) if $s=t$, then $s'=t'$;
(b) if the homomorphism of $s$ is supported in a set of weights, so is that of $s'$;
(c) if $s$ lies in $W_{\mathrm r}$, so does $s'$;
(d) if $s$ and $t$ lie in $W_{\mathrm r}$ and have the same $(0,-1)$-part, then $s'$ and $t'$ have the same $(0,-1)$-part;
(e) if the homomorphisms of $s$ and $t$ are supported in $W_{\mathrm m}$ and $s$ and $t$ have the same $(-1,0)$-part, then $s'$ and $t'$ have the same $(-1,0)$-part;
(f) if the homomorphism of $s$ is supported in $W_{\mathrm r}$ and has an entry of weight $(0,-1)$, so has that of $s'$;
(g) if the homomorphism of $s$ is supported in $W_{\mathrm m}$ and has an entry of weight $(-1,0)$, respectively $(1,-1)$, so has that of $s'$;
(h) if $s$ is over a flagged pair $s_0$ for the rational problem of $(\epsilon,g_0)$ along $a$, where $\epsilon$ is the relabelling of $T$ and $g_0\in G$, then every residual flagged pair for the problem of $T'$ with auxiliary primes $W'_1$ that is over $s_0$ is $s'$.
\end{enumerate}
\end{lemma}
\leanref{Development/AutomorphicArithmetic/AutomorphicVarieties/WeakShimuraTheory/CompactIntegralFields/RestrT.lean}{UpperTheory.V.Ws.restrT}
\leanref{Development/AutomorphicArithmetic/AutomorphicVarieties/WeakShimuraTheory/CompactIntegralFields/OverFacts.lean}{UpperTheory.V.Ws.overFacts}
\begin{proof}
Let $F\subseteq F'$ be the fields of $T$ and $T'$. As $a$ kills $\mathfrak m_{\mathcal O}$, $k$ has characteristic $\ell$, and $a(\chi_i(g))\ne0$ for all $g$.

(1) Write $s=(r,(n_v))$.
\begin{enumerate}
\item[(a)] \emph{Places.} Let $v'\in\Sigma_{T'}$. It is conjugate under $G_F$ to exactly one $v\in\Sigma_T$; choose $g_{v'}\in G_F$ carrying $v$ to $v'$. Then $g_{v'}^{-1}tg_{v'}\in D_v(F)$ for $t\in D_{v'}(F)$.
\item[(b)] \emph{The flags.} The matrix $X=r(g_{v'})$ is symplectic with multiplier $a(\theta(g_{v'}))\ne0$. As $s$ is residual, either $X$ is supported in $W_{\mathrm r}$ and $n_v$ in $W_{\mathrm n}$, or $X$ is supported in $W_{\mathrm m}$. By Lemma~\ref{sha-siegel-step-chains-rational-flags}(1) there is a lower unipotent $n'_{v'}$ with $n_{v'}'^*Xn_v$ upper triangular, which in the first case is supported in $W_{\mathrm n}$.
\item[(c)] \emph{The diagonal characters.} Let $t\in D_{v'}(F')$ and $t'=g_{v'}^{-1}tg_{v'}\in D_v(F)$. Then $r(t)=Xr(t')r(g_{v'}^{-1})$ with $Xr(g_{v'}^{-1})=1$, and $n_v^*r(t')n_v$ is upper triangular with diagonal entries $a(\chi_i(t'))=a(\chi_i(t))$. By Lemma~\ref{sha-siegel-step-chains-rational-flags}(2), $n_{v'}'^*r(t)n'_{v'}$ is upper triangular with diagonal entries $a(\chi_i(t))$.
\item[(d)] \emph{A flagged pair.} Let $r'$ be the restriction of $r$ to $G_{F'}$ and $s'=(r',(n'_{v'}))$. Each $r'(g)$ is symplectic with multiplier $a(\theta(g))$, the $n'_{v'}$ are lower unipotent, and (c) gives the flags. As $T\le T'$, $g_0^{T'}=(g_0^T)^{\ell^{j'}}$, so $r'(g_0^{T'})$ is a power of the diagonal matrix $r(g_0^T)$. For $t\in\mu_{v'}$, $a(\chi_i(t))=1$, as in step (1)(d) of the proof of Lemma~\ref{sha-siegel-step-chains-rational}, since $g\mapsto a(\chi_i(g))$ is a locally constant homomorphism $G_{F'}\to k^\times$; so the diagonal entries of $n_{v'}'^*r'(t)n'_{v'}$ are $1$. Let $g\in I_w(F')$ with $w$ not above $\ell$ and $\bar r_0(g)=1$. If $s$ is residual for the problem of $T$, then $g$ lies in its set $N_\Phi$, and $r(g)=1$. Otherwise, if $w$ lies above no prime of $W$, then $g$ lies in the set $N_\Phi$ of the problem of $T$ with auxiliary primes $W$, and $r(g)=1$; and if $w$ lies above $p\in W$, then $p$ splits in $T$, $T'$ kills $p$, $g\in I_w(F)\cap G_{F'}$, and $\ell\ge5$ or $r$ is supported in $W_{\mathrm n}$, so $r(g)=1$ by Lemma~\ref{sha-siegel-step-chains-restrict-kill}.
\item[(e)] \emph{Residual and over $s$.} The entries of $r'$ are locally constant, $r'(g)_{ii}=a(\chi_i(g))$, (c) gives $\psi'_{v',i}(t)=a(\chi_i(t))$ for $t\in D_{v'}(F')$, and by (b) either $r'$ is supported in $W_{\mathrm r}$ and the $n'_{v'}$ in $W_{\mathrm n}$, or $r'$ is supported in $W_{\mathrm m}$. So $s'$ is residual for the problem of $T'$, and it is over $s$ by (a) and (b).
\end{enumerate}

(2) Write $s=(r,(n_v))$, $t=(r_t,(m_v))$, $s'=(r',(n'_{v'}))$ and $t'=(r'_t,(m'_{v'}))$, so that $r'$ and $r'_t$ are the restrictions of $r$ and $r_t$ to $G_{F'}$. For $v'\in\Sigma_{T'}$ choose $v\in\Sigma_T$ and $g\in G_F$ carrying $v$ to $v'$ with $n_{v'}'^*r(g)n_v$ upper triangular; by Lemma~\ref{sha-siegel-step-chains-restrict-flags}(2), $m_{v'}'^*r_t(g)m_v$ is upper triangular too. For a lower unipotent $n'$ and a matrix $M$ with $U=n'^*M$ upper triangular, $M=n'U$ gives $M_{i1}=n'_{i1}M_{11}$ for all $i$.
\begin{enumerate}
\item[(a)] If $s=t$, then $s'$ and $t'$ are both over $s$, so $s'=t'$ by Lemma~\ref{sha-siegel-step-chains-restrict-flags}(3).
\item[(b)] The homomorphism $r'$ is a restriction of $r$.
\item[(c)] Suppose that $s$ lies in $W_{\mathrm r}$. Then $r'$ is supported in $W_{\mathrm r}$, and step (b) of the proof of Lemma~\ref{sha-siegel-step-chains-rational}(2), with $x=r(g)$, $n=n_v$ and $n'=n'_{v'}$, gives $(n'_{v'})_{32}=0$; so $n'_{v'}$ is supported in $W_{\mathrm n}$.
\item[(d)] Suppose that $s$ and $t$ lie in $W_{\mathrm r}$ and have the same $(0,-1)$-part. The entries $(3,1)$ of $r'$ and $r'_t$ agree. By the identity above with $M=r(g)n_v$, whose first row is $(r(g)_{11},0,0,0)$, $(n'_{v'})_{31}r(g)_{11}=r(g)_{31}+r(g)_{33}(n_v)_{31}$, as $r(g)_{32}=r(g)_{34}=0$; likewise $(m'_{v'})_{31}r_t(g)_{11}=r_t(g)_{31}+r_t(g)_{33}(m_v)_{31}$. The right sides agree, as $s$ and $t$ have the same $(0,-1)$-part and the same diagonal entries, and $r(g)_{11}=r_t(g)_{11}\ne0$. So $(n'_{v'})_{31}=(m'_{v'})_{31}$.
\item[(e)] Suppose that $r$ and $r_t$ are supported in $W_{\mathrm m}$, hence lower triangular, and that $s$ and $t$ have the same $(-1,0)$-part. In the same way $(n'_{v'})_{21}r(g)_{11}=r(g)_{21}+r(g)_{22}(n_v)_{21}$ and $(m'_{v'})_{21}r_t(g)_{11}=r_t(g)_{21}+r_t(g)_{22}(m_v)_{21}$, so $(n'_{v'})_{21}=(m'_{v'})_{21}$; and the entries $(2,1)$ of $r'$ and $r'_t$ agree.
\item[(f)] Suppose that $r$ is supported in $W_{\mathrm r}$ and has an entry of weight $(0,-1)$, at $(3,1)$ or $(4,2)$. By (f) of Definition~\ref{sha-siegel-tower-towers}, some $y\in G_F$ has $a(\chi_1(y))\ne a(\chi_3(y))$. The entries $(1,1)$ and $(2,2)$ of $r(y)^*r(y)=a(\theta(y))$ give $r(y)_{11}r(y)_{44}=a(\theta(y))=r(y)_{22}r(y)_{33}-r(y)_{23}r(y)_{32}$, and $r(y)_{32}=0$; so $r(y)_{22}\ne r(y)_{44}$ too. For $(i,j)=(3,1)$ and $(4,2)$, $r(xz)_{ij}=r(x)_{ij}r(z)_{jj}+r(x)_{ii}r(z)_{ij}$, as $r(x)_{32}$, $r(x)_{34}$, $r(z)_{12}$ and $r(z)_{32}$ vanish. Lemma~\ref{sha-siegel-step-chains-restrict-powers}, with $\Gamma=G_F$, gives $z\in G_F$ with $r(z^{\ell^e})_{ij}\ne0$ for all $e$. As $T\le T'$, $z^{\ell^e}\in G_{F'}$ for some $e$, so $r'$ has an entry of weight $(0,-1)$.
\item[(g)] Suppose that $r$ is supported in $W_{\mathrm m}$, hence lower triangular. The entries of weight $(-1,0)$ are $(2,1)$ and $(4,3)$, and that of weight $(1,-1)$ is $(3,2)$. By (f) of Definition~\ref{sha-siegel-tower-towers}, some $y\in G_F$ has $a(\chi_1(y))\ne a(\chi_2(y))$, and some $y_2\in G_F$ has $a(\chi_2(y_2))\ne a(\chi_3(y_2))$; as above $r(y)_{11}r(y)_{44}=r(y)_{22}r(y)_{33}$, as $r(y)_{23}=0$, so also $r(y)_{33}\ne r(y)_{44}$. For $(i,j)=(2,1)$, $(4,3)$ and $(3,2)$ the relation of Lemma~\ref{sha-siegel-step-chains-restrict-powers} holds, as $r$ is lower triangular. So as in (f), an entry of $r$ of weight $(-1,0)$, respectively $(1,-1)$, gives one of $r'$.
\item[(h)] Let $s_0=(r_0,n_0)$, with $s$ over $s_0$. Then $s'$ is over $s_0$: $r'$ is the restriction of $r_0$; and for $v'$, with $v$ and $g$ as above and $g_v\in G$ carrying $v_\ell$ to $v$ with $n_v^*r_0(g_v)n_0$ upper triangular, $gg_v$ carries $v_\ell$ to $v'$, and $n_{v'}'^*r_0(gg_v)n_0=\big(n_{v'}'^*r(g)n_v\big)\big(n_v^*r_0(g_v)n_0\big)$ is upper triangular. So every flagged pair for the problem of $T'$ with auxiliary primes $W'_1$ over $s_0$ is $s'$, by Lemma~\ref{sha-siegel-step-chains-restrict-flags}(3) for the rational problem.
\end{enumerate}
% (1) is restrT (RestrT.lean, 271 lines): (a) is T.cover and the private conj_dec, (b) is Rt.flagNew, (c) is Rt.flagMove (Lemma sha-siegel-step-chains-rational-flags), (d) is the kill case with the private killW (Lemma sha-siegel-step-chains-restrict-kill) and the mu case with Rt.muTriv (MuTriv.lean, 27 lines), (e) is the final refine. (2) is overFacts (OverFacts.lean, 548 lines): (a) is c1, (b) c2, (c) c3 with the private first_column_and_entry_two_one_of_isUp_adj_mul, (d) c4, (e) c5, (f) c6 with Ok.g3, diag_product_eq_of_isGSp_of_entries_zero and exists_entry_ne_zero_of_triangular_mul, (g) c7, (h) c8 with the private over_comp_of_over_of_isLow. About 820 lines.
\end{proof}

\begin{lemma}\label{sha-siegel-step-chains}
Assume the Galois representations as in Theorem~\ref{sha-residually-reducible-automorphy}.
Let $D$ be a Siegel datum of case II, let $T_0$ be a fit tame tower for $D$ (Definition~\ref{sha-siegel-tower-towers}), with relabelling $\epsilon$, let $\iota:\bar{\mathbf Q}_\ell\cong\mathbf C$, and let $V$ be an ambient ring with residue field $k_V$ (Definition~\ref{sha-siegel-tower-ambient}).
Let $c_0$ be a flagged pair of lattice shape for the rational problem of $(\epsilon,g_0^{T_0})$ along $\mathcal O\to k_V$ (Definition~\ref{sha-siegel-tower-pairs}).
Then there is a tame tower $T_1\ge T_0$ such that for every deep tower $T\ge T_1$ there are a ring $\mathcal O_1\supseteq\mathcal O$ of Definition~\ref{sha-character-parts} and an integral seed $Q$ for $T$ over $\mathcal O_1$ (Definition~\ref{sha-siegel-tower-ideal}) such that every residual flagged pair for the problem of $T$ along $\mathcal O\to k_V$ that is over $c_0$ is joined by a chain to the image of $Q$ under $\mathcal O_1\to V\to k_V$.
\end{lemma}
\leanref{Development/AutomorphicArithmetic/AutomorphicVarieties/ModularityArgument/ThresholdArguments/DenseThresholdCases/Chain2.lean}{UpperTheory.V.Ws.deep_tower_residue_chain_exists}
\begin{proof}
Flagged pairs are along $a:\mathcal O\to k_V$, which kills $\mathfrak m_{\mathcal O}$, and we use the notation of Definition~\ref{sha-siegel-step-chains-notions}. First, let $c=(r,(n_v))$ be a residual flagged pair. If $r$ is supported in $W_{\mathrm r}$, it has an entry of weight $(0,-1)$ if and only if $r(g)_{31}\ne0$ for some $g$; if $r$ is supported in $W_{\mathrm m}$, it has an entry of weight $(-1,0)$ if and only if $r(g)_{21}\ne0$ for some $g$. Indeed, the other entries of these weights are $(4,2)$ and $(4,3)$, and as $r(g^{-1})$ is symplectic with multiplier $a(\theta(g^{-1}))$, $r(g^{-1})^*=a(\theta(g^{-1}))\,r(g)$, whose entries $(4,2)$ and $(4,3)$ are $-r(g^{-1})_{31}$ and $r(g^{-1})_{21}$.
\begin{enumerate}
\item[(1)] \emph{A pair over $c_0$.} Being of lattice shape, $c_0$ is residual and lies in $W_{\mathrm r}$. Lemma~\ref{sha-siegel-step-chains-rational}(1) gives a residual flagged pair $s_0$ for the problem of $T_0$ over $c_0$. By (l), $a(\chi^w(g_0^{T_0}))\ne1$ for every root $w$; so by Lemma~\ref{sha-siegel-step-chains-rational}(2), $s_0$ is residual for the problem of $T_0$ with auxiliary primes $W$ for every $W$, and lies in $W_{\mathrm r}$.
\item[(2)] \emph{Auxiliary primes.} Lemma~\ref{sha-siegel-step-chains-gysin}, applied with $W=\emptyset$ and $p=p'=s_0$, gives a finite set $W$ of primes that split in $T_0$ and a residual flagged pair $s$ for the problem of $T_0$ with auxiliary primes $W$, lying in $W_{\mathrm n}$, with the same $(0,-1)$-part as $s_0$ and an entry of weight $(-1,0)$.
\item[(3)] \emph{The tower $T_1$.} Lemma~\ref{sha-siegel-step-chains-dies} gives a tame tower $T_1\ge T_0$ that kills every prime of $W$.
\end{enumerate}
Now let $T\ge T_1$ be a deep tower, with field $F$. By Lemma~\ref{sha-siegel-step-towers}(2), $T_0\le T$, and $T$ is fit and satisfies (n) to (r). As $G_F\subseteq G_{F_1}$ for the field $F_1$ of $T_1$, $T$ kills every prime of $W$.
\begin{enumerate}
\item[(4)] \emph{The seed.} Lemma~\ref{sha-siegel-step-seeds} gives $\mathcal O_1$, with residue field $k_1$, and an integral seed $Q$ for $T$ over $\mathcal O_1$ whose image in $k_1$ is of type V and has an entry of weight $(-1,0)$ or entries of weights $(-1,1)$ and $(-1,-1)$. Let $c_Q=(r_Q,(n_{Q,v}))$ be the image of $Q$ under $\mathcal O_1\to V\to k_V$. As $k_1\to k_V$ is an injective homomorphism of fields, $c_Q$ is residual and of type V, with the same entries of these weights as the image in $k_1$.
\item[(5)] \emph{Restriction to $T$.} Lemma~\ref{sha-siegel-step-chains-restrict}(1) gives residual flagged pairs $c_3$ over $s$ and $t$ over $s_0$ for the problem of $T$: for $s$ because its homomorphism is supported in $W_{\mathrm n}$, and for $s_0$ because it is residual for the problem of $T_0$. Let $c'_0$ be a residual flagged pair for the problem of $T$ over $c_0$. By Lemma~\ref{sha-siegel-step-chains-restrict}(2)(h), $c'_0=t$, since $s_0$ is over $c_0$. By (b), (c), (d) and (g) of the same lemma, as $s$ lies in $W_{\mathrm n}\subseteq W_{\mathrm r}$ and $W_{\mathrm n}\subseteq W_{\mathrm m}$, the homomorphism $r_3$ of $c_3$ is supported in $W_{\mathrm n}$, $c_3$ lies in $W_{\mathrm r}$, $c_3$ and $c'_0$ have the same $(0,-1)$-part, and $c_3$ has an entry of weight $(-1,0)$.
\item[(6)] \emph{The pair $c'_0$.} The towers $T$ and $T_0$ have the same relabelling, so Lemma~\ref{sha-siegel-step-chains-rational}(2), applied to $T$, $g_0^{T_0}$, $c_0$ and $c'_0$, shows that $c'_0$ lies in $W_{\mathrm r}$, and either has entries of weights $(0,-1)$ and $(-1,1)$, or is of type V and has an entry of weight $(-1,0)$.
\item[(7)] \emph{The Levi part.} Lemma~\ref{sha-siegel-step-chains-levi}(1), applied to $c_3$, gives $c'_4$ with the same $(-1,0)$-part as $c_3$, whose homomorphism and flags are supported in $\{(-1,0)\}$. By (o), $a(\chi^{(0,1)}(s))=1$ for $s\in D_v(F)$; so Lemma~\ref{sha-siegel-step-chains-levi}(2), with $y_v=(n_{Q,v})_{31}$, gives a residual flagged pair $c_4=(r_4,(n_{4,v}))$ with the homomorphism of $c'_4$, with $n_{4,v}$ supported in $W_{\mathrm n}$, with the same entries $(2,1)$ of the flags as $c'_4$, and with $(n_{4,v})_{31}=(n_{Q,v})_{31}$. So $r_4$ is supported in $\{(-1,0)\}$, hence in $W_{\mathrm m}$ and in $W_{\mathrm V}$, and $c_3$ and $c_4$ have the same $(-1,0)$-part. As $r_3$ has an entry of weight $(-1,0)$, some $r_3(g)_{21}=r_4(g)_{21}$ is nonzero, by the first paragraph; so $r_4$ has an entry of weight $(-1,0)$.
\item[(8)] \emph{Ampleness.} By (6), $c'_0$ is ample. If $r_3$ has an entry of weight $(0,-1)$, then $c_3$ is ample, as it also has one of weight $(-1,0)$. Otherwise $c'_0$ has no entry of weight $(0,-1)$, by the first paragraph, as $c_3$ and $c'_0$ have the same $(0,-1)$-part; so by (6) $c'_0$ is of type V. Then $c_3$ is of type V: $r_3$ is supported in $W_{\mathrm n}$ and has no entry of weight $(0,-1)$, so it is supported in $W_{\mathrm V}$, and the flags of $c_3$ are supported in $W_{\mathrm n}$ with the nonzero entries $(3,1)$ of those of $c'_0$. With its entry of weight $(-1,0)$, $c_3$ is ample. The pair $c_4$ is of type V, as $(n_{Q,v})_{31}\ne0$, and has an entry of weight $(-1,0)$, so it is ample. The pair $c_Q$ is of type V; if it has an entry of weight $(-1,0)$ it is ample, and otherwise it has entries of weights $(-1,1)$ and $(-1,-1)$, and by (r) each $v\in\Sigma_T$ has $g\in D_v(F)$ with $a(\chi^{(1,0)}(g))^4\ne1$, so it is ample.
\item[(9)] \emph{The chain.} The pairs $c'_0$ and $c_3$ lie in $W_{\mathrm r}$ and have the same $(0,-1)$-part, so they are linked by Lemma~\ref{sha-siegel-step-chains-lines}(1). The homomorphisms of $c_3$ and $c_4$ are supported in $W_{\mathrm m}$, and the pairs have the same $(-1,0)$-part, so they are linked by Lemma~\ref{sha-siegel-step-chains-lines}(2). The pairs $c_4$ and $c_Q$ lie in $W_{\mathrm r}$, as $W_{\mathrm V}\subseteq W_{\mathrm r}$, and have the same $(0,-1)$-part: the entries $(3,1)$ of $r_4$ and $r_Q$ vanish, as $(0,-1)\notin W_{\mathrm V}$, and $(n_{4,v})_{31}=(n_{Q,v})_{31}$. So they are linked by Lemma~\ref{sha-siegel-step-chains-lines}(1). Hence $c'_0,c_3,c_4,c_Q$ is a chain from $c'_0$ to $c_Q$.
\end{enumerate}
% deep_tower_residue_chain_exists (Chain2.lean, about 130 lines): the first paragraph is Chain.entIff; (1) is Chain.exists_residual_datum_over_restriction_to_rationals with Chain.rtFacts and goodQ; (2) gysinAl; (3) diesEx; the paragraph after (3) is leFit and Dies through fixingSubgroup_antitone; (4) is SeedIIAlt (Lemma sha-siegel-step-seeds) with Chain.isRes_map_red_and_inV_ent_map_red_of_map_residue (SeedRes.lean, 76 lines); (5) is restrT twice and overFacts (g8 for c'_0 = t, f2, f3, f4, f7); (6) is goodQ for T; (7) is Chain.levi and Chain.movePs; (8) is G_0, G_3, G_4 and Gs; (9) is lineP, lineM and lineP. Lemmas sha-siegel-step-chains-lines, -levi, -rational, -gysin (with -gysin-raw), -dies and -restrict (about 3,850 lines) are written out; Lemma sha-siegel-step-seeds is open.
\end{proof}

\begin{lemma}\label{sha-siegel-step-bounds-iwasawa}
Let $\ell$ be an odd prime, $G=\operatorname{Gal}(\bar{\mathbf Q}/\mathbf Q)$, and $U\subseteq G$ an open subgroup that contains the commutator subgroup of $G$ and on which $\chi_{\mathrm{cyc}}\equiv1$ modulo $\ell$. For $n\ge0$ let $U_n$ be the set of $g\in U$ with $\chi_{\mathrm{cyc}}(g)\equiv1$ modulo $\ell^n$. Then there is $r\ge0$ such that for every $n$ there are at most $\ell^r$ continuous homomorphisms $U_n\to\mathbf Z/\ell$ that vanish on $U_n\cap I_A$ for every place $A$ of $\bar{\mathbf Q}$.
\end{lemma}
\leanref{Development/AutomorphicArithmetic/AutomorphicVarieties/WeakShimuraTheory/RealPlaceConjugation/ExistsIwasawaRankBoundOfFerreroWashington.lean}{UpperTheory.V.Ws.exists_iwasawaRank_bound_of_ferreroWashington}
\leanref{Development/AutomorphicArithmetic/AutomorphicVarieties/GaloisLibrary/KummerTheory/PullOnto.lean}{UpperTheory.V.Kum.pullOnto}
\begin{proof}
\begin{enumerate}
\item[(a)] \emph{The fields.} The group $U$ is normal, as it contains the commutator subgroup, and open; so it is $\operatorname{Gal}(\bar{\mathbf Q}/K)$ for a number field $K$, Galois over $\mathbf Q$ with abelian Galois group. As $\chi_{\mathrm{cyc}}\equiv1$ modulo $\ell$ on $U$, $K$ contains a primitive $\ell$-th root of unity. Theorem~\ref{sha-ferrero-washington} gives $r$ such that $|\operatorname{Cl}(K_n)/\operatorname{Cl}(K_n)^\ell|\le\ell^r$ for the $\ell^n$-th cyclotomic extension $K_n$ of $K$, for every $n$; and $\operatorname{Gal}(\bar{\mathbf Q}/K_n)=U_n$.
\item[(b)] \emph{Characters of the class group.} Fix $n$, and let $\mathrm{art}:U_n\to\operatorname{Cl}(K_n)$ be the Artin map of the Hilbert class field of $K_n$, which is surjective. Every continuous homomorphism $\phi:U_n\to\mathbf Z/\ell$ that vanishes on the inertia groups is $\theta\circ\mathrm{art}$ for a homomorphism $\theta:\operatorname{Cl}(K_n)\to\mathbf Z/\ell$. Indeed, if $\tau\in U_n$ acts as complex conjugation for some embedding, then $\tau^2=1$, so $2\phi(\tau)=0$, and $\phi(\tau)=0$ as $\ell$ is odd. So Lemma~\ref{sha-descent-link-hilbert} applies to $\phi$, which vanishes on the kernel of $\mathrm{art}$, and factors through the surjection $\mathrm{art}$.
\item[(c)] \emph{The count.} By (b), there are at most as many such $\phi$ as homomorphisms $\operatorname{Cl}(K_n)\to\mathbf Z/\ell$. Each of these kills the $\ell$-th powers, so it factors through $\operatorname{Cl}(K_n)/\operatorname{Cl}(K_n)^\ell$, a finite $\mathbf F_\ell$-vector space, whose linear forms are as many as its elements, at most $\ell^r$ by (a).
\end{enumerate}
\end{proof}
% exists_iwasawaRank_bound_of_ferreroWashington (ExistsIwasawaRankBoundOfFerreroWashington.lean, 104 lines): (a) is the fixed field K of U, abelian by hc, with the l-th roots of unity by muIff, FW.main at K, and K' = Kum.Fn K Omega l n with Kum.galOfCyc; (b) is Kum.pullOnto (PullOnto.lean), whose inclusion of XDual 1 in the image of ClDual is Kum.artKer (Lemma sha-descent-link-hilbert) with artSurj, the condition at the complex conjugations coming from h2, that 2 is invertible in Z/l; (c) is cardLe_map_of_cardLe and Two.cardClDual (with the private cardLe_hom).

\begin{lemma}\label{sha-siegel-step-bounds-places}
Let $\ell$ be a prime, $G=\operatorname{Gal}(\bar{\mathbf Q}/\mathbf Q)$, $U\subseteq G$ an open normal subgroup and $N\ge1$ an integer. Then there is $g\ge0$ such that every number field $F\subset\bar{\mathbf Q}$ with $U\cap\ker\chi_{\mathrm{cyc}}\subseteq\operatorname{Gal}(\bar{\mathbf Q}/F)$ has at most $g$ places dividing $N$.
\end{lemma}
\leanref{Development/AutomorphicArithmetic/AutomorphicVarieties/WeakShimuraTheory/RealPlaceConjugation/PlacesBd.lean}{UpperTheory.V.Ws.exists_uniform_bound_ncard_placesDividing}
\leanref{Development/AutomorphicArithmetic/AutomorphicVarieties/WeakShimuraTheory/RealPlaceConjugation/DecIdx.lean}{UpperTheory.V.Ws.index_cyclotomicKer_sup_decompositionGroup_ne_zero}
\begin{proof}
Let $V=U\cap\ker\chi_{\mathrm{cyc}}$, a closed normal subgroup. A place of $F$ divides $N$ if and only if it lies above a prime $p$ dividing $N$, so it suffices to bound the places of $F$ above each such $p$. Let $A_0$ be a place of $\bar{\mathbf Q}$ above $p$, with decomposition group $D_0$.
\begin{enumerate}
\item[(a)] \emph{An element of infinite order.} Some $h\in D_0$ has $\chi_{\mathrm{cyc}}(h)$ of infinite order: for $p\ne\ell$ a Frobenius element, with $\chi_{\mathrm{cyc}}(h)=p$; for $p=\ell$ an element of the inertia group, which $\chi_{\mathrm{cyc}}$ maps onto $\mathbf Z_\ell^\times$, as its images modulo every $\ell^n$ are all of $(\mathbf Z/\ell^n)^\times$ by Lemma~\ref{sha-siegel-step-towers-kappa-dirichlet}(2).
\item[(b)] \emph{Finite index.} The image $\chi_{\mathrm{cyc}}(D_0)$ is a closed subgroup of $\mathbf Z_\ell^\times$ that contains an element of infinite order, so it has finite index in $\mathbf Z_\ell^\times$. Hence $\ker\chi_{\mathrm{cyc}}\cdot D_0$ has finite index in $G$, and as $U$ has finite index, so has $V\cdot D_0$; let $n_p$ be this index.
\item[(c)] \emph{Orbits.} The places of $\bar{\mathbf Q}$ above $p$ are the $\sigma A_0$ with $\sigma\in G$, and $\sigma A_0$ and $\sigma'A_0$ are in the same orbit under $V$ when $\sigma'\in V\sigma D_0$; as $V$ is normal, $V\sigma D_0=\sigma VD_0$, so there are at most $n_p$ orbits. If $V\subseteq\operatorname{Gal}(\bar{\mathbf Q}/F)$, places in one orbit lie above the same place of $F$, and every place of $F$ above $p$ lies below a place of $\bar{\mathbf Q}$ above $p$; so $F$ has at most $n_p$ places above $p$.
\end{enumerate}
Then $g=\sum_{p\mid N}n_p$ will do.
\end{proof}
% exists_uniform_bound_ncard_placesDividing (PlacesBd.lean, 78 lines): plDvd_eq for the reduction to primes; (a) and (b) are index_cyclotomicKer_sup_decompositionGroup_ne_zero (DecIdx.lean, with Base.decInf at l, Chr.decCyc away from l, and index_ne_zero_of_isClosed_of_not_isOfFinOrder) and supIndex; (c) is orbIndex and ncard_placesDividing_le_of_orbitBound (NcardPlacesDividingLeOfOrbitBound.lean, 58 lines).

\begin{lemma}\label{sha-siegel-step-bounds-strict}
Let $F\subset\bar{\mathbf Q}$ be a number field, $\ell$ a prime, $M$ a finite $\mathbf F_\ell$-vector space with a linear action of $\operatorname{Gal}(\bar{\mathbf Q}/F)$, and $K\subseteq\operatorname{Gal}(\bar{\mathbf Q}/F)$ a subgroup of finite index prime to $\ell$ that acts trivially on $M$. Let $r\ge0$ be such that there are at most $\ell^r$ continuous homomorphisms $K\to\mathbf Z/\ell$ that vanish on $K\cap I_A$ for every place $A$ of $\bar{\mathbf Q}$. Then there are at most $|M|^{r+1}$ locally constant maps $z:\operatorname{Gal}(\bar{\mathbf Q}/F)\to M$ with $z(gh)=z(g)+gz(h)$ that vanish on $K\cap I_A$ for every place $A$.
\end{lemma}
\leanref{Development/AutomorphicArithmetic/AutomorphicVarieties/WeakShimuraTheory/RealPlaceConjugation/Strict.lean}{UpperTheory.V.Ws.card_cocycles_vanishing_on_inertia_le_pow}
\begin{proof}
Write $G_F=\operatorname{Gal}(\bar{\mathbf Q}/F)$, let $c_1,\dots,c_e$ be the coordinates of $M\cong(\mathbf Z/\ell)^e$, and $Q\subset G_F$ a set of representatives of $G_F/K$. For such a map $z$, the restriction of $z$ to $K$ is a homomorphism, as $K$ acts trivially, so each $c_i\circ z$ restricted to $K$ is one of the at most $\ell^r$ homomorphisms of the statement. The map $z\mapsto\big((c_i\circ z|_K)_i,\sum_{q\in Q}z(q)\big)$ is injective. Indeed, let $w=z-z'$ be the difference of two maps with the same image; it vanishes on $K$, so $w(gk)=w(g)$ for $k\in K$, and $\sum_{q\in Q}w(q)=0$. For $h\in G_F$, the $hq$ run through representatives of $G_F/K$, so $\sum_qw(hq)=\sum_qw(q)=0$, while $\sum_qw(hq)=|G_F/K|\,w(h)+h\sum_qw(q)$. So $|G_F/K|\,w(h)=0$, and $w(h)=0$, as $|G_F/K|$ is prime to $\ell$. Hence there are at most $(\ell^r)^e|M|=|M|^{r+1}$ maps $z$.
\end{proof}
% card_cocycles_vanishing_on_inertia_le_pow (Strict.lean, 129 lines): the injection Phi into (Fin d -> XDual 1 K) x M, with key for the averaging argument and exists_natCard_eq_pow_and_separating_zmod_forms for the coordinates.

\begin{lemma}\label{sha-siegel-step-bounds-global}
Let $\ell$ be an odd prime, $F\subset\bar{\mathbf Q}$ a number field, $G_F=\operatorname{Gal}(\bar{\mathbf Q}/F)$, $\omega:G_F\to\mathbf F_\ell^\times$ the reduction of $\chi_{\mathrm{cyc}}$ modulo $\ell$, and $M$ a finite $\mathbf F_\ell$-vector space with a linear action of $G_F$. Let $N\ge1$ be divisible by $\ell$, such that $F$ has at most $g$ places dividing $N$ and the inertia groups at the places of $\bar{\mathbf Q}$ not above primes dividing $N$ act trivially on $M$. Let $K\subseteq G_F$ be an open normal subgroup of index prime to $\ell$, contained in $\ker\omega$, that acts trivially on $M$, and let $r$ be as in Lemma~\ref{sha-siegel-step-bounds-strict} for $K$. Let $Z$ be the set of locally constant maps $z:G_F\to M$ with $z(gh)=z(g)+gz(h)$ that vanish on the inertia groups at the places not above primes dividing $N$. Then $|Z|\le|M|^{r+2+g+[F:\mathbf Q]}$; and if $F$ is totally real and every complex conjugation in $G_F$ acts trivially on $M$, then $|Z|\le|M|^{r+2+g}$.
\end{lemma}
\leanref{Development/AutomorphicArithmetic/AutomorphicVarieties/WeakShimuraTheory/RealPlaceConjugation/CardGlobalCocyclesLeOfIwasawaRank.lean}{UpperTheory.V.Ws.card_globalCocycles_le_of_iwasawaRank}
\leanref{Development/AutomorphicArithmetic/AutomorphicVarieties/WeakShimuraTheory/RealPlaceConjugation/GwCount.lean}{UpperTheory.V.Ws.greenberg_wiles_count_mod_l}
\leanref{Development/AutomorphicArithmetic/AutomorphicVarieties/WeakShimuraTheory/RealPlaceConjugation/ShaCard.lean}{UpperTheory.V.Ws.tate_shafarevich_card_mod_l}
\begin{proof}
Let $S$ be the set of places of $F$ dividing $N$, which contains the places above $\ell$ and has at most $g$ elements, and $G_S$ the quotient of $G_F$ of Lemma~\ref{sha-descent-link-global}; for $v\in S$ let $D_v$ be the decomposition group of a place of $\bar{\mathbf Q}$ above $v$, and $d_v=[F_v:\mathbf Q_\ell]$ if $v$ lies above $\ell$ and $d_v=0$ otherwise. Let $r_1$ and $r_2$ be the numbers of real and complex places of $F$, and $c_w$ complex conjugations at the real places $w$.
\begin{enumerate}
\item[(a)] \emph{Descent.} Let $H$ be the closed normal subgroup of $G_F$ generated by the inertia groups at the places outside $S$; it acts trivially on $M$, and $G_S=G_F/H$. Each $z\in Z$ vanishes on $H$: it is a homomorphism on $H$, and it vanishes on the conjugates $\sigma I_A\sigma^{-1}=I_{\sigma A}$ of these inertia groups, so on a dense subgroup. So $z(gh)=z(g)$ for $h\in H$, and $z$ is the inflation of a continuous cocycle of $G_S$; this embeds $Z$ in the cocycles $Z^1(G_S,M)$.
\item[(b)] \emph{The Greenberg-Wiles formula.} Lemma~\ref{sha-descent-link-greenberg-wiles}, with $L_v=H^1(D_v,M)$ for all $v\in S$, so that $L_v^\perp=0$, gives
\[
|H^1(G_S,M)|\,|H^0(G_S,M^*(1))|\prod_{v\in S}|H^0(D_v,M)|\prod_w|M^{c_w}|\,|M|^{r_2}=|\text{Ш}^1_S(M^*(1))|\,|H^0(G_S,M)|\prod_{v\in S}|H^1(D_v,M)|,
\]
where $\text{Ш}^1_S(M^*(1))$ is the group of classes of $H^1(G_S,M^*(1))$ that vanish on every $D_v$.
\item[(c)] \emph{The dual Selmer group.} The group $K$ acts trivially on $M^*(1)$, as it lies in $\ker\omega$. Let $y$ be a cocycle of $G_F$ inflated from a class of $\text{Ш}^1_S(M^*(1))$. It vanishes on $K\cap I_A$ for every place $A$: for $A$ outside $S$ by (a); for $A=\sigma A_v$ above $v\in S$, because $y$ is a coboundary $s\mapsto sX-X$ on $D_v$, which vanishes on $K\cap I_{A_v}$, and $y(\sigma s\sigma^{-1})=\sigma y(s)$ for $s\in K\cap I_{A_v}$, as $\sigma s\sigma^{-1}$ acts trivially. Choosing one cocycle in each class, Lemma~\ref{sha-siegel-step-bounds-strict} gives $|\text{Ш}^1_S(M^*(1))|\le|M|^{r+1}$.
\item[(d)] \emph{Local terms.} By Lemma~\ref{sha-descent-link-local}, $\dim H^1(D_v,M)=\dim H^0(D_v,M)+\dim H^2(D_v,M)+d_v\dim M$, and $\dim H^2(D_v,M)\le\dim M$; so $|H^1(D_v,M)|\le|H^0(D_v,M)|\,|M|^{1+d_v}$. Also $\sum_vd_v\le[F:\mathbf Q]$.
\item[(e)] \emph{Cocycles.} The coboundaries form a group isomorphic to $M/H^0(G_S,M)$, so $|Z^1(G_S,M)|=|H^1(G_S,M)|\,|M|/|H^0(G_S,M)|$.
\end{enumerate}
By (b) to (e), $|Z^1(G_S,M)|\prod_w|M^{c_w}|\,|M|^{r_2}\le|M|\cdot|M|^{r+1}\cdot|M|^{|S|+[F:\mathbf Q]}$, and by (a) $|Z|\le|Z^1(G_S,M)|$. The first bound follows, as $|S|\le g$. If $F$ is totally real and the $c_w$ act trivially, then $r_2=0$ and $\prod_w|M^{c_w}|=|M|^{[F:\mathbf Q]}$, which gives the second.
\end{proof}
% card_globalCocycles_le_of_iwasawaRank (CardGlobalCocyclesLeOfIwasawaRank.lean, 66 lines): (a) is exists_descent_to_galoisGroupUnramifiedOutside (DescGS.lean, 135 lines) with exists_zmod_prime_representation_of_torsion; (b) to (e) are greenberg_wiles_count_mod_l (GwCount.lean, 65 lines): P.gw (the Greenberg-Wiles formula of CohomologyHyp) with L = ⊤, tate_shafarevich_card_mod_l (ShaCard.lean, 64 lines) for (c), local_h1_finite_mod_l (LocalH1FiniteModL.lean) for (d), cocycle_card_euler_formula and sum_local_degree_above_l_le_finrank for (e); archH0Card is the factor of the real and complex places.

\begin{lemma}\label{sha-siegel-step-bounds-base}
Let $k_0\subseteq k$ be fields, $Y$ a set, $V_0$ a $k_0$-subspace of the maps $Y\to k_0$ of dimension at most $C$, and $V$ a $k$-subspace of the maps $Y\to k$ such that $\phi\circ z\in V_0$ for every $z\in V$ and every $k_0$-linear map $\phi:k\to k_0$. Then $\dim_kV\le C$.
\end{lemma}
\leanref{Development/AutomorphicArithmetic/AutomorphicVarieties/WeakShimuraTheory/RealPlaceConjugation/BoundsExt.lean}{UpperTheory.V.Ws.cocycleBounds_of_ringHom_comp}
\leanref{Development/AutomorphicArithmetic/AutomorphicVarieties/WeakShimuraTheory/RealPlaceConjugation/BaseCh.lean}{UpperTheory.V.Ws.rank_le_of_linearForm_comp_mem_of_rank_le}
\begin{proof}
For $z\in V$ let $L_z\subseteq k$ be the $k_0$-span of the values of $z$.
\begin{enumerate}
\item[(a)] \emph{The span of the values.} For $z\in V$, $\dim_{k_0}L_z\le C$. Let $e_1,\dots,e_m\in L_z$ be linearly independent over $k_0$. Extending the coordinate forms of their span, there are $k_0$-linear maps $\phi_j:k\to k_0$ with $\phi_j(e_i)=\delta_{ij}$. The maps $\phi_j\circ z$ lie in $V_0$, and they are linearly independent: if $\sum_jc_j\phi_j\circ z=0$ with $c_j\in k_0$, then $\psi=\sum_jc_j\phi_j$ vanishes at the values of $z$, hence on $L_z$, and $c_j=\psi(e_j)=0$. So $m\le C$.
\item[(b)] \emph{Each element of $V$.} Let $z\in V$, let $e_1,\dots,e_m$ be a basis of $L_z$ over $k_0$, which is finite by (a), and let the $\phi_j$ be as in (a). For $y\in Y$, $z(y)\in L_z$, so $z(y)=\sum_j\phi_j(z(y))e_j$. Hence $z=\sum_je_j\,(\phi_j\circ z)$, a $k$-linear combination of elements of $V_0$ viewed as maps $Y\to k$.
\item[(c)] \emph{The bound.} Let $w_1,\dots,w_c$ be a basis of $V_0$, with $c\le C$. Every element of $V_0$ is a $k_0$-linear combination of the $w_i$, so by (b) $V$ lies in the $k$-span of the $w_i$, viewed as maps $Y\to k$. Hence $\dim_kV\le c\le C$.
\end{enumerate}
\end{proof}
% cocycleBounds_of_ringHom_comp (BoundsExt.lean, 64 lines) applies rank_le_of_linearForm_comp_mem_of_rank_le to each of the bounds, with zgl_phi for the cocycle conditions. rank_le_of_linearForm_comp_mem_of_rank_le (BaseCh.lean, 125 lines): (a) is the private rank_span_range_le_of_linearForm_comp_mem, (b) mem_span_image_algebraMap_of_linearForm_comp_mem (LinearMap.exists_extend for the phi_j), (c) the final calc with Module.Free.chooseBasis.

\begin{lemma}\label{sha-siegel-step-bounds}
Let $D$ be a Siegel datum and let $T$ be a fit tame tower for $D$ (Definition~\ref{sha-siegel-tower-towers}).
There is $C\ge0$ such that for every tower $T'\ge T$ that is a layer of $T$, and every field $k$ with a homomorphism $a:\mathcal O\to k$ that kills $\mathfrak m_{\mathcal O}$, the tower $T'$ satisfies the cohomology bounds $(d_{T'},C)$ along $a$ (Definition~\ref{sha-siegel-tower-ideal}).
The proof uses neither that $T$ is fit or tame nor that $T'\ge T$.
\end{lemma}
\leanref{Development/AutomorphicArithmetic/AutomorphicVarieties/WeakShimuraTheory/CompactIntegralFields/X2.lean}{UpperTheory.V.Ws.x2}
\begin{proof}
Let $F$ be the field of $T$, and for a number field $E\supseteq F$ let $U_E\subseteq G$ be the set of $g\in\operatorname{Gal}(\bar{\mathbf Q}/E)$ with $\bar r_0(g)=1$ and $\chi_{\mathrm{cyc}}(g)\equiv1$ modulo $\ell$. Let $N$ be the product of the primes of $S$, which is divisible by $\ell$.
\begin{enumerate}
\item[(a)] \emph{The constants.} The group $U_F$ is open, and it contains the commutator subgroup of $G$, as $F$ is abelian over $\mathbf Q$ and $\bar r_0$ and $\chi_{\mathrm{cyc}}$ modulo $\ell$ take values in commutative groups; so it is normal. Lemma~\ref{sha-siegel-step-bounds-iwasawa}, for $U=U_F$, gives $r$, and Lemma~\ref{sha-siegel-step-bounds-places}, for $U_F$ and $N$, gives $g$. Put $C=r+2+2g$.
\item[(b)] \emph{A layer.} Let $T'\ge T$ be a layer of $T$ of index $m$, with field $F'$, $d'=d_{T'}$, and $U'=U_{F'}$. An element $g\in U_F$ has $\chi_{\mathrm{cyc}}(g)\in1+\ell\mathbf Z_\ell$, on which raising to the power $\ell-1$ is bijective modulo each $\ell^{n}$; so $g\in\operatorname{Gal}(\bar{\mathbf Q}/F')$ if and only if $\chi_{\mathrm{cyc}}(g)\equiv1$ modulo $\ell^{m+1}$. Hence $U'$ is the group $U_{m+1}$ of Lemma~\ref{sha-siegel-step-bounds-iwasawa} for $U_F$, and $U_F\cap\ker\chi_{\mathrm{cyc}}\subseteq\operatorname{Gal}(\bar{\mathbf Q}/F')$. So there are at most $\ell^r$ continuous homomorphisms $U'\to\mathbf Z/\ell$ that vanish on the inertia groups, and $F'$ has at most $g$ places dividing $N$.
\item[(c)] \emph{The group $U'$.} It is open and normal in $G_{F'}=\operatorname{Gal}(\bar{\mathbf Q}/F')$, and $G_{F'}/U'$ embeds in $(\mathbf F^\times)^4\times\mathbf F_\ell^\times$, which has no element of order $\ell$; so its index is prime to $\ell$. Every character $\bar\chi^w$ is trivial on $U'$, as $\chi^w$ is a product of the $\delta_i^{\pm1}$. The characters $\delta_i$ are unramified outside $S$, so an element of the inertia group at a place $A$ not above a prime of $S$ has $\bar r_0=1$, lies in the set $N_{\Phi'}$ of the problem $\Phi'$ of $T'$, and acts trivially through every $\bar\chi^w$.
\item[(d)] \emph{The bounds along $\mathcal O\to\mathbf F$.} Let $w\in\mathbf Z^2$, and let $M_w$ be $\mathbf F$, viewed as an $\mathbf F_\ell$-vector space, with $G_{F'}$ acting through $\bar\chi^w$.
\begin{enumerate}
\item[(i)] Distinct elements of $\Sigma_{T'}$ lie above distinct places of $F'$ above $\ell$, as every place above $\ell$ is conjugate under $G_{F'}$ to exactly one of them; so $|\Sigma_{T'}|\le g\le C$.
\item[(ii)] Let $z\in S_w$, with the $m_v$ of Definition~\ref{sha-siegel-tower-ideal}(3). Then $z$ vanishes on $U'\cap I_A$ for every place $A$. For $A$ not above $\ell$, $U'\cap I_A(F')\subseteq N_{\Phi'}$. For $A$ above $\ell$, write $A=\sigma v$ with $v\in\Sigma_{T'}$ and $\sigma\in G_{F'}$, and let $h\in U'\cap I_A$. Then $\sigma^{-1}h\sigma\in U'\cap D_v(F')$, so $z(\sigma^{-1}h\sigma)=m_v-\bar\chi^w(\sigma^{-1}h\sigma)m_v=0$, and the cocycle relation gives $\bar\chi^w(\sigma^{-1})z(h)=0$. By Lemma~\ref{sha-siegel-step-bounds-strict}, for $K=U'$ and $M_w$, $|S_w|\le|\mathbf F|^{r+1}$, so $\dim S_w\le r+1\le C$.
\item[(iii)] Every $z\in Z_w$ vanishes on the inertia groups at the places not above primes of $S$, which lie in $N_{\Phi'}$ by (c). Lemma~\ref{sha-siegel-step-bounds-global}, for $F'$, $M_w$, $N$, $g$, $K=U'$ and $r$, gives $\dim Z_w\le d'+r+2+g\le d'+C$. For $w\in\{0,(0,1),(0,-1)\}$, $\bar\chi^w$ is $1$ or $\bar\Psi^{\pm1}$, and $\Psi(\tau)=\Psi(c)=1$ for every complex conjugation $\tau$, as complex conjugations are conjugate in $G$; as $F'$ is totally real, the same lemma gives $\dim Z_w\le r+2+g\le C$.
\item[(iv)] Let $v\in\Sigma_{T'}$ and $d_v$ the local degree of $F'$ at the place below $v$. By Lemma~\ref{sha-descent-link-local}, as in step (d) of the proof of Lemma~\ref{sha-siegel-step-bounds-global}, $|H^1(D_v(F'),M_w)|\le|H^0(D_v(F'),M_w)|\,|\mathbf F|^{1+d_v}$; the coboundaries form a group of order $|\mathbf F|/|H^0(D_v(F'),M_w)|$, so $\dim Z_{v,w}\le2+d_v$. As $\sum_vd_v\le d'$, $\sum_v\dim Z_{v,w}\le d'+2g\le d'+C$.
\end{enumerate}
So $T'$ satisfies the cohomology bounds $(d',C)$ along $\mathcal O\to\mathbf F$.
\item[(e)] \emph{Other fields.} Let $a:\mathcal O\to k$ kill $\mathfrak m_{\mathcal O}$; it factors as $\mathcal O\to\mathbf F\xrightarrow{j}k$. For a $\mathbf F$-linear map $\phi:k\to\mathbf F$ and $x\in\mathbf F$, $\phi(j(x)y)=x\phi(y)$; so composing with $\phi$ carries the spaces $Z_w$, $S_w$ and $Z_{v,w}$ along $a$ into the corresponding spaces along $\mathcal O\to\mathbf F$ (for $S_w$ with the elements $\phi(m_v)$). By Lemma~\ref{sha-siegel-step-bounds-base}, applied to each space (to the family $(z_v)_v$ on the disjoint union of the $D_v(F')$ for the last), the bounds of (d) hold along $a$.
\end{enumerate}
% x2 (X2.lean, 37 lines): (a) is kerOverRationals_open_commutator_layer (KerQFacts.lean, 97 lines; U_F is kerQ D F), exists_iwasawaRank_bound_of_ferreroWashington and exists_uniform_bound_ncard_placesDividing with prodS D; (b) is the layer part of kerOverRationals_open_commutator_layer; (c) is residualKernel_properties (KrFacts.lean, 95 lines; U' is T'.Kr); (d)(i) and (iv) are rank_localCocycleSpaces_le (RankLocalCocycleSpacesLe.lean, with card_localCocycles_le_of_place_above_prime); (d)(ii) and (iii) are rank_cocycleSpaces_le_of_towerIwasawaRank (BdGlob.lean, 91 lines, with str_residue_le_cocycles_vanishing_on_inertia_kernel and residualCharacter_complexConjugation_eq_one); (e) is cocycleBounds_of_ringHom_comp. The Lean takes C = r + 2 + 2g. It is stated for any Raw D and any tower T with T.Fit and TameL; it uses T.Fit only to get D.Good (hD), which the book's Siegel datum includes, and it uses neither TameL nor T.Le T'.
\end{proof}

\subsection{Modularity ideals at seeds}\label{subsection-sha-modularity-ideals}

\begin{lemma}\label{sha-siegel-step-ideal-links}
Let $D$ be a Siegel datum of case II or of case III, and let $T$ be a deep tower for $D$ that satisfies (a) to (m) of Definition~\ref{sha-siegel-tower-towers}.
Then there is a ring $\mathcal O_1\supseteq\mathcal O$ of Definition~\ref{sha-character-parts} such that for every ring $\mathcal O'\supseteq\mathcal O_1$ of Definition~\ref{sha-character-parts}, $T$ has two links over $\mathcal O'$ (Definition~\ref{sha-siegel-automorphy-link}) with the same character $\kappa$, whose places $v_0$ lie above distinct primes.
\end{lemma}
\leanref{Development/AutomorphicArithmetic/AutomorphicVarieties/ModularityArgument/ThresholdArguments/DenseThresholdCases/LinkTwo.lean}{UpperTheory.Gal.projectedFormsSystemPart.exists_two_links_same_kappa_not_dvd_of_fit_deep}
\leanref{Development/AutomorphicArithmetic/Foundations/GaloisRepresentations/PlainTorsionPairs/LinkTwo.lean}{UpperTheory.Gal.projectedFormsSystemPart.setAv}
\begin{proof}
In the proof of Lemma~\ref{sha-siegel-automorphy-link-exists}, the characters $\kappa$, $\beta_L$ and $t_L$ of step (1) and the ring $\mathcal O_1$ of step (2) do not depend on the places $p$, $v_0$ and $B_0$ chosen in step (3), which come from Lemma~\ref{sha-siegel-automorphy-link-places}(2). In the proof of that statement, the prime $p$ can be required not to divide a given integer $p_0\ge1$: apply the Chebotarev density theorem with $p_0\cdot2\prod_{q\in S}q\cdot\operatorname{disc}(L)$ in place of $2\prod_{q\in S}q\cdot\operatorname{disc}(L)$. Let $\mathcal O'\supseteq\mathcal O_1$. Steps (3) to (8) with $p_0=1$ give a link over $\mathcal O'$ whose place $v_0$ lies above a prime $p$, and with $p_0=p$ a second link whose place $v_0$ lies above a prime $p'\nmid p$, so $p'\ne p$. Both have the character $\kappa$ of step (1).
\end{proof}
% exists_two_links_same_kappa_not_dvd_of_fit_deep (LinkTwo.lean, about 50 lines): the characters are built once (pre_ex, kapCont, kapUnr, Pr.coefEx, liftO, ochOf), and the link of nonempty_link_of_fit_deep_unramified is built from each of two sets of places Y, Y' given by setAv T with p_0 = 1 and p_0 = Y.s.p; setAv is the proof of Lk.setEx with M = p_0 * 2 * (prod S) * |disc L| in v0Ex. The Lean concludes that K'.s.p does not divide K.s.p.

\begin{lemma}\label{sha-siegel-step-ideal-finite}
Let $D$ be a Siegel datum, let $T$ be a tower for $D$ that satisfies (d) of Definition~\ref{sha-siegel-tower-towers}, with fields $K$, $F$ and $L$, and let $\mathcal O'\supseteq\mathcal O$ be a ring of Definition~\ref{sha-character-parts}, with residue field $k'$. Let $X=(A,B,P_X)$ be a deformation space for $T$ over $\mathcal O'$, with the polynomials $c_g$ and the elements $\psi^B_{v,i}(s)$ of Definition~\ref{sha-siegel-tower-space}. Let a link of $T$ over $\mathcal O'$ be given (Definition~\ref{sha-siegel-automorphy-link}), let $B^{\mathrm{ord}}$ and $\Lambda$ be the universal ordinary determinant ring and the Iwasawa algebra of its residual datum (Situation~\ref{sha-hecke-situation}), let $\varphi:B^{\mathrm{ord}}\to B$ be as in Lemma~\ref{sha-siegel-automorphy-phi}, and let $J_0\subseteq B^{\mathrm{ord}}$ be an ideal such that $B^{\mathrm{ord}}/J_0$ is finite over $\Lambda$.
Then some power of $\mathfrak m_B$ lies in the ideal $\mathfrak a$ generated by $\varphi(J_0)$, the image of $\mathfrak m_{\mathcal O'}$ and the elements $\psi^B_{v,i}(s)-1$ with $v\in\Sigma_T$ and $s\in I_v(F)$.
\end{lemma}
\leanref{Development/AutomorphicArithmetic/AutomorphicVarieties/ModularityArgument/ThresholdArguments/DenseThresholdCases/TPmFin.lean}{UpperTheory.V.Ws.exists_pow_maximalIdeal_le_of_modular_phi}
\leanref{Development/AutomorphicArithmetic/AutomorphicVarieties/AlgebraTools/MonomialRepresentations/CharpolyPowInt.lean}{UpperTheory.V.AlgebraTools.charpolyPowInt}
\begin{proof}
As $B$ is noetherian, it suffices to show that every prime $\mathfrak q\supseteq\mathfrak a$ of $B$ contains $\mathfrak m_B$: then $\mathfrak m_B$ is the radical of $\mathfrak a$, and, being finitely generated, has a power in $\mathfrak a$. Let $\mathfrak q$ be such a prime, $R=B/\mathfrak q$ and $\pi:B\to R$ the quotient map. Then $R$ is a complete local noetherian $\mathcal O'$-algebra with residue field $k'$, a domain, and $\mathcal O'\to R$ kills $\mathfrak m_{\mathcal O'}$, so it factors through $k'$. Call an element of $R$ \emph{constant} if it is the image of an element of $\mathcal O'$. For a complete local $\mathcal O'$-algebra $C$ with residue field $k'$ let $s_C:C\to R$ be the residue map followed by $k'\to R$. A homomorphism of $\mathcal O'$-algebras $f:C\to R$ is compatible with the residue maps, as in the proof of Lemma~\ref{sha-siegel-step-linked}; so if $f(c)$ is the image of $a\in\mathcal O'$, then $c$ has the residue of $a$, and $f(c)=s_C(c)$.
\begin{enumerate}
\item[(a)] \emph{The Iwasawa algebra.} Let $f_\Lambda:\Lambda\to R$ be the map $\Lambda\to B^{\mathrm{ord}}$ followed by $\pi\circ\varphi$. Let $v\mid\ell$ be a place of $F$, the field $L^+$ of the link, let $\bar v\in\Sigma_T$ be the place above it, and let $g\in I_v$. By Lemma~\ref{sha-siegel-automorphy-phi}, $\varphi(\psi^{\mathrm u}_{v,i}(g))=\kappa(g)\psi^B_{\bar v,i}(g)$, and $g\in I_{\bar v}(F)$, so $\pi(\psi^B_{\bar v,i}(g))=1$. So $f_\Lambda$ carries the universal family of characters of the inertia groups to constant values, hence to the family that $s_\Lambda$ carries it to. By the universal property of $\Lambda$, $f_\Lambda=s_\Lambda$, and every element of $f_\Lambda(\Lambda)$ is constant.
\item[(b)] \emph{Integrality.} As $\varphi(J_0)\subseteq\mathfrak q$, $\pi\circ\varphi$ factors through $B^{\mathrm{ord}}/J_0$, which is finite over $\Lambda$. So every element of $\pi(\varphi(B^{\mathrm{ord}}))$ is a root of a monic polynomial whose coefficients lie in $f_\Lambda(\Lambda)$, which are constant by (a); it is integral over $\mathcal O'$.
\item[(c)] \emph{The polynomials $c_g$.} Let $g\in G_F$. As $K$ is quadratic over $\mathbf Q$, $g^2\in G_K\cap G_F=G_L$. By Lemma~\ref{sha-siegel-automorphy-phi}, $\varphi$ carries the coefficient of $Y^n$ in the characteristic polynomial of $D^{\mathrm u}$ at $\mathrm{pr}(g^2)$ to that of $c^\kappa_{g^2}$, which is $\kappa(g^2)^{4-n}$ times the coefficient of $Y^n$ in $c_{g^2}$. By (b), and as $\kappa(g^2)$ is a unit of $\mathcal O'$, $\pi$ maps the coefficients of $c_{g^2}$ to elements integral over $\mathcal O'$. Let $S\subseteq B$ be the inverse image under $\pi$ of the integral closure of $\mathcal O'$ in $R$. The polynomials $c_{g^2}$ and $c_g$ map to the characteristic polynomials of $r_X(g)^2$ and $r_X(g)$. For a square matrix $M$ over a commutative ring, the characteristic polynomial $\chi_M$ divides $\chi_{M^2}(Y^2)$, so the coefficients of $\chi_M$ are integral over every subring containing those of $\chi_{M^2}$. So the coefficients of $c_g$ satisfy, in $A$, monic polynomials with coefficients in the image of $S$, and since $B\to A$ is injective they satisfy them in $B$. Applying $\pi$, their images are integral over the integral closure of $\mathcal O'$ in $R$, hence over $\mathcal O'$.
\item[(d)] \emph{The elements $\psi^B_{v,i}(s)$.} Let $v\in\Sigma_T$ and $s\in D_v(F)$. As $\ell\in S$, (d) gives $D_v(\mathbf Q)\subseteq G_K$, so $s\in G_L$, and $s\in D_w$ for the place $w$ of $F$ below $v$, with $\bar w=v$. By Lemma~\ref{sha-siegel-automorphy-phi}, $\varphi(\psi^{\mathrm u}_{w,i}(s))=\kappa(s)\psi^B_{v,i}(s)$, so by (b) $\pi(\psi^B_{v,i}(s))$ is integral over $\mathcal O'$.
\item[(e)] \emph{Integral elements are constant.} Let $y\in R$ be integral over $\mathcal O'$, and let $a\in\mathcal O'$ have the residue of $y$. As $\mathcal O'$ acts on $R$ through $k'$, $y-a$ is integral over the field $k'$; if it were nonzero, it would be a unit of the domain $R$, but it lies in $\mathfrak m_R$. So $y=a$ is constant.
\item[(f)] \emph{Conclusion.} By (c), (d) and (e), $\pi$ maps the coefficients of the $c_g$ and the $\psi^B_{v,i}(s)$ to constants, so $\pi$ and $s_B$ agree there, by the first paragraph. By (4) of Definition~\ref{sha-siegel-tower-space}, $\pi=s_B$, so $\pi$ kills $\mathfrak m_B$, that is $\mathfrak m_B\subseteq\mathfrak q$.
\end{enumerate}
\end{proof}
% exists_pow_maximalIdeal_le_of_modular_phi (TPmFin.lean, about 170 lines): Ideal.exists_pow_le_of_le_radical_of_fg; (a) is the private pf_lam with pf_sect and Lambda.univ; (b) is pf_int with hM.ring.fin, the finiteness of B^ord/J_0 over Lambda; (c) is pf_lift and pf_sqK with the coefficient formula of Polynomial.scaleRoots, pf_unit, and pf_sq through AlgebraTools.charpolyPowInt; (d) is Ok.split at l; (e) is pf_const; (f) is pf_gen with Space.genB. The Lean statement asks for Fit, Deep and a modular J_0 (ModularAlt); its proof uses only (d) of Ok and the finiteness of B^ord/J_0 over Lambda.

\begin{situation}\label{sha-siegel-step-ideal-setting}
Assume the Galois representations as in Theorem~\ref{sha-residually-reducible-automorphy}.
Let $D$ be a Siegel datum, $T$ a fit deep tower for $D$ with fields $K$, $F$ and $L$ (Definition~\ref{sha-siegel-tower-towers}), $\iota:\bar{\mathbf Q}_\ell\cong\mathbf C$, $\mathcal O'\supseteq\mathcal O$ a ring of Definition~\ref{sha-character-parts}, and $X=(A,B,P_X)$ a deformation space for $T$ over $\mathcal O'$ (Definition~\ref{sha-siegel-tower-space}), with $P_X=(r_X,(n_{X,v}))$. Let a link of $T$ over $\mathcal O'$ be given (Definition~\ref{sha-siegel-automorphy-link}), with character $\kappa$ and place $v_0$; let $B^{\mathrm{ord}}$ and $\Lambda$ be the universal ordinary determinant ring and the Iwasawa algebra of its residual datum (Situation~\ref{sha-hecke-situation}), $J_0\subseteq B^{\mathrm{ord}}$ a Hecke ideal for $\iota$ (Proposition~\ref{sha-hecke-ideal}), $\varphi:B^{\mathrm{ord}}\to B$ as in Lemma~\ref{sha-siegel-automorphy-phi}, and $J_\varphi$ the ideal of $B$ generated by $\varphi(J_0)$. Let $m=(k+1,k,1,0)$ for the integer $k$ of $D$.
\end{situation}

\begin{lemma}\label{sha-siegel-step-ideal-patching-arc}
In Situation~\ref{sha-siegel-step-ideal-setting}, let $\kappa_1$ be a finite field with an $\mathcal O'$-algebra structure and $y:A\to\kappa_1[[T]]$ a homomorphism of $\mathcal O'$-algebras that maps $\mathfrak m_B$ into $T\kappa_1[[T]]$, such that for each $w\in\Sigma_T$ and $i\ne j$ some $s\in U_w$ has $y(\psi_{w,i}(s))\ne y(\psi_{w,j}(s))$, and the $y(r_X(g))$ with $g\in G_F$ span $M_4$ over the fraction field of $\kappa_1[[T]]$. Let $f:B^{\mathrm{ord}}\to\kappa_1[[T]]$ be the composite of $\varphi$, $B\to A$ and $y$. Then $f$ is attached to a nice arc of type I (Definition~\ref{sha-hecke-arcs}), whose representation $\rho$ has $\rho(\mathrm{pr}(g))=\kappa(g)y(r_X(g))$ for $g\in G_L$.
\end{lemma}
\leanref{Development/AutomorphicArithmetic/AutomorphicVarieties/WeakShimuraTheory/ProjectiveLineFlags/PsPairAlt.lean}{UpperTheory.V.Ws.P1.psPair}
\leanref{Development/AutomorphicArithmetic/AutomorphicVarieties/WeakShimuraTheory/ProjectiveLineFlags/AcClosed.lean}{UpperTheory.V.Ws.P1.rhoEx}
\leanref{Development/AutomorphicArithmetic/AutomorphicVarieties/WeakShimuraTheory/ProjectiveLineFlags/FormEx.lean}{UpperTheory.V.Ws.P1.formEx}
\leanref{Development/AutomorphicArithmetic/AutomorphicVarieties/WeakShimuraTheory/ProjectiveLineFlags/HasFlagPush.lean}{UpperTheory.V.Ws.P1.psiEx}
\leanref{Development/AutomorphicArithmetic/AutomorphicVarieties/WeakShimuraTheory/ProjectiveLineFlags/RedPsi.lean}{UpperTheory.V.Ws.P1.redPsi}
\leanref{Development/AutomorphicArithmetic/AutomorphicVarieties/WeakShimuraTheory/ProjectiveLineFlags/RedDet.lean}{UpperTheory.V.Ws.P1.redDet}
\leanref{Development/AutomorphicArithmetic/AutomorphicVarieties/WeakShimuraTheory/ProjectiveLineFlags/IneTriv.lean}{UpperTheory.V.Ws.P1.ineTriv}
\leanref{Development/AutomorphicArithmetic/AutomorphicVarieties/WeakShimuraTheory/ProjectiveLineFlags/MapsPt.lean}{UpperTheory.V.Ws.P1.mapsPt}
\leanref{Development/AutomorphicArithmetic/AutomorphicVarieties/WeakShimuraTheory/ProjectiveLineFlags/CondA.lean}{UpperTheory.V.Ws.P1.condA}
\leanref{Development/AutomorphicArithmetic/AutomorphicVarieties/WeakShimuraTheory/ProjectiveLineFlags/TraceOfCp.lean}{UpperTheory.V.Ws.P1.irrL}
\leanref{Development/AutomorphicArithmetic/AutomorphicVarieties/WeakShimuraTheory/CompactIntegralFields/H4.lean}{UpperTheory.V.Ws.h4}
\leanref{Development/AutomorphicArithmetic/AutomorphicVarieties/WeakShimuraTheory/ProjectiveLineFlags/ContinuousLinearClosed.lean}{UpperTheory.V.Ws.P1.n4V0}
\begin{proof}
Write $r=y\circ r_X$ and $y(P_X)=(r,(n_w))$, with diagonal characters $\psi''_{w,i}$, and let $\bar{\mathcal K}$ be an algebraic closure of the fraction field $\mathcal K$ of $\kappa_1[[T]]$, a field of characteristic $\ell$. For $g\in G$ write $\bar\chi_i(g)$ also for the image in $\kappa_1$ of the residue of $\chi_i(g)$.
\begin{enumerate}
\item[(a)] \emph{The pair.} The composite of $B\to A$, $y$ and the constant coefficient kills $\mathfrak m_B$, so $\mathcal O'\to\kappa_1$ kills $\mathfrak m_{\mathcal O'}$ and $\kappa_1$ is an extension of $k'$. The pair $y(P_X)$ is a flagged pair along $\mathcal O\to\kappa_1[[T]]$, and its entries are $T$-adically continuous, as those of $P_X$ are $\mathfrak m_B$-adically continuous and $y(\mathfrak m_B)\subseteq T\kappa_1[[T]]$. By (2) and (3) of Definition~\ref{sha-siegel-tower-space}, the characteristic polynomial of $r(g)$ reduces modulo $T$ to $\prod_i(X-\bar\chi_i(g))$, and $\psi''_{w,i}(s)$ reduces to $\bar\chi_i(s)$ for $s\in D_w(F)$.
\item[(b)] \emph{The representation.} Let $\rho_L(g)=\kappa(g)r(g)$ for $g\in G_L$. Its entries are continuous and $\kappa_1[[T]]$ is $T$-adically separated, so its kernel is closed. By steps (2) and (3) of the proof of Lemma~\ref{sha-siegel-automorphy-phi}, with $r$ in place of $r_X$, the kernel contains the kernel of $\mathrm{pr}$. So $\rho_L=\rho\circ\mathrm{pr}$ for a homomorphism $\rho:\Gamma_0\to GL_4(\kappa_1[[T]])$ with continuous entries, as $\mathrm{pr}$ is a quotient map.
\item[(c)] \emph{The form.} Let $\tau_F$ be $\tau$ viewed in $G_F$, a complex conjugation; by Lemma~\ref{sha-siegel-step-completion-facts}(2), $\tau_F^2=1$ and $r(\tau_F)$ is symplectic with multiplier $-1$. Let $Q_0$ be the matrix with rows $(0,0,0,1)$, $(0,0,1,0)$, $(0,-1,0,0)$ and $(-1,0,0,0)$, so that ${}^tXQ_0=Q_0X^*$ for every $X$ and ${}^tr(g)Q_0r(g)=\theta(g)Q_0$, and let $P=Q_0r(\tau_F)^{-1}=Q_0r(\tau_F)$. From ${}^tr(\tau_F)Q_0r(\tau_F)=-Q_0$ and ${}^tQ_0=-Q_0$ one gets ${}^tP=P$. For $g\in G_L$, $r(\tau g\tau^{-1})=r(\tau_F)r(g)r(\tau_F)^{-1}$, so ${}^t\rho_L(g)P\rho_L(\tau g\tau^{-1})=\kappa(g)\kappa(\tau g\tau^{-1})\theta(g)P=\chi_{\mathrm{cyc}}(g)^{-3}P$ by (a) of Definition~\ref{sha-siegel-automorphy-link}; that is, ${}^t\rho(\gamma)P\rho(c\gamma c^{-1})=\varepsilon(\gamma)^{-3}P$ for $\gamma\in\Gamma_0$.
\item[(d)] \emph{The characters.} For $v\mid\ell$ and $s\in D_v$, which lies in $D_{\bar v}(F)$, let $\psi_{v,i}(s)=\kappa(s)\psi''_{\bar v,i}(s)$. These are continuous characters, and by (a) and (3) of Definition~\ref{sha-siegel-automorphy-link} they lift the $\bar\psi_{v,i}$. As $y(P_X)$ is a flagged pair, $n_{\bar v}^{-1}\rho(\mathrm{pr}(s))n_{\bar v}$ is upper triangular with diagonal entries $\psi_{v,1}(s),\dots,\psi_{v,4}(s)$.
\item[(e)] \emph{Reduction and the auxiliary places.} By (a) and (3) of Definition~\ref{sha-siegel-automorphy-link}, the characteristic polynomial of $\rho(\mathrm{pr}(g))=\kappa(g)r(g)$ reduces modulo $T$ to $\prod_i(X-\bar\chi_i(\mathrm{pr}(g)))$, so the determinant of $\rho$ reduces to $\bar D$, as in step (6) of the proof of Lemma~\ref{sha-siegel-automorphy-phi}. Let $u\in S_{\mathrm a}$ and $g\in G_L$ in the inertia group at the place above $\tilde u$ with $\bar\chi_i(\mathrm{pr}(g))=1$ for all $i$. By (3) of Definition~\ref{sha-siegel-automorphy-link}, the residues of the $\kappa(g)\chi_i(g)$ are $1$; one of the $\chi_i$ is $\delta_4=1$, so $\kappa(g)\equiv1$, and by (c) there, $\kappa(g)^m=1$ for some $m$ prime to $\ell$, so $\kappa(g)=1$. Hence the residues of the $\delta_j(g)$ are $1$, $\bar r_0(g)=1$, and $g\in N_\Phi$, as $u$ is not above $\ell$; so $r(g)=1$ and $\rho(\mathrm{pr}(g))=1$. With (b) to (d), $(\rho,P,(\psi_{v,i}))$ is an arc.
\item[(f)] \emph{The map $f$.} By Lemma~\ref{sha-siegel-automorphy-phi}, $f$ carries the characteristic polynomial of $D^{\mathrm u}$ at $\mathrm{pr}(g)$ to the image under $y$ of $c^\kappa_g$, which is that of $\kappa(g)r(g)=\rho(\mathrm{pr}(g))$, and $\psi^{\mathrm u}_{v,i}(s)$ to $y(\kappa(s)\psi^B_{\bar v,i}(s))=\psi_{v,i}(s)$. So $f$ carries $D^{\mathrm u}$ to the determinant of $\rho$, as in (e), and the arc is attached to $f$.
\item[(g)] \emph{Condition (N1).} Let $v\mid\ell$ and $i\ne j$, and let $s\in U_{\bar v}$ with $\psi''_{\bar v,i}(s)\ne\psi''_{\bar v,j}(s)$. By (c) of Definition~\ref{sha-siegel-tower-towers}, $s\in I_{\bar v}(F)\subseteq G_L$, so $s\in D_v$; as $\chi_{\mathrm{cyc}}(s)=1$, $\mathrm{pr}(s)\in\Gamma'$. As $\bar r_0(s)=1$, the residues of the $\chi_i(s)$ are $1$, so $\bar\psi_{v,i}(s)=\bar\psi_{v,j}(s)$, while $\psi_{v,i}(s)\ne\psi_{v,j}(s)$. So $\psi_{v,i}\ne\psi_{v,j}$, and as $\varepsilon(s)=1$, $\psi_{v,i}\ne\psi_{v,j}\varepsilon$.
\item[(h)] \emph{A regular element.} Fix $v\mid\ell$. There is $g\in D_v$ with $\mathrm{pr}(g)\in\Gamma'$ such that no $\psi_{v,i}(g)/\psi_{v,j}(g)$ with $i\ne j$ has finite order. Indeed, a unit of $\kappa_1[[T]]$ of finite order whose constant coefficient is $1$ is $1$, as $\kappa_1$ has characteristic $\ell$; so, with $q=|\kappa_1|$, a unit $z$ of finite order has $z^{q-1}=1$. The homomorphisms $\lambda_{ij}=(\psi_{v,i}/\psi_{v,j})^{q-1}$ on the group $H$ of the $g\in D_v$ with $\mathrm{pr}(g)\in\Gamma'$ are nontrivial, by the element $s$ of (g), at which $\psi_{v,i}/\psi_{v,j}$ has constant coefficient $1$ and is not $1$. If no $g$ had all $\lambda_{ij}(g)\ne1$, $H$ would be the union of the finitely many subgroups $\ker\lambda_{ij}$, so one of them would have finite index $d$ \cite{Neu54}; then $\lambda_{ij}(h)^d=1$, so $\lambda_{ij}(h)=1$, for every $h\in H$, a contradiction. For such $g$, no $\psi_{v,i}(g)/\psi_{v,j}(g)$ has finite order.
\item[(i)] \emph{Irreducibility.} The representation $\rho$ is irreducible over $\bar{\mathcal K}$. By hypothesis, the $r(g)$ with $g\in G_F$ span $M_4(\mathcal K)$. The $\bar\chi_i\circ\mathrm{pr}$ are pairwise distinct on $G_L$, so, if $G_L\ne G_F$, there is $g\in G_F\setminus G_L$ with $\sum_i\bar\chi_i(g)\ne0$, by the linear independence of characters applied to the functions $h\mapsto\sum_i\bar\chi_i(gh)$ on $G_L$; by (a), $\operatorname{tr}r(g)\ne0$. As the $r(g)$ with $g\in G_F$ span $M_4(\bar{\mathcal K})$, $r$ is irreducible on $G_F$ over $\bar{\mathcal K}$, which settles the case $G_L=G_F$. Otherwise, suppose that $N$ is a nonzero proper subspace of $\bar{\mathcal K}^4$ stable under $r(G_L)$, and let $M=r(g)N$. As $G_L$ has index $2$ in $G_F$, $r(g')$ maps $N$ into $M$ and $M$ into $N$ for $g'\notin G_L$, and $M$ is stable under $r(G_L)$; so $N\cap M$ and $N+M$ are stable under $r(G_F)$, and $N\cap M=0$ and $N+M=\bar{\mathcal K}^4$. Then $r(g)$ exchanges the summands of $N\oplus M$, so $\operatorname{tr}r(g)=0$, a contradiction. So $r$ is irreducible on $G_L$ over $\bar{\mathcal K}$, and so is $\rho$, as $\rho(\mathrm{pr}(g))=\kappa(g)r(g)$.
\item[(j)] \emph{Conditions (N2) and (N3).} Let $g$ be as in (h) and $u=\mathrm{pr}(g)\in\Gamma'$. By (d), the characteristic polynomial of $\rho(u)$ is $\prod_i(X-\psi_{v,i}(g))$, and $\psi_{v,i}(g)^e\ne\psi_{v,j}(g)^e$ for $i\ne j$ and $e\ge1$. The proof of Lemma~\ref{sha-siegel-automorphy-rank-four}, steps (1) to (3), applies to $\rho$ over $\bar{\mathcal K}$ in place of $\bar{\mathbf Q}_\ell$, with $A=\kappa_1[[T]]$, $j:A\to\bar{\mathcal K}$ the inclusion, $\phi:A\to\kappa_1$ the constant coefficient, the characters $\bar\chi_i$ and the element $u$. Its hypotheses (2) and (3) are (i) and (ii) of step (F) of the proof of Proposition~\ref{sha-siegel-points}, which follow from (g) and (h) of Definition~\ref{sha-siegel-tower-towers}. Its other inputs hold over $\bar{\mathcal K}$: the lemmas it uses are stated for any algebraically closed field, $2$ and $4$ are nonzero in $\bar{\mathcal K}$ as $\ell$ is odd, $\rho$ is irreducible by (i), the characteristic polynomials of $\rho$ are those of (e), and the set of $\gamma\in\Gamma_0$ at which a polynomial in the entries of $\rho(\gamma)$ vanishes is closed, as the entries are $T$-adically continuous and $\kappa_1[[T]]\to\bar{\mathcal K}$ is injective. So $\rho$ is irreducible on $\Gamma'$ over $\bar{\mathcal K}$, which is (N2) of type I, and every open subset of $\Gamma_0$ that meets $\Gamma'$ contains an element of $\Gamma'$ at which the characteristic polynomial of $\rho$ is separable, which is (N3).
\item[(k)] \emph{Condition (N4).} By (3) of Definition~\ref{sha-siegel-automorphy-link}, there is $g\in D_{v_0}$ with $\varepsilon(g)\not\equiv1$ modulo $\ell$, and the $\bar\chi_i(g)$ are all equal to some $\chi\ne0$. Let $X$ be a matrix over $\mathcal K$ with $\varepsilon(g')\rho(g')X=X\rho(g')$ for all $g'\in D_{v_0}$. Write $X=b^{-1}Y$ with $Y$ over $\kappa_1[[T]]$ and $b\ne0$, and let $M=\rho(g)$ and $e=\varepsilon(g)$, so that $eMY=YM$. Suppose $Y\ne0$, and write $Y=T^aY_0$ with $Y_0\not\equiv0$ modulo $T$. Then $eMY_0=Y_0M$, and modulo $T$, with $\bar M=\chi+N$ and $N$ nilpotent by (e), $(\bar e-1)\chi\bar Y_0=\bar Y_0N-\bar eN\bar Y_0$. The map $Z\mapsto ZN-\bar eNZ$ is nilpotent, as left and right multiplication by $N$ commute and are nilpotent, while $(\bar e-1)\chi\ne0$; so $\bar Y_0=0$, a contradiction. Hence $X=0$.
\end{enumerate}
\end{proof}
% The Lean builds the arc as a W.Point K.r in modular_deformation_quotient_regular_at_point_alt. (a) is P1.psPair (PsPairAlt.lean, 64 lines, with P1.coefRes); (b) is P1.rhoEx (AcClosed.lean, through phiKer); (c) is P1.formEx (FormEx.lean, 64 lines: gspForm, ccEx, X1.thFacts, exists_gl4_form_mul_involution and transpose_form_mul_involution, levEps); (d) is P1.psiEx (HasFlagPush.lean) and P1.redPsi (RedPsi.lean); (e) is P1.redDet (RedDet.lean, with detExt) and P1.ineTriv (IneTriv.lean, 70 lines); (f) is P1.mapsPt; (g) is P1.condA (CondA.lean); (h) is exists_mem_powerSeries_char_ratios_not_isOfFinOrder (StabilizerIndexBounds), through Pr.commonFree (Subgroup.exists_finiteIndex_of_leftCoset_cover) and power_series_unit_eq_one_of_torsion; (i) is P1.irrL (TraceOfCp.lean: P1.trNe with linearIndependent_monoidHom, P1.irrTr in TraceSwap.lean, P1.towerGal); (j) is Ws.h4 (H4.lean, 117 lines), with the same lemmas as Lemma sha-siegel-automorphy-rank-four and P1.condC, P1.condD; (k) is P1.n4V0 (ContinuousLinearClosed.lean, with eq_zero_of_smul_mul_eq_mul_of_isNilpotent_constantCoeff). WeakShimuraTheory/ProjectiveLineFlags has about 1,470 lines.

\begin{lemma}\label{sha-siegel-step-ideal-patching}
In Situation~\ref{sha-siegel-step-ideal-setting}, the ideal $J_\varphi$ has the patching property of Definition~\ref{sha-siegel-tower-ideal}(2): for every finite field $\kappa_1$ with an $\mathcal O'$-algebra structure and every homomorphism of $\mathcal O'$-algebras $y:A\to\kappa_1[[\tau]]$ that maps $\mathfrak m_B$ into $\tau\kappa_1[[\tau]]$ and kills $J_\varphi$, such that for each $v\in\Sigma_T$ and $i\ne j$ some $s\in U_v$ has $y(\psi_{v,i}(s))\ne y(\psi_{v,j}(s))$ and the $y(r_X(g))$ span $M_4$ over the fraction field of $\kappa_1[[\tau]]$, the map $y$ detects $J_\varphi$.
\end{lemma}
\leanref{Development/AutomorphicArithmetic/AutomorphicVarieties/ModularityArgument/ThresholdArguments/DenseThresholdCases/ModularDeformationQuotientRegularAtPoint.lean}{UpperTheory.V.Ws.modular_deformation_quotient_regular_at_point_alt}
\begin{proof}
Let $f:B^{\mathrm{ord}}\to\kappa_1[[\tau]]$ be the composite of $\varphi$, $B\to A$ and $y$.
\begin{enumerate}
\item[(a)] \emph{An arc.} The map $f$ is attached to a nice arc of type I (Definition~\ref{sha-hecke-arcs}), by Lemma~\ref{sha-siegel-step-ideal-patching-arc}, with $T=\tau$.
% (a) is P1.psPair, P1.coefRes, P1.rhoEx, P1.formEx, P1.psiEx, P1.redPsi, P1.redDet, P1.ineTriv (the point x : W.Point K.r), P1.condA, P1.irrL, P1.condC, P1.condD with Ws.h4 (H4.lean, 117 lines) and P1.n4V0 for niceness, and P1.mapsPt for f attached to x (Lemma sha-siegel-step-ideal-patching-arc); WeakShimuraTheory/ProjectiveLineFlags has about 1,470 lines.
\item[(b)] \emph{The Hecke ideal.} As $y$ kills $J_\varphi$, $J_0\subseteq\ker f$. By Proposition~\ref{sha-hecke-ideal}(7), $J_0$ has Taylor-Wiles data, so by Proposition~\ref{sha-patching-bound}(1) it lies in every prime of $B^{\mathrm{ord}}$ contained in $\ker f$.
\item[(c)] \emph{Detection.} Let $\mathfrak q$ be a prime of $B$ whose image in $A$ is killed by $y$. Then $\varphi^{-1}(\mathfrak q)$ is a prime of $B^{\mathrm{ord}}$ contained in $\ker f$, so it contains $J_0$ by (b); hence $\mathfrak q$ contains $\varphi(J_0)$ and $J_\varphi$.
\end{enumerate}
\end{proof}
% modular_deformation_quotient_regular_at_point_alt (60 lines): (b) and (c) are hM.p1 (the field P1core of ModularAlt, Proposition sha-patching-bound (1)) applied to q.comap phi, and Ideal.map_le_iff_le_comap.

\begin{lemma}\label{sha-siegel-step-ideal-seeds-phi}
Let $D$, $T$, $\mathcal O'$, $X=(A,B,P_X)$, the link, $B^{\mathrm{ord}}$ and the polynomials $c_g$ be as in Lemma~\ref{sha-siegel-automorphy-phi}, with the characters $\kappa$ and $\beta_L$ of the link, and let $n\ge0$. For $g\in G_L$ let $\kappa_n(g)=\kappa(g)\beta_L(g)^n$ and $c^{\kappa_n}_g(Y)=\kappa_n(g)^4c_g(\kappa_n(g)^{-1}Y)\in B[Y]$.
Then there is a homomorphism of $\mathcal O'$-algebras $\varphi_n:B^{\mathrm{ord}}\to B$ that carries the characteristic polynomial of $D^{\mathrm u}$ at $\mathrm{pr}(g)$ to $c^{\kappa_n}_g$ for $g\in G_L$, and $\psi^{\mathrm u}_{v,i}(s)$ to $\kappa_n(s)\psi^B_{\bar v,i}(s)$ for $v\mid\ell$ and $s\in D_v$.
\end{lemma}
\leanref{Development/AutomorphicArithmetic/AutomorphicVarieties/ModularityArgument/ThresholdArguments/DenseThresholdCases/TPhiEx.lean}{UpperTheory.V.Ws.exists_phi_universal_ordinary_det_ring_of_link}
\leanref{Development/AutomorphicArithmetic/AutomorphicVarieties/WeakShimuraTheory/IsAlgClosureField.lean}{UpperTheory.V.Ws.Link.IsPhi}
\begin{proof}
The proof of Lemma~\ref{sha-siegel-automorphy-phi} applies with $\kappa_n$ in place of $\kappa$. It uses $\kappa$ only through the following properties, which $\kappa_n$ shares. First, $\kappa_n(g)\kappa_n(\tau g\tau^{-1})\theta(g)\chi_{\mathrm{cyc}}(g)^3=1$, as $\beta_L(g)\beta_L(\tau g\tau^{-1})=1$ by (a) of Definition~\ref{sha-siegel-automorphy-link}. Second, $\kappa_n$ is trivial on the inertia groups of (b) there and on the $\mu_v$, and $\kappa_n^m=\kappa^m$ on the inertia groups at the places not above $\ell$: by (f), $\beta_L$ is trivial on the kernel of $t_L$, which contains the inertia groups at the places not above $\ell$ by (f) and the $\mu_v$ by (d). Third, $\kappa_n$ is continuous and has the residue of $\kappa$, as $\beta_L\equiv1$ modulo $\mathfrak m_{\mathcal O'}$ by (f).
\end{proof}
% exists_phi_universal_ordinary_det_ring_of_link (TPhiEx.lean) for any n, with Link.IsPhi X n: phiUnr and phiIne (K.β_t with K.t_unr), phiPol (the relation conj for beta), phiCpRed and phiIsPsi (K.β_red, and K.β_t with K.mu), in WeakShimuraTheory/PhiCharacterResidues.

\begin{lemma}\label{sha-siegel-step-ideal-seeds-weights}
In Situation~\ref{sha-siegel-step-ideal-setting}, let $n\ge0$, let $\varphi_n$ and $\kappa_n$ be as in Lemma~\ref{sha-siegel-step-ideal-seeds-phi}, let $y:A\to\bar{\mathbf Q}_\ell$ be a homomorphism of $\mathcal O'$-algebras, and let $x_n$ be the composite of $\varphi_n$, $B\to A$ and $y$, a point of $B^{\mathrm{ord}}$ with characters $\psi_{x_n,v,i}$ (Definition~\ref{sha-hecke-arcs}(1)).
\begin{enumerate}
\item For every place $w$ of $\bar{\mathbf Q}$ above $\ell$ there are $e_w\in\mathbf Z$ with $e_w+n$ even and $k_w\ge1$ such that $(\kappa_n^2\chi_{\mathrm{cyc}}^{-e_w})^{k_w}=1$ on the inertia group $I_w(L)$ of $G_L$ at $w$.
\item Let $h\in\mathbf Z^4$ with $h_i+n$ even for every $i$, and suppose that $y(P_X)$ has inertial weights $h$ (Definition~\ref{sha-siegel-tower-ideal}). Then there are integers $m_{v,i}$, for the places $v\mid\ell$ of $L^+$ and $1\le i\le4$, with $2(m_{v,i}-m_{v,j})=h_i-h_j$, and $k\ge1$, such that $\big(\psi_{x_n,v,i}\chi_{\mathrm{cyc}}^{-m_{v,i}}\big)^k=1$ on $I_v$ for all $v$ and $i$.
\end{enumerate}
\end{lemma}
\leanref{Development/AutomorphicArithmetic/AutomorphicVarieties/WeakShimuraTheory/ResidueFlagPoints/TwWts.lean}{UpperTheory.V.Ws.Pt.twWts}
\leanref{Development/AutomorphicArithmetic/AutomorphicVarieties/WeakShimuraTheory/ResidueFlagPoints/PowMemOpen.lean}{UpperTheory.V.Ws.Pt.ptWts}
\begin{proof}
(1) By (e) of Definition~\ref{sha-siegel-automorphy-link} there are $a_w$, an odd $b_w$ and $k_w\ge1$ with $(\kappa\chi_{\mathrm{cyc}}^{-a_w})^{k_w}=1$ and $(\beta_L^2\chi_{\mathrm{cyc}}^{-b_w})^{k_w}=1$ on $I_w(L)$. Let $e_w=2a_w+nb_w$. Then $e_w+n=2a_w+n(b_w+1)$ is even, and $\kappa_n^2\chi_{\mathrm{cyc}}^{-e_w}=(\kappa\chi_{\mathrm{cyc}}^{-a_w})^2(\beta_L^2\chi_{\mathrm{cyc}}^{-b_w})^n$, whose $k_w$-th power is $1$ on $I_w(L)$.

(2) Write $\psi''_{w',i}$ for the diagonal characters of $y(P_X)$. As $y(P_X)$ has inertial weights $h$, there is an open subgroup $U$ of $G_F$ with $\psi''_{w',i}(s)^2=\chi_{\mathrm{cyc}}(s)^{h_i}$ for $w'\in\Sigma_T$ and $s\in U\cap I_{w'}(F)$; $U$ contains an open normal subgroup of $G_F$, of some index $N$, so $g^N\in U$ for every $g\in G_F$. For $v\mid\ell$, let $e_v$ and $k_v$ be those of (1) at the place of $\bar{\mathbf Q}$ at which $D_v$ is the decomposition group of $G_L$. As $h_i+n$ and $e_v+n$ are even, $m_{v,i}=(h_i+e_v)/2$ is an integer, and $2(m_{v,i}-m_{v,j})=h_i-h_j$. Let $k=2N\prod_vk_v$. Let $g\in I_v$. Then $g\in D_{\bar v}(F)$, and $g^N\in U\cap I_{\bar v}(F)$. By Lemma~\ref{sha-siegel-step-ideal-seeds-phi}, $\psi_{x_n,v,i}(g)=\kappa_n(g)\psi''_{\bar v,i}(g)$, so $\psi_{x_n,v,i}(g)^{2N}=\kappa_n(g)^{2N}\chi_{\mathrm{cyc}}(g)^{Nh_i}$, and
\[
\big(\psi_{x_n,v,i}(g)\chi_{\mathrm{cyc}}(g)^{-m_{v,i}}\big)^{2N}=\kappa_n(g^N)^2\chi_{\mathrm{cyc}}(g^N)^{-e_v}.
\]
By (1), applied to $g^N\in I_v$, the $k_v$-th power of the right side is $1$; so the $k$-th power of $\psi_{x_n,v,i}(g)\chi_{\mathrm{cyc}}(g)^{-m_{v,i}}$ is $1$.
\end{proof}
% twWts (TwWts.lean, 34 lines): K.cyc for a_w, b_w, k_w and the private aux. ptWts (PowMemOpen.lean, 75 lines): the private pow_mem_open (the normal core of U), m v i = (w i + e v)/2, IsPhi.psi_val, and the final computation hx with k = 2 N prod_v k_v.

\begin{lemma}\label{sha-siegel-step-ideal-seeds-point}
In Situation~\ref{sha-siegel-step-ideal-setting}, let $n\ge0$, $\varphi_n$, $\kappa_n$, $y$ and $x_n$ be as in Lemma~\ref{sha-siegel-step-ideal-seeds-weights}. Let $r=y\circ r_X:G_F\to M_4(\bar{\mathbf Q}_\ell)$ and $\rho_n(g)=\kappa_n(g)r(g)$ for $g\in G_L$, and suppose that the $r(g)$ with $g\in G_L\cap\ker\chi_{\mathrm{cyc}}$ span $M_4(\bar{\mathbf Q}_\ell)$. Then:
\begin{enumerate}
\item the $r(g)$ with $g\in G_F$ span $M_4(\bar{\mathbf Q}_\ell)$, and the restriction of $\rho_n$ to $G_L\cap\ker\chi_{\mathrm{cyc}}$ is irreducible;
\item $r$ and $\rho_n$ are continuous;
\item $x_n$ carries the characteristic polynomial of $D^{\mathrm u}$ at $\mathrm{pr}(g)$ to that of $\rho_n(g)$, for $g\in G_L$;
\item $\rho_n$ is trivial on the inertia group of $G_L$ at the place above $\tilde v_0$ whose image in $\Gamma_0$ is $I_{v_0}$;
\item for $v\mid\ell$, the restriction of $\rho_n$ to $D_v$ is conjugate to an upper triangular representation with diagonal characters $\psi_{x_n,v,1},\dots,\psi_{x_n,v,4}$;
\item there is $P\in GL_4(\bar{\mathbf Q}_\ell)$ with ${}^t\rho_n(g)P\rho_n(\tau g\tau^{-1})=\chi_{\mathrm{cyc}}(g)^{-3}P$ for $g\in G_L$.
\end{enumerate}
\end{lemma}
\leanref{Development/AutomorphicArithmetic/AutomorphicVarieties/WeakShimuraTheory/ResidueFlagPoints/ExistsResLe.lean}{UpperTheory.V.Ws.Pt.ptIrr}
\leanref{Development/AutomorphicArithmetic/AutomorphicVarieties/WeakShimuraTheory/ResidueFlagPoints/ResLeCont.lean}{UpperTheory.V.Ws.Pt.pt_cont}
\leanref{Development/AutomorphicArithmetic/AutomorphicVarieties/WeakShimuraTheory/ResidueFlagPoints/ResLeCont.lean}{UpperTheory.V.Ws.Pt.rho_cont}
\leanref{Development/AutomorphicArithmetic/AutomorphicVarieties/WeakShimuraTheory/AutomorphicCharacters/CharpolySmulScale.lean}{UpperTheory.V.Ws.auRep}
\leanref{Development/AutomorphicArithmetic/AutomorphicVarieties/WeakShimuraTheory/ResidueFlagPoints/ResOneOfMap.lean}{UpperTheory.V.Ws.Pt.ptUnr}
\leanref{Development/AutomorphicArithmetic/AutomorphicVarieties/WeakShimuraTheory/ResidueFlagPoints/PtFlag.lean}{UpperTheory.V.Ws.Pt.ptFlag}
\leanref{Development/AutomorphicArithmetic/AutomorphicVarieties/WeakShimuraTheory/ResidueFlagPoints/PtPair.lean}{UpperTheory.V.Ws.Pt.ptPair}
\begin{proof}
\begin{enumerate}
\item[(1)] As $G_L\cap\ker\chi_{\mathrm{cyc}}\subseteq G_F$, the first claim holds. A subspace of $\bar{\mathbf Q}_\ell^4$ stable under the $\rho_n(g)$ with $g\in G_L\cap\ker\chi_{\mathrm{cyc}}$ is stable under the $r(g)=\kappa_n(g)^{-1}\rho_n(g)$, hence under their span $M_4(\bar{\mathbf Q}_\ell)$, so it is $0$ or $\bar{\mathbf Q}_\ell^4$.
\item[(2)] By Lemma~\ref{sha-automorphy-unramified-near-tangent-points}, $|y(b)|<1$ for $b\in\mathfrak m_B$; as $\mathfrak m_B$ is finitely generated, there is $a<1$ with $|y(b)|\le a^N$ for $b\in\mathfrak m_B^N$. So $y\circ f$ is continuous for every $\mathfrak m_B$-adically continuous $f$ with values in $B$. The trace of $r(g)$ is $-y$ of the coefficient of $Y^3$ in $c_g$, so it is continuous in $g$. By (1) there are $h_1,\dots,h_M\in G_F$ and $z_1,\dots,z_M\in\bar{\mathbf Q}_\ell$ with $\sum_az_ar(h_a)=E_{ji}$, the matrix unit; then $r(g)_{ij}=\operatorname{tr}(r(g)E_{ji})=\sum_az_a\operatorname{tr}r(gh_a)$, which is continuous in $g$. Finally $\kappa_n$ is continuous, as $\kappa$ and $\beta_L$ are, so $\rho_n$ is.
\item[(3)] As in step (a)(i) of the proof of Lemma~\ref{sha-siegel-step-ideal-away}, with Lemma~\ref{sha-siegel-step-ideal-seeds-phi} in place of Lemma~\ref{sha-siegel-automorphy-phi}: $x_n$ carries the characteristic polynomial of $D^{\mathrm u}$ at $\mathrm{pr}(g)$ to $\kappa_n(g)^4\det(\kappa_n(g)^{-1}Y-r(g))=\det(Y-\rho_n(g))$.
\item[(4)] Let $g$ lie in that inertia group. As $v_0$ is not above $\ell$, $\beta_L(g)=1$ by (f) of Definition~\ref{sha-siegel-automorphy-link}. As in step (F) of the proof of Proposition~\ref{sha-siegel-points}, for (3): the $\bar\chi_i$ are trivial on $I_{v_0}$, so the residues of the $\kappa(g)\chi_i(g)$ are $1$; one of the $\chi_i$ is $\delta_4=1$, so $\kappa(g)\equiv1$; by (c), $\kappa(g)^m=1$ for some $m$ prime to $\ell$, so $\kappa(g)=1$. Hence the residues of the $\delta_j(g)$ are $1$, $\bar r_0(g)=1$, and $g$ lies in the set $N_\Phi$ of the problem of $T$, so $r(g)=y(r_X(g))=1$ and $\rho_n(g)=1$.
\item[(5)] Let $v\mid\ell$, $n_{\bar v}=y(n_{X,\bar v})$ and $s\in D_v$, which lies in $D_{\bar v}(F)$. As $y(P_X)$ is a flagged pair, $n_{\bar v}^*=n_{\bar v}^{-1}$ and $n_{\bar v}^{-1}r(s)n_{\bar v}$ is upper triangular with diagonal entries $\psi''_{\bar v,i}(s)$, the diagonal characters of $y(P_X)$. So $n_{\bar v}^{-1}\rho_n(s)n_{\bar v}$ is upper triangular with diagonal entries $\kappa_n(s)\psi''_{\bar v,i}(s)=\psi_{x_n,v,i}(s)$, by Lemma~\ref{sha-siegel-step-ideal-seeds-phi}.
\item[(6)] Let $\tau_F\in G_F$ be the restriction of $\tau$, let $Q_0$ be the matrix with rows $(0,0,0,1)$, $(0,0,1,0)$, $(0,-1,0,0)$ and $(-1,0,0,0)$, so that ${}^tXQ_0=Q_0X^*$ for every $X$, and let $P=Q_0r(\tau_F)^{-1}$. As $r(g)$ is symplectic with multiplier $\theta(g)$, ${}^tr(g)Q_0r(g)=\theta(g)Q_0$. For $g\in G_L$, $r(\tau g\tau^{-1})=r(\tau_F)r(g)r(\tau_F)^{-1}$, so ${}^t\rho_n(g)P\rho_n(\tau g\tau^{-1})=\kappa_n(g)\kappa_n(\tau g\tau^{-1})\theta(g)P$, which is $\chi_{\mathrm{cyc}}(g)^{-3}P$ by (a) of Definition~\ref{sha-siegel-automorphy-link}, as $\beta_L(g)\beta_L(\tau g\tau^{-1})=1$.
\end{enumerate}
\end{proof}
% (1) is ptIrr (ExistsResLe.lean, 45 lines: the private exists_resLe and auIrr); (2) is pt_cont (trCont, with adicPt through GradeZeroInterpolatingElement.coPt and X.cp_cont 3) and rho_cont (twc_cont); (3) is auRep (CharpolySmulScale.lean) with IsPhi n; (4) is ptUnr (ResOneOfMap.lean, 116 lines: the private not_valuation_natCast_prime_lt_one_link, K.β_t with K.t_unr, K.tame, K.v0, the label lab T.pr i = 3 of the trivial character, the private eq_one_of_res and res_one_of_map, then X.pair.kill); (5) is ptFlag (PtFlag.lean, 46 lines, with IsPhi.psi_val); (6) is ptPair (PtPair.lean, 82 lines: formJ, form_gsp, K.conj), Af.IsPairing s.c (K.rhoQ n r) (cyc^(1-4)) P.

\begin{lemma}\label{sha-siegel-step-ideal-seeds-automorphy}
In Situation~\ref{sha-siegel-step-ideal-setting}, let $y:A\to\bar{\mathbf Q}_\ell$ be a homomorphism of $\mathcal O'$-algebras that carries $P_X$ to a seed $(r,(n_w))$, with $h$, $\kappa'$, $a'$, $k'$ and $\pi$ as in Definition~\ref{sha-siegel-tower-ideal}(1) for the character, the integers and the automorphic representation there. Let $n\ge0$ with $h_1+n$ even, and let $\kappa_n$, $x_n$ and $\rho_n$ be as in Lemma~\ref{sha-siegel-step-ideal-seeds-point}.
Then $\rho_n$ is attached through $\iota$ to a regular algebraic cuspidal automorphic representation $\Pi$ of $GL_4(\mathbf A_L)$ with $\Pi^\vee\cong\Pi^c$, and it is windowed (Theorem~\ref{sha-residually-reducible-automorphy}).
\end{lemma}
\leanref{Development/AutomorphicArithmetic/AutomorphicVarieties/ModularityArgument/ThresholdArguments/DenseThresholdCases/BirthSeedB.lean}{UpperTheory.V.Z.Win.Birth.ptAutR'_alt}
\leanref{Development/AutomorphicArithmetic/AutomorphicVarieties/ModularityArgument/LeviSmoothLocalization/PrincipalSeriesMatching/PotOfPair.lean}{UpperTheory.V.Z.ImaginaryQuadraticEtaCharacter.potOfPair}
\begin{proof}
As $h_i-h_{i+1}$ is even, $h_i+n$ is even for every $i$, and Lemmas~\ref{sha-siegel-step-ideal-seeds-weights} and~\ref{sha-siegel-step-ideal-seeds-point} apply. Let $m_{v,i}$ and $k$ be as in Lemma~\ref{sha-siegel-step-ideal-seeds-weights}(2). As $h_i-h_{i+1}\ge4$, $m_{v,j}+2\le m_{v,i}$ for $i<j$.
\begin{enumerate}
\item[(a)] \emph{The character $\lambda=\kappa_n\kappa'^{-1}$.} It is continuous. Let $w$ be a place of $\bar{\mathbf Q}$ above $\ell$, and let $e_w$ and $k_w$ be as in Lemma~\ref{sha-siegel-step-ideal-seeds-weights}(1). As $a'+h_1$ and $h_1+n$ are even, so are $a'+n$ and $e_w-a'$; let $e_w-a'=2c_w$. On $I_w(L)$, $(\lambda\chi_{\mathrm{cyc}}^{-c_w})^2=(\kappa_n^2\chi_{\mathrm{cyc}}^{-e_w})(\kappa'^2\chi_{\mathrm{cyc}}^{-a'})^{-1}$, so $(\lambda\chi_{\mathrm{cyc}}^{-c_w})^{2k'k_w}=1$. So $\lambda$ is a power of $\chi_{\mathrm{cyc}}$ times a character of finite order on the inertia group at every place above $\ell$, and it is the character attached through $\iota$ to an algebraic Hecke character $\chi_\lambda$ of $L$ \cite{Ser68}.
\item[(b)] \emph{Automorphy.} For $g\in G_L$, $\rho_n(g)=\lambda(g)\kappa'(g)r(g)$. The twist of the restriction of $r$ to $G_L$ by $\kappa'$ is attached through $\iota$ to $\pi$, so $\rho_n$ is attached to $\Pi=\pi\otimes(\chi_\lambda\circ\det)$, which is regular algebraic and cuspidal.
\item[(c)] \emph{Polarization.} By Lemma~\ref{sha-siegel-step-ideal-seeds-point}(6), the eigenvalues of $\rho_n$ at a Frobenius element at $\tau w$ are $\chi_{\mathrm{cyc}}^{-3}$ times the inverses of those at $w$, for almost all places $w$; so $\Pi^c$ and $\Pi^\vee$ have the same Satake parameters at almost all places, and $\Pi^\vee\cong\Pi^c$ by strong multiplicity one (Theorem~\ref{lit-jacquet-shalika}(1)).
\item[(d)] \emph{The condition ($\alpha$).} Let $w$ be a place of $L$ above $\ell$, and $v$ the place of $L^+$ below it; as $v$ splits in $L$, $w$ is $\tilde v$ or its conjugate. At $\tilde v$, by Lemma~\ref{sha-siegel-step-ideal-seeds-point}(5), $\rho_n$ is triangular on $D_v$ with the characters $\xi_i=\psi_{x_n,v,i}$. On $I_v$, $\xi_i\xi_j^{-1}\chi_{\mathrm{cyc}}^{-1}$ is $\chi_{\mathrm{cyc}}^{m_{v,i}-m_{v,j}-1}$ times a character whose $k$-th power is $1$, by Lemma~\ref{sha-siegel-step-ideal-seeds-weights}(2), and $m_{v,i}-m_{v,j}-1\ne0$, as $|m_{v,i}-m_{v,j}|\ge2$ for $i\ne j$. As $\chi_{\mathrm{cyc}}$ has infinite order on $I_v$, $\xi_i^e\ne(\xi_j\chi_{\mathrm{cyc}})^e$ for all $i,j$ and $e\ge1$. At the conjugate place, whose decomposition group is $\tau D_v\tau^{-1}$, Lemma~\ref{sha-siegel-step-ideal-seeds-point}(6) gives $\rho_n(\tau g\tau^{-1})=\chi_{\mathrm{cyc}}(g)^{-3}P^{-1}\,{}^t\rho_n(g)^{-1}P$ for $g\in D_v$. So $\rho_n$ is triangular there with the characters $\tau g\tau^{-1}\mapsto\chi_{\mathrm{cyc}}(g)^{-3}\psi_{x_n,v,5-i}(g)^{-1}$, whose integers $-3-m_{v,5-i}$ have the same gaps, and the same argument applies. So ($\alpha$) holds at every place above $\ell$.
\item[(e)] \emph{The condition ($\beta$).} Let $\Pi'$ be attached to $\rho_n$ through $\iota$. Then $\Pi'$ and $\Pi$ have the same Satake parameters at almost all places, so $\Pi'\cong\Pi$ by strong multiplicity one. By Definition~\ref{sha-siegel-tower-ideal}(1) there is $t$ such that at every finite place $u$ of $L$ the supercuspidal support of $\pi_u$ consists of characters $\xi$ with $q_u^{t-1/2}<|\xi(\varpi_u)|<q_u^{t+1/2}$. As $L$ is a CM field, there is $b\in\mathbf R$ with $|\chi_{\lambda,u}(\varpi_u)|=q_u^b$ at every finite place $u$, as in step (F) of the proof of Proposition~\ref{sha-siegel-points}; so $\Pi$ and $\Pi'$ satisfy the bound of ($\beta$) with $s=t+b$ at every finite place.
\end{enumerate}
\end{proof}
% ptAutR'_alt (BirthSeedB.lean, 49 lines): (a) is Af.algOfCyc_s with twWts and the private half_pow (HalfPow.lean); (b) is automorphy_bridge_of_printed_hida and Af.twist_invariance_automorphic; (c) is GalN.V.csdOf0_alt with ptPair; (d) is the hypothesis hpot, supplied in local_zero_of_modular_p by Z.ImaginaryQuadraticEtaCharacter.potOfPair (PotOfPair.lean, 267 lines: potOfFlag with gap at the places s.A v, and FlagAt.dual through the pairing at the others); (e) is Z.Win.Birth.awOne (strong multiplicity one through hecke_eigenvalue_basic_properties and suppSmo) and Z.Win.Birth.awTw (winTw for the shift b, Af.Cq.attTwist), then WitA.

\begin{definition}\label{sha-siegel-step-ideal-seeds-inputs-twisting}
In Situation~\ref{sha-hecke-situation}, for $g\in\Gamma_0$ and $1\le i\le n$ let $a_i(g)\in B$ be the coefficient of $Y^{n-i}$ in the characteristic polynomial of $D^{\mathrm u}$ at $g$; for a function $\gamma$ on $\Gamma_0$ and $s\in D_v$, write $\gamma(s)$ for the value of $\gamma$ at the image of $s$ in $\Gamma_0$.
\begin{enumerate}
\item Let $R$ be a local ring with maximal ideal $\mathfrak m_R$. A \emph{twisting character} over $R$ is a homomorphism $\gamma:\Gamma_0\to R^\times$ such that: (a) $\gamma(cgc^{-1})\gamma(g)=1$ for $g\in\Gamma_0$; (b) $\gamma$ is trivial on $I_{v_0}$, on $I_u$ for $u\in S_{\mathrm a}$, and on the image of $\mu_v$ for $v\mid\ell$; (c) for every $N\ge1$ the map $g\mapsto\gamma(g)\bmod\mathfrak m_R^N$ is locally constant on $\Gamma_0$; (d) $\gamma(g)-1\in\mathfrak m_R$ for every $g\in\Gamma_0$.
\item Let $R$ be a ring, $f_0,f:B\to R$ ring homomorphisms and $\gamma:\Gamma_0\to R^\times$ a homomorphism. Then $f$ is a \emph{twist of $f_0$ by $\gamma$} if $f(a_i(g))=\gamma(g)^if_0(a_i(g))$ for $g\in\Gamma_0$ and $1\le i\le n$, and $f(\psi^{\mathrm u}_{v,i}(s))=f_0(\psi^{\mathrm u}_{v,i}(s))\gamma(s)$ for $v\mid\ell$, $1\le i\le n$ and $s\in D_v$.
\end{enumerate}
\end{definition}
\leanref{Development/AutomorphicArithmetic/AutomorphicVarieties/WeakShimuraTheory/TwistedFourierKernels/IsTwistTwD.lean}{UpperTheory.V.Ws.H5.TwCh}
\leanref{Development/AutomorphicArithmetic/AutomorphicVarieties/WeakShimuraTheory/TwistedFourierKernels/IsTwistTwD.lean}{UpperTheory.V.Ws.H5.TwMap}

\begin{lemma}\label{sha-siegel-step-ideal-seeds-inputs-twist-map}
In Situation~\ref{sha-hecke-situation}, let $R$ be a complete local noetherian $\mathcal O$-algebra with residue field $k$, and $f_0:B\to R$ a homomorphism of $\mathcal O$-algebras.
\begin{enumerate}
\item For every homomorphism $\gamma:\Gamma_0\to R^\times$, at most one homomorphism of $\mathcal O$-algebras $B\to R$ is a twist of $f_0$ by $\gamma$ (Definition~\ref{sha-siegel-step-ideal-seeds-inputs-twisting}).
\item For every twisting character $\gamma$ over $R$, some homomorphism of $\mathcal O$-algebras $B\to R$ is a twist of $f_0$ by $\gamma$.
\item Let $R'$ be a local ring and $\gamma$ a twisting character over $R'$. Then $\gamma^m$ for $m\ge0$, $\gamma^{-1}$, and $e\circ\gamma$ for every local homomorphism $e$ from $R'$ to a local ring, are twisting characters.
\item Let $R'$ be a ring, $f'_0,f':B\to R'$ ring homomorphisms and $\gamma:\Gamma_0\to R'^\times$ a homomorphism such that $f'$ is a twist of $f'_0$ by $\gamma$, and let $e$ be a ring homomorphism from $R'$. Then $e\circ f'$ is a twist of $e\circ f'_0$ by $e\circ\gamma$.
\end{enumerate}
\end{lemma}
\leanref{Development/AutomorphicArithmetic/AutomorphicVarieties/WeakShimuraTheory/TwistedFourierKernels/TwMapUniq.lean}{UpperTheory.V.Ws.H5.twMapUniq}
\leanref{Development/AutomorphicArithmetic/AutomorphicVarieties/WeakShimuraTheory/TwistedFourierKernels/TwMapEx.lean}{UpperTheory.V.Ws.H5.twMapEx}
\leanref{Development/AutomorphicArithmetic/AutomorphicVarieties/WeakShimuraTheory/TwistedFourierKernels/PairTw.lean}{UpperTheory.V.Ws.H5.pairTw}
\leanref{Development/AutomorphicArithmetic/AutomorphicVarieties/WeakShimuraTheory/TwistedFourierKernels/TwChOps.lean}{UpperTheory.V.Ws.H5.twChOps}
\leanref{Development/AutomorphicArithmetic/AutomorphicVarieties/WeakShimuraTheory/TwistedFourierKernels/IsTwistTwD.lean}{UpperTheory.V.Ws.H5.TwMap.comp}
\begin{proof}
Let $(D_0,\psi_0)$ be the image of $(D^{\mathrm u},\psi^{\mathrm u})$ under $f_0$, an ordinary pair over $R$.
\begin{enumerate}
\item[(1)] Let $f$ and $f'$ be twists of $f_0$ by $\gamma$. The determinants $f(D^{\mathrm u})$ and $f'(D^{\mathrm u})$ have the same coefficients $\gamma(g)^if_0(a_i(g))$ of characteristic polynomials at every $g\in\Gamma_0$, so they are equal, as a determinant of a group is determined by these coefficients at the elements of the group, by Amitsur's formula \cite[(1.5) and Corollary~1.14]{Che14}. Also $f(\psi^{\mathrm u}_{v,i}(s))=f'(\psi^{\mathrm u}_{v,i}(s))$ for all $v$, $i$ and $s$. So $f'$ carries $(D^{\mathrm u},\psi^{\mathrm u})$ to the ordinary pair $(f(D^{\mathrm u}),f\circ\psi^{\mathrm u})$, and $f=f'$ by the uniqueness in the universal property of $B$.
\item[(2)] Let $D_\gamma$ be the determinant of $\Gamma_0$ over $R$ obtained from $D_0$ by composing, for every commutative $R$-algebra $R'$, with the $R'$-algebra endomorphism $\sum_gr_gg\mapsto\sum_g\gamma(g)r_gg$ of $R'[\Gamma_0]$, and let $\psi_{\gamma,v,i}=\psi_{0,v,i}\gamma$ on $D_v$. For a determinant of dimension $n$, $r\in R'$ and $y\in R'[\Gamma_0]$, the coefficient of $Y^{n-i}$ in the characteristic polynomial at $ry$ is $r^i$ times that at $y$, by homogeneity; so the coefficient of $Y^{n-i}$ in the characteristic polynomial of $D_\gamma$ at $g$ is $\gamma(g)^if_0(a_i(g))$, and the trace of $D_\gamma$ at $y$ is the trace of $D_0$ at the image of $y$. The pair $(D_\gamma,\psi_\gamma)$ is an ordinary pair over $R$:
\begin{itemize}
\item the coefficients of the characteristic polynomials of $D_\gamma$ are continuous, by (c) of Definition~\ref{sha-siegel-step-ideal-seeds-inputs-twisting}(1) and the continuity of those of $D_0$; by (d), $D_\gamma$ reduces to $\bar D$, as $D_0$ does, and $\psi_{\gamma,v,i}$ lifts $\bar\psi_{v,i}$; $\psi_{\gamma,v,i}$ is continuous by (c), and equal to $\theta_{v,i}$ on $\mu_v$ by (b);
\item for (1) of Situation~\ref{sha-hecke-situation}, at $g^*=\varepsilon(g)^{1-n}(cgc^{-1})^{-1}$ the coefficients for $D_\gamma$ are those for $D_0$ times the powers of $\gamma((cgc^{-1})^{-1})=\gamma(g)$, by (a), and so they equal those at $g$;
\item for (2) there, the endomorphism carries $h\prod_i(g_i-\psi_{\gamma,v,i}(g_i))$ to $\gamma(h)\prod_i\gamma(g_i)$ times $h\prod_i(g_i-\psi_{0,v,i}(g_i))$, whose trace under $D_0$ vanishes;
\item for (3) there, $\gamma$ is trivial on $I_u$ by (b), so the endomorphism carries $h(g-1)$ to $\gamma(h)h(g-1)$.
\end{itemize}
The universal property gives a homomorphism of $\mathcal O$-algebras $f:B\to R$ that carries $(D^{\mathrm u},\psi^{\mathrm u})$ to $(D_\gamma,\psi_\gamma)$, which is a twist of $f_0$ by $\gamma$.
\item[(3)] Conditions (a) and (b) hold for every homomorphism on $\Gamma_0$ whose kernel contains that of $\gamma$. For (c), powers and inverses of a locally constant homomorphism to the units of $R'/\mathfrak m_{R'}^N$ are locally constant, and a local homomorphism $e:R'\to R''$ induces homomorphisms $R'/\mathfrak m_{R'}^N\to R''/\mathfrak m_{R''}^N$. For (d), $\gamma(g)^m-1$ and $\gamma(g)^{-1}-1$ lie in $\mathfrak m_{R'}$, and $e$ maps $\mathfrak m_{R'}$ into $\mathfrak m_{R''}$.
\item[(4)] Apply $e$ to the identities of Definition~\ref{sha-siegel-step-ideal-seeds-inputs-twisting}(2).
\end{enumerate}
\end{proof}
% twMapUniq (TwMapUniq.lean): exists_determinantPair_map_of_algHom for the image pair, CA.d_gext (a determinant is determined by the coefficients at group elements) and B.univ. twMapEx (TwMapEx.lean) with pairTw (PairTw.lean, 96 lines: twD and twA of IsTwistTwD.lean, lamSmul of LamSmul.lean for the homogeneity, then the fields cont, red, psi, P, Ord and I of an ordinary pair) for (2); twChOps (TwChOps.lean) for (3); TwMap.comp for (4). The Lean's TwMap uses the lam i = (-1)^i times the coefficients, for 0 <= i <= n; the signs cancel.

\begin{lemma}\label{sha-siegel-step-ideal-seeds-inputs-twist-quotient}
In Situation~\ref{sha-hecke-situation}, let $R$ be a complete local noetherian $\mathcal O$-algebra with residue field $k$, $\gamma$ a twisting character over $R$ (Definition~\ref{sha-siegel-step-ideal-seeds-inputs-twisting}) and $N\ge1$. Then there are a finite abelian $\ell$-group $Q$, a homomorphism $q:\Gamma_0\to Q$ with open kernel, and a homomorphism $j:Q\to(R/\mathfrak m_R^N)^\times$, such that $q(cgc^{-1})q(g)=1$ for $g\in\Gamma_0$, $q$ is trivial on $I_{v_0}$, on the $I_u$ and on the images of the $\mu_v$, and $j(q(g))$ is the image of $\gamma(g)$ for every $g\in\Gamma_0$. For every such $Q$ and $q$, the group ring $R[Q]$ is a complete local noetherian $\mathcal O$-algebra with residue field $k$, and $g\mapsto[q(g)]$ is a twisting character over $R[Q]$.
\end{lemma}
\leanref{Development/AutomorphicArithmetic/AutomorphicVarieties/WeakShimuraTheory/TwistedFourierKernels/FinQuot.lean}{UpperTheory.V.Ws.H5.finQuot}
\leanref{Development/AutomorphicArithmetic/AutomorphicVarieties/WeakShimuraTheory/TwistedFourierKernels/SingleSubMem.lean}{UpperTheory.V.Ws.H5.grpCO}
\begin{proof}
\begin{enumerate}
\item[(a)] \emph{The quotient.} The group $\Gamma_0$ is open in the compact group $\Gamma$, hence closed and compact. By (c) of Definition~\ref{sha-siegel-step-ideal-seeds-inputs-twisting}(1), $g\mapsto\gamma(g)\bmod\mathfrak m_R^N$ is a locally constant homomorphism to $(R/\mathfrak m_R^N)^\times$; so its image $Q$ is finite and its kernel is open. Let $q$ be this homomorphism onto $Q$ and $j$ the inclusion. The kernel of $q$ contains that of $\gamma$, so (a) and (b) of Definition~\ref{sha-siegel-step-ideal-seeds-inputs-twisting}(1) give the conditions on $q$.
\item[(b)] \emph{An $\ell$-group.} Let $g\in\Gamma_0$ and $x=\gamma(g)$. By (d), $x-1\in\mathfrak m_R$, and $\ell\in\mathfrak m_R$. If $x^{\ell^a}-1\in\mathfrak m_R^{a+1}$, then $x^{\ell^{a+1}}-1=(x^{\ell^a}-1)\sum_{b<\ell}x^{b\ell^a}$, and the sum is congruent to $\ell$ modulo $\mathfrak m_R$, so it lies in $\mathfrak m_R$; hence $x^{\ell^{a+1}}-1\in\mathfrak m_R^{a+2}$. So $x^{\ell^N}-1\in\mathfrak m_R^{N+1}$, and $q(g)^{\ell^N}=1$.
\item[(c)] \emph{The group ring.} Let $Q$ be a finite abelian $\ell$-group and $q$ as in the statement. The ring $R[Q]$ is finite and free over $R$, so it is noetherian. Let $M$ be a maximal ideal of $R[Q]$. As $R[Q]$ is integral over $R$, $M$ lies over $\mathfrak m_R$, so $\ell\in M$; for $z\in Q$ of order $\ell^a$, $([z]-1)^{\ell^a}=[z]^{\ell^a}-1=0$ in the field $R[Q]/M$, of characteristic $\ell$, so $[z]-1\in M$. So every element of $R[Q]$ is congruent modulo $M$ to its image under the augmentation $R[Q]\to R$, and $M$ contains the kernel of the augmentation followed by the residue map $R\to k$, which is a maximal ideal. So $R[Q]$ is local, with residue field $k$, and, being finite over the complete local ring $R$, it is complete \cite[\S8]{Mat86}. The homomorphism $g\mapsto[q(g)]$ satisfies (a) and (b) of Definition~\ref{sha-siegel-step-ideal-seeds-inputs-twisting}(1), as $q$ does; it is locally constant, as $q$ has open kernel, which gives (c); and (d) holds as $[z]-1$ lies in the maximal ideal.
\end{enumerate}
\end{proof}
% finQuot (FinQuot.lean, 39 lines: the image chi.range of gamma mod m^N, IsLocallyConstant.range_finite, and the private pow_pow_sub_one_mem for (b)) and grpCO (SingleSubMem.lean, 98 lines: the private comap_max, single_sub_mem and sub_aug_mem for (c), and co_finite). The twisting character g -> [q(g)] is hΓ in twJN (LcOfKer.lean, with the private lc_of_ker).

\begin{lemma}\label{sha-siegel-step-ideal-seeds-inputs-twist-stable}
In Situation~\ref{sha-hecke-situation}, let $J\subseteq B$ be an ideal with the density property (2) and the twist property (6) of Proposition~\ref{sha-hecke-ideal}, such as a Hecke ideal. Let $\gamma$ be a twisting character over $B$ (Definition~\ref{sha-siegel-step-ideal-seeds-inputs-twisting}) and $\tau:B\to B$ a homomorphism of $\mathcal O$-algebras that is a twist of the identity by $\gamma$ (Lemma~\ref{sha-siegel-step-ideal-seeds-inputs-twist-map}). Then $\tau(J)\subseteq J$.
\end{lemma}
\leanref{Development/AutomorphicArithmetic/AutomorphicVarieties/WeakShimuraTheory/TwistedFourierKernels/LcOfKer.lean}{UpperTheory.V.Ws.H5.twJN}
\leanref{Development/AutomorphicArithmetic/AutomorphicVarieties/WeakShimuraTheory/TwistedFourierKernels/LevelJ.lean}{UpperTheory.V.Ws.H5.levelJ}
\leanref{Development/AutomorphicArithmetic/AutomorphicVarieties/WeakShimuraTheory/TwistedFourierKernels/ClosedJ.lean}{UpperTheory.V.Ws.H5.closedJ}
\leanref{Development/AutomorphicArithmetic/AutomorphicVarieties/WeakShimuraTheory/TwistedFourierKernels/Fourier.lean}{UpperTheory.V.Ws.H5.fourier}
\begin{proof}
Let $b\in J$ and $N\ge1$. The ring $B$ is a complete local noetherian $\mathcal O$-algebra with residue field $k$.
\begin{enumerate}
\item[(a)] \emph{A twist over a group ring.} Let $Q$, $q$ and $j$ be as in Lemma~\ref{sha-siegel-step-ideal-seeds-inputs-twist-quotient} for $R=B$, $\gamma$ and $N$, and let $\gamma_Q(g)=[q(g)]$, a twisting character over $B[Q]$ by that lemma. By Lemma~\ref{sha-siegel-step-ideal-seeds-inputs-twist-map}(2), some homomorphism of $\mathcal O$-algebras $t:B\to B[Q]$ is a twist of the inclusion $\iota_Q:B\to B[Q]$ by $\gamma_Q$. Write $t(b)=\sum_{z\in Q}b_z[z]$ with $b_z\in B$.
\item[(b)] \emph{The coefficients lie in $J$.} Let $h_0$ be the increasing weight with $h_{0,v,j,\tau}=j$. By the density property (2), it suffices to show that $x(b_z)=0$ for every point $x$ of distinct type of weight $h_0$ with $J\subseteq\ker x$ and every $z\in Q$. Let $\chi:Q\to\bar{\mathbf Q}_\ell^\times$ be a character, and $e_\chi:B[Q]\to\bar{\mathbf Q}_\ell$ the ring homomorphism $\sum_za_z[z]\mapsto\sum_zx(a_z)\chi(z)$. Then $x'=e_\chi\circ t$ is a point of $B$, and by Lemma~\ref{sha-siegel-step-ideal-seeds-inputs-twist-map}(4) it is a twist of $e_\chi\circ\iota_Q=x$ by $\beta=e_\chi\circ\gamma_Q=\chi\circ q$. So the coefficients of the characteristic polynomials of $x'(D^{\mathrm u})$ are those of $x(D^{\mathrm u})\otimes\beta$, and as a determinant is determined by them \cite[(1.5) and Corollary~1.14]{Che14}, $x'(D^{\mathrm u})=x(D^{\mathrm u})\otimes\beta$; and $\psi_{x',v,j}=\psi_{x,v,j}\beta$. The character $\beta$ has open kernel, its order divides that of $Q$, a power of $\ell$, it is trivial on $I_{v_0}$ and on the $I_u$, and $\beta(cgc^{-1})=\beta(g)^{-1}$, by the properties of $q$. So the twist property (6) gives $J\subseteq\ker x'$, and $0=x'(b)=\sum_zx(b_z)\chi(z)$. As this holds for every $\chi$, and the characters of the finite abelian group $Q$ with values in the algebraically closed field $\bar{\mathbf Q}_\ell$ of characteristic $0$ separate its elements, the functions $\chi\mapsto\chi(z)$, $z\in Q$, are distinct characters of the group of characters of $Q$; by the linear independence of characters, $x(b_z)=0$ for every $z$. Hence $b_z\in J$.
\item[(c)] \emph{Modulo $\mathfrak m_B^N$.} Let $\pi_N:B\to B/\mathfrak m_B^N$ be the quotient map, and $e:B[Q]\to B/\mathfrak m_B^N$ the homomorphism of $B$-algebras with $e([z])=j(z)$. By Lemma~\ref{sha-siegel-step-ideal-seeds-inputs-twist-map}(4), $e\circ t$ is a twist of $e\circ\iota_Q=\pi_N$ by $e\circ\gamma_Q=j\circ q$, which is $\pi_N\circ\gamma$; and $\pi_N\circ\tau$ is a twist of $\pi_N$ by $\pi_N\circ\gamma$. The ring $B/\mathfrak m_B^N$ is a complete local noetherian $\mathcal O$-algebra with residue field $k$, so $e\circ t=\pi_N\circ\tau$ by Lemma~\ref{sha-siegel-step-ideal-seeds-inputs-twist-map}(1). Hence $\pi_N(\tau(b))=\sum_z\pi_N(b_z)j(z)$, which lies in the image of $J$ by (b); that is, $\tau(b)\in J+\mathfrak m_B^N$.
\item[(d)] \emph{Conclusion.} By (c), the image of $\tau(b)$ in $B/J$ lies in $\bigcap_N\mathfrak m_B^N(B/J)$, which is $0$ by the Krull intersection theorem \cite[Theorem~8.10]{Mat86}. So $\tau(b)\in J$.
\end{enumerate}
\end{proof}
% twJN (LcOfKer.lean, 63 lines) for (a) and (c): finQuot, grpCO, hΓ, twMapEx for tw, levelJ, then e, h1 = htw.comp e, h2 = hτ.comp, twMapUniq and the private lift_mem; with N + 1 in place of N. (b) is levelJ (LevelJ.lean, 78 lines): hD with h0 = fun _ j _ => j, fourier (Fourier.lean: linearIndependent_monoidHom), the private lift_map for e_chi, hT (the field TwistOk of the Hecke ideal, Proposition sha-hecke-ideal (6)) with isTwistOf (IsTwistOf.lean, through CA.d_gext). (d) is closedJ (ClosedJ.lean: Ideal.iInf_pow_smul_eq_bot_of_isLocalRing on B/J). The Lean asks l in the maximal ideal of B, which holds as k has characteristic l.

\begin{lemma}\label{sha-siegel-step-ideal-seeds-inputs-twist-descend}
Let $D$ be a Siegel datum, $T$ a tower for $D$ with fields $K$, $F$ and $L$, and $\mathcal O'\supseteq\mathcal O$ a ring of Definition~\ref{sha-character-parts}. Let a link of $T$ over $\mathcal O'$ be given (Definition~\ref{sha-siegel-automorphy-link}), with characters $\beta_L$ and $t_L$ and projection $\mathrm{pr}:G_L\to\Gamma_0$. Then there is a twisting character $\beta_0$ over $\mathcal O'$ (Definition~\ref{sha-siegel-step-ideal-seeds-inputs-twisting}, for the data of Situation~\ref{sha-hecke-situation} of the link) with $\beta_0(\mathrm{pr}(g))=\beta_L(g)$ for $g\in G_L$.
\end{lemma}
\leanref{Development/AutomorphicArithmetic/AutomorphicVarieties/WeakShimuraTheory/TwistedFourierKernels/Descend.lean}{UpperTheory.V.Ws.H5.descend}
\leanref{Development/AutomorphicArithmetic/AutomorphicVarieties/WeakShimuraTheory/TwistedFourierKernels/PrFacts.lean}{UpperTheory.V.Ws.H5.prFacts}
\begin{proof}
\begin{enumerate}
\item[(a)] \emph{Descent.} The kernel of $t_L$ is a closed subgroup of $G_L$, as $t_L$ is continuous, and it contains the inertia groups of $G_L$ at the places not above $\ell$, by (f) of Definition~\ref{sha-siegel-automorphy-link}; these include the places above the places of $L^+$ outside $S$, as $S$ contains the places above $\ell$. By step (2) of the proof of Lemma~\ref{sha-siegel-automorphy-phi}, the kernel of $\mathrm{pr}$ lies in the kernel of $t_L$, which lies in the kernel of $\beta_L$ by (f). As $\mathrm{pr}$ is surjective, $\beta_L=\beta_0\circ\mathrm{pr}$ for a homomorphism $\beta_0:\Gamma_0\to\mathcal O'^\times$.
\item[(b)] \emph{Condition (a).} As $c$ is the image of $\tau$, $c\,\mathrm{pr}(g)c^{-1}=\mathrm{pr}(\tau g\tau^{-1})$, so $\beta_0(c\,\mathrm{pr}(g)c^{-1})\beta_0(\mathrm{pr}(g))=\beta_L(\tau g\tau^{-1})\beta_L(g)=1$ by (a) of Definition~\ref{sha-siegel-automorphy-link}.
\item[(c)] \emph{Condition (b).} The groups $I_{v_0}$ and $I_u$ are the images under $\mathrm{pr}$ of inertia groups of $G_L$ at places above $\tilde v_0$ and $\tilde u$. These places are not above $\ell$: $u\in S_{\mathrm a}$ is prime to $\ell$, and the residue field of $v_0$ is $\mathbf F_p$ with $p\ne\ell$. So $t_L$, and hence $\beta_L$, is trivial on these inertia groups, by (f). By (d) of Definition~\ref{sha-siegel-automorphy-link}, $t_L$ is trivial on $\mu_v$, and so is $\beta_L$.
\item[(d)] \emph{Condition (c).} For $N\ge1$, the composite of $g\mapsto\beta_0(g)\bmod\mathfrak m_{\mathcal O'}^N$ with $\mathrm{pr}$ is $\beta_L\bmod\mathfrak m_{\mathcal O'}^N$, which is locally constant, as $\beta_L$ is continuous; as $\mathrm{pr}$ is a quotient map, $\beta_0\bmod\mathfrak m_{\mathcal O'}^N$ is locally constant.
\item[(e)] \emph{Condition (d).} $\beta_0(\mathrm{pr}(g))=\beta_L(g)$ is congruent to $1$ modulo $\mathfrak m_{\mathcal O'}$, by (f).
\end{enumerate}
\end{proof}
% descend (Descend.lean, 81 lines): prFacts (PrFacts.lean, through phiKer, which is step (2) of the proof of Lemma sha-siegel-automorphy-phi; prFacts also asks the subgroup to be normal and stable under conjugation by c, which ker t_L is, by K.conj) with K.t_cont, K.t_unr and K.β_t for (a); liftOfRightInverse for beta_0; K.conj for (b); K.t_unr with K.s.St_l, and with the private not_mem_of_norm and K.s.norm_v0_eq_one for v_0, and K.mu for (c); K.β_cont with hq.isOpen_preimage for (d); K.β_red for (e).

\begin{lemma}\label{sha-siegel-step-ideal-seeds-inputs-regular}
Let $\mathcal O$ be the ring of integers of a finite extension of $\mathbf Q_\ell$, $N\ge0$, $\Lambda=\mathcal O[[X_1,\dots,X_N]]$, $R$ a complete local noetherian $\mathcal O$-algebra with the residue field of $\mathcal O$, $f:\Lambda\to R$ a homomorphism of $\mathcal O$-algebras, and $J\subseteq R$ an ideal such that $R/J$ is finite over $\Lambda$ and such that $f(a)b\in J$ with $a\ne0$ implies $b\in J$. Let $x:R\to\bar{\mathbf Q}_\ell$ be a homomorphism of $\mathcal O$-algebras with $J\subseteq\ker x$, and suppose that every homomorphism of $\mathcal O$-algebras $x':R\to\bar{\mathbf Q}_\ell[\epsilon]/(\epsilon^2)$ that reduces to $x$ modulo $\epsilon$ and maps $f(\Lambda)$ into $\bar{\mathbf Q}_\ell$ is equal to $x$. Then every prime of $R$ contained in $\ker x$ contains $J$.
\end{lemma}
\leanref{Development/AutomorphicArithmetic/AutomorphicVarieties/WeakShimuraTheory/EigenbasisLifts/RegOfTan.lean}{UpperTheory.V.Ws.Eb.regOfTan}
\leanref{Development/AutomorphicArithmetic/AutomorphicVarieties/WeakShimuraTheory/EigenbasisLifts/TanEmb.lean}{UpperTheory.V.Ws.Eb.tanEmb}
\leanref{Development/AutomorphicArithmetic/AutomorphicVarieties/WeakShimuraTheory/EigenbasisLifts/DimLow.lean}{UpperTheory.V.Ws.Eb.dimLow}
\leanref{Development/AutomorphicArithmetic/AutomorphicVarieties/WeakShimuraTheory/EigenbasisLifts/DimLeOfInt.lean}{UpperTheory.V.Ws.Eb.htLow}
\leanref{Development/AutomorphicArithmetic/AutomorphicVarieties/AlgebraTools/MonomialRepresentations/RegOfDim.lean}{UpperTheory.V.AlgebraTools.regOfDim}
\begin{proof}
Let $\mathfrak p=\ker x$, a prime of $R$, and $R_{\mathfrak p}$ the localization. A \emph{derivation at} $x$ is an $\mathcal O$-linear map $D:R\to\bar{\mathbf Q}_\ell$ with $D(ab)=x(a)D(b)+x(b)D(a)$; they form a $\bar{\mathbf Q}_\ell$-vector space, and $D$ is one if and only if $a\mapsto x(a)+\epsilon D(a)$ is a homomorphism of $\mathcal O$-algebras.
\begin{enumerate}
\item[(a)] \emph{Derivations.} The derivations at $x$ form a space of dimension at most $N$. Indeed, a derivation $D$ at $x$ that vanishes on $f(\Lambda)$ gives a homomorphism $x+\epsilon D$ that maps $f(\Lambda)$ into $\bar{\mathbf Q}_\ell$, so $D=0$ by hypothesis; hence $D\mapsto D\circ f$ is injective into the derivations of $\Lambda$ at $x\circ f$. These form a space of dimension at most $N$, as $D'\mapsto(D'(X_1),\dots,D'(X_N))$ is injective on them. Indeed, let $D'$ be a derivation of $\Lambda$ at $x\circ f$ with $D'(X_i)=0$ for every $i$. The Leibniz rule at $1\cdot1$ gives $D'(1)=0$, so $D'$ vanishes on $\mathcal O$, and by the Leibniz rule and linearity it vanishes on the $\mathcal O$-subalgebra $\mathcal O[X_1,\dots,X_N]$. Let $a\in\Lambda$ and $k\ge0$, and let $a_k$ be the sum of the terms of $a$ of total degree less than $k$, a polynomial; as the $X_i$ lie in $\mathfrak m_\Lambda$, $a-a_k\in\mathfrak m_\Lambda^k$. The ring $\Lambda$ is a complete local noetherian $\mathcal O$-algebra with the residue field of $\mathcal O$, so by Lemma~\ref{sha-automorphy-unramified-near-tangent-small}, applied to $x\circ f$ and $D'$, $|D'(a)|=|D'(a-a_k)|$ is arbitrarily small for $k$ large. So $D'(a)=0$.
\item[(b)] \emph{The maximal ideal of $R_{\mathfrak p}$ is generated by at most $N$ elements.} The residue field $\kappa(\mathfrak p)$ of $R_{\mathfrak p}$ embeds into $\bar{\mathbf Q}_\ell$ through $x$, so it is algebraic over the fraction field of $\mathcal O$, of characteristic $0$. Hence there is a derivation $d$ of $R_{\mathfrak p}$ with values in the cotangent space $\mathfrak pR_{\mathfrak p}/\mathfrak p^2R_{\mathfrak p}$ over $\kappa(\mathfrak p)$ whose restriction to $\mathfrak pR_{\mathfrak p}$ is the projection. Composing $d$ with the coordinate forms of a basis of the cotangent space and the embedding of $\kappa(\mathfrak p)$, and restricting to $R$, gives derivations at $x$, which are linearly independent, as they take the values $\delta_{ij}$ at representatives of the basis. By (a) the cotangent space has dimension at most $N$.
\item[(c)] \emph{$\dim R_{\mathfrak p}/JR_{\mathfrak p}\ge N$.} Let $P\subseteq\mathfrak p$ be a prime minimal over $J$. The map $\Lambda\to R/P$ is injective: if $f(a)\in P$ with $a\ne0$, then, $P/J$ being a minimal prime of the noetherian ring $R/J$, some $b\notin J$ has $bf(a)\in J$ \cite[\S6]{Mat86}, against the hypothesis. So $T=R/P$ is a complete local domain, finite over $\Lambda$ and containing it, and $\dim T=\dim\Lambda=N+1$. The image of $T$ under $x$ is integral over $\mathcal O$, so $\dim T/(\mathfrak p/P)\le1$. As $T$ is a complete local domain, the height of $\mathfrak p/P$ plus $\dim T/(\mathfrak p/P)$ is $\dim T$ \cite[\S31]{Mat86}; so $\mathfrak p/P$ has height at least $N$. As $R_{\mathfrak p}/JR_{\mathfrak p}$ maps onto $T_{\mathfrak p/P}$, its dimension is at least $N$.
\item[(d)] \emph{Regularity.} By (b) and (c), the number of generators of the maximal ideal of the noetherian local ring $R_{\mathfrak p}$ is at most $\dim R_{\mathfrak p}/JR_{\mathfrak p}\le\dim R_{\mathfrak p}$, which is at most that number. So $R_{\mathfrak p}$ is regular, hence a domain \cite[\S14]{Mat86}, and $\dim R_{\mathfrak p}/JR_{\mathfrak p}=\dim R_{\mathfrak p}$ with $JR_{\mathfrak p}$ proper; so $JR_{\mathfrak p}=0$.
\item[(e)] \emph{Conclusion.} Let $a\in J$. By (d), $sa=0$ for some $s\notin\mathfrak p$. If $\mathfrak q\subseteq\mathfrak p$ is a prime, then $s\notin\mathfrak q$, so $a\in\mathfrak q$.
\end{enumerate}
\end{proof}
% regOfTan (RegOfTan.lean, 97 lines): (a) is the private rank_le_res with derLam (DerLam.lean, 87 lines: D' maps to (D'(X_i))_i, injective by the private ptDer_eq_zero_of_vanish_on_X, which forms x o f + eps D', uses Alg.co_mvps for Lambda, and applies derZero (coPt for the smallness, Lemma sha-automorphy-unramified-near-tangent-small) with the private derLam_dense, the truncation of GlobalTheory.algHom_sub_aeval_truncTotal_mem);(b) is tanEmb (TanEmb.lean, 118 lines: the private dual_der with cotDer, CotDer.lean); (c) is dimLow (DimLow.lean, 71 lines: Ideal.exists_mul_mem_of_mem_minimalPrimes and htLow, DimLeOfInt.lean, with Alg.co_caten for the dimension formula of the complete local domain and mvPowerSeries_cotInd for dim Lambda); (d) is AlgebraTools.regOfDim; (e) is IsLocalization.map_eq_zero_iff.

\begin{lemma}\label{sha-siegel-step-ideal-seeds-inputs-tangent-lift}
In Situation~\ref{sha-hecke-situation}, let $A=\bar{\mathbf Q}_\ell[\epsilon]/(\epsilon^2)$. Let $x$ be a point of $B$ and $\rho:\Gamma_0\to GL_n(\bar{\mathbf Q}_\ell)$ a homomorphism whose values span $M_n(\bar{\mathbf Q}_\ell)$, with $x(D^{\mathrm u})$ the determinant of $\rho$, and let $x':B\to A$ be a homomorphism of $\mathcal O$-algebras that reduces to $x$ modulo $\epsilon$. Then:
\begin{enumerate}
\item there are a homomorphism $\rho':\Gamma_0\to GL_n(A)$ whose determinant is $x'(D^{\mathrm u})$ and a continuous map $y:\Gamma_0\to M_n(\bar{\mathbf Q}_\ell)$ such that $\rho'(g)=(1+\epsilon y(g))\rho(g)$ and $y(gh)=y(g)+\rho(g)y(h)\rho(g)^{-1}$ for $g,h\in\Gamma_0$; and $\rho$ is continuous;
\item for $v\mid\ell$ and $g_1,\dots,g_n\in D_v$, the product $\big(\rho'(g_1)-x'(\psi^{\mathrm u}_{v,1}(g_1))\big)\cdots\big(\rho'(g_n)-x'(\psi^{\mathrm u}_{v,n}(g_n))\big)$ is $0$; and $\rho'(g)=1$ for $u\in S_{\mathrm a}$ and $g\in I_u\cap\bigcap_i\ker\bar\chi_i$;
\item there are $P\in GL_n(\bar{\mathbf Q}_\ell)$ with ${}^t\rho(g)P\rho(cgc^{-1})=\varepsilon(g)^{1-n}P$ for $g\in\Gamma_0$, an adjoint extension $R$ for $P$ (Definition~\ref{sha-automorphy-unramified-near-locally-generic}(2)), and a continuous map $z:\Gamma\to M_n(\bar{\mathbf Q}_\ell)$ with $z(gh)=z(g)+R(g)z(h)$ for $g,h\in\Gamma$ and $z(g)=y(g)$ for $g\in\Gamma_0$.
\end{enumerate}
\end{lemma}
\leanref{Development/AutomorphicArithmetic/AutomorphicVarieties/WeakShimuraTheory/EigenbasisLifts/LiftRep.lean}{UpperTheory.V.Ws.Eb.liftRep}
\leanref{Development/AutomorphicArithmetic/AutomorphicVarieties/WeakShimuraTheory/EigenbasisLifts/RelKill.lean}{UpperTheory.V.Ws.Eb.relKill}
\leanref{Development/AutomorphicArithmetic/AutomorphicVarieties/WeakShimuraTheory/EigenbasisLifts/RelForm.lean}{UpperTheory.V.Ws.Eb.relForm}
\leanref{Development/AutomorphicArithmetic/AutomorphicVarieties/WeakShimuraTheory/EigenbasisLifts/CocRel.lean}{UpperTheory.V.Ws.Eb.cocRel}
\leanref{Development/AutomorphicArithmetic/AutomorphicVarieties/WeakShimuraTheory/EigenbasisLifts/CocExt.lean}{UpperTheory.V.Ws.Eb.cocExt}
\begin{proof}
Two facts about a homomorphism $\sigma:\Gamma_0\to GL_n(A')$, over a commutative ring $A'$, whose values span $M_n(A')$ are used; they follow from the nondegeneracy of the trace form on $M_n(A')$. First, if an element $w$ of $A'[\Gamma_0]$ has $\operatorname{tr}(\sigma(h)\sigma(w))=0$ for all $h\in\Gamma_0$, then $\sigma(w)=0$. Second, if $A'$ is a topological ring and $g\mapsto\operatorname{tr}\sigma(g)$ is continuous, then $\sigma$ is continuous: writing a matrix unit $E_{ji}$ as $\sum_ka_k\sigma(h_k)$, $\sigma(g)_{ij}=\sum_ka_k\operatorname{tr}\sigma(gh_k)$.
\begin{enumerate}
\item[(1)] The ring $A$ is a complete local noetherian ring with residue field $\bar{\mathbf Q}_\ell$, and $x'(D^{\mathrm u})$ reduces to $x(D^{\mathrm u})$, the determinant of $\rho$, whose values span $M_n(\bar{\mathbf Q}_\ell)$. So $x'(D^{\mathrm u})$ is the determinant of a homomorphism $\rho':\Gamma_0\to GL_n(A)$ whose values span $M_n(A)$, by \cite[Theorem~2.22(i)]{Che14} applied to the Cayley-Hamilton quotient of $A[\Gamma_0]$; its reduction has the determinant of $\rho$, so it is conjugate to $\rho$ (Theorem~2.22(i) over $\bar{\mathbf Q}_\ell$), and after conjugating $\rho'$ by a lift of the conjugating matrix, $\rho'$ reduces to $\rho$. The traces of $\rho'$ are the images under $x'$ of the traces of $D^{\mathrm u}$, which are $\mathfrak m_B$-adically continuous on $\Gamma_0$; and $x'$ is continuous for the $\mathfrak m_B$-adic topology, by Lemmas~\ref{sha-automorphy-unramified-near-tangent-points} and~\ref{sha-automorphy-unramified-near-tangent-small}, applied to $x$ and to the coefficient of $\epsilon$ in $x'$, which is a derivation at $x$. So $\rho'$ is continuous, and so is $\rho$, its reduction. Let $y(g)$ be the coefficient of $\epsilon$ in $\rho'(g)\rho(g)^{-1}$; it is continuous, $\rho'(g)=(1+\epsilon y(g))\rho(g)$, and comparing the coefficients of $\epsilon$ in $\rho'(gh)=\rho'(g)\rho'(h)$ gives the cocycle relation.
\item[(2)] The elements $w=\prod_i(g_i-\psi^{\mathrm u}_{v,i}(g_i))$ and $w=g-1$ of (2) and (3) of Situation~\ref{sha-hecke-situation} satisfy $\operatorname{tr}_{D^{\mathrm u}}(hw)=0$ for every $h\in\Gamma_0$. Applying $x'$, $\operatorname{tr}(\rho'(h)\rho'(w'))=0$ for the image $w'$ of $w$ in $A[\Gamma_0]$, so $\rho'(w')=0$ by the first fact. This is the claim.
\item[(3)] \emph{The form.} Let $\rho'^\vee(g)=\varepsilon(g)^{1-n}\,{}^t\rho'\big((cgc^{-1})^{-1}\big)$, a homomorphism $\Gamma_0\to GL_n(A)$. By (1) of Situation~\ref{sha-hecke-situation}, applied through $x'$, $\rho'$ and $\rho'^\vee$ have the same trace at every element of $A[\Gamma_0]$. As the values of $\rho'$ span $M_n(A)$, there is $U\in GL_n(A)$ with $\rho'^\vee(g)U=U\rho'(g)$ for all $g$: the extension of $\rho'^\vee$ to $A[\Gamma_0]$ factors through that of $\rho'$, by the first fact, through an endomorphism of the $A$-algebra $M_n(A)$ that preserves traces; it is injective, as the trace form is nondegenerate, hence an automorphism, and inner, as $A$ is local. Transposing, $P'={}^tU$ satisfies ${}^t\rho'(g)P'\rho'(cgc^{-1})=\varepsilon(g)^{1-n}P'$.
\item[(4)] \emph{The relation for $y$.} Write $P'=P+\epsilon P_1$ with $P\in GL_n(\bar{\mathbf Q}_\ell)$, and $Z=P^{-1}P_1$. The part of the relation of (3) without $\epsilon$ is ${}^t\rho(g)P\rho(cgc^{-1})=\varepsilon(g)^{1-n}P$. Its coefficient of $\epsilon$, with $\rho'(g)=(1+\epsilon y(g))\rho(g)$, gives
\[
y(cgc^{-1})=-P^{-1}\,{}^ty(g)P+\rho(cgc^{-1})Z\rho(cgc^{-1})^{-1}-Z.
\]
As $\varepsilon(cgc^{-1})=\varepsilon(g)$ and the values of $\rho$ span $M_n(\bar{\mathbf Q}_\ell)$, Lemma~\ref{sha-automorphy-unramified-near-matrix-adjoint}, with $G=\Gamma$ and $G_0=\Gamma_0$, gives an adjoint extension $R$ for $P$, with $R(c)X=-P^{-1}\,{}^tXP$. The relation reads $y(cgc^{-1})=R(c)y(g)+R(cgc^{-1})Z-Z$.
\item[(5)] \emph{The cocycle $z$.} Let $w=\frac12(R(c)Z-Z)$, so that $R(c)w=-w$, as $R(c)^2=1$. Applying the relation of (4) to $cgc^{-1}$ and substituting it into itself gives $R(g)(Z+R(c)Z)=Z+R(c)Z$ for $g\in\Gamma_0$; hence $y(cgc^{-1})=R(c)y(g)+w-R(cgc^{-1})w$. Let $z(g)=y(g)$ for $g\in\Gamma_0$ and $z(g)=y(gc)+R(gc)w$ for $g\notin\Gamma_0$. As $\Gamma_0$ is open and $y$ and $g\mapsto R(g)w=\rho(g)w\rho(g)^{-1}$ are continuous on $\Gamma_0$, $z$ is continuous. The cocycle relation $z(gh)=z(g)+R(g)z(h)$ holds in the four cases according as $g$ and $h$ lie in $\Gamma_0$: for $g,h\in\Gamma_0$ it is (1); otherwise it follows from (1), $c^2=1$, the relation for $y(cgc^{-1})$ just shown and $R(c)w=-w$, writing $ghc=gc\cdot chc^{-1}$ when $g\notin\Gamma_0$ and $h\in\Gamma_0$, and $gh=gc\cdot c(hc)c^{-1}$ when $g,h\notin\Gamma_0$.
\end{enumerate}
\end{proof}
% liftRep (LiftRep.lean, 68 lines): detRep (UpperTheory.Reformulated.B.CayleyHamiltonSymmetric.detRep, CapelliIdentities/CayleyHamiltonSymmetric.lean, which the Lean proves itself; the text cites Chenevier), trPlus (TrPlus.lean, the two facts), the private lr_cont2 with coPt for the continuity of x', and tanOf (EbSMul.lean) for y. (2) is relKill (RelKill.lean). (3) is relForm (RelForm.lean, 97 lines: the private rf_dual for rho'^vee, B.u.P for the traces, and GlobalTheory.TraceConjugacy.exists_conj_of_surjective for U, which proves the factorization and inner automorphism through OrzechProperty). (4) is cocRel (CocRel.lean, 99 lines, the private cr_alg for the displayed identity) and adExt (Lemma sha-automorphy-unramified-near-matrix-adjoint). (5) is cocExt (CocExt.lean, 168 lines: hinv, w, hrel, and the private cocycle_extends_to_conj_group for z).

\begin{lemma}\label{sha-siegel-step-ideal-seeds-inputs-tangent-v0-local-model}
Let $\mathcal O\subset\bar{\mathbf Q}_\ell$ be the ring of integers of a finite extension of $\mathbf Q_\ell$, $H$ a compact Hausdorff group and $\rho:H\to GL_n(\bar{\mathbf Q}_\ell)$ a continuous homomorphism. Then there are the ring of integers $\mathcal O'\supseteq\mathcal O$ of a finite extension of $\mathbf Q_\ell$ in $\bar{\mathbf Q}_\ell$, a matrix $Q\in GL_n(\bar{\mathbf Q}_\ell)$ and a homomorphism $\rho_0:H\to GL_n(\mathcal O')$ with $\rho(g)Q=Q\rho_0(g)$ for $g\in H$, such that for every $N\ge0$ the entries of $\rho_0$ are locally constant modulo $\mathfrak m_{\mathcal O'}^N$.
\end{lemma}
\leanref{Development/AutomorphicArithmetic/AutomorphicVarieties/WeakShimuraTheory/EigenbasisLifts/IntModel.lean}{UpperTheory.V.Ws.Eb.intModel}
\leanref{Development/AutomorphicArithmetic/AutomorphicVarieties/PatchingModules/DeterminantPadicFree/LatInt.lean}{UpperTheory.V.Pa.Det.latInt}
\begin{proof}
\begin{enumerate}
\item[(a)] \emph{Coefficients.} As $H$ is compact and $\rho$ is continuous, the entries of $\rho$ lie in a finite extension $E_1$ of $\mathbf Q_\ell$, by the Baire category theorem. The ring $\mathcal O$ is finite over $\mathbf Z_\ell$, so it lies in the finite extension $E_2$ generated by finitely many generators, and its elements have absolute value at most $1$, being integral over $\mathbf Z_\ell$. Let $E=E_1E_2$ and $\mathcal O'=\{x\in E:|x|\le1\}$, its ring of integers, which contains $\mathcal O$.
\item[(b)] \emph{A stable lattice.} The entries of $\rho$ are continuous on the compact group $H$, so they are bounded by some $C$. Let $N_0\ge0$ with $|\ell|^{N_0}C\le1$, so that $\ell^{N_0}\rho(g)$ has entries in $\mathcal O'$ for every $g\in H$. Let $\Lambda\subseteq E^n$ be the $\mathcal O'$-module spanned by the columns of the matrices $\rho(g)$, $g\in H$. It lies in $\ell^{-N_0}\mathcal O'^n$, so it is finitely generated, and it contains the standard basis; so it is free of rank $n$ over the discrete valuation ring $\mathcal O'$. It is stable under $\rho(H)$. Let $Q$ be the matrix whose columns form a basis of $\Lambda$. Then $\rho_0(g)=Q^{-1}\rho(g)Q$ has entries in $\mathcal O'$, as the columns of $\rho(g)Q$ lie in $\Lambda$, and so does its inverse $\rho_0(g^{-1})$.
\item[(c)] \emph{Local constancy.} The entries of $\rho_0$ are continuous. Let $N\ge0$ and $g\in H$. For $g'$ near $g$ the entries of $\rho_0(g')-\rho_0(g)$, which lie in $E$, have absolute value at most $|\ell|^N$, so they are $\ell^N$ times elements of $\mathcal O'$, and lie in $\mathfrak m_{\mathcal O'}^N$, as $\ell\in\mathfrak m_{\mathcal O'}$.
\end{enumerate}
\end{proof}
% intModel (IntModel.lean, 130 lines): GlobalTheory.CompactRepField.exists_finiteDimensional_entries_mem (the Baire category argument) for E_1, the private im_field for E_2, X.So.coefOf for O'; the private im_model (the bound C by compactness, im_small for N_0, and Pa.Det.latInt, whose private exists_stable_lattice is the span of the columns of the r(g)) for (b); the private im_cont for (c).

\begin{lemma}\label{sha-siegel-step-ideal-seeds-inputs-tangent-v0-local-tame}
Let $p$ be a prime, $H$ a topological group with an open embedding $\phi$ into $\operatorname{Gal}(\bar{\mathbf Q}_p/\mathbf Q_p)$, through which $H$ acts on $\bar{\mathbf Q}_p$, and $I\subseteq H$ the inverse image of the inertia group. Let $m\ge1$ be prime to $p$, $U\subseteq H$ an open subgroup, and $c:I\to\bar{\mathbf Q}_p^\times$ a homomorphism with $c^m=1$ that is trivial on $I\cap U$. Then $c(gag^{-1})=g(c(a))$ for $g\in H$ and $a\in I$.
\end{lemma}
\leanref{Development/AutomorphicArithmetic/AutomorphicVarieties/PadicHodgeTheory/ExistsNontrivialCharacterPgroup.lean}{tame_character_conjugation}
\leanref{Development/AutomorphicArithmetic/AutomorphicVarieties/PadicHodgeTheory/ExistsNontrivialCharacterPgroup.lean}{exists_eigenvector_tame_inertia}
\begin{proof}
Three facts about $\bar{\mathbf Q}_p$ are used: its automorphisms preserve the absolute value; an element $\sigma$ of the inertia group has $|\sigma(y)-y|<1$ for $|y|\le1$; and two distinct $m$-th roots of unity $\zeta,\zeta'$, with $p\nmid m$, have $|\zeta-\zeta'|=1$, as $X^m-1$ is separable modulo $p$. The values of $c$ are $m$-th roots of unity, which lie in the maximal unramified extension of $\mathbf Q_p$, so they are fixed by $I$.
\begin{enumerate}
\item[(a)] \emph{An eigenvector.} There is $x\in\bar{\mathbf Q}_p$, $x\ne0$, with $a(x)=c(a)x$ for every $a\in I$. Indeed, the kernel of $c$ contains $I\cap U$, so it is open, hence closed, in $I$, which is closed in $H$. As the image of $\phi$ is an open, hence closed, subgroup, $\phi$ is a closed embedding, and $J=\phi(\ker c)$ is a closed subgroup of $\operatorname{Gal}(\bar{\mathbf Q}_p/\mathbf Q_p)$. Let $E$ be its fixed field. By infinite Galois theory, $J$ is the group of automorphisms that fix $E$; so for $a,b\in I$, $c(a)=c(b)$ if and only if $a$ and $b$ agree on $E$. The image of $c$ is finite; for each of its elements $q$ choose $s_q\in I$ with $c(s_q)=q$. The restrictions of the $s_q$ to $E$ are pairwise distinct homomorphisms $E^\times\to\bar{\mathbf Q}_p^\times$, so by the linear independence of characters some $z\in E$ has $x=\sum_qq^{-1}s_q(z)\ne0$. For $a\in I$, $a$ fixes each $q^{-1}$, and $as_q$ agrees with $s_{c(a)q}$ on $E$, as both have the value $c(a)q$ under $c$. So $a(x)=\sum_qq^{-1}s_{c(a)q}(z)=c(a)\sum_{q'}q'^{-1}s_{q'}(z)=c(a)x$, with $q'=c(a)q$.
\item[(b)] \emph{Conjugation.} Let $g\in H$, $a\in I$ and $a'=gag^{-1}$, which lies in the normal subgroup $I$. Then $a(g^{-1}(x))=g^{-1}(a'(x))=g^{-1}(c(a'))\,g^{-1}(x)$. So $y=g^{-1}(x)/x$ has $a(y)-y=\big(g^{-1}(c(a'))-c(a)\big)y/c(a)$, with $|y|=1$ and $|c(a)|=1$. As $a$ lies in the inertia group, $|a(y)-y|<1$, so $|g^{-1}(c(a'))-c(a)|<1$. Both are $m$-th roots of unity, so they are equal, and $c(a')=g(c(a))$.
\end{enumerate}
\end{proof}
% tame_character_conjugation (ExistsNontrivialCharacterPgroup.lean): the private exists_eigenvector_tame_inertia for (a) (hfix for the values of c, hJ for the closed subgroup J, InfiniteGalois.fixingSubgroup_fixedField, key, linearIndependent_monoidHom for F, and Fintype.sum_equiv for the reindexing), then e1, hy, Loc.iqBasic (the inertia group acts trivially on the residue field) and KStar.IdealValuations.norm_aut_and_roots for (b).

\begin{lemma}\label{sha-siegel-step-ideal-seeds-inputs-tangent-v0-local-integral}
Let $p\ne\ell$ be primes, $H$ a topological group with an open embedding $\phi$ into $\operatorname{Gal}(\bar{\mathbf Q}_p/\mathbf Q_p)$, and $I\subseteq H$ the inverse image of the inertia group. Let $A$ be a complete discrete valuation ring that is a $\mathbf Z_\ell$-algebra, with maximal ideal $\mathfrak m$ and finite residue field $k$ of characteristic $\ell$, and for $g\in H$ let $\bar\chi(g)\in k$ be the residue of the image in $A$ of $\chi_{\mathrm{cyc}}(\phi(g))\in\mathbf Z_\ell^\times$. Let $\kappa:H\to k^\times$ be a character trivial on $I$, and suppose that $\bar\chi(g)\ne1$ for some $g\in H$. Let $r:H\to GL_n(A)$ be a homomorphism such that for every $N\ge0$ the entries of $r$ are locally constant modulo $\mathfrak m^N$, and such that the reduction of $r(a)$ has characteristic polynomial $(X-\kappa(a))^n$ for every $a\in H$. Then $r$ is trivial on $I$.
\end{lemma}
\leanref{Development/AutomorphicArithmetic/AutomorphicVarieties/PadicHodgeTheory/ExistsNontrivialCharacterPgroup.lean}{UpperTheory.V.Loc.v0b}
\begin{proof}
Let $H$ act on $\bar{\mathbf Q}_p$ through $\phi$. Fix $g\in H$ with $\bar\chi(g)\ne1$, and $c\in A$ with residue $\kappa(g)$. For a matrix $M$ over $A$ and $b\in A$, write $M\equiv b$ modulo an ideal when the entries of $M-b\cdot1$ lie in it.
\begin{enumerate}
\item[(a)] \emph{Powers of $g$.} For $j\ge0$, $\bar\chi(g^{\ell^j})-1=(\bar\chi(g)-1)^{\ell^j}\ne0$, as $k$ has characteristic $\ell$.
\item[(b)] \emph{Congruences.} Let $a\in H$ and $b\in A$ with residue $\kappa(a)$. Then $r(a)^{\ell^{n+N}}\equiv b^{\ell^{n+N}}$ modulo $\mathfrak m^{N+1}$ for every $N\ge0$. For $N=0$: by the Cayley-Hamilton theorem the reduction $\bar M$ of $r(a)$ has $(\bar M-\bar b)^n=0$, so $(\bar M-\bar b)^{\ell^n}=0$, as $\ell^n\ge n$; as $\bar b$ is a scalar and $k$ has characteristic $\ell$, $(\bar M-\bar b)^{\ell^n}=\bar M^{\ell^n}-\bar b^{\ell^n}$. For the step from $N$ to $N+1$, let $M=r(a)^{\ell^{n+N}}$ and $b'=b^{\ell^{n+N}}$, so that $M\equiv b'$ modulo $\mathfrak m^{N+1}$. Then $M^\ell-b'^\ell=\big(\sum_{i<\ell}M^ib'^{\ell-1-i}\big)(M-b')$, and the sum is congruent to $\ell b'^{\ell-1}$ modulo $\mathfrak m$, so its entries lie in $\mathfrak m$, which contains $\ell$. So $M^\ell\equiv b'^\ell$ modulo $\mathfrak m^{N+2}$.
\item[(c)] \emph{Modulo $\mathfrak m^N$.} Let $N\ge0$ and $r_N:H\to GL_n(A/\mathfrak m^N)$ the reduction of $r$. The ring $A/\mathfrak m^N$ is finite, and the kernel $U$ of $r_N$ is open, by the local constancy. For $a\in I$, $\kappa(a)=1$, so (b) with $b=1$ gives $r_N(a)^{\ell^{n+N}}=1$; so the finite group $P=r_N(I)$ is an $\ell$-group. Suppose that $P\ne1$. As $P$ is nilpotent, its commutator subgroup is proper, and a character of its abelianization gives a character $\psi:P\to\bar{\mathbf Q}_p^\times$ and $a\in I$ with $\psi(r_N(a))\ne1$. The homomorphism $c_\psi=\psi\circ r_N$ on $I$ has $c_\psi^{\ell^{n+N}}=1$, with $\ell^{n+N}$ prime to $p$, and is trivial on $I\cap U$. Let $g'=g^{\ell^{n+N}}$ and $a'=g'ag'^{-1}\in I$. By (b), applied to $g$ and $c$, $r_N(g')$ is a scalar, so $r_N(a')=r_N(a)$ and $c_\psi(a')=c_\psi(a)$. By Lemma~\ref{sha-siegel-step-ideal-seeds-inputs-tangent-v0-local-tame}, $c_\psi(a')=g'(c_\psi(a))$. So $g'$ fixes $\zeta=c_\psi(a)$, a root of unity of order a power of $\ell$ other than $1$, and so a power $\zeta'$ of $\zeta$ that is a primitive $\ell$-th root of unity. As $g'(\zeta')=\zeta'^{\chi_{\mathrm{cyc}}(\phi(g'))}$, $\chi_{\mathrm{cyc}}(\phi(g'))\equiv1$ modulo $\ell$, and $\bar\chi(g')=1$, as $\ell\in\mathfrak m$. This contradicts (a) for $j=n+N$. So $P=1$, and $r(a)\equiv1$ modulo $\mathfrak m^N$ for $a\in I$.
\item[(d)] \emph{Conclusion.} For $a\in I$, the entries of $r(a)-1$ lie in every power of $\mathfrak m$, so they vanish.
\end{enumerate}
\end{proof}
% Loc.v0b (ExistsNontrivialCharacterPgroup.lean, 322 lines), its first conclusion (the second, on X with a twisted commutation relation, is not used here): hgk for (a); hpow for (b), through the private unipotent_residue_pow_mod_maximal (N = 0) and matrix_pow_l_iterate_congruence with matrix_pow_l_congruence_succ (the step); main for (c), with the private exists_nontrivial_character_pgroup, tame_character_conjugation (Lemma sha-siegel-step-ideal-seeds-inputs-tangent-v0-local-tame) and cyclotomic_congruence_fixed_root; IsHausdorff.haus for (d).

\begin{lemma}\label{sha-siegel-step-ideal-seeds-inputs-tangent-v0-local}
Let $p\ne\ell$ be primes, $H$ a compact group with an open embedding into $\operatorname{Gal}(\bar{\mathbf Q}_p/\mathbf Q_p)$, and $I\subseteq H$ the inverse image of the inertia group. Let $\mathcal O$ be the ring of integers of a finite extension of $\mathbf Q_\ell$, with residue field $k$, and $\kappa:H\to k^\times$ a character trivial on $I$, such that the reduction of $\chi_{\mathrm{cyc}}$ modulo $\ell$ is not trivial on $H$. Let $\rho:H\to GL_n(\bar{\mathbf Q}_\ell)$ be a continuous representation such that for every $g\in H$ and every $a\in\mathcal O$ with residue $\kappa(g)$, every eigenvalue $t$ of $\rho(g)$ has $|t-a|<1$. Then $\rho$ is trivial on $I$.
\end{lemma}
\leanref{Development/AutomorphicArithmetic/AutomorphicVarieties/WeakShimuraTheory/EigenbasisLifts/InertiaTrivialOfCharpolyRootsNearCharacter.lean}{UpperTheory.V.Ws.Eb.inertia_trivial_of_charpoly_roots_near_character}
\begin{proof}
Let $\phi$ be the open embedding.
\begin{enumerate}
\item[(a)] \emph{An integral model.} The group $H$ is Hausdorff, as $\phi$ embeds it in a Hausdorff group. Lemma~\ref{sha-siegel-step-ideal-seeds-inputs-tangent-v0-local-model} gives $\mathcal O'\supseteq\mathcal O$, $Q$ and $\rho_0:H\to GL_n(\mathcal O')$ with $\rho(g)=Q\rho_0(g)Q^{-1}$, whose entries are locally constant modulo every power of $\mathfrak m_{\mathcal O'}$. The ring $\mathcal O'$ is a complete discrete valuation ring with finite residue field $k'$ of characteristic $\ell$. An element of $\mathcal O$ or of $\mathcal O'$ lies in the maximal ideal if and only if its absolute value is less than $1$; so the inclusion $\mathcal O\subseteq\mathcal O'$ is local and induces an injective map $f:k\to k'$.
\item[(b)] \emph{Characteristic polynomials.} Let $g\in H$ and $a\in\mathcal O$ with residue $\kappa(g)$, and let $t_1,\dots,t_n$ be the eigenvalues of $\rho(g)$, with multiplicity. By hypothesis $|t_i-a|<1$, and $|a|\le1$. By induction on $j$, the coefficients of $\prod_{i\le j}(X-t_i)-(X-a)^j$ have absolute value less than $1$, and those of $(X-a)^j$ at most $1$: write $(X-t_{j+1})P-(X-a)^{j+1}=(X-t_{j+1})\big(P-(X-a)^j\big)+(a-t_{j+1})(X-a)^j$, use $|t_{j+1}|\le1$ and the ultrametric inequality. The characteristic polynomial of $\rho_0(g)=Q^{-1}\rho(g)Q$ is $\prod_i(X-t_i)$, with coefficients in $\mathcal O'$. Its coefficients differ from those of $(X-a)^n$ by elements of $\mathcal O'$ of absolute value less than $1$, which lie in $\mathfrak m_{\mathcal O'}$; so its reduction is $(X-f(\kappa(g)))^n$.
\item[(c)] \emph{Conclusion.} Lemma~\ref{sha-siegel-step-ideal-seeds-inputs-tangent-v0-local-integral} applies to $\phi$, $A=\mathcal O'$, the character $f\circ\kappa$, which is trivial on $I$, and $r=\rho_0$, by (a) and (b): the residue in $k'$ of $\chi_{\mathrm{cyc}}(\phi(g))$ is the image under $f$ of its residue in $k$, which is not $1$ for some $g\in H$. So $\rho_0$ is trivial on $I$, and so is $\rho=Q\rho_0Q^{-1}$.
\end{enumerate}
\end{proof}
% inertia_trivial_of_charpoly_roots_near_character (InertiaTrivialOfCharpolyRootsNearCharacter.lean, 113 lines): Ws.Eb.intModel (Lemma sha-siegel-step-ideal-seeds-inputs-tangent-v0-local-model) with abelian₂.dvr_finite and Fr.sc_norm for (a); the private vb_coeff (the induction of (b)) with e1 and Matrix.charpoly_map for (b); Loc.v0b (Lemma sha-siegel-step-ideal-seeds-inputs-tangent-v0-local-integral) with W.pushCh f κ for (c). With intModel (130 lines) and Loc.v0b (ExistsNontrivialCharacterPgroup.lean, 322 lines), about 560 lines.

\begin{lemma}\label{sha-siegel-step-ideal-seeds-inputs-tangent-v0}
In Situation~\ref{sha-hecke-situation}, suppose that the $\bar\chi_i$ are trivial on $I_{v_0}$ and take equal values on $D_{v_0}$, and that $\varepsilon\bmod\ell$ is nontrivial on $D_{v_0}$. Let $x$ be a point of $B$ and $\rho:\Gamma_0\to GL_n(\bar{\mathbf Q}_\ell)$ a continuous homomorphism with $x(D^{\mathrm u})$ the determinant of $\rho$, and let $y:\Gamma_0\to M_n(\bar{\mathbf Q}_\ell)$ be continuous with $y(gh)=y(g)+\rho(g)y(h)\rho(g)^{-1}$. Then $y$ vanishes on $I_{v_0}$.
\end{lemma}
\leanref{Development/AutomorphicArithmetic/AutomorphicVarieties/WeakShimuraTheory/EigenbasisLifts/CocycleVanishesOnInertiaOfPointPseudodeformation.lean}{UpperTheory.V.Ws.Eb.cocycle_vanishes_on_inertia_of_point_pseudodeformation}
\leanref{Development/AutomorphicArithmetic/AutomorphicVarieties/WeakShimuraTheory/EigenbasisLifts/DblRep.lean}{UpperTheory.V.Ws.Eb.dblRep}
\leanref{Development/AutomorphicArithmetic/AutomorphicVarieties/WeakShimuraTheory/EigenbasisLifts/RootsNear.lean}{UpperTheory.V.Ws.Eb.rootsNear}
\begin{proof}
\begin{enumerate}
\item[(a)] \emph{Doubling.} Let $\rho_2(g)=\begin{pmatrix}\rho(g)&y(g)\rho(g)\\0&\rho(g)\end{pmatrix}$. By the cocycle relation, $\rho_2$ is a continuous homomorphism $\Gamma_0\to GL_{2n}(\bar{\mathbf Q}_\ell)$; the characteristic polynomial of $\rho_2(g)$ is the square of that of $\rho(g)$; and $\rho_2(g)=1$ implies $y(g)=0$.
\item[(b)] \emph{Eigenvalues.} Let $g\in D_{v_0}$ and $a\in\mathcal O$ with residue $\bar\chi_1(g)$, which is also $\bar\chi_i(g)$ for every $i$. The characteristic polynomial of $\rho(g)$ is the image under $x$ of that of $D^{\mathrm u}$ at $g$, which reduces to $\prod_i(X-\bar\chi_i(g))=(X-\bar a)^n$; by Lemma~\ref{sha-automorphy-unramified-near-tangent-points}, the coefficients of its difference with $(X-a)^n$ have absolute value less than $1$. The roots of a monic polynomial whose coefficients are within distance less than $1$ of those of $(X-a)^n$, with $|a|\le1$, are within distance less than $1$ of $a$. So every eigenvalue $t$ of $\rho(g)$, and so of $\rho_2(g)$, has $|t-a|<1$.
\item[(c)] \emph{Conclusion.} Let $H$ be the decomposition group of $\operatorname{Gal}(\bar L/L)$ at a place above $\tilde v_0$ whose image in $\Gamma_0$ contains $I_{v_0}$ as the image of its inertia group; as the residue field of $v_0$ is $\mathbf F_p$ with $p\ne\ell$ unramified in $L$, $H$ has an open embedding into $\operatorname{Gal}(\bar{\mathbf Q}_p/\mathbf Q_p)$. By (b), Lemma~\ref{sha-siegel-step-ideal-seeds-inputs-tangent-v0-local} applies to the composite of $H\to\Gamma_0$ with $\rho_2$ and $\kappa=\bar\chi_1$: $\kappa$ is trivial on the inertia group, and $\chi_{\mathrm{cyc}}\bmod\ell$ is not trivial on $H$. So $\rho_2$ is trivial on $I_{v_0}$, and $y$ vanishes there by (a).
\end{enumerate}
\end{proof}
% cocycle_vanishes_on_inertia_of_point_pseudodeformation (67 lines): r.V0ok for the residual conditions, Loc.decEmb at s.A s.v0 for H, dblRep (DblRep.lean, 68 lines) for (a), rootsNear (RootsNear.lean, 89 lines) for (b), and inertia_trivial_of_charpoly_roots_near_character (Lemma sha-siegel-step-ideal-seeds-inputs-tangent-v0-local) for (c).

\begin{lemma}\label{sha-siegel-step-ideal-seeds-inputs-tangent-crystalline-descent}
Let $K$, $G_K$, $B_{\mathrm{cris}}$ and the crystalline cocycles be as in Lemma~\ref{sha-automorphy-unramified-tangents-unit-duality}. Let $V$ be a finite-dimensional $\mathbf Q_\ell$-vector space with a continuous linear action of $G_K$, $K'\subseteq\bar{\mathbf Q}_\ell$ a finite Galois extension of $K$, and $c$ a continuous cocycle of $V$. Then $c$ is a crystalline cocycle if and only if its restriction to $G_{K'}$ is one.
\end{lemma}
\leanref{Development/AutomorphicArithmetic/AutomorphicVarieties/PadicHodgeTheory/Z1fDescAvg.lean}{UpperTheory.V.Loc.z1fDesc}
\begin{proof}
If $c(g)\otimes1=gx-x$ for every $g\in G_K$, the same $x$ serves for $G_{K'}$. Conversely, let $x\in V\otimes_{\mathbf Q_\ell}B_{\mathrm{cris}}$ with $c(h)\otimes1=hx-x$ for $h\in G_{K'}$, and let $f(g)=c(g)\otimes1-(gx-x)$ for $g\in G_K$. Then $f(gh)=f(g)+gf(h)$, and $f$ vanishes on $G_{K'}$; so $f(gh)=f(g)$ for $h\in G_{K'}$, and $f$ defines a map $F$ on the set $G_K/G_{K'}$ of $d=[K':K]$ cosets, with $F(gq)=f(g)+gF(q)$. Let $S=\sum_qF(q)$. As $q\mapsto gq$ permutes the cosets, $S=d\,f(g)+gS$. So $c(g)\otimes1=f(g)+gx-x=gy-y$ for every $g\in G_K$, with $y=x-d^{-1}S$.
\end{proof}
% z1fDesc (Z1fDescAvg.lean, 62 lines): the private z1fDesc_avg, for a normal subgroup H of finite index, with f' = c - (gx - x), the map F on G/H and the sum S.

\begin{lemma}\label{sha-siegel-step-ideal-seeds-inputs-tangent-crystalline-extension}
Let $K$, $G_K$ and the crystalline representations and cocycles be as in Lemma~\ref{sha-automorphy-unramified-tangents-unit-crystalline}. Let $A$ be a finite-dimensional commutative $\mathbf Q_\ell$-algebra, $V$ and $V'$ finitely generated $A$-modules with continuous $A$-linear actions of $G_K$, and $\kappa:G_K\to A^\times$ a character such that $A(\kappa)$, the module $A$ on which $g$ acts by multiplication by $\kappa(g)$, is crystalline. Let $f:V\to V'$ be an injective $G_K$-equivariant $A$-linear map, $p:V'\to A$ an $A$-linear map with kernel $f(V)$, $s\in V'$ with $p(s)=1$, and $c:G_K\to V$ a map with $gs=\kappa(g)\big(s+f(c(g))\big)$ for every $g\in G_K$. Suppose that $V$ is crystalline. Then $c$ is a continuous cocycle of $V(\kappa^{-1})$, the module $V$ on which $g$ acts by $x\mapsto\kappa(g)^{-1}gx$, and $V'$ is crystalline if and only if $c$ is a crystalline cocycle of $V(\kappa^{-1})$.
\end{lemma}
\leanref{Development/AutomorphicArithmetic/AutomorphicVarieties/PadicHodgeTheory/MatrixExtendByZeroApply.lean}{UpperTheory.V.Loc.crExt}
\begin{proof}
\begin{enumerate}
\item[(a)] \emph{The twist.} In $V'(\kappa^{-1})$, $gs=s+f(c(g))$. As $f$ has a $\mathbf Q_\ell$-linear left inverse, $c$ is continuous, and comparing $(gh)s$ with $g(hs)$ gives $c(gh)=c(g)+\kappa(g)^{-1}gc(h)$. The proof of Lemma~\ref{sha-automorphy-unramified-tangents-unit-crystalline}(5) applies with $A$ in place of $E$: $A(\kappa^{-1})$ is crystalline, and so is the twist of a crystalline $A$-module by $\kappa$ or by $\kappa^{-1}$. So $V(\kappa^{-1})$ is crystalline, and $V'$ is crystalline if and only if $V'(\kappa^{-1})$ is.
\item[(b)] \emph{An extension of $\mathbf Q_\ell$.} Let $G_K$ act on $V\times\mathbf Q_\ell$ by $g(v,t)=(\kappa(g)^{-1}gv+tc(g),t)$, and call this representation $\sigma$. With $i(v)=(v,0)$, $p_2(v,t)=t$ and $s_0=(0,1)$, $gs_0-s_0=i(c(g))$. By Lemma~\ref{sha-automorphy-unramified-tangents-unit-crystalline}(7), applied to $V(\kappa^{-1})$, $\sigma$ is crystalline if and only if $c$ is a crystalline cocycle of $V(\kappa^{-1})$.
\item[(c)] \emph{Comparison.} The map $j(v,t)=f(v)+ts$ from $\sigma$ to $V'(\kappa^{-1})$ is $G_K$-equivariant, by (a). It is injective if $A\ne0$, as $p(j(v,t))=t$ and $f$ is injective; if $A=0$ all the spaces are $0$. So if $V'(\kappa^{-1})$ is crystalline, so is $\sigma$, by Lemma~\ref{sha-automorphy-unramified-tangents-unit-crystalline}(2). Conversely, give $A$ the trivial action, which is crystalline by Lemma~\ref{sha-automorphy-unramified-tangents-unit-crystalline}(6) with $w=0$. The $\mathbf Q_\ell$-bilinear map $A\times(V\times\mathbf Q_\ell)\to V'(\kappa^{-1})$, $(a,x)\mapsto a\,j(x)$, is $G_K$-equivariant, and its image spans $V'(\kappa^{-1})$, as every $y\in V'$ is $f(v)+p(y)s$ for some $v\in V$. So if $\sigma$ is crystalline, so is $V'(\kappa^{-1})$, by step (d) of the proof of Lemma~\ref{sha-automorphy-unramified-tangents-unit-crystalline}.
\end{enumerate}
The lemma follows from (a), (b) and (c).
\end{proof}
% crExt (MatrixExtendByZeroApply.lean, 273 lines in all), stated in the Lean for an open normal subgroup H, here H = G_K: hZ for the cocycle; isCr_inv and isCr_twist (part (5)) for (a); the private triv for (b) and (c), with ext0 for sigma, the private line (part (7)), isCrys_isDR_sub_quot (part (2)) for j, and admissible_of_bilinear_openSubgroup_galois with isCr_trivial (crUnr with w = 0) for the converse; z1fDesc for the passage to H.

\begin{lemma}\label{sha-siegel-step-ideal-seeds-inputs-tangent-crystalline-positive}
Let $K$, $G_K$, $\chi$, $B_{\mathrm{cris}}\subseteq B_{\mathrm{dR}}$ and the crystalline cocycles be as in Lemmas~\ref{sha-automorphy-unramified-tangents-unit-duality} and~\ref{sha-automorphy-unramified-tangents-unit-crystalline-periods}, and let $(\operatorname{Fil}^iB_{\mathrm{dR}})_{i\in\mathbf Z}$ be the filtration of $B_{\mathrm{dR}}$ by the powers of the maximal ideal of $B_{\mathrm{dR}}^+$ \cite[Definition~4.4.7]{BC09}. For a finite-dimensional $\mathbf Q_\ell$-vector space $W$ with a continuous linear action of $G_K$ write $D_{\mathrm{dR}}(W)=(W\otimes_{\mathbf Q_\ell}B_{\mathrm{dR}})^{G_K}$ and $\operatorname{Fil}^iD_{\mathrm{dR}}(W)=D_{\mathrm{dR}}(W)\cap(W\otimes_{\mathbf Q_\ell}\operatorname{Fil}^iB_{\mathrm{dR}})$, and call $W$ \emph{de Rham} if $\dim_{\mathbf Q_\ell}D_{\mathrm{dR}}(W)=[K:\mathbf Q_\ell]\dim_{\mathbf Q_\ell}W$. Let $E$ be a finite extension of $\mathbf Q_\ell$ and $V$ a finite-dimensional $E$-vector space with a continuous $E$-linear action of $G_K$. Suppose that:
\begin{enumerate}
\item[(i)] $V$ is de Rham;
\item[(ii)] there are $G_K$-stable $E$-subspaces $0=V_0\subseteq V_1\subseteq\dots\subseteq V_m=V$, vectors $x_i$ and characters $\kappa_i:G_K\to E^\times$ with $V_i=V_{i-1}+Ex_i$ and $gx_i-\kappa_i(g)x_i\in V_{i-1}$ for $g\in G_K$, such that each $\kappa_i$ agrees with $\chi^{w_i}$ on $U_i\cap I_K$ for some open subgroup $U_i$ of $G_K$ and some integer $w_i\ge1$;
\item[(iii)] the only $E$-linear $f:V\to E$ with $f(gx)=\chi(g)f(x)$ for all $g$ and $x$ is $0$.
\end{enumerate}
Then every continuous cocycle of $V$ is crystalline.
\end{lemma}
\leanref{Development/AutomorphicArithmetic/AutomorphicVarieties/PadicHodgeTheory/PrF.lean}{UpperTheory.V.Loc.prF}
\leanref{Development/AutomorphicArithmetic/AutomorphicVarieties/PadicHodgeTheory/PosBasic.lean}{UpperTheory.V.Loc.posBasic}
\leanref{Development/AutomorphicArithmetic/AutomorphicVarieties/PadicHodgeTheory/TateRat.lean}{UpperTheory.V.Loc.tateRat}
\leanref{Development/AutomorphicArithmetic/AutomorphicVarieties/PadicHodgeTheory/CharFil.lean}{UpperTheory.V.Loc.charFil}
\begin{proof}
\begin{enumerate}
\item[(a)] \emph{The filtration of $D_{\mathrm{dR}}(V)$.} $\operatorname{Fil}^0D_{\mathrm{dR}}(V)=0$ and $V^{G_K}=0$. For a continuous character $\kappa:G_K\to E^\times$ that agrees with $\chi^w$ on $U\cap I_K$ for an open subgroup $U$, $E(\kappa)$ is de Rham with $\operatorname{Fil}^iD_{\mathrm{dR}}(E(\kappa))=0$ for $i>-w$: the character $\kappa\chi^{-w}$ has open kernel on $I_K$, so $E(\kappa\chi^{-w})$ is de Rham with $D_{\mathrm{dR}}$ concentrated in filtration degree $0$ \cite[Example~6.3.9 and the remark after it]{BC09}, and the twist by $\chi^w$ moves the filtration jump to degree $-w$ \cite[Example~6.3.1]{BC09}. For $W\subseteq V$ let $T(W)$ be the $E$-span of the $v\otimes b$ with $v\in W$ and $b\in\operatorname{Fil}^0B_{\mathrm{dR}}$, so that $\operatorname{Fil}^0D_{\mathrm{dR}}(V)=D_{\mathrm{dR}}(V)\cap T(V_m)$. Let $1\le i\le m$, and suppose that $x_i\notin V_{i-1}$ (otherwise $V_i=V_{i-1}$). Let $\lambda:V_i\to E$ be $E$-linear with $\lambda(x_i)=1$ and $\lambda(V_{i-1})=0$; then $\lambda(gv)=\kappa_i(g)\lambda(v)$ for $v\in V_i$, so $\kappa_i(g)=\lambda(gx_i)$ is continuous, and $\operatorname{Fil}^0D_{\mathrm{dR}}(E(\kappa_i))=0$, as $w_i\ge1$. For $y\in D_{\mathrm{dR}}(V)\cap T(V_i)$, $(\lambda\otimes1)(y)$ is invariant in $E(\kappa_i)\otimes B_{\mathrm{dR}}$ and lies in $E\otimes\operatorname{Fil}^0B_{\mathrm{dR}}$, so it vanishes; hence $y=\big((1-x_i\lambda)\otimes1\big)(y)\in T(V_{i-1})$. By descending induction on $i$, $\operatorname{Fil}^0D_{\mathrm{dR}}(V)\subseteq T(V_0)=0$. If $v\in V^{G_K}$, then $v\otimes1\in\operatorname{Fil}^0D_{\mathrm{dR}}(V)=0$, so $v=0$.
\item[(b)] \emph{The finite part.} By Bloch and Kato's dimension formula \cite[Corollary~3.8.4]{BK90}, the subspace $H^1_f(K,V)$ of $H^1(K,V)$ of the classes of crystalline cocycles has dimension $\dim D_{\mathrm{dR}}(V)/\operatorname{Fil}^0D_{\mathrm{dR}}(V)+\dim V^{G_K}$ over $\mathbf Q_\ell$, which is $[K:\mathbf Q_\ell]\dim_{\mathbf Q_\ell}V$ by (i) and (a).
\item[(c)] \emph{The whole group.} By Tate's local Euler characteristic formula and local duality \cite[Chapter~VII, \S\S2--3]{NSW08}, $\dim H^1(K,V)=\dim V^{G_K}+\dim H^2(K,V)+[K:\mathbf Q_\ell]\dim V$, and $H^2(K,V)$ is dual to the space of the $f$ of (iii), which is $0$. With (a), $\dim H^1(K,V)=[K:\mathbf Q_\ell]\dim V$.
\item[(d)] \emph{Conclusion.} By (b) and (c), $H^1_f(K,V)=H^1(K,V)$. So for every continuous cocycle $c$ there is a crystalline cocycle $c'$ with $c-c'$ a coboundary $g\mapsto gv-v$; as $gv-v$ is $g(v\otimes1)-v\otimes1$ after tensoring with $B_{\mathrm{cris}}$, $c$ is crystalline.
\end{enumerate}
\end{proof}
% prF (PrF.lean, 23 lines): posBasic (PosBasic.lean: key, the descending induction with the private step and charFil, and H0_eq_bot_of_filD_zero_eq_bot) for (a); posBasic's dimension of finiteCohomology, from the field DimFormula of GlobalTheory.Published.BlochKatoLocal (LocalPublishedTheorems, field bk) through CocycleTransportH1fSurjective, for (b); tateRat (TateRat.lean, from the fields LocalEulerChar and LocalTateDuality) for (c); htop and z1fMap for (d). The Lean proves charFil (CharFil.lean, 193 lines), the text cites it.

\begin{lemma}\label{sha-siegel-step-ideal-seeds-inputs-tangent-crystalline-triangular}
Let $K$, $G_K$, $I_K$, $\chi$ and the crystalline representations be as in Lemma~\ref{sha-automorphy-unramified-tangents-unit-crystalline}. Let $E$ be a finite extension of $\mathbf Q_\ell$, $A$ a finite-dimensional commutative $E$-algebra with a homomorphism of $E$-algebras $\pi:A\to E$ whose kernel consists of nilpotent elements, $n\ge0$, and $M:G_K\to GL_n(A)$ a continuous homomorphism such that every $M(g)$ is upper triangular with diagonal entries $\kappa_1(g),\dots,\kappa_n(g)$, for characters $\kappa_i:G_K\to A^\times$. Let $w_1>\dots>w_n$ be integers with $\kappa_i(g)=\chi(g)^{w_i}$ for $g\in I_K$, and suppose that for all $i<j$ some $g\in G_K$ has $\pi(\kappa_i(g))\ne\pi(\kappa_j(g))\chi(g)$. Then $A^n$, on which $g$ acts by $M(g)$, is crystalline.
\end{lemma}
\leanref{Development/AutomorphicArithmetic/AutomorphicVarieties/PadicHodgeTheory/SnocZeroLastCoordinateLast.lean}{UpperTheory.V.Loc.prC}
\begin{proof}
By induction on $n$. For $n=0$ the space is $0$, which is crystalline by Lemma~\ref{sha-automorphy-unramified-tangents-unit-crystalline}(6). Suppose the lemma holds for $n$, and let $M:G_K\to GL_{n+1}(A)$ be as in the statement.
\begin{enumerate}
\item[(a)] \emph{The block.} Let $M_0(g)$ be the upper left $n\times n$ block of $M(g)$. As $M(g)$ is upper triangular, $M_0$ is a continuous homomorphism, with diagonal characters $\kappa_1,\dots,\kappa_n$; by induction $A^n$, with $g$ acting by $M_0(g)$, is crystalline. Let $k=\kappa_{n+1}$. As $k=\chi^{w_{n+1}}$ and $k^{-1}=\chi^{-w_{n+1}}$ on $I_K$, $A(k)$ and $A(k^{-1})$ are crystalline by Lemma~\ref{sha-automorphy-unramified-tangents-unit-crystalline}(6). Let $V$ be $A^n$ with $g$ acting by $k(g)^{-1}M_0(g)$. By step (d) of the proof of Lemma~\ref{sha-automorphy-unramified-tangents-unit-crystalline}, applied to $(a,x)\mapsto ax$ on $A(k^{-1})\times A^n$, $V$ is crystalline, hence de Rham \cite[Proposition~9.1.9]{BC09}.
\item[(b)] \emph{A filtration.} Let $\psi_i=\pi\circ(k^{-1}\kappa_i):G_K\to E^\times$ for $i\le n$. Let $\mathfrak m=\ker\pi$, and $N$ with $\mathfrak m^N=0$, which exists as $A$ is noetherian and $\mathfrak m$ consists of nilpotent elements. For $1\le i\le n$ and $0\le t\le N$ let $W_{i,t}$ be the set of $x\in A^n$ with $x_j=0$ for $j>i$ and $x_i\in\mathfrak m^t$, an $E$-subspace; then $W_{i,N}=W_{i-1,0}$ (with $W_{0,0}=0$) and $W_{n,0}=A^n$. For $x\in W_{i,t}$ and $g\in G_K$, $gx-\psi_i(g)x\in W_{i,t+1}$: the entries $j>i$ of $gx$ vanish, as $M_0(g)$ is upper triangular, and its entry $i$ is $k(g)^{-1}\kappa_i(g)x_i$, which differs from $\psi_i(g)x_i$ by an element of $\mathfrak m\cdot\mathfrak m^t$. Adding the elements of an $E$-basis of each $W_{i,t}$ modulo $W_{i,t+1}$ one at a time gives a filtration of $V$ as in (ii) of Lemma~\ref{sha-siegel-step-ideal-seeds-inputs-tangent-crystalline-positive}, with characters among the $\psi_i$. On $I_K$, $\psi_i=\chi^{w_i-w_{n+1}}$, and $w_i-w_{n+1}\ge1$.
\item[(c)] \emph{No quotient $\chi$.} Let $f:V\to E$ be $E$-linear with $f(gx)=\chi(g)f(x)$. By induction along the filtration of (b), $f$ vanishes on it: at a step adding $x$, with $gx-\psi_i(g)x$ in the previous subspace, $\chi(g)f(x)=f(gx)=\psi_i(g)f(x)$, and some $g$ has $\psi_i(g)\ne\chi(g)$, by the hypothesis for $i<n+1$; so $f(x)=0$. So $f=0$.
\item[(d)] \emph{The last column.} Let $c(g)=k(g)^{-1}\big(M(g)_{i,n+1}\big)_{i\le n}\in A^n$. Comparing the last columns of $M(gh)=M(g)M(h)$ shows that $c$ is a continuous cocycle of $V$. By (a) to (c), Lemma~\ref{sha-siegel-step-ideal-seeds-inputs-tangent-crystalline-positive} applies to $V$, and $c$ is a crystalline cocycle.
\item[(e)] \emph{Conclusion.} Lemma~\ref{sha-siegel-step-ideal-seeds-inputs-tangent-crystalline-extension} applies to $A$, the module $A^n$ with $M_0$, the module $A^{n+1}$ with $M$, the character $k$, the inclusion $f$ of the first $n$ coordinates, the last coordinate $p$ and $s=e_{n+1}$: indeed $M(g)e_{n+1}=k(g)\big(e_{n+1}+f(c(g))\big)$, and $A^n(k^{-1})=V$. By (d), $A^{n+1}$ with $M$ is crystalline.
\end{enumerate}
\end{proof}
% prC (SnocZeroLastCoordinateLast.lean, 352 lines): the private key, by induction on n, with crUnr for n = 0. (a) is topBlock, the induction hypothesis, crUnr for chRep k and chRep k^-1, and admissible_of_bilinear_openSubgroup_galois for the twist; isCrys_isDR_restrict_openSubgroup for de Rham. (b) is the private filt_tri (with ch_step and ch_step2) and hpos, with hπc for π(χ) = χ. (c) is the private h0dual_bot. (d) is hc and prF (Lemma sha-siegel-step-ideal-seeds-inputs-tangent-crystalline-positive). (e) is crExt (Lemma sha-siegel-step-ideal-seeds-inputs-tangent-crystalline-extension) with ext0, the last coordinate and Pi.single (Fin.last n) 1.

\begin{lemma}\label{sha-siegel-step-ideal-seeds-inputs-tangent-crystalline-deformation}
Let $K$, $G_K$ and the crystalline representations and cocycles be as in Lemma~\ref{sha-automorphy-unramified-tangents-unit-crystalline}, and let $E$ be a finite extension of $\mathbf Q_\ell$, $E[\epsilon]=E[\epsilon]/(\epsilon^2)$, $n\ge0$, $r:G_K\to GL_n(E)$ a continuous homomorphism, and $y:G_K\to M_n(E)$ a continuous cocycle of $\operatorname{ad}r$, the space $M_n(E)$ on which $g$ acts by $X\mapsto r(g)Xr(g)^{-1}$. Let $r'=(1+\epsilon y)r:G_K\to GL_n(E[\epsilon])$. Then $E[\epsilon]^n$, with $g$ acting by $r'(g)$, is crystalline if and only if $E^n$, with $g$ acting by $r(g)$, is crystalline and $y$ is a crystalline cocycle of $\operatorname{ad}r$.
\end{lemma}
\leanref{Development/AutomorphicArithmetic/AutomorphicVarieties/PadicHodgeTheory/DeformCr.lean}{UpperTheory.V.Loc.deformCr}
\begin{proof}
Let $W=M_n(E)\times E$, with $g(X,a)=\big(r(g)Xr(g)^{-1}+ay(g),a\big)$, and $\beta:W\times E^n\to E[\epsilon]^n$, $\beta((X,a),v)=av+\epsilon Xv$. Then $\beta(gw,r(g)v)=r'(g)\beta(w,v)$, as both are $ar(g)v+\epsilon\big(r(g)Xv+ay(g)r(g)v\big)$.
\begin{enumerate}
\item[(a)] Suppose that $E^n$ is crystalline. Then so is $\operatorname{ad}r$, by step (e) of the proof of Lemma~\ref{sha-automorphy-unramified-tangents-unit-crystalline}, as it maps injectively and equivariantly to $\operatorname{Hom}_{\mathbf Q_\ell}(E^n,E^n)$; and $E$ with the trivial action is crystalline by Lemma~\ref{sha-automorphy-unramified-tangents-unit-crystalline}(6). Lemma~\ref{sha-siegel-step-ideal-seeds-inputs-tangent-crystalline-extension}, with $A=E$, $\kappa=1$, $V=\operatorname{ad}r$, $V'=W$, $f(X)=(X,0)$, $p(X,a)=a$ and $s=(0,1)$, shows that $W$ is crystalline if and only if $y$ is a crystalline cocycle.
\item[(b)] Suppose that $E[\epsilon]^n$ is crystalline. Its quotient $E^n$, by reduction modulo $\epsilon$, is crystalline, by Lemma~\ref{sha-automorphy-unramified-tangents-unit-crystalline}(2). If $n=0$, $y=0$. Otherwise $w\mapsto\beta(w,\cdot)$ is an injective equivariant map from $W$ to $\operatorname{Hom}_{\mathbf Q_\ell}(E^n,E[\epsilon]^n)$, as $\beta((X,a),v)$ determines $a$ and $Xv$; so $W$ is crystalline by step (e) of the proof of Lemma~\ref{sha-automorphy-unramified-tangents-unit-crystalline}, and $y$ is a crystalline cocycle by (a).
\item[(c)] Conversely, suppose that $E^n$ is crystalline and $y$ is a crystalline cocycle. By (a), $W$ is crystalline. The image of $\beta$ spans $E[\epsilon]^n$, as $u+\epsilon u'=\beta((0,1),u)+\beta((1,0),u')$; so $E[\epsilon]^n$ is crystalline by step (d) of the proof of Lemma~\ref{sha-automorphy-unramified-tangents-unit-crystalline}.
\end{enumerate}
\end{proof}
% deformCr (DeformCr.lean, 177 lines), stated in the Lean for an open normal subgroup H, here H = G_K: extRep for W, pair for beta with hb; crM (isCrys_isDR_internalHom), cr1 (crUnr with w = 0) and key (crExt with kappa = 1) for (a); isCrys_isDR_sub_quot with the first component, and isCrys_isDR_internalHom with pair, for (b); admissible_of_bilinear_openSubgroup_galois with pair for (c).

\begin{lemma}\label{sha-siegel-step-ideal-seeds-inputs-tangent-crystalline}
Let $\ell$ be a prime, $H$ a compact group with an open embedding $\phi$ into $\operatorname{Gal}(\bar{\mathbf Q}_\ell/\mathbf Q_\ell)$, and $A=\bar{\mathbf Q}_\ell[\epsilon]/(\epsilon^2)$. Let $\rho:H\to GL_n(\bar{\mathbf Q}_\ell)$ be continuous, $y:H\to M_n(\bar{\mathbf Q}_\ell)$ a continuous cocycle for $g\mapsto\rho(g)(\cdot)\rho(g)^{-1}$, and $\rho'=(1+\epsilon y)\rho:H\to GL_n(A)$. Let $\kappa_1,\dots,\kappa_n:H\to A^\times$ be characters such that $\rho'$ is conjugate over $A$ to an upper triangular representation with diagonal characters $\kappa_1,\dots,\kappa_n$, let $m_1>\dots>m_n$ be integers with $m_j+2\le m_i$ for $i<j$, and let $U\subseteq H$ be an open subgroup with $\kappa_i(g)=\chi_{\mathrm{cyc}}(\phi(g))^{m_i}$ for every $i$ and every $g\in U$ whose image lies in the inertia group. Then $y$ satisfies the Bloch-Kato condition \cite{BK90}: its class in $H^1(H,\operatorname{ad}\rho)$, for $H$ viewed through $\phi$ as an open subgroup of $\operatorname{Gal}(\bar{\mathbf Q}_\ell/\mathbf Q_\ell)$, becomes a coboundary over $B_{\mathrm{cris}}$.
\end{lemma}
\leanref{Development/AutomorphicArithmetic/AutomorphicVarieties/PadicHodgeTheory/DualInclUnit.lean}{UpperTheory.V.Loc.pr6Bar}
\begin{proof}
Let $I\subseteq H$ be the inverse image of the inertia group, and $Q\in GL_n(A)$ with $Q\rho'(g)Q^{-1}$ upper triangular with diagonal entries $\kappa_1(g),\dots,\kappa_n(g)$ for $g\in H$. By infinite Galois theory, $\phi(H)$ and $\phi(U)$, open subgroups of $\operatorname{Gal}(\bar{\mathbf Q}_\ell/\mathbf Q_\ell)$, are $G_K$ and $G_{K_1}$ for finite extensions $K\subseteq K_1$ of $\mathbf Q_\ell$; identify $H$ with $G_K$ and $U$ with $G_{K_1}$ through $\phi$.
\begin{enumerate}
\item[(a)] \emph{Separation.} For $i<j$ some $g\in U$ has $\kappa_i(g)\not\equiv\kappa_j(g)\chi_{\mathrm{cyc}}(g)$ modulo $\epsilon$. Indeed, $\chi_{\mathrm{cyc}}$ maps the inertia group of $\operatorname{Gal}(\bar{\mathbf Q}_\ell/\mathbf Q_\ell)$ onto $\mathbf Z_\ell^\times$: for every $j$, $\mathbf Q_\ell(\mu_{\ell^j})$ is totally ramified over $\mathbf Q_\ell$ with group $(\mathbf Z/\ell^j)^\times$ \cite[Chapter~II, Proposition~7.13]{Neu99}, so the image of the inertia group, which is compact, is a closed subgroup of $\mathbf Z_\ell^\times$ that maps onto every $(\mathbf Z/\ell^j)^\times$; let $\sigma$ be an element of it with $\chi_{\mathrm{cyc}}(\sigma)=1+\ell$, which has infinite order. As $G_{K_1}$ has finite index, $\sigma^e\in G_{K_1}$ for some $e\ge1$; let $g\in U\cap I$ with $\phi(g)=\sigma^e$, so that $\chi_{\mathrm{cyc}}(g)$ has infinite order. If $\kappa_i(g)\equiv\kappa_j(g)\chi_{\mathrm{cyc}}(g)$, then $\chi_{\mathrm{cyc}}(g)^{m_i}=\chi_{\mathrm{cyc}}(g)^{m_j+1}$, while $m_i-m_j-1\ge1$; a contradiction.
\item[(b)] \emph{A finite extension.} The components of the entries of $\rho'$, those of $\rho$ and of $y\rho$, are continuous on the compact group $H$. By the Baire category theorem, as in step (a) of the proof of Lemma~\ref{sha-siegel-step-ideal-seeds-inputs-tangent-v0-local-model}, there is a finite extension $L$ of $\mathbf Q_\ell$ in $\bar{\mathbf Q}_\ell$ that contains the components of the entries of every $\rho'(g)$ and of $Q$. Then $\rho$ and $y$ take values in $GL_n(L)$ and $M_n(L)$, as the coefficient of $\epsilon$ in $\rho'(g)$ is $y(g)\rho(g)$, and the $\kappa_i$, the diagonal entries of $Q\rho'Q^{-1}$, take values in $L[\epsilon]^\times$.
\item[(c)] \emph{On $U$.} Lemma~\ref{sha-siegel-step-ideal-seeds-inputs-tangent-crystalline-triangular}, applied to $K_1$, $E=L$, $A=L[\epsilon]$, the reduction $\pi$ modulo $\epsilon$, $M(g)=Q\rho'(g)Q^{-1}$ for $g\in U$, the $\kappa_i$ and $w_i=m_i$, which are strictly decreasing, shows by (a) that $L[\epsilon]^n$ with $Q\rho'Q^{-1}$ is crystalline over $K_1$; so is $L[\epsilon]^n$ with $\rho'$, by Lemma~\ref{sha-automorphy-unramified-tangents-unit-crystalline}(2) applied to $x\mapsto Qx$. By Lemma~\ref{sha-siegel-step-ideal-seeds-inputs-tangent-crystalline-deformation}, applied to $K_1$, $E=L$ and the restrictions of $\rho$ and $y$ to $U$, the restriction of $y$ to $U$ is a crystalline cocycle of $\operatorname{ad}\rho$ over $L$, for $K_1$.
\item[(d)] \emph{Descent.} The normal core of $U$ in $H$, the intersection of its finitely many conjugates, is $G_{K_2}$ for a finite Galois extension $K_2$ of $K$ containing $K_1$. The restriction of $y$ to $G_{K_2}$ is crystalline, with the element that serves for $U$. By Lemma~\ref{sha-siegel-step-ideal-seeds-inputs-tangent-crystalline-descent}, $y$ is a crystalline cocycle of $\operatorname{ad}\rho$ over $L$, for $K$.
\item[(e)] \emph{Coefficients.} The inclusion $M_n(L)\subseteq M_n(\bar{\mathbf Q}_\ell)$ is $H$-equivariant. It carries an element $x$ of $M_n(L)\otimes_{\mathbf Q_\ell}B_{\mathrm{cris}}$ with $y(g)\otimes1=gx-x$ to an element of $M_n(\bar{\mathbf Q}_\ell)\otimes_{\mathbf Q_\ell}B_{\mathrm{cris}}$ with the same property. So $y$ satisfies the Bloch-Kato condition.
\end{enumerate}
\end{proof}
% Loc.pr6Bar (PadicHodgeTheory/DualInclUnit.lean, 291 lines, from Loc.LocalPublishedTheorems): hne with cycInertia (CycInertia.lean, 173 lines: ci_surj, ci_inf and a power in U) for (a); the private desc (the Baire category argument for the two components, with gl_range_iff) and s, z, κ₁ for (b); the private fin for (c) and (d): Loc.pr6 (CrystallineTransferSurjection.lean, 97 lines) on U with H = top, which is prC (Lemma sha-siegel-step-ideal-seeds-inputs-tangent-crystalline-triangular) with A = E[eps] and π = fst, crystalline_conj_invariant, and deformCr (Lemma sha-siegel-step-ideal-seeds-inputs-tangent-crystalline-deformation), then z1fDesc (Lemma sha-siegel-step-ideal-seeds-inputs-tangent-crystalline-descent) for the normal core of U; z1fCoef for (e). The Lean's hypothesis is, for i < j, m_j + 2 <= m_i or some g in U with fst(kappa_i(g)) != fst(kappa_j(g)) chi(g); the text uses the first alternative, which is the one the Lean uses here, and (a) turns it into the second. The published inputs are the fields DimFormula (Bloch-Kato), LocalEulerChar and LocalTateDuality, through Lemma sha-siegel-step-ideal-seeds-inputs-tangent-crystalline-positive; the facts on crystalline representations are those of Lemma sha-automorphy-unramified-tangents-unit-crystalline.

\begin{lemma}\label{sha-siegel-step-ideal-seeds-inputs-tangent-bloch-kato}
In Situation~\ref{sha-hecke-situation}, suppose that the $\bar\chi_i$ are trivial on $I_{v_0}$ and take equal values on $D_{v_0}$, and that $\varepsilon\bmod\ell$ is nontrivial on $D_{v_0}$. Let $x$, $\rho$, $x'$, and $\rho'$, $y$, $P$, $R$ and $z$ be as in Lemma~\ref{sha-siegel-step-ideal-seeds-inputs-tangent-lift}. Suppose that $x'$ maps the image of $\Lambda$ in $B$ into $\bar{\mathbf Q}_\ell$, and that for every $v\mid\ell$ the restriction of $\rho'$ to $D_v$ is conjugate over $\bar{\mathbf Q}_\ell[\epsilon]/(\epsilon^2)$ to an upper triangular representation with diagonal characters $x'\circ\psi^{\mathrm u}_{v,1},\dots,x'\circ\psi^{\mathrm u}_{v,n}$, and there are integers $m_{v,1}>\dots>m_{v,n}$ with $m_{v,j}+2\le m_{v,i}$ for $i<j$ and $k\ge1$ with $\big(\psi_{x,v,i}\chi_{\mathrm{cyc}}^{-m_{v,i}}\big)^k=1$ on $I_v$. Then the inflation of $z$ to $\operatorname{Gal}(\bar L/L^+)$ satisfies the Bloch-Kato conditions.
\end{lemma}
\leanref{Development/AutomorphicArithmetic/AutomorphicVarieties/WeakShimuraTheory/EigenbasisLifts/BkAbEmb.lean}{UpperTheory.V.Ws.Eb.bkAb}
\leanref{Development/AutomorphicArithmetic/AutomorphicVarieties/WeakShimuraTheory/EigenbasisLifts/AwAvg.lean}{UpperTheory.V.Ws.Eb.bkAw}
\leanref{Development/AutomorphicArithmetic/AutomorphicVarieties/WeakShimuraTheory/EigenbasisLifts/LocLOpen.lean}{UpperTheory.V.Ws.Eb.locL}
\begin{proof}
\begin{enumerate}
\item[(a)] \emph{Places not above $\ell$.} Let $w$ be a place of $L^+$ not above $\ell$. It suffices to show that $z$ restricted to the image in $\Gamma$ of the inertia group at one place above $w$ is a coboundary $g\mapsto R(g)X-X$. If $w\notin S$, that image is trivial, and $z(1)=0$. If $w=u\in S_{\mathrm a}$, the image is $I_u$, and $z$ vanishes on its subgroup $I_u\cap\bigcap_i\ker\bar\chi_i$, of finite index, by Lemma~\ref{sha-siegel-step-ideal-seeds-inputs-tangent-lift}(2), as $\rho'=1$ there forces $y=0$. A cocycle that vanishes on a normal subgroup $N'$ of finite index $d$ of a group $I'$, with values in a $\mathbf Q$-vector space, is a coboundary on $I'$: it is constant on the cosets of $N'$, and with $\Sigma$ the sum of its values on a set of representatives, $z(g)=R(g)X-X$ with $X=-d^{-1}\Sigma$, as in step (7) of the proof of Lemma~\ref{sha-siegel-automorphy-descent}. If $w=v_0$, $z$ vanishes on $I_{v_0}$ by Lemma~\ref{sha-siegel-step-ideal-seeds-inputs-tangent-v0}.
\item[(b)] \emph{Places above $\ell$.} Let $v\mid\ell$ and $H=D_v$, with its open embedding into $\operatorname{Gal}(\bar{\mathbf Q}_\ell/\mathbf Q_\ell)$, so that $I_v$ is the inverse image of inertia. Let $\kappa_i=x'\circ\psi^{\mathrm u}_{v,i}$. On $I_v$, $\psi^{\mathrm u}_{v,i}$ is the image of the universal character over $\Lambda$, so $\kappa_i$ has no $\epsilon$-part there, and $\kappa_i=\psi_{x,v,i}$ on $I_v$. The finitely many continuous functions $g\mapsto\psi_{x,v,i}(g)\chi_{\mathrm{cyc}}(g)^{-m_{v,i}}$ take the value $1$ at $1$, and the $k$-th roots of unity other than $1$ are finitely many; so there is an open subgroup $U$ of $D_v$ on which $\big(\psi_{x,v,i}\chi_{\mathrm{cyc}}^{-m_{v,i}}\big)^k=1$ implies $\psi_{x,v,i}=\chi_{\mathrm{cyc}}^{m_{v,i}}$. Then $\kappa_i=\chi_{\mathrm{cyc}}^{m_{v,i}}$ on $U\cap I_v$, and Lemma~\ref{sha-siegel-step-ideal-seeds-inputs-tangent-crystalline}, applied to $\rho$, $y$ and $\rho'$ restricted to $D_v$, shows that $z$ satisfies the Bloch-Kato condition at $v$. The places of $\operatorname{Gal}(\bar L/L^+)$ above $\ell$ are conjugate to these, so the condition holds at every place above $\ell$.
\end{enumerate}
\end{proof}
% bkAb (BkAbEmb.lean, 48 lines: the private bkAb_emb for the places above l, and Loc.bkAbove) from the field h1, which is the private at_l of EigenbasisRankFour.lean (with Lambda.toPair_spec for the vanishing epsilon-part of kappa_i on inertia) and locL (LocLOpen.lean, 58 lines: the private locL_open for U, then Loc.pr6Bar, Lemma sha-siegel-step-ideal-seeds-inputs-tangent-crystalline), and the field h2, which is bkAw (AwAvg.lean, 99 lines: the three cases w not in S, w in S_a with the private exists_coboundary_on_h_of_vanish_on_finiteIndex, and w = v_0 with cocycle_vanishes_on_inertia_of_point_pseudodeformation).

\begin{lemma}\label{sha-siegel-step-ideal-seeds-inputs-tangent-constant-separating}
Let $n\ge1$, let $C$ be a local ring whose residue field $k$ has at least $n$ elements, $H$ a group, and $\psi_1,\dots,\psi_n:H\to C^\times$ characters whose reductions $\bar\psi_i:H\to k^\times$ are pairwise distinct. Extend each $\psi_i$ linearly to a map $C[H]\to C$. Then there is $z\in C[H]$ with $\psi_i(z)-\psi_j(z)\in C^\times$ for $i\ne j$.
\end{lemma}
\leanref{Development/AutomorphicArithmetic/AutomorphicVarieties/AlgebraTools/SepChars.lean}{UpperTheory.V.AlgebraTools.sep_chars}
\begin{proof}
Extend each $\bar\psi_i$ linearly to a map $k[H]\to k$.
\begin{enumerate}
\item[(a)] \emph{Surjectivity.} The $k$-linear map $L:k[H]\to k^n$, $w\mapsto(\bar\psi_1(w),\dots,\bar\psi_n(w))$, is surjective. Otherwise a nonzero linear form $(c_i)\mapsto\sum_ia_ic_i$ on $k^n$ vanishes on the image of $L$, which contains $(\bar\psi_i(h))_i$ for every $h\in H$; so $\sum_ia_i\bar\psi_i(h)=0$ for every $h\in H$, with $(a_i)\ne0$. This contradicts the linear independence of the pairwise distinct characters $\bar\psi_i$ of $H$ with values in the field $k$.
\item[(b)] \emph{The element.} As $k$ has at least $n$ elements, there are pairwise distinct $b_1,\dots,b_n\in k$. By (a) some $w=\sum_hw_hh\in k[H]$ has $\bar\psi_i(w)=b_i$ for every $i$. Let $z=\sum_hz_hh$, where $h$ runs over the finitely many elements with $w_h\ne0$ and $z_h\in C$ lifts $w_h$. Then $\psi_i(z)$ reduces to $\bar\psi_i(w)=b_i$. So for $i\ne j$ the reduction of $\psi_i(z)-\psi_j(z)$ is $b_i-b_j\ne0$, and $\psi_i(z)-\psi_j(z)$ is a unit of the local ring $C$.
\end{enumerate}
\end{proof}
% sep_chars (SepChars.lean, 51 lines): hsurj for (a), through Submodule.exists_le_ker_of_lt_top and linearIndependent_monoidHom; for (b), the embedding a : Fin n -> k from the cardinality hypothesis, w from hsurj, z through Function.surjInv of the residue map on the support of w, and residue_ne_zero_iff_isUnit.

\begin{lemma}\label{sha-siegel-step-ideal-seeds-inputs-tangent-constant-flags}
Let $n\ge1$, let $C$ be a local ring whose residue field $k$ has at least $n$ elements, $H$ a group and $r:H\to GL_n(C)$ a homomorphism. Let $\psi_1,\dots,\psi_n$ and $\psi'_1,\dots,\psi'_n$ be characters $H\to C^\times$ such that the reductions $\bar\psi_i:H\to k^\times$ are pairwise distinct and $\bar\psi'_i=\bar\psi_i$ for every $i$. Let $g,g'\in GL_n(C)$ be such that for every $s\in H$, $gr(s)g^{-1}$ is upper triangular with diagonal entries $\psi_1(s),\dots,\psi_n(s)$, and $g'r(s)g'^{-1}$ is upper triangular with diagonal entries $\psi'_1(s),\dots,\psi'_n(s)$. Then $\psi_i=\psi'_i$ for every $i$, and $gg'^{-1}$ is upper triangular.
\end{lemma}
\leanref{Development/AutomorphicArithmetic/AutomorphicVarieties/AlgebraTools/FlagCharsUnique.lean}{UpperTheory.V.AlgebraTools.flag_chars_unique}
\begin{proof}
Let $V=gg'^{-1}$. For $w\in C[H]$ let $T(w)$ and $T'(w)$ be the images of $w$ under the linear extensions of $s\mapsto gr(s)g^{-1}$ and $s\mapsto g'r(s)g'^{-1}$, and extend the $\psi_i$ and $\psi'_i$ linearly to $C[H]$.
\begin{enumerate}
\item[(a)] \emph{The relation.} For $s\in H$, $gr(s)g^{-1}V=gr(s)g'^{-1}=V\,g'r(s)g'^{-1}$. By linearity, $T(w)V=VT'(w)$ for $w\in C[H]$. The matrices $T(w)$ and $T'(w)$ are upper triangular, with diagonal entries $\psi_i(w)$ and $\psi'_i(w)$, and $\psi'_i(w)$ has the same reduction as $\psi_i(w)$, as $\bar\psi'_i=\bar\psi_i$.
\item[(b)] \emph{Units.} By Lemma~\ref{sha-siegel-step-ideal-seeds-inputs-tangent-constant-separating}, some $z\in C[H]$ has $\psi_i(z)-\psi_j(z)\in C^\times$ for $i\ne j$. Let $P=T(z)$ and $P'=T'(z)$. For $j<i$, $P_{ii}-P'_{jj}$ has the same reduction as $\psi_i(z)-\psi_j(z)$, so it is a unit.
\item[(c)] \emph{Triangularity.} Let $j<i$, and suppose, by induction on $j$ and, for fixed $j$, by descending induction on $i$, that $V_{kj}=0$ for $k>i$ and $V_{ik}=0$ for $k<j$. The $(i,j)$ entry of $PV=VP'$ is $\sum_kP_{ik}V_{kj}=\sum_kV_{ik}P'_{kj}$. On the left, $P_{ik}=0$ for $k<i$ and $V_{kj}=0$ for $k>i$; on the right, $P'_{kj}=0$ for $k>j$ and $V_{ik}=0$ for $k<j$. So $(P_{ii}-P'_{jj})V_{ij}=0$, and $V_{ij}=0$ by (b). So $V$ is upper triangular.
\item[(d)] \emph{The characters.} By (c), the determinant of the invertible matrix $V$ is the product of its diagonal entries, so each $V_{ii}$ is a unit. Let $s\in H$. The $(i,i)$ entry of $gr(s)g^{-1}V=V\,g'r(s)g'^{-1}$, a relation between products of upper triangular matrices, is $\psi_i(s)V_{ii}=V_{ii}\psi'_i(s)$. So $\psi_i(s)=\psi'_i(s)$.
\end{enumerate}
\end{proof}
% flag_chars_unique (FlagCharsUnique.lean, 103 lines): sep_chars for z; hcomm for the relation at s and key (induction on the group algebra) for (a); the private exists_invariant_flag_mod_l (the double well-founded induction) for (c); Matrix.det_of_upperTriangular, the private diag_mul and the cancellation of the unit V_ii for (d).

\begin{lemma}\label{sha-siegel-step-ideal-seeds-inputs-tangent-constant}
In Situation~\ref{sha-hecke-situation}, let $x$, $\rho$, $x'$, $\rho'$ and $y$ be as in Lemma~\ref{sha-siegel-step-ideal-seeds-inputs-tangent-lift}(1), with $A=\bar{\mathbf Q}_\ell[\epsilon]/(\epsilon^2)$. Suppose that there is $X\in M_n(\bar{\mathbf Q}_\ell)$ with $y(g)=\rho(g)X\rho(g)^{-1}-X$ for $g\in\Gamma_0$, and that for every $v\mid\ell$ the restriction of $\rho$ to $D_v$ is conjugate to an upper triangular representation with pairwise distinct diagonal characters $\psi_{x,v,1},\dots,\psi_{x,v,n}$, and that of $\rho'$ is conjugate over $A$ to an upper triangular representation with diagonal characters $x'\circ\psi^{\mathrm u}_{v,1},\dots,x'\circ\psi^{\mathrm u}_{v,n}$. Then $x'=x$.
\end{lemma}
\leanref{Development/AutomorphicArithmetic/AutomorphicVarieties/WeakShimuraTheory/EigenbasisLifts/ConjOf.lean}{UpperTheory.V.Ws.Eb.conjOf}
\leanref{Development/AutomorphicArithmetic/AutomorphicVarieties/WeakShimuraTheory/EigenbasisLifts/ConstOf.lean}{UpperTheory.V.Ws.Eb.constOf}
\leanref{Development/AutomorphicArithmetic/AutomorphicVarieties/WeakShimuraTheory/EigenbasisLifts/DerZero.lean}{UpperTheory.V.Ws.Eb.derZero}
\leanref{Development/AutomorphicArithmetic/AutomorphicVarieties/AlgebraTools/FlagCharsUnique.lean}{UpperTheory.V.AlgebraTools.flag_chars_unique}
\begin{proof}
Let $D$ be the coefficient of $\epsilon$ in $x'$, a derivation at $x$.
\begin{enumerate}
\item[(a)] \emph{Conjugation.} With $Q=1-\epsilon X\in GL_n(A)$, $Q\rho(g)Q^{-1}=\rho(g)+\epsilon(\rho(g)X-X\rho(g))=(1+\epsilon y(g))\rho(g)=\rho'(g)$.
\item[(b)] \emph{Coefficients.} By (a), $x'(D^{\mathrm u})$, the determinant of $\rho'$, is that of $\rho$ viewed over $A$; so the coefficients $a_i(g)$ of the characteristic polynomials of $D^{\mathrm u}$ (Definition~\ref{sha-siegel-step-ideal-seeds-inputs-twisting}) have $x'(a_i(g))\in\bar{\mathbf Q}_\ell$, that is, $D(a_i(g))=0$.
\item[(c)] \emph{Characters.} Let $v\mid\ell$. By (a), $Q^{-1}\rho'Q=\rho$ on $D_v$, so $\rho'$ is conjugate over $A$ to upper triangular representations with the diagonal characters $\psi_{x,v,i}$, with values in $\bar{\mathbf Q}_\ell$, and with the diagonal characters $x'\circ\psi^{\mathrm u}_{v,i}$; these reduce to the same characters $\psi_{x,v,i}$, which are pairwise distinct. Lemma~\ref{sha-siegel-step-ideal-seeds-inputs-tangent-constant-flags}, with $C=A$, whose residue field $\bar{\mathbf Q}_\ell$ is infinite, $H=D_v$, $r=\rho'$, $\psi_i=x'\circ\psi^{\mathrm u}_{v,i}$ and $\psi'_i=\psi_{x,v,i}$, gives $x'(\psi^{\mathrm u}_{v,i}(s))=\psi_{x,v,i}(s)$, so $D(\psi^{\mathrm u}_{v,i}(s))=0$, for $s\in D_v$.
\item[(d)] \emph{Generators.} Every element of $B$ is congruent, modulo every power of $\mathfrak m_B$, to an element of the $\mathcal O$-subalgebra generated by the $a_i(g)$ and the $\psi^{\mathrm u}_{v,i}(s)$. The derivation $D$ vanishes on that subalgebra, by (b) and (c), and by Lemma~\ref{sha-automorphy-unramified-near-tangent-small} it is arbitrarily small on high powers of $\mathfrak m_B$; so $D=0$ and $x'=x$.
\end{enumerate}
\end{proof}
% conjOf (ConjOf.lean, 56 lines) for (a), with Q = dmk 1 (-X); constOf (ConstOf.lean, 70 lines) for (b) to (d): repApi for (b), the private flag_const (the flag g Q^-1 of rho' through the identity e, res_iff for the residues of the dual numbers, and the infinite residue field for the cardinality) with AlgebraTools.flag_chars_unique (Lemma sha-siegel-step-ideal-seeds-inputs-tangent-constant-flags) for (c), and derZero (DerZero.lean) with CA.exists_adjoin_coeff_character_sub_mem_maximalIdeal_pow for (d).

\begin{lemma}\label{sha-siegel-step-ideal-seeds-inputs}
In Situation~\ref{sha-siegel-step-ideal-setting}:
\begin{enumerate}
\item for every $n\ge0$ and $\varphi_n$ as in Lemma~\ref{sha-siegel-step-ideal-seeds-phi}, the ideal of $B$ generated by $\varphi_n(J_0)$ is $J_\varphi$;
\item let $x$ be a point of $B^{\mathrm{ord}}$ with $J_0\subseteq\ker x$ and $\rho:G_L\to GL_4(\bar{\mathbf Q}_\ell)$ a continuous representation such that $x$ carries the characteristic polynomial of $D^{\mathrm u}$ at $\mathrm{pr}(g)$ to that of $\rho(g)$ for $g\in G_L$; suppose that $\rho$ is attached through $\iota$ to a regular algebraic cuspidal automorphic representation $\Pi$ of $GL_4(\mathbf A_L)$ with $\Pi^\vee\cong\Pi^c$ and is windowed, that it is trivial on the inertia group of $G_L$ at the place above $\tilde v_0$ whose image in $\Gamma_0$ is $I_{v_0}$, and that its restriction to $G_L\cap\ker\chi_{\mathrm{cyc}}$ is irreducible; and suppose that there are integers $m_{v,i}$ with $m_{v,j}+2\le m_{v,i}$ for $i<j$ and $k\ge1$ such that for every $v\mid\ell$ the restriction of $\rho$ to $D_v$ is conjugate to an upper triangular representation with diagonal characters $\psi_{x,v,1},\dots,\psi_{x,v,4}$ and $\big(\psi_{x,v,i}\chi_{\mathrm{cyc}}^{-m_{v,i}}\big)^k=1$ on $I_v$. Then every prime of $B^{\mathrm{ord}}$ contained in $\ker x$ contains $J_0$.
\end{enumerate}
\end{lemma}
\leanref{Development/AutomorphicArithmetic/AutomorphicVarieties/ModularityArgument/ThresholdArguments/DenseThresholdCases/HidaModularIdealMap.lean}{UpperTheory.V.Ws.hida_modular_ideal_map}
\leanref{Development/AutomorphicArithmetic/AutomorphicVarieties/ModularityArgument/ThresholdArguments/DenseThresholdCases/EigenbasisRankFour.lean}{UpperTheory.V.Th.eigenbasis_rank_four}
\begin{proof}
\emph{Proof of (1).} Let $\beta_0$ be the twisting character over $\mathcal O'$ of Lemma~\ref{sha-siegel-step-ideal-seeds-inputs-twist-descend}, and $a_i(\gamma)$, for $\gamma\in\Gamma_0$ and $1\le i\le4$, as in Definition~\ref{sha-siegel-step-ideal-seeds-inputs-twisting}. For a twisting character $\gamma$ over $\mathcal O'$ write $\gamma_B$ and $\gamma^{\mathrm{ord}}$ for its images in $B$ and in $B^{\mathrm{ord}}$. The map $\varphi$ is the map $\varphi_0$ of Lemma~\ref{sha-siegel-step-ideal-seeds-phi}.
\begin{enumerate}
\item[(a)] \emph{Twists of $\varphi$.} The map $\varphi_n$ is a twist of $\varphi$ by $(\beta_0^n)_B$, and $\varphi$ is a twist of $\varphi_n$ by $(\beta_0^{-n})_B$. Indeed, let $g\in G_L$ and $1\le i\le4$. By Lemma~\ref{sha-siegel-step-ideal-seeds-phi}, $\varphi_n$ and $\varphi$ carry $a_i(\mathrm{pr}(g))$ to the coefficients of $Y^{4-i}$ in $c^{\kappa_n}_g$ and in $c^\kappa_g$, which are $\kappa_n(g)^i$ and $\kappa(g)^i$ times that in $c_g$; as $\kappa_n(g)=\kappa(g)\beta_0(\mathrm{pr}(g))^n$, $\varphi_n(a_i(\mathrm{pr}(g)))=\beta_0(\mathrm{pr}(g))^{ni}\varphi(a_i(\mathrm{pr}(g)))$. Likewise, for $v\mid\ell$ and $s\in D_v$, $\varphi_n(\psi^{\mathrm u}_{v,i}(s))=\kappa_n(s)\psi^B_{\bar v,i}(s)=\beta_0(s)^n\varphi(\psi^{\mathrm u}_{v,i}(s))$. As $\mathrm{pr}$ is surjective, this is the first claim; the second follows on multiplying by $\beta_0(\mathrm{pr}(g))^{-ni}$ and $\beta_0(s)^{-n}$.
\item[(b)] \emph{Comparison.} Let $\gamma$ be a twisting character over $\mathcal O'$, and $f,f':B^{\mathrm{ord}}\to B$ homomorphisms of $\mathcal O'$-algebras such that $f'$ is a twist of $f$ by $\gamma_B$. Then the ideal generated by $f'(J_0)$ lies in the ideal generated by $f(J_0)$. Indeed, $\mathcal O'\to B^{\mathrm{ord}}$ is a local homomorphism, so $\gamma^{\mathrm{ord}}$ is a twisting character over $B^{\mathrm{ord}}$ by Lemma~\ref{sha-siegel-step-ideal-seeds-inputs-twist-map}(3); let $\tau:B^{\mathrm{ord}}\to B^{\mathrm{ord}}$ be a twist of the identity by it, which exists by Lemma~\ref{sha-siegel-step-ideal-seeds-inputs-twist-map}(2). As $J_0$ is a Hecke ideal, Lemma~\ref{sha-siegel-step-ideal-seeds-inputs-twist-stable} gives $\tau(J_0)\subseteq J_0$. By Lemma~\ref{sha-siegel-step-ideal-seeds-inputs-twist-map}(4), $f\circ\tau$ is a twist of $f$ by $f\circ\gamma^{\mathrm{ord}}=\gamma_B$. As $B$ is a complete local noetherian $\mathcal O'$-algebra with residue field $k'$, Lemma~\ref{sha-siegel-step-ideal-seeds-inputs-twist-map}(1) gives $f'=f\circ\tau$. So $f'(J_0)=f(\tau(J_0))\subseteq f(J_0)$.
\item[(c)] \emph{Conclusion.} By Lemma~\ref{sha-siegel-step-ideal-seeds-inputs-twist-map}(3), $\beta_0^n$ and $\beta_0^{-n}$ are twisting characters over $\mathcal O'$. By (a), (b) applies to $\beta_0^n$, $\varphi$ and $\varphi_n$, and to $\beta_0^{-n}$, $\varphi_n$ and $\varphi$. So the ideal generated by $\varphi_n(J_0)$ and $J_\varphi$ each contain the other.
\end{enumerate}

\emph{Proof of (2).} Let $A=\bar{\mathbf Q}_\ell[\epsilon]/(\epsilon^2)$. The residual datum of the link is trivial on $I_{v_0}$, takes equal values on $D_{v_0}$, and $\varepsilon\bmod\ell$ is nontrivial on $D_{v_0}$, by (3) of Definition~\ref{sha-siegel-automorphy-link}.
\begin{enumerate}
\item[(a)] \emph{Reduction to tangent vectors.} The Iwasawa algebra $\Lambda$ is a power series ring over $\mathcal O'$ in $N=4[L^+:\mathbf Q]$ variables (proof of Lemma~\ref{sha-hecke-algebra}), and by Proposition~\ref{sha-hecke-ideal}(1), $B^{\mathrm{ord}}/J_0$ is finite over $\Lambda$ and $\lambda b\in J_0$ with $\lambda\in\Lambda$ nonzero implies $b\in J_0$. By Lemma~\ref{sha-siegel-step-ideal-seeds-inputs-regular}, applied to $R=B^{\mathrm{ord}}$, the map $\Lambda\to B^{\mathrm{ord}}$, $J_0$ and $x$, it suffices to show: every homomorphism of $\mathcal O'$-algebras $x':B^{\mathrm{ord}}\to A$ that reduces to $x$ and maps the image of $\Lambda$ into $\bar{\mathbf Q}_\ell$ is equal to $x$. Fix such an $x'$.
\item[(b)] \emph{The representation of $\Gamma_0$.} As $\rho$ is irreducible on $G_L\cap\ker\chi_{\mathrm{cyc}}$, the $\rho(g)$ with $g\in G_L$ span $M_4(\bar{\mathbf Q}_\ell)$. Let $g\in\ker\mathrm{pr}$. For every $h\in G_L$, $\rho(hg)$ and $\rho(h)$ have the same characteristic polynomial, the image under $x$ of that of $D^{\mathrm u}$ at $\mathrm{pr}(h)$; so $\operatorname{tr}(\rho(h)(\rho(g)-1))=0$ for all $h$, and $\rho(g)=1$ by the nondegeneracy of the trace form. So $\rho=\rho_0\circ\mathrm{pr}$ for a homomorphism $\rho_0:\Gamma_0\to GL_4(\bar{\mathbf Q}_\ell)$ whose values span $M_4(\bar{\mathbf Q}_\ell)$, and $x(D^{\mathrm u})$ is the determinant of $\rho_0$, as a determinant is determined by the characteristic polynomials at the elements of the group \cite[(1.5) and Corollary~1.14]{Che14}.
\item[(c)] \emph{The cocycle.} Lemma~\ref{sha-siegel-step-ideal-seeds-inputs-tangent-lift}, applied to $x$, $\rho_0$ and $x'$, gives $\rho'$, $y$, $P$, $R$ and $z$ as there, and $\rho_0$ is continuous.
\item[(d)] \emph{The flags over $A$.} Let $v\mid\ell$. The characters $\psi_{x,v,i}$ are pairwise distinct: if $\psi_{x,v,i}=\psi_{x,v,j}$, then $\chi_{\mathrm{cyc}}^{(m_{v,j}-m_{v,i})k}=1$ on $I_v$, on which $\chi_{\mathrm{cyc}}$ has infinite order, so $m_{v,i}=m_{v,j}$ and $i=j$. By Lemma~\ref{sha-siegel-step-ideal-seeds-inputs-tangent-lift}(2), the ordered products $\prod_i\big(\rho'(g_i)-x'(\psi^{\mathrm u}_{v,i}(g_i))\big)$ with $g_i\in D_v$ vanish, and modulo $\epsilon$, $\rho'=\rho_0$ is conjugate on $D_v$ to an upper triangular representation with the distinct diagonal characters $\psi_{x,v,i}$, the reductions of the $x'\circ\psi^{\mathrm u}_{v,i}$. By Lemma~\ref{sha-siegel-step-ideal-seeds-inputs-tangent-constant-separating}, with $C=A$, whose residue field $\bar{\mathbf Q}_\ell$ is infinite, $H=D_v$ and the characters $x'\circ\psi^{\mathrm u}_{v,i}$, whose reductions $\psi_{x,v,i}$ are distinct, some element $w$ of $A[D_v]$ has images under the $x'\circ\psi^{\mathrm u}_{v,i}$ whose differences are units. As $\rho_0$ is conjugate on $D_v$ to an upper triangular representation with diagonal characters $\psi_{x,v,i}$, the characteristic polynomial of the image of $w$ under $\rho'$ reduces to the product of the $X$ minus the images of $w$ under the $\psi_{x,v,i}$. The vanishing of the ordered products holds for $g_1,\dots,g_4$ in $A[D_v]$, by linearity in each $g_i$. So the proof of Lemma~\ref{sha-rigidity-lagrange}, with $C=A$, $r=\rho'$ on $D_v$, the characters $x'\circ\psi^{\mathrm u}_{v,i}$ and $w$ in place of $g_0$, shows that $\rho'$ is conjugate over $A$, on $D_v$, to an upper triangular representation with diagonal characters $x'\circ\psi^{\mathrm u}_{v,1},\dots,x'\circ\psi^{\mathrm u}_{v,4}$.
\item[(e)] \emph{The hypotheses of the adjoint lifting theorem.} The point $x$ and $\rho_0$ satisfy (1) to (6) of Theorem~\ref{sha-siegel-automorphy-adjoint}, with $n=4$, whose conditions at $v_0$ hold as noted above.
For (1), $x(D^{\mathrm u})$ is the determinant of $\rho_0$, by (b); the $\psi_{x,v,i}$ are continuous, as $x$ is (Lemma~\ref{sha-automorphy-unramified-near-tangent-points}), and $(\psi_{x,v,i}\chi_{\mathrm{cyc}}^{-m_{v,i}})^k=1$ on $I_v$; so, by local class field theory as for (1) in step (F) of the proof of Proposition~\ref{sha-siegel-points}, $x$ is arithmetic of the weight $h$ with $h_{v,i,\tau}=-m_{v,i}$, which is increasing.
(2) is the hypothesis on $\rho$, with $\Pi$. (3) is the hypothesis on $\rho$ at $\tilde v_0$ and on its restrictions to the $D_v$.
For (4), let $v\mid\ell$. For all $i,j$, $\psi_{x,v,i}\ne\psi_{x,v,j}\chi_{\mathrm{cyc}}$: otherwise, on $I_v$, $\chi_{\mathrm{cyc}}^{(1-m_{v,i}+m_{v,j})k}=1$, while $1-m_{v,i}+m_{v,j}\ne0$ by the gaps and $\chi_{\mathrm{cyc}}$ has infinite order on $I_v$. So the proof of Lemma~\ref{sha-rigidity-free}, with $F=L$, $a=0$, $r=\rho$, the integers $m_{v,i}$, and the characters $\psi_{x,v,i}$ at the place $A=A'$ at which $D_v$ is the decomposition group, gives $g\in D_v$ with $\varepsilon(g)=1$ and $\psi_{x,v,i}(g)^e\ne\psi_{x,v,j}(g)^e$ for $i\ne j$ and $e\ge1$. That proof uses the hypothesis only at the prime of $L$ below $A'$, and on the characters only through $(\psi_{x,v,i}\chi_{\mathrm{cyc}}^{-m_{v,i}})^k=1$ on the inertia group; and $(\Pi,\chi)$ is polarized for $\chi=\delta^4$, $\delta$ the quadratic character of $L/L^+$, by step (1) of the proof of Theorem~\ref{sha-siegel-automorphy-adjoint}.
For (5), Lemma~\ref{sha-siegel-automorphy-rank-four} applies as for (5) in step (F) of the proof of Proposition~\ref{sha-siegel-points}, with $\Gamma_0$, $\varepsilon$, the residue field $k'$ of $\mathcal O'$, $\rho_0$, which is irreducible as its values span $M_4(\bar{\mathbf Q}_\ell)$, the quotient $A_x$ of $B^{\mathrm{ord}}$ by $\ker x$ with its maps $j$ and $\phi$, the characters $\bar\chi_i$, the image in $\Gamma'$ of the element $g$ of (4), and the conditions (i) and (ii) there, which follow from (g) and (h) of Definition~\ref{sha-siegel-tower-towers}.
For (6), dividing the relations for $i$ and for $j$ gives $\big(\psi_{x,v,i}\psi_{x,v,j}^{-1}\varepsilon^{m_{v,j}-m_{v,i}}\big)^k=1$ on $I_v$.
\item[(f)] \emph{Bloch-Kato conditions.} By (d), Lemma~\ref{sha-siegel-step-ideal-seeds-inputs-tangent-bloch-kato} applies to $x$, $\rho_0$, $x'$ and the data of (c): the inflation of $z$ to $\operatorname{Gal}(\bar L/L^+)$ satisfies the Bloch-Kato conditions.
\item[(g)] \emph{A coboundary.} By (e) and (f), Theorem~\ref{sha-siegel-automorphy-adjoint} gives $X$ with $z(g)=R(g)X-X$ for $g\in\Gamma$; so $y(g)=\rho_0(g)X\rho_0(g)^{-1}-X$ for $g\in\Gamma_0$.
\item[(h)] \emph{Conclusion.} By (g), (d) and the flags of $\rho_0$, with the distinct characters $\psi_{x,v,i}$, Lemma~\ref{sha-siegel-step-ideal-seeds-inputs-tangent-constant} gives $x'=x$.
\end{enumerate}
% (1) is hida_modular_ideal_map (HidaModularIdealMap.lean, 52 lines; also Ws.h5), followed step by step: beta_0 is H5.descend K; (a) is H5.phiTw (PhiTw.lean) for phi_n, and the computation hu for the inverse; (b) is key, with stab (H5.twMapEx, H5.twJN from hM.dense and hM.twist, H5.closedJ), twChOps for the image in B^ord, TwMap.comp and H5.twMapUniq; (c) is le_antisymm with twChOps for the power and the inverse. WeakShimuraTheory/TwistedFourierKernels has about 860 lines, all written out in Lemmas sha-siegel-step-ideal-seeds-inputs-twist-*.
% (2) is eigenbasis_rank_four (EigenbasisRankFour.lean, 174 lines), followed step by step: (a) is Ws.Eb.regOfTan with hU.lam and lamIso (Lambda a power series ring) and hJ.ring.fin, hJ.ring.tf; (b) is Ws.Eb.factor (Factor.lean, through trPlus and CA.d_gext); (c) is liftRep, relKill, relForm, cocRel, adExt and cocExt; (d) is psiInj (PsiInj.lean) and flagDual (FlagDual.lean: the private charpoly_lift for the characteristic polynomial, and AlgebraTools.lemma_L, which is the proof of Lemma sha-rigidity-lagrange with sep_chars, Lemma sha-siegel-step-ideal-seeds-inputs-tangent-constant-separating, for w and ord_ext for the products over A[D_v]); (e) is hapt (the private arith with Loc.laArith), the private free (Z.Win.RR.exists_cycZ_eq_one_freeRatios_of_ne_mul_cycBar, Lemma sha-rigidity-free, with noRatio), Th.nice0four_alt (Lemma sha-siegel-automorphy-rank-four, with the hypotheses hC and hD4, which are P1.condC and P1.condD for a link of a fit tower) and the private par; (f) is bkAb with at_l, locL and bkAw; (g) is the hypothesis Th.AdjointLiftingTheorem (Theorem sha-siegel-automorphy-adjoint); (h) is conjOf and constOf. WeakShimuraTheory/EigenbasisLifts has about 2,560 lines; the local steps at v_0 (Lemmas sha-siegel-step-ideal-seeds-inputs-tangent-v0-local-*) and above l (Lemmas -tangent-crystalline-*) are written out.
\end{proof}

\begin{lemma}\label{sha-siegel-step-ideal-seeds}
In Situation~\ref{sha-siegel-step-ideal-setting}, every homomorphism of $\mathcal O'$-algebras $y:A\to\bar{\mathbf Q}_\ell$ that carries $P_X$ to a seed (Definition~\ref{sha-siegel-tower-ideal}) detects $J_\varphi$.
\end{lemma}
\leanref{Development/AutomorphicArithmetic/AutomorphicVarieties/ModularityArgument/ThresholdArguments/DenseThresholdCases/BirthPm.lean}{UpperTheory.V.Th.local_zero_of_modular_p}
\begin{proof}
Let $y(P_X)=(r,(n_w))$, a seed, with $h$ as in Definition~\ref{sha-siegel-tower-ideal}(1).
\begin{enumerate}
\item[(a)] \emph{The twist.} As $h_1-h_i$ is even for every $i$, there is $n\ge0$ with $h_i+n$ even for every $i$. Let $\varphi_n$ be as in Lemma~\ref{sha-siegel-step-ideal-seeds-phi}, $x_n$ and $\rho_n$ as in Lemma~\ref{sha-siegel-step-ideal-seeds-point}, and $m_{v,i}$ and $k$ as in Lemma~\ref{sha-siegel-step-ideal-seeds-weights}(2). As $h_i-h_{i+1}\ge4$, $m_{v,j}+2\le m_{v,i}$ for $i<j$.
\item[(b)] \emph{The point.} The $r(g)$ with $g\in G_L\cap\ker\chi_{\mathrm{cyc}}$ span $M_4(\bar{\mathbf Q}_\ell)$, so Lemma~\ref{sha-siegel-step-ideal-seeds-point} applies: $\rho_n$ is continuous and irreducible on $G_L\cap\ker\chi_{\mathrm{cyc}}$; $x_n$ carries the characteristic polynomial of $D^{\mathrm u}$ at $\mathrm{pr}(g)$ to that of $\rho_n(g)$; $\rho_n$ is trivial on the inertia group of $G_L$ at the place above $\tilde v_0$ whose image in $\Gamma_0$ is $I_{v_0}$; and for $v\mid\ell$, the restriction of $\rho_n$ to $D_v$ is triangular with diagonal characters $\psi_{x_n,v,1},\dots,\psi_{x_n,v,4}$. By Lemma~\ref{sha-siegel-step-ideal-seeds-automorphy}, $\rho_n$ is attached through $\iota$ to a regular algebraic cuspidal automorphic representation $\Pi$ of $GL_4(\mathbf A_L)$ with $\Pi^\vee\cong\Pi^c$, and it is windowed.
\item[(c)] \emph{The zero locus of $J_0$.} Proposition~\ref{sha-hecke-ideal}(3), in rank $4$, gives $J_0\subseteq\ker x_n$. Its hypotheses hold by (b), together with the following two facts. First, $x_n$ is arithmetic of the weight $h'$ with $h'_{v,j,\tau}=-m_{v,j}$, which is increasing: the $\psi_{x_n,v,j}$ are continuous, by the argument of step (2) of the proof of Lemma~\ref{sha-siegel-step-ideal-seeds-point} applied to $B^{\mathrm{ord}}$, and $(\psi_{x_n,v,j}\chi_{\mathrm{cyc}}^{-m_{v,j}})^k=1$ on $I_v$, so local class field theory applies as in step (a)(ii) of the proof of Lemma~\ref{sha-siegel-step-ideal-away}. Second, $\psi_{x_n,v,i}\ne\psi_{x_n,v,j}\varepsilon$ for all $i,j$. Indeed, $1-m_{v,i}+m_{v,j}\ne0$, and $\chi_{\mathrm{cyc}}(I_v)$ is infinite, so some $g\in I_v$ has $\chi_{\mathrm{cyc}}(g)^{(1-m_{v,i}+m_{v,j})k}\ne1$; if $\psi_{x_n,v,i}=\psi_{x_n,v,j}\chi_{\mathrm{cyc}}$, then $(\psi_{x_n,v,i}\chi_{\mathrm{cyc}}^{-m_{v,i}})^k=(\psi_{x_n,v,j}\chi_{\mathrm{cyc}}^{-m_{v,j}})^k\chi_{\mathrm{cyc}}^{(1-m_{v,i}+m_{v,j})k}$, and evaluating at $g$ gives a contradiction.
\item[(d)] \emph{Detection.} Let $\mathfrak q$ be a prime of $B$ whose image in $A$ is killed by $y$. Then $\varphi_n^{-1}(\mathfrak q)$ is a prime of $B^{\mathrm{ord}}$ contained in $\ker x_n$. By (a) to (c), Lemma~\ref{sha-siegel-step-ideal-seeds-inputs}(2) applies to $x_n$ and $\rho_n$, so $\varphi_n^{-1}(\mathfrak q)$ contains $J_0$. Hence $\mathfrak q$ contains $\varphi_n(J_0)$, and so the ideal it generates, which is $J_\varphi$ by Lemma~\ref{sha-siegel-step-ideal-seeds-inputs}(1).
\end{enumerate}
\end{proof}
% local_zero_of_modular_p (BirthPm.lean, 66 lines): (a) is the choice of n with Even (w 0 + n), exists_phi_universal_ordinary_det_ring_of_link for phi' = phi_n, and Pt.ptWts with hg for the gaps; (b) is Pt.ptIrr, Pt.rho_cont with Pt.pt_cont, auRep, Pt.ptUnr, Pt.ptFlag, Pt.ptPair and Z.Win.Birth.ptAutR'_alt with Z.ImaginaryQuadraticEtaCharacter.potOfPair; (c) is Pt.point_representation_weights (PointRepresentationWeights.lean, 56 lines: hM.pt, which is Proposition sha-hecke-ideal (3), Loc.laArith with psi_cont, and cycIne for psi_i != psi_j eps); (d) is hida_modular_ideal_map and Th.eigenbasis_rank_four (Lemma sha-siegel-step-ideal-seeds-inputs) applied to q.comap phi_n. The Lean asks for Ws.UnramifiedInputs and Th.AdjointLiftingTheorem (Theorem sha-siegel-automorphy-adjoint) as hypotheses. WeakShimuraTheory/ResidueFlagPoints has about 650 lines.

\subsection{Purity and descent}\label{subsection-sha-purity-descent}

\begin{lemma}\label{sha-siegel-step-ideal-purity-roots}
In Situation~\ref{sha-hecke-situation}, let $x$ be a point (Definition~\ref{sha-hecke-arcs}) and $\rho:\operatorname{Gal}(\bar L/L)\to GL_n(\bar{\mathbf Q}_\ell)$ a homomorphism such that $x(D^{\mathrm u})$ is the determinant of $\rho$. For $y\in\mathcal O[\operatorname{Gal}(\bar L/L)]$ write $\bar\chi_i(y)\in k$ for the image of $y$ under the linear extension of $\bar\chi_i\circ\mathrm{pr}$. Then for every $y\in\mathcal O[\operatorname{Gal}(\bar L/L)]$ and every eigenvalue $z$ of $\rho(y)$ there are $i$ and $a\in\mathcal O$ with residue $\bar\chi_i(y)$ such that $|z-a|<1$.
\end{lemma}
\leanref{Development/AutomorphicArithmetic/AutomorphicVarieties/WeakShimuraTheory/CharacteristicPolynomialRoots/LiftDiag.lean}{UpperTheory.V.Ws.Pu.ptRed}
\begin{proof}
Let $P\in B[X]$ be the characteristic polynomial of $D^{\mathrm u}$ at the image of $y$ in $B[\Gamma_0]$, a monic polynomial. As $x(D^{\mathrm u})$ is the determinant of $\rho$, $x(P)$ is the characteristic polynomial of $\rho(y)$. As $D^{\mathrm u}$ reduces to $\bar D$, the determinant of $\bar\chi_1\oplus\dots\oplus\bar\chi_n$, the reduction of $P$ is $\prod_i(X-\bar\chi_i(y))$. Choose $a_i\in\mathcal O$ with residue $\bar\chi_i(y)$; then the coefficients of $P-\prod_i(X-a_i)$ lie in $\mathfrak m_B$. By Lemma~\ref{sha-automorphy-unramified-near-tangent-points}, $|x(b)|\le1$ for $b\in B$ and $|x(b)|<1$ for $b\in\mathfrak m_B$. The eigenvalue $z$ is a root of the monic polynomial $x(P)$, whose coefficients have absolute value at most $1$, so $|z|\le1$. Hence $\prod_i(z-a_i)=-x\big(P-\prod_i(X-a_i)\big)(z)$ has absolute value less than $1$, and $|z-a_i|<1$ for some $i$.
\end{proof}
% ptRed (LiftDiag.lean, 201 lines): chi_map for x(P) and for the reduction (B.u.red, lift_diag), GradeZeroInterpolatingElement.coPt for the absolute values, and the private near for the last step.

\begin{lemma}\label{sha-siegel-step-ideal-purity-multiplicity}
Let $\mathcal O\subset\bar{\mathbf Q}_\ell$ be the ring of integers of a finite extension of $\mathbf Q_\ell$, with residue field $k$. Let $\Gamma$ be a group, $\rho:\Gamma\to GL_n(\bar{\mathbf Q}_\ell)$ a homomorphism, and $\bar\chi_1,\dots,\bar\chi_n:\Gamma\to k^\times$ and $\varepsilon:\Gamma\to\mathcal O^\times$ homomorphisms; extend the $\bar\chi_i$, and the $\bar\chi_i\bar\varepsilon$, linearly to ring homomorphisms $\mathcal O[\Gamma]\to k$. Suppose that:
\begin{enumerate}
\item[(i)] for every $y\in\mathcal O[\Gamma]$ and every eigenvalue $z$ of $\rho(y)$ there are $i$ and $a\in\mathcal O$ with residue $\bar\chi_i(y)$ such that $|z-a|<1$;
\item[(ii)] $\bar\chi_i\ne\bar\chi_j\bar\varepsilon$ for all $i$ and $j$.
\end{enumerate}
Let $\rho'_1,\dots,\rho'_a$ be representations of $\Gamma$ over $\bar{\mathbf Q}_\ell$ of positive dimensions, and $m_1,\dots,m_a\ge1$ integers, with $\det(X-\rho(g))=\prod_i\prod_{0\le j<m_i}\det\big(X-\varepsilon(g)^{-j}\rho'_i(g)\big)$ for $g\in\Gamma$. Then $m_i=1$ for every $i$.
\end{lemma}
\leanref{Development/AutomorphicArithmetic/AutomorphicVarieties/WeakShimuraTheory/CharacteristicPolynomialRoots/MemMax.lean}{UpperTheory.V.Ws.Pu.mOne}
\leanref{Development/AutomorphicArithmetic/AutomorphicVarieties/WeakShimuraTheory/CharacteristicPolynomialRoots/MemRootsIffSpec.lean}{UpperTheory.V.Ws.Pu.cpAlg}
\leanref{Development/AutomorphicArithmetic/AutomorphicVarieties/WeakShimuraTheory/CharacteristicPolynomialRoots/ExistsSeparatingElement.lean}{UpperTheory.V.Ws.Pu.sepEl}
\begin{proof}
\begin{enumerate}
\item[(a)] \emph{A separating element.} There is $y\in\mathcal O[\Gamma]$ with $\bar\chi_a(y)=1$ and $(\bar\chi_b\bar\varepsilon)(y)=0$ for all $a$ and $b$. Indeed, for $a$ and $b$, (ii) gives $\gamma\in\Gamma$ with $\bar\chi_a(\gamma)\ne(\bar\chi_b\bar\varepsilon)(\gamma)$; with $s,t\in\mathcal O$ having the residues $(\bar\chi_a(\gamma)-(\bar\chi_b\bar\varepsilon)(\gamma))^{-1}$ and $(\bar\chi_b\bar\varepsilon)(\gamma)$, the element $e_{ab}=s(\gamma-t)$ has $\bar\chi_a(e_{ab})=1$ and $(\bar\chi_b\bar\varepsilon)(e_{ab})=0$. Let $y_a=\prod_be_{ab}$ and $y=1-\prod_a(1-y_a)$. Then $\bar\chi_a(y_a)=1$ and $(\bar\chi_b\bar\varepsilon)(y_a)=0$ for all $a,b$, so $\bar\chi_a(y)=1$ and $(\bar\chi_b\bar\varepsilon)(y)=0$.
\item[(b)] \emph{Eigenvalues of the factors.} Let $\sigma_{ij}=\varepsilon^{-j}\rho'_i$ for $0\le j<m_i$. For $y\in\mathcal O[\Gamma]$, every eigenvalue of $\sigma_{ij}(y)$ is an eigenvalue of $\rho(y)$. Indeed, the hypothesis gives $\operatorname{tr}\rho(g)=\sum_{i,j}\operatorname{tr}\sigma_{ij}(g)$ for $g\in\Gamma$, hence for $g\in\mathcal O[\Gamma]$ by linearity, and so $\operatorname{tr}p(\rho(y))=\sum_{i,j}\operatorname{tr}p(\sigma_{ij}(y))$ for every polynomial $p$ over $\bar{\mathbf Q}_\ell$, applying it to the powers of $y$. Suppose that an eigenvalue $z$ of some $\sigma_{i_0j_0}(y)$ is not an eigenvalue of $\rho(y)$, and let $p$ vanish at every eigenvalue of $\rho(y)$ and of the $\sigma_{ij}(y)$ other than $z$, with $p(z)\ne0$. The eigenvalues of $p(M)$ are the $p(u)$ for the eigenvalues $u$ of $M$, so $\operatorname{tr}p(\rho(y))=0$, while $\sum_{i,j}\operatorname{tr}p(\sigma_{ij}(y))$ is $p(z)$ times the total multiplicity of $z$, which is positive. This contradicts $p(z)\ne0$ in characteristic $0$.
\item[(c)] \emph{Conclusion.} Suppose that $m_i\ge2$, and let $y=\sum_gc_gg$ be as in (a) and $Y=\sum_gc_g\varepsilon(g)g$. Then $\sigma_{i1}(Y)=\rho'_i(y)=\sigma_{i0}(y)$, and $(\bar\chi_b)(Y)=(\bar\chi_b\bar\varepsilon)(y)=0$ for every $b$. Let $z$ be an eigenvalue of $\rho'_i(y)$. By (b), $z$ is an eigenvalue of $\rho(y)$ and of $\rho(Y)$. By (i) there are $a_0,a_1\in\mathcal O$ with $|z-a_0|<1$, $|z-a_1|<1$, the residue of $a_0$ some $\bar\chi_{b_0}(y)=1$ and that of $a_1$ some $\bar\chi_{b_1}(Y)=0$. Then $|a_0-a_1|<1$, so $a_0-a_1\in\mathfrak m_{\mathcal O}$; but its residue is $1$.
\end{enumerate}
\end{proof}
% mOne (MemMax.lean, 100 lines): (a) is sepEl (ExistsSeparatingElement.lean, 55 lines, with the private exists_separating_element); (b) is cpAlg (MemRootsIffSpec.lean, 93 lines: traces, then a polynomial p vanishing at the other roots); (c) is the private twHom and lift_tw for Y, and mem_max.

\begin{lemma}\label{sha-siegel-step-ideal-purity-frobenius}
In Situation~\ref{sha-hecke-situation}, let $\iota:\bar{\mathbf Q}_\ell\cong\mathbf C$, $x$ a point and $\rho:\operatorname{Gal}(\bar L/L)\to GL_n(\bar{\mathbf Q}_\ell)$ a semisimple representation such that $x(D^{\mathrm u})$ is the determinant of $\rho$. Let $w$ be a place of $L$ not above a place of $S$.
\begin{enumerate}
\item $\rho$ is trivial on the inertia groups at the places of $\bar L$ above $w$.
\item Suppose that any two eigenvalues $z,z'$ of $\rho$ at the geometric Frobenius elements at $w$ satisfy $|\iota z|<\operatorname{N}w\,|\iota z'|$. Let $A$ be a place of $\bar L$ above $w$, $p$ the prime below $w$, $f\ge1$, and $\sigma$ an element of the decomposition group at $A$ that acts as $u\mapsto u^{p^f}$ on the residue field of $A$. Then any two eigenvalues $x',y'$ of $\rho(\sigma)$ satisfy $|\iota x'|<p^f|\iota y'|$.
\end{enumerate}
\end{lemma}
\leanref{Development/AutomorphicArithmetic/AutomorphicVarieties/WeakShimuraTheory/CharacteristicPolynomialRoots/LevFin.lean}{UpperTheory.V.Ws.Pu.levFin}
\leanref{Development/AutomorphicArithmetic/AutomorphicVarieties/WeakShimuraTheory/CharacteristicPolynomialRoots/RepKill.lean}{UpperTheory.V.Ws.Pu.repKill}
\leanref{Development/AutomorphicArithmetic/AutomorphicVarieties/WeakShimuraTheory/CharacteristicPolynomialRoots/FrobEigEx.lean}{UpperTheory.V.Ws.Pu.frobEigEx}
\leanref{Development/AutomorphicArithmetic/Foundations/GaloisRepresentations/AdmissibleResidues/PosOfMulPosOfPos.lean}{UpperTheory.V.Ws.Pu.frobAbs_w}
\begin{proof}
(1) Let $g$ lie in the inertia group at a place above $w$. As $w$ lies above no place of $S$, the image of $g$ in $\Gamma$, and so $\mathrm{pr}(g)$, is $1$. For every $h$, the characteristic polynomials of $\rho(gh)$ and $\rho(h)$ are the images under $x$ of those of $D^{\mathrm u}$ at $\mathrm{pr}(gh)=\mathrm{pr}(h)$, so they agree, and $\operatorname{tr}\big((\rho(g)-1)\rho(h)\big)=0$. As $\rho$ is semisimple, the trace form is nondegenerate on the span of the image of $\rho$, which contains $\rho(g)-1$; so $\rho(g)=1$.

(2) By (1), $\rho$ is unramified at $w$, and the characteristic polynomial of $\rho$ at the geometric Frobenius elements at the places above $w$ does not depend on the place. Let $\operatorname{N}w=p^e$ and let $\phi$ be an element of the decomposition group at $A$ that acts as $u\mapsto u^{\operatorname{N}w}$ on the residue field. Then $\sigma^e$ and $\phi^f$ both act as $u\mapsto u^{p^{ef}}$, so $\sigma^{-e}\phi^f$ lies in the inertia group at $A$, and $\rho(\sigma)^e=\rho(\phi)^f$. So for an eigenvalue $x'$ of $\rho(\sigma)$ there is an eigenvalue $u$ of $\rho(\phi)$ with $u^f=x'^e$; then $c=\iota(u)^{-1}$ is an eigenvalue of $\rho$ at the geometric Frobenius element $\phi^{-1}$ (through $\iota$), and $|\iota x'|^e=|c|^{-f}$. Likewise $|\iota y'|^e=|c'|^{-f}$ for such a $c'$. The hypothesis gives $|c|>0$ and $|c'|<\operatorname{N}w\,|c|$, so $|c|^{-1}<\operatorname{N}w\,|c'|^{-1}$, and $|\iota x'|^e=|c|^{-f}<p^{ef}|c'|^{-f}=p^{ef}|\iota y'|^e$. Taking $e$-th roots gives $|\iota x'|<p^f|\iota y'|$.
\end{proof}
% levFin (LevFin.lean, 34 lines: ine_dies, the places of L above V.S'); repKill (RepKill.lean, 33 lines: kill_of_charpoly with GlobalTheory.isSSRep_trace_nondegenerate, and repChi); frobEigEx (FrobEigEx.lean, 56 lines: frob_back and frob_ine) gives the multiset b of Af.FrobEig, the images under iota of the eigenvalues of rho(sigma^-1) for a Frobenius sigma; frobAbs_w (PosOfMulPosOfPos.lean: frobAbs_isFrob_pow, inv_mul_mem_Ine_of_isFrob, frobAbs_eig, exists_normV_eq_prime_pow_le_finrank).

\begin{lemma}\label{sha-siegel-step-ideal-purity}
In Situation~\ref{sha-siegel-step-ideal-setting}, let $y:A\to\bar{\mathbf Q}_\ell$ be a homomorphism of $\mathcal O'$-algebras that kills $J_\varphi$ and carries $P_X$ to a flagged pair with inertial weights $2m$ (Definition~\ref{sha-siegel-tower-ideal}). Then every homomorphism $r:G\to M_4(\bar{\mathbf Q}_\ell)$ whose restriction to $G_F$ is $y\circ r_X$ is pure through $\iota$ (Definition~\ref{sha-siegel-tower-ambient}).
\end{lemma}
\leanref{Development/AutomorphicArithmetic/AutomorphicVarieties/ModularityArgument/ThresholdArguments/DenseThresholdCases/VWinPm.lean}{UpperTheory.V.Ws.pmPur_w_alt}
\begin{proof}
Let $x$ be the composite of $\varphi$, $B\to A$ and $y$, a point of $B^{\mathrm{ord}}$ with $J_0\subseteq\ker x$, as $y$ kills $J_\varphi$.
\begin{enumerate}
\item[(a)] \emph{The weight.} The point $x$ is arithmetic of some weight. This is the argument of step (a)(ii) of the proof of Lemma~\ref{sha-siegel-step-ideal-away}, which uses only that $y(P_X)$ has inertial weights $2m$, with $m$ strictly decreasing.
\item[(b)] \emph{The representation.} By Proposition~\ref{sha-hecke-ideal}(4), $x(D^{\mathrm u})$ is the determinant of a continuous semisimple $\rho:G_L\to GL_4(\bar{\mathbf Q}_\ell)$, and there are continuous $\rho'_1,\dots,\rho'_a$ of positive dimensions and integers $m_1,\dots,m_a\ge1$ with $\det(X-\rho(g))=\prod_i\prod_{0\le j<m_i}\det\big(X-\rho'_i(g)\chi_{\mathrm{cyc}}(g)^{-j}\big)$ for $g\in G_L$.
\item[(c)] \emph{The multiplicities.} Every $m_i$ is $1$, by Lemma~\ref{sha-siegel-step-ideal-purity-multiplicity}, applied to $\Gamma=G_L$, $\rho$, the characters $\bar\chi_i\circ\mathrm{pr}$ and $\varepsilon=\chi_{\mathrm{cyc}}$ with values in $\mathcal O'^\times$. Hypothesis (i) is Lemma~\ref{sha-siegel-step-ideal-purity-roots}. For (ii), let $i,j$, and let $w'$ be the weight of the entry $(i,j)$ (or $0$ if $i=j$). By (e) of Definition~\ref{sha-siegel-tower-towers}, some $g\in G_L$ has $\bar\chi^{w'}(g)\ne\bar\omega(g)$. By (3) of Definition~\ref{sha-siegel-automorphy-link}, $\bar\chi_i(\mathrm{pr}(g))$ is the residue of $\kappa(g)\chi_i(g)$, so $\bar\chi_i(\mathrm{pr}(g))\bar\chi_j(\mathrm{pr}(g))^{-1}=\bar\chi^{w'}(g)$, while the residue of $\chi_{\mathrm{cyc}}(g)$ is $\bar\omega(g)$; so $\bar\chi_i\circ\mathrm{pr}\ne(\bar\chi_j\circ\mathrm{pr})\bar\chi_{\mathrm{cyc}}$ at $g$.
\item[(d)] \emph{The Frobenius eigenvalues.} As all $m_i=1$, Proposition~\ref{sha-hecke-ideal}(4) gives a finite set $Z_1$ of places of $L$ outside which any two Frobenius eigenvalues $z,z'$ of $\rho$ satisfy $|\iota z|<\operatorname{N}w\,|\iota z'|$. Let $Z$ be the union of $Z_1$ and the set of places of $L$ above the places of $S$, which is finite. By Lemma~\ref{sha-siegel-step-ideal-purity-frobenius}, for every place $A$ of $\bar{\mathbf Q}$ above a place $w\notin Z$ of $L$, every $f\ge1$ and every $\sigma\in G_L$ in the decomposition group at $A$ that acts as $u\mapsto u^{p^f}$ on its residue field, where $p$ is below $w$, any two eigenvalues $x',y'$ of $\rho(\sigma)$ satisfy $|\iota x'|<p^f|\iota y'|$.
\item[(e)] \emph{Purity.} Let $H=G_L$, an open subgroup of $G$, and $M=\prod_{w\in Z}\operatorname{N}w\ne0$. Let $p\nmid M$ be a prime, $f\ge1$, $\sigma\in H$ acting as $u\mapsto u^{p^f}$ on the residue field of a place $A$ above $p$, and $x',y'\ne0$ eigenvalues of $r(\sigma)$. The place $w$ of $L$ below $A$ is not in $Z$, as otherwise $p$ would divide $\operatorname{N}w$, which divides $M$. As in step (a)(i) of the proof of Lemma~\ref{sha-siegel-step-ideal-away}, the characteristic polynomial of $\rho(\sigma)$, the image under $x$ of that of $D^{\mathrm u}$ at $\mathrm{pr}(\sigma)$, is that of $\kappa(\sigma)r(\sigma)$, as $r(\sigma)=y(r_X(\sigma))$. So $\kappa(\sigma)x'$ and $\kappa(\sigma)y'$ are eigenvalues of $\rho(\sigma)$, and by (d) $|\iota(\kappa(\sigma)x')|<p^f|\iota(\kappa(\sigma)y')|$; dividing by $|\iota\kappa(\sigma)|$ gives $|\iota x'|<p^f|\iota y'|$. So $r$ is pure through $\iota$.
\end{enumerate}
% pmPur_w_alt (VWinPm.lean, 108 lines in all): (a) is Pu.arith (auArith with D.wts4, from D.Good); (b) is GalN.V.end1W' with hM.end1 and Link.le_ker_pt; (c) is Pu.mOne with Pu.ptRed and Pu.noOm (NoOmLift.lean: T.Ok.g2, which is (e), noOm_lift and noOm_rt, with D.Good.om_bar); (d) is the cofinite set of hpure joined to Pu.levFin, with Pu.frobEigEx (through Pu.repKill) and Pu.frobAbs_w; (e) is Pu.pureOf_w (H = L.fixingSubgroup, M the product of the norms of the places of Z) with Pu.ptCp. WeakShimuraTheory/CharacteristicPolynomialRoots has about 750 lines.
\end{proof}

\begin{lemma}\label{sha-siegel-step-ideal-away}
In Situation~\ref{sha-siegel-step-ideal-setting}, let $y:A\to\bar{\mathbf Q}_\ell$ be a homomorphism of $\mathcal O'$-algebras that kills $J_\varphi$ and carries $P_X$ to a flagged pair with inertial weights $2m$, and let $r:G\to M_4(\bar{\mathbf Q}_\ell)$ be a continuous homomorphism whose restriction to $G_F$ is $y\circ r_X$, with $r(g)$ symplectic with multiplier $\theta(g)$ for all $g$, and irreducible on $G_L$. Then the representation $g\mapsto\kappa(g)r(g)$ of $G_L$ is attached through $\iota$ to a regular algebraic cuspidal automorphic representation $\pi$ of $GL_4(\mathbf A_L)$ with $\pi^\vee\cong\pi^c$, and it is windowed at the places not above $v_0$.
\end{lemma}
\leanref{Development/AutomorphicArithmetic/AutomorphicVarieties/ModularityArgument/ThresholdArguments/DenseThresholdCases/PmAutO.lean}{UpperTheory.Gal.projectedFormsSystemPart.pmAutO_alt}
\begin{proof}
Let $x$ be the composite of $\varphi$, $B\to A$ and $y$, a point of $B^{\mathrm{ord}}$ with $J_0\subseteq\ker x$, as $y$ kills $J_\varphi$.
\begin{enumerate}
\item[(a)] \emph{The point.} Let $\rho_L(g)=\kappa(g)r(g)$ for $g\in G_L$, a continuous representation, as $r$ and $\kappa$ are continuous. The point $x$ is arithmetic of some weight, the composite of $x(D^{\mathrm u})$ with $\mathrm{pr}:G_L\to\Gamma_0$ is the determinant of $\rho_L$, and $\rho_L$ is irreducible. Write $\psi''_{\bar v,i}$ for the diagonal characters of the flagged pair $y(P_X)=(r|_{G_F},(y(n_{X,w})))$.
\begin{enumerate}
\item[(i)] \emph{The determinant.} By Lemma~\ref{sha-siegel-automorphy-phi}, $\varphi$ carries the characteristic polynomial of $D^{\mathrm u}$ at $\mathrm{pr}(g)$, for $g\in G_L$, to $c^\kappa_g(Y)=\kappa(g)^4c_g(\kappa(g)^{-1}Y)$. The image of $c_g$ in $A[Y]$ is the characteristic polynomial of $r_X(g)$ (Definition~\ref{sha-siegel-tower-space}), and $y(r_X(g))=r(g)$, as $G_L\subseteq G_F$. So $x$ carries the characteristic polynomial of $D^{\mathrm u}$ at $\mathrm{pr}(g)$ to $\kappa(g)^4\det(\kappa(g)^{-1}Y-r(g))=\det(Y-\rho_L(g))$.
\item[(ii)] \emph{The weight.} Let $v\mid\ell$. By Lemma~\ref{sha-siegel-automorphy-phi}, $\psi_{x,v,i}=\kappa\psi''_{\bar v,i}$ on $D_v$; so $\psi_{x,v,i}$ is continuous, as $\kappa$ and the entries of $r$ are. Let $a_v$ and $k_v$ be the integers of (e) of Definition~\ref{sha-siegel-automorphy-link} at the place of $\bar{\mathbf Q}$ at which $D_v$ is the decomposition group of $G_L$. As $y(P_X)$ has inertial weights $2m$, there is an open subgroup $U$ of $G_F$ with $\psi''_{\bar v,i}(s)^2=\chi_{\mathrm{cyc}}(s)^{2m_i}$ for $s\in U\cap I_{\bar v}(F)$; $U$ contains an open normal subgroup of $G_F$, of some index $N$, so $g^N\in U$ for every $g\in G_F$. Hence $\psi_{x,v,i}(g)^{2N}=\kappa(g)^{2N}\chi_{\mathrm{cyc}}(g)^{2Nm_i}$ for $g\in I_v$, and as $(\kappa\chi_{\mathrm{cyc}}^{-a_v})^{k_v}=1$ on $I_v$, $\big(\psi_{x,v,i}\chi_{\mathrm{cyc}}^{-(m_i+a_v)}\big)^{2Nk_v}=1$ on $I_v$. By local class field theory, as in step (F) of the proof of Proposition~\ref{sha-siegel-points}, $x$ is arithmetic of the weight $h$ with $h_{v,i,\tau}=-(m_i+a_v)$ (Definition~\ref{sha-hecke-arcs}(1)), which is increasing, as $m=(k+1,k,1,0)$ is strictly decreasing.
\item[(iii)] \emph{Irreducibility.} As $r$ is irreducible on $G_L$, the $r(g)$ with $g\in G_L$ span $M_4(\bar{\mathbf Q}_\ell)$. A subspace of $\bar{\mathbf Q}_\ell^4$ stable under $\rho_L$ is stable under every $r(g)=\kappa(g)^{-1}\rho_L(g)$ with $g\in G_L$, hence under $M_4(\bar{\mathbf Q}_\ell)$, so it is $0$ or $\bar{\mathbf Q}_\ell^4$.
\end{enumerate}
% (a) is auArith (AuArith.lean, 98 lines: the open subgroup U of Tower.Wts and its normal core, K.cyc for a_v and k_v, auPsi for psi_x = kappa psi'', mul_zpow_pow_eq_one_of_pow_eqs, then Loc.laArith from GlobalArtinReciprocity with the weight -(m_j + a_v)), auPsiCont (PsiCont.lean, through hφ.psi_val and auCont), auRep (CharpolySmulScale.lean: hφ.1, scale_point and CA.d_gext; Pt.IsRep is the equality of x(D^u) composed with pr and the determinant of the twist) and auIrr (IrrOfStab.lean: irr_of_stab), in WeakShimuraTheory/AutomorphicCharacters, about 420 lines.
\item[(b)] \emph{Classicality.} By Proposition~\ref{sha-hecke-ideal}(5), applied to $x$, which has $J_0\subseteq\ker x$, and to $\rho_L$, the representation $\rho_L$ is attached through $\iota$ to a regular algebraic cuspidal automorphic representation $\pi$ of $GL_4(\mathbf A_L)$ with $\pi^\vee\cong\pi^c$, and it is windowed at the places not above $v_0$.
\end{enumerate}
\end{proof}
% pmAutO_alt (PmAutO.lean, 47 lines): (b) is end2oN with hM.end2, the field IrrArithPointsAutomorphicAway of ModularAlt, which is Proposition sha-hecke-ideal (5); the conclusion is AutOffN K.s.v0 for Af.tw of r on G_L by κ.

\begin{lemma}\label{sha-siegel-step-ideal-descent-third}
Let $\ell$ be a prime, $n\ge1$, $G=\operatorname{Gal}(\bar{\mathbf Q}/\mathbf Q)$ and $r:G\to GL_n(\bar{\mathbf Q}_\ell)$ a homomorphism. Let $K$ and $K'$ be imaginary quadratic fields such that $r$ is irreducible on $G_K$ and on $G_{K'}$. Then there is an imaginary quadratic field $K''$ such that $r$ is irreducible on $G_K\cap G_{K''}$ and on $G_{K'}\cap G_{K''}$.
\end{lemma}
\leanref{Development/AutomorphicArithmetic/AutomorphicVarieties/ModularityArgument/BaseChangeUnits/CentralResidueGroups/EqOfRes.lean}{UpperTheory.V.Z.Bc.third}
\begin{proof}
For an imaginary quadratic field $E$ let $\eta_E:G\to\{\pm1\}$ be its quadratic character, with kernel $G_E$.
\begin{enumerate}
\item[(a)] \emph{Index two.} Let $E$ be one of $K$, $K'$, and $\delta$ a character of $G$ with values $\pm1$. If $r$ is irreducible on $G_E$ and reducible on $G_E\cap\ker\delta$, a subgroup of index at most $2$ of $G_E$, then the restriction of $r$ to $G_E$ is isomorphic to its twist by $\delta$. Indeed, let $H=G_E\cap\ker\delta$ and $W$ a nonzero proper subspace of $\bar{\mathbf Q}_\ell^n$ stable under $r(H)$. If $\delta$ is trivial on $G_E$, the identity is such an isomorphism. Otherwise let $t\in G_E$ with $\delta(t)=-1$. Let $s\in G_E$. If $\delta(s)=1$, then $r(s)$ maps $W$ into $W$, and $r(t)W$ into $r(t)W$, as $r(s)r(t)=r(t)r(t^{-1}st)$ with $t^{-1}st\in H$. If $\delta(s)=-1$, then $r(s)$ maps $W$ into $r(t)W$, as $r(s)=r(t)r(t^{-1}s)$ with $t^{-1}s\in H$, and $r(t)W$ into $W$, as $r(s)r(t)=r(st)$ with $st\in H$. So $W\cap r(t)W$ and $W+r(t)W$ are stable under $r(G_E)$. As $r$ is irreducible on $G_E$ and $W$ is nonzero and proper, $W\cap r(t)W=0$ and $W+r(t)W=\bar{\mathbf Q}_\ell^n$. Let $P$ be the map that is $1$ on $W$ and $-1$ on $r(t)W$. Then $P^2=1$, and by the two cases $Pr(s)=\delta(s)r(s)P$ for every $s\in G_E$. So $P$ is an isomorphism from $r|_{G_E}$ to $r|_{G_E}\otimes\delta|_{G_E}$.
\item[(b)] \emph{Counting twists.} Let $E$ be as in (a), and $\delta_1,\dots,\delta_N$ pairwise distinct characters of $G$ with values $\pm1$. At most $2n^2$ of them have $r|_{G_E}\cong r|_{G_E}\otimes\delta_i|_{G_E}$. Indeed, fix $t\in G\setminus G_E$. A character $\delta$ with values $\pm1$ is determined by its restriction to $G_E$ and its value at $t$: it is determined on $G_E$, and on the other coset by $\delta(s)=\delta(t)\delta(t^{-1}s)$. For each of the two values at $t$, the restrictions $\delta_i|_{G_E}$ with $r|_{G_E}\cong r|_{G_E}\otimes\delta_i|_{G_E}$ are therefore distinct characters of $G_E$; each gives a nonzero element of $\operatorname{End}(\bar{\mathbf Q}_\ell^n)$ on which $G_E$ acts, by conjugation through $r$, by that character, and eigenvectors for distinct characters are linearly independent; so there are at most $n^2$ of them.
\item[(c)] \emph{Choice of $K''$.} There are $4n^2+1$ distinct imaginary quadratic fields, with pairwise distinct quadratic characters. By (b), applied to $E=K$ and to $E=K'$, one of them, $K''$, has $r|_{G_K}\not\cong r|_{G_K}\otimes\eta_{K''}|_{G_K}$ and $r|_{G_{K'}}\not\cong r|_{G_{K'}}\otimes\eta_{K''}|_{G_{K'}}$. By (a), $r$ is irreducible on $G_K\cap G_{K''}$ and on $G_{K'}\cap G_{K''}$.
\end{enumerate}
\end{proof}
% Z.Bc.third (EqOfRes.lean, 120 lines): (a) is the private irr_inf with Br.resTwo (ResTwo.lean, its private part_one: the case δ = 1, the inclusions a1 to a4, hinf and hsup for the decomposition, and P = LinearMap.ofIsCompl with hPP and hPr);(b) is the private card_bad with eq_of_res and Br.twCard; (c) is Br.quadFam (4 n^2 + 1) with exists_quadratic_sign_character and Br.signNe, and Finset.exists_mem_notMem_of_card_lt_card.

\begin{lemma}\label{sha-siegel-step-ideal-descent-split}
Let $\ell$ be a prime, $n\ge1$, $G=\operatorname{Gal}(\bar{\mathbf Q}/\mathbf Q)$ and $r:G\to GL_n(\bar{\mathbf Q}_\ell)$ a homomorphism. Let $K$ and $K'$ be imaginary quadratic fields such that $r$ is irreducible on $G_K$ and on $G_{K'}$. Then there is an imaginary quadratic field $K''$, fixed by the decomposition group of $G$ at every place of $\bar{\mathbf Q}$ above $\ell$, such that $r$ is irreducible on $G_K\cap G_{K''}$ and on $G_{K'}\cap G_{K''}$.
\end{lemma}
\leanref{Development/AutomorphicArithmetic/AutomorphicVarieties/ModularityArgument/AutomorphyLifting/QuadraticCharpolySplitting/QuadFamSplit.lean}{UpperTheory.V.Z.Win.thirdSplit}
\leanref{Development/AutomorphicArithmetic/AutomorphicVarieties/ModularityArgument/AutomorphyLifting/QuadraticCharpolySplitting/QuadFamSplit.lean}{UpperTheory.V.Z.Win.quadFamSplit}
\begin{proof}
\begin{enumerate}
\item[(a)] \emph{The fields.} By Dirichlet's theorem there are $4n^2+1$ distinct primes $p\equiv-1$ modulo $4\ell$. For each of them, $\mathbf Q(\sqrt{-p})$ is an imaginary quadratic field of discriminant $-p$, as $p\equiv3$ modulo $4$; the discriminants are pairwise coprime, so the fields and their quadratic characters are pairwise distinct.
\item[(b)] \emph{The places above $\ell$.} Let $p$ be as in (a), $\delta\in\bar{\mathbf Q}$ with $\delta^2=-p$, and $c=(p+1)/(4\ell)$, an integer. Let $A$ be a place of $\bar{\mathbf Q}$ above $\ell$, $|\cdot|_A$ an absolute value defining it, and $g$ in the decomposition group at $A$, so that $|g(x)|_A<1$ if and only if $|x|_A<1$. Then $g(\delta)=\pm\delta$. Let $u=(1+\delta)/2$; then $u(1-u)=(1+p)/4=\ell c$, so $|u|_A<1$ or $|1-u|_A<1$. If $g(\delta)=-\delta$, then $g(u)=1-u$ and $g(1-u)=u$, so both $|u|_A$ and $|1-u|_A$ are less than $1$, which is impossible as $u+(1-u)=1$. So $g(\delta)=\delta$, and $g$ fixes $\mathbf Q(\delta)$.
\item[(c)] \emph{Choice of $K''$.} By (b) of the proof of Lemma~\ref{sha-siegel-step-ideal-descent-third}, applied to $E=K$ and to $E=K'$ and the quadratic characters of the fields of (a), one of these fields, $K''$, has $r|_{G_K}\not\cong r|_{G_K}\otimes\eta_{K''}|_{G_K}$ and $r|_{G_{K'}}\not\cong r|_{G_{K'}}\otimes\eta_{K''}|_{G_{K'}}$. By (a) of that proof, $r$ is irreducible on $G_K\cap G_{K''}$ and on $G_{K'}\cap G_{K''}$, and $K''$ is fixed by the decomposition groups above $\ell$, by (b).
\end{enumerate}
\end{proof}
% thirdSplit (QuadFamSplit.lean): quadFamSplit (Nat.infinite_setOfPred_prime_and_eq_mod for p = -1 mod 4l, the private fam_one with KStar.Consta.exists_quadField for discriminant -p) for (a); the private fix_root (u = (1 + δ)/2 with u(1 - u) = l c) and fix_field for (b); Br.signNe for the distinct characters, and the private card_bad and irr_inf of EqOfRes.lean for (c).

\begin{definition}\label{sha-siegel-step-ideal-descent-conditions}
Let $F\subseteq\bar{\mathbf Q}$ be a number field, $\ell$ a prime, $n\ge1$ and $r:\operatorname{Gal}(\bar{\mathbf Q}/F)\to GL_n(\bar{\mathbf Q}_\ell)$ a homomorphism. For a place $A$ of $\bar{\mathbf Q}$ let $D_A$ and $I_A$ be the decomposition and inertia groups of $\operatorname{Gal}(\bar{\mathbf Q}/F)$ at $A$; say that $r$ is \emph{triangular on $D_A$ with characters} $\kappa_1,\dots,\kappa_n:D_A\to\bar{\mathbf Q}_\ell^\times$ if some $Q\in GL_n(\bar{\mathbf Q}_\ell)$ has $Qr(g)Q^{-1}$ upper triangular with diagonal entries $\kappa_1(g),\dots,\kappa_n(g)$ for every $g\in D_A$.
\begin{enumerate}
\item $r$ satisfies \emph{condition (I)} if for every place $A$ of $\bar{\mathbf Q}$ above a finite place of $F$ not above $\ell$ there are continuous characters $\kappa_1,\dots,\kappa_n:D_A\to\bar{\mathbf Q}_\ell^\times$ with $\det(X-r(g))=\prod_i(X-\kappa_i(g))$ for $g\in I_A$.
\item $r$ satisfies \emph{condition (T)} if for every place $A$ of $\bar{\mathbf Q}$ above $\ell$ it is triangular on $D_A$ with some characters.
\item For a number field $L\supseteq F$ in $\bar{\mathbf Q}$, $r$ \emph{keeps $\ell$-distinctness in $L$} if for every place $w$ of $L$ above $\ell$ such that $r$ is $\ell$-distinct at the place of $F$ below $w$ (Definition~\ref{sha-eigensystem-distinctness}), its restriction to $\operatorname{Gal}(\bar{\mathbf Q}/L)$ is $\ell$-distinct at $w$. As the decomposition groups above a place are conjugate, $r$ is $\ell$-distinct at a place $v$ above $\ell$ if and only if for every place $A$ of $\bar{\mathbf Q}$ above $v$ it is triangular on $D_A$ with characters $\kappa_1,\dots,\kappa_n$ such that $\kappa_i\ne\kappa_j\chi_{\mathrm{cyc}}$ for all $i$ and $j$.
\end{enumerate}
\end{definition}
\leanref{Development/AutomorphicArithmetic/AutomorphicVarieties/ModularityArgument/AutomorphyLifting/WeakOfAtt.lean}{UpperTheory.V.Z.Win.IneAb}
\leanref{Development/AutomorphicArithmetic/AutomorphicVarieties/ModularityArgument/AutomorphyLifting/WeakOfAtt.lean}{UpperTheory.V.Z.Win.Flagged}
\leanref{Development/AutomorphicArithmetic/AutomorphicVarieties/ModularityArgument/AutomorphyLifting/WeakOfAtt.lean}{UpperTheory.V.Z.Win.Unl1}
\leanref{Development/AutomorphicArithmetic/AutomorphicVarieties/ModularityArgument/AutomorphyLifting/WeakOfAtt.lean}{UpperTheory.V.Z.Win.Keeps}

\begin{lemma}\label{sha-siegel-step-ideal-descent-updown-local-above}
Let $F\subseteq K$ be number fields, $\ell$ a prime, $\iota:\bar{\mathbf Q}_\ell\cong\mathbf C$, $v$ a place of $F$ above $\ell$, $w$ a place of $K$ above $v$, and $A$ a valuation ring of $\bar{\mathbf Q}$ over $w$, with decomposition groups $D_A(F)\supseteq D_A(K)$ in $\operatorname{Gal}(\bar{\mathbf Q}/F)$ and $\operatorname{Gal}(\bar{\mathbf Q}/K)$, Weil group $W_A\subseteq D_A(K)$ and local Artin map $a:W_A\to K_w^\times$ (Situation~\ref{gal-local-situation} and Definition~\ref{gal-local-artin}). A continuous character $\psi$ of $D_A(K)$ is \emph{locally algebraic with} $(\omega,h)$, for a smooth character $\omega$ of $K_w^\times$ and integers $h_\tau$ indexed by the embeddings $\tau:K\to\bar{\mathbf Q}_\ell$, if $\psi(s)=\iota^{-1}(\omega(a(s)))\prod_\tau\tau(a(s))^{-h_\tau}$ for $s\in W_A$, where the product runs over the $\tau$ that extend continuously to $K_w$, and $\tau$ also denotes that extension.
\begin{enumerate}
\item Let $r:D_A(K)\to GL_n(\bar{\mathbf Q}_\ell)$ be a continuous representation, conjugate to an upper triangular representation with diagonal characters $\psi_1,\dots,\psi_n$, and de Rham with Hodge-Tate multisets $\mathrm{HT}_\tau$. Then each $\psi_i$ is locally algebraic with some $(\omega_i,h_i)$.
\item If moreover $r$ is potentially semistable, then the characteristic polynomial of $\mathrm{WD}_\iota(r)(s)$ is $\prod_i(X-\omega_i(a(s)))$ for $s\in W_A$.
\item Let $\kappa$ be a continuous character of $D_A(F)$ whose restriction to $D_A(K)$ is locally algebraic with $(\omega,h)$. Then there is a smooth character $\chi$ of $F_v^\times$ such that $\omega\cdot(\chi\circ N_{K_w/F_v})$ is unramified.
\end{enumerate}
\end{lemma}
\leanref{Development/AutomorphicArithmetic/AutomorphicVarieties/HeckeAlgebras/Distinct.lean}{UpperTheory.V.Hi.OrdChar}
\leanref{Development/AutomorphicArithmetic/AutomorphicVarieties/ModularityArgument/OrdinaryJacquetModules/CoreOrdAlt.lean}{UpperTheory.V.Z.Ord.ExtChar}
\leanref{Development/AutomorphicArithmetic/AutomorphicVarieties/ModularityArgument/ThresholdArguments/DenseThresholdCases/OrdExtChar.lean}{UpperTheory.V.Z.Ord.ext_character_of_literal_automorphy_n}
\begin{proof}
\openstep{The proof is not yet written out: for (1), that the diagonal characters of a de Rham triangular representation are de Rham with one Hodge-Tate weight each, hence locally algebraic by local class field theory and Berger's theorem; for (2), the computation of the Weil-Deligne representation along the flag; for (3), the passage from the character of $D_A(K)$ to one of $F_v^\times$ through the norm, using that $\kappa$ is defined on $D_A(F)$.}
% (1) is the private ord_of_flag of OrdThruL.lean (with tri_dR, flag_jumps, cont_of_flag, Loc.laConv and Loc.ordCharIff); (2) is the private pol_of_flag (Loc.PstExact.flag and Loc.hecke_char_de_rham_of_local_artin_above_l); (3) is Z.Ord.ext_character_of_literal_automorphy_n (OrdExtChar.lean, 380 lines, localArtin_normU and the algebraic parts powU). All three use the literature hypothesis LiteratureRankNHypothesis (Berger's theorem and the l-adic characters of algebraic Hecke characters). Size: lean_size of ord_of_flag, pol_of_flag and ext_character_of_literal_automorphy_n together, stopping at the book's citations, 70,800 lines, of which about 5,400 in AutomorphicArithmetic/AutomorphicVarieties; the rest is mostly p-adic Hodge theory and local class field theory in AbelianVarietyTheory/Proofs (30,700) and GroundTheory/LocalFields (12,100).
\end{proof}

\begin{lemma}\label{sha-siegel-step-ideal-descent-updown-local-down}
Let $F\subset K$ be totally complex number fields, with $K$ Galois over $F$ of prime degree, $n\ge2$, and $\pi$ and $Q$ cuspidal automorphic representations of $GL_n(\mathbf A_F)$ and $GL_n(\mathbf A_K)$ such that $Q$ is a weak base change of $\pi$ (Definition~\ref{gal-depth-zero-lift-weak}). Let $v$ be a finite place of $F$.
\begin{enumerate}
\item Suppose that at every place $w$ of $K$ above $v$ the supercuspidal support of every local component of $Q$ at $w$ consists of characters $\zeta$, for each of which there is a smooth character $\chi'$ of $F_v^\times$ such that $\zeta\cdot(\chi'\circ N_{K_w/F_v})$ is unramified. Then the supercuspidal support of every local component of $\pi$ at $v$ consists of characters.
\item For every character $\xi$ in the supercuspidal support of a local component of $\pi$ at $v$ there is a place $w$ of $K$ above $v$ such that $\xi\circ N_{K_w/F_v}$ lies in the supercuspidal support of a local component of $Q$ at $w$.
\end{enumerate}
\end{lemma}
\leanref{Development/AutomorphicArithmetic/AutomorphicVarieties/ModularityArgument/GammaLaurentFactors/MatchIndOf.lean}{UpperTheory.V.Z.Gam.UpDown}
\leanref{Development/AutomorphicArithmetic/AutomorphicVarieties/ModularityArgument/GammaLaurentFactors/CrossG.lean}{UpperTheory.V.Z.Gam.upDown}
\leanref{Development/AutomorphicArithmetic/AutomorphicVarieties/ModularityArgument/LeviSmoothLocalization/WinUpDown.lean}{UpperTheory.V.Z.ImaginaryQuadraticEtaCharacter.UpDownX}
\leanref{Development/AutomorphicArithmetic/AutomorphicVarieties/ModularityArgument/LeviSmoothLocalization/PrincipalSeriesMatching/UpDownX.lean}{UpperTheory.V.Z.ImaginaryQuadraticEtaCharacter.upDownX}
\begin{proof}
\openstep{The proof is not yet written out. Alpöge compares the gamma factors of the local components of $\pi$ at $v$, twisted by characters of $F_v^\times$, with those of $Q$ at the places above $v$, twisted by the composites of these characters with the norm, and reads off the supercuspidal supports.}
% (1) is the down half of Z.Gam.upDown (upDownOf with normCnt and matchBCOf: pairGOps, crossG, sideG, normLoc, globTw, regroupBC); (2) is the second half of upDownX (matchBCX), which needs no hypothesis. Size: lean_size of upDown and upDownX, stopping at the book's citations, 125,000 lines, of which about 6,500 in AutomorphicArithmetic/AutomorphicVarieties (the gamma factors); the rest is mostly the local theory of GL_n in GroundTheory/AutomorphicForms (51,800) and AbelianVarietyTheory/Proofs (36,600); part of it is shared with the up half (Lemma sha-rigidity-windowed-local, which cites the published base change).
\end{proof}

\begin{lemma}\label{sha-siegel-step-ideal-descent-updown-local}
Assume the Galois representations as in Theorem~\ref{sha-residually-reducible-automorphy}. Let $\ell$ be a prime and $\iota:\bar{\mathbf Q}_\ell\cong\mathbf C$. For a finite place $u$ of a number field let $\varpi_u$ be a uniformizer and $q_u$ the order of the residue field.
\begin{enumerate}
\item Let $K$ be a CM field, $n\ge1$, and $r:\operatorname{Gal}(\bar{\mathbf Q}/K)\to GL_n(\bar{\mathbf Q}_\ell)$ a homomorphism. Suppose that there are a totally real subfield $F_0$ of $K$, a character $\mu$ of $\operatorname{Gal}(\bar{\mathbf Q}/F_0)$ attached through $\iota$ to an algebraic Hecke character of $F_0$ and taking the same value at all complex conjugations, a complex conjugation $c\in\operatorname{Gal}(\bar{\mathbf Q}/F_0)$ and $P\in GL_n(\bar{\mathbf Q}_\ell)$ with ${}^tr(s)\,P\,r(csc^{-1})=\mu(s)P$ for $s\in\operatorname{Gal}(\bar{\mathbf Q}/K)$. Then for every cuspidal automorphic representation $Q$ of $GL_n(\mathbf A_K)$ to which $r$ is attached through $\iota$ there is a Hecke character $\chi$ of the maximal totally real subfield of $K$ such that $(Q,\chi)$ is polarized (Definition~\ref{gal-polarized}).
\item Let $F\subseteq K$ be CM fields, $n\ge2$, $r:\operatorname{Gal}(\bar{\mathbf Q}/F)\to GL_n(\bar{\mathbf Q}_\ell)$ a continuous representation that satisfies condition (I) of Definition~\ref{sha-siegel-step-ideal-descent-conditions}, and $Q$ a regular algebraic cuspidal automorphic representation of $GL_n(\mathbf A_K)$, with $(Q,\chi)$ polarized for some $\chi$, to which the restriction of $r$ to $\operatorname{Gal}(\bar{\mathbf Q}/K)$ is attached through $\iota$. Let $v$ be a finite place of $F$ such that either $v\nmid\ell$, or $v\mid\ell$ and $r$ is potentially $\ell$-distinct at $v$ (Definition~\ref{sha-eigensystem-distinctness}). Let $w$ be a place of $K$ above $v$ and $t\in\mathbf R$ such that the supercuspidal support of every local component of $Q$ at $w$ consists of characters $\zeta$ with $q_w^{t-1/2}<|\zeta(\varpi_w)|<q_w^{t+1/2}$. Then for each such $\zeta$ there is a smooth character $\chi'$ of $F_v^\times$ such that $\zeta\cdot(\chi'\circ N_{K_w/F_v})$ is unramified.
\item Let $F\subset K$ be totally complex number fields, with $K$ Galois over $F$ of prime degree, $n\ge2$, and $\pi$ and $Q$ cuspidal automorphic representations of $GL_n(\mathbf A_F)$ and $GL_n(\mathbf A_K)$ such that $Q$ is a weak base change of $\pi$ (Definition~\ref{gal-depth-zero-lift-weak}). Let $t\in\mathbf R$ and $v$ a finite place of $F$ such that at every place $w$ of $K$ above $v$ the supercuspidal support of every local component of $Q$ at $w$ consists of characters $\zeta$ with $q_w^{t-1/2}<|\zeta(\varpi_w)|<q_w^{t+1/2}$, each of which has the property of the conclusion of (2). Then the supercuspidal support of every local component of $\pi$ at $v$ consists of characters $\xi$ with $q_v^{t-1/2}<|\xi(\varpi_v)|<q_v^{t+1/2}$.
\end{enumerate}
\end{lemma}
\leanref{Development/AutomorphicArithmetic/AutomorphicVarieties/ModularityArgument/ThresholdArguments/DenseThresholdCases/TAUpDownOf.lean}{UpperTheory.V.Z.Win.polarized_pair_of_attached}
\leanref{Development/AutomorphicArithmetic/AutomorphicVarieties/ModularityArgument/AutomorphyLifting/WeakOfAtt.lean}{UpperTheory.V.Z.Win.Thru}
\leanref{Development/AutomorphicArithmetic/AutomorphicVarieties/ModularityArgument/AutomorphyLifting/PolarizedLinkage/ThruAway.lean}{UpperTheory.V.Z.Win.thruAway}
\leanref{Development/AutomorphicArithmetic/AutomorphicVarieties/ModularityArgument/ThresholdArguments/DenseThresholdCases/OrdThruL.lean}{UpperTheory.V.Z.Ord.thruL_alt}
\leanref{Development/AutomorphicArithmetic/AutomorphicVarieties/ModularityArgument/GammaLaurentFactors/MatchIndOf.lean}{UpperTheory.V.Z.Gam.ThruAt}
\leanref{Development/AutomorphicArithmetic/AutomorphicVarieties/ModularityArgument/LeviSmoothLocalization/WinUpDown.lean}{UpperTheory.V.Z.ImaginaryQuadraticEtaCharacter.WinUpDown}
\begin{proof}
(1) Let $K^+$ be the maximal totally real subfield of $K$; it contains $F_0$, and $c$ lies in $\operatorname{Gal}(\bar{\mathbf Q}/K^+)$, as a complex conjugation fixes every totally real field.
\begin{enumerate}
\item[(a)] \emph{The polarization over $K^+$.} The cyclotomic character is attached through $\iota$ to an algebraic Hecke character of $\mathbf Q$ \cite{Ser68}, and the restriction of $\mu$ to $\operatorname{Gal}(\bar{\mathbf Q}/K^+)$ to the composite with the norm from $K^+$ to $F_0$ of the algebraic Hecke character of $\mu$ \cite[Chapter~VI]{Neu99}. So $\mu'=\chi_{\mathrm{cyc}}^{n-1}\mu$, on $\operatorname{Gal}(\bar{\mathbf Q}/K^+)$, is attached through $\iota$ to an algebraic Hecke character $\chi$ of $K^+$. It takes the same value at all complex conjugations, as $\mu$ does and $\chi_{\mathrm{cyc}}$ takes the value $-1$ at each of them. The relation reads ${}^tr(s)\,P\,r(csc^{-1})=\chi_{\mathrm{cyc}}(s)^{1-n}\mu'(s)P$ for $s\in\operatorname{Gal}(\bar{\mathbf Q}/K)$.
\item[(b)] \emph{Conjugate places.} Let $w$ be a finite place of $K$, not above $\ell$ and outside a finite set, such that $\chi\circ N_{K/K^+}$ is unramified at $w$, $Q$ has Satake parameters $B$ at $w$ and $B'$ at the conjugate place $w^c$, and $r$ has Frobenius eigenvalues $q_w^{(n-1)/2}B$ at $w$ and $q_w^{(n-1)/2}B'$ at $w^c$. Let $z$ be the value of $\chi\circ N_{K/K^+}$ at a uniformizer of $K_w$, which is the Frobenius eigenvalue at $w$ of the restriction of $\mu'$ to $\operatorname{Gal}(\bar{\mathbf Q}/K)$. Let $\sigma$ be a Frobenius element at a place $A$ above $w$; then $c\sigma c^{-1}$ is a Frobenius element at $c(A)$, which lies above $w^c$. The relation of (a) at $s=\sigma^{-1}$ gives $r(c\sigma^{-1}c^{-1})=\lambda P^{-1}\,{}^tr(\sigma^{-1})^{-1}P$ with $\lambda=\chi_{\mathrm{cyc}}(\sigma^{-1})^{1-n}\mu'(\sigma^{-1})$, and $\iota(\lambda)=q_w^{n-1}z$, as $\chi_{\mathrm{cyc}}(\sigma^{-1})=q_w^{-1}$. So the eigenvalues of $r(c\sigma^{-1}c^{-1})$ are the $\lambda/e$ for the eigenvalues $e$ of $r(\sigma^{-1})$, and $q_w^{(n-1)/2}B'=\{q_w^{n-1}z\,q_w^{-(n-1)/2}b^{-1}:b\in B\}$, that is, $B'=\{zb^{-1}:b\in B\}$.
\item[(c)] \emph{The pairing.} By (b), the composite $Q^c$ of $Q$ with $c$, which has Satake parameters $B'$ at $w$, and $Q^\vee\otimes(\chi\circ N_{K/K^+}\circ\det)$, which has Satake parameters $\{zb^{-1}:b\in B\}$ at $w$, have the same Satake parameters at all but finitely many places. The contragredient of $Q$ is realized on the functions $g\mapsto\phi({}^tg^{-1})$ for $\phi\in Q$ (Theorem~\ref{lit-bundle-unitary}(5)), so both are cuspidal automorphic representations, and by Theorem~\ref{lit-jacquet-shalika}(1) they are isomorphic. This gives the pairing of Definition~\ref{gal-polarized}(3) for $(Q,\chi)$.
\item[(d)] \emph{The signs.} The infinity type $(a_\tau)$ of $\chi$, indexed by the embeddings $\tau:K^+\to\mathbf R$, is parallel: as in the proof of Lemma~\ref{sha-satake-twist}, $\chi$ is trivial on the units of $K^+$, so $\sum_\tau a_\tau\log|\tau(\varepsilon)|=0$ for every unit $\varepsilon$, and by Dirichlet's unit theorem the $a_\tau$ are all equal to one integer $b$. For a real place $v$ of $K^+$, the global reciprocity map sends the idele that is $-1$ at $v$ and $1$ elsewhere to the class of a complex conjugation $c_v$, and as $\mu'$ is attached to $\chi$, $\mu'(c_v)=\iota^{-1}(\chi_v(-1))(-1)^b$. As $\mu'$ takes the same value at all complex conjugations, $\chi_v(-1)$ is the same at every real place $v$: $\chi_v(-1)=(-1)^m$ with $m\in\{0,1\}$.
\item[(e)] \emph{The sign condition.} Let $\eta$ be the quadratic character of $K/K^+$. Then $\eta\circ N_{K/K^+}=1$ and $\eta_v(-1)=-1$ at every real place $v$, so $(\chi\eta^{m+n})_v(-1)=(-1)^m(-1)^{m+n}=(-1)^n$ and $(\chi\eta^{m+n})\circ N_{K/K^+}=\chi\circ N_{K/K^+}$. So $(Q,\chi\eta^{m+n})$ is polarized.
\end{enumerate}

(2) Let $A$ be a valuation ring of $\bar{\mathbf Q}$ over $w$; it lies over $v$. Write $D_A(F)\supseteq D_A(K)$ and $I_A(F)\supseteq I_A(K)$ for the decomposition and inertia groups of $\operatorname{Gal}(\bar{\mathbf Q}/F)$ and $\operatorname{Gal}(\bar{\mathbf Q}/K)$ at $A$, $W_A(K)$ for the Weil group, and $r_K$ for the restriction of $r$. By Theorem~\ref{gal-llc-exists} there are a local Artin map $a$ of $K_w$ and a map $\operatorname{rc}$ that satisfies the laws with respect to $a$. Let $\sigma$ be a local component of $Q$ at $w$, with supercuspidal support $\{\xi_1,\dots,\xi_n\}$, and $\zeta=\xi_j$. Put $\xi'_i=\xi_i\nu^{(1-n)/2}$, which agrees with $\xi_i$ on $\mathcal O_w^\times$.
\begin{enumerate}
\item[(a)] \emph{The Weil-Deligne representation of $\sigma$.} By Lemma~\ref{sha-siegel-step-seeds-induced-automorphic-windows-plain}(1), $\sigma$ is plain and $y=\operatorname{rc}(\sigma\nu^{(1-n)/2})$ is abelian and monodromy-free. As $\sigma$ is a subquotient of $\mathrm{PS}(\xi')$ (Definition~\ref{sha-hecke-local-ordinary}), the characteristic polynomial of $y(s)$ is $\prod_i(X-\xi'_i(a(s)))$ for $s\in W_A(K)$, by Lemma~\ref{sha-hecke-local-rec}(1).
\item[(b)] \emph{Away from $\ell$.} Let $v\nmid\ell$. By Lemma~\ref{sha-hecke-ideal-attached}(1), applied to $Q$, which is regular algebraic and polarized, and $r_K$, $\mathrm{WD}_\iota(r_K|_{W_A(K)})=y$; so for $s\in I_A(K)$ the characteristic polynomial of $\iota r(s)$ is $\prod_i(X-\xi'_i(a(s)))$. By condition (I) at $A$ there are continuous characters $\kappa_1,\dots,\kappa_n$ of $D_A(F)$ with $\det(X-r(g))=\prod_k(X-\kappa_k(g))$ for $g\in I_A(F)$. Hence $\prod_k(X-\iota\kappa_k(s))=\prod_i(X-\xi'_i(a(s)))$ for $s\in I_A(K)$. Comparing the coefficients of $X^{n-1}$ at all $s$, the two families of characters of $I_A(K)$ have the same sums, and by the linear independence of distinct characters they agree up to order; so $\xi'_j\circ a=\iota\circ\kappa_k$ on $I_A(K)$ for some $k$. Let $\theta=\iota\circ\kappa_k$ on the Weil group $W_A(F)$. It agrees on $I_A(K)$, an open subgroup of $I_A(F)$, with $\xi'_j\circ a$, which has open kernel there, as $a$ is continuous on $I_A(K)$ and $\xi'_j$ is smooth. By (3) of Definition~\ref{gal-local-artin}, $\theta=\chi''\circ a_F$ for the local Artin map $a_F$ of $F_v$ and a smooth character $\chi''$ of $F_v^\times$. By (4) of that definition for $K$ and for $F$, and the compatibility of global reciprocity with the norm \cite[Chapter~VI]{Neu99}, $a_F(s)=N_{K_w/F_v}(a(s))$ for $s\in W_A(K)$, the reciprocity map being injective on the image of $F_v^\times$ as in the proof of Lemma~\ref{sha-hecke-local-rec}. For $x\in\mathcal O_w^\times$ write $x=a(s)$ with $s\in I_A(K)$; then $\zeta(x)=\xi'_j(x)=\theta(s)=\chi''(N_{K_w/F_v}(x))$. So $\chi'=(\chi'')^{-1}$ has $\zeta\cdot(\chi'\circ N_{K_w/F_v})$ unramified.
\item[(c)] \emph{Above $\ell$.} Let $v\mid\ell$, with $r$ potentially $\ell$-distinct at $v$. By Lemma~\ref{sha-hecke-ideal-attached}, $r_K|_{D_A(K)}$ is potentially semistable and de Rham, with $\mathrm{WD}_\iota(r_K|_{D_A(K)})=y$. As $r$ is potentially $\ell$-distinct at $v$, its restriction to $D_A(F)$ is conjugate to an upper triangular representation with diagonal characters $\kappa_1,\dots,\kappa_n$, so $r_K|_{D_A(K)}$ is conjugate to one with diagonal characters the restrictions $\psi_k$ of the $\kappa_k$. By Lemma~\ref{sha-siegel-step-ideal-descent-updown-local-above}(1) and~(2), each $\psi_k$ is locally algebraic with some $(\omega_k,h_k)$, and the characteristic polynomial of $y(s)$ is $\prod_k(X-\omega_k(a(s)))$ for $s\in W_A(K)$. Comparing with (a) as in (b), $\xi'_j\circ a=\omega_k\circ a$ on $W_A(K)$ for some $k$, so $\xi'_j=\omega_k$, as $a$ is surjective. By Lemma~\ref{sha-siegel-step-ideal-descent-updown-local-above}(3), applied to $\kappa_k$, there is a smooth character $\chi'$ of $F_v^\times$ with $\omega_k\cdot(\chi'\circ N_{K_w/F_v})$ unramified; as $\zeta$ agrees with $\xi'_j=\omega_k$ on $\mathcal O_w^\times$, so is $\zeta\cdot(\chi'\circ N_{K_w/F_v})$.
\end{enumerate}

(3) By Lemma~\ref{sha-siegel-step-ideal-descent-updown-local-down}(1), the supercuspidal support of every local component of $\pi$ at $v$ consists of characters. Let $\xi$ be one of them. By Lemma~\ref{sha-siegel-step-ideal-descent-updown-local-down}(2) there is a place $w$ of $K$ above $v$ such that $\xi\circ N_{K_w/F_v}$ lies in the supercuspidal support of a local component of $Q$ at $w$. Let $f$ be the residue degree of $w$ over $v$. Then $q_w=q_v^f$, and $N_{K_w/F_v}(\varpi_w)=u\varpi_v^f$ for a unit $u$ of $F_v$, with $|\xi(u)|=1$ as $\xi$ is continuous on the compact group of units; so $|\xi\circ N_{K_w/F_v}(\varpi_w)|=|\xi(\varpi_v)|^f$. The bound $q_v^{f(t-1/2)}<|\xi(\varpi_v)|^f<q_v^{f(t+1/2)}$ at $w$ gives, after taking $f$-th roots, $q_v^{t-1/2}<|\xi(\varpi_v)|<q_v^{t+1/2}$.
% (1) is polarized_pair_of_attached (TAUpDownOf.lean, section Pol, 25 lines), followed step by step: cyclotomicCharacter_plusPart (ToPolC, with Af.algOfCyc_s for the algebraicity and cyclotomic_char_conj_invariant, CycConj, for the sign) for (a); Af.self_dual_complex_of_printed_hida (TSelfDualC.lean, 129 lines: Bc.normGal, the private frob_dual with eig_dual and isFrob_conjGal) for (b), and polOfSat_of_thetaSat (theta_satake_parameters_inverse, Jacquet-Shalika strong multiplicity one and the Gelfand-Kazhdan description of the contragredient from HypothesesH', cuspContragredientIsCMConjTwist_of_paired_invTranspose) for (c); sign_at_infinity_of_parallel_weight (adelicSolution_parallel, heckeChar_parallel_weight; art_sign_eq and infPow_infIdele_neg_one_of_const_of_isReal) for (d); the twist by quadChar^(m + n) with compNorm_quadChar for (e).
% (2) is Thru false and Thru true. Thru false is thruAway plainRec (ThruAway.lean, 105 lines), followed step by step: recTExists, plainRec (Lemma sha-siegel-step-seeds-induced-automorphic-windows-plain (1)), plainTor_exists_chars, supp_nu_chars, charPol_of_chars and hasSupport_unique for (a); compatTOf (Lemma sha-hecke-ideal-attached) with charPol_wdOf, IneAb, the private exists_eq_of_prod (card_fiber_eq_of_sum_eq) and exists_char_of_ine (ThruLoc.lean: isSmooth_of_ine, IsLocalArtin.factor and localArtin_comp_norm_compat) for (b). Thru true is Z.Ord.thruL_alt (OrdThruL.lean, 220 lines), followed step by step for (c): compatTOf above l with Af.exists_bridge, PotUnl for the flag kappa, ord_of_flag and pol_of_flag (Lemma sha-siegel-step-ideal-descent-updown-local-above (1) and (2)), mem_of_prod, and ext_character_of_literal_automorphy_n (part (3) of that lemma).
% (3) is the down half of ImaginaryQuadraticEtaCharacter.winUpDown (PrincipalSeriesMatching/WinUpDown.lean, 40 lines): the down half of Z.Gam.upDown for the characters (Lemma sha-siegel-step-ideal-descent-updown-local-down (1)), the second half of upDownX (part (2) of that lemma), and normVal with the private win_pow_iff for the bound; the up half is Lemma sha-rigidity-windowed-local.
\end{proof}

\begin{lemma}\label{sha-siegel-step-ideal-descent-updown-auxiliary}
Let $K$ be a CM field and $N\ge0$ an integer. Then there is a CM field $M\supseteq K$, Galois over $K$, such that $[M:K]$ is a prime number greater than $N$.
\end{lemma}
\leanref{Development/AutomorphicArithmetic/AutomorphicVarieties/BuildingsResidueChars/SolvableCyclicSubextensionOverRat.lean}{UpperTheory.solvable_cyclic_subextension_over_rat}
\leanref{Development/AutomorphicArithmetic/AutomorphicVarieties/BuildingsResidueChars/SolvableAuxiliaryGaloisExtension.lean}{UpperTheory.solvable_auxiliary_galois_extension}
\leanref{Development/AutomorphicArithmetic/AutomorphicVarieties/BuildingsResidueChars/SolvableCmCompositumMaximalRealSubfield.lean}{UpperTheory.solvable_cm_compositum_maximalRealSubfield}
\begin{proof}
\begin{enumerate}
\item[(1)] \emph{A totally real field.} Let $p$ be a prime with $p\ge\max(N,[K:\mathbf Q])+3$; it is odd. By Dirichlet's theorem there is a prime $q\equiv1$ modulo $p$. The group $\operatorname{Gal}(\mathbf Q(\zeta_q)/\mathbf Q)$ is cyclic of order $q-1$, which $p$ divides; let $L$ be the fixed field of its subgroup of index $p$. Then $L$ is Galois over $\mathbf Q$ of degree $p$. As $p$ is odd, $L$ has a real place; as $\operatorname{Gal}(L/\mathbf Q)$ acts transitively on the infinite places of $L$ and preserves the real ones, $L$ is totally real.
\item[(2)] \emph{The compositum.} Embed $L$ in $\bar{\mathbf Q}$, let $\alpha$ be a primitive element of $L$ over $\mathbf Q$, $f$ its minimal polynomial, of degree $p$, and $M=K(\alpha)$. As $L$ is Galois over $\mathbf Q$, the roots of $f$ lie in $L$, so $M=KL$ is a splitting field of $f$ over $K$, and $M$ is Galois over $K$. Its degree $[M:K]$ is at most $p$, the degree of $f$. As $L\subseteq M$, $p$ divides $[M:\mathbf Q]=[M:K][K:\mathbf Q]$, and it does not divide $[K:\mathbf Q]<p$. So $[M:K]=p>N$.
\item[(3)] \emph{$M$ is CM.} As $K$ is CM, there is $\delta\in K$, $\delta\ne0$, with $\bar\delta=-\delta$ for the complex conjugation of $K$. Then $\varphi(\delta)$ is purely imaginary for every embedding $\varphi:M\to\mathbf C$, so $\delta^2$ lies in the maximal totally real subfield $M^+$ of $M$. Every element of $M$ is $a+b\delta$ with $a,b\in M^+$: this holds for the elements of $L$, which is totally real, and for $x\in K$, as $x=\frac{x+\bar x}2+\frac{x-\bar x}{2\delta}\delta$, and it is preserved by sums and products, as $\delta^2\in M^+$. If $s+t\delta=0$ with $s,t\in M^+$, then applying an embedding $\varphi$ and complex conjugation gives $\varphi(s)-\varphi(t)\varphi(\delta)=0$, so $s=t=0$. So $[M:M^+]=2$, and $M$, which is totally complex as it contains $K$, is a CM field.
\end{enumerate}
\end{proof}
% exists_cm_field_tower (TAUpDownOf.lean, private, 30 lines) with solvable_cyclic_subextension_over_rat (CycQ: a prime p >= N + 3 by Nat.exists_infinite_primes, q = 1 mod p by Nat.exists_prime_gt_modEq_one, the fixed field of zpowers (g^p) in CyclotomicField q, totally real through nrRealPlaces_pos_of_odd_finrank and exists_smul_eq_of_comap_eq) for (1); solvable_auxiliary_galois_extension (ExistsCompositumWithPrimeDegreeGalois: M = K(beta), the splitting field of the minimal polynomial, of degree [L:Q] by minpoly.degree_dvd and Nat.Prime.dvd_mul) for (2); solvable_cm_compositum_maximalRealSubfield (RCSup, with the private rcsup_cm_of_pair) for (3). The Lean states it for F and K with [K:F] = 1, and M is also Galois of degree p over F.

\begin{lemma}\label{sha-siegel-step-ideal-descent-updown-prime-index}
Let $\Gamma$ and $\Gamma'$ be groups, $k$ an algebraically closed field of characteristic $0$, $n\ge0$, $r:\Gamma\to GL_n(k)$ an irreducible representation, and $\phi:\Gamma'\to\Gamma$ a homomorphism whose image $H$ is a normal subgroup of $\Gamma$ of prime index $p>n$. Then $r\circ\phi$ is irreducible.
\end{lemma}
\leanref{Development/AutomorphicArithmetic/AutomorphicVarieties/BuildingsResidueChars/GaloisRepresentationIrreduciblePrimeDimension.lean}{UpperTheory.galois_representation_irreducible_prime_dimension}
\leanref{Development/AutomorphicArithmetic/AutomorphicVarieties/AlgebraTools/MonomialRepresentations/TwistOfRed.lean}{UpperTheory.V.AlgebraTools.twistOfRed}
\begin{proof}
The subspaces stable under $r\circ\phi$ are those stable under $r(H)$, so it suffices to show that the restriction of $r$ to $H$ is irreducible. As $r$ is irreducible, $n\ge1$. Suppose that the restriction is not irreducible. As $\Gamma/H$ is cyclic of order $p$, $H$ contains the commutators of $\Gamma$, and Lemma~\ref{sha-siegel-step-seeds-induced-irreducible-twist} gives a homomorphism $\eta:\Gamma\to k^\times$ with $\eta\ne1$ and $\eta(H)=1$, and $Q\in GL_n(k)$ with $\eta(g)Qr(g)=r(g)Q$ for $g\in\Gamma$. Taking determinants, $\eta(g)^n=1$; and $\eta(g)^p=\eta(g^p)=1$, as $g^p\in H$. As $p$ is a prime greater than $n\ge1$, $p$ and $n$ are coprime, so $\eta(g)=1$ for every $g$, which is false.
\end{proof}
% galois_representation_irreducible_prime_dimension (IrrPrime, 45 lines): isIrr_comp_of_range, hcomm (the quotient is cyclic of prime order), twistOfRed (TwistOfRed.lean, 140 lines: Clifford, a projection in the commutant of N, and a common eigenvector of the conjugation action; the same statement as exists_selfTwist of Lemma sha-siegel-step-seeds-induced-irreducible-twist, without eta^m = 1), then h1, h2 (the determinant) and Nat.coprime_of_lt_prime. For Galois groups, galois_representation_image_spec gives the image of Gal(Qbar/M) in Gal(Qbar/K) as a normal subgroup of index [M:K].

\begin{lemma}\label{sha-siegel-step-ideal-descent-updown}
Assume the Galois representations as in Theorem~\ref{sha-residually-reducible-automorphy}. Let $F$ be a CM field, $K\supseteq F$ a CM field Galois over $F$ with solvable Galois group, $\ell$ a prime, $\iota:\bar{\mathbf Q}_\ell\cong\mathbf C$, $n\ge1$, and $r:\operatorname{Gal}(\bar{\mathbf Q}/F)\to GL_n(\bar{\mathbf Q}_\ell)$ a continuous representation that is irreducible on $\operatorname{Gal}(\bar{\mathbf Q}/K)$, attached through $\iota$ to a regular algebraic cuspidal automorphic representation of $GL_n(\mathbf A_F)$, and satisfies condition (T) of Definition~\ref{sha-siegel-step-ideal-descent-conditions}. Suppose that there are a totally real subfield $F_0$ of $F$, a character $\mu$ of $\operatorname{Gal}(\bar{\mathbf Q}/F_0)$ attached through $\iota$ to an algebraic Hecke character of $F_0$ and taking the same value at all complex conjugations, a complex conjugation $c\in\operatorname{Gal}(\bar{\mathbf Q}/F_0)$ and $P\in GL_n(\bar{\mathbf Q}_\ell)$ with ${}^tr(s)\,P\,r(csc^{-1})=\mu(s)P$ for $s\in\operatorname{Gal}(\bar{\mathbf Q}/F)$. Then:
\begin{enumerate}
\item if $r$ is windowed (Theorem~\ref{sha-residually-reducible-automorphy}), so is its restriction to $\operatorname{Gal}(\bar{\mathbf Q}/K)$;
\item if $r$ satisfies condition (I) and keeps $\ell$-distinctness in $K$, and its restriction to $\operatorname{Gal}(\bar{\mathbf Q}/K)$ is windowed, then $r$ is windowed.
\end{enumerate}
\end{lemma}
\leanref{Development/AutomorphicArithmetic/AutomorphicVarieties/ModularityArgument/ThresholdArguments/DenseThresholdCases/EndsRankThreeOfLocal.lean}{UpperTheory.V.Z.Win.Dn.aUpDown_alt}
\leanref{Development/AutomorphicArithmetic/AutomorphicVarieties/ModularityArgument/ThresholdArguments/DenseThresholdCases/WinDefs.lean}{UpperTheory.V.Z.Win.AUpDownAlt}
\begin{proof}
For a field $E$ between $F$ and $K$ write $r_E$ for the restriction of $r$ to $\operatorname{Gal}(\bar{\mathbf Q}/E)$. Let $F^+$ be the maximal totally real subfield of $F$; it contains $F_0$, and $c$ lies in $\operatorname{Gal}(\bar{\mathbf Q}/F^+)$, as a complex conjugation fixes every totally real field.

(1) The restriction of $\mu$ to $\operatorname{Gal}(\bar{\mathbf Q}/F^+)$ is attached through $\iota$ to the composite with the norm from $F^+$ to $F_0$ of the algebraic Hecke character to which $\mu$ is attached \cite[Chapter~VI]{Neu99}, and it takes the same value at all complex conjugations. So the hypotheses (1) to (4) of Lemma~\ref{sha-rigidity-windowed-up} hold for $F\subseteq K$, $r$, this restriction, $c$ and $P$, and that lemma, with $E=K$, gives (1).

(2) The proof is by induction on $d=[K:F]$, for all fields $F\subseteq K$ as in the statement and all $\ell$, $\iota$, $n$ and $r$ at once, as in the proof of Lemma~\ref{sha-rigidity-windowed-up}, whose conventions are kept. Let $r$ be as in the statement, satisfy condition (I), keep $\ell$-distinctness in $K$, and have $r_K$ windowed.

\emph{Prime degree.} Suppose that $d$ is prime. For a finite place $w$ of $K$, $v$ denotes the place of $F$ below $w$.
\begin{enumerate}
\item[(d1)] \emph{Condition ($\alpha$) for $r$.} Let $v$ be a place of $F$ above $\ell$ at which $r$ is $\ell$-distinct, and $w$ a place of $K$ above $v$. As $r$ keeps $\ell$-distinctness in $K$, $r_K$ is $\ell$-distinct at $w$; by ($\alpha$) for $r_K$ it is potentially $\ell$-distinct at $w$; and by condition (T) and Lemma~\ref{sha-eigensystem-flags}(2), $r$ is potentially $\ell$-distinct at $v$.
\item[(d2)] \emph{Condition ($\beta$) for $n=1$.} As in step (p3) of the proof of Lemma~\ref{sha-rigidity-windowed-up}, Lemma~\ref{sha-hecke-ideal-plain}(1) gives ($\beta$) for $r$, with the bound at every finite place.
\item[(d3)] \emph{The data for $n\ge2$.} Let $n\ge2$ and $\pi$ a regular algebraic cuspidal automorphic representation of $GL_n(\mathbf A_F)$ attached through $\iota$ to $r$. By Theorem~\ref{sha-siegel-automorphy-base-change}, applied to $F\subseteq K$, $r$, the restriction of $\mu$ to $\operatorname{Gal}(\bar{\mathbf Q}/F^+)$ of (1), $c$ and $P$, the representation $r_K$ is attached through $\iota$ to a regular algebraic cuspidal automorphic representation $Q$ of $GL_n(\mathbf A_K)$. As $F_0\subseteq K$ and the relation between $r$, $\mu$, $c$ and $P$ holds on $\operatorname{Gal}(\bar{\mathbf Q}/K)$, Lemma~\ref{sha-siegel-step-ideal-descent-updown-local}(1), applied to $r_K$, $F_0$, $\mu$, $c$, $P$ and $Q$, gives a Hecke character $\chi$ of the maximal totally real subfield of $K$ with $(Q,\chi)$ polarized (Definition~\ref{gal-polarized}). By step (p4) of the proof of Lemma~\ref{sha-rigidity-windowed-up}, $Q$ is a weak base change of $\pi$. By ($\beta$) for $r_K$, applied to $Q$, there is $t\in\mathbf R$ such that at every finite place $w$ of $K$ with $w\nmid\ell$ or $r_K$ potentially $\ell$-distinct at $w$, the supercuspidal support of $Q_w$ consists of characters $\zeta$ with $q_w^{t-1/2}<|\zeta(\varpi_w)|<q_w^{t+1/2}$.
\item[(d4)] \emph{Condition ($\beta$) for $n\ge2$.} Let $v$ be a finite place of $F$ with $v\nmid\ell$ or $r$ potentially $\ell$-distinct at $v$, and $w$ a place of $K$ above $v$. Then $w\nmid\ell$, or $r_K$ is potentially $\ell$-distinct at $w$ by Lemma~\ref{sha-eigensystem-flags}(1); so the bound of (d3) holds at $w$. By Lemma~\ref{sha-siegel-step-ideal-descent-updown-local}(2), applied to $F\subseteq K$, $r$, which satisfies condition (I), $Q$, $v$, $w$ and $t$, every character $\zeta$ in the supercuspidal support of $Q_w$ is the product of an unramified character and of $\chi'\circ N_{K_w/F_v}$ for a smooth character $\chi'$ of $F_v^\times$. This holds at every $w$ above $v$, so by Lemma~\ref{sha-siegel-step-ideal-descent-updown-local}(3) the supercuspidal support of $\pi_v$ consists of characters $\xi$ with $q_v^{t-1/2}<|\xi(\varpi_v)|<q_v^{t+1/2}$. This is ($\beta$) for $r$, with $\pi$ and $t$.
\end{enumerate}

\emph{Degree one.} Suppose that $d=1$, so that $K=F$.
\begin{enumerate}
\item[(o1)] \emph{Condition ($\alpha$) for $r$, and ($\beta$) for $n=1$.} These follow as in (d1) and (d2).
\item[(o2)] \emph{An auxiliary field.} Let $n\ge2$. By Lemma~\ref{sha-siegel-step-ideal-descent-updown-auxiliary} there is a CM field $M\supseteq K$, Galois over $K$, of prime degree $p>n$; as $K=F$, it is Galois over $F$ of prime degree, with solvable Galois group. The group $\operatorname{Gal}(\bar{\mathbf Q}/M)$ is a normal subgroup of index $p$ of $\operatorname{Gal}(\bar{\mathbf Q}/K)$, and $r_K$ is irreducible, so $r_M$ is irreducible by Lemma~\ref{sha-siegel-step-ideal-descent-updown-prime-index}. By Theorem~\ref{sha-siegel-automorphy-base-change}, applied to $F\subseteq M$, $r$, the restriction of $\mu$ to $\operatorname{Gal}(\bar{\mathbf Q}/F^+)$ of (1), $c$ and $P$, the representation $r_M$ is attached through $\iota$ to a regular algebraic cuspidal automorphic representation of $GL_n(\mathbf A_M)$; and $r_K=r$ is attached through $\iota$ to one of $GL_n(\mathbf A_K)$.
\item[(o3)] \emph{Up from $K$ to $M$.} The representation $r_K$ satisfies condition (T), is attached through $\iota$ to a regular algebraic cuspidal automorphic representation by (o2), and satisfies ($\beta$). Steps (p2), (p4) and (p5) of the proof of Lemma~\ref{sha-rigidity-windowed-up}, applied to $K\subseteq M$ and $r_K$, show that $r_M$ satisfies ($\beta$).
\item[(o4)] \emph{Down from $M$ to $F$.} Steps (d3) and (d4), applied to $F\subseteq M$, which has prime degree, with $M$ in place of $K$ and ($\beta$) for $r_M$ from (o3), give ($\beta$) for $r$.
\end{enumerate}
As $K=F$ and $r_K=r$, (2) also holds in this case without the auxiliary field.

\emph{Other degrees.} Suppose that $d$ is neither $1$ nor prime. Step (t1) of the proof of Lemma~\ref{sha-rigidity-windowed-up} gives a CM field $E$ with $F\subsetneq E\subsetneq K$, Galois over $F$, with $K$ Galois over $E$, both with solvable Galois groups and of degrees less than $d$.
\begin{enumerate}
\item[(e1)] \emph{The hypotheses for $E\subseteq K$.} As in step (t3) of that proof, $r_E$ satisfies the hypotheses of the statement for $E\subseteq K$: it is irreducible on $\operatorname{Gal}(\bar{\mathbf Q}/K)$, polarized, attached through $\iota$ to a regular algebraic cuspidal automorphic representation of $GL_n(\mathbf A_E)$ by Theorem~\ref{sha-siegel-automorphy-base-change}, and satisfies condition (T) by Lemma~\ref{sha-rigidity-windowed-flags}(1). It satisfies condition (I), with the restrictions of the characters to the decomposition groups of $\operatorname{Gal}(\bar{\mathbf Q}/E)$.
\item[(e2)] \emph{$r_E$ keeps $\ell$-distinctness in $K$.} Let $u$ be a place of $K$ above $\ell$ such that $r_E$ is $\ell$-distinct at the place of $E$ below $u$. By condition (T) for $r$ and Lemma~\ref{sha-rigidity-windowed-flags}(2), applied to $F\subseteq E$, $r$ is $\ell$-distinct at the place of $F$ below $u$. As $r$ keeps $\ell$-distinctness in $K$, $r_K$ is $\ell$-distinct at $u$.
\item[(e3)] \emph{$r$ keeps $\ell$-distinctness in $E$.} Let $w$ be a place of $E$ above $\ell$ such that $r$ is $\ell$-distinct at the place of $F$ below $w$, and $u$ a place of $K$ above $w$. As $r$ keeps $\ell$-distinctness in $K$, $r_K$ is $\ell$-distinct at $u$. By (e1) and Lemma~\ref{sha-rigidity-windowed-flags}(2), applied to $E\subseteq K$ and $r_E$, $r_E$ is $\ell$-distinct at $w$.
\item[(e4)] \emph{Conclusion.} By (e1), (e2) and the statement for $E\subseteq K$, $r_E$ is windowed. By (e3) and the statement for $F\subseteq E$, $r$ is windowed.
\end{enumerate}
% Part (2) follows the down halves of aUpDownOf_alt (TAUpDownOf.lean, 360 lines): the private good_all, a strong induction on [K:F], with good_prime in prime degree and solvable_cm_extension_tower with good_trans otherwise. (d1) is first_condition_down (Keeps, then pot_down of FlagLaws, Lemma sha-eigensystem-flags (2)); (d2) is second_condition_trivial_extension (win_one); (d3) is second_condition_down with aut_up (Z.Bc.solC_alt), polarized_pair_of_attached (Lemma sha-siegel-step-ideal-descent-updown-local (1)) and weakOfAtt; (d4) is win_down with cond_up (pot_up), Thru false or Thru true (thruAway plainRec and Z.Ord.thruL_alt, Lemma sha-siegel-step-ideal-descent-updown-local (2)) and the down half of WinUpDown (part (3) of that lemma). (e1) to (e4) are the down half of good_trans: flagLaws (ineAb_res, flagged_res, unl_down) for (e1) to (e3), and d1.2 with d2.2 for (e4). Degree one is the down half of good_one: (o1) is first_condition_down and second_condition_trivial_extension; (o2) is exists_cm_field_tower (Lemma sha-siegel-step-ideal-descent-updown-auxiliary, with N = n), galois_representation_image_spec and galois_representation_irreducible_prime_dimension (Lemma sha-siegel-step-ideal-descent-updown-prime-index), and aut_up through Z.Bc.stepC_alt for F in M, with aut_up through solC_alt for haK; (o3) is second_condition_up for K in M, with flagged_res for Flagged r_K; (o4) is second_condition_down for F in M, with solvable_polarized_restriction. All of this is second_condition_down_trivial. In the Lean K is a field with an embedding of F of degree one, which it cannot identify with F; the last sentence of the degree-one case is not in the Lean.
% aUpDown_alt = Z.Win.aUpDownOf_alt (TAUpDownOf.lean, 360 lines) with weakOfAtt, flagLaws (FlagLaws.lean, 398 lines), thruAway plainRec, Z.Ord.thruL_alt (OrdThruL.lean), ImaginaryQuadraticEtaCharacter.winUpDown and winOf, and ends_rank_three_of_local (RowT, from the field litN.aut.gal of Af.HypothesesH'). The hypotheses: Sol.Pol ι r (F_0, μ with IsAlgChar and SameSign, c with IsConjEl, P with IsPairing), IsRegAlgAutomorphicAttached ι r, Flagged r; (2) asks IneAb r and Keeps r K. Down'_alt is (2) and automorphy_up_base_change is (1).
\end{proof}

\begin{lemma}\label{sha-siegel-step-ideal-descent-across}
Assume the Galois representations as in Theorem~\ref{sha-residually-reducible-automorphy}. Let $\ell$ be a prime, $\iota:\bar{\mathbf Q}_\ell\cong\mathbf C$, $n\ge1$, $G=\operatorname{Gal}(\bar{\mathbf Q}/\mathbf Q)$ and $r:G\to GL_n(\bar{\mathbf Q}_\ell)$ a continuous representation. Let $K$ be an imaginary quadratic field fixed by the decomposition group of $G$ at every place of $\bar{\mathbf Q}$ above $\ell$. Suppose that:
\begin{enumerate}
\item[(i)] for every imaginary quadratic field $E$ such that $r$ is irreducible on $G_E$, the restriction of $r$ to $G_E$ is attached through $\iota$ to a regular algebraic cuspidal automorphic representation of $GL_n(\mathbf A_E)$;
\item[(ii)] there are a character $\theta$ of $G$ attached through $\iota$ to an algebraic Hecke character of $\mathbf Q$ and $J\in GL_n(\bar{\mathbf Q}_\ell)$ with ${}^tr(s)Jr(s)=\theta(s)J$ for $s\in G$;
\item[(iii)] $r$ satisfies conditions (I) and (T) of Definition~\ref{sha-siegel-step-ideal-descent-conditions};
\item[(iv)] $r$ is irreducible on $G_K$, and its restriction to $G_K$ is windowed.
\end{enumerate}
Then for every imaginary quadratic field $K'$ such that $r$ is irreducible on $G_{K'}$, the restriction of $r$ to $G_{K'}$ is windowed.
\end{lemma}
\leanref{Development/AutomorphicArithmetic/AutomorphicVarieties/ModularityArgument/ThresholdArguments/DenseThresholdCases/EndsRankThreeOfLocal.lean}{UpperTheory.V.Z.Win.Dn.across'_alt}
\leanref{Development/AutomorphicArithmetic/AutomorphicVarieties/ModularityArgument/ThresholdArguments/DenseThresholdCases/TAAcrossOf.lean}{UpperTheory.V.Z.Win.a_across_of_statement}
\begin{proof}
\begin{enumerate}
\item[(a)] \emph{Restrictions.} For a number field $E$, the restriction of $r$ to $G_E$ satisfies conditions (I) and (T): restrict the characters to the decomposition groups of $G_E$, which lie in those of $G$, with the inertia groups. Let $c\in G$ be a complex conjugation and $P=Jr(c)^{-1}$; then ${}^tr(s)Pr(csc^{-1})=\theta(s)P$ for $s\in G$, and $\theta$ takes the same value at all complex conjugations, which are conjugate in $G$.
\item[(b)] \emph{A step.} Let $E$ and $E'$ be imaginary quadratic fields such that $E$ is fixed by the decomposition groups of $G$ above $\ell$, $r$ is irreducible on $G_E\cap G_{E'}$, and the restriction of $r$ to $G_E$ is windowed. Then so is its restriction to $G_{E'}$. Indeed, let $M=EE'$, a CM field, Galois over $E$ and over $E'$ with abelian Galois groups, with $G_M=G_E\cap G_{E'}$. As $r$ is irreducible on $G_M$, it is irreducible on $G_E$ and $G_{E'}$, so by (i) its restrictions to $G_E$ and $G_{E'}$ are attached through $\iota$ to regular algebraic cuspidal automorphic representations. By (a), Lemma~\ref{sha-siegel-step-ideal-descent-updown} applies to $F=E$, $K=M$, and to $F=E'$, $K=M$, with $F_0=\mathbf Q$, $\mu=\theta$, $c$ and $P$. By its part (1), the restriction of $r$ to $G_M$ is windowed. The restriction of $r$ to $G_{E'}$ keeps $\ell$-distinctness in $M$: for a place $A$ above $\ell$, the decomposition group of $G_{E'}$ at $A$ fixes $E$, so it is the decomposition group of $G_M$ at $A$, and the triangular forms and characters on it are the same. By part (2), with (a), the restriction of $r$ to $G_{E'}$ is windowed.
\item[(c)] \emph{Conclusion.} Let $K'$ be an imaginary quadratic field such that $r$ is irreducible on $G_{K'}$. Lemma~\ref{sha-siegel-step-ideal-descent-split} gives an imaginary quadratic field $K''$, fixed by the decomposition groups above $\ell$, such that $r$ is irreducible on $G_K\cap G_{K''}$ and on $G_{K''}\cap G_{K'}$. By (b), applied to $(E,E')=(K,K'')$ and then to $(K'',K')$, the restriction of $r$ to $G_{K'}$ is windowed.
\end{enumerate}
\end{proof}
% across'_alt = Z.Win.Dn.aAcross_alt = a_across_of_statement (TAAcrossOf.lean, 101 lines) with flagLaws, aUpDown_alt and isoTw_alt: thirdSplit (Lemma sha-siegel-step-ideal-descent-split) for (c); the private step for (b), with Z.Bc.biQuad for M, the private across_irr, across_pair (P = J r(c^-1)) and across_cc of BaseChangeUnits/CentralResidueGroups/Across.lean, auSign, the two halves of aUpDown_alt (Lemma sha-siegel-step-ideal-descent-updown) and the private keeps_of_split; flagLaws (ineAb_res, flagged_res) for (a).

\begin{lemma}\label{sha-siegel-step-ideal-descent-laws}
Let $D$ be a Siegel datum, $T$ a fit tower for $D$ with fields $K$, $F$ and $L$ (Definition~\ref{sha-siegel-tower-towers}), $\mathcal O'\supseteq\mathcal O$ a ring of Definition~\ref{sha-character-parts}, $X=(A,B,P_X)$ a deformation space for $T$ over $\mathcal O'$ (Definition~\ref{sha-siegel-tower-space}), with $P_X=(r_X,(n_{X,v}))$, and $m=(k+1,k,1,0)$ for the integer $k$ of $D$. Let $y:A\to\bar{\mathbf Q}_\ell$ be a homomorphism of $\mathcal O'$-algebras that carries $P_X$ to a flagged pair with inertial weights $2m$ (Definition~\ref{sha-siegel-tower-ideal}), and $r:G\to GL_4(\bar{\mathbf Q}_\ell)$ a homomorphism whose restriction to $G_F$ is $y\circ r_X$. Then $r$ satisfies conditions (I) and (T) of Definition~\ref{sha-siegel-step-ideal-descent-conditions}, and its restriction to $G_L$ is $\ell$-distinct at every place of $L$ above $\ell$ (Definition~\ref{sha-eigensystem-distinctness}).
\end{lemma}
\leanref{Development/AutomorphicArithmetic/AutomorphicVarieties/ModularityArgument/UpperTriangularPlaces/PtLaws.lean}{UpperTheory.V.Keep.ptLaws}
\leanref{Development/AutomorphicArithmetic/AutomorphicVarieties/ModularityArgument/UpperTriangularPlaces/PtLawsOf.lean}{UpperTheory.V.Keep.PtLaws}
\begin{proof}
Let $P=(r_P,(n_v))$ be the image of $P_X$ under $y$, so that $r_P$ is the restriction of $r$ to $G_F$, and write $\psi_{v,i}(s)$ for the diagonal entries of $n_v^*r(s)n_v$, $s\in D_v(F)$. Let $(\chi_1,\dots,\chi_4)$ be the relabelling characters of $T$; they are $(\delta_1,\dots,\delta_4)=(\theta,\Psi,\alpha,1)$ or $(\delta_3,\delta_4,\delta_1,\delta_2)$ (Definition~\ref{sha-siegel-datum}). Below, $A$ is a place of $\bar{\mathbf Q}$ above $\ell$.
\begin{enumerate}
\item[(a)] \emph{Inertia of infinite cyclotomic order.} For a number field $E$ and an open subgroup $U$ of $G_E$, some $s\in U\cap I_A(E)$ has $\chi_{\mathrm{cyc}}(s)$ of infinite order. Indeed, $\chi_{\mathrm{cyc}}$ maps $I_A(\mathbf Q)$ onto $\mathbf Z_\ell^\times$, as in step (a) of the proof of Lemma~\ref{sha-siegel-step-ideal-seeds-inputs-tangent-crystalline}; let $\sigma\in I_A(\mathbf Q)$ with $\chi_{\mathrm{cyc}}(\sigma)=1+\ell$, of infinite order. As $U$ is open in $G$, $\sigma^e\in U$ for some $e\ge1$, and $s=\sigma^e$ lies in $U\cap I_A(E)$.
\item[(b)] \emph{Triangular forms over $F$.} There are characters $\kappa_1,\dots,\kappa_4$ of $D_A(F)$ and an open subgroup $U_A$ of $G_F$ such that $r$ is triangular on $D_A(F)$ with characters $\kappa_i$ (Definition~\ref{sha-siegel-step-ideal-descent-conditions}), $|\kappa_i(s)-\chi_i(s)|<1$ for $s\in D_A(F)$, and $\kappa_i(s)^2=\chi_{\mathrm{cyc}}(s)^{2m_i}$ for $s\in U_A\cap I_A(F)$. Indeed, by Definition~\ref{sha-siegel-tower-towers}(3), $A=g(v)$ for some $v\in\Sigma_T$ and $g\in G_F$, so that $D_A(F)=gD_v(F)g^{-1}$ and $I_A(F)=gI_v(F)g^{-1}$. As $n_v$ is symplectic with multiplier $1$, $n_v^*n_v=1$, and $u=n_v^*r(g^{-1})$ is invertible, with inverse $r(g)n_v$. For $s\in D_A(F)$, $ur(s)u^{-1}=n_v^*r(g^{-1}sg)n_v$ is upper triangular with diagonal entries $\kappa_i(s)=\psi_{v,i}(g^{-1}sg)$; the $\kappa_i$ are characters, as the diagonal entries of a product of upper triangular matrices are the products of theirs. By (3) of Definition~\ref{sha-siegel-tower-space}, $\psi_{v,i}(g^{-1}sg)$ is the image under $y$ of an element of $B$ with residue $\bar\chi_i(g^{-1}sg)=\bar\chi_i(s)$, so $|\kappa_i(s)-\chi_i(s)|<1$ by Lemma~\ref{sha-automorphy-unramified-near-tangent-points}, applied to the restriction of $y$ to $B$. As $P$ has inertial weights $2m$ (Definition~\ref{sha-siegel-tower-ideal}), there is an open subgroup $U$ of $G_F$ with $\psi_{v,i}(s')^2=\chi_{\mathrm{cyc}}(s')^{2m_i}$ for $s'\in U\cap I_v(F)$; with $U_A=gUg^{-1}$, $\kappa_i(s)^2=\chi_{\mathrm{cyc}}(g^{-1}sg)^{2m_i}=\chi_{\mathrm{cyc}}(s)^{2m_i}$ for $s\in U_A\cap I_A(F)$.
\item[(c)] \emph{Distinctness over $F$.} For $\kappa_i$ as in (b), $\kappa_i\ne\kappa_j\chi_{\mathrm{cyc}}$ on $D_A(F)$ for all $i$ and $j$. Suppose otherwise. As $|\kappa_i-\chi_i|<1$, $|\kappa_j-\chi_j|<1$ and $\bar\chi_{\mathrm{cyc}}=\bar\omega$, $\bar\chi_i=\bar\chi_j\bar\omega$ on $D_A(F)$. By condition (m) of Definition~\ref{sha-siegel-tower-towers}(4) there are two cases.
\begin{enumerate}
\item[(c1)] Some $g\in I_A(F)$ has $\bar\omega(g)\ne1$. By condition (k), $\bar r_0(g)=1$, so $\bar\delta_1(g)=\dots=\bar\delta_4(g)=1$, and $\bar\chi_i(g)=\bar\chi_j(g)=1$. Then $\bar\omega(g)=1$, a contradiction.
\item[(c2)] Some $g\in D_A(F)$ has $\bar\omega(g)=1$, $\bar\chi^{(0,1)}(g)=1$ and $\bar\chi^{(1,0)}(g)\ne1$. By (a), with $E=F$ and $U=U_A$, some $s\in U_A\cap I_A(F)$ has $\chi_{\mathrm{cyc}}(s)$ of infinite order; squaring $\kappa_i(s)=\kappa_j(s)\chi_{\mathrm{cyc}}(s)$ gives $\chi_{\mathrm{cyc}}(s)^{2m_i}=\chi_{\mathrm{cyc}}(s)^{2m_j+2}$, so $m_i=m_j+1$. As $m=(k+1,k,1,0)$ with $k\ge2$, $(i,j)$ is $(1,2)$, $(2,3)$ or $(3,4)$. For both relabellings, $\chi^{(0,1)}=\chi_1\chi_3^{-1}$ is $\Psi^{\pm1}$ and $\chi^{(1,0)}=\chi_1\chi_2^{-1}$ is $\alpha$; so $\bar\Psi(g)=1$, $\bar\alpha(g)\ne1$, and $(\bar\chi_1(g),\dots,\bar\chi_4(g))=(\bar\alpha(g),1,\bar\alpha(g),1)$, as $\theta=\alpha\Psi$. For each of the three pairs, exactly one of $\bar\chi_i(g)$ and $\bar\chi_j(g)$ is $\bar\alpha(g)\ne1$, which contradicts $\bar\chi_i(g)=\bar\chi_j(g)\bar\omega(g)=\bar\chi_j(g)$.
\end{enumerate}
\item[(d)] \emph{$\ell$-distinctness over $L$.} By condition (d) of Definition~\ref{sha-siegel-tower-towers}(4), as $\ell\in S$, $D_A(\mathbf Q)\subseteq G_K$, so $D_A(F)\subseteq G_F\cap G_K=G_L$ and $D_A(L)=D_A(F)$. By (b) and (c), the restriction of $r$ to $G_L$ is triangular on $D_A(L)$ with characters $\kappa_i$ such that $\kappa_i\ne\kappa_j\chi_{\mathrm{cyc}}$ for all $i$ and $j$. As this holds at every $A$, the restriction of $r$ to $G_L$ is $\ell$-distinct at every place of $L$ above $\ell$.
\item[(e)] \emph{Condition (T).} Let $a\in GL_4(\bar{\mathbf Q}_\ell)$ conjugate $r$ on $D_A(F)$ to upper triangular form with diagonal characters $\kappa_i$, as in (b). Let $N\subseteq U_A$ be an open normal subgroup of $G$, which exists as $U_A$ is open in the profinite group $G$. By (a), with $E=\mathbf Q$ and $U=N$, some $s\in N\cap I_A(\mathbf Q)$ has $\chi_{\mathrm{cyc}}(s)$ of infinite order. Let $g\in D_A(\mathbf Q)$. Then $s$ and $gsg^{-1}$ lie in $N\cap I_A(\mathbf Q)\subseteq U_A\cap I_A(F)$, and $\chi_{\mathrm{cyc}}(gsg^{-1})=\chi_{\mathrm{cyc}}(s)$. So $X=ar(s)a^{-1}$ and $Y=ar(gsg^{-1})a^{-1}$ are upper triangular, and by (b) the diagonal entries of $X^2$ and $Y^2$ are both $d_i=\chi_{\mathrm{cyc}}(s)^{2m_i}$, which are pairwise distinct. The matrix $Z=ar(g)a^{-1}$ satisfies $ZX^2=Y^2Z$. For $j<i$, by induction on $j+(4-i)$, the entries $Z_{i'j'}$ with $j'<i'$ and $j'+(4-i')<j+(4-i)$ vanish, so the entry $(i,j)$ of $ZX^2$ is $Z_{ij}d_j$ and that of $Y^2Z$ is $d_iZ_{ij}$; as $d_i\ne d_j$, $Z_{ij}=0$. So $ar(g)a^{-1}$ is upper triangular for every $g\in D_A(\mathbf Q)$, and its diagonal entries are characters of $D_A(\mathbf Q)$. This is condition (T).
\item[(f)] \emph{Condition (I).} Let $H=G_F\cap\ker\bar r_0$, an open subgroup of $G$ that contains the commutator subgroup, as $F$ is abelian over $\mathbf Q$ and $\bar r_0$ is diagonal. Let $A'$ be a place of $\bar{\mathbf Q}$ above a prime other than $\ell$, and $g\in I_{A'}(\mathbf Q)\cap H$. Then $g\in I_{A'}(F)$ and $\bar r_0(g)=1$, so $g$ lies in the set $N_\Phi$ of the problem of $T$ (Definition~\ref{sha-siegel-tower-pairs}(1)), and $r(g)=y(r_X(g))=1$. So for $a,b\in I_{A'}(\mathbf Q)$, the commutator $aba^{-1}b^{-1}$, which lies in $I_{A'}(\mathbf Q)\cap H$, has $r(aba^{-1}b^{-1})=1$: the $r(a)$ commute. Commuting matrices over $\bar{\mathbf Q}_\ell$ can be put in upper triangular form simultaneously, which gives characters $\chi'_1,\dots,\chi'_4$ of $I_{A'}(\mathbf Q)$ with $\det(X-r(a))=\prod_i(X-\chi'_i(a))$. Each $\chi'_i$ is trivial on $I_{A'}(\mathbf Q)\cap H$, where $r=1$, so it is a character of the image of $I_{A'}(\mathbf Q)$ in the finite abelian group $G/H$, and it extends to a character $\psi_i$ of $G/H$, as $\bar{\mathbf Q}_\ell^\times$ is divisible. The restrictions of the $\psi_i$ to $D_{A'}(\mathbf Q)$ are continuous, as $H$ is open, and they give condition (I) at $A'$.
\end{enumerate}
% Keep.ptLaws = ptLaws_of (PtLawsOf.lean, 80 lines with the statements of its eight inputs), followed step by step, with the inputs in UpperTriangularPlaces/NormalValuationBounds (about 780 lines): (a) is infCyc (InfCyc.lean, through Z.Sym.Ga.ineInf); (b) is flagF (ValLtOneIff.lean, 193 lines: T.cover, conj_dec, conj_ine, flag_abs, X.ψ_red with coPt for the residues, and the conjugated U for the weights), whose conclusion is the structure OrdAt; (c) is unlF (NeOfPair.lean, 115 lines: res_step, bar_chi, LocL with Tower.Ok.ineL, exp_step and pairs, ne_of_pair); (d) is unlL (MemDec.lean, 44 lines: Tower.Ok.split with D.Good.l_mem, mem_dec); (e) is flagQ (ExistsNormal.lean, 134 lines: exists_normal, flag_of, tri_step with upTri, UpTri.lean, 39 lines); (f) is ineQ (CommLeRbar.lean, 44 lines: comm_le_fix, comm_le_rbar, the field kill of the flagged pair) with ineOf (CharpolyScalar.lean, 197 lines: common_eig, chars_of_comm, chars_of_gl, extend_char, descend). The hypothesis LocL T is condition (m) of the fit tower. The rest of UpperTriangularPlaces (BigOm, ConjLeftPos, CrOnTri and others) is not used here.
% Keep.ptLaws = ptLaws_of with flagF, infCyc, unlF, unlL, upTri, flagQ, ineOf and ineQ (ModularityArgument/UpperTriangularPlaces, about 2,300 lines); the hypotheses D.Good, T.Ok and LocL T come from the fit tower.
\end{proof}

\begin{lemma}\label{sha-siegel-step-ideal-descent}
Assume the Galois representations as in Theorem~\ref{sha-residually-reducible-automorphy}. Let $D$, $T$, $\iota$, $\mathcal O'$ and $X$ be as in Situation~\ref{sha-siegel-step-ideal-setting}, and let two links of $T$ over $\mathcal O'$ be given with the same character $\kappa$, whose places $v_0$ and $v'_0$ lie above distinct primes. Let $y:A\to\bar{\mathbf Q}_\ell$ be a homomorphism of $\mathcal O'$-algebras that carries $P_X$ to a flagged pair with inertial weights $2m$, and $r:G\to M_4(\bar{\mathbf Q}_\ell)$ a continuous homomorphism whose restriction to $G_F$ is $y\circ r_X$, with $r(g)$ symplectic with multiplier $\theta(g)$ for all $g$, and irreducible on $G_L$. Suppose that the conclusion of Lemma~\ref{sha-siegel-step-ideal-away} holds for both links. Then $r$ is automorphic through $\iota$ (Definition~\ref{sha-siegel-datum}).
\end{lemma}
\leanref{Development/AutomorphicArithmetic/Foundations/GaloisRepresentations/PlainTorsionPairs/IrrArithPointsAutomorphicAway.lean}{UpperTheory.Gal.projectedFormsSystemPart.glueN}
\leanref{Development/AutomorphicArithmetic/AutomorphicVarieties/ModularityArgument/ThresholdArguments/DenseThresholdCases/DnPmDown.lean}{UpperTheory.V.Z.Win.Dn.pmDown_alt}
\begin{proof}
Let $\rho_L(g)=\kappa(g)r(g)$ for $g\in G_L$.
\begin{enumerate}
\item[(a)] \emph{Gluing.} By hypothesis, $\rho_L$ is attached through $\iota$ to a regular algebraic cuspidal automorphic representation $\pi$ of $GL_4(\mathbf A_L)$ with $\pi^\vee\cong\pi^c$, and it is windowed at the places not above $v_0$ and at the places not above $v'_0$; in particular condition ($\alpha$) holds. Let $\pi'$ be a regular algebraic cuspidal automorphic representation attached to $\rho_L$ through $\iota$, and let $s$ and $s'$ be the exponents of ($\beta$) for $\pi'$ given by the two links. Let $w$ be a finite place of $L$ with $w\nmid\ell$ or satisfying the second condition of ($\alpha$). As $v_0$ and $v'_0$ lie above distinct primes, $w$ is not above one of them, so the bound of ($\beta$) holds at $w$ with $s$ or with $s'$. Then the supercuspidal support of $\pi'_w$ consists of characters, and the absolute values of their values at $\varpi_w$ have ratios strictly between $q_w^{-1}$ and $q_w$; so no two of them are linked, and $\pi'$ is plain at $w$ (Definition~\ref{gal-plain}). By Lemma~\ref{sha-hecke-ideal-plain}(1) there is $t\in\mathbf R$ such that the bound of ($\beta$) holds with $t$ at every finite place at which $\pi'$ is plain, in particular at every such $w$. So $\rho_L$ is windowed.
\item[(b)] \emph{Untwisting.} By (e) of Definition~\ref{sha-siegel-automorphy-link}, $\kappa$ is a power of $\chi_{\mathrm{cyc}}$ times a character of finite order on the inertia group at every place above $\ell$, so it is the character attached through $\iota$ to an algebraic Hecke character $\chi_\kappa$ of $L$ \cite{Ser68}. So the restriction of $r$ to $G_L$, which is the twist of $\rho_L$ by $\kappa^{-1}$, is attached through $\iota$ to $\pi\otimes(\chi_\kappa^{-1}\circ\det)$, which is regular algebraic and cuspidal. It is windowed, as in step (F) of the proof of Proposition~\ref{sha-siegel-points}: the twist multiplies the characters of ($\alpha$) by a common character, and as $L$ is a CM field there is $b\in\mathbf R$ with $|\chi_{\kappa,u}(\varpi_u)|=q_u^b$ at every finite place $u$, which shifts the bounds of ($\beta$) by $b$.
\item[(c)] \emph{Descent.} The field $L=FK$ is Galois over $\mathbf Q$ with abelian Galois group, as $F$ and $K$ are; $r$ is irreducible on $G_L$; and $r(g)$ is symplectic with multiplier $\theta(g)$, so with the matrix $Q_0$ of step (6) of the proof of Lemma~\ref{sha-siegel-step-ideal-seeds-point}, ${}^tr(g)Q_0r(g)=\theta(g)Q_0$ for $g\in G$. Let $c$ be the complex conjugation of $D$ and $P=Q_0r(c)^{-1}$; then ${}^tr(s)Pr(csc^{-1})=\theta(s)P$ for $s\in G$. The character $\theta=\chi_{\mathrm{cyc}}^{k+1}\omega^{-1}\xi\psi$ of Definition~\ref{sha-siegel-datum} is $\chi_{\mathrm{cyc}}^{k+1}$ times a character of finite order, so it is attached through $\iota$ to an algebraic Hecke character of $\mathbf Q$ \cite{Ser68}, and it takes the same value at all complex conjugations, which are conjugate in $G$. Let $K'$ be an imaginary quadratic field on which $r$ is irreducible. It remains to show that the restriction of $r$ to $\operatorname{Gal}(\bar{\mathbf Q}/K')$ is attached through $\iota$ to a regular algebraic cuspidal automorphic representation of $GL_4(\mathbf A_{K'})$ and is windowed.
\begin{enumerate}
\item[(i)] \emph{Over $K$.} By Theorem~\ref{sha-siegel-automorphy-base-change}, with $F=K$, $F_0=\mathbf Q$, $M=L$, which is Galois over $K$ with abelian Galois group, the restriction of $r$ to $G_K$, $\mu=\theta$, $c$ and $P$, and as $r$ is irreducible on $G_L$ and its restriction to $G_L$ is automorphic by (b), the restriction of $r$ to $G_K$ is attached through $\iota$ to a regular algebraic cuspidal automorphic representation of $GL_4(\mathbf A_K)$.
\item[(ii)] \emph{Across.} Let $E$ and $E'$ be imaginary quadratic fields such that $r$ is irreducible on $G_E\cap G_{E'}$ and its restriction to $G_E$ is attached through $\iota$ to a regular algebraic cuspidal automorphic representation. Let $M=EE'$, a CM field, Galois over $E$ and over $E'$ with abelian Galois groups, with $G_M=G_E\cap G_{E'}$. Theorem~\ref{sha-siegel-automorphy-base-change}, with $F=E$, $F_0=\mathbf Q$, $\mu=\theta$, $c$ and $P$, shows that the restriction of $r$ to $G_M$ is attached through $\iota$ to a regular algebraic cuspidal automorphic representation of $GL_4(\mathbf A_M)$; the same theorem with $F=E'$ then shows that the restriction of $r$ to $G_{E'}$ is.
\item[(iii)] \emph{Over $K'$.} As $r$ is irreducible on $G_L\subseteq G_K$, it is irreducible on $G_K$ and on $G_{K'}$. Lemma~\ref{sha-siegel-step-ideal-descent-third} gives an imaginary quadratic field $K''$ such that $r$ is irreducible on $G_K\cap G_{K''}$ and on $G_{K''}\cap G_{K'}$. By (i) and by (ii), applied to $(K,K'')$ and then to $(K'',K')$, the restriction of $r$ to $G_{K'}$ is attached through $\iota$ to a regular algebraic cuspidal automorphic representation of $GL_4(\mathbf A_{K'})$.
\end{enumerate}
\item[(d)] \emph{Windows.} By Lemma~\ref{sha-siegel-step-ideal-descent-laws}, applied to $D$, $T$, $\mathcal O'$, $X$, $y$ and $r$, whose values are invertible as they are symplectic, $r$ satisfies conditions (I) and (T) of Definition~\ref{sha-siegel-step-ideal-descent-conditions}, and its restriction to $G_L$ is $\ell$-distinct at every place of $L$ above $\ell$. So the restriction of $r$ to $G_K$ satisfies conditions (I) and (T), as in step (a) of the proof of Lemma~\ref{sha-siegel-step-ideal-descent-across}, and keeps $\ell$-distinctness in $L$.
\begin{enumerate}
\item[(i)] \emph{Over $K$.} The restriction of $r$ to $G_K$ is continuous, irreducible on $G_L$, and attached through $\iota$ to a regular algebraic cuspidal automorphic representation of $GL_4(\mathbf A_K)$ by (c)(i), and its restriction to $G_L$ is windowed by (b). So Lemma~\ref{sha-siegel-step-ideal-descent-updown}(2), applied to $F=K$, the field $L$, which is Galois over $K$ with abelian Galois group, $F_0=\mathbf Q$, $\mu=\theta$, $c$ and $P$ of (c), shows that the restriction of $r$ to $G_K$ is windowed.
\item[(ii)] \emph{Over $K'$.} As $\ell\in S$, by (d) of Definition~\ref{sha-siegel-tower-towers} the decomposition groups of $G$ at the places above $\ell$ lie in $G_K$, so they fix $K$. Lemma~\ref{sha-siegel-step-ideal-descent-across} applies to $r$, $K$, $\theta$ and $J=Q_0$: its hypothesis (i) is (c)(iii), (ii) is shown in (c), (iii) is shown above, and (iv) holds by (c)(iii) and (d)(i). So for every imaginary quadratic field $K'$ such that $r$ is irreducible on $G_{K'}$, the restriction of $r$ to $G_{K'}$ is windowed.
\end{enumerate}
\end{enumerate}
With (c), this is the automorphy of $r$ through $\iota$ (Definition~\ref{sha-siegel-datum}).
% (a) is glueN (IrrArithPointsAutomorphicAway.lean): AutCsd0 from either link, and witA_of with ImaginaryQuadraticEtaCharacter.winOf (Lemma sha-hecke-ideal-plain (1), the WinOf t) and winUnl (WinUnl.lean: a window at v makes the support unlinked, through quot_ne) from the two WitOff, with v0_ne. (b) is pmDown_alt's automorphy_bridge_of_printed_hida, isAlgChar_pushCh_link_kappa with Af.algOfCyc_s, Af.twist_invariance_automorphic and Z.Win.Dn.twist'_alt (with csdN'_alt). (c) is auGal (SupComm.lean), auPair (the polarization from D.Mult r), isAlgChar_pushCh_th_of_good and auSign for theta, auIrr, then isAutomorphicOverImagQuad_of_restriction_aut_irr (AutomorphyDescentSolvableCMToImaginaryQuadraticA.lean): (i) is the SolC of Theorem sha-siegel-automorphy-base-change (here Z.Bc.solC_of_solDescent of Af.solvable_descent_of_printed_hida, through Z.Bc.solC_alt) with the pairing P r(c^-1), (ii) is Z.Bc.across_alt (TAcross.lean, with Z.Bc.BiQuad for M = EE') and (iii) is GalN.V.ofOne0_alt (VWOfOne.lean) with Z.Bc.third (Lemma sha-siegel-step-ideal-descent-third) and automorphy_bridge_of_printed_hida. (d) is Keep.ptLaws (Lemma sha-siegel-step-ideal-descent-laws: IneAb, Flagged, and Unl1 over L from T.Wts (2 * D.wts4)) with flagLaws for the restriction to K; (i) is Z.Win.Dn.down'_alt (Lemma sha-siegel-step-ideal-descent-updown (2)) with hpol (Sol.Pol from auPair with D.cc and auSign) and the Keeps from Unl1 over L; (ii) is Z.Win.Dn.across'_alt (Lemma sha-siegel-step-ideal-descent-across) with hF.1.split l hG.l_mem for (d) of the tower, and the conclusion is isAutQA_of_isAutomorphicOverImagQuad_of_witAQ. pmDown_alt is 88 lines. What remains open is in Lemmas sha-siegel-step-ideal-descent-updown (Z.Win.aUpDownOf_alt, TAUpDownOf.lean, 360 lines, with thruAway, Ord.thruL_alt and ImaginaryQuadraticEtaCharacter.winUpDown) and -descent-laws (ModularityArgument/UpperTriangularPlaces, about 2,300 lines).
\end{proof}

\begin{lemma}\label{sha-siegel-step-ideal}
Assume the Galois representations as in Theorem~\ref{sha-residually-reducible-automorphy}.
Let $D$ be a Siegel datum, let $T$ be a fit deep tower for $D$ (Definition~\ref{sha-siegel-tower-towers}), and let $\iota:\bar{\mathbf Q}_\ell\cong\mathbf C$.
There is a ring $\mathcal O_1\supseteq\mathcal O$ of Definition~\ref{sha-character-parts} such that for every such ring $\mathcal O'\supseteq\mathcal O_1$, every deformation space for $T$ over $\mathcal O'$ (Definition~\ref{sha-siegel-tower-space}) has a modularity ideal (Definition~\ref{sha-siegel-tower-ideal}).
\end{lemma}
\leanref{Development/AutomorphicArithmetic/AutomorphicVarieties/ModularityArgument/ThresholdArguments/DenseThresholdCases/PotentialModularityTower.lean}{UpperTheory.V.Ws.PMExAlt}
\leanref{Development/AutomorphicArithmetic/AutomorphicVarieties/ModularityArgument/ThresholdArguments/DenseThresholdCases/AdjointLiftingOfCohomology.lean}{UpperTheory.V.HV.DescentVanishingReductions.pmEx_of_printsH}
\begin{proof}
Let $K$, $F$ and $L$ be the fields of $T$. As $T$ is fit, $D$ is of case II or of case III and $T$ satisfies (a) to (m), by (4) of Definition~\ref{sha-siegel-tower-towers}; let $\mathcal O_1$ be as in Lemma~\ref{sha-siegel-step-ideal-links}. Let $\mathcal O'\supseteq\mathcal O_1$ be a ring of Definition~\ref{sha-character-parts} and $X=(A,B,P_X)$ a deformation space for $T$ over $\mathcal O'$.
\begin{enumerate}
\item[(1)] \emph{Two links.} By Lemma~\ref{sha-siegel-step-ideal-links}, $T$ has two links over $\mathcal O'$ with the same character $\kappa$, whose places $v_0$ and $v'_0$ lie above distinct primes. By (1) and (3) of Definition~\ref{sha-siegel-automorphy-link}, each gives data of Situation~\ref{sha-hecke-situation} for $L$ over $\mathcal O'$ with $n=4$. Let $B^{\mathrm{ord}}$ and $\Lambda$ be the universal ordinary determinant ring and the Iwasawa algebra of the first, $J_0\subseteq B^{\mathrm{ord}}$ a Hecke ideal (Proposition~\ref{sha-hecke-ideal}), which satisfies the bounds of Proposition~\ref{sha-patching-bound}, and $\varphi:B^{\mathrm{ord}}\to B$ the homomorphism of Lemma~\ref{sha-siegel-automorphy-phi}; let $B'^{\mathrm{ord}}$, $\Lambda'$, $J'_0$ and $\varphi'$ be the same for the second.
\item[(2)] \emph{The ideal.} Let $J$ be the ideal of $B$ generated by $\varphi(J_0)$ and $\varphi'(J'_0)$. A prime of $B$, or the kernel of a homomorphism from $A$, contains $J$ if and only if it contains $\varphi(J_0)$ and $\varphi'(J'_0)$. So a homomorphism from $A$ kills $J$ if and only if it kills $\varphi(J_0)$ and $\varphi'(J'_0)$, and it detects $J$ if it detects the ideals generated by $\varphi(J_0)$ and by $\varphi'(J'_0)$.
\item[(3)] \emph{Finiteness.} By Proposition~\ref{sha-hecke-ideal}(1), $B^{\mathrm{ord}}/J_0$ is finite over $\Lambda$. Lemma~\ref{sha-siegel-step-ideal-finite}, which applies as $T$ satisfies (d), gives a power of $\mathfrak m_B$ in the ideal generated by $\varphi(J_0)$, $\mathfrak m_{\mathcal O'}$ and the $\psi^B_{v,i}(s)-1$ with $s\in I_v(F)$, which lies in the ideal generated by $J$, $\mathfrak m_{\mathcal O'}$ and these elements.
\item[(4)] \emph{Patching and seeds.} Each link, with $X$, $J_0$ and $\varphi$, or with $X$, $J'_0$ and $\varphi'$, is as in Situation~\ref{sha-siegel-step-ideal-setting}. By Lemmas~\ref{sha-siegel-step-ideal-patching} and~\ref{sha-siegel-step-ideal-seeds}, the ideals generated by $\varphi(J_0)$ and by $\varphi'(J'_0)$ have the patching and seed properties, and by (2) so has $J$: a homomorphism that kills $J$ kills both ideals, and one that detects both detects $J$.
\item[(5)] \emph{Purity.} A homomorphism $y$ that kills $J$ kills the ideal generated by $\varphi(J_0)$, so Lemma~\ref{sha-siegel-step-ideal-purity} gives the purity part of the last property of $J$.
\item[(6)] \emph{Automorphy.} Let $y$ and $r$ be as in the automorphy part of that property. As $y$ kills both ideals, Lemma~\ref{sha-siegel-step-ideal-away} applies to both links, which have the same character $\kappa$, and Lemma~\ref{sha-siegel-step-ideal-descent}, as $v_0$ and $v'_0$ lie above distinct primes, shows that $r$ is automorphic through $\iota$.
\end{enumerate}
% Per link (K, B, Lambda, J_0, phi): patching is modular_deformation_quotient_regular_at_point_alt (Lemma sha-siegel-step-ideal-patching), seeds Th.local_zero_of_modular_p (Lemma sha-siegel-step-ideal-seeds), purity pmPur_w_alt (Lemma sha-siegel-step-ideal-purity), automorphy away from v_0 pmAutO_alt (Lemma sha-siegel-step-ideal-away), and glueN with Z.Win.Dn.pmDown_alt (Lemma sha-siegel-step-ideal-descent). In more detail:
% per link (K, B, Lambda, J_0, phi): patching is modular_deformation_quotient_regular_at_point_alt (y composed with phi is a nice arc of type I through the WeakShimuraTheory/ProjectiveLineFlags lemmas P1.psPair, coefRes, rhoEx, formEx, psiEx, redPsi, condA, irrL, condC, condD, mapsPt, n4V0 and Ws.h4, then hM.p1, Proposition sha-patching-bound (1)); seeds are Th.local_zero_of_modular_p (BirthPm.lean: Pt.ptWts, Pt.ptIrr, auRep, Pt.ptUnr, Pt.ptFlag, Pt.ptPair, Z.Win.Birth.ptAutR'_alt, Pt.point_representation_weights, hida_modular_ideal_map, and Th.eigenbasis_rank_four through Th.AdjointLiftingTheorem); purity is pmPur_w_alt (VWinPm.lean: Pu.arith, GalN.V.end1W' from Proposition sha-hecke-ideal (4), Pu.mOne with Pu.noOm, Pu.levFin, Pu.pureOf_w, Pu.frobEigEx, Pu.frobAbs_w); automorphy away from v_0 is pmAutO_alt (auArith, auRep, auIrr and end2oN, Proposition sha-hecke-ideal (5)); the two are combined by glueN (witA_of) with v0_ne, and descended to Q by Z.Win.Dn.pmDown_alt (automorphy_bridge_of_printed_hida, twist_invariance_automorphic, isAutomorphicOverImagQuad_of_restriction_aut_irr, Keep.ptLaws, down'_alt, across'_alt).
% pmEx_of_printsH is pm_exists_of_unramified_adjoint_lifting (PmExTh.lean), that is promodular_existence_of_local_zero (PmTwo.lean: exists_two_links_same_kappa_not_dvd_of_fit_deep, exists_phi_universal_ordinary_det_ring_of_link, U.lam and U.modU for J_0, exists_pow_maximalIdeal_le_of_modular_phi for finiteness, modular_deformation_quotient_regular_at_point_alt for patching, pmPur_w_alt for purity, pmAutO_alt, glueN and Z.Win.Dn.pmDown_alt for automorphy) with Th.local_zero_of_modular_p (BirthPm.lean, through Th.AdjointLiftingTheorem) for seeds. Size: about 270 lines of Lean outside what the book states elsewhere, resting on the Hecke and patching results of this chapter.
\end{proof}

\subsection{Deformation spaces and their completions}\label{subsection-sha-deformation-spaces}

\begin{lemma}\label{sha-siegel-step-coefficients-extension}
Let $V$ be the valuation ring of an algebraically closed field $\Omega$, and let $A\subset V$ be a complete discrete valuation ring with finite residue field such that $\mathfrak m_A\subseteq\mathfrak m_V$ and $\Omega$ is algebraic over the fraction field $E$ of $A$. Let $s\subset\Omega$ be a finite set, and let $B=V\cap E(s)$.
Then $B$ is a complete discrete valuation ring with finite residue field, $\mathfrak m_B\subseteq\mathfrak m_V$, $\Omega$ is algebraic over the fraction field of $B$, and $B$ contains $A$ and every element of $V$ that lies in the subfield of $\Omega$ generated by $A$ and $s$.
\end{lemma}
\leanref{Development/AutomorphicArithmetic/AutomorphicVarieties/WeakShimuraTheory/CompactIntegralFields/SmallExt.lean}{UpperTheory.V.Ws.smallExt}
\begin{proof}
The elements of $s$ are algebraic over $E$, so $E'=E(s)$ is a finite extension of $E$, of degree $n$ say; it contains the subfield generated by $A$ and $s$, and $A\subseteq V\cap E'=B$. For $x,y\in B$ with $x\ne0$, either $y/x$ or $x/y$ lies in $V$, and both lie in $E'$; so $B$ is a valuation ring. A unit of $V$ that lies in $B$ has its inverse in $V\cap E'$, so $\mathfrak m_B\subseteq\mathfrak m_V$, and $A\to B$ is local, since $\mathfrak m_A\subseteq\mathfrak m_V$. Let $\varpi$ be a uniformizer of $A$.
\begin{enumerate}
\item[(a)] $B$ is $\varpi$-adically separated. Let $x\ne0$ lie in $\varpi^kB$ for every $k$. As $x$ is algebraic over $A$, the ideal $xB\cap A$ is nonzero; let $a=u\varpi^e$ be a nonzero element of it, with $u\in A^\times$. Writing $a=xb$ and $x=\varpi^{e+1}y$ with $b,y\in B$, we get $u=\varpi yb$, so $\varpi$ would be a unit of $B$, which it is not.
\item[(b)] $B/\varpi B$ has dimension at most $n$ over $A/\varpi A$. Let $b_1,\dots,b_m\in B$ have linearly independent images in $B/\varpi B$. If $\sum_ia_ib_i=0$ with $a_i\in A$ not all $0$, divide by an $a_j$ of least valuation; this gives a relation in which one coefficient is $1$, and its reduction contradicts the independence. So the $b_i$ are linearly independent over $A$, hence over $E$ in $E'$, and $m\le n$.
\item[(c)] By (a), (b) and the completeness of $A$, $B$ is a finite $A$-module \cite[Theorem~8.4]{Mat86}. So $B$ is noetherian, and as $\varpi$ is not a unit of $B$, it is a discrete valuation ring. As a finite module over the complete ring $A$ it is $\varpi$-adically complete, and $\varpi B$ is a power of $\mathfrak m_B$, so $B$ is complete. Its residue field is finite over that of $A$, and $\Omega$ is algebraic over $E\subseteq\operatorname{Frac}B$.
\end{enumerate}
\end{proof}
% smallExt (SmallExt.lean, 230 lines): B is the subring of V of the elements of IntermediateField.adjoin E s; (a) is the private haus, (b) the private fin_quot, (c) the private key with GlobalTheory.finite_of_isHausdorff_of_finite_quotient and GlobalTheory.FormalIsogeny.isAdicComplete_maximalIdeal_of_finite.

\begin{lemma}\label{sha-siegel-step-coefficients-subring}
Let $\mathcal O$ be a complete local noetherian ring, and let $R$ be a complete local noetherian ring with a local homomorphism $\mathcal O\to R$ such that the residue field of $R$ is finite over $\mathcal O$. Let $S$ be the set of $x\in R$ congruent modulo $\mathfrak m_R$ to the image of an element of $\mathcal O$.
Then $S$ is a subring of $R$ containing $\mathfrak m_R$, it is a complete local noetherian $\mathcal O$-algebra with maximal ideal $\mathfrak m_S=\mathfrak m_R$ and with the residue field of $\mathcal O$, $R$ is finite over $S$, and $\mathfrak m_S^N=S\cap\mathfrak m_R^N$ for every $N\ge1$.
\end{lemma}
\leanref{Development/AutomorphicArithmetic/AutomorphicVarieties/CommutativeAlgebraHeight/ClSub.lean}{UpperTheory.V.CA.cl_sub}
\begin{proof}
Clearly $S$ is a subring containing $\mathfrak m_R$ and the image of $\mathcal O$.
\begin{enumerate}
\item[(a)] $S$ is local with maximal ideal $\mathfrak m_R$. An element $x\in S\setminus\mathfrak m_R$ is congruent to the image of some $a\in\mathcal O$, which is then a unit of $\mathcal O$ as $\mathcal O\to R$ is local; $x$ is a unit of $R$, and $x^{-1}$ is congruent to the image of $a^{-1}$, so $x^{-1}\in S$. Elements of $\mathfrak m_R$ are not units. So the residue field of $S$ is that of $\mathcal O$.
\item[(b)] $\mathfrak m_S^N=S\cap\mathfrak m_R^N$. The inclusion $\subseteq$ is clear. Conversely, by induction on $N$, every element of $\mathfrak m_R^{N}=\mathfrak m_R\cdot\mathfrak m_R^{N-1}$ is a sum of products $ij$ with $i\in\mathfrak m_R=\mathfrak m_S$ and $j\in\mathfrak m_R^{N-1}\cap S=\mathfrak m_S^{N-1}$.
\item[(c)] $R$ is finite over $S$: it is generated by $1$ and lifts $y_1,\dots,y_t$ of generators of the residue field of $R$ over $\mathcal O$. Indeed, every $x\in R$ is congruent modulo $\mathfrak m_R$ to an $\mathcal O$-combination $y$ of the $y_i$, and $x=(x-y)\cdot1+y$ with $x-y\in\mathfrak m_R\subseteq S$.
\item[(d)] $R$ is a noetherian $S$-module. Let $N\subseteq R$ be an $S$-submodule, and $J=RN$. As $R$ is noetherian and finite over $S$, $J$ and $\mathfrak m_RJ$ are finitely generated $S$-modules, and $\mathfrak m_RJ\subseteq N\subseteq J$, since $\mathfrak m_R\subseteq S$. The quotient $J/\mathfrak m_RJ$ is a finite module over the field $S/\mathfrak m_S$, hence noetherian, so $N/\mathfrak m_RJ$ is finitely generated, and so is $N$.
\item[(e)] By (d), $S$, an $S$-submodule of $R$, is noetherian. It is complete: by (b), a sequence in $S$ that is Cauchy for the $\mathfrak m_S$-adic topology is Cauchy in $R$, and converges there to some $L$; $L$ is congruent modulo $\mathfrak m_R$ to an element of the sequence, so $L\in S$, and by (b) the sequence converges to $L$ in $S$; and $S$ is separated, since $R$ is.
\end{enumerate}
\end{proof}
% cl_sub (ClSub.lean, about 340 lines): (a) is the private residue_constant_subalgebra_isLocalRing, (b) residue_constant_subalgebra_mem_maximalIdeal_pow_iff, (c) residue_constant_subalgebra_module_finite, (d) residue_constant_subalgebra_isNoetherian with isNoetherian_of_finite_torsion_by_maximalIdeal, (e) isAdicComplete_residue_constant_subalgebra.

\begin{lemma}\label{sha-siegel-step-coefficients-standard}
Let $D$ be a Siegel datum, $T$ a fit tower for $D$ with field $F$ (Definition~\ref{sha-siegel-tower-towers}), and $\mathcal O'\supseteq\mathcal O$ a ring of Definition~\ref{sha-character-parts}, with residue field $k'$. Let $\Phi$ be the problem of $T$, with characters $\chi_i$ on $G_F$, places $v\in\Sigma_T$, groups $D_v=D_v(F)$, sets $\mu_v$ and set $N_\Phi$ (Definition~\ref{sha-siegel-tower-pairs}); bars denote images in $k'$.
\begin{enumerate}
\item $\chi_i\chi_{5-i}=\theta$ on $G_F$ for every $i$; $\bar\chi_i(g)=1$ for $g\in N_\Phi$; the $\bar\chi_i(g_0^T)$ are pairwise distinct; $2$ is a unit of $k'$; and the $\bar\chi_i$ are locally constant on $G_F$.
\item The field $k'$ has characteristic $\ell$, and the $\chi_i:G_F\to\mathcal O'$ are continuous for the $\mathfrak m_{\mathcal O'}$-adic topology.
\item For every open subgroup $H$ of $G_F$ that contains the normal subgroup generated by $N_\Phi$, there are only finitely many continuous homomorphisms $H\to\mathbf Z/\ell$ that vanish on the closure of that normal subgroup.
\item Each $D_v$ is closed in $G_F$, and for every open subgroup $H$ of $D_v$ there are only finitely many continuous homomorphisms $H\to\mathbf Z/\ell$.
\item Every $s\in\mu_v$ lies in $D_v$ and has $\bar\chi_i(s)=1$ for every $i$.
\end{enumerate}
\end{lemma}
\leanref{Development/AutomorphicArithmetic/AutomorphicVarieties/WeakShimuraTheory/CompactIntegralFields/ContResQ.lean}{UpperTheory.V.Ws.std}
\leanref{Development/AutomorphicArithmetic/AutomorphicVarieties/WeakShimuraTheory/RedIdeal.lean}{UpperTheory.V.Ws.Fr.Std}
\begin{proof}
Each $\delta_i$ is continuous with values in $\mathcal O^\times$, so its reduction modulo $\mathfrak m_{\mathcal O}^n$ is locally constant for every $n$; and the set of $g\in G$ with $\bar r_0(g)=1$ is an open subgroup.
\begin{enumerate}
\item[(1)] The relabelling permutes the $\delta_i$ compatibly with $i\mapsto5-i$, and $\delta_1\delta_4=\theta=\delta_2\delta_3$, as $\Psi\alpha=\theta$; so $\chi_i\chi_{5-i}=\theta$. An element $g$ of $N_\Phi$ has $\bar r_0(g)=1$, so $\bar\delta_i(g)=1$ for all $i$. The residues $\bar\chi_i(g_0^T)$ are pairwise distinct by (l). As $\ell$ is odd, $2$ is a unit of $k'$. The $\bar\chi_i$ are locally constant on $G_F$ by the first paragraph.
\item[(2)] As $\ell\in\mathfrak m_{\mathcal O}$, which maps into $\mathfrak m_{\mathcal O'}$, $k'$ has characteristic $\ell$. The reduction of $\chi_i$ modulo $\mathfrak m_{\mathcal O'}^n$ factors through its reduction modulo $\mathfrak m_{\mathcal O}^n$, which is locally constant.
\item[(3)] Let $H$ be as there, $M$ the closure of the normal subgroup generated by $N_\Phi$, and $U$ the open subgroup of the $g\in G_F$ with $\bar r_0(g)=1$. For a place $w$ not above $\ell$, $I_w(F)\cap U\subseteq N_\Phi\subseteq M$. Let $\zeta$ be a primitive $\ell$-th root of unity, and let $E$ be the fixed field of the open subgroup $H\cap U\cap G_{F(\zeta)}$, a number field containing $\zeta$ with $G_E\subseteq H$ of finite index. The restriction to $G_E$ of a continuous homomorphism $\phi:H\to\mathbf Z/\ell$ vanishing on $M$ vanishes on the inertia groups of $G_E$ at the places not above $\ell$; by Lemma~\ref{sha-descent-link-hermite}, with $n=\ell$ and $N=\ell$, there are only finitely many such restrictions. And $\phi$ is determined by its restriction to $G_E$ and its values at representatives of the finitely many cosets of $G_E$ in $H$.
\item[(4)] Decomposition groups are closed. The finiteness is the finiteness of the continuous characters of order $\ell$ of the absolute Galois group of a finite extension of $\mathbf Q_\ell$ \cite[Chapter~VII, \S2]{NSW08}, applied to the fixed field of $H$ in the completion.
\item[(5)] Let $s\in\mu_v$, so that $s\in D_v$ and $s^{\ell^m}$ lies in the closure of the commutator subgroup of $D_v$ for some $m$. The character $\bar\chi_i$ of $D_v$ is locally constant with values in the commutative group $k'^\times$, so its kernel is closed and contains that closure. Hence $\bar\chi_i(s)^{\ell^m}=1$, so $(\bar\chi_i(s)-1)^{\ell^m}=0$ in the field $k'$ of characteristic $\ell$, and $\bar\chi_i(s)=1$.
\end{enumerate}
\end{proof}
% std (ContResQ.lean, 143 lines): (1) is the private ok_of with facts (f6 for delta_i delta_{5-i} = theta, f11 for the local constancy), (2) is char_res and f11 through inc_max, (3) is the private fin_chars with Kum.finChar (Lemma sha-descent-link-hermite) and f12 for the openness of ker rbar, (4) is isClosed_Dec and Pa.phiDec from the published LocalFiniteness of CohomologyHyp, (5) is the last bullet of std. The Lean also proves (1) along every homomorphism from O to a field that kills the maximal ideal; Fr.Std asks for (1) along O -> O' -> k'.

\begin{lemma}\label{sha-siegel-step-coefficients-universal}
Let $D$, $T$, $F$, $\mathcal O'$, $k'$ and $\Phi$ be as in Lemma~\ref{sha-siegel-step-coefficients-standard}. For a complete local noetherian $\mathcal O'$-algebra $R$ with residue field $k'$, a \emph{near representation} over $R$ is a pair $(r,(\psi_{v,i}))$ of a homomorphism $r:G_F\to M_4(R)$ and multiplicative maps $\psi_{v,i}:D_v\to R$, for $v\in\Sigma_T$ and $1\le i\le4$, such that: $r(g)$ is symplectic with multiplier $\theta(g)$; $r(g)=1$ for $g\in N_\Phi$; $r(g_0^T)$ is diagonal; the entries of $r$ and the $\psi_{v,i}$ are continuous for the $\mathfrak m_R$-adic topology; $r(g)_{ii}$ has residue $\bar\chi_i(g)$ and $r(g)_{ij}\in\mathfrak m_R$ for $i\ne j$; and $\psi_{v,i}(s)$ has residue $\bar\chi_i(s)$.
Then there are a complete local noetherian $\mathcal O'$-algebra $R_\Phi$ with residue field $k'$ and a near representation $\mathcal D$ over $R_\Phi$ such that for every near representation $\mathcal D'$ over a complete local noetherian $\mathcal O'$-algebra $R'$ with residue field $k'$ there is a unique homomorphism of $\mathcal O'$-algebras $R_\Phi\to R'$ that carries $\mathcal D$ to $\mathcal D'$.
\end{lemma}
\leanref{Development/AutomorphicArithmetic/AutomorphicVarieties/WeakShimuraTheory/AdicExtensionMatrices/SxUnivD.lean}{UpperTheory.V.Ws.sxUnivD}
\leanref{Development/AutomorphicArithmetic/AutomorphicVarieties/WeakShimuraTheory/AdicExtensionMatrices/AdicContMap.lean}{UpperTheory.V.Ws.Fr.Near}
\leanref{Development/AutomorphicArithmetic/AutomorphicVarieties/WeakShimuraTheory/AdicExtensionMatrices/AdicContMap.lean}{UpperTheory.V.Ws.Fr.RUniv}
\begin{proof}
All rings below are complete local noetherian $\mathcal O'$-algebras with residue field $k'$, and homomorphisms between them are homomorphisms of $\mathcal O'$-algebras; such a homomorphism maps the maximal ideal into the maximal ideal, so it preserves continuity and residues. Let $N'$ be the closed normal subgroup of $G_F$ generated by $N_\Phi$, and let $\bar r_\Phi(g)=\operatorname{diag}(\bar\chi_1(g),\dots,\bar\chi_4(g))$ for $g\in G_F$.
\begin{enumerate}
\item[(a)] \emph{Lifts of $\bar r_\Phi$.} By Lemma~\ref{sha-siegel-step-coefficients-standard}(1), $\bar r_\Phi$ is a locally constant homomorphism, trivial on $N_\Phi$; its kernel is closed and normal, so it contains $N'$, and $\bar r_\Phi$ factors through the profinite group $G_F/N'$. Open subgroups of $G_F/N'$ are the images of the open subgroups $H$ of $G_F$ that contain $N'$, so by Lemma~\ref{sha-siegel-step-coefficients-standard}(3) each of them has only finitely many continuous homomorphisms to $\mathbf Z/\ell$. By Mazur's theorem \cite{Maz89}, the homomorphisms $G_F/N'\to M_4(R)$ with $\mathfrak m_R$-adically continuous entries that reduce to $\bar r_\Phi$ modulo $\mathfrak m_R$ are the images of a universal one, with values in a ring $R^\square$. A homomorphism $r:G_F\to M_4(R)$ with continuous entries that is trivial on $N_\Phi$ is trivial on $N'$, as the set of $g$ with $r(g)=1$ is a closed normal subgroup. So composing with $G_F\to G_F/N'$ gives a homomorphism $r^\square:G_F\to M_4(R^\square)$ of which every such $r$ reducing to $\bar r_\Phi$ is the image under a unique homomorphism $R^\square\to R$.
\item[(b)] \emph{The conditions on $r$.} Let $I\subseteq R^\square$ be the ideal generated by the entries of $r^\square(g)^*r^\square(g)-\theta(g)$ for $g\in G_F$, of $r^\square(g)-1$ for $g\in N_\Phi$, and the entries of $r^\square(g_0^T)$ off the diagonal. The image of $r^\square$ under a homomorphism $f:R^\square\to R$ is symplectic with multiplier $\theta$, trivial on $N_\Phi$ and diagonal at $g_0^T$ if and only if $f$ kills $I$. The homomorphism $\bar r_\Phi$ satisfies these conditions, as $\chi_i\chi_{5-i}=\theta$ by Lemma~\ref{sha-siegel-step-coefficients-standard}(1). So the map $R^\square\to k'$ that carries $r^\square$ to $\bar r_\Phi$ kills $I$, the quotient $A_r=R^\square/I$ is a ring of the kind above, and the image $r_A$ of $r^\square$ in $A_r$ has the following property: every homomorphism $r$ as in the definition of a near representation over $R$ is the image of $r_A$ under a unique homomorphism $A_r\to R$.
\item[(c)] \emph{The characters.} Let $v\in\Sigma_T$ and $1\le i\le4$. By Lemma~\ref{sha-siegel-step-coefficients-standard}(4), $D_v$ is closed, hence compact, and each open subgroup of $D_v$ has only finitely many continuous homomorphisms to $\mathbf Z/\ell$; and $\bar\chi_i$ is locally constant on $D_v$. By Mazur's theorem for lifts of dimension $1$, there are a ring $B_{v,i}$ and an $\mathfrak m_{B_{v,i}}$-adically continuous multiplicative map $\psi^{\mathrm u}_{v,i}:D_v\to B_{v,i}$ with residue $\bar\chi_i$ such that every continuous multiplicative map $D_v\to R$ with residue $\bar\chi_i$ is the image of $\psi^{\mathrm u}_{v,i}$ under a unique homomorphism $B_{v,i}\to R$.
\item[(d)] \emph{The product.} Each ring $C$ of the kind above is the quotient of a power series ring over $\mathcal O'$ in finitely many variables mapping to generators of $\mathfrak m_C$; with all variables of weight $0$ in the trivial group, this is a weighted presentation (Definition~\ref{sha-counts-vd-model}(1)). Lemma~\ref{sha-counts-vd-model-coproduct}(1), for $A_r$ and the $B_{v,i}$, gives a ring $R_\Phi$ with homomorphisms $j_r:A_r\to R_\Phi$ and $j_{v,i}:B_{v,i}\to R_\Phi$ such that, for every ring $R'$, homomorphisms $A_r\to R'$ and $B_{v,i}\to R'$ are induced by a unique homomorphism $R_\Phi\to R'$. Let $\mathcal D$ be the pair of $j_r\circ r_A$ and the $j_{v,i}\circ\psi^{\mathrm u}_{v,i}$, which is a near representation over $R_\Phi$.
\item[(e)] \emph{Universality.} Let $\mathcal D'=(r',(\psi'_{v,i}))$ be a near representation over $R'$. By (b) there is a unique homomorphism $A_r\to R'$ carrying $r_A$ to $r'$, and by (c) a unique homomorphism $B_{v,i}\to R'$ carrying $\psi^{\mathrm u}_{v,i}$ to $\psi'_{v,i}$ for each $(v,i)$. By (d) they are induced by a unique homomorphism $R_\Phi\to R'$, which carries $\mathcal D$ to $\mathcal D'$. A homomorphism $R_\Phi\to R'$ carrying $\mathcal D$ to $\mathcal D'$ induces these homomorphisms on $A_r$ and the $B_{v,i}$, by their uniqueness, so it is this one.
\end{enumerate}
% sxUnivD (SxUnivD.lean, 90 lines): (a) and (b) are sxUnivRep (UnivRep.lean, 113 lines): sxUnivKill (UnivKill.lean, 130 lines) through the quotient by the closure of the normal closure of N, with univ_lift (CommutativeAlgebra/UnivLift.lean, the Lean's proof of Mazur's theorem from the finiteness of lifts to the dual numbers, about 2,000 lines with its imports), then univ_cond (UnivCond.lean) with the ideal of rep_closed and the witness r_1 = splitRD. (c) is sxUnivChar (UnivChar.lean), univ_lift for 1 x 1 matrices with Fr.Std.phiv. (d) is sxCoprodOpt (CoprodOpt.lean) through CA.co_coprod (CoCoprod.lean), which applies Alg.coprod to the presentations co_coprod_pres by generators of the maximal ideals, with weights in Unit. (e) is the second half of sxUnivD.
\end{proof}

\begin{lemma}\label{sha-siegel-step-coefficients-space-core-weights}
Let $A$ be a commutative ring, $E\subseteq A$ a subring, and $\mathcal R\subseteq M_4(A)$ a subring that contains the matrix units $E_{11},\dots,E_{44}$, is stable under $X\mapsto X^*$ (Definition~\ref{sha-siegel-datum}), and satisfies $\operatorname{tr}x\in E$ for $x\in\mathcal R$. Let $x_1,\dots,x_n\in\mathcal R$ and let $(i_1,j_1),\dots,(i_n,j_n)$ be entries whose weights add up to $0$. Then $(x_1)_{i_1j_1}\cdots(x_n)_{i_nj_n}\in E$.
\end{lemma}
\leanref{Development/AutomorphicArithmetic/AutomorphicVarieties/WeakShimuraTheory/AdicExtensionMatrices/GdotSelf.lean}{UpperTheory.V.Ws.sxGma}
\begin{proof}
For $u,u'\in\mathbf Z^2$ let $\langle u,u'\rangle=u_1u'_1+u_2u'_2$, and write $w(p,q)$ for the weight of the entry $(p,q)$; then $w(p,q)=w(p,r)+w(r,q)$ and $w(5-q,5-p)=w(p,q)$.
\begin{enumerate}
\item[(a)] \emph{Diagonal entries.} For $x\in\mathcal R$, $x_{pp}=\operatorname{tr}(E_{pp}x)\in E$.
\item[(b)] \emph{Merging.} Let $x,y\in\mathcal R$ and let $(i,j)$, $(k,l)$ be entries with $\langle w(i,j),w(k,l)\rangle<0$. A check of the pairs of entries shows that $j=k$, $l=i$, $j=5-l$ or $i=5-k$. In each case some $z\in\mathcal R$ has an entry $(p,q)$ of weight $w(i,j)+w(k,l)$ with $z_{pq}=x_{ij}y_{kl}$: if $j=k$, take $z=xE_{jj}y$ and $(p,q)=(i,l)$; if $l=i$, take $z=yE_{ll}x$ and $(p,q)=(k,j)$; if $j=5-l$, the matrix $y'=\pm y^*\in\mathcal R$ with $y'_{j,5-k}=y_{kl}$ gives $z=xE_{jj}y'$ and $(p,q)=(i,5-k)$; if $i=5-k$, the matrix $x'=\pm x^*$ with $x'_{5-j,k}=x_{ij}$ gives $z=x'E_{kk}y$ and $(p,q)=(5-j,l)$.
\item[(c)] \emph{Induction on $n$.} For $n=0$ the product is $1$. Let $n\ge1$ and $w_1=w(i_1,j_1)$. If $i_1=j_1$, then $(x_1)_{i_1i_1}\in E$ by (a), $w_1=0$, and the product of the other factors lies in $E$ by induction. Otherwise $\langle w_1,w_1\rangle>0$, and as $\sum_kw(i_k,j_k)=0$, some $k\ge2$ has $\langle w_1,w(i_k,j_k)\rangle<0$. By (b), the factors $(x_1)_{i_1j_1}$ and $(x_k)_{i_kj_k}$ can be replaced by one entry of an element of $\mathcal R$ whose weight is the sum of theirs; the product is unchanged, the weights still add up to $0$, and there are $n-1$ factors, so the induction applies.
\end{enumerate}
\end{proof}
% sxGma (GdotSelf.lean, 138 lines): gdot is the pairing; (a) is hdiag; (b) is gmerge with gdot_neg (a decide over the entries), gcomp and gflip; (c) is the strong induction, with gdot_self and gdot_sum.

\begin{lemma}\label{sha-siegel-step-coefficients-space-core-traces}
Let $A$ be a commutative ring, $E\subseteq A$ a subring, $H$ a group, $h_0\in H$ and $r:H\to M_4(A)$ a homomorphism such that each $r(h)$ is symplectic with a multiplier in $E$, $\operatorname{tr}r(h)\in E$ for every $h$, and $r(h_0)$ is diagonal with diagonal entries $d_1,\dots,d_4\in E$ such that each $d_i-d_j$ with $i\ne j$ has an inverse in $E$. Then every product $r(h_1)_{i_1j_1}\cdots r(h_n)_{i_nj_n}$ of entries whose weights add up to $0$ lies in $E$.
\end{lemma}
\leanref{Development/AutomorphicArithmetic/AutomorphicVarieties/WeakShimuraTheory/AdicExtensionMatrices/MonoZero.lean}{UpperTheory.V.Ws.sxMonoZero}
\begin{proof}
Let $\mathcal R$ be the $E$-subalgebra of $M_4(A)$ generated by the $r(h)$; as $r$ is a homomorphism, its elements are the $E$-linear combinations of the $r(h)$. So the trace of every element of $\mathcal R$ lies in $E$. If $r(h)$ has multiplier $\nu\in E$, then $r(h)^*=\nu r(h^{-1})$, as $r(h)r(h^{-1})=1$; as $X\mapsto X^*$ is additive and commutes with scalars, $\mathcal R$ is stable under it. Let $u_{ij}\in E$ be the inverse of $d_i-d_j$; then $E_{ii}=\prod_{j\ne i}u_{ij}(r(h_0)-d_j)\in\mathcal R$. So Lemma~\ref{sha-siegel-step-coefficients-space-core-weights} applies.
\end{proof}
% sxMonoZero (MonoZero.lean, 142 lines): S is Algebra.adjoin E (range r), with hsp, htrSp and hadjSp; single_mem_of_diag and single_diag_mem_of_inverse_differences_mem for the matrix units; then sxGma.

\begin{lemma}\label{sha-siegel-step-coefficients-space-core-laurent}
Let $\mathcal O'$ be a complete local noetherian ring with residue field $k'$, $R'$ a complete local noetherian $\mathcal O'$-algebra with residue field $k'$ whose maximal ideal is nilpotent, and $c\ge0$ an integer. Let $R'((t))$ be the ring of Laurent series over $R'$. Then there are a complete local noetherian $\mathcal O'$-algebra $S$ with residue field $k'$ and an injective ring homomorphism $e:S\to R'((t))$ such that: $e(a)$ is the constant series $a$ for $a\in\mathcal O'$; for $n\in\mathbf Z$ and $x\in R'$ with $n>0$, or with $n\ge-c$ and $x\in\mathfrak m_{R'}$, there is $y\in\mathfrak m_S$ with $e(y)=xt^n$; and for every $n$ there is $M$ such that the coefficient of $t^n$ in $e(y)$ vanishes for $y\in\mathfrak m_S^M$.
\end{lemma}
\leanref{Development/AutomorphicArithmetic/AutomorphicVarieties/WeakShimuraTheory/AdicExtensionMatrices/NearOfAdjoin.lean}{UpperTheory.V.Ws.sxLaur}
\begin{proof}
Let $x_1,\dots,x_k$ generate $\mathfrak m_{R'}$, let $S\subseteq R'((t))$ be the $R'[[t]]$-subalgebra generated by the $t^{-c}x_i$, and let $e$ be the inclusion. Each $t^{-c}x_i$ is nilpotent, as $\mathfrak m_{R'}$ is; so $S$ is a finite $R'[[t]]$-module, and every $y\in S$ is the sum of a power series $p$ and a nilpotent element. Then $y$ is a unit if and only if the constant term of $p$ is a unit of $R'$; so $S$ is local, every element of $S$ is congruent modulo $\mathfrak m_S$ to an element of $\mathcal O'$, and, as a finite module over the complete local noetherian ring $R'[[t]]$ whose maximal ideal maps into $\mathfrak m_S$, $S$ is noetherian and complete. If $x=\sum_ia_ix_i\in\mathfrak m_{R'}$ and $n\ge-c$, then $xt^n=t^{n+c}\sum_ia_it^{-c}x_i\in S$ is nilpotent, so it lies in $\mathfrak m_S$; and $xt^n\in tR'[[t]]\subseteq\mathfrak m_S$ for $n>0$. Finally, $\mathfrak m_S$ lies in the radical of $tS$, as $y=p+(\text{nilpotent})\in\mathfrak m_S$ has $p$ with nilpotent constant term; as $S$ is noetherian, $\mathfrak m_S^K\subseteq tS$ for some $K$; and $S\subseteq t^{-N}R'[[t]]$ for some $N$, being generated by finitely many Laurent series. So for $y\in\mathfrak m_S^{K(n+N+1)}$, $e(y)\in t^{n+N+1}S\subseteq t^{n+1}R'[[t]]$, and its coefficient of $t^n$ vanishes.
\end{proof}
% sxLaur (NearOfAdjoin.lean): S = Algebra.adjoin (PowerSeries R) G with G the t^(-c) x_i, hfin (finite over PowerSeries R), near_of_adjoin and local_of_near, co_finite for IsCO, alg_monomial for the elements x t^n, bound_fin and hrad for the last property; Laur is the structure of the four properties.

\begin{lemma}\label{sha-siegel-step-coefficients-space-core}
Let $D$, $T$, $F$, $\mathcal O'$, $k'$ and $\Phi$ be as in Lemma~\ref{sha-siegel-step-coefficients-standard}, and let $R_\Phi$ and $\mathcal D=(r_{\mathcal D},(\psi_{\mathcal D,v,i}))$ be as in Lemma~\ref{sha-siegel-step-coefficients-universal}. Let $\pi$ be a weighted presentation of $R_\Phi$ by variables with weights in $\mathbf Z^2$ (Definition~\ref{sha-counts-vd-model}(1), with $\mathcal O'$ and $k'$ in place of $\mathcal O$ and $\mathbf F$) such that: each variable maps into $\mathfrak m_{R_\Phi}$, either to some $r_{\mathcal D}(g)_{ij}-c$ with weight the weight of $(i,j)$, where $c\in\mathcal O'$ and $c=0$ if $i\ne j$, or to some $\psi_{\mathcal D,v,i}(s)-c$ with weight $0$, where $c\in\mathcal O'$; each $r_{\mathcal D}(g)_{ij}$ lies in the part of the weight of $(i,j)$; and each $\psi_{\mathcal D,v,i}(s)$ lies in the part of weight $0$. Let $B$ be the part of weight $0$ and $A_0=R_\Phi^{\mathrm{gr}}$, so that $r_{\mathcal D}$ takes values in $M_4(A_0)$ and the $\psi_{\mathcal D,v,i}$ in $B$, and suppose that $B$ is a complete local noetherian $\mathcal O'$-algebra with residue field $k'$. For $g\in G_F$ let $c_g\in B[X]$ be the characteristic polynomial of $r_{\mathcal D}(g)$.
\begin{enumerate}
\item Let $Y$ be a topological space. A map $\phi:Y\to A_0$ with values in one part that is continuous for the $\mathfrak m_{R_\Phi}$-adic topology is continuous for the $\mathfrak m_BA_0$-adic topology, and a map $Y\to B$ that is continuous for the $\mathfrak m_{R_\Phi}$-adic topology is continuous for the $\mathfrak m_B$-adic topology.
\item Two homomorphisms of $\mathcal O'$-algebras from $B$ to a complete local noetherian $\mathcal O'$-algebra with residue field $k'$ that agree at the coefficients of the $c_g$ and at the $\psi_{\mathcal D,v,i}(s)$ are equal.
\item Let $\tau:\mathbf Z^2\to\mathbf Z$ be additive, $R'$ a complete local noetherian $\mathcal O'$-algebra with residue field $k'$, and $(r',(\psi'_{v,i}))$ a pair with all the properties of a near representation over $R'$, except that $r'(g)_{ij}\in\mathfrak m_{R'}$ is only required for $i\ne j$ with $\tau(w)\le0$, where $w$ is the weight of $(i,j)$. Then some homomorphism of $\mathcal O'$-algebras $A_0\to R'$ carries $(r_{\mathcal D},(\psi_{\mathcal D,v,i}))$ to $(r',(\psi'_{v,i}))$.
\end{enumerate}
\end{lemma}
\leanref{Development/AutomorphicArithmetic/AutomorphicVarieties/WeakShimuraTheory/AdicExtensionMatrices/ContGr.lean}{UpperTheory.V.Ws.sxContGr}
\leanref{Development/AutomorphicArithmetic/AutomorphicVarieties/WeakShimuraTheory/AdicExtensionMatrices/GenZero.lean}{UpperTheory.V.Ws.sxGenZero}
\leanref{Development/AutomorphicArithmetic/AutomorphicVarieties/WeakShimuraTheory/AdicExtensionMatrices/QuotPowNil.lean}{UpperTheory.V.Ws.sxNearA}
\begin{proof}
Write $R=R_\Phi$, $R_w$ for the parts and $z_s$ for the images of the variables of $\pi$, and $\mathfrak m_B=B\cap\mathfrak m_R$, the maximal ideal of $B$ (Lemma~\ref{sha-counts-vd-presentations}(3)). By Lemma~\ref{sha-counts-vd-presentations}(2), $A_0=\bigoplus_wR_w$, and $A_0$ is generated over $B$ by the $z_s$.

(1) Let $\phi:Y\to R_w$ be continuous for the $\mathfrak m_R$-adic topology, and $N\ge0$. Lemma~\ref{sha-automorphy-unramified-near-tangent-filtration} gives $M$ with $R_w\cap\mathfrak m_R^M\subseteq\mathfrak m_B^NR_w$. Let $y_0\in Y$; the $y$ with $\phi(y)\equiv\phi(y_0)$ modulo $\mathfrak m_R^M$ form a neighbourhood of $y_0$, and for them $\phi(y)-\phi(y_0)\in R_w\cap\mathfrak m_R^M\subseteq\mathfrak m_B^NR_w\subseteq(\mathfrak m_BA_0)^N$. So $\phi$ is continuous for the $\mathfrak m_BA_0$-adic topology. For $w=0$, $\mathfrak m_B^NR_0$ is the ideal $\mathfrak m_B^N$ of $B$, which gives the second claim.

(2) Let $f,f':B\to R'$ be as there, and let $E\subseteq R$ be the subring of the elements of $B$ at which $f$ and $f'$ agree.
\begin{enumerate}
\item[(a)] \emph{Monomials in the entries.} Every product of entries of the $r_{\mathcal D}(g)$ whose weights add up to $0$ lies in $E$. We apply Lemma~\ref{sha-siegel-step-coefficients-space-core-traces} to $r_{\mathcal D}$, $E$ and $h_0=g_0^T$. The multipliers $\theta(g)$ are constants, at which $f$ and $f'$ agree. The trace of $r_{\mathcal D}(g)$ is minus a coefficient of $c_g$, so it lies in $E$. The matrix $r_{\mathcal D}(g_0^T)$ is diagonal, and its diagonal entries $d_i$ lie in $B$, the part of weight $0$, with $c_{g_0^T}=\prod_i(X-d_i)$. As $f$ and $f'$ agree at the coefficients of $c_{g_0^T}$, $f'(d_i)$ is a root of $\prod_j(X-f(d_j))$; for $j\ne i$, $f'(d_i)-f(d_j)$ is a unit, as its residue is $\bar\chi_i(g_0^T)-\bar\chi_j(g_0^T)\ne0$ by Lemma~\ref{sha-siegel-step-coefficients-standard}(1); so $f'(d_i)=f(d_i)$, and $d_i\in E$. For the same reason $d_i-d_j$ is a unit of $B$ for $i\ne j$, and $f$ and $f'$ agree at its inverse, the inverse of their common value at $d_i-d_j$.
\item[(b)] \emph{Monomials in the variables.} Every $z^n=\prod_sz_s^{n_s}$ of weight $0$ lies in $E$. Each $z_s$ is $r_{\mathcal D}(g)_{ij}-a$, with the weight of $(i,j)$ and $a=0$ if $i\ne j$, or $\psi_{\mathcal D,v,i}(s')-a$, of weight $0$, with $a\in\mathcal O'$. In the second case, and in the first with $i=j$ by (a), $z_s\in E$. So, by induction on $n$, $z^n=e\,m$ with $e\in E$ and $m$ a product of entries of the $r_{\mathcal D}(g)$ whose weights add up to the weight of $n$. If that weight is $0$, then $m\in E$ by (a), and $z^n\in E$.
\item[(c)] \emph{Conclusion.} By Lemma~\ref{sha-counts-vd-presentations}(3), $\mathfrak m_B$ is generated by $\mathfrak m_{\mathcal O'}$ and finitely many $z^n\in B$ of weight $0$, at which $f$ and $f'$ agree by (b). As in step (ii) of the proof of Lemma~\ref{sha-counts-vd-model-coproduct}, two homomorphisms of $\mathcal O'$-algebras from $B$ to $R'$ that agree at generators of $\mathfrak m_B$ over $\mathfrak m_{\mathcal O'}$ are equal, as they map $\mathfrak m_B$ into $\mathfrak m_{R'}$ and $R'$ is separated. So $f=f'$.
\end{enumerate}

(3) Let $R'$, $\tau$ and $(r',(\psi'_{v,i}))$ be as there. For a commutative ring $C$ with such a pair, call it \emph{$\tau$-near}. By the hypothesis on $\pi$, each $z_s$ is $r_{\mathcal D}(g)_{ij}-a$ or $\psi_{\mathcal D,v,i}(s')-a$; let $y_s\in R'$ be $r'(g)_{ij}-a$, respectively $\psi'_{v,i}(s')-a$.
\begin{enumerate}
\item[(a)] \emph{Uniqueness.} Two homomorphisms of $\mathcal O'$-algebras from $A_0$ to a complete local noetherian $\mathcal O'$-algebra with residue field $k'$ that agree at the $z_s$ are equal: they agree at the $z^n$ of weight $0$, hence on $B$, as in step (c) of (2), and $A_0$ is generated over $B$ by the $z_s$.
\item[(b)] \emph{A model over Laurent series.} Suppose that $\mathfrak m_{R'}$ is nilpotent. Let $c\ge|\tau(w)|$ for the weights $w$ of the entries, and let $S$ and $e:S\to R'((t))$ be as in Lemma~\ref{sha-siegel-step-coefficients-space-core-laurent}. Let $h_i$ be the value of $\tau$ at the weight of the entry $(i,4)$; as weights of entries add, the entry $(i,j)$ has $\tau$ of its weight equal to $h_i-h_j$, and $h_i+h_{5-i}=\tau(1,1)$ does not depend on $i$. Every entry $r'(g)_{ij}t^{h_i-h_j}$ and every $\psi'_{v,i}(s')$ is $e$ of a unique element of $S$: for $i=j$, write the entry as $a+m$ with $a\in\mathcal O'$ and $m\in\mathfrak m_{R'}$; for $i\ne j$, either $h_i-h_j>0$, or $h_i-h_j\le0$ and $r'(g)_{ij}\in\mathfrak m_{R'}$. Let $(r'',(\psi''_{v,i}))$ be the pair over $S$ so obtained. As $e$ is injective and $e(r''(g))=\Delta r'(g)\Delta^{-1}$ with $\Delta=\operatorname{diag}(t^{h_1},\dots,t^{h_4})$, which is symplectic with multiplier $t^{h_1+h_4}$, the pair $(r'',(\psi''_{v,i}))$ is a near representation over $S$; its entries are locally constant, as those of $r'$ and $\psi'$ are, $\mathfrak m_{R'}$ being nilpotent, and its entries off the diagonal lie in $\mathfrak m_S$. By Lemma~\ref{sha-siegel-step-coefficients-universal}, a homomorphism of $\mathcal O'$-algebras $F:R\to S$ carries $\mathcal D$ to it. Then $e(F(z_s))=y_st^{\tau(w_s)}$ for the weight $w_s$ of $z_s$.
\item[(c)] \emph{The parts.} For $x\in R_w$, $e(F(x))=yt^{\tau(w)}$ for some $y\in R'$. Indeed, $x=\pi(g)$ for a series $g$ of weight $w$. Let $n\ne\tau(w)$, and $M$ as in Lemma~\ref{sha-siegel-step-coefficients-space-core-laurent} for $n$. The terms of $g$ of total degree at least $M$ map into $\mathfrak m_S^M$ under $F\circ\pi$, as $F$ maps the $z_s$ into $\mathfrak m_S$; the others form a polynomial whose monomials have weight $w$, and $e\circ F$ maps each of them to an element of $R'$ times $t^{\tau(w)}$, by (b) and the additivity of $\tau$. So the coefficient of $t^n$ in $e(F(x))$ is $0$.
\item[(d)] \emph{The Artinian case.} Still with $\mathfrak m_{R'}$ nilpotent, for each $w$ let $f_w:R_w\to R'$ send $x$ to the coefficient of $t^{\tau(w)}$ in $e(F(x))$. These maps are $\mathcal O'$-linear, $f_0(1)=1$, and $f_{w+w'}(xx')=f_w(x)f_{w'}(x')$ by (c). As $A_0=\bigoplus_wR_w$, they define a homomorphism of $\mathcal O'$-algebras $f:A_0\to R'$ with $e(F(x))=f(x)t^{\tau(w)}$ for $x\in R_w$. As $r_{\mathcal D}(g)_{ij}$ lies in the part of the weight of $(i,j)$ and $\psi_{\mathcal D,v,i}(s')$ in $R_0$, and $F$ maps them to the entries of $(r'',(\psi''_{v,i}))$, comparing coefficients gives $f(r_{\mathcal D}(g)_{ij})=r'(g)_{ij}$ and $f(\psi_{\mathcal D,v,i}(s'))=\psi'_{v,i}(s')$.
\item[(e)] \emph{The general case.} For $N\ge0$, the ring $R'_N=R'/\mathfrak m_{R'}^{N+1}$ has nilpotent maximal ideal, and the image of $(r',(\psi'_{v,i}))$ in it is $\tau$-near. By (d) there is $f_N:A_0\to R'_N$ carrying $(r_{\mathcal D},(\psi_{\mathcal D,v,i}))$ to it; so $f_N(z_s)$ is the image of $y_s$. By (a), $f_{N+1}$ followed by $R'_{N+1}\to R'_N$ is $f_N$. As $R'$ is complete, the $f_N$ define a homomorphism of $\mathcal O'$-algebras $f:A_0\to R'$, and it carries $(r_{\mathcal D},(\psi_{\mathcal D,v,i}))$ to $(r',(\psi'_{v,i}))$, as each $f_N$ does modulo $\mathfrak m_{R'}^{N+1}$.
\end{enumerate}
% (1) is sxContGr (ContGr.lean, 114 lines): wPres_maximalIdeal_gr_zero_eq_sup_span_monomials for m_B, part_filt (Lemma sha-automorphy-unramified-near-tangent-filtration) and sxContGr_cont. (2) is sxGenZero (GenZero.lean): (a) is sxEqMono (EqMono.lean, with eqMono_root for the roots d_i, Fr.Ok for the distinct residues) and sxMonoZero; (b) is sxZMono (SxZMono.lean, with Pm); (c) is sxExtGr0 (SxExtGr0.lean, with Alg.co_ext). (3) is sxNearA (QuotPowNil.lean): (a) is sxExtA; (b) is sxLaur and sxScaleRD (ScaleRD.lean) with hU.2, the universality of Fr.RUniv; (c) is sxLaurPart (LaurPart.lean, with truncTotal); (d) is sxLaurDesc (LaurDesc.lean, through DirectSum.toAlgebra) in sxNearArt (NearArt.lean); (e) is sxNearMap, hlev, hval and sxLimPt. About 1,500 lines in AdicExtensionMatrices and CommutativeAlgebra.
\end{proof}

\begin{lemma}\label{sha-siegel-step-coefficients-space}
Let $D$, $T$, $\mathcal O'$ and $\Phi$ be as in Lemma~\ref{sha-siegel-step-coefficients-standard}. Then there is a deformation space for $T$ over $\mathcal O'$ (Definition~\ref{sha-siegel-tower-space}). The same holds for the deformation problem $\Phi^\circ$ with the data of $\Phi$ except that every set $\mu_v$ is empty, with flagged pairs for $\Phi^\circ$ in Definition~\ref{sha-siegel-tower-space}.
\end{lemma}
\leanref{Development/AutomorphicArithmetic/AutomorphicVarieties/WeakShimuraTheory/CompactIntegralFields/SpaceEx.lean}{UpperTheory.V.Ws.spaceEx}
\leanref{Development/AutomorphicArithmetic/AutomorphicVarieties/WeakShimuraTheory/AdicExtensionMatrices/PreOfCore.lean}{UpperTheory.V.Ws.sxPreOfCore}
\leanref{Development/AutomorphicArithmetic/AutomorphicVarieties/WeakShimuraTheory/AdicExtensionMatrices/ChConj.lean}{UpperTheory.V.Ws.sxChart}
\leanref{Development/AutomorphicArithmetic/AutomorphicVarieties/WeakShimuraTheory/AdicExtensionMatrices/PairNear.lean}{UpperTheory.V.Ws.sxPairNear}
\leanref{Development/AutomorphicArithmetic/AutomorphicVarieties/WeakShimuraTheory/AdicExtensionMatrices/PreLine.lean}{UpperTheory.V.Ws.sxPreLine}
\begin{proof}
Let $\Psi$ be $\Phi$ or $\Phi^\circ$, with sets $\mu_v$. Both have the properties (1) to (5) of Lemma~\ref{sha-siegel-step-coefficients-standard}: $\Phi$ by that lemma, and $\Phi^\circ$ because (1) to (4) do not involve the sets $\mu_v$ and (5) is empty for it. Below, rings are complete local noetherian $\mathcal O'$-algebras with residue field $k'$ unless said otherwise, and homomorphisms between them are homomorphisms of $\mathcal O'$-algebras; such a homomorphism maps the maximal ideal into the maximal ideal.
\begin{enumerate}
\item[(a)] \emph{The universal near representation.} Near representations do not involve the sets $\mu_v$. Let $R_\Phi$ and $\mathcal D=(r_{\mathcal D},(\psi_{\mathcal D,v,i}))$ be as in Lemma~\ref{sha-siegel-step-coefficients-universal}.
\item[(b)] \emph{The grading.} Let $R'$ be a ring and $t=(t_1,t_2)\in(R'^\times)^2$ with $t_1\equiv t_2\equiv1$ modulo $\mathfrak m_{R'}$. Conjugation by $\tau(t)=\operatorname{diag}(t_1t_2,t_2,t_1,1)$ multiplies the entry $(i,j)$ by $t^w$, where $w$ is its weight (Definition~\ref{sha-counts-vd-model}), and $\tau(t)$ is symplectic with multiplier $t_1t_2$. So it carries a near representation $(r,(\psi_{v,i}))$ over $R'$ to the near representation $(\tau(t)r\tau(t)^{-1},(\psi_{v,i}))$: it keeps the multipliers, the diagonal entries, the diagonal matrix at $g_0^T$, the value $1$ on $N_\Phi$, the continuity and the residues. By the universal property, for every homomorphism $f:R_\Phi\to R'$ there is $f':R_\Phi\to R'$ with $f'(r_{\mathcal D}(g)_{ij})=t^wf(r_{\mathcal D}(g)_{ij})$ and $f'(\psi_{\mathcal D,v,i}(s))=f(\psi_{\mathcal D,v,i}(s))$, and two homomorphisms that agree at the entries of $r_{\mathcal D}$ and at the $\psi_{\mathcal D,v,i}$ are equal. Apply Lemma~\ref{sha-counts-vd-model-torus-grading}(1), with $\mathcal O'$ and $k'$ in place of $\mathcal O$ and $\mathbf F$, to the elements $r_{\mathcal D}(g)_{ij}-c$ of weight the weight of $(i,j)$, where $c$ is the image of $\chi_i(g)$ if $i=j$ and $c=0$ otherwise, and $\psi_{\mathcal D,v,i}(s)-\chi_i(s)$ of weight $0$; they lie in $\mathfrak m_{R_\Phi}$, as $\mathcal D$ is a near representation. This gives a weighted presentation $\pi$ of $R_\Phi$ as in Lemma~\ref{sha-siegel-step-coefficients-space-core}, constants lying in the part of weight $0$.
\item[(c)] \emph{The rings $B$ and $A_0$.} Let $B$ be the part of weight $0$ and $A_0=R_\Phi^{\mathrm{gr}}$. By Lemma~\ref{sha-counts-vd-presentations}(3), $B$ is a ring. By Lemma~\ref{sha-counts-vd-presentations}(2), $A_0$ is generated as an $\mathcal O'$-algebra by $B$ and the images of the finitely many variables, which are entries of $r_{\mathcal D}$ minus constants or lie in $B$; so $A_0$ is of finite type over $B$, generated by the entries of $r_{\mathcal D}$. The entries of $r_{\mathcal D}$ lie in $A_0$, and the $\psi_{\mathcal D,v,i}$ in $B$.
\item[(d)] \emph{The polynomials $c_g$.} Each term of the characteristic polynomial of $r_{\mathcal D}(g)$ is a product of entries $(i,\sigma(i))$ for a permutation $\sigma$, whose weights add up to $0$; and the product of elements of the parts of weights $w$ and $w'$ lies in the part of weight $w+w'$. So the characteristic polynomial is a polynomial $c_g\in B[X]$. Modulo $\mathfrak m_{R_\Phi}$, $r_{\mathcal D}(g)$ is $\operatorname{diag}(\bar\chi_1(g),\dots,\bar\chi_4(g))$, so $c_g$ reduces to $\prod_i(X-\bar\chi_i(g))$. Lemma~\ref{sha-siegel-step-coefficients-space-core} applies to $\pi$.
\item[(e)] \emph{The flags.} Let $\tilde A$ be the quotient of the polynomial ring over $A_0$ in variables $N_{v,ij}$, for $v\in\Sigma_T$ and $1\le i,j\le4$, by the ideal generated by: for each $v$, with $N_v=(N_{v,ij})$, the entries of $N_v^*N_v-1$, the $N_{v,ij}$ with $i<j$ and the $N_{v,ii}-1$; the entries $(i,j)$ with $i>j$ of $N_v^*r_{\mathcal D}(s)N_v$ and the $(N_v^*r_{\mathcal D}(s)N_v)_{ii}-\psi_{\mathcal D,v,i}(s)$, for $s\in D_v$; and the $\psi_{\mathcal D,v,i}(s)-1$ for $s\in\mu_v$. Let $n_{X,v}$ be the image of $N_v$ in $\tilde A$, and $P_X=(r_X,(n_{X,v}))$ with $r_X$ the image of $r_{\mathcal D}$. Then $\tilde A$ is of finite type over $B$, generated by the entries of $r_X$ and of the $n_{X,v}$, and $P_X$ is a flagged pair for $\Psi$ along $\mathcal O\to\tilde A$, whose diagonal characters are the images of the $\psi_{\mathcal D,v,i}$. For every ring homomorphism $f:A_0\to R$ to a commutative ring and matrices $n'_v\in M_4(R)$ such that $n'_v$ is symplectic with multiplier $1$ and lower triangular with diagonal entries $1$, $n_v'^*f(r_{\mathcal D}(s))n'_v$ is upper triangular with diagonal entries $f(\psi_{\mathcal D,v,i}(s))$ for $s\in D_v$, and $f(\psi_{\mathcal D,v,i}(s))=1$ for $s\in\mu_v$, there is a ring homomorphism $\tilde A\to R$ extending $f$ that carries $n_{X,v}$ to $n'_v$.
\item[(f)] \emph{Lifts of residual pairs.} Let $c=(r_c,(n_{c,v}))$ be a residual flagged pair for $\Psi$ along $\mathcal O\to k'$, and $Q=(r_Q,(n_{Q,v}))$ a lift of $c$ to a ring $R$, with diagonal characters $\psi_{Q,v,i}$. Let $\tau(w)=-2w_1-w_2$ if $r_c$ is supported in $\{(-1,1),(-1,0),(-1,-1),(0,-1)\}$, and $\tau(w)=-w_1-2w_2$ otherwise. A check of the two supports in Definition~\ref{sha-siegel-tower-pairs}(3) shows that $r_c(g)_{ij}=0$ whenever $i\ne j$ and $\tau(w)\le0$ for the weight $w$ of $(i,j)$. So $(r_Q,(\psi_{Q,v,i}))$ has the properties of Lemma~\ref{sha-siegel-step-coefficients-space-core}(3): it is symplectic with multiplier $\theta$, trivial on $N_\Phi$ and diagonal at $g_0^T$, as $Q$ is a flagged pair; its entries are continuous; $r_Q(g)_{ii}$ and $\psi_{Q,v,i}(s)$ have the residues $\bar\chi_i(g)$ and $\bar\chi_i(s)$, as $c$ is residual; and the entries $r_Q(g)_{ij}$ with $\tau(w)\le0$ reduce to $0$. That lemma gives $f:A_0\to R$ carrying $\mathcal D$ to $(r_Q,(\psi_{Q,v,i}))$, and (e), with $n'_v=n_{Q,v}$, extends it to $\tilde A\to R$ carrying $P_X$ to $Q$.
\item[(g)] \emph{$\tilde A$ is nonzero.} The pair $c_0=(\operatorname{diag}(\bar\chi_1,\dots,\bar\chi_4),(1))$ is a residual flagged pair for $\Psi$ along $\mathcal O\to k'$, by (1) and (5) of Lemma~\ref{sha-siegel-step-coefficients-standard}, and a lift of itself to $k'$. By (f) there is a homomorphism $\tilde A\to k'$, so $\tilde A\ne0$.
\item[(h)] \emph{The deformation space.} Let $\bar B$ be the image of $B$ in $\tilde A$. It is the quotient of $B$ by a proper ideal, by (g), so it is a ring of the kind above, and $\bar B\to\tilde A$ is injective. Let $X=(\tilde A,\bar B,P_X)$. Conditions (1) and (5) of Definition~\ref{sha-siegel-tower-space} hold by (e) and (f). The images of the $c_g$ satisfy (2), and the images of the $\psi_{\mathcal D,v,i}(s)$ satisfy (3), by (d) and (e), with the continuity of Lemma~\ref{sha-siegel-step-coefficients-space-core}(1); the same lemma gives the continuity of the entries of $r_X$ for the $\mathfrak m_{\bar B}\tilde A$-adic topology. For (4), two homomorphisms from $\bar B$ that agree at these elements give, composed with $B\to\bar B$, homomorphisms that agree at the coefficients of the $c_g$ and at the $\psi_{\mathcal D,v,i}(s)$, which are equal by Lemma~\ref{sha-siegel-step-coefficients-space-core}(2); as $B\to\bar B$ is surjective, they are equal.
\item[(i)] \emph{Condition (6).} Let $Q$ be a residual flagged pair for $\Psi$ along $\mathcal O\to k'\subset k'[X]$, and $j:k'[X]\to k'[[X]]$ the inclusion. The entries of $Q$ are locally constant, so $j(Q)$ is a lift to $k'[[X]]$ of its image modulo $X$, which is residual. By (f) there is $f:\tilde A\to k'[[X]]$ carrying $P_X$ to $j(Q)$. As $Q=(r_Q,(n_{Q,v}))$ is residual, each matrix $r_Q(g)$ is triangular for an ordering of the indices, with diagonal entries $\bar\chi_i(g)$, so the characteristic polynomial of $j(Q(g))$ is $\prod_i(X-\bar\chi_i(g))$, and the diagonal characters of $j(Q)$ are the $\bar\chi_i$. So $f$ restricted to $\bar B$ and the composite $\bar B\to k'\subset k'[[X]]$ agree at the generators of (4), and they are equal. Hence $f$ kills $\mathfrak m_{\bar B}$ and maps $\bar B$ into $k'$, and as $\tilde A$ is generated over $\bar B$ by the entries of $P_X$, which $f$ maps into $j(k'[X])$, $f=j\circ f_0$ for a homomorphism $f_0:\tilde A\to k'[X]$ that carries $P_X$ to $Q$ and kills $\mathfrak m_{\bar B}$.
\end{enumerate}
\end{proof}
% spaceEx (SpaceEx.lean, 108 lines), for a deformation problem with Fr.Std over a complete discrete valuation ring with finite residue field: (a) sxUnivD; (b) sxGradD (SubMemMax.lean, 101 lines: coord_grading, which is Lemma sha-counts-vd-model-torus-grading (1), with sxNearConj and sxNearMap); (c) wPres_gr_zero_isCO_of_mem_maximalIdeal, gr_noeth, sxSubRD and sxCoreGen; (d) sxCoreCp (CoreCp.lean); the structure Core collects (a) to (d) with sxContGr, sxGenZero and sxNearA (Lemma sha-siegel-step-coefficients-space-core); (e) is sxChart (ChConj.lean, 129 lines); (f) is sxPreOfCore (PreOfCore.lean, 90 lines) with sxPairNear (PairNear.lean, 61 lines: tau = lin (-2) (-1) or lin (-1) (-2)); (g) is sxSplitRes (SplitRes.lean) with Y.univ; (h) is sxPreImage (PreImage.lean, pre_of_surj) and Pre.toSpace; (i) is sxPreLine (PreLine.lean, 66 lines, with sxResLift and sxCpRes). The Lean proves it for every problem with Fr.Std, which Phi.noMu has (X1.defEx).

\begin{lemma}\label{sha-siegel-step-coefficients}
Let $D$ be a Siegel datum, let $V$ be an ambient ring with residue field $k_V$ (Definition~\ref{sha-siegel-tower-ambient}), and let $\Phi$ be a deformation problem for $D$ with field $E$ (Definition~\ref{sha-siegel-tower-pairs}). For a ring $\mathcal O'$ of Definition~\ref{sha-character-parts} write $k'$ for its residue field and $k'\to k_V$ for the map induced by $\mathcal O'\to V$.
\begin{enumerate}
\item If $\Phi$ is the problem of a fit tower $T$, then for every ring $\mathcal O'\supseteq\mathcal O$ of Definition~\ref{sha-character-parts} there is a deformation space for $T$ over $\mathcal O'$ (Definition~\ref{sha-siegel-tower-space}).
\item Let $c$ and $c'$ be flagged pairs for $\Phi$ along $\mathcal O\to k_V$, with $c$ residual and a chain from $c$ to $c'$. There is $\mathcal O_1\supseteq\mathcal O$ such that for every $\mathcal O'\supseteq\mathcal O_1$ there are flagged pairs $s$ and $t$ for $\Phi$ along $\mathcal O\to k'$ with images $c$ and $c'$ under $k'\to k_V$, with $s$ residual and a chain from $s$ to $t$.
\item Let $r:G_E\to M_4(V)$ be a homomorphism and $n_v\in M_4(V)$ matrices, with entries in a complete discrete valuation subring of $V$ with finite residue field through which $r$ factors continuously. There is $\mathcal O_1\supseteq\mathcal O$ such that for every $\mathcal O'\supseteq\mathcal O_1$ there are a complete local noetherian $\mathcal O'$-algebra $R$ with residue field $k'$, an injective ring homomorphism $f:R\to V$ that restricts to the given map on $\mathcal O'$ and induces $k'\to k_V$ on residue fields, and a homomorphism $r_R:G_E\to M_4(R)$ and matrices $n_{R,v}\in M_4(R)$ with $\mathfrak m_R$-adically continuous entries that $f$ carries to $r$ and the $n_v$.
\end{enumerate}
\end{lemma}
\leanref{Development/AutomorphicArithmetic/AutomorphicVarieties/WeakShimuraTheory/CompactIntegralFields/ContResQ.lean}{UpperTheory.V.Ws.std}
\leanref{Development/AutomorphicArithmetic/AutomorphicVarieties/WeakShimuraTheory/CompactIntegralFields/SpaceEx.lean}{UpperTheory.V.Ws.spaceEx}
\leanref{Development/AutomorphicArithmetic/AutomorphicVarieties/WeakShimuraTheory/CompactIntegralFields/ChainDesc.lean}{UpperTheory.V.Ws.chainDesc}
\leanref{Development/AutomorphicArithmetic/AutomorphicVarieties/WeakShimuraTheory/CompactIntegralFields/SmallCo.lean}{UpperTheory.V.Ws.smallCo}
\begin{proof}
Any two rings of Definition~\ref{sha-character-parts} lie in a third, and the maps from these rings to $V$ are compatible with the inclusions. So if an element of $k_V$ is the image of an element of the residue field of a ring $\mathcal O''$, it is also the image of an element of the residue field of every ring $\mathcal O'\supseteq\mathcal O''$; and by (1) of Definition~\ref{sha-siegel-tower-ambient}, every element of $k_V$ is the image of an element of the residue field of some ring $\mathcal O''$.

(1) Let $\mathcal O'\supseteq\mathcal O$. The problem of the fit tower $T$ has the properties of Lemma~\ref{sha-siegel-step-coefficients-standard} for $\mathcal O'$, and from them Lemma~\ref{sha-siegel-step-coefficients-space} gives a deformation space for $T$ over $\mathcal O'$.

(2) Let $c=c_0,c_1,\dots,c_m=c'$ be the chain, and for $1\le i\le m$ let $Q_i$ be a residual flagged pair for $\Phi$ along $\mathcal O\to k_V[X]$ whose images under $X\mapsto0$ and $X\mapsto1$ are $c_{i-1}$ and $c_i$. The entries of $c$, and the coefficients of the entries of the $Q_i$, are finitely many elements of $k_V$: the entries of the homomorphisms are locally constant on the compact group $G_E$, so they take finitely many values, and there are finitely many places $v$. By the first paragraph there is $\mathcal O_1\supseteq\mathcal O$ such that for every $\mathcal O'\supseteq\mathcal O_1$ all these elements lie in the image of $k'\to k_V$.
Fix such an $\mathcal O'$. The map $\iota:k'\to k_V$ is injective, as a homomorphism of fields, and so is the map $k'[X]\to k_V[X]$ it induces; the composite of $\mathcal O\to k'$ with $\iota$ is $\mathcal O\to k_V$. A pair $(r,(n_v))$ over $k_V$, or over $k_V[X]$, whose entries lie in the image of $k'$, or of $k'[X]$, is the image of a unique pair over $k'$, or $k'[X]$, and if $r$ is a homomorphism so is the pair it comes from, by the injectivity. The conditions that define flagged pairs, residual pairs, pairs of type V and ample pairs are identities between polynomials in the entries and in the images of the values of the characters, the vanishing or nonvanishing of entries, and the local constancy of entries; an injective ring homomorphism preserves and reflects all of them. So $c$ and the $c_i$ are the images of flagged pairs $s$ and $t_i$ along $\mathcal O\to k'$, and the $Q_i$ of residual flagged pairs $Q'_i$ along $\mathcal O\to k'[X]$, with $s=t_0$ residual and every $t_i$ ample. Evaluation at $0$ and at $1$ commutes with $\iota$, so by uniqueness the images of $Q'_i$ under $X\mapsto0$ and $X\mapsto1$ are $t_{i-1}$ and $t_i$. So $t_0,\dots,t_m$ is a chain from $s$ to $t=t_m$, and $s$ and $t$ have images $c$ and $c'$.

(3) Let $A\subset V$ be the complete discrete valuation ring of the statement, containing the entries of $r$ and the $n_v$, with $r$ continuous for the $\mathfrak m_A$-adic topology; $\mathfrak m_A\subseteq\mathfrak m_V$, and $\Omega_V$ is algebraic over the fraction field of $A$ (Definition~\ref{sha-siegel-tower-ambient}).
\begin{enumerate}
\item[(a)] \emph{The ring $\mathcal O_1$.} The residue field of $A$ is finite. Choosing a representative in $A$ of each residue class, the first paragraph gives $\mathcal O_1\supseteq\mathcal O$ such that for every $\mathcal O'\supseteq\mathcal O_1$ each representative is congruent modulo $\mathfrak m_V$ to the image of an element of $\mathcal O'$. As $\mathfrak m_A\subseteq\mathfrak m_V$, every $x\in A$ then has this property.
Fix $\mathcal O'\supseteq\mathcal O_1$, with residue field $k'$.
\item[(b)] \emph{The image of $\mathbf Z_\ell$ lies in $A$.} The image of $\ell$ in $V$ lies in $\mathfrak m_V$, as $\mathcal O'\to V$ is local, so $\ell$ is not a unit of $A$. Let $t\in\mathbf Z_\ell$, and let $t_n$ be integers with $t\equiv t_n\pmod{\ell^n}$. As $A$ is complete, the $t_n$ converge in $A$ to some $\mu$. For a uniformizer $\varpi$ of $A$ and every $k$, $\mu-t_k\in\varpi^kA$, and $t_k-t\in\ell^k\mathbf Z_\ell$ maps into $\ell^kV\subseteq\varpi^kV$; so the difference $d=\mu-t$, computed in $V$, lies in $\varpi^kV$ for every $k$. Suppose that $d\ne0$. As $d$ is algebraic over $A$, there is $p\in A[X]$ with $p(d)=0$ and $p(0)\ne0$; write $p(0)=u\varpi^e$ with $u\in A^\times$, $p(0)=-d\,q(d)$ with $q\in A[X]$, and $d=\varpi^{e+1}w$ with $w\in V$. Then $u=-\varpi wq(d)$, so $\varpi$ would be a unit of $V$, contrary to $\mathfrak m_A\subseteq\mathfrak m_V$. So $d=0$.
\item[(c)] \emph{A larger ring.} The ring $\mathcal O'$ is generated as a $\mathbf Z_\ell$-module by finitely many elements; let $s$ be the set of their images in $\Omega_V$. By (b), the image of $\mathcal O'$ lies in the subfield generated by $A$ and $s$, and it lies in $V$. Lemma~\ref{sha-siegel-step-coefficients-extension} gives a complete discrete valuation ring $A'\subset V$ with finite residue field, with $\mathfrak m_{A'}\subseteq\mathfrak m_V$, that contains $A$ and the image of $\mathcal O'$. The homomorphism $\mathcal O'\to A'$ is local, since $\mathcal O'\to V$ is, and $\mathfrak m_{A'}=A'\cap\mathfrak m_V$.
\item[(d)] \emph{The ring $R$.} Lemma~\ref{sha-siegel-step-coefficients-subring}, applied to $\mathcal O'\to A'$, gives the subring $R$ of the $x\in A'$ congruent modulo $\mathfrak m_{A'}$ to the image of an element of $\mathcal O'$; it is a complete local noetherian $\mathcal O'$-algebra with residue field $k'$, and $\mathfrak m_R^N=R\cap\mathfrak m_{A'}^N$. By (a), every $x\in A$ is congruent modulo $\mathfrak m_V$, hence modulo $A'\cap\mathfrak m_V=\mathfrak m_{A'}$, to the image of an element of $\mathcal O'$; so $A\subseteq R$, and $\mathfrak m_A^N\subseteq R\cap\mathfrak m_{A'}^N=\mathfrak m_R^N$.
\item[(e)] \emph{Conclusion.} Let $f:R\to V$ be the inclusion. It restricts to the given map on $\mathcal O'$. An element of $R$ is congruent modulo $\mathfrak m_R\subseteq\mathfrak m_V$ to the image of an element of $\mathcal O'$, so $f$ induces $k'\to k_V$ on residue fields. The entries of $r$ and of the $n_v$ lie in $A\subseteq R$; let $r_R$ and $n_{R,v}$ be $r$ and the $n_v$ viewed over $R$. Since $\mathfrak m_A^N\subseteq\mathfrak m_R^N$, the entries of $r_R$ are $\mathfrak m_R$-adically continuous, and $f$ carries $r_R$ and the $n_{R,v}$ to $r$ and the $n_v$.
\end{enumerate}
% (1) is std (Lemma sha-siegel-step-coefficients-standard) and spaceEx (Lemma sha-siegel-step-coefficients-space, with Lemmas sha-siegel-step-coefficients-universal and sha-siegel-step-coefficients-space-core), as in last_step: (std hC D T hF ∅ c' h).1 and spaceEx (T.fr ∅) l (inc h). These are about 2,600 and 2,200 lines of Lean.
% (2) is chainDesc (ChainDesc.lean): the private ev_pd, ev_k, ev_poly and ev_fin for the rings, desc_ex for the descent, and isRes_of, good_of, inV_of for the conditions; the rings are combined with HV.coefSup.
% (3) is smallCo (SmallCo.lean): (a) is the private exists_coef_forall_sub_mem_maximalIdeal_of_isDvr, (b) exists_map_eq_algebraMap_padicInt_of_isDvr, (c) exists_finset_coef_mem_subfield_closure_of_isDvr with smallExt, (d) CA.cl_sub. Size of the lemma: about 4,250 lines of Lean, of which (2) and (3) are about 600.
\end{proof}

\begin{lemma}\label{sha-siegel-step-linked}
Let $D$ be a Siegel datum, let $T$ be a fit tower for $D$ (Definition~\ref{sha-siegel-tower-towers}), let $X=(A,B,P_X)$ be a deformation space for $T$ over a ring $\mathcal O'$ with residue field $k'$, and let $J\subseteq B$ be an ideal (Definition~\ref{sha-siegel-tower-space}). That $T$ is fit is not used in the proof.
\begin{enumerate}
\item Let $x$ and $x'$ be points of $X$ whose images are linked. If $J$ lies in all minimal primes below $\ker x$, then it lies in some minimal prime below $\ker x'$; and the same holds with $x$ and $x'$ exchanged.
\item Let $y:A\to K'$ be a ring homomorphism to a domain that detects $J$, and let $\mathfrak p\supseteq\ker y$ be a prime of $A$. Then $J$ lies in some minimal prime below $\mathfrak p$.
\item Let $y:A\to K'$ be a ring homomorphism to a domain, and let $\mathfrak p\supseteq\ker y$ be a prime of $A$ such that $J$ lies in all minimal primes below $\mathfrak p$. Then $y$ kills $J$.
\item Suppose that for every point $x$ of $X$ whose image is residual and ample, if $J$ lies in some minimal prime below $\ker x$, then it lies in all of them. Let $x_0$ and $x_1$ be points of $X$ such that the image of $x_0$ is residual, there is a chain from the image of $x_0$ to the image of $x_1$, and $J$ lies in some minimal prime below $\ker x_1$. Then $J$ lies in all minimal primes below $\ker x_0$.
\end{enumerate}
\end{lemma}
\leanref{Development/AutomorphicArithmetic/AutomorphicVarieties/WeakShimuraTheory/CompactIntegralFields/AlgHomResUnique.lean}{UpperTheory.V.Ws.conn}
\leanref{Development/AutomorphicArithmetic/AutomorphicVarieties/WeakShimuraTheory/SomeInOfLocZero.lean}{UpperTheory.V.Ws.Asm.someIn_of_locZero}
\leanref{Development/AutomorphicArithmetic/AutomorphicVarieties/WeakShimuraTheory/SomeInOfLocZero.lean}{UpperTheory.V.Ws.Asm.zero_of_allIn}
\leanref{Development/AutomorphicArithmetic/AutomorphicVarieties/WeakShimuraTheory/SomeInOfLocZero.lean}{UpperTheory.V.Ws.Asm.allIn_of_chain}
\begin{proof}
First, two points $x$ and $x'$ with the same image are equal. A homomorphism of $\mathcal O'$-algebras $z:B\to k'$ is surjective, since $\mathcal O'\to k'$ is, so its kernel is the maximal ideal; as every $b\in B$ is congruent modulo $\mathfrak m_B$ to an element of $\mathcal O'$, $z(b)$ is the residue of $b$. So $x$ and $x'$ agree on $B$, and they agree on the entries of $P_X$, which generate $A$ over $B$.

(1) By (6) of Definition~\ref{sha-siegel-tower-space}, the residual flagged pair over $k'[X]$ that links the two images is the image of $P_X$ under a homomorphism of $\mathcal O'$-algebras $f:A\to k'[X]$. Composed with $X\mapsto0$, $f$ is a point with the same image as $x$, so it is $x$, and $\ker f\subseteq\ker x$; in the same way $\ker f\subseteq\ker x'$. The ideal $\ker f$ is prime, since $k'[X]$ is a domain, so it contains a minimal prime $\mathfrak q$ of $A$. If $J$ lies in all minimal primes below $\ker x$, then $JA\subseteq\mathfrak q$, and $\mathfrak q\subseteq\ker x'$; this is the first claim, and the second is the same argument.

(2) The ideal $\ker y$ is prime and contains a minimal prime $\mathfrak q$ of $A$. The inverse image of $\mathfrak q$ in $B$ is a prime whose image in $A$ is killed by $y$, so it contains $J$. Hence $JA\subseteq\mathfrak q\subseteq\ker y\subseteq\mathfrak p$.

(3) Let $\mathfrak q\subseteq\ker y$ be a minimal prime of $A$. Then $\mathfrak q\subseteq\mathfrak p$, so $JA\subseteq\mathfrak q\subseteq\ker y$.

(4) Let $c_0,\dots,c_m$ be the chain, with $c_0$ the image of $x_0$ and $c_m$ that of $x_1$. Each $c_i$ is ample, by the definition of a chain, and residual: $c_0$ by hypothesis, and $c_i$ with $i\ge1$ as the image under $X\mapsto1$ of a residual flagged pair over $k'[X]$. We show, by descending induction on $i$, that $J$ lies in all minimal primes below $\ker x$ for every point $x$ with image $c_i$.
For $i=m$, the hypothesis applied to $x_1$ shows that $J$ lies in all minimal primes below $\ker x_1$. The constant family links $c_m$ to itself, so by (1) $J$ lies in some minimal prime below $\ker x$, and by the hypothesis applied to $x$, in all of them.
For $i<m$, by (6) of Definition~\ref{sha-siegel-tower-space} the flagged pair over $k'[X]$ that links $c_i$ to $c_{i+1}$ is the image of $P_X$ under some $f:A\to k'[X]$, and $f$ composed with $X\mapsto1$ is a point $x'$ with image $c_{i+1}$. By induction $J$ lies in all minimal primes below $\ker x'$; by (1) it lies in some minimal prime below $\ker x$, and by the hypothesis applied to $x$, in all of them.
The case $i=0$, $x=x_0$, is the claim.
\end{proof}

\begin{lemma}\label{sha-siegel-step-completion-facts-bockstein}
Let $\ell$ be a prime, $\Gamma$ a profinite group, $H\subseteq\Gamma$ a closed subgroup and $C$ the closure of the commutator subgroup of $H$. Suppose that the group $H^2(H,\mathbf Z/\ell)$ of continuous cochains, for the trivial action of $H$ on $\mathbf Z/\ell$, is finite with at most $\ell$ elements. Let $s,t\in H$ with $s^\ell\in C$, $t^\ell\in C$ and $s\notin C$. Then $ts^{-n}\in C$ for some $n\ge0$.
\end{lemma}
\leanref{Development/AutomorphicArithmetic/AutomorphicVarieties/WeakShimuraTheory/AdjointCoefficients/MuBock.lean}{UpperTheory.V.Ws.X1.muBock}
\begin{proof}
The group $H$ is profinite, $C$ is a closed normal subgroup, and $H/C$ is commutative. A locally constant homomorphism from $H$ to a commutative group vanishes on $C$, as its kernel is open, hence closed, and contains the commutators. Suppose that there is no such $n$. As $s^\ell\in C$, the image of $s$ in $H/C$ generates a finite subgroup; let $K\subseteq H$ be its inverse image, a closed subgroup that contains $C$ and $s$. Then $t\notin K$: otherwise $t\in s^mC$ for an integer $m$, which can be taken $\ge0$ as $s^\ell\in C$.
\begin{enumerate}
\item[(a)] \emph{Characters.} Let $K'\supseteq C$ be a closed subgroup of $H$ and $g\in H\setminus K'$. Then some locally constant homomorphism $\chi:H\to\mathbf Q/\mathbf Z$ vanishes on $K'$ and has $\chi(g)\ne0$. Indeed, $K'$ is the intersection of the open subgroups of $H$ that contain it, so one of them, $N$, does not contain $g$; $N$ is normal and $H/N$ is finite and commutative, as $N\supseteq C$; and a finite commutative group has a homomorphism to $\mathbf Q/\mathbf Z$ that does not vanish at a given element other than $1$. Let $\chi_1$ vanish on $C$ with $\chi_1(s)\ne0$, and $\chi_2$ vanish on $K$ with $\chi_2(t)\ne0$. As $s^\ell$ and $t^\ell$ lie in $C\subseteq K$, $\chi_1(s)$ and $\chi_2(t)$ have order $\ell$.
\item[(b)] \emph{Bockstein cocycles.} Let $\sigma:\mathbf Q/\mathbf Z\to\mathbf Q$ be a section of the quotient map, and for $i=1,2$ let $a_i=\sigma\circ\chi_i:H\to\mathbf Q$ and $c_i(g,h)=a_i(g)+a_i(h)-a_i(gh)$. Then $a_i$ and $c_i$ are locally constant, $c_i$ takes integer values, as $\chi_i$ is a homomorphism, and $c_i(h,k)-c_i(gh,k)+c_i(g,hk)-c_i(g,h)=0$, by expanding. So the reduction $\bar c_i$ of $c_i$ modulo $\ell$ is a continuous $2$-cocycle with values in $\mathbf Z/\ell$.
\item[(c)] \emph{Independence.} Let $i,j\in\mathbf Z$ be such that $i\bar c_1+j\bar c_2$ is a coboundary: there is a continuous map $f:H\to\mathbf Z/\ell$ with $f(g)+f(h)-f(gh)\equiv ic_1(g,h)+jc_2(g,h)$ modulo $\ell$ for all $g,h$. Let $\tilde f(g)\in\{0,\dots,\ell-1\}$ represent $f(g)$, and let $\psi(g)$ be the class in $\mathbf Q/\mathbf Z$ of $(ia_1(g)+ja_2(g)-\tilde f(g))/\ell$. Then $\psi$ is locally constant, and it is a homomorphism: $\ell$ times a representative of $\psi(g)+\psi(h)-\psi(gh)$ is $ic_1(g,h)+jc_2(g,h)-(\tilde f(g)+\tilde f(h)-\tilde f(gh))$, which lies in $\ell\mathbf Z$. Also $\ell\psi=i\chi_1+j\chi_2$. As $\psi$ vanishes on $C$, every $x\in H$ with $x^\ell\in C$ has $i\chi_1(x)+j\chi_2(x)=\ell\psi(x)=\psi(x^\ell)=0$. For $x=s$, $\chi_2(s)=0$, so $i\chi_1(s)=0$, and $\ell$ divides $i$ by (a). For $x=t$, $i\chi_1(t)=(i/\ell)\chi_1(t^\ell)=0$, so $j\chi_2(t)=0$, and $\ell$ divides $j$.
\item[(d)] \emph{The count.} By (c), the classes of the cocycles $p\bar c_1+q\bar c_2$ with $p,q\in\mathbf Z/\ell$ are $\ell^2$ distinct elements of $H^2(H,\mathbf Z/\ell)$. As $\ell^2>\ell$, this contradicts the hypothesis.
\end{enumerate}
\end{proof}
% muBock (MuBock.lean, 221 lines): core, with C the closure of the commutators in H, K the inverse image of the subgroup generated by s; (a) is the private sep (ProfiniteGrp.closedSubgroup_eq_sInf_open and CharacterModule.exists_character_apply_ne_zero_of_ne_zero), (b) is lift and cocycle, (c) is divide, key, hi0 and hj0 with dvd_of_smul, (d) is hF and the final count.

\begin{lemma}\label{sha-siegel-step-completion-facts-cyclic}
Let $\ell$ be a prime, $\Gamma$ a profinite group, $H\subseteq\Gamma$ a closed subgroup and $C$ the closure of the commutator subgroup of $H$. Suppose that for all $s,t\in H$ with $s^\ell\in C$, $t^\ell\in C$ and $s\notin C$ there is $n\ge0$ with $ts^{-n}\in C$. Then there is $\zeta\in H$ with $\zeta^{\ell^a}\in C$ for some $a\ge0$ such that every $s\in H$ with $s^{\ell^b}\in C$ for some $b\ge0$ has $s\zeta^{-n}\in C$ for some $n\ge0$.
\end{lemma}
\leanref{Development/AutomorphicArithmetic/AutomorphicVarieties/WeakShimuraTheory/AdjointCoefficients/GenOfMax.lean}{UpperTheory.V.Ws.X1.muCyc}
\begin{proof}
The quotient $A=H/C$ is a commutative group. In $A$ the hypothesis says that if $x,y\in A$ satisfy $x^\ell=y^\ell=1$ and $x\ne1$, then $y=x^n$ for some $n\ge0$; and it suffices to find $z\in A$ of $\ell$-power order such that every element of $\ell$-power order of $A$ is a power $z^n$ with $n\ge0$, and to lift $z$ to $\zeta\in H$.
\begin{enumerate}
\item[(a)] \emph{Subgroups of finite index.} Every $x\ne1$ in $A$ lies outside some subgroup of finite index. Let $x$ be the image of $g\in H\setminus C$. As $C$ is closed, the $h\in H$ with $gh\notin C$ form an open neighbourhood of $1$, which contains an open normal subgroup $N$ of $H$. The image of $N$ in $A$ has finite index, and does not contain $x$: if $x$ were the image of $h\in N$, then $gh^{-1}\in C$, whereas $h^{-1}\in N$.
\item[(b)] \emph{No element of order $\ell$.} If $A$ has no element of order $\ell$, then $z=1$ will do: by induction on $b$, an element $y$ with $y^{\ell^b}=1$ is $1$, as $y^{\ell^{b-1}}$ has order dividing $\ell$.
\item[(c)] \emph{Bounded orders.} Otherwise, let $x\in A$ have order $\ell$, and let $U$ be a subgroup of finite index $m$ that does not contain $x$, by (a). If some $y\in A$ has order $\ell^{a+1}$, then $a\le m$. Indeed, suppose that $a>m$. Then $y^{\ell^a}$ has order $\ell$, so $x$ is a power of it. The image of $y$ in $A/U$ has order $\ell^j$ dividing $\ell^{a+1}$, and $\ell^j\le m$, as the order divides $m$. So $j<\ell^j\le m<a$, the order $\ell^j$ divides $\ell^a$, and $y^{\ell^a}\in U$; hence $x\in U$, a contradiction.
\item[(d)] \emph{The generator.} By (c), there is a largest $a$ such that $A$ has an element $z$ of order $\ell^{a+1}$; then every element of $\ell$-power order has order dividing $\ell^{a+1}$. We show, by induction on $b\le a+1$, that every $y$ with $y^{\ell^b}=1$ is a power of $z$. This is clear for $b=0$. For $b\ge1$, $y^{\ell^{b-1}}$ has order dividing $\ell$, and $z^{\ell^a}$ has order $\ell$, so $y^{\ell^{b-1}}=z^{\ell^an}$ for some $n\ge0$. Writing $a=c+b-1$, the element $y'=yz^{-\ell^cn}$ has $y'^{\ell^{b-1}}=1$, so $y'$ is a power of $z$ by induction, and so is $y$. As $z$ has finite order, every power of $z$ is a power $z^n$ with $n\ge0$.
\end{enumerate}
\end{proof}
% muCyc (GenOfMax.lean, 137 lines): (a) is the openness argument with ProfiniteGrp.exist_openNormalSubgroup_sub_open_nhds_of_one; (b) to (d) are the private cyc_abs (with Nat.findGreatest for the largest a and the bound hbd for (c)) and gen_of_max for the induction of (d).

\begin{lemma}\label{sha-siegel-step-completion-facts}
Let $D$ be a Siegel datum, let $T$ be a fit deep tower for $D$ with field $F$ (Definition~\ref{sha-siegel-tower-towers}), and let $\Phi$ be the problem of $T$, with sets $N_\Phi$ and $\mu_v$ for $v\in\Sigma_T$ (Definition~\ref{sha-siegel-tower-pairs}).
Write $G_F=\operatorname{Gal}(\bar{\mathbf Q}/F)$, let $\bar\omega_\ell:G_F\to\mathbf F_\ell^\times$ be the reduction of $\chi_{\mathrm{cyc}}$ modulo $\ell$, and let $K_T$ be the set of $g\in G_F$ with $\bar r_0(g)=1$.
\begin{enumerate}
\item $\theta(g)=1$ for every $g\in N_\Phi$, and $\theta(s)=1$ for every $s\in\mu_v$ and $v\in\Sigma_T$.
\item Every complex conjugation $\tau\in G_F$ has $\tau^2=1$ and $\theta(\tau)=-1$.
\item Let $k$ be a field of characteristic $\ell$ and $a:\mathcal O\to k$ a homomorphism that kills $\mathfrak m_{\mathcal O}$. For every $w\in\mathbf Z^2$ some $g\in G_F$ has $a(\chi^w(g))\ne\bar\omega_\ell(g)$.
\item The group $K_T$ is an open normal subgroup of $G_F$ of index prime to $\ell$. Let $S$ be the set of places of $F$ above $\ell$, $N_T$ the closed normal subgroup of $G_F$ generated by $N_\Phi$, $\Gamma_T=G_F/N_T$ and $\pi_T:G_F\to\Gamma_T$ the quotient map. Then $\ell$, $F$, $S$, $K_T$, $\Gamma_T$ and $\pi_T$ are as in Situation~\ref{sha-descent-link-setting}, and the map that sends $v\in\Sigma_T$ to the place of $F$ below it is a bijection from $\Sigma_T$ onto $S$.
\item For every $v\in\Sigma_T$ there is $\zeta_v\in\mu_v$ such that every $s\in\mu_v$ is the product of an element of the closure of the commutator subgroup of $D_v(F)$ and a power $\zeta_v^n$ with $n\ge0$.
\end{enumerate}
\end{lemma}
\leanref{Development/AutomorphicArithmetic/AutomorphicVarieties/WeakShimuraTheory/AdjointCoefficients/ThFacts.lean}{UpperTheory.V.Ws.X1.thFacts}
\leanref{Development/AutomorphicArithmetic/AutomorphicVarieties/WeakShimuraTheory/AdjointCoefficients/NotDvdIndex.lean}{UpperTheory.V.Ws.X1.linkT}
\leanref{Development/AutomorphicArithmetic/AutomorphicVarieties/WeakShimuraTheory/AdjointCoefficients/MuLoc.lean}{UpperTheory.V.Ws.X1.muLoc}
\begin{proof}
(1) Let $g\in N_\Phi$, so that $g\in I_w(F)$ for a place $w$ not above $\ell$ and $\bar r_0(g)=1$. Then $\chi_{\mathrm{cyc}}(g)=1$, as $\chi_{\mathrm{cyc}}$ is unramified at $w$, and $\psi(g)=1$ by (i). As $\omega^{\ell-1}=1$ and $\bar\omega(g)=\bar\chi_{\mathrm{cyc}}(g)=1$, Lemma~\ref{sha-siegel-step-towers-kappa-units}(1) gives $\omega(g)=1$. The entry $(3,3)$ of $\bar r_0(g)$ is $\bar\alpha(g)=\bar\xi(g)$, so $\bar\xi(g)=1$; as $\xi$ has order prime to $\ell$, the same lemma gives $\xi(g)=1$. Hence $\theta(g)=\chi_{\mathrm{cyc}}(g)^{k+1}\omega(g)^{-1}\xi(g)\psi(g)=1$.

Let now $s\in\mu_v$, so that $s^{\ell^a}$ lies in the closure $C_v$ of the commutator subgroup of $D_v(F)$ for some $a\ge0$. A continuous character of $G_F$ with values in a commutative group and with closed kernel is trivial on $C_v$, so its value at $s$ has order dividing $\ell^a$. As $\ell$ is odd, $\mathbf Z_\ell^\times$ has no element of order $\ell$, so $\chi_{\mathrm{cyc}}(s)=1$ and $\theta(s)=(\omega^{-1}\xi\psi)(s)$. The character $\omega^{-1}\xi\psi$ has open kernel, so $(\omega^{-1}\xi\psi)(s)^{\ell^a}=1$. Also $\omega^{\ell-1}=1$, $\xi^{m_1}=1$ for some $m_1$ prime to $\ell$, and, as $T$ is deep, $\psi^{m_2}=1$ on $G_F$ for some $m_2$ prime to $\ell$; so $(\omega^{-1}\xi\psi)(s)^{(\ell-1)m_1m_2}=1$. As $\ell^a$ and $(\ell-1)m_1m_2$ are coprime, $\theta(s)=1$.

(2) Clearly $\tau^2=1$. Any two complex conjugations are conjugate in $G$, and $\theta$ takes values in a commutative group, so $\theta(\tau)=\theta(c)=-1$ by Lemma~\ref{sha-flagged-residues}(3).

(3) By (e), some $g\in G_L\subseteq G_F$ has $\bar\chi^w(g)\ne\bar\omega(g)$ in $\mathbf F$. As $a$ kills $\mathfrak m_{\mathcal O}$, it induces an injective homomorphism $\mathbf F\to k$. So $a(\chi^w(g))\ne a(\omega(g))$, and $a(\omega(g))$ is the image of $\bar\omega(g)=\bar\chi_{\mathrm{cyc}}(g)$, which is $\bar\omega_\ell(g)$.

(4) As $\bar r_0$ is diagonal, $K_T$ is the kernel of the homomorphism $g\mapsto(\bar\delta_1(g),\dots,\bar\delta_4(g))$ to $(\mathbf F^\times)^4$, which is locally constant; so $K_T$ is open and normal in $G_F$. Its index is the order of the image, a finite subgroup of $(\mathbf F^\times)^4$. This group has no element of order $\ell$, as $x^\ell=1$ gives $(x-1)^\ell=0$ in $\mathbf F$; so $\ell$ does not divide the index. The set $N_\Phi$ is the union of the groups $I_A\cap K_T$ for the places $A$ of $\bar{\mathbf Q}$ above places of $F$ outside $S$, so $N_T$ is the closed normal subgroup generated by these groups. Every place of $\bar{\mathbf Q}$ above $\ell$ is conjugate under $G_F$ to exactly one element of $\Sigma_T$, by the definition of a tower, and two places of $\bar{\mathbf Q}$ above the same place of $F$ are conjugate under $G_F$; so the map $\Sigma_T\to S$ is bijective.

(5) Let $v\in\Sigma_T$, and let $C_v$ be the closure of the commutator subgroup of $D_v(F)$, which is a closed subgroup of the profinite group $G_F$. By Lemma~\ref{sha-descent-link-local}(1), applied to $D_v(F)$, $k=\mathbf F_\ell$ and $M=\mathbf F_\ell$ with the trivial action, $\dim H^2(D_v(F),\mathbf F_\ell)$ is the dimension of a space of linear forms on $\mathbf F_\ell$, which is at most $1$; so $H^2(D_v(F),\mathbf Z/\ell)$ has at most $\ell$ elements. By Lemma~\ref{sha-siegel-step-completion-facts-bockstein}, $H=D_v(F)$ satisfies the hypothesis of Lemma~\ref{sha-siegel-step-completion-facts-cyclic}, which gives $\zeta_v\in D_v(F)$ with $\zeta_v^{\ell^a}\in C_v$ for some $a$, so that $\zeta_v\in\mu_v$, such that every $s\in\mu_v$ has $s\zeta_v^{-n}\in C_v$ for some $n\ge0$. Then $s$ is the product of $s\zeta_v^{-n}\in C_v$ and $\zeta_v^n$.
\end{proof}
% thFacts (ThFacts.lean, 303 lines): (1) is nu_eq_one_of_mem_n_of_good_of_ok and nu_eq_one_of_mem_mu_of_good_of_deep (mu_pow, no_tors, one_of_cop), (2) is mul_self_eq_one_and_nu_neg_one_of_isCC with hom_eq_of_isCC_of_isCC, (3) is exists_rt_ne_cycF_of_good_of_ok. (4) is linkT (NotDvdIndex.lean, 118 lines) with not_dvd_index and clos_eq; Gamma_T is G_F modulo GC.tameN S K_T, and GC.gLink_quot gives GLink. (5) is muLoc (MuLoc.lean, 31 lines): Kum.placeDict for the place of F below v, GC.locEuler with M = ZMod l and the trivial action for |H^2(D_v, F_l)| <= l (e2 and Subspace.dual_finrank_eq), then muCyc (Lemma sha-siegel-step-completion-facts-cyclic) applied to muBock (Lemma sha-siegel-step-completion-facts-bockstein).

\begin{lemma}\label{sha-siegel-step-completion-adjoint}
Let $k$ be a field in which $2\ne0$, and let $\mathfrak{sp}_4(k)$ and $X\mapsto X^*$ be as in Definition~\ref{sha-siegel-datum}. For $n\in M_4(k)$ let $\mathfrak b(n)$ be the space of $X\in\mathfrak{sp}_4(k)$ with $n^*Xn$ upper triangular, and $\mathfrak u(n)$ the space of $X\in\mathfrak{sp}_4(k)$ with $(n^*Xn)_{ij}=0$ for $i\le j$. Let $\mathfrak t$ be the space of diagonal elements of $\mathfrak{sp}_4(k)$.
\begin{enumerate}
\item $\dim\mathfrak{sp}_4(k)=10$ and $\dim\mathfrak t=2$. If $n$ is symplectic with multiplier $1$, then $\dim\mathfrak b(n)=6$ and $\mathfrak{sp}_4(k)=\mathfrak b(n)\oplus\mathfrak u(n)$.
\item Let $H$ be a group and $r:H\to M_4(k)$ a homomorphism such that each $r(h)$ is symplectic with some multiplier. Then $h\cdot X=r(h)Xr(h)^{-1}$ is a linear action of $H$ on $\mathfrak{sp}_4(k)$. Let $h\in H$. If $n$ is symplectic with multiplier $1$ and $n^*r(h)n$ and $n^*r(h^{-1})n$ are upper triangular, then $h$ preserves $\mathfrak b(n)$. If $h^2=1$ and $r(h)$ is symplectic with multiplier $-1$, then the $X\in\mathfrak{sp}_4(k)$ with $h\cdot X=X$ form a space of dimension at most $4$.
\item Let $D$ be a Siegel datum, $T$ a fit tower for $D$ with field $F$ (Definition~\ref{sha-siegel-tower-towers}), $a:\mathcal O\to k$ a homomorphism, and $c=(r,(n_v))$ a residual flagged pair for the problem of $T$ along $a$ (Definition~\ref{sha-siegel-tower-pairs}). Let $\omega':G_F\to k$ be a map such that for every $w\in\mathbf Z^2$ some $g\in G_F$ has $a(\chi^w(g))\ne\omega'(g)$. Then every linear form $f$ on $\mathfrak{sp}_4(k)$ with $f(r(g)Xr(g)^{-1})=\omega'(g)f(X)$ for all $g\in G_F$ and $X\in\mathfrak{sp}_4(k)$ is $0$.
\end{enumerate}
\end{lemma}
\leanref{Development/AutomorphicArithmetic/AutomorphicVarieties/WeakShimuraTheory/AdjointCoefficients/SpDims.lean}{UpperTheory.V.Ws.X1.spDims}
\leanref{Development/AutomorphicArithmetic/AutomorphicVarieties/WeakShimuraTheory/AdjointCoefficients/AdFacts.lean}{UpperTheory.V.Ws.X1.adFacts}
\leanref{Development/AutomorphicArithmetic/AutomorphicVarieties/WeakShimuraTheory/AdjointCoefficients/RankAdjLe.lean}{UpperTheory.V.Ws.X1.ccDim}
\leanref{Development/AutomorphicArithmetic/AutomorphicVarieties/WeakShimuraTheory/AdjointCoefficients/DualZero.lean}{UpperTheory.V.Ws.X1.dualZero}
\begin{proof}
The map $X\mapsto X^*$ is linear, satisfies $(XY)^*=Y^*X^*$ and $X^{**}=X$, and carries the entry $(i,j)$ to the entry $(5-j,5-i)$ up to sign.

(1) Since $i\le j$ if and only if $5-j\le5-i$, the upper triangular part and the strictly lower triangular part of an element of $\mathfrak{sp}_4(k)$ lie in $\mathfrak{sp}_4(k)$. Counting the entries that $X^*=-X$ leaves free, $\mathfrak{sp}_4(k)$ has dimension $10$, its upper triangular elements $\mathfrak b(1)$ dimension $6$, its strictly lower triangular elements $\mathfrak u(1)$ dimension $4$, and its diagonal elements, $\operatorname{diag}(x,y,-y,-x)$, dimension $2$; and $\mathfrak{sp}_4(k)=\mathfrak b(1)\oplus\mathfrak u(1)$. If $n$ is symplectic with multiplier $1$, then $n^*=n^{-1}$, and $X\mapsto n^*Xn$ is an automorphism of $\mathfrak{sp}_4(k)$ that carries $\mathfrak b(n)$ and $\mathfrak u(n)$ onto $\mathfrak b(1)$ and $\mathfrak u(1)$.

(2) If $r(h)$ has multiplier $\nu$, then $r(h)^*=\nu r(h)^{-1}$ and $(r(h)^{-1})^*=\nu^{-1}r(h)$, so $(r(h)Xr(h)^{-1})^*=-r(h)Xr(h)^{-1}$ for $X\in\mathfrak{sp}_4(k)$. If $n^*=n^{-1}$, then $n^*(r(h)Xr(h)^{-1})n=(n^*r(h)n)(n^*Xn)(n^*r(h^{-1})n)$, a product of upper triangular matrices when $X\in\mathfrak b(n)$. Suppose now that $g=r(h)$ has $g^2=1$ and $g^*g=-1$, so that $g^*=-g$. Then $p=\frac12(1+g)$ and $q=\frac12(1-g)$ are idempotents with $p+q=1$, $pq=0$ and $p^*=q$. As $X\mapsto X^*$ preserves ranks and the image of $q$ lies in the kernel of $p$, both have rank $2$. Let $Y\in\mathfrak{sp}_4(k)$ commute with $g$, hence with $p$, and suppose that $pYp=0$. Then $Yp=0$; applying $X\mapsto X^*$ gives $qY=0$, so $Yq=0$ and $Y=Y(p+q)=0$. So $Y\mapsto pYp$, viewed as an endomorphism of the image of $p$, is injective on these $Y$, and their space has dimension at most $2^2=4$.

(3) Let $F'(Y)=f\big(\frac12(Y-Y^*)\big)$ for $Y\in M_4(k)$. As $(r(g)Yr(g)^{-1})^*=r(g)Y^*r(g)^{-1}$ by the computation in (2), $F'(r(g)Yr(g)^{-1})=\omega'(g)F'(Y)$, and $F'$ agrees with $f$ on $\mathfrak{sp}_4(k)$. Let $h=(-3,-1,-2,0)$ if $r$ is supported in $\{(-1,1),(-1,0),(-1,-1),(0,-1)\}$, and $h=(-3,-2,-1,0)$ otherwise. By the supports of Definition~\ref{sha-siegel-tower-pairs}(3), in both cases $r(g)_{ij}=0$ whenever $h_i<h_j$, and $|h_i-h_j|\le3$ for all $i,j$. We show that $F'(E_{ij})=0$ for the matrix units, by descending induction on $h_i-h_j$. Let $w$ be the weight of the entry $(i,j)$, and $g\in G_F$ with $a(\chi^w(g))\ne\omega'(g)$. In
\[
F'(r(g)E_{ij}r(g^{-1}))=\sum_{i',j'}r(g)_{i'i}\,r(g^{-1})_{jj'}F'(E_{i'j'}),
\]
a term with $(i',j')\ne(i,j)$ vanishes unless $h_{i'}\ge h_i$ and $h_{j'}\le h_j$, and then $h_{i'}-h_{j'}>h_i-h_j$, so $F'(E_{i'j'})=0$ by induction. As $c$ is residual, $r(g)_{ii}r(g^{-1})_{jj}=a(\chi_i(g)\chi_j(g)^{-1})=a(\chi^w(g))$ (Definition~\ref{sha-siegel-tower-towers}(2)). So $(a(\chi^w(g))-\omega'(g))F'(E_{ij})=0$, and $F'(E_{ij})=0$. Hence $F'=0$, and $f=0$.
\end{proof}
% (1) is spDims (SpDims.lean, 179 lines: tor2, bor1, compl1, and bor_eq, nil_eq through the conjugation cjE by n). (2) is adFacts (AdFacts.lean, 62 lines; it also gives the openness of the kernel of the action when the entries of r are locally constant) and ccDim (RankAdjLe.lean, 130 lines, with rank_adj_le and cmp_zero). (3) is dualZero (DualZero.lean, 164 lines): F is f composed with t(id - adj), t = 1/2; tri_zero is the induction, with the orders ![-3,-1,-2,0] and ![-3,-2,-1,0] for the two residual shapes; rt_wt is the weight identity. The Lean states (3) for any problem with Fr.Ok and Res.

\begin{lemma}\label{sha-siegel-step-completion-euler}
In Situation~\ref{sha-descent-link-setting}, suppose that $F$ is totally real and that $S$ is the set of places of $F$ above $\ell$. Let $V$ be a finite set with a bijection onto $S$, for each $v\in V$ let $A_v$ be a place of $\bar{\mathbf Q}$ above the image of $v$ in $S$, with decomposition group $G_v$ in $G_F$, and let $\mathrm{loc}_v:G_v\to\Gamma$ be the restriction of $\pi$.
Let $k$ be a finite field of characteristic $\ell$, $W$ a finite-dimensional $k$-vector space with a linear action of $\Gamma$ with open kernel, and $b_v\subseteq W$ a subspace stable under $\mathrm{loc}_v(G_v)$ for each $v$; define $H^i_b(\Gamma,W)$ by (1) to (3) of Definition~\ref{sha-descent-residual}, for these maps $\mathrm{loc}_v$.
Let $m,e\ge0$ be integers such that $\dim b_v=m$ for every $v$ and $\dim\ker(\pi(\tau)-1\mid W)\le e$ for every complex conjugation $\tau\in G_F$, and suppose that every linear form $\lambda$ on $W$ with $\lambda(\pi(g)X)=\omega(g)\lambda(X)$ for all $g\in G_F$ and $X\in W$ is $0$.
Then $H^1_b(\Gamma,W)$ and $H^2_b(\Gamma,W)$ are finite-dimensional, and
\[
\dim H^2_b(\Gamma,W)+\dim H^0_b(\Gamma,W)+[F:\mathbf Q]\,m\le\dim H^1_b(\Gamma,W)+[F:\mathbf Q]\,e.
\]
\end{lemma}
\leanref{Development/AutomorphicArithmetic/AutomorphicVarieties/WeakShimuraTheory/AdjointCoefficients/ExAbove.lean}{UpperTheory.V.Ws.X1.ecG}
\begin{proof}
For each infinite place $w$ of $F$ let $\tau_w\in G_F$ be a complex conjugation at $w$, and for $v\in V$ let $Q_v=W/b_v$, on which $G_v$ acts through $\mathrm{loc}_v$; let $d_v$ be the local degree of $F$ at the image of $v$.
\begin{enumerate}
\item[(a)] \emph{The global bound.} By the hypothesis on linear forms, Lemma~\ref{sha-descent-link-h2} shows that $H^1(\Gamma,W)$ and $H^2(\Gamma,W)$ are finite-dimensional and that
\[
\dim\text{Ш}^2_S(\Gamma,W)+\sum_v\dim H^2(G_v,W)+\dim H^0(\Gamma,W)+[F:\mathbf Q]\dim W\le\dim H^1(\Gamma,W)+\sum_w\dim\ker(\pi(\tau_w)-1\mid W).
\]
\item[(b)] \emph{Local quotients.} Let $X\subseteq H^2(\Gamma,W)$ be the subspace of classes whose image in $H^2(G_v,Q_v)$ vanishes for every $v$. Restriction maps $X$ to $\bigoplus_v\ker\big(H^2(G_v,W)\to H^2(G_v,Q_v)\big)$, with kernel $\text{Ш}^2_S(\Gamma,W)$, and the maps $H^2(G_v,W)\to H^2(G_v,Q_v)$ are surjective by Lemma~\ref{sha-descent-link-local}(3). So $\dim X+\sum_v\dim H^2(G_v,Q_v)\le\dim\text{Ш}^2_S(\Gamma,W)+\sum_v\dim H^2(G_v,W)$.
\item[(c)] \emph{The cone.} By Lemma~\ref{sha-descent-link-cone}, $H^1_b(\Gamma,W)$ and $H^2_b(\Gamma,W)$ are finite-dimensional, and
\[
\dim H^2_b+\dim H^0_b+\dim H^1(\Gamma,W)+\sum_v\dim H^0(G_v,Q_v)=\dim H^1_b+\dim H^0(\Gamma,W)+\sum_v\dim H^1(G_v,Q_v)+\dim X.
\]
\item[(d)] \emph{Local Euler characteristics.} By Lemma~\ref{sha-descent-link-local}(2), $\dim H^1(G_v,Q_v)=\dim H^0(G_v,Q_v)+\dim H^2(G_v,Q_v)+d_v(\dim W-m)$.
\item[(e)] \emph{Degrees.} The $d_v$ are the local degrees at distinct places above $\ell$, so $\sum_vd_v\le[F:\mathbf Q]$. As $F$ is totally real, it has $[F:\mathbf Q]$ infinite places, so $\sum_w\dim\ker(\pi(\tau_w)-1\mid W)\le[F:\mathbf Q]\,e$.
\end{enumerate}
Adding (a) to (d) gives $\dim H^2_b+\dim H^0_b+[F:\mathbf Q]\dim W\le\dim H^1_b+\sum_w\dim\ker(\pi(\tau_w)-1\mid W)+\sum_vd_v(\dim W-m)$, and (e) gives the claim.
\end{proof}
% ecG (ExAbove.lean, 135 lines): cc_ex for the tau_w; (a) and (b) together are GC.H2_galois_quotient_surjective (Lemma sha-descent-link-h2 with local_second_cohomology_surjective_at_place); (c) is GC.coneCount; (d) is GC.locEuler; (e) is sum_deg and the count of infinite places. This is the argument for condition (11) in Lemma sha-descent-link-cohomology, which the Lean proves through the same lemmas (GC.ec).

\begin{lemma}\label{sha-siegel-step-completion-tangent}
Let $D$ be a Siegel datum, $T$ a fit deep tower for $D$ with field $F$ (Definition~\ref{sha-siegel-tower-towers}), and $\mathcal O'\supseteq\mathcal O$ a ring of Definition~\ref{sha-character-parts}, with residue field $k'$. Let $\Phi^\circ$ be the deformation problem with the data of the problem of $T$, except that every set $\mu_v$ is empty; Definition~\ref{sha-siegel-tower-pairs}(2) to (5) apply to it. Let $c=(r_c,(n_{c,v}))$ be a residual flagged pair for $\Phi^\circ$ along $\mathcal O\to k'$, and $P'$ a universal lift of $c$ to $R'$ for $\Phi^\circ$. Let $\Gamma_T$, $N_T$ and $\pi_T$ be as in Lemma~\ref{sha-siegel-step-completion-facts}(4), and for $v\in\Sigma_T$ let $\mathrm{loc}_v:D_v(F)\to\Gamma_T$ be the restriction of $\pi_T$.
Then $r_c=\bar r\circ\pi_T$ for a homomorphism $\bar r:\Gamma_T\to M_4(k')$ with locally constant entries.
Let $\Gamma_T$ act on $W=\mathfrak{sp}_4(k')$ by $\gamma\cdot X=\bar r(\gamma)X\bar r(\gamma)^{-1}$, let $b_v=\mathfrak b(n_{c,v})$ (Lemma~\ref{sha-siegel-step-completion-adjoint}), and define $H^i_b(\Gamma_T,W)$ by (1) to (3) of Definition~\ref{sha-descent-residual}, for the maps $\mathrm{loc}_v$. Let $t_{R'}$ be the $k'$-vector space of $\mathcal O'$-linear maps $\delta:R'\to k'$ with $\delta(xy)=\bar x\delta(y)+\bar y\delta(x)$.
Then $H^1_b(\Gamma_T,W)$ is finite-dimensional, and
\[
\dim H^1_b(\Gamma_T,W)+2\le\dim t_{R'}+\dim H^0_b(\Gamma_T,W).
\]
\end{lemma}
\leanref{Development/AutomorphicArithmetic/AutomorphicVarieties/WeakShimuraTheory/AdjointCoefficients/Descend.lean}{UpperTheory.V.Ws.X1.descend}
\leanref{Development/AutomorphicArithmetic/AutomorphicVarieties/WeakShimuraTheory/AdjointCoefficients/TanInj.lean}{UpperTheory.V.Ws.X1.tanInj}
\leanref{Development/AutomorphicArithmetic/AutomorphicVarieties/WeakShimuraTheory/AdjointCoefficients/ConeTan.lean}{UpperTheory.V.Ws.X1.coneTan}
\begin{proof}
The set of $g\in G_F$ with $r_c(g)=1$ is a subgroup, as $r_c(g)r_c(g^{-1})=1$, and it is normal. It is closed, as the entries of $r_c$ are locally constant, and it contains $N_\Phi$; so it contains $N_T$. Hence $r_c=\bar r\circ\pi_T$, and $\bar r$ has locally constant entries, as $\pi_T$ is a quotient map. By Lemma~\ref{sha-siegel-step-completion-adjoint}(2), $\Gamma_T$ acts on $W$ and $b_v$ is stable under $\mathrm{loc}_v(D_v(F))$, as $n_{c,v}^*r_c(s)n_{c,v}$ is upper triangular for $s\in D_v(F)$. Write $\mathfrak u_v=\mathfrak u(n_{c,v})$, and $g_0=g_0^T$.
\begin{enumerate}
\item[(a)] \emph{Tangent pairs.} Let $Z$ be the space of pairs $(x,(Y_v))$ of a locally constant cocycle $x:\Gamma_T\to W$ and elements $Y_v\in\mathfrak u_v$ with $x(\mathrm{loc}_vs)-(\mathrm{loc}_v(s)\cdot Y_v-Y_v)\in b_v$ for all $v$ and $s\in D_v(F)$, and $Z_0\subseteq Z$ the subspace of the pairs with $x(\pi_T(g_0))$ diagonal.
\item[(b)] \emph{Lifts to the dual numbers.} Let $(x,(Y_v))\in Z_0$, and over $k'[\varepsilon]$ put $r'(g)=(1+\varepsilon x(\pi_Tg))r_c(g)$ and $n'_v=(1-\varepsilon Y_v)n_{c,v}$. Then $(r',(n'_v))$ is a lift of $c$ to $k'[\varepsilon]$ for $\Phi^\circ$. Indeed, $r'$ is a homomorphism by the cocycle condition, and $r'(g)$ is symplectic with multiplier $\theta(g)$, as $x(\pi_Tg)^*=-x(\pi_Tg)$. For $g\in N_\Phi$, $\pi_T(g)=1$ and $x(1)=0$, so $r'(g)=1$. The matrix $r'(g_0)$ is diagonal, as $x(\pi_Tg_0)$ and $r_c(g_0)$ are. The matrix $n'_v$ is symplectic with multiplier $1$, as $Y_v^*=-Y_v$, and lower triangular with diagonal entries $1$, as $n_{c,v}^{-1}Y_vn_{c,v}=n_{c,v}^*Y_vn_{c,v}$ is strictly lower triangular. For $s\in D_v(F)$, write $r=r_c(s)$, $n=n_{c,v}$, $Y=Y_v$ and $X=x(\mathrm{loc}_vs)$; then $n_v'^*=n^*(1+\varepsilon Y)$ and
\[
n_v'^*r'(s)n'_v=n^*rn+\varepsilon\,n^*\big(X-(rYr^{-1}-Y)\big)n\cdot n^*rn,
\]
which is upper triangular, as $n^*rn$ is and $X-(rYr^{-1}-Y)\in b_v$. The entries are locally constant, the image in $k'$ is $c$, and there is no condition at the empty sets $\mu_v$.
\item[(c)] \emph{Injectivity.} By the universal property, each element of $Z_0$ gives a homomorphism of $\mathcal O'$-algebras $R'\to k'[\varepsilon]$ that carries $P'$ to the lift of (b). Distinct elements give distinct lifts, as $\pi_T$ is surjective and $x$ and the $Y_v$ are read off from $r'$ and the $n'_v$. The homomorphisms of $\mathcal O'$-algebras $R'\to k'[\varepsilon]$ are the maps $y\mapsto\bar y+\varepsilon\delta(y)$ with $\delta\in t_{R'}$, and $t_{R'}$ is finite-dimensional, as $\mathfrak m_{R'}$ is finitely generated. As $k'$ is finite, $Z_0$ is finite-dimensional and $\dim Z_0\le\dim t_{R'}$.
\item[(d)] \emph{The count.} By Lemma~\ref{sha-siegel-step-completion-adjoint}(1), $W=b_v\oplus\mathfrak u_v$; let $p_v:W\to\mathfrak u_v$ be the projection along $b_v$. If $(x,(X_v))$ represents a class of $H^1_b(\Gamma_T,W)$, then $(x,(p_vX_v))$ lies in $Z$, as $X_v-p_vX_v\in b_v$ and $b_v$ is stable under $\mathrm{loc}_v(D_v(F))$, and it represents the same class, as $(0,(p_vX_v-X_v))$ is a coboundary. So $Z$ maps onto $H^1_b(\Gamma_T,W)$. The map $X\mapsto(\gamma\mapsto\gamma\cdot X-X,(p_vX))$ from $W$ to $Z$ lands in the kernel of this map, and its own kernel lies in $H^0_b(\Gamma_T,W)$: if $\gamma\cdot X=X$ for all $\gamma$ and $p_vX=0$, then $X\in b_v$ for all $v$. Finally $Z_0$ is the kernel of the map $Z\to W/\mathfrak t$, $(x,(Y_v))\mapsto x(\pi_Tg_0)\bmod\mathfrak t$, so $\dim Z\le\dim Z_0+10-2$. Hence $\dim H^1_b\le\dim Z-(10-\dim H^0_b)\le\dim Z_0-2+\dim H^0_b$, and (c) gives the claim.
\end{enumerate}
\end{proof}
% descend (Descend.lean, 47 lines) for rbar; tanInj (TanInj.lean, 256 lines): the pair P' of (b) is dRep and dm, checked by dm_gsp, dm_low, dm_flag, and injectivity by dm_inj; TanC (AdjointMemSp4.lean) is Z_0, with t = tor and n v = nil (c.n v). (d) is coneTan (ConeTan.lean, 135 lines): exProj for p_v, and dimLem for the count. In absPres the Lean counts t_{R'} as the number of maps to DualNumber, which obs_pres gives as |k'|^a.

\begin{lemma}\label{sha-siegel-step-completion-mu}
Let $D$, $T$, $\mathcal O'$, $k'$ and $\Phi^\circ$ be as in Lemma~\ref{sha-siegel-step-completion-tangent}, let $\Phi$ be the problem of $T$, and let the $\zeta_v$ be as in Lemma~\ref{sha-siegel-step-completion-facts}(5). Let $c$ be a residual flagged pair for $\Phi$ along $\mathcal O\to k'$, $P$ a universal lift of $c$ to $R$ for $\Phi$, and $P'$ a universal lift of $c$ to $R'$ for $\Phi^\circ$, whose diagonal characters (Definition~\ref{sha-siegel-tower-pairs}(2)) are the $\psi'_{v,i}$.
Then the homomorphism of $\mathcal O'$-algebras $f:R'\to R$ that carries $P'$ to $P$ is surjective, and its kernel is generated by the $2|\Sigma_T|$ elements $\psi'_{v,1}(\zeta_v)-1$ and $\psi'_{v,2}(\zeta_v)-1$ with $v\in\Sigma_T$.
\end{lemma}
\leanref{Development/AutomorphicArithmetic/AutomorphicVarieties/WeakShimuraTheory/AdjointCoefficients/ClosureComm.lean}{UpperTheory.V.Ws.X1.muIff}
\leanref{Development/AutomorphicArithmetic/AutomorphicVarieties/WeakShimuraTheory/AdjointCoefficients/MuQuot.lean}{UpperTheory.V.Ws.X1.muQuot}
\begin{proof}
The map $f$ exists, as $P$ is a lift of $c$ for $\Phi^\circ$.
\begin{enumerate}
\item[(a)] \emph{The condition at $\mu_v$.} Let $Q=(r,(n_v))$ be a lift of $c$ for $\Phi^\circ$ to a complete local noetherian $\mathcal O'$-algebra $R''$ with residue field $k'$, with diagonal characters $\psi_{v,i}$. Then $Q$ is a lift for $\Phi$ if and only if $\psi_{v,1}(\zeta_v)=\psi_{v,2}(\zeta_v)=1$ for every $v$. The condition is necessary, as $\zeta_v\in\mu_v$. Conversely, suppose that it holds, and let $v\in\Sigma_T$. As $n_v^*=n_v^{-1}$ and the matrices $n_v^*r(s)n_v$ with $s\in D_v(F)$ are upper triangular, each $\psi_{v,i}$ is a homomorphism $D_v(F)\to R''^\times$, continuous for the $\mathfrak m_{R''}$-adic topology. The diagonal entries of an upper triangular matrix that is symplectic with multiplier $\nu$ satisfy $\psi_i\psi_{5-i}=\nu$; at $\zeta_v$ the multiplier is $1$ by Lemma~\ref{sha-siegel-step-completion-facts}(1), so all four $\psi_{v,i}(\zeta_v)$ are $1$. As $R''$ is separated and $R''^\times$ is commutative, each $\psi_{v,i}$ is trivial on the closure of the commutator subgroup of $D_v(F)$. By Lemma~\ref{sha-siegel-step-completion-facts}(5), $\psi_{v,i}(s)=\psi_{v,i}(\zeta_v)^n=1$ for every $s\in\mu_v$.
\item[(b)] \emph{The quotient.} Let $I\subseteq R'$ be the ideal generated by the $2|\Sigma_T|$ elements. By (a), applied to $P$, $f$ kills $I$; so $I\ne R'$, and $R'/I$ is a complete local noetherian $\mathcal O'$-algebra with residue field $k'$. By (a), the image of $P'$ in $R'/I$ is a lift of $c$ for $\Phi$; let $h:R\to R'/I$ carry $P$ to it. Then $h\circ f$ carries $P'$ to the image of $P'$ in $R'/I$, so it is the quotient map, by uniqueness for $P'$. The map $\bar f:R'/I\to R$ induced by $f$ has $\bar f\circ h=\mathrm{id}_R$, by uniqueness for $P$. So $f$ is surjective, and if $f(y)=0$, then the class of $y$ is $h(f(y))=0$, so $y\in I$.
\end{enumerate}
\end{proof}
% (a) is muIff (ClosureComm.lean, 121 lines, with gsp_diag, up_mul_diag and closure_comm), with hnu = Lemma sha-siegel-step-completion-facts (1) at zeta_v; (b) is muQuot (MuQuot.lean, 97 lines) with co_quot.

\begin{lemma}\label{sha-siegel-step-completion-obstruction-lift}
Let $\mathcal O'$ be a complete local noetherian ring in which $2$ is a unit, with residue field $k'$, and let $p:A\to R$ be a surjective homomorphism of complete local noetherian $\mathcal O'$-algebras with residue field $k'$ whose kernel $K$ satisfies $\mathfrak m_AK=0$. Let $S$ be a topological space, $\nu:S\to\mathcal O'^\times$ a map that is continuous for the $\mathfrak m_{\mathcal O'}$-adic topology, and $x:S\to M_4(R)$ a map with $\mathfrak m_R$-adically continuous entries such that $x(s)$ is symplectic with multiplier $\nu(s)$ for every $s$ (Definition~\ref{sha-siegel-datum}). Then there is a map $y:S\to M_4(A)$ with $\mathfrak m_A$-adically continuous entries such that, for every $s$, $y(s)$ is symplectic with multiplier $\nu(s)$ and maps to $x(s)$; and if every $x(s)$ is upper triangular, so is every $y(s)$.
\end{lemma}
\leanref{Development/AutomorphicArithmetic/AutomorphicVarieties/WeakShimuraTheory/AdjointCoefficients/GspLift.lean}{UpperTheory.V.Ws.X1.gspLift}
\leanref{Development/AutomorphicArithmetic/AutomorphicVarieties/WeakShimuraTheory/AdjointCoefficients/ContLift.lean}{UpperTheory.V.Ws.X1.contLift}
\begin{proof}
As $p$ is surjective onto a nonzero ring, $K\subseteq\mathfrak m_A$, so $yy'=0$ for $y,y'\in K$, and a product of two matrices over $K$ is $0$.
\begin{enumerate}
\item[(a)] \emph{Continuous lifts.} The construction of step (1) of the proof of Lemma~\ref{sha-automorphy-unramified-near-completion-obstruction}, for $p:A\to R$, gives $N_0$ with $K\cap\mathfrak m_A^{N_0}=0$ and a map $\sigma:R\to A$ with $p\circ\sigma=\mathrm{id}$ such that $\sigma(c_1)-\sigma(c_2)\in\mathfrak m_A^N$ whenever $N\ge N_0$ and $c_1-c_2\in\mathfrak m_R^N$. Let $y_0(s)$ be the matrix whose entry $(i,j)$ is $0$ if $x(t)_{ij}=0$ for every $t\in S$, and $\sigma(x(s)_{ij})$ otherwise. Then $y_0(s)$ maps to $x(s)$, and its entries are continuous: modulo $\mathfrak m_A^N$, with $N\ge N_0$, $y_0(s)_{ij}$ depends only on $x(s)_{ij}$ modulo $\mathfrak m_R^N$, which is locally constant in $s$.
\item[(b)] \emph{The multiplier.} Let $E(s)=y_0(s)^*y_0(s)-\nu(s)$. Its entries lie in $K$, as $x(s)^*x(s)=\nu(s)$; so $E(s)^2=0$, and $E(s)^*=E(s)$. Let $y(s)=y_0(s)\big(1-(2\nu(s))^{-1}E(s)\big)$. Then $y(s)^*y(s)=\big(1-(2\nu(s))^{-1}E(s)\big)\big(\nu(s)+E(s)\big)\big(1-(2\nu(s))^{-1}E(s)\big)=\nu(s)+E(s)-E(s)=\nu(s)$, and $y(s)$ maps to $x(s)$, as $E(s)$ maps to $0$. Modulo $\mathfrak m_A^N$, $y(s)$ is a polynomial in the entries of $y_0(s)$ and in $\nu(s)^{-1}$, so its entries are continuous.
\item[(c)] \emph{Triangularity.} If every $x(s)$ is upper triangular, the entries of $x$ below the diagonal vanish identically, so $y_0(s)$ is upper triangular. As $X^*_{ij}=\pm X_{5-j,5-i}$, $y_0(s)^*$ is upper triangular, and so are $E(s)$ and $y(s)$.
\end{enumerate}
\end{proof}
% gspLift (GspLift.lean, 120 lines): (a) is contLift (ContLift.lean, 85 lines: inf_pow_bot, the Artin-Rees lemma, for N_0, and the section sigma) with the entries that vanish identically lifted by 0; (b) is corr with c = u (nu s)^-1, u an inverse of 2, through corr_gsp and corr_map; (c) is corr_up.

\begin{lemma}\label{sha-siegel-step-completion-obstruction-diagonal}
Let $\mathcal O'$, $p:A\to R$ and $K$ be as in Lemma~\ref{sha-siegel-step-completion-obstruction-lift}. Let $X\in M_4(A)$ be symplectic with a multiplier $\nu\in A^\times$, such that the image of $X$ in $M_4(R)$ is diagonal and $X_{ii}-X_{jj}$ is a unit for $i\ne j$. Then there is $u\in M_4(A)$, symplectic with multiplier $1$, that maps to $1$ in $M_4(R)$, such that $u^*Xu$ is diagonal.
\end{lemma}
\leanref{Development/AutomorphicArithmetic/AutomorphicVarieties/WeakShimuraTheory/AdjointCoefficients/DiagLift.lean}{UpperTheory.V.Ws.X1.diagLift}
\begin{proof}
A product of two matrices over $K$ is $0$, as in the proof of Lemma~\ref{sha-siegel-step-completion-obstruction-lift}. Let $D=\operatorname{diag}(X_{11},\dots,X_{44})$ and $E=X-D$, a matrix over $K$ with zero diagonal. As $E^*E=0$ and $D^*=\operatorname{diag}(X_{44},\dots,X_{11})$, the identity $X^*X=\nu$ gives $D^*D+D^*E+E^*D=\nu$; so $X_{5-i,5-i}X_{ii}=\nu$ for every $i$, and every $X_{ii}$ is a unit, and $X_{5-i,5-i}E_{ij}+(E^*)_{ij}X_{jj}=0$ for $i\ne j$. Let $N$ be the matrix with $N_{ii}=0$ and $N_{ij}=-E_{ij}(X_{ii}-X_{jj})^{-1}$ for $i\ne j$, a matrix over $K$. From these identities and $X_{5-j,5-j}-X_{5-i,5-i}=\nu(X_{ii}-X_{jj})X_{ii}^{-1}X_{jj}^{-1}$, a computation gives $N^*=-N$. Let $u=1+N$. Then $u^*=1-N$, $u^*u=1-N^2=1$, $u$ maps to $1$, and $u^*Xu=(1-N)(D+E)(1+N)=D+E+DN-ND$, whose entry $(i,j)$ with $i\ne j$ is $E_{ij}+(X_{ii}-X_{jj})N_{ij}=0$.
\end{proof}
% diagLift (DiagLift.lean, 87 lines): the private exists_symplectic_diagonalizing_lift, with key, k1, k2 for the two identities, aN for N^* = -N, and the final computation e.

\begin{lemma}\label{sha-siegel-step-completion-obstruction-flags}
Let $\mathcal O'$, $p:A\to R$ and $K$ be as in Lemma~\ref{sha-siegel-step-completion-obstruction-lift}, and let $n\in M_4(A)$ be symplectic with multiplier $1$, with image in $M_4(R)$ lower triangular with diagonal entries $1$. Then there are $m,b\in M_4(A)$, symplectic with multiplier $1$, such that $m$ is lower triangular with diagonal entries $1$, $b$ and $b^*$ are upper triangular, $m$ and $n$ have the same image in $M_4(R)$, and $m=nb$.
\end{lemma}
\leanref{Development/AutomorphicArithmetic/AutomorphicVarieties/WeakShimuraTheory/AdjointCoefficients/LowLift.lean}{UpperTheory.V.Ws.X1.lowLift}
\begin{proof}
A product of two matrices over $K$ is $0$, as in the proof of Lemma~\ref{sha-siegel-step-completion-obstruction-lift}.
\begin{enumerate}
\item[(a)] \emph{A unitriangular approximation.} For $a,b',c',d\in A$ let $\Lambda(a,b',c',d)$ be the lower triangular matrix with diagonal entries $1$ and entries $\Lambda_{21}=a$, $\Lambda_{31}=b'$, $\Lambda_{32}=c'$, $\Lambda_{41}=d$, $\Lambda_{42}=b'-ac'$, $\Lambda_{43}=-a$. A computation shows that it is symplectic with multiplier $1$, and that every lower triangular matrix with diagonal entries $1$ that is symplectic with multiplier $1$ is $\Lambda$ of its entries $(2,1)$, $(3,1)$, $(3,2)$, $(4,1)$. Let $m_0=\Lambda(n_{21},n_{31},n_{32},n_{41})$; it has the same image as $n$, as that image is of this form.
\item[(b)] \emph{The correction.} Let $Y=m_0^*n-1$, a matrix over $K$, as $m_0^*=m_0^{-1}$ and $m_0$ and $n$ have the same image. As $m_0^*n$ is symplectic with multiplier $1$ and $Y^*Y=0$, $Y^*+Y=0$. Let $Y_u$ be the matrix of the entries $(i,j)$ of $Y$ with $i\le j$, and $0$ elsewhere; as $X\mapsto X^*$ maps such entries to such entries, $Y_u^*+Y_u=0$. Let $b=1-Y_u$. It is upper triangular, $b^*=1+Y_u$ is upper triangular, $b^*b=1$, and $b$ maps to $1$.
\item[(c)] \emph{The matrix $m$.} Now $m_0^*nb=(1+Y)(1-Y_u)=1+(Y-Y_u)$, which is lower triangular with diagonal entries $1$, as $Y-Y_u$ is strictly lower triangular, and symplectic with multiplier $1$. Let $m=m_0(m_0^*nb)=nb$. It is a product of two lower triangular matrices with diagonal entries $1$ that are symplectic with multiplier $1$, so it is one too, and it has the same image as $n$, as $b$ maps to $1$.
\end{enumerate}
\end{proof}
% lowLift (LowLift.lean, 155 lines): (a) is lowOf, lowOf_low and eq_lowOf (the Lean's lowOf a b c d, with the indices from 0); (b) and (c) are the private aux, with low_mul and up_adj.

\begin{definition}\label{sha-siegel-step-completion-obstruction-generators}
Let $\mathcal O'$ be a complete local noetherian ring with residue field $k'$, $R$ a complete local noetherian $\mathcal O'$-algebra with residue field $k'$, and $H$ a $k'$-vector space. Let $C'$ be a complete local noetherian $\mathcal O'$-algebra with residue field $k'$ and $p:C'\to R$ a surjective $\mathcal O'$-algebra map with $\mathfrak m_{C'}\ker p=0$. A family $\varepsilon_1,\dots,\varepsilon_n$ of generators of the ideal $\ker p$ is \emph{minimal} if $c_1,\dots,c_n\in\mathfrak m_{\mathcal O'}$ whenever $c_1,\dots,c_n\in\mathcal O'$ and $\sum_jc_j\varepsilon_j=0$. Then $H$ is an \emph{obstruction space for minimal generators} of $R$ if for every such $C'$ and $p$ and every minimal family $\varepsilon_1,\dots,\varepsilon_n$ of generators of $\ker p$ there are $o_1,\dots,o_n\in H$ with the following property. Let $c_1,\dots,c_n\in\mathcal O'$ with $\sum_j\bar c_jo_j=0$, let $q:C'\to C''$ be a surjective $\mathcal O'$-algebra map onto a complete local noetherian $\mathcal O'$-algebra with residue field $k'$, let $p'':C''\to R$ be an $\mathcal O'$-algebra map with $p''\circ q=p$, and let $\eta\in C''$ with $\mathfrak m_{C''}\eta=0$, $p''(\eta)=0$ and $q(\varepsilon_j)=c_j\eta$ for every $j$. Then there is an $\mathcal O'$-algebra map $s:R\to C''$ with $p''\circ s=\mathrm{id}_R$.
\end{definition}
\leanref{Development/AutomorphicArithmetic/AutomorphicVarieties/WeakShimuraTheory/AdjointMemSp4.lean}{UpperTheory.V.Ws.Obs}

\begin{lemma}\label{sha-siegel-step-completion-obstruction}
Let $D$, $T$, $\mathcal O'$, $k'$, $\Phi^\circ$, $c$, $P'$, $R'$, $\Gamma_T$, $W$ and the $b_v$ be as in Lemma~\ref{sha-siegel-step-completion-tangent}. Then $H^2_b(\Gamma_T,W)$ is an obstruction space for minimal generators of $R'$ (Definition~\ref{sha-siegel-step-completion-obstruction-generators}).
\end{lemma}
\leanref{Development/AutomorphicArithmetic/AutomorphicVarieties/WeakShimuraTheory/AdjointCoefficients/ObWeak.lean}{UpperTheory.V.Ws.X1.obWeak}
\leanref{Development/AutomorphicArithmetic/AutomorphicVarieties/WeakShimuraTheory/AdjointCoefficients/ContConj.lean}{UpperTheory.V.Ws.X1.sect}
\leanref{Development/AutomorphicArithmetic/AutomorphicVarieties/WeakShimuraTheory/AdjointCoefficients/ContConjAlt.lean}{UpperTheory.V.Ws.X1.slEx}
\leanref{Development/AutomorphicArithmetic/AutomorphicVarieties/WeakShimuraTheory/AdjointCoefficients/MapAdjAlg.lean}{UpperTheory.V.Ws.X1.obCoc}
\leanref{Development/AutomorphicArithmetic/AutomorphicVarieties/WeakShimuraTheory/AdjointCoefficients/ExistsStrictHenselianAbsolutelyClosedLocalRing.lean}{UpperTheory.V.Ws.X1.obCor}
\begin{proof}
Write $P'=(r_{P'},(n_{P',v}))$, $c=(r_c,(n_{c,v}))$ and $r_c=\bar r\circ\pi_T$ as in Lemma~\ref{sha-siegel-step-completion-tangent}, $\mathrm{loc}_v:D_v(F)\to\Gamma_T$ for the restriction of $\pi_T$, and $\theta$ also for its image in $\mathcal O'^\times$. As $\ell$ is odd, $2$ is a unit of $\mathcal O'$. Let $C'$ be a complete local noetherian $\mathcal O'$-algebra with residue field $k'$ and $p:C'\to R'$ a surjective homomorphism of $\mathcal O'$-algebras with $\mathfrak m_{C'}K=0$, where $K=\ker p$; then $K\subseteq\mathfrak m_{C'}$, and $K$ is a $k'$-vector space.
\begin{enumerate}
\item[(0)] \emph{Descent to $\Gamma_T$.} By Lemma~\ref{sha-siegel-step-completion-facts}(1), $\theta(g)=1$ for $g\in N_\Phi$, and $r_{P'}(g)=1$ for $g\in N_\Phi$. The $g\in G_F$ with $r_{P'}(g)=1$ form a normal subgroup, which is closed, as the entries of $r_{P'}$ are continuous and $R'$ is separated; so it contains $N_T$, and in the same way $\ker\theta\supseteq N_T$. Hence $r_{P'}=r\circ\pi_T$ and $\theta=\theta_\Gamma\circ\pi_T$ for a homomorphism $r:\Gamma_T\to M_4(R')$ with continuous entries and a continuous character $\theta_\Gamma$, as $\pi_T$ is an open quotient map; $r$ reduces to $\bar r$.
\item[(1)] \emph{Lifts.} Lemma~\ref{sha-siegel-step-completion-obstruction-lift} gives lifts with $\mathfrak m_{C'}$-adically continuous entries: $\tilde r:\Gamma_T\to M_4(C')$ of $r$, with $\tilde r(\gamma)$ symplectic with multiplier $\theta_\Gamma(\gamma)$; $\tilde n_v$ of $n_{P',v}$, symplectic with multiplier $1$; and $\tilde t_v:D_v(F)\to M_4(C')$ of $s\mapsto n_{P',v}^*r(\mathrm{loc}_vs)n_{P',v}$, with $\tilde t_v(s)$ upper triangular and symplectic with multiplier $\theta(s)$, as $n_{P',v}^*=n_{P',v}^{-1}$ and $P'$ is a flagged pair.
\item[(2)] \emph{The defects.} Let $\mathrm{om}(\gamma,\delta)=\tilde r(\gamma)\tilde r(\delta)-\tilde r(\gamma\delta)$ and $\mathrm{de}_v(s)=\tilde r(\mathrm{loc}_vs)-\tilde n_v\tilde t_v(s)\tilde n_v^*$. Their entries lie in $K$, as $r$ is a homomorphism and $n_{P',v}n_{P',v}^*=1$. They are locally constant: as in step (2) of the proof of Lemma~\ref{sha-automorphy-unramified-near-completion-obstruction}, there is $N_0$ with $K\cap\mathfrak m_{C'}^{N_0}=0$, they are locally constant modulo $\mathfrak m_{C'}^{N_0}$, and two of their values that agree modulo $\mathfrak m_{C'}^{N_0}$ differ by a matrix over $K\cap\mathfrak m_{C'}^{N_0}=0$.
\item[(3)] \emph{Matrices over $K$.} Let $u:K\to k'$ be $\mathcal O'$-linear, applied to matrices entrywise; then $u(az)=\bar au(z)$ for $a\in C'$ and $z\in K$. For matrices $Z$ over $K$ and $M$ over $C'$: $u(MZ)=\bar Mu(Z)$, $u(ZM)=u(Z)\bar M$ and $u(Z^*)=u(Z)^*$. Let $Z_1,Z_2\in M_4(C')$ be symplectic with the same multiplier, with the same image in $M_4(R')$ and with invertible residue $g$. Then $Z_1-Z_2$ is over $K$, and $u(Z_1-Z_2)g^{-1}\in\mathfrak{sp}_4(k')$. Indeed, $(Z_1-Z_2)^*Z_2+Z_1^*(Z_1-Z_2)=Z_1^*Z_1-Z_2^*Z_2=0$; applying $u$ gives $u(Z_1-Z_2)^*g+g^*u(Z_1-Z_2)=0$, and multiplying by $(g^{-1})^*$ on the left and by $g^{-1}$ on the right gives the claim.
\item[(4)] \emph{The cocycle.} Let $\omega_u(\gamma,\delta)=u(\mathrm{om}(\gamma,\delta))\bar r(\gamma\delta)^{-1}$ and $\alpha_{u,v}(s)=u(\mathrm{de}_v(s))\bar r(\mathrm{loc}_vs)^{-1}$. They lie in $W$ by (3): $\tilde r(\gamma)\tilde r(\delta)$ and $\tilde r(\gamma\delta)$ are symplectic with multiplier $\theta_\Gamma(\gamma\delta)$ and map to $r(\gamma\delta)$, and $\tilde r(\mathrm{loc}_vs)$ and $\tilde n_v\tilde t_v(s)\tilde n_v^*$ are symplectic with multiplier $\theta(s)$ and map to $r(\mathrm{loc}_vs)$. They are locally constant by (2), as $\bar r$ is. Expanding gives $\tilde r(\gamma)\mathrm{om}(\delta,\epsilon)+\mathrm{om}(\gamma,\delta\epsilon)=\mathrm{om}(\gamma\delta,\epsilon)+\mathrm{om}(\gamma,\delta)\tilde r(\epsilon)$; applying $u$ with (3) and multiplying on the right by $\bar r(\gamma\delta\epsilon)^{-1}$ gives $\gamma\cdot\omega_u(\delta,\epsilon)+\omega_u(\gamma,\delta\epsilon)=\omega_u(\gamma\delta,\epsilon)+\omega_u(\gamma,\delta)$, so $\omega_u$ is a $2$-cocycle.
\item[(5)] \emph{The local condition.} Let $v\in\Sigma_T$ and $s,s'\in D_v(F)$, write $\gamma=\mathrm{loc}_vs$, $\gamma'=\mathrm{loc}_vs'$ and $\tilde n=\tilde n_v$, and let $T=\tilde t_v(s)\tilde t_v(s')-\tilde t_v(ss')$, an upper triangular matrix over $K$. Expanding, with $\tilde n^*\tilde n=1$,
\[
\mathrm{om}(\gamma,\gamma')+\mathrm{de}_v(ss')=\mathrm{de}_v(s)\tilde r(\gamma')+\tilde n\tilde t_v(s)\tilde n^*\mathrm{de}_v(s')+\tilde nT\tilde n^*.
\]
Apply $u$, using (3) and $\bar n\bar t_v(s)\bar n^*=\bar r(\gamma)$ for the residues $\bar n=n_{c,v}$ and $\bar t_v(s)$, and multiply on the right by $\bar r(\gamma\gamma')^{-1}$: then $\omega_u(\gamma,\gamma')+\alpha_{u,v}(ss')=\alpha_{u,v}(s)+\gamma\cdot\alpha_{u,v}(s')+\bar nu(T)\bar n^*\bar r(\gamma\gamma')^{-1}$. So $\omega_u(\gamma,\gamma')$ minus the coboundary of $\alpha_{u,v}$ at $(s,s')$ is $\bar nu(T)\bar n^*\bar r(\gamma\gamma')^{-1}$, and $\bar n^*$ times it times $\bar n$ is $u(T)\cdot\bar n^*\bar r((\gamma\gamma')^{-1})\bar n$, which is upper triangular, as $c$ is a flagged pair. So it lies in $b_v$. With (4), $(\omega_u,(\alpha_{u,v}))$ represents a class $\mathrm{ob}(u)\in H^2_b(\Gamma_T,W)$, and $u\mapsto\mathrm{ob}(u)$ is $k'$-linear.
\item[(6)] \emph{The classes and the vanishing.} Let $\varepsilon_1,\dots,\varepsilon_n$ be a minimal family of generators of $K$ (Definition~\ref{sha-siegel-step-completion-obstruction-generators}). Every $y\in K$ is $\sum_ja_j\varepsilon_j$ with $a_j\in\mathcal O'$, as an element of $C'$ acts on $K$ as any element of $\mathcal O'$ with the same residue; and the residues $\bar a_j$ depend only on $y$, as the family is minimal. Let $\mathrm{co}_j(y)=\bar a_j$; the maps $\mathrm{co}_j:K\to k'$ are $\mathcal O'$-linear. Let $o_j=\mathrm{ob}(\mathrm{co}_j)$. Now let $c_1,\dots,c_n\in\mathcal O'$ with $\sum_j\bar c_jo_j=0$, and let $q:C'\to C''$, $p'':C''\to R'$ and $\eta\in C''$ be as in Definition~\ref{sha-siegel-step-completion-obstruction-generators}. Let $u=\sum_j\bar c_j\mathrm{co}_j$. By (5), $\mathrm{ob}(u)=\sum_j\bar c_jo_j=0$: there are a continuous $h:\Gamma_T\to W$ and $Y_v\in W$ with $\omega_u(\gamma,\delta)=h(\gamma)+\gamma\cdot h(\delta)-h(\gamma\delta)$ and $B_v(s)=\alpha_{u,v}(s)-\big(h(\mathrm{loc}_vs)-(\mathrm{loc}_v(s)\cdot Y_v-Y_v)\big)\in b_v$. Let $\mu:k'\to C''$ be the map with $\mu(\bar a)=a\eta$ for $a\in\mathcal O'$. It is well defined and additive, as $\mathfrak m_{\mathcal O'}$ maps into $\mathfrak m_{C''}$, which kills $\eta$; and $\mu(\bar ac)=a\mu(c)$ for $a\in C''$, as $a$ is congruent modulo $\mathfrak m_{C''}$ to the image of an element of $\mathcal O'$. For $y=\sum_ja_j\varepsilon_j\in K$, $q(y)=\sum_ja_jc_j\eta=\mu(u(y))$. As $p''(\eta)=0$, $\eta\in\mathfrak m_{C''}$, so $\mu(k')\subseteq\mathfrak m_{C''}$, and $p''\circ\mu=0$. Apply $\mu$ to matrices entrywise. For matrices $Y,Y'$ over $k'$ and $M$ over $C''$: $\mu(Y)M=\mu(Y\bar M)$, $M\mu(Y)=\mu(\bar MY)$, $\mu(Y)\mu(Y')=0$, $\mu(Y)^*=\mu(Y^*)$, $\mu(Y)$ is upper triangular if $Y$ is, and $q(Z)=\mu(u(Z))$ for $Z$ over $K$. So $q(\mathrm{om}(\gamma,\delta))=\mu(\omega_u(\gamma,\delta)\bar r(\gamma\delta))$ and $q(\mathrm{de}_v(s))=\mu(\alpha_{u,v}(s)\bar r(\mathrm{loc}_vs))$. If $g\in M_4(C'')$ is symplectic with multiplier $\lambda$ and $Y\in\mathfrak{sp}_4(k')$, then $g+\mu(Y\bar g)$ is symplectic with multiplier $\lambda$, as $(g+\mu(Y\bar g))^*(g+\mu(Y\bar g))=g^*g+\mu(\bar g^*(Y^*+Y)\bar g)$.
\item[(7)] \emph{Corrected lifts.} Let $f(\gamma)=q(\tilde r(\gamma))-\mu(h(\gamma)\bar r(\gamma))$ and $n''_v=q(\tilde n_v)+\mu(Y_vn_{c,v})$. By (6), $f(\gamma)$ is symplectic with multiplier $\theta_\Gamma(\gamma)$, and $n''_v$ with multiplier $1$. By (6) and the formula for $\omega_u$,
\[
f(\gamma)f(\delta)=q(\tilde r(\gamma\delta))+\mu\big((\omega_u(\gamma,\delta)-\gamma\cdot h(\delta)-h(\gamma))\bar r(\gamma\delta)\big)=f(\gamma\delta),
\]
so $f$ is a homomorphism, as $f(1)$ is invertible. Its entries are continuous, $p''\circ f=r$, and $p''(n''_v)=n_{P',v}$. For $s\in D_v(F)$, with $\gamma=\mathrm{loc}_vs$, $\bar n=n_{c,v}$ and $t=q(\tilde t_v(s))$: $q(\tilde r(\gamma))=q(\tilde n_v)tq(\tilde n_v)^*+\mu(\alpha_{u,v}(s)\bar r(\gamma))$, and a computation with (6) gives
\[
n''^*_vf(\gamma)n''_v=t+\mu\big(\bar n^*(\alpha_{u,v}(s)-h(\gamma)+\gamma\cdot Y_v-Y_v)\bar r(\gamma)\bar n\big)=t+\mu\big((\bar n^*B_v(s)\bar n)(\bar n^*\bar r(\gamma)\bar n)\big),
\]
which is upper triangular, as $t$, $\bar n^*B_v(s)\bar n$ and $\bar n^*\bar r(\gamma)\bar n$ are.
\item[(8)] \emph{Normalization.} As $q$ is surjective, $\ker p''=q(K)$ and $\mathfrak m_{C''}\ker p''=q(\mathfrak m_{C'}K)=0$, so Lemmas~\ref{sha-siegel-step-completion-obstruction-diagonal} and~\ref{sha-siegel-step-completion-obstruction-flags} apply to $p''$. The image $r_{P'}(g_0^T)$ of $f(\pi_Tg_0^T)$ is diagonal, and the diagonal entries of $f(\pi_Tg_0^T)$ have the pairwise distinct residues $a(\chi_i(g_0^T))$, by Lemma~\ref{sha-siegel-step-coefficients-standard}(1), so their differences are units. Lemma~\ref{sha-siegel-step-completion-obstruction-diagonal} gives $u_0\in M_4(C'')$, symplectic with multiplier $1$ and mapping to $1$, with $u_0^*f(\pi_Tg_0^T)u_0$ diagonal. For each $v$, $u_0^*n''_v$ is symplectic with multiplier $1$ and maps to $n_{P',v}$, which is lower triangular with diagonal entries $1$; Lemma~\ref{sha-siegel-step-completion-obstruction-flags} gives $m_v=u_0^*n''_vb_v$, lower triangular with diagonal entries $1$ and mapping to $n_{P',v}$, with $b_v$ and $b_v^*$ upper triangular. Let $r''(g)=u_0^*f(\pi_Tg)u_0$ and $P''=(r'',(m_v))$. Then $P''$ is a lift of $c$ to $C''$ for $\Phi^\circ$: $r''$ is a homomorphism with continuous entries and $r''(g)$ is symplectic with multiplier $\theta(g)$; $r''(g)=1$ for $g\in N_\Phi$, as $\pi_T(g)=1$; $r''(g_0^T)$ is diagonal; as $u_0u_0^*=1$, $m_v^*r''(s)m_v=b_v^*\big(n''^*_vf(\pi_Ts)n''_v\big)b_v$ is upper triangular for $s\in D_v(F)$, by (7); the sets $\mu_v$ of $\Phi^\circ$ are empty; and $P''$ maps to $P'$ under $p''$, so its image in $k'$ is $c$.
\item[(9)] \emph{The section.} By the universal property of $P'$, some homomorphism of $\mathcal O'$-algebras $s:R'\to C''$ carries $P'$ to $P''$. Then $p''\circ s$ carries $P'$ to the image of $P''$ under $p''$, which is $P'$; by the uniqueness in the universal property, $p''\circ s=\mathrm{id}_{R'}$.
\end{enumerate}
\end{proof}
% obWeak (ObWeak.lean, 114 lines) and sect (ContConj.lean, 172 lines); absPres passes them to obs_pres. (0) is X1.descend in absPres (with hν from Lemma sha-siegel-step-completion-facts (1)); (1) is slEx (ContConjAlt.lean: gspLift three times, the structure SL with rt, nt, t); (2) is SL.om, SL.de with lc_ker (MapAdjAlg.lean); (3) is co_left, co_right, co_add and key_sp; (4) and (5) are obCoc, which puts (omega_u, alpha_u) in Fr.continuousCocycles2 with the cone condition bor; (6) is obWeak itself: coefEx (CoefEx.lean) for the maps co_j, o_j the class of the cocycle of obCoc at co_j, hx for ob(u) = 0 (the cocycle of u is the sum of the c_j times those of the co_j), ex_mu for mu, key for q(y) = mu(u(y)), keyM for the matrices, and the matrix map T of exists_strictHenselian_absolutely_closed_local_ring, with pert; (7) is obCor (f, fmul, fone, key) and the push of D to A' in slEx; (8) and (9) are sect with diagLift, lowLift and hDef.2 (obWeak supplies m_{C''} ker p'' = 0 as hs', through key). obWeak concludes Obs with the property surjective, m ker = 0 and WLift; absPres turns WLift into the section with sect. About 1,400 lines.

\begin{lemma}[Mazur {\cite{Maz89}}]\label{sha-siegel-step-completion-obstruction-presentation}
Let $\mathcal O'$ be a complete local noetherian ring with finite residue field $k'$, $R$ a complete local noetherian $\mathcal O'$-algebra with residue field $k'$, and $H$ a finite-dimensional obstruction space for minimal generators of $R$ (Definition~\ref{sha-siegel-step-completion-obstruction-generators}). Then there are $a\ge0$ and a surjective homomorphism of $\mathcal O'$-algebras $\mathcal O'[[Z_1,\dots,Z_a]]\to R$ whose kernel is generated by $\dim H$ elements, such that there are $|k'|^a$ homomorphisms of $\mathcal O'$-algebras from $R$ to $k'[\varepsilon]$, $\varepsilon^2=0$.
\end{lemma}
\leanref{Development/AutomorphicArithmetic/AutomorphicVarieties/CommutativeAlgebra/ObsPres.lean}{UpperTheory.V.Alg.obs_pres}

\begin{lemma}\label{sha-siegel-step-completion}
Let $D$ be a Siegel datum, let $T$ be a fit tower for $D$ (Definition~\ref{sha-siegel-tower-towers}), and let $\mathcal O'\supseteq\mathcal O$ be a ring of Definition~\ref{sha-character-parts}, with residue field $k'$.
\begin{enumerate}
\item Let $X=(A,B,P_X)$ be a deformation space for $T$ over $\mathcal O'$ and $x$ a point of $X$ whose image $c$ is residual (Definition~\ref{sha-siegel-tower-space}). There are a complete local noetherian $\mathcal O'$-algebra $R$ with residue field $k'$ and a flat homomorphism of $\mathcal O'$-algebras $f:A\to R$ such that $f$ carries $P_X$ to a universal lift of $c$ and $x$ is $f$ followed by the residue map of $R$. The proof of (1) does not use that $T$ is fit.
\item Suppose that $T$ is deep, let $c$ be a residual flagged pair for the problem of $T$ along $\mathcal O\to k'$, and let $P$ be a universal lift of $c$ to $R$ (Definition~\ref{sha-siegel-tower-pairs}). There are $a,b\ge0$ with $b+2+2d_T\le a+2|\Sigma_T|$ and a surjective homomorphism of $\mathcal O'$-algebras $\mathcal O'[[Z_1,\dots,Z_a]]\to R$ whose kernel is generated by $b$ elements.
\end{enumerate}
\end{lemma}
\leanref{Development/AutomorphicArithmetic/AutomorphicVarieties/WeakShimuraTheory/CompactIntegralFields/CmplEx.lean}{UpperTheory.V.Ws.cmplEx}
\leanref{Development/AutomorphicArithmetic/AutomorphicVarieties/WeakShimuraTheory/CompactIntegralFields/IncLoc.lean}{UpperTheory.V.Ws.x1}
\begin{proof}
(1) The ring $A$ is noetherian, being of finite type over the noetherian ring $B$. The point $x$ is surjective, since $\mathcal O'\to k'$ is, so its kernel $\mathfrak m$ is a maximal ideal, and every element of $A$ is congruent modulo $\mathfrak m$ to the image of an element of $\mathcal O'$. Let $R$ be the $\mathfrak m$-adic completion of $A$ and $f:A\to R$ the canonical map. Then $R$ is a complete local noetherian ring with maximal ideal $\mathfrak mR$, and $f$ is flat \cite[\S8]{Mat86}. Every element of $R$ is congruent modulo $\mathfrak mR$ to an element of $f(A)$, hence to the image of an element of $\mathcal O'$; an element of $\mathcal O'$ whose residue vanishes maps into $\mathfrak m$, hence into $\mathfrak mR$. So $R$ is a complete local noetherian $\mathcal O'$-algebra with residue field $k'$, and $x$ is $f$ followed by the residue map of $R$. By the universal property of completion, if $R'$ is a complete local noetherian $\mathcal O'$-algebra with residue field $k'$ and $g:A\to R'$ is a homomorphism of $\mathcal O'$-algebras whose composite with the residue map of $R'$ is $x$, then $g(\mathfrak m)\subseteq\mathfrak m_{R'}$, and there is a unique homomorphism of $\mathcal O'$-algebras $h:R\to R'$ with $h\circ f=g$.

Two homomorphisms of $\mathcal O'$-algebras from $A$ to a complete local noetherian $\mathcal O'$-algebra $R'$ with residue field $k'$ that carry $P_X$ to the same flagged pair are equal. Indeed, they agree at the coefficients of the polynomials $c_g$, which map to those of the characteristic polynomials of the images of the $r_X(g)$, and at the $\psi^B_{v,i}(s)$, which map to the diagonal entries $\psi_{v,i}(s)$ of the images; so they agree on $B$ by (4) of Definition~\ref{sha-siegel-tower-space}, and on the entries of $P_X$, which generate $A$ over $B$ by (1) there.

The image $f(P_X)$ is a lift of $c$ to $R$. It is a flagged pair along $\mathcal O\to R$, as the image of one. Its entries are continuous for the $\mathfrak m_R$-adic topology, since those of $P_X$ are continuous for the $\mathfrak m_BA$-adic topology and $f$ maps $\mathfrak m_B$ into $\mathfrak m_R$: the homomorphism of $\mathcal O'$-algebras $B\to R\to k'$ is the residue map of $B$, as in the proof of Lemma~\ref{sha-siegel-step-linked}. Its image in $k'$ is the image of $P_X$ under $x$, which is $c$.

The lift $f(P_X)$ is universal. Let $P'$ be a lift of $c$ to a complete local noetherian $\mathcal O'$-algebra $R'$ with residue field $k'$. By (5) of Definition~\ref{sha-siegel-tower-space}, a homomorphism of $\mathcal O'$-algebras $g:A\to R'$ carries $P_X$ to $P'$. Its composite with the residue map of $R'$ carries $P_X$ to $c$, as $x$ does, so it is $x$, by the previous paragraph applied to $k'$. So there is $h:R\to R'$ with $h\circ f=g$, and $h$ carries $f(P_X)$ to $P'$. If $h':R\to R'$ also carries $f(P_X)$ to $P'$, then $h'\circ f$ and $g$ both carry $P_X$ to $P'$, so $h'\circ f=g$ by the previous paragraph, and $h'=h$.

(2) Let $\Phi$ be the problem of $T$, and $\Phi^\circ$ the problem of Lemma~\ref{sha-siegel-step-completion-tangent}. The pair $c$ is residual for $\Phi^\circ$ too, as the conditions of a residual pair do not involve the sets $\mu_v$. Let $d=d_T=[F:\mathbf Q]$.
\begin{enumerate}
\item[(a)] \emph{A universal lift for $\Phi^\circ$.} By Lemma~\ref{sha-siegel-step-coefficients-space}, applied to $\Phi^\circ$, there is a deformation space $X^\circ$ for $\Phi^\circ$ over $\mathcal O'$. As $c$ is a lift of itself to $k'$, (5) of Definition~\ref{sha-siegel-tower-space} gives a point $x$ of $X^\circ$ with image $c$. The proof of (1) uses only Definition~\ref{sha-siegel-tower-space}, so it applies to $X^\circ$ and $x$, and gives a universal lift $P'$ of $c$ to a ring $R'$ for $\Phi^\circ$.
\item[(b)] \emph{The cohomology.} Let $\Gamma_T$, $\pi_T$, $\bar r$, $W=\mathfrak{sp}_4(k')$, the $b_v=\mathfrak b(n_{c,v})$ and $H^i_b=H^i_b(\Gamma_T,W)$ be as in Lemma~\ref{sha-siegel-step-completion-tangent}, for $c=(r_c,(n_{c,v}))$. Lemma~\ref{sha-siegel-step-completion-euler} applies with $k=k'$, $V=\Sigma_T$, the places $A_v=v$, $m=6$ and $e=4$. Indeed, $F$ is totally real. By Lemma~\ref{sha-siegel-step-completion-facts}(4), $\ell$, $F$, the set $S$ of places above $\ell$, $K_T$, $\Gamma_T$ and $\pi_T$ are as in Situation~\ref{sha-descent-link-setting}, and $\Sigma_T$ is in bijection with $S$. The action on $W$ has open kernel, as $\bar r$ has locally constant entries. Each $b_v$ is stable under $\mathrm{loc}_v(D_v(F))$, as $n_{c,v}^*r_c(s)n_{c,v}$ is upper triangular for $s\in D_v(F)$, and has dimension $6$, by Lemma~\ref{sha-siegel-step-completion-adjoint}(1) and~(2). For a complex conjugation $\tau\in G_F$, $\tau^2=1$ and $r_c(\tau)$ is symplectic with multiplier $-1$ by Lemma~\ref{sha-siegel-step-completion-facts}(2), so the invariants of $\pi_T(\tau)$ in $W$ have dimension at most $4$ by Lemma~\ref{sha-siegel-step-completion-adjoint}(2). A linear form $\lambda$ on $W$ with $\lambda(r_c(g)Xr_c(g)^{-1})=\bar\omega_\ell(g)\lambda(X)$ for all $g\in G_F$ is $0$, by Lemma~\ref{sha-siegel-step-completion-adjoint}(3), whose hypothesis holds by Lemma~\ref{sha-siegel-step-completion-facts}(3) for $\mathcal O\to k'$. So $H^1_b$ and $H^2_b$ are finite-dimensional, and $\dim H^2_b+\dim H^0_b+6d\le\dim H^1_b+4d$.
\item[(c)] \emph{The ring $R'$.} By Lemma~\ref{sha-siegel-step-completion-obstruction}, $H^2_b$ is an obstruction space for minimal generators of $R'$, and it is finite-dimensional by (b). By Lemma~\ref{sha-siegel-step-completion-obstruction-presentation}, there are $a\ge0$ and a surjective homomorphism of $\mathcal O'$-algebras $\mathcal O'[[Z_1,\dots,Z_a]]\to R'$ whose kernel is generated by $\dim H^2_b$ elements, with $|k'|^a$ homomorphisms of $\mathcal O'$-algebras $R'\to k'[\varepsilon]$. These are the maps $y\mapsto\bar y+\varepsilon\delta(y)$ with $\delta\in t_{R'}$, in the notation of Lemma~\ref{sha-siegel-step-completion-tangent}, so $a=\dim t_{R'}$. By Lemma~\ref{sha-siegel-step-completion-tangent}, $\dim H^1_b+2\le a+\dim H^0_b$.
\item[(d)] \emph{The ring $R$.} By Lemma~\ref{sha-siegel-step-completion-mu}, with the $\zeta_v$ of Lemma~\ref{sha-siegel-step-completion-facts}(5), there is a surjective homomorphism of $\mathcal O'$-algebras $R'\to R$ whose kernel is generated by $2|\Sigma_T|$ elements. So the composite $\mathcal O'[[Z_1,\dots,Z_a]]\to R$ is surjective, and its kernel is generated by $b=\dim H^2_b+2|\Sigma_T|$ elements: the generators of the kernel of the first map and lifts of the generators of the kernel of the second.
\item[(e)] \emph{The count.} By (b) and (c), $b+2+2d=\dim H^2_b+2d+2+2|\Sigma_T|\le\dim H^1_b-\dim H^0_b+2+2|\Sigma_T|\le a+2|\Sigma_T|$.
\end{enumerate}
% (1) is cmplEx through cmplPt (CmplPt.lean: the completion of A at ker x, with the private dense and lift_ex), Space.noeth, Space.pair_map, Space.cont_map, Space.univ and Space.hom_ext (SpaceMap.lean); the hypothesis Fr.Std of cmplEx is not used in its proof.
% (2) is x1 (IncLoc.lean, 117 lines): std, X1.thFacts, X1.linkT, X1.muLoc, X1.ecG and the private count. In count, (a) is X1.defEx (spaceEx for Phi.noMu, X.univ at c, cmplEx); (b) is X1.descend, X1.adFacts, X1.spDims, X1.dualZero, X1.ccDim and the hypothesis ec = X1.ecG; (c) is X1.absPres (obWeak and sect through obs_pres, tanInj and coneTan); (d) is X1.muQuot and X1.presQuot; (e) is the final omega. Lemma sha-siegel-step-completion-obstruction is about 1,400 lines of Lean.
\end{proof}

\subsection{Dimension bounds and propagation}\label{subsection-sha-propagation}

\begin{lemma}\label{sha-siegel-step-dimensions-pinning}
Let $F$ be a number field, let $u$ be a place of $\bar{\mathbf Q}$ above $\ell$, and let $d=[F_w:\mathbf Q_\ell]$ for the place $w$ of $F$ below $u$. Call a subset $\Sigma$ of $D_u(F)$ \emph{pinning} if the following holds: for every noetherian local domain $R$ in which $\ell=0$, and every homomorphism $\psi:D_u(F)\to R^\times$, continuous for the $\mathfrak m_R$-adic topology, with $\psi(g)\equiv1$ modulo $\mathfrak m_R$ for $g\in I_u(F)$ and $\psi(g)=1$ for $g\in\Sigma$, $\psi$ is trivial on $I_u(F)$.
\begin{enumerate}
\item There are $s_1,\dots,s_d\in I_u(F)$ that form a pinning set.
\item If $\ell$ is odd, there is $s_0\in I_u(F)$ such that $s_0$ and the $s\in I_u(F)$ with $\chi_{\mathrm{cyc}}(s)=1$ form a pinning set.
\end{enumerate}
\end{lemma}
\leanref{Development/AutomorphicArithmetic/AutomorphicVarieties/BuildingsResidueChars/ExistsPinningElementsOfLocAb.lean}{UpperTheory.exists_pinning_elements_of_locAb}
\leanref{Development/AutomorphicArithmetic/AutomorphicVarieties/BuildingsResidueChars/NonemptyLocalWeilDataDecompositionGroupOfMemPlac.lean}{UpperTheory.nonempty_localWeilData_decompositionGroup_of_mem_place}
\leanref{Development/AutomorphicArithmetic/AutomorphicVarieties/BuildingsResidueChars/ExistsMemPinsInsertKerOfPadicCharacter.lean}{UpperTheory.exists_mem_pins_insert_ker_of_padic_character}
\leanref{Development/AutomorphicArithmetic/AutomorphicVarieties/WeakShimuraTheory/NnDefs.lean}{UpperTheory.V.Ws.eq_one_of_pow}
\begin{proof}
First, let $R$ be a local domain in which $\ell=0$, and let $x\in R$ with $x\equiv1$ modulo $\mathfrak m_R$ and $x^N=1$ for some $N\ge1$; then $x=1$. Indeed, write $N=\ell^kN'$ with $\ell\nmid N'$ and put $y=x^{\ell^k}$. Then $0=y^{N'}-1=(y-1)(1+y+\dots+y^{N'-1})$, and the second factor is congruent to $N'$, a unit; so $y=1$. As $\ell=0$ in $R$, $(x-1)^{\ell^k}=x^{\ell^k}-1=0$, so $x=1$, $R$ being a domain.

(1) Let $\mu_u$ be the subgroup of the $g\in D_u(F)$ whose image in the quotient of $D_u(F)$ by the closure of its commutator subgroup has $\ell$-power order. By local class field theory for $F_w$ \cite[Chapter~VII]{NSW08}, the maximal abelian pro-$\ell$ quotient of $D_u(F)$ modulo the image of $\mu_u$ is $\mathbf Z_\ell^d\times\mathbf Z_\ell$, with inertia mapping onto $\mathbf Z_\ell^d\times\{0\}$, as the pro-$\ell$ completion of the units of $F_w$ is the product of a finite group and $\mathbf Z_\ell^d$. That is, there is a continuous surjective homomorphism $a:D_u(F)\to\mathbf Z_\ell^d\times\mathbf Z_\ell$ with $a(I_u(F))=\mathbf Z_\ell^d\times\{0\}$, such that every homomorphism from $D_u(F)$ to a finite abelian $\ell$-group with open kernel containing $\mu_u$ kills $\ker a$. Choose $s_k\in I_u(F)$ with $a(s_k)=(e_k,0)$, for the standard basis $e_1,\dots,e_d$ of $\mathbf Z_\ell^d$. Let $R$ and $\psi$ be as in the definition, with $\psi(s_k)=1$ for all $k$.
\begin{enumerate}
\item[(a)] The reduction of $\psi$ modulo $\mathfrak m_R$ is locally constant on the compact group $D_u(F)$, so its image is finite, of order $M$ say, and $\psi(g)^M\equiv1$ for all $g$.
\item[(b)] Let $g\in\mu_u$. As $R^\times$ is commutative and $\psi$ is continuous, $\psi$ factors through the quotient of $D_u(F)$ by the closure of its commutator subgroup, so $\psi(g)^{\ell^n}=1$ for some $n$. By the first paragraph, applied to $\psi(g)^M$, $\psi(g)^M=1$.
\item[(c)] For $N\ge1$ let $\varphi_N$ be $\psi^M$ modulo $\mathfrak m_R^N$, a locally constant homomorphism to $(R/\mathfrak m_R^N)^\times$. If $\psi(g)^M=1+y$, then $y\in\mathfrak m_R$ by (a), and $(1+y)^{\ell^N}=1+y^{\ell^N}$ as $\ell=0$ in $R$, with $y^{\ell^N}\in\mathfrak m_R^N$. So the image of $\varphi_N$ is a finite $\ell$-group, and $\varphi_N$ kills $\mu_u$ by (b). Hence $\ker a\subseteq\ker\varphi_N$.
\item[(d)] The subgroup $a(\ker\varphi_N)$ is closed, as the image of a compact set, and contains the $(e_k,0)$, as $\psi(s_k)=1$; so it contains $\mathbf Z^d\times\{0\}$, and by density $\mathbf Z_\ell^d\times\{0\}$.
\item[(e)] Let $g\in I_u(F)$. By (d), $a(g)=a(h)$ for some $h\in\ker\varphi_N$; then $h^{-1}g\in\ker a\subseteq\ker\varphi_N$ by (c), so $\varphi_N(g)=1$, that is $\psi(g)^M-1\in\mathfrak m_R^N$. This holds for every $N$, so $\psi(g)^M=1$ by Krull's intersection theorem \cite[Theorem~8.10]{Mat86}, and $\psi(g)=1$ by the first paragraph, as $\psi(g)\equiv1$.
\end{enumerate}

(2) As $\ell$ is odd, $\mathbf Z_\ell^\times\cong\mathbf Z/(\ell-1)\times\mathbf Z_\ell$ is procyclic, and so is its closed subgroup $\chi_{\mathrm{cyc}}(I_u(F))$, the image of a compact group. Let $s_0\in I_u(F)$ be such that $\chi_{\mathrm{cyc}}(s_0)$ generates it topologically. Let $R$ and $\psi$ be as in the definition, trivial at $s_0$ and at the $s\in I_u(F)$ with $\chi_{\mathrm{cyc}}(s)=1$. The subgroup $H=I_u(F)\cap\ker\psi$ is closed, since $\psi$ is continuous and $R$ is $\mathfrak m_R$-adically separated; so $\chi_{\mathrm{cyc}}(H)$ is closed, contains $\chi_{\mathrm{cyc}}(s_0)$, and hence contains $\chi_{\mathrm{cyc}}(I_u(F))$. For $g\in I_u(F)$ choose $h\in H$ with $\chi_{\mathrm{cyc}}(h)=\chi_{\mathrm{cyc}}(g)$; then $gh^{-1}\in I_u(F)$ and $\chi_{\mathrm{cyc}}(gh^{-1})=1$, so $\psi(g)=\psi(gh^{-1})\psi(h)=1$.
\end{proof}
% (1) is exists_pinning_elements_of_locAb (ExistsPinningElementsOfLocAb.lean), with the structure W.LocAb (the map a, inertia, and univ for finite l-groups killing muG) given for D_u(F) by nonempty_localWeilData_decompositionGroup_of_mem_place from W.lam_local, the Lean's own proof of local class field theory in this form (LamLocal.lean, about 740 lines); (a) is the order M of the range of rho_1, (b) exists_pow_l_pow_eq_one_of_mem_muG with eq_one_of_pow (the first paragraph), (c) hK through A.univ, (d) hmem (closedness and PadicInt.denseRange_intCast), (e) hgN and Ideal.iInf_pow_eq_bot_of_isLocalRing. (2) is exists_mem_pins_insert_ker_of_padic_character, for any compact group, closed subgroup and continuous character to Z_l^x, with Kum.existsGen for the procyclicity and isClosed_ker_of_adicCont.

\begin{definition}\label{sha-siegel-step-dimensions-tangent-pairs}
Let $D$ be a Siegel datum, $T$ a tower for $D$ with field $F$ (Definition~\ref{sha-siegel-tower-towers}), and $N_\Phi$ the set of the problem $\Phi$ of $T$ (Definition~\ref{sha-siegel-tower-pairs}); let $k$ be a field, and let $\mathfrak{sp}_4(k)$ and $X\mapsto X^*$ be as in Definition~\ref{sha-siegel-datum}. Let $c=(r,(n_v)_{v\in\Sigma_T})$ be a pair of a homomorphism $r:G_F\to M_4(k)$ and matrices $n_v\in M_4(k)$, and let $W$ be a set of weights. The space $\mathrm{Tan}_W(c)$ of \emph{tangent pairs} of $c$ supported in $W$ is the set of pairs $(x,(Y_v)_{v\in\Sigma_T})$ of a locally constant map $x:G_F\to\mathfrak{sp}_4(k)$ and elements $Y_v\in\mathfrak{sp}_4(k)$ such that: $x(gh)=x(g)+r(g)x(h)r(g^{-1})$; $x(g)=0$ for $g\in N_\Phi$; $x(g_0^T)$ is diagonal; $x(g)r(g)$ is supported in $W$ for every $g$; $(n_v^*Y_vn_v)_{ij}=0$ for $i\le j$; and $n_v^*\big(x(s)+r(s)Y_vr(s^{-1})-Y_v\big)n_v$ is upper triangular for $s\in D_v(F)$. These conditions are linear, so $\mathrm{Tan}_W(c)$ is a $k$-vector space.
\end{definition}
\leanref{Development/AutomorphicArithmetic/AutomorphicVarieties/WeakShimuraTheory/RedIdeal.lean}{UpperTheory.V.Ws.Fr.Tan}
\leanref{Development/AutomorphicArithmetic/AutomorphicVarieties/WeakShimuraTheory/RedTanDefs.lean}{UpperTheory.V.Ws.Fr.TanS}

\begin{lemma}\label{sha-siegel-step-dimensions-tangent-flow}
Let $D$, $T$, $F$, $k$ and $W$ be as in Definition~\ref{sha-siegel-step-dimensions-tangent-pairs}, and let $m_1,\dots,m_4$ be integers such that $m_i+m_{5-i}$ does not depend on $i$. For $p\in\mathbf Z$, call $X\in M_4(k)$ \emph{of degree} $\ge p$ if $X_{ij}=0$ whenever $m_i-m_j<p$, and let $X^{(p)}$ be the matrix whose entry $(i,j)$ is $X_{ij}$ if $m_i-m_j=p$ and $0$ otherwise. Let $c=(r,(n_v))$ and $c_0=(r_0,(n_{0,v}))$ be pairs as in that definition such that each $r(g)$ and each $n_v$ is of degree $\ge0$, with $r(g)^{(0)}=r_0(g)$ and $n_v^{(0)}=n_{0,v}$. Then $\dim\mathrm{Tan}_W(c)\le\dim\mathrm{Tan}_W(c_0)$.
\end{lemma}
\leanref{Development/AutomorphicArithmetic/AutomorphicVarieties/WeakShimuraTheory/RedTanFlow.lean}{UpperTheory.V.Ws.tan_flow}
\leanref{Development/AutomorphicArithmetic/AutomorphicVarieties/WeakShimuraTheory/RedTanFlow.lean}{UpperTheory.V.Ws.tan_lowest}
\leanref{Development/AutomorphicArithmetic/AutomorphicVarieties/WeakShimuraTheory/RedTanFlow.lean}{UpperTheory.V.Ws.rank_le_of_lowest}
\begin{proof}
For a pair $y=(x,(Y_v))$ of a map $x:G_F\to M_4(k)$ and matrices $Y_v$, write $y^{(p)}=(g\mapsto x(g)^{(p)},(Y_v^{(p)}))$, and say that $y$ is of degree $\ge p$ if every $x(g)$ and every $Y_v$ is.
\begin{enumerate}
\item[(a)] \emph{Rules.} If $X$ is of degree $\ge p$ and $X'$ of degree $\ge q$, then $XX'$ is of degree $\ge p+q$ and $(XX')^{(p+q)}=X^{(p)}X'^{(q)}$, as the entry $(i,j)$ of $XX'$ is a sum of products $X_{il}X'_{lj}$ with $m_i-m_j=(m_i-m_l)+(m_l-m_j)$. As $m_i-m_j=m_{5-j}-m_{5-i}$, the map $X\mapsto X^{(p)}$ commutes with $X\mapsto X^*$, so it preserves $\mathfrak{sp}_4(k)$, and $X^*$ is of degree $\ge p$ when $X$ is. The matrix $X^{(p)}$ vanishes wherever $X$ does.
\item[(b)] \emph{Lowest parts.} Let $y=(x,(Y_v))\in\mathrm{Tan}_W(c)$ be of degree $\ge p$. Then $y^{(p)}\in\mathrm{Tan}_W(c_0)$. Indeed, by (a), $r(g)x(h)r(g^{-1})$ is of degree $\ge p$, with degree-$p$ part $r_0(g)x(h)^{(p)}r_0(g^{-1})$, so the cocycle condition for $x$ gives that for $g\mapsto x(g)^{(p)}$ and $r_0$; this map is locally constant, vanishes on $N_\Phi$, and is diagonal at $g_0^T$; $x(g)^{(p)}r_0(g)=(x(g)r(g))^{(p)}$ is supported in $W$; and $n_{0,v}^*Y_v^{(p)}n_{0,v}=(n_v^*Y_vn_v)^{(p)}$ and $n_{0,v}^*\big(x(s)^{(p)}+r_0(s)Y_v^{(p)}r_0(s^{-1})-Y_v^{(p)}\big)n_{0,v}=\big(n_v^*(x(s)+r(s)Y_vr(s^{-1})-Y_v)n_v\big)^{(p)}$ vanish wherever the matrices on the right do.
\item[(c)] \emph{The count.} Let $N$ exceed every $|m_i-m_j|$, so that every pair is of degree $\ge-N$ and a pair of degree $\ge N$ is $0$. For $p\in\mathbf Z$ let $T_p\subseteq\mathrm{Tan}_W(c)$ and $U_p\subseteq\mathrm{Tan}_W(c_0)$ be the subspaces of pairs of degree $\ge p$. On $T_p$, the linear map $y\mapsto y^{(p)}$ has kernel $T_{p+1}$, and by (b) its image lies in the image of $U_p$ under the same map, as $y^{(p)}$ is of degree $\ge p$ and equal to its own part of degree $p$. On $U_p$ the kernel is $U_{p+1}$. So $\dim T_p\le\dim U_p^{(p)}+\dim T_{p+1}$ and $\dim U_p=\dim U_p^{(p)}+\dim U_{p+1}$, where $U_p^{(p)}$ is that image. By descending induction from $T_N=U_N=0$, $\dim T_p\le\dim U_p$ for every $p$; for $p=-N$ this is the claim.
\end{enumerate}
\end{proof}
% tan_flow (RedTanFlow.lean): FGe m p is "of degree >= p", cmp m p is X -> X^(p), Lim m A A_0 is "of degree >= 0 with degree-0 part A_0", and Symp m is the condition on m. (a) is FGe.mul, cmp_mul and cmp_adj (RedTanDefs.lean); (b) is tan_lowest; (c) is rank_le_of_lowest, with K p for the pairs of degree >= p.

\begin{lemma}\label{sha-siegel-step-dimensions-tangent-split}
Let $D$ be a Siegel datum, $T$ a tower for $D$ with field $F$ (Definition~\ref{sha-siegel-tower-towers}), $k$ a field, $a:\mathcal O\to k$ a homomorphism that kills $\mathfrak m_{\mathcal O}$, and $d,C\ge0$ such that $T$ satisfies the cohomology bounds $(d,C)$ along $a$, with the spaces $Z_w$ and $S_w$ of Definition~\ref{sha-siegel-tower-ideal}(3). Let $W$ be a set of weights, and $c=(r_1,(n_v))$ a pair as in Definition~\ref{sha-siegel-step-dimensions-tangent-pairs} with $r_1(g)=\operatorname{diag}(a(\chi_1(g)),\dots,a(\chi_4(g)))$ and each $n_v$ symplectic with multiplier $1$. Suppose that some entry $(i_0,j_0)$ whose weight is an odd root has the following property: for every $(x,(Y_v))\in\mathrm{Tan}_W(c)$ with $x(g)_{i_0j_0}=0$ for all $g$, and every entry $(i,j)$ with $i+j\le5$ whose weight $w$ is an odd root, there are $m_v\in k$, for $v\in\Sigma_T$, with $x(s)_{ij}=m_v-a(\chi^w(s))m_v$ for all $v$ and $s\in D_v(F)$. Then $\dim\mathrm{Tan}_W(c)\le d+20C$.
\end{lemma}
\leanref{Development/AutomorphicArithmetic/AutomorphicVarieties/WeakShimuraTheory/CompactIntegralFields/RedTan.lean}{rank_of_oneOdd}
\leanref{Development/AutomorphicArithmetic/AutomorphicVarieties/WeakShimuraTheory/RedTanDefs.lean}{UpperTheory.V.Ws.Fr.OneOdd}
\begin{proof}
\begin{enumerate}
\item[(a)] \emph{Entries.} Let $(x,(Y_v))\in\mathrm{Tan}_W(c)$ and let $(i,j)$ be an entry, of weight $w$. Then $z(g)=x(g)_{ij}$ lies in $Z_w$. It is locally constant. As $r_1$ is diagonal and $\chi_i\chi_j^{-1}=\chi^w$, the entry $(i,j)$ of $r_1(g)x(h)r_1(g^{-1})$ is $a(\chi^w(g))x(h)_{ij}$, so $z(gh)=z(g)+a(\chi^w(g))z(h)$. It vanishes on $N_\Phi$. If $w\ne0$, then $i\ne j$, and $z(g_0^T)=0$, as $x(g_0^T)$ is diagonal.
\item[(b)] \emph{The entry $(i_0,j_0)$.} Let $E$ be the linear map from $\mathrm{Tan}_W(c)$ to the maps $G_F\to k$ that sends $(x,(Y_v))$ to $g\mapsto x(g)_{i_0j_0}$. By (a), its image lies in $Z_{w_0}$ for the odd root $w_0$ of $(i_0,j_0)$, which has dimension at most $d+C$.
\item[(c)] \emph{The kernel of $E$.} Let $P$ be the set of the ten entries $(i,j)$ with $i+j\le5$, and $Q$ the set of the six entries with $j<i$. For $(i,j)\in P$ of weight $w$, let $M_{ij}=S_w$ if $w$ is an odd root, and $M_{ij}=Z_w$ otherwise, in which case $w\in\{0,(0,1),(0,-1)\}$; in both cases $\dim M_{ij}\le C$. By (a) and the hypothesis, for $(x,(Y_v))$ in the kernel of $E$ and $(i,j)\in P$, the map $g\mapsto x(g)_{ij}$ lies in $M_{ij}$. Consider the linear map from the kernel of $E$ to $\prod_{(i,j)\in P}M_{ij}\times k^{\Sigma_T\times Q}$ that sends $(x,(Y_v))$ to the maps $g\mapsto x(g)_{ij}$ for $(i,j)\in P$ and the entries $(n_v^*Y_vn_v)_{ij}$ for $v\in\Sigma_T$ and $(i,j)\in Q$. It is injective. Indeed, if $(x,(Y_v))$ maps to $0$, then $x(g)_{ij}=0$ for $i+j\le5$; as $x(g)\in\mathfrak{sp}_4(k)$, its entry $(i,j)$ is $\pm$ its entry $(5-j,5-i)$, so $x(g)=0$. And $n_v^*Y_vn_v$ vanishes at the entries with $j<i$ by assumption and at those with $i\le j$ by the definition of $\mathrm{Tan}_W(c)$; so it is $0$, and $Y_v=0$, as $n_v^*=n_v^{-1}$.
\item[(d)] \emph{The bound.} By (c), the kernel of $E$ has dimension at most $10C+6|\Sigma_T|\le16C$, as $|\Sigma_T|\le C$. With (b), $\dim\mathrm{Tan}_W(c)\le d+C+16C\le d+20C$.
\end{enumerate}
\end{proof}
% rank_of_oneOdd (RedTan.lean, about 85 lines): (a) is ent, with Ok.conj_split; (b) is E i_0 j_0 and r1 (hb.odd); (c) is F, hF (with the sp4 symmetry through adj) and r2, r3 (hb.str, hb.even, hb.card); (d) is r4 and the final count d + C + (10C + 6C) <= d + 20C. The hypothesis is Fr.OneOdd, with LocTriv for the elements m_v.

\begin{lemma}\label{sha-siegel-step-dimensions-tangent-cases}
Let $D$, $T$, $F$, $k$, $a$, $d$ and $C$ be as in Lemma~\ref{sha-siegel-step-dimensions-tangent-split}, with $2\ne0$ in $k$; let $W_1,\dots,W_9$ be the sets of weights of Lemma~\ref{sha-siegel-step-dimensions}, $r_1(g)=\operatorname{diag}(a(\chi_1(g)),\dots,a(\chi_4(g)))$, and $E_{ij}$ the matrix units.
\begin{enumerate}
\item If $W$ is $W_1$, $W_5$ or $W_7$, the pair $(r_1,(1))$ satisfies the hypothesis of Lemma~\ref{sha-siegel-step-dimensions-tangent-split}, with $(i_0,j_0)=(2,3)$ for $W_1$ and $W_7$ and $(i_0,j_0)=(1,2)$ for $W_5$.
\item Let $b_v\in k^\times$ for $v\in\Sigma_T$, and suppose that $a(\chi_1(s))=a(\chi_3(s))$ for every $v$ and $s\in D_v(F)$. Then $1+b_v(E_{31}+E_{42})$ is symplectic with multiplier $1$, and if $W$ is $W_1$, $W_3$, $W_5$, $W_6$, $W_7$ or $W_9$, the pair $(r_1,(1+b_v(E_{31}+E_{42})))$ satisfies the hypothesis of Lemma~\ref{sha-siegel-step-dimensions-tangent-split}, with $(i_0,j_0)=(2,3)$ for $W_1$, $W_3$ and $W_7$, $(1,2)$ for $W_5$ and $W_6$, and $(1,4)$ for $W_9$.
\end{enumerate}
\end{lemma}
\leanref{Development/AutomorphicArithmetic/AutomorphicVarieties/WeakShimuraTheory/CompactIntegralFields/RedTan.lean}{oneOdd_split}
\leanref{Development/AutomorphicArithmetic/AutomorphicVarieties/WeakShimuraTheory/CompactIntegralFields/RedTan.lean}{oneOdd_splitPs}
\begin{proof}
Let $c$ be the pair of (1) or (2), with matrices $n_v$, let $(x,(Y_v))\in\mathrm{Tan}_W(c)$ with $x(g)_{i_0j_0}=0$ for all $g$, and call an entry $(i,j)$ whose weight $w$ is an odd root \emph{trivial} if there are $m_v\in k$ with $x(s)_{ij}=m_v-a(\chi^w(s))m_v$ for all $v$ and $s\in D_v(F)$. We show that the entries $(i,j)$ with $i+j\le5$ whose weights are odd roots are trivial. The entry $(i_0,j_0)$ is trivial, with $m_v=0$. An entry $(i,j)$ with $i\ne j$ whose weight is not in $W$ is trivial, with $m_v=0$: the entry $(i,j)$ of $x(g)r_1(g)$, which is supported in $W$, is $x(g)_{ij}a(\chi_j(g))$, and $a(\chi_j(g))$ is a unit.

(1) Here $n_v=1$. Let $(i,j)$ be an entry with $j<i$ whose weight $w$ is an odd root. For $s\in D_v(F)$, the matrix $x(s)+r_1(s)Y_vr_1(s^{-1})-Y_v$ is upper triangular, and its entry $(i,j)$ is $x(s)_{ij}+a(\chi^w(s))(Y_v)_{ij}-(Y_v)_{ij}$; so $(i,j)$ is trivial, with $m_v=(Y_v)_{ij}$. The entries with $i<j$ and $i+j\le5$ whose weights are odd roots are $(1,2)$, $(1,4)$ and $(2,3)$, of weights $(1,0)$, $(1,1)$ and $(-1,1)$. Of these three weights only $(-1,1)$ lies in $W_1$ and in $W_7$, so $(1,2)$ and $(1,4)$ are trivial and $(2,3)$ is $(i_0,j_0)$; and only $(1,0)$ lies in $W_5$, so $(1,4)$ and $(2,3)$ are trivial and $(1,2)$ is $(i_0,j_0)$.

(2) Write $n_v=1+b_v(E_{31}+E_{42})$. A computation shows that $n_v^*n_v=1$, and that every $Z\in\mathfrak{sp}_4(k)$ for which $n_v^*Zn_v$ is upper triangular satisfies
\[
Z_{21}=-b_vZ_{23},\qquad Z_{41}=-b_v^2Z_{23},\qquad Z_{32}=2b_vZ_{12}+b_v^2Z_{14}.
\]
For $s\in D_v(F)$, $a(\chi^{(0,1)}(s))=a(\chi_1(s)\chi_3(s)^{-1})=1$, so $a(\chi^w(s))=a(\chi^{(w_1,0)}(s))$ for $w=(w_1,w_2)$. The entries with $i+j\le5$ whose weights are odd roots are $(2,1)$, $(2,3)$, $(4,1)$, of weights $(-1,0)$, $(-1,1)$, $(-1,-1)$, for which $a(\chi^w(s))=a(\chi^{(-1,0)}(s))$, and $(1,2)$, $(1,4)$, $(3,2)$, of weights $(1,0)$, $(1,1)$, $(1,-1)$, for which $a(\chi^w(s))=a(\chi^{(1,0)}(s))$. Let $Z_v(s)=x(s)+r_1(s)Y_vr_1(s^{-1})-Y_v$, which lies in $\mathfrak{sp}_4(k)$, with $n_v^*Z_v(s)n_v$ upper triangular for $s\in D_v(F)$. Its entry $(i,j)$ is $x(s)_{ij}-\big((Y_v)_{ij}-a(\chi^w(s))(Y_v)_{ij}\big)$, so $(i,j)$ is trivial if and only if there are $m'_v$ with $Z_v(s)_{ij}=m'_v-a(\chi^w(s))m'_v$ for all $v$ and $s$. As $b_v$ and $2b_v$ are units, these three relations show: if one of $(2,1)$, $(2,3)$, $(4,1)$ is trivial, then all three are; and if two of $(1,2)$, $(1,4)$, $(3,2)$ are trivial, then so is the third. For instance, if $(2,3)$ is trivial with the elements $m'_v$, then $(2,1)$ is trivial with the elements $-b_vm'_v$, and if $(1,2)$ and $(1,4)$ are trivial with $m'_v$ and $m''_v$, then $(3,2)$ is trivial with $2b_vm'_v+b_v^2m''_v$. It remains to find one trivial entry among the first three and two among the last three:
for $W_1$ and $W_7$, $(2,3)=(i_0,j_0)$, and the weights $(1,0)$ of $(1,2)$ and $(1,1)$ of $(1,4)$ are not in $W$;
for $W_3$, $(2,3)=(i_0,j_0)$, and the weights $(1,0)$ of $(1,2)$ and $(1,-1)$ of $(3,2)$ are not in $W_3$;
for $W_5$, the weight $(-1,1)$ of $(2,3)$ is not in $W_5$, $(1,2)=(i_0,j_0)$, and the weight $(1,1)$ of $(1,4)$ is not in $W_5$;
for $W_6$, the weight $(-1,-1)$ of $(4,1)$ is not in $W_6$, $(1,2)=(i_0,j_0)$, and the weight $(1,-1)$ of $(3,2)$ is not in $W_6$;
for $W_9$, the weight $(-1,0)$ of $(2,1)$ is not in $W_9$, $(1,4)=(i_0,j_0)$, and the weight $(1,0)$ of $(1,2)$ is not in $W_9$.
\end{proof}
% oneOdd_split (RedTan.lean): key, a decide over redTypes, and low (the entries below the diagonal, through Ok.conj_split), with zero_of_not_mem and triv_of_zero. oneOdd_splitPs: ps (a(chi_1) = a(chi_3) on D_v), rel through rel_ps (the three relations, checked by linear_combination), keyB (a decide over redTypes), z0, z1 and t1 to t6 for the propagation; gsp_ps is the symplecticity of n_v. The Lean indexes from 0, so its entries (1,0), (1,2), (3,0), (0,1), (0,3), (2,1) are the entries (2,1), (2,3), (4,1), (1,2), (1,4), (3,2) of the text. The choice of (i_0,j_0) in the text is one that keyB and key accept.

\begin{lemma}\label{sha-siegel-step-dimensions-tangent}
Let $D$ be a Siegel datum, $T$ a fit tower for $D$ with field $F$ (Definition~\ref{sha-siegel-tower-towers}), $k$ a field, $a:\mathcal O\to k$ a homomorphism, and $d,C\ge0$ such that $T$ satisfies the cohomology bounds $(d,C)$ along $a$ (Definition~\ref{sha-siegel-tower-ideal}). Let $c=(r_c,(n_{c,v}))$ be a residual ample flagged pair for the problem of $T$ along $a$, and let $W$ be one of the nine sets $W_1,\dots,W_9$ of Lemma~\ref{sha-siegel-step-dimensions} such that $r_c$ is supported in $W$ (Definition~\ref{sha-siegel-tower-pairs}). Let $\mathrm{Tan}_W(c)$ be the space of tangent pairs of $c$ supported in $W$ (Definition~\ref{sha-siegel-step-dimensions-tangent-pairs}): the pairs $(x,(Y_v)_{v\in\Sigma_T})$ of a locally constant map $x:G_F\to\mathfrak{sp}_4(k)$ and elements $Y_v\in\mathfrak{sp}_4(k)$ such that: $x(gh)=x(g)+r_c(g)x(h)r_c(g)^{-1}$; $x(g)=0$ for $g\in N_\Phi$; $x(g_0^T)$ is diagonal; $x(g)r_c(g)$ is supported in $W$ for every $g$; $(n_{c,v}^*Y_vn_{c,v})_{ij}=0$ for $i\le j$; and $n_{c,v}^*\big(x(s)+r_c(s)Y_vr_c(s)^{-1}-Y_v\big)n_{c,v}$ is upper triangular for $s\in D_v(F)$. Then $\mathrm{Tan}_W(c)$ has dimension at most $d+20C$. (The last clause of the fourth condition of ampleness, that for each $v$ some $s\in D_v(F)$ has $a(\chi^{(1,0)}(s))\ne1$, is not used in the proof.)
\end{lemma}
\leanref{Development/AutomorphicArithmetic/AutomorphicVarieties/WeakShimuraTheory/CompactIntegralFields/RedTan.lean}{UpperTheory.V.Ws.redTan}
\leanref{Development/AutomorphicArithmetic/AutomorphicVarieties/WeakShimuraTheory/RedIdeal.lean}{UpperTheory.V.Ws.Fr.Tan}
\begin{proof}
As $T$ satisfies the cohomology bounds along $a$, $a$ kills $\mathfrak m_{\mathcal O}$, so $k$ has the odd characteristic $\ell$ and $2\ne0$ in $k$. Let $r_1(g)=\operatorname{diag}(a(\chi_1(g)),\dots,a(\chi_4(g)))$, and $E_{ij}$ the matrix units. As $r_c$ is supported in $W$, every weight $w\ne0$ of an entry of $r_c$ lies in $W$.
\begin{enumerate}
\item[(a)] \emph{Split pairs.} Suppose that $W$ contains $(0,-1)$ and one of $(-1,1)$ and $(-1,0)$, or contains $(-1,0)$ and $(1,-1)$. Among the nine sets, these are $W_1$, $W_5$ and $W_7$. If $r_c$ is supported in $\{(-1,1),(-1,0),(-1,-1),(0,-1)\}$ and each $n_{c,v}$ in $\{(-1,0),(0,-1),(-1,-1)\}$, let $m=(0,2,1,3)$; otherwise, as $c$ is residual, $r_c$ is supported in $\{(-1,0),(1,-1),(0,-1),(-1,-1)\}$, and let $m=(0,1,2,3)$. In both cases $m_i+m_{5-i}=3$, and every off-diagonal entry $(i,j)$ of an $r_c(g)$ or an $n_{c,v}$ that may be nonzero has $m_i>m_j$. In the first case these are among the entries $(2,3)$, $(2,1)$, $(4,3)$, $(4,1)$, $(3,1)$ and $(4,2)$, of weights $(-1,1)$, $(-1,0)$, $(-1,0)$, $(-1,-1)$, $(0,-1)$ and $(0,-1)$. In the second case they lie below the diagonal, as $n_{c,v}$ is lower triangular and the entries of the four weights are $(2,1)$, $(4,3)$, $(3,2)$, $(3,1)$, $(4,2)$ and $(4,1)$. The diagonal entries are $a(\chi_i(g))$ and $1$, as $c$ is residual. So Lemma~\ref{sha-siegel-step-dimensions-tangent-flow}, with $c_0=(r_1,(1))$, gives $\dim\mathrm{Tan}_W(c)\le\dim\mathrm{Tan}_W(r_1,(1))$, which is at most $d+20C$ by Lemmas~\ref{sha-siegel-step-dimensions-tangent-cases}(1) and~\ref{sha-siegel-step-dimensions-tangent-split}.
\item[(b)] \emph{Type V.} Suppose that $c$ is of type V, and that $W$ contains $(-1,0)$, or contains $(-1,1)$ and $(-1,-1)$. Among the nine sets, these are $W_1$, $W_3$, $W_5$, $W_6$, $W_7$ and $W_9$. Let $b_v=(n_{c,v})_{31}$, which is nonzero. For $s\in D_v(F)$, $a(\chi_1(s))=a(\chi_3(s))$. Indeed, as $n_{c,v}^*=n_{c,v}^{-1}$, $r_c(s)n_{c,v}=n_{c,v}U$ with $U=n_{c,v}^*r_c(s)n_{c,v}$ upper triangular and $U_{11}=a(\chi_1(s))$. The entries $(3,1)$, $(3,2)$ and $(3,4)$ of $r_c(s)$ vanish, as their weights $(0,-1)$, $(1,-1)$ and $(1,0)$ are not in $\{(-1,1),(-1,0),(-1,-1)\}$, and $r_c(s)_{33}=a(\chi_3(s))$; so comparing the entries $(3,1)$ gives $a(\chi_3(s))b_v=b_va(\chi_1(s))$. Let $m=(0,1,0,1)$. The entries of $r_c(g)$ of weights $(-1,1)$, $(-1,0)$ and $(-1,-1)$ are $(2,3)$, $(2,1)$, $(4,3)$ and $(4,1)$, with $m_i>m_j$, so $r_c(g)$ is of degree $\ge0$ with $r_c(g)^{(0)}=r_1(g)$. The matrix $n_{c,v}$ is lower triangular with diagonal entries $1$, and its entry $(3,2)$, of weight $(1,-1)$, vanishes; so it is of degree $\ge0$, and as the entry $(3,1)$ of $n_{c,v}^*n_{c,v}=1$ is $(n_{c,v})_{31}-(n_{c,v})_{42}$, $n_{c,v}^{(0)}=1+b_v(E_{31}+E_{42})$. So Lemma~\ref{sha-siegel-step-dimensions-tangent-flow} gives $\dim\mathrm{Tan}_W(c)\le\dim\mathrm{Tan}_W(r_1,(1+b_v(E_{31}+E_{42})))$, which is at most $d+20C$ by Lemmas~\ref{sha-siegel-step-dimensions-tangent-cases}(2) and~\ref{sha-siegel-step-dimensions-tangent-split}.
\item[(c)] \emph{The cases.} As $c$ is ample, one of the four conditions of Definition~\ref{sha-siegel-tower-pairs}(3) holds. If $r_c$ has an entry of weight $(0,-1)$ and one of weight $(-1,1)$ or $(-1,0)$, or entries of weights $(-1,0)$ and $(1,-1)$, then these weights lie in $W$, and (a) applies. If $c$ is of type V and $r_c$ has an entry of weight $(-1,0)$, or entries of weights $(-1,1)$ and $(-1,-1)$, then (b) applies; the fourth condition is used only through these entries and the type V.
\end{enumerate}
\end{proof}
% redTan (RedTan.lean, 352 lines, with RedTanFlow.lean and RedTanDefs.lean; about 640 lines): ent for the first paragraph; (a) is A, with tan_flow for m = ![0,2,1,3] (shape nP) or ![0,1,2,3] (shape nM), limR and limN, then rank_of_oneOdd and oneOdd_split; (b) is B, with residual_chars_zero_two_eq, tan_flow for m = ![0,1,0,1], limR and limPs, then rank_of_oneOdd and oneOdd_splitPs; (c) is the final case analysis of Good. The last condition of ampleness (some s with a(chi^(1,0)(s)) != 1) is not used.

\begin{lemma}\label{sha-siegel-step-dimensions-torus}
Let $K$ be a field, $H$ a group, $r:H\to M_4(K)$ a homomorphism with $r(h)$ symplectic with multiplier $\nu(h)\in K$ for every $h$, and $h_0\in H$ with $r(h_0)$ diagonal with pairwise distinct diagonal entries, such that the $r(h)$ span $M_4(K)$.
\begin{enumerate}
\item There are $s,t\in K$ such that every subfield $E\subseteq K$ that contains $s$, $t$, the $\nu(h)$, the traces of the $r(h)$ and the diagonal entries of $r(h_0)$ contains every entry of every $r(h)$.
\item If $n\in M_4(K)$ is symplectic with multiplier $1$ and lower triangular with diagonal entries $1$, then a subfield of $K$ that contains $n_{21}$, $n_{31}$, $n_{41}$ and $n_{32}$ contains every entry of $n$.
\end{enumerate}
\end{lemma}
\leanref{Development/AutomorphicArithmetic/AutomorphicVarieties/WeakShimuraTheory/FibreTraces/FibTTrace.lean}{UpperTheory.V.Ws.fibTorus}
\begin{proof}
(1) Let $d_1,\dots,d_4$ be the diagonal entries of $r(h_0)$. If $E$ contains the $d_i$ and elements $\sum_id_i^by_i$ for $0\le b\le3$, then it contains the $y_i$, as the Vandermonde matrix of the $d_i$ is invertible. Let $E$ be as in (1). For $u,g,w\in H$ and indices $i,j$:
\begin{enumerate}
\item[(a)] $r(u)_{ii}\in E$, as $\operatorname{tr}r(h_0^bu)=\sum_id_i^br(u)_{ii}$;
\item[(b)] $r(u)_{1i}r(w)_{i1}\in E$, as $r(uh_0^bw)_{11}=\sum_id_i^br(u)_{1i}r(w)_{i1}$ lies in $E$ by (a);
\item[(c)] $r(u)_{1i}r(g)_{ij}r(w)_{j1}\in E$, as $\sum_jd_j^b\,r(u)_{1i}r(g)_{ij}r(w)_{j1}=r(u)_{1i}\,r(gh_0^bw)_{i1}$ lies in $E$ by (b).
\end{enumerate}
As the $r(h)$ span $M_4(K)$, for all $i,j$ some $g_{ij}\in H$ has $r(g_{ij})_{ij}\ne0$. Put $s=r(g_{12})_{12}$ and $t=r(g_{13})_{13}$. Then $x_j=r(g_{1j})_{1j}\in E$ for every $j$: for $j=1$ by (a), for $j=2,3$ by the choice of $s$ and $t$, and for $j=4$ as follows. Write $g_2=g_{13}$. As $r(g_2)$ is symplectic, $r(g_2)^*=\nu(g_2)r(g_2^{-1})$, and comparing the entries $(2,4)$ gives $r(g_2^{-1})_{24}=-t/\nu(g_2)\in E$, which is nonzero. By (b) and (c), $x_4\,r(g_{41})_{41}$ and $s\,r(g_2^{-1})_{24}\,r(g_{41})_{41}$ lie in $E$, so their quotient times $s\,r(g_2^{-1})_{24}$, which is $x_4$, does. Then $r(g_{j1})_{j1}=x_jr(g_{j1})_{j1}/x_j\in E$ by (b). Finally, for every $g$, $r(g)_{ij}=r(g_{1i})_{1i}r(g)_{ij}r(g_{j1})_{j1}/(x_ir(g_{j1})_{j1})\in E$ by (c).

(2) As $n^*n=1$, $n^*=n^{-1}$; comparing the entries $(2,1)$ and $(3,1)$ of $n^*$ and $n^{-1}$ gives $n_{43}=-n_{21}$ and $n_{42}=n_{31}-n_{21}n_{32}$.
\end{proof}
% fibTorus (FibTTrace.lean, 120 lines of its 454-line directory): mem_subfield_of_vandermonde_sums_mem, diag_entries_mem_of_trace_mem_of_regular_diag (a), entry_products_mem_of_trace_mem_of_regular_diag (b), triple_entry_products_mem_of_trace_regular_diag (c), exists_entry_ne_zero_of_irr and neg_entry_zero_two_eq_nu_mul_inv_entry_of_isGSp; the Lean indexes from 0, so its g2 = gs 0 2 is g_{13}. (2) is fib_low in FibLow.lean, with the entries ![1,2,3,2], ![0,0,0,1].

\begin{lemma}\label{sha-siegel-step-dimensions-generators}
Let $\mathcal O'$ be a discrete valuation ring with residue field $k'$, and $R$ a complete local noetherian $\mathcal O'$-algebra with residue field $k'$, the map $\mathcal O'\to R$ being local. Let $t_R$ be the $k'$-vector space of $\mathcal O'$-linear maps $\delta:R\to k'$ with $\delta(xy)=\bar x\delta(y)+\bar y\delta(x)$, where bars denote residues, and let $n\ge0$ with $\dim t_R\le n$. Then $\dim R\le n+1$.
\end{lemma}
\leanref{Development/AutomorphicArithmetic/AutomorphicVarieties/WeakShimuraTheory/RedDim.lean}{UpperTheory.V.Ws.redDim}
\begin{proof}
Write $\mathfrak m=\mathfrak m_R$.
\begin{enumerate}
\item[(a)] \emph{Elements.} For $y\in\mathfrak m$, the map $\delta\mapsto\delta(y)$ is a linear form on $t_R$. These forms span a space of dimension at most $n$, so there are $y_1,\dots,y_m\in\mathfrak m$, with $m\le n$, whose forms span it. So every $\delta\in t_R$ that vanishes at the $y_i$ vanishes on $\mathfrak m$.
\item[(b)] \emph{Generators modulo $\mathfrak m^2$.} Let $\varpi$ be a uniformizer of $\mathcal O'$, $\mathfrak a$ the ideal of $R$ generated by $\varpi$ and the $y_i$, which lies in $\mathfrak m$, and $I=\mathfrak a+\mathfrak m^2$. Suppose that some $x\in\mathfrak m$ is not in $I$, and let $N$ be the $\mathcal O'$-submodule $I+\mathcal O'\cdot1$ of $R$. Then $x\notin N$: if $x=y+a$ with $y\in I$ and $a\in\mathcal O'$, then $a$ maps into $\mathfrak m$, so $a\in\mathfrak m_{\mathcal O'}=\varpi\mathcal O'$ and $a\in I$, and $x\in I$. The quotient $R/N$ is killed by $\mathfrak m_{\mathcal O'}$, as $\varpi R\subseteq I$, so it is a $k'$-vector space, and some linear form on it does not vanish at the class of $x$. Composed with $R\to R/N$, it gives an $\mathcal O'$-linear map $\delta:R\to k'$ that vanishes on $N$, with $\delta(x)\ne0$. Then $\delta\in t_R$: for $y,z\in R$ choose $a,b\in\mathcal O'$ with $y-a$ and $z-b$ in $\mathfrak m$; then $yz=(y-a)(z-b)+az+by-ab$, where $(y-a)(z-b)\in\mathfrak m^2\subseteq N$ and $ab\in N$, so $\delta(yz)=\bar a\delta(z)+\bar b\delta(y)$. As $\delta$ vanishes at the $y_i$ but not at $x\in\mathfrak m$, this contradicts (a). So $\mathfrak m\subseteq\mathfrak a+\mathfrak m^2$.
\item[(c)] \emph{Conclusion.} By (b) and Nakayama's lemma, $\mathfrak m=\mathfrak a$, which is generated by at most $n+1$ elements. By Krull's height theorem \cite[Theorem~13.5]{Mat86}, $\dim R=\operatorname{ht}\mathfrak m\le n+1$.
\end{enumerate}
\end{proof}
% redDim (RedDim.lean, 123 lines): (a) is the private redDim_set, applied to the maps ev; (b) is hle with redDim_der (the map D through the quotient by N, with redDim_dual for the linear form); (c) is heq through Submodule.le_of_le_smul_of_le_jacobson_bot (Nakayama) and Mathlib's ringKrullDim_le_spanFinrank_maximalIdeal, which is the case I = m of Krull's height theorem.

\begin{lemma}\label{sha-siegel-step-dimensions}
Let $D$ be a Siegel datum, let $T$ be a fit tower for $D$ (Definition~\ref{sha-siegel-tower-towers}), and let $\mathcal O'\supseteq\mathcal O$ be a ring of Definition~\ref{sha-character-parts}, with residue field $k'$.
Let $W_1,\dots,W_9$ be the sets of weights
$\{\pm(-1,1),(-1,0),(-1,-1),(0,-1)\}$,
$\{\pm(-1,1),(1,0),(1,1),(0,1)\}$,
$\{\pm(1,1),(0,1),(-1,1),(-1,0)\}$,
$\{\pm(1,1),(0,-1),(1,-1),(1,0)\}$,
$\{\pm(1,0),(-1,-1),(0,-1),(1,-1)\}$,
$\{\pm(1,0),(1,1),(0,1),(-1,1)\}$,
$\{\pm(0,1),(-1,1),(-1,0),(-1,-1)\}$,
$\{\pm(0,1),(1,-1),(1,0),(1,1)\}$ and
$\{\pm(-1,1),\pm(1,1)\}$,
and for a homomorphism $r$ from $G_F$ to $M_4(R)$ let $I_W(r)\subseteq R$ be the ideal generated by the off-diagonal entries $r(g)_{ij}$ of weight outside $W$.
\begin{enumerate}
\item Let $d,C\ge0$ be such that $T$ satisfies the cohomology bounds $(d,C)$ along $\mathcal O\to k'$ (Definition~\ref{sha-siegel-tower-ideal}), let $c$ be a residual ample flagged pair for the problem of $T$ along $\mathcal O\to k'$, and let $P=(r_P,(n_v))$ be a universal lift of $c$ to $R$ (Definition~\ref{sha-siegel-tower-pairs}). Then $\dim R/I_{W_s}(r_P)\le d+20C+1$ for $s=1,\dots,9$.
\item Let $X=(A,B,P_X)$ be a deformation space for $T$ over $\mathcal O'$ (Definition~\ref{sha-siegel-tower-space}), let $c$ be residual, and let $f:A\to R$ be a homomorphism of $\mathcal O'$-algebras that carries $P_X$ to a universal lift of $c$. Let $\mathfrak q$ be a prime of $R$ such that the images of the $f(r_X(g))$ span $M_4$ over the fraction field of $R/\mathfrak q$, and let $\mathfrak q_B$ be its inverse image in $B$. Then $\dim R/\mathfrak q\le\dim B/\mathfrak q_B+2+4|\Sigma_T|$.
\item Suppose that $T$ is deep, let $\iota:\bar{\mathbf Q}_\ell\cong\mathbf C$, let $X$ be a deformation space for $T$ over $\mathcal O'$ with a modularity ideal $J$ (Definition~\ref{sha-siegel-tower-ideal}), and let $d_v\ge0$ with $d_v|\Sigma_T|=d_T$. Let $v\in\Sigma_T$, $i\ne j$, and let $\mathfrak q$ be a prime of $B$ that contains $J$, $\mathfrak m_{\mathcal O'}B$ and the elements $\psi^B_{v,i}(s)-\psi^B_{v,j}(s)$ with $s\in U_v$. Then $\dim B/\mathfrak q+d_v\le2d_T+1$. The proof of (3) does not use that $T$ is deep.
\end{enumerate}
\end{lemma}
\leanref{Development/AutomorphicArithmetic/AutomorphicVarieties/WeakShimuraTheory/CompactIntegralFields/HomOfDer.lean}{UpperTheory.V.Ws.red}
\leanref{Development/AutomorphicArithmetic/AutomorphicVarieties/WeakShimuraTheory/CompactIntegralFields/FibLow.lean}{UpperTheory.V.Ws.fib}
\leanref{Development/AutomorphicArithmetic/AutomorphicVarieties/ModularityArgument/ThresholdArguments/DenseThresholdCases/OffDiagonalNilpotentRankFour.lean}{UpperTheory.V.Ws.off_diagonal_nilpotent_rank_four}
\begin{proof}
(1) Let $s$ be given, $I=I_{W_s}(r_P)$, and $c=(r_c,(n_{c,v}))$. If $I=R$ there is nothing to prove. Otherwise $I\subseteq\mathfrak m_R$, and $R/I$ is a complete local noetherian $\mathcal O'$-algebra with residue field $k'$. Let $t$ be the $k'$-space of $\mathcal O'$-linear maps $\delta:R/I\to k'$ with $\delta(xy)=\bar x\delta(y)+\bar y\delta(x)$; for $y\in R$ write $\delta(y)$ for $\delta$ of its class, and apply $\delta$ to matrices entrywise.
\begin{enumerate}
\item[(a)] \emph{Lifts to the dual numbers.} For $\delta\in t$, $y\mapsto\bar y+\varepsilon\delta(y)$ is a homomorphism of $\mathcal O'$-algebras $R\to k'[\varepsilon]$; it carries $P$ to a lift $P_\delta$ of $c$ to $k'[\varepsilon]$, whose entries are locally constant, as $\mathfrak m_{k'[\varepsilon]}^2=0$. The entries of $r_P$ of weight outside $W_s$ lie in $I\subseteq\mathfrak m_R$, so $r_c$ is supported in $W_s$, and so is the homomorphism of $P_\delta$.
\item[(b)] \emph{Tangent pairs.} Put $x_\delta(g)=\delta(r_P(g))r_c(g)^{-1}$ and $Y_{\delta,v}=\delta(n_{P,v})n_{c,v}^*$. The homomorphism of $P_\delta$ is $g\mapsto(1+\varepsilon x_\delta(g))r_c(g)$, and its flags are the $(1+\varepsilon Y_{\delta,v})n_{c,v}$. As in step (b) of the proof of Lemma~\ref{sha-siegel-step-completion-tangent}, with $Y_v$ replaced by $-Y_{\delta,v}$, the conditions on $P_\delta$ say that $(x_\delta,(Y_{\delta,v}))$ lies in the set $\mathrm{Tan}_{W_s}(c)$ of Lemma~\ref{sha-siegel-step-dimensions-tangent}: the multiplicativity gives the cocycle condition, the multipliers give $x_\delta(g),Y_{\delta,v}\in\mathfrak{sp}_4(k')$, the value $1$ on $N_\Phi$ gives $x_\delta=0$ there, the diagonal matrix at $g_0^T$ gives a diagonal $x_\delta(g_0^T)$, $x_\delta(g)r_c(g)=\delta(r_P(g))$ is supported in $W_s$, $n_{c,v}^*Y_{\delta,v}n_{c,v}=n_{c,v}^*\delta(n_{P,v})$ is strictly lower triangular, and the triangular forms give the last condition.
\item[(c)] \emph{Injectivity.} The map $\delta\mapsto(x_\delta,(Y_{\delta,v}))$ is $k'$-linear. If it vanishes, $\delta$ vanishes on the entries of $P$, so the homomorphisms $R\to k'[\varepsilon]$ attached to $\delta$ and to $0$ carry $P$ to the same lift; by the universal property they are equal, and $\delta=0$.
\item[(d)] \emph{The bound.} By (b), (c) and Lemma~\ref{sha-siegel-step-dimensions-tangent}, for $a:\mathcal O\to k'$, $\dim t\le d+20C$. As $\mathcal O'$ is a discrete valuation ring and $R/I$ is a complete local noetherian $\mathcal O'$-algebra with residue field $k'$, Lemma~\ref{sha-siegel-step-dimensions-generators} gives $\dim R/I\le d+20C+1$.
\end{enumerate}

(2) Let $\mathfrak x=f^{-1}(\mathfrak m_R)$ and $\mathfrak p=f^{-1}(\mathfrak q)$, primes of $A$, so that $\mathfrak q_B$ is the inverse image of $\mathfrak p$ in $B$.
\begin{enumerate}
\item[(a)] \emph{The completion.} The composite of $f$ with the residue map of $R$ is a point of $X$, with image $c$, so $\mathfrak x$ is maximal. Let $\hat A$ be the $\mathfrak x$-adic completion of $A$. By the proof of Lemma~\ref{sha-siegel-step-completion}(1), $\hat A$ carries a universal lift of $c$, every homomorphism of $\mathcal O'$-algebras $g:A\to R'$ to a complete local noetherian $\mathcal O'$-algebra with residue field $k'$ whose composite with the residue map is that point factors uniquely through $\hat A$, and two homomorphisms from $A$ that carry $P_X$ to the same pair are equal. As $f$ carries $P_X$ to a universal lift, the same factorization property holds for $f$. So $R$ and $\hat A$ are isomorphic as $\mathcal O'$-algebras under $A$. As $\hat A$ is also the completion of the local ring $A_{\mathfrak x}$, $R$ is flat over $A_{\mathfrak x}$ \cite[Theorem~8.8]{Mat86}, and $\mathfrak xR=\mathfrak m_R$. By base change, $R/\mathfrak pR$ is flat over $(A/\mathfrak p)_{\mathfrak x}$, and the maximal ideal of $(A/\mathfrak p)_{\mathfrak x}$ generates that of $R/\mathfrak pR$, so the closed fibre is a field. By \cite[Theorem~15.1]{Mat86}, $\dim R/\mathfrak pR=\dim(A/\mathfrak p)_{\mathfrak x}$. As $\mathfrak pR\subseteq\mathfrak q$, $\dim R/\mathfrak q\le\dim(A/\mathfrak p)_{\mathfrak x}$.
\item[(b)] \emph{The dimension formula.} The ring $B/\mathfrak q_B$ is a complete local domain, so it is a quotient of a regular local ring \cite[Theorem~29.4(ii)]{Mat86}, hence of a Cohen-Macaulay ring, and it is universally catenary \cite[Theorem~17.9]{Mat86}; it embeds in the domain $A/\mathfrak p$, which is of finite type over it; $\mathfrak x\cap B=\mathfrak m_B$; and the residue fields of $\mathfrak x$ and of $\mathfrak m_B$ are both $k'$. By the dimension formula \cite[Theorem~15.6]{Mat86}, $\dim(A/\mathfrak p)_{\mathfrak x}=\dim B/\mathfrak q_B+\tau$, where $\tau$ is the transcendence degree of the fraction field of $A/\mathfrak p$ over that of $B/\mathfrak q_B$.
\item[(c)] \emph{Transcendence degree.} Let $\Omega_{\mathfrak q}$ be the fraction field of $R/\mathfrak q$ and $y:A\to\Omega_{\mathfrak q}$ the composite map, with kernel $\mathfrak p$. Lemma~\ref{sha-siegel-step-dimensions-torus}(1) applies to $g\mapsto y(r_X(g))$, $\nu(g)=y(\theta(g))$ and $h_0=g_0^T$: the diagonal entries of $r_X(g_0^T)$ have the pairwise distinct residues $\bar\chi_i(g_0^T)$, by (l), so their differences are units of $R$ and their images in $\Omega_{\mathfrak q}$ are distinct; and the $y(r_X(g))$ span $M_4(\Omega_{\mathfrak q})$ by hypothesis. Let $s,t$ be as there, and let $T_0$ consist of $s$, $t$ and the entries $(2,1)$, $(3,1)$, $(4,1)$, $(3,2)$ of the $y(n_{X,v})$, at most $2+4|\Sigma_T|$ elements. Let $E$ be the set of elements of $\Omega_{\mathfrak q}$ algebraic over the subfield generated by $y(B)$ and $T_0$; it is a subfield. It contains the images of the multipliers and the traces, which are coefficients of the $c_g$, and the diagonal entries of $y(r_X(g_0^T))$, which are roots of $y(c_{g_0^T})$. By Lemma~\ref{sha-siegel-step-dimensions-torus}, $E$ contains the entries of the $y(r_X(g))$ and of the $y(n_{X,v})$, which generate $y(A)$ over $y(B)$. So the fraction field of $A/\mathfrak p$ is algebraic over that of $B/\mathfrak q_B$ adjoined $T_0$, and $\tau\le2+4|\Sigma_T|$.
\end{enumerate}
By (a), (b) and (c), $\dim R/\mathfrak q\le\dim B/\mathfrak q_B+2+4|\Sigma_T|$.
% (1) is red (HomOfDer.lean, 253 lines): (a) is hom_of_der and lifts_map; (b) is the map T with redDual (RedDual.lean, 148 lines) for the image in the span of Tan; (c) is hT through hDef.2 (uniqueness); (d) is redTan (Lemma sha-siegel-step-dimensions-tangent) and redDim (Lemma sha-siegel-step-dimensions-generators).
% (2) is fib (FibLow.lean, 110 lines): (a) is fibCpl (FibreTraces/ComplOfNil.lean, the completion property from the universality of f, with X.hom_ext) and CA.c3c (KStableHighestWeightVector.lean: cpl_uniq for R as the completion Cpl x of A_x, flat over it by Mathlib's instance for adic completions, and dim_quot_eq_of_flat, the base change to A_x / p with dim_eq_of_flat, which is Matsumura's Theorem 15.1 (ii) through Ideal.height_eq_height_add_of_liesOver_of_hasGoingDown); (b) is CA.c3b (HighestWeightIntegral.lean: cl_quot and cl_cmquot, B / q_B as a quotient of a power series ring over O', which is Cohen-Macaulay, then g_dimform, the Lean's own proof of the dimension formula over a quotient of a Cohen-Macaulay local ring); (c) is fib_gen with fibTorus and fib_low (Lemma sha-siegel-step-dimensions-torus) and fibTrdeg.

(3) Let $v$, $i\ne j$ and $\mathfrak q$ be as there. As $F$ is Galois over $\mathbf Q$, its places above $\ell$ all have the same local degree; they correspond to the elements of $\Sigma_T$, and their local degrees add up to $d_T$, so each is $d_v$.
\begin{enumerate}
\item[(a)] \emph{The characters.} For $u\in\Sigma_T$ and $s\in D_u(F)$, the matrix $n_{X,u}^*r_X(s)n_{X,u}$ is upper triangular, with diagonal entries $\psi_{u,1}(s),\dots,\psi_{u,4}(s)$, and symplectic with multiplier $\theta(s)$. So the $\psi_{u,m}$ are characters of $D_u(F)$ with $\psi_{u,m}\psi_{u,5-m}=\theta$, and since $B\to A$ is injective, the same holds for the $\psi^B_{u,m}$. For $s\in I_u(F)$, $\bar r_0(s)=1$ by (k), so $\bar\theta(s)=1$, and the residue $\bar\chi_m(s)$ of $\psi^B_{u,m}(s)$ is $1$.
\item[(b)] \emph{The elements.} By Lemma~\ref{sha-siegel-step-dimensions-pinning}, for each $u\in\Sigma_T$ there are $s_{u,1},\dots,s_{u,d_v}\in I_u(F)$ that form a pinning set, and, as $\ell$ is odd, there is $s_0\in I_v(F)$ that forms a pinning set with the $s\in I_v(F)$ with $\chi_{\mathrm{cyc}}(s)=1$. Choose $a$ with $\{a,5-a,i,5-i\}=\{1,2,3,4\}$, and with $a=j$ unless $j=5-i$. Consider the elements $\psi^B_{u,a}(s_{u,k})-1$ for $u\in\Sigma_T$ and $1\le k\le d_v$, the elements $\psi^B_{u,i}(s_{u,k})-1$ for $u\ne v$ and $1\le k\le d_v$, and $\psi^B_{v,i}(s_0)-\psi^B_{v,j}(s_0)$. There are $d_T+(d_T-d_v)+1$ of them, and they lie in $\mathfrak m_B$ by (a).
\item[(c)] \emph{The primes.} Let $\mathfrak p\supseteq\mathfrak q$ be a prime of $B$ that contains the elements of (b), $R=B/\mathfrak p$, a noetherian local domain, and $f:B\to R$ the quotient map. As $\mathfrak m_{\mathcal O'}B\subseteq\mathfrak p$, $\ell=0$ in $R$, and $f(\theta(s))=1$ when $\bar\theta(s)=1$. Let $\kappa_{u,m}=f\circ\psi^B_{u,m}:D_u(F)\to R^\times$, a character, continuous for the $\mathfrak m_R$-adic topology since $\psi^B_{u,m}$ is continuous and $f$ maps $\mathfrak m_B$ into $\mathfrak m_R$. By (a), $\kappa_{u,m}(s)\equiv1$ and $\kappa_{u,m}(s)\kappa_{u,5-m}(s)=1$ for $s\in I_u(F)$. By the pinning sets of (b), $\kappa_{u,a}$ is trivial on $I_u(F)$ for every $u$, and so is $\kappa_{u,i}$ for $u\ne v$. The character $\kappa_{v,i}\kappa_{v,j}^{-1}$ is trivial at $s_0$, and at every $s\in I_v(F)$ with $\chi_{\mathrm{cyc}}(s)=1$, since such an $s$ lies in $U_v$ by (k) and $\mathfrak q\subseteq\mathfrak p$; so $\kappa_{v,i}=\kappa_{v,j}$ on $I_v(F)$. If $a=j$, then $\kappa_{v,i}$ is trivial on $I_v(F)$. Otherwise $j=5-i$, so $\kappa_{v,i}^2=\kappa_{v,i}\kappa_{v,5-i}=1$ on $I_v(F)$, and $\kappa_{v,i}$ is trivial there by the first paragraph of the proof of Lemma~\ref{sha-siegel-step-dimensions-pinning}. As $\{a,5-a,i,5-i\}=\{1,2,3,4\}$ and $\kappa_{u,5-m}=\kappa_{u,m}^{-1}$ on $I_u(F)$, every $\kappa_{u,m}$ is trivial on $I_u(F)$. So $\mathfrak p$ contains the $\psi^B_{u,m}(s)-1$ with $s\in I_u(F)$, as well as $J$ and $\mathfrak m_{\mathcal O'}B$. By the finiteness property of $J$, $\mathfrak p$ contains a power of $\mathfrak m_B$, hence $\mathfrak m_B$.
\item[(d)] \emph{The bound.} By (c), the maximal ideal of $B/\mathfrak q$ is minimal over the ideal generated by the images of the elements of (b). By Krull's height theorem \cite[Theorem~13.5]{Mat86}, $\dim B/\mathfrak q\le2d_T+1-d_v$, which is (3).
\end{enumerate}
% (3) is off_diagonal_nilpotent_rank_four (OffDiagonalNilpotentRankFour.lean, about 90 lines): the local degree is exists_place_over_l_locDeg_mul_card_eq; (a) is deformationSpace_local_characters_continuous with pseudoDeformation_diag_characters_mul_rev_eq_similitude and Ok.ineL; (b) is exists_pinning_elements_of_locAb and exists_mem_pins_insert_ker_of_padic_character, with the index a chosen by decide; (c) is hpin, hA, hE (Pins.eq), hI with eq_one_of_pow, hall, and hPM.fin; (d) is krullDim_quotient_le_card_of_radical_parameters. With its dependencies in BuildingsResidueChars, about 1,000 lines, most of them for local class field theory.
\end{proof}

\begin{lemma}\label{sha-siegel-step-irreducible}
Let $W_1,\dots,W_9$ and $I_W$ be as in Lemma~\ref{sha-siegel-step-dimensions}.
\begin{enumerate}
\item Let $K$ be a field, $H$ a group, and $r:H\to M_4(K)$ a homomorphism such that each $r(h)$ is symplectic with some multiplier (Definition~\ref{sha-siegel-datum}), and $r(h_0)$ is diagonal with pairwise distinct diagonal entries for some $h_0\in H$. Then the $r(h)$ span $M_4(K)$ if and only if $r$ is supported in none of $W_1,\dots,W_9$.
\item Let $D$ be a Siegel datum, $T$ a fit tower for $D$ (Definition~\ref{sha-siegel-tower-towers}), $\mathcal O'\supseteq\mathcal O$ a ring of Definition~\ref{sha-character-parts} with residue field $k'$, $c$ a residual flagged pair for the problem of $T$ along $\mathcal O\to k'$, and $P=(r_P,(n_v))$ a lift of $c$ to $R$ (Definition~\ref{sha-siegel-tower-pairs}). Let $\varphi:R\to K$ be a ring homomorphism to a field and $\mathfrak q$ a prime of $R$ containing $\ker\varphi$ but none of the $I_{W_s}(r_P)$. Then the $\varphi(r_P(g))$ span $M_4(K)$.
\end{enumerate}
\end{lemma}
\leanref{Development/AutomorphicArithmetic/AutomorphicVarieties/WeakShimuraTheory/CompactIntegralFields/RedList.lean}{UpperTheory.V.Ws.redList}
\leanref{Development/AutomorphicArithmetic/AutomorphicVarieties/WeakShimuraTheory/ValuationCharacters/ConstOfFinite.lean}{UpperTheory.V.Ws.Ch.irr}
\begin{proof}
(1) Since $r(h_0)$ is diagonal with pairwise distinct diagonal entries, the matrix units $E_{ii}$ are polynomials in $r(h_0)$, so they lie in the span of the $r(h)$, and Lemma~\ref{sha-flagged-relation} applies: with its relation $\succeq$, the $r(h)$ span $M_4(K)$ if and only if $i\succeq j$ for all $i,j$. A finite check over the reflexive and transitive relations with $i\succeq j\Rightarrow5-j\succeq5-i$ shows that if this fails, there is $s$ such that $i\not\succeq j$ for every entry $(i,j)$ of weight outside $W_s$, that is, $r$ is supported in $W_s$. Conversely, each $W_s$ omits the weight of some off-diagonal entry, so a representation supported in $W_s$ has $i\not\succeq j$ for that entry.
(2) The matrix $r_P(g_0^T)$ is diagonal, and its diagonal entries reduce to the $\bar\chi_i(g_0^T)$, which are pairwise distinct by (l); so their differences are units of $R$, and their images under $\varphi$ are pairwise distinct. Each $\varphi(r_P(g))$ is symplectic with multiplier $\varphi(\theta(g))$. If the $\varphi(r_P(g))$ did not span $M_4(K)$, then by (1) $\varphi\circ r_P$ would be supported in some $W_s$, so $\varphi$ would kill the generators of $I_{W_s}(r_P)$, and $I_{W_s}(r_P)\subseteq\ker\varphi\subseteq\mathfrak q$, a contradiction.
\end{proof}

\begin{lemma}\label{sha-siegel-step-commutative}
\begin{enumerate}
\item Let $f:A\to R$ be a flat ring homomorphism with $R$ local, and let $\mathfrak q$ be a minimal prime of $A$ contained in $f^{-1}(\mathfrak m_R)$. Then $\mathfrak q=f^{-1}(\mathfrak Q)$ for a minimal prime $\mathfrak Q$ of $R$.
\item Let $\mathcal O$ be a complete discrete valuation ring, let $a,b,N\ge0$ with $b+N\le a$, and let $R\ne0$ be the quotient of $\mathcal O[[Z_1,\dots,Z_a]]$ by an ideal generated by $b$ elements. Any two minimal primes $\mathfrak P$ and $\mathfrak Q$ of $R$ are joined by a sequence $\mathfrak P=\mathfrak P_0,\mathfrak P_1,\dots,\mathfrak P_m=\mathfrak Q$ of minimal primes with $\dim R/(\mathfrak P_{i-1}+\mathfrak P_i)\ge N$ for $1\le i\le m$.
\item Let $R$ be a noetherian local ring, $\varpi\in\mathfrak m_R$, $N\ge1$, $I\subseteq R$ an ideal with $\dim R/I\ge N+1$, $\mathfrak a_1,\dots,\mathfrak a_s$ ideals with $\dim R/\mathfrak a_i<N$, and $\mathfrak b_1,\dots,\mathfrak b_t$ ideals such that $\dim R/\mathfrak q<N$ for every prime $\mathfrak q$ that contains some $\mathfrak b_k$ and none of the $\mathfrak a_i$. Then there is a prime $\mathfrak p\supseteq I$ with $\dim R/\mathfrak p=1$ and $\varpi\in\mathfrak p$ that contains none of the $\mathfrak a_i$ and none of the $\mathfrak b_k$.
\item Let $\mathcal O$ be a local ring whose residue field is finite of characteristic $\ell$, let $R$ be a complete local noetherian $\mathcal O$-algebra with the same residue field, and let $\mathfrak p$ be a prime of $R$ with $\dim R/\mathfrak p=1$ that contains the image of $\mathfrak m_{\mathcal O}$. There are a finite field $\kappa$, an $\mathcal O$-algebra structure on $\kappa$, and a homomorphism of $\mathcal O$-algebras $j:R\to\kappa[[\tau]]$ with kernel $\mathfrak p$ that maps $\mathfrak m_R$ into $\tau\kappa[[\tau]]$.
\end{enumerate}
\end{lemma}
\leanref{Development/AutomorphicArithmetic/AutomorphicVarieties/WeakShimuraTheory/ValuationCharacters/ConstOfFinite.lean}{UpperTheory.V.Ws.Ch.flat_min}
\leanref{Development/AutomorphicArithmetic/AutomorphicVarieties/WeakShimuraTheory/ValuationCharacters/ConstOfFinite.lean}{UpperTheory.V.Ws.Ch.chain}
\leanref{Development/AutomorphicArithmetic/AutomorphicVarieties/WeakShimuraTheory/ValuationCharacters/ConstOfFinite.lean}{UpperTheory.V.Ws.Ch.nice}
\leanref{Development/AutomorphicArithmetic/AutomorphicVarieties/CurveDVRSpecialization/NicePrimeEx.lean}{UpperTheory.V.Pr.nicePrimeEx}
\leanref{Development/AutomorphicArithmetic/AutomorphicVarieties/WeakShimuraTheory/ValuationCharacters/ConstOfFinite.lean}{UpperTheory.V.Ws.Ch.curve}
\begin{proof}
(1) Flat homomorphisms have the going-down property \cite[Theorem~9.5]{Mat86}. Applied to $\mathfrak m_R$, which lies over $f^{-1}(\mathfrak m_R)$, and to $\mathfrak q\subseteq f^{-1}(\mathfrak m_R)$, it gives a prime of $R$ lying over $\mathfrak q$; a minimal prime $\mathfrak Q$ of $R$ inside it has $f^{-1}(\mathfrak Q)\subseteq\mathfrak q$, hence $f^{-1}(\mathfrak Q)=\mathfrak q$ by the minimality of $\mathfrak q$.

(2) The ring $\mathcal O[[Z_1,\dots,Z_a]]$ is a complete regular local domain of dimension $a+1$. By Grothendieck's connectedness theorem \cite[Expos\'e~XIII]{Gro68}, the spectrum of its quotient $R$ by an ideal with $b$ generators stays connected after removing any closed subset of dimension less than $a-b\ge N$; for the finitely many minimal primes of $R$ this is the existence of the sequence.

(3) Choose a chain of primes of length $N+1$ whose first prime contains $I$. Ending it at $\mathfrak m_R$ if necessary, its last prime contains $\varpi$, and by Krull's principal ideal theorem it may be changed into a chain $\mathfrak q_0\subsetneq\mathfrak q_1\subsetneq\mathfrak q_2\subsetneq\cdots$ of the same length with the same first prime and with $\varpi\in\mathfrak q_1$. The chain from $\mathfrak q_1$ has length at least $N$, so $\dim R/\mathfrak q_1\ge N$; hence $\mathfrak q_1$ contains no $\mathfrak a_i$, and no minimal prime over a $\mathfrak b_k$ that contains no $\mathfrak a_i$, since these have $\dim R/\mathfrak q<N$. Choose $g\in\mathfrak q_2\setminus\mathfrak q_1$, elements $x_i\in\mathfrak a_i\setminus\mathfrak q_1$, and elements $x'_j\in\mathfrak q'_j\setminus\mathfrak q_1$ for the finitely many minimal primes $\mathfrak q'_j$ over the $\mathfrak b_k$ that contain no $\mathfrak a_i$; their product $f$ lies in $\mathfrak m_R$ and not in $\mathfrak q_1$. Let $\mathfrak p$ be maximal among the primes that contain $\mathfrak q_1$ and not $f$. Every prime strictly containing $\mathfrak p$ contains $f$, so the maximal ideal of $R/\mathfrak p$ is minimal over the image of $f$, which is nonzero, and $\dim R/\mathfrak p=1$ by Krull's principal ideal theorem. The prime $\mathfrak p$ contains $I$ and $\varpi$. It contains no $\mathfrak a_i$, since $x_i\notin\mathfrak p$. If it contained some $\mathfrak b_k$, it would contain a minimal prime over $\mathfrak b_k$, which either contains some $\mathfrak a_i$ or is one of the $\mathfrak q'_j$; both are excluded.

(4) The ring $R/\mathfrak p$ is a complete local domain of dimension $1$ with finite residue field, and it has characteristic $\ell$, since $\ell\in\mathfrak m_{\mathcal O}$. Its normalization is finite over it and is a complete discrete valuation ring of characteristic $\ell$ with finite residue field $\kappa$, so it is $\kappa[[\tau]]$ by Cohen's structure theorem \cite{Mat86}; this gives an injective local homomorphism $R/\mathfrak p\to\kappa[[\tau]]$, and $j$ is its composite with $R\to R/\mathfrak p$. If $\ell^e$ is the order of the residue field of $\mathcal O$, every $a\in\mathcal O$ has $a^{\ell^e}-a\in\mathfrak m_{\mathcal O}$, so the image $y$ of $a$ in $\kappa[[\tau]]$ satisfies $y^{\ell^e}=y$. A power series $y$ with $y^{\ell^e}=y$ is constant: otherwise $y-y(0)$ would have some order $n\ge1$, and its $\ell^e$-th power, which is $y-y(0)$, would have order $\ell^en$. So $\mathcal O\to\kappa[[\tau]]$ takes values in $\kappa$, which makes $\kappa$ an $\mathcal O$-algebra and $j$ a homomorphism of $\mathcal O$-algebras. It maps $\mathfrak m_R$ into $\tau\kappa[[\tau]]$ because it is local.
\end{proof}

\begin{lemma}\label{sha-siegel-step-propagation}
Let $D$ be a Siegel datum, let $T$ be a fit deep tower for $D$ (Definition~\ref{sha-siegel-tower-towers}), let $\iota:\bar{\mathbf Q}_\ell\cong\mathbf C$, let $\mathcal O'\supseteq\mathcal O$ be a ring of Definition~\ref{sha-character-parts}, with residue field $k'$, and let $X=(A,B,P_X)$ be a deformation space for $T$ over $\mathcal O'$ with a modularity ideal $J$ (Definitions~\ref{sha-siegel-tower-space} and~\ref{sha-siegel-tower-ideal}).
Let $C\ge0$ be such that $T$ satisfies the cohomology bounds $(d_T,C)$ along $\mathcal O\to k'$ and $30(C+1)^2\le d_T$.
Let $x$ be a point of $X$ whose image is residual and ample. If $J$ lies in some minimal prime below $\ker x$, then it lies in all minimal primes below $\ker x$.
\end{lemma}
\leanref{Development/AutomorphicArithmetic/AutomorphicVarieties/ModularityArgument/ThresholdArguments/DenseThresholdCases/ClassicalFixedByIwahori.lean}{UpperTheory.V.Ws.residual_charpoly_rank_four}
\leanref{Development/AutomorphicArithmetic/AutomorphicVarieties/WeakShimuraTheory/ValuationCharacters/ValLtIff.lean}{UpperTheory.V.Ws.chDv}
\begin{proof}
Write $d=d_T$, let $c$ be the image of $x$, and let $\varphi:B\to A$ be the structure map.
\begin{enumerate}
\item[(1)] By Lemma~\ref{sha-siegel-step-completion}(1) there are $R$ and a flat $f:A\to R$ carrying $P_X$ to a universal lift of $c$, with $x$ equal to $f$ followed by the residue map; so $\ker x=f^{-1}(\mathfrak m_R)$. By Lemma~\ref{sha-siegel-step-completion}(2), $R$ is the quotient of $\mathcal O'[[Z_1,\dots,Z_a]]$ by an ideal generated by $b$ elements, with $b+2+2d\le a+2|\Sigma_T|$.
\item[(2)] The field $F$ is Galois over $\mathbf Q$, so its primes above $\ell$ all have the same ramification index and residue degree, with product $d_v$ say, and their number times $d_v$ is $d$; these primes correspond to the elements of $\Sigma_T$, so $d=d_v|\Sigma_T|$. As $|\Sigma_T|\le C$ and $d\ge30(C+1)^2$, we get $d_v\ge30(C+1)$ and $d\ge30(C+1)$.
\item[(3)] Put $N=2d+1-2|\Sigma_T|$, so that $b+N+1\le a$. Let $\mathfrak q_0$ and $\mathfrak q$ be minimal primes of $A$ contained in $\ker x$, with $JA\subseteq\mathfrak q_0$. By Lemma~\ref{sha-siegel-step-commutative}(1), $\mathfrak q_0=f^{-1}(\mathfrak P)$ and $\mathfrak q=f^{-1}(\mathfrak Q)$ for minimal primes $\mathfrak P$ and $\mathfrak Q$ of $R$, and by Lemma~\ref{sha-siegel-step-commutative}(2) they are joined by a sequence of minimal primes in which consecutive primes $\mathfrak P'$ and $\mathfrak Q'$ have $\dim R/(\mathfrak P'+\mathfrak Q')\ge N+1$. So it suffices to show: if $\mathfrak P'$ and $\mathfrak Q'$ are such primes and $f(\varphi(J))\subseteq\mathfrak P'$, then $f(\varphi(J))\subseteq\mathfrak Q'$.
\item[(4)] Let $\varpi$ be a uniformizer of $\mathcal O'$. Consider the ideals $\mathfrak a_s=I_{W_s}(f\circ r_X)$ of Lemma~\ref{sha-siegel-step-dimensions}, $1\le s\le9$, and, for $v\in\Sigma_T$ and $i\ne j$, the ideal $\mathfrak b_{v,i,j}$ generated by $f(\varphi(J))$, $\varpi$ and the $f(\varphi(\psi^B_{v,i}(s)-\psi^B_{v,j}(s)))$ with $s\in U_v$. By Lemma~\ref{sha-siegel-step-dimensions}(1), $\dim R/\mathfrak a_s\le d+20C+1<N$, the last inequality because $d\ge30(C+1)$ and $|\Sigma_T|\le C$.
\item[(5)] Let $\mathfrak q'$ be a prime of $R$ that contains some $\mathfrak b_{v,i,j}$ and no $\mathfrak a_s$. By Lemma~\ref{sha-siegel-step-irreducible}(2), applied to $R\to R/\mathfrak q'$ followed by the inclusion into its fraction field, the images of the $f(r_X(g))$ span $M_4$ over the fraction field of $R/\mathfrak q'$, so $\dim R/\mathfrak q'\le\dim B/\mathfrak q'_B+2+4|\Sigma_T|$ by Lemma~\ref{sha-siegel-step-dimensions}(2), where $\mathfrak q'_B$ is the inverse image of $\mathfrak q'$ in $B$. The prime $\mathfrak q'_B$ contains $J$, $\mathfrak m_{\mathcal O'}B$ and the $\psi^B_{v,i}(s)-\psi^B_{v,j}(s)$ with $s\in U_v$, so $\dim B/\mathfrak q'_B+d_v\le2d+1$ by Lemma~\ref{sha-siegel-step-dimensions}(3). Hence $\dim R/\mathfrak q'\le2d+3+4|\Sigma_T|-d_v<N$, because $d_v>2+6|\Sigma_T|$ by (2).
\item[(6)] By (4), (5) and Lemma~\ref{sha-siegel-step-commutative}(3), applied to $I=\mathfrak P'+\mathfrak Q'$, there is a prime $\mathfrak p\supseteq\mathfrak P'+\mathfrak Q'$ with $\dim R/\mathfrak p=1$ and $\varpi\in\mathfrak p$ that contains no $\mathfrak a_s$ and no $\mathfrak b_{v,i,j}$. Since $\varpi\in\mathfrak p$, $\mathfrak p$ contains the image of $\mathfrak m_{\mathcal O'}$, and Lemma~\ref{sha-siegel-step-commutative}(4) gives a finite field $\kappa$ and a homomorphism of $\mathcal O'$-algebras $j:R\to\kappa[[\tau]]$ with kernel $\mathfrak p$ that maps $\mathfrak m_R$ into $\tau\kappa[[\tau]]$.
\item[(7)] Let $y=j\circ f$. It maps $\mathfrak m_B$ into $\tau\kappa[[\tau]]$: the homomorphisms of $\mathcal O'$-algebras $B\to R\to k'$ and $B\to k'$ agree, as in the proof of Lemma~\ref{sha-siegel-step-linked}, so $f(\varphi(\mathfrak m_B))\subseteq\mathfrak m_R$. It kills $J$, since $f(\varphi(J))\subseteq\mathfrak P'\subseteq\mathfrak p$. For $v\in\Sigma_T$ and $i\ne j$, $\mathfrak p$ contains $f(\varphi(J))$ and $\varpi$ but not $\mathfrak b_{v,i,j}$, so some $s\in U_v$ has $y(\psi_{v,i}(s))\ne y(\psi_{v,j}(s))$, where the $\psi_{v,i}(s)$, the diagonal entries of $P_X$, are the images of the $\psi^B_{v,i}(s)$. By Lemma~\ref{sha-siegel-step-irreducible}(2), applied to $j$ followed by the inclusion into the fraction field of $\kappa[[\tau]]$, the $y(r_X(g))$ span $M_4$ over that field. By the patching property of $J$ (Definition~\ref{sha-siegel-tower-ideal}), $y$ detects $J$.
\item[(8)] The inverse image of $\mathfrak Q'$ in $B$ is a prime whose image in $A$ is killed by $y$, since $\mathfrak Q'\subseteq\mathfrak p=\ker j$. So it contains $J$, that is, $f(\varphi(J))\subseteq\mathfrak Q'$.
\end{enumerate}
\end{proof}

\begin{lemma}\label{sha-siegel-step-last}
Assume the Galois representations as in Theorem~\ref{sha-residually-reducible-automorphy}.
Let $D$ be a Siegel datum, let $T$ be a fit deep tower for $D$ (Definition~\ref{sha-siegel-tower-towers}), let $\iota:\bar{\mathbf Q}_\ell\cong\mathbf C$, and let $V$ be an ambient ring with residue field $k_V$ (Definition~\ref{sha-siegel-tower-ambient}).
Let $C\ge0$ be such that $30(C+1)^2\le d_T$ and $T$ satisfies the cohomology bounds $(d_T,C)$ along every homomorphism from $\mathcal O$ to a field that kills $\mathfrak m_{\mathcal O}$ (Definition~\ref{sha-siegel-tower-ideal}).
Let $Q$ be an integral seed for $T$ over a ring $\mathcal O_1\supseteq\mathcal O$ of Definition~\ref{sha-character-parts}, and let $P$ be a flagged pair for the problem of $T$ along $\mathcal O\to V$, with entries in a complete discrete valuation subring of $V$ with finite residue field through which its homomorphism factors continuously, whose image $c$ over $k_V$ is residual and joined by a chain to the image of $Q$ under $\mathcal O_1\to V\to k_V$ (Definition~\ref{sha-siegel-tower-pairs}).
Then there are a ring $\mathcal O'\supseteq\mathcal O$ of Definition~\ref{sha-character-parts}, a deformation space $X=(A,B,P_X)$ for $T$ over $\mathcal O'$ (Definition~\ref{sha-siegel-tower-space}), a modularity ideal $J$ for $X$ and a ring homomorphism $y:A\to V$ that restricts to the given map on $\mathcal O'$, maps $\mathfrak m_B$ into $\mathfrak m_V$, kills $J$, and carries $P_X$ to $P$.
\end{lemma}
\leanref{Development/AutomorphicArithmetic/AutomorphicVarieties/ModularityArgument/ThresholdArguments/DenseThresholdCases/Gsp4PromodularityCaseTwo.lean}{last_step}
\begin{proof}
Write $t_V$ for the image of $Q$ in $k_V$.
\begin{enumerate}
\item[(1)] Lemma~\ref{sha-siegel-step-coefficients}(2), applied to the chain from $c$ to $t_V$, Lemma~\ref{sha-siegel-step-coefficients}(3), applied to $P$, and Lemma~\ref{sha-siegel-step-ideal} each give a ring; let $\mathcal O'$, with residue field $k'$, contain these three rings and $\mathcal O_1$.
\item[(2)] By Lemma~\ref{sha-siegel-step-coefficients}(1) there is a deformation space $X=(A,B,P_X)$ for $T$ over $\mathcal O'$, and by Lemma~\ref{sha-siegel-step-ideal} it has a modularity ideal $J$.
\item[(3)] By Lemma~\ref{sha-siegel-step-coefficients}(2) there are flagged pairs $s$ and $t$ along $\mathcal O\to k'$, with images $c$ and $t_V$ in $k_V$, with $s$ residual and a chain from $s$ to $t$. By Lemma~\ref{sha-siegel-step-coefficients}(3) there are $R$, an injective $f:R\to V$ and a family $P_R$ over $R$ that $f$ carries to $P$.
\item[(4)] The pair $t$ is residual, as the last pair of a chain that starts at a residual pair (see the proof of Lemma~\ref{sha-siegel-step-linked}(4)). The image of $Q$ over $\mathcal O'$ is a lift of $t$: it is a flagged pair along $\mathcal O\to\mathcal O'$ with continuous entries, and its image in $k'$ is $t$, because both have the image $t_V$ in $k_V$ and $k'\to k_V$ is injective. By (5) of Definition~\ref{sha-siegel-tower-space}, a homomorphism of $\mathcal O'$-algebras $f_1:A\to\mathcal O'$ carries $P_X$ to the image of $Q$.
\item[(5)] Since the image of $Q$ in $\bar{\mathbf Q}_\ell$ is a seed, the seed property of $J$ shows that $f_1$ followed by $\mathcal O'\subset\bar{\mathbf Q}_\ell$ detects $J$. Its kernel is $\ker f_1$, which lies in the kernel of the point $x_1$, the composite of $f_1$ with the residue map; so by Lemma~\ref{sha-siegel-step-linked}(2), $J$ lies in some minimal prime below $\ker x_1$. The image of $x_1$ is $t$.
\item[(6)] The family $P_R$ is a lift of $s$: it is a flagged pair along $\mathcal O\to R$, because its image under the injective homomorphism $f$ is one, its entries are continuous, and its image in $k'$ is $s$, because both have the image $c$ in $k_V$. By (5) of Definition~\ref{sha-siegel-tower-space}, a homomorphism of $\mathcal O'$-algebras $f_0:A\to R$ carries $P_X$ to $P_R$; the point $x_0$, the composite of $f_0$ with the residue map, has image $s$, which is residual.
\item[(7)] By Lemma~\ref{sha-siegel-step-propagation}, the hypothesis of Lemma~\ref{sha-siegel-step-linked}(4) holds, so by that lemma and the chain from $s$ to $t$, $J$ lies in all minimal primes below $\ker x_0$.
\item[(8)] Let $y=f\circ f_0:A\to V$. Its kernel is $\ker f_0$, as $f$ is injective, and lies in $\ker x_0$; so $y$ kills $J$ by Lemma~\ref{sha-siegel-step-linked}(3). It restricts to the given map on $\mathcal O'$, it carries $P_X$ to $f(P_R)=P$, and it maps $\mathfrak m_B$ into $\mathfrak m_V$, because $x_0$ kills $\mathfrak m_B$ and $f$ is compatible with the residue maps.
\end{enumerate}
\end{proof}

\subsection{Adjoint vanishing and links of towers}\label{subsection-sha-tower-links}

\begin{lemma}\label{sha-siegel-automorphy-clifford}
Let $\Gamma$ be a group, let $H$ be a normal subgroup of $\Gamma$ that contains every commutator $aba^{-1}b^{-1}$, let $\kappa$ be an algebraically closed field of characteristic $0$, let $m\ge1$, and let $\rho:\Gamma\to GL_m(\kappa)$ be an irreducible representation whose restriction to $H$ is reducible.
Then there are a character $\phi:\Gamma\to\kappa^\times$ with $\phi\ne1$, $H\subseteq\ker\phi$ and $\phi^m=1$, and a matrix $M\in GL_m(\kappa)$ with $M\phi(g)\rho(g)=\rho(g)M$ for all $g\in\Gamma$.
\end{lemma}
\begin{proof}
Let $V=\kappa^m$, on which $\Gamma$ acts through $\rho$, and call a subspace of $V$ \emph{$H$-stable} if $\rho(h)$ preserves it for every $h\in H$.
\begin{enumerate}
\item[(1)] \emph{An $H$-stable decomposition.} Since the restriction of $\rho$ to $H$ is reducible, there are nonzero proper $H$-stable subspaces. Let $W$ be minimal among the nonzero $H$-stable subspaces, and $U$ maximal among the nonzero proper ones. For $g\in\Gamma$ the subspace $\rho(g)^{-1}U$ is $H$-stable, because $H$ is normal. Some $g$ has $W\not\subseteq\rho(g)^{-1}U$: otherwise the intersection of the $\rho(g)^{-1}U$ over all $g\in\Gamma$ would be $\Gamma$-stable, would contain $W\ne0$ and would be contained in $U\ne V$, which contradicts the irreducibility of $\rho$. Put $U'=\rho(g)^{-1}U$. Then $U'\cap W$ is $H$-stable and properly contained in $W$, so it is $0$ by the minimality of $W$. If $U'+W$ were not $V$, then $\rho(g)(U'+W)=U+\rho(g)W$ would be a proper $H$-stable subspace properly containing $U$, which contradicts the maximality of $U$. So $V=U'\oplus W$, with $U'$ and $W$ nonzero and $H$-stable.
\item[(2)] \emph{A traceless endomorphism commuting with $H$.} Let $Q$ be the projection of $V$ onto $U'$ along $W$; it commutes with $\rho(h)$ for $h\in H$. Then $P=mQ-\operatorname{tr}(Q)$ is traceless and commutes with $\rho(H)$. It is nonzero: it acts on $W$ by $-\operatorname{tr}(Q)$ and on $U'$ by $m-\operatorname{tr}(Q)$, and these do not both vanish, since $m\ne0$ in $\kappa$.
\item[(3)] \emph{A common eigenvector.} Let $\mathcal S$ be the space of traceless endomorphisms of $V$ that commute with $\rho(H)$; it contains $P\ne0$. For $g\in\Gamma$ the map $f_g(T)=\rho(g)T\rho(g)^{-1}$ preserves $\mathcal S$, since $H$ is normal, and $f_h$ is the identity on $\mathcal S$ for $h\in H$. As $f_af_b=f_{aba^{-1}b^{-1}}f_bf_a$ and $H$ contains the commutators, the $f_g$ commute with each other on $\mathcal S$. Let $\mathcal F$ be a minimal nonzero subspace of $\mathcal S$ stable under all $f_g$. For each $g$, $f_g$ has an eigenvalue on $\mathcal F$, since $\kappa$ is algebraically closed; the intersection of $\mathcal F$ with the corresponding eigenspace is nonzero and stable under all $f_{g'}$, as these commute with $f_g$, so it is $\mathcal F$. Hence a nonzero $T\in\mathcal F$ satisfies $\rho(g)T=\phi(g)T\rho(g)$ for all $g\in\Gamma$, with $\phi(g)\in\kappa$.
\item[(4)] \emph{The character.} Since $T\ne0$ and $\rho(g)$ is invertible, $\phi(g)\ne0$ and $\phi(g)$ is determined by $T$; so $\phi$ is a character of $\Gamma$, and it is trivial on $H$ because $T$ commutes with $\rho(H)$. The kernel of $T$ is $\Gamma$-stable and is not $V$, so it is $0$ and $T$ is invertible; its matrix $M$ satisfies $M\phi(g)\rho(g)=\rho(g)M$. If $\phi$ were trivial, $T$ would commute with $\rho(\Gamma)$; an eigenspace of $T$ would then be a nonzero $\Gamma$-stable subspace, hence $V$, so $T$ would be a scalar, and $T=0$ since $\operatorname{tr}T=0$ and $m\ne0$ in $\kappa$. So $\phi\ne1$. Taking determinants gives $\phi(g)^m\det\rho(g)=\det\rho(g)$, so $\phi^m=1$.
\end{enumerate}
\end{proof}
\leanref{Development/AutomorphicArithmetic/AutomorphicVarieties/UnitaryThreeGroups/Clifford.lean}{UpperTheory.V.U3.exists_selfTwist}

\begin{theorem}[\cite{AC89,BLGGT14}]\label{sha-siegel-automorphy-base-change}
Let $F$ be a CM field with maximal totally real subfield $F_0$, let $M\supseteq F$ be a CM field, Galois over $F$ with solvable Galois group, let $\ell$ be a prime, let $\iota:\bar{\mathbf Q}_\ell\cong\mathbf C$, let $n\ge1$, and fix an algebraic closure $\bar F$ of $F$.
Let $r:\operatorname{Gal}(\bar F/F)\to GL_n(\bar{\mathbf Q}_\ell)$ be a continuous representation whose restriction to $\operatorname{Gal}(\bar F/M)$ is irreducible, and let $\mu:\operatorname{Gal}(\bar F/F_0)\to\bar{\mathbf Q}_\ell^\times$ be the character attached through $\iota$ to an algebraic Hecke character of $F_0$, taking the same value at all complex conjugations.
Suppose that there are a complex conjugation $c\in\operatorname{Gal}(\bar F/F_0)$ and $P\in GL_n(\bar{\mathbf Q}_\ell)$ with ${}^tr(s)\,P\,r(csc^{-1})=\mu(s)P$ for all $s\in\operatorname{Gal}(\bar F/F)$.
Then $r$ is attached through $\iota$ to a regular algebraic cuspidal automorphic representation of $GL_n(\mathbf A_F)$ if and only if its restriction to $\operatorname{Gal}(\bar F/M)$ is attached through $\iota$ to a regular algebraic cuspidal automorphic representation of $GL_n(\mathbf A_M)$.
\end{theorem}
\leanref{Development/AutomorphicArithmetic/AutomorphicVarieties/ModularityArgument/BaseChangeUnits/SolCOfSolDescent.lean}{UpperTheory.V.Z.Bc.SolC}
\leanref{Development/AutomorphicArithmetic/AutomorphicVarieties/ModularityArgument/BaseChangeUnits/SolQ.lean}{UpperTheory.V.Z.Bc.solQ}

\begin{lemma}\label{sha-siegel-automorphy-polarized}
Let $K$ be an imaginary quadratic field, let $\ell$ be a prime, $\iota:\bar{\mathbf Q}_\ell\cong\mathbf C$, $n\ge1$ and $u\in\mathbf Z$.
Let $\theta:\operatorname{Gal}(\bar{\mathbf Q}/\mathbf Q)\to\bar{\mathbf Q}_\ell^\times$ be $\chi_{\mathrm{cyc}}^{u}$ times a character that is trivial on $\operatorname{Gal}(\bar{\mathbf Q}/\mathbf Q(\mu_M))$ for some $M\ge1$, and let $r:\operatorname{Gal}(\bar{\mathbf Q}/\mathbf Q)\to GL_n(\bar{\mathbf Q}_\ell)$ be a homomorphism with ${}^tr(s)\,P\,r(s)=\theta(s)P$ for all $s$ and some $P\in GL_n(\bar{\mathbf Q}_\ell)$.
Suppose that the restriction of $r$ to $\operatorname{Gal}(\bar{\mathbf Q}/K)$ is attached through $\iota$ to a cuspidal automorphic representation $\pi$ of $GL_n(\mathbf A_K)$.
Then there is an algebraic Hecke character $\chi$ of $\mathbf Q$ such that $(\pi,\chi)$ is polarized (Definition~\ref{gal-polarized}) and the character attached to $\chi$ through $\iota$ equals $\theta\chi_{\mathrm{cyc}}^{n-1}$ on $\operatorname{Gal}(\bar{\mathbf Q}/K)$.
\end{lemma}
\begin{proof}
The character $\theta\chi_{\mathrm{cyc}}^{n-1}$ is $\chi_{\mathrm{cyc}}^{u+n-1}$ times a character of finite order, so by class field theory it is the character attached through $\iota$ to an algebraic Hecke character $\chi_0$ of $\mathbf Q$ \cite{Ser68}.
The quadratic character $\eta_K$ of $K$ is trivial on the norms from $K$, its Galois character is trivial on $\operatorname{Gal}(\bar{\mathbf Q}/K)$, and $\eta_{K,\infty}(-1)=-1$; so one of $\chi_0$ and $\chi_0\eta_K$, call it $\chi$, has $\chi_\infty(-1)=(-1)^n$, and both have the same composite with the norm from $K$ and the same Galois character on $\operatorname{Gal}(\bar{\mathbf Q}/K)$.
For almost all primes $w$ of $K$, Frobenius elements at $w$ and at its conjugate $cw$ are conjugate in $\operatorname{Gal}(\bar{\mathbf Q}/\mathbf Q)$, so $\pi$ has the same Satake parameters $B_w$ at $w$ and at $cw$.
Since ${}^tr(s)Pr(s)=\theta(s)P$, the eigenvalues of $r$ at a Frobenius element $\sigma_w$ are stable under $b\mapsto\theta(\sigma_w)b^{-1}$; through $\iota$ this says that $B_w$ is stable under $b\mapsto(\chi\circ N_{K/\mathbf Q})(w)b^{-1}$.
So the representation $g\mapsto\pi({}^tc(g)^{-1})$, which has at $w$ the inverses of the Satake parameters of $\pi$ at $cw$, has the same Satake parameters as $\pi\otimes(\chi\circ N_{K/\mathbf Q}\circ\det)^{-1}$ at almost all places.
By strong multiplicity one (Theorem~\ref{lit-jacquet-shalika}(1)) the two are isomorphic, and since the contragredient of a cuspidal representation is its composite with the transpose inverse \cite{GK75}, this isomorphism gives the pairing of Definition~\ref{gal-polarized}; so $(\pi,\chi)$ is polarized.
\end{proof}
\leanref{Development/AutomorphicArithmetic/AutomorphicVarieties/ModularityArgument/BaseChangeUnits/SolCOfSolDescent.lean}{UpperTheory.V.Z.Bc.PolK}
\leanref{Development/AutomorphicArithmetic/AutomorphicVarieties/ModularityArgument/ThresholdArguments/DenseThresholdCases/UnderCPl.lean}{UpperTheory.V.Z.Bc.polarization_k_of_printed_results}

\begin{lemma}\label{sha-siegel-automorphy-free}
Assume the Galois representations as in Theorem~\ref{sha-residually-reducible-automorphy}.
Let $K$ be an imaginary quadratic field, $\ell$ an odd prime, $\iota:\bar{\mathbf Q}_\ell\cong\mathbf C$, $n\ge2$, $M\ge1$, and $m_1>m_2>\dots>m_n$ integers.
Let $r:\operatorname{Gal}(\bar{\mathbf Q}/K)\to GL_n(\bar{\mathbf Q}_\ell)$ be continuous, attached through $\iota$ to a regular algebraic cuspidal automorphic representation $\pi$ of $GL_n(\mathbf A_K)$ such that $(\pi,\chi)$ is polarized for some Hecke character $\chi$ of $\mathbf Q$ (Definition~\ref{gal-polarized}), and windowed (Theorem~\ref{sha-residually-reducible-automorphy}).
Suppose that for every prime $w$ of $K$ above $\ell$ there are a place $A$ of $\bar{\mathbf Q}$ above $w$ and characters $\kappa_1,\dots,\kappa_n$ of the decomposition group $D_A\subseteq\operatorname{Gal}(\bar{\mathbf Q}/K)$ such that the restriction of $r$ to $D_A$ is conjugate to an upper triangular representation with diagonal characters $\kappa_1,\dots,\kappa_n$, each $\kappa_i$ equals $\chi_{\mathrm{cyc}}^{m_i}\chi_i$ on the inertia group $I_A$ for a character $\chi_i$ of $\operatorname{Gal}(\bar{\mathbf Q}/K)$ trivial on $\operatorname{Gal}(\bar{\mathbf Q}/K(\mu_M))$, and $\kappa_i\ne\kappa_j\chi_{\mathrm{cyc}}$ on $D_A$ for all $i,j$.
Then for every place $A'$ of $\bar{\mathbf Q}$ above $\ell$ there is $g\in D_{A'}$ with $\chi_{\mathrm{cyc}}(g)=1$ such that the eigenvalues $\alpha_1,\dots,\alpha_n$ of $r(g)$ satisfy $\alpha_i^e\ne\alpha_j^e$ for $i\ne j$ and $e\ge1$.
\end{lemma}
\begin{proof}
The characters $\chi_i$ have finite order, so each $\kappa_i\chi_{\mathrm{cyc}}^{-m_i}$ has finite order on $I_A$.
The proof of Lemma~\ref{sha-rigidity-free}, with $F=K$ and $a=0$, then applies. It uses $\pi$ only through the windowed condition and through the triple attached to $\pi$ (Definition~\ref{gal-attached}), which exists because $(\pi,\chi)$ is polarized (Definition~\ref{gal-galois-representations}).
\end{proof}
\leanref{Development/AutomorphicArithmetic/AutomorphicVarieties/ModularityArgument/ThresholdArguments/DenseThresholdCases/RRFree.lean}{UpperTheory.V.Z.Win.RR.freeAt_of_autPA_of_unl_of_finrank_eq_two}
\leanref{Development/AutomorphicArithmetic/AutomorphicVarieties/ModularityArgument/ThresholdArguments/DenseThresholdCases/RRFreeW.lean}{UpperTheory.V.Z.Win.RR.freePUA_alt}

\begin{theorem}[\cite{NT23}]\label{sha-siegel-automorphy-adjoint}
In Situation~\ref{sha-hecke-situation}, suppose that $2\le n\le4$ and that the residual datum satisfies: the $\bar\chi_i$ are trivial on $I_{v_0}$ and take equal values on $D_{v_0}$, and $\varepsilon\bmod\ell$ is nontrivial on $D_{v_0}$.
Let $\iota:\bar{\mathbf Q}_\ell\cong\mathbf C$, let $x:B\to\bar{\mathbf Q}_\ell$ be a point (Definition~\ref{sha-hecke-arcs}), and let $\rho:\Gamma_0\to GL_n(\bar{\mathbf Q}_\ell)$ be a continuous representation such that:
\begin{enumerate}
\item $x$ carries $D^{\mathrm u}$ to the determinant of $\rho$, and $x$ is arithmetic of some increasing weight;
\item $\rho$, as a representation of $\operatorname{Gal}(\bar L/L)$, is attached through $\iota$ to a regular algebraic cuspidal automorphic representation $\Pi$ of $GL_n(\mathbf A_L)$ with $\Pi^\vee\cong\Pi^c$, and is windowed (Theorem~\ref{sha-residually-reducible-automorphy});
\item $\rho$ is trivial on $I_{v_0}$, and for each $v\mid\ell$ the restriction of $\rho$ to $D_v$ is conjugate to an upper triangular representation with diagonal characters $\psi_{x,v,1},\dots,\psi_{x,v,n}$;
\item for each $v\mid\ell$ some $g\in D_v$ with $\varepsilon(g)=1$ has $\psi_{x,v,i}(g)^e\ne\psi_{x,v,j}(g)^e$ for $i\ne j$ and $e\ge1$;
\item the restriction of $\rho$ to $\Gamma'$ is irreducible, and every open subset of $\Gamma_0$ that meets $\Gamma'$ contains an element $\sigma$ of $\Gamma'$ at which the characteristic polynomial of $\rho(\sigma)$ is separable;
\item for each $v\mid\ell$ there are integers $m_{v,1}>\dots>m_{v,n}$ and $k\ge1$ with $\big(\psi_{x,v,i}\psi_{x,v,j}^{-1}\varepsilon^{m_{v,j}-m_{v,i}}\big)^k=1$ on $I_v$ for all $i,j$.
\end{enumerate}
Let $P\in GL_n(\bar{\mathbf Q}_\ell)$ satisfy ${}^t\rho(g)P\rho(cgc^{-1})=\varepsilon(g)^{1-n}P$ for $g\in\Gamma_0$, and let $R$ be a representation of $\Gamma$ on $M_n(\bar{\mathbf Q}_\ell)$ with $R(g)X=\rho(g)X\rho(g)^{-1}$ for $g\in\Gamma_0$ and $R(c)X=-P^{-1}\,{}^tXP$.
Then every continuous cocycle of $\Gamma$ with values in $R$ whose inflation to $\operatorname{Gal}(\bar L/L^+)$ satisfies the Bloch-Kato conditions \cite{BK90} is a coboundary.
\end{theorem}
\begin{proof}
This is a case of \cite[Theorem~5.2]{NT23}, whose hypotheses follow from (2) and (5); the other hypotheses are not needed for it. Write $\mathrm{pr}:\operatorname{Gal}(\bar L/L)\to\Gamma_0$ for the projection and $\rho_L=\rho\circ\mathrm{pr}$.
\begin{enumerate}
\item[(1)] \emph{The representation.} Let $\delta$ be the quadratic character of $L/L^+$ and $\chi=\delta^n$. Since $\chi\circ N_{L/L^+}=1$ and $\delta_v(-1)=-1$ at every infinite place $v$ of $L^+$, the relation $\Pi^c\cong\Pi^\vee$ of (2) makes $(\Pi,\chi)$ polarized in the sense of \cite[\S2.1]{BLGGT14}. So $\Pi$ has a Galois representation $r_{\Pi,\iota}:\operatorname{Gal}(\bar L/L)\to GL_n(\bar{\mathbf Q}_\ell)$ \cite[Theorem~2.1.1]{BLGGT14}, compatible with the local Langlands correspondence at the finite places prime to $\ell$. At almost all finite places of $L$, both $r_{\Pi,\iota}$ and $\rho_L$ are unramified, with characteristic polynomials of Frobenius elements given by the Satake parameters of $\Pi$. By the Chebotarev density theorem and the Brauer-Nesbitt theorem \cite[(30.16)]{CR62}, the semisimplification of $r_{\Pi,\iota}$ is isomorphic to $\rho_L$. By (5), $\rho$ is irreducible, so $r_{\Pi,\iota}$ is irreducible and isomorphic to $\rho_L$: there is $A\in GL_n(\bar{\mathbf Q}_\ell)$ with $r_{\Pi,\iota}=A\rho_LA^{-1}$.
\item[(2)] \emph{The adjoint representation.} Lift $c$ to a complex conjugation $\tilde c\in\operatorname{Gal}(\bar L/L^+)$. By \cite[\S1]{NT23}, there is $Q\in GL_n(\bar{\mathbf Q}_\ell)$ with ${}^tr_{\Pi,\iota}(\sigma)\,Q\,r_{\Pi,\iota}(\tilde c\sigma\tilde c^{-1})=\varepsilon(\sigma)^{1-n}Q$ for $\sigma\in\operatorname{Gal}(\bar L/L)$, and $\operatorname{Gal}(\bar L/L^+)$ acts on $\operatorname{ad}r_{\Pi,\iota}=M_n(\bar{\mathbf Q}_\ell)$ by conjugation through $r_{\Pi,\iota}$ on $\operatorname{Gal}(\bar L/L)$ and by $Y\mapsto-Q^{-1}\,{}^tYQ$ at $\tilde c$. The matrix $P_1={}^tAQA$ satisfies ${}^t\rho(g)P_1\rho(cgc^{-1})=\varepsilon(g)^{1-n}P_1$ for $g\in\Gamma_0$. Comparing with the same relation for $P$ shows that $P^{-1}P_1$ commutes with every $\rho(cgc^{-1})$, $g\in\Gamma_0$, hence with $\rho(\Gamma_0)$; as $\rho$ is irreducible, $P_1$ is a scalar multiple of $P$. So $X\mapsto AXA^{-1}$ carries the inflation of $R$ to $\operatorname{Gal}(\bar L/L^+)$ isomorphically to $\operatorname{ad}r_{\Pi,\iota}$: on $\operatorname{Gal}(\bar L/L)$ both act by conjugation, and $A(-P^{-1}\,{}^tXP)A^{-1}=-Q^{-1}\,{}^t(AXA^{-1})Q$.
\item[(3)] \emph{Enormous image.} The projection maps $\operatorname{Gal}(\bar L/L(\zeta_{\ell^\infty}))$, the kernel of $\varepsilon$ on $\operatorname{Gal}(\bar L/L)$, onto $\Gamma'=\Gamma_0\cap\ker\varepsilon$; so $H=r_{\Pi,\iota}(\operatorname{Gal}(\bar L/L(\zeta_{\ell^\infty})))$ is $A\rho(\Gamma')A^{-1}$. The group $\operatorname{Gal}(\bar L/L)$ is compact, so after replacing $r_{\Pi,\iota}$ by a conjugate, and $A$ accordingly, $r_{\Pi,\iota}$ takes values in $GL_n(\mathcal O_{E_0})$ for a finite extension $E_0$ of $\mathbf Q_\ell$, and then in $GL_n(\mathcal O_E)$ for every finite extension $E\supseteq E_0$. Take $E$ to contain the finitely many extensions of $E_0$ of degree at most $n$; then the characteristic polynomial of every element of $H$ splits over $E$, as in \cite[Lemma~2.22]{NT23}. Then $H$ is enormous \cite[Definition~2.23]{NT23}. Indeed, let $V\subseteq M_n(E)$ be a nonzero $E[H]$-submodule for the action by conjugation, and suppose that $\operatorname{tr}(e_{h,\alpha}v)=0$ for all $v\in V$, all $h\in H$ with $n$ distinct eigenvalues and all eigenvalues $\alpha$ of $h$, where $e_{h,\alpha}$ is the projection onto the $\alpha$-eigenspace of $h$ along the others. Since such an $h$ equals $\sum_\alpha\alpha e_{h,\alpha}$, $\operatorname{tr}(hv)=0$ for these $h$. By the second half of (5), the $\sigma\in\Gamma'$ at which $\rho(\sigma)$, and so $A\rho(\sigma)A^{-1}$, has distinct eigenvalues are dense in $\Gamma'$, and $\sigma\mapsto\operatorname{tr}(A\rho(\sigma)A^{-1}v)$ is continuous; so $\operatorname{tr}(hv)=0$ for all $h\in H$. By the first half of (5) and Burnside's theorem \cite[\S27]{CR62}, the $\bar{\mathbf Q}_\ell$-span of $H$ is $M_n(\bar{\mathbf Q}_\ell)$, so the $E$-span of $H\subseteq M_n(E)$ is $M_n(E)$. As the trace form is nondegenerate, $V=0$, a contradiction. So every simple $E[H]$-submodule $V$ has some $h$ with $n$ distinct eigenvalues and an eigenvalue $\alpha$ with $\operatorname{tr}(e_{h,\alpha}V)\ne0$.
\item[(4)] \emph{Conclusion.} Let $z$ be a continuous cocycle of $\Gamma$ with values in $R$ whose inflation satisfies the Bloch-Kato conditions. By (2), $\sigma\mapsto Az(\bar\sigma)A^{-1}$, where $\bar\sigma$ is the image of $\sigma$ in $\Gamma$, is a continuous cocycle of $\operatorname{Gal}(\bar L/L^+)$ with values in $\operatorname{ad}r_{\Pi,\iota}$, and its class lies in the Bloch-Kato Selmer group $H^1_f(L^+,\operatorname{ad}r_{\Pi,\iota})$ of \cite[\S1]{NT23}: it is unramified at the places prime to $\ell$ and crystalline at the places above $\ell$. By (1) and (3), \cite[Theorem~5.2]{NT23} applies to $(\Pi,\chi)$, so this group vanishes, and there is $Y$ with $Az(\bar\sigma)A^{-1}=\sigma\cdot Y-Y$. With $X=A^{-1}YA$ this says $z(\bar\sigma)=R(\bar\sigma)X-X$ for all $\sigma$, and $\sigma\mapsto\bar\sigma$ is surjective.
\end{enumerate}
\end{proof}
\leanref{Development/AutomorphicArithmetic/AutomorphicVarieties/ModularityArgument/ThresholdArguments/DenseThresholdCases/AutomorphicPointGalois.lean}{UpperTheory.V.Th.AdjointLiftingTheorem}
\leanref{Development/AutomorphicArithmetic/AutomorphicVarieties/ModularityArgument/ThresholdArguments/DenseThresholdCases/AutomorphicPointGalois.lean}{UpperTheory.V.Th.AutomorphicPointGalois}
\leanref{Development/AutomorphicArithmetic/AutomorphicVarieties/ModularityArgument/ThresholdArguments/DenseThresholdCases/AutomorphicPointGalois.lean}{UpperTheory.V.Th.AutomorphicPointLocallyGeneric}
\leanref{Development/AutomorphicArithmetic/AutomorphicVarieties/TheoremsMuLinked/OrdZero.lean}{UpperTheory.V.Th.AdZero}
\leanref{Development/AutomorphicArithmetic/AutomorphicVarieties/TheoremsMuLinked/OrdZero.lean}{UpperTheory.V.Th.Par}
\leanref{Development/AutomorphicArithmetic/AutomorphicVarieties/TheoremsMuLinked/OrdZero.lean}{UpperTheory.V.Th.Free}
\leanref{Development/AutomorphicArithmetic/AutomorphicVarieties/TheoremsMuLinked/OrdZero.lean}{UpperTheory.V.Th.TypeI}
\leanref{Development/AutomorphicArithmetic/AutomorphicVarieties/TheoremsMuLinked/OrdZero.lean}{UpperTheory.V.Th.N3}
\leanref{Development/AutomorphicArithmetic/AutomorphicVarieties/WeilGroupCharacters/ArithmeticPointClassicality.lean}{UpperTheory.V.W.Resid.V0ok}

\begin{lemma}\label{sha-siegel-automorphy-descent}
Let $K$ be an imaginary quadratic field and $L\supseteq K$ a CM field with maximal totally real subfield $L^+$; let $\ell$ be a prime and $n\ge1$, and write $G=\operatorname{Gal}(\bar{\mathbf Q}/\mathbf Q)$ and $G_E=\operatorname{Gal}(\bar{\mathbf Q}/E)$.
Let $S'$ be a finite set of places of $L^+$ containing those above $\ell$, let $N'$ be the closed normal subgroup of $G_{L^+}$ generated by the inertia groups at the places of $L^+$ not in $S'$, let $\Gamma=G_{L^+}/N'$, let $\Gamma_0$ be the image of $G_L$ and $\mathrm{pr}:G_L\to\Gamma_0$ the projection, and let $\tau\in G_{L^+}$ be a complex conjugation, with image $c'\in\Gamma$.
Let $r:G_K\to GL_n(\bar{\mathbf Q}_\ell)$ be a continuous representation, $\mu:G_K\to\bar{\mathbf Q}_\ell^\times$ a character, and $R$ a representation of $G$ on $M_n(\bar{\mathbf Q}_\ell)$ such that, for some complex conjugation $c$ and some $P\in GL_n(\bar{\mathbf Q}_\ell)$, ${}^tr(s)Pr(csc^{-1})=\mu(s)P$ and $R(s)X=r(s)Xr(s)^{-1}$ for $s\in G_K$, and $R(c)X=-P^{-1}\,{}^tXP$.
Let $\rho:\Gamma_0\to GL_n(\bar{\mathbf Q}_\ell)$ and $\kappa:G_L\to\bar{\mathbf Q}_\ell^\times$ be homomorphisms with $\rho(\mathrm{pr}(g))=\kappa(g)r(g)$ and $\kappa(g)\kappa(\tau g\tau^{-1})\mu(g)=\chi_{\mathrm{cyc}}(g)^{1-n}$ for $g\in G_L$.
Suppose that for every $P'\in GL_n(\bar{\mathbf Q}_\ell)$ with ${}^t\rho(x)P'\rho(c'xc'^{-1})=\chi_{\mathrm{cyc}}(x)^{1-n}P'$ for $x\in\Gamma_0$, and every representation $R'$ of $\Gamma$ on $M_n(\bar{\mathbf Q}_\ell)$ with $R'(x)X=\rho(x)X\rho(x)^{-1}$ for $x\in\Gamma_0$ and $R'(c')X=-P'^{-1}\,{}^tXP'$, every continuous cocycle of $\Gamma$ with values in $R'$ whose inflation to $G_{L^+}$ satisfies the Bloch-Kato conditions \cite{BK90} is a coboundary.
Then every continuous cocycle of $G$ with values in $R$ that satisfies the Bloch-Kato conditions is a coboundary.
\end{lemma}
\begin{proof}
Let $y$ be a continuous cocycle of $G$ with values in $R$ that satisfies the Bloch-Kato conditions.
\begin{enumerate}
\item[(1)] The restriction of $y$ to $G_{L^+}$ satisfies the Bloch-Kato conditions, since these are local conditions and are preserved by restriction \cite{BK90}.
\item[(2)] The complex conjugations $\tau$ and $c$ both act on $K$ by its nontrivial automorphism, so $\tau=wc$ with $w\in G_K$.
\item[(3)] For $g\in G_L$, $R(g)X=r(g)Xr(g)^{-1}=\rho(\mathrm{pr}(g))X\rho(\mathrm{pr}(g))^{-1}$, the scalar $\kappa(g)$ cancelling. The subgroup $N'$ has trivial image in $\Gamma$, so it lies in $G_L$ and maps to $1$ in $\Gamma_0$; hence $R$ is trivial on $N'$.
\item[(4)] The cocycle $y$ vanishes on $N'$. Indeed, the elements of $N'$ at which $y$ vanishes form a closed normal subgroup of $G_{L^+}$, since $y$ is a continuous cocycle and $R$ is trivial on $N'$. Let $u\notin S'$ be a place of $L^+$; it is not above $\ell$. At one place above $u$ the Bloch-Kato condition makes $y$ on the inertia group a coboundary $s\mapsto R(s)X-X$, which is $0$ because $R$ is trivial there; the inertia groups at the other places above $u$ are conjugate to this one, so $y$ vanishes on them too, by normality.
\item[(5)] So $y$ induces a continuous cocycle $z$ of $\Gamma$ with values in the representation $R'$ of $\Gamma$ induced by $R$, and the inflation of $z$ satisfies the Bloch-Kato conditions by (1). By (3), $R'(x)X=\rho(x)X\rho(x)^{-1}$ for $x\in\Gamma_0$. Put $P'=P\,r(w)^{-1}$. Then $R'(c')X=R(w)R(c)X=-P'^{-1}\,{}^tXP'$. For $x=\mathrm{pr}(g)$ with $g\in G_L$, the element $c'xc'^{-1}$ is the image of $\tau g\tau^{-1}=wcgc^{-1}w^{-1}$, so
\[
 {}^t\rho(x)P'\rho(c'xc'^{-1})=\kappa(g)\kappa(\tau g\tau^{-1})\,{}^tr(g)Pr(cgc^{-1})r(w)^{-1}=\kappa(g)\kappa(\tau g\tau^{-1})\mu(g)P'=\chi_{\mathrm{cyc}}(g)^{1-n}P'.
\]
\item[(6)] By hypothesis $z$ is a coboundary, so there is $X$ with $y(s)=R(s)X-X$ for $s\in G_{L^+}$.
\item[(7)] The subgroup $G_{L^+}$ has finite index $d$ in $G$, and $M_n(\bar{\mathbf Q}_\ell)$ is a vector space over a field of characteristic $0$. The cocycle $y'(s)=y(s)-(R(s)X-X)$ vanishes on $G_{L^+}$, so it is constant on the left cosets of $G_{L^+}$. Let $\Sigma$ be the sum of its values on a set of representatives of these cosets. Since left multiplication by $s$ permutes the cosets, $\Sigma=d\,y'(s)+R(s)\Sigma$ for every $s\in G$. So $y(s)=R(s)Y-Y$ with $Y=X-d^{-1}\Sigma$.
\end{enumerate}
\end{proof}
\leanref{Development/AutomorphicArithmetic/AutomorphicVarieties/TheoremsMuLinked/FreeFlag.lean}{UpperTheory.V.Th.AdDown}
\leanref{Development/AutomorphicArithmetic/AutomorphicVarieties/TheoremsMuLinked/GaloisRepresentationTypes/AdDown.lean}{UpperTheory.V.Th.adDown}

\begin{definition}\label{sha-siegel-automorphy-link}
Let $D$ be a Siegel datum, with the notation of Definitions~\ref{sha-siegel-datum} and~\ref{sha-siegel-tower-towers}, let $T$ be a tower for $D$, with relabelling characters $\chi_1,\dots,\chi_4$ and fields $K$, $F$ and $L=FK$, and let $\mathcal O'\supseteq\mathcal O$ be a ring of Definition~\ref{sha-character-parts}. A \emph{link} of $T$ over $\mathcal O'$ consists of the following data.
\begin{enumerate}
\item The data of Situation~\ref{sha-hecke-situation} for the field $L$, with $\mathcal O'$ in place of $\mathcal O$ and $n=4$, in which $c$ is the image of a complex conjugation $\tau\in\operatorname{Gal}(\bar{\mathbf Q}/L^+)$; write $\mathrm{pr}:G_L\to\Gamma_0$ for the projection. With them, a bijection $v\mapsto\bar v$ from the places of $L^+$ above $\ell$ to $\Sigma_T$ such that $D_v$ is the decomposition group of $G_L$ at $\bar v$.
\item Characters $\kappa,\beta_L:G_L\to\mathcal O'^\times$, continuous for the $\mathfrak m_{\mathcal O'}$-adic topology, and a continuous homomorphism $t_L$ from $G_L$ to the additive group $\mathbf Z_\ell$, such that:
(a) $\kappa(g)\kappa(\tau g\tau^{-1})\theta(g)\chi_{\mathrm{cyc}}(g)^3=1$, $\beta_L(g)\beta_L(\tau g\tau^{-1})=1$ and $t_L(g)+t_L(\tau g\tau^{-1})=0$ for $g\in G_L$;
(b) $\kappa$ and $\bar r_0$ are trivial on the inertia groups of $G_L$ at the places above the places of $L^+$ outside $S$;
(c) for some $m$ prime to $\ell$, $\kappa^m$ is trivial on the inertia groups of $G_L$ at the places not above $\ell$;
(d) $\kappa$ and $t_L$ are trivial on $\mu_v$ for every $v\mid\ell$;
(e) for every place $w$ of $\bar{\mathbf Q}$ above $\ell$ there are $a_w\in\mathbf Z$, an odd $b_w\in\mathbf Z$ and $k_w\ge1$ with $(\kappa\chi_{\mathrm{cyc}}^{-a_w})^{k_w}=1$ and $(\beta_L^2\chi_{\mathrm{cyc}}^{-b_w})^{k_w}=1$ on $I_w(L)$;
(f) $\beta_L$ is congruent to $1$ modulo $\mathfrak m_{\mathcal O'}$ and trivial on the kernel of $t_L$, and $t_L$ is trivial on the inertia groups of $G_L$ at the places not above $\ell$.
\item A residual datum of rank $4$ such that $\bar\chi_i(\mathrm{pr}(g))$ is the residue of $\kappa(g)\chi_i(g)$ for $g\in G_L$, $\bar\psi_{v,i}$ is the restriction of $\bar\chi_i$ to $D_v$, and $\theta_{v,i}=1$; and such that the $\bar\chi_i$ are trivial on $I_{v_0}$ and take equal values on $D_{v_0}$, and $\varepsilon\bmod\ell$ is nontrivial on $D_{v_0}$.
\end{enumerate}
\end{definition}
\leanref{Development/AutomorphicArithmetic/AutomorphicVarieties/WeakShimuraTheory/IsAlgClosureField.lean}{UpperTheory.V.Ws.Link}

\begin{lemma}\label{sha-siegel-automorphy-phi}
Let $D$ be a Siegel datum, $T$ a tower for $D$, with fields $K$, $F$ and $L$, and $\mathcal O'\supseteq\mathcal O$ a ring of Definition~\ref{sha-character-parts}. Let $X=(A,B,P_X)$ be a deformation space for $T$ over $\mathcal O'$, with the polynomials $c_g$ and elements $\psi^B_{w,i}(s)$ of Definition~\ref{sha-siegel-tower-space}. Let a link of $T$ over $\mathcal O'$ be given (Definition~\ref{sha-siegel-automorphy-link}), with the notation there, and let $B^{\mathrm{ord}}$ be the universal ordinary determinant ring of its residual datum, with universal pair $(D^{\mathrm u},\psi^{\mathrm u})$ (Situation~\ref{sha-hecke-situation}). For $g\in G_L$ let $c^\kappa_g(Y)=\kappa(g)^4c_g(\kappa(g)^{-1}Y)\in B[Y]$.
Then there is a homomorphism of $\mathcal O'$-algebras $\varphi:B^{\mathrm{ord}}\to B$ that carries the characteristic polynomial of $D^{\mathrm u}$ at $\mathrm{pr}(g)$ to $c^\kappa_g$ for $g\in G_L$, and $\psi^{\mathrm u}_{v,i}(s)$ to $\kappa(s)\psi^B_{\bar v,i}(s)$ for $v\mid\ell$ and $s\in D_v$.
\end{lemma}
\begin{proof}
Let $\rho^\kappa:G_L\to M_4(A)$ be $g\mapsto\kappa(g)r_X(g)$, and write $\psi_{w,i}$ for the diagonal entries of $P_X$.
\begin{enumerate}
\item[(1)] \emph{Characteristic polynomials.} Since $c_g$ maps to the characteristic polynomial of $r_X(g)$ and $B\to A$ is injective, $c_g$ is monic of degree $4$; so $c^\kappa_g$ is monic and maps to the characteristic polynomial of $\rho^\kappa(g)$. Its coefficients are continuous in $g$, since those of $c_g$ and $\kappa$ are. By (2) of Definition~\ref{sha-siegel-tower-space} and (3) of Definition~\ref{sha-siegel-automorphy-link}, $c^\kappa_g$ reduces modulo $\mathfrak m_B$ to $\prod_i(Y-\bar\chi_i(\mathrm{pr}(g)))$.
\item[(2)] \emph{The kernel of $\mathrm{pr}$.} The projection $\mathrm{pr}$ is surjective and a quotient map. Its kernel lies in every closed subgroup $U$ of $G_L$ that contains the inertia groups of $G_L$ at the places above the places of $L^+$ outside $S$. Indeed, the kernel of $\operatorname{Gal}(\bar{\mathbf Q}/L^+)\to\Gamma$ is the closed normal subgroup generated by the inertia groups at these places; since $L/L^+$ is unramified at the finite places, these inertia groups lie in $G_L$, and their conjugates are the inertia groups at the conjugate places, which lie in $U$.
\item[(3)] \emph{Unramified places.} Let $g$ lie in such an inertia group. By (b) of Definition~\ref{sha-siegel-automorphy-link}, $\kappa(g)=1$ and $\bar r_0(g)=1$; the place is not above $\ell$, since $S$ contains the places above $\ell$, so $g$ lies in the set $N_\Phi$ of the problem of $T$ and $r_X(g)=1$. So $\rho^\kappa(g)=1$.
\item[(4)] \emph{The determinant.} The $k\in G_L$ with $c^\kappa_{hk}=c^\kappa_h$ for all $h$ form a subgroup, closed because $B$ is $\mathfrak m_B$-adically separated and the coefficients are continuous, and containing the kernel of $\rho^\kappa$, because $B\to A$ is injective. By (2) and (3), it contains the kernel of $\mathrm{pr}$, so $c^\kappa_g$ depends only on $\mathrm{pr}(g)$. Hence the determinant of $\rho^\kappa$ factors through $\Gamma_0$, and, since the characteristic polynomials of $\rho^\kappa$ at the elements of $G_L$ have coefficients in $B$ and $B\to A$ is injective, it descends by Amitsur's formula \cite{Che14}, as in the proof of Lemma~\ref{sha-hecke-glue}, to a continuous determinant $E$ of $\Gamma_0$ of dimension $4$ over $B$ whose characteristic polynomial at $\mathrm{pr}(g)$ is $c^\kappa_g$.
\item[(5)] \emph{The characters.} For $v\mid\ell$ and $s\in D_v$, which lies in $D_{\bar v}(F)$, put $\psi_{v,i}(s)=\kappa(s)\psi^B_{\bar v,i}(s)\in B^\times$. It is a continuous character, since its image in $A$ is $\kappa(s)\psi_{\bar v,i}(s)$ and $B\to A$ is injective. Its residue is $\bar\chi_i(\mathrm{pr}(s))=\bar\psi_{v,i}(s)$, by (3) of Definition~\ref{sha-siegel-tower-space} and (3) of Definition~\ref{sha-siegel-automorphy-link}. For $s\in\mu_v$, $\kappa(s)=1$ by (d) of Definition~\ref{sha-siegel-automorphy-link}, and $s$ lies in the corresponding subgroup $\mu_{\bar v}$ of $D_{\bar v}(F)$, where $\psi_{\bar v,i}=1$ since $P_X$ is a flagged pair; so $\psi_{v,i}(s)=1=\theta_{v,i}(s)$.
\item[(6)] \emph{An ordinary pair.} The determinant $E$ reduces to $\bar D$, the determinant of $\bar\chi_1\oplus\dots\oplus\bar\chi_4$, since by (1) the two have the same characteristic polynomials at the elements of $\Gamma_0$. The conditions (1) to (3) of Situation~\ref{sha-hecke-situation} are identities between traces of $E$, and since $B\to A$ is injective they can be checked for $\rho^\kappa$ over $A$.
For (1), let $g\in G_L$ and $g'=\tau g\tau^{-1}$. As $r_X(g)$ is symplectic with multiplier $\theta(g)$, $r_X(g)^{-1}=\theta(g)^{-1}r_X(g)^*$ has the characteristic polynomial of $\theta(g)^{-1}r_X(g)$, and $r_X(g'^{-1})$ is conjugate to $r_X(g)^{-1}$ under the image of $\tau$ in $\operatorname{Gal}(\bar{\mathbf Q}/F)$. So $\varepsilon(g)^{-3}\rho^\kappa(g'^{-1})$ has the characteristic polynomial of $\varepsilon(g)^{-3}\kappa(g')^{-1}\theta(g)^{-1}r_X(g)$, which is that of $\rho^\kappa(g)$ by (a) of Definition~\ref{sha-siegel-automorphy-link}; this is (1) at $\mathrm{pr}(g)$.
For (2), let $v\mid\ell$ and $n=n_{X,\bar v}$, which is symplectic with multiplier $1$. For $s\in D_v$, $n^{-1}\rho^\kappa(s)n=\kappa(s)n^*r_X(s)n$ is upper triangular with diagonal entries $\psi_{v,1}(s),\dots,\psi_{v,4}(s)$. So for $g_1,\dots,g_4\in D_v$ the product $\prod_i\big(\rho^\kappa(g_i)-\psi_{v,i}(g_i)\big)$ is conjugate to a product of upper triangular matrices of which the $i$-th has $i$-th diagonal entry $0$, which is $0$; this gives the ordinary relations.
For (3), let $u\in S_{\mathrm a}$ and let $g\in G_L$ lie in the inertia group at the place above $\tilde u$ with $\bar\chi_i(\mathrm{pr}(g))=1$ for all $i$. By (c) of Definition~\ref{sha-siegel-automorphy-link} there is $m$ prime to $\ell$ with $\kappa(g^m)=1$; then the residues of the $\chi_i(g^m)$ are $1$, so $\bar r_0(g^m)=1$, and as $u$ is not above $\ell$, $r_X(g^m)=1$ and $\rho^\kappa(g)^m=1$. By (1), $M=\rho^\kappa(g)$ is a root of the monic polynomial $c^\kappa_g$ over $B$, which reduces to $(Y-1)^4$. This polynomial and $1+Y+\dots+Y^{m-1}$ are coprime over $B$, as their reductions are, $m$ being nonzero in the residue field; so $1+M+\dots+M^{m-1}$ is invertible, and it kills $M-1$. Hence $\rho^\kappa(g)=1$, which gives (3).
\item[(7)] \emph{The map.} So $(E,(\psi_{v,i}))$ is an ordinary pair over $B$, and the universal property of $B^{\mathrm{ord}}$ gives $\varphi$ carrying $(D^{\mathrm u},\psi^{\mathrm u})$ to it. It carries the characteristic polynomial of $D^{\mathrm u}$ at $\mathrm{pr}(g)$ to that of $E$, which is $c^\kappa_g$, and $\psi^{\mathrm u}_{v,i}$ to $\psi_{v,i}$.
\end{enumerate}
\end{proof}
\leanref{Development/AutomorphicArithmetic/AutomorphicVarieties/ModularityArgument/ThresholdArguments/DenseThresholdCases/TPhiEx.lean}{UpperTheory.V.Ws.exists_phi_universal_ordinary_det_ring_of_link}
\leanref{Development/AutomorphicArithmetic/AutomorphicVarieties/WeakShimuraTheory/IsAlgClosureField.lean}{UpperTheory.V.Ws.Link.IsPhi}

\begin{lemma}\label{sha-siegel-automorphy-stabilizer}
Let $K$ be an algebraically closed field, $G$ a group, $n\ge1$ and $\rho:G\to GL_n(K)$ an irreducible representation. Let $N$ be a subgroup of $G$ that contains every commutator $aba^{-1}b^{-1}$, and suppose that the restriction of $\rho$ to $N$ is reducible and that $\rho(u)$ has separable characteristic polynomial for some $u\in N$.
Then there is a nonzero proper subspace $W\subset K^n$, stable under $\rho(N)$, such that the $g\in G$ with $\rho(g)W\subseteq W$ form a subgroup $S\supseteq N$ whose index divides $n$.
\end{lemma}
\begin{proof}
\begin{enumerate}
\item[(1)] The subgroup $N$ is normal: for $g\in G$ and $h\in N$, $g^{-1}hg$ is the product of $h$ and the commutator $h^{-1}g^{-1}hg$. So for $g\in G$ and a $\rho(N)$-stable subspace $X$, the subspace $\rho(g)X$ is $\rho(N)$-stable.
\item[(2)] Let $b_1,\dots,b_n$ be a basis of eigenvectors of $\rho(u)$, whose eigenvalues are distinct. A subspace stable under $\rho(u)$ is spanned by the $b_i$ it contains. So a $\rho(N)$-stable subspace $X$ has dimension equal to the number of $i$ with $b_i\in X$.
\item[(3)] Let $W$ be a nonzero $\rho(N)$-stable subspace of least dimension. Since $\rho$ is reducible on $N$, some nonzero proper subspace is $\rho(N)$-stable, so $W\ne K^n$. For $g\in G$ the subspace $\rho(g)W$ is $\rho(N)$-stable by (1), nonzero and of the same dimension. So $\rho(g)W\subseteq W$ implies $\rho(g)W=W$, and the set $S$ of such $g$ is a subgroup; it contains $N$.
\item[(4)] If $\rho(g_1)W$ and $\rho(g_2)W$ both contain some $b_i$, they are equal: their intersection is $\rho(N)$-stable and nonzero, so by the minimality of $\dim W$ it equals both.
\item[(5)] Every $b_i$ lies in some $\rho(g)W$. Otherwise the sum of the $\rho(g)W$, $g\in G$, is a nonzero $\rho(G)$-stable subspace, hence $K^n$ as $\rho$ is irreducible; but by (2) each $\rho(g)W$ is spanned by basis vectors other than $b_i$, so the $b_i$-coordinate vanishes on the sum, which is false at $b_i$.
\item[(6)] The subspaces $\rho(g)W$ correspond to the cosets $gS$, and by (4) and (5) they partition $\{b_1,\dots,b_n\}$ into sets of $\dim W$ elements. Hence $[G:S]\dim W=n$.
\end{enumerate}
\end{proof}
\leanref{Development/AutomorphicArithmetic/AutomorphicVarieties/WeakShimuraTheory/StabilizerIndexBounds/ExistsStabilizerIndexDvdOfNotIrreducibleRestrict.lean}{UpperTheory.V.Ws.exists_stabilizer_index_dvd_of_not_irreducible_restriction}
\leanref{Development/AutomorphicArithmetic/AutomorphicVarieties/WeakShimuraTheory/StabilizerIndexBounds/InvariantSubmoduleEqSpanEigenbasisOfSeparable.lean}{UpperTheory.V.Ws.invariant_submodule_eq_span_eigenbasis_of_separable}

\begin{lemma}\label{sha-siegel-automorphy-sign}
Let $\ell$ be an odd prime, $G$ a compact topological group, $\varepsilon:G\to\mathbf Z_\ell^\times$ a continuous homomorphism and $\eta:G\to\{\pm1\}$ a homomorphism whose kernel is open and contains $\ker\varepsilon$. Then $\eta(g)=1$ for every $g\in G$ with $\varepsilon(g)\equiv1\pmod\ell$.
\end{lemma}
\begin{proof}
Suppose that $\eta(g)=-1$. As $\ell$ is odd, $\eta(g^{\ell^m})=-1$ for every $m\ge0$, so $g^{\ell^m}$ lies in the complement $C$ of $\ker\eta$. Since $\ker\eta$ is open, $C$ is closed, hence compact, and so is its image $\varepsilon(C)\subset\mathbf Z_\ell$. Since $\varepsilon(g)\equiv1\pmod\ell$, $\varepsilon(g)^{\ell^m}\equiv1\pmod{\ell^{m+1}}$, so $\varepsilon(g^{\ell^m})$ tends to $1$, and $1\in\varepsilon(C)$. So some $h\in C$ has $\varepsilon(h)=1$; then $h\in\ker\varepsilon\subseteq\ker\eta$, a contradiction.
\end{proof}
\leanref{Development/AutomorphicArithmetic/AutomorphicVarieties/WeakShimuraTheory/StabilizerIndexBounds/SignCharEqOneOfCyclotomicSubOneNotIsUnit.lean}{UpperTheory.V.Ws.sign_char_eq_one_of_cyclotomic_sub_one_not_isUnit}

\begin{lemma}\label{sha-siegel-automorphy-sign-twist}
Let $K$ and $k$ be fields, $A$ a commutative ring with an injective homomorphism $j:A\to K$ and a homomorphism $\phi:A\to k$, $G$ a group, $n\ge1$, $\rho:G\to GL_n(K)$ a representation, and $\bar\chi_1,\dots,\bar\chi_n:G\to k^\times$ characters such that, for every $g\in G$, the characteristic polynomial of $\rho(g)$ is the image under $j$ of a polynomial $p_g$ over $A$ with $\phi(p_g)=\prod_i(X-\bar\chi_i(g))$.
Let $\eta:G\to\{\pm1\}$ be a homomorphism such that $\rho(g)$ and $\eta(g)\rho(g)$ have the same characteristic polynomial for every $g$, and let $\bar\eta:G\to k^\times$ be its composite with $\{\pm1\}\to k^\times$. Then for each $i$ there is $i'$ with $\bar\eta\bar\chi_i=\bar\chi_{i'}$.
\end{lemma}
\begin{proof}
Let $g\in G$. Then $\prod_i(X-\bar\chi_i(g))=\prod_i(X-\bar\eta(g)\bar\chi_i(g))$. This is clear if $\eta(g)=1$. If $\eta(g)=-1$, the characteristic polynomial $c$ of $\rho(g)$ equals that of $-\rho(g)$, which is $(-1)^nc(-X)$. As $j$ is injective, $p_g=(-1)^np_g(-X)$, and applying $\phi$ gives $\prod_i(X-\bar\chi_i(g))=(-1)^n\prod_i(-X-\bar\chi_i(g))=\prod_i(X+\bar\chi_i(g))$.
So Lemma~\ref{sha-rigidity-characters}, applied to the $\bar\chi_i$ and the $\bar\eta\bar\chi_i$, gives a permutation $\sigma$ with $\bar\eta\bar\chi_i=\bar\chi_{\sigma(i)}$ for all $i$.
\end{proof}
\leanref{Development/AutomorphicArithmetic/AutomorphicVarieties/TheoremsMuLinked/CharacteristicPolynomials/CharpolyCompNegRoots.lean}{UpperTheory.V.Th.Det.perm4}

\begin{lemma}\label{sha-siegel-automorphy-monomial}
Let $K$ be an algebraically closed field, $G$ a group, $n\ge1$ and $\rho:G\to GL_n(K)$ a representation whose values span $M_n(K)$. Let $N$ be a normal subgroup of $G$ such that $\rho(g)\rho(h)=\rho(h)\rho(g)$ for $g,h\in N$, and suppose that $\rho(u)$ has separable characteristic polynomial for some $u\in N$.
\begin{enumerate}
\item There are a basis $b_1,\dots,b_n$ of $K^n$ and a homomorphism $\pi$ from $G$ to the group of permutations of $\{1,\dots,n\}$, with $N\subseteq\ker\pi$, such that $\rho(g)b_i$ is a multiple of $b_{\pi(g)(i)}$ for all $g$ and $i$, and such that for all $i,j$ some $g$ has $\pi(g)(i)=j$.
\item Let $k$ be a field, $A$ a commutative ring with an injective homomorphism $j:A\to K$ and a homomorphism $\phi:A\to k$, and $\bar\chi_1,\dots,\bar\chi_n:G\to k^\times$ characters such that for every $g$ the characteristic polynomial of $\rho(g)$ is the image under $j$ of a polynomial over $A$ whose image under $\phi$ is $\prod_i(X-\bar\chi_i(g))$. Then, for $\pi$ as in (1), $\bar\chi_1(g)=\dots=\bar\chi_n(g)$ for every $g\in\ker\pi$.
\end{enumerate}
\end{lemma}
\begin{proof}
(1) Let $b_1,\dots,b_n$ be a basis of eigenvectors of $\rho(u)$, with distinct eigenvalues $\lambda_1,\dots,\lambda_n$; an eigenvector of $\rho(u)$ with eigenvalue $\lambda_i$ is a multiple of $b_i$, and $\rho(u)$ has no other eigenvalues.
For $h\in N$, $\rho(h)$ commutes with $\rho(u)$, so $\rho(h)b_i$ is a multiple of $b_i$.
Let $g\in G$ and $1\le i\le n$. As $N$ is normal, $g^{-1}ug\in N$, so $\rho(g^{-1}ug)b_i=cb_i$ for some $c$, that is, $\rho(u)\rho(g)b_i=c\rho(g)b_i$. As $\rho(g)b_i\ne0$, $c$ is some $\lambda_j$ and $\rho(g)b_i$ is a nonzero multiple of $b_j$; put $\pi(g)(i)=j$.
Then $\rho(gh)b_i=\rho(g)\rho(h)b_i$ is a multiple of $b_{\pi(g)(\pi(h)(i))}$, so $\pi(gh)=\pi(g)\pi(h)$; and $\pi(1)$ is the identity. So $\pi$ is a homomorphism to the permutations, and $N\subseteq\ker\pi$.
If no $g$ had $\pi(g)(i)=j$, the $b_j$-coordinate of $\rho(g)b_i$ would vanish for every $g$, hence that of $Mb_i$ for every $M$ in the span $M_n(K)$ of the $\rho(g)$, which is false for the matrix sending $b_i$ to $b_j$.

(2) Let $R$ be the integral closure of $j(A)$ in $K$. As in the proof of Lemma~\ref{sha-rigidity-normal}(c), the homomorphism $j(A)\to k\to\bar k$ induced by $\phi$, for an algebraic closure $\bar k$ of $k$, extends to a homomorphism $\Phi:R\to\bar k$ \cite[Chapter~5, Exercise~2]{AM69}.
For $y\in\ker\pi$ write $\rho(y)b_i=e_y(i)b_i$. The characteristic polynomial $\prod_i(X-e_y(i))$ of $\rho(y)$ is the image under $j$ of a polynomial $p$ over $A$, monic as $j$ is injective; so the $e_y(i)$ lie in $R$, and $\prod_i(X-\Phi(e_y(i)))$ is the image of $p$ under $\phi$, which is $\prod_s(X-\bar\chi_s(y))$. Hence the $\Phi(e_y(i))$ and the $\bar\chi_s(y)$ are the same elements of $\bar k$, counted with multiplicity.
Let $g\in\ker\pi$ and $1\le i,l\le n$. By (1) there is $h$ with $\pi(h)(l)=i$; write $\rho(h^{-1})b_i=c\,b_l$ and $\rho(h)b_l=d\,b_i$, so that $cd=1$. The commutator $y=ghg^{-1}h^{-1}$ lies in $\ker\pi$, and $\rho(y)b_i=cd\,e_{g^{-1}}(l)e_g(i)\,b_i$, so $e_y(i)=e_g(i)e_{g^{-1}}(l)$. Every $\bar\chi_s(y)$ is $1$, since the $\bar\chi_s$ are characters; so $\Phi(e_y(i))=1$. As $e_g(l)e_{g^{-1}}(l)=1$, this gives $\Phi(e_g(i))=\Phi(e_g(l))$. So the $\Phi(e_g(i))$ are all equal, and so are the $\bar\chi_s(g)$.
\end{proof}
\leanref{Development/AutomorphicArithmetic/AutomorphicVarieties/AlgebraTools/MonomialRepresentations/MonomialOfComm.lean}{UpperTheory.V.AlgebraTools.monomialOfComm}
\leanref{Development/AutomorphicArithmetic/AutomorphicVarieties/AlgebraTools/MonomialRepresentations/MonomialRed.lean}{UpperTheory.V.AlgebraTools.monomialRed}

\begin{lemma}\label{sha-siegel-automorphy-noncommuting}
Let $K$ be an algebraically closed field, $k$ a field, $A$ a commutative ring with an injective homomorphism $j:A\to K$ and a homomorphism $\phi:A\to k$, $G$ a group, $\rho:G\to GL_4(K)$ a representation whose values span $M_4(K)$, and $\bar\chi_1,\dots,\bar\chi_4:G\to k^\times$ characters such that for every $g$ the characteristic polynomial of $\rho(g)$ is the image under $j$ of a polynomial over $A$ whose image under $\phi$ is $\prod_i(X-\bar\chi_i(g))$. Suppose that $(\bar\chi_s/\bar\chi_t)^e\ne1$ for some $s,t$ and all $1\le e\le4$. Let $Z$ be a normal subgroup of $G$ and $u\in Z$ such that $\rho(u)^d$ has separable characteristic polynomial for every $d\ge1$.
Then every subgroup $N$ of finite index of $Z$ that is normal in $Z$ contains elements $g,h$ with $\rho(g)\rho(h)\ne\rho(h)\rho(g)$.
\end{lemma}
\begin{proof}
Suppose that the $\rho(g)$, $g\in N$, commute with each other. Let $d=[Z:N]$ and let $N_0$ be the intersection of the conjugates of $N$ in $G$, a normal subgroup of $G$ contained in $N$.
For $c\in G$, $cuc^{-1}$ lies in $Z$, as $Z$ is normal, and its $d$-th power lies in $N$, as $N$ is normal of index $d$ in $Z$; so $u^d\in c^{-1}Nc$ for every $c$, and $u^d\in N_0$.
So Lemma~\ref{sha-siegel-automorphy-monomial}, applied to $N_0$ and $u^d$, gives $\pi$ with $\bar\chi_s=\bar\chi_t$ on $\ker\pi$.
Every permutation $\tau$ of four letters has $\tau^3=1$ or $\tau^4=1$, so for each $g\in G$, $g^3$ or $g^4$ lies in $\ker\pi$, and $\chi=\bar\chi_s/\bar\chi_t$, which is trivial on $\ker\pi$, satisfies $\chi(g)^3=1$ or $\chi(g)^4=1$. One of the two holds for all $g$: if $\chi(g)^3\ne1$ and $\chi(g')^4\ne1$, then $\chi(g)^4=1$ and $\chi(g')^3=1$, and neither holds at $gg'$, since $\chi(gg')^3=\chi(g)^3$ and $\chi(gg')^4=\chi(g')^4$. So $\chi^3=1$ or $\chi^4=1$, contrary to the hypothesis.
\end{proof}
\leanref{Development/AutomorphicArithmetic/AutomorphicVarieties/WeakShimuraTheory/StabilizerIndexBounds/NonComm.lean}{UpperTheory.V.Ws.exists_noncommuting_images_in_finiteIndex_normal_subgroup}
\leanref{Development/AutomorphicArithmetic/AutomorphicVarieties/AlgebraTools/MonomialRepresentations/PermChar.lean}{UpperTheory.V.AlgebraTools.permChar}

\begin{lemma}\label{sha-siegel-automorphy-zariski}
Let $K$ be a field, $G$ a group, $n\ge1$ and $\rho:G\to GL_n(K)$ a representation. For a polynomial $f$ over $K$ in $n^2$ variables and $g\in G$, write $f(g)$ for the value of $f$ at the entries of $\rho(g)$.
There are a normal subgroup $N$ of finite index of $G$ and a set $I$ of polynomials such that $N$ is the set of $g$ with $f(g)=0$ for all $f\in I$, and such that the following holds: for every subgroup $U$ of finite index of $G$, every $\sigma\in G$ and all polynomials $f_1,\dots,f_r$ such that, for each $a$, $f_a(\sigma g)\ne0$ for some $g\in N$, there is $g\in U\cap N$ with $f_a(\sigma g)\ne0$ for all $a$.
\end{lemma}
\begin{proof}
Give $G$ the topology whose closed sets are the sets $\{g:f(g)=0\text{ for all }f\in J\}$, for sets $J$ of polynomials. It is Noetherian, by the Hilbert basis theorem, and left and right translations are homeomorphisms, since $f(hg)$ and $f(gh)$ are polynomials in the entries of $\rho(g)$. So $G$ is the union of its finitely many irreducible components, which are closed.
\begin{enumerate}
\item[(1)] \emph{Every point lies in only one component.} The component $Z$ through $1$ is not contained in the union of the other components, since it is irreducible and contained in none of them; so some $x\in Z$ lies in no other component. For $y\in G$, left translation by $xy^{-1}$ permutes the components and carries $y$ to $x$, so $y$ too lies in only one component.
\item[(2)] \emph{The component $N$ through $1$.} For $a\in N$, the translates $aN$ and $a^{-1}N$ are components containing $a$ and $1$, so $aN=N=a^{-1}N$ by (1); hence $N$ is a subgroup. For $g\in G$, $gNg^{-1}$ is a component containing $1$, so $N$ is normal. The cosets $gN$ are components, so there are finitely many. Being closed, $N$ is the set of $g$ with $f(g)=0$ for all $f$ in some set $I$; and it is irreducible.
\item[(3)] \emph{$N$ lies in the closure of $U\cap N$.} For each of the finitely many left cosets $q$ of $U$ let $S_q=N\cap q$. Then $N$ is the union of the closures of the $S_q$, and as $N$ is irreducible, it lies in the closure of one of them, $S_q$, which is then nonempty. Let $m\in S_q$. Then $S_q\subseteq m(U\cap N)$, so for $y\in N$, $my\in N$ lies in the closure of $m(U\cap N)$, which is $m$ times the closure of $U\cap N$; so $y$ lies in the closure of $U\cap N$.
\item[(4)] \emph{Conclusion.} The sets $O_a=\{g:f_a(\sigma g)\ne0\}$ are open, as left translation by $\sigma$ is continuous, and each meets $N$. Since $N$ is irreducible, the open set $O=O_1\cap\dots\cap O_r$ meets $N$, say at $x$. By (3), $x$ lies in the closure of $U\cap N$, so its open neighbourhood $O$ contains some $g\in U\cap N$.
\end{enumerate}
\end{proof}
\leanref{Development/AutomorphicArithmetic/AutomorphicVarieties/AlgebraTools/MonomialRepresentations/IdCompMat.lean}{UpperTheory.V.AlgebraTools.idCompMat}
\leanref{Development/AutomorphicArithmetic/AutomorphicVarieties/AlgebraTools/MonomialRepresentations/IdComp.lean}{UpperTheory.V.AlgebraTools.idComp}

\begin{lemma}\label{sha-siegel-automorphy-weight}
Let $H$ be a commutative group, written multiplicatively, $n\ge1$ and $a_1,\dots,a_n\in H$ such that $a_i/a_j$ has infinite order for $i\ne j$. Then there are pairwise distinct integers $w_1,\dots,w_n$ such that $\sum_im_iw_i=0$ for every $(m_1,\dots,m_n)\in\mathbf Z^n$ with $\prod_ia_i^{m_i}=1$.
\end{lemma}
\begin{proof}
Let $L\subseteq\mathbf Z^n$ be the set of $m$ such that $\prod_ia_i^{km_i}=1$ for some $k\ne0$. It is a subgroup, and $Q=\mathbf Z^n/L$ is torsion-free, hence free of finite rank. Let $q_i\in Q$ be the image of the $i$-th basis vector. For $i\ne j$, $q_i\ne q_j$, since $q_i=q_j$ would give $(a_i/a_j)^k=1$ with $k\ne0$.
Some homomorphism $f:Q\to\mathbf Z$ has $f(q_i)\ne f(q_j)$ for all $i\ne j$. Otherwise $\operatorname{Hom}(Q,\mathbf Z)$ is the union of the finitely many subgroups $\{f:f(q_i-q_j)=0\}$, $i\ne j$, and one of them has finite index $d$ \cite{Neu54}; then $d\,f(q_i-q_j)=0$ for every $f$, so $q_i=q_j$ as $Q$ is free, a contradiction.
Put $w_i=f(q_i)$. If $\prod_ia_i^{m_i}=1$, then $m\in L$, so $\sum_im_iq_i=0$ and $\sum_im_iw_i=0$.
\end{proof}
\leanref{Development/AutomorphicArithmetic/AutomorphicVarieties/AlgebraTools/MonomialRepresentations/RelWeight.lean}{UpperTheory.V.AlgebraTools.relWeight}

\begin{lemma}\label{sha-siegel-automorphy-torus}
Let $K$ be a field, $n\ge1$, $X\in M_n(K)$, $a_1,\dots,a_n\in K^\times$ and $w_1,\dots,w_n\in\mathbf Z$ such that $\sum_im_iw_i=0$ for every $m\in\mathbf Z^n$ with $\prod_ia_i^{m_i}=1$, and let $T$ be an indeterminate. If the characteristic polynomial of $X\operatorname{diag}(T^{w_1},\dots,T^{w_n})$ over $K(T)$ is separable, then that of $X\operatorname{diag}(a_1^m,\dots,a_n^m)$ is separable for some $m\ge0$.
\end{lemma}
\begin{proof}
A monic polynomial of degree $n$ over a field is separable if and only if its resultant with its derivative, in degrees $n$ and $n$, is nonzero. Let $F\in K[t_1,\dots,t_n]$ be this resultant for the characteristic polynomial of $X\operatorname{diag}(t_1,\dots,t_n)$; its value at a point $(s_1,\dots,s_n)$ of a field containing $K$ is the resultant for $X\operatorname{diag}(s_1,\dots,s_n)$.
Suppose that $F(a_1^m,\dots,a_n^m)=0$ for all $m\ge0$. Write $F=\sum_\mu c_\mu t^\mu$, over finitely many $\mu\in\mathbf N^n$, and $a^\mu=\prod_ia_i^{\mu_i}$. For each value $\beta$ of $a^\mu$ let $C_\beta$ be the sum of the $c_\mu$ with $a^\mu=\beta$. Then $\sum_\beta C_\beta\beta^m=0$ for all $m\ge0$, so $\sum_\beta C_\beta P(\beta)=0$ for every polynomial $P$; taking $P$ vanishing at every value other than $\beta$ gives $C_\beta=0$.
If $a^\mu=a^{\mu'}$, then $\prod_ia_i^{\mu_i-\mu'_i}=1$, so $\sum_i\mu_iw_i=\sum_i\mu'_iw_i$; call this common value $e_\beta$, for $\beta=a^\mu$. Hence $F(T^{w_1},\dots,T^{w_n})=\sum_\beta C_\beta T^{e_\beta}=0$, which contradicts the separability over $K(T)$.
\end{proof}
\leanref{Development/AutomorphicArithmetic/AutomorphicVarieties/AlgebraTools/MonomialRepresentations/TorusTransfer.lean}{UpperTheory.V.AlgebraTools.torusTransfer}

\begin{lemma}\label{sha-siegel-automorphy-quartic}
Let $K$ be a field in which $2\ne0$, and $T$ an indeterminate.
\begin{enumerate}
\item Let $X\in M_4(K)$ and let $w_1>w_2>w_3>w_4$ be integers. If $X_{11}\ne0$ and the principal minor of $X$ on the rows and columns $1,2,3$ is nonzero, then the characteristic polynomial of $X\operatorname{diag}(T^{w_1},\dots,T^{w_4})$ over $K(T)$ is separable.
\item Let $B,C\in GL_2(K)$ with $B_{11}\ne0$ and $C_{11}\ne0$, and let $w_1>w_2$ and $w'_1>w'_2$ be integers. Then the characteristic polynomial over $K(T)$ of
\[
 \begin{pmatrix}0&B\\C&0\end{pmatrix}\operatorname{diag}(T^{w_1},T^{w_2},T^{w'_1},T^{w'_2})
\]
is separable.
\end{enumerate}
\end{lemma}
\begin{proof}
Measure elements of $K(T)$ by their degree in $T$, the negative of the valuation at infinity.

(1) Write the characteristic polynomial as $Y^4+bY^3+cY^2+dY+e$. Up to sign, its coefficients are the sums of the principal minors of $M=X\operatorname{diag}(T^{w_i})$ of sizes $1$ to $4$, and the minor of $M$ on a set $J$ of indices is that of $X$ times $T^{\sum_{i\in J}w_i}$. As the $w_i$ decrease, $b$ has degree $w_1$, $d$ has degree $w_1+w_2+w_3$, $e$ has degree at most $w_1+w_2+w_3+w_4$, and $c$ has degree $w_1+w_2$ if the principal minor of $X$ on $1,2$ is nonzero and at most $w_1+w_3$ otherwise. In the discriminant of $Y^4+bY^3+cY^2+dY+e$, the term $b^2c^2d^2$ in the first case, and $-4b^3d^3$ in the second, has larger degree than each of the other terms. So the discriminant is nonzero.

(2) With $C'=C\operatorname{diag}(T^{w_1},T^{w_2})$ and $B'=B\operatorname{diag}(T^{w'_1},T^{w'_2})$, the matrix is $\begin{pmatrix}0&B'\\C'&0\end{pmatrix}$, whose characteristic polynomial is $Y^4-sY^2+p$ with $s$ the trace and $p$ the determinant of $C'B'$. Such a polynomial is separable if $2\ne0$, $p\ne0$ and $s^2-4p\ne0$. Here $p=\det C\det B\,T^{w_1+w_2+w'_1+w'_2}\ne0$, and $s^2-4p$ is $T^{2(w_2+w'_2)}$ times $q^2-4\det C\det B\,T^{n+n'}$, where $n=w_1-w_2$, $n'=w'_1-w'_2$ and $q$ is a polynomial in $T$ of degree $n+n'$ with leading coefficient $C_{11}B_{11}$. So $s^2-4p\ne0$.
\end{proof}
\leanref{Development/AutomorphicArithmetic/AutomorphicVarieties/AlgebraTools/MonomialRepresentations/AdjugateCornerEqMinor.lean}{UpperTheory.V.AlgebraTools.disc4}
\leanref{Development/AutomorphicArithmetic/AutomorphicVarieties/AlgebraTools/MonomialRepresentations/DiscSwap.lean}{UpperTheory.V.AlgebraTools.discSwap}

\begin{lemma}\label{sha-siegel-automorphy-coset-four}
Let $K$ be an algebraically closed field in which $2\ne0$, $G$ a topological group, and $\rho:G\to GL_4(K)$ a homomorphism such that, for every polynomial $f$ over $K$ in $16$ variables, the set of $g$ at which $f$ vanishes on the entries of $\rho(g)$ is closed. Suppose that:
\begin{enumerate}
\item the values of $\rho$ span $M_4(K)$;
\item some $u\in G$ has $\prod_i(X-a_i)$ as the characteristic polynomial of $\rho(u)$, where no $a_i/a_j$ with $i\ne j$ is a root of unity;
\item every open normal subgroup of finite index of $G$ contains elements $g,h$ with $\rho(g)\rho(h)\ne\rho(h)\rho(g)$.
\end{enumerate}
Then for every subgroup $U$ of finite index of $G$ and every $\sigma\in G$ there is $g\in U$ such that the characteristic polynomial of $\rho(\sigma g)$ is separable.
\end{lemma}
\begin{proof}
The $a_i$ are distinct, so after conjugating $\rho$ by a fixed matrix, which changes neither the hypotheses nor the conclusion, $\rho(u)=\operatorname{diag}(a_1,\dots,a_4)$. Write $\rho(g)_{ij}$ for the entries of $\rho(g)$.
\begin{enumerate}
\item[(1)] \emph{A normal subgroup.} Let $N$ be as in Lemma~\ref{sha-siegel-automorphy-zariski}. It is closed, by the hypothesis on $\rho$, and of finite index, so it is open; by (3) it contains $g,h$ with $\rho(g)\rho(h)\ne\rho(h)\rho(g)$.
\item[(2)] \emph{A diagonal element of $U\cap N$.} As $U\cap N$ has finite index, $t=u^k$ lies in it for some $k\ge1$. Then $\rho(t)=\operatorname{diag}(d_1,\dots,d_4)$ with $d_i=a_i^k$, and no $d_i/d_j$ with $i\ne j$ is a root of unity; in particular the $d_i$ are distinct. Lemma~\ref{sha-siegel-automorphy-weight} gives distinct integers $w_1,\dots,w_4$ with $\sum_im_iw_i=0$ whenever $\prod_id_i^{m_i}=1$.
\item[(3)] \emph{A relation.} For $g\in G$ write $i\to_gj$ if $\rho(gn)_{ij}\ne0$ for some $n\in N$.
\begin{itemize}
\item[(a)] If $\rho(g)_{ij}\ne0$, then $i\to_gj$. If $i\to_gj$ and $g^{-1}g'\in N$, then $i\to_{g'}j$.
\item[(b)] If $i\to_gj$ and $j\to_hk$, then $i\to_{gh}k$. Indeed, let $n,n'\in N$ with $\rho(gn)_{ij}\ne0$ and $\rho(hn')_{jk}\ne0$. For $p\ge0$ the element $m_p=h^{-1}nt^phn'$ lies in $N$, and $ghm_p=gn\,t^p\,hn'$. If $\rho(ghm)_{ik}=0$ for all $m\in N$, then $\sum_l\rho(gn)_{il}d_l^p\rho(hn')_{lk}=0$ for all $p\ge0$, and as the $d_l$ are distinct, the Vandermonde determinant gives $\rho(gn)_{il}\rho(hn')_{lk}=0$ for each $l$, which fails for $l=j$.
\item[(c)] For all $i,j$ some $g$ has $\rho(g)_{ij}\ne0$, since the $\rho(g)$ span $M_4(K)$. For all $g$ and $i$ some $j$ has $\rho(g)_{ij}\ne0$ and $\rho(g^{-1})_{ji}\ne0$, since $\sum_j\rho(g)_{ij}\rho(g^{-1})_{ji}=1$.
\item[(d)] If $i\to_gj$, then $j\to_{g^{-1}}i$. Indeed, by (c) and (a) there is $h$ with $j\to_hi$; by (b), $i\to_{gh}i$, and then $j\to_{h(gh)^m}i$ for all $m\ge0$. With $e=[G:N]$, $(gh)^e\in N$, as $N$ is normal, and $g\cdot h(gh)^{e-1}=(gh)^e$; so (a) gives $j\to_{g^{-1}}i$.
\item[(e)] For each $i$ there is $j\ne i$ with $i\to_1j$. First, $j\to_1k$ for some $j\ne k$: otherwise every $\rho(n)$, $n\in N$, is diagonal, and these commute, contrary to (1). Suppose that $i\to_1l$ only for $l=i$. By (c) there is $g$ with $\rho(g)_{ji}\ne0$, so $j\to_gi$; by (d) and (b), $k\to_1j$ and $k\to_gi$. If $a\to_gi$ and $\rho(g)_{al}\ne0$, then $l=i$: by (d), (a) and (b), $i\to_{g^{-1}}a$ and $a\to_gl$ give $i\to_1l$. By (c) there is $l$ with $\rho(g)_{kl}\ne0$ and $\rho(g^{-1})_{lk}\ne0$, and $l=i$. Since $j\ne k$, $0=\sum_b\rho(g)_{jb}\rho(g^{-1})_{bk}$, and every term with $b\ne i$ vanishes, as $\rho(g)_{jb}\ne0$ forces $b=i$; so $\rho(g)_{ji}\rho(g^{-1})_{ik}=0$, a contradiction.
\end{itemize}
So $\to_1$ is an equivalence relation, by (a), (d) and (b), whose classes have at least two elements.
\item[(4)] \emph{The pattern of $\sigma$.} Either $i\to_\sigma i$ for all $i$, or $\{1,2,3,4\}$ is the disjoint union of two pairs $P$ and $Q$ such that $i\to_\sigma j$ if and only if $i$ and $j$ lie in different pairs.
Suppose that $i_0\not\to_\sigma i_0$. By (b), (d) and (c): if $i\to_\sigma j$ and $j\to_1k$, or $i\to_1j$ and $j\to_\sigma k$, then $i\to_\sigma k$; if $i\to_\sigma k$ and $j\to_\sigma k$, then $i\to_1j$; and each $i$ has some $j$ with $i\to_\sigma j$.
By (e) take $i_1\ne i_0$ with $i_0\to_1i_1$, then $j_0$ with $i_0\to_\sigma j_0$, and $j_1\ne j_0$ with $j_0\to_1j_1$. If $i_0\to_1j_0$, then $i_0\to_\sigma i_0$; so $i_0\not\to_1j_0$, and $P=\{i_0,i_1\}$ and $Q=\{j_0,j_1\}$ are disjoint, since a common element would give $i_0\to_1j_0$.
No two elements of $P$ are related by $\to_\sigma$, since that would give $i_0\to_\sigma i_0$. No two elements of $Q$ are, since that would give $j_0\to_\sigma j_0$ and then, with $i_0\to_\sigma j_0$, $i_0\to_1j_0$. Each $i\in P$ is related to each $j\in Q$, through $i\to_1i_0$, $i_0\to_\sigma j_0$ and $j_0\to_1j$. For $j\in Q$ some $l$ has $j\to_\sigma l$, and $l\in P$; then $j\to_\sigma i$ for every $i\in P$, as $l\to_1i$.
\item[(5)] \emph{A separable twist.} There is $g\in U\cap N$ such that the characteristic polynomial of $\rho(\sigma g)\operatorname{diag}(T^{w_1},\dots,T^{w_4})$ over $K(T)$ is separable.
In the first case of (4), for each $i$, $\rho(\sigma n)_{ii}\ne0$ for some $n\in N$; and by (d), $\rho(\sigma^{-1}n')_{ii}\ne0$ for some $n'\in N$, so that for $n''=\sigma^{-1}n'^{-1}\sigma\in N$ the adjugate $\det\rho(\sigma n'')\,\rho(\sigma^{-1}n')$ of $\rho(\sigma n'')$ has nonzero entry $(i,i)$. Lemma~\ref{sha-siegel-automorphy-zariski}, with $\sigma$ and the eight polynomials giving the diagonal entries of a matrix and of its adjugate, gives $g\in U\cap N$ such that the diagonal entries of $\rho(\sigma g)$ and of its adjugate, which are its principal minors of size $3$, are nonzero. After renumbering so that the $w_i$ decrease, which does not change the characteristic polynomial, Lemma~\ref{sha-siegel-automorphy-quartic}(1) applies.
In the second case, Lemma~\ref{sha-siegel-automorphy-zariski}, with $\sigma$ and the polynomials giving the entries $(i,j)$ with $i\to_\sigma j$, gives $g\in U\cap N$ with $\rho(\sigma g)_{ij}\ne0$ whenever $i\to_\sigma j$; and $\rho(\sigma g)_{ij}=0$ otherwise, since $g\in N$. After renumbering so that $P=\{1,2\}$ and $Q=\{3,4\}$, with the $w_i$ decreasing on each pair, $\rho(\sigma g)=\begin{pmatrix}0&B\\C&0\end{pmatrix}$ with $B$ and $C$ invertible, as $\rho(\sigma g)$ is, and with nonzero entries; Lemma~\ref{sha-siegel-automorphy-quartic}(2) applies.
\item[(6)] \emph{Conclusion.} By (2), (5) and Lemma~\ref{sha-siegel-automorphy-torus}, with $X=\rho(\sigma g)$ and the $d_i$, there is $m\ge0$ such that the characteristic polynomial of $\rho(\sigma g)\operatorname{diag}(d_1^m,\dots,d_4^m)=\rho(\sigma gt^m)$ is separable, and $gt^m\in U$.
\end{enumerate}
\end{proof}
\leanref{Development/AutomorphicArithmetic/AutomorphicVarieties/AlgebraTools/MonomialRepresentations/RegCosetFour.lean}{UpperTheory.V.AlgebraTools.regCosetFour}

\begin{lemma}\label{sha-siegel-automorphy-link-self-twist}
Let $K$ be an algebraically closed field, $G$ a group, $V$ a finite-dimensional $K$-vector space and $\rho:G\to GL(V)$ an irreducible representation. Let $N$ be a normal subgroup of $G$ of finite index $d$, with $d\ne0$ in $K$, that contains every commutator $aba^{-1}b^{-1}$, and suppose that the restriction of $\rho$ to $N$ is reducible.
Then there are a character $\phi:G\to K^\times$ with $\phi\ne1$ and $N\subseteq\ker\phi$, and a linear automorphism $y$ of $V$ with $\rho(g)y=\phi(g)y\rho(g)$ for all $g\in G$.
\end{lemma}
\begin{proof}
Call a subspace of $V$ \emph{$N$-stable} if $\rho(h)$ preserves it for every $h\in N$. Let $G$ act on $\operatorname{End}(V)$ by $g\cdot f=\rho(g)f\rho(g)^{-1}$.
\begin{enumerate}
\item[(1)] \emph{An $N$-stable complement.} Let $W$ be a nonzero proper $N$-stable subspace, $W'$ maximal among the $N$-stable subspaces with $W\cap W'=0$, and $S$ minimal among the nonzero $N$-stable subspaces. Then $\rho(g)S\subseteq W+W'$ for every $g\in G$. Indeed, $S\cap\rho(g)^{-1}(W+W')$ is $N$-stable, as $N$ is normal. If it is nonzero, it is $S$ by minimality. If it is $0$, then $W'+\rho(g)S$ is $N$-stable and meets $W$ in $0$: if $x=x'+\rho(g)s\in W$ with $x'\in W'$ and $s\in S$, then $\rho(g)s\in W+W'$, so $s=0$ and $x\in W\cap W'=0$; by maximality, $\rho(g)S\subseteq W'$. So the intersection of the $\rho(g)^{-1}(W+W')$ over $g\in G$ is a $G$-stable subspace containing $S\ne0$; it is $V$, and $V=W\oplus W'$.
\item[(2)] \emph{A projection.} Let $p$ be the projection onto $W$ along $W'$. It is fixed by $N$. Some $g_0\in G$ has $g_0\cdot p\ne p$: otherwise $p$ commutes with $\rho(G)$, and $W=p(V)$ is $G$-stable, which contradicts the irreducibility of $\rho$.
\item[(3)] \emph{A common eigenvector.} Let $E\subseteq\operatorname{End}(V)$ be the subspace fixed by $N$. It is stable under $G$, as $N$ is normal, and on $E$ the actions of any $a,b\in G$ commute, since $ab=ba\,(a^{-1}b^{-1}ab)$ and $a^{-1}b^{-1}ab\in N$. Let $U=\{g_0\cdot x-x:x\in E\}\subseteq E$; it contains $g_0\cdot p-p\ne0$, and it is stable under $G$, whose elements commute with $g_0$ on $E$. A minimal nonzero $G$-stable subspace of $U$ is contained in an eigenspace of each $g\in G$, since that eigenspace meets it in a nonzero $G$-stable subspace, $K$ being algebraically closed. So some nonzero $y\in U$ has $g\cdot y=\phi(g)y$ for all $g$, with $\phi(g)\in K^\times$. Then $\phi$ is a character, trivial on $N$ since $y\in E$.
\item[(4)] \emph{$\phi\ne1$.} Otherwise $g_0\cdot y=y$. Write $y=g_0\cdot x-x$ with $x\in E$; then $g_0^j\cdot x=x+jy$ for $j\ge0$. As $N$ is normal of index $d$, $g_0^d\in N$, so $g_0^d\cdot x=x$ and $dy=0$, which is impossible as $d\ne0$ in $K$.
\item[(5)] \emph{The automorphism.} The relation $g\cdot y=\phi(g)y$ says $\rho(g)y=\phi(g)y\rho(g)$. So the kernel of $y$ is $G$-stable; it is not $V$, as $y\ne0$, so it is $0$, and $y$ is an automorphism.
\end{enumerate}
\end{proof}
\leanref{Development/AutomorphicArithmetic/AutomorphicVarieties/CurveDVRSpecialization/SelfTwistEigenvalues/SelfTwist.lean}{UpperTheory.V.Pr.T.selfTwist}
\leanref{Development/AutomorphicArithmetic/AutomorphicVarieties/CurveDVRSpecialization/SelfTwistEigenvalues/CliffordSS.lean}{UpperTheory.V.Pr.T.cliffordSS}

\begin{lemma}\label{sha-siegel-automorphy-rank-four}
Let $\ell$ be an odd prime, $\Gamma_0$ a profinite group, $\varepsilon:\Gamma_0\to\mathbf Z_\ell^\times$ a continuous homomorphism with kernel $\Gamma'$, and $k$ a field of characteristic $\ell$. Let $\rho:\Gamma_0\to GL_4(\bar{\mathbf Q}_\ell)$ be a continuous irreducible representation, $A$ a commutative ring with an injective homomorphism $j:A\to\bar{\mathbf Q}_\ell$ and a homomorphism $\phi:A\to k$, and $\bar\chi_1,\dots,\bar\chi_4:\Gamma_0\to k^\times$ characters such that for every $g$ the characteristic polynomial of $\rho(g)$ is the image under $j$ of a polynomial over $A$ whose image under $\phi$ is $\prod_i(X-\bar\chi_i(g))$. Suppose that:
\begin{enumerate}
\item some $u\in\Gamma'$ has $\prod_i(X-a_i)$ as the characteristic polynomial of $\rho(u)$, with $a_i^e\ne a_j^e$ for $i\ne j$ and $e\ge1$;
\item every character $\eta:\Gamma_0\to k^\times$ with $\eta^2=1$, trivial on the kernel of $\varepsilon\bmod\ell$, such that each $\eta\bar\chi_i$ is one of the $\bar\chi_j$, is trivial;
\item $(\bar\chi_s/\bar\chi_t)^e\ne1$ for some $s,t$ and all $1\le e\le4$.
\end{enumerate}
Then the restriction of $\rho$ to $\Gamma'$ is irreducible, and every open subset of $\Gamma_0$ that meets $\Gamma'$ contains an element $\sigma$ of $\Gamma'$ at which the characteristic polynomial of $\rho(\sigma)$ is separable.
\end{lemma}
\begin{proof}
No $a_i/a_j$ with $i\ne j$ is a root of unity, so for every $d\ge1$ the eigenvalues $a_i^d$ of $\rho(u)^d$ are distinct and its characteristic polynomial is separable. Every commutator of $\Gamma_0$ lies in $\Gamma'$, as $\mathbf Z_\ell^\times$ is commutative; so $\Gamma'$ is normal.
\begin{enumerate}
\item[(1)] \emph{Irreducibility on $\Gamma'$.} Suppose that $\rho$ is reducible on $\Gamma'$. Lemma~\ref{sha-siegel-automorphy-stabilizer}, with $N=\Gamma'$ and $u$, gives a nonzero proper subspace $W$, stable under $\rho(\Gamma')$, whose stabilizer $S\supseteq\Gamma'$ has index dividing $4$. The subgroup $S$ is closed, as the intersection of the sets of $g$ with $f(\rho(g)w)=0$ for $w\in W$ and linear forms $f$ vanishing on $W$; having finite index, it is open. It contains every commutator, and $\rho$ is reducible on $S$, as $W$ shows.
So $S$ is normal, and its index $d$ divides $4$, so $d\ne0$ in $\bar{\mathbf Q}_\ell$. By Lemma~\ref{sha-siegel-automorphy-link-self-twist}, with $K=\bar{\mathbf Q}_\ell$, $G=\Gamma_0$, $V=\bar{\mathbf Q}_\ell^4$ and $N=S$, there are a character $\lambda\ne1$ of $\Gamma_0$, trivial on $S$, and an invertible matrix $M$ with $M\lambda(g)\rho(g)=\rho(g)M$ for all $g$. As $S$ is normal of index $d$, $g^d\in S$ for every $g$, so $\lambda(g)^d=1$ and $\lambda^4=1$. Let $\eta=\lambda$ if $\lambda^2=1$, and $\eta=\lambda^2$, with $M^2$ in place of $M$, otherwise. Then $\eta$ takes values in $\{\pm1\}$, $\eta\ne1$, $S\subseteq\ker\eta$, and $\eta(g)\rho(g)$ is conjugate to $\rho(g)$, so the two have the same characteristic polynomial.
The kernel of $\eta$ contains the open subgroup $S$, so it is open, and it contains $\Gamma'=\ker\varepsilon$. By Lemma~\ref{sha-siegel-automorphy-sign}, $\eta(g)=1$ whenever $\varepsilon(g)\equiv1\pmod\ell$. Let $\bar\eta:\Gamma_0\to k^\times$ be the composite of $\eta$ with $\{\pm1\}\to k^\times$. Then $\bar\eta^2=1$, $\bar\eta$ is trivial on the kernel of $\varepsilon\bmod\ell$, and by Lemma~\ref{sha-siegel-automorphy-sign-twist} each $\bar\eta\bar\chi_i$ is one of the $\bar\chi_j$. By (2), $\bar\eta=1$. As $\ell$ is odd, $-1\ne1$ in $k$, so $\eta=1$, a contradiction.
\item[(2)] \emph{Noncommuting elements.} By Burnside's theorem \cite[\S27]{CR62}, the values of $\rho$ span $M_4(\bar{\mathbf Q}_\ell)$, as $\rho$ is irreducible, and so do its values on $\Gamma'$, by (1). By Lemma~\ref{sha-siegel-automorphy-noncommuting}, with $G=\Gamma_0$, $Z=\Gamma'$, $u$ and hypothesis (3), every normal subgroup of finite index of $\Gamma'$ contains $g,h$ with $\rho(g)\rho(h)\ne\rho(h)\rho(g)$.
\item[(3)] \emph{Separable elements.} Let $O$ be an open subset of $\Gamma_0$ and $\sigma\in O\cap\Gamma'$. As $\Gamma_0$ is profinite, there is an open normal subgroup $V$ of $\Gamma_0$ with $\sigma V\subseteq O$, and $V\cap\Gamma'$ has finite index in $\Gamma'$. Lemma~\ref{sha-siegel-automorphy-coset-four} applies to the restriction of $\rho$ to the topological group $\Gamma'$: the zero sets of polynomials in the entries are closed as $\rho$ is continuous, the values span $M_4(\bar{\mathbf Q}_\ell)$ by (2), the element $u$ satisfies (2) of that lemma by hypothesis (1), and (3) of that lemma holds by (2). With $U=V\cap\Gamma'$, it gives $g\in V\cap\Gamma'$ such that $\rho(\sigma g)$ has separable characteristic polynomial, and $\sigma g\in O\cap\Gamma'$.
\end{enumerate}
% The Lean proof is the private absolute_irreducible_rank_four of TNiceFour.lean, called by Th.nice0four_alt. (1) is TypeI: Ws.exists_stabilizer_index_dvd_of_not_irreducible_restriction, Ws.isClosed_setOf_matrix_stabilizes_submodule with isOpen_of_isClosed_of_finiteIndex, Ws.exists_sign_twist_charpoly_of_not_irreducible_restriction (through Pr.T.selfTwist, Lemma sha-siegel-automorphy-link-self-twist, which needs the index of S invertible in K, with lambda^4 = 1 from the index dividing 4 by S.pow_index_mem, and the private quadratic_twist_charpoly_eq for eta), Ws.sign_char_eq_one_of_cyclotomic_sub_one_not_isUnit and Th.Det.perm4. The separability of the powers is Ws.charpoly_pow_separable_of_eigenvalue_ratios_not_finOrder. (2) is KStar.Schema.span_eq_top_of_irreducible and Ws.exists_noncommuting_images_in_finiteIndex_normal_subgroup; (3) is Ws.exists_mem_open_inter_subgroup_of_cosetwise_finiteIndex with AlgebraTools.regCosetFour.
\end{proof}
\leanref{Development/AutomorphicArithmetic/AutomorphicVarieties/ModularityArgument/ThresholdArguments/DenseThresholdCases/TNiceFour.lean}{UpperTheory.V.Th.nice0four_alt}
\leanref{Development/AutomorphicArithmetic/AutomorphicVarieties/WeakShimuraTheory/StabilizerIndexBounds/ExistsSignTwistCharpolyOfNotIrreducibleRestricti.lean}{UpperTheory.V.Ws.exists_sign_twist_charpoly_of_not_irreducible_restriction}
\leanref{Development/AutomorphicArithmetic/AutomorphicVarieties/WeakShimuraTheory/StabilizerIndexBounds/ExistsMemOpenInterSubgroupOfCosetwiseFiniteIndex.lean}{UpperTheory.V.Ws.exists_mem_open_inter_subgroup_of_cosetwise_finiteIndex}

\begin{lemma}\label{sha-siegel-automorphy-link-root}
Let $\ell$ be an odd prime, $|\cdot|$ the absolute value of $\bar{\mathbf Q}_\ell$, $\Gamma$ a compact group, $\pi:\Gamma\to\mathbf Z_\ell^\times$ a continuous homomorphism with $\pi(s)\equiv1\pmod\ell$ for all $s\in\Gamma$, and $n\ge1$.
Then there are a continuous homomorphism $t$ from $\Gamma$ to the additive group $\mathbf Z_\ell$ with the same kernel as $\pi$, a ring $\mathcal O_1$ of Definition~\ref{sha-character-parts}, and a continuous homomorphism $\beta:\Gamma\to\bar{\mathbf Q}_\ell^\times$ with values in $\mathcal O_1$, such that $\beta$ is trivial on the kernel of $t$, $|\beta(s)-1|<1$ and $\beta(s)^n=\pi(s)$ for all $s$, and $\beta(s)\beta(s')=1$ whenever $\pi(s)\pi(s')=1$.
\end{lemma}
\begin{proof}
\begin{enumerate}
\item[(1)] \emph{A logarithm.} Since $\ell$ is odd, $(1+\ell)^{\ell^j}\equiv1+\ell^{j+1}\pmod{\ell^{j+2}}$ for every $j\ge0$, so $x\mapsto(1+\ell)^x$ is an isomorphism of topological groups from $\mathbf Z_\ell$ onto $1+\ell\mathbf Z_\ell$; let $e$ be its inverse. Put $t=e\circ\pi$, a continuous homomorphism with kernel $\ker\pi$.
\item[(2)] \emph{A root.} Let $u=1+\ell$. The product of the $1-b$ over the roots $b$ of $X^n-u$ in $\bar{\mathbf Q}_\ell$, counted with multiplicity, is $1-u$, and $|1-u|<1$; so some root $b$ has $|b-1|<1$.
\item[(3)] \emph{Its powers.} As $|b-1|<1$, the series $b^x=\sum_{j\ge0}\binom xj(b-1)^j$ converges for $x\in\mathbf Z_\ell$ and defines a continuous homomorphism from $\mathbf Z_\ell$ to the units of the ring of integers of $\mathbf Q_\ell(b)$ that extends $x\mapsto b^x$ on $\mathbf N$. Since $|b^x-1|\le|b-1|$ for $x\in\mathbf N$, which is dense, $|b^x-1|<1$ for all $x$. Let $\mathcal O_1$ be a ring of Definition~\ref{sha-character-parts} that contains this ring of integers.
\item[(4)] \emph{The character.} Put $\beta(s)=b^{t(s)}$. It is a continuous homomorphism with values in $\mathcal O_1$, trivial on $\ker t$, with $|\beta(s)-1|<1$. The continuous homomorphisms $x\mapsto b^{nx}$ and $x\mapsto(1+\ell)^x$ from $\mathbf Z_\ell$ agree at $1$, hence on the dense subgroup $\mathbf Z$, hence everywhere; so $\beta(s)^n=(1+\ell)^{t(s)}=\pi(s)$. If $\pi(s)\pi(s')=1$, then $t(ss')=0$, so $\beta(s)\beta(s')=\beta(ss')=1$.
\end{enumerate}
% (1) is rc_equiv (through princZl, which gives some isomorphism of 1 + l Z_l with Z_l); (2) is rc_root; (3) is rc_char (PadicInt.addChar_of_value_at_one, rc_coef) with rc_pow.
\end{proof}
\leanref{Development/AutomorphicArithmetic/AutomorphicVarieties/WeakShimuraTheory/FrobeniusDensity/RootChar.lean}{UpperTheory.V.Ws.Tw.rootChar}

\begin{lemma}\label{sha-siegel-automorphy-link-half}
Let $D$ be a Siegel datum, with the notation of Definition~\ref{sha-siegel-datum}, and let $|\cdot|$ be the absolute value of $\bar{\mathbf Q}_\ell$.
Then there is a continuous homomorphism $\gamma:G\to\bar{\mathbf Q}_\ell^\times$ such that $|\gamma(g)-1|<1$ and $\gamma(g)^2\omega(g)=\chi_{\mathrm{cyc}}(g)$ for all $g\in G$, and $\gamma(g)=1$ whenever $\chi_{\mathrm{cyc}}(g)=1$.
\end{lemma}
\begin{proof}
The character $\pi=\chi_{\mathrm{cyc}}^{\ell-1}$ is continuous with values in $1+\ell\mathbf Z_\ell$, and its kernel contains that of $\chi_{\mathrm{cyc}}$. Lemma~\ref{sha-siegel-automorphy-link-root}, with $\Gamma=G$ and $n=2(\ell-1)$, gives $t$ and a continuous $\gamma$ with $|\gamma(g)-1|<1$ and $\gamma^{2(\ell-1)}=\pi$, trivial on $\ker t=\ker\pi\supseteq\ker\chi_{\mathrm{cyc}}$.
Let $g\in G$ and $z=\gamma(g)^2\omega(g)\chi_{\mathrm{cyc}}(g)^{-1}$. Then $z^{\ell-1}=1$, since $\omega^{\ell-1}=1$. Also $|z-1|<1$: $|\gamma(g)^2-1|<1$, and $|\omega(g)-\chi_{\mathrm{cyc}}(g)|<1$ because $\bar\omega=\bar\chi_{\mathrm{cyc}}$. So $|1+z+\dots+z^{\ell-2}-(\ell-1)|<1$, and as $|\ell-1|=1$, $1+z+\dots+z^{\ell-2}\ne0$. It kills $z-1$, since $z^{\ell-1}=1$; so $z=1$.
% Tw.rootChar with pi = cycZ^(l-1), n = 2(l-1); the last step is sq_mul_eq_of_pow_two_mul_pred_eq_of_norm_sub_lt_one with padicAlgCl_eq_one_of_pow_eq_one_of_norm_sub_one_lt_one, using om_pow and om_bar of D.Good.
\end{proof}
\leanref{Development/AutomorphicArithmetic/AutomorphicVarieties/WeakShimuraTheory/UnipotentIdempotents/HalfEx.lean}{UpperTheory.V.Ws.Lk.halfEx}

\begin{lemma}\label{sha-siegel-automorphy-link-hecke}
Let $\ell$ be an odd prime, write $G=\operatorname{Gal}(\bar{\mathbf Q}/\mathbf Q)$ with the notation $G_E$, $D_B(E)$, $I_B(E)$ of Definition~\ref{sha-siegel-tower-towers}, let $K\subset\bar{\mathbf Q}$ be an imaginary quadratic field, and let $A$ be a place of $\bar{\mathbf Q}$ above $\ell$ with $D_A(\mathbf Q)\subseteq G_K$, above the prime $w$ of $K$.
Then there are $n\ge1$ and a continuous homomorphism $\mu:G_K\to\mathbf Z_\ell^\times$ such that:
\begin{enumerate}
\item $\mu(s)\equiv1\pmod\ell$ for all $s\in G_K$;
\item $\mu(s)\mu(csc^{-1})=\chi_{\mathrm{cyc}}(s)^n$ for all $s\in G_K$ and $c\in G\setminus G_K$;
\item $\mu$ is trivial on $I_B(K)$ for every place $B$ of $\bar{\mathbf Q}$ that does not lie above $w$.
\end{enumerate}
\end{lemma}
\begin{proof}
Let $\tau$ be the nontrivial automorphism of $K$.
\begin{enumerate}
\item[(1)] \emph{Units, classes and the completion at $w$.} The unit group of $\mathcal O_K$ is finite, of order $m$ say, and the class group of $K$ has order $h$. As $D_A(\mathbf Q)\subseteq G_K$, $\ell$ splits in $K$, so $\mathcal O_K/w^j\cong\mathbf Z/\ell^j$ for every $j$, and these isomorphisms give a ring homomorphism $\iota:\mathcal O_K\to\mathbf Z_\ell$ with $\iota(x)\in\ell^j\mathbf Z_\ell$ if and only if $x\in w^j$.
\item[(2)] \emph{A character of ideals.} Let $I$ be a nonzero ideal of $\mathcal O_K$ prime to $w$, and write $I^h=(a)$; then $\iota(a)$ is a unit, and we put $F(I)=\iota(a)^{m(\ell-1)}$. This does not depend on $a$, which is determined up to a unit $u$, and $\iota(u)^m=1$. So $F$ is multiplicative, $F(I)\equiv1\pmod\ell$, and $F((a))=\iota(a)^n$ for $a$ prime to $w$, with $n=hm(\ell-1)$.
\item[(3)] \emph{Ray class characters.} Let $j\ge0$. If $a\equiv1\pmod{w^{j+1}}$, then $\iota(a)\equiv1\pmod{\ell^{j+1}}$ and $F((a))\equiv1\pmod{\ell^j}$. By class field theory \cite[Chapter~VI]{Neu99}, there is a homomorphism $\kappa_j:G_K\to(\mathbf Z/\ell^j)^\times$ with open kernel that is trivial on $I_B(K)$ for every place $B$ above a prime $v\ne w$ of $K$ and maps every Frobenius element at such a place $B$ to $F(v)\bmod\ell^j$.
\item[(4)] \emph{The limit.} By the Chebotarev density theorem (Theorem~\ref{lit-chebotarev-degree-one}), the Frobenius elements at the places above the primes of $K$ prime to $\ell$ are dense in $G_K$. For $j'\le j$, the reduction of $\kappa_j$ modulo $\ell^{j'}$ and $\kappa_{j'}$ are locally constant and agree at these elements, so they are equal. Hence the $\kappa_j$ define a continuous homomorphism $\mu:G_K\to\mathbf Z_\ell^\times$ with $\mu\equiv\kappa_j\pmod{\ell^j}$ for all $j$, and $\mu(\sigma)=F(v)$ for every Frobenius element $\sigma$ at a place above a prime $v\ne w$.
\item[(5)] \emph{Property (1).} It holds at the Frobenius elements of (4), by (2), and the set where it holds is closed.
\item[(6)] \emph{Property (2).} Fix $c\in G\setminus G_K$. Both sides are continuous in $s$. Let $\sigma$ be a Frobenius element at a place $B$ above a prime $v$ of $K$ prime to $\ell$. Then $c\sigma c^{-1}$ is a Frobenius element at $cB$, which lies above $\tau(v)$, also prime to $\ell$; and $v\,\tau(v)=(Nv)$. So $\mu(\sigma)\mu(c\sigma c^{-1})=F(v\,\tau(v))=\iota(Nv)^n=(Nv)^n=\chi_{\mathrm{cyc}}(\sigma)^n$. By the density in (4), (2) holds everywhere.
\item[(7)] \emph{Property (3).} If $B$ lies above a prime $v\ne w$ of $K$, then every $\kappa_j$ is trivial on $I_B(K)$, so $\mu$ is.
\end{enumerate}
% (1) is unit_pow, class_pow, Tw.quadFacts (the split prime w with O_K -> Z/l^n) and exists_iota; (2) is exists_F; (3) is Tw.rayChar (RayChar.lean, about 180 lines, through the Artin map Modulus.rayClassArtinEquiv of the ray class field), for the modulus w^(k+1); (4) is Tw.frobDense (Published.ChebotarevDegreeOne) and exists_limit; (5) and (6) use Continuous.ext_on and the density, with exists_conj, the conjugate Frobenius and I tau(I) = (N I) of Tw.quadFacts, and cycZ_of_isFrob.
\end{proof}
\leanref{Development/AutomorphicArithmetic/AutomorphicVarieties/WeakShimuraTheory/FrobeniusDensity/MuChar.lean}{UpperTheory.V.Ws.Tw.muChar}

\begin{lemma}\label{sha-siegel-automorphy-link-beta}
Let $\ell$ be an odd prime, $|\cdot|$ the absolute value of $\bar{\mathbf Q}_\ell$, $G=\operatorname{Gal}(\bar{\mathbf Q}/\mathbf Q)$ with the notation of Definition~\ref{sha-siegel-tower-towers}, $K\subset\bar{\mathbf Q}$ an imaginary quadratic field, and $A$ a place of $\bar{\mathbf Q}$ above $\ell$ with $D_A(\mathbf Q)\subseteq G_K$.
Then there are a continuous homomorphism $t$ from $G_K$ to the additive group $\mathbf Z_\ell$, a ring $\mathcal O_1$ of Definition~\ref{sha-character-parts}, and a continuous homomorphism $\beta:G_K\to\bar{\mathbf Q}_\ell^\times$ with values in $\mathcal O_1$, such that:
\begin{enumerate}
\item $\beta$ is trivial on the kernel of $t$, and $|\beta(s)-1|<1$ for all $s$;
\item $\beta(s)\beta(csc^{-1})=1$ and $t(s)+t(csc^{-1})=0$ for $s\in G_K$ and $c\in G\setminus G_K$;
\item $t$ is trivial on $I_B(K)$ for every place $B$ of $\bar{\mathbf Q}$ not above $\ell$;
\item for every place $B$ of $\bar{\mathbf Q}$ above $\ell$ there are an odd integer $b$ and $k\ge1$ with $\big(\beta(s)^2\chi_{\mathrm{cyc}}(s)^{-b}\big)^k=1$ for $s\in I_B(K)$.
\end{enumerate}
\end{lemma}
\begin{proof}
Let $w$ be the prime of $K$ below $A$, and let $n$ and $\mu$ be as in Lemma~\ref{sha-siegel-automorphy-link-hecke}. For $s\in G_K$ and $c\in G$, $\chi_{\mathrm{cyc}}(csc^{-1})=\chi_{\mathrm{cyc}}(s)$.
\begin{enumerate}
\item[(1)] \emph{A character in $1+\ell\mathbf Z_\ell$.} Let $\pi=\mu^2\chi_{\mathrm{cyc}}^{-n}$ on $G_K$, a continuous homomorphism to $\mathbf Z_\ell^\times$. For $c\in G\setminus G_K$, by (2) of Lemma~\ref{sha-siegel-automorphy-link-hecke}, $\pi(s)=\mu(s)\mu(csc^{-1})^{-1}$, which is $\equiv1\pmod\ell$ by (1) of that lemma, and $\pi(s)\pi(csc^{-1})=\big(\mu(s)\mu(csc^{-1})\big)^2\chi_{\mathrm{cyc}}(s)^{-2n}=1$.
\item[(2)] \emph{The characters.} Lemma~\ref{sha-siegel-automorphy-link-root}, applied to $\pi$ and $2n$, gives $t$, $\mathcal O_1$ and $\beta$ with $\ker t=\ker\pi$ and $\beta^{2n}=\pi$. This gives (1), and (2) follows from $\pi(s)\pi(csc^{-1})=1$.
\item[(3)] \emph{Property (3).} If $B$ is not above $\ell$, it is not above $w$, so $\mu$ is trivial on $I_B(K)$, and so is $\chi_{\mathrm{cyc}}$; hence $\pi$ and $t$ are.
\item[(4)] \emph{Above $w$.} The function $f(s)=\beta(s)^2\chi_{\mathrm{cyc}}(s)^{-1}$ is continuous with $f(1)=1$, and the $n$-th roots of unity other than $1$ form a finite set; so there is an open subgroup $U$ of $G_K$ such that $f(s)=1$ whenever $s\in U$ and $f(s)^n=1$. Fix $c\in G\setminus G_K$. Let $B$ be a place above $w$ and $s\in U\cap I_B(K)$. Then $c^{-1}sc\in I_{c^{-1}B}(K)$, and $c^{-1}B$ lies above $\tau(w)$, where $\tau$ is the nontrivial automorphism of $K$; since $D_A(\mathbf Q)\subseteq G_K$, $\ell$ splits in $K$ and $\tau(w)\ne w$. So $\mu(c^{-1}sc)=1$ by (3) of Lemma~\ref{sha-siegel-automorphy-link-hecke}, and $\mu(s)=\chi_{\mathrm{cyc}}(s)^n$ by (2) of that lemma. Hence $\pi(s)=\chi_{\mathrm{cyc}}(s)^n$, $f(s)^n=\beta(s)^{2n}\chi_{\mathrm{cyc}}(s)^{-n}=1$, and $\beta(s)^2=\chi_{\mathrm{cyc}}(s)$.
\item[(5)] \emph{Property (4).} Let $k=e!$, where $e$ is the index of $U$ in $G_K$, so that $s^k\in U$ for every $s\in G_K$. Let $B$ be a place above $\ell$, and $\sigma\in G$ with $B=\sigma A$. If $\sigma\in G_K$, then $B$ lies above $w$, and for $s\in I_B(K)$, (4) applied to $s^k$ gives $\big(\beta(s)^2\chi_{\mathrm{cyc}}(s)^{-1}\big)^k=1$; take $b=1$. If $\sigma\notin G_K$, then for $s\in I_B(K)$ the element $s'=\sigma^{-1}s\sigma$ lies in $I_A(K)$, $\beta(s)=\beta(s')^{-1}$ by (2), with $c=\sigma^{-1}$, and $\chi_{\mathrm{cyc}}(s)=\chi_{\mathrm{cyc}}(s')$; so $\big(\beta(s)^2\chi_{\mathrm{cyc}}(s)\big)^k=\big(\beta(s'^k)^2\chi_{\mathrm{cyc}}(s'^k)^{-1}\big)^{-1}=1$ by (4); take $b=-1$.
\end{enumerate}
% Tw.bt0: muChar, exists_lift (a lift of tau), unit_div_near, Tw.rootChar with 2n, the open subgroup from exists_open, and differs with conjGal_mem_Ine in (4); (5) is Lk.betaCyc, with k = (index of U)!.
\end{proof}
\leanref{Development/AutomorphicArithmetic/AutomorphicVarieties/WeakShimuraTheory/FrobeniusDensity/UnitDivNear.lean}{UpperTheory.V.Ws.Tw.bt0}
\leanref{Development/AutomorphicArithmetic/AutomorphicVarieties/WeakShimuraTheory/UnipotentIdempotents/BetaCyc.lean}{UpperTheory.V.Ws.Lk.betaCyc}

\begin{lemma}\label{sha-siegel-automorphy-link-dirichlet}
Let $D$ be a Siegel datum of case II or of case III (Definition~\ref{sha-siegel-tower-towers}), with the notation of Definition~\ref{sha-siegel-datum}.
Then there is a character $\eta:G\to\bar{\mathbf Q}_\ell^\times$ such that: $\eta$ is trivial on $\operatorname{Gal}(\bar{\mathbf Q}/\mathbf Q(\mu_{N'}))$ for some $N'\ge1$; $\eta^m=1$ for some $m$ prime to $\ell$; $\eta$ is trivial at every complex conjugation; $\eta$ is trivial on the inertia groups at the places not above the primes of $S$; and $\eta(g)\omega(g)^{k+3}\xi(g)\psi(g)=1$ for every $g\in G$ such that $\psi(g)^{m'}=1$ for some $m'$ prime to $\ell$.
\end{lemma}
\begin{proof}
\begin{enumerate}
\item[(1)] \emph{An exponent.} Let $D$ be attached to $(h,k,\psi)$. By (4) of Definition~\ref{sha-siegel-datum}, $\psi$ is a Dirichlet character modulo $N_h$, and $\xi=\xi_0\xi'$ with $\xi'$ a Dirichlet character modulo $2^{n'}$ and $\xi_0$ one modulo $p'$, where $(n',p')=(n_{\mathrm{II}},1)$ in case II and $(n_{\mathrm{III}},p_1(\psi_0))$ in case III. Write the order of $(\mathbf Z/N_h)^\times$ as $\ell^am_0$ with $\ell\nmid m_0$, and choose $j\ge0$ with $j\equiv0\pmod{\ell^a}$ and $j\equiv1\pmod{m_0}$. As $\psi$ factors through $(\mathbf Z/N_h)^\times$, $\psi^{\ell^am_0}=1$. If $\psi(g)^{m'}=1$ with $m'$ prime to $\ell$, the order of $\psi(g)$ divides $\gcd(\ell^am_0,m')$, which divides $m_0$, so $\psi(g)^j=\psi(g)$.
\item[(2)] \emph{The character.} Put $\eta=(\omega^{k+3}\xi\psi^j)^{-1}$. The character $\omega$ is trivial on $\operatorname{Gal}(\bar{\mathbf Q}/\mathbf Q(\mu_\ell))$, since $\omega(g)$ is the $(\ell-1)$-st root of unity with residue $\bar\chi_{\mathrm{cyc}}(g)$; so $\eta$ is trivial on $\operatorname{Gal}(\bar{\mathbf Q}/\mathbf Q(\mu_{N'}))$ with $N'=\ell N_hp'2^{n'}$.
\item[(3)] \emph{Order.} $\omega^{\ell-1}=1$, $\xi$ has some order $m_\xi$ prime to $\ell$, and $(\psi^j)^{m_0}=1$ since $\ell^am_0$ divides $jm_0$. So $\eta^m=1$ with $m=(\ell-1)m_\xi m_0$.
\item[(4)] \emph{Complex conjugations.} A complex conjugation $\sigma$ acts on every group of roots of unity as $c$ does, by inversion. So $\omega(\sigma)=\omega(c)=-1$, as $\omega(c)$ is an $(\ell-1)$-st root of unity with residue $\bar\chi_{\mathrm{cyc}}(c)=-1$; $\xi(\sigma)=\xi(c)=-1$; and $\psi(\sigma)=\psi(c)=(-1)^k$, since $\Psi(c)=\chi_{\mathrm{cyc}}(c)^k\psi(c)=1$. By (1), $\psi(\sigma)^j=\psi(\sigma)$. So $\eta(\sigma)^{-1}=(-1)^{k+3}(-1)(-1)^k=1$.
\item[(5)] \emph{Ramification.} The characters $\xi$ and $\psi$ are unramified outside $S$, and $\omega$ outside $\ell\in S$.
\item[(6)] \emph{The relation.} If $\psi(g)^{m'}=1$ with $m'$ prime to $\ell$, then $\psi(g)^j=\psi(g)$ by (1), so $\eta(g)\omega(g)^{k+3}\xi(g)\psi(g)=1$.
\end{enumerate}
% Lk.etaEx is the private exists_character_of_dirichlet_data with exp_ex, cls_one, omT and cyclotomicClass_eq_of_isComplexConj; it takes the moduli of psi and xi (N, p_1 and 2^n') from Aut.OfChar, which Ctx2 D or Ctx3 D provides, and eta is (om^(e+3) xi psi^k)^(-1) for the exponent k of exp_ex.
\end{proof}
\leanref{Development/AutomorphicArithmetic/AutomorphicVarieties/WeakShimuraTheory/UnipotentIdempotents/ExistsCharacterOfDirichletData.lean}{UpperTheory.V.Ws.Lk.etaEx}

\begin{lemma}\label{sha-siegel-automorphy-link-kappa-zero}
Let $\ell$ be an odd prime, write $G=\operatorname{Gal}(\bar{\mathbf Q}/\mathbf Q)$ with the notation of Definition~\ref{sha-siegel-tower-towers}, and let $K\subset\bar{\mathbf Q}$ be an imaginary quadratic field with $D_B(\mathbf Q)\subseteq G_K$ for the places $B$ above $2$.
Let $\eta:G\to\bar{\mathbf Q}_\ell^\times$ be a character, trivial on $\operatorname{Gal}(\bar{\mathbf Q}/\mathbf Q(\mu_N))$ for some $N\ge1$ and at every complex conjugation, with $\eta^m=1$ for some $m$ prime to $\ell$, such that at every place $B$ of $\bar{\mathbf Q}$ either $\eta$ is trivial on $I_B(\mathbf Q)$ or $D_B(\mathbf Q)\subseteq G_K$.
Then there is a character $\kappa_0:G_K\to\bar{\mathbf Q}_\ell^\times$ with open kernel and of order prime to $\ell$ such that $\kappa_0(s)\kappa_0(csc^{-1})=\eta(s)$ for $s\in G_K$ and $c\in G\setminus G_K$, and $\kappa_0$ is trivial on $I_B(K)$ at every place $B$ at which $\eta$ is trivial on $I_B(\mathbf Q)$.
\end{lemma}
\begin{proof}
Let $\tau$ be the nontrivial automorphism of $K$.
\begin{enumerate}
\item[(1)] \emph{The conductor.} The character $\eta$ factors through $(\mathbf Z/N)^\times$; let $\chi$ be the primitive Dirichlet character, of conductor $M$, that induces it, so that $\eta(g)=\chi(a)$ when $g$ acts on $\mu_M$ by $\zeta\mapsto\zeta^a$. A primitive character is nontrivial on the inertia group at every prime dividing its conductor, since otherwise it would factor through the modulus $M/p^{v_p(M)}$. So $D_B(\mathbf Q)\subseteq G_K$ at the places $B$ above the primes $p\mid M$, and these primes split in $K$. Choosing a prime $w_p$ of $K$ above each such $p$, so that $\mathcal O_K/w_p^j\cong\mathbf Z/p^j$ for all $j$, the Chinese remainder theorem gives a ring homomorphism $\rho:\mathcal O_K\to\mathbf Z/M$.
\item[(2)] \emph{Parity.} $\chi(-1)=1$, since $\eta$ is trivial at a complex conjugation, which acts on $\mu_M$ by $\zeta\mapsto\zeta^{-1}$.
\item[(3)] \emph{Units.} Every unit $u$ of $\mathcal O_K$ is a root of unity in a quadratic field, so $u^{12}=1$. Since $2$ splits in $K$, there is a ring homomorphism $\mathcal O_K\to\mathbf Z/4$, and none of the factors $X^2+1$, $X^2\pm X+1$, $X^4-X^2+1$ of $X^{12}-1$ has a root in $\mathbf Z/4$. So $u=\pm1$.
\item[(4)] \emph{A character of ideals.} For $a\in\mathcal O_K$ with $(a)$ prime to $M$, $\rho(a)$ is a unit; put $g(a)=\chi(\rho(a))$. It is multiplicative and, by (2) and (3), trivial on units, so it defines a homomorphism on the group of principal fractional ideals prime to $M$. Since $\bar{\mathbf Q}_\ell^\times$ is divisible, it extends to the group of all fractional ideals prime to $M$; this gives a multiplicative function $f$ on the ideals prime to $M$ with $f((a))=g(a)$.
\item[(5)] \emph{Class field theory.} If $a\equiv1\pmod M$, then $\rho(a)=1$ and $f((a))=1$. By class field theory \cite[Chapter~VI]{Neu99}, there is a character $\kappa':G_K\to\bar{\mathbf Q}_\ell^\times$ with open kernel that is trivial on $I_B(K)$ for every place $B$ above a prime $v$ of $K$ prime to $M$, and maps every Frobenius element at such a place to $f(v)$.
\item[(6)] \emph{The relation.} Fix $c\in G\setminus G_K$. The map $s\mapsto\kappa'(s)\kappa'(csc^{-1})\eta(s)^{-1}$ is a homomorphism of $G_K$ whose kernel contains $\ker\kappa'\cap c^{-1}(\ker\kappa')c\cap\operatorname{Gal}(\bar{\mathbf Q}/K(\mu_M))$, so it is open and closed. Let $\sigma$ be a Frobenius element at a place $B$ above a prime $v$ of $K$ prime to $M$. Then $c\sigma c^{-1}$ is a Frobenius element at $cB$, above $\tau(v)$, which is prime to $M$; and $v\,\tau(v)=(Nv)$, with $Nv$ prime to $M$. So $\kappa'(\sigma)\kappa'(c\sigma c^{-1})=f(v\,\tau(v))=g(Nv)=\chi(Nv)=\eta(\sigma)$, as $\sigma$ acts on $\mu_M$ by $\zeta\mapsto\zeta^{Nv}$. These Frobenius elements are dense in $G_K$ (Theorem~\ref{lit-chebotarev-degree-one}), so the relation holds on $G_K$.
\item[(7)] \emph{Order prime to $\ell$.} The image of $\kappa'$ is finite, of order $\ell^bm''$ with $\ell\nmid m''$. Choose $j\ge0$ with $j\equiv0\pmod{\ell^b}$ and $j\equiv1\pmod m$, and put $\kappa_0=\kappa'^j$. It has open kernel, $\kappa_0^{m''}=1$, and $\kappa_0(s)\kappa_0(csc^{-1})=\eta(s)^j=\eta(s)$, as $\eta^m=1$.
\item[(8)] \emph{Ramification.} If $\eta$ is trivial on $I_B(\mathbf Q)$, then by (1) $B$ is not above a prime dividing $M$, so $\kappa'$ and $\kappa_0$ are trivial on $I_B(K)$.
\end{enumerate}
% (1) is exists_prim, ramified and exists_hom (with Tw.quadFacts for the split primes); (2) uses exists_cc; (3) is units_pm of Tw.quadFacts; (4) is Tw.extIdeal; (5) is Tw.rayChar for the modulus (M); (6) is the subgroup key, with Tw.frobDense and cls_frob; (7) is primeTo.
\end{proof}
\leanref{Development/AutomorphicArithmetic/AutomorphicVarieties/WeakShimuraTheory/FrobeniusDensity/CoprimeOfNotMem.lean}{UpperTheory.V.Ws.Tw.kap0}

\begin{lemma}\label{sha-siegel-automorphy-link-characters}
Let $D$ be a Siegel datum of case II or of case III, with the notation of Definitions~\ref{sha-siegel-datum} and~\ref{sha-siegel-tower-towers}, and let $|\cdot|$ be the absolute value of $\bar{\mathbf Q}_\ell$. Let $T$ be a deep tower for $D$ that satisfies (a) to (m) of Definition~\ref{sha-siegel-tower-towers}, with fields $K$, $F$ and $L$.
Then there are continuous homomorphisms $\kappa,\beta_L:G_L\to\bar{\mathbf Q}_\ell^\times$ and a continuous homomorphism $t_L$ from $G_L$ to the additive group $\mathbf Z_\ell$ such that:
\begin{enumerate}
\item $\kappa(g)\kappa(g')\theta(g)\chi_{\mathrm{cyc}}(g)^3=1$, $\beta_L(g)\beta_L(g')=1$ and $t_L(g)+t_L(g')=0$ for $g\in G_L$ and $g'=\sigma g\sigma^{-1}$ with $\sigma\in G\setminus G_K$;
\item $|\beta_L(g)-1|<1$ for all $g$, and $\beta_L$ is trivial on the kernel of $t_L$;
\item $\kappa$ is trivial on $I_B(L)$ for every place $B$ of $\bar{\mathbf Q}$ not above a prime of $S$; $t_L$ is trivial on $I_B(L)$ for every place $B$ not above $\ell$; and for some $m$ prime to $\ell$, $\kappa^m$ is trivial on $I_B(L)$ for every place $B$ not above $\ell$;
\item if $g\in G_L$ and some $g^{\ell^a}$ lies in the closure of the commutator subgroup of $G_L$, then $\kappa(g)=1$ and $t_L(g)=0$;
\item for every place $B$ above $\ell$ there are $a\in\mathbf Z$, an odd $b\in\mathbf Z$ and $k'\ge1$ with $\big(\kappa(g)\chi_{\mathrm{cyc}}(g)^{-a}\big)^{k'}=1$ and $\big(\beta_L(g)^2\chi_{\mathrm{cyc}}(g)^{-b}\big)^{k'}=1$ for $g\in I_B(L)$.
\end{enumerate}
\end{lemma}
\begin{proof}
By (d), $D_B(\mathbf Q)\subseteq G_K$ for the places $B$ above $2$ and above $\ell$, which lie in $S$.
\begin{enumerate}
\item[(1)] \emph{The ingredients.} Let $\eta$ be as in Lemma~\ref{sha-siegel-automorphy-link-dirichlet}. At a place $B$ not above a prime of $S$, $\eta$ is trivial on $I_B(\mathbf Q)$; at a place above a prime of $S$, $D_B(\mathbf Q)\subseteq G_K$ by (d). So Lemma~\ref{sha-siegel-automorphy-link-kappa-zero} gives $\kappa_0$ on $G_K$, of some order $m$ prime to $\ell$. Let $\gamma$ be as in Lemma~\ref{sha-siegel-automorphy-link-half}, and let $t_K$ and $\beta_K$ be the characters $t$ and $\beta$ of Lemma~\ref{sha-siegel-automorphy-link-beta}, with $A=v_\ell$.
\item[(2)] \emph{The characters.} Let $\beta_L$ and $t_L$ be the restrictions of $\beta_K$ and $t_K$ to $G_L$, let $e\in\{0,1\}$ with $e\equiv k\pmod2$, and put
\[
 \kappa(g)=\kappa_0(g)\,\gamma(g)^{-(k+4)}\,\beta_L(g)^{e}\qquad(g\in G_L).
\]
These are continuous, since $\kappa_0$ has open kernel. Property (2) is (1) of Lemma~\ref{sha-siegel-automorphy-link-beta}.
\item[(3)] \emph{Property (1).} Let $g\in G_L$ and $g'=\sigma g\sigma^{-1}$ with $\sigma\notin G_K$. As $T$ is deep and $G_L\subseteq G_F$, $\psi(g)$ has order prime to $\ell$, so $\kappa_0(g)\kappa_0(g')=\eta(g)=\big(\omega^{k+3}\xi\psi\big)(g)^{-1}$. Since $\gamma$ is a character of $G$, $\gamma(g')=\gamma(g)$, and $\gamma(g)^2=\chi_{\mathrm{cyc}}(g)\omega(g)^{-1}$. By (2) of Lemma~\ref{sha-siegel-automorphy-link-beta}, $\beta_L(g)\beta_L(g')=1$ and $t_L(g)+t_L(g')=0$. So
\[
 \kappa(g)\kappa(g')=\big(\omega^{k+3}\xi\psi\big)(g)^{-1}\big(\chi_{\mathrm{cyc}}\omega^{-1}\big)(g)^{-(k+4)}=\big(\chi_{\mathrm{cyc}}^{k+4}\omega^{-1}\xi\psi\big)(g)^{-1},
\]
and $\theta=\alpha\Psi=\chi_{\mathrm{cyc}}^{k+1}\omega^{-1}\xi\psi$, which gives (1).
\item[(4)] \emph{Property (3).} Let $B$ be a place not above $\ell$. Then $\chi_{\mathrm{cyc}}$ is trivial on $I_B(L)$, hence so are $\gamma$ (Lemma~\ref{sha-siegel-automorphy-link-half}), $t_L$ (Lemma~\ref{sha-siegel-automorphy-link-beta}(3)) and $\beta_L$, which is trivial on $\ker t_L$; so $\kappa^m=\kappa_0^m=1$ on $I_B(L)$. If moreover $B$ is not above a prime of $S$, then $\eta$ is trivial on $I_B(\mathbf Q)$, so $\kappa_0$ is trivial on $I_B(K)$, and $\kappa$ on $I_B(L)$.
\item[(5)] \emph{Property (4).} Let $g^{\ell^a}$ lie in the closure of the commutator subgroup of $G_L$, which lies in the kernel of every continuous homomorphism from $G_L$ to a commutative Hausdorff group. So $\ell^at_L(g)=0$, hence $t_L(g)=0$ and $\beta_L(g)=1$. Also $\chi_{\mathrm{cyc}}(g)^{\ell^a}=1$, and as $\ell$ is odd, $\mathbf Z_\ell^\times$ has no element of order $\ell$; so $\chi_{\mathrm{cyc}}(g)=1$ and $\gamma(g)=1$. Finally $\kappa_0(g)^{\ell^a}=1=\kappa_0(g)^m$ with $m$ prime to $\ell$, so $\kappa_0(g)=1$. Hence $\kappa(g)=1$.
\item[(6)] \emph{Property (5).} Let $B$ be above $\ell$, and let $b=2j+1$ and $k''$ be the integers $b$ and $k$ of (4) of Lemma~\ref{sha-siegel-automorphy-link-beta}. Write $k=2q+e$ and put $a=ej-q-2$ and $k'=2k''m(\ell-1)$. On $I_B(L)$,
\[
 \kappa\chi_{\mathrm{cyc}}^{-a}=\kappa_0\,(\gamma^2)^{-(q+2)}\chi_{\mathrm{cyc}}^{q+2}\,\big(\gamma^{-1}\beta_L\chi_{\mathrm{cyc}}^{-j}\big)^{e}=\kappa_0\,\omega^{q+2}\big(\gamma^{-1}\beta_L\chi_{\mathrm{cyc}}^{-j}\big)^{e},
\]
and $\big(\gamma^{-1}\beta_L\chi_{\mathrm{cyc}}^{-j}\big)^2=\omega\,\beta_L^2\chi_{\mathrm{cyc}}^{-b}$. As $\kappa_0^m=1$, $\omega^{\ell-1}=1$ and $(\beta_L^2\chi_{\mathrm{cyc}}^{-b})^{k''}=1$ on $I_B(L)$, both powers in (5) are $1$.
\end{enumerate}
% The characters are pre_ex (private, TLinkEx.lean), a Pre with Pre.κ = κ₀ γ^(-(e+4)) β^(e mod 2), Pre.βL and Pre.tL; then kapCont, kapConj (kapConj_alg), kapUnr, kapMu (kill, padic_pow) and kapCyc (mul_aux, with a = (e mod 2) j - e/2 - 2 and k = 2 k m (l-1)) give the QCh.
\end{proof}
\leanref{Development/AutomorphicArithmetic/AutomorphicVarieties/WeakShimuraTheory/UnipotentIdempotents/KapConj.lean}{UpperTheory.V.Ws.Lk.kapConj}
\leanref{Development/AutomorphicArithmetic/AutomorphicVarieties/WeakShimuraTheory/UnipotentIdempotents/KapUnr.lean}{UpperTheory.V.Ws.Lk.kapUnr}
\leanref{Development/AutomorphicArithmetic/AutomorphicVarieties/WeakShimuraTheory/UnipotentIdempotents/PadicPowL.lean}{UpperTheory.V.Ws.Lk.kapMu}
\leanref{Development/AutomorphicArithmetic/AutomorphicVarieties/WeakShimuraTheory/UnipotentIdempotents/KapCyc.lean}{UpperTheory.V.Ws.Lk.kapCyc}
\leanref{Development/AutomorphicArithmetic/AutomorphicVarieties/WeakShimuraTheory/UnipotentIdempotents/KapCont.lean}{UpperTheory.V.Ws.Lk.kapCont}

\begin{lemma}\label{sha-siegel-automorphy-link-places}
Let $D$ be a Siegel datum, with the notation of Definitions~\ref{sha-siegel-datum} and~\ref{sha-siegel-tower-towers}, and let $T$ be a tower for $D$ that satisfies (a) to (m), with fields $K$, $F$ and $L$, and relabelling characters $\chi_1,\dots,\chi_4$.
\begin{enumerate}
\item The field $F$ is the maximal totally real subfield of $L$, $L/F$ is unramified at every finite place, and every complex conjugation $\tau\in G_F$ acts nontrivially on $K$.
\item There are an odd prime $p\notin S$ unramified in $L$, a place $v_0$ of $F$ with residue field $\mathbf F_p$, and a place $B_0$ of $\bar{\mathbf Q}$ above $v_0$, such that $D_{B_0}(F)\subseteq G_L$, the residues $\bar\chi_1,\dots,\bar\chi_4$ take equal values on $D_{B_0}(L)$, and $\bar\omega$ is nontrivial on $D_{B_0}(L)$.
\item For such $v_0$, let $S_L$ be the set of places of $F$ above the primes of $S$, together with $v_0$. It is finite, $D_B(F)\subseteq G_L$ for every place $B$ of $\bar{\mathbf Q}$ above a place of $S_L$, and the places of $F$ outside $S_L$ are not above primes of $S$. The map that sends a place $v$ of $F$ above $\ell$ to the element of $\Sigma_T$ above it is a bijection onto $\Sigma_T$.
\end{enumerate}
\end{lemma}
\begin{proof}
(1) As $F$ is totally real and $K$ imaginary quadratic, $L=FK$ is a totally imaginary quadratic extension of $F$. By (c), $I_B(F)\subseteq G_L$ for every place $B$, so $L/F$ is unramified at the finite places. A complex conjugation acts on the imaginary quadratic field $K$ by its nontrivial automorphism.

(2) \emph{An element of $G_L$.} For all $i,j$, $\chi_i\chi_j^{-1}=\chi^{w}$ for the weight $w$ of the entry $(i,j)$, and $\chi^w\chi^{w''}=\chi^{w+w''}$. So by (e), $\bar\omega$ is not a product of the $\bar\chi_i\bar\chi_j^{-1}$ on $G_L$. The image of $G_L$ in $(\mathbf F^\times)^5$ under $(\bar\chi_1,\dots,\bar\chi_4,\bar\omega)$ is a finite abelian group, so by the duality of finite abelian groups some $g_0\in G_L$ has $\bar\chi_1(g_0)=\dots=\bar\chi_4(g_0)$ and $\bar\omega(g_0)\ne1$.

\emph{The place $v_0$.} Let $U\subseteq G_L$ be the open subgroup on which $\bar\delta_1,\dots,\bar\delta_4$ and $\bar\omega$ are trivial, and $U'\subseteq U$ an open subgroup normal in $G_F$. By the Chebotarev density theorem (Theorem~\ref{lit-chebotarev-degree-one}), applied to the fixed field of $U'$ over $F$, there are a place $v_0$ of $F$ of degree one over a prime $p$ not dividing $2\prod_{q\in S}q\cdot\operatorname{disc}(L)$, unramified in that field, a place $B_0$ above $v_0$, and a Frobenius element $\sigma\in D_{B_0}(F)$ with $\sigma\in g_0U'$. Then $p$ is odd, $p\notin S$ and $p$ is unramified in $L$. The group $D_{B_0}(F)$ is the closure of the subgroup generated by $\sigma$ and $I_{B_0}(F)\subseteq U'$; since $U'$ is normal and open, the subgroup generated by $g_0$ and $U'$ is closed, so $D_{B_0}(F)$ lies in it, and in particular in $G_L$. So every $d\in D_{B_0}(L)$ lies in $g_0^nU$ for some $n\in\mathbf Z$, and $\bar\chi_i(d)=\bar\chi_i(g_0)^n$ is independent of $i$, the $\bar\chi_i$ being the $\bar\delta_j$ in some order. And $\sigma\in g_0U$ has $\bar\omega(\sigma)=\bar\omega(g_0)\ne1$.

(3) There are finitely many places above the primes of $S$. For a place $B$ above a prime of $S$, $D_B(\mathbf Q)\subseteq G_K$ by (d), so $D_B(F)\subseteq G_F\cap G_K=G_L$; for $B$ above $v_0$ this holds by (2), as all places above $v_0$ are conjugate under $G_F$ and $G_L$ is normal. The statement on the places outside $S_L$ is the definition of $S_L$. The last statement holds because every place of $\bar{\mathbf Q}$ above $\ell$ is conjugate under $G_F$ to exactly one element of $\Sigma_T$.
% (1) is Lk.galEq, Tower.Ok.unr and Lk.ccOf; (2) is Lk.gEx (Tower.Ok.g2 with Loc.v0e and Loc.mem_closure_of_kernel_containment) and Lk.v0Ex (Loc.v00 through Published.cheb_holds, and TameQuot.closure_frob_sup_Ine), with the open subgroup U of setEx from Lk.barFacts and M = 2 (prod S) |disc L|; (3) is Lk.mkSet with Tower.Ok.split, and Lk.placesEq.
\end{proof}
\leanref{Development/AutomorphicArithmetic/AutomorphicVarieties/WeakShimuraTheory/UnipotentIdempotents/EqOfKer.lean}{UpperTheory.V.Ws.Lk.setEx}
\leanref{Development/AutomorphicArithmetic/AutomorphicVarieties/WeakShimuraTheory/UnipotentIdempotents/GExRtZero.lean}{UpperTheory.V.Ws.Lk.gEx}
\leanref{Development/AutomorphicArithmetic/AutomorphicVarieties/WeakShimuraTheory/UnipotentIdempotents/RangeOpen.lean}{UpperTheory.V.Ws.Lk.v0Ex}
\leanref{Development/AutomorphicArithmetic/AutomorphicVarieties/WeakShimuraTheory/UnipotentIdempotents/EqOfMem.lean}{UpperTheory.V.Ws.Lk.mkSet}

\begin{lemma}\label{sha-siegel-automorphy-link-exists}
Let $D$ be a Siegel datum of case II or of case III (Definition~\ref{sha-siegel-tower-towers}), with the notation of Definition~\ref{sha-siegel-datum}. Let $T$ be a deep tower for $D$ that satisfies (a) to (m) of Definition~\ref{sha-siegel-tower-towers}.
Then there is a ring $\mathcal O_1\supseteq\mathcal O$ of Definition~\ref{sha-character-parts} such that $T$ has a link over every ring $\mathcal O'\supseteq\mathcal O_1$ of Definition~\ref{sha-character-parts} (Definition~\ref{sha-siegel-automorphy-link}).
\end{lemma}
\begin{proof}
Let $K$, $F$ and $L$ be the fields of $T$, and $\chi_1,\dots,\chi_4$ its relabelling characters.
\begin{enumerate}
\item[(1)] \emph{Characters.} Let $\kappa$, $\beta_L$ and $t_L$ be as in Lemma~\ref{sha-siegel-automorphy-link-characters}.
\item[(2)] \emph{Coefficients.} As $G_L$ is compact, the continuous characters $\kappa$ and $\beta_L$ take values in a finite extension of $\mathbf Q_\ell$, by the Baire category theorem, and in the units of its ring of integers, since their images are compact. So there is a ring $\mathcal O_1\supseteq\mathcal O$ of Definition~\ref{sha-character-parts} that contains the values of $\kappa$ and $\beta_L$ and the images of all embeddings of $\mathcal O_L$ in $\bar{\mathbf Q}_\ell$. Let $\mathcal O'\supseteq\mathcal O_1$ be such a ring, with maximal ideal $\mathfrak m_{\mathcal O'}$ and residue field $k'$. Then $\kappa,\beta_L:G_L\to\mathcal O'^\times$ are continuous for the $\mathfrak m_{\mathcal O'}$-adic topology, which is the topology induced by $\bar{\mathbf Q}_\ell$, and $\beta_L\equiv1\pmod{\mathfrak m_{\mathcal O'}}$, as $|\beta_L-1|<1$.
\item[(3)] \emph{The Hida data.} Let $p$, $v_0$, $B_0$ and $S_L$ be as in Lemma~\ref{sha-siegel-automorphy-link-places}, and let $\tau\in G_F$ be a complex conjugation. The field $L$ contains no primitive $\ell$-th root of unity: an element $g$ of $G_L$ with $\bar\omega(g)\ne1$ has $\chi_{\mathrm{cyc}}(g)\not\equiv1\pmod\ell$, as $\bar\omega=\bar\chi_{\mathrm{cyc}}$, so it moves the $\ell$-th roots of unity. Take $S_L$ as the set $S$ of Situation~\ref{sha-hecke-situation}, with the places above the primes of $S$ other than $\ell$ as its auxiliary places; for $\tilde v$, when $v\mid\ell$, the place of $L$ below the element $\bar v$ of $\Sigma_T$ above $v$, and for $\tilde v_0$ the place of $L$ below $B_0$; $\mathcal O'$ in place of $\mathcal O$, and $n=4$. By Lemma~\ref{sha-siegel-automorphy-link-places} and (2), these are data of Situation~\ref{sha-hecke-situation} for $L$, with $L^+=F$ and $c$ the image of $\tau$, and $v\mapsto\bar v$ is a bijection as in (1) of Definition~\ref{sha-siegel-automorphy-link}, with $D_v=D_{\bar v}(L)$.
\item[(4)] \emph{Residual characters.} For each $i$, let $\varrho_i(g)$ be the residue of $\kappa(g)\chi_i(g)$, a character of $G_L$ with values in $k'^\times$ and open kernel. It is trivial on the inertia groups at the places above the places of $F$ outside $S_L$: these places are not above primes of $S$, so $\kappa$ is trivial there by (3) of Lemma~\ref{sha-siegel-automorphy-link-characters}, and so is $\chi_i$, which is unramified outside $S$. As in step (2) of the proof of Lemma~\ref{sha-siegel-automorphy-phi}, the kernel of $\mathrm{pr}$ lies in every closed subgroup of $G_L$ that contains these inertia groups; so $\varrho_i=\bar\chi'_i\circ\mathrm{pr}$ for a character $\bar\chi'_i$ of $\Gamma_0$ with open kernel.
\item[(5)] \emph{They are distinct.} If $\bar\chi'_i=\bar\chi'_j$ with $i\ne j$, then the residue of $\chi_i\chi_j^{-1}=\chi^{w'}$, where $w'$ is the weight of the entry $(i,j)$, a root, would be trivial on $G_L$, contrary to (f).
\item[(6)] \emph{Polarization.} Let $g\in G_L$ and $g'=\tau g\tau^{-1}$. Then $\chi_i(g')=\chi_i(g)$, since $\chi_i$ is a character of $G$, and $\chi_i\chi_{5-i}=\theta$. So (1) of Lemma~\ref{sha-siegel-automorphy-link-characters} gives $\varrho_i(g')\varrho_{5-i}(g)\bar\chi_{\mathrm{cyc}}(g)^3=1$. As $\varepsilon\circ\mathrm{pr}=\chi_{\mathrm{cyc}}$ on $G_L$, this says $\bar\chi'_i(y^*)=\bar\chi'_{5-i}(y)$ for $y\in\Gamma_0$, with $y^*=\varepsilon(y)^{-3}(c\,y\,c^{-1})^{-1}$, and by linearity for $y$ in the group algebra. So $g\mapsto g^*$ permutes the factors of $\bar D=\prod_i\bar\chi'_i$, and $\bar D$ takes the same value at $y$ and $y^*$.
\item[(7)] \emph{The residual datum.} With the $\bar\chi'_i$ as the characters $\bar\chi_i$ of (3) of Definition~\ref{sha-siegel-automorphy-link}, $\bar\psi_{v,i}$ the restriction of $\bar\chi'_i$ to $D_v$ and $\theta_{v,i}=1$, (4) to (6) give a residual datum of rank $4$ with the properties in (3) of Definition~\ref{sha-siegel-automorphy-link} other than those at $v_0$. At $v_0$: the place $v_0$ is not above a prime of $S$, so $\kappa$ and the $\chi_i$ are trivial on the inertia groups above it, and the $\bar\chi'_i$ are trivial on $I_{v_0}$; the residue of $\kappa$ is the same factor in every $\varrho_i$, so by (2) of Lemma~\ref{sha-siegel-automorphy-link-places} the $\bar\chi'_i$ take equal values on $D_{v_0}$; and some $d\in D_{B_0}(L)$ has $\bar\varepsilon(\mathrm{pr}(d))=\bar\chi_{\mathrm{cyc}}(d)=\bar\omega(d)\ne1$.
\item[(8)] \emph{The characters of the link.} Conditions (a) to (f) of Definition~\ref{sha-siegel-automorphy-link} hold for $\kappa$, $\beta_L$ and $t_L$. Condition (a) is (1) of Lemma~\ref{sha-siegel-automorphy-link-characters}, with $\sigma=\tau$, which is not in $G_K$ by (1) of Lemma~\ref{sha-siegel-automorphy-link-places}. In (b), the places above the places of $F$ outside $S_L$ are not above primes of $S$, so $\kappa$ is trivial on their inertia groups by (3) of Lemma~\ref{sha-siegel-automorphy-link-characters}, and so is $\bar r_0$, whose entries are unramified outside $S$. Condition (c) is (3) of that lemma. For (d), an element of $\mu_v$ has a power $g^{\ell^a}$ in the closure of the commutator subgroup of $D_v$, hence of $G_L$, and (4) of that lemma applies. Condition (e) is (5) of that lemma, and (f) is (2) and (3) of that lemma with (2) above.
\end{enumerate}
% (2) is Pr.coefEx (ce_char through CompactRepField.exists_finiteDimensional_entries_mem, ce_large for W.Large), Lk.liftO and Lk.ochOf; (3) is Lk.setEx, Lk.noMu and Lk.ccOf, with the group data V from the hypothesis W.LevEx of Ws.UnramifiedInputs, which W.levEx proves; (4) is Lk.descend (phiKer) with Lk.cbFacts; (5) is Lk.cbInj (Tower.Ok.g3); (6) is cbFacts and Lk.polar with the permutation i -> 5-i; (7) is Lk.residEx and Lk.v0ok; (8) is the construction of the Link in nonempty_link_of_fit_deep_unramified, with Lk.prFacts for (d).
\end{proof}
\leanref{Development/AutomorphicArithmetic/AutomorphicVarieties/ModularityArgument/ThresholdArguments/DenseThresholdCases/TLinkEx.lean}{UpperTheory.V.Ws.nonempty_link_of_fit_deep_unramified}
\leanref{Development/AutomorphicArithmetic/AutomorphicVarieties/CurveDVRSpecialization/CoefEx.lean}{UpperTheory.V.Pr.coefEx}
\leanref{Development/AutomorphicArithmetic/AutomorphicVarieties/WeakShimuraTheory/UnipotentIdempotents/ResidEx.lean}{UpperTheory.V.Ws.Lk.residEx}
\leanref{Development/AutomorphicArithmetic/AutomorphicVarieties/WeakShimuraTheory/UnipotentIdempotents/V0ok.lean}{UpperTheory.V.Ws.Lk.v0ok}
\leanref{Development/AutomorphicArithmetic/AutomorphicVarieties/WeakShimuraTheory/UnipotentIdempotents/NoMu.lean}{UpperTheory.V.Ws.Lk.noMu}

\subsection{The points of a Siegel datum}\label{subsection-sha-siegel-datum-points}

\begin{definition}\label{sha-siegel-points-three-types}
Let $D$ be a Siegel datum and $T$ a tower for $D$ with field $F$, with the notation of Definitions~\ref{sha-siegel-datum}, \ref{sha-siegel-tower-towers}, \ref{sha-siegel-tower-pairs} and~\ref{sha-siegel-step-chains-notions}.
\begin{enumerate}
\item The tower $T$ is \emph{of tensor type} if $\ell\ge5$, the set $S_b$ is empty, $\psi^m=1$ for some $m$ prime to $\ell$, and there is an imaginary quadratic field $K'$, a \emph{tensor field} of $T$, such that $\ell$ does not divide the discriminant of $K'$, $D_w(\mathbf Q)\subseteq G_{K'}$ for every place $w$ of $\bar{\mathbf Q}$ above $\ell$ or above a prime of $S$, and, for $g\in G_F$, $\bar\chi^{(1,0)}(g)=1$ if $g$ fixes $K'$ and $\bar\chi^{(1,0)}(g)=-1$ otherwise.
\item Put $W_{\mathrm A}=\{(1,-1),(0,-1),(-1,-1)\}$. A flagged pair $P=(r,(n_v))$ for a problem of $T$ is \emph{of type A} if $r$ is supported in $W_{\mathrm A}$ and every $n_v$ has a nonzero entry $(2,1)$. It is \emph{A-ample} if it satisfies one of the four conditions that define an ample pair in Definition~\ref{sha-siegel-tower-pairs}(3), or it is of type A and $r$ has an entry of weight $(0,-1)$, or it is of type A and $r$ has entries of weights $(1,-1)$ and $(-1,-1)$.
\item An \emph{A-chain} from a flagged pair $P$ to a flagged pair $P'$ is a sequence $P=P_0,P_1,\dots,P_m=P'$ of flagged pairs such that $P_0$ is A-ample, and each $P_i$ with $i\ge1$ is A-ample and linked to $P_{i-1}$.
\end{enumerate}
\end{definition}
\leanref{Development/AutomorphicArithmetic/AutomorphicVarieties/ModularityArgument/CuspidalTwists/PromodularityCaseThree.lean}{UpperTheory.V.Cu.QuadOn}
\leanref{Development/AutomorphicArithmetic/AutomorphicVarieties/ModularityArgument/CuspidalTwists/PromodularityCaseThree.lean}{UpperTheory.V.Cu.Four}
\leanref{Development/AutomorphicArithmetic/AutomorphicVarieties/ModularityArgument/CuspidalTwists/PromodularityCaseThree.lean}{UpperTheory.V.Cu.InA}
\leanref{Development/AutomorphicArithmetic/AutomorphicVarieties/ModularityArgument/CuspidalTwists/PromodularityCaseThree.lean}{UpperTheory.V.Cu.Good4}
\leanref{Development/AutomorphicArithmetic/AutomorphicVarieties/ModularityArgument/CuspidalTwists/PromodularityCaseThree.lean}{UpperTheory.V.Cu.Chain4}

\begin{lemma}\label{sha-siegel-points-three-tower-quadratic}
Let $S'$ be a finite set of primes and $B\ge0$. There are a prime $q>B$ not in $S'$ and a character $\eta:G\to\{\pm1\}$ with open kernel such that: $\eta(\sigma)=-1$ for every complex conjugation $\sigma$; $D_w(\mathbf Q)\subseteq\ker\eta$ for every place $w$ of $\bar{\mathbf Q}$ above a prime of $S'$; for every place $w$ above $q$ some $g\in I_w(\mathbf Q)$ has $\eta(g)=-1$; and $\ker\eta=G_{K'}$ for an imaginary quadratic field $K'$ of discriminant $-q$ (Definition~\ref{sha-siegel-tower-towers}).
\end{lemma}
\leanref{Development/AutomorphicArithmetic/AutomorphicVarieties/ModularityArgument/CuspidalTwists/ComplexConjugationsConjugateCaseThree.lean}{UpperTheory.V.Cu.RankFourGaloisLemmas.quad}
\leanref{Development/AutomorphicArithmetic/AutomorphicVarieties/GaloisLibrary/ClassFieldRamification/ExistsQuadraticGaussSum.lean}{UpperTheory.V.CF.gauss}
\leanref{Development/AutomorphicArithmetic/AutomorphicVarieties/HeightAndVanishing/SplittingFieldIsoQuadratic.lean}{UpperTheory.V.HV.kstar}
\begin{proof}
Use the characters $c_N$ of the first paragraph of the proof of Lemma~\ref{sha-siegel-step-towers-quadratic}, and the Legendre symbol $(\frac{\cdot}{q})$.
\begin{enumerate}
\item[(a)] \emph{The prime.} By Dirichlet's theorem there is a prime $q$, larger than $B$ and than every element of $S'$, with $q\equiv-1$ modulo $8\prod_{p\in S'}p$. Then $q\notin S'$ and $q\equiv3\pmod4$, and, as in step (a) of the proof of Lemma~\ref{sha-siegel-step-towers-quadratic} with $s=-1$, $(\frac{-1}q)=-1$ and $(\frac pq)=1$ for every $p\in S'$.
\item[(b)] \emph{The character.} Let $\eta(g)=(\frac{c_q(g)}q)$. As in steps (b) and (c) of that proof, $\eta$ has open kernel of index $2$, $\eta(\sigma)=(\frac{-1}q)=-1$ for every complex conjugation $\sigma$, $\eta$ is trivial on $D_w(\mathbf Q)$ for every place $w$ above a prime of $S'$, and for every place $w$ above $q$ some $g\in I_w(\mathbf Q)$ has $\eta(g)=-1$.
\item[(c)] \emph{The field.} The kernel of $\eta$ is open of index $2$ and contains the commutator subgroup, so it is $G_{K'}$ for a quadratic field $K'$, which is imaginary as no complex conjugation lies in $\ker\eta$. By the theory of quadratic Gauss sums \cite{Was97}, there is $s\in\bar{\mathbf Q}$ with $s^2=(-1)^{(q-1)/2}q=-q$ and $\sigma(s)=(\frac{c_q(\sigma)}q)s$ for every $\sigma\in G$. So $\ker\eta$ fixes $s$, and $K'=\mathbf Q(s)=\mathbf Q(\sqrt{-q})$. As $-q\equiv1\pmod4$, its discriminant is $-q$ \cite[Chapter~I]{Neu99}.
\end{enumerate}
\end{proof}
% RankFourGaloisLemmas.quad (ComplexConjugationsConjugateCaseThree.lean): (a) is the private exists_prime_quadratic_conditions (Nat.forall_exists_prime_gt_and_zmodEq for max (S.sup id) B, and the private jacobi_symbol_eq_one_of_congr for quadratic reciprocity); (b) is Chr.clsFacts, Chr.clsOnto for the index, hccη, Chr.decCyc with the private cyclotomic_class_eq_of_frobenius and valuation_prime_lt_one_exclusive for the splitting, and compIne for the ramification; (c) is fixAb and fixInf, CF.gauss (the Gauss sum, with CF.pStar p = -p as p is 3 mod 4), the private imag_quadratic_field_unique (K' is the splitting field of X^2 + p) and HV.kstar for the discriminant. It repeats quadEx (Lemma sha-siegel-step-towers-quadratic) for s = -1, with the bound B and the field K' added.

\begin{lemma}\label{sha-siegel-points-three-tower-elements}
Let $D$ be a Siegel datum of case III, attached to $(h,k,\psi)$, with $\xi=\xi_0\xi'$ as in Lemma~\ref{sha-siegel-step-towers-three-elements} (Definitions~\ref{sha-siegel-datum} and~\ref{sha-siegel-tower-towers}).
\begin{enumerate}
\item There are $g\in G$ and an odd integer $m$ prime to $\ell$ with $\omega(g)=1$, $\xi(g)=1$, $\psi(g)^m=1$ and $\psi(g)\ne1$.
\item There is $B\ge0$ such that $\xi$, $\psi$ and $\omega$ are trivial on $I_w(\mathbf Q)$ for every place $w$ above a prime $q>B$.
\end{enumerate}
\end{lemma}
\leanref{Development/AutomorphicArithmetic/AutomorphicVarieties/ModularityArgument/CuspidalTwists/ComplexConjugationsConjugateCaseThree.lean}{UpperTheory.V.Cu.RankFourGaloisLemmas.odd}
\leanref{Development/AutomorphicArithmetic/AutomorphicVarieties/ModularityArgument/CuspidalTwists/ComplexConjugationsConjugateCaseThree.lean}{UpperTheory.V.Cu.RankFourGaloisLemmas.unr}
\begin{proof}
Use the characters $c_M$ of Lemma~\ref{sha-siegel-step-towers-kappa-dirichlet}, write $p_1=p_1(\psi_0)$ and $n=n_{\mathrm{III}}$, so that $\psi$, $\xi_0$ and $\xi'$ are characters modulo $N_h$, $p_1$ and $2^n$, and recall that $p_1$ is an odd prime that divides neither $N_h$ nor $\ell$, and that $\ell\nmid N_h$ (Definition~\ref{sha-host-datum} and Lemma~\ref{sha-siegel-step-towers-three-host}).

(1) Let $m$ be the order of $\psi$, which is odd and prime to $\ell$ by Lemma~\ref{sha-siegel-step-towers-three-host}. As $\psi\ne1$, some $a\in(\mathbf Z/N_h)^\times$ has $\psi(a)\ne1$. As $N_h$ is prime to $\ell p_1$, step (a) of the proof of Lemma~\ref{sha-siegel-step-towers-three-elements} gives $u\in G$ with $c_{N_h}(u)=a$, $c_\ell(u)=1$ and $c_{p_1}(u)=1$. Let $g=u^{2^n}$. Then $c_{p_1}(g)=1$, and $c_{2^n}(g)=c_{2^n}(u)^{2^n}=1$, as the exponent of $(\mathbf Z/2^n)^\times$ divides $2^n$; so $\xi(g)=1$. As $c_\ell(g)=1$, $\omega(g)=1$ by Lemma~\ref{sha-siegel-step-towers-kappa-units}(3). Finally $\psi(g)=\psi(a)^{2^n}$, so $\psi(g)^m=1$, and $\psi(g)\ne1$, as $\psi(a)^m=1$, $\psi(a)\ne1$ and $2^n$ is prime to $m$.

(2) Let $B=2\ell p_1N_h$, let $q>B$ be a prime and $w$ a place above $q$. Then $q$ divides none of $2^n$, $p_1$, $N_h$ and $\ell$. So $\xi_0$, $\xi'$ and $\psi$ are trivial on $I_w(\mathbf Q)$ by Lemma~\ref{sha-siegel-step-towers-kappa-dirichlet}(2), and so is $\omega$, by Lemma~\ref{sha-siegel-step-towers-kappa-units}(3).
\end{proof}
% (1) is RankFourGaloisLemmas.odd: bchiFacts and pK_pOne_prime_conditions for the coprimality, hostOdd for the order of psi (hdo, hlm), omT for omega (hom), large₂.clsIndep for u, the private unit_pow_two_pow_eq_one and dirichlet_unit_pow_order_eq_one, and e1, e2. (2) is RankFourGaloisLemmas.unr, with B = N_h l p_1 2, the private dirichlet_galois_char_unramified (Chr.clsFacts) and omT (o6).

\begin{lemma}\label{sha-siegel-points-three-tower}
Let $D$ be a Siegel datum of case III and $\epsilon$ a relabelling satisfying (a) of Definition~\ref{sha-siegel-tower-towers}. There is a tame tower $T$ for $D$ with relabelling $\epsilon$ that satisfies (a) to (m) of Definition~\ref{sha-siegel-tower-towers} and is of tensor type (Definition~\ref{sha-siegel-points-three-types}). Such a tower is fit.
\end{lemma}
\leanref{Development/AutomorphicArithmetic/AutomorphicVarieties/ModularityArgument/CuspidalTwists/ComplexConjugationsConjugateCaseThree.lean}{UpperTheory.V.Cu.existsTensorSeedTower_holds}
\leanref{Development/AutomorphicArithmetic/AutomorphicVarieties/ModularityArgument/CuspidalTwists/Sp4In.lean}{UpperTheory.V.Cu.hFit}
\begin{proof}
A tower of tensor type for $D$ of case III that satisfies (a) to (m) is fit by (4) of Definition~\ref{sha-siegel-tower-towers}: $\ell\ge5$, $S_b$ is empty and $\psi^m=1$ for some $m$ prime to $\ell$.

Now let $D$ be attached to $(h,k,\psi)$, let $m$ be the order of $\psi$, and let $s=1$ if $\epsilon=0$ and $s=-1$ if $\epsilon=1$, so that $\chi_1\chi_2^{-1}=\alpha$, $\chi_1\chi_3^{-1}=\Psi^s$ and $\chi_1\chi_4^{-1}=\alpha\Psi^s$. Recall that $\bar\alpha=\bar\xi$ and $\bar\Psi=\bar\omega^k\bar\psi$, as $\bar\omega=\bar\chi_{\mathrm{cyc}}$, and that $\bar r_0(g)=1$ if and only if $\bar\alpha(g)=1$ and $\bar\Psi(g)=1$. We view characters with values $\pm1$ as characters with values in $\mathcal O^\times$.
\begin{enumerate}
\item[(a)] \emph{The tensor field.} By Lemma~\ref{sha-siegel-step-towers-three-host}, $\ell\ge5$, $\ell-1\mid2k$, $m$ is odd and prime to $\ell$, $S_b$ is empty, and $\xi$ and $\psi$ are trivial on the inertia groups above $\ell$; as $\omega^{\ell-1}=1$, $\omega^{2k}=1$. Let $B$ be as in Lemma~\ref{sha-siegel-points-three-tower-elements}(2). Lemma~\ref{sha-siegel-points-three-tower-quadratic}, applied to $S$ and $B+\ell$, gives a prime $q>B+\ell$ not in $S$, a character $\eta$ and an imaginary quadratic field $K'$ with $\ker\eta=G_{K'}$ and discriminant $-q$.
\item[(b)] \emph{The subgroup $H$.} Let $H$ be the set of $g\in G$ with $\xi(g)=\eta(g)$ and $\bar\Psi(g)=\bar\psi(g)$, the intersection of the kernels of two characters. It contains $\ker\xi\cap\ker\eta\cap\ker\psi\cap\ker\omega$, so it is open, and it contains the commutator subgroup. It contains $c$: $\psi(c)^2=1$ as $c^2=1$, and $\psi(c)^m=1$ with $m$ odd, so $\psi(c)=1=\Psi(c)$; and $\xi(c)=-1=\eta(c)$. For $w$ above $\ell$ and $g\in I_w(\mathbf Q)$, $g^2\in H$, as $\xi(g^2)=\psi(g^2)=1$, $\eta(g^2)=1$ and $\omega(g^2)^k=1$.
\item[(c)] \emph{The fields.} Lemma~\ref{sha-siegel-step-towers-three-fields}, with $S'=S\cup\{q\}$ and this $H$, gives $K$, $F$ and $J$. Put $L=FK$, so that $G_L=G_F\cap G_K=H\cap J$.
\item[(d)] \emph{Elements.} Lemma~\ref{sha-siegel-step-towers-three-elements}(1) gives $g_1\in I_2$ with $\xi(g_1)=-1$, $\psi(g_1)=1$ and $\omega(g_1)=1$; and $\eta(g_1)=1$, as $2\in S$. Lemma~\ref{sha-siegel-points-three-tower-elements}(1) gives $g_2\in G$ and an odd $m'$ prime to $\ell$ with $\omega(g_2)=\xi(g_2)=1$, $\psi(g_2)^{m'}=1$ and $\psi(g_2)\ne1$. Lemma~\ref{sha-siegel-points-three-tower-quadratic} gives $g_5\in I_q$ with $\eta(g_5)=-1$, and $\xi(g_5)=\psi(g_5)=\omega(g_5)=1$ as $q>B$. Let $g_0=g_1g_5g_2^2$. Then $\xi(g_0)=-1=\eta(g_0)$, $\omega(g_0)=1$ and $\psi(g_0)=\psi(g_2)^2$, so $\bar\Psi(g_0)=\bar\psi(g_0)$ and $g_0\in H$. As $2,q\in S'$, $g_1\in D_2$ and $g_5\in D_q$ lie in $J$, and so does $g_2^2$; so $g_0\in G_L$. Put $y=\bar\psi(g_2)^2$ and $z=y^s$. Then $y^{m'}=1$, and $y\ne1$: otherwise $\psi(g_2)^2=1$ by Lemma~\ref{sha-siegel-step-towers-kappa-units}(1), as $(\psi(g_2)^2)^{m'}=1$, and then $\psi(g_2)=1$, as $m'$ is odd. So $z^{m'}=1$, $z\ne1$, and $z\ne-1$, as $(-1)^{m'}=-1\ne1$ in $\mathbf F$. Finally, as in step (c) of the case III part of the proof of Lemma~\ref{sha-siegel-step-towers}(1), some $g_3\in I_\ell$ has $\omega(g_3^2)\ne1$; let $g_4=g_3^2$, which lies in $H$ by (b) and in $J$, so $g_4\in G_L$.
\item[(e)] \emph{The tower.} As $\ell\in S$, $D_w(\mathbf Q)\subseteq J\subseteq G_K$ for $w$ above $\ell$, and Lemma~\ref{sha-siegel-step-towers-make} gives a tower $T$ with relabelling $\epsilon$, fields $K$ and $F$, and $g_0^T=g_0$.
\item[(f)] \emph{Characters on $G_F$.} For $g\in G_F\subseteq H$, $\bar\chi^{(1,0)}(g)=\bar\alpha(g)=\bar\xi(g)=\eta(g)$ and $\bar\chi^{(0,1)}(g)=\bar\Psi(g)^s=\bar\psi(g)^s$; so $\bar\chi^{w'}(g)=\eta(g)^{w_1}\bar\psi(g)^{sw_2}$ for $w'=(w_1,w_2)$. At $g_0$ this is $(-1)^{w_1}z^{w_2}$, which is not $1$ when $w'$ is a root, as $-1\ne1$ and $z\ne\pm1$. The $\bar\chi_i(g_0)$ are $\bar\chi_1(g_0)$ times the inverses of $1$, $-1$, $z$ and $-z$, which are pairwise distinct as $-1\ne1$ and $z\ne\pm1$; so the $\bar\chi_i(g_0)$ are pairwise distinct.
\item[(g)] \emph{Conditions (a) to (m).} (a) is the hypothesis on $\epsilon$, and (b) holds as $F$ has even degree. (c): $I_w(F)\subseteq I_w(\mathbf Q)\cap G_F\subseteq G_K$. (d): $D_w(\mathbf Q)\subseteq J\subseteq G_K$ for $w$ above a prime of $S$. (e): $\bar\chi^{w'}(g_4)=1$ by (f), as $\xi(g_4)=\psi(g_4)=1$, while $\bar\omega(g_4)\ne1$ by Lemma~\ref{sha-siegel-step-towers-kappa-units}(1), as $\omega(g_4)\ne1$ and $\omega(g_4)^{\ell-1}=1$. (f) is in (f), as $g_0\in G_L$. (g): let $\eta'$ be as there, and $j$ with $\eta'\bar\chi_1=\bar\chi_j$ on $G_L$; as $\omega(g_0)=1$, $\eta'(g_0)=1$, so $\bar\chi_1(g_0)=\bar\chi_j(g_0)$, $j=1$ by (f), and $\eta'=1$. (h): for the root $(1,1)$ and $0<i\le4$, $\bar\chi^{(1,1)}(g_0)^i=(-z)^i\ne1$. Indeed, $z=x^2$ for $x=z^{(m'+1)/2}$, which has $x^{m'}=1$ and $x\ne1$; if $i$ is even, $(-x^2)^i=x^{2i}$ with $2i\mid8$ prime to $m'$, so $x^{2i}\ne1$; if $i$ is odd, $((-x^2)^i)^{m'}=-1\ne1$. (i): let $g\in I_w(F)$ with $w$ not above $\ell$ and $\bar r_0(g)=1$; then $\bar\Psi(g)=1$ and $\omega(g)=1$ (Lemma~\ref{sha-siegel-step-towers-kappa-units}(3)), so $\bar\psi(g)=1$, and $\psi(g)=1$ by Lemma~\ref{sha-siegel-step-towers-kappa-units}(1), as $\psi(g)^m=1$. (j): $\psi^m=1$. (k): let $w$ be above $\ell$ and $g\in I_w(F)$; then $\xi(g)=\psi(g)=1$, so $\bar\alpha(g)=1$ and $\bar\Psi(g)=\bar\psi(g)=1$, as $g\in H$; so $\bar r_0(g)=1$. (l) is in (f). (m): let $w$ be above $\ell$; if $\bar\omega$ were trivial on $I_w(F)$, then for $g\in I_w(\mathbf Q)$, $g^2\in H\cap J\subseteq G_F$ by (b), so $g^2\in I_w(F)$ and $\omega(g^2)=1$ by Lemma~\ref{sha-siegel-step-towers-kappa-units}(1); so $\omega^2$ would be trivial on $I_w(\mathbf Q)$, which is impossible, as in (d). So $\bar\omega$ is nontrivial on $I_w(F)$.
\item[(h)] \emph{Tensor type and tameness.} By (a), $\ell\ge5$, $S_b$ is empty and $\psi^m=1$ with $m$ prime to $\ell$. The field $K'$ is imaginary quadratic, and $\ell$ does not divide its discriminant $-q$, as $q>\ell$. As $\ell\in S$, $D_w(\mathbf Q)\subseteq\ker\eta=G_{K'}$ for every place $w$ above $\ell$ or above a prime of $S$. For $g\in G_F$, $\bar\chi^{(1,0)}(g)=\eta(g)$ by (f), which is $1$ if $g$ fixes $K'$ and $-1$ otherwise. So $T$ is of tensor type, with tensor field $K'$. It is tame: for $w$ above $\ell$ and $g\in I_w(\mathbf Q)$, $g^2\in H\cap J\subseteq G_F$ by (b), and $2$ is prime to $\ell$.
\end{enumerate}
% The fitness is Cu.hFit (Sp4In.lean): Fit is Ok with Ctx3 and Four.five, Four.sb, Four.psi. The construction is existsTensorSeedTower_holds (ComplexConjugationsConjugateCaseThree.lean, about 165 lines; with the lemmas cited, about 2,100 lines): (a) is large₂.host3, omT, RankFourGaloisLemmas.unr and RankFourGaloisLemmas.quad for S and B + l; (b) is memH, hopen (with reduction_psi_eq_omega_pow_case_three), hcomm, hψcc and hcc (complex_conjugation_sq, Psi_even, xi_odd), and hsq (with h2e); (c) is the private exists_cm_and_real_fields, a copy of the private fieldsEx (Lemma sha-siegel-step-towers-three-fields), and hL; (d) is large₂.eltTwo, RankFourGaloisLemmas.odd, hηram with unr for g_5, hx, hw, hp, hηx, then y, z, hym, hy1 (redOne), hz1 and hzn (odd_order_unit_ne_neg_one), and g_3 from omT (o7) with g_4 from hL; (e) is mkTower; (f) is residual_root_character_explicit with reduction_alpha_eq_xi_case_three (hrb, hrb0), reduced_character_quotient_explicit and distinct_four_eigenvalues_injective (hdist), and root_character_ne_one; (g) is the refine of Tower.Ok field by field, with neg_pow_ne_one_of_odd_order (through neg_sq_pow_ne_one_case_three) for (h), residual_rep_trivial_iff_case_three for (i) and (k), and Keep.loc3 (BigOm.lean) for (m); (h) is Four (h5, hSb, the order m, and QuadOn from hK'2, hK'c, hK'd with q > l, hηS, hK'fix and hξF) and TameL with 2 prime to l.
\end{proof}

\begin{lemma}\label{sha-siegel-points-three-seeds}
Assume the Galois representations as in Theorem~\ref{sha-residually-reducible-automorphy}.
Let $D$ be a Siegel datum of case III, let $T_0$ be a tower for $D$ of tensor type that satisfies (a) to (m) (Definitions~\ref{sha-siegel-tower-towers} and~\ref{sha-siegel-points-three-types}), let $\iota:\bar{\mathbf Q}_\ell\cong\mathbf C$, and let $V$ be an ambient ring with residue field $k_V$ (Definition~\ref{sha-siegel-tower-ambient}).
There are a prime $w$ that splits in $T_0$ and a residual flagged pair $s$ for the problem of $T_0$ with auxiliary primes $\{w\}$ along $\mathcal O\to k_V$ (Definition~\ref{sha-siegel-step-chains-notions}), whose homomorphism is supported in $W_{\mathrm m}$, with the following property. For every tower $T\ge T_0$ that satisfies (a) to (m) and kills $w$, there are a ring $\mathcal O_1\supseteq\mathcal O$ of Definition~\ref{sha-character-parts} and an integral seed $Q$ for $T$ over $\mathcal O_1$ (Definition~\ref{sha-siegel-tower-ideal}) whose image under $\mathcal O_1\to V\to k_V$ is over $s$, is of type A, and has entries of weights $(1,-1)$ and $(-1,-1)$.
\end{lemma}
\leanref{Development/AutomorphicArithmetic/AutomorphicVarieties/ModularityArgument/ThresholdArguments/DenseThresholdCases/PromodularityCaseThree.lean}{UpperTheory.V.Cu.SeedIVAlt}
\leanref{Development/AutomorphicArithmetic/AutomorphicVarieties/ModularityArgument/ThresholdArguments/DenseThresholdCases/BirthTop.lean}{UpperTheory.V.Z.Win.Birth.fourth_initial_of_prints_h}
\begin{proof}
Let $K'$ be a tensor field of $T_0$, let $\epsilon$ be its relabelling and $g_0=g_0^{T_0}$, put $\sigma=1$ if $\epsilon=0$ and $\sigma=-1$ if $\epsilon=1$, $\varphi=\xi\psi^\sigma$, and $\Phi_{k'}=\chi_{\mathrm{cyc}}^{k'-1}\varphi$ for $k'\ge1$. For a commutative ring $R$ and $a,b\in M_2(R)$ let
\[
a\boxtimes b=\begin{pmatrix}a_{11}b_{11}&a_{11}b_{12}&-a_{12}b_{12}&a_{12}b_{11}\\a_{11}b_{21}&a_{11}b_{22}&-a_{12}b_{22}&a_{12}b_{21}\\-a_{21}b_{21}&-a_{21}b_{22}&a_{22}b_{22}&-a_{22}b_{21}\\a_{21}b_{11}&a_{21}b_{12}&-a_{22}b_{12}&a_{22}b_{11}\end{pmatrix},
\]
so that $1\boxtimes1=1$ and $(aa')\boxtimes(bb')=(a\boxtimes b)(a'\boxtimes b')$.
\begin{enumerate}
\item[(a)] \emph{The element $g_0$.} It does not fix $K'$: otherwise $\bar\chi^{(1,0)}(g_0)=1$, as $T_0$ is of tensor type, so $\bar\chi_1(g_0)=\bar\chi_2(g_0)$, against (l).
\item[(b)] \emph{The data.} Alpöge constructs the following, over a ring $\mathcal O_1\supseteq\mathcal O$ of Definition~\ref{sha-character-parts}. First, a continuous character $\kappa$ of $G_{K'}$, attached through $\iota$ to an algebraic Hecke character of $K'$, and an integer $m\ge1$, such that for every $j\ge1$ there are a continuous $\tau_j:G\to M_2(\mathcal O_1)$ and signs $s_u=\pm1$ at the places $u$ above $\ell$, with: $\tau_j(g)$ preserves the form $\operatorname{diag}(1,-1)$; $\tau_j(g)=\big(\begin{smallmatrix}a&b\\b&a\end{smallmatrix}\big)$ for $g\in G_{K'}$ and $\tau_j(g)=\big(\begin{smallmatrix}a&b\\-b&-a\end{smallmatrix}\big)$ otherwise, with $a\equiv1$ and $b\equiv0$ modulo $\mathfrak m_{\mathcal O_1}$; $\tau_j(g_0)=\operatorname{diag}(1,-1)$; and $\kappa^j(h)=a+b$ for $h\in G_{K'}$. Second, a prime $w$ that splits in $T_0$, a character $\theta$ with $\theta^\ell=1$ that is trivial on $G_{\mathbf Q(\mu_w)}$, an integer $k'\ge18m+3$ with $k'-1\equiv k$ modulo $\ell-1$, an irreducible representation $\rho$ of $G$ automorphic through $\iota$, with determinant $\chi_{\mathrm{cyc}}^{k'-1}\varphi\theta$ and characteristic polynomials congruent to those of $1\oplus\Phi_{k'}$, and a lattice $\rho_1:G\to M_2(\mathcal O_1)$ in $\rho$ with $\rho_1\equiv\big(\begin{smallmatrix}\Phi_{k'}&0\\\ast&1\end{smallmatrix}\big)$, $\rho_1(g_0)$ diagonal, and elements $x_u\in\mathcal O_1$ at the places $u$ above $\ell$ such that the line spanned by $(1,x_u)$ is stable under $\rho_1(D_u(\mathbf Q))$. Third, an exponent $1\le j\le9$ such that $\rho_1\boxtimes\tau_j$ is automorphic through $\iota$ over the abelian CM fields that do not contain $K'$, at which $\rho$ satisfies a local condition and over which $\rho_1\boxtimes\tau_j$ is irreducible. Fourth, a continuous character $\delta:G\to\mathcal O^\times$ with $\delta^2\Phi_{k'}=\alpha\Psi$, the multiplier of $D$, and $\bar\delta=\bar\chi_4$. For a tower $T\ge T_0$ with field $F$, let $Q_T$ be the pair of the homomorphism $g\mapsto\delta(g)\,\rho_1(g)\boxtimes\tau_j(g)$ on $G_F$ and the flags $\Big(\big(\begin{smallmatrix}1&0\\x_v&1\end{smallmatrix}\big)\boxtimes1\Big)\big(1+s_v(E_{21}-E_{43})\big)$ for $v\in\Sigma_T$.
\item[(c)] \emph{The pair $s$.} Let $s$ be the image of $Q_{T_0}$ under $\mathcal O_1\to V\to k_V$. It is a residual flagged pair for the problem of $T_0$ with auxiliary primes $\{w\}$, whose homomorphism is supported in $W_{\mathrm m}$.
\item[(d)] \emph{The seeds.} Let $T\ge T_0$ satisfy (a) to (m) and kill $w$. Then $\theta$ is trivial on $G_F$, and $Q_T$ is an integral seed for $T$ over $\mathcal O_1$: it is a flagged pair along $\mathcal O\subseteq\mathcal O_1$ with continuous entries, its image in the residue field of $\mathcal O_1$ satisfies the last three conditions of a residual flagged pair, and its image in $\bar{\mathbf Q}_\ell$ is a seed, with inertial weights $h_0+2(k'-1+mj,\,k'-1-mj,\,mj,\,-mj)$ for an integer $h_0$ with $\delta^2=\chi_{\mathrm{cyc}}^{h_0}$ on an open subgroup of $G$. Its image under $\mathcal O_1\to V\to k_V$ is over $s$, is of type A, and has entries of weights $(1,-1)$ and $(-1,-1)$.
\end{enumerate}
\openstep{The construction of the data of (b), with the conditions they satisfy, and the proofs of the properties of the pairs $Q_T$ stated in (c) and (d), are not yet written out.}
% Size: about 4,700 lines of Lean in ModularityArgument and WeakShimuraTheory, with further published-theorem inputs.
% The written part is Cu.fourth_initial_of_rank_four_data (BirthSeedIV.lean, about 50 lines): (a) is hg0 (QuadOn.on and Tower.Ok.dist); (b) is the structure Cu.Dat (OverOf4.lean), with kappa, m and c_0 from Cu.TauExAlt (TauDat, TnOne.lean), w, theta, k', rho, rho_1 and x from Z.Win.Birth.RhoExP with B = 2(9m)+3 (Sd.EisRep for phi D pr and theta, IrrOpen, PR, Cu.Mod4), j from Z.Win.Birth.RankFourAutomorphyOfTensor (AutOffS and OffW for tenB rho_1 tau), tau and s for j, delta from Cu.TwEx4 (Cu.Tw4), and the unramifiedness of theta outside w from theta_unramified_of_kills_roots_of_unity (Sd.S3.ThetaUnr); Q_T is Dat.Q, rep4 (toC h delta) rho_1 tau on G_F with flags flg (x v) (s v) = tn (low x) 1 * nA s, where tn is the matrix a ⊠ b of the text; (c) is PairOf4 (with W = {w}), ResOf4 and ContOf4; (d) is PairOf4 (with W = empty and the determinant through Mod4.det and Dies), ContOf4, ResOf4, tower_is_initial_of_ok (SeedOf4P.lean, from AlNe, WtsOf4, IrrOf4, AutOf4 and OffW) for the seed, OverOf4 and EntOf4. The open parts are the Props TauExAlt (Cu.tauEx'), RhoExP (Cu2.rhoExP), RankFourAutomorphyOfTensor, TwEx4, PairOf4, ResOf4, ContOf4, EntOf4, OverOf4, WtsOf4, AlNe, IrrOf4 and AutOf4, with the hypotheses Hypotheses' and PropEGP (Proposition sha-automorphy-eisenstein-galois).
% SeedIVAlt is proved from HypothesesH' by Z.Win.Birth.fourth_initial_of_prints_h, through Cu.fourth_initial_of_rank_four_data with alNe, Cu.tauEx', Cu.rank_four_automorphy_of_tensor_of_base_change_of_induction (from Z.Bc.upQuadI_of_alt, Z.Bc.upSolI_of_alt, E.indAutCs_of_alt, Cu.polOf4_alt and badOfW), Cu2.rhoExP (Cu2.exists_weight_four_auxiliary_characters, Cu2.phiFacts, Cu2.modEx4), twEx4, pairOf4, resOf4, contOf4, entOf4, overOf4, wtsOf4, irrOf4_of, autOf4_of and propEGP_of_printsH.
\end{proof}

\begin{lemma}\label{sha-siegel-points-three-gysin}
Let $D$ be a Siegel datum, $T$ a fit tower for $D$, $k$ a field and $a:\mathcal O\to k$ a homomorphism that kills $\mathfrak m_{\mathcal O}$, and let $W$ be a finite set of primes that split in $T$ (Definitions~\ref{sha-siegel-tower-towers} and~\ref{sha-siegel-step-chains-notions}). Let $p$ be a residual flagged pair for the problem of $T$ with auxiliary primes $W$ along $a$ whose homomorphism is supported in $W_{\mathrm m}$.
Then there are a finite set $W'\supseteq W$ of primes that split in $T$ and a residual flagged pair $s$ for the problem of $T$ with auxiliary primes $W'$ along $a$, lying in $W_{\mathrm n}$, with the same $(-1,0)$-part as $p$ and with an entry of weight $(0,-1)$.
\end{lemma}
\leanref{Development/AutomorphicArithmetic/AutomorphicVarieties/WeakShimuraTheory/CompactIntegralFields/GysinPs.lean}{UpperTheory.V.Ws.gysinPs}
\leanref{Development/AutomorphicArithmetic/AutomorphicVarieties/WeakShimuraTheory/RamifiedLineCharacters/GysinPsOf.lean}{UpperTheory.V.Ws.Gys.gysinPs_of}
\leanref{Development/AutomorphicArithmetic/AutomorphicVarieties/WeakShimuraTheory/RamifiedLineCharacters/DataAl.lean}{UpperTheory.V.Ws.Gys.dataAl}
\leanref{Development/AutomorphicArithmetic/AutomorphicVarieties/WeakShimuraTheory/RamifiedLineCharacters/CupSwap.lean}{UpperTheory.V.Ws.Gys.cupSwap}
\begin{proof}
Let $F$ be the field of $T$; cocycles and cup cochains are as in Definition~\ref{sha-siegel-step-chains-gysin-cochains}. Write $p=(r_p,(n_{p,v}))$.
\begin{enumerate}
\item[(a)] \emph{The base point.} By (l), the $a(\chi_i(g_0^T))$ are pairwise distinct; so $a(\chi^{(0,-1)}(g_0^T))=a(\chi_3(g_0^T))/a(\chi_1(g_0^T))\ne1$ and $a(\chi^{(-1,-1)}(g_0^T))=a(\chi_4(g_0^T))/a(\chi_1(g_0^T))\ne1$.
\item[(b)] \emph{The $(-1,0)$-part.} Put $x(g)=r_p(g)_{21}a(\chi_1(g))^{-1}$. The entries $(2,3)$ and $(2,4)$ of $r_p$ have the weights $(-1,1)$ and $(0,1)$, outside $W_{\mathrm m}$, so they vanish, and $r_p(gh)_{21}=r_p(g)_{21}a(\chi_1(h))+a(\chi_2(g))r_p(h)_{21}$, as $p$ is residual. So $x$ is a cocycle of weight $(-1,0)$ for the problem of $T$ with auxiliary primes $W$: it is locally constant, it vanishes on $N_\Phi$, where $r_p=1$, and at $g_0^T$, where $r_p$ is diagonal. For $v\in\Sigma_T$ and $s\in D_v(F)$, the matrix $U=n_{p,v}^*r_p(s)n_{p,v}$ is upper triangular with $U_{11}=a(\chi_1(s))$, and $r_p(s)n_{p,v}=n_{p,v}U$, as $n_{p,v}n_{p,v}^*=1$; the entries $(2,1)$ of the two sides give $r_p(s)_{21}+a(\chi_2(s))(n_{p,v})_{21}=(n_{p,v})_{21}a(\chi_1(s))$. So $x$ is locally trivial with constants $m_v=(n_{p,v})_{21}$.
\item[(c)] \emph{Auxiliary primes.} Lemma~\ref{sha-siegel-step-chains-gysin-cup}, with $w=(0,-1)$, $w'=(-1,0)$ and the one cocycle $x$, which applies by (a), gives a finite set $W'\supseteq W$ of primes that split in $T$, a cocycle $y$ of weight $(0,-1)$ for the problem of $T$ with auxiliary primes $W'$, locally trivial with constants $(q_v)$, with $y(g_0^T)=0$ and $y(g)\ne0$ for some $g$, and a cup cochain $z$ for $(y,x)$, with constants $(o_v)$ say, such that $z(g_0^T)=0$; and $x$ is a cocycle for that problem.
\item[(d)] \emph{Exchanging the factors.} Put $z'=yx-z$. A direct computation from the cocycle relations of $x$ and $y$ and the relation of $z$ shows that $z'$ is a cup cochain for $(x,y)$, with constants $m_vq_v-o_v$; and $z'(g_0^T)=0$.
\item[(e)] \emph{The pair.} As in step (d) of the proof of Lemma~\ref{sha-siegel-step-chains-gysin}, $a(\chi_i(s))=1$ for $v\in\Sigma_T$ and $s\in\mu_v$. Lemma~\ref{sha-siegel-step-chains-gysin-matrix}, applied to $(x,y,z')$, gives a residual flagged pair $s=(r_s,(n_{s,v}))$ for the problem of $T$ with auxiliary primes $W'$, lying in $W_{\mathrm n}$, with $r_s(g)_{21}=x(g)a(\chi_1(g))=r_p(g)_{21}$, $(n_{s,v})_{21}=m_v=(n_{p,v})_{21}$ and $r_s(g)_{31}=y(g)a(\chi_1(g))$. So $s$ has the same $(-1,0)$-part as $p$, and, as $y(g)\ne0$ for some $g$, its homomorphism has an entry of weight $(0,-1)$.
\end{enumerate}
% gysinPs is Gys.gysinPs_of (GysinPsOf.lean, 20 lines), the mirror of gysinAl (Lemma sha-siegel-step-chains-gysin): (a) is Gys.rootCharacters_ne_one_at_basepoint; (b) is Gys.dataAl (DataAl.lean, 55 lines, the (-1,0) analogue of Gys.dataPs); (c) is Gys.cupEx with w = -wPs, w' = -wAl and the index set Unit; (d) is Gys.cupSwap (CupSwap.lean, 25 lines); (e) is Gys.mu_of and Gys.ob. About 230 lines with the helpers not already in the book.
\end{proof}

\begin{lemma}\label{sha-siegel-points-three-chains}
Assume the Galois representations as in Theorem~\ref{sha-residually-reducible-automorphy}.
Let $D$ be a Siegel datum of case III, let $T_0$ be a tame tower for $D$ of tensor type that satisfies (a) to (m) (Definitions~\ref{sha-siegel-tower-towers} and~\ref{sha-siegel-points-three-types}), with relabelling $\epsilon$, let $\iota:\bar{\mathbf Q}_\ell\cong\mathbf C$, and let $V$ be an ambient ring with residue field $k_V$ (Definition~\ref{sha-siegel-tower-ambient}).
Let $c_0$ be a flagged pair of lattice shape for the rational problem of $(\epsilon,g_0^{T_0})$ along $\mathcal O\to k_V$ (Definition~\ref{sha-siegel-tower-pairs}).
Then there is a tame tower $T_1\ge T_0$ such that for every deep tower $T\ge T_1$ there are a ring $\mathcal O_1\supseteq\mathcal O$ of Definition~\ref{sha-character-parts} and an integral seed $Q$ for $T$ over $\mathcal O_1$ (Definition~\ref{sha-siegel-tower-ideal}) such that every residual flagged pair for the problem of $T$ along $\mathcal O\to k_V$ that is over $c_0$ is joined by an A-chain (Definition~\ref{sha-siegel-points-three-types}) to the image of $Q$ under $\mathcal O_1\to V\to k_V$.
\end{lemma}
\leanref{Development/AutomorphicArithmetic/AutomorphicVarieties/ModularityArgument/ThresholdArguments/DenseThresholdCases/PromodularityCaseThree.lean}{UpperTheory.V.Cu.RankFourChainReductionDense}
\leanref{Development/AutomorphicArithmetic/AutomorphicVarieties/ModularityArgument/ThresholdArguments/DenseThresholdCases/CuPChainIV.lean}{UpperTheory.V.Cu.chainIV_alt}
\leanref{Development/AutomorphicArithmetic/AutomorphicVarieties/BuildingsResidueChars/ExistsResidualRepRestrictionOfLatticeShape.lean}{UpperTheory.exists_residualRep_restriction_of_latticeShape}
\begin{proof}
Flagged pairs are along $a:\mathcal O\to k_V$, which kills $\mathfrak m_{\mathcal O}$, and we use the notation of Definition~\ref{sha-siegel-step-chains-notions}. The tower $T_0$ is fit, by Lemma~\ref{sha-siegel-points-three-tower}. By the first paragraph of the proof of Lemma~\ref{sha-siegel-step-chains}, the entries $(4,2)$ and $(4,3)$ of the homomorphism $r$ of a residual pair are $-r(g^{-1})_{31}$ and $r(g^{-1})_{21}$ up to the factor $a(\theta(g))$; so $r$ has an entry of weight $(0,-1)$ if and only if some $r(g)_{31}$ is nonzero, and an entry of weight $(-1,0)$ if and only if some $r(g)_{21}$ is nonzero.
\begin{enumerate}
\item[(1)] \emph{A pair over $c_0$.} Being of lattice shape, $c_0$ is residual and lies in $W_{\mathrm r}$. Lemma~\ref{sha-siegel-step-chains-rational}(1) gives a residual flagged pair $p_0$ for the problem of $T_0$ over $c_0$, and by Lemma~\ref{sha-siegel-step-chains-rational}(2) $p_0$ is residual for the problem of $T_0$ with auxiliary primes $W$ for every $W$, and lies in $W_{\mathrm r}$.
\item[(2)] \emph{Seed data.} Lemma~\ref{sha-siegel-points-three-seeds} gives a prime $w$ that splits in $T_0$ and a residual flagged pair $s$ for the problem of $T_0$ with auxiliary primes $\{w\}$, whose homomorphism is supported in $W_{\mathrm m}$, with the property stated there.
\item[(3)] \emph{Auxiliary primes.} Lemma~\ref{sha-siegel-points-three-gysin}, applied to $\{w\}$ and $s$, gives a finite set $W_1\supseteq\{w\}$ of primes that split in $T_0$ and a residual flagged pair $p_2$ for the problem of $T_0$ with auxiliary primes $W_1$, lying in $W_{\mathrm n}$, with the same $(-1,0)$-part as $s$ and an entry of weight $(0,-1)$. Lemma~\ref{sha-siegel-step-chains-gysin}, applied to $W_1$, $p_0$ and $p_2$, which lie in $W_{\mathrm r}$ as $W_{\mathrm n}\subseteq W_{\mathrm r}$, gives a finite set $W_2\supseteq W_1$ of primes that split in $T_0$ and residual flagged pairs $p_3$ and $p_4$ for the problem of $T_0$ with auxiliary primes $W_2$, lying in $W_{\mathrm n}$, such that $p_3$ has the same $(0,-1)$-part as $p_0$, $p_4$ has the same $(0,-1)$-part as $p_2$, $p_3$ and $p_4$ have the same $(-1,0)$-part, and $p_3$ has an entry of weight $(-1,0)$.
\item[(4)] \emph{The tower $T_1$.} Lemma~\ref{sha-siegel-step-chains-dies} gives a tame tower $T_1\ge T_0$ that kills every prime of $W_2$.
\end{enumerate}
Now let $T\ge T_1$ be a deep tower, with field $F$. By Lemma~\ref{sha-siegel-step-towers}(2), $T_0\le T$, and $T$ satisfies (a) to (m) and is fit. As $G_F$ is contained in the group of the field of $T_1$, $T$ kills every prime of $W_2$.
\begin{enumerate}
\item[(5)] \emph{The seed.} By (2), applied to $T$, which kills $w\in W_2$, there are $\mathcal O_1$ and an integral seed $Q$ for $T$ over $\mathcal O_1$ whose image $c_Q$ under $\mathcal O_1\to V\to k_V$ is over $s$, of type A, with entries of weights $(1,-1)$ and $(-1,-1)$. As in step (4) of the proof of Lemma~\ref{sha-siegel-step-chains}, $c_Q$ is residual.
\end{enumerate}
Let $u$ be a residual flagged pair for the problem of $T$ that is over $c_0$.
\begin{enumerate}
\item[(6)] \emph{Restriction to $T$.} The set $N_\Phi$ of the problem of $T_0$ with auxiliary primes $W_2$ is contained in that of the problem with auxiliary primes $\{w\}$, or $W_1$, or $\emptyset$, and the other conditions are the same; so $s$, $p_0$ and $p_2$ are residual for the problem of $T_0$ with auxiliary primes $W_2$, as are $p_3$ and $p_4$. As $\ell\ge5$, Lemma~\ref{sha-siegel-step-chains-restrict}(1) gives residual flagged pairs $q_0$, $q_2$, $q_3$ and $q_4$ for the problem of $T$ over $p_0$, $p_2$, $p_3$ and $p_4$. By Lemma~\ref{sha-siegel-step-chains-restrict}(2)(h), $u=q_0$, since $p_0$ is over $c_0$.
\item[(7)] \emph{Supports and parts.} The towers $T$ and $T_0$ have the same relabelling, so Lemma~\ref{sha-siegel-step-chains-rational}(2), applied to $T$, $g_0^{T_0}$, $c_0$ and $u$, shows that $u$ lies in $W_{\mathrm r}$, and either has entries of weights $(0,-1)$ and $(-1,1)$, or is of type V and has an entry of weight $(-1,0)$. By Lemma~\ref{sha-siegel-step-chains-restrict}(2): $q_3$, $q_4$ and $q_2$ lie in $W_{\mathrm r}$, by (c), and their homomorphisms are supported in $W_{\mathrm n}\subseteq W_{\mathrm m}$, by (b); the homomorphism of $c_Q$ is supported in $W_{\mathrm m}$, by (b), as that of $s$ is; $u$ and $q_3$ have the same $(0,-1)$-part, and so have $q_4$ and $q_2$, by (d); $q_3$ and $q_4$ have the same $(-1,0)$-part, and so have $q_2$ and $c_Q$, by (e); $q_3$ has an entry of weight $(-1,0)$, by (g), and $q_2$ one of weight $(0,-1)$, by (f).
\item[(8)] \emph{A-ampleness.} The pair $u$ is ample, by (7). If $u$ has an entry of weight $(0,-1)$, so has $q_3$, by the first paragraph, and with its entry of weight $(-1,0)$, $q_3$ is ample. Otherwise $u$ is of type V; then $q_3$, which lies in $W_{\mathrm r}$ and has no entry of weight $(0,-1)$, has its homomorphism supported in $W_{\mathrm V}$ and its flags supported in $W_{\mathrm n}$ with the nonzero entries $(3,1)$ of those of $u$, so it is of type V, and with its entry of weight $(-1,0)$ it is ample. The pair $q_4$ has an entry of weight $(0,-1)$, as $q_2$ has, and one of weight $(-1,0)$, as $q_3$ has, so it is ample. The pair $q_2$ is of type A: its homomorphism is supported in $W_{\mathrm n}$, and its entries of weight $(-1,0)$ vanish, as the entries $(2,1)$ are those of $c_Q$, whose homomorphism is supported in $W_{\mathrm A}$; so it is supported in $W_{\mathrm A}$, and the entries $(2,1)$ of its flags are those of $c_Q$, which are nonzero. With its entry of weight $(0,-1)$, $q_2$ is A-ample. The pair $c_Q$ is A-ample, being of type A with entries of weights $(1,-1)$ and $(-1,-1)$.
\item[(9)] \emph{The A-chain.} The pairs $u$ and $q_3$ lie in $W_{\mathrm r}$ and have the same $(0,-1)$-part, so they are linked by Lemma~\ref{sha-siegel-step-chains-lines}(1). The homomorphisms of $q_3$ and $q_4$ are supported in $W_{\mathrm m}$ and the pairs have the same $(-1,0)$-part, so they are linked by Lemma~\ref{sha-siegel-step-chains-lines}(2). In the same way $q_4$ and $q_2$ are linked by Lemma~\ref{sha-siegel-step-chains-lines}(1), and $q_2$ and $c_Q$ by Lemma~\ref{sha-siegel-step-chains-lines}(2). Hence $u,q_3,q_4,q_2,c_Q$ is an A-chain from $u$ to $c_Q$.
\end{enumerate}
% chainIV_alt (CuPChainIV.lean, 88 lines): T_0 fit is hH (Cu.hFit); (1) is exists_residualRep_restriction_of_latticeShape, which Alpöge proves by the steps of Lemma sha-siegel-step-chains-rational (1) (Rq.placeMove, Rq.luG for the flags, Rq.conjFlag, Rq.killQ), with goodQ; (2) is SeedIVAlt (Lemma sha-siegel-points-three-seeds); (3) is gysinPs and gysinAl; (4) is diesEx; the paragraph after (4) is leFit and dies_mono; (5) is the seed with Chn.seed_res_alt; (6) is res_mono and restrT with Four.five, then the last part of overFacts; (7) is goodQ for T and the parts of overFacts; (8) is g0, g3 (samePs_ent, samePs_inV), g4, g2 (RankFourChainReduction.inA_of, with al_zero for the entries (4,3)) and gs; (9) is lineP, lineM, lineP, lineM. Lemma sha-siegel-points-three-seeds is open.
\end{proof}

\begin{lemma}\label{sha-siegel-points-three-descent}
Let $D$ be a Siegel datum.
\begin{enumerate}
\item Let $V$ be an ambient ring with residue field $k_V$ (Definition~\ref{sha-siegel-tower-ambient}), let $\Phi$ be a deformation problem for $D$ (Definition~\ref{sha-siegel-tower-pairs}), and let $c$ and $c'$ be flagged pairs for $\Phi$ along $\mathcal O\to k_V$, with $c$ residual and an A-chain from $c$ to $c'$ (Definition~\ref{sha-siegel-points-three-types}). There is a ring $\mathcal O_1\supseteq\mathcal O$ of Definition~\ref{sha-character-parts} such that for every such ring $\mathcal O'\supseteq\mathcal O_1$, with residue field $k'$, there are flagged pairs $s$ and $t$ for $\Phi$ along $\mathcal O\to k'$ with images $c$ and $c'$ under the map $k'\to k_V$ induced by $\mathcal O'\to V$, with $s$ residual and an A-chain from $s$ to $t$.
\item Let $c$ be a residual flagged pair for a deformation problem along a homomorphism to a ring, and let there be an A-chain from $c$ to $c'$. Then $c'$ is residual and A-ample.
\item Let $T$ be a fit tower for $D$ (Definition~\ref{sha-siegel-tower-towers}), let $X=(A,B,P_X)$ be a deformation space for $T$ over a ring $\mathcal O'$ of Definition~\ref{sha-character-parts}, and let $J\subseteq B$ be an ideal (Definition~\ref{sha-siegel-tower-space}). Suppose that for every point $x$ of $X$ whose image is residual and A-ample, if $J$ lies in some minimal prime below $\ker x$, then it lies in all of them. Let $x_0$ and $x_1$ be points of $X$ such that the image of $x_0$ is residual, there is an A-chain from the image of $x_0$ to the image of $x_1$, and $J$ lies in some minimal prime below $\ker x_1$. Then $J$ lies in all minimal primes below $\ker x_0$.
\end{enumerate}
\end{lemma}
\leanref{Development/AutomorphicArithmetic/AutomorphicVarieties/ModularityArgument/CuspidalTwists/PromodularityCaseThree.lean}{UpperTheory.V.Cu.ChainDesc4}
\leanref{Development/AutomorphicArithmetic/AutomorphicVarieties/ModularityArgument/CuspidalTwists/GoodPairOfInjective.lean}{UpperTheory.V.Cu.chainDesc4}
\leanref{Development/AutomorphicArithmetic/AutomorphicVarieties/ModularityArgument/ThresholdArguments/DenseThresholdCases/PromodularityCaseThreeOfLemmas.lean}{chain_rank_four_good}
\leanref{Development/AutomorphicArithmetic/AutomorphicVarieties/ModularityArgument/ThresholdArguments/DenseThresholdCases/PromodularityCaseThreeOfLemmas.lean}{allIn_of_chain4}
\begin{proof}
(1) The proof of Lemma~\ref{sha-siegel-step-coefficients}(2) applies, with A-chains in place of chains: the conditions that define pairs of type A and A-ample pairs are, like those that define ample pairs, the vanishing or nonvanishing of entries, identities between the values of the characters and the vanishing of entries of the flags, and an injective ring homomorphism preserves and reflects all of them.

(2) Let $c=c_0,\dots,c_m=c'$ be the A-chain. If $m=0$, then $c'=c$ is residual, and A-ample as the first pair of the A-chain. If $m\ge1$, then $c'$ is A-ample by the definition of an A-chain, and residual as the image under $X\mapsto1$ of a residual flagged pair over a polynomial ring.

(3) The proof of Lemma~\ref{sha-siegel-step-linked}(4), which is stated for fit towers, applies, with A-chains and A-ample pairs in place of chains and ample pairs: by (2), every pair of the A-chain is residual and A-ample, and the descending induction uses only Lemma~\ref{sha-siegel-step-linked}(1) and the hypothesis on points whose image is residual and A-ample.
% (1) is Cu.chainDesc4 (GoodPairOfInjective.lean, 205 lines), the copy of Ws.chainDesc with the private good_pair_of_injective and inA_of for the conditions of Good4. (2) is the private chain_rank_four_good, with Asm.linked_isRes. (3) is the private allIn_of_chain4, with conn (Lemma sha-siegel-step-linked (1)), Asm.end_pt and Asm.linked_self; it is stated for problems with Fr.Std, which the problem of a fit tower has (Lemma sha-siegel-step-coefficients-standard).
\end{proof}

\begin{lemma}\label{sha-siegel-points-three-bounds-rank}
Let $D$ be a Siegel datum, $T$ a tower for $D$ with field $F$ (Definition~\ref{sha-siegel-tower-towers}), $k$ a field, $a:\mathcal O\to k$ a homomorphism that kills $\mathfrak m_{\mathcal O}$, and $d,C\ge0$ such that $T$ satisfies the cohomology bounds $(d,C)$ along $a$, with the spaces $Z_w$ and $S_w$ of Definition~\ref{sha-siegel-tower-ideal}(3). Let $W$ be a set of weights, and $c=(r_1,(n_v))$ a pair as in Definition~\ref{sha-siegel-step-dimensions-tangent-pairs} with $r_1(g)=\operatorname{diag}(a(\chi_1(g)),\dots,a(\chi_4(g)))$ and each $n_v$ symplectic with multiplier $1$. Suppose that some entry $(i_0,j_0)$ whose weight is an odd root has the following property: for every $(x,(Y_v))\in\mathrm{Tan}_W(c)$ such that $x(g)_{i_0j_0}$, $x(g)_{11}$ and $x(g)_{22}$ vanish for all $g$, and every entry $(i,j)$ with $i+j\le5$ whose weight $w$ is an odd root, there are $m_v\in k$, for $v\in\Sigma_T$, with $x(s)_{ij}=m_v-a(\chi^w(s))m_v$ for all $v$ and $s\in D_v(F)$. Then $\dim\mathrm{Tan}_W(c)\le d+20C$.
\end{lemma}
\leanref{Development/AutomorphicArithmetic/AutomorphicVarieties/ModularityArgument/CuspidalTwists/PRedTanA.lean}{UpperTheory.V.Cu.ReductionTangentLimit.rank_A}
\begin{proof}
\begin{enumerate}
\item[(a)] \emph{Entries.} By step (a) of the proof of Lemma~\ref{sha-siegel-step-dimensions-tangent-split}, for $(x,(Y_v))\in\mathrm{Tan}_W(c)$ and an entry $(i,j)$ of weight $w$, the map $g\mapsto x(g)_{ij}$ lies in $Z_w$.
\item[(b)] \emph{Three entries.} The entries $(1,1)$ and $(2,2)$ have weight $0$; let $w_0$ be the weight of $(i_0,j_0)$. Let $E$ be the linear map from $\mathrm{Tan}_W(c)$ to $Z_{w_0}\times Z_0\times Z_0$ that sends $(x,(Y_v))$ to the maps $g\mapsto x(g)_{i_0j_0}$, $g\mapsto x(g)_{11}$ and $g\mapsto x(g)_{22}$. As $w_0$ is an odd root, its image has dimension at most $(d+C)+C+C$.
\item[(c)] \emph{The kernel of $E$.} Let $P$, $Q$ and the spaces $M_{ij}$ be as in step (c) of the proof of Lemma~\ref{sha-siegel-step-dimensions-tangent-split}, so that $\dim M_{ij}\le C$. By (a) and the hypothesis, for $(x,(Y_v))$ in the kernel of $E$ and $(i,j)\in P$, the map $g\mapsto x(g)_{ij}$ lies in $M_{ij}$. The linear map from the kernel of $E$ to $\prod_{(i,j)\in P}M_{ij}\times k^{\Sigma_T\times Q}$ that sends $(x,(Y_v))$ to the maps $g\mapsto x(g)_{ij}$ for $(i,j)\in P$ and the entries $(n_v^*Y_vn_v)_{ij}$ for $v\in\Sigma_T$ and $(i,j)\in Q$ is injective, by the argument of that step. So the kernel of $E$ has dimension at most $10C+6|\Sigma_T|\le16C$.
\item[(d)] \emph{The bound.} By (b) and (c), $\dim\mathrm{Tan}_W(c)\le d+3C+16C\le d+20C$.
\end{enumerate}
\end{proof}
% rank_A (PRedTanA.lean, about 95 lines): (a) is ent, with Ok.conj_split; (b) is E3, with r1 (hb.odd and hb.even twice); (c) is M, mem, F and hF (with the sp4 symmetry through adj), r2 and r3 (hb.str, hb.even, hb.card); (d) is r4 and the count d + C + (C + C) + (10C + 6C) <= d + 20C. It is rank_of_oneOdd (Lemma sha-siegel-step-dimensions-tangent-split) with the two factors for the entries (1,1) and (2,2) added to E3.

\begin{lemma}\label{sha-siegel-points-three-bounds-cases}
Let $D$, $T$, $F$, $k$, $a$, $d$ and $C$ be as in Lemma~\ref{sha-siegel-points-three-bounds-rank}; let $W_1,\dots,W_9$ be the sets of weights of Lemma~\ref{sha-siegel-step-dimensions}, $r_1(g)=\operatorname{diag}(a(\chi_1(g)),\dots,a(\chi_4(g)))$, and $E_{ij}$ the matrix units. Let $b_v\in k^\times$ for $v\in\Sigma_T$, and suppose that $a(\chi_1(s))=a(\chi_2(s))$ for every $v$ and $s\in D_v(F)$. Then $1+b_v(E_{21}-E_{43})$ is symplectic with multiplier $1$, and if $W$ is $W_1$, $W_4$, $W_5$, $W_7$ or $W_8$, the pair $(r_1,(1+b_v(E_{21}-E_{43})))$ satisfies the hypothesis of Lemma~\ref{sha-siegel-points-three-bounds-rank}, with $(i_0,j_0)=(2,3)$ for $W_1$ and $W_7$, $(1,4)$ for $W_4$ and $W_8$, and $(1,2)$ for $W_5$.
\end{lemma}
\leanref{Development/AutomorphicArithmetic/AutomorphicVarieties/ModularityArgument/CuspidalTwists/PRedTanA.lean}{UpperTheory.V.Cu.ReductionTangentLimit.oneOddA}
\leanref{Development/AutomorphicArithmetic/AutomorphicVarieties/ModularityArgument/CuspidalTwists/PRedTanA.lean}{UpperTheory.V.Cu.ReductionTangentLimit.keyA}
\leanref{Development/AutomorphicArithmetic/AutomorphicVarieties/ModularityArgument/CuspidalTwists/PRedTanA.lean}{UpperTheory.V.Cu.ReductionTangentLimit.rel_A}
\leanref{Development/AutomorphicArithmetic/AutomorphicVarieties/ModularityArgument/CuspidalTwists/PRedTanA.lean}{UpperTheory.V.Cu.ReductionTangentLimit.gsp_A}
\begin{proof}
Put $X=E_{21}-E_{43}$ and $n_v=1+b_vX$. Then $X^*=-X$ and $X^2=0$, so $n_v^*n_v=(1-b_vX)(1+b_vX)=1$.

Let $c=(r_1,(n_v))$, let $(x,(Y_v))\in\mathrm{Tan}_W(c)$ with $x(g)_{i_0j_0}=x(g)_{11}=x(g)_{22}=0$ for all $g$, and call an entry $(i,j)$ whose weight $w$ is an odd root \emph{trivial} if there are $m_v\in k$ with $x(s)_{ij}=m_v-a(\chi^w(s))m_v$ for all $v$ and $s\in D_v(F)$. The entries with $i+j\le5$ whose weights are odd roots are $(1,2)$, $(1,4)$, $(2,1)$, $(2,3)$, $(3,2)$ and $(4,1)$, of weights $(1,0)$, $(1,1)$, $(-1,0)$, $(-1,1)$, $(1,-1)$ and $(-1,-1)$. We show that they are trivial.
\begin{enumerate}
\item[(a)] \emph{Vanishing entries.} An entry $(i,j)$ with $i\ne j$ whose weight is not in $W$ has $x(g)_{ij}=0$ for all $g$, as in the proof of Lemma~\ref{sha-siegel-step-dimensions-tangent-cases}. So $x(g)_{14}=0$ for all $g$: $(i_0,j_0)=(1,4)$ for $W_4$ and $W_8$, and the weight $(1,1)$ is not in $W_1$, $W_5$ or $W_7$. Likewise $x(g)_{23}=0$ for all $g$: $(i_0,j_0)=(2,3)$ for $W_1$ and $W_7$, and $(-1,1)$ is not in $W_4$, $W_5$ or $W_8$. And $x(g)_{12}=0$ for all $g$, or $x(g)_{21}=0$ for all $g$: $(i_0,j_0)=(1,2)$ for $W_5$, $(1,0)$ is not in $W_1$ or $W_7$, and $(-1,0)$ is not in $W_4$ or $W_8$.
\item[(b)] \emph{Relations.} Let $v\in\Sigma_T$ and $s\in D_v(F)$, write $b=b_v$, and let $Z=x(s)+r_1(s)Y_vr_1(s^{-1})-Y_v$, so that $n_v^*Zn_v=(1-bX)Z(1+bX)$ is upper triangular. Its entries $(3,2)$, $(3,1)$, $(4,2)$, $(4,1)$ and $(2,1)$ are $Z_{32}$, $Z_{31}+bZ_{32}$, $Z_{42}+bZ_{32}$, $Z_{41}+b(Z_{31}+Z_{42})+b^2Z_{32}$ and $Z_{21}+b(Z_{22}-Z_{11})-b^2Z_{12}$. As they vanish,
\[
Z_{21}=b^2Z_{12}+b(Z_{11}-Z_{22}),\qquad Z_{32}=0,\qquad Z_{41}=0.
\]
The entry $(i,j)$ of $Z$ is $x(s)_{ij}-\big((Y_v)_{ij}-a(\chi^w(s))(Y_v)_{ij}\big)$ for the weight $w$ of $(i,j)$. As $a(\chi^{(\pm1,0)}(s))=a(\chi_1(s)\chi_2(s)^{-1})^{\pm1}=1$ and $\chi^0=1$, $Z_{12}=x(s)_{12}$, $Z_{21}=x(s)_{21}$, $Z_{11}=x(s)_{11}=0$ and $Z_{22}=x(s)_{22}=0$. So $x(s)_{21}=b_v^2x(s)_{12}$, and $x(s)_{ij}=(Y_v)_{ij}-a(\chi^w(s))(Y_v)_{ij}$ for $(i,j)$ equal to $(3,2)$ or $(4,1)$.
\item[(c)] \emph{Conclusion.} By (a) and the relation $x(s)_{21}=b_v^2x(s)_{12}$ of (b), with $b_v\ne0$, $x(s)_{12}=x(s)_{21}=0$ for every $v$ and $s\in D_v(F)$. So $(1,2)$, $(1,4)$, $(2,1)$ and $(2,3)$ are trivial with $m_v=0$, and $(3,2)$ and $(4,1)$ are trivial with $m_v=(Y_v)_{32}$ and $m_v=(Y_v)_{41}$, by (b).
\end{enumerate}
\end{proof}
% oneOddA (PRedTanA.lean): keyA, a decide over redTypes, gives (i_0,j_0); the text's choices are ones that keyA accepts. (a) is zeroA (through z1), g03, g12 and the case split of gα; (b) is rel_A (the five entries h10, h21, h20, h31, h30, checked by linear_combination), ent (Ok.conj_split) and rt_oneA; (c) is gα and the final case analysis pos. gsp_A is the symplecticity of n_v. The Lean indexes from 0, so its entries (1,0), (2,1), (3,0), (0,1), (0,3), (1,2) are the entries (2,1), (3,2), (4,1), (1,2), (1,4), (2,3) of the text, and its xMA is E_21 - E_43.

\begin{lemma}\label{sha-siegel-points-three-bounds-tangent}
Let $D$ be a Siegel datum, $T$ a fit tower for $D$ with field $F$ (Definition~\ref{sha-siegel-tower-towers}), $k$ a field, $a:\mathcal O\to k$ a homomorphism, and $d,C\ge0$ such that $T$ satisfies the cohomology bounds $(d,C)$ along $a$ (Definition~\ref{sha-siegel-tower-ideal}). Let $c=(r_c,(n_{c,v}))$ be a residual flagged pair for the problem of $T$ along $a$ that is of type A (Definition~\ref{sha-siegel-points-three-types}). Let $W$ be one of the sets $W_1,\dots,W_9$ of Lemma~\ref{sha-siegel-step-dimensions} that contains $(0,-1)$, or that contains $(1,-1)$ and $(-1,-1)$ and is not $W_9$, and suppose that $r_c$ is supported in $W$. Then the space $\mathrm{Tan}_W(c)$ of Definition~\ref{sha-siegel-step-dimensions-tangent-pairs} has dimension at most $d+20C$. (That $r_c$ is supported in $W$ is not used in the proof.)
\end{lemma}
\leanref{Development/AutomorphicArithmetic/AutomorphicVarieties/ModularityArgument/CuspidalTwists/PRedTanA.lean}{UpperTheory.V.Cu.redTanA}
\leanref{Development/AutomorphicArithmetic/AutomorphicVarieties/ModularityArgument/CuspidalTwists/PRedTanA.lean}{UpperTheory.V.Cu.ReductionTangentLimit.characters_agree_at_local_of_typeA}
\begin{proof}
As $T$ satisfies the cohomology bounds along $a$, $a$ kills $\mathfrak m_{\mathcal O}$. Let $r_1(g)=\operatorname{diag}(a(\chi_1(g)),\dots,a(\chi_4(g)))$, let $E_{ij}$ be the matrix units, and let $b_v=(n_{c,v})_{21}$, which is nonzero as $c$ is of type A.
\begin{enumerate}
\item[(a)] \emph{The characters at $v$.} Let $v\in\Sigma_T$ and $s\in D_v(F)$. Then $a(\chi_1(s))=a(\chi_2(s))$. Indeed, as $n_{c,v}^*=n_{c,v}^{-1}$, $r_c(s)n_{c,v}=n_{c,v}U$ with $U=n_{c,v}^*r_c(s)n_{c,v}$ upper triangular and $U_{11}=a(\chi_1(s))$, as $c$ is residual. The entries $(2,1)$, $(2,3)$ and $(2,4)$ of $r_c(s)$ vanish, as their weights $(-1,0)$, $(-1,1)$ and $(0,1)$ are not in $W_{\mathrm A}$, and $r_c(s)_{22}=a(\chi_2(s))$. So comparing the entries $(2,1)$ of $r_c(s)n_{c,v}$ and $n_{c,v}U$ gives $a(\chi_2(s))b_v=b_va(\chi_1(s))$, and $b_v\ne0$.
\item[(b)] \emph{The flow.} Let $m=(0,0,1,1)$, so that $m_i+m_{5-i}=1$, and use the notation of Lemma~\ref{sha-siegel-step-dimensions-tangent-flow}. The entries of weights in $W_{\mathrm A}=\{(1,-1),(0,-1),(-1,-1)\}$ are $(3,2)$, $(3,1)$, $(4,2)$ and $(4,1)$, with $m_i>m_j$. As $r_c$ is supported in $W_{\mathrm A}$ and $c$ is residual, $r_c(g)$ is of degree $\ge0$ with $r_c(g)^{(0)}=r_1(g)$. The entries $(i,j)$ with $m_i<m_j$ are above the diagonal, so they vanish in the lower triangular matrix $n_{c,v}$, which is of degree $\ge0$. Its entries of degree $0$ are the diagonal entries $1$, the entries $(1,2)$ and $(3,4)$, which vanish, the entry $(2,1)$, which is $b_v$, and the entry $(4,3)$, which is $-b_v$: comparing the entries $(2,1)$ of $n_{c,v}^*n_{c,v}=1$ gives $(n_{c,v})_{43}+(n_{c,v})_{21}=0$. So $n_{c,v}^{(0)}=1+b_v(E_{21}-E_{43})$, and Lemma~\ref{sha-siegel-step-dimensions-tangent-flow} gives $\dim\mathrm{Tan}_W(c)\le\dim\mathrm{Tan}_W(r_1,(1+b_v(E_{21}-E_{43})))$.
\item[(c)] \emph{The bound.} The sets $W_s$ that contain $(0,-1)$ are $W_1$, $W_4$, $W_5$, $W_7$ and $W_8$, and those that contain $(1,-1)$ and $(-1,-1)$ are $W_1$, $W_4$, $W_5$ and $W_9$; so $W$ is one of $W_1$, $W_4$, $W_5$, $W_7$ and $W_8$. By (a) and Lemma~\ref{sha-siegel-points-three-bounds-cases}, the pair $(r_1,(1+b_v(E_{21}-E_{43})))$ satisfies the hypothesis of Lemma~\ref{sha-siegel-points-three-bounds-rank}, which gives $\dim\mathrm{Tan}_W(r_1,(1+b_v(E_{21}-E_{43})))\le d+20C$. With (b), this is the claim.
\end{enumerate}
\end{proof}
% redTanA (PRedTanA.lean, 10 lines): (a) is characters_agree_at_local_of_typeA (hp.low for n* n = 1, hp.flag, h.char, h.diag and hA.1 for the vanishing entries); (b) is tan_flow with m = ![0,0,1,1] (Symp by decide), limRA with hA.1 and limA (h32 for the entry (4,3)); (c) is rank_A with gsp_A and oneOddA. The hypothesis hs (r_c supported in W) is introduced and not used.

\begin{lemma}\label{sha-siegel-points-three-bounds-torus}
Let $K$ be a field, $H$ a group, $r:H\to M_4(K)$ a homomorphism, and $h_0\in H$ with $r(h_0)$ diagonal with pairwise distinct diagonal entries. Suppose that $r$ is supported in $W_9=\{\pm(-1,1),\pm(1,1)\}$ and has entries of each of the four weights of $W_9$ (Definitions~\ref{sha-siegel-datum} and~\ref{sha-siegel-tower-pairs}). Then there are $s,t\in K$ such that every subfield $E\subseteq K$ that contains $s$, $t$, the traces of the $r(h)$ and the diagonal entries of $r(h_0)$ contains every entry of every $r(h)$.
\end{lemma}
\leanref{Development/AutomorphicArithmetic/AutomorphicVarieties/ModularityArgument/CuspidalTwists/PTorusY.lean}{UpperTheory.V.Cu.torusY}
\begin{proof}
The off-diagonal entries of weights $(-1,1)$, $(1,-1)$, $(1,1)$ and $(-1,-1)$ are $(2,3)$, $(3,2)$, $(1,4)$ and $(4,1)$. So the other off-diagonal entries of every $r(h)$ vanish, and for each of these four entries $(i,j)$ there is $g_{ij}\in H$ with $r(g_{ij})_{ij}\ne0$. Put $s=r(g_{14})_{14}$ and $t=r(g_{23})_{23}$. Let $E$ be as in the statement, $d_1,\dots,d_4$ the diagonal entries of $r(h_0)$, $E_{ij}$ the matrix units, and $\Lambda$ the $E$-subspace of $M_4(K)$ spanned by the $r(h)$. It contains $1=r(1)$, it is closed under products, as $r(g)r(h)=r(gh)$, and the traces of its elements lie in $E$.
\begin{enumerate}
\item[(a)] Each $E_{ii}$ lies in $\Lambda$: it is the product of the matrices $(d_i-d_j)^{-1}(r(h_0)-d_j)$ over $j\ne i$, which lie in $\Lambda$.
\item[(b)] For $h\in H$ and all $i,j$, $r(h)_{ij}E_{ij}=E_{ii}r(h)E_{jj}$ lies in $\Lambda$, by (a).
\item[(c)] If $c\in E$ is nonzero and $cE_{ij}\in\Lambda$, then $r(h)_{ji}\in E$ for every $h$: by (b), $cE_{ij}\cdot r(h)_{ji}E_{ji}$ lies in $\Lambda$, and its trace is $c\,r(h)_{ji}$.
\item[(d)] By (b) and (c), applied to $c=s$ and $(i,j)=(1,4)$, and to $c=t$ and $(i,j)=(2,3)$, the entries $r(h)_{41}$ and $r(h)_{32}$ lie in $E$ for every $h$. Applied then to $c=r(g_{41})_{41}$ and $(i,j)=(4,1)$, and to $c=r(g_{32})_{32}$ and $(i,j)=(3,2)$, which are nonzero elements of $E$, they show that the $r(h)_{14}$ and $r(h)_{23}$ lie in $E$. By (a) and (c), applied to $c=1$ and $(i,i)$, the $r(h)_{ii}$ lie in $E$. The other entries vanish.
\end{enumerate}
\end{proof}
% torusY (PTorusY.lean, 128 lines): the four entries are wt_be, wt_mbe, wt_th, wt_mth and wt_off, with gs, gt, g30, g21 from ent; Lambda is Submodule.span E (Set.range r), with span_mul and span_trace; (a) is eii; (b) is hent (single_mul_mul_single); (c) is hQE; (d) is e30, e21, e03, e12 and ed. Unlike Lemma sha-siegel-step-dimensions-torus (1), it does not use the multipliers.

\begin{lemma}\label{sha-siegel-points-three-bounds}
Let $D$ be a Siegel datum, let $T$ be a fit tower for $D$ (Definition~\ref{sha-siegel-tower-towers}), let $\mathcal O'\supseteq\mathcal O$ be a ring of Definition~\ref{sha-character-parts}, with residue field $k'$, and let $W_1,\dots,W_9$ and the ideals $I_W$ be as in Lemma~\ref{sha-siegel-step-dimensions}, so that $W_9=\{\pm(-1,1),\pm(1,1)\}$.
\begin{enumerate}
\item Let $d,C\ge0$ be such that $T$ satisfies the cohomology bounds $(d,C)$ along $\mathcal O\to k'$ (Definition~\ref{sha-siegel-tower-ideal}), let $c$ be a residual flagged pair for the problem of $T$ along $\mathcal O\to k'$ that is of type A (Definition~\ref{sha-siegel-points-three-types}) and whose homomorphism has an entry of weight $(0,-1)$ or entries of weights $(1,-1)$ and $(-1,-1)$, and let $P=(r_P,(n_v))$ be a universal lift of $c$ to $R$ (Definition~\ref{sha-siegel-tower-pairs}). Then $\dim R/I_{W_s}(r_P)\le d+20C+1$ for $s=1,\dots,8$.
\item Let $X=(A,B,P_X)$ be a deformation space for $T$ over $\mathcal O'$ (Definition~\ref{sha-siegel-tower-space}), let $c$ be residual, and let $f:A\to R$ be a homomorphism of $\mathcal O'$-algebras that carries $P_X$ to a universal lift of $c$. Let $\mathfrak q$ be a prime of $R$ such that the image $r'$ of $f\circ r_X$ over the fraction field of $R/\mathfrak q$ is supported in $W_9$ and has entries of the four weights $(-1,1)$, $(1,-1)$, $(1,1)$ and $(-1,-1)$, and let $\mathfrak q_B$ be the inverse image of $\mathfrak q$ in $B$. Then $\dim R/\mathfrak q\le\dim B/\mathfrak q_B+2+4|\Sigma_T|$.
\item Let $K$ be a field, $a:\mathcal O\to K$ a homomorphism, and $P=(r,(n_v))$ a flagged pair for the problem of $T$ along $a$ with $r$ supported in $W_9$, with diagonal entries $\psi_{v,i}$. Let $v\in\Sigma_T$ be such that $(n_v)_{21}\ne0$. Then $\psi_{v,1}(s)=\psi_{v,2}(s)$ for all $s\in D_v(F)$, or $\psi_{v,2}(s)=\psi_{v,4}(s)$ for all $s\in D_v(F)$, where $F$ is the field of $T$.
\end{enumerate}
\end{lemma}
\leanref{Development/AutomorphicArithmetic/AutomorphicVarieties/ModularityArgument/CuspidalTwists/PromodularityCaseThree.lean}{UpperTheory.V.Cu.RedA}
\leanref{Development/AutomorphicArithmetic/AutomorphicVarieties/ModularityArgument/CuspidalTwists/PRedA.lean}{UpperTheory.V.Cu.redA_of}
\leanref{Development/AutomorphicArithmetic/AutomorphicVarieties/ModularityArgument/CuspidalTwists/PromodularityCaseThree.lean}{UpperTheory.V.Cu.FibY}
\leanref{Development/AutomorphicArithmetic/AutomorphicVarieties/ModularityArgument/CuspidalTwists/PFibY.lean}{UpperTheory.V.Cu.fibY_of}
\leanref{Development/AutomorphicArithmetic/AutomorphicVarieties/ModularityArgument/CuspidalTwists/PromodularityCaseThree.lean}{UpperTheory.V.Cu.YCoin}
\leanref{Development/AutomorphicArithmetic/AutomorphicVarieties/ModularityArgument/CuspidalTwists/PYCoin.lean}{UpperTheory.V.Cu.yCoin}
\begin{proof}
(1) Let $s\le8$, $I=I_{W_s}(r_P)$ and $c=(r_c,(n_{c,v}))$. If $I=R$ there is nothing to prove. Otherwise let $t$ be the $k'$-space of $\mathcal O'$-linear maps $\delta:R/I\to k'$ with $\delta(xy)=\bar x\delta(y)+\bar y\delta(x)$. Steps (a) to (c) of the proof of Lemma~\ref{sha-siegel-step-dimensions}(1) use only that $P$ is a universal lift of $c$, and they apply without change: $r_c$ is supported in $W_s$, and $\delta\mapsto(x_\delta,(Y_{\delta,v}))$ is an injective $k'$-linear map from $t$ to $\mathrm{Tan}_{W_s}(c)$. As $r_c$ is supported in $W_s$, every nonzero weight of an entry of $r_c$ lies in $W_s$. So $W_s$ contains $(0,-1)$ if $r_c$ has an entry of that weight, and contains $(1,-1)$ and $(-1,-1)$ otherwise; and $W_s\ne W_9$, as $s\le8$. By Lemma~\ref{sha-siegel-points-three-bounds-tangent}, applied to $\mathcal O\to k'$, $\dim t\le d+20C$, and Lemma~\ref{sha-siegel-step-dimensions-generators} gives $\dim R/I\le d+20C+1$.

(2) Let $\mathfrak p=f^{-1}(\mathfrak q)$. Steps (a) and (b) of the proof of Lemma~\ref{sha-siegel-step-dimensions}(2) apply without change, and give $\dim R/\mathfrak q\le\dim B/\mathfrak q_B+\tau$, where $\tau$ is the transcendence degree of the fraction field of $A/\mathfrak p$ over that of $B/\mathfrak q_B$. Step (c) applies with Lemma~\ref{sha-siegel-points-three-bounds-torus} in place of Lemma~\ref{sha-siegel-step-dimensions-torus}(1). Indeed, let $\Omega_{\mathfrak q}$ be the fraction field of $R/\mathfrak q$ and $y:A\to\Omega_{\mathfrak q}$ the composite map, so that $r'=y\circ r_X$. The matrix $y(r_X(g_0^T))$ is diagonal with pairwise distinct diagonal entries, as there, and $r'$ is supported in $W_9$ and has entries of the four weights of $W_9$. Let $s,t$ be as in Lemma~\ref{sha-siegel-points-three-bounds-torus}, and let $T_0$ consist of $s$, $t$ and the entries $(2,1)$, $(3,1)$, $(4,1)$ and $(3,2)$ of the $y(n_{X,v})$, at most $2+4|\Sigma_T|$ elements. The field $E$ of the elements of $\Omega_{\mathfrak q}$ algebraic over the subfield generated by $y(B)$ and $T_0$ contains the traces of the $y(r_X(g))$, which are coefficients of the $y(c_g)$, and the diagonal entries of $y(r_X(g_0^T))$, which are roots of $y(c_{g_0^T})$. By Lemmas~\ref{sha-siegel-points-three-bounds-torus} and~\ref{sha-siegel-step-dimensions-torus}(2), $E$ contains the entries of the $y(r_X(g))$ and of the $y(n_{X,v})$, which generate $y(A)$ over $y(B)$. So $\tau\le2+4|\Sigma_T|$, which gives (2).

(3) Let $s\in D_v(F)$, put $U=n_v^*r(s)n_v$, which is upper triangular with diagonal entries $\psi_{v,i}(s)$, and $\nu=a(\theta(s))\ne0$. As $n_v$ is symplectic with multiplier $1$, $n_vn_v^*=1$, so $r(s)n_v=n_vU$, and $U$ is symplectic with multiplier $\nu$; so $U_{11}U_{44}=U_{22}U_{33}=\nu$, and $U_{11}\ne0$. Comparing the first columns of $r(s)n_v=n_vU$, with the vanishing entries of $r(s)$ of weight outside $W_9$, the lower unitriangular shape of $n_v$ and the diagonal entries of $r(s)^*r(s)=\nu$, gives
\[
U_{11}\big(U_{11}+U_{44}-r(s)_{11}-r(s)_{44}\big)=0,\qquad(n_v)_{21}U_{11}\big(U_{11}+U_{44}-r(s)_{22}-r(s)_{33}\big)=0.
\]
As $U_{11}\ne0$ and $(n_v)_{21}\ne0$, $U_{11}+U_{44}$ is $r(s)_{11}+r(s)_{44}$ and $r(s)_{22}+r(s)_{33}$; as $U$ and $r(s)$ have the same trace, $U_{22}+U_{33}=U_{11}+U_{44}$. Hence
\[
\big(U_{22}-U_{11}\big)\big(U_{22}-U_{44}\big)=U_{22}^2-U_{22}(U_{11}+U_{44})+U_{11}U_{44}=U_{22}^2-U_{22}(U_{22}+U_{33})+U_{22}U_{33}=0.
\]
So for every $s\in D_v(F)$, $\psi_{v,2}(s)=\psi_{v,1}(s)$ or $\psi_{v,2}(s)=\psi_{v,4}(s)$, and $\psi_{v,i}(s)\ne0$. The $\psi_{v,i}$ are multiplicative on $D_v(F)$, as the matrices $n_v^*r(s)n_v$ are upper triangular and $n_vn_v^*=1$. Suppose that $x,y\in D_v(F)$ have $\psi_{v,1}(x)\ne\psi_{v,2}(x)$ and $\psi_{v,2}(y)\ne\psi_{v,4}(y)$, so that $\psi_{v,2}(x)=\psi_{v,4}(x)$ and $\psi_{v,2}(y)=\psi_{v,1}(y)$. If $\psi_{v,2}(xy)=\psi_{v,1}(xy)$, then $\psi_{v,2}(x)\psi_{v,2}(y)=\psi_{v,1}(x)\psi_{v,2}(y)$, so $\psi_{v,2}(x)=\psi_{v,1}(x)$; if $\psi_{v,2}(xy)=\psi_{v,4}(xy)$, then $\psi_{v,4}(x)\psi_{v,2}(y)=\psi_{v,4}(x)\psi_{v,4}(y)$, so $\psi_{v,2}(y)=\psi_{v,4}(y)$. Both contradict the choice of $x$ and $y$; so no such pair exists, which is the claim.
% (1) is Cu.redA_of (PRedA.lean, 150 lines), the copy of Ws.red (Lemma sha-siegel-step-dimensions (1)) with redTanA (Lemma sha-siegel-points-three-bounds-tangent) in place of redTan: the trivial case is htop; (a) to (c) are lifts_map, lc_dual, hom_of_der, redDual (hmem) and hinj; the weights of W_s are hmemS and hcond; (d) is redTanA and redDim. It is stated for every problem with Fr.Std. (2) is Cu.fibY_of (PFibY.lean, 108 lines), the copy of Ws.fib (Lemma sha-siegel-step-dimensions (2)) with fibCpl, CA.c3b, CA.c3c and fibTrdeg as there, and the private fib_gen of PFibY.lean, which calls torusY (Lemma sha-siegel-points-three-bounds-torus) in place of fibTorus and its own copy of fib_low (Lemma sha-siegel-step-dimensions-torus (2)). With their shared dependencies, about 2,000 and 750 lines. (3) is Cu.yCoin (PYCoin.lean): the matrix computation is YCoincidenceMatrices.key (c0 to c3, dA, dB, p0, p1, hA, hB, htr), multiplicativity is psi_mul, and the last step is the final by_contra.
\end{proof}

\begin{lemma}\label{sha-siegel-points-three-propagation}
Let $D$ be a Siegel datum of case III, let $T$ be a fit deep tower for $D$ (Definition~\ref{sha-siegel-tower-towers}), let $\iota:\bar{\mathbf Q}_\ell\cong\mathbf C$, let $\mathcal O'\supseteq\mathcal O$ be a ring of Definition~\ref{sha-character-parts}, with residue field $k'$, and let $X=(A,B,P_X)$ be a deformation space for $T$ over $\mathcal O'$ with a modularity ideal $J$ (Definitions~\ref{sha-siegel-tower-space} and~\ref{sha-siegel-tower-ideal}).
Let $C\ge0$ be such that $T$ satisfies the cohomology bounds $(d_T,C)$ along $\mathcal O\to k'$ and $30(C+1)^2\le d_T$.
Let $x$ be a point of $X$ whose image is residual and A-ample (Definition~\ref{sha-siegel-points-three-types}). If $J$ lies in some minimal prime below $\ker x$, then it lies in all minimal primes below $\ker x$.
\end{lemma}
\leanref{Development/AutomorphicArithmetic/AutomorphicVarieties/ModularityArgument/ThresholdArguments/DenseThresholdCases/PromodularityCaseThree.lean}{UpperTheory.V.Cu.ChainRankFour}
\leanref{Development/AutomorphicArithmetic/AutomorphicVarieties/ModularityArgument/ThresholdArguments/DenseThresholdCases/ReductionTypeChainOfRedAYCoinFibY.lean}{UpperTheory.V.Cu.chain_rank_four_of_red_a_y_coin_fib_y}
\begin{proof}
Write $d=d_T$ and let $c$ be the image of $x$. Lemmas~\ref{sha-siegel-step-completion}, \ref{sha-siegel-step-dimensions}, \ref{sha-siegel-step-irreducible} and~\ref{sha-siegel-step-propagation} apply to $T$, which is fit and deep.
If $c$ is ample, Lemma~\ref{sha-siegel-step-propagation} is the claim. Otherwise $c$ is of type A, and its homomorphism has an entry of weight $(0,-1)$ or entries of weights $(1,-1)$ and $(-1,-1)$. Steps (1) to (3) of the proof of Lemma~\ref{sha-siegel-step-propagation} apply without change: they give $R$ and $f:A\to R$, and reduce the claim to the following, with $N=2d+1-2|\Sigma_T|$ and $\varphi:B\to A$ the structure map. Let $\mathfrak P'$ and $\mathfrak Q'$ be minimal primes of $R$ with $\dim R/(\mathfrak P'+\mathfrak Q')\ge N+1$ and $f(\varphi(J))\subseteq\mathfrak P'$; then $f(\varphi(J))\subseteq\mathfrak Q'$.
\begin{enumerate}
\item[(4)] Let $\varpi$ be a uniformizer of $\mathcal O'$, let $\mathfrak a_s=I_{W_s}(f\circ r_X)$ for $1\le s\le9$, and let $\mathfrak b_{v,i,j}$ be the ideals of step (4) of the proof of Lemma~\ref{sha-siegel-step-propagation}. By Lemma~\ref{sha-siegel-points-three-bounds}(1), $\dim R/\mathfrak a_s\le d+20C+1<N$ for $s\le8$, as there.
\item[(5)] Let $\mathfrak q'$ be a prime of $R$ that contains some $\mathfrak b_{v,i,j}$ and none of $\mathfrak a_1,\dots,\mathfrak a_8$. Then $\dim R/\mathfrak q'\le\dim B/\mathfrak q'_B+2+4|\Sigma_T|$, where $\mathfrak q'_B$ is the inverse image of $\mathfrak q'$ in $B$. Indeed, if $\mathfrak q'$ does not contain $\mathfrak a_9$, this follows from Lemmas~\ref{sha-siegel-step-irreducible}(2) and~\ref{sha-siegel-step-dimensions}(2), as in step (5) there. If $\mathfrak q'\supseteq\mathfrak a_9$, the image $r'$ of $f\circ r_X$ over the fraction field of $R/\mathfrak q'$ is supported in $W_9$. Each of the weights $(1,1)$, $(-1,-1)$, $(1,-1)$ and $(-1,1)$ of $W_9$ is the only one outside $W_1$, $W_2$, $W_3$ and $W_4$ respectively; as $\mathfrak q'$ does not contain $\mathfrak a_s$ for $s\le4$, some entry of $f\circ r_X$ of weight outside $W_s$ is nonzero modulo $\mathfrak q'$, and it has that weight. So $r'$ has entries of the four weights of $W_9$, and Lemma~\ref{sha-siegel-points-three-bounds}(2) gives the bound. As in step (5) there, Lemma~\ref{sha-siegel-step-dimensions}(3) then gives $\dim R/\mathfrak q'<N$.
\item[(6)] By (4), (5) and Lemma~\ref{sha-siegel-step-commutative}(3), with $I=\mathfrak P'+\mathfrak Q'$, the ideals $\mathfrak a_1,\dots,\mathfrak a_8$ and the $\mathfrak b_{v,i,j}$, there is a prime $\mathfrak p\supseteq\mathfrak P'+\mathfrak Q'$ with $\dim R/\mathfrak p=1$ and $\varpi\in\mathfrak p$ that contains none of $\mathfrak a_1,\dots,\mathfrak a_8$ and no $\mathfrak b_{v,i,j}$.
\item[(7)] The prime $\mathfrak p$ does not contain $\mathfrak a_9$. Suppose it does, and let $v\in\Sigma_T$. The image of $f(P_X)$ over the fraction field $K$ of $R/\mathfrak p$ is a flagged pair for the problem of $T$ whose homomorphism is supported in $W_9$. The entry $(2,1)$ of $f(n_{X,v})$ is a unit of $R$, since its residue is the entry $(2,1)$ of the flag of $c$ at $v$, which is nonzero as $c$ is of type A; so its image in $K$ is nonzero. By Lemma~\ref{sha-siegel-points-three-bounds}(3), $\psi_{v,1}=\psi_{v,2}$ or $\psi_{v,2}=\psi_{v,4}$ on $D_v(F)$ over $K$, where the $\psi_{v,i}(s)$ are the images of the $\psi^B_{v,i}(s)$. As $U_v\subseteq D_v(F)$, $\mathfrak p$ contains the $f(\varphi(\psi^B_{v,1}(s)-\psi^B_{v,2}(s)))$, or the $f(\varphi(\psi^B_{v,2}(s)-\psi^B_{v,4}(s)))$, for $s\in U_v$; it contains $f(\varphi(J))$ and $\varpi$, so it contains $\mathfrak b_{v,1,2}$ or $\mathfrak b_{v,2,4}$, a contradiction.
\item[(8)] So $\mathfrak p$ contains none of $\mathfrak a_1,\dots,\mathfrak a_9$, and steps (6) to (8) of the proof of Lemma~\ref{sha-siegel-step-propagation}, from the homomorphism $j:R\to\kappa[[\tau]]$ of Lemma~\ref{sha-siegel-step-commutative}(4) on, apply without change and give $f(\varphi(J))\subseteq\mathfrak Q'$.
\end{enumerate}
% ChainRankFour is chain_rank_four_of_red_a_y_coin_fib_y (ReductionTypeChainOfRedAYCoinFibY.lean): the ample case is Ws.residual_charpoly_rank_four (Lemma sha-siegel-step-propagation); the rest is cmplEx, x1, chDv, Ch.flat_min and Ch.chain as in steps (1) to (3) there, and the private ch_core4 for the induction step: (4) is the ideals A (redIdeal over redTypes minus yY) and B, with RedA; (5) is hb, with FibY through ent_of and the types t1 to t4, and with Ch.irr and fib otherwise, then ch_num with off_diagonal_nilpotent_rank_four; (6) is Ch.nice; (7) is hYp with YCoin; (8) is Ch.curve, hp1 (the patching property) and Ch.irr. The Lean behind the cited lemmas (cmplEx, x1, red, fib, off_diagonal_nilpotent_rank_four, Ch.irr, residual_charpoly_rank_four) takes T.Fit, Fr.Std or Fr.Ok, not Ctx2, and the book states those lemmas for fit towers.
\end{proof}

\begin{lemma}\label{sha-siegel-points-three-last}
Assume the Galois representations as in Theorem~\ref{sha-residually-reducible-automorphy}.
Let $D$ be a Siegel datum of case III, let $T$ be a fit deep tower for $D$ (Definition~\ref{sha-siegel-tower-towers}), let $\iota:\bar{\mathbf Q}_\ell\cong\mathbf C$, and let $V$ be an ambient ring with residue field $k_V$ (Definition~\ref{sha-siegel-tower-ambient}).
Let $C\ge0$ be such that $30(C+1)^2\le d_T$ and $T$ satisfies the cohomology bounds $(d_T,C)$ along every homomorphism from $\mathcal O$ to a field that kills $\mathfrak m_{\mathcal O}$ (Definition~\ref{sha-siegel-tower-ideal}).
Let $Q$ be an integral seed for $T$ over a ring $\mathcal O_1\supseteq\mathcal O$ of Definition~\ref{sha-character-parts}, and let $P$ be a flagged pair for the problem of $T$ along $\mathcal O\to V$, with entries in a complete discrete valuation subring of $V$ with finite residue field through which its homomorphism factors continuously, whose image $c$ over $k_V$ is residual and joined by an A-chain (Definition~\ref{sha-siegel-points-three-types}) to the image of $Q$ under $\mathcal O_1\to V\to k_V$.
Then there are a ring $\mathcal O'\supseteq\mathcal O$ of Definition~\ref{sha-character-parts}, a deformation space $X=(A,B,P_X)$ for $T$ over $\mathcal O'$ (Definition~\ref{sha-siegel-tower-space}), a modularity ideal $J$ for $X$ and a ring homomorphism $y:A\to V$ that restricts to the given map on $\mathcal O'$, maps $\mathfrak m_B$ into $\mathfrak m_V$, kills $J$, and carries $P_X$ to $P$.
\end{lemma}
\leanref{Development/AutomorphicArithmetic/AutomorphicVarieties/ModularityArgument/ThresholdArguments/DenseThresholdCases/PromodularityCaseThreeOfLemmas.lean}{last_step}
\begin{proof}
The proof of Lemma~\ref{sha-siegel-step-last} applies, with these changes. Lemmas~\ref{sha-siegel-step-ideal}, \ref{sha-siegel-step-linked} and~\ref{sha-siegel-step-propagation} apply to $T$, which is fit and deep. In (1) and (3), Lemma~\ref{sha-siegel-points-three-descent}(1), applied to the A-chain from $c$ to the image $t_V$ of $Q$, replaces Lemma~\ref{sha-siegel-step-coefficients}(2), and gives flagged pairs $s$ and $t$ along $\mathcal O\to k'$, with images $c$ and $t_V$, with $s$ residual and an A-chain from $s$ to $t$. In (4), $t$ is residual by Lemma~\ref{sha-siegel-points-three-descent}(2). In (7), Lemma~\ref{sha-siegel-points-three-propagation} gives the hypothesis of Lemma~\ref{sha-siegel-points-three-descent}(3), and that lemma, applied to the points $x_0$ and $x_1$ and the A-chain from $s$ to $t$, shows that $J$ lies in all minimal primes below $\ker x_0$. Steps (2), (5), (6) and (8) are unchanged.
% last_step of PromodularityCaseThreeOfLemmas.lean (ThresholdArguments): the same as the private last_step of Gsp4PromodularityCaseTwo.lean (Lemma sha-siegel-step-last) with hCD (ChainDesc4, Cu.chainDesc4) in place of chainDesc, chain_rank_four_good in place of Asm.chain_end, and allIn_of_chain4 with hCh (ChainRankFour) in place of Asm.allIn_of_chain with residual_charpoly_rank_four. PMExAlt (Lemma sha-siegel-step-ideal), smallCo, std and spaceEx take any fit tower.
\end{proof}

\begin{proposition}\label{sha-siegel-points}
Assume the Galois representations as in Theorem~\ref{sha-residually-reducible-automorphy}.
Let $\psi_0$ be a primitive Dirichlet character, let $h$ be a host datum for $\psi_0$ that is large for $t$, let $\ell$ be an odd prime, let $\mathcal O$ be as in Definition~\ref{sha-character-parts}, let $0<k<t$, and let $\psi\in X_h(k)$ be such that $(\ell,k,\psi)$ is in case II or case III (Definition~\ref{sha-host-datum}, with $\psi_0$ in place of $\psi$).
Let $D$ be a Siegel datum attached to $(h,k,\psi)$, with characters $\psi,\xi,\omega$ (Definition~\ref{sha-siegel-datum}).
\begin{enumerate}
\item (Finiteness.) $D$ is finite.
\item (No endoscopic points.) $D$ has no endoscopic $\bar{\mathbf Q}_\ell$-points.
\item (Density.) There is $\iota:\bar{\mathbf Q}_\ell\cong\mathbf C$ such that every irreducible $\bar{\mathbf Q}_\ell$-point of $D$ is automorphic through $\iota$.
\item (Rigidity.) Let $N\ge1$ be such that $\psi$, $\xi$ and $\omega$ are trivial on $\operatorname{Gal}(\bar{\mathbf Q}/\mathbf Q(\mu_N))$, and let $\iota:\bar{\mathbf Q}_\ell\cong\mathbf C$. Then $D$ is rigid at every irreducible $\bar{\mathbf Q}_\ell$-point that is automorphic through $\iota$.
\end{enumerate}
\end{proposition}
\begin{proof}
Use the notation of Definitions~\ref{sha-siegel-datum}, \ref{sha-siegel-tower-ambient} and~\ref{sha-siegel-tower-towers}, and fix $\iota:\bar{\mathbf Q}_\ell\cong\mathbf C$.
Below, $V$ is an ambient ring, with residue field $k_V$.
A \emph{$V$-point} of $D$ is the image of a Siegel point under a local homomorphism of $\mathcal O$-algebras $R\to V$, which factors continuously through a complete discrete valuation subring of $V$ with finite residue field.

\emph{Step 1: promodularity.}
Let $(r,u)$ be a flagged representation over $V$ with $r$ of type $\star$. Then $r$ is \emph{promodular}: there are a deep tower $T$ for $D$ satisfying (a) to (m) of Definition~\ref{sha-siegel-tower-towers}, with fields $K$, $F$ and $L=FK$; a ring $\mathcal O'\supseteq\mathcal O$ of Definition~\ref{sha-character-parts}; a deformation space $X=(A_F,B_F,P_X)$ for $T$ over $\mathcal O'$ (Definition~\ref{sha-siegel-tower-space}); a modularity ideal $J\subseteq B_F$ for $X$ (Definition~\ref{sha-siegel-tower-ideal}); a ring homomorphism $y:A_F\to V$; and a homomorphism $r_1:G\to M_4(V)$, such that: $y$ restricts to the given map on $\mathcal O'$, maps $\mathfrak m_{B_F}$ into $\mathfrak m_V$ and kills $J$; $r$ is conjugate to $r_1$ over $\Omega_V$; $r_1(g)$ is symplectic with multiplier $\theta(g)$ for all $g$; the restriction of $r_1$ to $G_F$ is $y\circ r_X$; and $y(\psi_{v,i}(s))=\delta_i(s)$ for $v\in\Sigma_T$ and $s\in I_v(F)$, where the $\psi_{v,i}$ are the diagonal entries of $P_X$.
By (k) and (f), $\bar r_0$ is trivial on the inertia groups of $F$ above $\ell$ and the $\bar\delta_i$ are distinct on $G_L$; and $G_L$ is normal in $G$, since $L$ is Galois over $\mathbf Q$.

\emph{Proof of Step 1 in case II.}
Suppose that $D$ is of case II (Definition~\ref{sha-siegel-tower-towers}).
\begin{enumerate}
\item[(1)] \emph{Relabelling.} Lemma~\ref{sha-siegel-step-labels} gives a relabelling $\epsilon$ satisfying (a), with the property stated there for the flagged representation $(r,u)$.
\item[(2)] \emph{A first tower.} By Lemma~\ref{sha-siegel-step-towers}(1) there is a fit tame tower $T_0$ with relabelling $\epsilon$. By (l), the $\bar\chi_i$ take pairwise distinct values at $g_0^{T_0}$.
\item[(3)] \emph{Lattice.} By (1), there is a near split rational flagged pair $P=(r_P,n)$ for $(\epsilon,g_0^{T_0})$ with $r_P$ conjugate to $r$ over $\Omega_V$; since $r$ is of type $\star$, $r_P$ has entries of weights $(0,1)$ and $(0,-1)$ and of one of the weights $(-1,1)$, $(-1,0)$ and $(-1,-1)$. Lemma~\ref{sha-siegel-step-lattice} gives a rational flagged pair $P'=(r_{P'},n')$ for $(\epsilon,g_0^{T_0})$ with $r_{P'}$ conjugate to $r_P$ over $\Omega_V$ and whose image $c_0$ over $k_V$ is of lattice shape.
\item[(4)] \emph{Seeds and chains.} Lemma~\ref{sha-siegel-step-chains}, applied to $T_0$ and $c_0$, gives a tame tower $T_1\ge T_0$ with the property stated there. By Lemma~\ref{sha-siegel-step-towers}(2), $T_1$ is fit.
\item[(5)] \emph{Bounds.} Lemma~\ref{sha-siegel-step-bounds}, which rests on Theorem~\ref{sha-ferrero-washington}, gives $C$ such that every tower $T'\ge T_1$ that is a layer of $T_1$ satisfies the cohomology bounds $(d_{T'},C)$ along every homomorphism from $\mathcal O$ to a field that kills $\mathfrak m_{\mathcal O}$.
\item[(6)] \emph{Depth.} By Lemma~\ref{sha-siegel-step-towers}(3) there is a deep tower $T\ge T_1$ that is a layer of $T_1$ with $d_T\ge30(C+1)^2$; by Lemma~\ref{sha-siegel-step-towers}(2), $T$ is fit and $T_0\le T$. By (4) there are $\mathcal O_1$ and an integral seed $Q$ for $T$ over $\mathcal O_1$ such that every residual flagged pair for $T$ along $\mathcal O\to k_V$ that is over $c_0$ is joined by a chain to the image of $Q$ in $k_V$.
\item[(7)] \emph{Restriction.} Lemma~\ref{sha-siegel-step-restriction}, applied to $T_0\le T$ and $P'$, gives a flagged pair $P''=(r'',(n''_v))$ for $T$ along $\mathcal O\to V$, with entries in a complete discrete valuation subring of $V$ with finite residue field through which $r''$ factors continuously, such that $r''$ is the restriction of $r_{P'}$ to $G_F$, the image of $P''$ over $k_V$ is residual and over $c_0$, and $\psi''_{v,i}(s)=\delta_i(s)$ for $s\in I_v(F)$. By (6), that image is joined by a chain to the image of $Q$ in $k_V$.
\item[(8)] \emph{Conclusion.} Lemma~\ref{sha-siegel-step-last}, applied to $T$, $C$, $Q$ and $P''$, gives $\mathcal O'$, a deformation space $X=(A_F,B_F,P_X)$ for $T$ over $\mathcal O'$, a modularity ideal $J$ for $X$ and $y:A_F\to V$ that restricts to the given map on $\mathcal O'$, maps $\mathfrak m_{B_F}$ into $\mathfrak m_V$, kills $J$ and carries $P_X$ to $P''$. Put $r_1=r_{P'}$. Then $r$ is conjugate to $r_P$, hence to $r_1$, over $\Omega_V$; $r_1(g)$ is symplectic with multiplier $\theta(g)$, as $P'$ is a flagged pair along $\mathcal O\to V$; the restriction of $r_1$ to $G_F$ is $r''=y\circ r_X$; and $y(\psi_{v,i}(s))=\psi''_{v,i}(s)=\delta_i(s)$ for $s\in I_v(F)$. This is the promodularity of $r$.
\end{enumerate}
\emph{Proof of Step 1 in case III.}
Suppose that $D$ is of case III.
\begin{enumerate}
\item[(1)] \emph{Relabelling.} As in case II, Lemma~\ref{sha-siegel-step-labels} gives a relabelling $\epsilon$ satisfying (a), with the property stated there for $(r,u)$.
\item[(2)] \emph{A first tower.} Lemma~\ref{sha-siegel-points-three-tower} gives a tame tower $T_0$ with relabelling $\epsilon$ that satisfies (a) to (m) and is of tensor type (Definition~\ref{sha-siegel-points-three-types}); it is fit. By (l), the $\bar\chi_i$ take pairwise distinct values at $g_0^{T_0}$.
\item[(3)] \emph{Lattice.} As in case II, (1) and Lemma~\ref{sha-siegel-step-lattice} give a near split rational flagged pair $P=(r_P,n)$ for $(\epsilon,g_0^{T_0})$ with $r_P$ conjugate to $r$ over $\Omega_V$, and a rational flagged pair $P'=(r_{P'},n')$ for $(\epsilon,g_0^{T_0})$ with $r_{P'}$ conjugate to $r_P$ over $\Omega_V$ and whose image $c_0$ over $k_V$ is of lattice shape.
\item[(4)] \emph{Seeds and A-chains.} Lemma~\ref{sha-siegel-points-three-chains}, applied to $T_0$ and $c_0$, gives a tame tower $T_1\ge T_0$ with the property stated there. By Lemma~\ref{sha-siegel-step-towers}(2), $T_1$ is fit.
\item[(5)] \emph{Bounds.} Lemma~\ref{sha-siegel-step-bounds}, applied to the fit tame tower $T_1$, gives $C$ such that every tower $T'\ge T_1$ that is a layer of $T_1$ satisfies the cohomology bounds $(d_{T'},C)$ along every homomorphism from $\mathcal O$ to a field that kills $\mathfrak m_{\mathcal O}$.
\item[(6)] \emph{Depth.} By Lemma~\ref{sha-siegel-step-towers}(3) there is a deep tower $T\ge T_1$ that is a layer of $T_1$ with $d_T\ge30(C+1)^2$; by Lemma~\ref{sha-siegel-step-towers}(2), $T$ is fit and $T_0\le T$. By (4) there are $\mathcal O_1$ and an integral seed $Q$ for $T$ over $\mathcal O_1$ such that every residual flagged pair for $T$ along $\mathcal O\to k_V$ that is over $c_0$ is joined by an A-chain to the image of $Q$ in $k_V$.
\item[(7)] \emph{Restriction.} As in case II, Lemma~\ref{sha-siegel-step-restriction}, applied to $T_0\le T$ and $P'$, gives a flagged pair $P''=(r'',(n''_v))$ for $T$ along $\mathcal O\to V$, with entries in a complete discrete valuation subring of $V$ with finite residue field through which $r''$ factors continuously, such that $r''$ is the restriction of $r_{P'}$ to $G_F$, the image of $P''$ over $k_V$ is residual and over $c_0$, and $\psi''_{v,i}(s)=\delta_i(s)$ for $s\in I_v(F)$. By (6), that image is joined by an A-chain to the image of $Q$ in $k_V$.
\item[(8)] \emph{Conclusion.} Lemma~\ref{sha-siegel-points-three-last}, applied to $T$, $C$, $Q$ and $P''$, gives $\mathcal O'$, a deformation space $X=(A_F,B_F,P_X)$ for $T$ over $\mathcal O'$, a modularity ideal $J$ for $X$ and $y:A_F\to V$ that restricts to the given map on $\mathcal O'$, maps $\mathfrak m_{B_F}$ into $\mathfrak m_V$, kills $J$ and carries $P_X$ to $P''$. With $r_1=r_{P'}$, the end of step (8) in case II shows that $r$ is promodular.
\end{enumerate}
% Case III of Step 1 is Cu.promodularity_case_three_of_chain_reduction (PromodularityCaseThreeOfLemmas.lean, ThresholdArguments, 24 lines): labels, ExistsTensorFactorTower for T_0 (existsTensorSeedTower_holds, Lemma sha-siegel-points-three-tower) with hH (Cu.hFit), RankFourChainReductionDense (chainIV_alt, Lemma sha-siegel-points-three-chains, with SeedIVAlt), lat, leFit, x2 (Lemma sha-siegel-step-bounds, applied to T_1 with hF₁ and tame), layerEx, restrQ, and its own last_step (Lemma sha-siegel-points-three-last) with ChainRankFour and ChainDesc4. The open part is the open step of Lemma sha-siegel-points-three-seeds.

\emph{Step 2: characteristic $\ell$.}
Let $\Omega_V$ have characteristic $\ell$, and let $(r,u)$ be a flagged representation over $V$ with $r$ promodular.
In characteristic $\ell$ the maps $\mathcal O'\to V$ factor through the residue fields, so $y$ kills $\mathfrak m_{\mathcal O}$ and the values minus $1$ of the inertial characters: $y$ carries these values to values of the $\delta_i$ on the inertia groups of $F$ above $\ell$, which are $1$ modulo $\mathfrak m_{\mathcal O}$ because $\bar r_0$ is trivial on these groups.
With $J$, this puts a power of the maximal ideal of $B_F$ in the kernel of $y$, which is prime, so $y$ maps $B_F$ through its residue field.
Hence for $g\in\operatorname{Gal}(\bar{\mathbf Q}/F)$ the characteristic polynomial of $r(g)$ is the image in $V$ of that of $\bar r_0(g)$.
Let $g\in G$. Some power of $g$ lies in $\operatorname{Gal}(\bar{\mathbf Q}/F)$, so the eigenvalues of $r(g)$ are roots of unity and $\operatorname{tr}r(g)-\sum_i\delta_i(g)$, computed in $V$, is integral over $\mathbf F_\ell$; it lies in $\mathfrak m_V$, since $r\equiv r_0$, so it is $0$.
Thus $\operatorname{tr}r(xgh)=\operatorname{tr}r(xhg)$ for all $x,g,h\in G$.
If $r$ were irreducible, the $r(x)$ would span $M_4(\Omega_V)$, so $r(g)r(h)=r(h)r(g)$ for all $g,h$, which is impossible for an irreducible representation of dimension $4$.
If $r$ were a non-split block extension, the same identity, used with an element of $\operatorname{Gal}(\bar{\mathbf Q}/F)$ at which the $\bar\delta_i$ are distinct, so that the characteristic polynomial of its image is separable, would make $r$ conjugate under some $1+Y$ to a block diagonal representation.
So $r$ is not of type $\star$.

\emph{Step 3: characteristic $0$.}
Let $V$ be the valuation ring of $\bar{\mathbf Q}_\ell$, and let $(r,u)$ be a flagged representation over $V$ with $r$ promodular.
On $I_\ell\cap\ker\psi\cap\ker\omega$, and on its conjugates, the characters $\delta_i$ are $\chi_{\mathrm{cyc}}^{m_i}$, with $m=(k+1,k,1,0)$, since $\xi$ is trivial on $I_\ell$; so the homomorphism $y$ of Step 1 carries $P_X$ to a flagged pair with inertial weights $2m$ (Definition~\ref{sha-siegel-tower-ideal}).
Hence, by the purity property of $J$, $r_1$, and with it $r$, which has the same characteristic polynomials, is pure through $\iota$.
If $r$ is irreducible, it stays irreducible on the normal subgroup $\operatorname{Gal}(\bar{\mathbf Q}/L)$: its reduction has trace $\sum_i\bar\delta_i$, and the $\bar\delta_i$ are distinct on $\operatorname{Gal}(\bar{\mathbf Q}/L)$.
Then $r_1$ is automorphic through $\iota$ by the automorphy property of $J$, and so is its conjugate $r$.

\emph{Step 4: points and flagged representations.}
\begin{enumerate}
\item[(a)] Every $V$-point $r$ has a flag $u$, that is, $(r,u)$ is a flagged representation.
The restriction of $r$ to $D_\ell$ comes from an ordinary lift $\rho$ over a domain $B$.
Since $\chi_{\mathrm{cyc}}(s_0)$ has infinite order, the $\delta_i(s_0)$ are distinct, so over the fraction field of $B$ the matrix $\rho(s_0)$ has distinct eigenvalues, and $\rho$ is upper triangular in an eigenbasis.
Hensel's lemma and the map to $V$ carry this basis to a triangularizing matrix over $V$, and an Iwasawa decomposition in the symplectic group makes it symplectic with multiplier $1$.
\item[(b)] Every $\bar{\mathbf Q}_\ell$-point is the image of a $V$-point for the valuation ring $V$ of $\bar{\mathbf Q}_\ell$: a homomorphism of $\mathcal O$-algebras from a complete local noetherian $\mathcal O$-algebra with residue field $\mathbf F$ to $\bar{\mathbf Q}_\ell$ takes values in one of the rings $\mathcal O'$ of Definition~\ref{sha-character-parts} and is local.
\item[(c)] A $V$-point $r$ is irreducible over $\Omega_V$; or it is of \emph{diagonal type}, meaning that for $i\ne j$ the entries $(i,j)$ and $(j,i)$ of the $r(g)$ are not both somewhere nonzero; or its entries $(1,3)$ and $(3,1)$ are somewhere nonzero and either all its entries $(i,j)$ with $\lambda_i<\lambda_j$ or all those with $\lambda_i>\lambda_j$ vanish.
\item[(d)] If $\Omega_V$ has characteristic $\ell$ and a $V$-point $r$ is of diagonal type, then every product of entries of $r$ of weight $0$ satisfies $x^{\ell^n}=x$ for some $n\ge1$.
\item[(e)] If a $V$-point $r$ with a flag is in the third case of (c), there is a flagged representation $(r_1,u_1)$ over $V$ with all $r_1(g)$ block diagonal and with first blocks spanning $M_2(\Omega_V)$.
\item[(f)] If $(r,u)$ is a flagged representation over $V$ with all $r(g)$ block diagonal and with first blocks spanning $M_2(\Omega_V)$, there is a flagged representation $(r_2,u_2)$ over $V$ with $r_2$ a non-split block extension and with the characteristic polynomials of $r$.
\item[(g)] If $D$ is not finite, there are $V$ with $\Omega_V$ of characteristic $\ell$, a $V$-point $r$ and a product $x$ of entries of $r$ of weight $0$ with $x^{\ell^n}\ne x$ for every $n\ge1$.
\item[(h)] If $V$ is the valuation ring of $\bar{\mathbf Q}_\ell$ and $(r,u)$ is a flagged representation over $V$ with all $r(g)$ block diagonal, then $r$ is not pure through $\iota$.
Indeed, on $I_\ell\cap\ker\psi\cap\ker\omega$ the flag shows that the two blocks of $r$ have determinants $\delta_1\delta_3=\chi_{\mathrm{cyc}}^{k+2}$ and $\delta_2\delta_4=\chi_{\mathrm{cyc}}^{k}$.
Each determinant is a character of $G$ that is trivial on the inertia groups on which $r$ is, so it equals $\chi_{\mathrm{cyc}}^{k+2}$, respectively $\chi_{\mathrm{cyc}}^{k}$, on an open subgroup $U$.
If $r$ were pure, with $H$ and $M$, there would be $\sigma\in H\cap U$ acting as $x\mapsto x^{p^f}$ at a place above a prime $p\nmid M\ell$, with $\chi_{\mathrm{cyc}}(\sigma)=p^f$.
The eigenvalues of the first block of $r(\sigma)$ have product $p^{f(k+2)}$, those of the second have product $p^{fk}$, and purity would give $p^{f(k+2)}<p^{2f}p^{fk}$.
\end{enumerate}

\emph{Proof of} (c).
Let $r$ be a $V$-point, with a flag $u$ by (a), and let $k_V$ be the residue field of $V$.
By Lemma~\ref{sha-flagged-idempotents}, the $V$-span of the $r(g)$ contains $E_{11},\dots,E_{44}$, so Lemma~\ref{sha-flagged-relation} applies to $r$ over $\Omega_V$; let $\succeq$ be its relation.
We first show:
\begin{itemize}
\item[($*$)] if $\lambda_i\ne\lambda_j$, $i\succeq j$ and $j\succeq i$, then $i\succeq k\succeq i$ for some $k\notin\{i,j\}$.
\end{itemize}
Suppose that ($*$) fails for some $i$ and $j$.
If $i\succeq a$ and $b\succeq i$, then $r(gh)_{ab}=r(g)_{ai}r(h)_{ib}+r(g)_{aj}r(h)_{jb}$ for all $g,h\in G$, since a nonzero term $r(g)_{ak}r(h)_{kb}$ with $k\notin\{i,j\}$ would give $i\succeq a\succeq k\succeq b\succeq i$.
Apply Lemma~\ref{sha-flagged-common-line} over $\Omega_V$ to the restriction of $r$ to $D_\ell$, with $u$ and $\Pi=\{k:k\succeq i\}$; its hypotheses hold by transitivity and since ($*$) fails.
It gives $(x,y)$; as $V$ is a valuation ring, either $y\ne0$ and $x/y\in V$, or $x\ne0$ and $y/x\in V$.
Put $(s,t)=(i,j)$ and $v=x/y$ in the first case, and $(s,t)=(j,i)$ and $v=y/x$ in the second. In both cases
\[
 r(h)_{ss}v+r(h)_{st}=\big(r(h)_{ts}v+r(h)_{tt}\big)v\quad(h\in D_\ell),\qquad r(gh)_{st}=r(g)_{ss}r(h)_{st}+r(g)_{st}r(h)_{tt}\quad(g,h\in G).
\]
Let $A\subset V$ be a complete discrete valuation ring through which $r$ factors continuously.
Since $s\succeq t$, there is $g_1$ for which $b=r(g_1)_{st}$ has the least valuation in $A$ among the nonzero $r(g)_{st}$; write $r(g)_{st}=b\,c(g)$ with $c(g)\in A$. The residue $\bar c(g)\in k_V$ is locally constant in $g$, by the continuity of $r$ modulo a power of $\mathfrak m_A$.

By Lemma~\ref{sha-flagged-idempotents}, $E_{ss}=\sum_ga_gr(g)$ for finitely many $a_g\in V$, so $\sum_ga_gr(gh)=E_{ss}r(h)$ for every $h$. Comparing the entries $(s,t)$, $(s,s)$ and $(t,t)$ gives $\sum_ga_gc(gh)=c(h)$, $\sum_ga_gr(gh)_{ss}=r(h)_{ss}$ and $\sum_ga_gr(gh)_{tt}=0$. Hence:
\begin{itemize}
\item[($**$)] there is no $m\in k_V$ with $\bar c(g)=m\big(\bar r(g)_{ss}-\bar r(g)_{tt}\big)$ for all $g\in G$,
\end{itemize}
where the bar denotes residues: otherwise $\bar c(h)=m\,\bar r(h)_{ss}$ for every $h$, and $h=1$, where $c(1)=0$, gives $m=0$, while $h=g_1$, where $c(g_1)=1$, gives $1=0$.

On the other hand, $\bar r(g)_{ss}$ and $\bar r(g)_{tt}$ are the images $\bar\delta_s(g)$ and $\bar\delta_t(g)$ in $k_V$, and $c(gh)=r(g)_{ss}c(h)+c(g)r(h)_{tt}$. So $z(g)=\bar c(g)/\bar\delta_t(g)$ is locally constant with $z(gh)=z(g)+\bar\gamma(g)z(h)$, where $\gamma=\delta_s\delta_t^{-1}$, which is one of $\alpha^{\pm1}$, $\beta^{\pm1}$, $\theta^{\pm1}$ since $\lambda_s\ne\lambda_t$. It satisfies the hypotheses of Lemma~\ref{sha-flagged-cocycles}~(1), with $a$ the map $\mathcal O\to k_V$:
\begin{itemize}
\item if $g$ lies in an inertia group at a place not above $T$ and $\bar r_0(g)=1$, then $r(g)=1$, so $c(g)=0$;
\item on $D_p$ with $p\in S_b$, $z$ is a coboundary by Lemma~\ref{sha-flagged-cocycles}~(2);
\item on $D_\ell$, write $\Delta(h)=r(h)_{tt}-r(h)_{ss}$. The first displayed identity gives $b\,c(h)=\Delta(h)v+r(h)_{ts}v^2$, and hence $c(h)\Delta(\sigma_\ell)-c(\sigma_\ell)\Delta(h)=v\big(c(\sigma_\ell)r(h)_{ts}-c(h)r(\sigma_\ell)_{ts}\big)$, which lies in $\mathfrak m_V$ since $r(h)_{ts}\equiv0$. As $\bar\delta_s(\sigma_\ell)\ne\bar\delta_t(\sigma_\ell)$ by Lemma~\ref{sha-flagged-residues}~(1), this gives $z(h)=\bar\gamma(h)m_\ell-m_\ell$ for $h\in D_\ell$, with $m_\ell=\bar c(\sigma_\ell)/\big(\bar\delta_s(\sigma_\ell)-\bar\delta_t(\sigma_\ell)\big)$.
\end{itemize}
So $z(g)=\bar\gamma(g)m-m$ for some $m\in k_V$, that is, $\bar c(g)=m\big(\bar r(g)_{ss}-\bar r(g)_{tt}\big)$ for all $g$, contrary to ($**$). This proves ($*$).

By Lemma~\ref{sha-flagged-relation}~(2) and ($*$), the relation $\succeq$ satisfies the hypotheses of Lemma~\ref{sha-flagged-relation-cases}.
In its first alternative, the span of the $r(g)$ is $M_4(\Omega_V)$ by Lemma~\ref{sha-flagged-relation}~(3), and $r$ is irreducible.
In its second alternative, $r$ is of diagonal type.
In its third, the entries $(1,3)$ and $(3,1)$ of $r$ are somewhere nonzero, and all entries $(i,j)$ with $\lambda_i<\lambda_j$ vanish, or all those with $\lambda_i>\lambda_j$ do.

\emph{Proof of} (d).
Let $\Omega_V$ have characteristic $\ell$, let $r$ be a $V$-point of diagonal type, with a flag $u$ by (a), and let $\succeq$ be the relation of Lemma~\ref{sha-flagged-relation}, which applies by Lemma~\ref{sha-flagged-idempotents}; as $r$ is of diagonal type, $\succeq$ is a partial order.
Let $a(i)$ be the number of $k$ with $i\succeq k$, and $h(i)=a(i)-a(5-i)$. If $i\ne j$ and $i\succeq j$, then $a(i)>a(j)$ and, as $5-j\succeq5-i$, $a(5-j)>a(5-i)$, so $h(i)>h(j)$.
Since $h(5-i)=-h(i)$, $\lambda_{5-i}=1-\lambda_i$ and $\mu_{5-i}=1-\mu_i$, we have $h(i)=(h(1)-h(2))\lambda_i+(h(1)-h(3))\mu_i-h(1)$. So the linear form $L(w)=(h(1)-h(2))w_1+(h(1)-h(3))w_2$ on $\mathbf Z^2$ takes the value $h(i)-h(j)>0$ at the weight of the entry $(i,j)$ whenever $i\ne j$ and $i\succeq j$.

Each diagonal entry $\kappa_i(g)=r(g)_{ii}$ is a character of $G$, since a nonzero term $r(g)_{ik}r(h)_{ki}$ with $k\ne i$ would give $i\succeq k\succeq i$.
Let $U=\ker\bar r_0\cap\ker\psi$, an open subgroup, and $N=[G:U]\,\ell^2$. Then $\kappa_i^N$ is trivial on every inertia group:
at a place not above $T$, $s^{[G:U]}\in\ker\bar r_0$ for $s$ in the inertia group, so $r(s^{[G:U]})=1$;
the inertia groups at the places above $p\in S_b$ are conjugate to $I_p$, and $r(s^{[G:U]})=1$ for $s\in I_p$;
for $s\in I_\ell\cap U$, the matrix $u^*r(s)u$ is upper triangular with diagonal entries the images of the $\delta_i(s)$, which are congruent to $1$ and hence equal to $1$, because the map $\mathcal O\to V$ kills $\mathfrak m_{\mathcal O}$ in characteristic $\ell$; so $(r(s)-1)^4=0$ and $r(s)^{\ell^2}=1$; the inertia groups at the other places above $\ell$ are conjugate to $I_\ell$, and $\kappa_i$ is a character.
Let $A\subset V$ be a complete discrete valuation ring through which $r$ factors continuously. Modulo each power of $\mathfrak m_A$, $\kappa_i^N$ is a character of $G$ with open kernel containing every inertia group, so it is trivial, since $\mathbf Q$ has no nontrivial extension unramified at every prime. Hence $\kappa_i^N=1$.

Now let $x$ be a nonzero product of entries of $r$ whose weights add up to $0$. The values of $L$ at the weights of its factors are $0$ for the diagonal factors and positive for the others, and they add up to $L(0)=0$, so all factors are diagonal entries, and $x$ is a root of unity. In characteristic $\ell$, $x^{\ell^am}=1$ implies $(x^m-1)^{\ell^a}=0$, so $x^m=1$ for some $m$ prime to $\ell$, and then $x^{\ell^n}=x$ with $n=\varphi(m)$. For $x=0$ every $n$ will do.

\emph{Proof of} (e).
Let $r$ be a $V$-point with a flag $u$ in the third case of (c), and let $\epsilon\in\{1,-1\}$ be such that the entries $(i,j)$ of $r$ with $\epsilon\lambda_i<\epsilon\lambda_j$ vanish. For a matrix $X$ let $X^\flat$ be its block diagonal part, obtained by replacing the entries $(i,j)$ with $\lambda_i\ne\lambda_j$ by $0$.
On the matrices whose entries $(i,j)$ with $\epsilon\lambda_i<\epsilon\lambda_j$ vanish, $X\mapsto X^\flat$ is multiplicative, and it commutes with $X\mapsto X^*$ since $\lambda_{5-i}=1-\lambda_i$. So $r_1(g)=r(g)^\flat$ is a homomorphism, symplectic with multiplier $\theta(g)$, congruent to $r_0$, factoring continuously through the same discrete valuation ring as $r$, and equal to $1$ wherever $r$ is.
The matrices $X$ with $u^*Xu$ upper triangular form a $V$-subalgebra of $M_4(V)$, since $u^*=u^{-1}$. It contains $r(\phi)$, which is block diagonal, with off-diagonal entries in $\mathfrak m_V$ and with residues of the diagonal entries distinct across the blocks, by Lemma~\ref{sha-flagged-residues}. By Lemma~\ref{sha-flagged-block-idempotent} it contains $E=E_{11}+E_{33}$, so $e=u^*Eu$ is an upper triangular idempotent.
Since $X^\flat=EXE+(1-E)X(1-E)$, for $s\in D_\ell$ and $y=u^*r(s)u$ we get $u^*r_1(s)u=eye+(1-e)y(1-e)$, which is upper triangular with diagonal entries $\big(e_{ii}^2+(1-e_{ii})^2\big)y_{ii}=y_{ii}$. So $(r_1,u)$ is a flagged representation.
Applying $X\mapsto X^\flat$ to expressions of the $E_{ii}$ as $V$-combinations of the $r(g)$, given by Lemma~\ref{sha-flagged-idempotents}, shows that the span of the $r_1(g)$ contains the $E_{ii}$. The entries $(1,3)$ and $(3,1)$ of $r_1$ are those of $r$, so by Lemma~\ref{sha-flagged-relation}~(4) the first blocks of the $r_1(g)$ span $M_2(\Omega_V)$.

\emph{Proof of} (f).
Let $(r,u)$ be as in (f), with entries in a complete discrete valuation ring $A\subset V$ with finite residue field through which $r$ factors continuously; $\ell$ lies in $\mathfrak m_A$, and $2$ is a unit.
\begin{enumerate}
\item[(1)] $r(c)=\operatorname{diag}(-1,1,-1,1)$. Indeed, $r(c)^2=1$, $r(c)$ is block diagonal, and $r(c)\equiv r_0(c)=\operatorname{diag}(-1,1,-1,1)$ by Lemma~\ref{sha-flagged-residues}~(3); its first block $X$ satisfies $(X-1)(X+1)=0$ with $X-1$ invertible, so $X=-1$, and in the same way the second block is $1$.
\item[(2)] Let $T$ be the $A$-module of the $Y\in M_4(A)$ with $Y^*=-Y$ whose entries $(i,j)$ vanish unless $\lambda_i<\lambda_j$; it is free of finite rank, and $YY'=0$ for $Y,Y'\in T$. Let $G$ act on $T$ by $\rho(g)Y=r(g)Yr(g)^{-1}$, which preserves $T$ since $r(g)$ is block diagonal and symplectic, and let $P$ be the set of $Y\in T$ with $u^*Yu$ strictly upper triangular.
The hypotheses of Lemma~\ref{sha-flagged-extension-cocycle} hold with $s=\sigma_\ell$, the set $S$ of $D$, and $N=\ker r$:
$\rho(c)=-1$ by (1), as conjugation by $r(c)$ multiplies the entries $(i,j)$ with $\lambda_i\ne\lambda_j$ by $-1$;
modulo $\mathfrak m_A$, $\rho(\sigma_\ell)-1$ multiplies the entry $(i,j)$ by $\bar\delta_i(\sigma_\ell)/\bar\delta_j(\sigma_\ell)-1\ne0$, by Lemma~\ref{sha-flagged-residues}~(1), so $\rho(\sigma_\ell)-1$ is bijective;
$\ker r$ contains the inertia groups at the places not above $S$, where $\psi$, $\xi$, $\omega$ and $\chi_{\mathrm{cyc}}$ are unramified, so that $\bar r_0$ is trivial there;
$P$ is saturated and stable under $D_\ell$, since the $u^*r(h)u$ with $h\in D_\ell$ are upper triangular;
and $P\ne0$: the first column of $u$ spans an $r(\sigma_\ell)$-stable line, on which $r(\sigma_\ell)$ acts by a scalar congruent to $\delta_1(\sigma_\ell)$, so by the separation of residues this column lies in the span of $e_1,e_3$, and the span of the first two columns of $u$ is spanned by it and a vector $w$ in the span of $e_2,e_4$; if $\varphi$ is the linear form vanishing on $e_2$, $e_4$ and on the first column of $u$, with $\varphi(e_1)$ the entry $(3,1)$ of $u$, a computation with $u^*u=1$ shows that $Y=w\varphi$ lies in $T$; it kills the first two columns of $u$ and has image in their span, so $Y\in P$, and $Y\ne0$.
\item[(3)] Lemma~\ref{sha-flagged-extension-cocycle} gives a cocycle $z:G\to T$, not identically $0$, with $z(c)=0$, and $y\in T$ with $z(h)-(\rho(h)y-y)\in P$ for $h\in D_\ell$.
\item[(4)] Let $\varpi$ be a uniformizer of $A$, and put $r_2(g)=(1+\varpi z(g))r(g)$ and $u_2=(1-\varpi y)u$. Since $YY'=0$ for $Y,Y'\in T$, $(1+\varpi z(g))r(g)(1+\varpi z(h))r(h)=(1+\varpi z(g)+\varpi\rho(g)z(h))r(gh)=r_2(gh)$, so $r_2$ is a homomorphism; and $(1+Y)^*(1+Y)=(1-Y)(1+Y)=1$ for $Y\in T$, so $r_2(g)$ is symplectic with multiplier $\theta(g)$ and $u_2$ with multiplier $1$. Also $r_2\equiv r_0$, $r_2$ factors continuously through $A$, and if $g$ lies in an inertia group at a place not above $\ell$ with $r(g)=1$, then $z(g)=0$ and $r_2(g)=1$.
For $h\in D_\ell$, $(1+\varpi y)(1+\varpi z(h))r(h)(1-\varpi y)=\big(1+\varpi(z(h)-(\rho(h)y-y))\big)r(h)$, so $u_2^*r_2(h)u_2=(1+N_h)\,u^*r(h)u$ with $N_h$ strictly upper triangular; it is upper triangular with the diagonal of $u^*r(h)u$. So $(r_2,u_2)$ is a flagged representation.
The matrix $\varpi z(g)r(g)$ is supported on the entries $(i,j)$ with $\lambda_i<\lambda_j$, so $r_2(g)$ is block triangular with the diagonal blocks of $r(g)$ and has the same characteristic polynomial.
\item[(5)] $r_2$ is a non-split block extension. Its entries $(i,j)$ with $\lambda_i>\lambda_j$ vanish. Its blocks in rows and columns $2$ and $4$ are those of $r$; for a block diagonal $X$ this block of $X^*$ is the adjugate of the first block of $X$, and $r(g)^*=\theta(g)r(g^{-1})$, so these blocks span the image under the adjugate, a bijection of $M_2(\Omega_V)$, of the span of the first blocks, which is $M_2(\Omega_V)$.
Suppose that $Y$ is supported on the entries $(i,j)$ with $\lambda_i<\lambda_j$ and that $(1-Y)r_2(g)(1+Y)$ is block diagonal for every $g$. The products $Yz(g)$, $Yr(g)Y$ and $z(g)r(g)Y$ vanish, since two matrices supported on these entries, with a block diagonal one between them, multiply to $0$; so $(1-Y)r_2(g)(1+Y)=r(g)+\big(\varpi z(g)-Y+r(g)Yr(g)^{-1}\big)r(g)$, and $\varpi z(g)r(g)=Yr(g)-r(g)Y$. At $g=c$ this gives $0=2Y$ by (1), so $Y=0$, and then $z=0$, a contradiction.
\end{enumerate}

\emph{Proof of} (g).
Suppose that $D$ is not finite, and let $U$, $\rho$, $\pi$, $z_s$, $w_s$ and $U_0$ be as in Lemma~\ref{sha-flagged-universal}. Lemma~\ref{sha-flagged-weight-zero} gives $\Omega$, $V$, $A'$, $Y:U\to A'$ and a monomial $v$ in the $z_s$ of weight $0$ with $Y(v)\ne0$. With the maps from the rings $\mathcal O'$ to $V$ through their residue fields, contained in $\bar{\mathbf F}_\ell\subset V$, the ring $V$ is as above with $\Omega_V=\Omega$ of characteristic $\ell$, and $r=Y\circ\rho$ is a $V$-point. Since $Y$ is local and $v\in\mathfrak m_U$, $Y(v)$ is a nonzero element of $\mathfrak m_V$.
Each $z_s$ is the entry $\rho(g_s)_{a_sb_s}$ if $a_s\ne b_s$, and $\rho(g_s)_{aa}-\delta_a(g_s)$, with $\rho(g_s)_{aa}$ of weight $0$, if $a_s=b_s=a$. Expanding the product, $v$ lies in the $\mathcal O$-algebra generated by the products of entries of $\rho$ whose weights add up to $0$.
Suppose that every product $x$ of entries of $r$ of weight $0$ satisfied $x^{\ell^n}=x$ for some $n\ge1$. Then so would every element of the $\mathcal O$-algebra they generate: the image of $\mathcal O$ lies in a finite field, $x\mapsto x^{\ell^n}$ is a ring homomorphism in characteristic $\ell$, and $x^{\ell^n}=x$ implies $x^{\ell^{nk}}=x$. In particular $Y(v)^{\ell^n}=Y(v)$ for some $n\ge1$, so $Y(v)\big(1-Y(v)^{\ell^n-1}\big)=0$, where $1-Y(v)^{\ell^n-1}$ is a unit since $Y(v)\in\mathfrak m_V$; so $Y(v)=0$, a contradiction. Hence some product $x$ of entries of $r$ of weight $0$ has $x^{\ell^n}\ne x$ for every $n\ge1$.

\emph{Proof of} (i).
Suppose that $D$ is not finite, and take $V$, $r$ and $x$ as in (g).
By (a), $r$ has a flag; consider the cases of (c).
If $r$ is irreducible, it is of type $\star$, hence promodular by Step 1, which contradicts Step 2.
If $r$ is of diagonal type, (d) contradicts the choice of $x$.
In the third case, (e) and (f) give a flagged representation with a non-split block extension, which is of type $\star$, and Steps 1 and 2 again give a contradiction.

\emph{Proof of} (ii).
Let $r$ be an endoscopic $\bar{\mathbf Q}_\ell$-point.
By (b) and (a), it is the image of a $V$-point with a flag, for the valuation ring $V$ of $\bar{\mathbf Q}_\ell$.
By (f), there is a flagged representation $(r_2,u_2)$ with $r_2$ a non-split block extension and with the characteristic polynomials of $r$.
By Steps 1 and 3, $r_2$ is pure through $\iota$, so $r$ is pure through $\iota$, which contradicts (h).

\emph{Proof of} (iii).
Let $r$ be an irreducible $\bar{\mathbf Q}_\ell$-point.
By (b) and (a), it is the image of a $V$-point with a flag, for the valuation ring $V$ of $\bar{\mathbf Q}_\ell$, and it is of type $\star$; by Steps 1 and 3 it is automorphic through $\iota$.

\emph{Proof of} (iv).
Let $r$ be an irreducible $\bar{\mathbf Q}_\ell$-point that is automorphic through $\iota$, and let $r'=r+\varepsilon r_1$ be a point with values in $\bar{\mathbf Q}_\ell[\varepsilon]$.
Let $\operatorname{ad}r$ be the action $g\cdot X=r(g)Xr(g)^{-1}$ of $G$ on $M_4(\bar{\mathbf Q}_\ell)$; it preserves $\mathfrak{sp}_4(\bar{\mathbf Q}_\ell)$ and the space of $X$ with $X^*=X$, since $r(g)^*=\theta(g)r(g)^{-1}$.
The function $z(g)=r_1(g)r(g)^{-1}$ is a continuous cocycle for $\operatorname{ad}r$, with values in $\mathfrak{sp}_4(\bar{\mathbf Q}_\ell)$ because $r'(g)$ is symplectic with multiplier $\theta(g)$.
It is a Bloch-Kato cocycle.
At a prime $p\ne\ell$, an inertia group at $p$ has an open subgroup on which $r'$ is trivial by the conditions on Siegel points, so $z$ vanishes there, and averaging over the finite quotient makes $z$ a coboundary on the inertia group.
At $\ell$, the ordinary condition puts the restriction of $z$ to $D_\ell$ in the finite part, by the Bloch-Kato local conditions.

There is an imaginary quadratic field $K$ such that the Bloch-Kato Selmer group of every polarized adjoint representation of the restriction of $r$ to $\operatorname{Gal}(\bar{\mathbf Q}/K)$ vanishes.
Here a \emph{polarized adjoint representation} is a representation $\rho$ of $G$ on $M_4(\bar{\mathbf Q}_\ell)$ that is $\operatorname{ad}r$ on $\operatorname{Gal}(\bar{\mathbf Q}/K)$, for which some complex conjugation $c'$ acts by $X\mapsto-P^{-1}\,{}^tX\,P$, where $P$ is invertible and ${}^tr(g)\,P\,r(c'gc'^{-1})=\theta(g)P$ for $g\in\operatorname{Gal}(\bar{\mathbf Q}/K)$.
The field $K$ comes from a tower. The datum $D$ is of case II or of case III (Definition~\ref{sha-siegel-tower-towers}), and the argument is the same in both cases.
Let $V$ be the valuation ring of $\bar{\mathbf Q}_\ell$. By (b) and (a), $r$ is the image of a $V$-point $r_V$ with a flag $u$; then $(r_V,u)$ is a flagged representation over $V$, and $r_V$ is irreducible, hence of type $\star$.
As in steps (1) to (3) of the proof of Step 1, Lemmas~\ref{sha-siegel-step-labels}, \ref{sha-siegel-step-towers}(1) and~\ref{sha-siegel-step-lattice} give a relabelling $\epsilon$, a fit tame tower $T_0$ with relabelling $\epsilon$, and a rational flagged pair $P'=(r_{P'},n')$ for $(\epsilon,g_0^{T_0})$ with $r_{P'}$ conjugate to $r$ and with image $c_0$ over $k_V$ of lattice shape; and Lemma~\ref{sha-siegel-step-towers}(3) and (2), applied to $T_0$, give a fit deep tower $T\ge T_0$, with fields $K$, $F$ and $L=FK$.
Lemma~\ref{sha-siegel-step-restriction} gives a flagged pair $P''$ for $T$ along $\mathcal O\to V$ whose homomorphism is the restriction of $r_{P'}$ to $G_F$, whose image over $k_V$ is residual, and with $\psi''_{v,i}(s)=\delta_i(s)$ for $s\in I_v(F)$.
For the ring $\mathcal O_1$ of step (F) below, Lemma~\ref{sha-siegel-step-coefficients}(3) gives a ring $\mathcal O'\supseteq\mathcal O_1$, a complete local noetherian $\mathcal O'$-algebra $R$ with residue field $k'$, an injective $f:R\to V$ and a family $P_R$ over $R$ that $f$ carries to $P''$, and Lemma~\ref{sha-siegel-step-coefficients}(1) a deformation space $X=(A,B,P_X)$ for $T$ over $\mathcal O'$.
The family $P_R$ is a flagged pair, since $f$ is injective, and it is a lift of its image $s$ in $k'$, which is residual because its image in $k_V$ is residual and $k'\to k_V$ is injective. By (5) of Definition~\ref{sha-siegel-tower-space} some $f_0:A\to R$ carries $P_X$ to $P_R$, and $y=f\circ f_0:A\to\bar{\mathbf Q}_\ell$ is a homomorphism of $\mathcal O'$-algebras that carries $P_X$ to $P''$.
Put $r_1=r_{P'}$. Then $r_1$ is conjugate to $r$, $r_1(g)$ is symplectic with multiplier $\theta(g)$, and the restriction of $r_1$ to $G_F$ is $y\circ r_X$.
The flagged pair $P''$ has inertial weights $2m$, with $m=(k+1,k,1,0)$ (Definition~\ref{sha-siegel-tower-ideal}): for $s$ in the open subgroup $G_F\cap\ker\psi\cap\ker\omega$ and in $I_v(F)$, $\psi''_{v,i}(s)=\delta_i(s)=\chi_{\mathrm{cyc}}(s)^{m_i}$, as in Step 3.
Finally, $r$ is irreducible on $G_L$. Every commutator of $G$ lies in $G_L$, since $F$ is abelian over $\mathbf Q$ and $K$ is quadratic, so $G_L$ is normal in $G$ with finite abelian quotient. If $r$ were reducible on $G_L$, then, as $r$ is irreducible, Lemma~\ref{sha-siegel-automorphy-clifford}, applied to $\Gamma=G$ and $H=G_L$, would give a character $\phi\ne1$ of $G$, trivial on $G_L$, and an invertible $M$ with $M\phi(g)r(g)=r(g)M$ for all $g\in G$. Let $g_0$ have $\phi(g_0)\ne1$. For $h\in G_L$, the matrix $\phi(g_0h)r(g_0h)=M^{-1}r(g_0h)M$ has the trace of $r(g_0h)$, and $\phi(g_0h)=\phi(g_0)\ne1$, so the trace of $r(g_0h)$ would vanish; reducing modulo $\mathfrak m_V$ would give $\sum_i\bar\delta_i(g_0)\bar\delta_i(h)=0$ for all $h\in G_L$. The $\bar\delta_i$ are pairwise distinct on $G_L$ by (f), so by the linear independence of characters every $\bar\delta_i(g_0)$ would vanish, which is absurd. So $r$, and with it its conjugate $r_1$, is irreducible on $G_L$.
So $r$ is irreducible on $G_K\supseteq G_L$.
With $m$ as above, write $(\nu_1,\nu_2,\nu_3,\nu_4)=(\omega^{-1}\xi\psi,\psi,\omega^{-1}\xi,1)$, so that $\delta_i=\chi_{\mathrm{cyc}}^{m_i}\nu_i$, with each $\nu_i$ trivial on $\operatorname{Gal}(\bar{\mathbf Q}/\mathbf Q(\mu_N))$; in particular $\theta=\chi_{\mathrm{cyc}}^{k+1}\nu_1$.
The following steps show that the field $K$ of $T$ has the property stated.
\begin{enumerate}
\item[(A)] \emph{Triangular forms above $\ell$.} For $s\in D_\ell$ the matrix $u^*r(s)u$ is upper triangular, with diagonal entries $\kappa_1(s),\dots,\kappa_4(s)$ congruent to the $\delta_i(s)$ modulo $\mathfrak m_V$ and equal to them for $s\in I_\ell$; so the $\kappa_i$ are characters of $D_\ell$, equal to $\chi_{\mathrm{cyc}}^{m_i}\nu_i$ on $I_\ell$. Moreover $\kappa_i\ne\kappa_j\chi_{\mathrm{cyc}}$ on $D_\ell$ for all $i,j$. For $i=j$ this holds at $s_0$, since $\chi_{\mathrm{cyc}}(s_0)\ne1$. Suppose that $\kappa_i=\kappa_j\chi_{\mathrm{cyc}}$ with $i\ne j$. At $s_0\in I_\ell\cap\ker\psi\cap\ker\omega$, where $\xi$ is trivial as well, $\delta_i(s_0)=\chi_{\mathrm{cyc}}(s_0)^{m_i}$; as $\chi_{\mathrm{cyc}}(s_0)$ has infinite order, $m_i=m_j+1$, which forces $\lambda_i\ne\lambda_j$ since $k\ge2$. At $\sigma_\ell$, where $\chi_{\mathrm{cyc}}$ is $1$, $\kappa_i(\sigma_\ell)=\kappa_j(\sigma_\ell)$, so $\bar\delta_i(\sigma_\ell)=\bar\delta_j(\sigma_\ell)$, contrary to Lemma~\ref{sha-flagged-residues}(1).
Let $w$ be a prime of $K$ above $\ell$ and $t\in G$ with $tv_\ell$ above $w$. By (d), $D_{tv_\ell}=tD_\ell t^{-1}$ lies in $G_K$, and $r$ is conjugate on it to an upper triangular representation with the characters $s\mapsto\kappa_i(t^{-1}st)$, which satisfy the same conditions, since $\chi_{\mathrm{cyc}}$ and the $\nu_i$ are characters of $G$.
\item[(B)] \emph{Automorphy over $K$.} Since $r$ is automorphic through $\iota$ and irreducible on $G_K$, its restriction to $G_K$ is attached through $\iota$ to a regular algebraic cuspidal automorphic representation $\pi$ of $GL_4(\mathbf A_K)$ and is windowed. Let $J_0$ be the matrix with rows $(0,0,0,1)$, $(0,0,1,0)$, $(0,-1,0,0)$ and $(-1,0,0,0)$, so that ${}^tX\,J_0=J_0X^*$ for every $X$. Since $r(g)^*r(g)=\theta(g)$, ${}^tr(g)J_0r(g)=\theta(g)J_0$ for all $g\in G$, so Lemma~\ref{sha-siegel-automorphy-polarized}, with $n=4$, $u=k+1$ and $M=N$, gives a Hecke character $\chi$ of $\mathbf Q$ with $(\pi,\chi)$ polarized.
\item[(C)] \emph{Free elements over $K$.} By (A) and (B), Lemma~\ref{sha-siegel-automorphy-free}, with $n=4$, $M=N$ and the weights $m$, applies to the restriction of $r$ to $G_K$: for every place $A'$ above $\ell$ some $g\in D_{A'}\cap G_K$ with $\chi_{\mathrm{cyc}}(g)=1$ has eigenvalues $\alpha_1,\dots,\alpha_4$ under $r$ with $\alpha_i^e\ne\alpha_j^e$ for $i\ne j$ and $e\ge1$.
\item[(D)] \emph{Automorphy over $L$.} The field $L=FK$ is Galois over $\mathbf Q$ with solvable Galois group, since $F$ and $K$ are abelian over $\mathbf Q$. The character $\theta=\chi_{\mathrm{cyc}}^{k+1}\nu_1$ is $\chi_{\mathrm{cyc}}^{k+1}$ times a character of finite order, so it is the character attached through $\iota$ to an algebraic Hecke character of $\mathbf Q$ \cite{Ser68}; it takes the same value at all complex conjugations, since they are conjugate in $G$. With $P=J_0r(c)^{-1}$, ${}^tr(s)Pr(csc^{-1})={}^tr(s)J_0r(s)r(c)^{-1}=\theta(s)P$ for all $s\in G$. As $r$ is irreducible on $G_L$, Theorem~\ref{sha-siegel-automorphy-base-change}, with $F_0=\mathbf Q$, $F=K$, $M=L$ and $\mu=\theta$, shows that the restriction of $r$ to $G_L$ is attached through $\iota$ to a regular algebraic cuspidal automorphic representation of $GL_4(\mathbf A_L)$.
\item[(E)] \emph{The windowed condition over $L$.} The restriction of $r$ to $G_L$ is windowed, by Lemma~\ref{sha-rigidity-windowed-up} with $F=K$, $E=L$, $n=4$, the restriction of $r$ to $G_K$, $\mu=\theta$, the complex conjugation $c$ and $P=J_0r(c)^{-1}$. Its hypotheses hold. The field $L$ is Galois over $K$, with Galois group a subgroup of the solvable group $\operatorname{Gal}(L/\mathbf Q)$ of (D). Hypothesis (1) is the irreducibility of $r$ on $G_L$ shown above. Hypothesis (2) is the pairing of (D), in which $\theta$ is attached through $\iota$ to an algebraic Hecke character of $\mathbf Q$, the maximal totally real subfield of $K$, and takes the same value at all complex conjugations. Hypothesis (3) is (B). For (4), by Lemma~\ref{sha-counts-vd-diag-local} there is $M_\ell\in GL_4(\bar{\mathbf Q}_\ell)$ such that $M_\ell r(s)M_\ell^{-1}$ is upper triangular for every $s\in D_\ell$. Every place $A$ of $\bar{\mathbf Q}$ above $\ell$ is $tv_\ell$ for some $t\in G$, so that $D_A(\mathbf Q)=tD_\ell t^{-1}$, and for $s\in D_A(\mathbf Q)$ the matrix $M_\ell r(t)^{-1}r(s)r(t)M_\ell^{-1}=M_\ell r(t^{-1}st)M_\ell^{-1}$ is upper triangular. So $r$ is conjugate to an upper triangular representation on $D_A(\mathbf Q)$, and on its subgroup $D_A(K)$. Finally, the restriction of $r$ to $G_K$ is windowed by (B).
\item[(F)] \emph{The vanishing over $L^+$.} There is a ring $\mathcal O_1$ of Definition~\ref{sha-character-parts}, the one used above, such that $T$ has a link over every such ring $\mathcal O'\supseteq\mathcal O_1$ (Definition~\ref{sha-siegel-automorphy-link}): this is Lemma~\ref{sha-siegel-automorphy-link-exists}, which applies since $D$ is of case II or of case III and the tower $T$ is fit and deep. Take the link over the ring $\mathcal O'$ of $X$, with its data of Situation~\ref{sha-hecke-situation} over $L$ (a finite set $S$ of places of $L^+$, the groups $\Gamma\supseteq\Gamma_0\supseteq\Gamma'$, the place $v_0$), its complex conjugation $\tau$, its projection $\mathrm{pr}:G_L\to\Gamma_0$, its bijection $v\mapsto\bar v$, its character $\kappa$ and its residual datum, which satisfies the conditions at $v_0$ of Theorem~\ref{sha-siegel-automorphy-adjoint}. Let $B^{\mathrm{ord}}$ be its universal ordinary determinant ring, and let $\rho_L(g)=\kappa(g)r(g)$ for $g\in G_L$.
\begin{itemize}
\item By Lemma~\ref{sha-siegel-automorphy-phi}, there is a homomorphism of $\mathcal O'$-algebras $\varphi:B^{\mathrm{ord}}\to B$ that carries the characteristic polynomial of $D^{\mathrm u}$ at $\mathrm{pr}(g)$, for $g\in G_L$, to the polynomial $c_g$ of Definition~\ref{sha-siegel-tower-space} with its roots multiplied by $\kappa(g)$, and $\psi^{\mathrm u}_{v,i}$ to $\kappa\psi^B_{\bar v,i}$ on $D_v$.
\item The point $x=y\circ\varphi$ of $B^{\mathrm{ord}}$ carries $D^{\mathrm u}$ to the determinant of $g\mapsto\kappa(g)r_1(g)$ on $G_L$, since $r_1|_{G_F}=y\circ r_X$, and so to the determinant of its conjugate $\rho_L$.
\item The representation $\rho_L$ factors through $\Gamma_0$. Its values span $M_4(\bar{\mathbf Q}_\ell)$, since $r$ is irreducible on $G_L$. Let $g\in\ker\mathrm{pr}$. As $x\circ D^{\mathrm u}$ is a determinant of $\Gamma_0$, $\operatorname{tr}\rho_L(hg)=\operatorname{tr}\rho_L(h)$ for all $h\in G_L$, so $\rho_L(g)-1$ is orthogonal to every $\rho_L(h)$ under the trace pairing, and $\rho_L(g)=1$. So $\rho_L=\rho\circ\mathrm{pr}$ for a representation $\rho$ of $\Gamma_0$, which is continuous since $\mathrm{pr}$ is a quotient map, and $x$ carries $D^{\mathrm u}$ to its determinant.
\item \emph{The inertial characters above $\ell$.} Let $v\mid\ell$, and let $a_v$ and $k_v$ be the integers of (e) of Definition~\ref{sha-siegel-automorphy-link} at the place of $\bar{\mathbf Q}$ at which $D_v$ is the decomposition group of $G_L$. Since $P''$ has inertial weights $2m$, there is an open subgroup $U$ of $G_F$ with $\psi''_{\bar v,i}(s)^2=\chi_{\mathrm{cyc}}(s)^{2m_i}$ for $s\in U\cap I_{\bar v}(F)$; $U$ contains an open normal subgroup of $G_F$, of some index $N$, so $g^N\in U$ for every $g\in G_F$. On $D_v$, $\psi_{x,v,i}=\kappa\psi''_{\bar v,i}$, by the property of $\varphi$ and since $y$ carries the $\psi^B_{\bar v,i}$ to the diagonal entries $\psi''_{\bar v,i}$ of $P''$. So $\psi_{x,v,i}(g)^{2N}=\kappa(g)^{2N}\chi_{\mathrm{cyc}}(g)^{2Nm_i}$ for $g\in I_v$, and since $(\kappa\chi_{\mathrm{cyc}}^{-a_v})^{k_v}=1$ on $I_v$,
\[
 \big(\psi_{x,v,i}(g)\,\chi_{\mathrm{cyc}}(g)^{-(m_i+a_v)}\big)^{2Nk_v}=1\qquad(g\in I_v).
\]
\item The point $x$ and $\rho$ satisfy hypotheses (1) to (6) of Theorem~\ref{sha-siegel-automorphy-adjoint}.

For (1), $x$ carries $D^{\mathrm u}$ to the determinant of $\rho$, as shown above. By local class field theory, if a continuous character $\psi$ of $D_v$ satisfies $(\psi\chi_{\mathrm{cyc}}^{-b})^{k'}=1$ on $I_v$ for some $k'\ge1$, then $\psi$ satisfies the condition of Definition~\ref{sha-hecke-arcs}(1) with the integers $-b$ at every embedding $\tau$. So, by the previous item, $x$ is arithmetic of the weight $h$ with $h_{v,i,\tau}=-(m_i+a_v)$, which is increasing because $m=(k+1,k,1,0)$ is strictly decreasing.

For (2), by (D) the restriction of $r$ to $G_L$ is attached through $\iota$ to a regular algebraic cuspidal automorphic representation $\pi$ of $GL_4(\mathbf A_L)$, and by (E) it is windowed. By (e) of Definition~\ref{sha-siegel-automorphy-link}, $\kappa$ is a power of $\chi_{\mathrm{cyc}}$ times a character of finite order on the inertia group at every place above $\ell$, so it is the character attached through $\iota$ to an algebraic Hecke character $\chi_\kappa$ of $L$ \cite{Ser68}. Hence $\rho_L$ is attached through $\iota$ to $\Pi=\pi\otimes(\chi_\kappa\circ\det)$, which is regular algebraic and cuspidal. With $P_L=J_0r(\tau)^{-1}$, (B) and (a) of Definition~\ref{sha-siegel-automorphy-link} give
\[
 {}^t\rho_L(g)\,P_L\,\rho_L(\tau g\tau^{-1})=\kappa(g)\kappa(\tau g\tau^{-1})\,{}^tr(g)J_0r(g)\,r(\tau)^{-1}=\chi_{\mathrm{cyc}}(g)^{-3}P_L\qquad(g\in G_L).
\]
So the eigenvalues of $\rho_L$ at a Frobenius element at $\tau w$ are $\chi_{\mathrm{cyc}}^{-3}$ times the inverses of those at $w$, and $\Pi^c$ and $\Pi^\vee$ have the same Satake parameters at almost all places; by strong multiplicity one (Theorem~\ref{lit-jacquet-shalika}(1)), $\Pi^\vee\cong\Pi^c$. Finally, $\rho_L$ is windowed. In ($\alpha$), the diagonal characters of $\rho_L$ above $\ell$ are those of $r$ multiplied by $\kappa$, which leaves the conditions unchanged. In ($\beta$), if $\Pi'$ is attached to $\rho_L$, then $\Pi'\otimes(\chi_\kappa^{-1}\circ\det)$ is attached to the restriction of $r$; as $L$ is a CM field, there is $b\in\mathbf R$ with $|\chi_{\kappa,u}(\varpi_u)|=q_u^{b}$ at every finite place $u$ of $L$, so the bound for $\Pi'\otimes(\chi_\kappa^{-1}\circ\det)$ with $s$ gives the bound for $\Pi'$ with $s+b$.

For (3), let $g\in G_L$ lie in the inertia group whose image in $\Gamma_0$ is $I_{v_0}$. The $\bar\chi_i$ are trivial on $I_{v_0}$, so by (3) of Definition~\ref{sha-siegel-automorphy-link} the residues of the $\kappa(g)\chi_i(g)$ are $1$. One of the $\chi_i$ is $\delta_4=1$, so $\kappa(g)\equiv1$ modulo $\mathfrak m_{\mathcal O'}$. By (c) of Definition~\ref{sha-siegel-automorphy-link}, $\kappa(g)^m=1$ for some $m$ prime to $\ell$; as the units of $\mathcal O'$ congruent to $1$ form a pro-$\ell$ group, $\kappa(g)=1$. Hence $\bar\delta_j(g)=1$ for all $j$, that is $\bar r_0(g)=1$. Since $v_0$ is not above $\ell$, $g$ lies in the set $N_\Phi$ of the problem of $T$, so $r_X(g)=1$, $r_1(g)=y(r_X(g))=1$, and $r(g)=1$, as $r$ is conjugate to $r_1$. So $\rho_L(g)=1$, and $\rho$ is trivial on $I_{v_0}$. Next let $v\mid\ell$ and $s\in D_v$, which lies in $D_{\bar v}(F)$. As $P''=(r_1|_{G_F},(n''_w))$ is a flagged pair and $n''_{\bar v}$ is symplectic with multiplier $1$, the matrix $(n''_{\bar v})^{-1}r_1(s)n''_{\bar v}$ is upper triangular with diagonal entries $\psi''_{\bar v,1}(s),\dots,\psi''_{\bar v,4}(s)$. Multiplying by $\kappa(s)$, and since $\psi_{x,v,i}=\kappa\psi''_{\bar v,i}$ on $D_v$, the restriction of $\kappa r_1$ to $D_v$ is conjugate to an upper triangular representation with diagonal characters $\psi_{x,v,1},\dots,\psi_{x,v,4}$; so is that of its conjugate $\rho_L$.

For (4), let $v\mid\ell$, and let $A'$ be the place of $\bar{\mathbf Q}$ at which $D_v$ is the decomposition group of $G_L$. By (C) some $g$ in the decomposition group of $G_K$ at $A'$ has $\chi_{\mathrm{cyc}}(g)=1$ and eigenvalues $\alpha_i$ under $r$ with $\alpha_i^e\ne\alpha_j^e$ for $i\ne j$ and $e\ge1$. A power $g^d$ with $d\ge1$ lies in $G_L$, which has finite index in $G_K$; then $g^d\in D_v$, $\varepsilon(g^d)=1$, and the eigenvalues of $\rho_L(g^d)=\kappa(g^d)r(g)^d$ are the $\kappa(g^d)\alpha_i^d$. By (3) they are the $\psi_{x,v,i}(g^d)$, and $(\kappa(g^d)\alpha_i^d)^e\ne(\kappa(g^d)\alpha_j^d)^e$ for $i\ne j$ and $e\ge1$.

For (5), two properties of the residual characters $\bar\chi_i$ of the link, with values in the residue field $k'$ of $\mathcal O'$, are used:
(i) every character $\eta:\Gamma_0\to k'^\times$ with $\eta^2=1$, trivial on the kernel of $\varepsilon\bmod\ell$, such that each $\eta\bar\chi_i$ is one of the $\bar\chi_j$, is trivial;
(ii) for some $i$ and $j$, $(\bar\chi_i/\bar\chi_j)^e\ne1$ for $1\le e\le4$.
For (i), such an $\eta$ takes the values $\pm1$, so $\eta\circ\mathrm{pr}$ is a character of $G_L$ with values in $\mathbf F^\times$ and square $1$. It is trivial on $G_L\cap\ker\omega$, since $\bar\omega=\bar\chi_{\mathrm{cyc}}$ and so $\varepsilon\equiv1$ modulo $\ell$ there. The residue of $\kappa$ cancels from the relations $\eta\bar\chi_i=\bar\chi_j$, which become $(\eta\circ\mathrm{pr})\bar\chi_i=\bar\chi_j$ on $G_L$ for the reductions of the relabelling characters. By (g) of Definition~\ref{sha-siegel-tower-towers}, $\eta\circ\mathrm{pr}=1$, and $\eta=1$ since $\mathrm{pr}$ is surjective. For (ii), by (h) there are a root $w'$ and $g\in G_L$ with $\bar\chi^{w'}(g)^e\ne1$ for $1\le e\le4$; $w'$ is the weight of an off-diagonal entry $(i,j)$, so $\chi_i\chi_j^{-1}=\chi^{w'}$, and $(\bar\chi_i/\bar\chi_j)(\mathrm{pr}(g))^e=\bar\chi^{w'}(g)^e\ne1$, the residue of $\kappa$ cancelling.
Let $A_x$ be the quotient of $B^{\mathrm{ord}}$ by the kernel of $x$, a local ring with residue field $k'$; let $j:A_x\to\bar{\mathbf Q}_\ell$ be the injective map induced by $x$, and $\phi:A_x\to k'$ the residue map. Since $x$ carries $D^{\mathrm u}$ to the determinant of $\rho$, and $D^{\mathrm u}$ reduces to the determinant of $\bar\chi_1\oplus\dots\oplus\bar\chi_4$, the characteristic polynomial of $\rho(y)$, for $y\in\Gamma_0$, is the image under $j$ of that of $D^{\mathrm u}$ at $y$ taken modulo $\ker x$, whose image under $\phi$ is $\prod_i(X-\bar\chi_i(y))$.
Fix $v\mid\ell$, and let $h=g^d\in D_v$ be the element found in (4). Its image $u_0$ in $\Gamma_0$ lies in $\Gamma'$, as $\varepsilon(h)=1$; by (3) the characteristic polynomial of $\rho(u_0)$ is $\prod_i(X-\psi_{x,v,i}(h))$; and $\psi_{x,v,i}(h)^e\ne\psi_{x,v,j}(h)^e$ for $i\ne j$ and $e\ge1$.
So Lemma~\ref{sha-siegel-automorphy-rank-four} applies, with the profinite group $\Gamma_0$, the character $\varepsilon$, the field $k'$, the continuous irreducible representation $\rho$, the ring $A_x$ with $j$ and $\phi$, the characters $\bar\chi_i$, the element $u_0$, and (i) and (ii); it gives (5).

For (6), let $m_{v,i}=m_i+a_v$, which is strictly decreasing in $i$, and $k=2N\prod_vk_v$. Dividing the relations of the item on the inertial characters for $i$ and for $j$ gives $\big(\psi_{x,v,i}\psi_{x,v,j}^{-1}\varepsilon^{m_{v,j}-m_{v,i}}\big)^k=1$ on $I_v$.
\item Theorem~\ref{sha-siegel-automorphy-adjoint} now shows: for every $P'$ and $R'$ as in Lemma~\ref{sha-siegel-automorphy-descent}, with $n=4$, $S'=S$ and $\tau$, every continuous cocycle of $\Gamma$ with values in $R'$ whose inflation to $G_{L^+}$ satisfies the Bloch-Kato conditions is a coboundary.
\end{itemize}
\item[(G)] \emph{Descent.} By (F), Lemma~\ref{sha-siegel-automorphy-descent}, with $n=4$, $\mu$ the restriction of $\theta$ to $G_K$, $S'=S$, the complex conjugation $\tau$, and $\rho$ and $\kappa$ of (F), shows that every Bloch-Kato cocycle of a polarized adjoint representation of the restriction of $r$ to $K$ is a coboundary: the relation $\kappa(g)\kappa(\tau g\tau^{-1})\theta(g)=\chi_{\mathrm{cyc}}(g)^{-3}$ of (F) is the one required there.
\end{enumerate}
% In (F), the link is Ws.nonempty_link_of_fit_deep_unramified (Lemma sha-siegel-automorphy-link-exists), from UnramifiedInputs, Fit (through Tower.Ok and Ctx2 or Ctx3, which the lemma states as (a) to (m) and the hypothesis that D is of case II or III) and Deep; the universal ring is from UnramifiedInputs.b0. The step uses only the twist Link.twc 0 = kappa, not beta and t.
% (E) is Z.Win.Dn.automorphy_up_base_change (Lemma sha-rigidity-windowed-up), applied in exists_cm_field_twist_h1 with F = T.K, K = T.L, the solvable group hsol' (injecting into Gal(L/Q) by AlgEquiv.restrictScalarsHom), the irreducibility hiL, Af.Sol.Pol with F0 = Q, mu = theta (isAlgChar_pushCh_th_of_good, auSign), c = D.cc and P from Ws.auPair, the automorphy over K from autKQ0_alt, Flagged from Z.Win.flagLaws (flagged_level from the level of GradeZeroInterpolatingElement.level, whose triangular form at D.pl l is Loc.spPtL, Lemma sha-counts-vd-diag-local, conjugated by r(gamma) with flag_pull and conj_eq, then flagged_res to K), and WitA over K from autQN'_alt.
% (5) is Th.nice0four_alt, through the private absolute_irreducible_rank_four of TNiceFour.lean, which is Lemma sha-siegel-automorphy-rank-four; A_x, j and phi are Th.Det.redB, the element u_0 is hF w with charpoly_eq_prod_flag_characters and hA.flag.
% Written in (F), all in adjoint_vanishing_gsp4_rigid_point: the determinant (Ws.auRep); the factorization (phiKer and the private factor_abs); the inertial characters (Pt.twWts from Link.cyc with n = 0, so e = 2a, and Tower.Wts with pow_mem_open, as in Ws.Pt.ptWts and Ws.auArith through auPsi); (1) is Ws.auArith with Loc.laArith (local class field theory, with GlobalArtinReciprocity); (2) is GalN.V.point_automorphy_w_of_gsp_cusp_sub_rank_four applied with the trivial twist, a = 0 and k = 1 (automorphy_bridge_of_printed_hida; Af.algOfCyc_s through Pt.twWts and half_pow; Af.twist_invariance_automorphic; Ws.Pt.ptPair with P = formJ r(c_F)^-1 and Link.conj; GalN.V.csdOf0_alt through self_dual_complex_of_printed_hida; WitA.witness_twist_algebraic_character, which is aIsoTwOf_alt with flagLaws for (alpha), Af.Cq.attTwist and winTw with charAbs for (beta)); (3) is Ws.Pt.ptUnr (Link.tame, Link.v0, Link.chib and the label of the character 1, then X.pair.kill) and Ws.Pt.ptFlag (X.pair.flag and X.pair.low, with IsPhi.psi_val); (4) is the private free_up with free_pow; (i) and (ii) of (5) are Ws.P1.condC from (g) of Fit (g3', with om_bar) and Ws.P1.condD from (h) (Fit.ord, condD_roots, condD_rt); (6) is Ws.Pt.ptWts with the private par'; the vanishing is the application of Th.AdjointLiftingTheorem.

Let $\eta$ be the quadratic character of $K$, and let $\rho$ act by $\operatorname{ad}r$ on $\mathfrak{sp}_4(\bar{\mathbf Q}_\ell)$ and by $\eta\cdot\operatorname{ad}r$ on the matrices $X$ with $X^*=X$.
Let $J_0$ be as in (B), so that ${}^tX\,J_0=J_0X^*$ for every $X$.
A computation from $r(g)^*r(g)=\theta(g)$ shows that $\rho$ is a polarized adjoint representation, with $c'=c$ and $P=J_0\,r(c)^{-1}$.
Since $z$ takes values in $\mathfrak{sp}_4(\bar{\mathbf Q}_\ell)$, on which $\rho$ is $\operatorname{ad}r$, it is a Bloch-Kato cocycle for $\rho$, so $z(g)=\rho(g)X_0-X_0$ for some $X_0$.
The projection $X\mapsto\frac12(X-X^*)$ onto $\mathfrak{sp}_4(\bar{\mathbf Q}_\ell)$ commutes with $\rho$ and fixes $z(g)$, so $z(g)=r(g)X_1r(g)^{-1}-X_1$ with $X_1=\frac12(X_0-X_0^*)$.
Hence $r_1(g)=Xr(g)-r(g)X$ with $X=-X_1\in\mathfrak{sp}_4(\bar{\mathbf Q}_\ell)$.
% Step 1 in case II is Ws.proMod_of_pmEx_of_ctx2_of_kPair_of_star with the private last_step of Gsp4PromodularityCaseTwo.lean; the propagation in last_step is Asm.allIn_of_chain with Ws.residual_charpoly_rank_four (cmplEx, x1, chDv, Ch.flat_min, Ch.chain, and the private ch_core with Ch.nice, Ch.curve, Ch.irr, red, fib and off_diagonal_nilpotent_rank_four).
% The open parts of the lemmas sha-siegel-step-* have their own open steps, with their Lean names and sizes. Case III of Step 1 is written out after case II, and its open part is the open step of Lemma sha-siegel-points-three-seeds; in the tower for r, case III uses exists_tame_tower_ok_three (Lemma sha-siegel-step-towers (1)).
% The field K in (iv) is Z.Win.Dn.exists_cm_field_twist_h1. The tower for r, written above for cases II and III, is Th.spPt (labels, exists_tame_tower_ok_two, lat, layerEx, restrQ, the private point with smallCo and spaceEx, ps_wt_any, and irr_L through irr_sub, comm_L and dist_L; irr_sub uses U3.exists_selfTwist for the Clifford step, Lemma sha-siegel-automorphy-clifford). Steps (A) to (G) follow exists_cm_field_twist_h1: (A) is Th.spUnl (with Loc.spPtL) and Th.unlRes; (B) is GalN.V.autKQ0_alt (Z.Bc.polarization_k_of_printed_results, Lemma sha-siegel-automorphy-polarized) with GradeZeroInterpolatingElement.level; (C) is Z.Win.RR.freePUA_alt; (D) is Ws.auGal, Ws.isAlgChar_pushCh_th_of_good, Ws.auSign, Ws.auPair and automorphic_of_gsp4_rigid_solvable_galois (Z.Bc.solQ); (E) is Z.Win.Dn.automorphy_up_base_change; (F) is Ws.nonempty_link_of_fit_deep_unramified and adjoint_vanishing_gsp4_rigid_point, through Th.AdjointLiftingTheorem (Theorem sha-siegel-automorphy-adjoint); (G) is Th.adDown with the field conj of the link (Lemma sha-siegel-automorphy-descent). Case III of the tower for r: spPt uses exists_tame_tower_ok_three (Ctx3, giving Tower.Ok and Tower.Three) in place of exists_tame_tower_ok_two, and is otherwise the same; Fit in case III is Ok with Ctx3 and Tower.Three0 (l >= 5, S_b empty, psi of order prime to l), and layerEx, leFit, restrQ, the private point, the link and steps (A) to (G) take any fit tower; the text states Definition sha-siegel-tower-towers and Lemmas sha-siegel-step-towers, sha-siegel-step-restriction and sha-siegel-step-coefficients (1) for fit towers in both cases.
% Step 1 is ProModAlt (with Tower.PotentialModularityTower: fin, p1, seed, pur, aut, and Tower.Ok: ineL, dist), proved in case II by Ws.proMod_of_pmEx_of_ctx2_of_kPair_of_star (seeds Ws.SeedIIAlt from Z.Bc.seedIIs_alt, Ws.PMExAlt from pmEx_of_printsH, Published.FerreroWashington) and in case III by Cu.PromodularityCaseThreeDense, from promodularity_case_three_of_chain_reduction (existsTensorSeedTower_holds, chain_rank_four_of_red_a_y_coin_fib_y, chainDesc4, chainIV_alt, Cu.SeedIVAlt from fourth_initial_of_prints_h); the assembly is Cu.modularity_rank_four_of_promodularity_case_three. A flagged representation is Ws.KPair; V is Ws.Amb; a V-point is Ws.VPt; type star is Ws.Star (Irr or BlockExt); pure is GalN.V.PureW.
% Step 2 is Ws.contradiction_star_promodular (con_const, con_trace, con_span, con_block); Step 3 is Ws.pureW_of_proModular_of_isZero and Ws.aut_n_alt (irr_L); (i) to (iii) are Ws.fin0_no_end_all_irr_of_good_con_purity_aut.
% Step 4: (a) Ws.ptPair (Pp.genTri, Pp.triSpec, Pp.triIwa); (b) Ws.eptV (coPt); (c) Ws.triV; (d) Ws.diagConst; (e) Ws.blockPart; (f) Ws.bl (Bl.cocA, from CohomologyHyp); (g) Ws.ptOver (Po.univGr, Po.ambPt, Po.ptA); (h) Ws.nePure_w (blockDet, wtChar, frobIn, pureDet_w).
% (iv) is Th.rigSp4_alt: tanBK for the cocycle (tb_place, tb_avg, tb_atL with Loc.bkAbove), rg_sign, rg_alg, rg_pair, rg_bk. The vanishing is Z.Win.Dn.exists_cm_field_twist_h1, through spPt, nonempty_link_of_fit_deep_unramified, autKQ0_alt, freePUA_alt, automorphic_of_gsp4_rigid_solvable_galois, automorphy_up_base_change, adjoint_vanishing_gsp4_rigid_point (from Th.AdjointLiftingTheorem) and adDown.
\end{proof}
\leanref{Development/AutomorphicArithmetic/AutomorphicVarieties/ModularityArgument/ThresholdArguments/DenseThresholdCases/IsAut4.lean}{UpperTheory.V.Aut.SymplecticFourImageDense}
\leanref{Development/AutomorphicArithmetic/AutomorphicVarieties/ModularityArgument/ThresholdArguments/DenseThresholdCases/Sp4In0.lean}{UpperTheory.V.HV.DescentVanishingReductions.sp4In0_alt}
\leanref{Development/AutomorphicArithmetic/AutomorphicVarieties/ModularityArgument/ThresholdArguments/DenseThresholdCases/RigEnd3.lean}{UpperTheory.V.Th.RigiditySymplecticFour}
\leanref{Development/AutomorphicArithmetic/AutomorphicVarieties/ModularityArgument/ThresholdArguments/DenseThresholdCases/RigSp4.lean}{UpperTheory.V.Th.rigSp4_alt}
\leanref{Development/AutomorphicArithmetic/AutomorphicVarieties/ModularityArgument/ThresholdArguments/DenseThresholdCases/VWinPm.lean}{UpperTheory.V.Ws.fin0_no_end_all_irr_of_good_con_purity_aut}

\subsection{Automorphic eigensystems and Satake parameters}\label{subsection-sha-eigensystems}

\begin{definition}\label{sha-eigensystem-level}
Let $F$ be a number field with algebraic closure $\bar F$, let $\ell$ be a prime, let $n,M\ge1$ and let $m=(m_1,\dots,m_n)\in\mathbf Z^n$.
Let $r:\operatorname{Gal}(\bar F/F)\to GL_n(\bar{\mathbf Q}_\ell)$ be a continuous representation.
\begin{enumerate}
\item The representation $r$ has \emph{level $M$ and weights $m$} if the following hold. At every prime $w$ of $F$ not above $\ell$, the restriction of $r$ to an inertia group at $w$ is conjugate to the diagonal representation with entries $\chi_1,\dots,\chi_n$, for characters $\chi_i$ of $\operatorname{Gal}(\bar F/F)$ that are trivial on $\operatorname{Gal}(\bar F/F(\mu_M))$. At every prime $w$ of $F$ above $\ell$ there are characters $\chi_1,\dots,\chi_n$ of $\operatorname{Gal}(\bar F/F)$, trivial on $\operatorname{Gal}(\bar F/F(\mu_M))$, and characters $\kappa_1,\dots,\kappa_n$ of a decomposition group $D_w$ at $w$, such that the restriction of $r$ to $D_w$ is conjugate to an upper triangular representation with diagonal entries $\kappa_1,\dots,\kappa_n$, the character $\kappa_i$ equals $\chi_{\mathrm{cyc}}^{m_i}\chi_i$ on the inertia group, and either $\kappa_i\ne\kappa_j\chi_{\mathrm{cyc}}$ for all $i,j$ or $r$ is potentially crystalline at $w$.
\item The representation $r$ is \emph{$\ell$-distinct} if at every prime $w$ of $F$ above $\ell$ its restriction to a decomposition group $D_w$ is conjugate to an upper triangular representation whose diagonal characters $\kappa_1,\dots,\kappa_n$ satisfy $\kappa_i\ne\kappa_j\chi_{\mathrm{cyc}}$ on $D_w$ for all $i,j$.
\end{enumerate}
\end{definition}
\leanref{Development/AutomorphicArithmetic/AutomorphicVarieties/AutomorphicForms/BlochKatoSelmerVanishes.lean}{UpperTheory.V.Af.Level}
\leanref{Development/AutomorphicArithmetic/AutomorphicVarieties/AutomorphicForms/BlochKatoSelmerVanishes.lean}{UpperTheory.V.Af.LevelAway}
\leanref{Development/AutomorphicArithmetic/AutomorphicVarieties/AutomorphicForms/BlochKatoSelmerVanishes.lean}{UpperTheory.V.Af.LevelAt}
\leanref{Development/AutomorphicArithmetic/AutomorphicVarieties/ModularityArgument/UpperTriangularPlaces/OddAOfUnl.lean}{UpperTheory.V.Keep.OddA}

\begin{definition}\label{sha-eigensystem-triple}
Let $K$ be an imaginary quadratic field of discriminant $d_K$, let $\ell$ be an odd prime, let $\iota:\bar{\mathbf Q}_\ell\cong\mathbf C$, and let $k,N\ge1$.
An \emph{automorphic triple of weight $k$ and level $N$} is a continuous representation $\rho:\operatorname{Gal}(\bar{\mathbf Q}/K)\to GL_3(\bar{\mathbf Q}_\ell)$ of one of the following two kinds.
\begin{enumerate}
\item \emph{Stable.} The representation $\rho$ is irreducible; it has level $N$ and weights $(k,0,-k)$ (Definition~\ref{sha-eigensystem-level}), with characters $\kappa_i$ at each prime above $\ell$ that satisfy $\kappa_i\ne\kappa_j\chi_{\mathrm{cyc}}$ for all $i,j$; and $\rho\otimes\chi_{\mathrm{cyc}}^{-1}$ is attached through $\iota$ to a regular algebraic cuspidal automorphic representation $\pi$ of $GL_3(\mathbf A_K)$ with $\pi^\vee\cong\pi^c$, and is windowed in the sense of Theorem~\ref{sha-residually-reducible-automorphy}.
\item \emph{Endoscopic.} There is an irreducible continuous representation $W:\operatorname{Gal}(\bar{\mathbf Q}/\mathbf Q)\to GL_2(\bar{\mathbf Q}_\ell)$ of level $N|d_K|$ and weights $(k,-k)$, with $\det W$ of finite order, such that $W\otimes\chi_{\mathrm{cyc}}^{-k}$ is attached through $\iota$ to a newform of weight $2k+1$, the restriction of $W$ to every imaginary quadratic field on which it is irreducible is windowed, and $\rho$ is conjugate to the representation
\[
 g\mapsto\begin{pmatrix}a(g)&0&b(g)\\0&1&0\\c(g)&0&d(g)\end{pmatrix},\qquad\text{where }W(g)=\begin{pmatrix}a(g)&b(g)\\c(g)&d(g)\end{pmatrix}.
\]
\end{enumerate}
\end{definition}
\leanref{Development/AutomorphicArithmetic/AutomorphicVarieties/ModularityArgument/ThresholdArguments/DenseThresholdCases/NumberU3US.lean}{UpperTheory.V.Z.J.Adm3U_s}
\leanref{Development/AutomorphicArithmetic/AutomorphicVarieties/ModularityArgument/ThresholdArguments/DenseThresholdCases/AfStmt.lean}{UpperTheory.V.Af.Adm3Alt}
\leanref{Development/AutomorphicArithmetic/AutomorphicVarieties/ModularityArgument/AutomorphyLifting/AutomorphyDescent/AutNewformW.lean}{UpperTheory.V.Z.Win.Dn.Adm2A}

\begin{definition}\label{sha-eigensystem-data}
Let $q,n\ge1$, let $m\in\mathbf Z^n$ be strictly decreasing and let $u\in\mathbf Z$.
\begin{enumerate}
\item An \emph{eigenvalue datum} of rank $n$ at $q$, with weights $m$ and weight $u$, is a finite set $S$ together with a multiset $\beta_s$ of $n$ complex numbers for each $s\in S$, such that:
\begin{itemize}
\item for every $\sigma\in\operatorname{Aut}(\mathbf C)$ there is a bijection $\tau$ of $S$ with $\beta_{\tau(s)}=\sigma(\beta_s)$ for all $s$;
\item for each $s$, every $b\in\beta_s$ satisfies $q^{(u-1)/2}<|b|<q^{(u+1)/2}$, and the product of the elements of $\beta_s$ has absolute value $q^{nu/2}$;
\item for each $s$, every $b\in\beta_s$ and every $a\in\mathbf Z$, the number $q^ab$ is an algebraic integer if $a\ge-m_i$ for all $i$, and $q^ab^{-1}$ is an algebraic integer if $a\ge m_i$ for all $i$.
\end{itemize}
\item Let $k\ge1$. A \emph{trace datum} of weight $k$ at $q$ is a finite set $S$ together with a complex number $t_s$ for each $s\in S$, such that for every $\sigma\in\operatorname{Aut}(\mathbf C)$ there is a bijection $\tau$ of $S$ with $t_{\tau(s)}=\sigma(t_s)$, and such that $|t_s|\le\sqrt q+2$ and $q^kt_s$ is an algebraic integer for every $s$.
\item Let $\ell$ be a prime, $\iota:\bar{\mathbf Q}_\ell\cong\mathbf C$, $\Gamma$ a group, and $\mathcal X$ a set of representations of $\Gamma$ in $GL_n(\bar{\mathbf Q}_\ell)$. A \emph{classifying map} for $\mathcal X$ into a finite set $S$ is a map $s:\mathcal X\to S$ such that $s(r)=s(r')$ implies $r\cong r'$. Let $F$ be a number field, $\Gamma=\operatorname{Gal}(\bar F/F)$ and $w$ a prime of $F$. An eigenvalue datum \emph{classifies $\mathcal X$ at $w$} through $\iota$ if there is a classifying map $s$ such that every $r\in\mathcal X$ is unramified at $w$ and, for an arithmetic Frobenius element $\sigma_w$ at $w$, the eigenvalues of $\iota r(\sigma_w)$ are the elements of $\beta_{s(r)}$. For $n=3$ and $g\in\Gamma$, a trace datum \emph{classifies $\mathcal X$ at $g$} through $\iota$ if there is a classifying map $s$ with $\iota\operatorname{tr}r(g)=t_{s(r)}$ for every $r\in\mathcal X$.
\end{enumerate}
\end{definition}
\leanref{Development/AutomorphicArithmetic/Foundations/GaloisRepresentations/AdmissibleResidues/WinDefs.lean}{UpperTheory.GalN.V.Win.EigData}
\leanref{Development/AutomorphicArithmetic/Foundations/GaloisRepresentations/AdmissibleResidues/WinDefs.lean}{UpperTheory.GalN.V.Win.EigData.Classifies}
\leanref{Development/AutomorphicArithmetic/Foundations/GaloisRepresentations/AdmissibleResidues/WinDefs.lean}{UpperTheory.GalN.V.Win.TrData}
\leanref{Development/AutomorphicArithmetic/Foundations/GaloisRepresentations/AdmissibleResidues/WinDefs.lean}{UpperTheory.GalN.V.Win.TrData.Classifies}
\leanref{Development/AutomorphicArithmetic/Foundations/GaloisRepresentations/AdmissibleResidues/SumThree.lean}{UpperTheory.GalN.V.Win}

\begin{lemma}\label{sha-eigensystem-cuspidal-count}
Let $K$ be a CM field and let $n\ge2$.
There are $c,A$ such that the following holds.
Let $m\in\mathbf Z^n$ be strictly decreasing, let $H\ge0$ with $|m_i|\le H$ for all $i$, let $M\ge1$, and let $K(M)$ be the kernel of $GL_n(\hat{\mathcal O}_K)\to GL_n(\hat{\mathcal O}_K/M\hat{\mathcal O}_K)$.
Let $S$ be a finite set of pairwise non-isomorphic cuspidal automorphic representations $\pi$ of $GL_n(\mathbf A_K)$ (Definition~\ref{gal-polarized}) such that, for every embedding $\tau:K\to\mathbf C$, the infinitesimal character of $\pi$ at $\tau$ is the multiset of the numbers $\frac{n-1}2+m_i$; such that $(\pi,\chi)$ is polarized for some Hecke character $\chi$; and such that $\pi$ has nonzero vectors fixed by $K(M)$.
Then $S$ has at most $c(M(H+1))^A$ elements.
\end{lemma}
\begin{proof}
There are $c_1,A_1$, depending only on $K$ and $n$, and an irreducible algebraic representation $\Xi$ of $GL_n(K\otimes_{\mathbf Q}\mathbf C)$ whose infinitesimal character at every $\tau$ is the negative of the multiset above, with $\dim\Xi\le c_1(H+1)^{A_1}$: at each $\tau$ the factor of $\Xi$ is built from the partition $\nu$ with $\nu_i=i+H-m_{n+1-i}$, whose entries lie between $1$ and $n+2H$, and Weyl's dimension formula \cite[Theorem~6.3]{FH91} bounds its dimension by a power of $n+2H+1$.
By Clozel's theorem (Theorem~\ref{lit-clozel-cohomological}), every $\pi\in S$ has nonzero cohomology with coefficients in $\Xi$.
Each $\pi\in S$ has nonzero vectors fixed by $K(3M)\subset K(M)$.
By Borel's theorem \cite{Bor83}, the cuspidal cohomology injects into the cohomology of the locally symmetric space of $GL_n$ over $K$ of level $K(3M)$ with coefficients in $\Xi$, so $|S|$ is at most the dimension of that cohomology.
That cohomology is the cohomology of $\Gamma=GL_n(K)$ with coefficients in the functions $GL_n(\mathbf A_{K,f})/K(3M)\to\Xi$, and its dimension is at most $c_2M^{A_2}\dim\Xi$, for $c_2,A_2$ depending only on $K$ and $n$.
Indeed, $\Gamma$ has finitely many orbits on $GL_n(\mathbf A_{K,f})/K(3)$, by the finiteness of the class number.
The stabilizer $\Gamma_g=\Gamma\cap gK(3)g^{-1}$ of $gK(3)$ is an arithmetic group, torsion free because $K(3)$ is, so the trivial module has a finite free resolution over it \cite{BS73}; if $c_g$ is the sum of the ranks, the cohomology of $\Gamma_g$ with coefficients in a finite-dimensional representation $A$ has total dimension at most $c_g\dim A$.
By Shapiro's lemma, the functions supported on the inverse image of the orbit of $gK(3)$ contribute the cohomology of $\Gamma_g$ with coefficients in the functions $gK(3)/K(3M)\to\Xi$, of dimension $[K(3):K(3M)]\dim\Xi$.
Finally $[K(3):K(3M)]\le(3M)^{n^2[K:\mathbf Q]}$.
So $|S|\le c_2c_1M^{A_2}(H+1)^{A_1}\le c(M(H+1))^A$ with $c=c_1c_2$ and $A=A_1+A_2$.
\end{proof}
\leanref{Development/AutomorphicArithmetic/AutomorphicVarieties/AutomorphicForms/IsoOfFrob.lean}{UpperTheory.V.Af.CardK}
\leanref{Development/AutomorphicArithmetic/AutomorphicVarieties/ModularityArgument/ThresholdArguments/DenseThresholdCases/TCardK.lean}{UpperTheory.V.Af.card_k_of_printed_results}
\leanref{Development/AutomorphicArithmetic/AutomorphicVarieties/ModularityArgument/ThresholdArguments/DenseThresholdCases/BrCardP.lean}{UpperTheory.V.Br.cardP_of_prints}
\leanref{Development/AutomorphicArithmetic/AutomorphicVarieties/ModularityArgument/ThresholdArguments/DenseThresholdCases/TWeylC.lean}{UpperTheory.V.Af.weyl_complex_of_printed_results}
\leanref{Development/AutomorphicArithmetic/AutomorphicVarieties/AutomorphicForms/TemperedLocalComponents/CohBoundFree.lean}{UpperTheory.V.Af.cohBoundFree}
\leanref{Development/AutomorphicArithmetic/AutomorphicVarieties/PublishedTheorems/GelfandKazhdanReal.lean}{UpperTheory.V.Published.BorelInjC}

\begin{definition}\label{sha-satake-slopes}
Let $\ell$ be a prime and $\iota:\bar{\mathbf Q}_\ell\cong\mathbf C$.
\begin{enumerate}
\item For $x\in\mathbf C$, put $|x|_\iota=|\iota^{-1}(x)|$, where $|\cdot|$ is the absolute value of $\bar{\mathbf Q}_\ell$ that extends the $\ell$-adic absolute value of $\mathbf Q_\ell$.
\item Let $E$ be a number field and $w$ a finite place of $E$. Let $\mathcal E_\iota(w)$ be the set of embeddings $\tau:E\to\mathbf C$ such that, for every $x\in\mathcal O_E$, $x\in w$ if and only if $|\tau(x)|_\iota<1$. Let $\varpi_w\in E^\times$ be an element of valuation $1$ at $w$, chosen once for each $w$.
\item Let moreover $n\ge1$, let $q$ be the norm of $w$, let $Z=(Z_\tau)_\tau$ be a family of multisets of complex numbers indexed by the embeddings $\tau:E\to\mathbf C$, and let $B$ be a multiset of nonzero complex numbers. Then $B$ has \emph{$\iota$-slopes bounded by $Z$ at $w$} if the following holds for all families of integers $(l_\tau)_\tau$ and $(h_\tau)_\tau$ such that every element of $Z_\tau$ is $\frac{n-1}2+k$ for an integer $k$ with $l_\tau\le k\le h_\tau$: every element $b$ of the multiset $q^{(n-1)/2}B$ satisfies
\[
 |b|_\iota\prod_{\tau\in\mathcal E_\iota(w)}|\tau(\varpi_w)|_\iota^{h_\tau}\le1
 \qquad\text{and}\qquad
 |b^{-1}|_\iota\prod_{\tau\in\mathcal E_\iota(w)}|\tau(\varpi_w)|_\iota^{-l_\tau}\le1.
\]
\end{enumerate}
\end{definition}
\leanref{Development/AutomorphicArithmetic/AutomorphicVarieties/ModularityArgument/IntegralDeterminants/PlaceConjAlt.lean}{UpperTheory.V.Z.Int.nu}
\leanref{Development/AutomorphicArithmetic/AutomorphicVarieties/ModularityArgument/IntegralDeterminants/PlaceConjAlt.lean}{UpperTheory.V.Z.Int.Ov}
\leanref{Development/AutomorphicArithmetic/AutomorphicVarieties/ModularityArgument/IntegralDeterminants/PlaceConjAlt.lean}{UpperTheory.V.Z.Int.Bnd}
\leanref{Development/AutomorphicArithmetic/AutomorphicVarieties/ModularityArgument/IntegralDeterminants/PlaceConjAlt.lean}{UpperTheory.V.Z.Int.Within}
\leanref{Development/AutomorphicArithmetic/AutomorphicVarieties/ModularityArgument/IntegralDeterminants/PlaceConjAlt.lean}{UpperTheory.V.Z.Int.Slope}

\begin{lemma}\label{sha-satake-integer-criterion}
Let $x$ be a complex number such that $|x|_\iota\le1$ for every prime $\ell$ and every $\iota:\bar{\mathbf Q}_\ell\cong\mathbf C$ (Definition~\ref{sha-satake-slopes}). Then $x$ is an algebraic integer.
\end{lemma}
\leanref{Development/AutomorphicArithmetic/AutomorphicVarieties/GaloisLibrary/IntCrit.lean}{UpperTheory.V.Lib.isIntegral_of_norm_le}
\begin{proof}
Two embeddings of a countable field, one into $\mathbf C$ and one into $\bar{\mathbf Q}_\ell$, are the restrictions of $\iota$ and $\iota^{-1}$ for some $\iota:\bar{\mathbf Q}_\ell\cong\mathbf C$, since both are algebraically closed of transcendence degree the continuum over the countable field.

First, $x$ is algebraic over $\mathbf Q$.
Otherwise, let $\iota_0:\bar{\mathbf Q}_2\cong\mathbf C$ and $t=\iota_0^{-1}(x)$.
Since $x$ and $t+\frac14$ are transcendental over $\mathbf Q$, applying this to $\mathbf Q(X)$ embedded by $X\mapsto x$ and $X\mapsto t+\frac14$ gives $\iota$ with $\iota^{-1}(x)=t+\frac14$.
Then $|\frac14|=4$ is at most $\max(|\iota^{-1}(x)|,|t|)\le1$, a contradiction.

Let $P$ be the minimal polynomial of $x$ over $\mathbf Q$, and let $\ell$ be a prime.
Every root $t$ of $P$ in $\bar{\mathbf Q}_\ell$ is $\iota^{-1}(x)$ for some $\iota$, by the same fact applied to $\mathbf Q(x)$ and the embedding that maps $x$ to $t$; so $|t|\le1$.
The coefficients of $P$ are, up to sign, the elementary symmetric functions of its roots, so they have $\ell$-adic absolute value at most $1$.
As this holds for every $\ell$, they lie in $\mathbf Z$.
\end{proof}

\begin{lemma}\label{sha-satake-place-embeddings}
Let $E$ be a number field, $\ell$ a prime, $\iota:\bar{\mathbf Q}_\ell\cong\mathbf C$ and $\tau:E\to\mathbf C$ an embedding.
\begin{enumerate}
\item There is a unique finite place $u$ of $E$ with $\tau\in\mathcal E_\iota(u)$ (Definition~\ref{sha-satake-slopes}), and there is a ring homomorphism $\psi:E_u\to\mathbf C$ that extends $\tau$ and satisfies $|\psi(x)|_\iota\le1$ if and only if $x\in\mathcal O_u$. Moreover $|\psi(x)|_\iota=|\tau(\varpi_u)|_\iota^{v_u(x)}$ for every $x\in E_u^\times$.
\item Let $E$ be a CM field with complex conjugation $c$ and maximal totally real subfield $F$, and let $w$ be a finite place of $E$ split over $F$. If $\tau\in\mathcal E_\iota(w)$, then $\tau\circ c\notin\mathcal E_\iota(w)$. If $\tau\in\mathcal E_\iota(u)$ for a place $u\ne w$ of $E$ over the same place of $F$ as $w$, then $\tau\circ c\in\mathcal E_\iota(w)$.
\end{enumerate}
\end{lemma}
\leanref{Development/AutomorphicArithmetic/AutomorphicVarieties/ModularityArgument/IntegralDeterminants/PlaceConjAlt.lean}{UpperTheory.V.Z.Int.PlaceHom}
\leanref{Development/AutomorphicArithmetic/AutomorphicVarieties/ModularityArgument/IntegralDeterminants/PlaceHom.lean}{UpperTheory.V.Z.Int.placeHom}
\leanref{Development/AutomorphicArithmetic/AutomorphicVarieties/ModularityArgument/IntegralDeterminants/PlaceConjAlt.lean}{UpperTheory.V.Z.Int.PlaceConj}
\leanref{Development/AutomorphicArithmetic/AutomorphicVarieties/ModularityArgument/IntegralDeterminants/PlaceConj.lean}{UpperTheory.V.Z.Int.placeConj}
\begin{proof}
(1) The function $x\mapsto|\tau(x)|_\iota$ is a nonarchimedean absolute value on $E$, at most $1$ on $\mathcal O_E$ and less than $1$ at $\ell$.
So $u=\{x\in\mathcal O_E:|\tau(x)|_\iota<1\}$ is a nonzero prime ideal, the only one with $\tau\in\mathcal E_\iota(u)$, and the absolute value is a power of the $u$-adic one \cite[Chapter~II, \S3]{Neu99}.
Hence $\iota^{-1}\circ\tau$ is continuous for the $u$-adic topology, with values in the finite extension of $\mathbf Q_\ell$ generated by its image, which is complete; it extends to $E_u$, and $\psi$ is $\iota$ composed with this extension.
The equivalence holds on $E$ by the definition of $u$, and on $E_u$ by continuity.
For $x\in E_u^\times$, $y=x\varpi_u^{-v_u(x)}$ and $y^{-1}$ lie in $\mathcal O_u$, so $|\psi(y)|_\iota=1$, which gives the last formula.

(2) For $x\in\mathcal O_E$, $|\tau(c(x))|_\iota<1$ if and only if $c(x)\in w$, that is $x\in c(w)$; so $\tau\circ c\in\mathcal E_\iota(c(w))$.
As $w$ is split over $F$, $c(w)\ne w$, and $c(w)$ is the only other place of $E$ over the place of $F$ under $w$.
The first claim follows from the uniqueness in (1).
In the second, $u=c(w)$, so $\tau\in\mathcal E_\iota(c(w))$ and $\tau\circ c\in\mathcal E_\iota(w)$.
\end{proof}

\begin{lemma}\label{sha-satake-norm-product}
Let $E$ be a number field, $w$ a finite place of $E$ of norm $q$, $\ell$ a prime and $\iota:\bar{\mathbf Q}_\ell\cong\mathbf C$. Then $0<|\tau(\varpi_w)|_\iota<1$ for every $\tau\in\mathcal E_\iota(w)$ (Definition~\ref{sha-satake-slopes}), and
\[
 \prod_{\tau\in\mathcal E_\iota(w)}|\tau(\varpi_w)|_\iota=|q|_\iota.
\]
\end{lemma}
\leanref{Development/AutomorphicArithmetic/AutomorphicVarieties/ModularityArgument/IntegralDeterminants/PlaceConjAlt.lean}{UpperTheory.V.Z.Int.ProdNorm}
\leanref{Development/AutomorphicArithmetic/AutomorphicVarieties/ModularityArgument/IntegralDeterminants/ProdNorm.lean}{UpperTheory.V.Z.Int.prodNorm}
\begin{proof}
For $\tau\in\mathcal E_\iota(w)$, $|\tau(\varpi_w)|_\iota<1$ by Lemma~\ref{sha-satake-place-embeddings}(1), and it is positive since $\varpi_w\ne0$.
As the class group of $E$ is finite, there are $h\ge1$ and $\alpha\in\mathcal O_E$ with $\alpha\mathcal O_E=w^h$.
By Lemma~\ref{sha-satake-place-embeddings}(1), $|\tau(\alpha)|_\iota=|\tau(\varpi_w)|_\iota^h$ for $\tau\in\mathcal E_\iota(w)$, and $|\tau(\alpha)|_\iota=1$ for every other $\tau$, which lies in $\mathcal E_\iota(u)$ for a place $u\ne w$ at which $\alpha$ is a unit.
Since $N_{E/\mathbf Q}(\alpha)=\pm q^h$,
\[
 \Big(\prod_{\tau\in\mathcal E_\iota(w)}|\tau(\varpi_w)|_\iota\Big)^h=\prod_\tau|\tau(\alpha)|_\iota=|N_{E/\mathbf Q}(\alpha)|_\iota=|q|_\iota^h,
\]
and taking $h$-th roots of these positive numbers gives the claim.
\end{proof}

\begin{lemma}\label{sha-satake-character-value}
Let $E$ be a totally complex number field, $w$ a finite place of $E$, $\ell$ a prime and $\iota:\bar{\mathbf Q}_\ell\cong\mathbf C$. Let $\kappa$ be a Hecke character of $E$, a continuous character of $\mathbf A_E^\times/E^\times$, that is unramified at $w$ and \emph{algebraic of infinity type $m$}: $m=(m_\tau)_\tau$ is a family of integers indexed by the embeddings $\tau:E\to\mathbf C$, and the component $\kappa_\infty$ of $\kappa$ on $(E\otimes\mathbf R)^\times$ satisfies $\kappa_\infty(x)=\prod_\tau\tau(x)^{m_\tau}$ for $x\in E^\times$. Then
\[
 |\kappa_w(\varpi_w)|_\iota\prod_{\tau\in\mathcal E_\iota(w)}|\tau(\varpi_w)|_\iota^{m_\tau}=1
\]
(Definition~\ref{sha-satake-slopes}).
\end{lemma}
\leanref{Development/AutomorphicArithmetic/AutomorphicVarieties/ModularityArgument/IntegralDeterminants/PlaceConjAlt.lean}{UpperTheory.V.Z.Int.CharVal}
\leanref{Development/AutomorphicArithmetic/AutomorphicVarieties/ModularityArgument/IntegralDeterminants/CharVal.lean}{UpperTheory.V.Z.Int.charVal}
\begin{proof}
Let $h\ge1$ and $\alpha\in\mathcal O_E$ with $\alpha\mathcal O_E=w^h$.
For each of the finitely many finite places $v$ at which $\kappa$ is ramified, $\kappa_v$ is trivial on an open subgroup of $\mathcal O_v^\times$, so there is $o_v\ge1$ with $\kappa_v(y)^{o_v}=1$ for all $y\in\mathcal O_v^\times$.
Let $M$ be the product of these $o_v$, and put $z=\alpha^M$.
Then $\kappa_v(z)=1$ for every finite place $v\ne w$, since $z$ is a unit at $v$, and an $o_v$-th power of a unit if $\kappa$ is ramified at $v$.
Since $\kappa$ is unramified at $w$ and $z\varpi_w^{-hM}$ is a unit at $w$, $\kappa_w(z)=\kappa_w(\varpi_w)^{hM}$.
As $\kappa$ is trivial on $E^\times$,
\[
 1=\kappa(z)=\kappa_w(\varpi_w)^{hM}\prod_\tau\tau(z)^{m_\tau}.
\]
By Lemma~\ref{sha-satake-place-embeddings}(1), $|\tau(z)|_\iota=|\tau(\varpi_w)|_\iota^{hM}$ for $\tau\in\mathcal E_\iota(w)$, and $|\tau(z)|_\iota=1$ for the other $\tau$.
Taking $|\cdot|_\iota$ shows that the $hM$-th power of the nonnegative number on the left of the claim is $1$.
\end{proof}

\begin{lemma}\label{sha-satake-newton-polygon}
Let $\ell$ be a prime and $\iota:\bar{\mathbf Q}_\ell\cong\mathbf C$, with $|\cdot|_\iota$ as in Definition~\ref{sha-satake-slopes}. Let $S$ be a multiset of $N\ge1$ complex numbers and $c\ge0$ a real number such that $|e_j(S)|_\iota\le c^j$ for $1\le j\le N$, where $e_j$ is the $j$-th elementary symmetric function. Then $|b|_\iota\le c$ for every $b\in S$.
\end{lemma}
\leanref{Development/AutomorphicArithmetic/AutomorphicVarieties/ModularityArgument/IntegralDeterminants/PlaceConjAlt.lean}{UpperTheory.V.Z.Int.RootBound}
\leanref{Development/AutomorphicArithmetic/AutomorphicVarieties/ModularityArgument/IntegralDeterminants/RootBound.lean}{UpperTheory.V.Z.Int.rootBound}
\begin{proof}
Suppose that $|b|_\iota>c$ for some $b\in S$.
Since $b$ is a root of $\prod_{s\in S}(X-s)=\sum_{j=0}^N(-1)^je_j(S)X^{N-j}$,
\[
 b^N=-\sum_{j=1}^N(-1)^je_j(S)b^{N-j}.
\]
The $j$-th term has $|\cdot|_\iota\le c^j|b|_\iota^{N-j}<|b|_\iota^N$, so by the ultrametric inequality the right side has $|\cdot|_\iota<|b|_\iota^N$, a contradiction.
\end{proof}

\begin{lemma}\label{sha-satake-weight-matrices}
Let $E$ be a number field, $u$ a finite place of $E$, $\ell$ a prime, $\iota:\bar{\mathbf Q}_\ell\cong\mathbf C$, and $\psi:E_u\to\mathbf C$ a ring homomorphism with $|\psi(x)|_\iota\le1$ for $x\in\mathcal O_u$ (Definition~\ref{sha-satake-slopes}). Let $n\ge1$, let $\lambda_0\ge\dots\ge\lambda_{n-1}$ be integers, let $\rho_\lambda$ be the action of $GL_n(\mathbf C)$ on its irreducible representation $\nabla(\lambda)$ attached to $\lambda$ (Definition~\ref{gal-definite-unitary}), and let $e$ be a basis of $\nabla(\lambda)$. Then there is $C\ge0$ such that the following holds. Let $l$ and $h$ be integers with $l\le\lambda_i\le h$ for all $i$, and let $g\in GL_n(E_u)$ have entries in $\mathcal O_u$. Then every entry $y$ of the matrix in the basis $e$ of $\rho_\lambda(\psi(g))$ satisfies $|y|_\iota\le C|\psi(\det g)|_\iota^l$, and every entry $y$ of the matrix of $\rho_\lambda(\psi(g^{-1}))$ satisfies $|y|_\iota\le C|\psi(\det g)|_\iota^{-h}$. Here $\psi$ is applied to matrices entrywise.
\end{lemma}
\leanref{Development/AutomorphicArithmetic/AutomorphicVarieties/ModularityArgument/IntegralDeterminants/PlaceConjAlt.lean}{UpperTheory.V.Z.Int.RepBound}
\leanref{Development/AutomorphicArithmetic/AutomorphicVarieties/ModularityArgument/IntegralDeterminants/RepBound.lean}{UpperTheory.V.Z.Int.repBound}
\leanref{Development/AutomorphicArithmetic/AutomorphicVarieties/HeckeAlgebras/OrdinaryHeckeCharacters/SqBasis.lean}{UpperTheory.V.Hi.sqBasis}
\begin{proof}
Let $\mathcal O\subset\mathbf C$ be the subring of the $x$ with $|x|_\iota\le1$; it contains $\psi(\mathcal O_u)$.
By the standard basis theorem for bitableaux \cite{DEP80}, $\nabla(\lambda)$ has a basis $b$ of weight vectors for the diagonal torus, with weights $\mu_k\in\mathbf Z^n$, such that the matrix in $b$ of every element of $GL_n(\mathcal O)$ has entries in $\mathcal O$.
Its lowest weight is $(\lambda_{n-1},\dots,\lambda_0)$, and $\sum_ia_i\mu_k(i)\ge\sum_ia_i\lambda_{n-1-i}$ for every $k$ and every non-increasing $a\in\mathbf Z^n$.
Let $c\ge0$ bound $|\cdot|_\iota$ on the entries of the two matrices that change the basis $b$ into $e$ and back, and put $C=c^2$.

Let $\varpi$ be a uniformizer of $E_u$.
By the Cartan decomposition, $g=k_1dk_2$ with $k_1,k_2\in GL_n(\mathcal O_u)$ and $d=\operatorname{diag}(\varpi^{a_0},\dots,\varpi^{a_{n-1}})$ for a non-increasing $a\in\mathbf Z^n$.
As $g$ has entries in $\mathcal O_u$, so does $d=k_1^{-1}gk_2^{-1}$, and $a_i\ge0$.
The determinants of $k_1$ and $k_2$ are units, so their images under $\psi$ have $|\cdot|_\iota=1$, and $|\psi(\det g)|_\iota=|\psi(\varpi)|_\iota^{\sum_ia_i}$.
In the basis $b$, $\rho_\lambda(\psi(k_1))$ and $\rho_\lambda(\psi(k_2))$ have entries in $\mathcal O$, and $\rho_\lambda(\psi(d))$ is diagonal with entries $\prod_i\psi(\varpi)^{a_i\mu_k(i)}$.
Since $0<|\psi(\varpi)|_\iota\le1$, $a_i\ge0$ and $\lambda_i\ge l$, these entries satisfy
\[
 |\psi(\varpi)|_\iota^{\sum_ia_i\mu_k(i)}\le|\psi(\varpi)|_\iota^{\sum_ia_i\lambda_{n-1-i}}\le|\psi(\varpi)|_\iota^{l\sum_ia_i}=|\psi(\det g)|_\iota^l.
\]
By the ultrametric inequality, the entries of the matrix of $\rho_\lambda(\psi(g))$ in $b$ satisfy the same bound, and those in $e$ the bound $C|\psi(\det g)|_\iota^l$.

For $g^{-1}=k_2^{-1}d^{-1}k_1^{-1}$, conjugating $d^{-1}$ by the permutation matrix that reverses the order of the coordinates gives $g^{-1}=k_3d'k_4$ with $k_3,k_4\in GL_n(\mathcal O_u)$ and $d'=\operatorname{diag}(\varpi^{a'_0},\dots,\varpi^{a'_{n-1}})$, where $a'_i=-a_{n-1-i}$ is non-increasing.
The same argument applies, with $\sum_ia'_i\lambda_{n-1-i}=-\sum_ia_i\lambda_i\ge-h\sum_ia_i$, and gives the bound $C|\psi(\det g)|_\iota^{-h}$.
\end{proof}

\begin{lemma}\label{sha-satake-eigenvalue-bound}
Let $\ell$ be a prime and $\iota:\bar{\mathbf Q}_\ell\cong\mathbf C$, with $|\cdot|_\iota$ as in Definition~\ref{sha-satake-slopes}. Let $G$ be a group with subgroups $\Gamma$ and $K$ such that $\Gamma\backslash G/K$ is finite, let $I$ be a finite set, $R:G\to GL_I(\mathbf C)$ a homomorphism, $s\subset G$ a finite subset, $a\in\mathbf C$, and $c,C\ge0$ real numbers. Let $F:G\to\mathbf C^I$ be a nonzero function with
\[
 F(\gamma g)=R(\gamma)F(g),\qquad F(gk)=F(g),\qquad\sum_{x\in s}F(gx)=aF(g)
\]
for all $\gamma\in\Gamma$, $k\in K$ and $g\in G$. Suppose that every entry $y$ of $R(k)$, for $k\in K$, has $|y|_\iota\le C$, and that every entry $y$ of $R(x_1\cdots x_r)$, for $r\ge0$ and $x_1,\dots,x_r\in s$, has $|y|_\iota\le Cc^r$. Then $|a|_\iota\le c$.
\end{lemma}
\leanref{Development/AutomorphicArithmetic/AutomorphicVarieties/ModularityArgument/IntegralDeterminants/PlaceConjAlt.lean}{UpperTheory.V.Z.Int.AbsBound}
\leanref{Development/AutomorphicArithmetic/AutomorphicVarieties/ModularityArgument/IntegralDeterminants/AbsBound.lean}{UpperTheory.V.Z.Int.absBound}
\begin{proof}
Put $N(g)=R(g)^{-1}F(g)$.
Then $N(\gamma g)=N(g)$ for $\gamma\in\Gamma$, $N(gk)=R(k^{-1})N(g)$ for $k\in K$, and
\[
 aN(g)=\sum_{x\in s}R(x)N(gx).
\]
Since $F\ne0$, there are $g_0$ and $i_0$ with $N(g_0)_{i_0}\ne0$.
The function $N$ is bounded: if $g_1,\dots,g_t$ represent $\Gamma\backslash G/K$, every $g$ is $\gamma g_ik$ with $\gamma\in\Gamma$ and $k\in K$, and $N(g)=R(k^{-1})N(g_i)$; by the ultrametric inequality its entries have $|\cdot|_\iota\le B$, where $B$ is $C$ times the largest $|N(g_i)_j|_\iota$.
Iterating the identity above $r$ times gives
\[
 a^rN(g_0)=\sum_{x_1,\dots,x_r\in s}R(x_1\cdots x_r)N(g_0x_1\cdots x_r),
\]
so $|a|_\iota^r|N(g_0)_{i_0}|_\iota\le Cc^rB$ for every $r\ge0$.
If $|a|_\iota>c$, the left side divided by $c^r$ is unbounded as $r\to\infty$ when $c>0$, and positive for $r=1$ when $c=0$; both contradict this inequality.
\end{proof}

\begin{lemma}\label{sha-satake-vector-forms}
Let $E$ be a CM field, $n\ge1$ and $U$ an automorphic representation of $U_n$ (Definition~\ref{gal-definite-unitary}). Write $U_n=U_{n,f}\times U_{n,\infty}$ for its finite part and its part at infinity, and for $g\in U_{n,\infty}$ and an embedding $\tau:E\to\mathbf C$ let $\tau(g)\in U(n)$ be the $\tau$-factor of $g$. For each infinite place $v$ of $E$ let $\varphi_v:E\to\mathbf C$ be an embedding over $v$, and suppose that $U$ has weight $\lambda_v$ at $\varphi_v$. Let $w$ be a finite place of $E$ at which $U$ is unramified. Then there are a nonzero multilinear map $f:\prod_v\nabla(\lambda_v)\to U$ and an open subgroup $K$ of $U_{n,f}$ such that:
\begin{enumerate}
\item for $g\in U_{n,\infty}$ and $x=(x_v)_v$, the right translate of $f(x)$ by $g$ is $f\big((\rho_{\lambda_v}(\varphi_v(g))x_v)_v\big)$;
\item for every $k\in K$ and every finite place $u$ of $E$, the component $k_u\in GL_n(E_u)$ lies in $GL_n(\mathcal O_u)$;
\item every $f(x)$ is fixed under right translation by $K$ and by $c_w^{-1}(GL_n(\mathcal O_w))$.
\end{enumerate}
\end{lemma}
\leanref{Development/AutomorphicArithmetic/AutomorphicVarieties/ModularityArgument/IntegralDeterminants/PlaceConjAlt.lean}{UpperTheory.V.Z.Int.VecForm}
\leanref{Development/AutomorphicArithmetic/AutomorphicVarieties/ModularityArgument/IntegralDeterminants/VecForm.lean}{UpperTheory.V.Z.Int.vecForm}
\begin{proof}
Let $N\subset U$ be the subspace of the vectors fixed by $c_w^{-1}(GL_n(\mathcal O_w))$.
It is nonzero, since $U$ is unramified at $w$, and stable under $U_{n,\infty}$, which commutes with $U_{n,w}$.
For each $v$, the sum of the images of the maps $\nabla(\lambda_v)\to U$ that intertwine $\rho_{\lambda_v}\circ\varphi_v$ with right translation by $U_{n,\infty}$ is nonzero and stable under $U_n$, since the $\varphi_v$-factor commutes with the other factors of $U_n$; so it is $U$, as $U$ is irreducible.
Since $\nabla(\lambda_v)$ stays irreducible under $U(n)$ \cite{FH91}, the $U(n)$-stable subspace $N$ is again a sum of copies of $\nabla(\lambda_v)$ for the $\varphi_v$-factor.
The factors at distinct $v$ commute, so $N$ contains the image of a nonzero map from $\bigotimes_v\nabla(\lambda_v)$ that intertwines $\bigotimes_v\rho_{\lambda_v}\circ\varphi_v$ with right translation by $U_{n,\infty}$; let $f$ be the corresponding multilinear map, which satisfies (1).
Its image spans a finite-dimensional subspace of $U$, and every vector of $U$ is fixed by an open subgroup of $U_{n,f}$, so this subspace is fixed by an open subgroup $K'$.
Then $K$, the intersection of $K'$ with the open subgroup of the $k\in U_{n,f}$ with $k_u\in GL_n(\mathcal O_u)$ for all $u$, satisfies (2) and (3).
\end{proof}

\begin{lemma}\label{sha-satake-hecke-operators}
Let $E$ be a CM field with maximal totally real subfield $F$, $n\ge1$, $U$ an automorphic representation of $U_n$, and $w$ a finite place of $E$ split over $F$, of norm $q$. Let $\sigma$ be a local component of $U$ at $w$ that is spherical with Satake parameters $B$ (Definition~\ref{gal-definite-unitary}). Then $U$ is unramified at $w$, $B$ has $n$ elements, and for every $0\le j\le n$ there are finite subsets $s_j,s'_j\subset U_{n,w}$ such that:
\begin{enumerate}
\item for $x\in s_j$ the matrix $c_w(x)$, and for $x\in s'_j$ the matrix $c_w(x)^{-1}$, has entries in $\mathcal O_w$ and determinant of valuation $j$;
\item for every $f\in U$ fixed by $c_w^{-1}(GL_n(\mathcal O_w))$, with $R$ denoting right translation,
\[
 q^{j(j-1)/2}\sum_{x\in s_j}R(x)f=e_j(q^{(n-1)/2}B)f,\qquad q^{j(j-1)/2}\sum_{x\in s'_j}R(x)f=e_j(q^{(n-1)/2}B^{-1})f,
\]
where $B^{-1}=\{b^{-1}:b\in B\}$ and $e_j$ is the $j$-th elementary symmetric function.
\end{enumerate}
\end{lemma}
\leanref{Development/AutomorphicArithmetic/AutomorphicVarieties/ModularityArgument/IntegralDeterminants/PlaceConjAlt.lean}{UpperTheory.V.Z.Int.EigAll}
\leanref{Development/AutomorphicArithmetic/AutomorphicVarieties/ModularityArgument/IntegralDeterminants/EigAll.lean}{UpperTheory.V.Z.Int.eigAll}
\begin{proof}
Since $w$ is split over $F$, $c_w:U_{n,w}\to GL_n(E_w)$ is an isomorphism; put $K_w=GL_n(\mathcal O_w)$.
Let $\varpi$ be a uniformizer of $E_w$ and $D_j=K_w\operatorname{diag}(\varpi,\dots,\varpi,1,\dots,1)K_w$, with $j$ entries $\varpi$.
It is a finite union of cosets $gK_w$; let $s_j$ be the preimages under $c_w$ of a set of representatives, and $s'_j$ the same for $D_j^{-1}$.
Every element of $D_j$ has entries in $\mathcal O_w$ and determinant of valuation $j$, which gives (1).
On the vectors fixed by $c_w^{-1}(K_w)$, $\sum_{x\in s_j}R(x)$ and $\sum_{x\in s'_j}R(x)$ are the Hecke operators $T_j$ and $T'_j$ of $D_j$ and $D_j^{-1}$.

Let $i:\sigma\to U$ be a nonzero map as in the definition of a local component.
It is injective, as $\sigma$ is irreducible, so the image of a nonzero spherical vector of $\sigma$ is a nonzero vector of $U$ fixed by $c_w^{-1}(K_w)$; hence $U$ is unramified at $w$.
The Satake parameters of $\sigma$ form a multiset of $n$ elements.
Let $C$ be the centralizer of $U_{n,w}$ in $U_n$.
Every $g\in U_n$ is $c_w^{-1}(g_w)z$ with $z\in C$, since $z=c_w^{-1}(g_w)^{-1}g$ is trivial at the place of $F$ under $w$.
So the sum of the subspaces $R(z)i(\sigma)$, $z\in C$, is stable under $U_n$, and it is $U$.
Each $R(z)\circ i$ is again a map as in the definition of a local component, and averaging over $c_w^{-1}(K_w)$ shows that a fixed vector $f\in U$ is a sum of images under these maps of spherical vectors of $\sigma$.
On the spherical vectors of $\sigma$, $T_j$ acts by $q^{j(n-j)/2}e_j(B)$, by the normalization of $B$.
The automorphism $\theta(g)={}^tg^{-1}$ of $GL_n(E_w)$ preserves $K_w$ and carries $D_j$ onto $D_j^{-1}$, so $T'_j$ on $\sigma$ is $T_j$ on $\sigma\circ\theta$, whose spherical vectors have Satake parameters $B^{-1}$ by the Satake isomorphism \cite{Car79}; so $T'_j$ acts by $q^{j(n-j)/2}e_j(B^{-1})$.
Finally $q^{j(j-1)/2}q^{j(n-j)/2}e_j(B)=q^{j(n-1)/2}e_j(B)=e_j(q^{(n-1)/2}B)$, and the same holds for $B^{-1}$.
\end{proof}

\begin{lemma}\label{sha-satake-operator-bound}
Let $E$ be a CM field with maximal totally real subfield $F$, $\ell$ a prime and $\iota:\bar{\mathbf Q}_\ell\cong\mathbf C$. Let $n\ge1$ and let $U$ be an automorphic representation of $U_n$ with weight $\lambda_\tau$ at every embedding $\tau:E\to\mathbf C$ (Definition~\ref{gal-definite-unitary}). Let $w$ be a finite place of $E$ split over $F$ at which $U$ is unramified, and let $(l_\tau)_\tau$ and $(h_\tau)_\tau$ be integers with $l_\tau\le\lambda_{\tau,i}\le h_\tau$ for all $\tau$ and $i$. Let $j\ge0$, let $s\subset U_{n,w}$ be finite and $a\in\mathbf C$, and suppose that $\sum_{x\in s}R(x)f=af$ for every $f\in U$ fixed by $c_w^{-1}(GL_n(\mathcal O_w))$, with $R$ denoting right translation.
\begin{enumerate}
\item If $c_w(x)$ has entries in $\mathcal O_w$ and determinant of valuation $j$ for every $x\in s$, then $|a|_\iota\le\big(\prod_{\tau\in\mathcal E_\iota(w)}|\tau(\varpi_w)|_\iota^{-h_\tau}\big)^j$.
\item If $c_w(x)^{-1}$ has entries in $\mathcal O_w$ and determinant of valuation $j$ for every $x\in s$, then $|a|_\iota\le\big(\prod_{\tau\in\mathcal E_\iota(w)}|\tau(\varpi_w)|_\iota^{l_\tau}\big)^j$.
\end{enumerate}
Here $|\cdot|_\iota$, $\mathcal E_\iota(w)$ and $\varpi_w$ are as in Definition~\ref{sha-satake-slopes}.
\end{lemma}
\leanref{Development/AutomorphicArithmetic/AutomorphicVarieties/ModularityArgument/IntegralDeterminants/PlaceConjAlt.lean}{UpperTheory.V.Z.Int.OpBound}
\leanref{Development/AutomorphicArithmetic/AutomorphicVarieties/ModularityArgument/IntegralDeterminants/OpBound.lean}{UpperTheory.V.Z.Int.opBound}
\leanref{Development/AbelianVarietyTheory/Proofs/DivisorProofs/FiniteDoubleCosetDefiniteUnitaryRationalOfIsOpen.lean}{GlobalTheory.BCT.finite_doubleCoset_ratFinU_isOpen}
\begin{proof}
For $\tau\in\mathcal E_\iota(w)$ put $\beta_\tau=|\tau(\varpi_w)|_\iota^{-jh_\tau}$ in case (1) and $\beta_\tau=|\tau(\varpi_w)|_\iota^{jl_\tau}$ in case (2).
The claim is $|a|_\iota\le\prod_{\tau\in\mathcal E_\iota(w)}\beta_\tau$.

\emph{Embeddings.} For each infinite place $v$ of $E$ choose an embedding $\varphi_v$ over $v$, in $\mathcal E_\iota(w)$ if one of the two embeddings over $v$ lies there.
By Lemma~\ref{sha-satake-place-embeddings}(1) there are a finite place $u_v$ with $\varphi_v\in\mathcal E_\iota(u_v)$ and an extension $\psi_v:E_{u_v}\to\mathbf C$ of $\varphi_v$ with $|\psi_v(x)|_\iota\le1$ if and only if $x\in\mathcal O_{u_v}$; and $\varphi_v\in\mathcal E_\iota(w)$ if and only if $u_v=w$.
By Lemma~\ref{sha-satake-place-embeddings}(2), $v\mapsto\varphi_v$ maps the places $v$ with $u_v=w$ bijectively onto $\mathcal E_\iota(w)$.
Also by Lemma~\ref{sha-satake-place-embeddings}(2), if $u_v\ne w$ then $u_v$ does not lie over the place of $F$ under $w$: otherwise $\varphi_v\circ c$, for the complex conjugation $c$ of $E$, would lie in $\mathcal E_\iota(w)$, and $\varphi_v$ would have been chosen there.

\emph{Vector-valued forms.} Lemma~\ref{sha-satake-vector-forms}, with $\lambda_v=\lambda_{\varphi_v}$, gives $f$ and $K$.
Choose a basis $e_v$ of each $\nabla(\lambda_{\varphi_v})$, and let $I$ be the set of families $i=(i_v)_v$ of indices of these bases.
For $g\in U_{n,f}$ let $F(g)\in\mathbf C^I$ have the entries $F(g)_i=f((e_{v,i_v})_v)(g)$.
For $h\in U_{n,f}$ let $M_v(h)$ be the matrix in $e_v$ of $\rho_{\lambda_{\varphi_v}}(\psi_v(h_{u_v}))$, and let $R'(h)$ be the matrix with entries $R'(h)_{ik}=\prod_vM_v(h^{-1})_{k_vi_v}$; then $R':U_{n,f}\to GL_I(\mathbf C)$ is a homomorphism.
The function $F$ is nonzero: every element of $U_n$ is $g_fg_\infty$ with $g_f\in U_{n,f}$ and $g_\infty\in U_{n,\infty}$, and $f(x)(g_fg_\infty)=f((\rho_{\lambda_{\varphi_v}}(\varphi_v(g_\infty))x_v)_v)(g_f)$ by Lemma~\ref{sha-satake-vector-forms}(1), so $F=0$ would force $f=0$.

\emph{Rational points.} Let $\Gamma\subset U_{n,f}$ be the group of the finite parts $\gamma_f$ of the elements $\gamma\in U_n(F)$, with parts $\gamma_\infty$ at infinity.
Since $f(x)$ is left invariant under $\gamma$ and $\gamma_\infty$ commutes with $U_{n,f}$, $f(x)(\gamma_fg)=f(x)(g\gamma_\infty^{-1})$ for $g\in U_{n,f}$, which is $f((\rho_{\lambda_{\varphi_v}}(\varphi_v(\gamma^{-1}))x_v)_v)(g)$ by Lemma~\ref{sha-satake-vector-forms}(1).
The matrix $\gamma^{-1}$ has entries in $E$ and $\psi_v$ extends $\varphi_v$, so $\varphi_v(\gamma^{-1})=\psi_v((\gamma_f^{-1})_{u_v})$.
Expanding in the bases $e_v$ gives $F(\gamma_fg)=R'(\gamma_f)F(g)$.

\emph{Hecke relation.} By Lemma~\ref{sha-satake-vector-forms}(3), $F(gk)=F(g)$ for $k\in K$, and the hypothesis on $s$, applied to the forms $f((e_{v,i_v})_v)$, gives $\sum_{x\in s}F(gx)=aF(g)$.

\emph{Bounds.} Let $C_v$ be the constant of Lemma~\ref{sha-satake-weight-matrices} for $u_v$, $\psi_v$, $\lambda_{\varphi_v}$ and $e_v$, and $D_v=\max(C_v,1)$.
For $k\in K$, $k_{u_v}\in GL_n(\mathcal O_{u_v})$ has a unit determinant, so $|\psi_v(\det k_{u_v})|_\iota=1$, and the second bound of that lemma bounds $|\cdot|_\iota$ of the entries of $M_v(k^{-1})$ by $C_v$; hence every entry of $R'(k)$ has $|\cdot|_\iota\le\prod_vD_v$.
Let $x_1,\dots,x_r\in s$ and $y=x_1\cdots x_r\in U_{n,w}$.
If $u_v\ne w$, then $u_v$ does not lie over the place of $F$ under $w$, so $y_{u_v}=1$ and $M_v(y^{-1})$ is the identity matrix.
If $u_v=w$ in case (1), then $c_w(y)$ has entries in $\mathcal O_w$ and determinant of valuation $jr$, so $|\psi_v(\det c_w(y))|_\iota=|\varphi_v(\varpi_w)|_\iota^{jr}$ by Lemma~\ref{sha-satake-place-embeddings}(1), and the second bound of Lemma~\ref{sha-satake-weight-matrices} bounds the entries of $M_v(y^{-1})$ by $C_v\beta_{\varphi_v}^r$.
If $u_v=w$ in case (2), the first bound of that lemma, applied to $c_w(y^{-1})$, which has entries in $\mathcal O_w$ and determinant of valuation $jr$, gives the same.
Hence every entry of $R'(y)$ has $|\cdot|_\iota\le\prod_vD_v\cdot\big(\prod_{\tau\in\mathcal E_\iota(w)}\beta_\tau\big)^r$.

\emph{Conclusion.} By Borel's finiteness of class numbers \cite{PR94}, $\Gamma\backslash U_{n,f}/K$ is finite.
Lemma~\ref{sha-satake-eigenvalue-bound}, with $G=U_{n,f}$, $R'$, $C=\prod_vD_v$ and $c=\prod_{\tau\in\mathcal E_\iota(w)}\beta_\tau$, gives $|a|_\iota\le c$.
\end{proof}

\begin{lemma}\label{sha-satake-unitary-slopes}
Let $E$ be a CM field with maximal totally real subfield $F$, $n\ge1$, and $U$ an automorphic representation of $U_n$ with weight $\lambda_\tau$ at every embedding $\tau:E\to\mathbf C$ (Definition~\ref{gal-definite-unitary}). Let $w$ be a finite place of $E$ split over $F$, and $\sigma$ a local component of $U$ at $w$ that is spherical with Satake parameters $B$. Then for every prime $\ell$ and every $\iota:\bar{\mathbf Q}_\ell\cong\mathbf C$, $B$ has $\iota$-slopes bounded at $w$ by the family of the multisets $\{\lambda_{\tau,k}+\frac{n-1}2-k:0\le k<n\}$ (Definition~\ref{sha-satake-slopes}).
\end{lemma}
\leanref{Development/AutomorphicArithmetic/AutomorphicVarieties/ModularityArgument/IntegralDeterminants/PlaceConjAlt.lean}{UpperTheory.V.Z.Int.CoreU}
\leanref{Development/AutomorphicArithmetic/AutomorphicVarieties/ModularityArgument/IntegralDeterminants/WeightBounds.lean}{UpperTheory.V.Z.Int.coreU}
\begin{proof}
Let $q$ be the norm of $w$, let $(l_\tau)$ and $(h_\tau)$ be as in Definition~\ref{sha-satake-slopes}, and put $P_\tau=|\tau(\varpi_w)|_\iota$ for $\tau\in\mathcal E_\iota(w)$; products over $\tau$ below run over $\mathcal E_\iota(w)$.
Then $l_\tau\le\lambda_{\tau,k}-k\le h_\tau$ for all $k$; as $\lambda_\tau$ is non-increasing, $l_\tau+n-1\le\lambda_{\tau,n-1}\le\lambda_{\tau,i}\le\lambda_{\tau,0}\le h_\tau$ for all $i$.
By Lemma~\ref{sha-satake-hecke-operators}, $U$ is unramified at $w$, $B$ has $n$ elements, and for $1\le j\le n$ the sums over $s_j$ and $s'_j$ act on the vectors fixed by $c_w^{-1}(GL_n(\mathcal O_w))$ as multiplication by
\[
 a_j=q^{-j(j-1)/2}e_j(q^{(n-1)/2}B)\qquad\text{and}\qquad a'_j=q^{-j(j-1)/2}e_j(q^{(n-1)/2}B^{-1}).
\]
Lemma~\ref{sha-satake-operator-bound}, with the bounds $l_\tau+n-1$ and $h_\tau$ for the weights, gives $|a_j|_\iota\le c^j$ and $|a'_j|_\iota\le c'^j$ for $c=\prod_\tau P_\tau^{-h_\tau}$ and $c'=\prod_\tau P_\tau^{l_\tau+n-1}$.
Since $|q|_\iota\le1$ by Lemma~\ref{sha-satake-norm-product}, $|e_j(q^{(n-1)/2}B)|_\iota\le c^j$ and $|e_j(q^{(n-1)/2}B^{-1})|_\iota\le c'^j$.
By Lemma~\ref{sha-satake-newton-polygon}, every $b\in q^{(n-1)/2}B$ has $|b|_\iota\le c$, which is the first inequality of Definition~\ref{sha-satake-slopes}, and every element of $q^{(n-1)/2}B^{-1}$ has $|\cdot|_\iota\le c'$.
For $b=q^{(n-1)/2}\beta$ with $\beta\in B$, the element $q^{(n-1)/2}\beta^{-1}$ is $q^{n-1}b^{-1}$, and $\prod_\tau P_\tau^{n-1}=|q|_\iota^{n-1}$ by Lemma~\ref{sha-satake-norm-product}.
So $|b^{-1}|_\iota|q|_\iota^{n-1}\le\prod_\tau P_\tau^{l_\tau}\cdot|q|_\iota^{n-1}$, and dividing by $|q|_\iota^{n-1}>0$ gives the second inequality.
\end{proof}

\begin{lemma}\label{sha-satake-descent}
Let $E$ be a CM field with maximal totally real subfield $F$ and $\eta$ the quadratic character of $E/F$, let $\ell$ be a prime and $a\ge1$, and let $\pi$ be a cuspidal automorphic representation of $GL_a(\mathbf A_E)$ such that $(\pi,\eta^a)$ is in self-dual scope for $\ell$ (Definition~\ref{gal-scope}). Let $Z$ be a finite set of finite places of $E$. Then there are a finite set $T$ of finite places of $E$ not split over $F$, disjoint from $Z$, and an automorphic representation $U$ of $U_a$ (Definition~\ref{gal-definite-unitary}) such that:
\begin{enumerate}
\item $U$ is unramified at every finite place of $E$ that is not split over $F$ and not in $T$;
\item at every finite place $w$ of $E$ split over $F$, every local component of $\pi$ at $w$ is a local component of $U$ at $w$;
\item if $U$ has weight $\lambda$ at an embedding $\tau:E\to\mathbf C$, then $Z_\tau(\pi)=\{\lambda_k+\frac{a-1}2-k:0\le k<a\}$.
\end{enumerate}
This is the descent of Labesse \cite{Lab11} from $GL_a$ to the definite unitary group.
\end{lemma}
\leanref{Development/AbelianVarietyTheory/Definitions/GaloisGeometry/MatchOfMatchOffTEmpty.lean}{GlobalTheory.Cong.DescOffT}
\leanref{Development/AbelianVarietyTheory/Definitions/GaloisGeometry/MatchOfMatchOffTEmpty.lean}{GlobalTheory.Cong.MatchOffT}
\leanref{Development/AutomorphicArithmetic/Foundations/GaloisRepresentations/BaseChangeDescent/BcBackOffT.lean}{UpperTheory.Gal.Desc.descOffT}

\begin{lemma}\label{sha-satake-translates}
Let $E$ be a number field, $w$ a finite place of $E$ with ring of integers $\mathcal O_w$, $a\ge1$ and $K_w=GL_a(\mathcal O_w)$. Let $\rho$ be a representation of $GL_a(E_w)$ on a complex vector space $V$, and $X\subset V$ the span of the vectors $\rho(k)z-z$ with $k\in K_w$ and $z\in V$. Let $\mu:E_w^\times\to\mathbf C^\times$ be a homomorphism and $u\in\mathcal O_w^\times$ with $\mu(u)\ne1$. If $y\in V$ satisfies $\rho(k)y=\mu(\det k)^{-1}y$ for every $k\in K_w$, then $\rho(g)y\in X$ for every $g\in GL_a(E_w)$.
\end{lemma}
\leanref{Development/AbelianVarietyTheory/Proofs/SatakeParameters/TranslateMemTwistedJacquetKernelRamifiedStatement.lean}{GlobalTheory.T.translate_mem_twistedJacquetKernel_ramified}
\leanref{Development/AbelianVarietyTheory/Proofs/SatakeParameters/TranslateMemTwistedJacquetKernelRamified.lean}{GlobalTheory.CartanCoinv.translate_mem_twKer}
\begin{proof}
The subspace $X$ is stable under $\rho(k)$ for $k\in K_w$, since $\rho(k)x=(\rho(k)x-x)+x$.
By the Cartan decomposition, $g=k_1Dk_2$ with $k_1,k_2\in K_w$ and $D$ diagonal.
Let $d\in K_w$ be the diagonal matrix with entries $u,1,\dots,1$, so that $\det d=u$ and $dD=Dd$.
Put $c=\mu(u)^{-1}$, so that $c\ne1$.
Then $\rho(d)\rho(D)y=\rho(D)\rho(d)y=c\,\rho(D)y$, so $\rho(D)y=(c-1)^{-1}(\rho(d)\rho(D)y-\rho(D)y)$ lies in $X$.
Hence $\rho(g)y=\mu(\det k_2)^{-1}\rho(k_1)\rho(D)y$ lies in $X$.
\end{proof}

\begin{lemma}\label{sha-satake-extension}
Let $E$ be a CM field with complex conjugation $c$ and maximal totally real subfield $F$, and let $s\in\mathbf Z$. Let $\psi$ be a Hecke character of $F$, a continuous character of $\mathbf A_F^\times/F^\times$, that is algebraic of parallel weight $s$: its component $\psi_\infty$ on $(F\otimes\mathbf R)^\times$ satisfies $\psi_\infty(x)=\prod_\phi\phi(x)^s$ for $x\in F^\times$, where $\phi$ runs over the embeddings $F\to\mathbf R$. Suppose that at every finite place $w$ of $E$ not split over $F$, $E$ is unramified over $F$ and $\psi$ is unramified at the place of $F$ under $w$. Let $w_0$ be a finite place of $E$ split over $F$. Then there is a Hecke character $\kappa$ of $E$, algebraic of some infinity type $m$ (Lemma~\ref{sha-satake-character-value}), such that $\kappa|_{\mathbf A_F^\times}=\psi$ and $\kappa$ is unramified at $w_0$ and at every finite place of $E$ not split over $F$. Compare \cite[Lemma~4.1.4]{CHT08}: for $\psi$ algebraic, with $\psi_v(-1)$ independent of the real place $v$, and unramified at the places of a finite set $S$ of places of $F$ split in $E$, it gives an algebraic $\kappa$, unramified at the places over $S$ and with exponents prescribed at the embeddings of a CM type, such that $\kappa\cdot(\kappa\circ c)=\psi\circ N_{E/F}$.
\end{lemma}
\leanref{Development/AbelianVarietyTheory/Proofs/SatakeParameters/ExistsAlgebraicHeckeCharSplitStatement.lean}{GlobalTheory.T.exists_algebraic_hecke_char_split}
\leanref{Development/AbelianVarietyTheory/Proofs/SatakeParameters/ExistsAlgebraicHeckeCharSplit.lean}{GlobalTheory.HeckeUnrExt.exists_ext_unr}
\begin{proof}
Let $v_0$ be the place of $F$ under $w_0$.
For a finite place $w$ of $E$ let $\mathcal O_w$ be the ring of integers of $E_w$, and for an idele $x$ let $x_w$ be its component at $w$.
For each infinite place $w$ of $E$ choose an embedding $\phi_w:E\to\mathbf C$ that induces it, and let $\iota_w:E_w\to\mathbf C$ be its continuous extension.

\emph{A level for $\psi$.}
The character $\psi$ is unramified at all but finitely many finite places $v$ of $F$, and at each of the others it is trivial on $1+\varpi_v^{n_v}\mathcal O_v$ for some $n_v\ge1$.
So there is an integer $M\ge1$ such that $\psi(a)=1$ for every idele $a$ of $F$ with $a_\infty=1$ and $a_v\in1+M\mathcal O_v$ for every finite place $v$: every $\psi_v(a_v)$ is $1$, and $\psi(a)$ is the limit of their finite products, as $\psi$ is continuous.

\emph{An auxiliary place.}
There is a finite place $w_1$ of $E$, split over $F$, that does not lie over $v_0$.
Indeed, there is a rational prime $p>N(w_0)$ that splits completely in $E$.
The places of $E$ over $p$ are split over $F$, and a place $w_1$ over $p$ does not lie over $v_0$: otherwise $w_0$ would contain $p$ and the integer $N(w_0)$, which are coprime.

\emph{A congruence subgroup.}
Let $K=[E:\mathbf Q]+1$, $q_1=N(w_1)$ and $L=3Mq_1^K$.
Let $T$ be the set of finite places of $E$ split over $F$ other than $w_0$, and $T'\subset T$ the set of $w$ such that every place of $E$ over the place of $F$ under $w$ lies in $T$; then $w_1\in T'$.
For a set $S$ of finite places of $E$, let $U_S$ be the group of ideles $u$ of $E$ with $u_w\in\mathcal O_w^\times$ for every finite place $w$ and $u_w\in1+L\mathcal O_w$ for $w\in S$, with no condition at the infinite places.
Then $U_T\subset U_{T'}$, and $U_T$ is open in $\mathbf A_E^\times$, since it contains $(E\otimes\mathbf R)^\times\times\prod_w(1+L\mathcal O_w)$.
Let $\kappa_0:\mathbf A_E^\times\to\mathbf C^\times$ be the continuous homomorphism $\kappa_0(x)=\prod_w\iota_w(x_w)^s$, the product over the infinite places $w$ of $E$.

\emph{Agreement on $\mathbf A_F^\times$.}
Let $b\in\mathbf A_F^\times$ lie in $U_T$. Then $\psi(b)=\kappa_0(b)$.
Indeed, write $b=b_\infty b_f$, where $b_\infty$ has the components of $b$ at the infinite places and $1$ at the finite places, and $b_f$ the other way round.
First, $\psi(b_f)=1$.
Every $b_v$, for $v$ finite, is a unit, since its image in $E_w$ for a place $w$ over $v$ is a unit; so $\psi_v(b_v)=1$ if $\psi$ is unramified at $v$.
Let $v$ be a finite place of $F$ at which $\psi$ is ramified.
By hypothesis the places of $E$ over $v$ are split over $F$, and one of them, $w'$, is not $w_0$; so $w'\in T$, and $b_v-1$, whose image in $E_{w'}$ lies in $L\mathcal O_{w'}$, lies in $L\mathcal O_v\subset M\mathcal O_v$.
The idele with component $b_v$ at $v$ and $1$ at every other place satisfies the condition of the first step, so $\psi_v(b_v)=1$.
As $\psi$ is continuous, $\psi(b_f)=1$.
Second, $\psi(b_\infty)=\kappa_0(b_\infty)$.
Both sides are continuous functions of $b_\infty\in(F\otimes\mathbf R)^\times$, and they agree on $F^\times$: for $x\in F^\times$, $\psi_\infty(x)=\prod_\phi\phi(x)^s$, and $\kappa_0$ of the idele $x_\infty$ is $\prod_w\phi_w(x)^s$, the same product, since $w\mapsto\phi_w|_F$ is a bijection from the infinite places of $E$ onto the embeddings $F\to\mathbf R$.
As $F$ is dense in $F\otimes\mathbf R$, $F^\times$ is dense in $(F\otimes\mathbf R)^\times$.
Since $\kappa_0$ depends only on the components at the infinite places, $\kappa_0(b)=\kappa_0(b_\infty)=\psi(b_\infty)=\psi(b)$.

\emph{Units.}
Let $\gamma\in E^\times$ lie in $U_{T'}$, with $|\tau(\gamma)|=1$ for every embedding $\tau:E\to\mathbf C$. Then $\gamma=1$.
Indeed, $\gamma$ is a unit at every finite place, so $\gamma\in\mathcal O_E$, and $\gamma-1\in L\mathcal O_{w_1}$ since $w_1\in T'$.
As $q_1\in w_1$, $L\in w_1^K$, so $\gamma-1\in w_1^K$.
If $\gamma\ne1$, then $q_1^K$ divides the nonzero integer $N_{E/\mathbf Q}(\gamma-1)$, and
\[
 2^K\le q_1^K\le|N_{E/\mathbf Q}(\gamma-1)|=\prod_\tau|\tau(\gamma)-1|\le2^{[E:\mathbf Q]},
\]
which contradicts $K>[E:\mathbf Q]$.

\emph{Gluing.}
Let $G=\mathbf A_F^\times\times U_T\times E^\times$, and let $p:G\to\mathbf A_E^\times$ and $f:G\to\mathbf C^\times$ be the homomorphisms $p(b,u,\gamma)=bu\gamma$ and $f(b,u,\gamma)=\psi(b)\kappa_0(u)$.
Then $f$ is trivial on the kernel of $p$.
Indeed, let $bu\gamma=1$, and apply $c$, which acts on $\mathbf A_E^\times$ through its action on $E$.
It fixes $b$ and maps $E^\times$ to itself, and $c(u)\in U_{T'}$: the component of $c(u)$ at $w$ is the image under $c$ of $u_{c(w)}$, and for $w\in T'$ the place $c(w)$ lies over the same place of $F$ as $w$, so $c(w)\in T$.
From $b\,c(u)\,c(\gamma)=1=bu\gamma$, $c(\gamma)/\gamma=u/c(u)$, which lies in $U_{T'}$.
For every embedding $\tau$, $\tau(c(\gamma))$ is the complex conjugate of $\tau(\gamma)$, as $E$ is a CM field, so $|\tau(c(\gamma)/\gamma)|=1$.
By the previous step $c(\gamma)=\gamma$, that is, $\gamma\in F^\times$.
Then $b\gamma\in\mathbf A_F^\times$ and $b\gamma=u^{-1}\in U_T$, so $\psi(b\gamma)=\kappa_0(u)^{-1}$ by the agreement on $\mathbf A_F^\times$; as $\psi$ is trivial on $F^\times$, $\psi(b)\kappa_0(u)=1$.
So $f$ induces a homomorphism on the image of $p$.
Since $\mathbf C^\times$ is a divisible group, it extends to a homomorphism $\kappa:\mathbf A_E^\times\to\mathbf C^\times$ with $\kappa(bu\gamma)=\psi(b)\kappa_0(u)$ for all $(b,u,\gamma)\in G$.
In particular $\kappa$ is trivial on $E^\times$, $\kappa|_{\mathbf A_F^\times}=\psi$, and $\kappa=\kappa_0$ on $U_T$.

\emph{Properties of $\kappa$.}
It agrees with the continuous $\kappa_0$ on the open subgroup $U_T$, so it is continuous at $1$, hence continuous: $\kappa$ is a Hecke character of $E$.
Let $w$ be a finite place of $E$ not in $T$, that is, $w=w_0$ or $w$ not split over $F$.
An idele with a unit of $\mathcal O_w$ at $w$ and $1$ at every other place lies in $U_T$, and $\kappa_0$ is $1$ on it; so $\kappa$ is unramified at $w$.
Finally, for $\gamma\in E^\times$ the idele $\gamma_\infty$, with components $\gamma$ at the infinite places and $1$ at the finite places, lies in $U_T$, so
\[
 \kappa_\infty(\gamma)=\kappa_0(\gamma_\infty)=\prod_w\phi_w(\gamma)^s=\prod_\tau\tau(\gamma)^{m_\tau},
\]
where $m_\tau=s$ if $\tau=\phi_w$ for the place $w$ induced by $\tau$, and $m_\tau=0$ otherwise.
\end{proof}

\begin{lemma}\label{sha-satake-twist}
Let $E$ be a CM field with complex conjugation $c$, maximal totally real subfield $F$, and $\eta$ the quadratic character of $E/F$. Let $\ell$ be a prime, $a\ge1$, $\pi$ a cuspidal automorphic representation of $GL_a(\mathbf A_E)$ and $\chi$ a Hecke character of $F$ such that $(\pi,\chi)$ is in good scope for $\ell$ (Definition~\ref{gal-scope}). Let $w_0$ be a finite place of $E$ split over $F$. Then there is a Hecke character $\kappa$ of $E$, unramified at $w_0$ and algebraic of some infinity type $m$ (Lemma~\ref{sha-satake-character-value}), such that $(\pi\otimes\kappa,\eta^a)$ is in self-dual scope for $\ell$, where $\pi\otimes\kappa$ is the twist of $\pi$ by $\kappa\circ\det$.
\end{lemma}
\leanref{Development/AbelianVarietyTheory/Proofs/SatakeParameters/ExistsGoodConjSelfDualTwistSplitUnramifiedStatement.lean}{GlobalTheory.T.exists_good_conjSelfDual_twist_split_unramified}
\leanref{Development/AbelianVarietyTheory/Proofs/SatakeParameters/ExistsGoodConjSelfDualTwistSplitUnramified.lean}{GlobalTheory.GoodCsdTwist.exists_goodCsd_twist}
\begin{proof}
\emph{The character $\chi$ is algebraic of parallel weight.}
For an embedding $\tau:E\to\mathbf C$, let $s_\tau$ be the exponent of $\chi\circ N_{E/F}$ at $\tau$: its component at the infinite place of $\tau$ is $z\mapsto(z\bar z)^{s_\tau}$, read through $\tau$, and twisting by it shifts $Z_\tau$ by $s_\tau$; also $s_{\tau\circ c}=s_\tau$.
The polarization pairs $\pi\otimes(\chi\circ N_{E/F}\circ\det)^{-1}$ with $\pi$ composed with $c$, so $Z_{\tau\circ c}(\pi)$ is the negative of $Z_\tau(\pi\otimes(\chi\circ N_{E/F}\circ\det)^{-1})$, that is, $s_\tau-z\in Z_{\tau\circ c}(\pi)$ for every $z\in Z_\tau(\pi)$.
As $\pi$ is regular algebraic, $z$ and $s_\tau-z$ lie in $\frac{a-1}2+\mathbf Z$, so $s_\tau\in\mathbf Z$.
For a totally positive $x\in F^\times$, $x$ is the norm of the element of $(E\otimes\mathbf R)^\times$ with positive real components $\sqrt x$, so $\chi_\infty(x)=\prod_\phi\phi(x)^{s_\phi}$, where $\phi$ runs over the embeddings $F\to\mathbf R$ and $s_\phi=s_\tau$ for $\tau$ over $\phi$.
For a unit $\varepsilon$ of $F$, $\chi_\infty(\varepsilon)^{-1}$ is the value of $\chi$ on the finite idele $\varepsilon$, which has absolute value $1$, as $\chi$ is continuous on the compact group $\prod_v\mathcal O_v^\times$.
Applied to the totally positive unit $\varepsilon^2$, this gives $\sum_\phi s_\phi\log|\phi(\varepsilon)|=0$ for every unit $\varepsilon$.
By Dirichlet's unit theorem the vectors $(\log|\phi(\varepsilon)|)_\phi$ span the hyperplane of the vectors with sum $0$, so the $s_\phi$ are all equal to one integer $b$.

\emph{A character of $F$.} Put $\psi=\chi^{-1}\eta^{a+|b|}$.
Since $\eta^2=1$, $\psi_\infty(x)=\prod_\phi\phi(x)^{-b}$ for totally positive $x$.
At a real place $v$ of $F$, $\chi_v(-1)=(-1)^a$ by the polarization and $\eta_v(-1)=-1$, so $\psi_v(-1)=(-1)^a(-1)^{a+|b|}=(-1)^b$.
Hence $\psi_\infty(x)=\prod_\phi\phi(x)^{-b}$ for every $x\in F^\times$.

\emph{Ramification.} Let $w$ be a finite place of $E$ not split over $F$, and $v$ the place of $F$ under $w$.
By good scope, $E$ is unramified over $F$ at $w$ and $\pi$ has Satake parameters at $w$.
As $c(w)=w$, the polarization identifies the local component of $\pi\otimes(\chi\circ N_{E/F}\circ\det)^{-1}$ at $w$ with the contragredient of that of $\pi$ at $w$ composed with $c$, so it is spherical.
Since $\pi$ and its twist by $\xi\circ\det$, with $\xi=(\chi\circ N_{E/F})^{-1}$, are both spherical at $w$, the character $\xi$ is unramified at $w$, and so is $\chi\circ N_{E/F}$.
Indeed, suppose that $\xi_w(u)\ne1$ for a unit $u$ of $E_w$, and let $K_w=GL_a(\mathcal O_w)$.
As the twist is spherical at $w$, there is a nonzero $y$ in the space of $\pi$ with $\pi_w(k)y=\xi_w(\det k)^{-1}y$ for $k\in K_w$.
Every translate $\pi_w(g)y$, $g\in GL_a(E_w)$, lies in the span $X$ of the vectors $\pi_w(k)z-z$ with $k\in K_w$, by Lemma~\ref{sha-satake-translates} for $\rho=\pi_w$ and $\mu=\xi_w$.
Let $W$ be the subspace of the vectors of $\pi$ all of whose translates under $GL_a(E_w)$ lie in $X$.
Every operator of the global representation is a translation at $w$ composed with an operator that commutes with the translations at $w$, so $W$ is stable under the whole representation; it contains $y\ne0$, so it is the whole space by irreducibility.
But $\pi$ has a nonzero $K_w$-fixed vector, which then lies in $X$; and the $K_w$-fixed vectors and $X$ are complementary, since $\pi_w$ is smooth and $K_w$ is compact open.
This contradiction shows that $\xi$ is unramified at $w$.
% The ramification step is T.heckeChar_unrAt_of_spherical_and_twist_spherical, proved as SphTwistUnr.unrAt_of_isSpherical_twist: the local computation is T.translate_mem_twistedJacquetKernel_ramified (CartanCoinv.translate_mem_twKer, Lemma sha-satake-translates), X is twKer (locRep w) (KGL w a) 1, the decomposition of operators is T.exists_localTranslation_comp_commuting_operator, irreducibility is pi.irred, and the complement is LocGL.isCompl_fixed_twKer.
As $E_w/F_v$ is unramified, the norm maps $\mathcal O_w^\times$ onto $\mathcal O_v^\times$, so $\chi$ is unramified at $v$, and so are $\eta$ and $\psi$.

\emph{The twist.} The character $\psi$ of $F$ is algebraic of parallel weight $-b$, and at every finite place $w$ of $E$ not split over $F$, $E$ is unramified over $F$ at $w$ and $\psi$ is unramified at the place under $w$; and $w_0$ is split over $F$. So Lemma~\ref{sha-satake-extension} gives a Hecke character $\kappa$ of $E$, algebraic of some infinity type $m$, with $\kappa|_{\mathbf A_F^\times}=\psi$, unramified at $w_0$ and at every finite place not split over $F$.
% The extension lemma is T.exists_algebraic_hecke_char_split, proved as HeckeUnrExt.exists_ext_unr (Lemma sha-satake-extension).

\emph{Self-dual scope.} The representation $\pi\otimes\kappa$ is regular algebraic, since $Z_\tau(\pi\otimes\kappa)=\{z+m_\tau:z\in Z_\tau(\pi)\}$.
At a real place $v$ of $F$, $\eta_v(-1)^a=(-1)^a$.
For $x\in\mathbf A_E^\times$, $\kappa(x)\kappa(c(x))=\psi(N_{E/F}(x))=\chi^{-1}(N_{E/F}(x))$, as $\eta\circ N_{E/F}=1$.
So $\kappa\circ c=\kappa^{-1}\cdot(\chi\circ N_{E/F})^{-1}$, and $(\pi\otimes\kappa)^c\cong\pi^\vee\otimes(\chi\circ N_{E/F}\circ\det)\otimes(\kappa\circ c\circ\det)\cong(\pi\otimes\kappa)^\vee$, so $(\pi\otimes\kappa,\eta^a)$ is polarized.
The conditions of good scope on $E$ are unchanged, and at a finite place $w$ not split over $F$, $\pi\otimes\kappa$ has the Satake parameters $\{\beta\kappa_w(\varpi_w):\beta\in B\}$, where $B$ are those of $\pi$, since $\kappa$ is unramified at $w$.
\end{proof}

\begin{remark}\label{sha-satake-twist-support}
The ramification step in the proof of Lemma~\ref{sha-satake-twist} also follows from the theory of the supercuspidal support \cite{BZ77}: the supercuspidal support of a spherical representation of $GL_a(E_w)$ consists of unramified characters, and twisting by $\xi\circ\det$ multiplies each of them by $\xi_w$, so if $\pi_w$ and its twist are both spherical, $\xi_w$ is unramified.
\end{remark}

\begin{lemma}\label{sha-satake-quadratic-extension}
Let $E$ be a CM field, $\Omega$ an algebraic closure of $E$, $\ell$ a prime, $a\ge1$, $\pi$ a cuspidal automorphic representation of $GL_a(\mathbf A_E)$ and $w_0$ a finite place of $E$. Then there are infinitely many quadratic extensions $Q\subset\Omega$ of $E$ that are CM fields, with maximal totally real subfield $Q^+$, such that:
\begin{enumerate}
\item $D_A\subset G_Q$ for every valuation ring $A$ of $\Omega$ over $w_0$;
\item every place of $Q$ over $w_0$ is split over $Q^+$;
\item $Q$ is ramified over $E$ at some place over a place of $E$ where $\pi$ has Satake parameters;
\item every base change $P$ of $\pi$ to $Q$ (Definition~\ref{gal-solvable-base-change}) satisfies the conditions that good scope for $\ell$ adds (Definition~\ref{gal-scope}): $Q$ contains $\sqrt{-d}$ for an integer $d>0$, $[Q^+:\mathbf Q]$ is even, and every place of $Q$ not split over $Q^+$ is prime to $\ell$, unramified over $Q^+$, and a place where $P$ has Satake parameters.
\end{enumerate}
\end{lemma}
\leanref{Development/AbelianVarietyTheory/Proofs/SatakeParameters/InfiniteQuadraticCMGoodBaseChangeSplitStatement.lean}{GlobalTheory.T.infinite_quadraticCM_goodBaseChange_split}
\leanref{Development/AbelianVarietyTheory/Proofs/SatakeParameters/InfiniteQuadraticCMGoodBaseChangeSplit.lean}{GlobalTheory.GalCut.infinite_qcm_split_good}
\begin{proof}
Let $F$ be the maximal totally real subfield of $E$, and write $E=F(\beta)$ with $\beta_0=\beta^2\in\mathcal O_F$ totally negative.
Put $N=|N_{F/\mathbf Q}(2\beta_0)|$.
For every CM field $K$ that contains $E$, $\beta$ does not lie in the totally real field $K^+$, so $K=K^+(\beta)$ with $\beta^2=\beta_0$, and every place of $K$ prime to $N$, hence to $2\beta_0$, is unramified over $K^+$.
Let $q_0$ be the prime under $w_0$, and $S$ the finite set of finite places of $E$ where $\pi$ has no Satake parameters.
Let $T$ be the set of primes made of $2$, $\ell$, $q_0$, the prime factors of $N$ and of the discriminant of $E$, and the primes under the places in $S$.
By Dirichlet's theorem there are infinitely many primes $p\notin T$ with $p\equiv-1$ modulo $8\prod_{t\in T}t$.

Let $p$ be such a prime, $\delta\in\Omega$ with $\delta^2=-p$, and $Q=E(\delta)$.
The prime $p$ is odd and unramified in $E$, so $-p$ has valuation $1$ at every place $u$ of $E$ over $p$; hence $\delta\notin E$, and $Q$ is ramified over $E$ at every place over $u$.
The element $\gamma=\delta\beta$ satisfies $\gamma^2=-p\beta_0$, which lies in $F$ and is totally positive, and $\gamma\notin E$; so $Q=E(\gamma)$, the field $F(\gamma)$ is totally real, and $Q=F(\gamma)(\beta)$ is a CM field, quadratic over $E$.
In particular $Q^+$ is totally real, $\delta^2=-p\in Q^+$, $\delta\notin Q^+$, and $Q=Q^+(\delta)$.

Let $t\in T$ and write $p+1=8tk$ with $k\in\mathbf Z$. Then $-p=1+4z$ with $z=-2tk$.
At every place $v$ of $Q^+$ over $t$, $z$ lies in the maximal ideal, so $y^2+y=z$ has a solution in $\mathcal O_v$ by Hensel's lemma, and $-p=(2y+1)^2$ is a square in $Q^+_v$.
Since $Q=Q^+(\delta)$ with $\delta^2=-p$, every place of $Q^+$ over $t$ splits in $Q$: every place of $Q$ over $t$ is split over $Q^+$.
Likewise $-p$ is a square in $\mathbf Q_{q_0}$, so every element of $D_A$, for $A$ over $w_0$, fixes $\delta$, which is a square root of $-p$ in the completion $E_{w_0}\supset\mathbf Q_{q_0}$; so $D_A\subset G_Q$.
This gives (1), and (2) since $q_0\in T$.
For (3), a place of $Q$ over $p$ is ramified over $E$, and $\pi$ has Satake parameters at the place of $E$ under it, which is not in $S$ since $p\notin T$.

For (4), $Q$ contains $\delta$ with $\delta^2=-p$, and $[Q^+:\mathbf Q]=[E:\mathbf Q]=2[F:\mathbf Q]$ is even.
Let $w$ be a place of $Q$ not split over $Q^+$, over the prime $t$.
Then $t\notin T$ by the previous paragraph, so $t\ne\ell$, $t$ is prime to $N$ and $w$ is unramified over $Q^+$, and the place $v$ of $E$ under $w$ is not in $S$.
So $\pi$ has Satake parameters $B$ at $v$, and $P$ has the Satake parameters $\{b^f:b\in B\}$ at $w$ by property (1)(a) of Definition~\ref{gal-solvable-base-change}, where $N(w)=N(v)^f$.

Finally, a given $Q$ is ramified over $E$ at finitely many places, while $E(\sqrt{-p})$ is ramified at the places over $p$; so distinct primes $p$ give distinct fields $Q$, and there are infinitely many of them.
\end{proof}

\begin{lemma}\label{sha-satake-unramified-places}
Let $K\subset M$ be totally complex number fields with $M$ a Galois extension of $K$ of prime degree, let $n\ge1$, and let $\pi$ and $P$ be cuspidal automorphic representations of $GL_n(\mathbf A_K)$ and $GL_n(\mathbf A_M)$ (Definition~\ref{gal-polarized}) such that $P$ is a weak base change of $\pi$ (Definition~\ref{gal-depth-zero-lift-weak}). Let $v$ be a finite place of $K$. For a finite extension $k$ of $\mathbf Q_p$ and smooth characters $\omega_1,\dots,\omega_n$ of $k^\times$, let $I(\omega_1,\dots,\omega_n)$ be the representation of $GL_n(k)$ obtained by normalized parabolic induction from the character $\operatorname{diag}(t_1,\dots,t_n)\mapsto\prod_i\omega_i(t_i)$ of the Borel subgroup of upper triangular matrices.
\begin{enumerate}
\item Suppose that every place of $M$ over $v$ is unramified over $K$, and that $\pi$ has Satake parameters $B$ at $v$ (Definition~\ref{gal-weakly-attached}). Then $P$ has Satake parameters $\{b^f:b\in B\}$ at every place $w$ of $M$ over $v$, where $N(w)=N(v)^f$.
\item Suppose that $P\circ\gamma\cong P$ for every $\gamma\in\operatorname{Gal}(M/K)$ (Lemma~\ref{gal-depth-zero-lift-stable-twist}), and that $v$ does not split completely in $M$, so that there is one place $w$ of $M$ over $v$. Let $\sigma$ and $\sigma'$ be local components of $\pi$ at $v$ and of $P$ at $w$, and $\omega_1,\dots,\omega_n$ smooth characters of $K_v^\times$. If $\sigma$ is a subquotient of $I(\omega_1,\dots,\omega_n)$, then $\sigma'$ is a subquotient of $I(\omega_1\circ N,\dots,\omega_n\circ N)$, where $N=N_{M_w/K_v}$.
\end{enumerate}
Both follow from the cyclic base change of Arthur and Clozel \cite[Chapter~3]{AC89}. By strong multiplicity one \cite{JS81a}, $P$ is the base change of $\pi$ constructed there, whose local component at every place $w$ is the local base change of that of $\pi$ at the place below. Local base change takes an unramified representation with Satake parameters $B$, at a place of $M$ unramified over $K$ of residue degree $f$, to the unramified representation with Satake parameters $\{b^f:b\in B\}$, and it commutes with parabolic induction, taking a character $\omega$ to $\omega\circ N$ \cite{AC89}.
\end{lemma}
\leanref{Development/AutomorphicArithmetic/Foundations/Convolution/SphericalIdeles/SphConj.lean}{UpperTheory.ConvolutionTheory.Sph.CBCsph}
\leanref{Development/AutomorphicArithmetic/Foundations/Convolution/SphericalIdeles/CbcSph.lean}{UpperTheory.ConvolutionTheory.Sph.cbcSph}
\leanref{Development/AutomorphicArithmetic/Foundations/Convolution/SphericalIdeles/TwoLeAbsNorm.lean}{UpperTheory.ConvolutionTheory.Sph.up}
\leanref{Development/AutomorphicArithmetic/Foundations/GaloisRepresentations/JacobianTraceQuotients/ElemPairJ.lean}{UpperTheory.Gal.Join.LiftGood}
\leanref{Development/AutomorphicArithmetic/Foundations/GaloisRepresentations/JacobianTraceQuotients/TorusCharComp.lean}{UpperTheory.Gal.Join.liftGood}

\begin{lemma}\label{sha-satake-unramified-principal}
Let $k$ be a finite extension of $\mathbf Q_p$ with ring of integers $\mathcal O$, uniformizer $\varpi$ and residue field of $q$ elements, and let $n\ge1$. For $b\in\mathbf C^\times$ let $\xi_b$ be the unramified character of $k^\times$ with $\xi_b(\varpi)=b$, and for $\beta\in(\mathbf C^\times)^n$ write $I(\beta)=I(\xi_{\beta_1},\dots,\xi_{\beta_n})$, as in Lemma~\ref{sha-satake-unramified-places}. An irreducible smooth representation $\sigma$ of $GL_n(k)$ is \emph{spherical} if it has a nonzero vector fixed by $GL_n(\mathcal O)$.
\begin{enumerate}
\item If $\sigma$ is spherical, then it is a subquotient of $I(\beta)$ for some $\beta$.
\item If $\sigma$ is spherical and a subquotient of $I(\beta)$, then $\sigma$ has Satake parameters $\{\beta_1,\dots,\beta_n\}$, normalized as in Definition~\ref{gal-definite-unitary}.
\item If $\beta_i\ne q\beta_j$ for all $i$ and $j$, then every irreducible subquotient of $I(\beta)$ is spherical.
\end{enumerate}
Parts (1) and (2) are the description of spherical representations by the Satake isomorphism \cite{Car79}. In (3), $I(\beta)$ is irreducible, since no two of the characters $\xi_{\beta_i}$ differ by a factor $|\cdot|$ \cite[Theorem~9.7]{Zel80}, and it has a nonzero $GL_n(\mathcal O)$-fixed vector by the Iwasawa decomposition.
\end{lemma}
\leanref{Development/AutomorphicArithmetic/Foundations/Convolution/SphericalIdeles/WeakStable.lean}{UpperTheory.ConvolutionTheory.Sph.SatDict}
\leanref{Development/AutomorphicArithmetic/Foundations/Convolution/SphericalIdeles/SatDict.lean}{UpperTheory.ConvolutionTheory.Sph.satDict}
\leanref{Development/AutomorphicArithmetic/Foundations/Convolution/SphericalIdeles/WeakStable.lean}{UpperTheory.ConvolutionTheory.Sph.UnlSph}
\leanref{Development/AutomorphicArithmetic/Foundations/Convolution/SphericalIdeles/UnlSph.lean}{UpperTheory.ConvolutionTheory.Sph.unlSph}

\begin{lemma}\label{sha-satake-ramified-places}
Let $K\subset M$ be totally complex number fields with $M$ a Galois extension of $K$ of prime degree, let $n\ge1$, and let $\pi$ and $P$ be cuspidal automorphic representations of $GL_n(\mathbf A_K)$ and $GL_n(\mathbf A_M)$ such that $P$ is a weak base change of $\pi$ (Definition~\ref{gal-depth-zero-lift-weak}). Let $w$ be a finite place of $M$ ramified over $K$, and $v$ the place of $K$ under $w$. If $\pi$ has Satake parameters $B$ at $v$ (Definition~\ref{gal-weakly-attached}), then $P$ has Satake parameters $B$ at $w$.
\end{lemma}
\leanref{Development/AutomorphicArithmetic/Foundations/Convolution/SphericalIdeles/RamDefs.lean}{UpperTheory.ConvolutionTheory.Sph.RamSph}
\leanref{Development/AutomorphicArithmetic/Foundations/Convolution/SphericalIdeles/QuadLift.lean}{UpperTheory.ConvolutionTheory.Sph.ramSph}
\leanref{Development/AutomorphicArithmetic/Foundations/Convolution/SphericalIdeles/RamSphOf.lean}{UpperTheory.ConvolutionTheory.Sph.ramSphOf}
\leanref{Development/AutomorphicArithmetic/Foundations/Convolution/SphericalIdeles/Unlinked.lean}{UpperTheory.ConvolutionTheory.Sph.unlinked}
\leanref{Development/AutomorphicArithmetic/Foundations/Convolution/SphericalIdeles/UnrNorm.lean}{UpperTheory.ConvolutionTheory.Sph.unrNorm}
\begin{proof}
Since $[M:K]$ is prime and $w$ is ramified over $K$, the ramification index of $w$ over $K$ is $[M:K]$.
So the residue degree of $w$ is $1$, $N(w)=N(v)=q$, and $w$ is the only place of $M$ over $v$; in particular $v$ does not split completely in $M$.
By the step \emph{stability} of the proof of Lemma~\ref{gal-depth-zero-lift-infinity}, $P\circ\gamma\cong P$ for every $\gamma\in\operatorname{Gal}(M/K)$.
Let $\sigma$ and $\sigma'$ be local components of $\pi$ at $v$ and of $P$ at $w$ \cite{Fla79}.
Write $\xi_b$ and $\xi'_b$ for the unramified characters of $K_v^\times$ and of $M_w^\times$ with value $b$ at a uniformizer, as in Lemma~\ref{sha-satake-unramified-principal}.

\emph{The parameters of $\pi$.}
Since $\pi$ has Satake parameters at $v$, $\sigma$ is spherical.
By Lemma~\ref{sha-satake-unramified-principal}(1), $\sigma$ is a subquotient of $I(\xi_{\beta_1},\dots,\xi_{\beta_n})$ for some $\beta\in(\mathbf C^\times)^n$, and by Lemma~\ref{sha-satake-unramified-principal}(2), $\pi$ has Satake parameters $\{\beta_1,\dots,\beta_n\}$ at $v$.
Satake parameters at a place are unique, so $B=\{\beta_1,\dots,\beta_n\}$.

\emph{Base change.}
By Lemma~\ref{sha-satake-unramified-places}(2), $\sigma'$ is a subquotient of $I(\xi_{\beta_1}\circ N,\dots,\xi_{\beta_n}\circ N)$, with $N=N_{M_w/K_v}$.
The norm satisfies $|N(x)|_v=|x|_w$ for the normalized absolute values, and $N(w)=N(v)$, so $N(x)$ and $x$ have the same valuation, and $\xi_{\beta_i}\circ N=\xi'_{\beta_i}$.

\emph{The parameters are unlinked.}
We have $\beta_i\ne q\beta_j$ for all $i$ and $j$.
Indeed, the absolute value of the central character of $\pi$ is trivial on the ideles of norm $1$, whose quotient by $K^\times$ is compact, so it is a real power of the idelic norm; hence for some real $r$ the twist $\pi\otimes|\det|^r$ has unitary central character.
It has Satake parameters $\{c\beta_1,\dots,c\beta_n\}$ at $v$, with $c=q^{-r}$, so $q^{-1/2}<|c\beta_i|<q^{1/2}$ for every $i$ by Theorem~\ref{lit-jacquet-shalika}(2).
If $\beta_i=q\beta_j$, then $|c\beta_i|=q|c\beta_j|>q^{1/2}$, a contradiction.

\emph{The parameters of $P$.}
By Lemma~\ref{sha-satake-unramified-principal}(3), applied to $M_w$, whose residue field also has $q$ elements, $\sigma'$ is spherical.
By Lemma~\ref{sha-satake-unramified-principal}(2), $P$ has Satake parameters $\{\beta_1,\dots,\beta_n\}=B$ at $w$.
\end{proof}

\begin{proposition}\label{sha-satake-integrality}
Let $E$ be a CM field with maximal totally real subfield $F$, let $n\ge1$, and let $m\in\mathbf Z^n$ be strictly decreasing. Let $\pi$ be a cuspidal automorphic representation of $GL_n(\mathbf A_E)$ such that $Z_\tau(\pi)=\{\frac{n-1}2+m_i:1\le i\le n\}$ for every embedding $\tau:E\to\mathbf C$, and such that $(\pi,\chi)$ is polarized for some Hecke character $\chi$ of $F$ (Definition~\ref{gal-polarized}). Let $w$ be a finite place of $E$, of norm $q$, at which $\pi$ has Satake parameters $B$. Then for every $b\in q^{(n-1)/2}B$ and every $a\in\mathbf Z$:
\begin{enumerate}
\item if $a\ge-m_i$ for all $i$, then $q^ab^{-1}$ is an algebraic integer;
\item if $a\ge m_i$ for all $i$, then $q^ab$ is an algebraic integer.
\end{enumerate}
\end{proposition}
\leanref{Development/AutomorphicArithmetic/AutomorphicVarieties/ModularityArgument/IntegralDeterminants/PlaceConjAlt.lean}{UpperTheory.V.Z.Int.SlopeAll}
\leanref{Development/AutomorphicArithmetic/AutomorphicVarieties/ModularityArgument/IntegralDeterminants/PlaceConjAlt.lean}{UpperTheory.V.Z.Int.SlopeOn}
\leanref{Development/AutomorphicArithmetic/AutomorphicVarieties/AutomorphicForms/IsoOfFrob.lean}{UpperTheory.V.Af.SlopesIn}
\leanref{Development/AutomorphicArithmetic/AutomorphicVarieties/ModularityArgument/IntegralDeterminants/SlopeAll.lean}{UpperTheory.V.Z.Int.slopeAll}
\leanref{Development/AutomorphicArithmetic/AutomorphicVarieties/ModularityArgument/IntegralDeterminants/SlopeAllOf.lean}{UpperTheory.V.Z.Int.slopeAllOf}
\leanref{Development/AutomorphicArithmetic/AutomorphicVarieties/ModularityArgument/IntegralDeterminants/SlopeGood.lean}{UpperTheory.V.Z.Int.slopeGood}
\leanref{Development/AutomorphicArithmetic/AutomorphicVarieties/ModularityArgument/IntegralDeterminants/SlopeCsd.lean}{UpperTheory.V.Z.Int.slopeCsd}
\leanref{Development/AutomorphicArithmetic/AutomorphicVarieties/ModularityArgument/IntegralDeterminants/IntPHolds.lean}{UpperTheory.V.Z.Int.intP_holds}
\begin{proof}
For a CM field $E'$ write $E'^+$ for its maximal totally real subfield, $\eta_{E'}$ for the quadratic character of $E'/E'^+$, and $Z(\pi')$ for the family $(Z_\tau(\pi'))_\tau$ of a cuspidal automorphic representation $\pi'$ of $GL_n(\mathbf A_{E'})$.

\emph{Step 1: from slopes to integrality.}
Let $E'$ be a number field, $w'$ a finite place of $E'$ of norm $q$, and suppose that for every prime $\ell$ and every $\iota:\bar{\mathbf Q}_\ell\cong\mathbf C$, $B$ has $\iota$-slopes bounded at $w'$ by the family all of whose members are $\{\frac{n-1}2+m_i:1\le i\le n\}$ (Definition~\ref{sha-satake-slopes}).
Then (1) and (2) hold.
Indeed, let $b\in q^{(n-1)/2}B$.
If $a\ge-m_i$ for all $i$, take $l_\tau=-a$ and $h_\tau=\max_im_i$ for all $\tau$; the second inequality of Definition~\ref{sha-satake-slopes} and Lemma~\ref{sha-satake-norm-product} give $|q^ab^{-1}|_\iota\le1$ for every $\ell$ and $\iota$, and $q^ab^{-1}$ is an algebraic integer by Lemma~\ref{sha-satake-integer-criterion}.
If $a\ge m_i$ for all $i$, take $l_\tau=\min_im_i$ and $h_\tau=a$; the first inequality gives $|q^ab|_\iota\le1$ in the same way.

\emph{Step 2: self-dual scope at split places.}
Let $E'$ be a CM field, $\ell_0$ a prime, $\pi'$ a cuspidal automorphic representation of $GL_n(\mathbf A_{E'})$ with $(\pi',\eta_{E'}^n)$ in self-dual scope for $\ell_0$, and $w'$ a finite place of $E'$ split over $E'^+$ at which $\pi'$ has Satake parameters $B'$.
Then $B'$ has $\iota$-slopes bounded by $Z(\pi')$ at $w'$, for every prime $\ell$ and every $\iota$.
Indeed, Lemma~\ref{sha-satake-descent} gives an automorphic representation $U$ of $U_n$.
It has a weight $\lambda_\tau$ at every embedding $\tau$, by the theorem of the highest weight \cite{FH91}, since $U$ is irreducible and the $\tau$-factor $U(n)$ is compact; and $Z_\tau(\pi')=\{\lambda_{\tau,k}+\frac{n-1}2-k:0\le k<n\}$ by Lemma~\ref{sha-satake-descent}(3).
The representation $\pi'$ has a local component $\sigma$ at $w'$, spherical with Satake parameters $B'$, and $\sigma$ is a local component of $U$ at $w'$ by Lemma~\ref{sha-satake-descent}(2).
Lemma~\ref{sha-satake-unitary-slopes} gives the claim.

\emph{Step 3: good scope at split places.}
Let $E'$, $\ell_0$, $w'$ and $B'$ be as in Step 2, and let $\pi'$ be such that $(\pi',\chi')$ is in good scope for $\ell_0$ for some Hecke character $\chi'$ of $E'^+$.
Then the conclusion of Step 2 holds.
Indeed, Lemma~\ref{sha-satake-twist} gives a Hecke character $\kappa$ of $E'$, algebraic of some infinity type $m'$ and unramified at $w'$, with $(\pi'\otimes\kappa,\eta_{E'}^n)$ in self-dual scope for $\ell_0$.
Since $\kappa$ is unramified at $w'$, $\pi'\otimes\kappa$ has the Satake parameters $\{\beta\kappa_{w'}(\varpi_{w'}):\beta\in B'\}$ at $w'$, and $Z_\tau(\pi'\otimes\kappa)=\{z+m'_\tau:z\in Z_\tau(\pi')\}$.
If $(l_\tau)$ and $(h_\tau)$ are as in Definition~\ref{sha-satake-slopes} for $Z(\pi')$, then $(l_\tau+m'_\tau)$ and $(h_\tau+m'_\tau)$ are as there for $Z(\pi'\otimes\kappa)$.
For $b\in q^{(n-1)/2}B'$, Step 2 for $\pi'\otimes\kappa$, applied to $b\kappa_{w'}(\varpi_{w'})$, gives
\[
 |b|_\iota\,|\kappa_{w'}(\varpi_{w'})|_\iota\prod_{\tau\in\mathcal E_\iota(w')}|\tau(\varpi_{w'})|_\iota^{h_\tau+m'_\tau}\le1,
\]
and $|\kappa_{w'}(\varpi_{w'})|_\iota\prod_{\tau\in\mathcal E_\iota(w')}|\tau(\varpi_{w'})|_\iota^{m'_\tau}=1$ by Lemma~\ref{sha-satake-character-value}; so $|b|_\iota\prod_\tau|\tau(\varpi_{w'})|_\iota^{h_\tau}\le1$.
The inequality for $b^{-1}$ follows in the same way.

\emph{Step 4: the general case.}
Since $m$ is strictly decreasing, $\pi$ is regular algebraic, so $(\pi,\chi)$ is in scope (Definition~\ref{gal-scope}).
Let $\Omega$ be an algebraic closure of $E$, and $Q\subset\Omega$ a quadratic extension of $E$ given by Lemma~\ref{sha-satake-quadratic-extension} for $\ell=2$ and $w_0=w$.
By its property (3), $Q$ is ramified over $E$ at a place $u$ over a place $v$ of $E$ at which $\pi$ has Satake parameters, and then $\pi$ has a base change $P$ to $Q$ (Definition~\ref{gal-solvable-base-change}).
Indeed, let $A$ be a valuation ring of $\Omega$ over $u$; as $u$ is ramified over $E$, some element of the inertia group of $A$ in $\operatorname{Gal}(\Omega/E)$ does not fix $Q$.
So the step \emph{no self-twist} of the proof of Lemma~\ref{gal-depth-zero-lift-step}, for $K=E$ and $M=Q$, shows that every norm character $\eta$ for $Q$ (Definition~\ref{gal-depth-zero-lift-weak}) is ramified at $v$ and that $\pi\otimes(\eta\circ\det)\not\cong\pi$: an isomorphism would carry a nonzero $GL_n(\mathcal O_v)$-fixed vector of $\pi$ to a nonzero vector on which $GL_n(\mathcal O_v)$ acts through $\eta_v\circ\det$, and the averaging projection there, with the Cartan decomposition and the irreducibility of $\pi$, shows that $\pi$ has no such vector.
So Theorem~\ref{gal-depth-zero-lift-base-change}(1) gives a weak base change $P$ of $\pi$ with $Z_\tau(P)=Z_{\tau|_E}(\pi)$ for every $\tau$, whose local components agree with those of $\pi$ at places of degree one; these are properties (1)(b) and (1)(c) of Definition~\ref{gal-solvable-base-change}.
For property (1)(a), let $u'$ be a finite place of $Q$ and $v'$ the place of $E$ under it, with $N(u')=N(v')^f$, and suppose that $\pi$ has Satake parameters $B'$ at $v'$.
If $u'$ is ramified over $E$, its ramification index is $2$, so $f=1$, and $P$ has Satake parameters $B'$ at $u'$ by Lemma~\ref{sha-satake-ramified-places}.
Otherwise every place of $Q$ over $v'$ is unramified over $E$, since a place of $Q$ ramified over $E$ is the only place of $Q$ over the place below it, as $[Q:E]=2$; and $P$ has Satake parameters $\{b^f:b\in B'\}$ at $u'$ by Lemma~\ref{sha-satake-unramified-places}(1).
% The base change is QuadLift, supplied by intK_of_desc as UpperTheory.ConvolutionTheory.Sph.quadLift = quadLiftOf ramSph. quadLiftUnr applies liftUnr (from cbcc) with T.cuspidal_twist_ramified_char_ne_self = TwistRamified.twist_ne (through no_eigvec) and the third part of Sph.normRam (UnramifiedFromSplittingCriterion.normRam, with the inertia element from InertFixVal.ramificationIdx_eq_one) for the ramification of each norm character; liftUnr gives P with the infinity types and the local components, and the Satake parameters at unramified places by Sph.cbcSph (Lemma sha-satake-unramified-places(1)) after the second part of normRam. quadLiftOf gets f = 1 at a ramified place from the first part of normRam, and the Satake parameters there from RamSph (ramSphOf with Gal.Join.locGal and liftGood, Lemma sha-satake-ramified-places).
By Lemma~\ref{gal-solvable-base-change-facts}(1), $(P,\chi\circ N_{Q^+/F})$ is in scope, and by property (4) of Lemma~\ref{sha-satake-quadratic-extension} it is in good scope for $2$.
Let $A$ be a valuation ring of $\Omega$ over $w$ and $w'$ the place of $Q$ under $A$.
Since $D_A\subset G_Q$ by property (1) of Lemma~\ref{sha-satake-quadratic-extension}, there is a continuous isomorphism $E_w\to Q_{w'}$ that extends the inclusion of $E$ in $Q$; so $w'$ has norm $q$, and $P$ has Satake parameters $B$ at $w'$ by property (1)(a) of Definition~\ref{gal-solvable-base-change}, with $f=1$.
The place $w'$ is split over $Q^+$ by property (2) of Lemma~\ref{sha-satake-quadratic-extension}, and $Z_\tau(P)=Z_{\tau|_E}(\pi)=\{\frac{n-1}2+m_i:1\le i\le n\}$ for every $\tau:Q\to\mathbf C$ by property (1)(b) of Definition~\ref{gal-solvable-base-change}.
By Step 3 for $P$ at $w'$, $B$ has $\iota$-slopes bounded by this family at $w'$ for every $\ell$ and $\iota$, and Step 1 for $Q$ and $w'$ gives (1) and (2).
\end{proof}

\begin{remark}\label{sha-satake-integrality-base-change}
The base change $P$ in Step 4 of the proof of Proposition~\ref{sha-satake-integrality} also comes directly from Theorem~\ref{gal-cyclic-base-change}, as in the proof of Lemma~\ref{gal-patch-descent}: a Hecke character $\eta$ as in its hypothesis is the quadratic character of $Q/E$, ramified at $v$, so $\pi\otimes(\eta\circ\det)$ is ramified at $v$, while it has the Satake parameters of $\pi$ at almost all places; strong multiplicity one \cite{JS81a} would give $\pi\cong\pi\otimes(\eta\circ\det)$, a contradiction. That theorem gives the Satake parameters at all places at once.
\end{remark}

\begin{definition}\label{sha-eigensystem-distinctness}
Let $F$ be a number field with algebraic closure $\bar F$, let $\ell$ be a prime, let $n,M\ge1$ and let $m=(m_1,\dots,m_n)\in\mathbf Z^n$.
Let $r:\operatorname{Gal}(\bar F/F)\to GL_n(\bar{\mathbf Q}_\ell)$ be a continuous representation, and let $w$ be a prime of $F$ above $\ell$, with a decomposition group $D_w$.
\begin{enumerate}
\item The representation $r$ is \emph{$\ell$-distinct at $w$} if its restriction to $D_w$ is conjugate to an upper triangular representation whose diagonal characters $\kappa_1,\dots,\kappa_n$ satisfy $\kappa_i\ne\kappa_j\chi_{\mathrm{cyc}}$ on $D_w$ for all $i,j$. So $r$ is $\ell$-distinct in the sense of Definition~\ref{sha-eigensystem-level} exactly when it is $\ell$-distinct at every prime of $F$ above $\ell$.
\item The representation $r$ is \emph{potentially $\ell$-distinct at $w$} if its restriction to $D_w$ is conjugate to an upper triangular representation whose diagonal characters $\kappa_1,\dots,\kappa_n$ satisfy $\kappa_i^a\ne(\kappa_j\chi_{\mathrm{cyc}})^a$ on $D_w$ for all $i,j$ and all integers $a\ge1$. It is \emph{potentially $\ell$-distinct} if it is potentially $\ell$-distinct at every prime of $F$ above $\ell$. Taking $a=1$ shows that a representation that is potentially $\ell$-distinct at $w$ is $\ell$-distinct at $w$.
\item The representation $r$ satisfies the \emph{distinctness condition at $\ell$} for $M$ and $m$ if either $\ell$ divides $M$ and $r$ is $\ell$-distinct, or $m_i\ne m_j+1$ for all $i,j$.
\item The representation $r$ satisfies the \emph{potential distinctness condition at $\ell$} for $M$ and $m$ if either $\ell$ divides $M$ and $r$ is potentially $\ell$-distinct, or $m_i\ne m_j+1$ for all $i,j$.
\end{enumerate}
\end{definition}
\leanref{Development/AutomorphicArithmetic/AutomorphicVarieties/ModularityArgument/AutomorphyLifting/WeakOfAtt.lean}{UpperTheory.V.Z.Win.Unl1}
\leanref{Development/AutomorphicArithmetic/AutomorphicVarieties/ModularityArgument/AutomorphyLifting/WeakOfAtt.lean}{UpperTheory.V.Z.Win.PotUnl}
\leanref{Development/AutomorphicArithmetic/AutomorphicVarieties/ModularityArgument/AutomorphyLifting/SplitIrr.lean}{UpperTheory.V.Z.Win.UnlAll}
\leanref{Development/AutomorphicArithmetic/AutomorphicVarieties/ModularityArgument/AutomorphyLifting/SplitIrr.lean}{UpperTheory.V.Z.Win.PotAll}
\leanref{Development/AutomorphicArithmetic/AutomorphicVarieties/ModularityArgument/AutomorphyLifting/AdmissibleLevelRestriction/UnlAllOf.lean}{UpperTheory.V.Z.Win.A4.SafeRA}
\leanref{Development/AutomorphicArithmetic/AutomorphicVarieties/ModularityArgument/AutomorphyLifting/AdmissibleLevelRestriction/UnlAllOf.lean}{UpperTheory.V.Z.Win.A4.SafePA}

\begin{lemma}\label{sha-eigensystem-unitary}
Let $K$ be a CM field with complex conjugation $c$, let $n\ge2$, and assume the Galois representations as in Theorem~\ref{sha-residually-reducible-automorphy}.
There are $c_0,A$ such that the following holds.
Let $m\in\mathbf Z^n$ be strictly decreasing, let $H\ge0$ with $|m_i|\le H$ for all $i$, let $M\ge1$, and let $w$ be a prime of $K$ above a prime $p\nmid Md_K$, of norm $q$.
\begin{enumerate}
\item Suppose that $m_i+m_{n+1-i}=1-n$ for all $i$. Then there is an eigenvalue datum of rank $n$ at $q$ with weights $m$ and weight $1-n$ (Definition~\ref{sha-eigensystem-data}), with at most $c_0(M(H+1))^A$ elements, that classifies at $w$, for every odd prime $\ell\ne p$ and every $\iota:\bar{\mathbf Q}_\ell\cong\mathbf C$, the set of irreducible continuous representations $r:\operatorname{Gal}(\bar K/K)\to GL_n(\bar{\mathbf Q}_\ell)$ that have level $M$ and weights $m$, are $\ell$-distinct (Definition~\ref{sha-eigensystem-level}), are attached through $\iota$ to a regular algebraic cuspidal automorphic representation $\pi$ of $GL_n(\mathbf A_K)$ with $\pi^\vee\cong\pi^c$, and are windowed in the sense of Theorem~\ref{sha-residually-reducible-automorphy}.
\item Suppose that $K$ is imaginary quadratic, and that $u\in\mathbf Z$ satisfies $m_i+m_{n+1-i}=u$ for all $i$. Then the same holds with weight $u$ for the set of irreducible continuous representations $r$ that satisfy the conditions of level $M$ and weights $m$ except possibly the last condition at the primes above $\ell$, satisfy the potential distinctness condition at $\ell$ for $M$ and $m$ (Definition~\ref{sha-eigensystem-distinctness}), are attached through $\iota$ to a regular algebraic cuspidal automorphic representation $\pi$ of $GL_n(\mathbf A_K)$ with $(\pi,\chi)$ polarized for some Hecke character $\chi$, and are windowed.
\end{enumerate}
\end{lemma}
\begin{proof}
We prove (1); the proof of (2) is the same, with polarized representations in place of conjugate self-dual ones and $u$ in place of $1-n$.
Let $\Sigma$ be the set of regular algebraic cuspidal automorphic representations $\pi$ of $GL_n(\mathbf A_K)$ with $\pi^\vee\cong\pi^c$, with infinitesimal character the multiset of the $\frac{n-1}2+m_i$ at every embedding, and with nonzero vectors fixed by $K(M)$, the group of Lemma~\ref{sha-eigensystem-cuspidal-count}.
For a prime $v\nmid M$ of $K$, the local component of $K(M)$ at $v$ is $GL_n(\mathcal O_v)$, so $\pi_v$ is unramified, with Satake parameters $\alpha_{v,1},\dots,\alpha_{v,n}$; let $B_v(\pi)$ be the multiset of the $q_v^{(n-1)/2}\alpha_{v,i}$, where $q_v$ is the norm of $v$.
Isomorphic representations have the same $B_v$ for all $v$.
Let $S$ be the set of the families $(B_v(\pi))_{v\nmid M}$ for $\pi\in\Sigma$, and for $s=(B_v(\pi))_v$ let $\beta_s$ be the multiset of the inverses of the elements of $B_w(\pi)$.

\emph{Size.} By Lemma~\ref{sha-eigensystem-cuspidal-count}, $\Sigma$ has at most $c(M(H+1))^A$ isomorphism classes, so $|S|\le c(M(H+1))^A$.

\emph{Stability.} Let $\sigma\in\operatorname{Aut}(\mathbf C)$ and $\pi\in\Sigma$. By Clozel's theorem \cite[Th\'eor\`eme~3.13]{Clo90} there is a cuspidal automorphic representation $\pi'$ whose finite part is the conjugate of that of $\pi$ by $\sigma$ and whose infinitesimal characters are those of $\pi$ with the embeddings composed with $\sigma$; so $\pi'$ has the infinitesimal character and the fixed vectors required by $\Sigma$, and $B_v(\pi')=\sigma(B_v(\pi))$ for every $v\nmid M$.
The conjugate dual of $\pi'$ has the same Satake parameters as $\pi'$ at almost all places, since this holds for $\pi$, so $\pi'^\vee\cong\pi'^c$ by strong multiplicity one (Theorem~\ref{lit-jacquet-shalika}(1)), and $\pi'\in\Sigma$.

\emph{Window.} The sum of the entries of the infinitesimal character of $\pi$ is $\sum_i\big(\frac{n-1}2+m_i\big)=\frac n2(n-1+u)$ with $u=1-n$, so the central character of $\pi\otimes|\det|^{-(n-1+u)/2}$ is unitary. By Theorem~\ref{lit-jacquet-shalika}(2) applied to this twist, every element $b$ of $B_w(\pi)$ satisfies $q^{(-u-1)/2}<|b|<q^{(-u+1)/2}$, and the product of the elements of $B_w(\pi)$ has absolute value $q^{-nu/2}$; so the inverses satisfy the window condition of Definition~\ref{sha-eigensystem-data} for the weight $u$.

\emph{Integrality.} Let $\eta$ be the quadratic character of $K/K^+$, where $K^+$ is the maximal totally real subfield of $K$.
Then $(\pi,\eta^n)$ is polarized (Definition~\ref{gal-polarized}): $\eta_v(-1)=-1$ at every real place $v$ of $K^+$, and $\eta\circ N_{K/K^+}=1$, so the condition on $\pi^c$ is $\pi^\vee\cong\pi^c$.
In (2), $(\pi,\chi)$ is polarized by hypothesis.
So Proposition~\ref{sha-satake-integrality} applies to $\pi$ at $w$: for $b\in B_w(\pi)$ and $a\in\mathbf Z$, $q^ab^{-1}$ is an algebraic integer if $a\ge-m_i$ for all $i$, and $q^ab$ is one if $a\ge m_i$ for all $i$.
For $\beta_s$, which consists of the inverses, this is the integrality condition of Definition~\ref{sha-eigensystem-data}.
% Integrality is Z.Int.intK_holds (and Z.Int.intP_holds in (2)), that is, Z.Int.slopeAll applied with chi = quadChar K ^ n.

\emph{Classification.} Let $\ell\ne p$ be odd, $\iota$ an isomorphism, and $r$ a representation in the set of (1), attached to $\pi$. By local-global compatibility \cite{Car12,Car14,BLGGT14}, the conditions of level $M$ and weights $m$ on $r$ imply that $\pi$ has nonzero vectors fixed by $K(M)$, so $\pi\in\Sigma$; and for every prime $v\nmid\ell M$, $r$ is unramified at $v$ and the eigenvalues of $\iota r(\sigma_v^{-1})$ are the elements of $B_v(\pi)$.
At a prime above $\ell$ this uses that $r$ is $\ell$-distinct or potentially $\ell$-distinct there (Definition~\ref{sha-eigensystem-distinctness}).
In (2) the potential distinctness condition at $\ell$ gives this. Under its first alternative $r$ is potentially $\ell$-distinct. Under the second, at a prime above $\ell$ the characters $\kappa_i$ of Definition~\ref{sha-eigensystem-level} satisfy $\kappa_i\ne\kappa_j\chi_{\mathrm{cyc}}$: on inertia $\kappa_i(\kappa_j\chi_{\mathrm{cyc}})^{-1}=\chi_{\mathrm{cyc}}^{m_i-m_j-1}\chi_i\chi_j^{-1}$, where $m_i-m_j-1\ne0$, $\chi_i\chi_j^{-1}$ has finite order and $\chi_{\mathrm{cyc}}$ has infinite order on inertia.
% In (2) this is Keep.levelPA_alt: the first alternative of SafePA gives PotUnl, the second gives Unl1 by Lev.unlGap and A4.unlOfOne; Keep.levelCore_alt uses Unl1 or PotUnl at each place above l.
Let $s(r)=(B_v(\pi))_v$. Since $w\nmid\ell M$, the eigenvalues of $\iota r(\sigma_w)$ are the elements of $\beta_{s(r)}$.
If $s(r)=s(r')$, then $r(\sigma_v)$ and $r'(\sigma_v)$ have the same characteristic polynomial for every $v\nmid\ell M$; by the Chebotarev density theorem (Theorem~\ref{lit-chebotarev-degree-one}) the characteristic polynomials of $r$ and $r'$ agree everywhere, and since both are irreducible, $r\cong r'$ by the Brauer-Nesbitt theorem \cite[(30.16)]{CR62}.
\end{proof}
\leanref{Development/AutomorphicArithmetic/AutomorphicVarieties/ModularityArgument/ThresholdArguments/DenseThresholdCases/FrobeniusTraceCountBoundOfCmEigensystemCount.lean}{UpperTheory.GalN.V.Win.N.CountKAlt}
\leanref{Development/AutomorphicArithmetic/AutomorphicVarieties/ModularityArgument/ThresholdArguments/DenseThresholdCases/FrobeniusTraceCountBoundOfCmEigensystemCount.lean}{UpperTheory.GalN.V.Win.N.CountPAlt}
\leanref{Development/AutomorphicArithmetic/AutomorphicVarieties/ModularityArgument/ThresholdArguments/DenseThresholdCases/FrobeniusTraceCountBoundOfCmEigensystemCount.lean}{UpperTheory.GalN.V.Win.N.countK_alt}
\leanref{Development/AutomorphicArithmetic/AutomorphicVarieties/ModularityArgument/ThresholdArguments/DenseThresholdCases/FrobeniusTraceCountBoundOfCmEigensystemCount.lean}{UpperTheory.GalN.V.Win.N.countP_alt}
\leanref{Development/AutomorphicArithmetic/Foundations/GaloisRepresentations/AdmissibleResidues/WinDefs.lean}{UpperTheory.GalN.V.Win.exists_eigData}
\leanref{Development/AutomorphicArithmetic/Foundations/GaloisRepresentations/AdmissibleResidues/WinSat.lean}{UpperTheory.GalN.V.win_algSat}
\leanref{Development/AutomorphicArithmetic/AutomorphicVarieties/ModularityArgument/ThresholdArguments/DenseThresholdCases/TConjK.lean}{UpperTheory.V.Af.conjK_alt}
\leanref{Development/AutomorphicArithmetic/AutomorphicVarieties/ModularityArgument/IntegralDeterminants/IntPHolds.lean}{UpperTheory.V.Z.Int.intK_holds}
\leanref{Development/AutomorphicArithmetic/AutomorphicVarieties/ModularityArgument/ThresholdArguments/DenseThresholdCases/KeepLevel.lean}{UpperTheory.V.Keep.levelKA_alt}
\leanref{Development/AutomorphicArithmetic/AutomorphicVarieties/ModularityArgument/ThresholdArguments/DenseThresholdCases/KeepLevel.lean}{UpperTheory.V.Keep.levelPA_alt}
\leanref{Development/AutomorphicArithmetic/AutomorphicVarieties/ModularityArgument/SemisimpleLevels/UnlGap.lean}{UpperTheory.V.Lev.unlGap}
\leanref{Development/AutomorphicArithmetic/AutomorphicVarieties/ModularityArgument/ThresholdArguments/DenseThresholdCases/TIsoOfFrob.lean}{UpperTheory.V.Af.iso_of_frobenius_of_printed_results}

\begin{lemma}\label{sha-eigensystem-twists}
Let $\Gamma$ be a group, let $k$ be an algebraically closed field of characteristic $0$, let $n\ge1$, and let $r:\Gamma\to GL_n(k)$ be an irreducible representation.
\begin{enumerate}
\item At most $n^2$ characters $\delta:\Gamma\to k^\times$ satisfy $r\otimes\delta\cong r$.
\item Let $\delta:\Gamma\to\{\pm1\}$ be a character with kernel $\Gamma'$. If $r\otimes\delta\not\cong r$, then the restriction of $r$ to $\Gamma'$ is irreducible. If the restriction of $r$ to $\Gamma'$ is irreducible, then every representation $r':\Gamma\to GL_n(k)$ whose restriction to $\Gamma'$ is isomorphic to that of $r$ is isomorphic to $r$ or to $r\otimes\delta$.
\end{enumerate}
\end{lemma}
\begin{proof}
(1) For each such $\delta$ choose $g_\delta\in GL_n(k)$ with $g_\delta r(s)=\delta(s)r(s)g_\delta$ for all $s\in\Gamma$.
Suppose that $\sum_\delta c_\delta g_\delta=0$ with $c_\delta\in k$.
Multiplying on the right by $r(s)$ gives $r(s)\sum_\delta c_\delta\delta(s)g_\delta=0$, so $\sum_\delta c_\delta\delta(s)g_\delta=0$ for every $s\in\Gamma$.
Entry by entry, this is a linear relation between the distinct characters $\delta$, so $c_\delta g_\delta=0$ by the linear independence of characters, and $c_\delta=0$.
So the $g_\delta$ are linearly independent in $M_n(k)$, which has dimension $n^2$.

(2) If $\delta$ is trivial, both statements are clear; otherwise let $t\in\Gamma$ with $\delta(t)=-1$.
Suppose that $W\subset k^n$ is a nonzero proper subspace stable under $r(\Gamma')$.
Then $r(s)$ preserves $W$ and $r(t)W$ for $s\in\Gamma'$ and exchanges them for $s\notin\Gamma'$, so $W\cap r(t)W$ and $W+r(t)W$ are stable under $r(\Gamma)$, and $k^n=W\oplus r(t)W$ since $r$ is irreducible.
The map $P$ that is $1$ on $W$ and $-1$ on $r(t)W$ satisfies $Pr(s)=\delta(s)r(s)P$ for all $s$, so $r\otimes\delta\cong r$.
For the second statement, let $g\in GL_n(k)$ with $r'(s)=gr(s)g^{-1}$ for $s\in\Gamma'$, and put $u=r(t)^{-1}g^{-1}r'(t)g$.
Then $u$ commutes with $r(\Gamma')$, since $tst^{-1}\in\Gamma'$ for $s\in\Gamma'$, so $u$ is a scalar by Schur's lemma; and $u^2=1$ because $t^2\in\Gamma'$.
Hence $r'(s)=gu^{(1-\delta(s))/2}r(s)g^{-1}$ for all $s\in\Gamma$, so $r'\cong r$ if $u=1$ and $r'\cong r\otimes\delta$ if $u=-1$.
\end{proof}
\leanref{Development/AutomorphicArithmetic/AutomorphicVarieties/GaloisLibrary/CrystallineBaseChange/CqStmt.lean}{UpperTheory.V.Br.TwCard}
\leanref{Development/AutomorphicArithmetic/AutomorphicVarieties/GaloisLibrary/CrystallineBaseChange/TwCard.lean}{UpperTheory.V.Br.twCard}
\leanref{Development/AutomorphicArithmetic/AutomorphicVarieties/GaloisLibrary/CrystallineBaseChange/CqStmt.lean}{UpperTheory.V.Br.ResTwo}
\leanref{Development/AutomorphicArithmetic/AutomorphicVarieties/GaloisLibrary/CrystallineBaseChange/ResTwo.lean}{UpperTheory.V.Br.resTwo}

\begin{lemma}\label{sha-eigensystem-split-field}
Let $\ell$ be a prime, let $n\ge1$, and let $r:\operatorname{Gal}(\bar{\mathbf Q}/\mathbf Q)\to GL_n(\bar{\mathbf Q}_\ell)$ be an irreducible representation.
Then there is an imaginary quadratic field $K\subset\bar{\mathbf Q}$ such that the restriction of $r$ to $\operatorname{Gal}(\bar{\mathbf Q}/K)$ is irreducible and the decomposition group in $\operatorname{Gal}(\bar{\mathbf Q}/\mathbf Q)$ of every place of $\bar{\mathbf Q}$ above $\ell$ is contained in $\operatorname{Gal}(\bar{\mathbf Q}/K)$.
\end{lemma}
\begin{proof}
By Dirichlet's theorem there are distinct primes $P_1,\dots,P_{n^2+1}$ with $P_i\equiv-1\pmod{4\ell}$.
Let $K_i=\mathbf Q(\sqrt{-P_i})\subset\bar{\mathbf Q}$, an imaginary quadratic field, and let $\delta_i$ be its quadratic character, with kernel $\operatorname{Gal}(\bar{\mathbf Q}/K_i)$.
The $\delta_i$ are pairwise distinct, since a quadratic field does not contain square roots of both $-P_i$ and $-P_j$ for $i\ne j$.
By Lemma~\ref{sha-eigensystem-twists}(1) there is an $i$ with $r\otimes\delta_i\not\cong r$, and by Lemma~\ref{sha-eigensystem-twists}(2) the restriction of $r$ to $\operatorname{Gal}(\bar{\mathbf Q}/K_i)$ is irreducible.
Let $A$ be the valuation ring of a place of $\bar{\mathbf Q}$ above $\ell$, let $g$ be in its decomposition group, let $y\in K_i$ with $y^2=-P_i$, and write $P_i+1=4\ell c$.
Suppose that $g(y)=-y$, and put $x=(1+y)/2$; then $x(1-x)=\ell c$ and $g(x)=1-x$.
If $x^{-1}\in A$, then $(1-x)^{-1}=g(x^{-1})\in A$, and $\ell^{-1}=c\,x^{-1}(1-x)^{-1}\in A$, which is false.
Otherwise neither $x^{-1}$ nor $(1-x)^{-1}=g(x^{-1})$ lies in $A$, so $x$ and $1-x$ lie in the maximal ideal of $A$, and so does their sum $1$, which is false.
So $g$ fixes $y$, hence $K_i$, and $K=K_i$ has the two properties.
\end{proof}
\leanref{Development/AutomorphicArithmetic/AutomorphicVarieties/ModularityArgument/AutomorphyLifting/SplitIrr.lean}{UpperTheory.V.Z.Win.SplitIrr}
\leanref{Development/AutomorphicArithmetic/AutomorphicVarieties/ModularityArgument/AutomorphyLifting/AdmissibleLevelRestriction/SplitIrr.lean}{UpperTheory.V.Z.Win.A4.splitIrr}

\begin{lemma}\label{sha-eigensystem-flags}
Let $F\subset L$ be number fields inside an algebraic closure $\bar F$ of $F$, let $\ell$ be a prime, let $n\ge1$, and let $r:\operatorname{Gal}(\bar F/F)\to GL_n(\bar{\mathbf Q}_\ell)$ be a continuous representation, with restriction $r_L$ to $\operatorname{Gal}(\bar F/L)$.
Let $w$ be a prime of $L$ above $\ell$, and $v$ the prime of $F$ below $w$.
\begin{enumerate}
\item If $r$ is potentially $\ell$-distinct at $v$ (Definition~\ref{sha-eigensystem-distinctness}), then $r_L$ is potentially $\ell$-distinct at $w$.
\item Suppose that the restriction of $r$ to a decomposition group at $v$ is conjugate to an upper triangular representation. If $r_L$ is potentially $\ell$-distinct at $w$, then $r$ is potentially $\ell$-distinct at $v$.
\end{enumerate}
\end{lemma}
\begin{proof}
Put $\Gamma=\operatorname{Gal}(\bar F/F)$ and $\Gamma_L=\operatorname{Gal}(\bar F/L)$, and for a place $A$ of $\bar F$ let $D_A\subset\Gamma$ be its decomposition group; if $A$ lies above $w$, its decomposition group in $\Gamma_L$ is $D_A\cap\Gamma_L$.
The cyclotomic character of $L$ is the restriction of that of $F$.

(1) Let $A$ be a place above $w$. Since $r$ is potentially $\ell$-distinct at $v$, it has on $D_A$ an upper triangular form whose diagonal characters $\kappa_i$ satisfy $\kappa_i^a\ne(\kappa_j\chi_{\mathrm{cyc}})^a$ for all $i,j$ and $a\ge1$; restricted to $D_A\cap\Gamma_L$, it is an upper triangular form of $r_L$ whose diagonal characters are the restrictions of the $\kappa_i$.
Let $N\ge1$ be the index of a normal subgroup of finite index of $\Gamma$ contained in $\Gamma_L$, so that $s^N\in\Gamma_L$ for every $s\in\Gamma$.
If $\kappa_i^a=(\kappa_j\chi_{\mathrm{cyc}})^a$ on $D_A\cap\Gamma_L$ for some $i,j$ and $a\ge1$, then evaluating at $s^N$ for $s\in D_A$ gives $\kappa_i^{Na}=(\kappa_j\chi_{\mathrm{cyc}})^{Na}$ on $D_A$, which is false.

(2) Let $A$ be a place above $v$ such that $r$ has on $D_A$ an upper triangular form, with diagonal characters $\kappa_1,\dots,\kappa_n$.
Let $B$ be a place above $w$ such that $r_L$ has on $D_B\cap\Gamma_L$ an upper triangular form whose diagonal characters $\theta_1,\dots,\theta_n$ satisfy $\theta_a^m\ne(\theta_b\chi_{\mathrm{cyc}})^m$ for all $a,b$ and $m\ge1$, and choose $\gamma\in\Gamma$ with $\gamma A=B$.
Then $s\mapsto\gamma^{-1}s\gamma$ maps $D_B\cap\Gamma_L$ into $D_A$, and $r(s)=r(\gamma)r(\gamma^{-1}s\gamma)r(\gamma)^{-1}$, so $r_L$ has on $D_B\cap\Gamma_L$ a second upper triangular form, with diagonal characters $\kappa'_i(s)=\kappa_i(\gamma^{-1}s\gamma)$.
Both forms give the trace of $r_L$, so $\sum_i\kappa'_i=\sum_a\theta_a$, and by the linear independence of characters each $\kappa'_i$ is one of the $\theta_a$.
If $\kappa_i^m=(\kappa_j\chi_{\mathrm{cyc}})^m$ on $D_A$ for some $i,j$ and $m\ge1$, then composing with $s\mapsto\gamma^{-1}s\gamma$, which preserves $\chi_{\mathrm{cyc}}$, gives $\theta_a^m=(\theta_b\chi_{\mathrm{cyc}})^m$ for the $a$ and $b$ with $\kappa'_i=\theta_a$ and $\kappa'_j=\theta_b$, which is false.
\end{proof}
\leanref{Development/AutomorphicArithmetic/AutomorphicVarieties/ModularityArgument/AutomorphyLifting/WeakOfAtt.lean}{UpperTheory.V.Z.Win.FlagLaws}
\leanref{Development/AutomorphicArithmetic/AutomorphicVarieties/ModularityArgument/AutomorphyLifting/PolarizedLinkage/FlagLaws.lean}{UpperTheory.V.Z.Win.flagLaws}

\begin{lemma}\label{sha-eigensystem-potential}
Let $\ell$ be a prime, let $\iota:\bar{\mathbf Q}_\ell\cong\mathbf C$, let $n,M\ge1$, let $m\in\mathbf Z^n$, and let $r:\operatorname{Gal}(\bar{\mathbf Q}/\mathbf Q)\to GL_n(\bar{\mathbf Q}_\ell)$ be an irreducible continuous representation whose restriction to every imaginary quadratic field on which it is irreducible is attached through $\iota$ to a regular algebraic cuspidal automorphic representation and is windowed (Theorem~\ref{sha-residually-reducible-automorphy}).
\begin{enumerate}
\item If $r$ is $\ell$-distinct, then $r$ is potentially $\ell$-distinct (Definition~\ref{sha-eigensystem-distinctness}).
\item If $r$ satisfies the distinctness condition at $\ell$ for $M$ and $m$, then for every number field $L\subset\bar{\mathbf Q}$ the restriction of $r$ to $\operatorname{Gal}(\bar{\mathbf Q}/L)$ satisfies the potential distinctness condition at $\ell$ for $M$ and $m$.
\end{enumerate}
\end{lemma}
\begin{proof}
(1) Let $K$ be the field of Lemma~\ref{sha-eigensystem-split-field} for $r$, and $r_K$ the restriction of $r$ to $\operatorname{Gal}(\bar{\mathbf Q}/K)$, which is irreducible.
Let $w$ be a prime of $K$ above $\ell$ and $A$ a place of $\bar{\mathbf Q}$ above $w$.
The decomposition group $D_A$ of $A$ in $\operatorname{Gal}(\bar{\mathbf Q}/\mathbf Q)$ lies in $\operatorname{Gal}(\bar{\mathbf Q}/K)$, so it is also the decomposition group of $A$ in $\operatorname{Gal}(\bar{\mathbf Q}/K)$, and an upper triangular form of $r$ on $D_A$ with $\kappa_i\ne\kappa_j\chi_{\mathrm{cyc}}$ for all $i,j$ is one of $r_K$.
So $r_K$ is $\ell$-distinct at $w$.
By hypothesis $r_K$ is windowed, and condition ($\alpha$) of Theorem~\ref{sha-residually-reducible-automorphy} says that a windowed representation is potentially $\ell$-distinct at every prime above $\ell$ at which it is $\ell$-distinct; so $r_K$ is potentially $\ell$-distinct at $w$.
Since $r$ is $\ell$-distinct, its restriction to a decomposition group at $\ell$ is conjugate to an upper triangular representation, and Lemma~\ref{sha-eigensystem-flags}(2) for $\mathbf Q\subset K$ shows that $r$ is potentially $\ell$-distinct at $\ell$, the only prime of $\mathbf Q$ above $\ell$.

(2) If $m_i\ne m_j+1$ for all $i,j$, every restriction of $r$ satisfies the potential distinctness condition.
Otherwise $\ell$ divides $M$ and $r$ is $\ell$-distinct.
By (1) $r$ is potentially $\ell$-distinct, and by Lemma~\ref{sha-eigensystem-flags}(1) so is its restriction to $\operatorname{Gal}(\bar{\mathbf Q}/L)$, at every prime of $L$ above $\ell$.
\end{proof}
\leanref{Development/AutomorphicArithmetic/AutomorphicVarieties/ModularityArgument/AutomorphyLifting/SplitIrr.lean}{UpperTheory.V.Z.Win.PotQ}
\leanref{Development/AutomorphicArithmetic/AutomorphicVarieties/ModularityArgument/AutomorphyLifting/AdmissibleLevelRestriction/PotQOf.lean}{UpperTheory.V.Z.Win.A4.potQ}
\leanref{Development/AutomorphicArithmetic/AutomorphicVarieties/ModularityArgument/ThresholdArguments/DenseThresholdCases/SafeUP.lean}{UpperTheory.V.Z.Win.A4.safeUP'_alt}
\leanref{Development/AutomorphicArithmetic/AutomorphicVarieties/ModularityArgument/AutomorphyLifting/AdmissibleLevelRestriction/UnlAllOf.lean}{UpperTheory.V.Z.Win.A4.safeRes'}

\begin{lemma}\label{sha-eigensystem-rational}
Let $n\ge2$ and assume the Galois representations as in Theorem~\ref{sha-residually-reducible-automorphy}.
There are $c_0,A$ such that the following holds.
Let $u\in\mathbf Z$, let $m\in\mathbf Z^n$ be strictly decreasing with $m_i+m_{n+1-i}=u$ for all $i$, let $H\ge0$ with $|m_i|\le H$ for all $i$, let $M\ge1$, and let $p\nmid M$ be a prime.
Then there is an eigenvalue datum of rank $n$ at $p$ with weights $m$ and weight $u$, with at most $c_0(M(H+1))^A$ elements, that classifies at $p$, for every odd prime $\ell\ne p$ and every $\iota:\bar{\mathbf Q}_\ell\cong\mathbf C$, the set of irreducible continuous representations $r:\operatorname{Gal}(\bar{\mathbf Q}/\mathbf Q)\to GL_n(\bar{\mathbf Q}_\ell)$ with the following properties: $r$ has level $M$ and weights $m$ (Definition~\ref{sha-eigensystem-level}); $r$ preserves a nondegenerate bilinear form up to a character $\theta$ that is $\chi_{\mathrm{cyc}}^{u}$ times a character of finite order; the restriction of $r$ to every imaginary quadratic field on which it is irreducible is attached through $\iota$ to a regular algebraic cuspidal automorphic representation and is windowed; and $r$ satisfies the distinctness condition at $\ell$ for $M$ and $m$ (Definition~\ref{sha-eigensystem-distinctness}).
\end{lemma}
\begin{proof}
Let $K_0,\dots,K_{n^2+1}$ be imaginary quadratic fields with pairwise coprime discriminants, for instance $\mathbf Q(\sqrt{-q_i})$ for distinct primes $q_i\equiv3\pmod4$.
Let $c_i,A_i$ be the constants of Lemma~\ref{sha-eigensystem-unitary} for $K_i$ and $n$, and put $c_0=2^{n+1}(n^2+1)\sum_ic_i$ and $A=\sum_iA_i$, both sums over $0\le i\le n^2+1$.
Let $u$, $m$, $H$, $M$ and $p$ be as in the statement.
The prime $p$ divides at most one of the discriminants; after renumbering, $p$ is unramified in $K_1,\dots,K_{n^2+1}$.
For $1\le i\le n^2+1$ let $v_i$ be a prime of $K_i$ above $p$, of residue degree $f_i\in\{1,2\}$, and let $\delta_i$ be the quadratic character of $K_i$; the $\delta_i$ are pairwise distinct.
Lemma~\ref{sha-eigensystem-unitary}(2) for $K_i$ at $v_i$ gives an eigenvalue datum $D_i=(\beta^{(i)}_s)_{s\in S_i}$ at $p^{f_i}$ with weights $m$ and weight $u$, with $|S_i|\le c_i(M(H+1))^{A_i}\le\big(\sum_jc_j\big)(M(H+1))^A$.

\emph{The datum.} Let $S$ be the set of quadruples $(i,s,B,\epsilon)$ with $1\le i\le n^2+1$, $s\in S_i$, $B$ a multiset of $n$ complex numbers whose $f_i$-th powers form $\beta^{(i)}_s$, and $\epsilon\in\{0,1\}$; put $\beta_{(i,s,B,\epsilon)}=B$.
Since $f_i\le2$, each element of $\beta^{(i)}_s$ has at most two $f_i$-th roots, so there are at most $2^n$ such $B$, and $|S|\le2^{n+1}\sum_{i=1}^{n^2+1}|S_i|\le c_0(M(H+1))^A$.
The family $(\beta_t)_{t\in S}$ is stable under $\operatorname{Aut}(\mathbf C)$ because the $D_i$ are; and the window and integrality conditions at $p$ for $B$ follow from those at $p^{f_i}$ for $\beta^{(i)}_s$ by taking $f_i$-th roots.

\emph{Restriction.} Let $\ell\ne p$ be odd, $\iota$ an isomorphism, $r$ a representation as in the statement, and $K$ an imaginary quadratic field such that the restriction $r_K$ of $r$ to $\operatorname{Gal}(\bar{\mathbf Q}/K)$ is irreducible.
Then $r_K$ lies in the set classified in Lemma~\ref{sha-eigensystem-unitary}(2) for $K$.
Indeed, by hypothesis $r_K$ is attached through $\iota$ to a regular algebraic cuspidal automorphic representation $\pi$ of $GL_n(\mathbf A_K)$ and is windowed.
The conditions of level $M$ and weights $m$, except the last condition at the primes above $\ell$, pass from $r$ to $r_K$, since the inertia and decomposition groups of $K$ are subgroups of those of $\mathbf Q$.
By Lemma~\ref{sha-eigensystem-potential}(2), $r_K$ satisfies the potential distinctness condition at $\ell$ for $M$ and $m$.
It remains to find a Hecke character $\chi$ of $\mathbf Q$ with $(\pi,\chi)$ polarized (Definition~\ref{gal-polarized}).
The character $\theta\chi_{\mathrm{cyc}}^{n-1}$ is $\chi_{\mathrm{cyc}}^{u+n-1}$ times a character of finite order, so by class field theory it is the character attached through $\iota$ to an algebraic Hecke character $\chi$ of $\mathbf Q$.
Multiplying $\chi$ by a power of the quadratic character of $K$, which is trivial on the norms from $K$ and whose Galois character is trivial on $\operatorname{Gal}(\bar{\mathbf Q}/K)$, makes $\chi_\infty(-1)=(-1)^n$.
For almost all primes $w$ of $K$, Frobenius elements at $w$ and at its conjugate $cw$ are conjugate in $\operatorname{Gal}(\bar{\mathbf Q}/\mathbf Q)$, so $\pi$ has the same Satake parameters $B_w$ at $w$ and at $cw$.
Since $r$ preserves a form up to $\theta$, its Frobenius eigenvalues at $w$ are stable under $b\mapsto\theta(\sigma_w)b^{-1}$; through $\iota$ this says that $B_w$ is stable under $b\mapsto(\chi\circ N_{K/\mathbf Q})(w)b^{-1}$.
The representation $g\mapsto\pi({}^tc(g)^{-1})$ has at $w$ the inverses of the Satake parameters of $\pi$ at $cw$, so it has the same Satake parameters as $\pi\otimes(\chi\circ N_{K/\mathbf Q}\circ\det)^{-1}$ at almost all places.
By strong multiplicity one (Theorem~\ref{lit-jacquet-shalika}(1)) the two are isomorphic, and since the contragredient of a cuspidal representation is its composite with the transpose inverse \cite{GK75}, this isomorphism gives the pairing of Definition~\ref{gal-polarized}; so $(\pi,\chi)$ is polarized.

\emph{Classification.} Let $\ell\ne p$ be odd, $\iota$ an isomorphism and $r$ a representation as in the statement.
By Lemma~\ref{sha-eigensystem-twists}(1) at most $n^2$ characters $\delta$ satisfy $r\otimes\delta\cong r$, so there is an $i$ with $r\otimes\delta_i\not\cong r$, and by Lemma~\ref{sha-eigensystem-twists}(2) the restriction $r_i$ of $r$ to $\operatorname{Gal}(\bar{\mathbf Q}/K_i)$ is irreducible.
By the restriction step, $D_i$ classifies $r_i$ at $v_i$, with a class $s_i(r_i)\in S_i$.
Since $p\nmid\ell M$, $r$ is unramified at $p$; let $B$ be the multiset of eigenvalues of $\iota r(\sigma_p)$. A Frobenius element at $v_i$ acts as $\sigma_p^{f_i}$, so the $f_i$-th powers of the elements of $B$ form $\beta^{(i)}_{s_i(r_i)}$.
If $r$, $r'$ and $r''$ give the same $(i,s,B)$, then their restrictions to $\operatorname{Gal}(\bar{\mathbf Q}/K_i)$ are isomorphic, since $D_i$ classifies them, and by Lemma~\ref{sha-eigensystem-twists}(2) each of $r'$ and $r''$ is isomorphic to $r$ or to $r\otimes\delta_i$: the representations with a given $(i,s,B)$ form at most two isomorphism classes.
Choosing $\epsilon\in\{0,1\}$ to separate these two classes gives a classifying map $r\mapsto(i,s_i(r_i),B,\epsilon)$, and the eigenvalues of $\iota r(\sigma_p)$ are the elements of $B$.
\end{proof}
\leanref{Development/AutomorphicArithmetic/AutomorphicVarieties/ModularityArgument/ThresholdArguments/DenseThresholdCases/FrobeniusTraceCountBoundOfCmEigensystemCount.lean}{UpperTheory.GalN.V.Win.N.CountQAlt}
\leanref{Development/AutomorphicArithmetic/AutomorphicVarieties/ModularityArgument/ThresholdArguments/DenseThresholdCases/AutQN.lean}{UpperTheory.GalN.V.AdmRNAlt}
\leanref{Development/AutomorphicArithmetic/AutomorphicVarieties/ModularityArgument/ThresholdArguments/DenseThresholdCases/FrobeniusTraceCountBoundOfCmEigensystemCount.lean}{UpperTheory.GalN.V.Win.N.countQ_alt}
\leanref{Development/AutomorphicArithmetic/AutomorphicVarieties/GaloisLibrary/CrystallineBaseChange/QuadFam.lean}{UpperTheory.V.Br.quadFam}
\leanref{Development/AutomorphicArithmetic/Foundations/GaloisRepresentations/AdmissibleResidues/FrobeniusTraceCountBoundOfRatEigensystemCount.lean}{UpperTheory.GalN.V.Win.rootSum}
\leanref{Development/AutomorphicArithmetic/AutomorphicVarieties/ModularityArgument/ThresholdArguments/DenseThresholdCases/AutQN.lean}{UpperTheory.GalN.V.AdmResNAlt}
\leanref{Development/AutomorphicArithmetic/AutomorphicVarieties/ModularityArgument/ThresholdArguments/DenseThresholdCases/LevelAdmissibleRestrictionHolds.lean}{UpperTheory.V.Z.Win.A4.admissible_restriction_level_n_of_printed_results}
\leanref{Development/AutomorphicArithmetic/AutomorphicVarieties/ModularityArgument/ThresholdArguments/DenseThresholdCases/UnderCPl.lean}{UpperTheory.V.Z.Bc.polarization_k_of_printed_results}
\leanref{Development/AutomorphicArithmetic/AutomorphicVarieties/ModularityArgument/SemisimpleLevels/LevelFRes.lean}{UpperTheory.V.Lev.levelFRes}

\begin{lemma}\label{sha-eigensystem-traces}
Let $K$ be an imaginary quadratic field of discriminant $d_K$, and assume the Galois representations as in Theorem~\ref{sha-residually-reducible-automorphy}.
There are $c,A$ such that the following holds.
Let $k,N\ge1$, let $w$ be a prime of $K$ above a prime $p\nmid Nd_K$, of norm $q$, and let $\sigma_w$ be an arithmetic Frobenius element at $w$.
Then there is a trace datum of weight $k$ at $q$ (Definition~\ref{sha-eigensystem-data}) with at most $c(N(k+1))^A$ elements that classifies at $\sigma_w^{-1}$, for every odd prime $\ell\ne p$ and every $\iota:\bar{\mathbf Q}_\ell\cong\mathbf C$, the set of automorphic triples of weight $k$ and level $N$ (Definition~\ref{sha-eigensystem-triple}).
\end{lemma}
\begin{proof}
\emph{Stable triples.} Apply Lemma~\ref{sha-eigensystem-unitary}(1) with $n=3$, $m=(k-1,-1,-k-1)$, which satisfies $m_i+m_{4-i}=-2$, $H=k+1$ and level $N$; it gives an eigenvalue datum $(\beta_s)_{s\in S_1}$ at $q$ of weight $-2$ with $|S_1|\le c_1(N(k+2))^{A_1}\le c_12^{A_1}(N(k+1))^{A_1}$.
For $s\in S_1$ put $t_s=\sum_{b\in\beta_s}(qb)^{-1}$.
The family $(t_s)$ is stable under $\operatorname{Aut}(\mathbf C)$.
By the window condition for the weight $-2$, the three numbers $y_b=|qb|^{-1}$, $b\in\beta_s$, lie between $q^{-1/2}$ and $q^{1/2}$ and have product $1$.
Two of the three numbers $y_b-1$ have the same sign, say those for $b_1$ and $b_2$; then $y_{b_1}+y_{b_2}\le1+y_{b_1}y_{b_2}=1+y_{b_3}^{-1}$, and $y_{b_3}+y_{b_3}^{-1}\le q^{1/2}+1$, so $|t_s|\le\sum_by_b\le\sqrt q+2$.
And $q^k(qb)^{-1}=q^{k-1}b^{-1}$ is an algebraic integer, since $k-1\ge m_i$ for all $i$.
So $(t_s)_{s\in S_1}$ is a trace datum of weight $k$ at $q$.
Let $\rho$ be a stable triple. Then $\rho\otimes\chi_{\mathrm{cyc}}^{-1}$ lies in the set classified in Lemma~\ref{sha-eigensystem-unitary}(1): the twist lowers the weights $(k,0,-k)$ by $1$, and keeps irreducibility, the condition $\kappa_i\ne\kappa_j\chi_{\mathrm{cyc}}$ and the other conditions.
Let $s(\rho)$ be the class of $\rho\otimes\chi_{\mathrm{cyc}}^{-1}$. Since $\chi_{\mathrm{cyc}}(\sigma_w)=q$, the eigenvalues of $\iota\rho(\sigma_w^{-1})$ are the $(qb)^{-1}$ for $b\in\beta_{s(\rho)}$, so $\iota\operatorname{tr}\rho(\sigma_w^{-1})=t_{s(\rho)}$; and $s(\rho)=s(\rho')$ implies $\rho\cong\rho'$.

\emph{Endoscopic triples.} Apply Lemma~\ref{sha-eigensystem-rational} with $n=2$, $m=(k,-k)$, which satisfies $m_1+m_2=0$, $H=k$, level $N|d_K|$ and the prime $p$; it gives an eigenvalue datum $(\beta_s)_{s\in S_2}$ at $p$ of weight $0$ with $|S_2|\le c_2|d_K|^{A_2}(N(k+1))^{A_2}$.
Let $h$ be the residue degree of $w$, so that $q=p^h$, and for $s\in S_2$ put $t_s=1+\sum_{b\in\beta_s}b^{-h}$.
The family is stable under $\operatorname{Aut}(\mathbf C)$. The two numbers $y_b=|b|^{-h}$ lie between $q^{-1/2}$ and $q^{1/2}$ and have product $1$, so $(q^{1/2}y_{b_1}-1)(q^{1/2}y_{b_2}-1)\ge0$ gives $y_{b_1}+y_{b_2}\le q^{1/2}+q^{-1/2}$, and $|t_s|\le\sqrt q+2$. And $q^kt_s=p^{hk}+\sum_b(p^kb^{-1})^h$ is an algebraic integer.
Let $\rho$ be an endoscopic triple with its representation $W$. Then $W$ lies in the set classified in Lemma~\ref{sha-eigensystem-rational}: it preserves the alternating form on $\bar{\mathbf Q}_\ell^2$ up to $\det W$, which has finite order, and its restriction to an imaginary quadratic field on which it is irreducible is attached to the base change of the newform \cite{AC89} and is windowed by hypothesis; and it satisfies the distinctness condition at $\ell$ for $N|d_K|$ and $m$ (Definition~\ref{sha-eigensystem-distinctness}), since $m_i\ne m_j+1$ for all $i,j$, as $2k\ne\pm1$.
Let $s(\rho)$ be the class of $W$. Since $p\nmid Nd_K$, $W$ is unramified at $p$, and $\sigma_w$ acts on $\bar{\mathbf Q}$ as $\sigma_p^h$, up to inertia, for an arithmetic Frobenius element $\sigma_p$ at $p$; so $\iota\operatorname{tr}\rho(\sigma_w^{-1})=1+\iota\operatorname{tr}W(\sigma_p^{-h})=t_{s(\rho)}$.
If $s(\rho)=s(\rho')$, then $W\cong W'$, and $\rho\cong\rho'$.

\emph{Union.} The disjoint union of the two trace data, with the maps $s$ above on stable and on endoscopic triples, classifies all automorphic triples of weight $k$ and level $N$ at $\sigma_w^{-1}$, and has at most $c(N(k+1))^A$ elements with $c=c_12^{A_1}+c_2|d_K|^{A_2}$ and $A=A_1+A_2$.
\end{proof}
\leanref{Development/AutomorphicArithmetic/AutomorphicVarieties/ModularityArgument/ThresholdArguments/DenseThresholdCases/FrobeniusTraceCountBoundOfCmEigensystemCount.lean}{UpperTheory.GalN.V.Win.N.frobenius_trace_count_bound_of_cm_eigensystem_count_alt}
\leanref{Development/AutomorphicArithmetic/AutomorphicVarieties/ModularityArgument/ThresholdArguments/DenseThresholdCases/FrobeniusTraceCountBoundOfCmEigensystemCount.lean}{UpperTheory.GalN.V.Win.N.frobenius_trace_count_bound_of_rat_eigensystem_count_alt}
\leanref{Development/AutomorphicArithmetic/Foundations/GaloisRepresentations/AdmissibleResidues/FrobeniusTraceCountBoundOfRatEigensystemCount.lean}{UpperTheory.GalN.V.Win.traces_from_eigenvalues_rank3}
\leanref{Development/AutomorphicArithmetic/Foundations/GaloisRepresentations/AdmissibleResidues/FrobeniusTraceCountBoundOfRatEigensystemCount.lean}{UpperTheory.GalN.V.Win.traces_from_eigenvalues_rank2}
\leanref{Development/AutomorphicArithmetic/Foundations/GaloisRepresentations/AdmissibleResidues/WinDefs.lean}{UpperTheory.GalN.V.Win.TrData.Classifies.sum}

\begin{lemma}\label{sha-eigensystem-eisenstein}
Let $q\ge7$ and $k\ge1$ be integers, let $z$ be a root of unity, and let $t\in\mathbf C$ with $|t|\le\sqrt q+2$ and $q^kt$ an algebraic integer.
Let $t_z=q^{-k}z+1+q^kz^{-1}$.
Then $t\ne t_z$, the number $f=q^k(t-t_z)$ is an algebraic integer, and $\log|f|\le2k\log q+1$.
\end{lemma}
\begin{proof}
Let $Q=q^k\ge7$. Then $f=q^kt-z-Q-Q^2z^{-1}$, a sum of algebraic integers.
Since $(\sqrt q-3)^2\ge0$, $\sqrt q\le(q+9)/6\le(Q+9)/6$.
So $|t_z|\ge Q-1-Q^{-1}>(Q+9)/6+2\ge|t|$, because $Q\ge7$, and $t\ne t_z$.
Finally $|f|\le Q(\sqrt q+2)+1+Q+Q^2\le2Q^2$, again because $\sqrt q\le(Q+9)/6$ and $Q\ge7$, and $\log2\le1$.
\end{proof}
\leanref{Development/AutomorphicArithmetic/Foundations/GaloisRepresentations/AdmissibleResidues/FrobeniusTraceCountBoundOfRatEigensystemCount.lean}{UpperTheory.GalN.V.Win.eisenstein_trace_difference_log_bound}

\begin{lemma}\label{sha-eigensystem-norm}
Let $q\ge7$ and $k\ge1$ be integers, let $(t_s)_{s\in S}$ be a trace datum of weight $k$ at $q$ (Definition~\ref{sha-eigensystem-data}), and let $(z_o)_{o\in O}$ be a finite family of roots of unity such that for every $\sigma\in\operatorname{Aut}(\mathbf C)$ there is a bijection $\tau$ of $O$ with $z_{\tau(o)}=\sigma(z_o)$.
Let $t_o=q^{-k}z_o+1+q^kz_o^{-1}$.
Then there is a nonzero integer $N_0$ with $\log|N_0|\le|O|\,|S|\,(2k\log q+1)$ such that the following holds.
Let $\ell\nmid q$ be a prime, $\iota:\bar{\mathbf Q}_\ell\cong\mathbf C$, $\Gamma$ a group, $g\in\Gamma$, and $\mathcal X$ a set of representations of $\Gamma$ in $GL_3(\bar{\mathbf Q}_\ell)$ that the trace datum classifies at $g$ through $\iota$. For each $o\in O$ let $X_o\subset\mathcal X$ be a finite set of pairwise non-isomorphic representations.
Then $\iota\operatorname{tr}\rho(g)\ne t_o$ for all $o$ and $\rho\in X_o$, and
\[
 |N_0|_\ell\le\prod_{o\in O}\prod_{\rho\in X_o}\big|\operatorname{tr}\rho(g)-\iota^{-1}(t_o)\big|_\ell.
\]
\end{lemma}
\begin{proof}
For $o\in O$ and $s\in S$ let $f(o,s)=q^k(t_s-t_o)$.
By Lemma~\ref{sha-eigensystem-eisenstein} each $f(o,s)$ is a nonzero algebraic integer with $\log|f(o,s)|\le2k\log q+1$.
Every $\sigma\in\operatorname{Aut}(\mathbf C)$ permutes the $t_s$ and the $z_o$, and $\sigma(t_o)=q^{-k}\sigma(z_o)+1+q^k\sigma(z_o)^{-1}$, so $\sigma$ permutes the $f(o,s)$.
Hence $N_0=\prod_{o,s}f(o,s)$ is an algebraic integer fixed by $\operatorname{Aut}(\mathbf C)$, that is, an integer; it is nonzero, and $\log|N_0|=\sum_{o,s}\log|f(o,s)|\le|O|\,|S|\,(2k\log q+1)$.
Let $\ell$, $\iota$, $\mathcal X$ and $X_o$ be as in the statement, and $s$ a classifying map.
For $\rho\in X_o$, $\iota\operatorname{tr}\rho(g)=t_{s(\rho)}\ne t_o$.
Since the elements of $X_o$ are pairwise non-isomorphic, $s$ is injective on $X_o$.
Each $\iota^{-1}f(o,s)$ is an algebraic integer, so $|\iota^{-1}f(o,s)|_\ell\le1$, and $|q|_\ell=1$. Hence
\[
 |N_0|_\ell=\prod_{o,s}|\iota^{-1}f(o,s)|_\ell\le\prod_o\prod_{\rho\in X_o}|\iota^{-1}f(o,s(\rho))|_\ell=\prod_o\prod_{\rho\in X_o}\big|\operatorname{tr}\rho(g)-\iota^{-1}(t_o)\big|_\ell.
\]
\end{proof}
\leanref{Development/AutomorphicArithmetic/Foundations/GaloisRepresentations/AdmissibleResidues/FrobeniusTraceCountBoundOfRatEigensystemCount.lean}{UpperTheory.GalN.V.Win.frobenius_trace_norm_nonzero_log_bound}
\leanref{Development/AutomorphicArithmetic/AutomorphicVarieties/BuildingsResidueChars/ExistsIntEqProdOfGaloisPermutedIntegral.lean}{UpperTheory.exists_int_eq_prod_of_galois_permuted_integral}
\leanref{Development/AutomorphicArithmetic/AutomorphicVarieties/BuildingsResidueChars/AdicNormProdLeProdOfInjectiveClasses.lean}{UpperTheory.adicNorm_prod_le_prod_of_injective_classes}

\subsection{Cotangent lengths of graded algebras}\label{subsection-sha-graded-algebras}

\begin{definition}\label{sha-algebra-graded}
Let $\mathcal O$ be a discrete valuation ring, let $X$ be an abelian group, let $A=\bigoplus_{j\in X}A_j$ be an $X$-graded commutative $\mathcal O$-algebra, and let $x:A\to\mathcal O$ be a homomorphism of $\mathcal O$-algebras.
\begin{enumerate}
\item The \emph{saturated square} $\Sigma_x$ of $x$ is the $\mathcal O$-submodule of the $f\in A$ such that $cf\in(\ker x)^2$ for some nonzero $c\in\mathcal O$.
For $j\in X$, the \emph{cotangent lattice of weight $j$} is the image $L_j$ of $A_j\cap\ker x$ in $A/\Sigma_x$.
For a family $F=(F_i)$ of elements of $A$, let $\langle F\rangle$ be the $\mathcal O$-span of the images of the $F_i$ in $A/\Sigma_x$; the \emph{cotangent length} $\lambda_j(F)$ is the $\mathcal O$-length, possibly infinite, of $L_j/(L_j\cap\langle F\rangle)$.
\item Let $\xi:A\to B$ be a homomorphism of $\mathcal O$-algebras and $Q$ a $B$-module. A \emph{derivation at $\xi$} with values in $Q$ is an $\mathcal O$-linear map $D:A\to Q$ with $D(fg)=\xi(f)D(g)+\xi(g)D(f)$. An $\mathcal O$-linear map on $A$ has \emph{weight $j$} if it vanishes on $A_i$ for every $i\ne j$.
\item Let $X=\mathbf Z$, let $\Omega$ be a field that is an $\mathcal O$-algebra, and let $\xi:A\to\Omega$ be a homomorphism of $\mathcal O$-algebras. For $t\in\Omega^\times$, let $t\cdot\xi$ be the homomorphism equal to $t^j\xi$ on $A_j$ for every $j$; the $t\cdot\xi$ form the \emph{orbit} of $\xi$. The homomorphism $\xi$ is \emph{two-signed} if it does not vanish on all the $A_j$ with $j<0$, and does not vanish on all the $A_j$ with $j>0$.
\end{enumerate}
\end{definition}
\leanref{Development/AutomorphicArithmetic/AutomorphicVarieties/CommutativeAlgebra/ResMapBijective.lean}{UpperTheory.V.Alg.satSq}
\leanref{Development/AutomorphicArithmetic/AutomorphicVarieties/CommutativeAlgebra/ResMapBijective.lean}{UpperTheory.V.Alg.cotLat}
\leanref{Development/AutomorphicArithmetic/AutomorphicVarieties/CommutativeAlgebra/ResMapBijective.lean}{UpperTheory.V.Alg.cotSpan}
\leanref{Development/AutomorphicArithmetic/AutomorphicVarieties/CommutativeAlgebra/ResMapBijective.lean}{UpperTheory.V.Alg.cotLen}
\leanref{Development/AutomorphicArithmetic/AutomorphicVarieties/CommutativeAlgebra/ResMapBijective.lean}{UpperTheory.V.Alg.PtDer}
\leanref{Development/AutomorphicArithmetic/AutomorphicVarieties/CommutativeAlgebra/ResMapBijective.lean}{UpperTheory.V.Alg.HasWt}
\leanref{Development/AutomorphicArithmetic/AutomorphicVarieties/CommutativeAlgebra/ResMapBijective.lean}{UpperTheory.V.Alg.act}
\leanref{Development/AutomorphicArithmetic/AutomorphicVarieties/CommutativeAlgebra/OrbitVanishingLocal.lean}{UpperTheory.V.Alg.TwoSigned}

\begin{lemma}\label{sha-cotangent-finite}
Let $\mathcal O$ be a complete discrete valuation ring of characteristic $0$ with fraction field $E$, and let $\Omega$ be an algebraically closed field containing $\mathcal O$.
\begin{enumerate}
\item Let $C$ be a finitely generated $\mathcal O$-algebra in which every nonzero element of $\mathcal O$ is a unit. Suppose that for every homomorphism of $\mathcal O$-algebras $\xi:C\to\Omega$, every derivation at $\xi$ with values in $\Omega$ (Definition~\ref{sha-algebra-graded}) is zero. Then there are only finitely many homomorphisms of $\mathcal O$-algebras $C\to\Omega$, and an element of $C$ at which all of them vanish is $0$.
\end{enumerate}
Let $A=\bigoplus_{j\in\mathbf Z}A_j$ be a $\mathbf Z$-graded commutative $\mathcal O$-algebra such that $A_0$ is a complete local noetherian $\mathcal O$-algebra with the residue field of $\mathcal O$ and $A$ is generated over $A_0$ by finitely many elements of degrees $-2$, $-1$, $1$ and $2$, and let $x:A\to\mathcal O$ be a homomorphism of $\mathcal O$-algebras that vanishes on $A_j$ for $j\ne0$. (The proofs of (2) to (4) use only the noetherianity of $A_0$ and the generators of $A$; the completeness of $A_0$, its residue field and $x$ are not used.)
\begin{enumerate}\setcounter{enumi}{1}
\item For every $j$, $A_j$ is a finitely generated $A_0$-module.
\item For $n\ge1$, $A_{-2n}$ and $A_{2n}$ are the $A_0$-modules spanned by the products of $n$ elements of $A_{-2}$, respectively of $A_2$. So a homomorphism from $A$ to a field vanishes on all the $A_j$ with $j<0$ if and only if it vanishes on $A_{-2}$, and on all the $A_j$ with $j>0$ if and only if it vanishes on $A_2$.
\item Let $I$ be an ideal of $A$ generated by homogeneous elements, such that the image of $A_0$ in $A/I$ is a finitely generated $\mathcal O$-module, and let $Y$ be the set of homomorphisms of $\mathcal O$-algebras $A\to\Omega$ that vanish on $I$. Suppose that at every two-signed $\xi\in Y$ any two derivations at $\xi$ with values in $\Omega$ that vanish on $I$ are linearly dependent over $\Omega$. Then for $f\in A_{-2}$ and $g\in A_2$, the set of $\xi\in Y$ with $\xi(f)=1$ and $\xi(g)\ne0$ is finite. Moreover, if $h\in A$ vanishes at every $\xi\in Y$ with $\xi(f)=1$ and $\xi(g)\ne0$, then $c\,g^nh\in I+(f-1)A$ for some $n\ge0$ and some nonzero $c\in\mathcal O$.
\end{enumerate}
\end{lemma}
\begin{proof}
(1) Since the nonzero elements of $\mathcal O$ are units in $C$, $C$ is a finitely generated $E$-algebra. Let $\mathfrak m$ be a maximal ideal of $C$. Then $C/\mathfrak m$ is a finite extension of $E$ \cite[Corollary~5.24]{AM69}, so it embeds over $\mathcal O$ in the algebraically closed field $\Omega$.
The module of differentials $\Omega_{C/\mathcal O}$ satisfies $\Omega_{C/\mathcal O}=\mathfrak m\Omega_{C/\mathcal O}$. Otherwise there is a nonzero $C/\mathfrak m$-linear form $\lambda$ on $\Omega_{C/\mathcal O}/\mathfrak m\Omega_{C/\mathcal O}$; composed with the universal derivation $d$ and with an embedding $\iota$ of $C/\mathfrak m$ in $\Omega$, it gives a derivation $\iota\circ\lambda\circ d$ at the point $C\to C/\mathfrak m\to\Omega$, which is zero by hypothesis, so $\lambda$ vanishes on the image of $d$, which generates; a contradiction.
Since $\Omega_{C/\mathcal O}$ is a finitely generated $C$-module, Nakayama's lemma gives, for each maximal ideal $\mathfrak m$, an element $1-u$ with $u\in\mathfrak m$ that kills it; so its annihilator lies in no maximal ideal, and $\Omega_{C/\mathcal O}=0$.
So $C$ is unramified over the field $E$, and being finitely generated it is a finite product of finite separable extensions of $E$ \cite[Chapter~I, Propositions~3.2 and~3.5]{Mil80}, hence a finite $E$-algebra. The homomorphisms of $\mathcal O$-algebras $C\to\Omega$ are $E$-algebra homomorphisms, which are linearly independent over $\Omega$ by Dedekind's lemma; so there are at most $\dim_EC$ of them.
Being a finite product of fields, $C$ is reduced and the intersection of its maximal ideals is $0$. For each maximal ideal $\mathfrak m$, an embedding $\iota$ of $C/\mathfrak m$ in $\Omega$ as above gives a homomorphism $C\to C/\mathfrak m\to\Omega$ with kernel $\mathfrak m$. So an element of $C$ at which every homomorphism to $\Omega$ vanishes lies in every maximal ideal, and is $0$.

(2) The ring $A$ is a finitely generated $A_0$-algebra, so it is noetherian by the Hilbert basis theorem. So the ideal of $A$ generated by $A_j$ is generated by finitely many $t_1,\dots,t_k\in A_j$. For $y\in A_j$, write $y=\sum_ia_it_i$ with $a_i\in A$; the components of degree $j$ give $y=\sum_i(a_i)_0t_i$.

(3) For $a,b\in\mathbf Z$ write $A_aA_b$ for the $A_0$-module spanned by the products of an element of $A_a$ and an element of $A_b$. Fix $\varepsilon\in\{1,-1\}$, and say that $f\in A_j$ has property $P$ if $f\in A_\varepsilon A_{j-\varepsilon}+A_{2\varepsilon}A_{j-2\varepsilon}$ when $\varepsilon j\ge1$, and $f\in A_{2\varepsilon}A_{j-2\varepsilon}$ when $\varepsilon j\ge2$.
Let $z\in A_w$ with $|w|\le2$, and let $f\in A_j$ have property $P$; then $zf\in A_{w+j}$ has property $P$. Indeed, if $w=2\varepsilon$, then $zf\in A_{2\varepsilon}A_j$. If $w=\varepsilon$, then $zf\in A_\varepsilon A_j$; and if $\varepsilon(w+j)\ge2$, then $\varepsilon j\ge1$, so $f$ is a sum of products $pq$ with $p\in A_\varepsilon$, $q\in A_{j-\varepsilon}$, and $p'q'$ with $p'\in A_{2\varepsilon}$, $q'\in A_{j-2\varepsilon}$, and $zpq=(zp)q$ and $zp'q'=p'(zq')$ lie in $A_{2\varepsilon}A_{w+j-2\varepsilon}$. Otherwise $\varepsilon w\le0$, so each condition on $w+j$ implies the same condition on $j$, and $z$ is absorbed into the second factor.
For each $j$ the elements of $A_j$ with property $P$ form an $\mathcal O$-submodule $P_j$. The sum $\sum_jP_j$ contains $1$ and is stable under multiplication by the elements of $A_0$ and by the generators of $A$, which generate $A$ over $\mathcal O$ together with $A_0$; so it is $A$. Taking components of degree $j$ gives $P_j=A_j$: every element of $A_j$ has property $P$.
In particular $A_{2\varepsilon n}=A_{2\varepsilon}A_{2\varepsilon(n-1)}$ for $n\ge1$, and induction on $n$ gives the first statement.
If a homomorphism $\xi$ from $A$ to a field vanishes on $A_{-2}$ and $f\in A_j$ with $j<0$, then $f^2\in A_{2j}$, which is spanned by products of $-j$ elements of $A_{-2}$, so $\xi(f)^2=0$ and $\xi(f)=0$. The case $j>0$ is the same, and the converses are clear.

(4) Let $J=I+(f-1)A$, let $\varpi$ be a uniformizer of $\mathcal O$, let $C$ be the localization of $A/J$ at the powers of the image $s$ of $\varpi g$, and let $\phi:A\to C$ be the canonical map.
The ring $A/I$ is a finitely generated $\mathcal O$-algebra, generated by the images of finitely many elements spanning the image of $A_0$ and of the generators of $A$; so $C$ is a finitely generated $\mathcal O$-algebra. Every nonzero element of $\mathcal O$ is a unit times a power of $\varpi$, hence a unit in $C$.
The map $\psi\mapsto\psi\circ\phi$ is a bijection from the homomorphisms of $\mathcal O$-algebras $C\to\Omega$ onto the $\xi\in Y$ with $\xi(f)=1$ and $\xi(g)\ne0$, since $\varpi\ne0$ in $\Omega$. By (1), it suffices to show that every derivation $D$ of $C$ at such a $\psi$ is zero.
Let $\xi=\psi\circ\phi$. Then $D'=D\circ\phi$ is a derivation at $\xi$ that vanishes on $I$. So is the Euler derivation $\mathcal E$, defined by $\mathcal E(h)=j\,\xi(h)$ for $h\in A_j$: it is a derivation because $(j+k)\xi(h)\xi(h')=j\xi(h)\xi(h')+k\xi(h)\xi(h')$ for $h\in A_j$ and $h'\in A_k$, and it vanishes on $I$ because the homogeneous components of the elements of $I$ lie in $I$.
The homomorphism $\xi$ is two-signed, since $\xi(f)=1$ with $f\in A_{-2}$ and $\xi(g)\ne0$ with $g\in A_2$. So $cD'+c'\mathcal E=0$ for some $(c,c')\ne(0,0)$.
At $f$: $D'(f)=D(1)=0$ and $\mathcal E(f)=-2\ne0$, since $\Omega$ has characteristic $0$; so $c'=0$, $c\ne0$ and $D'=0$.
Every element $z$ of $C$ satisfies $zu=\phi(h)$ for some $h\in A$ and a unit $u=\phi(\varpi g)^n$, and $D(\phi(h))=D(u)=0$; so $\psi(u)D(z)=0$ and $D(z)=0$.
For the second statement, let $h$ be as there. For every homomorphism $\psi:C\to\Omega$, $\psi\circ\phi$ is one of the $\xi$ above, so $\psi(\phi(h))=0$; by the second half of (1), $\phi(h)=0$. In the localization $C$ this means that $s^n$ kills the image of $h$ in $A/J$ for some $n\ge0$, that is, $\varpi^ng^nh\in J=I+(f-1)A$; take $c=\varpi^n$.
\end{proof}
\leanref{Development/AutomorphicArithmetic/AutomorphicVarieties/CommutativeAlgebra/GradedModulesLength/EtHolds.lean}{UpperTheory.V.Alg.M.et_holds}
\leanref{Development/AutomorphicArithmetic/AutomorphicVarieties/CommutativeAlgebra/GradedModulesLength/FgHolds.lean}{UpperTheory.V.Alg.M.fg_holds}
\leanref{Development/AutomorphicArithmetic/AutomorphicVarieties/CommutativeAlgebra/GradedModulesLength/OzTwoMulEqPow.lean}{UpperTheory.V.Alg.M.oz_two_mul_eq_pow}
\leanref{Development/AutomorphicArithmetic/AutomorphicVarieties/CommutativeAlgebra/GradedModulesLength/SignHolds.lean}{UpperTheory.V.Alg.M.sign_holds}
\leanref{Development/AutomorphicArithmetic/AutomorphicVarieties/CommutativeAlgebra/GradedModulesLength/SliceHolds.lean}{UpperTheory.V.Alg.M.slice_holds}

\begin{lemma}\label{sha-cotangent-implicit}
Let $E$ be a field, $\tau$ a finite set with a function $w:\tau\to\mathbf Z$, and $E[[Y]]$ the ring of formal power series in variables $Y_i$, $i\in\tau$, where $Y_i$ has weight $w(i)$; a power series has weight $j$ if all its monomials do. Let $\tau_0=\{i:w(i)=0\}$, and let $E[[Y']]$ be the ring of power series in the $Y_i$ with $i\notin\tau_0$. Let $z_1,\dots,z_m\in E[[Y]]$, with $m=|\tau_0|$, be power series of weight $0$ without constant term, such that for every $i\in\tau_0$ some $E$-linear combination of the $z_k$ has linear part $Y_i$.
Then there is a homomorphism of $E$-algebras $\psi:E[[Y]]\to E[[Y']]$ that maps power series of weight $j$ to power series of weight $j$, preserves constant terms, satisfies $\psi(z_k)=0$ for all $k$, and preserves the coefficient of $Y_i$ for every $i\notin\tau_0$.
\end{lemma}
\begin{proof}
For $i\in\tau_0$ choose $c_{i1},\dots,c_{im}\in E$ such that $\sum_kc_{ik}z_k$ has linear part $Y_i$, and put $g_i=\sum_kc_{ik}z_k-Y_i$ and $G_i=0$; for $i\notin\tau_0$ put $g_i=0$ and $G_i=Y_i$. Each $g_i$ has weight $w(i)$ and no terms of degree less than $2$, and each $G_i$ is an element of $E[[Y']]$ of weight $w(i)$ without constant term.
The iteration $\ell^{(0)}=0$, $\ell^{(n+1)}_i=G_i-g_i(\ell^{(n)})$, where $g_i(\ell)$ is the substitution of $\ell_j$ for $Y_j$, converges $Y'$-adically, because the $g_i$ have no terms of degree less than $2$; its limit $\ell$ satisfies $\ell_i=G_i-g_i(\ell)$, and $\ell_i$ has weight $w(i)$ and no constant term, since every $\ell^{(n)}_i$ has.
Let $\psi$ be the substitution $Y_i\mapsto\ell_i$. It preserves weights, since $\ell_i$ has weight $w(i)$.
Modulo terms of degree at least $2$, $\psi(h)$ is the constant term of $h$ plus $\sum_ia_iG_i$, where $a_i$ is the coefficient of $Y_i$ in $h$, since $\psi$ maps terms of degree at least $2$ to terms of degree at least $2$ and $\psi(Y_i)=G_i-\psi(g_i)$. So $\psi$ preserves constant terms and the coefficients of the $Y_i$ with $i\notin\tau_0$.
For $i\in\tau_0$, $\psi(Y_i)=\ell_i=-\psi(g_i)=\psi(Y_i)-\sum_kc_{ik}\psi(z_k)$, so $\sum_kc_{ik}\psi(z_k)=0$. The square matrix $(c_{ik})$ is invertible, since its product with the matrix of the coefficients of the $Y_{i'}$, $i'\in\tau_0$, in the $z_k$ is the identity. So $\psi(z_k)=0$ for all $k$.
\end{proof}
\leanref{Development/AutomorphicArithmetic/AutomorphicVarieties/CommutativeAlgebra/GradedModulesLength/IftSubstLow.lean}{UpperTheory.V.Alg.M.ift_holds}

\begin{definition}\label{sha-cotangent-euler}
Let $B$ be a noetherian local ring with maximal ideal $\mathfrak m_B$, and let $k\ge1$. Write $e_1,\dots,e_k$ for the unit vectors of $\mathbf Z^k$, and $m\ge m'$ for $m,m'\in\mathbf Z^k$ if all coordinates of $m-m'$ are nonnegative.
\begin{enumerate}
\item A \emph{standard $\mathbf Z^k$-graded $B$-algebra} is a commutative $B$-algebra $R=\bigoplus_{m\in\mathbf Z^k}R_m$, graded by $B$-submodules, such that $R_m=0$ unless $m\ge0$, $R_0$ is the image of $B$, each $R_{e_i}$ is a finitely generated $B$-module, and $R_{m+e_i}=R_mR_{e_i}$ for $m\ge0$ and $1\le i\le k$. Its \emph{irrelevant ideal} $\operatorname{Irr}$ is the product of the ideals $R_{e_i}R$.
\item A \emph{graded module} over such an $R$ is a finitely generated $R$-module $F$ with a decomposition $F=\bigoplus_{m\in\mathbf Z^k}F_m$ into $B$-submodules such that $R_dF_m\subseteq F_{d+m}$ for all $d,m$. For $d\in\mathbf Z^k$, the shift $F(-d)$ is $F$ with the pieces $F(-d)_m=F_{m-d}$. For $s\in R_d$, $F/sF$ has the images of the $F_m$ as pieces, and the kernel $F[s]$ of the multiplication by $s$ has the pieces $F[s]\cap F_m$.
\item A list $u=((d_1,s_1),\dots,(d_n,s_n))$ of elements of $\mathbf Z^k\times R$ is \emph{valid} if $d_i\ge0$ and $s_i\in R_{d_i}$ for every $i$. For a graded module $F$ the \emph{Euler characteristic} $\chi(u,F):\mathbf Z^k\to\mathbf Z$ is defined by induction on $n$: $\chi(\emptyset,F)(m)$ is the length of $F_m$ over $B$, or $0$ if this length is infinite; and $\chi(u,F)(m)=\chi(u',F/s_1F)(m)-\chi(u',F[s_1](-d_1))(m)$ with $u'=((d_2,s_2),\dots,(d_n,s_n))$.
\item The list $u$ is \emph{admitted} for an ideal $\mathfrak a$ of $R$ if some power of $\mathfrak m_BR\cdot\operatorname{Irr}$ lies in $\mathfrak a+(s_1,\dots,s_n)$. A statement about $m\in\mathbf Z^k$ holds \emph{for large $m$} if there is $m_0$ such that it holds for all $m\ge m_0$.
\end{enumerate}
\end{definition}
\leanref{Formalism/SharpRestriction/ValuationFrobeniusEpsilon/LeastPrimeTwists/QuotientShiftIrreducible/NoetherianQuotientCalc/EulerVanishingBelowDimension.lean}{Formalism.MainNotions.V.Pf.V2.Calc.Std}
\leanref{Formalism/SharpRestriction/ValuationFrobeniusEpsilon/LeastPrimeTwists/QuotientShiftIrreducible/NoetherianQuotientCalc/EulerVanishingBelowDimension.lean}{Formalism.MainNotions.V.Pf.V2.Calc.HidaSetting}
\leanref{Formalism/SharpRestriction/ValuationFrobeniusEpsilon/LeastPrimeTwists/QuotientShiftIrreducible/NoetherianQuotientCalc/EulerVanishingBelowDimension.lean}{Formalism.MainNotions.V.Pf.V2.Calc.GMod.IsGr}
\leanref{Formalism/SharpRestriction/ValuationFrobeniusEpsilon/LeastPrimeTwists/QuotientShiftIrreducible/NoetherianQuotientCalc/EulerVanishingBelowDimension.lean}{Formalism.MainNotions.V.Pf.V2.Calc.Valid}
\leanref{Formalism/SharpRestriction/ValuationFrobeniusEpsilon/LeastPrimeTwists/QuotientShiftIrreducible/NoetherianQuotientCalc/EulerVanishingBelowDimension.lean}{Formalism.MainNotions.V.Pf.V2.Calc.chi}
\leanref{Formalism/SharpRestriction/ValuationFrobeniusEpsilon/LeastPrimeTwists/QuotientShiftIrreducible/NoetherianQuotientCalc/EulerVanishingBelowDimension.lean}{Formalism.MainNotions.V.Pf.V2.Calc.admitted}

\begin{lemma}[\cite{BH93,Mat86,Eis95}]\label{sha-cotangent-calculus}
Let $B$ be a noetherian local ring, $R$ a standard $\mathbf Z^k$-graded $B$-algebra, $F$ a graded module over $R$, and $u=((d_1,s_1),\dots,(d_n,s_n))$ and $u'$ valid lists (Definition~\ref{sha-cotangent-euler}).
\begin{enumerate}
\item The submodules of $F$ generated by homogeneous elements, with the induced pieces, and the quotients of $F$ by them are graded modules; so are the shifts of $F$, and $F/sF$ and $F[s]$ for homogeneous $s$. Each $F_m$ is a finitely generated $B$-module, and if a power of $\mathfrak m_BR\cdot\operatorname{Irr}$ kills $F$, then $F_m$ has finite length for large $m$.
\item If $F$ and $G$ are graded modules with an isomorphism carrying each $F_m$ onto $G_m$, then $\chi(u,F)=\chi(u,G)$.
\item Let $0\to F'\to F\to F''\to0$ be an exact sequence of graded modules whose maps carry the pieces of degree $m$ into those of degree $m$, the pieces of $F'$ and $F''$ being the inverse images and the images of those of $F$, and let $u$ be admitted for the annihilator of $F$. Then $\chi(u,F)(m)=\chi(u,F')(m)+\chi(u,F'')(m)$ and $\chi(u,F'')(m)\le\chi(u,F)(m)$ for large $m$.
\item If $u$ is admitted for the annihilator of $F$ and $u'$ is a permutation of $u$, then $\chi(u,F)(m)=\chi(u',F)(m)$ for large $m$.
\item If $s\in R_d$ is a nonzerodivisor on $F$, then $\chi(((d,s))u,F)=\chi(u,F/sF)$, where $((d,s))u$ is $u$ preceded by $(d,s)$.
\item If $u$ is admitted for the annihilator of $F$, then $\chi(u,F)(m)\ge0$ for large $m$. If the concatenation $uu'$ is admitted for the annihilator of $F$, then $\chi(uu',F)(m)\le\chi(u',F/(s_1,\dots,s_n)F)(m)$ for large $m$.
\item If $u'=((d_1,s'_1),\dots,(d_n,s'_n))$ and each $s_i-s'_i$ kills $F$, then $\chi(u,F)=\chi(u',F)$.
\item Let $d\ge0$, $s\in R_0$ and $s'\in R_d$ be such that $((0,s),(d,s'))u$ is valid and $((d,ss'))u$ is admitted for the annihilator of $F$. Then $\chi(((d,ss'))u,F)(m)=\chi(((0,s))u,F)(m)+\chi(((d,s'))u,F)(m)$ for large $m$.
\item If the empty list is admitted for the annihilator of $F$, then $\chi(u,F)(m)=\sum_{I}(-1)^{|I|}\operatorname{len}_BF_{m-\sum_{i\in I}d_i}$ for large $m$, the sum over the subsets $I$ of $\{1,\dots,n\}$.
\end{enumerate}
\end{lemma}
\leanref{Formalism/SharpRestriction/ValuationFrobeniusEpsilon/LeastPrimeTwists/QuotientShiftIrreducible/NoetherianQuotientCalc/RuleGr.lean}{Formalism.MainNotions.V.Pf.V2.Calc.ruleGr}
\leanref{Formalism/SharpRestriction/ValuationFrobeniusEpsilon/LeastPrimeTwists/QuotientShiftIrreducible/NoetherianQuotientCalc/RuleFin.lean}{Formalism.MainNotions.V.Pf.V2.Calc.ruleFin}
\leanref{Formalism/SharpRestriction/ValuationFrobeniusEpsilon/LeastPrimeTwists/QuotientShiftIrreducible/NoetherianQuotientCalc/RuleIso.lean}{Formalism.MainNotions.V.Pf.V2.Calc.ruleIso}
\leanref{Formalism/SharpRestriction/ValuationFrobeniusEpsilon/LeastPrimeTwists/QuotientShiftIrreducible/NoetherianQuotientCalc/EulerAdditivityExactSequence.lean}{Formalism.MainNotions.V.Pf.V2.Calc.euler_additivity_of_rules}
\leanref{Formalism/SharpRestriction/ValuationFrobeniusEpsilon/LeastPrimeTwists/QuotientShiftIrreducible/NoetherianQuotientCalc/RuleMono.lean}{Formalism.MainNotions.V.Pf.V2.Calc.ruleMono}
\leanref{Formalism/SharpRestriction/ValuationFrobeniusEpsilon/LeastPrimeTwists/QuotientShiftIrreducible/NoetherianQuotientCalc/EulerIndependenceOfOrder.lean}{Formalism.MainNotions.V.Pf.V2.Calc.euler_rule_two_of_graded_finite_iso}
\leanref{Formalism/SharpRestriction/ValuationFrobeniusEpsilon/LeastPrimeTwists/QuotientShiftIrreducible/NoetherianQuotientCalc/RuleNzd.lean}{Formalism.MainNotions.V.Pf.V2.Calc.ruleNzd}
\leanref{Formalism/SharpRestriction/ValuationFrobeniusEpsilon/LeastPrimeTwists/QuotientShiftIrreducible/NoetherianQuotientCalc/RulePos.lean}{Formalism.MainNotions.V.Pf.V2.Calc.rulePos}
\leanref{Formalism/SharpRestriction/ValuationFrobeniusEpsilon/LeastPrimeTwists/QuotientShiftIrreducible/NoetherianQuotientCalc/RuleUB.lean}{Formalism.MainNotions.V.Pf.V2.Calc.ruleUB}
\leanref{Formalism/SharpRestriction/ValuationFrobeniusEpsilon/LeastPrimeTwists/QuotientShiftIrreducible/NoetherianQuotientCalc/RuleCong.lean}{Formalism.MainNotions.V.Pf.V2.Calc.Cong.chi_cong}
\leanref{Formalism/SharpRestriction/ValuationFrobeniusEpsilon/LeastPrimeTwists/QuotientShiftIrreducible/NoetherianQuotientCalc/EulerRuleProductZero.lean}{Formalism.MainNotions.V.Pf.V2.Calc.euler_rule_product_zero}
\leanref{Formalism/SharpRestriction/ValuationFrobeniusEpsilon/LeastPrimeTwists/QuotientShiftIrreducible/NoetherianQuotientCalc/RuleSmall.lean}{Formalism.MainNotions.V.Pf.V2.Calc.ruleSmall}

\begin{lemma}\label{sha-cotangent-nullstellensatz}
Let $\mathcal O\to B\to A$ be homomorphisms of commutative rings with $A$ a finitely generated $B$-algebra, let $\Omega$ be an algebraically closed field that is an $\mathcal O$-algebra, let $\sigma:B\to\Omega$ be a homomorphism of $\mathcal O$-algebras, let $J$ be an ideal of $A$ and let $f\in A$. Suppose that the image of $b\in B$ lies in $J$ whenever $\sigma(b)=0$, and that $bf^n\notin J$ whenever $\sigma(b)\ne0$ and $n\ge0$. Then there is a homomorphism of $\mathcal O$-algebras $\xi:A\to\Omega$ that equals $\sigma$ on the image of $B$, vanishes on $J$, and has $\xi(f)\ne0$.
\end{lemma}
\begin{proof}
The products $bf^n$ with $\sigma(b)\ne0$ form a multiplicative set disjoint from $J$, so there is a prime $P\supseteq J$ disjoint from it; then the image of $b\in B$ lies in $P$ if and only if $\sigma(b)=0$, and $f\notin P$.
Let $\mathfrak p=\ker\sigma$, the inverse image of $P$ in $B$. The residue field $\kappa(\mathfrak p)$ maps to $\Omega$ through $\sigma$ and to the residue field $\kappa(P)$ of $P$; the tensor product $\Omega\otimes_{\kappa(\mathfrak p)}\kappa(P)$ is nonzero, so it has a maximal ideal, whose inverse image $Q$ in $\Omega\otimes_BA$ is a prime with $1\otimes a\in Q$ if and only if $a\in P$.
The ring $\Omega\otimes_BA$ is a finitely generated $\Omega$-algebra, so by the Nullstellensatz $Q$ is the intersection of the maximal ideals containing it, and their residue fields are $\Omega$ \cite[Theorem~4.19]{Eis95}. As $1\otimes f\notin Q$, some maximal ideal $\mathfrak M\supseteq Q$ does not contain $1\otimes f$.
The composite of $a\mapsto1\otimes a$ with the quotient map by $\mathfrak M$ is the required $\xi$.
\end{proof}
\leanref{Development/AutomorphicArithmetic/AutomorphicVarieties/CommutativeAlgebra/GradedModulesLength/NsPoint.lean}{UpperTheory.V.Alg.M.ns_holds}

\begin{lemma}\label{sha-cotangent-herbrand}
Let $\mathcal O$ be a complete discrete valuation ring of characteristic $0$, let $\Omega$ be an algebraically closed field containing $\mathcal O$, and let $v:\Omega\to\mathbf R$ be a function such that, for nonzero $a,b$, $v(ab)=v(a)+v(b)$ and, if $a+b\ne0$, $v(a+b)\ge\min(v(a),v(b))$, and such that $v(c)$ is the normalized valuation of $c$ for nonzero $c\in\mathcal O$.
Let $B$ be a complete local noetherian $\mathcal O$-algebra with the residue field of $\mathcal O$, and $M$ a finitely generated $B$-module such that $B/\operatorname{Ann}_BM$ is finite over $\mathcal O$. For a homomorphism of $\mathcal O$-algebras $\sigma:B\to\Omega$, let $\ell_\sigma(M)$ be the length, possibly infinite, of the localization of $M$ at the prime $\ker\sigma$.
Let $a\in B$, let $S$ be a finite set of homomorphisms of $\mathcal O$-algebras $B\to\Omega$, and let $c\ge0$ be an integer, such that $\sigma(a)\ne0$, $v(\sigma(a))\ge0$ and $\ell_\sigma(M)\le c$ for $\sigma\in S$, and $\ell_\sigma(M)=0$ for $\sigma\notin S$.
Let $\theta(a,M)$ be the length of $M/aM$ minus the length of the kernel $M[a]$ of $a$ on $M$, each taken to be $0$ if it is infinite. Then
\[
 \theta(a,M)\le c\sum_{\sigma\in S}v(\sigma(a)).
\]
\end{lemma}
\begin{proof}
Let $I$ be an ideal of $B$ with $B/I$ finite over $\mathcal O$. We show, for every finitely generated $B$-module $N$ killed by $I$ with $\ell_\sigma(N)$ finite for $\sigma\in S$ and $0$ for $\sigma\notin S$: $N/aN$ and $N[a]$ have finite length, and
\[
 \operatorname{len}(N/aN)-\operatorname{len}(N[a])=\sum_{\sigma\in S}\ell_\sigma(N)\,v(\sigma(a)).
\]
Applied to $N=M$ and $I=\operatorname{Ann}_BM$, this gives the claim, since $v(\sigma(a))\ge0$ and $\ell_\sigma(M)\le c$ for $\sigma\in S$.
The module $N$ has a filtration with subquotients $B/\mathfrak p$ for primes $\mathfrak p$ \cite[Theorem~6.4]{Mat86}, and the statement is proved along it.
\begin{enumerate}
\item[(1)] For $N=0$ both sides vanish.
\item[(2)] For an exact sequence $0\to N_1\to N_2\to N_3\to0$, $\ell_\sigma(N_2)=\ell_\sigma(N_1)+\ell_\sigma(N_3)$, since localization is exact; and the snake lemma for the multiplication by $a$ gives an exact sequence $0\to N_1[a]\to N_2[a]\to N_3[a]\to N_1/aN_1\to N_2/aN_2\to N_3/aN_3\to0$. So the statement for $N_1$ and $N_3$ gives it for $N_2$.
\item[(3)] Let $N=B/\mathfrak p$, with $I\subseteq\mathfrak p$, so that $B/\mathfrak p$ is finite over $\mathcal O$. An $\mathcal O$-algebra map $\sigma$ whose kernel does not contain $\mathfrak p$ has $\ell_\sigma(B/\mathfrak p)=0$.
If $\mathfrak p$ contains a nonzero element of $\mathcal O$, then $\mathfrak p$ is maximal, since $B/\mathfrak p$ is a domain integral over the residue field of $\mathcal O$. Then $B/\mathfrak p$ is a simple module, on which $a$ acts by $0$ or invertibly, so the left side vanishes. So does the right side: no kernel of a $\sigma$ contains $\mathfrak p$, because $\mathcal O\to\Omega$ is injective.
Otherwise $D=B/\mathfrak p$ is a domain, finite and torsion free over $\mathcal O$. Every $\mathcal O$-algebra map $\tau:D\to\Omega$ is injective, and $\sigma=\tau\circ(B\to D)$ has kernel $\mathfrak p$ and $\ell_\sigma(B/\mathfrak p)=1$; so these $\sigma$ lie in $S$, and they are the $\sigma$ with $\ell_\sigma(B/\mathfrak p)\ne0$. As $\sigma(a)\ne0$ for them, the image of $a$ in $D$ is nonzero, so $D[a]=0$. The lengths of $D/aD$ over $B$ and over $\mathcal O$ agree, since $B$ has the residue field of $\mathcal O$. As $D$ is free of finite rank over $\mathcal O$, the length of $D/aD$ over $\mathcal O$ is the valuation of the determinant of the multiplication by $a$, the norm of $a$, which is $\prod_\tau\tau(a)$ over the $\mathcal O$-algebra maps $\tau:D\to\Omega$; so it is $\sum_\tau v(\tau(a))$, which is the right side.
\end{enumerate}
\end{proof}
\leanref{Development/AutomorphicArithmetic/AutomorphicVarieties/CommutativeAlgebra/GradedModulesLength/HerbrandFinLen.lean}{UpperTheory.V.Alg.M.th_holds}
\leanref{Development/AutomorphicArithmetic/AutomorphicVarieties/CommutativeAlgebra/GradedModulesLength/LengthQuotientSpanEqFinsumValuation.lean}{UpperTheory.V.Alg.M.length_quotient_span_eq_finsum_valuation}

\begin{lemma}\label{sha-algebra-cotangent}
Let $\mathcal O$ be a complete discrete valuation ring of characteristic $0$, let $\Omega$ be an algebraically closed field containing $\mathcal O$, and let $v:\Omega\to\mathbf R$ be a function such that, for nonzero $a,b$, $v(ab)=v(a)+v(b)$ and, if $a+b\ne0$, $v(a+b)\ge\min(v(a),v(b))$, and such that $v(c)$ is the normalized valuation of $c$ for nonzero $c\in\mathcal O$.
Let $A=\bigoplus_{j\in\mathbf Z}A_j$ be a $\mathbf Z$-graded commutative $\mathcal O$-algebra and $x:A\to\mathcal O$ a homomorphism of $\mathcal O$-algebras that vanishes on $A_j$ for $j\ne0$, such that $A_0$ is a complete local noetherian $\mathcal O$-algebra with the residue field of $\mathcal O$, and $A$ is generated over $A_0$ by finitely many elements of degrees $-2$, $-1$, $1$ and $2$.
Let $r\ge1$ and $s,r',s',m\ge0$, and let $F=(F_0,\dots,F_r)$ in $A_{-1}$, $G=(G_1,\dots,G_s)$ in $A_{-2}$, $(E_1,\dots,E_{r'})$ in $A_1$, $(E'_1,\dots,E'_{s'})$ in $A_2$ and $Z=(Z_1,\dots,Z_m)$ in $A_0\cap\ker x$.
Let $I$ be the ideal of $A$ that they generate, and $Y$ the set of homomorphisms of $\mathcal O$-algebras $A\to\Omega$ that vanish on $I$.
In the notation of Definition~\ref{sha-algebra-graded}, assume:
\begin{enumerate}
\item There is a graded homomorphism $\varphi$ from $A$ to a ring of formal power series over the fraction field of $\mathcal O$ in $r$ variables of weight $-1$, $s$ of weight $-2$, $r'+1$ of weight $1$, $s'+1$ of weight $2$ and $m$ of weight $0$, such that the constant term of $\varphi(f)$ is $x(f)$; every variable is the only linear term of $\varphi(f)$ for some $f$; and $f\in\Sigma_x$ whenever $x(f)=0$ and $\varphi(f)$ has no linear term. Moreover every element of $L_{-1}$, $L_{-2}$ and $L_0$ has a nonzero multiple in $\langle F\rangle$, $\langle G\rangle$ and $\langle Z\rangle$ respectively.
\item The image of $A_0$ in $A/I$ is a finitely generated $\mathcal O$-module.
\item Let $x_1\in Y$, different from $x$, vanish on $A_j$ for all $j\ne0$. Then either no $\xi\in Y$ that agrees with $x_1$ on $A_0$ vanishes on all $A_j$ with $j<0$ but not on all $A_j$ with $j>0$, or no such $\xi$ vanishes on all $A_j$ with $j>0$ but not on all $A_j$ with $j<0$.
\item At every two-signed $\xi\in Y$, any two derivations at $\xi$ with values in $\Omega$ that vanish on $I$ are linearly dependent over $\Omega$.
\item The element $a\in A_0$ satisfies $x(a)=0$, and $\xi(a)\ne0$ for every two-signed $\xi\in Y$.
\end{enumerate}
Then there is a finite set of two-signed elements of $Y$ that meets the orbit of every two-signed element of $Y$, the cotangent length $\lambda_{-1}(F)$ is finite, and for every such finite set $T$,
\[
 \lambda_{-1}(F)\le2\sum_{\xi\in T}v(\xi(a)).
\]
\end{lemma}
\begin{proof}
Let $E$ be the fraction field of $\mathcal O$ and $\varpi$ a uniformizer. The ideal $I$ is generated by homogeneous elements, and $x$ vanishes on the $F_i$, which lie in $A_{-1}$.

\emph{Step 1: the rank.} The cotangent lattice $L_{-1}$ is a free $\mathcal O$-module of rank $r$, the number of variables of weight $-1$ of the chart $\varphi$ of (1).
It is finitely generated: by Lemma~\ref{sha-cotangent-finite}(2), $A_{-1}$ is spanned over $A_0$ by a finite set $S$, and $x$ vanishes on $A_{-1}$; for $b\in A_0$ and $h\in A_{-1}$, $(b-x(b))h\in(\ker x)^2\subset\Sigma_x$, so $bh$ and $x(b)h$ have the same image, and $L_{-1}$ is spanned over $\mathcal O$ by the images of $S$.
For a variable $Y_i$ of weight $-1$, let $\lambda_i(f)$ be the coefficient of $Y_i$ in $\varphi(f)$. It vanishes on $(\ker x)^2$, since the power series $\varphi(h)$ with $h\in\ker x$ have no constant term and a product of two of them has no linear term; and so it vanishes on $\Sigma_x$, since $E$ is torsion free. So the $\lambda_i$ define an $\mathcal O$-linear map $\Lambda:L_{-1}\to E^r$.
It is injective: if $f\in A_{-1}$ and $\lambda_i(f)=0$ for all $i$, then $\varphi(f)$ has no linear term, since its coefficients at the variables of weight other than $-1$ vanish because $\varphi$ is graded, and $x(f)=0$; so $f\in\Sigma_x$ by (1).
So $L_{-1}$ is torsion free, hence free, and its rank is at most $r$. Conversely, for each variable $Y_i$ of weight $-1$ let $f_i$ be an element such that $Y_i$ is the only linear term of $\varphi(f_i)$. Its component of degree $-1$ has the same coefficients at the variables of weight $-1$, because $\varphi$ is graded; so the images of these components in $L_{-1}$ are mapped by $\Lambda$ to $r$ linearly independent vectors, and the rank is $r$.

\emph{Step 2: elementary divisors.} The quotient $A/\Sigma_x$ is torsion free: if $c\ne0$ and $cf\in\Sigma_x$, then $dcf\in(\ker x)^2$ for some $d\ne0$, so $f\in\Sigma_x$. By (1), $\langle F\rangle$ has finite index in $L_{-1}$, so it has rank $r$. By the elementary divisor theorem there are a basis $y_1,\dots,y_r$ of $L_{-1}$ and integers $c_1,\dots,c_r\ge0$ such that $\varpi^{c_1}y_1,\dots,\varpi^{c_r}y_r$ is a basis of $\langle F\rangle$; hence $\lambda_{-1}(F)=c_1+\dots+c_r$, which is finite. Since $\langle F\rangle$ is the image of $\mathcal O^{r+1}$ under $(b_i)\mapsto\sum_ib_iF_i$ modulo $\Sigma_x$, an invertible $(r+1)\times(r+1)$ matrix over $\mathcal O$ transforms $F$ into a family $F'=(F'_0,\dots,F'_r)$ in $A_{-1}$ whose images in $A/\Sigma_x$ are $\varpi^{c_1}y_1,\dots,\varpi^{c_r}y_r$ and $0$. The families $F$ and $F'$ have the same $\mathcal O$-span, so they generate the same ideal, have the same cotangent length, and both satisfy (1). Replace $F$ by $F'$, and lift each $y_i$ to an element of $A_{-1}$, still written $y_i$.

\emph{Step 3: removing the weight $0$.} There is a graded homomorphism $\varphi'$ from $A$ to a ring of formal power series over $E$ in $r$ variables of weight $-1$, $s$ of weight $-2$, $r'+1$ of weight $1$ and $s'+1$ of weight $2$, with no variable of weight $0$, such that the constant term of $\varphi'(f)$ is $x(f)$, every variable is the only linear term of $\varphi'(f)$ for some $f$, and $\varphi'(Z_i)=0$ for all $i$.
Indeed, the power series $\varphi(Z_k)$ have weight $0$ and no constant term, since $Z_k\in A_0\cap\ker x$. For a variable $Y_i$ of weight $0$, let $f$ be an element such that $Y_i$ is the only linear term of $\varphi(f)$, let $f_0$ be its component of degree $0$, and $h=f_0-x(f_0)\in A_0\cap\ker x$. Since $\varphi$ is graded, the linear part of $\varphi(h)$ is $\mu Y_i$ with $\mu\ne0$. By (1), $ch\equiv\sum_kb_kZ_k$ modulo $\Sigma_x$ for some nonzero $c\in\mathcal O$ and some $b_k\in\mathcal O$; the linear coefficients of $\varphi$ vanish on $\Sigma_x$, as in Step 1, so $\sum_k(c\mu)^{-1}b_k\varphi(Z_k)$ has linear part $Y_i$.
Lemma~\ref{sha-cotangent-implicit}, applied to the $m$ power series $\varphi(Z_k)$, gives a homomorphism $\psi$ to the power series in the variables of nonzero weight, and $\varphi'=\psi\circ\varphi$ has the required properties: it is graded and has constant terms $x(f)$, because $\psi$ preserves weights and constant terms; it kills the $Z_k$; and a variable that is the only linear term of $\varphi(f)$ is the only linear term of $\varphi'(f)$, because $\psi$ preserves the coefficients of the variables of nonzero weight.

\emph{Step 4: a representative of $a$.} There is $a'\in A_0$ with $a'-a\in I$, $x(a')=0$ and $\varphi'(a')\ne0$.
If $\varphi'(a)\ne0$, take $a'=a$. Otherwise choose a variable $X$ of weight $-1$ and a variable $X'$ of weight $1$ of $\varphi'$; since $\varphi'$ is graded, there are $f\in A_{-1}$ and $g\in A_1$ such that the coefficient of $X$ in $\varphi'(f)$ and that of $X'$ in $\varphi'(g)$ are nonzero.
Some $F_k$ has $\varphi'(F_k)\ne0$. Otherwise the coefficient $\lambda(h)$ of $X$ in $\varphi'(h)$ is an $\mathcal O$-linear form on $A$ that vanishes on $(\ker x)^2$, because the constant term of $\varphi'(h)$ is $x(h)$; on $\Sigma_x$, because $E$ is torsion free; and on the $F_k$. By (1), a nonzero multiple $cf$ lies in the span of the $F_k$ modulo $\Sigma_x$, so $c\lambda(f)=0$, which contradicts $\lambda(f)\ne0$.
Put $a'=a+gF_k$. Then $a'\in A_0$ and $a'-a\in I$; $x(a')=x(a)=0$, since $x(g)=0$; and $\varphi'(a')=\varphi'(g)\varphi'(F_k)\ne0$, since a power series ring over a field is a domain.
Every element of $Y$ vanishes on $I$, so it takes the same value at $a'$ as at $a$, and condition (5) holds for $a'$.

\emph{Step 5: a finite covering set.} By Lemma~\ref{sha-cotangent-finite}(2) and (3), the modules $A_{-2}$ and $A_2$ are generated over $A_0$ by finite sets $S_{-2}$ and $S_2$, and a homomorphism $\xi:A\to\Omega$ is two-signed if and only if it does not vanish on $A_{-2}$ and does not vanish on $A_2$. For $f\in S_{-2}$ and $g\in S_2$, the set of $\xi\in Y$ with $\xi(f)=1$ and $\xi(g)\ne0$ is finite, by Lemma~\ref{sha-cotangent-finite}(4), whose hypotheses are conditions (2) and (4) and the homogeneity of $I$.
Let $T_0$ be the set of the two-signed elements of these finitely many finite sets. Let $\xi\in Y$ be two-signed; then $\xi(f)\ne0$ and $\xi(g)\ne0$ for some $f\in S_{-2}$ and $g\in S_2$. Choose $t\in\Omega$ with $t^2=\xi(f)$. Then $t\cdot\xi$ vanishes on $I$, which is generated by homogeneous elements, is two-signed, takes the value $t^{-2}\xi(f)=1$ at $f$ and the value $t^2\xi(g)\ne0$ at $g$; so it lies in $T_0$, and $T_0$ meets the orbit of $\xi$.

\emph{Step 6: the bound.} For every finite set $T$ of two-signed elements of $Y$ that meets the orbit of every two-signed element of $Y$,
\[
 c_1+\dots+c_r\le2\sum_{\xi\in T}v(\xi(a')).
\]
Here $F$ is the family of Step 2, with $F_{i-1}-\varpi^{c_i}y_i\in\Sigma_x$ for $1\le i\le r$ and $F_r\in\Sigma_x$.
Let $N_1=A_{-2}$, $N_2=A_2$ and $N_3=A_0F_r^2+A_{-2}A_{-2}A_2$, which are $A_0$-submodules of $A$, and for $m\in\mathbf N^3$ let $N^m=N_1^{m_1}N_2^{m_2}N_3^{m_3}$, with $N^0=A_0$.
The \emph{Rees algebra} $\mathcal R$ is the $A_0$-subalgebra of the polynomial ring $A[T_1,T_2,T_3]$ formed by the polynomials $\sum_mf_mT^m$ with $f_m\in N^m$ for all $m$. It is graded by $\mathbf Z^3$, with $\mathcal R_m=N^mT^m$ for $m\in\mathbf N^3$, $\mathcal R_m=0$ for the other $m$, and $\mathcal R_0=A_0$.
Let $\mathcal J\subset\mathcal R$ be the ideal of the polynomials whose coefficients lie in the kernel of $\varphi'$; it is generated by homogeneous elements, and it contains the $Z_i$, since $\varphi'(Z_i)=0$.
Let $\theta$ be the list of homogeneous elements
\[
 F_0^2T_1,\dots,F_{r-1}^2T_1,\ G_1T_1,\dots,G_sT_1,\ E_1^2T_2,\dots,E_{r'}^2T_2,\ E'_1T_2,\dots,E'_{s'}T_2,\ F_r^2T_3,
\]
of degrees $e_1$, $e_2$ and $e_3$, the unit vectors of $\mathbf Z^3$, and let $(\theta)$ be the ideal they generate.
The algebra $\mathcal R$ is a standard $\mathbf Z^3$-graded $A_0$-algebra (Definition~\ref{sha-cotangent-euler}): $\mathcal R_{e_i}=N_iT_i$ is finitely generated over $A_0$ by Lemma~\ref{sha-cotangent-finite}(2), $N_3$ being a sum of products of such modules, and $\mathcal R_{m+e_i}=\mathcal R_m\mathcal R_{e_i}$ because $N^{m+e_i}=N^mN_i$. The list $\theta$, with the degrees above, followed by $(0,a')$, is valid. Euler characteristics, admitted lists and ``for large $m$'' are those of Definition~\ref{sha-cotangent-euler} for $B=A_0$, and $\chi(a',\cdot)$ is $\chi(((0,a')),\cdot)$, with $a'$ viewed in $A_0=\mathcal R_0$. Write $fT_i$ for an element $f$ of $N_i$ placed in degree $e_i$, and $h_j$ for the component of degree $j$ of $h\in A$. Then the following hold.
\begin{enumerate}
\item[(a)] Some power of $\mathfrak m_{A_0}\mathcal R\cdot\operatorname{Irr}$, where $\operatorname{Irr}=\mathcal R_{e_1}\mathcal R_{e_2}\mathcal R_{e_3}\mathcal R$, lies in $\mathcal J+(\theta)+a'\mathcal R$: the list $\theta$ followed by $(0,a')$ is admitted for $\mathcal J$.
\item[(b)] For large $m$, $(c_1+\dots+c_r)\,2^{r+r'}\le\chi(\theta,M)(m)$, where $M=(\mathcal R/\mathcal J)/a'(\mathcal R/\mathcal J)$.
\item[(c)] For large $m$, $\chi(\theta,M)(m)\le\chi(a',\mathcal R/(\mathcal J+(\theta)))(m)$.
\item[(d)] For large $m$, $\chi(a',\mathcal R/(\mathcal J+(\theta)))(m)\le2^{r+1+r'}\sum_{\xi\in T}v(\xi(a'))$.
\end{enumerate}
For one $m$ that is large for (b), (c) and (d), they give $(c_1+\dots+c_r)\,2^{r+r'}\le2^{r+1+r'}\sum_{\xi\in T}v(\xi(a'))$, which is the bound of Step 6.

\emph{Proof of (a).} The ring $\mathcal R$ is noetherian, being finitely generated over $A_0$, so the ideal $\mathfrak m_{A_0}\mathcal R\cdot\operatorname{Irr}$ is finitely generated, and it suffices to show that it lies in every prime $\mathfrak p$ of $\mathcal R$ containing $\mathcal J+(\theta)+a'\mathcal R$. Let $\mathfrak p$ be one. If $\operatorname{Irr}\subseteq\mathfrak p$ there is nothing to show, so suppose not: for each $i$ there is $f_i\in N_i$ with $f_iT_i\notin\mathfrak p$. We show that $\mathfrak m_{A_0}\subseteq\mathfrak p$. Let $\mathfrak d\subseteq A_0$ be the ideal generated by the products $fg$ with $f\in A_{-2}$ and $g\in A_2$. The $Z_k$ lie in $\mathcal J\subseteq\mathfrak p$.
\begin{enumerate}
\item[(a1)] \emph{The prime $\mathfrak p$ contains $F_i^2T_1$ for $0\le i\le r$, and every element of $A_0\cap I$.} For $i<r$, $F_i^2T_1$ lies in $\theta$. As $N_3\subseteq A_{-2}=N_1$, $(F_r^2T_1)(f_3T_3)=(f_3T_1)(F_r^2T_3)\in\mathfrak p$, and $f_3T_3\notin\mathfrak p$.
Let $u=(f_1T_1)(f_2T_2)\notin\mathfrak p$, let $y$ be one of the generators $F_i$, $G_i$, $E_i$, $E'_i$, $Z_i$ of $I$, of degree $j$, and let $c\in A_{-j}$. Then $(cy)^2u=(F_i^2T_1)(c^2f_1f_2T_2)$ if $y=F_i$, $(cy)u=(G_iT_1)(cf_1f_2T_2)$ if $y=G_i$, $(cy)^2u=(c^2f_1f_2T_1)(E_i^2T_2)$ if $y=E_i$, and $(cy)u=(cf_1f_2T_1)(E'_iT_2)$ if $y=E'_i$; these lie in $\mathfrak p$, and so does $cZ_i$. As $\mathfrak p$ is prime and $u\notin\mathfrak p$, $cy\in\mathfrak p$ in all cases. An element of $A_0\cap I$ is a sum of products $hy$ with $h\in A$ and $y$ a generator, hence the sum of the components of degree $0$ of these products, which are the $h_{-j}y$; so it lies in $\mathfrak p$.
\item[(a2)] \emph{Two cases.} Let $\mathfrak p_0$ be the inverse image of $\mathfrak p$ in $A_0$, a prime that contains $A_0\cap I$ by (a1). The domain $A_0/\mathfrak p_0$ is a quotient of the image of $A_0$ in $A/I$, so it is finite over $\mathcal O$ by condition (2). If $\mathfrak p_0$ contains a nonzero element of $\mathcal O$, then $A_0/\mathfrak p_0$ is a domain integral over the residue field of $\mathcal O$, hence a field, and $\mathfrak p_0=\mathfrak m_{A_0}$, as required. Otherwise $\mathcal O$ embeds in $A_0/\mathfrak p_0$, which is algebraic over the fraction field of $\mathcal O$ and so embeds in $\Omega$ over $\mathcal O$: $\mathfrak p_0=\ker\sigma$ for a homomorphism of $\mathcal O$-algebras $\sigma:A_0\to\Omega$. Steps (a3) to (a9) show that this case does not occur.
\item[(a3)] \emph{Extensions.} Every homomorphism of $\mathcal O$-algebras $\sigma':A_0\to\Omega$ that kills $A_0\cap I$ is the restriction to $A_0$ of some $\xi\in Y$. Indeed, $A$ is finitely generated over $A_0$, and Lemma~\ref{sha-cotangent-nullstellensatz} applies with $B=A_0$, $J=I+(\ker\sigma')A$ and $f=1$: as $I$ is generated by homogeneous elements, the component of degree $0$ of an element of $J$ lies in $(A_0\cap I)+\ker\sigma'$, which $\sigma'$ kills, so no $b\in A_0$ with $\sigma'(b)\ne0$ lies in $J$.
\item[(a4)] \emph{$\mathfrak d\subseteq\mathfrak p_0$.} Otherwise $\sigma(fg)\ne0$ for some $f\in A_{-2}$ and $g\in A_2$. Let $\xi\in Y$ extend $\sigma$, by (a3). Then $\xi(f)\xi(g)\ne0$, so $\xi$ is two-signed, and $\xi(a')\ne0$ by condition (5) for $a'$ (Step 4). But $a'\in\mathfrak p$, so $\xi(a')=\sigma(a')=0$.
\item[(a5)] \emph{A point $\xi_0\in Y$ that restricts to $\sigma$ on $A_0$ and vanishes on $A_j$ for all $j\ne0$.} Let $\xi_0(h)=\sigma(h_0)$. It is additive and $\xi_0(1)=1$. The $h$ with $\xi_0(hh')=\xi_0(h)\xi_0(h')$ for all $h'$ form a subring containing $A_0$, so for multiplicativity it suffices to take for $h$ a generator $z$ of $A$ over $A_0$, of degree $j\in\{-2,-1,1,2\}$; then $\xi_0(z)=0$ and $(zh')_0=zh'_{-j}$. For $j=\pm2$, $zh'_{-j}\in\mathfrak d\subseteq\ker\sigma$ by (a4); for $j=\pm1$, its square lies in $\mathfrak d$, so $\sigma(zh'_{-j})^2=0$. So $\xi_0$ is a homomorphism. It vanishes on the generators of $I$ of nonzero degree, and on the $Z_i$, which lie in $A_0\cap I\subseteq\mathfrak p_0$ by (a1).
\item[(a6)] \emph{A point $\xi_-\in Y$ that restricts to $\sigma$ on $A_0$, vanishes on $A_j$ for all $j>0$, and does not vanish on all $A_j$ with $j<0$.} For $n\ge0$ and $h\in A$, let $\iota_n(h)=h_{-2n}T_1^n$, which lies in $\mathcal R$ because $A_{-2n}=N_1^n$ by Lemma~\ref{sha-cotangent-finite}(3). Then $\iota_{n+n'}(hh')=\iota_n(h)\iota_{n'}(h')$ for $h\in A_{-2n}$ and $h'\in A_{-2n'}$, and $\iota_1(f_1)=f_1T_1\notin\mathfrak p$. Let $J_-$ be the ideal of the $h\in A$ with $\iota_n(h'h)\in\mathfrak p$ for all $h'\in A$ and $n\ge0$; a homogeneous $g$ of degree $j$ lies in $J_-$ as soon as $\iota_n(h'g)\in\mathfrak p$ for all $n$ and $h'\in A_{-2n-j}$. As $\mathfrak p$ is prime, this gives, for homogeneous $g$ and for $b$:
if $g^2\in J_-$, then $g\in J_-$, since $\iota_{2n}(h'^2g^2)=\iota_n(h'g)^2$;
if $f_1g\in J_-$, then $g\in J_-$, since $\iota_{n+1}(h'f_1g)=\iota_1(f_1)\iota_n(h'g)$;
if $b\in A_0\cap\mathfrak p_0$, then $b\in J_-$, since $\iota_n(h'b)=\iota_n(h')b$;
if $b\in A_{-2}$ and $bT_1\in\mathfrak p$, then $b\in J_-$, since $\iota_1(f_1)\iota_n(h'b)=\iota_n(f_1h')\,bT_1$;
if $b\in A_0\setminus\mathfrak p_0$, then $bf_1^n\notin J_-$, since $\iota_n(bf_1^n)=b(f_1T_1)^n\notin\mathfrak p$.
Hence $A_2\subseteq J_-$: for $g\in A_2$, $f_1g\in\mathfrak d\subseteq\mathfrak p_0$ by (a4), so $f_1g\in J_-$ and $g\in J_-$. By Lemma~\ref{sha-cotangent-finite}(3), $A_{2j}$ is spanned over $A_0$ by products of $j$ elements of $A_2$, so for $g\in A_j$ with $j>0$, $g^2\in J_-$ and $g\in J_-$. The $F_i$ lie in $J_-$, because $F_i^2T_1\in\mathfrak p$ by (a1); the $G_i$, because $G_iT_1$ lies in $\theta$; the $E_i$ and $E'_i$, which have positive degree; and the $Z_i$. So $I\subseteq J_-$, the image of $\ker\sigma=\mathfrak p_0$ lies in $J_-$, and $bf_1^n\notin J_-$ when $\sigma(b)\ne0$. Lemma~\ref{sha-cotangent-nullstellensatz}, with $B=A_0$, $J=J_-$ and $f=f_1$, gives $\xi_-$, with $\xi_-(f_1)\ne0$.
\item[(a7)] \emph{A point $\xi_+\in Y$ that restricts to $\sigma$ on $A_0$, vanishes on $A_j$ for all $j<0$, and does not vanish on all $A_j$ with $j>0$.} Let $J_+$ be the ideal of the $h\in A$ with $(h'h)_{2(n+1)}T_2^{n+1}\in\mathfrak p$ for all $h'\in A$ and $n\ge0$; this element lies in $\mathcal R$ because $A_{2(n+1)}=N_2^{n+1}$. As in (a6), using that $\mathfrak p$ is prime and $f_2T_2\notin\mathfrak p$:
$A_j\subseteq J_+$ for $j<0$, because for $g\in A_j$ and $h'\in A_{2(n+1)-j}$ the square of $(h'g)T_2^{n+1}$ is $g^2h'^2T_2^{2n+2}$, and by Lemma~\ref{sha-cotangent-finite}(3) $g^2h'^2$ is a sum of products $cwv$ with $c\in A_{-2}$, $w\in A_2$ and $v\in N_2^{2n+2}$, where $(cw)(vT_2^{2n+2})\in\mathfrak p$ since $cw\in\mathfrak d\subseteq\mathfrak p_0$;
$g\in J_+$ for $g\in A_2$ with $gT_2\in\mathfrak p$, since $(h'g)T_2^{n+1}=(gT_2)(h'T_2^n)$ for $h'\in A_{2n}$;
$g\in J_+$ for $g\in A_1$ with $g^2T_2\in\mathfrak p$, since $((h'g)T_2^{n+1})^2=(g^2T_2)(h'^2T_2^{2n+1})$ for $h'\in A_{2n+1}$;
$A_0\cap\mathfrak p_0\subseteq J_+$; and $bf_2^n\notin J_+$ for $b\in A_0\setminus\mathfrak p_0$, since $(f_2bf_2^n)T_2^{n+1}=b(f_2T_2)^{n+1}\notin\mathfrak p$.
As $E_i^2T_2$ and $E'_iT_2$ lie in $\theta$, $I\subseteq J_+$, and Lemma~\ref{sha-cotangent-nullstellensatz} with $J=J_+$ and $f=f_2$ gives $\xi_+$, with $\xi_+(f_2)\ne0$.
\item[(a8)] \emph{A point $\xi\in Y$ that agrees with $x$ on $A_0$ and vanishes on $A_j$ for all $j>0$ vanishes on $A_j$ for all $j<0$.} Such a $\xi$ vanishes on $A_j\cap\ker x$ for all $j\ge0$. The $\mathcal O$-linear form $\lambda(h)=\xi(h_{-1})$ on $A$ vanishes on $(\ker x)^2$: for $u,u'\in\ker x$, $(uu')_{-1}$ is the sum of the $u_au'_b$ with $a+b=-1$, where $u_a,u'_b\in\ker x$ because $x$ vanishes in nonzero degrees, and one of $a$, $b$ is nonnegative. As $\Omega$ is torsion free over $\mathcal O$, $\lambda$ vanishes on $\Sigma_x$; it vanishes on the $F_i\in I$. By condition (1), every element of $L_{-1}$ has a nonzero multiple in $\langle F\rangle$, so $\xi$ vanishes on $A_{-1}$. In the same way, with the $G_i$, $\xi$ vanishes on $A_{-2}$: for $a+b=-2$, one of $a$, $b$ is nonnegative or both are $-1$. By Lemma~\ref{sha-cotangent-finite}(3), $\xi$ then vanishes on all $A_j$ with $j<0$.
\item[(a9)] \emph{The contradiction.} If $\xi_0$ is $x$ followed by $\mathcal O\subseteq\Omega$, then $\xi_-$ agrees with $x$ on $A_0$, so it vanishes on all $A_j$ with $j<0$ by (a8), contrary to (a6). Otherwise $\xi_0\in Y$ is different from $x$ and vanishes on $A_j$ for all $j\ne0$, and condition (3) for $x_1=\xi_0$ says that $\xi_+$ or $\xi_-$, which agree with $\xi_0$ on $A_0$, does not exist. So $\mathfrak p_0=\mathfrak m_{A_0}$.
\end{enumerate}

\emph{Proof of (c).} Let $F=\mathcal R/\mathcal J$, a graded module by Lemma~\ref{sha-cotangent-calculus}(1), as $\mathcal J$ is generated by homogeneous elements; by (a), $\theta$ followed by $(0,a')$ is admitted for $\mathcal J$, which kills $F$. If $a'p\in\mathcal J$, then $p\in\mathcal J$: each coefficient $p_m$ has $\varphi'(a')\varphi'(p_m)=0$, with $\varphi'(a')\ne0$ by Step 4, in a power series ring over a field, which is a domain. So $a'$ is a nonzerodivisor on $F$, and by Lemma~\ref{sha-cotangent-calculus}(5), $\chi(((0,a'))\theta,F)=\chi(\theta,F/a'F)=\chi(\theta,M)$. By (4) of that lemma, $\chi(((0,a'))\theta,F)(m)=\chi(\theta((0,a')),F)(m)$ for large $m$, and by (6), $\chi(\theta((0,a')),F)(m)\le\chi(a',F/(\theta)F)(m)$ for large $m$. Finally $F/(\theta)F=\mathcal R/(\mathcal J+(\theta))$, with the same pieces, so (2) of that lemma gives (c).

\emph{Proof of (b).} The proof uses four facts.
\begin{enumerate}
\item[(b1)] For $\nu\in\mathbf N^{\tau'}$, where $\tau'$ is the set of variables of $\varphi'$ with weights $w$, call $\sum_{w(i)<0}|w(i)|\nu_i$ and $\sum_{w(i)>0}w(i)\nu_i$ the negative and positive degrees of the monomial $Y^\nu$. Let $\mathfrak q$ be the set of the $p=\sum_mp_mT^m\in\mathcal R$ such that, for every $m$, every monomial of $\varphi'(p_m)$ has negative degree at least $2m_1+4m_3+1$. Then, since $F_r\in\Sigma_x$, $\mathfrak q$ is an ideal generated by homogeneous elements; it contains $\mathcal J$ and the $b\in A_0$ with $x(b)=0$; $cp\in\mathfrak q$ with $c\in\mathcal O$ nonzero implies $p\in\mathfrak q$; and $f^2T_1-g^2T_1\in\mathfrak q$ for $f,g\in A_{-1}$ with $f-g\in\Sigma_x$.
\item[(b2)] For $m\ge0$, the length over $A_0$ of the piece of degree $m$ of $\mathcal R/(\mathfrak q+\varpi\mathcal R)$ is the number $h(m)$ of monomials in the variables of $\varphi'$ with negative degree $2m_1+4m_3$ and positive degree $2m_2+2m_3$.
\item[(b3)] For a list $\delta=(\delta_1,\dots,\delta_p)$ in $\mathbf Z^3$ and $g:\mathbf Z^3\to\mathbf Z$, let $\Delta_\delta g$ be defined by $\Delta_\emptyset g=g$ and $\Delta_\delta g(m)=\Delta_{\delta'}g(m)-\Delta_{\delta'}g(m-\delta_1)$, with $\delta'=(\delta_2,\dots,\delta_p)$. For the list $\delta$ formed by $r-1+s$ copies of $e_1$, $r'+s'$ copies of $e_2$ and $e_3$, $\Delta_\delta h(m)=2^{r-1+r'}$ for large $m$; here $\varphi'$ has $r$ variables of weight $-1$, $s$ of weight $-2$, $r'+1$ of weight $1$ and $s'+1$ of weight $2$.
\item[(b4)] Let $F$ be a graded module over a standard $\mathbf Z^k$-graded algebra $R$ over $B$, let $\beta\in B$ be a nonzerodivisor on $F$ such that the empty list is admitted for the annihilator of $F/\beta F$, let $d\ge0$, $n\ge0$, $y_1,\dots,y_n\in R_d$, $c_1,\dots,c_n\ge0$, and let $u$ be a valid list such that the list $u_\beta$ of the $(d,\beta^{c_i}y_i)$ followed by $u$ is admitted for the annihilator of $F$. Then, for large $m$, $(c_1+\dots+c_n)\,\Delta_\delta h_\beta(m)\le\chi(u_\beta,F)(m)$, where $h_\beta(m)$ is the length of the piece of degree $m$ of $F/\beta F$ and $\delta$ is formed by $n-1$ copies of $d$ followed by the degrees of $u$.
\end{enumerate}
Let $\mathfrak q$ be as in (b1). Since $\mathcal J\subseteq\mathfrak q$ and $a'\in\mathfrak q$, as $x(a')=0$, the module $\mathcal R/\mathfrak q$ is the quotient of $M$ by a submodule generated by homogeneous elements; and $\theta$ is admitted for the annihilator of $M$, which contains $\mathcal J+a'\mathcal R$, by (a). So $\chi(\theta,\mathcal R/\mathfrak q)(m)\le\chi(\theta,M)(m)$ for large $m$, by Lemma~\ref{sha-cotangent-calculus}(3).
Let $\theta'$ be $\theta$ with $F_{i-1}^2T_1$ replaced by $(\varpi^{c_i}y_i)^2T_1=\varpi^{2c_i}y_i^2T_1$ for $1\le i\le r$. Since $F_{i-1}-\varpi^{c_i}y_i\in\Sigma_x$, the differences lie in $\mathfrak q$ by (b1), so $\chi(\theta,\mathcal R/\mathfrak q)=\chi(\theta',\mathcal R/\mathfrak q)$ by Lemma~\ref{sha-cotangent-calculus}(7), and $\theta'$ is admitted for $\mathfrak q$.
By (b1), $\varpi$ is a nonzerodivisor on $\mathcal R/\mathfrak q$; and $\mathfrak m_{A_0}$ kills $\mathcal R/(\mathfrak q+\varpi\mathcal R)$, because for $b\in\mathfrak m_{A_0}$, $x(b)\in\varpi\mathcal O$ and $b-x(b)\in\mathfrak q$. So (b4) applies to $F=\mathcal R/\mathfrak q$, $\beta=\varpi$, $d=e_1$, $n=r$, the $y_i^2T_1$ with the exponents $2c_i$, and the rest of $\theta$; with (b2) and (b3) it gives, for large $m$,
\[
 (c_1+\dots+c_r)\,2^{r+r'}=(2c_1+\dots+2c_r)\,2^{r-1+r'}\le\chi(\theta',\mathcal R/\mathfrak q)(m)\le\chi(\theta,M)(m).
\]

\emph{Proof of (b1).} The variables of $\varphi'$ have weights $\pm1$ and $\pm2$. For $n\ge0$ let $\Phi_n$ be the $E$-vector space of the power series all of whose monomials have negative degree at least $n$; then $\Phi_n\Phi_{n'}\subseteq\Phi_{n+n'}$. The weight of a monomial is its positive degree minus its negative degree, so a power series of weight $-j\le0$ lies in $\Phi_j$, and every power series lies in $\Phi_0$.
\begin{enumerate}
\item[(1)] If $f\in A_{-1}\cap\Sigma_x$, then $\varphi'(f)\in\Phi_2$; if $b\in A_0$ and $x(b)=0$, then $\varphi'(b)\in\Phi_1$. Indeed, $cf\in(\ker x)^2$ for some nonzero $c\in\mathcal O$, and the power series $\varphi'(h)$ with $h\in\ker x$ have no constant term, so $\varphi'(f)$ has no terms of degree at most $1$. A monomial of weight $-1$ has negative degree at least $1$, and if its negative degree is $1$ it is one variable of weight $-1$, of degree $1$. A monomial of weight $0$ and negative degree $0$ is constant, and the constant term of $\varphi'(b)$ is $x(b)=0$.
\item[(2)] For $f\in N^m$, $f\in A_{2(m_2-m_1-m_3)}$ and $\varphi'(f)\in\Phi_{2m_1+4m_3}$. Indeed, $\varphi'$ maps $A_{-2}$ into $\Phi_2$, and $F_r^2$ into $\Phi_4$ by (1), as $F_r\in A_{-1}\cap\Sigma_x$; so it maps $N_3$ into $\Phi_4$, and the claim follows by multiplying.
\end{enumerate}
The set $\mathfrak q$ is an additive group, and it is stable under multiplication by $\mathcal R$: the coefficient at $T^m$ of $hp$ is a sum of products $h_{m'}p_{m''}$ with $m'+m''=m$, whose images lie in $\Phi_{2m'_1+4m'_3}\Phi_{2m''_1+4m''_3+1}\subseteq\Phi_{2m_1+4m_3+1}$ by (2). As the condition is imposed on each coefficient, $\mathfrak q$ is generated by homogeneous elements. It contains $\mathcal J$, and the $b\in A_0$ with $x(b)=0$ by (1). As each $\Phi_n$ is a vector space over $E$, $cp\in\mathfrak q$ with $c\in\mathcal O$ nonzero implies $p\in\mathfrak q$. Finally, for $f,g\in A_{-1}$ with $f-g\in\Sigma_x$, $\varphi'(f^2-g^2)=\varphi'(f-g)\varphi'(f+g)$ lies in $\Phi_2\Phi_1\subseteq\Phi_3$, by (1) and as $f+g$ has weight $-1$; so $f^2T_1-g^2T_1\in\mathfrak q$.

\emph{Proof of (b2).} Let $m\ge0$, let $I_m$ be the set of the monomials with negative degree $2m_1+4m_3$ and positive degree $2m_2+2m_3$, which is finite as no variable has weight $0$, and let $\lambda:N^m\to E^{I_m}$ map $f$ to the coefficients of $\varphi'(f)$ at the monomials of $I_m$.
\begin{enumerate}
\item[(1)] For $f\in N^m$, $\varphi'(f)\in\Phi_{2m_1+4m_3+1}$ if and only if $\lambda(f)=0$. Indeed, by (2) of the proof of (b1), $\varphi'(f)$ lies in $\Phi_{2m_1+4m_3}$ and has weight $2(m_2-m_1-m_3)$, so its monomials of negative degree $2m_1+4m_3$ have positive degree $2m_2+2m_3$.
\item[(2)] The piece of degree $m$ of $\mathcal R/(\mathfrak q+\varpi\mathcal R)$ is the image of $N^mT^m$, and since the condition defining $\mathfrak q$ is imposed on each coefficient, the image of $fT^m$ vanishes if and only if $\varphi'(f-\varpi g)\in\Phi_{2m_1+4m_3+1}$ for some $g\in N^m$. By (1) the piece is isomorphic to $\lambda(N^m)/\varpi\lambda(N^m)$.
\item[(3)] For $b\in A_0$ and $f\in N^m$, $\lambda(bf)=x(b)\lambda(f)$, as $\varphi'(b-x(b))\in\Phi_1$ by (1) of the proof of (b1). So $A_0$ acts on $\lambda(N^m)$ through $x$, and lengths over $A_0$ and over $\mathcal O$ agree, $A_0$ having the residue field of $\mathcal O$. By Lemma~\ref{sha-cotangent-finite}(2), $N^m$ is a finitely generated $A_0$-module, so $\lambda(N^m)$ is a finitely generated $\mathcal O$-submodule of $E^{I_m}$, and it is free.
\item[(4)] For each $\nu\in I_m$ some $f\in A_{-2}^{m_1+2m_3}A_2^{m_2+m_3}\subseteq N^m$ has $\lambda(f)$ a nonzero multiple of the unit vector at $\nu$. So $\lambda(N^m)$ has rank $|I_m|$, and by (2) and (3) the piece has length $|I_m|=h(m)$.
\end{enumerate}
For (4), say that a power series $P$ \emph{leads with} the monomial $Y^\nu$ of negative degree $n$ if $P\in\Phi_n$, the coefficient of $P$ at $Y^\nu$ is nonzero, and its coefficients at the other monomials of negative degree $n$ vanish. A product of power series that lead with $Y^\nu$ and $Y^{\nu'}$ leads with $Y^{\nu+\nu'}$.
If $Y_i$ is a variable of weight $\pm1$ and $Y_i$ is the only linear term of $\varphi'(f)$, then the component of $f$ of degree $\pm1$ leads with $Y_i$: a monomial of weight $\pm1$ of the negative degree of $Y_i$ is a variable of weight $\pm1$. Products of two such elements lead with the products of two variables.
If $Y_i$ has weight $\pm2$ and $Y_i$ is the only linear term of $\varphi'(f)$, let $f'$ be the component of $f$ of degree $\pm2$. The monomials of weight $\pm2$ of the negative degree of $Y_i$ are the variables of weight $\pm2$ and the products $Y_kY_l$ of two variables of weight $\pm1$; subtracting from $f'$ suitable multiples of the products leading with these $Y_kY_l$ gives an element of $A_{\pm2}$ that leads with $Y_i$.
Finally, a monomial $Y^\nu$ with $\nu\in I_m$ is a product of $m_1+2m_3$ monomials of negative degree $2$ and positive degree $0$ and $m_2+m_3$ monomials of negative degree $0$ and positive degree $2$, as every variable has negative or positive degree $1$ or $2$; each of them is a variable of weight $\pm2$ or a product of two variables of weight $\pm1$. The product of the elements leading with them lies in $A_{-2}^{m_1+2m_3}A_2^{m_2+m_3}$ and leads with $Y^\nu$, which gives (4).

\emph{Proof of (b3).} Call the pair of the negative and positive degrees of a monomial its bidegree; a variable $Y_i$ of weight $-1$, $-2$, $1$ or $2$ has bidegree $\gamma_i=(1,0)$, $(2,0)$, $(0,1)$ or $(0,2)$. For a set $X$ of variables of $\varphi'$ and $p\in\mathbf Z^2$, let $H_X(p)$ be the number of monomials in the variables of $X$ with bidegree $p$, which is finite as no variable has weight $0$, and let $H=H_{\tau'}$. For $v\in\mathbf Z^2$ and $g:\mathbf Z^2\to\mathbf Z$ let $(\rho_vg)(p)=g(p-v)$, $D_v=1-\rho_v$ and $D^+_v=1+\rho_v$; these operators commute, and $\rho_v\rho_{v'}=\rho_{v+v'}$.
Let $\omega(m)=(2m_1+4m_3,2m_2+2m_3)$, so that $h=H\circ\omega$, with $\omega(e_1)=(2,0)$, $\omega(e_2)=(0,2)$ and $\omega(e_3)=(4,2)$. As $\omega$ is additive, induction on the length of $\delta$ gives $\Delta_\delta h=(D_{\omega(\delta_1)}\cdots D_{\omega(\delta_p)}H)\circ\omega$; and both coordinates of $\omega(m)$ are large when $m$ is large.
\begin{enumerate}
\item[(1)] \emph{Removing a variable.} If $Y_i\notin X$, then $D_{\gamma_i}H_{X\cup\{Y_i\}}=H_X$: a monomial in the variables of $X\cup\{Y_i\}$ with bidegree $p$ either does not involve $Y_i$, or it is $Y_i$ times such a monomial with bidegree $p-\gamma_i$. So the product of the $D_{\gamma_i}$ over the variables outside $X$ maps $H$ to $H_X$.
\item[(2)] \emph{Two variables.} Let $Y_{i_0}$, $Y_{i_1}$ and $Y_{i_2}$ be variables of weights $-1$, $1$ and $2$, which exist as $r\ge1$. For $p_1,p_2\ge0$, $H_{\{Y_{i_0},Y_{i_1}\}}(p)=1$, the only monomial being $Y_{i_0}^{p_1}Y_{i_1}^{p_2}$. For $p_1\ge1$, $H_{\{Y_{i_1},Y_{i_2}\}}(p)=0$, as the monomials in $Y_{i_1}$ and $Y_{i_2}$ have negative degree $0$.
\item[(3)] \emph{The operator.} Since $\rho_{(2,0)}=\rho_{(1,0)}^2$, $\rho_{(0,2)}=\rho_{(0,1)}^2$ and $\rho_{(4,2)}=\rho_{(0,2)}\rho_{(2,0)}^2$,
\[
 D_{(2,0)}=D^+_{(1,0)}D_{(1,0)},\qquad D_{(0,2)}=D^+_{(0,1)}D_{(0,1)},\qquad D_{(4,2)}=D_{(0,2)}+\rho_{(0,2)}D^+_{(2,0)}D^+_{(1,0)}D_{(1,0)}.
\]
With $C=D_{(1,0)}^{r-1}D_{(2,0)}^{s}D_{(0,1)}^{r'}D_{(0,2)}^{s'}$, the operator $D_{\omega(\delta_1)}\cdots D_{\omega(\delta_p)}$ of the list $\delta$ of (b3) is therefore
\[
 D_{(2,0)}^{r-1+s}D_{(0,2)}^{r'+s'}D_{(4,2)}=(D^+_{(1,0)})^{r-1}(D^+_{(0,1)})^{r'}\big(CD_{(0,2)}+\rho_{(0,2)}D^+_{(2,0)}D^+_{(1,0)}\,CD_{(1,0)}\big).
\]
\item[(4)] \emph{Evaluation.} As $\varphi'$ has $r$ variables of weight $-1$, $s$ of weight $-2$, $r'+1$ of weight $1$ and $s'+1$ of weight $2$, $CD_{(0,2)}$ is the product of the $D_{\gamma_i}$ over all the variables except $Y_{i_0}$ and $Y_{i_1}$, and $CD_{(1,0)}$ is the product over all the variables except $Y_{i_1}$ and $Y_{i_2}$. By (1) and (2), at every $p$ with large coordinates, $CD_{(0,2)}H$ takes the value $1$, and $CD_{(1,0)}H$ and its image under $\rho_{(0,2)}D^+_{(2,0)}D^+_{(1,0)}$ take the value $0$. If $g(p)=a$ at every $p$ with large coordinates, then $(D^+_vg)(p)=2a$ at every such $p$. So $\Delta_\delta h(m)=2^{r-1+r'}$ for large $m$.
\end{enumerate}
% (b3) is num_holds, through the private NumericalFunctions of FrobeniusInclusionXEqMulQ.lean, which count the monomials of each bidegree directly with the shifts x(v) of the group algebra of Z^2: dif_eq (the differences of h as operators on the count N), N_insert and N_prod for (1), N_one and N_zero for (2), the identities hX2, hY2, h42 and hop of key for (3), and ev_pow with eventually_const_of_shift for (4).

\emph{Proof of (b4).} For a valid list $u'$ with degrees $\delta'$ let $\Delta_{\delta'}h_\beta$ be as in (b3); it does not depend on the order of $\delta'$.
\begin{enumerate}
\item[(1)] \emph{Peeling.} Let $c\ge0$, $y\in R_d$ and $u'$ a valid list such that $((d,\beta^cy))u'$ is admitted for the annihilator of $F$. Then $\chi(((d,\beta^cy))u',F)(m)=c\,\Delta_{\delta'}h_\beta(m)+\chi(((d,y))u',F)(m)$ for large $m$. By induction on $c$: the lists $((d,\beta^{c-1}y))u'$ and $((d,y))u'$ are admitted, as their ideals contain that of $((d,\beta^cy))u'$. By Lemma~\ref{sha-cotangent-calculus}(8), with $s=\beta$ of degree $0$ and $s'=\beta^{c-1}y$, the left side is $\chi(((0,\beta))u',F)(m)+\chi(((d,\beta^{c-1}y))u',F)(m)$ for large $m$. By (5) of that lemma, as $\beta$ is a nonzerodivisor on $F$, $\chi(((0,\beta))u',F)=\chi(u',F/\beta F)$, and by (9), as the empty list is admitted for the annihilator of $F/\beta F$, this is the alternating sum $\sum_I(-1)^{|I|}h_\beta(m-\sum_{i\in I}\delta'_i)=\Delta_{\delta'}h_\beta(m)$ for large $m$.
\item[(2)] \emph{Induction on $n$.} For $n=0$ the claim is $\chi(u,F)(m)\ge0$ for large $m$, which is Lemma~\ref{sha-cotangent-calculus}(6). For $n\ge1$, let $L$ be the list of the $(d,\beta^{c_i}y_i)$ with $i\ge2$. By (1), for large $m$, $\chi(u_\beta,F)(m)=c_1\Delta_\delta h_\beta(m)+\chi(((d,y_1))Lu,F)(m)$, since the degrees of $Lu$ are $n-1$ copies of $d$ followed by those of $u$. By Lemma~\ref{sha-cotangent-calculus}(4), $\chi(((d,y_1))Lu,F)(m)=\chi(L((d,y_1))u,F)(m)$ for large $m$, and by induction, applied to $y_2,\dots,y_n$ and the valid list $((d,y_1))u$, this is at least $(c_2+\dots+c_n)\Delta_\delta h_\beta(m)$, the degrees being again $n-1$ copies of $d$ and those of $u$ in another order.
\end{enumerate}

\emph{Proof of (d).} Let $M'=\mathcal R/(\mathcal J+(\theta))$, a graded module by Lemma~\ref{sha-cotangent-calculus}(1), whose pieces $M'_m$ are finitely generated $A_0$-modules. Since $a'$ has degree $0$, the definition gives $\chi(a',M')(m)=\theta(a',M'_m)$, in the notation of Lemma~\ref{sha-cotangent-herbrand}. For $\sigma:A_0\to\Omega$ let $\ell_\sigma$ be as there, let $I_2$ be the ideal of $A$ generated by the $F_i^2$, $G_i$, $E_i^2$, $E'_i$ and $Z_i$, and for an ideal $I'$ of $A$ and $w\in\mathbf Z$ let $A_w(I')$ be the image of $A_w$ in $A/I'$, an $A_0$-module. The proof uses four facts.
\begin{enumerate}
\item[(d1)] For large $m$, $A_0/\operatorname{Ann}(M'_m)$ is finite over $\mathcal O$, and $\ell_\sigma(M'_m)=0$ for every $\sigma$ with $\mathfrak d\subseteq\ker\sigma$ and for every $\sigma$ that does not kill $A_0\cap I$. This uses conditions (1) to (3).
\item[(d2)] For $m\ge(1,1,1)$ and $\sigma$ with $\mathfrak d\not\subseteq\ker\sigma$, $\ell_\sigma(M'_m)\le\ell_\sigma(A_w(I_2))$ with $w=2(m_2-m_1-m_3)$.
\item[(d3)] If $\ell_\sigma(A_w(I))\le1$ for all $w$, then $\ell_\sigma(A_w(I_2))\le2^{r+1+r'}$ for all $w$.
\item[(d4)] For $\sigma$ with $\mathfrak d\not\subseteq\ker\sigma$, $\ell_\sigma(A_w(I))\le1$ for all $w$. This uses conditions (2) and (4).
\end{enumerate}
A homomorphism $\xi:A\to\Omega$ is two-signed if and only if $\mathfrak d$ is not contained in the kernel of its restriction to $A_0$: by Lemma~\ref{sha-cotangent-finite}(3), $\xi$ is two-signed if and only if $\xi(f)\ne0$ and $\xi(g)\ne0$ for some $f\in A_{-2}$ and $g\in A_2$.
For $\xi\in Y$ and $b\in A_0$ with $\xi(b)\ne0$, $v(\xi(b))\ge0$: the $\mathcal O$-module $\xi(A_0)$ is an image of the image of $A_0$ in $A/I$, which is finitely generated by condition (2), and it contains $\mathcal O[\xi(b)]$, so $\xi(b)$ is integral over $\mathcal O$; if $v(\xi(b))<0$, then in an equation $\xi(b)^n+\sum_{i<n}a_i\xi(b)^i=0$ with $a_i\in\mathcal O$ the term $\xi(b)^n$ would have a smaller value of $v$ than all the others.
Let $S$ be the set of the restrictions to $A_0$ of the elements of $T$, and let $m$ be large for (d1), with $m\ge(1,1,1)$. For $\sigma\in S$, the restriction of a two-signed $\xi\in T$: $\sigma(a')=\xi(a')\ne0$ by condition (5) for $a'$, $v(\sigma(a'))\ge0$, and $\ell_\sigma(M'_m)\le2^{r+1+r'}$ by (d2), (d3) and (d4). For $\sigma\notin S$, $\ell_\sigma(M'_m)=0$: otherwise, by (d1), $\mathfrak d\not\subseteq\ker\sigma$ and $\sigma$ kills $A_0\cap I$; by (a3), $\sigma$ is the restriction of some $\xi\in Y$, which is two-signed; by the covering property, $T$ contains some $t\cdot\xi$, whose restriction to $A_0$ is that of $\xi$, so $\sigma\in S$.
Lemma~\ref{sha-cotangent-herbrand}, with $B=A_0$, $M=M'_m$ and $c=2^{r+1+r'}$, gives $\chi(a',M')(m)\le2^{r+1+r'}\sum_{\sigma\in S}v(\sigma(a'))\le2^{r+1+r'}\sum_{\xi\in T}v(\xi(a'))$, since each $\sigma\in S$ is the restriction of some $\xi\in T$ and all the terms are nonnegative.
\emph{Proof of (d1).} Let $\mathcal J'=\mathcal J+(\theta)$, so that $M'=\mathcal R/\mathcal J'$, and let $\mathfrak a=A_0\cap I$, the kernel of $A_0\to A/I$.
\begin{enumerate}
\item[(1)] Let $\mathfrak b,\mathfrak b'$ be ideals of $A_0$ such that every prime $\mathfrak p$ of $\mathcal R$ that contains $\mathcal J'+\mathfrak b'\mathcal R$ but not $\operatorname{Irr}$ contains $\mathfrak b$. Then for some $n$, $\mathfrak b^nM'_m\subseteq\mathfrak b'M'_m$ for large $m$. Indeed, $\mathfrak b\mathcal R\cdot\operatorname{Irr}$ lies in every prime containing $\mathcal J'+\mathfrak b'\mathcal R$, so, $\mathcal R$ being noetherian, $(\mathfrak b\mathcal R\cdot\operatorname{Irr})^n\subseteq\mathcal J'+\mathfrak b'\mathcal R$ for some $n$. As $\mathcal R$ is standard, $M'_m\subseteq\operatorname{Irr}^nM'$ for large $m$, so $\mathfrak b^nM'_m\subseteq\mathfrak b'M'$, and taking components of degree $m$ gives the claim.
\item[(2)] Step (a1) of the proof of (a) uses of $\mathfrak p$ only that it contains $\theta$ and the $Z_k$ and not $\operatorname{Irr}$; so every prime containing $\mathcal J'$ but not $\operatorname{Irr}$ contains $\mathfrak a$, and by (1), $\mathfrak a^n$ kills $M'_m$ for large $m$. The ring $A_0/\mathfrak a$ is finite over $\mathcal O$ by condition (2), and each $\mathfrak a^i/\mathfrak a^{i+1}$ is a finitely generated $A_0/\mathfrak a$-module, $A_0$ being noetherian; so $A_0/\mathfrak a^n$ is finite over $\mathcal O$, and so is its quotient $A_0/\operatorname{Ann}(M'_m)$. If $\sigma$ does not kill $\mathfrak a$, some $b\in\mathfrak a$ with $\sigma(b)\ne0$ has $b^nM'_m=0$, so the localization of $M'_m$ at $\ker\sigma$ is $0$.
\item[(3)] Steps (a2) to (a9), with (a4) replaced by the hypothesis $\mathfrak d\mathcal R\subseteq\mathfrak p$, show that every prime $\mathfrak p$ that contains $\mathcal J'+\mathfrak d\mathcal R$ but not $\operatorname{Irr}$ contains $\mathfrak m_{A_0}$; they use conditions (1) to (3). By (1), $\mathfrak m_{A_0}^nM'_m\subseteq\mathfrak dM'_m$ for large $m$. Let $\sigma$ kill $\mathfrak d$. If the localization of $M'_m$ at $\ker\sigma$ were nonzero, so would be that of $M'_m/\mathfrak dM'_m$, by Nakayama's lemma, as $\mathfrak d\subseteq\ker\sigma$; then the annihilator of $M'_m/\mathfrak dM'_m$, which contains $\varpi^n$, would lie in $\ker\sigma$. But $\sigma(\varpi)\ne0$.
\end{enumerate}

\emph{Proof of (d2).} Let $m\ge(1,1,1)$ and $w=2(m_2-m_1-m_3)$, and choose $f\in A_{-2}$ and $g\in A_2$ with $\sigma(fg)\ne0$; let $b=fg\in A_0$.
\begin{enumerate}
\item[(1)] For $\nu\in\mathbf N^3$ and $a\in A_{2(\nu_2-\nu_1-\nu_3)}$, some $b^na$ lies in $N^\nu$. Indeed, $N^\nu$ contains $A_{-2}^{\nu_1+2\nu_3}A_2^{\nu_2+\nu_3}$. If $\nu_1+2\nu_3\le\nu_2+\nu_3$, then $a\in A_{2c}=A_2^c$ with $c=\nu_2-\nu_1-\nu_3$, by Lemma~\ref{sha-cotangent-finite}(3), and $b^{\nu_1+2\nu_3}a\in A_{-2}^{\nu_1+2\nu_3}A_2^{\nu_2+\nu_3}$; the other case is symmetric.
\item[(2)] The piece $M'_m$ is the image of $N^mT^m$, and $N^m\subseteq A_w$. Let $a\in N^m\cap I_2$, and write $a=\sum_kc_ky_k$ with $c_k\in A$ and $y_k$ among the generators $F_i^2$, $G_i$, $E_i^2$, $E'_i$, $Z_i$ of $I_2$. Each $y_k$ lies in $N^{\nu_k}$ with $\nu_k=e_1$, $e_2$, $e_3$ or $0$, so that $\nu_k\le m$, and $y_kT^{\nu_k}$ lies in $\mathcal J'$: it is an element of $\theta$, or $F_r^2T_3$, or $Z_i\in\mathcal J$. Taking components of degree $w$, $a=\sum_k(c_k)_{w_k}y_k$ with $w_k=w-\deg y_k$, the degree of $N^{m-\nu_k}$. By (1) there is $n$ with $b^n(c_k)_{w_k}\in N^{m-\nu_k}$ for all $k$, and then $b^naT^m=\sum_k\big(b^n(c_k)_{w_k}T^{m-\nu_k}\big)(y_kT^{\nu_k})\in\mathcal J'$.
\item[(3)] So the surjection $N^m\to M'_m$ and the map $N^m\to A_w(I_2)$ have the property: an element of the kernel of the second is mapped by the first to an element killed by $b^n$, with $\sigma(b^n)\ne0$. After localization at $\ker\sigma$, $M'_m$ becomes a quotient of the image of $N^m$ in the localization of $A_w(I_2)$, so $\ell_\sigma(M'_m)\le\ell_\sigma(A_w(I_2))$.
\end{enumerate}

\emph{Proof of (d3).} Let $J'$ be an ideal of $A$ generated by homogeneous elements, $t\in A_d$ and $w\in\mathbf Z$. Then
\[
 \ell_\sigma(A_w(J'+t^2A))\le\ell_\sigma(A_{w-d}(J'+tA))+\ell_\sigma(A_w(J'+tA)).
\]
Indeed, $A_w(J'+t^2A)$ maps onto $A_w(J'+tA)$, and its kernel is the image of $A_w\cap(J'+tA)$, which is the image of $tA_{w-d}$, as $J'$ is generated by homogeneous elements; multiplication by $t$ maps $A_{w-d}(J'+tA)$ onto this kernel, since $t(J'+tA)\subseteq J'+t^2A$. Localization is exact, and lengths are additive on exact sequences.
So if $\ell_\sigma(A_w(J'+tA))\le c$ for all $w$, then $\ell_\sigma(A_w(J'+t^2A))\le2c$ for all $w$. Starting from $I$, with the bound $1$, and replacing $F_0,\dots,F_r$ and then $E_1,\dots,E_{r'}$ one at a time by their squares, the other generators being homogeneous, gives the bound $2^{r+1+r'}$ for $I_2$.

\emph{Proof of (d4).} Let $\sigma$ be as there, and choose $f\in A_{-2}$ and $g\in A_2$ with $\sigma(fg)\ne0$; let $t=fg\in A_0$ and $P=\ker\sigma$.
\begin{enumerate}
\item[(1)] If $p\in A_0$ and $\sigma(p)=0$, then $bp^n\in I$ for some $b\in A_0$ with $\sigma(b)\ne0$ and some $n$. Indeed, the image of $A_0$ in $A/I$ is finite over $\mathcal O$ by condition (2), so the image of $p$ is a root of a monic polynomial $q$ over $\mathcal O$. Write $q=X^nq_1$ with $q_1(0)\ne0$, and $b=q_1(p)$; then $bp^n\in I$, and $\sigma(b)=q_1(0)\ne0$, as $\mathcal O\to\Omega$ is injective.
\item[(2)] If $h\in A_j$, $b\in A_0$, $\sigma(b)\ne0$ and $bh^n\in I$, then $b'h\in I$ for some $b'\in A_0$ with $\sigma(b')\ne0$. Indeed, every $\xi\in Y$ with $\xi(f)=1$ and $\xi(g)\ne0$ has $\xi(b)\xi(h)^n=0$, so it vanishes at $bh$. By the second statement of Lemma~\ref{sha-cotangent-finite}(4), whose hypotheses are conditions (2) and (4), $cg^Nbh\in I+(f-1)A$ for some $N$ and some nonzero $c\in\mathcal O$. This element is homogeneous, of degree $2N+j$. If a homogeneous $h'$ lies in $I+(f-1)A$, say $h'=i+u(f-1)$, then comparing components gives $u_{k}\equiv fu_{k+2}$ modulo $I$ in every degree $k\ne\deg h'$ and $u_{\deg h'}\equiv fu_{\deg h'+2}-h'$; as only finitely many $u_k$ are nonzero, $u_k\in I$ for $k>\deg h'$, then $u_{\deg h'-2n}\equiv-f^nh'$ for $n\ge0$, and $f^nh'\in I$ for large $n$. So $f^Mcg^Nbh\in I$ for some $M$, and $b'=ct^{N+M}b$ satisfies $b'h=(f^Ng^M)(f^Mcg^Nbh)\in I$ and $\sigma(b')=c\,\sigma(t)^{N+M}\sigma(b)\ne0$.
\item[(3)] There is $h_0\in A_w$ such that for every $h\in A_w$ some $b\in A_0$ with $\sigma(b)\ne0$ and some $a\in A_0$ have $bh-ah_0\in I$. Let $e=f^w$ and $e'=g^w$ if $w\ge0$, and $e=g^{-w}$ and $e'=f^{-w}$ if $w<0$, so that $e\in A_{-2w}$ and $ee'=t^{|w|}$. If $\sigma(h_0h_0e)\ne0$ for some $h_0\in A_w$, then $b=h_0h_0e$ and $a=hh_0e$ satisfy $bh-ah_0=0$. Otherwise take $h_0=0$ and $a=0$: for $h\in A_w$, (1) gives $b$ with $\sigma(b)\ne0$ and $b(hhe)^n\in I$, so $bt^{|w|n}h^{2n}=e'^nb(hhe)^n\in I$, and (2) gives $b'h\in I$ with $\sigma(b')\ne0$.
\item[(4)] Every $p\in P$ kills the localization of $A_w(I)$ at $P$: for $h\in A_w$, (1) gives $bp^n\in I$ with $\sigma(b)\ne0$, so $b(ph)^n\in I$, and (2), applied to $ph\in A_w$, gives $b'ph\in I$ with $\sigma(b')\ne0$. By (3), this localization is generated by the image of $h_0$. So it is a quotient of the residue field of $(A_0)_P$, and $\ell_\sigma(A_w(I))\le1$.
\end{enumerate}

By Steps 2, 4 and 6, $\lambda_{-1}(F)=c_1+\dots+c_r\le2\sum_{\xi\in T}v(\xi(a))$ for every such $T$, and Step 5 gives one.
% The proof is Alg.M.main_holds. Step 1 is rk_holds (LinPartApply: rk_fin with fg_holds, linPart_sat, linPart_wt, linPart_proj); Step 2 is sm_holds (sm_smith and sm_lift) with main_span and main_cotSpan; Step 3 is cut_holds (lin_span, then ift_holds, which is Lemma sha-cotangent-implicit; IsChart0, weights mDims r s r' s' 0); Step 4 is nz_holds (nz_comp, nz_coeff_mul); Step 5 is fin_holds (fin_gen, fin_act, fin_onS), with fg_holds, sign_holds and slice_holds, which are Lemma sha-cotangent-finite. hom_holds gives the homogeneity of I.
% Step 6 is core_holds (AlgMemCoef): mR Q = rees (mN Q.F_last) with grading mGr Q (a Calc.HidaSetting 3 by std_holds), J = coefIdeal (ker phi'), theta = ts Q (tsF ++ tsR), chi = Calc.chi, (theta) = Calc.secs. (a) is admissibility_holds (Calc.admitted for ts ++ [(0, a')]); (b) is low_holds, applied to the admissibility of J + (a') obtained from (a) by E2.admissible_tail_of_cons_mem and E2.secs_perm; (c) is the private chain (ruleNzd, euler_rule_two_of_graded_finite_iso, ruleUB, qi_holds), with coef_hom and the nonzerodivisor hn; (d) is up_holds, with hom_sup. The final division is h4, h5.
% (a) is admissibility_holds: Ideal.exists_pow_le_of_le_radical_of_fg with irr_fg; (a1) is zi_holds (exists_sec_notMem, sec_mul_sec, algebraMap_mul_sec_mul_sec_eq_of_mul_eq, hgen); (a2) is emb_holds (prime_maximal_or_kernel); (a3) is tsp_holds (ns_holds with J = I + (ker sigma) A, through ft_holds); (a4) is the last part of admissibility_holds (tsp_holds, then M5); (a5) is fix_holds (kill, kill1); (a6) is zm_holds (zm_iota, zm_abs, zm_of_hom, zm_f, oz_two_mul_eq_pow, ns_holds); (a7) is zp_holds (InP, idJ, exists_not_inP, inP_neg, neg_mem_idJ, two_mem_idJ, one_mem_idJ, zero_mem_idJ, not_mem_idJ); (a8) is ct_holds (kill_sq, kill_part, sign_holds); (a9) is zii_holds with M3.
% (c) is the private chain of AlgMemCoef.lean: ruleNzd (h1), euler_rule_two_of_graded_finite_iso (h2), ruleUB (h3), qi_holds with ruleIso; the nonzerodivisor is hn (coeff_alg).
% (b) is low_holds: (b1) is lk_holds (leadKer, fil, negW); (b2) is hf_holds (cnt, posW); (b3) is num_holds; (b4) is la_holds (LaDifPerm.lean, with dif and hilb). The comparison with M is mono_quo (exists_surjective_graded_quot_to_quo, ruleMono), the congruence is Cong.chi_cong with h𝔮sq and hkey, and the admissibility for R/(q + pi R) is hadm0.
% (d) is up_holds: (d1) is eventually_mQ_gr_finite_annihilator; (d2) is locLen_mQ_gr_le_locLen_aPiece; (d3) is locLen_aPiece_ideal2_le_two_pow; (d4) is locLen_quotientWeightPiece_le_one; the nonnegativity of v is valuation_nonneg_apply_degreeZero_of_onS (int_holds); two-signed points are hD (sign_holds); chi of one element of degree 0 is Calc.chi_single0; the bound is th_holds (Lemma sha-cotangent-herbrand) with the restrictions res0 of T, and Finset.sum_image_le_of_nonneg.
% The facts: (b1) is lk_holds (satfil_holds is its (1), latfil_holds its (2), fil_holds the product of the Phi_n); (b2) is hf_holds (rees_graded_piece_surjection, the coefficient map lc, length_localization_bound, and exists_mem_oz_pow_coeff_ne_zero_unique of ChartLeadingWeightTwo.lean: Lead, sp_lin, chart_leading_weight_one, _sum, _two with sp_elim, sp_split, sp_main); (b3) is num_holds; (b4) is la_holds (la_peel through euler_rule_product_zero, ruleNzd and ruleSmall, then euler_rule_two_of_graded_finite_iso and la_dif_perm). (d1) is eventually_mQ_gr_finite_annihilator (power_ideal_kills_graded_pieces with kill_holds; zi_holds, which is (a1); zii_holds, which is (a2)-(a9) with D in p; finite_quotient_of_kernel_pow; localized_length_zero); (d2) is locLen_mQ_gr_le_locLen_aPiece (power_mul_mem_lattice, generated_component_mem, localized_length_le_of_surjective); (d3) is locLen_aPiece_ideal2_le_two_pow (locLen_aPiece_sup_span_sq_le_add, applied to the F_i and then the E_i); (d4) is locLen_quotientWeightPiece_le_one (nilpotent_modulo_kernel, red through slice_holds and deh_holds, cyc, localized_length_le_one).
\end{proof}
\leanref{Development/AutomorphicArithmetic/AutomorphicVarieties/CommutativeAlgebra/OrbitVanishingLocal.lean}{UpperTheory.V.Alg.LemmaM}
\leanref{Development/AutomorphicArithmetic/AutomorphicVarieties/CommutativeAlgebra/OrbitVanishingLocal.lean}{UpperTheory.V.Alg.IsWRing}
\leanref{Development/AutomorphicArithmetic/AutomorphicVarieties/CommutativeAlgebra/OrbitVanishingLocal.lean}{UpperTheory.V.Alg.Eqns}
\leanref{Development/AutomorphicArithmetic/AutomorphicVarieties/CommutativeAlgebra/OrbitVanishingLocal.lean}{UpperTheory.V.Alg.M1}
\leanref{Development/AutomorphicArithmetic/AutomorphicVarieties/CommutativeAlgebra/OrbitVanishingLocal.lean}{UpperTheory.V.Alg.M3}
\leanref{Development/AutomorphicArithmetic/AutomorphicVarieties/CommutativeAlgebra/OrbitVanishingLocal.lean}{UpperTheory.V.Alg.M4}
\leanref{Development/AutomorphicArithmetic/AutomorphicVarieties/CommutativeAlgebra/OrbitVanishingLocal.lean}{UpperTheory.V.Alg.M5}
\leanref{Development/AutomorphicArithmetic/AutomorphicVarieties/CommutativeAlgebra/OrbitVanishingLocal.lean}{UpperTheory.V.Alg.Covers}
\leanref{Development/AutomorphicArithmetic/AutomorphicVarieties/CommutativeAlgebra/GradedModulesLength/MainCotSpan.lean}{UpperTheory.V.Alg.M.main_holds}
\leanref{Development/AutomorphicArithmetic/AutomorphicVarieties/CommutativeAlgebra/GradedModulesLength/LkHolds.lean}{UpperTheory.V.Alg.M.lk_holds}
\leanref{Development/AutomorphicArithmetic/AutomorphicVarieties/CommutativeAlgebra/GradedModulesLength/ReesGradedPieceSurjection.lean}{UpperTheory.V.Alg.M.hf_holds}
\leanref{Development/AutomorphicArithmetic/AutomorphicVarieties/CommutativeAlgebra/GradedModulesLength/ChartLeadingWeightTwo.lean}{UpperTheory.V.Alg.M.exists_mem_oz_pow_coeff_ne_zero_unique}
\leanref{Development/AutomorphicArithmetic/AutomorphicVarieties/CommutativeAlgebra/GradedModulesLength/FrobeniusInclusionXEqMulQ.lean}{UpperTheory.V.Alg.M.num_holds}
\leanref{Development/AutomorphicArithmetic/AutomorphicVarieties/CommutativeAlgebra/GradedModulesLength/LaDifPerm.lean}{UpperTheory.V.Alg.M.la_holds}
\leanref{Development/AutomorphicArithmetic/AutomorphicVarieties/CommutativeAlgebra/GradedModulesLength/EventuallyMQGrFiniteAnnihilator.lean}{UpperTheory.V.Alg.M.eventually_mQ_gr_finite_annihilator}
\leanref{Development/AutomorphicArithmetic/AutomorphicVarieties/CommutativeAlgebra/GradedModulesLength/LocalizedLengthLeOfSurjective.lean}{UpperTheory.V.Alg.M.locLen_mQ_gr_le_locLen_aPiece}
\leanref{Development/AutomorphicArithmetic/AutomorphicVarieties/CommutativeAlgebra/GradedModulesLength/LocLenAPieceIdeal2LeTwoPow.lean}{UpperTheory.V.Alg.M.locLen_aPiece_ideal2_le_two_pow}
\leanref{Development/AutomorphicArithmetic/AutomorphicVarieties/CommutativeAlgebra/GradedModulesLength/LocLenQuotientWeightPieceLeOne.lean}{UpperTheory.V.Alg.M.locLen_quotientWeightPiece_le_one}

\begin{lemma}\label{sha-algebra-derivations}
Let $\mathcal O$ be a discrete valuation ring with finite residue field, let $X$ be an abelian group, let $A=\bigoplus_{j\in X}A_j$ be an $X$-graded commutative $\mathcal O$-algebra, let $x:A\to\mathcal O$ be a homomorphism of $\mathcal O$-algebras, and let $j\in X$ with $x(A_j)=0$.
Let $F=(F_i)$ be a family of elements of $A_j$, let $c\in\mathcal O$ be nonzero, and let $\mathcal M$ be an $\mathcal O$-module with an injective additive map $D:\mathcal M\to\operatorname{Hom}_{\mathcal O}(A,\mathcal O/c)$ such that, for every $\mu\in\mathcal M$, $D(\mu)$ vanishes on the $F_i$ and is the reduction modulo $c$ of a derivation at $x$ with values in $\mathcal O$ of weight $j$ (Definition~\ref{sha-algebra-graded}).
Then $\operatorname{len}_{\mathcal O}\mathcal M\le\lambda_j(F)$.
\end{lemma}
\begin{proof}
We may assume that $\lambda_j(F)$ is finite.
Let $\mu\in\mathcal M$, and let $D'$ be a derivation at $x$ with values in $\mathcal O$ of weight $j$ that reduces to $D(\mu)$.
It vanishes on $(\ker x)^2$; if $cf\in(\ker x)^2$ with $c\ne0$, then $cD'(f)=0$, so $D'(f)=0$.
Hence $D(\mu)$ vanishes on $\Sigma_x$ and on the $F_i$, and its restriction to $L_j$ gives an $\mathcal O$-linear map $\psi(\mu):P\to\mathcal O/c$ on $P=L_j/(L_j\cap\langle F\rangle)$.
The map $\mu\mapsto\psi(\mu)$ is injective: if $\psi(\mu)=\psi(\nu)$, then $D(\mu)$ and $D(\nu)$ agree on $A_j=A_j\cap\ker x$ and both vanish on $A_i$ for $i\ne j$, so $D(\mu)=D(\nu)$ and $\mu=\nu$.
Since $\operatorname{Hom}_{\mathcal O}(\mathcal O/\mathfrak m_{\mathcal O},\mathcal O/c)$ has length at most $1$, induction on the length of $P$ shows that $\operatorname{Hom}_{\mathcal O}(P,\mathcal O/c)$ has length at most $\operatorname{len}P=\lambda_j(F)$.
A module of finite length $\lambda$ over $\mathcal O$ has $|k|^\lambda$ elements, where $k$ is the finite residue field, so the injective map $\psi$ from $\mathcal M$ gives $\operatorname{len}\mathcal M\le\lambda_j(F)$.
\end{proof}
\leanref{Development/AutomorphicArithmetic/AutomorphicVarieties/CommutativeAlgebra/OrbitVanishingLocal.lean}{UpperTheory.V.Alg.LemmaM'}
\leanref{Development/AutomorphicArithmetic/AutomorphicVarieties/CommutativeAlgebra/GradedModulesLength/Mprime.lean}{UpperTheory.V.Alg.M.mprime_holds}

\subsection{Selmer groups and rigidity}\label{subsection-sha-rigidity}

\begin{theorem}[Kato \cite{Kat04}]\label{sha-kato-selmer}
Let $\ell$ be an odd prime and $k\ge1$ an integer with $\ell-1\nmid2k$; let $K$ be an imaginary quadratic field with $\ell\nmid d_Kh_K$, and $\eta$ its quadratic character; let $N\ge1$ be even and prime to $d_K$.
Let $\psi:\operatorname{Gal}(\bar{\mathbf Q}/\mathbf Q)\to\bar{\mathbf Q}_\ell^\times$ be a character, trivial on $\operatorname{Gal}(\bar{\mathbf Q}/\mathbf Q(\mu_N))$, of finite order prime to $\ell$, with $\psi^2$ trivial on the inertia groups at the places above $\ell$, and with $\chi_{\mathrm{cyc}}(c)^k\psi(c)=1$ for some complex conjugation $c$; put $\Psi=\chi_{\mathrm{cyc}}^k\psi$. Let $\iota:\bar{\mathbf Q}_\ell\cong\mathbf C$.
For an Eisenstein datum $(\ell,k,\psi,K)$ with $\psi$ modulo $N$ and deformation data $(\eta,c,\Psi)$ (Definition~\ref{sha-ordinary-deformation}) these conditions hold.
Let $W:\operatorname{Gal}(\bar{\mathbf Q}/\mathbf Q)\to GL_2(\bar{\mathbf Q}_\ell)$ be a continuous representation such that:
\begin{enumerate}
\item $\det W=\eta$, and $W$ has level $N|d_K|$ and weights $(k,-k)$ (Definition~\ref{sha-eigensystem-level});
\item on the decomposition group at a place above $\ell$ whose inertia group lies in $\ker\eta$, $W$ is conjugate to an upper triangular representation whose diagonal characters equal $\chi_{\mathrm{cyc}}^{k}\chi'$ and $\chi_{\mathrm{cyc}}^{-k}\chi'$ on inertia, for a character $\chi'$ trivial on $\operatorname{Gal}(\bar{\mathbf Q}/\mathbf Q(\mu_\ell))$; for some place above $\ell$ and some $M\ge1$, the characteristic polynomial of $W(s)$ is $(X-\chi_{\mathrm{cyc}}(s)^k)(X-\chi_{\mathrm{cyc}}(s)^{-k})$ for every $s$ in its inertia group that fixes $\mu_M$; and $\operatorname{tr}W=\Psi+\Psi^{-1}$ on the inertia groups at the places above $\ell$;
\item $\operatorname{tr}W\equiv\eta\Psi+\Psi^{-1}$ modulo the maximal ideal of the integers of $\bar{\mathbf Q}_\ell$, and the restriction of $W$ to $\operatorname{Gal}(\bar{\mathbf Q}/K(\mu_{\ell^\infty}))$ is irreducible;
\item $W\otimes\chi_{\mathrm{cyc}}^{-k}$ is attached through $\iota$ to a newform of weight $2k+1$, and the restriction of $W$ to every imaginary quadratic field on which it is irreducible is windowed (Theorem~\ref{sha-residually-reducible-automorphy}).
\end{enumerate}
Then the Bloch-Kato Selmer groups $H^1_f(\mathbf Q,W)$ and $H^1_f(\mathbf Q,W\otimes\eta)$ vanish.
\end{theorem}
\leanref{Development/AutomorphicArithmetic/AutomorphicVarieties/ModularityArgument/KatoQuadraticCharacters/SphericalDecomposition/KatoTwoW.lean}{UpperTheory.V.Z.Kato.WeilSphericalRecords.KatoTwoW}
\leanref{Development/AutomorphicArithmetic/AutomorphicVarieties/ModularityArgument/KatoQuadraticCharacters/KatoTwoOfUseN.lean}{UpperTheory.V.Z.Kato.Used}
\leanref{Development/AutomorphicArithmetic/AutomorphicVarieties/ModularityArgument/ThresholdArguments/DenseThresholdCases/KatoTwoWOfPublishedH.lean}{UpperTheory.V.Z.Kato.WeilSphericalRecords.katoTwoW_of_printsH}

\begin{lemma}\label{sha-rigidity-normalize}
Let $R$ be a complete local ring with maximal ideal $\mathfrak m_R$ and residue field $\kappa$, in which $2$ is a unit, and let $J$ be the antidiagonal $3\times3$ matrix with entries $1$.
\begin{enumerate}
\item Let $a,b,d\in R$ with $b\in R^\times$ and $d\in\mathfrak m_R$. Then $ax^2+bx+d=0$ for some $x\in\mathfrak m_R$.
\item Let $S$ be a symmetric $3\times3$ matrix over $R$ congruent to $J$ modulo $\mathfrak m_R$. Then there is a matrix $h$ over $R$, congruent to $1$, with $hS\,{}^th=J$.
\item Let $a\in R^\times$ and $b,d\in\mathfrak m_R$, and let $Z$ be the matrix with rows $(a,b,0)$, $(d,0,-b)$ and $(0,-d,-a)$. Then there is a matrix $h$ over $R$, congruent to $1$, with $hJ\,{}^thJ=1$, such that the entries $(1,2)$ and $(2,1)$ of $hZJ\,{}^thJ$ vanish.
\item Let $G$ be a group, let $\eta:G\to\{\pm1\}$ and $\bar\Psi:G\to\kappa^\times$ be homomorphisms, and let $c,g\in G$ with $c^2=1$, $\eta(c)=-1$, $\bar\Psi(c)=1$, $\eta(g)=1$ and $\bar\Psi(g)^2\ne1$. Let $r:G\to\mathcal G_3(R)$ be a homomorphism whose reduction modulo $\mathfrak m_R$ is $r_0(\eta,\bar\Psi)$ (Definition~\ref{sha-ordinary-deformation}). Then there is $h\in GL_3(R)$, congruent to $1$, such that $r'=hrh^{-1}$ is normalized relative to $c$ and $g$: the matrix part of $r'(c)$ is $1$, and the entries $(1,2)$ and $(2,1)$ of the difference of the matrix parts of $r'(g)$ and $r'(cg^{-1}c)$ vanish.
\end{enumerate}
\end{lemma}
\begin{proof}
(1) The ring $R$ is henselian, since it is complete. The monic polynomial $X^2+bX+ad$ takes at $X=-b$ the value $ad\in\mathfrak m_R$, and its derivative there is $-b$, a unit. By Hensel's lemma it has a root $w\equiv-b$, which is a unit. Then $x=dw^{-1}$ lies in $\mathfrak m_R$, and $ax^2+bx+d=dw^{-2}(w^2+bw+ad)=0$.

(2) Write $S$ with rows $(s,q,p)$, $(q,u,t)$ and $(p,t,v)$. Then $s,q,t,v\in\mathfrak m_R$ and $p\equiv u\equiv1$.
By (1), applied to $uy^2+2uy+(u-1)$, there is $y\in\mathfrak m_R$ with $u(1+y)^2=1$; put $w=1+y\equiv1$.
Let $h_1$ be the matrix with rows $(1,-qw^2,0)$, $(0,w,0)$ and $(0,-tw^2,1)$. Since $w^2u=1$, a direct computation gives that $h_1S\,{}^th_1$ has rows $(s',0,p')$, $(0,1,0)$ and $(p',0,v')$, with $s'=s-q^2w^2$, $p'=p-qtw^2$ and $v'=v-t^2w^2$.
Here $s',v'\in\mathfrak m_R$ and $p'\equiv1$. By (1), applied to $v'x^2+2p'x+s'$, there is $x\in\mathfrak m_R$ with $v'x^2+2p'x+s'=0$. Then $p'+xv'\equiv1$ is a unit; let $k$ be its inverse, and let $h_2$ be the matrix with rows $(1,0,x)$, $(0,1,0)$ and $(-v'k^2/2,0,k(1-v'kx/2))$.
A direct computation, using $v'x^2+2p'x+s'=0$ and $(p'+xv')k=1$, gives $h_2(h_1S\,{}^th_1)\,{}^th_2=J$.
So $h=h_2h_1$ satisfies $hS\,{}^th=J$, and $h\equiv1$ because $q,t,x,v'\in\mathfrak m_R$ and $w\equiv k\equiv1$.

(3) By (1), applied to $-(b/2)y^2+ay+d$, there is $y\in\mathfrak m_R$ with $-(b/2)y^2+ay+d=0$. Then $a-by$ is a unit; put $x=b(a-by)^{-1}\in\mathfrak m_R$.
Let $h$ be the product of the matrix with rows $(1,x,-x^2/2)$, $(0,1,-x)$ and $(0,0,1)$ by the matrix with rows $(1,0,0)$, $(y,1,0)$ and $(-y^2/2,-y,1)$. Then $h\equiv1$, and a direct computation, using the two relations satisfied by $x$ and $y$, gives $hJ\,{}^thJ=1$ and the vanishing of the entries $(1,2)$ and $(2,1)$ of $hZJ\,{}^thJ$.

(4) Write $M_x$ for the matrix part of $r(x)$. For $h\in GL_3(R)$ and $x\in G$, the matrix part of $hr(x)h^{-1}$ is $hM_xh^{-1}$ if $\eta(x)=1$, and $hM_xJ\,{}^thJ$ if $\eta(x)=-1$.
Since $r(c)^2=1$ and $r(c)$ has sign $-1$, $M_cJ\,{}^tM_c^{-1}J=1$, that is, $M_c=J\,{}^tM_cJ$; so $S=M_cJ$ is symmetric. It is congruent to $J$, since $M_c$ reduces to $\operatorname{diag}(\bar\Psi(c),1,\bar\Psi(c)^{-1})=1$.
By (2) there is $h_1\equiv1$ with $h_1S\,{}^th_1=J$. Then $h_1$ is invertible, and the matrix part of $h_1r(c)h_1^{-1}$ is $h_1M_cJ\,{}^th_1J=JJ=1$.
Let $r_1=h_1rh_1^{-1}$, with matrix parts $M'_x$. Since $M'_c=1$, $r_1(c)$ has sign $-1$ and $r_1(g)$ has sign $1$, the matrix part of $r_1(cg^{-1}c)$ is $J\,{}^tM'_gJ$. So the difference $Z=M'_g-J\,{}^tM'_gJ$ has rows $(a,b,0)$, $(d,0,-b)$ and $(0,-d,-a)$, with $a=M'_{g,11}-M'_{g,33}$, $b=M'_{g,12}-M'_{g,23}$ and $d=M'_{g,21}-M'_{g,32}$.
Since $h_1\equiv1$, $M'_g$ reduces to $\operatorname{diag}(\bar\Psi(g),1,\bar\Psi(g)^{-1})$. So $b,d\in\mathfrak m_R$, and $a$ reduces to $\bar\Psi(g)^{-1}(\bar\Psi(g)^2-1)\ne0$, so it is a unit.
By (3) there is $h_2\equiv1$ with $h_2J\,{}^th_2J=1$ such that the entries $(1,2)$ and $(2,1)$ of $h_2ZJ\,{}^th_2J$ vanish. Then $h_2$ is invertible with inverse $J\,{}^th_2J$. Conjugation by $h_2$ keeps the matrix part at $c$ equal to $h_2J\,{}^th_2J=1$, and it transforms $Z$, the difference of two matrix parts of sign $1$, into $h_2Zh_2^{-1}=h_2ZJ\,{}^th_2J$.
So $h=h_2h_1$, which is congruent to $1$, has the required property.
\end{proof}
\leanref{Development/AutomorphicArithmetic/AutomorphicVarieties/UnitaryThreeGroups/NormEx.lean}{UpperTheory.V.U3.exists_normalised}
\leanref{Development/AutomorphicArithmetic/AutomorphicVarieties/UnitaryThreeGroups/MatGstar.lean}{UpperTheory.V.U3.mat_gstar}

\begin{lemma}\label{sha-rigidity-free}
Assume the Galois representations as in Theorem~\ref{sha-residually-reducible-automorphy}.
Let $F$ be a CM field, $\ell$ an odd prime, $\iota:\bar{\mathbf Q}_\ell\cong\mathbf C$, $n\ge1$, $a\in\mathbf Z$, $M\ge1$, and $m_1>m_2>\dots>m_n$ integers.
Let $r:\operatorname{Gal}(\bar F/F)\to GL_n(\bar{\mathbf Q}_\ell)$ be continuous, such that $r\otimes\chi_{\mathrm{cyc}}^{a}$ is attached through $\iota$ to a regular algebraic cuspidal automorphic representation $\pi$ of $GL_n(\mathbf A_F)$ for which $(\pi,\chi)$ is polarized for some Hecke character $\chi$ of the maximal totally real subfield of $F$ (Definition~\ref{gal-polarized}), and is windowed (Theorem~\ref{sha-residually-reducible-automorphy}).
Suppose that for every prime $w$ of $F$ above $\ell$ there are a place $A$ of $\bar F$ above $w$ and characters $\kappa_1,\dots,\kappa_n$ of the decomposition group $D_A$ such that the restriction of $r$ to $D_A$ is conjugate to an upper triangular representation with diagonal characters $\kappa_1,\dots,\kappa_n$, each $\kappa_i$ equals $\chi_{\mathrm{cyc}}^{m_i}\chi_i$ on the inertia group $I_A$ for a character $\chi_i$ of $\operatorname{Gal}(\bar F/F)$ trivial on $\operatorname{Gal}(\bar F/F(\mu_M))$, and $\kappa_i\ne\kappa_j\chi_{\mathrm{cyc}}$ on $D_A$ for all $i,j$.
Then for every place $A'$ of $\bar F$ above $\ell$ there is $g\in D_{A'}$ with $\chi_{\mathrm{cyc}}(g)=1$ such that the eigenvalues $\alpha_1,\dots,\alpha_n$ of $r(g)$ satisfy $\alpha_i^m\ne\alpha_j^m$ for $i\ne j$ and $m\ge1$.
\end{lemma}
\begin{proof}
Write $\chi=\chi_{\mathrm{cyc}}$. Let $A'$ be a place above $\ell$, over the prime $w$ of $F$, and let $A$ and $\kappa_1,\dots,\kappa_n$ be as in the statement for $w$. Write $D=D_A$ and $I=I_A$, let $q$ be the norm of $w$, $W\subset D$ the Weil group, $\deg:W\to\mathbf Z$ its degree, and $\mathrm a:W\to F_w^\times$ the reciprocity map of local class field theory, with $\operatorname{ord}_w\mathrm a(s)=-\deg s$.
The group of $M$-th roots of unity in $\bar F$ is finite; let $k\ge1$ be the order of its group of permutations. Every $s\in\operatorname{Gal}(\bar F/F)$ permutes the $M$-th roots of unity, so $s^k$ fixes them, and $\chi_i(s^k)=1$. Hence $(\kappa_i\chi^{-m_i})(s)^k=\chi_i(s)^k=1$ for $s\in I$ and all $i$.
Put $r'=r\otimes\chi^a$, $\kappa'_i=\kappa_i\chi^a$ and $m'_i=m_i+a$. Then $r'$ restricted to $D$ is conjugate to an upper triangular representation with diagonal characters $\kappa'_1,\dots,\kappa'_n$; $m'_1>\dots>m'_n$; $(\kappa'_i\chi^{-m'_i})^k=1$ on $I$; and $\kappa'_i\ne\kappa'_j\chi$ for all $i,j$.

\emph{The window.} Let $B$ be a place above $w$, and $t$ with $B=tA$. For $s\in D_B$, $t^{-1}st\in D$ and $r'(s)=r'(t)r'(t^{-1}st)r'(t)^{-1}$; so $r'$ is triangular on $D_B$ with the characters $s\mapsto\kappa'_i(t^{-1}st)$, which still satisfy the condition $\kappa'_i\ne\kappa'_j\chi$, because $\chi(t^{-1}st)=\chi(s)$.
So the hypothesis of condition ($\alpha$) of the windowed condition holds at $w$, and with its conclusion, condition ($\beta$) applies at $w$: there is $s_\pi\in\mathbf R$, independent of $w$, such that $\pi_w$ has supercuspidal support consisting of characters $\xi_1,\dots,\xi_n$ of $F_w^\times$ with $q^{s_\pi-1/2}<|\xi_j(\varpi_w)|<q^{s_\pi+1/2}$.
The ratios $|\xi_i(\varpi_w)/\xi_j(\varpi_w)|$ lie strictly between $q^{-1}$ and $q$, so no $\xi_i$ is $\xi_j|\cdot|^{\pm1}$. Hence $\pi_w$ is the representation induced from $\xi_1,\dots,\xi_n$, which is irreducible \cite[Theorem~4.2]{BZ77}, and $y=\operatorname{rec}(\pi_w\nu^{(1-n)/2})$ is the sum of the characters $(\xi_j\nu^{(1-n)/2})\circ\mathrm a$ of $W$: it is abelian and monodromy-free.
For $s\in W$ of degree $1$, any two eigenvalues $\alpha,\beta$ of $y(s)$ satisfy $|\alpha|<q|\beta|$: since $\operatorname{ord}_w\mathrm a(s)=-1$ and the $\xi_j$ have absolute value $1$ on units, $|\alpha/\beta|$ is a ratio $|\xi_j(\varpi_w)/\xi_i(\varpi_w)|$.

\emph{Local-global compatibility.} Since $(\pi,\chi)$ is polarized, a triple $(R,m,\mathrm{HT})$ is attached to $\pi$ (Definition~\ref{gal-galois-representations}), and $R$ is isomorphic to $r'$, since both have the Frobenius eigenvalues of $\pi$ at almost all places. By Definition~\ref{gal-attached}(5)(b) for $R$, $r'$ is potentially semistable on $D$, and $\mathrm{WD}_\iota(r'|_D)$ is compatible with $y$; so $y(s)$ and $\mathrm{WD}_\iota(r'|_D)(s)$ have the same characteristic polynomial for $s\in W$.
For each $i$, $\mu_i(s)=q^{m'_i\deg s}\,\iota\big(\kappa'_i(s)\chi(s)^{-m'_i}\big)$ is a smooth character of $W$, because $\kappa'_i$ is continuous and $(\kappa'_i\chi^{-m'_i})^k=1$ on $I$. By local class field theory, $\kappa'_i$ is then a de Rham character with $\mathrm{WD}_\iota(\kappa'_i)=\mu_i$. Since $\mathrm{WD}_\iota$ is exact on potentially semistable representations \cite{Fon94} and $r'|_D$ is an iterated extension of the $\kappa'_i$, the characteristic polynomial of $y(s)$ is $\prod_i(X-\mu_i(s))$ for $s\in W$.

\emph{An element with trivial cyclotomic character.} Let $s_0\in W$ with $\mathrm a(s_0)=\ell$. Then $\deg s_0=-\operatorname{ord}_w(\ell)\ne0$. By local class field theory $\chi(s_0)$ is an integral power of $\ell$; it is an $\ell$-adic unit, so $\chi(s_0)=1$, and $\kappa'_i(s_0)=\kappa_i(s_0)$.

\emph{Free ratios.} Suppose that $\kappa_i(s_0)^N=\kappa_j(s_0)^N$ for some $i\ne j$ and $N\ge1$. Then $|\iota\kappa'_i(s_0)|=|\iota\kappa'_j(s_0)|$, and since $\chi(s_0)=1$, $|\mu_i(s_0)|=q^{(m'_i-m'_j)\deg s_0}|\mu_j(s_0)|$.
The function $f(s)=\log|\mu_i(s)|-\log|\mu_j(s)|$ is a homomorphism $W\to\mathbf R$. For $s$ of degree $1$, $\mu_i(s)$ and $\mu_j(s)$ are eigenvalues of $y(s)$, so $|f(s)|<\log q$.
Let $\varphi\in W$ have degree $1$, and put $e=m'_i-m'_j$, so that $f(s_0)=e\deg(s_0)\log q$. For every integer $m$, $s_0^m\varphi^{1-m\deg s_0}$ has degree $1$, so $|m\deg(s_0)(e\log q-f(\varphi))+f(\varphi)|<\log q$. Since $\deg s_0\ne0$, this forces $f(\varphi)=e\log q$, and then the case $m=0$ gives $|e|<1$. So $m'_i=m'_j$, which contradicts $i\ne j$.
Hence $\kappa_i(s_0)^N\ne\kappa_j(s_0)^N$ for $i\ne j$ and $N\ge1$.

\emph{Transport.} Since $A'$ lies over $w$, $A'=tA$ for some $t$. Then $g=ts_0t^{-1}\in D_{A'}$, $\chi(g)=\chi(s_0)=1$, and $r(g)$ is conjugate to $r(s_0)$, whose eigenvalues are $\kappa_1(s_0),\dots,\kappa_n(s_0)$, since $r$ is triangular on $D$ with these diagonal characters.
\end{proof}
\leanref{Development/AutomorphicArithmetic/AutomorphicVarieties/ModularityArgument/ThresholdArguments/DenseThresholdCases/FFreeCU.lean}{UpperTheory.V.Z.J.free_at_of_unlinked_of_automorphic_cyclotomic_twist}
\leanref{Development/AutomorphicArithmetic/AutomorphicVarieties/ModularityArgument/ThresholdArguments/DenseThresholdCases/RRFree.lean}{UpperTheory.V.Z.Win.RR.exists_cycZ_eq_one_freeRatios_of_ne_mul_cycBar}
\leanref{Development/AutomorphicArithmetic/AutomorphicVarieties/ModularityArgument/ThresholdArguments/DenseThresholdCases/RRFree.lean}{UpperTheory.V.Z.Win.RR.purLocA_alt}
\leanref{Development/AutomorphicArithmetic/AutomorphicVarieties/ModularityArgument/ThresholdArguments/DenseThresholdCases/RRFree.lean}{UpperTheory.V.Z.Win.RR.free_core_alt}
\leanref{Development/AutomorphicArithmetic/AutomorphicVarieties/ModularityArgument/ThresholdArguments/DenseThresholdCases/RRFree.lean}{UpperTheory.V.Z.Win.RR.freeAt_of_autPA_of_unl_of_finrank_eq_two}
\leanref{Development/AutomorphicArithmetic/AutomorphicVarieties/ModularityArgument/AutomorphyLifting/LocalConditionCharacters/LocalConditionOfFlagOfCharactersNotCyclotomicShi.lean}{UpperTheory.V.Z.Win.RR.local_condition_of_flag_of_characters_not_cyclotomic_shift}

\begin{lemma}\label{sha-rigidity-adjoint}
Assume the Galois representations as in Theorem~\ref{sha-residually-reducible-automorphy}.
Let $\ell$ be an odd prime, $\mathcal O$ as in Definition~\ref{sha-character-parts}, $N,e\ge1$, $\psi$ a Dirichlet character modulo $N$ with values in $\mathcal O$, and $\Psi=\chi_{\mathrm{cyc}}^{e}\psi$.
Let $K$ be an imaginary quadratic field with quadratic character $\eta$, let $c$ be a complex conjugation, $B$ a place above $\ell$ and $\iota:\bar{\mathbf Q}_\ell\cong\mathbf C$, and suppose that:
\begin{enumerate}
\item $\Psi(c)=1$, and $\eta$ is unramified at the primes not dividing an integer $M$ divisible by $\ell N$;
\item for some $g_2\in\operatorname{Gal}(\bar{\mathbf Q}/K)$ with $\Psi(g_2)^2-1\in\mathcal O^\times$ there are a link of $\eta$, $\Psi$, $c$, $g_2$ and $B$, for the integer $M$ of (1), to an abelian CM field $L\supset K$ that satisfies the residual conditions, and a seed datum for $L$, as in the proof of Theorem~\ref{sha-residually-reducible-automorphy}.
\end{enumerate}
Let $g_1\in\operatorname{Gal}(\bar{\mathbf Q}/\mathbf Q)$, let $A\subset\bar{\mathbf Q}_\ell$ be a discrete valuation ring finite over $\mathcal O$, and let $r$ be an ordinary point relative to $g_1$ and $B$ with values in $A$, in the sense of Definition~\ref{sha-ordinary-deformation}(c) for $\eta$, $c$ and $\Psi$; let $\rho$ be its restriction to $\operatorname{Gal}(\bar{\mathbf Q}/K)$. Suppose moreover that:
\begin{enumerate}\setcounter{enumi}{2}
\item $r$ is irreducible on every open subgroup of $\ker\eta$, and $\rho\otimes\chi_{\mathrm{cyc}}^{-1}$ is attached through $\iota$ to a regular algebraic cuspidal automorphic representation $\pi$ of $GL_3(\mathbf A_K)$ with $\pi^\vee\cong\pi^c$, and is windowed (Theorem~\ref{sha-residually-reducible-automorphy});
\item at every prime $w$ of $K$ above $\ell$, the restriction of $\rho$ to a decomposition group $D_w$ is conjugate to an upper triangular representation whose diagonal characters $\kappa_1,\kappa_2,\kappa_3$ equal $\chi_{\mathrm{cyc}}^{m_i}\chi_i$ on the inertia group $I_w$, for $(m_1,m_2,m_3)=(e,0,-e)$ and characters $\chi_1,\chi_2,\chi_3$ of $\operatorname{Gal}(\bar{\mathbf Q}/K)$ trivial on $\operatorname{Gal}(\bar{\mathbf Q}/K(\mu_N))$, and satisfy $\kappa_i\ne\kappa_j\chi_{\mathrm{cyc}}$ for all $i,j$.
\end{enumerate}
Then $H^1_f(\mathbf Q,\operatorname{ad}r)=0$, where $\operatorname{ad}r$ is $M_3(\bar{\mathbf Q}_\ell)$ with $\operatorname{Gal}(\bar{\mathbf Q}/\mathbf Q)$ acting through $r$ and the adjoint action of $\mathcal G_3$.
\end{lemma}
\begin{proof}
\emph{Step 1: a free element.} Let $B'$ be a place above $\ell$ and $D'\subset\operatorname{Gal}(\bar{\mathbf Q}/K)$ its decomposition group. There is $g\in D'$ with $\chi_{\mathrm{cyc}}(g)=1$ such that the eigenvalues $\alpha_1,\alpha_2,\alpha_3$ of $\rho(g)$ satisfy $\alpha_i^m\ne\alpha_j^m$ for $i\ne j$ and $m\ge1$. Indeed, $\rho$ is continuous, since $r$ is, and Lemma~\ref{sha-rigidity-free} applies to $\rho$ with $F=K$, $n=3$, $a=-1$, $M=N$ and $(m_1,m_2,m_3)=(e,0,-e)$: its automorphy and windowed hypotheses are (3), where $(\pi,\eta)$ is polarized because $\pi^\vee\cong\pi^c$, $\eta\circ N_{K/\mathbf Q}=1$ and $\eta(-1)=-1=(-1)^3$ at the real place, with $\eta$ viewed as a Hecke character of $\mathbf Q$; and its hypothesis at the primes above $\ell$ is (4), with the characters $\chi_i$ of (4).

\emph{Step 2: renormalization.} Let $g_2$ be the element of (2), and $J$ the antidiagonal $3\times3$ matrix with entries $1$.
Since $c$ is a complex conjugation, $c^2=1$. By Lemma~\ref{sha-descent-points}, applied to the link of (2), which satisfies the residual conditions, $\eta(c)=-1$, $\eta(g_2)=1$ and $\Psi(g_2)^2-1\in\mathcal O^\times$; and $\Psi(c)=1$ by (1). Since $\ell$ is odd, $2$ is a unit in $\mathcal O$, hence in $A$.
The reduction of $r$ modulo $\mathfrak m_A$ is $r_0(\eta,\Psi)$, and $A$ is complete. So Lemma~\ref{sha-rigidity-normalize}(4), for the reduction $\bar\Psi$ of $\Psi$, gives $h\in GL_3(A)$ congruent to $1$ such that $r'=hrh^{-1}$ is normalized relative to $c$ and $g_2$: its proof first makes the matrix part at $c$ equal to $1$ by part (2) of that lemma, and then clears the entries $(1,2)$ and $(2,1)$ at $g_2$ by part (3).
Conjugation by $h$ preserves the other conditions of an ordinary point: the multiplier; the reduction and the continuity, because $h\equiv1$; the value $1$ at the elements of the inertia groups away from $\ell$ at which $\eta$ and $\Psi$ are trivial; and the ordinarity at $B$, with $hh_B$ in place of the conjugating matrix $h_B$ of $r$. So $r'$ is an ordinary point relative to $g_2$ and $B$.
Conjugation by $h$ is an isomorphism from $\operatorname{ad}r$ to $\operatorname{ad}r'$, so it suffices to show $H^1_f(\mathbf Q,\operatorname{ad}r')=0$. On $\ker\eta$ the matrix parts of $r'$ are those of $r$ conjugated by $h$. So $r'$ is irreducible on every open subgroup of $\ker\eta$, by (3); and its restriction $\rho'=h\rho h^{-1}$ to $\operatorname{Gal}(\bar{\mathbf Q}/K)$ is isomorphic to $\rho$, so that $\rho'\otimes\chi_{\mathrm{cyc}}^{-1}$ is attached through $\iota$ to $\pi$ and is windowed, and the elements of Step 1 have the same eigenvalues under $\rho'$ as under $\rho$.

\emph{Step 3: over $L$.} The group $\operatorname{Gal}(\bar{\mathbf Q}/L)$ is an open subgroup of $\ker\eta$, so $r'$ is irreducible on it.
The restriction of $\rho'$ to the decomposition group $D_P\cap\operatorname{Gal}(\bar{\mathbf Q}/K)$ at every place $P$ of $\bar{\mathbf Q}$ above $\ell$ is conjugate to an upper triangular representation. Indeed, by (iv) of the link, $P=uA_v=u\gamma_vB$ for a place $v$ of $L^+$ above $\ell$ and some $u$; put $t=u\gamma_v$, so that $t^{-1}D_Pt=D_B$.
The ordinarity of $r'$ at $B$ gives $h_B\in GL_3(A)$ such that $h_B^{-1}M_xh_B$ is upper triangular for $x\in D_B\cap\ker\eta$, where $M_x$ is the matrix part of $r'(x)$; and for $s\in D_P\cap\operatorname{Gal}(\bar{\mathbf Q}/K)$, $t^{-1}st\in D_B\cap\ker\eta$ and $r'(s)=r'(t)r'(t^{-1}st)r'(t)^{-1}$.
If $\eta(t)=1$, then $\rho'(s)=M_tM_{t^{-1}st}M_t^{-1}$, and $(M_th_B)^{-1}\rho'(s)(M_th_B)$ is upper triangular.
If $\eta(t)=-1$, then $\rho'(s)=M_t\theta(M_{t^{-1}st})M_t^{-1}$ with $\theta(X)=J\,{}^tX^{-1}J$; since $\theta$ is a homomorphism that maps upper triangular matrices to upper triangular matrices, $(M_t\theta(h_B))^{-1}\rho'(s)(M_t\theta(h_B))=\theta(h_B^{-1}M_{t^{-1}st}h_B)$ is upper triangular.
Since $r'(c)$ has matrix part $1$ and sign $-1$, the matrix part of $r'(cgc^{-1})$ is $\theta(\rho'(g))$, so ${}^t\rho'(g)J\rho'(cgc^{-1})=J$ for $g\in\operatorname{Gal}(\bar{\mathbf Q}/K)$; thus $\rho'\otimes\chi_{\mathrm{cyc}}^{-1}$ is polarized, with multiplier $\chi_{\mathrm{cyc}}^{-2}$.
The extension $L/K$ is abelian, hence solvable. Base change along it \cite{AC89}, applied to the conjugate self-dual cuspidal representation $\pi$ attached to $\rho'\otimes\chi_{\mathrm{cyc}}^{-1}$, with this polarization, makes the restriction of $\rho'\otimes\chi_{\mathrm{cyc}}^{-1}$ to $\operatorname{Gal}(\bar{\mathbf Q}/L)$, which is irreducible, attached through $\iota$ to a regular algebraic conjugate self-dual cuspidal automorphic representation of $GL_3(\mathbf A_L)$. The windowed condition over $L$ follows from that over $K$ and the triangular forms at the places above $\ell$.

\emph{Step 4: linked points.} By the descent step of the proof of Theorem~\ref{sha-residually-reducible-automorphy}, $r'$ is a linked point of rank three for the link of (2): it factors through the quotient $\Gamma$ of the link, and its characters at the places of the link above $\ell$ are $\Psi$, $1$ and $\Psi^{-1}$ on inertia.

\emph{Step 5: the vanishing over $L^+$.} Let $\rho'_L$ be the restriction of $\rho'$ to $\operatorname{Gal}(\bar{\mathbf Q}/L)$. By Step 3, $\rho'_L$ is irreducible, and $\rho'_L\otimes\chi_{\mathrm{cyc}}^{-1}$ is attached through $\iota$ to a regular algebraic conjugate self-dual cuspidal automorphic representation of $GL_3(\mathbf A_L)$ and is windowed.
At every place $P$ above $\ell$, $\rho'_L$ has an element as in Step 1: if $g\in D_P\cap\operatorname{Gal}(\bar{\mathbf Q}/K)$ is the element of Step 1 for $\rho'$, some power $g^d$ with $d\ge1$ lies in the subgroup $\operatorname{Gal}(\bar{\mathbf Q}/L)$ of finite index; then $\chi_{\mathrm{cyc}}(g^d)=1$, and the eigenvalues $\alpha_i^d$ of $\rho'(g^d)$ satisfy $(\alpha_i^d)^m\ne(\alpha_j^d)^m$ for $i\ne j$ and $m\ge1$.
There are a finite set $S'$ of places of $L^+$ containing those above $\ell$, with the quotient $\Gamma'$ of $\operatorname{Gal}(\bar{\mathbf Q}/L^+)$ by the closed normal subgroup $N'$ generated by the inertia groups at the places outside $S'$, the image $\Gamma'_0$ of $\operatorname{Gal}(\bar{\mathbf Q}/L)$ and the image $c'$ of a complex conjugation $\tau\in\operatorname{Gal}(\bar{\mathbf Q}/L^+)$, and a representation $\rho_V$ of $\Gamma'_0$ with $\rho_V(s)=\chi_{\mathrm{cyc}}(s)^{-1}\rho'(s)$ for $s\in\operatorname{Gal}(\bar{\mathbf Q}/L)$, such that:
let $P\in GL_3(\bar{\mathbf Q}_\ell)$ satisfy ${}^t\rho_V(x)P\rho_V(c'xc'^{-1})=\chi_{\mathrm{cyc}}(x)^{-2}P$ for $x\in\Gamma'_0$, and let $R$ be a representation of $\Gamma'$ on $M_3(\bar{\mathbf Q}_\ell)$ with $R(x)X=\rho_V(x)X\rho_V(x)^{-1}$ for $x\in\Gamma'_0$ and $R(c')X=-P^{-1}\,{}^tXP$; then every continuous cocycle of $\Gamma'$ with values in $R$ whose inflation to $\operatorname{Gal}(\bar{\mathbf Q}/L^+)$ satisfies the Bloch-Kato conditions \cite{BK90} is a coboundary.
Alpöge obtains this from his own adjoint lifting theorem, applied through the universal ordinary determinant rings of the realizations of rank three of the link to the linked point of Step 4, whose characters on the inertia groups of the link are $\Psi$, $1$ and $\Psi^{-1}$. It is also a case of the theorem of Newton and Thorne \cite[Theorem~A]{NT23}, whose hypotheses hold here, as follows.

\emph{The representation $\rho_V$.} Let $S'$ be the set of places of $L^+$ above the primes dividing $M$; it contains those above $\ell$, since $\ell\mid M$. By Definition~\ref{sha-descent-link}(1), the inertia groups of $\operatorname{Gal}(\bar{\mathbf Q}/L^+)$ at the finite places lie in $\ker\eta$, hence in $\operatorname{Gal}(\bar{\mathbf Q}/L)$; so $N'\subset\operatorname{Gal}(\bar{\mathbf Q}/L)$.
Let $s$ lie in the inertia group of $\operatorname{Gal}(\bar{\mathbf Q}/L^+)$ at a place outside $S'$. Then $\eta(s)=1$, and $\chi_{\mathrm{cyc}}(s)=\psi(s)=1$ because the place does not divide $\ell N$; so $\Psi(s)=1$, and $\rho'(s)=1$ because $r'$ is an ordinary point (Step 2).
The elements $s\in\operatorname{Gal}(\bar{\mathbf Q}/L)$ with $\chi_{\mathrm{cyc}}(s)^{-1}\rho'(s)=1$ form a closed subgroup that contains these inertia groups, and with them their conjugates under $\operatorname{Gal}(\bar{\mathbf Q}/L^+)$, which are the inertia groups at other places outside $S'$; so it contains $N'$, and $\rho_V$ is well defined.

\emph{Enormous image.} Let $H\subset\operatorname{Gal}(\bar{\mathbf Q}/L)$ be the subgroup fixing $\mu_{\ell^\infty}$, and let $h_0\in H$ be one of the elements above, so that the eigenvalues $\beta_1,\beta_2,\beta_3$ of $\rho'(h_0)$ satisfy $\beta_i^m\ne\beta_j^m$ for $i\ne j$ and $m\ge1$.
First, $\rho'$ is irreducible on $U\cap H$ for every open subgroup $U$ of $\operatorname{Gal}(\bar{\mathbf Q}/L)$. We may assume that $U$ is normal, of index $n$. By Step 2, $\rho'$ is irreducible on $U$, which is an open subgroup of $\ker\eta$. Let $V_0\ne0$ be a subspace of $\bar{\mathbf Q}_\ell^3$ stable under $\rho'(U\cap H)$, of minimal dimension. Since $U\cap H$ is normal in $U$, each $\rho'(u)V_0$ with $u\in U$ is again such a subspace, and their sum, which is stable under $\rho'(U)$, is $\bar{\mathbf Q}_\ell^3$.
If $\dim V_0=2$, then $V_0\cap\rho'(u)V_0\ne0$ for every $u$, so $\rho'(u)V_0=V_0$ by minimality, and $V_0$ is stable under $\rho'(U)$, which is impossible.
If $\dim V_0=1$, each line $\rho'(u)V_0$ is stable under $\rho'(h_0^n)$, where $h_0^n\in U\cap H$ has the distinct eigenvalues $\beta_i^n$; so these lines are among the three eigenlines of $\rho'(h_0^n)$, and the stabilizer of $V_0$ in $U$, which is closed, has index at most $3$. It is therefore open, and $V_0$ is stable under it, which is impossible. So $V_0=\bar{\mathbf Q}_\ell^3$.
The representation $\rho'$ takes values in $GL_3(A)$. Let $E'$ be a finite extension of the fraction field of $A$ over which the characteristic polynomials of all $\rho'(g)$ split; it exists because the splitting field of a cubic has degree at most $6$, and the fraction field of $A$ has only finitely many extensions of degree at most $6$.
Let $G$ be the Zariski closure of $\rho'(H)$ and $G^\circ$ its identity component. The inverse image of $G^\circ$ in $H$ is closed of finite index, hence open in the profinite group $H$, so it contains $U\cap H$ for some open subgroup $U$ of $\operatorname{Gal}(\bar{\mathbf Q}/L)$. So $G^\circ$ acts irreducibly on $\bar{\mathbf Q}_\ell^3$; and it contains $\rho'(h_0^d)$ for some $d\ge1$, which is regular semisimple, since its eigenvalues $\beta_i^d$ are distinct. By \cite[Lemma~2.28(2)]{NT23}, applied with the coefficient field $E'$, $\rho'(H)$ is enormous in the sense of \cite[Definition~2.23]{NT23}. Since $\chi_{\mathrm{cyc}}$ is trivial on $H$, this is the image of $H$ under $s\mapsto\rho_V(s)$.

\emph{The vanishing.} By Step 3, $s\mapsto\rho_V(s)$ is irreducible and attached through $\iota$ to a regular algebraic cuspidal automorphic representation $\Pi$ of $GL_3(\mathbf A_L)$ with $\Pi^\vee\cong\Pi^c$. By the Chebotarev density and Brauer-Nesbitt theorems it has the same semisimplification as the Galois representation $r_{\Pi,\iota}$ attached to $\Pi$ in \cite{NT23}, so $r_{\Pi,\iota}$ is irreducible and isomorphic to it. Its image of $H$ is enormous, so \cite[Theorem~A]{NT23} gives $H^1_f(L^+,\operatorname{ad}r_{\Pi,\iota})=0$. Here $\operatorname{ad}r_{\Pi,\iota}$ is $M_3(\bar{\mathbf Q}_\ell)$, on which $\operatorname{Gal}(\bar{\mathbf Q}/L)$ acts by conjugation through $r_{\Pi,\iota}$ and $\tau$ acts by $X\mapsto-X^*$, where $X^*$ is the adjoint of $X$ for the pairing that expresses the conjugate self-duality of $r_{\Pi,\iota}$.
Let $P$ and $R$ be as in the statement of this step. Since $\rho_V$ is irreducible, its condition determines $P$ up to a scalar. The transpose ${}^tP$ satisfies it too: transpose the condition and replace $x$ by $c'xc'^{-1}$, using $c'^2=1$. So ${}^tP=\pm P$, and $X\mapsto P^{-1}\,{}^tXP$ is the adjoint for the bilinear form $(u,w)\mapsto{}^tuPw$, which expresses the conjugate self-duality of $\rho_V$. Hence the inflation of $R$ to $\operatorname{Gal}(\bar{\mathbf Q}/L^+)$ is isomorphic to $\operatorname{ad}r_{\Pi,\iota}$.
Let $z$ be a continuous cocycle of $\Gamma'$ with values in $R$ whose inflation $\tilde z$ satisfies the Bloch-Kato conditions. Then $\tilde z$ represents a class in $H^1_f(L^+,\operatorname{ad}r_{\Pi,\iota})=0$, so there is $X$ with $\tilde z(g)=R(g)X-X$ for all $g\in\operatorname{Gal}(\bar{\mathbf Q}/L^+)$, and $z(y)=R(y)X-X$ for all $y\in\Gamma'$.

\emph{Step 6: descent.} Let $R_0$ be the action of $\operatorname{Gal}(\bar{\mathbf Q}/\mathbf Q)$ on $\operatorname{ad}r'$, and $y$ a continuous cocycle with values in it that satisfies the Bloch-Kato conditions. We show that $y$ is a coboundary.
Its restriction to $\operatorname{Gal}(\bar{\mathbf Q}/L^+)$ satisfies the Bloch-Kato conditions, since these conditions are local and are preserved by restriction \cite{BK90}.
The elements $\tau$ and $c$ both act on $K$ by complex conjugation, so $\tau=wc$ with $w\in\operatorname{Gal}(\bar{\mathbf Q}/K)$.
For $s\in\operatorname{Gal}(\bar{\mathbf Q}/L)$, $R_0(s)X=\rho'(s)X\rho'(s)^{-1}=\rho_V(s)X\rho_V(s)^{-1}$, the scalar $\chi_{\mathrm{cyc}}(s)^{-1}$ cancelling. Since the matrix part of $r'(c)$ is $1$ and its sign is $-1$, $R_0(c)X=-J\,{}^tXJ$; so $R_0(\tau)X=-P^{-1}\,{}^tXP$ with $P=J\rho'(w)^{-1}$.
The subgroup $N'$ lies in $\operatorname{Gal}(\bar{\mathbf Q}/L)$ and maps to $1$ in $\Gamma'_0$, so $R_0$ is trivial on $N'$.
The cocycle $y$ vanishes on $N'$. Indeed, at a place outside $S'$, which is not above $\ell$, the Bloch-Kato condition makes $y$ on the inertia group a coboundary $s\mapsto R_0(s)X-X$, which is $0$ because $R_0$ is trivial there; and the elements of $N'$ at which $y$ vanishes form a closed normal subgroup of $\operatorname{Gal}(\bar{\mathbf Q}/L^+)$, since $y$ is a continuous cocycle and $R_0$ is trivial on $N'$.
So $y$ induces a continuous cocycle $z$ of $\Gamma'$ with values in the representation $R$ induced by $R_0$, and $R$ has the form of Step 5 for $P$.
The matrix $P$ satisfies the condition of Step 5: for $x$ the image of $s\in\operatorname{Gal}(\bar{\mathbf Q}/L)$, $c'xc'^{-1}$ is the image of $wcsc^{-1}w^{-1}$, so by Step 3
\[
 {}^t\rho_V(x)P\rho_V(c'xc'^{-1})=\chi_{\mathrm{cyc}}(s)^{-2}\,{}^t\rho'(s)J\rho'(csc^{-1})\rho'(w)^{-1}=\chi_{\mathrm{cyc}}(s)^{-2}P.
\]
By Step 5, $z$ is a coboundary, so there is $X$ with $y(s)=R_0(s)X-X$ for $s\in\operatorname{Gal}(\bar{\mathbf Q}/L^+)$.
Finally, $\operatorname{Gal}(\bar{\mathbf Q}/L^+)$ has finite index $d$ in $\operatorname{Gal}(\bar{\mathbf Q}/\mathbf Q)$, and $\operatorname{ad}r'$ is a vector space over $\bar{\mathbf Q}_\ell$, a field of characteristic $0$. The cocycle $y'(s)=y(s)-(R_0(s)X-X)$ vanishes on $\operatorname{Gal}(\bar{\mathbf Q}/L^+)$, so it is constant on the left cosets of this subgroup. Let $\Sigma$ be the sum of its values on a set of representatives of these cosets. Since left multiplication by $s$ permutes the cosets, $\Sigma=d\,y'(s)+R_0(s)\Sigma$ for every $s$. So $y(s)=R_0(s)Y-Y$ with $Y=X-d^{-1}\Sigma$, and $H^1_f(\mathbf Q,\operatorname{ad}r')=0$.
% Th.RigEnd3Alt, proved by Z.Win.Dn.rigEnd3_alt from HypothesesH' and Th.AdjointLiftingTheorem, with FreeAt from free_at_of_unlinked_of_automorphic_cyclotomic_twist (Th.Unl), which is Lemma sha-rigidity-free; then freeAt_iso (Step 2) and freeAt_res (Step 5).
% Step 2 is Pr.marks_of_link and Th.End3.ordPt_renorm (U3.exists_normalised: symm_norm, anti_norm, quad_root; ordPt_conj), with h1fZero_conjBy. Step 3 is Dn.localConditionRankThreeLink_of_inertia_ordinary (Flagged), pair_of and Th.End3.isAutCsdA_twist_comp_galEmb_of_repIso_of_flagged (autCsd0_comp_galEmb_of_isPair_of_isIrr_of_isAut, automorphy_up_base_change). Step 4 is Pr.exists_nrep3.
% Step 5 in the Lean is Th.linked_rigid_of_printed_adjoint_lifting (Pr.coefEx, Pr.cRealEx for n = 3, Th.rigPtOf_alt, with Th.AdjointLiftingTheorem, itself adjoint_lifting_of_cohomology), whose conclusion is Th.RigOver: a Hida setting s with level V (S' = V.S', Gamma' = V.frame.Gamma, c' the image of s.c, rho_V = rho with V.pr) and Th.AdZero V rho_V. The text proves the same conclusion, with S' the places above the primes dividing M, from [NT23, Theorem A], after checking its enormous-image hypothesis; see the correspondence note.
% Step 6 is the private ad_down, through Th.adDown: X.Kato.bkRes (restriction), zero_NS (vanishing on N'), hform and hext (the form P r(w)^-1 and the extension R'), AdZero, and avg (the averaging over the cosets).
\end{proof}
\leanref{Development/AutomorphicArithmetic/AutomorphicVarieties/ModularityArgument/ThresholdArguments/DenseThresholdCases/RigEnd3.lean}{UpperTheory.V.Th.RigEnd3Alt}
\leanref{Development/AutomorphicArithmetic/AutomorphicVarieties/ModularityArgument/ThresholdArguments/DenseThresholdCases/AutUpA.lean}{UpperTheory.V.Z.Win.Dn.rigEnd3_alt}

\begin{lemma}\label{sha-rigidity-two-by-two}
Let $\Gamma$ be a group, $k$ a field and $W:\Gamma\to GL_2(k)$ a homomorphism.
\begin{enumerate}
\item Let $H\subset\Gamma$ be a subgroup and $\Psi:\Gamma\to k^\times$ a character. For $g\in\Gamma$ let $B(g)$ be the $3\times3$ matrix whose rows and columns $1$ and $3$ form $W(g)$, whose entry $(2,2)$ is $1$, and whose other entries are $0$. If there is $X\in GL_3(k)$ with $XB(g)X^{-1}=\operatorname{diag}(\Psi(g),1,\Psi(g)^{-1})$ for all $g\in H$, then $W$ is conjugate on $H$ to $\operatorname{diag}(\Psi,\Psi^{-1})$.
\item Suppose that $k$ has characteristic $0$. Let $H\subset\Gamma$ be a subgroup and $\eta:\Gamma\to\{\pm1\}$ and $\Psi:\Gamma\to k^\times$ characters such that $W$ is conjugate on $H\cap\ker\eta$ to $\operatorname{diag}(\Psi,\Psi^{-1})$. Then there are characters $\chi_1,\chi_2$ of $\Gamma$, with $\chi_i(g)=1$ whenever $\Psi(g)=\eta(g)=1$, such that $W$ is conjugate on $H$ to $\operatorname{diag}(\chi_1,\chi_2)$.
\item Let $I\subset D$ be subgroups of $\Gamma$ with $I$ normal in $D$, and $\Psi:\Gamma\to k^\times$ a character. Suppose that there is $v\ne0$ in $k^2$ with $W(s)v=\Psi(s)v$ and $\det W(s)=1$ for all $s\in I$, and that $\Psi(s)^2\ne1$ for some $s\in I$. Then $W$ is conjugate on $D$ to an upper triangular representation whose diagonal characters $\kappa_1,\kappa_2$ of $D$ equal $\Psi$ and $\Psi^{-1}$ on $I$.
\end{enumerate}
\end{lemma}
\begin{proof}
(1) to (3) are direct computations with $2\times2$ and $3\times3$ matrices.
\end{proof}
\leanref{Development/AutomorphicArithmetic/AutomorphicVarieties/UnitaryThreeGroups/DiagonalOnSubgroupOfConjugate.lean}{UpperTheory.diagonal_on_subgroup_of_conjugate}
\leanref{Development/AutomorphicArithmetic/AutomorphicVarieties/UnitaryThreeGroups/ExistsIsDiagOnOfIsDiagOnInfKerSign.lean}{UpperTheory.exists_isDiagOn_of_isDiagOn_inf_ker_sign}
\leanref{Development/AutomorphicArithmetic/AutomorphicVarieties/UnitaryThreeGroups/ExistsIsFlagOnOfEigenvectorNormalSubgroup.lean}{UpperTheory.exists_isFlagOn_of_eigenvector_normal_subgroup}

\begin{lemma}\label{sha-rigidity-corner}
Let $k$ be a field and $H$ a set, and for $s\in H$ let $R(s)$ be a $3\times3$ matrix over $k$ whose entries $(1,2)$, $(2,1)$, $(2,3)$ and $(3,2)$ vanish and whose entry $(2,2)$ is $1$; let $C(s)$ be its corner block, formed by rows and columns $1$ and $3$. Let $g\in GL_3(k)$ and $a,c:H\to k$ be such that, for every $s\in H$, $gR(s)g^{-1}$ is upper triangular with entries $(1,1)$ and $(3,3)$ equal to $a(s)$ and $c(s)$. Suppose that $a(s_0)\ne1$ and $c(s_2)\ne1$ for some $s_0,s_2\in H$.
Then there is $Q\in GL_2(k)$ such that $QC(s)Q^{-1}$ is upper triangular with diagonal entries $a(s)$ and $c(s)$ for every $s\in H$.
\end{lemma}
\begin{proof}
Let $x=(x_1,x_2,x_3)$ be the first column of $g^{-1}$ and $p=(p_1,p_2,p_3)$ the last row of $g$. Since $gR(s)g^{-1}$ is upper triangular, $R(s)x=a(s)x$ and $pR(s)=c(s)p$ for every $s$.
By the shape of $R(s_0)$, the second coordinate of $R(s_0)x$ is $x_2$; so $x_2=a(s_0)x_2$, and $x_2=0$. In the same way $p_2=c(s_2)p_2$, and $p_2=0$.
So $x'=(x_1,x_3)$ and $p'=(p_1,p_3)$ satisfy $C(s)x'=a(s)x'$ and $p'C(s)=c(s)p'$, and $p'x'=px=0$, the entry $(3,1)$ of $gg^{-1}$. Both are nonzero, since the entries $(1,1)$ and $(3,3)$ of $gg^{-1}$ are $1$.
Choose $u,v\in k$ with $ux_1+vx_3=1$, and let $d=up_3-vp_1$. Using $p_1x_1+p_3x_3=0$, $p_3=dx_1$ and $p_1=-dx_3$; so $d\ne0$, as $p'\ne0$.
Let $Q$ be the matrix with rows $(u,v)$ and $(p_1,p_3)$. It is invertible, and $Q^{-1}$ has columns $x'$ and $d^{-1}(-v,u)$.
The entry $(2,1)$ of $QC(s)Q^{-1}$ is $p'C(s)x'=c(s)p'x'=0$; the entry $(1,1)$ is $(u,v)C(s)x'=a(s)(ux_1+vx_3)=a(s)$; and the entry $(2,2)$ is $d^{-1}p'C(s)(-v,u)=c(s)d^{-1}(up_3-vp_1)=c(s)$.
\end{proof}
\leanref{Development/AutomorphicArithmetic/AutomorphicVarieties/ModularityArgument/UpperTriangularPlaces/TriOfLines.lean}{UpperTheory.V.Z.Keep.OrdFlag2.flag_block}
\leanref{Development/AutomorphicArithmetic/AutomorphicVarieties/ModularityArgument/UpperTriangularPlaces/TriOfLines.lean}{UpperTheory.V.Z.Keep.OrdFlag2.tri_of_lines}

\begin{lemma}\label{sha-rigidity-lagrange}
Let $C$ be a local ring, $H$ a group, $n\ge1$, $r:H\to M_n(C)$ a map with $r(1)=1$ and $r(gh)=r(g)r(h)$, and $\psi_1,\dots,\psi_n:H\to C^\times$ characters. Write $\bar y$ for the reduction of $y$ modulo the maximal ideal of $C$. Suppose that
\[
 \big(r(g_1)-\psi_1(g_1)\big)\big(r(g_2)-\psi_2(g_2)\big)\cdots\big(r(g_n)-\psi_n(g_n)\big)=0\qquad\text{for all }g_1,\dots,g_n\in H,
\]
and that some $g_0\in H$ has $\psi_i(g_0)-\psi_j(g_0)\in C^\times$ for $i\ne j$ and the reduction of the characteristic polynomial of $r(g_0)$ equal to $\prod_i(X-\bar\psi_i(g_0))$.
Then there is $g\in GL_n(C)$ such that, for every $s\in H$, $gr(s)g^{-1}$ is upper triangular with diagonal entries $\psi_1(s),\dots,\psi_n(s)$.
\end{lemma}
\begin{proof}
Write $a_i=\psi_i(g_0)$ and $P=r(g_0)$, and let $e_i=\prod_{j\ne i}(a_i-a_j)^{-1}(P-a_j)$ be the Lagrange idempotents of $P$.
\begin{enumerate}
\item[(1)] \emph{Idempotents.} Taking all $g_i=g_0$ gives $\prod_i(P-a_i)=0$. Hence $Pe_c=a_ce_c$ for every $c$, so $(P-a_j)e_c=e_c(P-a_j)=(a_c-a_j)e_c$ for all $j,c$. For $i\ne j$, $e_ie_j$ is a multiple of $\prod_k(P-a_k)$, so $e_ie_j=0$. Lagrange interpolation at the $a_i$, whose differences are units, gives $\sum_ie_i=1$; so $e_i^2=e_i$.
\item[(2)] \emph{The blocks.} Let $s\in H$ and $1\le i\le n$, and let $c,d$ be indices with $c\ne j$ for every $j<i$ and $d\ne j$ for every $j>i$. Apply the relation with $g_i=s$ and $g_j=g_0$ for $j\ne i$, multiplied by $e_c$ on the left and by $e_d$ on the right. By (1), the factors with $j<i$ turn $e_c$ into $(a_c-a_j)e_c$ and those with $j>i$ turn $e_d$ into $(a_d-a_j)e_d$, with units $a_c-a_j$ and $a_d-a_j$; so $e_c\big(r(s)-\psi_i(s)\big)e_d=0$.
For $b<a$, taking $i=c=a$ and $d=b$ gives $e_ar(s)e_b=\psi_a(s)e_ae_b=0$. Taking $c=d=i$ gives $e_ir(s)e_i=\psi_i(s)e_i$.
\item[(3)] \emph{Nonvanishing.} Each $e_i$ is nonzero. Its reduction is the Lagrange idempotent $\prod_{j\ne i}(\bar a_i-\bar a_j)^{-1}(\bar P-\bar a_j)$ of $\bar P$. Since $\bar a_i$ is a root of the characteristic polynomial of $\bar P$, there is $w\ne0$ over the residue field with $\bar Pw=\bar a_iw$, and the reduction of $e_i$ maps $w$ to $w$.
\item[(4)] \emph{Conjugation.} By (1), $C^n=\bigoplus_ie_iC^n$. Each $e_iC^n$ is a direct summand, hence projective, hence free over the local ring $C$, and it is nonzero by (3); as there are $n$ of them, each is free of rank $1$. Let $f_i$ generate $e_iC^n$, and let $g$ be the inverse of the matrix with columns $f_1,\dots,f_n$. Then $e_if_j=\delta_{ij}f_j$, so $ge_ig^{-1}$ is the matrix unit $E_{ii}$.
\item[(5)] For $s\in H$ and indices $a,b$, $(gr(s)g^{-1})_{ab}=\big(E_{aa}\,gr(s)g^{-1}E_{bb}\big)_{ab}=\big(g\,e_ar(s)e_b\,g^{-1}\big)_{ab}$. By (2) this is $0$ for $b<a$, and $\psi_a(s)(E_{aa})_{aa}=\psi_a(s)$ for $a=b$.
\end{enumerate}
\end{proof}
\leanref{Development/AutomorphicArithmetic/AutomorphicVarieties/FrobeniusSmoothVectors/FlagLoc.lean}{UpperTheory.V.Fr.flag_loc}
\leanref{Development/AutomorphicArithmetic/AutomorphicVarieties/AlgebraTools/LagTriangular.lean}{UpperTheory.V.AlgebraTools.lag_triangular}
\leanref{Development/AutomorphicArithmetic/AutomorphicVarieties/AlgebraTools/LagNeZero.lean}{UpperTheory.V.AlgebraTools.lag_ne_zero}
\leanref{Development/AutomorphicArithmetic/AutomorphicVarieties/AlgebraTools/CoiConjSingle.lean}{UpperTheory.V.AlgebraTools.coi_conj_single}

\begin{lemma}[Clifford \cite{CR62}]\label{sha-rigidity-clifford}
Let $G$ be a group, $U$ a normal subgroup of $G$, $\kappa$ an algebraically closed field, $n$ a prime number and $r:G\to GL_n(\kappa)$ an irreducible representation. Then one of the following holds.
\begin{enumerate}
\item The restriction of $r$ to $U$ is irreducible.
\item $r(u)$ is a scalar matrix for every $u\in U$.
\item There are a basis $b_1,\dots,b_n$ of $\kappa^n$, pairwise distinct characters $\lambda_1,\dots,\lambda_n$ of $U$ with $r(u)b_i=\lambda_i(u)b_i$ for $u\in U$, and a homomorphism $\pi$ from $G$ to the group of permutations of $\{1,\dots,n\}$ whose image acts transitively, such that $r(g)b_i$ is a multiple of $b_{\pi(g)(i)}$ for all $g\in G$ and all $i$.
\end{enumerate}
\end{lemma}
\begin{proof}
By Clifford's theorem \cite[(49.2)]{CR62}, whose proof applies to any group acting irreducibly on a finite-dimensional space, the restriction of $r$ to $U$ is semisimple, and $G$ permutes its isotypic components $V_1,\dots,V_t$ transitively; so they all have the same dimension $em$, where $m$ is the dimension of the irreducible constituents and $e$ their multiplicity. Then $n=tem$, and since $n$ is prime, one of $t,e,m$ is $n$ and the other two are $1$.
If $m=n$, the restriction to $U$ is irreducible. If $e=n$, the restriction to $U$ is $n$ copies of one character, so $U$ acts by scalars. If $t=n$, the $V_i$ are lines $\kappa b_i$ on which $U$ acts through pairwise distinct characters $\lambda_i$, and $r(g)V_i=V_{\pi(g)(i)}$ defines $\pi$.
\end{proof}
\leanref{Development/AutomorphicArithmetic/AutomorphicVarieties/CurveDVRSpecialization/CliffPrime.lean}{UpperTheory.V.Pr.cliffPrime}

\begin{lemma}\label{sha-rigidity-characters}
Let $G$ be a group, $k$ a field, $n\ge1$, and $\chi_1,\dots,\chi_n$ and $\chi'_1,\dots,\chi'_n$ characters $G\to k^\times$ such that $\prod_i(X-\chi_i(g))=\prod_i(X-\chi'_i(g))$ for every $g\in G$. Then there is a permutation $\sigma$ of $\{1,\dots,n\}$ with $\chi'_i=\chi_{\sigma(i)}$ for all $i$.
\end{lemma}
\begin{proof}
For a character $\phi$ let $c_\phi$ be the number of $i$ with $\chi_i=\phi$ minus the number of $i$ with $\chi'_i=\phi$; it is nonzero for only finitely many $\phi$. For every $g$ the $\chi_i(g)$ and the $\chi'_i(g)$ are the same elements of $k$ counted with multiplicity, so $\sum_{\phi(g)=t}c_\phi=0$ in $\mathbf Z$ for every $t\in k$, and $\sum_\phi c_\phi\phi(g)=0$ in $k$.
By Artin's theorem on the independence of characters \cite[Chapter~VI, Theorem~4.1]{Lan02}, every $c_\phi$ is $0$ in $k$. If $k$ has characteristic $0$, the $c_\phi$ vanish. If it has characteristic $p>0$, they are divisible by $p$, and the integers $c_\phi/p$ satisfy the same relations; by induction every $c_\phi$ is divisible by every power of $p$, so again the $c_\phi$ vanish. So each character occurs as often among the $\chi_i$ as among the $\chi'_i$, which gives $\sigma$.
\end{proof}
\leanref{Development/AutomorphicArithmetic/AutomorphicVarieties/CurveDVRSpecialization/CharPerm.lean}{UpperTheory.V.Pr.charPerm}

\begin{lemma}\label{sha-rigidity-normal}
Let $G$ be a group, $\kappa$ an algebraically closed field, $k$ a field, $n$ a prime number and $r:G\to GL_n(\kappa)$ an irreducible representation.
Let $A$ be a commutative ring with an injective homomorphism $j:A\to\kappa$ and a homomorphism $\phi:A\to k$, and let $\bar\chi_1,\dots,\bar\chi_n:G\to k^\times$ be characters such that, for every $g\in G$, the characteristic polynomial of $r(g)$ is the image under $j$ of a polynomial over $A$ whose image under $\phi$ is $\prod_i(X-\bar\chi_i(g))$.
Let $D$ be a group, $\iota:D\to G$ a homomorphism and $\psi_1,\dots,\psi_n:D\to\kappa^\times$ characters such that the characteristic polynomial of $r(\iota(d))$ is $\prod_i(X-\psi_i(d))$ for every $d\in D$. Let $U$ be a normal subgroup of $G$, and suppose that:
\begin{enumerate}
\item for $a\ne b$ there is $d\in D$ with $\iota(d)\in U$ and $\psi_a(d)\ne\psi_b(d)$;
\item there are $a,b$ and $d\in D$ with $\bar\chi_a(\iota(d))\ne\bar\chi_b(\iota(d))$.
\end{enumerate}
Then the restriction of $r$ to $U$ is irreducible.
\end{lemma}
\begin{proof}
Lemma~\ref{sha-rigidity-clifford} applies to $r$ and $U$; we show that its cases (2) and (3) do not occur.

\emph{Case (2).} As $n\ge2$, (1) gives $d$ with $\iota(d)\in U$ and $\psi_1(d)\ne\psi_2(d)$. If $r(\iota(d))$ is a scalar $t$, its characteristic polynomial is $(X-t)^n$, so $\psi_1(d)=t=\psi_2(d)$, a contradiction.

\emph{Case (3).} Let $b_i$, $\lambda_i$ and $\pi$ be as there, and define $c(g,i)\in\kappa^\times$ by $r(g)b_i=c(g,i)b_{\pi(g)(i)}$.
\begin{enumerate}
\item[(a)] Then $c(1,i)=1$ and $c(gh,i)=c(g,\pi(h)(i))\,c(h,i)$. For $u\in U$, $r(u)b_i=\lambda_i(u)b_i$, so $\pi(u)=1$: $U$ lies in the kernel $G_0$ of $\pi$. On $G_0$, $h\mapsto c(h,i)$ is a character, and the $c(h,i)$ are the eigenvalues of $r(h)$. For $g\in G$ and $h\in G_0$, (a) applied to $(ghg^{-1})g=gh$ gives $c(ghg^{-1},\pi(g)(i))=c(h,i)$.
\item[(b)] \emph{$\pi(\iota(d))=1$ for every $d\in D$.} Let $D_U=\iota^{-1}(U)$. For $d\in D_U$ the eigenvalues of $r(\iota(d))$ are the $c(\iota(d),i)$ and the $\psi_i(d)$; by Lemma~\ref{sha-rigidity-characters} for the group $D_U$, there is a permutation $\sigma$ with $c(\iota(d),i)=\psi_{\sigma(i)}(d)$ for $d\in D_U$ and all $i$.
Suppose that $i'=\pi(\iota(d))(i)\ne i$ for some $d$ and $i$. By (1) there is $u\in D_U$ with $\psi_{\sigma(i')}(u)\ne\psi_{\sigma(i)}(u)$. But $dud^{-1}\in D_U$, as $U$ is normal, and by (a)
\[
 \psi_{\sigma(i')}(u)=\psi_{\sigma(i')}(dud^{-1})=c\big(\iota(d)\iota(u)\iota(d)^{-1},i'\big)=c(\iota(u),i)=\psi_{\sigma(i)}(u),
\]
a contradiction.
\item[(c)] \emph{Reduction.} Let $S$ be the integral closure of $j(A)$ in $\kappa$. The homomorphism $j(A)\to k\to\bar k$ induced by $\phi$, for an algebraic closure $\bar k$ of $k$, extends to a homomorphism $\Phi:S\to\bar k$ \cite[Chapter~5, Exercise~2]{AM69}. For $h\in G_0$ the characteristic polynomial $\prod_i(X-c(h,i))$ of $r(h)$ is $j(p)$ for a polynomial $p$ over $A$, monic because $j$ is injective; so the $c(h,i)$ lie in $S$, and $\prod_i(X-\Phi(c(h,i)))$ is the image of $p$ under $\phi$, which is $\prod_i(X-\bar\chi_i(h))$. The maps $h\mapsto\Phi(c(h,i))$ are characters of $G_0$, so by Lemma~\ref{sha-rigidity-characters} there is a permutation $\tau$ with $\Phi(c(h,i))=\bar\chi_{\tau(i)}(h)$ for $h\in G_0$ and all $i$.
\item[(d)] By (2) there are $a,a'$ and $d$ with $\bar\chi_a(\iota(d))\ne\bar\chi_{a'}(\iota(d))$. As $\pi(G)$ is transitive, some $g\in G$ has $\pi(g)(\tau^{-1}(a'))=\tau^{-1}(a)$. By (b), $\iota(d)$ and $h=g\iota(d)g^{-1}$ lie in $G_0$, and by (a) and (c)
\[
 \bar\chi_a(\iota(d))=\bar\chi_a(h)=\Phi\big(c(h,\pi(g)(\tau^{-1}(a')))\big)=\Phi\big(c(\iota(d),\tau^{-1}(a'))\big)=\bar\chi_{a'}(\iota(d)),
\]
a contradiction.
\end{enumerate}
So case (1) of Lemma~\ref{sha-rigidity-clifford} holds.
\end{proof}
\leanref{Development/AutomorphicArithmetic/AutomorphicVarieties/CurveDVRSpecialization/IrrNormal.lean}{UpperTheory.V.Pr.irrNormal}

\begin{lemma}\label{sha-rigidity-small}
Let $G$ be a compact topological group and $G'$ a normal subgroup. Let $\kappa$, $k$, the prime $n$, the irreducible representation $r:G\to GL_n(\kappa)$, $A$, $j$, $\phi$, the $\bar\chi_i$, $D$, $\iota$ and the $\psi_i$ be as in Lemma~\ref{sha-rigidity-normal}, and let $\bar\psi_1,\dots,\bar\psi_n:D\to k^\times$ be characters such that for all $i$ and $d$ some $y\in A$ has $j(y)=\psi_i(d)$ and $\phi(y)=\bar\psi_i(d)$. Suppose that some $d_0\in D$ has $\iota(d_0)\in G'$ and $\psi_a(d_0)^m\ne\psi_b(d_0)^m$ for $a\ne b$ and $m\ge1$, and that $\bar\psi_a(d_1)\ne\bar\psi_b(d_1)$ for some $a,b$ and $d_1\in D$.
Then for every open subgroup $U$ of $G$ the restriction of $r$ to $U\cap G'$ is irreducible.
\end{lemma}
\begin{proof}
As $G$ is compact, $U$ has finite index, and so has the intersection $N$ of its conjugates, which is normal; let $e$ be its index. Then $N\cap G'$ is normal, and it suffices to show that $r$ is irreducible on it, by Lemma~\ref{sha-rigidity-normal}.
For (1) of that lemma, let $a\ne b$ and $d=d_0^e$. Then $\iota(d)=\iota(d_0)^e$ lies in $N$, whose quotient has order $e$, and in $G'$; and $\psi_a(d)=\psi_a(d_0)^e\ne\psi_b(d_0)^e=\psi_b(d)$.
For (2), suppose that $\bar\chi_a(\iota(d_1))=\bar\chi_b(\iota(d_1))$ for all $a,b$. Choose $y_i\in A$ with $j(y_i)=\psi_i(d_1)$ and $\phi(y_i)=\bar\psi_i(d_1)$. The polynomial $p$ over $A$ with $j(p)$ the characteristic polynomial $\prod_i(X-\psi_i(d_1))$ of $r(\iota(d_1))$ is $\prod_i(X-y_i)$, since $j$ is injective. So $\prod_i(X-\bar\psi_i(d_1))=\phi(p)=\prod_i(X-\bar\chi_i(\iota(d_1)))=(X-\bar\chi_a(\iota(d_1)))^n$, and all the $\bar\psi_i(d_1)$ are equal, a contradiction.
\end{proof}
\leanref{Development/AutomorphicArithmetic/AutomorphicVarieties/TheoremsMuLinked/CharacteristicPolynomials/IrrSmall.lean}{UpperTheory.V.Th.Det.irrSmall}

\begin{lemma}\label{sha-rigidity-coset}
Let $G$ be a group, $\kappa$ a field, $1\le n\le3$ and $r:G\to GL_n(\kappa)$ a representation such that no subgroup of finite index of $G$ stabilizes a nonzero proper subspace of $\kappa^n$. Let $u\in G$ be such that the characteristic polynomial of $r(u)$ is $\prod_i(X-\alpha_i)$ with $\alpha_i^m\ne\alpha_j^m$ for $i\ne j$ and $m\ge1$. Then for every $y\in GL_n(\kappa)$ there is $h\in G$ such that the characteristic polynomial of $y\,r(h)$ is separable.
\end{lemma}
\begin{proof}
The $\alpha_i$ are distinct, so $r(u)$ is diagonalizable; after conjugating $r$ and $y$ by the same matrix, $r(u)=\Delta=\operatorname{diag}(\alpha_1,\dots,\alpha_n)$. The proof uses three facts.
\begin{enumerate}
\item[(F1)] \emph{Some $b\in G$ has all entries of $r(b)$ and of $r(b)^{-1}$ nonzero.}
Let $E=M_n(\kappa)^2$, $\varphi(g)=(r(g),r(g^{-1}))$ and $L_h(X,Y)=(r(h)X,Yr(h^{-1}))$, so that $\varphi(hg)=L_h\varphi(g)$ with $L_h$ linear. Call a subset $Z\subseteq G$ \emph{irreducible} if $Z=\varphi^{-1}(W)$ for a subspace $W$ of $E$, and $Z\subseteq\varphi^{-1}(B)$ for some $B\in T$ whenever $Z$ is covered by the sets $\varphi^{-1}(B)$, $B\in T$, for a finite set $T$ of subspaces. The empty set is not irreducible, and $hZ$ is irreducible with $Z$, through the $L_h$.
By induction on $W$, every $\varphi^{-1}(W)$ is a finite union of irreducible sets: if $\varphi^{-1}(W)$ is not irreducible, it is covered by finitely many $\varphi^{-1}(B)$, none containing it, and it is the union of the $\varphi^{-1}(W\cap B)$ with $W\cap B\subsetneq W$. Let $\mathcal C$ be such a finite family for $W=E$. Every irreducible set lies in a member of $\mathcal C$. Let $Z\in\mathcal C$ be maximal for inclusion. For $h\in G$, $hZ\subseteq Z_1\in\mathcal C$ and $h^{-1}Z_1\subseteq Z_2\in\mathcal C$, so $Z\subseteq Z_2$, $Z=Z_2$ and $hZ=Z_1\in\mathcal C$. So the stabilizer $H$ of $Z$ has finite index, and $Hx\subseteq Z$ for any $x\in Z$.
Suppose that (F1) fails. Then $G$ is covered by the finitely many sets $\varphi^{-1}(B)$ where $B$ is the subspace on which the entry $(i,j)$ of the first or of the second component vanishes; so $Z$ lies in one of them. In the first case $r(hx)_{ij}=0$ for $h\in H$: the span of the $r(h)r(x)e_j$, $h\in H$, is stable under $H$, nonzero and contained in the hyperplane of the vectors with $i$-th coordinate $0$, which contradicts the hypothesis on $r$. In the second case $(r(x^{-1})r(h))_{ij}=0$ for $h\in H$, as $(hx)^{-1}=x^{-1}h^{-1}$; so the span of the $r(h)e_j$, $h\in H$, is stable under $H$, nonzero and contained in the kernel of the nonzero linear form $v\mapsto(r(x^{-1})v)_i$, which is again impossible.
\item[(F2)] \emph{Let $\beta_{s,1},\dots,\beta_{s,n}\in\kappa^\times$ and $c_{s,1},\dots,c_{s,n}\in\kappa$, for $s$ in a finite set, be such that $\beta_{s,i}^m\ne\beta_{s,j}^m$ for $i\ne j$ and $m\ge1$, and for each $s$ some $c_{s,i}$ is nonzero. Then there is $m\ge0$ with $\sum_ic_{s,i}\beta_{s,i}^m\ne0$ for all $s$.}
Otherwise, colour each $m$ by an $s$ at which the sum vanishes. By van der Waerden's theorem \cite{Wae27} there are $s$, $a\ge1$ and $b\ge0$ with $\sum_ic_{s,i}\beta_{s,i}^b(\beta_{s,i}^a)^t=0$ for $0\le t<n$. The $\beta_{s,i}^a$ are distinct, so the Vandermonde matrix is invertible and every $c_{s,i}\beta_{s,i}^b$ vanishes, a contradiction.
\item[(F3)] \emph{If $Y\in M_n(\kappa)$ has all its principal minors nonzero, then the characteristic polynomial of $Y\Delta^m$ is separable for some $m\ge0$.}
Since no $\alpha_i/\alpha_j$ with $i\ne j$ is a root of unity, these ratios are nonzero in the $\mathbf Q$-vector space $\kappa^\times\otimes\mathbf Q$, and a linear form on it gives a homomorphism $w:\kappa^\times\to\mathbf Q$ with $w(\alpha_i)\ne w(\alpha_j)$ for $i\ne j$. A function $f(m)=\sum_\beta c_\beta\beta^m$, for finitely many distinct $\beta\in\kappa^\times$, is not identically zero if the $\beta$ with $c_\beta\ne0$ of largest $w(\beta)$ is unique, by the independence of the characters $m\mapsto\beta^m$ \cite[Chapter~VI, Theorem~4.1]{Lan02}. The discriminant of the characteristic polynomial of $Y\Delta^m$ is such a function of $m$, and the characteristic polynomial is separable where it does not vanish.
For $n\le1$ there is nothing to show. For $n=2$, with $w(\alpha_1)>w(\alpha_2)$ after renumbering, the discriminant is $(Y_{11}\alpha_1^m+Y_{22}\alpha_2^m)^2-4\det(Y)(\alpha_1\alpha_2)^m$, whose term of largest weight is $Y_{11}^2(\alpha_1^2)^m$. For $n=3$, with $w(\alpha_1)>w(\alpha_2)>w(\alpha_3)$, the characteristic polynomial is $X^3-e_1X^2+e_2X-e_3$ with $e_1=\sum_iY_{ii}\alpha_i^m$, $e_2=\sum_{i<j}Y_{\{i,j\}}(\alpha_i\alpha_j)^m$, where $Y_{\{i,j\}}$ is the principal minor of $Y$ on the rows and columns $i$ and $j$, and $e_3=\det(Y)(\alpha_1\alpha_2\alpha_3)^m$. In its discriminant $e_1^2e_2^2-4e_2^3-4e_1^3e_3+18e_1e_2e_3-27e_3^2$ the term of largest weight is $Y_{11}^2Y_{\{1,2\}}^2(\alpha_1^4\alpha_2^2)^m$, of weight $4w(\alpha_1)+2w(\alpha_2)$; all other terms have smaller weight.
\end{enumerate}
Let $b$ be as in (F1). The rows of $y$ and the columns of $y^{-1}$ are nonzero. By (F2), applied to the $2n$ sums $\sum_jy_{ij}r(b)_{ji}\alpha_j^m$ and $\sum_j(r(b)^{-1})_{ij}(y^{-1})_{ji}\alpha_j^{-m}$, which have nonzero coefficients by (F1), there is $m$ such that all diagonal entries of $Y=y\Delta^mr(b)$ and of $Y^{-1}=r(b)^{-1}\Delta^{-m}y^{-1}$ are nonzero.
Then every principal minor of $Y$, on a set $S$ of rows and columns, is nonzero. For $|S|=1$ it is a diagonal entry, and for $S$ empty it is $1$. Otherwise, as $n\le3$, $S$ omits at most one index. If the minor vanished, there would be $v\ne0$ supported on $S$ with $(Yv)_i=0$ for $i\in S$. If $S$ omits an index $k$, then $Yv$ is a multiple of the $k$-th unit vector, and $v=Y^{-1}(Yv)$ gives $0=v_k=(Y^{-1})_{kk}(Yv)_k$, so $(Yv)_k=0$. In either case $Yv=0$, which is impossible.
By (F3) there is $m'$ with the characteristic polynomial of $Y\Delta^{m'}$ separable. Then $h=u^mbu^{m'}$ has $y\,r(h)=y\Delta^mr(b)\Delta^{m'}=Y\Delta^{m'}$.
\end{proof}
\leanref{Development/AutomorphicArithmetic/AutomorphicVarieties/CurveDVRSpecialization/RegCoset.lean}{UpperTheory.V.Pr.regCoset}
\leanref{Development/AutomorphicArithmetic/IntegralCasts/CapelliIdentities/IrreducibleCoverForms.lean}{UpperTheory.Reformulated.B.StronglyIrreducibleFreeRatios.fullEx}
\leanref{Development/AutomorphicArithmetic/AutomorphicVarieties/CurveDVRSpecialization/ExpSums.lean}{UpperTheory.V.Pr.expSums}
\leanref{Development/AutomorphicArithmetic/AutomorphicVarieties/CurveDVRSpecialization/SepOfMinors.lean}{UpperTheory.V.Pr.sepOfMinors}

\begin{lemma}\label{sha-rigidity-separable}
Let $G$ be a profinite group, $G'$ a subgroup, $\ell$ a prime, $1\le n\le3$ and $r:G\to GL_n(\bar{\mathbf Q}_\ell)$ a continuous representation. Suppose that the restriction of $r$ to $U\cap G'$ is irreducible for every open subgroup $U$ of $G$, and that some $u\in G'$ has the characteristic polynomial of $r(u)$ equal to $\prod_i(X-\alpha_i)$ with $\alpha_i^m\ne\alpha_j^m$ for $i\ne j$ and $m\ge1$.
Then every open subset of $G$ that meets $G'$ contains an element $\sigma$ of $G'$ at which the characteristic polynomial of $r(\sigma)$ is separable.
\end{lemma}
\begin{proof}
\begin{enumerate}
\item[(1)] For every linear form $f$ on $M_n(\bar{\mathbf Q}_\ell)$ the set of $g$ with $f(r(g))=0$ is closed, since $f(r(g))$ is a linear combination of the entries of $r(g)$, which are continuous.
\item[(2)] Let $V$ be an open subgroup of $G$. It has finite index, so $u^e\in V$ for some $e\ge1$. The characteristic polynomial of $r(u^e)$ is $\prod_i(X-\alpha_i^e)$, and the $\alpha_i^e$ have the property of the $\alpha_i$.
\item[(3)] Let $O$ be open, and $s\in O\cap G'$. The open subgroups form a basis of neighbourhoods of $1$ in the profinite group $G$, so there is an open subgroup $V$ with $sV\subseteq O$.
\item[(4)] No subgroup $S$ of finite index of $V\cap G'$ stabilizes a nonzero proper subspace $W$. Indeed, we may take $S$ to be the stabilizer of $W$, which is the intersection of the sets of $g$ with $f(r(g)w)=0$ for $w\in W$ and linear forms $f$ vanishing on $W$; it is closed by (1). Being closed of finite index, $S$ is open in $V\cap G'$, so $S$ contains $V'\cap G'$ for an open subgroup $V'$ of $G$, on which $r$ is irreducible. So $W$ is $0$ or everything.
\item[(5)] By (2), (4) and Lemma~\ref{sha-rigidity-coset}, applied to the restriction of $r$ to $V\cap G'$, the element $u^e$ and $y=r(s)$, there is $h\in V\cap G'$ such that $r(s)r(h)=r(sh)$ has separable characteristic polynomial. Then $\sigma=sh$ lies in $O\cap G'$.
\end{enumerate}
\end{proof}
\leanref{Development/AutomorphicArithmetic/AutomorphicVarieties/TheoremsMuLinked/CharacteristicPolynomials/CharpolyPow.lean}{UpperTheory.V.Th.Det.n3small}
\leanref{Development/AutomorphicArithmetic/AutomorphicVarieties/CurveDVRSpecialization/StrongIrrOfOpen.lean}{UpperTheory.V.Pr.strongIrrOfOpen}

\begin{lemma}\label{sha-rigidity-windowed-flags}
Let $F\subseteq L$ be number fields inside an algebraic closure $\bar F$ of $F$, let $\ell$ be a prime, let $n\ge1$, and let $r:\operatorname{Gal}(\bar F/F)\to GL_n(\bar{\mathbf Q}_\ell)$ be a continuous representation, with restriction $r_L$ to $\Gamma_L=\operatorname{Gal}(\bar F/L)$. For a place $A$ of $\bar F$ let $D_A\subseteq\operatorname{Gal}(\bar F/F)$ be its decomposition group.
\begin{enumerate}
\item Suppose that for every place $A$ of $\bar F$ above $\ell$ the restriction of $r$ to $D_A$ is conjugate to an upper triangular representation. Then for every such $A$ the restriction of $r_L$ to $D_A\cap\Gamma_L$, the decomposition group of $A$ in $\Gamma_L$, is conjugate to an upper triangular representation.
\item Let $w$ be a prime of $L$ above $\ell$ and $v$ the prime of $F$ below $w$, and suppose that the restriction of $r$ to a decomposition group at $v$ is conjugate to an upper triangular representation. If $r_L$ is $\ell$-distinct at $w$ (Definition~\ref{sha-eigensystem-distinctness}), then $r$ is $\ell$-distinct at $v$.
\end{enumerate}
\end{lemma}
\begin{proof}
(1) An upper triangular form of $r$ on $D_A$ restricts to an upper triangular form of $r_L$ on $D_A\cap\Gamma_L$, whose diagonal characters are the restrictions of those of $r$.

(2) Let $A$ be a place above $v$ such that $r$ has on $D_A$ an upper triangular form, with diagonal characters $\kappa_1,\dots,\kappa_n$. Let $B$ be a place above $w$ such that $r_L$ has on $D_B\cap\Gamma_L$ an upper triangular form whose diagonal characters $\theta_1,\dots,\theta_n$ satisfy $\theta_a\ne\theta_b\chi_{\mathrm{cyc}}$ for all $a,b$, and choose $\gamma\in\operatorname{Gal}(\bar F/F)$ with $\gamma A=B$.
As in the proof of Lemma~\ref{sha-eigensystem-flags}(2), $r_L$ has on $D_B\cap\Gamma_L$ a second upper triangular form, with diagonal characters $\kappa'_i(s)=\kappa_i(\gamma^{-1}s\gamma)$, and each $\kappa'_i$ is one of the $\theta_a$.
If $\kappa_i=\kappa_j\chi_{\mathrm{cyc}}$ on $D_A$ for some $i,j$, then composing with $s\mapsto\gamma^{-1}s\gamma$, which preserves $\chi_{\mathrm{cyc}}$, gives $\theta_a=\theta_b\chi_{\mathrm{cyc}}$ for the $a$ and $b$ with $\kappa'_i=\theta_a$ and $\kappa'_j=\theta_b$, which is false.
So the form of $r$ on $D_A$ has $\kappa_i\ne\kappa_j\chi_{\mathrm{cyc}}$ for all $i,j$, and $r$ is $\ell$-distinct at $v$.
\end{proof}
\leanref{Development/AutomorphicArithmetic/AutomorphicVarieties/ModularityArgument/AutomorphyLifting/PolarizedLinkage/FlagLaws.lean}{UpperTheory.V.Z.Win.flagged_res}
\leanref{Development/AutomorphicArithmetic/AutomorphicVarieties/ModularityArgument/AutomorphyLifting/PolarizedLinkage/FlagLaws.lean}{UpperTheory.V.Z.Win.unl_down}

\begin{lemma}\label{sha-rigidity-windowed-local}
Let $K\subset M$ be totally complex number fields with $M$ a Galois extension of $K$ of prime degree, let $n\ge2$, and let $\pi$ and $P$ be cuspidal automorphic representations of $GL_n(\mathbf A_K)$ and $GL_n(\mathbf A_M)$ (Definition~\ref{gal-polarized}) such that $P$ is a weak base change of $\pi$ (Definition~\ref{gal-depth-zero-lift-weak}). For a finite place $u$ let $\varpi_u$ be a uniformizer and $q_u$ the order of the residue field.
Let $v$ be a finite place of $K$ such that the local components of $\pi$ at $v$ have supercuspidal support consisting of characters, let $w$ be a place of $M$ above $v$, and let $N=N_{M_w/K_v}$.
\begin{enumerate}
\item The local components of $P$ at $w$ have supercuspidal support consisting of characters, each of the form $\xi\circ N$ for a character $\xi$ in the supercuspidal support of a local component of $\pi$ at $v$.
\item Let $t\in\mathbf R$, and suppose that every character $\xi$ in the supercuspidal support of a local component of $\pi$ at $v$ satisfies $q_v^{t-1/2}<|\xi(\varpi_v)|<q_v^{t+1/2}$. Then every character $\zeta$ in the supercuspidal support of a local component of $P$ at $w$ satisfies $q_w^{t-1/2}<|\zeta(\varpi_w)|<q_w^{t+1/2}$.
\end{enumerate}
\end{lemma}
\begin{proof}
(1) This follows from the cyclic base change of Arthur and Clozel \cite[Chapter~3]{AC89}. By strong multiplicity one (Theorem~\ref{lit-jacquet-shalika}(1)), $P$ is the base change of $\pi$ constructed there, whose local component at $w$ is the local base change of that of $\pi$ at $v$. A local component of $\pi$ at $v$ is a subquotient of the representation induced from the characters $\xi_1,\dots,\xi_n$ of its supercuspidal support. Local base change commutes with parabolic induction, taking a character $\xi$ to $\xi\circ N$ \cite{AC89}; so the local components of $P$ at $w$ are subquotients of the representation induced from the $\xi_i\circ N$, which form their supercuspidal support \cite{Zel80}.

(2) By (1), $\zeta=\xi\circ N$ for such a character $\xi$. Let $f$ be the residue degree of $w$ over $v$. Then $q_w=q_v^f$, and $N(\varpi_w)=u\varpi_v^f$ for a unit $u$ of $K_v$. The character $\xi$ is continuous on the compact group of units, so $|\xi(u)|=1$ and $|\zeta(\varpi_w)|=|\xi(\varpi_v)|^f$. Raising the bounds for $\xi$ to the power $f$ gives those for $\zeta$.
% The up half of ImaginaryQuadraticEtaCharacter.winUpDown. (1) is the first halves of Z.Gam.upDown (TorAt) and upDownX (the characters are xi composed with the norm), which Alpöge proves through gamma factors (upDownOf, matchBCOf, matchBCX), in the directories GammaLaurentFactors and LeviSmoothLocalization/PrincipalSeriesMatching (about 7,700 lines together, not all used here); the text cites the published base change instead. (2) is normVal with the private win_pow_iff.
\end{proof}
\leanref{Development/AutomorphicArithmetic/AutomorphicVarieties/ModularityArgument/LeviSmoothLocalization/WinUpDown.lean}{UpperTheory.V.Z.ImaginaryQuadraticEtaCharacter.WinUpDown}
\leanref{Development/AutomorphicArithmetic/AutomorphicVarieties/ModularityArgument/LeviSmoothLocalization/PrincipalSeriesMatching/WinUpDown.lean}{UpperTheory.V.Z.ImaginaryQuadraticEtaCharacter.winUpDown}
\leanref{Development/AutomorphicArithmetic/AutomorphicVarieties/ModularityArgument/GammaLaurentFactors/CrossG.lean}{UpperTheory.V.Z.Gam.upDown}
\leanref{Development/AutomorphicArithmetic/AutomorphicVarieties/ModularityArgument/LeviSmoothLocalization/PrincipalSeriesMatching/UpDownX.lean}{UpperTheory.V.Z.ImaginaryQuadraticEtaCharacter.upDownX}
\leanref{Development/AutomorphicArithmetic/AutomorphicVarieties/ModularityArgument/LeviSmoothLocalization/PrincipalSeriesMatching/NormVal.lean}{UpperTheory.V.Z.ImaginaryQuadraticEtaCharacter.normVal}

\begin{lemma}\label{sha-rigidity-windowed-up}
Assume the Galois representations as in Theorem~\ref{sha-residually-reducible-automorphy}.
Let $F\subseteq E$ be CM fields inside an algebraic closure $\bar F$ of $F$, with $E$ Galois over $F$ with solvable Galois group. Let $\ell$ be a prime, let $\iota:\bar{\mathbf Q}_\ell\cong\mathbf C$, let $n\ge1$, and let $r:\operatorname{Gal}(\bar F/F)\to GL_n(\bar{\mathbf Q}_\ell)$ be a continuous representation, with restriction $r_E$ to $\operatorname{Gal}(\bar F/E)$. Suppose that:
\begin{enumerate}
\item $r_E$ is irreducible;
\item there are a totally real subfield $F_0$ of $F$, a character $\mu:\operatorname{Gal}(\bar F/F_0)\to\bar{\mathbf Q}_\ell^\times$, attached through $\iota$ to an algebraic Hecke character of $F_0$ and taking the same value at all complex conjugations, a complex conjugation $c\in\operatorname{Gal}(\bar F/F_0)$ and $P\in GL_n(\bar{\mathbf Q}_\ell)$ with ${}^tr(s)\,P\,r(csc^{-1})=\mu(s)P$ for all $s\in\operatorname{Gal}(\bar F/F)$;
\item $r$ is attached through $\iota$ to a regular algebraic cuspidal automorphic representation of $GL_n(\mathbf A_F)$;
\item for every place $A$ of $\bar F$ above $\ell$, the restriction of $r$ to the decomposition group $D_A\subseteq\operatorname{Gal}(\bar F/F)$ is conjugate to an upper triangular representation.
\end{enumerate}
If $r$ is windowed (Theorem~\ref{sha-residually-reducible-automorphy}), then so is $r_E$.
\end{lemma}
\begin{proof}
The proof is by induction on $d=[E:F]$, for all fields $F\subseteq E$ as in the statement and all $\ell$, $\iota$, $n$ and $r$ at once. For a finite place $w$ of $E$, $v$ denotes the place of $F$ below $w$. Condition ($\alpha$) of Theorem~\ref{sha-residually-reducible-automorphy} says that a representation is potentially $\ell$-distinct at every prime above $\ell$ at which it is $\ell$-distinct (Definition~\ref{sha-eigensystem-distinctness}).

\emph{Base change.} Let $F^+$ be the maximal totally real subfield of $F$; it contains $F_0$, and $c$ lies in $\operatorname{Gal}(\bar F/F^+)$, as a complex conjugation fixes every totally real field. The restriction $\mu^+$ of $\mu$ to $\operatorname{Gal}(\bar F/F^+)$ is attached through $\iota$ to the composite with the norm from $F^+$ to $F_0$ of the algebraic Hecke character to which $\mu$ is attached \cite[Chapter~VI]{Neu99}, and it takes the same value at all complex conjugations. Let $M\supseteq F$ be a CM field inside $\bar F$, Galois over $F$ with solvable Galois group, such that the restriction $r_M$ of $r$ to $\operatorname{Gal}(\bar F/M)$ is irreducible. Then Theorem~\ref{sha-siegel-automorphy-base-change}, applied to $F\subseteq M$, $r$, $\mu^+$, $c$ and $P$, and (3) show that $r_M$ is attached through $\iota$ to a regular algebraic cuspidal automorphic representation of $GL_n(\mathbf A_M)$.

\emph{Prime degree.} Suppose that $d$ is prime, and let $r$ satisfy (1) to (4) and be windowed.
\begin{enumerate}
\item[(p1)] \emph{Condition ($\alpha$) for $r_E$.} Let $w$ be a prime of $E$ above $\ell$ at which $r_E$ is $\ell$-distinct. By (4) and Lemma~\ref{sha-rigidity-windowed-flags}(2), $r$ is $\ell$-distinct at $v$. By ($\alpha$) for $r$, it is potentially $\ell$-distinct at $v$, and by Lemma~\ref{sha-eigensystem-flags}(1), $r_E$ is potentially $\ell$-distinct at $w$.
\item[(p2)] \emph{The places of ($\beta$).} Let $w$ be a finite place of $E$ such that $w\nmid\ell$ or $r_E$ is potentially $\ell$-distinct at $w$. Then $v\nmid\ell$ or $r$ is potentially $\ell$-distinct at $v$: if $w\mid\ell$, this is (4) and Lemma~\ref{sha-eigensystem-flags}(2).
\item[(p3)] \emph{Condition ($\beta$) for $n=1$.} A cuspidal automorphic representation of $GL_1(\mathbf A_E)$ is plain at every finite place (Lemma~\ref{sha-hecke-ideal-plain}): its local components are characters, each the only member of its supercuspidal support. So Lemma~\ref{sha-hecke-ideal-plain}(1) gives ($\beta$) for $r_E$, with the bound at every finite place.
\item[(p4)] \emph{A weak base change.} Let $n\ge2$, let $\pi$ be a regular algebraic cuspidal automorphic representation of $GL_n(\mathbf A_F)$ attached through $\iota$ to $r$, which exists by (3), and let $\Pi$ be a regular algebraic cuspidal automorphic representation of $GL_n(\mathbf A_E)$ attached through $\iota$ to $r_E$. Then $\Pi$ is a weak base change of $\pi$ (Definition~\ref{gal-depth-zero-lift-weak}).
Indeed, only finitely many places of $E$ lie above each place of $F$, so for all but finitely many finite places $w$ of $E$, with $N(w)=N(v)^f$, there are $B$ and $B'$ such that $\pi$ has Satake parameters $B$ at $v$ and $r$ has Frobenius eigenvalues $q_v^{(n-1)/2}B$ at $v$, and $\Pi$ has Satake parameters $B'$ at $w$ and $r_E$ has Frobenius eigenvalues $q_w^{(n-1)/2}B'$ at $w$ (Definition~\ref{gal-weakly-attached}).
Let $A$ be a place of $\bar F$ above $w$, let $\sigma\in D_A\cap\operatorname{Gal}(\bar F/E)$ act on the residue field by $z\mapsto z^{q_w}$, and let $\sigma_0\in D_A$ act by $z\mapsto z^{q_v}$. Then $\sigma$ and $\sigma_0^f$ act in the same way, so they differ by an element of the inertia group of $A$, on which $r$ is trivial; hence $r(\sigma^{-1})=r(\sigma_0^{-1})^f$, whose eigenvalues are, through $\iota$, the $(q_v^{(n-1)/2}b)^f=q_w^{(n-1)/2}b^f$ for $b\in B$. These eigenvalues determine $B'=\{b^f:b\in B\}$.
\item[(p5)] \emph{Condition ($\beta$) for $n\ge2$.} Let $\pi$ and $\Pi$ be as in (p4). By ($\beta$) for $r$, applied to $\pi$, there is $t\in\mathbf R$ such that at every finite place $v$ of $F$ with $v\nmid\ell$ or $r$ potentially $\ell$-distinct at $v$, the local components of $\pi$ at $v$ have supercuspidal support consisting of characters $\xi$ with $q_v^{t-1/2}<|\xi(\varpi_v)|<q_v^{t+1/2}$. Let $w$ be a place as in (p2). By (p2) this holds at $v$, and by (p4) and Lemma~\ref{sha-rigidity-windowed-local}, the local components of $\Pi$ at $w$ have supercuspidal support consisting of characters $\zeta$ with $q_w^{t-1/2}<|\zeta(\varpi_w)|<q_w^{t+1/2}$. This is ($\beta$) for $r_E$.
\end{enumerate}

\emph{Degree one.} Suppose that $d=1$, so that $E=F$, and let $r$ satisfy (1) to (4) and be windowed. Steps (p1) to (p4) do not use that $d$ is prime, and they apply to any finite extension of CM fields inside $\bar F$ in place of $F\subseteq E$.
\begin{enumerate}
\item[(o1)] \emph{Condition ($\alpha$) for $r_E$, and ($\beta$) for $n=1$.} These are (p1) and (p3).
\item[(o2)] \emph{An auxiliary field.} Let $n\ge2$. By Lemma~\ref{sha-siegel-step-ideal-descent-updown-auxiliary} there is a CM field $M\supseteq E$, Galois over $E$, of prime degree $p>n$; it may be taken inside $\bar F$. As $E=F$, $M$ is Galois over $F$ of degree $p$, with cyclic Galois group. The group $\operatorname{Gal}(\bar F/M)$ is a normal subgroup of index $p$ of $\operatorname{Gal}(\bar F/E)$, and $r_E$ is irreducible by (1), so the restriction $r_M$ of $r$ to $\operatorname{Gal}(\bar F/M)$ is irreducible by Lemma~\ref{sha-siegel-step-ideal-descent-updown-prime-index}. By the paragraph on base change, $r_M$ is attached through $\iota$ to a regular algebraic cuspidal automorphic representation $Q$ of $GL_n(\mathbf A_M)$.
\item[(o3)] \emph{Up from $F$ to $M$.} Let $\Pi$ be a regular algebraic cuspidal automorphic representation of $GL_n(\mathbf A_E)$ attached through $\iota$ to $r_E$, and $\pi$ one of $GL_n(\mathbf A_F)$ attached through $\iota$ to $r$, which exists by (3). By (p4), applied to $F\subseteq M$ with $r$, $\pi$ and $Q$, and to $E\subseteq M$ with $r_E$, $\Pi$ and $Q$, $Q$ is a weak base change of $\pi$ and of $\Pi$. By ($\beta$) for $r$, applied to $\pi$, there is $t\in\mathbf R$ such that at every finite place $v$ of $F$ with $v\nmid\ell$ or $r$ potentially $\ell$-distinct at $v$, the local components of $\pi$ at $v$ have supercuspidal support consisting of characters $\xi$ with $q_v^{t-1/2}<|\xi(\varpi_v)|<q_v^{t+1/2}$. Let $w$ be a place of $E$ as in (p2); by (p2) this holds at $v$. Let $u$ be a place of $M$ above $w$. By Lemma~\ref{sha-rigidity-windowed-local}, applied to $F\subseteq M$, $\pi$, $Q$, $v$ and $u$, every character $\zeta$ in the supercuspidal support of a local component of $Q$ at $u$ is $\xi\circ N_{M_u/F_v}$ for such a character $\xi$, and $q_u^{t-1/2}<|\zeta(\varpi_u)|<q_u^{t+1/2}$.
\item[(o4)] \emph{Down from $M$ to $E$.} Let $\zeta=\xi\circ N_{M_u/F_v}$ be as in (o3). As $N_{M_u/F_v}=N_{E_w/F_v}\circ N_{M_u/E_w}$, the smooth character $\chi'=\xi^{-1}\circ N_{E_w/F_v}$ of $E_w^\times$ has $\zeta\cdot(\chi'\circ N_{M_u/E_w})=1$, which is unramified. So at every place $u$ of $M$ above $w$ the supercuspidal support of every local component of $Q$ at $u$ consists of characters $\zeta$ with $q_u^{t-1/2}<|\zeta(\varpi_u)|<q_u^{t+1/2}$, each with the property of the conclusion of Lemma~\ref{sha-siegel-step-ideal-descent-updown-local}(2) for $E\subseteq M$. By Lemma~\ref{sha-siegel-step-ideal-descent-updown-local}(3), applied to $E\subseteq M$, which is Galois of prime degree, $\Pi$, $Q$, $t$ and $w$, the supercuspidal support of every local component of $\Pi$ at $w$ consists of characters $\xi'$ with $q_w^{t-1/2}<|\xi'(\varpi_w)|<q_w^{t+1/2}$. This is ($\beta$) for $r_E$, with $\Pi$ and $t$.
\end{enumerate}
As $E=F$ and $r_E=r$, the statement also holds in this case without the auxiliary field.

\emph{Other degrees.} Suppose that $d$ is neither $1$ nor prime.
\begin{enumerate}
\item[(t1)] \emph{An intermediate field.} The group $G=\operatorname{Gal}(E/F)$ is solvable of order $d$, so it has a normal subgroup $H$ with $1\ne H\ne G$: otherwise $G$ would be simple and solvable, hence cyclic of prime order. Let $E'$ be the fixed field of $H$. Then $F\subsetneq E'\subsetneq E$, $E'$ is Galois over $F$ with the solvable group $G/H$, $E$ is Galois over $E'$ with the solvable group $H$, and $[E':F]$ and $[E:E']$ are less than $d$.
The field $E'$ is CM. Let $c_E$ be the complex conjugation of $E$, so that $\varphi(c_Ex)=\overline{\varphi(x)}$ for every embedding $\varphi:E\to\mathbf C$. For $x\in E'$ and $h\in H$, applying this to $\varphi\circ h$ gives $\varphi(h(c_Ex))=\overline{\varphi(x)}=\varphi(c_Ex)$; so $c_E$ maps $E'$ to itself. The elements of $E'$ fixed by $c_E$ are real in every embedding, so they form a totally real field $E'_0$, and $E'$ is $E'_0$ or a totally imaginary quadratic extension of it. It is not totally real, since it contains $F$.
\item[(t2)] By induction, the statement holds for $F\subseteq E'$ and for $E'\subseteq E$. Let $r$ satisfy (1) to (4) and be windowed, and let $r_{E'}$ be its restriction to $\operatorname{Gal}(\bar F/E')$. Since $r_E$ is irreducible, so is $r_{E'}$, and (1) to (4) hold for $F\subseteq E'$ and $r$; so $r_{E'}$ is windowed.
\item[(t3)] Hypotheses (1) to (4) hold for $E'\subseteq E$ and $r_{E'}$, whose restriction to $\operatorname{Gal}(\bar F/E)$ is $r_E$. (1) is hypothesis (1).
(2) holds with the same $F_0$, $\mu$, $c$ and $P$: $F_0$ is a totally real subfield of $E'$, and the relation of (2) holds for $s\in\operatorname{Gal}(\bar F/E')$.
For (3), the paragraph on base change, with $M=E'$, shows that $r_{E'}$ is attached through $\iota$ to a regular algebraic cuspidal automorphic representation of $GL_n(\mathbf A_{E'})$, as $r_{E'}$ is irreducible by (t2).
(4) is Lemma~\ref{sha-rigidity-windowed-flags}(1).
So $r_E$ is windowed, by the statement for $E'\subseteq E$.
\end{enumerate}
% The Lean: automorphy_up_base_change is the first half of Z.Win.Dn.aUpDown_alt, that is Z.Win.aUpDownOf_alt (TAUpDownOf.lean) with weakOfAtt, flagLaws, winUpDown and winOf; its proof is good_all, a strong induction on [E:F] over all fields at once. Prime degree is good_prime: (p1) is first_condition_up (unl_down, pot_up); (p2) is cond_down (pot_down); (p3) is second_condition_trivial_extension (win_one, through winOf and suppFacts); (p4) is weakOfAtt (frobenius_power_compatibility, hecke_eigenvalue_basic_properties); (p5) is second_condition_up with winUpDown. Other degrees: (t1) is UpperTheory.solvable_cm_extension_tower (exists_proper_nontrivial_normal_of_solvable_not_prime_card, rc_of_intermediateField_stable_under_complexConj); (t2) and (t3) are good_trans (galoisRep_irreducibility_criterion, UpperTheory.solvable_polarized_restriction, aut_up through Z.Bc.solC_alt, flagged_res).
% Degree one is the up half of good_one: (o1) is first_condition_up and second_condition_trivial_extension; (o2) to (o4) are second_condition_up_trivial: exists_cm_field_tower with N = n (Lemma sha-siegel-step-ideal-descent-updown-auxiliary), galois_representation_image_spec and galois_representation_irreducible_prime_dimension (Lemma sha-siegel-step-ideal-descent-updown-prime-index), aut_up through Z.Bc.stepC_alt for F in M, weakOfAtt for F in M and for E in M, the up half of winUpDown for F in M, thruTr (the character chi' of (o4), through the transitivity of norms in exists_char and norm_through), and the down half of winUpDown for E in M. In the Lean E is a field with an embedding of F of degree one, which it cannot identify with F; the last sentence of the degree-one case is not in the Lean.
% Polarization: Sol.Pol allows any totally real subfield F0 of F, kept along the tower by UpperTheory.solvable_polarized_restriction in (t3); F0 = Q in (v)(f) of Lemma sha-rigidity-points. The paragraph on base change restricts mu to F^+ to apply Theorem sha-siegel-automorphy-base-change, which the book states for the maximal totally real subfield; the Lean's aut_up applies Z.Bc.solC_alt (and stepC_alt in degree one) with F0 directly. HypothesesH' enters only through aut_up. The down half of aUpDown_alt (Thru, IneAb, Keeps) is not used.
\end{proof}
\leanref{Development/AutomorphicArithmetic/AutomorphicVarieties/ModularityArgument/ThresholdArguments/DenseThresholdCases/EndsRankThreeOfLocal.lean}{UpperTheory.V.Z.Win.Dn.automorphy_up_base_change}
\leanref{Development/AutomorphicArithmetic/AutomorphicVarieties/ModularityArgument/ThresholdArguments/DenseThresholdCases/TAUpDownOf.lean}{UpperTheory.V.Z.Win.aUpDownOf_alt}
\leanref{Development/AutomorphicArithmetic/AutomorphicVarieties/ModularityArgument/AutomorphyLifting/PolarizedLinkage/WeakOfAtt.lean}{UpperTheory.V.Z.Win.weakOfAtt}

\begin{lemma}\label{sha-rigidity-points}
Assume the Galois representations as in Theorem~\ref{sha-residually-reducible-automorphy}.
Let $(\ell,k,\psi,K)$ be an Eisenstein datum with $\psi$ modulo $N$ and deformation data $(\eta,c,\Psi)$, let $g_1\in\operatorname{Gal}(\bar{\mathbf Q}/K)$ with $\Psi(g_1)^2-1\in\mathcal O^\times$, let $B$ be a place above $\ell$, and let $\iota:\bar{\mathbf Q}_\ell\cong\mathbf C$.
Let $R$ be the universal ring of the $\mathfrak a_B$-lifts relative to $g_1$ and $B$ (Lemma~\ref{sha-model-universal}), let $f$ be a point of $R$ and $r=r_f$ its representation (Lemma~\ref{sha-model-points}), and let $\rho$ be the restriction of $r$ to $\operatorname{Gal}(\bar{\mathbf Q}/K)$.
\begin{enumerate}
\item Suppose that the restriction of $r$ to $\ker\eta$ is irreducible, and that $\rho\otimes\chi_{\mathrm{cyc}}^{-1}$ is attached through $\iota$ to a regular algebraic cuspidal automorphic representation $\pi$ of $GL_3(\mathbf A_K)$ with $\pi^\vee\cong\pi^c$, and is windowed. Then $H^1_f(\mathbf Q,\operatorname{ad}r)=0$, where $\operatorname{ad}r$ is $M_3(\bar{\mathbf Q}_\ell)$ with $\operatorname{Gal}(\bar{\mathbf Q}/\mathbf Q)$ acting through $r$ and the adjoint action of $\mathcal G_3$.
\item Suppose that on $\ker\eta$ the entries $(1,2)$, $(2,1)$, $(2,3)$ and $(3,2)$ of $r$ vanish and the corner block formed by rows and columns $1$ and $3$ is irreducible. Let $W:\operatorname{Gal}(\bar{\mathbf Q}/\mathbf Q)\to GL_2(\bar{\mathbf Q}_\ell)$ be continuous, agree with the corner block on $\ker\eta$, and satisfy $\operatorname{tr}W(c)=0$ and $\det W=\eta$; suppose that $W\otimes\chi_{\mathrm{cyc}}^{-k}$ is attached through $\iota$ to a newform of weight $2k+1$ and that the restriction of $W$ to every imaginary quadratic field on which it is irreducible is windowed. Let $R$ be a representation of $\operatorname{Gal}(\bar{\mathbf Q}/\mathbf Q)$ on $M_2(\bar{\mathbf Q}_\ell)$ such that, for a complex conjugation $c'$ and a matrix $P$ with $W(g)^tPW(c'gc'^{-1})=P$ for $g\in\operatorname{Gal}(\bar{\mathbf Q}/K)$, the restriction of $R$ to $\operatorname{Gal}(\bar{\mathbf Q}/K)$ is the adjoint representation of $W$ and $R(c')X=-P^{-1}X^tP$. Then $H^1_f(\mathbf Q,R)=0$.
\end{enumerate}
\end{lemma}
\begin{proof}
As in the proof of Lemma~\ref{sha-model-tangent}(e), $r$ is an ordinary point relative to $g_1$ and $B$ with values in a discrete valuation ring $A\subset\bar{\mathbf Q}_\ell$ finite over $\mathcal O$.
Hypotheses (1) and (2) of Lemma~\ref{sha-rigidity-adjoint} hold with $e=k$ and $M=\ell N|d_K|$. Indeed, $\chi_{\mathrm{cyc}}(c)=-1$ and $\psi(c)=\psi(-1)=(-1)^k$, so $\Psi(c)=1$; $\eta$ is unramified at the primes not dividing $d_K$; and the descent step of the proof of Theorem~\ref{sha-residually-reducible-automorphy} gives the link and the seed datum.

(1) By Lemma~\ref{sha-admissibility-stable}, $\rho$ is a stable automorphic triple of weight $k$ and level $N$. So at every prime $w$ of $K$ above $\ell$ the restriction of $\rho$ to $D_w$ is conjugate to an upper triangular representation whose diagonal characters $\kappa_i$ satisfy $\kappa_i\ne\kappa_j\chi_{\mathrm{cyc}}$ and equal $\chi_{\mathrm{cyc}}^{m_i}\chi_i$ on $I_w$, for $(m_1,m_2,m_3)=(k,0,-k)$ and characters $\chi_i$ trivial on $\operatorname{Gal}(\bar{\mathbf Q}/K(\mu_N))$. This is hypothesis (4) of Lemma~\ref{sha-rigidity-adjoint}.
By the paragraph on irreducibility on open subgroups in the proof of Theorem~\ref{sha-residually-reducible-automorphy}, $r$ is irreducible on every open subgroup of $\ker\eta$; with the hypotheses on $\rho$, this is (3).
So Lemma~\ref{sha-rigidity-adjoint} gives $H^1_f(\mathbf Q,\operatorname{ad}r)=0$.

(2) The determinant of $W$ is $\eta$, which is trivial on $\operatorname{Gal}(\bar{\mathbf Q}/\mathbf Q(\mu_{|d_K|}))$; here $d_K$ is odd, being prime to the even integer $N$. The argument then follows that of Lemma~\ref{sha-rigidity-adjoint}, with the linked points of rank two of the link in place of those of rank three, in the following steps.
\begin{enumerate}
\item[(i)] \emph{Level and irreducibility.} $W$ has level $N|d_K|$ and weights $(k,-k)$, and its restriction $W_K$ to $\operatorname{Gal}(\bar{\mathbf Q}/K)$ is irreducible, even on the subgroup $H$ of the elements with $\chi_{\mathrm{cyc}}=1$.
Write $M_g$ for the matrix part of $r(g)$. By the argument of the proof of Lemma~\ref{sha-model-points}(c), the entry $(2,2)$ of $M_g$ is $1$ for $g\in\ker\eta$. So, for $g\in\ker\eta$, $M_g$ is the matrix $B(g)$ of Lemma~\ref{sha-rigidity-two-by-two}(1) for $W$.
\emph{Away from $\ell$.} Let $I$ be an inertia group at a prime $p\ne\ell$. As in the proof of Lemma~\ref{sha-admissibility-stable}, the $M_g$ with $g\in I\cap\ker\eta$ are simultaneously conjugate to $\operatorname{diag}(\Psi(g),1,\Psi(g)^{-1})$, and $\Psi=\psi$ on $I$, since $\chi_{\mathrm{cyc}}$ is trivial there. By Lemma~\ref{sha-rigidity-two-by-two}(1), $W$ is conjugate on $I\cap\ker\eta$ to $\operatorname{diag}(\psi,\psi^{-1})$, and by Lemma~\ref{sha-rigidity-two-by-two}(2) it is conjugate on $I$ to $\operatorname{diag}(\chi_1,\chi_2)$ for characters $\chi_i$ of $\operatorname{Gal}(\bar{\mathbf Q}/\mathbf Q)$ that are trivial wherever $\psi$ and $\eta$ are. Both $\psi$ and $\eta$ are trivial on $\operatorname{Gal}(\bar{\mathbf Q}/\mathbf Q(\mu_{N|d_K|}))$, so the $\chi_i$ are too.
\emph{Above $\ell$.} Let $A$ be a place above $\ell$, with decomposition group $D_A$ and inertia group $I_A\subset\ker\eta$. As in the proof of Lemma~\ref{sha-admissibility-stable}, there is $Y\in GL_3(\bar{\mathbf Q}_\ell)$ such that the $YM_gY^{-1}$ with $g\in D_A\cap\ker\eta$ are upper triangular, with diagonal $\Psi(g),1,\Psi(g)^{-1}$ for $g\in I_A$. So $v=Y^{-1}e_1\ne0$ satisfies $M_gv=\Psi(g)v$ for $g\in I_A$, and the characteristic polynomial of $M_g$ is $(X-\Psi(g))(X-1)(X-\Psi(g)^{-1})$ there.
The inertia group $I_A$ contains an element $s_0$ with $\Psi(s_0)^2\ne1$, since $\ell-1\nmid2k$ and $\psi_\ell^2=1$. The second coordinate of $M_{s_0}v$ is $v_2$, by the shape of $M_{s_0}$, and it is $\Psi(s_0)v_2$; so $v_2=0$. Hence $(v_1,v_3)\ne0$ is an eigenvector of $W(g)$ with eigenvalue $\Psi(g)$ for $g\in I_A$, where also $\det W(g)=\eta(g)=1$.
By Lemma~\ref{sha-rigidity-two-by-two}(3), with $I_A$ normal in $D_A$, $W$ is conjugate on $D_A$ to an upper triangular representation with diagonal characters $\kappa_1,\kappa_2$ equal to $\Psi=\chi_{\mathrm{cyc}}^k\psi$ and $\Psi^{-1}=\chi_{\mathrm{cyc}}^{-k}\psi^{-1}$ on $I_A$, where $\psi^{\pm1}$ are trivial on $\operatorname{Gal}(\bar{\mathbf Q}/\mathbf Q(\mu_N))$.
Moreover $\kappa_i\ne\kappa_j\chi_{\mathrm{cyc}}$ for all $i,j$. Let $s\in I_A$ with $\chi_{\mathrm{cyc}}(s)$ of infinite order, $(m_1,m_2)=(k,-k)$ and $n=|(\mathbf Z/N)^\times|$. If $\kappa_i=\kappa_j\chi_{\mathrm{cyc}}$, then raising to the power $n$ at $s$ kills the values of $\psi$ and gives $\chi_{\mathrm{cyc}}(s)^{n(m_i-m_j-1)}=1$, while $m_i-m_j-1$ is $-1$, $2k-1$ or $-2k-1$, which is not $0$.
So $W$ has level $N|d_K|$ and weights $(k,-k)$ (Definition~\ref{sha-eigensystem-level}).
\emph{Irreducibility.} The restriction $W_K$ is irreducible, since it is the corner block of $r$ on $\ker\eta=\operatorname{Gal}(\bar{\mathbf Q}/K)$. Suppose that it is reducible on $H$. Since $\operatorname{Gal}(\bar{\mathbf Q}/K)/H$ is abelian, $H$ contains every commutator of $\operatorname{Gal}(\bar{\mathbf Q}/K)$, and Lemma~\ref{sha-siegel-automorphy-clifford}, with $m=2$ and $\kappa=\bar{\mathbf Q}_\ell$, gives a character $\phi\ne1$ of $\operatorname{Gal}(\bar{\mathbf Q}/K)$, trivial on $H$, with $\phi^2=1$ and $W_K\otimes\phi\cong W_K$. Comparing traces, $\operatorname{tr}W(t)=0$ whenever $\phi(t)\ne1$.
Let $A$ be a place above $\ell$. For $t\in I_A$, $\operatorname{tr}W(t)=\operatorname{tr}M_t-1=\Psi(t)+\Psi(t)^{-1}$, by the characteristic polynomial above; so $\Psi(t)^4=1$ if moreover $\phi(t)\ne1$.
Let $g$ satisfy $\phi(g)\ne1$. Since $\chi_{\mathrm{cyc}}(I_A)=\mathbf Z_\ell^\times$, there are $s,u\in I_A$ with $\chi_{\mathrm{cyc}}(s)=\chi_{\mathrm{cyc}}(g)$ and $\chi_{\mathrm{cyc}}(u)=1+\ell$. Since $\phi$ factors through $\chi_{\mathrm{cyc}}$ and $\phi^2=1$, $\phi(s)=\phi(su^2)=\phi(g)\ne1$, so $\Psi(s)^4=\Psi(su^2)^4=1$ and $\Psi(u)^8=1$. But $\Psi(u)^{8\operatorname{ord}\psi}=(1+\ell)^{8k\operatorname{ord}\psi}\ne1$. So $W_K$ is irreducible on $H$.
\item[(ii)] \emph{Automorphy over $K$.} By (i) and the base change from $\mathbf Q$ to $K$ \cite{AC89} of the cuspidal representation of $GL_2(\mathbf A_{\mathbf Q})$ of the newform, twisted by a power of the absolute value, $W_K$ is attached through $\iota$ to a regular algebraic cuspidal automorphic representation $\pi_K$ of $GL_2(\mathbf A_K)$. The pair $(\pi_K,\chi)$ is polarized for some Hecke character $\chi$ of $\mathbf Q$, because $W$ preserves the alternating form on $\bar{\mathbf Q}_\ell^2$ up to $\det W=\eta$, a character of finite order. By hypothesis, $W_K$ is windowed.
\item[(iii)] \emph{Triangular forms.} Let $w$ be a prime of $K$ above $\ell$. By the level of $W$ at $\ell$, there are a place $A$ above $\ell$ and characters $\kappa_1,\kappa_2$ of $D_A$ such that $W$ is triangular on $D_A$ with diagonal $\kappa_1,\kappa_2$, and $\kappa_i=\chi_{\mathrm{cyc}}^{m_i}\chi_i$ on $I_A$ for $(m_1,m_2)=(k,-k)$ and characters $\chi_i$ trivial on $\operatorname{Gal}(\bar{\mathbf Q}/\mathbf Q(\mu_{N|d_K|}))$. Moving $A$ to a place $B=tA$ above $w$ by conjugation with $W(t)$, and restricting to $D_B\cap\operatorname{Gal}(\bar{\mathbf Q}/K)$, gives a triangular form of $W_K$ there, with characters $\kappa'_i=\kappa_i(t^{-1}\cdot t)$ that satisfy $\kappa'_i=\chi_{\mathrm{cyc}}^{m_i}\chi'_i$ on inertia, where $\chi'_i$ is the restriction of $\chi_i$ to $\operatorname{Gal}(\bar{\mathbf Q}/K)$, since the characters $\chi_{\mathrm{cyc}}$ and $\chi_i$ of $\operatorname{Gal}(\bar{\mathbf Q}/\mathbf Q)$ take the same value at $t^{-1}st$ as at $s$; so $\chi'_i$ is trivial on $\operatorname{Gal}(\bar{\mathbf Q}/K(\mu_{N|d_K|}))$. And $\kappa'_i\ne\kappa'_j\chi_{\mathrm{cyc}}$ for all $i,j$: otherwise $\chi_{\mathrm{cyc}}^{m_i-m_j-1}=\chi'_j/\chi'_i$ on the inertia group, a character trivial on $\operatorname{Gal}(\bar{\mathbf Q}/K(\mu_{N|d_K|}))$, which has finite order $d$, while $m_i-m_j-1\ne0$ because $2k\ne\pm1$, and $\chi_{\mathrm{cyc}}^{(m_i-m_j-1)d}$ is not trivial on the inertia group, on which $\chi_{\mathrm{cyc}}$ takes values of infinite order.
\item[(iv)] \emph{A free element.} By (ii) and (iii), Lemma~\ref{sha-rigidity-free}, with $F=K$, $n=2$, $a=0$, $M=N|d_K|$ and $(m_1,m_2)=(k,-k)$, gives at every place above $\ell$ an element of the decomposition group in $\operatorname{Gal}(\bar{\mathbf Q}/K)$ with trivial cyclotomic character at which the eigenvalues $\alpha_1,\alpha_2$ of $W$ satisfy $\alpha_1^m\ne\alpha_2^m$ for $m\ge1$.
\item[(v)] \emph{The vanishing.} Write $G_E=\operatorname{Gal}(\bar{\mathbf Q}/E)$ for a field $E$. Let $g_2$, the link of $\eta$, $\Psi$, $c$, $g_2$ and $B$ for $M=\ell N|d_K|$ to the abelian CM field $L\supset K$ that satisfies the residual conditions, and the seed datum for $L$ be those given by the descent step, as above. By Lemma~\ref{sha-descent-points}, $\eta(g_2)=1$ and $\Psi(g_2)^2-1\in\mathcal O^\times$; and $g_2\in G_{L^+}\cap\ker\eta=G_L$, by (iii) of Definition~\ref{sha-descent-link}.
\begin{enumerate}
\item[(a)] \emph{The base element.} The point $r$ is also an ordinary point relative to $g_2$ and $B$. Of the conditions of Definition~\ref{sha-ordinary-deformation}(c), only the normalization involves the base element; it asks that the entries $(1,2)$ and $(2,1)$ of the difference of the matrix parts of $r(g_2)$ and $r(cg_2^{-1}c)$ vanish. Both $g_2$ and $cg_2^{-1}c$ lie in $\ker\eta$, where these entries of the matrix part of $r$ vanish by hypothesis.
\item[(b)] \emph{The linked point.} By (a) and Lemma~\ref{sha-descent-points}(2), the rows and columns $1$ and $3$ give a linked point $(\rho_2,(\psi_{v,1},\psi_{v,3}))$ of rank two with values in $A$ for the link. Since $W$ agrees with the corner block on $\ker\eta\supseteq G_L$, $\rho_2(p(s))=W(s)$ for $s\in G_L$; and $\psi_{v,1}=\Psi$ and $\psi_{v,3}=\Psi^{-1}$ on $I_v$.
\item[(c)] \emph{Irreducibility over $L$.} The restriction of $W$ to $G_L$ is irreducible. Indeed, $W_K$ is irreducible by (i), and $G_L$ is normal in $G_K$. For $g\in G_K$, the trace of $W(g)$ is the sum of the entries $(1,1)$ and $(3,3)$ of the matrix part of $r(g)$, which reduce modulo $\mathfrak m_A$ to $\bar\Psi(g)$ and $\bar\Psi(g)^{-1}$, because $r$ reduces to $r_0(\eta,\Psi)$; so $\operatorname{tr}W(g)\equiv\Psi(g)+\Psi(g)^{-1}$. By Clifford's theorem \cite[(49.2)]{CR62}, applied to the irreducible representation $W_K$ of dimension $2$ and the normal subgroup $G_L$, either $W$ is irreducible on $G_L$, or $W(G_L)$ consists of scalars, or there are two lines in $\bar{\mathbf Q}_\ell^2$, each stable under $W(G_L)$, and some $t\in G_K$ with $W(t)$ exchanging them.
In the second case $W(g_2)=a$ with $a^2=\det W(g_2)=\eta(g_2)=1$, so $(\operatorname{tr}W(g_2))^2=4$; modulo $\mathfrak m_A$ this gives $(\bar\Psi(g_2)+\bar\Psi(g_2)^{-1})^2=4$, that is $(\bar\Psi(g_2)^2-1)^2=0$, which contradicts $\Psi(g_2)^2-1\in\mathcal O^\times$.
In the third case $W(ts)$ exchanges the two lines for every $s\in G_L$, so $\operatorname{tr}W(ts)=0$ and $\bar\Psi(ts)^2=-1$. For $s=1$ and $s=g_2$ this gives $\bar\Psi(t)^2=-1=\bar\Psi(t)^2\bar\Psi(g_2)^2$, so $\bar\Psi(g_2)^2=1$, the same contradiction.
\item[(d)] \emph{The pairing.} Let $J=\begin{pmatrix}0&1\\-1&0\end{pmatrix}$, so that ${}^tXJX=\det(X)J$ for every $2\times2$ matrix $X$, and let $P_0=JW(c)^{-1}$. For $s\in G_K$, $W(csc^{-1})=W(c)W(s)W(c)^{-1}$ and $\det W(s)=\eta(s)=1$, so ${}^tW(s)P_0W(csc^{-1})={}^tW(s)JW(s)W(c)^{-1}=P_0$.
\item[(e)] \emph{Automorphy over $L$.} The field $L$ is abelian over $\mathbf Q$, hence Galois over $K$ with abelian Galois group. By (ii), $W_K$ is attached through $\iota$ to a regular algebraic cuspidal automorphic representation of $GL_2(\mathbf A_K)$. Theorem~\ref{sha-siegel-automorphy-base-change}, with $F=K$, $F_0=\mathbf Q$, $M=L$, $\mu=1$, the complex conjugation $c$ and the matrix $P_0$ of (d), and with the irreducibility of (c), shows that the restriction $W_L$ of $W$ to $G_L$ is attached through $\iota$ to a regular algebraic cuspidal automorphic representation $\pi_L$ of $GL_2(\mathbf A_L)$.
\item[(f)] \emph{The twist.} Let $\kappa_2$ be the twist of rank two of (c)(1) of the proof of Theorem~\ref{sha-residually-reducible-automorphy}: the restriction to $G_L$ of an algebraic Hecke character $\kappa$ of $K$ of infinity type $(-1,0)$ with $\kappa\kappa^c=\chi_{\mathrm{cyc}}^{-1}$, which is moreover trivial on the inertia groups of $G_L$ at the places not dividing $M$. Every complex conjugation $\tau$ is $wc$ with $w\in G_K$, and $\kappa$ is a character of $G_K$, so $\kappa_2(g)\kappa_2(\tau g\tau^{-1})=\chi_{\mathrm{cyc}}(g)^{-1}$ for $g\in G_L$.
Then $W_L\otimes\kappa_2$ is attached through $\iota$ to a regular algebraic cuspidal automorphic representation $\Pi$ of $GL_2(\mathbf A_L)$ with $\Pi^\vee\cong\Pi^c$, and is windowed. Indeed, $\kappa_2$ is the character attached through $\iota$ to an algebraic Hecke character $\chi_\kappa$ of $L$, so $W_L\otimes\kappa_2$ is attached to $\Pi=\pi_L\otimes(\chi_\kappa\circ\det)$, which is regular algebraic and cuspidal. By (d), ${}^t(\kappa_2W)(g)\,P_0\,(\kappa_2W)(cgc^{-1})=\kappa_2(g)\kappa_2(cgc^{-1})P_0=\chi_{\mathrm{cyc}}(g)^{-1}P_0$ for $g\in G_L$; so the eigenvalues of $W_L\otimes\kappa_2$ at a Frobenius element at $cw$ are $\chi_{\mathrm{cyc}}^{-1}$ times the inverses of those at $w$, $\Pi^c$ and $\Pi^\vee$ have the same Satake parameters at almost all places, and $\Pi^\vee\cong\Pi^c$ by strong multiplicity one (Theorem~\ref{lit-jacquet-shalika}(1)). The representation $W_K$ is triangular on the decomposition group in $G_K$ at every place above $\ell$, as the corner of the triangular forms of $r$. Indeed, let $w$ be a prime of $K$ above $\ell$ and $P$ a place above $w$. By (a), $r$ is an ordinary point relative to $g_2$ and $B$. As in Step~3 of the proof of Lemma~\ref{sha-rigidity-adjoint}, with $P=tB$ for the element $t$ given by (iv) of the link, conjugation by the matrix part of $r(t)$, through $X\mapsto J\,{}^tX^{-1}J$ if $\eta(t)=-1$, carries the triangular form of $r$ on $D_B\cap\ker\eta$ to a matrix $Y\in GL_3(\bar{\mathbf Q}_\ell)$ such that $YM_sY^{-1}$ is upper triangular for $s\in D_P\cap G_K$, with diagonal entries $\Psi(s)$, $1$ and $\Psi(s)^{-1}$ for $s$ in the inertia group $I_P\cap G_K$. Since $\Psi=\chi_{\mathrm{cyc}}^k\psi$ with $k\ge1$ and $\psi$ of finite order, and $\chi_{\mathrm{cyc}}$ takes values of infinite order on $I_P\cap G_K$, some $s_0$ in this inertia group has $\Psi(s_0)\ne1$. On $\ker\eta\supseteq G_K$ the matrices $M_s$ have the shape of Lemma~\ref{sha-rigidity-corner}, by the hypothesis on $r$ and by (i), and their corner block is $W(s)$. So Lemma~\ref{sha-rigidity-corner}, with $H=D_P\cap G_K$, $g=Y$ and $s_2=s_0$, shows that $W$ is triangular on $D_P\cap G_K$. Conjugation by $W(g)$, $g\in G_K$, carries this triangular form to the decomposition groups at the other places $gP$ above $w$.
The restriction $W_L$ is windowed, by Lemma~\ref{sha-rigidity-windowed-up} with $F=K$, $E=L$, $n=2$, $r=W_K$, $F_0=\mathbf Q$, $\mu=1$, the complex conjugation $c$ and the matrix $P_0$ of (d): $L$ is Galois over $K$ with abelian Galois group, $W_L$ is irreducible by (c), $W_K$ is attached through $\iota$ to a regular algebraic cuspidal automorphic representation and is windowed by (ii), and it is triangular at the places above $\ell$, as shown above. The windowed condition passes from $W_L$ to its twist by the algebraic character $\kappa_2$: in ($\alpha$) the diagonal characters are multiplied by $\kappa_2$, which leaves the conditions unchanged, and in ($\beta$), as $L$ is a CM field, there is $b\in\mathbf R$ with $|\chi_{\kappa,u}(\varpi_u)|=q_u^b$ at every finite place $u$ of $L$, so the bound for the representation attached to $W_L$ with $s$ gives the bound for $\Pi$ with $s+b$.
\item[(g)] \emph{The vanishing over $L^+$.} By (c)(2) of the proof of Theorem~\ref{sha-residually-reducible-automorphy}, after enlarging $\mathcal O$ to a ring containing the values of $\rho_2$, there is a realization of rank two of the link, twisted by $\kappa_2$: the data of Situation~\ref{sha-hecke-situation} over $L$, with a finite set $S$ of places of $L^+$ containing those above $\ell$, groups $\Gamma\supseteq\Gamma_0\supseteq\Gamma'$, the projection $\mathrm{pr}:G_L\to\Gamma_0$, the image $c''$ of a complex conjugation $\tau\in G_{L^+}$, and a place $v_0$ not dividing $M$; a residual datum of rank two, which satisfies the conditions at $v_0$ of Theorem~\ref{sha-siegel-automorphy-adjoint}; its universal ordinary determinant ring $B_2$; and a map from $B_2$ to the universal ring of linked points of rank two, carrying the universal pair to the twist by $\kappa_2$ of the determinant and the characters of the universal linked point.
The linked point $\rho_2$ of (b) therefore gives a point $x$ of $B_2$. Let $\rho:\Gamma_0\to GL_2(\bar{\mathbf Q}_\ell)$ be the representation with $\rho(\mathrm{pr}(g))=\kappa_2(g)W(g)$ for $g\in G_L$; then $x$ carries $D^{\mathrm u}$ to the determinant of $\rho$, and $\psi_{x,v,1}=\kappa_2\psi_{v,1}$ and $\psi_{x,v,2}=\kappa_2\psi_{v,3}$ on $D_v$.
The point $x$ and $\rho$ satisfy hypotheses (1) to (6) of Theorem~\ref{sha-siegel-automorphy-adjoint}.
For (1), $x$ carries $D^{\mathrm u}$ to the determinant of $\rho$. On $I_v$, $\Psi^{\pm1}=\chi_{\mathrm{cyc}}^{\pm k}\psi^{\pm1}$, so $\psi_{x,v,i}=\chi_{\mathrm{cyc}}^{ke_{2,i}}\psi^{e_{2,i}}\kappa_2$ there. Since $\kappa_2$ is algebraic, local class field theory gives integers $a_{v,\tau}$ such that, for every $m\in\mathbf Z$ and every character $\chi$ of $G_L$ with open kernel, a character of $D_v$ equal to $\chi_{\mathrm{cyc}}^m\chi\kappa_2$ on $I_v$ satisfies the condition of Definition~\ref{sha-hecke-arcs}(1) with the integers $a_{v,\tau}-m$. So $x$ is arithmetic of the weight $h_{v,i,\tau}=a_{v,\tau}-ke_{2,i}$, which is increasing.
For (2), this is (f).
For (3), let $g\in G_L$ lie in an inertia group at a place above $v_0$. Since $v_0$ does not divide $M$, $\kappa_2(g)=1$, and $\eta$ and $\bar\Psi$ are trivial on that inertia group, so $p(g)=1$ by (i) of Definition~\ref{sha-descent-link}, and $W(g)=\rho_2(p(g))=1$. So $\rho$ is trivial on $I_{v_0}$. For $v\mid\ell$, condition (c) of a linked point gives $(\rho_2(\mathrm{loc}_vh_1)-\psi_{v,1}(h_1))(\rho_2(\mathrm{loc}_vh_2)-\psi_{v,3}(h_2))=0$ for $h_1,h_2\in G_v$. By (7) of Definition~\ref{sha-descent-residual} there is $h_0\in G_v$ with $\bar\Psi(h_0)^2\ne1$. Then $\psi_{v,1}(h_0)-\psi_{v,3}(h_0)$ is a unit, as it reduces to $\bar\Psi(h_0)-\bar\Psi(h_0)^{-1}$, and $\rho_2(\mathrm{loc}_vh_0)$ reduces to the diagonal matrix with entries $\bar\Psi(h_0)$ and $\bar\Psi(h_0)^{-1}$, so the reduction of its characteristic polynomial is $(X-\bar\Psi(h_0))(X-\bar\Psi(h_0)^{-1})$. By Lemma~\ref{sha-rigidity-lagrange}, with $C=A$, $H=G_v$, $r=\rho_2\circ\mathrm{loc}_v$, the characters $\psi_{v,1},\psi_{v,3}$ and $g_0=h_0$, the representation $\rho_2\circ\mathrm{loc}_v$ is conjugate over $A$ to an upper triangular representation with diagonal $\psi_{v,1},\psi_{v,3}$. Multiplying by $\kappa_2$, the restriction of $\rho$ to $D_v$ is conjugate to an upper triangular representation with diagonal $\psi_{x,v,1},\psi_{x,v,2}$.
For (4), let $v\mid\ell$. By (iv) there is $g$ in the decomposition group of $G_K$ at the place at which $D_v$ is the decomposition group of $G_L$, with $\chi_{\mathrm{cyc}}(g)=1$ and eigenvalues $\alpha_1,\alpha_2$ of $W(g)$ with $\alpha_1^m\ne\alpha_2^m$ for $m\ge1$. Some power $g^d$ with $d\ge1$ lies in the subgroup $G_L$ of finite index, so in $D_v$; $\varepsilon(g^d)=1$; and the eigenvalues $\alpha_i^d$ of $W(g^d)$ still have this property. By the triangular form of $\rho_2\circ\mathrm{loc}_v$ found for (3), the eigenvalues of $\rho_2(p(g^d))=W(g^d)$ are $\psi_{v,1}(g^d)$ and $\psi_{v,3}(g^d)$; multiplying them by $\kappa_2(g^d)$ does not change their ratio, so $\psi_{x,v,1}(g^d)^m\ne\psi_{x,v,2}(g^d)^m$ for $m\ge1$.
For (5), $\rho$ is irreducible: $\mathrm{pr}$ is surjective, $\rho(\mathrm{pr}(g))=\kappa_2(g)W(g)$ for $g\in G_L$, and $W$ is irreducible on $G_L$ by (c). It is continuous, as $W$ and $\kappa_2$ are.
Let $A_x$ be the quotient of $B_2$ by the kernel of $x$, a local ring with the residue field $k$ of $B_2$; let $j:A_x\to\bar{\mathbf Q}_\ell$ be the injective map induced by $x$, and $\phi:A_x\to k$ the residue map. Let $\bar\chi_1,\bar\chi_2$ be the characters of the residual datum, and $\bar\psi_{v,1},\bar\psi_{v,2}$ its characters at $v\mid\ell$. Since $x$ carries $D^{\mathrm u}$ to the determinant of $\rho$, and $D^{\mathrm u}$ reduces to the determinant of $\bar\chi_1\oplus\bar\chi_2$, the characteristic polynomial of $\rho(y)$, for $y\in\Gamma_0$, is the image under $j$ of that of $D^{\mathrm u}$ at $y$ taken modulo $\ker x$, whose image under $\phi$ is $(X-\bar\chi_1(y))(X-\bar\chi_2(y))$. In the same way, for $v\mid\ell$ and $g\in D_v$, the class $a$ of $\psi^{\mathrm u}_{v,i}(g)$ has $j(a)=\psi_{x,v,i}(g)$ and $\phi(a)=\bar\psi_{v,i}(g)$.
Fix $v\mid\ell$. By (3), the characteristic polynomial of $\rho$ at $g\in D_v$ is $(X-\psi_{x,v,1}(g))(X-\psi_{x,v,2}(g))$. By (4), some $g\in D_v$ with $\varepsilon(g)=1$, so with image in $\Gamma'$, has $\psi_{x,v,1}(g)^m\ne\psi_{x,v,2}(g)^m$ for $m\ge1$. The residual datum is obtained from the residual representation of the link by the twist, so $\bar\psi_{v,1}$ and $\bar\psi_{v,2}$ are the reductions of $\kappa_2$ times $\bar\Psi$ and $\bar\Psi^{-1}$ through $\mathrm{loc}_v$; at an element of $D_v$ that maps to the element $h_0$ of (3), their ratio is $\bar\Psi(h_0)^2\ne1$.
So Lemma~\ref{sha-rigidity-small} applies, with the compact group $\Gamma_0$, its normal subgroup $\Gamma'$, $n=2$, $r=\rho$, $D=D_v$ with its map to $\Gamma_0$, and the characters above: $\rho$ is irreducible on $U\cap\Gamma'$ for every open subgroup $U$ of $\Gamma_0$, and in particular on $\Gamma'$. With the image of $g$ in $\Gamma'$, Lemma~\ref{sha-rigidity-separable}, for the profinite group $\Gamma_0$, gives the second half of (5).
For (6), let $m_{v,i}=ke_{2,i}$ and $k_0=|(\mathbf Z/N)^\times|$. On $I_v$, $\psi_{x,v,i}\psi_{x,v,j}^{-1}\varepsilon^{m_{v,j}-m_{v,i}}=\psi^{e_{2,i}-e_{2,j}}$, whose $k_0$-th power is $1$.
Theorem~\ref{sha-siegel-automorphy-adjoint} then shows: for every $P'\in GL_2(\bar{\mathbf Q}_\ell)$ with ${}^t\rho(y)P'\rho(c''yc''^{-1})=\varepsilon(y)^{-1}P'$ for $y\in\Gamma_0$, and every representation $R'$ of $\Gamma$ on $M_2(\bar{\mathbf Q}_\ell)$ with $R'(y)X=\rho(y)X\rho(y)^{-1}$ for $y\in\Gamma_0$ and $R'(c'')X=-P'^{-1}\,{}^tXP'$, every continuous cocycle of $\Gamma$ with values in $R'$ whose inflation to $G_{L^+}$ satisfies the Bloch-Kato conditions is a coboundary.
\item[(h)] \emph{Descent.} Lemma~\ref{sha-siegel-automorphy-descent} applies with $n=2$, the restriction of $W$ to $G_K$, $\mu=1$, the representation $R$, the complex conjugation $c'$ and the matrix $P$ of the statement, the set $S'=S$, the complex conjugation $\tau$, and $\rho$ and $\kappa_2$ of (g): $\rho(\mathrm{pr}(g))=\kappa_2(g)W(g)$ and $\kappa_2(g)\kappa_2(\tau g\tau^{-1})=\chi_{\mathrm{cyc}}(g)^{1-2}$ for $g\in G_L$, by (f); and its hypothesis on the cocycles of $\Gamma$ is the conclusion of (g). So every continuous cocycle of $\operatorname{Gal}(\bar{\mathbf Q}/\mathbf Q)$ with values in $R$ that satisfies the Bloch-Kato conditions is a coboundary, and $H^1_f(\mathbf Q,R)=0$.
\end{enumerate}
\end{enumerate}
% Step (v) cannot be replaced by [NT23, Theorem A] for the twist of W over L as in Step 5 of Lemma sha-rigidity-adjoint: that needs the image of Gal(Qbar/L(mu_l^infty)) to be enormous, which would follow from the irreducibility of W on every open subgroup of Gal(Qbar/L(mu_l^infty)); the hypotheses of (2) and step (i) give irreducibility on Gal(Qbar/K(mu_l^infty)) and (in the Lean, irr2) on Gal(Qbar/L), not on its open subgroups. The text follows the Lean's route through Theorem sha-siegel-automorphy-adjoint instead.
% Hypothesis (5) in (v)(g) is Th.nice_zero_of_rank_two_or_three_irreducible_free, from rp_irr (irreducibility of rho on Gamma_0, from (c)), the free elements (rp_charpoly, rp_free) and R.ok.f3 with rp_w_ne (two residual characters differ at v): TypeI and the irreducibility on (open U) cap Gamma' are Th.Det.irrSmall (Lemma sha-rigidity-small: Pr.irrNormal, Lemma sha-rigidity-normal, through Pr.cliffPrime, Pr.charPerm and reduce_roots, applied to the normal core of U meet Gamma'), with Th.Det.redB (redB_abs: A_x = B/ker x) for the reduction; N3 is Th.Det.n3small (Lemma sha-rigidity-separable: cosetI, coset_in, strong_of, charpoly_pow, and Pr.regCoset, Lemma sha-rigidity-coset, with fullEx, expSums, minors_ne_zero and sepOfMinors).
% In (f), the windowed condition over L is Z.Win.Dn.automorphy_up_base_change (hWL, in hcsd), which is Lemma sha-rigidity-windowed-up, with F0 = Q, mu = 1, c = sigma, P = J W(sigma)^-1 (pair2, hpair), halg, haut, irreducibility hirr and the triangular forms hFl. The same Lean theorem gives (E) of Proposition sha-siegel-points.
% (1) is Z.J.rigIUE_alt: Pr.ordPt_of_gd (the ordinary point), LinkHyp with M = l N |d_K| (hram), Fr.gd_ex (the link), admissible_typeI_even_level_alt (Lemma sha-admissibility-stable), free_at_of_unlinked_of_automorphic_cyclotomic_twist and Pr.matIrr_open_gd, then Z.Win.Dn.rigEnd3_alt (Lemma sha-rigidity-adjoint).
% (2) is Z.Win.RR.rigII_a_alt: hwt (det W through mem_galOf_quadratic_of_fixes_discr_roots_of_unity); (ii) is witAQ_of_isModularA (windowed), mod0'_alt, newformAut0_alt (Sd.up_quad_of_printed_results) and autKQ0_alt (polarization_k_of_printed_results, with hadm); (iii) is unlRes, second half (Lev.LevelF.of_level, flag_move, Lev.unlGap with Ws.Pt.cycIne); (iv) is freePUA_alt (freeAt_of_autPA_of_unl_of_finrank_eq_two, Lemma sha-rigidity-free with a = 0).
% (i) is U3.typeII_ord (the level: typeII_facts for the entry (2,2), block_form, pt_inertia_away with diagonal_on_subgroup_of_conjugate and exists_isDiagOn_of_isDiagOn_inf_ker_sign away from l; pt_inertia_above, the eigenvector v with v 1 = 0, exists_isFlagOn_of_eigenvector_normal_subgroup and exists_two above l) and U3.typeII_irr_cyc (matIrr_ker, with exists_selfTwist, cycZ_Ine_surjective and exists_unit_pow_ne).
% (v) is Z.Win.Dn.rigid_end_two_w_of_printed_adjoint_lifting: (a) is hr' (the block form gives the normalization at g2); (b) is Pr.exists_nrep2; (c) is irr2 (Pr.cliffPrime, the scalar and monomial cases); (d) is pair2 and hpair; (e) is Z.Bc.solC_of_solDescent with halg; (f) is hka (Pr.IsKa from Pr.kaOf and Sd.twist_of_printed_results, in linked_rigid_of_printed_adjoint_lifting) and hcsd (GalN.V.csdOf0_alt with Af.twist_invariance_automorphic, Z.Win.Dn.twist'_alt, Z.Win.Dn.automorphy_up_base_change with hFl); the triangular forms hFl are Keep.ordFlag2 (Keep.ordFlag3 with the link's pl_all and pl_conj, OrdinaryFlagThree.flag_core for theta, OrdFlag2.ne_one, then OrdFlag2.flag_block and tri_of_lines, Lemma sha-rigidity-corner) moved to every place above w by Dn.flagAll_of. (g) is Th.linked_rigid_of_printed_adjoint_lifting for n = 2 (Pr.coefEx, Pr.cRealEx), that is Th.rigPtOf_alt: goodPtOf (the point x), hpr, ap_locAlg with rp_psi and hin ((1)), haut ((2)), rp_unr with hka.unr and R.v0M, rp_flag (Fr.flag_loc, Lemma sha-rigidity-lagrange) and rp_tw_flag ((3)), freeRes, rp_charpoly and rp_free ((4)), nice_zero_of_rank_two_or_three_irreducible_free ((5)), rp_par ((6)), then Th.AdjointLiftingTheorem (Theorem sha-siegel-automorphy-adjoint); the relation for kappa_2 with tau is rp_conj. (h) is Th.adDown (Lemma sha-siegel-automorphy-descent), with a1, a2 and a3.
\end{proof}
\leanref{Development/AutomorphicArithmetic/AutomorphicVarieties/ModularityArgument/ThresholdArguments/DenseThresholdCases/RigEnd3.lean}{UpperTheory.V.Th.RigidityFirstCase}
\leanref{Development/AutomorphicArithmetic/AutomorphicVarieties/ModularityArgument/ThresholdArguments/DenseThresholdCases/RigEnd3.lean}{UpperTheory.V.Th.RigiditySecondCase}
\leanref{Development/AutomorphicArithmetic/AutomorphicVarieties/TheoremsMuLinked/FreeFlag.lean}{UpperTheory.V.Th.AdZero2}
\leanref{Development/AutomorphicArithmetic/AutomorphicVarieties/ModularityArgument/ThresholdArguments/DenseThresholdCases/ERigIU.lean}{UpperTheory.V.Z.J.rigIUE_alt}
\leanref{Development/AutomorphicArithmetic/AutomorphicVarieties/ModularityArgument/ThresholdArguments/DenseThresholdCases/RRRigII.lean}{UpperTheory.V.Z.Win.RR.rigII_a_alt}

\begin{lemma}\label{sha-admissibility-stable}
Assume the Galois representations as in Theorem~\ref{sha-residually-reducible-automorphy}.
Let $(\ell,k,\psi,K)$ be an Eisenstein datum with $\psi$ modulo $N$ and deformation data $(\eta,c,\Psi)$, let $g_1\in\operatorname{Gal}(\bar{\mathbf Q}/K)$ with $\Psi(g_1)^2-1\in\mathcal O^\times$, let $B$ be a place above $\ell$, and let $\iota:\bar{\mathbf Q}_\ell\cong\mathbf C$.
Let $R$ be the universal ring of the $\mathfrak a_B$-lifts relative to $g_1$ and $B$ (Lemma~\ref{sha-model-universal}), and let $f$ be a point of $R$, with representation $r=r_f$ (Lemma~\ref{sha-model-points}). Suppose that the restriction of $r$ to $\ker\eta$ is irreducible, and that its restriction $\rho$ to $\operatorname{Gal}(\bar{\mathbf Q}/K)$ has $\rho\otimes\chi_{\mathrm{cyc}}^{-1}$ attached through $\iota$ to a regular algebraic cuspidal automorphic representation $\pi$ of $GL_3(\mathbf A_K)$ with $\pi^\vee\cong\pi^c$, and windowed.
Then $\rho$ is a stable automorphic triple of weight $k$ and level $N$ (Definition~\ref{sha-eigensystem-triple}).
\end{lemma}
\begin{proof}
Write $M_g$ for the matrix part of $r(g)$, and $(w_1,w_2,w_3)=(1,0,-1)$, $(m_1,m_2,m_3)=(k,0,-k)$.
The representation $\rho$ is irreducible, since $\ker\eta=\operatorname{Gal}(\bar{\mathbf Q}/K)$; it is continuous, since $r^{\mathrm u}$ is continuous and $f$ is a point; and $\rho\otimes\chi_{\mathrm{cyc}}^{-1}$ is automorphic and windowed by hypothesis. As in the proof of Lemma~\ref{sha-model-tangent}(e), $r$ is an ordinary point relative to $g_1$ and $B$ with values in a discrete valuation ring $A\subset\bar{\mathbf Q}_\ell$ finite over $\mathcal O$. It remains to check the conditions on the level and on the characters at $\ell$ in Definition~\ref{sha-eigensystem-triple}(1).

\emph{Away from $\ell$.} Let $w$ be a prime of $K$ not above $\ell$, and $I_w\subset\operatorname{Gal}(\bar{\mathbf Q}/K)$ an inertia group at $w$. Since $\chi_{\mathrm{cyc}}$ is trivial on $I_w$, the elements $s\in I_w$ with $\psi(s)=1$ have $\eta(s)=\Psi(s)=1$, so they lie in $\mathcal N$, where $r$ and $r_0$ are trivial. They form a normal subgroup of $I_w$ whose quotient is finite of order prime to $\ell$, because $\ell\nmid\operatorname{ord}\psi$. So the matrix parts of $r$ and of $r_0$ on $I_w$ factor through a finite group of order prime to $\ell$, and they are conjugate, since both have integral models with the same reduction, as in the second of the three facts in the proof of Lemma~\ref{sha-model-points}. On $I_w$ the matrix part of $r_0$ is $\operatorname{diag}(\psi,1,\psi^{-1})$, and the characters $\psi^{w_i}$ of $\operatorname{Gal}(\bar{\mathbf Q}/K)$ are trivial on $\operatorname{Gal}(\bar{\mathbf Q}/K(\mu_N))$.

\emph{The flag at $\ell$.} Let $w$ be a prime of $K$ above $\ell$. One of the places $B$ and $cB$ lies above $w$; call it $B_w$, and let $D_w$ and $I_w$ be its decomposition and inertia groups in $\operatorname{Gal}(\bar{\mathbf Q}/K)$.
By Lemma~\ref{sha-model-local}(c), applied to the homomorphism $R_B\to\bar{\mathbf Q}_\ell$ through which $r$ factors on $D_B$, which kills $\mathfrak a_B$, and since $\bar{\mathbf Q}_\ell$ is algebraically closed, there is $P\in GL_3(\bar{\mathbf Q}_\ell)$ such that the $P^{-1}M_gP$ with $g\in D_B\cap\ker\eta$ are upper triangular, with diagonal $\Psi(g),1,\Psi(g)^{-1}$ for $g\in I_B$.
Let $x\in\operatorname{Gal}(\bar{\mathbf Q}/\mathbf Q)$ with $B_w=xB$, so that the decomposition and inertia groups at $B_w$ are $xD_Bx^{-1}$ and $xI_Bx^{-1}$. For $g$ in $xD_Bx^{-1}\cap\ker\eta$, $M_g$ is $M_xM_{x^{-1}gx}M_x^{-1}$ if the sign of $r(x)$ is $1$, and $M_x\theta(M_{x^{-1}gx})M_x^{-1}$ with $\theta(X)=J\,{}^tX^{-1}J$ if it is $-1$. The map $\theta$ preserves upper triangular matrices and sends one with diagonal $(a,b,d)$ to one with diagonal $(d^{-1},b^{-1},a^{-1})$, and $\Psi(x^{-1}gx)=\Psi(g)$.
So in both cases the $M_g$ with $g\in D_w$ are simultaneously conjugate to upper triangular matrices whose diagonal entries are characters $\kappa_1,\kappa_2,\kappa_3$ of $D_w$, equal to $\Psi,1,\Psi^{-1}$ on $I_w$. On $I_w$, $\Psi=\chi_{\mathrm{cyc}}^k\psi$, so $\kappa_i=\chi_{\mathrm{cyc}}^{m_i}\psi^{w_i}$ there.

\emph{Potential crystallinity.} The image of the homomorphism $R_B\to\bar{\mathbf Q}_\ell$ above lies in a finite extension of $\mathbf Q_\ell$, as in the proof of Lemma~\ref{sha-model-tangent}(e), so by Lemma~\ref{sha-model-local}(e) the matrix parts of $r$ on the open subgroup $H_B$ of $D_B$ form a crystalline representation, and $\rho$ is potentially crystalline at $B$. On $cH_Bc$ the matrix parts of $r$ are $g\mapsto\theta(M_{cgc})$, because the matrix part of $r(c)$ is $1$ and its sign is $-1$; this is the contragredient of a crystalline representation conjugated by $J$, hence crystalline, and $\rho$ is potentially crystalline at $cB$. So $\rho$ has level $N$ and weights $(k,0,-k)$.

\emph{The condition $\kappa_i\ne\kappa_j\chi_{\mathrm{cyc}}$.} Suppose that $\kappa_i=\kappa_j\chi_{\mathrm{cyc}}$ on $D_w$ for some $i,j$; this leads to a contradiction.
Let $L$, $g_2$, the link and the seed datum be those of the descent step of the proof of Theorem~\ref{sha-residually-reducible-automorphy}, with $M=\ell N|d_K|$; $\eta$ is unramified at the primes not dividing $M$, $\Psi(c)=1$, and $N$ is even.
By the paragraph on irreducibility on open subgroups in that proof, $r$ is irreducible on every open subgroup of $\ker\eta$.
By Step 2 of the proof of Lemma~\ref{sha-rigidity-adjoint}, $r'=hrh^{-1}$ is an ordinary point relative to $g_2$ and $B$ for some $h\in GL_3(A)$ congruent to $1$, and $\rho'=h\rho h^{-1}$ has the flag at $B_w$, with the same characters, and is potentially crystalline there.
By the descent step, $r'$ is a linked point of rank three, with characters $\psi_{v,1},\psi_{v,2},\psi_{v,3}$ at the places $v$ of the link above $\ell$ equal to $\Psi,1,\Psi^{-1}$ on inertia, and with the ordinary condition $\prod_{a=1}^3\big(r'(g_a)-\psi_{v,a}(g_a)\big)=0$ for $g_1,g_2,g_3$ in the decomposition group $D_v$ of the link at $v$; it is irreducible on every open subgroup of $\operatorname{Gal}(\bar{\mathbf Q}/L)$.
The place $B_w$ is $uB_v$ for the place $B_v$ of some $v$ and some $u\in\operatorname{Gal}(\bar{\mathbf Q}/L^+)$. Conjugation by $r'(u)$, followed by $\theta$ and the reversal of the indices when the sign of $r'(u)$ is $-1$, carries the flag of $\rho'$ at $B_w$ to a flag at $B_v$ whose characters $\kappa'_1,\kappa'_2,\kappa'_3$ satisfy $\kappa'_{i'}=\kappa'_{j'}\chi_{\mathrm{cyc}}$ for some $i',j'$, and carries the potential crystallinity to $B_v$. The group $D_v$ maps into the decomposition group of $\operatorname{Gal}(\bar{\mathbf Q}/L)$ at $B_v$.
Each $\kappa'_b$ agrees on $D_v$ with some $\psi_{v,a}$. Otherwise choose, for each $a$, an element $g_a\in D_v$ with $\kappa'_b(g_a)\ne\psi_{v,a}(g_a)$; in a basis in which the $r'(g_a)$ are upper triangular, the entry $(b,b)$ of $\prod_a(r'(g_a)-\psi_{v,a}(g_a))$ is $\prod_a(\kappa'_b(g_a)-\psi_{v,a}(g_a))\ne0$, which contradicts the ordinary condition.
So $\psi_{v,a}=\psi_{v,a'}\chi_{\mathrm{cyc}}$ on $D_v$ for some $a,a'$. This contradicts proposition (d) of the proof of Theorem~\ref{sha-residually-reducible-automorphy}, since $r'$ is potentially crystalline at $B_v$ over $L$.
Hence $\kappa_i\ne\kappa_j\chi_{\mathrm{cyc}}$ for all $i,j$, and $\rho$ is a stable automorphic triple of weight $k$ and level $N$.
% admissible_typeI_even_level_alt: toAdm3 is U3.admissible_rank_three_of_type_one_of_automorphic_point (irr from TypeI; continuous_rho_of_point; level away by pt_inertia_away through conj_of_close; level at l by exists_place_over_conjugate_embedding, pt_inertia_above (Loc.lemma2b, move, sgn_up), exists_flag_characters_of_upper_triangular, and potentially_crystalline_conjugate_place through Loc.ptFin, Loc.rtPoints and crystalline_of_theta_twist_gl3). unl is UnlEndLogNorm, from Z.J.LinkedUnlE (proposition (d)) by Z.J.unlEnd3E: Fr.gd_ex (the link, with LinkHyp for M = l N |d_K|), Pr.matIrr_open_gd, ordPt_renorm, Pr.exists_nrep3, move_all (flag_pos, flag_neg, potentiallyCrystalline_comp_conjugation and potentiallyCrystalline_contragredient), k.pl_all, k.split, flag_match and Br.potRes.
\end{proof}
\leanref{Development/AutomorphicArithmetic/AutomorphicVarieties/ModularityArgument/ThresholdArguments/DenseThresholdCases/NumberU3US.lean}{UpperTheory.V.Z.J.AdmI_s}
\leanref{Development/AutomorphicArithmetic/AutomorphicVarieties/ModularityArgument/ThresholdArguments/DenseThresholdCases/AdmissibleTypeIEvenLevel.lean}{UpperTheory.V.Z.J.admissible_typeI_even_level_alt}

\begin{proposition}\label{sha-siegel-count}
For an odd prime $\ell$, an isomorphism $\iota:\bar{\mathbf Q}_\ell\cong\mathbf C$ and integers $M,k\ge1$, say that a continuous representation $r:\operatorname{Gal}(\bar{\mathbf Q}/\mathbf Q)\to GL_4(\bar{\mathbf Q}_\ell)$ is \emph{of Siegel type} with level $M$ and weight $k$ if: it is irreducible and automorphic through $\iota$ (Definition~\ref{sha-siegel-datum}); it has level $M$ and weights $(k+1,k,1,0)$ and is $\ell$-distinct (Definition~\ref{sha-eigensystem-level}); $\ell$ divides $M$; and it preserves a nondegenerate bilinear form up to a character that is $\chi_{\mathrm{cyc}}^{k+1}$ times a character of finite order.
Assume the Galois representations as in Theorem~\ref{sha-residually-reducible-automorphy}.
There are $c$ and $A$ with the following property.
Let $k\ge2$ and $M\ge1$, let $p\ge5$ be a prime not dividing $M$, let $n\ge1$, let $z_1,z_2,z_3\in\mathbf C$ be $n$-th roots of unity, and put $t=p^{k+1}z_1+p^kz_2+pz_3+1$.
There is a nonzero integer $N_0$ with $\log|N_0|\le c\,\varphi(n)\,(M(k+1))^{A}\log p$ such that for every odd prime $\ell\ne p$, every $\iota:\bar{\mathbf Q}_\ell\cong\mathbf C$, every Frobenius element $\sigma$ at a place above $p$, and every finite set $X$ of pairwise non-isomorphic representations of Siegel type with level $M$ and weight $k$, the traces $\operatorname{tr}r(\sigma)$ with $r\in X$ differ from $\iota^{-1}(t)$, and
\[
 |N_0|_\ell\le\prod_{r\in X}\big|\operatorname{tr}r(\sigma)-\iota^{-1}(t)\big|_\ell.
\]
\end{proposition}
\begin{proof}
Apply Lemma~\ref{sha-eigensystem-rational} with $n=4$, $u=k+1$, $m=(k+1,k,1,0)$, which satisfies $m_i+m_{5-i}=k+1$, and $H=k+1$.
A representation $r$ of Siegel type with level $M$ and weight $k$ lies in the set that this lemma classifies.
It is irreducible, has level $M$ and weights $m$, and preserves a nondegenerate bilinear form up to $\chi_{\mathrm{cyc}}^{k+1}$ times a character of finite order.
Being automorphic through $\iota$, its restriction to every imaginary quadratic field on which it is irreducible is attached through $\iota$ to a regular algebraic cuspidal automorphic representation and is windowed.
It satisfies the distinctness condition at $\ell$ for $M$ and $m$ (Definition~\ref{sha-eigensystem-distinctness}), since $\ell$ divides $M$ and $r$ is $\ell$-distinct.
% number_sp4_of_admissible_to_admissible_n applies countQ_alt with n = 4; a representation of Siegel type (Af.AdmRAlt) lies in AdmRNAlt by autQN'_alt, which is the identity, and its condition A4.SafeRA holds by the first alternative.
The lemma gives constants $c_1,A_1$ and an eigenvalue datum $(\beta_s)_{s\in\Sigma}$ of rank $4$ at $p$ with weights $m$ and weight $k+1$ (Definition~\ref{sha-eigensystem-data}), with $|\Sigma|\le c_1(M(k+2))^{A_1}$, that classifies at $p$, for every odd prime $\ell\ne p$ and every $\iota$, the representations of Siegel type with level $M$ and weight $k$, by a map $r\mapsto s(r)$.
Since all $m_i\ge0$, the elements of each $\beta_s$ are algebraic integers; they have absolute values between $p^{k/2}$ and $p^{(k+2)/2}$, and their product has absolute value $p^{2(k+1)}$.

For $s\in\Sigma$ and $x\in(\mathbf Z/n)^\times$ put
\[
 f(s,x)=\sum_{b\in\beta_s}b-\big(p^{k+1}z_1^{x}+p^kz_2^{x}+pz_3^{x}+1\big).
\]
Each $f(s,x)$ is an algebraic integer.
An automorphism of $\mathbf C$ permutes the multisets $\beta_s$ and raises every $n$-th root of unity to the same power $y\in(\mathbf Z/n)^\times$, so it permutes the numbers $f(s,x)$; hence $N_0=\prod_{s,x}f(s,x)$ is an integer.
One element of $\beta_s$ has absolute value at most $p^{(k+1)/2}$, since the product has absolute value $p^{2(k+1)}$, and the other three are less than $p^{(k+2)/2}=\sqrt p\,p^{(k+1)/2}$; as $p\ge5$ and $k\ge2$, $2p^{(k+1)/2}\le p^k$ and $\sqrt p\,p^{(k+1)/2}\le p^k$, so $\big|\sum_{b\in\beta_s}b\big|\le\tfrac72p^k$.
Since $\tfrac72p^k+p^k+p+1<p^{k+1}=|p^{k+1}z_1^x|$, each $f(s,x)$ is nonzero, and $|f(s,x)|\le p^{2(k+1)}$.
So $N_0\ne0$, and $\log|N_0|\le|\Sigma|\,\varphi(n)\cdot2(k+1)\log p\le c\,\varphi(n)(M(k+1))^{A}\log p$ with $c=2^{A_1+1}c_1$ and $A=A_1+1$.

Let $\ell$, $\iota$, $\sigma$ and $X$ be as in the statement.
For $r\in X$, the eigenvalues of $\iota r(\sigma)$ are the elements of $\beta_{s(r)}$, so $\operatorname{tr}r(\sigma)-\iota^{-1}(t)=\iota^{-1}f(s(r),1)$, which is nonzero.
The map $r\mapsto s(r)$ is injective on $X$, since the elements of $X$ are pairwise non-isomorphic.
Every $f(s,x)$ is an algebraic integer, so $|\iota^{-1}f(s,x)|_\ell\le1$; omitting from $N_0$ the factors $f(s,x)$ that are not of the form $f(s(r),1)$ with $r\in X$ gives $|N_0|_\ell\le\prod_{r\in X}|\iota^{-1}f(s(r),1)|_\ell$.
\end{proof}
\leanref{Development/AutomorphicArithmetic/AutomorphicVarieties/ModularityArgument/ThresholdArguments/DenseThresholdCases/AfStmt.lean}{UpperTheory.V.Af.NumberSp4Alt}
\leanref{Development/AutomorphicArithmetic/AutomorphicVarieties/ModularityArgument/ThresholdArguments/DenseThresholdCases/AfStmt.lean}{UpperTheory.V.Af.AdmRAlt}
\leanref{Development/AutomorphicArithmetic/AutomorphicVarieties/ModularityArgument/ThresholdArguments/DenseThresholdCases/VWLevel.lean}{UpperTheory.GalN.V.numberSp4_n_alt}
\leanref{Development/AutomorphicArithmetic/AutomorphicVarieties/ModularityArgument/ThresholdArguments/DenseThresholdCases/VWinNumN.lean}{UpperTheory.GalN.V.Win.N.number_sp4_of_admissible_to_admissible_n}

\subsection{Rings of lifts and their graded models}\label{subsection-sha-lift-rings}

\begin{definition}\label{sha-lifts-problem}
Let $\ell$ be an odd prime, let $\mathcal O$ be as in Definition~\ref{sha-character-parts}, with maximal ideal $\mathfrak m_{\mathcal O}$ and residue field $k$, let $e\ge1$, let $N\ge2$ be even, and let $\psi$ be a Dirichlet character modulo $N$ with values in $\mathcal O$ such that $\psi(-1)=(-1)^e$.
\begin{enumerate}
\item Let $S$ be the set of primes dividing $\ell N$, let $\mathbf Q_S\subset\bar{\mathbf Q}$ be the maximal extension of $\mathbf Q$ unramified at the primes outside $S$, and put $\Gamma=\operatorname{Gal}(\mathbf Q_S/\mathbf Q)$.
For each $p\in S$ fix a place of $\mathbf Q_S$ above $p$, with decomposition group $D_p\subset\Gamma$ and inertia group $I_p\subset D_p$.
Let $c\in\Gamma$ be a complex conjugation.
\item Let $\eta:\Gamma\to\{\pm1\}$ be the character with kernel $\operatorname{Gal}(\mathbf Q_S/\mathbf Q(i))$, and let $\Psi=\chi_{\mathrm{cyc}}^{e}\psi:\Gamma\to\mathcal O^\times$, with reduction $\bar\Psi:\Gamma\to k^\times$.
Put $N_p=I_p\cap\ker\eta\cap\ker\bar\Psi$ for $p\in S$.
\item Let $\kappa$ and $\theta_0$ be the parts of $\psi$ of $\ell$-power order and of order prime to $\ell$, let $\psi_2$ be the $2$-component of $\psi$, of conductor $2^{c_2}$, and put $\theta_1=\theta_0\psi_2^{-1}$.
Let $f_0$ be the order of $\ell$ modulo $\operatorname{lcm}(4,\operatorname{ord}\theta_1)$, and $n_0=v_2(\ell^{f_0}-1)$.
Let $S_b$ be the set of primes $p\in S$ other than $\ell$ at which the component of $\kappa$ is nontrivial, and put $T_b=S_b\cup\{\ell\}$.
\item Let $o,q\ge1$ be integers.
The data satisfy the \emph{dyadic condition} for $o$ if $c_2\ge4$, $8op^2<2^{c_2}$ for every $p\in T_b$, and the $r$-component of $\theta_0$ has order dividing $2o$ for every odd prime $r$.
They satisfy the \emph{Bernoulli condition} for $q$ if the conductor of $\theta_1^2$ divides $q$ and $2\ell q\,4^{n_0}\le2^{c_2}$.
\item For a ring $A$ and $t\in A^\times$ let $\tau(t)=\operatorname{diag}(t,1,t^{-1})$.
Let $r_0=r_0(\eta,\Psi)$ be as in Definition~\ref{sha-ordinary-deformation}, so that $r_0(g)=(\tau(\Psi(g)),\eta(g))$, and let $J$ be the antidiagonal $3\times3$ matrix with entries $1$.
The entry $(a,b)$ of a $3\times3$ matrix has \emph{weight} $b-a$.
\item Let $g_1\in\Gamma$ with $\eta(g_1)=1$ and $\Psi(g_1)^2-1\in\mathcal O^\times$, and let $A$ be a complete noetherian local $\mathcal O$-algebra with residue field $k$.
A \emph{unitary lift} with values in $A$, relative to $g_1$, is a homomorphism $r:\Gamma\to\mathcal G_3(A)$ with the following properties.
($\alpha$) It has multiplier $\eta$, it reduces to $r_0$ modulo $\mathfrak m_A$, and it is continuous for the $\mathfrak m_A$-adic topology.
($\beta$) The matrix part of $r(c)$ is $1$, and the $(1,2)$ and $(2,1)$ entries of the matrix part of $r(g_1)-r(cg_1^{-1}c)$ vanish.
($\gamma$) There is $h\in GL_3(A)$ with $h\equiv1\pmod{\mathfrak m_A}$ such that the matrix parts of $h^{-1}r(g)h$ for $g\in D_\ell\cap\ker\eta$ are upper triangular, with diagonal entries $\Psi(g),1,\Psi(g)^{-1}$ when $g\in I_\ell$.
($\delta$) For each $p\in S_b$ there is $h\in GL_3(A)$ with $h\equiv1\pmod{\mathfrak m_A}$ such that the matrix parts of $h^{-1}r(g)h$ for $g\in D_p\cap\ker\eta$ are diagonal, with diagonal entries $\Psi(g),1,\Psi(g)^{-1}$ when $g\in I_p$.
($\varepsilon$) $r(s)=1$ for $s\in N_p$ and $p\in S\setminus T_b$.
\end{enumerate}
\end{definition}
\leanref{Development/AutomorphicArithmetic/AutomorphicVarieties/UnitaryThreeDeformations/IsHostFrame.lean}{UpperTheory.V.UnitaryThreeInertia.IsHostFrame}
\leanref{Development/AutomorphicArithmetic/AutomorphicVarieties/UnitaryThreeDeformations/IsHostFrame.lean}{UpperTheory.V.UnitaryThreeInertia.HostData}
\leanref{Development/AutomorphicArithmetic/AutomorphicVarieties/UnitaryThreeDeformations/IsHostFrame.lean}{UpperTheory.V.UnitaryThreeInertia.XData}
\leanref{Development/AutomorphicArithmetic/AutomorphicVarieties/UnitaryThreeDeformations/AdjointDiagonalRepPreservesWeightPiece.lean}{UpperTheory.V.UnitaryThreeInertia.Frame}
\leanref{Development/AutomorphicArithmetic/AutomorphicVarieties/UnitaryThreeDeformations/AdjointDiagonalRepPreservesWeightPiece.lean}{UpperTheory.V.UnitaryThreeInertia.Frame.LocAt}
\leanref{Development/AutomorphicArithmetic/AutomorphicVarieties/UnitaryThreeDeformations/AdjointDiagonalRepPreservesWeightPiece.lean}{UpperTheory.V.UnitaryThreeInertia.Frame.IsPt}
\leanref{Development/AutomorphicArithmetic/AutomorphicVarieties/UnitaryThreeDeformations/HostEx.lean}{UpperTheory.V.UnitaryThreeInertia.exists_hostFrame}

\begin{definition}\label{sha-lifts-graded}
In Definition~\ref{sha-lifts-problem}, let $u:\Gamma\to\mathcal G_3(R)$ be a unitary lift.
\begin{enumerate}
\item An element $x\in R$ has \emph{weight} $j\in\mathbf Z$ if $f'(x)=t^{j}f(x)$ whenever $f,f':R\to A$ are homomorphisms of $\mathcal O$-algebras to a complete noetherian local $\mathcal O$-algebra $A$ with residue field $k$, and $t\in A^\times$, such that $f'\circ u=\tau(t)(f\circ u)\tau(t)^{-1}$.
The elements of weight $j$ form an $\mathcal O$-submodule $R_j$, with $R_iR_j\subset R_{i+j}$, and the entry $(a,b)$ of the matrix part of every $u(g)$ lies in $R_{b-a}$.
Let $R^{\mathrm{gr}}=\sum_jR_j$, an $\mathcal O$-subalgebra of $R$.
\item A \emph{point} of $u$ is a homomorphism of $\mathcal O$-algebras $\xi:R^{\mathrm{gr}}\to\bar{\mathbf Q}_\ell$.
Its \emph{representation} $r_\xi:\Gamma\to\mathcal G_3(\bar{\mathbf Q}_\ell)$ has the multiplier of $u$ and the matrix entries $\xi(u(g)_{ab})$.
\end{enumerate}
Let $r:\Gamma\to\mathcal G_3(\bar{\mathbf Q}_\ell)$ be a homomorphism with multiplier $\eta$.
\begin{enumerate}
\setcounter{enumi}{2}
\item $r$ is \emph{of characteristic zero} if for some $t\in\bar{\mathbf Q}_\ell^\times$ the homomorphism $\tau(t)r\tau(t)^{-1}$ is a unitary lift with values in a discrete valuation ring $A\subset\bar{\mathbf Q}_\ell$ finite over $\mathcal O$.
\item $r$ is \emph{irreducible} if no subspace of $\bar{\mathbf Q}_\ell^3$ other than $0$ and $\bar{\mathbf Q}_\ell^3$ is stable under the matrix parts of $r(g)$ for $g\in\ker\eta$.
\item $\Gamma$ acts on $M_3(\bar{\mathbf Q}_\ell)$ by $\operatorname{ad}r(g)X=MXM^{-1}$ if $r(g)=(M,1)$ and $\operatorname{ad}r(g)X=-MJ\,{}^tX\,JM^{-1}$ if $r(g)=(M,-1)$.
A \emph{cocycle} for $r$ is a continuous $y:\Gamma\to M_3(\bar{\mathbf Q}_\ell)$ with $y(gh)=y(g)+\operatorname{ad}r(g)y(h)$, and the \emph{coboundary} of $X$ is $g\mapsto X-\operatorname{ad}r(g)X$.
$r$ is \emph{rigid} if every cocycle $y$ for $r$ with the following properties is a coboundary: $y$ vanishes on $N_p$ for $p\in S\setminus T_b$, and the homomorphism $g\mapsto((1+\epsilon y(g))M_g,\eta(g))$, where $r(g)=(M_g,\eta(g))$, with values in $\mathcal G_3(\bar{\mathbf Q}_\ell[\epsilon]/(\epsilon^2))$, satisfies the conditions in ($\gamma$) and ($\delta$) of Definition~\ref{sha-lifts-problem} for some $h\in GL_3(\bar{\mathbf Q}_\ell[\epsilon]/(\epsilon^2))$.
\end{enumerate}
\end{definition}
\leanref{Development/AutomorphicArithmetic/AutomorphicVarieties/UnitaryThreeDeformations/IsCharZeroPt.lean}{UpperTheory.V.UnitaryThreeInertia.wpart}
\leanref{Development/AutomorphicArithmetic/AutomorphicVarieties/UnitaryThreeDeformations/IsCharZeroPt.lean}{UpperTheory.V.UnitaryThreeInertia.Sring}
\leanref{Development/AutomorphicArithmetic/AutomorphicVarieties/UnitaryThreeDeformations/IsCharZeroPt.lean}{UpperTheory.V.UnitaryThreeInertia.IsPtHom}
\leanref{Development/AutomorphicArithmetic/AutomorphicVarieties/UnitaryThreeDeformations/IsCharZeroPt.lean}{UpperTheory.V.UnitaryThreeInertia.Frame.IsCharZeroPt}
\leanref{Development/AutomorphicArithmetic/AutomorphicVarieties/UnitaryThreeGroups/BaseRepGL3NormalisedOfPsiEqOne.lean}{UpperTheory.V.U3.TypeI}
\leanref{Development/AutomorphicArithmetic/AutomorphicVarieties/UnitaryThreeDeformations/IsCharZeroPt.lean}{UpperTheory.V.UnitaryThreeInertia.Frame.Rigid}

\begin{definition}\label{sha-lifts-selmer}
In Definition~\ref{sha-lifts-problem}, let $\bar r_0$ be the reduction of $r_0$, and let $\Gamma$ act on $M_3(k)$ by $\operatorname{ad}\bar r_0$, defined as in Definition~\ref{sha-lifts-graded}(e); cocycles and coboundaries are as there.
\begin{enumerate}
\item For $j\in\mathbf Z$ let $\mathfrak g_j\subset M_3(k)$ be the subspace of matrices supported on the entries of weight $j$.
It is stable under $\Gamma$, and has dimension $3-|j|$ if $|j|\le2$ and $0$ otherwise.
\item For $p\in S$ let $L_{j,p}$ be the space of continuous cocycles $y:D_p\to\mathfrak g_j$ such that $y$ vanishes on $N_p$ if $p\ne\ell$, $y$ is a coboundary if $p=\ell$ and $j<0$, and $y$ vanishes on $I_\ell$ if $p=\ell$ and $j=0$.
\item The \emph{Selmer group} $H^1_L(\mathfrak g_j)$ is the space of continuous cocycles $y:\Gamma\to\mathfrak g_j$ with $y|_{D_p}\in L_{j,p}$ for every $p\in S$, modulo coboundaries.
\item The \emph{obstruction group} $\mathrm{Ob}_j$ is the quotient of the space of tuples $(z,(y_p)_{p\in S})$, where $z$ is a continuous $2$-cocycle of $\Gamma$ with values in $\mathfrak g_j$ and $y_p:D_p\to\mathfrak g_j$ is continuous with $dy_p=z|_{D_p}$, by the subspace spanned by the tuples $(dx,(x|_{D_p})_p)$ for continuous $x:\Gamma\to\mathfrak g_j$ and the tuples $(0,(\lambda_p)_p)$ with $\lambda_p\in L_{j,p}$.
\item Let $\delta(-1)=-1$, $\delta(1)=\delta(2)=1$ and $\delta(j)=0$ for the other $j\in\mathbf Z$.
\item For a nonzero $a\in\mathcal O$, let $\text{Ш}_a$ be the $\mathcal O$-submodule of $H^1(\Gamma,\mathcal O/a\mathcal O(\Psi^{-1}))$ generated by the classes of the continuous cocycles that vanish on $N_p$ for $p\in S\setminus T_b$ and are coboundaries on $D_p$ for $p\in T_b$.
Here $\sigma$ acts on $\mathcal O/a\mathcal O$ by $\Psi(\sigma)^{-1}$.
\end{enumerate}
\end{definition}
\leanref{Development/AutomorphicArithmetic/AutomorphicVarieties/UnitaryThreeDeformations/AdjointDiagonalRepPreservesWeightPiece.lean}{UpperTheory.V.UnitaryThreeInertia.Frame.ZL}
\leanref{Development/AutomorphicArithmetic/AutomorphicVarieties/UnitaryThreeDeformations/AdjointDiagonalRepPreservesWeightPiece.lean}{UpperTheory.V.UnitaryThreeInertia.Frame.selmerCohomologyQuotient}
\leanref{Development/AutomorphicArithmetic/AutomorphicVarieties/UnitaryThreeDeformations/AdjointDiagonalRepPreservesWeightPiece.lean}{UpperTheory.V.UnitaryThreeInertia.Frame.Ob}
\leanref{Development/AutomorphicArithmetic/AutomorphicVarieties/UnitaryThreeDeformations/IsCharZeroPt.lean}{UpperTheory.V.UnitaryThreeInertia.Frame.ShaO}

\begin{lemma}\label{sha-lifts-host-data}
Let $\psi$ be a primitive Dirichlet character, let $h$ be a host datum for $\psi$, let $\ell$ be an odd prime, let $\mathcal O$ be as in Definition~\ref{sha-character-parts}, and let $\psi'\in X_h(1)$.
Let $n_0$ and $2^{c_2}$ be attached to $\psi'$ as in Proposition~\ref{sha-counts}(d), and suppose that $16\,o_\psi(\ell N_\psi)^2\,4^{n_0}\le2^{c_2}$.
Then $N_h$ is even and $\psi'(-1)=-1$, and the data $(\ell,\mathcal O,1,\psi')$ of Definition~\ref{sha-lifts-problem}, with $N=N_h$, satisfy the dyadic condition for $o_\psi$ and the Bernoulli condition for $N_\psi$.
\end{lemma}
\begin{proof}
Since $\psi'\in X_h(1)$, $\psi'(-1)=-1$, and $16$ divides $N_h$.
Let $\theta_0$, $\psi'_2$ and $\theta_1$ be attached to $\psi'$ as in Definition~\ref{sha-lifts-problem}(c).
Since $\psi'$ is trivial on $H_h$, the $r$-component of $\theta_0$ has order dividing $2o_\psi$ for every odd prime $r$, and the conductor of $\theta_1^2$ divides $N_\psi$; this is a computation in $(\mathbf Z/N_h)^\times/H_h$.
Put $R_1=16\,o_\psi(\ell N_\psi)^2\,4^{n_0}$.
Since $\ell^{f_0}\equiv1\pmod4$, $n_0\ge2$, so $R_1\le2^{c_2}$ gives $c_2\ge12$.
A prime $p\in S_b$ divides $N_h$ and is not $2$, since a character of the $2$-group $(\mathbf Z/2^m)^\times$ of $\ell$-power order is trivial; it is not $q$ or $d$, where the components of $\psi'$ have order at most $2$, because $\psi'$ is trivial on the squares there.
So $p$ divides $N_\psi$, and $8o_\psi p^2<R_1\le2^{c_2}$; the same holds for $p=\ell$.
This is the dyadic condition for $o_\psi$.
Finally $2\ell N_\psi\,4^{n_0}\le R_1\le2^{c_2}$, which is the Bernoulli condition for $N_\psi$.
\end{proof}
\leanref{Development/AutomorphicArithmetic/AutomorphicVarieties/UnitaryThreeDeformations/HostDataOf.lean}{UpperTheory.V.UnitaryThreeInertia.hostData_of_host}
\leanref{Development/AutomorphicArithmetic/AutomorphicVarieties/HeightAndVanishing/OrdProjNDvd.lean}{UpperTheory.V.HV.hostTh1}
\leanref{Development/AutomorphicArithmetic/AutomorphicVarieties/HeightAndVanishing/CmDatumCharTwoAdicOrderBounds.lean}{UpperTheory.V.HV.cmDatum_char_two_adic_order_bounds}

\begin{lemma}\label{sha-lifts-local}
In Definition~\ref{sha-lifts-problem}, suppose that the data satisfy the dyadic condition for some $o\ge1$.
\begin{enumerate}
\item $\eta(c)=-1$, $\Psi(c)=1$, and $\eta$ is trivial on $I_p$ for every $p\in T_b$.
\item For every $p\in T_b$ there is $x_p\in D_p$ with $\eta(x_p)=1$, $\chi_{\mathrm{cyc}}(x_p)\equiv1\pmod\ell$ and $\Psi(x_p)^2-1\in\mathcal O^\times$.
\item For every $p\in S\setminus T_b$, $\Psi$ is trivial on $N_p$.
\item For every $p\in S$, the index of $N_p$ in $I_p$ is prime to $\ell$.
\item For every open subgroup $\Gamma'$ of $\Gamma$, the group $H^1(\Gamma',\mathbf Z/\ell)$ is finite.
\end{enumerate}
\end{lemma}
\begin{proof}
(a) Complex conjugation inverts the roots of unity, so $\eta(c)=-1$, $\chi_{\mathrm{cyc}}(c)=-1$ and $\Psi(c)=(-1)^e\psi(-1)=1$.
The primes of $T_b$ are odd: $\ell$ is odd, and the $2$-component of $\kappa$ is trivial because $\kappa$ has odd order.
Since $\mathbf Q(i)$ is unramified at the odd primes, $\eta$ is trivial on $I_p$ for $p\in T_b$.

(b) Let $p\in T_b$.
If $p\ne\ell$, then $p\equiv1\pmod\ell$, because the group $(\mathbf Z/p^m)^\times$ has a nontrivial character of $\ell$-power order.
Choose $\sigma\in D_p$ that acts on the roots of unity of order prime to $p$ by $\zeta\mapsto\zeta^p$ and satisfies $\chi_{\mathrm{cyc}}(\sigma)\equiv1\pmod\ell$: for $p\ne\ell$ a Frobenius element, and for $p=\ell$ a Frobenius element times an element of $I_\ell$ that corrects its cyclotomic character.
Put $x_p=\sigma^2$, so that $\eta(x_p)=1$ and $\chi_{\mathrm{cyc}}(x_p)\equiv1\pmod\ell$.
Let $u\in(\mathbf Z/N)^\times$ be the image of $\sigma$; it is congruent to $p$ modulo the $2$-part of $N$, and $\Psi(x_p)^2\equiv\psi(u)^4$ modulo $\mathfrak m_{\mathcal O}$.
Suppose that $\psi(u)^4\equiv1$.
Since $\kappa(u)\equiv1$, we get $\theta_0(u)^4\equiv1$, and $\theta_0(u)^4=1$ because a root of unity of order prime to $\ell$ is determined by its residue.
The $2$-component of $\theta_0$ is $\psi_2$, and its odd components have orders dividing $2o$, so $\theta_0^{4o}=\psi_2^{4o}$ and $\psi_2(u)^{4o}=1$.
This is impossible.
Indeed $v_2(p^{4o}-1)=v_2(p^2-1)+v_2(2o)<c_2-2$ because $2op^2<2^{c_2-2}$, so the powers of $p^{4o}$ in $(\mathbf Z/2^{c_2})^\times$ include $1+2^{c_2-1}$, on which the primitive character $\psi_2$ of conductor $2^{c_2}$ is nontrivial.
So $\Psi(x_p)^2-1$ is a unit.

(c) Let $p\in S\setminus T_b$.
Then $p\ne\ell$, so $\chi_{\mathrm{cyc}}$ is trivial on $I_p$ and $\Psi=\psi$ on $I_p$.
Since $p\notin S_b$, the $p$-component of $\kappa$ is trivial, so $\Psi$ has order prime to $\ell$ on $I_p$, and its values there are determined by their residues.
Hence $\Psi$ is trivial on $I_p\cap\ker\bar\Psi$, which contains $N_p$.

(d) $N_p$ is the kernel of $(\eta,\bar\Psi)$ on $I_p$, whose image has order dividing $2(\#k-1)$.

(e) This is the finiteness of the cohomology, with finite coefficients of order prime to the primes outside $S$, of the Galois group of the maximal extension of a number field unramified outside a finite set of places containing the primes above $\ell$ \cite[Chapter VIII, \S3]{NSW08}.
\end{proof}
\leanref{Development/AutomorphicArithmetic/AutomorphicVarieties/UnitaryThreeDeformations/HostOk.lean}{UpperTheory.V.UnitaryThreeInertia.hostFrame_ok}
\leanref{Development/AutomorphicArithmetic/AutomorphicVarieties/UnitaryThreeDeformations/AdjointDiagonalRepPreservesWeightPiece.lean}{UpperTheory.V.UnitaryThreeInertia.Frame.Ok}
\leanref{Development/AutomorphicArithmetic/AutomorphicVarieties/UnitaryThreeDeformations/DyadicCharacterPowerNeOne.lean}{UpperTheory.V.UnitaryThreeInertia.exists_sigma}
\leanref{Development/AutomorphicArithmetic/AutomorphicVarieties/UnitaryThreeDeformations/HostPhi.lean}{UpperTheory.V.UnitaryThreeInertia.hostFrame_phi}

\begin{lemma}\label{sha-lifts-universal}
In Definition~\ref{sha-lifts-problem}, suppose that the data satisfy the dyadic condition for some $o\ge1$, and let $g_1$ be as in (f).
Then there are a complete noetherian local $\mathcal O$-algebra $R$ with residue field $k$ and a unitary lift $u$ with values in $R$ that is \emph{universal}: for every unitary lift $r$ with values in $A$ there is a unique homomorphism of $\mathcal O$-algebras $f:R\to A$ with $f\circ u=r$.
Moreover, if $r$ is a unitary lift with values in $A$ and $t\in A^\times$, then $\tau(t)r\tau(t)^{-1}$ is a unitary lift.
\end{lemma}
\begin{proof}
By Lemma~\ref{sha-lifts-local}(e), $\Gamma$ satisfies Mazur's finiteness condition at $\ell$ \cite{Maz89}.
Let $\theta(M)=J\,{}^tM^{-1}J$, and embed $\mathcal G_3(A)$ in $GL_6(A)$ by sending $(M,1)$ to the block matrix with diagonal blocks $M$ and $\theta(M)$, and $(M,-1)$ to the block matrix with antidiagonal blocks $M$ (upper right) and $\theta(M)$ (lower left).
A matrix with blocks $X_{11},X_{12},X_{21},X_{22}$ is the image of some $(M,1)$ if and only if $X_{12}$, $X_{21}$ and $X_{11}J\,{}^tX_{22}J-1$ vanish, and of some $(M,-1)$ if and only if $X_{11}$, $X_{22}$ and $X_{12}J\,{}^tX_{21}J-1$ vanish.
The continuous lifts of the composite of $\bar r_0$ with this embedding form a functor represented by a complete noetherian local $\mathcal O$-algebra with residue field $k$: it satisfies Schlessinger's criteria, and its tangent space is finite by the finiteness condition, as in \cite{Maz89}.
Its quotient by the ideal generated by the entries of the equations above at $\eta(g)$, evaluated on the universal lift at every $g$, represents the continuous lifts of $\bar r_0$ to $\mathcal G_3$ with multiplier $\eta$.
It suffices to show that each of the conditions ($\beta$) to ($\varepsilon$) holds for the composite of the universal lift with $f$ if and only if $f$ kills a certain ideal, and to take the quotient by the sum of these ideals.
For ($\beta$) and ($\varepsilon$) the ideal is generated by matrix entries.
For ($\gamma$) and ($\delta$) at $p\in T_b$, let $x_p$ be as in Lemma~\ref{sha-lifts-local}(b).
The residues of $\Psi(x_p),1,\Psi(x_p)^{-1}$ are distinct, since $\Psi(x_p)^2-1$ and hence $\Psi(x_p)-1$ are units.
So for every lift $r$ there is a unique $h_r\equiv1$ with diagonal entries $1$ such that $h_r^{-1}r(x_p)h_r$ is diagonal, and its entries are given by formulas in the entries of $r(x_p)$.
A matrix $h$ as in ($\gamma$) or ($\delta$) diagonalizes $r(x_p)$ up to an upper triangular matrix congruent to $1$, so the condition holds for some $h$ if and only if it holds for $h_r$; for $h_r$ it is the vanishing of entries and the equality of diagonal entries with values of $\Psi$, again generated by matrix entries.
Finally, conjugation by $\tau(t)$ fixes $r_0$, preserves the triangular and diagonal shapes and the multiplier, and maps $(M,-1)$ to $(\tau(t)M\tau(t)^{-1},-1)$ because $J\,{}^t\tau(t)J=\tau(t)^{-1}$; it multiplies the entry $(a,b)$ by $t^{b-a}$, so it preserves ($\beta$).
\end{proof}
\leanref{Development/AutomorphicArithmetic/AutomorphicVarieties/UnitaryThreeDeformations/ExistsIsUniv.lean}{UpperTheory.V.UnitaryThreeInertia.exists_isUniv}
\leanref{Development/AutomorphicArithmetic/AutomorphicVarieties/UnitaryThreeDeformations/IsCharZeroPt.lean}{UpperTheory.V.UnitaryThreeInertia.Frame.IsUniv}

\begin{lemma}\label{sha-lifts-residual}
In Definition~\ref{sha-lifts-problem}, suppose that the data satisfy the dyadic condition for some $o\ge1$ and the Bernoulli condition for some $q\ge1$.
\begin{enumerate}
\item For $p\in T_b$ and $j\ne0$, $H^2(D_p,\mathfrak g_j)=0$; for $p\in S_b$ and $j\ne0$, also $H^1(D_p,\mathfrak g_j)=0$.
\item For every $j$, $H^1_L(\mathfrak g_j)$ and $\mathrm{Ob}_j$ are finite-dimensional, and $\dim H^1_L(\mathfrak g_j)-\dim\mathrm{Ob}_j=\delta(j)$.
\item Every continuous homomorphism from $\ker\eta$ to an abelian group that is trivial on every conjugate in $\ker\eta$ of every $I_p$, $p\in S$, is trivial.
\item Let $j\le0$, and let $\gamma:\Gamma\to M_3(\mathcal O)$ be a continuous cocycle for $\operatorname{ad}r_0$ with values in the matrices supported on the entries of weight $j$.
Suppose that $\gamma$ vanishes on $N_p$ for $p\in S\setminus T_b$, that $\gamma$ vanishes on $I_p$ for $p\in T_b$ if $j=0$, and that $\gamma$ agrees on $D_\ell\cap\ker\eta$ with the coboundary of a matrix of weight $j$ if $j<0$.
Then $a\gamma$ is a coboundary for some nonzero $a\in\mathcal O$.
\item $H^1_L(\mathfrak g_{-2})=0$.
\end{enumerate}
\end{lemma}
\begin{proof}
(a) Let $p\in T_b$, let $j\ne0$, and let $x=x_p$ be as in Lemma~\ref{sha-lifts-local}(b).
Since $\eta(x)=1$, $x$ acts on $\mathfrak g_j$ by the scalar $\bar\Psi(x)^{j}$, which is not $1$ because $|j|\le2$ and $\bar\Psi(x)^2\ne1$.
Since $\chi_{\mathrm{cyc}}(x)\equiv1\pmod\ell$, $x$ acts by a scalar other than $1$ on the dual of $\mathfrak g_j$ twisted by the cyclotomic character as well.
So that twisted dual has no invariants under $D_p$, and $H^2(D_p,\mathfrak g_j)=0$ by local Tate duality \cite[Chapter VII, \S2]{NSW08}.
For $p\in S_b$ we have $p\ne\ell$, $H^0(D_p,\mathfrak g_j)=0$ by the same element, and the local Euler characteristic formula \cite[Chapter VII, \S3]{NSW08} gives $\dim H^1(D_p,\mathfrak g_j)=\dim H^0+\dim H^2=0$.

(b) We may assume $|j|\le2$.
The groups $H^1_L(\mathfrak g_j)$ and $\mathrm{Ob}_j$ are the first and second cohomology groups of the Selmer complex of $\mathfrak g_j$ with the local conditions $L_{j,p}$, and they are finite-dimensional because the Galois cohomology of $\Gamma$ and of the $D_p$ with coefficients in $\mathfrak g_j$ is.
The map $H^2(\Gamma,\mathfrak g_j)\to\bigoplus_{p\in S}H^2(D_p,\mathfrak g_j)$ is surjective by Poitou-Tate duality \cite[Chapter VIII, \S6]{NSW08}, since some element of $\Gamma$ acts on $\mathfrak g_j$ by a scalar different from its cyclotomic character modulo $\ell$: $x_\ell$ if $j\ne0$, and an element fixing $i$ whose cyclotomic character is not $1$ modulo $\ell$ if $j=0$.
So the third cohomology of the Selmer complex vanishes, and its Euler characteristic, computed with the global and local Euler characteristic formulas \cite[Chapter VII, \S3 and Chapter VIII, \S7]{NSW08}, is the Greenberg-Wiles formula
\[
 \dim H^1_L(\mathfrak g_j)-\dim\mathrm{Ob}_j=\dim H^0(\Gamma,\mathfrak g_j)-\dim\mathfrak g_j^{c=1}+\sum_{p\in S}\big(\dim L_{j,p}/B_{j,p}-\dim H^0(D_p,\mathfrak g_j)\big),
\]
where $B_{j,p}\subset L_{j,p}$ is the space of coboundaries and $\mathfrak g_j^{c=1}$ is the space fixed by $c$.
For $p\ne\ell$, and for $p=\ell$ when $j=0$, the term of $p$ is $0$: the classes in $L_{j,p}/B_{j,p}$ vanish on $N_p$, which has index prime to $\ell$ in $I_p$ by Lemma~\ref{sha-lifts-local}(d), so they are unramified, and since $D_p/I_p$ is procyclic the unramified classes have dimension $\dim H^0(D_p,\mathfrak g_j)$.
For $p=\ell$ and $j<0$ the term is $0$, since $L_{j,\ell}=B_{j,\ell}$ and $H^0(D_\ell,\mathfrak g_j)=0$.
For $p=\ell$ and $j>0$, $L_{j,\ell}$ consists of all cocycles, $H^0(D_\ell,\mathfrak g_j)=H^2(D_\ell,\mathfrak g_j)=0$ by (a), and the local Euler characteristic formula makes the term $\dim\mathfrak g_j$.
Finally $H^0(\Gamma,\mathfrak g_j)=0$ for $j\ne0$, while $H^0(\Gamma,\mathfrak g_0)=\mathfrak g_0^{c=1}$ because $\ker\eta$ acts trivially on the diagonal matrices.
The element $c$ acts on $\mathfrak g_j$ by $X\mapsto-J\,{}^tXJ$, and $2$ is invertible in $k$, so $\dim\mathfrak g_j^{c=1}$ is $1$ if $|j|\le1$ and $0$ if $|j|=2$.
Adding the terms gives $\delta(j)$ in each case.

(c) Such a homomorphism has open kernel, so it factors through the Galois group of a finite abelian extension of $\mathbf Q(i)$ that is unramified at every finite place.
Since $\mathbf Q(i)$ has class number $1$, this extension is trivial.

(d) Let $j=0$, so that $\gamma$ takes values in the diagonal matrices, on which $\ker\eta$ acts trivially.
Then $\gamma$ restricts to a homomorphism on $\ker\eta$, and $\gamma(gsg^{-1})=\operatorname{ad}r_0(g)\gamma(s)$ for $s\in\ker\eta$ and $g\in\Gamma$.
It vanishes on $N_p$ for $p\in S\setminus T_b$, hence on $I_p\cap\ker\eta$, since the index is finite and $\mathcal O$ is torsion free; it vanishes on $I_p$ for $p\in T_b$ by hypothesis.
So $\gamma$ vanishes on $\ker\eta$ by (c).
Then $0=\gamma(c^2)=\gamma(c)+\operatorname{ad}r_0(c)\gamma(c)$, and $\gamma$ is the coboundary of $\gamma(c)/2$.
Now let $j=-n<0$.
Each entry of $\gamma$ of weight $j$ is a cocycle for a character $(\chi_{\mathrm{cyc}}^{ne}\theta)^{-1}$ with $\theta$ of finite order, and by hypothesis it is a coboundary on $D_\ell\cap\ker\eta$, hence becomes one on $D_\ell$ after multiplication by $2$.
At the primes $p\ne\ell$ the local $H^1$ with coefficients in $\mathcal O(\chi_{\mathrm{cyc}}^{ne}\theta)^{-1}$ is torsion, since a Frobenius element acts on the inertia invariants by an eigenvalue of absolute value $p^{-ne}\ne1$.
So the class of the entry lies, up to torsion, in the Selmer group with the conditions of being a coboundary at $\ell$.
By Poitou-Tate duality this Selmer group is torsion if $H^2(G_{L,S_\ell},\mathbf Z_\ell(ne+1))$ is finite for a totally complex field $L$ containing $\mu_\ell$ on whose absolute Galois group $\theta$ is trivial, and that holds by Theorem~\ref{sha-soule-bound}.
Hence a nonzero multiple of each entry, and so of $\gamma$, is a coboundary.

(e) The space $\mathfrak g_{-2}$ is the line of the entry $(3,1)$, on which $\Gamma$ acts by the character $\bar\Psi^{-2}\eta$.
A class in $H^1_L(\mathfrak g_{-2})$ is represented by a cocycle for this character that vanishes on $N_p$ for $p\ne\ell$ and is a coboundary on $D_\ell$.
The kernel of $(\eta,\bar\Psi)$ has index prime to $\ell$ in $\Gamma$, and on it the cocycle is a homomorphism that is unramified everywhere and trivial on the decomposition groups above $\ell$.
By the residual nonvanishing of Bernoulli numbers for the character $\bar\Psi^2\eta$, also used in the proof of Proposition~\ref{sha-counts}(b), whose hypotheses are given by the Bernoulli condition for $q$, such a cocycle is a coboundary.
\end{proof}
\leanref{Development/AutomorphicArithmetic/AutomorphicVarieties/UnitaryThreeDeformations/HostLocVan.lean}{UpperTheory.V.UnitaryThreeInertia.hostFrame_locVan}
\leanref{Development/AutomorphicArithmetic/AutomorphicVarieties/UnitaryThreeDeformations/HostCount.lean}{UpperTheory.V.UnitaryThreeInertia.hostFrame_count}
\leanref{Development/AutomorphicArithmetic/AutomorphicVarieties/UnitaryThreeDeformations/FrameClassOne.lean}{UpperTheory.V.UnitaryThreeInertia.frame_class_one}
\leanref{Development/AutomorphicArithmetic/AutomorphicVarieties/LevelCohomology/SoHostVanE.lean}{UpperTheory.V.X.So.hostVanE}
\leanref{Development/AutomorphicArithmetic/AutomorphicVarieties/LevelCohomology/SoLineNeg.lean}{UpperTheory.V.X.So.lineNeg}
\leanref{Development/AutomorphicArithmetic/AutomorphicVarieties/LevelCohomology/Xpp.lean}{UpperTheory.V.X.xpp}
\leanref{Development/AutomorphicArithmetic/AutomorphicVarieties/UnitaryThreeDeformations/AdjointDiagonalRepPreservesWeightPiece.lean}{UpperTheory.V.UnitaryThreeInertia.Frame.LocVan}
\leanref{Development/AutomorphicArithmetic/AutomorphicVarieties/UnitaryThreeDeformations/AdjointDiagonalRepPreservesWeightPiece.lean}{UpperTheory.V.UnitaryThreeInertia.Frame.Count}
\leanref{Development/AutomorphicArithmetic/AutomorphicVarieties/UnitaryThreeDeformations/IsCharZeroPt.lean}{UpperTheory.V.UnitaryThreeInertia.Frame.Cl1}
\leanref{Development/AutomorphicArithmetic/AutomorphicVarieties/UnitaryThreeDeformations/IsCharZeroPt.lean}{UpperTheory.V.UnitaryThreeInertia.Frame.VanE}
\leanref{Development/AutomorphicArithmetic/AutomorphicVarieties/UnitaryThreeDeformations/AdjointDiagonalRepPreservesWeightPiece.lean}{UpperTheory.V.UnitaryThreeInertia.Frame.Xpp}

\begin{lemma}\label{sha-lifts-presentation}
In Definition~\ref{sha-lifts-problem}, suppose that the data satisfy the dyadic condition for some $o\ge1$ and the Bernoulli condition for some $q\ge1$, let $g_1$ be as in (f), and let $u$ be the universal unitary lift, with values in $R$ (Lemma~\ref{sha-lifts-universal}).
\begin{enumerate}
\item There are elements $x_1,\dots,x_n\in R$, each of the form $u(g)_{ab}-r_0(g)_{ab}$ for some $g\in\ker\eta$ and of weight $w_i=b-a$, with the following properties.
The homomorphism of $\mathcal O$-algebras $\pi:P=\mathcal O[[X_1,\dots,X_n]]\to R$ with $\pi(X_i)=x_i$ is surjective.
Its kernel $I$ is homogeneous when $X_i$ is given the weight $w_i$, and $R_j$ is the image of the power series of weight $j$.
The coefficients of degree one of the elements of $I$ lie in $\mathfrak m_{\mathcal O}$.
There is a homomorphism of $\mathcal O$-algebras $o:R\to\mathcal O$ with $o(x_i)=0$ for all $i$, and $o\circ u=r_0$.
\item For every $j$, the number of $i$ with $w_i=j$ is at least $\dim H^1_L(\mathfrak g_j)$.
\item $I$ is generated by power series $f_{j,m}$, $1\le m\le\dim\mathrm{Ob}_j$, with $f_{j,m}$ homogeneous of weight $j$.
\item Let $j\le0$ and let $d\in\mathcal O^n$ with $d_i=0$ unless $w_i=j$. If $\sum_i\frac{\partial f}{\partial X_i}(0)\,d_i=0$ for every $f\in I$, then $d=0$.
\item For every nonzero $a\in\mathcal O$ there is an injective additive map $d_a:\text{Ш}_a\to(\mathcal O/a\mathcal O)^n$ such that $d_a(\mu)_i=0$ unless $w_i=-1$, and $\sum_i\frac{\partial f}{\partial X_i}(0)\,d_a(\mu)_i=0$ for every $f\in I$ and $\mu\in\text{Ш}_a$.
\end{enumerate}
\end{lemma}
\begin{proof}
(a) A homomorphism $f:R\to A$ is determined by $f\circ u$, which is determined by its values on $\ker\eta$ and at $c$, and the matrix part of $u(c)$ is $1$.
So the entries $u(g)_{ab}-r_0(g)_{ab}$ with $g\in\ker\eta$, which lie in $\mathfrak m_R$ and in $R_{b-a}$, generate $R$ topologically over $\mathcal O$.
By Lemma~\ref{sha-lifts-universal}, for $f$ and $t\in A^\times$ there is a unique $f'$ with $f'\circ u=\tau(t)(f\circ u)\tau(t)^{-1}$, and it multiplies each such entry of weight $w$ by $t^{w}$.
Choosing finitely many of these entries that give a basis of $\mathfrak m_R/(\mathfrak m_R^2+\mathfrak m_{\mathcal O}R)$, the map $\pi$ is surjective, the action of $t$ on $R$ lifts to $X_i\mapsto t^{w_i}X_i$ on $P$, and this gives the homogeneity of $I$, the description of $R_j$, and the condition on the coefficients of degree one.
Finally $r_0$ is a unitary lift with values in $\mathcal O$, so universality gives $o$ with $o\circ u=r_0$, and $o(x_i)=0$.

(b) Let $y$ be a cocycle representing a class in $H^1_L(\mathfrak g_j)$.
After adding a coboundary, the homomorphism $g\mapsto((1+\epsilon y(g))M_g,\eta(g))$, where $\bar r_0(g)=(M_g,\eta(g))$, with values in $\mathcal G_3(k[\epsilon]/(\epsilon^2))$, satisfies ($\beta$), because $\Psi(g_1)^2-1$ is a unit.
It satisfies ($\gamma$) to ($\varepsilon$): at $\ell$ by the definition of $L_{j,\ell}$ and because classes of positive weight are upper triangular; at $p\in S_b$ because $y|_{D_p}$ is a coboundary when $j\ne0$, by Lemma~\ref{sha-lifts-residual}(a), and vanishes on $I_p$ when $j=0$, since it vanishes on $N_p$, of index prime to $\ell$; and at the other primes because $y$ vanishes on $N_p$.
By universality it comes from a homomorphism $R\to k[\epsilon]/(\epsilon^2)$, whose derivative is a linear form on $\mathfrak m_R/(\mathfrak m_R^2+\mathfrak m_{\mathcal O}R)$ supported on the $x_i$ of weight $j$; distinct classes give distinct forms.

(c) By the Artin-Rees lemma there is $N\ge1$ with $\mathfrak m_P^N\cap I\subset\mathfrak m_PI$.
Let $T\subset P$ be the ideal generated by $\mathfrak m_{\mathcal O}^N$ and the monomials of degree $N$. It is homogeneous and contained in $\mathfrak m_P^N$, and $\bar P=P/T$ is finite; let $\bar I$ be the image of $I$ in $\bar P$, a homogeneous ideal.
Fix $j$, and let $K_j\subset\bar I$ be the ideal generated by $\mathfrak m_{\bar P}\bar I$ and the homogeneous elements of $\bar I$ of weight other than $j$.
Then $A'=\bar P/K_j\to A=\bar P/\bar I$ is surjective, and its kernel $V=\bar I/K_j$ is killed by $\mathfrak m_{A'}$ and has weight $j$. We show that $\dim_kV\le\dim\mathrm{Ob}_j$.
The map $\pi$ induces $\psi:R\to A$, and the matrix entries of the unitary lift $\psi\circ u$ lift to homogeneous elements of $A'$.
The obstruction to lifting $\psi\circ u$ to a unitary lift with values in $A'$ is a class in $\mathrm{Ob}_j\otimes_kV$; it is defined with local lifts at the primes of $S$, which exist by Lemma~\ref{sha-lifts-residual}(a).
If $\dim V>\dim\mathrm{Ob}_j$, some nonzero linear form $\lambda$ on $V$ kills this class, and $\psi\circ u$ lifts to a unitary lift with values in $A''=A'/\ker\lambda$, which is $s\circ u$ for a homomorphism $s:R\to A''$ by universality.
The homomorphisms $s\circ\pi$ and $P\to A'\to A''$ agree modulo the kernel of $A''\to A$, which is killed by $\mathfrak m_{A''}$; so their difference is a derivation that vanishes on $\mathfrak m_P^2$ and on $\mathfrak m_{\mathcal O}P$.
By (a), the coefficients of degree one of the elements of $I$ lie in $\mathfrak m_{\mathcal O}$, so the two homomorphisms agree on $I$, and $I$ maps to $0$ in $A''$; this contradicts $\lambda\ne0$.
So $\dim V\le\dim\mathrm{Ob}_j$, and there are at most $\dim\mathrm{Ob}_j$ homogeneous elements of $\bar I$ of weight $j$ that span $\bar I$ modulo $K_j$.
Taken over all $j$, these elements generate $\bar I$ modulo $\mathfrak m_{\bar P}\bar I$.
Lift each of them to a homogeneous element $f_{j,m}$ of $I$ of weight $j$, and let $I'$ be the ideal generated by the $f_{j,m}$.
Then $I\subset I'+\mathfrak m_PI+(T\cap I)$, and $T\cap I\subset\mathfrak m_P^N\cap I\subset\mathfrak m_PI$, so $I=I'$ by Nakayama's lemma.
This is a graded form of Mazur's bound on the number of relations of a deformation ring \cite{Maz89}.

(d) The vector $d$ defines a homomorphism $R\to\mathcal O[\epsilon]/(\epsilon^2)$ that sends $x_i$ to $\epsilon d_i$, and composed with $u$ it gives a unitary lift $(1+\epsilon\gamma)r_0$ with $\gamma$ a cocycle for $\operatorname{ad}r_0$ supported in weight $j$.
The conditions on $\gamma$ in Lemma~\ref{sha-lifts-residual}(d) follow from ($\gamma$), ($\delta$) and ($\varepsilon$), so $a\gamma$ is the coboundary of a matrix $Y$ of weight $j$ for some nonzero $a\in\mathcal O$.
The normalization ($\beta$) makes this coboundary vanish.
At $c$ it gives $Y=\operatorname{ad}r_0(c)Y=-J\,{}^tYJ$.
For $j=-2$ this forces $Y=0$.
For $j=0$ it forces $Y=\operatorname{diag}(y,0,-y)$, whose coboundary vanishes on $\ker\eta$ and at $c$, hence everywhere.
For $j=-1$ it forces $Y_{32}=-Y_{21}$, and at $g_1$ the $(2,1)$ entry of the normalization is $Y_{21}(\Psi(g_1)-\Psi(g_1)^{-1})=0$, so $Y=0$ because $\Psi(g_1)^2-1$ is a unit.
So $a\gamma=0$, hence $\gamma=0$ and $d=0$.

(e) Let $z$ be a cocycle as in the definition of $\text{Ш}_a$.
Let $Z(g)$ be the matrix with entries $z(g)$ at $(2,1)$ and $-z(g)$ at $(3,2)$, the two entries of weight $-1$; then $g\mapsto((1+\epsilon Z(g))M_g,\eta(g))$, with $r_0(g)=(M_g,\eta(g))$, is a homomorphism to $\mathcal G_3((\mathcal O/a\mathcal O)[\epsilon]/(\epsilon^2))$, since $\operatorname{ad}r_0$ acts on these entries through $\Psi^{-1}$ and $c$ exchanges them.
The local conditions on $z$ give ($\gamma$) to ($\varepsilon$), and adding to $z$ the coboundary of a constant gives ($\beta$).
By universality this unitary lift comes from a homomorphism $R\to(\mathcal O/a\mathcal O)[\epsilon]/(\epsilon^2)$, and $d_a$ sends the class of $z$ to the vector of $\epsilon$-coefficients of the images of the $x_i$.
It vanishes on coboundaries, satisfies the linearized relations because $I$ maps to $0$, and is injective because the class of $z$ is recovered from the entries $(2,1)$ of the lift on $\ker\eta$.
\end{proof}
\leanref{Development/AutomorphicArithmetic/AutomorphicVarieties/UnitaryThreeDeformations/WPresEx.lean}{UpperTheory.V.UnitaryThreeInertia.exists_adapted}
\leanref{Development/AutomorphicArithmetic/AutomorphicVarieties/UnitaryThreeDeformations/Origin.lean}{UpperTheory.V.UnitaryThreeInertia.hasOrigin}
\leanref{Development/AutomorphicArithmetic/AutomorphicVarieties/UnitaryThreeDeformations/CommAlgFunctorCoboundaryZeroMemCocyclesOne.lean}{UpperTheory.V.UnitaryThreeInertia.finrank_H1L_le}
\leanref{Development/AutomorphicArithmetic/AutomorphicVarieties/UnitaryThreeDeformations/ExistsRelations.lean}{UpperTheory.V.UnitaryThreeInertia.exists_relations}
\leanref{Development/AutomorphicArithmetic/AutomorphicVarieties/UnitaryThreeDeformations/TangentCoordEqZeroOfNonposWeight.lean}{UpperTheory.V.UnitaryThreeInertia.tangent_coord_eq_zero_of_nonpos_weight}
\leanref{Development/AutomorphicArithmetic/AutomorphicVarieties/UnitaryThreeDeformations/PrecompleteOfPowLe.lean}{UpperTheory.V.UnitaryThreeInertia.mprime_coord}

\begin{lemma}\label{sha-lifts-points}
In Lemma~\ref{sha-lifts-presentation}, suppose moreover that $R_0$ is finite over $\mathcal O$, and let $\xi$ be a point of $u$ with representation $r=r_\xi$.
\begin{enumerate}
\item The matrix parts of $r$ on $\ker\eta$ are upper triangular if and only if $\xi$ vanishes on $R_j$ for every $j<0$, and lower triangular if and only if $\xi$ vanishes on $R_j$ for every $j>0$.
\item $r$ is of characteristic zero.
\item If the matrix parts of $r$ on $\ker\eta$ are diagonal, then $r=r_0$.
\item If the matrix parts of $r$ on $\ker\eta$ are neither upper nor lower triangular, then $r$ is irreducible.
\item Let $\xi'$ be a point and $t\in\bar{\mathbf Q}_\ell^\times$. Then $r_{\xi'}=\tau(t)r\tau(t)^{-1}$ if and only if $\xi'(x)=t^j\xi(x)$ for every $j$ and $x\in R_j$.
\item If $r$ is rigid, then for every $\mathcal O$-linear $D:R^{\mathrm{gr}}\to\bar{\mathbf Q}_\ell$ with $D(xy)=\xi(x)D(y)+\xi(y)D(x)$ there is $a\in\bar{\mathbf Q}_\ell$ with $D(x)=aj\,\xi(x)$ for every $j$ and $x\in R_j$.
\item Let $r'$ be a homomorphism of characteristic zero. If $r$ is irreducible and the matrix parts of $r$ and $r'$ on $\ker\eta$ are conjugate by an element of $GL_3(\bar{\mathbf Q}_\ell)$, then $r'=\tau(t)r\tau(t)^{-1}$ for some $t\in\bar{\mathbf Q}_\ell^\times$.
\end{enumerate}
\end{lemma}
\begin{proof}
(a) Let $j<0$ and $x\in R_j$. By Lemma~\ref{sha-lifts-presentation}(a), $x=\pi(f)$ with $f$ of weight $j$, and every monomial of $f$ contains a variable of negative weight; so $x=\sum_ix_iy_i$, the sum over the $i$ with $w_i<0$, with $y_i\in R_{j-w_i}$.
The $x_i$ of negative weight are entries below the diagonal of matrix parts of $u(g)$ with $g\in\ker\eta$.
So $\xi$ vanishes on the $R_j$ with $j<0$ if and only if it vanishes on these entries, that is, if and only if $r$ is upper triangular on $\ker\eta$.
The other case is symmetric.

(b) An element of $R_0$ that is a unit of $R$ has its inverse in $R_0$, so $R_0$ is a local ring with maximal ideal $R_0\cap\mathfrak m_R$, finite over $\mathcal O$.
So, with $|\cdot|$ the absolute value of $\bar{\mathbf Q}_\ell$, $|\xi(x)|\le1$ for $x\in R_0$, and $|\xi(x)|<1$ if moreover $x\in\mathfrak m_R$.
Let $a_i=\xi(x_i)$, with $x_i$ and $w_i$ as in Lemma~\ref{sha-lifts-presentation}(a).
If $w_i=0$, then $x_i\in R_0\cap\mathfrak m_R$ and $|a_i|<1$; if $w_i>0>w_h$, then $x_i^{-w_h}x_h^{w_i}\in R_0\cap\mathfrak m_R$ and $|a_i|^{-w_h}|a_h|^{w_i}<1$.
By these inequalities the conditions $|t^{-w_i}a_i|<1$ on $\log|t|$, one for each $i$ with $w_i\ne0$, define a nonempty open interval, so some $t\in\bar{\mathbf Q}_\ell^\times$ satisfies $|t^{-w_i}a_i|<1$ for every $i$.
The $t^{-w_i}a_i$ lie in the maximal ideal of a discrete valuation ring $A\subset\bar{\mathbf Q}_\ell$ finite over $\mathcal O$; let $\Phi:P\to A$ be the continuous homomorphism of $\mathcal O$-algebras with $\Phi(X_i)=t^{-w_i}a_i$.
Then $\Phi(f)=t^{-j}\xi(\pi(f))$ for every power series $f$ of weight $j$.
For monomials this follows from multiplicativity.
By Dickson's lemma, a power series of weight $j$ is $\sum_mb_mm$, over a finite set of monomials $m$ of weight $j$, with $b_m$ of weight $0$; so it suffices to treat $j=0$.
A power series $f$ of weight $0$ without constant term is such a sum over a finite set of nonconstant monomials $m$ of weight $0$, for which $|\Phi(m)|\le\rho$ for some $\rho<1$; then $\Phi(f)-\xi(\pi(f))=\sum_m\big(\Phi(b_m)-\xi(\pi(b_m))\big)\Phi(m)$, and since $|\Phi|\le1$ on $P$ and $|\xi|\le1$ on $R_0$, induction gives $|\Phi(f)-\xi(\pi(f))|\le\rho^N$ for every $N$.
The ideal $I$ is generated by homogeneous power series (Lemma~\ref{sha-lifts-presentation}(c)), which $\Phi$ kills by this formula.
So $\Phi$ factors through a local homomorphism $\Xi:R\to A$ with $\Xi=t^{-j}\xi$ on $R_j$.
Since the entry $(a,b)$ of the matrix part of $u(g)$ lies in $R_{b-a}$, $\Xi\circ u=\tau(t)^{-1}r\tau(t)$; it is a unitary lift with values in $A$, so $r$ is of characteristic zero.

(c) Let $\tau(t)r\tau(t)^{-1}$ be a unitary lift with values in $A$.
Its diagonal entries on $\ker\eta$, divided by those of $r_0$, are characters of $\ker\eta$ with values in $1+\mathfrak m_A$.
By ($\gamma$), ($\delta$) and ($\varepsilon$) they are trivial on $I_p\cap\ker\eta$ for $p\in T_b$ and on $N_p$ for the other $p$, hence on these inertia groups and their conjugates, since $N_p$ has index prime to $\ell$ in $I_p$ by Lemma~\ref{sha-lifts-local}(d).
By Lemma~\ref{sha-lifts-residual}(c) they are trivial, so $\tau(t)r\tau(t)^{-1}=r_0$ on $\ker\eta$, and at $c$ by ($\beta$).
Since $\tau(t)$ commutes with $r_0$, $r=r_0$.

(d) Let $\zeta$ be the entry $(1,1)$ of the matrix part of $r(g_1)-r(cg_1^{-1}c)$; conjugation by $\tau(t)$ does not change it, and for the unitary lift $\tau(t)r\tau(t)^{-1}$ with values in $A$ it is congruent to $\Psi(g_1)-\Psi(g_1)^{-1}$ modulo $\mathfrak m_A$, so it is nonzero.
A homomorphism with multiplier $\eta$ that satisfies ($\beta$) and has $\zeta\ne0$ is either irreducible, or upper or lower triangular on $\ker\eta$, or block shaped: on $\ker\eta$ its entries of odd weight vanish and the corner block of the rows and columns $1$ and $3$ is irreducible.
Suppose that $r$ is block shaped, and let $r'=\tau(t)r\tau(t)^{-1}$ be a unitary lift with values in $A$.
The entry $(3,2)$ of $r'$ vanishes on $\ker\eta$, so modulo each power of a uniformizer of $A$ the entry $(3,1)$ of $r'$, divided by the largest possible power of the uniformizer, is a cocycle for $\bar\Psi^{-2}\eta$ that satisfies the local conditions of $H^1_L(\mathfrak g_{-2})$.
By Lemma~\ref{sha-lifts-residual}(e) it is a coboundary, and the normalization ($\beta$) makes it vanish; so the entry $(3,1)$ of $r$ vanishes.
Then the line of the first basis vector is stable under the corner block, which contradicts its irreducibility.
Hence $r$ is irreducible.

(e) The representation $\tau(t)r\tau(t)^{-1}$ has the entries $t^{b-a}r_{ab}$, and the entries of the matrix parts of the $u(g)$ generate $R^{\mathrm{gr}}$ by Lemma~\ref{sha-lifts-presentation}(a), with $u(g)_{ab}\in R_{b-a}$.

(f) Each $R_j$ is a finite $\mathcal O$-module, spanned by finitely many monomials in the $x_i$: by Dickson's lemma it is the sum of the $R_0m$ over finitely many monomials $m$, and $R_0$, which is finite over $\mathcal O$ and is topologically spanned by the monomials of weight $0$ (Lemma~\ref{sha-lifts-presentation}(a)), is spanned by finitely many of them.
% This is sring_eq_adjoin (part_zero_le and part_decomp, in Adjoin.lean).
The map $x\mapsto\xi(x)+\epsilon D(x)$ is a homomorphism of $\mathcal O$-algebras $R^{\mathrm{gr}}\to\bar{\mathbf Q}_\ell[\epsilon]/(\epsilon^2)$.
Applied to the entries of the matrix parts of the $u(g)$, which lie in $R^{\mathrm{gr}}$, it gives a homomorphism $g\mapsto((1+\epsilon y(g))M_g,\eta(g))$, where $r(g)=(M_g,\eta(g))$ and $y$ is a cocycle for $r$; $y$ is continuous because $D$ is $\mathcal O$-linear on each finite $\mathcal O$-module $R_j$.
The cocycle $y$ vanishes on $N_p$ for $p\in S\setminus T_b$, since $u$ is $1$ there.
The matrices $h$ of ($\gamma$) and ($\delta$) for $u$ can be taken with homogeneous entries, hence with entries in $R^{\mathrm{gr}}$, and their images give these conditions for $(1+\epsilon y)r$.
Since $r$ is rigid, $y$ is the coboundary of some $X\in M_3(\bar{\mathbf Q}_\ell)$.
By ($\beta$) for $u$, the matrix part of $r(c)$ is $1$, $y(c)=0$, and the entries $(1,2)$ and $(2,1)$ of the matrix part of $r(g_1)-r(cg_1^{-1}c)$, and of its coefficient of $\epsilon$, vanish.
So that matrix is $\operatorname{diag}(\zeta,0,-\zeta)$, with $\zeta=\xi(\zeta_u)$ for the corresponding entry $\zeta_u\in R_0$ of $u$; $\zeta_u$ is a unit of $R$, congruent to $\Psi(g_1)-\Psi(g_1)^{-1}$, and its inverse lies in $R_0$, so $\zeta\ne0$.
The vanishing of the coboundary of $X$ at $c$ gives $X=-J\,{}^tXJ$, and the vanishing of the entries $(1,2)$ and $(2,1)$ of the coefficient of $\epsilon$ gives $\zeta X_{12}=\zeta X_{21}=0$; since $2$ is invertible, $X=a\operatorname{diag}(1,0,-1)$ for some $a\in\bar{\mathbf Q}_\ell$.
For $g\in\ker\eta$ the coboundary of $X$ gives $y(g)M_g=XM_g-M_gX$, whose entry $(a',b')$ is $a(b'-a')$ times that of $M_g$; so $D(x)=aj\,\xi(x)$ for every entry $x\in R_j$ of the matrix parts of the $u(g)$ with $g\in\ker\eta$.
Both $D$ and the map equal to $x\mapsto aj\,\xi(x)$ on each $R_j$ are derivations at $\xi$, and these entries generate $R^{\mathrm{gr}}$ as an $\mathcal O$-algebra, since the $x_i$ are among them up to constants; so $D(x)=aj\,\xi(x)$ for every $j$ and $x\in R_j$.

(g) A homomorphism $r$ of characteristic zero satisfies ($\beta$), since $\tau(t)r\tau(t)^{-1}$ does for some $t$. So, if $r(g_1)=(M_1,1)$, the matrix part of $r(cg_1^{-1}c)$ is $J\,{}^tM_1J$, and that of $r(g_1)-r(cg_1^{-1}c)$ is $\operatorname{diag}(\zeta,0,-\zeta)$, with $\zeta=\zeta(r)$ the entry of (d).
Conjugation by $\tau(t)$ does not change $\zeta$, so $\zeta(r)$ and $\zeta(r')$ are congruent to $\Psi(g_1)-\Psi(g_1)^{-1}$, which is a unit; since $2$ is a unit, $\zeta(r)\ne0$, $\zeta(r')\ne0$ and $\zeta(r')\ne-\zeta(r)$.
Let $M\in GL_3(\bar{\mathbf Q}_\ell)$ conjugate the matrix parts of $r$ to those of $r'$ on $\ker\eta$.
Since $g_1$ and $cg_1^{-1}c$ lie in $\ker\eta$, $M\operatorname{diag}(\zeta,0,-\zeta)=\operatorname{diag}(\zeta',0,-\zeta')M$ with $\zeta'=\zeta(r')$, and the three conditions on $\zeta$ and $\zeta'$ force $M$ to be diagonal, $M=\operatorname{diag}(m_1,m_2,m_3)$.
For $g\in\ker\eta$, comparing $r$ and $r'$ at $g$ and at $cg^{-1}c$ gives $m_ar_{ab}=r'_{ab}m_b$ and $m_{4-b}r_{ab}=r'_{ab}m_{4-a}$ for the entries of the matrix parts.
If an entry of weight $\pm1$ of $r$ is nonzero somewhere on $\ker\eta$, these give $m_1m_3=m_2^2$, and $t=m_1/m_2$ satisfies $m_a/m_b=t^{b-a}$; otherwise only the diagonal and the entries $(1,3)$ and $(3,1)$ of $r$ can be nonzero, and a square root $t$ of $m_1/m_3$ will do.
In both cases $r'_{ab}=t^{b-a}r_{ab}$ on $\ker\eta$, that is, $r'=\tau(t)r\tau(t)^{-1}$ on $\ker\eta$.
At $c$ both have matrix part $1$, and $\tau(t)$ commutes with $(1,-1)$ because $J\,{}^t\tau(t)J=\tau(t)^{-1}$; so $r'=\tau(t)r\tau(t)^{-1}$.
The irreducibility of $r$ is not used.
\end{proof}
\leanref{Development/AutomorphicArithmetic/AutomorphicVarieties/UnitaryThreeDeformations/PtSign.lean}{UpperTheory.V.UnitaryThreeInertia.isPtHom_sign}
\leanref{Development/AutomorphicArithmetic/AutomorphicVarieties/UnitaryThreeDeformations/IsCharZeroPtOfIsPtHom.lean}{UpperTheory.V.UnitaryThreeInertia.isCharZeroPt_of_isPtHom}
\leanref{Development/AutomorphicArithmetic/AutomorphicVarieties/UnitaryThreeDeformations/UnitEqOneOfPowEqOneOfResidueOne.lean}{UpperTheory.V.UnitaryThreeInertia.eq_rz_of_diag}
\leanref{Development/AutomorphicArithmetic/AutomorphicVarieties/UnitaryThreeDeformations/TypeIOfTwoSigned.lean}{UpperTheory.V.UnitaryThreeInertia.typeI_of_twoSigned}
\leanref{Development/AutomorphicArithmetic/AutomorphicVarieties/UnitaryThreeDeformations/IsPtHomAct.lean}{UpperTheory.V.UnitaryThreeInertia.isPtHom_act}
\leanref{Development/AutomorphicArithmetic/AutomorphicVarieties/UnitaryThreeDeformations/ContinuousCompOfFilteredContinuity.lean}{UpperTheory.V.UnitaryThreeInertia.tangent_eq_orbit}
\leanref{Development/AutomorphicArithmetic/AutomorphicVarieties/UnitaryThreeDeformations/NormTwoEqOne.lean}{UpperTheory.V.UnitaryThreeInertia.conj_of_matIso}

\begin{lemma}\label{sha-lifts-coordinates}
Let $\ell$ be a prime, let $\mathcal O$ be as in Definition~\ref{sha-character-parts}, and let $v=e_{\mathcal O}v_\ell$ on $\bar{\mathbf Q}_\ell$, where $v_\ell(\ell)=1$.
Let $P=\mathcal O[[X_1,\dots,X_n]]$, graded by giving $X_i$ a weight $w_i\in\{-2,\dots,2\}$, and let $I\subset P$ be generated by power series $f_{j,m}$, $1\le m\le b_j$, with $f_{j,m}$ homogeneous of weight $j$ and $f_{j,m}(0)=0$, where $b_j+\delta(j)\le\#\{i:w_i=j\}$ for every $j$ (Definition~\ref{sha-lifts-selmer}).
Let $B=P/I$, let $B_j$ be the image of the power series of weight $j$, so that $B^{\mathrm{gr}}=\bigoplus_jB_j$ is a graded $\mathcal O$-algebra, and let $o:B\to\mathcal O$ be the homomorphism with $o(X_i)=0$.
Two-signed homomorphisms $B^{\mathrm{gr}}\to\bar{\mathbf Q}_\ell$ and their orbits are as in Definition~\ref{sha-algebra-graded}.
Suppose:
\begin{enumerate}
\item for $j\le0$, every $d\in\mathcal O^n$ with $d_i=0$ unless $w_i=j$ and $\sum_i\frac{\partial f}{\partial X_i}(0)d_i=0$ for all $f\in I$ is $0$;
\item $B_0$ is finite over $\mathcal O$;
\item a homomorphism $B^{\mathrm{gr}}\to\bar{\mathbf Q}_\ell$ that vanishes on every $B_j$ with $j\ne0$ is $o$;
\item for every two-signed $\xi$ and every $\mathcal O$-linear $D:B^{\mathrm{gr}}\to\bar{\mathbf Q}_\ell$ with $D(xy)=\xi(x)D(y)+\xi(y)D(x)$ there is $a\in\bar{\mathbf Q}_\ell$ with $D(x)=aj\xi(x)$ for $x\in B_j$;
\item $\alpha\in B_0$ satisfies $o(\alpha)=0$ and $\xi(\alpha)\ne0$ for every two-signed $\xi$.
\end{enumerate}
Then there are two-signed $\xi_1,\dots,\xi_s$ in pairwise distinct orbits and an integer $V$ with $0\le V\le2\sum_iv(\xi_i(\alpha))$ such that the following holds: for every nonzero $a\in\mathcal O$ and every $\mathcal O$-module $M$ with an injective additive map $d:M\to(\mathcal O/a\mathcal O)^n$ such that $d(\mu)_i=0$ unless $w_i=-1$ and $\sum_i\frac{\partial f}{\partial X_i}(0)d(\mu)_i=0$ for all $f\in I$ and $\mu\in M$, the length of $M$ is at most $V$.
\end{lemma}
\begin{proof}
If no $w_i$ is $-1$, then every such $M$ is $0$, and $s=V=0$ will do.
Otherwise let $n_{-1}\ge1$ be the number of $i$ with $w_i=-1$.
Complete the generators of $I$ by zeros to families $F$ of $n_{-1}+1$ power series of weight $-1$, $G$ of $\#\{i:w_i=-2\}$ of weight $-2$, $E$ and $E'$ of $\#\{i:w_i=j\}-1$ of weight $j$ for $j=1,2$, and $Z$ of $\#\{i:w_i=0\}$ of weight $0$; this is possible because $b_j+\delta(j)$ is at most the number of variables of weight $j$.
The graded algebra $P^{\mathrm{gr}}=\bigoplus_jP_j$ of the power series that are finite sums of homogeneous ones, with the evaluation at $0$, satisfies the hypotheses of Lemma~\ref{sha-algebra-cotangent} on $A$ and $x$, and the ideal generated by these families is $I\cap P^{\mathrm{gr}}$.
The conditions of that lemma hold: the first with the inclusion of $P^{\mathrm{gr}}$ into the power series over the fraction field of $\mathcal O$, and by (a), which says that the families of nonpositive weight span the corresponding cotangent lattices up to finite index; the second by (b); the third by (c); the fourth by (d), since two derivations that are multiples of $x\mapsto j\xi(x)$ are linearly dependent; and the fifth by (e).
So $\lambda_{-1}(F)$ is finite, and there is a finite set $T$ of two-signed homomorphisms meeting every orbit, with $\lambda_{-1}(F)\le2\sum_{\xi\in T}v(\xi(\alpha))$.
Keep one element of $T$ in each orbit; the set still meets every orbit, so the bound holds for it, and its elements are the $\xi_i$.
Put $V=\lambda_{-1}(F)$.
For $M$ and $d$ as in the statement, the map that sends $\mu$ to $f\mapsto\sum_i\frac{\partial f}{\partial X_i}(0)d(\mu)_i$ is injective, and its values are reductions modulo $a$ of derivations at the evaluation at $0$, of weight $-1$, that vanish on $F$; so $\operatorname{len}M\le V$ by Lemma~\ref{sha-algebra-derivations}.
\end{proof}
\leanref{Development/AutomorphicArithmetic/AutomorphicVarieties/UnitaryThreeDeformations/AbsCoord.lean}{UpperTheory.V.UnitaryThreeInertia.Abs.coord_count}

\begin{proposition}\label{sha-lifts-abstract-count}
In Lemma~\ref{sha-lifts-points}, let $\varphi\in\Gamma$, and let $\mathcal P$ be a property of homomorphisms $\Gamma\to\mathcal G_3(\bar{\mathbf Q}_\ell)$ such that:
the representation of every point of $u$ that is irreducible has $\mathcal P$; and every irreducible homomorphism of characteristic zero with $\mathcal P$ is rigid and has $\operatorname{tr}r(\varphi)\ne\operatorname{tr}r_0(\varphi)$, traces of matrix parts.
Then there are irreducible homomorphisms $r_1,\dots,r_s$ of characteristic zero with $\mathcal P$, whose matrix parts on $\ker\eta$ are pairwise not conjugate under $GL_3(\bar{\mathbf Q}_\ell)$, and an integer $V$ with $\operatorname{len}_{\mathcal O}\text{Ш}_a\le V$ for every nonzero $a\in\mathcal O$ and
\[
 V\le2e_{\mathcal O}\sum_{i=1}^{s}v_\ell\big(\operatorname{tr}r_i(\varphi)-\operatorname{tr}r_0(\varphi)\big).
\]
\end{proposition}
\begin{proof}
Let $x_i$, $w_i$, $P$, $I$ and $o$ be as in Lemma~\ref{sha-lifts-presentation}(a), and $f_{j,m}$ as in (c); identify $R$ with $P/I$.
By Lemma~\ref{sha-lifts-presentation}(b) and Lemma~\ref{sha-lifts-residual}(b), $\dim\mathrm{Ob}_j+\delta(j)\le\dim H^1_L(\mathfrak g_j)$ is at most the number of $i$ with $w_i=j$.
Let $\alpha=\operatorname{tr}u(\varphi)-\operatorname{tr}r_0(\varphi)$; the diagonal entries have weight $0$, so $\alpha\in R_0$, and $o(\alpha)=0$ because $o\circ u=r_0$.
For a point $\xi$ with representation $r$, $\xi(\alpha)=\operatorname{tr}r(\varphi)-\operatorname{tr}r_0(\varphi)$.
If $\xi$ is two-signed, then $r$ is of characteristic zero and irreducible by Lemma~\ref{sha-lifts-points}(a), (b) and (d), so it has $\mathcal P$, and it is rigid with $\xi(\alpha)\ne0$.

We check the hypotheses of Lemma~\ref{sha-lifts-coordinates}.
(a) is Lemma~\ref{sha-lifts-presentation}(d), and (b) is the finiteness of $R_0$.
For (c), a point that vanishes on every $R_j$ with $j\ne0$ has upper and lower triangular matrix parts on $\ker\eta$ by Lemma~\ref{sha-lifts-points}(a), so its representation is $r_0$ by Lemma~\ref{sha-lifts-points}(c); a point is determined by its representation, by Lemma~\ref{sha-lifts-points}(e) with $t=1$, so it is $o$.
(d) follows from Lemma~\ref{sha-lifts-points}(f), since the representation of a two-signed point is rigid, and (e) was shown above.
So Lemma~\ref{sha-lifts-coordinates} gives two-signed points $\xi_1,\dots,\xi_s$ and $V$.
Let $r_i$ be the representation of $\xi_i$; it is irreducible, of characteristic zero, and has $\mathcal P$.
If the matrix parts of $r_i$ and $r_{i'}$ on $\ker\eta$ are conjugate, then $r_{i'}=\tau(t)r_i\tau(t)^{-1}$ for some $t$ by Lemma~\ref{sha-lifts-points}(g), so $\xi_{i'}$ is in the orbit of $\xi_i$ by Lemma~\ref{sha-lifts-points}(e), and $i=i'$.
For a nonzero $a\in\mathcal O$, Lemma~\ref{sha-lifts-presentation}(e) gives the map $d_a$ on $M=\text{Ш}_a$ required in Lemma~\ref{sha-lifts-coordinates}, so $\operatorname{len}\text{Ш}_a\le V$.
Finally $V\le2\sum_iv(\xi_i(\alpha))=2e_{\mathcal O}\sum_iv_\ell(\operatorname{tr}r_i(\varphi)-\operatorname{tr}r_0(\varphi))$.
\end{proof}
\leanref{Development/AutomorphicArithmetic/AutomorphicVarieties/UnitaryThreeDeformations/AbstractCount.lean}{UpperTheory.V.UnitaryThreeInertia.abstract_count}

\begin{lemma}\label{sha-lifts-comparison}
In Definition~\ref{sha-lifts-problem}, suppose that the data satisfy the dyadic condition for some $o\ge1$, and let $V\ge0$ be an integer with $\operatorname{len}_{\mathcal O}\text{Ш}_a\le V$ for every nonzero $a\in\mathcal O$.
Then $V'_\ell(e,\psi)\le V$ (Definition~\ref{sha-character-parts}).
\end{lemma}
\begin{proof}
Let $M=\mathcal O\otimes\mathbf Q_\ell/\mathbf Z_\ell$ with $\sigma$ acting by $\Psi(\sigma)^{-1}$.
The classes defining $V'_\ell(e,\psi)$ are unramified outside $S$, so they are inflated from the subgroup $\text{Ш}$ of $H^1(\Gamma,M)$ of classes that vanish on $D_p$ for every $p\in S$.
Every element of $\text{Ш}$ is killed by a power of $\ell$, so $\text{Ш}$ is the increasing union of the submodules $\text{Ш}[\ell^m]$ of elements killed by $\ell^m$, and it suffices to bound their lengths by $V$.
Let $a=\ell^m$, so that $M[a]\cong\mathcal O/a\mathcal O(\Psi^{-1})$, and let $x\in\text{Ш}[a]$ be the class of a cocycle $z$.
Then $az$ is the coboundary of some element, which is $am_1$ since $M$ is divisible, and $z'=z-dm_1$ is a cocycle with values in $M[a]$ that represents $x$.
For $p\in S$, $z'$ is the coboundary on $D_p$ of some $m\in M$, and $a(\Psi(g)^{-1}-1)m=0$ for $g\in D_p$.
If $p\in S\setminus T_b$, then $\Psi$ is trivial on $N_p$ by Lemma~\ref{sha-lifts-local}(c), so $z'$ vanishes on $N_p$.
If $p\in T_b$, then $\Psi(x_p)-1$ is a unit by Lemma~\ref{sha-lifts-local}(b), so $am=0$, and $z'$ is a coboundary on $D_p$ with values in $M[a]$.
So $x$ comes from $\text{Ш}_a$, and $\operatorname{len}\text{Ш}[a]\le\operatorname{len}\text{Ш}_a\le V$.
\end{proof}
\leanref{Development/AutomorphicArithmetic/AutomorphicVarieties/UnitaryThreeDeformations/HostSha.lean}{UpperTheory.V.UnitaryThreeInertia.vlen_le}

\begin{lemma}\label{sha-lifts-automorphy-conditions}
In Definition~\ref{sha-lifts-problem}, suppose that the data satisfy the dyadic condition for some $o\ge1$.
View $\eta$, $\Psi$ and $\bar\Psi$ as characters of $\operatorname{Gal}(\bar{\mathbf Q}/\mathbf Q)$ through its quotient $\Gamma$, let $\omega:\operatorname{Gal}(\bar{\mathbf Q}/\mathbf Q)\to k^\times$ be the reduction of $\chi_{\mathrm{cyc}}$, let $c'\in\operatorname{Gal}(\bar{\mathbf Q}/\mathbf Q)$ be a complex conjugation, let $B$ be a place of $\bar{\mathbf Q}$ above $\ell$ with decomposition group $D_B$, and put $M=\ell N$.
Let $\kappa_b$ be the character $\kappa$ with its component at $\ell$ removed, and let $H_b\subset\operatorname{Gal}(\bar{\mathbf Q}/\mathbf Q)$ be its kernel.
\begin{enumerate}
\item The field $k$ is finite, and $2\ell$ divides $M$.
\item There are an integer $D\ge1$ whose prime factors divide $M$ and an element $\delta\in\bar{\mathbf Q}$ with $\delta^2=-D$ and $g(\delta)=\eta(g)\delta$ for every $g$.
\item There is $a\ge0$ such that $\eta$ and $\bar\Psi$ are trivial on $\operatorname{Gal}(\bar{\mathbf Q}/\mathbf Q(\mu_{M^a}))$.
\item The character $\eta$ is trivial on the inertia groups at the places above $\ell$, $\bar\Psi(c')=1$, and there is $j\ge0$ with $\bar\Psi=\omega^j$ on these inertia groups.
\item The subgroup $\ker\omega\cap\ker\bar\Psi$ is not contained in $\ker\eta$.
\item There is $s\in D_B$ with $\eta(s)=1$ and $\bar\Psi(s)^2\ne1$.
\item The subgroup $H_b$ is open and normal, with abelian quotient of $\ell$-power order, and contains the inertia groups at the places above $\ell$; and $\Psi(s)=1$ for every $s\in H_b$ with $\bar\Psi(s)=1$ that lies in an inertia group at a place not above $\ell$.
\end{enumerate}
\end{lemma}
\begin{proof}
(a) The ring $\mathcal O$ is finite over $\mathbf Z_\ell$, and $N$ is even.

(b) Take $D=1$ and $\delta=i$: $\eta$ is the character of $\mathbf Q(i)$.

(c) Take $a=2$.
Since $4$ divides $M^2$, $\eta$ is trivial on $\operatorname{Gal}(\bar{\mathbf Q}/\mathbf Q(\mu_{M^2}))$.
The characters $\omega$ and $\psi$ factor through $\operatorname{Gal}(\mathbf Q(\mu_\ell)/\mathbf Q)$ and $\operatorname{Gal}(\mathbf Q(\mu_N)/\mathbf Q)$, and $\ell$ and $N$ divide $M^2$; so $\bar\Psi=\omega^e\bar\psi$ is trivial there too.

(d) The field $\mathbf Q(i)$ is unramified at the odd prime $\ell$.
Since $\chi_{\mathrm{cyc}}(c')=-1$ and $\psi(c')=\psi(-1)=(-1)^e$, $\Psi(c')=1$.
An element $s$ of an inertia group above $\ell$ acts on the $N$-th roots of unity through the component at $\ell$ of $(\mathbf Z/N)^\times$, so $\psi(s)$ is a value of the component $\psi_\ell$ of $\psi$ at $\ell$, which is the product of the components of $\kappa$ and $\theta_0$ at $\ell$.
The first has $\ell$-power order, so its values are congruent to $1$.
The second has order prime to $\ell$, so it is trivial on the kernel of $(\mathbf Z/\ell^m)^\times\to(\mathbf Z/\ell)^\times$, an $\ell$-group, and its reduction is the reduction map $(\mathbf Z/\ell)^\times\to k^\times$ raised to some power $b\ge0$.
Hence $\bar\Psi(s)=\omega(s)^{e+b}$.

(e) There is $u\in(\mathbf Z/N)^\times$ with $u\equiv-1\pmod4$ and $\psi_2(u)=1$.
If $\psi_2(-1)=1$, take $u=-1$.
Otherwise, since the conductor $2^{c_2}$ of $\psi_2$ is at least $16$, $\psi_2$ is nontrivial on the elements congruent to $1$ modulo $4$; as $\psi_2$ has $2$-power order, some such element $v$ has $\psi_2(v)=-1$, and $u=-v$ will do.
Let $s$ be an element of an inertia group above $2$ that acts on the $N$-th roots of unity through the component of $u$ at $2$.
Then $\omega(s)=1$, since $\ell$ is odd, and $\psi(s)=\psi_2(u)=1$, so $s\in\ker\omega\cap\ker\bar\Psi$; but $\eta(s)=-1$, since $u\equiv-1\pmod4$.

(f) Lift the element $x_\ell$ of Lemma~\ref{sha-lifts-local}(b) to the decomposition group of a place above $\ell$, and conjugate it into $D_B$.

(g) The character $\kappa_b$ has $\ell$-power order and factors through $(\mathbf Z/N)^\times$, which gives the first assertion; its component at $\ell$ is trivial, so $H_b$ contains the inertia groups above $\ell$.
Let $s\in H_b$ lie in an inertia group at a place above a prime $p\ne\ell$, with $\bar\Psi(s)=1$.
Then $s$ fixes the $\ell$-power roots of unity, so $\chi_{\mathrm{cyc}}(s)=1$, $\Psi(s)=\psi(s)$, and the component of $\kappa$ at $\ell$ is trivial at $s$; with $\kappa_b(s)=1$ this gives $\kappa(s)=1$.
So $\psi(s)=\theta_0(s)$ is a root of unity of order prime to $\ell$ whose reduction is $\bar\Psi(s)=1$, and $\psi(s)=1$.
\end{proof}
\leanref{Development/AutomorphicArithmetic/AutomorphicVarieties/FrobeniusSmoothVectors/GenHost.lean}{UpperTheory.V.Fr.gen_host}
\leanref{Development/AutomorphicArithmetic/AutomorphicVarieties/FrobeniusSmoothVectors/FrobeniusSplitPrimeDatum.lean}{UpperTheory.V.Fr.Gen}

\begin{proposition}\label{sha-lifts-automorphy-form}
Assume that Galois representations exist as in Theorem~\ref{sha-residually-reducible-automorphy}.
In Definition~\ref{sha-lifts-problem}, suppose that the data satisfy the dyadic condition for some $o\ge1$, let $K=\mathbf Q(i)$, let $\iota:\bar{\mathbf Q}_\ell\cong\mathbf C$, let $c'\in\operatorname{Gal}(\bar{\mathbf Q}/\mathbf Q)$ be a complex conjugation and let $B$ be a place of $\bar{\mathbf Q}$ above $\ell$.
View $\eta$ and $\Psi$ as characters of $\operatorname{Gal}(\bar{\mathbf Q}/\mathbf Q)$, and define ordinary points relative to $g_1\in\operatorname{Gal}(\bar{\mathbf Q}/\mathbf Q)$ and $B$ as in Definition~\ref{sha-ordinary-deformation}(c), for these characters and with $c'$ in place of $c$.
Then there is $g_1\in\operatorname{Gal}(\bar{\mathbf Q}/K)$ with $\Psi(g_1)^2-1\in\mathcal O^\times$ such that the following hold for the ordinary points relative to $g_1$ and $B$.
\begin{enumerate}
\item (Finiteness.) There is a subgroup $U$ of finite index such that for every complete discrete valuation ring $A$ over $\mathcal O$ with finite residue field in which $\mathfrak m_{\mathcal O}$ maps to $0$, and every ordinary point $r$ with values in $A$, the characteristic polynomial of $r(g)$ equals that of $r_0(\eta,\Psi)(g)$ for all $g\in U\cap\ker\eta$.
\item (Stable points.) Let $A\subset\bar{\mathbf Q}_\ell$ be a discrete valuation ring finite over $\mathcal O$, let $r$ be an ordinary point with values in $A$ whose restriction to every open subgroup of $\ker\eta$ is irreducible over $\bar{\mathbf Q}_\ell$, and let $\rho:\operatorname{Gal}(\bar{\mathbf Q}/K)\to GL_3(\bar{\mathbf Q}_\ell)$ be its restriction. Then $\rho\otimes\chi_{\mathrm{cyc}}^{-1}$ is attached through $\iota$ to a regular algebraic cuspidal automorphic representation $\pi$ of $GL_3(\mathbf A_K)$ with $\pi^\vee\cong\pi^c$, and it is windowed.
\end{enumerate}
\end{proposition}
\begin{proof}
By Lemma~\ref{sha-lifts-automorphy-conditions}, whose conditions (a) to (g) are those of Definition~\ref{sha-descent-datum}, $(\eta,\Psi,c',B,\ell N,H_b)$ is a descent datum.
Proposition~\ref{sha-descent-construction} gives an abelian CM field $L$, an element $g_1\in\operatorname{Gal}(\bar{\mathbf Q}/L)$ with $\bar\Psi(g_1)^2\ne1$, a link of the datum to $L$ with base element $g_1$ that satisfies the residual conditions with a constant $C_0$ and is large for $C_0$, and a seed datum for $L$.
The quadratic field of the seed datum is $K$, because $\ker\eta=\operatorname{Gal}(\bar{\mathbf Q}/K)$; so $L\supset K$, $g_1\in\operatorname{Gal}(\bar{\mathbf Q}/K)$, and $\Psi(g_1)^2-1\in\mathcal O^\times$.
By Lemma~\ref{sha-descent-points}(1), the ordinary points relative to $g_1$ and $B$ give linked points of rank three.
The rest of the proof of the theorem uses of the Eisenstein datum only these objects and the facts that $e\ge1$, that $N$ is even, that $\Psi(c')=1$ and that $\eta$ is unramified at the primes not dividing $M$.
Here $\Psi(c')=1$ by Lemma~\ref{sha-lifts-automorphy-conditions}(d), and $\eta$ is ramified only at $2$.
So (i) follows from the first part of proposition (c) of that proof.
For (ii), proposition (c) makes the restriction of $\rho\otimes\chi_{\mathrm{cyc}}^{-1}$ to $L$ automorphic, the descent of conjugate self-dual automorphy along the solvable extension $L/K$ gives $\pi$, and the windowed conditions follow from the ordinary local conditions at $\ell$ and from proposition (d), as in that proof.
% The construction is Fr.host_ex, which applies Fr.seed_cl and Fr.frame_ex (Proposition sha-descent-construction) to Fr.gen_host (Lemma sha-lifts-automorphy-conditions).
% The rest is Z.Win.Dn.hostEnd_of_loc3_frame_linkedEnd_downCsdK, the proof of Z.Win.Dn.gdEnd_of_frame_linkedEnd_downCsdK_descent with this link: marks_of_link, finEnd_of_finPt (from the first part of LinkedEndEAlt), autEnd3_alt with DownCsdKAlt, Dn.Loc3 and Dn.LocalConditionRankThreeEvenLevel.
\end{proof}
\leanref{Development/AutomorphicArithmetic/AutomorphicVarieties/ModularityArgument/ThresholdArguments/DenseThresholdCases/AutEnd3AltAlt.lean}{UpperTheory.V.Pr.HostEndAlt}
\leanref{Development/AutomorphicArithmetic/AutomorphicVarieties/ModularityArgument/ThresholdArguments/DenseThresholdCases/DnHost.lean}{UpperTheory.V.Z.Win.Dn.hostEnd_of_loc3_frame_linkedEnd_downCsdK}
\leanref{Development/AutomorphicArithmetic/AutomorphicVarieties/FrobeniusSmoothVectors/HostEx.lean}{UpperTheory.V.Fr.host_ex}
\leanref{Development/AutomorphicArithmetic/AutomorphicVarieties/CurveDVRSpecialization/HostFrameEx.lean}{UpperTheory.V.Pr.HostFrameEx}

\begin{lemma}\label{sha-lifts-triangularity}
Let $G$ be a group, let $k$ be a finite field of characteristic other than $2$, let $\eta:G\to\{\pm1\}$ and $\bar\Psi:G\to k^\times$ be homomorphisms, and let $c,g_1\in G$ with $c^2=1$, $\eta(c)=-1$, $\bar\Psi(c)=1$, $\eta(g_1)=1$ and $\bar\Psi(g_1)^2\ne1$.
Let $A$ be a local domain with a homomorphism $k\to A$, and let $r:G\to\mathcal G_3(A)$ be a homomorphism with multiplier $\eta$ whose matrix parts are congruent modulo $\mathfrak m_A$ to those of $\bar r_0=r_0(\eta,\bar\Psi)$, such that the matrix part of $r(c)$ is $1$ and the $(1,2)$ and $(2,1)$ entries of the matrix part of $r(g_1)-r(cg_1^{-1}c)$ vanish.
Suppose that for some subgroup $U$ of finite index the characteristic polynomial of $r(g)$ equals that of $\bar r_0(g)$ for every $g\in U\cap\ker\eta$, characteristic polynomials being those of matrix parts.
Then the matrix parts of $r$ on $\ker\eta$ are all upper triangular or all lower triangular, and for every $g\in G$ the diagonal entries of the matrix parts of $r(g)$ and $\bar r_0(g)$ agree.
\end{lemma}
\begin{proof}
An element $a\in A$ that is algebraic over $k$ and congruent modulo $\mathfrak m_A$ to some $b\in k$ equals $b$: the ring $k[a]$ is a domain finite over $k$, hence a field, so it meets $\mathfrak m_A$ only in $0$, and $a-b\in k[a]\cap\mathfrak m_A$.

Let $g\in\ker\eta$, let $X$ be the matrix part of $r(g)$, with characteristic polynomial $P$, and choose $n\ge1$ with $g^n\in U$.
The matrix part of $r(g^n)$ is $X^n$, so the characteristic polynomial $P_n$ of $X^n$ is that of $\bar r_0(g^n)$, with coefficients in $k$; and $P(T)$ divides $P_n(T^n)$.
So the coefficients of $P$ are integral over $k$, and they are congruent to those of the characteristic polynomial of $\bar r_0(g)$; hence $P$ is the characteristic polynomial of $\bar r_0(g)$.
If a diagonal entry $X_{aa}$ is a root of $P$, it is algebraic over $k$ and congruent to the diagonal entry of $\bar r_0(g)$, so it equals it.

The entry $\zeta$ at $(1,1)$ of the matrix part of $r(g_1)-r(cg_1^{-1}c)$ is congruent to $\bar\Psi(g_1)-\bar\Psi(g_1)^{-1}\ne0$, so $\zeta\ne0$.
The shape argument recalled in the proof of Lemma~\ref{sha-lifts-points}(d) works over any field of characteristic other than $2$; over the fraction field of $A$ it shows that the matrix parts of $r$ on $\ker\eta$ are irreducible, or all upper triangular, or all lower triangular, or block shaped: their entries of odd weight vanish and the corner block of the rows and columns $1$ and $3$ is irreducible.
They are not irreducible: their characteristic polynomials are those of the sum of the characters $\bar\Psi$, $1$ and $\bar\Psi^{-1}$ of $\ker\eta$, so an irreducible representation with these characteristic polynomials would be isomorphic to that sum by the Brauer-Nesbitt theorem \cite[(30.16)]{CR62}.
If they were block shaped, the middle diagonal entry of $X$ would be a root of $P$, hence equal to $1$, and the corner block would have characteristic polynomial $(T-\bar\Psi(g))(T-\bar\Psi(g)^{-1})$ for every $g\in\ker\eta$; by the same theorem it would be reducible.
So they are all upper triangular or all lower triangular.
Then every diagonal entry of $X$ is a root of $P$, so the diagonal entries of $r(g)$ and $\bar r_0(g)$ agree for $g\in\ker\eta$.
Every other element of $G$ is $gc$ with $g\in\ker\eta$, and since the matrix parts of $r(c)$ and $\bar r_0(c)$ are $1$, the matrix parts of $r(gc)$ and $\bar r_0(gc)$ are those of $r(g)$ and $\bar r_0(g)$.
\end{proof}
\leanref{Development/AutomorphicArithmetic/AutomorphicVarieties/CurveDVRSpecialization/ConstOfChar.lean}{UpperTheory.V.Pr.constOfChar}
\leanref{Development/AutomorphicArithmetic/AutomorphicVarieties/CurveDVRSpecialization/IsUnitOrIsUnitOfAddOne.lean}{UpperTheory.V.Pr.ConstOfChar}
\leanref{Development/AutomorphicArithmetic/AutomorphicVarieties/UnitaryThreeGroups/ExistsCoordOfNotMatIrr.lean}{UpperTheory.V.U3.shape_cases}

\begin{lemma}\label{sha-lifts-finiteness-diagonal}
In Definition~\ref{sha-lifts-problem}, suppose that the data satisfy the dyadic condition for some $o\ge1$, and let $g_1$ be as in (f).
Let $A$ be a noetherian local domain with a local homomorphism $\mathcal O\to A$, and let $r:\Gamma\to\mathcal G_3(A)$ satisfy conditions ($\alpha$) to ($\varepsilon$) of Definition~\ref{sha-lifts-problem}(f).
If the matrix parts of $r$ on $\ker\eta$ are all upper triangular or all lower triangular, then for every $g\in\Gamma$ the diagonal entries of the matrix parts of $r(g)$ and $r_0(g)$ agree.
\end{lemma}
\begin{proof}
Write $M_g$ for the matrix part of $r(g)$ and $(w_1,w_2,w_3)=(1,0,-1)$, so that the diagonal entries of $r_0(g)$ are the $\Psi(g)^{w_a}$.
Since the $M_g$ with $g\in\ker\eta$ are triangular, $d_a(g)=(M_g)_{aa}\Psi(g)^{-w_a}$ defines a homomorphism $d_a:\ker\eta\to A^\times$, congruent to $1$ modulo $\mathfrak m_A$ and continuous for the $\mathfrak m_A$-adic topology.

First, $d_a$ is trivial on $I_p$ for $p\in T_b$, a subgroup of $\ker\eta$ by Lemma~\ref{sha-lifts-local}(a).
Let $h$ be as in ($\gamma$) or ($\delta$), and $V_g=h^{-1}M_gh$ for $g\in D_p\cap\ker\eta$; it is upper triangular, with diagonal entries $\Psi(g)^{w_b}$ if $g\in I_p$.
Let $g\in D_p\cap\ker\eta$ with $\Psi(g)^2-1\in\mathcal O^\times$.
The residues of the diagonal entries of $M_g$ and of $V_g$ are the $\bar\Psi(g)^{w_b}$, which are pairwise distinct, and
\[
 \prod_b\big((M_g)_{aa}-(V_g)_{bb}\big)=\det\big((M_g)_{aa}-V_g\big)=\det\big((M_g)_{aa}-M_g\big)=0,
\]
because $M_g$ is triangular; the factors with $b\ne a$ are units, so $(M_g)_{aa}=(V_g)_{aa}$.
Let $s\in I_p$.
If $\Psi(s)^2-1$ is a unit, this gives $(M_s)_{aa}=\Psi(s)^{w_a}$.
Otherwise $\bar\Psi(s)^2=1$, and it applies to the element $x_p$ of Lemma~\ref{sha-lifts-local}(b) and to $x_ps$; since the diagonal entries of $M$ and of $V$ are multiplicative on $D_p\cap\ker\eta$, again $(M_s)_{aa}=(V_s)_{aa}=\Psi(s)^{w_a}$.

Next, $d_a$ is trivial on $I_p\cap\ker\eta$ for $p\in S\setminus T_b$.
Let $s$ be in this group and $m=\#k^\times$.
Then $s^m\in N_p$, so $M_{s^m}=1$ by ($\varepsilon$) and $\Psi(s^m)=1$ by Lemma~\ref{sha-lifts-local}(c), and $d_a(s)^m=1$.
Write $d_a(s)=1+y$ with $y\in\mathfrak m_A$; then $0=(1+y)^m-1=y(m+yz)$ for some $z\in A$, and $m+yz$ is a unit since $\ell\nmid m$, so $d_a(s)=1$.

Finally, let $s\in I_p\cap\ker\eta$ for some $p\in S$, and $g\in\Gamma$; then $d_a(gsg^{-1})=1$.
If $g\in\ker\eta$, this is $d_a(s)=1$.
Otherwise $g=g'c$ with $g'\in\ker\eta$, and the matrix part of $r(csc^{-1})$ is $J\,{}^tM_s^{-1}J$, whose entry $(a,a)$ is the entry $(4-a,4-a)$ of $M_{s^{-1}}$; the same holds for $r_0$, so $d_a(csc^{-1})=d_{4-a}(s^{-1})=1$.
Hence, for every $n\ge1$, the reduction of $d_a$ modulo $\mathfrak m_A^n$ is a continuous homomorphism from $\ker\eta$ to an abelian group that is trivial on every conjugate in $\ker\eta$ of every $I_p$, and it is trivial by Lemma~\ref{sha-lifts-residual}(c).
By Krull's intersection theorem $d_a=1$.
So the diagonal entries of $M_g$ and $r_0(g)$ agree for $g\in\ker\eta$, and for $g=g'c$ they are those of $g'$, since the matrix parts of $r(c)$ and $r_0(c)$ are $1$.
\end{proof}
\leanref{Development/AutomorphicArithmetic/AutomorphicVarieties/UnitaryThreeDeformations/DiagEq.lean}{UpperTheory.V.UnitaryThreeInertia.diag_eq_of_tri}

\begin{lemma}\label{sha-lifts-finiteness-criterion}
Let $\mathcal O$ be as in Definition~\ref{sha-character-parts}, with a uniformizer $\varpi$, let $R$ be a complete noetherian local $\mathcal O$-algebra with finite residue field, and let $x_1,\dots,x_n\in\mathfrak m_R$ and $w_1,\dots,w_n\in\mathbf Z$.
Suppose that the homomorphism of $\mathcal O$-algebras $\pi:\mathcal O[[X_1,\dots,X_n]]\to R$ with $\pi(X_i)=x_i$ is surjective and that its kernel is homogeneous when $X_i$ is given the weight $w_i$, and let $R_0$ be the image of the power series of weight $0$.
Suppose that every local homomorphism of $\mathcal O$-algebras $f:R\to A$, to a complete discrete valuation ring $A$ with finite residue field in which $\varpi$ maps to $0$, kills the $x_i$ with $w_i=0$, and either all $x_i$ with $w_i<0$ or all $x_i$ with $w_i>0$.
Then $R_0$ is a finite $\mathcal O$-module.
\end{lemma}
\begin{proof}
By Dickson's lemma the nonconstant monomials of weight $0$ that are minimal for divisibility are finite in number, $m_1,\dots,m_t$, and every monomial of weight $0$ is a product of them.
So $R_0$ is the image of $\mathcal O[[Y_1,\dots,Y_t]]$ under $Y_j\mapsto b_j=\pi(m_j)$: a complete noetherian local ring with finite residue field, whose maximal ideal $\mathfrak m_0$ is generated by $\varpi$ and $b_1,\dots,b_t$.

Each $b_j$ lies in the radical of $\varpi R$.
Otherwise some prime of $R$ contains $\varpi$ and not $b_j$; let $\mathfrak p$ be maximal among such primes.
Every nonzero prime of the local domain $R/\mathfrak p$ contains the image of $b_j$, which is nonzero and not a unit, so $R/\mathfrak p$ has dimension $1$: a prime of height $2$ would contain infinitely many primes of height $1$, all minimal over $b_j$.
A valuation ring of the fraction field of $R/\mathfrak p$ that dominates $R/\mathfrak p$ is, by the Krull-Akizuki theorem, a discrete valuation ring $A$, finite over the complete local ring $R/\mathfrak p$; so $A$ is complete with finite residue field, and $f:R\to R/\mathfrak p\to A$ is a local homomorphism with kernel $\mathfrak p$ in which $\varpi$ maps to $0$.
The monomial $m_j$ is not constant and has weight $0$, so it contains a variable of weight $0$, or variables of both signs.
By hypothesis $f$ kills one of them, so $f(b_j)=0$ and $b_j\in\mathfrak p$, a contradiction.

So $b_j^{N}=\varpi y$ for some $N\ge1$ and $y\in R$.
Write $b_j^N=\pi(\mu)$ with $\mu$ of weight $0$ and $y=\pi(\nu)$; since $\mu-\varpi\nu$ lies in the homogeneous kernel of $\pi$, so does its component $\mu-\varpi\nu_0$ of weight $0$, and $b_j^N=\varpi\pi(\nu_0)\in\varpi R_0$.
Hence $\mathfrak m_0$ lies in the radical of $\varpi R_0$, and since it is finitely generated, $\mathfrak m_0^K\subset\varpi R_0$ for some $K$.
So $R_0/\varpi R_0$ is finite, and $R_0$, which is complete and separated for the $\varpi$-adic topology, is a finite $\mathcal O$-module by Nakayama's lemma.
\end{proof}
\leanref{Development/AutomorphicArithmetic/AutomorphicVarieties/CurveDVRSpecialization/FinCrit.lean}{UpperTheory.V.Pr.finCrit}
\leanref{Development/AutomorphicArithmetic/AutomorphicVarieties/CurveDVRSpecialization/IsUnitOrIsUnitOfAddOne.lean}{UpperTheory.V.Pr.FinCrit}

\begin{lemma}\label{sha-lifts-automorphy}
Assume that Galois representations exist as in Theorem~\ref{sha-residually-reducible-automorphy}.
In Definition~\ref{sha-lifts-problem}, suppose that the data satisfy the dyadic condition for some $o\ge1$ and the Bernoulli condition for some $q\ge1$, let $K=\mathbf Q(i)$ and let $\iota:\bar{\mathbf Q}_\ell\cong\mathbf C$.
Then there is $g_1$ as in (f) such that the universal unitary lift $u$ relative to $g_1$, with values in $R$, has the following properties.
\begin{enumerate}
\item $R_0$ is finite over $\mathcal O$.
\item Let $r$ be the representation of a point of $u$, and suppose that $r$ is irreducible.
Then the matrix parts of $r$, inflated to $\operatorname{Gal}(\bar{\mathbf Q}/\mathbf Q)$, have no common invariant subspace other than $0$ and $\bar{\mathbf Q}_\ell^3$ on any open subgroup of $\operatorname{Gal}(\bar{\mathbf Q}/K)$, nor on $\operatorname{Gal}(\bar{\mathbf Q}/K)\cap\ker\chi_{\mathrm{cyc}}$.
Their restriction $\rho:\operatorname{Gal}(\bar{\mathbf Q}/K)\to GL_3(\bar{\mathbf Q}_\ell)$, twisted by $\chi_{\mathrm{cyc}}^{-1}$, is attached through $\iota$ to a regular algebraic conjugate self-dual cuspidal automorphic representation of $GL_3(\mathbf A_K)$, and is windowed (Theorem~\ref{sha-residually-reducible-automorphy}).
\end{enumerate}
\end{lemma}
\begin{proof}
Let $c'\in\operatorname{Gal}(\bar{\mathbf Q}/\mathbf Q)$ be a complex conjugation that maps to $c$, and $B$ a place of $\bar{\mathbf Q}$ above the place of $\mathbf Q_S$ fixed above $\ell$.
Composed with $\operatorname{Gal}(\bar{\mathbf Q}/\mathbf Q)\to\Gamma$, a homomorphism that satisfies ($\alpha$) to ($\varepsilon$) is an ordinary point in the sense of Definition~\ref{sha-ordinary-deformation}(c) for the characters $\eta$ and $\Psi$, relative to $c'$, $g_1$ and $B$: conditions ($\alpha$) to ($\gamma$) are the conditions there, and at an inertia element $s$ outside $\ell$ with $\eta(s)=\Psi(s)=1$ the homomorphism is $1$, by ($\delta$) and ($\varepsilon$) at the primes of $S$ and because $s$ maps to $1$ in $\Gamma$ at the other primes.
Let $g_1$ and $U$ be given by Proposition~\ref{sha-lifts-automorphy-form} for $c'$ and $B$; the image of $g_1$ in $\Gamma$ is as in (f).

(a) Let $A$ be a complete discrete valuation ring with finite residue field, and $f:R\to A$ a local homomorphism of $\mathcal O$-algebras in which $\mathfrak m_{\mathcal O}$ maps to $0$.
Then $f\circ u$ satisfies ($\alpha$) to ($\varepsilon$), so by Proposition~\ref{sha-lifts-automorphy-form}(i) its characteristic polynomials on the image of $U\cap\ker\eta$ are those of $r_0$.
Lemma~\ref{sha-lifts-triangularity}, applied to $f\circ u$ composed with $\operatorname{Gal}(\bar{\mathbf Q}/\mathbf Q)\to\Gamma$, with $\bar\Psi$, $c'$ and $g_1$, shows that the matrix parts of $f\circ u$ on $\ker\eta$ are all upper triangular or all lower triangular, and Lemma~\ref{sha-lifts-finiteness-diagonal} that its diagonal entries are those of $r_0$.
Let $x_i$ and $w_i$ be as in Lemma~\ref{sha-lifts-presentation}(a).
Each $x_i$ is $u(g)_{ab}-r_0(g)_{ab}$ for some $g\in\ker\eta$, with $w_i=b-a$; so $f$ kills the $x_i$ of weight $0$, which come from diagonal entries, and either all those of negative weight or all those of positive weight.
By Lemma~\ref{sha-lifts-finiteness-criterion}, $R_0$ is finite over $\mathcal O$.

(b) By (a) and Lemma~\ref{sha-lifts-points}(b), $r$ is of characteristic zero: for some $t$, $r'=\tau(t)r\tau(t)^{-1}$ satisfies ($\alpha$) to ($\varepsilon$) with values in a discrete valuation ring $A\subset\bar{\mathbf Q}_\ell$ finite over $\mathcal O$, so it is an ordinary point.
It is irreducible on $\ker\eta$, as $r$ is.

\emph{Open subgroups.} Let $U'$ be an open subgroup of $\ker\eta$; we may assume that it is normal in $\operatorname{Gal}(\bar{\mathbf Q}/\mathbf Q)$.
Let $x$ be a lift of the element $x_\ell$ of Lemma~\ref{sha-lifts-local}(b) to the decomposition group at $B$, let $\sigma$ be an element of the inertia group at $B$ with $\chi_{\mathrm{cyc}}(\sigma)=1+\ell$, and let $u=\sigma^m$, where $m$ is the index of $U'$ times $\varphi(N)$.
Then $u\in U'$, $\psi(u)=1$ and $\Psi(u)^2=(1+\ell)^{2em}\ne1$.
By ($\gamma$), after conjugation $r'$ is upper triangular on the decomposition group at $B$ intersected with $\ker\eta$, with diagonal $\Psi,1,\Psi^{-1}$ on the inertia group, which lies in $\ker\eta$; so $r'(u)$ has three distinct eigenvalues, and $r'(x)$ is upper triangular, with $\eta(x)=1$ and $\bar\Psi(x)^2\ne1$.
These are the properties of $u$ and $x$ used in the paragraph on irreducibility on open subgroups in the proof of Theorem~\ref{sha-residually-reducible-automorphy}, and its argument shows that $r'$, and so $r$, is irreducible on $U'$.

\emph{The cyclotomic kernel.} Suppose that $r$ is reducible on $H=\ker\eta\cap\ker\chi_{\mathrm{cyc}}$, a normal subgroup of $\ker\eta$ with abelian quotient.
By Clifford's theorem the restriction of $r$ to $H$ is semisimple, so it has a nonzero endomorphism of trace $0$.
The endomorphisms of trace $0$ of the restriction to $H$ form a space on which $\ker\eta$ acts by conjugation through its abelian quotient by $H$, so they have a common eigenvector $T$: $r(g)Tr(g)^{-1}=\phi(g)^{-1}T$ for a character $\phi$ of $\ker\eta$ trivial on $H$.
The kernel of $T$ is stable under $\ker\eta$, so $T$ is invertible, and taking determinants gives $\phi^3=1$.
The kernel of $\phi$, which is the set of $g$ with $Tr(g)=r(g)T$, is closed of finite index, hence open, so $r$ is irreducible on it by the previous paragraph, and $T$ is a scalar by Schur's lemma.
This contradicts $T\ne0$ and $\operatorname{tr}T=0$.

\emph{Automorphy.} By Proposition~\ref{sha-lifts-automorphy-form}(ii), applied to $r'$, the restriction of $r'$ to $\operatorname{Gal}(\bar{\mathbf Q}/K)$, twisted by $\chi_{\mathrm{cyc}}^{-1}$, is attached through $\iota$ to a regular algebraic conjugate self-dual cuspidal automorphic representation of $GL_3(\mathbf A_K)$, and is windowed; $\rho$ is conjugate to it by $\tau(t)$.
\end{proof}
\leanref{Development/AutomorphicArithmetic/AutomorphicVarieties/ModularityArgument/ThresholdArguments/DenseThresholdCases/ThmS6CoreS.lean}{UpperTheory.V.Aut.ThmS6_s}
\leanref{Development/AutomorphicArithmetic/AutomorphicVarieties/ModularityArgument/ThresholdArguments/DenseThresholdCases/ThmS6CoreS.lean}{UpperTheory.V.Aut.ThmS6Core_s}
\leanref{Development/AutomorphicArithmetic/AutomorphicVarieties/ModularityArgument/ThresholdArguments/DenseThresholdCases/ThmS6CoreS.lean}{UpperTheory.V.Aut.BPrime_s}
\leanref{Development/AutomorphicArithmetic/AutomorphicVarieties/ModularityArgument/ThresholdArguments/DenseThresholdCases/UnitaryLiftTheoremOfReducedForm.lean}{UpperTheory.V.UnitaryThreeInertia.unitary_lift_of_reduced_core}
\leanref{Development/AutomorphicArithmetic/AutomorphicVarieties/ModularityArgument/ThresholdArguments/DenseThresholdCases/SymmetricCubeTheoremOfEndomorphisms.lean}{UpperTheory.V.Pr.symmetric_cube_theorem_of_endomorphisms_alt}
\leanref{Development/AutomorphicArithmetic/AutomorphicVarieties/UnitaryThreeDeformations/UniversalDiagonalOfTriangularPoints.lean}{UpperTheory.V.UnitaryThreeInertia.universal_diagonal_of_triangular_points}
\leanref{Development/AutomorphicArithmetic/AutomorphicVarieties/UnitaryThreeDeformations/IrrOpen.lean}{UpperTheory.V.UnitaryThreeInertia.matIrr_open}
\leanref{Development/AutomorphicArithmetic/AutomorphicVarieties/UnitaryThreeDeformations/IrrCyc.lean}{UpperTheory.V.UnitaryThreeInertia.matIrr_cyc}
\leanref{Development/AutomorphicArithmetic/AutomorphicVarieties/CurveDVRSpecialization/RestrictionToDecompositionGroup.lean}{UpperTheory.V.Pr.ordPt_of_isPt}

\begin{lemma}\label{sha-lifts-level}
In Definition~\ref{sha-lifts-problem}, suppose that the data satisfy the dyadic condition for some $o\ge1$, let $g_1$ be as in (f), and let $K=\mathbf Q(i)$.
Let $r:\Gamma\to\mathcal G_3(\bar{\mathbf Q}_\ell)$ be of characteristic zero, and let $\rho:\operatorname{Gal}(\bar{\mathbf Q}/K)\to GL_3(\bar{\mathbf Q}_\ell)$ be the restriction of its matrix parts, inflated to $\operatorname{Gal}(\bar{\mathbf Q}/\mathbf Q)$.
Then $\rho$ has level $N$ and weights $(e,0,-e)$ (Definition~\ref{sha-eigensystem-level}), and at every prime $w$ of $K$ above $\ell$ the characters $\kappa_1,\kappa_2,\kappa_3$ of that definition can be taken with $\kappa_i\ne\kappa_j\chi_{\mathrm{cyc}}$ on $D_w$ for all $i,j$.
\end{lemma}
\begin{proof}
Let $t\in\bar{\mathbf Q}_\ell^\times$ be such that $\tau(t)r\tau(t)^{-1}$ satisfies ($\alpha$) to ($\varepsilon$) with values in a discrete valuation ring $A\subset\bar{\mathbf Q}_\ell$ finite over $\mathcal O$.
On $\ker\eta=\operatorname{Gal}(\bar{\mathbf Q}/K)$ its matrix parts are those of $r$ conjugated by $\tau(t)$, and both conclusions are invariant under conjugation; so we may assume that $r$ itself satisfies ($\alpha$) to ($\varepsilon$).
Write $M_g$ for the matrix part of $r(g)$, $(w_1,w_2,w_3)=(1,0,-1)$, and $\theta$ for the restriction of $\psi$ to $\operatorname{Gal}(\bar{\mathbf Q}/K)$, which is trivial on $\operatorname{Gal}(\bar{\mathbf Q}/K(\mu_N))$; then $\Psi=\chi_{\mathrm{cyc}}^e\theta$ on $\operatorname{Gal}(\bar{\mathbf Q}/K)$.
Conjugation by $r(\gamma)$ acts on matrix parts by $X\mapsto M_\gamma XM_\gamma^{-1}$ if $\eta(\gamma)=1$ and by $X\mapsto M_\gamma J\,{}^tX^{-1}JM_\gamma^{-1}$ if $\eta(\gamma)=-1$; the map $X\mapsto J\,{}^tX^{-1}J$ preserves the upper triangular matrices and the diagonal matrices $\operatorname{diag}(z,1,z^{-1})$.
So the local statements below may be checked at the places of $\bar{\mathbf Q}$ that lie over the places of $\mathbf Q_S$ fixed in Definition~\ref{sha-lifts-problem}, the characters $\Psi$ and $\chi_{\mathrm{cyc}}$ being invariant under conjugation.

\emph{Primes not above $\ell$.} Let $w$ be such a prime of $K$, above $p$, with inertia group $I_w$.
If $p\notin S$, $I_w$ maps to $1$ in $\Gamma$ and $\rho$ is trivial on it.
Let $p\in S$, $p\ne\ell$; it suffices to show that the $M_s$ with $s\in I_p\cap\ker\eta$ are simultaneously conjugate to $\operatorname{diag}(\Psi(s),1,\Psi(s)^{-1})$, since $\chi_{\mathrm{cyc}}$ is trivial on $I_p$, so that $\Psi=\theta$ there.
For $p\in S_b$ this is ($\delta$).
For $p\notin T_b$, the homomorphisms $r$ and $r_0$ are trivial on $N_p$, by ($\varepsilon$) and Lemma~\ref{sha-lifts-local}(c), so on $I_p\cap\ker\eta$ they factor through the finite group $Q=(I_p\cap\ker\eta)/N_p$, which embeds in $k^\times$ through $\bar\Psi$, so that its order $m$ is prime to $\ell$.
Let $X=\sum_{q\in Q}M_qr_0(q)^{-1}$, with $r_0(q)$ standing for its matrix part.
Then $M_qX=Xr_0(q)$ for every $q\in Q$, and $X\equiv m\cdot1$ modulo $\mathfrak m_A$ because $M_q\equiv r_0(q)$, so $X$ is invertible.

\emph{Primes above $\ell$.} Let $w$ be a prime of $K$ above $\ell$; we may take its decomposition and inertia groups $D_w$ and $I_w$ to map onto $D_\ell\cap\ker\eta$ and $I_\ell$, and work with these images, through which $\rho$ factors.
Let $h$ be as in ($\gamma$).
Then $h^{-1}M_sh$ is upper triangular for $s\in D_w$, and its diagonal entries are characters $\kappa_1,\kappa_2,\kappa_3$ of $D_w$, congruent to $\bar\Psi^{w_i}$ since $h\equiv1$ and $M_s\equiv r_0(s)$, and equal to $\Psi^{w_i}=\chi_{\mathrm{cyc}}^{ew_i}\theta^{w_i}$ on $I_w$.
So $\rho$ has weights $(e,0,-e)$ at $w$, with the characters $\theta^{w_i}$.
Let $i,j\in\{1,2,3\}$.
If $i=j$, then $\kappa_i\ne\kappa_i\chi_{\mathrm{cyc}}$: for $\sigma\in I_\ell$ with $\chi_{\mathrm{cyc}}(\sigma)=1+\ell$, $\sigma^2\in D_w$ and $\chi_{\mathrm{cyc}}(\sigma^2)\ne1$.
If $i\ne j$, let $x=x_\ell$ be as in Lemma~\ref{sha-lifts-local}(b); then $x\in D_w$, $\chi_{\mathrm{cyc}}(x)\equiv1\pmod\ell$, and $\kappa_i(x)$ and $\kappa_j(x)\chi_{\mathrm{cyc}}(x)$ are congruent to $\bar\Psi(x)^{w_i}$ and $\bar\Psi(x)^{w_j}$, which differ because $0<|w_i-w_j|\le2$ and $\bar\Psi(x)^2\ne1$.
\end{proof}
\leanref{Development/AutomorphicArithmetic/AutomorphicVarieties/UnitaryThreeDeformations/LevelAway.lean}{UpperTheory.V.UnitaryThreeInertia.levelAway_of_pt}
\leanref{Development/AutomorphicArithmetic/AutomorphicVarieties/UnitaryThreeDeformations/LevelAt.lean}{UpperTheory.V.UnitaryThreeInertia.levelAt_of_pt}

\begin{lemma}\label{sha-lifts-rigidity}
Assume that Galois representations exist as in Theorem~\ref{sha-residually-reducible-automorphy}.
In Definition~\ref{sha-lifts-problem}, suppose that the data satisfy the dyadic condition for some $o\ge1$, let $g_1$ be as in (f), let $K=\mathbf Q(i)$ and $\iota:\bar{\mathbf Q}_\ell\cong\mathbf C$.
Let $r:\Gamma\to\mathcal G_3(\bar{\mathbf Q}_\ell)$ be of characteristic zero, such that its matrix parts, inflated to $\operatorname{Gal}(\bar{\mathbf Q}/\mathbf Q)$, have no common invariant subspace other than $0$ and $\bar{\mathbf Q}_\ell^3$ on any open subgroup of $\operatorname{Gal}(\bar{\mathbf Q}/K)$, and such that their restriction to $\operatorname{Gal}(\bar{\mathbf Q}/K)$, twisted by $\chi_{\mathrm{cyc}}^{-1}$, is attached through $\iota$ to a regular algebraic conjugate self-dual cuspidal automorphic representation of $GL_3(\mathbf A_K)$ and is windowed (Theorem~\ref{sha-residually-reducible-automorphy}).
Then $r$ is rigid.
\end{lemma}
\begin{proof}
\emph{Step 1.} The Bloch-Kato Selmer group of the adjoint representation of $r$ over $\mathbf Q$ vanishes.
Let $r'=\tau(t)r\tau(t)^{-1}$ satisfy ($\alpha$) to ($\varepsilon$) with values in a discrete valuation ring $A\subset\bar{\mathbf Q}_\ell$ finite over $\mathcal O$, and let $\rho'$ be its restriction to $\operatorname{Gal}(\bar{\mathbf Q}/K)$.
As in the proof of Lemma~\ref{sha-lifts-automorphy}, $r'$ is an ordinary point relative to $g_1$, a complex conjugation $c'$ that maps to $c$, and a place $B$ above the place fixed above $\ell$.
It is irreducible on every open subgroup of $\operatorname{Gal}(\bar{\mathbf Q}/K)$, and $\rho'\otimes\chi_{\mathrm{cyc}}^{-1}$ is attached through $\iota$ to a regular algebraic conjugate self-dual cuspidal automorphic representation and is windowed, since $\rho'$ is conjugate to the restriction of $r$ by $\tau(t)$.
Lemma~\ref{sha-rigidity-adjoint}, for these characters, $c'$ and $B$, derives the vanishing of $H^1_f(\mathbf Q,\operatorname{ad}r')$ from its hypotheses (1) to (4), which are the following properties of the characters and of the point.
\begin{itemize}
\item $e\ge1$, $\Psi(c')=1$, and $\eta$ is unramified at the primes not dividing an integer $M$ divisible by $\ell N$.
\item For some $g_2\in\operatorname{Gal}(\bar{\mathbf Q}/K)$ with $\Psi(g_2)^2-1\in\mathcal O^\times$, a link of $\eta$, $\Psi$, $c'$, $g_2$ and $B$ to an abelian CM field $L\supset K$ that satisfies the residual conditions, and a seed datum for $L$.
\item The irreducibility of $r'$ on every open subgroup of $\ker\eta$, and the automorphy and the windowed condition of $\rho'\otimes\chi_{\mathrm{cyc}}^{-1}$.
\item At every prime $w$ of $K$ above $\ell$, a conjugate of the restriction of $\rho'$ to $D_w$ that is upper triangular, with diagonal characters $\kappa_1,\kappa_2,\kappa_3$ equal to $\chi_{\mathrm{cyc}}^{m_i}\chi_i$ on $I_w$ for $(m_1,m_2,m_3)=(e,0,-e)$ and characters $\chi_i$ of $\operatorname{Gal}(\bar{\mathbf Q}/K)$ trivial on $\operatorname{Gal}(\bar{\mathbf Q}/K(\mu_N))$, and with $\kappa_i\ne\kappa_j\chi_{\mathrm{cyc}}$ for all $i,j$.
\end{itemize}
Here the first property holds with $M=\ell N$, since $\eta$ is ramified only at $2$, and $\Psi(c')=1$ by Lemma~\ref{sha-lifts-automorphy-conditions}(d); the second is given by the proof of Proposition~\ref{sha-lifts-automorphy-form}; the third was shown above; and the last is Lemma~\ref{sha-lifts-level}, applied to $r'$.
So $H^1_f(\mathbf Q,\operatorname{ad}r')=0$, and $H^1_f(\mathbf Q,\operatorname{ad}r)=0$, since conjugation by $\tau(t)$ is an isomorphism from $\operatorname{ad}r$ to $\operatorname{ad}r'$.
% Step 1 is selmer_h1_vanishes_at_rigid_point: the link is Fr.host_ex, the condition at l is Th.Unl from Z.J.hostUnlE_alt (the first alternative of levelAt_of_pt, returned by UnlEndLogNorm), the free element is free_at_of_unlinked_of_automorphic_cyclotomic_twist, and the vanishing is Z.Win.Dn.rigEnd3_alt (Th.RigEnd3Alt), which is Lemma sha-rigidity-adjoint.

\emph{Step 2.} Let $y$ be a cocycle for $r$ with the properties in the definition of rigidity.
Its inflation to $\operatorname{Gal}(\bar{\mathbf Q}/\mathbf Q)$ satisfies the Bloch-Kato local conditions.
At the primes outside $S$ it vanishes on inertia, which maps to $1$ in $\Gamma$.
At $p\in S\setminus\{\ell\}$ its restriction to $I_p$ is a coboundary: it vanishes on $N_p$ if $p\notin T_b$, and if $p\in S_b$ the diagonal shape on $D_p\cap\ker\eta$ with fixed diagonal on inertia makes it a coboundary on $I_p\cap\ker\eta$; in characteristic $0$ a cocycle that is a coboundary on a normal subgroup of finite index is a coboundary, by averaging.
At $\ell$ the deformation $(1+\epsilon y)r$ is upper triangular on $D_\ell\cap\ker\eta$ with the diagonal characters of $r$ on inertia, so $y$ lies in the Bloch-Kato subspace at $\ell$ \cite{BK90}.
By Step 1, $y$ is a coboundary on $\operatorname{Gal}(\bar{\mathbf Q}/\mathbf Q)$, hence on $\Gamma$.
\end{proof}
\leanref{Development/AutomorphicArithmetic/AutomorphicVarieties/ModularityArgument/ThresholdArguments/DenseThresholdCases/HostNumberS.lean}{UpperTheory.V.UnitaryThreeInertia.HostRigid_s}
\leanref{Development/AutomorphicArithmetic/AutomorphicVarieties/ModularityArgument/ThresholdArguments/DenseThresholdCases/EHostRigidU.lean}{UpperTheory.V.Z.J.hostRigidUE_alt}

\begin{lemma}\label{sha-lifts-number}
Assume that Galois representations exist as in Theorem~\ref{sha-residually-reducible-automorphy}.
There is $C$ with the following property.
In Definition~\ref{sha-lifts-problem}, suppose that the data satisfy the dyadic condition for some $o\ge1$, let $g_1$ be as in (f), let $K=\mathbf Q(i)$ and $\iota:\bar{\mathbf Q}_\ell\cong\mathbf C$.
Let $\mathcal A$ be the set of homomorphisms $r:\Gamma\to\mathcal G_3(\bar{\mathbf Q}_\ell)$ of characteristic zero that are irreducible and have the properties of the representations in Lemma~\ref{sha-lifts-automorphy}(b).
Then there is $\varphi\in\Gamma$ such that $\operatorname{tr}r(\varphi)\ne\operatorname{tr}r_0(\varphi)$ for every $r\in\mathcal A$, and
\[
 \sum_{i=1}^{s}v_\ell\big(\operatorname{tr}r_i(\varphi)-\operatorname{tr}r_0(\varphi)\big)\log\ell\le(\ell Ne)^{C}
\]
for all $r_1,\dots,r_s\in\mathcal A$ whose matrix parts on $\ker\eta$ are pairwise not conjugate.
\end{lemma}
\begin{proof}
Let $a$ and $b$ be the constants $c_5$ and $A$ of the counting proposition in the proof of Proposition~\ref{sha-counts}(a), for $K=\mathbf Q(i)$, whose discriminant is $-4$.
There is a prime $p\equiv3\pmod4$ with $7\le p\le12\ell N$ and $p\nmid\ell N$, by an elementary sieve.
It is inert in $K$; let $w$ be the prime of $K$ above it, with $\mathrm Nw=p^2$, let $\sigma\in\operatorname{Gal}(\bar{\mathbf Q}/K)$ be a Frobenius element at a place above $w$, and let $\varphi$ be the image of $\sigma^{-1}$ in $\Gamma$.
Then $\operatorname{tr}r_0(\varphi)=q^{-e}z+1+q^{e}z^{-1}$ with $q=p^2$ and $z=\psi(\sigma)^{-1}$, an $m$-th root of unity for $m=\#(\mathbf Z/N)^\times$, so it is the Eisenstein value $t_z$ of the counting proposition for weight $e$.
The restriction $\rho$ to $K$ of every $r\in\mathcal A$ is a stable automorphic triple of weight $e$ and level $N$ (Definition~\ref{sha-eigensystem-triple}): it is irreducible, since $r$ is irreducible and $\ker\eta=\operatorname{Gal}(\bar{\mathbf Q}/K)$; it has level $N$ and weights $(e,0,-e)$, with characters at each prime above $\ell$ that satisfy $\kappa_i\ne\kappa_j\chi_{\mathrm{cyc}}$, by Lemma~\ref{sha-lifts-level}; and $\rho\otimes\chi_{\mathrm{cyc}}^{-1}$ is attached through $\iota$ to a regular algebraic conjugate self-dual cuspidal automorphic representation and is windowed, by the properties of Lemma~\ref{sha-lifts-automorphy}(b).
Apply the counting proposition with the family of all $m$-th roots of unity, and with the given $r_i$ in the set attached to $z$.
It gives the inequality $\operatorname{tr}r_i(\varphi)\ne t_z$, and a nonzero integer $N_0$ with $\log|N_0|\le am(N(e+1))^{b}\log p$ and $|N_0|_\ell\le\prod_i|\operatorname{tr}r_i(\varphi)-t_z|_\ell$, so that $\sum_iv_\ell(\operatorname{tr}r_i(\varphi)-t_z)\log\ell\le\log|N_0|$.
The automorphy of each $r_i$ is relative to one of the two embeddings of $K$ in $\bar{\mathbf Q}$; applying this to the two subfamilies doubles the bound.
Since $\ell\ge3$, $N\ge2$ and $p\le12\ell N$, the result $2am(N(e+1))^b\log(12\ell N)$ is at most $(\ell Ne)^{2a+b+4}$.
\end{proof}
\leanref{Development/AutomorphicArithmetic/AutomorphicVarieties/ModularityArgument/ThresholdArguments/DenseThresholdCases/HostNumberS.lean}{UpperTheory.V.UnitaryThreeInertia.HostNumber_s}
\leanref{Development/AutomorphicArithmetic/AutomorphicVarieties/ModularityArgument/ThresholdArguments/DenseThresholdCases/EHostNumberU.lean}{UpperTheory.V.Z.J.hostNumberUE_alt}

\begin{lemma}\label{sha-model-local-component}
Use the notation of Lemma~\ref{sha-model-local}, which follows. Let $I_B$ be the ideal of the $a\in R_B$ with $\delta_B^na\in J_B$ for some $n\ge0$, and $\mathfrak c_B$ its radical, so that $\mathfrak a_B$ is the intersection of the minimal primes of $\mathfrak c_B$ contained in $\ker x_B$. Let $\mathfrak m$ be the maximal ideal of $E[[Y_1,\dots,Y_{12}]]$. Finite $\mathbf Q_\ell$-algebras $C$ below carry an $\mathcal O$-algebra structure compatible with their $\mathbf Z_\ell$-algebra structure, and the $\ell$-adic topology.
\begin{enumerate}
\item (The component.) The ideal $\mathfrak a_B$ is prime, it is the only minimal prime of $\mathfrak c_B$ contained in $\ker x_B$, $\ell\notin\mathfrak a_B$, and $R_B/\mathfrak a_B$ has Krull dimension $13$.
\item (Completion.) There is a homomorphism of $\mathcal O$-algebras $\varphi:R_B\to E[[Y_1,\dots,Y_{12}]]$ with kernel $\mathfrak a_B$ such that the constant term of $\varphi(f)$ is $x_B(f)$; for every $N\ge1$ and every power series $g$ there are $f\in R_B$ and $c\in\mathcal O\setminus\{0\}$ with $cg-\varphi(f)\in\mathfrak m^N$; and if $\varphi(f)\in\mathfrak m^N$, then $cf\in(\ker x_B)^N+\mathfrak a_B$ for some $c\in\mathcal O\setminus\{0\}$.
\item (Crystalline points.) Let $C$ be a finite $\mathbf Q_\ell$-algebra and $h:R_B\to C$ a homomorphism of $\mathcal O$-algebras that kills $I_B$. Then the entries of $h\circ u_B$ are continuous, and the matrix parts of $h\circ u_B$ on $H_B$ form a crystalline representation \cite{Fon94}.
\item (Ordinary points.) Let $C$ be a finite $\mathbf Q_\ell$-algebra that is an $E$-algebra, with a homomorphism of $E$-algebras $\pi:C\to E$ whose kernel consists of nilpotent elements, and let $r:D_B\to\mathcal G_3(C)$ be a homomorphism with multiplier $\eta$ and continuous entries, with $\pi\circ r=r_0$, whose matrix parts on $H_B$ form a crystalline representation. Then $r=nun^{-1}$ for a homomorphism $u:D_B\to\mathcal G_3(C)$ with upper triangular matrix parts and a matrix $n\in GL_3(C)$ congruent to $1$ modulo $\ker\pi$; and the diagonal entries of the matrix part of $u(s)$, for $s\in I_B$, are $\Psi(s),1,\Psi(s)^{-1}$.
\item (Conjugation.) The element $\delta_B$ lies in no minimal prime of $\mathfrak c_B$. Let $S$ be a domain that is an $\mathcal O$-algebra, let $i,\alpha:R_B\to S$ be homomorphisms of $\mathcal O$-algebras, and let $U\in GL_3(S)$ with $\alpha\circ u_B=U(i\circ u_B)U^{-1}$. If $i$ kills $I_B$ and $i(\delta_B)\ne0$, then $\alpha$ kills $\mathfrak c_B$.
\item (Dual numbers.) Let $K$ be a finite extension of $\mathbf Q_\ell$ that is an $\mathcal O$-algebra, let $\rho:D_B\to\mathcal G_3(K)$ be a homomorphism with multiplier $\eta$ and continuous matrix part $s\mapsto M_s$, let $\gamma$ be a continuous cocycle of $D_B$ for the adjoint action of $\rho$ (Definition~\ref{sha-lifts-graded}(e)), and let $\rho':D_B\to\mathcal G_3(K[\epsilon]/(\epsilon^2))$ be a homomorphism with the signs of $\rho$ and the matrix parts $M_s+\epsilon\gamma(s)M_s$. Then the matrix parts of $\rho'$ on $H_B$ form a crystalline representation if and only if those of $\rho$ do and $\gamma$ lies in the Bloch-Kato subspace \cite{BK90}, that is, $\gamma(s)\otimes1=sX-X$ for some $X\in M_3(K)\otimes_{\mathbf Q_\ell}B_{\mathrm{cris}}$ and all $s\in D_B$.
\item (Weights.) Let $\varphi'$ have the properties of $\varphi$ in (2), and let $w$ be a function on $\{1,\dots,12\}$ such that, with $Y_i$ of weight $w(i)$, $\varphi'$ maps the entry $(a,b)$ of the matrix part of every $u_B(h)$ to a power series of weight $b-a$. Suppose that for each $i$ some entry $e_i$ of the matrix part of some $u_B(h)$ has $\varphi'(e_i)$ with linear part $Y_i$. Then $w$ takes the values $-2,-1,0,1,2$ on at most $1,2,3,4,2$ elements respectively, and no other value.
\end{enumerate}
\end{lemma}
\leanref{Development/AutomorphicArithmetic/AutomorphicVarieties/ModularityArgument/CrystallineOrdinaryIdeals/OrdinaryComponentUniquePrimeDimThirteenNearBase.lean}{GlobalTheory.OrdLoc.component_basic}
\leanref{Development/AutomorphicArithmetic/AutomorphicVarieties/PadicHodgeTheory/RtCompl.lean}{UpperTheory.V.Loc.rtCompl}
\leanref{Development/AutomorphicArithmetic/AutomorphicVarieties/ModularityArgument/CrystallineOrdinaryIdeals/SaturatedOrdinaryIdealCutsOutCrystallineLocusStatement.lean}{GlobalTheory.T.saturatedOrdinaryIdeal_cuts_out_crystalline_locus}
\leanref{Development/AutomorphicArithmetic/AutomorphicVarieties/PadicHodgeTheory/ExistsEps.lean}{UpperTheory.V.Loc.ordEq}
\leanref{Development/AutomorphicArithmetic/AutomorphicVarieties/PadicHodgeTheory/TriDiag.lean}{UpperTheory.V.Loc.triDiag}
\leanref{Development/AutomorphicArithmetic/AutomorphicVarieties/ModularityArgument/CrystallineOrdinaryIdeals/SaturatingElementNotMemMinimalPrimesReducedOrdinStatement.lean}{GlobalTheory.T.saturatingElement_notMem_minimalPrimes_reducedOrdinaryIdeal}
\leanref{Development/AutomorphicArithmetic/AutomorphicVarieties/ModularityArgument/CrystallineOrdinaryIdeals/ReducedOrdinaryIdealLeKerOfConjOfSaturatingNeZerStatement.lean}{GlobalTheory.T.reducedOrdinaryIdeal_le_ker_of_conj_of_saturating_ne_zero}
\leanref{Development/AutomorphicArithmetic/AutomorphicVarieties/PadicHodgeTheory/DualCr.lean}{UpperTheory.V.Loc.dualCr}
\leanref{Development/AutomorphicArithmetic/AutomorphicVarieties/PadicHodgeTheory/ChartCount.lean}{UpperTheory.V.Loc.chartCount}
\begin{proof}
\openstep{The proofs of (1) to (7) are not yet written out: the smoothness of the ordinary component of the local deformation ring at the base point and its dimension, the description of its points over finite $\mathbf Q_\ell$-algebras as the crystalline ones, the triangular form of crystalline ordinary deformations, the stability of the saturated ideal under conjugation, the crystalline deformations to the dual numbers, and the count of the tangent directions of each weight.}
% Size: about 14,300 lines of Lean in PadicHodgeTheory and ModularityArgument/CrystallineOrdinaryIdeals, where (1), (2), (3), (4) and (7) take the structure Loc.LocalPublishedTheorems (local finiteness, local Tate duality and the local Euler characteristic, for the base field and its finite extensions, and the local Bloch-Kato theorem) as hypothesis, and (5), (6) and Loc.triDiag do not.
% (1) is GlobalTheory.OrdLoc.component_basic (from T.exists_unique_ordinaryPrime_dim_succ_near_base, T.saturatedOrdinaryIdeal_tangent_finrank_eq_twelve and finrank_pointDerivations_vanishing_eq_of_isCompl); I_B is Dat3.IY, c_B is Dat3.Icr, a_B is Dat3.It (the infimum of primesAt0). (2) is Loc.rtCompl (OrdLoc.rtCompl; the properties are Loc.IsCompl: ker, const, dense, inj). (3) is the first half (PtsCr) of T.saturatedOrdinaryIdeal_cuts_out_crystalline_locus (OrdLoc.ideal_cuts_crystalline_locus, about 6,400 lines), with ptCont for the continuity. (4) is Loc.ordEq (InDcr iff InDord, about 9,900 lines with (3)) and Loc.triDiag (about 2,100 lines). (5) is T.saturatingElement_notMem_minimalPrimes_reducedOrdinaryIdeal (45 lines) and T.reducedOrdinaryIdeal_le_ker_of_conj_of_saturating_ne_zero (OrdLoc.icr_le_ker_of_conj, about 680 lines, without the published theorems). (6) is Loc.dualCr (about 2,400 lines with DeformCr, AdmBilin, AdmHom). (7) is Loc.chartCount with locDims.
\end{proof}

\begin{lemma}\label{sha-model-local}
Let $(\ell,k,\psi,K)$ be an Eisenstein datum with deformation data $(\eta,c,\Psi)$ (Definition~\ref{sha-ordinary-deformation}), let $E$ and $\mathbf F$ be the fraction field and the residue field of $\mathcal O$, put $r_0=r_0(\eta,\Psi)$, and let $B$ be a place above $\ell$, with decomposition group $D_B$ and inertia group $I_B$.
Let $u_B:D_B\to\mathcal G_3(R_B)$ be the universal continuous lift of the reduction of $r_0$ to $\mathcal G_3$ with multiplier $\eta$, with values in a complete noetherian local $\mathcal O$-algebra $R_B$ with residue field $\mathbf F$, and let $x_B:R_B\to\mathcal O$ be the homomorphism with $x_B\circ u_B=r_0$.
For $s\in D_B$ with $\eta(s)=1$ let $M_s$ be the matrix part of $u_B(s)$.
Let $J_B\subset R_B$ be the ideal generated by the coefficients of $\det(X-M_s)-(X-\Psi(s))(X-1)(X-\Psi(s)^{-1})$ and by the entries of $(M_{s_1}-\Psi(s_1))(M_{s_2}-1)(M_{s_3}-\Psi(s_3)^{-1})$, for $s,s_1,s_2,s_3\in I_B$ at which $\eta$ and the $\ell$-component of $\psi$ are trivial.
Let $H_B$ be the subgroup of the elements of $D_B$ at which $\eta$ and the $\ell$-component of $\psi$ are trivial.
Fix $s_0\in I_B\cap H_B$ with $\chi_{\mathrm{cyc}}(s_0)$ of infinite order, and $\phi\in H_B$ that is a power of a Frobenius element times an element of $I_B$ and generates, with $I_B\cap H_B$, a dense subgroup of $H_B$.
Put $(\Psi_1,\Psi_2,\Psi_3)=(\Psi,1,\Psi^{-1})$, with indices modulo $3$, and
\[
 \beta_a=\big(\Psi_a(s_0)-\Psi_{a+1}(s_0)\big)\big(\Psi_a(s_0)-\Psi_{a+2}(s_0)\big),\qquad \gamma_a=\operatorname{tr}\big((M_{s_0}-\Psi_{a+1}(s_0))(M_{s_0}-\Psi_{a+2}(s_0))M_\phi\big),
\]
$\delta_a=\beta_{a+1}\gamma_a-\chi_{\mathrm{cyc}}(\phi)\beta_a\gamma_{a+1}$ and $\delta_B=\ell\delta_1\delta_2$.
The \emph{ordinary ideal} $\mathfrak a_B$ is the intersection of the minimal primes of $\{a\in R_B:\delta_B^na\in J_B\text{ for some }n\ge0\}$ that are contained in $\ker x_B$.
\begin{enumerate}
\item (Chart.) There are a function $w$ on $\{1,\dots,12\}$ that takes the values $-2,-1,0,1,2$ on $1,2,3,4,2$ elements respectively, and a homomorphism of $\mathcal O$-algebras $\varphi_B:R_B\to E[[Y_1,\dots,Y_{12}]]$ with kernel $\mathfrak a_B$, with the following properties, where $Y_i$ has weight $w(i)$: the entry $(a,b)$ of the matrix part of every $u_B(h)$ maps to a power series of weight $b-a$; the constant term of $\varphi_B(f)$ is $x_B(f)$; every $Y_i$ is the only linear term of $\varphi_B(f)$ for some $f$; and if $x_B(f)=0$ and $\varphi_B(f)$ has no linear term, then a nonzero multiple of $f$ lies in $(\ker x_B)^2+\mathfrak a_B$.
\item (Conjugation.) Let $A$ be a complete noetherian local $\mathcal O$-algebra with residue field $\mathbf F$, let $h:R_B\to A$ kill $\mathfrak a_B$, and let $P\in GL_3(A)$ be congruent to $1$ modulo $\mathfrak m_A$. Then the homomorphism $h'$ with $h'\circ u_B=P(h\circ u_B)P^{-1}$ kills $\mathfrak a_B$.
\item (Triangular form.) Let $L$ be a field over $\mathcal O$ and let $h:R_B\to L$ kill $\mathfrak a_B$. Over a finite extension of $L$, the matrix parts of $h\circ u_B$ on $D_B\cap\ker\eta$ are simultaneously conjugate to upper triangular matrices whose diagonal entries at the elements of $I_B\cap\ker\eta$ are $\Psi,1,\Psi^{-1}$.
\item (Tangent vectors.) Let $h:R_B\to\bar{\mathbf Q}_\ell[\epsilon]/(\epsilon^2)$ kill $\mathfrak a_B$, and suppose that $h_0=h\bmod\epsilon$ maps $\mathfrak m_{R_B}$ into the maximal ideal of the integers of $\bar{\mathbf Q}_\ell$. Write the matrix part of $h\circ u_B(s)$ as $M_s+\epsilon M'_s$. Then the cocycle $s\mapsto M'_sM_s^{-1}$ of $D_B$, for the adjoint action of $h_0\circ u_B$ defined as in Definition~\ref{sha-lifts-graded}(e), lies in the Bloch-Kato subspace \cite{BK90}.
\item (Crystalline points.) Let $E'$ be a finite extension of $\mathbf Q_\ell$ that is an $\mathcal O$-algebra, and let $h:R_B\to E'$ kill $\mathfrak a_B$. Then the matrix parts of $h\circ u_B$ on the open subgroup $H_B$ of $D_B$ form a crystalline representation \cite{Fon94}.
\end{enumerate}
\end{lemma}
\begin{proof}
(c) Let $h:R_B\to L$ kill $\mathfrak a_B$, and write $M_s$ also for the matrix part of $h\circ u_B(s)$. For $s,s_1,s_2,s_3\in I_B$:
\begin{enumerate}
\item[(i)] $\det(X-M_s)=(X-\Psi(s))(X-1)(X-\Psi(s)^{-1})$;
\item[(ii)] $(M_{s_1}-\Psi(s_1))(M_{s_2}-1)(M_{s_3}-\Psi(s_3)^{-1})=0$.
\end{enumerate}
Suppose first that $L$ has characteristic $0$. Since $\psi_\ell(s_0)=1$, $\Psi(s_0)=\chi_{\mathrm{cyc}}(s_0)^k$ has infinite order, so $d=(\Psi(s_0),1,\Psi(s_0)^{-1})$ has distinct entries in $L$. By (i) there is $P\in GL_3(L)$ with $P^{-1}M_{s_0}P=\Delta=\operatorname{diag}(d_1,d_2,d_3)$.
Let $s\in I_B$ and $N=P^{-1}M_sP$. By (ii) for $(s,s_0,s_0)$, $(s_0,s,s_0)$ and $(s_0,s_0,s)$,
\[
 (N-\Psi(s))(\Delta-d_2)(\Delta-d_3)=0,\qquad(\Delta-d_1)(N-1)(\Delta-d_3)=0,\qquad(\Delta-d_1)(\Delta-d_2)(N-\Psi(s)^{-1})=0.
\]
Since the $d_a$ are distinct, the first identity says that the first column of $N-\Psi(s)$ vanishes, the second that the entries $(a,b)$ of $N-1$ with $a\ne1$ and $b\ne3$ vanish, and the third that the last row of $N-\Psi(s)^{-1}$ vanishes. So $N$ is upper triangular with diagonal $\Psi(s),1,\Psi(s)^{-1}$.
Now let $g\in D_B\cap\ker\eta$ and $k=P^{-1}M_gP$. Then $gs_0g^{-1}\in I_B$ and $\Psi(gs_0g^{-1})=\Psi(s_0)$, so $U=P^{-1}M_{gs_0g^{-1}}P$ is upper triangular with diagonal $d_1,d_2,d_3$, and $Uk=k\Delta$. Comparing the entries $(3,1)$, $(3,2)$ and $(2,1)$ of the two sides, in this order, gives $(d_3-d_1)k_{31}=(d_3-d_2)k_{32}=(d_2-d_1)k_{21}=0$, so $k$ is upper triangular. This proves (c) for $L$ of characteristic $0$, with no extension of $L$.

In general, $R_B/\mathfrak a_B$ is a domain in which $\ell\ne0$, so its fraction field $F$ has characteristic $0$, and the case just treated gives $P\in GL_3(F)$ that puts the images of the matrix parts of $u_B$ on $D_B\cap\ker\eta$ in this triangular form. Let $\mathfrak p$ be the kernel of the map $R_B/\mathfrak a_B\to L$ induced by $h$, and $V$ a valuation ring of $F$ that dominates the local ring of $R_B/\mathfrak a_B$ at $\mathfrak p$.
Multiplying $P$ on the left by matrices of determinant $1$ over $V$ that clear its entries $(3,1)$, $(2,1)$ and $(3,2)$ in turn gives $k\in GL_3(V)$ with $kP$ upper triangular: for two entries $a,b$ of a column, $a/b$ or $b/a$ lies in $V$. Then $kM_gk^{-1}=(kP)(P^{-1}M_gP)(kP)^{-1}$ is upper triangular with the same diagonal as $P^{-1}M_gP$, and $k$ and $k^{-1}$ have entries in $V$. Reducing modulo the maximal ideal of $V$ gives the triangular form over the residue field $\kappa$ of $V$, for the map $R_B/\mathfrak a_B\to\kappa$, whose kernel is $\mathfrak p$.
The fields $\kappa$ and $L$ both contain the domain $(R_B/\mathfrak a_B)/\mathfrak p$, so their tensor product over it is a nonzero ring. Its $L$-subalgebra generated by the entries of the images of $k$ and $k^{-1}$ is finitely generated over $L$, so the residue field of any of its maximal ideals is finite over $L$, by Zariski's lemma \cite[Corollary~5.24]{AM69}, and the triangular form holds over it.

Let $I_B$, $\mathfrak c_B$ and $\mathfrak m$ be as in Lemma~\ref{sha-model-local-component}.

\emph{The domain.} By Lemma~\ref{sha-model-local-component}(1), $\mathfrak a_B$ is prime and $\ell\notin\mathfrak a_B$; so $R_B/\mathfrak a_B$ is a domain in which $\ell\ne0$. As $\mathfrak a_B$ is a minimal prime of $\mathfrak c_B$, $I_B\subseteq\mathfrak c_B\subseteq\mathfrak a_B$.

\emph{The identities (i) and (ii).} Let $Z\subseteq R_B$ be the set of the coefficients of $\det(X-M_s)-(X-\Psi(s))(X-1)(X-\Psi(s)^{-1})$ and of the entries of $(M_{s_1}-\Psi(s_1))(M_{s_2}-1)(M_{s_3}-\Psi(s_3)^{-1})$, for $s,s_1,s_2,s_3\in I_B$. For $h$ as in (c), (i) and (ii) say that $h$ kills $Z$; so it suffices to show that $Z\subseteq\mathfrak a_B$. Let $\varphi$ be as in Lemma~\ref{sha-model-local-component}(2), let $f\in Z$ and $n\ge0$. As $\ker\varphi=\mathfrak a_B$ and $\bigcap_n\mathfrak m^{n+1}=0$, it suffices to show that $\varphi(f)\in\mathfrak m^{n+1}$.
\begin{enumerate}
\item[(1)] The ring $C_n=E[[Y_1,\dots,Y_{12}]]/\mathfrak m^{n+1}$ is a finite $E$-algebra, hence a finite $\mathbf Q_\ell$-algebra, and the constant term $\pi:C_n\to E$ is a homomorphism of $E$-algebras whose kernel consists of nilpotent elements. Let $g_n:R_B\to C_n$ be $\varphi$ followed by the quotient map; it kills $\mathfrak a_B$, hence $I_B$.
\item[(2)] By Lemma~\ref{sha-model-local-component}(3), $g_n\circ u_B$ has continuous entries and its matrix parts on $H_B$ form a crystalline representation. It has multiplier $\eta$, and $\pi\circ g_n\circ u_B=r_0$, as $\pi\circ g_n$ is $x_B$ followed by $\mathcal O\subset E$.
\item[(3)] By Lemma~\ref{sha-model-local-component}(4), $g_n\circ u_B=nun^{-1}$ with $u$ upper triangular, with diagonal entries $\Psi(s),1,\Psi(s)^{-1}$ at $s\in I_B$. For $s\in I_B$, let $U_s$ be the matrix part of $u(s)$. Its characteristic polynomial is $(X-\Psi(s))(X-1)(X-\Psi(s)^{-1})$. The matrices $U_{s_1}-\Psi(s_1)$, $U_{s_2}-1$ and $U_{s_3}-\Psi(s_3)^{-1}$ are upper triangular with a zero diagonal entry in positions $1$, $2$ and $3$ respectively; so the product of the first two has its first two columns zero, and its product with the third is zero.
\item[(4)] As $\eta$ is trivial on $I_B$, the matrix part of $g_n(u_B(s))$ is $nU_sn^{-1}$ for $s\in I_B$, and both identities of (3) are preserved by this conjugation. So $g_n$ kills $Z$, that is, $\varphi(f)\in\mathfrak m^{n+1}$.
\end{enumerate}

(e) Let $h:R_B\to E'$ kill $\mathfrak a_B$. Then $h$ kills $I_B$, and Lemma~\ref{sha-model-local-component}(3), with $C=E'$, is the claim.

(b) Let $h:R_B\to A$ and $P$ be as in (b), let $q:R_B\to\bar R=R_B/\mathfrak a_B$ be the quotient map, and let $S=\bar R[[X_{ab}]]$, a power series ring in nine variables $X_{ab}$, $1\le a,b\le3$; it is a domain and a complete local noetherian $\mathcal O$-algebra with residue field $\mathbf F$.
\begin{enumerate}
\item[(1)] Let $i:R_B\to S$ be $q$ followed by $\bar R\subset S$, and $U=1+(X_{ab})\in GL_3(S)$. The homomorphism $U(i\circ u_B)U^{-1}$ is a continuous lift of the reduction of $r_0$ with multiplier $\eta$, as $U\equiv1$ modulo $\mathfrak m_S$; so by universality there is $\alpha:R_B\to S$ with $\alpha\circ u_B=U(i\circ u_B)U^{-1}$.
\item[(2)] The homomorphism $i$ kills $I_B\subseteq\mathfrak a_B$, and $i(\delta_B)\ne0$, as $\delta_B$ lies in no minimal prime of $\mathfrak c_B$, and $\mathfrak a_B$ is one. By Lemma~\ref{sha-model-local-component}(5), $\alpha$ kills $\mathfrak c_B$.
\item[(3)] Let $c:S\to\bar R$ set every $X_{ab}$ to $0$. Then $c\circ\alpha$ and $q$ carry $u_B$ to $q\circ u_B$, so $c\circ\alpha=q$ by universality, and $\ker\alpha\subseteq\mathfrak a_B$. The prime $\ker\alpha$ contains $\mathfrak c_B$ and lies in $\mathfrak a_B$, a minimal prime of $\mathfrak c_B$; so $\ker\alpha=\mathfrak a_B$.
\item[(4)] Let $F:S\to A$ be the homomorphism of $\mathcal O$-algebras that is the map $\bar R\to A$ induced by $h$ on $\bar R$ and sends $X_{ab}$ to the entry $(a,b)$ of $P-1$, which lies in $\mathfrak m_A$. Then $F\circ\alpha$ carries $u_B$ to $P(h\circ u_B)P^{-1}$, so $h'=F\circ\alpha$ by universality, and $h'$ kills $\mathfrak a_B$ by (3).
\end{enumerate}

(a) Let $\varphi_0$ be the homomorphism of Lemma~\ref{sha-model-local-component}(2).
\begin{enumerate}
\item[(1)] \emph{Uniqueness.} Two homomorphisms of $\mathcal O$-algebras from $R_B$ to a complete local noetherian $\mathcal O$-algebra with residue field $\mathbf F$ that agree at the entries of the matrix parts of $u_B$ are equal, by universality.
\item[(2)] \emph{A torus action.} Let $t\in1+\mathfrak m_{\mathcal O}$ have infinite order, and $d=\operatorname{diag}(t,1,t^{-1})$, so that conjugation by $d$ multiplies the entry $(a,b)$ by $t^{b-a}$. As $t\equiv1$, $d\,u_Bd^{-1}$ is a continuous lift of the reduction of $r_0$ with multiplier $\eta$, and universality gives an endomorphism $\theta$ of $R_B$ with $\theta\circ u_B=d\,u_Bd^{-1}$. Then $x_B\circ\theta=x_B$ by (1), as the matrix parts of $r_0$ are diagonal. And $\theta(\mathfrak a_B)\subseteq\mathfrak a_B$: by (b), applied to the quotient map $R_B\to R_B/\mathfrak a_B$ and $P=d$, the homomorphism carrying $u_B$ to $d\,u_Bd^{-1}$ modulo $\mathfrak a_B$ kills $\mathfrak a_B$, and by (1) it is $\theta$ followed by the quotient map.
\item[(3)] \emph{Spanning.} For $f\in R_B$ let $\lambda(f)\in E^{12}$ be the vector of the coefficients of $Y_1,\dots,Y_{12}$ in $\varphi_0(f)$. The $\lambda(e)$, for $e$ running over the entries of the matrix parts of $u_B$, span $E^{12}$. Otherwise a nonzero linear form $\mu$ on $E^{12}$ kills them, and $\delta=\mu\circ\lambda$ satisfies $\delta(fg)=x_B(f)\delta(g)+x_B(g)\delta(f)$ and vanishes on $\mathcal O$. As $\ker x_B$ is finitely generated and $\delta(gk)=x_B(g)\delta(k)$ for $k\in\ker x_B$, some $c\ne0$ in $\mathcal O$ has $c\delta(R_B)\subseteq\mathcal O$; then $f\mapsto x_B(f)+\epsilon c\delta(f)$ is a homomorphism $R_B\to\mathcal O[\epsilon]/(\epsilon^2)$ that agrees with $x_B$ at the entries, so $\delta=0$ by (1). But by Lemma~\ref{sha-model-local-component}(2) with $N=2$ and $g=Y_j$, some $f$ and $c'\ne0$ have $\lambda(f)=c'e_j$, where $e_j$ is the $j$-th unit vector; so $\mu(e_j)=0$ for all $j$, a contradiction.
\item[(4)] \emph{Coordinates.} Choose entries $e_1,\dots,e_{12}$ with $\lambda(e_1),\dots,\lambda(e_{12})$ a basis, and put $z_i=e_i-x_B(e_i)$, of weight $w(i)=b-a$ if $e_i$ is an entry $(a,b)$. The $\varphi_0(z_i)$ have no constant term and linearly independent linear parts, so there is an automorphism $\sigma$ of the $E$-algebra $E[[Y_1,\dots,Y_{12}]]$ with $\sigma(Y_i)=\varphi_0(z_i)$; it preserves the powers of $\mathfrak m$. Put $\varphi_B=\sigma^{-1}\circ\varphi_0$. It has kernel $\mathfrak a_B$, the properties of $\varphi$ in Lemma~\ref{sha-model-local-component}(2), and $\varphi_B(z_i)=Y_i$.
\item[(5)] \emph{Grading.} Let $\varrho$ be the automorphism of $E[[Y_1,\dots,Y_{12}]]$ with $\varrho(Y_i)=t^{w(i)}Y_i$. The homomorphisms $\varphi_B\circ\theta$ and $\varrho\circ\varphi_B$ have constant terms $x_B$, kill $\mathfrak a_B$, and agree at the $z_i$, as $\theta(z_i)=t^{w(i)}z_i$. They are equal: given $f$ and $N$, the density in Lemma~\ref{sha-model-local-component}(2) gives $c\ne0$ and a polynomial $p$ over $\mathcal O$ with $\varphi_B(cf-p(z))\in\mathfrak m^N$; by the injectivity there, $c'(cf-p(z))\in(\ker x_B)^N+\mathfrak a_B$ for some $c'\ne0$; both homomorphisms agree at $p(z)$, vanish on $\mathfrak a_B$, and map $(\ker x_B)^N$ into $\mathfrak m^N$; so $c'c$ times the difference of their values at $f$ lies in $\mathfrak m^N$ for every $N$. Now if $\theta(f)=t^jf$, then $\varrho(\varphi_B(f))=t^j\varphi_B(f)$, and comparing the coefficients of each monomial, which $\varrho$ multiplies by $t$ to the power of its weight, shows that $\varphi_B(f)$ has weight $j$, as $t$ has infinite order. This applies to every entry $(a,b)$ of a matrix part of $u_B$, with $j=b-a$.
\item[(6)] \emph{Counting.} By Lemma~\ref{sha-model-local-component}(7), applied to $\varphi_B$, $w$ and the $e_i$, the function $w$ takes the values $-2,-1,0,1,2$ on at most $1,2,3,4,2$ elements and no other value. As $1+2+3+4+2=12$, all five bounds are attained.
\item[(7)] \emph{Conclusion.} The constant term of $\varphi_B(f)$ is $x_B(f)$; $Y_i$ is the only linear term of $\varphi_B(z_i)$; and if $x_B(f)=0$ and $\varphi_B(f)$ has no linear term, then $\varphi_B(f)\in\mathfrak m^2$, as its constant term is $0$, so a nonzero multiple of $f$ lies in $(\ker x_B)^2+\mathfrak a_B$ by the injectivity in Lemma~\ref{sha-model-local-component}(2).
\end{enumerate}

(d) Let $h$ be as in (d), with $h_0=h\bmod\epsilon$.
\begin{enumerate}
\item[(1)] \emph{A finite extension.} Every element of $h_0(R_B)$ is integral over $\mathcal O$: it is $h_0(c)+h_0(f)$ with $c\in\mathcal O$ and $f\in\mathfrak m_{R_B}$, and $h_0(f)$ has absolute value less than $1$. So $h_0(R_B)$ is a local noetherian ring, integral over $\mathcal O$, whose residue field is that of $\mathcal O$; by Nakayama's lemma it is a finite $\mathcal O$-module. Choose $b_1,\dots,b_m\in R_B$ whose images generate it, generators $p_1,\dots,p_q$ of the ideal $\ker h_0$, and generators of $\mathcal O$ over $\mathbf Z_\ell$, and let $K\subset\bar{\mathbf Q}_\ell$ be the extension of $\mathbf Q_\ell$ generated by the images of the latter, the $h_0(b_i)$, and the coefficients of $\epsilon$ in the $h(b_i)$ and the $h(p_k)$. It is finite over $\mathbf Q_\ell$. Every $f\in R_B$ is $\sum_iu_ib_i+\sum_kg_kp_k$ with $u_i\in\mathcal O$ and $g_k\in R_B$, and $h(g_kp_k)=h_0(g_k)h(p_k)$ as $h_0(p_k)=0$; so $h$ takes values in $K[\epsilon]/(\epsilon^2)$.
\item[(2)] \emph{The cocycle.} Viewed over $K[\epsilon]/(\epsilon^2)$, a finite $\mathbf Q_\ell$-algebra, $h$ kills $I_B$, so by Lemma~\ref{sha-model-local-component}(3) $h\circ u_B$ has continuous entries and its matrix parts on $H_B$ form a crystalline representation. The map $\gamma(s)=M'_sM_s^{-1}$ is continuous, and multiplying out $(M_{st}+\epsilon M'_{st})=(M_s+\epsilon M'_s)\cdot(M_t+\epsilon M'_t)$, with the signs of $u_B$, shows that it is a cocycle for the adjoint action of $h_0\circ u_B$. By Lemma~\ref{sha-model-local-component}(6), applied over $K$ to $\rho=h_0\circ u_B$ and $\rho'=h\circ u_B$, $\gamma$ lies in the Bloch-Kato subspace over $K$.
\item[(3)] \emph{Extension of scalars.} The inclusion $K\subset\bar{\mathbf Q}_\ell$ is continuous and compatible with the adjoint actions; it carries an $X\in M_3(K)\otimes B_{\mathrm{cris}}$ with $\gamma(s)\otimes1=sX-X$ to such an element over $\bar{\mathbf Q}_\ell$. So $\gamma$ lies in the Bloch-Kato subspace over $\bar{\mathbf Q}_\ell$.
\end{enumerate}
% (c) is Loc.lemma2c: lemma2b (exists_inertia_distinct_characters, diagonalize_three_by_three_split, three_eigenvalue_product_zero, commuting_upper_triangular_distinct) over the fraction field of D.Rt U, then triangular_form_specialization (iwasawa_triangularization, triangular_form_borel_change, exists_common_ring_same_kernel, triangular_form_finite_extension). The identities (i) and (ii) are Loc.lemma2a, which Alpöge proves on all of the inertia group through rtCompl, rtPoints and ordEq (crystalline ordinary representations are triangular); the domain and l are Loc.rtBasic.
% The domain is Loc.rtBasic (OrdLoc.component_basic). The identities are Loc.lemma2a (ScalarMap.lean): Z is the set S of the private cond_eqs; (1) is trunc_data; (2) is rtPoints; (3) is ordEq, triDiag and the private cond_tri; (4) is the private cond_conj with eta_unr. (e) is Loc.rtPoints (OrdLoc.rtPoints: Icr <= It, then PtsCr). (b) is Loc.rtConj = OrdLoc.rtConj (ExistsConjLift.lean): (1) is exists_univMat, exists_conjLift and U.rep; (2) is saturatingElement_notMem and reducedOrdinaryIdeal_le_ker_of_conj; (3) is co_ev at 0 and the minimality of It; (4) is co_ev with gl_ext. (a) is Loc.rtChart: (1) is the uniq of torAut, (2) is Loc.torAut (exists_unit, conjBy_tor), (3) is Loc.chartSpan (with exists_den and der_zero), (4) is Loc.chartChange (exists_equiv), (5) is Loc.chartGraded (cg_ext, rescale), (6) is Loc.chartCount and count_weights_eq_local_dims, (7) is IsCompl.const and IsCompl.inj with N = 2. (d) is Loc.rtDual: (1) is Loc.ptFin (PtFinFinite.lean, with ptFin_int, ptFin_finite by Nakayama), (2) is the private fin_case (rtPoints, tc_mul, dualCr), (3) is the private map_mem_finiteCocycles_adRep_of_continuous. What remains open is the open step of Lemma sha-model-local-component.
\end{proof}
\leanref{Development/AutomorphicArithmetic/AutomorphicVarieties/PadicHodgeTheory/IsInertiaInfiniteCyclotomic.lean}{UpperTheory.V.Loc.Dat3.It}
\leanref{Development/AutomorphicArithmetic/AutomorphicVarieties/PadicHodgeTheory/LocGood.lean}{UpperTheory.V.Loc.locGood}
\leanref{Development/AutomorphicArithmetic/AutomorphicVarieties/PadicHodgeTheory/RtChart.lean}{UpperTheory.V.Loc.rtChart}
\leanref{Development/AutomorphicArithmetic/AutomorphicVarieties/PadicHodgeTheory/RtConj.lean}{UpperTheory.V.Loc.rtConj}
\leanref{Development/AutomorphicArithmetic/AutomorphicVarieties/PadicHodgeTheory/ValuationRingClearLastRowEntry.lean}{UpperTheory.V.Loc.lemma2c}
\leanref{Development/AutomorphicArithmetic/AutomorphicVarieties/PadicHodgeTheory/RtDual.lean}{UpperTheory.V.Loc.rtDual}
\leanref{Development/AutomorphicArithmetic/AutomorphicVarieties/PadicHodgeTheory/RtPoints.lean}{UpperTheory.V.Loc.rtPoints}
\leanref{Development/AutomorphicArithmetic/AutomorphicVarieties/PadicHodgeTheory/CommutingUpperTriangularDistinct.lean}{UpperTheory.V.Loc.lemma2b}

\begin{lemma}\label{sha-model-universal}
In Lemma~\ref{sha-model-local}, let $\psi$ be modulo $N$, let $g_1\in\operatorname{Gal}(\bar{\mathbf Q}/K)$ with $\Psi(g_1)^2-1\in\mathcal O^\times$, and let $\mathcal N$ be the set of the elements $s$ of the inertia groups at the primes $p\ne\ell$ with $\eta(s)=1$ and $\bar\Psi(s)=1$, where $\bar\Psi$ is the reduction of $\Psi$.
For a complete noetherian local $\mathcal O$-algebra $A$ with residue field $\mathbf F$, an \emph{$\mathfrak a_B$-lift} with values in $A$ is a homomorphism $r:\operatorname{Gal}(\bar{\mathbf Q}/\mathbf Q)\to\mathcal G_3(A)$ with multiplier $\eta$ that reduces to $r_0$ modulo $\mathfrak m_A$, is continuous for the $\mathfrak m_A$-adic topology, is normalized as in Definition~\ref{sha-ordinary-deformation}(c), satisfies $r(s)=1$ for $s\in\mathcal N$, and whose restriction to $D_B$ is $h\circ u_B$ for a homomorphism $h:R_B\to A$ that kills $\mathfrak a_B$.
\begin{enumerate}
\item There is an $\mathfrak a_B$-lift $r^{\mathrm u}$ with values in a complete noetherian local $\mathcal O$-algebra $R$ with residue field $\mathbf F$ that is universal: for every $\mathfrak a_B$-lift $r$ with values in $A$ there is a unique homomorphism of $\mathcal O$-algebras $f:R\to A$ with $f\circ r^{\mathrm u}=r$.
\item Let $A$ be a complete local domain over $\mathcal O$ and $f:R\to A$ a local homomorphism of $\mathcal O$-algebras. Then $f\circ r^{\mathrm u}$ is an ordinary point relative to $g_1$ and $B$.
\end{enumerate}
\end{lemma}
\begin{proof}
(a) Let $S$ be the set of primes dividing $\ell Nd_K$.
At a prime $p\notin S$ the characters $\eta$ and $\Psi$ are unramified, so $\mathcal N$ contains the inertia groups at $p$, and the quotient of $\operatorname{Gal}(\bar{\mathbf Q}/\mathbf Q)$ by the closed normal subgroup generated by $\mathcal N$ is a quotient of the Galois group of the maximal extension of $\mathbf Q$ unramified outside $S$.
That Galois group, and hence its quotient, satisfies Mazur's finiteness condition at $\ell$ \cite{Maz89}, by the finiteness of its cohomology with finite coefficients \cite[Chapter VIII, \S3]{NSW08}.
As in the proof of Lemma~\ref{sha-lifts-universal}, the continuous lifts of the reduction of $r_0$ to $\mathcal G_3$ with multiplier $\eta$ that are trivial on $\mathcal N$ are classified by a complete noetherian local $\mathcal O$-algebra $R^\square$ with residue field $\mathbf F$ and a universal lift $r^\square$.
The restriction of $r^\square$ to $D_B$ is $h^\square\circ u_B$ for a unique $h^\square:R_B\to R^\square$.
Let $R$ be the quotient of $R^\square$ by the ideal generated by $h^\square(\mathfrak a_B)$, by the entries of the matrix part of $r^\square(c)$ minus $1$, and by the entries $(1,2)$ and $(2,1)$ of the matrix part of $r^\square(g_1)-r^\square(cg_1^{-1}c)$, and let $r^{\mathrm u}$ be the image of $r^\square$.
A lift $r$ trivial on $\mathcal N$ is $f\circ r^\square$ for a unique $f$.
It is normalized if and only if $f$ kills these entries.
Its restriction to $D_B$ is $(f\circ h^\square)\circ u_B$, so, by the uniqueness in the universal property of $R_B$, it is $h\circ u_B$ for some $h$ killing $\mathfrak a_B$ if and only if $f$ kills $h^\square(\mathfrak a_B)$.

(b) Put $r=f\circ r^{\mathrm u}$.
The multiplier, the reduction, the continuity and the normalization hold by definition, and so does the condition away from $\ell$, since an element $s$ of an inertia group at $p\ne\ell$ with $\eta(s)=\Psi(s)=1$ lies in $\mathcal N$.
At $B$, the restriction of $r$ to $D_B$ is $h\circ u_B$ with $h$ killing $\mathfrak a_B$.
By Lemma~\ref{sha-model-local}(c), applied to the fraction field of $A$, the matrix parts of $r$ on $D_B\cap\ker\eta$ stabilize a full flag over a finite extension of that field, and inertia acts on its graded pieces by $\Psi,1,\Psi^{-1}$.
Since $\Psi=\chi_{\mathrm{cyc}}^k\psi_\ell$ on $I_B$ with $\psi_\ell^2=1$, and $\ell-1\nmid2k$, there is $x\in I_B\cap\ker\eta$ with $\Psi(x)^2-1\in\mathcal O^\times$.
The values $\Psi(x)$, $1$ and $\Psi(x)^{-1}$ are then pairwise distinct modulo $\mathfrak m_{\mathcal O}$, so the matrix part of $r(x)$, which reduces to $\operatorname{diag}(\bar\Psi(x),1,\bar\Psi(x)^{-1})$, has eigenvectors in $A^3$ congruent to the standard basis vectors, by Hensel's lemma.
The lines of the flag are spanned by eigenvectors of $r(x)$, so the matrix $Q\in GL_3(A)$ whose columns are these eigenvectors, which is congruent to $1$ modulo $\mathfrak m_A$, conjugates the matrix parts of $r$ on $D_B\cap\ker\eta$ to upper triangular matrices with diagonal entries $\Psi,1,\Psi^{-1}$ on $I_B\cap\ker\eta$.
\end{proof}
\leanref{Development/AutomorphicArithmetic/AutomorphicVarieties/UnitaryThreeGroups/RingEx.lean}{UpperTheory.V.U3.exists_ring}
\leanref{Development/AutomorphicArithmetic/AutomorphicVarieties/UnitaryThreeGroups/UnivQuot.lean}{UpperTheory.V.U3.exists_sRing}
\leanref{Development/AutomorphicArithmetic/AutomorphicVarieties/UnitaryThreeGroups/RhoKApply.lean}{UpperTheory.V.U3.Dat.Ring}
\leanref{Development/AutomorphicArithmetic/AutomorphicVarieties/CurveDVRSpecialization/TriangularOfReducibleRankThree.lean}{UpperTheory.V.Pr.ordPt_of_gd}

\begin{lemma}\label{sha-model-points}
In Lemma~\ref{sha-model-universal}, let $r^{\mathrm u}$ be the universal $\mathfrak a_B$-lift, with values in $R$.
Give the entry $(a,b)$ of a $3\times3$ matrix the weight $b-a$, and let $R^\natural_0$ be the $\mathcal O$-subalgebra of $R$ generated by the products of entries of the matrix parts of the $r^{\mathrm u}(g)$ whose weights add up to $0$.
A \emph{point} of $R$ is a homomorphism of $\mathcal O$-algebras $f:R\to\bar{\mathbf Q}_\ell$ with $|f(y)|<1$ for $y\in\mathfrak m_R$; its representation is $r_f=f\circ r^{\mathrm u}$, and $f$ is \emph{two-signed} if the matrix parts of $r_f$ on $\ker\eta$ are neither all upper triangular nor all lower triangular.
\begin{enumerate}
\item Suppose that there is a subgroup $U$ of finite index of $\operatorname{Gal}(\bar{\mathbf Q}/\mathbf Q)$ such that, for every complete discrete valuation ring $A$ over $\mathcal O$ with finite residue field in which $\mathfrak m_{\mathcal O}$ maps to $0$ and every ordinary point $r$ relative to $g_1$ and $B$ with values in $A$, the characteristic polynomials of $r(g)$ and $r_0(g)$ agree for $g\in U\cap\ker\eta$. Then $R^\natural_0$ is finite over $\mathcal O$.
\item Let $f$ be a point. If the matrix parts of $r_f$ on $\ker\eta$ are upper triangular and lower triangular, then $r_f=r_0$.
\item Let $f$ be two-signed. Then either the matrix parts of $r_f$ on $\ker\eta$ have no common invariant subspace other than $0$ and $\bar{\mathbf Q}_\ell^3$; or on $\ker\eta$ the entries $(1,2)$, $(2,1)$, $(2,3)$ and $(3,2)$ of $r_f$ vanish, the middle entry is $1$, and the corner block formed by rows and columns $1$ and $3$ is irreducible and has determinant $1$. In the second case there is a continuous $W:\operatorname{Gal}(\bar{\mathbf Q}/\mathbf Q)\to GL_2(\bar{\mathbf Q}_\ell)$ that agrees with the corner block on $\ker\eta$ and has $\operatorname{tr}W(c)=0$ and $\det W=\eta$.
\item Let $f$ be two-signed and $f'$ a point, and suppose that the matrix parts of $r_f$ and $r_{f'}$ on $\ker\eta$ are conjugate by an element of $GL_3(\bar{\mathbf Q}_\ell)$. Then $r_{f'}=\tau(t)r_f\tau(t)^{-1}$ for some $t\in\bar{\mathbf Q}_\ell^\times$, where $\tau(t)=\operatorname{diag}(t,1,t^{-1})$.
\end{enumerate}
\end{lemma}
\begin{proof}
For a homomorphism $r$ with multiplier $\eta$ that is normalized as in Definition~\ref{sha-ordinary-deformation}(c), let $\zeta(r)$ be the entry $(1,1)$ of the matrix part of $r(g_1)-r(cg_1^{-1}c)$; by the normalization, that matrix is $\operatorname{diag}(\zeta(r),0,-\zeta(r))$.
For $r_0$ it is $\Psi(g_1)-\Psi(g_1)^{-1}$, a unit.
For a point $f$, $\zeta(r_f)=f(\zeta(r^{\mathrm u}))$ is congruent to this unit, so it is nonzero; and $\zeta(r_{f'})\ne-\zeta(r_f)$ for any two points $f,f'$, since $2$ is a unit.
Let $\mathfrak m$ be the maximal ideal of the integers of $\bar{\mathbf Q}_\ell$.
Three facts about a point $f$ are used in (b) and (c).
First, for $s\in I_B\cap\ker\eta$ the characteristic polynomial of the matrix part of $r_f(s)$ is $(X-\Psi(s))(X-1)(X-\Psi(s)^{-1})$, by Lemma~\ref{sha-model-local}(c) applied to the homomorphism $R_B\to\bar{\mathbf Q}_\ell$ through which $r_f$ factors on $D_B$; the same holds on the conjugates of $I_B$, the inertia groups at the other places above $\ell$.
Second, at a prime $p\ne\ell$ the elements of the inertia group that lie in $\mathcal N$ form a normal subgroup of index prime to $\ell$, on which $r_f$ is trivial; so on the elements of the inertia group in $\ker\eta$ the matrix parts of $r_f$ factor through a finite group of order prime to $\ell$, and they are conjugate to those of $r_0$, since both have integral models with the same reduction.
Third, a continuous character of $\operatorname{Gal}(\bar{\mathbf Q}/K)=\ker\eta$ with values in $1+\mathfrak m$ that is trivial on every inertia group factors through the class group of $K$; its image is then finite of order prime to $\ell$, because $\ell\nmid h_K$, so it is trivial, since $1+\mathfrak m$ has no torsion of order prime to $\ell$.
Finally, let $x\in I_B\cap\ker\eta$ with $\Psi(x)^2-1\in\mathcal O^\times$, as in the proof of Lemma~\ref{sha-model-universal}(b).

(a) As in Lemma~\ref{sha-lifts-presentation}(a), with Lemma~\ref{sha-model-local}(b) for the action of $\tau(t)$ when $t\equiv1$, there is a surjection $\pi:\mathcal O[[X_1,\dots,X_n]]\to R$ with homogeneous kernel, where $X_i$ has a weight and maps to an entry of that weight of the matrix part of some $r^{\mathrm u}(g)$ minus its value at $r_0$; and the entry $(a,b)$ of every $r^{\mathrm u}(g)$ is the image of a power series of weight $b-a$.
Let $P_0\subset R$ be the image of the power series of weight $0$, a complete local noetherian $\mathcal O$-algebra; it contains $R^\natural_0$, so it suffices to show that $P_0$ is finite over $\mathcal O$.
Let $A$ be a complete discrete valuation ring with finite residue field in which $\mathfrak m_{\mathcal O}$ maps to $0$, and $f:R\to A$ a local homomorphism. Such an $f$ kills the $x_i=\pi(X_i)$ of weight $0$, and either all those of negative weight or all those of positive weight.
By Lemma~\ref{sha-model-universal}(b), $r=f\circ r^{\mathrm u}$ is an ordinary point.
Let $g\in\ker\eta$, and $n\ge1$ with $g^n\in U$.
The characteristic polynomial of $r(g)$ divides the polynomial obtained from that of $r(g^n)$ by substituting $X^n$ for $X$, whose coefficients lie in the copy of $\mathbf F$ in $A$; so the coefficients of the characteristic polynomial of $r(g)$ are algebraic over $\mathbf F$, and they lie in $\mathbf F$, since an element of $A$ algebraic over the finite field $\mathbf F$ is determined by its residue.
Their residues are those of $r_0(g)$, so $r(g)$ and $r_0(g)$ have the same characteristic polynomial for every $g\in\ker\eta$.
Over the fraction field of $A$, $r$ is normalized and $\zeta(r)$ is a unit, so, by the analysis in the proof of Lemma~\ref{sha-lifts-points}(d), on $\ker\eta$ it is irreducible, or upper triangular, or lower triangular, or block shaped: its entries of odd weight vanish and its corner block is irreducible.
It is not irreducible, by the Brauer-Nesbitt theorem, since its characteristic polynomials on $\ker\eta$ are those of $\operatorname{diag}(\Psi,1,\Psi^{-1})$.
If it were block shaped, its middle entry, a root of the characteristic polynomial congruent to $1$, would be $1$, so the characteristic polynomials of the corner block would be those of $\operatorname{diag}(\Psi,\Psi^{-1})$, and the Brauer-Nesbitt theorem would contradict the irreducibility of the corner block.
So $r$ is upper or lower triangular on $\ker\eta$, and its diagonal entries, roots of the characteristic polynomial congruent to those of $r_0$, are equal to those of $r_0$.
This proves the claim on $f$.

Let $\varpi$ be a uniformizer of $\mathcal O$.
The maximal ideal of $P_0$ is generated by $\varpi$ and by finitely many monomials $b$ of weight $0$ in the $x_i$ without constant term, and each such $b$ contains a variable of weight $0$ or variables of both signs.
If a prime $Q\supset\varpi R$ did not contain $b$, there would be a prime $\mathfrak p\supset Q$ with $R/\mathfrak p$ of dimension one and $b\notin\mathfrak p$, and the normalization of $R/\mathfrak p$ would give a local homomorphism $f$ as above with $f(b)\ne0$, which the claim excludes.
So every such $b$ lies in the radical of $\varpi R$, and, taking components of weight $0$, some power of $b$ lies in $\varpi P_0$.
Since the maximal ideal of $P_0$ is generated by $\varpi$ and these $b$, a power of it lies in $\varpi P_0$; hence $P_0/\varpi P_0$ is finite, and $P_0$, which is complete and separated for the $\varpi$-adic topology, is finite over $\mathcal O$.

(b) On $\ker\eta$ the matrix parts of $r_f$ are diagonal, and their diagonal entries are characters $a_1,a_2,a_3$; with $(w_1,w_2,w_3)=(1,0,-1)$, the values of $a_i\Psi^{-w_i}$ lie in $1+\mathfrak m$, since $r_f$ reduces to $r_0$.
Let $s\in I_B\cap\ker\eta$.
By the first fact, $a_i(s)=\Psi(s)^{w_j}$ for some $j$.
If $\Psi(s)^2-1$ is a unit, then $j=i$, since otherwise $\Psi(s)^{w_j-w_i}\in1+\mathfrak m$ with $w_j-w_i$ dividing $2$.
Otherwise $\Psi(sx)^2-1$ is a unit, and $a_i(s)=a_i(sx)a_i(x)^{-1}=\Psi(s)^{w_i}$.
So $a_i=\Psi^{w_i}$ on $I_B\cap\ker\eta$, and in the same way on its conjugates.
Let $s$ lie in the inertia group at a prime $p\ne\ell$, with $\eta(s)=1$.
By the second fact, $a_i(s)=\Psi(s)^{w_j}$ for some $j$; then $s'=s^{w_j-w_i}$ has $\Psi(s')\in1+\mathfrak m$, so $s'\in\mathcal N$ and $\Psi(s')=1$, because $\Psi=\psi$ on this inertia group and $\psi$ has order prime to $\ell$; so $a_i(s)=\Psi(s)^{w_i}$.
By the third fact, $a_i=\Psi^{w_i}$ on $\ker\eta$.
So $r_f=r_0$ on $\ker\eta$; the matrix part of $r_f(c)$ is $1$ by the normalization, and so $r_f=r_0$.

(c) By the analysis in the proof of Lemma~\ref{sha-lifts-points}(d), since $\zeta(r_f)\ne0$, $r_f$ is irreducible on $\ker\eta$, or upper or lower triangular there, or block shaped; since $f$ is two-signed, it is irreducible or block shaped.
Suppose it is block shaped.
The middle entry $m$ of its matrix parts on $\ker\eta$ is a character with values in $1+\mathfrak m$.
Let $s\in\ker\eta$ lie in an inertia group.
If $\bar\Psi(s)\ne1$, then $m(s)$, a root of the characteristic polynomial, which is $(X-\Psi(s))(X-1)(X-\Psi(s)^{-1})$ by the first and second facts, is congruent to $1$ and not to $\Psi(s)^{\pm1}$, so $m(s)=1$.
If $\bar\Psi(s)=1$ and $s$ lies above $\ell$, then $m(s)=m(sx)m(x)^{-1}=1$, replacing $x$ by a conjugate; and if $s$ lies above $p\ne\ell$, then $s\in\mathcal N$ and $r_f(s)=1$.
By the third fact $m=1$.
The determinant of the matrix parts on $\ker\eta$ is then the determinant of the corner block; it is a character with values in $1+\mathfrak m$, equal to $\Psi(s)\Psi(s)^{-1}=1$ on the inertia groups, so it is $1$ by the third fact.
For $g\in\operatorname{Gal}(\bar{\mathbf Q}/\mathbf Q)$ let $C_g$ be the corner block of the matrix part of $r_f(g)$, and put $W(g)=C_g$ if $\eta(g)=1$ and $W(g)=C_g\operatorname{diag}(-1,1)$ if $\eta(g)=-1$.
On a block shaped matrix of corner block $C$ with $\det C=1$, the automorphism $M\mapsto J\,{}^tM^{-1}J$ of $GL_3$ acts on the corner block by conjugation by $\operatorname{diag}(-1,1)$; since the matrix part of $r_f(c)$ is $1$, the matrix part of $r_f(g)$ for $\eta(g)=-1$ is that of $r_f(gc)$, which is block shaped, and this makes $W$ a homomorphism.
It is continuous, agrees with the corner block on $\ker\eta$, and $W(c)=\operatorname{diag}(-1,1)$ has trace $0$; $\det W=\eta$.

(d) This is the argument of the proof of Lemma~\ref{sha-lifts-points}(g), which uses only the normalization of $r_f$ and $r_{f'}$, the nonvanishing of $\zeta(r_f)$ and $\zeta(r_{f'})$, and $\zeta(r_{f'})\ne-\zeta(r_f)$.
\end{proof}
\leanref{Development/AutomorphicArithmetic/AutomorphicVarieties/CurveDVRSpecialization/GfResZpow.lean}{UpperTheory.V.Pr.finA_of_end}
\leanref{Development/AutomorphicArithmetic/AutomorphicVarieties/CurveDVRSpecialization/FinCrit.lean}{UpperTheory.V.Pr.finCrit}
\leanref{Development/AutomorphicArithmetic/AutomorphicVarieties/CurveDVRSpecialization/ConstOfChar.lean}{UpperTheory.V.Pr.constOfChar}
\leanref{Development/AutomorphicArithmetic/AutomorphicVarieties/UnitaryThreeGroups/RootOfCharpoly.lean}{UpperTheory.V.U3.pt_fixed}
\leanref{Development/AutomorphicArithmetic/AutomorphicVarieties/UnitaryThreeGroups/PtShape.lean}{UpperTheory.V.U3.twoPt_shape}
\leanref{Development/AutomorphicArithmetic/AutomorphicVarieties/UnitaryThreeGroups/CoboundaryLemmaNoColumn.lean}{UpperTheory.V.U3.typeII_facts}

\begin{lemma}\label{sha-model-tangent}
In Lemma~\ref{sha-model-points}, let $f$ be a point of $R$, and let $\tilde f:R\to\bar{\mathbf Q}_\ell[\epsilon]/(\epsilon^2)$ be a homomorphism of $\mathcal O$-algebras with $\tilde f\equiv f\pmod\epsilon$, so that $\tilde f=f+\epsilon D_{\tilde f}$ for a derivation $D_{\tilde f}$ at $f$ (Definition~\ref{sha-algebra-graded}).
Write the matrix part of $\tilde f\circ r^{\mathrm u}(g)$ as $M_g+\epsilon M'_g$, and let the \emph{tangent cocycle} of $\tilde f$ be $y_{\tilde f}(g)=M'_gM_g^{-1}$.
The action $\operatorname{ad}r_f$ of $\operatorname{Gal}(\bar{\mathbf Q}/\mathbf Q)$ on $M_3(\bar{\mathbf Q}_\ell)$, cocycles and coboundaries are as in Definition~\ref{sha-lifts-graded}(e).
\begin{enumerate}
\item $y_{\tilde f}$ is a continuous cocycle for $r_f$. It vanishes at $c$ and on $\mathcal N$, the entries $(1,2)$ and $(2,1)$ of $y_{\tilde f}(g_1)M_{g_1}-y_{\tilde f}(cg_1^{-1}c)M_{cg_1^{-1}c}$ vanish, and its class lies in the Bloch-Kato Selmer group $H^1_f(\mathbf Q,\operatorname{ad}r_f)$.
\item Let $\tilde f'$ be a second homomorphism with $\tilde f'\equiv f\pmod\epsilon$, and let $a,a'\in\bar{\mathbf Q}_\ell$ with $ay_{\tilde f}+a'y_{\tilde f'}=0$. Then $aD_{\tilde f}+a'D_{\tilde f'}=0$.
\item Let $X\in M_3(\bar{\mathbf Q}_\ell)$, and suppose that the coboundary of $X$ for $r_f$ vanishes at $c$ and has the vanishing property at $g_1$ of (a). Then $X$ is a multiple of $\operatorname{diag}(1,0,-1)$. If $r_f=r_0$, the coboundary of every multiple of $\operatorname{diag}(1,0,-1)$ is $0$.
\item Suppose that $r_f=r_0$ and that $y_{\tilde f}$ takes values in the matrices supported on the entries of weight $j$, for some $j\le0$. Then $y_{\tilde f}$ is a coboundary.
\item Assume the Galois representations as in Theorem~\ref{sha-residually-reducible-automorphy}, let $\iota:\bar{\mathbf Q}_\ell\cong\mathbf C$, and suppose that the conclusions (ii) and (iii) of that theorem hold for the ordinary points relative to $g_1$ and $B$. If $f$ is two-signed, then $H^1_f(\mathbf Q,\operatorname{ad}r_f)=0$, and $y_{\tilde f}$ is a coboundary.
\end{enumerate}
\end{lemma}
\begin{proof}
(a) Since $\tilde f\circ r^{\mathrm u}$ is a homomorphism, $y_{\tilde f}$ is a cocycle for $r_f$, and it is continuous because $r^{\mathrm u}$ is.
It vanishes at $c$ and on $\mathcal N$, where the matrix part of $r^{\mathrm u}$ is constant, and the condition at $g_1$ is the coefficient of $\epsilon$ in the normalization of $\tilde f\circ r^{\mathrm u}$.
For the Bloch-Kato conditions: at a prime $p\ne\ell$, the restriction of $y_{\tilde f}$ to the inertia group vanishes on the subgroup of finite index of its elements in $\mathcal N$, so it is inflated from a finite quotient, and averaging over that quotient, which is possible in characteristic $0$, makes it a coboundary there.
At the primes above $\ell$, the restriction of $\tilde f\circ r^{\mathrm u}$ to $D_B$ is $(\tilde f\circ h^{\mathrm u})\circ u_B$ for the homomorphism $h^{\mathrm u}:R_B\to R$ that kills $\mathfrak a_B$, and Lemma~\ref{sha-model-local}(d) puts the restriction of $y_{\tilde f}$ to $D_B$ in the Bloch-Kato subspace.

(b) Let $D=aD_{\tilde f}+a'D_{\tilde f'}$, a derivation at $f$.
For $g\in\operatorname{Gal}(\bar{\mathbf Q}/\mathbf Q)$ the matrix of the values of $D$ on the entries of the matrix part of $r^{\mathrm u}(g)$ is $(ay_{\tilde f}(g)+a'y_{\tilde f'}(g))M_g=0$.
So $D$ vanishes on these entries, which topologically generate $R$ over $\mathcal O$.
The image of $\tilde f$ is a quotient of $R$ whose reduction modulo $\epsilon$ is the image of $f$, which is finite over $\mathcal O$ as in the proof of (e), and whose intersection with $\epsilon\bar{\mathbf Q}_\ell$ is a finitely generated ideal on which it acts through that reduction; so it is finite over $\mathcal O$, and $D_{\tilde f}$ is continuous.
The same holds for $\tilde f'$, so $D$, which vanishes on the $\mathcal O$-algebra generated by the entries, vanishes on $R$.

(c) The coboundary of $X$ vanishes at $c$, where the matrix part of $r_f$ is $1$ and the multiplier is $-1$, so $X=-J\,{}^tXJ$.
At $g_1$, by the normalization of $r_f$, the matrix part of $r_f(g_1)-r_f(cg_1^{-1}c)$ is $\operatorname{diag}(\zeta,0,-\zeta)$ with $\zeta\ne0$ (proof of Lemma~\ref{sha-model-points}), and the vanishing of the entries $(1,2)$ and $(2,1)$ gives $\zeta X_{12}=\zeta X_{21}=0$; with $X=-J\,{}^tXJ$ and $2$ invertible, this gives $X=a\operatorname{diag}(1,0,-1)$ for some $a$, as in the proof of Lemma~\ref{sha-lifts-points}(f).
The matrix $\operatorname{diag}(1,0,-1)$ commutes with the matrix parts of $r_0$, and $-J\,{}^t\operatorname{diag}(1,0,-1)J=\operatorname{diag}(1,0,-1)$, so its coboundary for $r_0$ is $0$.

(d) Let $\theta\in\{1,\eta\}$ and $\epsilon_\theta=-1$ if $\theta=1$, $\epsilon_\theta=1$ if $\theta=\eta$.
For an entry $(a,b)$ of weight $j$, let $(a',b')=(4-b,4-a)$, also of weight $j$.
The map $Y\mapsto Y_{ab}+\epsilon_\theta Y_{a'b'}$ carries $\operatorname{ad}r_0$ on the matrices of weight $j$ to the character $\Psi^j\theta$, because $\operatorname{ad}r_0(g)$ multiplies the entry $(a,b)$ by $\Psi(g)^j$ when $\eta(g)=1$, and $\operatorname{ad}r_0(c)Y=-J\,{}^tYJ$.
So each $g\mapsto y_{\tilde f}(g)_{ab}+\epsilon_\theta y_{\tilde f}(g)_{a'b'}$ is a cocycle for $\Psi^j\theta$ that satisfies the Bloch-Kato conditions, by (a).
It suffices to show that each of these is a coboundary: then the entries of $y_{\tilde f}$, which are half the sums and differences of these combinations, are coboundaries, and $y_{\tilde f}$ is the coboundary of the matrix assembled from them.
Let $z$ be one of these cocycles.
For $j=0$, $z$ is a homomorphism on $\operatorname{Gal}(\bar{\mathbf Q}/K)=\ker\eta$, on which $\theta$ is trivial.
It vanishes on the inertia groups at the primes not above $\ell$, by the Bloch-Kato condition there, and on those above $\ell$, since $\theta$ is unramified at $\ell$ and a Bloch-Kato class of an unramified character of finite order is unramified \cite{BK90}.
So it factors through the class group of $K$, which is finite, and it vanishes, since $\bar{\mathbf Q}_\ell$ is torsion free.
For $g$ with $\eta(g)=-1$, $gc\in\ker\eta$, so $z(g)=z(gc\cdot c)=z(gc)+\theta(gc)z(c)=z(c)$.
If $\theta=1$, then $2z(c)=z(c^2)=0$ and $z=0$; if $\theta=\eta$, then $z$ is the coboundary $g\mapsto(\theta(g)-1)m$ of $m=-z(c)/2$.
For $j<0$, write $j=-n$; then $\Psi^j\theta=\chi_{\mathrm{cyc}}^{-nk}\psi^{-n}\theta$ with $nk\ge1$.
A Bloch-Kato class for $\chi_{\mathrm{cyc}}^{-nk}$ times a character of finite order is a coboundary on the decomposition groups above $\ell$ \cite{BK90}, and over a totally complex field $L$ containing $\mu_\ell$ on which the finite part is trivial it lies in a Selmer group that is controlled, through Poitou-Tate duality \cite[Chapter VIII, \S6]{NSW08} and local Tate duality \cite[Chapter VII, \S2]{NSW08}, by $H^2(G_{L,S_\ell},\mathbf Z_\ell(nk+1))$, which is finite by Theorem~\ref{sha-soule-bound}; since $\bar{\mathbf Q}_\ell$ is a vector space, the class vanishes, as in the proof of Lemma~\ref{sha-lifts-residual}(d).

(e) The image of $f$ is a complete local noetherian domain contained in $\bar{\mathbf Q}_\ell$, whose fraction field is algebraic over that of $\mathcal O$; so it is finite over $\mathcal O$, and its integral closure $A$ is a discrete valuation ring finite over $\mathcal O$; by Lemma~\ref{sha-model-universal}(b), $r_f$ is an ordinary point with values in $A$.
By Lemma~\ref{sha-model-points}(c), $r_f$ is irreducible on $\ker\eta$, or block shaped with the properties stated there.
In the first case, conclusion (ii) of Theorem~\ref{sha-residually-reducible-automorphy} shows that the restriction $\rho_f$ of $r_f$ to $\operatorname{Gal}(\bar{\mathbf Q}/K)$, twisted by $\chi_{\mathrm{cyc}}^{-1}$, is attached through $\iota$ to a regular algebraic cuspidal automorphic representation $\pi$ with $\pi^\vee\cong\pi^c$, and is windowed; Lemma~\ref{sha-rigidity-points}(1) gives $H^1_f(\mathbf Q,\operatorname{ad}r_f)=0$.
In the second case, let $W$ be as in Lemma~\ref{sha-model-points}(c).
Conclusion (iii) of Theorem~\ref{sha-residually-reducible-automorphy} makes $W\otimes\chi_{\mathrm{cyc}}^{-k}$ attached through $\iota$ to a newform of weight $2k+1$, with the windowed condition.
The block shape gives a decomposition of $\operatorname{ad}r_f$, stable under $\operatorname{Gal}(\bar{\mathbf Q}/\mathbf Q)$, into four pieces: the matrices supported on the corner, which form a representation $R$ as in Lemma~\ref{sha-rigidity-points}(2), with the pairing $P$ coming from the action of $c$; the entries $(1,2),(3,2)$ and the entries $(2,1),(2,3)$, combined with the signs $+$ and $-$, which form $W$ and $W\otimes\eta$; and the line of $\operatorname{diag}(0,1,0)$, on which $\operatorname{ad}r_f$ acts as $\operatorname{ad}r_0$, through $\eta$, and which has weight $0$.
The Bloch-Kato Selmer group of $R$ vanishes by Lemma~\ref{sha-rigidity-points}(2); those of $W$ and $W\otimes\eta$ by Theorem~\ref{sha-kato-selmer}, whose hypotheses on $W$ come from the conditions on $r_f$ away from $\ell$ and at $B$, its congruence to $r_0$, and conclusion (iii); and that of the last line by (d).
So $H^1_f(\mathbf Q,\operatorname{ad}r_f)=0$ in both cases, and $y_{\tilde f}$ is a coboundary by (a).
\end{proof}
\leanref{Development/AutomorphicArithmetic/AutomorphicVarieties/UnitaryThreeGroups/TanCocIsBK.lean}{UpperTheory.V.U3.tanCoc_isBK}
\leanref{Development/AutomorphicArithmetic/AutomorphicVarieties/UnitaryThreeGroups/TanDet.lean}{UpperTheory.V.U3.tanCoc_determines}
\leanref{Development/AutomorphicArithmetic/AutomorphicVarieties/UnitaryThreeGroups/BlochKatoNonpositiveWeight.lean}{UpperTheory.V.U3.bloch_kato_nonpositive_weight}
\leanref{Development/AutomorphicArithmetic/AutomorphicVarieties/UnitaryThreeGroups/Rigid0.lean}{UpperTheory.V.U3.PtHyp}
\leanref{Development/AutomorphicArithmetic/AutomorphicVarieties/ModularityArgument/ThresholdArguments/DenseThresholdCases/FPtHypU.lean}{UpperTheory.V.Z.J.ptHyp_of_autB_of_lemW_of_isBase}
\leanref{Development/AutomorphicArithmetic/AutomorphicVarieties/ModularityArgument/ThresholdArguments/DenseThresholdCases/TypeIIOfBlochKatoLocalPoint.lean}{UpperTheory.V.Z.Kato.WeilSphericalRecords.bloch_kato_selmer_vanishes_of_type_ii_modular_ext}

\begin{proposition}\label{sha-model-exists}
In Lemma~\ref{sha-model-points}, let $x_0:R\to\mathcal O$ be the homomorphism with $x_0\circ r^{\mathrm u}=r_0$.
There are a $\mathbf Z$-graded commutative $\mathcal O$-algebra $A=\bigoplus_jA_j$, a homomorphism $x:A\to\mathcal O$, integers $r\ge1$ and $s,r',s',m\ge0$, families $F$, $G$, $E$, $E'$ and $Z$, and a homomorphism of $\mathcal O$-algebras $q_A:A\to R$, with the following properties.
\begin{enumerate}
\item $A$, $x$ and the families satisfy the hypotheses of Lemma~\ref{sha-algebra-cotangent} that precede its conditions, and the first sentence of its condition (1); the kernel of $q_A$ is the ideal $I$ generated by the families; and $x=x_0\circ q_A$.
\item For every $g$ and every entry $(a,b)$, the entry $(a,b)$ of the matrix part of $r^{\mathrm u}(g)$ is $q_A(f)$ for some $f\in A_{b-a}$.
\item For every $j$, every $f\in A_j$ and every $n\ge1$, $q_A(f)$ is congruent modulo $\mathfrak m_R^n$ to an $\mathcal O$-linear combination of products of entries of the matrix parts of the $r^{\mathrm u}(g)$ whose weights add up to $j$.
\item Suppose that the image of $A_0$ in $A/I$ is a finitely generated $\mathcal O$-module. Then every homomorphism of $\mathcal O$-algebras $\xi:A\to\bar{\mathbf Q}_\ell$ that vanishes on $I$ is in the orbit (Definition~\ref{sha-algebra-graded}) of $f\circ q_A$ for some point $f$ of $R$.
\item Under the hypothesis of (d), let $f$ be a point of $R$ and $D$ a derivation at $f\circ q_A$ with values in $\bar{\mathbf Q}_\ell$ that vanishes on $I$. Then there is a homomorphism $\tilde f:R\to\bar{\mathbf Q}_\ell[\epsilon]/(\epsilon^2)$ with $\tilde f\equiv f\pmod\epsilon$ and $D=D_{\tilde f}\circ q_A$, in the notation of Lemma~\ref{sha-model-tangent}.
\item Let $a\in\mathcal O$ be nonzero, and let $D$ be a derivation at $x_0$ with values in $\mathcal O/a\mathcal O$ that vanishes on the entries of weight other than $-1$ of the matrix parts of the $r^{\mathrm u}(g)$. Suppose that there is a derivation $\partial$ at $x_B$ with values in $\mathcal O$ that vanishes on $\mathfrak a_B$ and on the entries of weight other than $-1$ of the matrix parts of the $u_B(h)$, and whose value on each entry of the matrix part of $u_B(h)$ reduces modulo $a$ to the value of $D$ on the same entry of $r^{\mathrm u}(h)$, for every $h\in D_B$. Then there is a derivation $D'$ at $x$ with values in $\mathcal O$, of weight $-1$, with $D\circ q_A\equiv D'\pmod a$.
\end{enumerate}
\end{proposition}
\begin{proof}
Let $\Gamma$ be the quotient of $\operatorname{Gal}(\bar{\mathbf Q}/\mathbf Q)$ by the closed normal subgroup generated by $\mathcal N$, and for $j\in\mathbf Z$ let $\mathfrak g_j\subset M_3(\mathbf F)$ be the matrices supported on the entries of weight $j$, with the action of $\Gamma$ through the reduction of $r_0$, as in Definition~\ref{sha-lifts-selmer}(a).

\emph{Step 1: the local presentation.}
By Lemma~\ref{sha-model-local}(b), for every $A'$ as there and every $t\in1+\mathfrak m_{A'}$ the conjugation by $\tau(t)$, which is congruent to $1$, preserves the homomorphisms $R_B\to A'$ that kill $\mathfrak a_B$.
As in the proof of Lemma~\ref{sha-lifts-presentation}(a), this gives a surjection $\pi_B:\mathcal O[[X_1,\dots,X_n]]\to R_B/\mathfrak a_B$ with homogeneous kernel, where $X_i$ has a weight in $\{-2,\dots,2\}$ and maps to the image of an entry of that weight of the image of some $u_B(h)$ under the embedding of $\mathcal G_3$ in $GL_6$ of the proof of Lemma~\ref{sha-lifts-universal}, minus its value at $r_0$; there the entry $(a,b)$ of the second diagonal block, which is $J\,{}^tM^{-1}J$ for the first block $M$, has the weight $b-a$, and every entry of the image of every $u_B(h)$ is the image of a power series of its weight.

\emph{Step 2: the relative presentation.}
Let $R^\flat$ be the quotient of the ring $R^\square$ of the proof of Lemma~\ref{sha-model-universal}(a) by the ideal generated by $h^\square(\mathfrak a_B)$; it classifies the lifts trivial on $\mathcal N$ whose restriction to $D_B$ factors through $R_B/\mathfrak a_B$.
Let $H^1_B(\mathfrak g_j)$ be the space of pairs $(z,y)$ of a continuous cocycle $z:\Gamma\to\mathfrak g_j$ and $y\in\mathfrak g_j$ with $z=dy$ on $D_B$, modulo the pairs $(dm,m)$, and let $\mathrm{Ob}^B_j$ be the analogous quotient of pairs of a continuous $2$-cocycle of $\Gamma$ and a $1$-cochain of $D_B$ whose coboundary is its restriction.
Then $R^\flat$ is the quotient of $(R_B/\mathfrak a_B)[[T_1,\dots,T_\tau]]$ by an ideal generated by $\kappa$ elements, each homogeneous of some weight, where $T_t$ has a weight and maps to an entry of that weight of the image of some $r^\flat(g)$ in $GL_6$ minus its value at $r_0$, and the numbers of variables $T_t$ and of generators of weight $j$ are $\dim H^1_B(\mathfrak g_j)$ and $\dim\mathrm{Ob}^B_j$.
This is the presentation of a global deformation ring over its local deformation ring \cite[\S2.2]{CHT08}, with weights obtained as in the proof of Lemma~\ref{sha-lifts-presentation}(a) and (c).
Moreover
\[
 \dim\mathrm{Ob}^B_j=\dim H^1_B(\mathfrak g_j)+\delta_B(j),\qquad\delta_B(j)=\dim\mathfrak g_j^{c=1},
\]
which is $1$ if $|j|\le1$ and $0$ otherwise, as in the proof of Lemma~\ref{sha-lifts-residual}(b).
Indeed, the map from $H^2(\Gamma,\mathfrak g_j)$ to the sum of the local $H^2$ is surjective, by Poitou-Tate duality \cite[Chapter VIII, \S6]{NSW08} and an element of $\Gamma$ on which the reduction of $\chi_{\mathrm{cyc}}$ is $-1$ and the action on $\mathfrak g_j$ is trivial; the local term at $B$ cancels because the cohomology at $B$ is taken relative to $D_B$; the local terms at the other primes $p\in S$ vanish because the unramified classes there have the dimension of $H^0$, the subgroup of the inertia group in $\mathcal N$ having index prime to $\ell$; and the global Euler characteristic formula \cite[Chapter VIII, \S7]{NSW08} leaves $\dim\mathfrak g_j^{c=1}$.

\emph{Step 3: the normalization.}
Since $c^2=1$ and the multiplier at $c$ is $-1$, the matrix part $M$ of $r^\flat(c)$ satisfies $MJ={}^t(MJ)$, so $M=1$ if and only if its entries $(1,1)$ and $(2,2)$ equal $1$ and its entries $(1,2)$, $(2,1)$, $(1,3)$ and $(3,1)$ vanish.
So $R$ is the quotient of $R^\flat$ by these six elements and the entries $(1,2)$ and $(2,1)$ of the matrix part of $r^\flat(g_1)-r^\flat(cg_1^{-1}c)$; the eight elements have weights $0,0,1,-1,2,-2,1,-1$.

\emph{Step 4: the model.}
Let $\mathcal P=(R_B/\mathfrak a_B)[[T_1,\dots,T_\tau]]$, the quotient of $\tilde{\mathcal P}=\mathcal O[[X_1,\dots,X_n,T_1,\dots,T_\tau]]$ by the homogeneous kernel of $\pi_B$, graded by Steps 1 and 2, and let $A$ be the $\mathcal O$-algebra of the elements of $\mathcal P$ that are finite sums of homogeneous ones, with $A_j$ its homogeneous elements of weight $j$.
Let $x:A\to\mathcal O$ be induced by $x_B$ and $T_t\mapsto0$, and $q_A$ the restriction to $A$ of the map $\mathcal P\to R^\flat\to R$.
The kernel of $\mathcal P\to R$ is generated by the homogeneous elements of Steps 2 and 3; since they lie in $A$, the kernel of $q_A$ is the ideal $I$ of $A$ that they generate, and the kernel of $\tilde{\mathcal P}\to R$ is generated by homogeneous elements.
Let $a(j)$ be the number of variables $T_t$ of weight $j$; then $\mathcal P$ has $\#\{i:w(i)=j\}+a(j)$ chart variables of weight $j$, through the chart $\varphi_B$ of Lemma~\ref{sha-model-local}(a) extended by $T_t\mapsto T_t$, and $I$ has, by Steps 2 and 3, $a(j)+\delta_B(j)$ plus $2,1,2,2,1$ generators of weights $j=-1,-2,0,1,2$.
With $r=a(-1)+2$, $s=a(-2)+1$, $r'=a(1)+3$, $s'=a(2)+1$ and $m=a(0)+3$, the numbers of chart variables of weights $-1,-2,1,2,0$ are $r$, $s$, $r'+1$, $s'+1$ and $m$, and the numbers of generators are $r+1$, $s$, $r'$, $s'$ and $m$; the generators form the families $F$, $G$, $E$, $E'$ and $Z$.
The ring $A_0$ is complete local noetherian with residue field $\mathbf F$, $A$ is generated over $A_0$ by the variables of nonzero weight, and $x$ vanishes on $A_j$ for $j\ne0$.
The chart has the properties of condition (1) of Lemma~\ref{sha-algebra-cotangent}, by those of $\varphi_B$, since $\mathfrak a_B$ is its kernel.
This proves (a).

(b) The entries of the matrix parts of the $r^\flat(g)$ are images of homogeneous elements of $\mathcal P$ of their weights, by the choice of the variables in Steps 1 and 2.

(c) Every variable maps to an entry of the matrix part of some $r^{\mathrm u}(g)$ minus a constant, with the same weight: for an entry of the second block of the embedding in $GL_6$ this holds because $r^{\mathrm u}(c)=(1,-1)$, so that the matrix part of $r^{\mathrm u}(cgc)$ is $J\,{}^tM_g^{-1}J$, where $M_g$ is that of $r^{\mathrm u}(g)$.
A power series of weight $j$ is congruent modulo $\mathfrak m_{\mathcal P}^n$ to a polynomial of weight $j$ in the variables, whose image is a combination of products of entries of total weight $j$.

(d) Let $\xi$ be as there. As in the proof of Lemma~\ref{sha-lifts-points}(b), the finiteness of the image of $A_0$ gives $|\xi(y)|\le1$ for $y\in A_0$ and $|\xi(y)|<1$ for $y\in A_0\cap\mathfrak m_{\mathcal P}$; these inequalities, applied to the products of powers of two variables of opposite weights, give $t\in\bar{\mathbf Q}_\ell^\times$ with $|t^{-w}\xi(X)|<1$ for every variable $X$ of weight $w$.
The homomorphism $\tilde{\mathcal P}\to\bar{\mathbf Q}_\ell$ that sends each variable $X$ of weight $w$ to $t^{-w}\xi(X)$ is continuous, and on the homogeneous elements of weight $j$ it is $t^{-j}$ times $\xi$ composed with $\tilde{\mathcal P}\to\mathcal P$, by the argument of that proof; so it kills the kernel of $\tilde{\mathcal P}\to R$, it factors through a point $f$ of $R$, and $\xi=t\cdot(f\circ q_A)$.

(e) The homomorphism $\tilde{\mathcal P}\to\bar{\mathbf Q}_\ell[\epsilon]/(\epsilon^2)$ that sends each variable $X$ to $f(q_A(X))+\epsilon D(X)$ is continuous, since $|f(q_A(X))|<1$.
On the homogeneous elements it agrees with $f\circ q_A+\epsilon D$ composed with $\tilde{\mathcal P}\to\mathcal P$: both are homomorphisms that agree on the variables, and on the elements of weight $j$ the second is continuous because $D$ vanishes on $I$ and is $\mathcal O$-linear on the image of $A_j$ in $A/I$, which is a finite $\mathcal O$-module, being finite over the image of $A_0$.
So it kills the kernel of $\tilde{\mathcal P}\to R$, and it factors through a homomorphism $\tilde f$ of $R$ with the required property.

(f) Let $D'$ be the derivation of $\mathcal P$ at the homomorphism induced by $x$ that is $\partial$ on $R_B/\mathfrak a_B$, and sends $T_t$ to a lift to $\mathcal O$ of $D(q_A(T_t))$ if $T_t$ has weight $-1$, and to $0$ otherwise.
It vanishes on the variables of weight other than $-1$, since $\partial$ does: on an entry of the second block of the embedding in $GL_6$, a derivation at $x_B$ is a multiple of its value on an entry of the matrix part of the same weight.
So its restriction to $A$ has weight $-1$.
The derivations $D\circ q_A$ and $D'\bmod a$ at $x$ agree on the variables, by the hypothesis on $\partial$ and the choice of the $D'(T_t)$. Both vanish on the square of the kernel of $x$ on $A$, and every homogeneous element of $A$ is congruent modulo that square to a constant plus an $\mathcal O$-linear combination of the variables, as one sees by taking homogeneous components; so they agree on $A$.
\end{proof}
\leanref{Development/AutomorphicArithmetic/AutomorphicVarieties/UnitaryThreeGroups/ModelEx.lean}{UpperTheory.V.U3.exists_model}
\leanref{Development/AutomorphicArithmetic/AutomorphicVarieties/UnitaryThreeGroups/ConjStable.lean}{UpperTheory.V.U3.Model}
\leanref{Development/AutomorphicArithmetic/AutomorphicVarieties/UnitaryThreeGroups/LocPres.lean}{UpperTheory.V.U3.exists_locPres}
\leanref{Development/AutomorphicArithmetic/AutomorphicVarieties/UnitaryThreeGroups/BoxPresOf.lean}{UpperTheory.V.U3.boxPres}
\leanref{Development/AutomorphicArithmetic/AutomorphicVarieties/CommutativeAlgebra/LemmaP.lean}{UpperTheory.V.Alg.lemmaP}
\leanref{Development/AutomorphicArithmetic/AutomorphicVarieties/UnitaryThreeGroups/BoxCount.lean}{UpperTheory.V.U3.box_count}
\leanref{Development/AutomorphicArithmetic/AutomorphicVarieties/UnitaryThreeGroups/ModelOfPres.lean}{UpperTheory.V.U3.model_of_pres}
\leanref{Development/AutomorphicArithmetic/AutomorphicVarieties/UnitaryThreeGroups/ExistsTorusTwistEqLocalPointOfWeightedPres.lean}{UpperTheory.exists_torus_twist_eq_localPoint_of_weightedPres}
\leanref{Development/AutomorphicArithmetic/AutomorphicVarieties/UnitaryThreeGroups/ExistsDualNumberPointOfPointDerivationVanishing.lean}{UpperTheory.exists_dualNumber_point_of_pointDerivation_vanishing}
\leanref{Development/AutomorphicArithmetic/AutomorphicVarieties/UnitaryThreeGroups/ExistsBasePointDerivationOfWeightHomogeneous.lean}{UpperTheory.exists_base_pointDerivation_of_weight_homogeneous}
\leanref{Development/AutomorphicArithmetic/AutomorphicVarieties/UnitaryThreeGroups/ExistsSpanEntMonApproxOfMemWtPart.lean}{UpperTheory.exists_span_entMon_approx_of_mem_wtPart}

\subsection{Existence of Siegel data}\label{subsection-sha-siegel-data-existence}

\begin{lemma}\label{sha-counts-siegel-dyadic}
Let $n\ge0$, let $\xi$ be a primitive Dirichlet character of conductor $2^n$ with values in a domain, and let $v>1$ be an odd integer.
If $a\ge0$ and $u$ is odd with $\xi(v)^{2^au}=1$, then $2^n<v^22^a$.
\end{lemma}
\begin{proof}
If $n\le1$, then $2^n\le2<9\le v^2$. Let $n\ge2$.
Write $v^2-1=2^sw$ with $w$ odd; since $v$ is odd, $s\ge3$, and $2^s<v^2$.
For odd $u$, $(v^2)^u=1+2^sw_1$ with $w_1$ odd, and since $(1+2^tx)^2=1+2^{t+1}(x+2^{t-1}x^2)$, squaring $a$ times gives $v^{2^{a+1}u}=1+2^{s+a}w_2$ with $w_2$ odd.
Suppose that $s+a<n$. Squaring $n-1-s-a$ more times gives $v^{2^{n-s}u}\equiv1+2^{n-1}\pmod{2^n}$, so $\xi(1+2^{n-1})=\big(\xi(v)^{2^au}\big)^{2^{n-s-a}}=1$.
The units congruent to $1$ modulo $2^{n-1}$ are $1$ and $1+2^{n-1}$, so $\xi$ would then factor through $(\mathbf Z/2^{n-1})^\times$, contradicting its primitivity.
Hence $n\le s+a$, and $2^n\le2^s2^a<v^22^a$.
\end{proof}
\leanref{Development/AutomorphicArithmetic/AutomorphicVarieties/SymplecticGSp4/TwoPowLtSqMulOfPrimitiveCharPowEqOne.lean}{UpperTheory.two_pow_lt_sq_mul_of_primitive_char_pow_eq_one}

\begin{lemma}\label{sha-counts-siegel-auxiliary}
Let $\ell$ be an odd prime, let $\mathcal O$ be as in Definition~\ref{sha-character-parts}, write $G=\operatorname{Gal}(\bar{\mathbf Q}/\mathbf Q)$ and $\bar x$ for the residue of $x\in\mathcal O$, let $k\ge0$, and let $\psi,\xi:G\to\mathcal O^\times$ be characters, the kernel of $\psi$ being open.
Put $\Psi=\chi_{\mathrm{cyc}}^k\psi$.
Then there are a character $\omega:G\to\mathcal O^\times$, a complex conjugation $c$, for each prime $p$ a place $v_p$ of $\bar{\mathbf Q}$ above $p$ and a Frobenius element $\sigma_p$ at $v_p$, and elements $\phi,g_0,s_0\in G$, with the following properties, where $\alpha=\chi_{\mathrm{cyc}}\omega^{-1}\xi$ and $I_\ell$ is the inertia group at $v_\ell$.
\begin{enumerate}
\item $\omega^{\ell-1}=1$ and $\bar\omega=\bar\chi_{\mathrm{cyc}}$.
\item $\chi_{\mathrm{cyc}}(\sigma_\ell)=1$, and $\phi$ is the limit of $\sigma_\ell^{\ell^{n!}}$ as $n\to\infty$.
\item If $\bar\alpha$ and $\bar\Psi$ are both nontrivial, then $\bar\alpha(g_0)\ne1$ and $\bar\Psi(g_0)\ne1$.
\item $s_0\in I_\ell$, $\psi(s_0)=\omega(s_0)=1$, and $\chi_{\mathrm{cyc}}(s_0)$ has infinite order.
\end{enumerate}
\end{lemma}
\begin{proof}
(1) Let $\omega$ be the Teichm\"uller lift of $\bar\chi_{\mathrm{cyc}}$: the $(\ell-1)$-st roots of unity in $\mathcal O$ map bijectively to $\mathbf F_\ell^\times$.
Complex conjugations exist, every prime has a place of $\bar{\mathbf Q}$ above it, and each decomposition group contains Frobenius elements.

(2) Let $g$ be a Frobenius element at $v_\ell$. Since $\chi_{\mathrm{cyc}}$ maps $I_\ell$ onto $\mathbf Z_\ell^\times$, there is $t\in I_\ell$ with $\chi_{\mathrm{cyc}}(t)=\chi_{\mathrm{cyc}}(g)^{-1}$, and $\sigma_\ell=tg$ is again a Frobenius element at $v_\ell$, with $\chi_{\mathrm{cyc}}(\sigma_\ell)=1$.
In a finite quotient of $G$ the image $x$ of $\sigma_\ell$ has order $\ell^am$ with $\ell\nmid m$, and $x^{\ell^{n!}}$ does not depend on $n$ once $n\ge a$ and $\varphi(m)\mid n!$.
By compactness the sequence $\sigma_\ell^{\ell^{n!}}$ has a cluster point $\phi$, and since its image in each finite quotient is eventually constant, the sequence converges to $\phi$.

(3) Let $\bar\alpha(x)\ne1$ and $\bar\Psi(y)\ne1$. If $\bar\Psi(x)\ne1$, take $g_0=x$; if $\bar\alpha(y)\ne1$, take $g_0=y$; otherwise $g_0=xy$ has $\bar\alpha(xy)=\bar\alpha(x)\ne1$ and $\bar\Psi(xy)=\bar\Psi(y)\ne1$.
If $\bar\alpha$ or $\bar\Psi$ is trivial, take $g_0=1$.

(4) Let $s\in I_\ell$ with $\chi_{\mathrm{cyc}}(s)=1+\ell$, which has infinite order in $\mathbf Z_\ell^\times$. Let $m$ be the index of $\ker\psi$, which is finite since $\ker\psi$ is open, and put $s_0=s^{m(\ell-1)}$.
Then $s_0\in I_\ell$, $\psi(s_0)=1$ because $g^m\in\ker\psi$ for every $g$, $\omega(s_0)=1$ because $\omega^{\ell-1}=1$, and $\chi_{\mathrm{cyc}}(s_0)=(1+\ell)^{m(\ell-1)}$ has infinite order.
\end{proof}
\leanref{Development/AutomorphicArithmetic/AutomorphicVarieties/SymplecticGSp4/ReFrob.lean}{UpperTheory.V.GradeZeroInterpolatingElement.rawEx}
\leanref{Development/AutomorphicArithmetic/AutomorphicVarieties/SymplecticGSp4/PrimeTo.lean}{UpperTheory.V.GradeZeroInterpolatingElement.primeTo}
\leanref{Development/AutomorphicArithmetic/AutomorphicVarieties/HeightAndVanishing/ExistsStrictHenselianAbsolutelyClosedLocalRing.lean}{UpperTheory.V.HV.omT}

\begin{lemma}\label{sha-counts-siegel-host}
Let $h$ be a host datum for a primitive Dirichlet character $\psi_0$, let $\ell$ be an odd prime, let $\mathcal O$ be as in Definition~\ref{sha-character-parts}, and let $\psi$ be a character of $(\mathbf Z/N_h)^\times$ with values in $\mathcal O$ that is trivial on $H_h$ (Definition~\ref{sha-host-datum}).
Let $\kappa$, $\theta_0$, $\psi_2$, $\theta_1$, $c_2$, $f_0$, $n_0$ and $T_b$ be attached to $\psi$ as in Definition~\ref{sha-lifts-problem}(c), with $N=N_h$, and put $\chi'=\chi'_\ell(\psi)$ (Definition~\ref{sha-ordinary-deformation}).
\begin{enumerate}
\item $\psi=\psi_2\theta_1\kappa$ and $\theta_0=\psi_2\theta_1$; $\kappa^{\ell^s}=1$ for some $s\ge0$; and $f_0\ge1$, $4\mid\ell^{f_0}-1$, $\ell-1\mid\ell^{f_0}-1$ and $\theta_1^{\ell^{f_0}-1}=1$.
\item $\theta_1$ has odd conductor and order dividing $2o_{\psi_0}$; the $\ell$-component of $\chi'$ is trivial, and the order of $\chi'$ is prime to $\ell$.
\item Every prime $p\in T_b$ other than $\ell$ is odd, satisfies $p\equiv1\pmod\ell$, and divides $N_{\psi_0}$; and a prime $p\ne\ell$ lies in $T_b$ if and only if $\ell$ divides the order of the component $\psi_p$ of $\psi$ at $p$.
\item If $\psi$ has odd order, then its order divides the odd part of $o_{\psi_0}$, $\psi_2=1$, $c_2=0$, and the conductor of $\psi$ divides $N_{\psi_0}$.
\end{enumerate}
\end{lemma}
\begin{proof}
These are computations in the group of characters of $(\mathbf Z/N_h)^\times/H_h$.
(1) By definition $\psi=\theta_0\kappa$; the $2$-component of $\kappa$ is trivial, since $\kappa$ has odd order and $(\mathbf Z/2^m)^\times$ is a $2$-group, so the $2$-component of $\theta_0$ is $\psi_2$ and $\theta_0=\psi_2\theta_1$. The order of $\theta_1$ divides that of $\theta_0$, and $f_0\ge1$ is the order of $\ell$ modulo $\operatorname{lcm}(4,\operatorname{ord}\theta_1)$.
(2) The $2$-component of $\theta_1$ is trivial. Since $\psi$ is trivial on $H_h$, the component of $\theta_0$ at every odd prime has order dividing $2o_{\psi_0}$, as in the proof of Lemma~\ref{sha-lifts-host-data}. The character $\chi'=\theta_0(\theta_0)_\ell^{-1}$ has trivial $\ell$-component, and its order divides that of $\theta_0$.
(3) Let $p\in T_b$, $p\ne\ell$, so that $\kappa_p\ne1$. Then $p\ne2$, as in (1); $p$ is neither of the primes of the prime pair of $h$, where the components of $\psi$ have order at most $2$; so $p$ divides $N_{\psi_0}$. A nontrivial character of $(\mathbf Z/p^m)^\times$ of $\ell$-power order exists only if $\ell\mid p-1$. Finally $\psi_p=(\theta_0)_p\kappa_p$, and $\ell$ divides the order of $\psi_p$ exactly when $\kappa_p\ne1$.
(4) If $\psi$ has odd order, its $2$-component and its components at the primes of the prime pair are trivial, since they have $2$-power order or order at most $2$; so the conductor of $\psi$ divides $N_{\psi_0}$. Since $\psi$ is trivial on $H_h$, its order divides a power of $2$ times $o_{\psi_0}$; being odd, it divides the odd part of $o_{\psi_0}$.
\end{proof}
\leanref{Development/AutomorphicArithmetic/AutomorphicVarieties/LevelCohomology/Th1Facts.lean}{UpperTheory.V.X.th1Facts}
\leanref{Development/AutomorphicArithmetic/AutomorphicVarieties/HeightAndVanishing/OrdProjNDvd.lean}{UpperTheory.V.HV.hostTh1}
\leanref{Development/AutomorphicArithmetic/AutomorphicVarieties/HeightAndVanishing/HostSb.lean}{UpperTheory.V.HV.hostSb}
\leanref{Development/AutomorphicArithmetic/AutomorphicVarieties/HeightAndVanishing/HostOdd.lean}{UpperTheory.V.HV.hostOdd}

\begin{lemma}\label{sha-counts-siegel-stickelberger}
Let $\ell$ be an odd prime, let $\mathcal O$ be as in Definition~\ref{sha-character-parts}, let $F$ be a field of characteristic $\ell$ with a ring homomorphism $\mathrm{red}:\mathcal O\to F$, and write $G=\operatorname{Gal}(\bar{\mathbf Q}/\mathbf Q)$.
Let $N\ge1$, let $\chi\ne1$ be a Dirichlet character modulo $N$ with values in $\mathcal O$ whose order is prime to $\ell$, viewed as a character of $G$ (Definition~\ref{sha-character-parts}), and put $\gamma=\mathrm{red}\circ\chi$.
For a Dirichlet character $\varepsilon$ modulo $N$ put
\[
 B_1(\varepsilon;N)=\sum_{a=1}^{N}\varepsilon(a)\Big(\frac aN-\frac12\Big).
\]
Suppose that $\gamma(g)\ne\mathrm{red}(\chi_{\mathrm{cyc}}(g))$ for some $g\in G$, and that $B_1(\chi^{-1};N)$ has $\ell$-adic absolute value $1$.
Let $z:G\to F$ be a locally constant function with $z(gh)=z(g)+\gamma(g)z(h)$ for all $g,h\in G$, such that at every place $A$ of $\bar{\mathbf Q}$, with inertia group $I_A$, there is a subgroup $H\subseteq G$ for which $H\cap I_A\cap\ker\gamma$ has index prime to $\ell$ in $I_A\cap\ker\gamma$ and $z$ vanishes on $H\cap I_A\cap\ker\gamma$.
Then there is $m\in F$ with $z(g)=\gamma(g)m-m$ for all $g\in G$.
\end{lemma}
\begin{proof}
Let $U\subseteq G$ be the subgroup fixing the $N$-th roots of unity, and for $b\in(\mathbf Z/N)^\times$ let $\sigma_b\in G$ act on them by $\zeta\mapsto\zeta^b$, so that $\chi(\sigma_b)=\chi(b)$. Since $\chi$ factors through $G/U$, $\gamma$ is trivial on $U$.
\begin{enumerate}
\item[(1)] Choose $b_0$ with $\chi(b_0)\ne1$; then $N>2$. Moreover $\gamma(\sigma_{b_0})\ne1$: if $x\ne1$ is a root of unity of order $r$ prime to $\ell$ with $\mathrm{red}(x)=1$, reducing $\sum_{i<r}x^i=0$ gives $r=0$ in $F$, which is impossible.
\item[(2)] There is an odd integer $a\ge1$, prime to $N$, with $\mathrm{red}(\chi(a))\ne a$ in $F$. If $\ell\mid N$, let $g$ be as in the hypothesis and let $a$ be an odd integer prime to $N$ such that $g$ acts on the $N$-th roots of unity by $\zeta\mapsto\zeta^a$; then $\mathrm{red}(\chi(a))=\gamma(g)$, and $a\equiv\chi_{\mathrm{cyc}}(g)\pmod\ell$ since $\ell\mid N$, so $\mathrm{red}(\chi(a))\ne a$. If $\ell\nmid N$, take $a=\ell$, which is $0$ in $F$, while $\chi(\ell)$ is a unit.
\item[(3)] For $b\in(\mathbf Z/N)^\times$ let $\langle b\rangle\in\{0,\dots,N-1\}$ be its representative and $E_a(b)=\frac{a-1}2-\big\lfloor a\langle a^{-1}b\rangle/N\big\rfloor$, an integer since $a$ is odd. For a Dirichlet character $\varepsilon\ne1$ modulo $N$ with values in a field of characteristic $0$,
\[
 \sum_{b\in(\mathbf Z/N)^\times}E_a(b)\,\varepsilon(b)=\big(1-a\varepsilon(a)\big)B_1(\varepsilon;N).
\]
Indeed $N\lfloor a\langle a^{-1}b\rangle/N\rfloor=a\langle a^{-1}b\rangle-\langle b\rangle$; since $\sum_b\varepsilon(b)=0$, $NB_1(\varepsilon;N)=\sum_b\langle b\rangle\varepsilon(b)$; and the substitution $b=ab'$ gives $\sum_b\langle a^{-1}b\rangle\varepsilon(b)=\varepsilon(a)NB_1(\varepsilon;N)$.
With $\varepsilon=\chi^{-1}$, the element $s=\sum_bE_a(b)\chi^{-1}(b)$ of $\mathcal O$ is the product of $1-a\chi^{-1}(a)$, whose residue is nonzero by (2), and of $B_1(\chi^{-1};N)$, a unit; so $\mathrm{red}(s)\ne0$.
\item[(4)] $z$ vanishes on $U$. Let $A$ be a place, with $H$ as in the hypothesis. On $I_A\cap\ker\gamma$ the cocycle relation makes $z$ a homomorphism to $F$; it vanishes on a subgroup of index prime to $\ell$, and $F$ has characteristic $\ell$, so it vanishes on $I_A\cap\ker\gamma$. Hence the restriction of $z$ to $U\subseteq\ker\gamma$ is a locally constant homomorphism that vanishes on every inertia group. By class field theory it factors through the class group of $\mathbf Q(\zeta_N)$, and Stickelberger's theorem \cite[Theorem 6.10]{Was97}, as in Step~2 of the proof of Lemma~\ref{sha-stickelberger-annihilator}, gives $\sum_bE_a(b)\,z(\sigma_b^{-1}u\sigma_b)=0$ for $u\in U$. For $u\in U$ the cocycle relation gives $z(\sigma_b^{-1}u\sigma_b)=\gamma(\sigma_b)^{-1}z(u)=\mathrm{red}(\chi^{-1}(b))\,z(u)$. So $\mathrm{red}(s)z(u)=0$, and $z(u)=0$ by (3).
\item[(5)] Put $m=z(\sigma_{b_0})/(\gamma(\sigma_{b_0})-1)$, using (1). For $g\in G$ the commutator $g\sigma_{b_0}g^{-1}\sigma_{b_0}^{-1}$ lies in $U$, since $\operatorname{Gal}(\mathbf Q(\zeta_N)/\mathbf Q)$ is abelian. Evaluating $z$ on $g\sigma_{b_0}=(g\sigma_{b_0}g^{-1}\sigma_{b_0}^{-1})\sigma_{b_0}g$ with the cocycle relation and (4) gives $z(g)+\gamma(g)z(\sigma_{b_0})=z(\sigma_{b_0})+\gamma(\sigma_{b_0})z(g)$, so $z(g)=\gamma(g)m-m$.
\end{enumerate}
\end{proof}
\leanref{Development/AutomorphicArithmetic/AutomorphicVarieties/LevelCohomology/UnramifiedCocycleIsCoboundaryOfBernoulliUnit.lean}{UpperTheory.V.X.UnramifiedCocycleIsCoboundaryOfBernoulliUnit}
\leanref{Development/AutomorphicArithmetic/AutomorphicVarieties/LevelCohomology/ResidualBernoulliNonvanishing.lean}{UpperTheory.V.X.residual_bernoulli_nonvanishing}

\begin{definition}\label{sha-counts-siegel-pair}
Let $\ell$ be a prime and $d\ge1$ with $4\mid\ell^d-1$, and put $e=v_2(\ell^d-1)$.
A \emph{Bernoulli pair} for $(\ell,d)$ is a pair $(\varphi,\xi)$ of Dirichlet characters with values in a commutative ring, where $\varphi$ is primitive of odd conductor $f$ and has order dividing $\ell^d-1$, and $\xi$ is primitive of conductor $2^n$, such that $f\,4^{e}\le2^n$ and $\varphi(-1)\xi(-1)=-1$.
Its \emph{product} is the Dirichlet character $\varphi\xi$ modulo $f2^n$.
\end{definition}
\leanref{Development/AutomorphicArithmetic/AutomorphicVarieties/LevelCohomology/UnramifiedCocycleIsCoboundaryOfBernoulliUnit.lean}{UpperTheory.V.X.XData}
\leanref{Development/AutomorphicArithmetic/AutomorphicVarieties/LevelCohomology/UnramifiedCocycleIsCoboundaryOfBernoulliUnit.lean}{UpperTheory.V.X.prodX}

\begin{lemma}\label{sha-counts-siegel-bernoulli}
Let $\ell$ be a prime and $d\ge1$ with $4\mid\ell^d-1$, and let $(\varphi,\xi)$ be a Bernoulli pair for $(\ell,d)$ with values in $\bar{\mathbf Q}_\ell$ (Definition~\ref{sha-counts-siegel-pair}), with product $\chi$ modulo $f2^n$.
Then $B_1(\chi;f2^n)$, as in Lemma~\ref{sha-counts-siegel-stickelberger}, has $\ell$-adic absolute value $1$.
\end{lemma}
\begin{proof}
This is a computation, whose steps are as follows. Let $e=v_2(\ell^d-1)$ and $\ell^d=1+2^ea$ with $a$ odd; the size condition gives $n=m+e$ with $f2^e\le2^m$.
The character $\chi$ is odd, so $f2^nB_1(\chi;f2^n)=\sum_{0\le a'<f2^n}\varphi(a')\xi(a')a'$.
Let $B$ be the ring of $\ell$-integral elements $x$ of $\bar{\mathbf Q}_\ell$ whose residue satisfies $\bar x^{\ell^d}=\bar x$; it contains the values of $\varphi$ and $2^{-n}$.
Grouping the terms by the residue of $a'$ modulo $2^n$ gives $B_1(\chi;f2^n)=\sum_{u\in\mathbf Z/2^n}\xi(u)g(u)$ for a function $g$ with values in $B$ and a constant $c\in B$ such that $g(u)-g(u-f2^e)=2^ec$, minus $\varphi(\langle u\rangle)$ when $\langle u\rangle<f2^e$, where $\langle u\rangle\in\{0,\dots,2^n-1\}$ represents $u$.
Reducing modulo the maximal ideal, a direct computation with these relations, with $\xi(1+2^{n-1})=-1$, $\varphi(1)=1$, $\varphi(f2^e-1)\xi(-1)=-1$ and $\ell^d=1+2^ea$, shows that the residue of $\sum_u\xi(u)g(u)$ is nonzero.
\end{proof}
\leanref{Development/AutomorphicArithmetic/AutomorphicVarieties/LevelCohomology/UnramifiedCocycleIsCoboundaryOfBernoulliUnit.lean}{UpperTheory.V.X.LemmaX}
\leanref{Development/AutomorphicArithmetic/AutomorphicVarieties/LevelCohomology/LemmaX.lean}{UpperTheory.V.X.lemmaX}
\leanref{Development/AutomorphicArithmetic/AutomorphicVarieties/LevelCohomology/LxPartial.lean}{UpperTheory.V.X.lx_partial}
\leanref{Development/AutomorphicArithmetic/AutomorphicVarieties/LevelCohomology/LxCore.lean}{UpperTheory.V.X.lx_core}

\begin{lemma}\label{sha-counts-siegel-residual}
Let $\ell$ be a prime and $d\ge1$ with $4\mid\ell^d-1$, let $\mathcal O$ be as in Definition~\ref{sha-character-parts}, let $F$ be a field of characteristic $\ell$ with a ring homomorphism $\mathrm{red}:\mathcal O\to F$, and write $G=\operatorname{Gal}(\bar{\mathbf Q}/\mathbf Q)$.
Let $(\varphi,\xi)$ be a Bernoulli pair for $(\ell,d)$ with values in $\mathcal O$ (Definition~\ref{sha-counts-siegel-pair}), with product $\chi$, viewed as a character of $G$, and let $\gamma$ be $\mathrm{red}\circ\chi$ or its inverse.
Let $z:G\to F$ be a locally constant function with $z(gh)=z(g)+\gamma(g)z(h)$ for all $g,h\in G$, such that at every place $A$ of $\bar{\mathbf Q}$, with inertia group $I_A$, there is a subgroup $H\subseteq G$ for which $H\cap I_A\cap\ker\gamma$ has index prime to $\ell$ in $I_A\cap\ker\gamma$ and $z$ vanishes on $H\cap I_A\cap\ker\gamma$.
Then there is $m\in F$ with $z(g)=\gamma(g)m-m$ for all $g\in G$.
\end{lemma}
\begin{proof}
The pair $(\varphi^{-1},\xi^{-1})$ is again a Bernoulli pair for $(\ell,d)$, with product $\chi^{-1}$, so we may assume that $\gamma=\mathrm{red}\circ\chi$. We check the hypotheses of Lemma~\ref{sha-counts-siegel-stickelberger} for $\chi$, modulo $f2^n$.
The prime $\ell$ is odd, since $4\mid\ell^d-1$.
The order of $\chi$ is prime to $\ell$: $\chi^{(\ell^d-1)2^n}=1$, since $\varphi^{\ell^d-1}=1$ and $(\mathbf Z/2^n)^\times$ has order $2^{n-1}$.
Since $2^n\ge f4^e\ge4$ and $\xi$ is primitive, $\xi\ne1$; its values are $2$-power roots of unity, so $\xi(b)=-1$ for some $b$. By the Chinese remainder theorem there is an integer $a'$ with $a'\equiv1\pmod{\ell f}$ and $a'\equiv b\pmod{2^n}$, and an element $g\in G$ acting on the $\ell f2^n$-th roots of unity by $\zeta\mapsto\zeta^{a'}$.
Then $g$ fixes the $\ell$-th roots of unity, so $\mathrm{red}(\chi_{\mathrm{cyc}}(g))=1$, while $\chi(g)=\varphi(1)\xi(b)=-1$ and $\gamma(g)=-1\ne1$. In particular $\chi\ne1$.
Finally, $(\varphi^{-1},\xi^{-1})$ is a Bernoulli pair with values in $\bar{\mathbf Q}_\ell$, through $\mathcal O\subset\bar{\mathbf Q}_\ell$, so $B_1(\chi^{-1};f2^n)$ has absolute value $1$ by Lemma~\ref{sha-counts-siegel-bernoulli}.
Lemma~\ref{sha-counts-siegel-stickelberger} gives the claim.
\end{proof}
\leanref{Development/AutomorphicArithmetic/AutomorphicVarieties/LevelCohomology/UnramifiedCocycleIsCoboundaryOfBernoulliUnit.lean}{UpperTheory.V.X.CorX}
\leanref{Development/AutomorphicArithmetic/AutomorphicVarieties/LevelCohomology/ChrInvNeg.lean}{UpperTheory.V.X.corX}
\leanref{Development/AutomorphicArithmetic/AutomorphicVarieties/LevelCohomology/ChrInvNeg.lean}{main}

\begin{lemma}\label{sha-counts-siegel-cocycles}
Let $\ell$ be an odd prime, let $\mathcal O$ be as in Definition~\ref{sha-character-parts}, with residue field $\mathbf F$, write $G=\operatorname{Gal}(\bar{\mathbf Q}/\mathbf Q)$, and let $d\ge1$ with $4\mid\ell^d-1$; put $e=v_2(\ell^d-1)$.
Let $N,M,n\ge1$ with $M$ odd, and let $\omega$, $\theta$, $\chi$, $\xi_0$ and $\xi$ be Dirichlet characters with values in $\mathcal O$, modulo $\ell$, $N$, $N$, $M$ and $2^n$ respectively, viewed as characters of $G$, such that:
$\omega$, $\theta$ and $\xi_0$ have order dividing $\ell^d-1$; $\theta$ has odd conductor $f_\theta$; $\chi$ has conductor $2^j$ with $j<n$; $\xi$ is primitive of conductor $2^n$; $M\ell f_\theta4^{e}\le2^n$; and $\omega(-1)\chi(-1)\theta(-1)\xi_0(-1)\xi(-1)=-1$.
Let $\gamma:G\to\mathcal O^\times$ be a character whose residue $\bar\gamma$ is that of $\omega\chi\theta\xi_0\xi$.
Let $T$ be a finite set of primes, with a place $v_p$ of $\bar{\mathbf Q}$ above each $p\in T$ and decomposition group $D_p$ at $v_p$, and let $K\subseteq G$ be a normal subgroup of finite index prime to $\ell$.
Let $\gamma'$ be $\bar\gamma$ or $\bar\gamma^{-1}$, and let $z:G\to\mathbf F$ be a locally constant function with $z(gh)=z(g)+\gamma'(g)z(h)$ for all $g,h\in G$, which vanishes on $I_A\cap K$ for every place $A$ of $\bar{\mathbf Q}$ not above a prime of $T$, and which for each $p\in T$ has the form $z(g)=\gamma'(g)m_p-m_p$ on $D_p$ for some $m_p\in\mathbf F$.
Then there is $m\in\mathbf F$ with $z(g)=\gamma'(g)m-m$ for all $g\in G$.
\end{lemma}
\begin{proof}
\emph{A Bernoulli pair.} Let $\varphi_1$ be the primitive character inducing $\omega\theta$; its conductor divides $\ell f_\theta$, and its order divides $\ell^d-1$. Let $\varphi$ be the primitive character inducing $\xi_0\varphi_1$; its conductor $f$ divides $M\ell f_\theta$, so $f$ is odd, and its order divides $\ell^d-1$.
Let $\xi_1$ be the character of $(\mathbf Z/2^n)^\times$ that agrees with $\chi\xi$ on the integers prime to $2N$; it is primitive of conductor $2^n$: $\chi$ factors through $(\mathbf Z/2^{n-1})^\times$ because $j\le n-1$, so if $\chi\xi$ did, then so would $\xi$, which is primitive.
Then $f4^e\le M\ell f_\theta4^e\le2^n$ and $\varphi(-1)\xi_1(-1)=-1$, so $(\varphi,\xi_1)$ is a Bernoulli pair for $(\ell,d)$ (Definition~\ref{sha-counts-siegel-pair}), and $\varphi(a)\xi_1(a)=\omega(a)\chi(a)\theta(a)\xi_0(a)\xi(a)$ for every integer $a$ prime to $M\ell N2^n$.

\emph{The residue of $\gamma$.} Let $g\in G$, and let $a$ be an integer prime to $M\ell N2^n$ such that $g$ acts on the $M\ell N2^n$-th roots of unity by $\zeta\mapsto\zeta^a$. Every Dirichlet character $\varepsilon$ whose modulus divides $M\ell N2^n$ satisfies $\varepsilon(g)=\varepsilon(a)$. So $\bar\gamma(g)$ is the residue of $(\varphi\xi_1)(a)=(\varphi\xi_1)(g)$, and $\gamma'$ is the residue of the product of the Bernoulli pair $(\varphi,\xi_1)$ or its inverse.

\emph{Unramified at every place.} By Lemma~\ref{sha-counts-siegel-residual} applied to $(\varphi,\xi_1)$, it suffices to find, at every place $A$, a subgroup $H\subseteq G$ such that $H\cap I_A\cap\ker\gamma'$ has index prime to $\ell$ in $I_A\cap\ker\gamma'$ and $z$ vanishes on it.
If $A$ lies above no prime of $T$, take $H=K$: $z$ vanishes on $K\cap I_A$, and since $K$ is normal, the index of $K\cap I_A\cap\ker\gamma'$ in $I_A\cap\ker\gamma'$ divides the index of $K$.
If $A$ lies above $p\in T$, take $H=G$. There is $g\in G$ with $g^{-1}I_Ag\subseteq I_{v_p}\subseteq D_p$, since the places above $p$ are conjugate. For $x\in I_A\cap\ker\gamma'$, the cocycle relation gives $z(g^{-1}xg)=\gamma'(g)^{-1}z(x)$, and $z(g^{-1}xg)=\gamma'(g^{-1}xg)m_p-m_p=0$; so $z(x)=0$.
\end{proof}
\leanref{Development/AutomorphicArithmetic/AutomorphicVarieties/SymplecticGSp4/DataEx.lean}{dirX}
\leanref{Development/AutomorphicArithmetic/AutomorphicVarieties/SymplecticGSp4/DataEx.lean}{xvOf}
\leanref{Development/AutomorphicArithmetic/AutomorphicVarieties/LevelCohomology/HcPrimProd.lean}{UpperTheory.V.X.primProd}
\leanref{Development/AutomorphicArithmetic/AutomorphicVarieties/SymplecticGSp4/ExistsPrimitiveTwoPowTwistMul.lean}{UpperTheory.exists_primitive_two_pow_twist_mul}
\leanref{Development/AutomorphicArithmetic/AutomorphicVarieties/LevelCohomology/HcDirVal.lean}{UpperTheory.V.X.dirVal}
\leanref{Development/AutomorphicArithmetic/AutomorphicVarieties/LevelCohomology/HcXvan.lean}{UpperTheory.V.X.xvan}
\leanref{Development/AutomorphicArithmetic/AutomorphicVarieties/LevelCohomology/HcVanOf.lean}{UpperTheory.V.X.vanOf}

\begin{lemma}\label{sha-counts-siegel-datum}
Let $\psi_0$ be a primitive Dirichlet character, let $h$ be a host datum for $\psi_0$ that is large for $t$, let $\ell$ be an odd prime, let $\mathcal O$ be as in Definition~\ref{sha-character-parts}, let $0<k<t$, and let $\psi\in X_h(k)$ be such that $(\ell,k,\psi)$ is in case II or case III (Definition~\ref{sha-host-datum}, with $\psi_0$ in place of $\psi$).
Let $n_{\mathrm{II}}$ and $n_{\mathrm{III}}$ be as in Definition~\ref{sha-siegel-datum}(d), and suppose that $\mathcal O$ contains a primitive $2^{n_{\mathrm{II}}}$-th and a primitive $2^{n_{\mathrm{III}}}$-th root of unity.
Then there is a Siegel datum attached to $(h,k,\psi)$ (Definition~\ref{sha-siegel-datum}).
\end{lemma}
\begin{proof}
Let $\kappa$, $\theta_0$, $\psi_2$, $\theta_1$, $c_2$, $f_0$, $n_0$ and $T_b$ be attached to $\psi$ as in Definition~\ref{sha-lifts-problem}(c), with $N=N_h$, let $\chi'=\chi'_\ell(\psi)$, write $o=o_{\psi_0}$, and let $s$ be as in Lemma~\ref{sha-counts-siegel-host}(1).
Recall that $\psi(-1)=(-1)^k$.

\emph{Step 1: the choices.} We choose an integer $n$, an odd integer $M$, a primitive odd Dirichlet character $\xi'$ of conductor $2^n$ and a character $\xi_0$ modulo $M$, both with values in $\mathcal O$, and integers $a\ge0$ and $u$ odd, such that:
\begin{enumerate}
\item[(H1)] $k\ge2$ and $\chi'\ne1$;
\item[(H2)] $\xi_0(-1)=1$, $\xi_0^{\ell^{f_0}-1}=1$, and no prime of $T_b$ divides $M$;
\item[(H3)] $c_2<n$ and $M\ell f_{\theta_1}4^{n_0}\le2^n$, where $f_{\theta_1}$ is the conductor of $\theta_1$;
\item[(H4)] $\psi^{2^au}=\xi_0^{2^au}=1$, and $v^22^a\le2^n$ for every $v\in T_b$.
\end{enumerate}
A primitive odd character of conductor $2^n$ with values in $\mathcal O$ exists when $n\ge3$ and $\mathcal O$ contains a primitive $2^n$-th root of unity, since $(\mathbf Z/2^n)^\times$ is the product of $\{\pm1\}$ and a cyclic group of order $2^{n-2}$.

\emph{Case II.} Here $k\ge2$ and $\chi'$ has even order, so (H1) holds. Take $n=n_{\mathrm{II}}$, $M=1$, $\xi_0=1$, $a=c_2+1+v_2(o)$ and $u=o\,2^{-v_2(o)}\ell^s$.
By the definition of $n_{\mathrm{II}}$, $\ell N_h4^{n_0}\le2^{n}$ and $2^{c_2+6}o\ell^3N_{\psi_0}^2\le2^n$; in particular $n\ge c_2+6\ge3$, which gives $\xi'$ and the first part of (H3), and the second part holds since $f_{\theta_1}\le N_h$.
(H2) is clear. For (H4): $2^au=2^{c_2}\cdot2o\cdot\ell^s$, and $\psi=\psi_2\theta_1\kappa$ with $\psi_2^{2^{c_2}}=1$, since $\psi_2$ factors through $(\mathbf Z/2^{c_2})^\times$, $\theta_1^{2o}=1$ and $\kappa^{\ell^s}=1$ (Lemma~\ref{sha-counts-siegel-host}(1) and (2)); and for $v\in T_b$, $v\le\ell N_{\psi_0}$ by Lemma~\ref{sha-counts-siegel-host}(3), so $v^22^a\le(\ell N_{\psi_0})^2\,2^{c_2+1}o\le2^n$.

\emph{Case III.} Here $\ell\notin B_{\psi_0}$, $\ell-1\mid2k$, and $\psi\ne1$ has odd order. As in the proof of Lemma~\ref{sha-cases}, $k\ge2$, $\ell\nmid o$ and $\ell\nmid N_h$; and $\ell\ne p_1(\psi_0)$, since $p_1(\psi_0)$ is odd and the odd primes dividing $p_1(\psi_0)$ lie in $B_{\psi_0}$.
By Lemma~\ref{sha-counts-siegel-host}(4), the order of $\psi$ divides the odd part of $o$, $\psi_2=1$, $c_2=0$, and $f_\psi\mid N_{\psi_0}$. So $\ell$ does not divide the order of $\psi$, hence $\kappa=1$ and $\theta_0=\theta_1=\psi$; since $\ell\nmid N_h$, the $\ell$-component of $\psi$ is trivial and $\chi'=\psi\ne1$. This is (H1), and $T_b=\{\ell\}$.
Take $n=n_{\mathrm{III}}$, $M=p_1(\psi_0)$, $a=0$ and $u$ the order of $\psi$. Since $p_1(\psi_0)\equiv1$ modulo the odd part of $o$, $u$ divides $M-1$, and there is a character $\xi_0$ modulo the prime $M$ of order $u$; then $\psi^u=\xi_0^u=1$.
Then $\xi_0(-1)=1$ since $u$ is odd; $u=\operatorname{ord}\theta_1$ divides $\ell^{f_0}-1$ by Lemma~\ref{sha-counts-siegel-host}(1); and $\ell\nmid M$. This is (H2).
By the definition of $n_{\mathrm{III}}$, $\ell N_{\psi_0}p_1(\psi_0)4^{n_0}\le2^n$ and $8\ell^2\le2^n$, so $n\ge3$, which gives $\xi'$; (H3) holds since $c_2=0$ and $f_{\theta_1}=f_\psi\le N_{\psi_0}$; and (H4) holds since $T_b=\{\ell\}$ and $\ell^2\le2^n$.

\emph{Step 2: the datum.} Let $S$ be the set of primes dividing $2\ell N_hM$, put $\xi=\xi_0\xi'$, and view $\psi$, $\xi_0$, $\xi'$ and $\xi$ as characters of $G=\operatorname{Gal}(\bar{\mathbf Q}/\mathbf Q)$.
Lemma~\ref{sha-counts-siegel-auxiliary}, applied to $k$, $\psi$ and $\xi$, gives $\omega$, $c$, the places $v_p$ with Frobenius elements $\sigma_p$, and $\phi$, $g_0$ and $s_0$.
Let $D$ be the datum formed by $S$, $k$, $\psi$, $\xi$, $\omega$, $c$, the $v_p$, $\phi$, $g_0$, $s_0$ and the $\sigma_p$, and use the notation of Definition~\ref{sha-siegel-datum} for it; we check its conditions.
By uniqueness of the Teichm\"uller lift, $\omega$ is the Dirichlet character modulo $\ell$ with $\omega^{\ell-1}=1$ whose residue is $x\mapsto x\bmod\ell$, viewed as a character of $G$.
Since $\chi_{\mathrm{cyc}}$ and $\omega$ have the same residue and $\kappa$ has $\ell$-power order,
\[
 \bar\alpha=\bar\xi=\bar\xi_0\bar\xi',\qquad\bar\Psi=\bar\omega^k\bar\psi=\bar\omega^k\bar\psi_2\bar\theta_1.
\]

\emph{The set $T$.} Let $p\in S_b$, so that $p\ne\ell$ and $\ell$ divides the order of $\psi(I_p)$. On $I_p$, $\psi$ agrees with its component $\psi_p$, so $\ell$ divides the order of $\psi_p$, and $p\in T_b$ by Lemma~\ref{sha-counts-siegel-host}(3). So $T\subseteq T_b$, and every $v\in T$ is an odd prime with $v=\ell$ or $v\equiv1\pmod\ell$.

\emph{The conditions on $S$, $k$, $\psi$, $\xi$ and $\omega$.} $S$ contains $2$ and $\ell$, and $k\ge2$ by (H1). The characters $\psi$ and $\xi$ have open kernels and are unramified outside $S$, since their moduli $N_h$ and $M2^n$ have their prime factors in $S$. $\Psi(c)=\chi_{\mathrm{cyc}}(c)^k\psi(-1)=(-1)^k(-1)^k=1$. $\omega$ is as required by Lemma~\ref{sha-counts-siegel-auxiliary}(1). The order of $\xi$ divides $(\ell^{f_0}-1)2^{n-1}$, which is prime to $\ell$; $\xi(c)=\xi_0(-1)\xi'(-1)=-1$. For $v\in T$, $v$ is odd and does not divide $M$, by (H2) and $T\subseteq T_b$, so $\xi$ is trivial on $I_v$.

\emph{The elements $\sigma_v$.} They are Frobenius elements, $\chi_{\mathrm{cyc}}(\sigma_\ell)=1$ and $\phi$ is the limit required (Lemma~\ref{sha-counts-siegel-auxiliary}(2)).
Let $v\in T$ and $i\in\mathbf Z$, and suppose that $\overline{\xi\Psi^i}(\sigma_v)=1$.
First $\bar\omega(\sigma_v)=1$: for $v=\ell$, $\bar\omega(\sigma_\ell)=\bar\chi_{\mathrm{cyc}}(\sigma_\ell)=1$; for $v\ne\ell$, $\omega(\sigma_v)=\omega(v)=1$ since $v\equiv1\pmod\ell$.
So $\bar\xi_0(\sigma_v)\bar\xi'(\sigma_v)\bar\psi(\sigma_v)^i=1$, and raising to the power $2^au$ gives $\bar\xi'(\sigma_v)^{2^au}=1$, by (H4).
The values of $\xi'$ are $2^{n-1}$-th roots of unity, and distinct roots of unity of order prime to $\ell$ have distinct residues, so $\xi'(\sigma_v)^{2^au}=1$.
As $v$ is odd, $\xi'(\sigma_v)=\xi'(v)$, and Lemma~\ref{sha-counts-siegel-dyadic} gives $2^n<v^22^a$, which contradicts (H4), since $v\in T_b$.
So $\overline{\xi\Psi^i}(\sigma_v)\ne1$.

\emph{The element $g_0$.} Since $\bar\alpha=\bar\xi$, $\bar\alpha(\sigma_\ell)\ne1$ by the case $v=\ell$, $i=0$ above.
And $\bar\Psi\ne1$. The $\ell$-component of $\theta_0$ has order prime to $\ell$ and factors through some $(\mathbf Z/\ell^r)^\times$, so it factors through $(\mathbf Z/\ell)^\times$ and is a power $\omega^b$; as $\theta_0=(\theta_0)_\ell\chi'$, $\bar\Psi=\bar\omega^{k+b}\bar\chi'$.
If $\bar\Psi=1$, then $\omega^{k+b}\chi'$ and $1$ are characters of order prime to $\ell$ with the same residue, so $\omega^{k+b}\chi'=1$. The character $\chi'$ is trivial on $I_\ell$, by Lemma~\ref{sha-counts-siegel-host}(2), so $\omega^{k+b}$ is trivial on $I_\ell$, on which $\omega$ has order $\ell-1$; hence $\omega^{k+b}=1$ and $\chi'=1$, contradicting (H1).
So Lemma~\ref{sha-counts-siegel-auxiliary}(3) gives $\bar\alpha(g_0)\ne1$ and $\bar\Psi(g_0)\ne1$, and Lemma~\ref{sha-counts-siegel-auxiliary}(4) gives the conditions on $s_0$.

\emph{Residual cocycles.} Let $K=\ker\bar r_0$, a normal subgroup whose index divides the order of $(\mathbf F^\times)^4$, which is prime to $\ell$.
Apply Lemma~\ref{sha-counts-siegel-cocycles} with $d=f_0$, $N=N_h$, the given $M$, $n$ and $\xi_0$, $\xi'$ in place of $\xi$, and $T$, the $v_p$ and $K$:
\begin{itemize}
\item to $\gamma=\alpha$, with $(\omega,\theta,\chi)$ trivial, since $\bar\alpha=\bar\xi_0\bar\xi'$;
\item to $\gamma=\theta$, with $(\omega^k,\theta_1,\psi_2)$ in place of $(\omega,\theta,\chi)$, since $\bar\theta=\bar\alpha\bar\Psi=\bar\omega^k\bar\psi_2\bar\theta_1\bar\xi_0\bar\xi'$;
\item to $\gamma=\beta$, with $(\omega^k,\theta_1,\psi_2)$ and $\xi_0^{-1}$, $\xi'^{-1}$ in place of $\xi_0$, $\xi'$, since $\bar\beta=\bar\Psi\bar\alpha^{-1}$.
\end{itemize}
Its hypotheses hold: $4\mid\ell^{f_0}-1$; $\omega^k$, $\theta_1$ and $\xi_0^{\pm1}$ have order dividing $\ell^{f_0}-1$, since $\ell-1\mid\ell^{f_0}-1$ (Lemma~\ref{sha-counts-siegel-host}(1)) and by (H2); $\theta_1$ has odd conductor (Lemma~\ref{sha-counts-siegel-host}(2)); $\psi_2$ has conductor $2^{c_2}$ with $c_2<n$, and $\xi'^{\pm1}$ is primitive; the size condition is (H3), and for $\alpha$ it follows from (H3) since $f_{\theta_1}\ge1$; and the sign is $-1$ in each case, because $\omega(-1)=-1$ gives $\omega^k(-1)\psi_2(-1)\theta_1(-1)=(-1)^k\psi(-1)\kappa(-1)^{-1}=1$, $\kappa$ having odd order, while $\xi_0^{\pm1}(-1)=1$ and $\xi'^{\pm1}(-1)=-1$.
The lemma gives the condition on residual cocycles for $\gamma\in\{\alpha^{\pm1},\beta^{\pm1},\theta^{\pm1}\}$.

So $D$ is a Siegel datum. Its integer is $k$, its character $\psi$ is the given one, and in case II, $S$ is the set of primes dividing $2\ell N_h$ and $\xi=\xi'$ is primitive and odd of conductor $2^{n_{\mathrm{II}}}$; in case III, $S$ is the set of primes dividing $2\ell N_hp_1(\psi_0)$ and $\xi=\xi_0\xi'$ with $\xi_0$ of order $\operatorname{ord}\psi$ modulo $p_1(\psi_0)$. So $D$ is attached to $(h,k,\psi)$.
\end{proof}
\leanref{Development/AutomorphicArithmetic/AutomorphicVarieties/SymplecticGSp4/DataEx.lean}{UpperTheory.V.GradeZeroInterpolatingElement.dataEx}
\leanref{Development/AutomorphicArithmetic/AutomorphicVarieties/SymplecticGSp4/DataEx.lean}{core}
\leanref{Development/AutomorphicArithmetic/AutomorphicVarieties/SymplecticGSp4/DataEx.lean}{t_sub}
\leanref{Development/AutomorphicArithmetic/AutomorphicVarieties/SymplecticGSp4/DataEx.lean}{psi_ne}
\leanref{Development/AutomorphicArithmetic/AutomorphicVarieties/AutomorphicCases/SymplecticGSp4.lean}{UpperTheory.V.Aut.OfChar}
\leanref{Development/AutomorphicArithmetic/AutomorphicVarieties/SymplecticGSp4/ExistsPrimitiveOddDirichletCharacterTwoPow.lean}{UpperTheory.exists_primitive_odd_dirichletCharacter_two_pow}
\leanref{Development/AutomorphicArithmetic/AutomorphicVarieties/SymplecticGSp4/ExistsDirichletCharacterPrimeOfOrderOfDvd.lean}{UpperTheory.exists_dirichletCharacter_prime_of_orderOf_dvd}
\leanref{Development/AutomorphicArithmetic/AutomorphicVarieties/LevelCohomology/HcOmDir.lean}{UpperTheory.V.X.omDir}
\leanref{Development/AutomorphicArithmetic/AutomorphicVarieties/HeightAndVanishing/PsiBar.lean}{UpperTheory.V.HV.psiBar}
\leanref{Development/AutomorphicArithmetic/AutomorphicVarieties/LevelCohomology/HcSpKer.lean}{UpperTheory.V.X.spKer}

\subsection{Graded models of Siegel data}\label{subsection-sha-graded-models}

\begin{definition}\label{sha-counts-vd-model}
Let $D$ be a Siegel datum, with the notation of Definition~\ref{sha-siegel-datum}, and let $E$ be the fraction field of $\mathcal O$.
A $\bar{\mathbf Q}_\ell$-point $r$ of $D$ is \emph{irreducible} if the $r(g)$ span $M_4(\bar{\mathbf Q}_\ell)$.
For a commutative ring $K$, $t=(t_1,t_2)\in(K^\times)^2$ and $w=(w_1,w_2)\in\mathbf Z^2$ put $t^w=t_1^{w_1}t_2^{w_2}$ and $\tau(t)=\operatorname{diag}(t_1t_2,t_2,t_1,1)$; conjugation by $\tau(t)$ multiplies the entry $(i,j)$ of a $4\times4$ matrix by $t^w$, where $w$ is the weight of $(i,j)$.
\begin{enumerate}
\item Let $X$ be a finitely generated free abelian group and $R$ a complete local noetherian $\mathcal O$-algebra with residue field $\mathbf F$. A \emph{weighted presentation} of $R$ is a surjective homomorphism of $\mathcal O$-algebras $\pi:\mathcal O[[X_s:s\in\Sigma]]\to R$, with $\Sigma$ finite and weights $w_s\in X$ given to the variables, such that for every $f\in\ker\pi$ and $w\in X$ the sum of the terms of $f$ of weight $w$ lies in $\ker\pi$. Its \emph{parts} $R_w$ are the images of the power series of weight $w$, $R^{\mathrm{gr}}$ is the $\mathcal O$-subalgebra of the finite sums of elements of the parts, and $z_s=\pi(X_s)$. It has an \emph{origin} if the series in $\ker\pi$ have constant term $0$; the constant term then defines a homomorphism of $\mathcal O$-algebras $x:R\to\mathcal O$.
\item A \emph{graded model} of $D$ consists of: a complete local noetherian $\mathcal O$-algebra $R$ with residue field $\mathbf F$ that is a domain, with a weighted presentation with origin $x$ by variables whose weights lie in $\{\pm(1,0),\pm(-1,1),\pm(1,1),\pm(0,1),(0,0)\}$; integers $r\ge1$ and $r',m\ge0$; elements $F_0,\dots,F_r$ of $R_{(0,-1)}$, $E_1,\dots,E_{r'}$ of $R_{(0,1)}$ and $Z_1,\dots,Z_m$ of $R_{(0,0)}$ with $x(Z_i)=0$; a complete local noetherian $\mathcal O$-algebra $S$ with residue field $\mathbf F$ and a surjective homomorphism of $\mathcal O$-algebras $q:R\to S$ whose kernel is the ideal generated by the $F_i$, $E_i$ and $Z_i$; and a Siegel point $r^{\mathrm u}$ of $D$ with values in $S$. These data satisfy the following conditions.
\begin{enumerate}
\item[(i)] (Chart.) There is a homomorphism of $\mathcal O$-algebras $\varphi$ from $R^{\mathrm{gr}}$ to a ring of formal power series over $E$ in variables with weights in $\mathbf Z^2$, one each of weights $(1,0)$, $(-1,1)$ and $(1,1)$, $r'$ of weight $(0,1)$, $r$ of weight $(0,-1)$, $m$ of weight $(0,0)$ and no others, such that $\varphi$ maps $R_w$ into the series of weight $w$, the constant term of $\varphi(f)$ is $x(f)$, every variable is the only linear term of $\varphi(f)$ for some $f$, and $f\in\Sigma_x$ (Definition~\ref{sha-algebra-graded}, for the algebra $R^{\mathrm{gr}}$) whenever $x(f)=0$ and $\varphi(f)$ has no linear term.
\item[(ii)] (Generation.) Two homomorphisms of $\mathcal O$-algebras from $S$ to a complete local noetherian $\mathcal O$-algebra with residue field $\mathbf F$ that agree on the entries of the $r^{\mathrm u}(g)$ are equal.
\item[(iii)] (Entries.) For all $g\in G$ and $i,j$, the entry $(i,j)$ of $r^{\mathrm u}(g)$ is $q(f)$ for some $f\in R_w$, where $w$ is the weight of $(i,j)$; and $x(f)=r_0(g)_{ij}$ for every $f\in R$ with $q(f)=r^{\mathrm u}(g)_{ij}$.
\item[(iv)] (Derivations of classes.) Let $a\in\mathcal O$ be nonzero, and let $z:G\to\mathcal O/a$ be a locally constant function with $z(gh)=z(g)+\Psi(g)^{-1}z(h)$ for all $g,h\in G$, which vanishes at $g_0$ and at the elements $g$ with $\bar r_0(g)=1$ of the inertia groups at the places not above $T$, and such that for each $p\in T$ there is $m_p\in\mathcal O/a$ with $z(g)=\Psi(g)^{-1}m_p-m_p$ for $g\in D_p$. Then there is a derivation $\delta:R\to\mathcal O$ at $x$ (Definition~\ref{sha-algebra-graded}) that vanishes on $R_w$ for $w\ne(0,-1)$, maps $\ker q$ into $a\mathcal O$, and satisfies $\delta(f)\equiv\theta(g)z(g)$ modulo $a$ whenever $q(f)=r^{\mathrm u}(g)_{31}$, and $\delta(f)\equiv\Psi(g)z(g)$ modulo $a$ whenever $q(f)=r^{\mathrm u}(g)_{42}$.
\end{enumerate}
\item For a graded model of $D$, let $A_n=R_{(0,n)}$ for $n\in\mathbf Z$, let $A$ be the $\mathcal O$-subalgebra of $R$ of the finite sums of elements of the $A_n$, write $x$ also for the restriction of $x$ to $A$, let $I$ be the ideal of $A$ generated by the $F_i$, $E_i$ and $Z_i$, and let $q_A:A\to S$ be the restriction of $q$. Put $S_w=q(R_w)$ for $w\in\mathbf Z^2$, and let $S^{\mathrm{gr}}$ be the $\mathcal O$-subalgebra of the finite sums of elements of the $S_w$; by (iii) it contains the entries of the $r^{\mathrm u}(g)$. For a homomorphism of $\mathcal O$-algebras $\xi$ from $S^{\mathrm{gr}}$ to a commutative ring $K$, let $r_\xi:G\to M_4(K)$ be the homomorphism with $r_\xi(g)_{ij}=\xi(r^{\mathrm u}(g)_{ij})$.
\end{enumerate}
\end{definition}
\leanref{Development/AutomorphicArithmetic/AutomorphicVarieties/SymplecticGSp4/EntryMem.lean}{UpperTheory.V.GradeZeroInterpolatingElement.Model}
\leanref{Development/AutomorphicArithmetic/AutomorphicVarieties/SymplecticGSp4/EntryMem.lean}{UpperTheory.V.GradeZeroInterpolatingElement.Raw.ShaCoc}
\leanref{Development/AutomorphicArithmetic/AutomorphicVarieties/SymplecticGSp4/EntryMem.lean}{UpperTheory.V.GradeZeroInterpolatingElement.tDims}
\leanref{Development/AutomorphicArithmetic/AutomorphicVarieties/CommutativeAlgebra/ResMapBijective.lean}{UpperTheory.V.Alg.WPres}
\leanref{Development/AutomorphicArithmetic/AutomorphicVarieties/CommutativeAlgebra/ResMapBijective.lean}{UpperTheory.V.Alg.WPres.HasOrigin}
\leanref{Development/AutomorphicArithmetic/AutomorphicVarieties/CommutativeAlgebra/OrbitVanishingLocal.lean}{UpperTheory.V.Alg.SmoothPt}
\leanref{Development/AutomorphicArithmetic/AutomorphicVarieties/CommutativeAlgebra/OrbitVanishingLocal.lean}{UpperTheory.V.Alg.IsChart}

\begin{lemma}\label{sha-counts-vd-presentations}
Let $X$ be a finitely generated free abelian group, and let $R$ be a complete local noetherian $\mathcal O$-algebra with residue field $\mathbf F$, with a weighted presentation $\pi$ by variables $X_s$ of weights $w_s\in X$ (Definition~\ref{sha-counts-vd-model}(1)) such that $z_s\in\mathfrak m_R$ for all $s$. For an exponent vector $n$ write $X^n=\prod_sX_s^{n_s}$ and $z^n=\prod_sz_s^{n_s}$.
\begin{enumerate}
\item Let $J$ be an ideal of $R$ generated by elements $g_1,\dots,g_k$ with $g_i\in R_{c_i}$, and let $q:R\to R/J$ be the quotient map. Then $q\circ\pi$ is a weighted presentation of $R/J$, with parts $q(R_w)$.
\item The sum of the $R_w$ is direct, so $R^{\mathrm{gr}}=\bigoplus_wR_w$ is an $X$-graded $\mathcal O$-algebra. Let $M$ be a set of exponent vectors of weight $w$. There is a finite subset $\mu\subseteq M$ such that every power series all of whose monomials have exponents in $M$ is $\sum_{n\in\mu}b_nX^n$ with power series $b_n$ of weight $0$. In particular $R_w=\sum_{n\in\mu}R_0z^n$ for a finite set $\mu$ of exponent vectors of weight $w$, and $R^{\mathrm{gr}}$ is generated as an $\mathcal O$-algebra by $R_0$ and the $z_s$.
\item $R_0$ is a complete local noetherian $\mathcal O$-algebra with residue field $\mathbf F$, whose maximal ideal $R_0\cap\mathfrak m_R$ is generated by $\mathfrak m_{\mathcal O}$ and finitely many $z^n$ with $n\ne0$ of weight $0$.
\item Let $B$ be a complete local noetherian $\mathcal O$-algebra with residue field $\mathbf F$, let $i:B\to R$ be a homomorphism of $\mathcal O$-algebras, and let $Z\subseteq\mathfrak m_R$ be such that two homomorphisms of $\mathcal O$-algebras $R\to\mathbf F[\varepsilon]$, $\varepsilon^2=0$, that agree on $i(B)$ and on $Z$ are equal. Then $\mathfrak m_R$ lies in the ideal generated by $Z$, $i(\mathfrak m_B)$ and $\mathfrak m_R^2$. If finitely many elements $z'_t\in\mathfrak m_R$, with $z'_t\in R_{w'_t}$, generate an ideal that, with $\mathfrak m_{\mathcal O}R$ and $\mathfrak m_R^2$, contains $\mathfrak m_R$, then $R$ has a weighted presentation by variables $X'_t$ of weights $w'_t$ with $X'_t\mapsto z'_t$.
\item A second weighted presentation of $R$ whose variables map into $\mathfrak m_R$ and into the parts $R_w$ of their weights $w$ has the same parts as the first.
\item Suppose that the presentation has an origin $x$. Every derivation of $R^{\mathrm{gr}}$ at $x$ with values in an $\mathcal O$-module extends uniquely to a derivation of $R$ at $x$, and a derivation of $R$ at $x$ that vanishes on the $z_s$ is $0$. For every $w$, the cotangent lattice $L_w$ of $R^{\mathrm{gr}}$ at $x$ (Definition~\ref{sha-algebra-graded}) is spanned by the images of the $z_s$ with $w_s=w$.
\end{enumerate}
\end{lemma}
\leanref{Development/AutomorphicArithmetic/AutomorphicVarieties/CommutativeAlgebra/HomogSpan.lean}{UpperTheory.V.Alg.homog_span}
\leanref{Development/AutomorphicArithmetic/AutomorphicVarieties/CommutativeAlgebra/GrInternal.lean}{UpperTheory.V.Alg.gr_internal}
\leanref{Development/AutomorphicArithmetic/AutomorphicVarieties/CommutativeAlgebra/PartSpan.lean}{UpperTheory.V.Alg.part_span}
\leanref{Development/AutomorphicArithmetic/AutomorphicVarieties/CommutativeAlgebra/GrGen.lean}{UpperTheory.V.Alg.gr_gen}
\leanref{Development/AutomorphicArithmetic/AutomorphicVarieties/CommutativeAlgebra/Gordan.lean}{UpperTheory.V.Alg.gordan}
\leanref{Development/AutomorphicArithmetic/AutomorphicVarieties/CommutativeAlgebra/WPresGrZeroIsCOOfMemMaximalIdeal.lean}{UpperTheory.V.Alg.wPres_gr_zero_isCO_of_mem_maximalIdeal}
\leanref{Development/AutomorphicArithmetic/AutomorphicVarieties/CommutativeAlgebra/WPresMaximalIdealGrZeroEqSupSpanMonomials.lean}{UpperTheory.V.Alg.wPres_maximalIdeal_gr_zero_eq_sup_span_monomials}
\leanref{Development/AutomorphicArithmetic/AutomorphicVarieties/CommutativeAlgebra/CoSpan.lean}{UpperTheory.V.Alg.co_span}
\leanref{Development/AutomorphicArithmetic/AutomorphicVarieties/CommutativeAlgebra/PresNew.lean}{UpperTheory.V.Alg.pres_new}
\leanref{Development/AutomorphicArithmetic/AutomorphicVarieties/CommutativeAlgebra/PresIndep.lean}{UpperTheory.V.Alg.pres_indep}
\leanref{Development/AutomorphicArithmetic/AutomorphicVarieties/CommutativeAlgebra/DerExtend.lean}{UpperTheory.V.Alg.der_extend}
\leanref{Development/AutomorphicArithmetic/AutomorphicVarieties/CommutativeAlgebra/CotFg.lean}{UpperTheory.V.Alg.cot_fg}
\leanref{Development/AutomorphicArithmetic/AutomorphicVarieties/CommutativeAlgebra/CoEvLow.lean}{UpperTheory.V.Alg.co_ev}
\leanref{Development/AutomorphicArithmetic/AutomorphicVarieties/CommutativeAlgebra/CoSurj.lean}{UpperTheory.V.Alg.co_surj}
\leanref{Development/AutomorphicArithmetic/AutomorphicVarieties/CommutativeAlgebra/PartClosed.lean}{UpperTheory.V.Alg.part_closed}
\leanref{Development/AutomorphicArithmetic/AutomorphicVarieties/CommutativeAlgebra/MapPart.lean}{UpperTheory.V.Alg.map_part}
\leanref{Development/AutomorphicArithmetic/IntegralCasts/CapelliIdentities/DensePiDecomposition.lean}{UpperTheory.Reformulated.B.AdicDensityOfPresentationSubalgebra.pi_dense}
\leanref{Development/AutomorphicArithmetic/AutomorphicVarieties/CommutativeAlgebra/KerPow.lean}{UpperTheory.V.Alg.ker_pow}
\begin{proof}
For a power series $h$ and $w\in X$ let $h_w$ be the sum of the terms of $h$ of weight $w$.

(a) Let $f$ be a power series with $q(\pi(f))=0$. Then $\pi(f)=\sum_ia_ig_i$ with $a_i=\pi(G_i)$, and $g_i=\pi(e_i)$ with $e_i$ of weight $c_i$. So $f-\sum_iG_ie_i\in\ker\pi$, and its terms of weight $w$, which form $f_w-\sum_i(G_i)_{w-c_i}e_i$, lie in $\ker\pi$. Hence $\pi(f_w)=\sum_i\pi((G_i)_{w-c_i})g_i\in J$, and $f_w$ lies in the kernel of $q\circ\pi$. The parts of $q\circ\pi$ are the images of the series of weight $w$, that is, the $q(R_w)$.

(b) If $\sum_w\pi(f^{(w)})=0$ for finitely many series $f^{(w)}$ of weight $w$, then $\sum_wf^{(w)}\in\ker\pi$, and its terms of weight $w$ form $f^{(w)}$; so each $\pi(f^{(w)})$ is $0$. By Dickson's lemma the elements of $M$ that are minimal for the componentwise order form a finite set $\mu$, and every $m\in M$ is $n+(m-n)$ with $n\in\mu$, $n\le m$, and $m-n$ of weight $0$; grouping the terms of a series according to a chosen $n\le m$ gives the second statement. The third follows by applying $\pi$ with $M$ the set of all exponent vectors of weight $w$.

(c) By Gordan's lemma the monoid of the exponent vectors of weight $0$ is generated by finitely many $n_1,\dots,n_k$, all nonzero. The homomorphism of $\mathcal O$-algebras $\mathcal O[[Y_1,\dots,Y_k]]\to\mathcal O[[X_s:s\in\Sigma]]$ with $Y_i\mapsto X^{n_i}$ maps onto the series of weight $0$. So $R_0$ is the image of a power series ring over $\mathcal O$ under a homomorphism of $\mathcal O$-algebras: it is a complete local noetherian $\mathcal O$-algebra with residue field $\mathbf F$, and its maximal ideal is generated by $\mathfrak m_{\mathcal O}$ and the $z^{n_i}$, which lie in $\mathfrak m_R$; so it is $R_0\cap\mathfrak m_R$.

(d) Let $J$ be the ideal generated by $Z$, $i(\mathfrak m_B)$ and $\mathfrak m_R^2$; it lies in $\mathfrak m_R$ and contains $\mathfrak m_{\mathcal O}R$, since $i$ is a homomorphism of $\mathcal O$-algebras. Suppose that $y\in\mathfrak m_R$ is not in $J$, and let $N=J+\mathcal O\cdot1$. Then $y\notin N$: if $y=y'+c$ with $y'\in J$ and $c\in\mathcal O$, then $c\in\mathfrak m_{\mathcal O}$, because $y,y'\in\mathfrak m_R$, and $y\in J$. The $\mathcal O$-module $R/N$ is killed by $\mathfrak m_{\mathcal O}$, so there is an $\mathbf F$-linear map $\lambda:R/N\to\mathbf F$ with $\lambda(y)\ne0$; let $D(f)=\lambda(f\bmod N)$ and let $\bar f$ be the residue of $f\in R$. For $f,f'\in R$ and $c,c'\in\mathcal O$ with $f-c,f'-c'\in\mathfrak m_R$, the product $(f-c)(f'-c')$ lies in $\mathfrak m_R^2\subseteq N$, and this gives $D(ff')=\bar fD(f')+\bar f'D(f)$. So $f\mapsto\bar f+\varepsilon D(f)$ is a homomorphism of $\mathcal O$-algebras $R\to\mathbf F[\varepsilon]$. It agrees with $f\mapsto\bar f$ on $Z$ and on $i(B)$, because $D$ vanishes on $J$ and on $\mathcal O\cdot1$, and differs from it at $y$, which contradicts the hypothesis.

For the second statement of (d), let $P=\mathcal O[[X_s:s\in\Sigma]]$, with maximal ideal $\mathfrak n$, and choose for each $t$ a series $f_t$ of weight $w'_t$ with $\pi(f_t)=z'_t$. The constant term of $f_t$ lies in $\mathfrak m_{\mathcal O}$, since otherwise $f_t$ and $z'_t$ would be units; so $f_t\in\mathfrak n$. As $P$ is $\mathfrak n$-adically complete, substituting $f_t$ for $X'_t$ is a homomorphism of $\mathcal O$-algebras $\Phi:\mathcal O[[X'_t]]\to P$, and $\Phi$ maps every series of order at least $N$ into $\mathfrak n^N$. Put $\pi'=\pi\circ\Phi$, so that $\pi'(X'_t)=z'_t$, and give $X'_t$ the weight $w'_t$.
\begin{enumerate}
\item[(i)] $\pi'$ is surjective. Let $\mathfrak n'$ be the maximal ideal of $\mathcal O[[X'_t]]$. The ideal $\pi'(\mathfrak n')R$ lies in $\mathfrak m_R$ and contains the $z'_t$ and $\mathfrak m_{\mathcal O}R$, so $\mathfrak m_R\subseteq\pi'(\mathfrak n')R+\mathfrak m_R^2$ by hypothesis. As $\mathfrak m_R$ is finitely generated, Nakayama's lemma \cite[Tag 00DV]{Sta} gives $\pi'(\mathfrak n')R=\mathfrak m_R$. So $R$, a module over $\mathcal O[[X'_t]]$ through $\pi'$, is separated for the $\mathfrak n'$-adic topology, because it is $\mathfrak m_R$-adically complete, and $R/\mathfrak n'R=\mathbf F$ is generated by the image of $1$. Since $\mathcal O[[X'_t]]$ is $\mathfrak n'$-adically complete, $R$ is generated by $1$ \cite[Theorem~8.4]{Mat86}, \cite[Tag 031D]{Sta}: $\pi'$ is surjective.
\item[(ii)] $\Phi(g_w)=\Phi(g)_w$ for every series $g$ in the $X'_t$ and every $w$, where $g_w$ is the sum of the terms of $g$ of weight $w$. For a polynomial $g$ this holds because $\Phi(X'^n)=\prod_tf_t^{n_t}$ has weight $\sum_tn_tw'_t$. In general let $g^{<N}$ be the sum of the terms of $g$ of degree less than $N$. Then $(g_w)^{<N}=(g^{<N})_w$, and $\Phi(g)-\Phi(g^{<N})$ and $\Phi(g_w)-\Phi((g_w)^{<N})$ lie in $\mathfrak n^N$. By induction on $N$, the coefficient of a monomial of degree $d$ in an element of $\mathfrak n^N$ lies in $\mathfrak m_{\mathcal O}^{N-d}$. So the coefficients of a monomial of degree $d$ in $\Phi(g_w)$ and in $\Phi(g)_w$ are congruent modulo $\mathfrak m_{\mathcal O}^{N-d}$ for every $N$, hence equal, since $\bigcap_k\mathfrak m_{\mathcal O}^k=0$.
\item[(iii)] Let $g\in\ker\pi'$. Then $\Phi(g)\in\ker\pi$, so $\Phi(g)_w\in\ker\pi$, and $\pi'(g_w)=\pi(\Phi(g)_w)=0$ by (ii). So $\pi'$ is a weighted presentation by the $X'_t$, with weights $w'_t$.
\end{enumerate}

(e) Let $\pi'$ be the second presentation, by variables $X'_t$ of weights $w'_t$ with $z'_t=\pi'(X'_t)\in\mathfrak m_R\cap R_{w'_t}$, and let $R'_w$ be its parts. For a series $h$ in the $X_s$, $h_w$ is as above.
\begin{enumerate}
\item[(i)] If $y\in R_w+\mathfrak m_R^N$ for every $N$, then $y\in R_w$; the same holds for the parts $R'_w$. Let $\mathfrak n_0$ be the ideal of $P$ generated by the $X_s$ and $\mathfrak m_{\mathcal O}$. Then $\pi(\mathfrak n_0)=\mathfrak m_R$, since $h-h(0)$ lies in the ideal generated by the $X_s$ for every $h\in P$, and $h(0)\in\mathfrak m_{\mathcal O}$ if $\pi(h)\in\mathfrak m_R$. The ideal $\mathfrak n_0^N$ is generated by products of $N$ of these generators, each a series of some weight, and if $g=\sum_ia_ie_i$ with $e_i$ of weight $c_i$, then $g_w=\sum_i(a_i)_{w-c_i}e_i$; so $g_w\in\mathfrak n_0^N$ whenever $g\in\mathfrak n_0^N$. Write $y=\pi(f)$, and for $N\ge0$ write $y=\pi(e)+\pi(g)$ with $e$ of weight $w$ and $g\in\mathfrak n_0^N$. The series $f-e-g$ lies in $\ker\pi$, and so does its component $f_w-e-g_w$ of weight $w$. Subtracting, $\pi(f-f_w)=\pi(g-g_w)\in\mathfrak m_R^N$ for every $N$. By Krull's intersection theorem \cite[Theorem~8.10]{Mat86}, \cite[Tag 00IP]{Sta}, $\pi(f-f_w)=0$, so $y=\pi(f_w)\in R_w$.
\item[(ii)] $R'_w\subseteq R_w$. Let $y=\pi'(g)$ with $g$ of weight $w$, and let $g^{<N}$ be as in (d). Each term of $g^{<N}$ is $cX'^n$ with $c\in\mathcal O$ and $n$ of weight $w$, and $\pi'(cX'^n)=c\prod_tz'^{n_t}_t$ lies in $R_w$, because $R_vR_{v'}\subseteq R_{v+v'}$. As the $z'_t$ lie in $\mathfrak m_R$, $\pi'(g-g^{<N})\in\mathfrak m_R^N$. So $y\in R_w+\mathfrak m_R^N$ for every $N$, and $y\in R_w$ by (i).
\item[(iii)] $R_w\subseteq R'_w$. Let $y=\pi(g)$ with $g$ of weight $w$, and let $N\ge0$; by (i) for $\pi'$ it suffices to find an element of $R'_w$ congruent to $y$ modulo $\mathfrak m_R^N$. First, $y\equiv a$ modulo $\mathfrak m_R^N$ for some $a\in R'^{\mathrm{gr}}$. This holds for every element of $R$, by induction on $N$: every series in the $X'_t$ is $c+\sum_tX'_th_t$ with $c\in\mathcal O$, so every element of $R$ is $c+\sum_tz'_t\pi'(h_t)$, and if $\pi'(h_t)\equiv a_t$ modulo $\mathfrak m_R^N$ with $a_t\in R'^{\mathrm{gr}}$, then $c+\sum_tz'_ta_t$ lies in $R'^{\mathrm{gr}}$ and is congruent to it modulo $\mathfrak m_R^{N+1}$. Write $a=\sum_va_v$, a finite sum with $a_v\in R'_v$; by (ii), $a_v=\pi(e^{(v)})$ with $e^{(v)}$ of weight $v$, and $e=\sum_ve^{(v)}$ has $\pi(e)=a$ and $\pi(e_w)=a_w$. Next, $\mathfrak m_R$ is generated by the $z_s$ and $\mathfrak m_{\mathcal O}$, as in (i), so $\mathfrak m_R^N$ is generated by elements $\pi(e')$ of $\mathfrak m_R^N$ with $e'$ a series of some weight $c$; for such $e'$ and every series $b$, $(be')_w=b_{w-c}e'$, so $\pi((be')_w)\in\mathfrak m_R^N$. Hence $y-a=\pi(h)$ for a series $h$ with $\pi(h_w)\in\mathfrak m_R^N$. The series $g-e-h$ lies in $\ker\pi$, and so does its component $g-e_w-h_w$; so $y-a_w=\pi(h_w)\in\mathfrak m_R^N$, with $a_w\in R'_w$.
\end{enumerate}

(f) Let $\mathfrak k$ be the kernel of $x$ on $R$, and $\mathfrak k^{\mathrm{gr}}=\mathfrak k\cap R^{\mathrm{gr}}$ its kernel on $R^{\mathrm{gr}}$. Since $x(\pi(h))$ is the constant term of $h$, $\mathfrak k$ is the image of the series with constant term $0$, and the $z_s$ lie in $\mathfrak k^{\mathrm{gr}}$.
\begin{enumerate}
\item[(i)] For $n\ge0$, $\mathfrak k^n\cap R^{\mathrm{gr}}=(\mathfrak k^{\mathrm{gr}})^n$, and every element of $R$ is congruent modulo $\mathfrak k^n$ to an element of $R^{\mathrm{gr}}$. Clearly $(\mathfrak k^{\mathrm{gr}})^n\subseteq\mathfrak k^n\cap R^{\mathrm{gr}}$. An element of $\mathfrak k^n$ is a sum of products of $n$ elements of $\mathfrak k$, so it is $\pi(h)$ for a series $h$ of order at least $n$. Let $y=\pi(h)$ be such an element that lies in $R^{\mathrm{gr}}$, so that also $y=\pi(e)$ for a finite sum $e$ of series of weights in a finite set $W_0$. Since $h-e\in\ker\pi$, $\pi(h_w)=\pi(e_w)$ for every $w$, and $y=\sum_{w\in W_0}\pi(h_w)$. Each $h_w$ has weight $w$ and only terms of degree at least $n$, so by (b), applied to the exponent vectors of weight $w$ and degree at least $n$, $h_w=\sum_{m\in\mu}b_mX^m$ with $b_m$ of weight $0$ and $m$ of degree at least $n$; then $\pi(h_w)=\sum_m\pi(b_m)z^m\in(\mathfrak k^{\mathrm{gr}})^n$. For the second statement let $f=\pi(h)$, and let $h^{<n}$ be the sum of the terms of $h$ of degree less than $n$. Then $\pi(h^{<n})\in R^{\mathrm{gr}}$, as each of its terms has a weight, and $h-h^{<n}$ lies in the ideal generated by the finitely many monomials of degree $n$, so $f-\pi(h^{<n})$ lies in the ideal generated by the products of $n$ of the $z_s$, which is contained in $\mathfrak k^n$.
\item[(ii)] Let $D$ be a derivation of $R^{\mathrm{gr}}$ at $x$ with values in an $\mathcal O$-module $Q$. From $D(1)=D(1\cdot1)=2D(1)$, $D$ vanishes on $\mathcal O\cdot1$, and by the product rule it vanishes on $(\mathfrak k^{\mathrm{gr}})^2$. For $f\in R$ choose $y_f\in R^{\mathrm{gr}}$ with $f-y_f\in\mathfrak k^2$, by (i), and put $D'(f)=D(y_f)$. If $a\in R^{\mathrm{gr}}$ and $f-a\in\mathfrak k^2$, then $y_f-a\in\mathfrak k^2\cap R^{\mathrm{gr}}=(\mathfrak k^{\mathrm{gr}})^2$, so $D(y_f)=D(a)$. Hence $D'$ is $\mathcal O$-linear and agrees with $D$ on $R^{\mathrm{gr}}$. For $f,f'\in R$, $ff'-y_fy_{f'}=(f-y_f)f'+y_f(f'-y_{f'})\in\mathfrak k^2$ and $x(y_f)=x(f)$, so $D'(ff')=D(y_fy_{f'})=x(f)D'(f')+x(f')D'(f)$. A derivation $E$ of $R$ at $x$ vanishes on $\mathfrak k^2$, so $E(f)=E(y_f)$; hence $D'$ is the only derivation of $R$ at $x$ extending $D$.
\item[(iii)] Let $D$ be a derivation of $R$ at $x$ that vanishes on the $z_s$, and let $f=\pi(h)$. Then $h-h(0)=\sum_sX_sh_s$ for some series $h_s$, so $f=h(0)+\sum_s\pi(h_s)z_s$, and $D(\pi(h_s)z_s)=x(\pi(h_s))D(z_s)+x(z_s)D(\pi(h_s))=0$, as $x(z_s)=0$. Since $D$ vanishes on $\mathcal O\cdot1$, $D(f)=0$.
\item[(iv)] The cotangent lattice $L_w$ is the image of $R_w\cap\mathfrak k$ in $R^{\mathrm{gr}}/\Sigma_x$, so it contains the images of the $z_s$ with $w_s=w$. Conversely let $y\in R_w\cap\mathfrak k$, and write $y=\pi(h)$ with $h$ of weight $w$; its constant term is $x(y)=0$. Let $c_s$ be the coefficient of $X_s$ in $h$; it is $0$ unless $w_s=w$. The series $h-\sum_{w_s=w}c_sX_s$ has only terms of weight $w$ and degree at least $2$, so by (b) it is $\sum_{m\in\mu}b_mX^m$ with $b_m$ of weight $0$ and $m$ of degree at least $2$. Hence $y-\sum_{w_s=w}c_sz_s=\sum_m\pi(b_m)z^m$ lies in $(\mathfrak k^{\mathrm{gr}})^2\subseteq\Sigma_x$, and the image of $y$ is $\sum_{w_s=w}c_s\bar z_s$, where $\bar z_s$ is the image of $z_s$.
\end{enumerate}
\end{proof}

\begin{lemma}\label{sha-counts-vd-extension}
Let $X$, $R$ and its weighted presentation be as in Lemma~\ref{sha-counts-vd-presentations}, with $z_s\in\mathfrak m_R$ for all $s$, and let $|\cdot|$ be the absolute value of $\bar{\mathbf Q}_\ell$.
Let $K$ be $\bar{\mathbf Q}_\ell$ or $\bar{\mathbf Q}_\ell[\varepsilon]$ with $\varepsilon^2=0$; write $y=y^\circ+\varepsilon y'$ for $y\in K$, with $y^\circ,y'\in\bar{\mathbf Q}_\ell$, and $y'=0$ if $K=\bar{\mathbf Q}_\ell$.
Let $\xi:R^{\mathrm{gr}}\to K$ be a homomorphism of $\mathcal O$-algebras such that $\xi(R_0)$ is a finitely generated $\mathcal O$-module.
\begin{enumerate}
\item If $|\xi(f)^\circ|\le1$ for every $f\in R^{\mathrm{gr}}$ and $|\xi(z_s)^\circ|<1$ for every $s$, then $\xi$ is the restriction of a homomorphism of $\mathcal O$-algebras $R\to K$.
\item If $K=\bar{\mathbf Q}_\ell$, there are a homomorphism $t:X\to\bar{\mathbf Q}_\ell^\times$ and a homomorphism of $\mathcal O$-algebras $\Xi:R\to\bar{\mathbf Q}_\ell$ whose image is finite over $\mathcal O$, such that $\Xi(f)=t(w)\xi(f)$ for every $w$ and $f\in R_w$.
\end{enumerate}
\end{lemma}
\leanref{Development/AutomorphicArithmetic/AutomorphicVarieties/CommutativeAlgebra/PointsFinite.lean}{UpperTheory.V.Alg.points_finite}
\leanref{Development/AutomorphicArithmetic/AutomorphicVarieties/CommutativeAlgebra/Carver.lean}{UpperTheory.V.Alg.carver}
\begin{proof}
Let $\varpi$ be a uniformizer of $\mathcal O$. An element $y$ of $K$ with $|y^\circ|\le1$ is integral over $\mathcal O$, since $y^\circ$ is integral over $\mathbf Z_\ell$ and $\varepsilon y'$ is nilpotent.

\emph{The torus element in (2).} Let $n\ne0$ be an exponent vector of weight $0$ with $u=\xi(z^n)\ne0$. Then $u^{-1}$ is not integral over $\mathcal O$. Otherwise $p(u^{-1})=0$ for a monic $p\in\mathcal O[T]$ of some degree $d$, and $p^*(T)=T^dp(1/T)$, whose constant term is $1$, has $p^*(u)=0$. The series $p^*(X^n)$ has weight $0$ and constant term $1$, so it has an inverse of weight $0$; hence $p^*(z^n)$ is a unit of $R^{\mathrm{gr}}$, while $\xi(p^*(z^n))=p^*(u)=0$, a contradiction. So $|\xi(z^n)|<1$ for every exponent vector $n\ne0$ of weight $0$. Let $v=-\log|\cdot|$, whose values on $\bar{\mathbf Q}_\ell^\times$ are dense in $\mathbf R$. By Fourier-Motzkin elimination, a theorem of the alternative for the linear inequalities $v(t(w_s))+v(\xi(z_s))>0$ in the homomorphism $v\circ t$, one for each $s$ with $\xi(z_s)\ne0$, as in the proof of Lemma~\ref{sha-flagged-weight-zero}, there is $t$ with $|t(w_s)\xi(z_s)|<1$ for all $s$: the alternative would be a relation $\sum_sn_sw_s=0$ with $n\ne0$ and $n_s\ge0$ supported on these $s$, and with $\sum_sn_sv(\xi(z_s))=v(\xi(z^n))\le0$.

\emph{The extension.} In case (1) put $t=1$. Let $y_s=t(w_s)\xi(z_s)$, so that $|y_s^\circ|<1$ for all $s$, and choose $e_s\ge1$ with $|(y_s^{e_s})^\circ|\le|\varpi|$. Let $B\subseteq K$ be the $\mathcal O$-algebra generated by a finite set of generators of the $\mathcal O$-module $\xi(R_0)$, the $y_s$ and the $y_s^{e_s}/\varpi$. These generators are integral over $\mathcal O$, so $B$ is finite over $\mathcal O$; it is complete for the radical $J$ of $\varpi B$, a power of which lies in $\varpi B$, and $y_s\in J$. Let $g:\mathcal O[[X_s:s\in\Sigma]]\to B$ be the continuous homomorphism of $\mathcal O$-algebras with $g(X_s)=y_s$.

Then $g(f)=t(w)\xi(\pi(f))$ for every $w$ and every series $f$ of weight $w$. For $w=0$, induction on $N$ gives $g(f)-\xi(\pi(f))\in J^N$; for $N=0$ both terms lie in $B$. By Lemma~\ref{sha-counts-vd-presentations}(b), applied to the nonzero exponent vectors of weight $0$, $f-f(0)=\sum_{n\in\mu}b_nX^n$ with $b_n$ of weight $0$ and $\mu$ finite. Since $n$ has weight $0$, $y^n=\xi(z^n)$, so $g(f)-\xi(\pi(f))=\sum_n\big(g(b_n)-\xi(\pi(b_n))\big)y^n\in J^{N}J$. As $B$ is $J$-adically separated, $g(f)=\xi(\pi(f))$. For general $w$, $f=\sum_nb_nX^n$ over a finite set of $n$ of weight $w$, with $b_n$ of weight $0$, so $g(f)=\sum_n\xi(\pi(b_n))y^n=t(w)\sum_n\xi(\pi(b_n))\xi(z^n)=t(w)\xi(\pi(f))$.

The homomorphism $g$ kills $\ker\pi$. Let $f\in\ker\pi$; its terms of each weight $w$ form a series $f_w\in\ker\pi$, so $g(f_w)=0$. For each $N$, $f$ minus the sum of the $f_w$ over the finitely many weights $w$ of its monomials of degree less than $N$ has only terms of degree at least $N$, whose images lie in $J^N$; so $g(f)\in\bigcap_NJ^N=0$. Hence $g=\Xi\circ\pi$ for a homomorphism of $\mathcal O$-algebras $\Xi:R\to B\subseteq K$ with $\Xi=t(w)\xi$ on $R_w$, whose image lies in the finite $\mathcal O$-module $B$. In case (1), $t=1$ and $\Xi$ extends $\xi$.
\end{proof}

\begin{lemma}\label{sha-counts-vd-structure}
Let $D$ be a Siegel datum with a graded model, with the notation of Definition~\ref{sha-counts-vd-model}.
\begin{enumerate}
\item The homomorphism $q\circ\pi$ is a weighted presentation of $S$, with parts $S_w$ and variables mapping into $\mathfrak m_S$. The sums of the $R_w$, of the $S_w$ and of the $A_n$ are direct; so $A=\bigoplus_nA_n$ is a $\mathbf Z$-graded $\mathcal O$-algebra, and $S^{\mathrm{gr}}=\bigoplus_wS_w$ is generated as an $\mathcal O$-algebra by $S_{(0,0)}$ and finitely many elements of the $S_w$.
\item The kernel of $q_A$ is $I$, and $q(A_n)=S_{(0,n)}$ for every $n$.
\item $A_0$ is a complete local noetherian $\mathcal O$-algebra with residue field $\mathbf F$, $A$ is generated over $A_0$ by finitely many elements of degrees $-2$, $-1$, $1$ and $2$, and $x$ vanishes on $A_n$ for $n\ne0$.
\item Let $\sigma\in G$. There is $a\in A_0$ with $x(a)=0$ and $\xi(q_A(a))=\operatorname{tr}r_\xi(\sigma)-\operatorname{tr}r_0(\sigma)$ for every homomorphism of $\mathcal O$-algebras $\xi:S^{\mathrm{gr}}\to\bar{\mathbf Q}_\ell$.
\end{enumerate}
\end{lemma}
\leanref{Development/AutomorphicArithmetic/AutomorphicVarieties/SymplecticGSp4/SInternal.lean}{UpperTheory.V.GradeZeroInterpolatingElement.sInternal}
\leanref{Development/AutomorphicArithmetic/AutomorphicVarieties/SymplecticGSp4/AGraded.lean}{UpperTheory.V.GradeZeroInterpolatingElement.aGraded}
\leanref{Development/AutomorphicArithmetic/AutomorphicVarieties/SymplecticGSp4/WcompMulHomog.lean}{UpperTheory.V.GradeZeroInterpolatingElement.kerA}
\leanref{Development/AutomorphicArithmetic/AutomorphicVarieties/SymplecticGSp4/WeightAlgebraGeneratedDegreeTwo.lean}{UpperTheory.V.GradeZeroInterpolatingElement.aW}
\leanref{Development/AutomorphicArithmetic/AutomorphicVarieties/SymplecticGSp4/AVal.lean}{UpperTheory.V.GradeZeroInterpolatingElement.aElt}
\begin{proof}
(a) Each $z_s$ has $x(z_s)=0$, as the constant term of $X_s$ is $0$; so $z_s\in\mathfrak m_R$, because the composite of $x$ with the residue map of $\mathcal O$ is a surjective homomorphism $R\to\mathbf F$, whose kernel is $\mathfrak m_R$. The kernel of $q$ is generated by the $F_i$, $E_i$ and $Z_i$, which lie in parts of $R$, so Lemma~\ref{sha-counts-vd-presentations}(a) gives the first statement, and $q(z_s)\in\mathfrak m_S$. Lemma~\ref{sha-counts-vd-presentations}(b), for $R$ and for $S$, gives the directness of the sums of the $R_w$, hence of the $A_n=R_{(0,n)}$, and of the $S_w$, and the generation of $S^{\mathrm{gr}}$ by $S_{(0,0)}$ and the $q(z_s)$.

(b) The ideal $I$ lies in the kernel, and $q(A_n)=q(R_{(0,n)})=S_{(0,n)}$ by definition. Let $f\in A$ with $q(f)=0$. Then $f=\sum_ia_ig_i$ with $a_i\in R$ and $g_i$ among the $F_j$, $E_j$ and $Z_j$, of weight $c_i\in\{(0,-1),(0,1),(0,0)\}$. Write $f=\pi(h)$ with $h$ a finite sum of series $h^{(n)}$ of weight $(0,n)$, $a_i=\pi(G_i)$, and $g_i=\pi(e_i)$ with $e_i$ of weight $c_i$. Since $h-\sum_iG_ie_i\in\ker\pi$, the terms of weight $(0,n)$ give $\pi(h^{(n)})=\sum_i\pi\big((G_i)_{(0,n)-c_i}\big)g_i$, in the notation of the proof of Lemma~\ref{sha-counts-vd-presentations}, and $\pi\big((G_i)_{(0,n)-c_i}\big)\in A$. Summing over $n$ gives $f\in I$.

(c) Since $A_0=R_{(0,0)}$, the first statement is Lemma~\ref{sha-counts-vd-presentations}(c). The weights of the variables have first coordinates $-1$, $0$ or $1$, and second coordinates $-1$, $0$ or $1$. An exponent vector whose weight has first coordinate $0$ is a sum of unit vectors of variables whose weight has first coordinate $0$ and of pairs of unit vectors of variables whose weights have first coordinates $1$ and $-1$: remove such a variable, or else such a pair, and otherwise all the first coordinates involved have the same sign, so the vector is $0$. By Lemma~\ref{sha-counts-vd-presentations}(b), $A_n=R_{(0,n)}$ is spanned over $R_{(0,0)}=A_0$ by monomials $z^m$ of weight $(0,n)$. So $A$ is generated over $A_0$ by the $z_s$ with weight of first coordinate $0$, of degrees $-1$, $0$ and $1$, and the products $z_sz_{s'}$ of variables with weights of first coordinates $1$ and $-1$, of degrees between $-2$ and $2$; those of degree $0$ lie in $A_0$. Finally $x$ is the constant term, and series of nonzero weight have constant term $0$.

(d) By Definition~\ref{sha-counts-vd-model}(iii) there are $f_1,\dots,f_4\in R_{(0,0)}=A_0$ with $q(f_i)=r^{\mathrm u}(\sigma)_{ii}$ and $x(f_i)=r_0(\sigma)_{ii}$. Put $a=\sum_if_i-\operatorname{tr}r_0(\sigma)$. Then $x(a)=0$, and $\xi(q_A(a))=\sum_i\xi(r^{\mathrm u}(\sigma)_{ii})-\operatorname{tr}r_0(\sigma)=\operatorname{tr}r_\xi(\sigma)-\operatorname{tr}r_0(\sigma)$.
\end{proof}

\begin{lemma}\label{sha-counts-vd-span}
Let $D$ be a finite Siegel datum (Definition~\ref{sha-siegel-datum}(3)) with a graded model, with the notation of Definition~\ref{sha-counts-vd-model}.
\begin{enumerate}
\item For every $w\in\mathbf Z^2$, $S_w$ is the $\mathcal O$-span of the products $r^{\mathrm u}(g_1)_{i_1j_1}\cdots r^{\mathrm u}(g_k)_{i_kj_k}$, $k\ge0$, whose weights add up to $w$.
\item The image of $A_0$ in $A/I$ is a finitely generated $\mathcal O$-module.
\end{enumerate}
\end{lemma}
\leanref{Development/AutomorphicArithmetic/AutomorphicVarieties/SymplecticGSp4/SpanMono.lean}{UpperTheory.V.GradeZeroInterpolatingElement.spanMono}
\leanref{Development/AutomorphicArithmetic/AutomorphicVarieties/SymplecticGSp4/FgMapMk.lean}{UpperTheory.V.GradeZeroInterpolatingElement.m2}
\begin{proof}
(a) Let $N_w$ be the span in question. By Definition~\ref{sha-counts-vd-model}(iii) the entry $(i,j)$ of $r^{\mathrm u}(g)$ lies in $S_w$, $w$ its weight, and $S_wS_{w'}\subseteq S_{w+w'}$; so $N_w\subseteq S_w$. Also $N_wN_{w'}\subseteq N_{w+w'}$, and $N_0$ is the $\mathcal O$-algebra generated by the products of entries of weight $0$, which is finite over $\mathcal O$ because $D$ is finite.

\emph{A presentation by entries.} For each $(g,i,j)$ let $z'_{g,i,j}$ be $r^{\mathrm u}(g)_{ij}$ if $i\ne j$, and $r^{\mathrm u}(g)_{ii}-c$ with $c\in\mathcal O$ of the same residue if $i=j$. These lie in $\mathfrak m_S$, since $r^{\mathrm u}\equiv r_0$ and $r_0$ is diagonal, and $z'_{g,i,j}\in S_w\cap N_w$, $w$ the weight of $(i,j)$. Two homomorphisms of $\mathcal O$-algebras $S\to\mathbf F[\varepsilon]$ that agree on the $z'$ agree on the entries of the $r^{\mathrm u}(g)$, so they are equal by Definition~\ref{sha-counts-vd-model}(ii), $\mathbf F[\varepsilon]$ being a complete local noetherian $\mathcal O$-algebra with residue field $\mathbf F$. So by Lemma~\ref{sha-counts-vd-presentations}(d), applied with $B=\mathcal O$, the $z'$ generate an ideal that, with $\mathfrak m_{\mathcal O}S$ and $\mathfrak m_S^2$, contains $\mathfrak m_S$. Since $S$ is noetherian, finitely many of them generate the same ideal, and by Lemma~\ref{sha-counts-vd-presentations}(d) they are the variables of a weighted presentation of $S$, which has the parts $S_w$ by Lemma~\ref{sha-counts-vd-presentations}(e) and Lemma~\ref{sha-counts-vd-structure}(a).

\emph{The part of weight $0$.} By Lemma~\ref{sha-counts-vd-presentations}(c), applied to this presentation, $S_0$ is a complete local noetherian $\mathcal O$-algebra whose maximal ideal is generated by $\mathfrak m_{\mathcal O}$ and finitely many monomials $u_1,\dots,u_k$ of weight $0$ in the $z'$, which lie in $N_0$. Each $u_i$ is integral over $\mathcal O$, as $N_0$ is finite over $\mathcal O$, and lies in $\mathfrak m_{S_0}$; so some power of $u_i$ lies in $\mathfrak m_{\mathcal O}S_0$. Indeed, let $u\in\mathfrak m_{S_0}$ and $u^jp(u)\in\mathfrak m_{\mathcal O}S_0$ with $p$ monic; if the constant term $p(0)$ lies in $\mathfrak m_{\mathcal O}$, then $u^{j+1}p_1(u)\in\mathfrak m_{\mathcal O}S_0$ with $p=Tp_1+p(0)$; otherwise $p(u)$ is a unit and $u^j\in\mathfrak m_{\mathcal O}S_0$; induction on the degree of $p$ concludes. So a power of the ideal $U$ generated by the $u_i$ lies in $\mathfrak m_{\mathcal O}S_0$. Since $S_0=\mathcal O+\mathfrak m_{\mathcal O}S_0+U$ and $u_iN_0\subseteq N_0$, induction gives $S_0=N_0+\mathfrak m_{\mathcal O}S_0+U^j$ for all $j$, hence $S_0=N_0+\mathfrak m_{\mathcal O}S_0$. As $N_0$ is finite over $\mathcal O$, hence $\mathfrak m_{\mathcal O}$-adically complete, and $S_0$ is $\mathfrak m_{\mathcal O}$-adically separated, this gives $S_0=N_0$.

\emph{All weights.} By Lemma~\ref{sha-counts-vd-presentations}(b) for the presentation by the $z'$, $S_w$ is the sum of the $S_0z'^n$ over finitely many $n$ of weight $w$; since $z'^n\in N_w$ and $S_0=N_0$, $S_w\subseteq N_0N_w\subseteq N_w$.

(b) By Lemma~\ref{sha-counts-vd-structure}(b), $q_A$ induces an injective map $A/I\to S$, which carries the image of $A_0$ into $S_{(0,0)}=N_0$, a finitely generated $\mathcal O$-module. Since $\mathcal O$ is noetherian, the image of $A_0$ is finitely generated.
\end{proof}

\begin{lemma}\label{sha-counts-vd-points}
Let $D$ be a Siegel datum with a graded model, with the notation of Definition~\ref{sha-counts-vd-model}, and let $K$ be a commutative $\mathcal O$-algebra.
\begin{enumerate}
\item Every homomorphism of $\mathcal O$-algebras $A\to\bar{\mathbf Q}_\ell$ that vanishes on $I$ is $\xi\circ q_A$ for a homomorphism of $\mathcal O$-algebras $\xi:S^{\mathrm{gr}}\to\bar{\mathbf Q}_\ell$.
\item Let $\xi:S^{\mathrm{gr}}\to K$ be a homomorphism of $\mathcal O$-algebras and $t\in(K^\times)^2$. The map $t\cdot\xi$ equal to $t^w\xi$ on each $S_w$ is a homomorphism of $\mathcal O$-algebras, and $r_{t\cdot\xi}=\tau(t)r_\xi\tau(t)^{-1}$.
\item Suppose that $D$ is finite. If $\xi,\xi':S^{\mathrm{gr}}\to K$ are homomorphisms of $\mathcal O$-algebras with $r_\xi=r_{\xi'}$, then $\xi=\xi'$.
\item Suppose that $D$ is finite, and let $\xi:S^{\mathrm{gr}}\to\bar{\mathbf Q}_\ell$ be a homomorphism of $\mathcal O$-algebras. There is $t\in(\bar{\mathbf Q}_\ell^\times)^2$ such that $\tau(t)r_\xi\tau(t)^{-1}$ is a $\bar{\mathbf Q}_\ell$-point of $D$.
\item Suppose that $D$ is finite, and let $\xi:S^{\mathrm{gr}}\to\bar{\mathbf Q}_\ell[\varepsilon]$, $\varepsilon^2=0$, be a homomorphism of $\mathcal O$-algebras such that the constant terms $c_{g,i,j}$ of the $r_\xi(g)_{ij}$ satisfy $|c_{g,i,j}-r_0(g)_{ij}|<1$, for the absolute value of $\bar{\mathbf Q}_\ell$. Then $r_\xi$ is a $\bar{\mathbf Q}_\ell[\varepsilon]$-point of $D$.
\end{enumerate}
\end{lemma}
\leanref{Development/AutomorphicArithmetic/AutomorphicVarieties/SymplecticGSp4/PtLift.lean}{UpperTheory.V.GradeZeroInterpolatingElement.ptLift}
\leanref{Development/AutomorphicArithmetic/AutomorphicVarieties/SymplecticGSp4/PtActInternal.lean}{UpperTheory.V.GradeZeroInterpolatingElement.ptAct}
\leanref{Development/AutomorphicArithmetic/AutomorphicVarieties/SymplecticGSp4/PtExt.lean}{UpperTheory.V.GradeZeroInterpolatingElement.ptExt}
\leanref{Development/AutomorphicArithmetic/AutomorphicVarieties/SymplecticGSp4/PtScale.lean}{UpperTheory.V.GradeZeroInterpolatingElement.ptScale}
\leanref{Development/AutomorphicArithmetic/AutomorphicVarieties/SymplecticGSp4/TanScale.lean}{UpperTheory.V.GradeZeroInterpolatingElement.tanScale}
\begin{proof}
(a) Let $x_1$ be such a homomorphism. By Lemma~\ref{sha-counts-vd-structure}(b), $x_1=\bar x_1\circ q_A$ for a homomorphism $\bar x_1$ on $q(A)=\bigoplus_nS_{(0,n)}$; make $\bar{\mathbf Q}_\ell$ a $q(A)$-algebra through $\bar x_1$. Let $P:S^{\mathrm{gr}}\to q(A)$ be the projection that keeps the components $S_w$ with $w_1=0$ (Lemma~\ref{sha-counts-vd-structure}(a)); it is $q(A)$-linear, since $S_{(0,n)}S_w\subseteq S_{w+(0,n)}$, and $P(1)=1$. By Lemma~\ref{sha-counts-vd-structure}(a), $S^{\mathrm{gr}}$ is a finitely generated $q(A)$-algebra. The $q(A)$-linear map $\bar x_1\circ P$ sends $1$ to $1$, so it induces a $\bar{\mathbf Q}_\ell$-linear map $\bar{\mathbf Q}_\ell\otimes_{q(A)}S^{\mathrm{gr}}\to\bar{\mathbf Q}_\ell$ sending $1$ to $1$; hence this tensor product is a nonzero finitely generated $\bar{\mathbf Q}_\ell$-algebra. The residue field of one of its maximal ideals is finite over $\bar{\mathbf Q}_\ell$ by Zariski's lemma \cite[Corollary~5.24]{AM69}, hence equal to it, and the composite $S^{\mathrm{gr}}\to\bar{\mathbf Q}_\ell\otimes_{q(A)}S^{\mathrm{gr}}\to\bar{\mathbf Q}_\ell$ is a homomorphism of $q(A)$-algebras $\xi$; so $\xi\circ q_A=x_1$.

(b) The $S_w$ form a grading of $S^{\mathrm{gr}}$ by Lemma~\ref{sha-counts-vd-structure}(a), and $t^wt^{w'}=t^{w+w'}$, so $t\cdot\xi$ is a homomorphism. The entry $(i,j)$ of $r^{\mathrm u}(g)$ lies in $S_w$, $w$ its weight, so $r_{t\cdot\xi}(g)_{ij}=t^wr_\xi(g)_{ij}$, which is the entry $(i,j)$ of $\tau(t)r_\xi(g)\tau(t)^{-1}$.

(c) The homomorphisms $\xi$ and $\xi'$ agree on the entries of the $r^{\mathrm u}(g)$, hence on their products, which span every $S_w$ by Lemma~\ref{sha-counts-vd-span}(a).

(d) By Lemma~\ref{sha-counts-vd-structure}(a), $S$ has the weighted presentation $q\circ\pi$, with parts $S_w$ and variables in $\mathfrak m_S$. By Lemma~\ref{sha-counts-vd-span}(a), $S_0$ lies in the $\mathcal O$-algebra generated by the products of entries of weight $0$, which is finite over $\mathcal O$ because $D$ is finite; so $\xi(S_0)$ is a finitely generated $\mathcal O$-module. Lemma~\ref{sha-counts-vd-extension}(2), with $X=\mathbf Z^2$, gives a homomorphism $\mathbf Z^2\to\bar{\mathbf Q}_\ell^\times$, that is, an element $t\in(\bar{\mathbf Q}_\ell^\times)^2$ with values $t^w$, and $\Xi:S\to\bar{\mathbf Q}_\ell$ with $\Xi=t^w\xi$ on $S_w$. The homomorphism $\Xi\circ r^{\mathrm u}$ is a $\bar{\mathbf Q}_\ell$-point of $D$, the image of the Siegel point $r^{\mathrm u}$, and its entries are $t^wr_\xi(g)_{ij}$; so it is $\tau(t)r_\xi\tau(t)^{-1}$ by (b).

(e) For $y\in\bar{\mathbf Q}_\ell[\varepsilon]$ let $y^\circ$ be its constant term. Let $N$ be the set of $f\in S^{\mathrm{gr}}$ for which there is $c\in\mathcal O$ with $f-c\in\mathfrak m_S$ and $|\xi(f)^\circ-c|<1$. It is an $\mathcal O$-subalgebra, since $|c|\le1$ for $c\in\mathcal O$, and it contains the entries of the $r^{\mathrm u}(g)$, with $c=r_0(g)_{ij}$. By Lemma~\ref{sha-counts-vd-span}(a), $N=S^{\mathrm{gr}}$. Hence $|\xi(f)^\circ|\le1$ for $f\in S^{\mathrm{gr}}$; and for a variable $q(z_s)\in\mathfrak m_S$ the corresponding $c$ lies in $\mathfrak m_{\mathcal O}$, so $|\xi(q(z_s))^\circ|<1$. As in (d), $\xi(S_0)$ is a finitely generated $\mathcal O$-module. Lemma~\ref{sha-counts-vd-extension}(1) gives $\Xi:S\to\bar{\mathbf Q}_\ell[\varepsilon]$ extending $\xi$, and $r_\xi=\Xi\circ r^{\mathrm u}$ is a $\bar{\mathbf Q}_\ell[\varepsilon]$-point of $D$.
\end{proof}

\begin{lemma}\label{sha-counts-vd-conjugate}
Let $D$ be a Siegel datum, with the notation of Definition~\ref{sha-siegel-datum}, and let $|\cdot|$ be the absolute value of $\bar{\mathbf Q}_\ell$.
\begin{enumerate}
\item Let $r$ be a point of $D$ with values in $\bar{\mathbf Q}_\ell$ or in $\bar{\mathbf Q}_\ell[\varepsilon]$, $\varepsilon^2=0$. Each $r(g)$ is symplectic with multiplier $\theta(g)$; $r(\phi)$ is block diagonal; the entries $(1,3)$ and $(3,1)$ of $r(g_0)+\theta(g_0)r(g_0^{-1})$ vanish; and the span of the $r(g)$ contains $E_{11},\dots,E_{44}$. If $r$ has values in $\bar{\mathbf Q}_\ell$, then $|r(g)_{ij}-r_0(g)_{ij}|<1$ for all $g$, $i$ and $j$.
\item Let $r$ and $r'$ be $\bar{\mathbf Q}_\ell$-points of $D$ with $r$ irreducible (Definition~\ref{sha-counts-vd-model}), and let $u\in GL_4(\bar{\mathbf Q}_\ell)$ satisfy $r'(g)=ur(g)u^{-1}$ for all $g$. Then $r'=\tau(t)r\tau(t)^{-1}$ for some $t\in(\bar{\mathbf Q}_\ell^\times)^2$.
\item Let $r$ be a $\bar{\mathbf Q}_\ell$-point of $D$ at which $D$ is rigid, and $r'$ a $\bar{\mathbf Q}_\ell[\varepsilon]$-point of $D$ with $r'\equiv r$ modulo $\varepsilon$. There are $a,b\in\bar{\mathbf Q}_\ell$ with $r'(g)_{ij}=\big(1+\varepsilon(aw_1+bw_2)\big)r(g)_{ij}$ for all $g$, $i$ and $j$, where $(w_1,w_2)$ is the weight of $(i,j)$.
\end{enumerate}
\end{lemma}
\leanref{Development/AutomorphicArithmetic/AutomorphicVarieties/SymplecticGSp4/PtInSmul.lean}{UpperTheory.V.GradeZeroInterpolatingElement.ptIn}
\leanref{Development/AutomorphicArithmetic/AutomorphicVarieties/SymplecticGSp4/EptInt.lean}{UpperTheory.V.GradeZeroInterpolatingElement.eptInt}
\leanref{Development/AutomorphicArithmetic/IntegralCasts/CapelliIdentities/NormLtOne.lean}{UpperTheory.Reformulated.B.NormLtOne.norm_lt_one}
\leanref{Development/AutomorphicArithmetic/AutomorphicVarieties/SymplecticGSp4/ConjTorus.lean}{UpperTheory.V.GradeZeroInterpolatingElement.conjTorus}
\leanref{Development/AutomorphicArithmetic/AutomorphicVarieties/SymplecticGSp4/RigidDiag.lean}{UpperTheory.V.GradeZeroInterpolatingElement.rigidDiag}
\begin{proof}
(a) Let $r=y\circ\rho$ with $\rho$ a Siegel point with values in $R'$ and $y$ a homomorphism of $\mathcal O$-algebras from $R'$. The first three properties hold for $\rho$ by Definition~\ref{sha-siegel-datum}(2), hence for $r$, and the $R'$-span of the $\rho(g)$ contains the $E_{ii}$ by Lemma~\ref{sha-flagged-idempotents}, so the span of the $r(g)$ contains them. Let $r$ have values in $\bar{\mathbf Q}_\ell$. Then $\rho(g)_{ij}-r_0(g)_{ij}\in\mathfrak m_{R'}$, and $|y(c)|<1$ for $c\in\mathfrak m_{R'}$. Otherwise $u=y(c)^{-1}$ has $|u|\le1$, so the minimal polynomial $T^d+\sum_{i<d}c_iT^i$ of $u$ over $\mathbf Q_\ell$ has coefficients in $\mathbf Z_\ell$, and multiplying the relation by $y(c)^d$ gives $y\big(1+\sum_ic_ic^{d-i}\big)=0$, although $1+\sum_ic_ic^{d-i}$ is a unit of $R'$.

(b) Let $d_i=\delta_i(\phi)$, so that $r_0(\phi)=\operatorname{diag}(d_1,\dots,d_4)$, and let $N=r(\phi)-r_0(\phi)$ and $N'=r'(\phi)-r_0(\phi)$. By (a), $N$ and $N'$ are block diagonal with entries of absolute value less than $1$; by Lemma~\ref{sha-flagged-residues}(1) and (2), $|d_i-d_j|=1$ when $\lambda_i\ne\lambda_j$.

\emph{$u$ is block diagonal.} From $ur(\phi)=r'(\phi)u$ we get $(d_j-d_i)u_{ij}=(N'u)_{ij}-(uN)_{ij}$. Suppose that $u_{ij}\ne0$ for some $(i,j)$ with $\lambda_i\ne\lambda_j$, and choose such $(i,j)$ with $|u_{ij}|$ largest. In $(uN)_{ij}=\sum_ku_{ik}N_{kj}$ only the $k$ with $\lambda_k=\lambda_j$ contribute, and for them $\lambda_i\ne\lambda_k$, so $|u_{ik}N_{kj}|<|u_{ij}|$; likewise for $(N'u)_{ij}$. So $|u_{ij}|=|(d_j-d_i)u_{ij}|<|u_{ij}|$, a contradiction.

\emph{The corner entries.} Let $H=r(g_0)+\theta(g_0)r(g_0^{-1})$ and $H'=r'(g_0)+\theta(g_0)r'(g_0^{-1})$, so that $uH=H'u$, and $H_{13}=H_{31}=H'_{13}=H'_{31}=0$ by (a). Put $h_i=\delta_i(g_0)+\theta(g_0)\delta_i(g_0)^{-1}$. By (a), and since $|\theta(g_0)|=1$, $|H_{ii}-h_i|<1$ and $|H'_{ii}-h_i|<1$. Now $h_1-h_3=\theta(g_0)+1-\alpha(g_0)-\Psi(g_0)=(\alpha(g_0)-1)(\Psi(g_0)-1)$, a unit since $\bar\alpha(g_0)\ne1$ and $\bar\Psi(g_0)\ne1$. So $H_{33}\ne H'_{11}$ and $H_{11}\ne H'_{33}$. As $u$ is block diagonal, the entries $(1,3)$ and $(3,1)$ of $uH=H'u$ give $u_{13}(H_{33}-H'_{11})=0$ and $u_{31}(H_{11}-H'_{33})=0$, so $u_{13}=u_{31}=0$.

\emph{$u$ is a torus element.} Since $r(g)^*r(g)=\theta(g)$, we have $r(g^{-1})=\theta(g)^{-1}r(g)^*$, and likewise for $r'$, with $X\mapsto X^*$ as in Definition~\ref{sha-siegel-datum}, which reverses products and satisfies $X^{**}=X$. Applying $*$ to $ur(g^{-1})=r'(g^{-1})u$ gives $r(g)u^*=u^*r'(g)$, so $u^*ur(g)=u^*r'(g)u=r(g)u^*u$, and $u^*u$ commutes with every $r(g)$, hence with $M_4(\bar{\mathbf Q}_\ell)$, as $r$ is irreducible: $u^*u=c\cdot1$ for some $c$. Row $3$ of $u$ has $u_{33}$ as its only possibly nonzero entry, so $u_{33}\ne0$; the entry $(1,3)$ of $u^*u$ is $u^*_{13}u_{33}=-u_{24}u_{33}$, so $u_{24}=0$, and in the same way $u_{42}=0$. So $u$ is diagonal, and the entries $(1,1)$ and $(3,3)$ of $u^*u$ give $u_{11}u_{44}=c=u_{22}u_{33}$. Hence $u=u_{44}\tau(t)$ with $t=(u_{33}/u_{44},u_{22}/u_{44})$, and $r'=\tau(t)r\tau(t)^{-1}$.

(c) Since $D$ is rigid at $r$, $r'=r+\varepsilon(Xr-rX)$ for some $X\in\mathfrak{sp}_4(\bar{\mathbf Q}_\ell)$. By (a), $r'(\phi)$ and $r(\phi)$ are block diagonal, so $Xr(\phi)-r(\phi)X$ is block diagonal. With $N$ and the $d_i$ as in (b), $(d_j-d_i)X_{ij}=(NX)_{ij}-(XN)_{ij}$ for $\lambda_i\ne\lambda_j$, and the argument of (b), with $X_{ij}$ of largest absolute value among the entries with $\lambda_i\ne\lambda_j$, shows that $X$ is block diagonal. With $H$ as in (b), the coefficient of $\varepsilon$ in the entries $(1,3)$ and $(3,1)$ of $r'(g_0)+\theta(g_0)r'(g_0^{-1})$ is the corresponding entry of $XH-HX$, and it vanishes by (a); as $H_{13}=H_{31}=0$ and $X$ is block diagonal, these entries are $X_{13}(H_{33}-H_{11})$ and $X_{31}(H_{11}-H_{33})$, and $H_{11}\ne H_{33}$ by (b). So $X_{13}=X_{31}=0$, and $X^*=-X$ gives $X_{24}=X_{42}=0$, $X_{33}=-X_{22}$ and $X_{44}=-X_{11}$. Thus $X=\operatorname{diag}(x_1,x_2,-x_2,-x_1)$, and $(XY-YX)_{ij}=(X_{ii}-X_{jj})Y_{ij}=(aw_1+bw_2)Y_{ij}$ with $a=x_1-x_2$ and $b=x_1+x_2$, for every matrix $Y$.
\end{proof}

\begin{lemma}\label{sha-counts-vd-flag}
Let $K$ be a field, $H$ a group, $I\subseteq H$ a subgroup and $s_0\in H$. Let $r:H\to M_4(K)$ be a homomorphism to the multiplicative monoid, and let $c_1,\dots,c_4:H\to K$ be functions such that the $c_i(s_0)$ are distinct and the characteristic polynomial of $r(s_0)$ is $\prod_j(X-c_j(s_0))$. Put $e_i=\prod_{j\ne i}(r(s_0)-c_j(s_0))$, and suppose that $e_jr(g)e_i=0$ for $i<j$ and $g\in H$, and that $e_i(r(s)-c_i(s))e_i=0$ for $s\in I$ and every $i$.
Then there are $u\in GL_4(K)$ and homomorphisms $\kappa_1,\dots,\kappa_4:H\to K^\times$ such that $ur(g)u^{-1}$ is upper triangular with diagonal entries $\kappa_1(g),\dots,\kappa_4(g)$ for every $g\in H$, $\kappa_i(s)=c_i(s)$ for $s\in I$, and $d_i\kappa_i(g)=\operatorname{tr}(e_ir(g))$ for all $i$ and $g$, where $d_i=\prod_{j\ne i}(c_i(s_0)-c_j(s_0))$.
\end{lemma}
\leanref{Development/AutomorphicArithmetic/AutomorphicVarieties/SymplecticGSp4/OrdFlag.lean}{UpperTheory.V.GradeZeroInterpolatingElement.ordFlag}
\begin{proof}
The $c_i(s_0)$ are distinct roots of the characteristic polynomial, so $r(s_0)P=P\operatorname{diag}(c_1(s_0),\dots,c_4(s_0))$ for some $P\in GL_4(K)$. Then $p(r(s_0))P=P\operatorname{diag}(p(c_1(s_0)),\dots,p(c_4(s_0)))$ for every polynomial $p$, so $e_i=d_iPE_{ii}P^{-1}$ with $d_i=\prod_{j\ne i}(c_i(s_0)-c_j(s_0))\ne0$. Let $M(g)=P^{-1}r(g)P$. For $i<j$, $0=P^{-1}e_jr(g)e_iP=d_jd_iE_{jj}M(g)E_{ii}$, so $M(g)_{ji}=0$, and $M(g)$ is upper triangular. Hence $\kappa_i(g)=M(g)_{ii}$ is multiplicative in $g$, and $\kappa_i(g)\kappa_i(g^{-1})=\kappa_i(1)=1$. For $s\in I$, $0=P^{-1}e_i(r(s)-c_i(s))e_iP=d_i^2(M(s)_{ii}-c_i(s))E_{ii}$, so $\kappa_i(s)=c_i(s)$. Finally $\operatorname{tr}(e_ir(g))=\operatorname{tr}(d_iE_{ii}M(g))=d_i\kappa_i(g)$. Take $u=P^{-1}$.
\end{proof}

\begin{lemma}\label{sha-counts-vd-lattice}
Let $\mathcal O'$ be a discrete valuation ring with residue field $k'$, and write $\bar a$ for the residue of $a\in\mathcal O'$. Let $H$ be a group, let $R:H\to M_4(\mathcal O')$ be a map with $R(1)=1$, and let $s\ne u$ be indices with $R(g)_{su}\ne0$ for some $g$. Suppose that there are $g_1,\dots,g_n\in H$ and $a_1,\dots,a_n\in\mathcal O'$ with $\sum_ka_kR(g_kh)=E_{ss}R(h)$ for every $h\in H$.
Then there are a nonzero $b\in\mathcal O'$ and a function $c:H\to\mathcal O'$ with $R(g)_{su}=bc(g)$ for all $g$, such that no $m\in k'$ satisfies $\bar c(g)=m(\bar R(g)_{ss}-\bar R(g)_{uu})$ for all $g\in H$.
\end{lemma}
\leanref{Development/AutomorphicArithmetic/AutomorphicVarieties/SymplecticGSp4/ExistsThirdIndexLinkedPair.lean}{lattice}
\begin{proof}
Choose $g_*$ with $R(g_*)_{su}$ of least valuation, and put $b=R(g_*)_{su}\ne0$ and $c(g)=R(g)_{su}/b\in\mathcal O'$. Then $c(g_*)=1$, and $c(1)=0$ since $R(1)=1$. Suppose that $m$ is as excluded. The entries $(s,u)$, $(s,s)$ and $(u,u)$ of the identity give $\sum_ka_kc(g_kh)=c(h)$, $\sum_ka_kR(g_kh)_{ss}=R(h)_{ss}$ and $\sum_ka_kR(g_kh)_{uu}=0$ for every $h$. Reducing, $\bar c(h)=\sum_k\bar a_km(\bar R(g_kh)_{ss}-\bar R(g_kh)_{uu})=m\bar R(h)_{ss}$. At $h=1$ this gives $m=0$, and then at $h=g_*$ it gives $1=0$.
\end{proof}

\begin{lemma}\label{sha-counts-vd-linked}
Let $D$ be a Siegel datum, with the notation of Definition~\ref{sha-siegel-datum}, let $r$ be a $\bar{\mathbf Q}_\ell$-point of $D$, and write $i\succeq j$ if $r(g)_{ij}\ne0$ for some $g\in G$.
If $\lambda_i\ne\lambda_j$, $i\succeq j$ and $j\succeq i$, then $i\succeq k\succeq i$ for some $k\notin\{i,j\}$.
\end{lemma}
\leanref{Development/AutomorphicArithmetic/AutomorphicVarieties/SymplecticGSp4/ExistsThirdIndexLinkedPair.lean}{UpperTheory.V.GradeZeroInterpolatingElement.exists_third_index_linked_pair}
\leanref{Development/AutomorphicArithmetic/AutomorphicVarieties/SymplecticGSp4/ExistsThirdIndexLinkedPair.lean}{core}
\begin{proof}
Suppose not. Let $|\cdot|$ be the absolute value of $\bar{\mathbf Q}_\ell$, and $\chi=\chi_{\mathrm{cyc}}$.
\begin{enumerate}
\item[(i)] \emph{Integrality.} Write $r=f\circ\rho$ with $\rho$ a Siegel point with values in $R'$ and $f:R'\to\bar{\mathbf Q}_\ell$. The image of $f$ is finite over $\mathcal O$, as in the proof of Lemma~\ref{sha-model-tangent}(e); its integral closure $\mathcal O'$ is a discrete valuation ring finite over $\mathbf Z_\ell$ that contains $\mathcal O$ and the entries of the $r(g)$. Let $k'$ be its residue field and $\bar x$ the residue of $x\in\mathcal O'$. By Lemma~\ref{sha-flagged-idempotents}, each $E_{aa}$ is an $R'$-linear combination of finitely many $\rho(g)$, so it is an $\mathcal O'$-linear combination of the $r(g)$. By Lemma~\ref{sha-counts-vd-conjugate}(a), $|r(g)_{ab}-r_0(g)_{ab}|<1$, so $\bar r(g)_{aa}=\bar\delta_a(g)$, and $|r(g)_{ab}|<1$ for $a\ne b$. As there, the finitely many generators of $\mathfrak m_{R'}$ have images of absolute value less than $1$, so $|f(t)|\le q_0^N$ for $t\in\mathfrak m_{R'}^N$, for some $q_0<1$; since $\rho$ is continuous, each function $g\mapsto r(g)_{ab}$ modulo a nonzero ideal of $\mathcal O'$ is locally constant.
\item[(ii)] \emph{Products.} If $i\succeq a$ and $b\succeq i$, then $r(gh)_{ab}=r(g)_{ai}r(h)_{ib}+r(g)_{aj}r(h)_{jb}$ for all $g,h\in G$. Indeed, by Lemma~\ref{sha-counts-vd-conjugate}(a) and Lemma~\ref{sha-flagged-relation}(2), $\succeq$ is transitive; if $k\notin\{i,j\}$ and $r(g)_{ak}r(h)_{kb}\ne0$, then $i\succeq a\succeq k\succeq b\succeq i$, which is excluded.
\item[(iii)] \emph{A common line at $\ell$.} Since $s_0\in I_\ell$, $\psi(s_0)=\omega(s_0)=1$ and $\xi$ is trivial on $I_\ell$, the $\delta_a(s_0)$ are $\chi(s_0)^{k+1}$, $\chi(s_0)^k$, $\chi(s_0)$ and $1$, which are distinct, as $\chi(s_0)$ has infinite order and $k\ge2$. The restriction of $r$ to $D_\ell$ is the image of an ordinary lift (Definition~\ref{sha-siegel-datum}(2)), so Lemma~\ref{sha-counts-vd-flag}, with $H=D_\ell$, $I=I_\ell$ and $c_a=\delta_a$, gives $u$ with $ur(g)u^{-1}$ upper triangular for $g\in D_\ell$. Apply Lemma~\ref{sha-flagged-common-line} to $r$ on $D_\ell$, the matrix $u^{-1}$, $\Pi=\{a:a\succeq i\}$ and the indices $i$ and $j$. Its hypotheses hold: if $r(g)_{ab}\ne0$ with $b\in\Pi$, then $a\succeq b\succeq i$, so $a\in\Pi$; and for $k\in\Pi$ other than $i$ and $j$, $r(g)_{ik}=0$, since otherwise $i\succeq k\succeq i$, and $r(g)_{jk}=0$, since otherwise $i\succeq j\succeq k\succeq i$. So there are $(x,y)\ne(0,0)$ and $\varphi$ with $r(g)_{ii}x+r(g)_{ij}y=\varphi(g)x$ and $r(g)_{ji}x+r(g)_{jj}y=\varphi(g)y$ for $g\in D_\ell$.
\item[(iv)] \emph{One index pair.} If $|x|\le|y|$ put $(a,b)=(i,j)$ and $v=x/y$; otherwise put $(a,b)=(j,i)$ and $v=y/x$. Then $|v|\le1$, $\lambda_a\ne\lambda_b$, $a\succeq b$, and eliminating $\varphi$ gives
\[
 r(g)_{aa}v+r(g)_{ab}=\big(r(g)_{ba}v+r(g)_{bb}\big)v\qquad(g\in D_\ell).\tag{$*$}
\]
By (ii), $r(gh)_{ab}=r(g)_{aa}r(h)_{ab}+r(g)_{ab}r(h)_{bb}$ for all $g,h\in G$.
\item[(v)] \emph{A cocycle.} By (i), $E_{aa}r(h)=\sum_kt_kr(g_kh)$ for every $h$, with $t_k\in\mathcal O'$, and $r(g)_{ab}\ne0$ for some $g$. Lemma~\ref{sha-counts-vd-lattice} gives a nonzero $b_0\in\mathcal O'$ and $c:G\to\mathcal O'$ with $r(g)_{ab}=b_0c(g)$, such that no $m\in k'$ has $\bar c(g)=m(\bar\delta_a(g)-\bar\delta_b(g))$ for all $g$. Let $\gamma=\delta_a\delta_b^{-1}$, one of $\alpha^{\pm1}$, $\beta^{\pm1}$, $\theta^{\pm1}$ since $\lambda_a\ne\lambda_b$, and put $z(g)=\bar c(g)/\bar\delta_b(g)$. By (iv), $c(gh)=r(g)_{aa}c(h)+c(g)r(h)_{bb}$; reducing, $\bar c(gh)=\bar\delta_a(g)\bar c(h)+\bar c(g)\bar\delta_b(h)$, so $z(gh)=z(g)+\bar\gamma(g)z(h)$. By (i), $z$ is locally constant.
\item[(vi)] \emph{Local conditions.} If $g$ lies in an inertia group at a place not above $T$ and $\bar r_0(g)=1$, then $r(g)=1$, so $z(g)=0$. Let $g\in D_\ell$ and $\sigma=\sigma_\ell\in D_\ell$. From $(*)$ at $g$ and at $\sigma$,
\[
 c(g)\big(r(\sigma)_{bb}-r(\sigma)_{aa}\big)-c(\sigma)\big(r(g)_{bb}-r(g)_{aa}\big)=v\big(c(\sigma)r(g)_{ba}-c(g)r(\sigma)_{ba}\big),
\]
since $b_0$ times either side is $v^2\big(r(g)_{ba}(r(\sigma)_{bb}-r(\sigma)_{aa})-r(\sigma)_{ba}(r(g)_{bb}-r(g)_{aa})\big)$. The right side has absolute value less than $1$, so $\bar c(g)(\bar\delta_b(\sigma)-\bar\delta_a(\sigma))=\bar c(\sigma)(\bar\delta_b(g)-\bar\delta_a(g))$. By Lemma~\ref{sha-flagged-residues}(1), $\bar\delta_b(\sigma)\ne\bar\delta_a(\sigma)$, and with $m_\ell=-\bar c(\sigma)/(\bar\delta_b(\sigma)-\bar\delta_a(\sigma))$ this says $z(g)=\bar\gamma(g)m_\ell-m_\ell$ for $g\in D_\ell$. For $p\in S_b$, the restriction of $z$ to $D_p$ is a coboundary by Lemma~\ref{sha-flagged-cocycles}(2).
\item[(vii)] By Lemma~\ref{sha-flagged-cocycles}(1), for the reduction $\mathcal O\to k'$, there is $m\in k'$ with $z(g)=\bar\gamma(g)m-m$ for all $g$. Then $\bar c(g)=\bar\delta_b(g)z(g)=m(\bar\delta_a(g)-\bar\delta_b(g))$ for all $g$, which (v) excludes.
\end{enumerate}
\end{proof}

\begin{lemma}\label{sha-counts-vd-soule}
Let $\ell$ be a prime, write $G=\operatorname{Gal}(\bar{\mathbf Q}/\mathbf Q)$ and $\chi=\chi_{\mathrm{cyc}}$, and let $|\cdot|$ be the absolute value of $\bar{\mathbf Q}_\ell$.
\begin{enumerate}
\item For every ring $\mathcal O'$ as in Definition~\ref{sha-character-parts}, every character $\psi':G\to\mathcal O'^\times$ with open kernel and every $k'\ge1$, the length $V_\ell(k',\psi')$ is finite.
\item Let $\psi':G\to\bar{\mathbf Q}_\ell^\times$ be a character with open kernel, let $e\ge1$, and put $\theta=\psi'^{-1}\chi^{-e}$. Let $c:G\to\bar{\mathbf Q}_\ell$ be continuous with $c(gh)=c(g)+\theta(g)c(h)$ for all $g,h\in G$. Suppose that for every inertia group $J$ of $G$, and for every decomposition group $J$ at a place above $\ell$, there is $m_J\in\bar{\mathbf Q}_\ell$ with $c(g)=\theta(g)m_J-m_J$ for $g\in J$. Then there is $m\in\bar{\mathbf Q}_\ell$ with $c(g)=\theta(g)m-m$ for all $g\in G$.
\end{enumerate}
\end{lemma}
\leanref{Development/AutomorphicArithmetic/AutomorphicVarieties/LevelCohomology/ExistsAb.lean}{UpperTheory.V.X.soFin}
\leanref{Development/AutomorphicArithmetic/AutomorphicVarieties/CharacterLevelBounds/LiftGens.lean}{UpperTheory.V.Chr.upV}
\leanref{Development/AutomorphicArithmetic/AutomorphicVarieties/GaloisLibrary/SelmerTransport/KerResAb.lean}{UpperTheory.V.Base.kerResAb}
\leanref{Development/AutomorphicArithmetic/AutomorphicVarieties/CharacterLevelBounds/H0Nu.lean}{UpperTheory.V.Chr.h0Nu}
\leanref{Development/AutomorphicArithmetic/AutomorphicVarieties/LevelCohomology/FdAdjoinRange.lean}{UpperTheory.V.X.soBar}
\leanref{Development/AutomorphicArithmetic/AutomorphicVarieties/LevelCohomology/IntIffNorm.lean}{UpperTheory.V.X.soVec}
\begin{proof}
(a) Let $N=\ell$ if $\ell$ is odd and $N=4$ if $\ell=2$, and let $L$ be the field fixed by the kernel of $g\mapsto(\psi'(g),\chi(g)\bmod N)$. It is a finite abelian extension of $\mathbf Q$ containing a primitive $N$-th root of unity, so it is totally complex. By Theorem~\ref{sha-soule-bound}, $H^2(G_{L,S_\ell},\mathbf Z_\ell(k'+1))$ is finite, and by Poitou-Tate duality, as in Step 2 of the proof of that theorem, $\text{Ш}^1_\ell(L,k')$ is finite.
Let $A=\mathcal O'\otimes\mathbf Q_\ell/\mathbf Z_\ell(\psi'^{-1}\chi^{-k'})$ and let $V\subseteq H^1(G,A)$ be the group whose length is $V_\ell(k',\psi')$, of the classes that vanish on the decomposition groups at $\ell$ and on the inertia groups. Let $G_L=\operatorname{Gal}(\bar{\mathbf Q}/L)$, and let $K$ be the kernel of restriction from $H^1(G,A)$ to $H^1(G_L,A)$.
\begin{enumerate}
\item[(i)] Restriction maps $V$ into the group of classes of $G_L$ with the same local conditions. Since $\psi'$ is trivial on $G_L$, this group is $\text{Ш}^1_\ell(L,k')\otimes\mathcal O'$, of length $e_{\mathcal O'}\operatorname{len}\text{Ш}^1_\ell(L,k')$. So $V_\ell(k',\psi')\le e_{\mathcal O'}\operatorname{len}\text{Ш}^1_\ell(L,k')+\operatorname{len}K$.
\item[(ii)] Let $T\subset G$ be a finite set whose image generates $\operatorname{Gal}(L/\mathbf Q)$. A class in $K$ is represented by a cocycle $z$ that vanishes on $G_L$, and such a cocycle is determined by its values on $T$, since the $g$ with $z(g)=0$ form a closed subgroup. Each $u\in G$ acts on $A$ by a scalar $c_u$; as $G/G_L$ is abelian and commutators act trivially on $A$, $z(ug)=z(gu)$ for all $g$, which says that $(c_u-1)z$ is the coboundary of $z(u)$. So $aK=0$ for a generator $a$ of the ideal of $\mathcal O'$ spanned by the $c_u-1$. Hence the length of $K$ is at most that of $Z/aZ$, with $Z$ the module of these cocycles, which is at most that of its $a$-torsion; and evaluation at $T$ embeds the $a$-torsion of $Z$ into $A[a]^T$, where $A[a]=A^G$ is the group of $G$-invariants. So $\operatorname{len}K\le|T|\operatorname{len}A^G$.
\item[(iii)] $A^G$ lies in the invariants of $\ker\psi'$, which act by $\chi^{-k'}$. As $\chi(\ker\psi')$ contains $\chi(G_L)$, which is open in $\mathbf Z_\ell^\times$, it contains an $x$ with $x^{k'}\ne1$, and $A^G$ lies in the finite group of elements of $A$ killed by $x^{k'}-1$.
\end{enumerate}
So $V_\ell(k',\psi')\le e_{\mathcal O'}\operatorname{len}\text{Ш}^1_\ell(L,k')+|T|\operatorname{len}A^G$ is finite.

(b) The character $\psi'$ has finite image, so its values and those of $\theta$ lie in a finite extension $K_0$ of $\mathbf Q_\ell$.
\begin{enumerate}
\item[(i)] The values of $c$ lie in a finite extension $E\supseteq K_0$. There is a countable family $(K_n)$ of finite extensions of $\mathbf Q_\ell$ in $\bar{\mathbf Q}_\ell$ such that every finite subset of $\bar{\mathbf Q}_\ell$ lies in some $K_n$. The sets $H_n$ of the $g$ with $c(g)\in K_0K_n$ are subgroups, by the cocycle relation, and closed, since $c$ is continuous and $K_0K_n$ is closed; they cover the compact group $G$. By the Baire category theorem some $H_n$ has nonempty interior, so it is open and of finite index, and $c$ takes values in the extension $E$ of $K_0K_n$ generated by the values of $c$ on representatives of $G/H_n$.
\item[(ii)] Let $\mathcal O'$ be the integral closure of $\mathbf Z_\ell$ in $E$, with a uniformizer $\varpi$; then $\psi'$ and $\theta$ take values in $\mathcal O'^\times$. By (a), the group $V$ of (a) for $\mathcal O'$, $\psi'$ and $k'=e$ has finite length, so $\varpi^M V=0$ for some $M$. Identify $A$ with $E/\mathcal O'$. For $n\ge0$, the reduction of $\varpi^{-(n+M)}c$ modulo $\mathcal O'$ is a cocycle with values in $A$ whose class lies in $V$, by the hypotheses on $c$. So $\varpi^M$ times it is a coboundary: there is $y_n\in E$ with $\varpi^{-n}c(g)-(\theta(g)y_n-y_n)\in\mathcal O'$ for all $g$, that is, $|c(g)-(\theta(g)m_n-m_n)|\le|\varpi|^n$ with $m_n=\varpi^ny_n$.
\item[(iii)] If $\theta=1$, this gives $|c(g)|\le|\varpi|^n$ for all $n$, so $c=0$. Otherwise choose $g_0$ with $\theta(g_0)\ne1$ and put $m=c(g_0)/(\theta(g_0)-1)$. Then
\[
 c(g)-(\theta(g)m-m)=\big(c(g)-(\theta(g)m_n-m_n)\big)-\big(c(g_0)-(\theta(g_0)m_n-m_n)\big)\frac{\theta(g)-1}{\theta(g_0)-1}
\]
has absolute value at most $|\varpi|^n\big(1+|(\theta(g)-1)/(\theta(g_0)-1)|\big)$ for every $n$, so it is $0$.
\end{enumerate}
\end{proof}

\begin{lemma}\label{sha-counts-vd-tangent}
Let $D$ be a Siegel datum, with the notation of Definition~\ref{sha-siegel-datum}, and let $\mathcal O[\varepsilon]$ be the ring of dual numbers, $\varepsilon^2=0$. Let $Z:D_\ell\to\mathfrak{sp}_4(\mathcal O)$ be a function whose entries are locally constant modulo every power of $\mathfrak m_{\mathcal O}$, such that $\rho(g)=(1+\varepsilon Z(g))r_0(g)$ is a homomorphism $D_\ell\to M_4(\mathcal O[\varepsilon])$ with the following properties, where $e_i=\prod_{j\ne i}(\rho(s_0)-\delta_j(s_0))$: each $\rho(g)$ is symplectic with multiplier $\theta(g)$; the characteristic polynomial of $\rho(s_0)$ is $\prod_j(X-\delta_j(s_0))$; $e_j\rho(g)e_i=0$ for $i<j$ and $g\in D_\ell$; and $e_i(\rho(s)-\delta_i(s))e_i=0$ for $s\in I_\ell$ and every $i$. (The proof uses neither the continuity of $Z$ nor the conditions on the multipliers and on the characteristic polynomial.)
\begin{enumerate}
\item If every $Z(g)$ is supported on the entries of weight $(0,-1)$, there are $a,m\in\mathcal O$ with $a\ne0$ and $aZ(g)_{31}=\Psi(g)^{-1}m-m$ for all $g\in D_\ell$.
\item If every $Z(g)$ is diagonal, then $Z(s)=0$ for $s\in I_\ell$.
\end{enumerate}
\end{lemma}
\leanref{Development/AutomorphicArithmetic/AutomorphicVarieties/SymplecticGSp4/SpWApply.lean}{UpperTheory.V.GradeZeroInterpolatingElement.tanLz}
\leanref{Development/AutomorphicArithmetic/AutomorphicVarieties/SymplecticGSp4/SpWApply.lean}{triple}
\begin{proof}
Let $c_k=\delta_k(s_0)$, which are distinct, as in step (iii) of the proof of Lemma~\ref{sha-counts-vd-linked}, and $d_i=\prod_{j\ne i}(c_i-c_j)\ne0$. Since $\rho(s_0)\equiv\operatorname{diag}(c_1,\dots,c_4)$ modulo $\varepsilon$, $e_i=d_iE_{ii}+\varepsilon e'_i$ for some $e'_i\in M_4(\mathcal O)$.
Let $A=\alpha_0E_{pp}+\varepsilon A'$, $B=\operatorname{diag}(\beta_1,\dots,\beta_4)+\varepsilon B'$ and $C=\gamma_0E_{qq}+\varepsilon C'$ be matrices over $\mathcal O[\varepsilon]$. Multiplying out, the coefficient of $\varepsilon$ in the entry $(p,q)$ of $ABC$ is $A'_{pq}\beta_q\gamma_0+\alpha_0B'_{pq}\gamma_0+\alpha_0\beta_pC'_{pq}$.

(a) Apply this to $e_3\rho(g)e_1=0$. The coefficient of $\varepsilon$ in $\rho(g)_{31}$ is $(Z(g)r_0(g))_{31}=Z(g)_{31}\theta(g)$, as $\delta_1=\theta$; so
\[
 0=e'_{3,31}\,\theta(g)d_1+d_3d_1\theta(g)Z(g)_{31}+d_3\alpha(g)\,e'_{1,31},
\]
where $e'_{i,31}$ is the entry $(3,1)$ of $e'_i$. At $g=1$, where $\rho(1)=1$ gives $Z(1)=0$, this reads $e'_{3,31}d_1+d_3e'_{1,31}=0$. Dividing by $\theta(g)$ and using $\alpha\theta^{-1}=\Psi^{-1}$, $d_1d_3Z(g)_{31}=\Psi(g)^{-1}m-m$ with $m=-d_3e'_{1,31}$. Take $a=d_1d_3$.

(b) Let $s\in I_\ell$, and apply the computation to $e_i(\rho(s)-\delta_i(s))e_i=0$, with $\beta_k=\delta_k(s)-\delta_i(s)$, so that $\beta_i=0$, and $B'=Z(s)r_0(s)$. The coefficient of $\varepsilon$ in the entry $(i,i)$ is $d_i^2\delta_i(s)Z(s)_{ii}=0$, so $Z(s)_{ii}=0$; and $Z(s)$ is diagonal.
\end{proof}

\begin{lemma}\label{sha-counts-vd-diag-henselian}
Let $\ell$ be a prime, let $|\cdot|$ be the absolute value of $\bar{\mathbf Q}_\ell$, and let $\mathcal O$ be a $\mathbf Z_\ell$-algebra with a homomorphism of $\mathbf Z_\ell$-algebras $\mathcal O\to\bar{\mathbf Q}_\ell$. Let $B$ be an $\mathcal O$-algebra that is a henselian local domain, with fraction field $F$, and let $x,y:B\to\bar{\mathbf Q}_\ell$ be homomorphisms of $\mathcal O$-algebras.
Let $d\in\mathcal O$ have nonzero image in $\bar{\mathbf Q}_\ell$, and let $t\in B$ with $x(t)=0$ be such that $dv=t$ for some $v\in F$ that is integral over $B$. If $\kappa\in\bar{\mathbf Q}_\ell$ satisfies $d\kappa=y(t)$, then $|\kappa|<1$.
\end{lemma}
\leanref{Development/AutomorphicArithmetic/IntegralCasts/CapelliIdentities/MajorantSquareBound.lean}{UpperTheory.Reformulated.B.MajorantSquareBound.spl_near}
\begin{proof}
Let $S$ be the integral closure of $B$ in an algebraic closure of $F$; it is a domain, integral over $B$.
\begin{enumerate}
\item[(i)] $S$ is a local ring. Every maximal ideal $\mathfrak n$ of $S$ meets $B$ in $\mathfrak m_B$, since $S$ is integral over $B$ \cite[Corollary~5.8]{AM69}. Let $\mathfrak n_1$ and $\mathfrak n_2$ be maximal ideals of $S$, and let $\theta\in S$ with $\theta\notin\mathfrak n_2$; it suffices to show that $\theta\notin\mathfrak n_1$. Let $Q\in B[X]$ be monic of degree $d$ with $Q(\theta)=0$. In the algebraic closure of $F$, $Q=\prod_{i=1}^d(X-\theta_i)$, where the $\theta_i$, being roots of $Q$, lie in $S$, and $\theta=\theta_{i_0}$ for some $i_0$. Let $a$ be the multiplicity of $0$ as a root of the reduction of $Q$ modulo $\mathfrak m_B$. For every maximal ideal $\mathfrak n$ of $S$, the reduction of $Q$ modulo $\mathfrak n$ is $\prod_i(X-\bar\theta_i)$ over the field $S/\mathfrak n$, which contains $B/\mathfrak m_B$; so $a$ of the $\theta_i$ lie in $\mathfrak n$.
Put $r=d-a$, and $\theta_I=\prod_{i\in I}\theta_i$ for $I\subseteq\{1,\dots,d\}$. The coefficients of $\prod_{|I|=r}(X-\theta_I)$ are symmetric polynomials with integer coefficients in the $\theta_i$, hence polynomials with integer coefficients in the coefficients of $Q$ \cite[Chapter~IV, Theorem~6.1]{Lan02}. This gives a monic $H\in B[X]$ with $H=\prod_{|I|=r}(X-\theta_I)$ in $S[X]$.
Let $I_0$ be the set of the $i$ with $\theta_i\notin\mathfrak n_2$, which has $r$ elements and contains $i_0$, and let $c\ne0$ be the image of $\theta_{I_0}$ in $S/\mathfrak n_2$. Modulo $\mathfrak n_2$, $a$ of the factors of $Q=\prod_i(X-\theta_i)$ become $X$, so the coefficient of $X^a$ in $Q$ reduces to $(-1)^rc$; let $A\in B$ be $(-1)^r$ times this coefficient, which reduces to $c$. Modulo $\mathfrak n_2$, $H$ becomes $(X-c)G$ with $G=\prod_{I\ne I_0}(X-\bar\theta_I)$, and $G(c)\ne0$, since every $I\ne I_0$ with $r$ elements contains some $i\notin I_0$, for which $\bar\theta_i=0$. As $B/\mathfrak m_B$ embeds in $S/\mathfrak n_2$, $H(A)\in\mathfrak m_B$ and $H'(A)\notin\mathfrak m_B$. Since $B$ is henselian, $H$ has a root $\tau\in B$ with $\tau-A\in\mathfrak m_B$ \cite[Tag 04GE]{Sta}, and $\tau\notin\mathfrak m_B$, as $c\ne0$.
In $S$, $H(\tau)=0$ gives $\tau=\theta_I$ for some $I$ with $r$ elements. For every maximal ideal $\mathfrak n$ of $S$, $\tau\notin\mathfrak n$, as $\mathfrak n\cap B=\mathfrak m_B$; so $\theta_i\notin\mathfrak n$ for $i\in I$, and $I$ is the set of the $i$ with $\theta_i\notin\mathfrak n$, both sets having $r$ elements. For $\mathfrak n=\mathfrak n_2$ this gives $i_0\in I$, and then for $\mathfrak n=\mathfrak n_1$ it gives $\theta=\theta_{i_0}\notin\mathfrak n_1$.
\item[(ii)] Since $S$ is integral over $B$ and $\bar{\mathbf Q}_\ell$ is algebraically closed, $x$ and $y$ extend to homomorphisms $x',y':S\to\bar{\mathbf Q}_\ell$ \cite[Chapter~5, Exercise~2]{AM69}. Let $v_S\in S$ be the element $v$, so that $dv_S=t$. Then $d\,x'(v_S)=x(t)=0$, so $x'(v_S)=0$, and $v_S$ lies in the maximal ideal $\mathfrak m_S$ of $S$, since $x'$ maps units to units. Also $d\,y'(v_S)=y(t)=d\kappa$, so $y'(v_S)=\kappa$.
\item[(iii)] Suppose that $|\kappa|\ge1$. Then $\kappa\ne0$ and $|\kappa^{-1}|\le1$, so the minimal polynomial of $\kappa^{-1}$ over $\mathbf Q_\ell$ has coefficients in $\mathbf Z_\ell$; reversing it gives $h\in\mathbf Z_\ell[X]$ with $h(\kappa)=0$ and $h(0)=1$. The element $w=h(v_S)$ of $S$ has $y'(w)=h(\kappa)=0$, so it is not a unit and lies in $\mathfrak m_S$. But $w-1$ is a multiple of $v_S\in\mathfrak m_S$, so $1\in\mathfrak m_S$, a contradiction.
\end{enumerate}
\end{proof}

\begin{lemma}\label{sha-counts-vd-diag-local}
Let $D$ be a Siegel datum, with the notation of Definition~\ref{sha-siegel-datum}, let $|\cdot|$ be the absolute value of $\bar{\mathbf Q}_\ell$, and let $r$ be a $\bar{\mathbf Q}_\ell$-point of $D$.
Then there are $u\in GL_4(\bar{\mathbf Q}_\ell)$ and homomorphisms $\kappa_1,\dots,\kappa_4:D_\ell\to\bar{\mathbf Q}_\ell^\times$ such that, for every $g\in D_\ell$, $ur(g)u^{-1}$ is upper triangular with diagonal entries $\kappa_1(g),\dots,\kappa_4(g)$ and $|\kappa_i(g)-\delta_i(g)|<1$ for every $i$, and $\kappa_i(g)=\delta_i(g)$ if $g\in I_\ell$.
\end{lemma}
\leanref{Development/AutomorphicArithmetic/AutomorphicVarieties/PadicHodgeTheory/SplAdjUnit.lean}{UpperTheory.V.Loc.spPtL}
\leanref{Development/AutomorphicArithmetic/AutomorphicVarieties/PadicHodgeTheory/SplAdjUnit.lean}{UpperTheory.V.Loc.spl_red}
\begin{proof}
By Definition~\ref{sha-siegel-datum}(2), the restriction of $r$ to $D_\ell$ is $y\circ\rho$, where $\rho:D_\ell\to M_4(B)$ is a lift that is symplectic with multiplier $\theta$ and ordinary, $B$ is a complete local noetherian $\mathcal O$-algebra with residue field $\mathbf F$ that is a domain, with a homomorphism of $\mathcal O$-algebras $x:B\to\mathcal O$ carrying $\rho$ to $r_0$, and $y:B\to\bar{\mathbf Q}_\ell$ is a homomorphism of $\mathcal O$-algebras.
The conditions of ordinarity are preserved by homomorphisms of $\mathcal O$-algebras, since the idempotents $e_i=\prod_{j\ne i}(\rho(s_0)-\delta_j(s_0))$ are polynomials in $\rho(s_0)$ with coefficients in $\mathcal O$.
The $\delta_i(s_0)$ are distinct, as in step (iii) of the proof of Lemma~\ref{sha-counts-vd-linked}. So Lemma~\ref{sha-counts-vd-flag}, with $K=\bar{\mathbf Q}_\ell$, $H=D_\ell$, $I=I_\ell$ and $c_i=\delta_i$, applied to $y\circ\rho$, gives $u$ and the $\kappa_i$ with the triangular form and $\kappa_i=\delta_i$ on $I_\ell$, and $d_i\kappa_i(g)=y\big(\operatorname{tr}(e_i\rho(g))\big)$ with $d_i=\prod_{j\ne i}(\delta_i(s_0)-\delta_j(s_0))\in\mathcal O$, which is nonzero.

Fix $i$ and $g\in D_\ell$, and put $t=\operatorname{tr}(e_i\rho(g))\in B$.
Since $x\circ\rho=r_0$ is diagonal, $x(e_i)$ is the diagonal matrix with entries $\prod_{j\ne i}(\delta_a(s_0)-\delta_j(s_0))$, which vanish for $a\ne i$; so $x(t)=d_i\delta_i(g)$.
The map $\mathcal O\to B$ is injective, since $x$ is a retraction of it. So Lemma~\ref{sha-counts-vd-flag}, applied in the same way to the image of $\rho$ over the fraction field $F$ of $B$, gives characters $\kappa^F_i$ with values in $F$ and $d_i\kappa^F_i(g)=t$. The element $\kappa^F_i(g)$ is integral over $B$: it is a diagonal entry of an upper triangular conjugate of $\rho(g)$, hence a root of the characteristic polynomial of $\rho(g)$, which is monic with coefficients in $B$.
The ring $B$ is henselian, being a complete local ring \cite[Tag 04GM]{Sta}. Lemma~\ref{sha-counts-vd-diag-henselian}, applied to $x$, $y$, $d=d_i$, $t-d_i\delta_i(g)$, which $x$ kills, $v=\kappa^F_i(g)-\delta_i(g)$ and $\kappa=\kappa_i(g)-\delta_i(g)$, gives $|\kappa_i(g)-\delta_i(g)|<1$.
\end{proof}

\begin{lemma}\label{sha-counts-vd-diag-characters}
Let $D$ be a Siegel datum, with the notation of Definition~\ref{sha-siegel-datum}, and put $(w_1,w_2,w_3,w_4)=(k+1,k,1,0)$. Let $r$ be a $\bar{\mathbf Q}_\ell$-point of $D$, write $i\succeq j$ if $r(g)_{ij}\ne0$ for some $g\in G$, and suppose that $i\succeq j$ and $j\succeq i$ imply $i=j$.
Then each $g\mapsto r(g)_{ii}$ is a homomorphism $\kappa_i:G\to\bar{\mathbf Q}_\ell^\times$, and there is a permutation $m$ of $\{1,2,3,4\}$, the identity or $(1\,3)(2\,4)$, with $r(s)_{ii}=\delta_{m(i)}(s)$ for $s\in I_\ell$ and every $i$. Moreover, for each $i$ there is a homomorphism $\nu_i:G\to\bar{\mathbf Q}_\ell^\times$ with open kernel such that $r(g)_{ii}=\chi_{\mathrm{cyc}}(g)^{w_{m(i)}}\nu_i(g)$ for all $g\in G$.
\end{lemma}
\leanref{Development/AutomorphicArithmetic/AutomorphicVarieties/SymplecticGSp4/DiagChar.lean}{UpperTheory.V.GradeZeroInterpolatingElement.diagChar}
\begin{proof}
Let $|\cdot|$ be the absolute value of $\bar{\mathbf Q}_\ell$.
\begin{enumerate}
\item[(i)] \emph{Characters.} For $g,h\in G$, $r(gh)_{ii}=\sum_kr(g)_{ik}r(h)_{ki}$, and a nonzero term with $k\ne i$ would give $i\succeq k\succeq i$. So $r(gh)_{ii}=r(g)_{ii}r(h)_{ii}$, and with $r(1)=1$ this makes $\kappa_i$ a homomorphism to $\bar{\mathbf Q}_\ell^\times$.
\item[(ii)] \emph{Matching.} Let $u$ and $\kappa'_1,\dots,\kappa'_4$ be as in Lemma~\ref{sha-counts-vd-diag-local}. For $g\in D_\ell$, $\sum_i\kappa_i(g)=\operatorname{tr}r(g)=\operatorname{tr}\big(ur(g)u^{-1}\big)=\sum_j\kappa'_j(g)$. The $\kappa'_j$ are pairwise distinct, since $\kappa'_j(s_0)=\delta_j(s_0)$. Distinct characters of $D_\ell$ are linearly independent over $\bar{\mathbf Q}_\ell$ \cite[Chapter~VI, Theorem~4.1]{Lan02}; collecting equal characters in $\sum_i\kappa_i-\sum_j\kappa'_j=0$, the coefficient of $\kappa'_j$ would be $-1$ if no $\kappa_i$ restricted to $\kappa'_j$. So for each $j$ some $\kappa_i$ restricts to $\kappa'_j$; the $i$ so chosen are distinct, and this defines a permutation, whose inverse $m$ satisfies $\kappa_i=\kappa'_{m(i)}$ on $D_\ell$. In particular $r(s)_{ii}=\delta_{m(i)}(s)$ for $s\in I_\ell$.
\item[(iii)] \emph{$\lambda_{m(i)}=\lambda_i$.} The Frobenius element $\sigma_\ell$ lies in $D_\ell$. By Lemma~\ref{sha-counts-vd-conjugate}(a), $|r(\sigma_\ell)_{ii}-\delta_i(\sigma_\ell)|<1$, and by Lemma~\ref{sha-counts-vd-diag-local}, $|r(\sigma_\ell)_{ii}-\delta_{m(i)}(\sigma_\ell)|=|\kappa'_{m(i)}(\sigma_\ell)-\delta_{m(i)}(\sigma_\ell)|<1$. So $|\delta_i(\sigma_\ell)-\delta_{m(i)}(\sigma_\ell)|<1$, and $\delta_i(\sigma_\ell)$ and $\delta_{m(i)}(\sigma_\ell)$ have the same residue; by Lemma~\ref{sha-flagged-residues}(1), $\lambda_i=\lambda_{m(i)}$.
\item[(iv)] \emph{$m(5-i)=5-m(i)$.} By Lemma~\ref{sha-counts-vd-conjugate}(a) and Lemma~\ref{sha-flagged-relation}(2), $i\succeq j$ implies $5-j\succeq5-i$. The entry $(i,i)$ of $r(g)^*r(g)=\theta(g)$ is $\sum_k\epsilon_i\epsilon_kr(g)_{5-k,5-i}r(g)_{ki}$, and a nonzero term with $k\ne i$ would give $5-k\succeq5-i$, hence $i\succeq k$, and $k\succeq i$. So $r(g)_{5-i,5-i}r(g)_{ii}=\theta(g)$. At $s_0$ this reads $\delta_{m(5-i)}(s_0)\delta_{m(i)}(s_0)=\theta(s_0)=\delta_{5-m(i)}(s_0)\delta_{m(i)}(s_0)$, since $\delta_a\delta_{5-a}=\theta$ for every $a$; the $\delta_a(s_0)$ being distinct, $m(5-i)=5-m(i)$.
\item[(v)] \emph{The two permutations.} By (iii), with $\lambda=(1,0,1,0)$, $m$ preserves $\{1,3\}$ and $\{2,4\}$. If $m(1)=1$, then $m(3)=3$, $m(4)=4$ by (iv), and $m(2)=2$. If $m(1)=3$, then $m(3)=1$, $m(4)=2$ by (iv), and $m(2)=4$.
\item[(vi)] \emph{The characters $\nu_i$.} Put $\nu_i=\chi_{\mathrm{cyc}}^{-w_{m(i)}}\kappa_i$, a continuous homomorphism, since $r$ is continuous (step (i) of the proof of Lemma~\ref{sha-counts-vd-linked}). Let $U$ be the intersection of the kernels of $\bar r_0$, $\psi$ and $\omega$, an open normal subgroup of $G$, of index $N$ say.
For $s\in I_\ell\cap U$, $\delta_a(s)=\chi_{\mathrm{cyc}}(s)^{w_a}$ for every $a$, since $\psi(s)=\omega(s)=1$ and $\xi$ is trivial on $I_\ell$; so $\nu_i(s)=1$.
For $p\in S_b$ and $g\in I_p\cap U$, $r(g)=1$ by Definition~\ref{sha-siegel-datum}(2), as $\psi(g)=1$, and $\chi_{\mathrm{cyc}}(g)=1$, as $p\ne\ell$; so $\nu_i(g)=1$.
Let $J$ be the inertia group at a place $v$ of $\bar{\mathbf Q}$, and $g\in J$. If $v$ lies above a prime $p\in T$, then $J=tI_pt^{-1}$ for some $t\in G$, and $t^{-1}g^Nt\in I_p\cap U$, as $U$ is normal of index $N$; so $\nu_i(g)^N=1$. Otherwise $g^N\in J$ and $\bar r_0(g^N)=1$, so $r(g^N)=1$ by Definition~\ref{sha-siegel-datum}(2), and $\chi_{\mathrm{cyc}}(g^N)=1$, as $v$ does not lie above $\ell$; so again $\nu_i(g)^N=1$.
So $\nu_i^N$ is trivial on every inertia group. Its kernel is a closed subgroup of $G$, so it is the intersection of the open subgroups that contain it. Each of these contains the commutators, so it is normal, and its fixed field is unramified over $\mathbf Q$ at every prime; by Minkowski's theorem \cite[Chapter~III, \S2]{Neu99}, which gives $|d_K|>1$ for every number field $K\ne\mathbf Q$, that field is $\mathbf Q$. So $\nu_i^N=1$, and $\nu_i$ takes values in the finite group of $N$-th roots of unity. Its kernel is therefore closed of finite index, hence open.
\end{enumerate}
\end{proof}

\begin{lemma}\label{sha-counts-vd-diag-order}
Let $D$ be a Siegel datum, with the notation of Definition~\ref{sha-siegel-datum}. Let $r$ be a $\bar{\mathbf Q}_\ell$-point of $D$, write $i\succeq j$ if $r(g)_{ij}\ne0$ for some $g\in G$, and suppose that $i\succeq j$ and $j\succeq i$ imply $i=j$.
Let $m$ be an injective map from $\{1,2,3,4\}$ to itself with $r(s)_{ii}=\delta_{m(i)}(s)$ for $s\in I_\ell$ and every $i$. If $i\ne j$ and $i\succeq j$, then $m(i)<m(j)$.
\end{lemma}
\leanref{Development/AutomorphicArithmetic/AutomorphicVarieties/SymplecticGSp4/ContentBoundLevel.lean}{UpperTheory.V.GradeZeroInterpolatingElement.diagonal_character_order_lt_of_stable_flag}
\leanref{Development/AutomorphicArithmetic/AutomorphicVarieties/SymplecticGSp4/ContentBoundLevel.lean}{pat_split}
\begin{proof}
Let $\kappa_i$, the permutation $m'$ and the $\nu_i$ be as in Lemma~\ref{sha-counts-vd-diag-characters}. Then $\delta_{m(i)}(s_0)=r(s_0)_{ii}=\delta_{m'(i)}(s_0)$, so $m=m'$, as the $\delta_a(s_0)$ are distinct. By Lemma~\ref{sha-counts-vd-conjugate}(a) and Lemma~\ref{sha-flagged-relation}(2), $\succeq$ is reflexive and transitive, so it is a partial order.
\begin{enumerate}
\item[(i)] \emph{Reduction.} It suffices to show $m(i)<m(j)$ when, moreover, no $k\notin\{i,j\}$ satisfies $i\succeq k\succeq j$. In general, argue by induction on the number of such $k$: if $k$ is one of them, then every such index for the pair $(i,k)$ or $(k,j)$ is one for $(i,j)$, and $k$ is not one for either pair; so $m(i)<m(k)<m(j)$.
\item[(ii)] \emph{Products.} Let $\Lambda$ be the set of $X\in M_4(\bar{\mathbf Q}_\ell)$ with $X_{ab}=0$ unless $a\succeq b$. It is a subalgebra, since $\succeq$ is reflexive and transitive, and it contains the $r(g)$. For $X,Y\in\Lambda$, $(XY)_{aa}=X_{aa}Y_{aa}$, since $\succeq$ is antisymmetric; and if $(i,j)$ is as in (i), $(XY)_{ij}=X_{ii}Y_{ij}+X_{ij}Y_{jj}$.
\end{enumerate}
Let $(i,j)$ be as in (i), and suppose that $m(j)<m(i)$. Put $\theta'=\kappa_i\kappa_j^{-1}$ and $z(g)=r(g)_{ij}/r(g)_{jj}$. By (ii), $z(gh)=z(g)+\theta'(g)z(h)$; $z$ is continuous, and $z(g)=0$ whenever $r(g)=1$.
\begin{enumerate}
\item[(iii)] \emph{At $\ell$.} For $a\in\{1,2,3,4\}$ the matrix $e_a=\prod_{c\ne a}(r(s_0)-\delta_c(s_0))$ lies in $\Lambda$, and by (ii) its entry $(b,b)$ is $\prod_{c\ne a}(\delta_{m(b)}(s_0)-\delta_c(s_0))$, which vanishes if and only if $m(b)\ne a$. Let $P=e_{m(i)}$ and $Q=e_{m(j)}$, so that $P_{ii}\ne0=P_{jj}$ and $Q_{jj}\ne0=Q_{ii}$. The restriction of $r$ to $D_\ell$ is the image of an ordinary lift, so $Pr(g)Q=0$ for $g\in D_\ell$, as $m(j)<m(i)$. By (ii), the entry $(i,j)$ of $Pr(g)Q$ is $P_{ii}r(g)_{ii}Q_{ij}+(P_{ii}r(g)_{ij}+P_{ij}r(g)_{jj})Q_{jj}$. Subtracting $r(g)_{ii}$ times its value at $g=1$ gives $P_{ii}Q_{jj}r(g)_{ij}=P_{ij}Q_{jj}\big(r(g)_{ii}-r(g)_{jj}\big)$. So $r(g)_{ij}=a\big(r(g)_{ii}-r(g)_{jj}\big)$ with $a=P_{ij}/P_{ii}$, and $z(g)=a(\theta'(g)-1)$, for $g\in D_\ell$.
\item[(iv)] \emph{Local conditions.} If $z(g)=\theta'(g)m_0-m_0$ for $g$ in a subgroup $J$, then $z(g)=\theta'(g)m_1-m_1$ for $g\in tJt^{-1}$, with $m_1=\theta'(t)m_0-z(t)$, by the cocycle relation. So by (iii) $z$ is a coboundary on the decomposition groups at all places above $\ell$, which are conjugate to $D_\ell$. Let $J$ be a subgroup of $G$ and $U$ a subgroup of finite index such that $z$ vanishes on $J\cap U$. Then $z$ is a coboundary on $J$: for $g\in J$ with $\theta'(g)=1$, some $g^n$ with $n\ge1$ lies in $U$, and $nz(g)=z(g^n)=0$; if $\theta'=1$ on $J$ this gives $z=0$ on $J$; otherwise, for $s\in J$ with $\theta'(s)\ne1$ and $g\in J$, $\theta'$ is $1$ at $c=(gs)(sg)^{-1}\in J$, so $z(c)=0$ and $z(gs)=z(c)+\theta'(c)z(sg)=z(sg)$, that is, $z(g)+\theta'(g)z(s)=z(s)+\theta'(s)z(g)$, and $z(g)=\theta'(g)m-m$ with $m=z(s)/(\theta'(s)-1)$. This applies to the inertia groups at places not above $T$, with $U=\ker\bar r_0$, and to $I_p$ for $p\in S_b$, with $U=\ker\psi$, since $r=1$ on $J\cap U$ in both cases by Definition~\ref{sha-siegel-datum}(2). With the conjugation, $z$ is a coboundary on every inertia group.
\item[(v)] \emph{Globally.} Since $k\ge2$, $w_1>w_2>w_3>w_4$, so $w_{m(j)}=w_{m(i)}+e$ with $e\ge1$. By Lemma~\ref{sha-counts-vd-diag-characters}, $\theta'=\psi'^{-1}\chi_{\mathrm{cyc}}^{-e}$ with $\psi'=\nu_j\nu_i^{-1}$, whose kernel contains $\ker\nu_i\cap\ker\nu_j$ and is open. Lemma~\ref{sha-counts-vd-soule}(b) gives $m_0$ with $z(g)=\theta'(g)m_0-m_0$ for all $g\in G$, that is, $r(g)_{ij}=m_0\big(r(g)_{ii}-r(g)_{jj}\big)$.
\item[(vi)] The identity $X_{ij}=m_0(X_{ii}-X_{jj})$ is linear in $X$, so it holds on the span of the $r(g)$, which contains $E_{ii}$ by Lemma~\ref{sha-counts-vd-conjugate}(a); at $X=E_{ii}$ it gives $m_0=0$. So $r(g)_{ij}=0$ for all $g$, contradicting $i\succeq j$.
\end{enumerate}
\end{proof}

\begin{lemma}\label{sha-counts-vd-chart-generators}
Let $\mathcal O$ be a domain with fraction field $E$, let $X$ be an abelian group, let $A'=\bigoplus_{j\in X}A'_j$ be an $X$-graded commutative $\mathcal O$-algebra, and let $x:A'\to\mathcal O$ be a homomorphism of $\mathcal O$-algebras. Let $\tau$ be a finite set with a map $w:\tau\to X$, and let $\varphi$ be a homomorphism of $\mathcal O$-algebras from $A'$ to the ring of formal power series over $E$ in variables $Y_i$, $i\in\tau$, where $Y_i$ has weight $w(i)$, such that $\varphi$ maps $A'_j$ into the series of weight $j$, the constant term of $\varphi(f)$ is $x(f)$, every variable is the only linear term of $\varphi(f)$ for some $f$, and $f\in\Sigma_x$ (Definition~\ref{sha-algebra-graded}) whenever $x(f)=0$ and $\varphi(f)$ has no linear term. Let $\lambda_i(f)$ be the coefficient of $Y_i$ in $\varphi(f)$.
\begin{enumerate}
\item There are $g_i\in A'_{w(i)}$ with $x(g_i)=0$, for $i\in\tau$, such that $\lambda_i(g_i)\ne0$ and $\lambda_{i'}(g_i)=0$ for $i'\ne i$.
\item For such $g_i$, every $j\in X$ and every $p\in A'_j$ with $x(p)=0$, there are a nonzero $c\in\mathcal O$ and $a_k\in\mathcal O$, with $a_k=0$ unless $w(k)=j$, such that $cp-\sum_ka_kg_k\in(\ker x)^2$.
\end{enumerate}
\end{lemma}
\leanref{Development/AutomorphicArithmetic/AutomorphicVarieties/SymplecticGSp4/ASmooth.lean}{homog_gens}
\leanref{Development/AutomorphicArithmetic/AutomorphicVarieties/SymplecticGSp4/ASmooth.lean}{expand}
\begin{proof}
The constant term of a series of weight $j\ne0$ is $0$, so $x$ vanishes on $A'_j$ for $j\ne0$. Each $\lambda_i$ is $\mathcal O$-linear; it vanishes on $A'_j$ for $j\ne w(i)$, since $\varphi$ is graded, on $\mathcal O\cdot1$, and on $(\ker x)^2$, since a product of two series without constant term has no linear term.

(a) Let $f_i$ be such that $Y_i$ is the only linear term of $\varphi(f_i)$, let $h_i$ be its component of degree $w(i)$, and put $g_i=h_i$ if $w(i)\ne0$ and $g_i=h_i-x(h_i)$ if $w(i)=0$. Then $g_i\in A'_{w(i)}$ and $x(g_i)=0$. For $i'$ with $w(i')=w(i)$, $\lambda_{i'}(g_i)=\lambda_{i'}(f_i)$, which is nonzero if $i'=i$ and $0$ otherwise; for the other $i'$, $\lambda_{i'}(g_i)=0$ because $\varphi$ is graded.

(b) For $k\in\tau$ let $a'_k=\lambda_k(p)/\lambda_k(g_k)\in E$ if $w(k)=j$ and $a'_k=0$ otherwise, and choose a nonzero $c'\in\mathcal O$ with $a_k=c'a'_k\in\mathcal O$ for all $k$. The element $q=c'p-\sum_ka_kg_k$ satisfies $x(q)=0$, and $\lambda_i(q)=c'\lambda_i(p)-a_i\lambda_i(g_i)=0$ for every $i$; when $w(i)\ne j$ both terms vanish. So $q\in\Sigma_x$, and $c''q\in(\ker x)^2$ for some nonzero $c''\in\mathcal O$. Replace $c'$ by $c''c'$ and each $a_k$ by $c''a_k$.
\end{proof}

\begin{lemma}\label{sha-counts-vd-chart-kernel}
Let $\mathcal O$ be a domain, let $R$ be a complete local noetherian $\mathcal O$-algebra with residue field $\mathbf F$ and with a weighted presentation $\pi$ with origin $x$ (Definition~\ref{sha-counts-vd-model}(1)) by variables $X_s$ whose weights $w_s\in\mathbf Z^2$ have first coordinates $-1$, $0$ or $1$, and put $z_s=\pi(X_s)$. Let $A\subseteq R$ be the $\mathcal O$-subalgebra of the finite sums of elements of the parts $R_{(0,n)}$, $n\in\mathbf Z$, and $K=\ker x$. Let $N$ be an $\mathcal O$-submodule of $R$ with the following properties.
\begin{enumerate}
\item $ab\in N$ for all $a,b\in A\cap K$.
\item If $c\in\mathcal O$ is nonzero and $cf\in N$, then $f\in N$.
\item For each $s$ with $w_s$ of first coordinate $0$ there are a nonzero $c\in\mathcal O$ and $y\in N\cap A$ with $cz_s-y\in K^2$.
\item For all $s,t$ with $w_s$ and $w_t$ of first coordinates $1$ and $-1$ there are a nonzero $c\in\mathcal O$ and $y\in N\cap A$ with $cz_sz_t-y\in K^3$.
\end{enumerate}
Then $A\cap K\subseteq N$.
\end{lemma}
\leanref{Development/AutomorphicArithmetic/AutomorphicVarieties/SymplecticGSp4/ASmooth.lean}{ker_le}
\leanref{Development/AutomorphicArithmetic/AutomorphicVarieties/SymplecticGSp4/ASmooth.lean}{pow_le}
\leanref{Development/AutomorphicArithmetic/AutomorphicVarieties/SymplecticGSp4/ASmooth.lean}{split_mono}
\begin{proof}
For an exponent vector $n$ write $z^n=\prod_sz_s^{n_s}$; it lies in $R_w$ for the weight $w$ of $n$, so in $A$ if the first coordinate of $w$ is $0$, and in $K$ if $n\ne0$, since $x(\pi(h))$ is the constant term of $h$.
\begin{enumerate}
\item[(i)] Let $k\ge1$, and suppose that $z^n\in N$ for every exponent vector $n$ of degree at least $k$ whose weight has first coordinate $0$. Then $A\cap K^k\subseteq N$. Let $y\in A\cap K^k$. As in the proof of Lemma~\ref{sha-counts-vd-presentations}(f)(i), $y=\pi(h)$ for a series $h$ of order at least $k$, and $y=\sum_{w\in W_0}\pi(h_w)$ for a finite set $W_0$ of weights with first coordinate $0$, where $h_w$ is the sum of the terms of $h$ of weight $w$. By Lemma~\ref{sha-counts-vd-presentations}(b), applied to the exponent vectors of weight $w$ and degree at least $k$, $h_w=\sum_{n\in\mu}b_nX^n$ with $b_n$ of weight $0$. Each $\pi(b_n)z^n$ is $x(\pi(b_n))z^n$, which lies in $N$ by hypothesis, plus $(\pi(b_n)-x(\pi(b_n)))z^n$, a product of two elements of $A\cap K$, which lies in $N$ by (1).
\item[(ii)] Let $n\ne0$ be an exponent vector whose weight has first coordinate $0$. If some $X_s$ with $w_s$ of first coordinate $0$ occurs in $n$, then $n=e_s+n'$. Otherwise the first coordinates, $1$ or $-1$, of the variables occurring in $n$ add up to $0$, so both signs occur, and $n=e_s+e_t+n'$ with $w_s$ and $w_t$ of first coordinates $1$ and $-1$. In both cases the weight of $n'$ has first coordinate $0$.
\item[(iii)] Let $n$ have degree at least $3$. In the first case of (ii), $z^n=z_s\cdot z^{n'}$ with $n'\ne0$; in the second, $z^n=(z_sz_t)\cdot z^{n'}$ with $n'\ne0$. Both are products of two elements of $A\cap K$, which lie in $N$ by (1). By (i), $A\cap K^3\subseteq N$.
\item[(iv)] Let $n$ have degree at least $2$. The argument of (iii) applies unless $n=e_s+e_t$, with $w_s$ and $w_t$ of first coordinates $1$ and $-1$. Then (4) gives $c$ and $y$ with $cz_sz_t-y\in A\cap K^3\subseteq N$, so $cz_sz_t\in N$ and $z_sz_t\in N$ by (2). By (i), $A\cap K^2\subseteq N$.
\item[(v)] Let $n$ have degree $1$. Then $n=e_s$ with $w_s$ of first coordinate $0$, and (3) gives $c$ and $y$ with $cz_s-y\in A\cap K^2\subseteq N$, so $z_s\in N$ by (2). With (iv) and (i), $A\cap K\subseteq N$.
\end{enumerate}
\end{proof}

\begin{lemma}\label{sha-counts-vd-inputs}
Let $D$ be a Siegel datum, with the notation of Definition~\ref{sha-siegel-datum}. For a homomorphism $r:G\to M_4(\bar{\mathbf Q}_\ell)$ write $i\succeq j$ if $r(g)_{ij}\ne0$ for some $g$, as in Lemma~\ref{sha-flagged-relation}, and say that $r$ is \emph{of diagonal type} if $i\succeq j$ and $j\succeq i$ imply $i=j$.
\begin{enumerate}
\item Let $r$ be a $\bar{\mathbf Q}_\ell$-point of $D$. Then $r$ is irreducible (Definition~\ref{sha-counts-vd-model}), or of diagonal type, or $1\succeq3$, $3\succeq1$ and either $i\not\succeq j$ whenever $\lambda_i<\lambda_j$ or $i\not\succeq j$ whenever $\lambda_i>\lambda_j$.
\item Let $r$ be a $\bar{\mathbf Q}_\ell$-point of $D$ of diagonal type. Then either $r(s)_{11}=\theta(s)$ for $s\in I_\ell$ and $i\not\succeq j$ whenever $\mu_i<\mu_j$, or $r(s)_{11}=\alpha(s)$ for $s\in I_\ell$ and $i\not\succeq j$ whenever $\mu_i>\mu_j$.
\item Let $r'$ be a point of $D$ with values in $\mathcal O[\varepsilon]$, $\varepsilon^2=0$, with $r'\equiv r_0$ modulo $\varepsilon$, and let $w$ be $(0,-1)$ or $(0,0)$. If the coefficient of $\varepsilon$ in $r'(g)_{ij}$ vanishes for all $g$ whenever the weight of $(i,j)$ is not $w$, then $r'=r_0$.
\item Let $D$ be finite, with a graded model, with the notation of Definition~\ref{sha-counts-vd-model}, and let $\xi:S^{\mathrm{gr}}\to\bar{\mathbf Q}_\ell$ be a homomorphism of $\mathcal O$-algebras such that either $r_\xi(g)_{ij}=0$ for all $g$ whenever $\lambda_i<\lambda_j$, or $r_\xi(g)_{ij}=0$ for all $g$ whenever $\lambda_i>\lambda_j$. Then there is a homomorphism of $\mathcal O$-algebras $\xi_0:S^{\mathrm{gr}}\to\bar{\mathbf Q}_\ell$ with $r_{\xi_0}(g)_{ij}=r_\xi(g)_{ij}$ if $\lambda_i=\lambda_j$ and $r_{\xi_0}(g)_{ij}=0$ otherwise.
\item For every ring $\mathcal O'$ as in Definition~\ref{sha-character-parts}, every character $\psi':G\to\mathcal O'^\times$ with open kernel and every $k'\ge1$, the length $V_\ell(k',\psi')$ is finite.
\item Let $D$ have a graded model, with the notation of Definition~\ref{sha-counts-vd-model}. There is a homomorphism of $\mathcal O$-algebras from $A$ to a ring of formal power series over $E$ in $r$ variables of weight $-1$, $r'+1$ of weight $1$, one of weight $2$, $m$ of weight $0$ and no others, with the properties of condition (1) of Lemma~\ref{sha-algebra-cotangent} for $A$ and $x$: it maps $A_n$ into the series of weight $n$, the constant term of the image of $f$ is $x(f)$, every variable is the only linear term of the image of some $f$, and $f\in\Sigma_x$ whenever $x(f)=0$ and the image of $f$ has no linear term. It is obtained from the chart $\varphi$ of Definition~\ref{sha-counts-vd-model}(i) by setting its variable of weight $(-1,1)$ equal to $1$; a variable of weight $(w_1,w_2)$ of $\varphi$ receives the weight $w_1+w_2$.
\end{enumerate}
\end{lemma}
\leanref{Development/AutomorphicArithmetic/AutomorphicVarieties/SymplecticGSp4/TriComb.lean}{UpperTheory.V.GradeZeroInterpolatingElement.tri}
\leanref{Development/AutomorphicArithmetic/AutomorphicVarieties/SymplecticGSp4/ExistsThirdIndexLinkedPair.lean}{UpperTheory.V.GradeZeroInterpolatingElement.exists_third_index_linked_pair}
\leanref{Development/AutomorphicArithmetic/AutomorphicVarieties/SymplecticGSp4/DiagCases.lean}{UpperTheory.V.GradeZeroInterpolatingElement.diagCases}
\leanref{Development/AutomorphicArithmetic/AutomorphicVarieties/SymplecticGSp4/TanGlob.lean}{UpperTheory.V.GradeZeroInterpolatingElement.tanGlob}
\leanref{Development/AutomorphicArithmetic/AutomorphicVarieties/SymplecticGSp4/PtLimit.lean}{UpperTheory.V.GradeZeroInterpolatingElement.ptLimit}
\leanref{Development/AutomorphicArithmetic/AutomorphicVarieties/LevelCohomology/ExistsAb.lean}{UpperTheory.V.X.soFin}
\leanref{Development/AutomorphicArithmetic/AutomorphicVarieties/SymplecticGSp4/ASmooth.lean}{UpperTheory.V.GradeZeroInterpolatingElement.aSmooth}
\leanref{Development/AutomorphicArithmetic/AutomorphicVarieties/SymplecticGSp4/PtLimit.lean}{limit_graded}
\leanref{Development/AutomorphicArithmetic/AutomorphicVarieties/PlacesOfRationals/Minkowski.lean}{UpperTheory.V.B3.minkowski}
\begin{proof}
(a) By Lemma~\ref{sha-counts-vd-conjugate}(a), each $r(g)$ is symplectic and the span of the $r(g)$ contains $E_{11},\dots,E_{44}$. So by Lemma~\ref{sha-flagged-relation}, $\succeq$ is reflexive and transitive, $i\succeq j$ implies $5-j\succeq5-i$, and $r$ is irreducible if and only if $i\succeq j$ for all $i$ and $j$. By Lemma~\ref{sha-counts-vd-linked}, $\succeq$ also has property (3) of Lemma~\ref{sha-flagged-relation-cases}, and that lemma gives the three alternatives; the second is that $r$ is of diagonal type.

(b) Lemma~\ref{sha-counts-vd-diag-characters} gives a permutation $m$, the identity or $(1\,3)(2\,4)$, with $r(s)_{ii}=\delta_{m(i)}(s)$ for $s\in I_\ell$, and by Lemma~\ref{sha-counts-vd-diag-order}, $m(i)<m(j)$ whenever $i\ne j$ and $i\succeq j$. Recall $\mu=(1,1,0,0)$.
If $m$ is the identity, then $r(s)_{11}=\delta_1(s)=\theta(s)$ for $s\in I_\ell$; and if $\mu_i<\mu_j$, then $i\in\{3,4\}$ and $j\in\{1,2\}$, so $m(i)>m(j)$ and $i\not\succeq j$.
If $m=(1\,3)(2\,4)$, then $r(s)_{11}=\delta_3(s)=\alpha(s)$ for $s\in I_\ell$; and if $\mu_i>\mu_j$, then $i\in\{1,2\}$ and $j\in\{3,4\}$, so $m(i)\in\{3,4\}$ and $m(j)\in\{1,2\}$, and again $i\not\succeq j$.

(c) Write $r'(g)=r_0(g)+\varepsilon Y(g)$ with $Y(g)\in M_4(\mathcal O)$, and put $Z(g)=Y(g)r_0(g)^{-1}$, so that $Z(g)_{ij}=Y(g)_{ij}\delta_j(g)^{-1}$.
\begin{enumerate}
\item[(i)] From $r'(gh)=r'(g)r'(h)$, $Z(gh)_{ij}=Z(g)_{ij}+\delta_i(g)\delta_j(g)^{-1}Z(h)_{ij}$. From $r'(g)^*r'(g)=\theta(g)=r_0(g)^*r_0(g)$, the coefficient of $\varepsilon$ gives $r_0(g)^*(Z(g)^*+Z(g))r_0(g)=0$, so $Z(g)\in\mathfrak{sp}_4(\mathcal O)$. Each $Z_{ij}$ is locally constant modulo every power of $\mathfrak m_{\mathcal O}$, so it is continuous as a function to $\bar{\mathbf Q}_\ell$. $Z(g)=0$ whenever $r'(g)=1$, in particular for $g$ with $\bar r_0(g)=1$ in an inertia group at a place not above $T$, and for $g\in I_p\cap\ker\psi$ with $p\in S_b$ (Definition~\ref{sha-siegel-datum}(2)). Every $Z(g)$ is supported on the entries of weight $w$, and the restriction of $r'$ to $D_\ell$ is the image of an ordinary lift, so Lemma~\ref{sha-counts-vd-tangent} applies to $Z$ on $D_\ell$.
\item[(ii)] Let $z:G\to\bar{\mathbf Q}_\ell$ and a character $\gamma:G\to\bar{\mathbf Q}_\ell^\times$ satisfy $z(gh)=z(g)+\gamma(g)z(h)$, let $z(g)=0$ whenever $r'(g)=1$, and let $z$ be a coboundary on $I_\ell$, that is, $z(g)=\gamma(g)m-m$ there for some $m$. Then $z$ is a coboundary on every inertia group. First, if $J$ is a subgroup of $G$ and $U$ a subgroup of finite index with $z=0$ on $J\cap U$, then $z$ is a coboundary on $J$: for $x\in J$ with $\gamma(x)=1$, some $x^n$ with $n\ge1$ lies in $U$, and $nz(x)=z(x^n)=0$; if $\gamma=1$ on $J$ this gives $z=0$ on $J$, and otherwise, for $s\in J$ with $\gamma(s)\ne1$ and $g\in J$, $(sg)^{-1}gs$ lies in $J$ and $\gamma$ is $1$ on it, so $z$ vanishes on it and $z(gs)=z(sg)$, that is, $z(g)+\gamma(g)z(s)=z(s)+\gamma(s)z(g)$, and $z(g)=\gamma(g)m-m$ with $m=z(s)/(\gamma(s)-1)$. This applies to the inertia groups at places not above $T$, with $U=\ker\bar r_0$, and to $I_p$ for $p\in S_b$, with $U=\ker\psi$. Second, a coboundary on $J$ is a coboundary on $tJt^{-1}$ for $t\in G$, of $\gamma(t)m-z(t)$; and the inertia groups at the places above a prime are conjugate in $G$.
\item[(iii)] Let $w=(0,0)$. The entries of weight $(0,0)$ are the diagonal ones, so $Z(g)$ is diagonal, and $z=Z_{ii}$ satisfies $z(gh)=z(g)+z(h)$. By Lemma~\ref{sha-counts-vd-tangent}(b), $z$ vanishes on $I_\ell$, and by (ii) with $\gamma=1$ it vanishes on every inertia group. Its kernel is a closed subgroup of $G$, so it is the intersection of the open subgroups that contain it. Each of these contains the commutators, so it is normal, and its fixed field is unramified over $\mathbf Q$ at every prime; by Minkowski's theorem \cite[Chapter~III, \S2]{Neu99}, which gives $|d_F|>1$ for every number field $F\ne\mathbf Q$, that field is $\mathbf Q$. So $z=0$, $Y=0$ and $r'=r_0$.
\item[(iv)] Let $w=(0,-1)$. The entries of weight $(0,-1)$ are $(3,1)$ and $(4,2)$, and $Z(g)_{42}=Z(g)_{31}$ since $Z(g)\in\mathfrak{sp}_4$. By (i), $z=Z_{31}$ satisfies $z(gh)=z(g)+\Psi(g)^{-1}z(h)$, as $\delta_3\delta_1^{-1}=\alpha\theta^{-1}=\Psi^{-1}$. By Lemma~\ref{sha-counts-vd-tangent}(a), $z$ is a coboundary on $D_\ell$, hence on $I_\ell$, so by (ii) on every inertia group, and on the decomposition groups at the places above $\ell$, which are conjugate to $D_\ell$. Lemma~\ref{sha-counts-vd-soule}(b), with $\psi'=\psi$, $e=k$ and $\Psi^{-1}=\psi^{-1}\chi_{\mathrm{cyc}}^{-k}$ in place of its $\theta$, gives $m$ with $z(g)=\Psi(g)^{-1}m-m$ for all $g$. The entry $(3,1)$ of $r'(g_0)+\theta(g_0)r'(g_0^{-1})$ vanishes, and its coefficient of $\varepsilon$ is $\theta(g_0)z(g_0)+\theta(g_0)\theta(g_0^{-1})z(g_0^{-1})=\theta(g_0)z(g_0)+z(g_0^{-1})$. Substituting, $\big(\theta(g_0)(\Psi(g_0)^{-1}-1)+\Psi(g_0)-1\big)m=-(\alpha(g_0)-1)(\Psi(g_0)-1)m=0$, since $\theta=\alpha\Psi$. As $\bar\alpha(g_0)\ne1$ and $\bar\Psi(g_0)\ne1$, $m=0$. So $Z_{31}=Z_{42}=0$, $Y=0$ and $r'=r_0$.
\end{enumerate}

(d) Let $p:\mathbf Z^2\to\mathbf Z$ be $(w_1,w_2)\mapsto w_1$ in the first case and $(w_1,w_2)\mapsto-w_1$ in the second. For the weight $w$ of $(i,j)$, $p(w)=\pm(\lambda_i-\lambda_j)$; so $p(w)=0$ if and only if $\lambda_i=\lambda_j$, and $r_\xi(g)_{ij}=0$ for all $g$ whenever $p(w)<0$.
By Lemma~\ref{sha-counts-vd-span}(a), $S_w$ is spanned by products of entries of the $r^{\mathrm u}(g)$ whose weights add up to $w$. If $p(w)<0$, one of the factors has a weight $w'$ with $p(w')<0$, so $\xi$ vanishes on $S_w$.
By Lemma~\ref{sha-counts-vd-structure}(a), $S^{\mathrm{gr}}=\bigoplus_wS_w$. Let $\xi_0$ be the $\mathcal O$-linear map equal to $\xi$ on the $S_w$ with $p(w)=0$ and to $0$ on the $S_w$ with $p(w)\ne0$; then $\xi_0(1)=1$.
For $f\in S_w$ and $f'\in S_{w'}$, $\xi_0(ff')=\xi_0(f)\xi_0(f')$: both sides are $\xi(f)\xi(f')$ if $p(w)=p(w')=0$; both vanish if exactly one of $p(w)$ and $p(w')$ is $0$, or if $p(w+w')\ne0$; and if $p(w),p(w')\ne0=p(w+w')$, one of $p(w)$ and $p(w')$ is negative, so $\xi(ff')=\xi(f)\xi(f')=0$.
So $\xi_0$ is a homomorphism of $\mathcal O$-algebras. As $r^{\mathrm u}(g)_{ij}\in S_w$ for the weight $w$ of $(i,j)$, $r_{\xi_0}(g)_{ij}$ is $r_\xi(g)_{ij}$ if $\lambda_i=\lambda_j$ and $0$ otherwise.

(e) This is Lemma~\ref{sha-counts-vd-soule}(a).

(f) Let $\varphi$ be the chart of Definition~\ref{sha-counts-vd-model}(i), with variables $Y_i$, $i\in\tau$, of weights $w(i)\in\mathbf Z^2$, and let $x_A$ be the restriction of $x$ to $A$.
\begin{enumerate}
\item[(i)] \emph{The weights.} By the counts in Definition~\ref{sha-counts-vd-model}(i), exactly one variable $Y_b$ has weight $(-1,1)$, and the first coordinates of the weights of the others are $0$ or $1$. Let $\tau'=\tau\setminus\{b\}$, and give $Y_i$, $i\in\tau'$, the weight $w'(i)=w(i)_1+w(i)_2$. The variables of weights $(0,-1)$, $(1,0)$, $(0,1)$, $(1,1)$ and $(0,0)$ receive the weights $-1$, $1$, $1$, $2$ and $0$, so there are $r$ variables of weight $-1$, $r'+1$ of weight $1$, one of weight $2$, $m$ of weight $0$ and no others.
\item[(ii)] \emph{Setting $Y_b=1$.} For an exponent vector $n$ on $\tau'$ let $\iota(n)$ be the exponent vector on $\tau$ with $\iota(n)_i=n_i$ for $i\in\tau'$ and $\iota(n)_b=\sum_{i:w(i)_1=1}n_i$. The map $\iota$ is additive and injective, the weight of $Y^{\iota(n)}$ is $(0,w'(n))$, and every exponent vector $n''$ on $\tau$ whose weight has first coordinate $0$ is in the image of $\iota$, since that coordinate is $\sum_{i:w(i)_1=1}n''_i-n''_b$.
For $f\in A$, $\varphi(f)$ is a finite sum of series of weights $(0,n)$, so its monomials are of the form $Y^{\iota(n)}$. Let $\psi(f)$ be the series in the $Y_i$, $i\in\tau'$, whose coefficient at $Y^n$ is that of $\varphi(f)$ at $Y^{\iota(n)}$; it is $\varphi(f)$ with $Y_b$ set to $1$. Then $\psi$ is a homomorphism of $\mathcal O$-algebras: in a product of two series supported on the monomials $Y^{\iota(n)}$, the coefficient at $Y^{\iota(n)}$ is the sum of the products of the coefficients at $Y^{\iota(n_1)}$ and $Y^{\iota(n_2)}$ over $n_1+n_2=n$, as $\iota$ is additive and injective.
It maps $A_n=R_{(0,n)}$ into the series of weight $n$, and the constant term of $\psi(f)$, that of $\varphi(f)$, is $x(f)$. The coefficient of each $Y_j$ in $\psi(f)$ vanishes on $(\ker x_A)^2$.
\item[(iii)] \emph{Generators.} Lemma~\ref{sha-counts-vd-chart-generators}(a), for $R^{\mathrm{gr}}$ and $\varphi$, gives $g_i\in R_{w(i)}$ with $x(g_i)=0$ such that the linear part of $\varphi(g_i)$ is $\lambda_iY_i$ with $\lambda_i\ne0$. For $k\in\tau'$ put $G_k=g_k$ if $w(k)_1=0$ and $G_k=g_kg_b$ if $w(k)_1=1$; then $G_k\in R_{(0,w'(k))}\subseteq A$ and $x(G_k)=0$. Let $c_{jk}$ be the coefficient of $Y_j$ in $\psi(G_k)$, that is, of $Y^{\iota(e_j)}$ in $\varphi(G_k)$.
If $w(j)_1=0$, then $Y^{\iota(e_j)}=Y_j$, and $c_{jk}$ is $\lambda_j$ if $k=j$ and $0$ for the other $k$: for $w(k)_1=1$, $\varphi(g_k)\varphi(g_b)$ is a product of two series without constant term.
If $w(j)_1=1$, then $Y^{\iota(e_j)}=Y_jY_b$, and for $w(k)_1=1$, $c_{jk}$, the coefficient of $Y_jY_b$ in $\varphi(g_k)\varphi(g_b)$, is $\lambda_k\lambda_b$ if $j=k$ and $0$ otherwise, since the terms of degree at most $1$ of the two factors are $\lambda_kY_k$ and $\lambda_bY_b$.
So if $v_k\in E$ satisfy $\sum_kv_kc_{jk}=0$ for every $j\in\tau'$, then $v_j=0$, first for the $j$ with $w(j)_1=0$ and then for the others: the matrix $(c_{jk})$ is invertible.
\item[(iv)] \emph{The span.} For every $f\in A$ with $x(f)=0$ there are a nonzero $c\in\mathcal O$ and $a_k\in\mathcal O$ with $cf-\sum_ka_kG_k\in(\ker x_A)^2$.
Let $L\subseteq R$ be the $\mathcal O$-module spanned by the $G_k$ and $(\ker x_A)^2$, and $N$ the set of $y\in R$ with $cy\in L$ for some nonzero $c\in\mathcal O$. The weights of the variables of the presentation of $R$ have first coordinates $-1$, $0$ or $1$, and it suffices to check the conditions of Lemma~\ref{sha-counts-vd-chart-kernel} for $N$. Conditions (1) and (2) hold by construction.
For (3), let $w_s$ have first coordinate $0$. Lemma~\ref{sha-counts-vd-chart-generators}(b), applied to $z_s\in R_{w_s}$, gives $c\ne0$ and $a_k$, zero unless $w(k)=w_s$, with $cz_s-\sum_ka_kg_k\in(\ker x)^2$; then $a_b=0$, and $g_k=G_k$ whenever $a_k\ne0$, so $y=\sum_ka_kg_k\in N\cap A$.
For (4), let $w_s$ and $w_t$ have first coordinates $1$ and $-1$. In the same way $c_1z_s-u$ and $c_2z_t-v$ lie in $(\ker x)^2$, where $u=\sum_ka_kg_k$ with $a_k=0$ unless $w(k)=w_s$, so that $k\ne b$ and $w(k)_1=1$, and $v=ag_b$ for some $a\in\mathcal O$, since $b$ is the only element of $\tau$ with $w(b)_1=-1$. As $u,v\in\ker x$,
\[
 c_1c_2z_sz_t-uv=(c_1z_s-u)(c_2z_t-v)+(c_1z_s-u)v+u(c_2z_t-v)\in(\ker x)^3,
\]
and $uv=\sum_kaa_kG_k\in N\cap A$. So Lemma~\ref{sha-counts-vd-chart-kernel} gives $f\in N$.
\item[(v)] \emph{The chart.} Let $i\in\tau'$. Since $(c_{jk})$ is invertible, there are $v_k\in E$ with $\sum_kc_{jk}v_k$ equal to $1$ for $j=i$ and to $0$ otherwise; with a nonzero $c\in\mathcal O$ such that the $a_k=cv_k$ lie in $\mathcal O$, the linear part of $\psi(\sum_ka_kG_k)$ is $cY_i$. Let $f\in A$ with $x(f)=0$ and no linear term in $\psi(f)$, and take $c$ and the $a_k$ of (iv). The coefficients of the $Y_j$ vanish on $(\ker x_A)^2$, so $\sum_ka_kc_{jk}=0$ for every $j$, and $a=0$ by (iii). Hence $cf\in(\ker x_A)^2$, and $f\in\Sigma_x$.
\end{enumerate}
So $\psi$ has the properties stated in (f).
\end{proof}

\begin{lemma}\label{sha-counts-vd-irreducible}
Let $D$ be a finite Siegel datum with a graded model, with the notation of Definition~\ref{sha-counts-vd-model}, and let $\xi:S^{\mathrm{gr}}\to\bar{\mathbf Q}_\ell$ be a homomorphism of $\mathcal O$-algebras. Use the relation $\succeq$ and the diagonal type of Lemma~\ref{sha-counts-vd-inputs}.
\begin{enumerate}
\item If $r_\xi(g)_{ij}=0$ for all $g$ whenever $\mu_i<\mu_j$, then $\xi\circ q_A$ vanishes on $A_n$ for all $n<0$. If $r_\xi(g)_{ij}=0$ for all $g$ whenever $\mu_i>\mu_j$, then $\xi\circ q_A$ vanishes on $A_n$ for all $n>0$.
\item If $\xi\circ q_A$ vanishes on $A_n$ for all $n>0$, then $r_\xi(g)_{13}=0$ for all $g$. If it vanishes on $A_n$ for all $n<0$, then $r_\xi(g)_{31}=0$ for all $g$.
\item Let $\xi':S^{\mathrm{gr}}\to\bar{\mathbf Q}_\ell$ be a homomorphism of $\mathcal O$-algebras that agrees with $\xi$ on $q(A_0)$. Then $r_\xi$ and $r_{\xi'}$ have the same diagonal entries, and if $r_\xi$ is of diagonal type, so is $r_{\xi'}$.
\item Suppose that $D$ has no endoscopic $\bar{\mathbf Q}_\ell$-points. If $\xi\circ q_A$ is two-signed (Definition~\ref{sha-algebra-graded}), then the $r_\xi(g)$ span $M_4(\bar{\mathbf Q}_\ell)$.
\end{enumerate}
\end{lemma}
\leanref{Development/AutomorphicArithmetic/AutomorphicVarieties/SymplecticGSp4/SpanZero.lean}{UpperTheory.V.GradeZeroInterpolatingElement.signs}
\leanref{Development/AutomorphicArithmetic/AutomorphicVarieties/SymplecticGSp4/TwoIrr.lean}{UpperTheory.V.GradeZeroInterpolatingElement.twoIrr}
\begin{proof}
(a) By Lemma~\ref{sha-counts-vd-structure}(b) and Lemma~\ref{sha-counts-vd-span}(a), $q_A(A_n)=S_{(0,n)}$ is spanned by products of entries of the $r^{\mathrm u}(g)$ whose weights add up to $(0,n)$. If $n<0$, the second coordinates of these weights add up to $n<0$, so one of the factors is an entry $(i,j)$ with $\mu_i<\mu_j$, which $\xi$ kills by hypothesis. The case $n>0$ is the same.

(b) The entry $(1,3)$ has weight $(0,1)$, so $r^{\mathrm u}(g)_{13}\in S_{(0,1)}=q_A(A_1)$; the entry $(3,1)$ has weight $(0,-1)$.

(c) The diagonal entries have weight $(0,0)$, so they lie in $S_{(0,0)}=q(A_0)$. Let $i\ne j$ with $r_{\xi'}(g)_{ij}\ne0$ and $r_{\xi'}(h)_{ji}\ne0$. The product $r^{\mathrm u}(g)_{ij}r^{\mathrm u}(h)_{ji}$ has weight $(0,0)$, so $\xi(r^{\mathrm u}(g)_{ij})\xi(r^{\mathrm u}(h)_{ji})=\xi'(r^{\mathrm u}(g)_{ij})\xi'(r^{\mathrm u}(h)_{ji})\ne0$, and $i\succeq j\succeq i$ for $r_\xi$.

(d) By Lemma~\ref{sha-counts-vd-points}(d) there is $t$ such that $r=\tau(t)r_\xi\tau(t)^{-1}$ is a $\bar{\mathbf Q}_\ell$-point of $D$. Conjugation by $\tau(t)$ multiplies the entries by units, so $r$ and $r_\xi$ have the same relation $\succeq$, and $r_\xi$ is irreducible if $r$ is. Apply Lemma~\ref{sha-counts-vd-inputs}(a) to $r$.
If $r$ is of diagonal type, Lemma~\ref{sha-counts-vd-inputs}(b) and (a) above show that $\xi\circ q_A$ vanishes on all $A_n$ with $n<0$, or on all $A_n$ with $n>0$, which is excluded since $\xi\circ q_A$ is two-signed.
Suppose that $1\succeq3$, $3\succeq1$, and $i\not\succeq j$ whenever $\lambda_i<\lambda_j$, or whenever $\lambda_i>\lambda_j$. Lemma~\ref{sha-counts-vd-inputs}(d), applied to $\xi$, gives $\xi_0$, and Lemma~\ref{sha-counts-vd-points}(d) gives $t_0$ such that $r_1=\tau(t_0)r_{\xi_0}\tau(t_0)^{-1}$ is a $\bar{\mathbf Q}_\ell$-point. Every $r_1(g)$ is block diagonal. Since $\lambda_1=\lambda_3$, the entries $(1,3)$ and $(3,1)$ of $r_{\xi_0}$ are those of $r_\xi$, so $1\succeq3$ and $3\succeq1$ for $r_1$. By Lemma~\ref{sha-counts-vd-conjugate}(a), each $r_1(g)$ is symplectic and the span of the $r_1(g)$ contains $E_{11},\dots,E_{44}$; so by Lemma~\ref{sha-flagged-relation}(4) the first blocks of the $r_1(g)$ span $M_2(\bar{\mathbf Q}_\ell)$, and $r_1$ is endoscopic, contrary to the hypothesis.
So $r$ is irreducible, and so is $r_\xi$.
\end{proof}

\begin{lemma}\label{sha-counts-vd-conditions}
Let $D$ be a finite Siegel datum without endoscopic $\bar{\mathbf Q}_\ell$-points that is rigid at each of its irreducible $\bar{\mathbf Q}_\ell$-points, and suppose that $D$ has a graded model; use the notation of Definition~\ref{sha-counts-vd-model}. Then $A$, $x$ and the families $F=(F_0,\dots,F_r)$, $E=(E_1,\dots,E_{r'})$ and $Z=(Z_1,\dots,Z_m)$, with $s=s'=0$ and $\Omega=\bar{\mathbf Q}_\ell$, satisfy the hypotheses of Lemma~\ref{sha-algebra-cotangent} on $A$, $x$ and the families, and its conditions (1) to (4). The ideal of that lemma is $I$, and its set $Y$ is the set of homomorphisms of $\mathcal O$-algebras $A\to\bar{\mathbf Q}_\ell$ that vanish on $I$.
\end{lemma}
\leanref{Development/AutomorphicArithmetic/AutomorphicVarieties/SymplecticGSp4/QuotientSaturatedSquareTorsionFree.lean}{UpperTheory.V.GradeZeroInterpolatingElement.m1b}
\leanref{Development/AutomorphicArithmetic/AutomorphicVarieties/CommutativeAlgebra/SpansOfDer.lean}{UpperTheory.V.Alg.spans_of_der}
\leanref{Development/AutomorphicArithmetic/AutomorphicVarieties/CommutativeAlgebra/ChartRank.lean}{UpperTheory.V.Alg.chart_rank}
\leanref{Development/AutomorphicArithmetic/AutomorphicVarieties/SymplecticGSp4/TorusConjugateCornerEntryEq.lean}{UpperTheory.V.GradeZeroInterpolatingElement.m3}
\leanref{Development/AutomorphicArithmetic/AutomorphicVarieties/SymplecticGSp4/ExistsLiftOfDerivationOnIdeal.lean}{UpperTheory.V.GradeZeroInterpolatingElement.m4}
\begin{proof}
\emph{Hypotheses.} The ring $\mathcal O$ is a complete discrete valuation ring of characteristic $0$. By Lemma~\ref{sha-counts-vd-structure}(a) and (c), $A$ is $\mathbf Z$-graded, $A_0$ is a complete local noetherian $\mathcal O$-algebra with residue field $\mathbf F$, $A$ is generated over $A_0$ by finitely many elements of degrees $-2$, $-1$, $1$ and $2$, and $x$ vanishes on $A_n$ for $n\ne0$. The $F_i$ lie in $R_{(0,-1)}=A_{-1}$, the $E_i$ in $A_1$, and the $Z_i$ in $A_0\cap\ker x$; and $r\ge1$.

\emph{Condition (2)} is Lemma~\ref{sha-counts-vd-span}(b).

\emph{Condition (1).} Its first sentence is Lemma~\ref{sha-counts-vd-inputs}(f), with $s=s'=0$. For the second, recall from Step 1 of the proof of Lemma~\ref{sha-algebra-cotangent} that the coefficients of the variables of weight $j$ in the image under a chart define an injective $\mathcal O$-linear map from the cotangent lattice of weight $j$ to $E^{n_j}$, $n_j$ the number of these variables. The chart of Lemma~\ref{sha-counts-vd-inputs}(f) has no variable of weight $-2$, so $L_{-2}=0$. Let $(j,w,H,n_j)$ be $(-1,(0,-1),F,r)$ or $(0,(0,0),Z,m)$.
\begin{enumerate}
\item[(i)] By the chart of Lemma~\ref{sha-counts-vd-inputs}(f), $L_j$ has rank at most $n_j$.
\item[(ii)] Let $L_w^{\mathrm{gr}}$ be the cotangent lattice of weight $w$ of $R^{\mathrm{gr}}$ at $x$. It is finitely generated by Lemma~\ref{sha-counts-vd-presentations}(f), and the chart $\varphi$ of Definition~\ref{sha-counts-vd-model}(i), with its $n_j$ variables of weight $w$, embeds it in $E^{n_j}$ with image containing $n_j$ linearly independent vectors, as in Step 1 of the proof of Lemma~\ref{sha-algebra-cotangent}; so it has rank $n_j$.
\item[(iii)] The inclusion $A\subseteq R^{\mathrm{gr}}$ maps $\Sigma_x$ of $A$ into $\Sigma_x$ of $R^{\mathrm{gr}}$, so it induces a map $L_j\to L_w^{\mathrm{gr}}$, which is surjective since $A_j=R_w$, and carries $\langle H\rangle$ onto the span of the images of $H$. Its kernel has rank $0$ by (i) and (ii), and it is $0$, since $A/\Sigma_x$ is torsion free (Step 2 of the proof of Lemma~\ref{sha-algebra-cotangent}). So it suffices to show that every element of $L_w^{\mathrm{gr}}$ has a nonzero multiple in the span of the images of $H$.
\item[(iv)] Suppose not. As $L_w^{\mathrm{gr}}$ is a finitely generated module over the discrete valuation ring $\mathcal O$, there is an $\mathcal O$-linear $\lambda:L_w^{\mathrm{gr}}\to\mathcal O$ that vanishes on the images of $H$ and is nonzero. The ideal $\ker x$ of $R^{\mathrm{gr}}$ is generated by homogeneous elements, and so is its square. Hence $D(f)=\lambda\big(\overline{f_w-x(f_w)}\big)$, with $f_w$ the component of weight $w$ of $f\in R^{\mathrm{gr}}$, is a derivation of $R^{\mathrm{gr}}$ at $x$ with values in $\mathcal O$; it vanishes on the $R_{w'}$ with $w'\ne w$ and on $H$, and it is nonzero.
\item[(v)] By Lemma~\ref{sha-counts-vd-presentations}(f), $D$ extends to a derivation $D'$ of $R$ at $x$. It vanishes on the $R_{w'}$ with $w'\ne w$, which lie in $R^{\mathrm{gr}}$, and so on the generators of $\ker q$ other than those of $H$, which have other weights; it vanishes on $H$; and all these generators lie in $\ker x$. So $D'$ vanishes on $\ker q$, and $f\mapsto x(f)+\varepsilon D'(f)$ factors through a homomorphism of $\mathcal O$-algebras $y:S\to\mathcal O[\varepsilon]$. Then $y\circ r^{\mathrm u}$ is a point of $D$ with values in $\mathcal O[\varepsilon]$; by Definition~\ref{sha-counts-vd-model}(iii) it is congruent to $r_0$ modulo $\varepsilon$, and the coefficient of $\varepsilon$ in its entry $(i,j)$ is $D'(f)$ for some $f\in R_{w''}$, $w''$ the weight of $(i,j)$, which vanishes when $w''\ne w$. By Lemma~\ref{sha-counts-vd-inputs}(c), $y\circ r^{\mathrm u}=r_0$. So $y$ and the homomorphism $S\to\mathcal O[\varepsilon]$ induced by $x$ agree on the entries of the $r^{\mathrm u}(g)$, and they are equal by Definition~\ref{sha-counts-vd-model}(ii). Hence $D'=0$, and $D=0$, a contradiction.
\end{enumerate}

\emph{Condition (3).} Let $x_1\in Y$ vanish on every $A_n$ with $n\ne0$. By Lemma~\ref{sha-counts-vd-points}(a), $x_1=\xi_1\circ q_A$, and by Lemma~\ref{sha-counts-vd-irreducible}(b), $r_{\xi_1}(g)_{13}=0$ for all $g$. By Lemma~\ref{sha-counts-vd-points}(d), $r=\tau(t)r_{\xi_1}\tau(t)^{-1}$ is a $\bar{\mathbf Q}_\ell$-point for some $t$, with the same relation $\succeq$ as $r_{\xi_1}$, so $1\not\succeq3$. By Lemma~\ref{sha-counts-vd-conjugate}(a) and Lemma~\ref{sha-flagged-relation}(3), $r$ is not irreducible; so by Lemma~\ref{sha-counts-vd-inputs}(a), $r$ and $r_{\xi_1}$ are of diagonal type.
Let $x'\in Y$ agree with $x_1$ on $A_0$. By Lemma~\ref{sha-counts-vd-points}(a), $x'=\xi\circ q_A$, and $\xi$ agrees with $\xi_1$ on $q(A_0)$; by Lemma~\ref{sha-counts-vd-irreducible}(c), $r_\xi$ is of diagonal type with the diagonal entries of $r_{\xi_1}$. Apply Lemma~\ref{sha-counts-vd-inputs}(b) to a $\bar{\mathbf Q}_\ell$-point $\tau(t')r_\xi\tau(t')^{-1}$ (Lemma~\ref{sha-counts-vd-points}(d)), which has the diagonal entries and the relation $\succeq$ of $r_\xi$, and then Lemma~\ref{sha-counts-vd-irreducible}(a): either $r_{\xi_1}(s_0)_{11}=\theta(s_0)$ and $x'$ vanishes on the $A_n$ with $n<0$, or $r_{\xi_1}(s_0)_{11}=\alpha(s_0)$ and $x'$ vanishes on the $A_n$ with $n>0$.
Since $s_0\in I_\ell$, $\psi(s_0)=\omega(s_0)=1$ and $\xi$ is trivial on $I_\ell$, we have $\alpha(s_0)=\chi_{\mathrm{cyc}}(s_0)$ and $\theta(s_0)=\chi_{\mathrm{cyc}}(s_0)^{k+1}$, which differ because $\chi_{\mathrm{cyc}}(s_0)$ has infinite order. So the alternative is determined by $x_1$. In the first case every such $x'$ vanishes on the $A_n$ with $n<0$, so none vanishes on all $A_n$ with $n>0$ but not on all $A_n$ with $n<0$; in the second case none vanishes on all $A_n$ with $n<0$ but not on all $A_n$ with $n>0$. This is condition (3).

\emph{Condition (4).} Let $x_1\in Y$ be two-signed and $\delta$ a derivation at $x_1$ with values in $\bar{\mathbf Q}_\ell$ that vanishes on $I$. It suffices to find $b\in\bar{\mathbf Q}_\ell$ with $\delta(f)=bn\,x_1(f)$ for every $n$ and $f\in A_n$: for a second such $\delta'$, with constant $b'$, $\delta$ and $\delta'$ are linearly dependent, as $1\cdot\delta+0\cdot\delta'=0$ if $b=0$, and $b'\delta-b\delta'=0$ otherwise.
\begin{enumerate}
\item[(i)] By Lemma~\ref{sha-counts-vd-points}(a), $x_1=\xi_1\circ q_A$, and $r_{\xi_1}$ is irreducible by Lemma~\ref{sha-counts-vd-irreducible}(d). So for every $(i,j)$ some $r_{\xi_1}(g)_{ij}$ is nonzero; choose $g_1,g_2$ with $\xi_1(u)\ne0$ and $\xi_1(v)\ne0$, where $u=r^{\mathrm u}(g_1)_{12}\in S_{(1,0)}$ and $v=r^{\mathrm u}(g_2)_{21}\in S_{(-1,0)}$. By Lemma~\ref{sha-counts-vd-structure}(b), $uv=q_A(a_0)$ for some $a_0\in A_0$, and $x_1(a_0)\ne0$.
\item[(ii)] The map $X=x_1+\varepsilon\delta:A\to\bar{\mathbf Q}_\ell[\varepsilon]$ is a homomorphism of $\mathcal O$-algebras that vanishes on $I=\ker q_A$, and $c=X(a_0)$ is a unit. For $w\in\mathbf Z^2$ and $f\in S_w$ put $p=\max(-w_1,0)$ and $p'=\max(w_1,0)$; then $fu^pv^{p'}\in S_{(0,w_2)}=q_A(A_{w_2})$, and $\Xi(f)=X(f')c^{-p'}$ for any $f'\in A$ with $q_A(f')=fu^pv^{p'}$ is well defined. The same formula holds with $(p+e,p'+e)$ in place of $(p,p')$ for every $e\ge0$, as $q_A(a_0)=uv$; so $\Xi$, extended additively, is a homomorphism of $\mathcal O$-algebras $S^{\mathrm{gr}}\to\bar{\mathbf Q}_\ell[\varepsilon]$ with $\Xi\circ q_A=X$.
\item[(iii)] Let $\xi'$ be the constant term of $\Xi$, so that $\xi'\circ q_A=x_1$. By Lemma~\ref{sha-counts-vd-points}(d) there is $t$ such that $r=\tau(t)r_{\xi'}\tau(t)^{-1}$ is a $\bar{\mathbf Q}_\ell$-point. Viewing $t$ in $(\bar{\mathbf Q}_\ell[\varepsilon]^\times)^2$, let $\Xi_1=t\cdot\Xi$ (Lemma~\ref{sha-counts-vd-points}(b)). Then $r_{\Xi_1}\equiv r$ modulo $\varepsilon$, and $\Xi_1(q_A(f))=t_2^n\big(x_1(f)+\varepsilon\delta(f)\big)$ for $f\in A_n$, since $q_A(A_n)\subseteq S_{(0,n)}$. The point $r$ is irreducible: the constant term of $\Xi_1$ is $t\cdot\xi'$, whose composite with $q_A$ is $t_2^nx_1$ on $A_n$ and so is two-signed, and Lemma~\ref{sha-counts-vd-irreducible}(d) applies to it.
\item[(iv)] By Lemma~\ref{sha-counts-vd-conjugate}(a), $|r(g)_{ij}-r_0(g)_{ij}|<1$, so $r_{\Xi_1}$ is a $\bar{\mathbf Q}_\ell[\varepsilon]$-point by Lemma~\ref{sha-counts-vd-points}(e). As $D$ is rigid at $r$, Lemma~\ref{sha-counts-vd-conjugate}(c) gives $a,b$ with $r_{\Xi_1}(g)_{ij}=\big(1+\varepsilon(aw_1+bw_2)\big)r(g)_{ij}$.
\item[(v)] Let $s'=(1+\varepsilon a,1+\varepsilon b)$, so that $s'^w=1+\varepsilon(aw_1+bw_2)$, and let $\Xi_2=s'\cdot(t\cdot\xi')$, with $t\cdot\xi'$ viewed with values in $\bar{\mathbf Q}_\ell[\varepsilon]$. By Lemma~\ref{sha-counts-vd-points}(b), $r_{\Xi_2}=r_{\Xi_1}$, so $\Xi_2=\Xi_1$ by Lemma~\ref{sha-counts-vd-points}(c). For $f\in A_n$, $q_A(f)\in S_{(0,n)}$, so $\Xi_2(q_A(f))=(1+\varepsilon bn)t_2^nx_1(f)$. Comparing with (iii), $\delta(f)=bn\,x_1(f)$.
\end{enumerate}
\end{proof}

\begin{lemma}\label{sha-counts-vd-sha}
Let $D$ be a Siegel datum, with the notation of Definition~\ref{sha-siegel-datum}, let $n_S$ be the product of the primes of $S$, and let $V_D=V_{\ell,n_S}(k,\psi)$ (Definition~\ref{sha-character-parts}), with $\mathcal H$ the group whose length it is.
\begin{enumerate}
\item $V_D$ is finite, and there are a nonzero $a\in\mathcal O$ and an injective additive map $h\mapsto z_h$ from $\mathcal H$ to the functions $G\to\mathcal O/a$ whose values satisfy the conditions on $z$ in Definition~\ref{sha-counts-vd-model}(iv).
\item If $D$ has a graded model, then $V_D\le\lambda_{-1}(F)$, the cotangent length of $F$ for $A$ and $x$ (Definition~\ref{sha-algebra-graded}), in the notation of Definition~\ref{sha-counts-vd-model}.
\end{enumerate}
\end{lemma}
\leanref{Development/AutomorphicArithmetic/AutomorphicVarieties/SymplecticGSp4/ShaCoc.lean}{UpperTheory.V.GradeZeroInterpolatingElement.shaCoc}
\leanref{Development/AutomorphicArithmetic/AutomorphicVarieties/SymplecticGSp4/MPrime.lean}{UpperTheory.V.GradeZeroInterpolatingElement.mPrime}
\begin{proof}
(a) The coefficients are $W=\mathcal O\otimes\mathbf Q_\ell/\mathbf Z_\ell(\Psi^{-1})$, since $\Psi^{-1}=\psi^{-1}\chi_{\mathrm{cyc}}^{-k}$. As $\ell\in S$, $\mathcal H$ lies in the group defining $V_\ell(k,\psi)$, which is finite by Lemma~\ref{sha-counts-vd-inputs}(e); so $\mathcal H$ has finite length and is killed by some nonzero $b\in\mathcal O$.
For each $p\in T$ choose $s_p\in D_p$ with $\Psi(s_p)\ne1$: for $p=\ell$ take $s_\ell=s_0$, at which $\Psi(s_0)=\chi_{\mathrm{cyc}}(s_0)^k$ has infinite order; for $p\in S_b$, $\psi(I_p)$ is nontrivial and $\Psi=\psi$ on $I_p$, as $\chi_{\mathrm{cyc}}$ is unramified at $p\ne\ell$. Put $a=b\prod_{p\in T}\big(\Psi(s_p)^{-1}-1\big)\ne0$.
For $h\in\mathcal H$ represented by a cocycle $c$, put $N_h(g)=c(g)-\big(\Psi(g)^{-1}-1\big)\big(\Psi(g_0)^{-1}-1\big)^{-1}c(g_0)$, using that $\Psi(g_0)^{-1}-1$ is a unit since $\bar\Psi(g_0)\ne1$. This is a locally constant cocycle in the class $h$ with $N_h(g_0)=0$; it depends only on $h$, and $h\mapsto N_h$ is $\mathcal O$-linear and injective, since a coboundary $g\mapsto(\Psi(g)^{-1}-1)m$ that vanishes at $g_0$ has $m=0$. Its values are killed by $b$. The elements of $W$ killed by $a$ form a copy of $\mathcal O/a$, so $N_h$ defines $z_h:G\to\mathcal O/a$, additive and injective in $h$, with $z_h(gg')=z_h(g)+\Psi(g)^{-1}z_h(g')$ and $z_h(g_0)=0$.
For $p\in T$, $D_p$ is a decomposition group at a prime dividing $n_S$, so $N_h(g)=(\Psi(g)^{-1}-1)m$ on $D_p$ for some $m\in W$; then $(\Psi(s_p)^{-1}-1)m=N_h(s_p)$ is killed by $b$, so $am=0$, and $m$ gives $m_p\in\mathcal O/a$. On every inertia group the class vanishes, so $N_h(g)=(\Psi(g)^{-1}-1)m'$ there for some $m'$, which is $0$ when $\Psi(g)=1$. Let $g$ lie in an inertia group at a place not above $T$, with $\bar r_0(g)=1$. If the place lies above a prime outside $S$, then $\psi$ and $\chi_{\mathrm{cyc}}$ are unramified there and $\Psi(g)=1$. Otherwise it lies above a prime $p\ne\ell$ of $S$ for which the order of $\psi(I_p)$ is prime to $\ell$; then $\Psi(g)=\psi(g)$ has order prime to $\ell$ and residue $\bar\Psi(g)=1$, so $\Psi(g)=1$. In both cases $z_h(g)=0$.

(b) For $h\in\mathcal H$ let $\delta_h$ be the derivation that Definition~\ref{sha-counts-vd-model}(iv) attaches to $z_h$, and $D_h:A\to\mathcal O/a$ the reduction modulo $a$ of its restriction to $A$.
\begin{enumerate}
\item[(i)] $h\mapsto D_h$ is additive. Let $d=\delta_{h+h'}-\delta_h-\delta_{h'}$, a derivation of $R$ at $x$ that maps $\ker q$ into $a\mathcal O$. For every entry of every $r^{\mathrm u}(g)$ there is $f$ with $q(f)$ equal to that entry and $d(f)\in a\mathcal O$: for the entries $(3,1)$ and $(4,2)$, the only ones of weight $(0,-1)$, by the congruences of Definition~\ref{sha-counts-vd-model}(iv) and the additivity of $h\mapsto z_h$; for the others, the three derivations vanish on an $f\in R_w$, $w\ne(0,-1)$, given by Definition~\ref{sha-counts-vd-model}(iii). The map $f\mapsto x(f)+\varepsilon(d(f)\bmod a)$ is a homomorphism of $\mathcal O$-algebras from $R$ to $\mathcal O\oplus\varepsilon\,\mathcal O/a$, $\varepsilon^2=0$, a complete local noetherian $\mathcal O$-algebra with residue field $\mathbf F$; it kills $\ker q$, so it factors through $S$, and there it agrees with the homomorphism induced by $x$ on the entries of the $r^{\mathrm u}(g)$. By Definition~\ref{sha-counts-vd-model}(ii) the two are equal, so $d(R)\subseteq a\mathcal O$.
\item[(ii)] $h\mapsto D_h$ is injective. If $D_h=0$, then for every $g$, choosing $f\in R_{(0,-1)}=A_{-1}$ with $q(f)=r^{\mathrm u}(g)_{31}$, we get $\theta(g)z_h(g)\equiv\delta_h(f)\equiv0$ modulo $a$, so $z_h=0$ and $h=0$.
\item[(iii)] Each $D_h$ vanishes on the $F_i$, which lie in $\ker q$, and it is the reduction of the derivation $\delta_h|_A$ at $x$ with values in $\mathcal O$, which vanishes on $A_n=R_{(0,n)}$ for $n\ne-1$. Also $x$ vanishes on $A_{-1}$.
\end{enumerate}
The residue field of $\mathcal O$ is finite, so Lemma~\ref{sha-algebra-derivations}, with $j=-1$, $c=a$ and $\mathcal M=\mathcal H$, gives $V_D=\operatorname{len}\mathcal H\le\lambda_{-1}(F)$.
\end{proof}

\subsection{Framed rings and charts}\label{subsection-sha-framed-rings}

\begin{definition}\label{sha-counts-vd-model-framing}
Let $D$ be a Siegel datum, with the notation of Definition~\ref{sha-siegel-datum}, and let $\mathcal W$ be the set of weights $\pm(1,0)$, $\pm(-1,1)$, $\pm(1,1)$, $\pm(0,1)$ and $(0,0)$. Below, $R$ denotes a complete local noetherian $\mathcal O$-algebra with residue field $\mathbf F$.
\begin{enumerate}
\item \emph{Local data} with values in $R$ are continuous lifts $\rho_v:D_v\to M_4(R)$ of $\bar r_0$, for $v\in T$, each symplectic with multiplier $\theta$, such that $\rho_\ell$ is ordinary (Definition~\ref{sha-siegel-datum}(2)) and $\rho_p$ is trivial on $I_p\cap\ker\psi$ for $p\in S_b$.
\item A \emph{local base} of $D$ is a complete local noetherian $\mathcal O$-algebra $B$ with residue field $\mathbf F$, with local data $(\rho^B_v)$ with values in $B$ that are universal (all local data with values in any $R$ are their image under a unique homomorphism of $\mathcal O$-algebras $B\to R$), and with a weighted presentation with origin $x_B$ (Definition~\ref{sha-counts-vd-model}(1)) by variables with weights in $\mathcal W$, such that: the entry $(a,b)$ of every $\rho^B_v(g)$ lies in the part of the weight of $(a,b)$; $x_B$ carries each $\rho^B_v$ to $r_0$; and $B^{\mathrm{gr}}$ has a chart at $x_B$, with the properties of Definition~\ref{sha-counts-vd-model}(i) other than the numbers of variables, with $n_B(j)=d(j)+|S_b|\,d'(j)$ variables of weight $j$. Here $d(j)=2$ for $j\in\{(1,0),(-1,1),(0,1),(1,1),(0,0)\}$, $d(j)=1$ for the other $j\in\mathcal W$, $d'(0,0)=2$, $d'(j)=1$ for the other $j\in\mathcal W$, and $d(j)=d'(j)=0$ for $j\notin\mathcal W$.
\item Let $B$ be a local base. A \emph{framing} with values in $R$ consists of a homomorphism of $\mathcal O$-algebras $i:B\to R$, a continuous lift $r:G\to M_4(R)$ of $\bar r_0$, and matrices $\alpha_v\in M_4(R)$ for $v\in T$, such that: each $r(g)$ is symplectic with multiplier $\theta(g)$; $r(g)=1$ for every $g$ with $\bar r_0(g)=1$ in an inertia group at a place not above $T$; each $\alpha_v$ is symplectic with multiplier $1$ and congruent to $1$ modulo $\mathfrak m_R$; $\alpha_\ell=1$; and $r(g)\alpha_v=\alpha_v\,i(\rho^B_v(g))$ for $v\in T$ and $g\in D_v$.
\item A \emph{framed ring} over $B$ is a framing $(i_U,r_U,\alpha)$ with values in some $U$ that is universal (every framing with values in any $R$ is its image under a unique homomorphism of $\mathcal O$-algebras $U\to R$), together with a weighted presentation of $U$ with origin $x_U$ with the following properties. Its variables are those of $B$, with the same weights and mapping to the images under $i_U$, and further variables $Y_t$, $t\in\tau$; each $Y_t$ has the weight of an index pair $(a,b)$ and maps to the entry $(a,b)$ of $r_U(g)-r_0(g)$ for some $g$ or of $\alpha_v-1$ for some $v$. Its kernel is generated by the image of the kernel of the presentation of $B$ and by power series $f_u$, $u\in\kappa$, of weights $w_f(u)$. The entry $(a,b)$ of every $r_U(g)$ and $\alpha_v$ lies in the part of the weight of $(a,b)$; $x_U\circ i_U=x_B$, and $x_U$ carries $r_U$ to $r_0$ and each $\alpha_v$ to $1$. No $Y_t$ has a weight outside $\{(0,0),\pm(0,1)\}$, and the number of $u$ with $w_f(u)=j$ is the number of $t$ with $Y_t$ of weight $j$ plus $2$ if $j=(0,0)$, plus $1$ if $j=\pm(0,1)$, and plus $0$ otherwise.
\end{enumerate}
\end{definition}
\leanref{Development/AutomorphicArithmetic/AutomorphicVarieties/SymplecticGSp4/ResidueCharMemLevelPrimes.lean}{UpperTheory.V.GradeZeroInterpolatingElement.LocBase}
\leanref{Development/AutomorphicArithmetic/AutomorphicVarieties/SymplecticGSp4/ExDefs.lean}{UpperTheory.V.GradeZeroInterpolatingElement.IsFr}
\leanref{Development/AutomorphicArithmetic/AutomorphicVarieties/SymplecticGSp4/ExDefs.lean}{UpperTheory.V.GradeZeroInterpolatingElement.Fr}

\begin{lemma}\label{sha-counts-vd-model-six}
Let $R$ be a local ring with maximal ideal $\mathfrak m$, let $\nu\in R^\times$, and let $g\in M_4(R)$ be symplectic with multiplier $\nu$ (Definition~\ref{sha-siegel-datum}), with $g_{ij}\in\mathfrak m$ for $i\ne j$ and $g_{12}=g_{14}=g_{32}=g_{21}=g_{41}=g_{23}=0$.
Then $g$ is block diagonal. If moreover $\nu=1$, $g_{13}=g_{31}=0$ and $g_{11}=g_{22}=1$, then $g=1$.
\end{lemma}
\leanref{Development/AutomorphicArithmetic/AutomorphicVarieties/SymplecticGSp4/SixBlock.lean}{UpperTheory.V.GradeZeroInterpolatingElement.sixBlock}
\begin{proof}
The entry $(i,j)$ of $g^*g=\nu$ is $\sum_k\epsilon_i\epsilon_kg_{5-k,5-i}g_{kj}$. With the given zeros, the entries $(1,1)$, $(2,2)$, $(1,2)$ and $(2,1)$ give
\[
 g_{11}g_{44}-g_{24}g_{31}=\nu,\qquad g_{22}g_{33}-g_{13}g_{42}=\nu,\qquad g_{34}g_{22}=0,\qquad g_{43}g_{11}=0.
\]
Since $g_{24},g_{13}\in\mathfrak m$ and $\nu$ is a unit, $g_{11}g_{44}$ and $g_{22}g_{33}$ are units, so $g_{11}$ and $g_{22}$ are units, and $g_{34}=g_{43}=0$. So $g_{ij}=0$ whenever $\lambda_i\ne\lambda_j$.
Under the further hypotheses, the first two identities give $g_{44}=g_{33}=1$, and the entries $(2,4)$ and $(3,1)$ of $g^*g=1$ give $g_{33}g_{24}=0$ and $-g_{42}g_{11}=0$, so $g_{24}=g_{42}=0$ and $g=1$.
\end{proof}

\begin{lemma}\label{sha-counts-vd-model-chart}
Let $\mathcal O$ be a domain with fraction field $E$, let $X$ be a finitely generated free abelian group, and let $R$ be a complete local noetherian $\mathcal O$-algebra with residue field $\mathbf F$, with a weighted presentation with origin $x$ (Definition~\ref{sha-counts-vd-model}(1)). Suppose that $R^{\mathrm{gr}}$ has a chart $\varphi$ at $x$ with the properties of Definition~\ref{sha-counts-vd-model}(i) other than the numbers of variables, with variables $Y_i$ of weights $w(i)\in X$.
Then there is a homomorphism of $\mathcal O$-algebras $\psi$ from $R$ to the power series over $E$ in the same variables that extends $\varphi$ and maps each part $R_j$ into the series of weight $j$, such that the constant term of $\psi(f)$ is $x(f)$, every variable is the only linear term of $\psi(f)$ for some $f$, and $cf\in(\ker x)^2$ for some nonzero $c\in\mathcal O$ whenever $x(f)=0$ and $\psi(f)$ has no linear term.
\end{lemma}
\leanref{Development/AutomorphicArithmetic/AutomorphicVarieties/SymplecticGSp4/SmoothPtPartOfSmoothPtAssociatedGraded.lean}{UpperTheory.smoothPt_part_of_smoothPt_associatedGraded}
\begin{proof}
Let $\mathfrak k=\ker x$ and $\mathfrak k^{\mathrm{gr}}=\mathfrak k\cap R^{\mathrm{gr}}$. By Lemma~\ref{sha-counts-vd-presentations}(f)(i), $\mathfrak k^n\cap R^{\mathrm{gr}}=(\mathfrak k^{\mathrm{gr}})^n$, and every $f\in R$ is congruent modulo $\mathfrak k^n$ to some $y\in R^{\mathrm{gr}}$. Since the series $\varphi(y)$, $y\in\mathfrak k^{\mathrm{gr}}$, have no constant term, $\varphi$ maps $(\mathfrak k^{\mathrm{gr}})^n$ into the series of order at least $n$.
For $f\in R$ and an exponent vector $m$ of degree $d$, let the coefficient of $Y^m$ in $\psi(f)$ be that of $\varphi(y)$ for any $y\in R^{\mathrm{gr}}$ with $f-y\in\mathfrak k^{d+1}$; two such $y$ differ by an element of $(\mathfrak k^{\mathrm{gr}})^{d+1}$, so this does not depend on $y$. If $f-y\in\mathfrak k^n$, the coefficients of $\psi(f)$ and $\varphi(y)$ agree in degrees less than $n$. Comparing $f+f'$ and $ff'$ with the sum and the product of approximations $y,y'$ modulo $\mathfrak k^{d+1}$, where $ff'-yy'=(f-y)f'+y(f'-y')$, shows that $\psi$ is a homomorphism of $\mathcal O$-algebras; it extends $\varphi$.
It maps $R_j\subseteq R^{\mathrm{gr}}$ into the series of weight $j$, since $\varphi$ does. For $f\in R$ choose $y\in R^{\mathrm{gr}}$ with $f-y\in\mathfrak k^2$; then $x(f)=x(y)$, and $\psi(f)$ and $\varphi(y)$ have the same terms of degree at most $1$. This gives the constant term; the variables are linear terms of images of elements of $R^{\mathrm{gr}}$; and if $x(f)=0$ and $\psi(f)$ has no linear term, then $\varphi(y)$ has none and $x(y)=0$, so $cy\in(\mathfrak k^{\mathrm{gr}})^2$ for some nonzero $c$, and $cf=c(f-y)+cy\in\mathfrak k^2$.
\end{proof}

\begin{lemma}\label{sha-counts-vd-model-tangent}
Let $D$ be a Siegel datum, with the notation of Definition~\ref{sha-siegel-datum}, let $B$ be a local base of $D$ and $(U,i_U,r_U,\alpha)$ a framed ring over $B$, with origins $x_B$ and $x_U$ (Definition~\ref{sha-counts-vd-model-framing}). Let $Y\in\mathfrak{sp}_4(\mathcal O)$, and let $A_v\in\mathfrak{sp}_4(\mathcal O)$ for $v\in T$, with $A_\ell=0$.
Then there is a derivation $\delta:U\to\mathcal O$ at $x_U$ (Definition~\ref{sha-algebra-graded}) with $\delta(r_U(g)_{ab})=(Yr_0(g)-r_0(g)Y)_{ab}$ and $\delta((\alpha_v)_{ab})=(A_v)_{ab}$ for all $g$, $v$, $a$ and $b$.
\end{lemma}
\leanref{Development/AutomorphicArithmetic/AutomorphicVarieties/SymplecticGSp4/ExistsPointDerivationFramedCommutatorTangent.lean}{UpperTheory.exists_pointDerivation_framed_commutator_tangent}
\leanref{Development/AutomorphicArithmetic/AutomorphicVarieties/SymplecticGSp4/ExistsLocalFirstOrderLifts.lean}{UpperTheory.exists_local_first_order_lifts}
\begin{proof}
Let $\mathcal O[\varepsilon]$ be the dual numbers, $\varepsilon^2=0$, a complete local noetherian $\mathcal O$-algebra with residue field $\mathbf F$.
\begin{enumerate}
\item[(i)] For $X\in\mathfrak{sp}_4(\mathcal O)$, $u=1+\varepsilon X$ is symplectic with multiplier $1$, since $u^*u=1+\varepsilon(X^*+X)=1$, and $u^{-1}=1-\varepsilon X$. The restrictions of $r_0$ to the $D_v$ are local data with values in $\mathcal O$: $r_0$ is diagonal, the idempotents $e_i$ of $r_0(s_0)$ are $d_iE_{ii}$, and $r_0$ is trivial on $I_p\cap\ker\psi$ for $p\in S_b$ by Lemma~\ref{sha-flagged-residues}(4). Conjugation by $u$ preserves the multiplier, the characteristic polynomial, the identities of ordinarity, in which the $e_i$ are conjugated too, and triviality on a subgroup. So, with $X_v=Y-A_v$, the lifts $\rho'_v(g)=(1+\varepsilon X_v)r_0(g)(1-\varepsilon X_v)=r_0(g)+\varepsilon(X_vr_0(g)-r_0(g)X_v)$ are local data with values in $\mathcal O[\varepsilon]$, and there is $i':B\to\mathcal O[\varepsilon]$ carrying $\rho^B_v$ to $\rho'_v$.
\item[(ii)] Let $r'(g)=(1+\varepsilon Y)r_0(g)(1-\varepsilon Y)$ and $\alpha'_v=1+\varepsilon A_v$. Then $(i',r',\alpha')$ is a framing: $r'$ is a continuous lift of $\bar r_0$, symplectic with multiplier $\theta$ and trivial where $r_0$ is; $\alpha'_v$ is symplectic with multiplier $1$ and congruent to $1$; $\alpha'_\ell=1$; and for $g\in D_v$ both $r'(g)\alpha'_v$ and $\alpha'_v\rho'_v(g)$ equal $r_0(g)+\varepsilon\big(Yr_0(g)-r_0(g)Y+r_0(g)A_v\big)$. By universality there is $y:U\to\mathcal O[\varepsilon]$ with $y\circ i_U=i'$, $y\circ r_U=r'$ and $y(\alpha_v)=\alpha'_v$.
\item[(iii)] Let $y_0$ be $y$ followed by $\varepsilon\mapsto0$. The homomorphisms $y_0\circ i_U$ and $x_B$ both carry $\rho^B_v$ to $r_0$, so they are equal by the universality of $B$; hence $y_0$ and $x_U$ both carry the framing $(i_U,r_U,\alpha)$ to $(x_B,r_0,(1)_v)$, and $y_0=x_U$ by the universality of $U$. So $y=x_U+\varepsilon\delta$ with $\delta$ a derivation at $x_U$, and its values on the entries of $r_U(g)$ and $\alpha_v$ are the coefficients of $\varepsilon$ in those of $r'(g)$ and $\alpha'_v$.
\end{enumerate}
\end{proof}

\begin{lemma}\label{sha-counts-vd-model-cocycle}
Let $D$, $B$ and $(U,i_U,r_U,\alpha)$ be as in Lemma~\ref{sha-counts-vd-model-tangent}. Let $a\in\mathcal O$ be nonzero and $z:G\to\mathcal O/a$ a function with the properties of condition (iv) of Definition~\ref{sha-counts-vd-model}, and put $X_m=E_{31}+E_{42}$.
Then there are a derivation $D_U:U\to\mathcal O/a$ at $x_U$ and a derivation $D_B:B\to\mathcal O$ at $x_B$ of weight $(0,-1)$ (Definition~\ref{sha-algebra-graded}) such that $D_U\circ i_U=D_B\bmod a$, $D_U(r_U(g)_{ij})=(X_mr_0(g))_{ij}\,z(g)$ and $D_U((\alpha_v)_{ij})=0$ for all $g$, $v$, $i$ and $j$.
\end{lemma}
\leanref{Development/AutomorphicArithmetic/AutomorphicVarieties/SymplecticGSp4/ExistsQuotientDerivationOfShaCocycle.lean}{UpperTheory.exists_quotientDerivation_of_shaCocycle}
\leanref{Development/AutomorphicArithmetic/AutomorphicVarieties/SymplecticGSp4/SmoothPtWeightedPresentationOfPointDerivations.lean}{UpperTheory.smoothPt_weightedPresentation_of_pointDerivations}
\begin{proof}
The matrix $X_m$ lies in $\mathfrak{sp}_4(\mathcal O)$, and $r_0(g)X_m=\Psi(g)^{-1}X_mr_0(g)$, since $\delta_3=\Psi^{-1}\delta_1$ and $\delta_4=\Psi^{-1}\delta_2$.
\begin{enumerate}
\item[(i)] For $v\in T$ choose $m_v\in\mathcal O$ with $z(g)\equiv(\Psi(g)^{-1}-1)m_v$ modulo $a$ for $g\in D_v$. As in step (i) of the proof of Lemma~\ref{sha-counts-vd-model-tangent}, with $X_v=-m_vX_m$, the lifts $\rho'_v(g)=r_0(g)+\varepsilon(X_vr_0(g)-r_0(g)X_v)=r_0(g)+\varepsilon(\Psi(g)^{-1}-1)m_vX_mr_0(g)$ are local data with values in $\mathcal O[\varepsilon]$, and there is $y_B:B\to\mathcal O[\varepsilon]$ carrying $\rho^B_v$ to $\rho'_v$. Its constant part carries $\rho^B_v$ to $r_0$, so it is $x_B$ by universality, and $y_B=x_B+\varepsilon D_B$ with $D_B$ a derivation at $x_B$.
\item[(ii)] $D_B$ has weight $(0,-1)$. Its value on the entry $(a',b')$ of $\rho^B_v(g)$ vanishes unless $(a',b')$ has weight $(0,-1)$, since $X_mr_0(g)$ is supported on $(3,1)$ and $(4,2)$. Let $w\ne(0,-1)$. Choose $u\in\mathcal O^\times$ of infinite order, such as $\ell+1$, and a homomorphism $s:\mathbf Z^2\to\mathcal O^\times$, $(w_1,w_2)\mapsto u^{w_1}$ or $u^{w_2}$, with $s(w)\ne s(0,-1)$. There is a homomorphism of $\mathcal O$-algebras $y_s:B\to\mathcal O[\varepsilon]$ equal to $s(w')y_B$ on each part $B_{w'}$: on power series it sends each variable $X_t$ to $s(w_t)y_B(z_t)$, where $z_t\in\mathfrak m_B$ is the image of $X_t$, and it kills the kernel of the presentation, whose homogeneous components lie in the kernel. Let $\phi:\mathcal O[\varepsilon]\to\mathcal O[\varepsilon]$ be $\varepsilon\mapsto s(0,-1)^{-1}\varepsilon$. On an entry $f\in B_{w'}$ of some $\rho^B_v(g)$, $\phi\circ y_s$ and $y_B$ agree: $x_B(f)=0$ unless $w'=0$, and $D_B(f)=0$ unless $w'=(0,-1)$. So $\phi\circ y_s=y_B$ by the universality of $B$, and for $f\in B_w$ this gives $s(0,-1)^{-1}s(w)D_B(f)=D_B(f)$, so $D_B(f)=0$.
\item[(iii)] Let $Q=\mathcal O/a$ and $\mathcal O\oplus\varepsilon Q$ the ring with $\varepsilon^2=0$, and let $r(g)=r_0(g)+\varepsilon z(g)X_mr_0(g)$. By the cocycle relation and $r_0(g)X_m=\Psi(g)^{-1}X_mr_0(g)$, $r$ is a homomorphism; it is a continuous lift of $\bar r_0$, as $z$ is locally constant; it is symplectic with multiplier $\theta$, as $X_m\in\mathfrak{sp}_4(\mathcal O)$; and $r(g)=1$ for $g$ with $\bar r_0(g)=1$ in an inertia group at a place not above $T$, where $z(g)=0$ and $r_0(g)=1$, each $\delta_i(g)$ being a root of unity of order prime to $\ell$ with residue $1$. Let $p:\mathcal O[\varepsilon]\to\mathcal O\oplus\varepsilon Q$ reduce the coefficient of $\varepsilon$ modulo $a$. Then $(p\circ y_B,r,(1)_v)$ is a framing: for $g\in D_v$, $p(\rho'_v(g))=r(g)$ by the choice of $m_v$.
\item[(iv)] By universality there is $y:U\to\mathcal O\oplus\varepsilon Q$ with $y\circ i_U=p\circ y_B$, $y\circ r_U=r$ and $y(\alpha_v)=1$. As in step (iii) of the proof of Lemma~\ref{sha-counts-vd-model-tangent}, its constant part is $x_U$, so $y=x_U+\varepsilon D_U$ with $D_U$ a derivation at $x_U$ with values in $Q$, and $D_U$ has the stated values.
\end{enumerate}
\end{proof}

\begin{definition}\label{sha-counts-vd-model-charted}
Let $\mathcal O$ be a domain with fraction field $E$, let $X$ be an abelian group, let $A$ be an $\mathcal O$-algebra with $\mathcal O$-submodules $A_j$, $j\in X$, its \emph{parts}, and let $x:A\to\mathcal O$ be a homomorphism of $\mathcal O$-algebras. Let $\Sigma_x$ be the set of the $f\in A$ with $cf\in(\ker x)^2$ for some nonzero $c\in\mathcal O$.
A \emph{chart} of $A$ at $x$ is a homomorphism of $\mathcal O$-algebras $\varphi:A\to E[[Y_i:i\in\tau]]$, where $\tau$ is a finite set and $Y_i$ has a weight $w(i)\in X$, with the following properties.
\begin{enumerate}
\item $\varphi$ maps $A_j$ into the series of weight $j$, for every $j$.
\item The constant term of $\varphi(f)$ is $x(f)$, for every $f\in A$.
\item For every $i\in\tau$ there is $f\in A$ such that the only linear term of $\varphi(f)$ is a nonzero multiple of $Y_i$.
\item If $x(f)=0$ and $\varphi(f)$ has no linear term, then $f\in\Sigma_x$.
\end{enumerate}
The chart has \emph{$n(j)$ variables of weight $j$} if $n(j)$ of the $Y_i$ have weight $j$, for every $j$. For a ring $R$ with a weighted presentation (Definition~\ref{sha-counts-vd-model}(1)), the parts of $R$ and of $R^{\mathrm{gr}}$ are the $R_j$; so condition (i) of Definition~\ref{sha-counts-vd-model} asks for a chart of $R^{\mathrm{gr}}$ at $x$.
\end{definition}
\leanref{Development/AutomorphicArithmetic/AutomorphicVarieties/CommutativeAlgebra/OrbitVanishingLocal.lean}{UpperTheory.V.Alg.IsChart}
\leanref{Development/AutomorphicArithmetic/AutomorphicVarieties/CommutativeAlgebra/OrbitVanishingLocal.lean}{UpperTheory.V.Alg.SmoothPt}
\leanref{Development/AutomorphicArithmetic/AutomorphicVarieties/CommutativeAlgebra/ResMapBijective.lean}{UpperTheory.V.Alg.satSq}

\begin{lemma}\label{sha-counts-vd-model-solve}
Let $E$ be a field, let $X$ be an abelian group, and let $\tau$ and $\kappa$ be finite sets. Let $Y_i$, $i\in\tau$, be variables with weights $w(i)\in X$, and for $f\in E[[Y_i:i\in\tau]]$ and $i\in\tau$ let $f_i$ be the coefficient of $Y_i$ in $f$. For $k\in\kappa$ let $z_k\in E[[Y_i:i\in\tau]]$ be a series of weight $c_k\in X$ with constant term $0$, and suppose that the vectors $(z_{k,i})_{i\in\tau}$, $k\in\kappa$, are linearly independent over $E$.
Then there are a finite set $\tau'$, an injective map $t:\tau'\to\tau$ and a homomorphism of $E$-algebras $\psi:E[[Y_i:i\in\tau]]\to E[[Y'_{i'}:i'\in\tau']]$, where $Y'_{i'}$ has the weight $w(t(i'))$, with the following properties.
\begin{enumerate}
\item $\psi$ maps the series of weight $j$ to series of weight $j$, for every $j$; it preserves constant terms; and $\psi(z_k)=0$ for every $k$.
\item Let $i\in\tau'$, and let $f$ have $f_{i''}=0$ for every $i''\ne t(i)$. Then $\psi(f)_{i'}=f_{t(i')}$ for every $i'\in\tau'$.
\item If $\psi(f)$ has no linear term, then $(f_i)_{i\in\tau}$ is a linear combination over $E$ of the vectors $(z_{k,i})_{i\in\tau}$.
\item For every $j\in X$, $\#\{i'\in\tau':w(t(i'))=j\}+\#\{k\in\kappa:c_k=j\}=\#\{i\in\tau:w(i)=j\}$.
\end{enumerate}
\end{lemma}
\leanref{Development/AutomorphicArithmetic/AutomorphicVarieties/CommutativeAlgebra/ChartCut.lean}{cc_series}
\leanref{Development/AutomorphicArithmetic/AutomorphicVarieties/CommutativeAlgebra/ChartCut.lean}{cc_pivot}
\leanref{Development/AutomorphicArithmetic/AutomorphicVarieties/CommutativeAlgebra/ChartCut.lean}{cc_ift}
\leanref{Development/AutomorphicArithmetic/AutomorphicVarieties/CommutativeAlgebra/ChartCut.lean}{cc_fix}
\begin{proof}
\begin{enumerate}
\item[(i)] \emph{Pivots.} Let $l$ be the $\kappa\times\tau$ matrix $(z_{k,i})$. Since $z_k$ has weight $c_k$, $z_{k,i}=0$ unless $w(i)=c_k$. The rows of $l$ are linearly independent, so some $|\kappa|$ of its columns form an invertible matrix, and some term $\pm\prod_kz_{k,s(k)}$ of the expansion of its determinant is nonzero, for a bijection $s$ of $\kappa$ onto these columns. So $s:\kappa\to\tau$ is injective, $M=(z_{k,s(k')})_{k,k'}$ is invertible, and $z_{k,s(k)}\ne0$ for every $k$. In particular $w(s(k))=c_k$, so $M_{kk'}=0$ unless $c_k=c_{k'}$, and the inverse $N$ of $M$ has the same property.
Put $z'_k=\sum_{k'}N_{kk'}z_{k'}$. It has weight $c_k$ and constant term $0$, and $z'_{k,s(k')}$ is $1$ if $k=k'$ and $0$ otherwise. Let $\tau'$ be the complement of $s(\kappa)$ in $\tau$ and $t$ its inclusion. Since $w(s(k))=c_k$, this gives (4).
\item[(ii)] \emph{Solving.} Put $b_{ki'}=z'_{k,t(i')}$, so that $z'_k=Y_{s(k)}+\sum_{i'}b_{ki'}Y_{t(i')}+h_k$ with $h_k$ of order at least $2$. Then $b_{ki'}=0$ unless $w(t(i'))=c_k$, and $h_k$ has weight $c_k$.
For a family $u=(u_k)_{k\in\kappa}$ of series in the $Y'_{i'}$ with constant term $0$, let $\psi_u$ be the homomorphism of $E$-algebras that substitutes $u_k$ for $Y_{s(k)}$ and $Y'_{i'}$ for $Y_{t(i')}$. Put $u^{(0)}=0$ and $u^{(n+1)}_k=-\sum_{i'}b_{ki'}Y'_{i'}-\psi_{u^{(n)}}(h_k)$. Since $h_k$ has order at least $2$, the terms of degree at most $d$ of $\psi_u(h_k)$ depend only on the terms of degree less than $d$ of the $u_k$. So, by induction on $d$, the terms of degree at most $d$ of $u^{(n)}_k$ do not depend on $n\ge d+1$. Let $u_k$ be the series whose terms of degree $d$ are those of $u^{(d+1)}_k$; then $u_k=-\sum_{i'}b_{ki'}Y'_{i'}-\psi_u(h_k)$.
A substitution of series of the weights of the variables maps series of weight $j$ to series of weight $j$. So, by induction on $n$, each $u^{(n)}_k$ has weight $c_k$, and so has $u_k$.
Let $\psi=\psi_u$. It maps series of weight $j$ to series of weight $j$, and it preserves constant terms, since the $u_k$ have constant term $0$. It kills each $z'_k$, as $\psi(z'_k)=u_k+\sum_{i'}b_{ki'}Y'_{i'}+\psi(h_k)=0$. Since $z_k=\sum_{k'}M_{kk'}z'_{k'}$, it kills each $z_k$. This proves (1).
\item[(iii)] \emph{Linear terms.} As $\psi(h_k)$ has order at least $2$, the linear part of $u_k$ is $-\sum_{i'}b_{ki'}Y'_{i'}$, and so, for every $f$ and every $i'\in\tau'$,
\[
 \psi(f)_{i'}=f_{t(i')}-\sum_kf_{s(k)}b_{ki'}.
\]
If $f_{i''}=0$ for every $i''\ne t(i)$, then $f_{s(k)}=0$ for every $k$, and (2) follows. If $\psi(f)$ has no linear term, then $f_{t(i')}=\sum_kf_{s(k)}z'_{k,t(i')}$ for every $i'$, and $f_{s(k')}=\sum_kf_{s(k)}z'_{k,s(k')}$ for every $k'$. So $(f_i)_i=\sum_kf_{s(k)}(z'_{k,i})_i=\sum_{k,k'}f_{s(k)}N_{kk'}(z_{k',i})_i$, which proves (3).
\end{enumerate}
\end{proof}

\begin{lemma}\label{sha-counts-vd-model-cut}
Let $\mathcal O$ be a domain with fraction field $E$, let $X$ be an abelian group, and let $A$ be an $\mathcal O$-algebra with parts $A_j$, $j\in X$, a homomorphism of $\mathcal O$-algebras $x:A\to\mathcal O$, and a chart $\varphi$ at $x$ with $n(j)$ variables of weight $j$ (Definition~\ref{sha-counts-vd-model-charted}). Let $\kappa$ be a finite set, and for $k\in\kappa$ let $g_k\in A_{c_k}$ with $x(g_k)=0$, such that the images of the $g_k$ in $A/\Sigma_x$ are linearly independent over $\mathcal O$. Let $p:A\to A'$ be a surjective homomorphism of $\mathcal O$-algebras whose kernel is the ideal generated by the $g_k$, and let $x':A'\to\mathcal O$ be a homomorphism of $\mathcal O$-algebras with $x'\circ p=x$.
Then $A'$, with the parts $p(A_j)$, has a chart at $x'$ with $n'(j)$ variables of weight $j$, where $n'(j)+\#\{k:c_k=j\}=n(j)$ for every $j$.
\end{lemma}
\leanref{Development/AutomorphicArithmetic/AutomorphicVarieties/CommutativeAlgebra/ChartCut.lean}{UpperTheory.V.Alg.chart_cut}
\begin{proof}
For a series $h$ in the variables $Y_i$ of $\varphi$ let $h_i$ be its coefficient of $Y_i$.
\begin{enumerate}
\item[(i)] The vectors $(\varphi(g_k)_i)_i$ are linearly independent over $E$. Let $\beta_k\in E$ with $\sum_k\beta_k\varphi(g_k)_i=0$ for every $i$, and choose a nonzero $d\in\mathcal O$ with $a_k=d\beta_k\in\mathcal O$ for every $k$. Then $f=\sum_ka_kg_k$ has $x(f)=0$, and $\varphi(f)$ has no linear term, so $f\in\Sigma_x$ by property (4) of $\varphi$. By hypothesis every $a_k$ vanishes, and so does every $\beta_k$.
\item[(ii)] The series $\varphi(g_k)$ has weight $c_k$ and constant term $x(g_k)=0$. Lemma~\ref{sha-counts-vd-model-solve}, applied to these series, gives $\tau'$, $t$ and $\psi$; put $n'(j)=\#\{i':w(t(i'))=j\}$, which satisfies the count by part (4) of that lemma.
\item[(iii)] The homomorphism $\psi\circ\varphi$ kills the $g_k$, so it kills $\ker p$, and it is $\varphi'\circ p$ for a homomorphism of $\mathcal O$-algebras $\varphi':A'\to E[[Y'_{i'}:i'\in\tau']]$.
\item[(iv)] $\varphi'$ is a chart of $A'$ at $x'$. It maps $p(A_j)$ into the series of weight $j$, by property (1) of $\varphi$ and part (1) of Lemma~\ref{sha-counts-vd-model-solve}. The constant term of $\varphi'(p(f))=\psi(\varphi(f))$ is that of $\varphi(f)$, which is $x(f)=x'(p(f))$. For $i'\in\tau'$ choose $f\in A$ such that the only linear term of $\varphi(f)$ is a nonzero multiple $cY_{t(i')}$; by part (2) of that lemma, and as $t$ is injective, the only linear term of $\varphi'(p(f))$ is $cY'_{i'}$.
Finally let $f\in A$ be such that $x'(p(f))=0$ and $\varphi'(p(f))$ has no linear term. By part (3), $\varphi(f)_i=\sum_k\beta_k\varphi(g_k)_i$ for every $i$, for some $\beta_k\in E$; choose a nonzero $d\in\mathcal O$ with $a_k=d\beta_k\in\mathcal O$. Then $x(df-\sum_ka_kg_k)=0$ and $\varphi(df-\sum_ka_kg_k)$ has no linear term, so $e(df-\sum_ka_kg_k)\in(\ker x)^2$ for some nonzero $e\in\mathcal O$. Applying $p$, which kills the $g_k$ and maps $\ker x$ into $\ker x'$, gives $ed\,p(f)\in(\ker x')^2$.
\end{enumerate}
\end{proof}

\begin{lemma}\label{sha-counts-vd-model-kernel}
Let $\mathcal O$ be a local domain with fraction field $E$ and residue field $\mathbf F$, let $X$ be a finitely generated free abelian group, and let $R$ be a complete local noetherian $\mathcal O$-algebra with residue field $\mathbf F$, with a weighted presentation $\pi$ with origin $x$ (Definition~\ref{sha-counts-vd-model}(1)) and a chart $\varphi$ of $R$ at $x$ (Definition~\ref{sha-counts-vd-model-charted}).
\begin{enumerate}
\item $\ker\varphi$ is a prime ideal contained in $\ker x$.
\item If $\varphi(\pi(f))=0$, then $\varphi(\pi(f_j))=0$ for every $j\in X$, where $f_j$ is the sum of the terms of $f$ of weight $j$.
\item Every derivation of $R$ at $x$ with values in $\mathcal O$ (Definition~\ref{sha-algebra-graded}) vanishes on $\ker\varphi$.
\item Let $p:R\to R'$ be a surjective homomorphism of $\mathcal O$-algebras, and let $x':R'\to\mathcal O$ and $\varphi':R'\to E[[Y_i]]$ be homomorphisms of $\mathcal O$-algebras with $x'\circ p=x$ and $\varphi'\circ p=\varphi$. Then $\varphi'$ is a chart of $R'$, with the parts $p(R_j)$, at $x'$.
\end{enumerate}
\end{lemma}
\leanref{Development/AutomorphicArithmetic/AutomorphicVarieties/CommutativeAlgebra/ChartKer.lean}{UpperTheory.V.Alg.chart_ker}
\leanref{Development/AutomorphicArithmetic/AutomorphicVarieties/CommutativeAlgebra/ChartKer.lean}{homog_ker}
\begin{proof}
(a) The ring $E[[Y_i]]$ is a domain, so $\ker\varphi$ is prime. If $\varphi(f)=0$, then the constant term $x(f)$ of $\varphi(f)$ vanishes.

(b) Let $z_s=\pi(X_s)$. The constant term of $\varphi(z_s)$ is $x(z_s)=0$, so $\varphi\circ\pi$ maps the ideal generated by the monomials of degree $N$ in the $X_s$ into the series of order at least $N$, for every $N$. Fix $j$ and a monomial $Y^n$ of degree $d$. If $Y^n$ has a weight other than $j$, its coefficient in $\varphi(\pi(f_j))$, a series of weight $j$, is $0$. Otherwise let $T\subseteq X$ be the finite set consisting of $j$ and of the weights of the monomials of degree at most $d$ in the $X_s$. The series $f-\sum_{i\in T}f_i$ has no terms of degree at most $d$, so it lies in the ideal generated by the monomials of degree $d+1$, and the coefficient of $Y^n$ in its image under $\varphi\circ\pi$ is $0$. As $\varphi(\pi(f))=0$, the coefficient of $Y^n$ in $\sum_{i\in T}\varphi(\pi(f_i))$ is $0$. The terms with $i\ne j$ contribute $0$, being series of weight $i$; so the coefficient of $Y^n$ in $\varphi(\pi(f_j))$ is $0$.

(c) Let $D$ be such a derivation, and let $\varphi(f)=0$. Then $x(f)=0$ by (a), and $\varphi(f)$ has no linear term, so $cf\in(\ker x)^2$ for some nonzero $c\in\mathcal O$. By the product rule $D$ vanishes on $(\ker x)^2$, so $cD(f)=0$, and $D(f)=0$.

(d) For $f\in R$, $\varphi'(p(f))=\varphi(f)$ and $x'(p(f))=x(f)$, so properties (1) to (3) pass from $\varphi$ to $\varphi'$. If $x'(p(f))=0$ and $\varphi'(p(f))$ has no linear term, then $cf\in(\ker x)^2$ for some nonzero $c$, and $c\,p(f)\in(\ker x')^2$, as $p$ maps $\ker x$ into $\ker x'$.
\end{proof}

\begin{lemma}\label{sha-counts-vd-model-graded}
Let $\mathcal O$ be a local domain with residue field $\mathbf F$, let $X$ be a finitely generated free abelian group, and let $R$ be a complete local noetherian $\mathcal O$-algebra with residue field $\mathbf F$, with a weighted presentation with origin $x$ and a chart $\varphi$ of $R$ at $x$ with $n(j)$ variables of weight $j$ (Definition~\ref{sha-counts-vd-model-charted}). Then the restriction of $\varphi$ to $R^{\mathrm{gr}}$ is a chart of $R^{\mathrm{gr}}$ at $x$ with $n(j)$ variables of weight $j$.
\end{lemma}
\leanref{Development/AutomorphicArithmetic/AutomorphicVarieties/CommutativeAlgebra/ChartGr.lean}{UpperTheory.V.Alg.chart_gr}
\begin{proof}
Let $\mathfrak k=\ker x$ and $\mathfrak k^{\mathrm{gr}}=\mathfrak k\cap R^{\mathrm{gr}}$. By Lemma~\ref{sha-counts-vd-presentations}(f)(i), $\mathfrak k^2\cap R^{\mathrm{gr}}=(\mathfrak k^{\mathrm{gr}})^2$, and every $g\in R$ is congruent modulo $\mathfrak k^2$ to some $y\in R^{\mathrm{gr}}$. By property (2) of $\varphi$, the series $\varphi(\mathfrak k)$ have constant term $0$, so the series $\varphi(\mathfrak k^2)$ have no constant and no linear terms, and $\varphi(g)$ and $\varphi(y)$ have the same linear terms.
Properties (1) and (2) of a chart pass to the restriction. For (3), let $g\in R$ be such that the only linear term of $\varphi(g)$ is a nonzero multiple of $Y_i$, and take $y$ as above; then the only linear term of $\varphi(y)$ is the same. For (4), if $f\in R^{\mathrm{gr}}$, $x(f)=0$ and $\varphi(f)$ has no linear term, then $cf\in\mathfrak k^2\cap R^{\mathrm{gr}}=(\mathfrak k^{\mathrm{gr}})^2$ for some nonzero $c\in\mathcal O$.
\end{proof}

\begin{lemma}\label{sha-counts-vd-model-quotient}
Let $\mathcal O$ be a complete local noetherian domain with residue field $\mathbf F$, and let $X$ be a finitely generated free abelian group. Let $P$ be a complete local noetherian $\mathcal O$-algebra with residue field $\mathbf F$, with a weighted presentation with origin $x$ (Definition~\ref{sha-counts-vd-model}(1)) and a chart of $P$ at $x$ with $n(j)$ variables of weight $j$ (Definition~\ref{sha-counts-vd-model-charted}). Let $\kappa$ be a finite set, and for $k\in\kappa$ let $g_k\in P_{w_k}$ with $x(g_k)=0$. Suppose that for every $k$ there is a derivation $\delta_k:P\to\mathcal O$ at $x$ (Definition~\ref{sha-algebra-graded}) with $\delta_k(g_k)\ne0$ and $\delta_k(g_{k'})=0$ for $k'\ne k$.
Then there are a complete local noetherian $\mathcal O$-algebra $R$ with residue field $\mathbf F$ that is a domain, a surjective homomorphism of $\mathcal O$-algebras $q:P\to R$ with $q(g_k)=0$ for every $k$, and a weighted presentation of $R$ by the variables of the presentation of $P$, with their weights, with parts $R_j=q(P_j)$ and with an origin $x_R$ such that $x_R\circ q=x$, with the following properties: every derivation of $P$ at $x$ with values in $\mathcal O$ that vanishes at the $g_k$ vanishes on $\ker q$; and $R^{\mathrm{gr}}$ has a chart at $x_R$ with $n'(j)$ variables of weight $j$, where $n'(j)+\#\{k:w_k=j\}=n(j)$ for every $j$.
\end{lemma}
\leanref{Development/AutomorphicArithmetic/AutomorphicVarieties/SymplecticGSp4/ExistsSmoothQuotientByWeightedEquations.lean}{UpperTheory.exists_smooth_quotient_by_weighted_equations}
\leanref{Development/AutomorphicArithmetic/AutomorphicVarieties/SymplecticGSp4/ExistsSmoothQuotientByWeightedEquations.lean}{indep}
\leanref{Development/AutomorphicArithmetic/AutomorphicVarieties/SymplecticGSp4/ExistsSmoothQuotientByWeightedEquations.lean}{der_desc}
\begin{proof}
Let $I$ be the ideal of $P$ generated by the $g_k$, $p:P\to P/I$ the quotient map, and $x':P/I\to\mathcal O$ the homomorphism induced by $x$, which kills the $g_k$.
\begin{enumerate}
\item[(i)] The images of the $g_k$ in $P/\Sigma_x$ are linearly independent over $\mathcal O$. By the product rule a derivation at $x$ vanishes on $(\ker x)^2$. If $c\sum_ka_kg_k\in(\ker x)^2$ with $a_k\in\mathcal O$ and $c\ne0$, applying $\delta_k$ gives $ca_k\delta_k(g_k)=0$, so $a_k=0$.
\item[(ii)] By Lemma~\ref{sha-counts-vd-model-cut}, $P/I$, with the parts $p(P_j)$, has a chart $\varphi'$ at $x'$ with $n'(j)$ variables of weight $j$, where $n'(j)+\#\{k:w_k=j\}=n(j)$. By Lemma~\ref{sha-counts-vd-presentations}(a), the presentation of $P$ followed by $p$ is a weighted presentation of $P/I$, with parts $p(P_j)$; it has the origin $x'$.
\item[(iii)] Let $J=\ker\varphi'$, $R=(P/I)/J$, and $q:P\to R$ the composite of the quotient maps. By Lemma~\ref{sha-counts-vd-model-kernel}(a), $J$ is prime and lies in $\ker x'$, so $R$ is a domain, and $x'$ and $\varphi'$ induce homomorphisms $x_R$ and $\varphi_R$ on $R$, with $x_R\circ q=x$. By Lemma~\ref{sha-counts-vd-model-kernel}(b), the presentation of $P/I$ followed by the quotient map is a weighted presentation of $R$; its parts are the $q(P_j)$, and it has the origin $x_R$. As a quotient of $P$ by a proper ideal, $R$ is a complete local noetherian $\mathcal O$-algebra with residue field $\mathbf F$.
\item[(iv)] By Lemma~\ref{sha-counts-vd-model-kernel}(d), $\varphi_R$ is a chart of $R$ at $x_R$, with $n'(j)$ variables of weight $j$, and by Lemma~\ref{sha-counts-vd-model-graded} its restriction to $R^{\mathrm{gr}}$ is a chart of $R^{\mathrm{gr}}$ at $x_R$.
\item[(v)] Let $\delta$ be a derivation of $P$ at $x$ with values in $\mathcal O$ that vanishes at the $g_k$. It vanishes on $I$, since $\delta(ag_k)=x(a)\delta(g_k)+x(g_k)\delta(a)=0$, so $\delta=\delta'\circ p$ for a derivation $\delta'$ of $P/I$ at $x'$. By Lemma~\ref{sha-counts-vd-model-kernel}(c), $\delta'$ vanishes on $J$; so $\delta$ vanishes on $\ker q=p^{-1}(J)$.
\end{enumerate}
\end{proof}

\begin{lemma}\label{sha-counts-vd-model-sum}
Let $\mathcal O$ be a local domain with fraction field $E$ and residue field $\mathbf F$, let $X$ be a finitely generated free abelian group, and let $T$ be a finite set. For $v\in T$ let $B_v$ be a complete local noetherian $\mathcal O$-algebra with residue field $\mathbf F$, with a weighted presentation $\pi_v$ by variables $X_{v,s}$, $s\in\Sigma_v$, with origin $x_v$ (Definition~\ref{sha-counts-vd-model}(1)), and with a chart $\varphi_v$ of $B_v$ at $x_v$ in variables $Y_{v,i}$, $i\in\tau_v$ (Definition~\ref{sha-counts-vd-model-charted}). Let $R$ be a complete local noetherian $\mathcal O$-algebra with residue field $\mathbf F$, with a weighted presentation $\pi$ by all the variables $X_{v,s}$, with the same weights, whose kernel is generated by the kernels of the $\pi_v$, viewed in $\mathcal O[[X_{v,s}:v\in T,s\in\Sigma_v]]$.
Then $\pi$ has an origin $x$, and $R$ has a chart $\psi$ at $x$ in the variables $Y_{v,i}$, $v\in T$, $i\in\tau_v$, with their weights, such that $\psi(\pi(X_{v,s}))=\varphi_v(\pi_v(X_{v,s}))$ for all $v$ and $s$.
\end{lemma}
\leanref{Development/AutomorphicArithmetic/AutomorphicVarieties/CommutativeAlgebra/ChartSum.lean}{UpperTheory.V.Alg.chart_sum}
\begin{proof}
Let $z_{v,s}=\pi_v(X_{v,s})$, and let $i_v:B_v\to R$ be the homomorphism induced by the inclusion of $\mathcal O[[X_{v,s}:s\in\Sigma_v]]$ in $\mathcal O[[X_{v,s}:v,s]]$, which maps $\ker\pi_v$ into $\ker\pi$.
\begin{enumerate}
\item[(i)] \emph{Origin.} The series of $\ker\pi$ lie in the ideal generated by the series of the $\ker\pi_v$, which have constant term $0$, so they have constant term $0$. Both $x\circ i_v$ and $x_v$ are given by constant terms, so they are equal; hence $i_v$ maps $\ker x_v$ into $\ker x$, and $\Sigma_{x_v}$ into $\Sigma_x$.
\item[(ii)] \emph{The homomorphism.} The series $a_{v,s}=\varphi_v(z_{v,s})$ has constant term $x_v(z_{v,s})=0$, and the weight of $X_{v,s}$. Let $\Phi:\mathcal O[[X_{v,s}]]\to E[[Y_{v,i}]]$ be the homomorphism of $\mathcal O$-algebras that substitutes $a_{v,s}$ for $X_{v,s}$; it maps series of weight $j$ to series of weight $j$.
On $\mathcal O[[X_{v,s}:s\in\Sigma_v]]$, $\Phi$ agrees with $\varphi_v\circ\pi_v$. Indeed both are homomorphisms of $\mathcal O$-algebras that agree on the variables and map them to series with constant term $0$; every series is a polynomial plus an element of the ideal generated by the monomials of degree $N$, which both map into the series of order at least $N$; so the two images of a series agree in degrees less than $N$, for every $N$.
Hence $\Phi$ kills $\ker\pi$, and $\Phi=\psi\circ\pi$ for a homomorphism of $\mathcal O$-algebras $\psi:R\to E[[Y_{v,i}]]$; then $\psi\circ i_v=\varphi_v$, and in particular $\psi(\pi(X_{v,s}))=\varphi_v(z_{v,s})$.
\item[(iii)] \emph{The chart.} Property (1) holds since $\Phi$ maps series of weight $j$ to series of weight $j$. For (2), $\Phi(g)-g(0)$ lies in the ideal generated by the $a_{v,s}$, so the constant term of $\psi(\pi(g))$ is $g(0)=x(\pi(g))$. For (3), given $v$ and $i$, choose $f\in B_v$ such that the only linear term of $\varphi_v(f)$ is a nonzero multiple of $Y_{v,i}$; then $\psi(i_v(f))=\varphi_v(f)$.
For (4), let $f=\pi(g)$ with $x(f)=g(0)=0$ and $\psi(f)$ without linear term, and let $\ell=\sum_{v,s}g_{v,s}X_{v,s}$ be the linear part of $g$. Then $g-\ell$ lies in the ideal generated by the products of two variables, so $\pi(g-\ell)\in(\ker x)^2$, and $\Phi(g-\ell)$ has no linear term; hence $\Phi(\ell)$ has no linear term. Write $\ell=\sum_v\ell_v$ with $\ell_v=\sum_sg_{v,s}X_{v,s}$. The series $\Phi(\ell_v)=\varphi_v(\pi_v(\ell_v))$ involves only the variables $Y_{v,i}$, so it has no linear term, for every $v$; and its constant term is $0$. By property (4) of $\varphi_v$, $\pi_v(\ell_v)\in\Sigma_{x_v}$, so $\pi(\ell_v)=i_v(\pi_v(\ell_v))\in\Sigma_x$. Hence $f=\sum_v\pi(\ell_v)+\pi(g-\ell)\in\Sigma_x$.
\end{enumerate}
\end{proof}

\begin{lemma}\label{sha-counts-vd-model-cover}
Let $\mathcal O$ be a complete local noetherian domain with fraction field $E$ and residue field $\mathbf F$, and let $X$ be a finitely generated free abelian group. Let $B$ and $U$ be complete local noetherian $\mathcal O$-algebras with residue field $\mathbf F$, and $i:B\to U$ a homomorphism of $\mathcal O$-algebras. Let $\pi_B$ be a weighted presentation of $B$ by variables $X_s$, $s\in\Sigma$, with origin $x_B$ (Definition~\ref{sha-counts-vd-model}(1)), such that $B$ has a chart at $x_B$ with $n(j)$ variables of weight $j$ (Definition~\ref{sha-counts-vd-model-charted}). Let $\pi_U$ be a weighted presentation of $U$ with origin $x_U$, by the variables $X_s$, with the same weights, and further variables $Y_t$, $t\in\tau$, such that $\pi_U$ restricts to $i\circ\pi_B$ on $\mathcal O[[X_s:s\in\Sigma]]$, and $\ker\pi_U$ is generated by $\ker\pi_B$ and by series $f_u$, $u\in\kappa$, with $\kappa$ finite and $f_u$ of weight $w_f(u)$. Let $z_t=\pi_U(Y_t)$, and let $a(j)$ be the number of $t$ with $Y_t$ of weight $j$.
Then there are a complete local noetherian $\mathcal O$-algebra $P$ with residue field $\mathbf F$, with a weighted presentation by the variables $X_s$ and $Y_t$, with their weights, and with an origin $x_P$; homomorphisms of $\mathcal O$-algebras $i_P:B\to P$ and $p:P\to U$; and elements $g_u\in P_{w_f(u)}$, $u\in\kappa$; with the following properties.
\begin{enumerate}
\item $p$ is surjective with kernel generated by the $g_u$, $p\circ i_P=i$, $x_U\circ p=x_P$, and $U_j=p(P_j)$ for every $j$.
\item $P$ has a chart at $x_P$ with $n(j)+a(j)$ variables of weight $j$.
\item Let $a\in\mathcal O$ be nonzero and $j\in X$. Let $D_U:U\to\mathcal O/a$ be a derivation at $x_U$, and $D_B:B\to\mathcal O$ a derivation at $x_B$ of weight $j$ (Definition~\ref{sha-algebra-graded}), such that $D_U\circ i=D_B\bmod a$ and $D_U(z_t)=0$ for every $t$ with $Y_t$ of weight other than $j$. Then there is a derivation $\delta:P\to\mathcal O$ at $x_P$ of weight $j$ with $\delta\bmod a=D_U\circ p$.
\end{enumerate}
\end{lemma}
\leanref{Development/AutomorphicArithmetic/AutomorphicVarieties/SymplecticGSp4/NonemptyCoverOfRelativePresentationSmoothPt.lean}{UpperTheory.nonempty_cover_of_relative_presentation_smoothPt}
\leanref{Development/AutomorphicArithmetic/AutomorphicVarieties/SymplecticGSp4/ExDefs.lean}{UpperTheory.V.GradeZeroInterpolatingElement.Cover}
\leanref{Development/AutomorphicArithmetic/AutomorphicVarieties/SymplecticGSp4/NonemptyCoverOfRelativePresentationSmoothPt.lean}{rp_lift}
\leanref{Development/AutomorphicArithmetic/AutomorphicVarieties/CommutativeAlgebra/ChartAdd.lean}{UpperTheory.V.Alg.chart_add}
\leanref{Development/AutomorphicArithmetic/AutomorphicVarieties/CommutativeAlgebra/ChartFree.lean}{UpperTheory.V.Alg.chart_free}
\leanref{Development/AutomorphicArithmetic/AutomorphicVarieties/CommutativeAlgebra/KerGen.lean}{UpperTheory.V.Alg.ker_gen}
\begin{proof}
Let $\mathcal P=\mathcal O[[X_s,Y_t:s\in\Sigma,t\in\tau]]$, let $J$ be the ideal of $\mathcal P$ generated by $\ker\pi_B$, let $P=\mathcal P/J$, and let $\pi_P:\mathcal P\to P$ be the quotient map.
\begin{enumerate}
\item[(i)] \emph{The presentation.} The ideal $\ker\pi_B$ is generated by the sums of the terms of a given weight of its elements, which lie in it. Indeed let $I'$ be the ideal generated by these sums, and $\mathfrak n$ the maximal ideal of the noetherian local ring $\mathcal O[[X_s]]$. Every $f\in\ker\pi_B$ is congruent, modulo the ideal generated by the monomials of degree $N$, which lies in $\mathfrak n^N$, to a finite sum of such sums, for every $N$; and $\bigcap_N(I'+\mathfrak n^N)=I'$, by Krull's intersection theorem \cite[Corollary~10.19]{AM69} for $\mathcal O[[X_s]]/I'$. As $\mathcal O[[X_s]]$ is noetherian, finitely many of these sums generate $\ker\pi_B$. So $J$ is generated by finitely many series of given weights, and by Lemma~\ref{sha-counts-vd-presentations}(a), applied to the identity of $\mathcal P$, $\pi_P$ is a weighted presentation.
As $J\subseteq\ker\pi_U$, there is a surjective homomorphism $p:P\to U$ with $p\circ\pi_P=\pi_U$. The series of $J$ have constant term $0$, as those of $\ker\pi_U$ do, so $\pi_P$ has an origin $x_P$, and $x_U\circ p=x_P$, both being given by constant terms. As a quotient of $\mathcal P$ by a proper ideal, $P$ is a complete local noetherian $\mathcal O$-algebra with residue field $\mathbf F$.
The kernel of $p$ is the image of $\ker\pi_U$, so it is generated by the elements $g_u=\pi_P(f_u)$, which lie in $P_{w_f(u)}$. The parts $U_j$ and $p(P_j)$ are both the images of the series of weight $j$. Let $i_P:B\to P$ be induced by the inclusion of $\mathcal O[[X_s]]$ in $\mathcal P$, which maps $\ker\pi_B$ into $J$. Then $p\circ i_P\circ\pi_B$ is the restriction of $\pi_U$ to $\mathcal O[[X_s]]$, which is $i\circ\pi_B$; so $p\circ i_P=i$. This proves (1).
\item[(ii)] \emph{The chart.} The identity of $\mathcal O[[Y_t]]$ is a weighted presentation whose origin is the constant term, and the inclusion of $\mathcal O[[Y_t]]$ in $E[[Y_t]]$ is a chart at it: a series with no constant and no linear term lies in the ideal generated by the products of two variables, which is the square of the kernel of the constant term. The presentation $\pi_P$ is by the variables of $\pi_B$ and of the identity of $\mathcal O[[Y_t]]$, and its kernel $J$ is generated by $\ker\pi_B$ and the kernel $0$ of the identity. So Lemma~\ref{sha-counts-vd-model-sum} gives a chart of $P$ at $x_P$ in the variables of the chart of $B$ and the $Y_t$, with $n(j)+a(j)$ variables of weight $j$. This proves (2).
\item[(iii)] \emph{Derivations.} For each variable $V$ of $\mathcal P$ choose $d_V\in\mathcal O$ as follows: $d_{X_s}=D_B(\pi_B(X_s))$; $d_{Y_t}$ is a lift of $D_U(z_t)$ if $Y_t$ has weight $j$; and $d_{Y_t}=0$ otherwise. Let $L:\mathcal P\to\mathcal O$ send a series $h$ to $\sum_Vh_Vd_V$, where $h_V$ is the coefficient of $V$ in $h$; it is a derivation at the constant term.
On $\mathcal O[[X_s]]$, $L$ and $D_B\circ\pi_B$ are derivations at the constant term that agree on the variables, so they are equal, by Lemma~\ref{sha-counts-vd-presentations}(f) applied to the identity of $\mathcal O[[X_s]]$. So $L$ vanishes on $\ker\pi_B$. It vanishes on $J$, since $L(hf)=h(0)L(f)+f(0)L(h)=0$ for $f\in\ker\pi_B$, which has $f(0)=0$. Hence $L=\delta\circ\pi_P$ for a derivation $\delta$ of $P$ at $x_P$.
If $X_s$ has a weight other than $j$, then $d_{X_s}=0$, as $D_B$ has weight $j$; and $d_{Y_t}=0$ if $Y_t$ has a weight other than $j$. A series of weight $j'\ne j$ has nonzero linear coefficients only at variables of weight $j'$, so $L$ kills it, and $\delta$ has weight $j$.
The maps $L\bmod a$ and $D_U\circ\pi_U=D_U\circ p\circ\pi_P$ are derivations of $\mathcal P$ at the constant term with values in $\mathcal O/a$, and they agree on the variables: at $X_s$, $D_U(i(\pi_B(X_s)))=D_B(\pi_B(X_s))\bmod a$; at $Y_t$ of weight $j$, by the choice of $d_{Y_t}$; and at $Y_t$ of another weight both vanish. By Lemma~\ref{sha-counts-vd-presentations}(f), applied to the identity of $\mathcal P$, they are equal; so $\delta\bmod a=D_U\circ p$. This proves (3).
\end{enumerate}
\end{proof}

\begin{lemma}\label{sha-counts-vd-model-coproduct}
Let $\mathcal O$ be a complete local noetherian ring with residue field $\mathbf F$, let $X$ be a finitely generated free abelian group, and let $T$ be a finite set. For $v\in T$ let $B_v$ be a complete local noetherian $\mathcal O$-algebra with residue field $\mathbf F$, with a weighted presentation $\pi_v$ by variables $X_{v,s}$, $s\in\Sigma_v$ (Definition~\ref{sha-counts-vd-model}(1)), such that $z_{v,s}=\pi_v(X_{v,s})$ lies in $\mathfrak m_{B_v}$.
Then there are a complete local noetherian $\mathcal O$-algebra $B$ with residue field $\mathbf F$, homomorphisms of $\mathcal O$-algebras $i_v:B_v\to B$, and a weighted presentation $\pi$ of $B$ by all the variables $X_{v,s}$, with the same weights, with the following properties.
\begin{enumerate}
\item For every complete local noetherian $\mathcal O$-algebra $R$ with residue field $\mathbf F$ and homomorphisms of $\mathcal O$-algebras $f_v:B_v\to R$, $v\in T$, there is a unique homomorphism of $\mathcal O$-algebras $f:B\to R$ with $f\circ i_v=f_v$ for every $v$.
\item $\pi(X_{v,s})=i_v(z_{v,s})$ for all $v$ and $s$, and $i_v$ maps each part of $B_v$ into the part of $B$ of the same weight.
\item $\ker\pi$ is generated by the kernels of the $\pi_v$, viewed in $\mathcal O[[X_{v,s}:v\in T,s\in\Sigma_v]]$.
\end{enumerate}
\end{lemma}
\leanref{Development/AutomorphicArithmetic/AutomorphicVarieties/CommutativeAlgebra/Coprod.lean}{UpperTheory.V.Alg.coprod}
\begin{proof}
Let $\mathcal P=\mathcal O[[X_{v,s}:v\in T,s\in\Sigma_v]]$, with maximal ideal $\mathfrak n$, let $J$ be the ideal of $\mathcal P$ generated by the $\ker\pi_v$, let $B=\mathcal P/J$, and let $\pi$ be the quotient map. Let $i_v:B_v\to B$ be induced by the inclusion of $\mathcal O[[X_{v,s}:s\in\Sigma_v]]$ in $\mathcal P$, which maps $\ker\pi_v$ into $J$. Then (2) and (3) hold, as the inclusion preserves weights.
\begin{enumerate}
\item[(i)] As in step (i) of the proof of Lemma~\ref{sha-counts-vd-model-cover}, each $\ker\pi_v$ is generated by finitely many series of given weights, so $J$ is, and $\pi$ is a weighted presentation by Lemma~\ref{sha-counts-vd-presentations}(a), applied to the identity of $\mathcal P$. A series $g\in\ker\pi_v$ is not a unit, as $\pi_v(g)=0$, so its constant term lies in $\mathfrak m_{\mathcal O}$; hence $J\subseteq\mathfrak n$, and $B$ is a complete local noetherian $\mathcal O$-algebra with residue field $\mathbf F$.
\item[(ii)] Let $R$ and the $f_v$ be as in (1). Two homomorphisms of $\mathcal O$-algebras from a power series ring over $\mathcal O$ in finitely many variables to $R$ that agree on the variables and map them into $\mathfrak m_R$ are equal: every series is a polynomial plus an element of the ideal generated by the monomials of degree $N$, which both map into $\mathfrak m_R^N$, and $\bigcap_N\mathfrak m_R^N=0$.
The composite of $f_v$ with $R\to\mathbf F$ is a homomorphism of $\mathcal O$-algebras onto $\mathbf F$, whose kernel is $\mathfrak m_{B_v}$; so $f_v(z_{v,s})\in\mathfrak m_R$. As $R$ is complete, there is a homomorphism of $\mathcal O$-algebras $F:\mathcal P\to R$ with $F(X_{v,s})=f_v(z_{v,s})$ for all $v$ and $s$. By the remark, $F$ agrees with $f_v\circ\pi_v$ on $\mathcal O[[X_{v,s}:s\in\Sigma_v]]$, so it kills $J$, and it induces $f:B\to R$ with $f\circ i_v=f_v$. If $f'$ is another such homomorphism, then $f'\circ\pi$ and $F$ agree on the variables, so they are equal by the remark, and $f'=f$.
\end{enumerate}
\end{proof}

\begin{lemma}\label{sha-counts-vd-model-universal}
Let $D$ be a Siegel datum, with the notation of Definitions~\ref{sha-siegel-datum} and~\ref{sha-counts-vd-model-framing}, let $E$ be the fraction field of $\mathcal O$, and let $v\in T$.
\begin{enumerate}
\item There are a complete local noetherian $\mathcal O$-algebra $B_v$ with residue field $\mathbf F$ and a continuous lift $\rho_v:D_v\to M_4(B_v)$ of $\bar r_0$ that satisfies the conditions at $v$ of Definition~\ref{sha-counts-vd-model-framing}(1), such that every continuous lift with values in a complete local noetherian $\mathcal O$-algebra $R$ with residue field $\mathbf F$ that satisfies these conditions is the image of $\rho_v$ under a unique homomorphism of $\mathcal O$-algebras $B_v\to R$. Moreover $B_v$ has a weighted presentation (Definition~\ref{sha-counts-vd-model}(1)) by variables $X_s$ with weights in $\mathcal W$, where $X_s$ maps to $z_s=\rho_v(g)_{ab}-r_0(g)_{ab}$ for some $g\in D_v$ and some index pair $(a,b)$ of the weight of $X_s$, so that $z_s\in\mathfrak m_{B_v}$; and the entry $(a,b)$ of every $\rho_v(g)$ lies in the part of $B_v$ of the weight of $(a,b)$.
\item Let $x:B_v\to\mathcal O$ be a homomorphism of $\mathcal O$-algebras carrying $\rho_v$ to $r_0$, let $\mathfrak p$ be an ideal of $B_v$, let $y\in B_v$ with $x(y)\ne0$ and $y\mathfrak p=0$, and let $\varphi:B_v\to E[[Y_1,\dots,Y_n]]$ be a homomorphism of $\mathcal O$-algebras, where $Y_i$ has a weight $w(i)\in\mathbf Z^2$, such that: $\ker\varphi=\mathfrak p$; the constant term of $\varphi(f)$ is $x(f)$ for every $f$; for every $N\ge1$ and every series $h$ there are $f\in B_v$ and a nonzero $c\in\mathcal O$ such that $ch-\varphi(f)$ has order at least $N$; for every $N\ge1$, if $\varphi(f)$ has order at least $N$, then $cf\in(\ker x)^N+\mathfrak p$ for some nonzero $c\in\mathcal O$; and $\varphi$ maps the entry $(a,b)$ of every $\rho_v(g)$ to a series of the weight of $(a,b)$. Then the presentation of (1) has the origin $x$, and $\varphi$ is a chart of $B_v$ at $x$ (Definition~\ref{sha-counts-vd-model-charted}).
\end{enumerate}
\end{lemma}
\leanref{Development/AutomorphicArithmetic/AutomorphicVarieties/SymplecticGSp4/LocBase.lean}{rep_ex}
\leanref{Development/AutomorphicArithmetic/AutomorphicVarieties/SymplecticGSp4/LocBase.lean}{pk_of}
\leanref{Development/AutomorphicArithmetic/AutomorphicVarieties/SymplecticGSp4/LocBase.lean}{lb_chart}
\leanref{Development/AutomorphicArithmetic/AutomorphicVarieties/SymplecticGSp4/LocBase.lean}{tstable}
\leanref{Development/AutomorphicArithmetic/AutomorphicVarieties/SymplecticGSp4/LocBase.lean}{phi_dec}
\leanref{Development/AutomorphicArithmetic/AutomorphicVarieties/CommutativeAlgebra/UnivLift.lean}{UpperTheory.V.Alg.univ_lift}
\leanref{Development/AutomorphicArithmetic/AutomorphicVarieties/CommutativeAlgebra/UnivCond.lean}{UpperTheory.V.Alg.univ_cond}
\leanref{Development/AutomorphicArithmetic/AutomorphicVarieties/PadicHodgeTheory/IsInertiaInfiniteCyclotomic.lean}{UpperTheory.V.Loc.IsCompl}
\begin{proof}
(1) The proof has three steps.
\begin{enumerate}
\item[(i)] \emph{Lifts.} The group $D_v$ is compact, and $\bar r_0$ is locally constant on it, its kernel being open. Every open subgroup $H$ of $D_v$ has only finitely many continuous homomorphisms to $\mathbf Z/\ell$: the place $v_p$ identifies $D_v$ with the absolute Galois group of $\mathbf Q_v$, so $H$ is the absolute Galois group of a finite extension $K$ of $\mathbf Q_v$, and $H^1(K,\mathbf Z/\ell)$ is finite \cite[Theorem~7.1.8]{NSW08}. So the continuous lifts of $\bar r_0$ on $D_v$ are the images of a universal lift $r^\square$ with values in a complete local noetherian $\mathcal O$-algebra $R^\square$ with residue field $\mathbf F$ \cite{Maz89}.
\item[(ii)] \emph{Conditions.} For a homomorphism of $\mathcal O$-algebras $f$ from $R^\square$ to a complete local noetherian $\mathcal O$-algebra with residue field $\mathbf F$, the image of $r^\square$ under $f$ satisfies the conditions at $v$ if and only if $f$ kills the following elements: the entries of $r^\square(g)^*r^\square(g)-\theta(g)$ for $g\in D_v$; if $v=\ell$, with $e_i=\prod_{j\ne i}(r^\square(s_0)-\delta_j(s_0))$, the coefficients of the characteristic polynomial of $r^\square(s_0)$ minus those of $\prod_j(X-\delta_j(s_0))$, and the entries of $e_jr^\square(g)e_i$ for $i<j$ and $g\in D_\ell$ and of $e_i(r^\square(s)-\delta_i(s))e_i$ for $s\in I_\ell$; and if $v\in S_b$, the entries of $r^\square(g)-1$ for $g\in I_v\cap\ker\psi$. The restriction of $r_0$ satisfies the conditions, as in step (i) of the proof of Lemma~\ref{sha-counts-vd-model-tangent}, so the homomorphism $R^\square\to\mathcal O$ that carries $r^\square$ to it kills these elements. So the quotient $B_v$ of $R^\square$ by the ideal they generate is a complete local noetherian $\mathcal O$-algebra with residue field $\mathbf F$, and the image $\rho_v$ of $r^\square$ has the universal property.
\item[(iii)] \emph{Weights.} Let $R$ be a complete local noetherian $\mathcal O$-algebra with residue field $\mathbf F$, and $t\in(R^\times)^2$ with both coordinates congruent to $1$ modulo $\mathfrak m_R$. Conjugation by $\tau(t)$ (Definition~\ref{sha-counts-vd-model}) preserves the multiplier, the identities of ordinarity, in which the $e_i$ are conjugated too, and triviality on a subgroup, and it does not change the reduction. So it carries the lifts satisfying the conditions at $v$ to such lifts, and by the universal property every image of $\rho_v$ conjugated by $\tau(t)$ is again an image of $\rho_v$. By the universal property, two homomorphisms from $B_v$ that agree on the entries of the $\rho_v(g)$ are equal. So Lemma~\ref{sha-counts-vd-model-torus-grading}(2), applied with $c(g,a,b)=r_0(g)_{ab}$, gives the weighted presentation by variables $X_s$ mapping to the $z_s$, with the entries of the $\rho_v(g)$ in the parts of their weights; $r_0(g)_{ab}=0$ for $a\ne b$, the weights of the index pairs lie in $\mathcal W$, and $z_s\in\mathfrak m_{B_v}$ since $\rho_v$ lifts $\bar r_0$.
\end{enumerate}

(2) Let $\pi$ be the presentation of (1).
\begin{enumerate}
\item[(i)] \emph{Origin.} Since $x$ carries $\rho_v$ to $r_0$, $x(z_s)=0$ for every $s$. So $x\circ\pi$ and the constant term are homomorphisms of $\mathcal O$-algebras from $\mathcal O[[X_s]]$ to $\mathcal O$ that agree on the variables and map them to $0$; by step (ii) of the proof of Lemma~\ref{sha-counts-vd-model-coproduct}, with $R=\mathcal O$, they are equal. So the series of $\ker\pi$ have constant term $0$, and $x$ is the origin of $\pi$.
\item[(ii)] \emph{Weights.} The series $\varphi(z_s)$ has constant term $x(z_s)=0$, and it has the weight $w_s$ of $X_s$: $\varphi$ maps $\rho_v(g)_{ab}$ to a series of weight $w_s$, and $r_0(g)_{ab}$ is $0$ if $a\ne b$ and a constant, of weight $0=w_s$, if $a=b$. As in step (ii) of the proof of Lemma~\ref{sha-counts-vd-model-sum}, $\varphi\circ\pi$ is the substitution of the $\varphi(z_s)$ for the $X_s$, so it maps series of weight $j$ to series of weight $j$. This is property (1) of a chart, and (2) is a hypothesis.
\item[(iii)] For (3), the hypothesis with $N=2$ and $h=Y_i$ gives $f$ and $c\ne0$ such that $cY_i-\varphi(f)$ has order at least $2$; so the only linear term of $\varphi(f)$ is $cY_i$.
\item[(iv)] For (4), let $x(f)=0$ with $\varphi(f)$ without linear term. Then $\varphi(f)$ has order at least $2$, so $cf=a+p$ with $c\ne0$, $a\in(\ker x)^2$ and $p\in\mathfrak p$. As $\varphi(p)=0$, $x(p)=0$. Since $yp=0$, $x(y)p=(x(y)-y)p$, a product of two elements of $\ker x$. So $x(y)cf=x(y)a+(x(y)-y)p\in(\ker x)^2$, with $x(y)c\ne0$.
\end{enumerate}
\end{proof}

\begin{lemma}\label{sha-counts-vd-model-base}
Let $D$ be a Siegel datum, with the notation of Definitions~\ref{sha-siegel-datum} and~\ref{sha-counts-vd-model-framing}, and put $n_\ell=d$ and $n_p=d'$ for $p\in S_b$. For $v\in T$ let $B_v$, $\rho_v$ and the presentation of $B_v$ be as in Lemma~\ref{sha-counts-vd-model-universal}(1), and suppose that there are $x_v=x$, $\mathfrak p$, $y$ and $\varphi_v=\varphi$ as in Lemma~\ref{sha-counts-vd-model-universal}(2), with $n_v(j)$ variables $Y_i$ of weight $j$ for every $j$.
Then $D$ has a local base (Definition~\ref{sha-counts-vd-model-framing}(2)).
\end{lemma}
\leanref{Development/AutomorphicArithmetic/AutomorphicVarieties/SymplecticGSp4/LocBase.lean}{UpperTheory.V.GradeZeroInterpolatingElement.locBase}
\leanref{Development/AutomorphicArithmetic/AutomorphicVarieties/SymplecticGSp4/LocBase.lean}{pk_sum}
\begin{proof}
By Lemma~\ref{sha-counts-vd-model-universal}(2), the presentation $\pi_v$ of $B_v$ has the origin $x_v$, and $\varphi_v$ is a chart of $B_v$ at $x_v$ with $n_v(j)$ variables of weight $j$. Lemma~\ref{sha-counts-vd-model-coproduct}, applied to the $B_v$, gives $B$, the $i_v$ and a weighted presentation $\pi$ of $B$, by variables with weights in $\mathcal W$. Put $\rho^B_v=i_v\circ\rho_v$.
\begin{enumerate}
\item[(i)] \emph{Local data.} The conditions at $v$ are identities between polynomial expressions in the entries, with coefficients in $\mathcal O$, so they pass to images; and $(\rho^B_v)$ are local data with values in $B$. Let $(\rho'_v)$ be local data with values in $R$. By Lemma~\ref{sha-counts-vd-model-universal}(1), $\rho'_v=f_v\circ\rho_v$ for a unique $f_v$, and by Lemma~\ref{sha-counts-vd-model-coproduct}(1) there is a unique $f$ with $f\circ i_v=f_v$ for every $v$. A homomorphism $f$ with $f\circ\rho^B_v=\rho'_v$ for every $v$ has $f\circ i_v=f_v$, by the uniqueness of $f_v$; so $f$ is unique, and the local data $(\rho^B_v)$ are universal.
\item[(ii)] \emph{Entries.} The entry $(a,b)$ of $\rho_v(g)$ lies in the part of $B_v$ of the weight of $(a,b)$, so by Lemma~\ref{sha-counts-vd-model-coproduct}(2) the entry $(a,b)$ of $\rho^B_v(g)$ lies in the part of $B$ of that weight.
\item[(iii)] \emph{Origin and chart.} By Lemma~\ref{sha-counts-vd-model-coproduct}(3), the kernel of $\pi$ is generated by those of the $\pi_v$. So Lemma~\ref{sha-counts-vd-model-sum} gives an origin $x_B$ of $\pi$ and a chart of $B$ at $x_B$ with $\sum_{v\in T}n_v(j)=d(j)+|S_b|\,d'(j)=n_B(j)$ variables of weight $j$. The homomorphisms $x_B\circ i_v$ and $x_v$ agree, both being given by constant terms on $\mathcal O[[X_{v,s}:s\in\Sigma_v]]$; so $x_B$ carries $\rho^B_v$ to $r_0$. By Lemma~\ref{sha-counts-vd-model-graded}, the restriction of the chart to $B^{\mathrm{gr}}$ is a chart of $B^{\mathrm{gr}}$ at $x_B$ with $n_B(j)$ variables of weight $j$.
\end{enumerate}
So $B$ is a local base of $D$.
\end{proof}

\begin{lemma}\label{sha-counts-vd-model-lift-chart-factor-integral}
Let $\ell$ be a prime and let $R$ be a local ring that is a $\mathbf Z_\ell$-algebra.
\begin{enumerate}
\item Let $K$ be a field that is finite over $\mathbf Q_\ell$, let $y:R\to K$ be a homomorphism of $\mathbf Z_\ell$-algebras, and let $a\in\mathfrak m_R$. Then $y(a)$ is integral over $\mathbf Z_\ell$, and $y(a)^n=\ell b$ for some $n\ge0$ and some $b\in K$ integral over $\mathbf Z_\ell$.
\item Suppose that every element of $R$ is congruent modulo $\mathfrak m_R$ to an element integral over $\mathbf Z_\ell$, and let $K$ and $y$ be as in (1). Then $y(a)$ is integral over $\mathbf Z_\ell$ for every $a\in R$.
\item Suppose moreover that $R$ is noetherian, let $B$ be a commutative finite-dimensional $\mathbf Q_\ell$-algebra, and let $x:R\to B$ be a homomorphism of $\mathbf Z_\ell$-algebras. Then $x(R)$ is a finitely generated $\mathbf Z_\ell$-module.
\end{enumerate}
\end{lemma}
\leanref{Development/AutomorphicArithmetic/IntegralCasts/CapelliIdentities/IntegralOfNormBounded.lean}{UpperTheory.Reformulated.B.IntegralOfNormBounded.ptc_field}
\leanref{Development/AutomorphicArithmetic/AutomorphicVarieties/PadicHodgeTheory/PtCont.lean}{ptc_field'}
\leanref{Development/AutomorphicArithmetic/AutomorphicVarieties/PadicHodgeTheory/PtCont.lean}{ptc_finite}
\begin{proof}
On a field $K$ finite over $\mathbf Q_\ell$ we use the absolute value extending that of $\mathbf Q_\ell$; an element of $K$ is integral over $\mathbf Z_\ell$ if and only if its absolute value is at most $1$.

(1) Let $z=y(a)$; if $z=0$, take $n=1$ and $b=0$. Suppose $|z|\ge1$. Then $t=z^{-1}$ is integral: $t^{d+1}+c_dt^d+\dots+c_0=0$ with $c_i\in\mathbf Z_\ell$. Multiplying by $z^d$ gives $t=-(c_d+c_{d-1}z+\dots+c_0z^d)$, so $t=y(r)$ for some $r\in R$. Then $1-ar$ is a unit of $R$, as $ar\in\mathfrak m_R$, while $y(1-ar)=1-zt=0$, a contradiction. So $|z|<1$, and $z$ is integral; for $n$ with $|z|^n\le|\ell|$, $b=z^n/\ell$ has $|b|\le1$, so it is integral.

(2) Write $a=c+m$ with $c$ integral and $m\in\mathfrak m_R$. Then $y(a)=y(c)+y(m)$ is integral by (1).

(3) The ring $B$ is artinian. For each maximal ideal $\mathfrak n$ of $B$, $B/\mathfrak n$ is a field finite over $\mathbf Q_\ell$, and by (2) the composite $R\to B/\mathfrak n$ takes values in the integral closure $O_{\mathfrak n}$ of $\mathbf Z_\ell$ in $B/\mathfrak n$, a finitely generated $\mathbf Z_\ell$-module. So $x(R)$ maps to $\prod_{\mathfrak n}O_{\mathfrak n}$, with an image that is finitely generated over the noetherian ring $\mathbf Z_\ell$. The kernel lies in the Jacobson radical of $B$, which is nilpotent; it is a nilpotent ideal of the noetherian ring $x(R)$, a quotient of $R$, so it is finitely generated. A ring with a finitely generated nilpotent ideal $J$, whose quotient by $J$ is a finitely generated $\mathbf Z_\ell$-module, is a finitely generated $\mathbf Z_\ell$-module, as each $J^k/J^{k+1}$ is a finitely generated module over the quotient. So $x(R)$ is a finitely generated $\mathbf Z_\ell$-module.
\end{proof}

\begin{lemma}\label{sha-counts-vd-model-lift-chart-factor}
Let $\mathcal O$ be a complete discrete valuation ring of characteristic $0$ that is finite over $\mathbf Z_\ell$, with fraction field $E$ and residue field $\mathbf F$. A \emph{pointed artinian $E$-algebra} is a commutative $E$-algebra $B$ of finite dimension over $E$ with a homomorphism of $E$-algebras $\pi_B:B\to E$ whose kernel consists of nilpotent elements; it has the topology of a finite-dimensional $\mathbf Q_\ell$-vector space. Let $\Gamma$ be a compact topological group, $\iota$ a finite set, $\bar r:\Gamma\to M_\iota(\mathbf F)$ a homomorphism, $R$ a complete local noetherian $\mathcal O$-algebra with residue field $\mathbf F$, and $\rho:\Gamma\to M_\iota(R)$ a continuous lift of $\bar r$. Let $x:R\to E$ be a homomorphism of $\mathcal O$-algebras, $B$ a pointed artinian $E$-algebra, and $r:\Gamma\to M_\iota(B)$ a homomorphism with continuous entries such that $\pi_B(r(g)_{ab})=x(\rho(g)_{ab})$ for all $g$, $a$ and $b$.
Then there are an $\mathcal O$-subalgebra $C$ of $B$ that is a complete local noetherian $\mathcal O$-algebra with residue field $\mathbf F$, and a continuous lift $\rho':\Gamma\to M_\iota(C)$ of $\bar r$, with $\rho'(g)=r(g)$ for every $g$.
\end{lemma}
\leanref{Development/AutomorphicArithmetic/AutomorphicVarieties/PadicHodgeTheory/ExistsLiftFactorCompactSubalgebra.lean}{UpperTheory.V.Loc.exists_lift_factor_compact_subalgebra}
\begin{proof}
Write $\pi=\pi_B$. Every element of $R$ is congruent modulo $\mathfrak m_R$ to an element of the image of $\mathcal O$, which is integral over $\mathbf Z_\ell$.
\begin{enumerate}
\item[(i)] By Lemma~\ref{sha-counts-vd-model-lift-chart-factor-integral}(2), $x(a)$ is integral over $\mathbf Z_\ell$, hence over $\mathcal O$, for every $a\in R$.
\item[(ii)] Let $C$ be the $\mathcal O$-subalgebra of $B$ generated by the entries $r(g)_{ab}$. As $\pi(r(g)_{ab})=x(\rho(g)_{ab})$, $\pi$ maps $C$ into $x(R)$.
\item[(iii)] \emph{$C$ is finite over $\mathcal O$.} Fix a basis $e_1,\dots,e_n$ of $B$ over $\mathbf Q_\ell$. The coordinates of the entries are continuous functions on the compact group $\Gamma$, so they are bounded, and all entries lie in the $\mathbf Z_\ell$-span of the $ce_i$ for a suitable $c\in\mathbf Q_\ell^\times$. So the $\mathbf Z_\ell$-span of the entries is generated by a finite set $T$, and $C=\mathcal O[T]$. Each $t\in T$ is integral over $\mathcal O$: by (i) and (ii), $P(\pi(t))=0$ for a monic $P$ over $\mathcal O$; then $P(t)\in\ker\pi$ is nilpotent, and $P^k(t)=0$ for some $k$. So $C$ is a finite $\mathcal O$-algebra.
\item[(iv)] \emph{$C$ is complete local noetherian with residue field $\mathbf F$.} The ring $x(R)\subseteq E$ is a quotient of $R$, so it is a complete local noetherian $\mathcal O$-algebra with residue field $\mathbf F$, and it is integral over $\mathcal O$ by (i). Let $p:C\to x(R)$ be the restriction of $\pi$, whose kernel consists of nilpotent elements. An element $a\in C$ with $p(a)$ a unit is a unit: if $u$ is the inverse of $p(a)$, then $u$ is integral over $\mathcal O$, and multiplying a monic equation of degree $d$ for $u$ by $p(a)^d$ gives $1+aq(a)\in\ker p$ for a polynomial $q$ over $\mathcal O$; so $aq(a)$ is $-1$ plus a nilpotent element, a unit. So $C$ is local, as $x(R)$ is. Every $c\in C$ is congruent modulo $\mathfrak m_C$ to an element of $\mathcal O$: if $p(c)-a\in\mathfrak m_{x(R)}$ with $a\in\mathcal O$, then $c-a$ is not a unit. A local ring that is finite over the complete local noetherian ring $\mathcal O$, with this property, is a complete local noetherian $\mathcal O$-algebra with residue field $\mathbf F$.
\item[(v)] \emph{The lift.} Let $\rho'(g)=r(g)$, with values in $M_\iota(C)$. The reduction of $\rho'(g)_{ab}$ is that of $p(r(g)_{ab})=x(\rho(g)_{ab})$, which is $\bar r(g)_{ab}$. For continuity, let $N\ge1$ and $g_0\in\Gamma$. The finitely generated torsion-free $\mathbf Z_\ell$-module $C$ is free, and the coordinates of $r(g)_{ab}$ in a basis of it extend to $\mathbf Q_\ell$-linear forms on $B$, so they are continuous in $g$. So for $g$ near $g_0$, $r(g)_{ab}-r(g_0)_{ab}\in\ell^NC\subseteq\mathfrak m_C^N$, as $\ell\in\mathfrak m_{\mathcal O}$ maps into $\mathfrak m_C$.
\end{enumerate}
\end{proof}

\begin{lemma}\label{sha-counts-vd-model-lift-chart-continuity}
Let $\mathcal O$ be as in Lemma~\ref{sha-counts-vd-model-lift-chart-factor}, let $R$ be a complete local noetherian $\mathcal O$-algebra with residue field $\mathbf F$, let $B$ be a commutative $\mathcal O$-algebra that is a finite-dimensional $\mathbf Q_\ell$-algebra, with its topology as a $\mathbf Q_\ell$-vector space, and let $f:R\to B$ be a homomorphism of $\mathcal O$-algebras.
\begin{enumerate}
\item Let $Y$ be a topological space and $h:Y\to R$ a function that is continuous for the $\mathfrak m_R$-adic topology, that is, $h$ modulo $\mathfrak m_R^N$ is locally constant for every $N$. Then $f\circ h$ is continuous.
\item If $B$ is a field, then $f(a)$ is integral over $\mathbf Z_\ell$ for every $a\in R$, and for every $a\in\mathfrak m_R$ there are $n\ge0$ and $b\in B$ integral over $\mathbf Z_\ell$ with $f(a)^n=\ell b$.
\end{enumerate}
\end{lemma}
\leanref{Development/AutomorphicArithmetic/AutomorphicVarieties/PadicHodgeTheory/PtCont.lean}{UpperTheory.V.Loc.ptCont}
\begin{proof}
Every element of $R$ is congruent modulo $\mathfrak m_R$ to an element of the image of $\mathcal O$, which is integral over $\mathbf Z_\ell$. So (2) is Lemma~\ref{sha-counts-vd-model-lift-chart-factor-integral}(1) and~(2).

(1) Let $A=f(R)$, a finitely generated $\mathbf Z_\ell$-module by Lemma~\ref{sha-counts-vd-model-lift-chart-factor-integral}(3).
\begin{enumerate}
\item[(i)] There is $N$ with $f(\mathfrak m_R^{Nk})\subseteq\ell^kA$ for every $k\ge0$. Let $P$ be a prime ideal of $A$ containing $\ell$. Then $P\cap\mathbf Z_\ell=\ell\mathbf Z_\ell$ is maximal and $A$ is integral over $\mathbf Z_\ell$, so $P$ is maximal, and its preimage in $R$ is a maximal ideal, that is $\mathfrak m_R$. So $f(\mathfrak m_R)A$ lies in the radical of $\ell A$, and as it is finitely generated, $(f(\mathfrak m_R)A)^N\subseteq\ell A$ for some $N$.
\item[(ii)] Fix a basis of $B$ over $\mathbf Q_\ell$, and let $\varphi:A\to\mathbf Q_\ell$ be the restriction of one of the coordinates. It is $\mathbf Z_\ell$-linear, so it is bounded, $|\varphi(b)|\le C$, as $A$ is finitely generated. Let $y\in Y$ and $k\ge0$. For $y'$ near $y$, $h(y')\equiv h(y)$ modulo $\mathfrak m_R^{Nk}$, so $f(h(y'))-f(h(y))\in\ell^kA$ by (i), and $|\varphi(f(h(y')))-\varphi(f(h(y)))|\le C\ell^{-k}$.
\end{enumerate}
So every coordinate of $f\circ h$ is continuous, and so is $f\circ h$.
\end{proof}

\begin{lemma}\label{sha-counts-vd-model-lift-chart-generation-span}
Let $\mathcal O$ be a local ring with residue field $\mathbf F$, let $B$ and $R$ be complete local noetherian $\mathcal O$-algebras with residue field $\mathbf F$, let $i:B\to R$ be a homomorphism of $\mathcal O$-algebras, and let $S\subseteq\mathfrak m_R$.
\begin{enumerate}
\item If any two homomorphisms of $\mathcal O$-algebras $R\to\mathbf F[\varepsilon]$, $\varepsilon^2=0$, that agree on $i(B)$ and on $S$ are equal, then $\mathfrak m_R$ is contained in the ideal generated by $S$, $i(\mathfrak m_B)$ and $\mathfrak m_R^2$.
\item Suppose that $\mathfrak m_R$ is contained in the ideal generated by $S$, $i(\mathfrak m_B)$ and $\mathfrak m_R^2$. Let $R'$ be an $\mathcal O$-algebra with an ideal $I$ such that $\bigcap_kI^k=0$, and let $f,f':R\to R'$ be homomorphisms of $\mathcal O$-algebras that agree on $i(B)$ and on $S$, such that $f(\mathfrak m_R^n)$ and $f'(\mathfrak m_R^n)$ lie in $I$ for some $n$. Then $f=f'$.
\end{enumerate}
\end{lemma}
\leanref{Development/AutomorphicArithmetic/AutomorphicVarieties/CommutativeAlgebra/CoSpan.lean}{UpperTheory.V.Alg.co_span}
\leanref{Development/AutomorphicArithmetic/AutomorphicVarieties/CommutativeAlgebra/CoExt.lean}{UpperTheory.V.Alg.co_ext}
\begin{proof}
Let $J$ be the ideal generated by $S$, $i(\mathfrak m_B)$ and $\mathfrak m_R^2$; it contains the image of $\mathfrak m_{\mathcal O}$, which lies in $i(\mathfrak m_B)$.

(1) Suppose that $x\in\mathfrak m_R$ is not in $J$, and let $N=J+\mathcal O\cdot1$, an $\mathcal O$-submodule of $R$. Then $x\notin N$: if $x=y+a$ with $y\in J$ and $a\in\mathcal O$, then $a\in\mathfrak m_{\mathcal O}$, as $x-y\in\mathfrak m_R$, and so $x\in J$. The quotient $R/N$ is killed by $\mathfrak m_{\mathcal O}$, so it is an $\mathbf F$-vector space, and there is an $\mathcal O$-linear $D:R\to\mathbf F$ that vanishes on $N$ with $D(x)\ne0$. Writing $y=a+y_1$ and $z=b+z_1$ with $a,b\in\mathcal O$ and $y_1,z_1\in\mathfrak m_R$, and using $y_1z_1\in\mathfrak m_R^2\subseteq N$, one gets $D(yz)=\bar y D(z)+D(y)\bar z$, where the bar is the residue. So $f(y)=\bar y+\varepsilon D(y)$ is a homomorphism of $\mathcal O$-algebras $R\to\mathbf F[\varepsilon]$. It agrees with $f'(y)=\bar y$ on $S$ and on $i(B)$, as $D$ vanishes on $S$, on $i(\mathfrak m_B)$ and on $\mathcal O$; but $f(x)\ne f'(x)$, a contradiction.

(2) For every $k$ and every $x\in R$ there is $a$ in the $\mathcal O$-subalgebra generated by $S$ and $i(\mathfrak m_B)$ with $x-a\in\mathfrak m_R^k$: by induction on $k$, as every element of $R$ is congruent to an element of $\mathcal O$ modulo $\mathfrak m_R$, and $\mathfrak m_R^j\subseteq(S,i(\mathfrak m_B))\mathfrak m_R^{j-1}+\mathfrak m_R^{j+1}$. As $f(a)=f'(a)$, $f(x)-f'(x)=f(x-a)-f'(x-a)$, which lies in $I^j$ if $k=nj$. So $f(x)-f'(x)\in\bigcap_jI^j=0$.
\end{proof}

\begin{lemma}\label{sha-counts-vd-model-lift-chart-generation}
Let $\mathcal O$ be as in Lemma~\ref{sha-counts-vd-model-lift-chart-factor}, let $\Gamma$ be a topological group, $\iota$ a finite set, $\bar r:\Gamma\to M_\iota(\mathbf F)$ a homomorphism, and $\rho:\Gamma\to M_\iota(R)$ a continuous lift of $\bar r$ with values in a complete local noetherian $\mathcal O$-algebra $R$ with residue field $\mathbf F$, such that two homomorphisms of $\mathcal O$-algebras from $R$ to a complete local noetherian $\mathcal O$-algebra with residue field $\mathbf F$ that agree on the entries of the $\rho(g)$ are equal. Let $B$ be as in Lemma~\ref{sha-counts-vd-model-lift-chart-continuity}. Then two homomorphisms of $\mathcal O$-algebras $R\to B$ that agree on the entries of the $\rho(g)$ are equal.
\end{lemma}
\leanref{Development/AutomorphicArithmetic/AutomorphicVarieties/PadicHodgeTheory/HomExt.lean}{UpperTheory.V.Loc.homExt}
\begin{proof}
Let $f,f':R\to B$ agree on the entries of the $\rho(g)$. For each entry $e=\rho(g)_{ab}$ choose $c(e)\in\mathcal O$ with $e-c(e)\in\mathfrak m_R$, and let $S\subseteq\mathfrak m_R$ be the set of the $e-c(e)$.
\begin{enumerate}
\item[(i)] $\mathfrak m_R$ lies in the ideal generated by $S$, $\mathfrak m_{\mathcal O}R$ and $\mathfrak m_R^2$. Indeed, $\mathbf F[\varepsilon]$ is a complete local noetherian $\mathcal O$-algebra with residue field $\mathbf F$, and two homomorphisms $R\to\mathbf F[\varepsilon]$ that agree on $S$ agree on the entries, so they are equal by the hypothesis; so Lemma~\ref{sha-counts-vd-model-lift-chart-generation-span}(1) applies, with $B=\mathcal O$.
\item[(ii)] By Lemma~\ref{sha-counts-vd-model-lift-chart-factor-integral}(3), $f(R)$ and $f'(R)$ are finitely generated $\mathbf Z_\ell$-modules, hence $\mathcal O$-modules, and so is the $\mathcal O$-subalgebra $A$ of $B$ they generate. As $A$ is a finitely generated module over the complete noetherian local ring $\mathcal O$, $\bigcap_k\mathfrak m_{\mathcal O}^kA=0$.
\item[(iii)] There is $N$ with $f(\mathfrak m_R^N)\subseteq\mathfrak m_{\mathcal O}A$, and the same for $f'$. Indeed, let $P$ be a prime ideal of $A$ containing $\mathfrak m_{\mathcal O}A$. As $A$ is integral over $\mathcal O$, hence over $f(R)$, $P$ is maximal, and so is $P\cap f(R)$, whose preimage in $R$ is therefore $\mathfrak m_R$. So $f(\mathfrak m_R)A$ lies in the radical of $\mathfrak m_{\mathcal O}A$, and being finitely generated, it has a power in $\mathfrak m_{\mathcal O}A$.
\end{enumerate}
By Lemma~\ref{sha-counts-vd-model-lift-chart-generation-span}(2), with $B=\mathcal O$, $R'=A$ and $I=\mathfrak m_{\mathcal O}A$, $f$ and $f'$, as maps to $A$, are equal: they agree on $\mathcal O$ and on $S$.
\end{proof}

\begin{definition}\label{sha-counts-vd-model-lift-chart-complete}
Let $\mathcal O$ be as in Lemma~\ref{sha-counts-vd-model-lift-chart-factor}, let $R$ be an $\mathcal O$-algebra, $\mathfrak p$ an ideal of $R$, $x:R\to\mathcal O$ a homomorphism of $\mathcal O$-algebras, and $\varphi:R\to E[[Y_1,\dots,Y_n]]$ a homomorphism of $\mathcal O$-algebras. Then $\varphi$ is a \emph{complete chart} of $(R,\mathfrak p)$ at $x$ if: $\ker\varphi=\mathfrak p$; the constant term of $\varphi(f)$ is $x(f)$ for every $f\in R$; for every $N\ge1$ and every series $h$ there are $f\in R$ and a nonzero $c\in\mathcal O$ such that $ch-\varphi(f)$ has order at least $N$; and for every $N\ge1$, if $\varphi(f)$ has order at least $N$, then $cf\in(\ker x)^N+\mathfrak p$ for some nonzero $c\in\mathcal O$.
\end{definition}
\leanref{Development/AutomorphicArithmetic/AutomorphicVarieties/PadicHodgeTheory/IsInertiaInfiniteCyclotomic.lean}{UpperTheory.V.Loc.IsCompl}

\begin{lemma}\label{sha-counts-vd-model-lift-chart-smooth-parameters}
Let $\mathcal O$ be a domain with fraction field $E$, let $R$ be a noetherian $\mathcal O$-algebra, $x:R\to\mathcal O$ a homomorphism of $\mathcal O$-algebras, and $I\subseteq\ker x$ an ideal. Then there are $n\ge0$, $\mathcal O$-linear maps $D_1,\dots,D_n:R\to E$ with $D_i(ab)=x(a)D_i(b)+x(b)D_i(a)$ that vanish on $I$, elements $t_1,\dots,t_n\in\ker x$, and a nonzero $c\in\mathcal O$ with $D_j(t_i)=c$ if $i=j$ and $0$ otherwise, such that for every $a\in\ker x$ there are a nonzero $e\in\mathcal O$ and $o_1,\dots,o_n\in\mathcal O$ with $ea-\sum_io_it_i\in(\ker x)^2+I$.
\end{lemma}
\leanref{Development/AutomorphicArithmetic/AutomorphicVarieties/PadicHodgeTheory/ExistsDerivationsDualToCotangentParameters.lean}{UpperTheory.exists_derivations_dual_to_cotangent_parameters}
\begin{proof}
Let $\mathfrak q=\ker x$ and $J=\mathfrak q^2+I$, an ideal contained in $\mathfrak q$. Let $N$ be the $\mathcal O$-submodule $\mathcal O\cdot1+J$ of $R$, let $V=E\otimes_{\mathcal O}(R/N)$, and let $\pi:R\to V$ be the $\mathcal O$-linear map $a\mapsto1\otimes(a\bmod N)$, which kills $N$.
\begin{enumerate}
\item[(i)] $\pi(ab)=x(a)\pi(b)+x(b)\pi(a)$ for $a,b\in R$. Indeed $a-x(a)$ and $b-x(b)$ lie in $\mathfrak q$, so their product lies in $\mathfrak q^2\subseteq N$, and $ab=(a-x(a))(b-x(b))-x(a)x(b)\cdot1+x(a)b+x(b)a$.
\item[(ii)] If $x(a)=0$ and $\pi(a)=0$, then $ea\in J$ for some nonzero $e\in\mathcal O$. Indeed $E$ is the localization of $\mathcal O$ at its nonzero elements, so $\pi(a)=0$ means that $e\cdot(a\bmod N)=0$ for some nonzero $e$, that is $ea=d\cdot1+v$ with $d\in\mathcal O$ and $v\in J$. Applying $x$, which vanishes on $J\subseteq\mathfrak q$, gives $0=d$, so $ea=v\in J$.
\item[(iii)] As $R$ is noetherian, $\mathfrak q$ is generated by finitely many elements $g_1,\dots,g_m$. The $\pi(g_k)$ span $V$ over $E$. Indeed, for $a=\sum_kr_kg_k\in\mathfrak q$, (i) gives $\pi(a)=\sum_kx(r_k)\pi(g_k)$, as $x(g_k)=0$; every $a\in R$ has $\pi(a)=\pi(a-x(a)\cdot1)$, with $a-x(a)\in\mathfrak q$; and $V$ is spanned over $E$ by the $\pi(a)$.
\item[(iv)] Choose among the $\pi(g_k)$ a basis $b_1,\dots,b_n$ of $V$, with $b_i=\pi(t_i)$ and each $t_i$ one of the $g_k$, so that $t_i\in\ker x$. Let $D_i$ be $\pi$ followed by the $i$-th coordinate in this basis, and $c=1$. Then $D_i$ is $\mathcal O$-linear, satisfies $D_i(ab)=x(a)D_i(b)+x(b)D_i(a)$ by (i), vanishes on $I\subseteq N$, and $D_j(t_i)$ is $1$ if $i=j$ and $0$ otherwise.
\item[(v)] Let $a\in\ker x$. Choose a nonzero $d\in\mathcal O$ such that $d$ times each coordinate of $\pi(a)$ is an element $o_i$ of $\mathcal O$. Then $\pi(da-\sum_io_it_i)=0$ and $x(da-\sum_io_it_i)=0$, so by (ii) there is a nonzero $e'\in\mathcal O$ with $e'(da-\sum_io_it_i)\in J$. So $e=e'd$ and the $e'o_i$ are as required.
\end{enumerate}
% exists_derivations_dual_to_cotangent_parameters (111 lines). The proof follows it: N, the tensor product E (x) R/N and pi; (i) is hder, (ii) is hker (IsLocalizedModule.eq_zero_iff), (iii) is hg and hspan, (iv) is exists_linearIndependent' and the basis b (with c = 1), (v) is the last item (exist_integer_multiples_of_finite).
\end{proof}

\begin{lemma}\label{sha-counts-vd-model-lift-chart-smooth-homogeneous}
Let $\mathcal O$ be a domain, $R$ an $\mathcal O$-algebra, $I$ an ideal of $R$, and $x:R\to\mathcal O$ a homomorphism of $\mathcal O$-algebras. Let $t_1,\dots,t_n\in\ker x$ be such that for every $a\in\ker x$ there are a nonzero $e\in\mathcal O$ and $o_i\in\mathcal O$ with $ea-\sum_io_it_i\in(\ker x)^2+I$. Then for every $N\ge0$ and $a\in(\ker x)^N$ there are a nonzero $e\in\mathcal O$ and a homogeneous polynomial $h\in\mathcal O[T_1,\dots,T_n]$ of degree $N$ with $ea-h(t_1,\dots,t_n)\in(\ker x)^{N+1}+I$.
\end{lemma}
\leanref{Development/AutomorphicArithmetic/AutomorphicVarieties/PadicHodgeTheory/ExistsHomogeneousLiftModKerPowSuccSup.lean}{UpperTheory.exists_homogeneous_lift_mod_ker_pow_succ_sup}
\begin{proof}
Let $\mathfrak q=\ker x$, and write $h(t)$ for $h(t_1,\dots,t_n)$. By induction on $N$, there are moreover such $e$ and $h$ with $h(t)\in\mathfrak q^N$. For $N=0$ take $e=1$ and $h$ the constant $x(a)$; then $a-x(a)\in\mathfrak q$.
Suppose the statement holds for $N$. The ideal $\mathfrak q^{N+1}$ consists of the sums of products $bc$ with $b\in\mathfrak q^N$ and $c\in\mathfrak q$.
\begin{enumerate}
\item[(i)] \emph{Products.} By induction there are a nonzero $e_1$ and $h_1$ homogeneous of degree $N$ with $h_1(t)\in\mathfrak q^N$ and $e_1b-h_1(t)\in\mathfrak q^{N+1}+I$; by hypothesis there are a nonzero $e_2$ and $L=\sum_io_iT_i$ with $e_2c-L(t)\in\mathfrak q^2+I$. Then $h_1L$ is homogeneous of degree $N+1$, $(h_1L)(t)\in\mathfrak q^{N+1}$ as the $t_i$ lie in $\mathfrak q$, and
\[
e_1e_2bc-(h_1L)(t)=(e_2c)\big(e_1b-h_1(t)\big)+h_1(t)\big(e_2c-L(t)\big),
\]
where the first term lies in $\mathfrak q(\mathfrak q^{N+1}+I)$ and the second in $\mathfrak q^N(\mathfrak q^2+I)$, both contained in $\mathfrak q^{N+2}+I$.
\item[(ii)] \emph{Sums.} If $e_1a_1-h_1(t)$ and $e_2a_2-h_2(t)$ lie in $\mathfrak q^{N+2}+I$, with $h_1(t),h_2(t)\in\mathfrak q^{N+1}$, then $e_1e_2(a_1+a_2)-(e_2h_1+e_1h_2)(t)=e_2\big(e_1a_1-h_1(t)\big)+e_1\big(e_2a_2-h_2(t)\big)$ lies there too, and $e_2h_1+e_1h_2$ is homogeneous of degree $N+1$.
\end{enumerate}
% exists_homogeneous_lift_mod_ker_pow_succ_sup (90 lines). The proof follows it: the stronger statement is key, (i) and (ii) are the two cases of Submodule.mul_induction_on, with the private mul_mem_pow_sup.
\end{proof}

\begin{lemma}\label{sha-counts-vd-model-lift-chart-smooth-vanishing}
Let $\mathcal O$ be a domain with fraction field $E$, let $c\in\mathcal O$ be nonzero, and let $u_1,\dots,u_n\in E[[Y_1,\dots,Y_n]]$ with $u_i-cY_i\in\mathfrak m^2$, where $\mathfrak m$ is the maximal ideal. If $h\in\mathcal O[T_1,\dots,T_n]$ is homogeneous of degree $N$ and $h(u_1,\dots,u_n)\in\mathfrak m^{N+1}$, then $h=0$.
\end{lemma}
\leanref{Development/AutomorphicArithmetic/AutomorphicVarieties/PadicHodgeTheory/EqZeroOfHomogeneousAevalMemMaximalIdealPowSucc.lean}{UpperTheory.eq_zero_of_homogeneous_aeval_mem_maximalIdeal_pow_succ}
\begin{proof}
Say that $(a,b)$ agree to order $k$ if $b\in\mathfrak m^k$ and $a-b\in\mathfrak m^{k+1}$. If $(a,b)$ agree to order $k$ and $(a',b')$ to order $k'$, then $(aa',bb')$ agree to order $k+k'$: $a'=b'+(a'-b')$ lies in $\mathfrak m^{k'}$, $bb'\in\mathfrak m^{k+k'}$, and $aa'-bb'=(a-b)a'+b(a'-b')\in\mathfrak m^{k+k'+1}$.
Let $v_i=cY_i$. The pairs $(u_i,v_i)$ agree to order $1$, as $v_i\in\mathfrak m$; so $(u^d,v^d)$ agree to order $|d|$ for every exponent vector $d$, where $u^d=\prod_iu_i^{d_i}$ and $|d|=\sum_id_i$. Every monomial of $h$ has degree $N$, so $h(u)-h(v)\in\mathfrak m^{N+1}$, and hence $h(v)\in\mathfrak m^{N+1}$. A series in $\mathfrak m^{N+1}$ has no terms of degree at most $N$. But $h(v)=\sum_dh_dc^NY^d$, where $h_d$ is the coefficient of $T^d$ in $h$ and $d$ runs over the exponent vectors of degree $N$. So $h_dc^N=0$ in $E$ for every $d$, and $h_d=0$, as $c\ne0$ and $\mathcal O$ embeds in $E$.
% eq_zero_of_homogeneous_aeval_mem_maximalIdeal_pow_succ (124 lines with its helpers). The proof follows it: agreement to order k is the private Cg, with Cg.mul, Cg.pow and Cg.prod; the order of an element of m^(N+1) is order_of_mem_pow; h(v) is hval.
\end{proof}

\begin{lemma}\label{sha-counts-vd-model-lift-chart-smooth-approximation}
Let $\mathcal O$ be a domain with fraction field $E$, $R$ an $\mathcal O$-algebra, $c\in\mathcal O$ nonzero, and $\varphi:R\to E[[Y_1,\dots,Y_n]]$ a homomorphism of $\mathcal O$-algebras with elements $t_i\in R$ such that $\varphi(t_i)-cY_i\in\mathfrak m^2$, where $\mathfrak m$ is the maximal ideal. Then for every $N\ge1$ and every series $f$ there are $a\in R$ and a nonzero $e\in\mathcal O$ with $ef-\varphi(a)\in\mathfrak m^N$.
\end{lemma}
\leanref{Development/AutomorphicArithmetic/AutomorphicVarieties/PadicHodgeTheory/PowerSeriesApproximationModMaximal.lean}{UpperTheory.power_series_approximation_mod_maximal}
\begin{proof}
For an ideal $J$ of $E[[Y_1,\dots,Y_n]]$, say that a series $f$ is \emph{approximable modulo $J$} if $ef-\varphi(a)\in J$ for some $a\in R$ and some nonzero $e\in\mathcal O$.
\begin{enumerate}
\item[(i)] \emph{Sums.} The series $0$ is approximable modulo $J$, with $a=0$ and $e=1$. If $ef-\varphi(a)$ and $e'g-\varphi(b)$ lie in $J$, then $ee'(f+g)-\varphi(e'a+eb)=e'\big(ef-\varphi(a)\big)+e\big(e'g-\varphi(b)\big)$ lies in $J$.
\item[(ii)] \emph{Products.} Let $f\in\mathfrak m^i$ be approximable modulo $\mathfrak m^{i+1}$, with $a$ and $e$, and $g\in\mathfrak m^j$ approximable modulo $\mathfrak m^{j+1}$, with $b$ and $e'$. Then $\varphi(b)=e'g-\big(e'g-\varphi(b)\big)\in\mathfrak m^j$, and $ee'fg-\varphi(ab)=ef\big(e'g-\varphi(b)\big)+\big(ef-\varphi(a)\big)\varphi(b)$ lies in $\mathfrak m^{i+j+1}$; so $fg$ is approximable modulo $\mathfrak m^{i+j+1}$.
\item[(iii)] Every series $g$ is approximable modulo $\mathfrak m$. Write its constant term as $a/e$ with $a,e\in\mathcal O$ and $e\ne0$; then $eg-\varphi(a\cdot1)$ has constant term $0$.
\item[(iv)] Every $g\in\mathfrak m$ is approximable modulo $\mathfrak m^2$. Write $g=\sum_ir_iY_i$ with series $r_i$. Each $r_i$ is approximable modulo $\mathfrak m$ by (iii), and $Y_i$ is approximable modulo $\mathfrak m^2$, with $t_i$ and $c$. So $g$ is approximable modulo $\mathfrak m^2$ by (ii), with $i=0$ and $j=1$, and (i).
\item[(v)] For every $N$, every $f\in\mathfrak m^N$ is approximable modulo $\mathfrak m^{N+1}$. By induction on $N$: the case $N=0$ is (iii), and $\mathfrak m^{N+1}$ consists of the sums of products $ab$ with $a\in\mathfrak m^N$ and $b\in\mathfrak m$, to which (ii), with $j=1$, and (iv) apply; sums are covered by (i).
\item[(vi)] For every $N$, every series $f$ is approximable modulo $\mathfrak m^N$. By induction on $N$, the case $N=0$ being clear. If $ef-\varphi(a)\in\mathfrak m^N$, then by (v) there are $b\in R$ and a nonzero $d\in\mathcal O$ with $d\big(ef-\varphi(a)\big)-\varphi(b)\in\mathfrak m^{N+1}$, that is $(de)f-\varphi(da+b)\in\mathfrak m^{N+1}$.
\end{enumerate}
% power_series_approximation_mod_maximal (107 lines with its helpers). The proof follows it: approximability is the private Appr; (i) is appr_zero and appr_add, (ii) is appr_mul, (iii) is h0, (iv) is h1 (with mem_span_X, UpperTheory.Reformulated.B.MemSpanVariable.mem_span_X), (v) is appr_step and (vi) is appr_all.
\end{proof}

\begin{lemma}\label{sha-counts-vd-model-lift-chart-smooth-prime}
Let $R$ be a noetherian ring, $I\subseteq\mathfrak q$ ideals of $R$ with $\mathfrak q$ prime, and $P$ an ideal such that for every $a\in P$ and $N\ge0$ there is $s\notin\mathfrak q$ with $sa\in\mathfrak q^N+I$. Then there is $y\notin\mathfrak q$ with $yP\subseteq I$.
\end{lemma}
\leanref{Development/AutomorphicArithmetic/AutomorphicVarieties/PadicHodgeTheory/ExistsNotMemMulMemIdealOfForallPrimePowSup.lean}{UpperTheory.exists_not_mem_mul_mem_ideal_of_forall_prime_pow_sup}
\begin{proof}
\begin{enumerate}
\item[(i)] \emph{One element.} Let $a\in P$. The image $\bar{\mathfrak q}$ of $\mathfrak q$ in $\bar R=R/I$ is a prime ideal, as $I\subseteq\mathfrak q$, and the localization $\bar R_{\bar{\mathfrak q}}$ is a noetherian local ring whose maximal ideal $\mathfrak n$ is generated by $\bar{\mathfrak q}$. For every $N$ there is $s\notin\mathfrak q$ with $sa\in\mathfrak q^N+I$, so the image of $sa$ in $\bar R_{\bar{\mathfrak q}}$ lies in $\mathfrak n^N$; as $s$ maps to a unit, so does the image of $a$. By Krull's intersection theorem \cite[Theorem~8.10]{Mat86}, $\bigcap_N\mathfrak n^N=0$, so $a$ maps to $0$ in $\bar R_{\bar{\mathfrak q}}$: there is $y\notin\mathfrak q$ with $ya\in I$.
\item[(ii)] \emph{The ideal.} As $R$ is noetherian, $P$ is generated by finitely many elements $a_1,\dots,a_k$. By (i) there are $y_i\notin\mathfrak q$ with $y_ia_i\in I$. Then $y=y_1\cdots y_k$ does not lie in the prime $\mathfrak q$, and $ya_i\in I$ for every $i$, so $yP\subseteq I$.
\end{enumerate}
% exists_not_mem_mul_mem_ideal_of_forall_prime_pow_sup (85 lines with its helpers). The proof follows it: (i) is the private a3one_of_factSplit_of_countCont (Ideal.iInf_pow_eq_bot_of_isLocalRing for the localization of R/I), which also treats the case where I is not contained in q; (ii) is hgen, by induction on a finite set of generators, and the span induction.
\end{proof}

\begin{lemma}\label{sha-counts-vd-model-lift-chart-smooth-height}
Let $\mathcal O$ be a domain with fraction field $E$, $S$ a noetherian $\mathcal O$-algebra, $x:S\to\mathcal O$ a homomorphism of $\mathcal O$-algebras, and $t_1,\dots,t_n\in\ker x$.
\begin{enumerate}
\item If for every $a\in\ker x$ there are a nonzero $e\in\mathcal O$ and $o_i\in\mathcal O$ with $ea-\sum_io_it_i\in(\ker x)^2$, then $\ker x$ has height at most $n$.
\item If there are a nonzero $c\in\mathcal O$ and a homomorphism of $\mathcal O$-algebras $\varphi:S\to E[[Y_1,\dots,Y_n]]$ such that the constant term of $\varphi(a)$ is $x(a)$ for every $a$ and $\varphi(t_i)-cY_i\in\mathfrak m^2$ for every $i$, where $\mathfrak m$ is the maximal ideal of $E[[Y_1,\dots,Y_n]]$, then $\ker x$ has height at least $n$.
\item Suppose that $\mathcal O$ is a complete discrete valuation ring and that $S$ is a domain and a complete local noetherian $\mathcal O$-algebra with residue field that of $\mathcal O$. Then $\dim S=\operatorname{ht}(\ker x)+1$.
\end{enumerate}
\end{lemma}
\leanref{Development/AutomorphicArithmetic/AutomorphicVarieties/PadicHodgeTheory/HeightKerSectionLeOfCotangentGenerators.lean}{UpperTheory.height_ker_section_le_of_cotangent_generators}
\leanref{Development/AutomorphicArithmetic/AutomorphicVarieties/PadicHodgeTheory/RingKrullDimMvPowerSeriesFinField.lean}{UpperTheory.ringKrullDim_mvPowerSeries_fin_field}
\leanref{Development/AutomorphicArithmetic/AutomorphicVarieties/PadicHodgeTheory/RingKrullDimEqHeightKerSectionAddOne.lean}{UpperTheory.ringKrullDim_eq_height_ker_section_add_one}
\begin{proof}
Let $\mathfrak q=\ker x$, a prime ideal, as $S/\mathfrak q$ embeds in $\mathcal O$. A nonzero $e\in\mathcal O$ maps to a unit of $S_{\mathfrak q}$, as $x(e)=e\ne0$.

(1) By the hypothesis, the images of the $t_i$ span $\mathfrak qS_{\mathfrak q}/\mathfrak q^2S_{\mathfrak q}$ over $S_{\mathfrak q}/\mathfrak qS_{\mathfrak q}$, so they generate $\mathfrak qS_{\mathfrak q}$ by the lemma of Nakayama, and $\operatorname{ht}\mathfrak q=\dim S_{\mathfrak q}\le n$ by the height theorem of Krull \cite[Theorem~13.5]{Mat86}.

(2) Let $P_n=E[[Y_1,\dots,Y_n]]$, a noetherian local ring whose maximal ideal $\mathfrak m$ is generated by the $Y_i$.
\begin{enumerate}
\item[(i)] For $a\in S$, $a\in\mathfrak q$ if and only if $\varphi(a)\in\mathfrak m$, since the constant term of $\varphi(a)$ is $x(a)$ and $\mathcal O$ embeds in $E$. So $\mathfrak m$ lies over $\mathfrak q$.
\item[(ii)] Let $J$ be the ideal of $P_n$ generated by $\varphi(\mathfrak q)$; by (i), $J\subseteq\mathfrak m$. Each $t_i$ lies in $\mathfrak q$, and $cY_i=\varphi(t_i)-\big(\varphi(t_i)-cY_i\big)$ lies in $J+\mathfrak m^2$; as $c$ is a unit of $E$, $\mathfrak m\subseteq J+\mathfrak m^2$, so $J=\mathfrak m$ by Nakayama's lemma.
\item[(iii)] By the dimension inequality \cite[Tag 00OM]{Sta}, for the homomorphism $\varphi$ of noetherian rings and the prime $\mathfrak m$ over $\mathfrak q$, the height of $\mathfrak m$ is at most $\operatorname{ht}\mathfrak q$ plus the height of the image of $\mathfrak m$ in $P_n/\mathfrak qP_n=P_n/J$. By (ii) that image is $0$ in the field $P_n/\mathfrak m$, of height $0$; so $\operatorname{ht}\mathfrak m\le\operatorname{ht}\mathfrak q$.
\item[(iv)] $\operatorname{ht}\mathfrak m=\dim P_n\ge n$. By induction on $n$, as $P_{n+1}\cong P_n[[Y_{n+1}]]$ and $\dim A[[Y]]\ge\dim A+1$ for every commutative ring $A$. Indeed $A$ is the quotient of $A[[Y]]$ by the ideal $(Y)$, so a chain of primes of $A$ gives a chain of primes of $A[[Y]]$ containing $Y$, of the same length. As $Y$ is a nonzerodivisor, the smallest member of this chain is not a minimal prime, so the chain extends by one prime below it.
\end{enumerate}
So $n\le\operatorname{ht}\mathfrak q$.
% (2) is ringKrullDim_mvPowerSeries_fin_field (89 lines). The proof follows it: (i) is hker (LiesOver), (ii) is hle (Nakayama, Submodule.le_of_le_smul_of_le_jacobson_bot, with the private max_eq_span), (iii) is Ideal.height_le_height_add_of_liesOver, (iv) is the private dim_ge (ringKrullDim_succ_le_ringKrullDim_powerseries, through ringKrullDim_succ_le_of_surjective for the constant term and X).

(3) By the structure theorem of Cohen \cite[Theorem~29.4]{Mat86}, $S$ is a quotient $A/Q$ of a regular, hence Cohen-Macaulay, local ring $A$, with $Q$ prime. Let $\tilde{\mathfrak q}$ be the preimage of $\mathfrak q$. By \cite[Theorem~17.4]{Mat86}, applied to $A$ and to $A_{\tilde{\mathfrak q}}$, $\operatorname{ht}\mathfrak q=\operatorname{ht}\tilde{\mathfrak q}-\operatorname{ht}Q$, $\dim S=\dim A-\operatorname{ht}Q$ and $\dim S/\mathfrak q=\dim A-\operatorname{ht}\tilde{\mathfrak q}$. As $S/\mathfrak q\cong\mathcal O$ has dimension $1$, $\dim S=\operatorname{ht}\mathfrak q+1$.
% (1) is height_ker_section_le_of_cotangent_generators (69 lines), whose proof is the one above: localization at ker x, Nakayama (le_of_le_smul_of_le_jacobson_bot) and Krull's height theorem (Ideal.height_le_card_of_mem_minimalPrimes_span_finset). (3) is ringKrullDim_eq_height_ker_section_add_one (63 lines), which the Lean proves its own way: a finite injective O[[Y_1..Y_d]] -> S compatible with x (exists_finite_injective_mvPowerSeries_compatible_section), going down over the normal ring O[[Y_1..Y_d]] (Ideal.height_eq_height_add_of_liesOver_of_hasGoingDown) and (2) for the constant term of O[[Y_1..Y_d]]; it is the case p = ker x of the published formula dim S = ht p + dim S/p, which the book cites.
\end{proof}

\begin{lemma}\label{sha-counts-vd-model-lift-chart-smooth}
Let $\mathcal O$, $E$ and pointed artinian $E$-algebras be as in Lemma~\ref{sha-counts-vd-model-lift-chart-factor}, let $R$ be a complete local noetherian $\mathcal O$-algebra with residue field $\mathbf F$, let $x:R\to\mathcal O$ be a homomorphism of $\mathcal O$-algebras, and let $I\subseteq\ker x$ be an ideal. Suppose that for every surjective homomorphism $p:B\to B'$ of pointed artinian $E$-algebras with $\pi_{B'}\circ p=\pi_B$, and every homomorphism of $\mathcal O$-algebras $f':R\to B'$ that kills $I$ and satisfies $\pi_{B'}(f'(a))=x(a)$ for all $a\in R$, there is a homomorphism of $\mathcal O$-algebras $f:R\to B$ that kills $I$ with $p\circ f=f'$.
Then there are $n\ge0$, an ideal $\mathfrak p\supseteq I$ of $R$ and a complete chart $\varphi:R\to E[[Y_1,\dots,Y_n]]$ of $(R,\mathfrak p)$ at $x$ (Definition~\ref{sha-counts-vd-model-lift-chart-complete}), and an element $y\in R$ with $x(y)\ne0$ and $y\mathfrak p\subseteq I$. Moreover $\mathfrak p$ is a minimal prime of the radical of $I$, it is the only one contained in $\ker x$, and $R/\mathfrak p$ has Krull dimension $n+1$.
\end{lemma}
\leanref{Development/AutomorphicArithmetic/AutomorphicVarieties/PadicHodgeTheory/SmoothCompl.lean}{UpperTheory.V.Loc.smoothCompl}
\begin{proof}
Write $P_n=E[[Y_1,\dots,Y_n]]$, with maximal ideal $\mathfrak m$, the series without constant term; it is generated by the $Y_i$, $P_n$ is $\mathfrak m$-adically complete, and it is a domain. For $k\ge2$, $P_n/\mathfrak m^k$ is a pointed artinian $E$-algebra, with the constant term as $\pi$, whose kernel $\mathfrak m/\mathfrak m^k$ consists of nilpotent elements.
\begin{enumerate}
\item[(i)] \emph{Parameters.} Lemma~\ref{sha-counts-vd-model-lift-chart-smooth-parameters} gives $n$, the $D_i$, the $t_i$ and $c$.
\item[(ii)] \emph{First order.} The map $f_0:R\to P_n/\mathfrak m^2$, $a\mapsto x(a)+\sum_iD_i(a)Y_i$, is a homomorphism of $\mathcal O$-algebras: it is $\mathcal O$-linear, $D_i(1)=0$ from the product rule, and $f_0(a)f_0(b)-f_0(ab)=(\sum_iD_i(a)Y_i)(\sum_iD_i(b)Y_i)\in\mathfrak m^2$. It kills $I$, as $x$ and the $D_i$ vanish on $I$.
\item[(iii)] \emph{The tower.} Suppose given $f_N:R\to P_n/\mathfrak m^{N+2}$ killing $I$ and reducing to $f_0$. Its constant terms are those of $f_0$, so $\pi(f_N(a))=x(a)$. The hypothesis, applied to the projection $P_n/\mathfrak m^{N+3}\to P_n/\mathfrak m^{N+2}$, gives $f_{N+1}$ killing $I$ that reduces to $f_N$. The compatible $f_N$ define a homomorphism of $\mathcal O$-algebras $\varphi:R\to P_n$, as $P_n$ is complete, which kills $I$, as $P_n$ is separated, and reduces to $f_0$ modulo $\mathfrak m^2$. So the constant term of $\varphi(a)$ is $x(a)$, $\varphi(t_i)-cY_i\in\mathfrak m^2$, and $\varphi$ maps $(\ker x)^k+I$ into $\mathfrak m^k$.
\item[(iv)] \emph{Injectivity up to constants.} For every $N$, if $\varphi(a)\in\mathfrak m^N$ then $ea\in(\ker x)^N+I$ for some nonzero $e\in\mathcal O$. By induction on $N$, the case $N=0$ being clear. If $\varphi(a)\in\mathfrak m^{N+1}$, write $ea=b+i'$ with $b\in(\ker x)^N$, $i'\in I$ and $e\ne0$. By Lemma~\ref{sha-counts-vd-model-lift-chart-smooth-homogeneous}, $e'b-h(t)\in(\ker x)^{N+1}+I$ for some $e'\ne0$ and $h$ homogeneous of degree $N$. Applying $\varphi$, with $\varphi(b)=e\varphi(a)\in\mathfrak m^{N+1}$, gives $h(\varphi(t_1),\dots,\varphi(t_n))\in\mathfrak m^{N+1}$, so $h=0$ by Lemma~\ref{sha-counts-vd-model-lift-chart-smooth-vanishing}. Then $e'ea=e'b+e'i'\in(\ker x)^{N+1}+I$.
\item[(v)] \emph{The prime.} Let $\mathfrak p=\ker\varphi$, a prime ideal, as $P_n$ is a domain. Then $I\subseteq\mathfrak p\subseteq\ker x$, the second by the constant terms. By (iv), for $a\in\mathfrak p$ and every $N$ there is a nonzero $e\in\mathcal O$, which does not lie in $\ker x$, with $ea\in(\ker x)^N+I$. So Lemma~\ref{sha-counts-vd-model-lift-chart-smooth-prime}, with $\mathfrak q=\ker x$, gives $y$ with $x(y)\ne0$ and $y\mathfrak p\subseteq I$.
\item[(vi)] \emph{The chart.} $\varphi$ is a complete chart of $(R,\mathfrak p)$ at $x$: $\ker\varphi=\mathfrak p$; the constant terms are given by $x$; the approximation property is Lemma~\ref{sha-counts-vd-model-lift-chart-smooth-approximation}; and the last property is (iv), as $I\subseteq\mathfrak p$.
\item[(vii)] \emph{Minimal primes.} The radical of $I$ lies in the prime $\mathfrak p$. Let $Q$ be a prime containing $I$ with $Q\subseteq\ker x$. Then $y\notin Q$, and $ya\in I\subseteq Q$ for $a\in\mathfrak p$; so $\mathfrak p\subseteq Q$. Hence $\mathfrak p$ is a minimal prime of the radical of $I$, and the only one contained in $\ker x$.
\item[(viii)] \emph{The dimension.} The ring $R/\mathfrak p$ is a domain and a complete local noetherian $\mathcal O$-algebra with residue field $\mathbf F$, and $x$ and $\varphi$ induce $\bar x:R/\mathfrak p\to\mathcal O$ and $\bar\varphi:R/\mathfrak p\to P_n$. The images of the $t_i$ satisfy the hypothesis of Lemma~\ref{sha-counts-vd-model-lift-chart-smooth-height}(1) for $\bar x$, since $(\ker x)^2+I$ maps into $(\ker\bar x)^2$, and that of (2) with $\bar\varphi$. So $\ker\bar x$ has height $n$, and $R/\mathfrak p$ has dimension $n+1$ by (3).
\end{enumerate}
\end{proof}

\begin{lemma}\label{sha-counts-vd-model-lift-chart-graded-change}
Let $\mathcal O$, $E$, $R$, $\mathfrak p$ and $x$ be as in Definition~\ref{sha-counts-vd-model-lift-chart-complete}, and let $\varphi:R\to E[[Y_1,\dots,Y_n]]$ be a complete chart of $(R,\mathfrak p)$ at $x$. For a series $f$ let $\ell(f)=(\ell_1(f),\dots,\ell_n(f))\in E^n$, where $\ell_i(f)$ is the coefficient of $Y_i$ in $f$. Let $z_1,\dots,z_n\in\ker x$ be such that $\ell(\varphi(z_1)),\dots,\ell(\varphi(z_n))$ span $E^n$.
Then there is a complete chart $\varphi'$ of $(R,\mathfrak p)$ at $x$ with $\varphi'(z_i)=Y_i$ for every $i$.
\end{lemma}
\leanref{Development/AutomorphicArithmetic/AutomorphicVarieties/PadicHodgeTheory/ChartChange.lean}{UpperTheory.V.Loc.chartChange}
\begin{proof}
Let $P=E[[Y_1,\dots,Y_n]]$, with maximal ideal $\mathfrak m$, the series with constant term $0$.
\begin{enumerate}
\item[(i)] For $f\in\mathfrak m$, $f-\sum_j\ell_j(f)Y_j\in\mathfrak m^2$. Indeed $f=\sum_jY_jg_j$ with series $g_j$; then $\ell_j(f)=g_j(0)$, and $f-\sum_jg_j(0)Y_j=\sum_jY_j\big(g_j-g_j(0)\big)$.
\item[(ii)] \emph{A substitution.} Let $y_i=\varphi(z_i)$, which lies in $\mathfrak m$, as its constant term is $x(z_i)=0$. As $P$ is complete, there is a homomorphism of $E$-algebras $s:P\to P$ with $s(Y_i)=y_i$, and $s(\mathfrak m)\subseteq\mathfrak m$. Let $J=s(\mathfrak m)P+\mathfrak m^2$. Each $y_i=s(Y_i)$ lies in $J$, so $\sum_j\ell_j(y_i)Y_j\in J$ by (i); as the $\ell(y_i)$ span $E^n$, $\sum_jv_jY_j\in J$ for every $v\in E^n$, and $\mathfrak m\subseteq J$ by (i) again. By Nakayama's lemma, $s(\mathfrak m)P=\mathfrak m$. So $P$, as a module over $P$ through $s$, is separated for the $\mathfrak m$-adic topology, and $P/s(\mathfrak m)P=E$ is generated by the image of $1$; as $P$ is $\mathfrak m$-adically complete, $P$ is generated by $1$ as a module through $s$ \cite[Theorem~8.4]{Mat86}, that is, $s$ is surjective.
\item[(iii)] \emph{$s$ is an automorphism.} A surjective endomorphism of a noetherian ring is injective. The kernels of the $s^k$ increase, so $\ker s^N=\ker s^{N+1}$ for some $N$. If $s(a)=0$, write $a=s^N(b)$; then $s^{N+1}(b)=0$, so $s^N(b)=0$, that is $a=0$.
\item[(iv)] $s$ and $s^{-1}$ preserve constant terms, so they map $\mathfrak m^N$ into $\mathfrak m^N$ for every $N$. For a series $f$, $f-f(0)$ is not a unit, so neither is $s(f-f(0))=s(f)-f(0)$, by (iii); that is, $s(f)$ has the constant term $f(0)$. The same holds for $s^{-1}$.
\item[(v)] \emph{The chart.} Let $\varphi'=s^{-1}\circ\varphi$, so that $\varphi'(z_i)=s^{-1}(y_i)=Y_i$. Its kernel is that of $\varphi$, which is $\mathfrak p$, and by (iv) the constant term of $\varphi'(a)$ is that of $\varphi(a)$, which is $x(a)$. Given $N$ and a series $f$, the approximation property of $\varphi$, for $s(f)$, gives $a\in R$ and a nonzero $c\in\mathcal O$ with $cs(f)-\varphi(a)\in\mathfrak m^N$, and applying $s^{-1}$ gives $cf-\varphi'(a)\in\mathfrak m^N$, by (iv). If $\varphi'(a)\in\mathfrak m^N$, then $\varphi(a)=s(\varphi'(a))\in\mathfrak m^N$ by (iv), and the last property of $\varphi$ gives a nonzero $c$ with $ca\in(\ker x)^N+\mathfrak p$.
\end{enumerate}
% chartChange (ChartChange.lean, 125 lines). The proof follows it: (i) is the private lin_part (with mem_span_X), (ii) is exists_equiv (co_ev for s, co_surj for its surjectivity), (iii) is inj_of_surj, (iv) is const_equiv and pow_equiv, (v) is the main proof.
\end{proof}

\begin{lemma}\label{sha-counts-vd-model-lift-chart-graded}
Let $\mathcal O$ and $E$ be as in Lemma~\ref{sha-counts-vd-model-lift-chart-factor}, let $X$ be an abelian group, and let $R$ be a complete local noetherian $\mathcal O$-algebra with residue field $\mathbf F$, with a weighted presentation with origin $x$ (Definition~\ref{sha-counts-vd-model}(1)), with parts $R_j$. Let $\mathfrak p$ be an ideal of $R$, $\varphi:R\to E[[Y_1,\dots,Y_n]]$ a complete chart of $(R,\mathfrak p)$ at $x$ (Definition~\ref{sha-counts-vd-model-lift-chart-complete}), and $y\in R$ with $x(y)\ne0$ and $y\mathfrak p=0$. For $j\in X$, let $\mathrm{Der}_j$ be the $E$-vector space of the $\mathcal O$-linear maps $\delta:R\to E$ with $\delta(fg)=x(f)\delta(g)+x(g)\delta(f)$ that vanish on the $R_i$ with $i\ne j$.
Then there are weights $u_1,\dots,u_n\in X$ and a complete chart $\varphi':R\to E[[Y_1,\dots,Y_n]]$ of $(R,\mathfrak p)$ at $x$, where $Y_i$ has the weight $u_i$, such that $\varphi'$ maps $R_j$ into the series of weight $j$, and the number of $i$ with $u_i=j$ is $\dim_E\mathrm{Der}_j$ for every $j$.
\end{lemma}
\leanref{Development/AutomorphicArithmetic/AutomorphicVarieties/PadicHodgeTheory/GradedChart.lean}{UpperTheory.V.Loc.gradedChart}
\leanref{Development/AutomorphicArithmetic/AutomorphicVarieties/PadicHodgeTheory/IsFirstOrderDeformationOfRZero.lean}{UpperTheory.V.Loc.wtDer}
\begin{proof}
Let $\pi:\mathcal O[[X_s:s\in\Sigma]]\to R$ be the presentation, with weights $w_s$, and $z_s=\pi(X_s)$, which lies in $R_{w_s}$; then $x(z_s)=0$, as $x\circ\pi$ is the constant term. For a series $h$ in the $X_s$ and $k\in X$ let $h_k$ be the sum of the terms of $h$ of weight $k$. Let $P=E[[Y_1,\dots,Y_n]]$, with maximal ideal $\mathfrak m$, and let $\ell=(\ell_1,\dots,\ell_n)$ be as in Lemma~\ref{sha-counts-vd-model-lift-chart-graded-change}. A series lies in $\mathfrak m^N$ if and only if its terms of degree less than $N$ vanish.
\begin{enumerate}
\item[(i)] \emph{Orders.} Let $\psi$ be a complete chart of $(R,\mathfrak p)$ at $x$, and $g$ a series in the $X_s$ of order at least $N$. Then $\psi(\pi(g))\in\mathfrak m^N$. Indeed $\psi\circ\pi$ maps each $X_s$ into $\mathfrak m$, as the constant term of $\psi(z_s)$ is $x(z_s)=0$, and $g$ lies in the ideal generated by the finitely many monomials of degree $N$.
\item[(ii)] \emph{Spanning generators.} There are $s_1,\dots,s_n\in\Sigma$ such that the $\ell(\varphi(z_{s_i}))$ span $E^n$. Let $V$ be the span of all the $\ell(\varphi(z_s))$, $s\in\Sigma$. The map $\ell$ kills $\mathfrak m^2$ and the constants. Let $g$ be a series in the $X_s$, with $g_s\in\mathcal O$ the coefficient of $X_s$. Then $g-g(0)-\sum_sg_sX_s$ has order at least $2$, so by (i) $\ell(\varphi(\pi(g)))=\sum_sg_s\ell(\varphi(z_s))$, which lies in $V$. For each $j$, the approximation property of $\varphi$ gives $a\in R$ and a nonzero $c\in\mathcal O$ with $cY_j-\varphi(a)\in\mathfrak m^2$; so $c$ times the $j$-th basis vector of $E^n$ is $\ell(\varphi(a))$, which lies in $V$. So $V=E^n$, and $n$ of the vectors $\ell(\varphi(z_s))$ form a basis of $E^n$.
\item[(iii)] \emph{The chart.} Let $z'_i=z_{s_i}$ and $u_i=w_{s_i}$, so that $z'_i\in R_{u_i}$, and give $Y_i$ the weight $u_i$. Lemma~\ref{sha-counts-vd-model-lift-chart-graded-change} gives a complete chart $\varphi'$ of $(R,\mathfrak p)$ at $x$ with $\varphi'(z'_i)=Y_i$. For an exponent vector $e$, the monomial $X^{\langle e\rangle}=\prod_iX_{s_i}^{e_i}$ has the weight $\sum_ie_iu_i$ of $Y^e$, and $\varphi'(\pi(X^{\langle e\rangle}))=Y^e$.
\item[(iv)] \emph{Components.} Let $f$ be a series in the $X_s$ with $\varphi'(\pi(f))\in\mathfrak m^N$. Then $\varphi'(\pi(f_k))\in\mathfrak m^N$ for every $k\in X$. By the last property of a complete chart, $c\pi(f)=b+b'$ with $c\in\mathcal O$ nonzero, $b\in(\ker x)^N$ and $b'\in\mathfrak p$. As $yb'=0$, $(x(y)-y)b'=x(y)b'$, so $x(y)^Nb'=(x(y)-y)^Nb'$, which lies in $(\ker x)^N$. So $e\pi(f)\in(\ker x)^N$, with $e=x(y)^Nc\ne0$. The ideal $\ker x$ is the image of the series with constant term $0$, so every element of $(\ker x)^N$ is the image of a series of order at least $N$, and $e\pi(f)=\pi(g)$ with $g$ of order at least $N$. Then $g-ef\in\ker\pi$, so $g_k-ef_k\in\ker\pi$, as $\pi$ is a weighted presentation. As $g_k$ has order at least $N$, $e\varphi'(\pi(f_k))=\varphi'(\pi(g_k))$ lies in $\mathfrak m^N$ by (i), and $e$ is a unit of $E$.
\item[(v)] \emph{$\varphi'$ respects the weights.} Let $a=\pi(f)$ with $f$ of weight $j$, let $Y^e$ be a monomial of weight $k\ne j$, and let $N=|e|+1$, where $|e|$ is the degree. Let $\Theta$ be a finite set of exponent vectors containing those of degree less than $N$, and choose a nonzero $d\in\mathcal O$ such that, for $e'\in\Theta$, $d$ times the coefficient of $Y^{e'}$ in $\varphi'(a)$ is an element $a_{e'}$ of $\mathcal O$. Let $F=df-\sum_{e'\in\Theta}a_{e'}X^{\langle e'\rangle}$. By (iii) the terms of $\varphi'(\pi(F))$ of degree less than $N$ vanish, so $\varphi'(\pi(F))\in\mathfrak m^N$, and $\varphi'(\pi(F_k))\in\mathfrak m^N$ by (iv). As $f$ has weight $j\ne k$, $F_k=-\sum a_{e'}X^{\langle e'\rangle}$, the sum over the $e'\in\Theta$ of weight $k$; so $\varphi'(\pi(F_k))=-\sum a_{e'}Y^{e'}$, over the same $e'$. Its coefficient of $Y^e$ is $-a_e$, and it vanishes, as $|e|<N$. So the coefficient of $Y^e$ in $\varphi'(a)$ is $0$.
\item[(vi)] \emph{Derivations.} Let $\delta:R\to E$ be $\mathcal O$-linear with $\delta(fg)=x(f)\delta(g)+x(g)\delta(f)$ and $\delta(z'_i)=0$ for every $i$. Then $\delta=0$. Indeed $\delta(1)=2\delta(1)$, so $\delta$ vanishes on $\mathcal O\cdot1$; it vanishes on $(\ker x)^2$ by the product rule; and it vanishes on $\mathfrak p$: for $b'\in\mathfrak p$, $x(b')=0$, as $\varphi'(b')=0$, and $0=\delta(yb')=x(y)\delta(b')$. Let $a\in R$, and choose a nonzero $d\in\mathcal O$ with $d\ell_i(\varphi'(a))=b_i\in\mathcal O$ for every $i$. With $F'=d\varphi'(a)$, $\varphi'\big(da-dx(a)-\sum_ib_iz'_i\big)=F'-F'(0)-\sum_i\ell_i(F')Y_i$, which lies in $\mathfrak m^2$. By the last property of a complete chart, $c\big(da-dx(a)-\sum_ib_iz'_i\big)\in(\ker x)^2+\mathfrak p$ for a nonzero $c\in\mathcal O$, so $cd\,\delta(a)=0$, and $\delta(a)=0$.
\item[(vii)] \emph{The count.} Let $j\in X$ and $\Lambda_j=\{i:u_i=j\}$. The $E$-linear map $\mathrm{Der}_j\to E^{\Lambda_j}$, $\delta\mapsto(\delta(z'_i))_{i\in\Lambda_j}$, is bijective, which gives the count. It is injective: if $\delta(z'_i)=0$ for $i\in\Lambda_j$, then also for $i\notin\Lambda_j$, as $z'_i\in R_{u_i}$ with $u_i\ne j$; so $\delta=0$ by (vi). It is surjective: for $v\in E^{\Lambda_j}$ let $\delta_v(a)=\sum_{i\in\Lambda_j}v_i\ell_i(\varphi'(a))$. It is $\mathcal O$-linear, and it satisfies the product rule at $x$, as $\ell_i(fg)=f(0)\ell_i(g)+g(0)\ell_i(f)$ for series $f,g$ and the constant term of $\varphi'(a)$ is $x(a)$. It vanishes on $R_k$ for $k\ne j$, as $\varphi'(R_k)$ consists of series of weight $k$ by (v) and $Y_i$ has the weight $j$ for $i\in\Lambda_j$. And $\delta_v(z'_i)=v_i$, as $\varphi'(z'_i)=Y_i$.
\end{enumerate}
So $\varphi'$, with the weights $u_i$, is as required.
% gradedChart (GradedChart.lean, 581 lines). The proof follows it: membership in m^N is gc_mem_pow_iff; (i) is gc_low (with origin_at_generator_eq_zero); (ii) is gc_span (gc_lin, exists_fun_fin_finrank_span_eq); (iii) is chartChange; (iv) is gc_homog (gc_sat, gc_exists_order, gc_unsmul); (v) is gc_graded (with the private finsupp_prod_X_pow_eq_monomial_mapDomain, wcomp_monomial_eq_ite_weight, gc_wcomp_ne); (vi) is gc_der; (vii) is gc_count.
\end{proof}

\begin{lemma}\label{sha-counts-vd-model-lift-chart}
Let $\mathcal O$ be a complete discrete valuation ring of characteristic $0$ that is finite over $\mathbf Z_\ell$, with fraction field $E$ and residue field $\mathbf F$. A \emph{pointed artinian $E$-algebra} is a commutative $E$-algebra $B$ of finite dimension over $E$ with a homomorphism of $E$-algebras $\pi_B:B\to E$ whose kernel consists of nilpotent elements; it has the topology of a finite-dimensional $\mathbf Q_\ell$-vector space. Let $\Gamma$ be a compact topological group, $\iota$ a finite set, $X$ a finitely generated free abelian group and $w:\iota\to X$ a function. Let $Q$ be a property of homomorphisms $\Gamma\to M_\iota(A)$, for commutative $\mathcal O$-algebras $A$, that is preserved by composition with every homomorphism of $\mathcal O$-algebras $A\to A'$ and reflected by every injective one. Let $R$ be a complete local noetherian $\mathcal O$-algebra with residue field $\mathbf F$, let $\bar r:\Gamma\to M_\iota(\mathbf F)$ be a homomorphism, and let $\rho:\Gamma\to M_\iota(R)$ be a continuous lift of $\bar r$ with the following properties.
\begin{enumerate}
\item[(a)] $\rho$ has the property $Q$, and every continuous lift of $\bar r$ with the property $Q$, with values in a complete local noetherian $\mathcal O$-algebra $R'$ with residue field $\mathbf F$, is the image of $\rho$ under a unique homomorphism of $\mathcal O$-algebras $R\to R'$.
\item[(b)] For every such $R'$, every homomorphism of $\mathcal O$-algebras $f:R\to R'$ and every torus element $t$ of $R'$ (Lemma~\ref{sha-counts-vd-model-torus-grading}) there is a homomorphism of $\mathcal O$-algebras $f':R\to R'$ with $f'(\rho(g)_{ab})=t(w_a-w_b)f(\rho(g)_{ab})$ for all $g$, $a$ and $b$.
\item[(c)] $\rho_0:\Gamma\to M_\iota(\mathcal O)$ is a continuous lift of $\bar r$ with the property $Q$, and $\rho_0(g)_{ab}=0$ whenever $w_a\ne w_b$.
\item[(d)] Let $p:B\to B'$ be a surjective homomorphism of pointed artinian $E$-algebras with $\pi_{B'}\circ p=\pi_B$, and let $r':\Gamma\to M_\iota(B')$ be a homomorphism with continuous entries and the property $Q$ such that $\pi_{B'}\circ r'=\rho_0$. Then $r'=p\circ r$ for a homomorphism $r:\Gamma\to M_\iota(B)$ with continuous entries and the property $Q$.
\end{enumerate}
Then there are a homomorphism of $\mathcal O$-algebras $x:R\to\mathcal O$ with $x\circ\rho=\rho_0$, an ideal $\mathfrak p$ of $R$, an element $y\in R$ with $x(y)\ne0$ and $y\mathfrak p=0$, and a homomorphism of $\mathcal O$-algebras $\varphi:R\to E[[Y_1,\dots,Y_n]]$, where $Y_i$ has a weight $u_i\in X$, with the following properties.
\begin{enumerate}
\item $\ker\varphi=\mathfrak p$, and the constant term of $\varphi(f)$ is $x(f)$ for every $f\in R$.
\item For every $N\ge1$ and every series $h$ there are $f\in R$ and a nonzero $c\in\mathcal O$ such that $ch-\varphi(f)$ has order at least $N$.
\item For every $N\ge1$, if $\varphi(f)$ has order at least $N$, then $cf\in(\ker x)^N+\mathfrak p$ for some nonzero $c\in\mathcal O$.
\item $\varphi(\rho(g)_{ab})$ is a series of weight $w_a-w_b$, for all $g$, $a$ and $b$.
\item For every $j\in X$, the number of $i$ with $u_i=j$ is the dimension over $E$ of the space of the functions $s:\Gamma\to M_\iota(E)$ with continuous entries and with $s(g)_{ab}=0$ whenever $w_a-w_b\ne j$, such that $g\mapsto\rho_0(g)+\varepsilon s(g)$ is a homomorphism $\Gamma\to M_\iota(E[\varepsilon])$, $\varepsilon^2=0$, with the property $Q$.
\end{enumerate}
\end{lemma}
\leanref{Development/AutomorphicArithmetic/AutomorphicVarieties/PadicHodgeTheory/LiftChart.lean}{UpperTheory.V.Loc.liftChart}
\leanref{Development/AutomorphicArithmetic/AutomorphicVarieties/CommutativeAlgebra/LiftGrading.lean}{UpperTheory.V.Alg.lift_grading}
\leanref{Development/AutomorphicArithmetic/AutomorphicVarieties/CommutativeAlgebra/DerHom.lean}{UpperTheory.V.Alg.der_hom}
\leanref{Development/AutomorphicArithmetic/AutomorphicVarieties/CommutativeAlgebra/DerExtend.lean}{UpperTheory.V.Alg.der_extend}
\begin{proof}
For a pointed artinian $E$-algebra $B$, a \emph{$B$-point} is a homomorphism $r:\Gamma\to M_\iota(B)$ with continuous entries and the property $Q$ such that $\pi_B\circ r=\rho_0$.
\begin{enumerate}
\item[(i)] \emph{Generation.} By the uniqueness in (a), two homomorphisms of $\mathcal O$-algebras from $R$ to a complete local noetherian $\mathcal O$-algebra with residue field $\mathbf F$ that agree on the entries of the $\rho(g)$ are equal. By (a), $\rho_0$ is the image of $\rho$ under a homomorphism of $\mathcal O$-algebras $x:R\to\mathcal O$.
\item[(ii)] \emph{Points.} Let $r$ be a $B$-point. Then $r=f\circ\rho$ for a homomorphism of $\mathcal O$-algebras $f:R\to B$. Indeed, $\pi_B(r(g)_{ab})=x(\rho(g)_{ab})$, so by Lemma~\ref{sha-counts-vd-model-lift-chart-factor} there are a subalgebra $C\subseteq B$ that is a complete local noetherian $\mathcal O$-algebra with residue field $\mathbf F$ and a continuous lift $\rho'$ of $\bar r$ with values in $C$ and $\rho'=r$. As the inclusion $C\to B$ is injective, $\rho'$ has the property $Q$; by (a), $\rho'=f_C\circ\rho$ for some $f_C:R\to C$, and $f$ is $f_C$ followed by the inclusion.
\item[(iii)] \emph{Continuity.} For every homomorphism of $\mathcal O$-algebras $f$ from $R$ to a commutative $\mathcal O$-algebra that is a finite-dimensional $\mathbf Q_\ell$-algebra, the entries of $f\circ\rho$ are continuous, by Lemma~\ref{sha-counts-vd-model-lift-chart-continuity}(1).
\item[(iv)] \emph{A complete chart.} Lemma~\ref{sha-counts-vd-model-lift-chart-smooth} applies to $R$, $x$ and $I=0$. Let $p:B\to B'$ be a surjective homomorphism of pointed artinian $E$-algebras with $\pi_{B'}\circ p=\pi_B$, and $f':R\to B'$ a homomorphism of $\mathcal O$-algebras with $\pi_{B'}\circ f'=x$. Then $f'\circ\rho$ has the property $Q$, which is preserved by $f'$, and continuous entries by (iii), and $\pi_{B'}\circ f'\circ\rho=x\circ\rho=\rho_0$. By (d), $f'\circ\rho=p\circ r$ for a homomorphism $r$ with continuous entries and the property $Q$; then $\pi_B\circ r=\pi_{B'}\circ p\circ r=\rho_0$, so $r$ is a $B$-point, and $r=f\circ\rho$ by (ii). The homomorphisms $p\circ f$ and $f'$ agree on the entries of the $\rho(g)$, so they are equal by Lemma~\ref{sha-counts-vd-model-lift-chart-generation}. So Lemma~\ref{sha-counts-vd-model-lift-chart-smooth} gives an ideal $\mathfrak p$, a complete chart $\varphi_0$ of $(R,\mathfrak p)$ at $x$, and $y\in R$ with $x(y)\ne0$ and $y\mathfrak p=0$.
\item[(v)] \emph{The weights.} By (i), (b), and since $\rho_0(g)_{ab}\in\mathcal O$ has the residue $\bar r(g)_{ab}$ and vanishes when $w_a\ne w_b$ by (c), Lemma~\ref{sha-counts-vd-model-torus-grading}(2), with $c(g,a,b)=\rho_0(g)_{ab}$, gives a weighted presentation of $R$ whose variables map to elements $z_s=\rho(g)_{ab}-\rho_0(g)_{ab}$ of weight $w_a-w_b$, with $\rho(g)_{ab}\in R_{w_a-w_b}$ for all $g$, $a$ and $b$. As $x(z_s)=0$, the presentation has the origin $x$, as in step (i) of the proof of Lemma~\ref{sha-counts-vd-model-universal}(2).
\item[(vi)] \emph{The chart.} Lemma~\ref{sha-counts-vd-model-lift-chart-graded}, applied to this presentation, $\mathfrak p$, $\varphi_0$ and $y$, gives weights $u_i$ and a complete chart $\varphi$ of $(R,\mathfrak p)$ at $x$ that maps $R_j$ into the series of weight $j$, with $\dim_E\mathrm{Der}_j$ variables of weight $j$. This gives (1), (2) and (3); and (4), as $\rho(g)_{ab}\in R_{w_a-w_b}$.
\item[(vii)] \emph{The count.} Let $V_j$ be the space of the functions $s$ of (5). For $\delta\in\mathrm{Der}_j$ let $s_\delta(g)_{ab}=\delta(\rho(g)_{ab})$. Since $\delta(1)=2\delta(1)$, $\delta$ vanishes on the image of $\mathcal O$, and $\delta(z_s)=\delta(\rho(g)_{ab})$ for $z_s=\rho(g)_{ab}-\rho_0(g)_{ab}$. The linear map $\delta\mapsto s_\delta$ is a bijection from $\mathrm{Der}_j$ onto $V_j$, which gives (5).
\begin{itemize}
\item It is injective: if $s_\delta=0$, then $\delta$ vanishes at every $z_s$, so $\delta=0$ by Lemma~\ref{sha-counts-vd-presentations}(f).
\item It maps into $V_j$. The map $F:f\mapsto x(f)+\varepsilon\delta(f)$ is a homomorphism of $\mathcal O$-algebras $R\to E[\varepsilon]$, by the product rule for $\delta$. So $F\circ\rho=\rho_0+\varepsilon s_\delta$ has the property $Q$; $s_\delta$ has continuous entries by (iii), applied to $F$, composed with the coefficient of $\varepsilon$; and $s_\delta(g)_{ab}=0$ when $w_a-w_b\ne j$, since then $\rho(g)_{ab}\in R_{w_a-w_b}$ and $\delta$ vanishes there.
\item It is onto $V_j$. Let $s\in V_j$. With $\pi(a+\varepsilon b)=a$, $E[\varepsilon]$ is a pointed artinian $E$-algebra, and $\rho_0+\varepsilon s$ is an $E[\varepsilon]$-point, its entries being continuous since both coordinates are. By (ii) it is $F\circ\rho$ for a homomorphism $F:R\to E[\varepsilon]$. The first coordinate of $F$ and $x$ agree on the entries of the $\rho(g)$, so they are equal by Lemma~\ref{sha-counts-vd-model-lift-chart-generation}; hence $F=x+\varepsilon\delta$ with $\delta$ satisfying the product rule at $x$, and $s=s_\delta$. Finally $\delta$ has weight $j$. Let $\pi$ be the presentation, with variables $X_s$, and let $\delta'(h)=\delta(\pi(h))-\sum_sh_s\delta(z_s)$, where $h_s$ is the coefficient of $X_s$ in $h$. Then $\delta'$ satisfies the product rule at the constant term and vanishes at the variables, so $\delta'=0$ by Lemma~\ref{sha-counts-vd-presentations}(f), applied to the power series ring. If $h$ has weight $i\ne j$, then $h_s=0$ unless $X_s$ has weight $i$, and $\delta(z_s)=s(g)_{ab}=0$ when $X_s$ has the weight $w_a-w_b\ne j$; so $\delta(\pi(h))=0$.
\end{itemize}
\end{enumerate}
% Written out after their statements: gradedChart (GradedChart.lean, 581 lines, with Lemma -lift-chart-graded-change for chartChange) and smoothCompl (SmoothCompl.lean, 174 lines: sc_first in (ii), sc_tower in (iii), sc_inj in (iv), sc_end in (v) to (vii) and sc_dim in (viii)), with its sub-lemmas -lift-chart-smooth-parameters, -homogeneous, -vanishing, -approximation, -prime and -height. The proofs of Lemmas -lift-chart-factor (exists_lift_factor_compact_subalgebra, with Lemma -lift-chart-factor-integral), -lift-chart-continuity (ptCont) and -lift-chart-generation (homExt, with Lemma -lift-chart-generation-span) are written out after their statements. Written out here: liftChart itself, with origin_eq in (v), der_hom in (vii) and the private hasWeight_of_generators in the last item of (vii).
\end{proof}

\subsection{Local lifts and framed cohomology}\label{subsection-sha-local-lifts}

\begin{lemma}\label{sha-counts-vd-model-local-lifts-flag}
Let $D$ be a Siegel datum, with the notation of Definitions~\ref{sha-siegel-datum} and~\ref{sha-counts-vd-model-framing}, let $E$ be the fraction field of $\mathcal O$, let $B$ be a pointed artinian $E$-algebra (Lemma~\ref{sha-counts-vd-model-lift-chart-factor}), and let $r:D_\ell\to M_4(B)$ be a homomorphism with continuous entries and $\pi_B\circ r=r_0$ on $D_\ell$.
Then $r$ satisfies the conditions at $\ell$ of Definition~\ref{sha-counts-vd-model-framing}(1) if and only if every $r(g)$ is symplectic with multiplier $\theta(g)$ and there are an invertible $h\in M_4(B)$, symplectic with multiplier $1$, with $\pi_B(h)=1$, and homomorphisms $\kappa_1,\dots,\kappa_4:D_\ell\to B^\times$, such that for every $g\in D_\ell$ the matrix $h^{-1}r(g)h$ is upper triangular with diagonal $(\kappa_1(g),\dots,\kappa_4(g))$, and $\kappa_i(s)=\delta_i(s)$ for $s\in I_\ell$.
\end{lemma}
\leanref{Development/AutomorphicArithmetic/AutomorphicVarieties/PadicHodgeTheory/SpFlag.lean}{UpperTheory.V.Loc.spFlag}
\leanref{Development/AutomorphicArithmetic/AutomorphicVarieties/PadicHodgeTheory/IsFirstOrderDeformationOfRZero.lean}{UpperTheory.V.Loc.FlagIn}
\begin{proof}
Write $\pi=\pi_B$, $c_i=\delta_i(s_0)$, viewed in $B$, and $d_i=\prod_{j\ne i}(c_i-c_j)$. An element $a\in B$ with $\pi(a)\ne0$ is a unit, as $a-\pi(a)$ is nilpotent; so a matrix over $B$ whose image under $\pi$ is invertible is invertible. The $c_i$ are units, and so are the $c_i-c_j$ for $i\ne j$ and the $d_i$, as the $\delta_i(s_0)$ are pairwise distinct (step (i) of the proof of Lemma~\ref{sha-counts-vd-model-local-tangent-conditions} at $\ell$); $c_{5-i}c_i=\theta(s_0)$; and $\pi(r(s_0))=\operatorname{diag}(\pi(c_1),\dots,\pi(c_4))$. Let $e_i=\prod_{j\ne i}(r(s_0)-c_j)$, as in Definition~\ref{sha-siegel-datum}(2), and let $E_{ii}$ be the matrix with entry $1$ at $(i,i)$ and $0$ elsewhere.
\begin{enumerate}
\item[(i)] \emph{Sandwiches.} Let $k\in M_4(B)$ be invertible with $r(s_0)k=k\operatorname{diag}(c_1,\dots,c_4)$. Then $e_i=k\,(d_iE_{ii})\,k^{-1}$, since $P(r(s_0))k=k\operatorname{diag}(P(c_1),\dots,P(c_4))$ for every polynomial $P$, and $\prod_{j\ne i}(c_a-c_j)$ is $d_i$ if $a=i$ and $0$ otherwise. So, for every $M\in M_4(B)$, $e_jMe_i=0$ if and only if $(k^{-1}Mk)_{ji}=0$.
\item[(ii)] \emph{Eigenvectors.} Let $X\in M_4(B)$ have the characteristic polynomial $\prod_j(T-c_j)$, with $\pi(X)=\operatorname{diag}(\pi(c_1),\dots,\pi(c_4))$. Then there is an invertible $k$ with $Xk=k\operatorname{diag}(c_1,\dots,c_4)$ and $\pi(k)=1$, which is upper triangular if $X$ is. Indeed, let $P_b=\prod_{j\ne b}(X-c_j)$; by the theorem of Cayley and Hamilton, $(X-c_b)P_b=0$. Let $k$ be the matrix whose column $b$ is $d_b^{-1}$ times the column $b$ of $P_b$. Then $Xk=k\operatorname{diag}(c)$; $\pi(P_b)=\pi(d_b)E_{bb}$, so $\pi(k)=1$ and $k$ is invertible; and if $X$ is upper triangular, so are the $P_b$ and $k$.
\item[(iii)] \emph{A symplectic change of basis.} Suppose that $r(s_0)$ is symplectic with multiplier $\theta(s_0)$, and let $k$ be as in (ii) for $X=r(s_0)$. Then there is an invertible $h$ with $r(s_0)h=h\operatorname{diag}(c)$, $\pi(h)=1$ and $h^*h=1$. Indeed, let $N=k^*k$. Applying ${}^*$ to $r(s_0)k=k\operatorname{diag}(c)$ gives $k^*r(s_0)^*=\operatorname{diag}(c_4,\dots,c_1)k^*$, so $\theta(s_0)N=\operatorname{diag}(c_4,\dots,c_1)N\operatorname{diag}(c)$; with $\theta(s_0)=c_{5-a}c_a$ this gives $c_{5-a}(c_a-c_b)N_{ab}=0$, so $N$ is diagonal. As $N^*=N$, $N_{5-a,5-a}=N_{aa}$, and as $\pi(N)=1$, the $N_{aa}$ are units. Let $h=k\operatorname{diag}(N_{11}^{-1},N_{22}^{-1},1,1)$. Then $h^*h=\operatorname{diag}(1,1,N_{22}^{-1},N_{11}^{-1})N\operatorname{diag}(N_{11}^{-1},N_{22}^{-1},1,1)=1$, $\pi(h)=1$, and $r(s_0)h=h\operatorname{diag}(c)$.
\end{enumerate}
Suppose that $r$ satisfies the conditions at $\ell$. Every $r(g)$ is symplectic with multiplier $\theta(g)$. By (ii) and (iii), applied with the condition on the characteristic polynomial of $r(s_0)$, there is $h$ as in (iii). By (i), the condition $e_jr(g)e_i=0$ for $i<j$ says that $h^{-1}r(g)h$ is upper triangular. Let $\kappa_i(g)=(h^{-1}r(g)h)_{ii}$; it is multiplicative, as the diagonal entries of a product of upper triangular matrices are the products of theirs, so $\kappa_i$ is a homomorphism to $B^\times$. By (i), the condition $e_i(r(s)-\delta_i(s))e_i=0$ says that $\kappa_i(s)=\delta_i(s)$ for $s\in I_\ell$.

Conversely, let $h$ and the $\kappa_i$ be as in the statement, and $T(g)=h^{-1}r(g)h$. As $s_0\in I_\ell$, $T(s_0)$ is upper triangular with diagonal $(c_1,\dots,c_4)$; so the characteristic polynomial of $r(s_0)$, which is that of $T(s_0)$, is $\prod_j(T-c_j)$, and $\pi(T(s_0))=\operatorname{diag}(\pi(c))$, as $\pi(h)=1$. By (ii), $T(s_0)k=k\operatorname{diag}(c)$ for an upper triangular invertible $k$, and then $r(s_0)hk=hk\operatorname{diag}(c)$. By (i), applied to $hk$, $e_jr(g)e_i=0$ if and only if $(k^{-1}T(g)k)_{ji}=0$, which holds for $j>i$, as $k^{-1}T(g)k$ is upper triangular; and $e_i(r(s)-\delta_i(s))e_i=0$ if and only if $(k^{-1}T(s)k)_{ii}=\delta_i(s)$, which holds for $s\in I_\ell$, as $(k^{-1}T(s)k)_{ii}=\kappa_i(s)$.
\end{proof}

\begin{lemma}\label{sha-counts-vd-model-local-lifts-symplectic}
Let $E$ be a field of characteristic $0$, let $B$ and $B'$ be commutative $E$-algebras with homomorphisms of $E$-algebras $\pi_B:B\to E$ and $\pi_{B'}:B'\to E$ whose kernels consist of nilpotent elements, and let $p:B\to B'$ be a surjective homomorphism of $E$-algebras with $\pi_{B'}\circ p=\pi_B$. Let $h'\in M_4(B')$ be invertible and symplectic with multiplier $1$, with $\pi_{B'}(h')=1$.
Then there is an invertible $h\in M_4(B)$, symplectic with multiplier $1$, with $\pi_B(h)=1$ and $p(h)=h'$.
\end{lemma}
\leanref{Development/AutomorphicArithmetic/AutomorphicVarieties/PadicHodgeTheory/LiftSymplecticSimilitudeNilpotent.lean}{UpperTheory.lift_symplectic_similitude_nilpotent}
\begin{proof}
A matrix over $B$ or $B'$ whose image under $\pi_B$ or $\pi_{B'}$ is invertible is invertible, as an element $a$ with $\pi(a)\ne0$ is a unit, $a-\pi(a)$ being nilpotent. Recall that $(YZ)^*=Z^*Y^*$ and $(Y^*)^*=Y$. Write $g=h'$, so that $g^*g=1$ and $\pi_{B'}(g)=1$.
\begin{enumerate}
\item[(i)] $1+g$ is invertible, as $\pi_{B'}(1+g)=2$ and $E$ has characteristic $0$. Let $Y'=(1+g)^{-1}(g-1)$; as $1+g$ and $g-1$ commute, also $Y'(1+g)=g-1$. Then $g(1-Y')=1+Y'$, since $g(1-Y')-(1+Y')=g-1-(1+g)Y'=0$; and $\pi_{B'}(Y')=0$. Moreover $Y'^*=-Y'$: $(1+g)^*(Y'^*+Y')(1+g)=(Y'(1+g))^*(1+g)+(1+g)^*Y'(1+g)=(g^*-1)(1+g)+(1+g^*)(g-1)=2(g^*g-1)=0$, and $1+g$ and $(1+g)^*$ are invertible.
\item[(ii)] As $p$ is surjective, $Y'=p(Z)$ for some $Z\in M_4(B)$. Let $Y=\frac12(Z-Z^*)$. Then $Y^*=-Y$, $p(Y)=\frac12(Y'-Y'^*)=Y'$, and $\pi_B(Y)=\pi_{B'}(Y')=0$.
\item[(iii)] So $1-Y$ and $1+Y$ are invertible. Let $h=(1+Y)(1-Y)^{-1}$. Then $\pi_B(h)=1$, and $p(h)(1-Y')=1+Y'=g(1-Y')$, so $p(h)=g$, as $1-Y'$ is invertible. Finally $h^*h=1$: as $(1-Y)^*=1+Y$ and $(1+Y)^*=1-Y$, $(1+Y)(h^*h-1)(1-Y)=(h(1-Y))^*(h(1-Y))-(1+Y)(1-Y)=(1-Y)(1+Y)-(1+Y)(1-Y)=0$.
\end{enumerate}
\end{proof}

\begin{lemma}\label{sha-counts-vd-model-local-lifts-unramified}
Let $v$ and $\ell$ be primes, let $B$ and $B'$ be commutative finite-dimensional $\mathbf Q_\ell$-algebras, with their topologies as $\mathbf Q_\ell$-vector spaces, and let $p:B\to B'$ be a surjective homomorphism of $\mathbf Q_\ell$-algebras whose kernel consists of nilpotent elements. Let $u':\operatorname{Gal}(\bar{\mathbf Q}_v/\mathbf Q_v)\to B'^\times$ be a continuous homomorphism, trivial on the inertia group, such that $u'(g)-1$ is nilpotent for every $g$.
Then there is a continuous homomorphism $u:\operatorname{Gal}(\bar{\mathbf Q}_v/\mathbf Q_v)\to B^\times$, trivial on the inertia group, with $p\circ u=u'$.
\end{lemma}
\leanref{Development/AutomorphicArithmetic/AutomorphicVarieties/PadicHodgeTheory/LiftUnramifiedUnipotentCharacter.lean}{UpperTheory.lift_unramified_unipotent_character}
\begin{proof}
Let $\sigma$ be a Frobenius element of $\operatorname{Gal}(\bar{\mathbf Q}_v/\mathbf Q_v)$. For a nilpotent element $a$ of a commutative $\mathbf Q$-algebra, $\exp(a)=\sum_{n\ge0}a^n/n!$ is a finite sum, $\exp(a+b)=\exp(a)\exp(b)$ for nilpotent $a$ and $b$, and every element $1+a$ with $a$ nilpotent is $\exp(\log(1+a))$, with $\log(1+a)=\sum_{n\ge1}(-1)^{n+1}a^n/n$ nilpotent.
\begin{enumerate}
\item[(i)] As $u'(\sigma)-1$ is nilpotent, $u'(\sigma)=\exp(x')$ for a nilpotent $x'\in B'$. Choose $x\in B$ with $p(x)=x'$; it is nilpotent, as $x^k\in\ker p$ for some $k$, and $\ker p$ consists of nilpotent elements.
\item[(ii)] By Lemma~\ref{sha-counts-vd-model-local-tangent-unramified-cocycle}, with $A=V=\mathbf Q_\ell$ and the trivial action, there is a continuous homomorphism $c:\operatorname{Gal}(\bar{\mathbf Q}_v/\mathbf Q_v)\to\mathbf Q_\ell$, trivial on the inertia group, with $c(\sigma)=1$.
\item[(iii)] Let $u(g)=\exp(c(g)x)$. It is a homomorphism to $B^\times$, as $c$ is additive; it is continuous, being a polynomial in $c(g)$ with coefficients in $B$; and it is trivial on the inertia group.
\item[(iv)] The homomorphisms $p\circ u$ and $u'$ agree on the inertia group, where both are $1$, and at $\sigma$, where $p(u(\sigma))=\exp(x')=u'(\sigma)$. The set where they agree is a closed subgroup, as $B'$ is Hausdorff, so it is everything, by fact (a) in the proof of Lemma~\ref{sha-counts-vd-model-local-tangent-unramified-cocycle}.
\end{enumerate}
\end{proof}

\begin{lemma}\label{sha-counts-vd-model-local-lifts-borel-cocycle}
Let $v$ and $\ell$ be primes, let $H$ be a topological group with an open embedding $\phi:H\to\operatorname{Gal}(\bar{\mathbf Q}_v/\mathbf Q_v)$ that is a homomorphism, let $B$ and $B'$ be commutative finite-dimensional $\mathbf Q_\ell$-algebras, with their topologies as $\mathbf Q_\ell$-vector spaces, and let $p:B\to B'$ be a surjective homomorphism of $\mathbf Q_\ell$-algebras. Let $\chi:H\to B$ be a continuous multiplicative homomorphism such that $\chi(g)-\chi_{\mathrm{cyc}}(\phi(g))$ is a unit of $B$ for some $g\in H$, and let $f:H\times H\to B$ be continuous, with $f(g,h)+f(gh,k)=\chi(g)f(h,k)+f(g,hk)$ for all $g,h,k\in H$. Let $x':H\to B'$ be continuous, with $x'(gh)=x'(g)+p(\chi(g))x'(h)+p(f(g,h))$ for all $g,h$.
Then there is a continuous $x:H\to B$ with $x(gh)=x(g)+\chi(g)x(h)+f(g,h)$ for all $g,h$, and $p\circ x=x'$.
\end{lemma}
\leanref{Development/AutomorphicArithmetic/AutomorphicVarieties/PadicHodgeTheory/SpBorLift.lean}{UpperTheory.V.Loc.cocLift}
\begin{proof}
The image of $\phi$ is an open subgroup of $\operatorname{Gal}(\bar{\mathbf Q}_v/\mathbf Q_v)$, the Galois group of a finite extension of $\mathbf Q_v$. Let $J=\ker p$, an ideal of $B$, on which $H$ acts continuously and $\mathbf Q_\ell$-linearly by $g\cdot b=\chi(g)b$.
\begin{enumerate}
\item[(i)] $H^2(H,J)=0$, for continuous cochains. By local Tate duality \cite[Chapter~VII, \S2]{NSW08}, $H^2(H,J)$ has the dimension of the space of the $\mathbf Q_\ell$-linear maps $F:J\to\mathbf Q_\ell$ with $F(\chi(g)b)=\chi_{\mathrm{cyc}}(\phi(g))F(b)$ for all $g$ and $b$. Such an $F$ vanishes: with $g$ as in the statement and $w=\chi(g)-\chi_{\mathrm{cyc}}(\phi(g))$, every $a\in J$ is $wb$ with $b=w^{-1}a\in J$, and $F(a)=F(\chi(g)b)-\chi_{\mathrm{cyc}}(\phi(g))F(b)=0$.
\item[(ii)] Let $s:B'\to B$ be a $\mathbf Q_\ell$-linear map with $p\circ s=\mathrm{id}$, which is continuous, and $X=s\circ x'$. Let $z(g,h)=X(gh)-X(g)-\chi(g)X(h)-f(g,h)$. Applying $p$ shows that $z$ takes values in $J$; it is continuous; and it is a $2$-cocycle for the action on $J$, being minus the coboundary of $X$ minus $f$, whose coboundary vanishes by the identity for $f$.
\end{enumerate}
By (i), $z(g,h)=\chi(g)y(h)-y(gh)+y(g)$ for a continuous $y:H\to J$. Then $x=X+y$ satisfies $x(gh)=x(g)+\chi(g)x(h)+f(g,h)$, and $p\circ x=x'$, as $p\circ y=0$.
\end{proof}

\begin{lemma}\label{sha-counts-vd-model-local-lifts-borel}
Let $v$ and $\ell$ be primes, let $H$ be a topological group with an open embedding $\phi:H\to\operatorname{Gal}(\bar{\mathbf Q}_v/\mathbf Q_v)$ that is a homomorphism, let $B$ and $B'$ be commutative finite-dimensional $\mathbf Q_\ell$-algebras, with their topologies as $\mathbf Q_\ell$-vector spaces, and let $p:B\to B'$ be a surjective homomorphism of $\mathbf Q_\ell$-algebras. Let $\alpha,\beta,u:H\to B$ be continuous multiplicative homomorphisms such that for each $\chi\in\{\alpha,\beta,\alpha\beta,\alpha^2\beta\}$ there is $g\in H$ with $\chi(g)-\chi_{\mathrm{cyc}}(\phi(g))$ a unit of $B$. Let $T':H\to M_4(B')$ be a homomorphism with continuous entries such that each $T'(g)$ is upper triangular with diagonal $p\big(\alpha(g)^2\beta(g)u(g),\alpha(g)\beta(g)u(g),\alpha(g)u(g),u(g)\big)$ and symplectic with some multiplier.
Then there is a homomorphism $T:H\to M_4(B)$ with continuous entries such that each $T(g)$ is upper triangular with diagonal $\big(\alpha(g)^2\beta(g)u(g),\alpha(g)\beta(g)u(g),\alpha(g)u(g),u(g)\big)$ and symplectic with multiplier $\alpha(g)^2\beta(g)u(g)^2$, and $p\circ T=T'$.
\end{lemma}
\leanref{Development/AutomorphicArithmetic/AutomorphicVarieties/PadicHodgeTheory/SpBorLift.lean}{UpperTheory.V.Loc.borLift}
\leanref{Development/AutomorphicArithmetic/AutomorphicVarieties/PadicHodgeTheory/LocalPublishedTheorems.lean}{UpperTheory.V.Loc.LocalPublishedTheorems}
\begin{proof}
For elements $\alpha,\beta,u,a,b,c,d$ of a commutative ring, let
\[
 T(\alpha,\beta,u;a,b,c,d)=\begin{pmatrix}\alpha^2\beta u&a\alpha\beta u&c\alpha u&du\\0&\alpha\beta u&b\alpha u&(c-ab)u\\0&0&\alpha u&-au\\0&0&0&u\end{pmatrix}.
\]
Three computations are used.
\begin{enumerate}
\item[(i)] $T(\alpha,\beta,u;a,b,c,d)\,T(\alpha',\beta',u';a',b',c',d')=T(\alpha\alpha',\beta\beta',uu';a'',b'',c'',d'')$ with $a''=a+\alpha a'$, $b''=b+\beta b'$, $c''=c+\alpha\beta c'+a\beta b'$ and $d''=d+\alpha^2\beta d'+a\alpha\beta(c'-a'b')-c\alpha a'$.
\item[(ii)] $T(\alpha,\beta,u;a,b,c,d)$ is symplectic with multiplier $\alpha^2\beta u^2$.
\item[(iii)] Let $\alpha$, $\beta$ and $u$ be units, and let $T$ be upper triangular with diagonal $(\alpha^2\beta u,\alpha\beta u,\alpha u,u)$ and symplectic with some multiplier. Then $T=T(\alpha,\beta,u;a,b,c,d)$ with $a=T_{12}/(\alpha\beta u)$, $b=T_{23}/(\alpha u)$, $c=T_{13}/(\alpha u)$ and $d=T_{14}/u$: the entries $(1,2)$ and $(1,3)$ of $T^*T$ vanish, which gives $T_{34}$ and $T_{24}$. And $T(\alpha,\beta,u;a,b,c,d)$ determines $a$, $b$, $c$ and $d$.
\end{enumerate}
The values $p(\alpha(g))$, $p(\beta(g))$ and $p(u(g))$ are units, with inverses at $g^{-1}$. By (iii), $T'(g)=T(p(\alpha(g)),p(\beta(g)),p(u(g));a'(g),b'(g),c'(g),d'(g))$ for continuous functions $a',b',c',d':H\to B'$. As $T'$ is a homomorphism, (i) and the uniqueness in (iii) give
\begin{align*}
 a'(gh)&=a'(g)+p(\alpha(g))a'(h),\qquad b'(gh)=b'(g)+p(\beta(g))b'(h),\\
 c'(gh)&=c'(g)+p(\alpha\beta(g))c'(h)+a'(g)p(\beta(g))b'(h),\\
 d'(gh)&=d'(g)+p(\alpha^2\beta(g))d'(h)+a'(g)p(\alpha\beta(g))\big(c'(h)-a'(h)b'(h)\big)-c'(g)p(\alpha(g))a'(h).
\end{align*}
Lemma~\ref{sha-counts-vd-model-local-lifts-borel-cocycle}, applied in turn with $(\chi,f)$ equal to $(\alpha,0)$, $(\beta,0)$, $(\alpha\beta,(g,h)\mapsto a(g)\beta(g)b(h))$ and $(\alpha^2\beta,(g,h)\mapsto a(g)\alpha(g)\beta(g)(c(h)-a(h)b(h))-c(g)\alpha(g)a(h))$, gives continuous $a,b,c,d:H\to B$ lifting $a',b',c',d'$ and satisfying the same identities with $\alpha$, $\beta$ in place of $p(\alpha)$, $p(\beta)$. The hypothesis on $\chi$ there is the hypothesis on $\alpha$, $\beta$, $\alpha\beta$ and $\alpha^2\beta$, and the identity for $f$ is a computation from the identities for $a$, $b$ and $c$ already obtained. Then $a$, $b$, $c$ and $d$ vanish at $1$, and $T(g)=T(\alpha(g),\beta(g),u(g);a(g),b(g),c(g),d(g))$ is a homomorphism by (i), with continuous entries, upper triangular with the required diagonal, symplectic with multiplier $\alpha(g)^2\beta(g)u(g)^2$ by (ii), and $p\circ T=T'$.
\end{proof}

\begin{lemma}\label{sha-counts-vd-model-local-lifts-commuting-diagonal}
Let $E$ be an infinite field, let $B'$ be a commutative $E$-algebra with a homomorphism of $E$-algebras $\pi':B'\to E$ whose kernel consists of nilpotent elements, and let $H$ be a group. Let $\chi_1,\dots,\chi_4:H\to E^\times$ be pairwise distinct homomorphisms and $\theta:H\to E$ a function with $\chi_{5-i}\chi_i=\theta$ for every $i$. Let $r':H\to M_4(B')$ be a homomorphism such that any two of the $r'(g)$ commute, $\pi'(r'(g))=\operatorname{diag}(\chi_1(g),\dots,\chi_4(g))$, and $r'(g)$ is symplectic with multiplier $\theta(g)$, for every $g$.
Then there is $u'\in M_4(B')$, symplectic with multiplier $1$, with $\pi'(u')=1$, such that $r'(g)u'=u'\operatorname{diag}(d(g))$ for every $g$, for some $d(g)\in B'^4$.
\end{lemma}
\leanref{Development/AutomorphicArithmetic/AutomorphicVarieties/PadicHodgeTheory/SpLiftP.lean}{diagAll}
\leanref{Development/AutomorphicArithmetic/AutomorphicVarieties/PadicHodgeTheory/SpLiftP.lean}{diagX}
\begin{proof}
An element $a\in B'$ with $\pi'(a)\ne0$ is a unit, as $a-\pi'(a)$ is nilpotent; so a matrix whose image under $\pi'$ is invertible is invertible.
\begin{enumerate}
\item[(i)] \emph{A generic element.} On the $E$-vector space of finitely supported functions $x:H\to E$, the six linear forms $x\mapsto\sum_gx(g)(\chi_i(g)-\chi_j(g))$, $i\ne j$, are nonzero. As $E$ is infinite, a vector space over $E$ is not a finite union of proper subspaces, so there is $x$ at which none of them vanishes. Let $X=\sum_gx(g)r'(g)$ and $c_i=\sum_gx(g)\chi_i(g)$. Then the $c_i$ are pairwise distinct, $\pi'(X)=\operatorname{diag}(c_1,\dots,c_4)$, and $X$ commutes with every $r'(g)$.
\item[(ii)] \emph{Eigenvectors of $X$.} Let $P$ be the characteristic polynomial of $X$; its image under $\pi'$ is $\prod_j(T-c_j)$, of which $c_i$ is a simple root. Since $\ker\pi'$ consists of nilpotent elements, there is $t_i\in B'$ with $P(t_i)=0$ and $\pi'(t_i)=c_i$ (a simple root lifts uniquely modulo a nil ideal, as in the lemma of Hensel). Let $d_i=\prod_{j\ne i}(c_i-c_j)\ne0$, and let $u$ be the matrix whose column $i$ is $d_i^{-1}$ times the column $i$ of the adjugate of $t_i-X$. As $(t_i-X)\operatorname{adj}(t_i-X)=P(t_i)=0$, $Xu=u\operatorname{diag}(t_1,\dots,t_4)$. The image of $\operatorname{adj}(t_i-X)$ under $\pi'$ is $\operatorname{adj}(\operatorname{diag}(c_i-c_1,\dots,c_i-c_4))$, whose column $i$ is $d_i$ times the column $i$ of $1$; so $\pi'(u)=1$, and $u$ is invertible.
\item[(iii)] \emph{Diagonalization.} For every $g$, $u^{-1}r'(g)u$ commutes with $\operatorname{diag}(t)=u^{-1}Xu$, and the $t_i-t_j$, $i\ne j$, are units, as $\pi'(t_i-t_j)=c_i-c_j\ne0$. So $u^{-1}r'(g)u$ is diagonal: $r'(g)u=u\operatorname{diag}(d(g))$.
\item[(iv)] \emph{The form.} Let $S=u^*u$. It is diagonal: for $i\ne j$ choose $g$ with $\chi_i(g)\ne\chi_j(g)$; from $r'(g)^*r'(g)=\theta(g)$ and $r'(g)u=uD$, $D=\operatorname{diag}(d(g))$, we get $D^*SD=\theta(g)S$, so $(d_{5-i}(g)d_j(g)-\theta(g))S_{ij}=0$, and $\pi'(d_{5-i}(g)d_j(g)-\theta(g))=\chi_{5-i}(g)(\chi_j(g)-\chi_i(g))\ne0$. So $S=\operatorname{diag}(s_1,\dots,s_4)$, with $\pi'(s_i)=1$, and $s_{5-i}=s_i$, as $S^*=S$.
\end{enumerate}
Let $u'=u\operatorname{diag}(1,1,s_2^{-1},s_1^{-1})$. Then $u'^*u'=\operatorname{diag}(s_1^{-1},s_2^{-1},1,1)\,S\,\operatorname{diag}(1,1,s_2^{-1},s_1^{-1})=1$, $\pi'(u')=1$, and $r'(g)u'=u'\operatorname{diag}(d(g))$, as diagonal matrices commute.
\end{proof}

\begin{lemma}\label{sha-counts-vd-model-local-lifts-commuting-characters}
Let $B$ and $B'$ be commutative $\mathbf Q$-algebras that are topological rings, with ideals $\mathfrak n\subseteq B$ and $\mathfrak n'\subseteq B'$ such that $\mathfrak n^m=0$ and $\mathfrak n'^{m'}=0$ for some $m,m'\ge1$. Let $p:B\to B'$ be a ring homomorphism and $s:B'\to B$ a continuous additive map with $p\circ s=\mathrm{id}$ and $s(\mathfrak n')\subseteq\mathfrak n$. Let $H$ be a topological group and $w:H\to B'$ a continuous multiplicative homomorphism with $w(g)-1\in\mathfrak n'$ for every $g$.
Then there is a continuous homomorphism $f:H\to B^\times$ with $p\circ f=w$ and $f(g)=1$ whenever $w(g)=1$.
\end{lemma}
\leanref{Development/AutomorphicArithmetic/AutomorphicVarieties/PadicHodgeTheory/SpLiftP.lean}{liftChar}
\leanref{Development/AutomorphicArithmetic/AutomorphicVarieties/PadicHodgeTheory/SpLiftP.lean}{log_ex}
\begin{proof}
For a nilpotent $a$, $\exp(a)=\sum_na^n/n!$ is a finite sum; $\exp(a+b)=\exp(a)\exp(b)$ for nilpotent $a$ and $b$; and $\exp$ is injective on nilpotent elements: if $\exp(a)=\exp(b)$, then $\exp(a-b)=1$, and $\exp(a-b)-1=(a-b)(1+(a-b)e)$ for some $e$, with $1+(a-b)e$ a unit.
\begin{enumerate}
\item[(i)] \emph{A logarithm.} There is a continuous $y:H\to\mathfrak n'$ with $\exp(y(g))=w(g)$. Let $y_0=0$ and $y_{k+1}=y_k+\big(w\exp(-y_k)-1\big)$. By induction, $y_k$ is continuous with values in $\mathfrak n'$, and $w\exp(-y_k)-1\in\mathfrak n'^{k+1}$: writing $e_k=w\exp(-y_k)-1$, $w\exp(-y_{k+1})-1=(1+e_k)\exp(-e_k)-1$, which lies in $e_k^2B'\subseteq\mathfrak n'^{2k+2}$, as $\exp(-e)=1-e+e^2(\cdots)$. Then $y=y_{m'}$ works.
\item[(ii)] By injectivity, $y(gh)=y(g)+y(h)$, both sides having the exponential $w(gh)$, and $y(g)=0$ when $w(g)=1$.
\end{enumerate}
Let $f(g)=\exp(s(y(g)))$, an exponential of an element of $\mathfrak n$. As $s$ is additive, $f$ is multiplicative; it is continuous, being a polynomial in $s(y(g))$; $p(f(g))=\exp(y(g))=w(g)$; and $f(g)=1$ when $w(g)=1$.
\end{proof}

\begin{lemma}\label{sha-counts-vd-model-local-lifts-commuting}
Let $\ell$ be a prime, $E$ a finite extension of $\mathbf Q_\ell$, $H$ a topological group, $B$ and $B'$ pointed artinian $E$-algebras (Lemma~\ref{sha-counts-vd-model-lift-chart-factor}), and $p:B\to B'$ a surjective homomorphism of $E$-algebras with $\pi_{B'}\circ p=\pi_B$. Let $\chi_1,\dots,\chi_4:H\to E^\times$ be homomorphisms, pairwise distinct, and $\theta:H\to E$ a function with $\chi_{5-i}\chi_i=\theta$ for every $i$. Let $r':H\to M_4(B')$ be a homomorphism with continuous entries, such that any two of the $r'(g)$ commute, $\pi_{B'}(r'(g))=\operatorname{diag}(\chi_1(g),\dots,\chi_4(g))$, and $r'(g)$ is symplectic with multiplier $\theta(g)$, for every $g$.
Then there is a homomorphism $r:H\to M_4(B)$ with continuous entries, with $r(g)$ symplectic with multiplier $\theta(g)$ for every $g$, $r(g)=1$ whenever $r'(g)=1$, and $p\circ r=r'$.
\end{lemma}
\leanref{Development/AutomorphicArithmetic/AutomorphicVarieties/PadicHodgeTheory/SpLiftP.lean}{core}
\begin{proof}
Write $\pi=\pi_B$ and $\pi'=\pi_{B'}$.
\begin{enumerate}
\item[(i)] The ideals $\ker\pi$ and $\ker\pi'$ consist of nilpotent elements and are finitely generated, $B$ and $B'$ being finite-dimensional; so they are nilpotent ideals. Let $s:B'\to B$ be an $E$-linear map with $p\circ s=\mathrm{id}$; it is continuous, and $s(\ker\pi')\subseteq\ker\pi$, as $\pi\circ s=\pi'\circ p\circ s=\pi'$. Each $\chi_i(g)=\pi'(r'(g)_{ii})$ is continuous in $g$.
\item[(ii)] By Lemma~\ref{sha-counts-vd-model-local-lifts-commuting-diagonal}, as $E$ is infinite, there is $u'\in M_4(B')$, symplectic with multiplier $1$, with $\pi'(u')=1$ and $r'(g)u'=u'\operatorname{diag}(d(g))$ for every $g$. By Lemma~\ref{sha-counts-vd-model-local-lifts-symplectic}, $u'=p(u)$ for some $u\in M_4(B)$ symplectic with multiplier $1$; so $u^{-1}=u^*$ and $u'^{-1}=u'^*$.
\item[(iii)] $d_i(g)=(u'^*r'(g)u')_{ii}$, so $d_i$ is multiplicative and continuous, $d(g)=(1,1,1,1)$ when $r'(g)=1$, $\pi'(d_i(g))=\chi_i(g)$, and $d_{5-i}(g)d_i(g)=\theta(g)$, as $\operatorname{diag}(d(g))^*\operatorname{diag}(d(g))=u'^*r'(g)^*r'(g)u'=\theta(g)$.
\item[(iv)] For $i=1,2$, the character $w_i(g)=d_i(g)\chi_i(g)^{-1}$ of $H$ is continuous, with $w_i(g)-1\in\ker\pi'$, and $w_i(g)=1$ when $r'(g)=1$ (then $\chi_i(g)=1$). By Lemma~\ref{sha-counts-vd-model-local-lifts-commuting-characters}, with $\mathfrak n=\ker\pi$ and $\mathfrak n'=\ker\pi'$, there is a continuous homomorphism $f_i:H\to B^\times$ with $p\circ f_i=w_i$ and $f_i(g)=1$ when $r'(g)=1$.
\item[(v)] Let $F(g)=(f_1(g),f_2(g),f_2(g)^{-1},f_1(g)^{-1})$. Then $F$ is multiplicative, $F_{5-i}F_i=1$, $F(g)=1$ when $r'(g)=1$, and $\chi_i(g)p(F_i(g))=d_i(g)$ for every $i$; for $i=3,4$ this follows from the case $5-i$, as $d_{5-i}d_i=\theta=\chi_{5-i}\chi_i$.
\end{enumerate}
Let $r(g)=u\operatorname{diag}\big(\chi_1(g)F_1(g),\dots,\chi_4(g)F_4(g)\big)u^*$. As $u^*u=1$ and $F$ and the $\chi_i$ are multiplicative, $r$ is a homomorphism; its entries are continuous; $r(g)^*r(g)=u\operatorname{diag}\big(\chi_{5-i}(g)\chi_i(g)F_{5-i}(g)F_i(g)\big)u^*=\theta(g)$; $r(g)=uu^*=1$ when $r'(g)=1$, as then $\chi_i(g)=1$ and $F(g)=1$; and $p(r(g))=u'\operatorname{diag}(d(g))u'^*=r'(g)$.
\end{proof}

\begin{lemma}\label{sha-counts-vd-model-local-lifts}
Let $D$ be a Siegel datum, with the notation of Definitions~\ref{sha-siegel-datum} and~\ref{sha-counts-vd-model-framing}, let $E$ be the fraction field of $\mathcal O$, let $v\in T$, and let pointed artinian $E$-algebras be as in Lemma~\ref{sha-counts-vd-model-lift-chart}. The conditions at $v$ of Definition~\ref{sha-counts-vd-model-framing}(1) make sense for homomorphisms $D_v\to M_4(A)$, for every commutative $\mathcal O$-algebra $A$. Let $p:B\to B'$ be a surjective homomorphism of pointed artinian $E$-algebras with $\pi_{B'}\circ p=\pi_B$, and let $r':D_v\to M_4(B')$ be a homomorphism with continuous entries and with $\pi_{B'}\circ r'=r_0$ on $D_v$ that satisfies the conditions at $v$.
Then $r'=p\circ r$ for a homomorphism $r:D_v\to M_4(B)$ with continuous entries and with $\pi_B\circ r=r_0$ on $D_v$ that satisfies the conditions at $v$.
\end{lemma}
\leanref{Development/AutomorphicArithmetic/AutomorphicVarieties/PadicHodgeTheory/SpLiftL.lean}{UpperTheory.V.Loc.spLiftL}
\leanref{Development/AutomorphicArithmetic/AutomorphicVarieties/PadicHodgeTheory/SpLiftP.lean}{UpperTheory.V.Loc.spLiftP}
\leanref{Development/AutomorphicArithmetic/AutomorphicVarieties/SymplecticGSp4/CommOfIntegral.lean}{UpperTheory.V.GradeZeroInterpolatingElement.abP}
\begin{proof}
Write $\pi=\pi_B$ and $\pi'=\pi_{B'}$, and view the $\delta_i$ as characters with values in $\mathcal O^\times\subseteq E^\times$; on $D_v$, $\pi'\circ r'=r_0=\operatorname{diag}(\delta_1,\dots,\delta_4)$.

\emph{The prime $\ell$.}
\begin{enumerate}
\item[(i)] The place $v_\ell$ identifies $D_\ell$ with $\operatorname{Gal}(\bar{\mathbf Q}_\ell/\mathbf Q_\ell)$, as topological groups, carrying $I_\ell$ to the inertia group and $\chi_{\mathrm{cyc}}$ to the cyclotomic character.
\item[(ii)] By Lemma~\ref{sha-counts-vd-model-local-lifts-flag}, applied to $r'$, every $r'(g)$ is symplectic with multiplier $\theta(g)$, and there are $h'$ and $\kappa'_1,\dots,\kappa'_4$ as there. By Lemma~\ref{sha-counts-vd-model-local-lifts-symplectic}, $h'=p(h)$ for an invertible $h\in M_4(B)$, symplectic with multiplier $1$, with $\pi(h)=1$. Let $T'(g)=h'^{-1}r'(g)h'$. It is upper triangular with diagonal $(\kappa'_i(g))$, it has continuous entries, it is symplectic with multiplier $\theta(g)$, and $\pi'\circ T'=r_0$, as $\pi'(h')=1$. In particular $\kappa'_4\kappa'_1=\kappa'_3\kappa'_2=\theta$, as the diagonal of an upper triangular matrix that is symplectic with multiplier $\nu$ satisfies $g_{ii}g_{5-i,5-i}=\nu$.
\item[(iii)] \emph{Two characters.} For $i=3,4$ let $w'_i=\kappa'_i\delta_i^{-1}:D_\ell\to B'^\times$. It is continuous, trivial on $I_\ell$, and $w'_i(g)-1$ is nilpotent, as $\pi'(w'_i(g))=1$. By (i) and Lemma~\ref{sha-counts-vd-model-local-lifts-unramified}, with the kernel of $p$ nilpotent since $\pi'\circ p=\pi$, there is a continuous homomorphism $u_i:D_\ell\to B^\times$, trivial on $I_\ell$, with $p\circ u_i=w'_i$. Put $k_i=\delta_iu_i$; then $p\circ k_i=\kappa'_i$ and $k_i=\delta_i$ on $I_\ell$.
\item[(iv)] Put $\kappa=(\theta k_4^{-1},\theta k_3^{-1},k_3,k_4)$. By (ii) and (iii), $p\circ\kappa_i=\kappa'_i$, and $\kappa_i=\delta_i$ on $I_\ell$, since $\delta_4=1$ and $\theta\delta_3^{-1}=\delta_2$. Let $\alpha'=k_3k_4^{-1}$, $\beta'=\theta k_3^{-2}$ and $u=k_4$; then $\kappa=(\alpha'^2\beta'u,\alpha'\beta'u,\alpha'u,u)$, and $\pi\circ\alpha'=\alpha$ and $\pi\circ\beta'=\beta$ on $D_\ell$, as $\pi\circ\kappa_i=\pi'\circ\kappa'_i=\delta_i$.
\item[(v)] \emph{Non-resonance.} There are elements of $D_\ell$ at which $\alpha$, $\beta$, $\alpha\beta=\Psi$ and $\alpha^2\beta=\theta$ differ from $\chi_{\mathrm{cyc}}$. As $\chi_{\mathrm{cyc}}(\sigma_\ell)=1$ and $\bar\omega(\sigma_\ell)=\bar\chi_{\mathrm{cyc}}(\sigma_\ell)=1$, the residues of $\alpha$, $\beta$ and $\theta$ at $\sigma_\ell$ are those of $\xi$, $\xi^{-1}\Psi$ and $\xi\Psi$, which differ from $1$ by the conditions on $\sigma_\ell$. At $s_0\in I_\ell$, $\Psi(s_0)=\chi_{\mathrm{cyc}}(s_0)^k$, as $\psi(s_0)=1$; if it were $\chi_{\mathrm{cyc}}(s_0)$, then $\chi_{\mathrm{cyc}}(s_0)^{k-1}=1$ with $k\ge2$, while $\chi_{\mathrm{cyc}}(s_0)$ has infinite order. So for each $\chi\in\{\alpha',\beta',\alpha'\beta',\alpha'^2\beta'\}$ there is $g$ with $\pi\big(\chi(g)-\chi_{\mathrm{cyc}}(g)\big)\ne0$, and $\chi(g)-\chi_{\mathrm{cyc}}(g)$ is a unit of $B$, as $\ker\pi$ consists of nilpotent elements.
\item[(vi)] By (i), (iv) and (v), Lemma~\ref{sha-counts-vd-model-local-lifts-borel} applies to $T'$, $\alpha'$, $\beta'$ and $u$; it gives $T$, upper triangular with diagonal $(\kappa_i(g))$, symplectic with multiplier $\alpha'^2\beta'u^2=\kappa_1\kappa_4=\theta$, with continuous entries and $p\circ T=T'$. Let $r(g)=hT(g)h^{-1}$. Then $r$ has continuous entries, $p\circ r=h'T'h'^{-1}=r'$, $\pi\circ r=\pi'\circ p\circ r=r_0$, and $r(g)$ is symplectic with multiplier $\theta(g)$, as $h$ is symplectic with multiplier $1$. By Lemma~\ref{sha-counts-vd-model-local-lifts-flag}, with $h$ and $\kappa$, $r$ satisfies the conditions at $\ell$.
\end{enumerate}

\emph{A prime $p_0\in S_b$.} Let $v=p_0$, and write $p$ only for the homomorphism $B\to B'$.
\begin{enumerate}
\item[(i)] The commutator of two elements of $D_{p_0}$ lies in $I_{p_0}\cap\ker\psi$: in $I_{p_0}$ because $D_{p_0}/I_{p_0}$ is abelian, and in $\ker\psi$ because $\psi$ is a character. Since $r'$ is trivial on $I_{p_0}\cap\ker\psi$, any two of the $r'(g)$ commute.
\item[(ii)] The group $\psi(I_{p_0})$ is finite, as $\psi$ has open kernel, and cyclic, as a finite group of units of a domain; its order is divisible by $\ell$, as $p_0\in S_b$. So it is generated by $\psi(\tau)\ne1$ for some $\tau\in I_{p_0}$.
\item[(iii)] On $I_{p_0}$, $\alpha=1$ and $\Psi=\theta=\psi$: $\chi_{\mathrm{cyc}}$ is unramified at $p_0\ne\ell$, $\xi$ is trivial on $I_{p_0}$, and $\omega$ is trivial there, its values being determined by their residues $\bar\chi_{\mathrm{cyc}}$. The restrictions $\chi_i$ of the $\delta_i$ to $D_{p_0}$ are pairwise distinct: if $\lambda_i\ne\lambda_j$, their residues at $\sigma_{p_0}$ differ, by Lemma~\ref{sha-flagged-residues}(1); and $\delta_1(\tau)=\delta_2(\tau)=\psi(\tau)\ne1=\delta_3(\tau)=\delta_4(\tau)$. Also $\delta_{5-i}\delta_i=\theta$.
\item[(iv)] By (i) and (iii), Lemma~\ref{sha-counts-vd-model-local-lifts-commuting} applies to $H=D_{p_0}$, the $\chi_i$, $\theta$ and $r'$, which is symplectic with multiplier $\theta$ by the conditions at $p_0$. It gives $r$ with continuous entries, symplectic with multiplier $\theta$, with $p\circ r=r'$, and trivial wherever $r'$ is, in particular on $I_{p_0}\cap\ker\psi$. So $r$ satisfies the conditions at $p_0$, and $\pi\circ r=\pi'\circ p\circ r=r_0$.
\end{enumerate}
% The proofs of Lemmas -local-lifts-flag (spFlag), -symplectic (lift_symplectic_similitude_nilpotent), -unramified (lift_unramified_unipotent_character), -borel (borLift, with cocLift as Lemma -local-lifts-borel-cocycle) and -commuting (the private core of SpLiftP.lean, with diagAll and diagX as Lemma -local-lifts-commuting-diagonal and liftChar with log_ex as Lemma -local-lifts-commuting-characters) are written out after their statements. Written out here: spLiftL itself (with the private exists_iso in (i), gsp_diag in (ii), hlift in (iii), hkappap and hkappaI in (iv), nonres and hun in (v), borLift and cj in (vi)), and spLiftP with abP and the facts f2, f5, f6.
\end{proof}

\begin{lemma}\label{sha-counts-vd-model-local-tangent-conditions}
Let $D$ be a Siegel datum, with the notation of Definitions~\ref{sha-siegel-datum} and~\ref{sha-counts-vd-model-framing}, let $E$ be the fraction field of $\mathcal O$, with its topology as a finite extension of $\mathbf Q_\ell$, and let $v\in T$. Let $\gamma:D_v\to M_4(E)$ have continuous entries, such that $r(g)=r_0(g)+\varepsilon\gamma(g)r_0(g)$ is a homomorphism $D_v\to M_4(E[\varepsilon])$, $\varepsilon^2=0$. Then $r$ satisfies the conditions at $v$ of Definition~\ref{sha-counts-vd-model-framing}(1) if and only if:
\begin{enumerate}
\item $\gamma(g)\in\mathfrak{sp}_4(E)$ for every $g\in D_v$;
\item there is $X\in\mathfrak{sp}_4(E)$ with $\gamma(g)_{ab}=\big(1-\delta_a(g)\delta_b(g)^{-1}\big)X_{ab}$ for all $g\in D_v$ and all $a,b$ with $b<a$ if $v=\ell$, and with $a\ne b$ if $v\in S_b$;
\item $\gamma(s)_{aa}=0$ for every $s\in I_v$ and every $a$.
\end{enumerate}
\end{lemma}
\leanref{Development/AutomorphicArithmetic/AutomorphicVarieties/PadicHodgeTheory/SpTanL.lean}{UpperTheory.V.Loc.spTanL}
\leanref{Development/AutomorphicArithmetic/AutomorphicVarieties/PadicHodgeTheory/SpTanP.lean}{UpperTheory.V.Loc.spTanP}
\begin{proof}
Let $\chi_a$ be the restriction of $\delta_a$ to $D_v$, with values in $\mathcal O^\times\subseteq E^\times$, so that $r_0=\operatorname{diag}(\chi_1,\dots,\chi_4)$ on $D_v$ and $\chi_{5-a}\chi_a=\theta$; put $\kappa_{ab}=\chi_a\chi_b^{-1}$, so that $\kappa_{5-b,5-a}=\kappa_{ab}$. Comparing the coefficients of $\varepsilon$, as $r$ is a homomorphism, $\gamma(gh)_{ab}=\gamma(g)_{ab}+\kappa_{ab}(g)\gamma(h)_{ab}$ for all $g$, $h$, $a$ and $b$.

\emph{The symplectic condition.} For $g\in D_v$, $r(g)$ is symplectic with multiplier $\theta(g)$ if and only if $\gamma(g)\in\mathfrak{sp}_4(E)$. Indeed, let $P=r_0(g)$. Then $P^*=\operatorname{diag}(\chi_4(g),\dots,\chi_1(g))$, $P^*P=\theta(g)$, and $(YZ)^*=Z^*Y^*$ for all $Y$ and $Z$; so $r(g)^*r(g)=\theta(g)+\varepsilon\big(P^*\gamma(g)P+P^*\gamma(g)^*P\big)$, whose coefficient of $\varepsilon$ has the entry $\chi_{5-i}(g)\chi_j(g)\big(\gamma(g)+\gamma(g)^*\big)_{ij}$ at $(i,j)$.

\emph{The prime $\ell$.} The conditions at $\ell$ are the symplectic condition and the three ordinary conditions of Definition~\ref{sha-siegel-datum}(2), on the characteristic polynomial of $r(s_0)$, on the $e_jr(g)e_i$ with $i<j$, and on the $e_i(r(s)-\delta_i(s))e_i$ with $s\in I_\ell$, where $e_i=\prod_{j\ne i}(r(s_0)-\delta_j(s_0))$.
\begin{enumerate}
\item[(i)] The $\chi_a(s_0)$ are pairwise distinct: $\chi_a(s_0)=\chi_{\mathrm{cyc}}(s_0)^{e_a}$ with $e=(k+1,k,1,0)$, as $\psi(s_0)=\omega(s_0)=1$ and $\xi$ is trivial on $I_\ell$, the $e_a$ are distinct, as $k\ge2$, and $\chi_{\mathrm{cyc}}(s_0)$ has infinite order.
\item[(ii)] For $Y\in M_4(E)$ let $u_Y=1+\varepsilon Y$, whose inverse is $1-\varepsilon Y$. The entry $(a,b)$ of $u_Y^{-1}r(g)u_Y$ is $P_{ab}+\varepsilon\big(\gamma(g)_{ab}\chi_b(g)-Y_{ab}\chi_b(g)+\chi_a(g)Y_{ab}\big)$, with $P=r_0(g)$.
\item[(iii)] Let $Y\in M_4(E)$ satisfy $\gamma(s_0)_{ab}=(1-\kappa_{ab}(s_0))Y_{ab}$ for $a\ne b$. Then $r$ satisfies the three ordinary conditions if and only if $\gamma(g)_{ab}=(1-\kappa_{ab}(g))Y_{ab}$ for all $g\in D_\ell$ and $b<a$, and $\gamma(s)_{aa}=0$ for all $s\in I_\ell$ and all $a$. By (ii) and the hypothesis on $Y$, $u_Y^{-1}r(s_0)u_Y$ is the diagonal matrix with entries $c_a=\chi_a(s_0)\big(1+\varepsilon\gamma(s_0)_{aa}\big)$. Both sides imply that $\gamma(s_0)_{aa}=0$ for every $a$: the right side as $s_0\in I_\ell$; the left side because, the characteristic polynomial of $r(s_0)$ being $\prod_j(X-c_j)$, it gives $\prod_j(c_a-\chi_j(s_0))=0$, where the factors with $j\ne a$ are units of $E[\varepsilon]$ by (i), so $c_a=\chi_a(s_0)$. Suppose then that $\gamma(s_0)_{aa}=0$ for every $a$, so that $r(s_0)=u_Y\operatorname{diag}(\chi_1(s_0),\dots,\chi_4(s_0))u_Y^{-1}$. Its characteristic polynomial is $\prod_j(X-\chi_j(s_0))$, and $e_i=u_Y\big(c'_iE_{ii}\big)u_Y^{-1}$, where $E_{ii}$ is the matrix with entry $1$ at $(i,i)$ and $0$ elsewhere and $c'_i=\prod_{j\ne i}(\chi_i(s_0)-\chi_j(s_0))$ is a unit by (i). So for every matrix $M$, $e_jMe_i=0$ if and only if the entry $(j,i)$ of $u_Y^{-1}Mu_Y$ vanishes. By (ii), for $i<j$ this entry for $M=r(g)$ is $\varepsilon\big(\gamma(g)_{ji}\chi_i(g)-Y_{ji}\chi_i(g)+\chi_j(g)Y_{ji}\big)$, which vanishes if and only if $\gamma(g)_{ji}=(1-\kappa_{ji}(g))Y_{ji}$; and for $M=r(s)-\chi_i(s)$ and $j=i$ it is $\varepsilon\gamma(s)_{ii}\chi_i(s)$, which vanishes if and only if $\gamma(s)_{ii}=0$.
\item[(iv)] Suppose that $r$ satisfies the conditions at $\ell$. Then (1) holds by the symplectic condition. Let $Y_{ab}=\gamma(s_0)_{ab}/(1-\kappa_{ab}(s_0))$ for $a\ne b$, which is defined by (i), and $Y_{aa}=0$. By (iii), (2) holds with $X=Y$, and (3) holds. And $Y\in\mathfrak{sp}_4(E)$, as $\gamma(s_0)\in\mathfrak{sp}_4(E)$ and $\kappa_{5-b,5-a}=\kappa_{ab}$.
\item[(v)] Conversely, suppose (1), (2) with some $X$, and (3). Let $Y_{ab}=X_{ab}$ if $b<a$, and $Y_{ab}=\gamma(s_0)_{ab}/(1-\kappa_{ab}(s_0))$ if $a<b$. Then $Y$ satisfies the hypothesis of (iii), by (2) at $s_0$ when $b<a$, and the right side of (iii) holds by (2) and (3). So $r$ satisfies the ordinary conditions, and by (1) the symplectic condition.
\end{enumerate}

\emph{A prime $v\in S_b$.} The conditions at $v$ are the symplectic condition and the triviality of $r$ on $N=I_v\cap\ker\psi$.
\begin{enumerate}
\item[(i)] On $I_v$, $(\chi_1,\chi_2,\chi_3,\chi_4)=(\psi,\psi,1,1)$, as in step (iii) of the proof of Lemma~\ref{sha-counts-vd-model-local-lifts} at the primes of $S_b$; so every $\chi_a$ is trivial on $N$, and for $n\in N$, $r(n)=1$ if and only if $\gamma(n)=0$.
\item[(ii)] For $s\in I_v$, $s^m\in N$ for some $m\ge1$: take $m$ to be the index of the open normal subgroup $\ker\psi$ of $G$.
\item[(iii)] For $a\ne b$ there is $g_0\in D_v$ with $\kappa_{ab}(g_0)\ne1$ such that, for every $g\in D_v$, the commutator $c=g_0gg_0^{-1}g^{-1}$ lies in $N$, or lies in $I_v$ while $\kappa_{ab}$ is trivial on $I_v$. If $\kappa_{ab}$ is trivial on $I_v$, then $\{a,b\}$ is $\{1,2\}$ or $\{3,4\}$ by (i), as $\psi$ is not trivial on $I_v$ ($\ell$ divides the order of $\psi(I_v)$); so $\kappa_{ab}=\alpha^{\pm1}$, and $g_0=\sigma_v$ works: $\alpha(\sigma_v)\ne1$, as the residues of $\delta_1(\sigma_v)$ and $\delta_2(\sigma_v)$ differ by Lemma~\ref{sha-flagged-residues}(1), and $c\in I_v$, as $D_v/I_v$ is abelian. Otherwise take $g_0=s\in I_v$ with $\kappa_{ab}(s)\ne1$: then $c=s\cdot gs^{-1}g^{-1}$ lies in $I_v$, which is normal in $D_v$, and in $\ker\psi$, as $\psi$ is a character.
\end{enumerate}
Suppose that $r$ satisfies the conditions at $v$. Then (1) holds, and $\gamma$ vanishes on $N$ by (i).
\begin{itemize}
\item If $s\in I_v$ and $\kappa_{ab}(s)=1$, then $\gamma(s)_{ab}=0$: with $m$ as in (ii), the identity for $\gamma$ gives $\gamma(s^m)_{ab}=m\gamma(s)_{ab}$, and $\gamma(s^m)=0$. In particular (3) holds.
\item Let $a\ne b$, let $g_0$ be as in (iii), and let $g\in D_v$ and $c=g_0gg_0^{-1}g^{-1}$. Then $\gamma(c)_{ab}=0$ and $\kappa_{ab}(c)=1$, by (i) if $c\in N$ and by the previous item otherwise. The identity for $\gamma$, applied to $c\cdot gg_0=g_0g$, gives $\gamma(g_0)_{ab}(1-\kappa_{ab}(g))=\gamma(g)_{ab}(1-\kappa_{ab}(g_0))$. So $\gamma(g)_{ab}=(1-\kappa_{ab}(g))X_{ab}$ with $X_{ab}=\gamma(g_0)_{ab}/(1-\kappa_{ab}(g_0))$. Put $X_{aa}=0$.
\item $X\in\mathfrak{sp}_4(E)$. Let $a\ne b$ and $g_0$ be as in (iii). As $\gamma(g_0)\in\mathfrak{sp}_4(E)$, $\epsilon_a\epsilon_b\gamma(g_0)_{5-b,5-a}=-\gamma(g_0)_{ab}$, and the previous item at $(a,b)$ and at $(5-b,5-a)$, with $\kappa_{5-b,5-a}=\kappa_{ab}$, gives $(1-\kappa_{ab}(g_0))(\epsilon_a\epsilon_bX_{5-b,5-a}+X_{ab})=0$.
\end{itemize}
So (2) holds. Conversely, suppose (1), (2) and (3). For $n\in N\subseteq I_v$, $\gamma(n)_{aa}=0$ by (3), and $\gamma(n)_{ab}=(1-\kappa_{ab}(n))X_{ab}=0$ for $a\ne b$ by (i); so $r(n)=1$, and with (1), $r$ satisfies the conditions at $v$.
\end{proof}

\begin{lemma}\label{sha-counts-vd-model-local-tangent-cocycles}
Let $D$ be a Siegel datum, with the notation of Definition~\ref{sha-siegel-datum}, let $E$ be the fraction field of $\mathcal O$, with its topology as a finite extension of $\mathbf Q_\ell$, and let $1\le a<b\le4$. Then the continuous functions $z:D_\ell\to E$ with $z(gh)=z(g)+\delta_a(g)\delta_b(g)^{-1}z(h)$ for all $g,h\in D_\ell$ form an $E$-vector space of dimension $2$.
\end{lemma}
\leanref{Development/AutomorphicArithmetic/AutomorphicVarieties/PadicHodgeTheory/ExistsFinrankTwoCocyclesCharRatio.lean}{UpperTheory.exists_finrank_two_cocycles_char_ratio}
\begin{proof}
The place $v_\ell$ identifies $D_\ell$ with $\operatorname{Gal}(\bar{\mathbf Q}_\ell/\mathbf Q_\ell)$, carrying $\chi_{\mathrm{cyc}}$ to the cyclotomic character. Let $\kappa=\delta_a\delta_b^{-1}$ on $D_\ell$, a continuous character with values in $\mathcal O^\times\subseteq E^\times$, and $E(\kappa)$ the line $E$ with $D_\ell$ acting through $\kappa$.
\begin{enumerate}
\item[(i)] $\kappa\ne1$: the $\delta_i(s_0)$ are pairwise distinct, as in step (i) of the proof of Lemma~\ref{sha-counts-vd-model-local-tangent-conditions} at $\ell$, so $\kappa(s_0)\ne1$.
\item[(ii)] $\kappa\ne\chi_{\mathrm{cyc}}$. If $\lambda_a=\lambda_b$, then $(a,b)$ is $(1,3)$ or $(2,4)$, and $\kappa(s_0)=\chi_{\mathrm{cyc}}(s_0)^{k}$, as $\delta_i(s_0)=\chi_{\mathrm{cyc}}(s_0)^{e_i}$ with $e=(k+1,k,1,0)$; this is not $\chi_{\mathrm{cyc}}(s_0)$, as $k\ge2$ and $\chi_{\mathrm{cyc}}(s_0)$ has infinite order. If $\lambda_a\ne\lambda_b$, the residue of $\kappa(\sigma_\ell)$ is not $1$, by Lemma~\ref{sha-flagged-residues}(1), while $\chi_{\mathrm{cyc}}(\sigma_\ell)=1$.
\item[(iii)] By the local Euler characteristic formula and local Tate duality for $\operatorname{Gal}(\bar{\mathbf Q}_\ell/\mathbf Q_\ell)$ \cite[Chapter~VII, \S\S2--3]{NSW08}, applied to the continuous representation $E(\kappa)$ over the finite extension $E$ of $\mathbf Q_\ell$, $\dim_EH^1(D_\ell,E(\kappa))=1+\dim_EH^0(D_\ell,E(\kappa))+\dim_EH^0(D_\ell,E(\kappa^{-1}\chi_{\mathrm{cyc}}))$. Both spaces $H^0$ vanish, by (i) and (ii); so $H^1(D_\ell,E(\kappa))$ has dimension $1$.
\item[(iv)] The cocycles of the statement are the continuous cocycles of $E(\kappa)$. The coboundaries $g\mapsto(\kappa(g)-1)m$, $m\in E$, are cocycles, and form a space of dimension $1$, as $\kappa\ne1$.
\end{enumerate}
So the cocycles form a space of dimension $1+1=2$.
\end{proof}

\begin{lemma}\label{sha-counts-vd-model-local-tangent-unramified-cocycle}
Let $v$ and $\ell$ be primes, let $G_v=\operatorname{Gal}(\bar{\mathbf Q}_v/\mathbf Q_v)$, with inertia group $I$, and let $\sigma\in G_v$ be a Frobenius element. Let $A$ be a commutative $\mathbf Q_\ell$-algebra and $V$ an $A$-module that is finite-dimensional over $\mathbf Q_\ell$, with $\mathbf Q_\ell$ acting through $A$ and with its topology as a $\mathbf Q_\ell$-vector space. Let $G_v$ act on $V$ by $A$-linear maps, trivially on $I$, such that $(g,y)\mapsto g\cdot y$ is continuous on $G_v\times V$.
Then for every $x\in V$ there is a unique continuous function $c:G_v\to V$ with $c(gh)=c(g)+g\cdot c(h)$ for all $g,h\in G_v$, $c(g)=0$ for $g\in I$, and $c(\sigma)=x$.
\end{lemma}
\leanref{Development/AutomorphicArithmetic/AutomorphicVarieties/PadicHodgeTheory/ExistsUniqueUnramifiedCocycleFrobeniusValue.lean}{UpperTheory.V.Loc.existsUnique_unramified_cocycle_frobenius_value}
\begin{proof}
Two published facts on $G_v$ are used \cite[Chapter~VII]{NSW08}: (a) $I$ and $\sigma$ generate a dense subgroup of $G_v$; (b) for every compact, totally disconnected Hausdorff group $K$ and every $y\in K$ there is a continuous homomorphism $G_v\to K$, trivial on $I$, that maps $\sigma$ to $y$. Fix a basis of $V$ over $\mathbf Q_\ell$, and for $y\in V$ let $|y|$ be the largest absolute value of its coordinates. Through the coordinates, $V$ is homeomorphic to $\mathbf Q_\ell^n$, so it is Hausdorff and totally disconnected.
\begin{enumerate}
\item[(i)] \emph{A stable lattice.} The map $g\mapsto g\cdot x$ is continuous on the compact group $G_v$, so there is $C\ge1$ with $|g\cdot x|\le C$ for every $g$. Let $L$ be the set of the $y\in V$ with $|g\cdot y|\le C$ for every $g\in G_v$. It is a subgroup of $V$, as the absolute value is ultrametric, and it is closed, as each $y\mapsto g\cdot y$ is continuous. It lies in the compact set of the $y$ with $|y|\le C$, the case $g=1$, so it is compact. It contains $x$, and it is stable under $G_v$: if $y\in L$, then $|h\cdot(g\cdot y)|=|(hg)\cdot y|\le C$ for every $h$.
\item[(ii)] \emph{The semidirect product.} On the set $V\times G_v$, the product $(y,g)(y',g')=(y+g\cdot y',gg')$ defines a group, with identity $(0,1)$ and inverse $(y,g)^{-1}=(-g^{-1}\cdot y,g^{-1})$. With the product topology it is a topological group, since the action is continuous. Its subset $L\times G_v$ is a subgroup, as $L$ is stable under $G_v$; and it is compact, totally disconnected and Hausdorff, as $L$ and $G_v$ are.
\item[(iii)] \emph{A homomorphism.} By (b), applied to $K=L\times G_v$ and $y=(x,\sigma)$, there is a continuous homomorphism $f:G_v\to L\times G_v$, trivial on $I$, with $f(\sigma)=(x,\sigma)$. Write $f(g)=(c(g),k(g))$. Then $k:G_v\to G_v$ is a continuous homomorphism, trivial on $I$, with $k(\sigma)=\sigma$; and $c$ is continuous, with $c(gh)=c(g)+k(g)\cdot c(h)$, $c(g)=0$ for $g\in I$, and $c(\sigma)=x$.
\item[(iv)] \emph{The action of $k(g)$.} Let $H$ be the set of the $g\in G_v$ with $k(g)\cdot y=g\cdot y$ for every $y\in V$. It is a subgroup, as $k$ is a homomorphism. It is closed: for each $y$ the maps $g\mapsto k(g)\cdot y$ and $g\mapsto g\cdot y$ are continuous, and $V$ is Hausdorff. It contains $I$, on which $k$ is trivial and $G_v$ acts trivially, and $\sigma$. By (a), $H=G_v$. So $c(gh)=c(g)+g\cdot c(h)$ for all $g,h$, and $c$ has the stated properties.
\item[(v)] \emph{Uniqueness.} A function $c'$ with $c'(gh)=c'(g)+g\cdot c'(h)$ has $c'(1)=0$, from $c'(1\cdot1)$, and $c'(g^{-1})=-g^{-1}\cdot c'(g)$, from $c'(g^{-1}g)$. So for two functions $c$ and $c'$ as in the statement, the set where they agree is a subgroup of $G_v$. It is closed, as $c$ and $c'$ are continuous and $V$ is Hausdorff, and it contains $I$ and $\sigma$. By (a) it is $G_v$.
\end{enumerate}
\end{proof}

\begin{lemma}\label{sha-counts-vd-model-local-tangent-unramified}
Let $v$ and $\ell$ be primes, let $D_v\supseteq I_v$ be the decomposition and inertia groups at a place of $\bar{\mathbf Q}$ above $v$, let $\sigma\in D_v$ be a Frobenius element, and let $V$ be a finite-dimensional $\mathbf Q_\ell$-vector space. For every $e\in V$ there is a unique continuous homomorphism $u:D_v\to V$ that vanishes on $I_v$ with $u(\sigma)=e$.
\end{lemma}
\leanref{Development/AutomorphicArithmetic/AutomorphicVarieties/PadicHodgeTheory/ExistsUniqueUnramifiedHomOfFrobeniusValue.lean}{UpperTheory.exists_unique_unramified_hom_of_frobenius_value}
\begin{proof}
The place identifies $D_v$ with $G_v=\operatorname{Gal}(\bar{\mathbf Q}_v/\mathbf Q_v)$, as topological groups, carrying $I_v$ to the inertia group and $\sigma$ to a Frobenius element $\sigma'$. Give $V$ its topology as a finite-dimensional $\mathbf Q_\ell$-vector space and the trivial action of $G_v$, which is continuous. For this action the functions $c$ of Lemma~\ref{sha-counts-vd-model-local-tangent-unramified-cocycle}, with $A=\mathbf Q_\ell$, are the continuous homomorphisms $G_v\to V$ that vanish on the inertia group.

\emph{Existence.} By Lemma~\ref{sha-counts-vd-model-local-tangent-unramified-cocycle} there is such a homomorphism $c$ with $c(\sigma')=e$. Its composite with the identification is a continuous homomorphism $u:D_v\to V$ that vanishes on $I_v$, with $u(\sigma)=e$.

\emph{Uniqueness.} Let $u$ be a continuous homomorphism $D_v\to V$ that vanishes on $I_v$ with $u(\sigma)=0$. Its composite with the inverse of the identification, and the function $0$, are homomorphisms as in Lemma~\ref{sha-counts-vd-model-local-tangent-unramified-cocycle} with value $0$ at $\sigma'$; by the uniqueness there, they are equal, so $u=0$. The difference of two homomorphisms as in the statement is such a $u$.
\end{proof}

\begin{lemma}\label{sha-counts-vd-model-local-tangent-root}
Let $E$ be a field of characteristic $0$, use the weights of Definition~\ref{sha-siegel-datum}, and let $a\ne b$. Then the entries of the weight of $(a,b)$ are $(a,b)$ and $(5-b,5-a)$, and there is $X\in\mathfrak{sp}_4(E)$ with $X_{ab}=1$ whose other entries vanish outside the weight of $(a,b)$, such that every $Y\in\mathfrak{sp}_4(E)$ satisfies $Y_{a'b'}=Y_{ab}X_{a'b'}$ for the $(a',b')$ of the weight of $(a,b)$.
\end{lemma}
\leanref{Development/AutomorphicArithmetic/AutomorphicVarieties/PadicHodgeTheory/ExistsSp4RootVectorUniqueOnWeight.lean}{UpperTheory.exists_sp4_root_vector_unique_on_weight}

\begin{lemma}\label{sha-counts-vd-model-local-tangent}
Let $D$ be a Siegel datum, with the notation of Definitions~\ref{sha-siegel-datum} and~\ref{sha-counts-vd-model-framing}, let $E$ be the fraction field of $\mathcal O$, with its topology as a finite extension of $\mathbf Q_\ell$, and let $v\in T$ and $j\in\mathbf Z^2$. Let $V_j$ be the $E$-vector space of the functions $s:D_v\to M_4(E)$ with continuous entries and with $s(g)_{ab}=0$ whenever $(a,b)$ does not have weight $j$, such that $g\mapsto r_0(g)+\varepsilon s(g)$ is a homomorphism $D_v\to M_4(E[\varepsilon])$, $\varepsilon^2=0$, that satisfies the conditions at $v$ of Definition~\ref{sha-counts-vd-model-framing}(1).
Then $\dim_EV_j=d(j)$ if $v=\ell$, and $\dim_EV_j=d'(j)$ if $v\in S_b$.
\end{lemma}
\leanref{Development/AutomorphicArithmetic/AutomorphicVarieties/PadicHodgeTheory/TangentSpaceDeformationAtL.lean}{UpperTheory.tangent_space_deformation_at_l}
\leanref{Development/AutomorphicArithmetic/AutomorphicVarieties/PadicHodgeTheory/TangentSpaceDeformationDualNumbers.lean}{UpperTheory.tangent_space_deformation_dual_numbers}
\begin{proof}
Let $\chi_a$ be the restriction of $\delta_a$ to $D_v$, with values in $\mathcal O^\times\subseteq E^\times$; the $\chi_a$ are continuous, and $\chi_{5-a}\chi_a=\theta$. They are pairwise distinct: if $v=\ell$, then $\delta_a(s_0)=\chi_{\mathrm{cyc}}(s_0)^{e_a}$ with $e=(k+1,k,1,0)$, since $\psi(s_0)=\omega(s_0)=1$ and $\xi$ is trivial on $I_\ell$, and $\chi_{\mathrm{cyc}}(s_0)$ has infinite order; if $v\in S_b$, as in step (iii) of the proof of Lemma~\ref{sha-counts-vd-model-local-lifts} at the primes of $S_b$. For $s:D_v\to M_4(E)$ let $\gamma_s(g)_{ab}=s(g)_{ab}\chi_b(g)^{-1}$, so that $r_0+\varepsilon s=r_0+\varepsilon\gamma_sr_0$; the map $s\mapsto\gamma_s$ is an $E$-linear bijection, which preserves continuity and the vanishing of entries.
\begin{enumerate}
\item[(i)] \emph{The conditions.} Comparing the coefficients of $\varepsilon$, $r_0+\varepsilon\gamma r_0$ is a homomorphism if and only if $\gamma(gh)_{ab}=\gamma(g)_{ab}+\chi_a(g)\chi_b(g)^{-1}\gamma(h)_{ab}$ for all $g$, $h$, $a$ and $b$. So by Lemma~\ref{sha-counts-vd-model-local-tangent-conditions}, $s\mapsto\gamma_s$ maps $V_j$ onto the space $Q_j$ of the $\gamma$ with continuous entries that satisfy this identity and the conditions (1) to (3) there, and whose entries outside the weight $j$ vanish.
\item[(ii)] \emph{Weight $(0,0)$.} The entries of weight $(0,0)$ are the diagonal ones. For $\gamma\in Q_{(0,0)}$, each $\gamma_{aa}$ is a continuous homomorphism $D_v\to E$ vanishing on $I_v$, and $\gamma_{44}=-\gamma_{11}$ and $\gamma_{33}=-\gamma_{22}$, since $\gamma(g)\in\mathfrak{sp}_4(E)$. By Lemma~\ref{sha-counts-vd-model-local-tangent-unramified}, the map $\gamma\mapsto(\gamma(\sigma_v)_{11},\gamma(\sigma_v)_{22})$ is injective, and it is surjective: for homomorphisms $u,u'$ vanishing on $I_v$ with given values at $\sigma_v$, $\gamma=\operatorname{diag}(u,u',-u',-u)$ lies in $Q_{(0,0)}$, with $X=0$ in (2). So $\dim Q_{(0,0)}=2=d(0,0)=d'(0,0)$.
\item[(iii)] \emph{A weight of a pair $(a,b)$ with $a<b$, at $\ell$.} Let $X$ be as in Lemma~\ref{sha-counts-vd-model-local-tangent-root}. The map $\gamma\mapsto\gamma_{ab}$ is a bijection from $Q_j$ onto the space of Lemma~\ref{sha-counts-vd-model-local-tangent-cocycles}. It is injective, since $\gamma(g)\in\mathfrak{sp}_4(E)$ vanishes outside the weight $j$, so $\gamma(g)=\gamma(g)_{ab}X$. For a cocycle $z$ in that space, $\gamma(g)=z(g)X$ lies in $Q_j$: the identity of (i) holds at $(a,b)$, and at $(5-b,5-a)$ because $\chi_{5-a}\chi_{5-b}^{-1}=\chi_a\chi_b^{-1}$; condition (2) holds with $0$, since no entry $(a',b')$ with $b'\le a'$ has the weight of $(a,b)$; and the diagonal entries vanish. So $\dim Q_j=2=d(j)$; the weights of the pairs $(a,b)$ with $a<b$ are $(1,0)$, $(0,1)$, $(1,1)$ and $(-1,1)$.
\item[(iv)] \emph{A weight of a pair $(a,b)$ with $b<a$ at $\ell$, or with $a\ne b$ at $v\in S_b$.} Let $E_{ab}$ be the matrix with entry $1$ at $(a,b)$ and $0$ elsewhere, so that $E_{ab}^*=\epsilon_a\epsilon_bE_{5-b,5-a}$. If $v=\ell$, let $X=E_{ab}-E_{ab}^*$. If $v\in S_b$, let $X=E_{ab}-E_{ab}^*$ if $(5-b,5-a)\ne(a,b)$, and $X=E_{ab}$ if $(5-b,5-a)=(a,b)$.
\begin{itemize}
\item $X\in\mathfrak{sp}_4(E)$: $(Y^*)^*=Y$ for every $Y$, and if $(5-b,5-a)=(a,b)$, then $\epsilon_b=\epsilon_{5-a}=-\epsilon_a$, so $E_{ab}^*=-E_{ab}$.
\item $X_{ab}\ne0$: it is $1$, except that it is $2$ if $v=\ell$ and $(5-b,5-a)=(a,b)$; and $E$ has characteristic $0$.
\item $X$ vanishes outside the weight of $(a,b)$, which is also the weight of $(5-b,5-a)$, by Lemma~\ref{sha-counts-vd-model-local-tangent-root}.
\end{itemize}
Let $w(g)_{a'b'}=(1-\chi_{a'}(g)\chi_{b'}(g)^{-1})X_{a'b'}$ for all $a'$ and $b'$; on the diagonal it vanishes. Then $w\in Q_j$: it satisfies the identity of (i), being a coboundary; $w(g)\in\mathfrak{sp}_4(E)$, as $X\in\mathfrak{sp}_4(E)$ and $\chi_{5-b'}\chi_{5-a'}^{-1}=\chi_{a'}\chi_{b'}^{-1}$; (2) holds with $X$; and its diagonal entries vanish. Also $w\ne0$: $\chi_a(g_0)\ne\chi_b(g_0)$ for some $g_0$, and $X_{ab}\ne0$, so $w(g_0)_{ab}\ne0$. Let $\gamma\in Q_j$, with $Y$ as in (2), and put $c=Y_{ab}/X_{ab}$. Then $\gamma(g)_{ab}=(1-\chi_a(g)\chi_b(g)^{-1})Y_{ab}=c\,w(g)_{ab}$. At $(5-b,5-a)$, if $v=\ell$, $\gamma(g)=c\,w(g)$ follows from the entry $(a,b)$, as $\gamma(g)$ and $w(g)$ lie in $\mathfrak{sp}_4(E)$; if $v\in S_b$, then $Y_{5-b,5-a}=-\epsilon_a\epsilon_bY_{ab}$ and $X_{5-b,5-a}=-\epsilon_a\epsilon_bX_{ab}$, as $X,Y\in\mathfrak{sp}_4(E)$, so (2) gives $\gamma(g)_{5-b,5-a}=(1-\chi_{5-b}(g)\chi_{5-a}(g)^{-1})Y_{5-b,5-a}=c\,w(g)_{5-b,5-a}$. Elsewhere both vanish. So $\gamma=c\,w$. So $\dim Q_j=1$, which is $d(j)$ at $\ell$, the weights of the pairs with $b<a$ being $(-1,0)$, $(0,-1)$, $(-1,-1)$ and $(1,-1)$, and $d'(j)$ at $v\in S_b$.
\item[(v)] For the other $j$, $Q_j=0$, and $d(j)=d'(j)=0$.
\end{enumerate}
% Lemma -local-tangent-root (exists_sp4_root_vector_unique_on_weight, 33 lines) is a computation. The proofs of Lemmas -local-tangent-conditions (spTanL, spTanP), -local-tangent-cocycles (exists_finrank_two_cocycles_char_ratio, through tateRat) and -local-tangent-unramified (exists_unique_unramified_hom_of_frobenius_value, through Lemma -local-tangent-unramified-cocycle, which is existsUnique_unramified_cocycle_frobenius_value with the private exists_lattice, SD, SDsub and all_of_closed, frobFree and frobDense) are written out after their statements. Written out here: tangent_space_deformation_at_l and tangent_space_deformation_dual_numbers with their private dim_core, hom_iff, dim_zero, dim_pos (at l), dim_root, ex_root and wt_fib.
\end{proof}

\begin{lemma}\label{sha-counts-vd-model-local-charts}
Let $D$ be a Siegel datum, with the notation of Definitions~\ref{sha-siegel-datum} and~\ref{sha-counts-vd-model-framing}, let $v\in T$, and let $B_v$ and $\rho_v$ be as in Lemma~\ref{sha-counts-vd-model-universal}(1).
Then there are $x$, $\mathfrak p$, $y$ and $\varphi$ as in Lemma~\ref{sha-counts-vd-model-universal}(2), with $d(j)$ variables of weight $j$ for every $j$ if $v=\ell$, and with $d'(j)$ variables of weight $j$ for every $j$ if $v\in S_b$.
\end{lemma}
\leanref{Development/AutomorphicArithmetic/AutomorphicVarieties/PadicHodgeTheory/SpChartL.lean}{UpperTheory.V.Loc.spChartL}
\leanref{Development/AutomorphicArithmetic/AutomorphicVarieties/PadicHodgeTheory/SpChartP.lean}{UpperTheory.V.Loc.spChartP}
\begin{proof}
Let $E$ be the fraction field of $\mathcal O$, a complete discrete valuation ring of characteristic $0$ that is finite over $\mathbf Z_\ell$ (Definition~\ref{sha-character-parts}). Let $Q$ be the property of satisfying the conditions at $v$ of Definition~\ref{sha-counts-vd-model-framing}(1), for homomorphisms $D_v\to M_4(A)$ with $A$ a commutative $\mathcal O$-algebra.
\begin{enumerate}
\item[(i)] $Q$ is preserved by homomorphisms of $\mathcal O$-algebras and reflected by injective ones. The conditions say that polynomial expressions in the entries, with coefficients in $\mathcal O$, vanish: the entries of $\rho(g)^*\rho(g)-\theta(g)$; at $\ell$, the coefficients of the characteristic polynomial of $\rho(s_0)$ minus those of $\prod_j(X-\delta_j(s_0))$, and the entries of $e_j\rho(g)e_i$ for $i<j$ and of $e_i(\rho(s)-\delta_i(s))e_i$ for $s\in I_\ell$, where $e_i=\prod_{j\ne i}(\rho(s_0)-\delta_j(s_0))$; at $p\in S_b$, the entries of $\rho(g)-1$ for $g\in I_p\cap\ker\psi$. A homomorphism of $\mathcal O$-algebras carries each of these expressions for $\rho$ to the same expression for its image, and an injective one reflects their vanishing.
\item[(ii)] The restriction $\rho_0$ of $r_0$ to $D_v$ satisfies the conditions at $v$, as in step (i) of the proof of Lemma~\ref{sha-counts-vd-model-tangent}, and $\rho_0(g)_{ab}=0$ for $a\ne b$, since $r_0$ is diagonal. The group $D_v$ is closed in the compact group $G$, so it is compact.
\item[(iii)] Let $X=\mathbf Z^2$ and $w_a=(\lambda_a,\mu_a)$, so that $w_a-w_b$ is the weight of $(a,b)$. For every $f:B_v\to R'$ and every torus element $t$ of $R'$ there is $f':B_v\to R'$ with $f'(\rho_v(g)_{ab})=t(w_a-w_b)f(\rho_v(g)_{ab})$ for all $g$, $a$ and $b$: this is the stability under conjugation by $\tau(t_1,t_2)$, with $t_1=t(1,0)$ and $t_2=t(0,1)$, of step (iii) of the proof of Lemma~\ref{sha-counts-vd-model-universal}(1).
\item[(iv)] Lemma~\ref{sha-counts-vd-model-lift-chart} applies to $\Gamma=D_v$, $R=B_v$, $\rho=\rho_v$ and $\rho_0$: the property (a) is Lemma~\ref{sha-counts-vd-model-universal}(1), (b) is (iii), (c) is (ii), and (d) is Lemma~\ref{sha-counts-vd-model-local-lifts}. It gives $x$, $\mathfrak p$, $y$ and $\varphi$, with $x\circ\rho_v=r_0$ on $D_v$, which have the properties required in Lemma~\ref{sha-counts-vd-model-universal}(2): $\ker\varphi=\mathfrak p$, the constant terms, the approximations (2) and (3) of Lemma~\ref{sha-counts-vd-model-lift-chart}, $y\mathfrak p=0$ with $x(y)\ne0$, and, by (4) there, the entry $(a,b)$ of every $\rho_v(g)$ maps to a series of the weight of $(a,b)$.
\item[(v)] By (5) of Lemma~\ref{sha-counts-vd-model-lift-chart}, the number of variables of $\varphi$ of weight $j$ is the dimension of the space $V_j$ of Lemma~\ref{sha-counts-vd-model-local-tangent}, which is $d(j)$ if $v=\ell$ and $d'(j)$ if $v\in S_b$.
\end{enumerate}
\end{proof}

\begin{lemma}\label{sha-counts-vd-model-descent}
Let $D$ be a Siegel datum, with the notation of Definitions~\ref{sha-siegel-datum} and~\ref{sha-counts-vd-model-framing}. Let $G_S$ be the Galois group of the maximal extension of $\mathbf Q$ in $\bar{\mathbf Q}$ unramified outside $S$, and $p_S:G\to G_S$ the quotient map; for a prime $q$, let $D_q\to G_S$ be the restriction of $p_S$. Let $Q=S\setminus T$, and for $q\in Q$ let $N_q=I_q\cap\ker\bar r_0$, a subgroup of $D_q$.
\begin{enumerate}
\item There are a homomorphism $\bar r_S:G_S\to M_4(\mathbf F)$ and a continuous lift $r_{0,S}:G_S\to M_4(\mathcal O)$ of $\bar r_S$ with $\bar r_S\circ p_S=\bar r_0$ and $r_{0,S}\circ p_S=r_0$. Their values are diagonal matrices, each symplectic with some multiplier.
\item Let $R$ be a complete local noetherian $\mathcal O$-algebra with residue field $\mathbf F$, let $r:G_S\to M_4(R)$ be a continuous lift of $\bar r_S$, and put $\rho=r\circ p_S$, a continuous lift of $\bar r_0$. Then $r(g)r_{0,S}(g)^{-1}$ is symplectic with multiplier $1$ for every $g\in G_S$ if and only if $\rho(g)$ is symplectic with multiplier $\theta(g)$ for every $g\in G$. And $r$ is trivial on the image of $N_q$ for every $q\in Q$ if and only if $\rho(g)=1$ for every $g$ with $\bar r_0(g)=1$ in an inertia group at a place not above $T$.
\item Every continuous lift $\rho$ of $\bar r_0$ with values in such an $R$, with $\rho(g)=1$ for every $g$ with $\bar r_0(g)=1$ in an inertia group at a place not above $T$, is $r\circ p_S$ for a continuous lift $r$ of $\bar r_S$.
\end{enumerate}
\end{lemma}
\leanref{Development/AutomorphicArithmetic/AutomorphicVarieties/SymplecticGSp4/Descend.lean}{UpperTheory.V.GradeZeroInterpolatingElement.descend}
\leanref{Development/AutomorphicArithmetic/AutomorphicVarieties/SymplecticGSp4/ResidueCharMemLevelPrimes.lean}{UpperTheory.V.GradeZeroInterpolatingElement.Raw.frame}
\begin{proof}
The group $G_S$ is the quotient of $G$ by the closed normal subgroup $N_S$ generated by the inertia groups at the places not above $S$, with the quotient topology. A homomorphism from $G_S$ is continuous if and only if its composite with $p_S$ is. We use four facts.
\begin{enumerate}
\item[(a)] The kernel of a continuous lift $\rho$ with values in a complete local noetherian $\mathcal O$-algebra $R$ is closed. It is the intersection, over the entries $(a,b)$ and $N\ge1$, of the sets where $\rho(g)_{ab}$ is congruent to $1_{ab}$ modulo $\mathfrak m_R^N$, which are closed as the entries are locally constant modulo $\mathfrak m_R^N$; and $\bigcap_N\mathfrak m_R^N=0$.
\item[(b)] A homomorphism from $G$ whose kernel is closed and contains the inertia groups at the places not above $S$ kills $N_S$, so it factors through $p_S$. Conversely $p_S$ kills these inertia groups.
\item[(c)] For a homomorphism $f$ from $G_S$ to a group, $f$ is trivial on the images of the $N_q$, $q\in Q$, if and only if $f(p_S(g))=1$ for every $g$ with $\bar r_0(g)=1$ in an inertia group at a place not above $T$. Suppose the first, and let $g$ be such an element, in the inertia group at a place $A$. If $A$ is not above $S$, then $p_S(g)=1$ by (b). Otherwise $A$ lies above some $q\in S$, and $q\notin T$, so $q\in Q$. The places $A$ and $v_q$ lie above $q$, so there is $u\in G$ with $u^{-1}gu\in I_q$, and $\bar r_0(u^{-1}gu)=\bar r_0(u)^{-1}\bar r_0(g)\bar r_0(u)=1$; so $u^{-1}gu\in N_q$, and $f(p_S(g))$ is conjugate to $f(p_S(u^{-1}gu))=1$. Conversely, for $q\in Q$ the place $v_q$ lies above no prime of $T$, since a place lies above only one prime; so the second condition applies to the elements of $N_q$.
\item[(d)] Let $s\in G_S$, $\nu\in\mathcal O$, and suppose that $r_{0,S}(s)$ is symplectic with multiplier $\nu$. Then $r(s)r_{0,S}(s)^{-1}$ is symplectic with multiplier $1$ if and only if $r(s)$ is symplectic with multiplier $\nu$. Indeed, for an invertible $Y$ with $Y^*Y=\nu$, $(XY^{-1})^*(XY^{-1})=(Y^{-1})^*(X^*X)Y^{-1}$, which is $1$ if and only if $X^*X=Y^*Y=\nu$. Moreover $r(s)r_{0,S}(s)^{-1}$ is congruent to $1$ modulo $\mathfrak m_R$, as $r$ and $r_{0,S}$ both lift $\bar r_S$.
\end{enumerate}
(1) The characters $\delta_i$ are continuous and unramified outside $S$: $\psi$ and $\xi$ are, by Definition~\ref{sha-siegel-datum}(1); $\chi_{\mathrm{cyc}}$ is unramified outside $\ell\in S$; and $\omega$ takes values in the $(\ell-1)$-st roots of unity, which are determined by their residues, so $\omega$ is determined by $\bar\omega=\bar\chi_{\mathrm{cyc}}$. So $r_0=\operatorname{diag}(\delta_1,\delta_2,\delta_3,\delta_4)$ is a continuous lift of $\bar r_0$ whose kernel contains the inertia groups at the places not above $S$, and it is closed by (a). By (b), $r_0=r_{0,S}\circ p_S$ for a continuous homomorphism $r_{0,S}$; let $\bar r_S$ be its reduction, so that $\bar r_S\circ p_S=\bar r_0$. The values $r_0(g)$ are diagonal, and symplectic with multiplier $\theta(g)$, as $\delta_i\delta_{5-i}=\theta$; so are their reductions.

(2) As $p_S$ is continuous, $\rho$ is a continuous lift of $\bar r_0$. For $g\in G$, $r_{0,S}(p_S(g))=r_0(g)$ is symplectic with multiplier $\theta(g)$, so by (d), $r(p_S(g))r_{0,S}(p_S(g))^{-1}$ is symplectic with multiplier $1$ if and only if $\rho(g)$ is symplectic with multiplier $\theta(g)$; as $p_S$ is surjective, this gives the first equivalence. The second is (c), applied to $f=r$.

(3) Let $A$ be a place not above $S$ and $g\in I_A$. Then $r_0(g)=1$ by (1), so $\bar r_0(g)=1$, and $A$ is not above $T\subseteq S$; so $\rho(g)=1$. The kernel of $\rho$ is closed by (a), so by (b) $\rho=r\circ p_S$ for a homomorphism $r$, which is continuous and lifts $\bar r_S$, since $\bar r_S\circ p_S=\bar r_0$ and $p_S$ is surjective.
\end{proof}

\begin{definition}\label{sha-counts-vd-model-framed-cohomology}
Let $D$ be a Siegel datum, and let $G_S$, $p_S$, the maps $D_v\to G_S$ for $v\in T$ and $D_q\to G_S$ for $q\in Q$, and the subgroups $N_q$ be as in Lemma~\ref{sha-counts-vd-model-descent}. Let $\bar r_S:G_S\to M_4(\mathbf F)$ be a homomorphism with $\bar r_S\circ p_S=\bar r_0$. For $j\in\mathbf Z^2$ let $\mathfrak g_j$ be the $\mathbf F$-vector space of the $Y\in\mathfrak{sp}_4(\mathbf F)$ whose entries $(a,b)$ vanish unless $(a,b)$ has weight $j$, with the action $g\cdot Y=\bar r_S(g)Y\bar r_S(g)^{-1}$ of $G_S$, and of $D_v$ and $D_q$ through their maps to $G_S$. All functions below are locally constant, and for a function $c$ on a group $H$ acting on $\mathfrak g_j$ let $dc(h,h')=c(h)+h\cdot c(h')-c(hh')$.
\begin{enumerate}
\item $Z^1_j$ is the $\mathbf F$-vector space of the pairs $(z,(y_v)_{v\in T})$ of a function $z:G_S\to\mathfrak g_j$ and elements $y_v\in\mathfrak g_j$ such that $z(gh)=z(g)+g\cdot z(h)$ for all $g,h\in G_S$; $z$ vanishes on the image of $N_q$ for $q\in Q$; $z(g)=g\cdot y_v-y_v$ for $v\in T$ and $g\in D_v$; and $y_\ell=0$.
\item $C^2_j$ is the $\mathbf F$-vector space of the families $(a,(b_v)_{v\in T},(b_q)_{q\in Q})$ of functions $a:G_S\times G_S\to\mathfrak g_j$, $b_v:D_v\to\mathfrak g_j$ and $b_q:D_q\to\mathfrak g_j$ such that $g\cdot a(h,k)+a(g,hk)=a(gh,k)+a(g,h)$ for all $g,h,k\in G_S$, and the restriction of $a$ to $D_v\times D_v$ is $db_v$ and to $D_q\times D_q$ is $db_q$. Let $B^2_j$ be the intersection with $C^2_j$ of the span of the families $\big(dc,(c|_{D_v}-(g\mapsto g\cdot y_v-y_v))_v,(c|_{D_q}-e_q)_q\big)$, where $c:G_S\to\mathfrak g_j$, $y_v\in\mathfrak g_j$, and $e_q:D_q\to\mathfrak g_j$ satisfies $e_q(gh)=e_q(g)+g\cdot e_q(h)$ and vanishes on $N_q$. Put $\mathrm{Ob}_j=C^2_j/B^2_j$.
\end{enumerate}
\end{definition}
\leanref{Development/AutomorphicArithmetic/AutomorphicVarieties/CommutativeAlgebra/FramedTwoCoboundaries.lean}{UpperTheory.V.Alg.FormData.gw}
\leanref{Development/AutomorphicArithmetic/AutomorphicVarieties/CommutativeAlgebra/FramedTwoCoboundaries.lean}{UpperTheory.V.Alg.Frame.H1c}
\leanref{Development/AutomorphicArithmetic/AutomorphicVarieties/CommutativeAlgebra/FramedTwoCoboundaries.lean}{UpperTheory.V.Alg.Frame.framedTwoCocycles}
\leanref{Development/AutomorphicArithmetic/AutomorphicVarieties/CommutativeAlgebra/FramedTwoCoboundaries.lean}{UpperTheory.V.Alg.Frame.framedTwoCoboundaries}
\leanref{Development/AutomorphicArithmetic/AutomorphicVarieties/CommutativeAlgebra/FramedTwoCoboundaries.lean}{UpperTheory.V.Alg.Frame.Ob}

\begin{definition}\label{sha-counts-vd-model-place-datum}
Let $\Gamma$ be a topological group.
\begin{enumerate}
\item A \emph{place datum} $F$ for $\Gamma$ consists of finite sets $T$ and $Q$, an element $v_0\in T$, topological groups $G_v$ for $v\in T$ and $G_q$ for $q\in Q$, continuous homomorphisms $\varphi_v:G_v\to\Gamma$ and $\psi_q:G_q\to\Gamma$, and subgroups $N_q\subseteq G_q$.
\item Let $F$ be a place datum for $\Gamma$, let $k$ be a field, and let $M$ be a $k$-vector space with a $k$-linear action of $\Gamma$, on which $G_v$ and $G_q$ act through $\varphi_v$ and $\psi_q$. For a function $c$ on a group $H$ acting on $M$ let $dc(h,h')=c(h)+h\cdot c(h')-c(hh')$. All functions below are locally constant. $Z^1_F(M)$ is the $k$-vector space of the pairs $(z,(y_v)_{v\in T})$ of a function $z:\Gamma\to M$ and elements $y_v\in M$ such that $z(gh)=z(g)+g\cdot z(h)$ for all $g,h\in\Gamma$; $z(\psi_q(n))=0$ for $q\in Q$ and $n\in N_q$; $z(\varphi_v(s))=\varphi_v(s)\cdot y_v-y_v$ for $v\in T$ and $s\in G_v$; and $y_{v_0}=0$. $C^2_F(M)$ is the $k$-vector space of the families $(a,(b_v)_{v\in T},(b_q)_{q\in Q})$ of functions $a:\Gamma\times\Gamma\to M$, $b_v:G_v\to M$ and $b_q:G_q\to M$ such that $g\cdot a(h,h')+a(g,hh')=a(gh,h')+a(g,h)$ for all $g,h,h'\in\Gamma$, and $a(\varphi_v(s),\varphi_v(t))=db_v(s,t)$ and $a(\psi_q(s),\psi_q(t))=db_q(s,t)$ for all $s$ and $t$. $B^2_F(M)$ is the intersection with $C^2_F(M)$ of the span of the families $\big(dc,(c\circ\varphi_v-(s\mapsto\varphi_v(s)\cdot y_v-y_v))_v,(c\circ\psi_q-e_q)_q\big)$, where $c:\Gamma\to M$, $y_v\in M$, and $e_q:G_q\to M$ satisfies $e_q(st)=e_q(s)+s\cdot e_q(t)$ and vanishes on $N_q$. Put $\mathrm{Ob}_F(M)=C^2_F(M)/B^2_F(M)$.
\item Let $M$ be as in (2), with the discrete topology, and suppose that the action $\Gamma\times M\to M$ is continuous. For a topological group $H$ acting continuously on $M$, $H^i(H,M)$ is the cohomology of the complex of locally constant cochains $H\times\dots\times H\to M$, and for a subgroup $N$ of $H$, $H^1_N(H,M)$ is the kernel of the restriction $H^1(H,M)\to H^1(N,M)$. The \emph{Selmer group} $\mathrm{Sel}_F(M)$ is the subspace of the classes in $H^1(\Gamma,M)$ whose restriction to $G_v$ vanishes for every $v\in T$ and whose restriction to $G_q$ lies in $H^1_{N_q}(G_q,M)$ for every $q\in Q$; and $\text{Ш}^2_F(M)$ is the kernel of the restriction $H^2(\Gamma,M)\to\bigoplus_{v\in T}H^2(G_v,M)\oplus\bigoplus_{q\in Q}H^2(G_q,M)$.
\end{enumerate}
With the notation of Lemma~\ref{sha-counts-vd-model-descent}, the sets $T$ and $Q=S\setminus T$, $v_0=\ell$, the maps $D_v\to G_S$ and $D_q\to G_S$ and the subgroups $N_q=I_q\cap\ker\bar r_0$ of $D_q$ form a place datum $F_S$ for $G_S$, and the spaces $Z^1_j$ and $\mathrm{Ob}_j$ of Definition~\ref{sha-counts-vd-model-framed-cohomology} are $Z^1_{F_S}(\mathfrak g_j)$ and $\mathrm{Ob}_{F_S}(\mathfrak g_j)$.
\end{definition}
\leanref{Development/AutomorphicArithmetic/AutomorphicVarieties/CommutativeAlgebra/FramedTwoCoboundaries.lean}{UpperTheory.V.Alg.Frame}
\leanref{Development/AutomorphicArithmetic/AutomorphicVarieties/LevelCohomology/HcDefs.lean}{UpperTheory.V.X.selF}
\leanref{Development/AutomorphicArithmetic/AutomorphicVarieties/LevelCohomology/HcDefs.lean}{UpperTheory.V.X.shaF}
\leanref{Development/AutomorphicArithmetic/AutomorphicVarieties/SymplecticGSp4/ResidueCharMemLevelPrimes.lean}{UpperTheory.V.GradeZeroInterpolatingElement.Raw.frame}

\begin{lemma}\label{sha-counts-vd-model-selmer}
Let $\Gamma$ be a topological group with a place datum $F$ (Definition~\ref{sha-counts-vd-model-place-datum}), let $k$ be a field, and let $M$ be a finite-dimensional $k$-vector space with a $k$-linear action of $\Gamma$ that is continuous for the discrete topology on $M$ and trivial on $\psi_q(N_q)$ for every $q\in Q$. Suppose that $H^1(\Gamma,M)$ is finite-dimensional.
Then $Z^1_F(M)$ is finite-dimensional, and
\[
 \dim_kH^0(\Gamma,M)+\dim_kZ^1_F(M)=\sum_{v\in T}\dim_kH^0(G_v,M)+\dim_k\mathrm{Sel}_F(M).
\]
\end{lemma}
\leanref{Development/AutomorphicArithmetic/AutomorphicVarieties/LevelCohomology/FiniteCompactlySupportedH1.lean}{UpperTheory.V.X.framed_h1_euler_formula}
\leanref{Development/AutomorphicArithmetic/AutomorphicVarieties/CommutativeAlgebra/KrullDimQuotientComapEq.lean}{UpperTheory.V.Alg.finrank_compact_support_H1}
\begin{proof}
The space $Z^1_F(M)$ lies in the product of $M^T$ and the space of locally constant cocycles $\Gamma\to M$, which is finite-dimensional, being an extension of $H^1(\Gamma,M)$ by the image of $M$ under the coboundary map. So $Z^1_F(M)$ is finite-dimensional.
For $q\in Q$, a class in $H^1(G_q,M)$ represented by a cocycle $e$ lies in $H^1_{N_q}(G_q,M)$ if and only if $e$ vanishes on $N_q$: on $N_q$ every coboundary $n\mapsto n\cdot m-m$ vanishes, as $N_q$ acts trivially.
Consider the maps
\[
 0\to H^0(\Gamma,M)\xrightarrow{\ \iota\ }\bigoplus_{v\in T}H^0(G_v,M)\xrightarrow{\ \beta\ }Z^1_F(M)\xrightarrow{\ p\ }\mathrm{Sel}_F(M)\to0,
\]
where $\iota(m)=(m)_v$, $\beta(b)=\big(g\mapsto b_{v_0}-g\cdot b_{v_0},(b_v-b_{v_0})_v\big)$ and $p(z,(y_v))$ is the class of $z$. They are well defined: $\beta(b)$ satisfies the conditions of $Z^1_F(M)$ because $b_v$ is fixed by $G_v$, and $p(z,(y_v))$ lies in $\mathrm{Sel}_F(M)$ because $z\circ\varphi_v$ is the coboundary of $y_v$ and $z\circ\psi_q$ vanishes on $N_q$. The sequence is exact.
\begin{enumerate}
\item[(i)] $\iota$ is injective, and $\beta\circ\iota=0$. If $\beta(b)=0$, then $b_v=b_{v_0}$ for every $v$, and $b_{v_0}$ is fixed by $\Gamma$; so $b=\iota(b_{v_0})$.
\item[(ii)] $p\circ\beta=0$, the first component of $\beta(b)$ being a coboundary. If $p(z,(y_v))=0$, then $z(g)=g\cdot m-m$ for some $m\in M$. Put $b_v=y_v-m$; it is fixed by $G_v$, since $\varphi_v(s)\cdot y_v-y_v=z(\varphi_v(s))=\varphi_v(s)\cdot m-m$. As $y_{v_0}=0$, $\beta(b)=(z,(y_v))$.
\item[(iii)] Let a class in $\mathrm{Sel}_F(M)$ be represented by $z$. Then $z(\varphi_v(s))=\varphi_v(s)\cdot m_v-m_v$ for some $m_v\in M$, and $z$ vanishes on the $\psi_q(N_q)$. So $\big(g\mapsto z(g)-(g\cdot m_{v_0}-m_{v_0}),(m_v-m_{v_0})_v\big)$ lies in $Z^1_F(M)$, and $p$ maps it to the class of $z$.
\end{enumerate}
The equality follows from the exactness.
\end{proof}

\begin{lemma}\label{sha-counts-vd-model-obstruction}
Let $\Gamma$, $F$, $k$ and $M$ be as in Lemma~\ref{sha-counts-vd-model-selmer}, and suppose that $H^1(\Gamma,M)$, $H^2(\Gamma,M)$, the $H^1(G_v,M)$ for $v\in T$ and the $H^1(G_q,M)$ for $q\in Q$ are finite-dimensional.
Then $\mathrm{Ob}_F(M)$ is finite-dimensional, and
\begin{multline*}
 \dim_k\mathrm{Ob}_F(M)+\dim_kH^1(\Gamma,M)+\sum_{q\in Q}\dim_kH^1_{N_q}(G_q,M)\\
 =\dim_k\text{Ш}^2_F(M)+\sum_{v\in T}\dim_kH^1(G_v,M)+\sum_{q\in Q}\dim_kH^1(G_q,M)+\dim_k\mathrm{Sel}_F(M).
\end{multline*}
\end{lemma}
\leanref{Development/AutomorphicArithmetic/AutomorphicVarieties/LevelCohomology/HcFrameOb.lean}{UpperTheory.V.X.frameOb}
\begin{proof}
The families of Definition~\ref{sha-counts-vd-model-place-datum}(2) that span $B^2_F(M)$ form a subspace and lie in $C^2_F(M)$: $d(dc)=0$, $d(c\circ\varphi_v)=(dc)\circ\varphi_v$, and coboundaries and the $e_q$ are cocycles. So $B^2_F(M)$ is the set of these families.
Let $W=\bigoplus_{v\in T}Z^1(G_v,M)\oplus\bigoplus_{q\in Q}Z^1(G_q,M)$, where $Z^1$ denotes the space of locally constant cocycles, let $L=\bigoplus_{v\in T}H^1(G_v,M)\oplus\bigoplus_{q\in Q}H^1(G_q,M)/H^1_{N_q}(G_q,M)$, and let $\pi:W\to L$ map cocycles to their classes, which is surjective. Let $f:H^1(\Gamma,M)\to L$ be the restriction, let $g:W\to\mathrm{Ob}_F(M)$ map $((a_v),(a_q))$ to the class of $(0,(a_v),(a_q))$, which lies in $C^2_F(M)$, and let $h:\mathrm{Ob}_F(M)\to H^2(\Gamma,M)$ map the class of $(a,(b_v),(b_q))$ to the class of $a$, which is well defined since the first component of a family in $B^2_F(M)$ is a coboundary.
\begin{enumerate}
\item[(i)] $\ker f=\mathrm{Sel}_F(M)$, by definition.
\item[(ii)] The image of $h$ is $\text{Ш}^2_F(M)$. The restriction of $a$ to $G_v$ is the coboundary of $b_v$, and to $G_q$ that of $b_q$; so the image lies in $\text{Ш}^2_F(M)$. Conversely, if the class of a cocycle $a$ lies in $\text{Ш}^2_F(M)$, its restrictions are coboundaries $db_v$ and $db_q$ of locally constant functions, and $(a,(b_v),(b_q))\in C^2_F(M)$ maps to it.
\item[(iii)] $\ker h$ is the image of $g$. The first component of $(0,(a_v),(a_q))$ is $0$. Conversely, if $h$ kills the class of $(a,(b_v),(b_q))$, then $a=dc$ for a locally constant $c$; the functions $b_v-c\circ\varphi_v$ and $b_q-c\circ\psi_q$ are cocycles, and $(a,(b_v),(b_q))$ minus their image $(0,(b_v-c\circ\varphi_v),(b_q-c\circ\psi_q))$ is $(dc,(c\circ\varphi_v),(c\circ\psi_q))$, which lies in $B^2_F(M)$, with $y_v=0$ and $e_q=0$.
\item[(iv)] $\ker g=\pi^{-1}(f(H^1(\Gamma,M)))$. If $g((a_v),(a_q))=0$, then $(0,(a_v),(a_q))=\big(dc,(c\circ\varphi_v-dy_v)_v,(c\circ\psi_q-e_q)_q\big)$ for some $c$, $y_v$ and $e_q$, where $dy_v$ denotes $s\mapsto\varphi_v(s)\cdot y_v-y_v$. Then $c$ is a cocycle, $a_v$ and $c\circ\varphi_v$ have the same class, and $c\circ\psi_q-a_q=e_q$ vanishes on $N_q$, so the class of $a_q$ is that of $c\circ\psi_q$ modulo $H^1_{N_q}(G_q,M)$; hence $\pi((a_v),(a_q))=f(\text{class of }c)$. Conversely, if $\pi((a_v),(a_q))=f(\text{class of }c)$ for a cocycle $c$, then $c\circ\varphi_v-a_v=dy_v$ for some $y_v$, and the class of $e_q=c\circ\psi_q-a_q$ lies in $H^1_{N_q}(G_q,M)$, so that $e_q$ vanishes on $N_q$, as shown in the proof of Lemma~\ref{sha-counts-vd-model-selmer}; then $(0,(a_v),(a_q))$ is the family of $c$, the $y_v$ and the $e_q$, and $g((a_v),(a_q))=0$.
\end{enumerate}
By (iii) and (iv), $\ker h$ is isomorphic to $W/\ker g$, which $\pi$ maps isomorphically onto $L/f(H^1(\Gamma,M))$. These spaces are finite-dimensional, and so is $\mathrm{Ob}_F(M)$, by (ii). By (i), $\dim f(H^1(\Gamma,M))=\dim H^1(\Gamma,M)-\dim\mathrm{Sel}_F(M)$, and $\dim L=\sum_v\dim H^1(G_v,M)+\sum_q\big(\dim H^1(G_q,M)-\dim H^1_{N_q}(G_q,M)\big)$. So $\dim\mathrm{Ob}_F(M)=\dim\text{Ш}^2_F(M)+\dim L-\dim H^1(\Gamma,M)+\dim\mathrm{Sel}_F(M)$, which is the stated equality.
\end{proof}

\begin{lemma}\label{sha-counts-vd-model-unramified}
Let $H$ be a topological group, let $k$ be a field, and let $M$ be a $k$-vector space with a $k$-linear action of $H$ that is continuous for the discrete topology on $M$; use the notation of Definition~\ref{sha-counts-vd-model-place-datum}(3).
\begin{enumerate}
\item Suppose that $k$ has characteristic $\ell>0$, and let $N\subseteq I$ be subgroups of $H$ such that $[I:N]$ is finite and prime to $\ell$. Then $H^1_N(H,M)=H^1_I(H,M)$.
\item Suppose that $M$ is finite, and let $I$ be a normal subgroup of $H$ and $\phi\in H$ such that $I$ and $\phi$ generate a dense subgroup of $H$ and, for every $n\ge1$, there is a locally constant homomorphism $d:H\to\mathbf Z/n$ that is trivial on $I$ with $d(\phi)=1$. Then $\dim_kH^1_I(H,M)=\dim_kH^0(H,M)$.
\end{enumerate}
\end{lemma}
\leanref{Development/AutomorphicArithmetic/AutomorphicVarieties/LevelCohomology/HcUrIndex.lean}{UpperTheory.V.X.urIndex}
\leanref{Development/AutomorphicArithmetic/AutomorphicVarieties/LevelCohomology/HcUrAbs.lean}{UpperTheory.V.X.urAbs}
\begin{proof}
A class lies in $H^1_N(H,M)$ if and only if a cocycle $z$ representing it satisfies $z(n)=n\cdot m-m$ for some $m\in M$ and all $n\in N$.

(1) Since $N\subseteq I$, $H^1_I(H,M)\subseteq H^1_N(H,M)$. Conversely, let $z$ be a cocycle with $z(n)=n\cdot m-m$ for $n\in N$, and let $c$ be the restriction to $I$ of $g\mapsto z(g)-(g\cdot m-m)$, a cocycle on $I$ that vanishes on $N$. Then $c(gn)=c(g)+g\cdot c(n)=c(g)$ for $n\in N$, so $c$ induces a function $c'$ on $I/N$, with $c'(hx)=c(h)+h\cdot c'(x)$ for $h\in I$. Let $\Sigma=\sum_{x\in I/N}c'(x)$. Summing over $x$ gives $\Sigma=[I:N]\,c(h)+h\cdot\Sigma$, and $[I:N]$ is invertible in $k$. So $c(h)=h\cdot m'-m'$ with $m'=-[I:N]^{-1}\Sigma$, and $z(h)=h\cdot(m+m')-(m+m')$ for $h\in I$.

(2) Let $Z_0$ be the space of the locally constant cocycles $z:H\to M$ that vanish on $I$, and $M^I$ the space of the elements fixed by $I$.
\begin{enumerate}
\item[(i)] For $z\in Z_0$ and $g\in H$, $z(g)\in M^I$: for $s\in I$, $z(sg)=s\cdot z(g)$, and $sg=g(g^{-1}sg)$ with $g^{-1}sg\in I$ gives $z(sg)=z(g)$.
\item[(ii)] The map $z\mapsto z(\phi)$ from $Z_0$ to $M^I$ is bijective. If $z(\phi)=0$, then $\{g:z(g)=0\}$ is a subgroup, closed as $z$ is locally constant, which contains $I$ and $\phi$; so it is $H$, and $z=0$. Let $m\in M^I$. The map $e:x\mapsto m+\phi\cdot x$ is a permutation of the finite set $M^I$, which is stable under $\phi$ as $I$ is normal. Composing a homomorphism $d$ for the order $n$ of $e$ with $\mathbf Z/n\to\operatorname{Sym}(M^I)$, $1\mapsto e$, gives a locally constant homomorphism $w:H\to\operatorname{Sym}(M^I)$, trivial on $I$, with $w(\phi)=e$. The set of $g\in H$ with $w(g)x-w(g)x'=g\cdot(x-x')$ for all $x,x'\in M^I$ is a closed subgroup containing $I$ and $\phi$, so it is $H$. Then $z(g)=w(g)(0)$ lies in $Z_0$: $z(gh)=w(g)(w(h)(0))=w(g)(0)+g\cdot w(h)(0)$; and $z(\phi)=m$.
\item[(iii)] The image of $Z_0$ in $H^1(H,M)$ is $H^1_I(H,M)$: a class whose cocycle $z$ satisfies $z(s)=s\cdot m-m$ on $I$ is the class of $g\mapsto z(g)-(g\cdot m-m)$, which lies in $Z_0$. The coboundaries in $Z_0$ are those of the elements of $M^I$, and the coboundary map $M^I\to Z_0$ has kernel $H^0(H,M)$.
\end{enumerate}
So $\dim H^1_I(H,M)=\dim Z_0-\big(\dim M^I-\dim H^0(H,M)\big)=\dim H^0(H,M)$, by (ii).
\end{proof}

\begin{lemma}\label{sha-counts-vd-model-euler}
Let $\ell$ be an odd prime, write $G=\operatorname{Gal}(\bar{\mathbf Q}/\mathbf Q)$, and let $S$ be a finite set of primes that is the disjoint union of sets $T$ and $Q$ with $\ell\in T$. For $p\in S$ let $D_p\supseteq I_p$ be the decomposition and inertia groups at a place of $\bar{\mathbf Q}$ above $p$. Let $G_S$ be the Galois group of the maximal extension of $\mathbf Q$ in $\bar{\mathbf Q}$ unramified outside $S$, $p_S:G\to G_S$ the quotient map, and $K$ a subgroup of $G$ such that $\ell$ does not divide $[I_q:I_q\cap K]$ for $q\in Q$. Let $F$ be the place datum for $G_S$ (Definition~\ref{sha-counts-vd-model-place-datum}) given by $T$, $Q$, $v_0=\ell$, the restrictions of $p_S$ to the $D_v$ and the $D_q$, and $N_q=I_q\cap K$. Let $k$ be a finite field of characteristic $\ell$, $\bar\chi:G\to k^\times$ the cyclotomic character modulo $\ell$, and $c$ a complex conjugation. Let $M$ be a finite-dimensional $k$-vector space with a $k$-linear action of $G_S$ such that $g\mapsto g\cdot m$ is locally constant for every $m\in M$ and $p_S(I_q\cap K)$ acts trivially for every $q\in Q$.
Then $\mathrm{Ob}_F(M)$ and $Z^1_F(M)$ are finite-dimensional, and
\[
 \dim_k\mathrm{Ob}_F(M)=\dim_kZ^1_F(M)+\dim_kM^{c=1}+\dim_k\operatorname{Hom}_G(M,k(\bar\chi)),
\]
where $M^{c=1}$ is the space of the elements fixed by $p_S(c)$, and $\operatorname{Hom}_G(M,k(\bar\chi))$ is the space of the $k$-linear maps $f:M\to k$ with $f(p_S(g)\cdot m)=\bar\chi(g)f(m)$ for all $g$ and $m$.
\end{lemma}
\leanref{Development/AutomorphicArithmetic/AutomorphicVarieties/LevelCohomology/HcGalCount.lean}{UpperTheory.V.X.galCount}
\leanref{Development/AutomorphicArithmetic/AutomorphicVarieties/LevelCohomology/HcGlobEuler.lean}{UpperTheory.V.X.globEuler}
\leanref{Development/AutomorphicArithmetic/AutomorphicVarieties/LevelCohomology/HcLocEuler.lean}{UpperTheory.V.X.locEuler}
\leanref{Development/AutomorphicArithmetic/AutomorphicVarieties/LevelCohomology/HcGalTail.lean}{UpperTheory.V.X.galTail}
\leanref{Development/AutomorphicArithmetic/AutomorphicVarieties/LevelCohomology/HcPtTail.lean}{UpperTheory.V.X.ptTail}
\leanref{Development/AutomorphicArithmetic/AutomorphicVarieties/LevelCohomology/HcDecZhat.lean}{UpperTheory.V.X.decZhat}
\begin{proof}
The action of $G_S$ on $M$ is continuous for the discrete topology on $M$, as $M$ is finite-dimensional and the maps $g\mapsto g\cdot m$ are locally constant. Write $h^i=\dim H^i(G_S,M)$ and $h^i_p=\dim H^i(D_p,M)$ for $p\in S$, and $u_q=\dim H^1_{N_q}(D_q,M)$ for $q\in Q$. Five published facts are used.
\begin{enumerate}
\item[(a)] The groups $H^i(G_S,M)$ and $H^i(D_p,M)$ are finite \cite[Chapter~VII, \S2 and Chapter~VIII, \S3]{NSW08}.
\item[(b)] $h^1+\dim M^{c=1}=h^0+h^2+\dim M$, by the global Euler characteristic formula \cite[Chapter~VIII, \S7]{NSW08}, since $\ell$ is odd and $\ell\in S$.
\item[(c)] $h^1_p=h^0_p+h^2_p+\dim M$ if $p=\ell$, and $h^1_p=h^0_p+h^2_p$ if $p\ne\ell$, by the local Euler characteristic formula \cite[Chapter~VII, \S3]{NSW08}.
\item[(d)] $\dim\text{Ш}^2_F(M)+\sum_{p\in S}h^2_p=h^2+\dim\operatorname{Hom}_G(M,k(\bar\chi))$. Since $S$ is the disjoint union of $T$ and $Q$, $\text{Ш}^2_F(M)$ is the kernel of the restriction $H^2(G_S,M)\to\bigoplus_{p\in S}H^2(D_p,M)$, and by Poitou-Tate duality \cite[Chapter~VIII, \S6]{NSW08} its cokernel is dual to $H^0(G_S,\operatorname{Hom}(M,\mu_\ell))$, which has the dimension of $\operatorname{Hom}_G(M,k(\bar\chi))$; the real place contributes nothing, as $\ell$ is odd.
\item[(e)] For $q\in Q$, $D_q/I_q$ is topologically generated by the image of a Frobenius element $\phi_q$, and for every $n\ge1$ there is a locally constant homomorphism $D_q\to\mathbf Z/n$, trivial on $I_q$, that maps $\phi_q$ to $1$ \cite[Chapter~VII]{NSW08}.
\end{enumerate}
For $q\in Q$, $q\ne\ell$, and $[I_q:N_q]$ is prime to $\ell$; so by Lemma~\ref{sha-counts-vd-model-unramified}(1) and~(2), with (e), $u_q=\dim H^1_{I_q}(D_q,M)=h^0_q$. By (a), Lemmas~\ref{sha-counts-vd-model-selmer} and~\ref{sha-counts-vd-model-obstruction} apply: $Z^1_F(M)$ and $\mathrm{Ob}_F(M)$ are finite-dimensional, $\dim Z^1_F(M)+h^0=\sum_{v\in T}h^0_v+\dim\mathrm{Sel}_F(M)$, and
\[
 \dim\mathrm{Ob}_F(M)+h^1+\sum_{q\in Q}h^0_q=\dim\text{Ш}^2_F(M)+\sum_{v\in T}h^1_v+\sum_{q\in Q}h^1_q+\dim\mathrm{Sel}_F(M).
\]
By (c), $\sum_{v\in T}h^1_v=\sum_{v\in T}(h^0_v+h^2_v)+\dim M$, as $\ell$ lies in $T$ once, and $\sum_{q\in Q}h^1_q=\sum_{q\in Q}(h^0_q+h^2_q)$. Substituting these and $\dim\mathrm{Sel}_F(M)$,
\[
 \dim\mathrm{Ob}_F(M)=\dim\text{Ш}^2_F(M)+\sum_{p\in S}h^2_p+\dim M+\dim Z^1_F(M)+h^0-h^1,
\]
which by (d) and (b) is $\dim Z^1_F(M)+\dim\operatorname{Hom}_G(M,k(\bar\chi))+\dim M^{c=1}$.
\end{proof}

\begin{lemma}\label{sha-counts-vd-model-counts}
Let $D$ be a Siegel datum, and let $\bar r_S$, $Z^1_j$ and $\mathrm{Ob}_j$ be as in Definition~\ref{sha-counts-vd-model-framed-cohomology}, for $j\in\mathbf Z^2$. Put $\delta(0,0)=2$, $\delta(0,1)=\delta(0,-1)=1$, and $\delta(j)=0$ for the other $j\in\mathbf Z^2$.
Then for every $j$, $\mathrm{Ob}_j$ is finite-dimensional, $\dim_{\mathbf F}\mathrm{Ob}_j=\dim_{\mathbf F}Z^1_j+\delta(j)$, and $Z^1_j=0$ if $j\notin\{(0,0),(0,1),(0,-1)\}$.
\end{lemma}
\leanref{Development/AutomorphicArithmetic/AutomorphicVarieties/SymplecticGSp4/NumsZero.lean}{UpperTheory.V.GradeZeroInterpolatingElement.nums}
\leanref{Development/AutomorphicArithmetic/AutomorphicVarieties/SymplecticGSp4/ResidueCharMemLevelPrimes.lean}{UpperTheory.V.GradeZeroInterpolatingElement.Raw.Nums}
\leanref{Development/AutomorphicArithmetic/AutomorphicVarieties/LevelCohomology/HcScalDim.lean}{UpperTheory.V.X.scalDim}
\leanref{Development/AutomorphicArithmetic/AutomorphicVarieties/LevelCohomology/H1cZero.lean}{UpperTheory.V.X.h1cZero}
\leanref{Development/AutomorphicArithmetic/AutomorphicVarieties/LevelCohomology/HcUnrOf.lean}{UpperTheory.V.X.unrOf}
\leanref{Development/AutomorphicArithmetic/AutomorphicVarieties/LevelCohomology/HcSpAct.lean}{UpperTheory.V.X.spAct}
\leanref{Development/AutomorphicArithmetic/AutomorphicVarieties/LevelCohomology/HcSpDim.lean}{UpperTheory.V.X.spDim}
\leanref{Development/AutomorphicArithmetic/AutomorphicVarieties/LevelCohomology/HcSpKer.lean}{UpperTheory.V.X.spKer}
\begin{proof}
Use the notation of Definition~\ref{sha-siegel-datum}. Let $j=(j_1,j_2)\in\mathbf Z^2$, and let $\gamma_j:G\to\mathbf F^\times$ be the residue of $\alpha^{j_1}\Psi^{j_2}$. Let $K=\ker\bar r_0$, and let $\bar\chi$ be the residue of $\chi_{\mathrm{cyc}}$.
\begin{enumerate}
\item[(a)] \emph{The action.} For $g\in G$ and $Y\in\mathfrak g_j$, $p_S(g)\cdot Y=\bar r_0(g)Y\bar r_0(g)^{-1}=\gamma_j(g)Y$. Indeed $\bar r_0(g)=\operatorname{diag}(\bar\delta_1(g),\dots,\bar\delta_4(g))$, so the entry $(a,b)$ of $\bar r_0(g)Y\bar r_0(g)^{-1}$ is $\bar\delta_a(g)\bar\delta_b(g)^{-1}Y_{ab}$; and $\delta_a\delta_b^{-1}=\alpha^{\lambda_a-\lambda_b}\Psi^{\mu_a-\mu_b}$, since $\delta_i=\alpha^{\lambda_i}\Psi^{\mu_i}$, while $Y_{ab}=0$ unless $(a,b)$ has weight $j$.
\item[(b)] \emph{The dimension.} The condition $Y^*=-Y$ relates the entries $(a,b)$ and $(5-b,5-a)$, which have the same weight. It leaves the entries $(1,4)$, $(2,3)$, $(3,2)$ and $(4,1)$ free, identifies the entries $(1,2)$ and $(3,4)$, $(1,3)$ and $(2,4)$, $(2,1)$ and $(4,3)$, $(3,1)$ and $(4,2)$ up to sign, and on the diagonal gives $Y_{44}=-Y_{11}$ and $Y_{33}=-Y_{22}$. So $\dim\mathfrak g_j$ is $2$ if $j=(0,0)$, $1$ if $j$ is one of $\pm(1,0)$, $\pm(-1,1)$, $\pm(1,1)$, $\pm(0,1)$, and $0$ otherwise.
\item[(c)] \emph{The kernel.} $K$ is an open normal subgroup of $G$ whose index divides the order of the group of invertible diagonal matrices over $\mathbf F$, which is prime to $\ell$; so $[I_q:I_q\cap K]$ is prime to $\ell$ for $q\in Q$. By (a), $\gamma_j$ is trivial on $K$, so $p_S(I_q\cap K)$ acts trivially on $\mathfrak g_j$ and the maps $g\mapsto g\cdot Y$ are locally constant.
\item[(d)] \emph{Values.} By Lemma~\ref{sha-flagged-residues}(3), $\gamma_j(c)=(-1)^{j_1}$; and $\bar\chi(c)=-1$, $\bar\chi(\sigma_\ell)=1$.
\end{enumerate}
By (c), Lemma~\ref{sha-counts-vd-model-euler} applies to $M=\mathfrak g_j$, with $K$, $k=\mathbf F$ and the place datum $F_S$ of Definition~\ref{sha-counts-vd-model-place-datum}. So $\mathrm{Ob}_j$ and $Z^1_j$ are finite-dimensional and $\dim\mathrm{Ob}_j=\dim Z^1_j+\dim\mathfrak g_j^{c=1}+\dim\operatorname{Hom}_G(\mathfrak g_j,\mathbf F(\bar\chi))$. By (a), $c$ acts on $\mathfrak g_j$ as $\gamma_j(c)$, so $\mathfrak g_j^{c=1}$ is $\mathfrak g_j$ if $\gamma_j(c)=1$ and $0$ otherwise; and an $f\in\operatorname{Hom}_G(\mathfrak g_j,\mathbf F(\bar\chi))$ satisfies $(\gamma_j(g)-\bar\chi(g))f=0$ for every $g$, so it vanishes if $\gamma_j(g)\ne\bar\chi(g)$ for some $g$.
\begin{enumerate}
\item[(i)] Let $j\in\{(0,0),(0,1),(0,-1)\}$. Then $\gamma_j(c)=1\ne-1=\bar\chi(c)$, so $\dim\mathfrak g_j^{c=1}=\dim\mathfrak g_j=\delta(j)$ by (b), and $\operatorname{Hom}_G(\mathfrak g_j,\mathbf F(\bar\chi))=0$.
\item[(ii)] Let $j$ be one of $\pm(1,0)$, $\pm(-1,1)$, $\pm(1,1)$. Then $\gamma_j(c)=-1\ne1$, as $\ell$ is odd, so $\mathfrak g_j^{c=1}=0$. The weight $j$ is that of a pair $(a,b)$ with $\lambda_a\ne\lambda_b$, and $\gamma_j=\bar\delta_a\bar\delta_b^{-1}$; so $\gamma_j(\sigma_v)\ne1$ for $v\in T$, by Lemma~\ref{sha-flagged-residues}(1), and $\gamma_j(\sigma_\ell)\ne\bar\chi(\sigma_\ell)$, so $\operatorname{Hom}_G(\mathfrak g_j,\mathbf F(\bar\chi))=0$. Hence $\dim\mathrm{Ob}_j=\dim Z^1_j$, and $\delta(j)=0$. It remains to show $Z^1_j=0$. Let $(z,(y_v))\in Z^1_j$ and let $\mu:\mathfrak g_j\to\mathbf F$ be linear. The function $z_\mu=\mu\circ z\circ p_S$ is locally constant, with $z_\mu(gh)=z_\mu(g)+\gamma_j(g)z_\mu(h)$ by (a).
\begin{itemize}
\item It vanishes at the elements $g\in K$ of the inertia group $I_A$ at a place $A$ not above $T$. If $A$ is not above $S$, then $p_S(g)=1$ and $z(1)=0$. Otherwise $A$ lies above some $q\in Q$, and $u^{-1}gu\in I_q\cap K=N_q$ for some $u\in G$, as $K$ is normal; so $z_\mu(u^{-1}gu)=0$, and comparing the two expansions of $z_\mu(gu)=z_\mu(u\cdot u^{-1}gu)$, with $\gamma_j(g)=1$, gives $z_\mu(g)=0$.
\item For $v\in T$ and $g\in D_v$, $z_\mu(g)=\gamma_j(g)\mu(y_v)-\mu(y_v)$.
\end{itemize}
The character $\alpha^{j_1}\Psi^{j_2}$ is one of $\alpha^{\pm1}$, $\beta^{\pm1}$, $\theta^{\pm1}$. So by the last condition of Definition~\ref{sha-siegel-datum}(1) there is $m\in\mathbf F$ with $z_\mu(g)=\gamma_j(g)m-m$ for all $g\in G$. As $y_\ell=0$, $z_\mu$ vanishes on $D_\ell$, so $(\gamma_j(\sigma_\ell)-1)m=0$, and $m=0$. So $\mu\circ z\circ p_S=0$ for every $\mu$, and $z=0$, as $p_S$ is surjective. Then $(\gamma_j(\sigma_v)-1)y_v=z(p_S(\sigma_v))=0$ gives $y_v=0$ for $v\in T$.
\item[(iii)] For the other $j$, $\mathfrak g_j=0$ by (b), so $Z^1_j=0$, $\mathfrak g_j^{c=1}=0$ and $\operatorname{Hom}_G(\mathfrak g_j,\mathbf F(\bar\chi))=0$, and $\dim\mathrm{Ob}_j=0=\dim Z^1_j+\delta(j)$.
\end{enumerate}
\end{proof}

\subsection{Presentation of the framed ring}\label{subsection-sha-framed-presentation}

\begin{definition}\label{sha-counts-vd-model-form-datum}
Let $\mathcal O$ be a local ring with residue field $\mathbf F$, let $\iota$ be a finite set, let $X$ be an abelian group, and let $w:\iota\to X$ be a function, $a\mapsto w_a$. For an $\mathcal O$-algebra $A$ and a matrix $Y$ over $\mathcal O$, write $Y_A$ for its image in $M_\iota(A)$, and $Y^{\mathrm t}$ for the transpose. Below, $R$ and $R'$ denote complete local noetherian $\mathcal O$-algebras with residue field $\mathbf F$.
\begin{enumerate}
\item A \emph{form datum} $\Phi$ for $w$ consists of an invertible matrix $J\in M_\iota(\mathcal O)$ with $J^{\mathrm t}=J$ or $J^{\mathrm t}=-J$, a set $\Lambda$ of \emph{blocks} with a map $\operatorname{bl}:\iota\to\Lambda$ and an involution $\sigma$ of $\Lambda$, and an element $c\in X$, such that $\operatorname{bl}(b)=\sigma(\operatorname{bl}(a))$ and $w_a+w_b=c$ whenever $J_{ab}\ne0$.
\item For an $\mathcal O$-algebra $A$, $\mathfrak g_\Phi(A)$ is the set of the $Y\in M_\iota(A)$ with $Y^{\mathrm t}J_A+J_AY=0$ and $Y_{ab}=0$ whenever $\operatorname{bl}(a)\ne\operatorname{bl}(b)$; for $j\in X$, $\mathfrak g_{\Phi,j}(A)$ is the set of the $Y\in\mathfrak g_\Phi(A)$ with $Y_{ab}=0$ whenever $w_a-w_b\ne j$.
\item A homomorphism $\rho:H\to M_\iota(A)$ from a group \emph{normalises} $\Phi$ if: for every $g\in H$ there is $\nu\in A$ with $\rho(g)^{\mathrm t}J_A\rho(g)=\nu J_A$; there is a homomorphism $\pi$ from $H$ to the group of permutations of $\Lambda$ with $\operatorname{bl}(a)=\pi(g)(\operatorname{bl}(b))$ whenever $\rho(g)_{ab}\ne0$; and $\rho(g)_{ab}=0$ whenever $w_a\ne w_b$. If $\bar r:H\to M_\iota(\mathbf F)$ normalises $\Phi$, then $Y\mapsto\bar r(g)Y\bar r(g)^{-1}$ maps $\mathfrak g_{\Phi,j}(\mathbf F)$ to itself, and $H$ acts on $\mathfrak g_{\Phi,j}(\mathbf F)$ in this way.
\item A matrix $h\in M_\iota(R)$ is \emph{$\Phi$-unitary} if $h^{\mathrm t}J_Rh=J_R$, $h_{ab}=0$ whenever $\operatorname{bl}(a)\ne\operatorname{bl}(b)$, and $h\equiv1$ modulo $\mathfrak m_R$.
\item Suppose that $\mathcal O$ is complete and noetherian. Let $\Gamma$ be a topological group, $\bar r:\Gamma\to M_\iota(\mathbf F)$ a homomorphism, and $r_0:\Gamma\to M_\iota(\mathcal O)$ a continuous lift of $\bar r$. A continuous lift $r$ of $\bar r$ with values in $R$ is \emph{$\Phi$-close to $r_0$} if $r(g)\,r_0(g^{-1})_R$ is $\Phi$-unitary for every $g\in\Gamma$. Let $N$ be a subgroup of $\Gamma$. A \emph{local condition} for $(\bar r,r_0)$ with \emph{tangent group} $N$ is a property $L$ of continuous lifts of $\bar r$ with values in the rings $R$ such that:
\begin{enumerate}
\item[(a)] (Closed.) For every lift $r$ with values in $R$ that is $\Phi$-close to $r_0$ there is an ideal $I$ of $R$ such that, for every homomorphism of $\mathcal O$-algebras $f:R\to R'$, the lift $f\circ r$ has $L$ if and only if $f(I)=0$.
\item[(b)] (Conjugation.) Let $r$ and $r'$ be lifts with values in $R$, with $r$ $\Phi$-close to $r_0$ and having $L$, and let $h\in M_\iota(R)$ be $\Phi$-unitary with $r'(g)h=hr(g)$ for all $g$. Then $r'$ has $L$.
\item[(c)] (Torus.) Let $r$ and $r'$ be lifts with values in $R$, with $r$ $\Phi$-close to $r_0$ and having $L$, and let $t:X\to R^\times$ be a homomorphism with $t(j)\equiv1$ modulo $\mathfrak m_R$ for every $j$, such that $r'(g)_{ab}=t(w_a-w_b)\,r(g)_{ab}$ for all $g$, $a$ and $b$. Then $r'$ has $L$.
\item[(d)] (Lifts.) Let $p:R\to R'$ be a surjective homomorphism of $\mathcal O$-algebras with $\mathfrak m_R\ker p=0$. Every lift $r'$ with values in $R'$ that is $\Phi$-close to $r_0$ and has $L$ is $p\circ r$ for a lift $r$ with values in $R$ that is $\Phi$-close to $r_0$ and has $L$.
\item[(e)] (Tangent.) Let $p$ be as in (d), and let $r_1$ and $r_2$ be lifts with values in $R$ that are $\Phi$-close to $r_0$, with $p\circ r_1=p\circ r_2$, such that $r_1$ has $L$. Then $r_2$ has $L$ if and only if $r_2(n)=r_1(n)$ for every $n\in N$.
\end{enumerate}
\item Let $\Gamma$, $\bar r$ and $r_0$ be as in (5), let $F$ be a place datum for $\Gamma$ (Definition~\ref{sha-counts-vd-model-place-datum}), let $L_q$ be a property of continuous lifts of $\bar r\circ\psi_q$ for $q\in Q$, let $B$ be a complete local noetherian $\mathcal O$-algebra with residue field $\mathbf F$, and let $\rho_v:G_v\to M_\iota(B)$ be a continuous lift of $\bar r\circ\varphi_v$ for $v\in T$. A \emph{framing} over $B$ with values in $R$ consists of a homomorphism of $\mathcal O$-algebras $i:B\to R$, a continuous lift $r$ of $\bar r$ with values in $R$, and matrices $\alpha_v\in M_\iota(R)$ for $v\in T$, such that: $r$ is $\Phi$-close to $r_0$; $r\circ\psi_q$ has $L_q$ for every $q\in Q$; each $\alpha_v$ is $\Phi$-unitary; $\alpha_{v_0}=1$; and $r(\varphi_v(s))\,\alpha_v=\alpha_v\,i(\rho_v(s))$ for $v\in T$ and $s\in G_v$. It is \emph{universal} if every framing over $B$ with values in any $R'$ is its image under a unique homomorphism of $\mathcal O$-algebras $R\to R'$.
\item Suppose that $\mathcal O$ is complete and noetherian and that $X$ is finitely generated and free. Let $\pi_B:\mathcal O[[X_s:s\in\Sigma_B]]\to B$ and $\pi:\mathcal O[[X_s:s\in\Sigma]]\to R$ be weighted presentations (Definition~\ref{sha-counts-vd-model}(1)), with $\Sigma_B\subseteq\Sigma$, let $i:B\to R$ be a homomorphism of $\mathcal O$-algebras, and let $f_u$, $u\in\kappa$, be finitely many power series in the variables of $\pi$, with weights $w_f(u)\in X$. Then $\pi$ is a \emph{presentation relative to $\pi_B$ and $i$ with relations $f_u$} if: the variables $X_s$ with $s\in\Sigma_B$ have the same weights for $\pi$ as for $\pi_B$; $\pi(g)=i(\pi_B(g))$ for every power series $g$ in these variables; each $f_u$ is a series of weight $w_f(u)$; and $\ker\pi$ is the ideal generated by $\ker\pi_B$ and the $f_u$.
\end{enumerate}
\end{definition}
\leanref{Development/AutomorphicArithmetic/AutomorphicVarieties/CommutativeAlgebra/FramedTwoCoboundaries.lean}{UpperTheory.V.Alg.FormData}
\leanref{Development/AutomorphicArithmetic/AutomorphicVarieties/CommutativeAlgebra/FramedTwoCoboundaries.lean}{UpperTheory.V.Alg.FormData.g}
\leanref{Development/AutomorphicArithmetic/AutomorphicVarieties/CommutativeAlgebra/FramedTwoCoboundaries.lean}{UpperTheory.V.Alg.FormData.gw}
\leanref{Development/AutomorphicArithmetic/AutomorphicVarieties/CommutativeAlgebra/FramedTwoCoboundaries.lean}{UpperTheory.V.Alg.FormData.Normalises}
\leanref{Development/AutomorphicArithmetic/AutomorphicVarieties/CommutativeAlgebra/FramedTwoCoboundaries.lean}{UpperTheory.V.Alg.FormData.adRep}
\leanref{Development/AutomorphicArithmetic/AutomorphicVarieties/CommutativeAlgebra/FramedTwoCoboundaries.lean}{UpperTheory.V.Alg.FormData.HH}
\leanref{Development/AutomorphicArithmetic/AutomorphicVarieties/CommutativeAlgebra/FramedTwoCoboundaries.lean}{UpperTheory.V.Alg.FormData.InH}
\leanref{Development/AutomorphicArithmetic/AutomorphicVarieties/CommutativeAlgebra/FramedTwoCoboundaries.lean}{UpperTheory.V.Alg.FormData.IsLocCond}
\leanref{Development/AutomorphicArithmetic/AutomorphicVarieties/CommutativeAlgebra/FramedTwoCoboundaries.lean}{UpperTheory.V.Alg.IsFramed}
\leanref{Development/AutomorphicArithmetic/AutomorphicVarieties/CommutativeAlgebra/FramedTwoCoboundaries.lean}{UpperTheory.V.Alg.IsUnivFramed}
\leanref{Development/AutomorphicArithmetic/AutomorphicVarieties/CommutativeAlgebra/FramedTwoCoboundaries.lean}{UpperTheory.V.Alg.IsRelPres}

\begin{situation}\label{sha-counts-vd-model-lemma-p-situation}
Let $\mathcal O$ be a complete local noetherian ring with residue field $\mathbf F$, in which $2$ is a unit. Let $\iota$ be a finite set, $X$ an abelian group, $w:\iota\to X$ a function, and $\Phi$ a form datum for $w$ (Definition~\ref{sha-counts-vd-model-form-datum}). Let $\Gamma$ be a topological group, $\bar r:\Gamma\to M_\iota(\mathbf F)$ a homomorphism and $r_0:\Gamma\to M_\iota(\mathcal O)$ a continuous lift of $\bar r$, both normalising $\Phi$; for $j\in X$ write $\mathfrak g_{\Phi,j}$ for $\mathfrak g_{\Phi,j}(\mathbf F)$ with the action of $\Gamma$ through $\bar r$. Let $F$ be a place datum for $\Gamma$ (Definition~\ref{sha-counts-vd-model-place-datum}) such that $\bar r$ is trivial on $\psi_q(N_q)$ for $q\in Q$, and for $q\in Q$ let $L_q$ be a local condition for $(\bar r\circ\psi_q,r_0\circ\psi_q)$ with tangent group $N_q$ that $r_0\circ\psi_q$ has. Let $B$ be a complete local noetherian $\mathcal O$-algebra with residue field $\mathbf F$, and for $v\in T$ let $\rho_v:G_v\to M_\iota(B)$ be a continuous lift of $\bar r\circ\varphi_v$ that is $\Phi$-close to $r_0\circ\varphi_v$. Framings are framings over $B$ for these data (Definition~\ref{sha-counts-vd-model-form-datum}(6)), and $R$, $R'$, $A$, $A'$ denote complete local noetherian $\mathcal O$-algebras with residue field $\mathbf F$.
\end{situation}
\leanref{Development/AutomorphicArithmetic/AutomorphicVarieties/CommutativeAlgebra/ResMapBijective.lean}{UpperTheory.V.Alg.PhiCond}
\leanref{Development/AutomorphicArithmetic/AutomorphicVarieties/CommutativeAlgebra/ResMapBijective.lean}{UpperTheory.V.Alg.Lift}
\leanref{Development/AutomorphicArithmetic/AutomorphicVarieties/CommutativeAlgebra/ResMapBijective.lean}{UpperTheory.V.Alg.Cond}

\begin{lemma}\label{sha-counts-vd-model-lemma-p-universal}
In Situation~\ref{sha-counts-vd-model-lemma-p-situation}, suppose that $\mathbf F$ is finite of characteristic a prime $\ell$, and that $\Gamma$ is compact, totally disconnected and Hausdorff and satisfies the finiteness condition of Mazur: every open subgroup of $\Gamma$ has only finitely many continuous homomorphisms to $\mathbf Z/\ell$.
Then there is a universal framing $(i,r,\alpha)$ with values in some $R$.
The hypotheses of the situation that $2$ is a unit of $\mathcal O$, that $\bar r$ and $r_0$ normalise $\Phi$, that $\bar r$ is trivial on the $\psi_q(N_q)$, and that the $\rho_v$ are $\Phi$-close to the $r_0\circ\varphi_v$ are not used in the proof.
\end{lemma}
\leanref{Development/AutomorphicArithmetic/AutomorphicVarieties/CommutativeAlgebra/FramedUniv.lean}{UpperTheory.V.Alg.framed_univ}
\leanref{Development/AutomorphicArithmetic/AutomorphicVarieties/CommutativeAlgebra/UnivLift.lean}{UpperTheory.V.Alg.univ_lift}
\begin{proof}
The proof has five steps.
\begin{enumerate}
\item[(i)] \emph{The universal lift.} The homomorphism $\bar r$ is locally constant, being the reduction of the continuous lift $r_0$. By the finiteness condition of Mazur, the continuous lifts of $\bar r$ are the images of a universal continuous lift $u:\Gamma\to M_\iota(R_0)$ \cite{Maz89}: every continuous lift with values in an $R$ is the image of $u$ under a unique homomorphism of $\mathcal O$-algebras $R_0\to R$.
\item[(ii)] \emph{A coproduct.} Choose generators $y_1,\dots,y_n$ of $\mathfrak m_{R_0}$. The homomorphism of $\mathcal O$-algebras $\mathcal O[[X_1,\dots,X_n]]\to R_0$ with $X_k\mapsto y_k$ is surjective by the complete Nakayama lemma \cite[Theorem~8.4]{Mat86}; let $K$ be its kernel. Let $S$ be the quotient of $B[[X_1,\dots,X_n]]$ by the ideal generated by $K$, with the homomorphisms $j:R_0\to S$ induced by $\mathcal O[[X_1,\dots,X_n]]\to B[[X_1,\dots,X_n]]$ and $B\to S$ induced by the constants. The residue maps give a homomorphism $S\to\mathbf F$, so $S$ is a complete local noetherian $\mathcal O$-algebra with residue field $\mathbf F$. A pair of homomorphisms $R_0\to R$ and $B\to R$ is the composite with $j$ and $B\to S$ of a unique homomorphism $S\to R$: the homomorphism $B[[X_1,\dots,X_n]]\to R$ that extends $B\to R$ and maps $X_k$ to the image of $y_k$, which lies in $\mathfrak m_R$, kills $K$.
\item[(iii)] \emph{Framing variables.} Let $A=S[[Y_{v,a,b}:v\in T,\ a,b\in\iota]]$, with $i:B\to A$, the lift $r=j\circ u$ and $\alpha_v=1+(Y_{v,a,b})_{a,b}$. Every triple $(i',r',\alpha')$ of a homomorphism $i':B\to R'$, a continuous lift $r'$ of $\bar r$ with values in $R'$, and matrices $\alpha'_v\equiv1$ modulo $\mathfrak m_{R'}$, is the image of $(i,r,\alpha)$ under a unique homomorphism $A\to R'$: by (i) and (ii) there is a unique homomorphism $S\to R'$ compatible with $i'$ and $r'$, and it extends uniquely with $Y_{v,a,b}\mapsto(\alpha'_v-1)_{ab}\in\mathfrak m_{R'}$.
\item[(iv)] \emph{The equations.} Each of the following conditions on a homomorphism $f:A\to R'$ says that $f$ kills a certain ideal of $A$: $f\circ r$ is $\Phi$-close to $r_0$; each $f(\alpha_v)$ is $\Phi$-unitary; $f(\alpha_{v_0})=1$; and $f(r(\varphi_v(s))\alpha_v)=f(\alpha_v\,i(\rho_v(s)))$ for all $v$ and $s$. Indeed each is the vanishing of the images of finitely many elements of $A$ for each $g$, $v$ or $s$: the entries of $h^{\mathrm t}J_Ah-J_A$ and the entries $h_{ab}$ with $\operatorname{bl}(a)\ne\operatorname{bl}(b)$, for $h=r(g)r_0(g^{-1})_A$ or $h=\alpha_v$, whose reduction modulo $\mathfrak m_A$ is $1$; the entries of $\alpha_{v_0}-1$; and those of $r(\varphi_v(s))\alpha_v-\alpha_v\,i(\rho_v(s))$. The residue map $A\to\mathbf F$ satisfies these conditions: it carries $r$ to $\bar r$, which is the residue of $r_0$, and each $\alpha_v$ to $1$, and $1$ is $\Phi$-unitary. So the quotient $A_1$ of $A$ by the sum of these ideals is a complete local noetherian $\mathcal O$-algebra with residue field $\mathbf F$. Write $r_1$ for the image of $r$ in $A_1$.
\item[(v)] \emph{The local conditions.} For $q\in Q$, $r_1\circ\psi_q$ is $\Phi$-close to $r_0\circ\psi_q$, so by property (a) of $L_q$ there is an ideal $I_q$ of $A_1$ such that $f\circ r_1\circ\psi_q$ has $L_q$ if and only if $f$ kills $I_q$. The residue map of $A_1$ kills $I_q$: property (a), applied to $r_0\circ\psi_q$ itself, gives an ideal $I_0$ of $\mathcal O$, which is $0$ since $r_0\circ\psi_q$ has $L_q$; so the residue of $r_0\circ\psi_q$, which is the residue of $r_1\circ\psi_q$, has $L_q$. Let $A_2$ be the quotient of $A_1$ by the sum of the $I_q$.
\end{enumerate}
The images of $i$, $r$ and $\alpha$ in $A_2$ form a framing. It is universal: a framing $(i',r',\alpha')$ with values in $R'$ is a triple as in (iii), since $\Phi$-unitary matrices are congruent to $1$; so it is the image of $(i,r,\alpha)$ under a unique homomorphism $A\to R'$, which kills the ideals of (iv) and (v) and so factors uniquely through $A_2$.
\end{proof}

\begin{lemma}\label{sha-counts-vd-model-lemma-p-grading}
In Situation~\ref{sha-counts-vd-model-lemma-p-situation}, suppose that $X$ is finitely generated and free. Let $\pi_B$ be a weighted presentation of $B$ (Definition~\ref{sha-counts-vd-model}(1)) whose variables map into $\mathfrak m_B$, such that the entry $(a,b)$ of every $\rho_v(s)$ lies in the part $B_{w_a-w_b}$. Let $(i,r,\alpha)$ be a universal framing with values in $R$.
Then $R$ has a weighted presentation whose variables map into $\mathfrak m_R$, such that $i$ maps the image in $B$ of each variable of $\pi_B$ into the part of $R$ of the weight of that variable, and the entry $(a,b)$ of every $r(g)$ and of every $\alpha_v$ lies in the part $R_{w_a-w_b}$.
The hypotheses of the situation that $2$ is a unit of $\mathcal O$, that $\bar r$ normalises $\Phi$, that $\bar r$ is trivial on the $\psi_q(N_q)$, that $r_0\circ\psi_q$ has $L_q$, and that the $\rho_v$ are $\Phi$-close to the $r_0\circ\varphi_v$ are not used in the proof.
\end{lemma}
\leanref{Development/AutomorphicArithmetic/AutomorphicVarieties/CommutativeAlgebra/FramedGrading.lean}{UpperTheory.V.Alg.framed_grading}
\begin{proof}
For a torus element $t$ of $R'$ (Lemma~\ref{sha-counts-vd-model-torus-grading}) and $M\in M_\iota(R')$, let $t\cdot M$ be the matrix with entries $t(w_a-w_b)M_{ab}$, that is $DMD^{-1}$ with $D=\operatorname{diag}(t(w_a))_a$; so $t\cdot(MN)=(t\cdot M)(t\cdot N)$.
\begin{enumerate}
\item[(i)] \emph{Images.} For a homomorphism of $\mathcal O$-algebras $f:R\to R'$, $(f\circ i,f\circ r,f(\alpha))$ is a framing. The $\Phi$-unitary matrices and $\alpha_{v_0}=1$ are preserved by $f$, and so are the identities $r(\varphi_v(s))\alpha_v=\alpha_vi(\rho_v(s))$. For $q\in Q$, property (a) of $L_q$ gives an ideal $I_q$ of $R$ such that $f'\circ r\circ\psi_q$ has $L_q$ if and only if $f'(I_q)=0$; as $r\circ\psi_q$ has $L_q$, $I_q=0$, so $f\circ r\circ\psi_q$ has $L_q$.
\item[(ii)] \emph{The torus.} Let $(i,r,\alpha)$ be a framing with values in $R'$, $t$ a torus element of $R'$, and $i':B\to R'$ a homomorphism of $\mathcal O$-algebras with $i'(\rho_v(s)_{ab})=t(w_a-w_b)\,i(\rho_v(s)_{ab})$. Then $(i',t\cdot r,t\cdot\alpha)$ is a framing. Indeed $t\cdot r_0(g)=r_0(g)$, as $r_0$ normalises $\Phi$, so $t\cdot\big(r(g)r_0(g^{-1})\big)=(t\cdot r(g))r_0(g^{-1})$; the $\Phi$-unitary matrices are stable under $h\mapsto t\cdot h$, as $DJD=t(c)J$, the entry $J_{ab}$ vanishing unless $w_a+w_b=c$; $L_q$ holds for $t\cdot r\circ\psi_q$ by property (c) of $L_q$; $t\cdot1=1$; and $t\cdot$ applied to $r(\varphi_v(s))\alpha_v=\alpha_vi(\rho_v(s))$ gives the identity for the new triple, as $t\cdot i(\rho_v(s))=i'(\rho_v(s))$.
\item[(iii)] \emph{Generators.} Let the $e$ be the elements $i(\pi_B(X_s))$, of weight that of $X_s$, the $r(g)_{ab}-r_0(g)_{ab}$, and the $(\alpha_v-1)_{ab}$, of weight $w_a-w_b$. They lie in $\mathfrak m_R$, as $\pi_B(X_s)\in\mathfrak m_B$, $r$ lifts $\bar r$, and $\alpha_v\equiv1$.
\item[(iv)] Two homomorphisms $f,f':R\to R'$ that agree on the $e$ are equal. They agree on $i(\pi_B(X_s))$, so $f\circ i=f'\circ i$, as in step (ii) of the proof of Lemma~\ref{sha-counts-vd-model-coproduct}; and $f\circ r=f'\circ r$ and $f(\alpha_v)=f'(\alpha_v)$. By (i) both carry $(i,r,\alpha)$ to the same framing, so $f=f'$ by universality.
\item[(v)] Let $f:R\to R'$ and let $t$ be a torus element of $R'$. Lemma~\ref{sha-counts-vd-model-torus-action}, applied to $\pi_B$ and $f\circ i$, gives $i':B\to R'$ with $i'(y)=t(j)f(i(y))$ for $y\in B_j$. As the entry $(a,b)$ of $\rho_v(s)$ lies in $B_{w_a-w_b}$, (ii) applies to the framing of (i), and by universality there is $f':R\to R'$ carrying $(i,r,\alpha)$ to $(i',t\cdot(f\circ r),t\cdot f(\alpha))$. Then $f'(e)=t(\operatorname{wt}(e))f(e)$ for every $e$ of (iii): for $i(\pi_B(X_s))$ as $\pi_B(X_s)$ lies in the part of the weight of $X_s$; for the others as $r_0(g)_{ab}$ and $1_{ab}$ are constants that vanish unless $w_a=w_b$, and $t(0)=1$.
\end{enumerate}
By (iii) to (v), Lemma~\ref{sha-counts-vd-model-torus-grading}(1) applies to the family of the $e$. It gives a weighted presentation of $R$ whose variables map to some of the $e$, hence into $\mathfrak m_R$, with each $e$ in the part of its weight. So $i(\pi_B(X_s))$ lies in the part of the weight of $X_s$; and $r(g)_{ab}$ and $(\alpha_v)_{ab}$ lie in $R_{w_a-w_b}$, as the constants $r_0(g)_{ab}$ and $1_{ab}$ lie in $R_0$ and vanish unless $w_a=w_b$.
\end{proof}

\begin{lemma}\label{sha-counts-vd-model-lemma-p-tangent-abstract}
Let $\mathcal O$ be a local ring with residue field $\mathbf F$, let $B$ and $R$ be complete local noetherian $\mathcal O$-algebras with residue field $\mathbf F$, and $i:B\to R$ a homomorphism of $\mathcal O$-algebras. Let $\mathbf F[\varepsilon]$ be the dual numbers, $\varepsilon^2=0$. A homomorphism of $\mathcal O$-algebras $f:R\to\mathbf F[\varepsilon]$ has the form $f(y)=\bar y+\varepsilon f'(y)$, with $\bar y$ the residue of $y$; it is \emph{trivial on $B$} if $f'\circ i=0$. Let $S$ be a set, $X$ an abelian group, $z:S\to\mathfrak m_R$ and $\operatorname{wt}:S\to X$ functions, and for $j\in X$ let $H_j$ be an $\mathbf F$-vector space with linear forms $l_{j,s}$, $s\in S$, such that:
\begin{enumerate}
\item[(a)] $l_{j,s}=0$ if $\operatorname{wt}(s)\ne j$;
\item[(b)] if $x\in H_j$ and $l_{j,s}(x)=0$ for every $s$, then $x=0$;
\item[(c)] for every $j$ and $x\in H_j$ there is $f$ trivial on $B$ with $f'(z(s))=l_{j,s}(x)$ for every $s$;
\item[(d)] for every $f$ trivial on $B$ and every $j$ there is $x\in H_j$ with $l_{j,s}(x)=f'(z(s))$ for every $s$ with $\operatorname{wt}(s)=j$;
\item[(e)] if $f$ is trivial on $B$ and $f'(z(s))=0$ for every $s$, then $f'=0$.
\end{enumerate}
Then there are a finite set $\tau$ and a map $u:\tau\to S$ such that $\mathfrak m_R$ is contained in the ideal generated by the $z(u(t))$, $i(\mathfrak m_B)$ and $\mathfrak m_R^2$; if $c_t\in\mathcal O$ and $\sum_tc_tz(u(t))$ lies in the ideal generated by $i(\mathfrak m_B)$ and $\mathfrak m_R^2$, then every $c_t\in\mathfrak m_{\mathcal O}$; and for every $j$ the number of $t$ with $\operatorname{wt}(u(t))=j$ is $\dim_{\mathbf F}H_j$.
\end{lemma}
\leanref{Development/AutomorphicArithmetic/AutomorphicVarieties/CommutativeAlgebra/FramedTangent.lean}{UpperTheory.V.Alg.framed_tangent}
\begin{proof}
For a subset $S'$ of $\mathfrak m_R$ let $N(S')$ be the ideal generated by $S'$, $i(\mathfrak m_B)$ and $\mathfrak m_R^2$.
\begin{enumerate}
\item[(i)] If $f$ is trivial on $B$ and $f'$ vanishes on $S'\subseteq\mathfrak m_R$, then $f'$ vanishes on $N(S')$: $f$ kills $S'$, $i(\mathfrak m_B)$, as $f(i(b))=\bar b$, and $\mathfrak m_R^2$, as $f(\mathfrak m_R)\subseteq\varepsilon\mathbf F$.
\item[(ii)] If every $f$ trivial on $B$ with $f'(S')=0$ has $f'=0$, then $\mathfrak m_R\subseteq N(S')$. Indeed, if $f_1$ and $f_2$ agree on $i(B)$ and on $S'$, then $y\mapsto\bar y+\varepsilon(f_1'(y)-f_2'(y))$ is a homomorphism of $\mathcal O$-algebras, trivial on $B$ and with derivative vanishing on $S'$, so $f_1=f_2$; and Lemma~\ref{sha-counts-vd-model-lift-chart-generation-span}(1) applies.
\item[(iii)] By (ii) and (e), $\mathfrak m_R\subseteq N(z(S))$. As $R$ is noetherian, the ideal generated by $z(S)$ is generated by $z(S_0)$ for a finite $S_0\subseteq S$; so $\mathfrak m_R\subseteq N(z(S_0))$.
\item[(iv)] For each $j$, the forms $l_{j,s}$ with $s\in S_0$ and $\operatorname{wt}(s)=j$ separate the points of $H_j$. Let $x\in H_j$ be killed by them, and take $f$ as in (c). Then $f'$ vanishes on $z(S_0)$, by (a) for the $s$ with $\operatorname{wt}(s)\ne j$, so on $N(z(S_0))\supseteq\mathfrak m_R$ by (i) and (iii). So $l_{j,s}(x)=f'(z(s))=0$ for every $s$, and $x=0$ by (b). Hence $H_j$ is finite-dimensional and these forms span its dual; choose a subset $\kappa_j$ of them that is a basis of the dual. Then $\#\kappa_j=\dim H_j$, and for $s\in\kappa_j$ there is $x_s\in H_j$ with $l_{j,s}(x_s)=1$ and $l_{j,s'}(x_s)=0$ for the other $s'\in\kappa_j$.
\item[(v)] Let $\tau$ be the disjoint union of the $\kappa_j$, which is finite, as it maps injectively to $S_0$, and $u$ the inclusion. The count holds.
\item[(vi)] \emph{Spanning.} Let $f$ be trivial on $B$ with $f'(z(u(t)))=0$ for every $t$, and let $s\in S$, $j=\operatorname{wt}(s)$. By (d) there is $x\in H_j$ with $l_{j,s'}(x)=f'(z(s'))$ whenever $\operatorname{wt}(s')=j$; it is killed by the basis $\kappa_j$, so $x=0$, and $f'(z(s))=0$. By (e), $f'=0$; so by (ii), $\mathfrak m_R\subseteq N(z(u(\tau)))$.
\item[(vii)] \emph{Independence.} Let $\sum_tc_tz(u(t))\in N(\emptyset)$, and let $t\in\kappa_j$. Take $f$ as in (c) for $x_{u(t)}$. By (i), $f'$ vanishes on $N(\emptyset)$, so $0=\sum_{t'}\bar c_{t'}l_{j,u(t')}(x_{u(t)})=\bar c_t$, by (a) for the $t'$ in the other $\kappa_{j'}$. So $c_t\in\mathfrak m_{\mathcal O}$.
\end{enumerate}
\end{proof}

\begin{lemma}\label{sha-counts-vd-model-lemma-p-tangent}
In Situation~\ref{sha-counts-vd-model-lemma-p-situation}, let $(i,r,\alpha)$ be a universal framing with values in $R$.
Then there are a finite set $\tau$, elements $z_t\in R$ and weights $\operatorname{wt}(t)\in X$ for $t\in\tau$, with the following properties.
\begin{enumerate}
\item For each $t$ there are $a,b\in\iota$ with $\operatorname{wt}(t)=w_a-w_b$ such that $z_t$ is the entry $(a,b)$ of $r(g)-r_0(g)_R$ for some $g\in\Gamma$ or of $\alpha_v-1$ for some $v\in T$.
\item $\mathfrak m_R$ is contained in the ideal generated by the $z_t$, $i(\mathfrak m_B)$ and $\mathfrak m_R^2$.
\item If $c_t\in\mathcal O$ and $\sum_tc_tz_t$ lies in the ideal generated by $i(\mathfrak m_B)$ and $\mathfrak m_R^2$, then $c_t\in\mathfrak m_{\mathcal O}$ for every $t$.
\item For every $j\in X$, the number of $t$ with $\operatorname{wt}(t)=j$ is $\dim_{\mathbf F}Z^1_F(\mathfrak g_{\Phi,j})$.
\end{enumerate}
The hypotheses of the situation that $2$ is a unit of $\mathcal O$, that $r_0$ normalises $\Phi$, that $\bar r$ is trivial on the $\psi_q(N_q)$, and that the $\rho_v$ are $\Phi$-close to the $r_0\circ\varphi_v$ are not used in the proof.
\end{lemma}
\leanref{Development/AutomorphicArithmetic/AutomorphicVarieties/CommutativeAlgebra/FramedTangent.lean}{UpperTheory.V.Alg.framed_tangent}
\begin{proof}
The dual numbers $\mathbf F[\varepsilon]$ form a complete local noetherian $\mathcal O$-algebra with residue field $\mathbf F$; use the notation of Lemma~\ref{sha-counts-vd-model-lemma-p-tangent-abstract}, and write a matrix over $\mathbf F[\varepsilon]$ as $M_0+\varepsilon M_1$. Let $i_0:B\to\mathbf F[\varepsilon]$ be the residue map. Write $Z^1_j=Z^1_F(\mathfrak g_{\Phi,j})$; an element of it is a pair $x=(c,(y_v))$.
\begin{enumerate}
\item[(i)] \emph{Framings over the dual numbers.} A matrix $M_0+\varepsilon M_1$ is $\Phi$-unitary if and only if $M_0=1$ and $M_1\in\mathfrak g_\Phi(\mathbf F)$, as $(1+\varepsilon M_1)^{\mathrm t}J(1+\varepsilon M_1)=J+\varepsilon(M_1^{\mathrm t}J+JM_1)$. A continuous lift of $\bar r$ with values in $\mathbf F[\varepsilon]$ is $s\mapsto\bar r(s)+\varepsilon e(s)$ with $e$ locally constant, as $\mathfrak m^2=0$ there. Let $q\in Q$. The image of $r_0\circ\psi_q$ in $\mathbf F[\varepsilon]$ has $L_q$, by property (a) of $L_q$ applied to $r_0\circ\psi_q$, which has $L_q$; so by the tangent property (e) of $L_q$ for the map $\mathbf F[\varepsilon]\to\mathbf F$, a lift $\bar r\circ\psi_q+\varepsilon e$ that is $\Phi$-close to $r_0\circ\psi_q$ has $L_q$ if and only if $e$ vanishes on $N_q$. Hence $(i_0,\bar r+\varepsilon e,(1+\varepsilon A_v)_v)$ is a framing if and only if: $e(s)\bar r(s^{-1})\in\mathfrak g_\Phi(\mathbf F)$ for every $s$; $e(\psi_q(n))=0$ for $n\in N_q$; $A_v\in\mathfrak g_\Phi(\mathbf F)$; $A_{v_0}=0$; and $\bar r(g)A_v+e(g)=A_v\bar r(g)$ for $g=\varphi_v(s)$, which is the coefficient of $\varepsilon$ in $r(g)\alpha_v=\alpha_vi_0(\rho_v(s))$.
\item[(ii)] For a homomorphism $f:R\to\mathbf F[\varepsilon]$, $(f\circ i,f\circ r,f(\alpha))$ is a framing: the $\Phi$-unitary matrices are preserved by Lemma~\ref{sha-counts-vd-model-lemma-p-unitary}(1), the identities by $f$, and $L_q$ by property (a) of $L_q$, as in step (i) of the proof of Lemma~\ref{sha-counts-vd-model-lemma-p-grading}.
\item[(iii)] \emph{The data of Lemma~\ref{sha-counts-vd-model-lemma-p-tangent-abstract}.} Let $S$ be the disjoint union of $\Gamma\times\iota\times\iota$ and $T\times\iota\times\iota$, with $z(s,a,b)=r(s)_{ab}-r_0(s)_{ab}$, $z(v,a,b)=(\alpha_v-1)_{ab}$, both of weight $w_a-w_b$; they lie in $\mathfrak m_R$, as $r$ and $r_0$ both lift $\bar r$ and $\alpha_v\equiv1$. Let $H_j=Z^1_j$, with $l_{j,(s,a,b)}(x)=(c(s)\bar r(s))_{ab}$ and $l_{j,(v,a,b)}(x)=-(y_v)_{ab}$. For $f$ trivial on $B$, $f'(z(s,a,b))$ is the entry $(a,b)$ of the coefficient of $\varepsilon$ in $f(r(s))$, as $f'$ vanishes on constants, and $f'(z(v,a,b))$ that of $f(\alpha_v)$.
\item[(iv)] Condition (a): $c(s)\in\mathfrak g_{\Phi,j}$ and the entries of $\bar r(s)$ vanish unless $w_a=w_b$, so the entry $(a,b)$ of $c(s)\bar r(s)$ vanishes unless $w_a-w_b=j$; and so does $(y_v)_{ab}$. Condition (b): if all $l_{j,s}(x)$ vanish, then $c(s)\bar r(s)=0$, so $c(s)=0$, and $y_v=0$.
\item[(v)] Condition (c). Let $x=(c,(y_v))\in Z^1_j$. Then $s\mapsto\bar r(s)+\varepsilon c(s)\bar r(s)$ is a homomorphism, by the cocycle identity $c(st)=c(s)+\bar r(s)c(t)\bar r(s)^{-1}$, and a continuous lift of $\bar r$, as $c$ is locally constant; with $A_v=-y_v$ it satisfies the conditions of (i): $c(s)\in\mathfrak g_\Phi(\mathbf F)$, $c$ vanishes on the $\psi_q(N_q)$, $y_{v_0}=0$, and $c(g)\bar r(g)=\bar r(g)y_v-y_v\bar r(g)$ for $g=\varphi_v(s)$, as $c(g)=\bar r(g)y_v\bar r(g)^{-1}-y_v$. By universality there is $f:R\to\mathbf F[\varepsilon]$ carrying $(i,r,\alpha)$ to this framing; it is trivial on $B$, and by (iii) $f'(z(s))=l_{j,s}(x)$ for every $s$.
\item[(vi)] Condition (d). Let $f$ be trivial on $B$, and write $f(r(s))=\bar r(s)+\varepsilon e(s)$ and $f(\alpha_v)=1+\varepsilon A_v$. By (ii), $(i_0,f\circ r,f(\alpha))$ is a framing, and (i) applies. For $Y\in M_\iota(\mathbf F)$ let $Y^{(j)}$ keep the entries of weight $j$ and replace the others by $0$; it commutes with multiplication by the values of $\bar r$, whose entries vanish unless $w_a=w_b$, and it maps $\mathfrak g_\Phi(\mathbf F)$ into $\mathfrak g_{\Phi,j}$, since the entry $(a,b)$ of $Y^{\mathrm t}J+JY$ involves only entries of $Y$ of weight $c-w_a-w_b$, as $J_{xy}\ne0$ implies $w_x+w_y=c$. Put $c(s)=(e(s)\bar r(s^{-1}))^{(j)}$ and $y_v=-A_v^{(j)}$. Then $c$ is locally constant, the homomorphism property of $f\circ r$ gives $c(st)=c(s)+\bar r(s)c(t)\bar r(s)^{-1}$, and (i) gives the other conditions of $Z^1_j$; so $x=(c,(y_v))\in Z^1_j$, and $l_{j,s}(x)=f'(z(s))$ for the $s$ of weight $j$, by (iii).
\item[(vii)] Condition (e). Let $f$ be trivial on $B$ with $f'(z(s))=0$ for every $s$, and let $f_0(y)=\bar y$. Then $f$ and $f_0$ carry $(i,r,\alpha)$ to the same framing, by (iii), so $f=f_0$ by the uniqueness in the universal property, and $f'=0$.
\end{enumerate}
Lemma~\ref{sha-counts-vd-model-lemma-p-tangent-abstract} gives $\tau$ and $u$; put $z_t=z(u(t))$ and $\operatorname{wt}(t)=\operatorname{wt}(u(t))$. Property (1) holds by the definition of $z$, and (2) to (4) are the conclusions of Lemma~\ref{sha-counts-vd-model-lemma-p-tangent-abstract}.
\end{proof}

\begin{lemma}\label{sha-counts-vd-model-lemma-p-relations-parts}
Let $\mathcal O$ be a complete local noetherian ring with residue field $\mathbf F$, let $X$ be a finitely generated free abelian group, and let $R$ be a complete local noetherian $\mathcal O$-algebra with residue field $\mathbf F$. Let $\pi$ and $\pi'$ be weighted presentations of $R$ (Definition~\ref{sha-counts-vd-model}(1)) whose variables map into $\mathfrak m_R$, such that each variable of $\pi'$ maps into the part of $\pi$ of its weight. Then $\pi$ and $\pi'$ have the same parts.
\end{lemma}
\leanref{Development/AutomorphicArithmetic/AutomorphicVarieties/CommutativeAlgebra/PresIndep.lean}{UpperTheory.V.Alg.pres_indep}
\begin{proof}
This is Lemma~\ref{sha-counts-vd-presentations}(e). Its proof, in steps (i) to (iii), uses of $\mathcal O$ only that it is a complete local noetherian ring with residue field $\mathbf F$: step (i) uses the theorem of Krull in $R$, step (ii) truncates power series, and step (iii) approximates elements of $R$ by finite sums of elements of the parts of $\pi'$.
\end{proof}

\begin{lemma}\label{sha-counts-vd-model-lemma-p-relations-kernel}
Let $\mathcal O$ be a complete local noetherian ring, let $X$ be an abelian group, and give the variables of $P=\mathcal O[[X_1,\dots,X_m]]$ weights in $X$. Let $J$ be an ideal of $P$ that contains, with every series, the sum of its terms of each weight. Then $J$ is generated by finitely many series, each of a single weight.
\end{lemma}
\leanref{Development/AutomorphicArithmetic/AutomorphicVarieties/CommutativeAlgebra/KerGen.lean}{UpperTheory.V.Alg.ker_gen}
\begin{proof}
For a series $f$ and $j\in X$ let $f_j$ be the sum of the terms of $f$ of weight $j$. Let $G$ be the set of the $f_j$ with $f\in J$; it lies in $J$. Conversely, let $f\in J$ and $N\ge1$. The monomials of degree less than $N$ have finitely many weights, and the sum of the $f_j$ over these weights lies in the ideal $(G)$ generated by $G$ and differs from $f$ by a series without terms of degree less than $N$, which lies in $\mathfrak m_P^N$. So $f$ lies in $\bigcap_N\big((G)+\mathfrak m_P^N\big)$, which is $(G)$ by the theorem of Krull, as $P$ is a noetherian local ring. So $J=(G)$, and as $P$ is noetherian, finitely many elements of $G$ generate it.
\end{proof}

\begin{lemma}\label{sha-counts-vd-model-lemma-p-relations}
Let $\mathcal O$ be a complete local noetherian ring with residue field $\mathbf F$, let $X$ be a finitely generated free abelian group, let $B$ and $R$ be complete local noetherian $\mathcal O$-algebras with residue field $\mathbf F$, and let $i:B\to R$ be a homomorphism of $\mathcal O$-algebras. Let $\pi_B$ and $\pi_0$ be weighted presentations of $B$ and $R$ (Definition~\ref{sha-counts-vd-model}(1)) whose variables map into the maximal ideals, such that $i$ maps the image of each variable $X_s$ of $\pi_B$ into the part of $\pi_0$ of the weight of $X_s$. Let $\tau$ be a finite set, and for $t\in\tau$ let $z_t\in\mathfrak m_R$ lie in the part of $\pi_0$ of a weight $\operatorname{wt}(t)$, such that $\mathfrak m_R$ is contained in the ideal generated by the $z_t$, $i(\mathfrak m_B)$ and $\mathfrak m_R^2$.
Then there is a weighted presentation $\pi$ of $R$, by the variables of $\pi_B$ and variables $Y_t$ for $t\in\tau$, relative to $\pi_B$ and $i$ with relations $f_u$ of weights $w_f(u)$, $u\in\kappa$ (Definition~\ref{sha-counts-vd-model-form-datum}(7)), such that $\pi(Y_t)=z_t$ and $Y_t$ has weight $\operatorname{wt}(t)$, $\pi$ and $\pi_0$ have the same parts, and the relations are minimal: if $c_u\in\mathcal O$ and $\sum_uc_uf_u$ lies in the ideal generated by $\ker\pi_B$ and $\mathfrak m\ker\pi$, where $\mathfrak m$ is the maximal ideal of the power series ring of $\pi$, then $c_u\in\mathfrak m_{\mathcal O}$ for every $u$.
\end{lemma}
\leanref{Development/AutomorphicArithmetic/AutomorphicVarieties/CommutativeAlgebra/RelPres.lean}{UpperTheory.V.Alg.rel_pres}
\begin{proof}
Write $X_s$ for the variables of $\pi_B$ and $P_0$ for the power series ring of $\pi_0$.
\begin{enumerate}
\item[(i)] $\mathfrak m_B$ is generated by the $\pi_B(X_s)$ and $\mathfrak m_{\mathcal O}$: if $\pi_B(f)\in\mathfrak m_B$, then $f-f(0)$ lies in the ideal of the variables and $f(0)\in\mathfrak m_{\mathcal O}$. So by the hypothesis, $\mathfrak m_R$ lies in the ideal generated by the $z_t$, the $i(\pi_B(X_s))$, $\mathfrak m_{\mathcal O}R$ and $\mathfrak m_R^2$.
\item[(ii)] \emph{The presentation.} Each of the elements $y=i(\pi_B(X_s))$ and $y=z_t$ lies in $\mathfrak m_R$ and in the part of $\pi_0$ of its weight $w(y)$, that of $X_s$ or $\operatorname{wt}(t)$; choose a series $f_y\in P_0$ of weight $w(y)$ with $\pi_0(f_y)=y$, which has no constant term. Let $\Phi$ map the power series ring $P$ in variables $X_s$ and $Y_t$ to $P_0$ by sending each variable to its $f_y$, and let $\pi=\pi_0\circ\Phi$. By (i) and the complete Nakayama lemma \cite[Theorem~8.4]{Mat86}, $\pi$ is surjective. As $\Phi$ maps series of weight $j$ to series of weight $j$, for $g\in\ker\pi$ the part of weight $j$ of $\Phi(g)$ is $\Phi(g_j)$, and it lies in $\ker\pi_0$; so $g_j\in\ker\pi$, and $\pi$ is a weighted presentation, with $\pi(Y_t)=z_t$ of weight $\operatorname{wt}(t)$.
\item[(iii)] By Lemma~\ref{sha-counts-vd-model-lemma-p-relations-parts}, $\pi$ and $\pi_0$ have the same parts.
\item[(iv)] On the series in the $X_s$, $\pi$ and $i\circ\pi_B$ agree: both are homomorphisms of $\mathcal O$-algebras that agree on the variables, which they map into $\mathfrak m_R$ (step (ii) of the proof of Lemma~\ref{sha-counts-vd-model-coproduct}). So $\ker\pi$ contains $\ker\pi_B$.
\item[(v)] By Lemma~\ref{sha-counts-vd-model-lemma-p-relations-kernel}, $\ker\pi$ is generated by finitely many series $g_k$ of single weights $c_k$; with (iv), it is generated by $\ker\pi_B$ and the $g_k$.
\item[(vi)] \emph{Minimal relations.} Choose a subset $\kappa$ of the indices $k$, minimal among those such that $\ker\pi$ is generated by $\ker\pi_B$ and the $g_u$, $u\in\kappa$, and put $f_u=g_u$ and $w_f(u)=c_u$. By (ii), (iv) and (v), $\pi$ is a presentation relative to $\pi_B$ and $i$ with relations $f_u$. Suppose that $\sum_uc_uf_u$ lies in the ideal generated by $\ker\pi_B$ and $\mathfrak m_P\ker\pi$, with $c_{u_0}$ a unit. Then $f_{u_0}$ lies in $N+\mathfrak m_P\ker\pi$, where $N$ is the ideal generated by $\ker\pi_B$ and the $f_u$ with $u\ne u_0$; so $\ker\pi\subseteq N+\mathfrak m_P\ker\pi$, and $N=\ker\pi$ by the lemma of Nakayama, as $\ker\pi$ is finitely generated. This contradicts the minimality of $\kappa$.
\end{enumerate}
\end{proof}

\begin{lemma}\label{sha-counts-vd-model-lemma-p-unitary}
Let $\mathcal O$ be a local ring with residue field $\mathbf F$, let $\iota$, $X$, $w$ and a form datum $\Phi$ be as in Definition~\ref{sha-counts-vd-model-form-datum}, and let $R$ and $R'$ be complete local noetherian $\mathcal O$-algebras with residue field $\mathbf F$.
\begin{enumerate}
\item The matrix $1$ is $\Phi$-unitary. A product of two $\Phi$-unitary matrices is $\Phi$-unitary, and a $\Phi$-unitary matrix has a $\Phi$-unitary inverse. The image of a $\Phi$-unitary matrix under a homomorphism of $\mathcal O$-algebras $R\to R'$ is $\Phi$-unitary. If $H$ is a group and $r_0:H\to M_\iota(\mathcal O)$ a homomorphism that normalises $\Phi$, then $r_0(s)_Rh\,r_0(s^{-1})_R$ is $\Phi$-unitary for every $s\in H$ and every $\Phi$-unitary $h\in M_\iota(R)$. If $t:X\to R^\times$ is a homomorphism and $h\in M_\iota(R)$ is $\Phi$-unitary, the matrix with entries $t(w_a-w_b)h_{ab}$ is $\Phi$-unitary.
\item Suppose that $\mathcal O$ is complete and noetherian and that $2$ is a unit of $\mathcal O$, and let $p:R\to R'$ be a surjective homomorphism of $\mathcal O$-algebras with $\mathfrak m_R\ker p=0$. Then: (a) every $\Phi$-unitary $h'\in M_\iota(R')$ is $p(h)$ for a $\Phi$-unitary $h\in M_\iota(R)$; (b) if $h_1,h_2\in M_\iota(R)$ are $\Phi$-unitary with $p(h_1)=p(h_2)$, then $h_2-h_1\in\mathfrak g_\Phi(R)$; (c) if $h\in M_\iota(R)$ is $\Phi$-unitary, $Y\in\mathfrak g_\Phi(R)$ and $p(Y)=0$, then $h+Y$ is $\Phi$-unitary.
\end{enumerate}
\end{lemma}
\leanref{Development/AutomorphicArithmetic/AutomorphicVarieties/CommutativeAlgebra/HhGroup.lean}{UpperTheory.V.Alg.hh_group}
\leanref{Development/AutomorphicArithmetic/AutomorphicVarieties/CommutativeAlgebra/HhSmall.lean}{UpperTheory.V.Alg.hh_small}
\begin{proof}
Write $J$ for $J_R$, which is invertible, with $J^{\mathrm t}=\epsilon J$ for $\epsilon=\pm1$. The conditions for $h$ to be $\Phi$-unitary are $h^{\mathrm t}Jh=J$, $h_{ab}=0$ whenever $\operatorname{bl}(a)\ne\operatorname{bl}(b)$, and $h\equiv1$ modulo $\mathfrak m_R$.

(1) The matrix $1$ satisfies them. Let $h$ and $h'$ be $\Phi$-unitary. Then $(hh')^{\mathrm t}J(hh')=h'^{\mathrm t}(h^{\mathrm t}Jh)h'=J$; an entry $(hh')_{ab}$ with $\operatorname{bl}(a)\ne\operatorname{bl}(b)$ is a sum of products $h_{ac}h'_{cb}$, each $0$ as $\operatorname{bl}(c)$ differs from $\operatorname{bl}(a)$ or from $\operatorname{bl}(b)$; and $hh'\equiv1$.
For the inverse, let $h'=J^{-1}h^{\mathrm t}J$. Then $h'h=1$, so $hh'=1$. So $h'^{\mathrm t}Jh'=h'^{\mathrm t}(h^{\mathrm t}Jh)h'=(hh')^{\mathrm t}J(hh')=J$, and $h'\equiv1$, its residue being the inverse of that of $h$. Let $a,b$ have $\operatorname{bl}(a)\ne\operatorname{bl}(b)$, and let $E$ be the diagonal matrix with $E_{xx}=1$ if $\operatorname{bl}(x)=\operatorname{bl}(a)$ and $E_{xx}=0$ otherwise. Then $Eh=hE$, by the zeros of $h$, so $h'E=h'(Eh)h'=Eh'$, and the entry $(a,b)$ of this identity is $0=h'_{ab}$.
A homomorphism $f:R\to R'$ carries $J_R$ to $J_{R'}$, and the three conditions are identities in the entries, which $f$ preserves, $f$ inducing the identity on the residue fields.
Let $A=r_0(s)_R$ and $A'=r_0(s^{-1})_R$, so that $AA'=1$ and $A^{\mathrm t}JA=\nu J$ for some $\nu$. Then $h^{\mathrm t}(A^{\mathrm t}JA)h=\nu J=A^{\mathrm t}JA$, so $(AhA')^{\mathrm t}J(AhA')=A'^{\mathrm t}A^{\mathrm t}JAA'=J$. If $(AhA')_{ab}\ne0$, then $A_{ac}$, $h_{cd}$ and $A'_{db}$ are nonzero for some $c,d$; with $\pi$ the homomorphism to the permutations of the blocks, $\operatorname{bl}(a)=\pi(s)(\operatorname{bl}(c))$, $\operatorname{bl}(c)=\operatorname{bl}(d)$ and $\operatorname{bl}(d)=\pi(s)^{-1}(\operatorname{bl}(b))$, so $\operatorname{bl}(a)=\operatorname{bl}(b)$. And $AhA'\equiv AA'=1$.
The matrix with entries $t(w_a-w_b)h_{ab}$ is $DhD^{-1}$ with $D=\operatorname{diag}(t(w_a))_a$, and $DJD=t(c)J$, as $J_{ab}\ne0$ only when $w_a+w_b=c$; so the computation above, with $A=D$ and $\nu=t(c)$, applies to it. Its entries vanish where those of $h$ do, and it is $\equiv1$, as its diagonal entries are those of $h$, since $t(0)=1$, and its other entries are multiples of those of $h$.

(2) The ideal $\ker p$ lies in $\mathfrak m_R$, and $\mathfrak m_R$ kills it.
\begin{enumerate}
\item[(i)] Let $h\equiv1$ modulo $\mathfrak m_R$, and let $Z\in M_\iota(R)$ have $p(Z)=0$. Then $(h+Z)^{\mathrm t}J(h+Z)=h^{\mathrm t}Jh+Z^{\mathrm t}J+JZ$ and $h+Z\equiv1$. Indeed the entries of $h-1$ lie in $\mathfrak m_R$ and those of $Z$ in $\ker p$, so $(h-1)^{\mathrm t}JZ$, $Z^{\mathrm t}J(h-1)$ and $Z^{\mathrm t}JZ$ vanish.
\item[(ii)] For (a), choose $h_0\in M_\iota(R)$ with $p(h_0)=h'$ and $(h_0)_{ab}=0$ whenever $\operatorname{bl}(a)\ne\operatorname{bl}(b)$, by lifting the other entries. Then $h_0\equiv1$, as $h'\equiv1$. Let $E=h_0^{\mathrm t}Jh_0-J$. Then $p(E)=0$, $E^{\mathrm t}=\epsilon E$, and $E_{ab}=0$ unless $\operatorname{bl}(b)=\sigma(\operatorname{bl}(a))$: a nonzero term $(h_0)_{da}J_{dc}(h_0)_{cb}$ of $(h_0^{\mathrm t}Jh_0)_{ab}$ has $\operatorname{bl}(d)=\operatorname{bl}(a)$, $\operatorname{bl}(c)=\operatorname{bl}(b)$ and $\operatorname{bl}(c)=\sigma(\operatorname{bl}(d))$, and $J_{ab}\ne0$ implies $\operatorname{bl}(b)=\sigma(\operatorname{bl}(a))$. Let $Z_0=-\tfrac12J^{-1}E$, and let $Z$ agree with $Z_0$ at the entries $(a,b)$ with $\operatorname{bl}(a)=\operatorname{bl}(b)$ and vanish at the others. Then $p(Z)=0$, and $JZ=-\tfrac12E$. Indeed, at an entry $(a,b)$, a term $J_{ac}Z_{cb}$ with $J_{ac}\ne0$ has $\operatorname{bl}(c)=\sigma(\operatorname{bl}(a))$; if $\operatorname{bl}(b)=\sigma(\operatorname{bl}(a))$, then $Z_{cb}=(Z_0)_{cb}$ for these terms, and the entry is $(JZ_0)_{ab}=-\tfrac12E_{ab}$; otherwise all these terms vanish, and so does $E_{ab}$. Transposing, $\epsilon Z^{\mathrm t}J=-\tfrac12\epsilon E$, so $Z^{\mathrm t}J=-\tfrac12E$. By (i), $h=h_0+Z$ has $h^{\mathrm t}Jh=J+E-E=J$ and $h\equiv1$; its entries vanish across blocks, and $p(h)=h'$.
\item[(iii)] For (b), (i) with $h=h_1$ and $Z=h_2-h_1$ gives $J=J+Z^{\mathrm t}J+JZ$, and $Z$ vanishes across blocks; so $Z\in\mathfrak g_\Phi(R)$. For (c), (i) gives $(h+Y)^{\mathrm t}J(h+Y)=J+Y^{\mathrm t}J+JY=J$ and $h+Y\equiv1$, and $h+Y$ vanishes across blocks.
\end{enumerate}
\end{proof}

\begin{lemma}\label{sha-counts-vd-model-lemma-p-classes-section-lift}
Let $\mathcal O$ be a complete local noetherian ring with residue field $\mathbf F$, and let $X$ be a finitely generated free abelian group. Let $A$ and $R$ be complete local noetherian $\mathcal O$-algebras with residue field $\mathbf F$, $p:A\to R$ a surjective homomorphism of $\mathcal O$-algebras with $\mathfrak m_A\ker p=0$, and let $A$ have a weighted presentation (Definition~\ref{sha-counts-vd-model}(1)) whose variables map into $\mathfrak m_A$, with parts $A_j$. Let $S$ be a topological space, $j\in X$, and $x:S\to R$ a function, continuous for the $\mathfrak m_R$-adic topology, with values in $p(A_j)$.
Then there is a function $y:S\to A_j$, continuous for the $\mathfrak m_A$-adic topology, with $p\circ y=x$.
\end{lemma}
\leanref{Development/AutomorphicArithmetic/AutomorphicVarieties/CommutativeAlgebra/ContSection.lean}{UpperTheory.V.Alg.cont_section}
\begin{proof}
Let $K=\ker p$, a vector space over $\mathbf F$, finite-dimensional as $A$ is noetherian; a decreasing sequence of subspaces of $K$ is eventually constant. The homomorphism $p$ maps $\mathfrak m_A^n$ onto $\mathfrak m_R^n$.
\begin{enumerate}
\item[(i)] As in Lemma~\ref{sha-counts-vd-model-lemma-p-kill-section}(1), there is $N_0$ with $K\cap\mathfrak m_A^{N_0}=0$; so $p$ is injective on $\mathfrak m_A^{N_0}$.
\item[(ii)] For $n\ge0$ let $F_n$ be the set of the $k\in K$ that lie in $A_j+\mathfrak m_A^n$. It is a subspace of $K$: for $a\in A$ and $k\in F_n$, $ak=ck$ for any $c\in\mathcal O$ with $a-c\in\mathfrak m_A$, and $A_j+\mathfrak m_A^n$ is an $\mathcal O$-module. The $F_n$ decrease, so they are equal to $F_{N_1}$ for $n\ge N_1$. Hence an element of $K\cap(A_j+\mathfrak m_A^{N_1})$ lies in $A_j+\mathfrak m_A^n$ for every $n$, so in $A_j$, by Lemma~\ref{sha-counts-vd-presentations}(e)(i).
\item[(iii)] Let $M=\max(N_0,N_1)$. If $a\in p(A_j)\cap\mathfrak m_R^M$, there is $w\in A_j\cap\mathfrak m_A^M$ with $p(w)=a$. Indeed $a=p(e)$ with $e\in A_j$, and $a=p(w)$ with $w\in\mathfrak m_A^M$; then $e-w\in K\cap(A_j+\mathfrak m_A^{N_1})$ lies in $A_j$ by (ii), and so does $w$.
\item[(iv)] For each class $C$ of $R$ modulo $\mathfrak m_R^M$ that meets $p(A_j)$, choose $e_C\in A_j$ with $p(e_C)\in C$. For $s\in S$, with $C$ the class of $x(s)$, choose by (iii) $w(s)\in A_j\cap\mathfrak m_A^M$ with $p(w(s))=x(s)-p(e_C)$, and put $y(s)=e_C+w(s)\in A_j$. Then $p(y(s))=x(s)$.
\item[(v)] $y$ is continuous. Let $n\ge0$ and $s_0\in S$. For $s$ near $s_0$, $x(s)-x(s_0)\in\mathfrak m_R^{\max(n,M)}$, so $x(s)$ and $x(s_0)$ have the same class $C$, and $x(s)-x(s_0)=p(v)$ with $v\in\mathfrak m_A^{\max(n,M)}$. Then $w(s)-w(s_0)$ and $v$ lie in $\mathfrak m_A^{N_0}$ and have the same image under $p$, so they are equal by (i); hence $y(s)-y(s_0)=v\in\mathfrak m_A^n$.
\end{enumerate}
\end{proof}

\begin{lemma}\label{sha-counts-vd-model-lemma-p-classes-section}
Let $\mathcal O$ be a complete local noetherian ring with residue field $\mathbf F$ in which $2$ is a unit, let $X$ be a finitely generated free abelian group, and let $\iota$, $w$ and a form datum $\Phi$ be as in Definition~\ref{sha-counts-vd-model-form-datum}. Let $A$ and $R$ be complete local noetherian $\mathcal O$-algebras with residue field $\mathbf F$, $p:A\to R$ a surjective homomorphism of $\mathcal O$-algebras with $\mathfrak m_A\ker p=0$, and let $A$ have a weighted presentation (Definition~\ref{sha-counts-vd-model}(1)) whose variables map into $\mathfrak m_A$, with parts $A_j$. Let $S$ be a topological space and $h:S\to M_\iota(R)$ a function with $\Phi$-unitary values, whose entries are continuous for the $\mathfrak m_R$-adic topology, with $h(s)_{ab}\in p(A_{w_a-w_b})$ for all $s$, $a$ and $b$.
Then there is a function $h':S\to M_\iota(A)$ with $\Phi$-unitary values, whose entries are continuous for the $\mathfrak m_A$-adic topology, with $h'(s)_{ab}\in A_{w_a-w_b}$ and $p(h'(s))=h(s)$ for all $s$, $a$ and $b$.
\end{lemma}
\leanref{Development/AutomorphicArithmetic/AutomorphicVarieties/CommutativeAlgebra/HhSection.lean}{UpperTheory.V.Alg.hh_section}
\begin{proof}
For $i\in X$ let $\mathcal C_i$ be the set of the functions $S\to A_i$ that are continuous for the $\mathfrak m_A$-adic topology. Then $\mathcal C_i$ is closed under sums, $\mathcal C_i\mathcal C_{i'}\subseteq\mathcal C_{i+i'}$, and a constant $c\in\mathcal O$ lies in $\mathcal C_i$ if $c=0$ or $i=0$.
\begin{enumerate}
\item[(i)] By Lemma~\ref{sha-counts-vd-model-lemma-p-classes-section-lift}, applied to each entry, there are $y_{ab}\in\mathcal C_{w_a-w_b}$ with $p(y_{ab}(s))=h(s)_{ab}$. Let $h_0(s)$ be the matrix with entries $y_{ab}(s)$ where $\operatorname{bl}(a)=\operatorname{bl}(b)$ and $0$ elsewhere. Then $p(h_0(s))=h(s)$, as $h(s)$ vanishes across blocks, and the entry $(a,b)$ of $h_0$ lies in $\mathcal C_{w_a-w_b}$.
\item[(ii)] Let $h'(s)=h_0(s)+Z(s)$, where $Z(s)$ is obtained from $h_0(s)$ as $Z$ from $h_0$ in step (ii) of the proof of Lemma~\ref{sha-counts-vd-model-lemma-p-unitary}(2): with $E(s)=h_0(s)^{\mathrm t}J_Ah_0(s)-J_A$, $Z(s)$ is the part of $-\tfrac12J_A^{-1}E(s)$ supported on the entries $(a,b)$ with $\operatorname{bl}(a)=\operatorname{bl}(b)$. By that step, $h'(s)$ is $\Phi$-unitary with $p(h'(s))=h(s)$.
\item[(iii)] The entry $(a,b)$ of $h'$ lies in $\mathcal C_{w_a-w_b}$. The constant $(J_A)_{ab}$ lies in $\mathcal C_{c-w_a-w_b}$, as $J_{ab}\ne0$ implies $w_a+w_b=c$. The constant $(J^{-1})_{ab}$ lies in $\mathcal C_{w_a+w_b-c}$: let $D$ and $D'$ be the diagonal matrices of the indicator functions of the sets of the $x$ with $\operatorname{bl}(x)=\operatorname{bl}(a)$ and $w_x=w_a$, and with $\operatorname{bl}(x)=\sigma(\operatorname{bl}(a))$ and $w_a+w_x=c$; the conditions on the nonzero entries of $J$, with $\sigma$ an involution, give $D'J=JD$, so $J^{-1}D'=DJ^{-1}$, and the entry $(a,b)$ of this identity shows that $(J^{-1})_{ab}\ne0$ implies $w_a+w_b=c$. So $s\mapsto E(s)_{ab}=\sum_{d,e}h_0(s)_{da}J_{de}h_0(s)_{eb}-J_{ab}$ lies in $\mathcal C_{c-w_a-w_b}$, and $s\mapsto Z(s)_{ab}$, which is $0$ or $-\tfrac12\sum_e(J^{-1})_{ae}E(s)_{eb}$, lies in $\mathcal C_{w_a-w_b}$. With (i), so does $s\mapsto h'(s)_{ab}$.
\end{enumerate}
\end{proof}

\begin{lemma}\label{sha-counts-vd-model-lemma-p-classes-local-split}
Let $\mathcal O$ be a complete local noetherian ring with residue field $\mathbf F$, let $X$ be a finitely generated free abelian group, and let torus elements be as in Lemma~\ref{sha-counts-vd-model-torus-grading}. Let $A$ and $R$ be complete local noetherian $\mathcal O$-algebras with residue field $\mathbf F$, $p:A\to R$ a homomorphism of $\mathcal O$-algebras with $\mathfrak m_A\ker p=0$, and $\pi$ a weighted presentation of $A$ (Definition~\ref{sha-counts-vd-model}(1)) whose variables map into $\mathfrak m_A$, with parts $A_j$, such that the kernel of $p\circ\pi$ contains, with every series, the sum of its terms of each weight. Let $D\subseteq X$ be a finite set.
Then there are complete local noetherian $\mathcal O$-algebras $S$ and $C$ with residue field $\mathbf F$, a torus element $T$ of $S$, and homomorphisms of $\mathcal O$-algebras $c:A\to S$, $\varphi:C\to S$, $q:C\to A$ and $i:A\to C$, with $\varphi$ injective and $q\circ i$ the identity, such that:
\begin{enumerate}
\item $\varphi(i(a))=T(j)c(a)$ for every $j\in X$ and every $a\in A_j$;
\item for every $d\in D$ and every $u\in\ker p$ there is $z\in C$ with $\varphi(z)=T(d)c(u)$, $q(z)\in A_d$ and $p(q(z))=0$.
\end{enumerate}
\end{lemma}
\leanref{Development/AutomorphicArithmetic/AutomorphicVarieties/CommutativeAlgebra/LocGraded.lean}{split_data}
\leanref{Development/AutomorphicArithmetic/AutomorphicVarieties/CommutativeAlgebra/LocGraded.lean}{torus_data}
\leanref{Development/AutomorphicArithmetic/AutomorphicVarieties/CommutativeAlgebra/LocGraded.lean}{ker_span}
\leanref{Development/AutomorphicArithmetic/AutomorphicVarieties/CommutativeAlgebra/LocGraded.lean}{square_zero_extension_ring}
\leanref{Development/AutomorphicArithmetic/AutomorphicVarieties/CommutativeAlgebra/CoSqz.lean}{UpperTheory.V.Alg.co_sqz}
\begin{proof}
Let $K=\ker p$. It lies in $\mathfrak m_A$, as $p(1)=1\ne0$; so $K^2\subseteq\mathfrak m_AK=0$.
\begin{enumerate}
\item[(i)] \emph{The torus.} Let $n$ be the rank of $X$, let $S=A[[U_1,\dots,U_n]]$, a complete local noetherian $\mathcal O$-algebra with residue field $\mathbf F$, and let $T(j)$ be the image in $S$ of the unit $t(j)$ of Lemma~\ref{sha-counts-vd-model-torus-characters}, for a basis of $X$. Then $T$ is a torus element of $S$, as $t(j)$ has constant term $1$; and by that lemma, if $s\subseteq X$ is finite and elements $a_j\in A$, $j\in s$, satisfy $\sum_{j\in s}T(j)a_j=0$, then every $a_j$ is $0$.
\item[(ii)] \emph{The homomorphism $c$.} Lemma~\ref{sha-counts-vd-model-torus-action}, applied to $\pi$, the inclusion of $A$ in $S$ as the constants, and $j\mapsto T(j)^{-1}$, gives $c:A\to S$ with $c(a)=T(j)^{-1}a$ for $j\in X$ and $a\in A_j$.
\item[(iii)] \emph{Generators of $K$.} By Lemma~\ref{sha-counts-vd-model-lemma-p-relations-kernel}, the kernel of $p\circ\pi$ is generated by finitely many series $f_1,\dots,f_m$, with $f_l$ of a single weight $\lambda_l$. Then $\varepsilon_l=\pi(f_l)$ lies in $A_{\lambda_l}\cap K$, and every $u\in K$ is $\sum_lo_l\varepsilon_l$ with $o_l\in\mathcal O$. Indeed $u=\pi(g)$ with $g=\sum_lg_lf_l$ for some series $g_l$; choose $o_l\in\mathcal O$ with $\pi(g_l)-o_l\in\mathfrak m_A$; then $(\pi(g_l)-o_l)\varepsilon_l\in\mathfrak m_AK=0$.
\item[(iv)] \emph{The ring $C$.} Let $F\subseteq X$ be the finite set of the nonzero differences $d-\lambda_l$ with $d\in D$ and $1\le l\le m$. Let $C=A\oplus K^F$, with the product $(a,k)(a',k')=(aa',ak'+a'k)$, the trivial square-zero extension of $A$ by the finitely generated $A$-module $K^F$. It is a local ring, finite over $A$, with residue field $\mathbf F$, so it is a complete local noetherian $\mathcal O$-algebra with residue field $\mathbf F$. Let $i(a)=(a,0)$, $q(a,k)=a$ and $\varphi(a,k)=a+\sum_{e\in F}T(e)k_e$. As the products of two elements of $K$ vanish, $\varphi$ is a homomorphism of $\mathcal O$-algebras. It is injective: if $\varphi(a,k)=0$, then $T(0)a+\sum_{e\in F}T(e)k_e=0$, where $T(0)=1$ and $0\notin F$, so $a=0$ and every $k_e=0$ by (i). For $e\in F$ and $u\in K$, the element $(0,k)$ with $k_e=u$ and $k_{e'}=0$ for $e'\ne e$ has image $T(e)u$ under $\varphi$ and $0$ under $q$.
\item[(v)] \emph{Property (1).} For $a\in A_j$, $\varphi(i(a))=a=T(j)c(a)$ by (ii).
\item[(vi)] \emph{Property (2).} Let $d\in D$ and $u=\sum_lo_l\varepsilon_l\in K$, as in (iii). Let $v=\sum_{\lambda_l=d}o_l\varepsilon_l$, which lies in $A_d\cap K$. For each $l$ with $\lambda_l\ne d$, $d-\lambda_l\in F$ and $o_l\varepsilon_l\in K$, so (iv) gives $z_l\in C$ with $\varphi(z_l)=T(d-\lambda_l)o_l\varepsilon_l$ and $q(z_l)=0$. Put $z=i(v)+\sum_{\lambda_l\ne d}z_l$. Then $q(z)=v\in A_d$ and $p(v)=0$; and, as $o_l\varepsilon_l\in A_{\lambda_l}$, (ii) gives $\varphi(z)=\sum_lT(d-\lambda_l)o_l\varepsilon_l=T(d)\sum_lT(\lambda_l)^{-1}o_l\varepsilon_l=T(d)c(u)$.
\end{enumerate}
% split_data (LocGraded.lean). The proof follows it: (i) is torus_data (tchar, with co_mvps for S), (ii) is tact applied to T^{-1}, (iii) is ker_span (with ker_gen), (iv) is square_zero_extension_ring (TrivSqZeroExt R (F -> ker p), with co_sqz), (v) and (vi) are the last two items.
\end{proof}

\begin{lemma}\label{sha-counts-vd-model-lemma-p-classes-local}
Let $\mathcal O$, $X$, $\iota$, $w$ and $\Phi$ be as in Lemma~\ref{sha-counts-vd-model-lemma-p-classes-section}. Let $\Gamma$ be a topological group, $\bar r:\Gamma\to M_\iota(\mathbf F)$ a homomorphism and $r_0$ a continuous lift of $\bar r$ with values in $\mathcal O$, both normalising $\Phi$, and $L$ a local condition for $(\bar r,r_0)$ with tangent group a subgroup $N$ of $\Gamma$ (Definition~\ref{sha-counts-vd-model-form-datum}(5)). Let $A$ and $R$ be complete local noetherian $\mathcal O$-algebras with residue field $\mathbf F$, $p:A\to R$ a surjective homomorphism of $\mathcal O$-algebras with $\mathfrak m_A\ker p=0$, and $\pi$ a weighted presentation of $A$ whose variables map into $\mathfrak m_A$, with parts $A_j$, such that $p\circ\pi$ is a weighted presentation of $R$. Let $r'$ be a continuous lift of $\bar r$ with values in $R$ that is $\Phi$-close to $r_0$ and has $L$, with $r'(s)_{ab}\in p(A_{w_a-w_b})$ for all $s$, $a$ and $b$.
Then there is a continuous lift $r$ of $\bar r$ with values in $A$ that is $\Phi$-close to $r_0$ and has $L$, with $p\circ r=r'$ and $r(s)_{ab}\in A_{w_a-w_b}$ for all $s$, $a$ and $b$.
The hypotheses that $2$ is a unit of $\mathcal O$ and that $\bar r$ normalises $\Phi$ are not used in the proof.
\end{lemma}
\leanref{Development/AutomorphicArithmetic/AutomorphicVarieties/CommutativeAlgebra/LocGradedTop.lean}{UpperTheory.V.Alg.loc_graded_top}
\leanref{Development/AutomorphicArithmetic/AutomorphicVarieties/CommutativeAlgebra/LocGraded.lean}{UpperTheory.V.Alg.loc_graded}
\begin{proof}
Let $K=\ker p$.
\begin{enumerate}
\item[(i)] \emph{Lifts.} By property (d) of $L$ there is a continuous lift $r_1$ of $\bar r$ with values in $A$, $\Phi$-close to $r_0$ and having $L$, with $p\circ r_1=r'$. By Lemma~\ref{sha-counts-vd-model-lemma-p-classes-section-lift}, for $a,b\in\iota$ there is a function $y_{ab}:\Gamma\to A_{w_a-w_b}$, continuous for the $\mathfrak m_A$-adic topology, with $p(y_{ab}(g))=r'(g)_{ab}$. Let $d_{ab}(g)=r_1(g)_{ab}-y_{ab}(g)$, which lies in $K$.
\item[(ii)] \emph{The $d_{ab}$ are locally constant.} By the Artin-Rees lemma \cite[Theorem~8.5]{Mat86} there is $k$ with $\mathfrak m_A^{k+1}\cap K=\mathfrak m_A(\mathfrak m_A^k\cap K)$, which lies in $\mathfrak m_AK=0$. The functions $g\mapsto r_1(g)_{ab}$ and $y_{ab}$ are locally constant modulo $\mathfrak m_A^{k+1}$, so for $g'$ near $g$, $d_{ab}(g')-d_{ab}(g)$ lies in $K\cap\mathfrak m_A^{k+1}=0$.
\item[(iii)] \emph{Auxiliary rings.} Lemma~\ref{sha-counts-vd-model-lemma-p-classes-local-split}, applied with $D$ the set of the $w_a-w_b$, gives $S$, $C$, $T$, $c$, $\varphi$, $q$ and $i$. For $a,b\in\iota$ and $u\in K$ choose $Z_{ab}(u)\in C$ with $\varphi(Z_{ab}(u))=T(w_a-w_b)c(u)$, $q(Z_{ab}(u))\in A_{w_a-w_b}$ and $p(q(Z_{ab}(u)))=0$.
\item[(iv)] \emph{A homomorphism over $C$.} Let $M(g)\in M_\iota(C)$ have the entries $M(g)_{ab}=i(y_{ab}(g))+Z_{ab}(d_{ab}(g))$, and let $E(g)\in M_\iota(S)$ have the entries $T(w_a-w_b)c(r_1(g)_{ab})$. Then $E(g)=\Delta\,c(r_1(g))\,\Delta^{-1}$, with $\Delta$ the diagonal matrix of the $T(w_a)$, so $E$ is a homomorphism $\Gamma\to M_\iota(S)$. And $\varphi(M(g))=E(g)$: by Lemma~\ref{sha-counts-vd-model-lemma-p-classes-local-split}(1), $\varphi(i(y_{ab}(g)))=T(w_a-w_b)c(y_{ab}(g))$, as $y_{ab}(g)\in A_{w_a-w_b}$; and $\varphi(Z_{ab}(d_{ab}(g)))=T(w_a-w_b)c(d_{ab}(g))$; their sum is $T(w_a-w_b)c(r_1(g)_{ab})$. As $\varphi$ is injective, $M$ is a homomorphism $\Gamma\to M_\iota(C)$.
\item[(v)] \emph{$M$ is a continuous lift of $\bar r$.} The function $g\mapsto i(y_{ab}(g))$ is continuous for the $\mathfrak m_C$-adic topology, as $i$ maps $\mathfrak m_A$ into $\mathfrak m_C$, and $g\mapsto Z_{ab}(d_{ab}(g))$ is locally constant by (ii). The residue of $M(g)_{ab}$ is that of $\varphi(M(g)_{ab})=T(w_a-w_b)c(r_1(g)_{ab})$, which is $\bar r(g)_{ab}$, as $T(w_a-w_b)\equiv1$ and $c$ is compatible with the maps to $\mathbf F$.
\item[(vi)] \emph{The image of $r_1$ in $S$.} The lift $c\circ r_1$ is $\Phi$-close to $r_0$ by Lemma~\ref{sha-counts-vd-model-lemma-p-unitary}(1), as $c(r_1(s))r_0(s^{-1})_S$ is the image of $r_1(s)r_0(s^{-1})_A$. It has $L$: property (a) of $L$ gives an ideal $I_0$ of $A$ such that, for every $f:A\to R'$, $f\circ r_1$ has $L$ if and only if $f(I_0)=0$; as $r_1$ has $L$, the identity of $A$ kills $I_0$, so $I_0=0$, and $c\circ r_1$ has $L$.
\item[(vii)] \emph{$\varphi\circ M$ has $L$.} By (iv), its entries are those of $c\circ r_1$ multiplied by $T(w_a-w_b)$, and $T$ is a torus element; so property (c) of $L$ applies.
\item[(viii)] \emph{$M$ is $\Phi$-close to $r_0$.} As $r_0$ normalises $\Phi$, $r_0(s^{-1})_{eb}=0$ unless $w_e=w_b$. So the entry $(a,b)$ of $\varphi(M(s))r_0(s^{-1})_S$, which is $\sum_eT(w_a-w_e)c(r_1(s))_{ae}r_0(s^{-1})_{eb}$, is $T(w_a-w_b)$ times that of the $\Phi$-unitary matrix $c(r_1(s))r_0(s^{-1})_S$, and it is $\Phi$-unitary by Lemma~\ref{sha-counts-vd-model-lemma-p-unitary}(1). It is the image under $\varphi$ of $h=M(s)r_0(s^{-1})_C$. As $\varphi$ is injective and carries $J_C$ to $J_S$, $h^{\mathrm t}J_Ch=J_C$ and $h_{ab}=0$ whenever $\operatorname{bl}(a)\ne\operatorname{bl}(b)$; and $h\equiv1$, its residue being that of $\varphi(h)$. So $h$ is $\Phi$-unitary.
\item[(ix)] \emph{The lift $r$.} Let $r=q\circ M$, a continuous lift of $\bar r$ with values in $A$, which is $\Phi$-close to $r_0$ by (viii) and Lemma~\ref{sha-counts-vd-model-lemma-p-unitary}(1). It has $L$: property (a) of $L$, for $M$, gives an ideal $I$ of $C$ such that $f\circ M$ has $L$ if and only if $f(I)=0$; by (vii), $\varphi(I)=0$, so $I=0$, as $\varphi$ is injective, and $q\circ M$ has $L$. As $q\circ i$ is the identity, $r(g)_{ab}=y_{ab}(g)+q(Z_{ab}(d_{ab}(g)))$, which lies in $A_{w_a-w_b}$, and $p(r(g)_{ab})=p(y_{ab}(g))=r'(g)_{ab}$.
\end{enumerate}
% loc_graded (LocGraded.lean, 468 lines). The proof follows it: (i) is hL.lift and cont_section (y); (ii) is the private lc_of_cont, with inf_pow_bot (Artin-Rees, Ideal.exists_pow_inf_eq_pow_smul); (iii) is split_data and the choice of Z; (iv) is M, E, hME, hE1, hEm and rAh; (v) is hcont (cont_map, hZlc) and hred; (vi) is hInHc (hh_group) and hLc; (vii) is hLA, the torus field of IsLocCond; (viii) is hInHA, with inH_map, the torus twist of hh_group and hh_of_inj; (ix) is loc_map and the last two items. loc_graded_top reduces to loc_graded by replacing the topology of Gamma with one that makes Gamma a topological group and has the same locally constant homomorphisms (the private exists_grpTop and tr); in the book Gamma is a topological group, so this step is not needed.
\end{proof}

\begin{lemma}\label{sha-counts-vd-model-lemma-p-classes-coefficients}
Let $\mathcal O$ be a local ring with residue field $\mathbf F$, let $X$ be an abelian group, and let $A$ be a complete local noetherian $\mathcal O$-algebra with residue field $\mathbf F$ with a weighted presentation $\pi$ (Definition~\ref{sha-counts-vd-model}(1)), with parts $A_j$. Let $K$ be an ideal of $A$ with $\mathfrak m_AK=0$, generated by elements $\varepsilon_u\in A_{w_f(u)}$, $u\in\kappa$, $\kappa$ finite, that are independent: if $c_u\in\mathcal O$ and $\sum_uc_u\varepsilon_u=0$, then every $c_u\in\mathfrak m_{\mathcal O}$.
Then there are functions $\mathrm{co}_u:A\to\mathbf F$, $u\in\kappa$, such that for $x,y\in K$ and $a\in A$: $\mathrm{co}_u(x+y)=\mathrm{co}_u(x)+\mathrm{co}_u(y)$; $\mathrm{co}_u(ax)=\bar a\,\mathrm{co}_u(x)$, with $\bar a$ the residue of $a$; $\mathrm{co}_u(x)=0$ if $x\in A_j$ for some $j\ne w_f(u)$; and $x=\sum_ua_u\varepsilon_u$ for some $a_u\in\mathcal O$ with residues $\mathrm{co}_u(x)$.
\end{lemma}
\leanref{Development/AutomorphicArithmetic/AutomorphicVarieties/CommutativeAlgebra/KerCoef.lean}{UpperTheory.V.Alg.ker_coef}
\begin{proof}
Every $x\in K$ is $\sum_ua_u\varepsilon_u$ with $a_u\in\mathcal O$: it is $\sum_ub_u\varepsilon_u$ with $b_u\in A$, and replacing $b_u$ by an element $a_u$ of $\mathcal O$ with $b_u-a_u\in\mathfrak m_A$ does not change the sum, as $\mathfrak m_A$ kills $\varepsilon_u$. By the independence, the residues of the $a_u$ depend only on $x$; let $\mathrm{co}_u(x)$ be the residue of $a_u$, and $\mathrm{co}_u(x)=0$ for $x\notin K$. Additivity is clear. For $a\in A$ choose $c\in\mathcal O$ with $a-c\in\mathfrak m_A$; then $ax=cx=\sum_u(ca_u)\varepsilon_u$, so $\mathrm{co}_u(ax)=\bar c\,\mathrm{co}_u(x)=\bar a\,\mathrm{co}_u(x)$.
Let $x\in K\cap A_j$ with $j\ne w_f(u)$, and write $x=\sum_va_v\varepsilon_v$ as above. Choose a power series $g$ of weight $j$ with $\pi(g)=x$ and series $f_v$ of weight $w_f(v)$ with $\pi(f_v)=\varepsilon_v$. Then $g-\sum_va_vf_v\in\ker\pi$, and so is its component of weight $w_f(u)$, which is $-\sum_{w_f(v)=w_f(u)}a_vf_v$, as $j\ne w_f(u)$. So $\sum_{w_f(v)=w_f(u)}a_v\varepsilon_v=0$, and by the independence $a_u\in\mathfrak m_{\mathcal O}$, that is $\mathrm{co}_u(x)=0$.
\end{proof}

\begin{lemma}\label{sha-counts-vd-model-lemma-p-classes-cocycle}
Let $\mathcal O$ be a complete local noetherian ring with residue field $\mathbf F$ in which $2$ is a unit, let $\iota$, $X$, $w$ and a form datum $\Phi$ be as in Definition~\ref{sha-counts-vd-model-form-datum}, let $\Gamma$ be a topological group, $\bar r:\Gamma\to M_\iota(\mathbf F)$ a homomorphism and $r_0$ a continuous lift of $\bar r$ with values in $\mathcal O$, both normalising $\Phi$, and let $F$ be a place datum for $\Gamma$ (Definition~\ref{sha-counts-vd-model-place-datum}); $\Gamma$, the $G_v$ and the $G_q$ act on $\mathfrak g_{\Phi,j}(\mathbf F)$ through $\bar r$. Let $A$ and $R$ be complete local noetherian $\mathcal O$-algebras with residue field $\mathbf F$, $p:A\to R$ a surjective homomorphism of $\mathcal O$-algebras with $\mathfrak m_A\ker p=0$, and let $A$ have a weighted presentation with parts $A_j$. Let $j\in X$, and let $\mathrm{co}:A\to\mathbf F$ satisfy, for $x,y\in\ker p$ and $a\in A$: $\mathrm{co}(x+y)=\mathrm{co}(x)+\mathrm{co}(y)$, $\mathrm{co}(ax)=\bar a\,\mathrm{co}(x)$, and $\mathrm{co}(x)=0$ if $x\in A_{j'}$ for some $j'\ne j$. Suppose given:
\begin{enumerate}
\item[(a)] $\tilde r:\Gamma\to M_\iota(A)$ with continuous entries in the parts $A_{w_a-w_b}$, such that $\tilde r(s)r_0(s^{-1})_A$ is $\Phi$-unitary and $p(\tilde r(s)\tilde r(t))=p(\tilde r(st))$ for all $s,t$;
\item[(b)] for $v\in T$, a $\Phi$-unitary $\tilde\alpha_v\in M_\iota(A)$ and a continuous lift $\rho^A_v$ of $\bar r\circ\varphi_v$ with values in $A$, $\Phi$-close to $r_0\circ\varphi_v$, both with entries in the parts $A_{w_a-w_b}$, such that $p(\tilde r(\varphi_v(s))\tilde\alpha_v)=p(\tilde\alpha_v\rho^A_v(s))$ for all $s\in G_v$;
\item[(c)] for $q\in Q$, a continuous lift $\rho^A_q$ of $\bar r\circ\psi_q$ with values in $A$, $\Phi$-close to $r_0\circ\psi_q$, with entries in the parts $A_{w_a-w_b}$, such that $p(\tilde r(\psi_q(s)))=p(\rho^A_q(s))$ for all $s\in G_q$.
\end{enumerate}
Then there is $(z,(z_v),(z_q))\in C^2_F(\mathfrak g_{\Phi,j}(\mathbf F))$ with $z(s,t)\bar r(st)=\mathrm{co}\big(\tilde r(s)\tilde r(t)-\tilde r(st)\big)$, $z_v(s)\bar r(\varphi_v(s))=\mathrm{co}\big(\tilde r(\varphi_v(s))\tilde\alpha_v-\tilde\alpha_v\rho^A_v(s)\big)$ and $z_q(s)\bar r(\psi_q(s))=\mathrm{co}\big(\tilde r(\psi_q(s))-\rho^A_q(s)\big)$, with $\mathrm{co}$ applied to the entries.
\end{lemma}
\leanref{Development/AutomorphicArithmetic/AutomorphicVarieties/CommutativeAlgebra/ObsCocycle.lean}{UpperTheory.V.Alg.obs_cocycle}
\begin{proof}
Let $K=\ker p$. For a matrix $Y$ with entries in $K$ write $\mathrm{co}(Y)$ for $\mathrm{co}$ applied to the entries. Then $\mathrm{co}(Y+Y')=\mathrm{co}(Y)+\mathrm{co}(Y')$, $\mathrm{co}(MY)=\bar M\mathrm{co}(Y)$ and $\mathrm{co}(YM)=\mathrm{co}(Y)\bar M$ for $M\in M_\iota(A)$ with residue $\bar M$. The residue of $\tilde r(s)$ is $\bar r(s)$, as that of $\tilde r(s)r_0(s^{-1})_A$ is $1$; the residues of $\rho^A_v(s)$, $\rho^A_q(s)$ and $\tilde\alpha_v$ are $\bar r(\varphi_v(s))$, $\bar r(\psi_q(s))$ and $1$.
\begin{enumerate}
\item[(i)] \emph{Values.} Let $g\in\Gamma$, and let $Y,Y'\in M_\iota(A)$ have entries in the parts $A_{w_a-w_b}$, with $Yr_0(g^{-1})_A$ and $Y'r_0(g^{-1})_A$ $\Phi$-unitary and $p(Y)=p(Y')$. Then $\mathrm{co}(Y-Y')\bar r(g^{-1})\in\mathfrak g_{\Phi,j}(\mathbf F)$. Indeed it is $\mathrm{co}(D)$ with $D=Yr_0(g^{-1})_A-Y'r_0(g^{-1})_A$. By Lemma~\ref{sha-counts-vd-model-lemma-p-unitary}(2)(b), $D\in\mathfrak g_\Phi(A)$; its entries lie in $K$ and in the parts $A_{w_a-w_b}$, as $r_0(g^{-1})_{ab}$ vanishes unless $w_a=w_b$. Applying $\mathrm{co}$ to $D^{\mathrm t}J_A+J_AD=0$ gives $\mathrm{co}(D)^{\mathrm t}J_{\mathbf F}+J_{\mathbf F}\mathrm{co}(D)=0$; $\mathrm{co}(D)$ vanishes across blocks; and its entry $(a,b)$ vanishes if $w_a-w_b\ne j$, by the last property of $\mathrm{co}$.
\item[(ii)] \emph{The functions.} Put
\[
 z(s,t)=\mathrm{co}\big(\tilde r(s)\tilde r(t)-\tilde r(st)\big)\bar r((st)^{-1}),\quad z_v(s)=\mathrm{co}\big(\tilde r(\varphi_v(s))\tilde\alpha_v-\tilde\alpha_v\rho^A_v(s)\big)\bar r(\varphi_v(s)^{-1}),
\]
and $z_q(s)=\mathrm{co}\big(\tilde r(\psi_q(s))-\rho^A_q(s)\big)\bar r(\psi_q(s)^{-1})$. They take values in $\mathfrak g_{\Phi,j}(\mathbf F)$ by (i): for $z$, with $g=st$, $Y=\tilde r(s)\tilde r(t)$ and $Y'=\tilde r(st)$, where $Yr_0((st)^{-1})_A$ is the product of $\tilde r(s)r_0(s^{-1})_A$ and of the conjugate by $r_0(s)_A$ of $\tilde r(t)r_0(t^{-1})_A$, $\Phi$-unitary by Lemma~\ref{sha-counts-vd-model-lemma-p-unitary}(1); for $z_v$, with $g=\varphi_v(s)$, $Y=\tilde r(g)\tilde\alpha_v$ and $Y'=\tilde\alpha_v\rho^A_v(s)$, similarly; for $z_q$, with $Y=\tilde r(\psi_q(s))$ and $Y'=\rho^A_q(s)$. The identities of the statement hold, as $\bar r(g^{-1})\bar r(g)=1$.
\item[(iii)] \emph{Local constancy.} By the theorem of Artin and Rees \cite[Theorem~8.5]{Mat86} there is $n$ with $\mathfrak m_A^{n+1}\cap K=\mathfrak m_A(\mathfrak m_A^n\cap K)$, which is $0$. The matrices $\tilde r(s)\tilde r(t)-\tilde r(st)$ have entries in $K$, and modulo $\mathfrak m_A^{n+1}$ they are locally constant in $(s,t)$, as the entries of $\tilde r$ are continuous; so they are locally constant. The same holds for the matrices defining $z_v$ and $z_q$, and $\bar r$ is locally constant, being the residue of $r_0$. So $z$, $z_v$ and $z_q$ are locally constant.
\item[(iv)] \emph{The cocycle identity.} The identity
\[
 \tilde r(s)\big(\tilde r(t)\tilde r(u)-\tilde r(tu)\big)+\big(\tilde r(s)\tilde r(tu)-\tilde r(stu)\big)=\big(\tilde r(st)\tilde r(u)-\tilde r(stu)\big)+\big(\tilde r(s)\tilde r(t)-\tilde r(st)\big)\tilde r(u),
\]
to which $\mathrm{co}$ applies term by term, gives, after multiplication by $\bar r((stu)^{-1})$ on the right, $\bar r(s)z(t,u)\bar r(s)^{-1}+z(s,tu)=z(st,u)+z(s,t)$.
\item[(v)] \emph{The places.} Let $\varphi=\varphi_v$, $\beta=\tilde\alpha_v$ and $\tau=\rho^A_v$, or $\varphi=\psi_q$, $\beta=1$ and $\tau=\rho^A_q$. The identity
\[
 \big(\tilde r(\varphi(s))\beta-\beta\tau(s)\big)\tau(t)+\tilde r(\varphi(s))\big(\tilde r(\varphi(t))\beta-\beta\tau(t)\big)=\big(\tilde r(\varphi(s))\tilde r(\varphi(t))-\tilde r(\varphi(st))\big)\beta+\big(\tilde r(\varphi(st))\beta-\beta\tau(st)\big),
\]
with $\mathrm{co}$ applied term by term and $\bar\beta=1$, gives $z(\varphi(s),\varphi(t))=y(s)+\varphi(s)\cdot y(t)-y(st)$ for $y=z_v$, respectively $y=z_q$.
\end{enumerate}
By (ii) to (v), $(z,(z_v),(z_q))\in C^2_F(\mathfrak g_{\Phi,j}(\mathbf F))$.
\end{proof}

\begin{lemma}\label{sha-counts-vd-model-lemma-p-classes-correction}
Let $\mathcal O$, $\iota$, $X$, $w$, $\Phi$, $\Gamma$, $\bar r$, $r_0$ and $F$ be as in Lemma~\ref{sha-counts-vd-model-lemma-p-classes-cocycle}, with $\bar r$ trivial on $\psi_q(N_q)$ for $q\in Q$, and for $q\in Q$ let $L_q$ be a local condition for $(\bar r\circ\psi_q,r_0\circ\psi_q)$ with tangent group $N_q$. Let $B$ be a complete local noetherian $\mathcal O$-algebra with residue field $\mathbf F$ and $\rho_v$ a continuous lift of $\bar r\circ\varphi_v$ with values in $B$ for $v\in T$; framings are over $B$ for these data (Definition~\ref{sha-counts-vd-model-form-datum}(6)). Let $A$ and $R$ be complete local noetherian $\mathcal O$-algebras with residue field $\mathbf F$, $p:A\to R$ a surjective homomorphism of $\mathcal O$-algebras with $\mathfrak m_A\ker p=0$, and $\mu:\mathbf F\to A$ an additive map with $\mu(\bar ac)=a\mu(c)$ and $p(\mu(c))=0$ for $a\in A$ and $c\in\mathbf F$. Let $i:B\to A$ be a homomorphism of $\mathcal O$-algebras, $r$ a continuous lift of $\bar r$ with values in $R$, and:
\begin{enumerate}
\item[(a)] $\tilde r:\Gamma\to M_\iota(A)$ with continuous entries, $\tilde r(s)r_0(s^{-1})_A$ $\Phi$-unitary and $p(\tilde r(s))=r(s)$ for every $s$;
\item[(b)] $\Phi$-unitary $\tilde\alpha_v\in M_\iota(A)$ for $v\in T$, with $\tilde\alpha_{v_0}=1$;
\item[(c)] for $q\in Q$, a continuous lift $\rho^A_q$ of $\bar r\circ\psi_q$ with values in $A$, $\Phi$-close to $r_0\circ\psi_q$, with $L_q$, and with $p\circ\rho^A_q=r\circ\psi_q$.
\end{enumerate}
Let $j\in X$ and $(z,(z_v),(z_q))\in B^2_F(\mathfrak g_{\Phi,j}(\mathbf F))$ with $\tilde r(s)\tilde r(t)-\tilde r(st)=\mu\big(z(s,t)\bar r(st)\big)$, $\tilde r(\varphi_v(s))\tilde\alpha_v-\tilde\alpha_vi(\rho_v(s))=\mu\big(z_v(s)\bar r(\varphi_v(s))\big)$ and $\tilde r(\psi_q(s))-\rho^A_q(s)=\mu\big(z_q(s)\bar r(\psi_q(s))\big)$, with $\mu$ applied to the entries.
Then there are a continuous lift $r'$ of $\bar r$ with values in $A$ and matrices $\alpha'_v$ such that $(i,r',\alpha')$ is a framing, $p\circ r'=r$ and $p(\alpha'_v)=p(\tilde\alpha_v)$ for every $v$.
\end{lemma}
\leanref{Development/AutomorphicArithmetic/AutomorphicVarieties/CommutativeAlgebra/ObsCorrect.lean}{UpperTheory.V.Alg.obs_correct}
\leanref{Development/AutomorphicArithmetic/AutomorphicVarieties/CommutativeAlgebra/CocycleNormalForm.lean}{UpperTheory.V.Alg.cocycle_normal_form}
\begin{proof}
For $Y\in M_\iota(\mathbf F)$ let $T(Y)\in M_\iota(A)$ be $\mu$ applied to the entries. Then $T$ is additive, and $MT(Y)=T(\bar MY)$ and $T(Y)M=T(Y\bar M)$ for $M\in M_\iota(A)$ with residue $\bar M$. As $p(\mu(c))=0$, $p(T(Y))=0$ and the entries of $T(Y)$ lie in $\mathfrak m_A$; so $T(Y)T(Y')=T(0\cdot Y')=0$. The residues of $\tilde r(s)$, $r_0(s)_A$ and $i(\rho_v(s))$ are $\bar r(s)$, $\bar r(s)$ and $\bar r(\varphi_v(s))$, and that of $\tilde\alpha_v$ is $1$. For $Y\in\mathfrak g_{\Phi,j}(\mathbf F)$, $T(Y)\in\mathfrak g_\Phi(A)$: $T(Y)^{\mathrm t}=T(Y^{\mathrm t})$, so $T(Y)^{\mathrm t}J_A+J_AT(Y)=T(Y^{\mathrm t}J_{\mathbf F}+J_{\mathbf F}Y)=0$, and $T(Y)$ vanishes where $Y$ does.
\begin{enumerate}
\item[(i)] \emph{Normal form.} By Definition~\ref{sha-counts-vd-model-place-datum}(2) and the proof of Lemma~\ref{sha-counts-vd-model-obstruction}, $(z,(z_v),(z_q))=\big(dc,(c\circ\varphi_v-dy'_v)_v,(c\circ\psi_q-e_q)_q\big)$ for some $c$, $y'_v$ and $e_q$, where $dy'_v(s)=\varphi_v(s)\cdot y'_v-y'_v$. Put $h(g)=c(g)-(g\cdot y'_{v_0}-y'_{v_0})$ and $y_v=y'_v-y'_{v_0}$. Then $h$ is locally constant, $dh=dc$, $y_{v_0}=0$, $z_v=h\circ\varphi_v-dy_v$, and $z_q(n)=h(\psi_q(n))$ for $n\in N_q$, since $e_q$ vanishes on $N_q$ and $\psi_q(N_q)$ acts trivially.
\item[(ii)] \emph{The lift.} Put $r'(s)=\tilde r(s)-T(h(s)\bar r(s))$. Then $r'(s)r_0(s^{-1})_A=\tilde r(s)r_0(s^{-1})_A-T(h(s))$ is $\Phi$-unitary, by Lemma~\ref{sha-counts-vd-model-lemma-p-unitary}(2)(c), and $p(r'(s))=r(s)$. Since $dh(s,t)\bar r(st)=h(s)\bar r(s)\bar r(t)+\bar r(s)h(t)\bar r(t)-h(st)\bar r(st)$ and $\tilde r(s)\tilde r(t)-\tilde r(st)=T(dh(s,t)\bar r(st))$,
\[
 r'(s)r'(t)=\tilde r(s)\tilde r(t)-T\big(\bar r(s)h(t)\bar r(t)\big)-T\big(h(s)\bar r(s)\bar r(t)\big)=\tilde r(st)-T\big(h(st)\bar r(st)\big)=r'(st).
\]
As $r'(1)$ is invertible and $r'(1)^2=r'(1)$, $r'(1)=1$, and $r'$ is a homomorphism. Its entries are continuous, as $h$ and $\bar r$ are locally constant, and its residue is $\bar r$. So $r'$ is a continuous lift of $\bar r$, $\Phi$-close to $r_0$, with $p\circ r'=r$.
\item[(iii)] \emph{The frames.} Put $\alpha'_v=\tilde\alpha_v+T(y_v)$. It is $\Phi$-unitary by Lemma~\ref{sha-counts-vd-model-lemma-p-unitary}(2)(c), $\alpha'_{v_0}=\tilde\alpha_{v_0}=1$, and $p(\alpha'_v)=p(\tilde\alpha_v)$. For $s\in G_v$ write $g=\varphi_v(s)$. Then
\[
 r'(g)\alpha'_v-\alpha'_vi(\rho_v(s))=\tilde r(g)\tilde\alpha_v-\tilde\alpha_vi(\rho_v(s))+T\big(\bar r(g)y_v-y_v\bar r(g)-h(g)\bar r(g)\big),
\]
and $\tilde r(g)\tilde\alpha_v-\tilde\alpha_vi(\rho_v(s))=T(z_v(s)\bar r(g))=T\big(h(g)\bar r(g)-\bar r(g)y_v+y_v\bar r(g)\big)$ by (i). So $r'(\varphi_v(s))\alpha'_v=\alpha'_vi(\rho_v(s))$.
\item[(iv)] \emph{The local conditions.} Let $q\in Q$. The lifts $\rho^A_q$ and $r'\circ\psi_q$ are $\Phi$-close to $r_0\circ\psi_q$ and have the same image $r\circ\psi_q$ under $p$, and $\rho^A_q$ has $L_q$. For $n\in N_q$, $\tilde r(\psi_q(n))-\rho^A_q(n)=T\big(h(\psi_q(n))\bar r(\psi_q(n))\big)$ by (i), so $r'(\psi_q(n))=\rho^A_q(n)$. By the tangent property (e) of $L_q$, $r'\circ\psi_q$ has $L_q$.
\end{enumerate}
So $(i,r',\alpha')$ is a framing with the stated properties.
\end{proof}

\begin{lemma}\label{sha-counts-vd-model-lemma-p-classes}
In Situation~\ref{sha-counts-vd-model-lemma-p-situation}, suppose that $X$ is finitely generated and free. Let $p:A\to R$ be a surjective homomorphism of $\mathcal O$-algebras with $\mathfrak m_A\ker p=0$, and let $A$ have a weighted presentation whose variables map into $\mathfrak m_A$. Let $\varepsilon_u\in A$, for $u$ in a finite set $\kappa$, lie in the parts $A_{w_f(u)}$, generate $\ker p$, and be independent: if $c_u\in\mathcal O$ and $\sum_uc_u\varepsilon_u=0$, then every $c_u\in\mathfrak m_{\mathcal O}$. Let $i:B\to A$ be a homomorphism of $\mathcal O$-algebras and $(p\circ i,r,\alpha)$ a framing with values in $R$, such that the entries $(a,b)$ of the $i(\rho_v(s))$ lie in $A_{w_a-w_b}$, and those of the $r(g)$ and the $\alpha_v$ lie in $p(A_{w_a-w_b})$.
Then for every $j\in X$ there are elements $o_j(u)\in\mathrm{Ob}_F(\mathfrak g_{\Phi,j})$, for the $u$ with $w_f(u)=j$, with the following property. Let $j\in X$, and let $c_u\in\mathcal O$ for the $u$ with $w_f(u)=j$, with $\sum_u\bar c_u\,o_j(u)=0$. Let $p':A\to A'$ be a surjective homomorphism of $\mathcal O$-algebras, $p'':A'\to R$ a homomorphism with $p''\circ p'=p$, and $\eta\in A'$ with $\mathfrak m_{A'}\eta=0$ and $p''(\eta)=0$, such that $p'(\varepsilon_u)=c_u\eta$ if $w_f(u)=j$ and $p'(\varepsilon_u)=0$ otherwise. Then there are a continuous lift $r'$ of $\bar r$ with values in $A'$ and matrices $\alpha'_v\in M_\iota(A')$ such that $(p'\circ i,r',\alpha')$ is a framing, $p''\circ r'=r$, and $p''(\alpha'_v)=\alpha_v$ for every $v$.
The hypothesis of the situation that $r_0\circ\psi_q$ has $L_q$ is not used in the proof.
\end{lemma}
\leanref{Development/AutomorphicArithmetic/AutomorphicVarieties/CommutativeAlgebra/ObsClass.lean}{UpperTheory.V.Alg.obs_class}
\begin{proof}
Write $A_j$ for the parts of $A$; then $p(A_j)$ are the parts of the presentation of $R$ obtained by composing with $p$, which is a weighted presentation by Lemma~\ref{sha-counts-vd-presentations}(a), as $\ker p$ is generated by the $\varepsilon_u$. The entries of $r_0(g)$ vanish unless $w_a=w_b$, so those of $r_0(g)_A$ lie in the parts $A_{w_a-w_b}$.
\begin{enumerate}
\item[(i)] \emph{Lifts of $r$ and $\alpha$.} The function $s\mapsto r(s)r_0(s^{-1})_R$ has $\Phi$-unitary values, continuous entries, and entries in the $p(A_{w_a-w_b})$. Lemma~\ref{sha-counts-vd-model-lemma-p-classes-section} lifts it to $h$ with the same properties over $A$; put $\tilde r(s)=h(s)r_0(s)_A$. Then $\tilde r(s)r_0(s^{-1})_A=h(s)$ is $\Phi$-unitary, the entries of $\tilde r$ are continuous and lie in the parts $A_{w_a-w_b}$, and $p(\tilde r(s))=r(s)$. Lemma~\ref{sha-counts-vd-model-lemma-p-classes-section}, applied to the finite set $T$ with the discrete topology, gives $\Phi$-unitary lifts $\beta_v$ of the $\alpha_v$ with entries in the parts $A_{w_a-w_b}$; put $\tilde\alpha_v=\beta_v$ for $v\ne v_0$ and $\tilde\alpha_{v_0}=1$.
\item[(ii)] \emph{Lifts at $Q$.} For $q\in Q$, $r\circ\psi_q$ is $\Phi$-close to $r_0\circ\psi_q$ and has $L_q$, and its entries lie in the $p(A_{w_a-w_b})$. So Lemma~\ref{sha-counts-vd-model-lemma-p-classes-local} gives a continuous lift $\rho^A_q$ of $\bar r\circ\psi_q$ with values in $A$, $\Phi$-close to $r_0\circ\psi_q$, with $L_q$, $p\circ\rho^A_q=r\circ\psi_q$, and entries in the parts $A_{w_a-w_b}$.
\item[(iii)] \emph{The classes.} Let $\mathrm{co}_u$ be as in Lemma~\ref{sha-counts-vd-model-lemma-p-classes-coefficients}, for $K=\ker p$. Since $(p\circ i,r,\alpha)$ is a framing, $p(\tilde r(s)\tilde r(t))=p(\tilde r(st))$, $p(\tilde r(\varphi_v(s))\tilde\alpha_v)=p(\tilde\alpha_vi(\rho_v(s)))$ and $p(\tilde r(\psi_q(s)))=p(\rho^A_q(s))$. The lifts $i\circ\rho_v$ are $\Phi$-close to $r_0\circ\varphi_v$ by Lemma~\ref{sha-counts-vd-model-lemma-p-unitary}(1). So for each $u$, Lemma~\ref{sha-counts-vd-model-lemma-p-classes-cocycle}, with $j=w_f(u)$, $\mathrm{co}=\mathrm{co}_u$ and $\rho^A_v=i\circ\rho_v$, gives $(z^u,(z^u_v),(z^u_q))\in C^2_F(\mathfrak g_{\Phi,j})$. Let $o_j(u)$ be its class in $\mathrm{Ob}_F(\mathfrak g_{\Phi,j})$.
\item[(iv)] Let $j$, the $c_u$, $p'$, $p''$ and $\eta$ be as in the statement, and let $x=\sum_{w_f(u)=j}\bar c_u(z^u,(z^u_v),(z^u_q))$, which lies in $B^2_F(\mathfrak g_{\Phi,j})$. Let $\mu:\mathbf F\to A'$ map the residue of $a\in\mathcal O$ to $a\eta$; it is well defined, as $\mathfrak m_{A'}\eta=0$, additive, $\mu(\bar bc)=b\mu(c)$ for $b\in A'$, and $p''(\mu(c))=0$. For $y\in\ker p$, write $y=\sum_ua_u\varepsilon_u$ with $a_u\in\mathcal O$ of residues $\mathrm{co}_u(y)$; then $p'(y)=\sum_{w_f(u)=j}a_uc_u\eta=\mu\big(\sum_{w_f(u)=j}\bar c_u\mathrm{co}_u(y)\big)$. Hence for matrices $Y$, $Y'$ over $A$ with $p(Y)=p(Y')$, $p'(Y)-p'(Y')=\mu\big(\sum_{w_f(u)=j}\bar c_u\mathrm{co}_u(Y-Y')\big)$. Applied to the three differences of (iii), with the identities of Lemma~\ref{sha-counts-vd-model-lemma-p-classes-cocycle}, this gives $p'(\tilde r(s))p'(\tilde r(t))-p'(\tilde r(st))=\mu(x_1(s,t)\bar r(st))$ and the two similar identities, where $x=(x_1,(x_v),(x_q))$.
\item[(v)] \emph{The lift.} The homomorphism $p''$ is surjective, and $\mathfrak m_{A'}\ker p''=0$: an element of $\ker p''$ is $p'(y)$ with $y\in\ker p$, which is $\mu(c)$ for some $c$ by (iv), and $\mathfrak m_{A'}$ kills $\mu(c)$. The lifts $p'\circ\rho^A_q$ have $L_q$: by property (a) of $L_q$ for $\rho^A_q$ there is an ideal $I$ of $A$ such that $f\circ\rho^A_q$ has $L_q$ if and only if $f(I)=0$; with $f$ the identity, $I=0$. By Lemma~\ref{sha-counts-vd-model-lemma-p-unitary}(1) the images under $p'$ of $\tilde r(s)r_0(s^{-1})_A$, of the $\tilde\alpha_v$ and of the $\rho^A_q(s)r_0(\psi_q(s)^{-1})_A$ are $\Phi$-unitary. So Lemma~\ref{sha-counts-vd-model-lemma-p-classes-correction}, applied to $p'':A'\to R$, $\mu$, $p'\circ i$, $r$, $p'\circ\tilde r$, the $p'(\tilde\alpha_v)$, the $p'\circ\rho^A_q$ and $x$, gives a continuous lift $r'$ of $\bar r$ with values in $A'$ and matrices $\alpha'_v$ such that $(p'\circ i,r',\alpha')$ is a framing, $p''\circ r'=r$, and $p''(\alpha'_v)=p''(p'(\tilde\alpha_v))=p(\tilde\alpha_v)=\alpha_v$.
\end{enumerate}
\end{proof}

\begin{lemma}\label{sha-counts-vd-model-lemma-p}
Let $\ell$ be an odd prime, and let $\mathcal O$ be a complete local noetherian ring whose residue field $\mathbf F$ is finite of characteristic $\ell$; then $2$ is a unit of $\mathcal O$. In Situation~\ref{sha-counts-vd-model-lemma-p-situation} for this $\mathcal O$, suppose that $X$ is finitely generated and free, that $\Gamma$ is compact, totally disconnected and Hausdorff and satisfies the finiteness condition of Mazur of Lemma~\ref{sha-counts-vd-model-lemma-p-universal}, and that $\mathrm{Ob}_F(\mathfrak g_{\Phi,j})$ is finite-dimensional for every $j\in X$. Let $\pi_B$ be a weighted presentation of $B$ (Definition~\ref{sha-counts-vd-model}(1)) whose variables map into $\mathfrak m_B$, such that the entry $(a,b)$ of every $\rho_v(s)$ lies in the part $B_{w_a-w_b}$.
Then there are an $R$, a universal framing $(i,r,\alpha)$ with values in $R$, and a weighted presentation $\pi$ of $R$ by the variables of $\pi_B$ and variables $Y_t$, $t\in\tau$, relative to $\pi_B$ and $i$ with relations $f_u$ of weights $w_f(u)$, $u\in\kappa$ (Definition~\ref{sha-counts-vd-model-form-datum}(7)), with the following properties.
\begin{enumerate}
\item The entry $(a,b)$ of every $r(g)$ and of every $\alpha_v$ lies in the part $R_{w_a-w_b}$.
\item For each $t$ there are $a,b\in\iota$ such that $Y_t$ has the weight $w_a-w_b$ and $\pi(Y_t)$ is the entry $(a,b)$ of $r(g)-r_0(g)_R$ for some $g\in\Gamma$ or of $\alpha_v-1$ for some $v\in T$.
\item For every $j\in X$, the number of $t$ with $Y_t$ of weight $j$ is $\dim_{\mathbf F}Z^1_F(\mathfrak g_{\Phi,j})$, and the number of $u$ with $w_f(u)=j$ is $\dim_{\mathbf F}\mathrm{Ob}_F(\mathfrak g_{\Phi,j})$.
\end{enumerate}
\end{lemma}
\leanref{Development/AutomorphicArithmetic/AutomorphicVarieties/CommutativeAlgebra/FramedTwoCoboundaries.lean}{UpperTheory.V.Alg.LemmaP}
\leanref{Development/AutomorphicArithmetic/AutomorphicVarieties/CommutativeAlgebra/LemmaP.lean}{UpperTheory.V.Alg.lemmaP}
\begin{proof}
Write $\mathfrak g_j=\mathfrak g_{\Phi,j}$, $Z^1_j=Z^1_F(\mathfrak g_j)$ and $\mathrm{Ob}_j=\mathrm{Ob}_F(\mathfrak g_j)$. Since $\ell$ is odd, $2$ is a unit of $\mathcal O$.
\begin{enumerate}
\item[(i)] \emph{The ring.} Lemma~\ref{sha-counts-vd-model-lemma-p-universal} gives a universal framing $(i,r,\alpha)$ with values in some $R$, and Lemma~\ref{sha-counts-vd-model-lemma-p-grading} a weighted presentation $\pi_0$ of $R$ with the properties there; let $R_j$ be its parts.
\item[(ii)] \emph{Tangent elements.} Lemma~\ref{sha-counts-vd-model-lemma-p-tangent} gives $\tau$, the $z_t$ and the $\operatorname{wt}(t)$. Each $z_t$ lies in $\mathfrak m_R$ and in $R_{\operatorname{wt}(t)}$. If $z_t=r(g)_{ab}-r_0(g)_{ab}$ with $\operatorname{wt}(t)=w_a-w_b$, then $r(g)_{ab}\in R_{w_a-w_b}$ by (i); $r_0(g)_{ab}$ is the image of a constant, which has weight $0$, and it vanishes unless $w_a=w_b$, as $r_0$ normalises $\Phi$; and $z_t\in\mathfrak m_R$ since $r$ and $r_0$ both lift $\bar r$. If $z_t=(\alpha_v-1)_{ab}$, the same holds, as $\alpha_v\equiv1$ modulo $\mathfrak m_R$.
\item[(iii)] \emph{The presentation.} By (i), (ii) and property (2) of Lemma~\ref{sha-counts-vd-model-lemma-p-tangent}, Lemma~\ref{sha-counts-vd-model-lemma-p-relations} applies to $\pi_B$, $\pi_0$, $i$ and the $z_t$. It gives a weighted presentation $\pi$ of $R$ relative to $\pi_B$ and $i$, by the variables of $\pi_B$ and variables $Y_t$ with $\pi(Y_t)=z_t$ of weight $\operatorname{wt}(t)$, with relations $f_u$ of weights $w_f(u)$, with the parts $R_j$, and with minimal relations. So (1) and (2) hold, and by property (4) of Lemma~\ref{sha-counts-vd-model-lemma-p-tangent} the number of $t$ with $Y_t$ of weight $j$ is $\dim Z^1_j$.
\item[(iv)] \emph{An extension.} Let $P$ be the power series ring of $\pi$, with maximal ideal $\mathfrak m_P$, let $N$ be the ideal of $P$ generated by $\ker\pi_B$ and $\mathfrak m_P\ker\pi$, let $A=P/N$ with the quotient map $\pi_A:P\to A$, and let $p:A\to R$ be the surjective homomorphism induced by $\pi$, as $N\subseteq\ker\pi$. As in step (i) of the proof of Lemma~\ref{sha-counts-vd-model-cover}, $\ker\pi_B$ and $\ker\pi$ are generated by finitely many series of given weights, and $\mathfrak m_P$ is generated by $\mathfrak m_{\mathcal O}$ and the variables; so $N$ is generated by finitely many series of given weights, and $\pi_A$ is a weighted presentation of $A$ by Lemma~\ref{sha-counts-vd-presentations}(a), whose variables map into $\mathfrak m_A$. Let $i_A:B\to A$ be induced by $\pi_A$ on the series in the variables of $\pi_B$, which is possible since $\ker\pi_B\subseteq N$. Then:
\begin{itemize}
\item $p\circ i_A=i$, $p$ maps $A_j$ onto $R_j$, and $i_A$ maps $B_j$ into $A_j$, so the entries $(a,b)$ of the $i_A(\rho_v(s))$ lie in $A_{w_a-w_b}$;
\item $\mathfrak m_A\ker p=0$, since $\mathfrak m_P\ker\pi\subseteq N$;
\item $\ker p$ is generated by the images $\varepsilon_u=\pi_A(f_u)\in A_{w_f(u)}$, as $\ker\pi$ is generated by $\ker\pi_B$ and the $f_u$;
\item the $\varepsilon_u$ are independent: if $\sum_uc_u\varepsilon_u=0$, then $\sum_uc_uf_u\in N$, and $c_u\in\mathfrak m_{\mathcal O}$ by the minimality in (iii);
\item $\ker p$ is contained in the ideal generated by $i_A(\mathfrak m_B)$ and $\mathfrak m_A^2$. Let $x=\pi_A(g)\in\ker p$. The constant term $g_0$ of $g$ lies in $\mathfrak m_{\mathcal O}$, as $p(x)\in\mathfrak m_R$. Write $g\equiv g_0+\sum_sc_sX_s+\sum_tc_tY_t$ modulo $\mathfrak m_P^2$, with $c_s,c_t\in\mathcal O$ and $X_s$ the variables of $\pi_B$. For a homomorphism of $\mathcal O$-algebras $q$ from $A$ to a complete local $\mathcal O$-algebra $A''$, the images under $q$ of $g_0$ and of the $\pi_A(X_s)$ lie in $q(i_A(\mathfrak m_B))$, so $q(x)-\sum_tc_tq(\pi_A(Y_t))$ lies in the ideal generated by $q(i_A(\mathfrak m_B))$ and $\mathfrak m_{A''}^2$. With $q=p$, as $p(x)=0$ and $p(\pi_A(Y_t))=z_t$, property (3) of Lemma~\ref{sha-counts-vd-model-lemma-p-tangent} gives $c_t\in\mathfrak m_{\mathcal O}$. With $q$ the identity, $x$ lies in the ideal generated by $i_A(\mathfrak m_B)$, $\mathfrak m_A^2$ and the $c_t\pi_A(Y_t)$, which lie in $\mathfrak m_A^2$.
\end{itemize}
\item[(v)] \emph{The bound.} By (iv), Lemma~\ref{sha-counts-vd-model-lemma-p-classes} applies to $p$, $\pi_A$, the $\varepsilon_u$, $i_A$ and the framing $(i,r,\alpha)=(p\circ i_A,r,\alpha)$, whose entries lie in the $p(A_j)=R_j$ by (i). Let $o_j(u)$ be its elements. Fix $j$, and let $c_u\in\mathcal O$ for the $u$ with $w_f(u)=j$, with $\sum_u\bar c_u\,o_j(u)=0$; put $c_u=0$ for the other $u$. We show that every $c_u$ lies in $\mathfrak m_{\mathcal O}$. Suppose that $c_{u_0}$ is a unit, and let $d=c_{u_0}^{-1}$. Let $I$ be the ideal of $A$ generated by the $\delta_u=\varepsilon_u-c_ud\,\varepsilon_{u_0}$, let $p':A\to A'=A/I$ be the quotient map, $p'':A'\to R$ the homomorphism induced by $p$, which kills $I\subseteq\ker p$, and $\eta=p'(d\varepsilon_{u_0})$. Then $p'(\varepsilon_u)=c_u\eta$ for every $u$, $p''(\eta)=0$, and $\mathfrak m_{A'}\eta=0$, as $\mathfrak m_A\ker p=0$. By Lemma~\ref{sha-counts-vd-model-lemma-p-classes} there is a framing $(p'\circ i_A,r',\alpha')$ with values in $A'$ that $p''$ carries to $(i,r,\alpha)$. By universality it is the image of $(i,r,\alpha)$ under a homomorphism $s:R\to A'$, so $s\circ i=p'\circ i_A$; and $p''\circ s$ carries $(i,r,\alpha)$ to itself, so $p''\circ s$ is the identity, by the uniqueness in the universal property.
The homomorphism $s$ is surjective. By the complete Nakayama lemma \cite[Theorem~8.4]{Mat86} it suffices that $\mathfrak m_{A'}$ lies in the ideal generated by $s(\mathfrak m_R)$ and $\mathfrak m_{A'}^2$. Let $x\in\mathfrak m_{A'}$; then $x-s(p''(x))=p'(y)$ with $y\in\ker p$, since $p''\circ s$ is the identity, and by (iv) $y$ lies in the ideal generated by $i_A(\mathfrak m_B)$ and $\mathfrak m_A^2$. As $p'(i_A(\mathfrak m_B))=s(i(\mathfrak m_B))\subseteq s(\mathfrak m_R)$ and $s(p''(x))\in s(\mathfrak m_R)$, the claim follows.
So $p'(\varepsilon_{u_0})=s(y)$ for some $y\in R$, and $y=p''(s(y))=p(\varepsilon_{u_0})=0$; hence $\varepsilon_{u_0}\in I$, say $\varepsilon_{u_0}=\sum_ua_u\delta_u$ with $a_u\in A$. Choose $b_u\in\mathcal O$ with $a_u-b_u\in\mathfrak m_A$; as $\delta_u\in\ker p$ and $\mathfrak m_A\ker p=0$, $\varepsilon_{u_0}=\sum_ub_u\delta_u=\sum_ub_u\varepsilon_u-\big(\sum_vb_vc_vd\big)\varepsilon_{u_0}$. By the independence in (iv), $b_u-e_u\in\mathfrak m_{\mathcal O}$ for every $u$, where $e_{u_0}=1+\sum_vb_vc_vd$ and $e_u=0$ for $u\ne u_0$. Multiplying by $c_ud$ and summing over $u$ gives $\sum_ub_uc_ud-e_{u_0}c_{u_0}d=-1\in\mathfrak m_{\mathcal O}$, a contradiction.
Hence the $o_j(u)$ with $w_f(u)=j$ are linearly independent over $\mathbf F$, and their number is at most $\dim\mathrm{Ob}_j$.
\item[(vi)] \emph{Equality.} If $j$ is not of the form $w_a-w_b$, then $\mathfrak g_j=0$ and $\mathrm{Ob}_j=0$. For each of the finitely many $j$ of the form $w_a-w_b$, add to the relations $\dim\mathrm{Ob}_j-\#\{u:w_f(u)=j\}$ copies of the series $0$, of weight $j$. This does not change the kernel, and it gives (3). Without weights, and for the groups $\mathcal G_n$ of Clozel, Harris and Taylor, compare \cite[Corollary~2.2.12]{CHT08}.
\end{enumerate}
% Written out after their statements: framed_tangent (FramedTangent.lean, 651 lines, through CoSpan, CoRes, HhGroup, the private framed_iff, liftOf and abstract) as Lemma -lemma-p-tangent with Lemma -lemma-p-tangent-abstract; obs_class (ObsClass.lean, 261 lines, through hh_section, ObsCocycle, ObsCorrect (with cocycle_normal_form and hh_small), HomogSpan, LocGradedTop and KerCoef) as Lemma -lemma-p-classes, with its sub-lemma -lemma-p-classes-local (loc_graded, with Lemma -lemma-p-classes-local-split for split_data). The proofs of Lemmas -lemma-p-grading (framed_grading), -lemma-p-relations-parts (pres_indep, as Lemma sha-counts-vd-presentations(e)) and -lemma-p-relations (rel_pres, with pres_new in step (ii), Lemma -lemma-p-relations-kernel for ker_gen and rel_pres_min in step (vi)) are written out after their statements. Written out here: framed_univ (Lemma -lemma-p-universal), and lemmaP itself, with the private lp_ker, lp_detect (with co_surj, co_quot) and lp_pad in steps (iv) to (vi).
\end{proof}

\begin{lemma}\label{sha-counts-vd-model-lemma-p-kill-section}
Let $\mathcal O$ be a complete local noetherian ring with residue field $\mathbf F$ in which $2$ is a unit, and let $\iota$, $X$, $w$ and a form datum $\Phi$ be as in Definition~\ref{sha-counts-vd-model-form-datum}. Let $R$ and $R'$ be complete local noetherian $\mathcal O$-algebras with residue field $\mathbf F$, and $p:R\to R'$ a surjective homomorphism of $\mathcal O$-algebras with $\mathfrak m_R\ker p=0$.
\begin{enumerate}
\item There is $M\ge0$ such that $0$ is the only element of $\ker p$ that lies in $\mathfrak m_R^M$.
\item There is a map $\sigma$ from the $\Phi$-unitary matrices in $M_\iota(R')$ to the $\Phi$-unitary matrices in $M_\iota(R)$ with $p(\sigma(h'))=h'$, such that for every $n\ge0$ and all $\Phi$-unitary $h',h''$ with $h'\equiv h''$ modulo $\mathfrak m_{R'}^{\max(n,M)}$, $\sigma(h')\equiv\sigma(h'')$ modulo $\mathfrak m_R^n$. In particular, if $f$ is a function from a topological space to the $\Phi$-unitary matrices in $M_\iota(R')$ whose entries are continuous for the $\mathfrak m_{R'}$-adic topology, then the entries of $\sigma\circ f$ are continuous for the $\mathfrak m_R$-adic topology.
\end{enumerate}
\end{lemma}
\leanref{Development/AutomorphicArithmetic/AutomorphicVarieties/CommutativeAlgebra/KillLoc.lean}{UpperTheory.V.Alg.kill_loc}
\begin{proof}
(1) The ideal $K=\ker p$ is killed by $\mathfrak m_R$, so it is a vector space over $R/\mathfrak m_R=\mathbf F$, finite-dimensional as $R$ is noetherian. So the decreasing sequence of subspaces $K\cap\mathfrak m_R^n$ is eventually constant, equal to $K\cap\mathfrak m_R^M$ say, and then $K\cap\mathfrak m_R^M=\bigcap_n(K\cap\mathfrak m_R^n)=0$ by the theorem of Krull \cite[Theorem~8.10]{Mat86}.

(2) First, if $a\in R$ and $p(a)\in\mathfrak m_{R'}^n$, then $a-k\in\mathfrak m_R^n$ for some $k\in K$, since $p$ maps $\mathfrak m_R^n$ onto $\mathfrak m_{R'}^n$. For each class $C$ of matrices over $R'$ modulo $\mathfrak m_{R'}^M$ that contains a $\Phi$-unitary matrix $h_C$, choose by Lemma~\ref{sha-counts-vd-model-lemma-p-unitary}(2)(a) a $\Phi$-unitary $H_C$ over $R$ with $p(H_C)=h_C$. Let $h'$ be $\Phi$-unitary, in the class $C$. For each entry choose $y_{ab}\in R$ with $p(y_{ab})=h'_{ab}$ and $y_{ab}-(H_C)_{ab}\in\mathfrak m_R^M$, by the first remark applied to a lift of $h'_{ab}$ minus $(H_C)_{ab}$; put $\sigma(h')=y$. Then $y$ is $\Phi$-unitary. The matrix $y^{\mathrm t}J_Ry-J_R$ maps to $h'^{\mathrm t}J_{R'}h'-J_{R'}=0$, and it is congruent modulo $\mathfrak m_R^M$ to $H_C^{\mathrm t}J_RH_C-J_R=0$, so it vanishes by (1). If $\operatorname{bl}(a)\ne\operatorname{bl}(b)$, then $y_{ab}\in K\cap\mathfrak m_R^M$, as $h'_{ab}=0=(H_C)_{ab}$, so $y_{ab}=0$. And $y\equiv1$, as $p(y)=h'\equiv1$.
Let $h'\equiv h''$ modulo $\mathfrak m_{R'}^{n'}$, with $n'=\max(n,M)$. They lie in the same class $C$, so $d=\sigma(h')-\sigma(h'')$ is $\equiv0$ modulo $\mathfrak m_R^M$. Its entries map to elements of $\mathfrak m_{R'}^{n'}$, so $d_{ab}-k\in\mathfrak m_R^{n'}$ for some $k\in K$. Then $k\in K\cap\mathfrak m_R^M$, so $k=0$, and $d_{ab}\in\mathfrak m_R^n$. The last statement follows, as the entries of $f$ are locally constant modulo $\mathfrak m_{R'}^{\max(n,M)}$ for every $n$.
\end{proof}

\begin{lemma}\label{sha-counts-vd-model-lemma-p-kill-coboundary}
Let $\mathcal O$ be a local ring with residue field $\mathbf F$, let $\iota$, $X$, $w$ and a form datum $\Phi$ be as in Definition~\ref{sha-counts-vd-model-form-datum}, let $\Gamma$ be a topological group, $\bar r:\Gamma\to M_\iota(\mathbf F)$ a homomorphism that normalises $\Phi$, and $N$ a subgroup of $\Gamma$. Suppose that for every $j\in X$, every locally constant $\kappa:\Gamma\times\Gamma\to\mathfrak g_{\Phi,j}(\mathbf F)$ with $g\cdot\kappa(h,h')+\kappa(g,hh')=\kappa(gh,h')+\kappa(g,h)$ and $\kappa(gn,hn')=\kappa(g,h)$ for all $g,h,h'\in\Gamma$ and $n,n'\in N$, where $\Gamma$ acts through $\bar r$, is $dc$ for a locally constant $c:\Gamma\to\mathfrak g_{\Phi,j}(\mathbf F)$ with $c(gn)=c(g)$ for $n\in N$.
Let $R$ be a complete local noetherian $\mathcal O$-algebra with residue field $\mathbf F$, $K$ an ideal of $R$ with $\mathfrak m_RK=0$, and $\rho:\Gamma\to M_\iota(R)$ a function whose residue is $\bar r$. Let $\kappa:\Gamma\times\Gamma\to\mathfrak g_\Phi(R)$ be locally constant, with entries in $K$, such that
\[
 \rho(s)\kappa(t,u)\rho(s^{-1})+\kappa(s,tu)=\kappa(st,u)+\kappa(s,t)\quad\text{and}\quad\kappa(sn,tn')=\kappa(s,t)
\]
for all $s,t,u\in\Gamma$ and $n,n'\in N$. Then there is a locally constant $a:\Gamma\to\mathfrak g_\Phi(R)$ with entries in $K$, with $a(sn)=a(s)$ for $s\in\Gamma$ and $n\in N$, and $a(s)+\rho(s)a(t)\rho(s^{-1})-a(st)=\kappa(s,t)$ for all $s,t\in\Gamma$.
\end{lemma}
\leanref{Development/AutomorphicArithmetic/AutomorphicVarieties/CommutativeAlgebra/KillLoc.lean}{UpperTheory.V.Alg.kill_loc}
\begin{proof}
\begin{enumerate}
\item[(i)] \emph{Coordinates.} As $\mathfrak m_R$ kills $K$, $K$ is a vector space over $R/\mathfrak m_R=\mathbf F$, finite-dimensional as $R$ is noetherian; choose a basis $k_1,\dots,k_d$. The map $\psi:\mathbf F^d\to K$, $(c_u)\mapsto\sum_u\tilde c_uk_u$, with $\tilde c_u\in R$ any lift of $c_u$, is an isomorphism of groups with $x\psi(c)=\psi(\bar xc)$ for $x\in R$ with residue $\bar x$. Applied to the entries, it gives a bijection $\Psi$ from $M_\iota(\mathbf F)^d$ onto the matrices with entries in $K$, with $A\Psi(Y)=\Psi((\bar AY_u)_u)$, $\Psi(Y)A=\Psi((Y_u\bar A)_u)$ and $\Psi(Y)^{\mathrm t}=\Psi((Y_u^{\mathrm t})_u)$ for $A\in M_\iota(R)$ with residue $\bar A$, and with $\Psi(Y)_{ab}=0$ if and only if $(Y_u)_{ab}=0$ for every $u$. As $J_R$ has the residue $J_{\mathbf F}$, $\Psi(Y)\in\mathfrak g_\Phi(R)$ if and only if every $Y_u\in\mathfrak g_\Phi(\mathbf F)$.
\item[(ii)] Write $\kappa=\Psi((\kappa_u)_u)$. Each $\kappa_u:\Gamma\times\Gamma\to\mathfrak g_\Phi(\mathbf F)$ is locally constant, invariant under $N$ as $\kappa$ is, and satisfies $\bar r(s)\kappa_u(t,v)\bar r(s)^{-1}+\kappa_u(s,tv)=\kappa_u(st,v)+\kappa_u(s,t)$, by (i), as $\rho(s^{-1})$ has the residue $\bar r(s)^{-1}$.
\item[(iii)] \emph{Weights.} For $j\in X$ and $Y\in M_\iota(\mathbf F)$, let $Y^{(j)}$ agree with $Y$ at the entries $(a,b)$ with $w_a-w_b=j$ and vanish at the others. If $Y\in\mathfrak g_\Phi(\mathbf F)$, then $Y^{(j)}\in\mathfrak g_{\Phi,j}(\mathbf F)$: the entry $(a,b)$ of $Y^{\mathrm t}J+JY$ involves only entries $Y_{ca}$ with $J_{cb}\ne0$ and $Y_{cb}$ with $J_{ac}\ne0$, all of weight $c-w_a-w_b$, as $J_{xy}\ne0$ implies $w_x+w_y=c$. If $A_{ab}=0$ whenever $w_a\ne w_b$, as for the values of $\bar r$, then $(AY)^{(j)}=AY^{(j)}$ and $(YA)^{(j)}=Y^{(j)}A$. And $Y=\sum_jY^{(j)}$, over the finitely many $j$ of the form $w_a-w_b$.
\item[(iv)] So for each $u$ and $j$, $\kappa_u^{(j)}$ is a locally constant function $\Gamma\times\Gamma\to\mathfrak g_{\Phi,j}(\mathbf F)$ that satisfies the hypotheses on $\kappa$ in the statement, and it is $de_{u,j}$ for a locally constant $e_{u,j}$ with $e_{u,j}(gn)=e_{u,j}(g)$. Put $a=\Psi\big((\sum_je_{u,j})_u\big)$. It is locally constant, invariant under $N$ on the right, with values in $\mathfrak g_\Phi(R)$ and entries in $K$ by (i), and by (i) and (iii), $a(s)+\rho(s)a(t)\rho(s^{-1})-a(st)=\Psi\big((\sum_jde_{u,j}(s,t))_u\big)=\Psi((\kappa_u(s,t))_u)=\kappa(s,t)$.
\end{enumerate}
\end{proof}

\begin{lemma}\label{sha-counts-vd-model-lemma-p-kill}
Let $\mathcal O$ be a complete local noetherian ring with residue field $\mathbf F$ in which $2$ is a unit, let $\iota$, $X$, $w$ and a form datum $\Phi$ be as in Definition~\ref{sha-counts-vd-model-form-datum}, and let $\Gamma$ be a topological group. Let $\bar r:\Gamma\to M_\iota(\mathbf F)$ be a homomorphism and $r_0$ a continuous lift of $\bar r$ with values in $\mathcal O$, both normalising $\Phi$, and let $N$ be a normal subgroup of $\Gamma$ with $r_0(n)=1$ for $n\in N$. Suppose that for every $j\in X$, every locally constant $\kappa:\Gamma\times\Gamma\to\mathfrak g_{\Phi,j}(\mathbf F)$ with $g\cdot\kappa(h,h')+\kappa(g,hh')=\kappa(gh,h')+\kappa(g,h)$ and $\kappa(gn,hn')=\kappa(g,h)$ for all $g,h,h'\in\Gamma$ and $n,n'\in N$, where $\Gamma$ acts through $\bar r$, is $dc$ for a locally constant $c:\Gamma\to\mathfrak g_{\Phi,j}(\mathbf F)$ with $c(gn)=c(g)$ for $n\in N$.
Then the property of a lift $r$ that $r(n)=1$ for every $n\in N$ is a local condition for $(\bar r,r_0)$ with tangent group $N$ (Definition~\ref{sha-counts-vd-model-form-datum}(5)).
\end{lemma}
\leanref{Development/AutomorphicArithmetic/AutomorphicVarieties/CommutativeAlgebra/KillLoc.lean}{UpperTheory.V.Alg.kill_loc}
\leanref{Development/AutomorphicArithmetic/AutomorphicVarieties/CommutativeAlgebra/FramedTwoCoboundaries.lean}{UpperTheory.V.Alg.killCond}
\begin{proof}
Let $L$ be the property; the properties (a) to (e) of Definition~\ref{sha-counts-vd-model-form-datum}(5) hold for $L$.
\begin{enumerate}
\item[(a)] \emph{Closed.} For a lift $r$ with values in $R$, let $I$ be the ideal generated by the entries of the $r(n)-1$, $n\in N$. Then $f\circ r$ has $L$ if and only if $f(I)=0$.
\item[(b)] \emph{Conjugation.} For $n\in N$, $r'(n)h=hr(n)=h$, and $h$ is invertible by Lemma~\ref{sha-counts-vd-model-lemma-p-unitary}(1); so $r'(n)=1$.
\item[(c)] \emph{Torus.} For $n\in N$, $r'(n)_{ab}=t(w_a-w_b)1_{ab}$, which is $1_{ab}$, as $t(0)=1$.
\item[(e)] \emph{Tangent.} As $r_1(n)=1$ for $n\in N$, $r_2$ has $L$ if and only if $r_2(n)=r_1(n)$ for every $n\in N$.
\item[(d)] \emph{Lifts.} Let $p:R\to R'$ be surjective with $\mathfrak m_R\ker p=0$, put $K=\ker p$, and let $r'$ be a lift with values in $R'$ that is $\Phi$-close to $r_0$ and has $L$. Write $\rho(s)=r_0(s)_R$ and $\rho'(s)=r_0(s)_{R'}$, so that $p\circ\rho=\rho'$, the residue of $\rho$ is $\bar r$, $\rho(n)=1$ and $\rho'(n)=1$ for $n\in N$, and $\rho$ is continuous. If $h\equiv1$ modulo $\mathfrak m_R$ and the entries of $Z$ lie in $K$, then $hZ=Z=Zh$, since $\mathfrak m_R$ kills $K$. A lift $r$ of $r'$ is constructed in (i) to (v).
\end{enumerate}
\begin{enumerate}
\item[(i)] \emph{A lift of the unitary part.} Let $M$ and $\sigma$ be as in Lemma~\ref{sha-counts-vd-model-lemma-p-kill-section}, and put $h(s)=\sigma\big(r'(s)\rho'(s^{-1})\big)$; the matrices $r'(s)\rho'(s^{-1})$ are $\Phi$-unitary, as $r'$ is $\Phi$-close to $r_0$, and their entries are continuous in $s$. So $h(s)$ is $\Phi$-unitary, $p(h(s))=r'(s)\rho'(s^{-1})$, and the entries of $h$ are continuous. For $n\in N$, $r'(sn)\rho'((sn)^{-1})=r'(s)r'(n)\rho'(n^{-1})\rho'(s^{-1})=r'(s)\rho'(s^{-1})$, so $h(sn)=h(s)$.
\item[(ii)] \emph{The defect.} Let $\kappa(s,t)=h(s)\rho(s)h(t)\rho(s^{-1})-h(st)$. Its two terms are $\Phi$-unitary, by Lemma~\ref{sha-counts-vd-model-lemma-p-unitary}(1), and they have the same image $r'(s)r'(t)\rho'((st)^{-1})$ under $p$. So $\kappa(s,t)\in\mathfrak g_\Phi(R)$ by Lemma~\ref{sha-counts-vd-model-lemma-p-unitary}(2)(b), and its entries lie in $K$. It is locally constant: its entries are continuous for the $\mathfrak m_R$-adic topology, so near $(s,t)$ the values of $\kappa$ are congruent to $\kappa(s,t)$ modulo $\mathfrak m_R^M$, and the differences lie in $K\cap\mathfrak m_R^M=0$.
\item[(iii)] \emph{The cocycle identity.} $\rho(s)\kappa(t,u)\rho(s^{-1})+\kappa(s,tu)=\kappa(st,u)+\kappa(s,t)$. By the last remark in (d), $\rho(s)\kappa(t,u)\rho(s^{-1})=h(s)\rho(s)\kappa(t,u)\rho(s^{-1})$ and $\kappa(s,t)=\kappa(s,t)\rho(st)h(u)\rho((st)^{-1})$. Expanding with $\rho(s)\rho(t)=\rho(st)$, both sides become $h(s)\rho(s)h(t)\rho(t)h(u)\rho((st)^{-1})-h(stu)$. Also $\kappa(sn,tn')=\kappa(s,t)$ for $n,n'\in N$: $h(sn)=h(s)$, $\rho(n)=1$, and $h(sntn')=h(st)$, as $sntn'=st(t^{-1}nt)n'$ with $t^{-1}nt\in N$.
\item[(iv)] By (ii), (iii) and the hypothesis on $2$-cocycles, Lemma~\ref{sha-counts-vd-model-lemma-p-kill-coboundary} gives a locally constant $a:\Gamma\to\mathfrak g_\Phi(R)$ with entries in $K$, with $a(sn)=a(s)$ for $n\in N$ and $a(s)+\rho(s)a(t)\rho(s^{-1})-a(st)=\kappa(s,t)$.
\item[(v)] \emph{The lift.} Put $r(s)=(h(s)-a(s))\rho(s)$. The matrix $h(s)-a(s)$ is $\Phi$-unitary, by Lemma~\ref{sha-counts-vd-model-lemma-p-unitary}(2)(c), as $p(a(s))=0$. By the last remark in (d), $a(s)\rho(s)h(t)\rho(s^{-1})=a(s)$ and $(h(s)-a(s))\rho(s)a(t)\rho(s^{-1})=\rho(s)a(t)\rho(s^{-1})$, so
\[
 r(s)r(t)=\big(h(s)\rho(s)h(t)\rho(s^{-1})-a(s)-\rho(s)a(t)\rho(s^{-1})\big)\rho(st)=\big(h(st)+\kappa(s,t)-a(s)-\rho(s)a(t)\rho(s^{-1})\big)\rho(st),
\]
which is $r(st)$ by (iv). As $r(1)$ is invertible and $r(1)=r(1)^2$, $r(1)=1$; so $r$ is a homomorphism. Its entries are continuous, by (i) and (iv) and the continuity of $\rho$; its residue is $\bar r$, as $h(s)-a(s)\equiv1$; $p(r(s))=r'(s)\rho'(s^{-1})\rho'(s)=r'(s)$; and $r(s)\rho(s^{-1})=h(s)-a(s)$ is $\Phi$-unitary, so $r$ is $\Phi$-close to $r_0$. For $n\in N$, $r(n)=(h(1\cdot n)-a(1\cdot n))\rho(n)=h(1)-a(1)=r(1)=1$, so $r$ has $L$.
\end{enumerate}
\end{proof}

\begin{lemma}\label{sha-counts-vd-model-lemma-p-siegel-transfer}
Let $\Gamma$ be a group, $k$ a field, and $M$ a $k$-vector space with a $k$-linear action of $\Gamma$. Let $V\subseteq H$ be subgroups of $\Gamma$, with $V$ normal and acting trivially on $M$, and $H$ of finite index $[\Gamma:H]$ that is invertible in $k$. Let $\kappa:\Gamma\times\Gamma\to M$ satisfy $g\cdot\kappa(h,h')+\kappa(g,hh')=\kappa(gh,h')+\kappa(g,h)$ and $\kappa(gv,h)=\kappa(g,h)=\kappa(g,hv)$ for all $g,h,h'\in\Gamma$ and $v\in V$. Let $a_0:\Gamma\to M$ satisfy $a_0(gv)=a_0(g)$ for $g\in\Gamma$ and $v\in V$, and $a_0(st)=a_0(s)+s\cdot a_0(t)-\kappa(s,t)$ for $s,t\in H$.
Then there is $c:\Gamma\to M$ with $c(gv)=c(g)$ for $g\in\Gamma$ and $v\in V$, and $\kappa(g,h)=c(g)+g\cdot c(h)-c(gh)$ for all $g,h\in\Gamma$.
\end{lemma}
\leanref{Development/AutomorphicArithmetic/AutomorphicVarieties/CommutativeAlgebra/ExistsLocallyConstantDOneEqCocycleOfTameInertia.lean}{second_cohomology_transfer}
\begin{proof}
For $g\in\Gamma$ let $\sigma(g)$ be a fixed element of the coset $gH$, depending only on $gH$; so $\sigma(g)^{-1}g\in H$, and $\sigma(gh)=\sigma(g)$ for $h\in H$. Put
\[
 f(g)=\kappa\big(\sigma(g),\sigma(g)^{-1}g\big)-\sigma(g)\cdot a_0\big(\sigma(g)^{-1}g\big).
\]
\begin{enumerate}
\item[(i)] For $g\in\Gamma$ and $h\in H$, $f(gh)=f(g)-g\cdot a_0(h)+\kappa(g,h)$. Indeed, let $s=\sigma(g)$ and $x=s^{-1}g\in H$. Then $f(gh)=\kappa(s,xh)-s\cdot a_0(x)-g\cdot a_0(h)+s\cdot\kappa(x,h)$, by the hypothesis on $a_0$, and $s\cdot\kappa(x,h)+\kappa(s,xh)=\kappa(g,h)+\kappa(s,x)$ by the cocycle identity.
\item[(ii)] For $v\in V$, $f(gv)=f(g)$, since $\sigma(gv)=\sigma(g)$, and $\kappa$ and $a_0$ are invariant under $V$ on the right.
\item[(iii)] Let $D(s,g)=s\cdot f(g)+\kappa(s,g)-f(sg)$ for $s,g\in\Gamma$. By the cocycle identity, $s\cdot D(t,g)+D(s,tg)=\kappa(s,t)+D(st,g)$. By (i) and the cocycle identity at $(s,g,h)$, $D(s,gh)=D(s,g)$ for $h\in H$. For $v\in V$, $D(sv,g)=D(s,g)$: $svg=sg(g^{-1}vg)$ with $g^{-1}vg\in V$, so $f(svg)=f(sg)$ by (ii); $\kappa(sv,g)=\kappa(s,g)$; and $v$ acts trivially.
\item[(iv)] Let $A(s)=\sum_{x\in\Gamma/H}D(s,\tilde x)$, for representatives $\tilde x$ of the cosets, which by (iii) does not depend on their choice. Left multiplication by $t$ permutes $\Gamma/H$, so $A(s)=\sum_xD(s,t\tilde x)$, and by (iii)
\[
 A(s)+s\cdot A(t)-A(st)=\sum_{x\in\Gamma/H}\big(D(s,t\tilde x)+s\cdot D(t,\tilde x)-D(st,\tilde x)\big)=[\Gamma:H]\,\kappa(s,t).
\]
\end{enumerate}
Put $c=[\Gamma:H]^{-1}A$. Then $c(g)+g\cdot c(h)-c(gh)=\kappa(g,h)$ by (iv), and $c(gv)=c(g)$ for $v\in V$ by (iii).
\end{proof}

\begin{lemma}\label{sha-counts-vd-model-lemma-p-siegel-tame}
Let $\ell$ be a prime, $k$ a field of characteristic $\ell$, and $\Gamma$ a compact, totally disconnected Hausdorff topological group. Let $M$ be a $k$-vector space that is a finite set, with a $k$-linear action of $\Gamma$ whose kernel is open. Let $N\subseteq I$ be closed normal subgroups of $\Gamma$ such that $N$ acts trivially on $M$ and the index $[I:N]$ is finite and prime to $\ell$. Suppose that there is $\phi\in\Gamma$ such that $I$ and $\phi$ generate a dense subgroup of $\Gamma$, and that for every $n\ge0$ there is an open subgroup of $\Gamma$ of index $\ell^n$ that contains $I$. Let $\kappa:\Gamma\times\Gamma\to M$ be locally constant, with $g\cdot\kappa(h,h')+\kappa(g,hh')=\kappa(gh,h')+\kappa(g,h)$ and $\kappa(gn,hn')=\kappa(g,h)$ for all $g,h,h'\in\Gamma$ and $n,n'\in N$.
Then there is a locally constant $c:\Gamma\to M$ with $c(gn)=c(g)$ for $g\in\Gamma$ and $n\in N$, and $\kappa(g,h)=c(g)+g\cdot c(h)-c(gh)$ for all $g,h\in\Gamma$.
\end{lemma}
\leanref{Development/AutomorphicArithmetic/AutomorphicVarieties/CommutativeAlgebra/ExistsLocallyConstantDOneEqCocycleOfTameInertia.lean}{UpperTheory.V.Alg.exists_locallyConstant_dOne_eq_cocycle_of_tame_inertia}
\begin{proof}
An open subgroup $W$ of $\Gamma$ that contains $I$ and $\phi$ is $\Gamma$: it is closed, and it contains the dense subgroup generated by $I$ and $\phi$. An open subgroup has finite index, as $\Gamma$ is compact, so it has finitely many conjugates, and its normal core, the intersection of its conjugates, is an open normal subgroup.
\begin{enumerate}
\item[(i)] \emph{Invariance of $\kappa$.} Let $V_0$ be the set of the $v\in\Gamma$ with $\kappa(gv,h)=\kappa(g,h)=\kappa(g,hv)$ for all $g,h\in\Gamma$; it is a subgroup. It is open. For each $(g,h)\in\Gamma\times\Gamma$, as $\kappa$ is locally constant, there are neighbourhoods $W_{g,h}$ of $1$ and $U_{g,h}$ of $(g,h)$ such that $\kappa(g'v,h')=\kappa(g',h')=\kappa(g',h'v)$ for $v\in W_{g,h}$ and $(g',h')\in U_{g,h}$. Finitely many $U_{g,h}$ cover the compact space $\Gamma\times\Gamma$, and the intersection of the corresponding $W_{g,h}$ is a neighbourhood of $1$ contained in $V_0$. By the hypothesis on $\kappa$, $N\subseteq V_0$.
\item[(ii)] \emph{Two open normal subgroups.} Let $V_1$ be the normal core of the intersection of $V_0$ with the kernel of the action; it is open and normal, and it contains $N$, as $N$ is normal. Let $r$ be the order of the image of $\phi$ in the finite group $\Gamma/V_1$, and write $r=\ell^em$ with $\ell\nmid m$. Let $U$ be an open subgroup of index $\ell^{e+1}$ that contains $I$, and $U_1$ its normal core, which contains $I$, as $I$ is normal. The preimage in $\Gamma$ of the subgroup of $\Gamma/U_1$ generated by the image of $\phi$ is open and contains $I$ and $\phi$, so it is $\Gamma$. So $\Gamma/U_1$ is cyclic, generated by the image of $\phi$, whose order $[\Gamma:U_1]$ is divisible by $[\Gamma:U]=\ell^{e+1}$. Let $V=V_1\cap U_1$, an open normal subgroup that contains $N$, is contained in $V_0$ and acts trivially on $M$.
\item[(iii)] \emph{A sequence.} Define $c_n\in M$ for $n\ge0$ by $c_0=\kappa(1,1)$ and $c_{n+1}=\phi\cdot c_n-\kappa(\phi,\phi^n)$. The cocycle identity at $(1,1,t)$ gives $\kappa(1,t)=\kappa(1,1)$ for every $t$, and by induction on $i$, with the cocycle identity at $(\phi,\phi^i,\phi^j)$,
\[
 c_{i+j}=c_i+\phi^i\cdot c_j-\kappa(\phi^i,\phi^j)\qquad(i,j\ge0).
\]
\item[(iv)] \emph{Periodicity.} If $\phi^i$ and $\phi^j$ have the same image in $\Gamma/V$, then $c_i=c_j$. Indeed $\phi^r\in V_1$, so $\phi^r$ acts trivially and $\kappa(\phi^r,t)=\kappa(1,t)=\kappa(1,1)$; by (iii) and induction on $n$, $c_{rn+j}=n(c_r-\kappa(1,1))+c_j$, so $c_i$ depends only on $i$ modulo $r\ell$, as $\ell=0$ in $k$. Now $\phi^{i-j}\in V_1$, so $i\equiv j$ modulo $r$; and $\phi^{i-j}\in U_1$, so $i\equiv j$ modulo $\ell^{e+1}$ by (ii). As $r\ell=\ell^{e+1}m$ with $\ell^{e+1}$ and $m$ coprime, $i\equiv j$ modulo $r\ell$.
\item[(v)] \emph{A trivialization on a subgroup.} Let $a_0(g)=c_n$ if $gV=\phi^nV$ for some $n\ge0$, and $a_0(g)=0$ if there is no such $n$; by (iv) this is well defined, and $a_0(gv)=a_0(g)$ for $v\in V$. Let $H$ be the preimage in $\Gamma$ of the subgroup of $\Gamma/V$ generated by the image of $\phi$. As $\Gamma/V$ is finite, every element of $H$ is $\phi^nv$ with $n\ge0$ and $v\in V$. For $s=\phi^iv$ and $t=\phi^jv'$ in $H$, with $v,v'\in V$, $stV=\phi^{i+j}V$, as $V$ is normal, and $\kappa(s,t)=\kappa(\phi^i,\phi^j)$, as $V\subseteq V_0$; so by (iii), $a_0(st)=c_{i+j}=a_0(s)+s\cdot a_0(t)-\kappa(s,t)$, as $v$ acts trivially.
\item[(vi)] \emph{The index.} The subgroup $HI$ is open and contains $I$ and $\phi$, so it is $\Gamma$. So $I$ acts transitively on $\Gamma/H$, with stabilizer $H\cap I$ at the coset $H$, and $[\Gamma:H]=[I:H\cap I]$, which divides $[I:N]$, as $N\subseteq V\subseteq H$. So $[\Gamma:H]$ is prime to $\ell$, and invertible in $k$.
\end{enumerate}
By (i), (ii), (v) and (vi), Lemma~\ref{sha-counts-vd-model-lemma-p-siegel-transfer} applies to $H$, $V$, $\kappa$ and $a_0$, and gives $c$ with $c(gv)=c(g)$ for $v\in V$ and $\kappa(g,h)=c(g)+g\cdot c(h)-c(gh)$. Then $c(gn)=c(g)$ for $n\in N\subseteq V$, and $c$ is locally constant, being constant on the cosets of the open subgroup $V$.
\end{proof}

\begin{corollary}\label{sha-counts-vd-model-lemma-p-siegel}
Let $D$ be a Siegel datum, with the notation of Definitions~\ref{sha-siegel-datum} and~\ref{sha-counts-vd-model-framing}, let $G_S$, $p_S$, $Q$, the $N_q$, $\bar r_S$ and $r_{0,S}$ be as in Lemma~\ref{sha-counts-vd-model-descent}, and let $Z^1_j$ and $\mathrm{Ob}_j$ be as in Definition~\ref{sha-counts-vd-model-framed-cohomology}; suppose that $\mathrm{Ob}_j$ is finite-dimensional for every $j\in\mathbf Z^2$. Let $B$ be a complete local noetherian $\mathcal O$-algebra with residue field $\mathbf F$, with a weighted presentation (Definition~\ref{sha-counts-vd-model}(1)) whose variables have weights in $\mathbf Z^2$ and map into $\mathfrak m_B$, and for $v\in T$ let $\rho_v:D_v\to M_4(B)$ be a continuous lift of $\bar r_0$ such that, for every $g\in D_v$, $\rho_v(g)$ is symplectic with multiplier $\theta(g)$ and its entry $(a,b)$ lies in the part of the weight of $(a,b)$.
For a complete local noetherian $\mathcal O$-algebra $R$ with residue field $\mathbf F$, an \emph{$S$-framing} with values in $R$ consists of a homomorphism of $\mathcal O$-algebras $i:B\to R$, a continuous lift $r:G_S\to M_4(R)$ of $\bar r_S$, and matrices $\alpha_v\in M_4(R)$ for $v\in T$, such that: $r(g)r_{0,S}(g)^{-1}$ is symplectic with multiplier $1$ for every $g\in G_S$; $r$ is trivial on the image of $N_q$ for $q\in Q$; each $\alpha_v$ is symplectic with multiplier $1$ and congruent to $1$ modulo $\mathfrak m_R$; $\alpha_\ell=1$; and $r(g)\alpha_v=\alpha_v\,i(\rho_v(g))$ for $v\in T$ and $g\in D_v$.
Then there are a complete local noetherian $\mathcal O$-algebra $U$ with residue field $\mathbf F$, an $S$-framing $(i_U,r,\alpha)$ with values in $U$, and a weighted presentation of $U$, with the following properties.
\begin{enumerate}
\item Every $S$-framing with values in any $R$ is the image of $(i_U,r,\alpha)$ under a unique homomorphism of $\mathcal O$-algebras $U\to R$.
\item The variables of the presentation are those of $B$, with the same weights, each mapping to the image under $i_U$ of its image in $B$, and further variables $Y_t$, $t\in\tau$; each $Y_t$ has the weight of an index pair $(a,b)$ and maps to the entry $(a,b)$ of $r(g)-r_{0,S}(g)$ for some $g$ or of $\alpha_v-1$ for some $v$. The kernel is generated by the image of the kernel of the presentation of $B$ and by power series $f_u$, $u\in\kappa$, with $f_u$ of weight $w_f(u)$.
\item The entry $(a,b)$ of every $r(g)$ and of every $\alpha_v$ lies in the part of $U$ of the weight of $(a,b)$.
\item For every $j$, the number of $t$ with $Y_t$ of weight $j$ is $\dim_{\mathbf F}Z^1_j$, and the number of $u$ with $w_f(u)=j$ is $\dim_{\mathbf F}\mathrm{Ob}_j$.
\end{enumerate}
\end{corollary}
\leanref{Development/AutomorphicArithmetic/AutomorphicVarieties/SymplecticGSp4/NonemptyLocalFramesOfGoodData.lean}{UpperTheory.nonempty_local_frames_of_good_data}
\leanref{Development/AutomorphicArithmetic/AutomorphicVarieties/SymplecticGSp4/ResidueCharMemLevelPrimes.lean}{UpperTheory.V.GradeZeroInterpolatingElement.form}
\leanref{Development/AutomorphicArithmetic/AutomorphicVarieties/SymplecticGSp4/PhiGam.lean}{UpperTheory.V.GradeZeroInterpolatingElement.phiGam}
\leanref{Development/AutomorphicArithmetic/AutomorphicVarieties/SymplecticGSp4/LocCondQ.lean}{UpperTheory.V.GradeZeroInterpolatingElement.locCondQ}
\leanref{Development/AutomorphicArithmetic/AutomorphicVarieties/SymplecticGSp4/LocCondQ.lean}{loc_diag}
\leanref{Development/AutomorphicArithmetic/AutomorphicVarieties/PatchingModules/ChebotarevFamilies/PhiDec.lean}{phi_condition_of_finite_cohomology}
\leanref{Development/AutomorphicArithmetic/AutomorphicVarieties/CommutativeAlgebra/ExistsLocallyConstantDOneEqCocycleOfTameInertia.lean}{UpperTheory.V.Alg.exists_locallyConstant_dOne_eq_cocycle_of_tame_inertia}
\begin{proof}
This is the case of Lemma~\ref{sha-counts-vd-model-lemma-p} described here. Let $\iota=\{1,2,3,4\}$, $X=\mathbf Z^2$ and $w_a=(\lambda_a,\mu_a)$, so that $w_a-w_b$ is the weight of $(a,b)$. Let $\Phi$ be the form datum with $J$ the matrix whose nonzero entries are $J_{14}=J_{23}=1$ and $J_{32}=J_{41}=-1$, a single block, and $c=(1,1)$.
\begin{enumerate}
\item[(a)] \emph{The form.} $J^{\mathrm t}=-J$ and $J^2=-1$, and $w_a+w_b=(1,1)$ at the four nonzero entries; so $\Phi$ is a form datum. For a matrix $Y$ over a commutative ring, $Y^*=J^{-1}Y^{\mathrm t}J$ (Definition~\ref{sha-siegel-datum}); so $g^{\mathrm t}Jg=\nu J$ if and only if $g$ is symplectic with multiplier $\nu$. Hence $\mathfrak g_{\Phi,j}(\mathbf F)=\mathfrak g_j$, with the same action, and a matrix is $\Phi$-unitary if and only if it is symplectic with multiplier $1$ and congruent to $1$. A lift $r$ of $\bar r_S$ is $\Phi$-close to $r_{0,S}$ if and only if every $r(g)r_{0,S}(g)^{-1}$ is symplectic with multiplier $1$; and, since $r_0(g)$ is symplectic with multiplier $\theta(g)$, $\rho_v$ is $\Phi$-close to the restriction of $r_{0,S}$ if and only if every $\rho_v(g)$ is symplectic with multiplier $\theta(g)$, by fact (d) in the proof of Lemma~\ref{sha-counts-vd-model-descent}.
\item[(b)] \emph{Normalising.} By Lemma~\ref{sha-counts-vd-model-descent}(1), the values of $\bar r_S$ and $r_{0,S}$ are diagonal and symplectic with some multiplier; the condition on blocks holds with the trivial permutation, and the entry $(a,b)$ of a diagonal matrix vanishes when $w_a\ne w_b$, since then $a\ne b$. So $\bar r_S$ and $r_{0,S}$ normalise $\Phi$.
\item[(c)] \emph{The condition of Mazur.} $G_S$ is a profinite group. Let $H$ be an open subgroup, and $M$ the finite module of functions $G_S/H\to\mathbf Z/\ell$, with $(g\cdot m)(x)=m(g^{-1}x)$, on which $G_S$ acts continuously. Fix a representative $\tilde x$ of each coset $x$, with $\tilde x\in H$ for $x=H$. A continuous homomorphism $\varphi:H\to\mathbf Z/\ell$ gives the locally constant cocycle $c_\varphi(g)(x)=\varphi\big(\tilde x^{-1}g\,\widetilde{g^{-1}x}\big)$, and $\varphi$ is determined by it: for $u\in H$ and $g=\tilde x u\tilde x^{-1}$ with $x=H$, $c_\varphi(g)(H)=\varphi(u)$. Since $H^1(G_S,M)$ is finite \cite[Chapter~VIII, \S3]{NSW08} and $M$ is finite, there are only finitely many locally constant cocycles, and so only finitely many $\varphi$.
\item[(d)] \emph{The conditions at $Q$.} For $q\in Q$, $\bar r_S$ is trivial on the image of $N_q=I_q\cap\ker\bar r_0$. Let $L_q$ be the property of a continuous lift of the restriction of $\bar r_S$ to $D_q$ that it is trivial on $N_q$. It is a local condition with tangent group $N_q$, by Lemma~\ref{sha-counts-vd-model-lemma-p-kill} applied to $\Gamma=D_q$ and $N=N_q$: $2$ is a unit of $\mathcal O$, as $\ell$ is odd; $N_q$ is normal in $D_q$; the restrictions of $\bar r_S$ and $r_{0,S}$ normalise $\Phi$ by (b); $r_0$ is trivial on $N_q$, each $\delta_i$ being trivial there as in step (ii) of the proof of Lemma~\ref{sha-counts-vd-model-framed}; and the hypothesis on $2$-cocycles holds by Lemma~\ref{sha-counts-vd-model-lemma-p-siegel-tame}, applied to $\Gamma=D_q$, $k=\mathbf F$, $M=\mathfrak g_j$, $I=I_q$, $N=N_q$ and a Frobenius element $\phi_q$ at $v_q$. Indeed:
\begin{itemize}
\item $D_q$ is closed in $G$, so it is compact, totally disconnected and Hausdorff; and $\mathfrak g_j$ is finite, as $\mathbf F$ is.
\item $D_q$ acts on $\mathfrak g_j$ through $\bar r_0$, so the kernel of the action contains $D_q\cap\ker\bar r_0$, which is open, as $\ker\bar r_0$ is (step (c) of the proof of Lemma~\ref{sha-counts-vd-model-counts}); and $N_q$ acts trivially.
\item $I_q$ and $N_q$ are closed normal subgroups of $D_q$, and $[I_q:N_q]$ is finite and prime to $\ell$, as in step (c) of the proof of Lemma~\ref{sha-counts-vd-model-counts}.
\item By fact (e) in the proof of Lemma~\ref{sha-counts-vd-model-euler}, $I_q$ and $\phi_q$ generate a dense subgroup of $D_q$: the closure of that subgroup contains $I_q$, and its image in $D_q/I_q$ is closed and contains the image of $\phi_q$. By the same fact, for $n\ge0$ the kernel of a locally constant homomorphism $D_q\to\mathbf Z/\ell^n$ that is trivial on $I_q$ and maps $\phi_q$ to $1$ is an open subgroup of $D_q$ of index $\ell^n$ that contains $I_q$.
\end{itemize}
The restriction of $r_{0,S}$ to $D_q$ has $L_q$.
\item[(e)] \emph{Finiteness.} By Definition~\ref{sha-counts-vd-model-place-datum}, $\mathrm{Ob}_{F_S}(\mathfrak g_j)=\mathrm{Ob}_j$ is finite-dimensional, and $Z^1_{F_S}(\mathfrak g_j)=Z^1_j$.
\end{enumerate}
By (a) to (e), $\mathcal O$, $\iota$, $X$, $w$, $\Phi$, $G_S$, $\bar r_S$, $r_{0,S}$, the place datum $F_S$, the $L_q$, $B$ and the $\rho_v$ form Situation~\ref{sha-counts-vd-model-lemma-p-situation}, and they satisfy the further hypotheses of Lemma~\ref{sha-counts-vd-model-lemma-p}: $\ell$ is odd, $\mathbf F$ is finite of characteristic $\ell$, $X=\mathbf Z^2$ is finitely generated and free, $G_S$ is as in (c), and the $\mathrm{Ob}_{F_S}(\mathfrak g_j)$ are finite-dimensional by (e). By (a) the framings over $B$ are the $S$-framings: $\Phi$-closeness is the first condition, $L_q$ the second, the $\Phi$-unitarity of the $\alpha_v$ the third, and $v_0=\ell$. So Lemma~\ref{sha-counts-vd-model-lemma-p}, applied with the presentation of $B$, gives $U$, a universal $S$-framing and a presentation with the properties (1) to (4).
% Lemma -lemma-p-kill (kill_loc) is written out after its statement, with Lemma -lemma-p-unitary (hh_group, hh_small), Lemma -lemma-p-kill-section (the private ker_inf, lift_pow, hh_near and hh_section of KillLoc.lean) and Lemma -lemma-p-kill-coboundary (the private ker_coord, ker_mat, mask, cob_all and cob_ker); absorb is the last remark of its step (d). Written out in (d): locCondQ through the private loc_diag (norm_diag, the index of I_q cap ker rbar from the units of F) and dec_unr (frobDense, frobFree), and Lemma -lemma-p-siegel-tame (exists_locallyConstant_dOne_eq_cocycle_of_tame_inertia, with second_cohomology_transfer as Lemma -lemma-p-siegel-transfer and the private cocycleTrivializationOnPowers, cocycle_sequence_add and cocycle_sequence_periodic in its steps (iii) and (iv)).
\end{proof}

\subsection{Proof of the plus bound}\label{subsection-sha-plus-bound}

\begin{lemma}\label{sha-counts-vd-model-framed}
Let $D$ be a Siegel datum, with the notation of Definitions~\ref{sha-siegel-datum} and~\ref{sha-counts-vd-model-framing}, and let $B$ be a local base of $D$, with local data $(\rho^B_v)$ and origin $x_B$.
Then there is a framed ring over $B$ (Definition~\ref{sha-counts-vd-model-framing}(4)).
\end{lemma}
\leanref{Development/AutomorphicArithmetic/AutomorphicVarieties/SymplecticGSp4/NonemptyLocalFramesOfGoodData.lean}{UpperTheory.nonempty_local_frames_of_good_data}
\leanref{Development/AutomorphicArithmetic/AutomorphicVarieties/CommutativeAlgebra/FramedOrigin.lean}{UpperTheory.V.Alg.framed_origin}
\begin{proof}
Let $G_S$, $p_S$, $Q$, the $N_q$, $\bar r_S$ and $r_{0,S}$ be as in Lemma~\ref{sha-counts-vd-model-descent}, and $Z^1_j$, $\mathrm{Ob}_j$ and $\delta$ as in Definition~\ref{sha-counts-vd-model-framed-cohomology} and Lemma~\ref{sha-counts-vd-model-counts}; by that lemma, each $\mathrm{Ob}_j$ is finite-dimensional.
\begin{enumerate}
\item[(i)] \emph{The relative presentation.} The variables of the presentation of $B$ map into $\mathfrak m_B$: their images $z$ have $x_B(z)=0$, and $x_B$ is compatible with the maps to $\mathbf F$. The local data $\rho^B_v$ are continuous lifts of $\bar r_0$, symplectic with multiplier $\theta$, whose entries lie in the parts of their weights. So Corollary~\ref{sha-counts-vd-model-lemma-p-siegel}, applied to $B$ and the $\rho^B_v$, gives $U$, an $S$-framing $(i_U,r,\alpha)$ with values in $U$, and a weighted presentation of $U$, with the properties (1) to (4) there.
\item[(ii)] \emph{The origin.} The triple $(x_B,r_{0,S},(1)_v)$ is an $S$-framing with values in $\mathcal O$. Indeed $r_{0,S}(g)r_{0,S}(g)^{-1}=1$. Next, $r_0(g)=1$ for every $g$ with $\bar r_0(g)=1$ in an inertia group at a place not above $T$, each $\delta_i(g)$ being then $1$, as in step (iii) of the proof of Lemma~\ref{sha-counts-vd-model-cocycle}; so $r_{0,S}$ is trivial on the images of the $N_q$, by Lemma~\ref{sha-counts-vd-model-descent}(2). Finally $r_{0,S}(g)=x_B(\rho^B_v(g))$ for $g\in D_v$, since $x_B$ carries $\rho^B_v$ to $r_0$. By (1) there is $x_U:U\to\mathcal O$ that carries $(i_U,r,\alpha)$ to $(x_B,r_{0,S},(1)_v)$. Then $x_U$ kills the images of the variables: those of the variables of $B$ are elements $i_U(z)$ with $x_B(z)=0$, and the others are entries of some $r(g)-r_{0,S}(g)$ or $\alpha_v-1$. So $x_U$ composed with the presentation and the constant term agree on the variables, which they map to $0$, and they are equal, by step (ii) of the proof of Lemma~\ref{sha-counts-vd-model-coproduct} with $R=\mathcal O$. Hence the series in the kernel of the presentation have constant term $0$, the presentation has the origin $x_U$, $x_U\circ i_U=x_B$, and $x_U$ carries $r$ to $r_{0,S}$ and each $\alpha_v$ to $1$.
\item[(iii)] \emph{The framing.} Put $r_U=r\circ p_S$. By Lemma~\ref{sha-counts-vd-model-descent}(2), each $r_U(g)$ is symplectic with multiplier $\theta(g)$, and $r_U(g)=1$ for every $g$ with $\bar r_0(g)=1$ in an inertia group at a place not above $T$. The other conditions of Definition~\ref{sha-counts-vd-model-framing}(3) are those of an $S$-framing, since $D_v\to G_S$ is the restriction of $p_S$. So $(i_U,r_U,\alpha)$ is a framing. Let $(i',r',\alpha')$ be a framing with values in $R$. By Lemma~\ref{sha-counts-vd-model-descent}(3), $r'=r''\circ p_S$ for a continuous lift $r''$ of $\bar r_S$, and by Lemma~\ref{sha-counts-vd-model-descent}(2), $(i',r'',\alpha')$ is an $S$-framing. As $p_S$ is surjective, a homomorphism $f:U\to R$ carries $r_U$ to $r'$ if and only if it carries $r$ to $r''$. So by (1) there is a unique $f$ that carries $(i_U,r_U,\alpha)$ to $(i',r',\alpha')$, and $(i_U,r_U,\alpha)$ is universal.
\item[(iv)] \emph{The presentation.} By (2) of Corollary~\ref{sha-counts-vd-model-lemma-p-siegel}, the presentation of $U$ has the variables of $B$, with their weights and mapping to their images under $i_U$, and variables $Y_t$ of the weight of an index pair $(a,b)$, mapping to the entry $(a,b)$ of $\alpha_v-1$ for some $v$ or of $r(p_S(g))-r_{0,S}(p_S(g))=r_U(g)-r_0(g)$ for some $g\in G$; its kernel is generated by the image of the kernel of the presentation of $B$ and by the series $f_u$ of weights $w_f(u)$. By (3) there, the entries of the $r_U(g)$ and of the $\alpha_v$ lie in the parts of their weights; and by (ii), $x_U$ carries $r_U$ to $r_0$.
\item[(v)] \emph{The counts.} By (4) of Corollary~\ref{sha-counts-vd-model-lemma-p-siegel} and Lemma~\ref{sha-counts-vd-model-counts}, the number of $u$ with $w_f(u)=j$ is $\dim\mathrm{Ob}_j=\dim Z^1_j+\delta(j)$, that is, the number of $t$ with $Y_t$ of weight $j$, plus $2$ if $j=(0,0)$, plus $1$ if $j=\pm(0,1)$, and plus $0$ otherwise. No $Y_t$ has a weight outside $\{(0,0),\pm(0,1)\}$, since $Z^1_j=0$ for those weights.
\end{enumerate}
So $(U,i_U,r_U,\alpha)$, with this presentation, is a framed ring over $B$.
\end{proof}

\begin{proposition}\label{sha-counts-vd-model-exists}
Every Siegel datum has a graded model (Definition~\ref{sha-counts-vd-model}).
\end{proposition}
\leanref{Development/AutomorphicArithmetic/AutomorphicVarieties/SymplecticGSp4/ExModel.lean}{UpperTheory.V.GradeZeroInterpolatingElement.exModel}
\leanref{Development/AutomorphicArithmetic/AutomorphicVarieties/SymplecticGSp4/ExModel.lean}{UpperTheory.V.GradeZeroInterpolatingElement.ExM.tower}
\begin{proof}
Let $D$ be a Siegel datum, with the notation of Definitions~\ref{sha-siegel-datum} and~\ref{sha-counts-vd-model-framing}, and let $\Pi$ be the list of the six index pairs $(1,2)$, $(1,4)$, $(3,2)$, $(2,1)$, $(4,1)$, $(2,3)$, of weights $(1,0)$, $(1,1)$, $(1,-1)$, $(-1,0)$, $(-1,-1)$, $(-1,1)$; for these pairs $\lambda_a\ne\lambda_b$.

\emph{Step 1: the local base and the framed ring.} For $v\in T$ let $B_v$ and $\rho_v$ be the universal ring and lift of Lemma~\ref{sha-counts-vd-model-universal}(1). Lemma~\ref{sha-counts-vd-model-local-charts} gives homomorphisms $x_v$ and $\varphi_v$ as in Lemma~\ref{sha-counts-vd-model-universal}(2), with $d(j)$ variables of weight $j$ if $v=\ell$ and $d'(j)$ if $v\in S_b$, so Lemma~\ref{sha-counts-vd-model-base} gives a local base $B$ of $D$, with local data $(\rho^B_v)$ and origin $x_B$. Lemma~\ref{sha-counts-vd-model-framed} gives a framed ring $(U,i_U,r_U,\alpha)$ over $B$, with origin $x_U$, variables $Y_t$ and series $f_u$ of weights $w_f(u)$; let $a(j)$ be the number of $t$ with $Y_t$ of weight $j$.
% Step 1 is written out down to the lemmas it uses. Their proofs are written out too, except for the narrower open steps in the proofs of Lemmas sha-counts-vd-model-lift-chart (five sub-lemmas: -lift-chart-factor, -continuity, -generation, -smooth, -graded), -local-lifts (-local-lifts-flag, -symplectic, -unramified, -borel, -commuting), -local-tangent (-local-tangent-conditions, -cocycles, -unramified) and -lemma-p (-lemma-p-grading, -tangent, -relations, -classes), and of Corollary -lemma-p-siegel (-lemma-p-kill). Lemma -counts is closed. Written out earlier: spChartL and spChartP (Lemma -local-charts), locBase (Lemmas -coproduct, -universal, -base), descend (Lemma -descent), and nonempty_local_frames_of_good_data with framed_origin, hh_iff and hh_sp (Lemma -framed).

\emph{Step 2: a cover.} By Lemma~\ref{sha-counts-vd-model-chart}, $B$ has a chart at $x_B$ with $n_B(j)$ variables of weight $j$ (Definition~\ref{sha-counts-vd-model-charted}). Lemma~\ref{sha-counts-vd-model-cover}, applied to it and to the presentation of $U$ relative to that of $B$ in Definition~\ref{sha-counts-vd-model-framing}(4), gives a complete local noetherian $\mathcal O$-algebra $P$ with residue field $\mathbf F$, with a weighted presentation with origin $x_P$ whose variables have weights among those of $U$, homomorphisms of $\mathcal O$-algebras $i_P:B\to P$ and $p:P\to U$, and elements $g_u\in P_{w_f(u)}$, $u\in\kappa$, such that: $p$ is surjective with kernel generated by the $g_u$, $p\circ i_P=i_U$, $x_U\circ p=x_P$, and $p$ maps each part of $P$ onto the part of $U$ of the same weight; $P$ has a chart at $x_P$, with the properties of Lemma~\ref{sha-counts-vd-model-chart}, with $n_B(j)+a(j)$ variables of weight $j$; and derivations lift: for a nonzero $a\in\mathcal O$, a weight $j$, a derivation $D_U:U\to\mathcal O/a$ at $x_U$ and a derivation $D_B:B\to\mathcal O$ at $x_B$ of weight $j$ with $D_U\circ i_U=D_B\bmod a$ and $D_U(z_t)=0$ for the images $z_t$ of the $Y_t$ of weight other than $j$, there is a derivation $\delta_P:P\to\mathcal O$ at $x_P$ of weight $j$ with $\delta_P\bmod a=D_U\circ p$.

\emph{Step 3: the equations.} For $v\in T$ and $k\in\Pi$ let $t_{v,k}\in U$ be the entry $k$ of $r_U(\phi)$ if $v=\ell$ and of $\alpha_v$ if $v\ne\ell$. For $p\in S_b$ and $k$ among $(1,3)$, $(3,1)$, $(1,1)$, $(2,2)$ let $s_{p,k}$ be the entry $k$ of $\alpha_p-1$, and let $s_1,s_2$ be the entries $(1,3)$ and $(3,1)$ of $H_U=r_U(g_0)+\theta(g_0)r_U(g_0^{-1})$. Each of these elements lies in the part of $U$ of the weight of its index pair, and in $\ker x_U$, since $x_U$ carries $r_U$ to the diagonal $r_0$ and $\alpha_v$ to $1$.
For each $(v,k)$ there is a derivation of $U$ at $x_U$ with values in $\mathcal O$ that is nonzero at $t_{v,k}$ and vanishes at the other $t_{v',k'}$. Let $X$ be $E_{12}-E_{34}$, $E_{14}$, $E_{32}$, $E_{21}-E_{43}$, $E_{41}$ or $E_{23}$ according as $k$ is the first, \dots, sixth pair of $\Pi$; it lies in $\mathfrak{sp}_4(\mathcal O)$, and its entries at the pairs of $\Pi$ are $1$ at $k$ and $0$ elsewhere. If $v=\ell$, Lemma~\ref{sha-counts-vd-model-tangent} with $Y=X$ and $A=0$ gives $\delta$ with $\delta(r_U(\phi)_{ab})=X_{ab}\big(\delta_b(\phi)-\delta_a(\phi)\big)$, which for $(a,b)\in\Pi$ vanishes unless $(a,b)=k$, and is a unit for $(a,b)=k$ by Lemma~\ref{sha-flagged-residues}(1) and~(2); and $\delta(\alpha_{v'})=0$. If $v\ne\ell$, the lemma with $Y=0$, $A_v=X$ and $A_{v'}=0$ for $v'\ne v$ gives $\delta$ with $\delta(r_U(g))=0$ and $\delta((\alpha_{v'})_{k'})$ equal to $1$ if $(v',k')=(v,k)$ and $0$ otherwise.
Choose $t'_{v,k}\in P$ in the part of the weight of $k$ with $p(t'_{v,k})=t_{v,k}$, and in the same way $s'_{p,k}$, $s'_1$, $s'_2$; they lie in $\ker x_P$. Composing with $p$, the derivations above separate the $t'_{v,k}$.

\emph{Step 4: the quotient.} Lemma~\ref{sha-counts-vd-model-quotient}, applied to $P$, its chart and the $t'_{v,k}$, which Step 3 separates by derivations, gives a complete local noetherian $\mathcal O$-algebra $R$ with residue field $\mathbf F$ that is a domain, a surjective homomorphism of $\mathcal O$-algebras $q:P\to R$ with $q(t'_{v,k})=0$, and a weighted presentation of $R$ with origin $x_R$, with $x_R\circ q=x_P$, variables with weights among those of $P$ and parts $q(P_j)$, such that every derivation of $P$ at $x_P$ with values in $\mathcal O$ that vanishes at the $t'_{v,k}$ vanishes on $\ker q$, and $R^{\mathrm{gr}}$ has a chart at $x_R$, with the properties of Definition~\ref{sha-counts-vd-model}(i) other than the numbers of variables, with $n'(j)$ variables of weight $j$, where $n'(j)+(|S_b|+1)\,\#\{k\in\Pi:k\text{ has weight }j\}=n_B(j)+a(j)$.
% Step 2 is nonempty_cover_of_relative_presentation_smoothPt (Lemma sha-counts-vd-model-cover). Step 4 is exists_smooth_quotient_by_weighted_equations (Lemma sha-counts-vd-model-quotient), through Alg.chart_cut (ChartCut.lean, 541 lines; Lemmas sha-counts-vd-model-solve and -cut), Alg.chart_ker (-kernel) and Alg.chart_gr (-graded).
Let $J$ be the ideal of $R$ generated by the $q(g_u)$, the $q(s'_{p,k})$, $q(s'_1)$ and $q(s'_2)$, let $S=R/J$, and let $\pi:R\to S$ be the quotient map. Since $p$ is surjective with kernel generated by the $g_u$, which $\pi\circ q$ kills, there is a surjective homomorphism $\sigma:U\to S$ with $\sigma\circ p=\pi\circ q$. Then $\sigma$ kills the $t_{v,k}$, $s_{p,k}$, $s_1$ and $s_2$.
\begin{enumerate}
\item[(a)] For $y\in U_j$ there is $f\in R_j$ with $\pi(f)=\sigma(y)$: lift $y$ to $P_j$ and apply $q$.
\item[(b)] If $\pi(f)=\sigma(y)$, then $x_R(f)=x_U(y)$: write $f=q(z)$ and $y=p(z')$; then $q(z)-q(z')\in J$, and the generators of $J$ lie in $\ker x_R$, so $x_R(f)=x_R(q(z'))=x_P(z')=x_U(y)$.
\item[(c)] Let $a\in\mathcal O$ be nonzero, $j$ a weight that is not the weight of a pair of $\Pi$, and $D_U$, $D_B$ as in Step 2, with $D_U$ vanishing moreover at the $s_{p,k}$, $s_1$ and $s_2$. Then there is a derivation $\delta:R\to\mathcal O$ at $x_R$ of weight $j$ with $\delta(f)\equiv D_U(y)$ modulo $a$ whenever $\pi(f)=\sigma(y)$. Indeed, Step 2 gives $\delta_P$ of weight $j$, which vanishes at the $t'_{v,k}$, as they lie in parts of weights other than $j$; by Step 4 it vanishes on $\ker q$, so $\delta_P=\delta\circ q$ for a derivation $\delta$ of $R$ at $x_R$, of weight $j$, since the parts of $R$ are the $q(P_j)$. Modulo $a$, $\delta$ vanishes at the generators of $J$, at which $D_U$ vanishes, and so on $J$, as $x_R$ kills $J$. If $\pi(f)=\sigma(y)$, with $f=q(z)$ and $y=p(z')$ as in (b), then $\delta(f)\equiv\delta(q(z'))=\delta_P(z')\equiv D_U(y)$.
\end{enumerate}

\emph{Step 5: the model.} The generators of $J$ have weights $(0,-1)$, $(0,1)$ or $(0,0)$: the $q(g_u)$ by Definition~\ref{sha-counts-vd-model-framing}(4), and $q(s'_{p,(1,3)})$, $q(s'_{p,(3,1)})$, $q(s'_{p,(1,1)})$, $q(s'_{p,(2,2)})$, $q(s'_1)$, $q(s'_2)$ have the weights $(0,1)$, $(0,-1)$, $(0,0)$, $(0,0)$, $(0,1)$, $(0,-1)$. Put $r=1+|S_b|+a(0,-1)$, $r'=2+|S_b|+a(0,1)$ and $m=2+2|S_b|+a(0,0)$. By the count in Definition~\ref{sha-counts-vd-model-framing}(4) there are $r+1$ generators of weight $(0,-1)$, $r'$ of weight $(0,1)$ and $m$ of weight $(0,0)$; call them $F_0,\dots,F_r$, $E_1,\dots,E_{r'}$ and $Z_1,\dots,Z_m$. Each lies in $\ker x_R$, and $r\ge1$.
The weights of the variables of $R$ lie in $\mathcal W$, being among those of $B$ and of the $Y_t$. By the formula of Step 4 and Definition~\ref{sha-counts-vd-model-framing}(2) and~(4), where $a(j)=0$ for $j\notin\{(0,0),\pm(0,1)\}$, the chart of $R^{\mathrm{gr}}$ has one variable of each of the weights $(1,0)$, $(1,1)$, $(-1,1)$, for which $n_B(j)=2+|S_b|$, none of the weights $(-1,0)$, $(-1,-1)$, $(1,-1)$, for which $n_B(j)=1+|S_b|$, and $r'$, $r$ and $m$ of the weights $(0,1)$, $(0,-1)$ and $(0,0)$. These are the counts of Definition~\ref{sha-counts-vd-model}(i).
Let $r^{\mathrm u}=\sigma\circ r_U$.
\begin{enumerate}
\item[(i)] \emph{$\sigma(\alpha_v)=1$ for $v\in T$.} For $v=\ell$, $\alpha_\ell=1$. For $p\in S_b$, $\sigma(\alpha_p)$ is symplectic with multiplier $1$, congruent to $1$, its entries at the pairs of $\Pi$ and at $(1,3)$ and $(3,1)$ vanish and its entries $(1,1)$ and $(2,2)$ are $1$, as $\sigma$ kills the $t_{p,k}$ and $s_{p,k}$; by Lemma~\ref{sha-counts-vd-model-six}, $\sigma(\alpha_p)=1$.
\item[(ii)] \emph{$r^{\mathrm u}$ is a Siegel point.} It is a continuous lift of $\bar r_0$, symplectic with multiplier $\theta$, and trivial on the elements $g$ with $\bar r_0(g)=1$ of the inertia groups at places not above $T$, since $r_U$ is. The matrix $\sigma(r_U(\phi))$ is symplectic with multiplier $\theta(\phi)$, a unit, its entries off the diagonal lie in $\mathfrak m_S$, since it reduces to the diagonal $\bar r_0(\phi)$, and its entries at the pairs of $\Pi$ vanish; by Lemma~\ref{sha-counts-vd-model-six}, $r^{\mathrm u}(\phi)$ is block diagonal. The entries $(1,3)$ and $(3,1)$ of $r^{\mathrm u}(g_0)+\theta(g_0)r^{\mathrm u}(g_0^{-1})$ are $\sigma(s_1)=\sigma(s_2)=0$. For $v\in T$ and $g\in D_v$, (i) and $r_U(g)\alpha_v=\alpha_vi_U(\rho^B_v(g))$ give $r^{\mathrm u}(g)=\sigma(i_U(\rho^B_v(g)))=\pi(\rho^R_v(g))$, with $\rho^R_v=q\circ i_P\circ\rho^B_v$. The $\rho^R_v$ are local data with values in the domain $R$, and $x_R\circ q\circ i_P=x_P\circ i_P=x_U\circ i_U=x_B$ carries them to $r_0$. So $r^{\mathrm u}$ satisfies the conditions at $\ell$ and at the $p\in S_b$ of Definition~\ref{sha-siegel-datum}(2).
\item[(iii)] \emph{Generation.} Let $y,y':S\to R'$ agree on the entries of the $r^{\mathrm u}(g)$. By (ii), the local data $y\sigma i_U(\rho^B_v)$ and $y'\sigma i_U(\rho^B_v)$ are the restrictions of $y\circ r^{\mathrm u}=y'\circ r^{\mathrm u}$, so $y\sigma i_U=y'\sigma i_U$ by the universality of $B$. Then the framings $(y\sigma i_U,y\circ r^{\mathrm u},(1)_v)$ and $(y'\sigma i_U,y'\circ r^{\mathrm u},(1)_v)$, images of $(i_U,r_U,\alpha)$ under $y\sigma$ and $y'\sigma$ by (i), are equal, so $y\sigma=y'\sigma$ by the universality of $U$, and $y=y'$ as $\sigma$ is surjective. This is condition (ii) of Definition~\ref{sha-counts-vd-model}.
\item[(iv)] \emph{Entries.} The entry $(i,j)$ of $r_U(g)$ lies in $U_w$, $w$ the weight of $(i,j)$, so by (a) $r^{\mathrm u}(g)_{ij}=\pi(f)$ for some $f\in R_w$; and by (b), $x_R(f)=x_U(r_U(g)_{ij})=r_0(g)_{ij}$ whenever $\pi(f)=r^{\mathrm u}(g)_{ij}$. This is condition (iii).
\item[(v)] \emph{Derivations.} Let $a$ and $z$ be as in condition (iv) of Definition~\ref{sha-counts-vd-model}, and let $X_m$, $D_U$ and $D_B$ be as in Lemma~\ref{sha-counts-vd-model-cocycle}. The matrix $X_mr_0(g)$ is supported on the pairs $(3,1)$ and $(4,2)$, of weight $(0,-1)$; so $D_U$ vanishes at the images of the $Y_t$ of other weights, which are entries of the $r_U(g)-r_0(g)$ at pairs of those weights or entries of the $\alpha_v-1$. It vanishes at the $s_{p,k}$, and at $s_1$ and $s_2$, since $z(g_0)=0$ and hence $z(g_0^{-1})=0$ by the cocycle relation. The weight $(0,-1)$ is not that of a pair of $\Pi$. So (c) gives $\delta:R\to\mathcal O$ at $x_R$ of weight $(0,-1)$ with $\delta(f)\equiv D_U(y)$ modulo $a$ whenever $\pi(f)=\sigma(y)$. It maps $\ker\pi$ into $a\mathcal O$, taking $y=0$; and if $\pi(f)$ is $r^{\mathrm u}(g)_{31}$ or $r^{\mathrm u}(g)_{42}$, then $\delta(f)$ is congruent to $(X_mr_0(g))_{31}z(g)=\theta(g)z(g)$ or to $(X_mr_0(g))_{42}z(g)=\Psi(g)z(g)$. This is condition (iv).
\end{enumerate}
So $R$ with its presentation, the integers $r,r',m$, the families $F$, $E$ and $Z$, the ring $S$ with the quotient map $\pi$ in the role of $q$, and the Siegel point $r^{\mathrm u}$ form a graded model of $D$.
\end{proof}

\begin{proposition}\label{sha-counts-vd-bound}
Let $D$ be a Siegel datum (Definition~\ref{sha-siegel-datum}) that is finite, has no endoscopic $\bar{\mathbf Q}_\ell$-points, is rigid at each of its irreducible $\bar{\mathbf Q}_\ell$-points (Definition~\ref{sha-counts-vd-model}), and has a graded model. Let $k$ and $\psi$ be its integer and its character, $n_S$ the product of the primes of $S$, and $V_D=V_{\ell,n_S}(k,\psi)$ (Definition~\ref{sha-character-parts}). Let $\sigma\in G$ be such that $\operatorname{tr}r(\sigma)\ne\operatorname{tr}r_0(\sigma)$ for every irreducible $\bar{\mathbf Q}_\ell$-point $r$ of $D$, and let $|\cdot|_\ell$ be the absolute value of $\bar{\mathbf Q}_\ell$ with $|\ell|_\ell=\ell^{-1}$.
Then $V_D$ is finite, and there is a finite set $X$ of irreducible $\bar{\mathbf Q}_\ell$-points of $D$, no two of which are conjugate under $GL_4(\bar{\mathbf Q}_\ell)$, with
\[
 \frac{V_D}{e_{\mathcal O}}\log\ell\le-2\sum_{r\in X}\log\big|\operatorname{tr}r(\sigma)-\operatorname{tr}r_0(\sigma)\big|_\ell.
\]
\end{proposition}
\leanref{Development/AutomorphicArithmetic/AutomorphicVarieties/SymplecticGSp4/GaloisRepIneqOfAllIrreducibleNonvanishing.lean}{UpperTheory.V.GradeZeroInterpolatingElement.galoisRep_ineq_of_allIrreducible_nonvanishing}
\leanref{Development/AutomorphicArithmetic/AutomorphicVarieties/SymplecticGSp4/DecompSigmaOrd.lean}{UpperTheory.V.GradeZeroInterpolatingElement.Raw.Ineq}
\leanref{Development/AutomorphicArithmetic/AutomorphicVarieties/SymplecticGSp4/Nval.lean}{UpperTheory.V.GradeZeroInterpolatingElement.nval}
\begin{proof}
Use the notation of Definition~\ref{sha-counts-vd-model} for a graded model of $D$, and let $Y$ be the set of homomorphisms of $\mathcal O$-algebras $A\to\bar{\mathbf Q}_\ell$ that vanish on $I$.

\emph{Step 1: the element $a$.} Lemma~\ref{sha-counts-vd-structure}(d) gives $a\in A_0$ with $x(a)=0$ and $\xi(q_A(a))=\operatorname{tr}r_\xi(\sigma)-\operatorname{tr}r_0(\sigma)$ for every $\xi:S^{\mathrm{gr}}\to\bar{\mathbf Q}_\ell$.

\emph{Step 2: the point of a two-signed element.} Let $x_1\in Y$ be two-signed. By Lemma~\ref{sha-counts-vd-points}(a) and (d), $x_1=\xi\circ q_A$ for some $\xi$, and $\rho(x_1)=\tau(t)r_\xi\tau(t)^{-1}$ is a $\bar{\mathbf Q}_\ell$-point of $D$ for some $t$; fix such choices. By Lemma~\ref{sha-counts-vd-irreducible}(d), $r_\xi$ is irreducible, and so is its conjugate $\rho(x_1)$. Since $\operatorname{tr}\rho(x_1)(\sigma)=\operatorname{tr}r_\xi(\sigma)$, Step 1 gives $x_1(a)=\operatorname{tr}\rho(x_1)(\sigma)-\operatorname{tr}r_0(\sigma)$, which is nonzero by hypothesis. So condition (5) of Lemma~\ref{sha-algebra-cotangent} holds for $a$.

\emph{Step 3: conjugate points.} Let $x_1,x'_1\in Y$ be two-signed, with $\xi,t$ and $\xi',t'$ as in Step 2, and suppose that $\rho(x_1)$ and $\rho(x'_1)$ are conjugate under $GL_4(\bar{\mathbf Q}_\ell)$. By Lemma~\ref{sha-counts-vd-conjugate}(b), $\rho(x'_1)=\tau(s)\rho(x_1)\tau(s)^{-1}$ for some $s$, so $r_{\xi'}=\tau(u)r_\xi\tau(u)^{-1}$ with $u=t'^{-1}st$. By Lemma~\ref{sha-counts-vd-points}(b), $r_{u\cdot\xi}=r_{\xi'}$, so $u\cdot\xi=\xi'$ by Lemma~\ref{sha-counts-vd-points}(c). As $q_A(A_n)\subseteq S_{(0,n)}$, this gives $x'_1=u_2^nx_1$ on $A_n$; so $x'_1$ is in the orbit of $x_1$ (Definition~\ref{sha-algebra-graded}).

\emph{Step 4: the commutative algebra.} Let $v(y)=-e_{\mathcal O}\log|y|_\ell/\log\ell$ for $y\ne0$. It satisfies $v(yy')=v(y)+v(y')$ and $v(y+y')\ge\min(v(y),v(y'))$, and for nonzero $c\in\mathcal O$, $v(c)$ is the normalized valuation of $c$, since a uniformizer of $\mathcal O$ has absolute value $\ell^{-1/e_{\mathcal O}}$. By Lemma~\ref{sha-counts-vd-conditions} and Step 2, Lemma~\ref{sha-algebra-cotangent} applies with $\Omega=\bar{\mathbf Q}_\ell$ and this $v$. It shows that $\lambda_{-1}(F)$ is finite, that there is a finite set of two-signed elements of $Y$ meeting the orbit of every two-signed element of $Y$, and that $\lambda_{-1}(F)\le2\sum_{x_1\in T}v(x_1(a))$ for every such set $T$.

\emph{Step 5: the set $X$.} Choose such a set $T$ with the least number of elements. No two elements of $T$ lie in the same orbit, since the orbits partition the two-signed elements of $Y$ and removing one of two elements in the same orbit would leave a smaller such set. Let $X=\{\rho(x_1):x_1\in T\}$. By Step 3, the points $\rho(x_1)$ for distinct $x_1\in T$ are not conjugate; in particular $x_1\mapsto\rho(x_1)$ is injective on $T$, and by Step 2,
\[
 \sum_{x_1\in T}v(x_1(a))=\frac{e_{\mathcal O}}{\log\ell}\sum_{r\in X}-\log\big|\operatorname{tr}r(\sigma)-\operatorname{tr}r_0(\sigma)\big|_\ell.
\]

\emph{Step 6: the bound.} By Lemma~\ref{sha-counts-vd-sha}, $V_D$ is finite and $V_D\le\lambda_{-1}(F)$. With Steps 4 and 5 this gives $V_D\le\frac{2e_{\mathcal O}}{\log\ell}\sum_{r\in X}-\log|\operatorname{tr}r(\sigma)-\operatorname{tr}r_0(\sigma)|_\ell$, which is the claim.
\end{proof}

\begin{proposition}\label{sha-counts}
Let $\psi$ be a primitive Dirichlet character, and in (a), (b) and (d) assume the Galois representations as in Theorem~\ref{sha-residually-reducible-automorphy}.
Below, $h$ runs over host data for $\psi$, $t$ over integers for which $h$ is large, $\ell$ over odd primes, $k$ over integers with $0<k<t$, $\mathcal O$ over rings as in Definition~\ref{sha-character-parts}, and $\psi'$ over $X_h(k)$; all the groups $V'_\ell(k,\psi')$ involved have finite length.
\begin{enumerate}
\item (Case I.) There are $C$ and primes $r_1,r_2$ such that if $r_1,r_2<q$ and $r_1,r_2<d$, then for every finite set $\Lambda$ of odd primes with $(\ell,k)$ in case I for all $\ell\in\Lambda$,
\[
 \sum_{\ell\in\Lambda}\log\ell\sum_{\psi'\in X_h(k)}\bar V_\ell(k,\psi')\le(N_ht)^{C}.
\]
\item (Cases II and III.) There is $C$ such that $\bar V_\ell(k,\psi')\log\ell\le(\ell N_hk)^{C}$ whenever $(\ell,k,\psi')$ is in case II or III.
\item (Case Z.) $\bar V_\ell(k,\psi')=0$ whenever $(\ell,k,\psi')$ is in case Z.
\item (Case V, small $2$-conductor.) There is $C$ such that $\bar V_\ell(k,\psi')\log\ell\le(\ell N_hk)^{C}$ whenever $(\ell,k,\psi')$ is in case V and $16\,o_\psi(\ell N_\psi)^2\,4^{n_0}\le2^{c_2}$. Here $2^{c_2}$ is the conductor of the $2$-component of $\psi'$, and $n_0=v_2(\ell^{f_0}-1)$ for the order $f_0$ of $\ell$ modulo $\operatorname{lcm}(4,\operatorname{ord}\theta_1)$, where $\theta_1$ is the part of $\psi'$ of order prime to $\ell$ with its $2$-component removed.
\item (Case V, large $2$-conductor.) The bound of (d) holds, without assuming the Galois representations, when $(\ell,k,\psi')$ is in case V and $2^{c_2}<16\,o_\psi(\ell N_\psi)^2\,4^{n_0}$.
\item (Case IV.) For every $\ell\in B_\psi$ and $j\ge0$ there is $C$ such that the following holds when $t\ge2$. Let $B$ be real and suppose that for every field $L_0$ of the cyclotomic datum of $\ell N_\psi$ and every layer $L_1$ of degree $\ell^{j}$ over $L_0$, the field $L_1$ satisfies the $\ell$-bound $(t,B)$. Then for every $\psi'$, trivial on $H_h$, for which $\chi'_\ell(\psi')$ has odd order, $\bar V_\ell(k,\psi')$ is finite for $0<k<t$ and $\sum_{k=1}^{t-1}\bar V_\ell(k,\psi')\le B+(N_ht)^{C}$.
\end{enumerate}
\end{proposition}
\begin{proof}
(c) is Theorem~\ref{sha-trivial-character}: in case Z, $\ell\ge5$ because $3\in B_\psi$, and $k$ is even because $\psi'=1$ lies in $X_h(k)$; the group defining $V'_\ell(k,1)$ is contained in the one defining $V_\ell(k,1)$.

(a) Let $p_K=p_K(\psi)$ and $K=\mathbf Q(\sqrt{-p_K})$, and take primes $r_2>r_1>30N_\psi p_K$.
The proof combines two propositions about $K$, on the automorphic triples of Definition~\ref{sha-eigensystem-triple} through an isomorphism $\iota:\bar{\mathbf Q}_\ell\cong\mathbf C$; $v_\ell$ is the valuation with $v_\ell(\ell)=1$.
\begin{itemize}
\item \emph{Counting.} There are $c_5,A$ such that for all $k,N\ge1$, every prime $p\ge7$ with $p\nmid Nd_K$, every prime $w$ of $K$ above $p$, of norm $q$, with arithmetic Frobenius element $\sigma_w$, and every finite family $(z_o)$ of roots of unity stable under $\operatorname{Aut}(\mathbf C)$, there is a nonzero integer $N_0$ with $\log|N_0|\le c_5\,\#\{o\}\,(N(k+1))^{A}\log p$ such that the following holds. Let $t_o=q^{-k}z_o+1+q^kz_o^{-1}$, the trace at $\sigma_w^{-1}$ of the Eisenstein representation $\operatorname{diag}(\Psi,1,\Psi^{-1})$ when $\Psi(\sigma_w)=q^kz_o^{-1}$. For every odd prime $\ell\ne p$, every $\iota$, and all finite sets $X_o$ of pairwise non-isomorphic automorphic triples of weight $k$ and level $N$, $\iota\operatorname{tr}\rho(\sigma_w^{-1})\ne t_o$ for $\rho\in X_o$, and $|N_0|_\ell\le\prod_o\prod_{\rho\in X_o}|\operatorname{tr}\rho(\sigma_w^{-1})-\iota^{-1}t_o|_\ell$.
\item \emph{One character.} Let $(\ell,k,\psi',K)$ be an Eisenstein datum (Definition~\ref{sha-ordinary-deformation}), $\sigma\in\operatorname{Gal}(\bar{\mathbf Q}/K)$ and $t=\Psi(\sigma)+1+\Psi(\sigma)^{-1}$. If no automorphic triple of weight $k$ and level $N$ has trace $t$ at $\sigma$, then $V'_\ell(k,\psi')$ is finite and there is a set $X$ of pairwise non-isomorphic automorphic triples with $\bar V_\ell(k,\psi')\le2\sum_{\rho\in X}v_\ell\big(\operatorname{tr}\rho(\sigma)-t\big)$.
\end{itemize}
\emph{Proof of the counting proposition.}
Let $c_1,A_1$ be the constants of Lemma~\ref{sha-eigensystem-traces} for $K$, and put $c_5=5c_1$ and $A=A_1+1$.
The norm $q$ of $w$ is $p$ or $p^2$, so $q\ge7$ and $\log q\le2\log p$.
Lemma~\ref{sha-eigensystem-traces} gives a trace datum $(t_s)_{s\in S}$ of weight $k$ at $q$, with $|S|\le c_1(N(k+1))^{A_1}$, that classifies the automorphic triples of weight $k$ and level $N$ at $\sigma_w^{-1}$ for every odd $\ell\ne p$ and every $\iota$.
Lemma~\ref{sha-eigensystem-norm}, applied to this datum and to $(z_o)$, gives a nonzero integer $N_0$ with $\log|N_0|\le\#\{o\}\,|S|\,(2k\log q+1)$, and both conclusions for every odd $\ell\ne p$, since such an $\ell$ does not divide $q$.
Since $p\ge7$, $\log p\ge1$, so $2k\log q+1\le4k\log p+1\le5N(k+1)\log p$, and $\log|N_0|\le c_5\,\#\{o\}\,(N(k+1))^{A}\log p$.

\emph{Proof of the one-character proposition.}
The character $\psi'$ has order prime to $\ell$, so its values lie in an unramified extension of $\mathbf Z_\ell$, and $V'_\ell(k,\psi')/e_{\mathcal O}$ does not depend on the choice of $\mathcal O$; so we may assume that $\mathcal O$ is unramified over $\mathbf Z_\ell$, and then $\bar V_\ell(k,\psi')=V'_\ell(k,\psi')$.
Theorem~\ref{sha-residually-reducible-automorphy} gives $g_1$ and the finiteness, stable-point and endoscopic-point statements for the ordinary points relative to $g_1$ and a place $B$ above $\ell$.
Let $R$ be the universal ring of the $\mathfrak a_B$-lifts relative to $g_1$ and $B$, with universal lift $r^{\mathrm u}$ (Lemma~\ref{sha-model-universal}), and let $x_0:R\to\mathcal O$ be the homomorphism with $x_0\circ r^{\mathrm u}=r_0(\eta,\Psi)$.
Points $f$ of $R$, their representations $r_f$ and two-signed points are as in Lemma~\ref{sha-model-points}; let $\rho_f$ be the restriction of $r_f$ to $\operatorname{Gal}(\bar{\mathbf Q}/K)$.
By the finiteness statement and Lemma~\ref{sha-model-points}(a), the algebra $R^\natural_0$ is finite over $\mathcal O$.

\emph{The model.} Let $A$, $x$, the families $F$, $G$, $E$, $E'$, $Z$, the ideal $I$ and $q_A:A\to R$ be as in Proposition~\ref{sha-model-exists}, and let $Y$ be the set of homomorphisms of $\mathcal O$-algebras $A\to\bar{\mathbf Q}_\ell$ that vanish on $I$.
Points of $R$ and elements of $Y$ are related as follows.
\begin{enumerate}
\item[(Y1)] For a point $f$, $f\circ q_A$ vanishes on every $A_j$ with $j<0$ if and only if the matrix parts of $r_f$ on $\ker\eta$ are upper triangular, and on every $A_j$ with $j>0$ if and only if they are lower triangular. So $f\circ q_A$ is two-signed if and only if $f$ is.
\item[(Y2)] For points $f,f'$ and $u\in\bar{\mathbf Q}_\ell^\times$, $r_{f'}=\tau(u)r_f\tau(u)^{-1}$ if and only if $f'\circ q_A=u\cdot(f\circ q_A)$.
\end{enumerate}
Indeed, a point $f$ is continuous: $|f(y)|\le\rho^n$ for $y\in\mathfrak m_R^n$, for some $\rho<1$.
By Proposition~\ref{sha-model-exists}(c), $q_A(A_j)$ lies in the closure of the $\mathcal O$-span of the products of entries of the matrix parts of the $r^{\mathrm u}(g)$ of total weight $j$.
For (Y1), such a product with $j<0$ contains an entry of negative weight, and the matrix part of $r^{\mathrm u}(g)$ with $\eta(g)=-1$ is that of $r^{\mathrm u}(gc)$, with $gc\in\ker\eta$, since $r^{\mathrm u}(c)$ has matrix part $1$; conversely the entries of negative weight are images of elements of the $A_j$ with $j<0$, by Proposition~\ref{sha-model-exists}(b).
For (Y2), if $r_{f'}=\tau(u)r_f\tau(u)^{-1}$, then $f'$ is $u^w$ times $f$ on each entry of weight $w$, so $f'\circ q_A=u^jf\circ q_A$ on $A_j$, by continuity; the converse holds because the entries of weight $w$ are images of elements of $A_w$.

\emph{Points.} Let $f$ be two-signed. By Lemma~\ref{sha-model-points}(c), either $r_f$ is irreducible on $\ker\eta$, or it is block shaped, with middle entry $1$ and irreducible corner block, and there is a representation $W$ as there.
As in the proof of Lemma~\ref{sha-model-tangent}(e), $r_f$ is an ordinary point with values in a discrete valuation ring finite over $\mathcal O$, by Lemma~\ref{sha-model-universal}(b).
In the first case, $\rho_f$ is a stable automorphic triple of weight $k$ and level $N$ by Theorem~\ref{sha-residually-reducible-automorphy}(ii) and Lemma~\ref{sha-admissibility-stable}; in the second, it is an endoscopic one by Theorem~\ref{sha-residually-reducible-automorphy}(iii), applied to $W$.
So, by hypothesis, $\operatorname{tr}\rho_f(\sigma)\ne t$ for every two-signed $f$.
By Lemma~\ref{sha-model-points}(d), two two-signed points $f,f'$ with $\rho_f\cong\rho_{f'}$ have $r_{f'}=\tau(u)r_f\tau(u)^{-1}$ for some $u$, and then $f'\circ q_A$ is in the orbit of $f\circ q_A$ by (Y2).

\emph{Conditions (1) to (4).} The model satisfies conditions (1) to (4) of Lemma~\ref{sha-algebra-cotangent}.

Condition (2). By Proposition~\ref{sha-model-exists}(c), $q_A(A_0)$ lies in the closure of $R^\natural_0$, which is closed in $R$ because it is a finitely generated $\mathcal O$-module. Since the kernel of $q_A$ is $I$, the image of $A_0$ in $A/I$ is isomorphic to a submodule of $R^\natural_0$, hence finitely generated over $\mathcal O$.

Condition (1). Its first sentence is Proposition~\ref{sha-model-exists}(a).
Let $j\le0$, and let $D$ be a derivation at $x$ with values in $\bar{\mathbf Q}_\ell$, of weight $j$, that vanishes on $I$. Then $D=0$.
Indeed, since $x=x_0\circ q_A$ and $x_0$ is a point, Proposition~\ref{sha-model-exists}(e), which applies by condition (2), gives $\tilde f$ with $\tilde f\equiv x_0\pmod\epsilon$ and $D=D_{\tilde f}\circ q_A$.
By Proposition~\ref{sha-model-exists}(b), $D_{\tilde f}$ vanishes on the entries of weight other than $j$ of the matrix parts of the $r^{\mathrm u}(g)$, so the tangent cocycle $y_{\tilde f}$ takes values in the matrices supported on the entries of weight $j$.
By Lemma~\ref{sha-model-tangent}(d), $y_{\tilde f}$ is the coboundary of some $X$; by Lemma~\ref{sha-model-tangent}(a) and (c), $X$ is a multiple of $\operatorname{diag}(1,0,-1)$ and its coboundary for $r_0$ is $0$.
So $y_{\tilde f}=0$, and $D_{\tilde f}=0$ by Lemma~\ref{sha-model-tangent}(b); hence $D=0$.
Now let $(j,H)$ be $(-1,F)$, $(-2,G)$ or $(0,Z)$, and suppose that some $p\in L_j$ has no nonzero multiple in $\langle H\rangle$.
Since $\bar{\mathbf Q}_\ell$ is divisible, hence injective as an $\mathcal O$-module, there is an $\mathcal O$-linear $\lambda:A/\Sigma_x\to\bar{\mathbf Q}_\ell$ with $\lambda(p)=1$ that vanishes on $\langle H\rangle$.
Let $D(f)=\lambda\big(\overline{f_j-x(f_j)}\big)$, where $f_j$ is the component of weight $j$ of $f$ and the bar denotes the class modulo $\Sigma_x$.
Since $(\ker x)^2\subset\Sigma_x$, $D$ is a derivation at $x$, of weight $j$; it vanishes on the members of $H$ by the choice of $\lambda$, and on the other generators of $I$, which have other weights and lie in $\ker x$.
So $D=0$, which contradicts $\lambda(p)=1$.

Condition (3). The homomorphism $x$ is the only element of $Y$ that vanishes on every $A_j$ with $j\ne0$, so the condition holds.
Indeed, let $x_1$ be such an element. By Proposition~\ref{sha-model-exists}(d), $x_1=u\cdot(f\circ q_A)$ for a point $f$, and $f\circ q_A$ also vanishes on these $A_j$.
By (Y1), the matrix parts of $r_f$ on $\ker\eta$ are upper and lower triangular, so $r_f=r_0=r_{x_0}$ by Lemma~\ref{sha-model-points}(b), and $f\circ q_A=x_0\circ q_A=x$ by (Y2) with $u=1$.
So $x_1$ and $x$ agree on $A_0$ and vanish on the other $A_j$, and $x_1=x$.

Condition (4). Let $\xi\in Y$ be two-signed and $D,D'$ derivations at $\xi$ with values in $\bar{\mathbf Q}_\ell$ that vanish on $I$.
By Proposition~\ref{sha-model-exists}(d), $\xi=u\cdot(f\circ q_A)$ for a point $f$, which is two-signed by (Y1).
The maps equal to $u^{-j}D$ and $u^{-j}D'$ on each $A_j$ are derivations at $f\circ q_A$ that vanish on $I$, which is generated by homogeneous elements; by Proposition~\ref{sha-model-exists}(e) they are $D_{\tilde f}\circ q_A$ and $D_{\tilde f'}\circ q_A$.
By Lemma~\ref{sha-model-tangent}(e), $y_{\tilde f}$ and $y_{\tilde f'}$ are the coboundaries of some $X$ and $X'$, which are multiples $b\operatorname{diag}(1,0,-1)$ and $b'\operatorname{diag}(1,0,-1)$ by Lemma~\ref{sha-model-tangent}(a) and (c).
Put $(a,a')=(0,1)$ if $b'=0$, so that $y_{\tilde f'}=0$, and $(a,a')=(b',-b)$ otherwise; in both cases $(a,a')\ne(0,0)$ and $ay_{\tilde f}+a'y_{\tilde f'}=0$.
Lemma~\ref{sha-model-tangent}(b) gives $aD_{\tilde f}+a'D_{\tilde f'}=0$, hence $aD+a'D'=0$ on each $A_j$, and $D$ and $D'$ are linearly dependent.

\emph{The first bound.} Let $e_{11},e_{22},e_{33}\in A_0$ map to the diagonal entries of the matrix part of $r^{\mathrm u}(\sigma)$ (Proposition~\ref{sha-model-exists}(b)), and $a=e_{11}+e_{22}+e_{33}-t$.
Since $\eta(\sigma)=1$, $t$ is the trace of $r_0(\eta,\Psi)(\sigma)$; so $x(a)=0$, and for $\xi$ in the orbit of $f\circ q_A$, $\xi(a)=\operatorname{tr}\rho_f(\sigma)-t$, since $a$ has weight $0$.
This is nonzero when $f$ is two-signed, which, with Proposition~\ref{sha-model-exists}(d) and (Y1), is condition (5).
Lemma~\ref{sha-algebra-cotangent}, with $v=v_\ell$ since $\mathcal O$ is unramified, shows that $\lambda_{-1}(F)$ is finite and gives a finite set $T_0$ of two-signed elements of $Y$ meeting every orbit of two-signed elements, with $\lambda_{-1}(F)\le2\sum_{\xi\in T_0}v_\ell(\xi(a))$.
Each $\xi\in T_0$ is in the orbit of $f\circ q_A$ for a two-signed point $f$.
Removing points one at a time, keep a set $T$ of such points $f$, no two of which are related by $r_{f'}=\tau(u)r_f\tau(u)^{-1}$, such that every two-signed point is so related to one of them.
By (Y2) and Proposition~\ref{sha-model-exists}(d), the $f\circ q_A$ with $f\in T$ meet the orbit of every two-signed element of $Y$, and they are pairwise distinct, since $f\circ q_A$ determines $r_f$ by Proposition~\ref{sha-model-exists}(b), and $r_f$ determines $f$.
By the last sentence of \emph{Points}, the $\rho_f$ with $f\in T$ are pairwise non-isomorphic, and the bound of Lemma~\ref{sha-algebra-cotangent} for this set gives
\[
 \lambda_{-1}(F)\le2\sum_{f\in T}v_\ell\big(\operatorname{tr}\rho_f(\sigma)-t\big).
\]

\emph{The second bound.} Let $b\in\mathcal O$ be nonzero and $\mathcal M$ the $b$-torsion of the group defining $V'_\ell(k,\psi')$, the group of classes in $H^1(\operatorname{Gal}(\bar{\mathbf Q}/\mathbf Q),\mathcal O\otimes\mathbf Q_\ell/\mathbf Z_\ell(\Psi^{-1}))$ that vanish on the inertia groups and on the decomposition groups at the primes dividing $\ell N$.
\begin{enumerate}
\item[(S1)] Let $\mu\in\mathcal M$. As in the proof of Lemma~\ref{sha-lifts-comparison}, $\mu$ is represented by a cocycle $z$ with values in $\mathcal O/b\mathcal O(\Psi^{-1})$. Let $Z(g)$ be the matrix with entries $z(g)$ at $(2,1)$ and $-z(g)$ at $(3,2)$, the two entries of weight $-1$. As in the proof of Lemma~\ref{sha-lifts-presentation}(e), $Z$ is a cocycle for the reduction of $r_0(\eta,\Psi)$ modulo $b$, and after adding the coboundary of a constant matrix of weight $-1$ it vanishes at $c$ and satisfies the condition at $g_1$ of Lemma~\ref{sha-model-tangent}(a), because $\Psi(c)=1$ and $\Psi(g_1)^2-1$ is a unit. It vanishes on $\mathcal N$, where $\Psi=1$ and the class is unramified. On $D_B$ the class vanishes, so $z=(\Psi^{-1}-1)m$ there for some $m$; evaluating at the element $x$ of the proof of Lemma~\ref{sha-model-universal}(b), for which $\Psi(x)-1$ is a unit, gives $bm=0$, so $Z$ is on $D_B$ the coboundary of a matrix $X_\mu$ of weight $-1$ with entries in $\mathcal O/b\mathcal O$. The class $\mu$ is recovered from the entry $(2,1)$ of $Z$, and $\mu\mapsto Z$ is additive.
\item[(S2)] Lift $X_\mu$ to a matrix $\tilde X_\mu$ over $\mathcal O$ of weight $-1$. The homomorphism $h\mapsto(1+\epsilon\tilde X_\mu)r_0(h)(1+\epsilon\tilde X_\mu)^{-1}$ of $D_B$, with values in $\mathcal G_3(\mathcal O[\epsilon]/(\epsilon^2))$, is the conjugate of $x_B\circ u_B$ by a matrix congruent to $1$, so by Lemma~\ref{sha-model-local}(b) it is $(x_B+\epsilon\partial_\mu)\circ u_B$ for a derivation $\partial_\mu$ at $x_B$ with values in $\mathcal O$ that vanishes on $\mathfrak a_B$; its values on the entries of the matrix parts are those of the coboundary of $\tilde X_\mu$ times $r_0$, so they vanish on the entries of weight other than $-1$. The homomorphism $g\mapsto((1+\epsilon Z(g))M_g,\eta(g))$, where $r_0(\eta,\Psi)(g)=(M_g,\eta(g))$, with values in $\mathcal G_3$ of the ring $\mathcal O\oplus\epsilon(\mathcal O/b\mathcal O)$ with $\epsilon^2=0$, is an $\mathfrak a_B$-lift: it is normalized and trivial on $\mathcal N$ by (S1), and on $D_B$ it is the reduction of the homomorphism above. So it is $(x_0+\epsilon D_\mu)\circ r^{\mathrm u}$ for a derivation $D_\mu$ at $x_0$ with values in $\mathcal O/b\mathcal O$, whose value on the entry $(a',b')$ of the matrix part of $r^{\mathrm u}(g)$ is $\Psi(g)^{w_{b'}}Z(g)_{a'b'}$ with $(w_1,w_2,w_3)=(1,0,-1)$. It vanishes on the entries of weight other than $-1$, and $\mu\mapsto D_\mu$ is additive and injective, by (S1).
\item[(S3)] Proposition~\ref{sha-model-exists}(f), applied to $D_\mu$ and $\partial_\mu$, gives a derivation $D'_\mu$ at $x$ with values in $\mathcal O$, of weight $-1$, with $D_\mu\circ q_A\equiv D'_\mu\pmod b$. The map $\mu\mapsto D_\mu\circ q_A$ is injective, by Proposition~\ref{sha-model-exists}(b), and $D_\mu\circ q_A$ vanishes on $F$, which lies in the kernel $I$ of $q_A$.
\end{enumerate}
So Lemma~\ref{sha-algebra-derivations} gives $\operatorname{len}\mathcal M\le\lambda_{-1}(F)$.
The group defining $V'_\ell(k,\psi')$ is torsion, so $V'_\ell(k,\psi')\le\lambda_{-1}(F)$.
With $X=\{\rho_f:f\in T\}$, the two bounds give the one-character proposition.
% Counting proposition: Z.J.NumberU3U_s, proved by GalN.V.Win.N.numberU3U_of_prints through numberU3U_of_adm2A_imp_adm2B (Lemmas sha-eigensystem-traces and sha-eigensystem-norm); written out.
% One-character proposition: OneCharU_s, proved by oneChar_unramU (FOneCharU) through exists_unram, coefIndep, exists_dat, exists_uLoc, exists_ring, exists_model (Proposition sha-model-exists), U3.model_upper (conditions (1)-(4): modelIdeal_finiteness_of_weightZero_finite, gradedModel_spansCot_of_rigid_weights, gradedModel_pointCondition_of_ptHyp, gradedModel_derivationCondition_of_ptHyp; model_signs for (Y1) and (Y2)) and U3.vlen_le (sha_inj_gen for (S1), model_lower with main_der for (S2) and (S3), LemmaM' for the last step). PtHyp is Lemma sha-model-tangent (ptHyp_of_autB_of_lemW_of_isBase); the points are twoPt_shape and GalN.V.adm2B_rhoK_of_typeII_of_modular_extensions.
% Lemmas sha-model-local, sha-algebra-cotangent, sha-rigidity-points (with sha-rigidity-adjoint) and sha-admissibility-stable are proved in place; what remains of them is in the open steps of their own proofs.

Group the characters $\psi'\in X_h(k)$ by their kernels $H'$; the characters with a given kernel form one orbit under $\operatorname{Aut}(\mathbf C)$, of size $\varphi([(\mathbf Z/N_h)^\times:H'])$.
For $\ell$ in case I, the data $(\ell,k,\psi',K)$ form an Eisenstein datum for every $\psi'\in X_h(k)$, by the reductions in the proof of Lemma~\ref{sha-cases}.
Fix $H'$ and a prime $p\in\{r_1,r_2\}$ with $p\ne\ell$, and apply the counting proposition with $w$ above $p$ and the family $z_\psi=\psi(\mathrm N w)^{-1}$ over the orbit; with the second proposition at $\sigma=\sigma_w^{-1}$ this gives $\sum_{\psi'}\bar V_\ell(k,\psi')\le2v_\ell(N_0)$, the sum over the orbit.
Both $r_1$ and $r_2$ are used so that one of them differs from $\ell$; adding over $\ell$ gives $\sum_\ell\log\ell\sum_{\psi'}\bar V_\ell(k,\psi')\le2\log|N_0(r_1)|+2\log|N_0(r_2)|$, which is at most $2c_5\varphi([(\mathbf Z/N_h)^\times:H'])(N_h(k+1))^A(\log r_1+\log r_2)$.
The sum of $\varphi([(\mathbf Z/N_h)^\times:H'])$ over the kernels is at most $N_h$, and $k<t$, which gives the bound $(N_ht)^{C}$.

(b) The proof has the shape of (a), with $GSp_4$ in place of $U(3)$.
Let $(\ell,k,\psi')$ be in case II or III; then $k\ge2$.
After enlarging $\mathcal O$, which changes neither side, $\mathcal O$ contains the $2^{n_{\mathrm{II}}+n_{\mathrm{III}}}$-th roots of unity (Definition~\ref{sha-siegel-datum}(d)).
Then there is a Siegel datum $D$ attached to $(h,k,\psi')$, by Lemma~\ref{sha-counts-siegel-datum}.
Let $V_D=V_{\ell,n_S}(k,\psi')$, where $n_S$ is the product of the primes of $S$ (Definition~\ref{sha-character-parts}).
Then $V'_\ell(k,\psi')\le V_D$: at a prime $p\in S$ not dividing $\ell N_h$, a class that is unramified at $p$ vanishes on the decomposition group at $p$, since Frobenius acts on the divisible coefficients by $\psi'(p)^{-1}p^{-k}\ne1$.
By Proposition~\ref{sha-siegel-points}, $D$ is finite and has no endoscopic $\bar{\mathbf Q}_\ell$-points, and there is $\iota$ such that every irreducible $\bar{\mathbf Q}_\ell$-point of $D$ is automorphic through $\iota$ and $D$ is rigid at it.
Let $(n',p')$ be $(n_{\mathrm{II}},1)$ in case II and $(n_{\mathrm{III}},p_1(\psi))$ in case III, and put $M=\ell2^{n'}N_hp'$ and $n=(\ell-1)2^{n'+1}o_\psi$.
The characters $\psi'$, $\xi$ and $\omega$ of $D$ have level $M$ and order dividing $n$, and every irreducible $\bar{\mathbf Q}_\ell$-point of $D$ is of Siegel type with level $M$ and weight $k$ (Proposition~\ref{sha-siegel-count}).
Let $p$ be the least prime $p\ge5$ not dividing $\ell N_hp'$, so that $p\nmid M$ and $p\notin S$, and let $\sigma$ be a Frobenius element at $v_p$.
Apply Proposition~\ref{sha-siegel-count} with the $n$-th roots of unity $z_1,z_2,z_3$ that $\iota$ attaches to $(\omega^{-1}\xi\psi')(\sigma)$, $\psi'(\sigma)$ and $(\omega^{-1}\xi)(\sigma)$, so that $\iota^{-1}(t)=\operatorname{tr}r_0(\sigma)$.
It gives $N_0\ne0$ with $\log|N_0|\le c\,n\,(M(k+1))^{A}\log p$ such that $\operatorname{tr}r(\sigma)\ne\operatorname{tr}r_0(\sigma)$ for every irreducible $\bar{\mathbf Q}_\ell$-point $r$, and $|N_0|_\ell\le\prod_{r\in X}|\operatorname{tr}r(\sigma)-\operatorname{tr}r_0(\sigma)|_\ell$ for every finite set $X$ of pairwise non-conjugate irreducible $\bar{\mathbf Q}_\ell$-points.
By Proposition~\ref{sha-counts-vd-model-exists}, $D$ has a graded model. So Proposition~\ref{sha-counts-vd-bound}, applied to $D$ and $\sigma$, shows that $V_D$ is finite and gives such a set $X$ with
\[
 \frac{V_D}{e_{\mathcal O}}\log\ell\le-2\sum_{r\in X}\log\big|\operatorname{tr}r(\sigma)-\operatorname{tr}r_0(\sigma)\big|_\ell.
\]
Since $-\log|N_0|_\ell\le\log|N_0|$, this gives $\bar V_\ell(k,\psi')\log\ell\le2c\,n\,(M(k+1))^{A}\log p$.
Here $M$, $n$ and $p$ are bounded by powers of $\ell N_hk$ with exponents depending only on $\psi$, which gives the bound $(\ell N_hk)^{C}$.
% (b) is countII_and_countIII_of_printsH_of_sp4In, through Th.sp4Main_alt and main_big. The enlargement of O is coefIndep with chHom and conductor_at_map_injective_coeff; D is GradeZeroInterpolatingElement.dataEx, written out as Lemma sha-counts-siegel-datum (core, with X.CorX as Lemma sha-counts-siegel-residual); V_D is Raw.sha, compared with V'_l by shaEq; Proposition sha-siegel-points is Aut.SymplecticFourImageDense (sp4In0_alt) and Th.RigiditySymplecticFour (rigSp4_alt); M and n come from levOf; Siegel type is GradeZeroInterpolatingElement.level, Z.Win.A4.safe4' and the density; the count is Z.Win.A4.exists_frob_nonvanishing_with_sha_bound from GalN.V.numberSp4_n_alt, with p = pZero (l N_h p'); the commutative algebra is galoisRep_ineq_of_allIrreducible_nonvanishing, written out as Proposition sha-counts-vd-bound (with Alg.M.main_holds and Alg.M.mprime_holds as Lemmas sha-algebra-cotangent and sha-algebra-derivations, and X.soBar and X.soFin), for the Model from exModel with Alg.lemmaP, phiGam, nums and locBase (Proposition sha-counts-vd-model-exists); the last step is GradeZeroInterpolatingElement.arith.
% The graded model (Proposition sha-counts-vd-model-exists) is exModel; what remains of it is in the open step of its own proof. Lemma sha-counts-vd-inputs, with (b) (diagCases) and (f) (aSmooth), and Lemma sha-counts-vd-presentations are written out.

(d) Let $(\ell,1,\psi')$ be in case V with $16\,o_\psi(\ell N_\psi)^2\,4^{n_0}\le2^{c_2}$, and let $K=\mathbf Q(i)$, the field of the character $\eta$ of Definition~\ref{sha-lifts-problem}; choose $\iota:\bar{\mathbf Q}_\ell\cong\mathbf C$.
By Lemma~\ref{sha-lifts-host-data}, the data $(\ell,\mathcal O,1,\psi')$, with $N=N_h$, satisfy the dyadic condition for $o_\psi$ and the Bernoulli condition for $N_\psi$.
Lemma~\ref{sha-lifts-automorphy} gives $g_1$ such that the universal unitary lift $u$ relative to $g_1$ (Lemma~\ref{sha-lifts-universal}) has $R_0$ finite over $\mathcal O$.
Let $C$ be the constant of Lemma~\ref{sha-lifts-number}, which depends only on $K$, let $\varphi\in\Gamma$ be given by that lemma, and let $\mathcal P$ be the property of lying in the set $\mathcal A$ of that lemma.
By Lemma~\ref{sha-lifts-automorphy}(b), the representation of every point of $u$ that is irreducible lies in $\mathcal A$.
Every $r\in\mathcal A$ is rigid by Lemma~\ref{sha-lifts-rigidity}, and has $\operatorname{tr}r(\varphi)\ne\operatorname{tr}r_0(\varphi)$ by Lemma~\ref{sha-lifts-number}.
So Proposition~\ref{sha-lifts-abstract-count} applies, and gives $r_1,\dots,r_s\in\mathcal A$ whose matrix parts on $\ker\eta$ are pairwise not conjugate, and $V$ with $\operatorname{len}\text{Ш}_a\le V$ for every nonzero $a$ and $V\le2e_{\mathcal O}\sum_iv_\ell(\operatorname{tr}r_i(\varphi)-\operatorname{tr}r_0(\varphi))$.
By Lemma~\ref{sha-lifts-comparison}, $V'_\ell(1,\psi')\le V$, so $V'_\ell(1,\psi')$ is finite, and by Lemma~\ref{sha-lifts-number}
\[
 \bar V_\ell(1,\psi')\log\ell=\frac{V'_\ell(1,\psi')}{e_{\mathcal O}}\log\ell\le2\sum_{i=1}^{s}v_\ell\big(\operatorname{tr}r_i(\varphi)-\operatorname{tr}r_0(\varphi)\big)\log\ell\le2(\ell N_h)^{C}\le(\ell N_hk)^{C+1},
\]
since $k=1$ and $\ell N_h\ge2$.
% (d) is count_v_of_printed_bound_k_reduced_unitary_lift, through UnitaryThreeInertia.countV_of_finiteness_duality_euler_host: Lemmas sha-lifts-host-data (hostData_of_host), sha-lifts-universal (exists_hostFrame, exists_isUniv), sha-lifts-residual (hostFrame_locVan, hostFrame_count, X.So.hostVanE, frame_class_one, X.xpp), Proposition sha-lifts-abstract-count (abstract_count), sha-lifts-comparison (vlen_le), sha-lifts-automorphy (unitary_lift_of_reduced_core of ends0_alt.2), sha-lifts-rigidity (Z.J.hostRigidUE_alt) and sha-lifts-number (Z.J.hostNumberUE_alt with numberU3U_of_prints); the final arithmetic is the private lemma arith in CountVOf.

For (e) and (f) we use the following comparison.
Let $L$ be a finite Galois extension of $\mathbf Q$ such that $\psi'$ is trivial on $\operatorname{Gal}(\bar{\mathbf Q}/L)$.
Restriction maps the group defining $V_\ell(k,\psi')$ into $\text{Ш}^1_\ell(L,k)\otimes\mathcal O$, with kernel contained in $H^1$ of $\operatorname{Gal}(L/\mathbf Q)$ with coefficients in the invariants of $\operatorname{Gal}(\bar{\mathbf Q}/L)$; by the bound for invariants in the proof of Lemma~\ref{sha-plus-by-characters} this gives
\[
 \bar V_\ell(k,\psi')\le V_\ell(k,\psi')/e_{\mathcal O}\le\operatorname{len}\text{Ш}^1_\ell(L,k)+[L:\mathbf Q]\,\nu(\ell,N_h,k),\tag{$\dagger$}
\]
where $\nu(\ell,N_h,k)\le2N_hk$ is the number $\nu$ of Lemma~\ref{sha-plus-by-characters}(b); the first inequality holds because the group defining $V'_\ell$ is contained in the one defining $V_\ell$.

(e) Since $(\ell,1,\psi')$ is in case V with large $2$-conductor, $\ell\in B_\psi$ and $\operatorname{ord}(\psi')\,\ell$ is at most a number $n_0$ depending only on $\psi$.
There is a Galois number field $L$ containing $\mu_\ell$, of degree at most $\operatorname{ord}(\psi')\,\ell\le n_0$, on whose absolute Galois group $\psi'$ is trivial; its discriminant is bounded in terms of $n_0$ and $\ell N_h$.
By Poitou-Tate duality, as at the end of the proof of Lemma~\ref{sha-tower-from-module}, $\operatorname{len}\text{Ш}^1_\ell(L,1)\le\operatorname{len}H^2(G_{L,S_\ell},\mathbf Z_\ell(2))$, and Theorem~\ref{sha-soule-bound} bounds $\log\ell$ times this length by $\exp\big((2n_0)^{C_K}\log(2|d_L|)\big)$.
With ($\dagger$) for $k=1$, and since $\ell$, $n_0$ and $\log|d_L|$ are bounded by powers of $\ell N_h$, this gives $\bar V_\ell(1,\psi')\log\ell\le(\ell N_h)^{C}$.

(f) Since $\ell\in B_\psi$, $\ell\ne q,d$, and since $\chi'_\ell(\psi')$ has odd order, the conductor of $\psi'$ divides $N_\psi$.
So there are a field $L_0$ of the cyclotomic datum of $\ell N_\psi$ and a layer $L_1$ of degree $\ell^{j}$ over $L_0$, Galois over $\mathbf Q$ of degree at most $\ell N_\psi\ell^{j}$, on whose absolute Galois group $\psi'$ is trivial.
By hypothesis $L_1$ satisfies the $\ell$-bound $(t,B)$.
Summing ($\dagger$) over $0<k<t$ gives $\sum_k\bar V_\ell(k,\psi')\le B+\sum_{k<t}\ell N_\psi\ell^{j}\cdot2N_hk$, and the last sum is at most $(N_ht)^{\ell N_\psi\ell^{j}+3}$.

\end{proof}
\leanref{Development/AutomorphicArithmetic/AutomorphicVarieties/HeightAndVanishing/Counts.lean}{UpperTheory.V.HV.CountI}
\leanref{Development/AutomorphicArithmetic/AutomorphicVarieties/HeightAndVanishing/Counts.lean}{UpperTheory.V.HV.CountII}
\leanref{Development/AutomorphicArithmetic/AutomorphicVarieties/HeightAndVanishing/Counts.lean}{UpperTheory.V.HV.CountIII}
\leanref{Development/AutomorphicArithmetic/AutomorphicVarieties/HeightAndVanishing/Counts.lean}{UpperTheory.V.HV.CountZ}
\leanref{Development/AutomorphicArithmetic/AutomorphicVarieties/HeightAndVanishing/Counts.lean}{UpperTheory.V.HV.CountV}
\leanref{Development/AutomorphicArithmetic/AutomorphicVarieties/HeightAndVanishing/Counts.lean}{UpperTheory.V.HV.CountVs}
\leanref{Development/AutomorphicArithmetic/AutomorphicVarieties/HeightAndVanishing/LayerLe.lean}{UpperTheory.V.HV.CountIV}
\leanref{Development/AutomorphicArithmetic/AutomorphicVarieties/ModularityArgument/ThresholdArguments/DenseThresholdCases/LinkedUnlE.lean}{UpperTheory.V.Z.J.count_i_of_printed_bound_k_kato_automorphy}
\leanref{Development/AutomorphicArithmetic/AutomorphicVarieties/ModularityArgument/ThresholdArguments/DenseThresholdCases/AdjointLiftingOfCohomology.lean}{UpperTheory.V.HV.DescentVanishingReductions.countII_and_countIII_of_printsH_of_sp4In}
\leanref{Development/AutomorphicArithmetic/AutomorphicVarieties/ModularityArgument/ThresholdArguments/DenseThresholdCases/LinkedUnlE.lean}{UpperTheory.V.Z.J.count_v_of_printed_bound_k_reduced_unitary_lift}
\leanref{Development/AutomorphicArithmetic/AutomorphicVarieties/HeightAndVanishing/CountVs.lean}{UpperTheory.V.HV.countVs}
\leanref{Development/AutomorphicArithmetic/AutomorphicVarieties/HeightAndVanishing/CountIV.lean}{UpperTheory.V.HV.countIV}
\leanref{Development/AutomorphicArithmetic/AutomorphicVarieties/HeightAndVanishing/CountZ.lean}{UpperTheory.V.HV.countZ}

\begin{theorem}\label{sha-plus-part}
Assume the Galois representations as in Theorem~\ref{sha-residually-reducible-automorphy}.
Let $K$ be a window constant and $C_t>0$ (Definition~\ref{sha-window-constants}).
For every primitive Dirichlet character $\psi$ there are $C$ and primes $r_1,r_2$ such that for every $0<\varepsilon<1/52$ and every real $A_1$ there is $X_0$ with the following property.
Let $(e,t)$ be admitted for $(A_1,C_t)$, let $X\ge X_0$, and let $(q,d)$ be a prime pair for $M_\psi$ such that $q$ and $d$ exceed $2X^{C_t}+1$ and each of $N_\psi,o_\psi,p_K(\psi),h^*(\psi),p_1(\psi),r_1,r_2$, and $\log(qd)\le X^{\varepsilon/(4J_\psi)}$.
Here $J_\psi$ is the number $J$ of Definition~\ref{sha-window-constants} for $M=M_\psi$, $B'=B_\psi$ and $L(\ell)$ the cyclotomic datum of $\ell N_\psi$.
Let $Y$ be a controlled window for $(K,C_t,\varepsilon,M_\psi,B_\psi,L,e,t,q,d,X)$, and let $x\in[Y,Y^2]$.
Then every CM host field $E$ of $(\psi,e(x),q,d)$ satisfies the plus bound
\[
 \Big(t(x),\;(N_Et(x))^{C}+C\exp\big(x^{4\varepsilon/5}\big)\Big).
\]
\end{theorem}
\begin{proof}
\emph{Primes outside $B_\psi$.} Let $h=(\psi,e,q,d)$ be a host datum large for $t\ge2$, $E$ a host field and $S$ a finite set of odd primes outside $B_\psi$.
By Lemma~\ref{sha-cases}, $\ell\nmid\operatorname{ord}\psi'$ for $\psi'\in X_h(k)$, and the odd primes dividing $[E:\mathbf Q]$ lie in $B_\psi$; so Lemma~\ref{sha-plus-by-characters}(a) gives $\operatorname{len}\text{Ш}^+_\ell(E,k)\le\sum_{\psi'\in X_h(k)}\bar V_\ell(k,\psi')$.
Put $M=N_ht$.
For fixed $k$, the primes of $S$ in case I contribute at most $M^{C_1}$ by Proposition~\ref{sha-counts}(a).
For the other primes of $S$, $\ell-1\mid2k$, so $\ell\le2t$; each $\psi'$ is in case Z, II or III and contributes at most $M^{3C_2}$ by Proposition~\ref{sha-counts}(b),(c), where $C_2$ is the larger of the constants there and $\ell N_hk\le M^3$; there are at most $2t$ such primes and at most $N_h$ characters.
Summing over $k<t$ gives at most $t(M^{C_1}+2MM^{3C_2})\le M^{C_1+3C_2+4}$.

\emph{Primes in $B_\psi$.} Fix $\ell\in B_\psi$ and $j\ge0$, and suppose that the $\ell$-bound $(t,B)$ holds for the layers of degree $\ell^{j}$ over the fields of the cyclotomic datum of $\ell N_\psi$, with $B\ge0$.
Lemma~\ref{sha-plus-by-characters}(b) bounds $\operatorname{len}\text{Ш}^+_\ell(E,k)$ by $\sum_{\psi'\in X_h(k)}\big(\bar V_\ell(k,\psi')+\lambda\big)$ with $\lambda\le(N_hk)^6$.
By Lemma~\ref{sha-cases} each $\psi'$ is in case II, IV or V.
The characters in case IV have conductor dividing $N_\psi$, so there are at most $N_\psi$ of them, and Proposition~\ref{sha-counts}(f) bounds the sum over $k$ of each by $B+M^{C_4}$.
The others satisfy $\bar V_\ell(k,\psi')\log\ell\le(M^2)^{C_m}$ by Proposition~\ref{sha-counts}(b),(d),(e), with $C_m$ the largest of their constants.
Adding these, with $\log\ell\le\ell\le M$ and $N_\psi\le M$, gives
\[
 \log\ell\sum_{k=1}^{t-1}\operatorname{len}\text{Ш}^+_\ell(E,k)\le M^{C_\ell}+C_\ell B
\]
for a constant $C_\ell$ depending on $\psi$, $\ell$ and $j$.

\emph{Assembly.} Let $j_\ell=\lfloor\log_\ell R_\ell(L(\ell))\rfloor+3$ for $\ell\in B_\psi$, let $D=\max_\ell C_\ell$, and let $C=C_0+D+|B_\psi|+|B_\psi|D$ with $C_0=C_1+3C_2+4$ and $r_1,r_2$ from Proposition~\ref{sha-counts}(a).
Combining the two parts, every host field of a host datum $h$ large for $t\ge2$ satisfies the plus bound $\big(t,(N_ht)^{C}+CB\big)$ whenever $r_1,r_2<q,d$ and the $\ell$-bounds $(t,B)$ hold for the layers of degree $\ell^{j_\ell}$, $\ell\in B_\psi$.

Now let $X_0=2^{J_\psi}$ and take the data of the statement.
For $x\in[Y,Y^2]$ we have $x\ge X^{1/J_\psi}\ge2$, so $e(x)\ge2$ and $2\le t(x)\le x^{C_t}\le X^{C_t}$.
Hence $h=(\psi,e(x),q,d)$ is a host datum that is large for $t(x)$.
Condition (ii) of a controlled window gives the $\ell$-bounds $\big(t(x),\exp(x^{4\varepsilon/5})\big)$ for the layers of degree $\ell^{j_\ell}$, and Lemma~\ref{sha-cases} gives $N_E=N_h$.
\end{proof}
\leanref{Development/AutomorphicArithmetic/AutomorphicVarieties/PrimeSetBoundDataOfPrimitiveChar.lean}{UpperTheory.V.ThmHV}
\leanref{Development/AutomorphicArithmetic/AutomorphicVarieties/HeightAndVanishing/JointWinOfCounts.lean}{UpperTheory.V.HV.thmHV_of_counts}
\leanref{Development/AutomorphicArithmetic/AutomorphicVarieties/HeightAndVanishing/ThmHV.lean}{UpperTheory.V.HV.thmHV}
\leanref{Development/AutomorphicArithmetic/AutomorphicVarieties/HeightAndVanishing/PointArith.lean}{UpperTheory.V.HV.point}
\leanref{Development/AutomorphicArithmetic/AutomorphicVarieties/HeightAndVanishing/SumOut.lean}{UpperTheory.V.HV.sumOut}
\leanref{Development/AutomorphicArithmetic/AutomorphicVarieties/HeightAndVanishing/FinalLe.lean}{UpperTheory.V.HV.sumIn}

\section{Proof of the weakening}\label{section-sha-assembly}

\begin{theorem}\label{sha-weakening-from-galois-representations}
Assume that Galois representations exist, in the sense of Definition~\ref{gal-galois-representations}, relative to the maps satisfying the laws of Definition~\ref{gal-llc-laws}.
Then the weakening of Vandiver's conjecture holds: every primitive Dirichlet character satisfies the weakening of Definition~\ref{sha-vandiver-weakening} with exponent $401$.
\end{theorem}
\begin{proof}
Let $K$ be the window constant of Theorem~\ref{sha-prime-two}, and let $C_b$ be the constant of Theorem~\ref{sha-minus-part}.
Fix a primitive character $\psi$, $0<\varepsilon<1/52$ and $e$ with $x^{2+8\varepsilon}\le2^{e(x)}\le2x^{2+8\varepsilon}$ for $x\ge1$.
Put $t(x)=\lfloor x^{401}\rfloor$.
For $x\ge2$ we have $2^{e(x)}\ge x^2\ge4$, $2^{e(x)}\le2x^{3}\le x^{4}$ and $2\le t(x)\le x^{401}$, so $(e,t)$ is admitted for $(4,401)$.

\emph{A joint window.} Apply Theorem~\ref{sha-prime-two} with $C_t=401$, $M=M_\psi$, $B'=B_\psi$, and $L(\ell)$ the cyclotomic datum of $\ell N_\psi$, which contains $\mu_\ell$; it gives $X_1$.
Apply Theorem~\ref{sha-plus-part} with the same data; it gives $C,r_1,r_2$ and $X_2$.
Let $J=J_\psi$ and let $P$ be the sum of $N_\psi,o_\psi,p_K(\psi),h^*(\psi),p_1(\psi),r_1,r_2$.
For $X\ge\max(X_1,X_2)$ and every prime pair $(q,d)$ for $M_\psi$ with $q,d>2X^{401}+1$, $q,d>P$ and $\log(qd)\le X^{\varepsilon/(4J)}$, there is $Y$ with $X^{1/J}\le Y$ and $Y^2\le X$ such that for every $x\in[Y,Y^2]$: every compositum of $(M_\psi,e(x),q,d)$ satisfies the $2$-bound $\big(t(x),\exp(x^{4\varepsilon/5})\big)$, and every CM host field $E$ of $(\psi,e(x),q,d)$ satisfies the plus bound $\big(t(x),(N_Et(x))^{C}+C\exp(x^{4\varepsilon/5})\big)$.

\emph{Choice of $X_0$.} Let $J'=\max(J,1)$.
If $(q,d)$ is admissible at $Y'=X^{1/J'}$, then $\log q>Y'^{\varepsilon/5}/2$ and $\log d>Y'^{\varepsilon/5}\log^2Y'/2$, so for $X$ large $q$ and $d$ exceed $2X^{401}+1$ and $\max(P,N_\psi)$; in particular $qd$ is prime to $N_\psi$, and $(q,d)$ is a prime pair for $M_\psi$.
Also $\log(qd)\le Y'^{\varepsilon/5}(1+\log^2Y')\le Y'^{\varepsilon/4}=X^{\varepsilon/(4J')}\le X^{\varepsilon/(4J)}$ for $X$ large.
So for $X\ge X_0$, with $X_0$ large, the joint window $Y$ exists, and $X^{1/J'}\le X^{1/J}\le Y$.

\emph{The bound at $x$.} Let $x\in[Y,Y^2]$ and let $E$ be a host field of $(\psi,e(x),q,d)$.
By Lemma~\ref{sha-composita}, $E$ is a CM field and a compositum of $(M_\psi,e(x),q,d)$, with $N_E\le2^{e(x)+2}N_\psi qd$ and $[E:\mathbf Q]\le N_E$.
Since $x\ge X^{1/J'}$, also $\log(qd)\le x^{\varepsilon/4}$, and with $2^{e(x)}\le2x^{2+8\varepsilon}$ and $t(x)\le x^{401}$ this gives $\log(N_Et(x))\le2x^{\varepsilon/4}$ for $x$ large.
Theorem~\ref{sha-minus-part} gives the minus bound $\big(t(x),B_-\big)$ with $B_-$ at most $K_1\exp(8x^{\varepsilon/4})$, where $K_1=3+|C_b|\sum_{\ell\mid o_\psi,\ \ell\ \mathrm{odd}}\ell^{3}$, because $[E:\mathbf Q]\le N_E$ and $\log(N_Et(x))\le2x^{\varepsilon/4}$.
By Lemma~\ref{sha-bound-from-parts}, $E$ satisfies the Tate-Shafarevich bound $(t(x),B)$ with
\[
 B=\exp(x^{4\varepsilon/5})\log2+K_1\exp(8x^{\varepsilon/4})+(N_Et(x))^{C}+C\exp(x^{4\varepsilon/5}).
\]
Since $(N_Et(x))^{C}\le\exp(2Cx^{\varepsilon/4})$, we have $B\le\sqrt{\exp(x^{\varepsilon})}$ for $x$ large, and $x\ge X^{1/J'}$ is large when $X_0$ is.
\end{proof}
\leanref{Development/AutomorphicArithmetic/AutomorphicVarieties/ModularityArgument/ThresholdArguments/DenseThresholdCases/VwinOfRow.lean}{UpperTheory.One.vwin_of_row}
\leanref{Development/AutomorphicArithmetic/AutomorphicVarieties/ModularityArgument/ThresholdArguments/DenseThresholdCases/HVCore.lean}{UpperTheory.V.HV.Core.of_four_alt}
\leanref{Development/AutomorphicArithmetic/AutomorphicVarieties/HeightAndVanishing/JointWinOfCounts.lean}{UpperTheory.V.HV.vwin_of_counts}
\leanref{Development/AutomorphicArithmetic/AutomorphicVarieties/PrimeSetBoundDataOfPrimitiveChar.lean}{UpperTheory.V.JointWin}
\leanref{Development/AutomorphicArithmetic/AutomorphicVarieties/PrimeSetBoundDataOfPrimitiveChar.lean}{UpperTheory.V.jointWin_of}
\leanref{Development/AutomorphicArithmetic/AutomorphicVarieties/CMFieldExponentialBounds/AdmissibleOfPageN.lean}{UpperTheory.V.vwinAt_of_joint}
